\documentclass[12pt]{article}
\def\unitnumber{1}
\def\pdftitle{Linear Algebra and Its Statistical Applications}
\def\runningtitle{Chapter~\unitnumber}
\usepackage[T1]{fontenc}
\usepackage{lmodern,microtype}
\usepackage[a4paper,margin=15mm,headheight=14pt,headsep=5mm]{geometry}
\usepackage{amsmath,amssymb,amsthm,mathtools}
\usepackage{enumitem,booktabs,needspace}
\usepackage{pifont}
\usepackage{tikz}
\usetikzlibrary{arrows.meta}
\usepackage{fancyhdr}
\usepackage{xr-hyper}
\usepackage[hidelinks]{hyperref}
\makeatletter
\newcommand{\notetarget}[1]{\Hy@raisedlink{\hypertarget{#1}{}}}
\makeatother
\hypersetup{pdftitle={\pdftitle},pdfauthor={Chang-Ye Tu}}
\pagestyle{fancy}
\fancyhf{}
\fancyhead[L]{\small Linear Algebra}
\fancyhead[R]{\small\runningtitle}
\fancyfoot[C]{\small\thepage}
\renewcommand{\headrulewidth}{0.3pt}
\setlength{\parindent}{0pt}
\setlength{\parskip}{4pt plus 1pt minus 1pt}
\setlength{\emergencystretch}{1.5em}
\setlength{\abovedisplayskip}{7pt plus 2pt minus 2pt}
\setlength{\belowdisplayskip}{7pt plus 2pt minus 2pt}
\setlist[enumerate]{label=(\alph*),leftmargin=*,itemsep=2pt,topsep=3pt}
\setlist[itemize]{leftmargin=*,itemsep=2pt,topsep=3pt}
\newcommand{\R}{\mathbb R}
\newcommand{\contradiction}{\mathord{\text{\large\ding{56}}}}
\newcommand{\norm}[1]{\left\lVert#1\right\rVert}
\newcommand{\ip}[2]{\left\langle#1,#2\right\rangle}
\newcommand{\one}{\mathbf 1}
\DeclareMathOperator{\tr}{tr}
\DeclareMathOperator{\row}{row}
\DeclareMathOperator{\col}{col}
\DeclareMathOperator{\diag}{diag}
\DeclareMathOperator{\rank}{rank}
\DeclareMathOperator\prb{\mathsf{P}}
\DeclareMathOperator\expc{\mathsf{E}}
\DeclareMathOperator\var{var}
\DeclareMathOperator\cov{cov}
\DeclareMathOperator\cor{cor}
\DeclareMathOperator\argmin{argmin}
\DeclareMathOperator\argmax{argmax}
\DeclareMathOperator{\Span}{span}
\DeclareMathOperator{\nullity}{nullity}
\DeclareMathOperator{\im}{im}
\DeclareMathOperator{\id}{id}
\DeclareMathOperator{\sgn}{sgn}
\DeclareMathOperator{\inv}{inv}
\theoremstyle{definition}
\newtheorem{theorem}{Theorem}
\renewcommand{\thetheorem}{\unitnumber.\arabic{theorem}}
\newtheorem{proposition}[theorem]{Proposition}
\newtheorem{lemma}[theorem]{Lemma}
\newtheorem{corollary}[theorem]{Corollary}
\newtheorem{definition}[theorem]{Definition}
\newtheorem{example}[theorem]{Example}
\newtheorem{remark}[theorem]{Remark}
\newtheorem{exercise}{Exercise}
\renewcommand{\theexercise}{\unitnumber.\arabic{exercise}}
\newtheorem{problem}{Problem}
\newcommand{\citeat}[2]{\if\relax\detokenize{#2}\relax\cite{#1}\else\cite[#2]{#1}\fi}
\renewcommand{\theproblem}{\unitnumber.\arabic{problem}}
\newenvironment{solution}{\par\noindent\textit{Solution.}\ }{\par\smallskip}
\AddToHook{env/exercise/before}{\needspace{8\baselineskip}}
\AddToHook{env/exercise/after}{\nopagebreak[4]}
\AddToHook{env/problem/before}{\needspace{8\baselineskip}}
\AddToHook{env/problem/after}{\nopagebreak[4]}
\AddToHook{cmd/section/before}{\setcounter{exercise}{0}}
\def\theHtheorem{\unitnumber.\arabic{theorem}}
\def\theHexercise{\unitnumber.\arabic{exercise}}
\def\theHproblem{\unitnumber.\arabic{problem}}

\begin{document}
\pagenumbering{alph}
\begin{titlepage}
\definecolor{coveraccent}{RGB}{31,73,125}
\centering
\vspace*{24mm}
{\color{coveraccent}\rule{\linewidth}{1.6pt}\par}
\vspace{10mm}
{\fontsize{46}{54}\selectfont\bfseries Linear Algebra\par}
\vspace{7mm}
{\LARGE\scshape\textls[120]{And Its Statistical Applications}\par}
\vspace{9mm}
{\color{coveraccent}\rule{\linewidth}{0.6pt}\par}
\vspace{\stretch{1}}
\begin{tikzpicture}[x={(1.05cm,0cm)},y={(0.6cm,0.36cm)},z={(0cm,1.1cm)},
  >={Stealth[length=3.6mm,width=2.5mm]},font=\LARGE]
  \coordinate (O) at (0,0,0);
  \coordinate (P) at (6,4,0);
  \coordinate (Y) at (6,4,5);
  \filldraw[fill=coveraccent!7,draw=coveraccent!45,line width=0.7pt]
    (-1.5,-2,0) -- (9,-2,0) -- (9,6.5,0) -- (-1.5,6.5,0) -- cycle;
  \node[coveraccent] at (8,4.3,0) {\(W\)};
  \draw[line width=0.6pt] (5.626,3.75,0) -- (5.626,3.75,0.45) -- (6,4,0.45);
  \draw[->,line width=1.3pt] (O) -- (Y) node[pos=0.55,above left] {\(y\)};
  \draw[->,line width=1.3pt,coveraccent] (O) -- (P) node[pos=0.55,below right] {\(P_Wy\)};
  \draw[->,line width=1.3pt,densely dashed] (P) -- (Y) node[pos=0.5,right] {\(y-P_Wy\)};
  \fill (O) circle (1.6pt) node[below left] {\(0\)};
\end{tikzpicture}
\vspace{\stretch{1.4}}
\end{titlepage}
\pagenumbering{arabic}

\clearpage
\def\unitnumber{1}
\setcounter{theorem}{0}\setcounter{problem}{0}
\section*{Chapter 1. Matrices}

\subsection*{Matrices and products}

\begin{definition}[\citeat{hoffman-kunze}{\S1.3}, \citeat{searle-khuri}{Ch.~2}, \citeat{bapat}{\S1.1}]\label{r1-def-matrix}
Let \(m,n\) be positive integers. An \(m\times n\) \textbf{matrix} \(A=(a_{ij})\) is a
rectangular array of real numbers, its \textbf{entries}; \(a_{ij}\) stands in row \(i\) and
column \(j\) (\(1\leqslant i\leqslant m\), \(1\leqslant j\leqslant n\)). The set of all
\(m\times n\) matrices is \(\R^{m\times n}\).
\begin{itemize}
\item \(A=B\) means that \(A\) and \(B\) have the same size and \(a_{ij}=b_{ij}\) for all
\(i,j\).
\item \(\R^n=\R^{n\times1}\) is the set of \textbf{column vectors} \(x=(x_i)\), and
\(\R^{1\times n}\) the set of \textbf{row vectors}. A \(1\times1\) matrix is identified with
its entry.
\item \(\col_j(A)\in\R^m\) is the \(j\)th column and
\(\row_i(A)=(a_{i1},\dots,a_{in})\in\R^{1\times n}\) the \(i\)th row of \(A\). We write
\(a_j=\col_j(A)\), so \(A=[a_1\ \cdots\ a_n]\).
\item Sums and scalar multiples are taken entrywise:
\((A+B)_{ij}=a_{ij}+b_{ij}\), \((cA)_{ij}=ca_{ij}\) (\(c\in\R\)); \(-A=(-1)A\) and
\(A-B=A+(-B)\). They obey the laws of numbers, such as \(A+B=B+A\) and
\(c(A+B)=cA+cB\), because these hold entry by entry.
\item The \textbf{zero matrix} \(0\) has all entries \(0\). The \textbf{identity matrix} is
\(I_n=(\delta_{ij})\in\R^{n\times n}\), where \(\delta_{ij}=1\) if \(i=j\) and \(\delta_{ij}=0\)
if \(i\ne j\); its columns \(e_1,\dots,e_n\) are the \textbf{unit vectors} of \(\R^n\).
\end{itemize}
\end{definition}

In statistics, measurements of \(k\) variables on \(n\) units are stored in an \(n\times k\)
\textbf{data matrix} \((x_{ij})\): row \(i\) holds the measurements on unit \(i\), column \(j\)
holds the \(n\) values of variable \(j\), and \(x_{ij}\) is the value of variable \(j\) on unit
\(i\).

\begin{definition}[\citeat{hoffman-kunze}{\S1.5}, \citeat{lang}{Ch.~II, \S3}]\label{r1-def-product}
Let \(A\in\R^{m\times n}\) and \(B\in\R^{n\times p}\). The \textbf{product}
\(AB\in\R^{m\times p}\) is defined by
\[
(AB)_{ij}=\sum_{r=1}^n a_{ir}b_{rj}=a_{i1}b_{1j}+a_{i2}b_{2j}+\dots+a_{in}b_{nj}
\qquad(1\leqslant i\leqslant m,\ 1\leqslant j\leqslant p)
\]
It is defined only when \(A\) has as many columns as \(B\) has rows.
\end{definition}

\begin{theorem}[\citeat{hoffman-kunze}{\S1.5, Theorem~8}, \citeat{lang}{Ch.~II, Theorem~3.1}]\label{r1-thm-rules}
Whenever the sizes allow the operations,
\[
A(BC)=(AB)C,\quad A(B+C)=AB+AC,\quad (A+B)C=AC+BC,\quad A(cB)=c(AB)=(cA)B
\]
\[
I_mA=A=AI_n,\qquad 0A=0,\qquad A0=0\qquad(A\in\R^{m\times n})
\]
\end{theorem}
\begin{proof}
Let \(A\in\R^{m\times n}\). In each rule both sides have the same size, so by
Definition~\ref{r1-def-matrix} it suffices to show that they have the same \((i,j)\) entry.
Every entry below is computed from Definition~\ref{r1-def-product} and from the entrywise
definition of sums and scalar multiples.
\begin{itemize}
\item \(A(BC)=(AB)C\) for \(B\in\R^{n\times p}\), \(C\in\R^{p\times q}\). Since
\((BC)_{rj}=\sum_{s=1}^p b_{rs}c_{sj}\),
\[
[A(BC)]_{ij}=\sum_{r=1}^n a_{ir}(BC)_{rj}=\sum_{r=1}^n a_{ir}\Bigl(\sum_{s=1}^p b_{rs}c_{sj}\Bigr)
=\sum_{r=1}^n\sum_{s=1}^p a_{ir}b_{rs}c_{sj}
\]
Since \((AB)_{is}=\sum_{r=1}^n a_{ir}b_{rs}\),
\[
[(AB)C]_{ij}=\sum_{s=1}^p (AB)_{is}c_{sj}=\sum_{s=1}^p\Bigl(\sum_{r=1}^n a_{ir}b_{rs}\Bigr)c_{sj}
=\sum_{s=1}^p\sum_{r=1}^n a_{ir}b_{rs}c_{sj}
\]
Both double sums add the same \(np\) numbers \(a_{ir}b_{rs}c_{sj}\)
(\(1\leqslant r\leqslant n\), \(1\leqslant s\leqslant p\)), only in a different order, so they
are equal.
\item \(A(B+C)=AB+AC\) for \(B,C\in\R^{n\times p}\):
\[
[A(B+C)]_{ij}=\sum_{r=1}^n a_{ir}(b_{rj}+c_{rj})=\sum_{r=1}^n a_{ir}b_{rj}+\sum_{r=1}^n a_{ir}c_{rj}
=(AB)_{ij}+(AC)_{ij}=(AB+AC)_{ij}
\]
\item \((A+B)C=AC+BC\) for \(B\in\R^{m\times n}\), \(C\in\R^{n\times p}\):
\[
[(A+B)C]_{ij}=\sum_{r=1}^n (a_{ir}+b_{ir})c_{rj}=\sum_{r=1}^n a_{ir}c_{rj}+\sum_{r=1}^n b_{ir}c_{rj}
=(AC+BC)_{ij}
\]
\item \(A(cB)=c(AB)=(cA)B\) for \(B\in\R^{n\times p}\), \(c\in\R\):
\[
\begin{gathered}
[A(cB)]_{ij}=\sum_{r=1}^n a_{ir}(cb_{rj})=c\sum_{r=1}^n a_{ir}b_{rj}=[c(AB)]_{ij}\\
[(cA)B]_{ij}=\sum_{r=1}^n (ca_{ir})b_{rj}=c\sum_{r=1}^n a_{ir}b_{rj}=[c(AB)]_{ij}
\end{gathered}
\]
\item \(I_mA=A=AI_n\). In \((I_mA)_{ij}=\sum_{r=1}^m\delta_{ir}a_{rj}\) every term with
\(r\ne i\) is \(0\), because \(\delta_{ir}=0\), and the term with \(r=i\) is
\(1\cdot a_{ij}\); so \((I_mA)_{ij}=a_{ij}\). In the same way
\((AI_n)_{ij}=\sum_{r=1}^n a_{ir}\delta_{rj}=a_{ij}\).
\item \(0A=0\) and \(A0=0\): every term of \((0A)_{ij}=\sum_r 0\cdot a_{rj}\) and of
\((A0)_{ij}=\sum_r a_{ir}\cdot0\) is \(0\).\qedhere
\end{itemize}
\end{proof}
By the first rule, a product of several matrices needs no parentheses. A \(1\times1\) factor in a product acts as
a number: for \(M\in\R^{m\times1}\), \(N\in\R^{1\times p}\) and a \(1\times1\) matrix \(c\), identified with its entry
(Definition~\ref{r1-def-matrix}), Definition~\ref{r1-def-product} gives \((Mc)_{i1}=m_{i1}c\) and
\((cN)_{1j}=cn_{1j}\), so the products \(Mc\) and \(cN\) are the scalar multiples \(cM\) and \(cN\), and hence
\(McN=M(cN)=c(MN)\) by the rules above. For \(A\in\R^{n\times n}\) we put \(A^0=I_n\) and \(A^{k+1}=A^kA\)
(\(k\geqslant0\)).

\begin{theorem}[\citeat{harville}{(2.2.7)--(2.2.11)}, \citeat{hoffman-kunze}{\S1.5}, \citeat{bapat}{\S1.1}]\label{r1-thm-views}
Let \(A\in\R^{m\times n}\) have columns \(a_1,\dots,a_n\) and \(B\in\R^{n\times p}\) columns
\(b_1,\dots,b_p\).
\begin{enumerate}
\item \(AB=[Ab_1\ \cdots\ Ab_p]\), that is, \(\col_j(AB)=Ab_j\).
\item For \(x\in\R^n\), \(Ax=x_1a_1+x_2a_2+\dots+x_na_n\). In particular \(Ae_j=a_j\).
\item \(\row_i(AB)=\row_i(A)B=a_{i1}\row_1(B)+\dots+a_{in}\row_n(B)\).
\item \((AB)_{ij}=\row_i(A)\,b_j\), and \(AB=a_1\row_1(B)+\dots+a_n\row_n(B)\), a sum of \(n\)
products of a column and a row (the \textbf{outer product expansion}).
\end{enumerate}
\end{theorem}
\begin{proof}
Write \(B=(b_{rj})\). The column \(a_r\) of \(A\) has entries
\(a_{1r},\dots,a_{mr}\), the column \(b_j\) of \(B\) has entries \(b_{1j},\dots,b_{nj}\), and
\(\row_r(B)=(b_{r1},\dots,b_{rp})\). By Definition~\ref{r1-def-product},
\[
(AB)_{ij}=\sum_{r=1}^n a_{ir}b_{rj}=a_{i1}b_{1j}+a_{i2}b_{2j}+\dots+a_{in}b_{nj}
\]
Each part compares entries with this sum (Definition~\ref{r1-def-matrix}).
\begin{itemize}
\item (a) Fix \(j\). The product \(Ab_j\) of \(A\in\R^{m\times n}\) and
\(b_j\in\R^{n\times1}\) lies in \(\R^m\), and by Definition~\ref{r1-def-product} its
\(i\)th entry is
\[
(Ab_j)_i=\sum_{r=1}^n a_{ir}(b_j)_r=\sum_{r=1}^n a_{ir}b_{rj}=(AB)_{ij}\qquad(1\leqslant i\leqslant m)
\]
So the vector \(Ab_j\) and \(\col_j(AB)\), whose entries are
\((AB)_{1j},\dots,(AB)_{mj}\), have the same entries:
\[
\col_j(AB)=\begin{pmatrix}a_{11}b_{1j}+\dots+a_{1n}b_{nj}\\\vdots\\a_{m1}b_{1j}+\dots+a_{mn}b_{nj}\end{pmatrix}
=\begin{pmatrix}a_{11}&\cdots&a_{1n}\\\vdots&&\vdots\\a_{m1}&\cdots&a_{mn}\end{pmatrix}
\begin{pmatrix}b_{1j}\\\vdots\\b_{nj}\end{pmatrix}=Ab_j
\]
\item (b) By Definition~\ref{r1-def-product} (with \(p=1\)), the \(i\)th entry of
\(Ax\) is \(a_{i1}x_1+\dots+a_{in}x_n\). By the entrywise definition of scalar multiples and
sums (Definition~\ref{r1-def-matrix}), the \(i\)th entry of \(x_1a_1+\dots+x_na_n\) is
\(x_1a_{i1}+\dots+x_na_{in}\), the same number. Hence
\[
Ax=\begin{pmatrix}a_{11}x_1+\dots+a_{1n}x_n\\\vdots\\a_{m1}x_1+\dots+a_{mn}x_n\end{pmatrix}
=x_1\begin{pmatrix}a_{11}\\\vdots\\a_{m1}\end{pmatrix}+\dots+x_n\begin{pmatrix}a_{1n}\\\vdots\\a_{mn}\end{pmatrix}
=x_1a_1+\dots+x_na_n
\]
For \(x=e_j\) we have \(x_j=1\) and \(x_r=0\) for \(r\ne j\), so \(Ae_j=1\cdot a_j=a_j\).
\item (c) Fix \(i\). The product of \(\row_i(A)=(a_{i1},\dots,a_{in})\in\R^{1\times n}\) and
\(B\in\R^{n\times p}\) lies in \(\R^{1\times p}\), and by Definition~\ref{r1-def-product} its
\(j\)th entry is \(\sum_{r=1}^n a_{ir}b_{rj}=(AB)_{ij}\), the \(j\)th entry of \(\row_i(AB)\).
So \(\row_i(AB)=\row_i(A)B\). The \(j\)th entry of the row vector
\(a_{i1}\row_1(B)+\dots+a_{in}\row_n(B)\) is \(a_{i1}b_{1j}+\dots+a_{in}b_{nj}\), the same
sum. Hence
\[
\row_i(AB)=\Bigl(\sum_{r=1}^n a_{ir}b_{r1},\ \dots,\ \sum_{r=1}^n a_{ir}b_{rp}\Bigr)
=\sum_{r=1}^n a_{ir}\,(b_{r1},\dots,b_{rp})=\sum_{r=1}^n a_{ir}\row_r(B)
\]
\item (d) The product of \(\row_i(A)\in\R^{1\times n}\) and \(b_j\in\R^{n\times1}\) is the
\(1\times1\) matrix with entry \(\sum_{r=1}^n a_{ir}b_{rj}=(AB)_{ij}\). Next fix \(r\). The
product of \(a_r\in\R^{m\times1}\) and \(\row_r(B)\in\R^{1\times p}\) is the \(m\times p\)
matrix whose \((i,j)\) entry is the one-term sum \(a_{ir}b_{rj}\)
(Definition~\ref{r1-def-product} with inner size \(1\)):
\[
a_r\row_r(B)=\begin{pmatrix}a_{1r}\\\vdots\\a_{mr}\end{pmatrix}(b_{r1},\ \dots,\ b_{rp})
=\begin{pmatrix}a_{1r}b_{r1}&\cdots&a_{1r}b_{rp}\\\vdots&&\vdots\\a_{mr}b_{r1}&\cdots&a_{mr}b_{rp}\end{pmatrix}
\]
Adding these \(n\) matrices entry by entry, the \((i,j)\) entry of
\(a_1\row_1(B)+\dots+a_n\row_n(B)\) is \(a_{i1}b_{1j}+\dots+a_{in}b_{nj}=(AB)_{ij}\).\qedhere
\end{itemize}
\end{proof}

\begin{example}\label{r1-exm-views}
Let
\[
A=\begin{pmatrix}1&2&0&1\\0&1&3&1\\2&0&1&1\end{pmatrix}\in\R^{3\times4},\qquad
B=\begin{pmatrix}1&0\\0&1\\1&1\\2&-1\end{pmatrix}\in\R^{4\times2}
\]
By Definition~\ref{r1-def-product}, each entry of \(AB\in\R^{3\times2}\) is a sum of four
products:
\[
\begin{aligned}
(AB)_{11}&=1\cdot1+2\cdot0+0\cdot1+1\cdot2=3, &\quad (AB)_{12}&=1\cdot0+2\cdot1+0\cdot1+1\cdot(-1)=1\\
(AB)_{21}&=0\cdot1+1\cdot0+3\cdot1+1\cdot2=5, &\quad (AB)_{22}&=0\cdot0+1\cdot1+3\cdot1+1\cdot(-1)=3\\
(AB)_{31}&=2\cdot1+0\cdot0+1\cdot1+1\cdot2=5, &\quad (AB)_{32}&=2\cdot0+0\cdot1+1\cdot1+1\cdot(-1)=0
\end{aligned}
\]
so
\[
AB=\begin{pmatrix}3&1\\5&3\\5&0\end{pmatrix}
\]
The three views of Theorem~\ref{r1-thm-views} produce the same matrix.
\begin{itemize}
\item Columns, by (a) and (b): column \(j\) of \(AB\) is the combination of the four columns
of \(A\) with coefficients from column \(j\) of \(B\),
\[
\begin{aligned}
Ab_1&=1\begin{pmatrix}1\\0\\2\end{pmatrix}+0\begin{pmatrix}2\\1\\0\end{pmatrix}
+1\begin{pmatrix}0\\3\\1\end{pmatrix}+2\begin{pmatrix}1\\1\\1\end{pmatrix}=\begin{pmatrix}3\\5\\5\end{pmatrix}\\
Ab_2&=0\begin{pmatrix}1\\0\\2\end{pmatrix}+1\begin{pmatrix}2\\1\\0\end{pmatrix}
+1\begin{pmatrix}0\\3\\1\end{pmatrix}-1\begin{pmatrix}1\\1\\1\end{pmatrix}=\begin{pmatrix}1\\3\\0\end{pmatrix}
\end{aligned}
\]
\item Rows, by (c): row \(i\) of \(AB\) is the combination of the four rows of \(B\) with
coefficients from row \(i\) of \(A\),
\[
\begin{aligned}
\row_1(AB)&=1\,(1,0)+2\,(0,1)+0\,(1,1)+1\,(2,-1)=(3,\ 1)\\
\row_2(AB)&=0\,(1,0)+1\,(0,1)+3\,(1,1)+1\,(2,-1)=(5,\ 3)\\
\row_3(AB)&=2\,(1,0)+0\,(0,1)+1\,(1,1)+1\,(2,-1)=(5,\ 0)
\end{aligned}
\]
\item Outer products, by (d): \(AB\) is the sum of four \(3\times2\) matrices
\(a_r\row_r(B)\),
\[
\begin{aligned}
AB&=\begin{pmatrix}1\\0\\2\end{pmatrix}(1,\ 0)+\begin{pmatrix}2\\1\\0\end{pmatrix}(0,\ 1)
+\begin{pmatrix}0\\3\\1\end{pmatrix}(1,\ 1)+\begin{pmatrix}1\\1\\1\end{pmatrix}(2,\ -1)\\
&=\begin{pmatrix}1&0\\0&0\\2&0\end{pmatrix}+\begin{pmatrix}0&2\\0&1\\0&0\end{pmatrix}
+\begin{pmatrix}0&0\\3&3\\1&1\end{pmatrix}+\begin{pmatrix}2&-1\\2&-1\\2&-1\end{pmatrix}
=\begin{pmatrix}3&1\\5&3\\5&0\end{pmatrix}
\end{aligned}
\]
\end{itemize}
\end{example}

\begin{example}[\citeat{hoffman-kunze}{\S1.5, Example~10}]\label{r1-exm-products}
\begin{enumerate}
\item By Definition~\ref{r1-def-product},
\[
\begin{pmatrix}1&0&0&0\\-3&1&0&0\\2&0&1&0\\-1&0&0&1\end{pmatrix}
\begin{pmatrix}1&-1&2&1\\3&4&8&5\\-2&3&-1&0\\1&1&3&3\end{pmatrix}
=\begin{pmatrix}1&-1&2&1\\0&7&2&2\\0&1&3&2\\0&2&1&2\end{pmatrix}
\]
By Theorem~\ref{r1-thm-views}(c), its second row is \(-3\,(1,-1,2,1)+1\,(3,4,8,5)=(0,7,2,2)\), its third row is
\(2\,(1,-1,2,1)+1\,(-2,3,-1,0)=(0,1,3,2)\), and its fourth row is \(-1\,(1,-1,2,1)+1\,(1,1,3,3)=(0,2,1,2)\).
\item By Theorem~\ref{r1-thm-views}(b),
\[
\begin{pmatrix}2&1&0\\5&4&1\\1&3&2\\0&1&3\end{pmatrix}\begin{pmatrix}2\\-1\\1\end{pmatrix}
=2\begin{pmatrix}2\\5\\1\\0\end{pmatrix}-1\begin{pmatrix}1\\4\\3\\1\end{pmatrix}+1\begin{pmatrix}0\\1\\2\\3\end{pmatrix}
=\begin{pmatrix}3\\7\\1\\2\end{pmatrix}
\]
\item For \(N=\begin{pmatrix}0&1&0\\0&0&0\\0&0&0\end{pmatrix}\) and
\(M=\begin{pmatrix}1&-5&2\\2&3&4\\9&-1&3\end{pmatrix}\),
\[
NM=\begin{pmatrix}2&3&4\\0&0&0\\0&0&0\end{pmatrix},\qquad
MN=\begin{pmatrix}0&1&0\\0&2&0\\0&9&0\end{pmatrix}
\]
Row \(1\) of \(NM\) is row \(2\) of \(M\), and column \(2\) of \(MN\) is column \(1\) of \(M\).
So \(NM\ne MN\), although both are \(3\times3\).
\end{enumerate}
\end{example}

\begin{remark}[\citeat{hoffman-kunze}{\S1.5, Example~10}, \citeat{searle-khuri}{\S4.9}]\label{r1-rem-noncomm}
Matrix products differ from products of numbers. In general \(AB\ne BA\)
(Example~\ref{r1-exm-products}(c)); square matrices with \(AB=BA\) are said to
\textbf{commute}. A product can be \(0\) although no factor is \(0\): for
\[
X=\begin{pmatrix}1&1&3\\2&2&6\\-1&-1&-3\end{pmatrix},\qquad X^2=0
\]
To see this, let \(u=\col_1(X)\), the column with entries \(u_1=1\), \(u_2=2\), \(u_3=-1\), and let \(w_1=1\),
\(w_2=1\), \(w_3=3\) be the entries of \(\row_1(X)=(1,1,3)\). Rows \(1,2,3\) of \(X\) are \(1,2,-1\) times
\(\row_1(X)\), that is, \(\row_i(X)=u_i\row_1(X)\); columns \(1,2,3\) of \(X\) are \(1,1,3\) times \(u\), that
is, \(\col_j(X)=w_ju\). By Theorem~\ref{r1-thm-views}(d) and the rule \(A(cB)=c(AB)=(cA)B\) of
Theorem~\ref{r1-thm-rules},
\[
\begin{gathered}
(X^2)_{ij}=\row_i(X)\col_j(X)=\bigl(u_i\row_1(X)\bigr)(w_ju)=u_iw_j\,\bigl(\row_1(X)\,u\bigr)\\
\row_1(X)\,u=1\cdot1+1\cdot2+3\cdot(-1)=0
\end{gathered}
\]
so every entry of \(X^2\) is \(0\). Hence \(AB=0\) does not force \(A=0\) or \(B=0\) (take
\(A=B=X\)), and \(AB=AC\) does not force \(B=C\) (take \(A=B=X\) and \(C=0\)).
\end{remark}

\subsection*{Transpose and trace}

\begin{definition}\label{r1-def-transpose}
The \textbf{transpose} of \(A\in\R^{m\times n}\) is \(A'\in\R^{n\times m}\) with
\((A')_{ij}=a_{ji}\): the rows of \(A'\) are the transposed columns of \(A\). We write
column vectors as \(x=(x_1,\dots,x_n)'\). A square \(A\) is \textbf{symmetric} if
\(A'=A\).
\end{definition}

\begin{theorem}[\citeat{lang}{Ch.~II, Theorem~3.3}, \citeat{searle-khuri}{Example~5.6}]\label{r1-thm-transpose}
For matrices of sizes for which the operations are defined and \(c\in\R\),
\[
(A')'=A,\quad(A+B)'=A'+B',\quad(cA)'=cA',\quad
(AB)'=B'A',\quad(ABC)'=C'B'A'
\]
For \(x,y\in\R^n\), \(x'y=y'x=\sum_{i=1}^nx_iy_i\); in particular
\(x'x=\sum_ix_i^2\). The identity matrix is symmetric, \(I_n'=I_n\), and for every \(B\in\R^{m\times n}\), \(B'B\)
and \(BB'\) are symmetric.
\end{theorem}
\begin{proof}
In each identity both sides have the same size; we compare their entries, using
\((A')_{ij}=a_{ji}\) (Definition~\ref{r1-def-transpose}).
\begin{itemize}
\item For \(A,B\in\R^{m\times n}\) and \(c\in\R\), Definitions~\ref{r1-def-transpose}
and~\ref{r1-def-matrix} give
\[
\begin{gathered}
((A')')_{ij}=(A')_{ji}=a_{ij},\qquad ((cA)')_{ij}=(cA)_{ji}=ca_{ji}=(cA')_{ij}\\
((A+B)')_{ij}=(A+B)_{ji}=a_{ji}+b_{ji}=(A'+B')_{ij}
\end{gathered}
\]
\item Let \(A\in\R^{m\times n}\) and \(B\in\R^{n\times p}\). Then \((AB)'\) and \(B'A'\) are
both \(p\times m\). For \(1\leqslant k\leqslant p\) and \(1\leqslant i\leqslant m\),
Definition~\ref{r1-def-product} gives
\[
(B'A')_{ki}=\sum_{j=1}^n(B')_{kj}(A')_{ji}=\sum_{j=1}^n b_{jk}a_{ij}=\sum_{j=1}^n a_{ij}b_{jk}
=(AB)_{ik}=((AB)')_{ki}
\]
\item For a third matrix \(C\) with \(ABC\) defined, the previous bullet applied first to the
pair \(AB,C\) and then to the pair \(A,B\) gives
\[
(ABC)'=((AB)C)'=C'(AB)'=C'B'A'
\]
\item For \(x,y\in\R^n=\R^{n\times1}\), \(x'\in\R^{1\times n}\) has entries
\((x')_{1i}=x_i\), so \(x'y\) is \(1\times1\) with entry \(\sum_{i=1}^n x_iy_i\). In the same
way \(y'x=\sum_{i=1}^n y_ix_i\), the same number since \(y_ix_i=x_iy_i\). With \(y=x\),
\(x'x=\sum_ix_i^2\).
\item \((I_n')_{ij}=\delta_{ji}=\delta_{ij}\) for all \(i,j\) (Definition~\ref{r1-def-matrix}), so \(I_n'=I_n\).
\item For \(B\in\R^{m\times n}\), the product rule and \((B')'=B\) give
\((B'B)'=B'(B')'=B'B\) and \((BB')'=(B')'B'=BB'\).\qedhere
\end{itemize}
\end{proof}

\begin{definition}[\citeat{harville}{\S5.1}]\label{r1-def-trace}
The \textbf{trace} of \(A\in\R^{n\times n}\) is \(\tr(A)=a_{11}+a_{22}+\dots+a_{nn}\).
\end{definition}

\begin{theorem}[\citeat{harville}{(5.1.3)--(5.1.5), Lemma~5.2.1, (5.2.9), (5.2.5), Lemma~5.3.1}]\label{r1-thm-trace}
\begin{enumerate}
\item \(\tr(A+B)=\tr(A)+\tr(B)\), \(\tr(cA)=c\tr(A)\) and \(\tr(A')=\tr(A)\) for
\(A,B\in\R^{n\times n}\).
\item \(\tr(AB)=\tr(BA)\) for \(A\in\R^{m\times n}\) and \(B\in\R^{n\times m}\).
\item \(\tr(ABC)=\tr(CAB)=\tr(BCA)\) for \(A\in\R^{m\times n}\), \(B\in\R^{n\times p}\),
\(C\in\R^{p\times m}\).
\item \(\tr(A'A)=\sum_{i=1}^m\sum_{j=1}^na_{ij}^2\) for \(A\in\R^{m\times n}\); hence
\(\tr(A'A)=0\) only if \(A=0\).
\end{enumerate}
\end{theorem}
\begin{proof}
\begin{itemize}
\item (a) The \(i\)th diagonal entries of \(A+B\), \(cA\) and \(A'\) are \(a_{ii}+b_{ii}\),
\(ca_{ii}\) and \((A')_{ii}=a_{ii}\) (Definitions~\ref{r1-def-matrix}
and~\ref{r1-def-transpose}). Hence, by Definition~\ref{r1-def-trace},
\[
\tr(A+B)=\sum_{i=1}^n(a_{ii}+b_{ii})=\sum_{i=1}^na_{ii}+\sum_{i=1}^nb_{ii},\qquad
\tr(cA)=\sum_{i=1}^nca_{ii}=c\sum_{i=1}^na_{ii},\qquad
\tr(A')=\sum_{i=1}^na_{ii}
\]
\item (b) Both \(AB\in\R^{m\times m}\) and \(BA\in\R^{n\times n}\) are square. By
Definitions~\ref{r1-def-trace} and~\ref{r1-def-product},
\[
\tr(AB)=\sum_{i=1}^m(AB)_{ii}=\sum_{i=1}^m\sum_{j=1}^na_{ij}b_{ji},\qquad
\tr(BA)=\sum_{j=1}^n(BA)_{jj}=\sum_{j=1}^n\sum_{i=1}^mb_{ji}a_{ij}
\]
Both double sums add the same \(mn\) numbers \(a_{ij}b_{ji}\)
(\(1\leqslant i\leqslant m\), \(1\leqslant j\leqslant n\)), so they are equal.
\item (c) Since \(AB\in\R^{m\times p}\) and \(C\in\R^{p\times m}\), (b) applied to the pair
\(AB,C\) gives \(\tr((AB)C)=\tr(C(AB))\), that is, \(\tr(ABC)=\tr(CAB)\). Since
\(A\in\R^{m\times n}\) and \(BC\in\R^{n\times m}\), (b) applied to the pair \(A,BC\) gives
\(\tr(A(BC))=\tr((BC)A)\), that is, \(\tr(ABC)=\tr(BCA)\). Parentheses are dropped by
Theorem~\ref{r1-thm-rules}.
\item (d) The matrix \(A'A\) is \(n\times n\), and its \(j\)th diagonal entry is
\((A'A)_{jj}=\sum_{i=1}^m(A')_{ji}a_{ij}=\sum_{i=1}^ma_{ij}^2\). Summing over
\(j=1,\dots,n\) gives \(\tr(A'A)=\sum_{i=1}^m\sum_{j=1}^na_{ij}^2\). Every term
\(a_{ij}^2\) is \(\geqslant0\). If some \(a_{kl}\ne0\), then
\(\tr(A'A)\geqslant a_{kl}^2>0\). So \(\tr(A'A)=0\) forces every \(a_{ij}=0\), that is,
\(A=0\).\qedhere
\end{itemize}
\end{proof}

\subsection*{Special matrices; means and sums of squares}

\begin{definition}[\citeat{searle-khuri}{\S4.7.8}, \citeat{harville}{\S1.3g}]\label{r1-def-diag}
For \(d_1,\dots,d_n\in\R\), \(\diag(d_1,\dots,d_n)\) is the \(n\times n\) matrix with diagonal entries
\(d_1,\dots,d_n\) and all other entries \(0\); a matrix of this form is \textbf{diagonal}. A square matrix \(A\) is
\textbf{upper triangular} if \(a_{ij}=0\) for \(i>j\), and \textbf{lower triangular} if \(a_{ij}=0\) for \(i<j\).
\end{definition}

\begin{proposition}[\citeat{searle-khuri}{\S4.7.8}, \citeat{harville}{Lemma~1.3.1}]\label{r1-prop-diag}
Let \(D=\diag(d_1,\dots,d_n)\) and \(F=\diag(f_1,\dots,f_n)\).
\begin{enumerate}
\item Left multiplication by a diagonal matrix scales the rows, right multiplication scales the columns: for a
matrix \(A\) with \(n\) rows, \((DA)_{ij}=d_ia_{ij}\) and \(\row_i(DA)=d_i\row_i(A)\); for a matrix \(A\) with \(n\)
columns, \((AD)_{ij}=a_{ij}d_j\) and \(\col_j(AD)=d_j\col_j(A)\).
\item \(\diag(d_1,\dots,d_n)\diag(f_1,\dots,f_n)=\diag(d_1f_1,\dots,d_nf_n)\).
\item If \(A,B\in\R^{n\times n}\) are upper (lower) triangular, then \(AB\) is upper (lower) triangular, with
diagonal entries \(a_{11}b_{11},\dots,a_{nn}b_{nn}\).
\end{enumerate}
\end{proposition}
\begin{proof}
(a) The \((i,r)\) entry of \(D\) is \(d_i\) if \(r=i\) and \(0\) if \(r\ne i\), so in each sum of
Definition~\ref{r1-def-product} only one term survives:
\[
(DA)_{ij}=\sum_{r=1}^n D_{ir}a_{rj}=d_ia_{ij},\qquad (AD)_{ij}=\sum_{r=1}^n a_{ir}D_{rj}=a_{ij}d_j
\]
So row \(i\) of \(DA\) is \(d_i\) times row \(i\) of \(A\), and column \(j\) of \(AD\) is \(d_j\) times column \(j\)
of \(A\).

(b) By (a), \((DF)_{ij}=d_iF_{ij}\) is \(d_if_i\) for \(i=j\) and \(0\) for \(i\ne j\), that is,
\(DF=\diag(d_1f_1,\dots,d_nf_n)\).

(c) For upper triangular \(A,B\in\R^{n\times n}\), the term \(a_{ik}b_{kj}\) is \(0\) if \(k<i\) (then
\(a_{ik}=0\)) and if \(k>j\) (then \(b_{kj}=0\)); so only the terms with \(i\leqslant k\leqslant j\) remain:
\[
(AB)_{ij}=\sum_{k=1}^n a_{ik}b_{kj}=\sum_{k=i}^{j}a_{ik}b_{kj}
\]
For \(i>j\) no \(k\) satisfies \(i\leqslant k\leqslant j\), so \((AB)_{ij}=0\); for \(i=j\) only \(k=i\) remains,
so \((AB)_{ii}=a_{ii}b_{ii}\). For lower triangular \(A,B\) the term \(a_{ik}b_{kj}\) is \(0\) if \(k>i\) or
\(k<j\), so only the terms with \(j\leqslant k\leqslant i\) remain, and \((AB)_{ij}=0\) for \(i<j\),
\((AB)_{ii}=a_{ii}b_{ii}\).
\end{proof}

\begin{definition}[\citeat{searle-khuri}{\S\S5.2, 5.3}]\label{r1-def-ones}
The \textbf{summing vector} \(\one_n\in\R^n\) has all entries \(1\). Put
\[
J_n=\one_n\one_n'\in\R^{n\times n},\qquad \bar J_n=\tfrac1nJ_n,\qquad C_n=I_n-\bar J_n
\]
Every entry of \(J_n\) is \(1\), and \(C_n\) is the \textbf{centering matrix}. Subscripts are
dropped when the size is clear. A square matrix \(A\) with \(A^2=A\) is \textbf{idempotent}.
\end{definition}

\begin{proposition}[\citeat{searle-khuri}{\S5.2, (5.13)--(5.18)}]\label{r1-prop-centering}
\(\one'\one=n\), \(J_n^2=nJ_n\), and \(\bar J_n\) and \(C=C_n\) are symmetric and
idempotent:
\[
\bar J_n'=\bar J_n=\bar J_n^2,\qquad C=C'=C^2,\qquad C\one=0
\]
\end{proposition}
\begin{proof}
Write \(\one=\one_n\), \(J=J_n\), \(\bar J=\bar J_n\), \(I=I_n\).
\begin{itemize}
\item By Theorem~\ref{r1-thm-transpose}, \(\one'\one=\sum_{i=1}^n1\cdot1=n\).
\item By Theorem~\ref{r1-thm-rules} (associativity) and the paragraph after Theorem~\ref{r1-thm-rules} (the
\(1\times1\) factor \(\one'\one=n\) acts as a number),
\[
J^2=(\one\one')(\one\one')=\one(\one'\one)\one'=\one\,n\,\one'=n\,\one\one'=nJ,\qquad
\bar J^2=\tfrac1{n^2}J^2=\tfrac1{n^2}\,nJ=\tfrac1nJ=\bar J
\]
\item By Theorem~\ref{r1-thm-transpose}, \(J'=(\one\one')'=(\one')'\one'=\one\one'=J\), so
\(\bar J'=\frac1nJ'=\bar J\); and \(I'=I\), by the same theorem. Hence \(C'=I'-\bar J'=I-\bar J=C\).
\item By Theorem~\ref{r1-thm-rules} and \(\bar J^2=\bar J\),
\[
C^2=(I-\bar J)(I-\bar J)=I\cdot I-I\bar J-\bar JI+\bar J^2=I-2\bar J+\bar J=I-\bar J=C
\]
\item By Theorem~\ref{r1-thm-rules}, the paragraph after it and \(\one'\one=n\),
\(C\one=\one-\frac1n\one(\one'\one)=\one-\frac1n\,n\,\one=0\).\qedhere
\end{itemize}
\end{proof}

\begin{example}[\citeat{searle-khuri}{Example~5.6}]\label{r1-exm-centering}
For data \(x=(x_1,\dots,x_n)'\), \(\one'x=x_1+\dots+x_n\) (Theorem~\ref{r1-thm-transpose}), so
the mean is \(\bar x=\frac1n\one'x\), and \(\one'x=n\bar x\). By Theorem~\ref{r1-thm-rules},
\[
Cx=(I-\tfrac1n\one\one')x=x-\tfrac1n\one(\one'x)=x-\bar x\one=(x_1-\bar x,\ \dots,\ x_n-\bar x)'
\]
the vector of \textbf{deviations} from the mean. Moreover \(x'\one=\one'x=n\bar x\)
(Theorem~\ref{r1-thm-transpose}), so
\[
x'Cx=x'x-\tfrac1n(x'\one)(\one'x)=\sum_{i=1}^nx_i^2-\tfrac1n(n\bar x)(n\bar x)=\sum_{i=1}^nx_i^2-n\bar x^2
\]
Expanding each square and using \(\sum_ix_i=n\bar x\),
\[
\sum_{i=1}^n(x_i-\bar x)^2=\sum_{i=1}^n(x_i^2-2\bar xx_i+\bar x^2)
=\sum_{i=1}^nx_i^2-2\bar x\sum_{i=1}^nx_i+n\bar x^2
=\sum_{i=1}^nx_i^2-2n\bar x^2+n\bar x^2=\sum_{i=1}^nx_i^2-n\bar x^2
\]
Hence the sum of squares about the mean is
\[
\sum_{i=1}^n(x_i-\bar x)^2=x'x-n\bar x^2=x'Cx
\]
\end{example}

\begin{example}[cf.\ \citeat{searle-khuri}{Example~11.1}]\label{r1-exm-data}
Four units (\(n=4\)) are measured on two variables (\(k=2\)); \(x_{ij}\) is the value of
variable \(j\) on unit \(i\):
\[
X_0=(x_{ij})=\begin{pmatrix}2&3\\0&3\\2&1\\0&1\end{pmatrix}\in\R^{4\times2},\qquad
\one'X_0=(4,\ 8),\qquad
\bar x=\tfrac14X_0'\one=\begin{pmatrix}1\\2\end{pmatrix}
\]
The vector \(\bar x=(\bar x_1,\bar x_2)'\) of column means \(\bar x_j=\frac1n\sum_ix_{ij}\) is
the \textbf{mean vector}. By Theorem~\ref{r1-thm-views}(a),
\(\col_j(CX_0)=C\col_j(X_0)\), which by Example~\ref{r1-exm-centering} is the vector of
deviations of variable \(j\) from its mean \(\bar x_j\). In matrix form, by
Theorem~\ref{r1-thm-rules} and \(\one'X_0=(X_0'\one)'=n\bar x'\)
(Theorem~\ref{r1-thm-transpose}),
\[
CX_0=X_0-\tfrac1n\one(\one'X_0)=X_0-\one\bar x'
\]
By Proposition~\ref{r1-prop-centering}, \(C=C^2=C'C\), so
\(X_0'CX_0=X_0'C'CX_0=(CX_0)'(CX_0)\) (Theorem~\ref{r1-thm-transpose}). The centered data
matrix and the matrix of \textbf{corrected sums of squares and products} are therefore
\[
CX_0=\begin{pmatrix}2-1&3-2\\0-1&3-2\\2-1&1-2\\0-1&1-2\end{pmatrix}=\begin{pmatrix}1&1\\-1&1\\1&-1\\-1&-1\end{pmatrix},\qquad
X_0'CX_0=(CX_0)'(CX_0)=\begin{pmatrix}4&0\\0&4\end{pmatrix}
\]
Here the \((j,l)\) entry of \((CX_0)'(CX_0)\) is \(\col_j(CX_0)'\col_l(CX_0)\)
(Theorem~\ref{r1-thm-views}(d)): \(1+1+1+1=4\) for \((1,1)\) and \((2,2)\), and
\(1-1-1+1=0\) for \((1,2)\) and \((2,1)\). The same matrix comes from the uncorrected sums: in
general, by Theorem~\ref{r1-thm-rules} and \(X_0'\one=n\bar x\),
\[
X_0'CX_0=X_0'(I-\tfrac1n\one\one')X_0=X_0'X_0-\tfrac1n(X_0'\one)(\one'X_0)
=X_0'X_0-\tfrac1n(n\bar x)(n\bar x')=X_0'X_0-n\bar x\bar x'
\]
and here
\[
X_0'X_0-4\bar x\bar x'=\begin{pmatrix}8&8\\8&20\end{pmatrix}-\begin{pmatrix}4&8\\8&16\end{pmatrix}
=\begin{pmatrix}4&0\\0&4\end{pmatrix}
\]
\end{example}

\subsection*{Unit vectors and the matrices \(E_{ij}\)}

\begin{proposition}[\citeat{searle-khuri}{\S5.1.5}, \citeat{harville}{(2.4.1)--(2.4.5), Lemma~2.3.2}]\label{r1-prop-units}
Let \(A\in\R^{m\times n}\). Write \(e_i\) for the unit vectors of \(\R^m\) or of \(\R^n\), whichever the product
requires, and put \(E_{ij}=e_ie_j'\in\R^{m\times n}\) (\(1\leqslant i\leqslant m\), \(1\leqslant j\leqslant n\)).
\begin{enumerate}
\item \(e_i'e_k=\delta_{ik}\), \(Ae_j=\col_j(A)\), \(e_i'A=\row_i(A)\) and \(e_i'Ae_j=a_{ij}\).
\item \(E_{ij}\) has \((i,j)\) entry \(1\) and all other entries \(0\), and
\[
A=\sum_{i=1}^me_i\row_i(A)=\sum_{j=1}^n\col_j(A)\,e_j'=\sum_{i=1}^m\sum_{j=1}^na_{ij}E_{ij}
\]
In particular \(I_m=e_1e_1'+\dots+e_me_m'\).
\item \(E_{ij}E_{kl}=\delta_{jk}E_{il}\), where \(E_{kl}=e_ke_l'\in\R^{n\times p}\) and \(E_{il}=e_ie_l'\in\R^{m\times p}\).
\item For \(B\in\R^{m\times n}\), \(A=B\) if and only if \(Ax=Bx\) for every \(x\in\R^n\).
\end{enumerate}
\end{proposition}
\begin{proof}
(a)
\begin{itemize}
\item \(e_i'e_k=\sum_r(e_i)_r(e_k)_r\) (Theorem~\ref{r1-thm-transpose}), and the term \((e_i)_r(e_k)_r\) is \(1\) if
\(r=i=k\) and \(0\) otherwise. So \(e_i'e_k=\delta_{ik}\).
\item \(Ae_j=\col_j(A)\) is Theorem~\ref{r1-thm-views}(b). By Theorem~\ref{r1-thm-views}(c),
\(e_i'A=(e_i)_1\row_1(A)+\dots+(e_i)_m\row_m(A)=\row_i(A)\).
\item Applied to the \(m\times1\) matrix \(\col_j(A)\), the last statement gives \(e_i'Ae_j=e_i'\col_j(A)=a_{ij}\).
\end{itemize}
(b)
\begin{itemize}
\item The \((k,l)\) entry of \(e_ie_j'\) is \((e_i)_k(e_j)_l\) (Definition~\ref{r1-def-product} with one term), which is
\(1\) for \(k=i\), \(l=j\) and \(0\) otherwise.
\item The outer product expansion (Theorem~\ref{r1-thm-views}(d)) of \(A=I_mA\), whose first factor has the columns
\(e_1,\dots,e_m\), gives \(A=\sum_ie_i\row_i(A)\). As \(I_n'=I_n\) (Theorem~\ref{r1-thm-transpose}), the rows of
\(I_n\) are \(e_1',\dots,e_n'\), and
the expansion of \(A=AI_n\) gives \(A=\sum_j\col_j(A)e_j'\).
\item By Theorem~\ref{r1-thm-views}(b), \(\col_j(A)=I_m\col_j(A)=\sum_ia_{ij}e_i\). Substituting and using
Theorem~\ref{r1-thm-rules}, \(A=\sum_j\sum_ia_{ij}e_ie_j'\). For \(A=I_m\), \(\row_i(I_m)=e_i'\), and the first
expansion reads \(I_m=\sum_ie_ie_i'\).
\end{itemize}
(c) \(E_{ij}E_{kl}=e_i(e_j'e_k)e_l'=\delta_{jk}e_ie_l'\) by Theorem~\ref{r1-thm-rules} and (a), the \(1\times1\)
factor \(e_j'e_k\) being a number.

(d) If \(A=B\), then \(Ax=Bx\) for every \(x\). Conversely, \(x=e_j\) gives \(\col_j(A)=Ae_j=Be_j=\col_j(B)\) for
every \(j\), by (a); so \(A=B\).
\end{proof}

For a data matrix \(X_0\in\R^{n\times k}\), \(e_i'X_0\) is the row of measurements on unit \(i\), \(X_0e_j\) the column
of values of variable \(j\), and \(x_{ij}=e_i'X_0e_j\). Several unit vectors select several variables or units: for
instance, if \(k\geqslant3\) and \(n\geqslant4\), then by
Theorem~\ref{r1-thm-views}(a), \(X_0[e_3\ e_1]=[\col_3(X_0)\ \col_1(X_0)]\) keeps variables \(3\) and \(1\), in this
order; by Theorem~\ref{r1-thm-views}(c), the rows of \([e_2\ e_4]'X_0\) are \(\row_2(X_0)\) and \(\row_4(X_0)\).

\subsection*{Powers and polynomials of a square matrix}

\begin{proposition}[\citeat{lang}{Ch.~II, \S3}, \citeat{abadir-magnus}{Exercises~2.33, 2.34}, \citeat{shakarchi}{Ch.~II, \S3, Exercise~30}, \citeat{harville}{Lemma~1.3.1}]\label{r1-prop-powers}
Let \(A\in\R^{n\times n}\), and let \(r,s\geqslant0\) be integers.
\begin{enumerate}
\item \(A^rA^s=A^{r+s}=A^sA^r\).
\item \((A^r)^s=A^{rs}\).
\item \((A')^r=(A^r)'\). In particular, every power of a symmetric matrix is symmetric.
\item \(\diag(d_1,\dots,d_n)^r=\diag(d_1^r,\dots,d_n^r)\) for all \(d_1,\dots,d_n\in\R\); here \(t^0=1\) for every
\(t\in\R\), also for \(t=0\).
\item If \(A\) is upper (lower) triangular, then \(A^r\) is upper (lower) triangular, with diagonal entries
\(a_{11}^r,\dots,a_{nn}^r\) (with \(t^0=1\) for every \(t\in\R\), as in (d)).
\end{enumerate}
\end{proposition}
\begin{proof}
(a) Induction on \(s\), with \(r\) fixed.
\begin{itemize}
\item \(s=0\): \(A^rA^0=A^rI_n=A^r\) (Theorem~\ref{r1-thm-rules}).
\item If \(A^rA^s=A^{r+s}\), then \(A^{r+s+1}=A^{r+s}A=(A^rA^s)A=A^r(A^sA)=A^rA^{s+1}\) (Theorem~\ref{r1-thm-rules}).
\item Exchanging the roles of \(r\) and \(s\), \(A^sA^r=A^{s+r}=A^{r+s}\).
\end{itemize}
(b) Induction on \(s\): \((A^r)^0=I_n=A^0\), and if \((A^r)^s=A^{rs}\), then
\((A^r)^{s+1}=(A^r)^sA^r=A^{rs}A^r=A^{rs+r}\) by (a).

(c)
\begin{itemize}
\item Induction on \(r\): \((A')^0=I_n=I_n'=(A^0)'\) (Theorem~\ref{r1-thm-transpose}). If \((A')^r=(A^r)'\), then
by Theorem~\ref{r1-thm-transpose}
and (a), \((A^{r+1})'=(A^rA)'=A'(A^r)'=A'(A')^r=(A')^{r+1}\).
\item If \(A'=A\), this reads \((A^r)'=A^r\).
\end{itemize}
(d) Induction on \(r\): \(\diag(d_1,\dots,d_n)^0=I_n=\diag(1,\dots,1)=\diag(d_1^0,\dots,d_n^0)\), and if the formula
holds for \(r\), then \(\diag(d_i)^{r+1}=\diag(d_i^r)\diag(d_i)=\diag(d_i^{r+1})\) by
Proposition~\ref{r1-prop-diag}(b).

(e) Induction on \(r\): \(A^0=I_n\) is upper triangular with diagonal entries \(1=a_{ii}^0\). If \(A^r\) is upper
triangular with diagonal entries \(a_{ii}^r\), then \(A^{r+1}=A^rA\) is a product of two upper triangular matrices,
so it is upper triangular with diagonal entries \(a_{ii}^ra_{ii}=a_{ii}^{r+1}\) (Proposition~\ref{r1-prop-diag}(c)).
The lower triangular case is the same.
\end{proof}

\begin{definition}[\citeat{lang}{Ch.~II, \S3, Exercise~39}, \citeat{searle-khuri}{\S5.3}, \citeat{harville}{Ch.~8, Exercise~2}]\label{r1-def-nilpotent}
A matrix \(A\in\R^{n\times n}\) is \textbf{nilpotent} if \(A^k=0\) for some integer \(k\geqslant1\), and
\textbf{involutory} if \(A^2=I_n\).
\end{definition}
For instance, the matrix \(X\) of Remark~\ref{r1-rem-noncomm} is nilpotent (\(X^2=0\)), and \(I_n\), \(-I_n\) and
\(\diag(1,-1,1,-1)\) are involutory.

\needspace{12\baselineskip}
\begin{proposition}[\citeat{harville-solutions}{Ch.~1, Exercise~8}]\label{r1-prop-strictly-triangular}
Let \(A\in\R^{n\times n}\) be \textbf{strictly upper triangular}, that is, \(a_{ij}=0\) for \(i\geqslant j\), and let
\(p\geqslant1\) be an integer.
\begin{enumerate}
\item The \((i,j)\) entry of \(A^p\) is \(0\) for \(1\leqslant i\leqslant n\) and \(1\leqslant j\leqslant\min(n,i+p-1)\).
\item For \(i\geqslant n-p+1\), the \(i\)th row of \(A^p\) is zero.
\item \(A^p=0\) for \(p\geqslant n\); so \(A\) is nilpotent.
\end{enumerate}
The same holds for a \textbf{strictly lower triangular} \(A\) (\(a_{ij}=0\) for \(i\leqslant j\)), with the
\((j,i)\) entry in (a) and the \(i\)th column in (b).
\end{proposition}
\begin{proof}
Write \(b_{ik}\) for the \((i,k)\) entry of \(A^p\).

(a) Induction on \(p\).
\begin{itemize}
\item \(p=1\): the range is \(j\leqslant\min(n,i)=i\), and \(a_{ij}=0\) for \(j\leqslant i\).
\item Suppose (a) holds for \(p\), and let \(j\leqslant\min(n,i+p)\). Since \(A^{p+1}=A^pA\),
\[
(A^{p+1})_{ij}=\sum_{k=1}^nb_{ik}a_{kj}=\sum_{k=1}^{\min(n,i+p-1)}0\cdot a_{kj}+\sum_{k=i+p}^nb_{ik}a_{kj}=0
\]
The first sum vanishes by the hypothesis on \(A^p\). The second is empty if \(i+p>n\); in it
\(k\geqslant i+p\geqslant j\), so \(a_{kj}=0\). This is (a) for \(p+1\).
\end{itemize}
(b) For \(i\geqslant n-p+1\), \(i+p-1\geqslant n\), so \(\min(n,i+p-1)=n\), and by (a) all \(n\) entries of row \(i\)
of \(A^p\) are \(0\).

(c) For \(p\geqslant n\), \(n-p+1\leqslant1\), so by (b) all rows of \(A^p\) are zero.

If \(A\) is strictly lower triangular, then \(A'\) is strictly upper triangular and \(A^p=\bigl((A')^p\bigr)'\)
(Proposition~\ref{r1-prop-powers}(c)); transposing (a)--(c) for \(A'\) gives the statements for \(A\).
\end{proof}

\begin{example}[\citeat{lang}{Ch.~II, \S3, Exercise~31}]\label{r1-exm-powers}
\begin{enumerate}
\item \(N=\begin{pmatrix}0&1&6\\0&0&4\\0&0&0\end{pmatrix}\) is strictly upper triangular. By
Theorem~\ref{r1-thm-views}(c), \(\row_1(N^2)=1\,\row_2(N)+6\,\row_3(N)=(0,0,4)\),
\(\row_2(N^2)=4\,\row_3(N)=0\) and \(\row_3(N^2)=0\); then \(\row_1(N^3)=4\,\row_3(N)=0\). So
\[
N^2=\begin{pmatrix}0&0&4\\0&0&0\\0&0&0\end{pmatrix},\qquad N^3=N^2N=\begin{pmatrix}0&0&0\\0&0&0\\0&0&0\end{pmatrix}
\]
as Proposition~\ref{r1-prop-strictly-triangular} says for \(n=3\). Here \(N^2\ne0\), so the exponent \(n\) in
Proposition~\ref{r1-prop-strictly-triangular}(c) cannot be lowered in general.
\item For the upper triangular
\[
U=\begin{pmatrix}1&1&0&1\\0&2&1&0\\0&0&-1&1\\0&0&0&2\end{pmatrix},\qquad
U^2=\begin{pmatrix}1&3&1&3\\0&4&1&1\\0&0&1&1\\0&0&0&4\end{pmatrix},\qquad
U^3=U^2U=\begin{pmatrix}1&7&2&8\\0&8&3&3\\0&0&-1&3\\0&0&0&8\end{pmatrix}
\]
for instance \(\row_1(U^2)=1\,\row_1(U)+1\,\row_2(U)+1\,\row_4(U)=(1,3,1,3)\) and
\(\row_1(U^3)=1\,\row_1(U)+3\,\row_2(U)+1\,\row_3(U)+3\,\row_4(U)=(1,7,2,8)\). Both powers are upper triangular, with
the diagonal entries \(1^2,2^2,(-1)^2,2^2\) and \(1^3,2^3,(-1)^3,2^3\) (Proposition~\ref{r1-prop-powers}(e)); and
\((U')^3=(U^3)'\) is lower triangular (Proposition~\ref{r1-prop-powers}(c)).
\end{enumerate}
\end{example}

\begin{definition}[\citeat{hoffman-kunze}{\S4.2}]\label{r1-def-poly}
A (real) \textbf{polynomial} \(f(t)=c_0+c_1t+\dots+c_mt^m\) is given by its coefficients \(c_0,\dots,c_m\in\R\). For
\(A\in\R^{n\times n}\),
\[
f(A)=c_0I_n+c_1A+\dots+c_mA^m
\]
Polynomials are added and multiplied as usual: for \(g(t)=d_0+d_1t+\dots+d_lt^l\), the product
\((fg)(t)=\sum_{i=0}^m\sum_{j=0}^lc_id_jt^{i+j}\) has the coefficient \(\sum_{i+j=k}c_id_j\) at \(t^k\); in
particular \(fg=gf\).
\end{definition}

\begin{proposition}[\citeat{hoffman-kunze}{\S4.2, Theorem~2}]\label{r1-prop-poly}
Let \(f\) and \(g\) be polynomials, \(c\in\R\) and \(A\in\R^{n\times n}\).
\begin{enumerate}
\item \((cf+g)(A)=cf(A)+g(A)\).
\item \((fg)(A)=f(A)g(A)\). Hence \(f(A)g(A)=g(A)f(A)\): polynomials in the same matrix commute.
\item \(f(A)'=f(A')\); in particular \(f(A)\) is symmetric if \(A\) is.
\end{enumerate}
\end{proposition}
\begin{proof}
Write \(f(t)=\sum_ic_it^i\) and \(g(t)=\sum_jd_jt^j\), with coefficients \(0\) beyond the degrees.

(a) \(cf+g\) has the coefficients \(cc_i+d_i\), so \((cf+g)(A)=\sum_i(cc_i+d_i)A^i=c\sum_ic_iA^i+\sum_id_iA^i\) by the
laws of Definition~\ref{r1-def-matrix}.

(b)
\begin{itemize}
\item By Definition~\ref{r1-def-poly} and (a), \((fg)(A)=\sum_i\sum_jc_id_jA^{i+j}\).
\item \(c_id_jA^{i+j}=(c_iA^i)(d_jA^j)\) by Proposition~\ref{r1-prop-powers}(a) and Theorem~\ref{r1-thm-rules}.
\item By the distributive laws (Theorem~\ref{r1-thm-rules}),
\(\sum_i\sum_j(c_iA^i)(d_jA^j)=\bigl(\sum_ic_iA^i\bigr)\bigl(\sum_jd_jA^j\bigr)=f(A)g(A)\).
\item As \(fg=gf\), \(f(A)g(A)=(fg)(A)=(gf)(A)=g(A)f(A)\).
\end{itemize}
(c) By Theorem~\ref{r1-thm-transpose} and Proposition~\ref{r1-prop-powers}(c),
\(f(A)'=\sum_ic_i(A^i)'=\sum_ic_i(A')^i=f(A')\).
\end{proof}

\begin{example}\label{r1-exm-poly}
Matrices \(aI_n+bJ_n\), with all diagonal entries \(a+b\) and all other entries \(b\), are frequent in statistics;
\(\bar J_n\) and \(C_n=I_n-\bar J_n\) are of this form. They are polynomials of degree at most \(1\) in \(J_n\), so any
two of them commute (Proposition~\ref{r1-prop-poly}(b)). Since \(J_n^2=nJ_n\) (Proposition~\ref{r1-prop-centering}),
\[
(aI_n+bJ_n)(cI_n+dJ_n)=acI_n+(ad+bc)J_n+bdJ_n^2=acI_n+(ad+bc+nbd)J_n
\]
is again of this form. For instance, \(C_n=I_n-\frac1nJ_n\) and \(J_n\) (\(a=1\), \(b=-\frac1n\), \(c=0\), \(d=1\))
give \(C_nJ_n=\bigl(1+0-n\cdot\frac1n\bigr)J_n=0\), and \(J_nC_n=C_nJ_n=0\). For \(n=4\) and \(a=b=c=d=1\),
\[
A=I_4+J_4=\begin{pmatrix}2&1&1&1\\1&2&1&1\\1&1&2&1\\1&1&1&2\end{pmatrix},\qquad
A^2=I_4+6J_4=\begin{pmatrix}7&6&6&6\\6&7&6&6\\6&6&7&6\\6&6&6&7\end{pmatrix}
\]
For \(f(t)=t^2-6t+5=(t-1)(t-5)\), Proposition~\ref{r1-prop-poly}(b) gives
\[
f(A)=(A-I_4)(A-5I_4)=J_4(J_4-4I_4)=J_4^2-4J_4=0
\]
So \(A^2=6A-5I_4\), and every power of \(A\) reduces to a combination of \(I_4\) and \(A\), by induction on \(k\):
\(A^1=A\), and if \(A^k=\alpha A+\beta I_4\), then \(A^{k+1}=A^kA=\alpha A^2+\beta A=(6\alpha+\beta)A-5\alpha I_4\)
(Theorem~\ref{r1-thm-rules}). For instance
\(A^3=AA^2=6A^2-5A=6(6A-5I_4)-5A=31A-30I_4\) (Proposition~\ref{r1-prop-powers}(a)), that is,
\[
A^3=\begin{pmatrix}32&31&31&31\\31&32&31&31\\31&31&32&31\\31&31&31&32\end{pmatrix}
\]
\end{example}

\subsection*{Symmetric, skew-symmetric and idempotent matrices}

\begin{definition}[\citeat{searle-khuri}{\S5.1.6}, \citeat{abadir-magnus}{Ch.~2}]\label{r1-def-skew}
A matrix \(A\in\R^{n\times n}\) is \textbf{skew-symmetric} if \(A'=-A\), that is, \(a_{ji}=-a_{ij}\) for all \(i,j\).
\end{definition}

\needspace{16\baselineskip}
\begin{proposition}[\citeat{abadir-magnus}{Exercises~2.22, 2.23}]\label{r1-prop-sym-skew}
Let \(A\in\R^{n\times n}\).
\begin{enumerate}
\item \(A+A'\) is symmetric and \(A-A'\) is skew-symmetric.
\item \(A=\frac12(A+A')+\frac12(A-A')\), the sum of a symmetric and a skew-symmetric matrix.
\item If \(A\) is skew-symmetric, then \(a_{ii}=0\) for every \(i\).
\item If \(B\in\R^{m\times m}\) is symmetric and \(P\in\R^{m\times n}\), then \(P'BP\in\R^{n\times n}\) is symmetric.
\end{enumerate}
\end{proposition}
\begin{proof}
(a) By Theorem~\ref{r1-thm-transpose}, \((A+A')'=A'+(A')'=A'+A\) and \((A-A')'=A'-A=-(A-A')\).

(b) By (a) and \((cM)'=cM'\) (Theorem~\ref{r1-thm-transpose}), \(\frac12(A+A')\) is symmetric and \(\frac12(A-A')\)
skew-symmetric; their sum is \(A\), since the terms \(\pm\frac12A'\) cancel.

(c) \(A'\) has the same diagonal entries as \(A\), so \(A'=-A\) gives \(a_{ii}=-a_{ii}\), that is, \(a_{ii}=0\).

(d) \((P'BP)'=P'B'(P')'=P'BP\) (Theorem~\ref{r1-thm-transpose}).
\end{proof}
For instance, the matrix \(X_0'C_nX_0\) of corrected sums of squares and products (Example~\ref{r1-exm-data}) is
symmetric by (d), since \(C_n\) is symmetric (Proposition~\ref{r1-prop-centering}).

\begin{example}\label{r1-exm-sym-skew}
For
\[
A=\begin{pmatrix}1&4&-1&2\\0&2&3&-1\\3&-1&5&3\\0&1&1&3\end{pmatrix},\qquad
A'=\begin{pmatrix}1&0&3&0\\4&2&-1&1\\-1&3&5&1\\2&-1&3&3\end{pmatrix}
\]
Proposition~\ref{r1-prop-sym-skew}(b) gives \(A=S+K\) with
\[
S=\tfrac12(A+A')=\begin{pmatrix}1&2&1&1\\2&2&1&0\\1&1&5&2\\1&0&2&3\end{pmatrix},\qquad
K=\tfrac12(A-A')=\begin{pmatrix}0&2&-2&1\\-2&0&2&-1\\2&-2&0&1\\-1&1&-1&0\end{pmatrix}
\]
\(S\) is symmetric, \(K\) is skew-symmetric with zero diagonal, and \(S+K=A\) entry by entry.
\end{example}

\begin{proposition}[\citeat{harville}{Lemma~10.1.2}, \citeat{searle-khuri}{\S5.3, Exercise~5.10(b)}, \citeat{harville-solutions}{Ch.~10, Exercise~1(b)}]\label{r1-prop-idempotent}
Let \(A,B\in\R^{n\times n}\).
\begin{enumerate}
\item \(A\) is idempotent if and only if \(A'\) is idempotent.
\item \(A\) is idempotent if and only if \(I_n-A\) is idempotent; and then \(A(I_n-A)=(I_n-A)A=0\).
\item If \(A\) is idempotent, then \(A^k=A\) for every integer \(k\geqslant1\).
\item If \(A\) and \(B\) are idempotent and \(AB=BA\), then \(AB\) is idempotent.
\end{enumerate}
\end{proposition}
\begin{proof}
(a)
\begin{itemize}
\item If \(A^2=A\), then \((A')^2=(AA)'=A'\) (Theorem~\ref{r1-thm-transpose}).
\item Conversely, if \(A'\) is idempotent, so is \((A')'=A\), by the first step.
\end{itemize}
(b)
\begin{itemize}
\item If \(A^2=A\), then \((I-A)^2=I-2A+A^2=I-A\) (Theorem~\ref{r1-thm-rules}).
\item Conversely, if \(I-A\) is idempotent, so is \(I-(I-A)=A\), by the first step.
\item Then \(A(I-A)=A-A^2=0\) and \((I-A)A=A-A^2=0\).
\end{itemize}
(c) Induction on \(k\): \(A^1=A\), and if \(A^k=A\), then \(A^{k+1}=A^kA=A^2=A\).

(d) \((AB)(AB)=A(BA)B=A(AB)B=A^2B^2=AB\) (Theorem~\ref{r1-thm-rules}).
\end{proof}

\subsection*{Traces of products; a cancellation rule}

\begin{proposition}[\citeat{harville}{(5.1.1), (5.2.4), (5.2.6)}, \citeat{abadir-magnus}{Exercise~2.26(b)}]\label{r1-prop-trace-more}
\begin{enumerate}
\item \(\tr(I_n)=n\), \(\tr(J_n)=n\), \(\tr(\bar J_n)=1\) and \(\tr(C_n)=n-1\).
\item For \(A,B\in\R^{m\times n}\) with columns \(a_1,\dots,a_n\) and \(b_1,\dots,b_n\),
\[
\tr(A'B)=\tr(AB')=\sum_{j=1}^na_j'b_j=\sum_{i=1}^m\sum_{j=1}^na_{ij}b_{ij}
\]
In particular \(\tr(ab')=a'b\) for \(a,b\in\R^m\).
\item \(x'Ax=\tr(Axx')\) for \(A\in\R^{n\times n}\) and \(x\in\R^n\).
\end{enumerate}
\end{proposition}
\begin{proof}
(a) The diagonal entries of \(I_n\) and of \(J_n\) are all \(1\). By Theorem~\ref{r1-thm-trace}(a),
\(\tr(\bar J_n)=\frac1n\tr(J_n)=1\) and \(\tr(C_n)=\tr(I_n)-\tr(\bar J_n)=n-1\).

(b)
\begin{itemize}
\item The \(j\)th diagonal entry of \(A'B\) is \(\row_j(A')\,b_j=a_j'b_j=\sum_ia_{ij}b_{ij}\)
(Theorems~\ref{r1-thm-views}(d) and~\ref{r1-thm-transpose}); summing over \(j\) gives the two sums.
\item \(\tr(AB')=\tr(B'A)\) (Theorem~\ref{r1-thm-trace}(b)), and the diagonal entries of \(B'A\) are
\(b_j'a_j=a_j'b_j\), so \(\tr(AB')\) is the same sum.
\item For \(n=1\), \(A=a\) and \(B=b\), this reads \(a'b=\tr(ab')\).
\end{itemize}
(c) By Theorem~\ref{r1-thm-trace}(b), applied to the \(1\times n\) matrix \(x'\) and the \(n\times1\) matrix \(Ax\),
\(\tr\bigl(x'(Ax)\bigr)=\tr\bigl((Ax)x'\bigr)\); the trace of the \(1\times1\) matrix \(x'Ax\) is this number.
\end{proof}

\needspace{14\baselineskip}
\begin{corollary}[\citeat{harville}{Corollaries~5.3.2, 5.3.3}]\label{r1-cor-cancel}
Let \(A\in\R^{m\times n}\).
\begin{enumerate}
\item \(A'A=0\) if and only if \(A=0\).
\item For \(B_1,B_2\in\R^{n\times p}\), \(A'AB_1=A'AB_2\) if and only if \(AB_1=AB_2\).
\item For \(B_1,B_2\in\R^{p\times n}\), \(B_1A'A=B_2A'A\) if and only if \(B_1A'=B_2A'\).
\end{enumerate}
\end{corollary}
\begin{proof}
(a) If \(A=0\), then \(A'A=0\). If \(A'A=0\), then \(\tr(A'A)=0\), so \(A=0\) by Theorem~\ref{r1-thm-trace}(d).

(b)
\begin{itemize}
\item If \(AB_1=AB_2\), multiplying by \(A'\) on the left gives \(A'AB_1=A'AB_2\).
\item Conversely, let \(A'AB_1=A'AB_2\). By Theorems~\ref{r1-thm-rules} and~\ref{r1-thm-transpose},
\((AB_1-AB_2)'=(B_1'-B_2')A'\), so
\[
(AB_1-AB_2)'(AB_1-AB_2)=(B_1'-B_2')(A'AB_1-A'AB_2)=0
\]
\item By (a), applied to the matrix \(AB_1-AB_2\), \(AB_1-AB_2=0\).
\end{itemize}
(c) By Theorem~\ref{r1-thm-transpose}, \(B_1A'A=B_2A'A\) is equivalent to \(A'AB_1'=A'AB_2'\), and \(B_1A'=B_2A'\)
to \(AB_1'=AB_2'\); these two are equivalent by (b).
\end{proof}

\begin{example}[\citeat{searle-khuri}{Example~5.2}]\label{r1-exm-cancel-law}
Let
\[
A=\begin{pmatrix}1&2&-1\\3&0&1\end{pmatrix},\qquad A'A=\begin{pmatrix}10&2&2\\2&4&-2\\2&-2&2\end{pmatrix}
\]
(the \((j,l)\) entry of \(A'A\) is \(\col_j(A)'\col_l(A)\), Theorems~\ref{r1-thm-views}(d) and~\ref{r1-thm-transpose};
for instance the \((1,1)\) entry \(1\cdot1+3\cdot3=10\) and the \((2,3)\) entry \(2\cdot(-1)+0\cdot1=-2\)), and let
\(b_1=(1,0,0)'\) and \(b_2=(0,2,3)'\). By Theorem~\ref{r1-thm-views}(b),
\[
A'Ab_1=\col_1(A'A)=\begin{pmatrix}10\\2\\2\end{pmatrix},\qquad
A'Ab_2=2\col_2(A'A)+3\col_3(A'A)=\begin{pmatrix}4+6\\8-6\\-4+6\end{pmatrix}=\begin{pmatrix}10\\2\\2\end{pmatrix}
\]
So \(A'Ab_1=A'Ab_2\), although \(b_1\ne b_2\): the factor \(A'A\) cannot be cancelled. The factor \(A'\) can: by
Corollary~\ref{r1-cor-cancel}(b), with \(B_1=b_1\) and \(B_2=b_2\) (\(p=1\)), \(Ab_1=Ab_2\). Indeed, by
Theorem~\ref{r1-thm-views}(b),
\[
Ab_1=\col_1(A)=\begin{pmatrix}1\\3\end{pmatrix},\qquad
Ab_2=2\col_2(A)+3\col_3(A)=\begin{pmatrix}4-3\\0+3\end{pmatrix}=\begin{pmatrix}1\\3\end{pmatrix}
\]
\end{example}

\needspace{20\baselineskip}
\subsection*{Sums of squares and products; linear combinations of variables}

\begin{proposition}[\citeat{searle-khuri}{\S5.1.4, \S10.3}, \citeat{abadir-magnus}{Exercise~2.17(a)}]\label{r1-prop-ssp}
Let \(X_0\in\R^{n\times k}\) be a data matrix with the rows \(x_1',\dots,x_n'\) (\(x_i\in\R^k\) holds the
measurements on unit \(i\)), let \(y\in\R^n\), \(\one=\one_n\), \(C=C_n\), and let \(\bar x=\frac1nX_0'\one\) be the
mean vector of \(X_0\) (Example~\ref{r1-exm-data}).
\begin{enumerate}
\item \(X_0'X_0=\sum_{i=1}^nx_ix_i'\) and \(X_0'y=\sum_{i=1}^nx_iy_i\); the \((j,l)\) entry of \(X_0'X_0\) is
\(\sum_ix_{ij}x_{il}\).
\item \(CX_0=X_0-\one\bar x'\) has the rows \((x_1-\bar x)',\dots,(x_n-\bar x)'\), and
\[
X_0'CX_0=(CX_0)'(CX_0)=\sum_{i=1}^n(x_i-\bar x)(x_i-\bar x)'=X_0'X_0-n\bar x\bar x'
\]
Its \((j,l)\) entry is \(\sum_i(x_{ij}-\bar x_j)(x_{il}-\bar x_l)\).
\end{enumerate}
\end{proposition}
\begin{proof}
(a)
\begin{itemize}
\item The columns of \(X_0'\) are \(x_1,\dots,x_n\) and the rows of \(X_0\) are \(x_1',\dots,x_n'\), so the outer
product expansion (Theorem~\ref{r1-thm-views}(d)) gives \(X_0'X_0=\sum_ix_ix_i'\).
\item By Theorem~\ref{r1-thm-views}(b), \(X_0'y=y_1x_1+\dots+y_nx_n\).
\item The \((j,l)\) entry of \(X_0'X_0\) is \(\col_j(X_0)'\col_l(X_0)=\sum_ix_{ij}x_{il}\)
(Theorems~\ref{r1-thm-views}(d) and~\ref{r1-thm-transpose}).
\end{itemize}
(b)
\begin{itemize}
\item \(\one'X_0=n\bar x'\), so \(CX_0=X_0-\frac1n\one(\one'X_0)=X_0-\one\bar x'\). The \(i\)th row of \(\one\bar x'\) is
\(\bar x'\) (Theorem~\ref{r1-thm-views}(c)), so the \(i\)th row of \(CX_0\) is \((x_i-\bar x)'\).
\item \(C'=C=C^2\) (Proposition~\ref{r1-prop-centering}), so \(X_0'CX_0=X_0'C'CX_0=(CX_0)'(CX_0)\).
\item (a), applied to the centered data matrix \(CX_0\), gives the sum of outer products and the entries.
\item As \(X_0'\one=n\bar x\), \(X_0'CX_0=X_0'X_0-\frac1n(X_0'\one)(\one'X_0)=X_0'X_0-n\bar x\bar x'\).
\end{itemize}
\end{proof}

\begin{proposition}[cf.\ \citeat{bapat}{\S10.1}, \citeat{tu-asm}{Definition~5.19}]\label{r1-prop-lincomb}
Let \(X_0\in\R^{n\times k}\) be a data matrix with mean vector \(\bar x\), let \(C=C_n\), and let
\(A\in\R^{k\times q}\). The \(j\)th column of \(Z=X_0A\in\R^{n\times q}\) is
\(X_0\col_j(A)=a_{1j}\col_1(X_0)+\dots+a_{kj}\col_k(X_0)\) (Theorem~\ref{r1-thm-views}(a),(b)): it holds the values,
on the \(n\) units, of the linear combination of the \(k\) variables with the coefficients \(a_{1j},\dots,a_{kj}\).
\begin{enumerate}
\item The mean vector of \(Z\) is \(A'\bar x\), and \(CZ=(CX_0)A\).
\item The matrix of corrected sums of squares and products of \(Z\) is \(Z'CZ=A'(X_0'CX_0)A\).
\item For \(a,b\in\R^k\), the combinations \(z=X_0a\) and \(w=X_0b\) have the means \(\bar z=a'\bar x\) and
\(\bar w=b'\bar x\), and
\[
\sum_{i=1}^n(z_i-\bar z)(w_i-\bar w)=a'(X_0'CX_0)b,\qquad \sum_{i=1}^n(z_i-\bar z)^2=a'(X_0'CX_0)a
\]
\end{enumerate}
\end{proposition}
\begin{proof}
(a) \(\frac1nZ'\one=\frac1nA'X_0'\one=A'\bar x\) (Theorem~\ref{r1-thm-transpose}), and \(CZ=C(X_0A)=(CX_0)A\)
(Theorem~\ref{r1-thm-rules}).

(b) \(Z'CZ=(X_0A)'C(X_0A)=A'(X_0'CX_0)A\) (Theorems~\ref{r1-thm-transpose} and~\ref{r1-thm-rules}); by
Proposition~\ref{r1-prop-ssp}(b), applied to \(Z\), it is the matrix of corrected sums of squares and products of the
columns of \(Z\).

(c)
\begin{itemize}
\item Take \(A=[a\ b]\in\R^{k\times2}\), so that \(Z=[z\ w]\). By (a), the means are \(a'\bar x\) and \(b'\bar x\).
\item By (b) and Proposition~\ref{r1-prop-ssp}(b), the \((1,2)\) and \((1,1)\) entries of \(A'(X_0'CX_0)A\) are the
two sums; by Theorem~\ref{r1-thm-views}(d) they are \(a'(X_0'CX_0)b\) and \(a'(X_0'CX_0)a\).
\end{itemize}
\end{proof}

\begin{example}\label{r1-exm-lincomb}
Four students take three tests; \(x_{ij}\) is the score of student \(i\) on test \(j\):
\[
X_0=\begin{pmatrix}5&4&6\\3&4&2\\5&3&6\\3&1&6\end{pmatrix},\qquad
\bar x=\tfrac14X_0'\one_4=\begin{pmatrix}4\\3\\5\end{pmatrix},\qquad
C_4X_0=X_0-\one_4\bar x'=\begin{pmatrix}1&1&1\\-1&1&-3\\1&0&1\\-1&-2&1\end{pmatrix}
\]
By Proposition~\ref{r1-prop-ssp}(b), \(M=X_0'C_4X_0\) is the sum of the outer products of the rows of \(C_4X_0\):
\[
M=\begin{pmatrix}1&1&1\\1&1&1\\1&1&1\end{pmatrix}+\begin{pmatrix}1&-1&3\\-1&1&-3\\3&-3&9\end{pmatrix}
+\begin{pmatrix}1&0&1\\0&0&0\\1&0&1\end{pmatrix}+\begin{pmatrix}1&2&-1\\2&4&-2\\-1&-2&1\end{pmatrix}
=\begin{pmatrix}4&2&4\\2&6&-4\\4&-4&12\end{pmatrix}
\]
Its trace, \(\tr(M)=4+6+12=22\), is the sum of the squares of all entries of \(C_4X_0\), since
\(M=(C_4X_0)'(C_4X_0)\) (Theorem~\ref{r1-thm-trace}(d)); row by row, \(3+11+2+6=22\).

The total score \(z=X_0\one_3=(15,9,14,10)'\) and the difference of the first two tests
\(w=X_0b=(1,-1,2,2)'\), \(b=(1,-1,0)'\), are linear combinations of the tests. By
Proposition~\ref{r1-prop-lincomb}(c), \(\bar z=\one_3'\bar x=12\), \(\bar w=b'\bar x=1\), and
\[
\sum_i(z_i-\bar z)^2=\one_3'M\one_3=26,\qquad \sum_i(w_i-\bar w)^2=b'Mb=6,\qquad
\sum_i(z_i-\bar z)(w_i-\bar w)=\one_3'Mb=6
\]
(\(\one_3'M=(10,4,12)\), the column sums of \(M\), and \(Mb=(2,-4,8)'\)). Directly,
\(z-\bar z\one_4=(3,-3,2,-2)'\) and \(w-\bar w\one_4=(0,-2,1,1)'\) give \(9+9+4+4=26\), \(0+4+1+1=6\) and
\(0+6+2-2=6\). The coefficients of \(w\) add up to \(0\) (\(\one_3'b=0\)): \(w\) is a \textbf{contrast} of the tests,
and it does not change if the same number \(t_i\) is added to the three scores of each student \(i\), since
\((X_0+t\one_3')b=X_0b+t(\one_3'b)=X_0b\) for every \(t\in\R^4\).
\end{example}

\subsection*{Exercises}

\begin{exercise}\label{r1-dr-three-views}
Let
\[
A=\begin{pmatrix}1&2&0&0\\0&1&-1&1\\2&0&1&0\\0&1&0&1\end{pmatrix},\qquad
B=\begin{pmatrix}1&0&2&0\\-1&1&0&1\\0&3&1&0\\1&0&0&2\end{pmatrix}
\]
\begin{enumerate}
\item Find the \((2,3)\) entry of \(AB\) from \(\row_2(A)\) and \(\col_3(B)\).
\item Find \(\col_2(AB)\) as a combination of the columns of \(A\).
\item Find \(\row_3(AB)\) as a combination of the rows of \(B\).
\item Compute \(AB\) and \(BA\).
\end{enumerate}
\end{exercise}
\begin{solution}
(a) By Theorem~\ref{r1-thm-views}(d), \((AB)_{23}=\row_2(A)\col_3(B)=0\cdot2+1\cdot0+(-1)\cdot1+1\cdot0=-1\).

(b) By Theorem~\ref{r1-thm-views}(a),(b), \(\col_2(AB)=A\col_2(B)=0\,a_1+1\,a_2+3\,a_3+0\,a_4\), with
\(a_j=\col_j(A)\):
\[
\col_2(AB)=\begin{pmatrix}2\\1\\0\\1\end{pmatrix}+3\begin{pmatrix}0\\-1\\1\\0\end{pmatrix}
=\begin{pmatrix}2\\-2\\3\\1\end{pmatrix}
\]
(c) By Theorem~\ref{r1-thm-views}(c),
\(\row_3(AB)=2\,\row_1(B)+0\,\row_2(B)+1\,\row_3(B)+0\,\row_4(B)=(2,0,4,0)+(0,3,1,0)=(2,3,5,0)\).

(d) Row by row, \(\row_1(AB)=1\,(1,0,2,0)+2\,(-1,1,0,1)=(-1,2,2,2)\),
\(\row_2(AB)=1\,(-1,1,0,1)-1\,(0,3,1,0)+1\,(1,0,0,2)=(0,-2,-1,3)\) and
\(\row_4(AB)=1\,(-1,1,0,1)+1\,(1,0,0,2)=(0,1,0,3)\); with (c),
\[
AB=\begin{pmatrix}-1&2&2&2\\0&-2&-1&3\\2&3&5&0\\0&1&0&3\end{pmatrix}
\]
in agreement with (a) and (b). In the same way \(\row_1(BA)=1\,(1,2,0,0)+2\,(2,0,1,0)=(5,2,2,0)\),
\(\row_2(BA)=-1\,(1,2,0,0)+1\,(0,1,-1,1)+1\,(0,1,0,1)=(-1,0,-1,2)\),
\(\row_3(BA)=3\,(0,1,-1,1)+1\,(2,0,1,0)=(2,3,-2,3)\) and \(\row_4(BA)=1\,(1,2,0,0)+2\,(0,1,0,1)=(1,4,0,2)\), so
\[
BA=\begin{pmatrix}5&2&2&0\\-1&0&-1&2\\2&3&-2&3\\1&4&0&2\end{pmatrix}\ne AB
\]
\end{solution}

\begin{exercise}[cf.\ \citeat{abadir-magnus}{Exercise~2.31}]\label{r1-ex-commuting}
Let
\[
N=\begin{pmatrix}0&1&0&0\\0&0&1&0\\0&0&0&1\\0&0&0&0\end{pmatrix}
\]
\begin{enumerate}
\item Compute \(N^2\), \(N^3\) and \(N^4\).
\item Show that the matrices \(B\in\R^{4\times4}\) with \(NB=BN\) are exactly
\[
B=\alpha I_4+\beta N+\gamma N^2+\delta N^3
=\begin{pmatrix}\alpha&\beta&\gamma&\delta\\0&\alpha&\beta&\gamma\\0&0&\alpha&\beta\\0&0&0&\alpha\end{pmatrix}
\qquad(\alpha,\beta,\gamma,\delta\in\R)
\]
\end{enumerate}
\end{exercise}
\begin{solution}
(a) By Theorem~\ref{r1-thm-views}(a),(b), \(Ne_1=0\) and \(Ne_j=e_{j-1}\) for \(j=2,3,4\).
Hence \(N^2e_j=e_{j-2}\) for \(j=3,4\) and \(N^2e_j=0\) for \(j=1,2\); \(N^3e_4=e_1\) and
\(N^3e_j=0\) for \(j\leqslant3\); and \(N^4e_j=0\) for all \(j\). Since \(\col_j(N^k)=N^ke_j\),
\[
N^2=\begin{pmatrix}0&0&1&0\\0&0&0&1\\0&0&0&0\\0&0&0&0\end{pmatrix},\qquad
N^3=\begin{pmatrix}0&0&0&1\\0&0&0&0\\0&0&0&0\\0&0&0&0\end{pmatrix},\qquad N^4=0
\]
This gives the matrix form of \(\alpha I_4+\beta N+\gamma N^2+\delta N^3\) in (b).

(b) Let \(B=(b_{ij})\). Since \(\row_i(N)=e_{i+1}'\) for \(i\leqslant3\) and \(\row_4(N)=0\),
Theorem~\ref{r1-thm-views}(c) gives \(\row_i(NB)=\row_{i+1}(B)\) for \(i\leqslant3\) and
\(\row_4(NB)=0\). By (a), \(\col_j(BN)=BNe_j=Be_{j-1}=\col_{j-1}(B)\) for \(j\geqslant2\) and
\(\col_1(BN)=0\). So
\[
(NB)_{ij}=\begin{cases}b_{i+1,j}&i\leqslant3\\0&i=4\end{cases},\qquad
(BN)_{ij}=\begin{cases}b_{i,j-1}&j\geqslant2\\0&j=1\end{cases}
\]
Suppose \(NB=BN\). Comparing the entries:
\begin{itemize}
\item row \(4\), \(j\geqslant2\): \(b_{4,j-1}=0\), so \(b_{41}=b_{42}=b_{43}=0\);
\item column \(1\), \(i\leqslant3\): \(b_{i+1,1}=0\), so \(b_{21}=b_{31}=b_{41}=0\);
\item \(i\leqslant3\), \(j\geqslant2\): \(b_{i+1,j}=b_{i,j-1}\), so every entry equals the entry
diagonally above and to its left.
\end{itemize}
The last condition gives \(b_{32}=b_{21}=0\); so all entries below the diagonal vanish. It
also gives \(b_{11}=b_{22}=b_{33}=b_{44}\), \(b_{12}=b_{23}=b_{34}\) and \(b_{13}=b_{24}\).
With \(\alpha=b_{11}\), \(\beta=b_{12}\), \(\gamma=b_{13}\), \(\delta=b_{14}\), \(B\) has the
stated form. Conversely, by Theorem~\ref{r1-thm-rules} and \(NN^k=N^{k+1}=N^kN\)
(Proposition~\ref{r1-prop-powers}(a)),
\[
N(\alpha I+\beta N+\gamma N^2+\delta N^3)=\alpha N+\beta N^2+\gamma N^3+\delta N^4
=(\alpha I+\beta N+\gamma N^2+\delta N^3)N
\]
for all \(\alpha,\beta,\gamma,\delta\).
\end{solution}

\begin{exercise}\label{r1-ex-products-drill}
Let
\[
A=\begin{pmatrix}1&1&0&0\\0&1&1&0\\0&0&1&1\end{pmatrix},\qquad
B=\begin{pmatrix}1&-1&1\\0&1&-1\\0&0&1\\0&0&0\end{pmatrix}
\]
\begin{enumerate}
\item Compute \(AB\) and \(BA\).
\item Write \(AB\) as the sum of the four products \(\col_j(A)\row_j(B)\) of
Theorem~\ref{r1-thm-views}(d).
\item Compute \(B'A'\), \(\tr(AB)\), \(\tr(BA)\) and \(\tr(A'A)\), and show that
\((BA)^2=BA\).
\end{enumerate}
\end{exercise}
\begin{solution}
(a) By Theorem~\ref{r1-thm-views}(c), \(\row_i(AB)=\row_i(A)B\) is a combination of the rows
\((1,-1,1)\), \((0,1,-1)\), \((0,0,1)\), \((0,0,0)\) of \(B\):
\[
\begin{aligned}
\row_1(AB)&=(1,-1,1)+(0,1,-1)=(1,\ 0,\ 0)\\
\row_2(AB)&=(0,1,-1)+(0,0,1)=(0,\ 1,\ 0)\\
\row_3(AB)&=(0,0,1)+(0,0,0)=(0,\ 0,\ 1)
\end{aligned}
\]
Likewise each row of \(BA\in\R^{4\times4}\) is a combination of the rows of \(A\):
\[
\begin{aligned}
\row_1(BA)&=(1,1,0,0)-(0,1,1,0)+(0,0,1,1)=(1,\ 0,\ 0,\ 1)\\
\row_2(BA)&=(0,1,1,0)-(0,0,1,1)=(0,\ 1,\ 0,\ -1)\\
\row_3(BA)&=(0,0,1,1),\qquad \row_4(BA)=(0,0,0,0)
\end{aligned}
\]
Hence
\[
AB=I_3,\qquad
BA=\begin{pmatrix}1&0&0&1\\0&1&0&-1\\0&0&1&1\\0&0&0&0\end{pmatrix}\ne I_4
\]

(b) \(\col_4(A)\row_4(B)=0\), and
\[
\begin{pmatrix}1\\0\\0\end{pmatrix}(1,\ -1,\ 1)+\begin{pmatrix}1\\1\\0\end{pmatrix}(0,\ 1,\ -1)
+\begin{pmatrix}0\\1\\1\end{pmatrix}(0,\ 0,\ 1)
=\begin{pmatrix}1&-1&1\\0&0&0\\0&0&0\end{pmatrix}+\begin{pmatrix}0&1&-1\\0&1&-1\\0&0&0\end{pmatrix}
+\begin{pmatrix}0&0&0\\0&0&1\\0&0&1\end{pmatrix}=I_3
\]

(c) \(B'A'=(AB)'=I_3'=I_3\), by the rules \((AB)'=B'A'\) and \(I_3'=I_3\) of Theorem~\ref{r1-thm-transpose}.
Next \(\tr(AB)=3\) and
\(\tr(BA)=1+1+1+0=3\), equal by Theorem~\ref{r1-thm-trace}(b); and by
Theorem~\ref{r1-thm-trace}(d), \(\tr(A'A)\) is the sum of the squares of the entries of \(A\),
which is \(6\). Finally, by Theorem~\ref{r1-thm-rules},
\((BA)^2=B(AB)A=BI_3A=BA\).
\end{solution}

\begin{exercise}[cf.\ \citeat{abadir-magnus}{Exercise~2.5(a)}]\label{r1-ex-products}
Let \(A=\begin{pmatrix}1&2&1&0\\0&1&2&1\\2&1&0&1\end{pmatrix}\), and let \(B\) be the matrix
obtained from \(A\) by interchanging its first two rows. Compute \(AB'\), \(BA'\), \(A'B\) and
\(B'A\), and compare their traces.
\end{exercise}
\begin{solution}
Write \(r_1'=(1,2,1,0)\), \(r_2'=(0,1,2,1)\), \(r_3'=(2,1,0,1)\) for the rows of \(A\) (\(r_1,r_2,r_3\in\R^4\)). The rows
of \(B\) are \(r_2',r_1',r_3'\), so
\(B=\left(\begin{smallmatrix}0&1&2&1\\1&2&1&0\\2&1&0&1\end{smallmatrix}\right)\). Since
\(\col_j(B')=\row_j(B)'\), the \((i,j)\) entry of \(AB'\) is \(\row_i(A)\row_j(B)'\)
(Theorem~\ref{r1-thm-views}(d)). The six numbers \(r_i'r_k\) needed are
\[
\begin{gathered}
r_1'r_1=1+4+1+0=6,\quad r_2'r_2=0+1+4+1=6,\quad r_3'r_3=4+1+0+1=6\\
r_1'r_2=0+2+2+0=4,\quad r_1'r_3=2+2+0+0=4,\quad r_2'r_3=0+1+0+1=2
\end{gathered}
\]
and \(r_k'r_i=r_i'r_k\) (Theorem~\ref{r1-thm-transpose}). Hence
\[
AB'=\begin{pmatrix}r_1'r_2&r_1'r_1&r_1'r_3\\r_2'r_2&r_2'r_1&r_2'r_3\\r_3'r_2&r_3'r_1&r_3'r_3\end{pmatrix}
=\begin{pmatrix}4&6&4\\6&4&2\\2&4&6\end{pmatrix},\qquad
BA'=(AB')'=\begin{pmatrix}4&6&2\\6&4&4\\4&2&6\end{pmatrix}
\]
by Theorem~\ref{r1-thm-transpose}. Since \(\col_k(A')=\row_k(A)'=r_k\), the outer product expansion (Theorem~\ref{r1-thm-views}(d)) gives
\[
A'B=\col_1(A')\row_1(B)+\col_2(A')\row_2(B)+\col_3(A')\row_3(B)=r_1r_2'+r_2r_1'+r_3r_3'
\]
This sum is symmetric, since \((r_1r_2')'=r_2r_1'\) and \((r_3r_3')'=r_3r_3'\)
(Theorem~\ref{r1-thm-transpose}). Its \((i,j)\) entry is \(\col_i(A)'\col_j(B)\); for instance
\((A'B)_{23}=2\cdot2+1\cdot1+1\cdot0=5\). Computing all sixteen entries in this way,
\[
A'B=\begin{pmatrix}4&3&2&3\\3&5&5&3\\2&5&4&1\\3&3&1&1\end{pmatrix}=(A'B)'=B'A
\]
So \(A'B\) is symmetric while \(AB'\) is not. The reason: \(B=SA\) with the symmetric matrix
\(S=\left(\begin{smallmatrix}0&1&0\\1&0&0\\0&0&1\end{smallmatrix}\right)\), since
\(\row_1(SA)=\row_2(A)\), \(\row_2(SA)=\row_1(A)\) and \(\row_3(SA)=\row_3(A)\)
(Theorem~\ref{r1-thm-views}(c)). Then \(A'B=A'SA\) and \((A'SA)'=A'S'A=A'SA\), whereas
\(AB'=AA'S\) need not be symmetric.

The four products have the same trace \(14\) (\(4+4+6\) and \(4+5+4+1\)); this is no
coincidence: \(\tr(AB')=\tr(B'A)\) and \(\tr(BA')=\tr(A'B)\) by
Theorem~\ref{r1-thm-trace}(b), and \(\tr(BA')=\tr\bigl((AB')'\bigr)=\tr(AB')\) by
Theorems~\ref{r1-thm-transpose} and~\ref{r1-thm-trace}(a).
\end{solution}

\begin{exercise}\label{r1-ex-outer}
Let \(u=(1,2,3,4)'\), \(v=(1,1,1,-1)'\), \(w=(1,1,-1,0)'\), and put \(A=uv'\) and
\(B=uw'\).
\begin{enumerate}
\item Write out \(A\) and \(B\), and compute \(v'u\) and \(w'u\).
\item Show that \(A^2=2A\), \(A^3=4A\), \(B^2=0\), \(AB=2B\) and \(BA=0\).
\item Compute \(\tr(A)\) and \(\tr(B)\).
\end{enumerate}
\end{exercise}
\begin{solution}
(a) Row \(i\) of \(uv'\) is \(u_iv'\) (Theorem~\ref{r1-thm-views}(c)), so
\[
A=\begin{pmatrix}1&1&1&-1\\2&2&2&-2\\3&3&3&-3\\4&4&4&-4\end{pmatrix},\qquad
B=\begin{pmatrix}1&1&-1&0\\2&2&-2&0\\3&3&-3&0\\4&4&-4&0\end{pmatrix}
\]
and \(v'u=1+2+3-4=2\), \(w'u=1+2-3+0=0\).

(b) By Theorem~\ref{r1-thm-rules}, a product of two such matrices collapses to a number
times a column times a row: for instance \(A^2=(uv')(uv')=u(v'u)v'=(v'u)\,uv'\), because
\(v'u\) is a number. In this way
\[
\begin{gathered}
A^2=u(v'u)v'=2A,\qquad A^3=A^2A=2A^2=4A,\qquad B^2=u(w'u)w'=0\\
AB=u(v'u)w'=2B,\qquad BA=u(w'u)v'=0
\end{gathered}
\]
So \(A\) and \(B\) do not commute, and \(B\ne0\) has \(B^2=0\).

(c) \(\tr(A)=1+2+3-4=2\) and \(\tr(B)=1+2-3+0=0\). In general
\(\tr(uv')=\tr(v'u)=v'u\) by Theorem~\ref{r1-thm-trace}(b).
\end{solution}

\begin{exercise}\label{r1-ex-triangular}
Let \(L\in\R^{4\times4}\) be the lower triangular matrix with all entries on and below the
diagonal equal to \(1\), and \(B=I_4-Z\), where \(Z\) has entries \(z_{21}=z_{32}=z_{43}=1\) and
all others \(0\):
\[
L=\begin{pmatrix}1&0&0&0\\1&1&0&0\\1&1&1&0\\1&1&1&1\end{pmatrix},\qquad
B=\begin{pmatrix}1&0&0&0\\-1&1&0&0\\0&-1&1&0\\0&0&-1&1\end{pmatrix}
\]
\begin{enumerate}
\item For \(x\in\R^4\), compute \(Lx\) and \(Bx\). Deduce that \(BL=I_4=LB\).
\item Show that \(LL'=(\min(i,j))\) and \(L'L=(5-\max(i,j))\), and write out both.
\item Compute \(B'B\) and show that \((LL')(B'B)=I_4\).
\end{enumerate}
\end{exercise}
\begin{solution}
(a) \(\row_i(L)x=x_1+\dots+x_i\), so \(Lx=(x_1,\ x_1+x_2,\ x_1+x_2+x_3,\ x_1+x_2+x_3+x_4)'\)
is the vector of partial sums, and \(Bx=x-Zx=(x_1,\ x_2-x_1,\ x_3-x_2,\ x_4-x_3)'\) is the
vector of differences. The differences of the partial sums of \(x\) are the entries of \(x\),
and the partial sums of the differences of \(x\) are again the entries of \(x\):
\(B(Lx)=x=L(Bx)\) for every \(x\). With \(x=e_j\), \(\col_j(BL)=BLe_j=e_j\) and
\(\col_j(LB)=e_j\) (Theorem~\ref{r1-thm-views}(a),(b)); so \(BL=I_4=LB\).

(b) \((LL')_{ij}=\row_i(L)\row_j(L)'\) counts the indices \(k\) with \(k\leqslant i\) and
\(k\leqslant j\), which is \(\min(i,j)\) (Theorems~\ref{r1-thm-views}(d) and~\ref{r1-thm-transpose}). \((L'L)_{ij}=\col_i(L)'\col_j(L)\) counts the
indices \(k\leqslant4\) with \(k\geqslant i\) and \(k\geqslant j\), which is \(5-\max(i,j)\).
So
\[
LL'=\begin{pmatrix}1&1&1&1\\1&2&2&2\\1&2&3&3\\1&2&3&4\end{pmatrix},\qquad
L'L=\begin{pmatrix}4&3&2&1\\3&3&2&1\\2&2&2&1\\1&1&1&1\end{pmatrix}
\]
both with trace \(10\).

(c) The columns of \(B\) are \(b_1=e_1-e_2\), \(b_2=e_2-e_3\), \(b_3=e_3-e_4\), \(b_4=e_4\),
and \((B'B)_{ij}=b_i'b_j\) (Theorems~\ref{r1-thm-views}(d) and~\ref{r1-thm-transpose}). Here
\(b_1'b_1=b_2'b_2=b_3'b_3=1+1=2\), \(b_4'b_4=1\), \(b_1'b_2=b_2'b_3=b_3'b_4=-1\), and
\(b_i'b_j=0\) for \(|i-j|\geqslant2\), since \(b_i\) and \(b_j\) then have no nonzero entry in
the same position:
\[
B'B=\begin{pmatrix}2&-1&0&0\\-1&2&-1&0\\0&-1&2&-1\\0&0&-1&1\end{pmatrix}
\]
By Theorems~\ref{r1-thm-rules} and~\ref{r1-thm-transpose} and (a),
\((LL')(B'B)=L(L'B')B=L(BL)'B=LI_4B=LB=I_4\).
\end{solution}

\begin{exercise}\label{r1-ex-data-drill}
Four units are measured on two variables (\(n=4\), \(k=2\)), with data matrix
\[
X_0=\begin{pmatrix}1&0\\3&2\\2&0\\2&2\end{pmatrix}\in\R^{4\times2}
\]
Find
\begin{enumerate}
\item the mean vector \(\bar x\) and the centered data matrix \(C_4X_0\);
\item the matrix \(X_0'C_4X_0\) of corrected sums of squares and products, both as
\((C_4X_0)'(C_4X_0)\) and as \(X_0'X_0-4\bar x\bar x'\).
\end{enumerate}
\end{exercise}
\begin{solution}
(a) \(\one'X_0=(1+3+2+2,\ 0+2+0+2)=(8,\ 4)\), so \(\bar x=\frac14X_0'\one=(2,\ 1)'\). As in
Example~\ref{r1-exm-data}, \(C_4X_0=X_0-\one\bar x'\):
\[
C_4X_0=\begin{pmatrix}1-2&0-1\\3-2&2-1\\2-2&0-1\\2-2&2-1\end{pmatrix}
=\begin{pmatrix}-1&-1\\1&1\\0&-1\\0&1\end{pmatrix}
\]

(b) The \((j,l)\) entry of \((C_4X_0)'(C_4X_0)\) is \(\col_j(C_4X_0)'\col_l(C_4X_0)\)
(Theorems~\ref{r1-thm-views}(d) and~\ref{r1-thm-transpose}): \(1+1+0+0=2\) for \((1,1)\),
\(1+1+0+0=2\) for \((1,2)\) and \((2,1)\), and \(1+1+1+1=4\) for \((2,2)\). On the other hand
\[
X_0'X_0-4\bar x\bar x'=\begin{pmatrix}18&10\\10&8\end{pmatrix}-4\begin{pmatrix}4&2\\2&1\end{pmatrix}
=\begin{pmatrix}18&10\\10&8\end{pmatrix}-\begin{pmatrix}16&8\\8&4\end{pmatrix}
=\begin{pmatrix}2&2\\2&4\end{pmatrix}
\]
(for instance \(18=1+9+4+4\) and \(10=0+6+0+4\)). So
\(X_0'C_4X_0=\left(\begin{smallmatrix}2&2\\2&4\end{smallmatrix}\right)\): the sums of squares
about the means are \(\sum_i(x_{i1}-\bar x_1)^2=2\) and \(\sum_i(x_{i2}-\bar x_2)^2=4\), and the
corrected sum of products is \(\sum_i(x_{i1}-\bar x_1)(x_{i2}-\bar x_2)=2\).
\end{solution}

\begin{exercise}\label{r1-ex-data-five}
Five units are measured on three variables (\(n=5\), \(k=3\)), with data matrix
\[
X_0=\begin{pmatrix}2&2&5\\3&2&2\\1&4&2\\3&6&8\\1&6&8\end{pmatrix}\in\R^{5\times3}
\]
\begin{enumerate}
\item Find the mean vector \(\bar x\) and the centered data matrix \(C_5X_0\).
\item Find \(S=X_0'C_5X_0\), both as \((C_5X_0)'(C_5X_0)\) and as
\(X_0'X_0-5\bar x\bar x'\).
\item Let \(D=\diag(\frac12,\frac14,\frac16)\). Compute \(DSD\), and explain why its diagonal
entries equal \(1\).
\end{enumerate}
\end{exercise}
\begin{solution}
(a) \(\one'X_0=(10,\ 20,\ 25)\), so \(\bar x=\frac15X_0'\one=(2,\ 4,\ 5)'\). As in
Example~\ref{r1-exm-data}, \(C_5X_0=X_0-\one\bar x'\):
\[
C_5X_0=\begin{pmatrix}0&-2&0\\1&-2&-3\\-1&0&-3\\1&2&3\\-1&2&3\end{pmatrix}
\]

(b) Write \(c_1,c_2,c_3\) for the columns of \(C_5X_0\). The \((j,l)\) entry of
\((C_5X_0)'(C_5X_0)\) is \(c_j'c_l\) (Theorems~\ref{r1-thm-views}(d) and~\ref{r1-thm-transpose}):
\[
\begin{gathered}
c_1'c_1=0+1+1+1+1=4,\quad c_2'c_2=4+4+0+4+4=16,\quad c_3'c_3=0+9+9+9+9=36\\
c_1'c_2=0-2+0+2-2=-2,\quad c_1'c_3=0-3+3+3-3=0,\quad c_2'c_3=0+6+0+6+6=18
\end{gathered}
\]
On the other hand
\[
X_0'X_0-5\bar x\bar x'
=\begin{pmatrix}24&38&50\\38&96&118\\50&118&161\end{pmatrix}
-\begin{pmatrix}20&40&50\\40&80&100\\50&100&125\end{pmatrix}
=\begin{pmatrix}4&-2&0\\-2&16&18\\0&18&36\end{pmatrix}=S
\]
(for instance \(38=4+6+4+18+6\) and \(161=25+4+4+64+64\)).

(c) By Proposition~\ref{r1-prop-diag}(a), left multiplication by \(D\) scales the rows and right multiplication
scales the columns, so \((DSD)_{ij}=d_is_{ij}d_j\):
\[
DSD=\begin{pmatrix}1&-\frac14&0\\-\frac14&1&\frac34\\0&\frac34&1\end{pmatrix}
\]
for instance \((DSD)_{23}=\frac14\cdot18\cdot\frac16=\frac34\). The diagonal entries are
\(d_i^2s_{ii}=\frac44,\frac{16}{16},\frac{36}{36}\), all equal to \(1\), because
\(d_i=1/\sqrt{s_{ii}}\) for \(i=1,2,3\).
\end{solution}

\begin{exercise}\label{r1-dr-units}
Four units are measured on three variables, with the data matrix
\(X_0=\begin{pmatrix}2&5&1\\4&0&3\\1&2&2\\3&1&4\end{pmatrix}\).
\begin{enumerate}
\item Compute \(e_3'X_0\), \(X_0e_2\) and \(e_3'X_0e_2\).
\item Compute \(X_0[e_3\ e_1]\), \([e_2\ e_4]'X_0\) and \([e_2\ e_4]'X_0[e_3\ e_1]\).
\item For the matrices \(E_{ij}\in\R^{4\times4}\), compute \(E_{12}E_{23}\) and \(E_{23}E_{12}\).
\end{enumerate}
\end{exercise}
\begin{solution}
(a) By Proposition~\ref{r1-prop-units}(a), \(e_3'X_0=\row_3(X_0)=(1,2,2)\), \(X_0e_2=\col_2(X_0)=(5,0,2,1)'\) and
\(e_3'X_0e_2=x_{32}=2\).

(b) By Theorem~\ref{r1-thm-views}(a) and (c), the first product has the columns \(3\) and \(1\) of \(X_0\), the second
the rows \(2\) and \(4\), and the third the entries of rows \(2,4\) in columns \(3,1\):
\[
X_0[e_3\ e_1]=\begin{pmatrix}1&2\\3&4\\2&1\\4&3\end{pmatrix},\qquad
[e_2\ e_4]'X_0=\begin{pmatrix}4&0&3\\3&1&4\end{pmatrix},\qquad
[e_2\ e_4]'X_0[e_3\ e_1]=\begin{pmatrix}3&4\\4&3\end{pmatrix}
\]
(c) With
\[
E_{12}=\begin{pmatrix}0&1&0&0\\0&0&0&0\\0&0&0&0\\0&0&0&0\end{pmatrix},\qquad
E_{23}=\begin{pmatrix}0&0&0&0\\0&0&1&0\\0&0&0&0\\0&0&0&0\end{pmatrix}
\]
Theorem~\ref{r1-thm-views}(c) gives \(\row_1(E_{12}E_{23})=1\,\row_2(E_{23})=(0,0,1,0)\), the other rows being zero,
and \(\row_2(E_{23}E_{12})=1\,\row_3(E_{12})=0\), the other rows being zero as well. So
\[
E_{12}E_{23}=\begin{pmatrix}0&0&1&0\\0&0&0&0\\0&0&0&0\\0&0&0&0\end{pmatrix}=E_{13},\qquad E_{23}E_{12}=0
\]
as Proposition~\ref{r1-prop-units}(c) says (\(\delta_{22}=1\), \(\delta_{31}=0\)).
\end{solution}

\begin{exercise}\label{r1-ex-trace-order}
For \(i,j\in\{1,2,3\}\) let \(E_{ij}=e_ie_j'\in\R^{3\times3}\), as in Proposition~\ref{r1-prop-units}. Compute
\(\tr(E_{12}E_{23}E_{31})\) and \(\tr(E_{12}E_{31}E_{23})\).
\end{exercise}
\begin{solution}
By Proposition~\ref{r1-prop-units}(c), \(E_{12}E_{23}E_{31}=E_{13}E_{31}=E_{11}\), with trace \(1\), whereas
\(E_{12}E_{31}=\delta_{23}E_{11}=0\), so \(E_{12}E_{31}E_{23}=0\), with trace \(0\). Thus the trace of a product
of three matrices is unchanged when the factors are moved cyclically (Theorem~\ref{r1-thm-trace}(c)), but it can
change when two factors are interchanged: \(\tr(E_{12}E_{23}E_{31})=1\ne0=\tr(E_{12}E_{31}E_{23})\).
\end{solution}

\begin{exercise}\label{r1-dr-powers}
Let
\[
U=\begin{pmatrix}1&2&-1&0\\0&-1&3&1\\0&0&2&1\\0&0&0&-1\end{pmatrix},\qquad
N=\begin{pmatrix}0&1&2&0\\0&0&1&3\\0&0&0&1\\0&0&0&0\end{pmatrix}
\]
\begin{enumerate}
\item Compute \(U^2\) and \(U^3\), and compare their diagonal entries with those of \(U\).
\item Compute \(N^2\), \(N^3\) and \(N^4\).
\item Write down \((U')^2\) without multiplying.
\end{enumerate}
\end{exercise}
\begin{solution}
(a) By rows (Theorem~\ref{r1-thm-views}(c)): \(\row_1(U^2)=1\,(1,2,-1,0)+2\,(0,-1,3,1)-1\,(0,0,2,1)=(1,0,3,1)\),
\(\row_2(U^2)=-1\,(0,-1,3,1)+3\,(0,0,2,1)+1\,(0,0,0,-1)=(0,1,3,1)\),
\(\row_3(U^2)=2\,(0,0,2,1)+1\,(0,0,0,-1)=(0,0,4,1)\) and \(\row_4(U^2)=-1\,(0,0,0,-1)=(0,0,0,1)\). Then, for
\(U^3=U^2U\), \(\row_1(U^3)=1\,(1,2,-1,0)+3\,(0,0,2,1)+1\,(0,0,0,-1)=(1,2,5,2)\),
\(\row_2(U^3)=1\,(0,-1,3,1)+3\,(0,0,2,1)+1\,(0,0,0,-1)=(0,-1,9,3)\),
\(\row_3(U^3)=4\,(0,0,2,1)+1\,(0,0,0,-1)=(0,0,8,3)\) and \(\row_4(U^3)=1\,(0,0,0,-1)=(0,0,0,-1)\):
\[
U^2=\begin{pmatrix}1&0&3&1\\0&1&3&1\\0&0&4&1\\0&0&0&1\end{pmatrix},\qquad
U^3=\begin{pmatrix}1&2&5&2\\0&-1&9&3\\0&0&8&3\\0&0&0&-1\end{pmatrix}
\]
Both are upper triangular; their diagonal entries \(1,1,4,1\) and \(1,-1,8,-1\) are the squares and the cubes of the
diagonal entries \(1,-1,2,-1\) of \(U\) (Proposition~\ref{r1-prop-powers}(e)).

(b) \(\row_1(N^2)=1\,\row_2(N)+2\,\row_3(N)=(0,0,1,5)\), \(\row_2(N^2)=1\,\row_3(N)+3\,\row_4(N)=(0,0,0,1)\), and
rows \(3\), \(4\) of \(N^2\) are zero. Then \(\row_1(N^3)=1\,\row_3(N)+5\,\row_4(N)=(0,0,0,1)\),
\(\row_2(N^3)=1\,\row_4(N)=0\), and \(\row_1(N^4)=1\,\row_4(N)=0\):
\[
N^2=\begin{pmatrix}0&0&1&5\\0&0&0&1\\0&0&0&0\\0&0&0&0\end{pmatrix},\qquad
N^3=\begin{pmatrix}0&0&0&1\\0&0&0&0\\0&0&0&0\\0&0&0&0\end{pmatrix},\qquad N^4=0
\]
With each power one more diagonal becomes zero, as in Proposition~\ref{r1-prop-strictly-triangular}(a), and
\(N^4=0\) (Proposition~\ref{r1-prop-strictly-triangular}(c), \(n=4\)).

(c) By Proposition~\ref{r1-prop-powers}(c),
\[
(U')^2=(U^2)'=\begin{pmatrix}1&0&0&0\\0&1&0&0\\3&3&4&0\\1&1&1&1\end{pmatrix}
\]
\end{solution}

\begin{exercise}\label{r1-ex-nilpotent}
Let
\[
N=\begin{pmatrix}0&1&1&1\\0&0&2&1\\0&0&0&3\\0&0&0&0\end{pmatrix}
\]
\begin{enumerate}
\item Compute \(N^2\), \(N^3\) and \(N^4\).
\item Let \(G=I_4-N+N^2-N^3\). Show that \((I_4+N)G=I_4=G(I_4+N)\), and write out \(G\).
\end{enumerate}
\end{exercise}
\begin{solution}
(a) By Theorem~\ref{r1-thm-views}(c), \(\row_i(N^{k+1})=\row_i(N^k)N\) is the combination of
the rows of \(N\) with coefficients \(\row_i(N^k)\). The rows of \(N\) are
\((0,1,1,1)\), \((0,0,2,1)\), \((0,0,0,3)\), \(0\). So
\[
\begin{aligned}
\row_1(N^2)&=1\,(0,0,2,1)+1\,(0,0,0,3)=(0,0,2,4), &\quad \row_2(N^2)&=2\,(0,0,0,3)=(0,0,0,6)\\
\row_1(N^3)&=2\,(0,0,0,3)=(0,0,0,6), &\quad \row_2(N^3)&=6\,\row_4(N)=0
\end{aligned}
\]
while \(\row_3(N^2)=3\,\row_4(N)=0\), \(\row_4(N^2)=\row_4(N)N=0\), and rows \(3,4\) of
\(N^3\) are \(\row_3(N^2)N=0\) and \(\row_4(N^2)N=0\). Hence
\[
N^2=\begin{pmatrix}0&0&2&4\\0&0&0&6\\0&0&0&0\\0&0&0&0\end{pmatrix},\qquad
N^3=6\,e_1e_4',\qquad N^4=6\,e_1e_4'N=6\,e_1\row_4(N)=0
\]

(b) By Theorem~\ref{r1-thm-rules} and \(NN^k=N^{k+1}=N^kN\) (Proposition~\ref{r1-prop-powers}(a)), the terms
cancel in pairs:
\[
(I+N)(I-N+N^2-N^3)=I-N+N^2-N^3+N-N^2+N^3-N^4=I-N^4=I
\]
and the product in the other order expands to the same sum. Adding the entries,
\[
G=\begin{pmatrix}1&-1&1&-3\\0&1&-2&5\\0&0&1&-3\\0&0&0&1\end{pmatrix}
\]
For instance \(g_{14}=-1+4-6=-3\) and \(g_{24}=-1+6=5\). As a check,
\(\row_1(I+N)\col_4(G)=1\cdot(-3)+1\cdot5+1\cdot(-3)+1\cdot1=0\).
\end{solution}

\begin{exercise}\label{r1-dr-polynomial}
Let
\[
A=\begin{pmatrix}2&-1&0&-1\\-1&2&0&1\\1&1&2&0\\1&-1&0&0\end{pmatrix}
\]
and \(f(t)=t^2-3t+2=(t-1)(t-2)\).
\begin{enumerate}
\item Compute \(A^2\) and \(f(A)\).
\item Compute \((A-I_4)(A-2I_4)\) and \((A-2I_4)(A-I_4)\).
\item Use (a) to write \(A^3\) as \(\alpha A+\beta I_4\), and check the result by computing \(AA^2\).
\end{enumerate}
\end{exercise}
\begin{solution}
(a) By rows (Theorem~\ref{r1-thm-views}(c)), \(\row_1(A^2)=2\,(2,-1,0,-1)-(-1,2,0,1)-(1,-1,0,0)=(4,-3,0,-3)\),
\(\row_2(A^2)=-(2,-1,0,-1)+2\,(-1,2,0,1)+(1,-1,0,0)=(-3,4,0,3)\),
\(\row_3(A^2)=(2,-1,0,-1)+(-1,2,0,1)+2\,(1,1,2,0)=(3,3,4,0)\) and
\(\row_4(A^2)=(2,-1,0,-1)-(-1,2,0,1)=(3,-3,0,-2)\). So
\[
A^2=\begin{pmatrix}4&-3&0&-3\\-3&4&0&3\\3&3&4&0\\3&-3&0&-2\end{pmatrix}=3A-2I_4
\]
since \(3A-2I_4\) has the same entries; that is, \(f(A)=A^2-3A+2I_4=0\).

(b) Here
\[
A-I_4=\begin{pmatrix}1&-1&0&-1\\-1&1&0&1\\1&1&1&0\\1&-1&0&-1\end{pmatrix},\qquad
A-2I_4=\begin{pmatrix}0&-1&0&-1\\-1&0&0&1\\1&1&0&0\\1&-1&0&-2\end{pmatrix}
\]
Write \(r_1',\dots,r_4'\) for the rows of \(A-2I_4\) and \(s_1',\dots,s_4'\) for those of \(A-I_4\)
(\(r_i,s_i\in\R^4\)). By Theorem~\ref{r1-thm-views}(c), the rows of \((A-I_4)(A-2I_4)\) are \(r_1'-r_2'-r_4'=0\),
\(-r_1'+r_2'+r_4'=0\), \(r_1'+r_2'+r_3'=0\) and \(r_1'-r_2'-r_4'=0\); the rows of \((A-2I_4)(A-I_4)\) are
\(-s_2'-s_4'=(0,0,0,0)\), \(-s_1'+s_4'=(0,0,0,0)\), \(s_1'+s_2'=(0,0,0,0)\) and \(s_1'-s_2'-2s_4'=(0,0,0,0)\). Both
products are \(0\), as Proposition~\ref{r1-prop-poly}(b) says: each equals \(f(A)\).

(c) By (a), \(A^2=3A-2I_4\), so \(A^3=AA^2=3A^2-2A=3(3A-2I_4)-2A=7A-6I_4\) (Proposition~\ref{r1-prop-powers}(a),
Theorem~\ref{r1-thm-rules}):
\[
A^3=7A-6I_4=\begin{pmatrix}8&-7&0&-7\\-7&8&0&7\\7&7&8&0\\7&-7&0&-6\end{pmatrix}
\]
Check: \(\row_1(AA^2)=2\,(4,-3,0,-3)-(-3,4,0,3)-(3,-3,0,-2)=(8,-7,0,-7)\),
\(\row_2(AA^2)=-(4,-3,0,-3)+2\,(-3,4,0,3)+(3,-3,0,-2)=(-7,8,0,7)\),
\(\row_3(AA^2)=(4,-3,0,-3)+(-3,4,0,3)+2\,(3,3,4,0)=(7,7,8,0)\) and
\(\row_4(AA^2)=(4,-3,0,-3)-(-3,4,0,3)=(7,-7,0,-6)\).
\end{solution}

\begin{exercise}\label{r1-ex-cyclic}
Let
\[
S=\begin{pmatrix}0&0&0&1\\1&0&0&0\\0&1&0&0\\0&0&1&0\end{pmatrix}
\]
so that \(Se_1=e_2\), \(Se_2=e_3\), \(Se_3=e_4\), \(Se_4=e_1\).
\begin{enumerate}
\item Compute \(S^2\), \(S^3\) and \(S^4\), and show that \(S'=S^3\) and \(S'S=I_4\).
\item Let \(T=I_4+2S+3S^2+4S^3\). Write out \(T\) and \(T'\), show that \(TT'=T'T\), and
compute \(TT'\).
\end{enumerate}
\end{exercise}
\begin{solution}
(a) Applying \(S\) repeatedly moves each unit vector one step along the cycle
\(e_1\to e_2\to e_3\to e_4\to e_1\). So \(S^2e_j=e_{j+2}\) and \(S^3e_j=e_{j+3}\), indices taken
cyclically (\(e_5=e_1\), \(e_6=e_2\), \dots), and \(S^4e_j=e_j\). With
\(\col_j(S^k)=S^ke_j\) (Theorem~\ref{r1-thm-views}(b)),
\[
S^2=\begin{pmatrix}0&0&1&0\\0&0&0&1\\1&0&0&0\\0&1&0&0\end{pmatrix},\qquad
S^3=\begin{pmatrix}0&1&0&0\\0&0&1&0\\0&0&0&1\\1&0&0&0\end{pmatrix},\qquad S^4=I_4
\]
The entries \(1\) of \(S\) stand in positions \((2,1)\), \((3,2)\), \((4,3)\), \((1,4)\), so
by \((S')_{ij}=s_{ji}\) the entries \(1\) of \(S'\) stand in positions \((1,2)\), \((2,3)\),
\((3,4)\), \((4,1)\), all other entries being \(0\). These are exactly the positions of the
entries \(1\) of \(S^3\), so \(S'=S^3\), and then \(S'S=S^3S=S^4=I_4\).

(b) Adding the entries of \(I\), \(2S\), \(3S^2\), \(4S^3\),
\[
T=\begin{pmatrix}1&4&3&2\\2&1&4&3\\3&2&1&4\\4&3&2&1\end{pmatrix},\qquad
T'=\begin{pmatrix}1&2&3&4\\4&1&2&3\\3&4&1&2\\2&3&4&1\end{pmatrix}
\]
By Theorem~\ref{r1-thm-transpose}, Proposition~\ref{r1-prop-powers}(a)--(c) and part (a),
\((S^2)'=(SS)'=S'S'=(S^3)^2=S^6=S^4S^2=S^2\) and \((S^3)'=(S')^3=S^9=(S^4)^2S=S\), so \(T'=I+4S+3S^2+2S^3\).
Both \(T\) and \(T'\) are polynomials in \(S\) (Definition~\ref{r1-def-poly}), so \(TT'=T'T\) by
Proposition~\ref{r1-prop-poly}(b). Expanding by Theorem~\ref{r1-thm-rules}, and using \(S^aS^b=S^{a+b}\)
(Proposition~\ref{r1-prop-powers}(a)) and \(S^4=I\), the coefficient of \(S^k\) in
\(TT'=(I+2S+3S^2+4S^3)(I+4S+3S^2+2S^3)\) collects the pairs of exponents with sum \(k\) or
\(k+4\):
\[
\begin{aligned}
I&:\ 1\cdot1+2\cdot2+3\cdot3+4\cdot4=30, &\quad S&:\ 1\cdot4+2\cdot1+3\cdot2+4\cdot3=24\\
S^2&:\ 1\cdot3+2\cdot4+3\cdot1+4\cdot2=22, &\quad S^3&:\ 1\cdot2+2\cdot3+3\cdot4+4\cdot1=24
\end{aligned}
\]
Hence
\[
TT'=30I+24S+22S^2+24S^3=\begin{pmatrix}30&24&22&24\\24&30&24&22\\22&24&30&24\\24&22&24&30\end{pmatrix}
\]
Its \((1,1)\) entry is \(\row_1(T)\row_1(T)'=1+16+9+4=30\).
\end{solution}

\begin{exercise}\label{r1-dr-sym-skew}
Let
\[
A=\begin{pmatrix}2&3&-1&1\\1&0&4&0\\5&-2&1&-1\\3&2&1&2\end{pmatrix},\qquad
P=\begin{pmatrix}1&0\\-1&1\\0&2\\1&0\end{pmatrix}
\]
\begin{enumerate}
\item Compute \(S=\frac12(A+A')\) and \(K=\frac12(A-A')\), and check that \(S'=S\), \(K'=-K\) and \(S+K=A\).
\item Compute \(P'SP\) and check that it is symmetric.
\end{enumerate}
\end{exercise}
\begin{solution}
(a)
\[
A'=\begin{pmatrix}2&1&5&3\\3&0&-2&2\\-1&4&1&1\\1&0&-1&2\end{pmatrix},\qquad
S=\tfrac12\begin{pmatrix}4&4&4&4\\4&0&2&2\\4&2&2&0\\4&2&0&4\end{pmatrix}
=\begin{pmatrix}2&2&2&2\\2&0&1&1\\2&1&1&0\\2&1&0&2\end{pmatrix}
\]
\[
K=\tfrac12\begin{pmatrix}0&2&-6&-2\\-2&0&6&-2\\6&-6&0&-2\\2&2&2&0\end{pmatrix}
=\begin{pmatrix}0&1&-3&-1\\-1&0&3&-1\\3&-3&0&-1\\1&1&1&0\end{pmatrix}
\]
\(S\) equals its transpose, \(K'=-K\) with zero diagonal (Proposition~\ref{r1-prop-sym-skew}(c)), and
\(S+K=A\) entry by entry.

(b) By columns (Theorem~\ref{r1-thm-views}(a)), \(SP=[S(1,-1,0,1)'\ \ S(0,1,2,0)']\), and then \(P'(SP)\) by rows:
\[
SP=\begin{pmatrix}2&6\\3&2\\1&3\\3&1\end{pmatrix},\qquad
P'SP=\begin{pmatrix}1&-1&0&1\\0&1&2&0\end{pmatrix}\begin{pmatrix}2&6\\3&2\\1&3\\3&1\end{pmatrix}
=\begin{pmatrix}2&5\\5&8\end{pmatrix}
\]
which is symmetric, as Proposition~\ref{r1-prop-sym-skew}(d) says.
\end{solution}

\begin{exercise}\label{r1-dr-special}
Let
\[
M_1=\frac12\begin{pmatrix}1&0&1&0\\0&1&0&-1\\1&0&1&0\\0&-1&0&1\end{pmatrix},\qquad
M_2=\begin{pmatrix}0&1&0&1\\-1&-1&1&0\\-1&0&1&1\\-1&-1&1&0\end{pmatrix},\qquad
M_3=\begin{pmatrix}2&0&-1&0\\0&1&0&0\\3&0&-2&0\\0&1&0&-1\end{pmatrix}
\]
\begin{enumerate}
\item Compute \(M_1^2\), \(M_2^2\), \(M_2^3\) and \(M_3^2\), and decide which of the three matrices is idempotent,
which nilpotent and which involutory.
\item For the idempotent one, compute \(I_4-M_1\), \((I_4-M_1)^2\) and \(M_1(I_4-M_1)\).
\end{enumerate}
\end{exercise}
\begin{solution}
(a) By rows (Theorem~\ref{r1-thm-views}(c)):
\begin{itemize}
\item Let \(B=2M_1\). Then \(\row_1(B^2)=1\,(1,0,1,0)+1\,(1,0,1,0)=(2,0,2,0)\),
\(\row_2(B^2)=1\,(0,1,0,-1)-1\,(0,-1,0,1)=(0,2,0,-2)\), \(\row_3(B^2)=1\,(1,0,1,0)+1\,(1,0,1,0)=(2,0,2,0)\) and
\(\row_4(B^2)=-1\,(0,1,0,-1)+1\,(0,-1,0,1)=(0,-2,0,2)\); so \(B^2=2B\), and \(M_1^2=\frac14B^2=\frac12B=M_1\)
(Theorem~\ref{r1-thm-rules}).
\item \(\row_1(M_2^2)=\row_2(M_2)+\row_4(M_2)=(-2,-2,2,0)\),
\(\row_2(M_2^2)=-\row_1(M_2)-\row_2(M_2)+\row_3(M_2)=(0,0,0,0)\),
\(\row_3(M_2^2)=-\row_1(M_2)+\row_3(M_2)+\row_4(M_2)=(-2,-2,2,0)\) and
\(\row_4(M_2^2)=-\row_1(M_2)-\row_2(M_2)+\row_3(M_2)=(0,0,0,0)\). Then rows \(1\) and \(3\) of \(M_2^3=M_2^2M_2\)
are \(-2\,\row_1(M_2)-2\,\row_2(M_2)+2\,\row_3(M_2)=(0,0,0,0)\), and rows \(2\) and \(4\) are \(0\).
\item \(\row_1(M_3^2)=2\,(2,0,-1,0)-(3,0,-2,0)=(1,0,0,0)\), \(\row_2(M_3^2)=(0,1,0,0)\),
\(\row_3(M_3^2)=3\,(2,0,-1,0)-2\,(3,0,-2,0)=(0,0,1,0)\) and \(\row_4(M_3^2)=(0,1,0,0)-(0,1,0,-1)=(0,0,0,1)\).
\end{itemize}
\[
\begin{gathered}
M_1^2=\frac12\begin{pmatrix}1&0&1&0\\0&1&0&-1\\1&0&1&0\\0&-1&0&1\end{pmatrix},\qquad
M_2^2=\begin{pmatrix}-2&-2&2&0\\0&0&0&0\\-2&-2&2&0\\0&0&0&0\end{pmatrix},\qquad M_2^3=0\\
M_3^2=\begin{pmatrix}1&0&0&0\\0&1&0&0\\0&0&1&0\\0&0&0&1\end{pmatrix}
\end{gathered}
\]
So \(M_1\) is idempotent (and symmetric), \(M_2\) is nilpotent (\(M_2^3=0\), \(M_2^2\ne0\); it is not
triangular), and \(M_3\) is involutory. Each lacks the other two properties: \(M_1^2=M_1\ne I_4\), and
\(M_1^k=M_1\ne0\) for every \(k\geqslant1\) (Proposition~\ref{r1-prop-idempotent}(c)); \(M_2^2\) differs from \(M_2\)
and from \(I_4\) (its second row is \(0\), theirs are not); and \(M_3^2=I_4\ne M_3\), while \(M_3^k\) is \(I_4\) for
even \(k\) and \(M_3\) for odd \(k\) (Proposition~\ref{r1-prop-powers}(a),(b)), never \(0\).

(b)
\[
\begin{gathered}
I_4-M_1=\frac12\begin{pmatrix}1&0&-1&0\\0&1&0&1\\-1&0&1&0\\0&1&0&1\end{pmatrix},\qquad
(I_4-M_1)^2=\frac12\begin{pmatrix}1&0&-1&0\\0&1&0&1\\-1&0&1&0\\0&1&0&1\end{pmatrix}\\
M_1(I_4-M_1)=\begin{pmatrix}0&0&0&0\\0&0&0&0\\0&0&0&0\\0&0&0&0\end{pmatrix}
\end{gathered}
\]
By rows, with \(F=2(I_4-M_1)\): \(\row_1(F^2)=1\,(1,0,-1,0)-1\,(-1,0,1,0)=(2,0,-2,0)\),
\(\row_2(F^2)=1\,(0,1,0,1)+1\,(0,1,0,1)=(0,2,0,2)\), \(\row_3(F^2)=-\row_1(F^2)\) and \(\row_4(F^2)=\row_2(F^2)\)
(\(\row_3(F)=-\row_1(F)\), \(\row_4(F)=\row_2(F)\); Theorem~\ref{r1-thm-views}(c)), so \(F^2=2F\) and
\((I_4-M_1)^2=\frac14F^2=\frac12F=I_4-M_1\); and \(\row_1(BF)=1\,(1,0,-1,0)+1\,(-1,0,1,0)=0\),
\(\row_2(BF)=1\,(0,1,0,1)-1\,(0,1,0,1)=0\), \(\row_3(BF)=\row_1(BF)\) and \(\row_4(BF)=-\row_2(BF)\)
(\(\row_3(B)=\row_1(B)\), \(\row_4(B)=-\row_2(B)\)), so \(M_1(I_4-M_1)=\frac14BF=0\). So \(I_4-M_1\) is idempotent and
\(M_1(I_4-M_1)=0\), as Proposition~\ref{r1-prop-idempotent}(b) says.
\end{solution}

\begin{exercise}\label{r1-ex-hadamard}
Let
\[
H=[h_1\ h_2\ h_3\ h_4]=\begin{pmatrix}1&1&1&1\\1&-1&1&-1\\1&1&-1&-1\\1&-1&-1&1\end{pmatrix}
\]
\begin{enumerate}
\item Show that \(H'=H\) and \(H^2=4I_4\).
\item Let \(E=\diag(1,1,0,0)\). Compute \(HE\), and use the outer product expansion
(Theorem~\ref{r1-thm-views}(d)) to write \(HEH\) and \(H^2\) as sums of the matrices \(h_jh_j'\).
\item Let \(P=\frac14HEH\), with \(E\) as in (b). Write out \(P\), and show that \(P^2=P\), \(\tr(P)=2\) and
\(\frac14H(I_4-E)H=I_4-P\).
\end{enumerate}
\end{exercise}
\begin{solution}
(a) Row \(i\) of \(H\) equals \(h_i'\) for each \(i\), so \(H'=H\). Then
\((H^2)_{ij}=\row_i(H)\col_j(H)=h_i'h_j\) (Theorem~\ref{r1-thm-views}(d)). Each \(h_i'h_i=1+1+1+1=4\), and the six products
with \(i<j\) vanish:
\[
\begin{gathered}
h_1'h_2=1-1+1-1=0,\quad h_1'h_3=1+1-1-1=0,\quad h_1'h_4=1-1-1+1=0\\
h_2'h_3=1-1-1+1=0,\quad h_2'h_4=1+1-1-1=0,\quad h_3'h_4=1-1+1-1=0
\end{gathered}
\]
So \(H^2=4I_4\).

(b) By Proposition~\ref{r1-prop-diag}(a), right multiplication by \(E\) multiplies the columns of \(H\) by
\(1,1,0,0\):
\[
HE=[h_1\ h_2\ 0\ 0]=\begin{pmatrix}1&1&0&0\\1&-1&0&0\\1&1&0&0\\1&-1&0&0\end{pmatrix}
\]
By (a), \(\row_j(H)=h_j'\). So the outer product expansions (Theorem~\ref{r1-thm-views}(d)) of \((HE)H\) and of
\(HH\) are
\[
\begin{gathered}
HEH=\sum_{j=1}^4\col_j(HE)\row_j(H)=h_1h_1'+h_2h_2'\\
H^2=\sum_{j=1}^4\col_j(H)\row_j(H)=h_1h_1'+h_2h_2'+h_3h_3'+h_4h_4'
\end{gathered}
\]

(c) By (b), \(P=\frac14HEH=\frac14(h_1h_1'+h_2h_2')\). Here
\(h_1h_1'=J_4\). The \(i\)th entry of \(h_2\) is \((-1)^{i+1}\), so the \((i,j)\) entry of
\(h_2h_2'\) is \((-1)^{i+1}(-1)^{j+1}=(-1)^{i+j}\): \(1\) where \(i+j\) is even and \(-1\)
where \(i+j\) is odd. Adding, \(h_1h_1'+h_2h_2'\) has entry \(2\) where \(i+j\) is even and
\(0\) where \(i+j\) is odd; so
\[
P=\frac12\begin{pmatrix}1&0&1&0\\0&1&0&1\\1&0&1&0\\0&1&0&1\end{pmatrix}
\]
By Theorem~\ref{r1-thm-rules}, \(H^2=4I\) from (a), and \(E^2=E\),
\[
P^2=\tfrac1{16}(HEH)(HEH)=\tfrac1{16}HE(H^2)EH=\tfrac1{16}HE(4I)EH=\tfrac14HE^2H=\tfrac14HEH=P
\]
Next \(\tr(P)=\frac12(1+1+1+1)=2\). Finally, by Theorem~\ref{r1-thm-rules} and (a),
\(\frac14H(I-E)H=\frac14H^2-\frac14HEH=\frac14\,4I-P=I-P\).
\end{solution}

\begin{exercise}\label{r1-ex-gram}
With \(h_1,\dots,h_4\) as in Exercise~\ref{r1-ex-hadamard}, let
\[
A=[h_2\ h_3\ h_4]=\begin{pmatrix}1&1&1\\-1&1&-1\\1&-1&-1\\-1&-1&1\end{pmatrix}\in\R^{4\times3}
\]
\begin{enumerate}
\item Compute \(A'A\) and \(AA'\).
\item Show that \(AA'=4C_4\), \(\one_4'A=0\) and \(C_4A=A\).
\item Compute \(\tr(A'A)\) and \(\tr(AA')\).
\end{enumerate}
\end{exercise}
\begin{solution}
(a) \((A'A)_{jl}=h_{j+1}'h_{l+1}\), so \(A'A=4I_3\) by the products listed in
Exercise~\ref{r1-ex-hadamard}(a). By Theorem~\ref{r1-thm-views}(d),
\(AA'=h_2h_2'+h_3h_3'+h_4h_4'\), and by Exercise~\ref{r1-ex-hadamard}(a),(b),
\(h_1h_1'+h_2h_2'+h_3h_3'+h_4h_4'=H^2=4I_4\). Since \(h_1h_1'=J_4\),
\[
AA'=4I_4-J_4=\begin{pmatrix}3&-1&-1&-1\\-1&3&-1&-1\\-1&-1&3&-1\\-1&-1&-1&3\end{pmatrix}
\]

(b) \(4I_4-J_4=4(I_4-\frac14J_4)=4C_4\) (Definition~\ref{r1-def-ones}). Next
\(\one_4'A=(h_1'h_2,\ h_1'h_3,\ h_1'h_4)=0\), so
\(C_4A=A-\frac14\one_4(\one_4'A)=A\): each column of \(A\) already has mean \(0\).

(c) \(\tr(A'A)=12\) and \(\tr(AA')=4\cdot3=12\), equal by Theorem~\ref{r1-thm-trace}(b); both
are the sum of the squares of the twelve entries \(\pm1\) of \(A\)
(Theorem~\ref{r1-thm-trace}(d)).
\end{solution}

\begin{exercise}\label{r1-dr-trace}
Let \(A=\begin{pmatrix}1&-2\\0&3\\2&1\end{pmatrix}\), \(B=\begin{pmatrix}2&1\\1&0\\-1&4\end{pmatrix}\),
\[
M=\begin{pmatrix}1&2&0&1\\-1&3&1&0\\2&0&-2&1\\0&1&-1&2\end{pmatrix}
\]
and \(x=(1,-1,2,1)'\).
\begin{enumerate}
\item Compute \(\sum_i\sum_ja_{ij}b_{ij}\), \(A'B\) and \(AB'\), and compare \(\tr(A'B)\) and \(\tr(AB')\) with the sum.
\item Compute \(x'Mx\), \(xx'\) and \(Mxx'\), and compare \(\tr(Mxx')\) with \(x'Mx\).
\end{enumerate}
\end{exercise}
\begin{solution}
(a) \(\sum_i\sum_ja_{ij}b_{ij}=1\cdot2+(-2)\cdot1+0\cdot1+3\cdot0+2\cdot(-1)+1\cdot4=2\), and
\[
A'B=\begin{pmatrix}1&0&2\\-2&3&1\end{pmatrix}\begin{pmatrix}2&1\\1&0\\-1&4\end{pmatrix}=\begin{pmatrix}0&9\\-2&2\end{pmatrix},\qquad
AB'=\begin{pmatrix}1&-2\\0&3\\2&1\end{pmatrix}\begin{pmatrix}2&1&-1\\1&0&4\end{pmatrix}=\begin{pmatrix}0&1&-9\\3&0&12\\5&2&2\end{pmatrix}
\]
So \(\tr(A'B)=0+2=2\) and \(\tr(AB')=0+0+2=2\), the same sum (Proposition~\ref{r1-prop-trace-more}(b)).

(b) \(Mx=(1-2+0+1,\ -1-3+2+0,\ 2+0-4+1,\ 0-1-2+2)'=(0,-2,-1,-1)'\), so \(x'Mx=0+2-2-1=-1\). Further
\[
xx'=\begin{pmatrix}1&-1&2&1\\-1&1&-2&-1\\2&-2&4&2\\1&-1&2&1\end{pmatrix},\qquad
Mxx'=(Mx)x'=\begin{pmatrix}0&0&0&0\\-2&2&-4&-2\\-1&1&-2&-1\\-1&1&-2&-1\end{pmatrix}
\]
and \(\tr(Mxx')=0+2-2-1=-1=x'Mx\) (Proposition~\ref{r1-prop-trace-more}(c)).
\end{solution}

\begin{exercise}\label{r1-dr-lincomb}
Five units are measured on three variables, with the data matrix
\[
X_0=\begin{pmatrix}0&1&1\\3&1&0\\2&4&2\\3&3&1\\2&1&1\end{pmatrix}\in\R^{5\times3}
\]
\begin{enumerate}
\item Find the mean vector \(\bar x\) and the centered data matrix \(C_5X_0\).
\item Find \(M=X_0'C_5X_0\), both as \((C_5X_0)'(C_5X_0)\) and as \(X_0'X_0-5\bar x\bar x'\).
\item For \(a=(1,1,1)'\) and \(b=(1,0,-1)'\), let \(z=X_0a\) and \(w=X_0b\). Find \(\bar z\), \(\bar w\),
\(\sum_i(z_i-\bar z)^2\), \(\sum_i(w_i-\bar w)^2\) and \(\sum_i(z_i-\bar z)(w_i-\bar w)\), first by
Proposition~\ref{r1-prop-lincomb}(c) and then directly.
\end{enumerate}
\end{exercise}
\begin{solution}
(a) \(\one_5'X_0=(10,10,5)\), so \(\bar x=\frac15X_0'\one_5=(2,2,1)'\), and \(C_5X_0=X_0-\one_5\bar x'\)
(Proposition~\ref{r1-prop-ssp}(b)):
\[
C_5X_0=\begin{pmatrix}0-2&1-2&1-1\\3-2&1-2&0-1\\2-2&4-2&2-1\\3-2&3-2&1-1\\2-2&1-2&1-1\end{pmatrix}
=\begin{pmatrix}-2&-1&0\\1&-1&-1\\0&2&1\\1&1&0\\0&-1&0\end{pmatrix}
\]
(b) The \((j,l)\) entry of \((C_5X_0)'(C_5X_0)\) is the product of columns \(j\) and \(l\) of \(C_5X_0\)
(Proposition~\ref{r1-prop-ssp}(a)): \(4+1+1=6\), \(1+1+4+1+1=8\) and \(1+1=2\) on the diagonal, and
\(2-1+1=2\), \(-1\), \(1+2=3\) off it. On the other hand, \(X_0'X_0\) has the entries
\(\sum_ix_{ij}x_{il}\) (for instance \(0+9+4+9+4=26\)), and
\[
X_0'X_0-5\bar x\bar x'=\begin{pmatrix}26&22&9\\22&28&13\\9&13&7\end{pmatrix}-\begin{pmatrix}20&20&10\\20&20&10\\10&10&5\end{pmatrix}
=\begin{pmatrix}6&2&-1\\2&8&3\\-1&3&2\end{pmatrix}=M
\]
(c) By Proposition~\ref{r1-prop-lincomb}(c), \(\bar z=a'\bar x=5\) and \(\bar w=b'\bar x=1\). With
\(Mb=\col_1(M)-\col_3(M)=(7,-1,-3)'\),
\[
a'Ma=24,\qquad b'Mb=7+3=10,\qquad a'Mb=7-1-3=3
\]
(\(a'Ma\) is the sum of all entries of \(M\)). Directly, \(z=(2,4,8,7,4)'\) and \(w=(-1,3,0,2,1)'\) have the means
\(\frac{25}5=5\) and \(\frac55=1\), the deviations are \(z-5\,\one_5=(-3,-1,3,2,-1)'\) and
\(w-\one_5=(-2,2,-1,1,0)'\), and
\[
\sum_i(z_i-\bar z)^2=9+1+9+4+1=24,\qquad \sum_i(w_i-\bar w)^2=4+4+1+1+0=10
\]
and \(\sum_i(z_i-\bar z)(w_i-\bar w)=6-2-3+2+0=3\).
\end{solution}

\needspace{14\baselineskip}
\subsection*{Problems}

\begin{problem}[\citeat{abadir-magnus}{Exercise~2.8(a), (c), (d)}]\label{r1-ex-unlike-scalars}
Give examples of real \(2\times2\) matrices such that:
(a) \(A^2=-I\);
(b) \(CD=-DC\) (\(CD\ne0\));
(c) \(EF=0\) (no entries of \(E\) or \(F\) are zero).
\end{problem}
\begin{solution}
(a) No real number \(a\) has \(a^2=-1\), but
\(A=\left(\begin{smallmatrix}0&1\\-1&0\end{smallmatrix}\right)\) gives
\(A^2=\left(\begin{smallmatrix}0\cdot0+1\cdot(-1)&0\cdot1+1\cdot0\\(-1)\cdot0+0\cdot(-1)&(-1)\cdot1+0\cdot0\end{smallmatrix}\right)=-I\).

(b) For \(C=\left(\begin{smallmatrix}1&1\\1&-1\end{smallmatrix}\right)\) and
\(D=\left(\begin{smallmatrix}1&-1\\-1&-1\end{smallmatrix}\right)\),
\[
CD=\begin{pmatrix}1-1&-1-1\\1+1&-1+1\end{pmatrix}=\begin{pmatrix}0&-2\\2&0\end{pmatrix},\qquad
DC=\begin{pmatrix}1-1&1+1\\-1-1&-1+1\end{pmatrix}=\begin{pmatrix}0&2\\-2&0\end{pmatrix}=-CD\ne0
\]

(c) For \(E=\left(\begin{smallmatrix}1&1\\1&1\end{smallmatrix}\right)\) and
\(F=\left(\begin{smallmatrix}1&1\\-1&-1\end{smallmatrix}\right)\), every entry of \(EF\) is
\(1\cdot1+1\cdot(-1)=0\), so \(EF=0\).
\end{solution}

\begin{problem}[\citeat{harville-solutions}{Ch.~1, Exercise~5}]\label{r1-ex-difference-squares}
Let \(A,B\in\R^{n\times n}\). Show that \((A+B)(A-B)=A^2-B^2\) if and only if \(A\) and \(B\)
commute.
\end{problem}
\begin{solution}
By Theorem~\ref{r1-thm-rules}, \((A+B)(A-B)=A(A-B)+B(A-B)=A^2-AB+BA-B^2\). Thus
\((A+B)(A-B)=A^2-B^2\) if and only if \(-AB+BA=0\), that is, if and only if \(AB=BA\).
\end{solution}

\begin{problem}[\citeat{harville-solutions}{Ch.~1, Exercise~6}]\label{r1-ex-symmetric-product}
\begin{enumerate}
\item Show that the product \(AB\) of two symmetric matrices \(A,B\in\R^{n\times n}\) is
symmetric if and only if \(A\) and \(B\) commute.
\item Give an example of two symmetric matrices of the same size whose product is not
symmetric.
\end{enumerate}
\end{problem}
\begin{solution}
(a) Since \(A\) and \(B\) are symmetric, \((AB)'=B'A'=BA\)
(Theorem~\ref{r1-thm-transpose}). So \(AB=(AB)'\) holds if and only if \(AB=BA\).

(b) Take \(A=\left(\begin{smallmatrix}1&0\\0&2\end{smallmatrix}\right)\) and
\(B=\left(\begin{smallmatrix}0&1\\1&0\end{smallmatrix}\right)\). Then
\(AB=\left(\begin{smallmatrix}0&1\\2&0\end{smallmatrix}\right)\ne\left(\begin{smallmatrix}0&2\\1&0\end{smallmatrix}\right)=BA\),
so by (a) \(AB\) is not symmetric.
\end{solution}

\begin{problem}[\citeat{harville-solutions}{Ch.~5, Exercise~1}]\label{r1-ex-trace-transposes}
Show that for any \(A\in\R^{m\times n}\), \(B\in\R^{n\times p}\) and \(C\in\R^{p\times m}\),
\(\tr(ABC)=\tr(B'A'C')=\tr(A'C'B')\).
\end{problem}
\begin{solution}
By Theorem~\ref{r1-thm-trace}(c), then (a), then Theorem~\ref{r1-thm-transpose}, and finally
(c) again,
\[
\tr(ABC)=\tr(CAB)=\tr\bigl((CAB)'\bigr)=\tr(B'A'C')=\tr(A'C'B')
\]
In the last step the factors \(B'\in\R^{p\times n}\), \(A'\in\R^{n\times m}\),
\(C'\in\R^{m\times p}\) are moved cyclically.
\end{solution}

\begin{problem}[\citeat{lang}{Ch.~III, \S3, Exercise~11(e)}, \citeat{shakarchi}{}]\label{r1-pr-commutator}
Prove that there are no matrices \(A,B\in\R^{n\times n}\) such that \(AB-BA=I_n\).
\end{problem}
\begin{solution}
Suppose that \(A,B\in\R^{n\times n}\) satisfy \(AB-BA=I_n\). Then \(\tr(AB-BA)=\tr(I_n)=n\), since the \(n\)
diagonal entries of \(I_n\) are \(1\) (Proposition~\ref{r1-prop-trace-more}(a)). On the other hand, by
Theorem~\ref{r1-thm-trace}(a),(b),
\[
\tr(AB-BA)=\tr(AB)-\tr(BA)=0
\]
So \(0=n\geqslant1\) \(\contradiction\) Hence no such \(A\) and \(B\) exist.
\end{solution}

\begin{problem}[\citeat{abadir-magnus}{Exercise~2.27}]\label{r1-ex-sum-entries}
Let \(A\in\R^{n\times n}\) and let \(D=\diag(a_{11},\dots,a_{nn})\). Show that the sum of all
entries of \(A\) is
\[
\one_n'A\one_n=\one_n'D\one_n+\tr\bigl((\one_n\one_n'-I_n)A\bigr)
\]
\end{problem}
\begin{solution}
By Theorem~\ref{r1-thm-trace}(a),(b),
\[
\tr\bigl((\one\one'-I_n)A\bigr)=\tr(\one\one'A)-\tr(A)=\tr(\one'A\one)-\tr(A)
=\one'A\one-\one'D\one
\]
Here \(\tr(\one\one'A)=\tr(\one'A\one)\) moves the \(n\times1\) factor \(\one\) past
the \(1\times n\) factor \(\one'A\); the trace of the \(1\times1\) matrix
\(\one'A\one\) is that number; and \(\one'D\one=\sum_ia_{ii}=\tr(A)\), because
\(D\one=(a_{11},\dots,a_{nn})'\). Rearranging gives the identity. Also
\(\one'A\one=\sum_i\sum_ja_{ij}\): by Theorem~\ref{r1-thm-views}(b),
\(A\one=a_1+\dots+a_n\) has \(i\)th entry \(\sum_ja_{ij}\), and by
Theorem~\ref{r1-thm-transpose}, \(\one'(A\one)\) is the sum of these \(n\) entries.
\end{solution}

\begin{problem}[\citeat{abadir-magnus}{Exercise~2.37(a), (c)}]\label{r1-ex-idempotent}
\begin{enumerate}
\item Show that the only symmetric idempotent \(2\times2\) matrices are \(I_2\), \(0\), and
\(uu'\) with \(u\in\R^2\), \(u'u=1\).
\item Give an example of an \(n\times n\) idempotent matrix (\(n\geqslant2\)) that is not
symmetric.
\end{enumerate}
\end{problem}
\begin{solution}
(a) The symmetric matrix \(A=\left(\begin{smallmatrix}a&b\\b&d\end{smallmatrix}\right)\) is
idempotent if and only if
\(\left(\begin{smallmatrix}a^2+b^2&b(a+d)\\b(a+d)&b^2+d^2\end{smallmatrix}\right)=\left(\begin{smallmatrix}a&b\\b&d\end{smallmatrix}\right)\),
that is, \(a^2+b^2=a\), \(b(a+d)=b\), \(d^2+b^2=d\).
\begin{itemize}
\item If \(a+d\ne1\), the second equation gives \(b=0\); then \(a^2=a\) and \(d^2=d\), so
\(a,d\in\{0,1\}\), and \(a+d\ne1\) leaves \(a=d=0\) or \(a=d=1\): \(A=0\) or \(A=I_2\).
\item If \(a+d=1\), then \(d=1-a\) and the first equation gives \(b^2=a(1-a)\), so
\(0\leqslant a\leqslant1\) and \(b=\pm\sqrt{a(1-a)}\); the third equation then holds, since
\((1-a)^2+a(1-a)=1-a\). Hence
\[
A=\begin{pmatrix}a&\pm\sqrt{a(1-a)}\\\pm\sqrt{a(1-a)}&1-a\end{pmatrix}
=\begin{pmatrix}\sqrt a\\\pm\sqrt{1-a}\end{pmatrix}\begin{pmatrix}\sqrt a&\pm\sqrt{1-a}\end{pmatrix}=uu'
\]
with the same sign throughout, and \(u'u=a+(1-a)=1\).
\end{itemize}
Conversely, \(I_2\) and \(0\) are symmetric idempotent, and if \(u'u=1\), then \(uu'\)
is symmetric and \((uu')(uu')=u(u'u)u'=uu'\).

(b) Let \(a,b\in\R^n\) with \(b'a=1\), and \(A=ab'\). Then
\(A^2=a(b'a)b'=ab'=A\). This \(A\) is symmetric only when \(b\) is a multiple of
\(a\): if \(ab'=ba'\), multiplying on the right by \(a\) gives \(a(b'a)=b(a'a)\),
that is, \(a=(a'a)b\); here \(a\ne0\) because \(b'a=1\), so \(a'a>0\) and
\(b=a/(a'a)\). Conversely, if \(b=ca\), then \(ab'=c\,aa'\) is symmetric. For a
concrete example take \(a=e_1\) and \(b=e_1+e_2\): \(b'a=1\), and \(A=e_1e_1'+e_1e_2'\)
has \((1,1)\) and \((1,2)\) entries \(1\) and all other entries \(0\); for \(n=2\),
\(A=\left(\begin{smallmatrix}1&1\\0&0\end{smallmatrix}\right)\), with \(A^2=A\ne A'\).
\end{solution}

\begin{problem}[\citeat{searle-khuri}{Exercise~4.15}]\label{r1-ex-idempotent-difference}
Let \(A,B\in\R^{n\times n}\) satisfy \(A^2=A\) and \(B^2=B\). Show that if \((A-B)^2=A-B\), then \(AB=BA=B\).
\end{problem}
\begin{solution}
Expanding by Theorem~\ref{r1-thm-rules}, \((A-B)^2=A-B\) reads \(A^2-AB-BA+B^2=A-B\); with \(A^2=A\) and \(B^2=B\)
this is \(A-AB-BA+B=A-B\), that is,
\[
2B=AB+BA
\]
\begin{itemize}
\item Multiplying on the left by \(A\): \(2AB=A^2B+ABA=AB+ABA\), hence \(AB=ABA\).
\item Multiplying on the right by \(A\): \(2BA=ABA+BA^2=ABA+BA\), hence \(ABA=BA\).
\item Therefore \(AB=BA\), and \(2B=AB+BA\) gives \(2B=2AB=2BA\), that is, \(AB=BA=B\).
\end{itemize}
\end{solution}

\begin{problem}[\citeat{harville-solutions}{Ch.~10, Exercise~6}]\label{r1-ex-involutory-idempotent}
Let \(A\in\R^{n\times n}\). Show that \(\frac12(I_n+A)\) is idempotent if and only if \(A\) is involutory.
\end{problem}
\begin{solution}
By Theorem~\ref{r1-thm-rules},
\[
\bigl[\tfrac12(I+A)\bigr]^2=\tfrac14I+\tfrac12A+\tfrac14A^2
\]
So \(\frac12(I+A)\) is idempotent if and only if \(\frac14I+\frac12A+\frac14A^2=\frac12I+\frac12A\), that is, if and
only if \(\frac14A^2=\frac14I\), that is, if and only if \(A^2=I\): \(A\) is involutory.
\end{solution}

\begin{problem}[\citeat{harville-solutions}{Ch.~5, Exercise~3}]\label{r1-ex-ata-a2}
Let \(A\in\R^{n\times n}\) satisfy \(A'A=A^2\).
\begin{enumerate}
\item Show that \(\tr\bigl[(A-A')'(A-A')\bigr]=0\).
\item Show that \(A\) is symmetric.
\end{enumerate}
\end{problem}
\begin{solution}
(a) Since \((A-A')'=A'-A\) and \(A'A'=(AA)'\) (Theorem~\ref{r1-thm-transpose}),
\[
\begin{aligned}
\tr\bigl[(A-A')'(A-A')\bigr]&=\tr\bigl[A'A-A'A'-AA+AA'\bigr]\\
&=\tr(A'A)-\tr\bigl[(AA)'\bigr]-\tr(A^2)+\tr(AA')\\
&=\tr(AA')-\tr\bigl[(AA)'\bigr]\\
&=\tr(A'A)-\tr\bigl[(AA)'\bigr]\\
&=\tr(A'A)-\tr(AA)=0
\end{aligned}
\]
The second line uses Theorem~\ref{r1-thm-trace}(a) (linearity); the third, \(A'A=A^2\); the fourth,
\(\tr(AA')=\tr(A'A)\) (Theorem~\ref{r1-thm-trace}(b)); the last, \(\tr(M')=\tr(M)\) (Theorem~\ref{r1-thm-trace}(a))
and \(A'A=A^2\) again.

(b) By (a) and Theorem~\ref{r1-thm-trace}(d), \(A-A'=0\), that is, \(A'=A\).
\end{solution}

\begin{problem}[\citeat{searle-khuri}{Exercise~5.8}]\label{r1-ex-trace-quadratic}
Let \(X=[x_1\ x_2\ \cdots\ x_c]\in\R^{r\times c}\) and let \(A\in\R^{r\times r}\) be symmetric. Prove that
\[
\tr\bigl[(AXX')^2\bigr]=\sum_{j=1}^c\sum_{k=1}^c(x_j'Ax_k)^2
\]
\end{problem}
\begin{solution}
Moving factors cyclically (Theorem~\ref{r1-thm-trace}(c), applied twice),
\[
\tr\bigl[AXX'AXX'\bigr]=\tr\bigl[XX'AXX'A\bigr]=\tr\bigl(X'AXX'AX\bigr)
\]
Let \(Y=X'AX\in\R^{c\times c}\); its \(k\)th column is \(y_k=X'Ax_k\) (Theorem~\ref{r1-thm-views}(a)).
\begin{itemize}
\item \(Y'=X'A'X=Y\), because \(A'=A\). So \(\row_k(Y)=y_k'\), and \(\tr(YY)=\sum_k\row_k(Y)\,y_k=\sum_ky_k'y_k\)
(Theorem~\ref{r1-thm-views}(d)).
\item \(y_k'=x_k'A'X=x_k'AX\), again because \(A'=A\). By the outer product expansion \(XX'=\sum_jx_jx_j'\)
(Theorem~\ref{r1-thm-views}(d)),
\[
y_k'y_k=x_k'AXX'Ax_k=\sum_{j=1}^c(x_k'Ax_j)(x_j'Ax_k)
\]
\item The number \(x_j'Ax_k\) equals its transpose \(x_k'A'x_j=x_k'Ax_j\), so each term is \((x_k'Ax_j)^2\).
\end{itemize}
Hence \(\tr\bigl[(AXX')^2\bigr]=\tr(YY)=\sum_k\sum_j(x_k'Ax_j)^2\), which is the stated double sum (the summand is
symmetric in \(j\) and \(k\)).
\end{solution}

\begin{problem}[\citeat{prasolov}{Problems~39.1, 39.2}]\label{r1-ex-pr-commute-diag}
Let \(d_1,\dots,d_n\) be distinct real numbers and \(D=\diag(d_1,\dots,d_n)\).
\begin{enumerate}
\item Show that every \(X\in\R^{n\times n}\) with \(DX=XD\) is diagonal.
\item Suppose, moreover, that \(d_1,\dots,d_n\) are nonzero, and let \(N\in\R^{n\times n}\) have the \((i,i+1)\) entries
\(1\) (\(1\leqslant i\leqslant n-1\)) and all other entries \(0\). Show that every \(X\in\R^{n\times n}\) that commutes
with \(D\) and with \(ND\) is \(X=cI_n\) for some \(c\in\R\).
\item Show that every \(X\in\R^{n\times n}\) that commutes with all matrices in \(\R^{n\times n}\) is \(X=cI_n\) for
some \(c\in\R\).
\end{enumerate}
\end{problem}
\begin{solution}
Write \(X=(x_{ij})\). By Proposition~\ref{r1-prop-diag}(a), left multiplication by a diagonal matrix scales the
rows, and right multiplication scales the columns:
\[
(DX)_{ij}=d_ix_{ij},\qquad (XD)_{ij}=x_{ij}d_j\qquad(1\leqslant i,j\leqslant n)
\]

(a) If \(DX=XD\), then \(d_ix_{ij}=x_{ij}d_j\), that is, \((d_i-d_j)x_{ij}=0\), for all \(i,j\). For \(i\ne j\),
\(d_i-d_j\ne0\), so \(x_{ij}=0\): \(X\) is diagonal.

(b)
\begin{itemize}
\item By (a), \(X=\diag(x_1,\dots,x_n)\), with \(x_i=x_{ii}\).
\item Let \(M=ND=(m_{ij})\). By the same rule, column \(j\) of \(M\) is \(d_j\col_j(N)\), so \(M\) has the
\((i,i+1)\) entries
\(m_{i,i+1}=d_{i+1}\) (\(1\leqslant i\leqslant n-1\)) and all other entries \(0\).
\item By the same rule for the diagonal matrix \(X\), \((MX)_{i,i+1}=m_{i,i+1}x_{i+1}=d_{i+1}x_{i+1}\) and
\((XM)_{i,i+1}=x_im_{i,i+1}=x_id_{i+1}\).
\item As \(MX=XM\), \(d_{i+1}(x_{i+1}-x_i)=0\), and \(d_{i+1}\ne0\); so \(x_{i+1}=x_i\) for \(i=1,\dots,n-1\). Hence
\(x_1=\dots=x_n=c\), and \(X=cI_n\).
\end{itemize}
(c) Let \(X\) commute with every matrix in \(\R^{n\times n}\), and take \(D=\diag(1,2,\dots,n)\), whose diagonal entries
are distinct and nonzero. Then \(X\) commutes with \(D\) and with \(ND\), and by (b), \(X=cI_n\) for some \(c\in\R\).
\end{solution}

\begin{problem}[\citeat{prasolov}{Problem~39.3}]\label{r1-ex-pr-commute-ones}
Find all matrices \(X\in\R^{n\times n}\) that commute with \(J_n=\one_n\one_n'\), the \(n\times n\) matrix with all
entries \(1\).
\end{problem}
\begin{solution}
Let \(\one=\one_n\) and \(J=J_n\), and put \(p=X\one\in\R^n\) and \(q'=\one'X\in\R^{1\times n}\). By
Theorem~\ref{r1-thm-views}(b),(c), the entries \(p_i=\sum_jx_{ij}\) of \(p\) are the row sums of \(X\), and the entries
\(q_j=\sum_ix_{ij}\) of \(q'\) are its column sums. By Theorem~\ref{r1-thm-rules},
\[
JX=\one(\one'X)=\one q',\qquad XJ=(X\one)\one'=p\one'
\]
so \((JX)_{ij}=q_j\) and \((XJ)_{ij}=p_i\) (Definition~\ref{r1-def-product}, with one term): every row of \(JX\) is
\((q_1,\dots,q_n)\), and every column of \(XJ\) is \(p\).
\begin{itemize}
\item If \(JX=XJ\), then \(q_j=p_i\) for all \(i,j\). Taking \(j=1\), every \(p_i\) equals \(q_1\); taking \(i=1\),
every \(q_j\) equals \(p_1=q_1\). So all row sums and all column sums of \(X\) equal the same number \(c=q_1\).
\item Conversely, if every row sum and every column sum of \(X\) equals \(c\), then all entries of \(JX\) and of \(XJ\)
are \(c\), so \(JX=XJ\).
\end{itemize}
Hence \(X\) commutes with \(J_n\) if and only if \(X\one_n=c\one_n\) and \(\one_n'X=c\one_n'\) for some \(c\in\R\), that
is, all row sums and all column sums of \(X\) are equal. For instance, \(aI_n+bJ_n\) has all row and column sums
\(a+nb\), in accordance with Example~\ref{r1-exm-poly}.
\end{solution}

\begin{problem}[\citeat{prasolov}{Problems~5.1, 40.7(a)}]\label{r1-ex-pr-trace-scalar}
Let \(A\in\R^{n\times n}\).
\begin{enumerate}
\item Suppose that \(\tr(AX)=0\) for every \(X\in\R^{n\times n}\) with \(\tr(X)=0\). Show that \(A=cI_n\) for some
\(c\in\R\).
\item Suppose that \(\tr(AXY)=\tr(AYX)\) for all \(X,Y\in\R^{n\times n}\). Show that \(A=cI_n\) for some \(c\in\R\).
\end{enumerate}
\end{problem}
\begin{solution}
Let \(E_{ij}=e_ie_j'\in\R^{n\times n}\) be the matrices of Proposition~\ref{r1-prop-units}. For all \(i,j\), by
Theorem~\ref{r1-thm-rules}, Theorem~\ref{r1-thm-trace}(b), applied to the \(n\times1\) matrix \(Ae_i\) and the
\(1\times n\) matrix \(e_j'\), and Proposition~\ref{r1-prop-units}(a),
\[
\tr(AE_{ij})=\tr\bigl((Ae_i)e_j'\bigr)=\tr\bigl(e_j'Ae_i\bigr)=e_j'Ae_i=a_{ji}
\]

(a)
\begin{itemize}
\item Let \(i\ne j\). The only nonzero entry of \(E_{ij}\) is its \((i,j)\) entry (Proposition~\ref{r1-prop-units}(b)),
which lies off the diagonal; so \(\tr(E_{ij})=0\), and by hypothesis \(a_{ji}=\tr(AE_{ij})=0\). Thus \(A\) is diagonal.
\item Let \(i\ne j\). Then \(\tr(E_{ii}-E_{jj})=1-1=0\) (Theorem~\ref{r1-thm-trace}(a)), and by hypothesis, with
Theorems~\ref{r1-thm-rules} and~\ref{r1-thm-trace}(a),
\(0=\tr\bigl(A(E_{ii}-E_{jj})\bigr)=\tr(AE_{ii})-\tr(AE_{jj})=a_{ii}-a_{jj}\).
\item So \(A\) is diagonal with all diagonal entries equal to \(a_{11}\), that is, \(A=a_{11}I_n\).
\end{itemize}
(b) By Theorems~\ref{r1-thm-rules} and~\ref{r1-thm-trace}(a), the hypothesis says that
\(\tr\bigl(A(XY-YX)\bigr)=\tr(AXY)-\tr(AYX)=0\) for all \(X,Y\in\R^{n\times n}\). Let \(i\ne j\).
\begin{itemize}
\item By Proposition~\ref{r1-prop-units}(c), \(E_{ij}E_{jj}=E_{ij}\) and \(E_{jj}E_{ij}=0\), as \(j\ne i\). With
\(X=E_{ij}\) and \(Y=E_{jj}\), \(a_{ji}=\tr(AE_{ij})=\tr(AE_{ij}E_{jj})=\tr(AE_{jj}E_{ij})=\tr(A0)=0\).
\item Likewise \(E_{ij}E_{ji}=E_{ii}\) and \(E_{ji}E_{ij}=E_{jj}\). With \(X=E_{ij}\) and \(Y=E_{ji}\),
\(a_{ii}=\tr(AE_{ii})=\tr(AE_{jj})=a_{jj}\).
\item As in (a), \(A=a_{11}I_n\). (So the traceless matrices \(E_{ij}=E_{ij}E_{jj}-E_{jj}E_{ij}\) and
\(E_{ii}-E_{jj}=E_{ij}E_{ji}-E_{ji}E_{ij}\) used in (a) are of the form \(XY-YX\).)
\end{itemize}
\end{solution}

\begin{problem}[\citeat{prasolov}{Problem~39.6(a)}]\label{r1-ex-pr-nilpotent-sum}
Let \(A,B\in\R^{n\times n}\) be nilpotent, with \(AB=BA\). Show that \(A+B\) is nilpotent.
\end{problem}
\begin{solution}
\begin{itemize}
\item \(A^jB^l=B^lA^j\) for all integers \(j,l\geqslant0\). First, \(AB^l=B^lA\), by induction on \(l\):
\(AB^0=A=B^0A\), and if \(AB^l=B^lA\), then \(AB^{l+1}=(AB^l)B=B^lAB=B^lBA=B^{l+1}A\) (Theorem~\ref{r1-thm-rules}).
Then, by induction on \(j\): \(A^0B^l=B^l=B^lA^0\), and if \(A^jB^l=B^lA^j\), then
\(A^{j+1}B^l=A^j(AB^l)=A^j(B^lA)=(A^jB^l)A=B^lA^jA=B^lA^{j+1}\).
\item The binomial formula: for every integer \(k\geqslant0\),
\[
(A+B)^k=\sum_{j=0}^k\binom kjA^jB^{k-j}
\]
by induction on \(k\). For \(k=0\) both sides are \(I_n\). If the formula holds for \(k\), then, by
Theorem~\ref{r1-thm-rules} and \(A^jB^{k-j}A=A^jAB^{k-j}=A^{j+1}B^{k-j}\) (first item),
\[
(A+B)^{k+1}=(A+B)^k(A+B)=\sum_{j=0}^k\binom kjA^{j+1}B^{k-j}+\sum_{j=0}^k\binom kjA^jB^{k+1-j}
=\sum_{j=0}^{k+1}\binom{k+1}jA^jB^{k+1-j}
\]
In the last step, the coefficient of \(A^jB^{k+1-j}\) is \(1\) for \(j=0\) and for \(j=k+1\), and for
\(1\leqslant j\leqslant k\) it is, by Pascal's rule,
\[
\binom k{j-1}+\binom kj=\frac{k!\,j+k!\,(k+1-j)}{j!\,(k+1-j)!}=\frac{(k+1)!}{j!\,(k+1-j)!}=\binom{k+1}j
\]
\item Let \(A^p=0\) and \(B^q=0\) with \(p,q\geqslant1\) (Definition~\ref{r1-def-nilpotent}), and let \(m=\max\{p,q\}\).
For \(j\geqslant m\), \(A^j=A^pA^{j-p}=0\) (Proposition~\ref{r1-prop-powers}(a), Theorem~\ref{r1-thm-rules}); likewise
\(B^l=0\) for \(l\geqslant m\).
\item In \((A+B)^{2m-1}=\sum_{j=0}^{2m-1}\binom{2m-1}jA^jB^{2m-1-j}\), either \(j\geqslant m\), or \(j\leqslant m-1\)
and then \(2m-1-j\geqslant m\). So every term has a zero factor and is \(0\) (Theorem~\ref{r1-thm-rules}). Hence
\((A+B)^{2m-1}=0\), with \(2m-1\geqslant1\): \(A+B\) is nilpotent.
\end{itemize}
\end{solution}

\begin{problem}[\citeat{prasolov}{Problem~34.4}]\label{r1-ex-pr-trace-ineq}
Let \(A,B\in\R^{n\times n}\) be symmetric. Show that \(\tr\bigl((AB)^2\bigr)\leqslant\tr(A^2B^2)\).
\end{problem}
\begin{solution}
Let \(K=AB-BA\).
\begin{itemize}
\item \(K\) is skew-symmetric: as \(A'=A\) and \(B'=B\), \(K'=B'A'-A'B'=BA-AB=-K\) (Theorem~\ref{r1-thm-transpose}).
\item Hence \(K'K=-K^2\), and by Theorem~\ref{r1-thm-trace}(a),(d),
\(\tr(K^2)=-\tr(K'K)=-\sum_{i,j}k_{ij}^2\leqslant0\).
\item By Theorem~\ref{r1-thm-rules}, \(K^2=ABAB-ABBA-BAAB+BABA\). By Theorem~\ref{r1-thm-trace}(b), applied to the
factors \(B\) and \(ABA\), to \(A\) and \(B^2A\), to \(B^2\) and \(A^2\), and to \(B\) and \(A^2B\),
\[
\tr(BABA)=\tr(ABAB),\qquad \tr(ABBA)=\tr(B^2A^2)=\tr(A^2B^2),\qquad \tr(BAAB)=\tr(A^2B^2)
\]
\item So, by Theorem~\ref{r1-thm-trace}(a),
\[
0\geqslant\tr(K^2)=2\tr\bigl((AB)^2\bigr)-2\tr(A^2B^2)
\]
that is, \(\tr\bigl((AB)^2\bigr)\leqslant\tr(A^2B^2)\).
\end{itemize}
\end{solution}

\begin{problem}[\citeat{hoffman-kunze}{\S8.5, Exercise~11}, \citeat{searle-khuri}{Exercise~5.7(a)}, \citeat{harville-solutions}{Ch.~10, Exercise~5}]\label{r1-pr-symmetric-nilpotent}
Let \(A\in\R^{n\times n}\) be symmetric and nilpotent: \(A^k=0\) for some integer \(k\geqslant1\)
(Definition~\ref{r1-def-nilpotent}). Show that \(A=0\). (Compare the matrix \(X\) of Remark~\ref{r1-rem-noncomm}:
\(X\ne0\) and \(X^2=0\), but \(X\) is not symmetric, since \(x_{12}=1\ne2=x_{21}\).)
\end{problem}
\begin{solution}
Since \(A^1=A^0A=I_nA=A\) (Theorem~\ref{r1-thm-rules}), the case \(k=1\) is \(A=A^1=0\). Let \(k\geqslant2\). It
suffices to show that, for every integer \(j\) with \(2\leqslant j\leqslant k\), \(A^j=0\) implies \(A^{j-1}=0\): then
\(A^k=0\) gives, in turn, \(A^{k-1}=0\), \dots, \(A^1=0\), that is, \(A=0\).

So let \(2\leqslant j\leqslant k\) and \(A^j=0\). Since \(A'=A\), Proposition~\ref{r1-prop-powers}(a) and
Theorem~\ref{r1-thm-rules} give
\[
A'A\,A^{j-2}=A^2A^{j-2}=A^j=0=A'A\,0
\]
where \(A^{j-2}=A^0=I_n\) if \(j=2\), and \(0\) is the \(n\times n\) zero matrix. By
Corollary~\ref{r1-cor-cancel}(b), with \(B_1=A^{j-2}\) and \(B_2=0\) (\(m=p=n\)), \(A\,A^{j-2}=A\,0\), that is,
\[
A^{j-1}=A\,A^{j-2}=A\,0=0
\]
again by Proposition~\ref{r1-prop-powers}(a) and Theorem~\ref{r1-thm-rules}.
\end{solution}

\clearpage
\def\unitnumber{2}
\setcounter{theorem}{0}\setcounter{problem}{0}
\section*{Chapter 2. Gauss elimination, RREF, and inverses}

\subsection*{Linear systems and augmented matrices}

\begin{definition}[\citeat{hoffman-kunze}{\S\S1.2--1.5, Example~12}, \citeat{abadir-magnus}{Exercise~6.41}]\label{r2-def-system}
Let \(A=(a_{ij})\in\R^{m\times n}\) and \(b\in\R^m\). The system of \(m\) linear equations
in \(n\) unknowns
\[
\begin{aligned}
a_{11}x_1+a_{12}x_2+\dots+a_{1n}x_n&=b_1\\
&\ \,\vdots\\
a_{m1}x_1+a_{m2}x_2+\dots+a_{mn}x_n&=b_m
\end{aligned}
\]
is the matrix equation \(Ax=b\), since \((Ax)_i=\sum_ja_{ij}x_j\). \(A\) is the
\textbf{coefficient matrix} and \([A\mid b]\in\R^{m\times(n+1)}\), the matrix \(A\) with
the column \(b\) appended, is the \textbf{augmented matrix}. A \textbf{solution} is an
\(x\in\R^n\) with \(Ax=b\). The system is \textbf{consistent} if it has a solution and
\textbf{inconsistent} otherwise. It is \textbf{homogeneous} if \(b=0\); then \(x=0\) is a
solution, the \textbf{trivial solution}.
\end{definition}

\subsection*{Elementary row operations}

\begin{definition}[\citeat{hoffman-kunze}{\S1.3}]\label{r2-def-rowops}
The \textbf{elementary row operations} on a matrix with \(m\) rows are, for row indices
\(r\ne s\):
\begin{enumerate}[label=(\arabic*)]
\item multiply row \(r\) by a number \(c\ne0\), written \(R_r\leftarrow cR_r\);
\item add \(c\) times row \(s\) to row \(r\) (\(c\in\R\)), written \(R_r\leftarrow R_r+cR_s\);
\item interchange rows \(r\) and \(s\), written \(R_r\leftrightarrow R_s\).
\end{enumerate}
When several operations label one arrow, they are performed one after another, in the order
written (upper label first). \(B\) is \textbf{row-equivalent} to \(A\) if \(B\) is obtained
from \(A\) by finitely many elementary row operations.
\end{definition}

\begin{theorem}[\citeat{hoffman-kunze}{\S\S1.2--1.4, Theorems~1--3}]\label{r2-thm-rowops}
\begin{enumerate}
\item Each elementary row operation is undone by an operation of the same type:
\(R_r\leftarrow cR_r\) by \(R_r\leftarrow c^{-1}R_r\), \(R_r\leftarrow R_r+cR_s\) by
\(R_r\leftarrow R_r-cR_s\), and \(R_r\leftrightarrow R_s\) by itself.
\item If row operations take \([A\mid b]\) to \([B\mid z]\), then the systems \(Ax=b\) and
\(Bx=z\) have exactly the same solutions.
\end{enumerate}
\end{theorem}
\begin{proof}
(a) is checked row by row. (b) It suffices to consider one operation. Each equation of the
new system is a multiple of an equation of the old one, or a sum of an old equation and a
multiple of another, so every solution of \(Ax=b\) solves \(Bx=z\). By (a), the old system
is obtained back from the new one by a row operation, so every solution of \(Bx=z\) solves
\(Ax=b\).
\end{proof}
So we solve a system by simplifying its augmented matrix with row operations; the next
definition says how far to simplify.

\subsection*{Echelon form and RREF}

\begin{definition}[\citeat{hoffman-kunze}{\S\S1.3--1.4}, \citeat{abadir-magnus}{Ch.~6}]\label{r2-def-rref}
A matrix \(R\) is in \textbf{reduced row echelon form} (\textbf{RREF}) if
\begin{enumerate}
\item the first nonzero entry of each nonzero row is \(1\), its \textbf{pivot};
\item each column that contains a pivot has all its other entries \(0\);
\item every zero row lies below every nonzero row;
\item if rows \(1,\dots,r\) are the nonzero rows and the pivot of row \(i\) lies in column
\(k_i\), then \(k_1<k_2<\dots<k_r\).
\end{enumerate}
The columns \(k_1,\dots,k_r\) are the \textbf{pivot columns}; column \(k_i\) of \(R\) is the
unit vector \(e_i\). A matrix that satisfies (a), (c), (d), but not necessarily (b), is in
\textbf{echelon form}; then every entry below a pivot is \(0\).
\end{definition}
For example, \(I_n\) and \(0\) are in RREF, and so is
\[
R=\begin{pmatrix}0&1&-3&0&2\\0&0&0&1&2\\0&0&0&0&0\end{pmatrix}\qquad(r=2,\ k_1=2,\ k_2=4)
\]
The matrix \(\left(\begin{smallmatrix}1&0&0&0\\0&1&-1&0\\0&0&1&0\end{smallmatrix}\right)\) is
in echelon form but not in RREF (column \(3\) violates (b)), and
\(\left(\begin{smallmatrix}0&2&1\\1&0&-3\\0&0&0\end{smallmatrix}\right)\) violates (a) and (d).

\begin{definition}[Gauss--Jordan elimination; \citeat{abadir-magnus}{Exercise~6.13, \S6.3}, \citeat{hoffman-kunze}{\S\S1.3--1.4, Theorems~4, 5}]\label{r2-def-gauss-jordan}
Let \(M\) be any matrix (for instance an augmented matrix). Start with \(i=1\).
\begin{enumerate}[label=\textit{Step \arabic*.},leftmargin=*]
\item Find the first column that has a nonzero entry in rows \(i,i+1,\dots\); if there is
none, go to Step 5. This is the next pivot column.
\item If the entry of row \(i\) in this column is \(0\), interchange row \(i\) with a lower row
that has a nonzero entry there.
\item Multiply row \(i\) by the reciprocal of this entry, so that the pivot is \(1\).
\item Add suitable multiples of row \(i\) to the rows below it, so that every entry below the
pivot is \(0\). Increase \(i\) by \(1\) and return to Step 1.
\item (\textit{Backward elimination.}) For each pivot, from the last one to the first, add
suitable multiples of its row to the rows above it, so that every entry above the pivot is
\(0\).
\end{enumerate}
Steps 1--4 (\textbf{forward elimination}) produce an echelon form; Step 5 produces a matrix in
RREF. Indeed, when Step 1 finds the pivot column of row \(i\), rows \(i,i+1,\dots\) are \(0\) in every column before
it, and Step 4 makes rows \(i+1,i+2,\dots\) \(0\) in this column too; so each nonzero row begins with its pivot \(1\),
the pivot columns are taken from left to right, and the rows left when Step 1 finds no column are zero. Thus forward
elimination ends in an echelon form, with pivots in columns \(k_1<\dots<k_r\), say. In Step 5, when the pivot of row
\(i\) is processed, row \(i\) is \(0\) before column \(k_i\) and, the pivots of rows \(i+1,\dots,r\) being done,
\(0\) in columns \(k_{i+1},\dots,k_r\); so the multiples of row \(i\) that are added change no entry before column
\(k_i\) and no other pivot column, and in the end every column \(k_i\) is \(e_i\). The scaling in Step 3 may also be
postponed, to the end of forward elimination or even to the very end (the leading entries become \(1\) only then);
this often avoids fractions. So for every matrix \(M\), some matrix in RREF is row-equivalent to \(M\).
\end{definition}

\begin{example}[cf.\ \citeat{hoffman-kunze}{\S1.3, Example~5}]\label{r2-exm-gauss-jordan}
Solve the homogeneous system
\[
x_1+x_2+x_3-2x_4=0,\qquad x_1+x_2+2x_3-3x_4=0,\qquad 2x_1+x_2+5x_3-4x_4=0
\]
\textit{Solution.} Row-reduce the coefficient matrix \(A\) (the column \(b=0\) never changes,
so it is not written).
\[
\begin{gathered}
\begin{pmatrix}1&1&1&-2\\1&1&2&-3\\2&1&5&-4\end{pmatrix}
\xrightarrow[R_3\leftarrow R_3-2R_1]{R_2\leftarrow R_2-R_1}
\begin{pmatrix}1&1&1&-2\\0&0&1&-1\\0&-1&3&0\end{pmatrix}
\xrightarrow{R_2\leftrightarrow R_3}
\begin{pmatrix}1&1&1&-2\\0&-1&3&0\\0&0&1&-1\end{pmatrix}
\\
\xrightarrow{R_2\leftarrow-R_2}
\begin{pmatrix}1&1&1&-2\\0&1&-3&0\\0&0&1&-1\end{pmatrix}
\xrightarrow[R_1\leftarrow R_1-R_3]{R_2\leftarrow R_2+3R_3}
\begin{pmatrix}1&1&0&-1\\0&1&0&-3\\0&0&1&-1\end{pmatrix}
\xrightarrow{R_1\leftarrow R_1-R_2}
\begin{pmatrix}1&0&0&2\\0&1&0&-3\\0&0&1&-1\end{pmatrix}
\end{gathered}
\]
After the first arrow the entry of row \(2\) in column \(2\) is \(0\), so Step 2 interchanges
rows \(2\) and \(3\), and Step 3 multiplies the new row \(2\) by \(-1\). The first matrix of the
second line is an echelon form (end of forward elimination); the last matrix, \(R\), is in RREF,
with pivot columns \(1,2,3\). The system \(Rx=0\) reads \(x_1+2x_4=0\), \(x_2-3x_4=0\),
\(x_3-x_4=0\), so \(x_4=t\) is free and all solutions are
\[
x=t\,(-2,\ 3,\ 1,\ 1)'\qquad(t\in\R)
\]
Check: \(-2+3+1-2=0\), \(-2+3+2-3=0\), \(-4+3+5-4=0\).
\end{example}

\subsection*{Reading off the solutions}

\begin{proposition}[\citeat{hoffman-kunze}{\S1.4}]\label{r2-prop-solutions}
Let row operations take \([A\mid b]\) (\(A\in\R^{m\times n}\)) to \([R\mid z]\) with \(R\) in
RREF, nonzero rows \(1,\dots,r\) and pivot columns \(k_1<\dots<k_r\). Call \(x_{k_1},\dots,x_{k_r}\)
the \textbf{pivot unknowns} and the other \(n-r\) unknowns the \textbf{free unknowns}.
\begin{enumerate}
\item \(Ax=b\) is consistent if and only if \(z_i=0\) for every \(i>r\), that is, no row of
\([R\mid z]\) reads \((0\ \cdots\ 0\mid z_i)\) with \(z_i\ne0\).
\item If it is consistent, all solutions are obtained by choosing arbitrary values for the free
unknowns and computing each pivot unknown from its row:
\(x_{k_i}=z_i-\sum_{j\text{ free}}R_{ij}x_j\) (\(1\leqslant i\leqslant r\)).
\end{enumerate}
\end{proposition}
\begin{proof}
By Theorem~\ref{r2-thm-rowops}, \(Ax=b\) and \(Rx=z\) have the same solutions. (a) Row
\(i>r\) of \(Rx=z\) reads \(0=z_i\). So there is no solution if one of these \(z_i\) is not \(0\);
if all of them are \(0\), these rows hold for every \(x\), and the argument for (b) gives a
solution. (b) For \(i\leqslant r\), column \(k_l\) of \(R\) is
\(e_l\), so the pivot unknown \(x_{k_i}\) occurs only in equation \(i\), with coefficient \(1\);
each choice of the free unknowns therefore gives exactly one solution.
\end{proof}

\begin{theorem}[\citeat{hefferon}{One.III.2, Theorem~2.6}]\label{r2-thm-rref-unique}
For every matrix \(A\), exactly one matrix in RREF is row-equivalent to \(A\); it is called \textbf{the RREF of
\(A\)}.
\end{theorem}
\begin{proof}
By Definition~\ref{r2-def-gauss-jordan}, for every \(A\in\R^{m\times n}\) some matrix in RREF is row-equivalent to
\(A\). Let \(R\) and \(S\) be matrices in RREF that are row-equivalent to \(A\); we show \(R=S\) by induction on
\(n\), for fixed \(m\).
\begin{itemize}
\item \(n=1\): an \(m\times1\) matrix in RREF is \(0\) or \(e_1\). Row operations take \(0\) to \(0\), and a nonzero
matrix to a nonzero one, since they can be undone (Theorem~\ref{r2-thm-rowops}(a)). So \(R=S=0\) if \(A=0\), and
\(R=S=e_1\) otherwise.
\item Let \(n>1\), and let \(\hat A,\hat R,\hat S\) consist of the first \(n-1\) columns of \(A,R,S\). Row
operations act on each column separately: by Definition~\ref{r2-def-rowops}, each of them computes the new entries
of a column from the entries of that column only. So the operations that take \(A\) to \(R\) take \(\hat A\) to
\(\hat R\). Moreover \(\hat R\) is in RREF: a row of \(R\)
whose pivot lies in column \(n\) is the last nonzero row, and it becomes a zero row of \(\hat R\). Likewise for
\(\hat S\); so \(\hat R=\hat S\) by the induction hypothesis.
\item Suppose \(R\ne S\); then they differ in column \(n\), say in row \(i\). The systems \(Ax=0\), \(Rx=0\) and
\(Sx=0\) have the same solutions (Theorem~\ref{r2-thm-rowops}(b)). Subtracting equation \(i\) of \(Sx=0\) from
equation \(i\) of \(Rx=0\) leaves \((r_{in}-s_{in})x_n=0\), so \(x_n=0\) for every solution.
\item Hence \(x_n\) is not a free unknown of \(Rx=0\), since a free unknown can take the value \(1\)
(Proposition~\ref{r2-prop-solutions}(b)): column \(n\) of \(R\) is a pivot column, and so is column \(n\) of \(S\).
\item Let \(r\) be the number of nonzero rows of \(\hat R=\hat S\). In \(R\), rows \(1,\dots,r\) have their pivots
before column \(n\), and rows \(r+1,\dots,m\) are zero before column \(n\); as zero rows come last, the pivot of
column \(n\) lies in row \(r+1\), and column \(n\) of \(R\) is \(e_{r+1}\) (Definition~\ref{r2-def-rref}). The same
holds for \(S\). So \(R=S\ \contradiction\)\qedhere
\end{itemize}
\end{proof}
In particular, the pivot columns, and hence the free unknowns of \(Ax=b\), do not depend on the row operations used to
reach the RREF.

\begin{example}[cf.\ \citeat{hoffman-kunze}{\S1.4, Example~8}]\label{r2-exm-read}
For the RREF \(R\) displayed after Definition~\ref{r2-def-rref}, the system \(Rx=0\) is
\(x_2-3x_3+2x_5=0\), \(x_4+2x_5=0\); the pivot unknowns are \(x_2,x_4\) and the free
unknowns \(x_1,x_3,x_5\). With \(x_1=a\), \(x_3=b\), \(x_5=c\), all solutions are
\[
x=\begin{pmatrix}a\\3b-2c\\b\\-2c\\c\end{pmatrix}
=a\begin{pmatrix}1\\0\\0\\0\\0\end{pmatrix}+b\begin{pmatrix}0\\3\\1\\0\\0\end{pmatrix}+c\begin{pmatrix}0\\-2\\0\\-2\\1\end{pmatrix}
\qquad(a,b,c\in\R)
\]
\end{example}

\begin{example}[cf.\ \citeat{hoffman-kunze}{\S1.4, Example~9}]\label{r2-exm-consistency}
Let
\[
A=\begin{pmatrix}1&2&0&1\\2&3&1&1\\1&1&1&1\\2&2&2&1\end{pmatrix}
\]
For which \(y=(y_1,y_2,y_3,y_4)'\) is \(Ax=y\) consistent? Solve it for \(y=(4,7,4,7)'\).

\textit{Solution.} Row-reduce the augmented matrix, keeping \(y_1,\dots,y_4\) as letters:
\[
\begin{gathered}
\left(\begin{array}{cccc|c}1&2&0&1&y_1\\2&3&1&1&y_2\\1&1&1&1&y_3\\2&2&2&1&y_4\end{array}\right)
\xrightarrow[R_4\leftarrow R_4-2R_1]{R_2\leftarrow R_2-2R_1,\ R_3\leftarrow R_3-R_1}
\left(\begin{array}{cccc|c}1&2&0&1&y_1\\0&-1&1&-1&y_2-2y_1\\0&-1&1&0&y_3-y_1\\0&-2&2&-1&y_4-2y_1\end{array}\right)
\\
\xrightarrow{R_2\leftarrow-R_2}
\left(\begin{array}{cccc|c}1&2&0&1&y_1\\0&1&-1&1&2y_1-y_2\\0&-1&1&0&y_3-y_1\\0&-2&2&-1&y_4-2y_1\end{array}\right)
\xrightarrow[R_4\leftarrow R_4+2R_2]{R_3\leftarrow R_3+R_2}
\left(\begin{array}{cccc|c}1&2&0&1&y_1\\0&1&-1&1&2y_1-y_2\\0&0&0&1&y_1-y_2+y_3\\0&0&0&1&2y_1-2y_2+y_4\end{array}\right)
\\
\xrightarrow{R_4\leftarrow R_4-R_3}
\left(\begin{array}{cccc|c}1&2&0&1&y_1\\0&1&-1&1&2y_1-y_2\\0&0&0&1&y_1-y_2+y_3\\0&0&0&0&y_1-y_2-y_3+y_4\end{array}\right)
\xrightarrow[R_1\leftarrow R_1-R_3]{R_2\leftarrow R_2-R_3}
\left(\begin{array}{cccc|c}1&2&0&0&y_2-y_3\\0&1&-1&0&y_1-y_3\\0&0&0&1&y_1-y_2+y_3\\0&0&0&0&y_1-y_2-y_3+y_4\end{array}\right)
\\
\xrightarrow{R_1\leftarrow R_1-2R_2}
\left(\begin{array}{cccc|c}1&0&2&0&-2y_1+y_2+y_3\\0&1&-1&0&y_1-y_3\\0&0&0&1&y_1-y_2+y_3\\0&0&0&0&y_1-y_2-y_3+y_4\end{array}\right)
\end{gathered}
\]
The left part of the last matrix is the RREF of \(A\), with \(r=3\) and pivot columns
\(1,2,4\), and \(z_4=y_1-y_2-y_3+y_4\). By Proposition~\ref{r2-prop-solutions}, the system is
consistent if and only if \(y_1-y_2-y_3+y_4=0\); then \(x_3=c\) is free and
\[
x_1=-2y_1+y_2+y_3-2c,\qquad x_2=y_1-y_3+c,\qquad x_3=c,\qquad x_4=y_1-y_2+y_3\qquad(c\in\R)
\]
For \(y=(4,7,4,7)'\), \(4-7-4+7=0\), and the solutions are
\(x=(3,0,0,1)'+c\,(-2,1,1,0)'\): the solution \((3,0,0,1)'\) plus the multiples of
\((-2,1,1,0)'\), which solve \(Ax=0\). Check: \(A(3,0,0,1)'=(4,7,4,7)'\), and
\(A(-2,1,1,0)'=0\) says \(\col_3(A)=2\col_1(A)-\col_2(A)\). For \(y=(1,0,0,0)'\),
\(y_1-y_2-y_3+y_4=1\ne0\), so \(Ax=y\) has no solution.
\end{example}

\subsection*{Homogeneous systems}

\begin{theorem}[\citeat{hoffman-kunze}{\S1.4, Theorems~6, 7}]\label{r2-thm-homogeneous}
\begin{enumerate}
\item If \(A\in\R^{m\times n}\) with \(m<n\) (fewer equations than unknowns), then \(Ax=0\)
has a nontrivial solution.
\item For \(A\in\R^{n\times n}\), \(Ax=0\) has only the trivial solution if and only if the
RREF of \(A\) is \(I_n\).
\end{enumerate}
\end{theorem}
\begin{proof}
Let \(R\) be the RREF of \(A\), with \(r\) nonzero rows. (a) \(r\leqslant m<n\), so there is
a free unknown; giving it the value \(1\) and the other free unknowns the value \(0\) gives a
solution \(x\ne0\) (Proposition~\ref{r2-prop-solutions}(b)). (b) If \(R=I\), then \(Rx=x\),
so \(x=0\) is the only solution (Theorem~\ref{r2-thm-rowops}(b)). If \(Ax=0\) has only the trivial
solution, there is no free unknown, so \(r=n\); then \(1\leqslant k_1<\dots<k_n\leqslant n\) forces \(k_i=i\), so column
\(i\) of \(R\) is \(e_i\) and \(R=I\).
\end{proof}

\subsection*{Inverses}

\begin{definition}[\citeat{hoffman-kunze}{\S1.6}, \citeat{lang}{Ch.~II, \S3, Exercise~34}]\label{r2-def-inverse}
\(A\in\R^{n\times n}\) is \textbf{invertible} if \(AB=BA=I_n\) for some
\(B\in\R^{n\times n}\). Such a \(B\) is unique: if \(BA=I\) and \(AB_1=I\), then
\(B=B(AB_1)=(BA)B_1=B_1\). It is called the \textbf{inverse} of \(A\), written \(A^{-1}\). For
example, \(\diag(d_1,\dots,d_n)\) with all \(d_i\ne0\) has inverse
\(\diag(d_1^{-1},\dots,d_n^{-1})\), since by Proposition~\ref{r1-prop-diag}(b) the products in both orders are
\(\diag(1,\dots,1)=I\).
\end{definition}

\begin{theorem}[\citeat{hoffman-kunze}{\S1.6, Theorem~10}, \citeat{lang}{Ch.~II, \S3, Exercise~32}, \citeat{shakarchi}{}]\label{r2-thm-inverse-props}
Let \(A,B\in\R^{n\times n}\) be invertible.
\begin{enumerate}
\item \(A^{-1}\) is invertible and \((A^{-1})^{-1}=A\).
\item \(AB\) is invertible and \((AB)^{-1}=B^{-1}A^{-1}\); more generally
\((A_1\cdots A_k)^{-1}=A_k^{-1}\cdots A_1^{-1}\) for invertible \(A_1,\dots,A_k\).
\item \(A'\) is invertible and \((A')^{-1}=(A^{-1})'\).
\end{enumerate}
\end{theorem}
\begin{proof}
(a) \(AA^{-1}=A^{-1}A=I\) says that \(A\) is the inverse of \(A^{-1}\). (b)
\((AB)(B^{-1}A^{-1})=A(BB^{-1})A^{-1}=AA^{-1}=I\), and similarly
\((B^{-1}A^{-1})(AB)=I\); for \(k\) factors repeat this. (c)
\(A'(A^{-1})'=(A^{-1}A)'=I\) and \((A^{-1})'A'=(AA^{-1})'=I\) by
Theorem~\ref{r1-thm-transpose}.
\end{proof}

\subsection*{Elementary matrices}

\begin{definition}[\citeat{hoffman-kunze}{\S1.5}]\label{r2-def-elementary}
An \textbf{elementary matrix} is a matrix \(E=\omega(I_m)\) obtained from \(I_m\) by one
elementary row operation \(\omega\). For \(r\ne s\) we write:
\begin{itemize}
\item \(D_r(c)\) (\(c\ne0\)): \(I_m\) with its \((r,r)\) entry replaced by \(c\)
(from \(R_r\leftarrow cR_r\));
\item \(T_{rs}(c)\): \(I_m\) with its \((r,s)\) entry replaced by \(c\)
(from \(R_r\leftarrow R_r+cR_s\));
\item \(P_{rs}\): \(I_m\) with rows \(r\) and \(s\) interchanged (from \(R_r\leftrightarrow R_s\)).
\end{itemize}
For example, for \(m=4\):
\[
P_{23}=\begin{pmatrix}1&0&0&0\\0&0&1&0\\0&1&0&0\\0&0&0&1\end{pmatrix},\qquad
T_{31}(c)=\begin{pmatrix}1&0&0&0\\0&1&0&0\\c&0&1&0\\0&0&0&1\end{pmatrix},\qquad
D_3(c)=\begin{pmatrix}1&0&0&0\\0&1&0&0\\0&0&c&0\\0&0&0&1\end{pmatrix}
\]
\end{definition}

\begin{theorem}[\citeat{hoffman-kunze}{\S\S1.5--1.6, Theorems~9, 11}]\label{r2-thm-elementary}
Let \(\omega\) be an elementary row operation on matrices with \(m\) rows and \(E=\omega(I_m)\).
\begin{enumerate}
\item \(\omega(A)=EA\) for every \(A\in\R^{m\times n}\): a row operation is a left
multiplication.
\item \(E\) is invertible, and \(E^{-1}\) is the elementary matrix of the inverse operation:
\(D_r(c)^{-1}=D_r(c^{-1})\), \(T_{rs}(c)^{-1}=T_{rs}(-c)\), \(P_{rs}^{-1}=P_{rs}\).
\item If row operations \(\omega_1,\dots,\omega_k\) (in this order) take \(A\) to \(B\), then
\(B=PA\) with \(P=E_k\cdots E_2E_1\), \(E_j=\omega_j(I_m)\); and \(P\) is invertible.
\end{enumerate}
\end{theorem}
\begin{proof}
(a) By Theorem~\ref{r1-thm-views}(c), \(\row_i(EA)=\row_i(E)A\). An unchanged row of \(E\) is
\(e_i'\), and \(e_i'A=\row_i(A)\) (Proposition~\ref{r1-prop-units}(a)). For type (2), \(\row_r(E)=e_r'+ce_s'\),
so \(\row_r(EA)=\row_r(A)+c\row_s(A)\); types (1) and (3) are the same computation. (b) Let
\(\tilde\omega\) be the inverse operation (Theorem~\ref{r2-thm-rowops}(a)) and \(\tilde E=\tilde\omega(I)\). By (a),
\(E\tilde E=\omega(\tilde\omega(I))=I\) and \(\tilde EE=\tilde\omega(\omega(I))=I\). (c) By (a), \(B=E_k(\cdots(E_1A))=(E_k\cdots E_1)A\),
and a product of invertible matrices is invertible (Theorem~\ref{r2-thm-inverse-props}(b)).
\end{proof}

\subsection*{Deciding invertibility and computing inverses}

\begin{theorem}[\citeat{hoffman-kunze}{\S1.6, Theorems~12, 13}, \citeat{bapat}{Result~2.11}]\label{r2-thm-invertible}
For \(A\in\R^{n\times n}\) the following are equivalent:
\begin{enumerate}[label=(\roman*)]
\item \(A\) is invertible;
\item the RREF of \(A\) is \(I_n\);
\item \(Ax=0\) has only the trivial solution;
\item \(Ax=b\) has a solution for every \(b\in\R^n\);
\item \(A\) is a product of elementary matrices.
\end{enumerate}
Moreover, if \(BA=I\) or \(AB=I\) for some \(B\in\R^{n\times n}\), then \(A\) is invertible
and \(B=A^{-1}\).
\end{theorem}
\begin{proof}
Let \(R=PA\) be the RREF of \(A\), with \(P=E_k\cdots E_1\) (Theorem~\ref{r2-thm-elementary}(c)).
\begin{itemize}
\item If \(R=I\), then \(A=P^{-1}=E_1^{-1}\cdots E_k^{-1}\) is a product of elementary
matrices, hence invertible (Theorem~\ref{r2-thm-inverse-props}(b)); \(Ax=0\) has only the trivial
solution (Theorem~\ref{r2-thm-homogeneous}(b)); and \(x=A^{-1}b\) solves \(Ax=b\).
\item If \(R\ne I\), then \(R\) has a zero row, hence a zero last row (by (c) of
Definition~\ref{r2-def-rref}); otherwise \(r=n\), so \(1\leqslant k_1<\dots<k_n\leqslant n\) forces \(k_i=i\),
column \(i\) of \(R\) is \(e_i\) for every \(i\), and \(R=I\). Then \(Ax=0\) has a solution \(x\ne0\)
(Theorem~\ref{r2-thm-homogeneous}(b)), so \(A\) has no inverse (else \(x=A^{-1}Ax=0\)
\(\contradiction\)), and \(A\) is not a product of elementary matrices (such a product is
invertible, Theorems~\ref{r2-thm-elementary}(b) and~\ref{r2-thm-inverse-props}(b)). For
\(b=P^{-1}e_n\), \(Ax=b\) holds if and only if \(Rx=PAx=Pb=e_n\) does, as \(P\) is invertible
(Theorem~\ref{r2-thm-elementary}(c)); the last equation of \(Rx=e_n\) reads \(0=1\), so (iv) fails.
\end{itemize}
Thus each of (i)--(v) holds exactly when \(R=I\). If \(BA=I\), then \(Ax=0\) gives
\(x=BAx=0\), so \(A\) is invertible by (iii), and \(B=BAA^{-1}=A^{-1}\). If \(AB=I\), then
\(B\) is invertible with \(B^{-1}=A\) by the case just shown, so \(A=B^{-1}\) is invertible
with \(A^{-1}=B\).
\end{proof}

\begin{corollary}[Computing \(A^{-1}\); \citeat{hoffman-kunze}{\S1.6, Corollaries to Theorem~12}, \citeat{abadir-magnus}{Exercises~6.22(b), 6.23}]\label{r2-cor-computing-inverse}
Let \(A\in\R^{n\times n}\), and row-reduce the \(n\times2n\) matrix \([A\mid I_n]\) by Gauss--Jordan elimination.
\begin{enumerate}
\item If the left half becomes \(I_n\), the right half is \(A^{-1}\).
\item If a zero row appears in the left half, \(A\) is not invertible.
\end{enumerate}
\end{corollary}
\begin{proof}
\begin{enumerate}
\item \([A\mid I]\to[I\mid B]\) means \(E[A\mid I]=[EA\mid E]=[I\mid B]\) for the product \(E\) of
the elementary matrices used (Theorem~\ref{r2-thm-elementary}(c); \(E\) multiplies each column
separately, Theorem~\ref{r1-thm-views}(a)), so \(EA=I\) and \(B=E=A^{-1}\) by the last sentence of
Theorem~\ref{r2-thm-invertible}.
\item The left half is then
\(EA\), with \(E\) the (invertible) product of the elementary matrices used so far
(Theorem~\ref{r2-thm-elementary}(c)). For every \(X\in\R^{n\times n}\), \(EAX\) has the same zero
row (Theorem~\ref{r1-thm-views}(c)), so \(EAX\ne I_n\) and \(EA\) is not invertible; if \(A\) were
invertible, so would be \(EA\) (Theorem~\ref{r2-thm-inverse-props}(b)).\qedhere
\end{enumerate}
\end{proof}

\begin{proposition}[\citeat{hefferon}{Three.IV.4, Corollary~4.11, Exercise~4.21}]\label{r2-prop-inverse-2x2}
Let \(A=\begin{pmatrix}a&b\\c&d\end{pmatrix}\in\R^{2\times2}\). Then \(A\) is invertible if and only if
\(ad-bc\ne0\), and in that case
\[
A^{-1}=\frac1{ad-bc}\begin{pmatrix}d&-b\\-c&a\end{pmatrix}
\]
\end{proposition}
\begin{proof}
We row-reduce \([A\mid I_2]\) and read the result by Corollary~\ref{r2-cor-computing-inverse}. Put \(\Delta=ad-bc\).
\begin{itemize}
\item Let \(a\ne0\). Then
\[
\left(\begin{array}{cc|cc}a&b&1&0\\c&d&0&1\end{array}\right)
\xrightarrow{R_2\leftarrow R_2-\frac caR_1}
\left(\begin{array}{cc|cc}a&b&1&0\\0&\frac\Delta a&-\frac ca&1\end{array}\right)
\]
since \(d-\frac{cb}a=\frac\Delta a\). If \(\Delta=0\), the left half has a zero row, so \(A\) is not invertible
(Corollary~\ref{r2-cor-computing-inverse}(b)).
\item Let \(a\ne0\) and \(\Delta\ne0\). Then
\[
\xrightarrow[R_2\leftarrow\frac a\Delta R_2]{R_1\leftarrow\frac1aR_1}
\left(\begin{array}{cc|cc}1&\frac ba&\frac1a&0\\0&1&-\frac c\Delta&\frac a\Delta\end{array}\right)
\xrightarrow{R_1\leftarrow R_1-\frac baR_2}
\left(\begin{array}{cc|cc}1&0&\frac d\Delta&-\frac b\Delta\\0&1&-\frac c\Delta&\frac a\Delta\end{array}\right)
\]
using \(\frac1a+\frac{bc}{a\Delta}=\frac{\Delta+bc}{a\Delta}=\frac d\Delta\). The left half is \(I_2\), so \(A\) is
invertible and the right half is the stated \(A^{-1}\) (Corollary~\ref{r2-cor-computing-inverse}(a)).
\item Let \(a=0\), so that \(\Delta=-bc\). If \(c=0\), the first column of \(A\) is zero, so \(Ae_1=0\) and \(A\) is not
invertible (Theorem~\ref{r2-thm-invertible}, (i)\,\(\Leftrightarrow\)\,(iii)). Let \(c\ne0\). Then
\[
\left(\begin{array}{cc|cc}0&b&1&0\\c&d&0&1\end{array}\right)
\xrightarrow{R_1\leftrightarrow R_2}
\left(\begin{array}{cc|cc}c&d&0&1\\0&b&1&0\end{array}\right)
\]
If \(b=0\), the left half has a zero row, and \(A\) is not invertible (Corollary~\ref{r2-cor-computing-inverse}(b)).
So for \(a=0\), \(A\) is invertible only if
\(b\ne0\) and \(c\ne0\), that is, only if \(\Delta=-bc\ne0\).
\item Let \(a=0\) and \(b,c\ne0\). Then
\[
\xrightarrow[R_2\leftarrow\frac1bR_2]{R_1\leftarrow\frac1cR_1}
\left(\begin{array}{cc|cc}1&\frac dc&0&\frac1c\\0&1&\frac1b&0\end{array}\right)
\xrightarrow{R_1\leftarrow R_1-\frac dcR_2}
\left(\begin{array}{cc|cc}1&0&-\frac d{bc}&\frac1c\\0&1&\frac1b&0\end{array}\right)
\]
and with \(\Delta=-bc\) the right half is \(\frac1\Delta\left(\begin{smallmatrix}d&-b\\-c&0\end{smallmatrix}\right)\),
the stated \(A^{-1}\) for \(a=0\) (Corollary~\ref{r2-cor-computing-inverse}(a)).
\end{itemize}
In both cases \(A\) is invertible exactly when \(\Delta\ne0\), with the stated inverse.
\end{proof}
The formula is easy to check by multiplication:
\(\left(\begin{smallmatrix}a&b\\c&d\end{smallmatrix}\right)\left(\begin{smallmatrix}d&-b\\-c&a\end{smallmatrix}\right)
=\left(\begin{smallmatrix}ad-bc&0\\0&ad-bc\end{smallmatrix}\right)
=\left(\begin{smallmatrix}d&-b\\-c&a\end{smallmatrix}\right)\left(\begin{smallmatrix}a&b\\c&d\end{smallmatrix}\right)\).

\begin{example}\label{r2-exm-inverse-2x2}
For \(A=\left(\begin{smallmatrix}3&5\\1&2\end{smallmatrix}\right)\), \(ad-bc=6-5=1\), so
Proposition~\ref{r2-prop-inverse-2x2} gives
\(A^{-1}=\left(\begin{smallmatrix}2&-5\\-1&3\end{smallmatrix}\right)\). For
\(B=\left(\begin{smallmatrix}2&6\\1&3\end{smallmatrix}\right)\), \(ad-bc=6-6=0\), so \(B\) is
not invertible; indeed \(x=(3,-1)'\ne0\) satisfies \(Bx=(6-6,\ 3-3)'=0\)
(Theorem~\ref{r2-thm-invertible}, (i)\,\(\Leftrightarrow\)\,(iii)).
\end{example}

\begin{example}[cf.\ \citeat{hoffman-kunze}{\S1.6, Example~15}]\label{r2-exm-inverse-gj}
Find the inverse of
\[
A=\begin{pmatrix}1&-1&1&-1\\-1&1&-2&1\\2&-1&2&-1\\1&-1&1&0\end{pmatrix}
\]
and write \(A\) and \(A^{-1}\) as products of elementary matrices.

\textit{Solution.} Forward elimination of \([A\mid I_4]\):
\[
\begin{gathered}
\left(\begin{array}{cccc|cccc}1&-1&1&-1&1&0&0&0\\-1&1&-2&1&0&1&0&0\\2&-1&2&-1&0&0&1&0\\1&-1&1&0&0&0&0&1\end{array}\right)
\xrightarrow[R_4\leftarrow R_4-R_1]{R_2\leftarrow R_2+R_1,\ R_3\leftarrow R_3-2R_1}
\left(\begin{array}{cccc|cccc}1&-1&1&-1&1&0&0&0\\0&0&-1&0&1&1&0&0\\0&1&0&1&-2&0&1&0\\0&0&0&1&-1&0&0&1\end{array}\right)
\\
\xrightarrow{R_2\leftrightarrow R_3}
\left(\begin{array}{cccc|cccc}1&-1&1&-1&1&0&0&0\\0&1&0&1&-2&0&1&0\\0&0&-1&0&1&1&0&0\\0&0&0&1&-1&0&0&1\end{array}\right)
\xrightarrow{R_3\leftarrow-R_3}
\left(\begin{array}{cccc|cccc}1&-1&1&-1&1&0&0&0\\0&1&0&1&-2&0&1&0\\0&0&1&0&-1&-1&0&0\\0&0&0&1&-1&0&0&1\end{array}\right)
\end{gathered}
\]
After the first arrow the entry of row \(2\) in column \(2\) is \(0\), so Step 2 interchanges
rows \(2\) and \(3\), and Step 3 multiplies the new row \(3\) by \(-1\). Backward elimination,
clearing columns \(4\), \(3\), \(2\) in turn:
\[
\begin{gathered}
\xrightarrow[R_2\leftarrow R_2-R_4]{R_1\leftarrow R_1+R_4}
\left(\begin{array}{cccc|cccc}1&-1&1&0&0&0&0&1\\0&1&0&0&-1&0&1&-1\\0&0&1&0&-1&-1&0&0\\0&0&0&1&-1&0&0&1\end{array}\right)
\xrightarrow{R_1\leftarrow R_1-R_3}
\left(\begin{array}{cccc|cccc}1&-1&0&0&1&1&0&1\\0&1&0&0&-1&0&1&-1\\0&0&1&0&-1&-1&0&0\\0&0&0&1&-1&0&0&1\end{array}\right)
\\
\xrightarrow{R_1\leftarrow R_1+R_2}
\left(\begin{array}{cccc|cccc}1&0&0&0&0&1&1&0\\0&1&0&0&-1&0&1&-1\\0&0&1&0&-1&-1&0&0\\0&0&0&1&-1&0&0&1\end{array}\right)
\end{gathered}
\]
The left half has become \(I_4\), so by Corollary~\ref{r2-cor-computing-inverse}(a)
\[
A^{-1}=\begin{pmatrix}0&1&1&0\\-1&0&1&-1\\-1&-1&0&0\\-1&0&0&1\end{pmatrix}
\]
Check: \(\row_3(A)A^{-1}=(2,-1,2,-1)A^{-1}=(0+1-2+1,\ 2+0-2+0,\ 2-1+0+0,\ 0+1+0-1)=(0,0,1,0)\),
and similarly for the other rows. The nine operations have the elementary matrices
\[
\begin{gathered}
E_1=T_{21}(1),\ E_2=T_{31}(-2),\ E_3=T_{41}(-1),\ E_4=P_{23},\ E_5=D_3(-1)\\
E_6=T_{14}(1),\ E_7=T_{24}(-1),\ E_8=T_{13}(-1),\ E_9=T_{12}(1)
\end{gathered}
\]
By Theorem~\ref{r2-thm-elementary}(c), \(E_9E_8\cdots E_1A=I_4\), so \(A^{-1}=E_9E_8\cdots E_1\) by the last
sentence of Theorem~\ref{r2-thm-invertible}. By Theorem~\ref{r2-thm-inverse-props}(a), \(A=(A^{-1})^{-1}\); so, by
Theorems~\ref{r2-thm-inverse-props}(b) and~\ref{r2-thm-elementary}(b),
\[
A=E_1^{-1}E_2^{-1}\cdots E_9^{-1}
=T_{21}(-1)\,T_{31}(2)\,T_{41}(1)\,P_{23}\,D_3(-1)\,T_{14}(-1)\,T_{24}(1)\,T_{13}(1)\,T_{12}(-1)
\]
\end{example}

\begin{example}\label{r2-exm-inverse-4x4}
Find the inverse of
\[
P=\begin{pmatrix}1&1&1&1\\1&2&3&4\\1&3&6&10\\1&4&10&20\end{pmatrix}
\]
(each entry outside row \(1\) and column \(1\) is the sum of the entry above it and the entry
to its left; \(P\) is the \(4\times4\) \textbf{Pascal matrix}).

\textit{Solution.} Every pivot is \(1\), so no scaling is needed. Forward elimination:
\[
\begin{gathered}
\left(\begin{array}{cccc|cccc}1&1&1&1&1&0&0&0\\1&2&3&4&0&1&0&0\\1&3&6&10&0&0&1&0\\1&4&10&20&0&0&0&1\end{array}\right)
\xrightarrow[R_4\leftarrow R_4-R_1]{R_2\leftarrow R_2-R_1,\ R_3\leftarrow R_3-R_1}
\left(\begin{array}{cccc|cccc}1&1&1&1&1&0&0&0\\0&1&2&3&-1&1&0&0\\0&2&5&9&-1&0&1&0\\0&3&9&19&-1&0&0&1\end{array}\right)
\\
\xrightarrow[R_4\leftarrow R_4-3R_2]{R_3\leftarrow R_3-2R_2}
\left(\begin{array}{cccc|cccc}1&1&1&1&1&0&0&0\\0&1&2&3&-1&1&0&0\\0&0&1&3&1&-2&1&0\\0&0&3&10&2&-3&0&1\end{array}\right)
\xrightarrow{R_4\leftarrow R_4-3R_3}
\left(\begin{array}{cccc|cccc}1&1&1&1&1&0&0&0\\0&1&2&3&-1&1&0&0\\0&0&1&3&1&-2&1&0\\0&0&0&1&-1&3&-3&1\end{array}\right)
\end{gathered}
\]
Backward elimination, clearing columns \(4\), \(3\), \(2\) in turn:
\[
\begin{gathered}
\xrightarrow[R_1\leftarrow R_1-R_4]{R_3\leftarrow R_3-3R_4,\ R_2\leftarrow R_2-3R_4}
\left(\begin{array}{cccc|cccc}1&1&1&0&2&-3&3&-1\\0&1&2&0&2&-8&9&-3\\0&0&1&0&4&-11&10&-3\\0&0&0&1&-1&3&-3&1\end{array}\right)
\\
\xrightarrow[R_1\leftarrow R_1-R_3]{R_2\leftarrow R_2-2R_3}
\left(\begin{array}{cccc|cccc}1&1&0&0&-2&8&-7&2\\0&1&0&0&-6&14&-11&3\\0&0&1&0&4&-11&10&-3\\0&0&0&1&-1&3&-3&1\end{array}\right)
\xrightarrow{R_1\leftarrow R_1-R_2}
\left(\begin{array}{cccc|cccc}1&0&0&0&4&-6&4&-1\\0&1&0&0&-6&14&-11&3\\0&0&1&0&4&-11&10&-3\\0&0&0&1&-1&3&-3&1\end{array}\right)
\end{gathered}
\]
The left half has become \(I_4\), so by Corollary~\ref{r2-cor-computing-inverse}(a)
\[
P^{-1}=\begin{pmatrix}4&-6&4&-1\\-6&14&-11&3\\4&-11&10&-3\\-1&3&-3&1\end{pmatrix}
\]
Check: \(\row_2(P)P^{-1}=(4-12+12-4,\ -6+28-33+12,\ 4-22+30-12,\ -1+6-9+4)=(0,1,0,0)\), and
similarly for the other rows. Every operation added an integer multiple of one row to another,
so the right half stayed an integer matrix throughout; this is why \(P^{-1}\) has integer
entries.
\end{example}

\begin{example}\label{r2-exm-singular}
The matrix \(A\) of Example~\ref{r2-exm-consistency} is not invertible: as the left part of the last matrix there
shows, its RREF is
\[
\begin{pmatrix}1&0&2&0\\0&1&-1&0\\0&0&0&1\\0&0&0&0\end{pmatrix}\ne I_4
\]
Accordingly \(Ax=0\) has the nontrivial solution
\(x=(-2,1,1,0)'\), and \(Ax=y\) has no solution for \(y=(1,0,0,0)'\)
(Theorem~\ref{r2-thm-invertible}, (iii) and (iv)).
\end{example}

\begin{corollary}[\citeat{hoffman-kunze}{\S1.6}, \citeat{harville}{\S8.2}]\label{r2-cor-products}
\begin{enumerate}
\item Let \(A_1,\dots,A_k\in\R^{n\times n}\). The product \(A_1A_2\cdots A_k\) is invertible if and only if every
\(A_j\) is invertible.
\item If \(A\in\R^{n\times n}\) is symmetric and invertible, then \(A^{-1}\) is symmetric.
\end{enumerate}
\end{corollary}
\begin{proof}
\begin{enumerate}
\item
\begin{itemize}
\item If every \(A_j\) is invertible, so is the product (Theorem~\ref{r2-thm-inverse-props}(b)).
\item Let \(M=A_1\cdots A_k\) be invertible. If \(A_kx=0\), then \(Mx=(A_1\cdots A_{k-1})A_kx=0\), so
\(x=M^{-1}Mx=0\). Hence \(A_k\) is invertible (Theorem~\ref{r2-thm-invertible}, (iii)\,\(\Rightarrow\)\,(i)).
\item Then \(A_1\cdots A_{k-1}=MA_k^{-1}\) is invertible (Theorem~\ref{r2-thm-inverse-props}(b)), and the same
argument, applied to these \(k-1\) factors, shows that \(A_{k-1}\) is invertible. Continuing in this way, every
\(A_j\) is invertible.
\end{itemize}
\item As \(A'=A\), Theorem~\ref{r2-thm-inverse-props}(c) gives \((A^{-1})'=(A')^{-1}=A^{-1}\).\qedhere
\end{enumerate}
\end{proof}
For instance, the Pascal matrix \(P\) of Example~\ref{r2-exm-inverse-4x4} is symmetric, and so is its inverse
\[
P^{-1}=\begin{pmatrix}4&-6&4&-1\\-6&14&-11&3\\4&-11&10&-3\\-1&3&-3&1\end{pmatrix}
\]
By (a), a product of square matrices with a factor that is not invertible is not invertible.

\subsection*{Row equivalence}

\begin{corollary}[\citeat{hoffman-kunze}{\S1.3, \S1.6, Corollaries to Theorem~12}, \citeat{abadir-magnus}{Exercise~6.21}, \citeat{hefferon}{One.III.2}]\label{r2-cor-row-equivalent}
Let \(A,B,G\in\R^{m\times n}\).
\begin{enumerate}
\item \(B\) is row-equivalent to \(A\) if and only if \(B=PA\) for some invertible \(P\in\R^{m\times m}\).
Moreover, if row operations take the \(m\times(n+m)\) matrix \([A\mid I_m]\) to \([B\mid P]\), then \(P\) is
invertible and \(B=PA\).
\item If \(B\) is row-equivalent to \(A\), then \(A\) is row-equivalent to \(B\). If, in addition, \(G\) is
row-equivalent to \(B\), then \(G\) is row-equivalent to \(A\).
\item \(B\) is row-equivalent to \(A\) if and only if \(A\) and \(B\) have the same RREF.
\end{enumerate}
\end{corollary}
\begin{proof}
\begin{enumerate}
\item
\begin{itemize}
\item If \(B\) is row-equivalent to \(A\), then \(B=PA\) with \(P=E_k\cdots E_1\) invertible
(Theorem~\ref{r2-thm-elementary}(c)).
\item Conversely, let \(B=PA\) with \(P\) invertible. By Theorem~\ref{r2-thm-invertible}
((i)\,\(\Rightarrow\)\,(v)), \(P=F_1F_2\cdots F_l\) with elementary matrices \(F_j\). Then
\(B=F_1(F_2(\cdots(F_lA)))\) is obtained from \(A\) by the row operations of \(F_l,F_{l-1},\dots,F_1\), in this
order (Theorem~\ref{r2-thm-elementary}(a)).
\item Let row operations take \([A\mid I_m]\) to \([B\mid P]\), and let \(E\) be the product of their elementary
matrices (Theorem~\ref{r2-thm-elementary}(c)). Since \(E\) multiplies each column separately
(Theorem~\ref{r1-thm-views}(a)), \([B\mid P]=E[A\mid I_m]=[EA\mid E]\). So \(P=E\) is invertible and
\(B=EA=PA\).
\end{itemize}
\item
\begin{itemize}
\item Let row operations \(\omega_1,\dots,\omega_k\) take \(A\) to \(B\), and let \(\tilde\omega_j\) be the
operation that undoes \(\omega_j\) (Theorem~\ref{r2-thm-rowops}(a)). Applying
\(\tilde\omega_k,\dots,\tilde\omega_1\), in this order, to \(B\) undoes \(\omega_k,\dots,\omega_1\) in turn and
gives \(A\).
\item If row operations \(\omega_1',\dots,\omega_l'\) take \(B\) to \(G\), then
\(\omega_1,\dots,\omega_k,\omega_1',\dots,\omega_l'\) take \(A\) to \(G\).
\end{itemize}
\item
\begin{itemize}
\item Let \(B\) be row-equivalent to \(A\), and let \(S\) be the RREF of \(B\). Operations taking \(A\) to \(B\),
followed by operations taking \(B\) to \(S\), take \(A\) to \(S\); so \(S\) is the RREF of \(A\)
(Theorem~\ref{r2-thm-rref-unique}).
\item Let \(A\) and \(B\) have the same RREF \(R\). Then \(R\) is row-equivalent to \(A\), and \(B\) is
row-equivalent to \(R\) (by (b)); so \(B\) is row-equivalent to \(A\) (by (b)).\qedhere
\end{itemize}
\end{enumerate}
\end{proof}

\needspace{14\baselineskip}
\begin{example}[cf.\ \citeat{hoffman-kunze}{\S1.6}]\label{r2-exm-row-equivalent}
Let
\[
A=\begin{pmatrix}1&2&0&1\\2&3&1&1\\1&1&1&1\\2&2&2&1\end{pmatrix}
\]
be the matrix of Example~\ref{r2-exm-consistency}, and \(R\) its RREF. Find an invertible \(P\) with \(PA=R\).

\textit{Solution.} Row-reduce \([A\mid I_4]\) with the operations of Example~\ref{r2-exm-consistency}:
\[
\begin{gathered}
\left(\begin{array}{cccc|cccc}1&2&0&1&1&0&0&0\\2&3&1&1&0&1&0&0\\1&1&1&1&0&0&1&0\\2&2&2&1&0&0&0&1\end{array}\right)
\xrightarrow[R_4\leftarrow R_4-2R_1]{R_2\leftarrow R_2-2R_1,\ R_3\leftarrow R_3-R_1}
\left(\begin{array}{cccc|cccc}1&2&0&1&1&0&0&0\\0&-1&1&-1&-2&1&0&0\\0&-1&1&0&-1&0&1&0\\0&-2&2&-1&-2&0&0&1\end{array}\right)
\\
\xrightarrow{R_2\leftarrow-R_2}
\left(\begin{array}{cccc|cccc}1&2&0&1&1&0&0&0\\0&1&-1&1&2&-1&0&0\\0&-1&1&0&-1&0&1&0\\0&-2&2&-1&-2&0&0&1\end{array}\right)
\xrightarrow[R_4\leftarrow R_4+2R_2]{R_3\leftarrow R_3+R_2}
\left(\begin{array}{cccc|cccc}1&2&0&1&1&0&0&0\\0&1&-1&1&2&-1&0&0\\0&0&0&1&1&-1&1&0\\0&0&0&1&2&-2&0&1\end{array}\right)
\\
\xrightarrow{R_4\leftarrow R_4-R_3}
\left(\begin{array}{cccc|cccc}1&2&0&1&1&0&0&0\\0&1&-1&1&2&-1&0&0\\0&0&0&1&1&-1&1&0\\0&0&0&0&1&-1&-1&1\end{array}\right)
\xrightarrow[R_1\leftarrow R_1-R_3]{R_2\leftarrow R_2-R_3}
\left(\begin{array}{cccc|cccc}1&2&0&0&0&1&-1&0\\0&1&-1&0&1&0&-1&0\\0&0&0&1&1&-1&1&0\\0&0&0&0&1&-1&-1&1\end{array}\right)
\\
\xrightarrow{R_1\leftarrow R_1-2R_2}
\left(\begin{array}{cccc|cccc}1&0&2&0&-2&1&1&0\\0&1&-1&0&1&0&-1&0\\0&0&0&1&1&-1&1&0\\0&0&0&0&1&-1&-1&1\end{array}\right)
\end{gathered}
\]
The left half is \(R\). By Corollary~\ref{r2-cor-row-equivalent}(a),
\[
P=\begin{pmatrix}-2&1&1&0\\1&0&-1&0\\1&-1&1&0\\1&-1&-1&1\end{pmatrix}\qquad\text{is invertible, and}\qquad
PA=R=\begin{pmatrix}1&0&2&0\\0&1&-1&0\\0&0&0&1\\0&0&0&0\end{pmatrix}
\]
\begin{itemize}
\item Check: \(\row_4(P)A=(1,2,0,1)-(2,3,1,1)-(1,1,1,1)+(2,2,2,1)=(0,0,0,0)\)
(Theorem~\ref{r1-thm-views}(c)), and likewise \(\row_1(P)A=(1,0,2,0)\), \(\row_2(P)A=(0,1,-1,0)\),
\(\row_3(P)A=(0,0,0,1)\).
\item For \(y\in\R^4\), \(Py=(-2y_1+y_2+y_3,\ y_1-y_3,\ y_1-y_2+y_3,\ y_1-y_2-y_3+y_4)'\) is the last column of the
last matrix in Example~\ref{r2-exm-consistency}: the same operations take \([A\mid y]\) to \([PA\mid Py]=[R\mid Py]\).
\item As \(P\) is invertible, \(Ax=y\) holds if and only if \(Rx=PAx=Py\) does. Row \(4\) of \(R\) is zero, so by
Proposition~\ref{r2-prop-solutions}(a), \(Ax=y\) is consistent exactly when \(\row_4(P)\,y=y_1-y_2-y_3+y_4=0\), the
condition found in Example~\ref{r2-exm-consistency}.
\end{itemize}
\end{example}

\subsection*{Solving \(Ax=b\)}

In practice \(Ax=b\) is solved by forward elimination of \([A\mid b]\) alone: an echelon form
already shows whether there is a solution, and how many.

\begin{theorem}[\citeat{hoffman-kunze}{\S1.4}]\label{thm:r2-consistency}
Let \(A\in\R^{m\times n}\) and \(b\in\R^m\), and let row operations take \([A\mid b]\) to a
matrix \([H\mid d]\) in echelon form (\(H\in\R^{m\times n}\), \(d\in\R^m\)). Then \(Ax=b\) is
consistent if and only if \([H\mid d]\) has no row \((0\ \cdots\ 0\mid c)\) with \(c\ne0\),
that is, if and only if the last column of \([H\mid d]\) is not a pivot column. Moreover, if row operations take
\([A\mid b]\) to any matrix, in echelon form or not, with a row \((0\ \cdots\ 0\mid c)\), \(c\ne0\), then \(Ax=b\)
has no solution. If \([H\mid d]\) has no such row, call \(x_j\) a pivot unknown when column \(j\) of \(H\) is a
pivot column and a free unknown otherwise, as in Proposition~\ref{r2-prop-solutions}; then each choice of values of
the free unknowns gives, by \textbf{back substitution} (solving the equations of the nonzero rows of \(H\), from the
last one up, for their pivot unknowns), exactly one solution, and every solution arises in this way.
\end{theorem}
\begin{proof}
\begin{itemize}
\item By Theorem~\ref{r2-thm-rowops}, \(Ax=b\) and \(Hx=d\) have the same solutions. In an
echelon form, a row \((0\ \cdots\ 0\mid c)\) with \(c\ne0\) is exactly a row whose pivot lies
in the last column.
\item Let row operations take \([A\mid b]\) to a matrix \([B\mid z]\) with a row \((0\ \cdots\ 0\mid c)\),
\(c\ne0\). The corresponding equation of \(Bx=z\) reads \(0=c\), and no \(x\) satisfies it; as \(Bx=z\) and
\(Ax=b\) have the same solutions (Theorem~\ref{r2-thm-rowops}(b)), \(Ax=b\) has no solution. This applies to
\([H\mid d]\), and to the matrix reached at any stage of the elimination, so one may stop there.
\item Suppose \([H\mid d]\) has no such row. Then every zero row of \(H\) is a zero row of \([H\mid d]\),
and its equation reads \(0=0\). Let rows \(1,\dots,r\) be the nonzero rows of \(H\), with pivots
in columns \(k_1<\dots<k_r\); so \(x_{k_1},\dots,x_{k_r}\) are the pivot unknowns, and the others are free.
\item Give the free unknowns any values. Row \(i\) of \(H=(h_{ij})\) is \(0\) before column
\(k_i\) and \(1\) there, so for \(i=r,r-1,\dots,1\) equation \(i\) reads
\(x_{k_i}+\sum_{j>k_i}h_{ij}x_j=d_i\). Each \(x_j\) with \(j>k_i\) is free or the pivot unknown
of a lower row, already found; so the equation determines \(x_{k_i}\) (back substitution). The resulting \(x\)
solves \(Hx=d\), hence \(Ax=b\).
\item Every solution arises in this way from the values of its free unknowns, since the
equations force the pivot unknowns; so each choice of values of the free unknowns gives exactly one
solution.\qedhere
\end{itemize}
\end{proof}

\begin{theorem}[\citeat{harville}{Theorem~11.2.3}, \citeat{bapat}{\S2.3, Ch.~2, Exercise~13}, \citeat{abadir-magnus}{Exercise~6.45}]\label{thm:r2-solution-structure}
Let \(A\in\R^{m\times n}\) and \(b\in\R^m\).
\begin{enumerate}
\item If \(x_p\) is one solution of \(Ax=b\) (a \textbf{particular solution}), then the
solutions of \(Ax=b\) are exactly the vectors \(x_p+z\) with \(Az=0\).
\item \(Ax=b\) has no solution, exactly one solution, or infinitely many solutions. If it is
consistent, it has exactly one solution if and only if there are no free unknowns, that is,
every column of \(H\) in Theorem~\ref{thm:r2-consistency} is a pivot column.
\item If \(A\) is square, \(Ax=b\) has exactly one solution if and only if \(A\) is invertible;
this solution is \(x=A^{-1}b\).
\end{enumerate}
\end{theorem}
\begin{proof}
\begin{enumerate}
\item If \(Az=0\), then \(A(x_p+z)=Ax_p+Az=b\). Conversely, if \(Ax=b\), then \(z=x-x_p\)
satisfies \(Az=b-b=0\), and \(x=x_p+z\).
\item Suppose \(Ax=b\) is consistent, and take \([H\mid d]\) as in
Theorem~\ref{thm:r2-consistency}.
\begin{itemize}
\item By that theorem, each choice of values of the free unknowns gives exactly one solution, and every solution
arises in this way; different choices give different solutions (they differ in a free unknown, an entry of \(x\)).
\item Without free unknowns there is therefore exactly one solution; with a free unknown, each
of its values gives a different solution, so there are infinitely many.
\end{itemize}
\item
\begin{itemize}
\item If \(A\) is invertible, \(x=A^{-1}b\) is a solution, since \(A(A^{-1}b)=b\); and every
solution satisfies \(x=A^{-1}(Ax)=A^{-1}b\).
\item If \(Ax=b\) has exactly one solution \(x_p\), then \(Az=0\) only for \(z=0\), since
otherwise \(x_p+z\) would be a second solution by (a). So \(A\) is invertible by
Theorem~\ref{r2-thm-invertible} ((iii)\,\(\Rightarrow\)\,(i)).\qedhere
\end{itemize}
\end{enumerate}
\end{proof}
So, to solve \(Ax=b\): reduce \([A\mid b]\) by forward elimination to an echelon form; if a row
\((0\ \cdots\ 0\mid c)\) with \(c\ne0\) appears, there is no solution; otherwise choose the free
unknowns arbitrarily and find the pivot unknowns by back substitution (or continue to the RREF
and use Proposition~\ref{r2-prop-solutions}).

\begin{example}\label{r2-exm-unique}
Solve the system
\(x_1-x_2+x_3-x_4=1\), \(-x_1+x_2-2x_3+x_4=0\), \(2x_1-x_2+2x_3-x_4=3\), \(x_1-x_2+x_3=3\).

\textit{Solution.} The coefficient matrix is the matrix \(A\) of
Example~\ref{r2-exm-inverse-gj}, and \(b=(1,0,3,3)'\). Forward elimination of \([A\mid b]\)
uses the operations of the forward elimination there:
\[
\begin{gathered}
\left(\begin{array}{cccc|c}1&-1&1&-1&1\\-1&1&-2&1&0\\2&-1&2&-1&3\\1&-1&1&0&3\end{array}\right)
\xrightarrow[R_4\leftarrow R_4-R_1]{R_2\leftarrow R_2+R_1,\ R_3\leftarrow R_3-2R_1}
\left(\begin{array}{cccc|c}1&-1&1&-1&1\\0&0&-1&0&1\\0&1&0&1&1\\0&0&0&1&2\end{array}\right)
\\
\xrightarrow{R_2\leftrightarrow R_3}
\left(\begin{array}{cccc|c}1&-1&1&-1&1\\0&1&0&1&1\\0&0&-1&0&1\\0&0&0&1&2\end{array}\right)
\xrightarrow{R_3\leftarrow-R_3}
\left(\begin{array}{cccc|c}1&-1&1&-1&1\\0&1&0&1&1\\0&0&1&0&-1\\0&0&0&1&2\end{array}\right)
\end{gathered}
\]
\begin{itemize}
\item The last matrix is an echelon form with no row \((0\ 0\ 0\ 0\mid c)\), \(c\ne0\), so the
system is consistent (Theorem~\ref{thm:r2-consistency}).
\item Every column of the left part is a pivot column, so there are no free unknowns and the
solution is unique (Theorem~\ref{thm:r2-solution-structure}(b)).
\item Back substitution: \(x_4=2\); \(x_3=-1\); \(x_2=1-x_4=-1\);
\(x_1=1+x_2-x_3+x_4=1-1+1+2=3\). So \(x=(3,-1,-1,2)'\).
\item Check, with \(A^{-1}\) from Example~\ref{r2-exm-inverse-gj}
(Theorem~\ref{thm:r2-solution-structure}(c)):
\(A^{-1}b=(0+0+3+0,\ -1+0+3-3,\ -1+0+0+0,\ -1+0+0+3)'=(3,-1,-1,2)'\).
\end{itemize}
\end{example}

\begin{example}\label{r2-exm-infinite}
Solve the system
\(x_1+x_2+x_3-2x_4=2\), \(x_1+x_2+2x_3-3x_4=3\), \(2x_1+x_2+5x_3-4x_4=7\).

\textit{Solution.} The coefficient matrix is the matrix \(A\) of
Example~\ref{r2-exm-gauss-jordan}, and \(b=(2,3,7)'\). Forward elimination of \([A\mid b]\)
uses the same operations as there:
\[
\begin{gathered}
\left(\begin{array}{cccc|c}1&1&1&-2&2\\1&1&2&-3&3\\2&1&5&-4&7\end{array}\right)
\xrightarrow[R_3\leftarrow R_3-2R_1]{R_2\leftarrow R_2-R_1}
\left(\begin{array}{cccc|c}1&1&1&-2&2\\0&0&1&-1&1\\0&-1&3&0&3\end{array}\right)
\\
\xrightarrow{R_2\leftrightarrow R_3}
\left(\begin{array}{cccc|c}1&1&1&-2&2\\0&-1&3&0&3\\0&0&1&-1&1\end{array}\right)
\xrightarrow{R_2\leftarrow-R_2}
\left(\begin{array}{cccc|c}1&1&1&-2&2\\0&1&-3&0&-3\\0&0&1&-1&1\end{array}\right)
\end{gathered}
\]
\begin{itemize}
\item The last matrix is an echelon form with no row \((0\ 0\ 0\ 0\mid c)\), \(c\ne0\): the
system is consistent (Theorem~\ref{thm:r2-consistency}).
\item The pivot columns are \(1,2,3\), so \(x_4\) is free and there are infinitely many
solutions (Theorem~\ref{thm:r2-solution-structure}(b)).
\item Back substitution with \(x_4=t\): \(x_3=1+t\); \(x_2=-3+3x_3=3t\);
\(x_1=2-x_2-x_3+2x_4=2-3t-(1+t)+2t=1-2t\). So
\[
x=\begin{pmatrix}1-2t\\3t\\1+t\\t\end{pmatrix}=\begin{pmatrix}1\\0\\1\\0\end{pmatrix}+t\begin{pmatrix}-2\\3\\1\\1\end{pmatrix}\qquad(t\in\R)
\]
\item This is the particular solution \(x_p=(1,0,1,0)'\) plus the solutions \(t(-2,3,1,1)'\)
of \(Ax=0\) found in Example~\ref{r2-exm-gauss-jordan}, as
Theorem~\ref{thm:r2-solution-structure}(a) says.
\end{itemize}
\end{example}

\begin{example}[cf.\ \citeat{hoffman-kunze}{\S1.4, Exercise~6}]\label{r2-exm-inconsistent}
Show that the system
\(x_1-2x_2+x_3+2x_4=1\), \(x_1-x_2-x_3+x_4=2\), \(x_1+x_2-5x_3-x_4=3\) has no solution.

\textit{Solution.} Forward elimination of the augmented matrix \([A\mid b]\), \(b=(1,2,3)'\):
\[
\begin{gathered}
\left(\begin{array}{cccc|c}1&-2&1&2&1\\1&-1&-1&1&2\\1&1&-5&-1&3\end{array}\right)
\xrightarrow[R_3\leftarrow R_3-R_1]{R_2\leftarrow R_2-R_1}
\left(\begin{array}{cccc|c}1&-2&1&2&1\\0&1&-2&-1&1\\0&3&-6&-3&2\end{array}\right)
\\
\xrightarrow{R_3\leftarrow R_3-3R_2}
\left(\begin{array}{cccc|c}1&-2&1&2&1\\0&1&-2&-1&1\\0&0&0&0&-1\end{array}\right)
\xrightarrow{R_3\leftarrow-R_3}
\left(\begin{array}{cccc|c}1&-2&1&2&1\\0&1&-2&-1&1\\0&0&0&0&1\end{array}\right)
\end{gathered}
\]
The last matrix is an echelon form whose last column is a pivot column: row \(3\) reads
\(0=1\). So the system has no solution (Theorem~\ref{thm:r2-consistency}). Indeed, the
coefficient matrix \(A\) has \(\row_3(A)=3\row_2(A)-2\row_1(A)\), while
\(3b_2-2b_1=4\ne3=b_3\): the third equation contradicts the first two.
\end{example}

\subsection*{Triangular matrices}

\begin{theorem}[\citeat{hoffman-kunze}{\S1.6, Exercise~9}, \citeat{harville}{Corollaries~8.5.6, 8.5.8, 8.5.9}, \citeat{shakarchi}{Ch.~VI, \S5, Exercise~2}, \citeat{abadir-magnus}{Exercise~7.60(a)}]\label{r2-thm-triangular}
Let \(A=(a_{ij})\in\R^{n\times n}\) be upper triangular or lower triangular.
\begin{enumerate}
\item \(A\) is invertible if and only if \(a_{ii}\ne0\) for \(i=1,\dots,n\).
\item If \(A\) is invertible, then \(A^{-1}\) is triangular of the same kind as \(A\), with diagonal entries
\(1/a_{11},\dots,1/a_{nn}\).
\end{enumerate}
In particular, \(\diag(d_1,\dots,d_n)\) is invertible if and only if every \(d_i\ne0\), and the inverse of a
\textbf{unit triangular} matrix (a triangular matrix with all diagonal entries \(1\)) is unit triangular of the
same kind.
\end{theorem}
\begin{proof}
\begin{enumerate}
\item By Theorem~\ref{r2-thm-invertible} ((i)\,\(\Leftrightarrow\)\,(iii)), \(A\) is invertible if and only if
\(Ax=0\) has only the trivial solution.
\begin{itemize}
\item Let \(A\) be lower triangular, so that equation \(i\) of \(Ax=0\) reads \(a_{i1}x_1+\dots+a_{ii}x_i=0\). If
every \(a_{ii}\ne0\) and \(Ax=0\), then equation \(1\) gives \(x_1=0\), equation \(2\) then gives
\(a_{22}x_2=0\), so \(x_2=0\), and so on: \(x=0\).
\item Let \(A\) be lower triangular with some \(a_{kk}=0\), and let \(k\) be the smallest such index. Put
\(x_1=\dots=x_{k-1}=0\). Then equations \(1,\dots,k-1\) hold, and so does equation \(k\), which reads
\(a_{kk}x_k=0\).
\item Equations \(k+1,\dots,n\) are \(n-k\) equations in the \(n-k+1\) unknowns \(x_k,\dots,x_n\). If \(k<n\),
they have a nontrivial solution (Theorem~\ref{r2-thm-homogeneous}(a)); if \(k=n\), there are none, and we take
\(x_n=1\). Either way \(x\ne0\) and \(Ax=0\), so \(A\) is not invertible.
\item If \(A\) is upper triangular, then \(A'\) is lower triangular with the same diagonal entries, and \(A\) is
invertible if and only if \(A'\) is (Theorem~\ref{r2-thm-inverse-props}(c), with \((A')'=A\)). So the lower
triangular case gives the upper one.
\end{itemize}
\item
\begin{itemize}
\item Let \(A\) be upper triangular and invertible, so that every \(a_{ii}\ne0\) by (a). For an upper triangular
\(B\in\R^{n\times n}\), \(AB\) is upper triangular with diagonal entries \(a_{ii}b_{ii}\)
(Proposition~\ref{r1-prop-diag}(c)), and \((AB)_{ij}=\sum_{k=i}^ja_{ik}b_{kj}\) for \(i<j\), since \(a_{ik}=0\) for
\(k<i\) and \(b_{kj}=0\) for \(k>j\).
\item Hence an upper triangular \(B\) satisfies \(AB=I_n\) if and only if
\[
b_{jj}=\frac1{a_{jj}},\qquad b_{ij}=-\frac1{a_{ii}}\sum_{k=i+1}^ja_{ik}b_{kj}\qquad(1\leqslant i<j\leqslant n)
\]
\item For each column \(j\), these formulas give \(b_{jj}\) and then \(b_{j-1,j},\dots,b_{1j}\) in turn, since
\(b_{ij}\) uses only the \(b_{kj}\) with \(k>i\). So such a \(B\) exists, and \(B=A^{-1}\)
(Theorem~\ref{r2-thm-invertible}): \(A^{-1}\) is upper triangular with diagonal entries \(1/a_{ii}\).
\item If \(A\) is lower triangular and invertible, then \(A'\) is upper triangular and invertible, with the same
diagonal; so \((A')^{-1}\) is upper triangular with diagonal entries \(1/a_{ii}\), and
\(A^{-1}=((A')^{-1})'\) (Theorem~\ref{r2-thm-inverse-props}(c)) is lower triangular with the same diagonal.
\end{itemize}
\end{enumerate}
A diagonal matrix is both upper and lower triangular, and a unit triangular matrix has \(1/a_{ii}=1\); this gives
the last sentence.
\end{proof}

\begin{example}\label{r2-exm-triangular}
Find the inverse of
\[
U=\begin{pmatrix}1&-2&-2&3\\0&-1&-1&1\\0&0&1&1\\0&0&0&1\end{pmatrix}
\]

\textit{Solution.} \(U\) is upper triangular with diagonal entries \(1,-1,1,1\), all nonzero, so it is invertible,
and \(B=U^{-1}=(b_{ij})\) is upper triangular with \(b_{11}=1\), \(b_{22}=-1\), \(b_{33}=b_{44}=1\)
(Theorem~\ref{r2-thm-triangular}). The other entries come from
\(b_{ij}=-\frac1{u_{ii}}\sum_{k=i+1}^ju_{ik}b_{kj}\), column by column, from the bottom up:
\begin{itemize}
\item column \(2\): \(b_{12}=-u_{12}b_{22}=-(-2)(-1)=-2\);
\item column \(3\): \(b_{23}=u_{23}b_{33}=(-1)(1)=-1\) (here \(-1/u_{22}=1\)), and
\(b_{13}=-(u_{12}b_{23}+u_{13}b_{33})=-\bigl((-2)(-1)+(-2)(1)\bigr)=0\);
\item column \(4\): \(b_{34}=-u_{34}b_{44}=-1\), \(b_{24}=u_{23}b_{34}+u_{24}b_{44}=(-1)(-1)+(1)(1)=2\) (again
\(-1/u_{22}=1\)), and
\(b_{14}=-(u_{12}b_{24}+u_{13}b_{34}+u_{14}b_{44})=-\bigl((-2)(2)+(-2)(-1)+(3)(1)\bigr)=-1\).
\end{itemize}
So
\[
U^{-1}=\begin{pmatrix}1&-2&0&-1\\0&-1&-1&2\\0&0&1&-1\\0&0&0&1\end{pmatrix}
\]
Check: \(\row_1(U)\,U^{-1}=(1,\ -2+2,\ 2-2,\ -1-4+2+3)=(1,0,0,0)\), and similarly for the other rows.
\end{example}

\begin{theorem}[\citeat{harville}{\S8.5d, Corollary~8.5.9}, \citeat{abadir-magnus}{Exercise~7.60(a)}, \citeat{lang}{Ch.~II, \S3, Exercise~18(c)}]\label{r2-thm-unit-lower-inverse}
Let \(L=(l_{ij})\in\R^{n\times n}\) be unit lower triangular, and let \(N=I_n-L\); so \(N\) is strictly lower
triangular, with entries \(-l_{ij}\) for \(i>j\). Let \(M=L^{-1}=(m_{ij})\).
\begin{enumerate}
\item \(M\) is unit lower triangular, and for \(i>j\)
\[
m_{ij}=-\sum_{k=j}^{i-1}l_{ik}m_{kj}
\]
so column \(j\) of \(M\) is found from the top down: \(m_{jj}=1\), then \(m_{j+1,j},\dots,m_{nj}\) in turn.
\item \(M=I_n+NM\), and repeating this recursion gives
\[
L^{-1}=I_n+N+N^2+\dots+N^{n-1}
\]
\end{enumerate}
\end{theorem}
\begin{proof}
(a)
\begin{itemize}
\item \(L\) is invertible, and \(M\) is unit lower triangular (Theorem~\ref{r2-thm-triangular}(a),(b)).
\item Let \(i>j\). The \((i,j)\) entry of \(LM=I_n\) is \(\sum_kl_{ik}m_{kj}=0\). Here \(l_{ik}=0\) for \(k>i\),
\(l_{ii}=1\), and \(m_{kj}=0\) for \(k<j\), so the sum runs over \(j\leqslant k\leqslant i\):
\(m_{ij}+\sum_{k=j}^{i-1}l_{ik}m_{kj}=0\).
\item The right side of the formula uses only \(m_{jj},\dots,m_{i-1,j}\), entries of column \(j\) above row \(i\).
\end{itemize}
(b)
\begin{itemize}
\item \(LM=I_n\) reads \((I_n-N)M=I_n\), that is, \(M=I_n+NM\) (Theorem~\ref{r1-thm-rules}). (Entry \((i,j)\) of
this equation, \(i>j\), is the formula of (a).)
\item Substituting \(I_n+NM\) for \(M\) on the right gives \(M=I_n+N+N^2M\), and after \(p\) substitutions
\[
M=I_n+N+\dots+N^{p-1}+N^pM
\]
by induction on \(p\) (Theorem~\ref{r1-thm-rules}).
\item \(N\) is strictly lower triangular, so \(N^n=0\) (Proposition~\ref{r1-prop-strictly-triangular}(c)). With
\(p=n\), \(M=I_n+N+\dots+N^{n-1}\).\qedhere
\end{itemize}
\end{proof}

\begin{example}\label{r2-exm-unit-lower-inverse}
\begin{enumerate}
\item Find the inverse of
\[
L=\begin{pmatrix}1&0&0&0\\1&1&0&0\\1&2&1&0\\2&0&-1&1\end{pmatrix}
\]
\item Let \(S\in\R^{4\times4}\) have the entries \(1\) in the positions \((2,1),(3,2),(4,3)\) and \(0\) elsewhere. The
first-difference matrix \(I_4-S\) turns \(x=(x_1,x_2,x_3,x_4)'\) into \((x_1,x_2-x_1,x_3-x_2,x_4-x_3)'\). Find its
inverse.
\end{enumerate}

\textit{Solution.}
(a) \(L\) is unit lower triangular. By Theorem~\ref{r2-thm-unit-lower-inverse}(a), column by column from the top down,
with \(m_{jj}=1\):
\begin{itemize}
\item column \(1\): \(m_{21}=-l_{21}m_{11}=-1\), \(m_{31}=-(l_{31}m_{11}+l_{32}m_{21})=-(1-2)=1\),
\(m_{41}=-(l_{41}m_{11}+l_{42}m_{21}+l_{43}m_{31})=-(2+0-1)=-1\);
\item column \(2\): \(m_{32}=-l_{32}m_{22}=-2\), \(m_{42}=-(l_{42}m_{22}+l_{43}m_{32})=-(0+2)=-2\);
\item column \(3\): \(m_{43}=-l_{43}m_{33}=1\).
\end{itemize}
So
\[
L^{-1}=\begin{pmatrix}1&0&0&0\\-1&1&0&0\\1&-2&1&0\\-1&-2&1&1\end{pmatrix}
\]
The formula of Theorem~\ref{r2-thm-unit-lower-inverse}(b) gives the same: with
\[
N=I_4-L=\begin{pmatrix}0&0&0&0\\-1&0&0&0\\-1&-2&0&0\\-2&0&1&0\end{pmatrix},\qquad
N^2=\begin{pmatrix}0&0&0&0\\0&0&0&0\\2&0&0&0\\-1&-2&0&0\end{pmatrix},\qquad
N^3=\begin{pmatrix}0&0&0&0\\0&0&0&0\\0&0&0&0\\2&0&0&0\end{pmatrix}
\]
the sum \(I_4+N+N^2+N^3\) has the entries \(-1\), \(-1+2=1\), \(-2-1+2=-1\) in column \(1\) below the diagonal,
\(-2\) and \(0-2=-2\) in column \(2\), and \(1\) in column \(3\): it is \(L^{-1}\).

(b) \(I_4-S\) is unit lower triangular with \(N=S\). By Theorem~\ref{r2-thm-unit-lower-inverse}(b),
\((I_4-S)^{-1}=I_4+S+S^2+S^3\), where
\[
S=\begin{pmatrix}0&0&0&0\\1&0&0&0\\0&1&0&0\\0&0&1&0\end{pmatrix},\qquad
S^2=\begin{pmatrix}0&0&0&0\\0&0&0&0\\1&0&0&0\\0&1&0&0\end{pmatrix},\qquad
S^3=\begin{pmatrix}0&0&0&0\\0&0&0&0\\0&0&0&0\\1&0&0&0\end{pmatrix}
\]
So
\[
(I_4-S)^{-1}=\begin{pmatrix}1&0&0&0\\1&1&0&0\\1&1&1&0\\1&1&1&1\end{pmatrix}
\]
the matrix that turns \(y\) into its partial sums \((y_1,y_1+y_2,y_1+y_2+y_3,y_1+y_2+y_3+y_4)'\): summing undoes
differencing.
\end{example}

\subsection*{Column operations}

\begin{theorem}[\citeat{abadir-magnus}{Exercises~6.4(a), 6.10}]\label{r2-thm-column-ops}
The \textbf{elementary column operations} on matrices with \(n\) columns are, for column indices \(r\ne s\):
multiply column \(s\) by \(c\ne0\), written \(\col_s\leftarrow c\col_s\); add \(c\) times column \(r\) to column \(s\)
(\(c\in\R\)), written \(\col_s\leftarrow\col_s+c\col_r\); and interchange columns \(r\) and \(s\), written
\(\col_r\leftrightarrow\col_s\). Let \(\gamma\) be one of them and \(F=\gamma(I_n)\).
\begin{enumerate}
\item \(F\) is \(D_s(c)\), \(T_{rs}(c)\) or \(P_{rs}\), respectively. Moreover \(D_s(c)'=D_s(c)\),
\(T_{sr}(c)'=T_{rs}(c)\) and \(P_{rs}'=P_{rs}\).
\item \(\gamma(A)=AF\) for every \(A\in\R^{m\times n}\): a column operation is a right multiplication.
\item If column operations \(\gamma_1,\dots,\gamma_l\) (in this order) take \(A\) to \(B\), then \(B=AQ\) with
\(Q=F_1F_2\cdots F_l\), \(F_j=\gamma_j(I_n)\), and \(Q\) is invertible.
\end{enumerate}
\end{theorem}
\begin{proof}
\begin{enumerate}
\item
\begin{itemize}
\item Column \(s\) of \(I_n\) is \(e_s\). The three operations replace it by \(ce_s\), by \(e_s+ce_r\), or exchange it
with column \(r\). The results are \(I_n\) with its \((s,s)\) entry replaced by \(c\), \(I_n\) with its \((r,s)\)
entry replaced by \(c\), and \(I_n\) with rows \(r\) and \(s\) interchanged (exchanging columns \(r\) and \(s\) of
\(I_n\), or its rows \(r\) and \(s\), gives in both cases the matrix whose \((r,s)\) and \((s,r)\) entries are \(1\),
whose \((r,r)\) and \((s,s)\) entries are \(0\), and whose other entries are those of \(I_n\)); that is, \(D_s(c)\),
\(T_{rs}(c)\) and \(P_{rs}\) (Definition~\ref{r2-def-elementary}).
\item \(D_s(c)\) and \(P_{rs}\) are symmetric, by inspection, and \(T_{sr}(c)'\) is \(I_n\) with its \((r,s)\)
entry replaced by \(c\), that is, \(T_{rs}(c)\).
\end{itemize}
\item
\begin{itemize}
\item The columns of \(\gamma(A)\) are the transposed rows of \(\gamma(A)'\) (Definition~\ref{r1-def-transpose}).
So \(\gamma(A)'\) is obtained from \(A'\) by the row operation \(\omega\) with the same indices:
\(R_s\leftarrow cR_s\), \(R_s\leftarrow R_s+cR_r\) or \(R_r\leftrightarrow R_s\).
\item By Theorem~\ref{r2-thm-elementary}(a), \(\gamma(A)'=EA'\) with \(E=\omega(I_n)\), which is \(D_s(c)\),
\(T_{sr}(c)\) or \(P_{rs}\). Transposing (Theorem~\ref{r1-thm-transpose}), \(\gamma(A)=(EA')'=AE'\), and \(E'=F\)
by (a).
\end{itemize}
\item By (b), \(B=\gamma_l(\cdots\gamma_1(A))=(\cdots(AF_1)\cdots)F_l=AF_1\cdots F_l\)
(Theorem~\ref{r1-thm-rules}). Each \(F_j\) is invertible (Theorem~\ref{r2-thm-elementary}(b)), and so is \(Q\)
(Theorem~\ref{r2-thm-inverse-props}(b)).\qedhere
\end{enumerate}
\end{proof}

\begin{example}[cf.\ \citeat{tu-clm}{Corollary~11.5(c)}]\label{r2-exm-centering}
Let \(X_0\) be the data matrix of Example~\ref{r1-exm-data}, and \(M=[\one_4\ X_0]\in\R^{4\times3}\). Write
\([\one_4\ CX_0]=MQ\) with a product \(Q\) of elementary matrices, and find \(Q^{-1}\).

\textit{Solution.} Centering the columns of \(X_0\) is a sequence of column operations on \(M\).
\begin{itemize}
\item Since \(\bar x=(1,2)'\), the columns of \(CX_0=X_0-\one\bar x'\) are \(\col_1(X_0)-\one\) and
\(\col_2(X_0)-2\cdot\one\). So \([\one_4\ CX_0]\) arises from \(M\) by \(\col_2\leftarrow\col_2-\col_1\) and then
\(\col_3\leftarrow\col_3-2\col_1\), and Theorem~\ref{r2-thm-column-ops} gives \([\one_4\ CX_0]=MQ\) with
\[
Q=T_{12}(-1)\,T_{13}(-2)=\begin{pmatrix}1&-1&0\\0&1&0\\0&0&1\end{pmatrix}\begin{pmatrix}1&0&-2\\0&1&0\\0&0&1\end{pmatrix}
=\begin{pmatrix}1&-1&-2\\0&1&0\\0&0&1\end{pmatrix}
\]
\item \(Q\) is unit upper triangular, so it is invertible and \(Q^{-1}\) is unit upper triangular
(Theorem~\ref{r2-thm-triangular}); indeed \(Q^{-1}=T_{13}(2)\,T_{12}(1)\) (Theorems~\ref{r2-thm-inverse-props}(b)
and~\ref{r2-thm-elementary}(b)), that is,
\[
Q^{-1}=\begin{pmatrix}1&1&2\\0&1&0\\0&0&1\end{pmatrix}
\]
\item With \(CX_0\) from Example~\ref{r1-exm-data},
\[
M=\begin{pmatrix}1&2&3\\1&0&3\\1&2&1\\1&0&1\end{pmatrix},\qquad
MQ=[\one_4\ CX_0]=\begin{pmatrix}1&1&1\\1&-1&1\\1&1&-1\\1&-1&-1\end{pmatrix}
\]
\item So \(M=[\one_4\ CX_0]\,Q^{-1}\): by Theorem~\ref{r2-thm-column-ops}(c), the column operations of
\(Q^{-1}=T_{13}(2)\,T_{12}(1)\), \(\col_3\leftarrow\col_3+2\col_1\) and then \(\col_2\leftarrow\col_2+\col_1\), undo
the centering; they add \(2\cdot\one\) back to \(\col_2(X_0)-2\cdot\one\) and \(\one\) back to \(\col_1(X_0)-\one\).
\end{itemize}
\end{example}

\subsection*{Exercises}

\begin{exercise}\label{r2-dr-rref}
Find the RREF and the pivot columns of each matrix by Gauss--Jordan elimination
(Definition~\ref{r2-def-gauss-jordan}).
\begin{enumerate}
\item \(A=\begin{pmatrix}1&1&-1&1\\-1&-1&1&0\\2&3&1&-3\end{pmatrix}\)
\item \(B=\begin{pmatrix}1&0&3&0&-1\\1&-1&2&1&-2\\0&-1&-1&1&-1\\-2&3&-3&-2&2\end{pmatrix}\)
\end{enumerate}
\end{exercise}
\begin{solution}
(a) After the first arrow, row \(2\) is \(0\) in column \(2\) and row \(3\) is not, so Step 2 interchanges rows
\(2\) and \(3\); the matrix after the interchange is an echelon form (end of forward elimination).
\[
\begin{gathered}
\begin{pmatrix}1&1&-1&1\\-1&-1&1&0\\2&3&1&-3\end{pmatrix}
\xrightarrow[R_3\leftarrow R_3-2R_1]{R_2\leftarrow R_2+R_1}
\begin{pmatrix}1&1&-1&1\\0&0&0&1\\0&1&3&-5\end{pmatrix}
\\
\xrightarrow{R_2\leftrightarrow R_3}
\begin{pmatrix}1&1&-1&1\\0&1&3&-5\\0&0&0&1\end{pmatrix}
\xrightarrow[R_1\leftarrow R_1-R_3]{R_2\leftarrow R_2+5R_3}
\begin{pmatrix}1&1&-1&0\\0&1&3&0\\0&0&0&1\end{pmatrix}
\\
\xrightarrow{R_1\leftarrow R_1-R_2}
\begin{pmatrix}1&0&-4&0\\0&1&3&0\\0&0&0&1\end{pmatrix}
\end{gathered}
\]
The RREF is the last matrix; its pivot columns are \(1\), \(2\) and \(4\).

(b) After the third arrow, column \(3\) has no pivot, and row \(3\) is zero while row \(4\) is not in column \(4\);
so Step 2 interchanges rows \(3\) and \(4\).
\[
\begin{gathered}
\begin{pmatrix}1&0&3&0&-1\\1&-1&2&1&-2\\0&-1&-1&1&-1\\-2&3&-3&-2&2\end{pmatrix}
\xrightarrow[R_4\leftarrow R_4+2R_1]{R_2\leftarrow R_2-R_1}
\begin{pmatrix}1&0&3&0&-1\\0&-1&-1&1&-1\\0&-1&-1&1&-1\\0&3&3&-2&0\end{pmatrix}
\\
\xrightarrow{R_2\leftarrow-R_2}
\begin{pmatrix}1&0&3&0&-1\\0&1&1&-1&1\\0&-1&-1&1&-1\\0&3&3&-2&0\end{pmatrix}
\xrightarrow[R_4\leftarrow R_4-3R_2]{R_3\leftarrow R_3+R_2}
\begin{pmatrix}1&0&3&0&-1\\0&1&1&-1&1\\0&0&0&0&0\\0&0&0&1&-3\end{pmatrix}
\\
\xrightarrow{R_3\leftrightarrow R_4}
\begin{pmatrix}1&0&3&0&-1\\0&1&1&-1&1\\0&0&0&1&-3\\0&0&0&0&0\end{pmatrix}
\xrightarrow{R_2\leftarrow R_2+R_3}
\begin{pmatrix}1&0&3&0&-1\\0&1&1&0&-2\\0&0&0&1&-3\\0&0&0&0&0\end{pmatrix}
\end{gathered}
\]
The RREF is the last matrix; its pivot columns are \(1\), \(2\) and \(4\).
\end{solution}

\begin{exercise}\label{r2-dr-consistency}
Let \(A=\begin{pmatrix}1&-2&-2\\2&-3&-2\\0&1&2\\1&-1&0\end{pmatrix}\).
\begin{enumerate}
\item For which \(b=(b_1,b_2,b_3,b_4)'\) is \(Ax=b\) consistent?
\item Find all solutions of \(Ax=b\) for \(b=(-2,-1,3,1)'\).
\item Does \(Ax=b\) have a solution for \(b=(1,0,0,0)'\)?
\end{enumerate}
\end{exercise}
\begin{solution}
(a) Row-reduce \([A\mid b]\), keeping \(b_1,\dots,b_4\) as letters:
\[
\begin{gathered}
\left(\begin{array}{ccc|c}1&-2&-2&b_1\\2&-3&-2&b_2\\0&1&2&b_3\\1&-1&0&b_4\end{array}\right)
\xrightarrow[R_4\leftarrow R_4-R_1]{R_2\leftarrow R_2-2R_1}
\left(\begin{array}{ccc|c}1&-2&-2&b_1\\0&1&2&b_2-2b_1\\0&1&2&b_3\\0&1&2&b_4-b_1\end{array}\right)
\\
\xrightarrow[R_4\leftarrow R_4-R_2]{R_3\leftarrow R_3-R_2}
\left(\begin{array}{ccc|c}1&-2&-2&b_1\\0&1&2&b_2-2b_1\\0&0&0&2b_1-b_2+b_3\\0&0&0&b_1-b_2+b_4\end{array}\right)
\xrightarrow{R_1\leftarrow R_1+2R_2}
\left(\begin{array}{ccc|c}1&0&2&2b_2-3b_1\\0&1&2&b_2-2b_1\\0&0&0&2b_1-b_2+b_3\\0&0&0&b_1-b_2+b_4\end{array}\right)
\end{gathered}
\]
The left part of the last matrix is the RREF of \(A\), with \(r=2\). By Proposition~\ref{r2-prop-solutions}(a),
\(Ax=b\) is consistent if and only if
\[
2b_1-b_2+b_3=0\qquad\text{and}\qquad b_1-b_2+b_4=0
\]
(b) Here \(2(-2)-(-1)+3=0\) and \(-2-(-1)+1=0\), so the system is consistent, and the last matrix becomes
\[
\left(\begin{array}{ccc|c}1&0&2&4\\0&1&2&3\\0&0&0&0\\0&0&0&0\end{array}\right)
\]
With \(x_3=t\) free, \(x_1=4-2t\) and \(x_2=3-2t\):
\(x=(4,3,0)'+t\,(-2,-2,1)'\) (\(t\in\R\)). Check: \(A(4,3,0)'=(4-6,\ 8-9,\ 3,\ 4-3)'=(-2,-1,3,1)'\) and
\(A(-2,-2,1)'=(-2+4-2,\ -4+6-2,\ -2+2,\ -2+2)'=0\).

(c) No: \(2b_1-b_2+b_3=2\ne0\).
\end{solution}

\begin{exercise}[cf.\ \citeat{hoffman-kunze}{\S1.4, Exercises~6--9}]\label{r2-ex-which-b}
Let \(A\) be the coefficient matrix of Example~\ref{r2-exm-inconsistent}.
\begin{enumerate}
\item For which \(b=(b_1,b_2,b_3)'\) is \(Ax=b\) consistent?
\item Find all solutions of \(Ax=b\) for \(b=(1,2,4)'\).
\end{enumerate}
\end{exercise}
\begin{solution}
(a) Row-reduce \([A\mid b]\), keeping \(b_1,b_2,b_3\) as letters:
\[
\begin{gathered}
\left(\begin{array}{cccc|c}1&-2&1&2&b_1\\1&-1&-1&1&b_2\\1&1&-5&-1&b_3\end{array}\right)
\xrightarrow[R_3\leftarrow R_3-R_1]{R_2\leftarrow R_2-R_1}
\left(\begin{array}{cccc|c}1&-2&1&2&b_1\\0&1&-2&-1&b_2-b_1\\0&3&-6&-3&b_3-b_1\end{array}\right)
\\
\xrightarrow{R_3\leftarrow R_3-3R_2}
\left(\begin{array}{cccc|c}1&-2&1&2&b_1\\0&1&-2&-1&b_2-b_1\\0&0&0&0&b_3-3b_2+2b_1\end{array}\right)
\xrightarrow{R_1\leftarrow R_1+2R_2}
\left(\begin{array}{cccc|c}1&0&-3&0&2b_2-b_1\\0&1&-2&-1&b_2-b_1\\0&0&0&0&b_3-3b_2+2b_1\end{array}\right)
\end{gathered}
\]
The left part of the last matrix is the RREF of \(A\), with \(r=2\). By
Proposition~\ref{r2-prop-solutions}(a), \(Ax=b\) is consistent if and only if
\(2b_1-3b_2+b_3=0\). (In Example~\ref{r2-exm-inconsistent}, \(2-6+3=-1\ne0\).)

(b) For \(b=(1,2,4)'\), \(2-6+4=0\), so the system is consistent, and the last matrix becomes
\[
\left(\begin{array}{cccc|c}1&0&-3&0&3\\0&1&-2&-1&1\\0&0&0&0&0\end{array}\right)
\]
The free unknowns are \(x_3=s\), \(x_4=t\), and by Proposition~\ref{r2-prop-solutions}(b)
\(x_1=3+3s\), \(x_2=1+2s+t\). So
\[
x=\begin{pmatrix}3\\1\\0\\0\end{pmatrix}+s\begin{pmatrix}3\\2\\1\\0\end{pmatrix}+t\begin{pmatrix}0\\1\\0\\1\end{pmatrix}\qquad(s,t\in\R)
\]
infinitely many solutions (Theorem~\ref{thm:r2-solution-structure}(b)): the particular solution
\((3,1,0,0)'\) plus the solutions of \(Ax=0\) (Theorem~\ref{thm:r2-solution-structure}(a)).
Check: \(A(3,1,0,0)'=(3-2,\ 3-1,\ 3+1)'=(1,2,4)'\).
\end{solution}

\begin{exercise}\label{r2-ex-six-unknowns}
Let
\[
A=\begin{pmatrix}1&2&1&2&1&1\\2&4&3&5&3&3\\1&2&2&3&3&1\\3&6&4&7&3&5\end{pmatrix}\in\R^{4\times6}
\]
\begin{enumerate}
\item Show that \(Ax=b\), \(b=(b_1,b_2,b_3,b_4)'\), is consistent if and only if \(2b_2=b_3+b_4\).
\item Find all solutions of \(Ax=b\) for \(b=(2,4,3,5)'\).
\end{enumerate}
\end{exercise}
\begin{solution}
(a) Forward elimination of \([A\mid b]\), keeping \(b_1,\dots,b_4\) as letters:
\[
\begin{gathered}
\left(\begin{array}{cccccc|c}1&2&1&2&1&1&b_1\\2&4&3&5&3&3&b_2\\1&2&2&3&3&1&b_3\\3&6&4&7&3&5&b_4\end{array}\right)
\xrightarrow[R_4\leftarrow R_4-3R_1]{R_2\leftarrow R_2-2R_1,\ R_3\leftarrow R_3-R_1}
\left(\begin{array}{cccccc|c}1&2&1&2&1&1&b_1\\0&0&1&1&1&1&b_2-2b_1\\0&0&1&1&2&0&b_3-b_1\\0&0&1&1&0&2&b_4-3b_1\end{array}\right)
\\
\xrightarrow[R_4\leftarrow R_4-R_2]{R_3\leftarrow R_3-R_2}
\left(\begin{array}{cccccc|c}1&2&1&2&1&1&b_1\\0&0&1&1&1&1&b_2-2b_1\\0&0&0&0&1&-1&b_1-b_2+b_3\\0&0&0&0&-1&1&b_4-b_1-b_2\end{array}\right)
\\
\xrightarrow{R_4\leftarrow R_4+R_3}
\left(\begin{array}{cccccc|c}1&2&1&2&1&1&b_1\\0&0&1&1&1&1&b_2-2b_1\\0&0&0&0&1&-1&b_1-b_2+b_3\\0&0&0&0&0&0&-2b_2+b_3+b_4\end{array}\right)
\end{gathered}
\]
The left part of the last matrix is in echelon form, with pivots in columns \(1,3,5\) of \(A\), and the last row
reads \(0=-2b_2+b_3+b_4\).
By Theorem~\ref{thm:r2-consistency}, \(Ax=b\) is consistent if and only if \(-2b_2+b_3+b_4=0\), that is,
\(2b_2=b_3+b_4\).

(b) For \(b=(2,4,3,5)'\), \(2\cdot4=3+5\), and the last matrix has rows
\((1,2,1,2,1,1\mid2)\), \((0,0,1,1,1,1\mid0)\), \((0,0,0,0,1,-1\mid1)\) and a zero row.
Backward elimination:
\[
\xrightarrow[R_1\leftarrow R_1-R_3]{R_2\leftarrow R_2-R_3}
\left(\begin{array}{cccccc|c}1&2&1&2&0&2&1\\0&0&1&1&0&2&-1\\0&0&0&0&1&-1&1\\0&0&0&0&0&0&0\end{array}\right)
\xrightarrow{R_1\leftarrow R_1-R_2}
\left(\begin{array}{cccccc|c}1&2&0&1&0&0&2\\0&0&1&1&0&2&-1\\0&0&0&0&1&-1&1\\0&0&0&0&0&0&0\end{array}\right)
\]
The free unknowns are \(x_2=\lambda\), \(x_4=\mu\), \(x_6=\nu\), and
Proposition~\ref{r2-prop-solutions}(b) gives \(x_1=2-2\lambda-\mu\), \(x_3=-1-\mu-2\nu\),
\(x_5=1+\nu\):
\[
x=\begin{pmatrix}2\\0\\-1\\0\\1\\0\end{pmatrix}
+\lambda\begin{pmatrix}-2\\1\\0\\0\\0\\0\end{pmatrix}
+\mu\begin{pmatrix}-1\\0\\-1\\1\\0\\0\end{pmatrix}
+\nu\begin{pmatrix}0\\0\\-2\\0\\1\\1\end{pmatrix}
\qquad(\lambda,\mu,\nu\in\R)
\]
the particular solution \((2,0,-1,0,1,0)'\) plus the solutions of \(Ax=0\)
(Theorem~\ref{thm:r2-solution-structure}(a)). Check:
\(A(2,0,-1,0,1,0)'=2\col_1(A)-\col_3(A)+\col_5(A)=(2-1+1,\ 4-3+3,\ 2-2+3,\ 6-4+3)'=(2,4,3,5)'\).
\end{solution}

\begin{exercise}\label{r2-dr-homogeneous}
\begin{enumerate}
\item Find all solutions of \(Ax=0\) for
\(A=\begin{pmatrix}1&1&-1&0&0\\2&1&-3&0&3\\2&1&-3&1&1\end{pmatrix}\).
\item Does \(Bx=0\) have a nontrivial solution, for
\(B=\begin{pmatrix}1&-2&-1&-1\\-1&3&3&1\\-1&2&1&1\\-1&2&1&0\end{pmatrix}\)? Find all solutions.
\end{enumerate}
\end{exercise}
\begin{solution}
(a) There are fewer equations than unknowns, so nontrivial solutions exist (Theorem~\ref{r2-thm-homogeneous}(a)).
Row-reduce \(A\):
\[
\begin{gathered}
\begin{pmatrix}1&1&-1&0&0\\2&1&-3&0&3\\2&1&-3&1&1\end{pmatrix}
\xrightarrow[R_3\leftarrow R_3-2R_1]{R_2\leftarrow R_2-2R_1}
\begin{pmatrix}1&1&-1&0&0\\0&-1&-1&0&3\\0&-1&-1&1&1\end{pmatrix}
\\
\xrightarrow{R_2\leftarrow-R_2}
\begin{pmatrix}1&1&-1&0&0\\0&1&1&0&-3\\0&-1&-1&1&1\end{pmatrix}
\xrightarrow{R_3\leftarrow R_3+R_2}
\begin{pmatrix}1&1&-1&0&0\\0&1&1&0&-3\\0&0&0&1&-2\end{pmatrix}
\\
\xrightarrow{R_1\leftarrow R_1-R_2}
\begin{pmatrix}1&0&-2&0&3\\0&1&1&0&-3\\0&0&0&1&-2\end{pmatrix}
\end{gathered}
\]
The pivot columns are \(1,2,4\); the free unknowns are \(x_3=s\), \(x_5=t\), and the rows give \(x_1=2s-3t\),
\(x_2=-s+3t\), \(x_4=2t\):
\[
x=s\begin{pmatrix}2\\-1\\1\\0\\0\end{pmatrix}+t\begin{pmatrix}-3\\3\\0\\2\\1\end{pmatrix}\qquad(s,t\in\R)
\]
Check: \(A(2,-1,1,0,0)'=(2-1-1,\ 4-1-3,\ 4-1-3)'=0\) and \(A(-3,3,0,2,1)'=(-3+3,\ -6+3+3,\ -6+3+2+1)'=0\).

(b) Row-reduce \(B\); after the first arrow, column \(3\) has no pivot and row \(3\) is zero, so Step 2
interchanges rows \(3\) and \(4\):
\[
\begin{gathered}
\begin{pmatrix}1&-2&-1&-1\\-1&3&3&1\\-1&2&1&1\\-1&2&1&0\end{pmatrix}
\xrightarrow[R_4\leftarrow R_4+R_1]{R_2\leftarrow R_2+R_1,\ R_3\leftarrow R_3+R_1}
\begin{pmatrix}1&-2&-1&-1\\0&1&2&0\\0&0&0&0\\0&0&0&-1\end{pmatrix}
\\
\xrightarrow{R_3\leftrightarrow R_4}
\begin{pmatrix}1&-2&-1&-1\\0&1&2&0\\0&0&0&-1\\0&0&0&0\end{pmatrix}
\xrightarrow{R_3\leftarrow-R_3}
\begin{pmatrix}1&-2&-1&-1\\0&1&2&0\\0&0&0&1\\0&0&0&0\end{pmatrix}
\\
\xrightarrow{R_1\leftarrow R_1+R_3}
\begin{pmatrix}1&-2&-1&0\\0&1&2&0\\0&0&0&1\\0&0&0&0\end{pmatrix}
\xrightarrow{R_1\leftarrow R_1+2R_2}
\begin{pmatrix}1&0&3&0\\0&1&2&0\\0&0&0&1\\0&0&0&0\end{pmatrix}
\end{gathered}
\]
The RREF has pivot columns \(1,2,4\), so \(x_3=t\) is free: yes, there are nontrivial solutions, and all
solutions are
\[
x=t\,(-3,\ -2,\ 1,\ 0)'\qquad(t\in\R)
\]
Check: \(B(-3,-2,1,0)'=(-3+4-1,\ 3-6+3,\ 3-4+1,\ 3-4+1)'=0\). In particular \(B\) is not invertible
(Theorem~\ref{r2-thm-invertible}).
\end{solution}

\begin{exercise}[cf.\ \citeat{abadir-magnus}{Exercise~6.37}]\label{r2-ex-homogeneous-two}
Find the solutions of the equations
\[
\begin{aligned}
x_1+x_2+x_3+x_4+2x_5&=0\\
2x_1+3x_2+x_3+3x_4+5x_5&=0\\
x_1+2x_3+x_4+2x_5&=0
\end{aligned}
\]
\end{exercise}
\begin{solution}
Row-reduce the coefficient matrix:
\[
\begin{gathered}
\begin{pmatrix}1&1&1&1&2\\2&3&1&3&5\\1&0&2&1&2\end{pmatrix}
\xrightarrow[R_3\leftarrow R_3-R_1]{R_2\leftarrow R_2-2R_1}
\begin{pmatrix}1&1&1&1&2\\0&1&-1&1&1\\0&-1&1&0&0\end{pmatrix}
\xrightarrow{R_3\leftarrow R_3+R_2}
\begin{pmatrix}1&1&1&1&2\\0&1&-1&1&1\\0&0&0&1&1\end{pmatrix}
\\
\xrightarrow[R_1\leftarrow R_1-R_3]{R_2\leftarrow R_2-R_3}
\begin{pmatrix}1&1&1&0&1\\0&1&-1&0&0\\0&0&0&1&1\end{pmatrix}
\xrightarrow{R_1\leftarrow R_1-R_2}
\begin{pmatrix}1&0&2&0&1\\0&1&-1&0&0\\0&0&0&1&1\end{pmatrix}
\end{gathered}
\]
The pivot columns are \(1,2,4\), so \(x_3\) and \(x_5\) are free unknowns
(Proposition~\ref{r2-prop-solutions}). With \(x_3=\lambda\) and \(x_5=\mu\), the rows of the RREF
give \(x_1=-2\lambda-\mu\), \(x_2=\lambda\), \(x_4=-\mu\). All solutions are
\[
x=\lambda(-2,1,1,0,0)'+\mu(-1,0,0,-1,1)'\qquad(\lambda,\mu\in\R)
\]
Check: \((-2,1,1,0,0)'\) gives \(-2+1+1=0\), \(-4+3+1=0\), \(-2+2=0\), and \((-1,0,0,-1,1)'\)
gives \(-1-1+2=0\), \(-2-3+5=0\), \(-1-1+2=0\). Three equations in five unknowns always have
nontrivial solutions (Theorem~\ref{r2-thm-homogeneous}(a)); here there are two free unknowns.
\end{solution}

\begin{exercise}[cf.\ \citeat{abadir-magnus}{Exercise~6.40}]\label{r2-ex-homogeneous-four}
Solve the homogeneous system
\[
\begin{aligned}
x_1-x_2+2x_3+x_4&=0\\
x_1-x_3+3x_4&=0\\
3x_1-2x_2+3x_3+5x_4&=0\\
x_1+x_2-4x_3+5x_4&=0
\end{aligned}
\]
\end{exercise}
\begin{solution}
Transform the coefficient matrix to echelon form:
\[
\begin{pmatrix}1&-1&2&1\\1&0&-1&3\\3&-2&3&5\\1&1&-4&5\end{pmatrix}
\xrightarrow[R_4\leftarrow R_4-R_1]{R_2\leftarrow R_2-R_1,\ R_3\leftarrow R_3-3R_1}
\begin{pmatrix}1&-1&2&1\\0&1&-3&2\\0&1&-3&2\\0&2&-6&4\end{pmatrix}
\xrightarrow[R_4\leftarrow R_4-2R_2]{R_3\leftarrow R_3-R_2}
\begin{pmatrix}1&-1&2&1\\0&1&-3&2\\0&0&0&0\\0&0&0&0\end{pmatrix}
\]
By Theorem~\ref{r2-thm-rowops}, the system has the same solutions as the two equations
\(x_1-x_2+2x_3+x_4=0\) and \(x_2-3x_3+2x_4=0\). For any values of \(x_3\) and \(x_4\), the
second equation determines \(x_2\) and then the first determines \(x_1\); so \(x_3\) and \(x_4\)
are free. Let \(x_3=\lambda\) and \(x_4=\mu\). Then \(x_2=3x_3-2x_4=3\lambda-2\mu\) and
\(x_1=x_2-2x_3-x_4=\lambda-3\mu\). Hence
\[
x=\lambda(1,3,1,0)'+\mu(-3,-2,0,1)'\qquad(\lambda,\mu\in\R)
\]
(One more arrow, \(R_1\leftarrow R_1+R_2\), gives the RREF, with rows \((1,0,-1,3)\) and
\((0,1,-3,2)\), from which the same solutions are read off.)
\end{solution}

\begin{exercise}[\citeat{abadir-magnus}{Exercise~6.1}]\label{r2-ex-elementary}
Let \(A=\begin{pmatrix}1&2\\3&4\\5&6\end{pmatrix}\). Find an elementary matrix \(E\) with
\(B=EA\) when:
(a) \(B\) is obtained from \(A\) by interchanging the first and third rows;
(b) \(B\) is obtained from \(A\) by multiplying the second row by \(7\);
(c) \(B\) is obtained from \(A\) by adding five times the second row to the third row.
\end{exercise}
\begin{solution}
Perform each operation on \(I_3\) (Theorem~\ref{r2-thm-elementary}(a)):
\[
P_{13}=\begin{pmatrix}0&0&1\\0&1&0\\1&0&0\end{pmatrix},\quad
D_2(7)=\begin{pmatrix}1&0&0\\0&7&0\\0&0&1\end{pmatrix},\quad
T_{32}(5)=\begin{pmatrix}1&0&0\\0&1&0\\0&5&1\end{pmatrix}
\]
and
\[
P_{13}A=\begin{pmatrix}5&6\\3&4\\1&2\end{pmatrix},\quad
D_2(7)A=\begin{pmatrix}1&2\\21&28\\5&6\end{pmatrix},\quad
T_{32}(5)A=\begin{pmatrix}1&2\\3&4\\20&26\end{pmatrix}
\]
\end{solution}

\begin{exercise}[cf.\ \citeat{abadir-magnus}{Exercise~6.8}]\label{r2-ex-order}
Let \(A=\begin{pmatrix}1&2&3&4\\0&1&2&3\\1&0&1&0\\2&1&0&1\end{pmatrix}\).
\begin{enumerate}
\item Find the matrix \(E\) such that \(EA\) is obtained from \(A\) by first interchanging rows
\(1\) and \(2\) and then multiplying row \(2\) by \(2\). Does the order of the two operations
matter?
\item Find the matrix \(E\) such that \(EA\) is obtained from \(A\) by first interchanging rows
\(1\) and \(2\) and then adding \(3\) times row \(1\) to row \(2\). Does the order matter?
\item Do elementary matrices commute?
\end{enumerate}
\end{exercise}
\begin{solution}
By Theorem~\ref{r2-thm-elementary}, the operation done first is the factor next to \(A\).

(a) \(E=D_2(2)P_{12}=\left(\begin{smallmatrix}0&1&0&0\\2&0&0&0\\0&0&1&0\\0&0&0&1\end{smallmatrix}\right)\),
while \(P_{12}D_2(2)=\left(\begin{smallmatrix}0&2&0&0\\1&0&0&0\\0&0&1&0\\0&0&0&1\end{smallmatrix}\right)\), and
\[
D_2(2)P_{12}A=\begin{pmatrix}0&1&2&3\\2&4&6&8\\1&0&1&0\\2&1&0&1\end{pmatrix}
\ne\begin{pmatrix}0&2&4&6\\1&2&3&4\\1&0&1&0\\2&1&0&1\end{pmatrix}=P_{12}D_2(2)A
\]
The order matters.

(b) \(E=T_{21}(3)P_{12}=\left(\begin{smallmatrix}0&1&0&0\\1&3&0&0\\0&0&1&0\\0&0&0&1\end{smallmatrix}\right)\),
while \(P_{12}T_{21}(3)=\left(\begin{smallmatrix}3&1&0&0\\1&0&0&0\\0&0&1&0\\0&0&0&1\end{smallmatrix}\right)\), and
\[
T_{21}(3)P_{12}A=\begin{pmatrix}0&1&2&3\\1&5&9&13\\1&0&1&0\\2&1&0&1\end{pmatrix}
\ne\begin{pmatrix}3&7&11&15\\1&2&3&4\\1&0&1&0\\2&1&0&1\end{pmatrix}=P_{12}T_{21}(3)A
\]
The order matters.

(c) No: (a) and (b) give pairs of elementary matrices that do not commute.
\end{solution}

\begin{exercise}[\citeat{abadir-magnus}{Exercise~6.14}]\label{r2-ex-echelon}
Reduce the matrix
\(A=\begin{pmatrix}0&0&1&2&3\\2&4&-2&0&4\\2&4&-1&2&7\end{pmatrix}\)
to echelon form by elementary row operations, and record the elementary matrices used.
\end{exercise}
\begin{solution}
Follow Definition~\ref{r2-def-gauss-jordan}, postponing the scaling (Step 3) to the end.
Column \(1\) is the first pivot column; since \(a_{11}=0\), interchange rows \(1\) and \(2\);
then clear column \(1\) below the pivot, and then column \(3\), the next pivot column:
\[
A\xrightarrow{R_1\leftrightarrow R_2}
\begin{pmatrix}2&4&-2&0&4\\0&0&1&2&3\\2&4&-1&2&7\end{pmatrix}
\xrightarrow{R_3\leftarrow R_3-R_1}
\begin{pmatrix}2&4&-2&0&4\\0&0&1&2&3\\0&0&1&2&3\end{pmatrix}
\xrightarrow{R_3\leftarrow R_3-R_2}
\begin{pmatrix}2&4&-2&0&4\\0&0&1&2&3\\0&0&0&0&0\end{pmatrix}
\]
Finally divide row \(1\) by \(2\):
\[
H=\begin{pmatrix}1&2&-1&0&2\\0&0&1&2&3\\0&0&0&0&0\end{pmatrix}
=D_1(\tfrac12)\,T_{32}(-1)\,T_{31}(-1)\,P_{12}\,A
\]
(Theorem~\ref{r2-thm-elementary}(c)). \(H\) is in echelon form with pivot columns \(1\) and
\(3\). One more operation, \(R_1\leftarrow R_1+R_2\), clears the entry above the second pivot
and gives the RREF
\(\left(\begin{smallmatrix}1&2&0&2&5\\0&0&1&2&3\\0&0&0&0&0\end{smallmatrix}\right)\).
\end{solution}

\begin{exercise}\label{r2-dr-inverse3}
Row-reduce \([M\mid I_4]\) to find \(M^{-1}\), or show that \(M\) is not invertible, for
\begin{enumerate}
\item \(M=\begin{pmatrix}1&1&0&-1\\1&1&-1&0\\-2&-1&3&0\\-1&-1&2&0\end{pmatrix}\)
\item \(M=\begin{pmatrix}1&-1&0&0\\-3&2&2&-1\\-2&2&1&-1\\-1&1&1&0\end{pmatrix}\)
\item \(M=\begin{pmatrix}1&2&-1&1\\2&3&1&-1\\1&1&2&-2\\1&2&0&-1\end{pmatrix}\)
\end{enumerate}
\end{exercise}
\begin{solution}
(a) After the first arrow, row \(2\) is \(0\) in column \(2\), so Step 2 interchanges rows \(2\) and \(3\):
\[
\begin{gathered}
\left(\begin{array}{cccc|cccc}1&1&0&-1&1&0&0&0\\1&1&-1&0&0&1&0&0\\-2&-1&3&0&0&0&1&0\\-1&-1&2&0&0&0&0&1\end{array}\right)
\xrightarrow[R_4\leftarrow R_4+R_1]{R_2\leftarrow R_2-R_1,\ R_3\leftarrow R_3+2R_1}
\left(\begin{array}{cccc|cccc}1&1&0&-1&1&0&0&0\\0&0&-1&1&-1&1&0&0\\0&1&3&-2&2&0&1&0\\0&0&2&-1&1&0&0&1\end{array}\right)
\\
\xrightarrow{R_2\leftrightarrow R_3}
\left(\begin{array}{cccc|cccc}1&1&0&-1&1&0&0&0\\0&1&3&-2&2&0&1&0\\0&0&-1&1&-1&1&0&0\\0&0&2&-1&1&0&0&1\end{array}\right)
\xrightarrow{R_3\leftarrow-R_3}
\left(\begin{array}{cccc|cccc}1&1&0&-1&1&0&0&0\\0&1&3&-2&2&0&1&0\\0&0&1&-1&1&-1&0&0\\0&0&2&-1&1&0&0&1\end{array}\right)
\\
\xrightarrow{R_4\leftarrow R_4-2R_3}
\left(\begin{array}{cccc|cccc}1&1&0&-1&1&0&0&0\\0&1&3&-2&2&0&1&0\\0&0&1&-1&1&-1&0&0\\0&0&0&1&-1&2&0&1\end{array}\right)
\xrightarrow[R_3\leftarrow R_3+R_4]{R_1\leftarrow R_1+R_4,\ R_2\leftarrow R_2+2R_4}
\left(\begin{array}{cccc|cccc}1&1&0&0&0&2&0&1\\0&1&3&0&0&4&1&2\\0&0&1&0&0&1&0&1\\0&0&0&1&-1&2&0&1\end{array}\right)
\\
\xrightarrow{R_2\leftarrow R_2-3R_3}
\left(\begin{array}{cccc|cccc}1&1&0&0&0&2&0&1\\0&1&0&0&0&1&1&-1\\0&0&1&0&0&1&0&1\\0&0&0&1&-1&2&0&1\end{array}\right)
\xrightarrow{R_1\leftarrow R_1-R_2}
\left(\begin{array}{cccc|cccc}1&0&0&0&0&1&-1&2\\0&1&0&0&0&1&1&-1\\0&0&1&0&0&1&0&1\\0&0&0&1&-1&2&0&1\end{array}\right)
\end{gathered}
\]
So
\[
M^{-1}=\begin{pmatrix}0&1&-1&2\\0&1&1&-1\\0&1&0&1\\-1&2&0&1\end{pmatrix}
\]
Check: \(\row_1(M)M^{-1}=(0+0+1,\ 1+1-2,\ -1+1,\ 2-1-1)=(1,0,0,0)\), and similarly for rows \(2,3,4\).

(b)
\[
\begin{gathered}
\left(\begin{array}{cccc|cccc}1&-1&0&0&1&0&0&0\\-3&2&2&-1&0&1&0&0\\-2&2&1&-1&0&0&1&0\\-1&1&1&0&0&0&0&1\end{array}\right)
\xrightarrow[R_4\leftarrow R_4+R_1]{R_2\leftarrow R_2+3R_1,\ R_3\leftarrow R_3+2R_1}
\left(\begin{array}{cccc|cccc}1&-1&0&0&1&0&0&0\\0&-1&2&-1&3&1&0&0\\0&0&1&-1&2&0&1&0\\0&0&1&0&1&0&0&1\end{array}\right)
\\
\xrightarrow{R_2\leftarrow-R_2}
\left(\begin{array}{cccc|cccc}1&-1&0&0&1&0&0&0\\0&1&-2&1&-3&-1&0&0\\0&0&1&-1&2&0&1&0\\0&0&1&0&1&0&0&1\end{array}\right)
\xrightarrow{R_4\leftarrow R_4-R_3}
\left(\begin{array}{cccc|cccc}1&-1&0&0&1&0&0&0\\0&1&-2&1&-3&-1&0&0\\0&0&1&-1&2&0&1&0\\0&0&0&1&-1&0&-1&1\end{array}\right)
\\
\xrightarrow[R_3\leftarrow R_3+R_4]{R_2\leftarrow R_2-R_4}
\left(\begin{array}{cccc|cccc}1&-1&0&0&1&0&0&0\\0&1&-2&0&-2&-1&1&-1\\0&0&1&0&1&0&0&1\\0&0&0&1&-1&0&-1&1\end{array}\right)
\xrightarrow{R_2\leftarrow R_2+2R_3}
\left(\begin{array}{cccc|cccc}1&-1&0&0&1&0&0&0\\0&1&0&0&0&-1&1&1\\0&0&1&0&1&0&0&1\\0&0&0&1&-1&0&-1&1\end{array}\right)
\\
\xrightarrow{R_1\leftarrow R_1+R_2}
\left(\begin{array}{cccc|cccc}1&0&0&0&1&-1&1&1\\0&1&0&0&0&-1&1&1\\0&0&1&0&1&0&0&1\\0&0&0&1&-1&0&-1&1\end{array}\right)
\end{gathered}
\]
So
\[
M^{-1}=\begin{pmatrix}1&-1&1&1\\0&-1&1&1\\1&0&0&1\\-1&0&-1&1\end{pmatrix}
\]
Check: \(\row_2(M)M^{-1}=(-3+2+1,\ 3-2,\ -3+2+1,\ -3+2+2-1)=(0,1,0,0)\), and similarly for rows \(1,3,4\).

(c)
\[
\begin{gathered}
\left(\begin{array}{cccc|cccc}1&2&-1&1&1&0&0&0\\2&3&1&-1&0&1&0&0\\1&1&2&-2&0&0&1&0\\1&2&0&-1&0&0&0&1\end{array}\right)
\xrightarrow[R_4\leftarrow R_4-R_1]{R_2\leftarrow R_2-2R_1,\ R_3\leftarrow R_3-R_1}
\left(\begin{array}{cccc|cccc}1&2&-1&1&1&0&0&0\\0&-1&3&-3&-2&1&0&0\\0&-1&3&-3&-1&0&1&0\\0&0&1&-2&-1&0&0&1\end{array}\right)
\\
\xrightarrow{R_2\leftarrow-R_2}
\left(\begin{array}{cccc|cccc}1&2&-1&1&1&0&0&0\\0&1&-3&3&2&-1&0&0\\0&-1&3&-3&-1&0&1&0\\0&0&1&-2&-1&0&0&1\end{array}\right)
\xrightarrow{R_3\leftarrow R_3+R_2}
\left(\begin{array}{cccc|cccc}1&2&-1&1&1&0&0&0\\0&1&-3&3&2&-1&0&0\\0&0&0&0&1&-1&1&0\\0&0&1&-2&-1&0&0&1\end{array}\right)
\end{gathered}
\]
A zero row has appeared in the left half (row \(3\)), so \(M\) is not invertible
(Corollary~\ref{r2-cor-computing-inverse}(b)). Indeed, \(Mx=0\) has the same solutions as the system of the left half
(Theorem~\ref{r2-thm-rowops}(b)): with \(x_4=t\), \(x_3-2x_4=0\) gives \(x_3=2t\), \(x_2-3x_3+3x_4=0\) gives
\(x_2=3t\), and then \(x_1=-2x_2+x_3-x_4=-5t\). So \(Mx=0\) has the nontrivial solution \(x=(-5,3,2,1)'\):
\(M(-5,3,2,1)'=(-5+6-2+1,\ -10+9+2-1,\ -5+3+4-2,\ -5+6+0-1)'=0\).
\end{solution}

\begin{exercise}[\citeat{abadir-magnus}{Exercise~6.23}]\label{r2-ex-inverse}
Use row reduction of \([A\mid I_3]\) to find the inverse of
\(A=\begin{pmatrix}0&1&1\\1&1&1\\3&4&3\end{pmatrix}\).
\end{exercise}
\begin{solution}
\[
\begin{gathered}
\left(\begin{array}{ccc|ccc}0&1&1&1&0&0\\1&1&1&0&1&0\\3&4&3&0&0&1\end{array}\right)
\xrightarrow{R_1\leftrightarrow R_2}
\left(\begin{array}{ccc|ccc}1&1&1&0&1&0\\0&1&1&1&0&0\\3&4&3&0&0&1\end{array}\right)
\xrightarrow{R_3\leftarrow R_3-3R_1}
\left(\begin{array}{ccc|ccc}1&1&1&0&1&0\\0&1&1&1&0&0\\0&1&0&0&-3&1\end{array}\right)
\\
\xrightarrow{R_3\leftarrow R_3-R_2}
\left(\begin{array}{ccc|ccc}1&1&1&0&1&0\\0&1&1&1&0&0\\0&0&-1&-1&-3&1\end{array}\right)
\xrightarrow{R_3\leftarrow -R_3}
\left(\begin{array}{ccc|ccc}1&1&1&0&1&0\\0&1&1&1&0&0\\0&0&1&1&3&-1\end{array}\right)
\\
\xrightarrow[R_2\leftarrow R_2-R_3]{R_1\leftarrow R_1-R_3}
\left(\begin{array}{ccc|ccc}1&1&0&-1&-2&1\\0&1&0&0&-3&1\\0&0&1&1&3&-1\end{array}\right)
\xrightarrow{R_1\leftarrow R_1-R_2}
\left(\begin{array}{ccc|ccc}1&0&0&-1&1&0\\0&1&0&0&-3&1\\0&0&1&1&3&-1\end{array}\right)
\end{gathered}
\]
Hence
\(A^{-1}=\left(\begin{smallmatrix}-1&1&0\\0&-3&1\\1&3&-1\end{smallmatrix}\right)\). Check:
\(\row_3(A)A^{-1}=(3,4,3)A^{-1}=(-3+3,\ 3-12+9,\ 4-3)=(0,0,1)\), and similarly for the other
rows.
\end{solution}

\begin{exercise}[cf.\ \citeat{abadir-magnus}{Exercise~4.19}]\label{r2-ex-two-inverses}
Do \(A=\begin{pmatrix}1&2&1&1\\2&4&3&1\\1&3&2&2\\3&7&5&3\end{pmatrix}\) and
\(B=\begin{pmatrix}1&1&0&0\\1&2&1&0\\0&1&2&1\\0&0&1&2\end{pmatrix}\) have inverses? Find the
inverse when it exists.
\end{exercise}
\begin{solution}
(\(A\)) Row-reduce \(A\); after the first arrow the entry of row \(2\) in column \(2\) is \(0\),
so rows \(2\) and \(3\) are interchanged:
\[
\begin{gathered}
A\xrightarrow[R_4\leftarrow R_4-3R_1]{R_2\leftarrow R_2-2R_1,\ R_3\leftarrow R_3-R_1}
\begin{pmatrix}1&2&1&1\\0&0&1&-1\\0&1&1&1\\0&1&2&0\end{pmatrix}
\xrightarrow{R_2\leftrightarrow R_3}
\begin{pmatrix}1&2&1&1\\0&1&1&1\\0&0&1&-1\\0&1&2&0\end{pmatrix}
\\
\xrightarrow{R_4\leftarrow R_4-R_2}
\begin{pmatrix}1&2&1&1\\0&1&1&1\\0&0&1&-1\\0&0&1&-1\end{pmatrix}
\xrightarrow{R_4\leftarrow R_4-R_3}
\begin{pmatrix}1&2&1&1\\0&1&1&1\\0&0&1&-1\\0&0&0&0\end{pmatrix}
\end{gathered}
\]
A zero row appears, so \(A\) has no inverse (Corollary~\ref{r2-cor-computing-inverse}(b): the same operations
give this left half in \([A\mid I_4]\)). (Indeed row \(4\) of \(A\) is the sum of rows \(2\) and \(3\).)

(\(B\)) Row-reduce \([B\mid I_4]\). Every pivot is \(1\); forward elimination subtracts each row from the next one,
and backward elimination each row from the one above it:
\[
\begin{gathered}
\left(\begin{array}{cccc|cccc}1&1&0&0&1&0&0&0\\1&2&1&0&0&1&0&0\\0&1&2&1&0&0&1&0\\0&0&1&2&0&0&0&1\end{array}\right)
\xrightarrow{R_2\leftarrow R_2-R_1}
\left(\begin{array}{cccc|cccc}1&1&0&0&1&0&0&0\\0&1&1&0&-1&1&0&0\\0&1&2&1&0&0&1&0\\0&0&1&2&0&0&0&1\end{array}\right)
\\
\xrightarrow{R_3\leftarrow R_3-R_2}
\left(\begin{array}{cccc|cccc}1&1&0&0&1&0&0&0\\0&1&1&0&-1&1&0&0\\0&0&1&1&1&-1&1&0\\0&0&1&2&0&0&0&1\end{array}\right)
\xrightarrow{R_4\leftarrow R_4-R_3}
\left(\begin{array}{cccc|cccc}1&1&0&0&1&0&0&0\\0&1&1&0&-1&1&0&0\\0&0&1&1&1&-1&1&0\\0&0&0&1&-1&1&-1&1\end{array}\right)
\\
\xrightarrow{R_3\leftarrow R_3-R_4}
\left(\begin{array}{cccc|cccc}1&1&0&0&1&0&0&0\\0&1&1&0&-1&1&0&0\\0&0&1&0&2&-2&2&-1\\0&0&0&1&-1&1&-1&1\end{array}\right)
\xrightarrow{R_2\leftarrow R_2-R_3}
\left(\begin{array}{cccc|cccc}1&1&0&0&1&0&0&0\\0&1&0&0&-3&3&-2&1\\0&0&1&0&2&-2&2&-1\\0&0&0&1&-1&1&-1&1\end{array}\right)
\\
\xrightarrow{R_1\leftarrow R_1-R_2}
\left(\begin{array}{cccc|cccc}1&0&0&0&4&-3&2&-1\\0&1&0&0&-3&3&-2&1\\0&0&1&0&2&-2&2&-1\\0&0&0&1&-1&1&-1&1\end{array}\right)
\end{gathered}
\]
The left block is \(I_4\), so \(B\) is invertible and
\[
B^{-1}=\begin{pmatrix}4&-3&2&-1\\-3&3&-2&1\\2&-2&2&-1\\-1&1&-1&1\end{pmatrix}
\]
(Corollary~\ref{r2-cor-computing-inverse}(a)). Check: the second row of \(BB^{-1}\) is
\((4-6+2,\ -3+6-2,\ 2-4+2,\ -1+2-1)=(0,1,0,0)\), and similarly for the other rows.
\end{solution}

\begin{exercise}[cf.\ \citeat{bapat}{Ch.~2, Exercise~15}]\label{r2-ex-inverses-b}
Use Proposition~\ref{r2-prop-inverse-2x2} to decide which of
\(\left(\begin{smallmatrix}5&3\\3&2\end{smallmatrix}\right)\),
\(\left(\begin{smallmatrix}4&6\\6&9\end{smallmatrix}\right)\) and
\(\left(\begin{smallmatrix}7&-3\\-5&2\end{smallmatrix}\right)\) are invertible, and find
the inverses that exist. Then find the inverse of
\[
\begin{pmatrix}0&1&1&1\\1&1&1&1\\0&1&1&2\\0&0&1&1\end{pmatrix}
\]
\end{exercise}
\begin{solution}
(\(2\times2\)) By Proposition~\ref{r2-prop-inverse-2x2}:
\begin{itemize}
\item \(5\cdot2-3\cdot3=1\), so
\(\left(\begin{smallmatrix}5&3\\3&2\end{smallmatrix}\right)^{-1}
=\left(\begin{smallmatrix}2&-3\\-3&5\end{smallmatrix}\right)\).
\item \(4\cdot9-6\cdot6=0\), so
\(\left(\begin{smallmatrix}4&6\\6&9\end{smallmatrix}\right)\) is not invertible.
\item \(7\cdot2-(-3)(-5)=-1\), so
\(\left(\begin{smallmatrix}7&-3\\-5&2\end{smallmatrix}\right)^{-1}
=-\left(\begin{smallmatrix}2&3\\5&7\end{smallmatrix}\right)
=\left(\begin{smallmatrix}-2&-3\\-5&-7\end{smallmatrix}\right)\).
\end{itemize}
Check: \(\left(\begin{smallmatrix}7&-3\\-5&2\end{smallmatrix}\right)
\left(\begin{smallmatrix}-2&-3\\-5&-7\end{smallmatrix}\right)
=\left(\begin{smallmatrix}-14+15&-21+21\\10-10&15-14\end{smallmatrix}\right)=I_2\).

(\(4\times4\)) Call the \(4\times4\) matrix \(K\) and row-reduce \([K\mid I_4]\). The entry in
position \((1,1)\) is \(0\), so Step 2 interchanges rows \(1\) and \(2\); after clearing column
\(2\), the entry in position \((3,3)\) is \(0\), so rows \(3\) and \(4\) are interchanged:
\[
\begin{gathered}
\left(\begin{array}{cccc|cccc}0&1&1&1&1&0&0&0\\1&1&1&1&0&1&0&0\\0&1&1&2&0&0&1&0\\0&0&1&1&0&0&0&1\end{array}\right)
\xrightarrow{R_1\leftrightarrow R_2}
\left(\begin{array}{cccc|cccc}1&1&1&1&0&1&0&0\\0&1&1&1&1&0&0&0\\0&1&1&2&0&0&1&0\\0&0&1&1&0&0&0&1\end{array}\right)
\\
\xrightarrow{R_3\leftarrow R_3-R_2}
\left(\begin{array}{cccc|cccc}1&1&1&1&0&1&0&0\\0&1&1&1&1&0&0&0\\0&0&0&1&-1&0&1&0\\0&0&1&1&0&0&0&1\end{array}\right)
\xrightarrow{R_3\leftrightarrow R_4}
\left(\begin{array}{cccc|cccc}1&1&1&1&0&1&0&0\\0&1&1&1&1&0&0&0\\0&0&1&1&0&0&0&1\\0&0&0&1&-1&0&1&0\end{array}\right)
\\
\xrightarrow[R_1\leftarrow R_1-R_4]{R_3\leftarrow R_3-R_4,\ R_2\leftarrow R_2-R_4}
\left(\begin{array}{cccc|cccc}1&1&1&0&1&1&-1&0\\0&1&1&0&2&0&-1&0\\0&0&1&0&1&0&-1&1\\0&0&0&1&-1&0&1&0\end{array}\right)
\\
\xrightarrow[R_1\leftarrow R_1-R_3]{R_2\leftarrow R_2-R_3}
\left(\begin{array}{cccc|cccc}1&1&0&0&0&1&0&-1\\0&1&0&0&1&0&0&-1\\0&0&1&0&1&0&-1&1\\0&0&0&1&-1&0&1&0\end{array}\right)
\xrightarrow{R_1\leftarrow R_1-R_2}
\left(\begin{array}{cccc|cccc}1&0&0&0&-1&1&0&0\\0&1&0&0&1&0&0&-1\\0&0&1&0&1&0&-1&1\\0&0&0&1&-1&0&1&0\end{array}\right)
\end{gathered}
\]
So
\[
K^{-1}=\begin{pmatrix}-1&1&0&0\\1&0&0&-1\\1&0&-1&1\\-1&0&1&0\end{pmatrix}
\]
Check: \(\row_3(K)K^{-1}=(0,1,1,2)K^{-1}=(1+1-2,\ 0,\ -1+2,\ -1+1)=(0,0,1,0)\), and similarly
for the other rows.
\end{solution}

\begin{exercise}\label{r2-dr-inverse4}
Find the inverse of \(A=\begin{pmatrix}1&2&0&0\\1&2&0&1\\1&2&1&1\\1&1&0&1\end{pmatrix}\) by row reduction of
\([A\mid I_4]\), and check \(AA^{-1}=I_4\).
\end{exercise}
\begin{solution}
After the first arrow, rows \(2\) and \(3\) are \(0\) in column \(2\), so Step 2 interchanges rows \(2\) and \(4\):
\[
\begin{gathered}
\left(\begin{array}{cccc|cccc}1&2&0&0&1&0&0&0\\1&2&0&1&0&1&0&0\\1&2&1&1&0&0&1&0\\1&1&0&1&0&0&0&1\end{array}\right)
\xrightarrow[R_4\leftarrow R_4-R_1]{R_2\leftarrow R_2-R_1,\ R_3\leftarrow R_3-R_1}
\left(\begin{array}{cccc|cccc}1&2&0&0&1&0&0&0\\0&0&0&1&-1&1&0&0\\0&0&1&1&-1&0&1&0\\0&-1&0&1&-1&0&0&1\end{array}\right)
\\
\xrightarrow{R_2\leftrightarrow R_4}
\left(\begin{array}{cccc|cccc}1&2&0&0&1&0&0&0\\0&-1&0&1&-1&0&0&1\\0&0&1&1&-1&0&1&0\\0&0&0&1&-1&1&0&0\end{array}\right)
\xrightarrow{R_2\leftarrow-R_2}
\left(\begin{array}{cccc|cccc}1&2&0&0&1&0&0&0\\0&1&0&-1&1&0&0&-1\\0&0&1&1&-1&0&1&0\\0&0&0&1&-1&1&0&0\end{array}\right)
\\
\xrightarrow[R_2\leftarrow R_2+R_4]{R_3\leftarrow R_3-R_4}
\left(\begin{array}{cccc|cccc}1&2&0&0&1&0&0&0\\0&1&0&0&0&1&0&-1\\0&0&1&0&0&-1&1&0\\0&0&0&1&-1&1&0&0\end{array}\right)
\xrightarrow{R_1\leftarrow R_1-2R_2}
\left(\begin{array}{cccc|cccc}1&0&0&0&1&-2&0&2\\0&1&0&0&0&1&0&-1\\0&0&1&0&0&-1&1&0\\0&0&0&1&-1&1&0&0\end{array}\right)
\end{gathered}
\]
So
\[
A^{-1}=\begin{pmatrix}1&-2&0&2\\0&1&0&-1\\0&-1&1&0\\-1&1&0&0\end{pmatrix}
\]
Check: \(\row_1(A)A^{-1}=(1,\ -2+2,\ 0,\ 2-2)=(1,0,0,0)\),
\(\row_2(A)A^{-1}=(1-1,\ -2+2+1,\ 0,\ 2-2)=(0,1,0,0)\),
\(\row_3(A)A^{-1}=(1-1,\ -2+2-1+1,\ 1,\ 2-2)=(0,0,1,0)\) and
\(\row_4(A)A^{-1}=(1-1,\ -2+1+1,\ 0,\ 2-1)=(0,0,0,1)\).
\end{solution}

\begin{exercise}\label{r2-dr-elementary}
Let
\[
A=\begin{pmatrix}1&-1&0&0\\0&0&1&0\\-2&1&2&-1\\0&0&-1&1\end{pmatrix}
\]
\begin{enumerate}
\item Row-reduce \([A\mid I_4]\) to \([I_4\mid A^{-1}]\), one elementary operation per arrow.
\item Write \(A^{-1}\) as a product of elementary matrices.
\item Write \(A\) as a product of elementary matrices, and check the product.
\end{enumerate}
\end{exercise}
\begin{solution}
(a) After the first arrow, row \(2\) is \(0\) in column \(2\), so Step 2 interchanges rows \(2\) and \(3\):
\[
\begin{gathered}
\left(\begin{array}{cccc|cccc}1&-1&0&0&1&0&0&0\\0&0&1&0&0&1&0&0\\-2&1&2&-1&0&0&1&0\\0&0&-1&1&0&0&0&1\end{array}\right)
\xrightarrow{R_3\leftarrow R_3+2R_1}
\left(\begin{array}{cccc|cccc}1&-1&0&0&1&0&0&0\\0&0&1&0&0&1&0&0\\0&-1&2&-1&2&0&1&0\\0&0&-1&1&0&0&0&1\end{array}\right)
\\
\xrightarrow{R_2\leftrightarrow R_3}
\left(\begin{array}{cccc|cccc}1&-1&0&0&1&0&0&0\\0&-1&2&-1&2&0&1&0\\0&0&1&0&0&1&0&0\\0&0&-1&1&0&0&0&1\end{array}\right)
\xrightarrow{R_2\leftarrow-R_2}
\left(\begin{array}{cccc|cccc}1&-1&0&0&1&0&0&0\\0&1&-2&1&-2&0&-1&0\\0&0&1&0&0&1&0&0\\0&0&-1&1&0&0&0&1\end{array}\right)
\\
\xrightarrow{R_4\leftarrow R_4+R_3}
\left(\begin{array}{cccc|cccc}1&-1&0&0&1&0&0&0\\0&1&-2&1&-2&0&-1&0\\0&0&1&0&0&1&0&0\\0&0&0&1&0&1&0&1\end{array}\right)
\xrightarrow{R_2\leftarrow R_2-R_4}
\left(\begin{array}{cccc|cccc}1&-1&0&0&1&0&0&0\\0&1&-2&0&-2&-1&-1&-1\\0&0&1&0&0&1&0&0\\0&0&0&1&0&1&0&1\end{array}\right)
\\
\xrightarrow{R_2\leftarrow R_2+2R_3}
\left(\begin{array}{cccc|cccc}1&-1&0&0&1&0&0&0\\0&1&0&0&-2&1&-1&-1\\0&0&1&0&0&1&0&0\\0&0&0&1&0&1&0&1\end{array}\right)
\xrightarrow{R_1\leftarrow R_1+R_2}
\left(\begin{array}{cccc|cccc}1&0&0&0&-1&1&-1&-1\\0&1&0&0&-2&1&-1&-1\\0&0&1&0&0&1&0&0\\0&0&0&1&0&1&0&1\end{array}\right)
\end{gathered}
\]
So
\[
A^{-1}=\begin{pmatrix}-1&1&-1&-1\\-2&1&-1&-1\\0&1&0&0\\0&1&0&1\end{pmatrix}
\]
(b) The seven operations have the elementary matrices \(E_1=T_{31}(2)\), \(E_2=P_{23}\), \(E_3=D_2(-1)\),
\(E_4=T_{43}(1)\), \(E_5=T_{24}(-1)\), \(E_6=T_{23}(2)\), \(E_7=T_{12}(1)\), and \(E_7E_6E_5E_4E_3E_2E_1A=I_4\)
(Theorem~\ref{r2-thm-elementary}(c)). Hence, by the last sentence of Theorem~\ref{r2-thm-invertible},
\(A^{-1}=E_7E_6E_5E_4E_3E_2E_1\), the product, in this order, of
\[
\begin{gathered}
E_7=\begin{pmatrix}1&1&0&0\\0&1&0&0\\0&0&1&0\\0&0&0&1\end{pmatrix},\qquad E_6=\begin{pmatrix}1&0&0&0\\0&1&2&0\\0&0&1&0\\0&0&0&1\end{pmatrix},\qquad E_5=\begin{pmatrix}1&0&0&0\\0&1&0&-1\\0&0&1&0\\0&0&0&1\end{pmatrix}\\
E_4=\begin{pmatrix}1&0&0&0\\0&1&0&0\\0&0&1&0\\0&0&1&1\end{pmatrix},\qquad E_3=\begin{pmatrix}1&0&0&0\\0&-1&0&0\\0&0&1&0\\0&0&0&1\end{pmatrix}\\
E_2=\begin{pmatrix}1&0&0&0\\0&0&1&0\\0&1&0&0\\0&0&0&1\end{pmatrix},\qquad E_1=\begin{pmatrix}1&0&0&0\\0&1&0&0\\2&0&1&0\\0&0&0&1\end{pmatrix}
\end{gathered}
\]
(c) By Theorem~\ref{r2-thm-inverse-props}(a), \(A=(A^{-1})^{-1}\); so, by Theorems~\ref{r2-thm-inverse-props}(b)
and~\ref{r2-thm-elementary}(b),
\[
\begin{aligned}
A&=E_1^{-1}E_2^{-1}E_3^{-1}E_4^{-1}E_5^{-1}E_6^{-1}E_7^{-1}\\
&=T_{31}(-2)\,P_{23}\,D_2(-1)\,T_{43}(-1)\,T_{24}(1)\,T_{23}(-2)\,T_{12}(-1)
\end{aligned}
\]
Check: the product applied to \(I_4\) performs, from the right, the operations of \(T_{12}(-1)\), \(T_{23}(-2)\),
\(T_{24}(1)\), \(T_{43}(-1)\), \(D_2(-1)\), \(P_{23}\), \(T_{31}(-2)\) (Theorem~\ref{r2-thm-elementary}(a)):
\[
\begin{gathered}
\begin{pmatrix}1&0&0&0\\0&1&0&0\\0&0&1&0\\0&0&0&1\end{pmatrix}
\xrightarrow{R_1\leftarrow R_1-R_2}
\begin{pmatrix}1&-1&0&0\\0&1&0&0\\0&0&1&0\\0&0&0&1\end{pmatrix}
\xrightarrow{R_2\leftarrow R_2-2R_3}
\begin{pmatrix}1&-1&0&0\\0&1&-2&0\\0&0&1&0\\0&0&0&1\end{pmatrix}
\\
\xrightarrow{R_2\leftarrow R_2+R_4}
\begin{pmatrix}1&-1&0&0\\0&1&-2&1\\0&0&1&0\\0&0&0&1\end{pmatrix}
\xrightarrow{R_4\leftarrow R_4-R_3}
\begin{pmatrix}1&-1&0&0\\0&1&-2&1\\0&0&1&0\\0&0&-1&1\end{pmatrix}
\xrightarrow{R_2\leftarrow-R_2}
\begin{pmatrix}1&-1&0&0\\0&-1&2&-1\\0&0&1&0\\0&0&-1&1\end{pmatrix}
\\
\xrightarrow{R_2\leftrightarrow R_3}
\begin{pmatrix}1&-1&0&0\\0&0&1&0\\0&-1&2&-1\\0&0&-1&1\end{pmatrix}
\xrightarrow{R_3\leftarrow R_3-2R_1}
\begin{pmatrix}1&-1&0&0\\0&0&1&0\\-2&1&2&-1\\0&0&-1&1\end{pmatrix}
\end{gathered}
\]
which is \(A\).
\end{solution}

\begin{exercise}\label{r2-dr-row-equivalent}
Let \(A=\begin{pmatrix}1&-1&2&0\\-1&1&-1&1\\-1&1&0&2\end{pmatrix}\).
\begin{enumerate}
\item Row-reduce \([A\mid I_3]\) to \([R\mid P]\), where \(R\) is the RREF of \(A\).
\item Check that \(PA=R\).
\item For which \(y\in\R^3\) is \(Ax=y\) consistent?
\end{enumerate}
\end{exercise}
\begin{solution}
(a)
\[
\begin{gathered}
\left(\begin{array}{cccc|ccc}1&-1&2&0&1&0&0\\-1&1&-1&1&0&1&0\\-1&1&0&2&0&0&1\end{array}\right)
\xrightarrow[R_3\leftarrow R_3+R_1]{R_2\leftarrow R_2+R_1}
\left(\begin{array}{cccc|ccc}1&-1&2&0&1&0&0\\0&0&1&1&1&1&0\\0&0&2&2&1&0&1\end{array}\right)
\\
\xrightarrow{R_3\leftarrow R_3-2R_2}
\left(\begin{array}{cccc|ccc}1&-1&2&0&1&0&0\\0&0&1&1&1&1&0\\0&0&0&0&-1&-2&1\end{array}\right)
\xrightarrow{R_1\leftarrow R_1-2R_2}
\left(\begin{array}{cccc|ccc}1&-1&0&-2&-1&-2&0\\0&0&1&1&1&1&0\\0&0&0&0&-1&-2&1\end{array}\right)
\end{gathered}
\]
So \(R=\left(\begin{smallmatrix}1&-1&0&-2\\0&0&1&1\\0&0&0&0\end{smallmatrix}\right)\) and
\(P=\left(\begin{smallmatrix}-1&-2&0\\1&1&0\\-1&-2&1\end{smallmatrix}\right)\), which is invertible
(Corollary~\ref{r2-cor-row-equivalent}(a)).

(b) Row by row (Theorem~\ref{r1-thm-views}(c)):
\begin{itemize}
\item \(\row_1(P)A=-(1,-1,2,0)-2(-1,1,-1,1)=(1,-1,0,-2)\);
\item \(\row_2(P)A=(1,-1,2,0)+(-1,1,-1,1)=(0,0,1,1)\);
\item \(\row_3(P)A=-(1,-1,2,0)-2(-1,1,-1,1)+(-1,1,0,2)=(0,0,0,0)\).
\end{itemize}
(c) As \(P\) is invertible, \(Ax=y\) holds if and only if \(Rx=Py\) does. Row \(3\) of \(R\) is zero, so
\(Ax=y\) is consistent if and only if \(\row_3(P)\,y=-y_1-2y_2+y_3=0\) (Proposition~\ref{r2-prop-solutions}(a)).
\end{solution}

\begin{exercise}\label{r2-dr-solve}
\begin{enumerate}
\item Solve \(x_1-2x_3-x_4=2\), \(2x_1-x_2-3x_3=3\), \(-x_1+x_2+2x_3=-1\), \(-x_1+x_2+x_3=0\) by forward elimination
and back substitution.
\item Find all solutions of \(x_1+2x_2+x_3+x_4=-2\), \(-2x_1-3x_2-3x_3-2x_4=3\), \(-x_1-3x_2=1\), in vector form,
by Gauss--Jordan elimination.
\item Show that the system \(x_1-x_3=-1\), \(-x_1+x_2+x_3=1\), \(x_1-x_2=-1\), \(2x_1-x_2=-1\) has no solution.
\end{enumerate}
\end{exercise}
\begin{solution}
(a) Forward elimination of \([A\mid b]\):
\[
\begin{gathered}
\left(\begin{array}{cccc|c}1&0&-2&-1&2\\2&-1&-3&0&3\\-1&1&2&0&-1\\-1&1&1&0&0\end{array}\right)
\xrightarrow[R_4\leftarrow R_4+R_1]{R_2\leftarrow R_2-2R_1,\ R_3\leftarrow R_3+R_1}
\left(\begin{array}{cccc|c}1&0&-2&-1&2\\0&-1&1&2&-1\\0&1&0&-1&1\\0&1&-1&-1&2\end{array}\right)
\\
\xrightarrow{R_2\leftarrow-R_2}
\left(\begin{array}{cccc|c}1&0&-2&-1&2\\0&1&-1&-2&1\\0&1&0&-1&1\\0&1&-1&-1&2\end{array}\right)
\xrightarrow[R_4\leftarrow R_4-R_2]{R_3\leftarrow R_3-R_2}
\left(\begin{array}{cccc|c}1&0&-2&-1&2\\0&1&-1&-2&1\\0&0&1&1&0\\0&0&0&1&1\end{array}\right)
\end{gathered}
\]
\begin{itemize}
\item The last matrix is an echelon form with pivots in columns \(1,2,3,4\) and none in the last column: there is
exactly one solution (Theorems~\ref{thm:r2-consistency} and~\ref{thm:r2-solution-structure}(b)).
\item Back substitution: \(x_4=1\); \(x_3=-x_4=-1\); \(x_2=1+x_3+2x_4=2\); \(x_1=2+2x_3+x_4=1\).
So \(x=(1,2,-1,1)'\).
\item Check: \(1+2-1=2\), \(2-2+3=3\), \(-1+2-2=-1\), \(-1+2-1=0\).
\end{itemize}
(b) Gauss--Jordan elimination of \([A\mid b]\):
\[
\begin{gathered}
\left(\begin{array}{cccc|c}1&2&1&1&-2\\-2&-3&-3&-2&3\\-1&-3&0&0&1\end{array}\right)
\xrightarrow[R_3\leftarrow R_3+R_1]{R_2\leftarrow R_2+2R_1}
\left(\begin{array}{cccc|c}1&2&1&1&-2\\0&1&-1&0&-1\\0&-1&1&1&-1\end{array}\right)
\\
\xrightarrow{R_3\leftarrow R_3+R_2}
\left(\begin{array}{cccc|c}1&2&1&1&-2\\0&1&-1&0&-1\\0&0&0&1&-2\end{array}\right)
\xrightarrow{R_1\leftarrow R_1-R_3}
\left(\begin{array}{cccc|c}1&2&1&0&0\\0&1&-1&0&-1\\0&0&0&1&-2\end{array}\right)
\\
\xrightarrow{R_1\leftarrow R_1-2R_2}
\left(\begin{array}{cccc|c}1&0&3&0&2\\0&1&-1&0&-1\\0&0&0&1&-2\end{array}\right)
\end{gathered}
\]
\begin{itemize}
\item The left part of the last matrix is the RREF of \(A\), with pivot columns \(1,2,4\), and the last column has
no pivot: the system is consistent, and \(x_3=t\) is free (Proposition~\ref{r2-prop-solutions}).
\item The rows give \(x_1=2-3t\), \(x_2=-1+t\), \(x_4=-2\), so
\[
x=\begin{pmatrix}2\\-1\\0\\-2\end{pmatrix}+t\begin{pmatrix}-3\\1\\1\\0\end{pmatrix}\qquad(t\in\R)
\]
\item Check: \(A(2,-1,0,-2)'=(2-2-2,\ -4+3+4,\ -2+3)'=(-2,3,1)'=b\) and
\(A(-3,1,1,0)'=(-3+2+1,\ 6-3-3,\ 3-3)'=0\).
\end{itemize}
(c) Forward elimination of \([A\mid b]\):
\[
\begin{gathered}
\left(\begin{array}{ccc|c}1&0&-1&-1\\-1&1&1&1\\1&-1&0&-1\\2&-1&0&-1\end{array}\right)
\xrightarrow[R_4\leftarrow R_4-2R_1]{R_2\leftarrow R_2+R_1,\ R_3\leftarrow R_3-R_1}
\left(\begin{array}{ccc|c}1&0&-1&-1\\0&1&0&0\\0&-1&1&0\\0&-1&2&1\end{array}\right)
\\
\xrightarrow[R_4\leftarrow R_4+R_2]{R_3\leftarrow R_3+R_2}
\left(\begin{array}{ccc|c}1&0&-1&-1\\0&1&0&0\\0&0&1&0\\0&0&2&1\end{array}\right)
\xrightarrow{R_4\leftarrow R_4-2R_3}
\left(\begin{array}{ccc|c}1&0&-1&-1\\0&1&0&0\\0&0&1&0\\0&0&0&1\end{array}\right)
\end{gathered}
\]
The last matrix is an echelon form whose last column is a pivot column: row \(4\) reads \(0=1\). So the system has
no solution (Theorem~\ref{thm:r2-consistency}). Indeed, rows \(1\)--\(3\) give \(x_3=0\), \(x_2=0\),
\(x_1=-1+x_3=-1\): the first three equations have only the solution \((-1,0,0)'\), and it violates the fourth,
\(2(-1)-0=-2\ne-1\).
\end{solution}

\begin{exercise}[cf.\ \citeat{abadir-magnus}{Exercise~6.27}]\label{r2-ex-two-systems}
Solve the equations
\[
\begin{pmatrix}1&-3&-2&-1\\-1&2&1&1\\-1&2&2&1\\0&1&1&1\end{pmatrix}\begin{pmatrix}x_1\\x_2\\x_3\\x_4\end{pmatrix}
=\begin{pmatrix}1\\-1\\-2\\1\end{pmatrix}
\]
and
\[
\begin{pmatrix}1&1&2&1\\2&2&3&1\\1&2&2&2\\0&1&1&3\end{pmatrix}\begin{pmatrix}x_1\\x_2\\x_3\\x_4\end{pmatrix}=\begin{pmatrix}4\\6\\3\\1\end{pmatrix}
\]
by Gaussian elimination.
\end{exercise}
\begin{solution}
First system. Forward elimination:
\[
\begin{gathered}
\left(\begin{array}{cccc|c}1&-3&-2&-1&1\\-1&2&1&1&-1\\-1&2&2&1&-2\\0&1&1&1&1\end{array}\right)
\xrightarrow[R_3\leftarrow R_3+R_1]{R_2\leftarrow R_2+R_1}
\left(\begin{array}{cccc|c}1&-3&-2&-1&1\\0&-1&-1&0&0\\0&-1&0&0&-1\\0&1&1&1&1\end{array}\right)
\\
\xrightarrow{R_2\leftarrow-R_2}
\left(\begin{array}{cccc|c}1&-3&-2&-1&1\\0&1&1&0&0\\0&-1&0&0&-1\\0&1&1&1&1\end{array}\right)
\xrightarrow[R_4\leftarrow R_4-R_2]{R_3\leftarrow R_3+R_2}
\left(\begin{array}{cccc|c}1&-3&-2&-1&1\\0&1&1&0&0\\0&0&1&0&-1\\0&0&0&1&1\end{array}\right)
\end{gathered}
\]
Backward elimination:
\[
\begin{gathered}
\xrightarrow{R_1\leftarrow R_1+R_4}
\left(\begin{array}{cccc|c}1&-3&-2&0&2\\0&1&1&0&0\\0&0&1&0&-1\\0&0&0&1&1\end{array}\right)
\xrightarrow[R_2\leftarrow R_2-R_3]{R_1\leftarrow R_1+2R_3}
\left(\begin{array}{cccc|c}1&-3&0&0&0\\0&1&0&0&1\\0&0&1&0&-1\\0&0&0&1&1\end{array}\right)
\\
\xrightarrow{R_1\leftarrow R_1+3R_2}
\left(\begin{array}{cccc|c}1&0&0&0&3\\0&1&0&0&1\\0&0&1&0&-1\\0&0&0&1&1\end{array}\right)
\end{gathered}
\]
By Proposition~\ref{r2-prop-solutions}, the unique solution is \((x_1,x_2,x_3,x_4)=(3,1,-1,1)\).
Check: \(3-3+2-1=1\), \(-3+2-1+1=-1\), \(-3+2-2+1=-2\), \(1-1+1=1\).

Second system. Forward elimination; after the first arrow the entry of row \(2\) in column \(2\)
is \(0\), so rows \(2\) and \(3\) are interchanged (Step 2):
\[
\begin{gathered}
\left(\begin{array}{cccc|c}1&1&2&1&4\\2&2&3&1&6\\1&2&2&2&3\\0&1&1&3&1\end{array}\right)
\xrightarrow[R_3\leftarrow R_3-R_1]{R_2\leftarrow R_2-2R_1}
\left(\begin{array}{cccc|c}1&1&2&1&4\\0&0&-1&-1&-2\\0&1&0&1&-1\\0&1&1&3&1\end{array}\right)
\xrightarrow{R_2\leftrightarrow R_3}
\left(\begin{array}{cccc|c}1&1&2&1&4\\0&1&0&1&-1\\0&0&-1&-1&-2\\0&1&1&3&1\end{array}\right)
\\
\xrightarrow[R_3\leftarrow -R_3]{R_4\leftarrow R_4-R_2}
\left(\begin{array}{cccc|c}1&1&2&1&4\\0&1&0&1&-1\\0&0&1&1&2\\0&0&1&2&2\end{array}\right)
\xrightarrow{R_4\leftarrow R_4-R_3}
\left(\begin{array}{cccc|c}1&1&2&1&4\\0&1&0&1&-1\\0&0&1&1&2\\0&0&0&1&0\end{array}\right)
\end{gathered}
\]
Every column of the left part is a pivot column and the last column is not, so there is exactly
one solution (Theorems~\ref{thm:r2-consistency} and~\ref{thm:r2-solution-structure}(b)). Back
substitution: \(x_4=0\); \(x_3=2-x_4=2\); \(x_2=-1-x_4=-1\);
\(x_1=4-x_2-2x_3-x_4=4+1-4-0=1\). The unique solution is \((x_1,x_2,x_3,x_4)=(1,-1,2,0)\).
Check: \(1-1+4+0=4\), \(2-2+6+0=6\), \(1-2+4+0=3\), \(0-1+2+0=1\).
\end{solution}

\begin{exercise}\label{r2-ex-system-drill}
Let
\[
A=\begin{pmatrix}1&3&5&-1\\1&2&2&0\\0&1&3&-1\\0&0&1&-1\end{pmatrix},\qquad
b=\begin{pmatrix}6\\4\\2\\0\end{pmatrix},\qquad c=\begin{pmatrix}6\\4\\3\\0\end{pmatrix}
\]
\begin{enumerate}
\item Find all solutions of \(Ax=b\).
\item Show that \(Ax=c\) has no solution.
\end{enumerate}
\end{exercise}
\begin{solution}
(a) Row-reduce \([A\mid b]\):
\[
\begin{gathered}
\left(\begin{array}{cccc|c}1&3&5&-1&6\\1&2&2&0&4\\0&1&3&-1&2\\0&0&1&-1&0\end{array}\right)
\xrightarrow{R_2\leftarrow R_2-R_1}
\left(\begin{array}{cccc|c}1&3&5&-1&6\\0&-1&-3&1&-2\\0&1&3&-1&2\\0&0&1&-1&0\end{array}\right)
\xrightarrow{R_2\leftarrow-R_2}
\left(\begin{array}{cccc|c}1&3&5&-1&6\\0&1&3&-1&2\\0&1&3&-1&2\\0&0&1&-1&0\end{array}\right)
\\
\xrightarrow{R_3\leftarrow R_3-R_2}
\left(\begin{array}{cccc|c}1&3&5&-1&6\\0&1&3&-1&2\\0&0&0&0&0\\0&0&1&-1&0\end{array}\right)
\xrightarrow{R_3\leftrightarrow R_4}
\left(\begin{array}{cccc|c}1&3&5&-1&6\\0&1&3&-1&2\\0&0&1&-1&0\\0&0&0&0&0\end{array}\right)
\\
\xrightarrow[R_1\leftarrow R_1-5R_3]{R_2\leftarrow R_2-3R_3}
\left(\begin{array}{cccc|c}1&3&0&4&6\\0&1&0&2&2\\0&0&1&-1&0\\0&0&0&0&0\end{array}\right)
\xrightarrow{R_1\leftarrow R_1-3R_2}
\left(\begin{array}{cccc|c}1&0&0&-2&0\\0&1&0&2&2\\0&0&1&-1&0\\0&0&0&0&0\end{array}\right)
\end{gathered}
\]
After the third arrow the entry of row \(3\) in column \(3\) is \(0\), so Step 2 interchanges rows
\(3\) and \(4\). The left part of the last matrix is the RREF \(R\) of \(A\), with \(r=3\) and pivot
columns \(1,2,3\); the unknown \(x_4\) is free. The last matrix is \([R\mid z]\) with
\(z=(0,2,0,0)'\). As \(z_4=0\), the system is consistent (Proposition~\ref{r2-prop-solutions}(a)),
and with \(x_4=t\), Proposition~\ref{r2-prop-solutions}(b) gives \(x_1=2t\), \(x_2=2-2t\), \(x_3=t\):
\[
x=\begin{pmatrix}0\\2\\0\\0\end{pmatrix}+t\begin{pmatrix}2\\-2\\1\\1\end{pmatrix}\qquad(t\in\R)
\]
infinitely many solutions (Theorem~\ref{thm:r2-solution-structure}(a),(b)). Check:
\(A(0,2,0,0)'=2\col_2(A)=(6,4,2,0)'=b\) and
\(A(2,-2,1,1)'=(2-6+5-1,\ 2-4+2,\ -2+3-1,\ 1-1)'=0\).

(b) Row-reduce \([A\mid c]\) by the same first three operations:
\[
\begin{gathered}
\left(\begin{array}{cccc|c}1&3&5&-1&6\\1&2&2&0&4\\0&1&3&-1&3\\0&0&1&-1&0\end{array}\right)
\xrightarrow{R_2\leftarrow R_2-R_1}
\left(\begin{array}{cccc|c}1&3&5&-1&6\\0&-1&-3&1&-2\\0&1&3&-1&3\\0&0&1&-1&0\end{array}\right)
\\
\xrightarrow{R_2\leftarrow-R_2}
\left(\begin{array}{cccc|c}1&3&5&-1&6\\0&1&3&-1&2\\0&1&3&-1&3\\0&0&1&-1&0\end{array}\right)
\xrightarrow{R_3\leftarrow R_3-R_2}
\left(\begin{array}{cccc|c}1&3&5&-1&6\\0&1&3&-1&2\\0&0&0&0&1\\0&0&1&-1&0\end{array}\right)
\end{gathered}
\]
The third row of the last matrix reads \((0\ 0\ 0\ 0\mid1)\), that is, \(0=1\); so \(Ax=c\) has no
solution (Theorem~\ref{thm:r2-consistency}). Indeed \(\row_3(A)=\row_1(A)-\row_2(A)\), while
\(c_3=3\ne2=c_1-c_2\).
\end{solution}

\begin{exercise}\label{r2-dr-general}
Consider the system
\[
\begin{aligned}
x_1-2x_3+3x_5&=-1\\
2x_1+x_2-x_3-x_4+3x_5&=-4\\
-2x_1+4x_3+x_4-4x_5&=3\\
x_1+x_2+x_3+2x_5&=-2
\end{aligned}
\]
\begin{enumerate}
\item Reduce its augmented matrix \([A\mid b]\) to RREF.
\item Write all solutions in the form \(x=x_p+sv_1+tv_2\) (\(s,t\in\R\)).
\item Check that \(Ax_p=b\) and \(Av_1=Av_2=0\).
\end{enumerate}
\end{exercise}
\begin{solution}
(a)
\[
\begin{gathered}
\left(\begin{array}{ccccc|c}1&0&-2&0&3&-1\\2&1&-1&-1&3&-4\\-2&0&4&1&-4&3\\1&1&1&0&2&-2\end{array}\right)
\xrightarrow[R_4\leftarrow R_4-R_1]{R_2\leftarrow R_2-2R_1,\ R_3\leftarrow R_3+2R_1}
\left(\begin{array}{ccccc|c}1&0&-2&0&3&-1\\0&1&3&-1&-3&-2\\0&0&0&1&2&1\\0&1&3&0&-1&-1\end{array}\right)
\\
\xrightarrow{R_4\leftarrow R_4-R_2}
\left(\begin{array}{ccccc|c}1&0&-2&0&3&-1\\0&1&3&-1&-3&-2\\0&0&0&1&2&1\\0&0&0&1&2&1\end{array}\right)
\xrightarrow{R_4\leftarrow R_4-R_3}
\left(\begin{array}{ccccc|c}1&0&-2&0&3&-1\\0&1&3&-1&-3&-2\\0&0&0&1&2&1\\0&0&0&0&0&0\end{array}\right)
\\
\xrightarrow{R_2\leftarrow R_2+R_3}
\left(\begin{array}{ccccc|c}1&0&-2&0&3&-1\\0&1&3&0&-1&-1\\0&0&0&1&2&1\\0&0&0&0&0&0\end{array}\right)
\end{gathered}
\]
The last matrix is \([R\mid z]\), with \(R\) the RREF of \(A\), pivot columns \(1,2,4\) and \(z=(-1,-1,1,0)'\).

(b) Since \(z_4=0\), the system is consistent (Proposition~\ref{r2-prop-solutions}(a)). The free unknowns are
\(x_3=s\) and \(x_5=t\), and the rows give \(x_1=-1+2s-3t\), \(x_2=-1-3s+t\), \(x_4=1-2t\)
(Proposition~\ref{r2-prop-solutions}(b)):
\[
x=\begin{pmatrix}-1\\-1\\0\\1\\0\end{pmatrix}+s\begin{pmatrix}2\\-3\\1\\0\\0\end{pmatrix}+t\begin{pmatrix}-3\\1\\0\\-2\\1\end{pmatrix}
=x_p+sv_1+tv_2
\]
(c) Row by row:
\begin{itemize}
\item \(Ax_p=(-1,\ -2-1-1,\ 2+1,\ -1-1)'=(-1,-4,3,-2)'=b\);
\item \(Av_1=(2-2,\ 4-3-1,\ -4+4,\ 2-3+1)'=0\);
\item \(Av_2=(-3+3,\ -6+1+2+3,\ 6-2-4,\ -3+1+2)'=0\).
\end{itemize}
So \(x_p\) is a particular solution, and \(sv_1+tv_2\) runs through the solutions of \(Ax=0\)
(Theorem~\ref{thm:r2-solution-structure}(a)).
\end{solution}

\begin{exercise}\label{r2-dr-parameter}
For which values of \(k\) does the system
\[
\begin{aligned}
x_1+2x_2-x_3&=1\\
2x_1+5x_2+x_3+x_4&=3\\
x_1+3x_2+3x_3+2x_4&=2\\
x_1+2x_2+k^2x_4&=k+2
\end{aligned}
\]
have
\begin{enumerate}
\item no solution,
\item exactly one solution,
\item infinitely many solutions?
\end{enumerate}
Give the solutions when they exist.
\end{exercise}
\begin{solution}
Forward elimination:
\[
\begin{gathered}
\left(\begin{array}{cccc|c}1&2&-1&0&1\\2&5&1&1&3\\1&3&3&2&2\\1&2&0&k^2&k+2\end{array}\right)
\xrightarrow[R_4\leftarrow R_4-R_1]{R_2\leftarrow R_2-2R_1,\ R_3\leftarrow R_3-R_1}
\left(\begin{array}{cccc|c}1&2&-1&0&1\\0&1&3&1&1\\0&1&4&2&1\\0&0&1&k^2&k+1\end{array}\right)
\\
\xrightarrow{R_3\leftarrow R_3-R_2}
\left(\begin{array}{cccc|c}1&2&-1&0&1\\0&1&3&1&1\\0&0&1&1&0\\0&0&1&k^2&k+1\end{array}\right)
\xrightarrow{R_4\leftarrow R_4-R_3}
\left(\begin{array}{cccc|c}1&2&-1&0&1\\0&1&3&1&1\\0&0&1&1&0\\0&0&0&k^2-1&k+1\end{array}\right)
\end{gathered}
\]
Row \(4\) reads \((k^2-1)x_4=k+1\), where \(k^2-1=(k-1)(k+1)\).
\begin{enumerate}
\item \(k=1\): row \(4\) reads \((0\ 0\ 0\ 0\mid2)\), that is, \(0=2\); so there is no solution
(Theorem~\ref{thm:r2-consistency}).
\item \(k\ne\pm1\): dividing row \(4\) by \(k^2-1\) gives an echelon form with pivots in columns \(1,2,3,4\) and
none in the last column, so there is exactly one solution
(Theorems~\ref{thm:r2-consistency} and~\ref{thm:r2-solution-structure}(b)). Back
substitution: \(x_4=\frac{k+1}{k^2-1}=\frac1{k-1}\); \(x_3=-x_4=-\frac1{k-1}\);
\(x_2=1-3x_3-x_4=\frac{k+1}{k-1}\); \(x_1=1-2x_2+x_3=-\frac{k+4}{k-1}\).
So \(x=\frac1{k-1}(-k-4,\ k+1,\ -1,\ 1)'\); for \(k=2\), \(x=(-6,3,-1,1)'\), and indeed \(-6+6+1=1\),
\(-12+15-1+1=3\), \(-6+9-3+2=2\), \(-6+6+4=4=k+2\).
\item \(k=-1\): row \(4\) is zero, and \(x_4=t\) is free, so there are infinitely many solutions
(Theorems~\ref{thm:r2-consistency} and~\ref{thm:r2-solution-structure}(b)): \(x_3=-t\),
\(x_2=1-3x_3-x_4=1+2t\), \(x_1=1-2x_2+x_3=-1-5t\), that is, \(x=(-1,1,0,0)'+t\,(-5,2,-1,1)'\) (\(t\in\R\)).
\end{enumerate}
\end{solution}

\begin{exercise}[cf.\ \citeat{abadir-magnus}{Exercise~6.47}]\label{r2-ex-parameters}
For which values of \(k\) and \(h\) does the system
\[
\begin{aligned}
x_1+x_2+x_3+x_4&=1\\
x_1+2x_2+2x_3+2x_4&=2\\
x_1+2x_2+3x_3+3x_4&=4\\
x_1+2x_2+3x_3+kx_4&=h
\end{aligned}
\]
have no solution, exactly one solution, or infinitely many solutions? Give the solutions when
they exist.
\end{exercise}
\begin{solution}
Forward elimination; each arrow subtracts a row from all rows below it:
\[
\begin{gathered}
\left(\begin{array}{cccc|c}1&1&1&1&1\\1&2&2&2&2\\1&2&3&3&4\\1&2&3&k&h\end{array}\right)
\xrightarrow[R_4\leftarrow R_4-R_1]{R_2\leftarrow R_2-R_1,\ R_3\leftarrow R_3-R_1}
\left(\begin{array}{cccc|c}1&1&1&1&1\\0&1&1&1&1\\0&1&2&2&3\\0&1&2&k-1&h-1\end{array}\right)
\\
\xrightarrow[R_4\leftarrow R_4-R_2]{R_3\leftarrow R_3-R_2}
\left(\begin{array}{cccc|c}1&1&1&1&1\\0&1&1&1&1\\0&0&1&1&2\\0&0&1&k-2&h-2\end{array}\right)
\xrightarrow{R_4\leftarrow R_4-R_3}
\left(\begin{array}{cccc|c}1&1&1&1&1\\0&1&1&1&1\\0&0&1&1&2\\0&0&0&k-3&h-4\end{array}\right)
\end{gathered}
\]
\begin{itemize}
\item \(k\ne3\): \(R_4\leftarrow\frac1{k-3}R_4\) gives an echelon form with pivots in columns
\(1,2,3,4\) and none in the last column. So there is exactly one solution
(Theorems~\ref{thm:r2-consistency} and~\ref{thm:r2-solution-structure}(b)). Back substitution
with \(s=\frac{h-4}{k-3}\): \(x_4=s\), \(x_3=2-s\), \(x_2=1-x_3-x_4=-1\),
\(x_1=1-x_2-x_3-x_4=0\), that is, \(x=(0,-1,2,0)'+s\,(0,0,-1,1)'\).
\item \(k=3\), \(h\ne4\): row \(4\) reads \(0=h-4\ne0\), so there is no solution.
\item \(k=3\), \(h=4\): row \(4\) is \(0\), and the first three rows form an echelon form with
pivots in columns \(1,2,3\); \(x_4=t\) is free, so there are infinitely many solutions:
\(x=(0,-1,2,0)'+t\,(0,0,-1,1)'\) (\(t\in\R\)).
\end{itemize}
Check: \((0,-1,2,0)'\) satisfies the first three equations (\(0-1+2+0=1\),
\(0-2+4+0=2\), \(0-2+6+0=4\)), and \((0,0,-1,1)'\) satisfies their homogeneous versions
(\(-1+1=0\), \(-2+2=0\), \(-3+3=0\)). In the fourth equation, \((0,-1,2,0)'\) gives \(0-2+6+k\cdot0=4\) and
\((0,0,-1,1)'\) gives \(-3+k\); so for \(k\ne3\) its left side is \(4+s(k-3)=4+(h-4)=h\),
and for \(k=3\), \(h=4\) it is \(4+t(3-3)=4=h\).
\end{solution}

\begin{exercise}\label{r2-dr-triangular}
\begin{enumerate}
\item Find \(U^{-1}\) for \(U=\begin{pmatrix}1&1&2&-1\\0&-1&-1&1\\0&0&-1&1\\0&0&0&1\end{pmatrix}\) by the formulas
of Theorem~\ref{r2-thm-triangular}(b).
\item Find \(L^{-1}\) for \(L=\begin{pmatrix}1&0&0&0\\2&1&0&0\\-1&2&-1&0\\0&1&1&1\end{pmatrix}\) by row reduction of \([L\mid I_4]\).
\item Show that \(G=\begin{pmatrix}1&2&-1&1\\0&0&3&1\\0&0&-1&2\\0&0&0&1\end{pmatrix}\) is not invertible by finding \(x\ne0\) with \(Gx=0\).
\end{enumerate}
\end{exercise}
\begin{solution}
(a) The diagonal entries of \(B=U^{-1}\) are \(b_{11}=1\), \(b_{22}=-1\), \(b_{33}=-1\), \(b_{44}=1\), and
\(b_{ij}=-\frac1{u_{ii}}\sum_{k=i+1}^ju_{ik}b_{kj}\) for \(i<j\), where \(-1/u_{11}=-1\), \(-1/u_{22}=1\) and
\(-1/u_{33}=1\):
\begin{itemize}
\item column \(2\): \(b_{12}=-u_{12}b_{22}=-(1)(-1)=1\);
\item column \(3\): \(b_{23}=u_{23}b_{33}=(-1)(-1)=1\), \(b_{13}=-(u_{12}b_{23}+u_{13}b_{33})=-(1-2)=1\);
\item column \(4\): \(b_{34}=u_{34}b_{44}=1\), \(b_{24}=u_{23}b_{34}+u_{24}b_{44}=-1+1=0\),
\(b_{14}=-(u_{12}b_{24}+u_{13}b_{34}+u_{14}b_{44})=-(0+2-1)=-1\).
\end{itemize}
So
\[
U^{-1}=\begin{pmatrix}1&1&1&-1\\0&-1&1&0\\0&0&-1&1\\0&0&0&1\end{pmatrix}
\]
Check: \(\row_1(U)U^{-1}=(1,\ 1-1,\ 1+1-2,\ -1+0+2-1)=(1,0,0,0)\), and similarly for the other rows.

(b) \(L\) is lower triangular, so forward elimination and one scaling already turn the left half into \(I_4\):
\[
\begin{gathered}
\left(\begin{array}{cccc|cccc}1&0&0&0&1&0&0&0\\2&1&0&0&0&1&0&0\\-1&2&-1&0&0&0&1&0\\0&1&1&1&0&0&0&1\end{array}\right)
\xrightarrow[R_3\leftarrow R_3+R_1]{R_2\leftarrow R_2-2R_1}
\left(\begin{array}{cccc|cccc}1&0&0&0&1&0&0&0\\0&1&0&0&-2&1&0&0\\0&2&-1&0&1&0&1&0\\0&1&1&1&0&0&0&1\end{array}\right)
\\
\xrightarrow[R_4\leftarrow R_4-R_2]{R_3\leftarrow R_3-2R_2}
\left(\begin{array}{cccc|cccc}1&0&0&0&1&0&0&0\\0&1&0&0&-2&1&0&0\\0&0&-1&0&5&-2&1&0\\0&0&1&1&2&-1&0&1\end{array}\right)
\xrightarrow{R_4\leftarrow R_4+R_3}
\left(\begin{array}{cccc|cccc}1&0&0&0&1&0&0&0\\0&1&0&0&-2&1&0&0\\0&0&-1&0&5&-2&1&0\\0&0&0&1&7&-3&1&1\end{array}\right)
\\
\xrightarrow{R_3\leftarrow-R_3}
\left(\begin{array}{cccc|cccc}1&0&0&0&1&0&0&0\\0&1&0&0&-2&1&0&0\\0&0&1&0&-5&2&-1&0\\0&0&0&1&7&-3&1&1\end{array}\right)
\end{gathered}
\]
So
\[
L^{-1}=\begin{pmatrix}1&0&0&0\\-2&1&0&0\\-5&2&-1&0\\7&-3&1&1\end{pmatrix}
\]
lower triangular with diagonal entries \(1,1,-1,1\), the reciprocals of those of \(L\)
(Theorem~\ref{r2-thm-triangular}(b)). Check: \(\row_4(L)L^{-1}=(-2-5+7,\ 1+2-3,\ -1+1,\ 1)=(0,0,0,1)\).

(c) \(Gx=0\) reads \(x_1+2x_2-x_3+x_4=0\), \(3x_3+x_4=0\), \(-x_3+2x_4=0\), \(x_4=0\). So \(x_4=0\), \(x_3=0\) and
\(x_1=-2x_2\); with \(x_2=1\), \(x=(-2,1,0,0)'\ne0\) and \(Gx=(-2+2,\ 0,\ 0,\ 0)'=0\). Hence \(G\) is not invertible
(Theorem~\ref{r2-thm-invertible}), as Theorem~\ref{r2-thm-triangular}(a) predicts: \(g_{22}=0\).
\end{solution}

\begin{exercise}\label{r2-dr-centering}
Six units are measured on two variables:
\(X_0=\begin{pmatrix}2&2&0&0&2&0\\2&0&2&0&1&1\end{pmatrix}'\in\R^{6\times2}\). Let \(M=[\one_6\ X_0]\in\R^{6\times3}\).
Find the mean vector \(\bar x\), the centered data \(CX_0\), a product \(Q\) of elementary matrices with
\([\one_6\ CX_0]=MQ\), and \(Q^{-1}\).
\end{exercise}
\begin{solution}
\begin{itemize}
\item \(\one'X_0=(6,6)\), so \(\bar x=\frac16X_0'\one=(1,1)'\), and \(CX_0=X_0-\one\bar x'\)
(Example~\ref{r1-exm-data}) has the columns \((1,1,-1,-1,1,-1)'\) and \((1,-1,1,-1,0,0)'\).
\item \([\one_6\ CX_0]\) arises from \(M\) by \(\col_2\leftarrow\col_2-\col_1\) and \(\col_3\leftarrow\col_3-\col_1\), so
(Theorem~\ref{r2-thm-column-ops})
\[
Q=T_{12}(-1)\,T_{13}(-1)=\begin{pmatrix}1&-1&-1\\0&1&0\\0&0&1\end{pmatrix},\qquad
Q^{-1}=T_{13}(1)\,T_{12}(1)=\begin{pmatrix}1&1&1\\0&1&0\\0&0&1\end{pmatrix}
\]
by Theorems~\ref{r2-thm-inverse-props}(b) and~\ref{r2-thm-elementary}(b); \(Q^{-1}\) is unit upper triangular, as
Theorem~\ref{r2-thm-triangular} says.
\end{itemize}
\end{solution}

\needspace{14\baselineskip}
\subsection*{Problems}

\begin{problem}[\citeat{abadir-magnus}{Exercise~6.6}]\label{r2-ex-product-elementary}
Show that the product of two elementary matrices is, in general, not elementary.
\end{problem}
\begin{solution}
For example,
\[
P_{13}D_2(7)=\begin{pmatrix}0&0&1\\0&1&0\\1&0&0\end{pmatrix}\begin{pmatrix}1&0&0\\0&7&0\\0&0&1\end{pmatrix}
=\begin{pmatrix}0&0&1\\0&7&0\\1&0&0\end{pmatrix}
\]
This is not elementary: a matrix \(D_r(c)\) or \(T_{rs}(c)\) differs from \(I_3\) in one entry
only, while this product differs from \(I_3\) in five entries; and every \(P_{rs}\) has all
entries \(0\) or \(1\), while this product has the entry \(7\).
\end{solution}

\begin{problem}[\citeat{abadir-magnus}{Exercise~4.22(a), (e)}]\label{r2-ex-inverse-props}
Let \(A\in\R^{n\times n}\) be invertible. Show that:
\begin{enumerate}
\item \((\lambda A)^{-1}=\frac1\lambda A^{-1}\) (\(\lambda\ne0\));
\item if \(A\) and \(B\) commute and \(B\) is invertible, then \(A\) and \(B^{-1}\) commute.
\end{enumerate}
\end{problem}
\begin{solution}
(a) We check both products (Definition~\ref{r2-def-inverse}):
\((\lambda A)(\frac1\lambda A^{-1})=AA^{-1}=I\) and \((\frac1\lambda A^{-1})(\lambda A)=A^{-1}A=I\).
(b) Multiply \(AB=BA\) on the left and on the right by \(B^{-1}\):
\(B^{-1}ABB^{-1}=B^{-1}BAB^{-1}\), that is, \(B^{-1}A=AB^{-1}\).
\end{solution}

\begin{problem}[\citeat{lang}{Ch.~II, \S3, Exercise~18}, \citeat{shakarchi}{}]\label{r2-ex-nilpotent}
Let \(A\in\R^{n\times n}\).
\begin{enumerate}
\item If \(A^2=0\), show that \(I-A\) is invertible.
\item If \(A^3=0\), show that \(I-A\) is invertible.
\item In general, if \(A^k=0\) for some positive integer \(k\), show that \(I-A\) is invertible.
\item Suppose that \(A^2+2A+I=0\). Show that \(A\) is invertible.
\item Suppose that \(A^3-A+I=0\). Show that \(A\) is invertible.
\end{enumerate}
\end{problem}
\begin{solution}
(a) By Theorem~\ref{r1-thm-rules}, \((I-A)(I+A)=I+A-A-A^2=I-A^2\) and \((I+A)(I-A)=I-A^2\). If \(A^2=0\),
both products equal \(I\), so \(I-A\) is invertible with \((I-A)^{-1}=I+A\).

(b) By Theorem~\ref{r1-thm-rules}, \((I-A)(I+A+A^2)=(I+A+A^2)-(A+A^2+A^3)=I-A^3\) and in the same way
\((I+A+A^2)(I-A)=I-A^3\). If \(A^3=0\), both products equal \(I\), so \((I-A)^{-1}=I+A+A^2\).

(c) For every positive integer \(k\), by Theorem~\ref{r1-thm-rules},
\[
(I-A)(I+A+A^2+\dots+A^{k-1})=(I+A+\dots+A^{k-1})-(A+A^2+\dots+A^k)=I-A^k
\]
and in the same way \((I+A+\dots+A^{k-1})(I-A)=I-A^k\). If \(A^k=0\), both products equal
\(I\), so \(I-A\) is invertible with inverse \(I+A+\dots+A^{k-1}\).

(d) \(A(-A-2I)=(-A-2I)A=-A^2-2A=I\), since \(A^2+2A+I=0\). So \(A^{-1}=-A-2I\).

(e) \(A(-A^2+I)=(-A^2+I)A=-A^3+A=I\), since \(A^3-A+I=0\). So \(A^{-1}=I-A^2\).
\end{solution}

\begin{problem}[\citeat{hoffman-kunze}{\S1.6, Exercise~10}]\label{r2-pr-shape}
Let \(A\in\R^{m\times n}\) and \(B\in\R^{n\times m}\), with \(n<m\). Show that \(AB\) is not invertible.
\end{problem}
\begin{solution}
\(AB\) is an \(m\times m\) matrix. By Theorem~\ref{r2-thm-invertible} ((i)\,\(\Leftrightarrow\)\,(iii)), it suffices
to find \(x\ne0\) with \((AB)x=0\).
\begin{itemize}
\item The system \(Bx=0\) has \(n\) equations in the \(m\) unknowns \(x_1,\dots,x_m\), and \(n<m\). So it has a
nontrivial solution \(x\in\R^m\) (Theorem~\ref{r2-thm-homogeneous}(a)).
\item Then \((AB)x=A(Bx)=A0=0\) (Theorem~\ref{r1-thm-rules}), with \(x\ne0\). Hence \(AB\) is not invertible.
\end{itemize}
\end{solution}

\begin{problem}[\citeat{prasolov}{Problem~7.3}]\label{r2-ex-pr-diag-dominant}
Let \(A=(a_{ij})\in\R^{n\times n}\) satisfy \(\lvert a_{ii}\rvert>\sum_{k\ne i}\lvert a_{ik}\rvert\) for
\(i=1,\dots,n\). Show that \(A\) is invertible.
\end{problem}
\begin{solution}
By Theorem~\ref{r2-thm-invertible} ((iii)\,\(\Rightarrow\)\,(i)), it suffices to show that \(Ax=0\) only for \(x=0\).
Suppose that \(Ax=0\) with \(x\ne0\), and choose an index \(s\) with \(\lvert x_s\rvert\geqslant\lvert x_k\rvert\) for
\(k=1,\dots,n\); then \(\lvert x_s\rvert>0\).
\begin{itemize}
\item Equation \(s\) of \(Ax=0\) reads \(\sum_ka_{sk}x_k=0\), that is, \(a_{ss}x_s=-\sum_{k\ne s}a_{sk}x_k\).
\item By the triangle inequality, the choice of \(s\), and the hypothesis for \(i=s\) (with \(\lvert x_s\rvert>0\)),
\[
\lvert a_{ss}\rvert\,\lvert x_s\rvert=\Bigl\lvert\sum_{k\ne s}a_{sk}x_k\Bigr\rvert
\leqslant\sum_{k\ne s}\lvert a_{sk}\rvert\,\lvert x_k\rvert\leqslant\lvert x_s\rvert\sum_{k\ne s}\lvert a_{sk}\rvert
<\lvert x_s\rvert\,\lvert a_{ss}\rvert\ \contradiction
\]
\end{itemize}
Hence \(Ax=0\) only for \(x=0\), and \(A\) is invertible.
\end{solution}

\begin{problem}[\citeat{prasolov}{Problems~21.1, 21.2}]\label{r2-ex-pr-skew-inverse}
Let \(A\in\R^{n\times n}\) be skew-symmetric (Definition~\ref{r1-def-skew}).
\begin{enumerate}
\item Show that \(I_n+A\) is invertible.
\item Show that if \(A\) is invertible, then \(A^{-1}\) is skew-symmetric.
\end{enumerate}
\end{problem}
\begin{solution}
(a) By Theorem~\ref{r2-thm-invertible} ((iii)\,\(\Rightarrow\)\,(i)), it suffices to show that \((I_n+A)x=0\) only for
\(x=0\).
\begin{itemize}
\item For every \(x\in\R^n\), the \(1\times1\) matrix \(x'Ax\) equals its transpose, so by
Theorem~\ref{r1-thm-transpose} and \(A'=-A\), \(x'Ax=(x'Ax)'=x'A'x=-x'Ax\). Hence \(x'Ax=0\).
\item If \((I_n+A)x=0\), then by Theorems~\ref{r1-thm-rules} and~\ref{r1-thm-transpose},
\(0=x'(I_n+A)x=x'x+x'Ax=x'x=\sum_ix_i^2\), so \(x=0\).
\end{itemize}
(b)
\begin{itemize}
\item \(-A\) is invertible with \((-A)^{-1}=-A^{-1}\), since \((-A)(-A^{-1})=AA^{-1}=I_n\) and
\((-A^{-1})(-A)=A^{-1}A=I_n\) (Theorem~\ref{r1-thm-rules}, Definition~\ref{r2-def-inverse}).
\item By Theorem~\ref{r2-thm-inverse-props}(c) and \(A'=-A\), \((A^{-1})'=(A')^{-1}=(-A)^{-1}=-A^{-1}\). So \(A^{-1}\)
is skew-symmetric.
\end{itemize}
\end{solution}

\begin{problem}[\citeat{prasolov}{Problem~2.12}]\label{r2-ex-pr-nonneg-inverse}
Find all invertible matrices \(A\in\R^{n\times n}\) such that all entries of \(A\) and all entries of \(A^{-1}\) are
nonnegative.
\end{problem}
\begin{solution}
These are exactly the matrices \(A=(a_{ij})\) with nonnegative entries that have exactly one nonzero entry in each row
and in each column; the inverse of such an \(A\) is \(A'\) with every nonzero entry replaced by its reciprocal.
\begin{itemize}
\item Rows. Let \(A\) be invertible, with nonnegative entries, and let \(B=A^{-1}=(b_{ij})\) have nonnegative entries.
Suppose that row \(i\) of \(A\) has two positive entries, \(a_{ir}>0\) and \(a_{is}>0\) with \(r\ne s\). For \(j\ne i\),
the \((i,j)\) entry of \(AB=I_n\) is \(\sum_ka_{ik}b_{kj}=0\); its terms are \(\geqslant0\), so all of them vanish, and
\(a_{ir}b_{rj}=0\), \(a_{is}b_{sj}=0\) give \(b_{rj}=b_{sj}=0\). Thus \(\row_r(B)=b_{ri}e_i'\) and
\(\row_s(B)=b_{si}e_i'\): rows \(r\) and \(s\) of \(B\) are multiples of \(e_i'\).
\item This is impossible. By Theorem~\ref{r1-thm-views}(c), \(\row_r(BA)=\row_r(B)A=b_{ri}\row_i(A)\) and
\(\row_s(BA)=b_{si}\row_i(A)\), and as \(BA=I_n\), these rows are \(e_r'\) and \(e_s'\). So \(b_{ri}\ne0\), and
\(e_s'=b_{si}\row_i(A)=(b_{si}/b_{ri})\,e_r'\); but the \(s\)th entry of \(e_s'\) is \(1\), while that of
\((b_{si}/b_{ri})\,e_r'\) is \(0\), as \(r\ne s\) \(\contradiction\) Hence every row of \(A\) has at most one nonzero
entry.
\item No row of \(A\) is zero: if \(\row_i(A)=0\), then \(\row_i(AB)=\row_i(A)B=0\), while \(\row_i(I_n)=e_i'\ne0\)
(Theorem~\ref{r1-thm-views}(c)). So every row of \(A\) has exactly one nonzero entry.
\item Columns. \(A'\) has nonnegative entries and is invertible, with \((A')^{-1}=(A^{-1})'=B'\), which has nonnegative
entries (Theorem~\ref{r2-thm-inverse-props}(c)). By the first three items, applied to \(A'\), every row of \(A'\), that
is, every column of \(A\), has exactly one nonzero entry.
\item Converse. Let \(A\) have nonnegative entries, with exactly one nonzero entry in each row and in each column, and
let \(a_{i\sigma(i)}\) be the nonzero (hence positive) entry of row \(i\). If \(i\ne l\), then
\(\sigma(i)\ne\sigma(l)\), as column \(\sigma(i)\) has only one nonzero entry. Let \(B=(b_{kj})\in\R^{n\times n}\) have
the entries \(b_{\sigma(i)i}=1/a_{i\sigma(i)}\) (\(i=1,\dots,n\)) and all other entries \(0\); it is \(A'\) with every
nonzero entry replaced by its reciprocal, and its entries are nonnegative.
\item In \((AB)_{ij}=\sum_ka_{ik}b_{kj}\) only the term \(k=\sigma(i)\) can be nonzero, so
\((AB)_{ij}=a_{i\sigma(i)}b_{\sigma(i)j}\). Here \(b_{\sigma(i)j}\ne0\) only if \(\sigma(i)=\sigma(j)\), that is,
\(j=i\), and \(a_{i\sigma(i)}b_{\sigma(i)i}=1\). So \(AB=I_n\), and by Theorem~\ref{r2-thm-invertible}, \(A\) is
invertible with \(A^{-1}=B\), whose entries are nonnegative.
\end{itemize}
\end{solution}

\clearpage
\def\unitnumber{3}
\setcounter{theorem}{0}\setcounter{problem}{0}
\section*{Chapter 3. Vector spaces, subspaces, span, and bases}

\subsection*{Vector spaces}

\begin{definition}[\citeat{hoffman-kunze}{\S2.1}, \citeat{lang}{Ch.~I, \S1}, \citeat{searle-khuri}{Definition~1.1}, \citeat{bapat}{\S1.2}]\label{r3-def-vector-space}
A (real) \textbf{vector space} is a set \(V\), whose elements are called \textbf{vectors}, with
an \textbf{addition} that gives \(u+v\in V\) for \(u,v\in V\) and a \textbf{scalar
multiplication} that gives \(cv\in V\) for \(c\in\R\) and \(v\in V\), such that for all
\(u,v,w\in V\) and \(c,d\in\R\):
\begin{enumerate}[label=(V\arabic*)]
\item \(u+v=v+u\) and (V2) \(u+(v+w)=(u+v)+w\);
\setcounter{enumi}{2}
\item there is a \textbf{zero vector} \(0\in V\) with \(v+0=v\) for every \(v\in V\);
\item for each \(v\in V\) there is \(-v\in V\) with \(v+(-v)=0\);
\item \(1v=v\) and (V6) \((cd)v=c(dv)\);
\setcounter{enumi}{6}
\item \(c(u+v)=cu+cv\) and (V8) \((c+d)v=cv+dv\).
\end{enumerate}
A vector \(c_1v_1+\dots+c_kv_k\) (\(c_i\in\R\)) is a \textbf{linear combination} of
\(v_1,\dots,v_k\).
\end{definition}
From these rules, \(0v=0\) and \((-1)v=-v\): indeed \(0v=(0+0)v=0v+0v\), and adding \(-(0v)\)
to both sides gives \(0=0v\); then \(v+(-1)v=(1+(-1))v=0v=0\), and \(-v\) is the only vector \(w\)
with \(v+w=0\), since such a \(w\) equals \(w+(v+(-v))=(w+v)+(-v)=0+(-v)=-v\). Likewise \(c0=0\), and
\(cv=0\) only if \(c=0\) or \(v=0\) (if \(c\ne0\), then \(v=c^{-1}(cv)=c^{-1}0=0\)).

\begin{example}[\citeat{hoffman-kunze}{\S2.1, Example~2}]\label{r3-exm-spaces}
\(\R^n\), \(\R^{1\times n}\) and \(\R^{m\times n}\), with the entrywise sum and scalar multiple
of Definition~\ref{r1-def-matrix}, are vector spaces: (V1)--(V8) hold entry by entry, the zero vector is the zero
matrix, and \(-A=(-1)A\). In this course every vector space is one of these or a subspace of
one of them.
\end{example}

\subsection*{Subspaces}

\begin{definition}[\citeat{hoffman-kunze}{\S2.2}, \citeat{searle-khuri}{Definition~1.2}]\label{r3-def-subspace}
A \textbf{subspace} of a vector space \(V\) is a subset \(W\subseteq V\) that is itself a
vector space with the addition and scalar multiplication of \(V\).
\end{definition}

\begin{theorem}[\citeat{hoffman-kunze}{\S2.2, Theorem~1}]\label{r3-thm-subspace-test}
For a subset \(W\) of a vector space \(V\) the following are equivalent:
\begin{enumerate}
\item \(W\) is a subspace of \(V\);
\item \(0\in W\), and \(u+v\in W\) and \(cu\in W\) for all \(u,v\in W\) and \(c\in\R\)
(\(W\) is \textbf{closed under addition and scalar multiplication});
\item \(0\in W\), and \(cu+v\in W\) for all \(u,v\in W\) and \(c\in\R\).
\end{enumerate}
\end{theorem}
\begin{proof}
(a)\,\(\Rightarrow\)\,(b): the operations of the vector space \(W\) are those of \(V\) and give
vectors of \(W\); \(W\) has a zero vector, so it contains some \(w\), and then
\(0=0w\in W\). (b)\,\(\Rightarrow\)\,(c): \(cu\in W\), so \(cu+v\in W\).
(c)\,\(\Rightarrow\)\,(a): taking \(v=0\) gives \(cu\in W\), and \(c=1\) gives \(u+v\in W\);
so the operations of \(V\) give operations on \(W\). The rules (V1), (V2), (V5)--(V8) hold in
\(W\) because they hold in \(V\); (V3) holds since \(0\in W\), and (V4) since
\(-w=(-1)w\in W\).
\end{proof}
So, to show that \(W\) is a subspace, check that it contains \(0\) and is closed under the two
operations; to show that it is not, exhibit one failure.

\begin{example}[\citeat{hoffman-kunze}{\S2.2, Examples~6, 7}, \citeat{searle-khuri}{\S1.1}, \citeat{abadir-magnus}{Exercise~3.10(b)}]\label{r3-exm-subspaces}
\begin{enumerate}
\item \(\{0\}\) and \(V\) itself are subspaces of \(V\).
\item In \(\R^n\), \(W=\{x:x_1=0\}\) is a subspace: \(0\in W\), and
\((cx+y)_1=cx_1+y_1=0\) for \(x,y\in W\). For \(n\geqslant2\), \(\{x:x_1=1+x_2\}\) is not,
since it does not contain \(0\).
\item In \(\R^2\), \(\{(x,y)':2x-5y=0\}\) is a subspace: it contains \(0\), and
\(2(cx_1+x_2)-5(cy_1+y_2)=c(2x_1-5y_1)+(2x_2-5y_2)\) for \((x_1,y_1)'\) and \((x_2,y_2)'\) in
it; but for \(k\ne0\) the set
\(\{(x,y)':ax+by=k\}\) is not, since it does not contain \(0\).
\item The symmetric \(n\times n\) matrices form a subspace of \(\R^{n\times n}\): \(0'=0\),
and \((cA+B)'=cA'+B'=cA+B\) for symmetric \(A,B\). So do the skew-symmetric matrices
(Definition~\ref{r1-def-skew}), those with \(A'=-A\): \(0'=-0\), and \((cA+B)'=cA'+B'=-(cA+B)\) if \(A'=-A\) and
\(B'=-B\).
\item For \(A\in\R^{m\times n}\), the solution set \(\{x\in\R^n:Ax=0\}\) of a homogeneous
system is a subspace of \(\R^n\): \(A0=0\), and \(A(cu+v)=cAu+Av=0\) if \(Au=Av=0\). For
\(b\ne0\), the solutions of \(Ax=b\) do not form a subspace (\(A0=0\ne b\)).
\item In statistics, the \textbf{centered} vectors \(W_0=\{x\in\R^n:\one'x=0\}\), whose
entries add up to \(0\), form a subspace: this is (e) with the \(1\times n\) matrix
\(A=\one'\). For every \(x\), the deviation vector \(Cx\) of Example~\ref{r1-exm-centering} lies in \(W_0\),
since \(\one'Cx=\one'C'x=(C\one)'x=0\) (Proposition~\ref{r1-prop-centering}: \(C'=C\), \(C\one=0\)).
\item The set \(\{(x,y,z)'\in\R^3:xy=0\}\) is not a subspace: \((1,0,0)'\) and
\((0,1,0)'\) lie in it, but their sum \((1,1,0)'\) does not. It is the union of the
subspaces \(\{x=0\}\) and \(\{y=0\}\); a union of subspaces need not be a subspace.
\end{enumerate}
\end{example}

\subsection*{Span}

\begin{definition}[\citeat{lang}{Ch.~I, \S1}]\label{r3-def-span}
Let \(v_1,\dots,v_k\) be vectors of a vector space \(V\). Their \textbf{span} is the set of all
their linear combinations,
\[
\Span\{v_1,\dots,v_k\}=\{c_1v_1+\dots+c_kv_k:\ c_1,\dots,c_k\in\R\}
\]
and \(\Span\varnothing=\{0\}\). If \(\Span\{v_1,\dots,v_k\}=W\), we say that \(v_1,\dots,v_k\)
\textbf{span} \(W\).
\end{definition}

\begin{theorem}[\citeat{hoffman-kunze}{\S2.2, Theorem~3}]\label{r3-thm-span}
\(\Span\{v_1,\dots,v_k\}\) is a subspace of \(V\) that contains \(v_1,\dots,v_k\), and it is
contained in every subspace that contains \(v_1,\dots,v_k\): it is the smallest subspace
containing them.
\end{theorem}
\begin{proof}
\(0=0v_1+\dots+0v_k\) lies in the span, and
\(c\sum_ia_iv_i+\sum_ib_iv_i=\sum_i(ca_i+b_i)v_i\); so the span is a subspace
(Theorem~\ref{r3-thm-subspace-test}(c)). It contains \(v_j=1v_j+\sum_{i\ne j}0v_i\). A subspace
that contains \(v_1,\dots,v_k\) contains every \(c_iv_i\) and every sum of these, by closure.
\end{proof}

For vectors \(v_1,\dots,v_k\in\R^m\), put \(V=[v_1\ \cdots\ v_k]\in\R^{m\times k}\). Since
\(Vc=c_1v_1+\dots+c_kv_k\) (Theorem~\ref{r1-thm-views}(b)),
\[
\Span\{v_1,\dots,v_k\}=\{Vc:\ c\in\R^k\},\qquad
b\in\Span\{v_1,\dots,v_k\}\iff\text{the system }Vc=b\text{ is consistent}
\]
So membership in a span is decided by row-reducing \([V\mid b]\)
(Proposition~\ref{r2-prop-solutions}, Theorem~\ref{thm:r2-consistency}). For instance, by
Example~\ref{r2-exm-consistency}, a vector \(y\in\R^4\) lies in the span of the columns of the
matrix \(A\) there if and only if \(y_1-y_2-y_3+y_4=0\).

\begin{example}[\citeat{hoffman-kunze}{\S2.2, Example~8}]\label{r3-exm-span}
Let \(W=\Span\{v_1,v_2,v_3\}\subseteq\R^5\), where \(v_1=(1,2,0,3,0)'\),
\(v_2=(0,0,1,4,0)'\) and \(v_3=(0,0,0,0,1)'\). Which \(y=(y_1,\dots,y_5)'\) lie in \(W\)?

\textit{Solution.} Row-reduce \([v_1\ v_2\ v_3\mid y]\), keeping the \(y_i\) as letters:
\[
\left(\begin{array}{ccc|c}1&0&0&y_1\\2&0&0&y_2\\0&1&0&y_3\\3&4&0&y_4\\0&0&1&y_5\end{array}\right)
\xrightarrow[R_4\leftarrow R_4-3R_1]{R_2\leftarrow R_2-2R_1}
\left(\begin{array}{ccc|c}1&0&0&y_1\\0&0&0&y_2-2y_1\\0&1&0&y_3\\0&4&0&y_4-3y_1\\0&0&1&y_5\end{array}\right)
\xrightarrow{R_4\leftarrow R_4-4R_3}
\left(\begin{array}{ccc|c}1&0&0&y_1\\0&0&0&y_2-2y_1\\0&1&0&y_3\\0&0&0&y_4-3y_1-4y_3\\0&0&1&y_5\end{array}\right)
\]
\[
\xrightarrow{R_2\leftrightarrow R_3}
\left(\begin{array}{ccc|c}1&0&0&y_1\\0&1&0&y_3\\0&0&0&y_2-2y_1\\0&0&0&y_4-3y_1-4y_3\\0&0&1&y_5\end{array}\right)
\xrightarrow{R_3\leftrightarrow R_5}
\left(\begin{array}{ccc|c}1&0&0&y_1\\0&1&0&y_3\\0&0&1&y_5\\0&0&0&y_4-3y_1-4y_3\\0&0&0&y_2-2y_1\end{array}\right)
\]
By Proposition~\ref{r2-prop-solutions}, \(y\in W\) if and only if
\[
y_2=2y_1,\qquad y_4=3y_1+4y_3
\]
and then \(y=y_1v_1+y_3v_2+y_5v_3\). For example, \((-3,-6,1,-5,2)'\in W\), since
\(-6=2(-3)\) and \(-5=3(-3)+4\cdot1\); indeed it equals \(-3v_1+v_2+2v_3\). But
\((2,4,6,7,8)'\notin W\), since \(3\cdot2+4\cdot6=30\ne7\).
\end{example}

\begin{example}\label{r3-exm-constant}
\(\Span\{\one_n\}=\{c\one_n:c\in\R\}\) is the subspace of constant vectors. Every
\(x\in\R^n\) is the sum of a constant vector and a centered vector:
\[
x=\bar x\one+Cx,\qquad \bar x\one\in\Span\{\one\},\quad Cx\in W_0
\]
since \(Cx=x-\bar x\one\) (Example~\ref{r1-exm-centering}) and \(Cx\in W_0\)
(Example~\ref{r3-exm-subspaces}(f)). For \(x=(1,3,5,7,9)'\):
\(x=5\one+(-4,-2,0,2,4)'\).
\end{example}

\begin{theorem}[\citeat{harville}{Lemmas~4.2.1, 4.2.2, 4.3.1}, \citeat{harville-solutions}{Ch.~4, Exercise~6}]\label{r3-thm-span-inclusion}
Let \(A=[a_1\ \cdots\ a_n]\in\R^{m\times n}\) and \(B=[b_1\ \cdots\ b_p]\in\R^{m\times p}\).
\begin{enumerate}
\item \(\Span\{b_1,\dots,b_p\}\subseteq\Span\{a_1,\dots,a_n\}\) if and only if every \(b_j\) lies in
\(\Span\{a_1,\dots,a_n\}\), and if and only if \(B=AF\) for some \(F\in\R^{n\times p}\).
\item \(\Span\{b_1,\dots,b_p\}=\Span\{a_1,\dots,a_n\}\) if and only if \(B=AF\) and \(A=BG\) for some
\(F\in\R^{n\times p}\) and \(G\in\R^{p\times n}\).
\item If every \(b_j\) lies in \(\Span\{a_1,\dots,a_n\}\), then
\(\Span\{a_1,\dots,a_n,b_1,\dots,b_p\}=\Span\{a_1,\dots,a_n\}\). In particular, deleting from a list of
vectors a vector that is a linear combination of the others does not change their span.
\end{enumerate}
\end{theorem}
\begin{proof}
(a)
\begin{itemize}
\item If \(\Span\{b_1,\dots,b_p\}\subseteq\Span\{a_1,\dots,a_n\}\), then each \(b_j\), which lies in
\(\Span\{b_1,\dots,b_p\}\) (Theorem~\ref{r3-thm-span}), lies in \(\Span\{a_1,\dots,a_n\}\).
\item If every \(b_j\) lies in \(\Span\{a_1,\dots,a_n\}\), this subspace contains \(b_1,\dots,b_p\), so it
contains \(\Span\{b_1,\dots,b_p\}\) (Theorem~\ref{r3-thm-span}).
\item For \(F=[f_1\ \cdots\ f_p]\in\R^{n\times p}\), \(\col_j(AF)=Af_j\) (Theorem~\ref{r1-thm-views}(a)), and
\(Af_j\) is the combination of \(a_1,\dots,a_n\) with the coefficients \(f_j\) (Theorem~\ref{r1-thm-views}(b)).
So \(B=AF\) for some \(F\) if and only if each \(b_j\) is a linear combination of \(a_1,\dots,a_n\).
\end{itemize}
(b) The two spans are equal if and only if each is contained in the other; apply (a) to both
inclusions.

(c)
\begin{itemize}
\item Let \(b_j=k_{1j}a_1+\dots+k_{nj}a_n\) (\(1\leqslant j\leqslant p\)). For all scalars \(x_i\) and
\(y_j\),
\[
\sum_{i=1}^nx_ia_i+\sum_{j=1}^py_jb_j=\sum_{i=1}^n\Bigl(x_i+\sum_{j=1}^pk_{ij}y_j\Bigr)a_i
\]
so every linear combination of \(a_1,\dots,a_n,b_1,\dots,b_p\) is one of \(a_1,\dots,a_n\).
\item Conversely, \(\sum_ix_ia_i=\sum_ix_ia_i+\sum_j0\,b_j\).
\item For the last sentence, let the \(a_i\) be the remaining vectors and \(b_1\) the deleted one
(\(p=1\)); the order of the vectors does not affect their span.\qedhere
\end{itemize}
\end{proof}

\subsection*{Linear independence}

\begin{definition}[\citeat{hoffman-kunze}{\S2.3}, \citeat{lang}{Ch.~I, \S2}, \citeat{searle-khuri}{Definition~1.3}, \citeat{bapat}{\S1.3}]\label{r3-def-independence}
Vectors \(v_1,\dots,v_k\) of a vector space are \textbf{linearly dependent} if
\(c_1v_1+\dots+c_kv_k=0\) for some scalars \(c_1,\dots,c_k\), not all \(0\) (a \textbf{linear
relation}). Otherwise they are \textbf{linearly independent}: \(c_1v_1+\dots+c_kv_k=0\) only
for \(c_1=\dots=c_k=0\). The empty list is independent.
\end{definition}

\begin{proposition}[cf.\ \citeat{hoffman-kunze}{\S2.3, Example~14}]\label{r3-prop-independence-test}
Let \(v_1,\dots,v_k\in\R^m\) and \(V=[v_1\ \cdots\ v_k]\in\R^{m\times k}\). Then
\(v_1,\dots,v_k\) are linearly independent if and only if \(Vc=0\) has only the solution
\(c=0\), that is, if and only if every column of the RREF of \(V\) is a pivot column (no free
unknowns). In particular, more than \(m\) vectors of \(\R^m\) are linearly dependent.
\end{proposition}
\begin{proof}
\(Vc=c_1v_1+\dots+c_kv_k\), so the linear relations among the \(v_i\) are the solutions of
\(Vc=0\). This system is consistent (\(c=0\)), so by
Theorem~\ref{thm:r2-solution-structure}(b) \(c=0\) is its only solution exactly when there is no
free unknown. If \(k>m\), then \(Vc=0\) has a nontrivial solution (Theorem~\ref{r2-thm-homogeneous}(a)).
\end{proof}

\begin{proposition}[\citeat{hoffman-kunze}{\S2.3}, \citeat{harville}{\S3.2, Lemma~3.2.1}]\label{r3-prop-dependence}
Let \(v_1,\dots,v_k\) be vectors of a vector space.
\begin{enumerate}
\item If some \(v_i=0\), or \(v_i=v_j\) for some \(i\ne j\), then \(v_1,\dots,v_k\) are
dependent. A single vector \(v\) is independent if and only if \(v\ne0\).
\item A sublist of an independent list is independent.
\item For \(k\geqslant2\), \(v_1,\dots,v_k\) are dependent if and only if some \(v_j\) is a
linear combination of the others. In particular, two vectors are dependent if and only if one
of them is a multiple of the other.
\end{enumerate}
\end{proposition}
\begin{proof}
(a) \(1\cdot0=0\) and \(v_i-v_j=0\) are linear relations with a coefficient \(1\); and if
\(v\ne0\), then \(cv=0\) forces \(c=0\). (b) A relation among some of the vectors is a
relation among all of them, with coefficient \(0\) for the others. (c) If
\(v_j=\sum_{i\ne j}x_iv_i\), then \(v_j-\sum_{i\ne j}x_iv_i=0\) is a relation with coefficient
\(1\) on \(v_j\). Conversely, if \(\sum_ix_iv_i=0\) with \(x_j\ne0\), then
\(v_j=\sum_{i\ne j}(-x_i/x_j)v_i\).
\end{proof}

\begin{proposition}[\citeat{harville}{Lemma~3.2.2}, \citeat{hoffman-kunze}{\S2.3, Lemma before Theorem~5}]\label{r3-prop-adjoin}
Let \(v_1,\dots,v_k\) be vectors of a vector space \(V\).
\begin{enumerate}
\item If \(k\geqslant2\) and \(v_1\ne0\), then \(v_1,\dots,v_k\) are linearly dependent if and only if
some \(v_j\) (\(2\leqslant j\leqslant k\)) is a linear combination of the preceding vectors
\(v_1,\dots,v_{j-1}\).
\item If \(v_1,\dots,v_k\) are linearly independent and \(v\in V\) does not lie in
\(\Span\{v_1,\dots,v_k\}\), then \(v_1,\dots,v_k,v\) are linearly independent.
\end{enumerate}
\end{proposition}
\begin{proof}
(a)
\begin{itemize}
\item If \(v_j=\sum_{i<j}x_iv_i\), then \(v_j\) is a combination of the other vectors (with
coefficient \(0\) for \(i>j\)), so the list is dependent (Proposition~\ref{r3-prop-dependence}(c)).
\item Conversely, let the list be dependent, and let \(j\) be the least index such that
\(v_1,\dots,v_j\) are dependent. Since \(v_1\ne0\) is independent
(Proposition~\ref{r3-prop-dependence}(a)), \(j\geqslant2\).
\item There are \(x_1,\dots,x_j\), not all \(0\), with \(x_1v_1+\dots+x_jv_j=0\). Here \(x_j\ne0\),
since otherwise \(v_1,\dots,v_{j-1}\) would be dependent, contrary to the choice of \(j\).
\item Hence \(v_j=(-x_1/x_j)v_1+\dots+(-x_{j-1}/x_j)v_{j-1}\).
\end{itemize}
(b)
\begin{itemize}
\item Let \(c_1v_1+\dots+c_kv_k+bv=0\). Then \(b=0\): otherwise
\(v=(-c_1/b)v_1+\dots+(-c_k/b)v_k\) would lie in \(\Span\{v_1,\dots,v_k\}\).
\item So \(c_1v_1+\dots+c_kv_k=0\), and every \(c_i=0\), as \(v_1,\dots,v_k\) are
independent.\qedhere
\end{itemize}
\end{proof}

\begin{example}\label{r3-exm-dependent}
Let \(a_1,a_2,a_3,a_4\) be the columns of the matrix \(A\) of Example~\ref{r2-exm-gauss-jordan}:
\[
a_1=\begin{pmatrix}1\\1\\2\end{pmatrix},\quad a_2=\begin{pmatrix}1\\1\\1\end{pmatrix},\quad
a_3=\begin{pmatrix}1\\2\\5\end{pmatrix},\quad a_4=\begin{pmatrix}-2\\-3\\-4\end{pmatrix}
\]
Four vectors of \(\R^3\) are dependent (Proposition~\ref{r3-prop-independence-test}). A relation
comes from the solutions \(x=t(-2,3,1,1)'\) of \(Ax=0\) found there:
\[
-2a_1+3a_2+a_3+a_4=0,\qquad\text{so}\qquad a_4=2a_1-3a_2-a_3
\]
On the other hand \(a_1,a_2,a_3\) are independent: the same row operations take
\([a_1\ a_2\ a_3]\) to the first three columns of the RREF found there (row operations act on each column separately:
Theorems~\ref{r2-thm-elementary}(c) and~\ref{r1-thm-views}(a)), which form \(I_3\), so
every column is a pivot column.
\end{example}

\begin{example}[\citeat{hoffman-kunze}{\S2.3, Exercises~2, 4}]\label{r3-exm-adjoin}
\begin{enumerate}
\item Let \(u_1=(1,1,2,4)'\), \(u_2=(2,-1,-5,2)'\), \(u_3=(1,-1,-4,0)'\) and \(u_4=(2,1,1,6)'\). Here \(u_1\ne0\), and
\(u_2\) is not a linear combination of the preceding vector \(u_1\), that is, not a multiple of \(u_1\): the only
multiple \(cu_1\) with first entry \(2\) is \(2u_1\), whose second entry is \(2\ne-1\). So \(u_1,u_2\) are independent
(Proposition~\ref{r3-prop-adjoin}(a) with \(k=2\)). But \(u_3\) is a linear combination of the preceding vectors:
\[
-\tfrac13u_1+\tfrac23u_2=\tfrac13\begin{pmatrix}-1+4\\-1-2\\-2-10\\-4+4\end{pmatrix}=\begin{pmatrix}1\\-1\\-4\\0\end{pmatrix}=u_3
\]
By Proposition~\ref{r3-prop-adjoin}(a) (\(j=3\); with \(k=3\) and with \(k=4\)), \(u_1,u_2,u_3\) and \(u_1,\dots,u_4\)
are linearly dependent. Since the single vector \(u_1\ne0\) is independent (Proposition~\ref{r3-prop-dependence}(a)),
and so are \(u_1,u_2\), \(3\) is the least \(j\) for which \(u_1,\dots,u_j\) are dependent.
\item Let \(v_1=(1,0,-1)'\), \(v_2=(1,2,1)'\) and \(v_3=(0,-3,2)'\). The single vector \(v_1\ne0\) is independent
(Proposition~\ref{r3-prop-dependence}(a)).
\begin{itemize}
\item \(v_2\notin\Span\{v_1\}\): every multiple of \(v_1\) has second entry \(0\), and \(v_2\) has second entry \(2\).
So \(v_1,v_2\) are independent (Proposition~\ref{r3-prop-adjoin}(b)).
\item \(v_3\notin\Span\{v_1,v_2\}\): if \(v_3=xv_1+yv_2\), the second entries give \(2y=-3\), so \(y=-\frac32\); the
first entries give \(x+y=0\), so \(x=\frac32\); but then the third entry of \(xv_1+yv_2\) is \(-x+y=-3\ne2\). So
\(v_1,v_2,v_3\) are independent (Proposition~\ref{r3-prop-adjoin}(b) again).
\end{itemize}
\end{enumerate}
\end{example}

\subsection*{Bases}

\begin{definition}[\citeat{hoffman-kunze}{\S2.3}, \citeat{lang}{Ch.~I, \S2}]\label{r3-def-basis}
Vectors \(v_1,\dots,v_k\) of a vector space \(V\) form a \textbf{basis} of a subspace \(W\) if
they lie in \(W\), are linearly independent, and span \(W\).
\begin{itemize}
\item The unit vectors \(e_1,\dots,e_n\) form a basis of \(\R^n\), the \textbf{standard basis}:
\(x=x_1e_1+\dots+x_ne_n\) for every \(x\in\R^n\), and this is \(0\) only for \(x=0\).
\item Let \(E_{ij}\in\R^{m\times n}\) have \((i,j)\) entry \(1\) and all other entries \(0\)
(two indices; not the elementary matrices \(E_1,E_2,\dots\) of Theorem~\ref{r2-thm-elementary}(c)). The \(mn\) matrices
\(E_{ij}\) form a basis of \(\R^{m\times n}\): \(A=\sum_i\sum_ja_{ij}E_{ij}\), and
\(\sum_i\sum_jc_{ij}E_{ij}=(c_{ij})\) is \(0\) only if every \(c_{ij}=0\).
\item The empty list is a basis of \(\{0\}\).
\end{itemize}
\end{definition}

\begin{theorem}[\citeat{lang}{Ch.~I, Theorem~2.1}]\label{r3-thm-unique-coef}
If \(v_1,\dots,v_k\) form a basis of \(W\), every \(w\in W\) can be written as
\(w=c_1v_1+\dots+c_kv_k\) in exactly one way; the numbers \(c_1,\dots,c_k\) are the
\textbf{coordinates} of \(w\) in this basis.
\end{theorem}
\begin{proof}
Such an expression exists because the \(v_i\) span \(W\). If
\(\sum_ix_iv_i=\sum_iy_iv_i\), then \(\sum_i(x_i-y_i)v_i=0\), and independence gives
\(x_i=y_i\) for every \(i\).
\end{proof}

\begin{example}[\citeat{hoffman-kunze}{\S2.3, Example~14}]\label{r3-exm-invertible-basis}
If \(P\in\R^{n\times n}\) is invertible, its columns \(p_1,\dots,p_n\) form a basis of \(\R^n\),
and the coordinates of \(y\in\R^n\) are \(x=P^{-1}y\). Indeed \(Px=\sum_jx_jp_j\); \(Px=0\) only
for \(x=0\) (Theorem~\ref{r2-thm-invertible}), so the columns are independent; and
\(y=P(P^{-1}y)\), so they span. For example, the matrix \(A\) of Example~\ref{r2-exm-inverse-gj} is invertible, its
inverse being computed there; so its columns \((1,-1,2,1)'\), \((-1,1,-1,-1)'\), \((1,-2,2,1)'\),
\((-1,1,-1,0)'\) form a basis of \(\R^4\).
The coordinates of \(y=(1,0,3,3)'\) are \(A^{-1}y=(3,-1,-1,2)'\), the solution of \(Ax=y\) found in
Example~\ref{r2-exm-unique}:
\(y=3(1,-1,2,1)'-(-1,1,-1,-1)'-(1,-2,2,1)'+2(-1,1,-1,0)'\).
\end{example}

\begin{example}[cf.\ \citeat{searle-khuri}{Example~1.6}]\label{r3-exm-two-of-three}
Let \(t=(1,2,3,4)'\) and \(W=\Span\{\one_4,t\}\), a subspace of \(\R^4\). Any two of the vectors
\(\one_4\), \(t\), \(\one_4+t=(2,3,4,5)'\) form a basis of \(W\).
\begin{itemize}
\item Each pair is independent: none of the three vectors is \(0\) or a multiple of another,
since their first two entries are \((1,1)\), \((1,2)\), \((2,3)\), no two of them proportional
(Proposition~\ref{r3-prop-dependence}(c)).
\item Each pair spans \(W\): the pair \(\one_4,t\) by definition, and for \(w=a\one_4+bt\in W\),
\[
w=(a-b)\,\one_4+b\,(\one_4+t)=(b-a)\,t+a\,(\one_4+t)
\]
\end{itemize}
All three vectors together also span \(W\), but they are dependent:
\(\one_4+t-(\one_4+t)=0\) is a linear relation with coefficients \(1,1,-1\). So they do not form
a basis. A subspace has many bases.
\end{example}

\begin{example}[\citeat{harville}{\S4.3b}]\label{r3-exm-symmetric-basis}
The symmetric \(2\times2\) matrices (for instance \(2\times2\) covariance matrices) form a
subspace of \(\R^{2\times2}\) (Example~\ref{r3-exm-subspaces}(d)). The matrices
\(E_{11}\), \(E_{12}+E_{21}\), \(E_{22}\) form a basis of it: they are symmetric, every
symmetric matrix is
\[
\begin{pmatrix}a&b\\b&d\end{pmatrix}=a\begin{pmatrix}1&0\\0&0\end{pmatrix}+b\begin{pmatrix}0&1\\1&0\end{pmatrix}+d\begin{pmatrix}0&0\\0&1\end{pmatrix}
\]
and this combination is \(0\) only if \(a=b=d=0\).
\end{example}

\begin{example}\label{r3-exm-centered-basis}
In \(W_0=\{x\in\R^4:x_1+x_2+x_3+x_4=0\}\), the vectors \(w_1=(1,-1,0,0)'\),
\(w_2=(0,1,-1,0)'\) and \(w_3=(0,0,1,-1)'\) form a basis. They lie in \(W_0\). They span
\(W_0\): for \(x\in W_0\), \(x_4=-x_1-x_2-x_3\), so
\[
x_1w_1+(x_1+x_2)w_2+(x_1+x_2+x_3)w_3=\bigl(x_1,\ x_2,\ x_3,\ -x_1-x_2-x_3\bigr)'=x
\]
They are independent: \(c_1w_1+c_2w_2+c_3w_3=(c_1,\ c_2-c_1,\ c_3-c_2,\ -c_3)'=0\) forces
\(c_1=0\), then \(c_2=0\), then \(c_3=0\).
\end{example}

\subsection*{Coordinates and change of basis}

\begin{definition}[\citeat{hoffman-kunze}{\S2.4}, \citeat{lang}{Ch.~I, \S2}]\label{r3-def-coordinates}
Let \(W\) be a subspace of a vector space \(V\). An \textbf{ordered basis} of \(W\) is a list
\(\mathcal B=(v_1,\dots,v_k)\) of vectors that form a basis of \(W\), taken in this order. Each
\(w\in W\) is \(w=c_1v_1+\dots+c_kv_k\) for exactly one choice of \(c_1,\dots,c_k\)
(Theorem~\ref{r3-thm-unique-coef}), and
\[
[w]_{\mathcal B}=(c_1,\dots,c_k)'\in\R^k
\]
is the \textbf{coordinate vector} of \(w\) relative to \(\mathcal B\). For the \textbf{standard ordered
basis} \(\mathcal E=(e_1,\dots,e_n)\) of \(\R^n\), \([x]_{\mathcal E}=x\), since
\(x=x_1e_1+\dots+x_ne_n\).
\end{definition}
If \(W\subseteq\R^n\) and \(\mathcal B\) consists of the columns of \(X=[v_1\ \cdots\ v_k]\in\R^{n\times k}\),
then \(Xc=c_1v_1+\dots+c_kv_k\) (Theorem~\ref{r1-thm-views}(b)); so \([w]_{\mathcal B}\) is the unique
solution \(c\) of the system \(Xc=w\), found by row-reducing \([X\mid w]\).

\begin{proposition}[\citeat{hoffman-kunze}{\S2.4}]\label{r3-prop-coordinates}
Let \(\mathcal B=(v_1,\dots,v_k)\) be an ordered basis of a subspace \(W\).
\begin{enumerate}
\item \([u+w]_{\mathcal B}=[u]_{\mathcal B}+[w]_{\mathcal B}\) and \([cw]_{\mathcal B}=c[w]_{\mathcal B}\)
for \(u,w\in W\) and \(c\in\R\).
\item Every \(c\in\R^k\) is the coordinate vector of exactly one vector of \(W\), namely
\(c_1v_1+\dots+c_kv_k\).
\end{enumerate}
\end{proposition}
\begin{proof}
(a) If \(u=\sum_ix_iv_i\) and \(w=\sum_iy_iv_i\), then \(u+w=\sum_i(x_i+y_i)v_i\) and
\(cw=\sum_i(cy_i)v_i\). By the uniqueness in Theorem~\ref{r3-thm-unique-coef}, these are the
expressions that give \([u+w]_{\mathcal B}\) and \([cw]_{\mathcal B}\).

(b) \(w=\sum_ic_iv_i\) lies in \(W\), and \([w]_{\mathcal B}=c\) by definition. If also
\([w_1]_{\mathcal B}=c\), then \(w_1=\sum_ic_iv_i=w\).
\end{proof}

\begin{example}[\citeat{hoffman-kunze}{\S2.2, Example~8}]\label{r3-exm-coordinates}
In Example~\ref{r3-exm-span}, the row operations shown take \([v_1\ v_2\ v_3]\) to the first three
columns of the last matrix (row operations act on each column separately,
Theorems~\ref{r2-thm-elementary}(c) and~\ref{r1-thm-views}(a)). These columns are \(e_1,e_2,e_3\in\R^5\);
they form the RREF of \([v_1\ v_2\ v_3]\), in which every column is a pivot column. So \(v_1,v_2,v_3\) are independent
(Proposition~\ref{r3-prop-independence-test}), and \(\mathcal B=(v_1,v_2,v_3)\) is an ordered basis of
\(W\). Every \(y\in W\) is \(y=y_1v_1+y_3v_2+y_5v_3\), so
\[
[y]_{\mathcal B}=(y_1,y_3,y_5)'
\]
the first, third and fifth entries of \(y\). For instance \([(-3,-6,1,-5,2)']_{\mathcal B}=(-3,1,2)'\),
and by Proposition~\ref{r3-prop-coordinates}(b), \((1,1,1)'\) is the coordinate vector of
\(v_1+v_2+v_3=(1,2,1,7,1)'\).
\end{example}

\begin{example}[\citeat{harville}{\S4.3b}]\label{r3-exm-symmetric-coordinates}
For \(1\leqslant j\leqslant i\leqslant n\) let \(U_{ii}=E_{ii}\) and \(U_{ij}=E_{ij}+E_{ji}\) (\(j<i\)),
with the matrices \(E_{ij}\in\R^{n\times n}\) of Definition~\ref{r3-def-basis}. Each \(U_{ij}\) is
symmetric. For \(k\geqslant l\), \(U_{kl}\) is the only one of these matrices with a nonzero
\((k,l)\) entry, and this entry is \(1\); for \(k<l\), the only one is \(U_{lk}\). Hence the \((k,l)\)
entry of \(\sum_{j\leqslant i}c_{ij}U_{ij}\) is \(c_{kl}\) for \(k\geqslant l\), and \(c_{lk}\) for
\(k<l\). So every symmetric \(A=(a_{ij})\) is
\[
A=\sum_{j=1}^n\sum_{i=j}^na_{ij}U_{ij}
\]
and \(\sum_{j\leqslant i}c_{ij}U_{ij}=0\) only if every \(c_{ij}=0\). Thus the matrices \(U_{ij}\)
(\(j\leqslant i\)) form a basis of the subspace of symmetric \(n\times n\) matrices
(Example~\ref{r3-exm-subspaces}(d)); for column \(j\) there are \(n-j+1\) of them, \(n(n+1)/2\) in all.
In the ordered basis \(\mathcal B=(U_{11},U_{21},\dots,U_{n1},U_{22},\dots,U_{n2},\dots,U_{nn})\), the
coordinate vector of \(A\) lists the entries of \(A\) on and below the diagonal, column by column. For instance, a
covariance matrix of four variables is determined by ten numbers, its variances and covariances:
\[
\left[\begin{pmatrix}4&2&-1&1\\2&5&0&-2\\-1&0&3&1\\1&-2&1&6\end{pmatrix}\right]_{\mathcal B}=(4,2,-1,1,5,0,-2,3,1,6)'
\]
For \(n=2\) this is Example~\ref{r3-exm-symmetric-basis}.
\end{example}

\begin{theorem}[\citeat{hoffman-kunze}{\S2.4, Theorem~7}]\label{r3-thm-change-basis}
Let \(\mathcal B=(v_1,\dots,v_k)\) and \(\tilde{\mathcal B}=(\tilde v_1,\dots,\tilde v_k)\) be ordered
bases of a subspace \(W\), each with \(k\) vectors, and let \(P=(p_{ij})\in\R^{k\times k}\) have \(j\)th
column \([\tilde v_j]_{\mathcal B}\), that is, \(\tilde v_j=p_{1j}v_1+\dots+p_{kj}v_k\)
(\(1\leqslant j\leqslant k\)).
\begin{enumerate}
\item \([w]_{\mathcal B}=P[w]_{\tilde{\mathcal B}}\) for every \(w\in W\), and \(P\) is the only
\(k\times k\) matrix with this property.
\item \(P\) is invertible, and \([w]_{\tilde{\mathcal B}}=P^{-1}[w]_{\mathcal B}\) for every \(w\in W\).
\end{enumerate}
\(P\) is the \textbf{change-of-basis matrix} from \(\tilde{\mathcal B}\) to \(\mathcal B\).
\end{theorem}
\begin{proof}
(a)
\begin{itemize}
\item Let \(w\in W\) and \(\tilde x=[w]_{\tilde{\mathcal B}}\). Then
\[
w=\sum_{j=1}^k\tilde x_j\tilde v_j=\sum_{j=1}^k\tilde x_j\sum_{i=1}^kp_{ij}v_i
=\sum_{i=1}^k\Bigl(\sum_{j=1}^kp_{ij}\tilde x_j\Bigr)v_i
\]
\item By the uniqueness of coordinates (Theorem~\ref{r3-thm-unique-coef}), the \(i\)th coordinate of
\(w\) relative to \(\mathcal B\) is \(\sum_jp_{ij}\tilde x_j=(P\tilde x)_i\). So
\([w]_{\mathcal B}=P\tilde x\).
\item If also \([w]_{\mathcal B}=Q[w]_{\tilde{\mathcal B}}\) for every \(w\), take \(w=\tilde v_j\), for
which \([\tilde v_j]_{\tilde{\mathcal B}}=e_j\): then \(\col_j(Q)=Qe_j=[\tilde v_j]_{\mathcal B}=\col_j(P)\)
(Theorem~\ref{r1-thm-views}(b)) for every \(j\), so \(Q=P\).
\end{itemize}
(b)
\begin{itemize}
\item Let \(P\tilde x=0\), and let \(w\in W\) be the vector with \([w]_{\tilde{\mathcal B}}=\tilde x\)
(Proposition~\ref{r3-prop-coordinates}(b)). By (a), \([w]_{\mathcal B}=0\), so \(w=0v_1+\dots+0v_k=0\),
and \(\tilde x=[0]_{\tilde{\mathcal B}}=0\).
\item So \(P\tilde x=0\) only for \(\tilde x=0\), and \(P\) is invertible
(Theorem~\ref{r2-thm-invertible}).
\item Then \([w]_{\tilde{\mathcal B}}=P^{-1}P[w]_{\tilde{\mathcal B}}=P^{-1}[w]_{\mathcal B}\).\qedhere
\end{itemize}
\end{proof}
If \(W\subseteq\R^n\) and the two ordered bases are the columns of \(X\) and of \(\tilde X\) in
\(\R^{n\times k}\), then \(\col_j(\tilde X)=\tilde v_j=X\col_j(P)\), that is, \(\tilde X=XP\)
(Theorem~\ref{r1-thm-views}(a),(b)). Hence \(\tilde X\tilde c=(XP)\tilde c=X(P\tilde c)\) for every \(\tilde c\in\R^k\)
(Theorem~\ref{r1-thm-rules}): the coefficients \(\tilde c\) on the columns of \(\tilde X\) and the coefficients
\(c=P\tilde c\) on the columns of \(X\) give the same vector, which is (a) in matrix form.

\needspace{12\baselineskip}
\begin{theorem}[\citeat{hoffman-kunze}{\S2.4, Theorem~8}, \citeat{tu-clm}{Proposition~11.4(i)}]\label{r3-thm-basis-from-P}
Let \(\mathcal B=(v_1,\dots,v_k)\) be an ordered basis of a subspace \(W\), and let
\(P=(p_{ij})\in\R^{k\times k}\) be invertible. Put \(\tilde v_j=p_{1j}v_1+\dots+p_{kj}v_k\)
(\(1\leqslant j\leqslant k\)). Then \(\tilde{\mathcal B}=(\tilde v_1,\dots,\tilde v_k)\) is an ordered
basis of \(W\), and it is the only ordered basis of \(W\) with \([w]_{\mathcal B}=P[w]_{\tilde{\mathcal B}}\)
for every \(w\in W\). In particular, if the columns of \(X\in\R^{n\times k}\) form a basis of a subspace
\(W\) of \(\R^n\), so do the columns of \(XP\).
\end{theorem}
\begin{proof}
\begin{itemize}
\item At most one: if \(\mathcal Y=(y_1,\dots,y_k)\) is an ordered basis of \(W\) with
\([w]_{\mathcal B}=P[w]_{\mathcal Y}\) for every \(w\), then \(w=y_j\), with \([y_j]_{\mathcal Y}=e_j\),
gives \([y_j]_{\mathcal B}=Pe_j=\col_j(P)\), that is, \(y_j=\sum_ip_{ij}v_i=\tilde v_j\).
\item Spanning: the \(\tilde v_j\) lie in \(W\). Let \(Q=(q_{jl})=P^{-1}\). For each \(l\),
\[
\sum_{j=1}^kq_{jl}\tilde v_j=\sum_{j=1}^kq_{jl}\sum_{i=1}^kp_{ij}v_i
=\sum_{i=1}^k\Bigl(\sum_{j=1}^kp_{ij}q_{jl}\Bigr)v_i=\sum_{i=1}^k(PQ)_{il}v_i=v_l
\]
So \(\Span\{\tilde v_1,\dots,\tilde v_k\}\) contains \(v_1,\dots,v_k\), hence contains
\(W=\Span\{v_1,\dots,v_k\}\) (Theorem~\ref{r3-thm-span}).
\item Independent: if \(\sum_jc_j\tilde v_j=0\), the same computation gives \(\sum_i(Pc)_iv_i=0\); so
\(Pc=0\), as the \(v_i\) are independent, and \(c=P^{-1}Pc=0\).
\item So \(\tilde{\mathcal B}\) is an ordered basis with \([\tilde v_j]_{\mathcal B}=\col_j(P)\), and
Theorem~\ref{r3-thm-change-basis}(a) gives \([w]_{\mathcal B}=P[w]_{\tilde{\mathcal B}}\).
\item For the last sentence take \(v_j=\col_j(X)\); then \(\tilde v_j=X\col_j(P)=\col_j(XP)\)
(Theorem~\ref{r1-thm-views}(a),(b)).\qedhere
\end{itemize}
\end{proof}

\begin{example}\label{r3-exm-contrasts}
In \(W_0=\{x\in\R^4:x_1+x_2+x_3+x_4=0\}\) take the ordered basis \(\mathcal B=(w_1,w_2,w_3)\) of
Example~\ref{r3-exm-centered-basis}. There \(x=x_1w_1+(x_1+x_2)w_2+(x_1+x_2+x_3)w_3\) for \(x\in W_0\), so
\([x]_{\mathcal B}=(x_1,\ x_1+x_2,\ x_1+x_2+x_3)'\), the partial sums of the entries of \(x\). Let
\(h_1=(1,0,0,-1)'\), \(h_2=(1,1,0,-2)'\) and \(h_3=(1,1,1,-3)'\), three vectors of \(W_0\) that compare the first
entry, the first two entries and the first three entries with the last.
\begin{itemize}
\item \([h_1]_{\mathcal B}=(1,1,1)'\), \([h_2]_{\mathcal B}=(1,2,2)'\) and \([h_3]_{\mathcal B}=(1,2,3)'\), the
partial sums; that is, \(h_1=w_1+w_2+w_3\), \(h_2=w_1+2w_2+2w_3\) and \(h_3=w_1+2w_2+3w_3\). With these columns,
\[
P=\begin{pmatrix}1&1&1\\1&2&2\\1&2&3\end{pmatrix},\qquad P^{-1}=\begin{pmatrix}2&-1&0\\-1&2&-1\\0&-1&1\end{pmatrix}
\]
since \(PP^{-1}=I_3\) (Theorem~\ref{r2-thm-invertible}).
\item As \(P\) is invertible, \(\tilde{\mathcal B}=(h_1,h_2,h_3)\) is an ordered basis of \(W_0\)
(Theorem~\ref{r3-thm-basis-from-P}), and \(P\) is the change-of-basis matrix from \(\tilde{\mathcal B}\)
to \(\mathcal B\); in matrix form \([h_1\ h_2\ h_3]=[w_1\ w_2\ w_3]P\).
\item By Theorem~\ref{r3-thm-change-basis}(b), for \(x\in W_0\),
\[
[x]_{\tilde{\mathcal B}}=P^{-1}[x]_{\mathcal B}
=\begin{pmatrix}2x_1-(x_1+x_2)\\-x_1+2(x_1+x_2)-(x_1+x_2+x_3)\\-(x_1+x_2)+(x_1+x_2+x_3)\end{pmatrix}
=\begin{pmatrix}x_1-x_2\\x_2-x_3\\x_3\end{pmatrix}
\]
For \(x=(3,1,2,-6)'\): \([x]_{\mathcal B}=(3,4,6)'\) and \([x]_{\tilde{\mathcal B}}=(2,-1,2)'\); indeed
\(3w_1+4w_2+6w_3=2h_1-h_2+2h_3=(3,1,2,-6)'\).
\end{itemize}
\end{example}

\subsection*{Sums and direct sums of subspaces}

\begin{definition}[\citeat{harville}{\S17.1}, \citeat{lang}{Ch.~I, \S4}, \citeat{abadir-magnus}{Exercise~3.12}]\label{r3-def-sum}
Let \(U\) and \(W\) be subspaces of a vector space \(V\). Their \textbf{intersection}
\(U\cap W\) is the set of vectors that lie in both \(U\) and \(W\), and their \textbf{sum} is
\[
U+W=\{u+w:\ u\in U,\ w\in W\}
\]
The sum is \textbf{direct}, written \(U\oplus W\), if \(U\cap W=\{0\}\), that is, if \(0\) is the
only vector that \(U\) and \(W\) have in common. Thus \(V=U\oplus W\) means that \(V=U+W\) and
\(U\cap W=\{0\}\).
\end{definition}

\begin{theorem}[\citeat{harville}{\S17.1, Lemma~17.1.1}, \citeat{abadir-magnus}{Exercise~3.12(a), (c)}, \citeat{harville-solutions}{Ch.~17, Exercise~1(a), (b)}, \citeat{hoffman-kunze}{\S2.2}]\label{r3-thm-sum}
Let \(U\) and \(W\) be subspaces of a vector space \(V\).
\begin{enumerate}
\item \(U\cap W\) and \(U+W\) are subspaces of \(V\).
\item \(U+W\) contains \(U\) and \(W\), and it is contained in every subspace of \(V\) that
contains \(U\cup W\): it is the smallest subspace of \(V\) containing \(U\cup W\). Moreover
\(U+W\) is the set of all linear combinations \(c_1x_1+\dots+c_kx_k\) of finitely many vectors
\(x_1,\dots,x_k\in U\cup W\) (this set is written \(\Span(U\cup W)\)).
\item If \(U=\Span\{u_1,\dots,u_p\}\) and \(W=\Span\{w_1,\dots,w_q\}\), then
\(U+W=\Span\{u_1,\dots,u_p,w_1,\dots,w_q\}\).
\end{enumerate}
\end{theorem}
\begin{proof}
(a)
\begin{itemize}
\item Both sets contain \(0\): \(0\in U\), \(0\in W\), and \(0=0+0\).
\item Let \(x,y\in U\cap W\) and \(c\in\R\). Then \(cx+y\in U\) and \(cx+y\in W\), since \(U\)
and \(W\) are subspaces; so \(cx+y\in U\cap W\).
\item Let \(x=u_1+w_1\) and \(y=u_2+w_2\) with \(u_1,u_2\in U\), \(w_1,w_2\in W\), and let
\(c\in\R\). Then \(cx+y=(cu_1+u_2)+(cw_1+w_2)\) with \(cu_1+u_2\in U\) and \(cw_1+w_2\in W\);
so \(cx+y\in U+W\).
\item By Theorem~\ref{r3-thm-subspace-test}(c), \(U\cap W\) and \(U+W\) are subspaces.
\end{itemize}
(b)
\begin{itemize}
\item Every \(u\in U\) is \(u=u+0\), and every \(w\in W\) is \(w=0+w\); so \(U\) and \(W\) lie
in \(U+W\).
\item If a subspace \(Z\) of \(V\) contains \(U\cup W\), then it contains \(u\) and \(w\) for all
\(u\in U\), \(w\in W\), hence also \(u+w\); so \(U+W\subseteq Z\).
\item In a combination \(c_1x_1+\dots+c_kx_k\) with all \(x_i\in U\cup W\), the terms with
\(x_i\in U\) add up to a vector \(u\in U\), and the other terms, whose \(x_i\) lie in \(W\), add
up to a vector \(w\in W\) (an empty sum is \(0\)). So the combination is \(u+w\in U+W\).
Conversely, \(u+w=1u+1w\) is such a combination.
\end{itemize}
(c)
\begin{itemize}
\item If \(u=\sum_ia_iu_i\in U\) and \(w=\sum_jb_jw_j\in W\), then
\(u+w=\sum_ia_iu_i+\sum_jb_jw_j\) is a combination of \(u_1,\dots,u_p,w_1,\dots,w_q\).
\item Conversely, \(\sum_ia_iu_i+\sum_jb_jw_j=u+w\) with \(u=\sum_ia_iu_i\in U\) and
\(w=\sum_jb_jw_j\in W\).\qedhere
\end{itemize}
\end{proof}

For matrices: let \(A\in\R^{m\times p}\) have columns \(a_1,\dots,a_p\), let \(B\in\R^{m\times q}\)
have columns \(b_1,\dots,b_q\), and let \([A\ B]\in\R^{m\times(p+q)}\) be the matrix with columns
\(a_1,\dots,a_p,b_1,\dots,b_q\). By (c), the columns of \([A\ B]\) span the sum of the subspaces
spanned by the columns of \(A\) and by the columns of \(B\),
\[
\Span\{a_1,\dots,a_p,b_1,\dots,b_q\}=\Span\{a_1,\dots,a_p\}+\Span\{b_1,\dots,b_q\}
\]
and, by Theorem~\ref{r1-thm-views}(b), the vectors of this subspace are the vectors \(Ax+By\)
(\(x\in\R^p\), \(y\in\R^q\)).

\begin{theorem}[\citeat{abadir-magnus}{Exercise~3.13(a), (b)}]\label{r3-thm-direct-sum}
Let \(U\) and \(W\) be subspaces of a vector space \(V\). Then \(U\cap W=\{0\}\) if and only if
every \(x\in U+W\) can be written in exactly one way as \(x=u+w\) with \(u\in U\) and \(w\in W\),
that is: if \(u+w=u_1+w_1\) with \(u,u_1\in U\) and \(w,w_1\in W\), then \(u=u_1\) and \(w=w_1\).
\end{theorem}
\begin{proof}
\begin{itemize}
\item Let \(U\cap W=\{0\}\), and let \(u+w=u_1+w_1\) as above. Then \(u-u_1=w_1-w\). The left
side lies in \(U\) and the right side in \(W\), so this vector lies in \(U\cap W=\{0\}\). Hence
\(u=u_1\) and \(w=w_1\).
\item Let \(U\cap W\ne\{0\}\), and take \(x\ne0\) in \(U\cap W\). Then \(x=x+0\) (with \(x\in U\),
\(0\in W\)) and \(x=0+x\) (with \(0\in U\), \(x\in W\)) are two different ways of writing
\(x\in U+W\).\qedhere
\end{itemize}
\end{proof}
So, if \(V=U\oplus W\), each \(v\in V\) is \(v=u+w\) for exactly one \(u\in U\) and one
\(w\in W\) (the parts of \(v\) in \(U\) and in \(W\)).

\begin{theorem}[\citeat{harville}{Theorem~17.1.5}]\label{r3-thm-direct-sum-basis}
Let \(U\) and \(W\) be subspaces of a vector space \(V\), let \(u_1,\dots,u_p\) form a basis of
\(U\), and let \(w_1,\dots,w_q\) form a basis of \(W\). If \(U\cap W=\{0\}\), then
\(u_1,\dots,u_p,w_1,\dots,w_q\) form a basis of \(U\oplus W\). If \(U\cap W\ne\{0\}\), then
\(u_1,\dots,u_p,w_1,\dots,w_q\) are linearly dependent. Hence the sum \(U+W\) is direct if and
only if \(u_1,\dots,u_p,w_1,\dots,w_q\) are linearly independent.
\end{theorem}
\begin{proof}
\begin{itemize}
\item The vectors \(u_1,\dots,u_p,w_1,\dots,w_q\) lie in \(U+W\) and span it, by
Theorem~\ref{r3-thm-sum}(b),(c).
\item Let \(U\cap W=\{0\}\) and \(\sum_ia_iu_i+\sum_jb_jw_j=0\). Then
\(x=\sum_ia_iu_i=\sum_j(-b_j)w_j\) lies in \(U\) and in \(W\), so \(x=0\). Hence
\(\sum_ia_iu_i=0\) and \(\sum_jb_jw_j=0\), and since each list is independent, every \(a_i\) and
every \(b_j\) is \(0\). So the vectors form a basis of \(U+W=U\oplus W\).
\item Let \(U\cap W\ne\{0\}\), take \(x\ne0\) in \(U\cap W\), and write
\(x=\sum_ia_iu_i=\sum_jb_jw_j\). Then \(\sum_ia_iu_i+\sum_j(-b_j)w_j=0\), and not every \(a_i\)
is \(0\), since \(x\ne0\): a linear relation among \(u_1,\dots,u_p,w_1,\dots,w_q\).\qedhere
\end{itemize}
\end{proof}

\begin{example}\label{r3-exm-sum-r3}
In \(\R^4\) let \(U=\{x:x_4=0\}=\Span\{e_1,e_2,e_3\}\) and
\(W_0=\{x:x_1+x_2+x_3+x_4=0\}=\Span\{w_1,w_2,w_3\}\), with \(w_1=(1,-1,0,0)'\),
\(w_2=(0,1,-1,0)'\) and \(w_3=(0,0,1,-1)'\) (Example~\ref{r3-exm-centered-basis}).
\begin{enumerate}
\item \(U+W_0=\R^4\), but the sum is not direct.
\begin{itemize}
\item Every \(x\in\R^4\) is \(x=(x_1,\ x_2,\ x_3+x_4,\ 0)'+(0,\ 0,\ -x_4,\ x_4)'\), with the
first vector in \(U\) and the second, \(-x_4w_3\), in \(W_0\). So \(U+W_0=\R^4\).
\item \(x\in U\cap W_0\) means \(x_4=0\) and \(x_1+x_2+x_3=0\), that is,
\(x=(a,\ b,\ -a-b,\ 0)'=a\,w_1+(a+b)\,w_2\) for some \(a,b\in\R\). So
\(U\cap W_0=\Span\{w_1,w_2\}\ne\{0\}\), and the sum is not direct.
\item Accordingly (Theorem~\ref{r3-thm-direct-sum}) the parts are not unique: for example
\[
(1,2,3,6)'=(1,2,9,0)'+(0,0,-6,6)'=(7,2,3,0)'+(-6,0,0,6)'
\]
and the two parts in \(U\) differ by \((-6,0,6,0)'=-6(w_1+w_2)\in U\cap W_0\). Likewise
\(e_1,e_2,e_3,w_1,w_2,w_3\) are dependent, since \(w_1=e_1-e_2\)
(Theorem~\ref{r3-thm-direct-sum-basis}).
\end{itemize}
\item \(\R^4=W_0\oplus\Span\{e_4\}\).
\begin{itemize}
\item \(ce_4=(0,0,0,c)'\) lies in \(W_0\) only for \(c=0\); so \(W_0\cap\Span\{e_4\}=\{0\}\).
\item For \(x\in\R^4\) and \(c\in\R\), \(x-ce_4=(x_1,\ x_2,\ x_3,\ x_4-c)'\) lies in \(W_0\)
exactly when \(c=x_1+x_2+x_3+x_4\). So every \(x\) is
\(x=(x_1,\ x_2,\ x_3,\ -x_1-x_2-x_3)'+(x_1+x_2+x_3+x_4)\,e_4\) with the first vector in \(W_0\),
and in no other way. For example, \((1,2,3,6)'=(1,2,3,-6)'+12e_4\).
\item The nonzero vector \(e_4\) is a basis of \(\Span\{e_4\}\) (Proposition~\ref{r3-prop-dependence}(a)); so, by
Theorem~\ref{r3-thm-direct-sum-basis}, \(w_1,w_2,w_3,e_4\) form a basis of \(\R^4\).
\end{itemize}
\end{enumerate}
\end{example}

\begin{example}\label{r3-exm-constant-centered}
The constant vectors and the centered vectors \(W_0=\{x\in\R^n:\one'x=0\}\)
(Example~\ref{r3-exm-subspaces}(f)) give a direct sum:
\[
\R^n=\Span\{\one_n\}\oplus W_0,\qquad x=\bar x\one+Cx
\]
\begin{itemize}
\item Sum: every \(x\in\R^n\) is \(x=\bar x\one+Cx\) with \(\bar x\one\in\Span\{\one\}\) and
\(Cx\in W_0\) (Example~\ref{r3-exm-constant}).
\item Intersection: if \(c\one\in W_0\), then \(0=\one'(c\one)=cn\)
(Proposition~\ref{r1-prop-centering}), so \(c=0\).
\item Hence \(\bar x\one\) is the only constant vector whose difference from \(x\) is centered
(Theorem~\ref{r3-thm-direct-sum}). In matrix form, \(x=\bar J_nx+C_nx\) with
\(\bar J_nx=\frac1n\one(\one'x)=\bar x\one\), and \(\bar J_n+C_n=I_n\)
(Definition~\ref{r1-def-ones}). For \(x=(1,3,5,7,9)'\) the two parts are \(\bar x\one=5\one\)
and \(Cx=(-4,-2,0,2,4)'\).
\item For \(n=4\), the nonzero vector \(\one_4\) is a basis of \(\Span\{\one_4\}\)
(Proposition~\ref{r3-prop-dependence}(a)); so \(\one_4\) together with the basis \(w_1,w_2,w_3\) of \(W_0\)
(Example~\ref{r3-exm-centered-basis}) is a basis of \(\R^4\) (Theorem~\ref{r3-thm-direct-sum-basis}).
\item The second subspace matters: also \(\R^4=W_0\oplus\Span\{e_4\}\)
(Example~\ref{r3-exm-sum-r3}(b)), and the parts of \(x=(1,2,3,6)'\) in \(W_0\) differ: here
\(\bar x=3\) and \(Cx=(-2,-1,0,3)'\), there \((1,2,3,-6)'\).
\end{itemize}
\end{example}

\begin{example}[cf.\ \citeat{harville-solutions}{Ch.~17, Exercise~2}]\label{r3-exm-sum-intersection-r4}
In \(\R^4\) let \(U=\Span\{u_1,u_2\}\) and \(W=\Span\{w_1,w_2\}\), where
\[
u_1=\begin{pmatrix}1\\0\\1\\2\end{pmatrix},\quad u_2=\begin{pmatrix}1\\1\\0\\1\end{pmatrix},\quad
w_1=\begin{pmatrix}0\\1\\0\\1\end{pmatrix},\quad w_2=\begin{pmatrix}2\\0\\1\\2\end{pmatrix}
\]
Find bases of \(U+W\) and of \(U\cap W\). Is the sum direct?

\textit{Solution.} By Theorem~\ref{r3-thm-sum}(c), \(U+W=\Span\{u_1,u_2,w_1,w_2\}\). Row-reduce
\(M=[u_1\ u_2\ w_1\ w_2]\):
\[
\begin{gathered}
\begin{pmatrix}1&1&0&2\\0&1&1&0\\1&0&0&1\\2&1&1&2\end{pmatrix}
\xrightarrow[R_4\leftarrow R_4-2R_1]{R_3\leftarrow R_3-R_1}
\begin{pmatrix}1&1&0&2\\0&1&1&0\\0&-1&0&-1\\0&-1&1&-2\end{pmatrix}
\xrightarrow[R_4\leftarrow R_4+R_2]{R_3\leftarrow R_3+R_2}
\begin{pmatrix}1&1&0&2\\0&1&1&0\\0&0&1&-1\\0&0&2&-2\end{pmatrix}
\\
\xrightarrow{R_4\leftarrow R_4-2R_3}
\begin{pmatrix}1&1&0&2\\0&1&1&0\\0&0&1&-1\\0&0&0&0\end{pmatrix}
\xrightarrow{R_2\leftarrow R_2-R_3}
\begin{pmatrix}1&1&0&2\\0&1&0&1\\0&0&1&-1\\0&0&0&0\end{pmatrix}
\xrightarrow{R_1\leftarrow R_1-R_2}
\begin{pmatrix}1&0&0&1\\0&1&0&1\\0&0&1&-1\\0&0&0&0\end{pmatrix}=R
\end{gathered}
\]
\begin{itemize}
\item Relations: \(Mc=0\) has the same solutions as \(Rc=0\) (Theorem~\ref{r2-thm-rowops}(b)), that is,
\(c_1+c_4=0\), \(c_2+c_4=0\), \(c_3-c_4=0\), with \(c_4=t\) free. So \(c=t(-1,-1,1,1)'\)
(Proposition~\ref{r2-prop-solutions}(b)), and \(t=1\) gives \(-u_1-u_2+w_1+w_2=0\), that is,
\(w_2=u_1+u_2-w_1\).
\item \(U+W\): by Theorem~\ref{r3-thm-span-inclusion}(c), \(U+W=\Span\{u_1,u_2,w_1\}\). The same row
operations take \([u_1\ u_2\ w_1]\) to the first three columns of \(R\) (row operations act on each
column separately, Theorems~\ref{r2-thm-elementary}(c) and~\ref{r1-thm-views}(a)), which are \(e_1,e_2,e_3\in\R^4\); every column is a pivot column, so
\(u_1,u_2,w_1\) are independent (Proposition~\ref{r3-prop-independence-test}) and form a basis of
\(U+W\).
\item \(U\cap W\): \(x\in U\cap W\) if and only if \(x=a_1u_1+a_2u_2=b_1w_1+b_2w_2\) for some
scalars, that is, \(M(a_1,a_2,-b_1,-b_2)'=0\). By the first item,
\((a_1,a_2,-b_1,-b_2)'=t(-1,-1,1,1)'\), so \(x=-t(u_1+u_2)\). Hence \(U\cap W=\Span\{u_1+u_2\}\), and the
nonzero vector \(u_1+u_2=(2,1,1,3)'=w_1+w_2\) is a basis of \(U\cap W\)
(Proposition~\ref{r3-prop-dependence}(a)).
\item The sum is not direct, since \(U\cap W\ne\{0\}\) (Definition~\ref{r3-def-sum}). Here \(u_1,u_2\) form a
basis of \(U\) (a sublist of the independent \(u_1,u_2,w_1\), Proposition~\ref{r3-prop-dependence}(b)), and
\(w_1,w_2\) a basis of \(W\) (neither is a multiple of the other, Proposition~\ref{r3-prop-dependence}(c));
accordingly \(u_1,u_2,w_1,w_2\) are dependent (Theorem~\ref{r3-thm-direct-sum-basis}).
\end{itemize}
\end{example}

\needspace{12\baselineskip}
\subsection*{Exercises}

\begin{exercise}\label{r3-dr-subspaces}
Decide whether each set is a subspace; for each subspace, find vectors (or matrices) that span it.
\begin{enumerate}
\item \(W_1=\{x\in\R^4:\ x_1+x_2+x_3=0,\ x_2-x_4=0,\ x_1+2x_2+x_3-x_4=0\}\).
\item \(W_2=\{x\in\R^3:\ x_1+x_2+x_3=1\}\).
\item \(W_3=\{x\in\R^3:\ x_1=x_2^2\}\).
\item \(W_4=\{A\in\R^{2\times3}:\ A\one_3=0\}\), the \(2\times3\) matrices whose rows add up to \(0\).
\end{enumerate}
\end{exercise}
\begin{solution}
(a) \(W_1\) is the set of solutions of \(Ax=0\) for the coefficient matrix \(A\) below, so it is a
subspace (Example~\ref{r3-exm-subspaces}(e)). Row-reduce \(A\):
\[
\begin{pmatrix}1&1&1&0\\0&1&0&-1\\1&2&1&-1\end{pmatrix}
\xrightarrow{R_3\leftarrow R_3-R_1}
\begin{pmatrix}1&1&1&0\\0&1&0&-1\\0&1&0&-1\end{pmatrix}
\xrightarrow{R_3\leftarrow R_3-R_2}
\begin{pmatrix}1&1&1&0\\0&1&0&-1\\0&0&0&0\end{pmatrix}
\xrightarrow{R_1\leftarrow R_1-R_2}
\begin{pmatrix}1&0&1&1\\0&1&0&-1\\0&0&0&0\end{pmatrix}
\]
The RREF gives \(x_1+x_3+x_4=0\) and \(x_2-x_4=0\), with \(x_3=s\) and \(x_4=t\) free
(Proposition~\ref{r2-prop-solutions}(b)):
\[
x=\begin{pmatrix}-s-t\\t\\s\\t\end{pmatrix}=s\begin{pmatrix}-1\\0\\1\\0\end{pmatrix}+t\begin{pmatrix}-1\\1\\0\\1\end{pmatrix}
\qquad(s,t\in\R)
\]
So \(W_1=\Span\{(-1,0,1,0)',(-1,1,0,1)'\}\). Check: \(-1+0+1=0\), \(0-0=0\), \(-1+0+1-0=0\), and
\(-1+1+0=0\), \(1-1=0\), \(-1+2+0-1=0\).

(b) Not a subspace: \(0\notin W_2\), since \(0+0+0\ne1\) (Theorem~\ref{r3-thm-subspace-test}(b)).

(c) Not a subspace: \((1,1,0)'\in W_3\), since \(1=1^2\), but \(2(1,1,0)'=(2,2,0)'\notin W_3\), since
\(2\ne2^2\); so \(W_3\) is not closed under scalar multiplication.

(d) A subspace: the zero matrix lies in \(W_4\), and if \(A\one=0\) and \(B\one=0\), then
\((cA+B)\one=cA\one+B\one=0\) (Theorem~\ref{r1-thm-rules}); apply Theorem~\ref{r3-thm-subspace-test}(c).
A matrix lies in \(W_4\) if and only if \(a_{13}=-a_{11}-a_{12}\) and \(a_{23}=-a_{21}-a_{22}\), that is,
\[
\begin{pmatrix}a&b&-a-b\\p&q&-p-q\end{pmatrix}=a\begin{pmatrix}1&0&-1\\0&0&0\end{pmatrix}+b\begin{pmatrix}0&1&-1\\0&0&0\end{pmatrix}
+p\begin{pmatrix}0&0&0\\1&0&-1\end{pmatrix}+q\begin{pmatrix}0&0&0\\0&1&-1\end{pmatrix}\qquad(a,b,p,q\in\R)
\]
So \(W_4\) is spanned by the four matrices on the right, \(E_{11}-E_{13}\), \(E_{12}-E_{13}\),
\(E_{21}-E_{23}\) and \(E_{22}-E_{23}\) (Definition~\ref{r3-def-basis}).
\end{solution}

\begin{exercise}\label{r3-ex-which-subspaces}
Which of the following subsets of \(\R^4\) are subspaces?
(a) \(\{x:x_1+x_2=x_3+x_4\}\);
(b) \(\{x:x_1x_4=x_2x_3\}\);
(c) \(\{x:x_1+x_2+x_3+x_4=1\}\);
(d) \(\{x:x_1\leqslant x_2\}\);
(e) \(\{x:x_1=x_2=x_3=x_4\}\).
\end{exercise}
\begin{solution}
(a) Yes: it is the solution set of \(Ax=0\) for the \(1\times4\) matrix \(A=(1,1,-1,-1)\),
a subspace by Example~\ref{r3-exm-subspaces}(e).

(b) No: \(e_1\) and \(e_4\) lie in the set (\(1\cdot0=0\cdot0\) and \(0\cdot1=0\cdot0\)), but
\(e_1+e_4=(1,0,0,1)'\) does not, since \(1\cdot1\ne0\cdot0\).

(c) No: \(0\) is not in the set, since \(0+0+0+0\ne1\).

(d) No: \(e_2=(0,1,0,0)'\) lies in the set, but \((-1)e_2=(0,-1,0,0)'\) does not, since
\(0>-1\).

(e) Yes: the set is \(\{c\one_4:c\in\R\}=\Span\{\one_4\}\), a subspace by
Theorem~\ref{r3-thm-span}.
\end{solution}

\begin{exercise}\label{r3-ex-span-drill}
Let \(W=\Span\{v_1,v_2,v_3\}\subseteq\R^4\), where \(v_1=(1,2,1,0)'\), \(v_2=(0,1,1,1)'\) and
\(v_3=(1,3,3,2)'\).
\begin{enumerate}
\item Which \(y=(y_1,y_2,y_3,y_4)'\) lie in \(W\)?
\item Show that \((3,6,4,1)'\in W\), and write it as a combination of \(v_1,v_2,v_3\).
\item Does \((1,1,1,1)'\) lie in \(W\)?
\end{enumerate}
\end{exercise}
\begin{solution}
(a) With \(V=[v_1\ v_2\ v_3]\), \(y\in W\) if and only if \(Vc=y\) is consistent (the paragraph
after Theorem~\ref{r3-thm-span}). Row-reduce \([V\mid y]\), keeping the \(y_i\) as letters:
\[
\begin{gathered}
\left(\begin{array}{ccc|c}1&0&1&y_1\\2&1&3&y_2\\1&1&3&y_3\\0&1&2&y_4\end{array}\right)
\xrightarrow[R_3\leftarrow R_3-R_1]{R_2\leftarrow R_2-2R_1}
\left(\begin{array}{ccc|c}1&0&1&y_1\\0&1&1&y_2-2y_1\\0&1&2&y_3-y_1\\0&1&2&y_4\end{array}\right)
\\
\xrightarrow[R_4\leftarrow R_4-R_2]{R_3\leftarrow R_3-R_2}
\left(\begin{array}{ccc|c}1&0&1&y_1\\0&1&1&y_2-2y_1\\0&0&1&y_1-y_2+y_3\\0&0&1&2y_1-y_2+y_4\end{array}\right)
\xrightarrow{R_4\leftarrow R_4-R_3}
\left(\begin{array}{ccc|c}1&0&1&y_1\\0&1&1&y_2-2y_1\\0&0&1&y_1-y_2+y_3\\0&0&0&y_1-y_3+y_4\end{array}\right)
\end{gathered}
\]
The left part of the last matrix is in echelon form, and its last row is \((0\ 0\ 0\mid y_1-y_3+y_4)\). If
\(y_1-y_3+y_4\ne0\), this row leaves no solution; if \(y_1-y_3+y_4=0\), the last matrix is in echelon form without
such a row. So, by Theorem~\ref{thm:r2-consistency}, \(Vc=y\) is consistent if and only if \(y_1-y_3+y_4=0\):
\[
y\in W\iff y_1-y_3+y_4=0
\]
Then every column of the left part is a pivot column, so the coefficients are unique
(Theorem~\ref{thm:r2-solution-structure}(b)); back substitution gives \(c_3=y_1-y_2+y_3\),
\(c_2=(y_2-2y_1)-c_3=-3y_1+2y_2-y_3\) and \(c_1=y_1-c_3=y_2-y_3\).

(b) For \(y=(3,6,4,1)'\), \(3-4+1=0\), so \(y\in W\), with \(c_3=3-6+4=1\), \(c_2=-9+12-4=-1\) and
\(c_1=6-4=2\): \(y=2v_1-v_2+v_3\). Check: \((2,4,2,0)'-(0,1,1,1)'+(1,3,3,2)'=(3,6,4,1)'\).

(c) No: \(y_1-y_3+y_4=1-1+1=1\ne0\).
\end{solution}

\begin{exercise}\label{r3-dr-representations}
Let \(v_1=(1,2,0,1)'\), \(v_2=(0,1,1,-1)'\), \(v_3=(1,4,2,-1)'\) and \(b=(2,3,-1,3)'\).
\begin{enumerate}
\item Show that \(b\in\Span\{v_1,v_2,v_3\}\), and find all \(c\in\R^3\) with \(b=c_1v_1+c_2v_2+c_3v_3\).
\item Give the representation with \(c_3=0\) and the one with \(c_1=0\).
\item Are \(v_1,v_2,v_3\) linearly independent? If not, give a linear relation among them.
\end{enumerate}
\end{exercise}
\begin{solution}
With \(V=[v_1\ v_2\ v_3]\), \(b=c_1v_1+c_2v_2+c_3v_3\) means \(Vc=b\) (Theorem~\ref{r1-thm-views}(b)).
Row-reduce \([V\mid b]\):
\[
\left(\begin{array}{ccc|c}1&0&1&2\\2&1&4&3\\0&1&2&-1\\1&-1&-1&3\end{array}\right)
\xrightarrow[R_4\leftarrow R_4-R_1]{R_2\leftarrow R_2-2R_1}
\left(\begin{array}{ccc|c}1&0&1&2\\0&1&2&-1\\0&1&2&-1\\0&-1&-2&1\end{array}\right)
\xrightarrow[R_4\leftarrow R_4+R_2]{R_3\leftarrow R_3-R_2}
\left(\begin{array}{ccc|c}1&0&1&2\\0&1&2&-1\\0&0&0&0\\0&0&0&0\end{array}\right)
\]
(a) The last matrix is in RREF and has no row \((0\ 0\ 0\mid d)\) with \(d\ne0\), so \(Vc=b\) is
consistent (Proposition~\ref{r2-prop-solutions}(a)) and \(b\in\Span\{v_1,v_2,v_3\}\). The pivot
columns are \(1,2\), and \(c_3=t\) is free; then \(c_1=2-t\) and \(c_2=-1-2t\)
(Proposition~\ref{r2-prop-solutions}(b)):
\[
c=\begin{pmatrix}2\\-1\\0\end{pmatrix}+t\begin{pmatrix}-1\\-2\\1\end{pmatrix}\qquad(t\in\R)
\]
(b) \(t=0\): \(b=2v_1-v_2\). \(t=2\): \(c=(0,-5,2)'\), so \(b=-5v_2+2v_3\). Check:
\((2,4,0,2)'-(0,1,1,-1)'=(2,3,-1,3)'\) and \(-5(0,1,1,-1)'+2(1,4,2,-1)'=(2,3,-1,3)'\).

(c) No: the left part of the last matrix is the RREF of \(V\), and its column \(3\) is not a pivot
column (Proposition~\ref{r3-prop-independence-test}). The difference of the solutions for \(t=1\) and
\(t=0\), \(c=(-1,-2,1)'\), solves \(Vc=0\) (Theorem~\ref{thm:r2-solution-structure}(a)):
\(-v_1-2v_2+v_3=0\), that is, \(v_3=v_1+2v_2\).
\end{solution}

\begin{exercise}\label{r3-dr-same-span}
Let \(a_1=(1,0,1,1)'\), \(a_2=(0,1,1,2)'\), \(b_1=(1,1,2,3)'\) and \(b_2=(1,2,3,5)'\).
\begin{enumerate}
\item Row-reduce \([a_1\ a_2\ b_1\ b_2]\), write \(b_1\) and \(b_2\) as combinations of \(a_1,a_2\), and
find \(F\in\R^{2\times2}\) with \([b_1\ b_2]=[a_1\ a_2]F\).
\item Write \(a_1\) and \(a_2\) as combinations of \(b_1,b_2\), and find \(G\in\R^{2\times2}\) with
\([a_1\ a_2]=[b_1\ b_2]G\).
\item Conclude that \(\Span\{a_1,a_2\}=\Span\{b_1,b_2\}\), and check that \(FG=I_2\).
\end{enumerate}
\end{exercise}
\begin{solution}
(a)
\[
\begin{pmatrix}1&0&1&1\\0&1&1&2\\1&1&2&3\\1&2&3&5\end{pmatrix}
\xrightarrow[R_4\leftarrow R_4-R_1]{R_3\leftarrow R_3-R_1}
\begin{pmatrix}1&0&1&1\\0&1&1&2\\0&1&1&2\\0&2&2&4\end{pmatrix}
\xrightarrow[R_4\leftarrow R_4-2R_2]{R_3\leftarrow R_3-R_2}
\begin{pmatrix}1&0&1&1\\0&1&1&2\\0&0&0&0\\0&0&0&0\end{pmatrix}
\]
The last matrix is in RREF; the solutions of \(c_1a_1+c_2a_2+c_3b_1+c_4b_2=0\) are \(c_1=-s-t\),
\(c_2=-s-2t\), \(c_3=s\), \(c_4=t\) (Theorem~\ref{r2-thm-rowops}(b),
Proposition~\ref{r2-prop-solutions}(b)). For \(s=1\), \(t=0\): \(-a_1-a_2+b_1=0\); for \(s=0\),
\(t=1\): \(-a_1-2a_2+b_2=0\). So
\[
b_1=a_1+a_2,\qquad b_2=a_1+2a_2,\qquad F=\begin{pmatrix}1&1\\1&2\end{pmatrix}
\]
(Theorem~\ref{r1-thm-views}(a),(b)). Check: \((1,0,1,1)'+(0,1,1,2)'=(1,1,2,3)'\) and
\((1,0,1,1)'+2(0,1,1,2)'=(1,2,3,5)'\).

(b) Subtracting, \(b_2-b_1=a_2\), and then \(a_1=b_1-a_2=2b_1-b_2\). So
\(G=\left(\begin{smallmatrix}2&-1\\-1&1\end{smallmatrix}\right)\).

(c) By (a), (b) and Theorem~\ref{r3-thm-span-inclusion}(b), the two spans are equal. Finally
\(FG=\left(\begin{smallmatrix}2-1&-1+1\\2-2&-1+2\end{smallmatrix}\right)=I_2\).
\end{solution}

\begin{exercise}[\citeat{abadir-magnus}{Exercise~3.19(c)}]\label{r3-ex-same-span}
Consider \(e_1=(1,0,0)'\), \(e_2=(0,1,0)'\) and \(v=(0,-2,0)'\) in \(\R^3\). Show
that \(e_1,e_2,v\) span \(W=\{(x,y,0)':x,y\in\R\}\), that \(e_1\) and \(e_2\) alone span the
same subspace \(W\), and that \(e_2\) and \(v\) span only \(\{(0,y,0)':y\in\R\}\).
\end{exercise}
\begin{solution}
For scalars \(a_1,a_2,a_3\) put \(\lambda=a_2-2a_3\). Since \(v=-2e_2\),
\[
a_1e_1+a_2e_2+a_3v=a_1e_1+\lambda e_2=(a_1,\lambda,0)',\qquad
a_2e_2+a_3v=\lambda e_2=(0,\lambda,0)'
\]
So \(\Span\{e_1,e_2,v\}\) consists of the vectors \((a_1,\lambda,0)'\). Each lies in \(W\),
and every \((x,y,0)'\in W\) is obtained with \(a_1=x\), \(a_2=y\), \(a_3=0\); so
\(\Span\{e_1,e_2,v\}=W\). The same computation with \(a_3=0\) gives \(\Span\{e_1,e_2\}=W\), and
\(\Span\{e_2,v\}=\{\lambda e_2:\lambda\in\R\}=\{(0,y,0)':y\in\R\}\).
\end{solution}

\begin{exercise}\label{r3-ex-independence-drill}
Let \(u_1=(1,2,3,1)'\), \(u_2=(0,1,1,1)'\), \(u_3=(1,2,3,0)'\) and \(u_4=(0,2,2,1)'\).
\begin{enumerate}
\item Are \(u_1,u_2,u_3\) linearly independent?
\item Are \(u_1,u_2,u_3,u_4\) linearly independent? If not, find a linear relation among them, and
write \(u_4\) as a combination of \(u_1,u_2,u_3\).
\end{enumerate}
\end{exercise}
\begin{solution}
Row-reduce \(V=[u_1\ u_2\ u_3\ u_4]\):
\[
\begin{gathered}
\begin{pmatrix}1&0&1&0\\2&1&2&2\\3&1&3&2\\1&1&0&1\end{pmatrix}
\xrightarrow[R_4\leftarrow R_4-R_1]{R_2\leftarrow R_2-2R_1,\ R_3\leftarrow R_3-3R_1}
\begin{pmatrix}1&0&1&0\\0&1&0&2\\0&1&0&2\\0&1&-1&1\end{pmatrix}
\xrightarrow[R_4\leftarrow R_4-R_2]{R_3\leftarrow R_3-R_2}
\begin{pmatrix}1&0&1&0\\0&1&0&2\\0&0&0&0\\0&0&-1&-1\end{pmatrix}
\\
\xrightarrow{R_3\leftrightarrow R_4}
\begin{pmatrix}1&0&1&0\\0&1&0&2\\0&0&-1&-1\\0&0&0&0\end{pmatrix}
\xrightarrow{R_3\leftarrow-R_3}
\begin{pmatrix}1&0&1&0\\0&1&0&2\\0&0&1&1\\0&0&0&0\end{pmatrix}
\xrightarrow{R_1\leftarrow R_1-R_3}
\begin{pmatrix}1&0&0&-1\\0&1&0&2\\0&0&1&1\\0&0&0&0\end{pmatrix}
\end{gathered}
\]
After the second arrow the entry of row \(3\) in column \(3\) is \(0\), so Step 2 interchanges rows
\(3\) and \(4\). The last matrix is the RREF of \(V\), with pivot columns \(1,2,3\).

(a) Yes. Row operations act on each column separately (Theorem~\ref{r2-thm-elementary}(c) and
Theorem~\ref{r1-thm-views}(a)), so the same operations take \([u_1\ u_2\ u_3]\) to the first three
columns of the last matrix, whose columns are the unit vectors \(e_1,e_2,e_3\) of \(\R^4\). This is
the RREF of \([u_1\ u_2\ u_3]\), and every column is a pivot column; so \(u_1,u_2,u_3\) are
independent (Proposition~\ref{r3-prop-independence-test}).

(b) No. Column \(4\) is not a pivot column, so \(Vc=0\) has a free unknown and a solution
\(c\ne0\) (Proposition~\ref{r3-prop-independence-test}). The RREF gives \(c_1-c_4=0\),
\(c_2+2c_4=0\), \(c_3+c_4=0\); with \(c_4=t\), \(c=t\,(1,-2,-1,1)'\). For \(t=1\),
\[
u_1-2u_2-u_3+u_4=0,\qquad\text{so}\qquad u_4=-u_1+2u_2+u_3
\]
Check: \(-(1,2,3,1)'+2(0,1,1,1)'+(1,2,3,0)'=(0,2,2,1)'=u_4\).
\end{solution}

\begin{exercise}\label{r3-dr-independence-parameter}
Let \(u_1=(1,0,1,2)'\), \(u_2=(2,1,3,k)'\) and \(u_3=(0,1,1,1)'\), where \(k\) is a number.
\begin{enumerate}
\item For which \(k\) are \(u_1,u_2,u_3\) linearly dependent?
\item For this \(k\), find a linear relation, and write \(u_2\) as a combination of \(u_1\) and \(u_3\).
\end{enumerate}
\end{exercise}
\begin{solution}
Row-reduce \(U=[u_1\ u_2\ u_3]\), keeping \(k\) as a letter:
\[
\begin{pmatrix}1&2&0\\0&1&1\\1&3&1\\2&k&1\end{pmatrix}
\xrightarrow[R_4\leftarrow R_4-2R_1]{R_3\leftarrow R_3-R_1}
\begin{pmatrix}1&2&0\\0&1&1\\0&1&1\\0&k-4&1\end{pmatrix}
\xrightarrow[R_4\leftarrow R_4-(k-4)R_2]{R_3\leftarrow R_3-R_2}
\begin{pmatrix}1&2&0\\0&1&1\\0&0&0\\0&0&5-k\end{pmatrix}
\]
(a) The vectors are independent exactly when \(Uc=0\) has only the solution \(c=0\)
(Proposition~\ref{r3-prop-independence-test}).
\begin{itemize}
\item If \(k\ne5\), interchange rows \(3\) and \(4\), and then scale row \(3\) by \(1/(5-k)\):
\[
\xrightarrow{R_3\leftrightarrow R_4}
\begin{pmatrix}1&2&0\\0&1&1\\0&0&5-k\\0&0&0\end{pmatrix}
\xrightarrow{R_3\leftarrow\frac1{5-k}R_3}
\begin{pmatrix}1&2&0\\0&1&1\\0&0&1\\0&0&0\end{pmatrix}
\]
The last matrix is an echelon form in which every column is a pivot column; so
\(c=0\) is the only solution (Theorem~\ref{thm:r2-solution-structure}(b)), and \(u_1,u_2,u_3\) are
independent.
\item If \(k=5\), the matrix after the second arrow has two zero rows; one more arrow gives the
RREF:
\[
\begin{pmatrix}1&2&0\\0&1&1\\0&0&0\\0&0&0\end{pmatrix}
\xrightarrow{R_1\leftarrow R_1-2R_2}
\begin{pmatrix}1&0&-2\\0&1&1\\0&0&0\\0&0&0\end{pmatrix}
\]
Column \(3\) is not a pivot column, so \(c_3\) is free and \(Uc=0\) has solutions \(c\ne0\): the
vectors are dependent.
\end{itemize}
So \(u_1,u_2,u_3\) are dependent only for \(k=5\).

(b) For \(k=5\), the RREF in (a) gives \(c_1-2c_3=0\) and \(c_2+c_3=0\), with \(c_3=t\) free: \(c=t(2,-1,1)'\). For \(t=1\),
\(2u_1-u_2+u_3=0\), that is, \(u_2=2u_1+u_3\). Check: \(2(1,0,1,2)'+(0,1,1,1)'=(2,1,3,5)'\), which is
\(u_2\) for \(k=5\).
\end{solution}

\begin{exercise}[\citeat{lang}{Ch.~I, \S2, Exercise~1(f)--(h)}, \citeat{shakarchi}{}]\label{r3-ex-independent-examples}
Show that the following vectors are linearly independent.
(a) \((1,2)'\) and \((1,3)'\).
(b) \((1,1,0)'\), \((1,1,1)'\) and \((0,1,-1)'\).
(c) \((0,1,1)'\), \((0,2,1)'\) and \((1,5,3)'\).
\end{exercise}
\begin{solution}
(a) If \(a(1,2)'+b(1,3)'=0\), then \(a+b=0\) and \(2a+3b=0\). The second equation minus
twice the first gives \(b=0\), and then \(a=0\).

(b) If \(a(1,1,0)'+b(1,1,1)'+c(0,1,-1)'=0\), then \(a+b=0\), \(a+b+c=0\),
\(b-c=0\). The second equation minus the first gives \(c=0\); then the third gives \(b=0\), and
the first \(a=0\).

(c) If \(a(0,1,1)'+b(0,2,1)'+c(1,5,3)'=0\), then \(c=0\), \(a+2b+5c=0\),
\(a+b+3c=0\). Subtracting the third equation from the second gives \(b+2c=0\), so \(b=0\) (as
\(c=0\)), and then \(a=0\).
\end{solution}

\begin{exercise}[\citeat{abadir-magnus}{Exercise~3.16(b)}]\label{r3-ex-column-row-relations}
Show that the columns of the matrix
\(A=\begin{pmatrix}1&3&2&0\\1&-5&-3&1\\5&-1&0&2\end{pmatrix}\) are linearly dependent. How
about the rows?
\end{exercise}
\begin{solution}
The columns satisfy, for example, \(3\col_1(A)-\col_2(A)-8\col_4(A)=0\):
\[
3\begin{pmatrix}1\\1\\5\end{pmatrix}-\begin{pmatrix}3\\-5\\-1\end{pmatrix}-8\begin{pmatrix}0\\1\\2\end{pmatrix}
=\begin{pmatrix}3-3-0\\3+5-8\\15+1-16\end{pmatrix}=0
\]
The coefficients \(3,-1,0,-8\) are not all \(0\), so the columns are dependent; equivalently,
\(Ax=0\) for \(x=(3,-1,0,-8)'\ne0\) (Theorem~\ref{r1-thm-views}(b)). The rows satisfy
\(3\row_1(A)+2\row_2(A)-\row_3(A)=0\):
\((3,9,6,0)+(2,-10,-6,2)-(5,-1,0,2)=(0,0,0,0)\), that is, \(y'A=0\) for
\(y=(3,2,-1)'\) (Theorem~\ref{r1-thm-views}(c)). So the rows are dependent too.
\end{solution}

\begin{exercise}[\citeat{harville-solutions}{Ch.~3, Exercise~1}]\label{r3-ex-dependent-k}
For what values of the scalar \(k\) are the three row vectors \((k,1,0)\), \((1,k,1)\) and
\((0,1,k)\) linearly dependent, and for what values are they linearly independent?
\end{exercise}
\begin{solution}
Let \(x_1,x_2,x_3\) be scalars with \(x_1(k,1,0)+x_2(1,k,1)+x_3(0,1,k)=0\), that is,
\[
kx_1+x_2=0,\qquad x_1+kx_2+x_3=0,\qquad x_2+kx_3=0
\]
or equivalently (the first and third equations give (1), the second gives (2))
\[
x_2=-kx_3=-kx_1\quad(1),\qquad kx_2=-x_1-x_3\quad(2)
\]
\begin{itemize}
\item \(k=0\): (1) and (2) become \(x_2=0\) and \(x_3=-x_1\); for instance
\((x_1,x_2,x_3)=(1,0,-1)\) is a solution other than \(0\).
\item \(k\ne0\): (1) gives \(kx_3=kx_1\), so \(x_3=x_1\); then (1) and (2) give
\(-k^2x_1=kx_2=-2x_1\), so \((2-k^2)x_1=0\). Hence \(k^2=2\), or else \(x_1=0\) and then
\(x_3=x_2=0\).
\item If \(k=\sqrt2\), conditions (1), (2) are equivalent to \(x_3=x_1\) and
\(x_2=-\sqrt2\,x_1\); if \(k=-\sqrt2\), to \(x_3=x_1\) and \(x_2=\sqrt2\,x_1\). In both cases
\(x_1=1\) gives a solution other than \(0\).
\end{itemize}
Thus there are solutions other than \(x_1=x_2=x_3=0\) if and only if \(k=0\) or
\(k=\pm\sqrt2\). The three vectors are linearly dependent if \(k\in\{0,\sqrt2,-\sqrt2\}\) and
linearly independent otherwise.
\end{solution}

\begin{exercise}\label{r3-dr-basis-of-span}
Let \(W=\Span\{v_1,\dots,v_5\}\subseteq\R^4\), where
\[
v_1=\begin{pmatrix}1\\1\\0\\2\end{pmatrix},\quad v_2=\begin{pmatrix}1\\0\\1\\1\end{pmatrix},\quad
v_3=\begin{pmatrix}1\\2\\-1\\3\end{pmatrix},\quad v_4=\begin{pmatrix}1\\0\\2\\1\end{pmatrix},\quad
v_5=\begin{pmatrix}1\\1\\-1\\2\end{pmatrix}
\]
\begin{enumerate}
\item Row-reduce \(V=[v_1\ \cdots\ v_5]\) to its RREF, and find all linear relations among
\(v_1,\dots,v_5\).
\item Find a basis of \(W\) that consists of some of the \(v_j\).
\item Write the remaining vectors as combinations of this basis.
\end{enumerate}
\end{exercise}
\begin{solution}
(a)
\[
\begin{gathered}
\begin{pmatrix}1&1&1&1&1\\1&0&2&0&1\\0&1&-1&2&-1\\2&1&3&1&2\end{pmatrix}
\xrightarrow[R_4\leftarrow R_4-2R_1]{R_2\leftarrow R_2-R_1}
\begin{pmatrix}1&1&1&1&1\\0&-1&1&-1&0\\0&1&-1&2&-1\\0&-1&1&-1&0\end{pmatrix}
\xrightarrow{R_2\leftarrow-R_2}
\begin{pmatrix}1&1&1&1&1\\0&1&-1&1&0\\0&1&-1&2&-1\\0&-1&1&-1&0\end{pmatrix}
\\
\xrightarrow[R_4\leftarrow R_4+R_2]{R_3\leftarrow R_3-R_2}
\begin{pmatrix}1&1&1&1&1\\0&1&-1&1&0\\0&0&0&1&-1\\0&0&0&0&0\end{pmatrix}
\xrightarrow[R_2\leftarrow R_2-R_3]{R_1\leftarrow R_1-R_3}
\begin{pmatrix}1&1&1&0&2\\0&1&-1&0&1\\0&0&0&1&-1\\0&0&0&0&0\end{pmatrix}
\xrightarrow{R_1\leftarrow R_1-R_2}
\begin{pmatrix}1&0&2&0&1\\0&1&-1&0&1\\0&0&0&1&-1\\0&0&0&0&0\end{pmatrix}
\end{gathered}
\]
The last matrix is the RREF \(R\) of \(V\). The relations \(c_1v_1+\dots+c_5v_5=0\) are the solutions
of \(Vc=0\), that is, of \(Rc=0\) (Theorem~\ref{r2-thm-rowops}(b)): \(c_1+2c_3+c_5=0\),
\(c_2-c_3+c_5=0\), \(c_4-c_5=0\), with \(c_3=s\) and \(c_5=t\) free
(Proposition~\ref{r2-prop-solutions}(b)):
\[
c=s\begin{pmatrix}-2\\1\\1\\0\\0\end{pmatrix}+t\begin{pmatrix}-1\\-1\\0\\1\\1\end{pmatrix}\qquad(s,t\in\R)
\]
For \(s=1\), \(t=0\): \(-2v_1+v_2+v_3=0\); for \(s=0\), \(t=1\): \(-v_1-v_2+v_4+v_5=0\).

(b) By these relations, \(v_3\) and \(v_5\) lie in \(\Span\{v_1,v_2,v_4\}\), so
\(W=\Span\{v_1,v_2,v_4\}\) (Theorem~\ref{r3-thm-span-inclusion}(c)). The same row operations take
\([v_1\ v_2\ v_4]\) to columns \(1,2,4\) of \(R\) (Theorems~\ref{r2-thm-elementary}(c)
and~\ref{r1-thm-views}(a)), which are \(e_1,e_2,e_3\in\R^4\); every column is a pivot column, so
\(v_1,v_2,v_4\) are independent (Proposition~\ref{r3-prop-independence-test}) and form a basis of \(W\).

(c) \(v_3=2v_1-v_2\) and \(v_5=v_1+v_2-v_4\). Check: \(2(1,1,0,2)'-(1,0,1,1)'=(1,2,-1,3)'\) and
\((1,1,0,2)'+(1,0,1,1)'-(1,0,2,1)'=(1,1,-1,2)'\).
\end{solution}

\begin{exercise}\label{r3-ex-basis-equations}
Let \(W=\{x\in\R^5:\ x_1+x_2+x_3+x_4+x_5=0,\ x_1+2x_2+3x_3+4x_4+5x_5=0\}\).
\begin{enumerate}
\item Find a basis of \(W\).
\item Show that \(w=(0,0,1,-2,1)'\in W\), and find its coordinates in this basis.
\end{enumerate}
\end{exercise}
\begin{solution}
(a) \(W\) is the solution set of \(Ax=0\) with
\(A=\left(\begin{smallmatrix}1&1&1&1&1\\1&2&3&4&5\end{smallmatrix}\right)\). Row-reduce:
\[
\begin{pmatrix}1&1&1&1&1\\1&2&3&4&5\end{pmatrix}
\xrightarrow{R_2\leftarrow R_2-R_1}
\begin{pmatrix}1&1&1&1&1\\0&1&2&3&4\end{pmatrix}
\xrightarrow{R_1\leftarrow R_1-R_2}
\begin{pmatrix}1&0&-1&-2&-3\\0&1&2&3&4\end{pmatrix}
\]
The free unknowns are \(x_3,x_4,x_5\), and \(x_1=x_3+2x_4+3x_5\), \(x_2=-2x_3-3x_4-4x_5\)
(Proposition~\ref{r2-prop-solutions}(b)). So
\[
W=\{x_3b_1+x_4b_2+x_5b_3:\ x_3,x_4,x_5\in\R\},\qquad
b_1=\begin{pmatrix}1\\-2\\1\\0\\0\end{pmatrix},\ b_2=\begin{pmatrix}2\\-3\\0\\1\\0\end{pmatrix},\
b_3=\begin{pmatrix}3\\-4\\0\\0\\1\end{pmatrix}
\]
The \(b_i\) lie in \(W\) and span it. Entries \(3,4,5\) of \(c_1b_1+c_2b_2+c_3b_3\) are
\(c_1,c_2,c_3\), so the combination is \(0\) only for \(c_1=c_2=c_3=0\). Hence \(b_1,b_2,b_3\)
form a basis of \(W\).

(b) \(0+0+1-2+1=0\) and \(0+0+3-8+5=0\), so \(w\in W\). By (a) its coordinates are its entries
\(3,4,5\), namely \((1,-2,1)\), and they are unique (Theorem~\ref{r3-thm-unique-coef}):
\(w=b_1-2b_2+b_3=(1-4+3,\ -2+6-4,\ 1,\ -2,\ 1)'\).
\end{solution}

\begin{exercise}\label{r3-ex-trend}
Let \(t=(1,2,3,4)'\) and \(W=\Span\{\one_4,t\}\subseteq\R^4\).
\begin{enumerate}
\item Show that \(\one_4,t\) form a basis of \(W\).
\item Which \(y=(y_1,y_2,y_3,y_4)'\) lie in \(W\)? Give the coordinates of such \(y\).
\item Find the coordinates of \((3,5,7,9)'\). Does \((1,2,4,8)'\) lie in \(W\)?
\end{enumerate}
\end{exercise}
\begin{solution}
(a) The vectors lie in \(W\) and span it. If \(a\one_4+bt=(a+b,\ a+2b,\ a+3b,\ a+4b)'=0\), the
second entry minus the first gives \(b=0\), and then \(a=0\). So they are independent.

(b) With \(c=(a,b)'\), \(y\in W\) if and only if \([\one_4\ t]c=y\) is consistent (the paragraph after
Theorem~\ref{r3-thm-span}). Row-reduce \([\one_4\ t\mid y]\), keeping the \(y_i\) as letters:
\[
\left(\begin{array}{cc|c}1&1&y_1\\1&2&y_2\\1&3&y_3\\1&4&y_4\end{array}\right)
\xrightarrow[R_4\leftarrow R_4-R_1]{R_2\leftarrow R_2-R_1,\ R_3\leftarrow R_3-R_1}
\left(\begin{array}{cc|c}1&1&y_1\\0&1&y_2-y_1\\0&2&y_3-y_1\\0&3&y_4-y_1\end{array}\right)
\xrightarrow[R_4\leftarrow R_4-3R_2]{R_3\leftarrow R_3-2R_2}
\left(\begin{array}{cc|c}1&1&y_1\\0&1&y_2-y_1\\0&0&y_1-2y_2+y_3\\0&0&2y_1-3y_2+y_4\end{array}\right)
\]
The left part of the last matrix is in echelon form, and its last two rows are \((0\ 0\mid y_1-2y_2+y_3)\) and
\((0\ 0\mid 2y_1-3y_2+y_4)\). So, by Theorem~\ref{thm:r2-consistency},
\[
y\in W\iff y_1-2y_2+y_3=0\ \text{ and }\ 2y_1-3y_2+y_4=0
\]
Then back substitution gives the coordinates \(b=y_2-y_1\) and \(a=y_1-b=2y_1-y_2\):
\(y=(2y_1-y_2)\one_4+(y_2-y_1)t\).

(c) For \((3,5,7,9)'\): \(3-10+7=0\) and \(6-15+9=0\), so it lies in \(W\), with \(a=1\),
\(b=2\): \((3,5,7,9)'=\one_4+2t\). For \((1,2,4,8)'\): \(1-4+4=1\ne0\), so it does not lie in
\(W\).
\end{solution}

\begin{exercise}\label{r3-dr-coordinates}
Let \(w_1=(1,1,0,1)'\), \(w_2=(1,2,1,0)'\) and \(W=\Span\{w_1,w_2\}\subseteq\R^4\).
\begin{enumerate}
\item Row-reduce \([w_1\ w_2\mid y]\) for \(y=(y_1,y_2,y_3,y_4)'\), keeping the \(y_i\) as letters. Show that
\(\mathcal B=(w_1,w_2)\) is an ordered basis of \(W\), find the conditions for \(y\in W\), and find
\([y]_{\mathcal B}\) for \(y\in W\).
\item Which of \((3,4,1,2)'\) and \((1,1,1,1)'\) lie in \(W\)? Give the coordinate vector of the one
that does.
\item Find \(w\in W\) with \([w]_{\mathcal B}=(3,-2)'\), and the coordinate vector of \(y+w\) for
\(y=(3,4,1,2)'\).
\end{enumerate}
\end{exercise}
\begin{solution}
(a)
\[
\begin{gathered}
\left(\begin{array}{cc|c}1&1&y_1\\1&2&y_2\\0&1&y_3\\1&0&y_4\end{array}\right)
\xrightarrow[R_4\leftarrow R_4-R_1]{R_2\leftarrow R_2-R_1}
\left(\begin{array}{cc|c}1&1&y_1\\0&1&y_2-y_1\\0&1&y_3\\0&-1&y_4-y_1\end{array}\right)
\xrightarrow[R_4\leftarrow R_4+R_2]{R_3\leftarrow R_3-R_2}
\left(\begin{array}{cc|c}1&1&y_1\\0&1&y_2-y_1\\0&0&y_1-y_2+y_3\\0&0&-2y_1+y_2+y_4\end{array}\right)
\\
\xrightarrow{R_1\leftarrow R_1-R_2}
\left(\begin{array}{cc|c}1&0&2y_1-y_2\\0&1&y_2-y_1\\0&0&y_1-y_2+y_3\\0&0&-2y_1+y_2+y_4\end{array}\right)
\end{gathered}
\]
The left part of the last matrix is the RREF of \([w_1\ w_2]\), with both columns pivot columns; so
\(w_1,w_2\) are independent (Proposition~\ref{r3-prop-independence-test}), and \(\mathcal B\) is an
ordered basis of \(W\). By Proposition~\ref{r2-prop-solutions}(a), the system \([w_1\ w_2]c=y\) is
consistent, that is, \(y\in W\), if and only if
\[
y_1-y_2+y_3=0,\qquad -2y_1+y_2+y_4=0
\]
and then its unique solution is \([y]_{\mathcal B}\) (Definition~\ref{r3-def-coordinates} and the
paragraph after it):
\[
[y]_{\mathcal B}=\begin{pmatrix}2y_1-y_2\\y_2-y_1\end{pmatrix}
\]
(b) For \((3,4,1,2)'\): \(3-4+1=0\) and \(-6+4+2=0\), so it lies in \(W\), with coordinate vector
\((6-4,\ 4-3)'=(2,1)'\). Check: \(2(1,1,0,1)'+(1,2,1,0)'=(3,4,1,2)'\). For \((1,1,1,1)'\):
\(1-1+1=1\ne0\), so it does not lie in \(W\).

(c) \(w=3w_1-2w_2=(3,3,0,3)'-(2,4,2,0)'=(1,-1,-2,3)'\) (Proposition~\ref{r3-prop-coordinates}(b)). By
Proposition~\ref{r3-prop-coordinates}(a), \([y+w]_{\mathcal B}=(2,1)'+(3,-2)'=(5,-1)'\). Check:
\(y+w=(4,3,-1,5)'=5(1,1,0,1)'-(1,2,1,0)'\).
\end{solution}

\begin{exercise}\label{r3-dr-change-basis}
Let \(u_1=(1,0,1,1)'\), \(u_2=(1,1,1,1)'\), \(u_3=(1,1,2,1)'\), \(u_4=(0,0,1,1)'\), \(U=[u_1\ u_2\ u_3\ u_4]\), and
\(\tilde v_1=(2,1,3,2)'\), \(\tilde v_2=(3,2,3,2)'\), \(\tilde v_3=(2,2,2,1)'\), \(\tilde v_4=(1,1,3,2)'\),
\(\tilde V=[\tilde v_1\ \tilde v_2\ \tilde v_3\ \tilde v_4]\).
\begin{enumerate}
\item Find \(U^{-1}\), and conclude that \(\mathcal B=(u_1,u_2,u_3,u_4)\) is an ordered basis of \(\R^4\).
\item Find the matrix \(P\) with columns \([\tilde v_1]_{\mathcal B},\dots,[\tilde v_4]_{\mathcal B}\), and check that
\(UP=\tilde V\).
\item Check that
\[
Q=\begin{pmatrix}0&-2&1&-1\\1&2&-1&1\\-1&-1&1&-1\\0&1&0&1\end{pmatrix}
\]
satisfies \(PQ=I_4\), and conclude that \(\tilde{\mathcal B}=(\tilde v_1,\tilde v_2,\tilde v_3,\tilde v_4)\) is an
ordered basis of \(\R^4\). For the vector \(w\) with \([w]_{\mathcal B}=(1,1,1,1)'\), find \(w\) and
\([w]_{\tilde{\mathcal B}}\).
\end{enumerate}
\end{exercise}
\begin{solution}
(a) Row-reduce \([U\mid I_4]\):
\[
\begin{gathered}
\left(\begin{array}{cccc|cccc}1&1&1&0&1&0&0&0\\0&1&1&0&0&1&0&0\\1&1&2&1&0&0&1&0\\1&1&1&1&0&0&0&1\end{array}\right)
\xrightarrow[R_4\leftarrow R_4-R_1]{R_3\leftarrow R_3-R_1}
\left(\begin{array}{cccc|cccc}1&1&1&0&1&0&0&0\\0&1&1&0&0&1&0&0\\0&0&1&1&-1&0&1&0\\0&0&0&1&-1&0&0&1\end{array}\right)
\\
\xrightarrow{R_3\leftarrow R_3-R_4}
\left(\begin{array}{cccc|cccc}1&1&1&0&1&0&0&0\\0&1&1&0&0&1&0&0\\0&0&1&0&0&0&1&-1\\0&0&0&1&-1&0&0&1\end{array}\right)
\xrightarrow[R_2\leftarrow R_2-R_3]{R_1\leftarrow R_1-R_3}
\left(\begin{array}{cccc|cccc}1&1&0&0&1&0&-1&1\\0&1&0&0&0&1&-1&1\\0&0&1&0&0&0&1&-1\\0&0&0&1&-1&0&0&1\end{array}\right)
\\
\xrightarrow{R_1\leftarrow R_1-R_2}
\left(\begin{array}{cccc|cccc}1&0&0&0&1&-1&0&0\\0&1&0&0&0&1&-1&1\\0&0&1&0&0&0&1&-1\\0&0&0&1&-1&0&0&1\end{array}\right)
\end{gathered}
\]
The left half has become \(I_4\), so \(U\) is invertible and \(U^{-1}\) is the right half
(Corollary~\ref{r2-cor-computing-inverse}(a)). By Example~\ref{r3-exm-invertible-basis}, \(\mathcal B\) is an
ordered basis of \(\R^4\), and \([y]_{\mathcal B}=U^{-1}y\) for every \(y\in\R^4\).

(b) Column \(j\) of \(P\) is \([\tilde v_j]_{\mathcal B}=U^{-1}\tilde v_j\), so \(P=U^{-1}\tilde V\)
(Theorem~\ref{r1-thm-views}(a)):
\[
P=\begin{pmatrix}1&-1&0&0\\0&1&-1&1\\0&0&1&-1\\-1&0&0&1\end{pmatrix}
\begin{pmatrix}2&3&2&1\\1&2&2&1\\3&3&2&3\\2&2&1&2\end{pmatrix}
=\begin{pmatrix}1&1&0&0\\0&1&1&0\\1&1&1&1\\0&-1&-1&1\end{pmatrix}
\]
Check:
\[
UP=\begin{pmatrix}1&1&1&0\\0&1&1&0\\1&1&2&1\\1&1&1&1\end{pmatrix}
\begin{pmatrix}1&1&0&0\\0&1&1&0\\1&1&1&1\\0&-1&-1&1\end{pmatrix}
=\begin{pmatrix}2&3&2&1\\1&2&2&1\\3&3&2&3\\2&2&1&2\end{pmatrix}=\tilde V
\]
(c) The product
\[
PQ=\begin{pmatrix}1&1&0&0\\0&1&1&0\\1&1&1&1\\0&-1&-1&1\end{pmatrix}
\begin{pmatrix}0&-2&1&-1\\1&2&-1&1\\-1&-1&1&-1\\0&1&0&1\end{pmatrix}
=\begin{pmatrix}1&0&0&0\\0&1&0&0\\0&0&1&0\\0&0&0&1\end{pmatrix}
\]
shows that \(P\) is invertible with \(P^{-1}=Q\) (Theorem~\ref{r2-thm-invertible}). As
\(\tilde v_j=U\col_j(P)=p_{1j}u_1+p_{2j}u_2+p_{3j}u_3+p_{4j}u_4\), Theorem~\ref{r3-thm-basis-from-P} shows that
\(\tilde{\mathcal B}\) is an ordered basis of \(\R^4\), and \(P\) is the change-of-basis matrix from
\(\tilde{\mathcal B}\) to \(\mathcal B\). Then \(w=u_1+u_2+u_3+u_4=(3,2,5,4)'\)
(Proposition~\ref{r3-prop-coordinates}(b)), and by Theorem~\ref{r3-thm-change-basis}(b)
\[
[w]_{\tilde{\mathcal B}}=Q[w]_{\mathcal B}
=\begin{pmatrix}0&-2&1&-1\\1&2&-1&1\\-1&-1&1&-1\\0&1&0&1\end{pmatrix}\begin{pmatrix}1\\1\\1\\1\end{pmatrix}
=\begin{pmatrix}-2\\3\\-2\\2\end{pmatrix}
\]
Check: \(-2\tilde v_1+3\tilde v_2-2\tilde v_3+2\tilde v_4=(-4,-2,-6,-4)'+(9,6,9,6)'+(-4,-4,-4,-2)'+(2,2,6,4)'
=(3,2,5,4)'=w\).
\end{solution}

\begin{exercise}\label{r3-dr-sum-intersection}
In \(\R^4\) let \(U=\Span\{u_1,u_2\}\) and \(W=\Span\{w_1,w_2\}\), where \(u_1=(1,0,1,1)'\),
\(u_2=(0,1,1,-1)'\), \(w_1=(1,1,2,0)'\) and \(w_2=(1,0,0,1)'\).
\begin{enumerate}
\item Row-reduce \(M=[u_1\ u_2\ w_1\ w_2]\) and find the linear relations among its columns.
\item Find a basis of \(U+W\).
\item Find a basis of \(U\cap W\). Is the sum \(U+W\) direct?
\end{enumerate}
\end{exercise}
\begin{solution}
(a)
\[
\begin{gathered}
\begin{pmatrix}1&0&1&1\\0&1&1&0\\1&1&2&0\\1&-1&0&1\end{pmatrix}
\xrightarrow[R_4\leftarrow R_4-R_1]{R_3\leftarrow R_3-R_1}
\begin{pmatrix}1&0&1&1\\0&1&1&0\\0&1&1&-1\\0&-1&-1&0\end{pmatrix}
\xrightarrow[R_4\leftarrow R_4+R_2]{R_3\leftarrow R_3-R_2}
\begin{pmatrix}1&0&1&1\\0&1&1&0\\0&0&0&-1\\0&0&0&0\end{pmatrix}
\\
\xrightarrow{R_3\leftarrow-R_3}
\begin{pmatrix}1&0&1&1\\0&1&1&0\\0&0&0&1\\0&0&0&0\end{pmatrix}
\xrightarrow{R_1\leftarrow R_1-R_3}
\begin{pmatrix}1&0&1&0\\0&1&1&0\\0&0&0&1\\0&0&0&0\end{pmatrix}=R
\end{gathered}
\]
\(Rc=0\) reads \(c_1+c_3=0\), \(c_2+c_3=0\), \(c_4=0\), with \(c_3=t\) free; so the relations are
\(c=t(-1,-1,1,0)'\) (Theorem~\ref{r2-thm-rowops}(b), Proposition~\ref{r2-prop-solutions}(b)), the
multiples of \(-u_1-u_2+w_1=0\): \(w_1=u_1+u_2\).

(b) By Theorem~\ref{r3-thm-sum}(c) and Theorem~\ref{r3-thm-span-inclusion}(c),
\(U+W=\Span\{u_1,u_2,w_1,w_2\}=\Span\{u_1,u_2,w_2\}\). The same row operations take
\([u_1\ u_2\ w_2]\) to columns \(1,2,4\) of \(R\) (Theorems~\ref{r2-thm-elementary}(c)
and~\ref{r1-thm-views}(a)), which are \(e_1,e_2,e_3\in\R^4\); so \(u_1,u_2,w_2\) are independent
(Proposition~\ref{r3-prop-independence-test}) and form a basis of \(U+W\).

(c) \(x\in U\cap W\) if and only if \(x=a_1u_1+a_2u_2=b_1w_1+b_2w_2\) for some scalars, that is,
\(M(a_1,a_2,-b_1,-b_2)'=0\). By (a), \((a_1,a_2,-b_1,-b_2)'=t(-1,-1,1,0)'\), so \(b_1=-t\), \(b_2=0\) and
\(x=-tw_1\). Hence \(U\cap W=\Span\{w_1\}\), with the basis \(w_1=(1,1,2,0)'\)
(Proposition~\ref{r3-prop-dependence}(a)). The sum is not direct, since \(U\cap W\ne\{0\}\)
(Definition~\ref{r3-def-sum}).
\end{solution}

\begin{exercise}[\citeat{harville-solutions}{Ch.~17, Exercise~2}]\label{r3-ex-sum-intersection}
Let \(U\) and \(W\) be the subspaces of \(\R^3\) spanned by the columns of
\[
A=\begin{pmatrix}1&0\\0&1\\0&0\end{pmatrix}\qquad\text{and}\qquad
B=\begin{pmatrix}0&2\\1&1\\2&3\end{pmatrix}
\]
respectively. Find (a) a basis of \(U+W\); (b) a basis of \(U\cap W\) (is the sum \(U+W\)
direct?); (c) a vector of \(U+W\) that is not in \(U\cup W\).
\end{exercise}
\begin{solution}
Let \(b_1=(0,1,2)'\) and \(b_2=(2,1,3)'\) be the columns of \(B\); the columns of \(A\) are
\(e_1,e_2\).

(a)
\begin{itemize}
\item By Theorem~\ref{r3-thm-sum}(c), \(U+W=\Span\{e_1,e_2,b_1,b_2\}\).
\item \(e_1,e_2,b_1\) are independent: \(c_1e_1+c_2e_2+c_3b_1=(c_1,\ c_2+c_3,\ 2c_3)'=0\) gives
\(c_3=0\), then \(c_2=0\) and \(c_1=0\).
\item So \(M=[e_1\ e_2\ b_1]\in\R^{3\times3}\) has \(Mc=0\) only for \(c=0\). By
Theorem~\ref{r2-thm-invertible}, \(M\) is invertible, and by
Example~\ref{r3-exm-invertible-basis} its columns form a basis of \(\R^3\).
\item These columns lie in the subspace \(U+W\), so every vector of \(\R^3\), a combination of
them, lies in \(U+W\). Hence \(U+W=\R^3\), and \(e_1,e_2,b_1\) form a basis of \(U+W\).
\end{itemize}
(b)
\begin{itemize}
\item \(U\) consists of the vectors \((x_1,x_2,0)'\), and \(W\) of the vectors
\(y_1b_1+y_2b_2=(2y_2,\ y_1+y_2,\ 2y_1+3y_2)'\) (\(x_1,x_2,y_1,y_2\in\R\)).
\item So \(U\cap W\) consists of the vectors \(y_1b_1+y_2b_2\) with \(2y_1+3y_2=0\). This
equation holds exactly when \(y_1=3t\) and \(y_2=-2t\) for some \(t\in\R\): these values give
\(6t-6t=0\); conversely, if \(2y_1+3y_2=0\), then \(t=y_1+y_2\) has \(3t=3y_1-2y_1=y_1\) and
\(-2t=3y_2-2y_2=y_2\).
\item For \(y_1=3t\), \(y_2=-2t\): \(y_1b_1+y_2b_2=(-4t,\ t,\ 0)'=t\,(-4,1,0)'\). Hence
\(U\cap W=\Span\{(-4,1,0)'\}\), and the single nonzero vector \((-4,1,0)'\) is a basis of it
(Proposition~\ref{r3-prop-dependence}(a)).
\item Since \(U\cap W\ne\{0\}\), the sum \(U+W\) is not direct (Definition~\ref{r3-def-sum}).
\end{itemize}
(c) Take \(x=(0,2,2)'\); it lies in \(U+W=\R^3\).
\begin{itemize}
\item \(x\notin U\), since its third entry is not \(0\).
\item If \(x=(2y_2,\ y_1+y_2,\ 2y_1+3y_2)'\), the first two entries give \(y_2=0\) and \(y_1=2\),
and then the third entry would be \(4\ne2\). So \(x\notin W\).
\end{itemize}
\end{solution}

\begin{exercise}[\citeat{lang}{Ch.~I, \S4, Exercise~2}, \citeat{shakarchi}{}]\label{r3-ex-direct-sum-r3}
In \(\R^3\) let \(W=\Span\{e_1\}\) and \(U=\Span\{u_1,u_2\}\), where \(u_1=(1,1,0)'\) and
\(u_2=(0,1,1)'\). Show that \(\R^3=W\oplus U\), and write \(x=(1,2,3)'\) as \(x=w+u\) with
\(w\in W\) and \(u\in U\).
\end{exercise}
\begin{solution}
The vectors \(e_1,u_1,u_2\) are linearly independent: if \(ae_1+bu_1+cu_2=0\), then
\[
a+b=0,\qquad b+c=0,\qquad c=0
\]
so \(c=0\), then \(b=0\) and \(a=0\).
\begin{itemize}
\item \(W\cap U=\{0\}\): \(e_1\) is a basis of \(W\), and \(u_1,u_2\) (independent by
Proposition~\ref{r3-prop-dependence}(b)) form a basis of \(U\). Their combined list
\(e_1,u_1,u_2\) is independent, so the sum \(W+U\) is direct
(Theorem~\ref{r3-thm-direct-sum-basis}).
\item \(W+U=\R^3\): the matrix \(M=[e_1\ u_1\ u_2]\in\R^{3\times3}\) has \(Mc=0\) only for
\(c=0\), so it is invertible (Theorem~\ref{r2-thm-invertible}), and its columns span \(\R^3\)
(Example~\ref{r3-exm-invertible-basis}). By Theorem~\ref{r3-thm-sum}(c),
\(W+U=\Span\{e_1,u_1,u_2\}=\R^3\).
\end{itemize}
Hence \(\R^3=W\oplus U\). To write \(x=(1,2,3)'\) as \(ae_1+bu_1+cu_2\), solve the system whose
augmented matrix is already in echelon form,
\[
\left(\begin{array}{ccc|c}1&1&0&1\\0&1&1&2\\0&0&1&3\end{array}\right)
\]
Back substitution (Theorem~\ref{thm:r2-consistency}) gives \(c=3\), \(b=2-3=-1\),
\(a=1-(-1)=2\). So \(w=2e_1=(2,0,0)'\) and \(u=-u_1+3u_2=(-1,2,3)'\); indeed
\(w+u=(1,2,3)'\). By Theorem~\ref{r3-thm-direct-sum}, this is the only such decomposition.
\end{solution}

\needspace{14\baselineskip}
\subsection*{Problems}

\begin{problem}[\citeat{abadir-magnus}{Exercise~3.9(c)}]\label{r3-ex-subsets-r2}
Is the set of vectors \((x,y)'\in\R^2\) with \(x\geqslant0\) and \(y\geqslant0\) a subspace of
\(\R^2\)?
\end{problem}
\begin{solution}
No: \((1,1)'\) lies in the set, but \((-1)(1,1)'=(-1,-1)'\) does not; so the set is not closed
under scalar multiplication (Theorem~\ref{r3-thm-subspace-test}(b)).
\end{solution}

\begin{problem}[\citeat{abadir-magnus}{Exercise~3.11}]\label{r3-ex-matrix-subspaces}
In the vector space \(\R^{3\times3}\), show that the set of lower triangular matrices is a
subspace.
\end{problem}
\begin{solution}
Let \(A,B\in\R^{3\times3}\) be lower triangular, that is, \(a_{ij}=b_{ij}=0\) for \(i<j\), and
let \(c\in\R\). For \(i<j\), \((A+B)_{ij}=a_{ij}+b_{ij}=0\) and \((cA)_{ij}=ca_{ij}=0\); so
\(A+B\) and \(cA\) are lower triangular, and so is \(0\). By
Theorem~\ref{r3-thm-subspace-test}(b), the lower triangular matrices form a subspace.
\end{solution}

\begin{problem}[\citeat{hoffman-kunze}{\S2.2, Exercise~7}]\label{r3-pr-union}
Let \(U\) and \(W\) be subspaces of a vector space \(V\) such that the union \(U\cup W\) is also a subspace. Show that
one of the spaces \(U\), \(W\) is contained in the other.
\end{problem}
\begin{solution}
Suppose that neither space is contained in the other: there are \(u\in U\) with \(u\notin W\), and \(w\in W\) with
\(w\notin U\). Both \(u\) and \(w\) lie in \(U\cup W\).
\begin{itemize}
\item If \(u+w\in U\), then, as \(U\) is a subspace, \(w=(u+w)+(-1)u\in U\ \contradiction\).
\item If \(u+w\in W\), then, as \(W\) is a subspace, \(u=(u+w)+(-1)w\in W\ \contradiction\).
\item So \(u+w\notin U\cup W\): the set \(U\cup W\) is not closed under addition, and it is not a subspace
(Theorem~\ref{r3-thm-subspace-test}(b)).
\end{itemize}
Hence, if \(U\cup W\) is a subspace, then \(U\subseteq W\) or \(W\subseteq U\).
\end{solution}

\begin{problem}[\citeat{harville-solutions}{Ch.~4, Exercise~7}]\label{r3-ex-nonunique}
Suppose that \(v_1,\dots,v_k\) span a vector space \(V\) but do not form a basis of \(V\). Show
that, for any \(v\in V\), the representation of \(v\) as a linear combination of
\(v_1,\dots,v_k\) is not unique.
\end{problem}
\begin{solution}
Let \(v\in V\). Since \(v_1,\dots,v_k\) span \(V\), there are scalars \(x_1,\dots,x_k\) with
\(v=\sum_ix_iv_i\). Since \(v_1,\dots,v_k\) span \(V\) but do not form a basis, they are
linearly dependent (Definition~\ref{r3-def-basis}): there are scalars \(z_1,\dots,z_k\), not
all \(0\), with \(\sum_iz_iv_i=0\). Let \(y_i=x_i+z_i\). Then
\(\sum_iy_iv_i=\sum_ix_iv_i+\sum_iz_iv_i=v+0=v\), and
\((y_1,\dots,y_k)\ne(x_1,\dots,x_k)\) because some \(z_i\ne0\). So \(v\) has two different
representations. (Compare Theorem~\ref{r3-thm-unique-coef}.)
\end{solution}

\begin{problem}[\citeat{lang}{Ch.~III, \S3, Exercise~14}, \citeat{shakarchi}{}]\label{r3-ex-sym-skew}
Show that every \(A\in\R^{n\times n}\) can be written in exactly one way as \(A=B+C\) with \(B\)
symmetric and \(C\) skew-symmetric, and conclude that \(\R^{n\times n}\) is the direct sum of the
subspace of symmetric matrices and the subspace of skew-symmetric matrices.
\end{problem}
\begin{solution}
Both sets are subspaces of \(\R^{n\times n}\) (Example~\ref{r3-exm-subspaces}(d)).
\begin{itemize}
\item Existence: let \(B=\frac12(A+A')\) and \(C=\frac12(A-A')\). Then \(B+C=A\), and, since
\((A')'=A\) (Theorem~\ref{r1-thm-transpose}), \(B'=\frac12(A'+A)=B\) and
\(C'=\frac12(A'-A)=-C\). (This is Proposition~\ref{r1-prop-sym-skew}(b).)
\item Uniqueness: let also \(A=B_1+C_1\) with \(B_1\) symmetric and \(C_1\) skew-symmetric. Then
\(B_1-B=C-C_1\), where \(B_1-B\) is symmetric and \(C-C_1\) is skew-symmetric. So
\(B_1-B=(B_1-B)'=(C-C_1)'=-(C-C_1)=-(B_1-B)\), hence \(2(B_1-B)=0\). Thus \(B_1=B\), and then
\(C_1=A-B_1=C\).
\item So the sum of the two subspaces is all of \(\R^{n\times n}\), and by
Theorem~\ref{r3-thm-direct-sum} their intersection is \(\{0\}\): \(\R^{n\times n}\) is their
direct sum.
\end{itemize}
For a numerical instance, see Example~\ref{r1-exm-sym-skew}, where the two parts of a \(4\times4\) matrix are
called \(S\) and \(K\).
\end{solution}

\begin{problem}[\citeat{harville-solutions}{Ch.~17, Exercise~4}]\label{r3-pr-modular}
Let \(U\), \(W\) and \(X\) be subspaces of a vector space \(V\).
\begin{enumerate}
\item Show that \((U\cap W)+(U\cap X)\subseteq U\cap(W+X)\).
\item Show by an example that \(U\cap W=\{0\}\) and \(U\cap X=\{0\}\) do not imply \(U\cap(W+X)=\{0\}\).
\item Show that if \(W\subseteq U\), then \(U+W=U\) and \(U\cap(W+X)=W+(U\cap X)\).
\end{enumerate}
\end{problem}
\begin{solution}
By Theorem~\ref{r3-thm-sum}(a), all the intersections and sums of subspaces below are subspaces of \(V\).

(a) Let \(y\in(U\cap W)+(U\cap X)\), that is, \(y=w+x\) with \(w\in U\cap W\) and \(x\in U\cap X\). Since \(w\) and
\(x\) lie in the subspace \(U\), so does \(y\); since \(w\in W\) and \(x\in X\), \(y\in W+X\). So \(y\in U\cap(W+X)\).

(b) In \(V=\R^2\) let \(U\), \(W\) and \(X\) be the spans of \((1,1)'\), \((1,0)'\) and \((0,1)'\), respectively.
\begin{itemize}
\item \(U\cap W=\{0\}\): a vector of \(U\cap W\) is \(a(1,1)'=b(1,0)'\) for some \(a,b\in\R\); then
\((a,a)'=(b,0)'\), so \(a=0\), and the vector is \(0\). Likewise \(a(1,1)'=b(0,1)'\) gives \((a,a)'=(0,b)'\), so
\(a=0\), and \(U\cap X=\{0\}\).
\item \(W+X=\Span\{e_1,e_2\}=\R^2\) (Theorem~\ref{r3-thm-sum}(c) and the standard basis,
Definition~\ref{r3-def-basis}). So \(U\cap(W+X)=U\ne\{0\}\).
\end{itemize}
(c) Let \(W\subseteq U\).
\begin{itemize}
\item \(U+W=U\): \(U\subseteq U+W\) by Theorem~\ref{r3-thm-sum}(b). Conversely, if \(y=u+w\) with \(u\in U\) and
\(w\in W\), then \(w\in U\), so \(y\in U\).
\item \(W+(U\cap X)\subseteq U\cap(W+X)\): this is (a), since \(U\cap W=W\).
\item \(U\cap(W+X)\subseteq W+(U\cap X)\): let \(y\in U\cap(W+X)\), say \(y=w+x\) with \(w\in W\) and \(x\in X\).
Then \(x=y-w\), where \(y\in U\) and \(w\in U\) (as \(W\subseteq U\)); so \(x\in U\), as \(U\) is a subspace.
Hence \(x\in U\cap X\), and \(y=w+x\in W+(U\cap X)\).
\end{itemize}
\end{solution}

\begin{problem}[\citeat{prasolov}{Problem~5.3}]\label{r3-ex-pr-orthants}
Let \(v\in\R^n\) have all entries nonzero, and let \(W=\{w\in\R^n:v'w=0\}\), a subspace of \(\R^n\)
(Example~\ref{r3-exm-subspaces}(e), with the \(1\times n\) matrix \(v'\)). For
\(\varepsilon=(\varepsilon_1,\dots,\varepsilon_n)'\) with every \(\varepsilon_i\in\{1,-1\}\), the \textbf{orthant}
\(O_\varepsilon\) is the set of the \(w\in\R^n\) with \(\varepsilon_iw_i>0\) for \(i=1,\dots,n\). Show that \(W\)
contains a vector of every orthant except the two orthants that contain \(v\) and \(-v\), and that \(W\) contains no
vector of these two orthants.
\end{problem}
\begin{solution}
Here \(v'w=\sum_iv_iw_i\) (Theorem~\ref{r1-thm-transpose}), and \(\varepsilon_i^2=1\) for every \(i\).
\begin{itemize}
\item Let \(v\in O_\varepsilon\). For \(w\in O_\varepsilon\), \(v_iw_i=(\varepsilon_iv_i)(\varepsilon_iw_i)>0\) for
every \(i\), so \(v'w>0\) and \(w\notin W\). Likewise, if \(-v\in O_\varepsilon\) and \(w\in O_\varepsilon\), then
\(-v_iw_i=(-\varepsilon_iv_i)(\varepsilon_iw_i)>0\) for every \(i\), so \(v'w<0\) and \(w\notin W\).
\item Let \(O_\varepsilon\) contain neither \(v\) nor \(-v\). As \(v\notin O_\varepsilon\), some
\(\varepsilon_jv_j\leqslant0\), hence \(\varepsilon_jv_j<0\) since \(v_j\ne0\); as \(-v\notin O_\varepsilon\), some
\(-\varepsilon_iv_i\leqslant0\), hence \(\varepsilon_iv_i>0\), again since \(v_i\ne0\).
\item The vector \(\varepsilon\) itself lies in \(O_\varepsilon\), since \(\varepsilon_i\varepsilon_i=1>0\). Let
\(s=v'\varepsilon\). If \(s=0\), then \(\varepsilon\in W\cap O_\varepsilon\).
\item If \(s>0\), enlarge entry \(j\): let \(w=\varepsilon+t\varepsilon_je_j\) with \(t=s/(-\varepsilon_jv_j)>0\). Entry
\(j\) of \(w\) is \((1+t)\varepsilon_j\), so \(\varepsilon_jw_j=1+t>0\), and the other entries are those of
\(\varepsilon\); so \(w\in O_\varepsilon\). Moreover \(v'w=s+t\varepsilon_jv_j=s-s=0\).
\item If \(s<0\), enlarge entry \(i\): let \(w=\varepsilon+t\varepsilon_ie_i\) with \(t=-s/(\varepsilon_iv_i)>0\). As
before, \(w\in O_\varepsilon\), and \(v'w=s+t\varepsilon_iv_i=s-s=0\).
\end{itemize}
In each case \(W\) contains a vector of \(O_\varepsilon\). For \(v=\one_n\), \(W\) is the subspace \(W_0\) of centered
vectors (Example~\ref{r3-exm-subspaces}(f)): a centered vector without zero entries can have any pattern of signs except
all entries positive and all entries negative.
\end{solution}

\begin{problem}[\citeat{prasolov}{Problems~7.1, 7.2}]\label{r3-ex-pr-caratheodory}
Let \(v_1,\dots,v_m\in\R^n\).
\begin{enumerate}
\item Show that if \(m\geqslant n+2\), then there are numbers \(\alpha_1,\dots,\alpha_m\), not all \(0\), with
\(\alpha_1v_1+\dots+\alpha_mv_m=0\) and \(\alpha_1+\dots+\alpha_m=0\).
\item A \textbf{convex combination} of \(v_1,\dots,v_m\) is a vector \(x=t_1v_1+\dots+t_mv_m\) with
\(t_1,\dots,t_m\geqslant0\) and \(t_1+\dots+t_m=1\). Show that every convex combination of \(v_1,\dots,v_m\) is a convex
combination of at most \(n+1\) of the vectors \(v_1,\dots,v_m\).
\end{enumerate}
\end{problem}
\begin{solution}
(a) For \(i=1,\dots,m\), let \(u_i\in\R^{n+1}\) be \(v_i\) with the entry \(1\) appended: its first \(n\) entries are
those of \(v_i\), and its last entry is \(1\). For numbers \(\alpha_1,\dots,\alpha_m\), the first \(n\) entries of
\(\sum_i\alpha_iu_i\) form the vector \(\sum_i\alpha_iv_i\), and its last entry is \(\sum_i\alpha_i\). The \(m\) vectors
\(u_1,\dots,u_m\) of \(\R^{n+1}\) are more than \(n+1\), so they are linearly dependent
(Proposition~\ref{r3-prop-independence-test}): \(\sum_i\alpha_iu_i=0\) for some \(\alpha_1,\dots,\alpha_m\), not all
\(0\). These numbers satisfy \(\sum_i\alpha_iv_i=0\) and \(\sum_i\alpha_i=0\).

(b) We show, by induction on \(m\), that a convex combination \(x\) of any \(m\) vectors of \(\R^n\) is a convex
combination of at most \(n+1\) of them.
\begin{itemize}
\item If \(m\leqslant n+1\), there is nothing to prove.
\item Let \(m\geqslant n+2\), assume the statement for \(m-1\) vectors, and let \(x=\sum_it_iv_i\) with all
\(t_i\geqslant0\) and \(\sum_it_i=1\). Take \(\alpha_1,\dots,\alpha_m\) as in (a). Some \(\alpha_i\) is positive:
otherwise all \(\alpha_i\leqslant0\), and as \(\sum_i\alpha_i=0\), all \(\alpha_i\) would be \(0\).
\item For every \(\lambda\in\R\), \(x=\sum_it_iv_i+\lambda\sum_i\alpha_iv_i=\sum_i\tilde t_iv_i\) with
\(\tilde t_i=t_i+\lambda\alpha_i\), and \(\sum_i\tilde t_i=\sum_it_i+\lambda\sum_i\alpha_i=1\).
\item Among the indices \(i\) with \(\alpha_i>0\), choose \(k\) with \(t_k/\alpha_k\leqslant t_i/\alpha_i\) for all of
them, and take \(\lambda=-t_k/\alpha_k\leqslant0\). Then \(\tilde t_k=t_k-t_k=0\). If \(\alpha_i>0\), then
\(\tilde t_i=\alpha_i\bigl(t_i/\alpha_i-t_k/\alpha_k\bigr)\geqslant0\); if \(\alpha_i\leqslant0\), then
\(\lambda\alpha_i\geqslant0\) and \(\tilde t_i\geqslant t_i\geqslant0\).
\item So \(x=\sum_{i\ne k}\tilde t_iv_i\) with all \(\tilde t_i\geqslant0\) and \(\sum_{i\ne k}\tilde t_i=1\): \(x\) is a convex
combination of the \(m-1\) vectors \(v_i\) (\(i\ne k\)). By the induction hypothesis, it is a convex combination of at
most \(n+1\) of them.
\end{itemize}
\end{solution}

\clearpage
\def\unitnumber{4}
\setcounter{theorem}{0}\setcounter{problem}{0}
\section*{Chapter 4. Bases, dimension, and rank}

\subsection*{Bases and dimension of subspaces of \(\R^n\)}

\begin{theorem}[cf.\ \citeat{hoffman-kunze}{\S2.3, Example~14}]\label{thm:r4-pivot-test}
Let \(R\) be the RREF of \(A\in\R^{m\times n}\), with \(r\) nonzero rows. The following are
equivalent:
\begin{enumerate}[label=(\roman*)]
\item the columns \(a_1,\dots,a_n\) of \(A\) are linearly independent;
\item the system \(Ax=0\) has only the solution \(x=0\);
\item \(r=n\), that is, every column of \(R\) is a pivot column.
\end{enumerate}
In particular, if \(n>m\), then any \(n\) vectors of \(\R^m\) are linearly dependent.
\end{theorem}
We call this the \textbf{pivot test}.
\begin{proof}
This is Proposition~\ref{r3-prop-independence-test} with \(V=A\) and \(k=n\): each nonzero row of \(R\)
contains exactly one pivot, and the pivots lie in different columns (Definition~\ref{r2-def-rref}), so
\(R\) has \(r\) pivot columns, and every column of \(R\) is a pivot column exactly when \(r=n\).
\end{proof}

\begin{remark}[cf.\ \citeat{hoffman-kunze}{\S1.3, proof of Theorem~4}]\label{rem:r4-echelon}
For the pivot test one may stop at an echelon form \(E\) of \(A\) (Definitions~\ref{r2-def-rref}
and~\ref{r2-def-gauss-jordan}).
Backward elimination adds multiples of a row \(i\), whose leading \(1\) is in column \(k_i\), to rows
\(l<i\), whose leading \(1\) is in a column \(k_l<k_i\). Row \(i\) is \(0\) before column \(k_i\), so
row \(l\) changes only in columns \(\geqslant k_i\) and keeps its leading \(1\) in column \(k_l\).
Hence the RREF of \(A\) reached from \(E\) has the same leading columns as \(E\). As the RREF of \(A\)
is unique (Theorem~\ref{r2-thm-rref-unique}), the pivot columns of \(A\), and their
number \(r\), can be read off from any echelon form of \(A\). In particular,
(iii) of Theorem~\ref{thm:r4-pivot-test} says that every column of \(E\) contains a leading \(1\).
The same holds before the leading entries are scaled to \(1\): scaling a row does not move its
leading entry.
\end{remark}

\begin{theorem}[cf.\ \citeat{lang-ila}{Ch.~III, \S6, Theorem~6.1}]\label{thm:r4-pivot-columns}
Let \(R\) be the RREF of \(A\in\R^{m\times n}\), with \(r\) nonzero rows, pivot columns
\(k_1<\dots<k_r\) and entries \(R_{ij}\).
\begin{enumerate}
\item The columns of \(A\) and those of \(R\) satisfy the same linear relations: for
\(x\in\R^n\), \(x_1a_1+\dots+x_na_n=0\) if and only if \(x_1\col_1(R)+\dots+x_n\col_n(R)=0\).
\item If \(r\geqslant1\), the pivot columns \(a_{k_1},\dots,a_{k_r}\) of \(A\) form a basis of
\(\Span\{a_1,\dots,a_n\}\), and for every \(j=1,\dots,n\)
\[
a_j=R_{1j}a_{k_1}+R_{2j}a_{k_2}+\dots+R_{rj}a_{k_r}
\]
If \(r=0\), then \(A=0\) and \(\Span\{a_1,\dots,a_n\}=\{0\}\).
\end{enumerate}
In particular, every finite list that spans a subspace \(W\) contains a basis of \(W\).
\end{theorem}
\begin{proof}
(a) The two sums are \(Ax\) and \(Rx\), and \(Ax=0\) if and only if \(Rx=0\) (Theorem~\ref{r2-thm-rowops}(b)).

(b) If \(r=0\), then \(R=0\), and \(A=P^{-1}R=0\), where \(R=PA\) with \(P\in\R^{m\times m}\)
invertible (Theorem~\ref{r2-thm-elementary}(c)). Let \(r\geqslant1\).
\begin{itemize}
\item Columns of \(R\): \(\col_{k_i}(R)=e_i\), and rows \(r+1,\dots,m\) of \(R\) are zero, so
\(\col_j(R)=\sum_{i=1}^rR_{ij}e_i=\sum_{i=1}^rR_{ij}\col_{k_i}(R)\) for every \(j\).
\item Relations: for a column \(j\notin\{k_1,\dots,k_r\}\) let \(x\in\R^n\) have \(x_j=1\),
\(x_{k_i}=-R_{ij}\) (\(1\leqslant i\leqslant r\)) and all other entries \(0\). By the first item,
\(Rx=\col_j(R)-\sum_iR_{ij}\col_{k_i}(R)=0\); by (a), \(Ax=a_j-\sum_iR_{ij}a_{k_i}=0\). For
\(j=k_l\) the formula reads \(a_{k_l}=a_{k_l}\), since \(R_{ik_l}\) is \(1\) for \(i=l\) and \(0\)
otherwise.
\item Spanning: each \(a_j\) lies in the subspace \(U=\Span\{a_{k_1},\dots,a_{k_r}\}\), so
\(\Span\{a_1,\dots,a_n\}\subseteq U\) (Theorem~\ref{r3-thm-span}); the reverse inclusion is clear.
\item Independence: if \(c_1a_{k_1}+\dots+c_ra_{k_r}=0\), then by (a)
\(c_1\col_{k_1}(R)+\dots+c_r\col_{k_r}(R)=c_1e_1+\dots+c_re_r=0\), so every \(c_i=0\).
\end{itemize}
For the last sentence, write the spanning vectors as the columns of a matrix and apply (b).
\end{proof}

By Theorem~\ref{thm:r4-pivot-columns}(a), \(a_j\) is a linear combination of \(a_1,\dots,a_{j-1}\) if
and only if \(\col_j(R)\) is a linear combination of \(\col_1(R),\dots,\col_{j-1}(R)\) (take
\(x_j=-1\) and \(x_l=0\) for \(l>j\)). In \(R\) this happens exactly when \(j\) is not a pivot column:
\(\col_{k_i}(R)=e_i\), while the columns before \(k_i\) are \(0\) in rows \(i,\dots,m\); and a
non-pivot column \(j\) equals \(\sum_{k_i<j}R_{ij}\col_{k_i}(R)\). So the pivot columns of \(A\) are
the columns that are not combinations of the earlier ones; they do not depend on the row operations
used.

\begin{theorem}[\citeat{hoffman-kunze}{\S2.3, Theorems~4, 5}]\label{thm:r4-dimension}
Every subspace \(W\) of \(\R^n\) has a basis, and any two bases of \(W\) have the same number of
elements. This number is the \textbf{dimension} of \(W\), written \(\dim W\). For example,
\(\dim\R^n=n\) (standard basis) and \(\dim\{0\}=0\) (the empty list is a basis of \(\{0\}\)). Conversely,
\(\dim W=0\) only for \(W=\{0\}\): a basis of \(W\) with no vectors is the empty list, which spans \(\{0\}\)
(Definition~\ref{r3-def-span}). Moreover, if \(W\) is spanned by \(m\) vectors, then any \(k>m\) vectors of \(W\)
are linearly dependent.
\end{theorem}
\begin{proof}
\begin{itemize}
\item If \(W\) is spanned by \(w_1,\dots,w_m\), then any \(k>m\) vectors \(v_1,\dots,v_k\) of \(W\)
are linearly dependent. Indeed, write \(v_j=a_{1j}w_1+\dots+a_{mj}w_m\) (\(1\leqslant j\leqslant k\))
and \(A=(a_{ij})\in\R^{m\times k}\). As \(k>m\), the columns of \(A\) are dependent
(Theorem~\ref{thm:r4-pivot-test}): \(Ax=0\) for some \(x\ne0\). Then
\[
x_1v_1+\dots+x_kv_k=\sum_{j=1}^kx_j\sum_{i=1}^ma_{ij}w_i=\sum_{i=1}^m(Ax)_iw_i=0
\]
\item Same size: if \(v_1,\dots,v_k\) and \(w_1,\dots,w_m\) are bases of \(W\), each list spans
\(W\) and is independent, so the first item gives \(k\leqslant m\) and \(m\leqslant k\).
\item Existence: start with the empty list and, as long as the vectors chosen do not span \(W\),
adjoin a vector \(v\in W\) outside their span. The list stays independent: in a relation
\(c_1v_1+\dots+c_kv_k+cv=0\), \(c\ne0\) would put \(v\) in the span of \(v_1,\dots,v_k\), so \(c=0\),
and then every \(c_i=0\). More than \(n\) vectors of \(\R^n\) are dependent
(Theorem~\ref{thm:r4-pivot-test}), so the process stops after at most \(n\) steps, at a basis of
\(W\).\qedhere
\end{itemize}
\end{proof}

By Theorem~\ref{thm:r4-pivot-columns}, \(\dim\Span\{a_1,\dots,a_n\}=r\), the number of pivots of
the RREF of \([a_1\ \cdots\ a_n]\).

\begin{theorem}[\citeat{hoffman-kunze}{\S2.3, Theorems~4, 5}, \citeat{lang}{Ch.~I, Theorem~3.4}]\label{thm:r4-counting}
Let \(W\) be a subspace of \(\R^n\) with \(\dim W=d\).
\begin{enumerate}
\item If \(U\) is a subspace of \(W\), then \(\dim U\leqslant d\), and \(\dim U=d\) only for
\(U=W\). In particular \(d\leqslant n\), and \(d=n\) only for \(W=\R^n\).
\item Every linearly independent list of vectors of \(W\) can be extended to a basis of \(W\).
\item More than \(d\) vectors of \(W\) are linearly dependent; fewer than \(d\) vectors do not span
\(W\).
\item \(d\) vectors of \(W\) are linearly independent \(\Leftrightarrow\) they span \(W\)
\(\Leftrightarrow\) they form a basis of \(W\).
\end{enumerate}
\end{theorem}
\begin{proof}
If \(d=0\), then \(W=\{0\}\) (Theorem~\ref{thm:r4-dimension}), and (a)--(c) are immediate: every subspace \(U\) of
\(W\) is \(\{0\}=W\); every vector of \(W\) is \(0\), so a nonempty list of vectors of \(W\) is dependent
(Proposition~\ref{r3-prop-dependence}(a)); and the empty list, the only independent one, is a basis of
\(W\) (Definition~\ref{r3-def-basis}). This also gives (d): the only list of \(0\) vectors is the empty
list, which is independent, spans \(\{0\}=W\) (Definition~\ref{r3-def-span}) and is a basis of \(W\). Let \(d\geqslant1\), and let \(w_1,\dots,w_d\) be a basis of \(W\)
(Theorem~\ref{thm:r4-dimension}). We prove (c) first.

\needspace{3\baselineskip}
(c)
\begin{itemize}
\item The vectors \(w_1,\dots,w_d\) span \(W\), so any \(k>d\) vectors of \(W\) are linearly dependent, by
Theorem~\ref{thm:r4-dimension} (with \(m=d\)).
\item If \(k<d\) vectors spanned \(W\), Theorem~\ref{thm:r4-dimension}, applied to them (with \(m=k\)), would make the
\(d>k\) vectors \(w_1,\dots,w_d\) of \(W\) dependent \(\contradiction\)
\end{itemize}
(b)
\begin{itemize}
\item Let \(v_1,\dots,v_k\) be linearly independent vectors of \(W\). If they span \(W\), they form a basis
of \(W\).
\item Otherwise take \(v\in W\) outside \(\Span\{v_1,\dots,v_k\}\). Then \(v_1,\dots,v_k,v\) are independent:
in a relation \(c_1v_1+\dots+c_kv_k+cv=0\), \(c\ne0\) would put \(v\) in the span, so \(c=0\), and then
every \(c_i=0\).
\item Repeat this. Each list obtained consists of independent vectors of \(W\), so it has at most \(d\)
vectors by (c); hence the process stops, at a basis of \(W\) that contains \(v_1,\dots,v_k\).
\end{itemize}
(a)
\begin{itemize}
\item A basis of \(U\) (Theorem~\ref{thm:r4-dimension}) consists of \(\dim U\) independent vectors of \(W\), so
\(\dim U\leqslant d\) by (c).
\item If \(U\ne W\), take \(v\in W\) outside \(U\). Adjoining \(v\) to a basis of \(U\) gives, as in (b),
\(\dim U+1\) independent vectors of \(W\), so \(\dim U+1\leqslant d\) by (c). Hence \(\dim U=d\) only for
\(U=W\).
\item For the last sentence, apply this to \(W\) as a subspace of \(\R^n\), which has dimension \(n\)
(Theorem~\ref{thm:r4-dimension}).
\end{itemize}
(d) Let \(v_1,\dots,v_d\in W\). A basis of \(W\) is independent and spans \(W\)
(Definition~\ref{r3-def-basis}).
\begin{itemize}
\item Let \(v_1,\dots,v_d\) be independent, and let \(w\in W\). By (c), the \(d+1\) vectors
\(w,v_1,\dots,v_d\) are dependent: \(x_0w+x_1v_1+\dots+x_dv_d=0\) with not all \(x_i=0\). Here
\(x_0\ne0\), since \(v_1,\dots,v_d\) are independent; so \(w=-(x_1/x_0)v_1-\dots-(x_d/x_0)v_d\). Thus
\(v_1,\dots,v_d\) span \(W\), and they form a basis.
\item Let \(v_1,\dots,v_d\) span \(W\). By Theorem~\ref{thm:r4-pivot-columns}, some of them form a basis of
\(W\), which has \(d\) vectors (Theorem~\ref{thm:r4-dimension}); so it consists of all of
\(v_1,\dots,v_d\), and they form a basis.\qedhere
\end{itemize}
\end{proof}

Transposition carries a subspace of \(\R^{1\times n}\) onto a subspace of \(\R^n\) and preserves linear
combinations, hence independence, spanning and bases; so Theorems~\ref{thm:r4-dimension}
and~\ref{thm:r4-counting} also hold for subspaces of \(\R^{1\times n}\), such as the row space
\(\mathcal R(A)\) below.

\begin{example}[cf.\ \citeat{searle-khuri}{Example~1.5}]\label{exm:r4-span}
Let \(W=\Span\{u_1,\dots,u_5\}\subseteq\R^4\), where
\[
u_1=\begin{pmatrix}1\\2\\1\\1\end{pmatrix},\quad
u_2=\begin{pmatrix}1\\3\\2\\1\end{pmatrix},\quad
u_3=\begin{pmatrix}1\\1\\0\\1\end{pmatrix},\quad
u_4=\begin{pmatrix}1\\2\\2\\3\end{pmatrix},\quad
u_5=\begin{pmatrix}1\\3\\1\\-1\end{pmatrix}
\]
(a) Are \(u_1,\dots,u_5\) linearly independent? (b) Find a basis of \(W\) among
\(u_1,\dots,u_5\), and \(\dim W\). (c) Express the remaining vectors in this basis.

\textit{Solution.} Row-reduce \(A=[u_1\ \cdots\ u_5]\):
\[
\begin{gathered}
\begin{pmatrix}1&1&1&1&1\\2&3&1&2&3\\1&2&0&2&1\\1&1&1&3&-1\end{pmatrix}
\xrightarrow[R_4\leftarrow R_4-R_1]{R_2\leftarrow R_2-2R_1,\ R_3\leftarrow R_3-R_1}
\begin{pmatrix}1&1&1&1&1\\0&1&-1&0&1\\0&1&-1&1&0\\0&0&0&2&-2\end{pmatrix}
\\
\xrightarrow{R_3\leftarrow R_3-R_2}
\begin{pmatrix}1&1&1&1&1\\0&1&-1&0&1\\0&0&0&1&-1\\0&0&0&2&-2\end{pmatrix}
\xrightarrow{R_4\leftarrow R_4-2R_3}
\begin{pmatrix}1&1&1&1&1\\0&1&-1&0&1\\0&0&0&1&-1\\0&0&0&0&0\end{pmatrix}
\\
\xrightarrow{R_1\leftarrow R_1-R_3}
\begin{pmatrix}1&1&1&0&2\\0&1&-1&0&1\\0&0&0&1&-1\\0&0&0&0&0\end{pmatrix}
\xrightarrow{R_1\leftarrow R_1-R_2}
\begin{pmatrix}1&0&2&0&1\\0&1&-1&0&1\\0&0&0&1&-1\\0&0&0&0&0\end{pmatrix}=R
\end{gathered}
\]
So \(r=3\), with pivot columns \(k_1=1\), \(k_2=2\), \(k_3=4\).
\begin{enumerate}
\item No: \(r=3<5\) (Theorem~\ref{thm:r4-pivot-test}); indeed any five vectors of \(\R^4\) are
dependent.
\item By Theorem~\ref{thm:r4-pivot-columns}, \(u_1,u_2,u_4\) form a basis of \(W\), and
\(\dim W=3\). As \(3<4\), \(W\neq\R^4\) (Theorem~\ref{thm:r4-counting}(a)).
\item Columns \(3\) and \(5\) of \(R\) give \(u_3=2u_1-u_2\) and \(u_5=u_1+u_2-u_4\)
(Theorem~\ref{thm:r4-pivot-columns}(b)). Check:
\(2(1,2,1,1)'-(1,3,2,1)'=(1,1,0,1)'=u_3\) and
\((1,2,1,1)'+(1,3,2,1)'-(1,2,2,3)'=(1,3,1,-1)'=u_5\). Equivalently,
\(2u_1-u_2-u_3=0\) and \(u_1+u_2-u_4-u_5=0\).
\end{enumerate}
\end{example}

\subsection*{The four subspaces of a matrix}

\begin{definition}[\citeat{harville}{\S4.1}, \citeat{hoffman-kunze}{\S2.2, Examples~7, 10}, \citeat{bapat}{\S2.1}]\label{def:r4-spaces}
For \(A\in\R^{m\times n}\),
\[
\begin{aligned}
\mathcal C(A)&=\{Ax:x\in\R^n\}\subseteq\R^m, &\qquad
\mathcal R(A)&=\{y'A:y\in\R^m\}\subseteq\R^{1\times n}\\
\mathcal N(A)&=\{x\in\R^n:Ax=0\}\subseteq\R^n, &
\mathcal N(A')&=\{y\in\R^m:A'y=0\}\subseteq\R^m
\end{aligned}
\]
are the \textbf{column space}, \textbf{row space}, \textbf{null space} and \textbf{left null space}
of \(A\). The columns of \(A\) span \(\mathcal C(A)\), the rows of \(A\) span \(\mathcal R(A)\) (a
space of row vectors), \(\mathcal N(A)\) is the set of solutions of \(Ax=0\), and
\(y\in\mathcal N(A')\) if and only if \(y'A=0\) (transpose both sides). All four are subspaces
(Theorem~\ref{r3-thm-span} for \(\mathcal C(A)\) and \(\mathcal R(A)\), Example~\ref{r3-exm-subspaces}(e) for
\(\mathcal N(A)\) and \(\mathcal N(A')\)).
\end{definition}

Let \(R\) be the RREF of \(A\in\R^{m\times n}\), with \(r\) nonzero rows, pivot columns
\(k_1<\dots<k_r\) and other columns \(f_1<\dots<f_{n-r}\), so that \(x_{f_1},\dots,x_{f_{n-r}}\) are
the free unknowns of \(Ax=0\) (Proposition~\ref{r2-prop-solutions}(b)).
For \(j=1,\dots,n-r\), the \textbf{special solution} \(s_j\in\R^n\) is the solution of \(Ax=0\)
with free unknowns \(x_{f_j}=1\) and \(x_{f_l}=0\) (\(l\ne j\))~\citeat{hoffman-kunze}{\S2.3, Example~15}. By
Proposition~\ref{r2-prop-solutions}(b), its entries are
\[
(s_j)_{f_l}=\delta_{jl}\quad(1\leqslant l\leqslant n-r),\qquad
(s_j)_{k_i}=-R_{if_j}\quad(1\leqslant i\leqslant r)
\]

\begin{theorem}[\citeat{hoffman-kunze}{\S2.5, Theorems~9, 10; \S2.3, Example~15}]\label{thm:r4-four-bases}
Let \(A\in\R^{m\times n}\), and let row operations take \([A\mid I_m]\) to \([R\mid P]\), where
\(R\) is the RREF of \(A\), with nonzero rows \(\rho_1',\dots,\rho_r'\) (\(\rho_i\in\R^n\)), pivot columns
\(k_1<\dots<k_r\) and free columns \(f_1<\dots<f_{n-r}\). Then \(P\in\R^{m\times m}\) is invertible,
\(PA=R\), and:
\begin{enumerate}
\item if \(r\geqslant1\), the pivot columns \(a_{k_1},\dots,a_{k_r}\) of \(A\) form a basis of
\(\mathcal C(A)\);
\item \(\mathcal R(R)=\mathcal R(A)\), and if \(r\geqslant1\), the nonzero rows \(\rho_1',\dots,\rho_r'\) of \(R\) form
a basis of \(\mathcal R(A)\);
\item the special solutions \(s_1,\dots,s_{n-r}\) form a basis of \(\mathcal N(A)\) (if \(r=n\),
then \(\mathcal N(A)=\{0\}\), with the empty basis);
\item the transposed rows \(\row_{r+1}(P)',\dots,\row_m(P)'\) form a basis of \(\mathcal N(A')\) (if
\(r=m\), then \(\mathcal N(A')=\{0\}\), with the empty basis).
\end{enumerate}
Hence \(\dim\mathcal C(A)=\dim\mathcal R(A)=r\), \(\dim\mathcal N(A)=n-r\) and
\(\dim\mathcal N(A')=m-r\) (if \(r=0\), then \(A=0\) and \(\mathcal C(A)=\mathcal R(A)=\{0\}\)).
\end{theorem}
\begin{proof}
The row operations multiply \([A\mid I_m]\) on the left by an invertible matrix \(E\)
(Theorem~\ref{r2-thm-elementary}(c)), and \(E[A\mid I_m]=[EA\mid E]\), column by column
(Theorem~\ref{r1-thm-views}(a)). So \(P=E\) is invertible and
\(PA=R\).

(a) \(\mathcal C(A)=\Span\{a_1,\dots,a_n\}\); apply Theorem~\ref{thm:r4-pivot-columns}.

(b)
\begin{itemize}
\item Same row space: with \(P=(p_{ij})\), row \(i\) of \(R=PA\) is
\(\row_i(R)=\row_i(P)A=p_{i1}\row_1(A)+\dots+p_{im}\row_m(A)\) (Theorem~\ref{r1-thm-views}(c)), a linear
combination of the rows of \(A\). So every row of \(R\) lies in \(\mathcal R(A)\), and so does their span
\(\mathcal R(R)\) (Theorem~\ref{r3-thm-span}): \(\mathcal R(R)\subseteq\mathcal R(A)\). Likewise \(A=P^{-1}R\) gives
\(\mathcal R(A)\subseteq\mathcal R(R)\). Hence \(\mathcal R(R)=\mathcal R(A)\).
\item Spanning: the rows of \(R\) span \(\mathcal R(R)\), and rows \(r+1,\dots,m\) are zero, so
\(\rho_1',\dots,\rho_r'\) span \(\mathcal R(R)=\mathcal R(A)\).
\item Independent: \(\col_{k_i}(R)=e_i\), so entry \(k_i\) of \(c_1\rho_1'+\dots+c_r\rho_r'\) is
\(c_i\); if this combination is \(0\), every \(c_i=0\).
\end{itemize}
(c) The solutions of \(Ax=0\) are those of \(Rx=0\) (Theorem~\ref{r2-thm-rowops}(b)).
\begin{itemize}
\item Independent: by the definition of the special solutions before the theorem, \((s_j)_{f_l}=\delta_{jl}\),
so the \(f_l\)th entry of \(c_1s_1+\dots+c_{n-r}s_{n-r}\), that is, its value of the unknown \(x_{f_l}\), is
\(\sum_jc_j\delta_{jl}=c_l\). If this combination is \(0\), all its entries are \(0\), so every \(c_l=0\).
\item Spanning: let \(t\in\mathcal N(A)\), with free unknowns \(t_{f_1},\dots,t_{f_{n-r}}\), and put
\(\tilde t=t_{f_1}s_1+\dots+t_{f_{n-r}}s_{n-r}\). Then \(\tilde t\in\mathcal N(A)\), as \(\mathcal N(A)\) is a
subspace (Definition~\ref{def:r4-spaces}).
\item By the previous item, with \(c_l=t_{f_l}\), the \(f_l\)th entry of \(\tilde t\) is \(t_{f_l}\): \(\tilde t\)
has the same free unknowns as \(t\).
\item By Proposition~\ref{r2-prop-solutions}(b), with \(b=0\) and so \(z=0\), the pivot unknowns of a solution are
computed from its free unknowns: \(x_{k_i}=-\sum_lR_{if_l}x_{f_l}\). So two solutions with the same free unknowns are
equal, and \(t=\tilde t\in\Span\{s_1,\dots,s_{n-r}\}\).
\item If \(r=n\), there are no free unknowns, and \(x=0\) is the only solution.
\end{itemize}
(d)
\begin{itemize}
\item In \(\mathcal N(A')\): for \(i>r\), \(\row_i(P)A=\row_i(PA)=\row_i(R)=0\) (Theorem~\ref{r1-thm-views}(c)), so
\(A'\row_i(P)'=(\row_i(P)A)'=0\).
\item Independent: if \(c'P=c_1\row_1(P)+\dots+c_m\row_m(P)=0\), then \(c'=c'PP^{-1}=0\); so the
rows of \(P\) are independent, and so are rows \(r+1,\dots,m\)
(Proposition~\ref{r3-prop-dependence}(b)).
\item Spanning: let \(A'y=0\), so \(y'A=0\), and put \(z'=y'P^{-1}\), so that
\(y'=z'P=\sum_{i=1}^mz_i\row_i(P)\). Then \(0=y'A=z'PA=z'R=z_1\rho_1'+\dots+z_r\rho_r'\), and (b)
gives \(z_1=\dots=z_r=0\). So \(y'=\sum_{i>r}z_i\row_i(P)\) (and \(y=0\) if \(r=m\)).\qedhere
\end{itemize}
\end{proof}

\begin{example}[cf.\ \citeat{searle-khuri}{Example~4.5}]\label{exm:r4-four-bases}
Find bases of the four subspaces of
\[
A=\begin{pmatrix}1&2&1&3\\1&3&0&5\\2&5&1&8\end{pmatrix}
\]
\textit{Solution.} Row-reduce \([A\mid I_3]\):
\[
\begin{gathered}
\left(\begin{array}{cccc|ccc}1&2&1&3&1&0&0\\1&3&0&5&0&1&0\\2&5&1&8&0&0&1\end{array}\right)
\xrightarrow[R_3\leftarrow R_3-2R_1]{R_2\leftarrow R_2-R_1}
\left(\begin{array}{cccc|ccc}1&2&1&3&1&0&0\\0&1&-1&2&-1&1&0\\0&1&-1&2&-2&0&1\end{array}\right)
\\
\xrightarrow{R_3\leftarrow R_3-R_2}
\left(\begin{array}{cccc|ccc}1&2&1&3&1&0&0\\0&1&-1&2&-1&1&0\\0&0&0&0&-1&-1&1\end{array}\right)
\xrightarrow{R_1\leftarrow R_1-2R_2}
\left(\begin{array}{cccc|ccc}1&0&3&-1&3&-2&0\\0&1&-1&2&-1&1&0\\0&0&0&0&-1&-1&1\end{array}\right)
=[R\mid P]
\end{gathered}
\]
So \(r=2\), with pivot columns \(1,2\) and free unknowns \(x_3,x_4\). (The pivots are already
visible in the echelon form after the second arrow; Remark~\ref{rem:r4-echelon}.)
\begin{itemize}
\item \(\mathcal C(A)\): \(a_1=(1,1,2)'\) and \(a_2=(2,3,5)'\) form a basis;
\(\dim\mathcal C(A)=2\). Columns \(3\), \(4\) of \(R\) give
\[
a_3=3a_1-a_2,\qquad a_4=-a_1+2a_2
\]
Check: \(3(1,1,2)'-(2,3,5)'=(1,0,1)'\) and \(-(1,1,2)'+2(2,3,5)'=(3,5,8)'\).
\item \(\mathcal R(A)\): \(\rho_1'=(1,0,3,-1)\) and \(\rho_2'=(0,1,-1,2)\) form a basis;
\(\dim\mathcal R(A)=2\).
\item \(\mathcal N(A)\): \(Rx=0\) reads \(x_1+3x_3-x_4=0\), \(x_2-x_3+2x_4=0\). The special
solutions
\[
s_1=\begin{pmatrix}-3\\1\\1\\0\end{pmatrix}\ (x_3=1,\ x_4=0),\qquad
s_2=\begin{pmatrix}1\\-2\\0\\1\end{pmatrix}\ (x_3=0,\ x_4=1)
\]
form a basis of \(\mathcal N(A)\); \(\dim\mathcal N(A)=2=4-2\). Note \(As_1=-3a_1+a_2+a_3=0\) is the
relation \(a_3=3a_1-a_2\): each special solution is a column relation.
\item \(\mathcal N(A')\): row \(3\) of \(P\) gives the basis \(y=(-1,-1,1)'\);
\(\dim\mathcal N(A')=1=3-2\). Indeed
\(y'A=-\row_1(A)-\row_2(A)+\row_3(A)=0\), that is, \(\row_3(A)=\row_1(A)+\row_2(A)\): each
vector of \(\mathcal N(A')\) is a row relation.
\item Warning: every column of \(R\) has third entry \(0\), so every vector of \(\mathcal C(R)\)
does; but \(a_1\) has third entry \(2\). Thus \(\mathcal C(R)\ne\mathcal C(A)\), and the pivot
columns \(e_1,e_2\) of \(R\) do not span \(\mathcal C(A)\).
\end{itemize}
\end{example}

\begin{remark}\label{rem:r4-warning}
Row operations do not change \(\mathcal R(A)\) or \(\mathcal N(A)\): \(\mathcal R(R)=\mathcal R(A)\)
(Theorem~\ref{thm:r4-four-bases}(b)) and \(\mathcal N(R)=\mathcal N(A)\) (Theorem~\ref{r2-thm-rowops}(b)), so \(R\) may
be used for these two spaces. They do change \(\mathcal C(A)\) in general: a basis of \(\mathcal C(A)\) consists of the
pivot columns of \(A\) itself, not of \(R\) (Example~\ref{exm:r4-four-bases}).
\end{remark}

\subsection*{Rank}

\begin{definition}[\citeat{bapat}{\S2.1}, \citeat{harville}{\S\S4.4a, 4.4c}]\label{def:r4-rank}
By Theorem~\ref{thm:r4-four-bases}, \(\dim\mathcal C(A)=\dim\mathcal R(A)=r\), the number of
pivots of the RREF of \(A\) (also when \(r=0\)). This common number is the \textbf{rank} of \(A\),
written \(\rank A\); by Remark~\ref{rem:r4-echelon} it is also the
number of nonzero rows of any echelon form of \(A\). \(A\) has \textbf{full column rank} if
\(\rank A=n\), and \textbf{full row rank} if \(\rank A=m\). The number \(\dim\mathcal N(A)\) is the
\textbf{nullity} of \(A\).
\end{definition}

\begin{theorem}[\citeat{bapat}{Result~2.10}, \citeat{knapp-basic-algebra}{Ch.~II, Corollary~2.9}]\label{thm:r4-rank-nullity}
For every \(A\in\R^{m\times n}\),
\[
\rank A+\dim\mathcal N(A)=n,\qquad\rank A+\dim\mathcal N(A')=m
\]
\end{theorem}
\begin{proof}
Each of the \(n\) columns of \(R\) is either a pivot column (\(r\) of them) or a free column
(\(n-r\) of them), and by Theorem~\ref{thm:r4-four-bases}(c), \(\dim\mathcal N(A)=n-r\), the number
of free unknowns. Likewise \(\dim\mathcal N(A')=m-r\), the number of zero rows of \(R\)
(Theorem~\ref{thm:r4-four-bases}(d)). Since \(\rank A=r\), both equalities follow.
\end{proof}

\begin{theorem}[\citeat{bapat}{\S2.1; Result~2.11}, \citeat{harville}{Lemma~4.4.3}, \citeat{abadir-magnus}{Exercise~4.7}]\label{thm:r4-rank-rules}
Let \(A\in\R^{m\times n}\).
\begin{enumerate}
\item \(\rank A'=\rank A\), and \(\rank A\leqslant\min\{m,n\}\).
\item \(A\) has full column rank \(\Leftrightarrow\) its columns are linearly independent
\(\Leftrightarrow\) \(\mathcal N(A)=\{0\}\).
\item \(A\) has full row rank \(\Leftrightarrow\) its rows are linearly independent
\(\Leftrightarrow\) \(\mathcal N(A')=\{0\}\) \(\Leftrightarrow\) \(\mathcal C(A)=\R^m\).
\item If \(m=n\), then \(A\) is invertible \(\Leftrightarrow\) \(\rank A=n\).
\end{enumerate}
\end{theorem}
\begin{proof}
Let \(R\) be the RREF of \(A\), with \(r\) nonzero rows and pivot columns \(k_1<\dots<k_r\), so that
\(\rank A=r\) (Definition~\ref{def:r4-rank}).

(a) For \(y\in\R^m\), \((y'A)'=A'y\), so transposing maps \(\mathcal R(A)\) onto
\(\mathcal C(A')\). It preserves linear combinations and is undone by transposing again, so it
carries a basis of \(\mathcal R(A)\) to a basis of \(\mathcal C(A')\), and
\(\rank A'=\dim\mathcal C(A')=\dim\mathcal R(A)=\rank A\). Also \(\rank A=r\leqslant m\), as
\(R\) has \(m\) rows, and \(r\leqslant n\), as \(k_1<\dots<k_r\) lie in \(\{1,\dots,n\}\).

(b) \(A\) has full column rank if and only if \(r=n\) (Definition~\ref{def:r4-rank}); so (b) is
Theorem~\ref{thm:r4-pivot-test}: (iii)\,\(\Leftrightarrow\)\,(i)\,\(\Leftrightarrow\)\,(ii).

(c)
\begin{itemize}
\item \(\rank A=m\) if and only if \(\dim\mathcal N(A')=m-\rank A=0\)
(Theorem~\ref{thm:r4-rank-nullity}), that is, \(\mathcal N(A')=\{0\}\) (Theorem~\ref{thm:r4-dimension}).
\item \(y'A=y_1\row_1(A)+\dots+y_m\row_m(A)\) (Theorem~\ref{r1-thm-views}(c)), and \(y'A=0\) if and only if
\(A'y=0\). So \(\mathcal N(A')=\{0\}\) says exactly that the rows of \(A\) are independent.
\item \(\rank A=\dim\mathcal C(A)\), and a subspace of \(\R^m\) has dimension \(m\) only if it is
\(\R^m\) (Theorem~\ref{thm:r4-counting}(a)); conversely, \(\dim\R^m=m\) (Theorem~\ref{thm:r4-dimension}); so
\(\rank A=m\) if and only if \(\mathcal C(A)=\R^m\).
\end{itemize}
(d) A square matrix \(A\) is invertible if and only if \(Ax=0\) has only the solution \(x=0\)
(Theorem~\ref{r2-thm-invertible}), that is, \(\mathcal N(A)=\{0\}\); by (b), this means \(\rank A=n\).
\end{proof}

\begin{corollary}[\citeat{hoffman-kunze}{\S2.6}, \citeat{harville}{Lemma~4.5.1, Corollary~4.5.2}]\label{cor:r4-ax-b}
Let \(A\in\R^{m\times n}\), \(b\in\R^m\), and let \([A\mid b]\in\R^{m\times(n+1)}\) be the augmented
matrix.
\begin{enumerate}
\item \(Ax=b\) is consistent \(\Leftrightarrow\) \(b\in\mathcal C(A)\) \(\Leftrightarrow\)
\(\rank[A\mid b]=\rank A\).
\item If \(Ax=b\) is consistent, its solution is unique \(\Leftrightarrow\) \(\rank A=n\).
\end{enumerate}
\end{corollary}
\begin{proof}
(a)
\begin{itemize}
\item \(Ax=x_1a_1+\dots+x_na_n\) (Theorem~\ref{r1-thm-views}(b)), so \(Ax=b\) has a solution if and only if \(b\) is
a linear combination of the columns of \(A\), that is, \(b\in\mathcal C(A)\).
\item \(\mathcal C(A)\subseteq\mathcal C([A\mid b])\): a combination of \(a_1,\dots,a_n\) is one of
\(a_1,\dots,a_n,b\), with coefficient \(0\) on \(b\).
\item If \(b=Ac\), then \(x_1a_1+\dots+x_na_n+tb=A(x+tc)\in\mathcal C(A)\); so
\(\mathcal C([A\mid b])=\mathcal C(A)\), and the ranks are equal.
\item If the ranks are equal, the subspace \(\mathcal C(A)\) of \(\mathcal C([A\mid b])\) has the
same dimension, so the two are equal (Theorem~\ref{thm:r4-counting}(a)), and
\(b\in\mathcal C([A\mid b])=\mathcal C(A)\).
\end{itemize}
(b) By Theorem~\ref{thm:r2-solution-structure}(a), the solutions are \(x_p+z\), where \(x_p\) is
one solution and \(z\) runs through the solutions of \(Ax=0\), that is, through \(\mathcal N(A)\). So
the solution is unique if and only if \(\mathcal N(A)=\{0\}\), that is, \(\rank A=n\)
(Theorem~\ref{thm:r4-rank-rules}(b)).
\end{proof}
In practice, consistency is read off from an echelon form of \([A\mid b]\)
(Theorem~\ref{thm:r2-consistency}).

\subsection*{Rank equalities and inequalities}

\begin{theorem}[\citeat{bapat}{Result~2.2; \S2.4}, \citeat{abadir-magnus}{Exercise~4.24}]\label{thm:r4-rank-product}
Let \(A\in\R^{m\times n}\) and \(B\in\R^{n\times p}\).
\begin{enumerate}
\item \(\mathcal C(AB)\subseteq\mathcal C(A)\) and \(\mathcal N(B)\subseteq\mathcal N(AB)\); and
\(\rank(AB)\leqslant\min\{\rank A,\rank B\}\).
\item If \(A\) has full column rank, then \(\mathcal N(AB)=\mathcal N(B)\) and \(\rank(AB)=\rank B\).
\item If \(B\) has full row rank, then \(\mathcal C(AB)=\mathcal C(A)\) and \(\rank(AB)=\rank A\).
\item \(\rank(PAQ)=\rank A\) for all invertible \(P\in\R^{m\times m}\) and \(Q\in\R^{n\times n}\).
\end{enumerate}
\end{theorem}
\begin{proof}
(a)
\begin{itemize}
\item \(ABx=A(Bx)\in\mathcal C(A)\), and \(Bx=0\) gives \(ABx=A0=0\).
\item So \(\rank(AB)=\dim\mathcal C(AB)\leqslant\dim\mathcal C(A)=\rank A\)
(Theorem~\ref{thm:r4-counting}(a)).
\item By this, \((AB)'=B'A'\) (Theorem~\ref{r1-thm-transpose}) and Theorem~\ref{thm:r4-rank-rules}(a),
\(\rank(AB)=\rank(B'A')\leqslant\rank B'=\rank B\).
\end{itemize}
(b)
\begin{itemize}
\item If \(ABx=0\), then \(Bx\in\mathcal N(A)=\{0\}\) (Theorem~\ref{thm:r4-rank-rules}(b)), so
\(x\in\mathcal N(B)\). With (a), \(\mathcal N(AB)=\mathcal N(B)\).
\item \(AB\) and \(B\) both have \(p\) columns, so by Theorem~\ref{thm:r4-rank-nullity},
\(\rank(AB)=p-\dim\mathcal N(AB)=p-\dim\mathcal N(B)=\rank B\).
\end{itemize}
(c)
\begin{itemize}
\item \(\mathcal C(B)=\R^n\) (Theorem~\ref{thm:r4-rank-rules}(c)): every \(x\in\R^n\) is \(Bw\) for
some \(w\in\R^p\), and then \(Ax=(AB)w\in\mathcal C(AB)\).
\item So \(\mathcal C(A)\subseteq\mathcal C(AB)\); with (a), \(\mathcal C(AB)=\mathcal C(A)\), and the
ranks are equal.
\end{itemize}
(d) \(Px=0\) gives \(x=P^{-1}Px=0\), so \(P\) has full column rank
(Theorem~\ref{thm:r4-rank-rules}(b)); and \(\rank Q=n\) (Theorem~\ref{thm:r4-rank-rules}(d)), so
\(Q\) has full row rank. By (b) and (c), \(\rank(P(AQ))=\rank(AQ)=\rank A\).
\end{proof}

\begin{theorem}[\citeat{harville}{Lemma~4.5.1, Corollary~4.5.2}, \citeat{zhang}{Example~2.3.1}]\label{thm:r4-partitioned}
\begin{enumerate}
\item For \(A\in\R^{m\times n}\) and \(B\in\R^{m\times p}\), let \([A\ B]\in\R^{m\times(n+p)}\) have
the columns of \(A\) followed by those of \(B\). Then
\[
\max\{\rank A,\rank B\}\leqslant\rank[A\ B]\leqslant\rank A+\rank B
\]
\item For \(A\in\R^{m\times n}\) and \(C\in\R^{q\times n}\), let
\(\left(\begin{smallmatrix}A\\C\end{smallmatrix}\right)\in\R^{(m+q)\times n}\) have the rows of \(A\)
followed by those of \(C\). Then \(\left(\begin{smallmatrix}A\\C\end{smallmatrix}\right)'=[A'\ C']\), and
\[
\max\{\rank A,\rank C\}\leqslant\rank\begin{pmatrix}A\\C\end{pmatrix}\leqslant\rank A+\rank C
\]
\item For \(A\in\R^{m\times n}\) and \(B\in\R^{q\times p}\), let
\(\diag(A,B)=\left(\begin{smallmatrix}A&0\\0&B\end{smallmatrix}\right)\in\R^{(m+q)\times(n+p)}\), with
zero blocks of sizes \(m\times p\) and \(q\times n\). Then \(\rank\diag(A,B)=\rank A+\rank B\).
\end{enumerate}
\end{theorem}
\begin{proof}
(a)
\begin{itemize}
\item A combination of the columns of \(A\), or of those of \(B\), is a combination of the columns
of \([A\ B]\) (the others with coefficient \(0\)). So \(\mathcal C(A)\) and \(\mathcal C(B)\) are
subspaces of \(\mathcal C([A\ B])\), and Theorem~\ref{thm:r4-counting}(a) gives the lower bound.
\item Let \(r=\rank A\), \(s=\rank B\), and take bases \(v_1,\dots,v_r\) of \(\mathcal C(A)\) and
\(w_1,\dots,w_s\) of \(\mathcal C(B)\). Every column of \([A\ B]\) is a combination of these
\(r+s\) vectors of \(\mathcal C([A\ B])\), so they span \(\mathcal C([A\ B])\) (Theorem~\ref{r3-thm-span}), and
\(\rank[A\ B]\leqslant r+s\) (Theorem~\ref{thm:r4-counting}(c)).
\end{itemize}
(b) Transposing turns the rows of \(\left(\begin{smallmatrix}A\\C\end{smallmatrix}\right)\) into the
columns of \([A'\ C']\), so \(\left(\begin{smallmatrix}A\\C\end{smallmatrix}\right)'=[A'\ C']\).
Apply (a) to \(A'\) and \(C'\), and use Theorem~\ref{thm:r4-rank-rules}(a).

(c) Let \(M=\diag(A,B)\), \(r=\rank A\) and \(s=\rank B\).
\begin{itemize}
\item Upper bound: \(M=[M_1\ M_2]\) with \(M_1=\left(\begin{smallmatrix}A\\0\end{smallmatrix}\right)\)
and \(M_2=\left(\begin{smallmatrix}0\\B\end{smallmatrix}\right)\). By (b), and as a zero matrix has
rank \(0\), \(r\leqslant\rank M_1\leqslant r+0\); so \(\rank M_1=r\), and likewise
\(\rank M_2=s\). By (a), \(\rank M\leqslant r+s\).
\item Lower bound: let \(a_{k_1},\dots,a_{k_r}\) and \(b_{l_1},\dots,b_{l_s}\) be the pivot columns of
\(A\) and of \(B\), bases of \(\mathcal C(A)\) and \(\mathcal C(B)\)
(Theorem~\ref{thm:r4-four-bases}(a)). The corresponding columns
\(\left(\begin{smallmatrix}a_{k_i}\\0\end{smallmatrix}\right)\) and
\(\left(\begin{smallmatrix}0\\b_{l_j}\end{smallmatrix}\right)\) of \(M\) are independent: a combination
of them is \(\left(\begin{smallmatrix}\sum_i\alpha_ia_{k_i}\\\sum_j\beta_jb_{l_j}\end{smallmatrix}\right)\),
which is \(0\) only if all \(\alpha_i\) and \(\beta_j\) are \(0\). These \(r+s\) vectors lie in
\(\mathcal C(M)\), so \(\rank M\geqslant r+s\) (Theorem~\ref{thm:r4-counting}(c)).\qedhere
\end{itemize}
\end{proof}

\begin{corollary}[\citeat{bapat}{Result~2.4}, \citeat{zhang}{Example~2.3.1(2)}]\label{cor:r4-rank-sum}
For \(A,B\in\R^{m\times n}\), \(\rank(A+B)\leqslant\rank A+\rank B\).
\end{corollary}
\begin{proof}
For \(x\in\R^n\), \((A+B)x=Ax+Bx\) is a combination of the columns of \(A\) and of \(B\), so it
lies in \(\mathcal C([A\ B])\). Hence \(\mathcal C(A+B)\subseteq\mathcal C([A\ B])\), so
\(\rank(A+B)\leqslant\rank[A\ B]\) (Theorem~\ref{thm:r4-counting}(a)), and
\(\rank[A\ B]\leqslant\rank A+\rank B\) (Theorem~\ref{thm:r4-partitioned}(a)).
\end{proof}

\begin{theorem}[\citeat{tu-clm}{Theorem~2.4, Step~3}]\label{thm:r4-rank-ata}
Let \(A\in\R^{m\times n}\).
\begin{enumerate}
\item \(\mathcal N(A'A)=\mathcal N(A)\).
\item \(\rank(A'A)=\rank A=\rank(AA')\).
\item \(\mathcal C(A'A)=\mathcal C(A')\) and \(\mathcal C(AA')=\mathcal C(A)\).
\end{enumerate}
\end{theorem}
\begin{proof}
(a)
\begin{itemize}
\item If \(Ax=0\), then \(A'Ax=A'0=0\).
\item If \(A'Ax=0\), then \((Ax)'(Ax)=x'A'Ax=0\). But \((Ax)'(Ax)\) is the sum of the squares of
the entries of \(Ax\) (Theorem~\ref{r1-thm-transpose}), so \(Ax=0\).
\end{itemize}
(b) \(A'A\) and \(A\) both have \(n\) columns, so by (a) and Theorem~\ref{thm:r4-rank-nullity},
\(\rank(A'A)=n-\dim\mathcal N(A'A)=n-\dim\mathcal N(A)=\rank A\). Applied to \(A'\) in place of
\(A\), this gives \(\rank(AA')=\rank A'=\rank A\) (Theorem~\ref{thm:r4-rank-rules}(a)).

(c) \(\mathcal C(A'A)\subseteq\mathcal C(A')\) (Theorem~\ref{thm:r4-rank-product}(a)), and both have
dimension \(\rank A\), by (b) and Theorem~\ref{thm:r4-rank-rules}(a); so they are equal
(Theorem~\ref{thm:r4-counting}(a)). Likewise \(\mathcal C(AA')\subseteq\mathcal C(A)\), both of
dimension \(\rank A\).
\end{proof}

\begin{theorem}[Sylvester's inequality; \citeat{abadir-magnus}{Exercise~5.47(b); Ch.~5, Notes}, \citeat{harville}{(17.5.9)}]\label{thm:r4-sylvester}
For \(A\in\R^{m\times n}\) and \(B\in\R^{n\times p}\),
\[
\dim\mathcal N(AB)\leqslant\dim\mathcal N(A)+\dim\mathcal N(B),\qquad\text{equivalently}\qquad
\rank(AB)\geqslant\rank A+\rank B-n
\]
In particular, if \(AB=0\), then \(\rank A+\rank B\leqslant n\).
\end{theorem}
\begin{proof}
\begin{itemize}
\item Let \(z_1,\dots,z_s\) be a basis of \(\mathcal N(B)\). Since
\(\mathcal N(B)\subseteq\mathcal N(AB)\) (Theorem~\ref{thm:r4-rank-product}(a)), it extends to a
basis \(z_1,\dots,z_s,w_1,\dots,w_t\) of \(\mathcal N(AB)\) (Theorem~\ref{thm:r4-counting}(b));
so \(\dim\mathcal N(AB)=s+t\).
\item The vectors \(Bw_1,\dots,Bw_t\) lie in \(\mathcal N(A)\), since \(A(Bw_j)=(AB)w_j=0\).
\item They are independent: if \(c_1Bw_1+\dots+c_tBw_t=0\), then \(B(c_1w_1+\dots+c_tw_t)=0\), so
\(c_1w_1+\dots+c_tw_t\in\mathcal N(B)\) equals some \(\lambda_1z_1+\dots+\lambda_sz_s\); as
\(z_1,\dots,z_s,w_1,\dots,w_t\) are independent, every \(c_j=0\).
\item Hence \(t\leqslant\dim\mathcal N(A)\) (Theorem~\ref{thm:r4-counting}(c)), and
\(\dim\mathcal N(AB)=s+t\leqslant\dim\mathcal N(B)+\dim\mathcal N(A)\).
\item By Theorem~\ref{thm:r4-rank-nullity}, \(\dim\mathcal N(AB)=p-\rank(AB)\),
\(\dim\mathcal N(B)=p-\rank B\) and \(\dim\mathcal N(A)=n-\rank A\). Substituting gives
\(p-\rank(AB)\leqslant(p-\rank B)+(n-\rank A)\), which is the rank form. If \(AB=0\), then
\(\rank(AB)=0\).\qedhere
\end{itemize}
\end{proof}

\begin{example}\label{exm:r4-bounds}
Let \(A=\diag(1,1,1,0)\) and \(B=\diag(0,0,1,1)\) in \(\R^{4\times4}\). Their columns are
\(e_1,e_2,e_3,0\) and \(0,0,e_3,e_4\). A zero column adds nothing to the span
(Theorem~\ref{r3-thm-span-inclusion}(c)), so \(\mathcal C(A)=\Span\{e_1,e_2,e_3\}\) and
\(\mathcal C(B)=\Span\{e_3,e_4\}\); unit vectors are independent, as part of the standard basis
(Definition~\ref{r3-def-basis} and Proposition~\ref{r3-prop-dependence}(b)). So \(\rank A=3\) and \(\rank B=2\).
\begin{itemize}
\item \(AB=\diag(0,0,1,0)\) (row \(i\) of \(AB\) is \(a_{ii}\row_i(B)\), Theorem~\ref{r1-thm-views}(c)) has the
columns \(0,0,e_3,0\); in the same way, \(\mathcal C(AB)=\Span\{e_3\}\) and \(\rank(AB)=1\): less than
\(\min\{3,2\}\) (Theorem~\ref{thm:r4-rank-product}(a)), and equal to \(\rank A+\rank B-4=1\), so Sylvester's bound
is attained (Theorem~\ref{thm:r4-sylvester}).
\item \(A+B=\diag(1,1,2,1)\) is invertible, with inverse \(\diag(1,1,\frac12,1)\) (Definition~\ref{r2-def-inverse}); so
\(\rank(A+B)=4<3+2\) (Theorem~\ref{thm:r4-rank-rules}(d) and Corollary~\ref{cor:r4-rank-sum}). And \(A+(-A)=0\) has
rank \(0\), as \(\mathcal C(0)=\{0\}\) and \(\dim\{0\}=0\) (Theorem~\ref{thm:r4-dimension}).
\item \([A\ B]\in\R^{4\times8}\) has the columns \(e_1,e_2,e_3,0,0,0,e_3,e_4\), so
\(\mathcal C([A\ B])=\R^4\) and \(\rank[A\ B]=\dim\R^4=4\) (Theorem~\ref{thm:r4-dimension}), strictly between
\(\max\{3,2\}\) and \(3+2\) (Theorem~\ref{thm:r4-partitioned}(a)).
\item \(\diag(A,B)=\diag(1,1,1,0,0,0,1,1)\in\R^{8\times8}\) has rank \(5=3+2\)
(Theorem~\ref{thm:r4-partitioned}(c)).
\item \(A'A=A^2=A\) has rank \(3=\rank A\) (Theorem~\ref{thm:r4-rank-ata}).
\end{itemize}
\end{example}

\begin{theorem}[\citeat{bapat}{Results~2.3, 2.13}, \citeat{harville}{Theorem~4.4.8}, \citeat{searle-khuri}{(4.42)}]\label{thm:r4-full-rank-factorization}
Let \(A\in\R^{m\times n}\) have rank \(r\). If \(r\geqslant1\), a product \(A=BC\) with
\(B\in\R^{m\times r}\) of full column rank and \(C\in\R^{r\times n}\) of full row rank is a
\textbf{full-rank factorization} of \(A\).
\begin{enumerate}
\item If \(r\geqslant1\), let \(R\) be the RREF of \(A\), with pivot columns \(k_1<\dots<k_r\); let
\(B=[a_{k_1}\ \cdots\ a_{k_r}]\in\R^{m\times r}\) consist of the pivot columns of \(A\), and
\(C\in\R^{r\times n}\) of the \(r\) nonzero rows of \(R\). Then \(A=BC\) is a full-rank
factorization. If \(r=0\), then \(A=0\), and there is nothing to factor.
\item For every full-rank factorization \(A=BC\), \(\mathcal C(A)=\mathcal C(B)\) and
\(\mathcal N(A)=\mathcal N(C)\).
\item For every full-rank factorization \(A=BC\), the \(r\times r\) matrices \(B'B\) and \(CC'\) are
invertible, and
\[
L=(B'B)^{-1}B'\in\R^{r\times m},\qquad K=C'(CC')^{-1}\in\R^{n\times r}
\]
satisfy \(LB=I_r\) and \(CK=I_r\): \(B\) has a left inverse and \(C\) has a right inverse.
\end{enumerate}
\end{theorem}
\begin{proof}
(a)
\begin{itemize}
\item Column \(j\) of \(BC\) is \(B\col_j(C)=R_{1j}a_{k_1}+\dots+R_{rj}a_{k_r}\)
(Theorem~\ref{r1-thm-views}(a),(b)), which
is \(a_j\) by Theorem~\ref{thm:r4-pivot-columns}(b). So \(BC=A\).
\item The columns of \(B\) form a basis of \(\mathcal C(A)\), so they are independent: \(B\) has full
column rank (Theorem~\ref{thm:r4-rank-rules}(b)).
\item The rows of \(C\) form a basis of \(\mathcal R(A)\) (Theorem~\ref{thm:r4-four-bases}(b)), so
they are independent: \(C\) has full row rank (Theorem~\ref{thm:r4-rank-rules}(c)).
\item If \(r=0\), then \(A=0\) (Theorem~\ref{thm:r4-four-bases}).
\end{itemize}
(b) Theorem~\ref{thm:r4-rank-product}(b) gives \(\mathcal N(BC)=\mathcal N(C)\), as \(B\) has full
column rank, and Theorem~\ref{thm:r4-rank-product}(c) gives \(\mathcal C(BC)=\mathcal C(B)\), as
\(C\) has full row rank.

(c)
\begin{itemize}
\item \(\rank(B'B)=\rank B=r\) and \(\rank(CC')=\rank C=r\) (Theorem~\ref{thm:r4-rank-ata}(b)), so
the \(r\times r\) matrices \(B'B\) and \(CC'\) are invertible (Theorem~\ref{thm:r4-rank-rules}(d)).
\item \(LB=(B'B)^{-1}(B'B)=I_r\) and \(CK=(CC')(CC')^{-1}=I_r\).\qedhere
\end{itemize}
\end{proof}

\begin{example}\label{exm:r4-factorization}
For \(A\) of Example~\ref{exm:r4-four-bases}, Theorem~\ref{thm:r4-full-rank-factorization}(a) gives
\[
\begin{pmatrix}1&2&1&3\\1&3&0&5\\2&5&1&8\end{pmatrix}
=\begin{pmatrix}1&2\\1&3\\2&5\end{pmatrix}\begin{pmatrix}1&0&3&-1\\0&1&-1&2\end{pmatrix}
\]
\begin{itemize}
\item Columns \(1\), \(2\) of \(BC\) are \(a_1\), \(a_2\), and columns \(3\), \(4\) are
\(3a_1-a_2=a_3\) and \(-a_1+2a_2=a_4\) (Example~\ref{exm:r4-four-bases}).
\item By Theorem~\ref{thm:r4-full-rank-factorization}(b),
\(\mathcal C(A)=\mathcal C(B)=\Span\{a_1,a_2\}\) and
\(\mathcal N(A)=\mathcal N(C)=\Span\{s_1,s_2\}\): the system \(Cx=0\) is \(Rx=0\) without its zero
row.
\item The \(4\times2\) matrix \([s_1\ s_2]\) has \(A[s_1\ s_2]=0\) and \(\rank[s_1\ s_2]=2\), and Sylvester's
inequality \(\rank A+\rank[s_1\ s_2]\leqslant4\) holds with equality: \(2+2=4\).
\end{itemize}
\end{example}

\subsection*{Sums and intersections of subspaces}

\begin{theorem}[\citeat{hoffman-kunze}{\S2.3, Theorem~6}]\label{thm:r4-dim-sum}
Let \(U\) and \(W\) be subspaces of \(\R^n\). Then
\[
\dim(U+W)+\dim(U\cap W)=\dim U+\dim W
\]
\end{theorem}
\begin{proof}
\(U\cap W\) and \(U+W\) are subspaces of \(\R^n\) (Theorem~\ref{r3-thm-sum}(a)). The lists below may be empty
(the empty list is a basis of \(\{0\}\), Definition~\ref{r3-def-basis}).
\begin{itemize}
\item Bases: let \(x_1,\dots,x_k\) be a basis of \(U\cap W\) (Theorem~\ref{thm:r4-dimension}). It is an independent
list of vectors of \(U\), and also of \(W\); so it extends to a basis \(x_1,\dots,x_k,u_1,\dots,u_p\) of \(U\) and to a
basis \(x_1,\dots,x_k,w_1,\dots,w_q\) of \(W\) (Theorem~\ref{thm:r4-counting}(b)). Thus \(\dim U=k+p\) and
\(\dim W=k+q\).
\item Spanning: a vector of \(U+W\) is \(u+w\) with \(u\in U\) and \(w\in W\), that is, a combination of
\(x_1,\dots,x_k,u_1,\dots,u_p\) plus a combination of \(x_1,\dots,x_k,w_1,\dots,w_q\); so it is a combination of the
\(k+p+q\) vectors
\[
x_1,\dots,x_k,\ u_1,\dots,u_p,\ w_1,\dots,w_q
\]
These vectors lie in \(U+W\) (Theorem~\ref{r3-thm-sum}(b)), so they span \(U+W\).
\item Independent: let \(\sum_ia_ix_i+\sum_jb_ju_j+\sum_lc_lw_l=0\). Then
\(\sum_lc_lw_l=-\sum_ia_ix_i-\sum_jb_ju_j\) lies in \(U\) as well as in \(W\), hence in \(U\cap W\); so
\(\sum_lc_lw_l=\sum_i\lambda_ix_i\) for some scalars \(\lambda_i\). Since \(x_1,\dots,x_k,w_1,\dots,w_q\) are
independent, the relation \(\sum_i\lambda_ix_i-\sum_lc_lw_l=0\) forces every \(c_l=0\). Then
\(\sum_ia_ix_i+\sum_jb_ju_j=0\), and the
independence of \(x_1,\dots,x_k,u_1,\dots,u_p\) gives every \(a_i=0\) and every \(b_j=0\).
\item So these \(k+p+q\) vectors form a basis of \(U+W\), and
\[
\dim(U+W)+\dim(U\cap W)=(k+p+q)+k=(k+p)+(k+q)=\dim U+\dim W\qedhere
\]
\end{itemize}
\end{proof}

\begin{example}[\citeat{hoffman-kunze}{\S2.3, Exercise~7(b)}]\label{exm:r4-dim-sum}
Write a \(2\times2\) matrix \(\left(\begin{smallmatrix}a_{11}&a_{12}\\a_{21}&a_{22}\end{smallmatrix}\right)\) as the
column \((a_{11},a_{12},a_{21},a_{22})'\) of its entries. The matrices
\(\left(\begin{smallmatrix}x&-x\\y&z\end{smallmatrix}\right)\) and
\(\left(\begin{smallmatrix}a&b\\-a&c\end{smallmatrix}\right)\) (\(a,b,c,x,y,z\in\R\)) then form the subspaces
\[
U=\{(x,-x,y,z)':x,y,z\in\R\},\qquad W=\{(a,b,-a,c)':a,b,c\in\R\}
\]
of \(\R^4\): \(U\) consists of the vectors whose second entry is minus the first, and \(W\) of those whose third entry
is minus the first; so \(U\) and \(W\) are the solution sets of \(x_1+x_2=0\) and of \(x_1+x_3=0\), subspaces of
\(\R^4\) by Example~\ref{r3-exm-subspaces}(e).
\begin{itemize}
\item \((x,-x,y,z)'=x\,u_1+y\,e_3+z\,e_4\) with \(u_1=(1,-1,0,0)'\), and this vector is \(0\) only for
\(x=y=z=0\); so \(u_1,e_3,e_4\) form a basis of \(U\) (Definition~\ref{r3-def-basis}), and \(\dim U=3\)
(Theorem~\ref{thm:r4-dimension}). Likewise
\((a,b,-a,c)'=a\,w_1+b\,e_2+c\,e_4\) with \(w_1=(1,0,-1,0)'\), so \(w_1,e_2,e_4\) form a basis of \(W\), and
\(\dim W=3\).
\item A vector \((x,-x,y,z)'\) of \(U\) lies in \(W\) if and only if \(y=-x\). So
\(U\cap W=\{(x,-x,-x,z)':x,z\in\R\}\), and \((x,-x,-x,z)'=x\,(1,-1,-1,0)'+z\,e_4\) is \(0\) only for \(x=z=0\);
so \((1,-1,-1,0)'\) and \(e_4\) form a basis of \(U\cap W\), a subspace by Theorem~\ref{r3-thm-sum}(a),
and \(\dim(U\cap W)=2\) (Theorem~\ref{thm:r4-dimension}).
\item By Theorem~\ref{thm:r4-dim-sum}, \(\dim(U+W)=\dim U+\dim W-\dim(U\cap W)=3+3-2=4\), so \(U+W=\R^4\)
(Theorem~\ref{thm:r4-counting}(a)). Directly: \(e_3,e_4\in U\) and \(e_2\in W\) lie in \(U+W\)
(Theorem~\ref{r3-thm-sum}(b)), and \(e_1=u_1+e_2\in U+W\); so the
subspace \(U+W\) (Theorem~\ref{r3-thm-sum}(a)) contains every \(x=x_1e_1+x_2e_2+x_3e_3+x_4e_4\) of \(\R^4\).
\end{itemize}
With this direct computation of \(U+W\), \(\dim(U+W)+\dim(U\cap W)=4+2=3+3=\dim U+\dim W\).
\end{example}

\begin{corollary}[\citeat{harville}{Corollary~17.1.6}, \citeat{bapat}{Ch.~1, Exercise~19}, \citeat{lang}{Ch.~I, Theorem~4.2}]\label{cor:r4-dim-sum}
Let \(U\) and \(W\) be subspaces of \(\R^n\).
\begin{enumerate}
\item \(\dim(U+W)\leqslant\dim U+\dim W\), with equality if and only if \(U\cap W=\{0\}\). In particular
\(\dim(U\oplus W)=\dim U+\dim W\).
\item If \(\dim U+\dim W>n\), then \(U\cap W\ne\{0\}\).
\item \(W\) has a \textbf{complement} in \(\R^n\), that is, a subspace \(V\) with \(\R^n=W\oplus V\); and every
complement \(V\) of \(W\) has \(\dim V=n-\dim W\).
\end{enumerate}
\end{corollary}
\begin{proof}
(a) Let \(u_1,\dots,u_p\) and \(w_1,\dots,w_q\) be bases of \(U\) and of \(W\) (Theorem~\ref{thm:r4-dimension}).
\begin{itemize}
\item If \(U\cap W=\{0\}\), then \(u_1,\dots,u_p,w_1,\dots,w_q\) form a basis of \(U\oplus W\)
(Theorem~\ref{r3-thm-direct-sum-basis}), so \(\dim(U\oplus W)=p+q\).
\item If \(U\cap W\ne\{0\}\), these \(p+q\) vectors are dependent (Theorem~\ref{r3-thm-direct-sum-basis}). They span
\(U+W\) (Theorem~\ref{r3-thm-sum}(c)), so some of them form a basis of \(U+W\)
(Theorem~\ref{thm:r4-pivot-columns}), and not all of them, as a basis is independent. Hence \(\dim(U+W)<p+q\).
\end{itemize}
(b) If \(U\cap W=\{0\}\), then (a) and Theorem~\ref{thm:r4-counting}(a), applied to the subspace \(U+W\) of \(\R^n\),
give \(n<\dim U+\dim W=\dim(U+W)\leqslant n\) \(\contradiction\)

\needspace{4\baselineskip}
(c)
\begin{itemize}
\item Let \(w_1,\dots,w_d\) be a basis of \(W\) (Theorem~\ref{thm:r4-dimension}). By
Theorem~\ref{thm:r4-counting}(b), applied to \(\R^n\), it extends to a basis \(w_1,\dots,w_d,v_1,\dots,v_{n-d}\) of
\(\R^n\), which has \(n\) vectors (Theorem~\ref{thm:r4-dimension}). Let \(V=\Span\{v_1,\dots,v_{n-d}\}\)
(\(V=\{0\}\) if \(d=n\)).
\item \(W+V=\Span\{w_1,\dots,w_d,v_1,\dots,v_{n-d}\}=\R^n\) (Theorem~\ref{r3-thm-sum}(c)).
\item \(W\cap V=\{0\}\): if \(x=\sum_ia_iw_i=\sum_jb_jv_j\), then \(\sum_ia_iw_i-\sum_jb_jv_j=0\), so every \(a_i\)
and every \(b_j\) is \(0\), and \(x=0\).
\item Hence \(\R^n=W\oplus V\) (Definition~\ref{r3-def-sum}). If \(V\) is any complement of \(W\), then
\(n=\dim(W\oplus V)=\dim W+\dim V\) by (a).\qedhere
\end{itemize}
\end{proof}
A complement is not unique in general: in \(\R^4\), both \(\Span\{\one_4\}\) and \(\Span\{e_4\}\) are complements of
the centered vectors \(W_0\) (Examples~\ref{r3-exm-constant-centered} and~\ref{r3-exm-sum-r3}(b)).

\needspace{12\baselineskip}
\subsection*{Column spaces and the system \(AX=B\)}

\begin{proposition}[\citeat{harville}{Lemmas~4.2.1, 4.2.2}]\label{prop:r4-inclusion}
Let \(A\in\R^{m\times n}\).
\begin{enumerate}
\item For \(B\in\R^{m\times p}\) the following are equivalent: \(\mathcal C(B)\subseteq\mathcal C(A)\); every column
of \(B\) lies in \(\mathcal C(A)\); \(B=AF\) for some \(F\in\R^{n\times p}\).
\item For \(C\in\R^{q\times n}\) the following are equivalent: \(\mathcal R(C)\subseteq\mathcal R(A)\); every row of
\(C\) lies in \(\mathcal R(A)\); \(C=LA\) for some \(L\in\R^{q\times m}\).
\end{enumerate}
\end{proposition}
\begin{proof}
(a) The columns of \(A\) span \(\mathcal C(A)\), and those of \(B\) span \(\mathcal C(B)\)
(Definition~\ref{def:r4-spaces}); so (a) is Theorem~\ref{r3-thm-span-inclusion}(a).

(b)
\begin{itemize}
\item If \(\mathcal R(C)\subseteq\mathcal R(A)\), then \(\row_i(C)=e_i'C\in\mathcal R(C)\subseteq\mathcal R(A)\)
(Theorem~\ref{r1-thm-views}(c), with \(e_i\in\R^q\)). Conversely, if every row of \(C\) lies in the subspace
\(\mathcal R(A)\), so does every \(y'C=y_1\row_1(C)+\dots+y_q\row_q(C)\) (Theorem~\ref{r1-thm-views}(c)).
\item For \(L\in\R^{q\times m}\) with rows \(l_1',\dots,l_q'\), \(\row_i(LA)=l_i'A\)
(Theorem~\ref{r1-thm-views}(c)). So \(C=LA\) for some \(L\) if and only if every \(\row_i(C)\) has the form \(l'A\)
with \(l\in\R^m\), that is, lies in \(\mathcal R(A)=\{y'A:y\in\R^m\}\).\qedhere
\end{itemize}
\end{proof}

\begin{theorem}[\citeat{harville}{Theorem~17.4.2, Lemma~4.5.1, Corollaries~4.5.2, 17.1.6}]\label{thm:r4-rank-concat}
Let \(A\in\R^{m\times n}\) and \(B\in\R^{m\times p}\), and let \([A\ B]\) be as in
Theorem~\ref{thm:r4-partitioned}(a).
\begin{enumerate}
\item \(\mathcal C([A\ B])=\mathcal C(A)+\mathcal C(B)\), and
\[
\rank[A\ B]=\rank A+\rank B-\dim\bigl(\mathcal C(A)\cap\mathcal C(B)\bigr)
\]
\item \(\rank[A\ B]=\rank A\) \(\Leftrightarrow\) \(\mathcal C([A\ B])=\mathcal C(A)\) \(\Leftrightarrow\)
\(\mathcal C(B)\subseteq\mathcal C(A)\).
\item \(\rank[A\ B]=\rank A+\rank B\) \(\Leftrightarrow\) \(\mathcal C(A)\cap\mathcal C(B)=\{0\}\).
\item For \(C\in\R^{q\times n}\) and \(\left(\begin{smallmatrix}A\\C\end{smallmatrix}\right)\) as in
Theorem~\ref{thm:r4-partitioned}(b),
\[
\rank\begin{pmatrix}A\\C\end{pmatrix}=\rank A+\rank C-\dim\bigl(\mathcal R(A)\cap\mathcal R(C)\bigr)
\]
and \(\rank\left(\begin{smallmatrix}A\\C\end{smallmatrix}\right)=\rank A\) if and only if
\(\mathcal R(C)\subseteq\mathcal R(A)\).
\end{enumerate}
\end{theorem}
\begin{proof}
(a) The columns of \([A\ B]\) are those of \(A\) followed by those of \(B\), so
\(\mathcal C([A\ B])=\mathcal C(A)+\mathcal C(B)\) (Theorem~\ref{r3-thm-sum}(c) and Definition~\ref{def:r4-spaces}).
Theorem~\ref{thm:r4-dim-sum} with \(U=\mathcal C(A)\) and \(W=\mathcal C(B)\) gives the formula, since
\(\rank M=\dim\mathcal C(M)\) for every matrix \(M\) (Definition~\ref{def:r4-rank}).

(b)
\begin{itemize}
\item By (a), \(\mathcal C(A)\) and \(\mathcal C(B)\) lie in \(\mathcal C([A\ B])\) (Theorem~\ref{r3-thm-sum}(b)).
\item If \(\mathcal C(B)\subseteq\mathcal C(A)\), then \(B=AF\) for some \(F=[f_1\ \cdots\ f_p]\)
(Proposition~\ref{prop:r4-inclusion}(a)), and the columns of \([A\ B]\) are \(a_1,\dots,a_n,Af_1,\dots,Af_p\)
(Theorem~\ref{r1-thm-views}(a)), all in \(\mathcal C(A)\). So \(\mathcal C([A\ B])\subseteq\mathcal C(A)\)
(Proposition~\ref{prop:r4-inclusion}(a)), and with the first item the two spaces are equal. Conversely, if
\(\mathcal C([A\ B])=\mathcal C(A)\), then \(\mathcal C(B)\subseteq\mathcal C([A\ B])=\mathcal C(A)\).
\item Equal column spaces have equal dimensions, so \(\mathcal C([A\ B])=\mathcal C(A)\) gives
\(\rank[A\ B]=\rank A\). Conversely, if \(\rank[A\ B]=\rank A\), the subspace \(\mathcal C(A)\) of
\(\mathcal C([A\ B])\) has the same dimension, so the two are equal (Theorem~\ref{thm:r4-counting}(a)).
\end{itemize}
(c) This is Corollary~\ref{cor:r4-dim-sum}(a) for \(U=\mathcal C(A)\) and \(W=\mathcal C(B)\), by (a).

\needspace{4\baselineskip}
(d)
\begin{itemize}
\item \(\left(\begin{smallmatrix}A\\C\end{smallmatrix}\right)'=[A'\ C']\)
(Theorem~\ref{thm:r4-partitioned}(b)), and a matrix and its transpose have the same rank
(Theorem~\ref{thm:r4-rank-rules}(a)).
\item For every matrix \(G\) with \(n\) columns, \((y'G)'=G'y\), so a row vector \(z\) lies in \(\mathcal R(G)\) if and
only if \(z'\in\mathcal C(G')\). Hence transposition maps \(\mathcal R(A)\cap\mathcal R(C)\) onto
\(\mathcal C(A')\cap\mathcal C(C')\), and \(\mathcal R(C)\subseteq\mathcal R(A)\) holds if and only if
\(\mathcal C(C')\subseteq\mathcal C(A')\). Transposition preserves linear combinations and is undone by transposing
again, so it carries bases to bases, and the two intersections have the same dimension.
\item Now apply (a) and (b) to \(A'\in\R^{n\times m}\) and \(C'\in\R^{n\times q}\).\qedhere
\end{itemize}
\end{proof}

For \(A\in\R^{m\times n}\) and \(B\in\R^{m\times p}\), the \textbf{system} \(AX=B\) asks for the matrices
\(X\in\R^{n\times p}\) with \(AX=B\), its \textbf{solutions}; it is \textbf{consistent} if it has a solution~\citeat{harville}{\S\S7.1, 7.2}. By
Theorem~\ref{r1-thm-views}(a), \(X=[x_1\ \cdots\ x_p]\) is a solution if and only if \(Ax_j=\col_j(B)\) for
\(j=1,\dots,p\): the system collects \(p\) systems with the same coefficient matrix. For \(p=1\) it is the system
\(Ax=b\) of Definition~\ref{r2-def-system}.

\begin{corollary}[\citeat{harville}{Theorems~7.2.1, 7.2.2}]\label{cor:r4-matrix-system}
Let \(A\in\R^{m\times n}\) and \(B\in\R^{m\times p}\). The following are equivalent:
\begin{enumerate}[label=(\roman*)]
\item the system \(AX=B\) is consistent;
\item \(\mathcal C(B)\subseteq\mathcal C(A)\), that is, every column of \(B\) lies in \(\mathcal C(A)\);
\item \(\mathcal C([A\ B])=\mathcal C(A)\);
\item \(\rank[A\ B]=\rank A\).
\end{enumerate}
If \(A\) has full row rank, then \(AX=B\) is consistent for every \(B\in\R^{m\times p}\).
\end{corollary}
\begin{proof}
\(AX=B\) is consistent if and only if \(B=AF\) for some \(F\in\R^{n\times p}\), which is (ii) by
Proposition~\ref{prop:r4-inclusion}(a); and (ii) \(\Leftrightarrow\) (iii) \(\Leftrightarrow\) (iv) is
Theorem~\ref{thm:r4-rank-concat}(b). If \(\rank A=m\), then
\(m=\rank A\leqslant\rank[A\ B]\leqslant m\) (Theorems~\ref{thm:r4-partitioned}(a) and~\ref{thm:r4-rank-rules}(a)),
so (iv) holds.
\end{proof}

\begin{theorem}[\citeat{harville}{Theorem~7.3.1}]\label{thm:r4-compatibility}
Let \(A\in\R^{m\times n}\) and \(B\in\R^{m\times p}\). The system \(AX=B\) is consistent if and only if \(y'B=0\) for
every \(y\in\mathcal N(A')\), that is, if and only if every linear relation
\(y_1\row_1(A)+\dots+y_m\row_m(A)=0\) among the rows of \(A\) also holds among the rows of \(B\). In particular,
\(Ax=b\) is consistent if and only if \(y'b=0\) for every \(y\in\mathcal N(A')\).
\end{theorem}
\begin{proof}
If \(AX_0=B\) and \(y\in\mathcal N(A')\), that is, \(y'A=0\) (Definition~\ref{def:r4-spaces}), then
\(y'B=y'AX_0=0\). Conversely, let \(y'B=0\) for every \(y\in\mathcal N(A')\).
\begin{itemize}
\item If \(A=0\), every \(y\in\R^m\) lies in \(\mathcal N(A')\); with \(y=e_i\), \(\row_i(B)=e_i'B=0\) for every \(i\)
(Theorem~\ref{r1-thm-views}(c)), so \(B=0\) and \(X=0\) is a solution. Let \(A\ne0\) and \(r=\rank A\geqslant1\).
\item The rows of \(A\) span \(\mathcal R(A)\), so \(r\) of them, \(\row_{h_1}(A),\dots,\row_{h_r}(A)\) with
\(h_1<\dots<h_r\), form a basis of \(\mathcal R(A)\) (Theorem~\ref{thm:r4-pivot-columns}, last sentence, for row
vectors by transposition, as in the paragraph after Theorem~\ref{thm:r4-counting}). Let \(A_h\in\R^{r\times n}\) and
\(B_h\in\R^{r\times p}\) consist of rows \(h_1,\dots,h_r\) of \(A\) and of \(B\). The rows of \(A_h\) are independent,
so \(A_h\) has full row rank (Theorem~\ref{thm:r4-rank-rules}(c)), and \(A_hX=B_h\) has a solution \(X^*\)
(Corollary~\ref{cor:r4-matrix-system}): \(\row_{h_j}(A)X^*=\row_{h_j}(B)\) for \(j=1,\dots,r\)
(Theorem~\ref{r1-thm-views}(c)).
\item Let \(i\notin\{h_1,\dots,h_r\}\). Then \(\row_i(A)=\sum_j\lambda_{ij}\row_{h_j}(A)\) for some scalars
\(\lambda_{ij}\), so \(y=e_i-\sum_j\lambda_{ij}e_{h_j}\in\R^m\) has \(y'A=\row_i(A)-\sum_j\lambda_{ij}\row_{h_j}(A)=0\),
that is, \(y\in\mathcal N(A')\). Hence \(0=y'B=\row_i(B)-\sum_j\lambda_{ij}\row_{h_j}(B)\), and
\[
\row_i(A)X^*=\sum_j\lambda_{ij}\row_{h_j}(A)X^*=\sum_j\lambda_{ij}\row_{h_j}(B)=\row_i(B)
\]
\item Thus \(\row_i(AX^*)=\row_i(A)X^*=\row_i(B)\) for every \(i\) (Theorem~\ref{r1-thm-views}(c)), that is,
\(AX^*=B\). The last sentence is the case \(p=1\).\qedhere
\end{itemize}
\end{proof}

In practice, \(\mathcal C(A)\) is described by equations obtained from the reduction of \([A\mid I_m]\) to
\([R\mid P]\) (Theorem~\ref{thm:r4-four-bases}). The same row operations take \([A\mid b]\) to \([R\mid Pb]\): by
Theorem~\ref{r2-thm-elementary}(c) they multiply on the left by an invertible \(E\), and \(E[A\mid I_m]=[EA\mid E]\),
column by column (Theorem~\ref{r1-thm-views}(a)), so \(E=P\), \(EA=R\) and \(E[A\mid b]=[R\mid Pb]\). As
\(b\in\mathcal C(A)\) if and only if \(Ax=b\) is consistent (Corollary~\ref{cor:r4-ax-b}(a)),
Proposition~\ref{r2-prop-solutions}(a) gives, with \(r=\rank A\),
\[
b\in\mathcal C(A)\iff\row_i(P)\,b=0\quad\text{for }i=r+1,\dots,m
\]
(no condition if \(r=m\); Example~\ref{r2-exm-row-equivalent} is an instance)~\citeat{hoffman-kunze}{\S2.6, Example~22(h)}. The vectors \(\row_i(P)'\)
(\(i>r\)) form a basis of \(\mathcal N(A')\)
(Theorem~\ref{thm:r4-four-bases}(d)), so these \(m-r\) equations are the conditions \(y'b=0\) of
Theorem~\ref{thm:r4-compatibility}.

\begin{example}\label{exm:r4-compatibility}
For \(A\) of Example~\ref{exm:r4-four-bases}, \(r=2\) and \(\row_3(P)=(-1,-1,1)\), so
\[
b\in\mathcal C(A)\iff -b_1-b_2+b_3=0
\]
This is the row relation \(\row_3(A)=\row_1(A)+\row_2(A)\) found there, which every right side must share
(Theorem~\ref{thm:r4-compatibility}).
\begin{enumerate}
\item \(Ax=(1,t,4)'\) is consistent if and only if \(-1-t+4=0\), that is, \(t=3\). For \(t=3\) and \(b=(1,3,4)'\),
\(Pb=(3-6,\ -1+3,\ -1-3+4)'=(-3,2,0)'\), and \(Rx=Pb\) reads \(x_1+3x_3-x_4=-3\), \(x_2-x_3+2x_4=2\). The free
unknowns \(x_3=x_4=0\) give the solution \(x_p=(-3,2,0,0)'\), and all solutions are \(x_p+c_1s_1+c_2s_2\)
(\(c_1,c_2\in\R\)), by Theorems~\ref{thm:r2-solution-structure}(a) and~\ref{thm:r4-four-bases}(c). Check:
\(Ax_p=-3a_1+2a_2=(-3+4,\ -3+6,\ -6+10)'=(1,3,4)'\).
\item For the two right sides
\[
B_1=\begin{pmatrix}1&0\\1&1\\2&1\end{pmatrix},\qquad B_2=\begin{pmatrix}1&0\\0&1\\0&0\end{pmatrix}
\]
and \(y=(-1,-1,1)'\), \(y'B_1=(-1-1+2,\ 0-1+1)=(0,0)\) and \(y'B_2=(-1,-1)\ne(0,0)\). Every vector of
\(\mathcal N(A')\) is a multiple of \(y\), so \(AX=B_1\) is consistent and \(AX=B_2\) is not. Indeed the columns of
\(B_1\) are \(a_1\) and \(a_2-2a_1\), so
\[
X=\begin{pmatrix}1&-2\\0&1\\0&0\\0&0\end{pmatrix}
\]
solves \(AX=B_1\) (Theorem~\ref{r1-thm-views}(a),(b)).
\end{enumerate}
\end{example}

\subsection*{The normal form of a matrix}

\begin{theorem}[\citeat{bapat}{Result~2.15}, \citeat{abadir-magnus}{Exercise~6.17}, \citeat{harville}{Theorem~4.4.9}]\label{thm:r4-normal-form}
Let \(A\in\R^{m\times n}\) have rank \(r\), and let \(N_r\in\R^{m\times n}\) be the matrix whose first \(r\) columns
are the unit vectors \(e_1,\dots,e_r\in\R^m\) and whose other \(n-r\) columns are \(0\) (\(N_0=0\)):
\[
N_r=\begin{pmatrix}I_r&0\\0&0\end{pmatrix}
\]
There are invertible matrices \(P\in\R^{m\times m}\) and \(Q\in\R^{n\times n}\) such that \(PAQ=N_r\). This \(N_r\) is
the \textbf{normal form} of \(A\).
\end{theorem}
\begin{proof}
If \(r=0\), then \(A=0\) (Theorem~\ref{thm:r4-four-bases}); take \(P=I_m\) and \(Q=I_n\). Let \(r\geqslant1\).
\begin{itemize}
\item By Theorem~\ref{thm:r4-full-rank-factorization}(a), \(A=B_1C_1\) with \(B_1\in\R^{m\times r}\) of full column
rank and \(C_1\in\R^{r\times n}\) of full row rank. So the columns \(b_1,\dots,b_r\) of \(B_1\) are independent, and so
are the rows \(c_1',\dots,c_r'\) of \(C_1\) (\(c_j\in\R^n\); Theorem~\ref{thm:r4-rank-rules}(b),(c)).
\item By Theorem~\ref{thm:r4-counting}(b), \(b_1,\dots,b_r\) extend to a basis \(b_1,\dots,b_m\) of \(\R^m\), and
\(c_1',\dots,c_r'\) extend to a basis \(c_1',\dots,c_n'\) of \(\R^{1\times n}\) (the paragraph after
Theorem~\ref{thm:r4-counting}; bases of \(\R^m\) and of \(\R^{1\times n}\) have \(m\) and \(n\) vectors, and nothing is
added if \(r=m\), or \(r=n\)).
\item Let \(B=[b_1\ \cdots\ b_m]\in\R^{m\times m}\), and let \(C\in\R^{n\times n}\) have the rows
\(c_1',\dots,c_n'\). The columns of \(B\) are independent, so \(\rank B=m\) and \(B\) is invertible
(Theorem~\ref{thm:r4-rank-rules}(b),(d)); the rows of \(C\) are independent, so \(\rank C=n\) and \(C\) is invertible
(Theorem~\ref{thm:r4-rank-rules}(c),(d)).
\item By Theorem~\ref{r1-thm-views}(a),(b), the columns of \(BN_r\) are \(Be_j=b_j\) (\(j\leqslant r\)) and \(0\)
(\(j>r\)). By the outer product expansion (Theorem~\ref{r1-thm-views}(d)), applied to \((BN_r)C\) and to \(B_1C_1\),
\[
BN_rC=\sum_{j=1}^rb_jc_j'=B_1C_1=A
\]
\item Hence \(P=B^{-1}\) and \(Q=C^{-1}\) are invertible (Theorem~\ref{r2-thm-inverse-props}(a)), and
\(PAQ=B^{-1}(BN_rC)C^{-1}=N_r\).\qedhere
\end{itemize}
\end{proof}

\begin{remark}[cf.\ \citeat{abadir-magnus}{Exercises~6.15, 6.16}]\label{rem:r4-normal-form}
The matrices \(P\) and \(Q\) of Theorem~\ref{thm:r4-normal-form} can be read off from the computation of
Theorem~\ref{thm:r4-four-bases}. Let \([A\mid I_m]\) be reduced to \([R\mid P]\), with pivot columns
\(k_1<\dots<k_r\), free columns \(f_1<\dots<f_{n-r}\) and special solutions \(s_1,\dots,s_{n-r}\), and let
\[
Q=[e_{k_1}\ \cdots\ e_{k_r}\ s_1\ \cdots\ s_{n-r}]\in\R^{n\times n}
\]
Then \(P\) and \(Q\) are invertible and \(PAQ=N_r\).
\begin{itemize}
\item \(P\) is invertible and \(PA=R\) (Theorem~\ref{thm:r4-four-bases}). The columns of \(RQ\) are
\(Re_{k_i}=\col_{k_i}(R)=e_i\) (\(i\leqslant r\); Definition~\ref{r2-def-rref}) and \(Rs_j=0\), as
\(s_j\in\mathcal N(A)=\mathcal N(R)\) (Theorem~\ref{r2-thm-rowops}(b)). So \(PAQ=RQ=N_r\)
(Theorem~\ref{r1-thm-views}(a)).
\item Let \(Qx=0\), that is, \(\sum_{i=1}^rx_ie_{k_i}+\sum_{j=1}^{n-r}x_{r+j}s_j=0\). Entry \(f_l\) of the left side is
\(x_{r+l}\), since \(e_{k_i}\) has entry \(0\) there and \(s_j\) has entry \(\delta_{jl}\); so
\(x_{r+1}=\dots=x_n=0\), and then \(\sum_ix_ie_{k_i}=0\) gives \(x_1=\dots=x_r=0\). By
Theorem~\ref{r2-thm-invertible}, \(Q\) is invertible.
\end{itemize}
\end{remark}

\begin{example}\label{exm:r4-normal-form}
For \(A\) of Example~\ref{exm:r4-four-bases}, the pivot columns are \(1,2\) and the special solutions are
\(s_1=(-3,1,1,0)'\), \(s_2=(1,-2,0,1)'\). Remark~\ref{rem:r4-normal-form} gives
\[
P=\begin{pmatrix}3&-2&0\\-1&1&0\\-1&-1&1\end{pmatrix},\qquad
Q=[e_1\ e_2\ s_1\ s_2]=\begin{pmatrix}1&0&-3&1\\0&1&1&-2\\0&0&1&0\\0&0&0&1\end{pmatrix}
\]
and
\[
PAQ=RQ=\begin{pmatrix}1&0&3&-1\\0&1&-1&2\\0&0&0&0\end{pmatrix}\begin{pmatrix}1&0&-3&1\\0&1&1&-2\\0&0&1&0\\0&0&0&1\end{pmatrix}
=\begin{pmatrix}1&0&0&0\\0&1&0&0\\0&0&0&0\end{pmatrix}=N_2
\]
These are also the matrices of the proof of Theorem~\ref{thm:r4-normal-form} for the factorization \(A=BC\) of
Example~\ref{exm:r4-factorization}, completed by \(e_3\) and by the rows \(e_3',e_4'\):
\begin{itemize}
\item \(\tilde B=[a_1\ a_2\ e_3]\) has \(P\tilde B=[Pa_1\ Pa_2\ Pe_3]=[e_1\ e_2\ e_3]=I_3\), since \(Pa_1\), \(Pa_2\) are
columns \(1\), \(2\) of \(PA=R\) and \(Pe_3=\col_3(P)=e_3\); so \(\tilde B^{-1}=P\)
(Theorem~\ref{r2-thm-invertible});
\item \(\tilde C\in\R^{4\times4}\), with rows \((1,0,3,-1)\), \((0,1,-1,2)\), \(e_3'\), \(e_4'\), has \(\tilde CQ=I_4\): the
first two rows of \(\tilde C\) form the matrix \(C\) with \(Ce_1=e_1\), \(Ce_2=e_2\), \(Cs_1=Cs_2=0\), and
\(e_3'Q\), \(e_4'Q\) are rows \(3\), \(4\) of \(Q\), that is, \(e_3'\) and \(e_4'\). So \(\tilde C^{-1}=Q\).
\end{itemize}
\end{example}

\needspace{12\baselineskip}
\subsection*{One-sided inverses}

\begin{definition}[\citeat{harville}{\S8.1}]\label{def:r4-one-sided}
Let \(A\in\R^{m\times n}\). A \textbf{right inverse} of \(A\) is a matrix \(K\in\R^{n\times m}\) with \(AK=I_m\), and a
\textbf{left inverse} of \(A\) is a matrix \(L\in\R^{n\times m}\) with \(LA=I_n\) (as for the factors of a full-rank
factorization in Theorem~\ref{thm:r4-full-rank-factorization}(c)).
\end{definition}

\begin{theorem}[\citeat{harville}{Lemma~8.1.1, Corollary~8.1.2, proof of Lemma~8.1.3}]\label{thm:r4-one-sided}
Let \(A\in\R^{m\times n}\).
\begin{enumerate}
\item \(A\) has a right inverse if and only if \(A\) has full row rank, \(\rank A=m\).
\item \(A\) has a left inverse if and only if \(A\) has full column rank, \(\rank A=n\).
\item \(A\) has both a right and a left inverse if and only if \(m=n\) and \(A\) is invertible; then every right inverse
and every left inverse of \(A\) equals \(A^{-1}\).
\end{enumerate}
\end{theorem}
\begin{proof}
(a)
\begin{itemize}
\item If \(\rank A=m\), the system \(AX=I_m\) is consistent (Corollary~\ref{cor:r4-matrix-system}, last sentence): it has a
solution \(K\in\R^{n\times m}\), a right inverse of \(A\).
\item Conversely, if \(AK=I_m\), then \(m=\rank I_m=\rank(AK)\leqslant\rank A\leqslant m\)
(Theorems~\ref{thm:r4-rank-rules}(d),~\ref{thm:r4-rank-product}(a) and~\ref{thm:r4-rank-rules}(a)); so \(\rank A=m\).
\end{itemize}
(b) \(LA=I_n\) if and only if \(A'L'=I_n\) (Theorem~\ref{r1-thm-transpose}), so \(L\) is a left inverse of \(A\) if and only
if \(L'\) is a right inverse of \(A'\in\R^{n\times m}\). By (a), \(A'\) has a right inverse if and only if \(\rank A'=n\), and
\(\rank A'=\rank A\) (Theorem~\ref{thm:r4-rank-rules}(a)).

(c)
\begin{itemize}
\item Let \(AK=I_m\) and \(LA=I_n\). By (a) and (b), \(\rank A=m\) and \(\rank A=n\), so \(m=n\), and \(A\) is invertible
(Theorem~\ref{thm:r4-rank-rules}(d)).
\item Then \(L=LI_m=L(AK)=(LA)K=I_nK=K\) (Theorem~\ref{r1-thm-rules}), so \(K\) is a two-sided inverse of \(A\), that is,
\(K=L=A^{-1}\) (Definition~\ref{r2-def-inverse}). This holds for every right inverse \(K\) and every left inverse \(L\).
\item Conversely, if \(m=n\) and \(A\) is invertible, \(A^{-1}\) is a right and a left inverse of \(A\).\qedhere
\end{itemize}
\end{proof}
For instance, if \(A\) has full column rank, then \(A'A\) is invertible
(Theorems~\ref{thm:r4-rank-ata}(b) and~\ref{thm:r4-rank-rules}(d)), and \((A'A)^{-1}A'\) is a left inverse of \(A\). A
one-sided inverse of a matrix that is not square is never unique: if \(AK_0=I_m\) and \(m\ne n\), then \(\rank A=m\) by (a),
so \(m<n\) (Theorem~\ref{thm:r4-rank-rules}(a)) and \(\dim\mathcal N(A)=n-m\geqslant1\) (Theorem~\ref{thm:r4-rank-nullity}), and for a nonzero \(s\in\mathcal N(A)\) every
\(K=K_0+st'\) (\(t\in\R^m\)) satisfies \(AK=AK_0+(As)t'=I_m\); for left inverses apply this to \(A'\) (\(L\) is a left
inverse of \(A\) if and only if \(L'\) is a right inverse of \(A'\), Theorem~\ref{r1-thm-transpose}). The next example
finds all right inverses of a \(3\times4\) matrix.

\begin{example}\label{exm:r4-one-sided}
Find all right inverses of
\[
A=\begin{pmatrix}1&0&1&2\\0&1&1&1\\1&1&3&3\end{pmatrix}
\]
and show that \(A\) has no left inverse.

\textit{Solution.}
\begin{itemize}
\item Row-reduce \([A\mid I_3]\):
\[
\begin{gathered}
\left(\begin{array}{cccc|ccc}1&0&1&2&1&0&0\\0&1&1&1&0&1&0\\1&1&3&3&0&0&1\end{array}\right)
\xrightarrow{R_3\leftarrow R_3-R_1}
\left(\begin{array}{cccc|ccc}1&0&1&2&1&0&0\\0&1&1&1&0&1&0\\0&1&2&1&-1&0&1\end{array}\right)
\\
\xrightarrow{R_3\leftarrow R_3-R_2}
\left(\begin{array}{cccc|ccc}1&0&1&2&1&0&0\\0&1&1&1&0&1&0\\0&0&1&0&-1&-1&1\end{array}\right)
\xrightarrow[R_2\leftarrow R_2-R_3]{R_1\leftarrow R_1-R_3}
\left(\begin{array}{cccc|ccc}1&0&0&2&2&1&-1\\0&1&0&1&1&2&-1\\0&0&1&0&-1&-1&1\end{array}\right)
\end{gathered}
\]
\item The RREF of \(A\), the left part, has three pivots, so \(\rank A=3=m\): \(A\) has right inverses
(Theorem~\ref{thm:r4-one-sided}(a)). As \(\rank A=3<4=n\), \(A\) has no left inverse (Theorem~\ref{thm:r4-one-sided}(b)).
\item Column \(j\) of a right inverse solves \(Ax=e_j\), and the same row operations take \([A\mid e_j]\) to the left part
and column \(j\) of the right part of the last matrix (Theorems~\ref{r2-thm-elementary}(c) and~\ref{r1-thm-views}(a)).
Let \(K_0\in\R^{4\times3}\) have this right part as its first three rows and a zero fourth row:
\[
K_0=\begin{pmatrix}2&1&-1\\1&2&-1\\-1&-1&1\\0&0&0\end{pmatrix},\qquad s=\begin{pmatrix}-2\\-1\\0\\1\end{pmatrix}
\]
The unknown \(x_4\) is free; with \(x_4=t_j\), Proposition~\ref{r2-prop-solutions}(b) gives \(x_1=k_{1j}-2t_j\),
\(x_2=k_{2j}-t_j\), \(x_3=k_{3j}\), where \(k_{ij}\) are the entries of \(K_0\). So the solutions are
\(x=\col_j(K_0)+t_js\) (\(t_j\in\R\)).
\item Hence the right inverses of \(A\) are the matrices \(K=K_0+st'\) with \(t=(t_1,t_2,t_3)'\in\R^3\), as column \(j\) of
\(st'\) is \(t_js\) (Theorem~\ref{r1-thm-views}(a),(d)); there are infinitely many. Check: \(As=(-2+2,\ -1+1,\ -2-1+3)'=0\),
and by rows, \(\row_3(AK_0)=1\,(2,1,-1)+1\,(1,2,-1)+3\,(-1,-1,1)=(0,0,1)\); likewise \(\row_1(AK_0)=(1,0,0)\) and
\(\row_2(AK_0)=(0,1,0)\), so \(AK_0=I_3\).
\end{itemize}
\end{example}

\needspace{14\baselineskip}
\subsection*{Exercises}

\begin{exercise}\label{ex:r4-extend-drill}
Let \(u_1=(1,-2,1,-1)'\) and \(u_2=(0,1,-1,1)'\) in \(\R^4\), and let \(e_1,\dots,e_4\) be the
unit vectors.
\begin{enumerate}
\item Row-reduce \([u_1\ u_2\ e_1\ e_2\ e_3\ e_4]\in\R^{4\times6}\), and find a basis of \(\R^4\) that
consists of \(u_1\), \(u_2\) and two of the unit vectors.
\item Do \(u_1,u_2,e_1,e_2\) form a basis of \(\R^4\)?
\end{enumerate}
\end{exercise}
\begin{solution}
(a) Call the matrix \(A=[u_1\ u_2\ e_1\ e_2\ e_3\ e_4]\); its columns span \(\R^4\), since they
include \(e_1,\dots,e_4\).
\[
\begin{gathered}
\left(\begin{array}{cc|cccc}1&0&1&0&0&0\\-2&1&0&1&0&0\\1&-1&0&0&1&0\\-1&1&0&0&0&1\end{array}\right)
\xrightarrow[R_4\leftarrow R_4+R_1]{R_2\leftarrow R_2+2R_1,\ R_3\leftarrow R_3-R_1}
\left(\begin{array}{cc|cccc}1&0&1&0&0&0\\0&1&2&1&0&0\\0&-1&-1&0&1&0\\0&1&1&0&0&1\end{array}\right)
\\
\xrightarrow[R_4\leftarrow R_4-R_2]{R_3\leftarrow R_3+R_2}
\left(\begin{array}{cc|cccc}1&0&1&0&0&0\\0&1&2&1&0&0\\0&0&1&1&1&0\\0&0&-1&-1&0&1\end{array}\right)
\xrightarrow{R_4\leftarrow R_4+R_3}
\left(\begin{array}{cc|cccc}1&0&1&0&0&0\\0&1&2&1&0&0\\0&0&1&1&1&0\\0&0&0&0&1&1\end{array}\right)
\end{gathered}
\]
The last matrix is an echelon form, with leading \(1\)s in columns \(1,2,3,5\); these are the pivot
columns of \(A\) (Remark~\ref{rem:r4-echelon}). By Theorem~\ref{thm:r4-pivot-columns}(b),
\(u_1,u_2,e_1,e_3\) form a basis of the span of the columns, which is \(\R^4\).

(b) No. Backward elimination finishes the RREF:
\[
\xrightarrow{R_3\leftarrow R_3-R_4}
\left(\begin{array}{cc|cccc}1&0&1&0&0&0\\0&1&2&1&0&0\\0&0&1&1&0&-1\\0&0&0&0&1&1\end{array}\right)
\xrightarrow[R_2\leftarrow R_2-2R_3]{R_1\leftarrow R_1-R_3}
\left(\begin{array}{cc|cccc}1&0&0&-1&0&1\\0&1&0&-1&0&2\\0&0&1&1&0&-1\\0&0&0&0&1&1\end{array}\right)
\]
Column \(4\) of the RREF gives \(e_2=-u_1-u_2+e_1\) (Theorem~\ref{thm:r4-pivot-columns}(b)). So
\(u_1,u_2,e_1,e_2\) are linearly dependent and do not form a basis. Check:
\(-(1,-2,1,-1)'-(0,1,-1,1)'+(1,0,0,0)'=(0,1,0,0)'=e_2\).
\end{solution}

\begin{exercise}[\citeat{abadir-magnus}{Exercise~3.23}]\label{ex:r4-which-bases}
Let \(a_1,\dots,a_9\) be the columns of
\[
A=\begin{pmatrix}1&1&1&1&1&1&2&2&5\\0&1&1&2&2&3&-1&3&3\\1&1&2&3&5&5&1&0&4\end{pmatrix}
\]
Which of the following form a basis of \(\R^3\)? (a) \(a_1,a_2\); (b) \(a_1,a_4,a_6,a_8\);
(c) \(a_2,a_4,a_7\); (d) \(a_3,a_5,a_9\).
\end{exercise}
\begin{solution}
(a) No: \(\dim\R^3=3\), so two vectors cannot span \(\R^3\) (Theorem~\ref{thm:r4-counting}(c)).

(b) No: \(\dim\R^3=3\), so four vectors of \(\R^3\) are dependent (Theorem~\ref{thm:r4-counting}(c)).

(c) Yes. Three vectors of \(\R^3\) form a basis if and only if they are independent
(Theorem~\ref{thm:r4-counting}(d)). Row-reduce \([a_2\ a_4\ a_7]\):
\[
\begin{gathered}
\begin{pmatrix}1&1&2\\1&2&-1\\1&3&1\end{pmatrix}
\xrightarrow[R_3\leftarrow R_3-R_1]{R_2\leftarrow R_2-R_1}
\begin{pmatrix}1&1&2\\0&1&-3\\0&2&-1\end{pmatrix}
\xrightarrow{R_3\leftarrow R_3-2R_2}
\begin{pmatrix}1&1&2\\0&1&-3\\0&0&5\end{pmatrix}
\\
\xrightarrow{R_3\leftarrow\frac15R_3}
\begin{pmatrix}1&1&2\\0&1&-3\\0&0&1\end{pmatrix}
\xrightarrow[R_2\leftarrow R_2+3R_3]{R_1\leftarrow R_1-2R_3}
\begin{pmatrix}1&1&0\\0&1&0\\0&0&1\end{pmatrix}
\xrightarrow{R_1\leftarrow R_1-R_2}
\begin{pmatrix}1&0&0\\0&1&0\\0&0&1\end{pmatrix}
\end{gathered}
\]
All three columns are pivot columns, so \(a_2,a_4,a_7\) are independent
(Theorem~\ref{thm:r4-pivot-test}).

(d) No: the three vectors are dependent, since
\[
7a_3-2a_5-a_9=7\begin{pmatrix}1\\1\\2\end{pmatrix}-2\begin{pmatrix}1\\2\\5\end{pmatrix}-\begin{pmatrix}5\\3\\4\end{pmatrix}
=\begin{pmatrix}7-2-5\\7-4-3\\14-10-4\end{pmatrix}=0
\]
\end{solution}

\begin{exercise}[\citeat{abadir-magnus}{Exercise~3.31(a), (b)}]\label{ex:r4-column-basis}
Let \(W=\mathcal C(A)\subseteq\R^4\) for
\[
A=\begin{pmatrix}1&2&3\\-2&3&8\\5&1&-3\\-3&-4&-5\end{pmatrix}
\]
(a) Find a basis of \(W\). (b) What is \(\dim W\)?
\end{exercise}
\begin{solution}
Let \(a_1,a_2,a_3\) be the columns of \(A\); they span \(W\).

(a) \(a_1,a_2\) are independent: otherwise, as \(a_1\ne0\), \(a_2=ca_1\) for some \(c\)
(Proposition~\ref{r3-prop-dependence}(c)); the first entries give \(c=2\), and then the second entries give \(3=2(-2)\)
\(\contradiction\). But
\[
a_1-2a_2+a_3=\begin{pmatrix}1-4+3\\-2-6+8\\5-2-3\\-3+8-5\end{pmatrix}=0
\]
so \(a_3=2a_2-a_1\), and every combination \(x_1a_1+x_2a_2+x_3a_3=(x_1-x_3)a_1+(x_2+2x_3)a_2\) is a
combination of \(a_1,a_2\). Thus \(a_1,a_2\) span \(W\) and form a basis of \(W\).

(b) \(\dim W=2\).
\end{solution}

\begin{exercise}\label{ex:r4-four-subspaces-drill}
Let
\[
A=\begin{pmatrix}1&2&0&1&-1\\1&2&1&0&1\\2&4&1&1&0\\1&2&2&-1&3\end{pmatrix}
\]
\begin{enumerate}
\item Find \(\rank A\) and bases of \(\mathcal C(A)\), \(\mathcal R(A)\), \(\mathcal N(A)\) and
\(\mathcal N(A')\), and write the non-pivot columns of \(A\) as combinations of the pivot columns.
\item For which \(b\in\R^4\) is \(Ax=b\) consistent? Give all solutions for \(b=(1,2,3,3)'\).
\end{enumerate}
\end{exercise}
\begin{solution}
(a) Row-reduce \([A\mid I_4]\) (Theorem~\ref{thm:r4-four-bases}):
\[
\begin{gathered}
\left(\begin{array}{ccccc|cccc}1&2&0&1&-1&1&0&0&0\\1&2&1&0&1&0&1&0&0\\2&4&1&1&0&0&0&1&0\\1&2&2&-1&3&0&0&0&1\end{array}\right)
\xrightarrow[R_4\leftarrow R_4-R_1]{R_2\leftarrow R_2-R_1,\ R_3\leftarrow R_3-2R_1}
\left(\begin{array}{ccccc|cccc}1&2&0&1&-1&1&0&0&0\\0&0&1&-1&2&-1&1&0&0\\0&0&1&-1&2&-2&0&1&0\\0&0&2&-2&4&-1&0&0&1\end{array}\right)
\\
\xrightarrow[R_4\leftarrow R_4-2R_2]{R_3\leftarrow R_3-R_2}
\left(\begin{array}{ccccc|cccc}1&2&0&1&-1&1&0&0&0\\0&0&1&-1&2&-1&1&0&0\\0&0&0&0&0&-1&-1&1&0\\0&0&0&0&0&1&-2&0&1\end{array}\right)
=[R\mid P]
\end{gathered}
\]
Row \(1\) has entry \(0\) in column \(3\), so no backward arrow is needed. Thus \(\rank A=2\), with
pivot columns \(1,3\) and free unknowns \(x_2,x_4,x_5\).
\begin{itemize}
\item \(\mathcal C(A)\): \(a_1=(1,1,2,1)'\) and \(a_3=(0,1,1,2)'\). Columns \(2\), \(4\), \(5\) of \(R\)
give \(a_2=2a_1\), \(a_4=a_1-a_3\) and \(a_5=-a_1+2a_3\) (Theorem~\ref{thm:r4-pivot-columns}(b)).
Check: \((1,1,2,1)'-(0,1,1,2)'=(1,0,1,-1)'=a_4\) and
\(-(1,1,2,1)'+2(0,1,1,2)'=(-1,1,0,3)'=a_5\).
\item \(\mathcal R(A)\): \(\rho_1'=(1,2,0,1,-1)\) and \(\rho_2'=(0,0,1,-1,2)\).
\item \(\mathcal N(A)\): \(Rx=0\) reads \(x_1+2x_2+x_4-x_5=0\), \(x_3-x_4+2x_5=0\). The special
solutions
\[
s_1=(-2,1,0,0,0)',\qquad s_2=(-1,0,1,1,0)',\qquad s_3=(1,0,-2,0,1)'
\]
form a basis; \(\dim\mathcal N(A)=3=5-2\).
\item \(\mathcal N(A')\): rows \(3\), \(4\) of \(P\) give the basis \(y_1=(-1,-1,1,0)'\),
\(y_2=(1,-2,0,1)'\); \(\dim\mathcal N(A')=2=4-2\). Indeed \(\row_3(A)=\row_1(A)+\row_2(A)\) and
\(\row_4(A)=-\row_1(A)+2\row_2(A)\).
\end{itemize}
(b) The row operations of (a) multiply on the left by the invertible matrix \(P\) of
\([R\mid P]\) (Theorem~\ref{r2-thm-elementary}(c)), so they take \([A\mid b]\) to
\(P[A\mid b]=[R\mid Pb]\) (Theorem~\ref{r1-thm-views}(a)), where
\[
Pb=(b_1,\ b_2-b_1,\ b_3-b_1-b_2,\ b_4+b_1-2b_2)'
\]
By Proposition~\ref{r2-prop-solutions}(a), \(Ax=b\) is consistent if and only if
\(b_3=b_1+b_2\) and \(b_4=2b_2-b_1\), that is, \(y_1'b=y_2'b=0\). For \(b=(1,2,3,3)'\) both hold, and
\(Pb=(1,1,0,0)'\). By Proposition~\ref{r2-prop-solutions}(b), \(x_1=1-2x_2-x_4+x_5\) and
\(x_3=1+x_4-2x_5\), so the solutions are
\[
x=(1,0,1,0,0)'+x_2s_1+x_4s_2+x_5s_3\qquad(x_2,x_4,x_5\in\R)
\]
Check: \(A(1,0,1,0,0)'=a_1+a_3=(1,2,3,3)'\).
\end{solution}

\begin{exercise}\label{r4-dr-four-subspaces}
Let
\[
A=\begin{pmatrix}1&1&1&0&1\\2&3&1&1&5\\1&2&0&2&5\\0&1&-1&1&3\end{pmatrix}
\]
\begin{enumerate}
\item Row-reduce \([A\mid I_4]\). Find \(\rank A\), bases of \(\mathcal C(A)\) and \(\mathcal R(A)\), and the non-pivot
columns of \(A\) as combinations of the pivot columns.
\item Find bases of \(\mathcal N(A)\) and \(\mathcal N(A')\).
\item Find an equation describing \(\mathcal C(A)\). Which of \(b=(1,3,0,1)'\) and \(u=(1,1,1,1)'\) lie in
\(\mathcal C(A)\)? For each one that does, find an \(x\) with \(Ax\) equal to it.
\end{enumerate}
\end{exercise}
\begin{solution}
(a)
\[
\begin{gathered}
\left(\begin{array}{ccccc|cccc}1&1&1&0&1&1&0&0&0\\2&3&1&1&5&0&1&0&0\\1&2&0&2&5&0&0&1&0\\0&1&-1&1&3&0&0&0&1\end{array}\right)
\xrightarrow[R_3\leftarrow R_3-R_1]{R_2\leftarrow R_2-2R_1}
\left(\begin{array}{ccccc|cccc}1&1&1&0&1&1&0&0&0\\0&1&-1&1&3&-2&1&0&0\\0&1&-1&2&4&-1&0&1&0\\0&1&-1&1&3&0&0&0&1\end{array}\right)
\\
\xrightarrow[R_4\leftarrow R_4-R_2]{R_3\leftarrow R_3-R_2}
\left(\begin{array}{ccccc|cccc}1&1&1&0&1&1&0&0&0\\0&1&-1&1&3&-2&1&0&0\\0&0&0&1&1&1&-1&1&0\\0&0&0&0&0&2&-1&0&1\end{array}\right)
\xrightarrow{R_2\leftarrow R_2-R_3}
\left(\begin{array}{ccccc|cccc}1&1&1&0&1&1&0&0&0\\0&1&-1&0&2&-3&2&-1&0\\0&0&0&1&1&1&-1&1&0\\0&0&0&0&0&2&-1&0&1\end{array}\right)
\\
\xrightarrow{R_1\leftarrow R_1-R_2}
\left(\begin{array}{ccccc|cccc}1&0&2&0&-1&4&-2&1&0\\0&1&-1&0&2&-3&2&-1&0\\0&0&0&1&1&1&-1&1&0\\0&0&0&0&0&2&-1&0&1\end{array}\right)=[R\mid P]
\end{gathered}
\]
So \(\rank A=3\), with pivot columns \(1,2,4\) and free unknowns \(x_3,x_5\).
\begin{itemize}
\item \(\mathcal C(A)\): \(a_1=(1,2,1,0)'\), \(a_2=(1,3,2,1)'\) and \(a_4=(0,1,2,1)'\) form a basis. Columns \(3\) and
\(5\) of \(R\) give \(a_3=2a_1-a_2\) and \(a_5=-a_1+2a_2+a_4\) (Theorem~\ref{thm:r4-pivot-columns}(b)). Check:
\(2(1,2,1,0)'-(1,3,2,1)'=(1,1,0,-1)'\) and \(-(1,2,1,0)'+2(1,3,2,1)'+(0,1,2,1)'=(1,5,5,3)'\).
\item \(\mathcal R(A)\): \(\rho_1'=(1,0,2,0,-1)\), \(\rho_2'=(0,1,-1,0,2)\) and \(\rho_3'=(0,0,0,1,1)\) form a basis.
\end{itemize}
(b)
\begin{itemize}
\item \(\mathcal N(A)\): \(Rx=0\) reads \(x_1+2x_3-x_5=0\), \(x_2-x_3+2x_5=0\), \(x_4+x_5=0\). The special solutions
\(s_1=(-2,1,1,0,0)'\) (\(x_3=1\), \(x_5=0\)) and \(s_2=(1,-2,0,-1,1)'\) (\(x_3=0\), \(x_5=1\)) form a basis;
\(\dim\mathcal N(A)=2=5-3\).
\item \(\mathcal N(A')\): row \(4\) of \(P\) gives the basis \(y=(2,-1,0,1)'\) (Theorem~\ref{thm:r4-four-bases}(d));
\(\dim\mathcal N(A')=1=4-3\). Indeed \(y'A=2\row_1(A)-\row_2(A)+\row_4(A)=0\).
\end{itemize}
(c)
\begin{itemize}
\item By the paragraph after Theorem~\ref{thm:r4-compatibility}, \(b\in\mathcal C(A)\) if and only if
\(\row_4(P)\,b=0\), that is, \(2b_1-b_2+b_4=0\).
\item For \(b=(1,3,0,1)'\): \(2-3+1=0\), so \(b\in\mathcal C(A)\). Then \(Pb=(4-6,\ -3+6,\ 1-3,\ 2-3+1)'=(-2,3,-2,0)'\),
and \(Rx=Pb\) reads \(x_1+2x_3-x_5=-2\), \(x_2-x_3+2x_5=3\), \(x_4+x_5=-2\). The free unknowns \(x_3=x_5=0\) give the
solution \(x=(-2,3,0,-2,0)'\). Check: \(Ax=-2a_1+3a_2-2a_4=(-2+3,\ -4+9-2,\ -2+6-4,\ 3-2)'=(1,3,0,1)'\).
\item For \(u=(1,1,1,1)'\): \(2-1+1=2\ne0\), so \(u\notin\mathcal C(A)\).
\end{itemize}
\end{solution}

\begin{exercise}\label{ex:r4-rank-drill}
Find \(\rank A\), bases of \(\mathcal C(A)\), \(\mathcal R(A)\) and \(\mathcal N(A)\), and
\(\dim\mathcal N(A')\) for
\[
A=\begin{pmatrix}1&2&1&1\\0&1&1&-1\\1&1&0&2\\0&-2&-2&3\end{pmatrix}
\]
and write the non-pivot column of \(A\) as a combination of the pivot columns.
\end{exercise}
\begin{solution}
Row-reduce \(A\):
\[
\begin{gathered}
A\xrightarrow{R_3\leftarrow R_3-R_1}
\begin{pmatrix}1&2&1&1\\0&1&1&-1\\0&-1&-1&1\\0&-2&-2&3\end{pmatrix}
\xrightarrow[R_4\leftarrow R_4+2R_2]{R_3\leftarrow R_3+R_2}
\begin{pmatrix}1&2&1&1\\0&1&1&-1\\0&0&0&0\\0&0&0&1\end{pmatrix}
\xrightarrow{R_3\leftrightarrow R_4}
\begin{pmatrix}1&2&1&1\\0&1&1&-1\\0&0&0&1\\0&0&0&0\end{pmatrix}
\\
\xrightarrow[R_1\leftarrow R_1-R_3]{R_2\leftarrow R_2+R_3}
\begin{pmatrix}1&2&1&0\\0&1&1&0\\0&0&0&1\\0&0&0&0\end{pmatrix}
\xrightarrow{R_1\leftarrow R_1-2R_2}
\begin{pmatrix}1&0&-1&0\\0&1&1&0\\0&0&0&1\\0&0&0&0\end{pmatrix}=R
\end{gathered}
\]
After the second arrow column \(3\) has no nonzero entry in rows \(3,4\), and the next pivot column,
\(4\), has its nonzero entry in row \(4\); so rows \(3\) and \(4\) are interchanged. The echelon form
after the third arrow has three nonzero rows: \(\rank A=3\), with pivot columns \(1,2,4\).
\begin{itemize}
\item \(\mathcal C(A)\): \(a_1=(1,0,1,0)'\), \(a_2=(2,1,1,-2)'\) and \(a_4=(1,-1,2,3)'\) form a basis
(Theorem~\ref{thm:r4-four-bases}(a)). Column \(3\) of \(R\) gives \(a_3=-a_1+a_2\); check:
\(-(1,0,1,0)'+(2,1,1,-2)'=(1,1,0,-2)'=a_3\).
\item \(\mathcal R(A)\): \(\rho_1'=(1,0,-1,0)\), \(\rho_2'=(0,1,1,0)\) and \(\rho_3'=(0,0,0,1)\) form a basis
(Theorem~\ref{thm:r4-four-bases}(b)).
\item \(\mathcal N(A)\): \(Rx=0\) reads \(x_1-x_3=0\), \(x_2+x_3=0\), \(x_4=0\); the special solution
\(s=(1,-1,1,0)'\) (\(x_3=1\)) is a basis (Theorem~\ref{thm:r4-four-bases}(c)), and \(\dim\mathcal N(A)=1=4-3\).
\item \(\dim\mathcal N(A')=4-\rank A=1\) (Theorem~\ref{thm:r4-rank-nullity}).
\end{itemize}
\end{solution}

\begin{exercise}\label{r4-dr-basis-check}
Let \(W_0=\{x\in\R^4:x_1+x_2+x_3+x_4=0\}\), the centered vectors of \(\R^4\), and let
\[
v_1=\begin{pmatrix}1\\0\\0\\-1\end{pmatrix},\quad v_2=\begin{pmatrix}0\\1\\0\\-1\end{pmatrix},\quad
v_3=\begin{pmatrix}1\\1\\-1\\-1\end{pmatrix},\quad w=\begin{pmatrix}2\\-1\\3\\-4\end{pmatrix}
\]
\begin{enumerate}
\item Find \(\dim W_0\).
\item Show that \(v_1,v_2,v_3\) form a basis of \(W_0\), and find the coordinates of \(w\) in this basis.
\item Does \((1,2,3,4)'\) lie in \(\Span\{v_1,v_2,v_3\}\)?
\end{enumerate}
\end{exercise}
\begin{solution}
(a) \(W_0=\mathcal N(\one_4')\), where the \(1\times4\) matrix \(\one_4'=(1,1,1,1)\) is in RREF with one nonzero row. So
\(\rank\one_4'=1\) and \(\dim W_0=4-1=3\) (Theorem~\ref{thm:r4-rank-nullity}).

(b) The entries of each of \(v_1,v_2,v_3,w\) add up to \(0\), so these vectors lie in \(W_0\). Row-reduce
\([v_1\ v_2\ v_3\mid w]\):
\[
\begin{gathered}
\left(\begin{array}{ccc|c}1&0&1&2\\0&1&1&-1\\0&0&-1&3\\-1&-1&-1&-4\end{array}\right)
\xrightarrow{R_4\leftarrow R_4+R_1}
\left(\begin{array}{ccc|c}1&0&1&2\\0&1&1&-1\\0&0&-1&3\\0&-1&0&-2\end{array}\right)
\xrightarrow{R_4\leftarrow R_4+R_2}
\left(\begin{array}{ccc|c}1&0&1&2\\0&1&1&-1\\0&0&-1&3\\0&0&1&-3\end{array}\right)
\\
\xrightarrow{R_3\leftarrow-R_3}
\left(\begin{array}{ccc|c}1&0&1&2\\0&1&1&-1\\0&0&1&-3\\0&0&1&-3\end{array}\right)
\xrightarrow{R_4\leftarrow R_4-R_3}
\left(\begin{array}{ccc|c}1&0&1&2\\0&1&1&-1\\0&0&1&-3\\0&0&0&0\end{array}\right)
\xrightarrow[R_2\leftarrow R_2-R_3]{R_1\leftarrow R_1-R_3}
\left(\begin{array}{ccc|c}1&0&0&5\\0&1&0&2\\0&0&1&-3\\0&0&0&0\end{array}\right)
\end{gathered}
\]
\begin{itemize}
\item Row operations act on each column separately (Theorems~\ref{r2-thm-elementary}(c) and~\ref{r1-thm-views}(a)), so
the left part shows the RREF of \([v_1\ v_2\ v_3]\): every column is a pivot column, and \(v_1,v_2,v_3\) are
independent (Theorem~\ref{thm:r4-pivot-test}). Three independent vectors of the \(3\)-dimensional space \(W_0\) form a
basis of it (Theorem~\ref{thm:r4-counting}(d)).
\item The last column is not a pivot column, and the coordinates are read off: \(c=(5,2,-3)'\)
(Proposition~\ref{r2-prop-solutions}), that is, \(w=5v_1+2v_2-3v_3\). Check:
\(5(1,0,0,-1)'+2(0,1,0,-1)'-3(1,1,-1,-1)'=(5-3,\ 2-3,\ 3,\ -5-2+3)'=(2,-1,3,-4)'\).
\end{itemize}
(c) No. By (b), \(\Span\{v_1,v_2,v_3\}=W_0\), and the entries of \((1,2,3,4)'\) add up to \(10\ne0\).
\end{solution}

\begin{exercise}[\citeat{lang}{Ch.~V, \S3, Exercise~4(b), (d)}, \citeat{shakarchi}{}]\label{ex:r4-solution-spaces}
Find the dimension of the space of solutions \((x,y,z)'\in\R^3\) of each system, and a basis of
it: (a) \(x-y+z=0\); (b) \(x+y+z=0\), \(x-y=0\), \(y+z=0\).
\end{exercise}
\begin{solution}
The space of solutions is \(\mathcal N(A)\) for the coefficient matrix \(A\), of dimension
\(3-\rank A\) (Theorem~\ref{thm:r4-rank-nullity}).

(a) \(A=(1,-1,1)\) is a nonzero RREF with one row, so \(\rank A=1\) and the dimension is \(2\). The
vectors \((1,1,0)'\) and \((0,1,1)'\) are solutions (\(1-1+0=0\), \(0-1+1=0\)) and are
independent (\(c_1(1,1,0)'+c_2(0,1,1)'=(c_1,c_1+c_2,c_2)'=0\) forces \(c_1=c_2=0\)); so
they form a basis (Theorem~\ref{thm:r4-counting}(d)).

(b) Row-reduce the coefficient matrix:
\[
\begin{pmatrix}1&1&1\\1&-1&0\\0&1&1\end{pmatrix}
\xrightarrow{R_2\leftarrow R_2-R_1}
\begin{pmatrix}1&1&1\\0&-2&-1\\0&1&1\end{pmatrix}
\xrightarrow{R_2\leftrightarrow R_3}
\begin{pmatrix}1&1&1\\0&1&1\\0&-2&-1\end{pmatrix}
\xrightarrow{R_3\leftarrow R_3+2R_2}
\begin{pmatrix}1&1&1\\0&1&1\\0&0&1\end{pmatrix}
\]
This echelon form has three nonzero rows, so \(\rank A=3\) (Definition~\ref{def:r4-rank}): the
space of solutions has dimension \(3-3=0\), so it is \(\{0\}\) (Theorem~\ref{thm:r4-dimension}), with the empty basis.
\end{solution}

\begin{exercise}\label{r4-dr-rank-parameter}
For \(t\in\R\) let
\[
A=\begin{pmatrix}1&2&1&1\\2&5&1&4\\1&3&1&3\\1&1&2&t\end{pmatrix}
\]
\begin{enumerate}
\item Find \(\rank A\) for every \(t\).
\item For the value of \(t\) with \(\rank A<4\), find a basis of \(\mathcal C(A)\), write the non-pivot column as a
combination of the pivot columns, and find a basis of \(\mathcal N(A)\).
\end{enumerate}
\end{exercise}
\begin{solution}
(a) Forward elimination, keeping \(t\) as a letter:
\[
\begin{gathered}
\begin{pmatrix}1&2&1&1\\2&5&1&4\\1&3&1&3\\1&1&2&t\end{pmatrix}
\xrightarrow[R_4\leftarrow R_4-R_1]{R_2\leftarrow R_2-2R_1,\ R_3\leftarrow R_3-R_1}
\begin{pmatrix}1&2&1&1\\0&1&-1&2\\0&1&0&2\\0&-1&1&t-1\end{pmatrix}
\\
\xrightarrow[R_4\leftarrow R_4+R_2]{R_3\leftarrow R_3-R_2}
\begin{pmatrix}1&2&1&1\\0&1&-1&2\\0&0&1&0\\0&0&0&t+1\end{pmatrix}
\end{gathered}
\]
The last row is \((0,0,0,t+1)\). If \(t\ne-1\), scaling row \(4\) by \(1/(t+1)\) gives an echelon form with four
nonzero rows, so \(\rank A=4\) (Remark~\ref{rem:r4-echelon} and Definition~\ref{def:r4-rank}). If \(t=-1\), the last row
is zero, and the last matrix is an echelon form with three nonzero rows: \(\rank A=3\).

(b) Let \(t=-1\). Backward elimination finishes the RREF:
\[
\begin{gathered}
\begin{pmatrix}1&2&1&1\\0&1&-1&2\\0&0&1&0\\0&0&0&0\end{pmatrix}
\xrightarrow[R_1\leftarrow R_1-R_3]{R_2\leftarrow R_2+R_3}
\begin{pmatrix}1&2&0&1\\0&1&0&2\\0&0&1&0\\0&0&0&0\end{pmatrix}
\xrightarrow{R_1\leftarrow R_1-2R_2}
\begin{pmatrix}1&0&0&-3\\0&1&0&2\\0&0&1&0\\0&0&0&0\end{pmatrix}=R
\end{gathered}
\]
\begin{itemize}
\item The pivot columns are \(1,2,3\): \(a_1=(1,2,1,1)'\), \(a_2=(2,5,3,1)'\) and \(a_3=(1,1,1,2)'\) form a basis of
\(\mathcal C(A)\). Column \(4\) of \(R\) gives \(a_4=-3a_1+2a_2\) (Theorem~\ref{thm:r4-pivot-columns}(b)); check:
\(-3(1,2,1,1)'+2(2,5,3,1)'=(1,4,3,-1)'\), the last column of \(A\) for \(t=-1\).
\item \(Rx=0\) reads \(x_1-3x_4=0\), \(x_2+2x_4=0\), \(x_3=0\). The special solution \(s=(3,-2,0,1)'\) (\(x_4=1\)) is a
basis of \(\mathcal N(A)\), and \(\dim\mathcal N(A)=1=4-3\).
\end{itemize}
\end{solution}

\begin{exercise}\label{ex:r4-rank-parameter}
For \(t\in\R\) let
\[
A_t=\begin{pmatrix}t&1&1&1\\1&t&1&1\\1&1&t&1\\1&1&1&t\end{pmatrix}
\]
Find \(\rank A_t\) for every \(t\), and a basis of \(\mathcal N(A_t)\) whenever \(A_t\) is not
invertible.
\end{exercise}
\begin{solution}
Interchange the rows so that the first row is \((1,1,1,t)\), and eliminate, keeping \(t\) as a letter:
\[
\begin{gathered}
A_t\xrightarrow[R_2\leftrightarrow R_3]{R_1\leftrightarrow R_4}
\begin{pmatrix}1&1&1&t\\1&1&t&1\\1&t&1&1\\t&1&1&1\end{pmatrix}
\xrightarrow[R_4\leftarrow R_4-tR_1]{R_2\leftarrow R_2-R_1,\ R_3\leftarrow R_3-R_1}
\begin{pmatrix}1&1&1&t\\0&0&t-1&1-t\\0&t-1&0&1-t\\0&1-t&1-t&1-t^2\end{pmatrix}
\\
\xrightarrow{R_4\leftarrow R_4+R_2}
\begin{pmatrix}1&1&1&t\\0&0&t-1&1-t\\0&t-1&0&1-t\\0&1-t&0&-(t-1)(t+2)\end{pmatrix}
\xrightarrow{R_4\leftarrow R_4+R_3}
\begin{pmatrix}1&1&1&t\\0&0&t-1&1-t\\0&t-1&0&1-t\\0&0&0&-(t-1)(t+3)\end{pmatrix}
\\
\xrightarrow{R_2\leftrightarrow R_3}
\begin{pmatrix}1&1&1&t\\0&t-1&0&1-t\\0&0&t-1&1-t\\0&0&0&-(t-1)(t+3)\end{pmatrix}=E_t
\end{gathered}
\]
since \((1-t^2)+(1-t)=-(t-1)(t+2)\) and \(-(t-1)(t+2)+(1-t)=-(t-1)(t+3)\). The pivot columns can be read off
before the leading entries are scaled to \(1\) (Remark~\ref{rem:r4-echelon}).
\begin{itemize}
\item \(t\ne1\), \(t\ne-3\): all four diagonal entries of \(E_t\) are nonzero, so \(\rank A_t=4\)
(Definition~\ref{def:r4-rank}) and
\(A_t\) is invertible (Theorem~\ref{thm:r4-rank-rules}(d)).
\item \(t=1\): rows \(2\), \(3\), \(4\) of \(E_1\) are zero, so \(\rank A_1=1\). Here \(\mathcal N(A_1)\)
is given by \(x_1+x_2+x_3+x_4=0\), with special solutions \((-1,1,0,0)'\), \((-1,0,1,0)'\),
\((-1,0,0,1)'\) as a basis (Theorem~\ref{thm:r4-four-bases}(c)); \(\dim\mathcal N(A_1)=3\).
\item \(t=-3\): \(E_{-3}\) has nonzero rows \((1,1,1,-3)\), \((0,-4,0,4)\), \((0,0,-4,4)\) and a zero
row, so \(\rank A_{-3}=3\) and \(\dim\mathcal N(A_{-3})=1\) (Theorem~\ref{thm:r4-rank-nullity}).
Every row of \(A_t\) has sum \(t+3\), so \(A_{-3}\one_4=0\); the nonzero vector \(\one_4=(1,1,1,1)'\)
is a basis of \(\mathcal N(A_{-3})\) (Proposition~\ref{r3-prop-dependence}(a),
Theorem~\ref{thm:r4-counting}(d)).
\end{itemize}
\end{solution}

\begin{exercise}\label{r4-dr-consistency}
Let
\[
A=\begin{pmatrix}1&2&1\\2&5&3\\1&3&2\\3&7&4\end{pmatrix},\qquad b=\begin{pmatrix}1\\t\\2\\4\end{pmatrix}\qquad(t\in\R)
\]
\begin{enumerate}
\item For which \(t\) is \(Ax=b\) consistent? Compare \(\rank A\) and \(\rank[A\mid b]\).
\item For this \(t\), find all solutions. Is the solution unique?
\end{enumerate}
\end{exercise}
\begin{solution}
(a) Row-reduce \([A\mid b]\), keeping \(t\) as a letter:
\[
\begin{gathered}
\left(\begin{array}{ccc|c}1&2&1&1\\2&5&3&t\\1&3&2&2\\3&7&4&4\end{array}\right)
\xrightarrow[R_4\leftarrow R_4-3R_1]{R_2\leftarrow R_2-2R_1,\ R_3\leftarrow R_3-R_1}
\left(\begin{array}{ccc|c}1&2&1&1\\0&1&1&t-2\\0&1&1&1\\0&1&1&1\end{array}\right)
\\
\xrightarrow[R_4\leftarrow R_4-R_2]{R_3\leftarrow R_3-R_2}
\left(\begin{array}{ccc|c}1&2&1&1\\0&1&1&t-2\\0&0&0&3-t\\0&0&0&3-t\end{array}\right)
\xrightarrow{R_4\leftarrow R_4-R_3}
\left(\begin{array}{ccc|c}1&2&1&1\\0&1&1&t-2\\0&0&0&3-t\\0&0&0&0\end{array}\right)
\end{gathered}
\]
\begin{itemize}
\item The left part of the last matrix is an echelon form with two nonzero rows, so \(\rank A=2\)
(Remark~\ref{rem:r4-echelon}).
\item If \(t\ne3\), row \(3\) reads \((0\ 0\ 0\mid3-t)\) with \(3-t\ne0\): the equation \(0=3-t\) has no solution
(Theorem~\ref{thm:r2-consistency}). Scaling row \(3\) by \(1/(3-t)\) gives an echelon form of \([A\mid b]\) with three
nonzero rows, so \(\rank[A\mid b]=3>\rank A\).
\item If \(t=3\), rows \(3\) and \(4\) are zero: \(\rank[A\mid b]=2=\rank A\), and the system is consistent
(Corollary~\ref{cor:r4-ax-b}(a)).
\end{itemize}
(b) Let \(t=3\). Backward elimination:
\[
\begin{gathered}
\left(\begin{array}{ccc|c}1&2&1&1\\0&1&1&1\\0&0&0&0\\0&0&0&0\end{array}\right)
\xrightarrow{R_1\leftarrow R_1-2R_2}
\left(\begin{array}{ccc|c}1&0&-1&-1\\0&1&1&1\\0&0&0&0\\0&0&0&0\end{array}\right)
\end{gathered}
\]
With the free unknown \(x_3=c\): \(x_1=-1+c\), \(x_2=1-c\), that is,
\[
x=\begin{pmatrix}-1\\1\\0\end{pmatrix}+c\begin{pmatrix}1\\-1\\1\end{pmatrix}\qquad(c\in\R)
\]
The solution is not unique, since \(\rank A=2<3\) (Corollary~\ref{cor:r4-ax-b}(b)); \((1,-1,1)'\), the special
solution of the left part, spans \(\mathcal N(A)\) (Theorem~\ref{thm:r4-four-bases}(c)). Check: \(A(-1,1,0)'=-a_1+a_2=(1,3,2,4)'\) and \(A(1,-1,1)'=a_1-a_2+a_3=(1-2+1,\ 2-5+3,\ 1-3+2,\ 3-7+4)'=0\).
\end{solution}

\needspace{24\baselineskip}
\begin{exercise}\label{r4-dr-product-sum}
Let
\[
A=\begin{pmatrix}1&1&0&0\\0&1&1&1\\1&2&1&1\\1&1&1&1\end{pmatrix},\qquad
B=\begin{pmatrix}0&1&0&1\\0&-1&0&-1\\1&1&1&-1\\0&1&-1&1\end{pmatrix}
\]
\begin{enumerate}
\item Find \(\rank A\) and \(\rank B\).
\item Compute \(AB\), \(BA\) and their ranks, and compare them with the bounds of
Theorems~\ref{thm:r4-rank-product}(a) and~\ref{thm:r4-sylvester}.
\item Compute \(A+B\) and its rank, and compare it with the bound of Corollary~\ref{cor:r4-rank-sum}.
\end{enumerate}
\end{exercise}
\begin{solution}
(a) Row-reduce \(A\) and \(B\). For \(A\), Step 2 interchanges rows \(3\) and \(4\) after the second arrow; for
\(B\), Step 2 interchanges rows \(1\) and \(3\) at the first arrow and rows \(3\) and \(4\) at the fourth:
\[
\begin{gathered}
\begin{pmatrix}1&1&0&0\\0&1&1&1\\1&2&1&1\\1&1&1&1\end{pmatrix}
\xrightarrow[R_4\leftarrow R_4-R_1]{R_3\leftarrow R_3-R_1}
\begin{pmatrix}1&1&0&0\\0&1&1&1\\0&1&1&1\\0&0&1&1\end{pmatrix}
\xrightarrow{R_3\leftarrow R_3-R_2}
\begin{pmatrix}1&1&0&0\\0&1&1&1\\0&0&0&0\\0&0&1&1\end{pmatrix}
\\
\xrightarrow{R_3\leftrightarrow R_4}
\begin{pmatrix}1&1&0&0\\0&1&1&1\\0&0&1&1\\0&0&0&0\end{pmatrix}
\xrightarrow{R_2\leftarrow R_2-R_3}
\begin{pmatrix}1&1&0&0\\0&1&0&0\\0&0&1&1\\0&0&0&0\end{pmatrix}
\xrightarrow{R_1\leftarrow R_1-R_2}
\begin{pmatrix}1&0&0&0\\0&1&0&0\\0&0&1&1\\0&0&0&0\end{pmatrix}
\end{gathered}
\]
\[
\begin{gathered}
\begin{pmatrix}0&1&0&1\\0&-1&0&-1\\1&1&1&-1\\0&1&-1&1\end{pmatrix}
\xrightarrow{R_1\leftrightarrow R_3}
\begin{pmatrix}1&1&1&-1\\0&-1&0&-1\\0&1&0&1\\0&1&-1&1\end{pmatrix}
\xrightarrow{R_2\leftarrow-R_2}
\begin{pmatrix}1&1&1&-1\\0&1&0&1\\0&1&0&1\\0&1&-1&1\end{pmatrix}
\\
\xrightarrow[R_4\leftarrow R_4-R_2]{R_3\leftarrow R_3-R_2}
\begin{pmatrix}1&1&1&-1\\0&1&0&1\\0&0&0&0\\0&0&-1&0\end{pmatrix}
\xrightarrow{R_3\leftrightarrow R_4}
\begin{pmatrix}1&1&1&-1\\0&1&0&1\\0&0&-1&0\\0&0&0&0\end{pmatrix}
\xrightarrow{R_3\leftarrow-R_3}
\begin{pmatrix}1&1&1&-1\\0&1&0&1\\0&0&1&0\\0&0&0&0\end{pmatrix}
\\
\xrightarrow{R_1\leftarrow R_1-R_3}
\begin{pmatrix}1&1&0&-1\\0&1&0&1\\0&0&1&0\\0&0&0&0\end{pmatrix}
\xrightarrow{R_1\leftarrow R_1-R_2}
\begin{pmatrix}1&0&0&-2\\0&1&0&1\\0&0&1&0\\0&0&0&0\end{pmatrix}
\end{gathered}
\]
So \(\rank A=\rank B=3\).

(b) By Theorem~\ref{r1-thm-views}(c), \(\row_i(AB)=\row_i(A)B\) and \(\row_i(BA)=\row_i(B)A\):
\(\row_1(AB)=\row_1(B)+\row_2(B)=(0,0,0,0)\), \(\row_2(AB)=\row_2(B)+\row_3(B)+\row_4(B)=(1,1,0,-1)\),
\(\row_3(AB)=\row_1(B)+2\row_2(B)+\row_3(B)+\row_4(B)=(1,1,0,-1)\),
\(\row_4(AB)=\row_1(B)+\row_2(B)+\row_3(B)+\row_4(B)=(1,2,0,0)\); and \(\row_1(BA)=\row_2(A)+\row_4(A)\),
\(\row_2(BA)=-\row_2(A)-\row_4(A)\), \(\row_3(BA)=\row_1(A)+\row_2(A)+\row_3(A)-\row_4(A)\),
\(\row_4(BA)=\row_2(A)-\row_3(A)+\row_4(A)\):
\[
AB=\begin{pmatrix}0&0&0&0\\1&1&0&-1\\1&1&0&-1\\1&2&0&0\end{pmatrix},\qquad
BA=\begin{pmatrix}1&2&2&2\\-1&-2&-2&-2\\1&3&1&1\\0&0&1&1\end{pmatrix}
\]
The rows of \(AB\) span \(\Span\{(1,1,0,-1),(1,2,0,0)\}\), and these two rows are independent (neither is a multiple
of the other, Proposition~\ref{r3-prop-dependence}(c)); so \(\dim\mathcal R(AB)=2=\rank(AB)\)
(Definition~\ref{def:r4-rank}). For \(BA\):
\[
\begin{gathered}
\begin{pmatrix}1&2&2&2\\-1&-2&-2&-2\\1&3&1&1\\0&0&1&1\end{pmatrix}
\xrightarrow[R_3\leftarrow R_3-R_1]{R_2\leftarrow R_2+R_1}
\begin{pmatrix}1&2&2&2\\0&0&0&0\\0&1&-1&-1\\0&0&1&1\end{pmatrix}
\xrightarrow{R_2\leftrightarrow R_3}
\begin{pmatrix}1&2&2&2\\0&1&-1&-1\\0&0&0&0\\0&0&1&1\end{pmatrix}
\\
\xrightarrow{R_3\leftrightarrow R_4}
\begin{pmatrix}1&2&2&2\\0&1&-1&-1\\0&0&1&1\\0&0&0&0\end{pmatrix}
\xrightarrow[R_2\leftarrow R_2+R_3]{R_1\leftarrow R_1-2R_3}
\begin{pmatrix}1&2&0&0\\0&1&0&0\\0&0&1&1\\0&0&0&0\end{pmatrix}
\xrightarrow{R_1\leftarrow R_1-2R_2}
\begin{pmatrix}1&0&0&0\\0&1&0&0\\0&0&1&1\\0&0&0&0\end{pmatrix}
\end{gathered}
\]
so \(\rank(BA)=3\).
\begin{itemize}
\item Theorem~\ref{thm:r4-rank-product}(a): \(\rank(AB)=2<\min\{3,3\}\) and \(\rank(BA)=3=\min\{3,3\}\).
\item Theorem~\ref{thm:r4-sylvester}: both ranks are at least \(3+3-4=2\); \(AB\) attains the bound. The products
\(AB\) and \(BA\) have different ranks.
\end{itemize}
(c) \(A+B\) is the first matrix below; row-reduce it (Step 2 interchanges rows \(2\) and \(3\)):
\[
\begin{gathered}
\begin{pmatrix}1&2&0&1\\0&0&1&0\\2&3&2&0\\1&2&0&2\end{pmatrix}
\xrightarrow[R_4\leftarrow R_4-R_1]{R_3\leftarrow R_3-2R_1}
\begin{pmatrix}1&2&0&1\\0&0&1&0\\0&-1&2&-2\\0&0&0&1\end{pmatrix}
\xrightarrow{R_2\leftrightarrow R_3}
\begin{pmatrix}1&2&0&1\\0&-1&2&-2\\0&0&1&0\\0&0&0&1\end{pmatrix}
\\
\xrightarrow{R_2\leftarrow-R_2}
\begin{pmatrix}1&2&0&1\\0&1&-2&2\\0&0&1&0\\0&0&0&1\end{pmatrix}
\xrightarrow[R_2\leftarrow R_2-2R_4]{R_1\leftarrow R_1-R_4}
\begin{pmatrix}1&2&0&0\\0&1&-2&0\\0&0&1&0\\0&0&0&1\end{pmatrix}
\xrightarrow{R_2\leftarrow R_2+2R_3}
\begin{pmatrix}1&2&0&0\\0&1&0&0\\0&0&1&0\\0&0&0&1\end{pmatrix}
\\
\xrightarrow{R_1\leftarrow R_1-2R_2}
\begin{pmatrix}1&0&0&0\\0&1&0&0\\0&0&1&0\\0&0&0&1\end{pmatrix}
\end{gathered}
\]
So \(\rank(A+B)=4<3+3\) (Corollary~\ref{cor:r4-rank-sum}).
\end{solution}

\begin{exercise}\label{ex:r4-shift-ranks}
Let \(N\in\R^{4\times4}\) be the matrix of Exercise~\ref{r1-ex-commuting}: its entries are \(n_{12}=n_{23}=n_{34}=1\)
and \(0\) elsewhere, so that \(Ne_1=0\) and \(Ne_j=e_{j-1}\) (\(j=2,3,4\)).
\begin{enumerate}
\item Compute \(N^2\), \(N^3\), \(N^4\), and \(\rank N^k\) for \(k=0,\dots,4\).
\item Show that Sylvester's inequality for the product \(N^k\cdot N\) is an equality for
\(k=1,2,3\).
\item Compute \(N'N\), \(NN'\) and \(\rank(N+N')\).
\end{enumerate}
\end{exercise}
\begin{solution}
(a) \(\col_j(N^k)=N^ke_j\) (Theorem~\ref{r1-thm-views}(b)), and applying \(N\) repeatedly gives
\(N^ke_j=e_{j-k}\) for \(j>k\) and \(N^ke_j=0\) for \(j\leqslant k\). So
\[
N^2=\begin{pmatrix}0&0&1&0\\0&0&0&1\\0&0&0&0\\0&0&0&0\end{pmatrix},\qquad
N^3=\begin{pmatrix}0&0&0&1\\0&0&0&0\\0&0&0&0\\0&0&0&0\end{pmatrix},\qquad N^4=0
\]
The nonzero columns of \(N^k\) are \(e_1,\dots,e_{4-k}\). They are independent, as part of the standard basis
(Definition~\ref{r3-def-basis} and Proposition~\ref{r3-prop-dependence}(b)), and they span \(\mathcal C(N^k)\), since
the zero columns add nothing to the span (Theorem~\ref{r3-thm-span-inclusion}(c)); so \(\rank N^k=4-k\)
(\(k=0,\dots,4\)).

(b) Theorem~\ref{thm:r4-sylvester} gives
\(\rank(N^k\cdot N)\geqslant\rank N^k+\rank N-4=(4-k)+3-4=3-k\), and \(N^k\cdot N=N^{k+1}\) has rank
\(4-(k+1)=3-k\) by (a).

(c) \(N'\) has entries \(1\) in positions \((2,1)\), \((3,2)\), \((4,3)\) and \(0\) elsewhere, so \(N'e_j=e_{j+1}\)
(\(j=1,2,3\)) and \(N'e_4=0\). Hence \(N'Ne_1=0\), \(N'Ne_j=N'e_{j-1}=e_j\) (\(j\geqslant2\)), and
\(NN'e_j=Ne_{j+1}=e_j\) (\(j\leqslant3\)), \(NN'e_4=0\):
\[
N'N=\diag(0,1,1,1),\qquad NN'=\diag(1,1,1,0)
\]
both of rank \(3=\rank N\) (Theorem~\ref{thm:r4-rank-ata}(b)). For \(N+N'\):
\[
\begin{pmatrix}0&1&0&0\\1&0&1&0\\0&1&0&1\\0&0&1&0\end{pmatrix}
\xrightarrow{R_1\leftrightarrow R_2}
\begin{pmatrix}1&0&1&0\\0&1&0&0\\0&1&0&1\\0&0&1&0\end{pmatrix}
\xrightarrow{R_3\leftarrow R_3-R_2}
\begin{pmatrix}1&0&1&0\\0&1&0&0\\0&0&0&1\\0&0&1&0\end{pmatrix}
\xrightarrow{R_3\leftrightarrow R_4}
\begin{pmatrix}1&0&1&0\\0&1&0&0\\0&0&1&0\\0&0&0&1\end{pmatrix}
\]
an echelon form with four nonzero rows; so \(\rank(N+N')=4\), below the bound
\(\rank N+\rank N'=6\) of Corollary~\ref{cor:r4-rank-sum}.
\end{solution}

\begin{exercise}\label{r4-dr-factorization}
Four variables are measured on four units: \(x_1=(1,1,2,0)'\), \(x_2=2x_1\) (the first variable in other units),
\(x_3=(1,2,3,1)'\) and \(x_4=x_1-x_3\). Let \(A=[x_1\ x_2\ x_3\ x_4]\in\R^{4\times4}\).
\begin{enumerate}
\item Row-reduce \(A\) and find \(\rank A\).
\item Find the full-rank factorization \(A=BC\) of Theorem~\ref{thm:r4-full-rank-factorization}(a), and check it.
\item Write \(A\) as a sum of \(\rank A\) matrices of rank \(1\).
\end{enumerate}
\end{exercise}
\begin{solution}
(a)
\[
\begin{gathered}
\begin{pmatrix}1&2&1&0\\1&2&2&-1\\2&4&3&-1\\0&0&1&-1\end{pmatrix}
\xrightarrow[R_3\leftarrow R_3-2R_1]{R_2\leftarrow R_2-R_1}
\begin{pmatrix}1&2&1&0\\0&0&1&-1\\0&0&1&-1\\0&0&1&-1\end{pmatrix}
\\
\xrightarrow[R_4\leftarrow R_4-R_2]{R_3\leftarrow R_3-R_2}
\begin{pmatrix}1&2&1&0\\0&0&1&-1\\0&0&0&0\\0&0&0&0\end{pmatrix}
\xrightarrow{R_1\leftarrow R_1-R_2}
\begin{pmatrix}1&2&0&1\\0&0&1&-1\\0&0&0&0\\0&0&0&0\end{pmatrix}
\end{gathered}
\]
So \(\rank A=2\), with pivot columns \(1\) and \(3\).

(b) \(B\) consists of the pivot columns of \(A\), and \(C\) of the nonzero rows of the RREF:
\[
B=[x_1\ x_3]=\begin{pmatrix}1&1\\1&2\\2&3\\0&1\end{pmatrix},\qquad C=\begin{pmatrix}1&2&0&1\\0&0&1&-1\end{pmatrix}
\]
Check by columns (Theorem~\ref{r1-thm-views}(a),(b)): \(B\col_1(C)=x_1\), \(B\col_2(C)=2x_1=x_2\),
\(B\col_3(C)=x_3\), \(B\col_4(C)=x_1-x_3=x_4\). So \(BC=A\), with \(\rank B=\rank C=2\).

(c) By the outer product expansion (Theorem~\ref{r1-thm-views}(d)), \(A=x_1\row_1(C)+x_3\row_2(C)\):
\[
A=\begin{pmatrix}1&2&0&1\\1&2&0&1\\2&4&0&2\\0&0&0&0\end{pmatrix}+\begin{pmatrix}0&0&1&-1\\0&0&2&-2\\0&0&3&-3\\0&0&1&-1\end{pmatrix}
\]
Each term is nonzero and its rows are multiples of one row, so each has rank \(1\). Check: the sum has rows
\((1,2,1,0)\), \((1,2,2,-1)\), \((2,4,3,-1)\), \((0,0,1,-1)\), the rows of \(A\).
\end{solution}

\begin{exercise}\label{ex:r4-factorization-drill}
Let
\[
A=\begin{pmatrix}1&1&0&2\\-1&0&-1&-1\\1&-1&2&0\\0&1&-1&1\\1&0&1&1\end{pmatrix}
\]
\begin{enumerate}
\item Find the RREF \(R\) and \(\rank A\), and the full-rank factorization \(A=BC\) of
Theorem~\ref{thm:r4-full-rank-factorization}(a).
\item Compute \(B'B\), \(CC'\) and the matrices \(L\), \(K\) of
Theorem~\ref{thm:r4-full-rank-factorization}(c), and verify \(LB=I_2\) and \(CK=I_2\).
\item Find a basis of \(\mathcal N(A)\).
\end{enumerate}
\end{exercise}
\begin{solution}
(a) Row-reduce \(A\):
\[
A\xrightarrow[R_5\leftarrow R_5-R_1]{R_2\leftarrow R_2+R_1,\ R_3\leftarrow R_3-R_1}
\begin{pmatrix}1&1&0&2\\0&1&-1&1\\0&-2&2&-2\\0&1&-1&1\\0&-1&1&-1\end{pmatrix}
\xrightarrow[R_4\leftarrow R_4-R_2,\ R_5\leftarrow R_5+R_2]{R_1\leftarrow R_1-R_2,\ R_3\leftarrow R_3+2R_2}
\begin{pmatrix}1&0&1&1\\0&1&-1&1\\0&0&0&0\\0&0&0&0\\0&0&0&0\end{pmatrix}=R
\]
So \(\rank A=2\), with pivot columns \(1,2\), and
\[
B=\begin{pmatrix}1&1\\-1&0\\1&-1\\0&1\\1&0\end{pmatrix},\qquad
C=\begin{pmatrix}1&0&1&1\\0&1&-1&1\end{pmatrix}
\]
Columns \(3\), \(4\) of \(R\) say \(a_3=a_1-a_2\) and \(a_4=a_1+a_2\); check:
\(a_1-a_2=(0,-1,2,-1,1)'\) and \(a_1+a_2=(2,-1,0,1,1)'\).

(b) With \(a_1=(1,-1,1,0,1)'\) and \(a_2=(1,0,-1,1,0)'\), the entries of \(B'B\) are the products \(a_j'a_l\)
(Theorem~\ref{r1-thm-views}(d)):
\[
B'B=\begin{pmatrix}a_1'a_1&a_1'a_2\\a_2'a_1&a_2'a_2\end{pmatrix}=\begin{pmatrix}4&0\\0&3\end{pmatrix},\qquad
CC'=\begin{pmatrix}3&0\\0&3\end{pmatrix}
\]
(\(a_1'a_2=1+0-1+0+0=0\) and \(\row_1(C)\row_2(C)'=0+0-1+1=0\)). Hence, as \((B'B)^{-1}=\diag(\frac14,\frac13)\) and
\((CC')^{-1}=\frac13I_2\) (Definition~\ref{r2-def-inverse}),
\[
L=\begin{pmatrix}\frac14&0\\0&\frac13\end{pmatrix}B'
=\begin{pmatrix}\frac14&-\frac14&\frac14&0&\frac14\\[2pt]
\frac13&0&-\frac13&\frac13&0\end{pmatrix},\qquad
K=\frac13C'=\frac13\begin{pmatrix}1&0\\0&1\\1&-1\\1&1\end{pmatrix}
\]
Then \(LB=\diag(\frac14,\frac13)B'B=\diag(\frac14,\frac13)\diag(4,3)=I_2\) and
\(CK=\frac13CC'=I_2\).

(c) \(\mathcal N(A)=\mathcal N(C)\) (Theorem~\ref{thm:r4-full-rank-factorization}(b)); \(Cx=0\) reads
\(x_1=-x_3-x_4\), \(x_2=x_3-x_4\), and the special solutions \((-1,1,1,0)'\) and \((-1,-1,0,1)'\)
form a basis (Theorem~\ref{thm:r4-four-bases}(c)). They are the relations \(a_3=a_1-a_2\) and \(a_4=a_1+a_2\).
\end{solution}

\begin{exercise}\label{ex:r4-gram-drill}
For \(A=BC\) of Exercise~\ref{ex:r4-factorization-drill}:
\begin{enumerate}
\item compute \(A'A\) from \(A'A=C'(B'B)C\);
\item row-reduce \(A'A\), and check that \(\rank(A'A)=\rank A\) and \(\mathcal N(A'A)=\mathcal N(A)\);
\item find \(\rank(AA')\) without computing the \(5\times5\) matrix \(AA'\).
\end{enumerate}
\end{exercise}
\begin{solution}
(a) \(A'A=C'B'BC\) (Theorem~\ref{r1-thm-transpose}), with \(B'B=\diag(4,3)\)
(Exercise~\ref{ex:r4-factorization-drill}(b)). The rows of \(C\) are \(\gamma_1'\) and \(\gamma_2'\), with
\(\gamma_1=(1,0,1,1)'\) and \(\gamma_2=(0,1,-1,1)'\) in \(\R^4\). The columns of \(C'\diag(4,3)\) are \(4\gamma_1\) and
\(3\gamma_2\) (Proposition~\ref{r1-prop-diag}(a)), so the outer product expansion (Theorem~\ref{r1-thm-views}(d)) gives
\[
A'A=4\gamma_1\gamma_1'+3\gamma_2\gamma_2'
=\begin{pmatrix}4&0&4&4\\0&0&0&0\\4&0&4&4\\4&0&4&4\end{pmatrix}
+\begin{pmatrix}0&0&0&0\\0&3&-3&3\\0&-3&3&-3\\0&3&-3&3\end{pmatrix}
=\begin{pmatrix}4&0&4&4\\0&3&-3&3\\4&-3&7&1\\4&3&1&7\end{pmatrix}
\]
(b) Row-reduce \(A'A\):
\[
\begin{gathered}
\begin{pmatrix}4&0&4&4\\0&3&-3&3\\4&-3&7&1\\4&3&1&7\end{pmatrix}
\xrightarrow[R_2\leftarrow\frac13R_2]{R_1\leftarrow\frac14R_1}
\begin{pmatrix}1&0&1&1\\0&1&-1&1\\4&-3&7&1\\4&3&1&7\end{pmatrix}
\xrightarrow[R_4\leftarrow R_4-4R_1]{R_3\leftarrow R_3-4R_1}
\begin{pmatrix}1&0&1&1\\0&1&-1&1\\0&-3&3&-3\\0&3&-3&3\end{pmatrix}
\\
\xrightarrow[R_4\leftarrow R_4-3R_2]{R_3\leftarrow R_3+3R_2}
\begin{pmatrix}1&0&1&1\\0&1&-1&1\\0&0&0&0\\0&0&0&0\end{pmatrix}
\end{gathered}
\]
The RREF of \(A'A\) is the matrix \(R\) of Exercise~\ref{ex:r4-factorization-drill} without its last
row. So \(\rank(A'A)=2=\rank A\), and \(A'Ax=0\) and \(Ax=0\) both reduce to \(Cx=0\):
\(\mathcal N(A'A)=\mathcal N(A)\), as Theorem~\ref{thm:r4-rank-ata}(a),(b) asserts.

(c) \(\rank(AA')=\rank A=2\) (Theorem~\ref{thm:r4-rank-ata}(b)).
\end{solution}

\begin{exercise}\label{r4-dr-sum-intersection}
In \(\R^4\) let \(U=\Span\{u_1,u_2,u_3\}\) and \(W=\Span\{w_1,w_2\}\), where
\[
u_1=\begin{pmatrix}1\\1\\-1\\1\end{pmatrix},\quad u_2=\begin{pmatrix}0\\1\\0\\1\end{pmatrix},\quad
u_3=\begin{pmatrix}1\\3\\-1\\3\end{pmatrix},\qquad w_1=\begin{pmatrix}0\\0\\1\\-1\end{pmatrix},\quad
w_2=\begin{pmatrix}1\\2\\1\\0\end{pmatrix}
\]
\begin{enumerate}
\item Find \(\dim U\), \(\dim W\) and \(\dim(U+W)\) by one row reduction.
\item Find \(\dim(U\cap W)\) by Theorem~\ref{thm:r4-dim-sum}, and a basis of \(U\cap W\).
\end{enumerate}
\end{exercise}
\begin{solution}
(a) \(U+W=\Span\{u_1,u_2,u_3,w_1,w_2\}\) (Theorem~\ref{r3-thm-sum}(c)). Row-reduce \(M=[u_1\ u_2\ u_3\ w_1\ w_2]\):
\[
\begin{gathered}
\begin{pmatrix}1&0&1&0&1\\1&1&3&0&2\\-1&0&-1&1&1\\1&1&3&-1&0\end{pmatrix}
\xrightarrow[R_4\leftarrow R_4-R_1]{R_2\leftarrow R_2-R_1,\ R_3\leftarrow R_3+R_1}
\begin{pmatrix}1&0&1&0&1\\0&1&2&0&1\\0&0&0&1&2\\0&1&2&-1&-1\end{pmatrix}
\\
\xrightarrow{R_4\leftarrow R_4-R_2}
\begin{pmatrix}1&0&1&0&1\\0&1&2&0&1\\0&0&0&1&2\\0&0&0&-1&-2\end{pmatrix}
\xrightarrow{R_4\leftarrow R_4+R_3}
\begin{pmatrix}1&0&1&0&1\\0&1&2&0&1\\0&0&0&1&2\\0&0&0&0&0\end{pmatrix}
\end{gathered}
\]
The last matrix is the RREF of \(M\), with pivot columns \(1,2,4\).
\begin{itemize}
\item \(\dim U=2\): row operations act on each column separately (Theorems~\ref{r2-thm-elementary}(c)
and~\ref{r1-thm-views}(a)), so the first three columns of the last matrix form the RREF of \([u_1\ u_2\ u_3]\), with pivot
columns \(1,2\). So \(u_1,u_2\) form a basis of \(U\), and column \(3\) gives \(u_3=u_1+2u_2\)
(Theorem~\ref{thm:r4-pivot-columns}(b)).
\item \(\dim W=2\): \(w_1\) and \(w_2\) are not multiples of each other (the first entry of \(w_1\) is \(0\), that of
\(w_2\) is not, and \(w_1\ne0\)), so they are independent (Proposition~\ref{r3-prop-dependence}(c)).
\item \(\dim(U+W)=3\), with the basis \(u_1,u_2,w_1\) (Theorem~\ref{thm:r4-pivot-columns}(b)).
\end{itemize}
(b)
\begin{itemize}
\item \(\dim(U\cap W)=\dim U+\dim W-\dim(U+W)=2+2-3=1\) (Theorem~\ref{thm:r4-dim-sum}).
\item Column \(5\) gives \(w_2=u_1+u_2+2w_1\), so \(v=w_2-2w_1=u_1+u_2=(1,2,-1,2)'\) lies in \(W\) and in \(U\). This
nonzero vector is a basis of the one-dimensional space \(U\cap W\) (Theorem~\ref{thm:r4-counting}(d)). Check:
\(w_2-2w_1=(1,\ 2,\ 1-2,\ 0+2)'=(1,2,-1,2)'\) and \(u_1+u_2=(1,2,-1,2)'\).
\end{itemize}
\end{solution}

\begin{exercise}\label{r4-dr-normal-form}
Let
\[
A=\begin{pmatrix}1&2&1&0\\2&4&3&-1\\1&2&2&0\\0&0&-1&1\end{pmatrix}
\]
Find invertible matrices \(P\) and \(Q\) with \(PAQ=N_r\), where \(r=\rank A\) (Theorem~\ref{thm:r4-normal-form}), by
Remark~\ref{rem:r4-normal-form}, and check the product.
\end{exercise}
\begin{solution}
Row-reduce \([A\mid I_4]\); after the first arrow column \(2\) has no nonzero entry in rows \(2\)--\(4\), so the next
pivot column is \(3\):
\[
\begin{gathered}
\left(\begin{array}{cccc|cccc}1&2&1&0&1&0&0&0\\2&4&3&-1&0&1&0&0\\1&2&2&0&0&0&1&0\\0&0&-1&1&0&0&0&1\end{array}\right)
\xrightarrow[R_3\leftarrow R_3-R_1]{R_2\leftarrow R_2-2R_1}
\left(\begin{array}{cccc|cccc}1&2&1&0&1&0&0&0\\0&0&1&-1&-2&1&0&0\\0&0&1&0&-1&0&1&0\\0&0&-1&1&0&0&0&1\end{array}\right)
\\
\xrightarrow[R_4\leftarrow R_4+R_2]{R_3\leftarrow R_3-R_2}
\left(\begin{array}{cccc|cccc}1&2&1&0&1&0&0&0\\0&0&1&-1&-2&1&0&0\\0&0&0&1&1&-1&1&0\\0&0&0&0&-2&1&0&1\end{array}\right)
\xrightarrow{R_2\leftarrow R_2+R_3}
\left(\begin{array}{cccc|cccc}1&2&1&0&1&0&0&0\\0&0&1&0&-1&0&1&0\\0&0&0&1&1&-1&1&0\\0&0&0&0&-2&1&0&1\end{array}\right)
\\
\xrightarrow{R_1\leftarrow R_1-R_2}
\left(\begin{array}{cccc|cccc}1&2&0&0&2&0&-1&0\\0&0&1&0&-1&0&1&0\\0&0&0&1&1&-1&1&0\\0&0&0&0&-2&1&0&1\end{array}\right)
\end{gathered}
\]
So \(\rank A=3\), with pivot columns \(1,3,4\), the free column \(2\), and
\[
R=\begin{pmatrix}1&2&0&0\\0&0&1&0\\0&0&0&1\\0&0&0&0\end{pmatrix},\qquad
P=\begin{pmatrix}2&0&-1&0\\-1&0&1&0\\1&-1&1&0\\-2&1&0&1\end{pmatrix}
\]
\(Rx=0\) reads \(x_1+2x_2=0\), \(x_3=0\), \(x_4=0\), and the special solution is \(s=(-2,1,0,0)'\). By
Remark~\ref{rem:r4-normal-form}, \(P\) and
\[
Q=[e_1\ e_3\ e_4\ s]=\begin{pmatrix}1&0&0&-2\\0&0&0&1\\0&1&0&0\\0&0&1&0\end{pmatrix}
\]
are invertible and \(PAQ=N_3\). Check: the columns of \(AQ\) are \(a_1\), \(a_3\), \(a_4\) and \(As=-2a_1+a_2=0\), so
\[
PAQ=\begin{pmatrix}2&0&-1&0\\-1&0&1&0\\1&-1&1&0\\-2&1&0&1\end{pmatrix}
\begin{pmatrix}1&1&0&0\\2&3&-1&0\\1&2&0&0\\0&-1&1&0\end{pmatrix}
=\begin{pmatrix}1&0&0&0\\0&1&0&0\\0&0&1&0\\0&0&0&0\end{pmatrix}=N_3
\]
\end{solution}

\needspace{14\baselineskip}
\subsection*{Problems}

\begin{problem}[\citeat{harville-solutions}{Ch.~18, Exercise~26}]\label{r4-pr-rank-difference}
Show that, for all \(A,B\in\R^{m\times n}\), \(\rank(A+B)\geqslant\lvert\rank A-\rank B\rvert\).
\end{problem}
\begin{solution}
\begin{itemize}
\item \(\rank(-B)=\rank B\): \(-B=(-I_m)BI_n\) (Theorem~\ref{r1-thm-rules}), and \(-I_m\) is invertible, as
\((-I_m)(-I_m)=I_m\) (Definition~\ref{r2-def-inverse}); so Theorem~\ref{thm:r4-rank-product}(d) applies. Likewise
\(\rank(-A)=\rank A\).
\item By Corollary~\ref{cor:r4-rank-sum}, applied to \(A+B\) and \(-B\),
\[
\rank A=\rank\bigl((A+B)+(-B)\bigr)\leqslant\rank(A+B)+\rank(-B)=\rank(A+B)+\rank B
\]
so \(\rank(A+B)\geqslant\rank A-\rank B\).
\item Similarly, \(\rank B=\rank\bigl((A+B)+(-A)\bigr)\leqslant\rank(A+B)+\rank A\), so
\(\rank(A+B)\geqslant\rank B-\rank A=-(\rank A-\rank B)\).
\end{itemize}
Together, \(\rank(A+B)\geqslant\lvert\rank A-\rank B\rvert\).
\end{solution}

\begin{problem}[\citeat{abadir-magnus}{Exercise~4.18}]\label{ex:r4-rank-atab}
Show that \(\rank(A'AB)=\rank(AB)=\rank(ABB')\) for all \(A\in\R^{m\times n}\) and
\(B\in\R^{n\times p}\).
\end{problem}
\begin{solution}
Apply Theorem~\ref{thm:r4-rank-ata}(b) to \(AB\), with \((AB)'=B'A'\) (Theorem~\ref{r1-thm-transpose}),
and Theorem~\ref{thm:r4-rank-product}(a) twice:
\[
\begin{aligned}
\rank(AB)&=\rank\bigl((AB)'AB\bigr)=\rank\bigl(B'(A'AB)\bigr)\leqslant\rank(A'AB)
=\rank\bigl(A'(AB)\bigr)\leqslant\rank(AB)\\
\rank(AB)&=\rank\bigl(AB(AB)'\bigr)=\rank\bigl((ABB')A'\bigr)\leqslant\rank(ABB')
=\rank\bigl((AB)B'\bigr)\leqslant\rank(AB)
\end{aligned}
\]
So equality holds throughout.
\end{solution}

\begin{problem}[\citeat{bapat}{Ch.~11, Exercise~37}]\label{ex:r4-involutory-rank}
Let \(A\in\R^{n\times n}\) satisfy \(A^2=I\). Show that \(\rank(I-A)+\rank(I+A)=n\).
\end{problem}
\begin{solution}
\begin{itemize}
\item \(2I=(I-A)+(I+A)\), and \(2I=\diag(2,\dots,2)\) is invertible (Definition~\ref{r2-def-inverse}). By
Theorem~\ref{thm:r4-rank-rules}(d) and Corollary~\ref{cor:r4-rank-sum},
\(n=\rank(2I)\leqslant\rank(I-A)+\rank(I+A)\).
\item \((I-A)(I+A)=I+A-A-A^2=I-A^2=0\) (Theorem~\ref{r1-thm-rules}), so
\(\rank(I-A)+\rank(I+A)\leqslant n\) by Theorem~\ref{thm:r4-sylvester}.
\end{itemize}
Together, \(\rank(I-A)+\rank(I+A)=n\).
\end{solution}

\begin{problem}[\citeat{harville-solutions}{Ch.~4, Exercise~10}]\label{ex:r4-sum-rank-one}
Let \(A\in\R^{m\times n}\) have rank \(r\). Show that \(A\) is the sum of \(r\) matrices of
rank \(1\).
\end{problem}
\begin{solution}
If \(r=0\), then \(A=0\), the empty sum. Let \(r\geqslant1\), and let \(A=BC\) be the full-rank
factorization of Theorem~\ref{thm:r4-full-rank-factorization}(a), where \(B\in\R^{m\times r}\) has columns
\(b_1,\dots,b_r\) and \(C\in\R^{r\times n}\). By the outer product expansion
(Theorem~\ref{r1-thm-views}(d)),
\[
A=\sum_{j=1}^rA_j,\qquad A_j=b_j\row_j(C)\in\R^{m\times n}
\]
\begin{itemize}
\item The columns of \(B\) are independent, and so are the rows of \(C\)
(Theorem~\ref{thm:r4-rank-rules}(b),(c)); hence every \(b_j\) and every \(\row_j(C)\) is nonzero,
since a list containing \(0\) is dependent (Proposition~\ref{r3-prop-dependence}(a)).
\item If the \(l\)th entry \(c_{jl}\) of \(\row_j(C)\) is nonzero, then
\(\col_l(A_j)=c_{jl}b_j\ne0\); so \(A_j\ne0\), and \(\rank A_j\geqslant1\).
\item \(\rank A_j\leqslant\rank b_j\leqslant1\) (Theorems~\ref{thm:r4-rank-product}(a)
and~\ref{thm:r4-rank-rules}(a)), as \(b_j\) has one column.
\end{itemize}
So each \(A_j\) has rank \(1\).
\end{solution}

\begin{problem}[\citeat{bapat}{Ch.~2, Exercise~17}]\label{r4-pr-frobenius}
Let \(A,X,B\in\R^{n\times n}\). Show the following generalization of Sylvester's inequality
(Theorem~\ref{thm:r4-sylvester}) for square matrices:
\[
\rank(AXB)\geqslant\rank(AX)+\rank(XB)-\rank X
\]
\end{problem}
\begin{solution}
If \(X=0\), then \(AX\), \(XB\) and \(AXB\) are zero, and both sides are \(0\). Let \(r=\rank X\geqslant1\).
\begin{itemize}
\item Let \(X=UV\) be a full-rank factorization (Theorem~\ref{thm:r4-full-rank-factorization}(a)), with
\(U\in\R^{n\times r}\) of full column rank and \(V\in\R^{r\times n}\) of full row rank. Then
\(AXB=(AU)(VB)\), \(AX=(AU)V\) and \(XB=U(VB)\) (Theorem~\ref{r1-thm-rules}).
\item As \(V\) has full row rank, \(\rank(AX)=\rank\bigl((AU)V\bigr)=\rank(AU)\)
(Theorem~\ref{thm:r4-rank-product}(c)); as \(U\) has full column rank,
\(\rank(XB)=\rank\bigl(U(VB)\bigr)=\rank(VB)\) (Theorem~\ref{thm:r4-rank-product}(b)).
\item Sylvester's inequality (Theorem~\ref{thm:r4-sylvester}) for \(AU\in\R^{n\times r}\) and \(VB\in\R^{r\times n}\),
whose product has the inner size \(r\), gives
\[
\rank(AXB)=\rank\bigl((AU)(VB)\bigr)\geqslant\rank(AU)+\rank(VB)-r=\rank(AX)+\rank(XB)-\rank X
\]
\end{itemize}
\end{solution}

\begin{problem}[\citeat{bapat}{Ch.~2, Exercise~2}]\label{r4-ex-exchange}
Let \(x_1,\dots,x_p\) and \(y_1,\dots,y_q\) be linearly independent lists of vectors of \(\R^n\), where
\(p<q\leqslant n\). Show that there is an \(i\in\{1,\dots,q\}\) such that \(x_1,\dots,x_p,y_i\) are linearly
independent.
\end{problem}
\begin{solution}
Let \(A=[x_1\ \cdots\ x_p]\in\R^{n\times p}\) and \(B=[y_1\ \cdots\ y_q]\in\R^{n\times q}\), and suppose that
\(x_1,\dots,x_p,y_i\) are dependent for every \(i\).
\begin{itemize}
\item Then each \(y_i\) lies in \(\Span\{x_1,\dots,x_p\}=\mathcal C(A)\): otherwise \(x_1,\dots,x_p,y_i\) would be
independent, by Proposition~\ref{r3-prop-adjoin}(b), as \(x_1,\dots,x_p\) are independent.
\item By Proposition~\ref{prop:r4-inclusion}(a), \(B=AC\) for some \(C\in\R^{p\times q}\).
\item The columns of \(A\) and of \(B\) are independent, so \(\rank A=p\) and \(\rank B=q\)
(Theorem~\ref{thm:r4-rank-rules}(b)). By Theorem~\ref{thm:r4-rank-product}(a),
\[
q=\rank B=\rank(AC)\leqslant\rank A=p<q\ \contradiction
\]
\end{itemize}
\end{solution}

\begin{problem}[\citeat{harville-solutions}{Ch.~7, Exercise~1}]\label{r4-ex-cancellation}
\begin{enumerate}
\item Let \(A\in\R^{m\times n}\), \(C\in\R^{n\times q}\) and \(B\in\R^{q\times p}\). Show that if
\(\rank(AC)=\rank C\), then \(\mathcal R(ACB)=\mathcal R(CB)\) and \(\rank(ACB)=\rank(CB)\); and that if
\(\rank(CB)=\rank C\), then \(\mathcal C(ACB)=\mathcal C(AC)\) and \(\rank(ACB)=\rank(AC)\).
\item Let \(A,B\in\R^{m\times n}\). Show that (1) if \(C\in\R^{s\times q}\) and \(D\in\R^{q\times m}\) satisfy
\(\rank(CD)=\rank D\), then \(CDA=CDB\) implies \(DA=DB\); and (2) if \(C\in\R^{n\times q}\) and \(D\in\R^{q\times p}\)
satisfy \(\rank(CD)=\rank C\), then \(ACD=BCD\) implies \(AC=BC\).
\end{enumerate}
\end{problem}
\begin{solution}
(a)
\begin{itemize}
\item \(\mathcal R(ACB)=\mathcal R(A(CB))\subseteq\mathcal R(CB)\) and
\(\mathcal C(ACB)=\mathcal C((AC)B)\subseteq\mathcal C(AC)\) (Proposition~\ref{prop:r4-inclusion}(b),(a)).
\item Let \(\rank(AC)=\rank C\). Then \(\mathcal R(AC)\subseteq\mathcal R(C)\) (Proposition~\ref{prop:r4-inclusion}(b)),
and both have dimension \(\rank C\) (Definition~\ref{def:r4-rank}), so \(\mathcal R(AC)=\mathcal R(C)\)
(Theorem~\ref{thm:r4-counting}(a), for row spaces). Hence \(C=LAC\) for some \(L\)
(Proposition~\ref{prop:r4-inclusion}(b)), and
\(\mathcal R(CB)=\mathcal R(L(ACB))\subseteq\mathcal R(ACB)\). With the first item,
\(\mathcal R(ACB)=\mathcal R(CB)\), and the two ranks, the dimensions of these spaces, are equal.
\item Let \(\rank(CB)=\rank C\). Then \(\mathcal C(CB)\subseteq\mathcal C(C)\) (Theorem~\ref{thm:r4-rank-product}(a)),
both of dimension \(\rank C\), so \(\mathcal C(CB)=\mathcal C(C)\) (Theorem~\ref{thm:r4-counting}(a)). Hence \(C=CBG\)
for some \(G\) (Proposition~\ref{prop:r4-inclusion}(a)), and
\(\mathcal C(AC)=\mathcal C((ACB)G)\subseteq\mathcal C(ACB)\). With the first item,
\(\mathcal C(ACB)=\mathcal C(AC)\), and \(\rank(ACB)=\rank(AC)\).
\end{itemize}
(b) Let \(F=A-B\in\R^{m\times n}\).
\begin{itemize}
\item For (1): by (a), applied to \(C\), \(D\), \(F\) in place of \(A\), \(C\), \(B\) (so that the hypothesis is
\(\rank(CD)=\rank D\)), \(\rank(DF)=\rank(CDF)=\rank(CDA-CDB)=\rank0=0\). So \(DF=0\) (Theorem~\ref{thm:r4-four-bases}),
that is, \(DA=DB\).
\item For (2): by (a), applied to \(F\), \(C\), \(D\) in place of \(A\), \(C\), \(B\) (so that the hypothesis is
\(\rank(CD)=\rank C\)), \(\rank(FC)=\rank(FCD)=\rank(ACD-BCD)=0\). So \(FC=0\), that is, \(AC=BC\).
\end{itemize}
\end{solution}

\begin{problem}[\citeat{abadir-magnus}{Exercise~6.19}]\label{r4-ex-equivalent}
Let \(A,B\in\R^{m\times n}\). Show that \(B=PAQ\) for some invertible \(P\in\R^{m\times m}\) and \(Q\in\R^{n\times n}\)
if and only if \(\rank A=\rank B\).
\end{problem}
\begin{solution}
If \(P\) and \(Q\) are invertible, then \(\rank(PAQ)=\rank A\) (Theorem~\ref{thm:r4-rank-product}(d)). Conversely, let
\(\rank A=\rank B=r\). By Theorem~\ref{thm:r4-normal-form} there are invertible \(P_1,P_2\in\R^{m\times m}\) and \(Q_1,Q_2\in\R^{n\times n}\) with
\(P_1AQ_1=N_r=P_2BQ_2\). Hence
\[
B=P_2^{-1}P_1AQ_1Q_2^{-1}=PAQ,\qquad P=P_2^{-1}P_1,\quad Q=Q_1Q_2^{-1}
\]
and \(P\), \(Q\) are invertible (Theorem~\ref{r2-thm-inverse-props}(a),(b)).
\end{solution}

\begin{problem}[\citeat{prasolov}{Problem~8.1}]\label{r4-ex-pr-rank-two}
Let \(x_1,\dots,x_n,y_1,\dots,y_n\in\R\), and let \(A=(a_{ij})\in\R^{n\times n}\) have the entries \(a_{ij}=x_i+y_j\).
Show that \(\rank A\leqslant2\).
\end{problem}
\begin{solution}
Let
\[
X=\begin{pmatrix}x_1&1\\\vdots&\vdots\\x_n&1\end{pmatrix}\in\R^{n\times2},\qquad
Y=\begin{pmatrix}1&\cdots&1\\y_1&\cdots&y_n\end{pmatrix}\in\R^{2\times n}
\]
By Theorem~\ref{r1-thm-views}(d), \((XY)_{ij}=\row_i(X)\col_j(Y)=x_i\cdot1+1\cdot y_j=a_{ij}\), so \(A=XY\). Hence, by
Theorems~\ref{thm:r4-rank-product}(a) and~\ref{thm:r4-rank-rules}(a),
\(\rank A\leqslant\rank X\leqslant\min\{n,2\}\leqslant2\).
\end{solution}

\begin{problem}[\citeat{prasolov}{Problem~8.7}]\label{r4-ex-pr-odd-order}
Let \(A,B\in\R^{n\times n}\), where \(n\) is odd, satisfy \(AB=0\). Show that at least one of the matrices \(A+A'\) and
\(B+B'\) is not invertible.
\end{problem}
\begin{solution}
Write \(n=2m+1\).
\begin{itemize}
\item By Sylvester's inequality (Theorem~\ref{thm:r4-sylvester}, the case \(AB=0\)),
\(\rank A+\rank B\leqslant n=2m+1\). So \(\rank A\leqslant m\) or \(\rank B\leqslant m\): otherwise
\(\rank A+\rank B\geqslant2(m+1)>2m+1\).
\item If \(\rank A\leqslant m\), then, as \(\rank A'=\rank A\) (Theorem~\ref{thm:r4-rank-rules}(a)),
Corollary~\ref{cor:r4-rank-sum} gives
\[
\rank(A+A')\leqslant\rank A+\rank A'=2\rank A\leqslant2m<n
\]
so \(A+A'\) is not invertible (Theorem~\ref{thm:r4-rank-rules}(d)).
\item If \(\rank B\leqslant m\), the same argument, with \(B\) in place of \(A\), shows that \(B+B'\) is not invertible.
\end{itemize}
\end{solution}

\begin{problem}[cf.\ \citeat{prasolov}{Problems~8.5, 8.6}]\label{r4-ex-pr-rank-additive}
Let \(A_1,A_2\in\R^{m\times n}\).
\begin{enumerate}
\item Show that \(\rank(A_1+A_2)=\rank A_1+\rank A_2\) if and only if \(\mathcal R(A_1)\cap\mathcal R(A_2)=\{0\}\) and
\(\mathcal C(A_1)\cap\mathcal C(A_2)=\{0\}\).
\item Show that if \(A_2'A_1=0\) and \(A_1A_2'=0\), then \(\rank(A_1+A_2)=\rank A_1+\rank A_2\).
\end{enumerate}
\end{problem}
\begin{solution}
(a) Let \(r_i=\rank A_i\) (\(i=1,2\)).
\begin{itemize}
\item If \(r_1=0\), then \(A_1=0\) (Theorem~\ref{thm:r4-four-bases}), so \(\rank(A_1+A_2)=r_2=r_1+r_2\), and
\(\mathcal R(A_1)=\{0\}\) and \(\mathcal C(A_1)=\{0\}\) meet \(\mathcal R(A_2)\) and \(\mathcal C(A_2)\) only in \(0\):
both sides of the equivalence hold. The case \(r_2=0\) is the same. Let \(r_1,r_2\geqslant1\).
\item Let \(A_i=B_iC_i\) be full-rank factorizations (Theorem~\ref{thm:r4-full-rank-factorization}(a)), with
\(B_i\in\R^{m\times r_i}\) of full column rank and \(C_i\in\R^{r_i\times n}\) of full row rank. Then
\(\mathcal C(B_i)=\mathcal C(A_i)\) (Theorem~\ref{thm:r4-full-rank-factorization}(b)). Also
\(\mathcal R(A_i)\subseteq\mathcal R(C_i)\), as \(A_i=B_iC_i\) (Proposition~\ref{prop:r4-inclusion}(b)), and both spaces
have dimension \(r_i\) (Definition~\ref{def:r4-rank}); so \(\mathcal R(C_i)=\mathcal R(A_i)\)
(Theorem~\ref{thm:r4-counting}(a), for row spaces by the paragraph after it).
\item Let \(B=[B_1\ B_2]\in\R^{m\times(r_1+r_2)}\) and
\(C=\left(\begin{smallmatrix}C_1\\C_2\end{smallmatrix}\right)\in\R^{(r_1+r_2)\times n}\), as in
Theorem~\ref{thm:r4-partitioned}(a),(b). The first \(r_1\) columns of \(B\) and rows of \(C\) are those of \(B_1\) and
\(C_1\), and the last \(r_2\) are those of \(B_2\) and \(C_2\). So, by the outer product expansion
(Theorem~\ref{r1-thm-views}(d)), applied to \(BC\), \(B_1C_1\) and \(B_2C_2\),
\[
BC=\sum_{j=1}^{r_1}\col_j(B_1)\row_j(C_1)+\sum_{j=1}^{r_2}\col_j(B_2)\row_j(C_2)=B_1C_1+B_2C_2=A_1+A_2
\]
\item By Theorem~\ref{thm:r4-rank-concat}(a),(d) and the second item,
\[
\rank B=r_1+r_2-\dim\bigl(\mathcal C(A_1)\cap\mathcal C(A_2)\bigr),\qquad
\rank C=r_1+r_2-\dim\bigl(\mathcal R(A_1)\cap\mathcal R(A_2)\bigr)
\]
\item If both intersections are \(\{0\}\), then \(\rank B=r_1+r_2\), so \(B\) has full column rank, and by
Theorem~\ref{thm:r4-rank-product}(b), \(\rank(A_1+A_2)=\rank(BC)=\rank C=r_1+r_2\).
\item Conversely, let \(\rank(A_1+A_2)=r_1+r_2\). By Theorem~\ref{thm:r4-rank-product}(a),
\(r_1+r_2=\rank(BC)\leqslant\min\{\rank B,\rank C\}\), so by the fourth item both intersections have dimension \(0\); a
subspace of dimension \(0\) is \(\{0\}\) (Theorem~\ref{thm:r4-dimension}).
\end{itemize}
(b)
\begin{itemize}
\item \(\mathcal C(A_1)\cap\mathcal C(A_2)=\{0\}\): if \(z=A_1u=A_2w\), then \(z'z=(A_2w)'(A_1u)=w'(A_2'A_1)u=0\), so
\(z=0\), as \(z'z=\sum_iz_i^2\) (Theorem~\ref{r1-thm-transpose}).
\item \(\mathcal R(A_1)\cap\mathcal R(A_2)=\{0\}\): if \(h=y_1'A_1=y_2'A_2\) with \(y_1,y_2\in\R^m\), then
\(hh'=(y_1'A_1)(y_2'A_2)'=y_1'(A_1A_2')y_2=0\), so \(h=0\), as \(hh'\) is the sum of the squares of the entries of
\(h\) (Theorem~\ref{r1-thm-transpose}, for the column \(h'\)).
\item By (a), \(\rank(A_1+A_2)=\rank A_1+\rank A_2\).
\end{itemize}
Neither condition in (a) suffices alone, and the hypothesis \(A_1A_2'=0\) cannot be omitted in (b). For
\(A_1=E_{11}=\left(\begin{smallmatrix}1&0\\0&0\end{smallmatrix}\right)\) and
\(A_2=E_{21}=\left(\begin{smallmatrix}0&0\\1&0\end{smallmatrix}\right)\),
\(\mathcal C(A_1)\cap\mathcal C(A_2)=\Span\{e_1\}\cap\Span\{e_2\}=\{0\}\) and, as \(A_2'=E_{12}\),
\(A_2'A_1=E_{12}E_{11}=\delta_{21}E_{11}=0\) (Proposition~\ref{r1-prop-units}(b),(c)), but
\(\mathcal R(A_1)=\mathcal R(A_2)=\Span\{e_1'\}\) and
\(\rank(A_1+A_2)=\rank\left(\begin{smallmatrix}1&0\\1&0\end{smallmatrix}\right)=1<2\). The transposes \(A_1'\), \(A_2'\)
have \(\mathcal R(A_1')\cap\mathcal R(A_2')=\{0\}\), and again \(\rank(A_1'+A_2')=1<2\).
\end{solution}

\clearpage
\def\unitnumber{5}
\setcounter{theorem}{0}\setcounter{problem}{0}
\section*{Chapter 5. Block matrices and determinants}
\subsection*{Block matrices}

\begin{definition}[\citeat{harville}{\S2.1}, \citeat{horn-johnson}{\S0.7.2}]\label{r5b-def-partition}
Let \(A\in\R^{m\times n}\). A \textbf{partition} of \(A\) splits the rows into consecutive groups
of \(m_1,\dots,m_r\) rows (\(m_1+\dots+m_r=m\)) and the columns into consecutive groups of
\(n_1,\dots,n_s\) columns (\(n_1+\dots+n_s=n\)). Put \(M_0=0\), \(M_i=m_1+\dots+m_i\) and
\(N_0=0\), \(N_j=n_1+\dots+n_j\): row group \(i\) consists of the rows \(M_{i-1}+1,\dots,M_i\),
and column group \(j\) of the columns \(N_{j-1}+1,\dots,N_j\). The \textbf{block}
\(A_{ij}\in\R^{m_i\times n_j}\) consists of the entries in row group \(i\) and column group \(j\):
\[
(A_{ij})_{xy}=a_{M_{i-1}+x,\,N_{j-1}+y}\qquad(1\leqslant x\leqslant m_i,\ 1\leqslant y\leqslant n_j)
\]
We write
\[
A=\begin{pmatrix}A_{11}&\cdots&A_{1s}\\\vdots&&\vdots\\A_{r1}&\cdots&A_{rs}\end{pmatrix}
\]
Lines between rows and between columns may mark the partition. The group sizes are arbitrary
positive integers, so a block \(A_{ij}\) is square only when \(m_i=n_j\). Two matrices of the same
size, partitioned in the same way, are equal if and only if their corresponding blocks are equal (both statements say that all entries agree,
Definition~\ref{r1-def-matrix}).
\end{definition}

\begin{definition}[\citeat{harville}{\S2.2}]\label{r5b-def-conformal}
Let \(A\in\R^{m\times n}\) be partitioned with row groups \(m_1,\dots,m_r\) and column groups
\(n_1,\dots,n_s\).
\begin{enumerate}
\item A partition of \(B\in\R^{m\times n}\) is \textbf{conformal for addition} with that of \(A\)
if it has the same row groups \(m_1,\dots,m_r\) and the same column groups \(n_1,\dots,n_s\).
\item A partition of \(B\in\R^{n\times p}\) is \textbf{conformal for multiplication} with that
of \(A\) (in the product \(AB\)) if its row groups are \(n_1,\dots,n_s\), the column groups of
\(A\). Its column groups \(p_1,\dots,p_t\) (\(p_1+\dots+p_t=p\)) may be any.
\end{enumerate}
\end{definition}
Conformality for multiplication restricts only the column groups of \(A\) and the row groups of
\(B\); the row groups \(m_i\) of \(A\) and the column groups \(p_k\) of \(B\) are free. It means
exactly that every block \(A_{ij}\in\R^{m_i\times n_j}\) has as many columns as every block
\(B_{jk}\in\R^{n_j\times p_k}\) has rows, so that the product \(A_{ij}B_{jk}\in\R^{m_i\times p_k}\)
is defined (Definition~\ref{r1-def-product}).

\begin{theorem}[\citeat{harville}{(2.2.1)--(2.2.6)}]\label{r5b-thm-block}
Let \(A\in\R^{m\times n}\) be partitioned into blocks \(A_{ij}\in\R^{m_i\times n_j}\)
(\(1\leqslant i\leqslant r\), \(1\leqslant j\leqslant s\)).
\begin{enumerate}
\item If \(B\in\R^{m\times n}\) is partitioned conformally for addition, into blocks
\(B_{ij}\in\R^{m_i\times n_j}\), then \(A+B\) has the blocks \(A_{ij}+B_{ij}\), and for
\(c\in\R\) the matrix \(cA\) has the blocks \(cA_{ij}\) (same partition).
\item Partition \(A'\in\R^{n\times m}\) with row groups \(n_1,\dots,n_s\) and column groups
\(m_1,\dots,m_r\). Its \((j,i)\) block is \(A_{ij}'\in\R^{n_j\times m_i}\): transpose the array
of blocks and also each block.
\item If \(B\in\R^{n\times p}\) is partitioned conformally for multiplication, into blocks
\(B_{jk}\in\R^{n_j\times p_k}\) (\(1\leqslant j\leqslant s\), \(1\leqslant k\leqslant t\)), then
\(AB\in\R^{m\times p}\), partitioned with row groups \(m_1,\dots,m_r\) and column groups
\(p_1,\dots,p_t\), has the blocks
\[
(AB)_{ik}=\sum_{j=1}^sA_{ij}B_{jk}=A_{i1}B_{1k}+A_{i2}B_{2k}+\dots+A_{is}B_{sk}\in\R^{m_i\times p_k}
\]
where the term \(A_{ij}B_{jk}\) is the product of an \(m_i\times n_j\) and an \(n_j\times p_k\)
matrix.
\end{enumerate}
\end{theorem}
So conformal block matrices are added and multiplied as if the blocks were numbers, but the order
of the factors in each \(A_{ij}B_{jk}\) must be kept. No group size has to equal another, and no
block has to be square. For \(r=s=t=2\),
\[
\begin{pmatrix}A_{11}&A_{12}\\A_{21}&A_{22}\end{pmatrix}
\begin{pmatrix}B_{11}&B_{12}\\B_{21}&B_{22}\end{pmatrix}
=\begin{pmatrix}A_{11}B_{11}+A_{12}B_{21}&A_{11}B_{12}+A_{12}B_{22}\\
A_{21}B_{11}+A_{22}B_{21}&A_{21}B_{12}+A_{22}B_{22}\end{pmatrix}
\]
whenever \(A_{11}\) and \(A_{21}\) have as many columns as \(B_{11}\) and \(B_{12}\) have rows
(the second groups then match as well, since \(n_1+n_2=n\) is the number of columns of \(A\) and
of rows of \(B\)).
The cases \(r=t=1,\ s=2\); \(r=t=2,\ s=1\); and \(r=s=1,\ t=2\) read
\[
[A_1\ A_2]\begin{pmatrix}B_1\\B_2\end{pmatrix}=A_1B_1+A_2B_2,\qquad
\begin{pmatrix}A_1\\A_2\end{pmatrix}[B_1\ B_2]=\begin{pmatrix}A_1B_1&A_1B_2\\A_2B_1&A_2B_2\end{pmatrix},\qquad
A[B_1\ B_2]=[AB_1\ \ AB_2]
\]
\begin{proof}
Let \(M_i\) and \(N_j\) be as in Definition~\ref{r5b-def-partition}. Each row index \(u\) of \(A\) lies in exactly one row group \(i\), and then
\(u=M_{i-1}+x\) with \(1\leqslant x\leqslant m_i\); likewise for columns.

(a) Fix \(i,j\) and \(1\leqslant x\leqslant m_i\), \(1\leqslant y\leqslant n_j\), and put
\(u=M_{i-1}+x\), \(v=N_{j-1}+y\). Since \(B\) has the same groups as \(A\), entry \((x,y)\) of
\(B_{ij}\) is \(b_{uv}\). By Definition~\ref{r1-def-matrix}, entry \((x,y)\) of block \((i,j)\) of
\(A+B\) is
\[
(A+B)_{uv}=a_{uv}+b_{uv}=(A_{ij})_{xy}+(B_{ij})_{xy}=(A_{ij}+B_{ij})_{xy}
\]
and in the same way \((cA)_{uv}=ca_{uv}=(cA_{ij})_{xy}\).

(b) Row group \(j\) of \(A'\) consists of the rows \(N_{j-1}+1,\dots,N_j\) and column group \(i\) of
the columns \(M_{i-1}+1,\dots,M_i\). By Definition~\ref{r1-def-transpose}, entry \((y,x)\) of the
\((j,i)\) block of \(A'\) is
\[
(A')_{N_{j-1}+y,\,M_{i-1}+x}=a_{M_{i-1}+x,\,N_{j-1}+y}=(A_{ij})_{xy}=(A_{ij}')_{yx}
\]

(c) For a row group \(i\) of \(A\) and a column group \(k\) of \(B\), let \(A_{i\ast}\in\R^{m_i\times n}\) consist of
the rows of \(A\) in row group \(i\), and \(B_{\ast k}\in\R^{n\times p_k}\) of the columns of \(B\) in column group
\(k\). By Definition~\ref{r5b-def-partition}, the columns of \(A_{i\ast}\) in column group \(j\) are, in order, the
columns of \(A_{ij}\); since the row groups of \(B\) are \(n_1,\dots,n_s\) (conformality), the rows of \(B_{\ast k}\)
in row group \(j\) are, in order, the rows of \(B_{jk}\). In blocks,
\[
A_{i\ast}=[A_{i1}\ \cdots\ A_{is}],\qquad
B_{\ast k}=\begin{pmatrix}B_{1k}\\\vdots\\B_{sk}\end{pmatrix}
\]
\begin{itemize}
\item Columns. By Theorem~\ref{r1-thm-views}(a), each column of \(AB\) is \(A\) times the column of \(B\) with the
same index. The columns of \(B\) in column group \(k\) are, in order, the columns of \(B_{\ast k}\), so the columns of
\(AB\) in column group \(k\) are, in order, \(A\) times the columns of \(B_{\ast k}\), that is, the columns of the
product \(AB_{\ast k}\) (again Theorem~\ref{r1-thm-views}(a)).
\item Rows. By Theorem~\ref{r1-thm-views}(c), each row of \(AB_{\ast k}\) is the row of \(A\) with the same index
times \(B_{\ast k}\). The rows of \(A\) in row group \(i\) are, in order, the rows of \(A_{i\ast}\), so the rows of
\(AB_{\ast k}\) in row group \(i\) are, in order, the rows of \(A_{i\ast}B_{\ast k}\) (again
Theorem~\ref{r1-thm-views}(c)). These are the entries of \(AB\) in row group \(i\) and column group \(k\), that is,
the block \((AB)_{ik}\) (Definition~\ref{r5b-def-partition}): \((AB)_{ik}=A_{i\ast}B_{\ast k}\in\R^{m_i\times p_k}\).
\item Inner sum. By the outer product expansion (Theorem~\ref{r1-thm-views}(d)),
\[
A_{i\ast}B_{\ast k}=\sum_{l=1}^n\col_l(A_{i\ast})\row_l(B_{\ast k})
\]
The column groups of \(A\) cut \(l=1,\dots,n\) into \(s\) consecutive stretches of \(n_1,\dots,n_s\) indices. In
stretch \(j\), the columns \(\col_l(A_{i\ast})\) are \(\col_1(A_{ij}),\dots,\col_{n_j}(A_{ij})\) and the rows
\(\row_l(B_{\ast k})\) are \(\row_1(B_{jk}),\dots,\row_{n_j}(B_{jk})\), in this order (see above). So the terms of
stretch \(j\) add up to
\[
\sum_{y=1}^{n_j}\col_y(A_{ij})\row_y(B_{jk})=A_{ij}B_{jk}
\]
the outer product expansion of the product of \(A_{ij}\in\R^{m_i\times n_j}\) and \(B_{jk}\in\R^{n_j\times p_k}\)
(Theorem~\ref{r1-thm-views}(d)). Adding the \(s\) stretches,
\[
(AB)_{ik}=A_{i\ast}B_{\ast k}=\sum_{j=1}^sA_{ij}B_{jk}\qedhere
\]
\end{itemize}
\end{proof}
In entries: let \(u=M_{i-1}+x\) (row \(x\) of row group \(i\)) and let column \(v\) be the \(z\)th column of
column group \(k\) of \(B\). By Definition~\ref{r1-def-product}, \((AB)_{uv}=\sum_{l=1}^na_{ul}b_{lv}\). Split the
sum into the stretches of the column groups of \(A\) and write \(l=N_{j-1}+y\) in stretch \(j\); by
Definition~\ref{r5b-def-partition} (for \(A\), and for \(B\) with row groups \(n_1,\dots,n_s\)),
\[
\sum_{l=N_{j-1}+1}^{N_j}a_{ul}b_{lv}=\sum_{y=1}^{n_j}(A_{ij})_{xy}(B_{jk})_{yz}=(A_{ij}B_{jk})_{xz}
\]
so the part of stretch \(j\) is the entry \((x,z)\) of \(A_{ij}B_{jk}\): this is (c), entry by entry.

\begin{example}\label{r5b-exm-block}
Partition \(A\in\R^{5\times6}\) with row groups \((2,3)\) and column groups \((4,2)\), and
\(B\in\R^{6\times5}\) with row groups \((4,2)\) and column groups \((3,2)\):
\[
A=\left(\begin{array}{cccc|cc}1&1&1&1&1&0\\1&-1&1&-1&0&1\\\hline1&1&1&1&2&0\\1&1&1&1&0&1\\1&1&1&1&1&1\end{array}\right)
=\begin{pmatrix}A_{11}&A_{12}\\A_{21}&A_{22}\end{pmatrix},\qquad
B=\left(\begin{array}{ccc|cc}1&0&0&1&0\\0&1&0&0&1\\0&0&1&1&0\\1&1&1&0&1\\\hline1&0&-1&1&0\\0&1&-1&0&1\end{array}\right)
=\begin{pmatrix}B_{11}&B_{12}\\B_{21}&B_{22}\end{pmatrix}
\]
The blocks are
\[
\begin{gathered}
A_{11}=\begin{pmatrix}1&1&1&1\\1&-1&1&-1\end{pmatrix}\in\R^{2\times4},\quad A_{12}=I_2,\quad
A_{21}=\one_3\one_4'\in\R^{3\times4},\quad A_{22}=\begin{pmatrix}2&0\\0&1\\1&1\end{pmatrix}\in\R^{3\times2}\\
B_{11}=\begin{pmatrix}I_3\\\one_3'\end{pmatrix}\in\R^{4\times3},\quad
B_{12}=\begin{pmatrix}1&0\\0&1\\1&0\\0&1\end{pmatrix}\in\R^{4\times2},\quad
B_{21}=\begin{pmatrix}1&0&-1\\0&1&-1\end{pmatrix}\in\R^{2\times3},\quad B_{22}=I_2
\end{gathered}
\]
Only \(A_{12}\) and \(B_{22}\) are square.
\begin{enumerate}
\item The column groups \((4,2)\) of \(A\) are the row groups of \(B\), so the partitions are
conformal for multiplication (Definition~\ref{r5b-def-conformal}). By
Theorem~\ref{r5b-thm-block}(c), \(AB\in\R^{5\times5}\) has row groups \((2,3)\), those of \(A\),
and column groups \((3,2)\), those of \(B\). Every term is a product of a block with \(n_j\) columns
and a block with \(n_j\) rows:
\[
\begin{aligned}
(AB)_{11}&=A_{11}B_{11}+A_{12}B_{21}:&&(2\times4)(4\times3)+(2\times2)(2\times3)&&\text{is }2\times3\\
(AB)_{12}&=A_{11}B_{12}+A_{12}B_{22}:&&(2\times4)(4\times2)+(2\times2)(2\times2)&&\text{is }2\times2\\
(AB)_{21}&=A_{21}B_{11}+A_{22}B_{21}:&&(3\times4)(4\times3)+(3\times2)(2\times3)&&\text{is }3\times3\\
(AB)_{22}&=A_{21}B_{12}+A_{22}B_{22}:&&(3\times4)(4\times2)+(3\times2)(2\times2)&&\text{is }3\times2
\end{aligned}
\]
\item The column sums of \(B_{11}\) and \(B_{12}\) are \(\one_4'B_{11}=(2,2,2)\) and
\(\one_4'B_{12}=(2,2)\), so \(A_{21}B_{11}=(\one_3\one_4')B_{11}=\one_3(\one_4'B_{11})\) (Theorem~\ref{r1-thm-rules}) has all entries \(2\), and so
has \(A_{21}B_{12}\). With \(A_{12}=B_{22}=I_2\),
\[
\begin{aligned}
(AB)_{11}&=A_{11}B_{11}+B_{21}=\begin{pmatrix}2&2&2\\0&-2&0\end{pmatrix}+\begin{pmatrix}1&0&-1\\0&1&-1\end{pmatrix}
=\begin{pmatrix}3&2&1\\0&-1&-1\end{pmatrix}\\
(AB)_{12}&=A_{11}B_{12}+I_2=\begin{pmatrix}2&2\\2&-2\end{pmatrix}+I_2=\begin{pmatrix}3&2\\2&-1\end{pmatrix}\\
(AB)_{21}&=A_{21}B_{11}+A_{22}B_{21}=\begin{pmatrix}2&2&2\\2&2&2\\2&2&2\end{pmatrix}+\begin{pmatrix}2&0&-2\\0&1&-1\\1&1&-2\end{pmatrix}
=\begin{pmatrix}4&2&0\\2&3&1\\3&3&0\end{pmatrix}\\
(AB)_{22}&=A_{21}B_{12}+A_{22}=\begin{pmatrix}2&2\\2&2\\2&2\end{pmatrix}+\begin{pmatrix}2&0\\0&1\\1&1\end{pmatrix}
=\begin{pmatrix}4&2\\2&3\\3&3\end{pmatrix}
\end{aligned}
\]
So
\[
AB=\left(\begin{array}{ccc|cc}3&2&1&3&2\\0&-1&-1&2&-1\\\hline4&2&0&4&2\\2&3&1&2&3\\3&3&0&3&3\end{array}\right)
\]
\item Check of one entry. Entry \((4,2)\) of \(AB\) lies in row group \(2\) (row \(u=4=M_1+2\),
so \(x=2\)) and column group \(1\) (\(z=2\)). As in the paragraph after the proof of Theorem~\ref{r5b-thm-block},
the sum over \(l=1,\dots,6\) splits at \(N_1=4\):
\[
(AB)_{42}=\underbrace{1\cdot0+1\cdot1+1\cdot0+1\cdot1}_{l=1,\dots,4}+\underbrace{0\cdot0+1\cdot1}_{l=5,6}=2+1=3
\]
The first part is row \(2\) of \(A_{21}\) times column \(2\) of \(B_{11}\), entry \((2,2)\) of
\(A_{21}B_{11}\); the second is row \(2\) of \(A_{22}\) times column \(2\) of \(B_{21}\), entry
\((2,2)\) of \(A_{22}B_{21}\).
\item A partition of \(B\) that is not conformal. Give \(B\) the row groups \((3,3)\), with any
column groups. The product \(AB\) is still defined, but \(A_{11}\in\R^{2\times4}\) has \(4\) columns
while the new \(B_{11}\) has \(3\) rows, so \(A_{11}B_{11}\) is not defined, and the blocks of
\(AB\) cannot be computed from these blocks. By contrast, the column groups of \(B\) are free: with its row groups
\((4,2)\) and the column groups changed, say to \((1,4)\), the partitions are still conformal
(Definition~\ref{r5b-def-conformal}).
\item By Theorem~\ref{r5b-thm-block}(b), \(A'\in\R^{6\times5}\) with row groups \((4,2)\) and
column groups \((2,3)\) is
\(A'=\left(\begin{smallmatrix}A_{11}'&A_{21}'\\A_{12}'&A_{22}'\end{smallmatrix}\right)\), where
\(A_{11}'\in\R^{4\times2}\), \(A_{21}'=\one_4\one_3'\in\R^{4\times3}\), \(A_{12}'=I_2\) and
\(A_{22}'\in\R^{2\times3}\).
\end{enumerate}
\end{example}

\begin{definition}[\citeat{harville}{\S2.1}]\label{r5b-def-block-tri}
Let \(A\in\R^{n\times n}\) be partitioned with the same groups \(n_1,\dots,n_r\) for rows and
columns. Then the \textbf{diagonal blocks} \(A_{ii}\in\R^{n_i\times n_i}\) are square, while an
off-diagonal block \(A_{ij}\in\R^{n_i\times n_j}\) (\(i\ne j\)) is rectangular when
\(n_i\ne n_j\). \(A\) is \textbf{block diagonal}, written \(A=\diag(A_{11},\dots,A_{rr})\), if
\(A_{ij}=0\) for \(i\ne j\); \textbf{upper block triangular} if \(A_{ij}=0\) for \(i>j\); and
\textbf{lower block triangular} if \(A_{ij}=0\) for \(i<j\).
\end{definition}

\begin{proposition}[\citeat{searle-khuri}{\S\S4.10, 4.12(f)}, \citeat{harville}{\S2.1}]\label{r5b-prop-block-tri}
Let \(A,B\in\R^{n\times n}\) be partitioned as in Definition~\ref{r5b-def-block-tri}, with the
same groups.
\begin{enumerate}
\item If \(A\) and \(B\) are upper (lower) block triangular, so is \(AB\), and its diagonal
blocks are \(A_{ii}B_{ii}\).
\item \(\diag(A_{11},\dots,A_{rr})\diag(B_{11},\dots,B_{rr})=\diag(A_{11}B_{11},\dots,A_{rr}B_{rr})\).
If every \(A_{ii}\) is invertible, then
\(\diag(A_{11},\dots,A_{rr})^{-1}=\diag(A_{11}^{-1},\dots,A_{rr}^{-1})\).
\item If \(A_{11}\in\R^{p\times p}\) and \(A_{22}\in\R^{q\times q}\) are invertible and
\(A_{12}\in\R^{p\times q}\), then
\[
\begin{pmatrix}A_{11}&A_{12}\\0&A_{22}\end{pmatrix}^{-1}
=\begin{pmatrix}A_{11}^{-1}&-A_{11}^{-1}A_{12}A_{22}^{-1}\\0&A_{22}^{-1}\end{pmatrix}
\]
\end{enumerate}
\end{proposition}
In (c) only the diagonal blocks must be square; for \(p\ne q\) the blocks \(A_{12}\) and
\(-A_{11}^{-1}A_{12}A_{22}^{-1}\) are \(p\times q\).
\begin{proof}
(a)
\begin{itemize}
\item The column groups \(n_1,\dots,n_r\) of \(A\) are the row groups of \(B\), so by
Theorem~\ref{r5b-thm-block}(c) the \((i,k)\) block of \(AB\) is
\(\sum_{j=1}^rA_{ij}B_{jk}\in\R^{n_i\times n_k}\).
\item Let \(A\) and \(B\) be upper block triangular. Then \(A_{ij}=0\) unless \(i\leqslant j\), and
\(B_{jk}=0\) unless \(j\leqslant k\). So the term \(A_{ij}B_{jk}\) has a zero factor, and is \(0\),
unless \(i\leqslant j\leqslant k\).
\item If \(i>k\), no \(j\) satisfies \(i\leqslant j\leqslant k\), so \((AB)_{ik}=0\): \(AB\) is
upper block triangular. If \(i=k\), only \(j=i\) remains, so \((AB)_{ii}=A_{ii}B_{ii}\).
\item For lower block triangular \(A,B\), the term \(j\) can be nonzero only if
\(i\geqslant j\geqslant k\); the same argument applies with the inequalities reversed.
\end{itemize}
(b) A block diagonal matrix is upper and lower block triangular, so by (a) the product is upper
and lower block triangular, that is, block diagonal, with diagonal blocks \(A_{ii}B_{ii}\). With
\(B_{ii}=A_{ii}^{-1}\) this gives
\(\diag(A_{11},\dots,A_{rr})\diag(A_{11}^{-1},\dots,A_{rr}^{-1})=\diag(I_{n_1},\dots,I_{n_r})=I_n\)
(both have entry \(1\) at every \((u,u)\) and \(0\) elsewhere), and likewise in the other order; so the second matrix is the inverse of the first
(Definition~\ref{r2-def-inverse}).

(c) Call the two matrices \(T\) and \(X\); both have row and column groups \((p,q)\), so the
\(2\times2\) block formula of Theorem~\ref{r5b-thm-block}(c) applies. The four blocks of \(TX\) are
\[
\begin{aligned}
(TX)_{11}&=A_{11}A_{11}^{-1}+A_{12}\cdot0=I_p,&
(TX)_{12}&=A_{11}\bigl(-A_{11}^{-1}A_{12}A_{22}^{-1}\bigr)+A_{12}A_{22}^{-1}=0\\
(TX)_{21}&=0\cdot A_{11}^{-1}+A_{22}\cdot0=0,&
(TX)_{22}&=0\cdot\bigl(-A_{11}^{-1}A_{12}A_{22}^{-1}\bigr)+A_{22}A_{22}^{-1}=I_q
\end{aligned}
\]
and those of \(XT\) are
\[
\begin{aligned}
(XT)_{11}&=A_{11}^{-1}A_{11}+\bigl(-A_{11}^{-1}A_{12}A_{22}^{-1}\bigr)\cdot0=I_p,&
(XT)_{12}&=A_{11}^{-1}A_{12}-A_{11}^{-1}A_{12}A_{22}^{-1}A_{22}=0\\
(XT)_{21}&=0\cdot A_{11}+A_{22}^{-1}\cdot0=0,&
(XT)_{22}&=0\cdot A_{12}+A_{22}^{-1}A_{22}=I_q
\end{aligned}
\]
So \(TX=XT=I_{p+q}\), and \(X=T^{-1}\) (Definition~\ref{r2-def-inverse}).
\end{proof}

\begin{example}\label{r5b-exm-block-inverse}
Find the inverse of
\[
T=\left(\begin{array}{cc|ccc}3&2&1&0&2\\1&1&1&1&1\\\hline0&0&1&1&0\\0&0&0&1&1\\0&0&0&0&1\end{array}\right)
=\begin{pmatrix}A_{11}&A_{12}\\0&A_{22}\end{pmatrix}\in\R^{5\times5}
\]

\textit{Solution.} The groups are \((2,3)\): \(A_{11}=\left(\begin{smallmatrix}3&2\\1&1\end{smallmatrix}\right)\)
is \(2\times2\), \(A_{12}=\left(\begin{smallmatrix}1&0&2\\1&1&1\end{smallmatrix}\right)\) is \(2\times3\), and
\[
A_{22}=\begin{pmatrix}1&1&0\\0&1&1\\0&0&1\end{pmatrix}
\]
is \(3\times3\). By Proposition~\ref{r2-prop-inverse-2x2} (\(3\cdot1-2\cdot1=1\)),
\(A_{11}^{-1}=\left(\begin{smallmatrix}1&-2\\-1&3\end{smallmatrix}\right)\). For
\[
Y=\begin{pmatrix}1&-1&1\\0&1&-1\\0&0&1\end{pmatrix},\qquad
A_{22}Y=\begin{pmatrix}1&-1+1&1-1\\0&1&-1+1\\0&0&1\end{pmatrix}=I_3
\]
so \(A_{22}^{-1}=Y\) (Theorem~\ref{r2-thm-invertible}). Then
\[
A_{11}^{-1}A_{12}=\begin{pmatrix}1&-2\\-1&3\end{pmatrix}\begin{pmatrix}1&0&2\\1&1&1\end{pmatrix}
=\begin{pmatrix}-1&-2&0\\2&3&1\end{pmatrix}
\]
\[
-A_{11}^{-1}A_{12}A_{22}^{-1}=-\begin{pmatrix}-1&-2&0\\2&3&1\end{pmatrix}\begin{pmatrix}1&-1&1\\0&1&-1\\0&0&1\end{pmatrix}
=\begin{pmatrix}1&1&-1\\-2&-1&0\end{pmatrix}
\]
This block is \(2\times3\). By Proposition~\ref{r5b-prop-block-tri}(c) with \(p=2\), \(q=3\),
\[
T^{-1}=\left(\begin{array}{cc|ccc}1&-2&1&1&-1\\-1&3&-2&-1&0\\\hline0&0&1&-1&1\\0&0&0&1&-1\\0&0&0&0&1\end{array}\right)
\]
\end{example}

\subsection*{Block elimination and the Schur complement}

\begin{proposition}[\citeat{abadir-magnus}{Exercises~5.6, 5.7}, \citeat{harville}{Lemma~8.5.2}]\label{r5b-prop-block-elementary}
Let \(M=\begin{pmatrix}A&B\\C&D\end{pmatrix}\) with \(A\in\R^{m\times p}\),
\(B\in\R^{m\times q}\), \(C\in\R^{n\times p}\), \(D\in\R^{n\times q}\), and let
\(V\in\R^{n\times m}\), \(F\in\R^{m\times n}\), \(H\in\R^{p\times q}\). Then
\[
\begin{pmatrix}I_m&0\\V&I_n\end{pmatrix}M=\begin{pmatrix}A&B\\C+VA&D+VB\end{pmatrix},\qquad
\begin{pmatrix}I_m&F\\0&I_n\end{pmatrix}M=\begin{pmatrix}A+FC&B+FD\\C&D\end{pmatrix}
\]
\[
M\begin{pmatrix}I_p&H\\0&I_q\end{pmatrix}=\begin{pmatrix}A&B+AH\\C&D+CH\end{pmatrix}
\]
These \textbf{block elementary matrices} are invertible:
\(\left(\begin{smallmatrix}I_m&0\\V&I_n\end{smallmatrix}\right)^{-1}=\left(\begin{smallmatrix}I_m&0\\-V&I_n\end{smallmatrix}\right)\),
\(\left(\begin{smallmatrix}I_m&F\\0&I_n\end{smallmatrix}\right)^{-1}=\left(\begin{smallmatrix}I_m&-F\\0&I_n\end{smallmatrix}\right)\)
and
\(\left(\begin{smallmatrix}I_p&H\\0&I_q\end{smallmatrix}\right)^{-1}=\left(\begin{smallmatrix}I_p&-H\\0&I_q\end{smallmatrix}\right)\).
\end{proposition}
None of \(A,B,C,D,V,F,H\) needs to be square.
\begin{proof}
\begin{itemize}
\item \(M\) has row groups \((m,n)\) and column groups \((p,q)\). The matrices
\(\left(\begin{smallmatrix}I_m&0\\V&I_n\end{smallmatrix}\right)\) and
\(\left(\begin{smallmatrix}I_m&F\\0&I_n\end{smallmatrix}\right)\) have column groups \((m,n)\), and
\(\left(\begin{smallmatrix}I_p&H\\0&I_q\end{smallmatrix}\right)\) has row groups \((p,q)\); so all
three products are conformal, and Theorem~\ref{r5b-thm-block}(c) applies.
\item First product: block row \(1\) is \((I_mA+0C,\ I_mB+0D)=(A,\ B)\), and block row \(2\) is
\((VA+I_nC,\ VB+I_nD)=(C+VA,\ D+VB)\), where \(VA\in\R^{n\times p}\) and \(VB\in\R^{n\times q}\) have
the sizes of \(C\) and \(D\).
\item Second product: block row \(1\) is \((I_mA+FC,\ I_mB+FD)=(A+FC,\ B+FD)\), and block row \(2\)
is \((0A+I_nC,\ 0B+I_nD)=(C,\ D)\).
\item Third product: block column \(1\) is \((AI_p+B0,\ CI_p+D0)=(A,\ C)\), and block column \(2\)
is \((AH+BI_q,\ CH+DI_q)=(B+AH,\ D+CH)\).
\item Thus left multiplication by \(\left(\begin{smallmatrix}I&0\\V&I\end{smallmatrix}\right)\) adds
\(V\) times the first block row to the second, the block analogue of the row operation
\(R_2\leftarrow R_2+cR_1\) of Definition~\ref{r2-def-rowops}; right multiplication by
\(\left(\begin{smallmatrix}I&H\\0&I\end{smallmatrix}\right)\) adds the first block column times
\(H\) to the second.
\item Inverses. The first identity with \(A=I_m\), \(B=0\), \(C=V\), \(D=I_n\) (so \(p=m\),
\(q=n\)) and with \(-V\) in place of \(V\) gives
\[
\begin{pmatrix}I_m&0\\-V&I_n\end{pmatrix}\begin{pmatrix}I_m&0\\V&I_n\end{pmatrix}
=\begin{pmatrix}I_m&0\\V-VI_m&I_n-V\cdot0\end{pmatrix}=I_{m+n}
\]
and with \(C=-V\) and \(V\) itself it gives the product in the other order,
\(\left(\begin{smallmatrix}I_m&0\\-V+V&I_n\end{smallmatrix}\right)=I_{m+n}\). In the same way the
second identity gives
\(\left(\begin{smallmatrix}I_m&-F\\0&I_n\end{smallmatrix}\right)\left(\begin{smallmatrix}I_m&F\\0&I_n\end{smallmatrix}\right)
=\left(\begin{smallmatrix}I_m&F-F\\0&I_n\end{smallmatrix}\right)=I_{m+n}\) and the product in the
other order; the third case is the second with \(p,q,H\) in place of \(m,n,F\).\qedhere
\end{itemize}
\end{proof}

\begin{theorem}[\citeat{abadir-magnus}{Exercise~5.10}]\label{r5b-thm-schur}
Let \(M=\begin{pmatrix}A&B\\C&D\end{pmatrix}\) with \(A\in\R^{p\times p}\) invertible,
\(B\in\R^{p\times q}\), \(C\in\R^{n\times p}\), \(D\in\R^{n\times q}\). Then
\[
\begin{pmatrix}I_p&0\\-CA^{-1}&I_n\end{pmatrix}\begin{pmatrix}A&B\\C&D\end{pmatrix}
=\begin{pmatrix}A&B\\0&D-CA^{-1}B\end{pmatrix},\qquad
\begin{pmatrix}A&B\\C&D\end{pmatrix}\begin{pmatrix}I_p&-A^{-1}B\\0&I_q\end{pmatrix}
=\begin{pmatrix}A&0\\C&D-CA^{-1}B\end{pmatrix}
\]
and hence
\(\left(\begin{smallmatrix}I_p&0\\-CA^{-1}&I_n\end{smallmatrix}\right)M\left(\begin{smallmatrix}I_p&-A^{-1}B\\0&I_q\end{smallmatrix}\right)
=\left(\begin{smallmatrix}A&0\\0&D-CA^{-1}B\end{smallmatrix}\right)\). The matrix
\(D-CA^{-1}B\in\R^{n\times q}\) is the \textbf{Schur complement} of \(A\) in \(M\).
\end{theorem}
Only \(A\) has to be square: \(B\), \(C\), \(D\) and \(M\) itself may be rectangular.
\begin{proof}
\begin{itemize}
\item Apply Proposition~\ref{r5b-prop-block-elementary} with \(m=p\). The choice
\(V=-CA^{-1}\in\R^{n\times p}\) makes the \((2,1)\) block \(C+VA=C-CA^{-1}A=0\), and then
\(D+VB=D-CA^{-1}B\). This is the first identity.
\item The choice \(H=-A^{-1}B\in\R^{p\times q}\) makes \(B+AH=B-AA^{-1}B=0\), and then
\(D+CH=D-CA^{-1}B\). This is the second identity.
\item Put \(S=D-CA^{-1}B\). Applying the column step to the result of the row step, by the third
identity of Proposition~\ref{r5b-prop-block-elementary} with \(C\) replaced by \(0\) and \(D\) by
\(S\),
\(\left(\begin{smallmatrix}A&B\\0&S\end{smallmatrix}\right)\left(\begin{smallmatrix}I_p&-A^{-1}B\\0&I_q\end{smallmatrix}\right)
=\left(\begin{smallmatrix}A&B-AA^{-1}B\\0&S-0\cdot A^{-1}B\end{smallmatrix}\right)=\left(\begin{smallmatrix}A&0\\0&S\end{smallmatrix}\right)\).\qedhere
\end{itemize}
\end{proof}
In Theorem~\ref{r5b-thm-schur} the left factor adds \(-CA^{-1}\) times the first block row of \(M\) to
the second, and the right factor adds the first block column times \(-A^{-1}B\) to the second
(Proposition~\ref{r5b-prop-block-elementary}). When \(D\) is invertible, the matrix \(A-BD^{-1}C\) is
likewise called the \textbf{Schur complement} of \(D\) in \(M\).

Let \(R_1,R_2\) be the block rows of a matrix with row groups \((p,q)\). The \textbf{block row
operations} are \(R_1\leftarrow XR_1\) and \(R_2\leftarrow YR_2\) with \(X\in\R^{p\times p}\) and
\(Y\in\R^{q\times q}\) invertible, \(R_2\leftarrow R_2+VR_1\) with \(V\in\R^{q\times p}\), and
\(R_1\leftarrow R_1+WR_2\) with \(W\in\R^{p\times q}\). They are the left multiplications by
\(\diag(X,I_q)\), \(\diag(I_p,Y)\), \(\left(\begin{smallmatrix}I_p&0\\V&I_q\end{smallmatrix}\right)\) and
\(\left(\begin{smallmatrix}I_p&W\\0&I_q\end{smallmatrix}\right)\) (Theorem~\ref{r5b-thm-block}(c) and
Proposition~\ref{r5b-prop-block-elementary}), and these matrices are invertible
(Propositions~\ref{r5b-prop-block-tri}(b) and~\ref{r5b-prop-block-elementary}). \textbf{Block
elimination} applies the steps of Gauss--Jordan elimination (Definition~\ref{r2-def-gauss-jordan}) to
the block rows of \([M\mid I]\), with the inverse of a pivot block in place of the reciprocal of a pivot
entry: the first pivot is a diagonal block of \(M\), the second its Schur complement, and the result is
\([I\mid M^{-1}]\), as in Corollary~\ref{r2-cor-computing-inverse}.

\begin{corollary}[Block elimination; \citeat{abadir-magnus}{Exercise~5.10}, \citeat{harville}{\S8.5}]\label{r5b-cor-schur-inverse}
Let \(M=\begin{pmatrix}A&B\\C&D\end{pmatrix}\) with \(A\in\R^{p\times p}\), \(B\in\R^{p\times q}\),
\(C\in\R^{q\times p}\) and \(D\in\R^{q\times q}\).
\begin{enumerate}
\item If \(A\) and its Schur complement \(S=D-CA^{-1}B\) are invertible, then \(M\) is invertible, the block
row operations \(R_2\leftarrow R_2-CA^{-1}R_1\), then \(R_1\leftarrow A^{-1}R_1\) and \(R_2\leftarrow S^{-1}R_2\),
then \(R_1\leftarrow R_1-A^{-1}BR_2\) take \([M\mid I_{p+q}]\) to \([I_{p+q}\mid M^{-1}]\),
\[
M^{-1}=\begin{pmatrix}I_p&-A^{-1}B\\0&I_q\end{pmatrix}\begin{pmatrix}A^{-1}&0\\0&S^{-1}\end{pmatrix}
\begin{pmatrix}I_p&0\\-CA^{-1}&I_q\end{pmatrix}
\]
and \(M=\left(\begin{smallmatrix}I_p&0\\CA^{-1}&I_q\end{smallmatrix}\right)\diag(A,S)
\left(\begin{smallmatrix}I_p&A^{-1}B\\0&I_q\end{smallmatrix}\right)\).
\item If \(D\) and its Schur complement \(F=A-BD^{-1}C\) are invertible, then \(M\) is invertible, the block
row operations \(R_1\leftarrow R_1-BD^{-1}R_2\), then \(R_1\leftarrow F^{-1}R_1\) and \(R_2\leftarrow D^{-1}R_2\),
then \(R_2\leftarrow R_2-D^{-1}CR_1\) take \([M\mid I_{p+q}]\) to \([I_{p+q}\mid M^{-1}]\),
\[
M^{-1}=\begin{pmatrix}I_p&0\\-D^{-1}C&I_q\end{pmatrix}\begin{pmatrix}F^{-1}&0\\0&D^{-1}\end{pmatrix}
\begin{pmatrix}I_p&-BD^{-1}\\0&I_q\end{pmatrix}
\]
and \(M=\left(\begin{smallmatrix}I_p&BD^{-1}\\0&I_q\end{smallmatrix}\right)\diag(F,D)
\left(\begin{smallmatrix}I_p&0\\D^{-1}C&I_q\end{smallmatrix}\right)\).
\end{enumerate}
\end{corollary}
For \(p\ne q\) the blocks \(B\) and \(C\) are rectangular.
\begin{proof}
(a)
\begin{itemize}
\item The three steps are the left multiplications by
\(E_1=\left(\begin{smallmatrix}I_p&0\\-CA^{-1}&I_q\end{smallmatrix}\right)\),
\(E_2=\diag(I_p,S^{-1})\diag(A^{-1},I_q)=\diag(A^{-1},S^{-1})\) (Proposition~\ref{r5b-prop-block-tri}(b))
and \(E_3=\left(\begin{smallmatrix}I_p&-A^{-1}B\\0&I_q\end{smallmatrix}\right)\).
\item By Theorem~\ref{r5b-thm-schur} (with \(n=q\)),
\(E_1M=\left(\begin{smallmatrix}A&B\\0&S\end{smallmatrix}\right)\). By the block product formula
(Theorem~\ref{r5b-thm-block}(c)),
\(E_2E_1M=\left(\begin{smallmatrix}A^{-1}A&A^{-1}B\\S^{-1}\cdot0&S^{-1}S\end{smallmatrix}\right)
=\left(\begin{smallmatrix}I_p&A^{-1}B\\0&I_q\end{smallmatrix}\right)\). By
Proposition~\ref{r5b-prop-block-elementary} (second identity, with upper right block \(-A^{-1}B\)),
\(E_3E_2E_1M=\left(\begin{smallmatrix}I_p&A^{-1}B-A^{-1}B\\0&I_q\end{smallmatrix}\right)=I_{p+q}\).
\item Put \(E=E_3E_2E_1\), the displayed product. Then \(EM=I_{p+q}\), so \(M\) is invertible with
\(M^{-1}=E\) (Theorem~\ref{r2-thm-invertible}), and
\(E[M\mid I_{p+q}]=[EM\mid E]=[I_{p+q}\mid M^{-1}]\) (Theorem~\ref{r1-thm-views}(a)).
\item \(E_1^{-1}=\left(\begin{smallmatrix}I_p&0\\CA^{-1}&I_q\end{smallmatrix}\right)\) and
\(E_3^{-1}=\left(\begin{smallmatrix}I_p&A^{-1}B\\0&I_q\end{smallmatrix}\right)\)
(Proposition~\ref{r5b-prop-block-elementary}), and \(E_2^{-1}=\diag(A,S)\)
(Proposition~\ref{r5b-prop-block-tri}(b)). So \(M=E^{-1}=E_1^{-1}E_2^{-1}E_3^{-1}\)
(Theorem~\ref{r2-thm-inverse-props}(a),(b)).
\end{itemize}
(b) The three steps are the left multiplications by
\(E_1'=\left(\begin{smallmatrix}I_p&-BD^{-1}\\0&I_q\end{smallmatrix}\right)\),
\(E_2'=\diag(F^{-1},D^{-1})\) (as in (a)) and \(E_3'=\left(\begin{smallmatrix}I_p&0\\-D^{-1}C&I_q\end{smallmatrix}\right)\).
By Proposition~\ref{r5b-prop-block-elementary} and Theorem~\ref{r5b-thm-block}(c),
\[
E_1'M=\begin{pmatrix}A-BD^{-1}C&B-BD^{-1}D\\C&D\end{pmatrix}=\begin{pmatrix}F&0\\C&D\end{pmatrix},\qquad
E_2'E_1'M=\begin{pmatrix}I_p&0\\D^{-1}C&I_q\end{pmatrix}
\]
and \(E_3'E_2'E_1'M=I_{p+q}\). The last two bullets of (a) apply to \(E'=E_3'E_2'E_1'\), the displayed
product, with \(E_1'^{-1}=\left(\begin{smallmatrix}I_p&BD^{-1}\\0&I_q\end{smallmatrix}\right)\),
\(E_2'^{-1}=\diag(F,D)\) and \(E_3'^{-1}=\left(\begin{smallmatrix}I_p&0\\D^{-1}C&I_q\end{smallmatrix}\right)\)
(Propositions~\ref{r5b-prop-block-elementary} and~\ref{r5b-prop-block-tri}(b)).
\end{proof}

\begin{example}\label{r5b-exm-schur}
Partition \(M\in\R^{5\times5}\) with groups \((2,3)\) for rows and columns. Find the Schur complement \(S\) of
\(A\), the factorization of \(M\) in Corollary~\ref{r5b-cor-schur-inverse}(a), and \(M^{-1}\).
\[
M=\left(\begin{array}{cc|ccc}2&1&2&1&3\\1&1&1&1&2\\\hline2&1&3&2&3\\1&1&2&3&3\\3&2&3&3&7\end{array}\right)
=\begin{pmatrix}A&B\\C&D\end{pmatrix}
\]
\[
A=\begin{pmatrix}2&1\\1&1\end{pmatrix},\quad B=\begin{pmatrix}2&1&3\\1&1&2\end{pmatrix},\quad C=B',\quad
D=\begin{pmatrix}3&2&3\\2&3&3\\3&3&7\end{pmatrix}
\]
Here \(A\) is \(2\times2\), \(D\) is \(3\times3\), \(B\) is \(2\times3\) and \(C\) is \(3\times2\).
By Proposition~\ref{r2-prop-inverse-2x2} (\(2\cdot1-1\cdot1=1\)),
\(A^{-1}=\left(\begin{smallmatrix}1&-1\\-1&2\end{smallmatrix}\right)\). Then
\[
A^{-1}B=\begin{pmatrix}1&-1\\-1&2\end{pmatrix}\begin{pmatrix}2&1&3\\1&1&2\end{pmatrix}=\begin{pmatrix}1&0&1\\0&1&1\end{pmatrix},\qquad
CA^{-1}=\begin{pmatrix}2&1\\1&1\\3&2\end{pmatrix}\begin{pmatrix}1&-1\\-1&2\end{pmatrix}=\begin{pmatrix}1&0\\0&1\\1&1\end{pmatrix}
\]
These are \(2\times3\) and \(3\times2\), and the Schur complement of \(A\) is the \(3\times3\) matrix
\[
S=D-CA^{-1}B=\begin{pmatrix}3&2&3\\2&3&3\\3&3&7\end{pmatrix}-\begin{pmatrix}1&0\\0&1\\1&1\end{pmatrix}\begin{pmatrix}2&1&3\\1&1&2\end{pmatrix}
=\begin{pmatrix}3&2&3\\2&3&3\\3&3&7\end{pmatrix}-\begin{pmatrix}2&1&3\\1&1&2\\3&2&5\end{pmatrix}
=\begin{pmatrix}1&1&0\\1&2&1\\0&1&2\end{pmatrix}
\]
By Theorem~\ref{r5b-thm-schur},
\[
\left(\begin{array}{cc|ccc}1&0&0&0&0\\0&1&0&0&0\\\hline-1&0&1&0&0\\0&-1&0&1&0\\-1&-1&0&0&1\end{array}\right)M
\left(\begin{array}{cc|ccc}1&0&-1&0&-1\\0&1&0&-1&-1\\\hline0&0&1&0&0\\0&0&0&1&0\\0&0&0&0&1\end{array}\right)
=\left(\begin{array}{cc|ccc}2&1&0&0&0\\1&1&0&0&0\\\hline0&0&1&1&0\\0&0&1&2&1\\0&0&0&1&2\end{array}\right)
\]
All matrices here have integer entries, because \(A^{-1}\) does. Multiplying on the left and on the
right by the inverses of the two outer factors (Proposition~\ref{r5b-prop-block-elementary}),
\[
\begin{aligned}
M&=\begin{pmatrix}I_2&0\\CA^{-1}&I_3\end{pmatrix}\begin{pmatrix}A&0\\0&S\end{pmatrix}
\begin{pmatrix}I_2&A^{-1}B\\0&I_3\end{pmatrix}\\
&=\left(\begin{array}{cc|ccc}1&0&0&0&0\\0&1&0&0&0\\\hline1&0&1&0&0\\0&1&0&1&0\\1&1&0&0&1\end{array}\right)
\left(\begin{array}{cc|ccc}2&1&0&0&0\\1&1&0&0&0\\\hline0&0&1&1&0\\0&0&1&2&1\\0&0&0&1&2\end{array}\right)
\left(\begin{array}{cc|ccc}1&0&1&0&1\\0&1&0&1&1\\\hline0&0&1&0&0\\0&0&0&1&0\\0&0&0&0&1\end{array}\right)
\end{aligned}
\]
so \(M\) is the product of a lower block triangular, a block diagonal and an upper block triangular
factor, whose off-diagonal blocks \(CA^{-1}\in\R^{3\times2}\) and \(A^{-1}B\in\R^{2\times3}\) are
rectangular. Since \(A'=A\) and \(C=B'\), \(A^{-1}B=(A^{-1})'C'=(CA^{-1})'\)
(Theorems~\ref{r2-thm-inverse-props}(c) and~\ref{r1-thm-transpose}), and the right factor is the
transpose of the left (Theorem~\ref{r5b-thm-block}(b)).

Finally, \(M^{-1}\) is found by Corollary~\ref{r5b-cor-schur-inverse}(a). The Schur complement \(S\) is invertible, with
\[
S^{-1}=\begin{pmatrix}3&-2&1\\-2&2&-1\\1&-1&1\end{pmatrix}
\]
since \(SS^{-1}=I_3\) (Theorem~\ref{r2-thm-invertible}). By the corollary and the block product formula
(Theorem~\ref{r5b-thm-block}(c)),
\[
M^{-1}=\begin{pmatrix}I_2&-A^{-1}B\\0&I_3\end{pmatrix}\begin{pmatrix}A^{-1}&0\\0&S^{-1}\end{pmatrix}
\begin{pmatrix}I_2&0\\-CA^{-1}&I_3\end{pmatrix}
=\begin{pmatrix}A^{-1}+A^{-1}BS^{-1}CA^{-1}&-A^{-1}BS^{-1}\\-S^{-1}CA^{-1}&S^{-1}\end{pmatrix}
\]
With \(A^{-1}B\) and \(CA^{-1}\) from above,
\[
\begin{gathered}
-A^{-1}BS^{-1}=-\begin{pmatrix}1&0&1\\0&1&1\end{pmatrix}\begin{pmatrix}3&-2&1\\-2&2&-1\\1&-1&1\end{pmatrix}
=\begin{pmatrix}-4&3&-2\\1&-1&0\end{pmatrix}\\
-S^{-1}CA^{-1}=-\begin{pmatrix}3&-2&1\\-2&2&-1\\1&-1&1\end{pmatrix}\begin{pmatrix}1&0\\0&1\\1&1\end{pmatrix}
=\begin{pmatrix}-4&1\\3&-1\\-2&0\end{pmatrix}
\end{gathered}
\]
and
\[
\begin{aligned}
A^{-1}+A^{-1}BS^{-1}CA^{-1}&=A^{-1}-(A^{-1}B)(-S^{-1}CA^{-1})
=\begin{pmatrix}1&-1\\-1&2\end{pmatrix}-\begin{pmatrix}1&0&1\\0&1&1\end{pmatrix}\begin{pmatrix}-4&1\\3&-1\\-2&0\end{pmatrix}\\
&=\begin{pmatrix}1&-1\\-1&2\end{pmatrix}-\begin{pmatrix}-6&1\\1&-1\end{pmatrix}=\begin{pmatrix}7&-2\\-2&3\end{pmatrix}
\end{aligned}
\]
So
\[
M^{-1}=\left(\begin{array}{cc|ccc}7&-2&-4&3&-2\\-2&3&1&-1&0\\\hline-4&1&3&-2&1\\3&-1&-2&2&-1\\-2&0&1&-1&1\end{array}\right)
\]
Check: \(\row_1(M)M^{-1}=2(7,-2,-4,3,-2)+(-2,3,1,-1,0)+2(-4,1,3,-2,1)+(3,-1,-2,2,-1)+3(-2,0,1,-1,1)
=(1,0,0,0,0)\) (Theorem~\ref{r1-thm-views}(c)).
\end{example}

\begin{theorem}[\citeat{abadir-magnus}{Exercise~5.16}, \citeat{harville}{Theorem~8.5.11, Corollary~8.5.12}, \citeat{searle-khuri}{\S4.12(h)}]\label{r5b-thm-partitioned-inverse}
Let \(M=\begin{pmatrix}A&B\\C&D\end{pmatrix}\) with \(A\in\R^{p\times p}\), \(B\in\R^{p\times q}\),
\(C\in\R^{q\times p}\) and \(D\in\R^{q\times q}\).
\begin{enumerate}
\item If \(A\) and its Schur complement \(E=D-CA^{-1}B\) are invertible, then \(M\) is invertible and
\[
M^{-1}=\begin{pmatrix}A^{-1}+A^{-1}BE^{-1}CA^{-1}&-A^{-1}BE^{-1}\\-E^{-1}CA^{-1}&E^{-1}\end{pmatrix}
\]
\item If \(D\) and its Schur complement \(F=A-BD^{-1}C\) are invertible, then \(M\) is invertible and
\[
M^{-1}=\begin{pmatrix}F^{-1}&-F^{-1}BD^{-1}\\-D^{-1}CF^{-1}&D^{-1}+D^{-1}CF^{-1}BD^{-1}\end{pmatrix}
\]
\end{enumerate}
In particular, in (a) the lower right block of \(M^{-1}\) is the inverse of the Schur complement of \(A\).
\end{theorem}
\begin{proof}
(a)
\begin{itemize}
\item By Theorem~\ref{r5b-thm-schur}, \(LM=U\) with
\(L=\left(\begin{smallmatrix}I_p&0\\-CA^{-1}&I_q\end{smallmatrix}\right)\) and
\(U=\left(\begin{smallmatrix}A&B\\0&E\end{smallmatrix}\right)\).
\item \(L\) is invertible with inverse \(\left(\begin{smallmatrix}I_p&0\\CA^{-1}&I_q\end{smallmatrix}\right)\)
(Proposition~\ref{r5b-prop-block-elementary}), so \(M=L^{-1}U\); and \(U\) is invertible with
\(U^{-1}=\left(\begin{smallmatrix}A^{-1}&-A^{-1}BE^{-1}\\0&E^{-1}\end{smallmatrix}\right)\)
(Proposition~\ref{r5b-prop-block-tri}(c)).
\item Hence \(M\) is invertible and \(M^{-1}=U^{-1}(L^{-1})^{-1}=U^{-1}L\)
(Theorem~\ref{r2-thm-inverse-props}(a),(b)). By the block product formula
(Theorem~\ref{r5b-thm-block}(c)),
\[
U^{-1}L=\begin{pmatrix}A^{-1}&-A^{-1}BE^{-1}\\0&E^{-1}\end{pmatrix}
\begin{pmatrix}I_p&0\\-CA^{-1}&I_q\end{pmatrix}
=\begin{pmatrix}A^{-1}+A^{-1}BE^{-1}CA^{-1}&-A^{-1}BE^{-1}\\-E^{-1}CA^{-1}&E^{-1}\end{pmatrix}
\]
\end{itemize}

(b) Look for \(X=\left(\begin{smallmatrix}P&Q\\R&T\end{smallmatrix}\right)\) with
\(MX=I_{p+q}\); by the block product formula this means
\[
AP+BR=I_p,\qquad CP+DR=0,\qquad AQ+BT=0,\qquad CQ+DT=I_q
\]
\begin{itemize}
\item The second equation gives \(R=-D^{-1}CP\); then the first becomes
\((A-BD^{-1}C)P=FP=I_p\), so \(P=F^{-1}\) and \(R=-D^{-1}CF^{-1}\).
\item The fourth equation gives \(T=D^{-1}-D^{-1}CQ\); then the third becomes
\(AQ+BD^{-1}-BD^{-1}CQ=FQ+BD^{-1}=0\), so \(Q=-F^{-1}BD^{-1}\) and
\(T=D^{-1}+D^{-1}CF^{-1}BD^{-1}\).
\item Conversely, these blocks satisfy the four equations:
\(AP+BR=(A-BD^{-1}C)F^{-1}=I_p\), \(CP+DR=CF^{-1}-CF^{-1}=0\),
\(AQ+BT=-(A-BD^{-1}C)F^{-1}BD^{-1}+BD^{-1}=0\) and
\(CQ+DT=-CF^{-1}BD^{-1}+I_q+CF^{-1}BD^{-1}=I_q\).
\item So \(MX=I_{p+q}\), and since \(M\) is square, \(M\) is invertible with \(M^{-1}=X\)
(Theorem~\ref{r2-thm-invertible}).\qedhere
\end{itemize}
\end{proof}

\subsection*{Determinants: definition and computation}

\begin{definition}[\citeat{hefferon}{Four.I.2, Definition~2.1}, \citeat{hoffman-kunze}{\S5.2}]\label{r5d-def-det}
Let \(n\geqslant1\). A function \(D\) that assigns a number \(D(A)\) to each \(A\in\R^{n\times n}\) is a
\textbf{determinant function} if, for every \(A\in\R^{n\times n}\), all row indices \(r\ne s\) and every
\(c\in\R\):
\begin{enumerate}
\item multiplying row \(r\) of \(A\) by \(c\) multiplies \(D(A)\) by \(c\);
\item adding \(c\) times row \(s\) to row \(r\) does not change \(D(A)\);
\item interchanging rows \(r\) and \(s\) changes the sign of \(D(A)\);
\item \(D(I_n)=1\).
\end{enumerate}
We show that for each \(n\) there is exactly one determinant function
(Theorems~\ref{r5d-thm-gauss} and~\ref{r5d-thm-expansion}). It is written \(\det\), and \(\det A\) is the
\textbf{determinant} of \(A\); many books write \(\lvert A\rvert\), but in these notes \(\lvert\cdot\rvert\) always
means absolute value. Until its existence is proved, every statement
about \(\det\) holds for every determinant function. For \(n=1\), (a) and (d) give \(\det(a)=a\det(1)=a\).
\end{definition}

For \(n=2\), conditions (a)--(d) are the rules for the area of the parallelogram whose sides are the two rows
\(u'\), \(v'\) of \(A\) (\(u,v\in\R^2\)), counted with a sign: stretching the side \(u'\) by \(c\) multiplies the area
by \(c\); replacing \(u'\) by \(u'+cv'\) slides that side along \(v'\) and keeps both the base \(v'\) and the height;
and the unit square, with sides \(e_1'\) and \(e_2'\), has area \(1\). The sign records the orientation of the pair
\(u',v'\), which an interchange reverses. For \(n\geqslant3\), \(\lvert\det A\rvert\) is accordingly taken as the volume of
the box spanned by the rows~\citeat{hefferon}{Four.II.1}. This is why determinants appear in the change-of-variables formula for probability
densities; we do not use this here.

\begin{theorem}[\citeat{hefferon}{Four.I.2, Lemma~2.4; Four.II.1, proof of Theorem~1.5}, \citeat{searle-khuri}{\S3.4}]\label{r5d-thm-basic}
Let \(A\in\R^{n\times n}\).
\begin{enumerate}
\item If two rows of \(A\) are equal, then \(\det A=0\).
\item If \(A\) has a zero row, then \(\det A=0\).
\item If \(A\) is upper triangular, then \(\det A=a_{11}a_{22}\cdots a_{nn}\). In particular
\(\det\diag(d_1,\dots,d_n)=d_1d_2\cdots d_n\).
\item For the elementary matrices of Definition~\ref{r2-def-elementary}, \(\det D_r(c)=c\),
\(\det T_{rs}(c)=1\) and \(\det P_{rs}=-1\); and \(\det(EA)=\det E\cdot\det A\) for every elementary matrix
\(E\in\R^{n\times n}\).
\end{enumerate}
\end{theorem}
\begin{proof}
(a) Interchanging the two equal rows does not change \(A\), but it changes the sign of \(\det A\)
(Definition~\ref{r5d-def-det}(c)). So \(\det A=-\det A\), that is, \(\det A=0\).

(b) Multiplying the zero row by \(2\) does not change \(A\), but it doubles \(\det A\)
(Definition~\ref{r5d-def-det}(a)). So \(\det A=2\det A\), that is, \(\det A=0\).

(c)
\begin{itemize}
\item Let \(a_{11},\dots,a_{nn}\) be nonzero. Multiplying row \(i\) by \(1/a_{ii}\) for \(i=1,\dots,n\) gives an upper
triangular matrix \(B\) with diagonal entries \(1\), and \(\det B=\det A/(a_{11}\cdots a_{nn})\)
(Definition~\ref{r5d-def-det}(a)).
\item Backward elimination (Step 5 of Definition~\ref{r2-def-gauss-jordan}, with the diagonal entries of \(B\) as
pivots) turns \(B\) into \(I_n\) by adding multiples of rows to other rows; so \(\det B=\det I_n=1\)
(Definition~\ref{r5d-def-det}(b),(d)), and \(\det A=a_{11}\cdots a_{nn}\).
\item Now let some \(a_{ii}\) be \(0\), with \(i\) the largest such index. Row \(i\) of \(A\) is \(0\) in columns
\(1,\dots,i\). For \(j=i+1,\dots,n\) in turn, subtract a multiple of row \(j\) from row \(i\) so that entry
\((i,j)\) becomes \(0\) (possible as \(a_{jj}\ne0\)); row \(j\) is \(0\) before column \(j\), so the zeros already
made stay \(0\). These additions do not change \(\det A\) and turn row \(i\) into a zero row; so \(\det A=0\) by (b),
and also \(a_{11}\cdots a_{nn}=0\).
\end{itemize}

(d)
\begin{itemize}
\item \(D_r(c)\), \(T_{rs}(c)\) and \(P_{rs}\) are obtained from \(I_n\) by \(R_r\leftarrow cR_r\),
\(R_r\leftarrow R_r+cR_s\) and \(R_r\leftrightarrow R_s\), so by Definition~\ref{r5d-def-det} their determinants
are \(c\), \(1\) and \(-1\).
\item \(EA\) is obtained from \(A\) by the row operation of \(E\) (Theorem~\ref{r2-thm-elementary}(a)), so
\(\det(EA)\) is \(c\det A\), \(\det A\) or \(-\det A\), that is, \(\det E\cdot\det A\).\qedhere
\end{itemize}
\end{proof}

\begin{theorem}[\citeat{hefferon}{Four.I.2, Lemma~2.7, Examples~2.5, 2.6}, \citeat{searle-khuri}{\S3.3.5}]\label{r5d-thm-gauss}
Let interchanges of two rows and additions of a multiple of one row to another take \(A\in\R^{n\times n}\) to an
upper triangular matrix \(U\), with \(t\) interchanges. Then
\[
\det A=(-1)^tu_{11}u_{22}\cdots u_{nn}
\]
Forward elimination (Steps 1, 2 and 4 of Definition~\ref{r2-def-gauss-jordan}, without scaling the leading
entries to \(1\)) always ends with such a matrix \(U\). In particular, for each \(n\) there is at most one
determinant function.
\end{theorem}
\begin{proof}
\begin{itemize}
\item \(\det U=u_{11}\cdots u_{nn}\) (Theorem~\ref{r5d-thm-basic}(c)). Each addition leaves the determinant
unchanged and each interchange changes its sign (Definition~\ref{r5d-def-det}(b),(c)); so
\(\det U=(-1)^t\det A\), and the formula follows.
\item In the result of forward elimination, the leading entries of the nonzero rows lie in columns
\(k_1<k_2<\cdots\), so \(k_i\geqslant i\), and the zero rows lie at the bottom
(Definition~\ref{r2-def-gauss-jordan}); hence the result is upper triangular.
\item For each \(A\), fix one forward elimination of \(A\) (Step 2 may allow several; any one will do). By the
first item, every determinant function takes the value \((-1)^tu_{11}\cdots u_{nn}\) of this reduction at \(A\); so
two determinant functions are equal.\qedhere
\end{itemize}
\end{proof}
If a row is multiplied by \(c\ne0\) on the way, the determinant is multiplied by \(c\)
(Definition~\ref{r5d-def-det}(a)); so divide by \(c\) at the end.

\begin{example}\label{r5d-exm-rowred-swap}
Compute \(\det A\) for
\[
A=\begin{pmatrix}0&2&1&1\\1&2&1&0\\2&6&4&4\\1&0&2&3\end{pmatrix}
\]

\textit{Solution.} Since \(a_{11}=0\), begin with an interchange:
\[
\begin{gathered}
A\xrightarrow{R_1\leftrightarrow R_2}
\begin{pmatrix}1&2&1&0\\0&2&1&1\\2&6&4&4\\1&0&2&3\end{pmatrix}
\xrightarrow[R_4\leftarrow R_4-R_1]{R_3\leftarrow R_3-2R_1}
\begin{pmatrix}1&2&1&0\\0&2&1&1\\0&2&2&4\\0&-2&1&3\end{pmatrix}
\\
\xrightarrow[R_4\leftarrow R_4+R_2]{R_3\leftarrow R_3-R_2}
\begin{pmatrix}1&2&1&0\\0&2&1&1\\0&0&1&3\\0&0&2&4\end{pmatrix}
\xrightarrow{R_4\leftarrow R_4-2R_3}
\begin{pmatrix}1&2&1&0\\0&2&1&1\\0&0&1&3\\0&0&0&-2\end{pmatrix}=U
\end{gathered}
\]
The interchange changes the sign and the additions change nothing (Theorem~\ref{r5d-thm-gauss}, \(t=1\)), so
\(\det A=-\det U=-\bigl(1\cdot2\cdot1\cdot(-2)\bigr)=4\).
\end{example}

\begin{example}[\citeat{hoffman-kunze}{\S5.4, Example~6}]\label{r5d-exm-rowred4}
Compute \(\det A\) for
\[
A=\begin{pmatrix}1&-1&2&3\\2&2&0&2\\4&1&-1&-1\\1&2&3&0\end{pmatrix}
\]

\textit{Solution.}
\[
\begin{gathered}
A\xrightarrow[R_4\leftarrow R_4-R_1]{R_2\leftarrow R_2-2R_1,\ R_3\leftarrow R_3-4R_1}
\begin{pmatrix}1&-1&2&3\\0&4&-4&-4\\0&5&-9&-13\\0&3&1&-3\end{pmatrix}
\xrightarrow{R_2\leftarrow\frac14R_2}
\begin{pmatrix}1&-1&2&3\\0&1&-1&-1\\0&5&-9&-13\\0&3&1&-3\end{pmatrix}
\\
\xrightarrow[R_4\leftarrow R_4-3R_2]{R_3\leftarrow R_3-5R_2}
\begin{pmatrix}1&-1&2&3\\0&1&-1&-1\\0&0&-4&-8\\0&0&4&0\end{pmatrix}
\xrightarrow{R_4\leftarrow R_4+R_3}
\begin{pmatrix}1&-1&2&3\\0&1&-1&-1\\0&0&-4&-8\\0&0&0&-8\end{pmatrix}=U
\end{gathered}
\]
Only the second arrow changes the determinant: it multiplies it by \(\frac14\)
(Definition~\ref{r5d-def-det}(a)). Hence
\(\det A=4\det U=4\cdot1\cdot1\cdot(-4)\cdot(-8)=128\).
\end{example}

\subsection*{Main properties}

\begin{theorem}[\citeat{harville}{Theorems~13.3.4, 13.3.7, (13.3.10), (13.3.12)}, \citeat{searle-khuri}{Theorem~3.4, Corollary~3.3}]\label{r5d-thm-product}
Let \(A,B\in\R^{n\times n}\).
\begin{enumerate}
\item \(A\) is invertible \(\Leftrightarrow\) \(\det A\ne0\) \(\Leftrightarrow\) \(\rank A=n\).
\item \(\det(AB)=\det A\cdot\det B\).
\item If \(A\) is invertible, then \(\det A^{-1}=1/\det A\). Moreover \(\det(cA)=c^n\det A\) for
\(c\in\R\), and \(\det A^k=(\det A)^k\) for \(k=1,2,\dots\)
\end{enumerate}
\end{theorem}
\begin{proof}
For \(n=1\), \(\det(a)=a\): (a) holds since \((a)\) is invertible, and \(\rank(a)=1\), if and only if \(a\ne0\)
(Definition~\ref{def:r4-rank}), and (b), (c) are identities between numbers; let \(n\geqslant2\).

\needspace{4\baselineskip}
(a)
\begin{itemize}
\item Let \(R=PA\) be the RREF of \(A\), where \(P=E_k\cdots E_1\) is a product of elementary matrices
(Theorem~\ref{r2-thm-elementary}(c)). By Theorem~\ref{r5d-thm-basic}(d), applied \(k\) times,
\(\det R=\det E_k\cdots\det E_1\cdot\det A\), and each \(\det E_i\) is \(c\ne0\), \(1\) or \(-1\).
\item If \(A\) is invertible, then \(R=I_n\) (Theorem~\ref{r2-thm-invertible}), so \(\det A\ne0\). If
not, then \(\rank A<n\) (Theorem~\ref{thm:r4-rank-rules}(d)), so \(R\) has fewer than \(n\) nonzero
rows (Definition~\ref{def:r4-rank}), that is, a zero row; then \(\det R=0\)
(Theorem~\ref{r5d-thm-basic}(b)) and \(\det A=0\). Finally, \(A\) is invertible if and only if
\(\rank A=n\) (Theorem~\ref{thm:r4-rank-rules}(d)).
\end{itemize}

(b)
\begin{itemize}
\item If \(A\) is not invertible, neither is \(AB\): an inverse \(G\) of \(AB\) would give
\(A(BG)=I_n\), and \(A\) would be invertible (Theorem~\ref{r2-thm-invertible}). By (a),
\(\det(AB)=0=\det A\cdot\det B\).
\item If \(A\) is invertible, then \(A=F_1\cdots F_k\) with elementary matrices \(F_i\)
(Theorem~\ref{r2-thm-invertible}). By Theorem~\ref{r5d-thm-basic}(d), applied \(k\) times,
\(\det(AB)=\det F_1\cdots\det F_k\cdot\det B\); for \(B=I_n\) this reads
\(\det A=\det F_1\cdots\det F_k\), so \(\det(AB)=\det A\cdot\det B\).
\end{itemize}

(c) By (b), \(\det A\cdot\det A^{-1}=\det(AA^{-1})=\det I_n=1\). Next, \(cA=(cI_n)A\) and \(\det(cI_n)=c^n\)
(Theorem~\ref{r5d-thm-basic}(c)), so \(\det(cA)=c^n\det A\) by (b). Finally
\(\det A^{k+1}=\det(A^kA)=\det A^k\cdot\det A\) by (b), which gives the powers by induction on \(k\).
\end{proof}

\subsection*{Cofactor expansion and the transpose}

\begin{lemma}[\citeat{hefferon}{Four.I.3, Lemma~3.3}]\label{r5d-lem-multilinear}
Let \(A_1,A_2,A_3\in\R^{n\times n}\) agree in all rows except row \(i\), and let \(\row_i(A_3)=\row_i(A_1)+\row_i(A_2)\).
Then \(\det A_3=\det A_1+\det A_2\). Together with Definition~\ref{r5d-def-det}(a): \(\det\) is linear in each row when
the other rows are held fixed.
\end{lemma}
\begin{proof}
Let \(\rho_l'\) (\(l\ne i\), \(\rho_l\in\R^n\)) be the common rows, and let \(u'=\row_i(A_1)\) and \(v'=\row_i(A_2)\)
(\(u,v\in\R^n\)).
\begin{itemize}
\item Let the \(n-1\) vectors \(\rho_l\) be linearly dependent. The columns of \(A_1'\) are \(u\) and the \(\rho_l\)
(Definition~\ref{r1-def-transpose}), so they are dependent too (Proposition~\ref{r3-prop-dependence}(b)). So \(A_1'\) is
not invertible (Theorems~\ref{thm:r4-pivot-test} and~\ref{r2-thm-invertible}), neither is \(A_1\)
(Theorem~\ref{r2-thm-inverse-props}(c)), and \(\det A_1=0\) (Theorem~\ref{r5d-thm-product}(a)). Likewise
\(\det A_2=\det A_3=0\).
\item Otherwise extend the \(n-1\) vectors \(\rho_l\) by one vector \(\beta\) to a basis of \(\R^n\)
(Theorem~\ref{thm:r4-counting}(b); a basis of \(\R^n\) has \(n\) vectors, Theorem~\ref{thm:r4-dimension}). Write
\(u=\sum_{l\ne i}x_l\rho_l+x_i\beta\) and \(v=\sum_{l\ne i}y_l\rho_l+y_i\beta\); then
\(u+v=\sum_{l\ne i}(x_l+y_l)\rho_l+(x_i+y_i)\beta\).
\item Let \(G\) have row \(i\) equal to \(\beta'\) and the rows \(\rho_l'\) elsewhere. Subtracting \(x_l\) times row \(l\)
from row \(i\) for every \(l\ne i\) turns \(A_1\) into \(G\) with row \(i\) multiplied by \(x_i\); so
\(\det A_1=x_i\det G\) (Definition~\ref{r5d-def-det}(b),(a)). Likewise \(\det A_2=y_i\det G\) and
\(\det A_3=(x_i+y_i)\det G\). Hence \(\det A_3=\det A_1+\det A_2\).\qedhere
\end{itemize}
\end{proof}

\begin{definition}[\citeat{hoffman-kunze}{\S\S5.2, 5.4}, \citeat{hefferon}{Four.III.1, Definition~1.2}, \citeat{searle-khuri}{\S3.1.2}, \citeat{lang-ila}{Ch.~VII, \S2}]\label{r5d-def-cofactor}
Let \(n\geqslant2\) and \(A\in\R^{n\times n}\). For \(1\leqslant i,j\leqslant n\), \(A(i|j)\in\R^{(n-1)\times(n-1)}\) is the
matrix obtained from \(A\) by deleting row \(i\) and column \(j\). The number \(\det A(i|j)\) is the \textbf{minor}
of \(a_{ij}\), and \(c_{ij}=(-1)^{i+j}\det A(i|j)\) is its \textbf{cofactor}.
\end{definition}
The signs \((-1)^{i+j}\) form a checkerboard with \(+\) on the diagonal, for \(n=3\)
\(\left(\begin{smallmatrix}+&-&+\\-&+&-\\+&-&+\end{smallmatrix}\right)\).

\begin{theorem}[\citeat{hoffman-kunze}{\S5.2, Theorem~1; \S5.4}, \citeat{lang}{Ch.~VI, Theorem~2.1}, \citeat{hefferon}{Four.III.1, Theorem~1.5}, \citeat{searle-khuri}{(3.1)--(3.2)}]\label{r5d-thm-expansion}
For each \(n\geqslant1\) there is exactly one determinant function on \(\R^{n\times n}\). Let \(n\geqslant2\) and
\(A\in\R^{n\times n}\), with cofactors \(c_{ij}\). For every column \(j\) and every row \(i\),
\[
\det A=a_{1j}c_{1j}+a_{2j}c_{2j}+\dots+a_{nj}c_{nj}=a_{i1}c_{i1}+a_{i2}c_{i2}+\dots+a_{in}c_{in}
\]
(\textbf{expansion along column \(j\)} and \textbf{along row \(i\)}).
\end{theorem}
\begin{proof}
\textit{Existence and the expansion along a column}, by induction on \(n\).
\begin{itemize}
\item \(n=1\): the function \((a)\mapsto a\) satisfies Definition~\ref{r5d-def-det}. Multiplying the only row by \(c\)
multiplies the value by \(c\); conditions (b) and (c) concern two different rows and say nothing for \(n=1\); and
\((1)\mapsto1\). By Theorem~\ref{r5d-thm-gauss} it is the only determinant function.
\item Let \(n\geqslant2\), and let \(\det\) be the determinant function on \(\R^{(n-1)\times(n-1)}\), which exists by the
induction hypothesis. Fix a column \(j\) and let \(E(A)=\sum_{i=1}^na_{ij}c_{ij}=\sum_{i=1}^n(-1)^{i+j}a_{ij}\det A(i|j)\) for
\(A\in\R^{n\times n}\). We show that \(E\) is a determinant function on \(\R^{n\times n}\). By Theorem~\ref{r5d-thm-gauss} it
is then the only one, so \(E=\det\): a determinant function exists for \(n\), and the expansion along column \(j\) holds.
\item \(E\) is linear in each row \(k\). For \(i\ne k\), \(a_{ij}\) does not depend on row \(k\), and \(\det A(i|j)\) is
linear in row \(k\) (Lemma~\ref{r5d-lem-multilinear}; row \(k\) of \(A\) without its entry in column \(j\) is a row
of \(A(i|j)\)). For \(i=k\), \(\det A(k|j)\) does not depend on row \(k\), and \(a_{kj}\) is linear in it. In
particular \(E\) satisfies Definition~\ref{r5d-def-det}(a).
\item If two adjacent rows \(k\), \(k+1\) of \(A\) are equal, then \(E(A)=0\). For \(i\ne k,k+1\), \(A(i|j)\) has two
equal rows, so \(\det A(i|j)=0\) (Theorem~\ref{r5d-thm-basic}(a)). The two remaining terms
\((-1)^{k+j}a_{kj}\det A(k|j)\) and \((-1)^{k+1+j}a_{k+1,j}\det A(k+1|j)\) cancel, as \(a_{kj}=a_{k+1,j}\) and
\(A(k|j)=A(k+1|j)\).
\item An interchange of two adjacent rows changes the sign of \(E\). Write \(E(u,v)\) for the value of \(E\) when
rows \(k\), \(k+1\) are \(u'\), \(v'\) (\(u,v\in\R^n\)) and the other rows are fixed. By the first two items,
\(0=E(u+v,u+v)=E(u,u)+E(u,v)+E(v,u)+E(v,v)=E(u,v)+E(v,u)\).
\item An interchange of rows \(k<l\) is a succession of \(2(l-k)-1\) interchanges of adjacent rows (move row \(l\) up
to position \(k\), then the former row \(k\) down to position \(l\)), an odd number; so it changes the sign of
\(E\), which is Definition~\ref{r5d-def-det}(c). In particular \(E(A)=0\) if two rows of \(A\) are equal.
\item If \(B\) is \(A\) with \(\lambda\) times row \(s\) added to row \(r\), then by linearity in row \(r\),
\(E(B)=E(A)+\lambda E(A_s)\), where \(A_s\) is \(A\) with row \(r\) replaced by row \(s\). As \(A_s\) has two equal rows,
\(E(A_s)=0\) and \(E(B)=E(A)\): Definition~\ref{r5d-def-det}(b).
\item Column \(j\) of \(I_n\) is \(e_j\), and \(I_n(j|j)=I_{n-1}\); so \(E(I_n)=\det I_{n-1}=1\):
Definition~\ref{r5d-def-det}(d).
\end{itemize}
\textit{The expansion along row \(i\).}
\begin{itemize}
\item Row \(i\) of \(A\) is \(\sum_{j=1}^na_{ij}e_j'\). As \(\det\) is linear in row \(i\) (Lemma~\ref{r5d-lem-multilinear} and
Definition~\ref{r5d-def-det}(a)), \(\det A=\sum_{j=1}^na_{ij}\det G_j\), where \(G_j\) is \(A\) with row \(i\) replaced by
\(e_j'\).
\item In \(G_j\), subtract \(a_{kj}\) times row \(i\) from row \(k\), for every \(k\ne i\). As row \(i\) is \(e_j'\), this makes
entry \((k,j)\) zero and leaves the other entries of row \(k\) unchanged; it does not change the determinant
(Definition~\ref{r5d-def-det}(b)). The resulting matrix \(H_j\) has column \(j\) equal to \(e_i\), and \(H_j(i|j)=A(i|j)\).
\item Expanding \(\det H_j\) along column \(j\) (proved above), whose only nonzero entry is the \(1\) in row \(i\):
\(\det G_j=\det H_j=(-1)^{i+j}\det A(i|j)=c_{ij}\). Hence \(\det A=\sum_{j=1}^na_{ij}c_{ij}\).\qedhere
\end{itemize}
\end{proof}
For \(n=2\), the expansion along row \(1\) gives
\(\det\left(\begin{smallmatrix}a_{11}&a_{12}\\a_{21}&a_{22}\end{smallmatrix}\right)=a_{11}a_{22}-a_{12}a_{21}\), the number \(ad-bc\)
of the \(2\times2\) inverse formula of Proposition~\ref{r2-prop-inverse-2x2}; for \(n=3\) it gives~\citeat{lang}{Ch.~VI, \S\S1, 2}
\[
\det A=a_{11}(a_{22}a_{33}-a_{23}a_{32})-a_{12}(a_{21}a_{33}-a_{23}a_{31})+a_{13}(a_{21}a_{32}-a_{22}a_{31})
\]
In practice one expands along a row or column with many zeros; a row or column with a single nonzero entry is
expanded at once. Since a determinant function exists, Theorem~\ref{r5d-thm-gauss} shows that every reduction of \(A\) to
an upper triangular matrix by interchanges and additions gives the same number \((-1)^tu_{11}\cdots u_{nn}\), namely
\(\det A\).

\begin{theorem}[\citeat{hefferon}{Four.I.3, Theorem~3.15, Corollary~3.16}, \citeat{hoffman-kunze}{\S5.4}, \citeat{searle-khuri}{Theorem~3.1}]\label{r5d-thm-transpose}
For every \(A\in\R^{n\times n}\), \(\det A'=\det A\). Consequently Definition~\ref{r5d-def-det}(a)--(c),
Theorem~\ref{r5d-thm-basic}(a),(b) and Lemma~\ref{r5d-lem-multilinear} hold with columns in place of rows, and the
determinant of a lower triangular matrix is the product of its diagonal entries.
\end{theorem}
\begin{proof}
\begin{itemize}
\item Induction on \(n\). For \(n=1\), \(A'=A\). Let \(n\geqslant2\), and assume \(\det B'=\det B\) for every
\(B\in\R^{(n-1)\times(n-1)}\).
\item The \((k,1)\) entry of \(A'\) is \(a_{1k}\), and deleting row \(k\) and column \(1\) of \(A'\) leaves the transpose of
\(A(1|k)\): \(A'(k|1)=A(1|k)'\). By the expansion of \(\det A'\) along column \(1\) and the induction hypothesis,
\(\det A'=\sum_k(-1)^{k+1}a_{1k}\det A(1|k)'=\sum_k(-1)^{1+k}a_{1k}\det A(1|k)\), which is the expansion of \(\det A\) along
row \(1\) (Theorem~\ref{r5d-thm-expansion}).
\item An operation on the columns of \(A\) is the same operation on the rows of \(A'\), and two equal columns or a
zero column of \(A\) are two equal rows or a zero row of \(A'\). As \(\det A=\det A'\), the statements for rows
give those for columns. If \(A\) is lower triangular, \(A'\) is upper triangular with the same diagonal, so
\(\det A=\det A'=a_{11}\cdots a_{nn}\) (Theorem~\ref{r5d-thm-basic}(c)).\qedhere
\end{itemize}
\end{proof}
In computing a determinant one may therefore use column operations as well as row operations.

\begin{example}[cf.\ \citeat{searle-khuri}{\S3.1.2, Example~3.4}]\label{r5d-exm-expansion}
Let
\[
A=\begin{pmatrix}2&0&1&0\\1&3&0&2\\0&1&2&0\\1&0&0&1\end{pmatrix}
\]
Compute \(\det A\) by expansion along the first row, and check the value by expansion along the third column.

\textit{Solution.} Row \(1\) has the nonzero entries \(a_{11}=2\) and \(a_{13}=1\), so
\(\det A=2\det A(1|1)+1\det A(1|3)\). Each \(3\times3\) minor has a row with a single nonzero
entry, and expanding along that row (Theorem~\ref{r5d-thm-expansion}),
\[
\begin{aligned}
\det A(1|1)&=\det\begin{pmatrix}3&0&2\\1&2&0\\0&0&1\end{pmatrix}=(-1)^{3+3}\cdot1\cdot\det\begin{pmatrix}3&0\\1&2\end{pmatrix}=6\\
\det A(1|3)&=\det\begin{pmatrix}1&3&2\\0&1&0\\1&0&1\end{pmatrix}=(-1)^{2+2}\cdot1\cdot\det\begin{pmatrix}1&2\\1&1\end{pmatrix}=-1
\end{aligned}
\]
so \(\det A=2\cdot6+1\cdot(-1)=11\). Column \(3\) has the nonzero entries \(a_{13}=1\) and
\(a_{33}=2\), with cofactors \(c_{13}=\det A(1|3)=-1\) and
\[
c_{33}=(-1)^{3+3}\det A(3|3)=\det\begin{pmatrix}2&0&0\\1&3&2\\1&0&1\end{pmatrix}
=(-1)^{1+1}\cdot2\cdot\det\begin{pmatrix}3&2\\0&1\end{pmatrix}=6
\]
(expansion along row \(1\)). By Theorem~\ref{r5d-thm-expansion},
\(\det A=1\cdot(-1)+2\cdot6=11\).
\end{example}

\begin{theorem}[\citeat{hefferon}{Four.III.1, Definition~1.8, Theorem~1.9}, \citeat{searle-khuri}{\S3.3.3}, \citeat{harville}{Theorem~13.5.3, Corollary~13.5.4}]\label{r5d-thm-adjugate}
Let \(A\in\R^{n\times n}\) with \(n\geqslant2\), and let \(c_{ij}\) be the cofactors of \(A\). The
\textbf{adjugate} of \(A\) is the transpose of the matrix of cofactors: \(\operatorname{adj}A=(c_{ij})'\),
that is, \((\operatorname{adj}A)_{ij}=c_{ji}\). Then
\[
A\operatorname{adj}A=(\det A)\,I_n
\]
Hence, if \(\det A\ne0\), then \(A^{-1}=\frac1{\det A}\operatorname{adj}A\).
\end{theorem}
\begin{proof}
\begin{itemize}
\item The \((i,k)\) entry of \(A\operatorname{adj}A\) is
\(\sum_{j=1}^na_{ij}(\operatorname{adj}A)_{jk}=\sum_{j=1}^na_{ij}c_{kj}\).
\item For \(k=i\), this is the expansion of \(\det A\) along row \(i\) (Theorem~\ref{r5d-thm-expansion}).
\item For \(k\ne i\), let \(\tilde A\) be \(A\) with row \(k\) replaced by row \(i\). Since
\(\tilde A(k|j)=A(k|j)\), the cofactors of row \(k\) of \(\tilde A\) are \(c_{k1},\dots,c_{kn}\), and the
sum is the expansion of \(\det\tilde A\) along row \(k\). As \(\tilde A\) has two equal rows, the sum is
\(0\) (Theorem~\ref{r5d-thm-basic}(a)).
\item If \(\det A\ne0\), then \(A\bigl(\frac1{\det A}\operatorname{adj}A\bigr)=I_n\), so
\(\frac1{\det A}\operatorname{adj}A\) is the inverse of \(A\) (Theorem~\ref{r2-thm-invertible}).\qedhere
\end{itemize}
\end{proof}
For \(n=2\), \(\operatorname{adj}\left(\begin{smallmatrix}a&b\\c&d\end{smallmatrix}\right)
=\left(\begin{smallmatrix}d&-b\\-c&a\end{smallmatrix}\right)\), and the theorem gives again the formula of
Proposition~\ref{r2-prop-inverse-2x2}.

\begin{remark}[\citeat{harville-solutions}{Ch.~13, Exercise~6}]\label{r5d-rem-symmetric-cofactors}
Let \(A\in\R^{n\times n}\), \(n\geqslant2\), be symmetric. Deleting row \(j\) and column \(i\) of
\(A=A'\) gives the transpose of \(A(i|j)\): \(A(j|i)=A'(j|i)=A(i|j)'\). By
Theorem~\ref{r5d-thm-transpose}, \(\det A(j|i)=\det A(i|j)\), so the cofactors satisfy
\(c_{ji}=c_{ij}\).
\end{remark}

\begin{example}\label{r5d-exm-adjugate}
Find \(A^{-1}\) by Theorem~\ref{r5d-thm-adjugate} for
\[
A=\begin{pmatrix}1&1&0&0\\0&1&1&0\\0&0&1&1\\-1&0&0&1\end{pmatrix}
\]

\textit{Solution.} The cofactors are
\[
\begin{aligned}
c_{11}&=\det\left(\begin{smallmatrix}1&1&0\\0&1&1\\0&0&1\end{smallmatrix}\right)=1,&
c_{12}&=-\det\left(\begin{smallmatrix}0&1&0\\0&1&1\\-1&0&1\end{smallmatrix}\right)=1\\
c_{13}&=\det\left(\begin{smallmatrix}0&1&0\\0&0&1\\-1&0&1\end{smallmatrix}\right)=-1,&
c_{14}&=-\det\left(\begin{smallmatrix}0&1&1\\0&0&1\\-1&0&0\end{smallmatrix}\right)=1\\
c_{21}&=-\det\left(\begin{smallmatrix}1&0&0\\0&1&1\\0&0&1\end{smallmatrix}\right)=-1,&
c_{22}&=\det\left(\begin{smallmatrix}1&0&0\\0&1&1\\-1&0&1\end{smallmatrix}\right)=1\\
c_{23}&=-\det\left(\begin{smallmatrix}1&1&0\\0&0&1\\-1&0&1\end{smallmatrix}\right)=1,&
c_{24}&=\det\left(\begin{smallmatrix}1&1&0\\0&0&1\\-1&0&0\end{smallmatrix}\right)=-1\\
c_{31}&=\det\left(\begin{smallmatrix}1&0&0\\1&1&0\\0&0&1\end{smallmatrix}\right)=1,&
c_{32}&=-\det\left(\begin{smallmatrix}1&0&0\\0&1&0\\-1&0&1\end{smallmatrix}\right)=-1\\
c_{33}&=\det\left(\begin{smallmatrix}1&1&0\\0&1&0\\-1&0&1\end{smallmatrix}\right)=1,&
c_{34}&=-\det\left(\begin{smallmatrix}1&1&0\\0&1&1\\-1&0&0\end{smallmatrix}\right)=1\\
c_{41}&=-\det\left(\begin{smallmatrix}1&0&0\\1&1&0\\0&1&1\end{smallmatrix}\right)=-1,&
c_{42}&=\det\left(\begin{smallmatrix}1&0&0\\0&1&0\\0&1&1\end{smallmatrix}\right)=1\\
c_{43}&=-\det\left(\begin{smallmatrix}1&1&0\\0&1&0\\0&0&1\end{smallmatrix}\right)=-1,&
c_{44}&=\det\left(\begin{smallmatrix}1&1&0\\0&1&1\\0&0&1\end{smallmatrix}\right)=1
\end{aligned}
\]
Each \(3\times3\) minor here is triangular, so that its determinant is the product of its diagonal entries
(Theorems~\ref{r5d-thm-basic}(c) and~\ref{r5d-thm-transpose}), or has a row with a single nonzero entry, along
which it is expanded (Theorem~\ref{r5d-thm-expansion}). Along row \(1\) (Theorem~\ref{r5d-thm-expansion}),
\(\det A=1\cdot c_{11}+1\cdot c_{12}+0\cdot c_{13}+0\cdot c_{14}=2\). Hence
\[
A^{-1}=\frac12\operatorname{adj}A
=\frac12\begin{pmatrix}c_{11}&c_{21}&c_{31}&c_{41}\\c_{12}&c_{22}&c_{32}&c_{42}\\c_{13}&c_{23}&c_{33}&c_{43}\\c_{14}&c_{24}&c_{34}&c_{44}\end{pmatrix}
=\frac12\begin{pmatrix}1&-1&1&-1\\1&1&-1&1\\-1&1&1&-1\\1&-1&1&1\end{pmatrix}
\]
Check: \(\row_4(A)\operatorname{adj}A=(-1,0,0,1)\operatorname{adj}A=(-1+1,\ 1-1,\ -1+1,\ 1+1)=(0,0,0,2)\),
the last row of \(2I_4\).
\end{example}

\begin{theorem}[\citeat{harville-solutions}{Ch.~13, Exercise~11}, \citeat{hoffman-kunze}{\S5.4}, \citeat{lang}{Ch.~VI, Theorem~4.1}]\label{r5d-thm-cramer}
Let \(A\in\R^{n\times n}\) (\(n\geqslant2\)) with \(\det A\ne0\), and let \(b\in\R^n\). The system \(Ax=b\) has exactly
one solution \(x\), and
\[
x_j=\frac{\det A_j}{\det A}\qquad(1\leqslant j\leqslant n)
\]
where \(A_j\) is \(A\) with column \(j\) replaced by \(b\) (\textbf{Cramer's rule}).
\end{theorem}
\begin{proof}
\begin{itemize}
\item By Theorem~\ref{r5d-thm-adjugate}, \(A\) is invertible with \(A^{-1}=\frac1{\det A}\operatorname{adj}A\).
If \(Ax=b\), then \(x=A^{-1}(Ax)=A^{-1}b\), so there is at most one solution; and \(A(A^{-1}b)=b\), so
\(A^{-1}b\) is a solution. Hence the only solution is \(x=\frac1{\det A}(\operatorname{adj}A)b\).
\item Let \(c_{ij}\) be the cofactors of \(A\). Entry \(j\) of \((\operatorname{adj}A)b\) is \(\sum_{i=1}^n(\operatorname{adj}A)_{ji}b_i=\sum_{i=1}^nc_{ij}b_i\).
\item \(A_j\) differs from \(A\) only in column \(j\), so \(A_j(i|j)=A(i|j)\), and the \((i,j)\) entry
of \(A_j\) is \(b_i\). Expanding \(\det A_j\) along column \(j\) (Theorem~\ref{r5d-thm-expansion}) gives
\[
\det A_j=\sum_{i=1}^nb_i(-1)^{i+j}\det A_j(i|j)=\sum_{i=1}^nb_i(-1)^{i+j}\det A(i|j)=\sum_{i=1}^nc_{ij}b_i
\]
Hence \(x_j=\det A_j/\det A\).\qedhere
\end{itemize}
\end{proof}

\begin{example}\label{r5d-exm-cramer}
For the matrix \(A\) of Example~\ref{r5d-exm-adjugate} and \(b=(3,1,0,0)'\), find \(x_3\) in \(Ax=b\) by
Cramer's rule, and check it against \(x=A^{-1}b\) computed with that example.

\textit{Solution.} Here \(\det A=2\), and \(A_3\) is \(A\) with column \(3\) replaced by \(b\). Row \(3\) of \(A_3\)
has the single nonzero entry \(1\), at position \((3,4)\); expanding along it and then along the last row of the
\(3\times3\) minor (Theorem~\ref{r5d-thm-expansion}),
\[
\det A_3=\det\begin{pmatrix}1&1&3&0\\0&1&1&0\\0&0&0&1\\-1&0&0&1\end{pmatrix}
=(-1)^{3+4}\det\begin{pmatrix}1&1&3\\0&1&1\\-1&0&0\end{pmatrix}
=-(-1)(-1)^{3+1}\det\begin{pmatrix}1&3\\1&1\end{pmatrix}=-2
\]
So \(x_3=-2/2=-1\) by Theorem~\ref{r5d-thm-cramer}. With \(\operatorname{adj}A\) of
Example~\ref{r5d-exm-adjugate} and \(b=3e_1+e_2\), \((\operatorname{adj}A)b=3\col_1(\operatorname{adj}A)+\col_2(\operatorname{adj}A)\)
(Theorem~\ref{r1-thm-views}(b)), so
\[
x=\frac12(\operatorname{adj}A)b=\frac12\bigl(3(1,1,-1,1)'+(-1,1,1,-1)'\bigr)=\frac12(2,4,-2,2)'=(1,2,-1,1)'
\]
Its third entry is \(-1\); indeed \(Ax=(1+2,\ 2-1,\ -1+1,\ -1+1)'=b\).
\end{example}

\subsection*{Determinants of block matrices and the Schur complement}

\begin{theorem}[\citeat{abadir-magnus}{Exercises~5.22, 5.25}, \citeat{hoffman-kunze}{(5-20)}, \citeat{harville}{Theorem~13.3.1, (13.3.5)}]\label{r5d-thm-block-tri}
Let \(A\in\R^{p\times p}\), \(B\in\R^{p\times q}\), \(C\in\R^{q\times p}\) and \(D\in\R^{q\times q}\). Then
\[
\det\begin{pmatrix}A&B\\0&D\end{pmatrix}=\det A\cdot\det D=\det\begin{pmatrix}A&0\\C&D\end{pmatrix}
\]
In particular \(\det\diag(A,D)=\det A\cdot\det D\).
\end{theorem}
The diagonal blocks \(A\) and \(D\) are square, of possibly different sizes; for \(p\ne q\) the
off-diagonal blocks \(B\) and \(C\) are rectangular.
\begin{proof}
\begin{itemize}
\item The first row of \(\left(\begin{smallmatrix}I_p&0\\0&D\end{smallmatrix}\right)\) is
\((1,0,\dots,0)\), and deleting row \(1\) and column \(1\) leaves
\(\left(\begin{smallmatrix}I_{p-1}&0\\0&D\end{smallmatrix}\right)\); expanding along the first row (Theorem~\ref{r5d-thm-expansion}) \(p\)
times gives \(\det\left(\begin{smallmatrix}I_p&0\\0&D\end{smallmatrix}\right)=\det D\). Likewise,
expanding \(\left(\begin{smallmatrix}A&0\\0&I_q\end{smallmatrix}\right)\) \(q\) times along its last row
\((0,\dots,0,1)\), whose entry \(1\) is a diagonal entry (sign \(+1\)), gives \(\det A\).
\item \(\left(\begin{smallmatrix}I_p&B\\0&I_q\end{smallmatrix}\right)\) is upper triangular with
diagonal entries \(1\), so its determinant is \(1\) (Theorem~\ref{r5d-thm-basic}(c)).
\item By the block product formula (Theorem~\ref{r5b-thm-block}),
\[
\begin{pmatrix}I_p&0\\0&D\end{pmatrix}\begin{pmatrix}I_p&B\\0&I_q\end{pmatrix}\begin{pmatrix}A&0\\0&I_q\end{pmatrix}
=\begin{pmatrix}I_p&0\\0&D\end{pmatrix}\begin{pmatrix}A&B\\0&I_q\end{pmatrix}=\begin{pmatrix}A&B\\0&D\end{pmatrix}
\]
and by Theorem~\ref{r5d-thm-product}(b) its determinant is \(\det D\cdot1\cdot\det A\).
\item The transpose of \(\left(\begin{smallmatrix}A&0\\C&D\end{smallmatrix}\right)\) is
\(\left(\begin{smallmatrix}A'&C'\\0&D'\end{smallmatrix}\right)\) (Theorem~\ref{r5b-thm-block}(b)); by
Theorem~\ref{r5d-thm-transpose} and the case just proved, its determinant is
\(\det A'\cdot\det D'=\det A\cdot\det D\). Finally, \(\diag(A,D)\) is the case \(B=0\).\qedhere
\end{itemize}
\end{proof}

\begin{theorem}[\citeat{abadir-magnus}{Exercise~5.30}, \citeat{harville}{Theorem~13.3.8}]\label{r5d-thm-schur-det}
Let \(M=\begin{pmatrix}A&B\\C&D\end{pmatrix}\) with \(A\in\R^{p\times p}\), \(B\in\R^{p\times q}\),
\(C\in\R^{q\times p}\) and \(D\in\R^{q\times q}\).
\begin{enumerate}
\item If \(A\) is invertible, then \(\det M=\det A\cdot\det(D-CA^{-1}B)\).
\item If \(D\) is invertible, then \(\det M=\det D\cdot\det(A-BD^{-1}C)\).
\end{enumerate}
In particular, if \(A\) is invertible, then \(M\) is invertible if and only if its Schur complement
\(D-CA^{-1}B\) is invertible.
\end{theorem}
\begin{proof}
(a) By Theorem~\ref{r5b-thm-schur},
\(\left(\begin{smallmatrix}I_p&0\\-CA^{-1}&I_q\end{smallmatrix}\right)M
=\left(\begin{smallmatrix}A&B\\0&D-CA^{-1}B\end{smallmatrix}\right)\). The left factor is lower
triangular with diagonal entries \(1\), so its determinant is \(1\) (Theorem~\ref{r5d-thm-transpose}). By
Theorems~\ref{r5d-thm-product}(b) and~\ref{r5d-thm-block-tri},
\(\det M=\det A\cdot\det(D-CA^{-1}B)\).

(b) By the block product formula,
\(\left(\begin{smallmatrix}I_p&-BD^{-1}\\0&I_q\end{smallmatrix}\right)M
=\left(\begin{smallmatrix}A-BD^{-1}C&B-BD^{-1}D\\C&D\end{smallmatrix}\right)
=\left(\begin{smallmatrix}A-BD^{-1}C&0\\C&D\end{smallmatrix}\right)\). The left factor is upper
triangular with diagonal entries \(1\) (Theorem~\ref{r5d-thm-basic}(c)); as in (a), \(\det M=\det(A-BD^{-1}C)\cdot\det D\).

If \(A\) is invertible, then \(\det A\ne0\), so by (a) \(\det M\ne0\) if and only if
\(\det(D-CA^{-1}B)\ne0\); apply Theorem~\ref{r5d-thm-product}(a) to \(M\) and to \(D-CA^{-1}B\).
\end{proof}
Here \(A\) and \(D\) are square, of possibly different sizes \(p\) and \(q\);
\(B\in\R^{p\times q}\) and \(C\in\R^{q\times p}\) are rectangular when \(p\ne q\). For \(p=q=1\) and
\(a\ne0\), (a) reads
\(\det\left(\begin{smallmatrix}a&b\\c&d\end{smallmatrix}\right)=a\bigl(d-\frac{cb}a\bigr)=ad-bc\).

\begin{theorem}[\citeat{harville}{Theorem~18.1.1, Corollary~18.1.3}]\label{r5d-thm-det-lemma}
Let \(A\in\R^{n\times n}\) be invertible and \(u,v\in\R^n\). Then
\[
\det(A+uv')=\det A\cdot(1+v'A^{-1}u)
\]
In particular \(\det(I_n+uv')=1+v'u\). (This is the \textbf{matrix determinant lemma}.)
\end{theorem}
\begin{proof}
Apply Theorem~\ref{r5d-thm-schur-det} to
\(M=\left(\begin{smallmatrix}A&-u\\v'&1\end{smallmatrix}\right)\in\R^{(n+1)\times(n+1)}\) (\(p=n\), \(q=1\)).
\begin{itemize}
\item By (a), with the invertible block \(A\), whose Schur complement is the \(1\times1\) matrix
\(1-v'A^{-1}(-u)\): \(\det M=\det A\cdot(1+v'A^{-1}u)\).
\item By (b), with the invertible block \(D=(1)\): \(\det M=\det(1)\cdot\det\bigl(A-(-u)(1)^{-1}v'\bigr)=\det(A+uv')\).
\item Comparing gives the formula; for \(A=I_n\), \(A^{-1}=I_n\).\qedhere
\end{itemize}
\end{proof}

\begin{theorem}[\citeat{harville}{Theorems~18.1.1, 18.2.8, Corollary~18.2.10}]\label{r5d-thm-woodbury}
Let \(A\in\R^{n\times n}\) and \(T\in\R^{m\times m}\) be invertible, and let \(U\in\R^{n\times m}\) and
\(V\in\R^{m\times n}\).
\begin{enumerate}
\item \(\det(A+UTV)=\det A\cdot\det T\cdot\det(T^{-1}+VA^{-1}U)\). Hence \(A+UTV\) is invertible if and only if
\(T^{-1}+VA^{-1}U\) is invertible, and then
\[
(A+UTV)^{-1}=A^{-1}-A^{-1}U(T^{-1}+VA^{-1}U)^{-1}VA^{-1}
\]
(the \textbf{Sherman--Morrison--Woodbury formula}).
\item For \(u,v\in\R^n\), \(A+uv'\) is invertible if and only if \(1+v'A^{-1}u\ne0\), and then
\[
(A+uv')^{-1}=A^{-1}-\frac1{1+v'A^{-1}u}A^{-1}uv'A^{-1}
\]
\end{enumerate}
\end{theorem}
\begin{proof}
(a) Let \(M=\left(\begin{smallmatrix}A&-U\\V&T^{-1}\end{smallmatrix}\right)\in\R^{(n+m)\times(n+m)}\).
\begin{itemize}
\item The Schur complement of \(A\) in \(M\) is \(E=T^{-1}-VA^{-1}(-U)=T^{-1}+VA^{-1}U\), and that of \(T^{-1}\) is
\(F=A-(-U)(T^{-1})^{-1}V=A+UTV\).
\item By Theorem~\ref{r5d-thm-schur-det}(a),(b), \(\det M=\det A\cdot\det E=\det T^{-1}\cdot\det F\). As
\(\det T^{-1}=1/\det T\) (Theorem~\ref{r5d-thm-product}(c)), \(\det F=\det A\cdot\det T\cdot\det E\).
\item Since \(\det A\ne0\) and \(\det T\ne0\), \(\det F\ne0\) if and only if \(\det E\ne0\); so \(F\) is invertible if
and only if \(E\) is (Theorem~\ref{r5d-thm-product}(a)).
\item Let \(E\) and \(F\) be invertible. Theorem~\ref{r5b-thm-partitioned-inverse}, with the blocks \(A\), \(-U\), \(V\),
\(T^{-1}\) of \(M\) in place of \(A\), \(B\), \(C\), \(D\), gives \(M^{-1}\) twice. Its upper left block is
\(A^{-1}+A^{-1}(-U)E^{-1}VA^{-1}\) by (a) of that theorem and \(F^{-1}\) by (b). An invertible matrix has only one
inverse (Definition~\ref{r2-def-inverse}), so \(F^{-1}=A^{-1}-A^{-1}UE^{-1}VA^{-1}\).
\end{itemize}

(b) Take \(m=1\), \(U=u\), \(T=(1)\) and \(V=v'\) in (a). Then \(T^{-1}+VA^{-1}U\) is the \(1\times1\) matrix
\(1+v'A^{-1}u\), which is invertible if and only if it is not \(0\), with inverse \(1/(1+v'A^{-1}u)\).
\end{proof}
For instance, let \(G\in\R^{q\times n}\) with \(G'G\) invertible, and let \(x\in\R^n\). Appending the row \(x'\) to \(G\)
changes \(G'G\) into \(G'G+xx'\): by the block transpose and the block product formula
(Theorem~\ref{r5b-thm-block}(b),(c)),
\(\left(\begin{smallmatrix}G\\x'\end{smallmatrix}\right)'\left(\begin{smallmatrix}G\\x'\end{smallmatrix}\right)
=[G'\ \ x]\left(\begin{smallmatrix}G\\x'\end{smallmatrix}\right)=G'G+xx'\). So (b), with \(A=G'G\) and \(u=v=x\), gives
\((G'G+xx')^{-1}\) from \((G'G)^{-1}\) without a new inversion. The condition of (b) holds: with \(w=(G'G)^{-1}x\),
\(x'(G'G)^{-1}x=(G'Gw)'w=(Gw)'(Gw)\geqslant0\) (Theorem~\ref{r1-thm-transpose}), so \(1+x'(G'G)^{-1}x\geqslant1\).

\needspace{7\baselineskip}
\begin{example}[\citeat{searle-khuri}{Example~10.1}]\label{r5d-exm-equicorrelation}
Let \(n\geqslant2\) and \(\rho\in\R\). The \textbf{equicorrelation matrix} is
\(R=(1-\rho)I_n+\rho J_n\in\R^{n\times n}\), with diagonal entries \(1\) and all other entries
\(\rho\): if \(n\) variables have correlation \(\rho\) for every pair, their correlation matrix is \(R\).
Find \(\det R\), the values of \(\rho\) for which \(R\) is invertible, and \(R^{-1}\) for these \(\rho\).

\textit{Solution.}
\begin{itemize}
\item If \(\rho=1\), then \(R=J_n\) has equal rows (\(n\geqslant2\)), so \(\det R=0\)
(Theorem~\ref{r5d-thm-basic}(a)).
\item Let \(\rho\ne1\). As \(J_n=\one_n\one_n'\), \(R=A+uv'\) with \(A=(1-\rho)I_n\), \(u=\rho\one_n\) and
\(v=\one_n\). Here \(\det A=(1-\rho)^n\) (Theorem~\ref{r5d-thm-basic}(c)),
\(A^{-1}=(1-\rho)^{-1}I_n\) and \(v'A^{-1}u=\frac{\rho}{1-\rho}\one'\one=\frac{n\rho}{1-\rho}\)
(Proposition~\ref{r1-prop-centering}). By Theorem~\ref{r5d-thm-det-lemma},
\[
\det R=(1-\rho)^n\Bigl(1+\frac{n\rho}{1-\rho}\Bigr)=(1-\rho)^{n-1}(1-\rho+n\rho)
=(1-\rho)^{n-1}\bigl(1+(n-1)\rho\bigr)
\]
This formula also gives \(0\) for \(\rho=1\).
\item By Theorem~\ref{r5d-thm-product}(a), \(R\) is invertible if and only if \(\det R\ne0\), that is,
\(\rho\ne1\) and \(\rho\ne-1/(n-1)\). For instance, for \(n=4\) and \(\rho=-\frac13\) every row of
\(R\) has sum \(1+3\cdot(-\frac13)=0\), so \(R\one_4=0\).
\item For these \(\rho\), \(1+v'A^{-1}u=1+\frac{n\rho}{1-\rho}=\frac{1+(n-1)\rho}{1-\rho}\ne0\) and
\(A^{-1}uv'A^{-1}=\frac{\rho}{(1-\rho)^2}\one\one'=\frac{\rho}{(1-\rho)^2}J_n\). By Theorem~\ref{r5d-thm-woodbury}(b),
\[
R^{-1}=\frac1{1-\rho}I_n-\frac{1-\rho}{1+(n-1)\rho}\cdot\frac{\rho}{(1-\rho)^2}J_n
=\frac1{1-\rho}\Bigl(I_n-\frac{\rho}{1+(n-1)\rho}J_n\Bigr)
\]
For instance, for \(n=4\) and \(\rho=\frac13\), \(R^{-1}=\frac32\bigl(I_4-\frac16J_4\bigr)=\frac32I_4-\frac14J_4\). Check,
with \(J_4^2=4J_4\) (Proposition~\ref{r1-prop-centering}): \(\bigl(\frac23I_4+\frac13J_4\bigr)\bigl(\frac32I_4-\frac14J_4\bigr)
=I_4-\frac16J_4+\frac12J_4-\frac13J_4=I_4\).
\end{itemize}
\end{example}

\subsection*{Rank and the Schur complement}

\begin{theorem}[\citeat{harville}{Theorem~8.5.10}, \citeat{abadir-magnus}{Exercise~5.48}]\label{r5d-thm-schur-rank}
Let \(M=\begin{pmatrix}A&B\\C&D\end{pmatrix}\).
\begin{enumerate}
\item If \(A\in\R^{p\times p}\) is invertible, and \(B\in\R^{p\times q}\), \(C\in\R^{n\times p}\),
\(D\in\R^{n\times q}\), then \(\rank M=p+\rank(D-CA^{-1}B)\).
\item If \(D\in\R^{q\times q}\) is invertible, and \(A\in\R^{m\times p}\), \(B\in\R^{m\times q}\),
\(C\in\R^{q\times p}\), then \(\rank M=q+\rank(A-BD^{-1}C)\).
\end{enumerate}
\end{theorem}
\begin{proof}
(a)
\begin{itemize}
\item By Theorem~\ref{r5b-thm-schur},
\(\left(\begin{smallmatrix}I_p&0\\-CA^{-1}&I_n\end{smallmatrix}\right)M
\left(\begin{smallmatrix}I_p&-A^{-1}B\\0&I_q\end{smallmatrix}\right)
=\left(\begin{smallmatrix}A&0\\0&D-CA^{-1}B\end{smallmatrix}\right)\), and the two outer factors are
invertible (Proposition~\ref{r5b-prop-block-elementary}).
\item Invertible factors do not change the rank (Theorem~\ref{thm:r4-rank-product}(d)), and the rank
of a block-diagonal matrix is the sum of the ranks of its blocks
(Theorem~\ref{thm:r4-partitioned}(c)); so \(\rank M=\rank A+\rank(D-CA^{-1}B)\), where \(\rank A=p\)
(Theorem~\ref{thm:r4-rank-rules}(d)).
\end{itemize}

(b)
\begin{itemize}
\item By the block product formula,
\(\left(\begin{smallmatrix}I_m&-BD^{-1}\\0&I_q\end{smallmatrix}\right)M
=\left(\begin{smallmatrix}A-BD^{-1}C&0\\C&D\end{smallmatrix}\right)\) (as in the proof of
Theorem~\ref{r5d-thm-schur-det}(b)), and multiplying on the right by
\(\left(\begin{smallmatrix}I_p&0\\-D^{-1}C&I_q\end{smallmatrix}\right)\) gives
\(\left(\begin{smallmatrix}A-BD^{-1}C&0\\C-DD^{-1}C&D\end{smallmatrix}\right)
=\left(\begin{smallmatrix}A-BD^{-1}C&0\\0&D\end{smallmatrix}\right)\).
\item Both outer factors are invertible (Proposition~\ref{r5b-prop-block-elementary}); as in (a),
\(\rank M=\rank(A-BD^{-1}C)+\rank D=\rank(A-BD^{-1}C)+q\).\qedhere
\end{itemize}
\end{proof}

\begin{example}\label{r5d-exm-schur-rank}
Find the rank of
\[
M=\left(\begin{array}{cc|ccc}1&1&1&0&2\\1&2&0&1&1\\\hline2&3&2&0&3\\1&2&1&0&1\end{array}\right)
=\begin{pmatrix}A&B\\C&D\end{pmatrix}\in\R^{4\times5}
\]

\textit{Solution.} Here \(A=\left(\begin{smallmatrix}1&1\\1&2\end{smallmatrix}\right)\) has
\(\det A=1\) and \(A^{-1}=\left(\begin{smallmatrix}2&-1\\-1&1\end{smallmatrix}\right)\)
(Proposition~\ref{r2-prop-inverse-2x2}), \(B\in\R^{2\times3}\), \(C\in\R^{2\times2}\) and
\(D\in\R^{2\times3}\). Then
\[
CA^{-1}=\begin{pmatrix}2&3\\1&2\end{pmatrix}\begin{pmatrix}2&-1\\-1&1\end{pmatrix}=\begin{pmatrix}1&1\\0&1\end{pmatrix},\qquad
CA^{-1}B=\begin{pmatrix}1&1\\0&1\end{pmatrix}\begin{pmatrix}1&0&2\\0&1&1\end{pmatrix}=\begin{pmatrix}1&1&3\\0&1&1\end{pmatrix}
\]
and the Schur complement of \(A\) is
\[
D-CA^{-1}B=\begin{pmatrix}2&0&3\\1&0&1\end{pmatrix}-\begin{pmatrix}1&1&3\\0&1&1\end{pmatrix}
=\begin{pmatrix}1&-1&0\\1&-1&0\end{pmatrix}
\]
It is nonzero and its two rows are equal, so its row space is spanned by the nonzero row \((1,-1,0)\) and its rank
is \(1\) (Definition~\ref{def:r4-rank} and Proposition~\ref{r3-prop-dependence}(a)). By
Theorem~\ref{r5d-thm-schur-rank}(a),
\(\rank M=2+1=3\). The equal rows of the Schur complement give a relation among the rows of \(M\): by
Theorem~\ref{r5b-thm-schur}, the last two rows of
\(\left(\begin{smallmatrix}I_2&0\\-CA^{-1}&I_2\end{smallmatrix}\right)M\) are those of
\([0\ \ D-CA^{-1}B]\), that is (Theorem~\ref{r1-thm-views}(c)),
\[
\row_3(M)-\row_1(M)-\row_2(M)=(0,0,1,-1,0)=\row_4(M)-\row_2(M)
\]
so \(\row_3(M)=\row_1(M)+\row_4(M)\).
\end{example}

\needspace{12\baselineskip}
\subsection*{Rank and subdeterminants}

\begin{definition}[\citeat{bapat}{\S2.3}]\label{r5d-def-minor}
Let \(A\in\R^{m\times n}\) and \(1\leqslant k\leqslant\min\{m,n\}\). For rows \(i_1<\dots<i_k\) and columns \(j_1<\dots<j_k\)
of \(A\), the \(k\times k\) matrix with \((s,t)\) entry \(a_{i_sj_t}\) is a \(k\times k\) \textbf{submatrix} of \(A\), and its
determinant is a \(k\times k\) \textbf{minor} of \(A\).
\end{definition}

\begin{theorem}[\citeat{bapat}{\S2.3}]\label{r5d-thm-rank-minors}
Let \(A\in\R^{m\times n}\) have rank \(r\).
\begin{enumerate}
\item Every \(k\times k\) minor of \(A\) with \(k>r\) is \(0\).
\item If \(r\geqslant1\), then \(A\) has an \(r\times r\) submatrix with nonzero determinant.
\end{enumerate}
So if \(A\ne0\), then \(\rank A\) is the largest \(k\) for which \(A\) has a nonzero \(k\times k\) minor.
\end{theorem}
\begin{proof}
(a)
\begin{itemize}
\item Let \(A_*\) be the \(k\times k\) submatrix in rows \(i_1<\dots<i_k\) and columns \(j_1<\dots<j_k\), with \(k>r\). The
\(k\) columns \(a_{j_1},\dots,a_{j_k}\) of \(A\) lie in \(\mathcal C(A)\), of dimension \(r<k\), so they are linearly dependent
(Theorem~\ref{thm:r4-counting}(c)): \(c_1a_{j_1}+\dots+c_ka_{j_k}=0\) for some \(c=(c_1,\dots,c_k)'\ne0\).
\item Entries \(i_1,\dots,i_k\) of this vector equation read \(A_*c=0\). So the columns of \(A_*\) are linearly dependent,
\(\rank A_*<k\) (Theorem~\ref{thm:r4-rank-rules}(b)), and \(\det A_*=0\) (Theorem~\ref{r5d-thm-product}(a)).
\end{itemize}
(b)
\begin{itemize}
\item \(\rank A'=r\) (Theorem~\ref{thm:r4-rank-rules}(a)), so the pivot columns of \(A'\) are \(r\) linearly independent
columns of \(A'\) (Theorem~\ref{thm:r4-four-bases}(a)): \(A\) has \(r\) linearly independent rows, say rows
\(i_1<\dots<i_r\).
\item The matrix \(B\in\R^{r\times n}\) formed by these rows has full row rank, \(\rank B=r\)
(Theorem~\ref{thm:r4-rank-rules}(c)); so it has \(r\) linearly independent columns, say columns \(j_1<\dots<j_r\) (its pivot
columns, Theorem~\ref{thm:r4-four-bases}(a)).
\item They form an \(r\times r\) submatrix \(A_*\) of \(B\), and so of \(A\), with independent columns: \(\rank A_*=r\)
(Theorem~\ref{thm:r4-rank-rules}(b)), and \(\det A_*\ne0\) (Theorem~\ref{r5d-thm-product}(a)).
\end{itemize}
By (a) and (b), no \(k\times k\) minor with \(k>r\) is nonzero, and some \(r\times r\) minor is.
\end{proof}

\begin{example}[\citeat{lang}{Ch.~VI, \S9, Examples~1, 2}]\label{r5d-exm-rank-minors}
Find, by minors, the ranks of
\[
B=\begin{pmatrix}3&1&2&5\\1&2&-1&2\\4&3&1&1\end{pmatrix},\qquad C=\begin{pmatrix}3&1&2&5\\1&2&-1&2\\4&3&1&7\end{pmatrix}
\]

\textit{Solution.} Both are \(3\times4\), so their ranks are at most \(3\) (Theorem~\ref{thm:r4-rank-rules}(a)).
\begin{itemize}
\item \(B\): in the \(3\times3\) submatrix in columns \(1,2,3\), row \(3\), \((4,3,1)\), is the sum of rows \(1\) and \(2\). The
operations \(R_3\leftarrow R_3-R_1\), \(R_3\leftarrow R_3-R_2\) do not change the determinant (Definition~\ref{r5d-def-det}(b))
and give a zero row; so this minor is \(0\) (Theorem~\ref{r5d-thm-basic}(b)). The minor in columns \(2,3,4\) is, along row
\(1\) (Theorem~\ref{r5d-thm-expansion}),
\[
\det\begin{pmatrix}1&2&5\\2&-1&2\\3&1&1\end{pmatrix}=1\cdot(-1-2)-2\cdot(2-6)+5\cdot(2+3)=-3+8+25=30\ne0
\]
So \(\rank B=3\) (Theorem~\ref{r5d-thm-rank-minors}).
\item \(C\): row \(3\) of \(C\) is the sum of rows \(1\) and \(2\), and so is row \(3\) of each of its four \(3\times3\)
submatrices; as for \(B\), all four \(3\times3\) minors are \(0\). The \(2\times2\) minor in rows \(1,2\) and columns \(1,2\) is
\(3\cdot2-1\cdot1=5\ne0\). So \(\rank C=2\) (Theorem~\ref{r5d-thm-rank-minors}).
\end{itemize}
\end{example}

\subsection*{Permutations}

\begin{definition}[\citeat{hoffman-kunze}{\S5.3}, \citeat{hefferon}{Four.I.3, Definition~3.7, Example~3.9}]\label{r5d-def-permutation}
A \textbf{permutation} of degree \(n\) is a function \(\varphi\) from \(\{1,\dots,n\}\) to itself that is one-to-one
and onto: each of \(1,\dots,n\) equals \(\varphi(i)\) for exactly one \(i\). We write
\(\varphi=\langle\varphi(1),\dots,\varphi(n)\rangle\); the \textbf{identity permutation} is
\(\id=\langle1,2,\dots,n\rangle\). The \textbf{product} of permutations \(\varphi\), \(\psi\) of degree \(n\) is the
composite \(\varphi\psi\), \((\varphi\psi)(i)=\varphi(\psi(i))\), again a permutation, and the \textbf{inverse}
\(\varphi^{-1}\), also a permutation, is given by \(\varphi^{-1}(k)=i\) when \(\varphi(i)=k\); thus
\(\varphi\varphi^{-1}=\varphi^{-1}\varphi=\id\).
\end{definition}
There are \(n!=n(n-1)\cdots2\cdot1\) permutations of degree \(n\): \(\varphi(1)\) can be chosen in \(n\) ways, then
\(\varphi(2)\) in \(n-1\) ways, and so on. For \(n=3\) they are \(\langle1,2,3\rangle\), \(\langle1,3,2\rangle\),
\(\langle2,1,3\rangle\), \(\langle2,3,1\rangle\), \(\langle3,1,2\rangle\) and \(\langle3,2,1\rangle\).

\begin{definition}[\citeat{hefferon}{Four.I.3, Example~3.11}, \citeat{hoffman-kunze}{\S5.3}]\label{r5d-def-permutation-matrix}
The \textbf{permutation matrix} of a permutation \(\varphi\) of degree \(n\) is the matrix \(P_\varphi\in\R^{n\times n}\)
with rows \(e_{\varphi(1)}',\dots,e_{\varphi(n)}'\): its \((i,j)\) entry is \(1\) if \(j=\varphi(i)\) and \(0\)
otherwise.
\end{definition}
So \(P_\varphi\) has exactly one entry \(1\) in each row and in each column, and \(P_{\id}=I_n\). The interchange
matrix \(P_{rs}\) of Definition~\ref{r2-def-elementary} is \(P_\varphi\) for the permutation \(\varphi\) that exchanges
\(r\) and \(s\) and leaves the other numbers in place.

\begin{proposition}[\citeat{hoffman-kunze}{\S5.3, (5-16), Exercise~12(a), (c)}, \citeat{hefferon}{Four.I.4, Exercise~4.17; Four.I.3, Exercise~3.33}]\label{r5d-prop-permutation-matrix}
Let \(\varphi\), \(\psi\) be permutations of degree \(n\).
\begin{enumerate}
\item For \(A\in\R^{n\times p}\), row \(i\) of \(P_\varphi A\) is row \(\varphi(i)\) of \(A\): \(P_\varphi A\) lists the rows of
\(A\) in the order \(\varphi(1),\dots,\varphi(n)\).
\item \(P_\psi P_\varphi=P_{\varphi\psi}\).
\item \(P_\varphi'=P_{\varphi^{-1}}\), and \(P_\varphi'P_\varphi=P_\varphi P_\varphi'=I_n\); so \(P_\varphi\) is invertible, with
\(P_\varphi^{-1}=P_\varphi'\).
\end{enumerate}
\end{proposition}
\begin{proof}
(a) \(\row_i(P_\varphi A)=\row_i(P_\varphi)A=e_{\varphi(i)}'A=\row_{\varphi(i)}(A)\) (Theorem~\ref{r1-thm-views}(c)).

(b) By (a), row \(i\) of \(P_\psi P_\varphi\) is row \(\psi(i)\) of \(P_\varphi\), which is
\(e_{\varphi(\psi(i))}'=e_{(\varphi\psi)(i)}'\), row \(i\) of \(P_{\varphi\psi}\).

(c)
\begin{itemize}
\item The \((i,j)\) entry of \(P_\varphi'\) is the \((j,i)\) entry of \(P_\varphi\). It is \(1\) if \(i=\varphi(j)\), that is,
\(j=\varphi^{-1}(i)\), and \(0\) otherwise; this is the \((i,j)\) entry of \(P_{\varphi^{-1}}\).
\item By (b), \(P_\varphi'P_\varphi=P_{\varphi^{-1}}P_\varphi=P_{\varphi\varphi^{-1}}=P_{\id}=I_n\), and likewise
\(P_\varphi P_\varphi'=P_\varphi P_{\varphi^{-1}}=P_{\varphi^{-1}\varphi}=I_n\).\qedhere
\end{itemize}
\end{proof}
If the rows of a data matrix \(X_0\) hold the measurements on \(n\) units, \(P_\varphi X_0\) lists the same units in
another order, for instance sorted by one variable. By Theorem~\ref{r1-thm-transpose} and
Proposition~\ref{r5d-prop-permutation-matrix}(c), \((P_\varphi X_0)'(P_\varphi X_0)=X_0'P_\varphi'P_\varphi X_0=X_0'X_0\):
reordering the units does not change \(X_0'X_0\).

\begin{example}\label{r5d-exm-permutations}
\begin{enumerate}
\item The data matrix \(X_0\) below holds the measurements of two variables on four units. Sorting the units by the
first variable puts them in the order \(3,1,4,2\); so for \(\varphi=\langle3,1,4,2\rangle\)
(Proposition~\ref{r5d-prop-permutation-matrix}(a))
\[
X_0=\begin{pmatrix}2&7\\5&1\\1&4\\3&6\end{pmatrix},\qquad
P_\varphi=\begin{pmatrix}0&0&1&0\\1&0&0&0\\0&0&0&1\\0&1&0&0\end{pmatrix},\qquad
P_\varphi X_0=\begin{pmatrix}1&4\\2&7\\3&6\\5&1\end{pmatrix}
\]
\item Let \(\varphi=\langle2,3,1\rangle\) and \(\psi=\langle2,1,3\rangle\). Then
\(\varphi\psi=\langle\varphi(2),\varphi(1),\varphi(3)\rangle=\langle3,2,1\rangle\) and
\(\psi\varphi=\langle\psi(2),\psi(3),\psi(1)\rangle=\langle1,3,2\rangle\), so \(\varphi\psi\ne\psi\varphi\). Since
\(\varphi(3)=1\), \(\varphi(1)=2\) and \(\varphi(2)=3\), the inverse is \(\varphi^{-1}=\langle3,1,2\rangle\).
\end{enumerate}
\end{example}

\subsection*{Exercises}

\begin{exercise}\label{r5b-dr-block-product}
Let
\[
A=\left(\begin{array}{cc|cc}1&2&0&-1\\0&1&3&2\\\hline2&-1&1&0\end{array}\right),\qquad
B=\begin{pmatrix}1&0&2\\-1&1&0\\2&0&1\\0&3&-1\end{pmatrix}
\]
where the rows of \(A\) are split into groups of \(2\) and \(1\) rows and its columns into groups of \(2\) and \(2\)
columns.
\begin{enumerate}
\item Split the columns of \(B\) into groups of \(1\) and \(2\) columns, and split its rows so that the partitions of
\(A\) and \(B\) are conformal for multiplication. Write down the blocks \(A_{ij}\) and \(B_{jk}\).
\item Compute the blocks of \(AB\) by Theorem~\ref{r5b-thm-block}(c).
\item Check the result against the ordinary product \(AB\).
\end{enumerate}
\end{exercise}
\begin{solution}
(a) The row groups of \(B\) must be the column groups of \(A\), of \(2\) and \(2\) (Definition~\ref{r5b-def-conformal}):
\[
B=\left(\begin{array}{c|cc}1&0&2\\-1&1&0\\\hline2&0&1\\0&3&-1\end{array}\right)
\]
So \(A_{11}=\left(\begin{smallmatrix}1&2\\0&1\end{smallmatrix}\right)\),
\(A_{12}=\left(\begin{smallmatrix}0&-1\\3&2\end{smallmatrix}\right)\), \(A_{21}=(2,\ -1)\), \(A_{22}=(1,\ 0)\), and
\(B_{11}=\left(\begin{smallmatrix}1\\-1\end{smallmatrix}\right)\),
\(B_{12}=\left(\begin{smallmatrix}0&2\\1&0\end{smallmatrix}\right)\),
\(B_{21}=\left(\begin{smallmatrix}2\\0\end{smallmatrix}\right)\),
\(B_{22}=\left(\begin{smallmatrix}0&1\\3&-1\end{smallmatrix}\right)\).

(b) \(AB\) is partitioned with the row groups of \(A\) (\(2\) and \(1\) rows) and the column groups of \(B\) (\(1\) and \(2\)
columns), so its blocks are \(2\times1\), \(2\times2\), \(1\times1\) and \(1\times2\):
\[
\begin{aligned}
(AB)_{11}&=A_{11}B_{11}+A_{12}B_{21}=\begin{pmatrix}-1\\-1\end{pmatrix}+\begin{pmatrix}0\\6\end{pmatrix}=\begin{pmatrix}-1\\5\end{pmatrix}\\
(AB)_{12}&=A_{11}B_{12}+A_{12}B_{22}=\begin{pmatrix}2&2\\1&0\end{pmatrix}+\begin{pmatrix}-3&1\\6&1\end{pmatrix}=\begin{pmatrix}-1&3\\7&1\end{pmatrix}\\
(AB)_{21}&=A_{21}B_{11}+A_{22}B_{21}=3+2=5\\
(AB)_{22}&=A_{21}B_{12}+A_{22}B_{22}=(-1,\ 4)+(0,\ 1)=(-1,\ 5)
\end{aligned}
\]
For instance, the first row of \(A_{12}B_{22}\) is \(0\cdot(0,\ 1)+(-1)(3,\ -1)=(-3,\ 1)\) (Theorem~\ref{r1-thm-views}(c)).
Hence
\[
AB=\left(\begin{array}{c|cc}-1&-1&3\\5&7&1\\\hline5&-1&5\end{array}\right)
\]

(c) By Theorem~\ref{r1-thm-views}(c), \(\row_i(AB)=\row_i(A)B\):
\begin{itemize}
\item \(\row_1(A)B=(1,0,2)+2(-1,1,0)-(0,3,-1)=(-1,-1,3)\);
\item \(\row_2(A)B=(-1,1,0)+3(2,0,1)+2(0,3,-1)=(5,7,1)\);
\item \(\row_3(A)B=2(1,0,2)-(-1,1,0)+(2,0,1)=(5,-1,5)\).
\end{itemize}
These are the rows found in (b).
\end{solution}

\begin{exercise}[cf.\ \citeat{abadir-magnus}{Exercise~2.16}]\label{r5b-ex-partitioned}
Let \(A\) and \(B\) be the \(3\times5\) matrices and \(C\) the \(5\times4\) matrix partitioned
as
\[
A=\left(\begin{array}{ccc|cc}1&3&-2&1&2\\2&1&0&-1&3\\\hline0&0&1&2&1\end{array}\right),\quad
B=\left(\begin{array}{ccc|cc}1&-3&-2&2&1\\1&2&3&2&0\\\hline1&0&2&0&1\end{array}\right),\quad
C=\left(\begin{array}{cc|cc}1&0&2&1\\0&2&0&0\\-1&0&3&1\\\hline3&1&0&2\\2&-1&1&1\end{array}\right)
\]
With blocks \(A=(A_{st})\), \(B=(B_{st})\), \(C=(C_{st})\) (\(s,t=1,2\)), show that
\[
A+B=(A_{st}+B_{st}),\qquad
A'=\begin{pmatrix}A_{11}'&A_{21}'\\A_{12}'&A_{22}'\end{pmatrix},\qquad
AC=\begin{pmatrix}A_{11}C_{11}+A_{12}C_{21}&A_{11}C_{12}+A_{12}C_{22}\\A_{21}C_{11}+A_{22}C_{21}&A_{21}C_{12}+A_{22}C_{22}\end{pmatrix}
\]
\end{exercise}
\begin{solution}
Adding entrywise,
\[
A+B=\left(\begin{array}{ccc|cc}2&0&-4&3&3\\3&3&3&1&3\\\hline1&0&3&2&2\end{array}\right)
\]
and each block is the sum of the corresponding blocks of \(A\) and \(B\). Transposing,
\[
A'=\left(\begin{array}{cc|c}1&2&0\\3&1&0\\-2&0&1\\\hline1&-1&2\\2&3&1\end{array}\right)
=\begin{pmatrix}A_{11}'&A_{21}'\\A_{12}'&A_{22}'\end{pmatrix}
\]
Multiplying in the ordinary way,
\[
AC=\left(\begin{array}{cc|cc}10&5&-2&3\\5&-2&7&3\\\hline7&1&4&6\end{array}\right)
\]
The blockwise products give the same four blocks:
\[
\begin{aligned}
A_{11}C_{11}+A_{12}C_{21}&=\begin{pmatrix}3&6\\2&2\end{pmatrix}+\begin{pmatrix}7&-1\\3&-4\end{pmatrix}=\begin{pmatrix}10&5\\5&-2\end{pmatrix}\\
A_{11}C_{12}+A_{12}C_{22}&=\begin{pmatrix}-4&-1\\4&2\end{pmatrix}+\begin{pmatrix}2&4\\3&1\end{pmatrix}=\begin{pmatrix}-2&3\\7&3\end{pmatrix}\\
A_{21}C_{11}+A_{22}C_{21}&=(-1,\ 0)+(8,\ 1)=(7,\ 1)\\
A_{21}C_{12}+A_{22}C_{22}&=(3,\ 1)+(1,\ 5)=(4,\ 6)
\end{aligned}
\]
For instance,
\(A_{12}C_{21}=\left(\begin{smallmatrix}1&2\\-1&3\end{smallmatrix}\right)\left(\begin{smallmatrix}3&1\\2&-1\end{smallmatrix}\right)=\left(\begin{smallmatrix}7&-1\\3&-4\end{smallmatrix}\right)\).
This agrees with Theorem~\ref{r5b-thm-block}.
\end{solution}

\begin{exercise}\label{r5b-ex-block-numeric}
Let
\[
N=\left(\begin{array}{ccc|cc}1&0&0&0&0\\1&1&0&0&0\\0&1&1&0&0\\\hline1&0&1&3&1\\0&1&1&2&1\end{array}\right)
=\begin{pmatrix}A&0\\C&D\end{pmatrix}
\]
Here \(A\in\R^{3\times3}\), \(C\in\R^{2\times3}\) and \(D\in\R^{2\times2}\).
Find \(N^{-1}\) by applying Proposition~\ref{r5b-prop-block-tri}(c) to \(N'\), and check the
answer by multiplication.
\end{exercise}
\begin{solution}
By Theorem~\ref{r5b-thm-block}(b), \(N'\) with groups \((3,2)\) is
\(N'=\left(\begin{smallmatrix}A'&C'\\0&D'\end{smallmatrix}\right)\), upper block triangular, with
\[
A'=\begin{pmatrix}1&1&0\\0&1&1\\0&0&1\end{pmatrix},\qquad
C'=\begin{pmatrix}1&0\\0&1\\1&1\end{pmatrix}\in\R^{3\times2},\qquad
D'=\begin{pmatrix}3&2\\1&1\end{pmatrix}
\]
The matrix \(A'\) is the block \(A_{22}\) of Example~\ref{r5b-exm-block-inverse}, so
\[
(A')^{-1}=\begin{pmatrix}1&-1&1\\0&1&-1\\0&0&1\end{pmatrix}
\]
and by Proposition~\ref{r2-prop-inverse-2x2} (\(3\cdot1-2\cdot1=1\)),
\((D')^{-1}=\left(\begin{smallmatrix}1&-2\\-1&3\end{smallmatrix}\right)\). Then
\[
(A')^{-1}C'=\begin{pmatrix}2&0\\-1&0\\1&1\end{pmatrix},\qquad
-(A')^{-1}C'(D')^{-1}=-\begin{pmatrix}2&0\\-1&0\\1&1\end{pmatrix}\begin{pmatrix}1&-2\\-1&3\end{pmatrix}
=\begin{pmatrix}-2&4\\1&-2\\0&-1\end{pmatrix}
\]
so by Proposition~\ref{r5b-prop-block-tri}(c) with \(p=3\), \(q=2\)
\[
(N')^{-1}=\left(\begin{array}{ccc|cc}1&-1&1&-2&4\\0&1&-1&1&-2\\0&0&1&0&-1\\\hline0&0&0&1&-2\\0&0&0&-1&3\end{array}\right),\qquad
N^{-1}=\bigl((N')^{-1}\bigr)'=\left(\begin{array}{ccc|cc}1&0&0&0&0\\-1&1&0&0&0\\1&-1&1&0&0\\\hline-2&1&0&1&-1\\4&-2&-1&-2&3\end{array}\right)
\]
by Theorem~\ref{r2-thm-inverse-props}(c). Check, with the blocks
\[
A^{-1}=\begin{pmatrix}1&0&0\\-1&1&0\\1&-1&1\end{pmatrix},\qquad
R=\begin{pmatrix}-2&1&0\\4&-2&-1\end{pmatrix}\in\R^{2\times3},\qquad
D^{-1}=\begin{pmatrix}1&-1\\-2&3\end{pmatrix}
\]
of \(N^{-1}\): the
\((2,1)\) block of \(NN^{-1}\) is
\[
CA^{-1}+DR=\begin{pmatrix}2&-1&1\\0&0&1\end{pmatrix}+\begin{pmatrix}-2&1&-1\\0&0&-1\end{pmatrix}=0
\]
The \((1,2)\) block is \(A\cdot0+0\cdot D^{-1}=0\), and the diagonal blocks are \(AA^{-1}=I_3\)
and \(DD^{-1}=I_2\). So \(NN^{-1}=I_5\).
\end{solution}

\begin{exercise}\label{r5b-ex-partitioned-inverse-drill}
Let
\[
M=\left(\begin{array}{ccc|cc}1&1&0&1&0\\0&1&1&0&1\\0&0&1&1&1\\\hline1&1&1&4&2\\0&1&2&2&3\end{array}\right)
=\begin{pmatrix}A&B\\C&D\end{pmatrix}
\]
Here \(A\in\R^{3\times3}\), \(B\in\R^{3\times2}\), \(C\in\R^{2\times3}\), \(D\in\R^{2\times2}\).
Find the Schur complement \(E\) of \(A\) and \(\det M\). Then compute \(M^{-1}\) with
Corollary~\ref{r5b-cor-schur-inverse}(a).
\end{exercise}
\begin{solution}
Here \(A\) is the block \(A_{22}\) of Example~\ref{r5b-exm-block-inverse}, so
\[
A=\begin{pmatrix}1&1&0\\0&1&1\\0&0&1\end{pmatrix},\qquad A^{-1}=\begin{pmatrix}1&-1&1\\0&1&-1\\0&0&1\end{pmatrix}
\]
and \(\det A=1\) (Theorem~\ref{r5d-thm-basic}(c)). With
\(B=\left(\begin{smallmatrix}1&0\\0&1\\1&1\end{smallmatrix}\right)\),
\(C=\left(\begin{smallmatrix}1&1&1\\0&1&2\end{smallmatrix}\right)\) and
\(D=\left(\begin{smallmatrix}4&2\\2&3\end{smallmatrix}\right)\),
\[
A^{-1}B=\begin{pmatrix}1&-1&1\\0&1&-1\\0&0&1\end{pmatrix}\begin{pmatrix}1&0\\0&1\\1&1\end{pmatrix}=\begin{pmatrix}2&0\\-1&0\\1&1\end{pmatrix},\qquad
CA^{-1}=\begin{pmatrix}1&1&1\\0&1&2\end{pmatrix}\begin{pmatrix}1&-1&1\\0&1&-1\\0&0&1\end{pmatrix}=\begin{pmatrix}1&0&1\\0&1&1\end{pmatrix}
\]
\[
E=D-CA^{-1}B=\begin{pmatrix}4&2\\2&3\end{pmatrix}-\begin{pmatrix}1&0&1\\0&1&1\end{pmatrix}\begin{pmatrix}1&0\\0&1\\1&1\end{pmatrix}
=\begin{pmatrix}4&2\\2&3\end{pmatrix}-\begin{pmatrix}2&1\\1&2\end{pmatrix}=\begin{pmatrix}2&1\\1&1\end{pmatrix}
\]
So \(\det M=\det A\cdot\det E=1\cdot1=1\) (Theorem~\ref{r5d-thm-schur-det}(a)), and
\(E^{-1}=\left(\begin{smallmatrix}1&-1\\-1&2\end{smallmatrix}\right)\) by
Proposition~\ref{r2-prop-inverse-2x2}.
\begin{itemize}
\item With \(L=\left(\begin{smallmatrix}I_3&0\\-CA^{-1}&I_2\end{smallmatrix}\right)\) and
\(U=\left(\begin{smallmatrix}I_3&-A^{-1}B\\0&I_2\end{smallmatrix}\right)\),
Corollary~\ref{r5b-cor-schur-inverse}(a), with \(E\) in place of \(S\), gives \(M^{-1}=U\diag(A^{-1},E^{-1})L\).
\item By the block product formula,
\[
U\diag(A^{-1},E^{-1})=\begin{pmatrix}A^{-1}&-A^{-1}BE^{-1}\\0&E^{-1}\end{pmatrix}
=\left(\begin{array}{ccc|cc}1&-1&1&-2&2\\0&1&-1&1&-1\\0&0&1&0&-1\\\hline0&0&0&1&-1\\0&0&0&-1&2\end{array}\right)
\]
Here \(-A^{-1}BE^{-1}=-\left(\begin{smallmatrix}2&0\\-1&0\\1&1\end{smallmatrix}\right)\left(\begin{smallmatrix}1&-1\\-1&2\end{smallmatrix}\right)
=\left(\begin{smallmatrix}-2&2\\1&-1\\0&-1\end{smallmatrix}\right)\in\R^{3\times2}\). Multiplying on the right by \(L\) adds
to block column \(1\) the block column \(2\) times \(-CA^{-1}\in\R^{2\times3}\):
\[
M^{-1}=\left(\begin{array}{ccc|cc}1&-1&1&-2&2\\0&1&-1&1&-1\\0&0&1&0&-1\\\hline0&0&0&1&-1\\0&0&0&-1&2\end{array}\right)
\left(\begin{array}{ccc|cc}1&0&0&0&0\\0&1&0&0&0\\0&0&1&0&0\\\hline-1&0&-1&1&0\\0&-1&-1&0&1\end{array}\right)
=\left(\begin{array}{ccc|cc}3&-3&1&-2&2\\-1&2&-1&1&-1\\0&1&2&0&-1\\\hline-1&1&0&1&-1\\1&-2&-1&-1&2\end{array}\right)
\]
\end{itemize}
Check: \(\row_1(M)M^{-1}\) is the sum of rows \(1\), \(2\), \(4\) of \(M^{-1}\), namely
\((3-1-1,\ -3+2+1,\ 1-1+0,\ -2+1+1,\ 2-1-1)=(1,0,0,0,0)\); and
\(\row_4(M)M^{-1}=(1,1,1,4,2)M^{-1}=(3-1+0-4+2,\ -3+2+1+4-4,\ 1-1+2+0-2,\ -2+1+0+4-2,\ 2-1-1-4+4)=(0,0,0,1,0)\).
\end{solution}

\begin{exercise}\label{r5b-dr-partitioned-inverse}
Let
\[
M=\left(\begin{array}{cc|cc}4&2&1&2\\2&1&0&1\\\hline1&0&1&1\\2&1&1&2\end{array}\right)
=\begin{pmatrix}A&B\\C&D\end{pmatrix}
\]
\begin{enumerate}
\item Show that \(A\) is not invertible, so that Theorem~\ref{r5b-thm-partitioned-inverse}(a) does not apply.
\item Find \(D^{-1}\), the Schur complement \(F=A-BD^{-1}C\) of \(D\), and \(\det M\).
\item Use Theorem~\ref{r5b-thm-partitioned-inverse}(b) to find \(M^{-1}\), and check the answer.
\end{enumerate}
\end{exercise}
\begin{solution}
(a) Here \(A=\left(\begin{smallmatrix}4&2\\2&1\end{smallmatrix}\right)\),
\(B=\left(\begin{smallmatrix}1&2\\0&1\end{smallmatrix}\right)\),
\(C=\left(\begin{smallmatrix}1&0\\2&1\end{smallmatrix}\right)\) and
\(D=\left(\begin{smallmatrix}1&1\\1&2\end{smallmatrix}\right)\). Since \(\det A=4\cdot1-2\cdot2=0\), \(A\) is not
invertible (Theorem~\ref{r5d-thm-product}(a)).

(b)
\begin{itemize}
\item \(\det D=1\cdot2-1\cdot1=1\), so \(D^{-1}=\left(\begin{smallmatrix}2&-1\\-1&1\end{smallmatrix}\right)\) (the
\(2\times2\) case of Theorem~\ref{r5d-thm-adjugate}).
\item Then
\[
BD^{-1}=\begin{pmatrix}1&2\\0&1\end{pmatrix}\begin{pmatrix}2&-1\\-1&1\end{pmatrix}=\begin{pmatrix}0&1\\-1&1\end{pmatrix},\qquad
BD^{-1}C=\begin{pmatrix}0&1\\-1&1\end{pmatrix}\begin{pmatrix}1&0\\2&1\end{pmatrix}=\begin{pmatrix}2&1\\1&1\end{pmatrix}
\]
\item So \(F=A-BD^{-1}C=\left(\begin{smallmatrix}2&1\\1&0\end{smallmatrix}\right)\). Its determinant is
\(2\cdot0-1\cdot1=-1\), so \(F\) is invertible, and by the same formula
\(F^{-1}=\left(\begin{smallmatrix}0&1\\1&-2\end{smallmatrix}\right)\).
\item By Theorem~\ref{r5d-thm-schur-det}(b), \(\det M=\det D\cdot\det F=1\cdot(-1)=-1\).
\end{itemize}

(c)
\begin{itemize}
\item The products in the formula of Theorem~\ref{r5b-thm-partitioned-inverse}(b):
\[
\begin{aligned}
F^{-1}BD^{-1}&=\begin{pmatrix}0&1\\1&-2\end{pmatrix}\begin{pmatrix}0&1\\-1&1\end{pmatrix}=\begin{pmatrix}-1&1\\2&-1\end{pmatrix}\\
D^{-1}C&=\begin{pmatrix}2&-1\\-1&1\end{pmatrix}\begin{pmatrix}1&0\\2&1\end{pmatrix}=\begin{pmatrix}0&-1\\1&1\end{pmatrix}\\
D^{-1}CF^{-1}&=\begin{pmatrix}0&-1\\1&1\end{pmatrix}\begin{pmatrix}0&1\\1&-2\end{pmatrix}=\begin{pmatrix}-1&2\\1&-1\end{pmatrix}\\
D^{-1}CF^{-1}BD^{-1}&=\begin{pmatrix}-1&2\\1&-1\end{pmatrix}\begin{pmatrix}0&1\\-1&1\end{pmatrix}=\begin{pmatrix}-2&1\\1&0\end{pmatrix}
\end{aligned}
\]
\item So \(D^{-1}+D^{-1}CF^{-1}BD^{-1}=\left(\begin{smallmatrix}0&0\\0&1\end{smallmatrix}\right)\), and by
Theorem~\ref{r5b-thm-partitioned-inverse}(b)
\[
M^{-1}=\begin{pmatrix}F^{-1}&-F^{-1}BD^{-1}\\-D^{-1}CF^{-1}&D^{-1}+D^{-1}CF^{-1}BD^{-1}\end{pmatrix}
=\left(\begin{array}{cc|cc}0&1&1&-1\\1&-2&-2&1\\\hline1&-2&0&0\\-1&1&0&1\end{array}\right)
\]
\item Check: \(\row_1(M)M^{-1}=4(0,1,1,-1)+2(1,-2,-2,1)+(1,-2,0,0)+2(-1,1,0,1)=(1,0,0,0)\) and
\(\row_3(M)M^{-1}=(0,1,1,-1)+(1,-2,0,0)+(-1,1,0,1)=(0,0,1,0)\); rows \(2\) and \(4\) likewise give \((0,1,0,0)\) and
\((0,0,0,1)\).
\end{itemize}
\end{solution}

\begin{exercise}\label{r5d-dr-elimination}
Let
\[
A=\begin{pmatrix}0&0&2&1\\1&2&0&-1\\2&3&1&0\\-1&0&0&1\end{pmatrix},\qquad
B=\begin{pmatrix}1&-1&2&1\\2&-1&3&4\\-1&3&-4&2\\3&-2&5&5\end{pmatrix}
\]
\begin{enumerate}
\item Compute \(\det A\) by forward elimination (Theorem~\ref{r5d-thm-gauss}), keeping track of the interchanges.
\item Compute \(\det B\) in the same way.
\item Which of \(A\) and \(B\) is invertible?
\end{enumerate}
\end{exercise}
\begin{solution}
(a)
\[
\begin{gathered}
A\xrightarrow{R_1\leftrightarrow R_2}
\begin{pmatrix}1&2&0&-1\\0&0&2&1\\2&3&1&0\\-1&0&0&1\end{pmatrix}
\xrightarrow[R_4\leftarrow R_4+R_1]{R_3\leftarrow R_3-2R_1}
\begin{pmatrix}1&2&0&-1\\0&0&2&1\\0&-1&1&2\\0&2&0&0\end{pmatrix}
\xrightarrow{R_2\leftrightarrow R_3}
\begin{pmatrix}1&2&0&-1\\0&-1&1&2\\0&0&2&1\\0&2&0&0\end{pmatrix}
\\
\xrightarrow{R_4\leftarrow R_4+2R_2}
\begin{pmatrix}1&2&0&-1\\0&-1&1&2\\0&0&2&1\\0&0&2&4\end{pmatrix}
\xrightarrow{R_4\leftarrow R_4-R_3}
\begin{pmatrix}1&2&0&-1\\0&-1&1&2\\0&0&2&1\\0&0&0&3\end{pmatrix}=U
\end{gathered}
\]
\begin{itemize}
\item Since \(a_{11}=0\), Step 2 of Definition~\ref{r2-def-gauss-jordan} interchanges rows \(1\) and \(2\). After the
first column is cleared, the entry in position \((2,2)\) is \(0\), and rows \(2\) and \(3\) are interchanged.
\item There are two interchanges, and the additions do not change the determinant. By Theorem~\ref{r5d-thm-gauss},
\(\det A=(-1)^2\cdot1\cdot(-1)\cdot2\cdot3=-6\).
\end{itemize}

(b)
\[
B\xrightarrow[R_4\leftarrow R_4-3R_1]{R_2\leftarrow R_2-2R_1,\ R_3\leftarrow R_3+R_1}
\begin{pmatrix}1&-1&2&1\\0&1&-1&2\\0&2&-2&3\\0&1&-1&2\end{pmatrix}
\xrightarrow[R_4\leftarrow R_4-R_2]{R_3\leftarrow R_3-2R_2}
\begin{pmatrix}1&-1&2&1\\0&1&-1&2\\0&0&0&-1\\0&0&0&0\end{pmatrix}=U
\]
There are no interchanges, and \(U\) is upper triangular with \(u_{33}=u_{44}=0\). So
\(\det B=1\cdot1\cdot0\cdot0=0\) by Theorem~\ref{r5d-thm-gauss}.

(c) By Theorem~\ref{r5d-thm-product}(a), \(A\) is invertible, as \(\det A=-6\ne0\), and \(B\) is not. Indeed
\(\row_4(B)=\row_1(B)+\row_2(B)\).
\end{solution}

\begin{exercise}[\citeat{searle-khuri}{Exercise~3.1}]\label{r5d-ex-three-by-three}
Let
\[
A=\begin{pmatrix}1&5&-5\\3&2&-5\\6&-2&-5\end{pmatrix},\qquad B=\begin{pmatrix}2&6&5\\-2&-3&-5\\2&3&9\end{pmatrix}
\]
Show that \(\det A=-5\) and \(\det B=24\), in each case (i) by expansion along the first row and (ii) by row and
column operations.
\end{exercise}
\begin{solution}
(\(A\))
\begin{itemize}
\item (i) Along the first row (Theorem~\ref{r5d-thm-expansion}),
\(\det A=1\det\left(\begin{smallmatrix}2&-5\\-2&-5\end{smallmatrix}\right)
-5\det\left(\begin{smallmatrix}3&-5\\6&-5\end{smallmatrix}\right)
-5\det\left(\begin{smallmatrix}3&2\\6&-2\end{smallmatrix}\right)
=(-10-10)-5(-15+30)-5(-6-12)=-20-75+90=-5\).
\item (ii) Column \(3\) is \(-5\) times \((1,1,1)'\). Taking the factor \(-5\) out of column \(3\)
(Definition~\ref{r5d-def-det}(a) for columns, Theorem~\ref{r5d-thm-transpose}) and then subtracting row \(1\) from rows \(2\) and \(3\)
(Definition~\ref{r5d-def-det}(b)),
\[
\det A=-5\det\begin{pmatrix}1&5&1\\3&2&1\\6&-2&1\end{pmatrix},\qquad
\begin{pmatrix}1&5&1\\3&2&1\\6&-2&1\end{pmatrix}
\xrightarrow[R_3\leftarrow R_3-R_1]{R_2\leftarrow R_2-R_1}
\begin{pmatrix}1&5&1\\2&-3&0\\5&-7&0\end{pmatrix}
\]
The last matrix has the single nonzero entry \(1\) in column \(3\), at position \((1,3)\); expanding
along column \(3\) (Theorem~\ref{r5d-thm-expansion}), its determinant is \((-1)^{1+3}\det\left(\begin{smallmatrix}2&-3\\5&-7\end{smallmatrix}\right)=-14+15=1\).
Hence \(\det A=-5\cdot1=-5\).
\end{itemize}

(\(B\))
\begin{itemize}
\item (i) Along the first row,
\(\det B=2\det\left(\begin{smallmatrix}-3&-5\\3&9\end{smallmatrix}\right)
-6\det\left(\begin{smallmatrix}-2&-5\\2&9\end{smallmatrix}\right)
+5\det\left(\begin{smallmatrix}-2&-3\\2&3\end{smallmatrix}\right)
=2(-27+15)-6(-18+10)+5(-6+6)=-24+48+0=24\).
\item (ii)
\[
B\xrightarrow[R_3\leftarrow R_3-R_1]{R_2\leftarrow R_2+R_1}
\begin{pmatrix}2&6&5\\0&3&0\\0&-3&4\end{pmatrix}
\xrightarrow{R_3\leftarrow R_3+R_2}
\begin{pmatrix}2&6&5\\0&3&0\\0&0&4\end{pmatrix}
\]
No arrow changes the determinant (Theorem~\ref{r5d-thm-gauss} with \(t=0\)), so \(\det B=2\cdot3\cdot4=24\).
\end{itemize}
\end{solution}

\begin{exercise}\label{r5d-ex-two-dets}
Compute, by row and column operations,
\[
\det\begin{pmatrix}1&2&-1&3\\2&1&0&3\\3&1&2&3\\1&0&1&3\end{pmatrix}\qquad\text{and}\qquad
\det\begin{pmatrix}1&2&0&1\\2&5&1&2\\-1&-2&3&1\\1&2&6&7\end{pmatrix}
\]
\end{exercise}
\begin{solution}
Call the two matrices \(A\) and \(B\).

(\(A\)) Column \(4\) is \(3\) times \(\one_4\). Taking the factor \(3\) out of column \(4\)
(Definition~\ref{r5d-def-det}(a) for columns, Theorem~\ref{r5d-thm-transpose}) and then subtracting row \(1\)
from rows \(2\), \(3\), \(4\) (Definition~\ref{r5d-def-det}(b)),
\[
\det A=3\det\begin{pmatrix}1&2&-1&1\\2&1&0&1\\3&1&2&1\\1&0&1&1\end{pmatrix}
=3\det\begin{pmatrix}1&2&-1&1\\1&-1&1&0\\2&-1&3&0\\0&-2&2&0\end{pmatrix}
\]
The last matrix has the single nonzero entry \(1\) in column \(4\), at position \((1,4)\). Expanding
along column \(4\) and then the \(3\times3\) determinant along its first row (Theorem~\ref{r5d-thm-expansion}),
\[
\det A=3\cdot(-1)^{1+4}\det\begin{pmatrix}1&-1&1\\2&-1&3\\0&-2&2\end{pmatrix}
=-3\bigl(1(-2+6)+1(4-0)+1(-4-0)\bigr)=-3\cdot4=-12
\]

(\(B\))
\[
B\xrightarrow[R_4\leftarrow R_4-R_1]{R_2\leftarrow R_2-2R_1,\ R_3\leftarrow R_3+R_1}
\begin{pmatrix}1&2&0&1\\0&1&1&0\\0&0&3&2\\0&0&6&6\end{pmatrix}
\xrightarrow{R_4\leftarrow R_4-2R_3}
\begin{pmatrix}1&2&0&1\\0&1&1&0\\0&0&3&2\\0&0&0&2\end{pmatrix}
\]
No arrow changes the determinant (Theorem~\ref{r5d-thm-gauss} with \(t=0\)), so \(\det B=1\cdot1\cdot3\cdot2=6\).
\end{solution}

\begin{exercise}\label{r5d-dr-det-rules}
Let \(A,B\in\R^{4\times4}\) with \(\det A=2\) and \(\det B=-3\). Find
\begin{enumerate}
\item \(\det(AB)\), \(\det A'\) and \(\det(A'B)\);
\item \(\det(2A)\), \(\det(-B)\) and \(\det(-2AB)\);
\item \(\det A^{-1}\), \(\det(A^{-1}B)\) and \(\det(B^{-1}AB)\);
\item \(\det(A^3B^2)\), and the real numbers \(c\) with \(\det(cA)=162\).
\end{enumerate}
\end{exercise}
\begin{solution}
(a) By Theorem~\ref{r5d-thm-product}(b), \(\det(AB)=\det A\cdot\det B=2\cdot(-3)=-6\). By
Theorem~\ref{r5d-thm-transpose}, \(\det A'=\det A=2\), so \(\det(A'B)=\det A'\cdot\det B=-6\).

(b) By Theorem~\ref{r5d-thm-product}(c) with \(n=4\), \(\det(2A)=2^4\det A=32\) and
\(\det(-B)=(-1)^4\det B=-3\). With (a) of this exercise, \(\det(-2AB)=(-2)^4\det(AB)=16\cdot(-6)=-96\).

(c) \(\det A\ne0\) and \(\det B\ne0\), so \(A\) and \(B\) are invertible (Theorem~\ref{r5d-thm-product}(a)). By
Theorem~\ref{r5d-thm-product}(b),(c),
\[
\begin{gathered}
\det A^{-1}=\frac1{\det A}=\frac12,\qquad\det(A^{-1}B)=\det A^{-1}\cdot\det B=\frac12\cdot(-3)=-\frac32\\
\det(B^{-1}AB)=\det B^{-1}\cdot\det A\cdot\det B=-\frac13\cdot2\cdot(-3)=2
\end{gathered}
\]

(d) By Theorem~\ref{r5d-thm-product}(b),(c), \(\det(A^3B^2)=(\det A)^3(\det B)^2=8\cdot9=72\). Also
\(\det(cA)=c^4\det A=2c^4\), and \(2c^4=162\) gives \(c^4=81\), that is, \((c^2-9)(c^2+9)=0\). As \(c^2+9>0\),
\(c^2=9\): \(c=3\) or \(c=-3\).
\end{solution}

\begin{exercise}\label{r5d-ex-det-rules}
Let \(A,B\in\R^{4\times4}\) with \(\det A=2\) and \(\det B=-3\). Compute \(\det(3A)\),
\(\det(A'B^{-1})\), \(\det(A^2B')\), \(\det(B^{-1}AB)\), \(\det(-B)\) and \(\det(\operatorname{adj}A)\).
\end{exercise}
\begin{solution}
By Theorems~\ref{r5d-thm-product}(b),(c) and~\ref{r5d-thm-transpose} (\(B\) is invertible since
\(\det B\ne0\)):
\[
\begin{aligned}
\det(3A)&=3^4\cdot2=162,&\det(A'B^{-1})&=\det A\cdot\frac1{\det B}=-\frac23\\
\det(A^2B')&=2^2\cdot(-3)=-12,&\det(B^{-1}AB)&=\frac1{-3}\cdot2\cdot(-3)=2\\
\det(-B)&=(-1)^4(-3)=-3
\end{aligned}
\]
By Theorem~\ref{r5d-thm-adjugate}, \(A\operatorname{adj}A=2I_4\), so
\(\det A\cdot\det(\operatorname{adj}A)=\det(2I_4)=2^4\) (Theorems~\ref{r5d-thm-product}(b)
and~\ref{r5d-thm-basic}(c)) and \(\det(\operatorname{adj}A)=16/2=8\).
\end{solution}

\begin{exercise}\label{r5d-dr-cofactor}
Let
\[
A=\begin{pmatrix}3&1&0&0\\2&-1&4&1\\1&2&-2&0\\0&1&0&2\end{pmatrix},\qquad
B=\begin{pmatrix}2&-1&2&3\\-2&3&0&-2\\2&0&0&-2\\1&0&3&2\end{pmatrix}
\]
\begin{enumerate}
\item Compute \(\det A\) by expansion along the third column, and check the value by expansion along the first row;
expand each \(3\times3\) minor along a row or a column that contains a zero.
\item Compute \(\det B\) by expansion along the third row, expanding each \(3\times3\) minor along a row or a column
that contains a zero.
\item Check the value of \(\det B\) by expansion along the second column.
\end{enumerate}
\end{exercise}
\begin{solution}
(a) Column \(3\) of \(A\) is \((0,4,-2,0)'\), so by Theorem~\ref{r5d-thm-expansion}, \(\det A=4c_{23}+(-2)c_{33}\).
Expanding the first \(3\times3\) minor along its third column \((0,0,2)'\) and the second along its first row
\((3,1,0)\),
\[
\begin{aligned}
c_{23}&=(-1)^{2+3}\det\begin{pmatrix}3&1&0\\1&2&0\\0&1&2\end{pmatrix}=-2\det\begin{pmatrix}3&1\\1&2\end{pmatrix}=-2\cdot5=-10\\
c_{33}&=(-1)^{3+3}\det\begin{pmatrix}3&1&0\\2&-1&1\\0&1&2\end{pmatrix}
=3\det\begin{pmatrix}-1&1\\1&2\end{pmatrix}-1\cdot\det\begin{pmatrix}2&1\\0&2\end{pmatrix}=3\cdot(-3)-4=-13
\end{aligned}
\]
Hence \(\det A=4\cdot(-10)+(-2)\cdot(-13)=-40+26=-14\). Along the first row \((3,1,0,0)\), \(\det A=3c_{11}+1\cdot c_{12}\);
expanding the first minor along its second row \((2,-2,0)\) and the second along its third row \((0,0,2)\),
\[
\begin{aligned}
c_{11}&=(-1)^{1+1}\det\begin{pmatrix}-1&4&1\\2&-2&0\\1&0&2\end{pmatrix}
=-2\det\begin{pmatrix}4&1\\0&2\end{pmatrix}+(-2)\det\begin{pmatrix}-1&1\\1&2\end{pmatrix}=-16+6=-10\\
c_{12}&=(-1)^{1+2}\det\begin{pmatrix}2&4&1\\1&-2&0\\0&0&2\end{pmatrix}=-2\det\begin{pmatrix}2&4\\1&-2\end{pmatrix}=-2\cdot(-8)=16
\end{aligned}
\]
So \(\det A=3\cdot(-10)+1\cdot16=-14\), as before.

(b)
\begin{itemize}
\item Row \(3\) of \(B\) is \((2,0,0,-2)\), so only the cofactors \(c_{31}\) and \(c_{34}\) are needed:
\(\det B=2c_{31}+(-2)c_{34}\), where
\[
c_{31}=(-1)^{3+1}\det\begin{pmatrix}-1&2&3\\3&0&-2\\0&3&2\end{pmatrix},\qquad
c_{34}=(-1)^{3+4}\det\begin{pmatrix}2&-1&2\\-2&3&0\\1&0&3\end{pmatrix}
\]
\item The first minor, along its first column \((-1,3,0)'\):
\((-1)\det\left(\begin{smallmatrix}0&-2\\3&2\end{smallmatrix}\right)-3\det\left(\begin{smallmatrix}2&3\\3&2\end{smallmatrix}\right)+0
=-1\cdot6-3\cdot(-5)=9\). So \(c_{31}=9\).
\item The second minor, along its second row \((-2,3,0)\):
\(-(-2)\det\left(\begin{smallmatrix}-1&2\\0&3\end{smallmatrix}\right)+3\det\left(\begin{smallmatrix}2&2\\1&3\end{smallmatrix}\right)-0
=2\cdot(-3)+3\cdot4=6\). So \(c_{34}=-6\).
\item Hence \(\det B=2\cdot9+(-2)\cdot(-6)=18+12=30\).
\end{itemize}

(c) Column \(2\) of \(B\) is \((-1,3,0,0)'\), so \(\det B=(-1)c_{12}+3c_{22}\).
\begin{itemize}
\item The two minors are
\[
B(1|2)=\begin{pmatrix}-2&0&-2\\2&0&-2\\1&3&2\end{pmatrix},\qquad B(2|2)=\begin{pmatrix}2&2&3\\2&0&-2\\1&3&2\end{pmatrix}
\]
\item Along the second column \((0,0,3)'\) of \(B(1|2)\),
\(\det B(1|2)=3\cdot(-1)^{3+2}\det\left(\begin{smallmatrix}-2&-2\\2&-2\end{smallmatrix}\right)=-3\cdot8=-24\),
so \(c_{12}=(-1)^{1+2}\cdot(-24)=24\).
\item Along the second row \((2,0,-2)\) of \(B(2|2)\),
\(\det B(2|2)=-2\det\left(\begin{smallmatrix}2&3\\3&2\end{smallmatrix}\right)+0-(-2)\det\left(\begin{smallmatrix}2&2\\1&3\end{smallmatrix}\right)
=-2\cdot(-5)+2\cdot4=18\), so \(c_{22}=(-1)^{2+2}\cdot18=18\).
\item Hence \(\det B=(-1)\cdot24+3\cdot18=30\), as in (b).
\end{itemize}
\end{solution}

\begin{exercise}[\citeat{harville-solutions}{Ch.~13, Exercise~5(b)}]\label{r5d-ex-four-by-four}
Compute the determinant of
\[
A=\begin{pmatrix}0&4&0&5\\1&0&-1&2\\0&3&0&-2\\0&0&-6&0\end{pmatrix}
\]
(a) by repeated cofactor expansion: expand \(\det A\) in terms of determinants of \(3\times3\) matrices, these in terms
of determinants of \(2\times2\) matrices, and finally these in terms of determinants of \(1\times1\) matrices; (b) by
row reduction to triangular form.
\end{exercise}
\begin{solution}
(a)
\begin{itemize}
\item Column \(1\) has the single nonzero entry \(a_{21}=1\); expanding along column \(1\)
(Theorem~\ref{r5d-thm-expansion}),
\[
\det A=1\cdot(-1)^{2+1}\det\begin{pmatrix}4&0&5\\3&0&-2\\0&-6&0\end{pmatrix}
\]
\item Row \(3\) of this \(3\times3\) matrix has the single nonzero entry \(-6\), in column \(2\):
\[
\det\begin{pmatrix}4&0&5\\3&0&-2\\0&-6&0\end{pmatrix}=(-6)(-1)^{3+2}\det\begin{pmatrix}4&5\\3&-2\end{pmatrix}
\]
\item Expanding the \(2\times2\) determinant along its first row,
\(4(-1)^{1+1}\det(-2)+5(-1)^{1+2}\det(3)=-8-15=-23\).
\item Hence \(\det A=(-1)\cdot6\cdot(-23)=138\).
\end{itemize}
(b)
\[
\begin{gathered}
A\xrightarrow{R_1\leftrightarrow R_2}
\begin{pmatrix}1&0&-1&2\\0&4&0&5\\0&3&0&-2\\0&0&-6&0\end{pmatrix}
\xrightarrow{R_2\leftarrow R_2-R_3}
\begin{pmatrix}1&0&-1&2\\0&1&0&7\\0&3&0&-2\\0&0&-6&0\end{pmatrix}
\\
\xrightarrow{R_3\leftarrow R_3-3R_2}
\begin{pmatrix}1&0&-1&2\\0&1&0&7\\0&0&0&-23\\0&0&-6&0\end{pmatrix}
\xrightarrow{R_3\leftrightarrow R_4}
\begin{pmatrix}1&0&-1&2\\0&1&0&7\\0&0&-6&0\\0&0&0&-23\end{pmatrix}
\end{gathered}
\]
There are two interchanges, and the additions do not change the determinant (Theorem~\ref{r5d-thm-gauss}, \(t=2\)); so
\(\det A=(-1)^2\cdot1\cdot1\cdot(-6)\cdot(-23)=138\).
\end{solution}

\begin{exercise}\label{r5d-dr-parameter}
For \(t\in\R\) let
\[
A(t)=\begin{pmatrix}1&1&1&1\\1&t&1&2-t\\1&1&t&2-t\\1&t&t&5-t\end{pmatrix}
\]
\begin{enumerate}
\item Compute \(\det A(t)\) by row operations, and factor it.
\item For which \(t\) is \(A(t)\) invertible?
\item Find \(\rank A(t)\) for each \(t\) for which \(A(t)\) is not invertible.
\end{enumerate}
\end{exercise}
\begin{solution}
(a)
\[
\begin{gathered}
A(t)\xrightarrow[R_4\leftarrow R_4-R_1]{R_2\leftarrow R_2-R_1,\ R_3\leftarrow R_3-R_1}
\begin{pmatrix}1&1&1&1\\0&t-1&0&1-t\\0&0&t-1&1-t\\0&t-1&t-1&4-t\end{pmatrix}
\xrightarrow{R_4\leftarrow R_4-R_2}
\begin{pmatrix}1&1&1&1\\0&t-1&0&1-t\\0&0&t-1&1-t\\0&0&t-1&3\end{pmatrix}
\\
\xrightarrow{R_4\leftarrow R_4-R_3}
\begin{pmatrix}1&1&1&1\\0&t-1&0&1-t\\0&0&t-1&1-t\\0&0&0&t+2\end{pmatrix}
\end{gathered}
\]
using \((4-t)-(1-t)=3\) and \(3-(1-t)=t+2\). Every operation adds a multiple of a row to another row, for every value
of \(t\), and there are no interchanges; the last matrix is upper triangular. By Theorem~\ref{r5d-thm-gauss},
\[
\det A(t)=1\cdot(t-1)\cdot(t-1)\cdot(t+2)=(t-1)^2(t+2)
\]

(b) By Theorem~\ref{r5d-thm-product}(a), \(A(t)\) is invertible if and only if \(\det A(t)\ne0\), that is,
\(t\ne1\) and \(t\ne-2\).

(c) The operations of (a) do not depend on \(t\), so for \(t=1\) and for \(t=-2\) they take \(A(t)\) to the last
matrix of (a) with that value of \(t\).
\begin{itemize}
\item For \(t=1\) its rows \(2\) and \(3\) are zero, and row \(4\) is not (the second and third arrows of (a) subtract
these zero rows and change nothing). The next pivot column is column \(4\) (Step 1 of
Definition~\ref{r2-def-gauss-jordan}), and Step 2 interchanges rows \(2\) and \(4\):
\[
\begin{pmatrix}1&1&1&1\\0&0&0&0\\0&0&0&0\\0&0&0&3\end{pmatrix}
\xrightarrow{R_2\leftrightarrow R_4}
\begin{pmatrix}1&1&1&1\\0&0&0&3\\0&0&0&0\\0&0&0&0\end{pmatrix}
\]
\item For \(t=-2\) it is
\[
\begin{pmatrix}1&1&1&1\\0&-3&0&3\\0&0&-3&3\\0&0&0&0\end{pmatrix}
\]
\item The matrix after the interchange (\(t=1\)) and this matrix (\(t=-2\)) are the results of forward elimination of
\(A(1)\) and \(A(-2)\) (Definition~\ref{r2-def-gauss-jordan}, with the scaling of Step 3 postponed), and
\(\rank A(t)\) is their number of nonzero rows (Definition~\ref{def:r4-rank} and Remark~\ref{rem:r4-echelon}). So
\(\rank A(1)=2\) and \(\rank A(-2)=3\).
\end{itemize}
\end{solution}

\begin{exercise}[\citeat{abadir-magnus}{Exercise~4.31(a)}]\label{r5d-ex-alpha}
Let \(\alpha\in\R\) and
\[
A=\begin{pmatrix}1&\alpha&0\\4&5&3\\1&0&2\end{pmatrix}
\]
For which \(\alpha\) is \(\det A=0\)? Show that the columns of \(A\) are linearly dependent in that case.
\end{exercise}
\begin{solution}
Expanding along the first row (Theorem~\ref{r5d-thm-expansion}),
\(\det A=1\det\left(\begin{smallmatrix}5&3\\0&2\end{smallmatrix}\right)
-\alpha\det\left(\begin{smallmatrix}4&3\\1&2\end{smallmatrix}\right)=10-5\alpha\), so \(\det A=0\) exactly
when \(\alpha=2\) (and \(A\) is invertible exactly when \(\alpha\ne2\),
Theorem~\ref{r5d-thm-product}(a)). For \(\alpha=2\) and \(x=(2,-1,-1)'\),
\(Ax=(2-2+0,\ 8-5-3,\ 2-0-2)'=0\); that is, \(2a_1-a_2-a_3=0\), a nontrivial relation among the columns.
\end{solution}

\begin{exercise}\label{r5d-dr-adjugate-cramer}
Let
\[
A=\begin{pmatrix}3&-1&2&0\\1&0&1&0\\4&-2&3&0\\0&0&1&1\end{pmatrix},\qquad b=\begin{pmatrix}-1\\0\\-3\\0\end{pmatrix}
\]
\begin{enumerate}
\item Compute the sixteen cofactors \(c_{ij}\) of \(A\), and \(\det A\) by expansion along the second row.
\item Find \(\operatorname{adj}A\) and \(A^{-1}\) (Theorem~\ref{r5d-thm-adjugate}), and check that \(AA^{-1}=I_4\).
\item Solve \(Ax=b\) by Cramer's rule (Theorem~\ref{r5d-thm-cramer}), and check the solution with \(A^{-1}b\).
\end{enumerate}
\end{exercise}
\begin{solution}
(a) By Definition~\ref{r5d-def-cofactor},
\[
\begin{aligned}
c_{11}&=\det\left(\begin{smallmatrix}0&1&0\\-2&3&0\\0&1&1\end{smallmatrix}\right)=2,&
c_{12}&=-\det\left(\begin{smallmatrix}1&1&0\\4&3&0\\0&1&1\end{smallmatrix}\right)=1\\
c_{13}&=\det\left(\begin{smallmatrix}1&0&0\\4&-2&0\\0&0&1\end{smallmatrix}\right)=-2,&
c_{14}&=-\det\left(\begin{smallmatrix}1&0&1\\4&-2&3\\0&0&1\end{smallmatrix}\right)=2\\
c_{21}&=-\det\left(\begin{smallmatrix}-1&2&0\\-2&3&0\\0&1&1\end{smallmatrix}\right)=-1,&
c_{22}&=\det\left(\begin{smallmatrix}3&2&0\\4&3&0\\0&1&1\end{smallmatrix}\right)=1\\
c_{23}&=-\det\left(\begin{smallmatrix}3&-1&0\\4&-2&0\\0&0&1\end{smallmatrix}\right)=2,&
c_{24}&=\det\left(\begin{smallmatrix}3&-1&2\\4&-2&3\\0&0&1\end{smallmatrix}\right)=-2\\
c_{31}&=\det\left(\begin{smallmatrix}-1&2&0\\0&1&0\\0&1&1\end{smallmatrix}\right)=-1,&
c_{32}&=-\det\left(\begin{smallmatrix}3&2&0\\1&1&0\\0&1&1\end{smallmatrix}\right)=-1\\
c_{33}&=\det\left(\begin{smallmatrix}3&-1&0\\1&0&0\\0&0&1\end{smallmatrix}\right)=1,&
c_{34}&=-\det\left(\begin{smallmatrix}3&-1&2\\1&0&1\\0&0&1\end{smallmatrix}\right)=-1\\
c_{41}&=-\det\left(\begin{smallmatrix}-1&2&0\\0&1&0\\-2&3&0\end{smallmatrix}\right)=0,&
c_{42}&=\det\left(\begin{smallmatrix}3&2&0\\1&1&0\\4&3&0\end{smallmatrix}\right)=0\\
c_{43}&=-\det\left(\begin{smallmatrix}3&-1&0\\1&0&0\\4&-2&0\end{smallmatrix}\right)=0,&
c_{44}&=\det\left(\begin{smallmatrix}3&-1&2\\1&0&1\\4&-2&3\end{smallmatrix}\right)=1
\end{aligned}
\]
\begin{itemize}
\item Each of the twelve minors with \(i\leqslant3\) has a third row or a third column whose only nonzero entry is the
\(1\) in position \((3,3)\); expanding along it (Theorem~\ref{r5d-thm-expansion}) leaves the \(2\times2\) determinant of
its first two rows and columns. For instance, \(\det\left(\begin{smallmatrix}0&1\\-2&3\end{smallmatrix}\right)=0+2=2\) for
\(c_{11}\).
\item The minors \(\det A(4|j)\) with \(j\leqslant3\) have a zero third column, so they are \(0\)
(Theorem~\ref{r5d-thm-basic}(b) for columns, Theorem~\ref{r5d-thm-transpose}).
\item \(c_{44}=\det A(4|4)\), along its second row \((1,0,1)\), is
\(-\det\left(\begin{smallmatrix}-1&2\\-2&3\end{smallmatrix}\right)-\det\left(\begin{smallmatrix}3&-1\\4&-2\end{smallmatrix}\right)=-1+2=1\).
\item Along row \(2\) of \(A\) (Theorem~\ref{r5d-thm-expansion}), \(\det A=1\cdot c_{21}+0\cdot c_{22}+1\cdot c_{23}+0\cdot c_{24}
=-1+2=1\).
\end{itemize}

(b) Since \(\det A=1\), Theorem~\ref{r5d-thm-adjugate} gives
\[
A^{-1}=\operatorname{adj}A=\begin{pmatrix}c_{11}&c_{21}&c_{31}&c_{41}\\c_{12}&c_{22}&c_{32}&c_{42}\\c_{13}&c_{23}&c_{33}&c_{43}\\c_{14}&c_{24}&c_{34}&c_{44}\end{pmatrix}
=\begin{pmatrix}2&-1&-1&0\\1&1&-1&0\\-2&2&1&0\\2&-2&-1&1\end{pmatrix}
\]
Check: \(\row_1(A)A^{-1}=3(2,-1,-1,0)-(1,1,-1,0)+2(-2,2,1,0)=(1,0,0,0)\),
\(\row_2(A)A^{-1}=(2,-1,-1,0)+(-2,2,1,0)=(0,1,0,0)\), \(\row_3(A)A^{-1}=4(2,-1,-1,0)-2(1,1,-1,0)+3(-2,2,1,0)=(0,0,1,0)\)
and \(\row_4(A)A^{-1}=(-2,2,1,0)+(2,-2,-1,1)=(0,0,0,1)\) (Theorem~\ref{r1-thm-views}(c)).

(c)
\begin{itemize}
\item As \(\det A=1\ne0\), Theorem~\ref{r5d-thm-cramer} applies. Replacing column \(j\) of \(A\) by \(b\) gives
\[
\begin{gathered}
A_1=\begin{pmatrix}-1&-1&2&0\\0&0&1&0\\-3&-2&3&0\\0&0&1&1\end{pmatrix},\qquad
A_2=\begin{pmatrix}3&-1&2&0\\1&0&1&0\\4&-3&3&0\\0&0&1&1\end{pmatrix}\\
A_3=\begin{pmatrix}3&-1&-1&0\\1&0&0&0\\4&-2&-3&0\\0&0&0&1\end{pmatrix},\qquad
A_4=\begin{pmatrix}3&-1&2&-1\\1&0&1&0\\4&-2&3&-3\\0&0&1&0\end{pmatrix}
\end{gathered}
\]
\item Each second row contains zeros; expanding along it (Theorem~\ref{r5d-thm-expansion}),
\[
\begin{aligned}
\det A_1&=1\cdot(-1)^{2+3}\det\begin{pmatrix}-1&-1&0\\-3&-2&0\\0&0&1\end{pmatrix}=-(2-3)=1\\
\det A_2&=1\cdot(-1)^{2+1}\det\begin{pmatrix}-1&2&0\\-3&3&0\\0&1&1\end{pmatrix}
+1\cdot(-1)^{2+3}\det\begin{pmatrix}3&-1&0\\4&-3&0\\0&0&1\end{pmatrix}=-(-3+6)-(-9+4)=2\\
\det A_3&=1\cdot(-1)^{2+1}\det\begin{pmatrix}-1&-1&0\\-2&-3&0\\0&0&1\end{pmatrix}=-(3-2)=-1\\
\det A_4&=1\cdot(-1)^{2+1}\det\begin{pmatrix}-1&2&-1\\-2&3&-3\\0&1&0\end{pmatrix}
+1\cdot(-1)^{2+3}\det\begin{pmatrix}3&-1&-1\\4&-2&-3\\0&0&0\end{pmatrix}=-(-1)-0=1
\end{aligned}
\]
Here the \(3\times3\) determinants with the third column \((0,0,1)'\) are \(2\times2\) determinants, as in (a); the last
one has a zero row, so it is \(0\) (Theorem~\ref{r5d-thm-basic}(b)); and the one before it, along its third row
\((0,1,0)\), is \((-1)^{3+2}\det\left(\begin{smallmatrix}-1&-1\\-2&-3\end{smallmatrix}\right)=-(3-2)=-1\).
\item Hence \(x_1=\frac11\), \(x_2=\frac21\), \(x_3=\frac{-1}1\), \(x_4=\frac11\), that is, \(x=(1,2,-1,1)'\). Check:
\(A^{-1}b=(-2+0+3+0,\ -1+0+3+0,\ 2+0-3+0,\ -2+0+3+0)'=(1,2,-1,1)'\).
\end{itemize}
\end{solution}

\begin{exercise}[cf.\ \citeat{harville-solutions}{Ch.~13, Exercise~9}]\label{r5d-ex-cofactors}
Let
\[
T=\begin{pmatrix}2&-1&0&0\\-1&2&-1&0\\0&-1&2&-1\\0&0&-1&1\end{pmatrix}
\]
\begin{enumerate}
\item Compute all cofactors of \(T\).
\item Compute \(\det T\) by expansion along the second row, and check the answer by expansion along the
fourth column.
\item Use Theorem~\ref{r5d-thm-adjugate} to show that \(T^{-1}\) is the matrix \(K\) with entries
\(k_{ij}=\min(i,j)\).
\item For \(b=(1,2,3,4)'\), find \(x_2\) in \(Tx=b\) by Cramer's rule, and check it with \(x=Kb\).
\end{enumerate}
\end{exercise}
\begin{solution}
(a) \(T\) is symmetric, so \(c_{ji}=c_{ij}\) (Remark~\ref{r5d-rem-symmetric-cofactors}). The ten cofactors with \(i\leqslant j\) are (each minor is expanded along a row or
column with at most two nonzero entries, Theorem~\ref{r5d-thm-expansion})
\[
\begin{aligned}
c_{11}&=\det\left(\begin{smallmatrix}2&-1&0\\-1&2&-1\\0&-1&1\end{smallmatrix}\right)=2\cdot1+1\cdot(-1)=1,&
c_{12}&=-\det\left(\begin{smallmatrix}-1&-1&0\\0&2&-1\\0&-1&1\end{smallmatrix}\right)=-(-1)=1\\
c_{13}&=\det\left(\begin{smallmatrix}-1&2&0\\0&-1&-1\\0&0&1\end{smallmatrix}\right)=(-1)(-1)(1)=1,&
c_{14}&=-\det\left(\begin{smallmatrix}-1&2&-1\\0&-1&2\\0&0&-1\end{smallmatrix}\right)=-(-1)=1\\
c_{22}&=\det\left(\begin{smallmatrix}2&0&0\\0&2&-1\\0&-1&1\end{smallmatrix}\right)=2(2-1)=2,&
c_{23}&=-\det\left(\begin{smallmatrix}2&-1&0\\0&-1&-1\\0&0&1\end{smallmatrix}\right)=-(-2)=2\\
c_{24}&=\det\left(\begin{smallmatrix}2&-1&0\\0&-1&2\\0&0&-1\end{smallmatrix}\right)=2(-1)(-1)=2,&
c_{33}&=\det\left(\begin{smallmatrix}2&-1&0\\-1&2&0\\0&0&1\end{smallmatrix}\right)=4-1=3\\
c_{34}&=-\det\left(\begin{smallmatrix}2&-1&0\\-1&2&-1\\0&0&-1\end{smallmatrix}\right)=-(-3)=3,&
c_{44}&=\det\left(\begin{smallmatrix}2&-1&0\\-1&2&-1\\0&-1&2\end{smallmatrix}\right)=2\cdot3+1\cdot(-2)=4
\end{aligned}
\]
(the triangular minors by Theorem~\ref{r5d-thm-basic}(c)). So \(c_{ij}=\min(i,j)\) for all
\(i,j\).

(b) Along row \(2\): \(\det T=(-1)c_{21}+2c_{22}+(-1)c_{23}=-1+4-2=1\). Along column \(4\):
\(\det T=(-1)c_{34}+1\cdot c_{44}=-3+4=1\) (Theorem~\ref{r5d-thm-expansion}).

(c) Since \(\det T=1\), Theorem~\ref{r5d-thm-adjugate} gives \(T^{-1}=\operatorname{adj}T=(c_{ij})'\),
and \((c_{ij})'=(c_{ji})=(\min(i,j))=K\) by (a):
\[
T^{-1}=\begin{pmatrix}1&1&1&1\\1&2&2&2\\1&2&3&3\\1&2&3&4\end{pmatrix}
\]
Check: \(\row_4(T)K=-(1,2,3,3)+(1,2,3,4)=(0,0,0,1)\) and
\(\row_2(T)K=-(1,1,1,1)+2(1,2,2,2)-(1,2,3,3)=(0,1,0,0)\).

(d) \(T_2\) is \(T\) with column \(2\) replaced by \(b\). Expanding along column \(1\) (Theorem~\ref{r5d-thm-expansion}),
\[
\begin{aligned}
\det T_2&=\det\begin{pmatrix}2&1&0&0\\-1&2&-1&0\\0&3&2&-1\\0&4&-1&1\end{pmatrix}
=2\det\begin{pmatrix}2&-1&0\\3&2&-1\\4&-1&1\end{pmatrix}+\det\begin{pmatrix}1&0&0\\3&2&-1\\4&-1&1\end{pmatrix}\\
&=2\bigl(2(2-1)+1\cdot(3+4)\bigr)+1\cdot(2-1)=19
\end{aligned}
\]
(the \(3\times3\) determinants along their first rows), and \(\det T=1\) by (b), so \(x_2=19\) (Theorem~\ref{r5d-thm-cramer}). Check: \(x_2=\row_2(K)b=1+4+6+8=19\);
indeed \(x=Kb=(10,19,26,30)'\).
\end{solution}

\begin{exercise}[cf.\ \citeat{abadir-magnus}{Exercise~6.29}]\label{r5d-ex-cramer}
Use Cramer's rule (Theorem~\ref{r5d-thm-cramer}) to solve the first system of Exercise~\ref{r2-ex-two-systems},
\(Ax=b\) with
\[
A=\begin{pmatrix}1&-3&-2&-1\\-1&2&1&1\\-1&2&2&1\\0&1&1&1\end{pmatrix},\qquad
b=\begin{pmatrix}1\\-1\\-2\\1\end{pmatrix}
\]
\end{exercise}
\begin{solution}
\begin{itemize}
\item Let \(A_j\) be \(A\) with column \(j\) replaced by \(b\). Subtracting row \(2\) from row \(3\) does not change a
determinant (Definition~\ref{r5d-def-det}(b)). It turns row \(3\) of \(A\) into \((0,0,1,0)\), and row \(3\) of \(A_j\)
into the same row with its \(j\)th entry replaced by \(b_3-b_2=-1\): the five matrices become
\[
\begin{gathered}
\begin{pmatrix}1&-3&-2&-1\\-1&2&1&1\\0&0&1&0\\0&1&1&1\end{pmatrix}\ (A),\qquad
\begin{pmatrix}1&-3&-2&-1\\-1&2&1&1\\-1&0&1&0\\1&1&1&1\end{pmatrix}\ (A_1),\qquad
\begin{pmatrix}1&1&-2&-1\\-1&-1&1&1\\0&-1&1&0\\0&1&1&1\end{pmatrix}\ (A_2)\\
\begin{pmatrix}1&-3&1&-1\\-1&2&-1&1\\0&0&-1&0\\0&1&1&1\end{pmatrix}\ (A_3),\qquad
\begin{pmatrix}1&-3&-2&1\\-1&2&1&-1\\0&0&1&-1\\0&1&1&1\end{pmatrix}\ (A_4)
\end{gathered}
\]
\item Expanding along the new row \(3\) (Theorem~\ref{r5d-thm-expansion}), and then the \(3\times3\) determinants along
their first rows,
\[
\begin{aligned}
\det A&=\det\begin{pmatrix}1&-3&-1\\-1&2&1\\0&1&1\end{pmatrix}=1(2-1)+3(-1-0)-1(-1-0)=-1\\
\det A_1&=-\det\begin{pmatrix}-3&-2&-1\\2&1&1\\1&1&1\end{pmatrix}+\det\begin{pmatrix}1&-3&-1\\-1&2&1\\1&1&1\end{pmatrix}
=-(0+2-1)+(1-6+3)=-3\\
\det A_2&=\det\begin{pmatrix}1&-2&-1\\-1&1&1\\0&1&1\end{pmatrix}+\det\begin{pmatrix}1&1&-1\\-1&-1&1\\0&1&1\end{pmatrix}
=(0-2+1)+(-2+1+1)=-1\\
\det A_3&=-\det\begin{pmatrix}1&-3&-1\\-1&2&1\\0&1&1\end{pmatrix}=1\\
\det A_4&=\det\begin{pmatrix}1&-3&1\\-1&2&-1\\0&1&1\end{pmatrix}+\det\begin{pmatrix}1&-3&-2\\-1&2&1\\0&1&1\end{pmatrix}
=(3-3-1)+(1-3+2)=-1
\end{aligned}
\]
(the minor of \(\det A_3\) is the one of \(\det A\), as columns \(1\), \(2\), \(4\) of \(A_3\) are those of \(A\)).
\item So \(x=\bigl(\frac{-3}{-1},\ \frac{-1}{-1},\ \frac1{-1},\ \frac{-1}{-1}\bigr)'=(3,1,-1,1)'\)
(Theorem~\ref{r5d-thm-cramer}), the solution found by elimination in Exercise~\ref{r2-ex-two-systems}.
\end{itemize}
\end{solution}

\begin{exercise}\label{r5d-ex-block-tri-det}
Let
\[
M=\left(\begin{array}{cc|ccc}3&1&4&-1&7\\5&2&2&9&3\\\hline0&0&2&1&0\\0&0&1&2&1\\0&0&0&1&2\end{array}\right)
\]
Compute \(\det M\), \(\det M^3\) and \(\det(2M^{-1})\).
\end{exercise}
\begin{solution}
\(M=\left(\begin{smallmatrix}A&B\\0&D\end{smallmatrix}\right)\) with
\(A=\left(\begin{smallmatrix}3&1\\5&2\end{smallmatrix}\right)\), \(\det A=6-5=1\), and
\[
D=\begin{pmatrix}2&1&0\\1&2&1\\0&1&2\end{pmatrix}
\]
whose expansion along the first row (Theorem~\ref{r5d-thm-expansion}) gives \(\det D=2(4-1)-1(2-0)=4\). By Theorem~\ref{r5d-thm-block-tri},
\(\det M=1\cdot4=4\); the block \(B\) plays no role. As \(\det M\ne0\), \(M\) is invertible, and by
Theorem~\ref{r5d-thm-product}(a),(c),
\(\det M^3=4^3=64\) and \(\det(2M^{-1})=2^5\cdot\frac14=8\).
\end{solution}

\begin{exercise}\label{r5d-ex-schur-numeric}
For \(t\in\R\) let
\[
M=\left(\begin{array}{ccc|cc}2&1&0&1&0\\1&1&0&1&1\\0&0&1&0&1\\\hline1&1&0&t&2\\0&1&1&2&t+2\end{array}\right)
=\begin{pmatrix}A&B\\C&D\end{pmatrix}
\]
Here \(A\in\R^{3\times3}\), \(B\in\R^{3\times2}\), \(C\in\R^{2\times3}\), \(D\in\R^{2\times2}\).
(a) Find \(\det M\) with the Schur complement of \(A\). (b) For which \(t\) is \(M\) invertible?
(c) Find \(\rank M\) for these exceptional values of \(t\).
\end{exercise}
\begin{solution}
(a) Here \(A=\diag(A_1,1)\) with \(A_1=\left(\begin{smallmatrix}2&1\\1&1\end{smallmatrix}\right)\),
\(B=\left(\begin{smallmatrix}1&0\\1&1\\0&1\end{smallmatrix}\right)\), \(C=B'\) and
\(D=\left(\begin{smallmatrix}t&2\\2&t+2\end{smallmatrix}\right)\). By
Theorem~\ref{r5d-thm-block-tri}, \(\det A=\det A_1\cdot1=1\), as \(\det A_1=2\cdot1-1\cdot1=1\); by
Propositions~\ref{r2-prop-inverse-2x2} and~\ref{r5b-prop-block-tri}(b),
\[
A^{-1}=\diag(A_1^{-1},1)=\begin{pmatrix}1&-1&0\\-1&2&0\\0&0&1\end{pmatrix}
\]
Then
\[
CA^{-1}=\begin{pmatrix}1&1&0\\0&1&1\end{pmatrix}\begin{pmatrix}1&-1&0\\-1&2&0\\0&0&1\end{pmatrix}=\begin{pmatrix}0&1&0\\-1&2&1\end{pmatrix},\qquad
CA^{-1}B=\begin{pmatrix}0&1&0\\-1&2&1\end{pmatrix}\begin{pmatrix}1&0\\1&1\\0&1\end{pmatrix}=\begin{pmatrix}1&1\\1&3\end{pmatrix}
\]
so the Schur complement is
\(S=D-CA^{-1}B=\left(\begin{smallmatrix}t-1&1\\1&t-1\end{smallmatrix}\right)\). By
Theorem~\ref{r5d-thm-schur-det}(a),
\(\det M=\det A\cdot\det S=(t-1)^2-1=t(t-2)\).

(b) By Theorem~\ref{r5d-thm-product}(a), \(M\) is invertible if and only if \(t\ne0\) and \(t\ne2\).

(c) For \(t=0\), \(S=\left(\begin{smallmatrix}-1&1\\1&-1\end{smallmatrix}\right)\); for \(t=2\),
\(S=\left(\begin{smallmatrix}1&1\\1&1\end{smallmatrix}\right)\). In both cases the second row of \(S\) is
a multiple of the first and \(S\ne0\), so \(\rank S=1\) (the row space is spanned by the nonzero first row;
Definition~\ref{def:r4-rank} and Proposition~\ref{r3-prop-dependence}(a)). By Theorem~\ref{r5d-thm-schur-rank}(a)
with \(p=3\), \(\rank M=3+1=4\).
\end{solution}

\begin{exercise}\label{r5d-ex-schur-nonsquare}
Let
\[
B=\begin{pmatrix}1&0\\0&1\\1&1\end{pmatrix},\qquad C=\begin{pmatrix}1&1&0\\0&1&1\end{pmatrix},\qquad
M=\begin{pmatrix}I_3&B\\C&I_2\end{pmatrix}\in\R^{5\times5}
\]
\begin{enumerate}
\item Compute \(\det M\) with the Schur complement of \(I_3\), and again with the Schur complement of
\(I_2\).
\item Compute \(M^{-1}\) with Corollary~\ref{r5b-cor-schur-inverse}(a).
\end{enumerate}
\end{exercise}
\begin{solution}
(a) The Schur complement of \(I_3\) is \(S=I_2-CI_3^{-1}B=I_2-CB\) and that of \(I_2\) is
\(I_3-BC\), with
\[
CB=\begin{pmatrix}1&1\\1&2\end{pmatrix},\quad S=\begin{pmatrix}0&-1\\-1&-1\end{pmatrix},\qquad
BC=\begin{pmatrix}1&1&0\\0&1&1\\1&2&1\end{pmatrix},\quad
I_3-BC=\begin{pmatrix}0&-1&0\\0&0&-1\\-1&-2&0\end{pmatrix}
\]
By Theorem~\ref{r5d-thm-schur-det}(a), \(\det M=\det I_3\cdot\det S=0-1=-1\). Expanding
\(\det(I_3-BC)\) along the first row (Theorem~\ref{r5d-thm-expansion}) gives
\((-1)^{1+2}(-1)\det\left(\begin{smallmatrix}0&-1\\-1&0\end{smallmatrix}\right)=0-1=-1\), so
Theorem~\ref{r5d-thm-schur-det}(b) gives again \(\det M=\det I_2\cdot\det(I_3-BC)=-1\).

(b) The Schur complement \(S=I_2-CB\) of \(I_3\) is invertible (\(\det S=-1\)), with
\(S^{-1}=\left(\begin{smallmatrix}1&-1\\-1&0\end{smallmatrix}\right)\) by the \(2\times2\) formula of
Proposition~\ref{r2-prop-inverse-2x2}. Corollary~\ref{r5b-cor-schur-inverse}(a) with \(A=I_3\) gives
\(M^{-1}=\left(\begin{smallmatrix}I_3&-B\\0&I_2\end{smallmatrix}\right)
\diag(I_3,S^{-1})\left(\begin{smallmatrix}I_3&0\\-C&I_2\end{smallmatrix}\right)\).
By the block product formula,
\[
M^{-1}=\begin{pmatrix}I_3&-BS^{-1}\\0&S^{-1}\end{pmatrix}\begin{pmatrix}I_3&0\\-C&I_2\end{pmatrix}
=\begin{pmatrix}I_3+BS^{-1}C&-BS^{-1}\\-S^{-1}C&S^{-1}\end{pmatrix}
\]
where
\[
-BS^{-1}=\begin{pmatrix}-1&1\\1&0\\0&1\end{pmatrix},\qquad
-S^{-1}C=\begin{pmatrix}-1&0&1\\1&1&0\end{pmatrix},\qquad
BS^{-1}C=-B(-S^{-1}C)=\begin{pmatrix}1&0&-1\\-1&-1&0\\0&-1&-1\end{pmatrix}
\]
Hence
\[
M^{-1}=\left(\begin{array}{ccc|cc}2&0&-1&-1&1\\-1&0&0&1&0\\0&-1&0&0&1\\\hline-1&0&1&1&-1\\1&1&0&-1&0\end{array}\right)
\]
Check: \(\row_4(M)M^{-1}=(1,1,0,1,0)M^{-1}=(2-1-1,\ 0+0+0,\ -1+0+1,\ -1+1+1,\ 1+0-1)=(0,0,0,1,0)\).
\end{solution}

\begin{exercise}\label{r5d-dr-schur}
Let
\[
M=\left(\begin{array}{cc|cc}1&1&2&2\\1&2&2&1\\\hline2&2&3&2\\2&1&2&1\end{array}\right)
=\begin{pmatrix}A&B\\C&D\end{pmatrix}
\]
\begin{enumerate}
\item Find \(D^{-1}\) and the Schur complement \(F=A-BD^{-1}C\) of \(D\).
\item Find \(\det M\) and \(\rank M\) from \(F\).
\item Find a basis of \(\mathcal N(M)\).
\end{enumerate}
\end{exercise}
\begin{solution}
(a)
\begin{itemize}
\item Here \(A=\left(\begin{smallmatrix}1&1\\1&2\end{smallmatrix}\right)\),
\(B=C=\left(\begin{smallmatrix}2&2\\2&1\end{smallmatrix}\right)\) and
\(D=\left(\begin{smallmatrix}3&2\\2&1\end{smallmatrix}\right)\). As \(\det D=3\cdot1-2\cdot2=-1\),
\(D^{-1}=\left(\begin{smallmatrix}-1&2\\2&-3\end{smallmatrix}\right)\) (the \(2\times2\) case of
Theorem~\ref{r5d-thm-adjugate}).
\item Then
\[
BD^{-1}=\begin{pmatrix}2&2\\2&1\end{pmatrix}\begin{pmatrix}-1&2\\2&-3\end{pmatrix}=\begin{pmatrix}2&-2\\0&1\end{pmatrix},\qquad
BD^{-1}C=\begin{pmatrix}2&-2\\0&1\end{pmatrix}\begin{pmatrix}2&2\\2&1\end{pmatrix}=\begin{pmatrix}0&2\\2&1\end{pmatrix}
\]
\item So \(F=A-BD^{-1}C=\left(\begin{smallmatrix}1&-1\\-1&1\end{smallmatrix}\right)\).
\end{itemize}

(b) \(\det F=1-1=0\), so \(\det M=\det D\cdot\det F=0\) (Theorem~\ref{r5d-thm-schur-det}(b)), and \(M\) is not
invertible. The second row of \(F\) is \(-1\) times the first and \(F\ne0\), so \(\rank F=1\) (the row space is
spanned by the nonzero first row; Definition~\ref{def:r4-rank} and Proposition~\ref{r3-prop-dependence}(a)); by
Theorem~\ref{r5d-thm-schur-rank}(b) with \(q=2\), \(\rank M=2+1=3\).

(c)
\begin{itemize}
\item By Proposition~\ref{r5b-prop-block-elementary}, with \(-BD^{-1}\) in the upper right block of the block
elementary matrix,
\[
L=\begin{pmatrix}I_2&-BD^{-1}\\0&I_2\end{pmatrix}\text{ is invertible, and}\quad
LM=\begin{pmatrix}A-BD^{-1}C&B-BD^{-1}D\\C&D\end{pmatrix}=\begin{pmatrix}F&0\\C&D\end{pmatrix}
\]
\item So for \(x=\left(\begin{smallmatrix}y\\z\end{smallmatrix}\right)\) with \(y,z\in\R^2\): \(Mx=0\) if and only if
\(LMx=0\), that is, \(Fy=0\) and \(Cy+Dz=0\).
\item \(Fy=0\) means \(y_1-y_2=0\), so \(y=s(1,1)'\) with \(s\in\R\). Then \(Cy=s(4,3)'\) and
\(z=-D^{-1}Cy=-s(-4+6,\ 8-9)'=s(-2,1)'\).
\item Hence \(\mathcal N(M)=\{s(1,1,-2,1)':s\in\R\}\), and \((1,1,-2,1)'\) is a basis. This agrees with
\(\dim\mathcal N(M)=4-\rank M=1\) (Theorem~\ref{thm:r4-rank-nullity}). Check:
\(M(1,1,-2,1)'=(1+1-4+2,\ 1+2-4+1,\ 2+2-6+2,\ 2+1-4+1)'=0\).
\end{itemize}
\end{solution}

\begin{exercise}\label{r5d-ex-det-lemma}
Use Theorem~\ref{r5d-thm-det-lemma} to compute
\[
\det\begin{pmatrix}2&1&1&1\\1&3&1&1\\1&1&4&1\\1&1&1&7\end{pmatrix}\qquad\text{and}\qquad
\det\begin{pmatrix}2&-1&1&-1\\2&-1&2&-2\\3&-3&4&-3\\4&-4&4&-3\end{pmatrix}
\]
\end{exercise}
\begin{solution}
The first matrix is \(D+\one_4\one_4'\) with \(D=\diag(1,2,3,6)\), \(\det D=36\) (Theorem~\ref{r5d-thm-basic}(c)) and
\(D^{-1}=\diag(1,\frac12,\frac13,\frac16)\) (Proposition~\ref{r5b-prop-block-tri}(b)). By
Theorem~\ref{r5d-thm-det-lemma},
\[
\det(D+\one\one')=\det D\cdot(1+\one'D^{-1}\one)=36\Bigl(1+1+\frac12+\frac13+\frac16\Bigr)=36\cdot3=108
\]
The second matrix is \(I_4+uv'\) with \(u=(1,2,3,4)'\) and \(v=(1,-1,1,-1)'\): its \((i,j)\) entry
is \(\delta_{ij}+u_iv_j\), where \(\delta_{ij}=1\) if \(i=j\) and \(0\) otherwise. Hence, by
Theorem~\ref{r5d-thm-det-lemma}, its determinant is \(1+v'u=1+(1-2+3-4)=-1\).
\end{solution}

\begin{exercise}\label{r5d-dr-sherman-morrison}
Let
\[
A=\begin{pmatrix}1&2&2&0\\0&1&1&0\\2&2&3&0\\0&0&1&1\end{pmatrix},\qquad
A^{-1}=\begin{pmatrix}1&-2&0&0\\2&-1&-1&0\\-2&2&1&0\\2&-2&-1&1\end{pmatrix},\qquad
u=\begin{pmatrix}0\\0\\1\\1\end{pmatrix},\qquad v=\begin{pmatrix}1\\2\\0\\0\end{pmatrix}
\]
\begin{enumerate}
\item Compute \(A^{-1}u\), \(v'A^{-1}\) and \(\det A\), and find \(\det(A+uv')\) by Theorem~\ref{r5d-thm-det-lemma}.
\item Find \((A+uv')^{-1}\) by Theorem~\ref{r5d-thm-woodbury}(b), and check the answer by multiplication.
\item For which \(s\in\R\) is \(A+s\,uv'\) invertible?
\end{enumerate}
\end{exercise}
\begin{solution}
(a)
\begin{itemize}
\item By Theorem~\ref{r1-thm-views}(b),(c), \(A^{-1}u=\col_3(A^{-1})+\col_4(A^{-1})=(0,-1,1,0)'\) and
\(v'A^{-1}=\row_1(A^{-1})+2\row_2(A^{-1})=(5,-4,-2,0)\). So \(v'A^{-1}u=5\cdot0+(-4)\cdot0+(-2)\cdot1+0\cdot1=-2\).
\item Along the fourth column, whose only nonzero entry is \(a_{44}=1\), and then along the second row \((0,1,1)\) of the
\(3\times3\) minor (Theorem~\ref{r5d-thm-expansion}),
\[
\det A=\det\begin{pmatrix}1&2&2\\0&1&1\\2&2&3\end{pmatrix}=\det\begin{pmatrix}1&2\\2&3\end{pmatrix}-\det\begin{pmatrix}1&2\\2&2\end{pmatrix}
=-1+2=1
\]
\item By Theorem~\ref{r5d-thm-det-lemma}, \(\det(A+uv')=\det A\cdot(1+v'A^{-1}u)=1\cdot(1-2)=-1\).
\end{itemize}

(b)
\begin{itemize}
\item \(1+v'A^{-1}u=-1\ne0\), so Theorem~\ref{r5d-thm-woodbury}(b) applies, with
\[
A^{-1}uv'A^{-1}=\begin{pmatrix}0\\-1\\1\\0\end{pmatrix}(5,\ -4,\ -2,\ 0)
=\begin{pmatrix}0&0&0&0\\-5&4&2&0\\5&-4&-2&0\\0&0&0&0\end{pmatrix}
\]
\item Hence
\[
(A+uv')^{-1}=A^{-1}-\frac1{-1}A^{-1}uv'A^{-1}=A^{-1}+A^{-1}uv'A^{-1}
=\begin{pmatrix}1&-2&0&0\\-3&3&1&0\\3&-2&-1&0\\2&-2&-1&1\end{pmatrix}
\]
\item Check:
\[
A+uv'=A+\begin{pmatrix}0&0&0&0\\0&0&0&0\\1&2&0&0\\1&2&0&0\end{pmatrix}=\begin{pmatrix}1&2&2&0\\0&1&1&0\\3&4&3&0\\1&2&1&1\end{pmatrix}
\]
and, with the rows of \((A+uv')^{-1}\), \((1,-2,0,0)+2(-3,3,1,0)+2(3,-2,-1,0)=(1,0,0,0)\),
\((-3,3,1,0)+(3,-2,-1,0)=(0,1,0,0)\), \(3(1,-2,0,0)+4(-3,3,1,0)+3(3,-2,-1,0)=(0,0,1,0)\) and
\((1,-2,0,0)+2(-3,3,1,0)+(3,-2,-1,0)+(2,-2,-1,1)=(0,0,0,1)\). So \((A+uv')(A+uv')^{-1}=I_4\).
\end{itemize}

(c) \(A+s\,uv'=A+(su)v'\). By Theorem~\ref{r5d-thm-det-lemma} with \(su\) in place of \(u\),
\(\det(A+s\,uv')=\det A\cdot(1+s\,v'A^{-1}u)=1-2s\). By Theorem~\ref{r5d-thm-product}(a), \(A+s\,uv'\) is invertible if
and only if \(s\ne\frac12\).
\end{solution}

\begin{exercise}\label{r5d-ex-permutations}
Let \(\varphi=\langle2,4,1,3\rangle\) and \(\psi=\langle4,1,3,2\rangle\).
\begin{enumerate}
\item Find \(\varphi\psi\), \(\psi\varphi\) and \(\varphi^{-1}\).
\item Write down \(P_\varphi\) and \(P_\psi\), and compute \(P_\psi P_\varphi\).
\item Turn \(P_\varphi\) into \(I_4\) by interchanges of rows, and find \(\det P_\varphi\).
\end{enumerate}
\end{exercise}
\begin{solution}
(a) By Definition~\ref{r5d-def-permutation},
\begin{itemize}
\item \(\varphi\psi=\langle\varphi(4),\varphi(1),\varphi(3),\varphi(2)\rangle=\langle3,2,1,4\rangle\) and
\(\psi\varphi=\langle\psi(2),\psi(4),\psi(1),\psi(3)\rangle=\langle1,2,4,3\rangle\), so \(\varphi\psi\ne\psi\varphi\);
\item since \(\varphi(3)=1\), \(\varphi(1)=2\), \(\varphi(4)=3\) and \(\varphi(2)=4\), \(\varphi^{-1}=\langle3,1,4,2\rangle\).
\end{itemize}

(b) \(P_\varphi\) has rows \(e_2',e_4',e_1',e_3'\), and \(P_\psi\) has rows \(e_4',e_1',e_3',e_2'\). By
Proposition~\ref{r5d-prop-permutation-matrix}(a), the rows of \(P_\psi P_\varphi\) are rows \(4,1,3,2\) of \(P_\varphi\):
\[
P_\varphi=\begin{pmatrix}0&1&0&0\\0&0&0&1\\1&0&0&0\\0&0&1&0\end{pmatrix},\qquad
P_\psi=\begin{pmatrix}0&0&0&1\\1&0&0&0\\0&0&1&0\\0&1&0&0\end{pmatrix},\qquad
P_\psi P_\varphi=\begin{pmatrix}0&0&1&0\\0&1&0&0\\1&0&0&0\\0&0&0&1\end{pmatrix}
\]
This is \(P_{\langle3,2,1,4\rangle}=P_{\varphi\psi}\), as Proposition~\ref{r5d-prop-permutation-matrix}(b) says.

(c)
\[
\begin{gathered}
P_\varphi=\begin{pmatrix}0&1&0&0\\0&0&0&1\\1&0&0&0\\0&0&1&0\end{pmatrix}
\xrightarrow{R_1\leftrightarrow R_3}
\begin{pmatrix}1&0&0&0\\0&0&0&1\\0&1&0&0\\0&0&1&0\end{pmatrix}
\\
\xrightarrow{R_2\leftrightarrow R_3}
\begin{pmatrix}1&0&0&0\\0&1&0&0\\0&0&0&1\\0&0&1&0\end{pmatrix}
\xrightarrow{R_3\leftrightarrow R_4}
\begin{pmatrix}1&0&0&0\\0&1&0&0\\0&0&1&0\\0&0&0&1\end{pmatrix}=I_4
\end{gathered}
\]
Each interchange changes the sign of the determinant, and \(\det I_4=1\) (Definition~\ref{r5d-def-det}(c),(d)). Three
interchanges turn \(P_\varphi\) into \(I_4\), so \(\det P_\varphi=(-1)^3=-1\).
\end{solution}

\begin{exercise}\label{r5d-dr-permutations}
\begin{enumerate}
\item The data matrix
\[
X_0=\begin{pmatrix}3&8\\1&5\\4&2\\2&6\end{pmatrix}
\]
holds the measurements of two variables on four units. Find the permutation \(\varphi\) for which \(P_\varphi X_0\)
lists the units in increasing order of the second variable. Write down \(P_\varphi\) and \(P_\varphi X_0\), and find
\(\det P_\varphi\) by turning \(P_\varphi\) into \(I_4\) by interchanges of rows.
\item Find \(\varphi^{-1}\) and \(P_{\varphi^{-1}}\), and check that \(P_{\varphi^{-1}}=P_\varphi'\) and
\(P_\varphi'(P_\varphi X_0)=X_0\).
\end{enumerate}
\end{exercise}
\begin{solution}
(a)
\begin{itemize}
\item The second variable takes the values \(8,5,2,6\) on units \(1,2,3,4\). In increasing order these are
\(2,5,6,8\), the values of units \(3,2,4,1\). Row \(i\) of \(P_\varphi X_0\) is row \(\varphi(i)\) of \(X_0\)
(Proposition~\ref{r5d-prop-permutation-matrix}(a)), so \(\varphi=\langle3,2,4,1\rangle\).
\item \(P_\varphi\) has the rows \(e_3',e_2',e_4',e_1'\) (Definition~\ref{r5d-def-permutation-matrix}):
\[
P_\varphi=\begin{pmatrix}0&0&1&0\\0&1&0&0\\0&0&0&1\\1&0&0&0\end{pmatrix},\qquad
P_\varphi X_0=\begin{pmatrix}4&2\\1&5\\2&6\\3&8\end{pmatrix}
\]
\item Interchanging rows \(1\) and \(4\) of \(P_\varphi\), and then rows \(3\) and \(4\), gives \(I_4\):
\[
P_\varphi\xrightarrow{R_1\leftrightarrow R_4}
\begin{pmatrix}1&0&0&0\\0&1&0&0\\0&0&0&1\\0&0&1&0\end{pmatrix}
\xrightarrow{R_3\leftrightarrow R_4}
\begin{pmatrix}1&0&0&0\\0&1&0&0\\0&0&1&0\\0&0&0&1\end{pmatrix}=I_4
\]
Each interchange changes the sign of the determinant, and \(\det I_4=1\) (Definition~\ref{r5d-def-det}(c),(d)); so
\(\det P_\varphi=(-1)^2=1\).
\end{itemize}

(b) Since \(\varphi(4)=1\), \(\varphi(2)=2\), \(\varphi(1)=3\) and \(\varphi(3)=4\), the inverse is
\(\varphi^{-1}=\langle4,2,1,3\rangle\) (Definition~\ref{r5d-def-permutation}), and \(P_{\varphi^{-1}}\) has the rows
\(e_4',e_2',e_1',e_3'\):
\[
P_{\varphi^{-1}}=\begin{pmatrix}0&0&0&1\\0&1&0&0\\1&0&0&0\\0&0&1&0\end{pmatrix}=P_\varphi',\qquad
P_\varphi'(P_\varphi X_0)=\begin{pmatrix}3&8\\1&5\\4&2\\2&6\end{pmatrix}=X_0
\]
The last product lists rows \(4,2,1,3\) of \(P_\varphi X_0\) (Proposition~\ref{r5d-prop-permutation-matrix}(a)): the
units return to their original order, as Proposition~\ref{r5d-prop-permutation-matrix}(c) says.
\end{solution}

\needspace{14\baselineskip}
\subsection*{Problems}

\begin{problem}[\citeat{harville-solutions}{Ch.~8, Exercise~4(a)}]\label{r5b-ex-inverse-blocks}
Let \(A\in\R^{n\times n}\) be invertible and partitioned by columns as \(A=[A_1\ A_2]\), where
\(A_1\) has \(k\) columns, and let \(A^{-1}\) be partitioned by rows into \(B_1\) (\(k\) rows)
and \(B_2\). Show that
\[
B_1A_1=I,\quad B_1A_2=0,\quad B_2A_1=0,\quad B_2A_2=I,\qquad
A_1B_1=I-A_2B_2,\quad A_2B_2=I-A_1B_1
\]
\end{problem}
\begin{solution}
With the partition \((k,n-k)\) of \(I_n\), Theorem~\ref{r5b-thm-block}(c) gives
\[
\begin{pmatrix}B_1A_1&B_1A_2\\B_2A_1&B_2A_2\end{pmatrix}
=\begin{pmatrix}B_1\\B_2\end{pmatrix}[A_1\ A_2]=A^{-1}A=I_n=\begin{pmatrix}I_k&0\\0&I_{n-k}\end{pmatrix}
\]
and comparing blocks gives the first four identities. Also
\(A_1B_1+A_2B_2=[A_1\ A_2]\left(\begin{smallmatrix}B_1\\B_2\end{smallmatrix}\right)=AA^{-1}=I_n\),
which gives the last two.
\end{solution}

\begin{problem}[\citeat{abadir-magnus}{Exercise~5.8}]\label{r5b-ex-unipotent}
Show that the matrix
\(Z=\begin{pmatrix}-I_m&B\\C&I_n\end{pmatrix}\) (\(B\in\R^{m\times n}\), \(C\in\R^{n\times m}\))
satisfies \(Z^2=I_{m+n}\) if \(B=0\) or \(C=0\). Is this condition necessary?
\end{problem}
\begin{solution}
By the \(2\times2\) block formula (Theorem~\ref{r5b-thm-block}),
\[
Z^2=\begin{pmatrix}(-I)(-I)+BC&-B+B\\-C+C&CB+I\end{pmatrix}=\begin{pmatrix}I_m+BC&0\\0&I_n+CB\end{pmatrix}
\]
so \(Z^2=I_{m+n}\) if and only if \(BC=0\) and \(CB=0\). This holds if \(B=0\) or \(C=0\).
The condition is not necessary when \(m,n\geqslant2\): for \(m=n=2\) and
\(B=C=\left(\begin{smallmatrix}0&1\\0&0\end{smallmatrix}\right)\ne0\) we have \(BC=CB=0\)
(Proposition~\ref{r1-prop-strictly-triangular}(c)); for larger \(m,n\) put this \(2\times2\) matrix in
the top left corner of zero matrices. If \(m=1\), then \(CB\in\R^{n\times n}\) has entries
\(c_ib_j\), so \(CB=0\) forces \(B=0\) or \(C=0\); the case \(n=1\) is the same with \(BC\).
\end{solution}

\begin{problem}[\citeat{abadir-magnus}{Exercise~5.9}]\label{r5b-ex-commuting}
Show that
\(Z_1=\begin{pmatrix}I_m&B_1\\0&I_n\end{pmatrix}\) and
\(Z_2=\begin{pmatrix}I_m&B_2\\0&I_n\end{pmatrix}\) (\(B_1,B_2\in\R^{m\times n}\)) commute.
\end{problem}
\begin{solution}
Write \(Z_1=I_{m+n}+X_1\) and \(Z_2=I_{m+n}+X_2\) with
\(X_i=\left(\begin{smallmatrix}0&B_i\\0&0\end{smallmatrix}\right)\). By the block formula,
\(X_1X_2=\left(\begin{smallmatrix}0\cdot0+B_1\cdot0&0\cdot B_2+B_1\cdot0\\0&0\end{smallmatrix}\right)=0\),
and likewise \(X_2X_1=0\). Hence
\[
Z_1Z_2=I+X_1+X_2+X_1X_2=I+X_1+X_2=I+X_2+X_1+X_2X_1=Z_2Z_1
\]
\end{solution}

\begin{problem}[\citeat{hoffman-kunze}{\S5.3, Exercise~2}]\label{r5d-ex-vandermonde}
Let \(a,b,c\in\R\). Show that the \textbf{Vandermonde matrix}
\[
A=\begin{pmatrix}1&a&a^2\\1&b&b^2\\1&c&c^2\end{pmatrix}
\]
has \(\det A=(b-a)(c-a)(c-b)\), so that \(A\) is invertible if and only if \(a\), \(b\), \(c\) are distinct.
\end{problem}
\begin{solution}
\begin{itemize}
\item Subtracting row \(1\) from rows \(2\) and \(3\) does not change the determinant
(Definition~\ref{r5d-def-det}(b)); it gives the rows \((0,\,b-a,\,b^2-a^2)\) and \((0,\,c-a,\,c^2-a^2)\).
\item Expanding along column \(1\) (Theorem~\ref{r5d-thm-expansion}),
\[
\det A=\det\begin{pmatrix}b-a&b^2-a^2\\c-a&c^2-a^2\end{pmatrix}=(b-a)(c^2-a^2)-(c-a)(b^2-a^2)
\]
\item As \(c^2-a^2=(c-a)(c+a)\) and \(b^2-a^2=(b-a)(b+a)\), this is
\((b-a)(c-a)\bigl((c+a)-(b+a)\bigr)=(b-a)(c-a)(c-b)\).
\item The product is nonzero exactly when \(a\), \(b\), \(c\) are distinct, so \(A\) is invertible exactly then
(Theorem~\ref{r5d-thm-product}(a)).\qedhere
\end{itemize}
\end{solution}

\begin{problem}[\citeat{searle-khuri}{Exercise~3.13}]\label{r5d-ex-k-power}
Let \(K\in\R^{12\times12}\) satisfy \(K^5=3K\). What is the numerical value of \(\det K\)? Is there more
than one possibility? If so, give all possible values.
\end{problem}
\begin{solution}
\begin{itemize}
\item By Theorem~\ref{r5d-thm-product}(c), \(\det K^5=(\det K)^5\) and \(\det(3K)=3^{12}\det K\). So
\((\det K)^5=3^{12}\det K\), that is, \(\det K\cdot\bigl((\det K)^4-3^{12}\bigr)=0\).
\item Hence \(\det K=0\) or \((\det K)^4=3^{12}\), that is,
\(\bigl((\det K)^2-3^6\bigr)\bigl((\det K)^2+3^6\bigr)=0\). The second factor is positive, so
\((\det K)^2=3^6\) and \(\det K=\pm3^3=\pm27\).
\item All three values occur. Let \(t\) be the positive number with \(t^4=3\), so \(t^5=3t\). For
\(K=0\), \(\det K=0\). For \(K=tI_{12}\), \(K^5=t^5I_{12}=3K\) and \(\det K=t^{12}=3^3=27\) (Theorem~\ref{r5d-thm-basic}(c)). For
\(K=\diag(-t,t,\dots,t)\), \(K^5=\diag(-t^5,t^5,\dots,t^5)=3K\) and \(\det K=-t^{12}=-27\).
\end{itemize}
So \(\det K\) is \(0\), \(27\) or \(-27\).
\end{solution}

\begin{problem}[\citeat{bapat}{Ch.~2, Exercise~14}, \citeat{harville-solutions}{Ch.~13, Exercise~3}]\label{r5d-ex-integer}
Let \(A\in\R^{n\times n}\) have integer entries.
\begin{enumerate}
\item Show that if \(A\) is invertible and all entries of \(A^{-1}\) are integers, then \(\det A=\pm1\).
\item Conversely, show that if \(\det A=\pm1\), then \(A\) is invertible and all entries of \(A^{-1}\) are integers.
\end{enumerate}
\end{problem}
\begin{solution}
(a)
\begin{itemize}
\item A matrix with integer entries has an integer determinant, by induction on its size. For \(n=1\),
\(\det(a)=a\) (Definition~\ref{r5d-def-det}). For \(n\geqslant2\), the expansion along the first row
(Theorem~\ref{r5d-thm-expansion}) gives \(\det A=\sum_{k=1}^n(-1)^{1+k}a_{1k}\det A(1|k)\), where every
\(A(1|k)\in\R^{(n-1)\times(n-1)}\) has integer entries; so every \(\det A(1|k)\) is an integer, and so is \(\det A\).
Hence \(d=\det A\) and \(\det A^{-1}\) are integers.
\item By Theorem~\ref{r5d-thm-product}(a),(c), \(d\ne0\) and \(\det A^{-1}=1/d\). So \(d\) and \(1/d\) are nonzero
integers.
\item A nonzero integer has absolute value at least \(1\): \(\lvert d\rvert\geqslant1\) and \(\lvert1/d\rvert\geqslant1\), that
is, \(\lvert d\rvert\leqslant1\). Hence \(\lvert d\rvert=1\), and \(\det A=\pm1\).
\end{itemize}
(b) As \(\det A=\pm1\ne0\), \(A\) is invertible (Theorem~\ref{r5d-thm-product}(a)). For \(n=1\), \(A=(\pm1)\) and
\(A^{-1}=(\pm1)\). For \(n\geqslant2\), each cofactor \(c_{ij}=(-1)^{i+j}\det A(i|j)\) is an integer by the first item
of (a), applied to \(A(i|j)\). By Theorem~\ref{r5d-thm-adjugate},
\(A^{-1}=\frac1{\det A}\operatorname{adj}A=\pm\operatorname{adj}A\), whose entries \(\pm c_{ji}\) are integers.
\end{solution}

\begin{problem}[\citeat{searle-khuri}{Exercise~3.14}]\label{r5d-ex-sylvester}
Let \(A\in\R^{m\times n}\) and \(B\in\R^{n\times m}\). Show that \(\det(I_m-AB)=\det(I_n-BA)\).
\end{problem}
\begin{solution}
By the block product formula (Theorem~\ref{r5b-thm-block}),
\[
\begin{gathered}
\begin{pmatrix}I_m-AB&A\\0&I_n\end{pmatrix}\begin{pmatrix}I_m&0\\B&I_n\end{pmatrix}
=\begin{pmatrix}I_m-AB+AB&A\\B&I_n\end{pmatrix}=\begin{pmatrix}I_m&A\\B&I_n\end{pmatrix}\\
\begin{pmatrix}I_m&0\\B&I_n\end{pmatrix}\begin{pmatrix}I_m&A\\0&I_n-BA\end{pmatrix}
=\begin{pmatrix}I_m&A\\B&BA+I_n-BA\end{pmatrix}=\begin{pmatrix}I_m&A\\B&I_n\end{pmatrix}
\end{gathered}
\]
So the two products are equal. By Theorems~\ref{r5d-thm-product}(b) and~\ref{r5d-thm-block-tri}, the
first has determinant \(\det(I_m-AB)\det I_n\cdot\det I_m\det I_n=\det(I_m-AB)\), and the second
\(\det I_m\det I_n\cdot\det I_m\det(I_n-BA)=\det(I_n-BA)\).
\end{solution}

\begin{problem}[\citeat{bapat}{Ch.~2, Exercise~18}]\label{r5d-pr-commuting-blocks}
Let \(A,B,C,D\in\R^{n\times n}\) with \(A\) invertible and \(AC=CA\). Show that
\[
\det\begin{pmatrix}A&B\\C&D\end{pmatrix}=\det(AD-CB)
\]
\end{problem}
\begin{solution}
\begin{itemize}
\item Subtracting \(CA^{-1}\) times the first block row from the second turns the matrix into
\(\left(\begin{smallmatrix}A&B\\0&D-CA^{-1}B\end{smallmatrix}\right)\) (Theorem~\ref{r5b-thm-schur}). By
Theorems~\ref{r5d-thm-schur-det}(a) and~\ref{r5d-thm-product}(b),
\[
\det\begin{pmatrix}A&B\\C&D\end{pmatrix}=\det A\cdot\det(D-CA^{-1}B)=\det\bigl(A(D-CA^{-1}B)\bigr)=\det(AD-ACA^{-1}B)
\]
the last equality by Theorem~\ref{r1-thm-rules}.
\item From \(AC=CA\), \(ACA^{-1}=CAA^{-1}=C\), so the last determinant is \(\det(AD-CB)\).
\end{itemize}
\end{solution}

\begin{problem}[\citeat{abadir-magnus}{Exercise~5.45}]\label{r4-pr-block-triangular}
Let \(A\in\R^{m\times p}\), \(B\in\R^{m\times q}\), \(C\in\R^{n\times p}\) and \(D\in\R^{n\times q}\).
\begin{enumerate}
\item Show that
\[
\rank\begin{pmatrix}A&B\\0&D\end{pmatrix}\geqslant\rank A+\rank D\qquad\text{and}\qquad
\rank\begin{pmatrix}A&0\\C&D\end{pmatrix}\geqslant\rank A+\rank D
\]
\item Show that \(\rank\left(\begin{smallmatrix}A&B\\C&D\end{smallmatrix}\right)\geqslant\rank A+\rank D\) is not true in
general.
\end{enumerate}
\end{problem}
\begin{solution}
(a)
\begin{itemize}
\item Let \(Z=\left(\begin{smallmatrix}A&B\\0&D\end{smallmatrix}\right)\), \(r=\rank A\) and \(s=\rank D\). By
Theorem~\ref{thm:r4-four-bases}(a), the pivot columns of \(A\), say columns \(k_1<\dots<k_r\), are linearly
independent, and so are the pivot columns of \(D\), say columns \(l_1<\dots<l_s\) (a list is empty if the rank is
\(0\)). Column \(k_i\) of \(Z\) is \(\col_{k_i}(A)\) followed by zeros, and column \(p+l_j\) of \(Z\) is \(\col_{l_j}(B)\)
followed by \(\col_{l_j}(D)\).
\item These \(r+s\) columns of \(Z\) are linearly independent. Let
\[
\sum_{i=1}^r\alpha_i\begin{pmatrix}\col_{k_i}(A)\\0\end{pmatrix}+\sum_{j=1}^s\beta_j\begin{pmatrix}\col_{l_j}(B)\\\col_{l_j}(D)\end{pmatrix}=0
\]
The last \(n\) entries give \(\sum_j\beta_j\col_{l_j}(D)=0\), so every \(\beta_j=0\); then the first \(m\) entries give
\(\sum_i\alpha_i\col_{k_i}(A)=0\), so every \(\alpha_i=0\).
\item The \(r+s\) columns lie in \(\mathcal C(Z)\), whose dimension is \(\rank Z\) (Definition~\ref{def:r4-rank}); so
\(r+s\leqslant\rank Z\) (Theorem~\ref{thm:r4-counting}(c)).
\item For the second matrix, \(\left(\begin{smallmatrix}A&0\\C&D\end{smallmatrix}\right)'
=\left(\begin{smallmatrix}A'&C'\\0&D'\end{smallmatrix}\right)\) (Theorem~\ref{r5b-thm-block}(b)). By
Theorem~\ref{thm:r4-rank-rules}(a) and the first inequality, applied to \(A'\), \(C'\), \(D'\),
\[
\rank\begin{pmatrix}A&0\\C&D\end{pmatrix}=\rank\begin{pmatrix}A'&C'\\0&D'\end{pmatrix}\geqslant\rank A'+\rank D'
=\rank A+\rank D
\]
\end{itemize}
(b) Let \(m=n=p=q\geqslant1\) and \(A=B=C=D=I_m\). Then \(\rank A+\rank D=2m\) (Theorem~\ref{thm:r4-rank-rules}(d)). Every
column of \(Z=\left(\begin{smallmatrix}I_m&I_m\\I_m&I_m\end{smallmatrix}\right)\) is a column of its first block column
\(G=\left(\begin{smallmatrix}I_m\\I_m\end{smallmatrix}\right)\), so \(\mathcal C(Z)=\mathcal C(G)\)
(Theorem~\ref{r3-thm-span-inclusion}(c)). If \(Gx=0\), the first \(m\) entries give \(x=I_mx=0\); so \(G\) has full column
rank (Theorem~\ref{thm:r4-rank-rules}(b)), and \(\rank Z=\dim\mathcal C(G)=m<2m\) (Definition~\ref{def:r4-rank}).
\end{solution}

\begin{problem}[\citeat{prasolov}{Problem~1.5}]\label{r5d-ex-pr-ar1}
Let \(n\geqslant1\) and \(\rho\in\R\), and let \(A\in\R^{n\times n}\) have the entries \(a_{ij}=\rho^{\lvert i-j\rvert}\), with
\(\rho^0=1\) also for \(\rho=0\):
\[
A=\begin{pmatrix}1&\rho&\rho^2&\cdots&\rho^{n-1}\\\rho&1&\rho&\cdots&\rho^{n-2}\\\rho^2&\rho&1&\cdots&\rho^{n-3}\\
\vdots&\vdots&\vdots&\ddots&\vdots\\\rho^{n-1}&\rho^{n-2}&\rho^{n-3}&\cdots&1\end{pmatrix}
\]
Compute \(\det A\).
\end{problem}
\begin{solution}
For \(n=1\), \(A=(1)\) and \(\det A=1\). Let \(n\geqslant2\).
\begin{itemize}
\item For \(i=n,n-1,\dots,2\) in turn, subtract \(\rho\) times row \(i-1\) from row \(i\). Row \(i-1\) is changed only after
row \(i\), so it is still row \(i-1\) of \(A\) when it is used; and no step changes the determinant
(Definition~\ref{r5d-def-det}(b)).
\item Row \(1\) is unchanged, and the new row \(i\) (\(i\geqslant2\)) has the entries
\(\rho^{\lvert i-j\rvert}-\rho\cdot\rho^{\lvert i-1-j\rvert}\). For \(j\leqslant i-1\) this is
\(\rho^{i-j}-\rho\cdot\rho^{i-1-j}=0\), and for \(j\geqslant i\) it is \(\rho^{j-i}-\rho\cdot\rho^{j-i+1}=\rho^{j-i}(1-\rho^2)\).
So the result is the upper triangular matrix
\[
U=\begin{pmatrix}1&\rho&\rho^2&\cdots&\rho^{n-1}\\0&1-\rho^2&\rho(1-\rho^2)&\cdots&\rho^{n-2}(1-\rho^2)\\
0&0&1-\rho^2&\cdots&\rho^{n-3}(1-\rho^2)\\\vdots&\vdots&\vdots&\ddots&\vdots\\0&0&0&\cdots&1-\rho^2\end{pmatrix}
\]
\item By Theorem~\ref{r5d-thm-basic}(c), \(\det A=\det U=1\cdot(1-\rho^2)^{n-1}\).
\end{itemize}
So \(\det A=(1-\rho^2)^{n-1}\) for every \(n\geqslant1\). In particular, for \(n\geqslant2\), \(A\) is invertible if and only
if \(\rho\ne\pm1\) (Theorem~\ref{r5d-thm-product}(a)).
\end{solution}

\begin{problem}[\citeat{prasolov}{Problem~39.4}]\label{r5d-ex-pr-perm-commute}
Let \(n\geqslant2\) and \(A\in\R^{n\times n}\), and suppose that \(AP_\varphi=P_\varphi A\) for every permutation \(\varphi\) of
degree \(n\) (Definition~\ref{r5d-def-permutation-matrix}). Show that \(A=cI_n+dJ_n\) for some \(c,d\in\R\).
\end{problem}
\begin{solution}
\begin{itemize}
\item Let \(\varphi\) be a permutation of degree \(n\). Multiplying \(AP_\varphi=P_\varphi A\) on the right by \(P_\varphi'\) and
using \(P_\varphi P_\varphi'=I_n\) (Proposition~\ref{r5d-prop-permutation-matrix}(c)) gives \(A=P_\varphi AP_\varphi'\).
\item Column \(i\) of \(P_\varphi'\) is the transpose of row \(i\) of \(P_\varphi\), which is \(e_{\varphi(i)}'\)
(Definition~\ref{r5d-def-permutation-matrix}); so \(P_\varphi'e_i=e_{\varphi(i)}\) (Proposition~\ref{r1-prop-units}(a)). Hence,
by Theorem~\ref{r1-thm-transpose} and Proposition~\ref{r1-prop-units}(a), the \((i,j)\) entry of \(P_\varphi AP_\varphi'\) is
\[
e_i'P_\varphi AP_\varphi'e_j=(P_\varphi'e_i)'A(P_\varphi'e_j)=e_{\varphi(i)}'Ae_{\varphi(j)}=a_{\varphi(i)\varphi(j)}
\]
and the first item gives \(a_{ij}=a_{\varphi(i)\varphi(j)}\) for all \(i,j\) and every \(\varphi\).
\item Diagonal entries: let \(p\ne q\), and let \(\varphi\) interchange \(p\) and \(q\) and leave the other numbers in place.
Then \(a_{qq}=a_{\varphi(q)\varphi(q)}=a_{pp}\). So all diagonal entries of \(A\) are equal, say to \(\alpha\).
\item Off-diagonal entries: let \(i\ne j\) and \(p\ne q\). Put \(\varphi(p)=i\) and \(\varphi(q)=j\), and let \(\varphi\) map the
other \(n-2\) numbers of \(\{1,\dots,n\}\), in increasing order, to the \(n-2\) numbers other than \(i\) and \(j\), in
increasing order. Then \(\varphi\) is one-to-one and onto, a permutation (Definition~\ref{r5d-def-permutation}), and
\(a_{pq}=a_{\varphi(p)\varphi(q)}=a_{ij}\). So all off-diagonal entries of \(A\) are equal, say to \(\beta\).
\item Every entry of \(J_n\) is \(1\) (Definition~\ref{r1-def-ones}), so
\(A=\alpha I_n+\beta(J_n-I_n)=(\alpha-\beta)I_n+\beta J_n\): take \(c=\alpha-\beta\) and \(d=\beta\).
\end{itemize}
Conversely, every row and every column of \(P_\varphi\) has exactly one entry \(1\)
(Definition~\ref{r5d-def-permutation-matrix}, as \(\varphi\) is one-to-one and onto), so \(P_\varphi\one_n=\one_n\) and
\(P_\varphi'\one_n=\one_n\) (Theorem~\ref{r1-thm-views}(b)); hence \(P_\varphi J_n=\one_n\one_n'=J_nP_\varphi\), and every \(cI_n+dJ_n\) commutes with every
\(P_\varphi\). In statistical terms: listing \(n\) variables in the order \(\varphi(1),\dots,\varphi(n)\) replaces their
correlation matrix \(R\) by \(P_\varphi RP_\varphi'\), with \((i,j)\) entry \(r_{\varphi(i)\varphi(j)}\), and
\(P_\varphi RP_\varphi'=R\) is the same as \(RP_\varphi=P_\varphi R\); so \(R\) is unchanged by every reordering exactly when
\(R=(1-\rho)I_n+\rho J_n\), the equicorrelation matrix of Example~\ref{r5d-exm-equicorrelation} (the diagonal entries \(1\)
give \(c+d=1\)).
\end{solution}

\begin{problem}[\citeat{prasolov}{Problem~1.23, \S2.6}]\label{r5d-ex-pr-checkerboard}
Let \(A\in\R^{n\times n}\), and let \(B\in\R^{n\times n}\) have the entries \(b_{ij}=(-1)^{i+j}a_{ij}\). Show that
\(\det B=\det A\).
\end{problem}
\begin{solution}
Let \(k\) be the integer part of \(n/2\), the number of even numbers among \(1,\dots,n\).
\begin{itemize}
\item Multiply the rows \(2,4,\dots,2k\) of \(B\) by \(-1\). An entry in an odd row \(i\) stays
\((-1)^{i+j}a_{ij}=(-1)^{j+1}a_{ij}\), and an entry in an even row \(i\) becomes \(-(-1)^{i+j}a_{ij}=(-1)^{j+1}a_{ij}\): every
entry \((i,j)\) is now \((-1)^{j+1}a_{ij}\).
\item Multiply the columns \(2,4,\dots,2k\) of this matrix by \(-1\). An entry in an odd column \(j\) stays
\((-1)^{j+1}a_{ij}=a_{ij}\), and an entry in an even column \(j\) becomes \(-(-1)^{j+1}a_{ij}=a_{ij}\). The result is \(A\).
\item Each of these \(2k\) steps multiplies the determinant by \(-1\) (Definition~\ref{r5d-def-det}(a), for columns by
Theorem~\ref{r5d-thm-transpose}). So \(\det A=(-1)^{2k}\det B=\det B\).
\end{itemize}
For example, for \(n\geqslant2\) the matrix of the cofactors \(c_{ij}=(-1)^{i+j}\det A(i|j)\) of \(A\) and the matrix of the
minors \(\det A(i|j)\) (Definition~\ref{r5d-def-cofactor}) have the same determinant.
\end{solution}

\begin{problem}[\citeat{prasolov}{Problems~1.1, 1.2}]\label{r5d-ex-pr-skew-det}
Let \(A\in\R^{n\times n}\) be skew-symmetric (Definition~\ref{r1-def-skew}).
\begin{enumerate}
\item Show that \(\det A=0\) if \(n\) is odd.
\item Let \(n\) be even. Show that adding the same number \(x\in\R\) to all entries of \(A\) does not change the
determinant: \(\det(A+xJ_n)=\det A\).
\end{enumerate}
\end{problem}
\begin{solution}
(a)
\begin{itemize}
\item \(A'=-A\), so by Theorems~\ref{r5d-thm-transpose} and~\ref{r5d-thm-product}(c),
\(\det A=\det A'=\det(-A)=(-1)^n\det A=-\det A\), as \(n\) is odd.
\item Hence \(2\det A=0\), and \(\det A=0\).
\end{itemize}

(b) Let \(\one=\one_n\). The matrices below have order \(n+1\) and are partitioned after the first row and column.
\begin{itemize}
\item Let \(M=\left(\begin{smallmatrix}0&\one'\\-\one&A\end{smallmatrix}\right)\). By Theorem~\ref{r5b-thm-block}(b),
\(M'=\left(\begin{smallmatrix}0&-\one'\\\one&A'\end{smallmatrix}\right)=\left(\begin{smallmatrix}0&-\one'\\\one&-A\end{smallmatrix}\right)=-M\).
So \(M\) is skew-symmetric of the odd order \(n+1\), and \(\det M=0\) by (a).
\item The matrix \(\left(\begin{smallmatrix}0&\one'\\-x\one&A\end{smallmatrix}\right)\) is \(M\) with its first column
multiplied by \(x\); so its determinant is \(x\det M=0\) (Definition~\ref{r5d-def-det}(a) for columns,
Theorem~\ref{r5d-thm-transpose}). Also
\(\det\left(\begin{smallmatrix}1&0\\-x\one&A\end{smallmatrix}\right)=\det(1)\cdot\det A=\det A\)
(Theorem~\ref{r5d-thm-block-tri}).
\item These two matrices differ only in the first row, and their first rows \((0,1,\dots,1)\) and \((1,0,\dots,0)\) add up
to \((1,1,\dots,1)\). By Lemma~\ref{r5d-lem-multilinear},
\[
\det\begin{pmatrix}1&\one'\\-x\one&A\end{pmatrix}=0+\det A=\det A
\]
\item Subtract column \(1\) from each of the columns \(2,\dots,n+1\); this does not change the determinant
(Definition~\ref{r5d-def-det}(b) for columns, Theorem~\ref{r5d-thm-transpose}). Column \(j+1\) (\(1\leqslant j\leqslant n\))
is \(\left(\begin{smallmatrix}1\\\col_j(A)\end{smallmatrix}\right)\) and column \(1\) is
\(\left(\begin{smallmatrix}1\\-x\one\end{smallmatrix}\right)\), so column \(j+1\) becomes
\(\left(\begin{smallmatrix}0\\\col_j(A)+x\one\end{smallmatrix}\right)\). As every column of \(xJ_n=x\one\one'\) is \(x\one\)
(Definition~\ref{r1-def-ones}), Theorem~\ref{r5d-thm-block-tri} gives
\[
\det A=\det\begin{pmatrix}1&0\\-x\one&A+xJ_n\end{pmatrix}=\det(A+xJ_n)
\]
\end{itemize}
\end{solution}

\begin{problem}[\citeat{prasolov}{Problem~1.26}]\label{r5d-ex-pr-row-sums}
Let \(A\in\R^{n\times n}\), let \(s_i=a_{i1}+\dots+a_{in}\) be the sum of row \(i\) of \(A\), and let \(B\in\R^{n\times n}\)
have the entries \(b_{ij}=s_i-a_{ij}\). Show that \(\det B=(-1)^{n-1}(n-1)\det A\).
\end{problem}
\begin{solution}
Let \(s=(s_1,\dots,s_n)'\) and \(\one=\one_n\). The matrices below have order \(n+1\) and are partitioned after the \(n\)th
row and column.
\begin{itemize}
\item \(\det\left(\begin{smallmatrix}B&0\\-\one'&1\end{smallmatrix}\right)=\det B\cdot\det(1)=\det B\)
(Theorem~\ref{r5d-thm-block-tri}).
\item Add the columns \(1,\dots,n\) to column \(n+1\); this does not change the determinant
(Definition~\ref{r5d-def-det}(b) for columns, Theorem~\ref{r5d-thm-transpose}). Entry \(i\leqslant n\) of the new last column
is \(0+\sum_{j=1}^n(s_i-a_{ij})=ns_i-s_i=(n-1)s_i\), and its last entry is \(1+n\cdot(-1)=-(n-1)\); so the new last column is
\(n-1\) times \(\left(\begin{smallmatrix}s\\-1\end{smallmatrix}\right)\). Taking out this factor
(Definition~\ref{r5d-def-det}(a) for columns),
\[
\det B=\det\begin{pmatrix}B&(n-1)s\\-\one'&-(n-1)\end{pmatrix}=(n-1)\det\begin{pmatrix}B&s\\-\one'&-1\end{pmatrix}
\]
\item Subtract column \(n+1\) from each of the columns \(1,\dots,n\); this does not change the determinant either
(Definition~\ref{r5d-def-det}(b) for columns, Theorem~\ref{r5d-thm-transpose}). Entry \((i,j)\) with \(i\leqslant n\) becomes
\((s_i-a_{ij})-s_i=-a_{ij}\), and the last entry of column \(j\) becomes \(-1-(-1)=0\). So, by
Theorems~\ref{r5d-thm-block-tri} and~\ref{r5d-thm-product}(c),
\[
\det\begin{pmatrix}B&s\\-\one'&-1\end{pmatrix}=\det\begin{pmatrix}-A&s\\0&-1\end{pmatrix}=\det(-A)\cdot(-1)
=(-1)^{n+1}\det A
\]
\item Hence \(\det B=(n-1)(-1)^{n+1}\det A=(-1)^{n-1}(n-1)\det A\).
\end{itemize}
\end{solution}

\begin{problem}[\citeat{prasolov}{Problem~2.7}]\label{r5d-ex-pr-adj-one-entry}
Let \(n\geqslant2\) and \(1\leqslant i,j\leqslant n\). Give an example of a matrix \(A\in\R^{n\times n}\) whose adjugate
\(\operatorname{adj}A\) (Theorem~\ref{r5d-thm-adjugate}) has exactly one nonzero entry, in position \((i,j)\).
\end{problem}
\begin{solution}
Insert into \(I_{n-1}\) a column of zeros so that it becomes column \(i\), and then a row of zeros so that it becomes
row \(j\). The result \(A\in\R^{n\times n}\) has the zero row \(j\) and the zero column \(i\), and deleting them gives back
\(I_{n-1}\): \(A(j|i)=I_{n-1}\). For instance, for \(n=4\) and \((i,j)=(1,2)\),
\[
I_3\to\begin{pmatrix}0&1&0&0\\0&0&1&0\\0&0&0&1\end{pmatrix}\to A=\begin{pmatrix}0&1&0&0\\0&0&0&0\\0&0&1&0\\0&0&0&1\end{pmatrix}
\]
By Theorem~\ref{r5d-thm-adjugate}, the \((k,l)\) entry of \(\operatorname{adj}A\) is the cofactor
\(c_{lk}=(-1)^{l+k}\det A(l|k)\) (Definition~\ref{r5d-def-cofactor}).
\begin{itemize}
\item \((k,l)=(i,j)\): \(c_{ji}=(-1)^{j+i}\det I_{n-1}=(-1)^{i+j}\ne0\) (Definition~\ref{r5d-def-det}(d)).
\item \(l\ne j\): row \(j\) of \(A\), without its entry in column \(k\), is a zero row of \(A(l|k)\); so \(\det A(l|k)=0\)
(Theorem~\ref{r5d-thm-basic}(b)).
\item \(l=j\) and \(k\ne i\): column \(i\) of \(A\), without its entry in row \(j\), is a zero column of \(A(j|k)\); so
\(\det A(j|k)=0\) (Theorem~\ref{r5d-thm-basic}(b) for columns, Theorem~\ref{r5d-thm-transpose}).
\end{itemize}
So \(\operatorname{adj}A\) has exactly one nonzero entry, \((-1)^{i+j}\) in position \((i,j)\). In the instance,
\(\operatorname{adj}A=-E_{12}\), with \(E_{12}\) as in Proposition~\ref{r1-prop-units}(b).
\end{solution}

\begin{problem}[\citeat{prasolov}{Problem~2.9}]\label{r5d-ex-pr-adj-skew}
Let \(n\geqslant2\) and let \(A\in\R^{n\times n}\) be skew-symmetric. Show that \(\operatorname{adj}A\) is symmetric if
\(n\) is odd and skew-symmetric if \(n\) is even.
\end{problem}
\begin{solution}
Let \(c_{ij}\) be the cofactors of \(A\) (Definition~\ref{r5d-def-cofactor}).
\begin{itemize}
\item Deleting row \(j\) and column \(i\) of \(A\) and then transposing gives the same matrix as deleting row \(i\) and
column \(j\) of \(A'\): \(A(j|i)'=A'(i|j)\). As \(A'=-A\), \(A'(i|j)=-A(i|j)\).
\item \(A(i|j)\) has order \(n-1\), so by Theorems~\ref{r5d-thm-transpose} and~\ref{r5d-thm-product}(c),
\(\det A(j|i)=\det A(j|i)'=\det\bigl(-A(i|j)\bigr)=(-1)^{n-1}\det A(i|j)\). Hence
\[
c_{ji}=(-1)^{j+i}\det A(j|i)=(-1)^{n-1}(-1)^{i+j}\det A(i|j)=(-1)^{n-1}c_{ij}
\]
\item By Theorem~\ref{r5d-thm-adjugate}, the \((i,j)\) entry of \(\operatorname{adj}A\) is \(c_{ji}\), and that of
\((\operatorname{adj}A)'\) is \(c_{ij}\). So \(\operatorname{adj}A=(-1)^{n-1}(\operatorname{adj}A)'\), that is,
\((\operatorname{adj}A)'=(-1)^{n-1}\operatorname{adj}A\).
\item For odd \(n\), \((\operatorname{adj}A)'=\operatorname{adj}A\): \(\operatorname{adj}A\) is symmetric. For even \(n\),
\((\operatorname{adj}A)'=-\operatorname{adj}A\): \(\operatorname{adj}A\) is skew-symmetric (Definition~\ref{r1-def-skew}).
\end{itemize}
\end{solution}

\begin{problem}[\citeat{prasolov}{Problems~2.2, 3.1}]\label{r5d-ex-pr-bordered}
Let \(n\geqslant2\), let \(A\in\R^{n\times n}\) have the cofactors \(c_{ij}\) (Definition~\ref{r5d-def-cofactor}), and let
\(x,y,u,v\in\R^n\).
\begin{enumerate}
\item Show that
\[
\det\begin{pmatrix}A&x\\y'&0\end{pmatrix}=-\sum_{i=1}^n\sum_{j=1}^nx_iy_jc_{ij}=-y'(\operatorname{adj}A)x
\]
\item Show that \(\det(A+uv')=\det A+v'(\operatorname{adj}A)u\), also when \(A\) is not invertible.
\end{enumerate}
\end{problem}
\begin{solution}
(a) Let \(M\) be the matrix of order \(n+1\) on the left.
\begin{itemize}
\item Column \(n+1\) of \(M\) is \(\left(\begin{smallmatrix}x\\0\end{smallmatrix}\right)\). Expanding along it
(Theorem~\ref{r5d-thm-expansion}),
\[
\det M=\sum_{i=1}^nx_i(-1)^{i+n+1}\det M(i|n+1)
\]
\item \(M(i|n+1)\in\R^{n\times n}\) consists of the rows of \(A\) other than row \(i\), followed by the row \(y'\); deleting
this last row and column \(j\) leaves \(A(i|j)\). Expanding along the last row (Theorem~\ref{r5d-thm-expansion}),
\(\det M(i|n+1)=\sum_{j=1}^ny_j(-1)^{n+j}\det A(i|j)\).
\item As \((-1)^{i+n+1}(-1)^{n+j}=-(-1)^{i+j}\),
\[
\det M=-\sum_{i=1}^n\sum_{j=1}^nx_iy_j(-1)^{i+j}\det A(i|j)=-\sum_{i=1}^n\sum_{j=1}^nx_iy_jc_{ij}
\]
\item By Theorem~\ref{r5d-thm-adjugate}, \((\operatorname{adj}A)_{ji}=c_{ij}\); so entry \(j\) of \((\operatorname{adj}A)x\) is
\(\sum_{i=1}^nc_{ij}x_i\), and \(y'(\operatorname{adj}A)x=\sum_{j=1}^ny_j\sum_{i=1}^nc_{ij}x_i\) is the double sum.
\end{itemize}

(b) Let \(N=\left(\begin{smallmatrix}A&-u\\v'&1\end{smallmatrix}\right)\in\R^{(n+1)\times(n+1)}\). We compute \(\det N\) in
two ways.
\begin{itemize}
\item By Theorem~\ref{r5d-thm-schur-det}(b), with the invertible block \(D=(1)\),
\(\det N=\det(1)\cdot\det\bigl(A-(-u)(1)^{-1}v'\bigr)=\det(A+uv')\).
\item The last column of \(N\) is \(\left(\begin{smallmatrix}-u\\0\end{smallmatrix}\right)
+\left(\begin{smallmatrix}0\\1\end{smallmatrix}\right)\), and the other columns are kept. By
Lemma~\ref{r5d-lem-multilinear} for columns (Theorem~\ref{r5d-thm-transpose}), then by (a) with \(x=-u\), \(y=v\), and by
Theorem~\ref{r5d-thm-block-tri},
\[
\det N=\det\begin{pmatrix}A&-u\\v'&0\end{pmatrix}+\det\begin{pmatrix}A&0\\v'&1\end{pmatrix}
=-v'(\operatorname{adj}A)(-u)+\det A\cdot\det(1)
\]
\item Comparing the two values, \(\det(A+uv')=\det A+v'(\operatorname{adj}A)u\).
\end{itemize}
If \(A\) is invertible, then \(\operatorname{adj}A=(\det A)A^{-1}\) (Theorem~\ref{r5d-thm-adjugate}), and (b) is
Theorem~\ref{r5d-thm-det-lemma}. If \(A\) is not invertible, then \(\det A=0\) (Theorem~\ref{r5d-thm-product}(a)), and (b) reads
\(\det(A+uv')=v'(\operatorname{adj}A)u\).
\end{solution}

\begin{problem}[\citeat{prasolov}{Problem~19.5}]\label{r5d-ex-pr-sym-minor}
Let \(A\in\R^{n\times n}\) be symmetric with \(\rank A=r\geqslant1\). By Theorem~\ref{r5d-thm-rank-minors}(b), \(A\) has an
\(r\times r\) submatrix with nonzero determinant. Show that there is one whose rows and columns have the same indices:
for some \(i_1<\dots<i_r\), the \(r\times r\) submatrix \(A_*\) of \(A\) in the rows \(i_1,\dots,i_r\) and the columns
\(i_1,\dots,i_r\) (Definition~\ref{r5d-def-minor}) has \(\det A_*\ne0\).
\end{problem}
\begin{solution}
\begin{itemize}
\item Let \(R\) be the RREF of \(A\); it has \(r\) nonzero rows (Definition~\ref{def:r4-rank}). Let \(i_1<\dots<i_r\) be its
pivot columns. By Theorem~\ref{thm:r4-pivot-columns}(b), the columns \(a_{i_1},\dots,a_{i_r}\) of \(A\) form a basis of
\(\Span\{a_1,\dots,a_n\}\), so they are linearly independent, and \(a_j=R_{1j}a_{i_1}+\dots+R_{rj}a_{i_r}\) for every \(j\).
\item Let \(S=[e_{i_1}\ \cdots\ e_{i_r}]\in\R^{n\times r}\), and let \(X\in\R^{r\times n}\) consist of the first \(r\) rows of
\(R\). Then \(AS=[a_{i_1}\ \cdots\ a_{i_r}]\) (Proposition~\ref{r1-prop-units}(a)), and the first item reads \(A=(AS)X\)
(Theorem~\ref{r1-thm-views}(a),(b)).
\item Multiplying on the left by \(S'\) gives \(S'A=(S'AS)X\). The \((s,t)\) entry of \(S'AS\) is
\(e_{i_s}'Ae_{i_t}=a_{i_si_t}\) (Proposition~\ref{r1-prop-units}(a)), so \(S'AS=A_*\).
\item As \(A'=A\), \(S'A=S'A'=(AS)'\) (Theorem~\ref{r1-thm-transpose}). The \(r\) columns of \(AS\) are linearly independent,
so \(\rank(S'A)=\rank\bigl((AS)'\bigr)=\rank(AS)=r\) (Theorem~\ref{thm:r4-rank-rules}(a),(b)). So \(S'A\) has full row
rank, and its rows, which are the rows \(i_1,\dots,i_r\) of \(A\) (Proposition~\ref{r1-prop-units}(a)), are linearly
independent too (Theorem~\ref{thm:r4-rank-rules}(c)).
\item By Theorems~\ref{thm:r4-rank-product}(a) and~\ref{thm:r4-rank-rules}(a),
\(r=\rank(S'A)=\rank(A_*X)\leqslant\rank A_*\leqslant r\). So \(\rank A_*=r\), \(A_*\) is invertible
(Theorem~\ref{thm:r4-rank-rules}(d)), and \(\det A_*\ne0\) (Theorem~\ref{r5d-thm-product}(a)).
\end{itemize}
\end{solution}

\clearpage
\def\unitnumber{6}
\setcounter{theorem}{0}\setcounter{problem}{0}
\section*{Chapter 6. Orthogonality and projections}

\subsection*{Inner products, lengths, and orthogonality}

\begin{definition}[\citeat{hoffman-kunze}{\S8.1}, \citeat{lang}{Ch.~V, \S1}, \citeat{tu-clm}{\S1.3}]\label{r6-def-inner}
For \(x,y\in\R^n\), the \textbf{inner product} of \(x\) and \(y\) is the number
\(x'y=x_1y_1+\dots+x_ny_n\); it equals \(y'x\) (Theorem~\ref{r1-thm-transpose}). The \textbf{length}
(or norm) of \(x\) is \(\lVert x\rVert=\sqrt{x'x}=(x_1^2+\dots+x_n^2)^{1/2}\), and \(\lVert x-y\rVert\) is the
\textbf{distance} between \(x\) and \(y\). The vectors \(x\) and \(y\) are \textbf{orthogonal}, written
\(x\perp y\), if \(x'y=0\).
\end{definition}
Since \(x'x\) is a sum of squares, \(\lVert x\rVert\geqslant0\), and \(\lVert x\rVert=0\) only for \(x=0\); also
\(\lVert cx\rVert=\lvert c\rvert\,\lVert x\rVert\) for \(c\in\R\). A vector of length \(1\) is a \textbf{unit vector};
for \(x\ne0\), \(x/\lVert x\rVert\) is one. The zero vector is orthogonal to every vector, and it is the only
vector orthogonal to itself.

\begin{proposition}[\citeat{hoffman-kunze}{\S8.1}, \citeat{lang}{Ch.~V, \S1}, \citeat{tu-clm}{proof of Theorem~2.5}]\label{r6-prop-pythagoras}
For \(x,y\in\R^n\), \(\lVert x+y\rVert^2=\lVert x\rVert^2+2x'y+\lVert y\rVert^2\). In particular, if \(x\perp y\), then
\(\lVert x+y\rVert^2=\lVert x\rVert^2+\lVert y\rVert^2\) (\textbf{Pythagoras}).
\end{proposition}
\begin{proof}
\((x+y)'(x+y)=x'x+x'y+y'x+y'y\), and \(y'x=x'y\).
\end{proof}

\begin{proposition}[\citeat{hoffman-kunze}{\S8.1, (8-3)}, \citeat{lang}{Ch.~V, \S1}]\label{r6-prop-polarization}
Let \(x,y\in\R^n\).
\begin{enumerate}
\item \(\lVert x-y\rVert^2=\lVert x\rVert^2-2x'y+\lVert y\rVert^2\), and
\(x'y=\frac14\lVert x+y\rVert^2-\frac14\lVert x-y\rVert^2\) (\textbf{polarization}).
\item (\textbf{Pythagoras for \(k\) vectors}) If \(u_1,\dots,u_k\in\R^n\) are pairwise orthogonal
(\(u_i'u_j=0\) for \(i\ne j\)), then \(\lVert u_1+\dots+u_k\rVert^2=\lVert u_1\rVert^2+\dots+\lVert u_k\rVert^2\).
\end{enumerate}
\end{proposition}
\begin{proof}
(a)
\begin{itemize}
\item Proposition~\ref{r6-prop-pythagoras} with \(-y\) in place of \(y\) gives the first identity, as
\(x'(-y)=-x'y\) and \(\lVert-y\rVert=\lVert y\rVert\).
\item Subtracting it from \(\lVert x+y\rVert^2=\lVert x\rVert^2+2x'y+\lVert y\rVert^2\) gives
\(\lVert x+y\rVert^2-\lVert x-y\rVert^2=4x'y\).
\end{itemize}
(b) Induction on \(k\); for \(k=1\) there is nothing to prove.
\begin{itemize}
\item Let \(k\geqslant2\) and \(s=u_1+\dots+u_{k-1}\). Then \(s'u_k=u_1'u_k+\dots+u_{k-1}'u_k=0\), so
\(s\perp u_k\).
\item By Proposition~\ref{r6-prop-pythagoras}, \(\lVert s+u_k\rVert^2=\lVert s\rVert^2+\lVert u_k\rVert^2\), and
\(\lVert s\rVert^2=\lVert u_1\rVert^2+\dots+\lVert u_{k-1}\rVert^2\) by the induction hypothesis.\qedhere
\end{itemize}
\end{proof}

\begin{theorem}[\citeat{tu-clm}{Lemma~2.16}, \citeat{hoffman-kunze}{\S8.2, Theorem~1(iv)}, \citeat{harville-solutions}{Ch.~6, Exercise~1}]\label{r6-thm-cauchy-schwarz}
Let \(x,y\in\R^n\).
\begin{enumerate}
\item (\textbf{Cauchy--Schwarz inequality}) \(\lvert x'y\rvert\leqslant\lVert x\rVert\,\lVert y\rVert\), with equality if and
only if one of \(x,y\) is a multiple of the other.
\item (\textbf{Triangle inequality}) \(\lVert x+y\rVert\leqslant\lVert x\rVert+\lVert y\rVert\), with equality if and only if
\(y=0\) or \(x=cy\) for some \(c\geqslant0\).
\end{enumerate}
\end{theorem}
\begin{proof}
(a) If \(x=0\), both sides are \(0\), and \(x=0y\). Let \(x\ne0\).
\begin{itemize}
\item Put \(c=x'y/\lVert x\rVert^2\) and \(w=y-cx\). Then \(x'w=x'y-c\,x'x=0\), so \(cx\perp w\).
\item Pythagoras for \(y=cx+w\) gives
\(\lVert y\rVert^2=c^2\lVert x\rVert^2+\lVert w\rVert^2\geqslant c^2\lVert x\rVert^2=(x'y)^2/\lVert x\rVert^2\). Multiply by
\(\lVert x\rVert^2\) and take square roots.
\item Equality holds if and only if \(w=0\), that is, \(y=cx\). Conversely, if \(y=ax\), then \(c=a\) and
\(w=0\); if \(x=ay\) with \(x\ne0\), then \(a\ne0\) and \(y=a^{-1}x\).
\end{itemize}
(b)
\begin{itemize}
\item By Proposition~\ref{r6-prop-pythagoras} and (a),
\[
\lVert x+y\rVert^2=\lVert x\rVert^2+2x'y+\lVert y\rVert^2\leqslant\lVert x\rVert^2+2\lvert x'y\rvert+\lVert y\rVert^2
\leqslant\lVert x\rVert^2+2\lVert x\rVert\,\lVert y\rVert+\lVert y\rVert^2=(\lVert x\rVert+\lVert y\rVert)^2
\]
and both \(\lVert x+y\rVert\) and \(\lVert x\rVert+\lVert y\rVert\) are \(\geqslant0\).
\item Equality holds if and only if both steps are equalities: \(x'y=\lvert x'y\rvert\), that is, \(x'y\geqslant0\);
and \(\lvert x'y\rvert=\lVert x\rVert\,\lVert y\rVert\).
\item If \(y=0\), both hold. Let \(y\ne0\). By (a), the second holds if and only if \(x=cy\) for some \(c\) (if
\(y=ax\), then \(a\ne0\) and \(x=a^{-1}y\)); and then \(x'y=c\lVert y\rVert^2\geqslant0\) if and only if
\(c\geqslant0\).\qedhere
\end{itemize}
\end{proof}

\begin{definition}[\citeat{tu-clm}{\S1.3}, \citeat{hoffman-kunze}{\S8.2}, \citeat{harville}{\S12.1}]\label{r6-def-complement}
Let \(W\) be a subspace of \(\R^n\). A vector \(y\in\R^n\) is \textbf{orthogonal to} \(W\), written
\(y\perp W\), if \(y'w=0\) for every \(w\in W\). The \textbf{orthogonal complement} of \(W\) is
\(W^\perp=\{y\in\R^n:y\perp W\}\). Subspaces \(U\) and \(W\) of \(\R^n\) are \textbf{orthogonal}, written
\(U\perp W\), if \(u'w=0\) for all \(u\in U\) and \(w\in W\); their sum \(U+W\) is then called an
\textbf{orthogonal sum}.
\end{definition}

\begin{theorem}[\citeat{harville}{Lemma~12.1.1, Corollaries~12.1.2, 12.5.3}, \citeat{harville-solutions}{Ch.~12, Exercise~1}, \citeat{abadir-magnus}{Exercise~3.48(c)}]\label{r6-thm-complement}
Let \(W\) be a subspace of \(\R^n\).
\begin{enumerate}
\item If \(w_1,\dots,w_k\) span \(W\) and \(A=[w_1\ \cdots\ w_k]\in\R^{n\times k}\), then \(y\perp W\) if and
only if \(y'w_1=\dots=y'w_k=0\), that is, \(A'y=0\). So \(W^\perp=\mathcal N(A')\), a subspace of \(\R^n\).
\item \(W\cap W^\perp=\{0\}\).
\item \(\dim W+\dim W^\perp=n\) and \(\R^n=W\oplus W^\perp\): every \(y\in\R^n\) is \(y=w+v\) with \(w\in W\)
and \(v\in W^\perp\) in exactly one way.
\end{enumerate}
\end{theorem}
\begin{proof}
(a) If \(y\perp W\), then \(y'w_j=0\), as \(w_j\in W\). Conversely, let all \(y'w_j=0\). Every \(w\in W\)
is \(w=Ac\) for some \(c\in\R^k\) (Definition~\ref{r3-def-span} and Theorem~\ref{r1-thm-views}(b)), and
\(y'w=(A'y)'c=0\). The \(j\)th entry of \(A'y\) is \(w_j'y\) (Theorem~\ref{r1-thm-views}(d)).

(b) If \(y\in W\cap W^\perp\), then \(y\perp y\), so \(y=0\).

(c) Let \(d=\dim W\). If \(d=0\), then \(W=\{0\}\) and \(W^\perp=\R^n\). Let \(d\geqslant1\).
\begin{itemize}
\item Let \(A\in\R^{n\times d}\) have a basis of \(W\) as its columns, so \(\rank A=d\) (Definition~\ref{def:r4-rank}, as \(\mathcal C(A)=W\)). By (a) and
Theorem~\ref{thm:r4-rank-nullity}, \(\dim W^\perp=\dim\mathcal N(A')=n-\rank A=n-d\).
\item By (b), the sum \(W+W^\perp\) is direct (Definition~\ref{r3-def-sum}), and a basis of \(W\)
together with a basis of \(W^\perp\), \(d+(n-d)=n\) vectors, is a basis of \(W\oplus W^\perp\)
(Theorem~\ref{r3-thm-direct-sum-basis}). So \(W\oplus W^\perp\) is a subspace of \(\R^n\)
(Theorem~\ref{r3-thm-sum}(a)) of dimension \(n\), hence equal to \(\R^n\)
(Theorem~\ref{thm:r4-counting}(a)). The decomposition \(y=w+v\) is unique by
Theorem~\ref{r3-thm-direct-sum}.\qedhere
\end{itemize}
\end{proof}

\begin{corollary}[\citeat{bapat}{Result~2.9}, \citeat{harville}{Lemma~12.5.2}]\label{r6-cor-four-subspaces}
For \(A\in\R^{m\times n}\),
\[
\mathcal N(A)=\mathcal C(A')^\perp,\quad \R^n=\mathcal C(A')\oplus\mathcal N(A);\qquad
\mathcal N(A')=\mathcal C(A)^\perp,\quad \R^m=\mathcal C(A)\oplus\mathcal N(A')
\]
So \(Ax=0\) if and only if \(x\) is orthogonal to every (transposed) row of \(A\), and the four
subspaces of \(A\) (Definition~\ref{def:r4-spaces}) are orthogonal in pairs: the transposed row space \(\mathcal C(A')\) and
the null space in \(\R^n\), the column space and the left null space in \(\R^m\).
\end{corollary}
\begin{proof}
The columns of \(A'\) span \(\mathcal C(A')\), so Theorem~\ref{r6-thm-complement}(a) gives
\(\mathcal C(A')^\perp=\mathcal N((A')')=\mathcal N(A)\); applied to the columns of \(A\), it gives
\(\mathcal C(A)^\perp=\mathcal N(A')\). The direct sums follow from
Theorem~\ref{r6-thm-complement}(c).
\end{proof}

\begin{theorem}[\citeat{harville}{Theorem~12.5.4, Corollaries~12.5.5, 12.5.6}]\label{r6-thm-perp-perp}
Let \(W\), \(W_1\) and \(W_2\) be subspaces of \(\R^n\).
\begin{enumerate}
\item \((W^\perp)^\perp=W\).
\item \(W_1\subseteq W_2\) if and only if \(W_2^\perp\subseteq W_1^\perp\).
\item \(W_1=W_2\) if and only if \(W_1^\perp=W_2^\perp\).
\end{enumerate}
\end{theorem}
\begin{proof}
(a)
\begin{itemize}
\item If \(w\in W\), then \(w'v=v'w=0\) for every \(v\in W^\perp\), so \(w\in(W^\perp)^\perp\). Thus
\(W\subseteq(W^\perp)^\perp\).
\item Conversely, let \(y\in(W^\perp)^\perp\), and write \(y=w+v\) with \(w\in W\) and \(v\in W^\perp\)
(Theorem~\ref{r6-thm-complement}(c)).
\item Then \(y'v=0\), as \(y\perp W^\perp\), and \(w'v=0\), as \(v\perp W\); so
\(v'v=(y-w)'v=y'v-w'v=0\).
\item Hence \(v=0\) and \(y=w\in W\).
\end{itemize}
(b)
\begin{itemize}
\item For subspaces \(Y\), \(U\) of \(\R^n\): \(Y\subseteq U\) if and only if \(Y\perp U^\perp\). Indeed,
\(Y\perp U^\perp\) says that every vector of \(Y\) lies in \((U^\perp)^\perp\), which is \(U\) by (a).
\item By this, (a) for \(W_1\), the symmetry of \(\perp\) (\(u'w=w'u\)), and this again with
\(Y=W_2^\perp\) and \(U=W_1^\perp\),
\[
W_1\subseteq W_2\iff W_1\perp W_2^\perp\iff(W_1^\perp)^\perp\perp W_2^\perp\iff W_2^\perp\perp(W_1^\perp)^\perp
\iff W_2^\perp\subseteq W_1^\perp
\]
\end{itemize}
(c) If \(W_1^\perp=W_2^\perp\), then \(W_1^\perp\subseteq W_2^\perp\) and \(W_2^\perp\subseteq W_1^\perp\), so
\(W_2\subseteq W_1\) and \(W_1\subseteq W_2\) by (b). Conversely, if \(W_1=W_2\), then \(W_1\subseteq W_2\) and
\(W_2\subseteq W_1\), so \(W_2^\perp\subseteq W_1^\perp\) and \(W_1^\perp\subseteq W_2^\perp\) by (b).
\end{proof}

\begin{corollary}[\citeat{harville}{Lemma~12.5.2, Theorem~12.5.4, Corollary~12.5.7}]\label{r6-cor-consistency}
For \(A\in\R^{m\times n}\), \(\mathcal C(A')=\mathcal N(A)^\perp\) and \(\mathcal C(A)=\mathcal N(A')^\perp\). Hence,
for \(b\in\R^m\), the system \(Ax=b\) has a solution if and only if \(z'b=0\) for every \(z\in\R^m\) with
\(A'z=0\); if \(z_1,\dots,z_s\) span \(\mathcal N(A')\), it suffices that \(z_1'b=\dots=z_s'b=0\).
\end{corollary}
\begin{proof}
\begin{itemize}
\item By Corollary~\ref{r6-cor-four-subspaces}, \(\mathcal N(A)=\mathcal C(A')^\perp\) and
\(\mathcal N(A')=\mathcal C(A)^\perp\). By Theorem~\ref{r6-thm-perp-perp}(a),
\(\mathcal N(A)^\perp=(\mathcal C(A')^\perp)^\perp=\mathcal C(A')\) and \(\mathcal N(A')^\perp=\mathcal C(A)\).
\item \(Ax=b\) has a solution if and only if \(b\in\mathcal C(A)\) (Definition~\ref{def:r4-spaces}), that is,
\(b\perp\mathcal N(A')\). By Theorem~\ref{r6-thm-complement}(a), applied to \(\mathcal N(A')\) and its spanning
vectors \(z_1,\dots,z_s\), this means \(z_1'b=\dots=z_s'b=0\).\qedhere
\end{itemize}
\end{proof}

\begin{example}\label{r6-exm-four-subspaces}
Let
\[
A=\begin{pmatrix}1&2&1&3\\1&3&0&5\\2&5&1&8\end{pmatrix}
\]
be the matrix of Example~\ref{exm:r4-four-bases}. The nonzero rows of its RREF are \(\rho_1'\) and \(\rho_2'\), with
\(\rho_1=(1,0,3,-1)'\) and \(\rho_2=(0,1,-1,2)'\), and they form a basis of \(\mathcal R(A)\); \(s_1=(-3,1,1,0)'\) and
\(s_2=(1,-2,0,1)'\) form a basis of \(\mathcal N(A)\) (Example~\ref{exm:r4-four-bases}). Transposing maps
\(\mathcal R(A)\) onto \(\mathcal C(A')\), as \((z'A)'=A'z\) (\(z\in\R^3\)), preserves linear combinations and is undone
by transposing again; so \(\rho_1,\rho_2\) form a basis of \(\mathcal C(A')\). Indeed
\[
\rho_1's_1=-3+0+3+0=0,\quad \rho_1's_2=1+0+0-1=0,\quad \rho_2's_1=0+1-1+0=0,\quad \rho_2's_2=0-2+0+2=0
\]
and \(2+2=4\). In \(\R^3\), \(y=(-1,-1,1)'\) is a basis of \(\mathcal N(A')\), orthogonal to the basis
\(a_1=(1,1,2)'\), \(a_2=(2,3,5)'\) of \(\mathcal C(A)\): \(-1-1+2=0\) and \(-2-3+5=0\); and \(2+1=3\).
By Corollary~\ref{r6-cor-consistency}, \(Ax=b\) has a solution exactly when \(y'b=-b_1-b_2+b_3=0\). For
\(b=(1,2,3)'\), \(-1-2+3=0\), and indeed \(A(-1,1,0,0)'=-a_1+a_2=(1,2,3)'\)
(Theorem~\ref{r1-thm-views}(b)); for \(b=e_1\), \(y'e_1=-1\ne0\), so \(Ax=e_1\) has no solution.
\end{example}

\subsection*{Orthonormal bases, Gram--Schmidt, and QR}

\begin{proposition}[\citeat{hoffman-kunze}{\S8.2, Theorem~2, (8-8)}, \citeat{bapat}{Result~2.5}]\label{r6-prop-orthogonal-list}
Let \(u_1,\dots,u_k\in\R^n\) be nonzero and pairwise orthogonal (\(u_i'u_j=0\) for \(i\ne j\)).
\begin{enumerate}
\item If \(y=c_1u_1+\dots+c_ku_k\), then \(c_j=u_j'y/\lVert u_j\rVert^2\) for \(j=1,\dots,k\).
\item \(u_1,\dots,u_k\) are linearly independent.
\end{enumerate}
\end{proposition}
\begin{proof}
(a) \(u_j'y=\sum_ic_iu_j'u_i=c_j\lVert u_j\rVert^2\), and \(\lVert u_j\rVert^2>0\). (b) If \(\sum_ic_iu_i=0\), then (a)
with \(y=0\) gives every \(c_j=0\).
\end{proof}

\begin{definition}[\citeat{hoffman-kunze}{\S8.2}, \citeat{searle-khuri}{\S5.4}, \citeat{harville}{\S8.4a}]\label{r6-def-orthonormal}
Vectors \(q_1,\dots,q_p\in\R^n\) are \textbf{orthonormal} if \(q_i'q_j=0\) for \(i\ne j\) and \(\lVert q_j\rVert=1\)
for every \(j\). A matrix \(Q=[q_1\ \cdots\ q_p]\in\R^{n\times p}\) has \textbf{orthonormal columns} if
\(Q'Q=I_p\), since \((Q'Q)_{ij}=q_i'q_j\) (Theorem~\ref{r1-thm-views}(d)). A square matrix \(Q\in\R^{n\times n}\) with
\(Q'Q=I_n\) is an \textbf{orthogonal matrix}. A basis of a subspace consisting of orthonormal vectors is
an \textbf{orthonormal basis}.
\end{definition}

\begin{proposition}[\citeat{lang}{Ch.~VII, Theorems~3.1, 3.2}, \citeat{abadir-magnus}{Exercise~8.4}, \citeat{harville}{\S8.4a}]\label{r6-prop-orthonormal}
Let \(Q\in\R^{n\times p}\) satisfy \(Q'Q=I_p\).
\begin{enumerate}
\item \((Qx)'(Qy)=x'y\) and \(\lVert Qx\rVert=\lVert x\rVert\) for all \(x,y\in\R^p\): \(Q\) preserves inner products and
lengths.
\item If \(p=n\), then \(Q\) is invertible, \(Q^{-1}=Q'\) and \(QQ'=I_n\): the transposed rows of \(Q\) are
orthonormal too.
\end{enumerate}
\end{proposition}
\begin{proof}
(a) \((Qx)'(Qy)=x'Q'Qy=x'y\); take \(y=x\). (b) \(Q'\) is a left inverse of the square matrix \(Q\), so
\(Q^{-1}=Q'\) (Theorem~\ref{r2-thm-invertible}), and \(QQ'=QQ^{-1}=I_n\).
\end{proof}

\begin{theorem}[Gram--Schmidt; \citeat{hoffman-kunze}{\S8.2, Theorem~3 and its Corollary}, \citeat{abadir-magnus}{Exercise~3.49}]\label{r6-thm-gram-schmidt}
Let \(a_1,\dots,a_p\in\R^n\) be linearly independent. Put \(b_1=a_1\) and, for \(j=2,\dots,p\),
\[
b_j=a_j-t_{1j}b_1-t_{2j}b_2-\dots-t_{j-1,j}b_{j-1},\qquad t_{ij}=\frac{a_j'b_i}{\lVert b_i\rVert^2}\quad(i<j)
\]
Then \(b_1,\dots,b_p\) are nonzero and pairwise orthogonal, and
\(\Span\{b_1,\dots,b_j\}=\Span\{a_1,\dots,a_j\}\) for every \(j\). Hence the vectors
\(q_j=b_j/\lVert b_j\rVert\) (\(1\leqslant j\leqslant p\)) form an orthonormal basis of
\(\Span\{a_1,\dots,a_p\}\).
\end{theorem}
\begin{proof}
Induction on \(j\): \(b_1=a_1\ne0\), since an independent list does not contain \(0\). Let
\(j\geqslant2\), and let \(b_1,\dots,b_{j-1}\) be nonzero, pairwise orthogonal, with span
\(W_{j-1}=\Span\{a_1,\dots,a_{j-1}\}\); so \(b_j\) is defined.
\begin{itemize}
\item Orthogonal: for \(l<j\), \(b_l'b_j=b_l'a_j-\sum_{i<j}t_{ij}b_l'b_i=b_l'a_j-t_{lj}\lVert b_l\rVert^2=0\).
\item Nonzero: otherwise \(a_j=\sum_{i<j}t_{ij}b_i\in W_{j-1}\), and \(a_1,\dots,a_j\) would be dependent
(Proposition~\ref{r3-prop-dependence}(c)) \(\contradiction\)
\item Span: \(b_j\) is a combination of \(a_j\) and \(b_1,\dots,b_{j-1}\in W_{j-1}\), and
\(a_j=b_j+\sum_{i<j}t_{ij}b_i\); so each of the lists \(a_1,\dots,a_j\) and \(b_1,\dots,b_j\) lies in the
span of the other, and the two spans are equal (Theorem~\ref{r3-thm-span}).
\end{itemize}
Finally, the \(q_j\) are pairwise orthogonal unit vectors, hence independent
(Proposition~\ref{r6-prop-orthogonal-list}), and they span \(\Span\{b_1,\dots,b_p\}=\Span\{a_1,\dots,a_p\}\).
\end{proof}

\begin{corollary}[\citeat{hoffman-kunze}{\S8.2, Corollary to Theorem~3}, \citeat{harville}{Theorem~6.4.5}, \citeat{lang}{Ch.~V, Theorems~2.1, 2.3}]\label{r6-cor-onb}
Let \(W\) be a subspace of \(\R^n\) with \(\dim W=d\geqslant1\).
\begin{enumerate}
\item \(W\) has an orthonormal basis.
\item If \(q_1,\dots,q_p\in W\) are orthonormal and \(p<d\), there are \(q_{p+1},\dots,q_d\in W\) such that
\(q_1,\dots,q_d\) form an orthonormal basis of \(W\).
\item If \(Q\in\R^{n\times p}\) has orthonormal columns and \(p<n\), there is \(T\in\R^{n\times(n-p)}\) such that
\([Q\ T]\) is an orthogonal matrix. For every such \(T\), the columns of \(T\) form an orthonormal basis of
\(\mathcal C(Q)^\perp\), and \(QQ'+TT'=I_n\).
\end{enumerate}
\end{corollary}
\begin{proof}
(a) \(W\) has a basis \(a_1,\dots,a_d\) (Theorem~\ref{thm:r4-dimension}). By
Theorem~\ref{r6-thm-gram-schmidt}, the vectors \(q_j=b_j/\lVert b_j\rVert\) obtained from it form an orthonormal
basis of \(\Span\{a_1,\dots,a_d\}=W\).

(b)
\begin{itemize}
\item \(q_1,\dots,q_p\) are independent (Proposition~\ref{r6-prop-orthogonal-list}(b)), so they extend to a
basis \(q_1,\dots,q_p,a_{p+1},\dots,a_d\) of \(W\) (Theorem~\ref{thm:r4-counting}(b)).
\item Apply Theorem~\ref{r6-thm-gram-schmidt} to this list. Its first \(p\) outputs are \(b_j=q_j\): \(b_1=q_1\),
and if \(b_i=q_i\) for all \(i<j\leqslant p\), then \(t_{ij}=q_j'q_i/\lVert q_i\rVert^2=0\), so \(b_j=q_j\).
\item Hence the normalized outputs \(q_1,\dots,q_p,b_{p+1}/\lVert b_{p+1}\rVert,\dots,b_d/\lVert b_d\rVert\) form an
orthonormal basis of \(W\) (Theorem~\ref{r6-thm-gram-schmidt}), and its first \(p\) vectors are the given
ones.
\end{itemize}
(c)
\begin{itemize}
\item The columns \(q_1,\dots,q_p\) of \(Q\) are orthonormal (Definition~\ref{r6-def-orthonormal}). By (b) with
\(W=\R^n\), there are \(q_{p+1},\dots,q_n\) such that \(q_1,\dots,q_n\) form an orthonormal basis of \(\R^n\);
then \(T=[q_{p+1}\ \cdots\ q_n]\) makes \([Q\ T]\) an \(n\times n\) matrix with orthonormal columns, that is,
an orthogonal matrix.
\item Let \(T\) be any such matrix. Each column of \(T\) is orthogonal to every column of \(Q\), so it lies in
\(\mathcal C(Q)^\perp\) (Theorem~\ref{r6-thm-complement}(a)). Orthonormal vectors are independent
(Proposition~\ref{r6-prop-orthogonal-list}(b)); so \(\rank Q=p\) (Theorem~\ref{thm:r4-rank-rules}(b)),
\(\dim\mathcal C(Q)^\perp=n-p\) (Theorem~\ref{r6-thm-complement}(c)), and the \(n-p\) independent columns of
\(T\) form a basis of \(\mathcal C(Q)^\perp\) (Theorem~\ref{thm:r4-counting}(d)).
\item \([Q\ T][Q\ T]'=I_n\) (Proposition~\ref{r6-prop-orthonormal}(b)), and by the outer product expansion
(Theorem~\ref{r1-thm-views}(d)), \([Q\ T][Q\ T]'=q_1q_1'+\dots+q_nq_n'=QQ'+TT'\).\qedhere
\end{itemize}
\end{proof}

\begin{theorem}[QR decomposition; \citeat{searle-khuri}{Theorem~5.2}, \citeat{harville}{\S6.4d}, \citeat{abadir-magnus}{Exercises~7.34, 7.35}]\label{r6-thm-qr}
Let \(A=[a_1\ \cdots\ a_p]\in\R^{n\times p}\) have full column rank \(p\). Then \(A=QR\), where
\(Q\in\R^{n\times p}\) has orthonormal columns (\(Q'Q=I_p\)) and \(R\in\R^{p\times p}\) is upper triangular with
positive diagonal entries; moreover \(R=Q'A\).
\end{theorem}
\begin{proof}
\begin{itemize}
\item The columns of \(A\) are independent (Theorem~\ref{thm:r4-rank-rules}(b)). Apply
Theorem~\ref{r6-thm-gram-schmidt} to them, and let \(Q=[q_1\ \cdots\ q_p]\).
\item The construction gives \(a_j=b_j+\sum_{i<j}t_{ij}b_i=\sum_{i\leqslant j}r_{ij}q_i\), where
\(r_{jj}=\lVert b_j\rVert>0\) and \(r_{ij}=t_{ij}\lVert b_i\rVert\) for \(i<j\). Let \(R=(r_{ij})\), with \(r_{ij}=0\) for
\(i>j\). Then \(\col_j(QR)=Q\col_j(R)=\sum_{i\leqslant j}r_{ij}q_i=a_j\) (Theorem~\ref{r1-thm-views}(a),(b)), so
\(A=QR\).
\item \(Q'Q=I_p\), as the \(q_j\) are orthonormal; hence \(Q'A=Q'QR=R\).\qedhere
\end{itemize}
\end{proof}

\begin{example}[cf.\ \citeat{hoffman-kunze}{\S8.2, Example~12}]\label{r6-exm-qr}
Find the QR decomposition of
\[
A=\begin{pmatrix}1&3&2\\2&1&3\\0&2&4\\2&2&5\end{pmatrix}
\]

\textit{Solution.} Apply Gram--Schmidt (Theorem~\ref{r6-thm-gram-schmidt}) to the columns
\(a_1,a_2,a_3\).
\begin{itemize}
\item \(b_1=a_1=(1,2,0,2)'\), \(\lVert b_1\rVert^2=1+4+0+4=9\).
\item \(a_2'b_1=3+2+0+4=9\), so \(t_{12}=1\) and \(b_2=a_2-b_1=(3,1,2,2)'-(1,2,0,2)'=(2,-1,2,0)'\), with
\(\lVert b_2\rVert^2=4+1+4+0=9\).
\item \(a_3'b_1=2+6+0+10=18\) and \(a_3'b_2=4-3+8+0=9\), so \(t_{13}=2\), \(t_{23}=1\) and
\[
b_3=a_3-2b_1-b_2=\begin{pmatrix}2\\3\\4\\5\end{pmatrix}-\begin{pmatrix}2\\4\\0\\4\end{pmatrix}-\begin{pmatrix}2\\-1\\2\\0\end{pmatrix}
=\begin{pmatrix}-2\\0\\2\\1\end{pmatrix},\qquad\lVert b_3\rVert^2=4+0+4+1=9
\]
\item Check: \(b_1'b_2=2-2+0+0=0\), \(b_1'b_3=-2+0+0+2=0\), \(b_2'b_3=-4+0+4+0=0\).
\item All three lengths are \(3\), so \(q_j=\frac13b_j\), \(r_{jj}=3\), \(r_{12}=3t_{12}=3\), \(r_{13}=3t_{13}=6\)
and \(r_{23}=3t_{23}=3\):
\[
Q=\frac13\begin{pmatrix}1&2&-2\\2&-1&0\\0&2&2\\2&0&1\end{pmatrix},\qquad
R=\begin{pmatrix}3&3&6\\0&3&3\\0&0&3\end{pmatrix}
\]
\end{itemize}
As a check, \(R=Q'A\); for instance \(r_{23}=q_2'a_3=\frac13(4-3+8+0)=3\).
\end{example}

\begin{example}\label{r6-exm-extend}
Extend the orthonormal columns \(q_1,q_2,q_3\) of \(Q\) in Example~\ref{r6-exm-qr} to an orthonormal basis of
\(\R^4\).

\textit{Solution.} By Corollary~\ref{r6-cor-onb}(c) we need a unit vector \(q_4\) orthogonal to \(q_1,q_2,q_3\),
that is, a normalized \(v\ne0\) with \(b_1'v=b_2'v=b_3'v=0\):
\[
v_1+2v_2+2v_4=0,\qquad 2v_1-v_2+2v_3=0,\qquad -2v_1+2v_3+v_4=0
\]
\begin{itemize}
\item The third equation gives \(v_4=2v_1-2v_3\), and the second \(v_2=2v_1+2v_3\).
\item Substituting both into the first: \(v_1+(4v_1+4v_3)+(4v_1-4v_3)=9v_1=0\), so \(v_1=0\), \(v_2=2v_3\) and
\(v_4=-2v_3\).
\item With \(v_3=1\): \(v=(0,2,1,-2)'\), \(v'v=0+4+1+4=9\), so \(q_4=\frac13v\).
\end{itemize}
Hence
\[
[Q\ q_4]=\frac13\begin{pmatrix}1&2&-2&0\\2&-1&0&2\\0&2&2&1\\2&0&1&-2\end{pmatrix}
\]
is an orthogonal matrix. By Corollary~\ref{r6-cor-onb}(c), \(q_4\) alone is an orthonormal basis of
\(\mathcal C(Q)^\perp=\mathcal C(A)^\perp\) (\(\mathcal C(Q)=\mathcal C(A)\) by Theorem~\ref{r6-thm-gram-schmidt}),
and \(QQ'+q_4q_4'=I_4\).
\end{example}

\subsection*{Orthogonal matrices}

\begin{theorem}[\citeat{harville}{\S8.4}, \citeat{lang}{Ch.~VII, \S3}, \citeat{shakarchi}{Ch.~VII, \S3}]\label{r6-thm-orthogonal}
\begin{enumerate}
\item For \(Q\in\R^{n\times n}\) the following are equivalent:
\begin{enumerate}[label=(\roman*)]
\item \(Q'Q=I_n\), that is, \(Q\) is orthogonal;
\item \(QQ'=I_n\);
\item \(Q\) is invertible and \(Q^{-1}=Q'\);
\item the columns of \(Q\) are orthonormal;
\item the transposed rows of \(Q\) are orthonormal;
\item \((Qx)'(Qy)=x'y\) for all \(x,y\in\R^n\);
\item \(\lVert Qx\rVert=\lVert x\rVert\) for all \(x\in\R^n\).
\end{enumerate}
\item If \(P,Q\in\R^{n\times n}\) are orthogonal, so are \(PQ\), \(Q'\) and \(Q^{-1}\).
\item If \(Q\) is orthogonal, then \(\det Q=1\) or \(\det Q=-1\).
\item Every permutation matrix \(P_\varphi\) (Definition~\ref{r5d-def-permutation-matrix}) is orthogonal, and
\(\det P_\varphi=1\) or \(\det P_\varphi=-1\).
\end{enumerate}
\end{theorem}
\begin{proof}
(a)
\begin{itemize}
\item (i)\(\Rightarrow\)(ii),(iii): Proposition~\ref{r6-prop-orthonormal}(b).
\item (ii)\(\Rightarrow\)(iii): \(Q'\) is a right inverse of the square matrix \(Q\), so \(Q\) is invertible and
\(Q^{-1}=Q'\) (Theorem~\ref{r2-thm-invertible}). (iii)\(\Rightarrow\)(i): \(Q'Q=Q^{-1}Q=I_n\).
\item (i)\(\Leftrightarrow\)(iv): \((Q'Q)_{ij}=q_i'q_j\) for the columns \(q_1,\dots,q_n\) of \(Q\)
(Theorem~\ref{r1-thm-views}(d)); so \(Q'Q=I_n\) exactly when \(q_i'q_j\) is \(1\) for \(i=j\) and \(0\) for
\(i\ne j\) (Definition~\ref{r6-def-orthonormal}).
\item (ii)\(\Leftrightarrow\)(v): likewise, \((QQ')_{ij}=\row_i(Q)\row_j(Q)'\).
\item (i)\(\Rightarrow\)(vi): Proposition~\ref{r6-prop-orthonormal}(a). (vi)\(\Rightarrow\)(vii): take \(y=x\).
\item (vi)\(\Rightarrow\)(i): with \(x=e_i\) and \(y=e_j\), \((Q'Q)_{ij}=q_i'q_j=(Qe_i)'(Qe_j)=e_i'e_j\), which is
\(1\) for \(i=j\) and \(0\) otherwise (Theorem~\ref{r1-thm-views}(b),(d)).
\item (vii)\(\Rightarrow\)(vi): by polarization (Proposition~\ref{r6-prop-polarization}(a)) and
\(Q(x\pm y)=Qx\pm Qy\),
\[
4(Qx)'(Qy)=\lVert Q(x+y)\rVert^2-\lVert Q(x-y)\rVert^2=\lVert x+y\rVert^2-\lVert x-y\rVert^2=4x'y
\]
\end{itemize}
(b)
\begin{itemize}
\item By (iii), \(P^{-1}=P'\) and \(Q^{-1}=Q'\), so \((PQ)^{-1}=Q^{-1}P^{-1}=Q'P'=(PQ)'\)
(Theorems~\ref{r2-thm-inverse-props}(b) and~\ref{r1-thm-transpose}); by (iii), \(PQ\) is orthogonal.
\item \((Q')'Q'=QQ'=I_n\) by (ii), so \(Q'\) is orthogonal; and \(Q^{-1}=Q'\).
\end{itemize}
(c) \(1=\det I_n=\det(Q'Q)=\det Q'\cdot\det Q=(\det Q)^2\) (Definition~\ref{r5d-def-det}(d),
Theorems~\ref{r5d-thm-product}(b) and~\ref{r5d-thm-transpose}).

(d) \(P_\varphi'P_\varphi=I_n\) (Proposition~\ref{r5d-prop-permutation-matrix}(c)), so \(P_\varphi\) is orthogonal, and
\(\det P_\varphi=\pm1\) by (c).
\end{proof}
So multiplying by an orthogonal matrix keeps inner products, lengths and distances:
\(\lVert Qx-Qy\rVert=\lVert Q(x-y)\rVert=\lVert x-y\rVert\). It only changes the coordinates in which a data vector
is written.

\begin{example}[\citeat{harville}{\S8.4b; Ch.~8, Exercise~4(b)}, \citeat{searle-khuri}{\S5.4.1}]\label{r6-exm-helmert}
Let \(n\geqslant2\). For \(k=2,\dots,n\) let \(u_k=(1,\dots,1,-(k-1),0,\dots,0)'\in\R^n\): \(k-1\) entries \(1\),
then \(-(k-1)\) in position \(k\), then zeros. The \textbf{Helmert matrix} of order \(n\) is the \(n\times n\)
matrix \(K_n\) with rows
\[
p_1'=\frac1{\sqrt n}\,\one_n',\qquad p_k'=\frac1{\sqrt{k(k-1)}}\,u_k'\quad(k=2,\dots,n)
\]
For instance
\[
K_4=\begin{pmatrix}1/2&1/2&1/2&1/2\\1/\sqrt2&-1/\sqrt2&0&0\\1/\sqrt6&1/\sqrt6&-2/\sqrt6&0\\
1/\sqrt{12}&1/\sqrt{12}&1/\sqrt{12}&-3/\sqrt{12}\end{pmatrix}
\]
\begin{enumerate}
\item \(K_n\) is orthogonal. By Theorem~\ref{r6-thm-orthogonal}(a) it suffices that \(p_1,\dots,p_n\) are
orthonormal.
\begin{itemize}
\item \(\one_n'u_k=(k-1)-(k-1)=0\) for \(k\geqslant2\).
\item For \(2\leqslant j<k\): \(u_j\) is \(0\) outside positions \(1,\dots,j\), and these positions of \(u_k\)
hold \(1\) (as \(j\leqslant k-1\)); so \(u_j'u_k\) is the sum of the entries of \(u_j\), which is
\(\one_n'u_j=0\).
\item \(\one_n'\one_n=n\) and \(u_k'u_k=(k-1)\cdot1+(k-1)^2=k(k-1)\), so every \(p_k\) has length \(1\).
\end{itemize}
\item For data \(y\in\R^n\) with mean \(\bar y\), the first entry of \(K_ny\) is \(p_1'y=\frac1{\sqrt n}\one'y=\sqrt n\,\bar y\),
and for \(k\geqslant2\) its \(k\)th entry is the \textbf{contrast}
\(p_k'y=\bigl(y_1+\dots+y_{k-1}-(k-1)y_k\bigr)/\sqrt{k(k-1)}\), a combination of \(y_1,\dots,y_n\) whose
coefficients add up to \(0\).
\begin{itemize}
\item By (a) and Theorem~\ref{r6-thm-orthogonal}(a)(vii), \(\sum_iy_i^2=\lVert K_ny\rVert^2=n\bar y^2+\sum_{k=2}^n(p_k'y)^2\).
As \(\sum_i(y_i-\bar y)^2=\sum_iy_i^2-n\bar y^2\) (Example~\ref{r1-exm-centering}), the sum of squared
deviations is the sum of the squares of the \(n-1\) contrasts:
\[
\sum_{i=1}^n(y_i-\bar y)^2=\sum_{k=2}^n(p_k'y)^2
\]
\item In matrix form: \(K_n'K_n=I_n\) reads \(p_1p_1'+\dots+p_np_n'=I_n\) (Theorem~\ref{r1-thm-views}(d)), and
\(p_1p_1'=\frac1n\one\one'=\bar J_n\); hence \(p_2p_2'+\dots+p_np_n'=I_n-\bar J_n=C_n\).
\end{itemize}
\item For \(n=4\) and \(y=(5,3,1,7)'\): \(\bar y=4\), \(p_1'y=16/2=8=\sqrt4\cdot4\), and
\[
p_2'y=\frac{5-3}{\sqrt2}=\sqrt2,\qquad p_3'y=\frac{5+3-2}{\sqrt6}=\sqrt6,\qquad p_4'y=\frac{5+3+1-21}{\sqrt{12}}=-\sqrt{12}
\]
so the contrasts give \(2+6+12=20=1+1+9+9=\sum_i(y_i-4)^2\), and \(\lVert y\rVert^2=25+9+1+49=84=64+20\).
\end{enumerate}
\end{example}

\begin{example}[\citeat{harville-solutions}{Ch.~10, Exercise~3}]\label{r6-exm-reflection}
\begin{enumerate}
\item If \(A\in\R^{n\times n}\) is symmetric and idempotent, then \(I_n-2A\) is symmetric and orthogonal:
\((I-2A)'(I-2A)=(I-2A)^2=I-4A+4A^2=I\).
\item Let \(h\in\R^n\), \(h\ne0\). The matrix \(A=hh'/(h'h)\) is symmetric and idempotent, since \((hh')'=hh'\) and
\(hh'hh'=h(h'h)h'\). So the \textbf{reflection}
\[
F_h=I_n-\frac2{h'h}\,hh'
\]
is symmetric and orthogonal, and \(F_h^2=F_h'F_h=I_n\). It reverses \(h\) and fixes every vector orthogonal to
\(h\): \(F_hh=h-2h=-h\), and \(F_hx=x-\frac{2h'x}{h'h}h=x\) if \(h'x=0\). By Theorem~\ref{r5d-thm-det-lemma}
with \(u=-\frac2{h'h}h\) and \(v=h\), \(\det F_h=1+v'u=1-2=-1\).
\item For \(h=e_r-e_s\) (\(r\ne s\)), \(h'h=2\) and \(F_h=I_n-(e_r-e_s)(e_r-e_s)'\) is the interchange matrix
\(P_{rs}\) of Definition~\ref{r2-def-elementary}: its entries \((r,r)\) and \((s,s)\) are \(1-1=0\), its entries
\((r,s)\) and \((s,r)\) are \(1\), and the others are those of \(I_n\). Interchanging two coordinates is a
reflection.
\item For \(h=(1,1,2,2)'\), \(h'h=10\) and
\[
F_h=I_4-\frac15\begin{pmatrix}1&1&2&2\\1&1&2&2\\2&2&4&4\\2&2&4&4\end{pmatrix}
=\frac15\begin{pmatrix}4&-1&-2&-2\\-1&4&-2&-2\\-2&-2&1&-4\\-2&-2&-4&1\end{pmatrix}
\]
Its columns are orthonormal (for instance \(16+1+4+4=25\) and \(-4-4+4+4=0\)), and
\(F_hh=\frac15(4-1-4-4,\ -1+4-4-4,\ -2-2+2-8,\ -2-2-8+2)'=-(1,1,2,2)'\).
\end{enumerate}
\end{example}

\subsection*{Orthogonal projections}

\begin{definition}[\citeat{hoffman-kunze}{\S8.2}, \citeat{tu-clm}{Remark~2.10}]\label{r6-def-projection}
Let \(W\) be a subspace of \(\R^n\) and \(y\in\R^n\). By Theorem~\ref{r6-thm-complement}(c),
\(y=\hat w+(y-\hat w)\) with \(\hat w\in W\) and \(y-\hat w\in W^\perp\) for exactly one vector \(\hat w\), the
\textbf{orthogonal projection} of \(y\) onto \(W\). Thus \(\hat w\) is the only vector of \(W\) with
\(y-\hat w\perp W\).
\end{definition}

\begin{theorem}[\citeat{hoffman-kunze}{\S8.2, Theorem~4}, \citeat{tu-clm}{Theorem~2.5}, \citeat{harville}{Theorem~12.4.1}]\label{r6-thm-best-approximation}
Let \(\hat w\) be the orthogonal projection of \(y\in\R^n\) onto a subspace \(W\) of \(\R^n\). Then
\[
\lVert y-w\rVert^2=\lVert y-\hat w\rVert^2+\lVert \hat w-w\rVert^2\qquad(w\in W)
\]
Hence \(\hat w\) is the vector of \(W\) closest to \(y\): \(\lVert y-w\rVert>\lVert y-\hat w\rVert\) for every \(w\in W\)
with \(w\ne\hat w\). Equivalently, a vector \(w_0\in W\) satisfies \(\lVert y-w_0\rVert\leqslant\lVert y-w\rVert\) for
every \(w\in W\) if and only if \(y-w_0\perp W\), and only \(w_0=\hat w\) does.
\end{theorem}
\begin{proof}
\(y-w=(y-\hat w)+(\hat w-w)\), where \(\hat w-w\in W\) and \(y-\hat w\perp W\); apply Pythagoras
(Proposition~\ref{r6-prop-pythagoras}). If \(w\ne\hat w\), then \(\lVert \hat w-w\rVert^2>0\). So \(w_0\in W\) is a
closest vector if and only if \(w_0=\hat w\), and by Definition~\ref{r6-def-projection} this holds if and only if
\(y-w_0\perp W\).
\end{proof}
The number \(\lVert y-\hat w\rVert=\min_{w\in W}\lVert y-w\rVert\) is the \textbf{distance} from \(y\) to \(W\).

\begin{theorem}[\citeat{hoffman-kunze}{\S8.2, Theorem~4(iii), Corollary after Theorem~5}, \citeat{harville}{Theorem~12.1.3, \S12.2a, Theorem~12.5.8}]\label{r6-thm-projection-onb}
Let \(W\) be a subspace of \(\R^n\), \(y\in\R^n\), and \(\hat w\) the orthogonal projection of \(y\) onto \(W\).
\begin{enumerate}
\item If \(u_1,\dots,u_k\) (\(k\geqslant1\)) form an orthogonal basis of \(W\), that is, a basis of pairwise
orthogonal vectors, then
\[
\hat w=\frac{u_1'y}{u_1'u_1}\,u_1+\dots+\frac{u_k'y}{u_k'u_k}\,u_k
\]
In particular, if \(q_1,\dots,q_k\) form an orthonormal basis of \(W\) and \(Q=[q_1\ \cdots\ q_k]\), then
\(\hat w=(q_1'y)q_1+\dots+(q_k'y)q_k=QQ'y\).
\item \(\hat w=y\) if and only if \(y\in W\).
\item (\textbf{Bessel's inequality}) In (a),
\(\lVert\hat w\rVert^2=\sum_{j=1}^k(u_j'y)^2/u_j'u_j\leqslant\lVert y\rVert^2\), with equality if and only if \(y\in W\);
for an orthonormal basis, \(\sum_{j=1}^k(q_j'y)^2\leqslant\lVert y\rVert^2\).
\item \(y-\hat w\) is the orthogonal projection of \(y\) onto \(W^\perp\).
\end{enumerate}
\end{theorem}
\begin{proof}
By Definition~\ref{r6-def-projection}, \(\hat w\) is the only vector of \(W\) with \(y-\hat w\perp W\).

(a)
\begin{itemize}
\item The \(u_j\) are nonzero (Proposition~\ref{r3-prop-dependence}(a)); put \(c_j=u_j'y/u_j'u_j\) and
\(w_0=c_1u_1+\dots+c_ku_k\in W\).
\item For \(i=1,\dots,k\), \(u_i'(y-w_0)=u_i'y-\sum_jc_ju_i'u_j=u_i'y-c_iu_i'u_i=0\). As \(u_1,\dots,u_k\) span \(W\),
\(y-w_0\perp W\) (Theorem~\ref{r6-thm-complement}(a)); so \(w_0=\hat w\).
\item For an orthonormal basis, \(q_j'q_j=1\), and \(\sum_j(q_j'y)q_j=Q(Q'y)\), as \(Q'y\) has the entries
\(q_j'y\) (Theorem~\ref{r1-thm-views}(b),(d)).
\end{itemize}
(b) If \(\hat w=y\), then \(y\in W\). Conversely, if \(y\in W\), then \(y-y=0\perp W\), so \(\hat w=y\).

(c)
\begin{itemize}
\item \(y=\hat w+(y-\hat w)\) with \(\hat w\perp y-\hat w\), so
\(\lVert y\rVert^2=\lVert\hat w\rVert^2+\lVert y-\hat w\rVert^2\geqslant\lVert\hat w\rVert^2\)
(Proposition~\ref{r6-prop-pythagoras}).
\item The vectors \(c_ju_j\) are pairwise orthogonal, so
\(\lVert\hat w\rVert^2=\sum_jc_j^2\,u_j'u_j=\sum_j(u_j'y)^2/u_j'u_j\) (Proposition~\ref{r6-prop-polarization}(b)).
\item Equality holds if and only if \(\lVert y-\hat w\rVert=0\), that is, \(y=\hat w\), that is, \(y\in W\) by (b).
\end{itemize}
(d) \(y-\hat w\in W^\perp\), and \(y-(y-\hat w)=\hat w\in W=(W^\perp)^\perp\) (Theorem~\ref{r6-thm-perp-perp}(a)),
that is, \(y-(y-\hat w)\perp W^\perp\). By Definition~\ref{r6-def-projection} for the subspace \(W^\perp\),
\(y-\hat w\) is the orthogonal projection of \(y\) onto \(W^\perp\).
\end{proof}

\begin{example}[cf.\ \citeat{hoffman-kunze}{\S8.2, Example~14}, \citeat{harville}{\S\S12.2b, 12.5c}]\label{r6-exm-projection-onb}
Let \(W=\Span\{a_1,a_2\}\subseteq\R^4\), where \(a_1=(1,2,0,2)'\) and \(a_2=(3,1,2,2)'\) are the first two columns
of \(A\) in Example~\ref{r6-exm-qr}, and let \(y=(-2,7,1,3)'\). Find the orthogonal projections of \(y\) onto
\(W\) and onto \(W^\perp\), and the distance from \(y\) to the nearest vector of \(W\).

\textit{Solution.} By Example~\ref{r6-exm-qr} and Theorem~\ref{r6-thm-gram-schmidt}, \(b_1=(1,2,0,2)'\) and
\(b_2=(2,-1,2,0)'\) form an orthogonal basis of \(W\), with \(b_1'b_1=b_2'b_2=9\).
\begin{itemize}
\item \(b_1'y=-2+14+0+6=18\) and \(b_2'y=-4-7+2+0=-9\). By Theorem~\ref{r6-thm-projection-onb}(a),
\[
\hat w=\frac{18}9\,b_1-\frac99\,b_2=\begin{pmatrix}2\\4\\0\\4\end{pmatrix}-\begin{pmatrix}2\\-1\\2\\0\end{pmatrix}
=\begin{pmatrix}0\\5\\-2\\4\end{pmatrix}
\]
\item By Theorem~\ref{r6-thm-projection-onb}(d), the projection of \(y\) onto \(W^\perp\) is
\(y-\hat w=(-2,2,3,-1)'\). Check: \(b_1'(y-\hat w)=-2+4+0-2=0\) and \(b_2'(y-\hat w)=-4-2+6+0=0\).
\item The vector of \(W\) nearest to \(y\) is \(\hat w\) (Theorem~\ref{r6-thm-best-approximation}), at distance
\(\lVert y-\hat w\rVert=\sqrt{4+4+9+1}=\sqrt{18}=3\sqrt2\).
\item Bessel's inequality (Theorem~\ref{r6-thm-projection-onb}(c)):
\(\lVert\hat w\rVert^2=18^2/9+9^2/9=36+9=45\leqslant\lVert y\rVert^2=4+49+1+9=63\), and \(63=45+18\).
\end{itemize}
\end{example}

\begin{theorem}[\citeat{tu-clm}{Definition~2.6, Theorem~2.9(i)--(iii), Remark~2.10}, \citeat{harville}{Theorem~12.3.1}]\label{r6-thm-hat}
Let \(A\in\R^{n\times p}\) have full column rank \(p\). Then \(A'A\) is invertible. Put
\[
P=A(A'A)^{-1}A'\in\R^{n\times n}
\]
\begin{enumerate}
\item For every \(y\in\R^n\), \(Py\) is the orthogonal projection of \(y\) onto \(\mathcal C(A)\), and
\(y=Py+(I_n-P)y\) is the decomposition of \(y\) along \(\R^n=\mathcal C(A)\oplus\mathcal C(A)^\perp\).
\item \(P'=P=P^2\) and \((I_n-P)'=I_n-P=(I_n-P)^2\): both are symmetric and idempotent.
\item \(PA=A\), \((I_n-P)A=0\) and \(P(I_n-P)=(I_n-P)P=0\).
\end{enumerate}
\end{theorem}
\begin{proof}
\(\rank(A'A)=\rank A=p\) (Theorem~\ref{thm:r4-rank-ata}(b)), so the \(p\times p\) matrix \(A'A\) is
invertible (Theorem~\ref{thm:r4-rank-rules}(d)).

(a) \(Py=A\bigl((A'A)^{-1}A'y\bigr)\in\mathcal C(A)\), and \(A'(y-Py)=A'y-(A'A)(A'A)^{-1}A'y=0\), so
\((I_n-P)y=y-Py\in\mathcal C(A)^\perp\) (Theorem~\ref{r6-thm-complement}(a)). By
Definition~\ref{r6-def-projection}, \(Py\) is the projection.

(b) \(A'A\) is symmetric, hence so is its inverse:
\(((A'A)^{-1})'=((A'A)')^{-1}=(A'A)^{-1}\) (Theorem~\ref{r2-thm-inverse-props}(c)). So
\(P'=A((A'A)^{-1})'A'=P\) and \(P^2=A(A'A)^{-1}(A'A)(A'A)^{-1}A'=P\). Then \((I_n-P)'=I_n-P'=I_n-P\) and
\((I_n-P)^2=I_n-2P+P^2=I_n-P\).

(c) \(PA=A(A'A)^{-1}A'A=A\), \((I_n-P)A=A-PA=0\), and \(P(I_n-P)=P-P^2=0=(I_n-P)P\).
\end{proof}

\begin{example}[\citeat{tu-clm}{Definition~6.1, Remark~6.3}, \citeat{abadir-magnus}{Exercise~8.71(a)}, \citeat{hoffman-kunze}{\S8.2, Example~14}]\label{r6-exm-centering}
For \(A=\one_n\) (\(p=1\)), \(A'A=\one'\one=n\), so the matrix \(P=A(A'A)^{-1}A'\) of Theorem~\ref{r6-thm-hat} is
\(\frac1n\one\one'=\bar J_n\), and \(I_n-P=I_n-\bar J_n=C_n\), the centering matrix of Definition~\ref{r1-def-ones}:
by Theorem~\ref{r6-thm-hat}(a), \(\bar J_ny=\bar y\one\) (Example~\ref{r3-exm-constant-centered}) is the orthogonal
projection of \(y\) onto \(\Span\{\one\}\), and \(C_ny=y-\bar y\one\) (Example~\ref{r1-exm-centering}). The centered vectors form \(W_0=\{x:\one'x=0\}=\Span\{\one\}^\perp\)
(Theorem~\ref{r6-thm-complement}(a)), so the decomposition \(y=\bar y\one+C_ny\) along \(\R^n=\Span\{\one\}\oplus W_0\)
of Example~\ref{r3-exm-constant-centered} is the orthogonal one, and \(\bar y\one\) is the constant vector closest to
\(y\) (Theorem~\ref{r6-thm-best-approximation}): \(\sum_i(y_i-c)^2=\lVert y-c\one\rVert^2>\lVert y-\bar y\one\rVert^2
=\sum_i(y_i-\bar y)^2\) for every \(c\ne\bar y\). For instance, \(y=(9,3,5,3)'\) has \(\bar y=5\) and
\(y=5\,\one_4+(4,-2,0,-2)'\), with \(\one_4'(4,-2,0,-2)'=0\). Likewise, for any nonzero \(a\in\R^n\) the projection of
\(y\) onto \(\Span\{a\}\) is \(a(a'a)^{-1}a'y=\frac{a'y}{a'a}\,a\) (Theorem~\ref{r6-thm-hat}(a) with the matrix \(a\), which
has full column rank \(1\) as \(a\ne0\)).
\end{example}

\begin{example}\label{r6-exm-projection}
Let \(a_1=\one_4\), \(a_2=(2,0,2,0)'\), \(a_3=(3,3,1,1)'\), \(A=[a_1\ a_2\ a_3]\in\R^{4\times3}\), \(W=\mathcal C(A)\) and
\(y=(9,3,5,3)'\). Find \((A'A)^{-1}\), the matrix \(P=A(A'A)^{-1}A'\) of Theorem~\ref{r6-thm-hat}, the vector of \(W\)
closest to \(y\), and the distance from \(y\) to \(W\).

\textit{Solution.}
\begin{itemize}
\item The \((j,l)\) entry of \(A'A\) is \(a_j'a_l\) (Theorem~\ref{r1-thm-views}(d)), so
\[
A'A=\begin{pmatrix}4&4&8\\4&8&8\\8&8&20\end{pmatrix}
\]
It is symmetric (Theorem~\ref{r1-thm-transpose}), so its cofactors (Definition~\ref{r5d-def-cofactor}) satisfy \(c_{ji}=c_{ij}\)
(Remark~\ref{r5d-rem-symmetric-cofactors}), and six of them are needed (each \(2\times2\) determinant along its
row \(1\), Theorem~\ref{r5d-thm-expansion}):
\[
\begin{aligned}
c_{11}&=\det\begin{pmatrix}8&8\\8&20\end{pmatrix}=96,&
c_{12}&=-\det\begin{pmatrix}4&8\\8&20\end{pmatrix}=-16,&
c_{13}&=\det\begin{pmatrix}4&8\\8&8\end{pmatrix}=-32\\
c_{22}&=\det\begin{pmatrix}4&8\\8&20\end{pmatrix}=16,&
c_{23}&=-\det\begin{pmatrix}4&4\\8&8\end{pmatrix}=0,&
c_{33}&=\det\begin{pmatrix}4&4\\4&8\end{pmatrix}=16
\end{aligned}
\]
Along row \(1\) (Theorem~\ref{r5d-thm-expansion}), \(\det(A'A)=4\cdot96+4\cdot(-16)+8\cdot(-32)=64\). Hence, by
Theorem~\ref{r5d-thm-adjugate},
\[
(A'A)^{-1}=\frac1{64}\begin{pmatrix}96&-16&-32\\-16&16&0\\-32&0&16\end{pmatrix}
=\frac14\begin{pmatrix}6&-1&-2\\-1&1&0\\-2&0&1\end{pmatrix}
\]
As \(A'A\) is invertible, \(\rank A=\rank(A'A)=3\) (Theorems~\ref{thm:r4-rank-ata}(b) and~\ref{thm:r4-rank-rules}(d)):
\(A\) has full column rank, and Theorem~\ref{r6-thm-hat} applies.
\item Then
\[
(A'A)^{-1}A'=\frac14\begin{pmatrix}6&-1&-2\\-1&1&0\\-2&0&1\end{pmatrix}\begin{pmatrix}1&1&1&1\\2&0&2&0\\3&3&1&1\end{pmatrix}
=\frac14\begin{pmatrix}-2&0&2&4\\1&-1&1&-1\\1&1&-1&-1\end{pmatrix}
\]
\[
P=A(A'A)^{-1}A'=\begin{pmatrix}1&2&3\\1&0&3\\1&2&1\\1&0&1\end{pmatrix}\frac14\begin{pmatrix}-2&0&2&4\\1&-1&1&-1\\1&1&-1&-1\end{pmatrix}
=\frac14\begin{pmatrix}3&1&1&-1\\1&3&-1&1\\1&-1&3&1\\-1&1&1&3\end{pmatrix}
\]
\item By Theorem~\ref{r6-thm-hat}(a), \(Py\) is the orthogonal projection of \(y\) onto \(W\), so it is the vector of
\(W\) closest to \(y\) (Theorem~\ref{r6-thm-best-approximation}):
\[
Py=\frac14\begin{pmatrix}27+3+5-3\\9+9-5+3\\9-3+15+3\\-9+3+5+9\end{pmatrix}=\begin{pmatrix}8\\4\\6\\2\end{pmatrix},\qquad
(I_4-P)y=y-Py=\begin{pmatrix}1\\-1\\-1\\1\end{pmatrix}
\]
The distance from \(y\) to \(W\) is \(\lVert y-Py\rVert=\sqrt{1+1+1+1}=2\). Check: \(A'(I_4-P)y=(1-1-1+1,\ 2-2,\ 3-3-1+1)'=0\),
and \((Py)'(I_4-P)y=8-4-6+2=0\).
\item Check through \(W^\perp\): \(h=(1,-1,-1,1)'\) is orthogonal to \(a_1\), \(a_2\), \(a_3\) (\(1-1-1+1=0\), \(2-2=0\),
\(3-3-1+1=0\)), so \(h\in W^\perp\) (Theorem~\ref{r6-thm-complement}(a)). As \(\dim W=\rank A=3\)
(Definition~\ref{def:r4-rank}), \(\dim W^\perp=4-3=1\) (Theorem~\ref{r6-thm-complement}(c)), and the nonzero \(h\)
spans \(W^\perp\) (Theorem~\ref{thm:r4-counting}(d)). By Theorem~\ref{r6-thm-projection-onb}(d) and the last sentence
of Example~\ref{r6-exm-centering}, \((I_4-P)y=y-Py\) is the projection \(\frac14(h'y)h=\frac14hh'y\) of \(y\) onto
\(W^\perp=\Span\{h\}\), for every \(y\in\R^4\); so (Theorem~\ref{r1-thm-views}(b), with \(y=e_j\))
\[
I_4-P=\frac14\begin{pmatrix}1&-1&-1&1\\-1&1&1&-1\\-1&1&1&-1\\1&-1&-1&1\end{pmatrix}=\frac14hh'
\]
which agrees with \(P\) above. Pythagoras (Proposition~\ref{r6-prop-pythagoras}): \(\lVert y\rVert^2=81+9+25+9=124\) and
\(\lVert Py\rVert^2+\lVert(I_4-P)y\rVert^2=(64+16+36+4)+4=124\).
\end{itemize}
\end{example}

\subsection*{Idempotent matrices: rank equals trace}

\begin{theorem}[\citeat{tu-clm}{proof of Lemma~2.8}, \citeat{abadir-magnus}{Exercise~8.53}, \citeat{bapat}{Result~3.11}, \citeat{harville}{Theorems~10.2.1, 12.3.9}]\label{r6-thm-idempotent}
Let \(A\in\R^{n\times n}\) be idempotent, \(A^2=A\).
\begin{enumerate}
\item \(Av=v\) for every \(v\in\mathcal C(A)\), and \(\R^n=\mathcal N(A)\oplus\mathcal C(A)\): every \(v\in\R^n\) is
\(v=(v-Av)+Av\) with \(v-Av\in\mathcal N(A)\) and \(Av\in\mathcal C(A)\).
\item \(I_n-A\) is idempotent.
\item \(\rank A=\tr A\).
\item If moreover \(A'=A\), then \(Av\) is the orthogonal projection of \(v\) onto \(\mathcal C(A)\), for every
\(v\in\R^n\).
\end{enumerate}
\end{theorem}
\begin{proof}
(a)
\begin{itemize}
\item If \(v=Aw\), then \(Av=A^2w=Aw=v\). For every \(v\), \(A(v-Av)=Av-A^2v=0\).
\item If \(v\in\mathcal N(A)\cap\mathcal C(A)\), then \(v=Av=0\). So the sum \(\mathcal N(A)+\mathcal C(A)\) is direct
(Definition~\ref{r3-def-sum}), and it is \(\R^n\) by the displayed decomposition.
\end{itemize}
(b) \((I-A)^2=I-2A+A^2=I-A\).

(c) If \(A=0\), then \(\rank A=0=\tr A\). Let \(r=\rank A\geqslant1\).
\begin{itemize}
\item Take a full-rank factorization \(A=BC\), with \(B\in\R^{n\times r}\), \(C\in\R^{r\times n}\), and
\(L\in\R^{r\times n}\), \(K\in\R^{n\times r}\) such that \(LB=I_r\) and \(CK=I_r\)
(Theorem~\ref{thm:r4-full-rank-factorization}(a),(c)).
\item \(A^2=A\) reads \(BCBC=BC\). Multiplying on the left by \(L\) and on the right by \(K\),
\(L(BCBC)K=(LB)(CB)(CK)=CB\) and \(L(BC)K=(LB)(CK)=I_r\); so \(CB=I_r\).
\item Hence \(\tr A=\tr(BC)=\tr(CB)=\tr I_r=r\) (Theorem~\ref{r1-thm-trace}(b)).
\end{itemize}
(d) By Corollary~\ref{r6-cor-four-subspaces}, \(\mathcal N(A)=\mathcal C(A')^\perp=\mathcal C(A)^\perp\). So in (a),
\(Av\in\mathcal C(A)\) and \(v-Av\in\mathcal C(A)^\perp\): by Definition~\ref{r6-def-projection}, \(Av\) is the
orthogonal projection of \(v\) onto \(\mathcal C(A)\).
\end{proof}

\begin{corollary}[\citeat{tu-clm}{Theorem~2.9(iv)}, \citeat{abadir-magnus}{Exercise~8.68(c)}]\label{r6-cor-trace-hat}
Let \(A\in\R^{n\times p}\) have full column rank, and let \(P=A(A'A)^{-1}A'\) (Theorem~\ref{r6-thm-hat}). Then
\(\rank P=\tr P=p=\rank A\) and \(\rank(I_n-P)=\tr(I_n-P)=n-p\). In particular \(\rank C_n=\tr C_n=n-1\).
\end{corollary}
\begin{proof}
By Theorem~\ref{r1-thm-trace}(a),(b),
\(\tr P=\tr\bigl(A(A'A)^{-1}A'\bigr)=\tr\bigl((A'A)^{-1}A'A\bigr)=\tr I_p=p\) and \(\tr(I_n-P)=\tr I_n-\tr P=n-p\).
Both matrices are idempotent (Theorem~\ref{r6-thm-hat}(b)), so their ranks equal their traces
(Theorem~\ref{r6-thm-idempotent}(c)). For \(A=\one_n\), \(I_n-P=C_n\) (Example~\ref{r6-exm-centering}).
\end{proof}
In Example~\ref{r6-exm-projection}, \(\tr P=\frac14(3+3+3+3)=3\) and \(\tr(I_4-P)=\frac14(1+1+1+1)=1\); indeed every
column of \(I_4-P=\frac14hh'\) is a nonzero multiple of \(h=(1,-1,-1,1)'\), so \(\mathcal C(I_4-P)=\Span\{h\}\) and
\(\rank(I_4-P)=1\) (Definition~\ref{def:r4-rank} and Proposition~\ref{r3-prop-dependence}(a)).

\subsection*{Projection matrices}

\begin{theorem}[\citeat{harville}{Theorem~12.3.1, Corollary~12.3.2, \S12.2a, Theorems~12.3.4, 12.3.8, Corollary~12.5.10}]\label{r6-thm-projection-matrix}
Let \(W\) be a subspace of \(\R^n\).
\begin{enumerate}
\item There is exactly one matrix \(P\in\R^{n\times n}\) such that \(Py\) is the orthogonal projection of \(y\) onto
\(W\) for every \(y\in\R^n\). It is the \textbf{(orthogonal) projection matrix} for \(W\), written \(P_W\). For
\(W=\{0\}\), \(P_W=0\).
\item If \(A\in\R^{n\times p}\) has full column rank and \(\mathcal C(A)=W\), then \(P_W=A(A'A)^{-1}A'\). If the
columns of \(Q\in\R^{n\times k}\) form an orthonormal basis of \(W\), then \(P_W=QQ'\).
\item \(P_W\) is symmetric and idempotent; \(P_Ww=w\) for \(w\in W\) and \(P_Wv=0\) for \(v\in W^\perp\);
\(\mathcal C(P_W)=W\) and \(\rank P_W=\tr P_W=\dim W\); \(\lVert P_Wy\rVert^2=y'P_Wy\) for \(y\in\R^n\); and
\(P_{W^\perp}=I_n-P_W\).
\end{enumerate}
\end{theorem}
\begin{proof}
(a)
\begin{itemize}
\item Uniqueness: if \(Py=\tilde Py\) for every \(y\), then \(\col_j(P)=Pe_j=\tilde Pe_j=\col_j(\tilde P)\) for every
\(j\) (Theorem~\ref{r1-thm-views}(b)), so \(P=\tilde P\).
\item Existence for \(W=\{0\}\): \(y=0+y\) with \(0\in W\) and \(y\perp W\), so the projection of every \(y\) is
\(0=0y\), and \(P=0\) has the property.
\item Existence for \(W\ne\{0\}\): \(W\) has a basis \(a_1,\dots,a_p\) (\(p\geqslant1\); Theorem~\ref{thm:r4-dimension}),
so \(A=[a_1\ \cdots\ a_p]\) has full column rank (Theorem~\ref{thm:r4-rank-rules}(b)) and \(\mathcal C(A)=W\).
By Theorem~\ref{r6-thm-hat}(a), \(P=A(A'A)^{-1}A'\) has the property.
\end{itemize}
(b) By Theorem~\ref{r6-thm-hat}(a), every \(A\) of full column rank with \(\mathcal C(A)=W\) has the property in (a), so
\(A(A'A)^{-1}A'=P_W\) by the uniqueness in (a). If the columns of \(Q\) form an orthonormal basis of \(W\), then
\(\mathcal C(Q)=W\), \(Q\) has full column rank (Proposition~\ref{r6-prop-orthogonal-list}(b),
Theorem~\ref{thm:r4-rank-rules}(b)) and \(Q'Q=I_k\) (Definition~\ref{r6-def-orthonormal}), so
\(P_W=Q(Q'Q)^{-1}Q'=QQ'\).

(c)
\begin{itemize}
\item For \(w\in W\), \(w=w+0\) with \(0\in W^\perp\); for \(v\in W^\perp\), \(v=0+v\) with \(0\in W\). By the
uniqueness in Definition~\ref{r6-def-projection}, \(P_Ww=w\) and \(P_Wv=0\).
\item If \(W=\{0\}\), then \(P_W=0\) is symmetric and idempotent, \(\mathcal C(0)=\{0\}\) and
\(\rank0=\tr0=0\). Let \(W\ne\{0\}\), and \(P_W=A(A'A)^{-1}A'\) as in (b).
\item \(P_W\) is symmetric and idempotent, and \(P_WA=A\) (Theorem~\ref{r6-thm-hat}(b),(c)).
\item \(\mathcal C(P_W)=\mathcal C\bigl(A(A'A)^{-1}A'\bigr)\subseteq\mathcal C(A)\) and
\(\mathcal C(A)=\mathcal C(P_WA)\subseteq\mathcal C(P_W)\) (Theorem~\ref{thm:r4-rank-product}(a)); so
\(\mathcal C(P_W)=W\), and \(\rank P_W=\dim W=\tr P_W\) (Definition~\ref{def:r4-rank},
Theorem~\ref{r6-thm-idempotent}(c)).
\item For every \(W\) and \(y\in\R^n\), as \(P_W\) is symmetric and idempotent,
\(\lVert P_Wy\rVert^2=(P_Wy)'(P_Wy)=y'P_W'P_Wy=y'P_W^2y=y'P_Wy\) (Theorem~\ref{r1-thm-transpose}).
\item For every \(W\): \((I_n-P_W)y=y-P_Wy\) is the orthogonal projection of \(y\) onto \(W^\perp\) for every \(y\)
(Theorem~\ref{r6-thm-projection-onb}(d)); so \(P_{W^\perp}=I_n-P_W\) by (a) for \(W^\perp\).\qedhere
\end{itemize}
\end{proof}
By (a) and (b), \(QQ'\) does not depend on the choice of the orthonormal basis of \(W\).

\begin{corollary}\label{r6-cor-qr-projection}
Let \(A\in\R^{n\times p}\) have full column rank \(p\), and let \(A=QR\) as in Theorem~\ref{r6-thm-qr}. Then
\(\mathcal C(A)=\mathcal C(Q)\) and \(P_{\mathcal C(A)}=QQ'\).
\end{corollary}
\begin{proof}
\begin{itemize}
\item \(\col_j(A)=Q\col_j(R)\in\mathcal C(Q)\) (Theorem~\ref{r1-thm-views}(a),(b)), so \(\mathcal C(A)\subseteq\mathcal C(Q)\).
\item \(R\) is upper triangular with positive diagonal entries, so \(R\) is invertible
(Theorem~\ref{r2-thm-triangular}(a)). Then \(Q=AR^{-1}\), so \(\col_j(Q)=A\col_j(R^{-1})\in\mathcal C(A)\) for every \(j\),
and \(\mathcal C(Q)\subseteq\mathcal C(A)\).
\item The columns of \(Q\) are orthonormal, hence linearly independent
(Proposition~\ref{r6-prop-orthogonal-list}(b)), and they span \(\mathcal C(Q)=\mathcal C(A)\). So they form an
orthonormal basis of \(\mathcal C(A)\), and \(P_{\mathcal C(A)}=QQ'\) by
Theorem~\ref{r6-thm-projection-matrix}(b).\qedhere
\end{itemize}
\end{proof}

\begin{example}\label{r6-exm-projection-matrix}
For \(W=\Span\{a_1,a_2\}\) of Example~\ref{r6-exm-projection-onb}, find \(P_W\) in two ways, and \(P_{W^\perp}\).

\textit{Solution.}
\begin{itemize}
\item By Example~\ref{r6-exm-qr}, \(q_1=\frac13b_1\) and \(q_2=\frac13b_2\) form an orthonormal basis of \(W\).
By Theorem~\ref{r6-thm-projection-matrix}(b) and the outer product expansion (Theorem~\ref{r1-thm-views}(d)),
\(P_W=q_1q_1'+q_2q_2'=\frac19(b_1b_1'+b_2b_2')\):
\[
P_W=\frac19\left[\begin{pmatrix}1&2&0&2\\2&4&0&4\\0&0&0&0\\2&4&0&4\end{pmatrix}
+\begin{pmatrix}4&-2&4&0\\-2&1&-2&0\\4&-2&4&0\\0&0&0&0\end{pmatrix}\right]
=\frac19\begin{pmatrix}5&0&4&2\\0&5&-2&4\\4&-2&4&0\\2&4&0&4\end{pmatrix}
\]
\item With \(B=[a_1\ a_2]\): \(B'B=\left(\begin{smallmatrix}9&9\\9&18\end{smallmatrix}\right)\), of determinant
\(81\), so \((B'B)^{-1}=\frac1{81}\left(\begin{smallmatrix}18&-9\\-9&9\end{smallmatrix}\right)
=\frac19\left(\begin{smallmatrix}2&-1\\-1&1\end{smallmatrix}\right)\) (Theorem~\ref{r5d-thm-adjugate}), and
\[
B(B'B)^{-1}=\frac19[2a_1-a_2\ \ \ -a_1+a_2]=\frac19\begin{pmatrix}-1&2\\3&-1\\-2&2\\2&0\end{pmatrix},\qquad
B(B'B)^{-1}B'=\frac19\begin{pmatrix}5&0&4&2\\0&5&-2&4\\4&-2&4&0\\2&4&0&4\end{pmatrix}
\]
the same matrix, as Theorem~\ref{r6-thm-projection-matrix}(a),(b) asserts. Also
\(\tr P_W=\frac19(5+5+4+4)=2=\dim W\), and \(P_Wy=(0,5,-2,4)'=\hat w\) for \(y\) of
Example~\ref{r6-exm-projection-onb}.
\item By Theorem~\ref{r6-thm-projection-matrix}(c),
\[
P_{W^\perp}=I_4-P_W=\frac19\begin{pmatrix}4&0&-4&-2\\0&4&2&-4\\-4&2&5&0\\-2&-4&0&5\end{pmatrix}
\]
The columns \(q_3=\frac13b_3\), \(q_4=\frac13v\) of the orthogonal matrix of Example~\ref{r6-exm-extend} form an
orthonormal basis of \(W^\perp\) (Corollary~\ref{r6-cor-onb}(c)), and indeed
\(\frac19(b_3b_3'+vv')\) is this matrix.
\end{itemize}
\end{example}

\begin{theorem}[\citeat{harville}{Theorems~12.3.8, 12.3.9}, \citeat{harville-solutions}{Ch.~10, Exercise~7}]\label{r6-thm-symmetric-idempotent}
A matrix \(A\in\R^{n\times n}\) is the projection matrix for some subspace of \(\R^n\) if and only if
\(A'=A=A^2\). In that case \(A=P_{\mathcal C(A)}\). In particular, two symmetric idempotent matrices with the same
column space are equal.
\end{theorem}
\begin{proof}
\begin{itemize}
\item If \(A=P_W\), then \(A\) is symmetric and idempotent and \(\mathcal C(A)=W\)
(Theorem~\ref{r6-thm-projection-matrix}(c)).
\item Conversely, let \(A'=A=A^2\). By Theorem~\ref{r6-thm-idempotent}(d), \(Av\) is the orthogonal projection of
\(v\) onto \(\mathcal C(A)\) for every \(v\in\R^n\); so \(A=P_{\mathcal C(A)}\)
(Theorem~\ref{r6-thm-projection-matrix}(a)).
\item If \(A\) and \(B\) are symmetric and idempotent with \(\mathcal C(A)=\mathcal C(B)=W\), then
\(A=P_W=B\).\qedhere
\end{itemize}
\end{proof}

\begin{corollary}[\citeat{tu-clm}{Definition~6.1}, \citeat{abadir-magnus}{Exercise~8.71(a)}]\label{r6-cor-centering}
Let \(W_0=\{x\in\R^n:\one_n'x=0\}\).
\begin{enumerate}
\item \(P_{\Span\{\one_n\}}=\bar J_n\) and \(P_{W_0}=C_n\).
\item \(\rank C_n=\tr C_n=n-1\).
\end{enumerate}
\end{corollary}
\begin{proof}
(a) The vector \(q=n^{-1/2}\one_n\) satisfies \(q'q=n^{-1}\one'\one=1\) (Proposition~\ref{r1-prop-centering}), so it
is an orthonormal basis of \(\Span\{\one_n\}\), and \(P_{\Span\{\one_n\}}=qq'=\frac1n\one\one'=\bar J_n\)
(Theorem~\ref{r6-thm-projection-matrix}(b) and Definition~\ref{r1-def-ones}). By
Theorem~\ref{r6-thm-complement}(a), \(W_0=\mathcal N(\one')=\Span\{\one_n\}^\perp\), so
\(P_{W_0}=I_n-\bar J_n=C_n\) (Theorem~\ref{r6-thm-projection-matrix}(c)).

(b) By (a) and Theorem~\ref{r6-thm-projection-matrix}(c), \(\rank C_n=\tr C_n\). Every diagonal entry of
\(\bar J_n\) is \(\frac1n\), so \(\tr C_n=\tr I_n-\tr\bar J_n=n-1\) (Theorem~\ref{r1-thm-trace}(a)).
\end{proof}

\subsection*{Projections onto orthogonal sums and nested subspaces}

\begin{theorem}[\citeat{harville}{Corollary~12.3.7, Theorem~12.3.5}, \citeat{tu-clm}{Exercise~B3(b), Remark~7.7}]\label{r6-thm-nested}
Let \(U\subseteq W\) be subspaces of \(\R^n\), with projection matrices \(P_U\) and \(P_W\).
\begin{enumerate}
\item \(P_WP_U=P_UP_W=P_U\).
\item \(P_W-P_U\) is symmetric and idempotent, of rank \(\dim W-\dim U\), and it is the projection matrix for
\(W\cap U^\perp\). Moreover \((P_W-P_U)P_U=0\) and \((P_W-P_U)(I_n-P_W)=0\), and \(W\) is the orthogonal sum
of \(U\) and \(W\cap U^\perp\).
\item Every \(y\in\R^n\) splits as \(y=P_Uy+(P_W-P_U)y+(I_n-P_W)y\) into three pairwise orthogonal parts, so that
\[
\lVert y\rVert^2=\lVert P_Uy\rVert^2+\lVert(P_W-P_U)y\rVert^2+\lVert(I_n-P_W)y\rVert^2,\qquad
\lVert P_Wy\rVert^2=\lVert P_Uy\rVert^2+\lVert(P_W-P_U)y\rVert^2
\]
\end{enumerate}
\end{theorem}
\begin{proof}
(a) If \(U=\{0\}\), then \(P_U=0\) and (a) is clear. Otherwise let \(Z\) have full column rank with
\(\mathcal C(Z)=U\) (a basis of \(U\) as columns; Theorems~\ref{thm:r4-dimension} and~\ref{thm:r4-rank-rules}(b)),
so that \(P_U=Z(Z'Z)^{-1}Z'\) (Theorem~\ref{r6-thm-projection-matrix}(b)).
\begin{itemize}
\item Each column of \(Z\) lies in \(U\subseteq W\), so it is its own projection onto \(W\)
(Theorem~\ref{r6-thm-projection-onb}(b)): \(P_WZ=Z\) (Theorem~\ref{r1-thm-views}(a)).
\item Hence \(P_WP_U=P_WZ(Z'Z)^{-1}Z'=Z(Z'Z)^{-1}Z'=P_U\). As \(P_W\) and \(P_U\) are symmetric
(Theorem~\ref{r6-thm-projection-matrix}(c)), \(P_UP_W=P_U'P_W'=(P_WP_U)'=P_U'=P_U\).
\end{itemize}
(b)
\begin{itemize}
\item \(P_W\) and \(P_U\) are symmetric and idempotent (Theorem~\ref{r6-thm-projection-matrix}(c)), so \(P_W-P_U\)
is symmetric, and by (a), \((P_W-P_U)^2=P_W^2-P_WP_U-P_UP_W+P_U^2=P_W-P_U-P_U+P_U=P_W-P_U\).
\item Its rank is its trace (Theorem~\ref{r6-thm-idempotent}(c)):
\(\tr P_W-\tr P_U=\dim W-\dim U\) (Theorems~\ref{r1-thm-trace}(a) and~\ref{r6-thm-projection-matrix}(c)).
\item \((P_W-P_U)P_U=P_WP_U-P_U^2=P_U-P_U=0\), and \((P_W-P_U)(I-P_W)=P_W-P_W^2-P_U+P_UP_W=0\).
\item By Theorem~\ref{r6-thm-symmetric-idempotent}, \(P_W-P_U\) is the projection matrix for
\(\mathcal C(P_W-P_U)\). If \(v=(P_W-P_U)y\), then \(P_Wv=(P_W^2-P_WP_U)y=v\), so \(v\in\mathcal C(P_W)=W\); and
\(P_Uv=(P_UP_W-P_U^2)y=0\), so \(v\in\mathcal N(P_U)=\mathcal C(P_U')^\perp=U^\perp\)
(Corollary~\ref{r6-cor-four-subspaces}, with \(P_U'=P_U\), and \(\mathcal C(P_W)=W\), \(\mathcal C(P_U)=U\) by
Theorem~\ref{r6-thm-projection-matrix}(c)). Thus \(\mathcal C(P_W-P_U)\subseteq W\cap U^\perp\).
\item Conversely, if \(v\in W\cap U^\perp\), then \(P_Wv=v\) (Theorem~\ref{r6-thm-projection-onb}(b)) and
\(P_Uv=0\) (as \(v=0+v\) with \(v\perp U\)), so \(v=(P_W-P_U)v\in\mathcal C(P_W-P_U)\).
\item \(W\cap U^\perp\) is a subspace (Theorems~\ref{r6-thm-complement}(a) and~\ref{r3-thm-sum}(a)), orthogonal to
\(U\), as it lies in \(U^\perp\). For \(w\in W\), \(w=P_Uw+(w-P_Uw)\), where \(P_Uw\in U\), and \(w-P_Uw\) lies
in \(W\) (as \(P_Uw\in U\subseteq W\)) and in \(U^\perp\) (Definition~\ref{r6-def-projection}). So
\(W\subseteq U+(W\cap U^\perp)\); conversely \(U+(W\cap U^\perp)\subseteq W\), as \(U\) and \(W\cap U^\perp\) lie
in the subspace \(W\) (Theorem~\ref{r3-thm-sum}(b)). Hence \(W=U+(W\cap U^\perp)\), an orthogonal sum
(Definition~\ref{r6-def-complement}).
\end{itemize}
(c)
\begin{itemize}
\item \(P_U+(P_W-P_U)+(I-P_W)=I\) gives the splitting.
\item For symmetric \(A\), \(B\) with \(AB=0\), \((Ay)'(By)=y'ABy=0\). Here \(P_U(P_W-P_U)=P_UP_W-P_U^2=0\),
\(P_U(I-P_W)=P_U-P_UP_W=0\) and \((P_W-P_U)(I-P_W)=0\), by (a) and (b); so the three parts are pairwise
orthogonal.
\item Proposition~\ref{r6-prop-polarization}(b) gives the equation for \(\lVert y\rVert^2\). As
\(P_Wy=P_Uy+(P_W-P_U)y\), with \(P_Uy\perp(P_W-P_U)y\), Pythagoras (Proposition~\ref{r6-prop-pythagoras}) gives the
equation for \(\lVert P_Wy\rVert^2\).\qedhere
\end{itemize}
\end{proof}

\begin{corollary}[cf.\ \citeat{harville-solutions}{Ch.~18, Exercise~23(a)}, \citeat{abadir-magnus}{Exercise~8.64}]\label{r6-cor-orthogonal-sum}
Let \(W_1\perp W_2\) be orthogonal subspaces of \(\R^n\), and \(P_1=P_{W_1}\), \(P_2=P_{W_2}\). Then
\(W_1\cap W_2=\{0\}\), so the orthogonal sum \(W_1+W_2\) (Definition~\ref{r6-def-complement}) is direct, written
\(W_1\oplus W_2\), and
\[
P_{W_1\oplus W_2}=P_1+P_2,\qquad P_1P_2=P_2P_1=0,\qquad
\lVert P_{W_1\oplus W_2}y\rVert^2=\lVert P_1y\rVert^2+\lVert P_2y\rVert^2\quad(y\in\R^n)
\]
\end{corollary}
\begin{proof}
\begin{itemize}
\item If \(w\in W_1\cap W_2\), then \(w'w=0\), as \(W_1\perp W_2\); so \(w=0\), and the sum is direct
(Definition~\ref{r3-def-sum}). Let \(W=W_1\oplus W_2\), a subspace that contains \(W_1\) and \(W_2\)
(Theorem~\ref{r3-thm-sum}(a),(b)).
\item Apply Theorem~\ref{r6-thm-nested} to \(U=W_1\subseteq W\). Every vector of \(W_2\) lies in \(W\) and is
orthogonal to \(W_1\), so \(W_2\subseteq W\cap W_1^\perp\). By Theorem~\ref{r6-thm-nested}(b), \(P_W-P_1\) is the
projection matrix for \(W\cap W_1^\perp\), of rank \(\dim W-\dim W_1\); so
\(\dim(W\cap W_1^\perp)=\dim W-\dim W_1=\dim W_2\) (Theorem~\ref{r6-thm-projection-matrix}(c) and
Corollary~\ref{cor:r4-dim-sum}(a)). Hence \(W\cap W_1^\perp=W_2\) (Theorem~\ref{thm:r4-counting}(a)), and
\(P_W-P_1=P_2\), that is, \(P_W=P_1+P_2\).
\item By Theorem~\ref{r6-thm-nested}(b), \(P_2P_1=(P_W-P_1)P_1=0\). As \(P_1\) and \(P_2\) are symmetric
(Theorem~\ref{r6-thm-projection-matrix}(c)), \(P_1P_2=P_1'P_2'=(P_2P_1)'=0\) (Theorem~\ref{r1-thm-transpose}).
\item The second equation of Theorem~\ref{r6-thm-nested}(c), with \(P_W-P_1=P_2\), gives
\(\lVert P_Wy\rVert^2=\lVert P_1y\rVert^2+\lVert P_2y\rVert^2\).\qedhere
\end{itemize}
\end{proof}

\begin{example}\label{r6-exm-nested}
In Example~\ref{r6-exm-projection-onb} take \(U=\Span\{a_1\}\subseteq W=\Span\{a_1,a_2\}\). Then
\(P_U=q_1q_1'=\frac19b_1b_1'\) (Theorem~\ref{r6-thm-projection-matrix}(b)), and by
Example~\ref{r6-exm-projection-matrix}, \(P_W-P_U=q_2q_2'=\frac19b_2b_2'\). By Theorem~\ref{r6-thm-nested}(b),
this is the projection matrix for \(W\cap U^\perp\), of dimension \(2-1=1\) (Theorem~\ref{r6-thm-projection-matrix}(c)); as \(0\ne b_2\in W\) and \(b_2\perp U\)
(\(b_2'a_1=0\); Theorem~\ref{r6-thm-complement}(a)), \(W\cap U^\perp=\Span\{b_2\}\) (Theorem~\ref{thm:r4-counting}(a)). For \(y=(-2,7,1,3)'\),
\[
P_Uy=\frac{b_1'y}9\,b_1=\begin{pmatrix}2\\4\\0\\4\end{pmatrix},\qquad
(P_W-P_U)y=\frac{b_2'y}9\,b_2=\begin{pmatrix}-2\\1\\-2\\0\end{pmatrix},\qquad
(I_4-P_W)y=\begin{pmatrix}-2\\2\\3\\-1\end{pmatrix}
\]
and \(\lVert y\rVert^2=63=36+9+18\), as Theorem~\ref{r6-thm-nested}(c) asserts. Here \(U\perp\Span\{b_2\}\) and
\(W=U\oplus\Span\{b_2\}\), so Corollary~\ref{r6-cor-orthogonal-sum} gives \(P_W=P_U+\frac19b_2b_2'\) and
\(\lVert P_Wy\rVert^2=36+9=45\).
\end{example}

\needspace{14\baselineskip}
\subsection*{Exercises}

\begin{exercise}\label{r6-dr-qr}
Let
\[
A=\begin{pmatrix}1&3&1\\1&1&-2\\1&2&-2\\0&1&0\end{pmatrix}
\]
with columns \(a_1,a_2,a_3\).
\begin{enumerate}
\item Apply the Gram--Schmidt process to \(a_1,a_2,a_3\).
\item Find the QR decomposition of \(A\), and check that \(R=Q'A\).
\item The vector \(y=(5,0,1,1)'\) lies in \(\mathcal C(A)\). Find its coordinates \(q_j'y\) in the orthonormal basis
\(q_1,q_2,q_3\) of (b), and check that \(\lVert y\rVert^2=\sum_j(q_j'y)^2\).
\end{enumerate}
\end{exercise}
\begin{solution}
(a) By Theorem~\ref{r6-thm-gram-schmidt}:
\begin{itemize}
\item \(b_1=a_1=(1,1,1,0)'\), \(b_1'b_1=1+1+1+0=3\).
\item \(a_2'b_1=3+1+2+0=6\), so \(t_{12}=2\) and \(b_2=a_2-2b_1=(3,1,2,1)'-(2,2,2,0)'=(1,-1,0,1)'\), with
\(b_2'b_2=1+1+0+1=3\).
\item \(a_3'b_1=1-2-2+0=-3\) and \(a_3'b_2=1+2+0+0=3\), so \(t_{13}=-1\), \(t_{23}=1\) and
\[
b_3=a_3+b_1-b_2=\begin{pmatrix}1\\-2\\-2\\0\end{pmatrix}+\begin{pmatrix}1\\1\\1\\0\end{pmatrix}-\begin{pmatrix}1\\-1\\0\\1\end{pmatrix}
=\begin{pmatrix}1\\0\\-1\\-1\end{pmatrix},\qquad b_3'b_3=1+0+1+1=3
\]
\item Check: \(b_1'b_2=1-1+0+0=0\), \(b_1'b_3=1+0-1+0=0\), \(b_2'b_3=1+0+0-1=0\).
\end{itemize}
(b) All three lengths are \(\sqrt3\), so \(q_j=\frac1{\sqrt3}b_j\), \(r_{jj}=\sqrt3\), \(r_{12}=2\sqrt3\),
\(r_{13}=-\sqrt3\) and \(r_{23}=\sqrt3\) (Theorem~\ref{r6-thm-qr}):
\[
Q=\frac1{\sqrt3}\begin{pmatrix}1&1&1\\1&-1&0\\1&0&-1\\0&1&-1\end{pmatrix},\qquad
R=\sqrt3\begin{pmatrix}1&2&-1\\0&1&1\\0&0&1\end{pmatrix}
\]
Check: the entries of \(\sqrt3\,Q'A\) are the products \(b_i'a_j\), so
\[
Q'A=\frac1{\sqrt3}\begin{pmatrix}3&6&-3\\0&3&3\\0&0&3\end{pmatrix}=R
\]

(c) \(b_1'y=5+0+1+0=6\), \(b_2'y=5-0+0+1=6\) and \(b_3'y=5+0-1-1=3\), so \(q_1'y=2\sqrt3\), \(q_2'y=2\sqrt3\) and
\(q_3'y=\sqrt3\): \(y=2\sqrt3\,q_1+2\sqrt3\,q_2+\sqrt3\,q_3=2b_1+2b_2+b_3\)
(Theorem~\ref{r6-thm-projection-onb}(a),(b)). Check: \(2(1,1,1,0)'+2(1,-1,0,1)'+(1,0,-1,-1)'=(5,0,1,1)'\), and
\(\sum_j(q_j'y)^2=12+12+3=27=25+0+1+1=\lVert y\rVert^2\) (Theorem~\ref{r6-thm-projection-onb}(c)).
\end{solution}

\begin{exercise}[cf.\ \citeat{searle-khuri}{Exercise~5.13(a)}]\label{r6-ex-qr}
Let \(A=[a_1\ a_2\ a_3]\in\R^{4\times3}\) have the columns \(a_1=\one_4\), \(a_2=(2,0,2,0)'\) and
\(a_3=(3,3,1,1)'\).
\begin{enumerate}
\item Find the QR decomposition of \(A\) by the Gram--Schmidt process.
\item Show that \(Q'Q=I_3\) but \(QQ'\ne I_4\), and that \(QQ'\) is the matrix \(P\) of
Example~\ref{r6-exm-projection}.
\end{enumerate}
\end{exercise}
\begin{solution}
(a) By Theorem~\ref{r6-thm-gram-schmidt}:
\begin{itemize}
\item \(b_1=a_1\), \(\lVert b_1\rVert^2=4\).
\item \(a_2'b_1=4\), so \(t_{12}=1\) and \(b_2=a_2-b_1=(1,-1,1,-1)'\), with \(\lVert b_2\rVert^2=4\).
\item \(a_3'b_1=8\) and \(a_3'b_2=3-3+1-1=0\), so \(t_{13}=2\), \(t_{23}=0\) and \(b_3=a_3-2b_1=(1,1,-1,-1)'\),
with \(\lVert b_3\rVert^2=4\).
\item All three lengths are \(2\): \(q_j=\frac12b_j\), \(r_{jj}=2\), \(r_{12}=2t_{12}=2\), \(r_{13}=2t_{13}=4\) and
\(r_{23}=0\):
\[
Q=\frac12\begin{pmatrix}1&1&1\\1&-1&1\\1&1&-1\\1&-1&-1\end{pmatrix},\qquad
R=\begin{pmatrix}2&2&4\\0&2&0\\0&0&2\end{pmatrix}
\]
Check: the third column of \(QR\) is \(Q(4,0,2)'=\frac12(4+2,\ 4+2,\ 4-2,\ 4-2)'=(3,3,1,1)'=a_3\).
\end{itemize}
(b) The columns of \(Q\) are orthonormal, so \(Q'Q=I_3\). By the outer product expansion
(Theorem~\ref{r1-thm-views}(d)),
\[
QQ'=\tfrac14\bigl(b_1b_1'+b_2b_2'+b_3b_3'\bigr)=\frac14\begin{pmatrix}3&1&1&-1\\1&3&-1&1\\1&-1&3&1\\-1&1&1&3\end{pmatrix}\ne I_4
\]
For instance, its \((1,1)\) entry is \(\frac14(1+1+1)=\frac34\). The transposed rows of \(Q\) are not orthonormal. By
Corollary~\ref{r6-cor-qr-projection}, \(QQ'=P_{\mathcal C(A)}\), and \(P_{\mathcal C(A)}=A(A'A)^{-1}A'\) by
Theorem~\ref{r6-thm-projection-matrix}(b): \(QQ'\) is the matrix \(P\) of Example~\ref{r6-exm-projection}.
\end{solution}

\begin{exercise}\label{r6-ex-orthonormal-basis}
Let
\[
Q=\frac12\begin{pmatrix}1&1&1&1\\1&-1&1&-1\\1&1&-1&-1\\1&-1&-1&1\end{pmatrix}
\]
with columns \(q_1,\dots,q_4\), and let \(y=(4,2,0,6)'\).
\begin{enumerate}
\item Show that \(Q\) is a symmetric orthogonal matrix, and find \(Q^{-1}\).
\item Find \(c\in\R^4\) with \(y=c_1q_1+\dots+c_4q_4\), and check that \(\lVert c\rVert=\lVert y\rVert\).
\item Find the vector of \(\Span\{q_1,q_2\}\) closest to \(y\), and its distance from \(y\).
\end{enumerate}
\end{exercise}
\begin{solution}
(a) \(Q'=Q\) by inspection. Each column of \(2Q\) has squared length \(4\), and any two of them are orthogonal
(each pair of columns agrees in two entries and differs in the other two), for
instance \((1,-1,1,-1)'\) and \((1,1,-1,-1)'\) have inner product \(1-1-1+1=0\); so \(Q'Q=I_4\) (Definition~\ref{r6-def-orthonormal}), and
\(Q^{-1}=Q'=Q\) (Proposition~\ref{r6-prop-orthonormal}(b)).

(b) \(y=Qc\) (Theorem~\ref{r1-thm-views}(b)) gives \(c=Q^{-1}y=Q'y\):
\[
c=\tfrac12\bigl(4+2+0+6,\ 4-2+0-6,\ 4+2-0-6,\ 4-2-0+6\bigr)'=(6,-2,0,4)'
\]
(equivalently, \(c_j=q_j'y\) by Proposition~\ref{r6-prop-orthogonal-list}(a)). Then
\(\lVert c\rVert^2=36+4+0+16=56=16+4+0+36=\lVert y\rVert^2\), as Proposition~\ref{r6-prop-orthonormal}(a) says.

(c) \(q_1,q_2\) are orthonormal by (a), hence linearly independent (Proposition~\ref{r6-prop-orthogonal-list}(b)),
so they form an orthonormal basis of \(V=\Span\{q_1,q_2\}\), and
\(P_Vy=q_1(q_1'y)+q_2(q_2'y)=6q_1-2q_2\) (Theorems~\ref{r6-thm-projection-matrix}(b) and~\ref{r1-thm-views}(d)):
\[
P_Vy=3\begin{pmatrix}1\\1\\1\\1\end{pmatrix}-\begin{pmatrix}1\\-1\\1\\-1\end{pmatrix}=\begin{pmatrix}2\\4\\2\\4\end{pmatrix},\qquad
y-P_Vy=\begin{pmatrix}2\\-2\\-2\\2\end{pmatrix}=4q_4
\]
By Theorem~\ref{r6-thm-best-approximation}, \((2,4,2,4)'\) is the closest vector, at distance
\(\lVert y-P_Vy\rVert=\sqrt{16}=4\); by (b), \(y-P_Vy=c_3q_3+c_4q_4\) and \(\lVert y-P_Vy\rVert^2=c_3^2+c_4^2=16\).
\end{solution}

\begin{exercise}\label{r6-ex-qr-solve}
Let \(B=[a_1\ a_2\ a_3\ a_4]\in\R^{4\times4}\), with \(a_1,a_2,a_3\) of Exercise~\ref{r6-ex-qr} and
\(a_4=(4,0,0,0)'\).
\begin{enumerate}
\item Continue the Gram--Schmidt process of Exercise~\ref{r6-ex-qr}(a) to find the QR decomposition of \(B\).
Which matrix of Exercise~\ref{r6-ex-orthonormal-basis} is its \(Q\)?
\item Solve \(Bx=b\) for \(b=(9,3,5,3)'\) by solving \(Rx=Q'b\).
\item Show that \(\lvert\det B\rvert=\det R=16\), without expanding \(\det B\).
\end{enumerate}
\end{exercise}
\begin{solution}
(a) \(a_4'b_1=4\), \(a_4'b_2=4\), \(a_4'b_3=4\), so \(t_{14}=t_{24}=t_{34}=1\) and
\[
b_4=a_4-b_1-b_2-b_3=(4-1-1-1,\ 0-1+1-1,\ 0-1-1+1,\ 0-1+1+1)'=(1,-1,-1,1)'
\]
with \(\lVert b_4\rVert^2=4\). So \(q_4=\frac12b_4\), \(r_{44}=2\) and \(r_{14}=r_{24}=r_{34}=2\cdot1=2\): \(B=QR\) with
\[
Q=\frac12\begin{pmatrix}1&1&1&1\\1&-1&1&-1\\1&1&-1&-1\\1&-1&-1&1\end{pmatrix},\qquad
R=\begin{pmatrix}2&2&4&2\\0&2&0&2\\0&0&2&2\\0&0&0&2\end{pmatrix}
\]
This \(Q\) is the orthogonal matrix of Exercise~\ref{r6-ex-orthonormal-basis}.

(b) \(Bx=b\) is \(QRx=b\); multiplying by \(Q'\), and back by \(Q\), it is equivalent to \(Rx=Q'b\), as \(Q'Q=QQ'=I_4\).
Here \(Q'b=\frac12(20,8,4,4)'=(10,4,2,2)'\), and back substitution gives
\[
2x_4=2,\quad 2x_3+2x_4=2,\quad 2x_2+2x_4=4,\quad 2x_1+2x_2+4x_3+2x_4=10:\qquad x=(3,1,0,1)'
\]
Check: \(3a_1+a_2+a_4=(3+2+4,\ 3,\ 3+2,\ 3)'=(9,3,5,3)'\).

(c) \(\det B=\det Q\det R\) (Theorem~\ref{r5d-thm-product}(b)). As \(Q\) is orthogonal, \(\det Q=\pm1\)
(Theorem~\ref{r6-thm-orthogonal}(c)); and \(R\) is upper triangular, so \(\det R=2\cdot2\cdot2\cdot2=16\)
(Theorem~\ref{r5d-thm-basic}(c)). Hence \(\lvert\det B\rvert=\lvert\det Q\rvert\det R=16\).
\end{solution}

\begin{exercise}\label{r6-dr-complement}
Let \(W=\Span\{v_1,v_2,v_3\}\subseteq\R^4\) with \(v_1=(1,1,-2,3)'\), \(v_2=(2,1,-3,5)'\) and \(v_3=(1,2,-3,4)'\).
\begin{enumerate}
\item Find a basis of \(W^\perp\), and check that \(\dim W+\dim W^\perp=4\).
\item Find an orthonormal basis of \(W^\perp\).
\item Using (b), write \(y=(1,2,0,4)'\) as \(y=w+z\) with \(w\in W\) and \(z\in W^\perp\), and check that
\(w'z=0\).
\end{enumerate}
\end{exercise}
\begin{solution}
(a) \(W^\perp=\mathcal N(V')\) for \(V=[v_1\ v_2\ v_3]\) (Theorem~\ref{r6-thm-complement}(a)). Row-reduce \(V'\):
\[
V'=\begin{pmatrix}1&1&-2&3\\2&1&-3&5\\1&2&-3&4\end{pmatrix}
\xrightarrow[R_3\leftarrow R_3-R_1]{R_2\leftarrow R_2-2R_1}
\begin{pmatrix}1&1&-2&3\\0&-1&1&-1\\0&1&-1&1\end{pmatrix}
\xrightarrow{R_2\leftarrow-R_2}
\begin{pmatrix}1&1&-2&3\\0&1&-1&1\\0&1&-1&1\end{pmatrix}
\]
\[
\xrightarrow[R_3\leftarrow R_3-R_2]{R_1\leftarrow R_1-R_2}
\begin{pmatrix}1&0&-1&2\\0&1&-1&1\\0&0&0&0\end{pmatrix}
\]
So \(\dim W=\rank V'=2\) (Definition~\ref{def:r4-rank}, Theorem~\ref{thm:r4-rank-rules}(a)). With \(x_3=s\) and
\(x_4=t\) free, \(x_1=s-2t\) and \(x_2=s-t\): \(s_1=(1,1,1,0)'\) and \(s_2=(-2,-1,0,1)'\) form a basis of \(W^\perp\)
(Theorem~\ref{thm:r4-four-bases}(c)), and \(2+2=4\) (Theorem~\ref{r6-thm-complement}(c)). Check:
\(s_1'v_1=1+1-2+0=0\), \(s_1'v_2=2+1-3+0=0\), \(s_1'v_3=1+2-3+0=0\), \(s_2'v_1=-2-1+0+3=0\), \(s_2'v_2=-4-1+0+5=0\),
\(s_2'v_3=-2-2+0+4=0\).

(b) Gram--Schmidt: \(s_1's_1=3\) and \(s_2's_1=-2-1+0+0=-3\), so \(u_2=s_2+s_1=(-1,0,1,1)'\), with
\(u_2'u_2=3\) and \(s_1'u_2=-1+0+1+0=0\). An orthonormal basis of \(W^\perp\) is
\(\frac1{\sqrt3}(1,1,1,0)'\), \(\frac1{\sqrt3}(-1,0,1,1)'\).

(c) \(s_1'y=1+2+0+0=3\) and \(u_2'y=-1+0+0+4=3\). By Theorem~\ref{r6-thm-projection-onb}(a) for \(W^\perp\), the
projection of \(y\) onto \(W^\perp\) is \(z=\frac33s_1+\frac33u_2=(0,1,2,1)'\). By
Theorem~\ref{r6-thm-projection-onb}(d) for \(W^\perp\), with \((W^\perp)^\perp=W\)
(Theorem~\ref{r6-thm-perp-perp}(a)), \(w=y-z=(1,1,-2,3)'\) is the projection of \(y\) onto \(W\); indeed \(w=v_1\).
Check: \(w'z=0+1-4+3=0\).
\end{solution}

\begin{exercise}[\citeat{abadir-magnus}{Exercise~3.44(b)}, \citeat{lang}{Ch.~V, \S2, Exercise~0}, \citeat{shakarchi}{}]\label{r6-ex-complements}
\begin{enumerate}
\item Find all vectors of \(\R^3\) that are orthogonal to \(x=(2,2,-1)'\).
\item What is the dimension of the subspace of \(\R^6\) consisting of the vectors orthogonal to both
\((1,1,-2,3,4,5)'\) and \((0,0,1,1,0,7)'\)?
\end{enumerate}
\end{exercise}
\begin{solution}
(a) \(y=(y_1,y_2,y_3)'\) is orthogonal to \(x\) if and only if \(x'y=2y_1+2y_2-y_3=0\), that is,
\(y_3=2y_1+2y_2\). So these vectors are
\[
y=\begin{pmatrix}y_1\\y_2\\2y_1+2y_2\end{pmatrix}=y_1\begin{pmatrix}1\\0\\2\end{pmatrix}+y_2\begin{pmatrix}0\\1\\2\end{pmatrix}\qquad(y_1,y_2\in\R)
\]
Hence \(\Span\{x\}^\perp=\mathcal N(x')=\Span\{(1,0,2)',(0,1,2)'\}\), of dimension \(2=3-1\)
(Theorem~\ref{r6-thm-complement}).

(b) Let \(a_1,a_2\) be the two vectors and \(A=[a_1\ a_2]\in\R^{6\times2}\). The subspace is
\(\Span\{a_1,a_2\}^\perp=\mathcal N(A')\) (Theorem~\ref{r6-thm-complement}(a)). The vectors are independent:
if \(c_1a_1+c_2a_2=0\), the first entry gives \(c_1=0\), and then \(c_2a_2=0\) with \(a_2\ne0\) gives \(c_2=0\).
So \(\dim\Span\{a_1,a_2\}=2\), and the answer is \(6-2=4\) (Theorem~\ref{r6-thm-complement}(c)).
\end{solution}

\begin{exercise}\label{r6-dr-consistency}
Let \(b=(2,3,1,5)'\), \(c=\one_4\) and
\[
A=\begin{pmatrix}1&1&2&1\\2&3&5&1\\1&2&3&0\\3&4&7&2\end{pmatrix}
\]
\begin{enumerate}
\item Find a basis of \(\mathcal N(A')\).
\item Use Corollary~\ref{r6-cor-consistency} to decide whether \(Ax=b\) and \(Ax=c\) have solutions.
\item Solve the consistent system, and check that \(\mathcal N(A)\) is orthogonal to the transposed rows of \(A\).
\end{enumerate}
\end{exercise}
\begin{solution}
(a) Row-reduce \(A'\):
\[
\begin{gathered}
A'=\begin{pmatrix}1&2&1&3\\1&3&2&4\\2&5&3&7\\1&1&0&2\end{pmatrix}
\xrightarrow[R_4\leftarrow R_4-R_1]{R_2\leftarrow R_2-R_1,\ R_3\leftarrow R_3-2R_1}
\begin{pmatrix}1&2&1&3\\0&1&1&1\\0&1&1&1\\0&-1&-1&-1\end{pmatrix}
\\
\xrightarrow[R_4\leftarrow R_4+R_2]{R_1\leftarrow R_1-2R_2,\ R_3\leftarrow R_3-R_2}
\begin{pmatrix}1&0&-1&1\\0&1&1&1\\0&0&0&0\\0&0&0&0\end{pmatrix}
\end{gathered}
\]
With \(w_3=s\) and \(w_4=t\) free, \(A'w=0\) gives \(w_1=s-t\) and \(w_2=-s-t\); so \(z_1=(1,-1,1,0)'\) and
\(z_2=(-1,-1,0,1)'\) form a basis of \(\mathcal N(A')\) (Theorem~\ref{thm:r4-four-bases}(c)), and
\(\rank A=\rank A'=2\) (Definition~\ref{def:r4-rank}, Theorem~\ref{thm:r4-rank-rules}(a)).

(b) \(z_1'b=2-3+1+0=0\) and \(z_2'b=-2-3+0+5=0\), while \(z_1'c=1-1+1+0=1\ne0\). By
Corollary~\ref{r6-cor-consistency}, \(Ax=b\) has a solution and \(Ax=c\) has none.

(c) Row-reduce \([A\mid b]\):
\[
\begin{gathered}
\left(\begin{array}{cccc|c}1&1&2&1&2\\2&3&5&1&3\\1&2&3&0&1\\3&4&7&2&5\end{array}\right)
\xrightarrow[R_4\leftarrow R_4-3R_1]{R_2\leftarrow R_2-2R_1,\ R_3\leftarrow R_3-R_1}
\left(\begin{array}{cccc|c}1&1&2&1&2\\0&1&1&-1&-1\\0&1&1&-1&-1\\0&1&1&-1&-1\end{array}\right)
\\
\xrightarrow[R_4\leftarrow R_4-R_2]{R_1\leftarrow R_1-R_2,\ R_3\leftarrow R_3-R_2}
\left(\begin{array}{cccc|c}1&0&1&2&3\\0&1&1&-1&-1\\0&0&0&0&0\\0&0&0&0&0\end{array}\right)
\end{gathered}
\]
With \(x_3=s\) and \(x_4=t\) free:
\[
x=\begin{pmatrix}3-s-2t\\-1-s+t\\s\\t\end{pmatrix}=\begin{pmatrix}3\\-1\\0\\0\end{pmatrix}+s\begin{pmatrix}-1\\-1\\1\\0\end{pmatrix}
+t\begin{pmatrix}-2\\1\\0\\1\end{pmatrix}\qquad(s,t\in\R)
\]
Check: \(A(3,-1,0,0)'=(3-1,\ 6-3,\ 3-2,\ 9-4)'=(2,3,1,5)'=b\). So \(\mathcal N(A)=\Span\{s_1,s_2\}\) with
\(s_1=(-1,-1,1,0)'\) and \(s_2=(-2,1,0,1)'\) (Theorem~\ref{thm:r4-four-bases}(c)), and both are orthogonal to the four
transposed rows of \(A\): \(-1-1+2+0=0\), \(-2-3+5+0=0\), \(-1-2+3+0=0\), \(-3-4+7+0=0\) for \(s_1\), and \(-2+1+0+1=0\),
\(-4+3+0+1=0\), \(-2+2+0+0=0\), \(-6+4+0+2=0\) for \(s_2\) (Corollary~\ref{r6-cor-four-subspaces}).
\end{solution}

\begin{exercise}\label{r6-dr-extend}
Let \(q_1=\frac13(1,0,2,2)'\), \(q_2=\frac13(0,1,2,-2)'\) and \(q_3=\frac13(2,2,-1,0)'\).
\begin{enumerate}
\item Check that \(q_1,q_2,q_3\) are orthonormal, and find all vectors of \(\R^4\) orthogonal to them.
\item Extend \(q_1,q_2,q_3\) to an orthonormal basis \(q_1,\dots,q_4\) of \(\R^4\), and write down the orthogonal
matrix \(U=[q_1\ \cdots\ q_4]\) and \(U^{-1}\).
\item Find the coordinates \(q_j'y\) of \(y=(3,3,3,3)'\) in this basis, and check that
\(\lVert y\rVert^2=\sum_j(q_j'y)^2\).
\end{enumerate}
\end{exercise}
\begin{solution}
(a) Put \(u_1=(1,0,2,2)'\), \(u_2=(0,1,2,-2)'\), \(u_3=(2,2,-1,0)'\), so \(q_j=\frac13u_j\). Then
\(u_1'u_2=0+0+4-4=0\), \(u_1'u_3=2+0-2+0=0\), \(u_2'u_3=0+2-2+0=0\), and \(u_1'u_1=1+0+4+4=9\),
\(u_2'u_2=0+1+4+4=9\), \(u_3'u_3=4+4+1+0=9\); so the \(q_j\) are orthonormal. The vectors orthogonal to them form
\(\mathcal N(B)\), where \(B\) has the rows \(u_j'\) (Theorem~\ref{r6-thm-complement}(a)):
\[
B=\begin{pmatrix}1&0&2&2\\0&1&2&-2\\2&2&-1&0\end{pmatrix}
\xrightarrow{R_3\leftarrow R_3-2R_1}
\begin{pmatrix}1&0&2&2\\0&1&2&-2\\0&2&-5&-4\end{pmatrix}
\xrightarrow{R_3\leftarrow R_3-2R_2}
\begin{pmatrix}1&0&2&2\\0&1&2&-2\\0&0&-9&0\end{pmatrix}
\]
\[
\xrightarrow{R_3\leftarrow-\frac19R_3}
\begin{pmatrix}1&0&2&2\\0&1&2&-2\\0&0&1&0\end{pmatrix}
\xrightarrow[R_2\leftarrow R_2-2R_3]{R_1\leftarrow R_1-2R_3}
\begin{pmatrix}1&0&0&2\\0&1&0&-2\\0&0&1&0\end{pmatrix}
\]
With \(x_4=t\) free, \(x_1=-2t\), \(x_2=2t\), \(x_3=0\): the vectors are \(t(-2,2,0,1)'\) (\(t\in\R\)).

(b) \(v=(-2,2,0,1)'\) has \(v'v=4+4+0+1=9\), so \(q_4=\frac13(-2,2,0,1)'\) completes an orthonormal basis
(Corollary~\ref{r6-cor-onb}(b)), and by Theorem~\ref{r6-thm-orthogonal}(a),
\[
U=\frac13\begin{pmatrix}1&0&2&-2\\0&1&2&2\\2&2&-1&0\\2&-2&0&1\end{pmatrix},\qquad
U^{-1}=U'=\frac13\begin{pmatrix}1&0&2&2\\0&1&2&-2\\2&2&-1&0\\-2&2&0&1\end{pmatrix}
\]

(c) \(u_1'y=3(1+0+2+2)=15\), \(u_2'y=3(0+1+2-2)=3\), \(u_3'y=3(2+2-1+0)=9\) and \(v'y=3(-2+2+0+1)=3\); so
\(q_1'y=5\), \(q_2'y=1\), \(q_3'y=3\), \(q_4'y=1\), and \(y=5q_1+q_2+3q_3+q_4\)
(Theorem~\ref{r6-thm-projection-onb}(a),(b) with \(W=\R^4\)). Check: \(25+1+9+1=36=9+9+9+9=\lVert y\rVert^2\).
\end{solution}

\begin{exercise}\label{r6-dr-orthogonal}
Let \(Q=\frac13G\) with
\[
G=\begin{pmatrix}1&0&-2&-2\\0&1&-2&2\\2&2&1&0\\2&-2&0&1\end{pmatrix}
\]
\begin{enumerate}
\item Show that \(Q\) is orthogonal, and write down \(Q^{-1}\).
\item Find \(\det Q\) by row-reducing \(G\).
\item For \(x=(3,3,0,0)'\) and \(y=(0,3,3,0)'\), check that \(\lVert Qx\rVert=\lVert x\rVert\) and \((Qx)'(Qy)=x'y\).
\end{enumerate}
\end{exercise}
\begin{solution}
(a) The product of \(G'\) and \(G\) is
\[
\begin{pmatrix}1&0&2&2\\0&1&2&-2\\-2&-2&1&0\\-2&2&0&1\end{pmatrix}\begin{pmatrix}1&0&-2&-2\\0&1&-2&2\\2&2&1&0\\2&-2&0&1\end{pmatrix}
=\begin{pmatrix}9&0&0&0\\0&9&0&0\\0&0&9&0\\0&0&0&9\end{pmatrix}
\]
so \(Q'Q=\frac19G'G=I_4\): \(Q\) is orthogonal, and \(Q^{-1}=Q'=\frac13G'\) (Theorem~\ref{r6-thm-orthogonal}(a)).

(b)
\[
G\xrightarrow[R_4\leftarrow R_4-2R_1]{R_3\leftarrow R_3-2R_1}
\begin{pmatrix}1&0&-2&-2\\0&1&-2&2\\0&2&5&4\\0&-2&4&5\end{pmatrix}
\xrightarrow[R_4\leftarrow R_4+2R_2]{R_3\leftarrow R_3-2R_2}
\begin{pmatrix}1&0&-2&-2\\0&1&-2&2\\0&0&9&0\\0&0&0&9\end{pmatrix}
\]
Adding multiples of rows does not change the determinant (Definition~\ref{r5d-def-det}(b)), and the last matrix
is upper triangular; so \(\det G=1\cdot1\cdot9\cdot9=81\) (Theorem~\ref{r5d-thm-basic}(c)), and
\(\det Q=\bigl(\frac13\bigr)^4\cdot81=1\) (Theorem~\ref{r5d-thm-product}(c)), as
Theorem~\ref{r6-thm-orthogonal}(c) allows.

(c) \(Qx=\frac13(3+0,\ 0+3,\ 6+6,\ 6-6)'=(1,1,4,0)'\), so \(\lVert Qx\rVert^2=1+1+16+0=18=9+9+0+0=\lVert x\rVert^2\). Next
\(Qy=\frac13(0+0-6+0,\ 0+3-6+0,\ 0+6+3+0,\ 0-6+0+0)'=(-2,-1,3,-2)'\), and
\((Qx)'(Qy)=-2-1+12+0=9=0+9+0+0=x'y\).
\end{solution}

\begin{exercise}\label{r6-dr-helmert}
Let \(K_4\) be the Helmert matrix of order \(4\) printed in Example~\ref{r6-exm-helmert}, with rows
\(p_1',\dots,p_4'\), and let \(y=(7,3,2,0)'\).
\begin{enumerate}
\item Check that \(p_1,\dots,p_4\) are orthonormal.
\item Compute \(K_4y\), and check that \(\sum_i(y_i-\bar y)^2\) is the sum of the squares of the three contrasts.
\item Check that \(p_2p_2'+p_3p_3'+p_4p_4'=C_4\).
\end{enumerate}
\end{exercise}
\begin{solution}
(a) The rows are \(p_1'=\frac12(1,1,1,1)\), \(p_2'=\frac1{\sqrt2}(1,-1,0,0)\), \(p_3'=\frac1{\sqrt6}(1,1,-2,0)\) and
\(p_4'=\frac1{\sqrt{12}}(1,1,1,-3)\). Then
\[
\begin{gathered}
p_1'p_2=\frac{1-1+0+0}{2\sqrt2}=0,\quad p_1'p_3=\frac{1+1-2+0}{2\sqrt6}=0,\quad p_1'p_4=\frac{1+1+1-3}{2\sqrt{12}}=0\\
p_2'p_3=\frac{1-1+0+0}{\sqrt{12}}=0,\quad p_2'p_4=\frac{1-1+0+0}{\sqrt{24}}=0,\quad p_3'p_4=\frac{1+1-2+0}{\sqrt{72}}=0
\end{gathered}
\]
and \(p_1'p_1=\frac44=1\), \(p_2'p_2=\frac22=1\), \(p_3'p_3=\frac{1+1+4}6=1\), \(p_4'p_4=\frac{1+1+1+9}{12}=1\).

(b) \(K_4y=\bigl(\frac{7+3+2+0}2,\ \frac{7-3}{\sqrt2},\ \frac{7+3-4}{\sqrt6},\ \frac{7+3+2-0}{\sqrt{12}}\bigr)'
=(6,\ 2\sqrt2,\ \sqrt6,\ 2\sqrt3)'\). Here \(\bar y=3\), and \(p_1'y=6=\sqrt4\,\bar y\). The contrasts give
\((2\sqrt2)^2+(\sqrt6)^2+(2\sqrt3)^2=8+6+12=26\), and \(\sum_i(y_i-3)^2=16+0+1+9=26\)
(Example~\ref{r6-exm-helmert}(b)).

(c)
\[
\begin{gathered}
p_2p_2'+p_3p_3'+p_4p_4'=\frac12\begin{pmatrix}1&-1&0&0\\-1&1&0&0\\0&0&0&0\\0&0&0&0\end{pmatrix}
+\frac16\begin{pmatrix}1&1&-2&0\\1&1&-2&0\\-2&-2&4&0\\0&0&0&0\end{pmatrix}
+\frac1{12}\begin{pmatrix}1&1&1&-3\\1&1&1&-3\\1&1&1&-3\\-3&-3&-3&9\end{pmatrix}\\
=\frac14\begin{pmatrix}3&-1&-1&-1\\-1&3&-1&-1\\-1&-1&3&-1\\-1&-1&-1&3\end{pmatrix}=I_4-\bar J_4=C_4
\end{gathered}
\]
(the \((1,1)\) entry, for instance, is \(\frac12+\frac16+\frac1{12}=\frac34\)).
\end{solution}

\begin{exercise}\label{r6-dr-reflection}
Let \(h=\one_4\) and \(F_h=I_4-\frac2{h'h}hh'\) (Example~\ref{r6-exm-reflection}).
\begin{enumerate}
\item Compute \(F_h\), and check that \(F_h^2=I_4\).
\item Check that \(F_hh=-h\) and that \(F_hx=x\) for \(x=(1,-1,0,0)'\), and find \(\det F_h\).
\item For \(y=(3,1,0,4)'\), compute \(F_hy\) and check that \(\lVert F_hy\rVert=\lVert y\rVert\).
\end{enumerate}
\end{exercise}
\begin{solution}
(a) \(h'h=4\), so \(F_h=I_4-\frac12J_4=\frac12N\) with
\[
N=\begin{pmatrix}1&-1&-1&-1\\-1&1&-1&-1\\-1&-1&1&-1\\-1&-1&-1&1\end{pmatrix},\qquad
N^2=\begin{pmatrix}4&0&0&0\\0&4&0&0\\0&0&4&0\\0&0&0&4\end{pmatrix}
\]
(each diagonal entry of \(N^2\) is \(1+1+1+1\), each other entry is, for instance, \(-1-1+1+1\)); so
\(F_h^2=\frac14N^2=I_4\).

(b) \(F_hh=\frac12(1-1-1-1,\ -1+1-1-1,\ -1-1+1-1,\ -1-1-1+1)'=(-1,-1,-1,-1)'=-h\), and
\(F_hx=\frac12(1+1,\ -1-1,\ -1+1,\ -1+1)'=(1,-1,0,0)'=x\) (here \(h'x=0\)). By Theorem~\ref{r5d-thm-det-lemma}
with \(u=-\frac12h\) and \(v=h\), \(\det F_h=1+v'u=1-\frac12\cdot4=-1\).

(c) \(h'y=3+1+0+4=8\), so \(F_hy=y-\frac{2\cdot8}4\,h=y-4h=(-1,-3,-4,0)'\), and
\(\lVert F_hy\rVert^2=1+9+16+0=26=9+1+0+16=\lVert y\rVert^2\). (Here \(F_hy=y-2\bar y\one\): the mean \(\bar y=2\) is
reflected to \(-2\), and the deviations \(y-\bar y\one\) are kept.)
\end{solution}

\begin{exercise}\label{r6-dr-projection-basis}
Let \(W=\Span\{w_1,w_2\}\subseteq\R^4\) with \(w_1=(1,0,1,1)'\) and \(w_2=(2,1,1,0)'\), and let \(y=(1,0,1,4)'\).
\begin{enumerate}
\item Find an orthogonal basis of \(W\) by the Gram--Schmidt process.
\item Find the orthogonal projection \(\hat w\) of \(y\) onto \(W\) and the projection of \(y\) onto \(W^\perp\), and
check that the latter is orthogonal to \(w_1\) and \(w_2\).
\item Find the distance from \(y\) to the nearest vector of \(W\), and check Bessel's inequality.
\end{enumerate}
\end{exercise}
\begin{solution}
(a) \(u_1=w_1\), \(u_1'u_1=1+0+1+1=3\); \(w_2'u_1=2+0+1+0=3\), so \(u_2=w_2-u_1=(1,1,0,-1)'\), with
\(u_2'u_2=3\) and \(u_1'u_2=1+0+0-1=0\).

(b) \(u_1'y=1+0+1+4=6\) and \(u_2'y=1+0+0-4=-3\). By Theorem~\ref{r6-thm-projection-onb}(a),
\[
\hat w=\frac63\,u_1-\frac33\,u_2=\begin{pmatrix}2\\0\\2\\2\end{pmatrix}-\begin{pmatrix}1\\1\\0\\-1\end{pmatrix}
=\begin{pmatrix}1\\-1\\2\\3\end{pmatrix}
\]
and the projection onto \(W^\perp\) is \(y-\hat w=(0,1,-1,1)'\) (Theorem~\ref{r6-thm-projection-onb}(d)). Check:
\(w_1'(y-\hat w)=0+0-1+1=0\) and \(w_2'(y-\hat w)=0+1-1+0=0\).

(c) The distance is \(\lVert y-\hat w\rVert=\sqrt{0+1+1+1}=\sqrt3\) (Theorem~\ref{r6-thm-best-approximation}).
Bessel: \(\lVert\hat w\rVert^2=6^2/3+3^2/3=12+3=15\leqslant\lVert y\rVert^2=1+0+1+16=18\), and \(18=15+3\).
\end{solution}

\begin{exercise}\label{r6-ex-two-vector-span}
Let \(a_1=(1,2,2,0)'\), \(a_2=(3,2,1,2)'\), \(W=\Span\{a_1,a_2\}\subseteq\R^4\) and \(y=(4,2,5,3)'\).
\begin{enumerate}
\item Find an orthonormal basis \(q_1,q_2\) of \(W\) by the Gram--Schmidt process.
\item Compute \(P_W\), and check that \(\tr P_W=2\).
\item Find the vector of \(W\) closest to \(y\), and write it as a combination of \(a_1\) and \(a_2\). Find the
distance from \(y\) to \(W\).
\end{enumerate}
\end{exercise}
\begin{solution}
(a) \(b_1=a_1\), with \(\lVert b_1\rVert^2=1+4+4+0=9\); \(a_2'b_1=3+4+2+0=9\), so \(t_{12}=1\) and
\(b_2=a_2-b_1=(2,0,-1,2)'\), with \(\lVert b_2\rVert^2=4+0+1+4=9\) (Theorem~\ref{r6-thm-gram-schmidt}). Hence
\(q_1=\frac13(1,2,2,0)'\) and \(q_2=\frac13(2,0,-1,2)'\).

(b) With \(Q=[q_1\ q_2]\), Theorems~\ref{r6-thm-projection-matrix}(b) and~\ref{r1-thm-views}(d) give
\[
P_W=QQ'=\frac19\bigl(b_1b_1'+b_2b_2'\bigr)
=\frac19\left[\begin{pmatrix}1&2&2&0\\2&4&4&0\\2&4&4&0\\0&0&0&0\end{pmatrix}
+\begin{pmatrix}4&0&-2&4\\0&0&0&0\\-2&0&1&-2\\4&0&-2&4\end{pmatrix}\right]
=\frac19\begin{pmatrix}5&2&0&4\\2&4&4&0\\0&4&5&-2\\4&0&-2&4\end{pmatrix}
\]
and \(\tr P_W=\frac19(5+4+5+4)=2=\dim W\) (Theorem~\ref{r6-thm-projection-matrix}(c)).

(c) The closest vector is the projection (Theorem~\ref{r6-thm-best-approximation}):
\[
P_Wy=\frac19\begin{pmatrix}20+4+0+12\\8+8+20+0\\0+8+25-6\\16+0-10+12\end{pmatrix}=\begin{pmatrix}4\\4\\3\\2\end{pmatrix}
=a_1+a_2
\]
The distance is \(\lVert y-P_Wy\rVert=\lVert(0,-2,2,1)'\rVert=\sqrt{0+4+4+1}=3\). Check:
\(a_1'(y-P_Wy)=0-4+4+0=0\) and \(a_2'(y-P_Wy)=0-4+2+2=0\).
\end{solution}

\begin{exercise}\label{r6-ex-qr-projection}
Let \(A\) be the matrix of Example~\ref{r6-exm-qr}, \(W=\mathcal C(A)\) and \(y=(0,6,3,3)'\).
\begin{enumerate}
\item Show that \(h=(0,2,1,-2)'\) spans \(W^\perp\).
\item Compute \(P_W=QQ'\), with \(Q\) of Example~\ref{r6-exm-qr}, and check that \(P_W=I_4-\frac19hh'\).
\item Find the vector of \(W\) closest to \(y\) and its distance from \(y\), and write the closest vector as \(Ax\) by
solving \(Rx=Q'y\).
\end{enumerate}
\end{exercise}
\begin{solution}
(a) By Example~\ref{r6-exm-qr}, \(W=\Span\{b_1,b_2,b_3\}\) with \(b_1=(1,2,0,2)'\), \(b_2=(2,-1,2,0)'\),
\(b_3=(-2,0,2,1)'\) (Theorem~\ref{r6-thm-gram-schmidt}). Now \(b_1'h=0+4+0-4=0\), \(b_2'h=0-2+2+0=0\) and
\(b_3'h=0+0+2-2=0\), so \(h\in W^\perp\) (Theorem~\ref{r6-thm-complement}(a)). As \(\dim W=3\) (the
nonzero, pairwise orthogonal \(b_j\) are linearly independent, Proposition~\ref{r6-prop-orthogonal-list}(b)),
\(\dim W^\perp=1\) (Theorem~\ref{r6-thm-complement}(c)), and the nonzero vector \(h\) spans \(W^\perp\)
(Theorem~\ref{thm:r4-counting}(d)). (This \(h\) is the vector \(v\) found in Example~\ref{r6-exm-extend}.)

(b) By Corollary~\ref{r6-cor-qr-projection} and Theorem~\ref{r1-thm-views}(d),
\[
P_W=QQ'=\tfrac19\bigl(b_1b_1'+b_2b_2'+b_3b_3'\bigr)=\frac19\begin{pmatrix}9&0&0&0\\0&5&-2&4\\0&-2&8&2\\0&4&2&5\end{pmatrix}
\]
For instance, the \((2,3)\) entry is \(\frac19(0-2+0)\) and the \((3,3)\) entry is \(\frac19(0+4+4)\). Since
\(h'h=9\), \(\frac19hh'=h(h'h)^{-1}h'\) is the projection matrix for \(\Span\{h\}=W^\perp\)
(Theorem~\ref{r6-thm-projection-matrix}(b) with the matrix \(h\)), and \(P_W=I_4-\frac19hh'\)
(Theorem~\ref{r6-thm-projection-matrix}(c)); entry by entry, \(I_4-\frac19hh'\) is the matrix
above, as \(hh'\) has rows \(0\), \((0,4,2,-4)\), \((0,2,1,-2)\), \((0,-4,-2,4)\).

(c) \(h'y=0+12+3-6=9\), so \((I_4-P_W)y=\frac{h'y}{h'h}h=h\) and \(P_Wy=y-h=(0,4,2,5)'\). This is the closest vector
(Theorem~\ref{r6-thm-best-approximation}), at distance \(\lVert h\rVert=3\). Next
\(Q'y=\frac13(b_1'y,b_2'y,b_3'y)'=\frac13(18,0,9)'=(6,0,3)'\), and with \(R\) of Example~\ref{r6-exm-qr},
\[
3x_3=3,\quad 3x_2+3x_3=0,\quad 3x_1+3x_2+6x_3=6:\qquad x=(1,-1,1)'
\]
Then \(Ax=QRx=QQ'y=P_Wy\) (Corollary~\ref{r6-cor-qr-projection}); directly,
\(Ax=a_1-a_2+a_3=(1-3+2,\ 2-1+3,\ 0-2+4,\ 2-2+5)'=(0,4,2,5)'\).
\end{solution}

\begin{exercise}\label{r6-ex-projection-drill}
Let \(W=\mathcal C(A)\) and \(y=(2,1,1,4,8)'\), where
\[
A=\begin{pmatrix}1&1\\1&0\\1&0\\2&0\\3&1\end{pmatrix}
\]
\begin{enumerate}
\item Find a basis of \(W^\perp\).
\item Find the QR decomposition of \(A\), and compute \(P_W=QQ'\). Check that \(P_WA=A\) and \(\tr P_W=2\).
\item Write \(y=P_Wy+(I_5-P_W)y\), the decomposition along \(\R^5=W\oplus W^\perp\), check that the two parts are
orthogonal, and find the distance from \(y\) to \(W\).
\end{enumerate}
\end{exercise}
\begin{solution}
(a) \(W^\perp=\mathcal N(A')\) (Theorem~\ref{r6-thm-complement}(a)). Row-reduce \(A'\):
\[
\begin{pmatrix}1&1&1&2&3\\1&0&0&0&1\end{pmatrix}
\xrightarrow{R_1\leftrightarrow R_2}
\begin{pmatrix}1&0&0&0&1\\1&1&1&2&3\end{pmatrix}
\xrightarrow{R_2\leftarrow R_2-R_1}
\begin{pmatrix}1&0&0&0&1\\0&1&1&2&2\end{pmatrix}
\]
With \(x_3,x_4,x_5\) free, \(x_1=-x_5\) and \(x_2=-x_3-2x_4-2x_5\); so
\(v_1=(0,-1,1,0,0)'\), \(v_2=(0,-2,0,1,0)'\), \(v_3=(-1,-2,0,0,1)'\) form a basis of \(W^\perp\)
(Theorem~\ref{thm:r4-four-bases}(c)), of dimension \(5-2=3\).

(b) By Theorem~\ref{r6-thm-gram-schmidt}, \(b_1=a_1=(1,1,1,2,3)'\), \(\lVert b_1\rVert^2=1+1+1+4+9=16\);
\(a_2'b_1=1+3=4\), so \(t_{12}=\frac14\) and
\[
b_2=a_2-\tfrac14b_1=\tfrac14(3,-1,-1,-2,1)',\qquad\lVert b_2\rVert^2=\tfrac1{16}(9+1+1+4+1)=1
\]
So \(q_1=\frac14b_1\), \(q_2=b_2\), \(r_{11}=4\), \(r_{12}=t_{12}\lVert b_1\rVert=1\), \(r_{22}=1\)
(Theorem~\ref{r6-thm-qr}):
\[
Q=\frac14\begin{pmatrix}1&3\\1&-1\\1&-1\\2&-2\\3&1\end{pmatrix},\qquad R=\begin{pmatrix}4&1\\0&1\end{pmatrix},\qquad
P_W=QQ'=\frac18\begin{pmatrix}5&-1&-1&-2&3\\-1&1&1&2&1\\-1&1&1&2&1\\-2&2&2&4&2\\3&1&1&2&5\end{pmatrix}
\]
(Corollary~\ref{r6-cor-qr-projection}); for instance, the \((1,5)\) entry is \(\frac1{16}(1\cdot3+3\cdot1)=\frac38\). Row by
row, \(P_Wa_1=\frac18(5-1-1-4+9,\ -1+1+1+4+3,\ -1+1+1+4+3,\ -2+2+2+8+6,\ 3+1+1+4+15)'=(1,1,1,2,3)'\) and
\(P_Wa_2=\frac18(5+3,\ -1+1,\ -1+1,\ -2+2,\ 3+5)'=(1,0,0,0,1)'\), so \(P_WA=A\); and
\(\tr P_W=\frac18(5+1+1+4+5)=2=\dim W\) (Theorem~\ref{r6-thm-projection-matrix}(c)).

(c) \(Q'y=(\frac14(2+1+1+8+24),\ \frac14(6-1-1-8+8))'=(9,1)'\), so
\(P_Wy=Q(Q'y)=9q_1+q_2=\frac14(9+3,\ 9-1,\ 9-1,\ 18-2,\ 27+1)'=(3,2,2,4,7)'\), and
\[
\begin{pmatrix}2\\1\\1\\4\\8\end{pmatrix}=\begin{pmatrix}3\\2\\2\\4\\7\end{pmatrix}+\begin{pmatrix}-1\\-1\\-1\\0\\1\end{pmatrix}
\]
The second part is \(v_3-v_1\in W^\perp\), and \((P_Wy)'(I_5-P_W)y=-3-2-2+0+7=0\). The distance is
\(\lVert(-1,-1,-1,0,1)'\rVert=2\) (Theorem~\ref{r6-thm-best-approximation}).
\end{solution}

\begin{exercise}\label{r6-ex-idempotent-drill}
Let \(v=\one_4\) and
\[
E=\begin{pmatrix}1&1&-1&1\\1&0&1&0\\0&-1&2&-1\\-1&-1&1&-1\end{pmatrix}
\]
\begin{enumerate}
\item Show that \(E\) is idempotent but not symmetric.
\item Find \(\rank E\), \(\tr E\), \(\rank(I_4-E)\) and \(\tr(I_4-E)\).
\item Decompose \(v\) along \(\R^4=\mathcal N(E)\oplus\mathcal C(E)\). Are the two parts orthogonal?
\end{enumerate}
\end{exercise}
\begin{solution}
Let \(\rho_1',\dots,\rho_4'\) be the rows of \(E\) (\(\rho_i\in\R^4\)); then \(\rho_3=\rho_2-\rho_1\) and \(\rho_4=-\rho_1\). Write \(E=(e_{ij})\).

(a) By Theorem~\ref{r1-thm-views}(c), \(\row_i(E^2)=\rho_i'E=e_{i1}\rho_1'+e_{i2}\rho_2'+e_{i3}\rho_3'+e_{i4}\rho_4'\). So
\(\rho_1'E=\rho_1'+\rho_2'-\rho_3'+\rho_4'=\rho_1'+\rho_2'-(\rho_2'-\rho_1')-\rho_1'=\rho_1'\) and
\(\rho_2'E=\rho_1'+\rho_3'=\rho_2'\); then \(\rho_3'E=\rho_2'E-\rho_1'E=\rho_3'\) and \(\rho_4'E=-\rho_1'E=\rho_4'\). Hence
\(E^2=E\). It is not symmetric: \(e_{13}=-1\ne0=e_{31}\).

(b) Row-reduce \(E\):
\[
E\xrightarrow[R_4\leftarrow R_4+R_1]{R_2\leftarrow R_2-R_1}
\begin{pmatrix}1&1&-1&1\\0&-1&2&-1\\0&-1&2&-1\\0&0&0&0\end{pmatrix}
\xrightarrow{R_3\leftarrow R_3-R_2}
\begin{pmatrix}1&1&-1&1\\0&-1&2&-1\\0&0&0&0\\0&0&0&0\end{pmatrix}
\]
Scaling row \(2\) by \(-1\) gives an echelon form with two nonzero rows, so \(\rank E=2=1+0+2-1=\tr E\)
(Remark~\ref{rem:r4-echelon} and Definition~\ref{def:r4-rank}). Next
\[
I_4-E=\begin{pmatrix}0&-1&1&-1\\-1&1&-1&0\\0&1&-1&1\\1&1&-1&2\end{pmatrix}
\]
whose rows \(\sigma_1',\dots,\sigma_4'\) (\(\sigma_i\in\R^4\)) satisfy \(\sigma_3=-\sigma_1\) and \(\sigma_4=-2\sigma_1-\sigma_2\).
Row-reduce it:
\[
\begin{gathered}
I_4-E\xrightarrow[R_4\leftarrow R_4+2R_1]{R_3\leftarrow R_3+R_1}
\begin{pmatrix}0&-1&1&-1\\-1&1&-1&0\\0&0&0&0\\1&-1&1&0\end{pmatrix}
\xrightarrow{R_4\leftarrow R_4+R_2}
\begin{pmatrix}0&-1&1&-1\\-1&1&-1&0\\0&0&0&0\\0&0&0&0\end{pmatrix}
\\
\xrightarrow{R_1\leftrightarrow R_2}
\begin{pmatrix}-1&1&-1&0\\0&-1&1&-1\\0&0&0&0\\0&0&0&0\end{pmatrix}
\end{gathered}
\]
Scaling rows \(1\) and \(2\) by \(-1\) gives an echelon form with two nonzero rows, so
\(\rank(I_4-E)=2=0+1-1+2=\tr(I_4-E)\) (Remark~\ref{rem:r4-echelon} and Definition~\ref{def:r4-rank}), as Theorem~\ref{r6-thm-idempotent}(b),(c) say.

(c) \(Ev=(2,2,0,-2)'\in\mathcal C(E)\) (the row sums of \(E\)) and \(v-Ev=(-1,-1,1,3)'\), which lies in \(\mathcal N(E)\):
\(E(-1,-1,1,3)'=(-1-1-1+3,\ -1+0+1+0,\ 0+1+2-3,\ 1+1+1-3)'=0\). So \(v=(-1,-1,1,3)'+(2,2,0,-2)'\)
(Theorem~\ref{r6-thm-idempotent}(a)). The parts are not orthogonal: \((v-Ev)'(Ev)=-2-2+0-6=-10\). (By
Theorem~\ref{r6-thm-idempotent}(d), they would be if \(E\) were symmetric.)
\end{solution}

\begin{exercise}\label{r6-dr-projection-matrix}
Let \(W=\mathcal C(A)\) and \(y=(2,1,2,0)'\), where
\[
A=[a_1\ a_2]=\begin{pmatrix}1&3\\1&1\\0&1\\1&2\end{pmatrix}
\]
\begin{enumerate}
\item Compute \(P_W=A(A'A)^{-1}A'\).
\item Find an orthonormal basis \(q_1,q_2\) of \(W\) by the Gram--Schmidt process, and check that
\(q_1q_1'+q_2q_2'=P_W\).
\item Find \(P_{W^\perp}\), check that \(\tr P_W=\tr P_{W^\perp}=2\), and split \(y\) along \(\R^4=W\oplus W^\perp\).
\end{enumerate}
\end{exercise}
\begin{solution}
(a) \(A'A=\left(\begin{smallmatrix}3&6\\6&15\end{smallmatrix}\right)\), with determinant \(45-36=9\), so
\((A'A)^{-1}=\frac19\left(\begin{smallmatrix}15&-6\\-6&3\end{smallmatrix}\right)
=\frac13\left(\begin{smallmatrix}5&-2\\-2&1\end{smallmatrix}\right)\) (Theorem~\ref{r5d-thm-adjugate}). Then
\[
(A'A)^{-1}A'=\frac13\begin{pmatrix}5&-2\\-2&1\end{pmatrix}\begin{pmatrix}1&1&0&1\\3&1&1&2\end{pmatrix}
=\frac13\begin{pmatrix}-1&3&-2&1\\1&-1&1&0\end{pmatrix}
\]
\[
P_W=A(A'A)^{-1}A'=\begin{pmatrix}1&3\\1&1\\0&1\\1&2\end{pmatrix}\frac13\begin{pmatrix}-1&3&-2&1\\1&-1&1&0\end{pmatrix}
=\frac13\begin{pmatrix}2&0&1&1\\0&2&-1&1\\1&-1&1&0\\1&1&0&1\end{pmatrix}
\]

(b) \(u_1=a_1=(1,1,0,1)'\), \(u_1'u_1=3\); \(a_2'u_1=3+1+0+2=6\), so \(u_2=a_2-2u_1=(1,-1,1,0)'\), with
\(u_2'u_2=3\) and \(u_1'u_2=0\). So \(q_j=\frac1{\sqrt3}u_j\), and
\[
q_1q_1'+q_2q_2'=\frac13\left[\begin{pmatrix}1&1&0&1\\1&1&0&1\\0&0&0&0\\1&1&0&1\end{pmatrix}
+\begin{pmatrix}1&-1&1&0\\-1&1&-1&0\\1&-1&1&0\\0&0&0&0\end{pmatrix}\right]
=\frac13\begin{pmatrix}2&0&1&1\\0&2&-1&1\\1&-1&1&0\\1&1&0&1\end{pmatrix}=P_W
\]
as Theorem~\ref{r6-thm-projection-matrix}(b) asserts.

(c) By Theorem~\ref{r6-thm-projection-matrix}(c),
\[
P_{W^\perp}=I_4-P_W=\frac13\begin{pmatrix}1&0&-1&-1\\0&1&1&-1\\-1&1&2&0\\-1&-1&0&2\end{pmatrix}
\]
and \(\tr P_W=\frac13(2+2+1+1)=2\), \(\tr P_{W^\perp}=\frac13(1+1+2+2)=2\). Next
\(P_Wy=\frac13(4+0+2+0,\ 0+2-2+0,\ 2-1+2+0,\ 2+1+0+0)'=(2,0,1,1)'\) and
\(P_{W^\perp}y=y-P_Wy=(0,1,1,-1)'\); check: \((2,0,1,1)(0,1,1,-1)'=0+0+1-1=0\).
\end{solution}

\begin{exercise}\label{r6-dr-nested}
Let \(a_1=(0,1,1,1)'\), \(a_2=(1,2,0,1)'\), \(U=\Span\{a_1\}\), \(W=\Span\{a_1,a_2\}\), and \(y=(2,2,1,3)'\).
\begin{enumerate}
\item Compute \(P_U\) and \(P_W\), and check that \(P_Wa_1=a_1\) and hence \(P_WP_U=P_U\).
\item Compute \(P_W-P_U\), check that it is idempotent of rank \(1\), and find a vector spanning
\(W\cap U^\perp\).
\item Split \(y=P_Uy+(P_W-P_U)y+(I_4-P_W)y\), and check the equation of Theorem~\ref{r6-thm-nested}(c).
\end{enumerate}
\end{exercise}
\begin{solution}
(a) \(a_1'a_1=3\), so \(P_U=\frac13a_1a_1'\) (Theorem~\ref{r6-thm-projection-matrix}(b)). For \(A=[a_1\ a_2]\),
\(A'A=\left(\begin{smallmatrix}3&3\\3&6\end{smallmatrix}\right)\) (\(a_1'a_2=0+2+0+1=3\), \(a_2'a_2=1+4+0+1=6\)),
with determinant \(9\), so \((A'A)^{-1}=\frac13\left(\begin{smallmatrix}2&-1\\-1&1\end{smallmatrix}\right)\)
(Theorem~\ref{r5d-thm-adjugate}) and \(A(A'A)^{-1}=\frac13[2a_1-a_2\ \ \ -a_1+a_2]\); by
Theorem~\ref{r6-thm-projection-matrix}(b), \(P_W=A(A'A)^{-1}A'\):
\[
P_U=\frac13\begin{pmatrix}0&0&0&0\\0&1&1&1\\0&1&1&1\\0&1&1&1\end{pmatrix},\qquad
P_W=\frac13\begin{pmatrix}-1&1\\0&1\\2&-1\\1&0\end{pmatrix}\begin{pmatrix}0&1&1&1\\1&2&0&1\end{pmatrix}
=\frac13\begin{pmatrix}1&1&-1&0\\1&2&0&1\\-1&0&2&1\\0&1&1&1\end{pmatrix}
\]
Then \(P_Wa_1=\frac13(0+1-1+0,\ 0+2+0+1,\ 0+0+2+1,\ 0+1+1+1)'=(0,1,1,1)'=a_1\), so
\(P_WP_U=\frac13(P_Wa_1)a_1'=P_U\).

(b)
\[
P_W-P_U=\frac13\begin{pmatrix}1&1&-1&0\\1&1&-1&0\\-1&-1&1&0\\0&0&0&0\end{pmatrix}=\frac13uu',\qquad u=(1,1,-1,0)'
\]
As \(u'u=3\), \(\bigl(\frac13uu'\bigr)^2=\frac19u(u'u)u'=\frac13uu'\): it is idempotent, of rank
\(\tr(P_W-P_U)=\frac13(1+1+1+0)=1\) (Theorem~\ref{r6-thm-idempotent}(c)). By Theorem~\ref{r6-thm-nested}(b), it
is the projection matrix for \(W\cap U^\perp=\mathcal C(P_W-P_U)=\Span\{u\}\); indeed \(u=a_2-a_1\in W\) and
\(u'a_1=0+1-1+0=0\).

(c) \(P_Uy=\frac{a_1'y}3\,a_1=\frac63\,a_1=(0,2,2,2)'\), \((P_W-P_U)y=\frac{u'y}3\,u=\frac33\,u=(1,1,-1,0)'\), and
\((I_4-P_W)y=y-(0,2,2,2)'-(1,1,-1,0)'=(1,-1,0,1)'\). Then
\(\lVert y\rVert^2=4+4+1+9=18=12+3+3\).
\end{solution}

\needspace{14\baselineskip}
\subsection*{Problems}

\begin{problem}[\citeat{harville-solutions}{Ch.~17, Exercise~3}]\label{r6-ex-complement-sum}
Let \(U\), \(W\) and \(Z\) be subspaces of \(\R^n\), and \(y\in\R^n\).
\begin{enumerate}
\item Show (1) that if \(y\perp W\) and \(y\perp Z\), then \(y\perp W+Z\); and (2) that if \(U\perp W\) and
\(U\perp Z\), then \(U\perp W+Z\).
\item Show (1) that \((U+W)^\perp=U^\perp\cap W^\perp\); and (2) that \((U\cap W)^\perp=U^\perp+W^\perp\).
\end{enumerate}
\end{problem}
\begin{solution}
(a)
\begin{itemize}
\item (1) Let \(t\in W+Z\), say \(t=w+z\) with \(w\in W\) and \(z\in Z\) (Definition~\ref{r3-def-sum}). Then
\(y't=y'w+y'z=0+0=0\). So \(y\perp W+Z\).
\item (2) Every \(u\in U\) satisfies \(u\perp W\) and \(u\perp Z\), so \(u\perp W+Z\) by (1). Hence \(U\perp W+Z\).
\end{itemize}
(b)
\begin{itemize}
\item (1) \(U\subseteq U+W\) and \(W\subseteq U+W\) (Theorem~\ref{r3-thm-sum}(b)). So, using (a)(1) for the
converse,
\(y\in(U+W)^\perp\iff y\perp U+W\iff y\perp U\text{ and }y\perp W\iff y\in U^\perp\cap W^\perp\).
\item (2) \(U^\perp+W^\perp\) is a subspace of \(\R^n\) (Theorems~\ref{r6-thm-complement}(a)
and~\ref{r3-thm-sum}(a)). By Theorem~\ref{r6-thm-perp-perp}(a), used three times, and (1) for the subspaces
\(U^\perp\) and \(W^\perp\),
\[
U^\perp+W^\perp=\bigl((U^\perp+W^\perp)^\perp\bigr)^\perp=\bigl((U^\perp)^\perp\cap(W^\perp)^\perp\bigr)^\perp
=(U\cap W)^\perp
\]
\end{itemize}
\end{solution}

\begin{problem}[\citeat{harville-solutions}{Ch.~12, Exercise~8}]\label{r6-ex-gs-projection}
Let \(a_1,\dots,a_k\in\R^n\) (\(k\geqslant2\)) be linearly independent, and let \(b_1,\dots,b_k\) be obtained from
them by the Gram--Schmidt process (Theorem~\ref{r6-thm-gram-schmidt}). Show that for \(j=2,\dots,k\) the vector
\(b_j\) is the orthogonal projection of \(a_j\) onto some subspace \(U_j\) of \(\R^n\), and describe \(U_j\).
\end{problem}
\begin{solution}
Let \(q_i=b_i/\lVert b_i\rVert\) (\(i=1,\dots,k\)) and \(W_j=\Span\{q_1,\dots,q_j\}\).
\begin{itemize}
\item For \(j=2,\dots,k\), since \(t_{ij}b_i=\frac{a_j'b_i}{\lVert b_i\rVert^2}\,b_i=(a_j'q_i)q_i\),
\[
b_j=a_j-\sum_{i=1}^{j-1}t_{ij}b_i=a_j-\sum_{i=1}^{j-1}(a_j'q_i)q_i
\]
\item \(q_1,\dots,q_{j-1}\) are orthonormal (Theorem~\ref{r6-thm-gram-schmidt}), so they form an orthonormal basis
of \(W_{j-1}\), and \(\sum_{i<j}(a_j'q_i)q_i\) is the orthogonal projection of \(a_j\) onto \(W_{j-1}\)
(Theorem~\ref{r6-thm-projection-onb}(a)).
\item By Theorem~\ref{r6-thm-projection-onb}(d), \(b_j\) is the orthogonal projection of \(a_j\) onto
\(W_{j-1}^\perp\).
\item Each \(q_i\) is a multiple of \(b_i\) and vice versa, so \(W_{j-1}=\Span\{b_1,\dots,b_{j-1}\}\)
(Theorem~\ref{r3-thm-span}), and this is \(\Span\{a_1,\dots,a_{j-1}\}\) (Theorem~\ref{r6-thm-gram-schmidt}).
\end{itemize}
Hence \(b_j\) is the orthogonal projection of \(a_j\) onto \(U_j=\Span\{a_1,\dots,a_{j-1}\}^\perp\), the orthogonal
complement of the span of the preceding vectors.
\end{solution}

\begin{problem}[\citeat{harville}{\S6.4d}]\label{r6-pr-qr-unique}
Let \(A=[a_1\ \cdots\ a_p]\in\R^{n\times p}\) have full column rank, and let \(A=QR\) be a QR decomposition as in
Theorem~\ref{r6-thm-qr}: \(Q\in\R^{n\times p}\) has orthonormal columns and \(R\in\R^{p\times p}\) is upper triangular
with positive diagonal entries. Show that this decomposition is unique: if \(Q^*\in\R^{n\times p}\) has orthonormal
columns, \(R^*\in\R^{p\times p}\) is upper triangular with positive diagonal entries, and \(A=Q^*R^*\), then \(Q^*=Q\)
and \(R^*=R\).
\end{problem}
\begin{solution}
The columns \(a_1,\dots,a_p\) of \(A\) are linearly independent (Theorem~\ref{thm:r4-rank-rules}(b)). Let
\(b_1,\dots,b_p\) and \(t_{ij}\) (\(i<j\)) be the vectors and coefficients obtained from them by the Gram--Schmidt
process (Theorem~\ref{r6-thm-gram-schmidt}), so that \(b_1,\dots,b_p\) are nonzero and pairwise orthogonal.
\begin{itemize}
\item Step 1. Let \(A=BX\), where the columns \(c_1,\dots,c_p\) of \(B\in\R^{n\times p}\) are nonzero and pairwise
orthogonal, and \(X=(x_{ij})\in\R^{p\times p}\) is unit upper triangular (upper triangular with diagonal entries
\(1\)). We show that \(c_j=b_j\) for every \(j\) and \(x_{ij}=t_{ij}\) for \(i<j\). By Theorem~\ref{r1-thm-views}(a),(b),
\(a_j=\col_j(BX)=B\col_j(X)=c_j+\sum_{i<j}x_{ij}c_i\), that is,
\[
c_j=a_j-\sum_{i<j}x_{ij}c_i\qquad(j=1,\dots,p)
\]
We use induction on \(j\). For \(j=1\), \(c_1=a_1=b_1\). Let \(j\geqslant2\), and let \(c_i=b_i\) for \(i<j\). For
\(l<j\), as the \(c_i\) are pairwise orthogonal,
\(0=c_l'c_j=c_l'a_j-\sum_{i<j}x_{ij}c_l'c_i=b_l'a_j-x_{lj}\lVert b_l\rVert^2\); so
\(x_{lj}=a_j'b_l/\lVert b_l\rVert^2=t_{lj}\), and \(c_j=a_j-\sum_{i<j}t_{ij}b_i=b_j\).
\item Step 2. Let \(A=Q^*R^*\) be as in the statement, with columns \(q_1^*,\dots,q_p^*\) of \(Q^*\) and entries
\(r_{ij}^*\) of \(R^*\). Put \(E^*=\diag(r_{11}^*,\dots,r_{pp}^*)\) and \(D^*=\diag(1/r_{11}^*,\dots,1/r_{pp}^*)\), so that
\(E^*D^*=I_p\), and let \(B^*=Q^*E^*\) and \(X^*=D^*R^*\). The columns of \(B^*\) are \(r_{jj}^*q_j^*\)
(Theorem~\ref{r1-thm-views}(a),(b)): nonzero, as \(r_{jj}^*>0\) and \(\lVert q_j^*\rVert=1\), and pairwise orthogonal. The
rows of \(X^*\) are \(\row_i(R^*)/r_{ii}^*\) (Theorem~\ref{r1-thm-views}(c)), so \(X^*\) is unit upper triangular, with
entries \(r_{ij}^*/r_{ii}^*\). And \(B^*X^*=Q^*E^*D^*R^*=Q^*R^*=A\).
\item By Step 1, \(r_{jj}^*q_j^*=b_j\) for every \(j\), and \(r_{ij}^*/r_{ii}^*=t_{ij}\) for \(i<j\). Taking lengths,
\(r_{jj}^*=\lVert b_j\rVert\), as \(r_{jj}^*>0\) and \(\lVert q_j^*\rVert=1\). Hence
\[
q_j^*=\frac{b_j}{\lVert b_j\rVert},\qquad r_{jj}^*=\lVert b_j\rVert,\qquad r_{ij}^*=t_{ij}\lVert b_i\rVert\quad(i<j),
\qquad r_{ij}^*=0\quad(i>j)
\]
\item Step 2 applies to the decomposition \(A=QR\) as well, with the same right sides. So \(Q\) and \(Q^*\) have the
same columns, and \(R\) and \(R^*\) the same entries: \(Q^*=Q\) and \(R^*=R\).
\end{itemize}
\end{solution}

\begin{problem}[\citeat{abadir-magnus}{Exercise~8.64}]\label{r6-ex-idempotent-sum}
Let \(A,B\in\R^{n\times n}\) be idempotent.
\begin{enumerate}
\item Show that \(A+B\) is idempotent if and only if \(AB=BA=0\).
\item If \(A+B\) is idempotent, show that \(\rank(A+B)=\rank A+\rank B\).
\end{enumerate}
\end{problem}
\begin{solution}
(a)
\begin{itemize}
\item If \(AB=BA=0\), then \((A+B)^2=A^2+AB+BA+B^2=A^2+B^2=A+B\).
\item Conversely, if \((A+B)^2=A+B\), expanding as above gives \(AB+BA=0\). Multiplying on the left by \(A\) gives
\(AB+ABA=0\); multiplying on the right by \(A\) gives \(ABA+BA=0\).
\item Subtracting, \(AB-BA=0\), so \(AB=BA\); then \(AB+BA=2AB=0\), and \(AB=BA=0\).
\end{itemize}
(b) By Theorem~\ref{r6-thm-idempotent}(c), applied to \(A+B\), \(A\) and \(B\), and
Theorem~\ref{r1-thm-trace}(a), \(\rank(A+B)=\tr(A+B)=\tr A+\tr B=\rank A+\rank B\).
\end{solution}

\begin{problem}[\citeat{abadir-magnus}{Exercise~8.67}]\label{r6-ex-complementary-projections}
Let \(A=[A_1\ A_2]\in\R^{n\times n}\), with \(A_1\in\R^{n\times n_1}\) and \(A_2\in\R^{n\times n_2}\)
(\(n_1,n_2\geqslant1\), \(n_1+n_2=n\)). Suppose that \(\rank A_1+\rank A_2=n\) and \(A_1'A_2=0\). Show that \(A_1\)
and \(A_2\) have full column rank and that
\[
A_1(A_1'A_1)^{-1}A_1'+A_2(A_2'A_2)^{-1}A_2'=I_n
\]
\end{problem}
\begin{solution}
\begin{itemize}
\item \(\rank A_1\leqslant n_1\) and \(\rank A_2\leqslant n_2\) (Theorem~\ref{thm:r4-rank-rules}(a)), and the
ranks add up to \(n=n_1+n_2\); so \(\rank A_1=n_1\) and \(\rank A_2=n_2\).
\item The columns of \(A\) are independent: let \(A_1c+A_2d=0\) with \(c\in\R^{n_1}\), \(d\in\R^{n_2}\), and put
\(v=A_1c=-A_2d\). Then \(v'v=-c'A_1'A_2d=0\), so \(v=0\); hence \(A_1c=0\) and \(A_2d=0\), and \(c=0\), \(d=0\)
because \(\mathcal N(A_1)=\{0\}\) and \(\mathcal N(A_2)=\{0\}\) (Theorem~\ref{thm:r4-rank-rules}(b)). So \(A\) is
invertible (Theorem~\ref{thm:r4-rank-rules}(b),(d)).
\item \(A_2'A_1=(A_1'A_2)'=0\), so by Theorem~\ref{r5b-thm-block}(c),
\[
A'A=\begin{pmatrix}A_1'A_1&A_1'A_2\\A_2'A_1&A_2'A_2\end{pmatrix}=\begin{pmatrix}A_1'A_1&0\\0&A_2'A_2\end{pmatrix}
\]
Here \(A_1'A_1\) and \(A_2'A_2\) are invertible (Theorems~\ref{thm:r4-rank-ata}(b) and~\ref{thm:r4-rank-rules}(d)),
so \((A'A)^{-1}\) is block diagonal with the diagonal blocks \((A_1'A_1)^{-1}\) and \((A_2'A_2)^{-1}\)
(Proposition~\ref{r5b-prop-block-tri}(b)).
\item As \(A\) is invertible, \((A'A)^{-1}=A^{-1}(A')^{-1}\) (Theorem~\ref{r2-thm-inverse-props}(b),(c)), so
\(A(A'A)^{-1}A'=AA^{-1}(A')^{-1}A'=I_n\). By Theorem~\ref{r5b-thm-block}(c), therefore,
\[
\begin{aligned}
I_n&=[A_1\ A_2]\begin{pmatrix}(A_1'A_1)^{-1}&0\\0&(A_2'A_2)^{-1}\end{pmatrix}\begin{pmatrix}A_1'\\A_2'\end{pmatrix}\\
&=A_1(A_1'A_1)^{-1}A_1'+A_2(A_2'A_2)^{-1}A_2'
\end{aligned}
\]
\end{itemize}
\end{solution}

\begin{problem}[\citeat{searle-khuri}{Exercise~5.5}, \citeat{harville-solutions}{Ch.~17, Exercise~19(a)}]\label{r6-pr-commuting-projections}
Let \(U,W\) be subspaces of \(\R^n\). Show that \(P_UP_W\) is symmetric if and only if \(P_UP_W=P_WP_U\), and that
then \(P_UP_W=P_{U\cap W}\).
\end{problem}
\begin{solution}
\((P_UP_W)'=P_W'P_U'=P_WP_U\) (Theorems~\ref{r1-thm-transpose} and~\ref{r6-thm-projection-matrix}(c)), so
\(P_UP_W\) is symmetric if and only if \(P_UP_W=P_WP_U\). Let them commute, and put \(P=P_UP_W\).
\begin{itemize}
\item \(P^2=P_UP_WP_UP_W=P_UP_UP_WP_W=P_UP_W=P\). So \(P\) is symmetric idempotent, and \(P=P_{\mathcal C(P)}\)
(Theorem~\ref{r6-thm-symmetric-idempotent}).
\item \(\mathcal C(P)=U\cap W\): for \(y\in\R^n\), \(Py=P_U(P_Wy)\in U\) and \(Py=P_W(P_Uy)\in W\); conversely, for
\(x\in U\cap W\), \(Px=P_U(P_Wx)=P_Ux=x\) (Theorem~\ref{r6-thm-projection-matrix}(c)), so \(x\in\mathcal C(P)\). Hence
\(P=P_{U\cap W}\).
\end{itemize}
\end{solution}

\begin{problem}[\citeat{prasolov}{Problem~7.4}]\label{r6-ex-pr-negative-products}
Let \(v_1,\dots,v_m\in\R^n\) satisfy \(v_i'v_j<0\) for all \(i\ne j\).
\begin{enumerate}
\item Show that if \(m=n+1\), then any \(n\) of the vectors \(v_1,\dots,v_m\) form a basis of \(\R^n\).
\item Show that \(m\leqslant n+1\).
\end{enumerate}
\end{problem}
\begin{solution}
Both parts use the following fact.

(\(*\)) If \(c_1v_1+\dots+c_mv_m=0\), then the nonzero numbers among \(c_1,\dots,c_m\) all have the same sign.
\begin{itemize}
\item Suppose not: the sets \(S_+=\{i:c_i>0\}\) and \(S_-=\{j:c_j<0\}\) are both nonempty. They are disjoint, and the
terms with \(c_i=0\) vanish, so
\[
z=\sum_{i\in S_+}c_iv_i=\sum_{j\in S_-}(-c_j)v_j
\]
\item Multiplying the transpose of the first expression by the second (Theorems~\ref{r1-thm-transpose}
and~\ref{r1-thm-rules}),
\[
z'z=\sum_{i\in S_+}\sum_{j\in S_-}c_i(-c_j)\,v_i'v_j
\]
Every term is \(<0\), as \(c_i>0\), \(-c_j>0\) and \(v_i'v_j<0\) (\(i\ne j\), since \(S_+\) and \(S_-\) are disjoint), and
there is at least one term. So \(z'z<0\), while \(z'z=\lVert z\rVert^2\geqslant0\) (Definition~\ref{r6-def-inner})
\(\contradiction\)
\end{itemize}
(a) Let \(m=n+1\), fix \(r\in\{1,\dots,n+1\}\), and let \(\sum_{i\ne r}c_iv_i=0\), the sum over the \(n\) indices
\(i\ne r\).
\begin{itemize}
\item With \(c_r=0\), this is a relation \(c_1v_1+\dots+c_{n+1}v_{n+1}=0\). Suppose that some \(c_i\ne0\). By (\(*\)),
the nonzero \(c_i\) have one sign; replacing every \(c_i\) by \(-c_i\) if necessary, every \(c_i\geqslant0\) and some
\(c_l>0\), where \(l\ne r\), as \(c_r=0\).
\item Multiplying the relation on the left by \(v_r'\) gives \(\sum_{i\ne r}c_i\,v_r'v_i=0\). But every term
\(c_i\,v_r'v_i\) is \(\leqslant0\), as \(v_r'v_i<0\), and the term with \(i=l\) is \(<0\); so the sum is \(<0\)
\(\contradiction\)
\item Hence every \(c_i=0\): the \(n\) vectors \(v_i\) (\(i\ne r\)) are linearly independent. As \(\dim\R^n=n\)
(Theorem~\ref{thm:r4-dimension}), they form a basis of \(\R^n\) (Theorem~\ref{thm:r4-counting}(d)).
\end{itemize}
(b) Suppose that \(m\geqslant n+2\).
\begin{itemize}
\item Consider the homogeneous linear equations
\[
c_1v_1+c_2v_2+\dots+c_{n+2}v_{n+2}=0,\qquad c_1+c_2+\dots+c_{n+2}=0
\]
in the unknowns \(c_1,\dots,c_{n+2}\). The first is a system of \(n\) equations, one for each entry, whose coefficient
matrix is \([v_1\ \cdots\ v_{n+2}]\) (Theorem~\ref{r1-thm-views}(b)); so there are \(n+1\) equations in \(n+2\) unknowns,
and by Theorem~\ref{r2-thm-homogeneous}(a) they have a solution with some \(c_i\ne0\).
\item Put \(c_i=0\) for \(n+2<i\leqslant m\). Then \(c_1v_1+\dots+c_mv_m=0\), so by (\(*\)) the nonzero \(c_i\) all
have one sign. As some \(c_i\ne0\), the sum \(c_1+\dots+c_{n+2}\), which is the sum of the nonzero \(c_i\), is then
\(\ne0\) \(\contradiction\)
\end{itemize}
Hence \(m\leqslant n+1\).
\end{solution}

\begin{problem}[\citeat{prasolov}{Problem~9.1}]\label{r6-ex-pr-perp-pair}
Let \(U\) and \(W\) be subspaces of \(\R^n\) with \(\dim U=\dim W=m\). Suppose that some \(u\in U\), \(u\ne0\), is
orthogonal to \(W\). Show that some \(w\in W\), \(w\ne0\), is orthogonal to \(U\).
\end{problem}
\begin{solution}
As \(u\ne0\), \(m\geqslant1\). Put \(q_1=u/\lVert u\rVert\), a unit vector of \(U\).
\begin{itemize}
\item If \(m=1\), then \(q_1\) alone is an orthonormal basis of \(U\) (Proposition~\ref{r3-prop-dependence}(a),
Theorem~\ref{thm:r4-counting}(d)). If \(m\geqslant2\), there are \(q_2,\dots,q_m\in U\) such that \(q_1,\dots,q_m\)
form an orthonormal basis of \(U\) (Corollary~\ref{r6-cor-onb}(b)).
\item Let \(Z=\Span\{q_2,\dots,q_m\}\), so \(Z=\{0\}\) if \(m=1\) (Definition~\ref{r3-def-span}). The vectors
\(q_2,\dots,q_m\) are orthonormal, hence linearly independent (Proposition~\ref{r6-prop-orthogonal-list}(b)); so they
form a basis of \(Z\), and \(\dim Z=m-1\). By Theorem~\ref{r6-thm-complement}(c), \(\dim Z^\perp=n-m+1\).
\item Hence \(\dim Z^\perp+\dim W=(n-m+1)+m=n+1>n\), and \(Z^\perp\cap W\ne\{0\}\) (Corollary~\ref{cor:r4-dim-sum}(b)).
Take \(w\in Z^\perp\cap W\) with \(w\ne0\).
\item Then \(w'q_j=0\) for \(j=2,\dots,m\), as \(w\in Z^\perp\) and \(q_j\in Z\); and \(w'q_1=u'w/\lVert u\rVert=0\), as
\(u\perp W\) and \(w\in W\) (Definition~\ref{r6-def-complement}). Since \(q_1,\dots,q_m\) span \(U\), \(w\perp U\)
(Theorem~\ref{r6-thm-complement}(a)).
\end{itemize}
\end{solution}

\begin{problem}[\citeat{prasolov}{Problem~9.2}]\label{r6-ex-pr-gram}
Let \(x_1,\dots,x_n\) and \(y_1,\dots,y_n\) be two bases of \(\R^n\) such that \(x_i'x_j=y_i'y_j\) for all \(i,j\).
Show that there is an orthogonal matrix \(U\in\R^{n\times n}\) with \(Ux_i=y_i\) for \(i=1,\dots,n\).
\end{problem}
\begin{solution}
Let \(A=[x_1\ \cdots\ x_n]\) and \(B=[y_1\ \cdots\ y_n]\).
\begin{itemize}
\item The \((i,j)\) entries of \(A'A\) and \(B'B\) are \(x_i'x_j\) and \(y_i'y_j\) (Theorem~\ref{r1-thm-views}(d)), so
\(A'A=B'B\).
\item The columns of \(A\) are linearly independent, so \(A\) has full column rank and is invertible
(Theorem~\ref{thm:r4-rank-rules}(b),(d)). By Theorem~\ref{r6-thm-qr}, \(A=Q_1R\), where \(Q_1'Q_1=I_n\) and \(R\) is upper
triangular with positive diagonal entries. So \(Q_1\) is orthogonal (Definition~\ref{r6-def-orthonormal}), and
\(R=Q_1'A=Q_1^{-1}A\) is invertible (Theorems~\ref{r6-thm-orthogonal}(a) and~\ref{r2-thm-inverse-props}(a),(b)), with
\(AR^{-1}=Q_1\).
\item Apply the same combinations to \(y_1,\dots,y_n\): let \(Q_2=BR^{-1}\). By Theorem~\ref{r1-thm-transpose} and
\(B'B=A'A\),
\[
Q_2'Q_2=(R^{-1})'B'BR^{-1}=(R^{-1})'A'AR^{-1}=(AR^{-1})'(AR^{-1})=Q_1'Q_1=I_n
\]
so \(Q_2\) is orthogonal too.
\item Let \(U=Q_2Q_1'\). It is orthogonal (Theorem~\ref{r6-thm-orthogonal}(b)), and \(UA=Q_2(Q_1'Q_1)R=Q_2R=B\). Column
\(i\) of this equation reads \(Ux_i=y_i\) (Theorem~\ref{r1-thm-views}(a)).
\end{itemize}
In fact \(U=BR^{-1}Q_1'=B(Q_1R)^{-1}=BA^{-1}\) (Theorem~\ref{r2-thm-inverse-props}(b)).
\end{solution}

\clearpage
\def\unitnumber{7}
\setcounter{theorem}{0}\setcounter{problem}{0}
\section*{Chapter 7. Eigenvalues and diagonalization}

\subsection*{Eigenvalues and eigenvectors}

\begin{definition}[\citeat{hoffman-kunze}{\S6.2}, \citeat{hefferon}{Five.II.3, Definition~3.5, Lemma~3.13}, \citeat{lang}{Ch.~VIII, \S1}, \citeat{harville}{\S21.1}]\label{r7-def-eigen}
Let \(A\in\R^{n\times n}\). A real number \(\lambda\) is an \textbf{eigenvalue} of \(A\) if \(Ax=\lambda x\) for some
\(x\in\R^n\) with \(x\ne0\); such a vector \(x\) is an \textbf{eigenvector} of \(A\) \textbf{belonging to}
\(\lambda\). Since \(Ax=\lambda x\) if and only if \((A-\lambda I)x=0\), the set
\[
\mathcal N(A-\lambda I)=\{x\in\R^n:Ax=\lambda x\}
\]
is called the \textbf{eigenspace} of \(A\) for \(\lambda\).
\end{definition}
Eigenvalues need not be real numbers: once the characteristic polynomial of \(A\) is defined below, the eigenvalues of
\(A\) are defined as the \(n\) complex roots of that polynomial, listed with multiplicity, and the eigenvalues of this
definition are exactly the real ones among them.

The eigenspace \(\mathcal N(A-\lambda I)\) is a subspace of \(\R^n\) (Example~\ref{r3-exm-subspaces}(e)), and its
nonzero vectors are exactly the eigenvectors belonging to \(\lambda\). So a nonzero multiple of an eigenvector, and a
nonzero sum of eigenvectors belonging to the same \(\lambda\), are again eigenvectors belonging to \(\lambda\). An
eigenvector belongs to only one eigenvalue: if \(Ax=\lambda_1x\) and \(Ax=\lambda_2x\), then
\((\lambda_1-\lambda_2)x=0\) with \(x\ne0\), so \(\lambda_1=\lambda_2\) (Proposition~\ref{r3-prop-dependence}(a)).

\begin{example}\label{r7-exm-two}
Let
\[
A=\begin{pmatrix}1&2&-2&-2\\-1&0&-2&2\\0&-1&-1&1\\-1&2&-2&0\end{pmatrix},\qquad
v_1=\begin{pmatrix}1\\1\\0\\1\end{pmatrix},\qquad v_2=\begin{pmatrix}1\\1\\1\\1\end{pmatrix},\qquad
v_3=\begin{pmatrix}0\\1\\0\\1\end{pmatrix},\qquad v_4=\begin{pmatrix}0\\1\\1\\0\end{pmatrix}
\]
Multiplying out,
\[
Av_1=\begin{pmatrix}1\\1\\0\\1\end{pmatrix},\qquad Av_2=\begin{pmatrix}-1\\-1\\-1\\-1\end{pmatrix},\qquad
Av_3=\begin{pmatrix}0\\2\\0\\2\end{pmatrix},\qquad Av_4=\begin{pmatrix}0\\-2\\-2\\0\end{pmatrix}
\]
So \(Av_1=v_1\), \(Av_2=-v_2\), \(Av_3=2v_3\) and \(Av_4=-2v_4\): the numbers \(1\), \(-1\), \(2\) and \(-2\) are
eigenvalues of \(A\), and \(v_1\), \(v_2\), \(v_3\), \(v_4\) are eigenvectors belonging to them
(Definition~\ref{r7-def-eigen}). The vector \(e_1\) is not an eigenvector: \(Ae_1=(1,-1,0,-1)'\) is not a multiple of
\(e_1\).
\end{example}

\begin{theorem}[\citeat{lang}{Ch.~VIII, Theorems~2.1, 2.2}, \citeat{hoffman-kunze}{\S6.2, Theorem~1}, \citeat{harville}{\S21.1}]\label{r7-thm-eig-singular}
Let \(A\in\R^{n\times n}\) and \(\lambda\in\R\). The following are equivalent:
\begin{enumerate}[label=(\roman*)]
\item \(\lambda\) is an eigenvalue of \(A\);
\item \(A-\lambda I\) is not invertible;
\item \(\det(\lambda I-A)=0\);
\item \(\rank(A-\lambda I)<n\).
\end{enumerate}
\end{theorem}
\begin{proof}
\begin{itemize}
\item (i)\(\Leftrightarrow\)(ii): \(\lambda\) is an eigenvalue if and only if \((A-\lambda I)x=0\) has a solution
\(x\ne0\), that is, if and only if \(A-\lambda I\) is not invertible (Theorem~\ref{r2-thm-invertible},
(i)\(\Leftrightarrow\)(iii)).
\item (ii)\(\Leftrightarrow\)(iii): \(\det(\lambda I-A)=\det\bigl((-1)(A-\lambda I)\bigr)=(-1)^n\det(A-\lambda I)\)
(Theorem~\ref{r5d-thm-product}(c)), and \(A-\lambda I\) is invertible if and only if \(\det(A-\lambda I)\ne0\)
(Theorem~\ref{r5d-thm-product}(a)).
\item (ii)\(\Leftrightarrow\)(iv): Theorem~\ref{thm:r4-rank-rules}(d).\qedhere
\end{itemize}
\end{proof}
So the eigenvalues of \(A\) are the real solutions \(\lambda\) of \(\det(\lambda I-A)=0\), and the eigenvectors
belonging to \(\lambda\) are the nonzero solutions of \((A-\lambda I)x=0\), found by row reduction.

\subsection*{Polynomials}

\textbf{Complex numbers.} A \textbf{complex number} is \(z=a+bi\) with \(a,b\in\R\), where \(i^2=-1\); \(\mathbb C\) is
the set of complex numbers, and a real number \(a\) is the complex number \(a+0i\). Sums and products are
\((a+bi)+(c+di)=(a+c)+(b+d)i\) and \((a+bi)(c+di)=(ac-bd)+(ad+bc)i\); they obey the rules of arithmetic of real
numbers (commutative, associative and distributive laws), as one checks by multiplying out both sides~\citeat{hefferon}{Five.I.1}, \citeat{hoffman-kunze}{\S1.1}. The
\textbf{conjugate} of \(z=a+bi\) is \(\bar z=a-bi\); multiplying out likewise gives: \(\bar z=z\) if and only if \(z\) is
real, \(\overline{z+w}=\bar z+\bar w\), \(\overline{zw}=\bar z\,\bar w\), and \(z\bar z=a^2+b^2>0\) for \(z\ne0\). Hence every \(z\ne0\) has the inverse \(z^{-1}=\bar z/(a^2+b^2)\), and \(zw=0\)
only if \(z=0\) or \(w=0\) (if \(z\ne0\), then \(w=z^{-1}zw=0\)).

\begin{definition}[\citeat{hoffman-kunze}{\S\S4.1, 4.2}, \citeat{hefferon}{Five.I.1}, \citeat{lang}{Ch.~IX, \S\S1, 2}]\label{r7-def-polynomial}
A \textbf{polynomial} (in the variable \(t\)) is an expression \(f=c_0+c_1t+\dots+c_mt^m\) with \textbf{coefficients}
\(c_0,\dots,c_m\in\mathbb C\); we put \(c_k=0\) for \(k>m\), and \(f\) equals \(g=d_0+d_1t+\cdots\) if \(c_k=d_k\) for every
\(k\geqslant0\). We write \(f(t)\) for \(f\) when the variable is to be shown.
\begin{itemize}
\item \(f\) is \textbf{real} if every \(c_k\in\R\); the \textbf{zero polynomial} \(0\) has all coefficients \(0\). For
\(f\ne0\), the \textbf{degree} \(\deg f\) is the largest \(k\) with \(c_k\ne0\), and \(c_{\deg f}\) is the
\textbf{leading coefficient}; \(f\) is \textbf{monic} if its leading coefficient is \(1\). \(c_0\) is the
\textbf{constant term}, and a polynomial with \(c_k=0\) for all \(k\geqslant1\) is written as the number \(c_0\).
\item The sum \(f+g\), the multiple \(af\) (\(a\in\mathbb C\)) and the product \(fg\) have the coefficients \(c_k+d_k\),
\(ac_k\) and \(\sum_{i+j=k}c_id_j=c_0d_k+c_1d_{k-1}+\dots+c_kd_0\); thus \((ct^i)(dt^j)=cd\,t^{i+j}\). These operations
obey the rules of arithmetic of numbers (\(fg=gf\), \((fg)h=f(gh)\), \(f(g+h)=fg+fh\), and so on), because these hold
coefficient by coefficient; for instance, with \(h=\sum_lh_lt^l\), the coefficient of \(t^k\) in \((fg)h\) and in
\(f(gh)\) is \(\sum_{i+j+l=k}c_id_jh_l\).
\item The \textbf{value} of \(f\) at \(z\in\mathbb C\) is \(f(z)=c_0+c_1z+\dots+c_mz^m\), and \(z\) is a \textbf{root} of \(f\)
if \(f(z)=0\). For a real polynomial \(f\) and \(A\in\R^{n\times n}\), \(f(A)=c_0I_n+c_1A+\dots+c_mA^m\) (with
\(A^0=I_n\); Definition~\ref{r1-def-poly}); for instance, for \(f=t-c\), \(f(A)=A-cI\).
\end{itemize}
\end{definition}

\begin{proposition}[\citeat{hoffman-kunze}{\S4.2, Theorems~1, 2; \S4.4, Corollaries~1, 2}, \citeat{hefferon}{Five.I.1, Corollaries~1.4, 1.5}]\label{r7-prop-poly}
Let \(f\), \(g\), \(h\) be polynomials.
\begin{enumerate}
\item If \(f\ne0\) and \(g\ne0\), then \(fg\ne0\), \(\deg(fg)=\deg f+\deg g\), and the leading coefficient of \(fg\) is the
product of those of \(f\) and \(g\). If \(f\ne0\) and \(fg=fh\), then \(g=h\).
\item \((af+g)(z)=af(z)+g(z)\) and \((fg)(z)=f(z)g(z)\) for \(a,z\in\mathbb C\); the same holds for real \(f\), \(g\),
\(a\in\R\) and \(z=A\in\R^{n\times n}\).
\item (\textbf{Factor theorem}) For \(c\in\mathbb C\), \(f(c)=0\) if and only if \(f=(t-c)q\) for some polynomial \(q\).
\item A nonzero polynomial of degree \(m\) has at most \(m\) roots. Consequently, if \(f(z)=g(z)\) for infinitely many
\(z\) (for instance, for all real \(z\), or for all real \(z\ne0\)), then \(f=g\).
\item For \(\mu_1,\dots,\mu_n\in\mathbb C\) (\(n\geqslant1\)),
\((t-\mu_1)\cdots(t-\mu_n)=t^n-(\mu_1+\dots+\mu_n)t^{n-1}+r\), where \(r=0\) or \(\deg r\leqslant n-2\); and the constant
term of this product is \((-1)^n\mu_1\cdots\mu_n\).
\end{enumerate}
\end{proposition}
\begin{proof}
Write \(f=\sum_kc_kt^k\) and \(g=\sum_kd_kt^k\).

(a)
\begin{itemize}
\item Let \(\deg f=m\) and \(\deg g=l\). A term \(c_id_{m+l+j-i}\) (\(j\geqslant0\)) of the coefficient of \(t^{m+l+j}\) in
\(fg\) can be nonzero only if \(i\leqslant m\) and \(m+l+j-i\leqslant l\), that is, \(m+j\leqslant i\leqslant m\); so
\(j=0\) and \(i=m\). Hence the coefficient of \(t^{m+l}\) is \(c_md_l\ne0\), and all higher coefficients are \(0\).
\item If \(fg=fh\), then \(f(g-h)=0\); if \(g\ne h\), the first item would give \(f(g-h)\ne0\) \(\contradiction\) So
\(g=h\).
\end{itemize}
(b)
\begin{itemize}
\item \(af+g\) has the coefficients \(ac_k+d_k\), so \((af+g)(z)=\sum_k(ac_k+d_k)z^k=af(z)+g(z)\).
\item By the product rule of Definition~\ref{r7-def-polynomial} and the distributive laws,
\((fg)(z)=\sum_k\sum_{i+j=k}c_id_jz^{i+j}=\sum_{i,j}(c_iz^i)(d_jz^j)=f(z)g(z)\).
\item For real \(f\), \(g\), \(a\) and a matrix \(z=A\), both rules are Proposition~\ref{r1-prop-poly}(a),(b).
\end{itemize}
(c)
\begin{itemize}
\item First, \(f=(t-c)q+r\) for a polynomial \(q\) and a number \(r\), by induction on the degree: if \(f=0\) or
\(\deg f=0\), take \(q=0\) and \(r=f\); if \(\deg f=m\geqslant1\) and \(a\) is the leading coefficient, then
\(f-at^{m-1}(t-c)\) has coefficient \(0\) at \(t^m\), so it is \(0\) or of degree less than \(m\); by the induction
hypothesis it equals \((t-c)q_1+r\), and \(f=(t-c)(q_1+at^{m-1})+r\).
\item By (b), \(f(c)=(c-c)q(c)+r=r\). So if \(f(c)=0\), then \(r=0\) and \(f=(t-c)q\). Conversely, if \(f=(t-c)q\), then
\(f(c)=(c-c)q(c)=0\) by (b).
\end{itemize}
(d)
\begin{itemize}
\item Induction on \(m\). If \(m=0\), \(f\) is a nonzero number and has no root. Let \(m\geqslant1\), and let \(f\) have a
root \(a\) (otherwise there is nothing to prove). By (c), \(f=(t-a)q\); here \(q\ne0\), as \(f\ne0\), so \(\deg q=m-1\)
by (a).
\item If \(f(b)=0\), then \((b-a)q(b)=0\) by (b), so \(b=a\) or \(q(b)=0\). By the induction hypothesis \(q\) has at most
\(m-1\) roots, so \(f\) has at most \(m\).
\item If \(f\ne g\), then \(f-g\) is a nonzero polynomial, and every \(z\) with \(f(z)=g(z)\) is a root of it (by (b)); so
there are only finitely many such \(z\).
\end{itemize}
(e)
\begin{itemize}
\item Induction on \(n\); for \(n=1\), \(t-\mu_1\) has this form with \(r=0\).
\item If \((t-\mu_1)\cdots(t-\mu_n)=t^n-s_nt^{n-1}+r_n\) with \(s_n=\mu_1+\dots+\mu_n\), then multiplying by \(t-\mu_{n+1}\)
gives \(t^{n+1}-(s_n+\mu_{n+1})t^n+r_{n+1}\) with \(r_{n+1}=s_n\mu_{n+1}t^{n-1}+r_n(t-\mu_{n+1})\), which is \(0\) or of
degree at most \(n-1\) by (a).
\item The constant term of a polynomial is its value at \(0\), here \((0-\mu_1)\cdots(0-\mu_n)=(-1)^n\mu_1\cdots\mu_n\) by
(b).\qedhere
\end{itemize}
\end{proof}

\subsection*{The characteristic polynomial}

\begin{lemma}[cf.\ \citeat{lang}{Ch.~VIII, \S2}]\label{r7-lem-degree}
Let \(M\) be an \(m\times m\) matrix whose entries have the form \(\alpha+\beta t\) with \(\alpha,\beta\in\R\), and let
\(\det M\) be computed by cofactor expansion (Theorem~\ref{r5d-thm-expansion}) with \(t\) treated as a number. If at
most \(d\) rows of \(M\) contain an entry with \(\beta\ne0\), then \(\det M\) is \(0\) or a real polynomial of degree at
most \(d\), and its value at every \(c\in\R\) is \(\det M(c)\), the determinant of the real matrix \(M(c)\) given by
\(t=c\).
\end{lemma}
\begin{proof}
Induction on \(m\). For \(m=1\), \(\det M=\alpha+\beta t\) (Definition~\ref{r5d-def-det}(a),(d)), which is \(0\) or of
degree at most \(d\), since \(\beta=0\) if \(d=0\); its value at \(c\) is \(\alpha+\beta c=\det M(c)\). Let \(m\geqslant2\).
\begin{itemize}
\item If no entry has \(\beta\ne0\), then \(M=M(c)\) for every \(c\), and \(\det M\) is the number \(\det M(c)\).
\item Otherwise \(d\geqslant1\); let row \(i\) contain an entry with \(\beta\ne0\), write
\(m_{ij}=\alpha_{ij}+\beta_{ij}t\) for its entries, and expand along row \(i\) (Theorem~\ref{r5d-thm-expansion}):
\[
\det M=\sum_{j=1}^m(-1)^{i+j}m_{ij}\det\bigl(M(i|j)\bigr)
\]
\item Each \(M(i|j)\) is an \((m-1)\times(m-1)\) matrix with entries of the same form, and its rows are the rows of \(M\)
other than row \(i\), each shortened by one entry; so at most \(d-1\) of them contain an entry with \(\beta\ne0\). By the
induction hypothesis, \(q_j=\det\bigl(M(i|j)\bigr)\) is \(0\) or a real polynomial of degree at most \(d-1\), and its
value at \(c\) is \(\det\bigl(M(c)(i|j)\bigr)\).
\item Each term \((-1)^{i+j}(\alpha_{ij}+\beta_{ij}t)q_j\) is \(0\) or of degree at most \(1+(d-1)=d\)
(Proposition~\ref{r7-prop-poly}(a)), so its coefficients of \(t^k\), \(k>d\), are \(0\), and so are those of \(\det M\)
(Definition~\ref{r7-def-polynomial}). By Proposition~\ref{r7-prop-poly}(b), the value of \(\det M\) at \(c\) is
\(\sum_j(-1)^{i+j}(\alpha_{ij}+\beta_{ij}c)\det\bigl(M(c)(i|j)\bigr)=\det M(c)\)
(Theorem~\ref{r5d-thm-expansion}).\qedhere
\end{itemize}
\end{proof}

\begin{definition}[\citeat{lang}{Ch.~VIII, \S2}, \citeat{hoffman-kunze}{\S6.2}, \citeat{harville}{\S21.1, (21.1.3)--(21.1.8)}, \citeat{hefferon}{Five.II.3, Definition~3.9}]\label{r7-def-charpoly}
Let \(A=(a_{ij})\in\R^{n\times n}\). The \textbf{characteristic polynomial} of \(A\) is
\[
p_A(t)=\det(tI-A)
\]
computed by cofactor expansion with \(t\) treated as a number. The \((i,j)\) entry of \(tI-A\) is
\(\delta_{ij}t-a_{ij}\), so \(p_A\) is a real polynomial (Lemma~\ref{r7-lem-degree} with \(m=d=n\)). The equation
\(p_A(t)=0\) is the \textbf{characteristic equation} of \(A\). Moreover:
\begin{enumerate}
\item \(p_A(c)=\det(cI-A)\) for every \(c\in\R\); so the eigenvalues of \(A\) are the real roots of \(p_A\)
(Theorem~\ref{r7-thm-eig-singular});
\item \(p_A\) is monic of degree \(n\), its coefficient of \(t^{n-1}\) is \(-\tr A\), and its constant term is
\((-1)^n\det A\) (for \(n=1\) both say \(p_A(t)=t-a_{11}\)). For \(n=2\), \(p_A(t)=t^2-(\tr A)t+\det A\).
\end{enumerate}
\end{definition}
\begin{proof}
(a) By Lemma~\ref{r7-lem-degree}, \(p_A(c)=\det(cI-A)\) for every real \(c\). By Theorem~\ref{r7-thm-eig-singular}, a real
\(\lambda\) is an eigenvalue of \(A\) if and only if \(\det(\lambda I-A)=0\), that is, \(p_A(\lambda)=0\).

(b) Induction on \(n\). For \(n=1\), \(p_A=\det(tI_1-A)=t-a_{11}\) (Definition~\ref{r5d-def-det}(a),(d)), and
\(-a_{11}=-\tr A=(-1)^1\det A\). Let \(n\geqslant2\).
\begin{itemize}
\item Expand \(\det(tI-A)\) along row \(1\), whose entries are \(t-a_{11},-a_{12},\dots,-a_{1n}\)
(Theorem~\ref{r5d-thm-expansion}); as \((tI-A)(1|1)=tI_{n-1}-A(1|1)\),
\[
p_A=(t-a_{11})\,p_{A(1|1)}+\sum_{j=2}^n(-1)^{1+j}(-a_{1j})\det\bigl((tI-A)(1|j)\bigr)
\]
\item For \(j\geqslant2\), the rows of \((tI-A)(1|j)\) are rows \(2,\dots,n\) of \(tI-A\) without their \(j\)th entry.
Row \(i\) of \(tI-A\) contains \(t\) only in column \(i\), so the row that comes from row \(j\) does not contain \(t\),
and at most \(n-2\) rows contain \(t\). By Lemma~\ref{r7-lem-degree}, \(q_j=\det\bigl((tI-A)(1|j)\bigr)\) is \(0\) or a
real polynomial of degree at most \(n-2\).
\item By the induction hypothesis, \(p_{A(1|1)}=t^{n-1}-st^{n-2}+r\) with \(s=\tr A(1|1)=a_{22}+\dots+a_{nn}\) and
\(r=0\) or \(\deg r\leqslant n-3\) (\(r=0\) if \(n=2\)). Multiplying out,
\((t-a_{11})p_{A(1|1)}=t^n-(a_{11}+s)t^{n-1}+r_1\) with \(r_1=a_{11}st^{n-2}+(t-a_{11})r\), which is \(0\) or of degree
at most \(n-2\) (Proposition~\ref{r7-prop-poly}(a)); and \(a_{11}+s=\tr A\) (Definition~\ref{r1-def-trace}).
\item So \(p_A=(t-a_{11})p_{A(1|1)}+\sum_{j=2}^n(-1)^{1+j}(-a_{1j})q_j\) (Definition~\ref{r7-def-polynomial}). The
coefficients of \(t^k\), \(k\geqslant n-1\), of \(r_1\) and of every \(q_j\) are \(0\); so \(p_A\) has the coefficient
\(1\) at \(t^n\), the coefficient \(-\tr A\) at \(t^{n-1}\), and \(0\) above \(t^n\).
\item The constant term is the value at \(0\): by (a), \(p_A(0)=\det(-A)=(-1)^n\det A\)
(Theorem~\ref{r5d-thm-product}(c)).\qedhere
\end{itemize}
\end{proof}
In examples, \(\det(tI-A)\) is also computed by row and column operations with \(t\) treated as a number
(Definition~\ref{r5d-def-det}, Theorem~\ref{r5d-thm-transpose}); the result is \(p_A\), since its value at every real
\(t\) is \(\det(tI-A)\), by (a) (Proposition~\ref{r7-prop-poly}(d)). Many books call \(\det(A-tI)=(-1)^np_A(t)\) the characteristic polynomial; it has
the same roots. By (a), (b) and Proposition~\ref{r7-prop-poly}(d), a matrix \(A\in\R^{n\times n}\) has at most \(n\)
eigenvalues; so the four eigenvalues \(1\), \(-1\), \(2\), \(-2\) of the \(4\times4\) matrix of
Example~\ref{r7-exm-two} are all its eigenvalues.

\needspace{10\baselineskip}
\begin{example}[\citeat{hefferon}{Five.II.3, Examples~3.7, 3.14}]\label{r7-exm-charpoly}
Find the eigenvalues of
\[
A=\begin{pmatrix}1&2&1\\2&0&-2\\-1&2&3\end{pmatrix}
\]
and a basis of each eigenspace.

\textit{Solution.}
\begin{itemize}
\item For \(t\in\R\), expanding along row \(1\) (Theorem~\ref{r5d-thm-expansion}),
\[
\begin{aligned}
\det(tI-A)&=\det\begin{pmatrix}t-1&-2&-1\\-2&t&2\\1&-2&t-3\end{pmatrix}
=(t-1)\bigl(t(t-3)+4\bigr)+2\bigl(-2(t-3)-2\bigr)-(4-t)\\
&=(t^3-4t^2+7t-4)+(-4t+8)+(t-4)=t^3-4t^2+4t=t(t-2)^2
\end{aligned}
\]
So \(p_A(t)=t(t-2)^2\), and the eigenvalues are \(0\) and \(2\) (Definition~\ref{r7-def-charpoly}).
\item \(\lambda=0\): row-reduce \(A-0I=A\).
\[
\begin{gathered}
A\xrightarrow[R_3\leftarrow R_3+R_1]{R_2\leftarrow R_2-2R_1}
\begin{pmatrix}1&2&1\\0&-4&-4\\0&4&4\end{pmatrix}
\xrightarrow{R_3\leftarrow R_3+R_2}
\begin{pmatrix}1&2&1\\0&-4&-4\\0&0&0\end{pmatrix}
\\
\xrightarrow{R_2\leftarrow-\frac14R_2}
\begin{pmatrix}1&2&1\\0&1&1\\0&0&0\end{pmatrix}
\xrightarrow{R_1\leftarrow R_1-2R_2}
\begin{pmatrix}1&0&-1\\0&1&1\\0&0&0\end{pmatrix}
\end{gathered}
\]
The free unknown is \(x_3\); the special solution \((1,-1,1)'\) is a basis of \(\mathcal N(A)\)
(Theorem~\ref{thm:r4-four-bases}(c)).
\item \(\lambda=2\): row-reduce \(A-2I\).
\[
\begin{pmatrix}-1&2&1\\2&-2&-2\\-1&2&1\end{pmatrix}
\xrightarrow[R_3\leftarrow R_3-R_1]{R_2\leftarrow R_2+2R_1}
\begin{pmatrix}-1&2&1\\0&2&0\\0&0&0\end{pmatrix}
\xrightarrow[R_2\leftarrow\frac12R_2]{R_1\leftarrow-R_1}
\begin{pmatrix}1&-2&-1\\0&1&0\\0&0&0\end{pmatrix}
\xrightarrow{R_1\leftarrow R_1+2R_2}
\begin{pmatrix}1&0&-1\\0&1&0\\0&0&0\end{pmatrix}
\]
The special solution \((1,0,1)'\) is a basis of \(\mathcal N(A-2I)\).
\end{itemize}
Both eigenspaces are lines, although \(p_A\) has the factor \((t-2)^2\).
\end{example}

\begin{theorem}[\citeat{hoffman-kunze}{\S4.5, Theorem~9}, \citeat{lang}{Ch.~IX, \S1}, \citeat{hefferon}{Five.I.1}]\label{r7-thm-factor}
\begin{enumerate}
\item (\textbf{Fundamental theorem of algebra}; not proved here.) Every polynomial of degree at least \(1\) has a root
in \(\mathbb C\).
\item Every monic polynomial \(f\) of degree \(n\geqslant1\) is a product
\[
f=(t-\lambda_1)(t-\lambda_2)\cdots(t-\lambda_n)\qquad(\lambda_1,\dots,\lambda_n\in\mathbb C)
\]
and the list \(\lambda_1,\dots,\lambda_n\) is unique up to order: if also \(f=(t-\mu_1)\cdots(t-\mu_n)\), then
\(\mu_1,\dots,\mu_n\) are \(\lambda_1,\dots,\lambda_n\) in some order.
\item In (b), the roots of \(f\) are exactly \(\lambda_1,\dots,\lambda_n\).
\end{enumerate}
\end{theorem}
\begin{proof}
(b)
\begin{itemize}
\item Existence, by induction on \(n\). For \(n=1\), \(f=t+c_0=t-(-c_0)\). Let \(n\geqslant2\). By (a), \(f\) has a root
\(\lambda_n\), so \(f=(t-\lambda_n)g\) (Proposition~\ref{r7-prop-poly}(c)); here \(g\ne0\), as \(f\ne0\), so \(g\) is monic of
degree \(n-1\) (Proposition~\ref{r7-prop-poly}(a)). By the induction hypothesis, \(g=(t-\lambda_1)\cdots(t-\lambda_{n-1})\).
\item Uniqueness, by induction on \(n\). For \(n=1\), \(t-\lambda_1=t-\mu_1\) gives \(\lambda_1=\mu_1\) (constant terms).
Let \(n\geqslant2\). The value of both products at \(\lambda_n\) (Proposition~\ref{r7-prop-poly}(b)) gives
\(0=(\lambda_n-\mu_1)\cdots(\lambda_n-\mu_n)\), so \(\lambda_n=\mu_k\) for some \(k\).
\item Renumber the \(\mu_i\) so that \(k=n\), and cancel the factor \(t-\lambda_n\ne0\) (Proposition~\ref{r7-prop-poly}(a)):
\((t-\lambda_1)\cdots(t-\lambda_{n-1})=(t-\mu_1)\cdots(t-\mu_{n-1})\). By the induction hypothesis,
\(\mu_1,\dots,\mu_{n-1}\) are \(\lambda_1,\dots,\lambda_{n-1}\) in some order, and \(\mu_n=\lambda_n\).
\end{itemize}
(c) By Proposition~\ref{r7-prop-poly}(b), \(f(z)=(z-\lambda_1)\cdots(z-\lambda_n)\), which is \(0\) if and only if
\(z=\lambda_i\) for some \(i\).
\end{proof}

\begin{definition}[\citeat{bapat}{\S3.1}, \citeat{hefferon}{Five.II.3, Definition~3.17}, \citeat{hoffman-kunze}{\S6.2, Example~1}, \citeat{harville}{\S21.1}]\label{r7-def-eigenvalues}
Let \(A\in\R^{n\times n}\). By Definition~\ref{r7-def-charpoly}(b) and Theorem~\ref{r7-thm-factor}(b),
\[
p_A(t)=(t-\lambda_1)(t-\lambda_2)\cdots(t-\lambda_n)
\]
for complex numbers \(\lambda_1,\dots,\lambda_n\), unique up to order. These \(n\) numbers are the \textbf{eigenvalues}
of \(A\), listed \textbf{with multiplicity}; by Theorem~\ref{r7-thm-factor}(c) they are the roots of \(p_A\). For
\(c\in\mathbb C\), the \textbf{algebraic multiplicity} \(a_A(c)\) is the number of indices \(i\) with \(\lambda_i=c\)
(\(a_A(c)=0\) if \(c\) is not an eigenvalue); so if \(c_1,\dots,c_k\) are the distinct eigenvalues, then
\(p_A=(t-c_1)^{a_A(c_1)}\cdots(t-c_k)^{a_A(c_k)}\) and \(a_A(c_1)+\dots+a_A(c_k)=n\). An eigenvalue \(c\) is
\textbf{simple} if \(a_A(c)=1\) and \textbf{multiple} if \(a_A(c)\geqslant2\). A real number is an eigenvalue in this
sense if and only if it is one in the sense of Definition~\ref{r7-def-eigen}, since both mean that it is a real root
of \(p_A\) (Definition~\ref{r7-def-charpoly}(a)); these are the \textbf{real eigenvalues}. From now on
``eigenvalue'' has this meaning; eigenvectors and eigenspaces belong to real eigenvalues only.
\end{definition}

\begin{proposition}[\citeat{hoffman-kunze}{\S4.5}, \citeat{abadir-magnus}{Exercise~7.17(b)}]\label{r7-prop-conjugate}
Let \(A\in\R^{n\times n}\).
\begin{enumerate}
\item If \(z\) is a non-real eigenvalue of \(A\), then so is \(\bar z\), and \(a_A(\bar z)=a_A(z)\): the non-real
eigenvalues come in conjugate pairs.
\item If \(n\) is odd, then \(A\) has a real eigenvalue.
\end{enumerate}
\end{proposition}
\begin{proof}
(a)
\begin{itemize}
\item For a polynomial \(f=\sum_kc_kt^k\) put \(\bar f=\sum_k\bar c_kt^k\). Conjugation respects sums and products of
numbers, so the coefficient \(\overline{\sum_{i+j=k}c_id_j}\) of \(\overline{fg}\) is \(\sum_{i+j=k}\bar c_i\bar d_j\):
\(\overline{fg}=\bar f\,\bar g\). Hence
\(\overline{(t-\lambda_1)\cdots(t-\lambda_n)}=(t-\bar\lambda_1)\cdots(t-\bar\lambda_n)\).
\item \(p_A\) is real, so \(p_A=\overline{p_A}\), that is,
\((t-\lambda_1)\cdots(t-\lambda_n)=(t-\bar\lambda_1)\cdots(t-\bar\lambda_n)\) (Definition~\ref{r7-def-eigenvalues}). By
Theorem~\ref{r7-thm-factor}(b), \(\bar\lambda_1,\dots,\bar\lambda_n\) are \(\lambda_1,\dots,\lambda_n\) in some order.
\item So the number of indices \(i\) with \(\lambda_i=\bar z\) equals the number of \(i\) with \(\bar\lambda_i=\bar z\),
that is, with \(\lambda_i=z\): \(a_A(\bar z)=a_A(z)\geqslant1\).
\end{itemize}
(b) By (a), the non-real numbers in the list \(\lambda_1,\dots,\lambda_n\) fall into pairs \(z\ne\bar z\) that occur
equally often, so there is an even number of them. If \(n\) is odd, some \(\lambda_i\) is real: a real eigenvalue.
\end{proof}

\begin{example}[cf.\ \citeat{hoffman-kunze}{\S6.2, Example~1}]\label{r7-exm-complex}
Find the eigenvalues of
\[
A=\begin{pmatrix}2&-1&0\\1&2&0\\1&1&1\end{pmatrix}
\]
and the eigenspace of each real eigenvalue.

\textit{Solution.}
\begin{itemize}
\item For \(t\in\R\), column \(3\) of \(tI-A\) has the entry \(t-1\) in position \((3,3)\) and \(0\) elsewhere; expanding
along it (Theorem~\ref{r5d-thm-expansion}),
\[
\det(tI-A)=\det\begin{pmatrix}t-2&1&0\\-1&t-2&0\\-1&-1&t-1\end{pmatrix}
=(t-1)\det\begin{pmatrix}t-2&1\\-1&t-2\end{pmatrix}=(t-1)\bigl((t-2)^2+1\bigr)
\]
\item Since \((t-2-i)(t-2+i)=(t-2)^2-i^2=(t-2)^2+1\), \(p_A(t)=(t-1)\bigl(t-(2+i)\bigr)\bigl(t-(2-i)\bigr)\): the
eigenvalues are \(1\), \(2+i\) and \(2-i\), each simple, and \(2\pm i\) is a conjugate pair
(Proposition~\ref{r7-prop-conjugate}(a)). The only real eigenvalue is \(1\).
\item Row-reduce \(A-I\):
\[
\begin{gathered}
\begin{pmatrix}1&-1&0\\1&1&0\\1&1&0\end{pmatrix}
\xrightarrow[R_3\leftarrow R_3-R_1]{R_2\leftarrow R_2-R_1}
\begin{pmatrix}1&-1&0\\0&2&0\\0&2&0\end{pmatrix}
\xrightarrow{R_3\leftarrow R_3-R_2}
\begin{pmatrix}1&-1&0\\0&2&0\\0&0&0\end{pmatrix}
\\
\xrightarrow{R_2\leftarrow\frac12R_2}
\begin{pmatrix}1&-1&0\\0&1&0\\0&0&0\end{pmatrix}
\xrightarrow{R_1\leftarrow R_1+R_2}
\begin{pmatrix}1&0&0\\0&1&0\\0&0&0\end{pmatrix}
\end{gathered}
\]
So \(x_1=x_2=0\) and \(x_3\) is free: \(\mathcal N(A-I)=\Span\{e_3\}\). The eigenvalues \(2\pm i\) have no eigenvectors in
\(\R^3\) (Definition~\ref{r7-def-eigenvalues}).
\end{itemize}
\end{example}

\begin{remark}[Stationary autoregressions]\label{r7-rem-ar}
In time series, the \textbf{autoregressive model} of order \(p\) is
\(y_k=\beta_0+\beta_1y_{k-1}+\dots+\beta_py_{k-p}+\varepsilon_k\). It describes a stationary process when every root
\(z\) of the real polynomial \(\Phi(t)=1-\beta_1t-\dots-\beta_pt^p\) lies outside the unit circle, that is,
\(\lvert z\rvert>1\), where \(\lvert a+bi\rvert=\sqrt{a^2+b^2}\) (not proved
here). The recursion \(y_k=\beta_0+\beta_1y_{k-1}+\dots+\beta_py_{k-p}\), written for the vector
\((y_{k-p+1},\dots,y_k)'\), uses the \(p\times p\) \textbf{companion matrix}
\[
C=\begin{pmatrix}0&1&0&\cdots&0\\0&0&1&\cdots&0\\\vdots&\vdots&&\ddots&\vdots\\0&0&0&\cdots&1\\
\beta_p&\beta_{p-1}&\beta_{p-2}&\cdots&\beta_1\end{pmatrix}
\]
whose eigenvalues are the roots of \(f=t^p-\beta_1t^{p-1}-\dots-\beta_p\)~\citeat{abadir-magnus}{Appendix~A.4.1}. As
\(\Phi(0)=1\) and \(z^p\Phi(1/z)=f(z)\) for \(z\ne0\), the nonzero roots of \(f\) are the numbers \(1/z\) for the roots
\(z\) of \(\Phi\), and \(\lvert1/z\rvert=1/\lvert z\rvert\) (as \(1/z=\bar z/(a^2+b^2)\) for \(z=a+bi\ne0\)); so the condition says that every eigenvalue \(\lambda\) of \(C\), real or not, has
\(\lvert\lambda\rvert<1\).
\end{remark}

\subsection*{Similar and triangular matrices}

\begin{definition}[\citeat{hoffman-kunze}{\S3.4}, \citeat{hefferon}{Five.II.1, Definition~1.2}]\label{r7-def-similar}
A matrix \(B\in\R^{n\times n}\) is \textbf{similar} to \(A\in\R^{n\times n}\) if \(B=P^{-1}AP\) for some invertible
\(P\in\R^{n\times n}\). Then \(A=PBP^{-1}=(P^{-1})^{-1}BP^{-1}\), so \(A\) is similar to \(B\).
\end{definition}

\begin{theorem}[\citeat{hoffman-kunze}{\S6.2}, \citeat{harville}{Theorem~21.3.1}]\label{r7-thm-similar}
Let \(A,P\in\R^{n\times n}\) with \(P\) invertible, and let \(B=P^{-1}AP\).
\begin{enumerate}
\item \(p_B=p_A\). Hence \(A\) and \(B\) have the same eigenvalues with the same algebraic multiplicities, \(\tr B=\tr A\)
and \(\det B=\det A\).
\item For every \(c\in\R\), \(x\mapsto P^{-1}x\) maps \(\mathcal N(A-cI)\) onto \(\mathcal N(B-cI)\), with inverse
\(y\mapsto Py\); so \(x\) is an eigenvector of \(A\) belonging to \(c\) if and only if \(P^{-1}x\) is an eigenvector of \(B\)
belonging to \(c\). Moreover \(\dim\mathcal N(B-cI)=\dim\mathcal N(A-cI)\), and \(\rank B=\rank A\).
\end{enumerate}
\end{theorem}
\begin{proof}
(a)
\begin{itemize}
\item For \(c\in\R\), \(cI-B=P^{-1}(cI)P-P^{-1}AP=P^{-1}(cI-A)P\) (Theorem~\ref{r1-thm-rules}), so
\(\det(cI-B)=\det P^{-1}\cdot\det(cI-A)\cdot\det P=\det(cI-A)\) (Theorem~\ref{r5d-thm-product}(b),(c)).
\item So \(p_B(c)=p_A(c)\) for every real \(c\) (Definition~\ref{r7-def-charpoly}(a)), and \(p_B=p_A\)
(Proposition~\ref{r7-prop-poly}(d)).
\item The eigenvalues of \(A\) and \(B\) come from the same factorization (Definition~\ref{r7-def-eigenvalues}); \(-\tr B\)
and \(-\tr A\) are the same coefficient of \(t^{n-1}\), and \((-1)^n\det B\), \((-1)^n\det A\) the same constant term
(Definition~\ref{r7-def-charpoly}(b)).
\end{itemize}
(b)
\begin{itemize}
\item \(B-cI=P^{-1}(A-cI)P\) and \(A-cI=P(B-cI)P^{-1}\). If \((A-cI)x=0\), then \((B-cI)P^{-1}x=P^{-1}(A-cI)x=0\); if
\((B-cI)y=0\), then \((A-cI)Py=P(B-cI)y=0\).
\item The two maps are inverse to each other (\(PP^{-1}x=x\), \(P^{-1}Py=y\)), and \(P^{-1}x\ne0\) if and only if \(x\ne0\).
\item \(\rank(B-cI)=\rank\bigl(P^{-1}(A-cI)P\bigr)=\rank(A-cI)\) (Theorem~\ref{thm:r4-rank-product}(d)), so the two null
spaces have the same dimension (Theorem~\ref{thm:r4-rank-nullity}). In the same way,
\(\rank B=\rank(P^{-1}AP)=\rank A\).\qedhere
\end{itemize}
\end{proof}

\begin{proposition}[\citeat{hoffman-kunze}{\S6.4}, \citeat{harville}{Lemma~21.2.1, Corollary~21.2.2}]\label{r7-prop-block-triangular}
\begin{enumerate}
\item Let \(A=\begin{pmatrix}B&C\\0&D\end{pmatrix}\) or \(A=\begin{pmatrix}B&0\\C&D\end{pmatrix}\) with \(B\in\R^{p\times p}\)
and \(D\in\R^{q\times q}\). Then \(p_A=p_Bp_D\): the eigenvalues of \(A\), listed with multiplicity, are those of \(B\)
together with those of \(D\).
\item If \(A\in\R^{n\times n}\) is upper or lower triangular, then \(p_A=(t-a_{11})\cdots(t-a_{nn})\): the eigenvalues of
\(A\), listed with multiplicity, are \(a_{11},\dots,a_{nn}\).
\end{enumerate}
\end{proposition}
\begin{proof}
(a)
\begin{itemize}
\item For \(c\in\R\), \(cI-A\) is block triangular of the same kind, with diagonal blocks \(cI_p-B\) and \(cI_q-D\)
(Theorem~\ref{r5b-thm-block}(a)). By Theorem~\ref{r5d-thm-block-tri},
\(\det(cI-A)=\det(cI_p-B)\det(cI_q-D)=p_B(c)p_D(c)\), the value of \(p_Bp_D\) at \(c\)
(Proposition~\ref{r7-prop-poly}(b)). So \(p_A=p_Bp_D\) (Proposition~\ref{r7-prop-poly}(d)).
\item If \(p_B=(t-\beta_1)\cdots(t-\beta_p)\) and \(p_D=(t-\delta_1)\cdots(t-\delta_q)\) (Definition~\ref{r7-def-eigenvalues}),
then \(p_A\) is the product of these \(p+q\) linear factors, and by Theorem~\ref{r7-thm-factor}(b) the eigenvalues of
\(A\) are \(\beta_1,\dots,\beta_p,\delta_1,\dots,\delta_q\).
\end{itemize}
(b)
\begin{itemize}
\item For \(c\in\R\), \(cI-A\) is triangular of the same kind with diagonal entries \(c-a_{ii}\), so
\(\det(cI-A)=(c-a_{11})\cdots(c-a_{nn})\) (Theorem~\ref{r5d-thm-basic}(c) for upper, Theorem~\ref{r5d-thm-transpose}
for lower triangular matrices), the value of \((t-a_{11})\cdots(t-a_{nn})\) at \(c\).
\item As in (a), \(p_A=(t-a_{11})\cdots(t-a_{nn})\), and Definition~\ref{r7-def-eigenvalues} gives the
eigenvalues.\qedhere
\end{itemize}
\end{proof}

\begin{example}\label{r7-exm-three}
Find the eigenvalues of
\[
A=\begin{pmatrix}-1&2&3&-1\\2&-1&1&1\\0&0&1&2\\0&0&2&1\end{pmatrix}=\begin{pmatrix}B&F\\0&D\end{pmatrix},\qquad
B=\begin{pmatrix}-1&2\\2&-1\end{pmatrix},\qquad D=\begin{pmatrix}1&2\\2&1\end{pmatrix}
\]
and a basis of each eigenspace.

\textit{Solution.}
\begin{itemize}
\item By Proposition~\ref{r7-prop-block-triangular}(a) and the formula for \(n=2\) in
Definition~\ref{r7-def-charpoly}(b), with \(\tr B=-2\), \(\det B=-3\), \(\tr D=2\) and \(\det D=-3\),
\[
p_A=p_Bp_D=(t^2+2t-3)(t^2-2t-3)=(t+3)(t-1)(t-3)(t+1)
\]
So the eigenvalues are \(-3\), \(-1\), \(1\) and \(3\), each simple (Definition~\ref{r7-def-eigenvalues}). Check:
\(p_A=(t^2-9)(t^2-1)=t^4-10t^2+9\); the coefficient \(0\) of \(t^3\) is \(-\tr A=-(-1-1+1+1)\), and the constant term
\(9\) is \((-1)^4\det A\), as \(\det A=\det B\det D=(-3)(-3)=9\) (Theorem~\ref{r5d-thm-block-tri},
Definition~\ref{r7-def-charpoly}(b)).
\item \(\lambda=-3\):
\[
\begin{gathered}
A+3I=\begin{pmatrix}2&2&3&-1\\2&2&1&1\\0&0&4&2\\0&0&2&4\end{pmatrix}
\xrightarrow{R_2\leftarrow R_2-R_1}
\begin{pmatrix}2&2&3&-1\\0&0&-2&2\\0&0&4&2\\0&0&2&4\end{pmatrix}
\xrightarrow[R_4\leftarrow R_4+R_2]{R_3\leftarrow R_3+2R_2}
\begin{pmatrix}2&2&3&-1\\0&0&-2&2\\0&0&0&6\\0&0&0&6\end{pmatrix}
\\
\xrightarrow{R_4\leftarrow R_4-R_3}
\begin{pmatrix}2&2&3&-1\\0&0&-2&2\\0&0&0&6\\0&0&0&0\end{pmatrix}
\end{gathered}
\]
An echelon form with three nonzero rows: \(\rank(A+3I)=3\) (Definition~\ref{def:r4-rank}), so
\(\dim\mathcal N(A+3I)=4-3=1\) (Theorem~\ref{thm:r4-rank-nullity}). Back substitution gives \(x_4=0\), \(x_3=0\) and
\(x_1=-x_2\); so \((-1,1,0,0)'\) is a basis of \(\mathcal N(A+3I)\) (Theorem~\ref{thm:r4-counting}(d)).
\item \(\lambda=-1\):
\[
A+I=\begin{pmatrix}0&2&3&-1\\2&0&1&1\\0&0&2&2\\0&0&2&2\end{pmatrix}
\xrightarrow[R_4\leftarrow R_4-R_3]{R_1\leftrightarrow R_2}
\begin{pmatrix}2&0&1&1\\0&2&3&-1\\0&0&2&2\\0&0&0&0\end{pmatrix}
\]
Again \(\rank(A+I)=3\) and \(\dim\mathcal N(A+I)=1\). With \(x_4=1\), back substitution gives \(x_3=-1\), \(x_2=2\)
and \(x_1=0\); so \((0,2,-1,1)'\) is a basis of \(\mathcal N(A+I)\).
\item \(\lambda=1\):
\[
\begin{gathered}
A-I=\begin{pmatrix}-2&2&3&-1\\2&-2&1&1\\0&0&0&2\\0&0&2&0\end{pmatrix}
\xrightarrow{R_2\leftarrow R_2+R_1}
\begin{pmatrix}-2&2&3&-1\\0&0&4&0\\0&0&0&2\\0&0&2&0\end{pmatrix}
\xrightarrow{R_2\leftrightarrow R_4}
\begin{pmatrix}-2&2&3&-1\\0&0&2&0\\0&0&0&2\\0&0&4&0\end{pmatrix}
\\
\xrightarrow{R_4\leftarrow R_4-2R_2}
\begin{pmatrix}-2&2&3&-1\\0&0&2&0\\0&0&0&2\\0&0&0&0\end{pmatrix}
\end{gathered}
\]
So \(\rank(A-I)=3\), \(\dim\mathcal N(A-I)=1\), and back substitution gives \(x_3=x_4=0\) and \(x_1=x_2\): the vector
\((1,1,0,0)'\) is a basis of \(\mathcal N(A-I)\).
\item \(\lambda=3\):
\[
A-3I=\begin{pmatrix}-4&2&3&-1\\2&-4&1&1\\0&0&-2&2\\0&0&2&-2\end{pmatrix}
\xrightarrow{R_1\leftrightarrow R_2}
\begin{pmatrix}2&-4&1&1\\-4&2&3&-1\\0&0&-2&2\\0&0&2&-2\end{pmatrix}
\xrightarrow[R_4\leftarrow R_4+R_3]{R_2\leftarrow R_2+2R_1}
\begin{pmatrix}2&-4&1&1\\0&-6&5&1\\0&0&-2&2\\0&0&0&0\end{pmatrix}
\]
So \(\rank(A-3I)=3\), \(\dim\mathcal N(A-3I)=1\), and back substitution gives \(x_3=x_4\), \(x_2=x_4\) and
\(x_1=x_4\): the vector \(\one_4\) is a basis of \(\mathcal N(A-3I)\). Indeed every row of \(A\) adds up to \(3\).
\end{itemize}
Check: \(A(-1,1,0,0)'=(3,-3,0,0)'\) and \(A(0,2,-1,1)'=(0,-2,1,-1)'\).
\end{example}

\subsection*{Eigenvalues of related matrices; trace and determinant}

\begin{proposition}[\citeat{harville}{Lemmas~21.1.2, 21.1.3, \S21.10}, \citeat{searle-khuri}{\S\S6.2.2, 6.2.3}, \citeat{hoffman-kunze}{\S6.2}]\label{r7-prop-eig-poly}
Let \(A\in\R^{n\times n}\).
\begin{enumerate}
\item If \(Ax=\lambda x\) with \(\lambda\in\R\) and \(x\ne0\), then \(A^kx=\lambda^kx\) (\(k=0,1,2,\dots\)) and
\(f(A)x=f(\lambda)x\) for every real polynomial \(f\); in particular \((cA)x=(c\lambda)x\) and \((A+cI)x=(\lambda+c)x\)
(\(c\in\R\)). So \(x\) is an eigenvector of \(A^k\), \(f(A)\), \(cA\) and \(A+cI\), belonging to \(\lambda^k\),
\(f(\lambda)\), \(c\lambda\) and \(\lambda+c\).
\item If, moreover, \(A\) is invertible, then \(\lambda\ne0\) and \(A^{-1}x=\lambda^{-1}x\).
\item If \(A\) has the eigenvalues \(\lambda_1,\dots,\lambda_n\) (with multiplicity) and \(c\in\R\), then \(A+cI\) has the
eigenvalues \(\lambda_1+c,\dots,\lambda_n+c\), and, if \(c\ne0\), \(cA\) has the eigenvalues \(c\lambda_1,\dots,c\lambda_n\).
\item \(p_{A'}=p_A\): \(A'\) and \(A\) have the same eigenvalues with the same algebraic multiplicities.
\end{enumerate}
\end{proposition}
\begin{proof}
(a)
\begin{itemize}
\item By induction on \(k\): \(A^0x=x=\lambda^0x\), and if \(A^kx=\lambda^kx\), then
\(A^{k+1}x=A^k(Ax)=A^k(\lambda x)=\lambda A^kx=\lambda^{k+1}x\) (Theorem~\ref{r1-thm-rules}).
\item If \(f=\sum_kc_kt^k\), then \(f(A)x=\sum_kc_kA^kx=\sum_kc_k\lambda^kx=f(\lambda)x\). For \(f=ct\) and \(f=t+c\) this
gives \((cA)x=c\lambda x\) and \((A+cI)x=(\lambda+c)x\). In each case \(x\ne0\).
\end{itemize}
(b) If \(\lambda=0\), then \(Ax=0\) with \(x\ne0\), and \(A\) would not be invertible (Theorem~\ref{r2-thm-invertible})
\(\contradiction\) So \(\lambda\ne0\), and \(x=A^{-1}(Ax)=A^{-1}(\lambda x)=\lambda A^{-1}x\), that is,
\(A^{-1}x=\lambda^{-1}x\).

(c)
\begin{itemize}
\item For \(s\in\R\), \(\det\bigl(sI-(A+cI)\bigr)=\det\bigl((s-c)I-A\bigr)=p_A(s-c)=(s-c-\lambda_1)\cdots(s-c-\lambda_n)\)
(Definition~\ref{r7-def-charpoly}(a), Proposition~\ref{r7-prop-poly}(b)), the value at \(s\) of
\(\bigl(t-(\lambda_1+c)\bigr)\cdots\bigl(t-(\lambda_n+c)\bigr)\). By Proposition~\ref{r7-prop-poly}(d) this product is
\(p_{A+cI}\), and Definition~\ref{r7-def-eigenvalues} gives the eigenvalues.
\item Let \(c\ne0\). For \(s\in\R\), \(\det(sI-cA)=\det\bigl(c(\tfrac scI-A)\bigr)=c^np_A(\tfrac sc)
=c^n(\tfrac sc-\lambda_1)\cdots(\tfrac sc-\lambda_n)=(s-c\lambda_1)\cdots(s-c\lambda_n)\)
(Theorem~\ref{r5d-thm-product}(c)); conclude as before.
\end{itemize}
(d) For \(c\in\R\), \((cI-A)'=cI-A'\) (Theorem~\ref{r1-thm-transpose}), so \(\det(cI-A')=\det(cI-A)\)
(Theorem~\ref{r5d-thm-transpose}). By Definition~\ref{r7-def-charpoly}(a) and Proposition~\ref{r7-prop-poly}(d),
\(p_{A'}=p_A\).
\end{proof}
In Example~\ref{r7-exm-two}, by (a), \(v_1\) and \(v_2\) are eigenvectors of \(A^2\) belonging to \(1\), \(v_3\) and
\(v_4\) eigenvectors of \(A^2\) belonging to \(4\), and \(v_1,\dots,v_4\) eigenvectors of \(A-2I\) belonging to \(-1\),
\(-3\), \(0\), \(-4\). The eigenvalues of that \(A\) are \(1\), \(-1\), \(2\), \(-2\) (the remark after
Definition~\ref{r7-def-charpoly}), so \(A\) is invertible (Theorem~\ref{r7-thm-eig-singular} with \(\lambda=0\)), and
by (b), \(v_1,\dots,v_4\) are eigenvectors of \(A^{-1}\) belonging to \(1\), \(-1\), \(\frac12\), \(-\frac12\).

\begin{theorem}[\citeat{bapat}{\S3.1}, \citeat{abadir-magnus}{Exercise~7.16}]\label{r7-thm-trace-det}
Let \(A\in\R^{n\times n}\) have the eigenvalues \(\lambda_1,\dots,\lambda_n\), listed with multiplicity. Then
\[
\tr A=\lambda_1+\lambda_2+\dots+\lambda_n,\qquad\det A=\lambda_1\lambda_2\cdots\lambda_n
\]
Hence \(A\) is invertible if and only if \(0\) is not an eigenvalue of \(A\).
\end{theorem}
\begin{proof}
\begin{itemize}
\item \(p_A=(t-\lambda_1)\cdots(t-\lambda_n)\) (Definition~\ref{r7-def-eigenvalues}). By Proposition~\ref{r7-prop-poly}(e),
its coefficient of \(t^{n-1}\) is \(-(\lambda_1+\dots+\lambda_n)\) and its constant term is \((-1)^n\lambda_1\cdots\lambda_n\).
\item By Definition~\ref{r7-def-charpoly}(b), the same two coefficients are \(-\tr A\) and \((-1)^n\det A\).
\item \(A\) is invertible if and only if \(\det A\ne0\) (Theorem~\ref{r5d-thm-product}(a)), that is, if and only if no
\(\lambda_i\) is \(0\).\qedhere
\end{itemize}
\end{proof}
For the matrix of Example~\ref{r7-exm-charpoly}, \(\tr A=4=0+2+2\) and \(\det A=0=0\cdot2\cdot2\), so \(A\) is not
invertible; for that of Example~\ref{r7-exm-complex}, \(\tr A=5=1+(2+i)+(2-i)\) and \(\det A=5=1\cdot(2+i)(2-i)\).

\needspace{12\baselineskip}
\subsection*{Eigenvectors of distinct eigenvalues; multiplicities}

\begin{theorem}[\citeat{hefferon}{Five.II.3, Theorem~3.18, Corollary~3.20}, \citeat{lang}{Ch.~VIII, \S1, Remark~1}]\label{r7-thm-indep}
Let \(v_1,\dots,v_m\) be eigenvectors of \(A\in\R^{n\times n}\) belonging to real eigenvalues
\(\lambda_1,\dots,\lambda_m\) that are pairwise distinct (\(\lambda_i\ne\lambda_j\) for \(i\ne j\)). Then
\(v_1,\dots,v_m\) are linearly independent. In particular, if \(A\) has \(n\) distinct real eigenvalues, then
eigenvectors belonging to them form a basis of \(\R^n\).
\end{theorem}
\begin{proof}
Induction on \(m\). For \(m=0\) the empty list is independent (Definition~\ref{r3-def-independence}). Let
\(m\geqslant1\), assume the statement for \(m-1\) eigenvectors, and let \(c_1v_1+\dots+c_mv_m=0\).
\begin{itemize}
\item Multiplying by \(\lambda_m\) gives \(\sum_ic_i\lambda_mv_i=0\); multiplying by \(A\) gives \(\sum_ic_i\lambda_iv_i=0\),
since \(Av_i=\lambda_iv_i\) (Theorem~\ref{r1-thm-rules}).
\item Subtracting, the \(v_m\) term vanishes:
\(c_1(\lambda_m-\lambda_1)v_1+\dots+c_{m-1}(\lambda_m-\lambda_{m-1})v_{m-1}=0\).
\item By the induction hypothesis, \(c_i(\lambda_m-\lambda_i)=0\) for \(i<m\), and \(c_i=0\) since \(\lambda_m\ne\lambda_i\).
\item Then \(c_mv_m=0\) with \(v_m\ne0\), so \(c_m=0\) (Proposition~\ref{r3-prop-dependence}(a)).
\end{itemize}
If \(A\) has \(n\) distinct real eigenvalues, eigenvectors belonging to them are \(n\) independent vectors of \(\R^n\),
which has dimension \(n\) (Theorem~\ref{thm:r4-dimension}); so they form a basis (Theorem~\ref{thm:r4-counting}(d)).
\end{proof}

\begin{definition}[\citeat{harville}{\S21.1}, \citeat{hefferon}{Five.II.3, Definition~3.17}, \citeat{bapat}{\S3.2}]\label{r7-def-multiplicity}
Let \(A\in\R^{n\times n}\) and \(c\in\R\). The \textbf{geometric multiplicity} of \(c\) is
\[
g_A(c)=\dim\mathcal N(A-cI)=n-\rank(A-cI)
\]
(Theorem~\ref{thm:r4-rank-nullity}). So \(g_A(c)\geqslant1\) if \(c\) is a real eigenvalue of \(A\), and \(g_A(c)=0\)
otherwise (Definition~\ref{r7-def-eigen}, Theorem~\ref{thm:r4-dimension}).
\end{definition}
By Theorem~\ref{r7-thm-similar}(a),(b), similar matrices have the same algebraic and geometric multiplicities. In
Example~\ref{r7-exm-charpoly}, \(g_A(0)=a_A(0)=1\), but \(g_A(2)=1<2=a_A(2)\).

\begin{theorem}[\citeat{hoffman-kunze}{\S6.4, Example~10}, \citeat{harville}{Corollary~21.3.5}, \citeat{hefferon}{Five.II.3}]\label{r7-thm-geo-alg}
If \(c\) is a real eigenvalue of \(A\in\R^{n\times n}\), then \(1\leqslant g_A(c)\leqslant a_A(c)\). In particular, if
\(c\) is simple, then \(g_A(c)=1\), that is, \(\rank(A-cI)=n-1\).
\end{theorem}
\begin{proof}
\begin{itemize}
\item \(g=g_A(c)\geqslant1\) (Definition~\ref{r7-def-multiplicity}). Extend a basis \(x_1,\dots,x_g\) of \(\mathcal N(A-cI)\)
to a basis \(x_1,\dots,x_n\) of \(\R^n\) (Theorem~\ref{thm:r4-counting}(b)). Then \(P=[x_1\ \cdots\ x_n]\) has
independent columns, so it is invertible (Theorem~\ref{thm:r4-rank-rules}(b),(d)).
\item For \(j\leqslant g\), \(Ax_j=cx_j\); put \(m_{ij}=c\delta_{ij}\). For \(j>g\), write \(Ax_j=\sum_{i=1}^nm_{ij}x_i\), as
\(x_1,\dots,x_n\) span \(\R^n\). With \(M=(m_{ij})\), \(\col_j(AP)=Ax_j=\sum_im_{ij}x_i=P\col_j(M)\) for every \(j\)
(Theorem~\ref{r1-thm-views}(a),(b)); so \(AP=PM\) and \(M=P^{-1}AP\).
\item The first \(g\) columns of \(M\) are \(ce_1,\dots,ce_g\): \(M=cI_n\) if \(g=n\), and
\(M=\left(\begin{smallmatrix}cI_g&C\\0&D\end{smallmatrix}\right)\) with \(D\in\R^{(n-g)\times(n-g)}\) if \(g<n\). By
Theorem~\ref{r7-thm-similar}(a) and Proposition~\ref{r7-prop-block-triangular}(a),(b), the eigenvalues of \(A\) are those of
\(M\): \(c\) (\(g\) times), together with those of \(D\) if \(g<n\). Hence \(a_A(c)\geqslant g\).
\item If \(c\) is simple, then \(1\leqslant g_A(c)\leqslant a_A(c)=1\), and \(\rank(A-cI)=n-g_A(c)=n-1\).\qedhere
\end{itemize}
\end{proof}

\needspace{16\baselineskip}
\subsection*{Diagonalization}

\begin{definition}[\citeat{hefferon}{Five.II.2, Definition~2.1, Lemma~2.4}, \citeat{hoffman-kunze}{\S6.2}, \citeat{harville}{\S21.5, Theorems~21.5.1, 21.5.2}]\label{r7-def-diagonalizable}
A matrix \(A\in\R^{n\times n}\) is \textbf{diagonalizable} if \(P^{-1}AP\) is diagonal for some invertible
\(P\in\R^{n\times n}\), that is, if \(A\) is similar to a diagonal matrix; \(P\) is then said to \textbf{diagonalize}
\(A\). Let \(P=[p_1\ \cdots\ p_n]\in\R^{n\times n}\) be invertible and \(\Lambda=\diag(\lambda_1,\dots,\lambda_n)\) with
\(\lambda_j\in\R\). Then
\[
P^{-1}AP=\Lambda\iff AP=P\Lambda\iff A=P\Lambda P^{-1}\iff Ap_j=\lambda_jp_j\quad(1\leqslant j\leqslant n)
\]
Consequently \(A\) is diagonalizable if and only if \(\R^n\) has a basis consisting of eigenvectors of \(A\). If
\(P^{-1}AP=\Lambda\), then \(p_j\) is an eigenvector belonging to \(\lambda_j\), and \(\lambda_1,\dots,\lambda_n\) are the
eigenvalues of \(A\), listed with multiplicity; in particular all eigenvalues of a diagonalizable matrix are real.
\end{definition}
\begin{proof}
\begin{itemize}
\item If \(P^{-1}AP=\Lambda\), then \(AP=P(P^{-1}AP)=P\Lambda\) and \(A=P\Lambda P^{-1}\); if \(AP=P\Lambda\) or
\(A=P\Lambda P^{-1}\), then \(P^{-1}AP=\Lambda\).
\item \(\col_j(AP)=Ap_j\) and \(\col_j(P\Lambda)=P(\lambda_je_j)=\lambda_jp_j\) (Theorem~\ref{r1-thm-views}(a),(b)); two
matrices are equal if and only if their columns are.
\item If \(P^{-1}AP=\Lambda\), the columns of \(P\) form a basis of \(\R^n\) (Example~\ref{r3-exm-invertible-basis}), so they
are nonzero (Proposition~\ref{r3-prop-dependence}(a)) and, as \(Ap_j=\lambda_jp_j\), eigenvectors. Conversely, if
eigenvectors \(p_1,\dots,p_n\) with \(Ap_j=\lambda_jp_j\) form a basis of \(\R^n\), then \(P=[p_1\ \cdots\ p_n]\) is invertible
(Theorem~\ref{thm:r4-rank-rules}(b),(d)) and \(P^{-1}AP=\Lambda\).
\item By Theorem~\ref{r7-thm-similar}(a) and Proposition~\ref{r7-prop-block-triangular}(b),
\(p_A=p_\Lambda=(t-\lambda_1)\cdots(t-\lambda_n)\); so \(\lambda_1,\dots,\lambda_n\) are the eigenvalues of \(A\)
(Definition~\ref{r7-def-eigenvalues}).\qedhere
\end{itemize}
\end{proof}
If \(P^{-1}AP=\Lambda\) and \(c=P^{-1}x\) is the vector of coordinates of \(x\in\R^n\) in the basis \(p_1,\dots,p_n\)
(Example~\ref{r3-exm-invertible-basis}), then the coordinates of \(Ax\) are
\(P^{-1}Ax=\Lambda P^{-1}x=\Lambda c=(\lambda_1c_1,\dots,\lambda_nc_n)'\) (as \(P^{-1}A=\Lambda P^{-1}\); Proposition~\ref{r1-prop-diag}(a)):
in a basis of eigenvectors, \(A\) multiplies the \(j\)th coordinate by \(\lambda_j\)~\citeat{hoffman-kunze}{\S6.2}.

\begin{lemma}[\citeat{hoffman-kunze}{\S6.2}]\label{r7-lem-eigenspaces}
Let \(c_1,\dots,c_k\) (\(k\geqslant1\)) be distinct real eigenvalues of \(A\in\R^{n\times n}\), and
\(W_i=\mathcal N(A-c_iI)\).
\begin{enumerate}
\item If \(w_i\in W_i\) (\(1\leqslant i\leqslant k\)) and \(w_1+\dots+w_k=0\), then \(w_i=0\) for every \(i\).
\item If \(v_{i1},\dots,v_{id_i}\) is a basis of \(W_i\) for each \(i\), then the \(d_1+\dots+d_k\) vectors
\(v_{11},\dots,v_{1d_1},\dots,v_{k1},\dots,v_{kd_k}\) are linearly independent.
\end{enumerate}
\end{lemma}
\begin{proof}
(a)
\begin{itemize}
\item For a real polynomial \(f\) and \(w\in W_i\), \(f(A)w=f(c_i)w\) (Proposition~\ref{r7-prop-eig-poly}(a) if \(w\ne0\);
both sides are \(0\) if \(w=0\)). Applying \(f(A)\) to \(w_1+\dots+w_k=0\) gives \(f(c_1)w_1+\dots+f(c_k)w_k=0\)
(Theorem~\ref{r1-thm-rules}).
\item For a fixed \(i\), let \(f_i=\prod_{j\ne i}\frac{t-c_j}{c_i-c_j}\), the product over \(j\in\{1,\dots,k\}\), \(j\ne i\)
(\(f_1=1\) if \(k=1\)). By Proposition~\ref{r7-prop-poly}(b), \(f_i(c_l)=\prod_{j\ne i}\frac{c_l-c_j}{c_i-c_j}\), which is
\(1\) for \(l=i\) and \(0\) for \(l\ne i\) (the factor \(j=l\) vanishes).
\item With \(f=f_i\), the relation of the first item reads \(w_i=0\).
\end{itemize}
(b) A linear relation among these vectors has the form \(w_1+\dots+w_k=0\) with \(w_i=\sum_la_{il}v_{il}\in W_i\). By (a),
every \(w_i=0\), and then every \(a_{il}=0\), since \(v_{i1},\dots,v_{id_i}\) are independent.
\end{proof}

\needspace{15\baselineskip}
\begin{theorem}[\citeat{hoffman-kunze}{\S6.2, Theorem~2, Exercise~7}, \citeat{searle-khuri}{Theorem~7.1}, \citeat{lang}{Ch.~VIII, \S1, Remark~1}, \citeat{hefferon}{Five.II.3, Corollary~3.20}]\label{r7-thm-diag-criterion}
Let \(A\in\R^{n\times n}\), let \(c_1,\dots,c_k\) be the distinct real eigenvalues of \(A\) (\(k\geqslant0\); an empty sum
is \(0\), an empty product \(1\)), and let \(W_i=\mathcal N(A-c_iI)\). The following are equivalent:
\begin{enumerate}[label=(\roman*)]
\item \(A\) is diagonalizable;
\item \(p_A=(t-c_1)^{d_1}\cdots(t-c_k)^{d_k}\) for some positive integers \(d_1,\dots,d_k\), and \(\dim W_i=d_i\) for each
\(i\);
\item \(\dim W_1+\dots+\dim W_k=n\).
\end{enumerate}
In other words, \(A\) is diagonalizable if and only if all its eigenvalues are real and \(g_A(c)=a_A(c)\) for each of
them, that is, \(\rank(A-c_iI)=n-a_A(c_i)\) for each \(i\). In particular, an \(n\times n\) matrix with \(n\) distinct real
eigenvalues is diagonalizable.
\end{theorem}
\begin{proof}
\begin{itemize}
\item (i)\(\Rightarrow\)(ii): let \(P^{-1}AP=\Lambda=\diag(\lambda_1,\dots,\lambda_n)\). By
Definition~\ref{r7-def-diagonalizable}, \(\lambda_1,\dots,\lambda_n\) are the eigenvalues of \(A\); they are real, so each is
one of \(c_1,\dots,c_k\). Let \(d_i\) be the number of \(j\) with \(\lambda_j=c_i\); \(d_i\geqslant1\), as \(c_i\) is an
eigenvalue. Then \(p_A=(t-\lambda_1)\cdots(t-\lambda_n)=(t-c_1)^{d_1}\cdots(t-c_k)^{d_k}\).
\item Still (i)\(\Rightarrow\)(ii): \(\dim W_i=\dim\mathcal N(\Lambda-c_iI)\) (Theorem~\ref{r7-thm-similar}(b)). The \(j\)th entry
of \((\Lambda-c_iI)x\) is \((\lambda_j-c_i)x_j\), so \(\mathcal N(\Lambda-c_iI)\) consists of the \(x\) with \(x_j=0\) whenever
\(\lambda_j\ne c_i\); it has the basis \(\{e_j:\lambda_j=c_i\}\) of \(d_i\) vectors (Definition~\ref{r3-def-basis},
Proposition~\ref{r3-prop-dependence}(b)).
\item (ii)\(\Rightarrow\)(iii): \(\dim W_1+\dots+\dim W_k=d_1+\dots+d_k=\deg p_A=n\) (Proposition~\ref{r7-prop-poly}(a),
Definition~\ref{r7-def-charpoly}(b)).
\item (iii)\(\Rightarrow\)(i): as \(n\geqslant1\), \(k\geqslant1\). By Lemma~\ref{r7-lem-eigenspaces}(b), bases of
\(W_1,\dots,W_k\) together give \(n\) independent eigenvectors of \(A\); they form a basis of \(\R^n\)
(Theorem~\ref{thm:r4-counting}(d)), and \(A\) is diagonalizable (Definition~\ref{r7-def-diagonalizable}).
\item In (ii), \(d_i=a_A(c_i)\) (Theorem~\ref{r7-thm-factor}(b)), and \(p_A\) has no other roots (Theorem~\ref{r7-thm-factor}(c)); so (ii) says that all
eigenvalues are real and \(g_A(c_i)=a_A(c_i)\). Conversely, if all eigenvalues are real, then
\(p_A=(t-c_1)^{a_A(c_1)}\cdots(t-c_k)^{a_A(c_k)}\) (Definition~\ref{r7-def-eigenvalues}), and (ii) holds when
\(g_A(c_i)=a_A(c_i)\) for each \(i\). By Definition~\ref{r7-def-multiplicity}, \(g_A(c_i)=a_A(c_i)\) means
\(\rank(A-c_iI)=n-a_A(c_i)\).
\item If \(A\) has \(n\) distinct real eigenvalues, eigenvectors belonging to them form a basis of \(\R^n\)
(Theorem~\ref{r7-thm-indep}), so \(A\) is diagonalizable (Definition~\ref{r7-def-diagonalizable}).\qedhere
\end{itemize}
\end{proof}
For example, the matrix of Example~\ref{r7-exm-charpoly} is not diagonalizable: its eigenvalues \(0\), \(2\), \(2\) are
real, but \(g_A(2)=1<2=a_A(2)\). Nor is the matrix of Example~\ref{r7-exm-complex}, which has non-real eigenvalues.
The matrix \(A\) of Example~\ref{r7-exm-three} has the four distinct real eigenvalues \(-3\), \(-1\), \(1\), \(3\), so it is
diagonalizable. The eigenvectors found there are independent (Theorem~\ref{r7-thm-indep}), so
\[
P=\begin{pmatrix}-1&0&1&1\\1&2&1&1\\0&-1&0&1\\0&1&0&1\end{pmatrix}
\]
is invertible (Theorem~\ref{thm:r4-rank-rules}(b),(d)), and \(P^{-1}AP=\Lambda=\diag(-3,-1,1,3)\)
(Definition~\ref{r7-def-diagonalizable}).

To diagonalize \(A\in\R^{n\times n}\): find the eigenvalues, the roots of \(p_A\); if one of them is not real, \(A\) is not
diagonalizable (Definition~\ref{r7-def-diagonalizable}). For each real eigenvalue \(c\), the special solutions of \((A-cI)x=0\) form a basis of
\(\mathcal N(A-cI)\) (Theorem~\ref{thm:r4-four-bases}(c)), and \(g_A(c)\) is their number, the number of free unknowns.
If \(g_A(c)=a_A(c)\) for every eigenvalue \(c\), put these bases into the columns of \(P\) and the corresponding
eigenvalues, in the same order, on the diagonal of \(\Lambda\); then \(P^{-1}AP=\Lambda\)
(Lemma~\ref{r7-lem-eigenspaces}(b), Theorem~\ref{thm:r4-counting}(d), Definition~\ref{r7-def-diagonalizable}). By
Theorem~\ref{r7-thm-geo-alg}, only the eigenvalues with \(a_A(c)\geqslant2\) need to be checked.

\begin{example}[\citeat{hoffman-kunze}{\S6.2, Example~3}]\label{r7-exm-diagonalize}
Show that
\[
A=\begin{pmatrix}5&-6&-6\\-1&4&2\\3&-6&-4\end{pmatrix}
\]
is diagonalizable, and find an invertible \(P\) and a diagonal \(\Lambda\) with \(P^{-1}AP=\Lambda\).

\textit{Solution.}
\begin{itemize}
\item For \(t\in\R\), subtract column \(3\) from column \(2\) and take the factor \(t-2\) out of column \(2\)
(Definition~\ref{r5d-def-det}(b),(a) for columns, Theorem~\ref{r5d-thm-transpose}):
\[
\begin{aligned}
\det(tI-A)&=\det\begin{pmatrix}t-5&6&6\\1&t-4&-2\\-3&6&t+4\end{pmatrix}
=\det\begin{pmatrix}t-5&0&6\\1&t-2&-2\\-3&2-t&t+4\end{pmatrix}\\
&=(t-2)\det\begin{pmatrix}t-5&0&6\\1&1&-2\\-3&-1&t+4\end{pmatrix}
\end{aligned}
\]
Add row \(2\) to row \(3\) (Definition~\ref{r5d-def-det}(b)) and expand along column \(2\), whose only nonzero entry is
the \((2,2)\) entry \(1\) (Theorem~\ref{r5d-thm-expansion}):
\[
\det(tI-A)=(t-2)\det\begin{pmatrix}t-5&0&6\\1&1&-2\\-2&0&t+2\end{pmatrix}
=(t-2)\det\begin{pmatrix}t-5&6\\-2&t+2\end{pmatrix}=(t-2)(t^2-3t+2)
\]
As \(t^2-3t+2=(t-1)(t-2)\), \(p_A=(t-1)(t-2)^2\): the eigenvalues are \(1\), \(2\), \(2\), with \(a_A(1)=1\) and
\(a_A(2)=2\).
\item \(\lambda=1\):
\[
\begin{gathered}
A-I=\begin{pmatrix}4&-6&-6\\-1&3&2\\3&-6&-5\end{pmatrix}
\xrightarrow{R_1\leftrightarrow R_2}
\begin{pmatrix}-1&3&2\\4&-6&-6\\3&-6&-5\end{pmatrix}
\xrightarrow[R_3\leftarrow R_3+3R_1]{R_2\leftarrow R_2+4R_1}
\begin{pmatrix}-1&3&2\\0&6&2\\0&3&1\end{pmatrix}
\\
\xrightarrow{R_2\leftrightarrow R_3}
\begin{pmatrix}-1&3&2\\0&3&1\\0&6&2\end{pmatrix}
\xrightarrow{R_3\leftarrow R_3-2R_2}
\begin{pmatrix}-1&3&2\\0&3&1\\0&0&0\end{pmatrix}
\end{gathered}
\]
So \(\rank(A-I)=2\) (Definition~\ref{def:r4-rank}) and \(g_A(1)=1\) (Definition~\ref{r7-def-multiplicity}). The
equations \(-x_1+3x_2+2x_3=0\),
\(3x_2+x_3=0\) with \(x_3=3\) give \(x_2=-1\) and \(x_1=3\): \(\alpha_1=(3,-1,3)'\) is a basis of \(\mathcal N(A-I)\)
(Theorem~\ref{thm:r4-counting}(d)).
\item \(\lambda=2\):
\[
\begin{gathered}
A-2I=\begin{pmatrix}3&-6&-6\\-1&2&2\\3&-6&-6\end{pmatrix}
\xrightarrow{R_1\leftrightarrow R_2}
\begin{pmatrix}-1&2&2\\3&-6&-6\\3&-6&-6\end{pmatrix}
\xrightarrow[R_3\leftarrow R_3+3R_1]{R_2\leftarrow R_2+3R_1}
\begin{pmatrix}-1&2&2\\0&0&0\\0&0&0\end{pmatrix}
\\
\xrightarrow{R_1\leftarrow-R_1}
\begin{pmatrix}1&-2&-2\\0&0&0\\0&0&0\end{pmatrix}
\end{gathered}
\]
So \(g_A(2)=3-1=2\), and the special solutions \(\alpha_2=(2,1,0)'\), \(\alpha_3=(2,0,1)'\) form a basis of
\(\mathcal N(A-2I)\) (Theorem~\ref{thm:r4-four-bases}(c)).
\item \(g_A(1)+g_A(2)=3\), so \(A\) is diagonalizable (Theorem~\ref{r7-thm-diag-criterion}), and
\(\alpha_1,\alpha_2,\alpha_3\) form a basis of \(\R^3\) (Lemma~\ref{r7-lem-eigenspaces}(b),
Theorem~\ref{thm:r4-counting}(d)). With
\[
P=[\alpha_1\ \alpha_2\ \alpha_3]=\begin{pmatrix}3&2&2\\-1&1&0\\3&0&1\end{pmatrix},\qquad\Lambda=\diag(1,2,2),\qquad
AP=\begin{pmatrix}3&4&4\\-1&2&0\\3&0&2\end{pmatrix}=P\Lambda
\]
we have \(P^{-1}AP=\Lambda\) (Definition~\ref{r7-def-diagonalizable}).
\end{itemize}
\end{example}

\subsection*{Powers and rank of a diagonalizable matrix}

\begin{proposition}[\citeat{searle-khuri}{\S6.4.2}]\label{r7-prop-diag-powers}
Let \(A=P\Lambda P^{-1}\) with \(P=[p_1\ \cdots\ p_n]\in\R^{n\times n}\) invertible and
\(\Lambda=\diag(\lambda_1,\dots,\lambda_n)\), \(\lambda_i\in\R\).
\begin{enumerate}
\item \(A^k=P\Lambda^kP^{-1}=P\diag(\lambda_1^k,\dots,\lambda_n^k)P^{-1}\) for \(k=0,1,2,\dots\); hence, for \(x\in\R^n\)
and \(c=P^{-1}x\),
\[
A^kx=c_1\lambda_1^kp_1+c_2\lambda_2^kp_2+\dots+c_n\lambda_n^kp_n
\]
\item \(f(A)=Pf(\Lambda)P^{-1}=P\diag\bigl(f(\lambda_1),\dots,f(\lambda_n)\bigr)P^{-1}\) for every real polynomial \(f\).
\item \(A\) is invertible if and only if every \(\lambda_i\ne0\). Then \(A^{-1}=P\Lambda^{-1}P^{-1}\), and for \(k\geqslant1\) the
power \(A^{-k}=(A^{-1})^k\) equals \(P\diag(\lambda_1^{-k},\dots,\lambda_n^{-k})P^{-1}\).
\end{enumerate}
\end{proposition}
\begin{proof}
(a)
\begin{itemize}
\item By induction on \(k\): \(A^0=I=P\Lambda^0P^{-1}\), and if \(A^k=P\Lambda^kP^{-1}\), then
\(A^{k+1}=A^kA=P\Lambda^kP^{-1}P\Lambda P^{-1}=P\Lambda^{k+1}P^{-1}\) (Theorem~\ref{r1-thm-rules}).
\item \(\Lambda^k=\diag(\lambda_1^k,\dots,\lambda_n^k)\) (Proposition~\ref{r1-prop-powers}(d)).
\item \(A^kx=P\Lambda^kP^{-1}x=P\Lambda^kc\), and \(\Lambda^kc=(\lambda_1^kc_1,\dots,\lambda_n^kc_n)'\)
(Proposition~\ref{r1-prop-diag}(a)); so \(A^kx=\sum_j\lambda_j^kc_jp_j\) (Theorem~\ref{r1-thm-views}(b)).
\end{itemize}
(b) If \(f=\sum_kd_kt^k\), then by (a) and Theorem~\ref{r1-thm-rules}, \(f(A)=\sum_kd_kP\Lambda^kP^{-1}=Pf(\Lambda)P^{-1}\);
and \(f(\Lambda)=\sum_kd_k\diag(\lambda_1^k,\dots,\lambda_n^k)=\diag\bigl(f(\lambda_1),\dots,f(\lambda_n)\bigr)\), entry by
entry.

\needspace{3\baselineskip}
(c)
\begin{itemize}
\item The eigenvalues of \(A\) are \(\lambda_1,\dots,\lambda_n\) (Definition~\ref{r7-def-diagonalizable}), so \(A\) is
invertible if and only if every \(\lambda_i\ne0\) (Theorem~\ref{r7-thm-trace-det}).
\item Then \(\Lambda^{-1}=\diag(\lambda_1^{-1},\dots,\lambda_n^{-1})\) (Definition~\ref{r2-def-inverse}), and
\(A^{-1}=(P^{-1})^{-1}\Lambda^{-1}P^{-1}=P\Lambda^{-1}P^{-1}\) (Theorem~\ref{r2-thm-inverse-props}(a),(b)). Applying (a) to
\(A^{-1}=P\Lambda^{-1}P^{-1}\) gives \(A^{-k}=P(\Lambda^{-1})^kP^{-1}=P\diag(\lambda_1^{-k},\dots,\lambda_n^{-k})P^{-1}\).\qedhere
\end{itemize}
\end{proof}

\begin{example}\label{r7-exm-diag-powers}
Let
\[
B=\begin{pmatrix}1&2&-2&2\\-2&3&-2&0\\0&-2&3&-2\\2&-2&2&1\end{pmatrix}
\]
Show that \(3\) and \(1\) are eigenvalues of \(B\), diagonalize \(B\), and find \(B^k\) for \(k=0,1,2,\dots\)

\textit{Solution.}
\begin{itemize}
\item \(\lambda=3\):
\[
\begin{gathered}
B-3I=\begin{pmatrix}-2&2&-2&2\\-2&0&-2&0\\0&-2&0&-2\\2&-2&2&-2\end{pmatrix}
\xrightarrow[R_4\leftarrow R_4+R_1]{R_2\leftarrow R_2-R_1}
\begin{pmatrix}-2&2&-2&2\\0&-2&0&-2\\0&-2&0&-2\\0&0&0&0\end{pmatrix}
\xrightarrow{R_3\leftarrow R_3-R_2}
\begin{pmatrix}-2&2&-2&2\\0&-2&0&-2\\0&0&0&0\\0&0&0&0\end{pmatrix}
\\
\xrightarrow[R_2\leftarrow-\frac12R_2]{R_1\leftarrow-\frac12R_1}
\begin{pmatrix}1&-1&1&-1\\0&1&0&1\\0&0&0&0\\0&0&0&0\end{pmatrix}
\xrightarrow{R_1\leftarrow R_1+R_2}
\begin{pmatrix}1&0&1&0\\0&1&0&1\\0&0&0&0\\0&0&0&0\end{pmatrix}
\end{gathered}
\]
So \(\rank(B-3I)=2<4\) (Definition~\ref{def:r4-rank}): \(3\) is an eigenvalue (Theorem~\ref{r7-thm-eig-singular}), with
\(g_B(3)=4-2=2\) (Definition~\ref{r7-def-multiplicity}). The special solutions \(p_1=(-1,0,1,0)'\) and
\(p_2=(0,-1,0,1)'\) form a basis of \(\mathcal N(B-3I)\) (Theorem~\ref{thm:r4-four-bases}(c)).
\item \(\lambda=1\):
\[
\begin{gathered}
B-I=\begin{pmatrix}0&2&-2&2\\-2&2&-2&0\\0&-2&2&-2\\2&-2&2&0\end{pmatrix}
\xrightarrow{R_1\leftrightarrow R_2}
\begin{pmatrix}-2&2&-2&0\\0&2&-2&2\\0&-2&2&-2\\2&-2&2&0\end{pmatrix}
\xrightarrow[R_4\leftarrow R_4+R_1]{R_3\leftarrow R_3+R_2}
\begin{pmatrix}-2&2&-2&0\\0&2&-2&2\\0&0&0&0\\0&0&0&0\end{pmatrix}
\\
\xrightarrow[R_2\leftarrow\frac12R_2]{R_1\leftarrow-\frac12R_1}
\begin{pmatrix}1&-1&1&0\\0&1&-1&1\\0&0&0&0\\0&0&0&0\end{pmatrix}
\xrightarrow{R_1\leftarrow R_1+R_2}
\begin{pmatrix}1&0&0&1\\0&1&-1&1\\0&0&0&0\\0&0&0&0\end{pmatrix}
\end{gathered}
\]
So \(\rank(B-I)=2\): \(1\) is an eigenvalue with \(g_B(1)=2\), and the special solutions \(p_3=(0,1,1,0)'\) and
\(p_4=(-1,-1,0,1)'\) form a basis of \(\mathcal N(B-I)\).
\item By Lemma~\ref{r7-lem-eigenspaces}(b), \(p_1,\dots,p_4\) are four independent eigenvectors, a basis of \(\R^4\)
(Theorem~\ref{thm:r4-counting}(d)). So \(P^{-1}BP=\Lambda\) (Definition~\ref{r7-def-diagonalizable}) with
\[
P=[p_1\ p_2\ p_3\ p_4]=\begin{pmatrix}-1&0&0&-1\\0&-1&1&-1\\1&0&1&0\\0&1&0&1\end{pmatrix},\qquad
\Lambda=\diag(3,3,1,1),\qquad
P^{-1}=\begin{pmatrix}0&-1&1&-1\\1&-1&1&0\\0&1&0&1\\-1&1&-1&1\end{pmatrix}
\]
Check: \(P^{-1}P=I_4\), so this is the inverse (Theorem~\ref{r2-thm-invertible}); for instance, row \(1\) of \(P^{-1}\)
times \(p_4\) is \(0+1+0-1=0\). The eigenvalues of \(B\), listed with multiplicity, are \(3\), \(3\), \(1\), \(1\)
(Definition~\ref{r7-def-diagonalizable}), so \(a_B(3)=g_B(3)=2\) and \(a_B(1)=g_B(1)=2\).
\item By Proposition~\ref{r7-prop-diag-powers}(a), \(B^k=P\diag(3^k,3^k,1,1)P^{-1}\). As
\(\diag(3^k,3^k,1,1)=3^k\diag(1,1,0,0)+\bigl(I_4-\diag(1,1,0,0)\bigr)\) and \(PP^{-1}=I_4\),
\[
B^k=3^kE+(I_4-E),\qquad E=P\diag(1,1,0,0)P^{-1}
\]
For \(k=1\), \(B=3E+I_4-E=I_4+2E\), so \(E=\frac12(B-I_4)\) and
\[
B^k=I_4+\tfrac12(3^k-1)(B-I_4),\qquad
B^2=I_4+4(B-I_4)=\begin{pmatrix}1&8&-8&8\\-8&9&-8&0\\0&-8&9&-8\\8&-8&8&1\end{pmatrix}
\]
\end{itemize}
\end{example}

\needspace{14\baselineskip}
\begin{proposition}[\citeat{searle-khuri}{\S6.5.4}, \citeat{harville}{Theorems~21.4.3, 21.5.1}]\label{r7-prop-diag-rank}
Let \(A=P\Lambda P^{-1}\) with \(P=[p_1\ \cdots\ p_n]\in\R^{n\times n}\) invertible and
\(\Lambda=\diag(\lambda_1,\dots,\lambda_n)\), \(\lambda_i\in\R\).
\begin{enumerate}
\item \(\rank A\) is the number of indices \(i\) with \(\lambda_i\ne0\), and the columns \(p_i\) with \(\lambda_i\ne0\) form a
basis of \(\mathcal C(A)\). So \(\rank A\) is the number of nonzero eigenvalues of \(A\), listed with multiplicity.
\item For every real polynomial \(f\), \(f(A)\) is diagonalizable, and its eigenvalues, listed with multiplicity, are
\(f(\lambda_1),\dots,f(\lambda_n)\); hence \(\tr f(A)=f(\lambda_1)+\dots+f(\lambda_n)\) and
\(\det f(A)=f(\lambda_1)\cdots f(\lambda_n)\). In particular, \(A^k\) has the eigenvalues \(\lambda_1^k,\dots,\lambda_n^k\)
(\(k=0,1,2,\dots\)), and, if \(A\) is invertible, \(A^{-1}\) has the eigenvalues \(\lambda_1^{-1},\dots,\lambda_n^{-1}\).
\end{enumerate}
\end{proposition}
\begin{proof}
(a) Let \(S\) be the set of indices \(i\) with \(\lambda_i\ne0\), and \(r\) the number of its elements.
\begin{itemize}
\item \(\rank A=\rank(P\Lambda P^{-1})=\rank\Lambda\) (Theorem~\ref{thm:r4-rank-product}(d)).
\item Column \(i\) of \(\Lambda\) is \(\lambda_ie_i\) (Definition~\ref{r1-def-diag}), which is \(0\) for \(i\notin S\); so the \(r\)
columns \(\lambda_ie_i\), \(i\in S\), span \(\mathcal C(\Lambda)\) (Definition~\ref{def:r4-spaces}). They are independent: if
\(\sum_{i\in S}c_i\lambda_ie_i=0\), then every \(c_i\lambda_i=0\), as \(e_1,\dots,e_n\) are independent
(Definition~\ref{r3-def-basis}), and so \(c_i=0\). Hence \(\rank\Lambda=\dim\mathcal C(\Lambda)=r\)
(Definition~\ref{def:r4-rank}).
\item If \(r=0\), then \(\Lambda=0\), so \(A=0\) and \(\mathcal C(A)=\{0\}\), with the empty basis. Let \(r\geqslant1\), let \(X\in\R^{n\times r}\)
consist of the columns \(p_i\), \(i\in S\), and let \(\Delta\) be the \(r\times r\) diagonal matrix of the corresponding
\(\lambda_i\). As \(Ap_i=\lambda_ip_i\) (Definition~\ref{r7-def-diagonalizable}), \(AX=X\Delta\)
(Theorem~\ref{r1-thm-views}(a),(b)). \(\Delta\) is invertible (Definition~\ref{r2-def-inverse}), so
\(X=AX\Delta^{-1}=A(X\Delta^{-1})\), and \(\mathcal C(X)\subseteq\mathcal C(A)\) (Theorem~\ref{thm:r4-rank-product}(a)).
\item The columns of \(X\) are columns of the invertible \(P\), hence independent (Example~\ref{r3-exm-invertible-basis},
Proposition~\ref{r3-prop-dependence}(b)). So \(\dim\mathcal C(X)=r=\rank A=\dim\mathcal C(A)\), and
\(\mathcal C(X)=\mathcal C(A)\) (Theorem~\ref{thm:r4-counting}(a)): the columns of \(X\) form a basis of \(\mathcal C(A)\).
\item The eigenvalues of \(A\), listed with multiplicity, are \(\lambda_1,\dots,\lambda_n\)
(Definition~\ref{r7-def-diagonalizable}).
\end{itemize}
(b)
\begin{itemize}
\item By Proposition~\ref{r7-prop-diag-powers}(b), \(f(A)=Pf(\Lambda)P^{-1}\), that is, \(P^{-1}f(A)P=f(\Lambda)\), with
\(f(\Lambda)=\diag\bigl(f(\lambda_1),\dots,f(\lambda_n)\bigr)\). So \(f(A)\) is diagonalizable, and its eigenvalues, listed
with multiplicity, are \(f(\lambda_1),\dots,f(\lambda_n)\) (Definition~\ref{r7-def-diagonalizable}).
\item Theorem~\ref{r7-thm-trace-det} gives the trace and the determinant; \(f=t^k\) gives the statement on \(A^k\).
\item If \(A\) is invertible, then \(P^{-1}A^{-1}P=\Lambda^{-1}=\diag(\lambda_1^{-1},\dots,\lambda_n^{-1})\)
(Proposition~\ref{r7-prop-diag-powers}(c)), and Definition~\ref{r7-def-diagonalizable} gives the eigenvalues of
\(A^{-1}\).\qedhere
\end{itemize}
\end{proof}

\needspace{14\baselineskip}
\subsection*{Orthogonal triangularization}

\begin{theorem}[Schur triangularization; \citeat{harville}{Theorem~21.5.11}]\label{r7-thm-schur}
Let \(A\in\R^{n\times n}\), and let all eigenvalues of \(A\) be real, listed with multiplicity in any order as
\(\lambda_1,\dots,\lambda_n\). Then \(A=QTQ'\), where \(Q\in\R^{n\times n}\) is orthogonal and \(T\in\R^{n\times n}\) is
upper triangular with the diagonal entries \(\lambda_1,\dots,\lambda_n\), in this order.
\end{theorem}
\begin{proof}
Induction on \(n\); for \(n=1\), \(A=(\lambda_1)=I_1(\lambda_1)I_1'\). Let \(n\geqslant2\).
\begin{itemize}
\item \(\lambda_1\) is a real eigenvalue, so \(Ax=\lambda_1x\) for some \(x\ne0\) (Definition~\ref{r7-def-eigenvalues}).
Then \(u=x/\lVert x\rVert\) has \(u'u=1\) and \(Au=\lambda_1u\). By Corollary~\ref{r6-cor-onb}(c) (with \(p=1\)), there is
\(V\in\R^{n\times(n-1)}\) such that \(Q_1=[u\ V]\) is orthogonal; comparing blocks in \(Q_1'Q_1=I_n\)
(Theorem~\ref{r5b-thm-block}(b),(c)) gives \(V'u=0\).
\item By Theorem~\ref{r5b-thm-block}(b),(c) and \(Au=\lambda_1u\),
\[
Q_1'AQ_1=\begin{pmatrix}u'Au&u'AV\\V'Au&V'AV\end{pmatrix}=\begin{pmatrix}\lambda_1&b'\\0&A_1\end{pmatrix},\qquad
b'=u'AV,\qquad A_1=V'AV\in\R^{(n-1)\times(n-1)}
\]
\item \(Q_1'=Q_1^{-1}\) (Proposition~\ref{r6-prop-orthonormal}(b)), so \(A\) and \(Q_1'AQ_1\) have the same eigenvalues
(Theorem~\ref{r7-thm-similar}(a)); by Proposition~\ref{r7-prop-block-triangular}(a) these are \(\lambda_1\) together with
the eigenvalues of \(A_1\). By Theorem~\ref{r7-thm-factor}(b), the eigenvalues of \(A_1\), listed with multiplicity, are
\(\lambda_2,\dots,\lambda_n\), all real.
\item By the induction hypothesis, \(A_1=Q_2T_2Q_2'\) with \(Q_2\in\R^{(n-1)\times(n-1)}\) orthogonal and \(T_2\) upper
triangular with the diagonal entries \(\lambda_2,\dots,\lambda_n\); so \(Q_2'A_1Q_2=T_2\). Let \(S=\diag(1,Q_2)\). Then
\(S'S=\diag(1,Q_2'Q_2)=I_n\) (Theorem~\ref{r5b-thm-block}(b), Proposition~\ref{r5b-prop-block-tri}(b)), so \(S\) is
orthogonal, and so is \(Q=Q_1S\) (Theorem~\ref{r6-thm-orthogonal}(b)). As \(Q'=S'Q_1'\)
(Theorem~\ref{r1-thm-transpose}), Theorem~\ref{r5b-thm-block}(c) gives
\[
Q'AQ=S'(Q_1'AQ_1)S=\begin{pmatrix}1&0\\0&Q_2'\end{pmatrix}\begin{pmatrix}\lambda_1&b'\\0&A_1\end{pmatrix}
\begin{pmatrix}1&0\\0&Q_2\end{pmatrix}=\begin{pmatrix}\lambda_1&b'Q_2\\0&T_2\end{pmatrix}
\]
an upper triangular matrix \(T\) with the diagonal entries \(\lambda_1,\dots,\lambda_n\). As \(QQ'=I_n\),
\(A=Q(Q'AQ)Q'=QTQ'\).\qedhere
\end{itemize}
\end{proof}
Conversely, if \(A=QTQ'\) with \(Q\) orthogonal and \(T\) upper triangular, then \(A\) is similar to \(T\), as
\(Q'=Q^{-1}\), and its eigenvalues are the diagonal entries of \(T\) (Theorem~\ref{r7-thm-similar}(a),
Proposition~\ref{r7-prop-block-triangular}(b)), which are real. So such a factorization exists exactly when all
eigenvalues of \(A\) are real; the matrix of Example~\ref{r7-exm-complex} has none.

\begin{example}\label{r7-exm-schur}
Let
\[
A=\begin{pmatrix}1&2&1&0\\1&2&0&1\\1&0&1&2\\0&1&1&2\end{pmatrix},\qquad
G=\begin{pmatrix}1&1&1&1\\1&1&-1&-1\\1&-1&1&-1\\1&-1&-1&1\end{pmatrix},\qquad Q=\tfrac12G
\]
Every row of \(A\) adds up to \(4\), so \(A\one_4=4\one_4\), and \(u=\frac12\one_4\), the first column of \(Q\), has
\(u'u=1\) and \(Au=4u\). The columns of \(G\) are orthogonal and have squared length \(4\) (for instance, columns \(2\)
and \(3\) give \(1-1-1+1=0\)), so \(G'G=4I_4\) and \(Q'Q=I_4\): \(Q=[u\ V]\) is an orthogonal matrix as in the first
step of the proof of Theorem~\ref{r7-thm-schur}. As \(G'=G\),
\[
AG=\begin{pmatrix}4&2&0&-2\\4&2&-2&0\\4&-2&0&2\\4&-2&-2&0\end{pmatrix},\qquad
Q'AQ=\tfrac14G(AG)=\begin{pmatrix}4&0&-1&0\\0&2&0&-1\\0&0&1&0\\0&0&0&-1\end{pmatrix}=T
\]
In the notation of that proof, \(\lambda_1=4\), \(b'=(0,-1,0)\), and \(A_1=V'AV\) is already upper triangular, so the
induction stops after one step. Hence \(A=QTQ'\), and the eigenvalues of \(A\) are \(4\), \(2\), \(1\), \(-1\)
(Theorem~\ref{r7-thm-similar}(a), Proposition~\ref{r7-prop-block-triangular}(b)); \(\tr T=6=\tr A\) and
\(\det T=-8=\det A\). As these four eigenvalues are real and distinct, \(A\) is diagonalizable
(Theorem~\ref{r7-thm-diag-criterion}); but no orthogonal \(Q_0\) makes \(Q_0'AQ_0\) diagonal: if \(Q_0'AQ_0=\Lambda\),
then \(A=Q_0\Lambda Q_0'\) (as \(Q_0Q_0'=I\), Proposition~\ref{r6-prop-orthonormal}(b)) and \((Q_0\Lambda Q_0')'=Q_0\Lambda'Q_0'=Q_0\Lambda Q_0'\) (Theorem~\ref{r1-thm-transpose}), while
\(A\) is not symmetric (\(a_{12}=2\ne1=a_{21}\)).
\end{example}

\begin{corollary}[\citeat{lang}{Ch.~X, \S1, Exercise~1}, \citeat{shakarchi}{}, \citeat{abadir-magnus}{Exercise~7.65}]\label{r7-cor-schur}
Let \(A\in\R^{n\times n}\), and let all eigenvalues \(\lambda_1,\dots,\lambda_n\) of \(A\), listed with multiplicity, be
real.
\begin{enumerate}
\item For every real polynomial \(f\), the eigenvalues of \(f(A)\), listed with multiplicity, are
\(f(\lambda_1),\dots,f(\lambda_n)\). In particular \(\tr A^k=\lambda_1^k+\dots+\lambda_n^k\) for \(k=0,1,2,\dots\)
\item If \(A\) is invertible, the eigenvalues of \(A^{-1}\), listed with multiplicity, are
\(\lambda_1^{-1},\dots,\lambda_n^{-1}\).
\end{enumerate}
\end{corollary}
\begin{proof}
Let \(A=QTQ'\) as in Theorem~\ref{r7-thm-schur}: \(Q'=Q^{-1}\) (Proposition~\ref{r6-prop-orthonormal}(b)), and \(T\) is
upper triangular with the diagonal entries \(\lambda_1,\dots,\lambda_n\).

(a)
\begin{itemize}
\item \(A^k=QT^kQ'\) for \(k=0,1,2,\dots\), by induction: \(A^0=I=QQ'\), and \(A^{k+1}=A^kA=QT^kQ'QTQ'=QT^{k+1}Q'\)
(Theorem~\ref{r1-thm-rules}). So, for \(f=\sum_kc_kt^k\), \(f(A)=\sum_kc_kQT^kQ'=Qf(T)Q'\).
\item \(T^k\) is upper triangular with the diagonal entries \(\lambda_1^k,\dots,\lambda_n^k\)
(Proposition~\ref{r1-prop-powers}(e)); adding multiples of these matrices entry by entry, \(f(T)\) is upper triangular
with the diagonal entries \(f(\lambda_1),\dots,f(\lambda_n)\).
\item \(f(A)\) is similar to \(f(T)\), so by Theorem~\ref{r7-thm-similar}(a) and
Proposition~\ref{r7-prop-block-triangular}(b) its eigenvalues are \(f(\lambda_1),\dots,f(\lambda_n)\). For \(f=t^k\),
Theorem~\ref{r7-thm-trace-det} gives \(\tr A^k=\lambda_1^k+\dots+\lambda_n^k\).
\end{itemize}
(b) \(0\) is not an eigenvalue of \(A\) (Theorem~\ref{r7-thm-trace-det}), so every \(\lambda_i\ne0\); hence \(T\) is
invertible, and \(T^{-1}\) is upper triangular with the diagonal entries \(\lambda_1^{-1},\dots,\lambda_n^{-1}\)
(Theorem~\ref{r2-thm-triangular}). Then \(A^{-1}=(Q')^{-1}T^{-1}Q^{-1}=QT^{-1}Q'\)
(Theorem~\ref{r2-thm-inverse-props}(a),(b)), so \(A^{-1}\) is similar to \(T^{-1}\); conclude as in (a).
\end{proof}
Proposition~\ref{r7-prop-diag-rank}(b) gives (a) for a diagonalizable \(A\); here \(A\) need not be diagonalizable.
For the matrix \(A\) of Example~\ref{r7-exm-three}, with the eigenvalues \(-3\), \(-1\), \(1\), \(3\),
\[
A^2=\begin{pmatrix}5&-4&0&8\\-4&5&8&0\\0&0&5&4\\0&0&4&5\end{pmatrix},\qquad\tr A^2=20=(-3)^2+(-1)^2+1^2+3^2
\]
and, as \(\det A=9\ne0\) (Example~\ref{r7-exm-three}), \(A^{-1}\) has the eigenvalues \(-\frac13\), \(-1\), \(1\), \(\frac13\) (part (b)).

\needspace{14\baselineskip}
\subsection*{The products \(AB\) and \(BA\)}

\begin{theorem}[\citeat{harville}{Theorem~21.10.1}]\label{r7-thm-AB-BA}
Let \(A\in\R^{m\times n}\) and \(B\in\R^{n\times m}\) with \(m\leqslant n\).
\begin{enumerate}
\item \(p_{BA}(t)=t^{n-m}p_{AB}(t)\); in particular \(p_{BA}=p_{AB}\) if \(m=n\).
\item The eigenvalues of \(BA\), listed with multiplicity, are those of \(AB\) together with \(n-m\) zeros. So \(AB\) and
\(BA\) have the same nonzero eigenvalues, with the same algebraic multiplicities.
\item \(g_{AB}(c)=g_{BA}(c)\) for every real \(c\ne0\). More precisely, if \(c\ne0\) is a real eigenvalue of \(AB\) and
the columns of \(X\in\R^{m\times g}\) form a basis of \(\mathcal N(AB-cI_m)\), then the \(g\) columns of \(BX\) are linearly
independent eigenvectors of \(BA\) belonging to \(c\).
\end{enumerate}
\end{theorem}
\begin{proof}
(a)
\begin{itemize}
\item Let \(t\in\R\), \(t\ne0\). Theorem~\ref{r5d-thm-woodbury}(a), with \(I_n\), \(I_m\), \(-t^{-1}B\) and \(A\) in
place of its \(A\), \(T\), \(U\) and \(V\), gives
\[
\det(I_n-t^{-1}BA)=\det I_n\cdot\det I_m\cdot\det(I_m-t^{-1}AB)=\det(I_m-t^{-1}AB)
\]
\item By Theorem~\ref{r5d-thm-product}(c), \(\det(tI_n-BA)=t^n\det(I_n-t^{-1}BA)\) and
\(\det(tI_m-AB)=t^m\det(I_m-t^{-1}AB)\). Hence \(\det(tI_n-BA)=t^{n-m}\det(tI_m-AB)\).
\item So \(p_{BA}\) and \(t^{n-m}p_{AB}\) have the same value at every real \(t\ne0\) (Definition~\ref{r7-def-charpoly}(a),
Proposition~\ref{r7-prop-poly}(b)), and they are equal (Proposition~\ref{r7-prop-poly}(d)).
\end{itemize}
(b) If \(p_{AB}=(t-\mu_1)\cdots(t-\mu_m)\) (Definition~\ref{r7-def-eigenvalues}), then by (a)
\(p_{BA}=(t-0)^{n-m}(t-\mu_1)\cdots(t-\mu_m)\), a product of \(n\) linear factors. By Theorem~\ref{r7-thm-factor}(b), the
eigenvalues of \(BA\) are \(\mu_1,\dots,\mu_m\) and \(n-m\) zeros; so \(a_{BA}(c)=a_{AB}(c)\) for every \(c\ne0\).

(c)
\begin{itemize}
\item If \(c\) is not an eigenvalue of \(AB\), then by (b) it is not one of \(BA\), and both geometric multiplicities are
\(0\) (Definition~\ref{r7-def-multiplicity}).
\item Let \(c\ne0\) be a real eigenvalue of \(AB\), \(g=g_{AB}(c)\), and let \(X\in\R^{m\times g}\) have a basis of
\(\mathcal N(AB-cI_m)\) as its columns. Then \(\rank X=g\) (Theorem~\ref{thm:r4-rank-rules}(b)) and \(ABX=cX\)
(Theorem~\ref{r1-thm-views}(a)).
\item By Theorem~\ref{thm:r4-rank-product}(a),(d), \(g=\rank X\geqslant\rank(BX)\geqslant\rank(ABX)=\rank(cX)=g\), as
\(cX=X(cI_g)\) with \(cI_g\) invertible. So the \(g\) columns of \(BX\) are independent
(Theorem~\ref{thm:r4-rank-rules}(b)).
\item \(BA(BX)=B(ABX)=cBX\): the columns of \(BX\) lie in \(\mathcal N(BA-cI_n)\). Being independent, they are
nonzero (Proposition~\ref{r3-prop-dependence}(a)), so they are eigenvectors of \(BA\) belonging to \(c\), and
\(g_{BA}(c)\geqslant g\) (Theorem~\ref{thm:r4-counting}(c)).
\item The argument does not use \(m\leqslant n\); with the roles of \(A\) and \(B\) exchanged it gives
\(g_{AB}(c)\geqslant g_{BA}(c)\).\qedhere
\end{itemize}
\end{proof}

\begin{example}\label{r7-exm-ones}
For \(n\geqslant2\), take \(A=\one_n'\in\R^{1\times n}\) and \(B=\one_n\in\R^{n\times1}\) in Theorem~\ref{r7-thm-AB-BA}
(\(m=1\)): \(AB=\one'\one=n\) and \(BA=\one\one'=J_n\) (Definition~\ref{r1-def-ones}, Proposition~\ref{r1-prop-centering}).
So \(p_{J_n}(t)=t^{n-1}(t-n)\): the eigenvalues of \(J_n\) are \(n\) (simple) and \(0\) (\(n-1\) times). Indeed
\(J_n\one=\one(\one'\one)=n\one\), and \(J_nx=\one(\one'x)=0\) for every centered \(x\), that is, \(\one'x=0\)
(Example~\ref{r3-exm-subspaces}(f)).
\end{example}

\begin{example}\label{r7-exm-ata}
Let
\[
A=\begin{pmatrix}1&1&0\\1&0&1\\0&1&1\\-1&0&1\end{pmatrix},\qquad
A'A=\begin{pmatrix}3&1&0\\1&2&1\\0&1&3\end{pmatrix},\qquad
AA'=\begin{pmatrix}2&1&1&-1\\1&2&1&0\\1&1&2&1\\-1&0&1&2\end{pmatrix}
\]
Find the eigenvalues of \(A'A\) and \(AA'\), and an eigenvector for each.

\textit{Solution.}
\begin{itemize}
\item For \(t\in\R\), expanding along row \(1\) (Theorem~\ref{r5d-thm-expansion}),
\[
\det(tI_3-A'A)=\det\begin{pmatrix}t-3&-1&0\\-1&t-2&-1\\0&-1&t-3\end{pmatrix}=(t-3)\bigl((t-2)(t-3)-1\bigr)-(t-3)
=(t-3)(t^2-5t+4)
\]
So \(p_{A'A}=(t-1)(t-3)(t-4)\) (Definition~\ref{r7-def-charpoly}(a), Proposition~\ref{r7-prop-poly}(d)).
Theorem~\ref{r7-thm-AB-BA}(a), with \(A'\in\R^{3\times4}\) and \(A\in\R^{4\times3}\) in place of \(A\) and \(B\) (\(m=3\),
\(n=4\)), gives \(p_{AA'}=t^{4-3}p_{A'A}=t(t-1)(t-3)(t-4)\).
\item Multiplying out, \(A'A(1,1,1)'=(4,4,4)'\), \(A'A(-1,0,1)'=(-3,0,3)'\) and \(A'A(1,-2,1)'=(1,-2,1)'\): these are
eigenvectors of \(A'A\) belonging to \(4\), \(3\) and \(1\). Each of these eigenvalues is simple, so its eigenspace has
dimension \(1\) (Theorem~\ref{r7-thm-geo-alg}), and the vector found is a basis of it (Theorem~\ref{thm:r4-counting}(d)).
By Theorem~\ref{r7-thm-AB-BA}(c), the vectors
\[
A\begin{pmatrix}1\\1\\1\end{pmatrix}=\begin{pmatrix}2\\2\\2\\0\end{pmatrix},\qquad
A\begin{pmatrix}-1\\0\\1\end{pmatrix}=\begin{pmatrix}-1\\0\\1\\2\end{pmatrix},\qquad
A\begin{pmatrix}1\\-2\\1\end{pmatrix}=\begin{pmatrix}-1\\2\\-1\\0\end{pmatrix}
\]
are eigenvectors of \(AA'\) belonging to \(4\), \(3\) and \(1\); for instance \(AA'(-1,0,1,2)'=(-3,0,3,6)'\).
\item The eigenvalue \(0\) of \(AA'\) has the eigenspace \(\mathcal N(AA')=\mathcal N(A')\)
(Theorem~\ref{thm:r4-rank-ata}(a), applied to \(A'\)). Row-reduce \(A'\):
\[
A'=\begin{pmatrix}1&1&0&-1\\1&0&1&0\\0&1&1&1\end{pmatrix}
\xrightarrow{R_2\leftarrow R_2-R_1}
\begin{pmatrix}1&1&0&-1\\0&-1&1&1\\0&1&1&1\end{pmatrix}
\xrightarrow{R_3\leftarrow R_3+R_2}
\begin{pmatrix}1&1&0&-1\\0&-1&1&1\\0&0&2&2\end{pmatrix}
\]
So \(\rank A'=3\) and \(\dim\mathcal N(A')=1\) (Definition~\ref{def:r4-rank}, Theorem~\ref{thm:r4-rank-nullity}). With
\(x_4=1\), back substitution gives \(x_3=-1\), \(x_2=0\) and \(x_1=1\); so \((1,0,-1,1)'\) is a basis of
\(\mathcal N(AA')\) (Theorem~\ref{thm:r4-counting}(d)).
\end{itemize}
\end{example}

\needspace{6\baselineskip}
\begin{corollary}[\citeat{searle-khuri}{Lemma~7.1, \S6.5.4}]\label{r7-cor-rank-eigen}
Let \(A\in\R^{n\times n}\) have rank \(r\). Then \(a_A(0)\geqslant n-r\): at most \(r\) of the eigenvalues of \(A\), listed
with multiplicity, are nonzero. If \(A\) is diagonalizable, exactly \(r\) of them are nonzero.
\end{corollary}
\begin{proof}
\begin{itemize}
\item If \(r=0\), then \(A=0\) and \(p_A=t^n\) (Proposition~\ref{r7-prop-block-triangular}(b)), so \(a_A(0)=n\).
\item Let \(r\geqslant1\), and let \(A=BC\) be a full-rank factorization, with \(B\in\R^{n\times r}\) and \(C\in\R^{r\times n}\)
(Theorem~\ref{thm:r4-full-rank-factorization}(a)); here \(r\leqslant n\) (Theorem~\ref{thm:r4-rank-rules}(a)).
Theorem~\ref{r7-thm-AB-BA}(a), applied to \(C\in\R^{r\times n}\) and \(B\in\R^{n\times r}\) in place of \(A\) and \(B\), gives
\(p_A=p_{BC}=t^{n-r}p_{CB}\).
\item By Theorem~\ref{r7-thm-AB-BA}(b), the eigenvalues of \(A\), listed with multiplicity, are the \(r\) eigenvalues of
the \(r\times r\) matrix \(CB\) together with \(n-r\) zeros. So at most \(r\) of them are nonzero, and \(a_A(0)\geqslant n-r\).
\item If \(A\) is diagonalizable, exactly \(r\) of them are nonzero, by Proposition~\ref{r7-prop-diag-rank}(a).\qedhere
\end{itemize}
\end{proof}

\needspace{14\baselineskip}
\begin{example}[\citeat{searle-khuri}{Example~7.1}]\label{r7-exm-rank-eigen}
Let
\[
A=\begin{pmatrix}-1&1&2\\2&2&0\\0&4&4\end{pmatrix}
\]
Find \(\rank A\), a full-rank factorization \(A=BC\), the eigenvalues of \(A\), and compare \(\rank A\) with the
number of nonzero eigenvalues.

\textit{Solution.}
\begin{itemize}
\item Row-reduce \(A\) to its RREF \(R\):
\[
\begin{gathered}
\begin{pmatrix}-1&1&2\\2&2&0\\0&4&4\end{pmatrix}
\xrightarrow{R_2\leftarrow R_2+2R_1}
\begin{pmatrix}-1&1&2\\0&4&4\\0&4&4\end{pmatrix}
\xrightarrow{R_3\leftarrow R_3-R_2}
\begin{pmatrix}-1&1&2\\0&4&4\\0&0&0\end{pmatrix}
\\
\xrightarrow[R_2\leftarrow\frac14R_2]{R_1\leftarrow-R_1}
\begin{pmatrix}1&-1&-2\\0&1&1\\0&0&0\end{pmatrix}
\xrightarrow{R_1\leftarrow R_1+R_2}
\begin{pmatrix}1&0&-1\\0&1&1\\0&0&0\end{pmatrix}
\end{gathered}
\]
The pivot columns are columns \(1\) and \(2\), so \(\rank A=2\) (Definition~\ref{def:r4-rank}). By
Theorem~\ref{thm:r4-full-rank-factorization}(a), \(A=BC\) with
\[
B=\begin{pmatrix}-1&1\\2&2\\0&4\end{pmatrix},\qquad C=\begin{pmatrix}1&0&-1\\0&1&1\end{pmatrix},\qquad
CB=\begin{pmatrix}-1&-3\\2&6\end{pmatrix}
\]
\item By Definition~\ref{r7-def-charpoly}(b), \(p_{CB}=t^2-5t+(-6+6)=t(t-5)\), and, as in the proof of
Corollary~\ref{r7-cor-rank-eigen}, \(p_A=t^{3-2}p_{CB}=t^2(t-5)\). So the eigenvalues of \(A\) are \(0\), \(0\), \(5\):
\(a_A(0)=2\), and only one eigenvalue is nonzero, although \(\rank A=2\).
\item Hence \(A\) is not diagonalizable (Proposition~\ref{r7-prop-diag-rank}(a)). Indeed
\(g_A(0)=3-\rank A=1<2=a_A(0)\) (Definition~\ref{r7-def-multiplicity}, Theorem~\ref{r7-thm-diag-criterion}).
\item Check: \(\tr A=-1+2+4=5=0+0+5\) and \(\det A=0=0\cdot0\cdot5\) (Theorem~\ref{r7-thm-trace-det}), as \(\rank A<3\)
(Theorem~\ref{r5d-thm-product}(a)).
\end{itemize}
\end{example}

\subsection*{The Cayley--Hamilton theorem}

\begin{theorem}[Cayley--Hamilton; \citeat{searle-khuri}{Theorem~7.7}, \citeat{harville-solutions}{Ch.~21, Exercise~3}]\label{r7-thm-cayley-hamilton}
Let \(A\in\R^{n\times n}\) and \(p_A(t)=c_0+c_1t+\dots+c_nt^n\) (so \(c_n=1\)). Then
\[
p_A(A)=c_0I+c_1A+\dots+c_nA^n=0
\]
\end{theorem}
\begin{proof}
For \(n=1\), \(A=(a)\), \(p_A=t-a\) and \(p_A(A)=A-aI_1=0\). Let \(n\geqslant2\).
\begin{itemize}
\item For \(t\in\R\), the \((i,j)\) entry of \(\operatorname{adj}(tI-A)\) is the cofactor
\((-1)^{i+j}\det\bigl((tI-A)(j|i)\bigr)\) (Definition~\ref{r5d-def-cofactor}, Theorem~\ref{r5d-thm-adjugate}). Every entry
of the \((n-1)\times(n-1)\) matrix \((tI-A)(j|i)\) is an entry \(\delta_{rs}t-a_{rs}\) of \(tI-A\); by
Lemma~\ref{r7-lem-degree} (with \(m=d=n-1\)) and Proposition~\ref{r7-prop-poly}(b), the cofactor is the value at \(t\)
of a real polynomial \(h_{ij}\) that is \(0\) or of degree at most \(n-1\). Writing
\(h_{ij}=\sum_{s=0}^{n-1}k^{(s)}_{ij}t^s\) and \(K_s=(k^{(s)}_{ij})\in\R^{n\times n}\),
\[
\operatorname{adj}(tI-A)=K_0+tK_1+\dots+t^{n-1}K_{n-1}\qquad(t\in\R)
\]
\item By Theorem~\ref{r5d-thm-adjugate}, which holds for every matrix, and Definition~\ref{r7-def-charpoly}(a),
\((tI-A)\operatorname{adj}(tI-A)=p_A(t)I\) for every real \(t\). By Theorem~\ref{r1-thm-rules} the left side is
\[
(tI-A)\sum_{s=0}^{n-1}t^sK_s=T_0+tT_1+\dots+t^nT_n
\]
with \(T_0=-AK_0\), \(T_s=K_{s-1}-AK_s\) (\(1\leqslant s\leqslant n-1\)) and \(T_n=K_{n-1}\). So for each \((i,j)\) the polynomials \(\sum_s(T_s)_{ij}t^s\) and \(\delta_{ij}p_A\) have the same value at every real
\(t\); by Proposition~\ref{r7-prop-poly}(d), their coefficients agree: \(T_s=c_sI\) for \(0\leqslant s\leqslant n\).
\item Hence
\[
\begin{aligned}
p_A(A)&=\sum_{s=0}^nA^s(c_sI)=\sum_{s=0}^nA^sT_s=-AK_0+\sum_{s=1}^{n-1}\bigl(A^sK_{s-1}-A^{s+1}K_s\bigr)+A^nK_{n-1}\\
&=\sum_{s=1}^nA^sK_{s-1}-\sum_{s=0}^{n-1}A^{s+1}K_s=0
\end{aligned}
\]
since the two sums are the same (replace \(s\) by \(s-1\) in the second).\qedhere
\end{itemize}
\end{proof}

\needspace{13\baselineskip}
\begin{corollary}[\citeat{searle-khuri}{\S7.3, Example~7.2}, \citeat{harville-solutions}{Ch.~21, Exercise~4}]\label{r7-cor-cayley-hamilton}
Let \(A\in\R^{n\times n}\) and \(p_A(t)=c_0+c_1t+\dots+c_nt^n\) (so \(c_n=1\)).
\begin{enumerate}
\item \(A^n=-(c_0I+c_1A+\dots+c_{n-1}A^{n-1})\), and every power \(A^k\) (\(k=0,1,2,\dots\)) is a linear combination of
\(I,A,\dots,A^{n-1}\).
\item If \(A\) is invertible, then \(c_0\ne0\) and
\[
A^{-1}=-\frac1{c_0}\bigl(c_1I+c_2A+\dots+c_nA^{n-1}\bigr)
\]
\end{enumerate}
\end{corollary}
\begin{proof}
(a)
\begin{itemize}
\item By Theorem~\ref{r7-thm-cayley-hamilton}, \(c_0I+c_1A+\dots+c_{n-1}A^{n-1}+A^n=0\), which is the formula for \(A^n\).
\item For \(k\leqslant n-1\), \(A^k\) is one of \(I,A,\dots,A^{n-1}\). Let \(k\geqslant n-1\) and
\(A^k=b_0I+b_1A+\dots+b_{n-1}A^{n-1}\). Then \(A^{k+1}=A^kA=b_0A+\dots+b_{n-2}A^{n-1}+b_{n-1}A^n\)
(Theorem~\ref{r1-thm-rules}), and replacing \(A^n\) by the first item gives a linear combination of
\(I,A,\dots,A^{n-1}\). By induction on \(k\), this holds for every \(k\).
\end{itemize}
(b)
\begin{itemize}
\item \(c_0=(-1)^n\det A\) (Definition~\ref{r7-def-charpoly}(b)), which is nonzero, as \(A\) is invertible
(Theorem~\ref{r5d-thm-product}(a)).
\item Multiplying \(c_0I+c_1A+\dots+c_nA^n=0\) (Theorem~\ref{r7-thm-cayley-hamilton}) by \(A^{-1}\) gives
\(c_0A^{-1}+c_1I+c_2A+\dots+c_nA^{n-1}=0\); dividing by \(c_0\) gives the formula.\qedhere
\end{itemize}
\end{proof}

\begin{example}[cf.\ \citeat{searle-khuri}{Example~7.2}]\label{r7-exm-cayley-hamilton}
Let \(A\) be the matrix of Example~\ref{r7-exm-two},
\[
A=\begin{pmatrix}1&2&-2&-2\\-1&0&-2&2\\0&-1&-1&1\\-1&2&-2&0\end{pmatrix}
\]
Find \(p_A\), check Theorem~\ref{r7-thm-cayley-hamilton}, and use Corollary~\ref{r7-cor-cayley-hamilton} to find
\(A^{-1}\) and \(A^4\).

\textit{Solution.}
\begin{itemize}
\item By Example~\ref{r7-exm-two}, \(1\), \(-1\), \(2\) and \(-2\) are eigenvalues of \(A\), that is, roots of \(p_A\)
(Definition~\ref{r7-def-charpoly}(a)). The eigenvalues \(\lambda_1,\dots,\lambda_4\) of \(A\), listed with multiplicity,
are the roots of \(p_A\) (Definition~\ref{r7-def-eigenvalues}, Theorem~\ref{r7-thm-factor}(c)), so each of the four
distinct numbers \(1\), \(-1\), \(2\), \(-2\) occurs among them, and there is no room for more. Hence
\[
p_A=(t-1)(t+1)(t-2)(t+2)=(t^2-1)(t^2-4)=t^4-5t^2+4
\]
and \(c_0=4\), \(c_1=0\), \(c_2=-5\), \(c_3=0\), \(c_4=1\).
\item Multiplying out,
\[
A^2=\begin{pmatrix}1&0&0&0\\-3&4&0&0\\0&3&1&-3\\-3&0&0&4\end{pmatrix},\qquad
A^4=(A^2)^2=\begin{pmatrix}1&0&0&0\\-15&16&0&0\\0&15&1&-15\\-15&0&0&16\end{pmatrix}
\]
and entry by entry, \(p_A(A)=A^4-5A^2+4I=0\).
\item As \(c_0=(-1)^4\det A\) (Definition~\ref{r7-def-charpoly}(b)), \(\det A=4\ne0\), and \(A\) is invertible
(Theorem~\ref{r5d-thm-product}(a)). With
\[
A^3=A^2A=\begin{pmatrix}1&2&-2&-2\\-7&-6&-2&14\\0&-7&-1&7\\-7&2&-2&6\end{pmatrix}
\]
Corollary~\ref{r7-cor-cayley-hamilton}(b) gives
\[
A^{-1}=-\frac14\bigl(c_1I+c_2A+c_3A^2+c_4A^3\bigr)=\frac14(5A-A^3)
=\frac12\begin{pmatrix}2&4&-4&-4\\1&3&-4&-2\\0&1&-2&-1\\1&4&-4&-3\end{pmatrix}
\]
Check: row \(1\) of \(AA^{-1}\) is \(\frac12(2+2+0-2,\ 4+6-2-8,\ -4-8+4+8,\ -4-4+2+6)=(1,0,0,0)\), and likewise rows
\(2\), \(3\), \(4\) are \(e_2'\), \(e_3'\), \(e_4'\).
\item By Corollary~\ref{r7-cor-cayley-hamilton}(a), \(A^4=-(c_0I+c_1A+c_2A^2+c_3A^3)=5A^2-4I\), as found above; then
\(A^5=5A^3-4A\), and every power of \(A\) is a linear combination of \(I\), \(A\), \(A^2\), \(A^3\).
\end{itemize}
\end{example}

\needspace{14\baselineskip}
\subsection*{Exercises}

\begin{exercise}\label{r7-dr-charpoly}
Find the characteristic polynomial and the eigenvalues of each matrix, and check the trace and the determinant against
the eigenvalues (Theorem~\ref{r7-thm-trace-det}).
\[
\text{(a)}\ \ A=\begin{pmatrix}1&-2&-3&-1\\0&-1&-3&-1\\0&0&2&1\\2&-2&0&3\end{pmatrix},\qquad
\text{(b)}\ \ B=\begin{pmatrix}0&1&0&1\\-1&1&1&1\\-1&3&-1&1\\-1&2&-3&4\end{pmatrix},\qquad
\text{(c)}\ \ M=\begin{pmatrix}2&1&1&-1\\4&-1&2&0\\0&0&0&1\\0&0&-2&3\end{pmatrix}
\]
\end{exercise}
\begin{solution}
(a)
\begin{itemize}
\item For \(t\in\R\), row \(3\) of \(tI-A\) is \((0,\ 0,\ t-2,\ -1)\); expanding along it (Theorem~\ref{r5d-thm-expansion}),
\[
\begin{aligned}
\det(tI-A)&=\det\begin{pmatrix}t-1&2&3&1\\0&t+1&3&1\\0&0&t-2&-1\\-2&2&0&t-3\end{pmatrix}\\
&=(t-2)\det\begin{pmatrix}t-1&2&1\\0&t+1&1\\-2&2&t-3\end{pmatrix}+\det\begin{pmatrix}t-1&2&3\\0&t+1&3\\-2&2&0\end{pmatrix}
\end{aligned}
\]
\item Expanding the first \(3\times3\) determinant along column \(1\) and the second along row \(3\),
\[
\begin{aligned}
\det\begin{pmatrix}t-1&2&1\\0&t+1&1\\-2&2&t-3\end{pmatrix}&=(t-1)\bigl((t+1)(t-3)-2\bigr)-2\bigl(2-(t+1)\bigr)\\
&=(t-1)(t^2-2t-5)+2(t-1)\\
&=(t-1)(t^2-2t-3)=(t-1)(t+1)(t-3)\\
\det\begin{pmatrix}t-1&2&3\\0&t+1&3\\-2&2&0\end{pmatrix}&=-2\bigl(6-3(t+1)\bigr)-2\cdot3(t-1)=(6t-6)-(6t-6)=0
\end{aligned}
\]
So \(\det(tI-A)=(t-2)(t-1)(t+1)(t-3)\), and \(p_A=(t+1)(t-1)(t-2)(t-3)\) (Definition~\ref{r7-def-charpoly}(a),
Proposition~\ref{r7-prop-poly}(d)): the eigenvalues are \(-1\), \(1\), \(2\), \(3\).
\item \(\tr A=1-1+2+3=5=-1+1+2+3\). Expanding \(\det A\) along row \(3\) in the same way, and the two \(3\times3\)
determinants along column \(1\),
\[
\det A=2\det\begin{pmatrix}1&-2&-1\\0&-1&-1\\2&-2&3\end{pmatrix}-\det\begin{pmatrix}1&-2&-3\\0&-1&-3\\2&-2&0\end{pmatrix}
=2\bigl((-3-2)+2(2-1)\bigr)-\bigl((0-6)+2(6-3)\bigr)=-6
\]
and \(-6=(-1)\cdot1\cdot2\cdot3\).
\end{itemize}
(b)
\begin{itemize}
\item Every row of \(B\) has the sum \(2\). For \(t\in\R\), add columns \(2\), \(3\) and \(4\) of \(tI-B\) to column \(1\),
and take the factor \(t-2\) out of column \(1\) (Definition~\ref{r5d-def-det}(b),(a) for columns,
Theorem~\ref{r5d-thm-transpose}):
\[
\begin{aligned}
\det(tI-B)&=\det\begin{pmatrix}t&-1&0&-1\\1&t-1&-1&-1\\1&-3&t+1&-1\\1&-2&3&t-4\end{pmatrix}
=\det\begin{pmatrix}t-2&-1&0&-1\\t-2&t-1&-1&-1\\t-2&-3&t+1&-1\\t-2&-2&3&t-4\end{pmatrix}\\
&=(t-2)\det\begin{pmatrix}1&-1&0&-1\\1&t-1&-1&-1\\1&-3&t+1&-1\\1&-2&3&t-4\end{pmatrix}
\end{aligned}
\]
\item Subtract row \(1\) from rows \(2\), \(3\) and \(4\) (Definition~\ref{r5d-def-det}(b)), expand along column \(1\),
and then along column \(3\) (Theorem~\ref{r5d-thm-expansion}):
\[
\begin{aligned}
\det(tI-B)&=(t-2)\det\begin{pmatrix}1&-1&0&-1\\0&t&-1&0\\0&-2&t+1&0\\0&-1&3&t-3\end{pmatrix}
=(t-2)\det\begin{pmatrix}t&-1&0\\-2&t+1&0\\-1&3&t-3\end{pmatrix}\\
&=(t-2)(t-3)(t^2+t-2)
\end{aligned}
\]
So \(p_B=(t-2)(t-3)(t+2)(t-1)\), and the eigenvalues are \(-2\), \(1\), \(2\), \(3\).
\item \(\tr B=0+1-1+4=4=-2+1+2+3\). Expanding \(\det B\) along row \(1\), and the two \(3\times3\) determinants along
column \(1\),
\[
\det B=-\det\begin{pmatrix}-1&1&1\\-1&-1&1\\-1&-3&4\end{pmatrix}-\det\begin{pmatrix}-1&1&1\\-1&3&-1\\-1&2&-3\end{pmatrix}
=-(1+7-2)-(7-5+4)=-12
\]
and \(-12=(-2)\cdot1\cdot2\cdot3\).
\end{itemize}
(c)
\begin{itemize}
\item \(M=\left(\begin{smallmatrix}M_1&M_2\\0&M_3\end{smallmatrix}\right)\) is block upper triangular, with
\(M_1=\left(\begin{smallmatrix}2&1\\4&-1\end{smallmatrix}\right)\) and \(M_3=\left(\begin{smallmatrix}0&1\\-2&3\end{smallmatrix}\right)\).
By Definition~\ref{r7-def-charpoly}(b) (\(n=2\)), \(p_{M_1}=t^2-t-6=(t-3)(t+2)\) and \(p_{M_3}=t^2-3t+2=(t-1)(t-2)\).
\item By Proposition~\ref{r7-prop-block-triangular}(a), \(p_M=p_{M_1}p_{M_3}=(t-3)(t+2)(t-1)(t-2)\): the eigenvalues are
\(-2\), \(1\), \(2\), \(3\).
\item \(\tr M=2-1+0+3=4=-2+1+2+3\), and \(\det M=\det M_1\det M_3=(-6)\cdot2=-12=(-2)\cdot1\cdot2\cdot3\)
(Theorem~\ref{r5d-thm-block-tri}).
\end{itemize}
\end{solution}

\begin{exercise}\label{r7-ex-eigenpairs}
Find the eigenvalues of
\[
B=\begin{pmatrix}3&1&0&-1\\0&1&0&0\\0&1&3&0\\-1&0&1&3\end{pmatrix}
\]
and a basis of each eigenspace. Check that the coefficient of \(t^3\) in \(p_B(t)\) is \(-\tr B\) and that the constant
term is \(\det B\).
\end{exercise}
\begin{solution}
For \(t\in\R\), row \(2\) of \(tI_4-B\) is \((0,\ t-1,\ 0,\ 0)\). Expanding along it, and then the remaining
\(3\times3\) determinant (rows and columns \(1\), \(3\), \(4\)) along its second row \((0,\ t-3,\ 0)\)
(Theorem~\ref{r5d-thm-expansion}),
\[
\begin{aligned}
\det(tI_4-B)&=(t-1)\det\begin{pmatrix}t-3&0&1\\0&t-3&0\\1&-1&t-3\end{pmatrix}
=(t-1)(t-3)\det\begin{pmatrix}t-3&1\\1&t-3\end{pmatrix}\\
&=(t-1)(t-3)(t-2)(t-4)
\end{aligned}
\]
So \(p_B=(t-1)(t-2)(t-3)(t-4)\) (Definition~\ref{r7-def-charpoly}(a), Proposition~\ref{r7-prop-poly}(d)), and
multiplying out, \(p_B=t^4-10t^3+35t^2-50t+24\). Here \(\tr B=3+1+3+3=10\), and expanding \(\det B\) in the same
way, \(\det B=1\cdot3\cdot(9-1)=24\); so the coefficient of \(t^3\) is \(-\tr B\) and the constant term is
\((-1)^4\det B\), as Definition~\ref{r7-def-charpoly}(b) says. The eigenvalues are \(1\), \(2\), \(3\), \(4\)
(Definition~\ref{r7-def-charpoly}(a)).
\begin{itemize}
\item \(\lambda=1\):
\[
\begin{gathered}
B-I_4=\begin{pmatrix}2&1&0&-1\\0&0&0&0\\0&1&2&0\\-1&0&1&2\end{pmatrix}
\xrightarrow{R_1\leftrightarrow R_4}
\begin{pmatrix}-1&0&1&2\\0&0&0&0\\0&1&2&0\\2&1&0&-1\end{pmatrix}
\xrightarrow{R_1\leftarrow-R_1}
\begin{pmatrix}1&0&-1&-2\\0&0&0&0\\0&1&2&0\\2&1&0&-1\end{pmatrix}
\\
\xrightarrow{R_4\leftarrow R_4-2R_1}
\begin{pmatrix}1&0&-1&-2\\0&0&0&0\\0&1&2&0\\0&1&2&3\end{pmatrix}
\xrightarrow{R_2\leftrightarrow R_4}
\begin{pmatrix}1&0&-1&-2\\0&1&2&3\\0&1&2&0\\0&0&0&0\end{pmatrix}
\\
\xrightarrow{R_3\leftarrow R_3-R_2}
\begin{pmatrix}1&0&-1&-2\\0&1&2&3\\0&0&0&-3\\0&0&0&0\end{pmatrix}
\xrightarrow{R_3\leftarrow-\frac13R_3}
\begin{pmatrix}1&0&-1&-2\\0&1&2&3\\0&0&0&1\\0&0&0&0\end{pmatrix}
\\
\xrightarrow[R_2\leftarrow R_2-3R_3]{R_1\leftarrow R_1+2R_3}
\begin{pmatrix}1&0&-1&0\\0&1&2&0\\0&0&0&1\\0&0&0&0\end{pmatrix}
\end{gathered}
\]
The special solution \((1,-2,1,0)'\) is a basis of \(\mathcal N(B-I_4)\) (Theorem~\ref{thm:r4-four-bases}(c)).
\item \(\lambda=2\):
\[
\begin{gathered}
B-2I_4=\begin{pmatrix}1&1&0&-1\\0&-1&0&0\\0&1&1&0\\-1&0&1&1\end{pmatrix}
\xrightarrow{R_4\leftarrow R_4+R_1}
\begin{pmatrix}1&1&0&-1\\0&-1&0&0\\0&1&1&0\\0&1&1&0\end{pmatrix}
\\
\xrightarrow[R_4\leftarrow R_4+R_2]{R_1\leftarrow R_1+R_2,\ R_3\leftarrow R_3+R_2}
\begin{pmatrix}1&0&0&-1\\0&-1&0&0\\0&0&1&0\\0&0&1&0\end{pmatrix}
\xrightarrow{R_4\leftarrow R_4-R_3}
\begin{pmatrix}1&0&0&-1\\0&-1&0&0\\0&0&1&0\\0&0&0&0\end{pmatrix}
\\
\xrightarrow{R_2\leftarrow-R_2}
\begin{pmatrix}1&0&0&-1\\0&1&0&0\\0&0&1&0\\0&0&0&0\end{pmatrix}
\end{gathered}
\]
The special solution \((1,0,0,1)'\) is a basis of \(\mathcal N(B-2I_4)\).
\item \(\lambda=3\):
\[
\begin{gathered}
B-3I_4=\begin{pmatrix}0&1&0&-1\\0&-2&0&0\\0&1&0&0\\-1&0&1&0\end{pmatrix}
\xrightarrow[R_2\leftrightarrow R_3]{R_1\leftrightarrow R_4}
\begin{pmatrix}-1&0&1&0\\0&1&0&0\\0&-2&0&0\\0&1&0&-1\end{pmatrix}
\xrightarrow[R_4\leftarrow R_4-R_2]{R_3\leftarrow R_3+2R_2}
\begin{pmatrix}-1&0&1&0\\0&1&0&0\\0&0&0&0\\0&0&0&-1\end{pmatrix}
\\
\xrightarrow{R_3\leftrightarrow R_4}
\begin{pmatrix}-1&0&1&0\\0&1&0&0\\0&0&0&-1\\0&0&0&0\end{pmatrix}
\xrightarrow[R_3\leftarrow-R_3]{R_1\leftarrow-R_1}
\begin{pmatrix}1&0&-1&0\\0&1&0&0\\0&0&0&1\\0&0&0&0\end{pmatrix}
\end{gathered}
\]
The special solution \((1,0,1,0)'\) is a basis of \(\mathcal N(B-3I_4)\).
\item \(\lambda=4\):
\[
\begin{gathered}
B-4I_4=\begin{pmatrix}-1&1&0&-1\\0&-3&0&0\\0&1&-1&0\\-1&0&1&-1\end{pmatrix}
\xrightarrow{R_4\leftarrow R_4-R_1}
\begin{pmatrix}-1&1&0&-1\\0&-3&0&0\\0&1&-1&0\\0&-1&1&0\end{pmatrix}
\\
\xrightarrow{R_2\leftarrow-\frac13R_2}
\begin{pmatrix}-1&1&0&-1\\0&1&0&0\\0&1&-1&0\\0&-1&1&0\end{pmatrix}
\xrightarrow[R_4\leftarrow R_4+R_2]{R_1\leftarrow R_1-R_2,\ R_3\leftarrow R_3-R_2}
\begin{pmatrix}-1&0&0&-1\\0&1&0&0\\0&0&-1&0\\0&0&1&0\end{pmatrix}
\\
\xrightarrow{R_4\leftarrow R_4+R_3}
\begin{pmatrix}-1&0&0&-1\\0&1&0&0\\0&0&-1&0\\0&0&0&0\end{pmatrix}
\xrightarrow[R_3\leftarrow-R_3]{R_1\leftarrow-R_1}
\begin{pmatrix}1&0&0&1\\0&1&0&0\\0&0&1&0\\0&0&0&0\end{pmatrix}
\end{gathered}
\]
The special solution \((-1,0,0,1)'\) is a basis of \(\mathcal N(B-4I_4)\).
\end{itemize}
Check: \(B(1,-2,1,0)'=(1,-2,1,0)'\) and \(B(-1,0,0,1)'=(-4,0,0,4)'\). Also \(1+2+3+4=10=\tr B\) and
\(1\cdot2\cdot3\cdot4=24=\det B\) (Theorem~\ref{r7-thm-trace-det}).
\end{solution}

\begin{exercise}\label{r7-dr-eigenspaces}
Let
\[
A=\begin{pmatrix}0&0&2&-1\\0&3&0&0\\1&0&1&-1\\-2&-2&2&1\end{pmatrix},\qquad
B=\begin{pmatrix}1&-2&-2&-2\\2&1&2&2\\0&0&1&0\\-2&2&0&1\end{pmatrix}
\]
\begin{enumerate}
\item Find the eigenvalues of \(A\) and a basis of each eigenspace.
\item Do the same for \(B\).
\item Find an invertible \(P\) and a diagonal \(\Lambda\) with \(AP=P\Lambda\), and an invertible \(Q\) and a diagonal \(\Delta\)
with \(BQ=Q\Delta\); check both equations by multiplying out.
\end{enumerate}
\end{exercise}
\begin{solution}
(a)
\begin{itemize}
\item For \(t\in\R\), row \(2\) of \(tI-A\) is \((0,\ t-3,\ 0,\ 0)\); expanding along it, and then along row \(1\)
(Theorem~\ref{r5d-thm-expansion}),
\[
\begin{aligned}
\det(tI-A)&=\det\begin{pmatrix}t&0&-2&1\\0&t-3&0&0\\-1&0&t-1&1\\2&2&-2&t-1\end{pmatrix}
=(t-3)\det\begin{pmatrix}t&-2&1\\-1&t-1&1\\2&-2&t-1\end{pmatrix}\\
&=(t-3)\Bigl(t\bigl((t-1)^2+2\bigr)+2\bigl(-(t-1)-2\bigr)+\bigl(2-2(t-1)\bigr)\Bigr)\\
&=(t-3)\bigl(t^3-2t^2+3t-2t-2+4-2t\bigr)\\
&=(t-3)(t^3-2t^2-t+2)=(t-3)(t-2)(t^2-1)
\end{aligned}
\]
So \(p_A=(t+1)(t-1)(t-2)(t-3)\): the eigenvalues are \(-1\), \(1\), \(2\), \(3\).
\item \(\lambda=-1\):
\[
\begin{gathered}
A+I=\begin{pmatrix}1&0&2&-1\\0&4&0&0\\1&0&2&-1\\-2&-2&2&2\end{pmatrix}
\xrightarrow[R_4\leftarrow R_4+2R_1]{R_3\leftarrow R_3-R_1}
\begin{pmatrix}1&0&2&-1\\0&4&0&0\\0&0&0&0\\0&-2&6&0\end{pmatrix}
\xrightarrow{R_2\leftarrow\frac14R_2}
\begin{pmatrix}1&0&2&-1\\0&1&0&0\\0&0&0&0\\0&-2&6&0\end{pmatrix}
\\
\xrightarrow{R_4\leftarrow R_4+2R_2}
\begin{pmatrix}1&0&2&-1\\0&1&0&0\\0&0&0&0\\0&0&6&0\end{pmatrix}
\xrightarrow{R_3\leftrightarrow R_4}
\begin{pmatrix}1&0&2&-1\\0&1&0&0\\0&0&6&0\\0&0&0&0\end{pmatrix}
\\
\xrightarrow{R_3\leftarrow\frac16R_3}
\begin{pmatrix}1&0&2&-1\\0&1&0&0\\0&0&1&0\\0&0&0&0\end{pmatrix}
\xrightarrow{R_1\leftarrow R_1-2R_3}
\begin{pmatrix}1&0&0&-1\\0&1&0&0\\0&0&1&0\\0&0&0&0\end{pmatrix}
\end{gathered}
\]
The special solution \((1,0,0,1)'\) is a basis of \(\mathcal N(A+I)\) (Theorem~\ref{thm:r4-four-bases}(c)).
\item \(\lambda=1\):
\[
\begin{gathered}
A-I=\begin{pmatrix}-1&0&2&-1\\0&2&0&0\\1&0&0&-1\\-2&-2&2&0\end{pmatrix}
\xrightarrow[R_4\leftarrow R_4-2R_1]{R_3\leftarrow R_3+R_1}
\begin{pmatrix}-1&0&2&-1\\0&2&0&0\\0&0&2&-2\\0&-2&-2&2\end{pmatrix}
\xrightarrow{R_4\leftarrow R_4+R_2}
\begin{pmatrix}-1&0&2&-1\\0&2&0&0\\0&0&2&-2\\0&0&-2&2\end{pmatrix}
\\
\xrightarrow[R_4\leftarrow R_4+R_3]{R_1\leftarrow R_1-R_3}
\begin{pmatrix}-1&0&0&1\\0&2&0&0\\0&0&2&-2\\0&0&0&0\end{pmatrix}
\xrightarrow[R_3\leftarrow\frac12R_3]{R_1\leftarrow-R_1,\ R_2\leftarrow\frac12R_2}
\begin{pmatrix}1&0&0&-1\\0&1&0&0\\0&0&1&-1\\0&0&0&0\end{pmatrix}
\end{gathered}
\]
The special solution \((1,0,1,1)'\) is a basis of \(\mathcal N(A-I)\).
\item \(\lambda=2\):
\[
\begin{gathered}
A-2I=\begin{pmatrix}-2&0&2&-1\\0&1&0&0\\1&0&-1&-1\\-2&-2&2&-1\end{pmatrix}
\xrightarrow{R_1\leftrightarrow R_3}
\begin{pmatrix}1&0&-1&-1\\0&1&0&0\\-2&0&2&-1\\-2&-2&2&-1\end{pmatrix}
\xrightarrow[R_4\leftarrow R_4+2R_1]{R_3\leftarrow R_3+2R_1}
\begin{pmatrix}1&0&-1&-1\\0&1&0&0\\0&0&0&-3\\0&-2&0&-3\end{pmatrix}
\\
\xrightarrow{R_4\leftarrow R_4+2R_2}
\begin{pmatrix}1&0&-1&-1\\0&1&0&0\\0&0&0&-3\\0&0&0&-3\end{pmatrix}
\xrightarrow{R_3\leftarrow-\frac13R_3}
\begin{pmatrix}1&0&-1&-1\\0&1&0&0\\0&0&0&1\\0&0&0&-3\end{pmatrix}
\\
\xrightarrow[R_4\leftarrow R_4+3R_3]{R_1\leftarrow R_1+R_3}
\begin{pmatrix}1&0&-1&0\\0&1&0&0\\0&0&0&1\\0&0&0&0\end{pmatrix}
\end{gathered}
\]
The special solution \((1,0,1,0)'\) is a basis of \(\mathcal N(A-2I)\).
\item \(\lambda=3\):
\[
\begin{gathered}
A-3I=\begin{pmatrix}-3&0&2&-1\\0&0&0&0\\1&0&-2&-1\\-2&-2&2&-2\end{pmatrix}
\xrightarrow{R_1\leftrightarrow R_3}
\begin{pmatrix}1&0&-2&-1\\0&0&0&0\\-3&0&2&-1\\-2&-2&2&-2\end{pmatrix}
\xrightarrow[R_4\leftarrow R_4+2R_1]{R_3\leftarrow R_3+3R_1}
\begin{pmatrix}1&0&-2&-1\\0&0&0&0\\0&0&-4&-4\\0&-2&-2&-4\end{pmatrix}
\\
\xrightarrow{R_2\leftrightarrow R_4}
\begin{pmatrix}1&0&-2&-1\\0&-2&-2&-4\\0&0&-4&-4\\0&0&0&0\end{pmatrix}
\xrightarrow[R_3\leftarrow-\frac14R_3]{R_2\leftarrow-\frac12R_2}
\begin{pmatrix}1&0&-2&-1\\0&1&1&2\\0&0&1&1\\0&0&0&0\end{pmatrix}
\\
\xrightarrow[R_2\leftarrow R_2-R_3]{R_1\leftarrow R_1+2R_3}
\begin{pmatrix}1&0&0&1\\0&1&0&1\\0&0&1&1\\0&0&0&0\end{pmatrix}
\end{gathered}
\]
The special solution \((-1,-1,-1,1)'\) is a basis of \(\mathcal N(A-3I)\).
\end{itemize}
(b)
\begin{itemize}
\item For \(t\in\R\), row \(3\) of \(tI-B\) is \((0,\ 0,\ t-1,\ 0)\); expanding along it, and then along row \(1\),
\[
\begin{aligned}
\det(tI-B)&=\det\begin{pmatrix}t-1&2&2&2\\-2&t-1&-2&-2\\0&0&t-1&0\\2&-2&0&t-1\end{pmatrix}
=(t-1)\det\begin{pmatrix}t-1&2&2\\-2&t-1&-2\\2&-2&t-1\end{pmatrix}\\
&=(t-1)\Bigl((t-1)\bigl((t-1)^2-4\bigr)-2\bigl(-2(t-1)+4\bigr)+2\bigl(4-2(t-1)\bigr)\Bigr)\\
&=(t-1)\bigl((t-1)(t-3)(t+1)+4(t-1)-8+8-4(t-1)\bigr)\\
&=(t-1)^2(t-3)(t+1)
\end{aligned}
\]
So \(p_B=(t-1)^2(t-3)(t+1)\): the eigenvalues are \(1\), \(1\), \(3\), \(-1\), with \(a_B(1)=2\).
\item \(\lambda=1\):
\[
\begin{gathered}
B-I=\begin{pmatrix}0&-2&-2&-2\\2&0&2&2\\0&0&0&0\\-2&2&0&0\end{pmatrix}
\xrightarrow{R_1\leftrightarrow R_2}
\begin{pmatrix}2&0&2&2\\0&-2&-2&-2\\0&0&0&0\\-2&2&0&0\end{pmatrix}
\xrightarrow{R_4\leftarrow R_4+R_1}
\begin{pmatrix}2&0&2&2\\0&-2&-2&-2\\0&0&0&0\\0&2&2&2\end{pmatrix}
\\
\xrightarrow{R_4\leftarrow R_4+R_2}
\begin{pmatrix}2&0&2&2\\0&-2&-2&-2\\0&0&0&0\\0&0&0&0\end{pmatrix}
\xrightarrow[R_2\leftarrow-\frac12R_2]{R_1\leftarrow\frac12R_1}
\begin{pmatrix}1&0&1&1\\0&1&1&1\\0&0&0&0\\0&0&0&0\end{pmatrix}
\end{gathered}
\]
The free unknowns are \(x_3\), \(x_4\); the special solutions \((-1,-1,1,0)'\), \((-1,-1,0,1)'\) form a basis of
\(\mathcal N(B-I)\), and \(g_B(1)=2\).
\item \(\lambda=3\):
\[
\begin{gathered}
B-3I=\begin{pmatrix}-2&-2&-2&-2\\2&-2&2&2\\0&0&-2&0\\-2&2&0&-2\end{pmatrix}
\xrightarrow[R_4\leftarrow R_4-R_1]{R_2\leftarrow R_2+R_1}
\begin{pmatrix}-2&-2&-2&-2\\0&-4&0&0\\0&0&-2&0\\0&4&2&0\end{pmatrix}
\\
\xrightarrow{R_2\leftarrow-\frac14R_2}
\begin{pmatrix}-2&-2&-2&-2\\0&1&0&0\\0&0&-2&0\\0&4&2&0\end{pmatrix}
\xrightarrow[R_4\leftarrow R_4-4R_2]{R_1\leftarrow R_1+2R_2}
\begin{pmatrix}-2&0&-2&-2\\0&1&0&0\\0&0&-2&0\\0&0&2&0\end{pmatrix}
\\
\xrightarrow[R_4\leftarrow R_4+R_3]{R_1\leftarrow R_1-R_3}
\begin{pmatrix}-2&0&0&-2\\0&1&0&0\\0&0&-2&0\\0&0&0&0\end{pmatrix}
\xrightarrow[R_3\leftarrow-\frac12R_3]{R_1\leftarrow-\frac12R_1}
\begin{pmatrix}1&0&0&1\\0&1&0&0\\0&0&1&0\\0&0&0&0\end{pmatrix}
\end{gathered}
\]
The special solution \((-1,0,0,1)'\) is a basis of \(\mathcal N(B-3I)\).
\item \(\lambda=-1\):
\[
\begin{gathered}
B+I=\begin{pmatrix}2&-2&-2&-2\\2&2&2&2\\0&0&2&0\\-2&2&0&2\end{pmatrix}
\xrightarrow{R_1\leftarrow\frac12R_1}
\begin{pmatrix}1&-1&-1&-1\\2&2&2&2\\0&0&2&0\\-2&2&0&2\end{pmatrix}
\xrightarrow[R_4\leftarrow R_4+2R_1]{R_2\leftarrow R_2-2R_1}
\begin{pmatrix}1&-1&-1&-1\\0&4&4&4\\0&0&2&0\\0&0&-2&0\end{pmatrix}
\\
\xrightarrow{R_2\leftarrow\frac14R_2}
\begin{pmatrix}1&-1&-1&-1\\0&1&1&1\\0&0&2&0\\0&0&-2&0\end{pmatrix}
\xrightarrow{R_1\leftarrow R_1+R_2}
\begin{pmatrix}1&0&0&0\\0&1&1&1\\0&0&2&0\\0&0&-2&0\end{pmatrix}
\\
\xrightarrow{R_3\leftarrow\frac12R_3}
\begin{pmatrix}1&0&0&0\\0&1&1&1\\0&0&1&0\\0&0&-2&0\end{pmatrix}
\xrightarrow[R_4\leftarrow R_4+2R_3]{R_2\leftarrow R_2-R_3}
\begin{pmatrix}1&0&0&0\\0&1&0&1\\0&0&1&0\\0&0&0&0\end{pmatrix}
\end{gathered}
\]
The special solution \((0,-1,0,1)'\) is a basis of \(\mathcal N(B+I)\).
\end{itemize}
(c)
\begin{itemize}
\item The four eigenvectors of \(A\) belong to distinct eigenvalues, so they are independent
(Theorem~\ref{r7-thm-indep}), and \(P\) below is invertible (Theorem~\ref{thm:r4-rank-rules}(b),(d)):
\[
P=\begin{pmatrix}1&1&1&-1\\0&0&0&-1\\0&1&1&-1\\1&1&0&1\end{pmatrix},\qquad\Lambda=\diag(-1,1,2,3),\qquad
AP=\begin{pmatrix}-1&1&2&-3\\0&0&0&-3\\0&1&2&-3\\-1&1&0&3\end{pmatrix}=P\Lambda
\]
\item The bases of \(\mathcal N(B-I)\), \(\mathcal N(B-3I)\) and \(\mathcal N(B+I)\) together are independent
(Lemma~\ref{r7-lem-eigenspaces}(b)), so \(Q\) below is invertible (Theorem~\ref{thm:r4-rank-rules}(b),(d)):
\[
Q=\begin{pmatrix}-1&-1&-1&0\\-1&-1&0&-1\\1&0&0&0\\0&1&1&1\end{pmatrix},\qquad\Delta=\diag(1,1,3,-1),\qquad
BQ=\begin{pmatrix}-1&-1&-3&0\\-1&-1&0&1\\1&0&0&0\\0&1&3&-1\end{pmatrix}=Q\Delta
\]
\end{itemize}
\end{solution}

\begin{exercise}[\citeat{lang}{Ch.~VIII, \S2, Exercise~4(a), (b)}, \citeat{shakarchi}{}]\label{r7-ex-lang-eigen}
Find the characteristic polynomial, the eigenvalues, and a basis of each eigenspace of
\[
\text{(a)}\ \ A=\begin{pmatrix}4&0&1\\-2&1&0\\-2&0&1\end{pmatrix},\qquad
\text{(b)}\ \ B=\begin{pmatrix}1&-3&3\\3&-5&3\\6&-6&4\end{pmatrix}
\]
\end{exercise}
\begin{solution}
(a)
\begin{itemize}
\item For \(t\in\R\), expanding along row \(1\) (Theorem~\ref{r5d-thm-expansion}),
\[
\det(tI-A)=\det\begin{pmatrix}t-4&0&-1\\2&t-1&0\\2&0&t-1\end{pmatrix}
=(t-4)(t-1)^2-\det\begin{pmatrix}2&t-1\\2&0\end{pmatrix}=(t-4)(t-1)^2+2(t-1)
\]
So
\(p_A=(t-1)\bigl((t-4)(t-1)+2\bigr)=(t-1)(t^2-5t+6)=(t-1)(t-2)(t-3)\): the eigenvalues are \(1\), \(2\), \(3\).
\item \(Ax=\lambda x\) reads \(4x_1+x_3=\lambda x_1\), \(-2x_1+x_2=\lambda x_2\), \(-2x_1+x_3=\lambda x_3\).
\item \(\lambda=1\): \(3x_1+x_3=0\) and \(-2x_1=0\); so \(x_1=x_3=0\), \(x_2\) is free, and \((0,1,0)'\) is a basis of
\(\mathcal N(A-I)\).
\item \(\lambda=2\): \(2x_1+x_3=0\), \(-2x_1-x_2=0\), \(-2x_1-x_3=0\); so \(x_3=-2x_1\) and \(x_2=-2x_1\), and \((1,-2,-2)'\) is a
basis of \(\mathcal N(A-2I)\).
\item \(\lambda=3\): \(x_1+x_3=0\), \(-2x_1-2x_2=0\), \(-2x_1-2x_3=0\); so \(x_3=-x_1\) and \(x_2=-x_1\), and \((1,-1,-1)'\) is a
basis of \(\mathcal N(A-3I)\).
\end{itemize}
(b)
\begin{itemize}
\item For \(t\in\R\), expanding along row \(1\),
\[
\begin{aligned}
\det(tI-B)&=\det\begin{pmatrix}t-1&3&-3\\-3&t+5&-3\\-6&6&t-4\end{pmatrix}\\
&=(t-1)\bigl((t+5)(t-4)+18\bigr)-3\bigl(-3(t-4)-18\bigr)-3\bigl(-18+6(t+5)\bigr)\\
&=(t-1)(t^2+t-2)+9(t+2)-18(t+2)=(t+2)\bigl((t-1)^2-9\bigr)=(t+2)^2(t-4)
\end{aligned}
\]
So \(p_B=t^3-12t-16=(t+2)^2(t-4)\): the eigenvalues are \(-2\), \(-2\), \(4\).
\item \(\lambda=-2\): \(Bx=-2x\) reads \(3x_1-3x_2+3x_3=0\), \(3x_1-3x_2+3x_3=0\), \(6x_1-6x_2+6x_3=0\), that is,
\(x_1-x_2+x_3=0\). Every row of \(B+2I\) is a multiple of \((1,-1,1)\), so the row space of \(B+2I\) is
\(\Span\{(1,-1,1)\}\), \(\rank(B+2I)=1\) (Definition~\ref{def:r4-rank}) and \(\dim\mathcal N(B+2I)=2\)
(Theorem~\ref{thm:r4-rank-nullity}). The solutions \((1,0,-1)'\) and \((0,1,1)'\) are not multiples
of each other, so they are independent (Proposition~\ref{r3-prop-dependence}(c)) and form a basis of
\(\mathcal N(B+2I)\) (Theorem~\ref{thm:r4-counting}(d)).
\item \(\lambda=4\): \(-3x_1-3x_2+3x_3=0\), \(3x_1-9x_2+3x_3=0\), \(6x_1-6x_2=0\); so \(x_2=x_1\), then \(x_3=2x_1\) by the
first equation, and the second holds: \((1,1,2)'\) is a basis of \(\mathcal N(B-4I)\).
\end{itemize}
\end{solution}

\begin{exercise}\label{r7-ex-complex-drill}
Let
\[
A=\begin{pmatrix}3&0&0&2\\-2&1&0&-2\\0&-3&2&1\\-2&2&0&-1\end{pmatrix}
\]
\begin{enumerate}
\item Find \(p_A\) and all eigenvalues of \(A\), and check Theorem~\ref{r7-thm-trace-det}.
\item Find the eigenspace of each real eigenvalue. Is \(A\) diagonalizable?
\end{enumerate}
\end{exercise}
\begin{solution}
(a)
\begin{itemize}
\item For \(t\in\R\), column \(3\) of \(tI-A\) has the entry \(t-2\) in position \((3,3)\) and \(0\) elsewhere; expanding
along it, and then along row \(1\) (Theorem~\ref{r5d-thm-expansion}),
\[
\begin{aligned}
\det(tI-A)&=\det\begin{pmatrix}t-3&0&0&-2\\2&t-1&0&2\\0&3&t-2&-1\\2&-2&0&t+1\end{pmatrix}
=(t-2)\det\begin{pmatrix}t-3&0&-2\\2&t-1&2\\2&-2&t+1\end{pmatrix}\\
&=(t-2)\Bigl((t-3)\bigl((t-1)(t+1)+4\bigr)-2\bigl(-4-2(t-1)\bigr)\Bigr)\\
&=(t-2)(t^3-3t^2+7t-5)
\end{aligned}
\]
\item Let \(q=t^3-3t^2+7t-5\). Then \(q(1)=1-3+7-5=0\), so \(q=(t-1)h\) for a polynomial \(h\)
(Proposition~\ref{r7-prop-poly}(c)). Multiplying out, \((t-1)(t^2-2t+5)=t^3-3t^2+7t-5=q\); so \(h=t^2-2t+5\)
(Proposition~\ref{r7-prop-poly}(a)). Since \(t^2-2t+5=(t-1)^2+4=(t-1-2i)(t-1+2i)\),
\(p_A=(t-2)(t-1)\bigl(t-(1+2i)\bigr)\bigl(t-(1-2i)\bigr)\) (Definition~\ref{r7-def-charpoly}(a),
Proposition~\ref{r7-prop-poly}(d)): the eigenvalues are \(2\), \(1\), \(1+2i\), \(1-2i\).
\item \(\tr A=3+1+2-1=5=2+1+(1+2i)+(1-2i)\). Expanding \(\det A\) along column \(3\), and then along row \(1\)
(Theorem~\ref{r5d-thm-expansion}),
\[
\det A=2\det\begin{pmatrix}3&0&2\\-2&1&-2\\-2&2&-1\end{pmatrix}=2\bigl(3(-1+4)+2(-4+2)\bigr)=10
\]
and \(10=2\cdot1\cdot(1+2i)(1-2i)\) (Theorem~\ref{r7-thm-trace-det}).
\end{itemize}
(b) The real eigenvalues are \(1\) and \(2\). Row-reduce \(A-I\) and \(A-2I\):
\[
\begin{gathered}
A-I=\begin{pmatrix}2&0&0&2\\-2&0&0&-2\\0&-3&1&1\\-2&2&0&-2\end{pmatrix}
\xrightarrow{R_1\leftarrow\frac12R_1}
\begin{pmatrix}1&0&0&1\\-2&0&0&-2\\0&-3&1&1\\-2&2&0&-2\end{pmatrix}
\xrightarrow[R_4\leftarrow R_4+2R_1]{R_2\leftarrow R_2+2R_1}
\begin{pmatrix}1&0&0&1\\0&0&0&0\\0&-3&1&1\\0&2&0&0\end{pmatrix}
\\
\xrightarrow{R_2\leftrightarrow R_4}
\begin{pmatrix}1&0&0&1\\0&2&0&0\\0&-3&1&1\\0&0&0&0\end{pmatrix}
\xrightarrow{R_2\leftarrow\frac12R_2}
\begin{pmatrix}1&0&0&1\\0&1&0&0\\0&-3&1&1\\0&0&0&0\end{pmatrix}
\\
\xrightarrow{R_3\leftarrow R_3+3R_2}
\begin{pmatrix}1&0&0&1\\0&1&0&0\\0&0&1&1\\0&0&0&0\end{pmatrix}
\end{gathered}
\]
\[
\begin{gathered}
A-2I=\begin{pmatrix}1&0&0&2\\-2&-1&0&-2\\0&-3&0&1\\-2&2&0&-3\end{pmatrix}
\xrightarrow[R_4\leftarrow R_4+2R_1]{R_2\leftarrow R_2+2R_1}
\begin{pmatrix}1&0&0&2\\0&-1&0&2\\0&-3&0&1\\0&2&0&1\end{pmatrix}
\xrightarrow{R_2\leftarrow-R_2}
\begin{pmatrix}1&0&0&2\\0&1&0&-2\\0&-3&0&1\\0&2&0&1\end{pmatrix}
\\
\xrightarrow[R_4\leftarrow R_4-2R_2]{R_3\leftarrow R_3+3R_2}
\begin{pmatrix}1&0&0&2\\0&1&0&-2\\0&0&0&-5\\0&0&0&5\end{pmatrix}
\xrightarrow{R_3\leftarrow-\frac15R_3}
\begin{pmatrix}1&0&0&2\\0&1&0&-2\\0&0&0&1\\0&0&0&5\end{pmatrix}
\\
\xrightarrow[R_4\leftarrow R_4-5R_3]{R_1\leftarrow R_1-2R_3,\ R_2\leftarrow R_2+2R_3}
\begin{pmatrix}1&0&0&0\\0&1&0&0\\0&0&0&1\\0&0&0&0\end{pmatrix}
\end{gathered}
\]
The special solutions give \(\mathcal N(A-I)=\Span\{(-1,0,-1,1)'\}\) and \(\mathcal N(A-2I)=\Span\{e_3\}\)
(Theorem~\ref{thm:r4-four-bases}(c)). \(A\) is not diagonalizable, since it has non-real eigenvalues
(Definition~\ref{r7-def-diagonalizable}).
\end{solution}

\begin{exercise}\label{r7-dr-similar}
Let
\[
A=\begin{pmatrix}1&1&0&-1\\0&2&0&1\\0&0&3&0\\0&0&0&4\end{pmatrix},\qquad
P=\begin{pmatrix}1&0&0&0\\-1&1&0&0\\0&1&1&0\\-1&1&0&1\end{pmatrix}
\]
\begin{enumerate}
\item Find \(P^{-1}\) by Gauss--Jordan elimination, and compute \(B=P^{-1}AP\).
\item Find \(p_B\), and check that \(p_B=p_A\), \(\tr B=\tr A\) and \(\det B=\det A\).
\item The vector \(x=(1,1,0,0)'\) is an eigenvector of \(A\) belonging to \(2\). Find an eigenvector of \(B\) belonging
to \(2\).
\item Show that
\[
F=\begin{pmatrix}1&1&1&1\\0&1&1&1\\0&0&1&1\\0&0&0&2\end{pmatrix},\qquad
G=\begin{pmatrix}1&0&1&1\\0&1&1&1\\0&0&1&1\\0&0&0&2\end{pmatrix}
\]
have the same characteristic polynomial, but are not similar.
\end{enumerate}
\end{exercise}
\begin{solution}
(a) Row-reduce \([P\mid I_4]\) (Corollary~\ref{r2-cor-computing-inverse}(a)):
\[
\begin{gathered}
\left(\begin{array}{cccc|cccc}1&0&0&0&1&0&0&0\\-1&1&0&0&0&1&0&0\\0&1&1&0&0&0&1&0\\-1&1&0&1&0&0&0&1\end{array}\right)
\xrightarrow[R_4\leftarrow R_4+R_1]{R_2\leftarrow R_2+R_1}
\left(\begin{array}{cccc|cccc}1&0&0&0&1&0&0&0\\0&1&0&0&1&1&0&0\\0&1&1&0&0&0&1&0\\0&1&0&1&1&0&0&1\end{array}\right)
\\
\xrightarrow[R_4\leftarrow R_4-R_2]{R_3\leftarrow R_3-R_2}
\left(\begin{array}{cccc|cccc}1&0&0&0&1&0&0&0\\0&1&0&0&1&1&0&0\\0&0&1&0&-1&-1&1&0\\0&0&0&1&0&-1&0&1\end{array}\right)
\end{gathered}
\]
So
\[
P^{-1}=\begin{pmatrix}1&0&0&0\\1&1&0&0\\-1&-1&1&0\\0&-1&0&1\end{pmatrix},\qquad
AP=\begin{pmatrix}1&0&0&-1\\-3&3&0&1\\0&3&3&0\\-4&4&0&4\end{pmatrix},\qquad
B=P^{-1}(AP)=\begin{pmatrix}1&0&0&-1\\-2&3&0&0\\2&0&3&0\\-1&1&0&3\end{pmatrix}
\]
(b)
\begin{itemize}
\item \(A\) is upper triangular, so \(p_A=(t-1)(t-2)(t-3)(t-4)\) (Proposition~\ref{r7-prop-block-triangular}(b)).
\item For \(t\in\R\), column \(3\) of \(tI-B\) has the entry \(t-3\) in position \((3,3)\) and \(0\) elsewhere; expanding
along it, and then along row \(1\),
\[
\begin{aligned}
\det(tI-B)&=\det\begin{pmatrix}t-1&0&0&1\\2&t-3&0&0\\-2&0&t-3&0\\1&-1&0&t-3\end{pmatrix}
=(t-3)\det\begin{pmatrix}t-1&0&1\\2&t-3&0\\1&-1&t-3\end{pmatrix}\\
&=(t-3)\bigl((t-1)(t-3)^2+(-2-(t-3))\bigr)=(t-3)\bigl((t-1)(t-3)^2-(t-1)\bigr)
\end{aligned}
\]
So \(p_B=(t-3)(t-1)\bigl((t-3)^2-1\bigr)=(t-1)(t-2)(t-3)(t-4)=p_A\), as Theorem~\ref{r7-thm-similar}(a) asserts.
\item \(\tr B=1+3+3+3=10=\tr A\). Expanding \(\det B\) along column \(3\) and then along row \(1\),
\[
\det B=3\det\begin{pmatrix}1&0&-1\\-2&3&0\\-1&1&3\end{pmatrix}=3\bigl(9-(-2+3)\bigr)=24
\]
and \(24=1\cdot2\cdot3\cdot4=\det A\) (Theorem~\ref{r5d-thm-basic}(c)).
\end{itemize}
(c) By Theorem~\ref{r7-thm-similar}(b), \(P^{-1}x=(1,2,-2,-1)'\) is an eigenvector of \(B\) belonging to \(2\). Check:
\(B(1,2,-2,-1)'=(2,4,-4,-2)'=2(1,2,-2,-1)'\).

(d)
\begin{itemize}
\item \(F\) and \(G\) are upper triangular with the diagonal entries \(1\), \(1\), \(1\), \(2\), so
\(p_F=p_G=(t-1)^3(t-2)\) (Proposition~\ref{r7-prop-block-triangular}(b)).
\item Row-reduce \(F-I\) and \(G-I\):
\[
F-I=\begin{pmatrix}0&1&1&1\\0&0&1&1\\0&0&0&1\\0&0&0&1\end{pmatrix}
\xrightarrow{R_4\leftarrow R_4-R_3}
\begin{pmatrix}0&1&1&1\\0&0&1&1\\0&0&0&1\\0&0&0&0\end{pmatrix}
\]
\[
G-I=\begin{pmatrix}0&0&1&1\\0&0&1&1\\0&0&0&1\\0&0&0&1\end{pmatrix}
\xrightarrow[R_4\leftarrow R_4-R_3]{R_2\leftarrow R_2-R_1}
\begin{pmatrix}0&0&1&1\\0&0&0&0\\0&0&0&1\\0&0&0&0\end{pmatrix}
\xrightarrow{R_2\leftrightarrow R_3}
\begin{pmatrix}0&0&1&1\\0&0&0&1\\0&0&0&0\\0&0&0&0\end{pmatrix}
\]
So \(\rank(F-I)=3\) and \(\rank(G-I)=2\) (Definition~\ref{def:r4-rank}), that is, \(\dim\mathcal N(F-I)=1\) and
\(\dim\mathcal N(G-I)=2\) (Theorem~\ref{thm:r4-rank-nullity}).
\item Similar matrices have eigenspaces of equal dimension (Theorem~\ref{r7-thm-similar}(b)), so \(F\) and \(G\) are not
similar.
\end{itemize}
\end{solution}

\begin{exercise}\label{r7-ex-similar}
Let \(B\) be the matrix of Example~\ref{r7-exm-diag-powers}, and
\[
F=\begin{pmatrix}3&0&0&0\\0&3&2&-2\\0&0&1&0\\0&0&0&1\end{pmatrix},\qquad
G=\begin{pmatrix}3&0&0&0\\0&3&2&-2\\0&0&1&1\\0&0&0&1\end{pmatrix}
\]
\begin{enumerate}
\item Show that \(p_B=p_F=p_G\).
\item Diagonalize \(F\), and find an invertible \(S\) with \(F=S^{-1}BS\).
\item Show that \(G\) is not similar to \(B\).
\end{enumerate}
\end{exercise}
\begin{solution}
(a) By Example~\ref{r7-exm-diag-powers}, the eigenvalues of \(B\), listed with multiplicity, are \(3\), \(3\), \(1\),
\(1\), so \(p_B=(t-3)^2(t-1)^2\) (Definition~\ref{r7-def-eigenvalues}). \(F\) and \(G\) are upper triangular with the
diagonal entries \(3\), \(3\), \(1\), \(1\), so \(p_F=p_G=(t-3)^2(t-1)^2\) (Proposition~\ref{r7-prop-block-triangular}(b)).

(b)
\begin{itemize}
\item \(\lambda=3\):
\[
\begin{gathered}
F-3I_4=\begin{pmatrix}0&0&0&0\\0&0&2&-2\\0&0&-2&0\\0&0&0&-2\end{pmatrix}
\xrightarrow{R_1\leftrightarrow R_3}
\begin{pmatrix}0&0&-2&0\\0&0&2&-2\\0&0&0&0\\0&0&0&-2\end{pmatrix}
\xrightarrow{R_2\leftarrow R_2+R_1}
\begin{pmatrix}0&0&-2&0\\0&0&0&-2\\0&0&0&0\\0&0&0&-2\end{pmatrix}
\\
\xrightarrow{R_4\leftarrow R_4-R_2}
\begin{pmatrix}0&0&-2&0\\0&0&0&-2\\0&0&0&0\\0&0&0&0\end{pmatrix}
\xrightarrow[R_2\leftarrow-\frac12R_2]{R_1\leftarrow-\frac12R_1}
\begin{pmatrix}0&0&1&0\\0&0&0&1\\0&0&0&0\\0&0&0&0\end{pmatrix}
\end{gathered}
\]
The free unknowns are \(x_1\), \(x_2\); the special solutions \(e_1\), \(e_2\) form a basis of \(\mathcal N(F-3I_4)\)
(Theorem~\ref{thm:r4-four-bases}(c)).
\item \(\lambda=1\):
\[
F-I_4=\begin{pmatrix}2&0&0&0\\0&2&2&-2\\0&0&0&0\\0&0&0&0\end{pmatrix}
\xrightarrow[R_2\leftarrow\frac12R_2]{R_1\leftarrow\frac12R_1}
\begin{pmatrix}1&0&0&0\\0&1&1&-1\\0&0&0&0\\0&0&0&0\end{pmatrix}
\]
The free unknowns are \(x_3\), \(x_4\); the special solutions \((0,-1,1,0)'\) and \((0,1,0,1)'\) form a basis of
\(\mathcal N(F-I_4)\).
\item So \(g_F(3)=a_F(3)=2\) and \(g_F(1)=a_F(1)=2\), and \(P_F^{-1}FP_F=\Lambda=\diag(3,3,1,1)\) (the paragraph after
Theorem~\ref{r7-thm-diag-criterion}), where
\[
P_F=\begin{pmatrix}1&0&0&0\\0&1&-1&1\\0&0&1&0\\0&0&0&1\end{pmatrix},\qquad
P_F^{-1}=\begin{pmatrix}1&0&0&0\\0&1&1&-1\\0&0&1&0\\0&0&0&1\end{pmatrix}
\]
Check: row \(2\) of the second matrix times columns \(1\), \(2\), \(3\), \(4\) of \(P_F\) gives \(0\), \(1\), \(-1+1=0\),
\(1-1=0\), and the other rows are rows of \(I_4\) times \(P_F\); so the product is \(I_4\), and the second matrix is
\(P_F^{-1}\) (Theorem~\ref{r2-thm-invertible}).
\item With \(P\) of Example~\ref{r7-exm-diag-powers}, \(P^{-1}BP=\Lambda\) with the same \(\Lambda\), so
\(F=P_F\Lambda P_F^{-1}=P_FP^{-1}BPP_F^{-1}=S^{-1}BS\) for \(S=PP_F^{-1}\) (Theorem~\ref{r2-thm-inverse-props}(a),(b)):
\[
S=PP_F^{-1}=\begin{pmatrix}-1&0&0&-1\\0&-1&0&0\\1&0&1&0\\0&1&1&0\end{pmatrix},\qquad
BS=SF=\begin{pmatrix}-3&0&0&-1\\0&-3&-2&2\\3&0&1&0\\0&3&3&-2\end{pmatrix}
\]
The equality \(BS=SF\) confirms \(F=S^{-1}BS\).
\end{itemize}
(c) Scaling rows \(1\) and \(2\) of
\[
G-I_4=\begin{pmatrix}2&0&0&0\\0&2&2&-2\\0&0&0&1\\0&0&0&0\end{pmatrix}
\]
by \(\frac12\) gives an echelon form with three nonzero rows, so \(\rank(G-I_4)=3\) (Remark~\ref{rem:r4-echelon},
Definition~\ref{def:r4-rank}) and \(g_G(1)=4-3=1\) (Definition~\ref{r7-def-multiplicity}). As \(g_B(1)=2\)
(Example~\ref{r7-exm-diag-powers}), Theorem~\ref{r7-thm-similar}(b) shows that \(G\) is not similar to \(B\); by (a), the
characteristic polynomial does not detect this.
\end{solution}

\begin{exercise}\label{r7-dr-multiplicities}
For each matrix, find the eigenvalues with their algebraic and geometric multiplicities, decide whether the matrix is
diagonalizable, and compare its rank with the number of its nonzero eigenvalues, listed with multiplicity.
\[
\text{(a)}\ \ A=\begin{pmatrix}2&1&-1&0\\0&2&3&1\\0&0&-1&-1\\0&0&0&3\end{pmatrix},\qquad
\text{(b)}\ \ B=\begin{pmatrix}-1&0&-1&-2\\1&3&-2&-1\\1&0&1&2\\-1&0&-1&-1\end{pmatrix}
\]
\[
\text{(c)}\ \ M=\begin{pmatrix}3&1&1&2\\0&2&0&0\\-1&-1&1&-2\\0&0&0&2\end{pmatrix}
\]
\end{exercise}
\begin{solution}
(a)
\begin{itemize}
\item \(A\) is upper triangular, so \(p_A=(t-2)^2(t+1)(t-3)\) (Proposition~\ref{r7-prop-block-triangular}(b)):
\(a_A(2)=2\) and \(a_A(-1)=a_A(3)=1\). As \(-1\) and \(3\) are simple, \(g_A(-1)=g_A(3)=1\)
(Theorem~\ref{r7-thm-geo-alg}).
\item Row-reduce \(A-2I\):
\[
A-2I=\begin{pmatrix}0&1&-1&0\\0&0&3&1\\0&0&-3&-1\\0&0&0&1\end{pmatrix}
\xrightarrow{R_3\leftarrow R_3+R_2}
\begin{pmatrix}0&1&-1&0\\0&0&3&1\\0&0&0&0\\0&0&0&1\end{pmatrix}
\xrightarrow{R_3\leftrightarrow R_4}
\begin{pmatrix}0&1&-1&0\\0&0&3&1\\0&0&0&1\\0&0&0&0\end{pmatrix}
\]
Scaling row \(2\) by \(\frac13\) gives an echelon form with three nonzero rows, so \(\rank(A-2I)=3\)
(Remark~\ref{rem:r4-echelon}, Definition~\ref{def:r4-rank}) and \(g_A(2)=4-3=1<2=a_A(2)\)
(Definition~\ref{r7-def-multiplicity}): \(A\) is not diagonalizable (Theorem~\ref{r7-thm-diag-criterion}).
\item \(\det A=2\cdot2\cdot(-1)\cdot3=-12\ne0\) (Theorem~\ref{r5d-thm-basic}(c)), so \(\rank A=4\)
(Theorem~\ref{r5d-thm-product}(a), Theorem~\ref{thm:r4-rank-rules}(d)): all four eigenvalues \(2\), \(2\), \(-1\), \(3\)
are nonzero, and their number is \(\rank A\), although \(A\) is not diagonalizable.
\end{itemize}
(b)
\begin{itemize}
\item For \(t\in\R\), column \(2\) of \(tI-B\) has the entry \(t-3\) in position \((2,2)\) and \(0\) elsewhere; expanding
along it, and then along row \(1\),
\[
\begin{aligned}
\det(tI-B)&=\det\begin{pmatrix}t+1&0&1&2\\-1&t-3&2&1\\-1&0&t-1&-2\\1&0&1&t+1\end{pmatrix}
=(t-3)\det\begin{pmatrix}t+1&1&2\\-1&t-1&-2\\1&1&t+1\end{pmatrix}\\
&=(t-3)\Bigl((t+1)\bigl((t-1)(t+1)+2\bigr)-\bigl(-(t+1)+2\bigr)+2\bigl(-1-(t-1)\bigr)\Bigr)\\
&=(t-3)\bigl(t^3+t^2+t+1-1+t-2t\bigr)=(t-3)(t^3+t^2)=(t-3)t^2(t+1)
\end{aligned}
\]
So \(a_B(0)=2\) and \(a_B(3)=a_B(-1)=1\), and \(g_B(3)=g_B(-1)=1\) (Theorem~\ref{r7-thm-geo-alg}).
\item Row-reduce \(B-0I=B\):
\[
\begin{gathered}
B=\begin{pmatrix}-1&0&-1&-2\\1&3&-2&-1\\1&0&1&2\\-1&0&-1&-1\end{pmatrix}
\xrightarrow[R_4\leftarrow R_4-R_1]{R_2\leftarrow R_2+R_1,\ R_3\leftarrow R_3+R_1}
\begin{pmatrix}-1&0&-1&-2\\0&3&-3&-3\\0&0&0&0\\0&0&0&1\end{pmatrix}
\\
\xrightarrow{R_3\leftrightarrow R_4}
\begin{pmatrix}-1&0&-1&-2\\0&3&-3&-3\\0&0&0&1\\0&0&0&0\end{pmatrix}
\end{gathered}
\]
Scaling rows \(1\) and \(2\) by \(-1\) and \(\frac13\) gives an echelon form with three nonzero rows, so \(\rank B=3\)
and \(g_B(0)=4-3=1<2=a_B(0)\): \(B\) is not diagonalizable.
\item \(\rank B=3\), but only two eigenvalues, \(3\) and \(-1\), are nonzero: the inequality of
Corollary~\ref{r7-cor-rank-eigen} is strict.
\end{itemize}
(c)
\begin{itemize}
\item For \(t\in\R\), rows \(2\) and \(4\) of \(tI-M\) have the entry \(t-2\) on the diagonal and \(0\) elsewhere;
expanding along row \(2\), and then along the last row,
\[
\begin{aligned}
\det(tI-M)&=\det\begin{pmatrix}t-3&-1&-1&-2\\0&t-2&0&0\\1&1&t-1&2\\0&0&0&t-2\end{pmatrix}
=(t-2)\det\begin{pmatrix}t-3&-1&-2\\1&t-1&2\\0&0&t-2\end{pmatrix}\\
&=(t-2)^2\det\begin{pmatrix}t-3&-1\\1&t-1\end{pmatrix}=(t-2)^2(t^2-4t+4)=(t-2)^4
\end{aligned}
\]
So \(2\) is the only eigenvalue, with \(a_M(2)=4\).
\item Row-reduce \(M-2I\):
\[
M-2I=\begin{pmatrix}1&1&1&2\\0&0&0&0\\-1&-1&-1&-2\\0&0&0&0\end{pmatrix}
\xrightarrow{R_3\leftarrow R_3+R_1}
\begin{pmatrix}1&1&1&2\\0&0&0&0\\0&0&0&0\\0&0&0&0\end{pmatrix}
\]
So \(\rank(M-2I)=1\) and \(g_M(2)=3<4=a_M(2)\): \(M\) is not diagonalizable.
\item \(\det M=(-1)^4p_M(0)=16\ne0\), so \(\rank M=4\), the number of the nonzero eigenvalues \(2\), \(2\), \(2\), \(2\).
\end{itemize}
\end{solution}

\begin{exercise}[\citeat{abadir-magnus}{Exercises~7.3, 7.11}]\label{r7-ex-am-eigen}
Let
\[
A=\begin{pmatrix}1&3&0\\0&1&0\\2&1&5\end{pmatrix}
\]
\begin{enumerate}
\item Find the characteristic equation and the eigenvalues of \(A\).
\item What is the rank of \(A\)?
\item What are the eigenvectors of \(A\)? Are they linearly independent? Are they orthogonal?
\end{enumerate}
\end{exercise}
\begin{solution}
(a) For \(t\in\R\), column \(3\) of \(tI-A\) has the entry \(t-5\) at \((3,3)\) and \(0\) elsewhere; expanding along
it (Theorem~\ref{r5d-thm-expansion}),
\[
\det(tI-A)=\det\begin{pmatrix}t-1&-3&0\\0&t-1&0\\-2&-1&t-5\end{pmatrix}
=(t-5)\det\begin{pmatrix}t-1&-3\\0&t-1\end{pmatrix}=(t-1)^2(t-5)
\]
The characteristic equation is \((t-1)^2(t-5)=0\), and the eigenvalues are \(1\) (twice) and \(5\).

(b) The constant term of \(p_A\) is \((0-1)^2(0-5)=-5=(-1)^3\det A\) (Definition~\ref{r7-def-charpoly}(b)), so
\(\det A=5\ne0\): \(A\) is invertible and \(\rank A=3\) (Theorem~\ref{r5d-thm-product}(a)).

(c)
\begin{itemize}
\item \(\lambda=1\): \(Ax=x\) reads \(x_1+3x_2=x_1\), \(x_2=x_2\), \(2x_1+x_2+5x_3=x_3\); so \(x_2=0\) and \(x_1+2x_3=0\): the
eigenvectors are the nonzero multiples of \((2,0,-1)'\).
\item \(\lambda=5\): \(Ax=5x\) reads \(x_1+3x_2=5x_1\), \(x_2=5x_2\), \(2x_1+x_2+5x_3=5x_3\); so \(x_2=0\), then \(x_1=0\), and
\(x_3\) is free: the eigenvectors are the nonzero multiples of \((0,0,1)'\).
\item \((2,0,-1)'\) and \((0,0,1)'\) are independent (Theorem~\ref{r7-thm-indep}) but not orthogonal: their inner product
(Definition~\ref{r6-def-inner}) is \(2\cdot0+0\cdot0+(-1)\cdot1=-1\). The double eigenvalue \(1\) has only one independent eigenvector: \(g_A(1)=1<2=a_A(1)\)
(Definitions~\ref{r7-def-eigenvalues} and~\ref{r7-def-multiplicity}), as Theorem~\ref{r7-thm-geo-alg} allows.
\end{itemize}
\end{solution}

\begin{exercise}\label{r7-ex-which-diagonalizable}
Which of the following matrices are diagonalizable?
\[
A_1=\begin{pmatrix}3&1&0&-1\\0&3&0&0\\0&1&3&0\\-1&0&1&3\end{pmatrix},\qquad
A_2=\begin{pmatrix}1&-2\\1&-1\end{pmatrix},\qquad
A_3=\begin{pmatrix}1&1&1&0\\0&2&1&1\\0&0&1&-1\\0&0&0&2\end{pmatrix}
\]
\end{exercise}
\begin{solution}
(\(A_1\)) \(A_1\) differs from the matrix \(B\) of Exercise~\ref{r7-ex-eigenpairs} only in the \((2,2)\) entry. The same
two expansions give \(p_{A_1}=(t-3)\cdot(t-3)(t-2)(t-4)=(t-3)^2(t-2)(t-4)\), so \(a_{A_1}(3)=2\). But
\[
\begin{gathered}
A_1-3I_4=\begin{pmatrix}0&1&0&-1\\0&0&0&0\\0&1&0&0\\-1&0&1&0\end{pmatrix}
\xrightarrow{R_1\leftrightarrow R_4}
\begin{pmatrix}-1&0&1&0\\0&0&0&0\\0&1&0&0\\0&1&0&-1\end{pmatrix}
\xrightarrow{R_2\leftrightarrow R_4}
\begin{pmatrix}-1&0&1&0\\0&1&0&-1\\0&1&0&0\\0&0&0&0\end{pmatrix}
\\
\xrightarrow{R_3\leftarrow R_3-R_2}
\begin{pmatrix}-1&0&1&0\\0&1&0&-1\\0&0&0&1\\0&0&0&0\end{pmatrix}
\end{gathered}
\]
and scaling row \(1\) by \(-1\) gives an echelon form with three nonzero rows; so \(\rank(A_1-3I_4)=3\)
(Remark~\ref{rem:r4-echelon}, Definition~\ref{def:r4-rank}), \(g_{A_1}(3)=1<2\) (Definition~\ref{r7-def-multiplicity}),
and \(A_1\) is not diagonalizable (Theorem~\ref{r7-thm-diag-criterion}). The matrix \(B\) has four distinct real
eigenvalues, so it is diagonalizable (Theorem~\ref{r7-thm-diag-criterion}); changing one entry has destroyed
diagonalizability.

(\(A_2\)) \(p_{A_2}=t^2-0\cdot t+(-1+2)=t^2+1=(t-i)(t+i)\) (Definition~\ref{r7-def-charpoly}(b)). So the eigenvalues
\(i\) and \(-i\) of \(A_2\) are not real, and \(A_2\) is not diagonalizable (Definition~\ref{r7-def-diagonalizable}).

(\(A_3\)) \(A_3\) is upper triangular, so \(p_{A_3}=(t-1)^2(t-2)^2\) (Proposition~\ref{r7-prop-block-triangular}(b)).
\[
A_3-I_4=\begin{pmatrix}0&1&1&0\\0&1&1&1\\0&0&0&-1\\0&0&0&1\end{pmatrix}
\xrightarrow{R_2\leftarrow R_2-R_1}
\begin{pmatrix}0&1&1&0\\0&0&0&1\\0&0&0&-1\\0&0&0&1\end{pmatrix}
\xrightarrow[R_4\leftarrow R_4-R_2]{R_3\leftarrow R_3+R_2}
\begin{pmatrix}0&1&1&0\\0&0&0&1\\0&0&0&0\\0&0&0&0\end{pmatrix}
\]
The free unknowns are \(x_1\), \(x_3\); the special solutions \((1,0,0,0)'\), \((0,-1,1,0)'\) form a basis of
\(\mathcal N(A_3-I_4)\) (Theorem~\ref{thm:r4-four-bases}(c)), and \(g_{A_3}(1)=2\).
\[
A_3-2I_4=\begin{pmatrix}-1&1&1&0\\0&0&1&1\\0&0&-1&-1\\0&0&0&0\end{pmatrix}
\xrightarrow[R_3\leftarrow R_3+R_2]{R_1\leftarrow R_1-R_2}
\begin{pmatrix}-1&1&0&-1\\0&0&1&1\\0&0&0&0\\0&0&0&0\end{pmatrix}
\xrightarrow{R_1\leftarrow-R_1}
\begin{pmatrix}1&-1&0&1\\0&0&1&1\\0&0&0&0\\0&0&0&0\end{pmatrix}
\]
The free unknowns are \(x_2\), \(x_4\); the special solutions \((1,1,0,0)'\), \((-1,0,-1,1)'\) form a basis of
\(\mathcal N(A_3-2I_4)\), and \(g_{A_3}(2)=2\). All eigenvalues of \(A_3\) are real and both geometric multiplicities
equal the algebraic ones, so \(A_3\) is diagonalizable (Theorem~\ref{r7-thm-diag-criterion}).
\end{solution}

\begin{exercise}\label{r7-ex-diag-drill}
Let
\[
A=\begin{pmatrix}2&0&-1&-1\\0&3&-2&-2\\2&0&2&-2\\-2&0&-1&3\end{pmatrix}
\]
\begin{enumerate}
\item Find the eigenvalues of \(A\), a basis of each eigenspace, and an invertible \(P\) with \(P^{-1}AP\) diagonal.
\item Find \(P^{-1}\) by Gauss--Jordan elimination.
\item Find \(A^k\) for \(k\geqslant1\), and check the result for \(k=1\).
\end{enumerate}
\end{exercise}
\begin{solution}
(a)
\begin{itemize}
\item For \(t\in\R\), column \(2\) of \(tI-A\) has the entry \(t-3\) in position \((2,2)\) and \(0\) elsewhere; expanding
along it, and then along row \(1\),
\[
\begin{aligned}
\det(tI-A)&=\det\begin{pmatrix}t-2&0&1&1\\0&t-3&2&2\\-2&0&t-2&2\\2&0&1&t-3\end{pmatrix}
=(t-3)\det\begin{pmatrix}t-2&1&1\\-2&t-2&2\\2&1&t-3\end{pmatrix}\\
&=(t-3)\Bigl((t-2)\bigl((t-2)(t-3)-2\bigr)-\bigl(-2(t-3)-4\bigr)+\bigl(-2-2(t-2)\bigr)\Bigr)\\
&=(t-3)\bigl((t-2)(t^2-5t+4)+(2t-2)+(2-2t)\bigr)\\
&=(t-3)(t-2)(t-1)(t-4)
\end{aligned}
\]
So \(p_A=(t-1)(t-2)(t-3)(t-4)\): the eigenvalues \(1\), \(2\), \(3\), \(4\) are distinct, and \(A\) is diagonalizable
(Theorem~\ref{r7-thm-diag-criterion}).
\item \(\lambda=1\):
\[
\begin{gathered}
A-I=\begin{pmatrix}1&0&-1&-1\\0&2&-2&-2\\2&0&1&-2\\-2&0&-1&2\end{pmatrix}
\xrightarrow[R_4\leftarrow R_4+2R_1]{R_3\leftarrow R_3-2R_1}
\begin{pmatrix}1&0&-1&-1\\0&2&-2&-2\\0&0&3&0\\0&0&-3&0\end{pmatrix}
\xrightarrow[R_3\leftarrow\frac13R_3]{R_2\leftarrow\frac12R_2}
\begin{pmatrix}1&0&-1&-1\\0&1&-1&-1\\0&0&1&0\\0&0&-3&0\end{pmatrix}
\\
\xrightarrow[R_4\leftarrow R_4+3R_3]{R_1\leftarrow R_1+R_3,\ R_2\leftarrow R_2+R_3}
\begin{pmatrix}1&0&0&-1\\0&1&0&-1\\0&0&1&0\\0&0&0&0\end{pmatrix}
\end{gathered}
\]
The special solution \(p_1=(1,1,0,1)'\) is a basis of \(\mathcal N(A-I)\) (Theorem~\ref{thm:r4-four-bases}(c)).
\item \(\lambda=2\):
\[
\begin{gathered}
A-2I=\begin{pmatrix}0&0&-1&-1\\0&1&-2&-2\\2&0&0&-2\\-2&0&-1&1\end{pmatrix}
\xrightarrow{R_1\leftrightarrow R_3}
\begin{pmatrix}2&0&0&-2\\0&1&-2&-2\\0&0&-1&-1\\-2&0&-1&1\end{pmatrix}
\xrightarrow{R_4\leftarrow R_4+R_1}
\begin{pmatrix}2&0&0&-2\\0&1&-2&-2\\0&0&-1&-1\\0&0&-1&-1\end{pmatrix}
\\
\xrightarrow[R_4\leftarrow R_4-R_3]{R_2\leftarrow R_2-2R_3}
\begin{pmatrix}2&0&0&-2\\0&1&0&0\\0&0&-1&-1\\0&0&0&0\end{pmatrix}
\xrightarrow[R_3\leftarrow-R_3]{R_1\leftarrow\frac12R_1}
\begin{pmatrix}1&0&0&-1\\0&1&0&0\\0&0&1&1\\0&0&0&0\end{pmatrix}
\end{gathered}
\]
The special solution \(p_2=(1,0,-1,1)'\) is a basis of \(\mathcal N(A-2I)\).
\item \(\lambda=3\):
\[
\begin{gathered}
A-3I=\begin{pmatrix}-1&0&-1&-1\\0&0&-2&-2\\2&0&-1&-2\\-2&0&-1&0\end{pmatrix}
\xrightarrow[R_4\leftarrow R_4-2R_1]{R_3\leftarrow R_3+2R_1}
\begin{pmatrix}-1&0&-1&-1\\0&0&-2&-2\\0&0&-3&-4\\0&0&1&2\end{pmatrix}
\xrightarrow{R_2\leftarrow-\frac12R_2}
\begin{pmatrix}-1&0&-1&-1\\0&0&1&1\\0&0&-3&-4\\0&0&1&2\end{pmatrix}
\\
\xrightarrow[R_4\leftarrow R_4-R_2]{R_1\leftarrow R_1+R_2,\ R_3\leftarrow R_3+3R_2}
\begin{pmatrix}-1&0&0&0\\0&0&1&1\\0&0&0&-1\\0&0&0&1\end{pmatrix}
\xrightarrow[R_4\leftarrow R_4+R_3]{R_2\leftarrow R_2+R_3}
\begin{pmatrix}-1&0&0&0\\0&0&1&0\\0&0&0&-1\\0&0&0&0\end{pmatrix}
\\
\xrightarrow[R_3\leftarrow-R_3]{R_1\leftarrow-R_1}
\begin{pmatrix}1&0&0&0\\0&0&1&0\\0&0&0&1\\0&0&0&0\end{pmatrix}
\end{gathered}
\]
The free unknown is \(x_2\); the special solution \(p_3=e_2\) is a basis of \(\mathcal N(A-3I)\).
\item \(\lambda=4\):
\[
\begin{gathered}
A-4I=\begin{pmatrix}-2&0&-1&-1\\0&-1&-2&-2\\2&0&-2&-2\\-2&0&-1&-1\end{pmatrix}
\xrightarrow{R_1\leftrightarrow R_3}
\begin{pmatrix}2&0&-2&-2\\0&-1&-2&-2\\-2&0&-1&-1\\-2&0&-1&-1\end{pmatrix}
\xrightarrow[R_4\leftarrow R_4+R_1]{R_3\leftarrow R_3+R_1}
\begin{pmatrix}2&0&-2&-2\\0&-1&-2&-2\\0&0&-3&-3\\0&0&-3&-3\end{pmatrix}
\\
\xrightarrow[R_3\leftarrow-\frac13R_3]{R_2\leftarrow-R_2}
\begin{pmatrix}2&0&-2&-2\\0&1&2&2\\0&0&1&1\\0&0&-3&-3\end{pmatrix}
\xrightarrow[R_4\leftarrow R_4+3R_3]{R_1\leftarrow R_1+2R_3,\ R_2\leftarrow R_2-2R_3}
\begin{pmatrix}2&0&0&0\\0&1&0&0\\0&0&1&1\\0&0&0&0\end{pmatrix}
\\
\xrightarrow{R_1\leftarrow\frac12R_1}
\begin{pmatrix}1&0&0&0\\0&1&0&0\\0&0&1&1\\0&0&0&0\end{pmatrix}
\end{gathered}
\]
The special solution \(p_4=(0,0,-1,1)'\) is a basis of \(\mathcal N(A-4I)\).
\item The eigenvectors \(p_1,\dots,p_4\) are independent (Theorem~\ref{r7-thm-indep}). With
\[
P=[p_1\ p_2\ p_3\ p_4]=\begin{pmatrix}1&1&0&0\\1&0&1&0\\0&-1&0&-1\\1&1&0&1\end{pmatrix}
\]
\(P^{-1}AP=\diag(1,2,3,4)\) (Definition~\ref{r7-def-diagonalizable}).
\end{itemize}
(b) Row-reduce \([P\mid I_4]\) (Corollary~\ref{r2-cor-computing-inverse}(a)):
\[
\begin{gathered}
\left(\begin{array}{cccc|cccc}1&1&0&0&1&0&0&0\\1&0&1&0&0&1&0&0\\0&-1&0&-1&0&0&1&0\\1&1&0&1&0&0&0&1\end{array}\right)
\xrightarrow[R_4\leftarrow R_4-R_1]{R_2\leftarrow R_2-R_1}
\left(\begin{array}{cccc|cccc}1&1&0&0&1&0&0&0\\0&-1&1&0&-1&1&0&0\\0&-1&0&-1&0&0&1&0\\0&0&0&1&-1&0&0&1\end{array}\right)
\\
\xrightarrow{R_2\leftrightarrow R_3}
\left(\begin{array}{cccc|cccc}1&1&0&0&1&0&0&0\\0&-1&0&-1&0&0&1&0\\0&-1&1&0&-1&1&0&0\\0&0&0&1&-1&0&0&1\end{array}\right)
\\
\xrightarrow[R_3\leftarrow R_3-R_2]{R_1\leftarrow R_1+R_2}
\left(\begin{array}{cccc|cccc}1&0&0&-1&1&0&1&0\\0&-1&0&-1&0&0&1&0\\0&0&1&1&-1&1&-1&0\\0&0&0&1&-1&0&0&1\end{array}\right)
\\
\xrightarrow[R_3\leftarrow R_3-R_4]{R_1\leftarrow R_1+R_4,\ R_2\leftarrow R_2+R_4}
\left(\begin{array}{cccc|cccc}1&0&0&0&0&0&1&1\\0&-1&0&0&-1&0&1&1\\0&0&1&0&0&1&-1&-1\\0&0&0&1&-1&0&0&1\end{array}\right)
\\
\xrightarrow{R_2\leftarrow-R_2}
\left(\begin{array}{cccc|cccc}1&0&0&0&0&0&1&1\\0&1&0&0&1&0&-1&-1\\0&0&1&0&0&1&-1&-1\\0&0&0&1&-1&0&0&1\end{array}\right)
\end{gathered}
\]
So
\[
P^{-1}=\begin{pmatrix}0&0&1&1\\1&0&-1&-1\\0&1&-1&-1\\-1&0&0&1\end{pmatrix}
\]
(c) By Proposition~\ref{r7-prop-diag-powers}(a), \(A^k=P\diag(1,2^k,3^k,4^k)P^{-1}\). With \(r_i'\) the \(i\)th row of
\(P^{-1}\), the outer product expansion (Theorem~\ref{r1-thm-views}(d)) gives
\(A^k=p_1r_1'+2^kp_2r_2'+3^kp_3r_3'+4^kp_4r_4'\), where
\[
\begin{gathered}
p_1r_1'=\begin{pmatrix}0&0&1&1\\0&0&1&1\\0&0&0&0\\0&0&1&1\end{pmatrix},\qquad
p_2r_2'=\begin{pmatrix}1&0&-1&-1\\0&0&0&0\\-1&0&1&1\\1&0&-1&-1\end{pmatrix}\\
p_3r_3'=\begin{pmatrix}0&0&0&0\\0&1&-1&-1\\0&0&0&0\\0&0&0&0\end{pmatrix},\qquad
p_4r_4'=\begin{pmatrix}0&0&0&0\\0&0&0&0\\1&0&0&-1\\-1&0&0&1\end{pmatrix}
\end{gathered}
\]
So
\[
A^k=\begin{pmatrix}2^k&0&1-2^k&1-2^k\\0&3^k&1-3^k&1-3^k\\4^k-2^k&0&2^k&2^k-4^k\\2^k-4^k&0&1-2^k&1-2^k+4^k\end{pmatrix}
\qquad(k\geqslant1)
\]
For \(k=1\) this is
\[
\begin{pmatrix}2&0&-1&-1\\0&3&-2&-2\\2&0&2&-2\\-2&0&-1&3\end{pmatrix}=A
\]
\end{solution}

\begin{exercise}\label{r7-ex-diagonalize}
Let
\[
A=\begin{pmatrix}2&-2&1&-1\\-1&1&1&-1\\0&0&3&0\\-1&-2&1&2\end{pmatrix}
\]
Show that \(\rank(A-3I_4)=1\), diagonalize \(A\), find a formula for \(A^k\) (\(k\geqslant0\)), and compute \(A^3\).
\end{exercise}
\begin{solution}
\begin{itemize}
\item Subtracting \(3\) from the diagonal entries of \(A\),
\[
A-3I_4=\begin{pmatrix}-1&-2&1&-1\\-1&-2&1&-1\\0&0&0&0\\-1&-2&1&-1\end{pmatrix}=uv'
\]
with \(u=(1,1,0,1)'\) and \(v'=(-1,-2,1,-1)\): rows \(1\), \(2\) and \(4\) equal \(v'\), and row
\(3\) is \(0\). So the row space of \(A-3I_4\) is \(\Span\{v'\}\), of dimension \(1\) as \(v'\ne0\)
(Proposition~\ref{r3-prop-dependence}(a)), and \(\rank(A-3I_4)=1\) (Definition~\ref{def:r4-rank}). Hence
\(g_A(3)=4-1=3\) (Definition~\ref{r7-def-multiplicity}) and \(a_A(3)\geqslant3\) (Theorem~\ref{r7-thm-geo-alg}).
\item So at least three of the four eigenvalues of \(A\), listed with multiplicity, are \(3\)
(Definition~\ref{r7-def-eigenvalues}); let \(\mu\) be the fourth. By Theorem~\ref{r7-thm-trace-det},
\(3+3+3+\mu=\tr A=2+1+3+2=8\), so \(\mu=-1\) and \(p_A=(t-3)^3(t+1)\).
\item \(\mathcal N(A-3I_4)=\mathcal N(uv')=\mathcal N(v')\), as \(uv'x=(v'x)u\) and \(u\ne0\). The equation \(v'x=0\), that
is, \(x_1+2x_2-x_3+x_4=0\), with \(x_2,x_3,x_4\) free, has the special solutions \(p_1=(-2,1,0,0)'\),
\(p_2=(1,0,1,0)'\) and \(p_3=(-1,0,0,1)'\), a basis of \(\mathcal N(A-3I_4)\) (Theorem~\ref{thm:r4-four-bases}(c)).
Next \(Au=3u+u(v'u)=3u-4u=-u\), so \(p_4=u\) spans \(\mathcal N(A+I_4)\) (Theorem~\ref{r7-thm-geo-alg}:
\(g_A(-1)\leqslant a_A(-1)=1\)).
\item By Lemma~\ref{r7-lem-eigenspaces}(b) and Theorem~\ref{thm:r4-counting}(d), \(p_1,\dots,p_4\) form a basis of
\(\R^4\), so \(P^{-1}AP=\Lambda\) with \(P=[p_1\ p_2\ p_3\ p_4]\) and \(\Lambda=\diag(3,3,3,-1)\)
(Definition~\ref{r7-def-diagonalizable}). The row \(w'=-\frac14v'=\frac14(1,2,-1,1)\) satisfies \(w'p_j=0\) for
\(j\leqslant3\) (as \(v'p_j=0\)) and \(w'p_4=-\frac14v'u=1\); so \(w'P=e_4'\), and \(w'=e_4'P^{-1}=\row_4(P^{-1})\).
\end{itemize}
By Proposition~\ref{r7-prop-diag-powers}(a),
\[
A^k=P\diag\bigl(3^k,3^k,3^k,(-1)^k\bigr)P^{-1}=3^kM_1+(-1)^kM_2
\]
with \(M_1=P\diag(1,1,1,0)P^{-1}\) and \(M_2=P\diag(0,0,0,1)P^{-1}\). By the outer product expansion
(Theorem~\ref{r1-thm-views}(d)), \(M_2=p_4\row_4(P^{-1})=-\frac14uv'\), and \(M_1=I_4-M_2\), as \(M_1+M_2=PP^{-1}\). So
\[
A^k=3^kI_4+\tfrac14\bigl(3^k-(-1)^k\bigr)uv'\qquad(k\geqslant0)
\]
For \(k=1\) this gives back \(A=3I_4+uv'\), and
\[
A^3=27I_4+7uv'=\begin{pmatrix}20&-14&7&-7\\-7&13&7&-7\\0&0&27&0\\-7&-14&7&20\end{pmatrix}
\]
\end{solution}

\begin{exercise}\label{r7-ex-powers}
Let
\[
A=\begin{pmatrix}0&2&-1&0\\1&1&-1&0\\-2&2&1&0\\1&-1&0&-2\end{pmatrix},\qquad
p_1=\begin{pmatrix}1\\1\\0\\0\end{pmatrix},\quad p_2=\begin{pmatrix}1\\1\\1\\0\end{pmatrix},\quad
p_3=\begin{pmatrix}1\\0\\1\\1\end{pmatrix},\quad p_4=\begin{pmatrix}0\\0\\0\\1\end{pmatrix}
\]
\begin{enumerate}
\item Show that \(p_1,\dots,p_4\) are eigenvectors of \(A\), and find \(p_A\), \(\tr A\) and \(\det A\).
\item Diagonalize \(A\) and compute \(A^{-1}\).
\item For \(x_0=(3,2,2,1)'\), find \(A^kx_0\) (\(k\geqslant0\)), and check \(A^3x_0\) by computing \(Ax_0\), \(A^2x_0\),
\(A^3x_0\) in turn.
\end{enumerate}
\end{exercise}
\begin{solution}
(a) Multiplying out, \(Ap_1=(2,2,0,0)'=2p_1\), \(Ap_2=(1,1,1,0)'=p_2\), \(Ap_3=(-1,0,-1,-1)'=-p_3\) and
\(Ap_4=\col_4(A)=(0,0,0,-2)'=-2p_4\). The eigenvalues \(2\), \(1\), \(-1\), \(-2\) are distinct, so \(p_1,\dots,p_4\)
form a basis of \(\R^4\) (Theorem~\ref{r7-thm-indep}), and \(P^{-1}AP=\Lambda\) with \(P=[p_1\ p_2\ p_3\ p_4]\) and
\(\Lambda=\diag(2,1,-1,-2)\) (Definition~\ref{r7-def-diagonalizable}). By the same definition, \(2\), \(1\), \(-1\), \(-2\)
are the eigenvalues of \(A\), listed with multiplicity; so \(p_A=(t-2)(t-1)(t+1)(t+2)\)
(Definition~\ref{r7-def-eigenvalues}), \(\tr A=2+1-1-2=0\) and \(\det A=2\cdot1\cdot(-1)\cdot(-2)=4\)
(Theorem~\ref{r7-thm-trace-det}). Directly, \(\tr A=0+1+1-2=0\).

(b)
\[
P^{-1}=\begin{pmatrix}1&0&-1&0\\-1&1&1&0\\1&-1&0&0\\-1&1&0&1\end{pmatrix}
\]
Check: by Theorem~\ref{r1-thm-views}(d), the \((i,j)\) entry of the product of this matrix and \(P\) is its row \(i\) times
\(p_j\). For row \(1\), \((1,0,-1,0)p_j\) is \(1\), \(1-1=0\), \(1-1=0\), \(0\) (\(j=1,2,3,4\)); for row \(2\),
\((-1,1,1,0)p_j\) is \(-1+1=0\), \(-1+1+1=1\), \(-1+1=0\), \(0\); for row \(3\), \((1,-1,0,0)p_j\) is \(1-1=0\),
\(1-1=0\), \(1\), \(0\); for row \(4\), \((-1,1,0,1)p_j\) is \(-1+1=0\), \(-1+1=0\), \(-1+1=0\), \(1\). So the product
is \(I_4\), and the matrix is \(P^{-1}\) (Theorem~\ref{r2-thm-invertible}). By
Proposition~\ref{r7-prop-diag-powers}(c), \(A^{-1}=P\Lambda^{-1}P^{-1}=(P\Lambda^{-1})P^{-1}\) with
\(\Lambda^{-1}=\diag(\frac12,1,-1,-\frac12)\); the columns of \(P\Lambda^{-1}\) are those of \(P\) times \(\frac12\), \(1\), \(-1\),
\(-\frac12\) (Proposition~\ref{r1-prop-diag}(a)):
\[
P\Lambda^{-1}=\frac12\begin{pmatrix}1&2&-2&0\\1&2&0&0\\0&2&-2&0\\0&0&-2&-1\end{pmatrix},\qquad
A^{-1}=\frac12\begin{pmatrix}-3&4&1&0\\-1&2&1&0\\-4&4&2&0\\-1&1&0&-1\end{pmatrix}
\]
Check: \(\row_1(A)\col_1(A^{-1})=\frac12(-2+4)=1\) and \(\row_1(A)\col_2(A^{-1})=\frac12(4-4)=0\).

(c) \(c=P^{-1}x_0=(1,1,1,0)'\), that is, \(x_0=p_1+p_2+p_3\). By Proposition~\ref{r7-prop-diag-powers}(a),
\[
A^kx_0=2^kp_1+p_2+(-1)^kp_3=\bigl(2^k+1+(-1)^k,\ 2^k+1,\ 1+(-1)^k,\ (-1)^k\bigr)'
\]
For \(k=3\) this is \((8,9,0,-1)'\). Directly, \(Ax_0=(2,3,0,-1)'\), \(A^2x_0=A(2,3,0,-1)'=(6,5,2,1)'\) and
\(A^3x_0=A(6,5,2,1)'=(8,9,0,-1)'\).
\end{solution}

\begin{exercise}\label{r7-dr-related}
The matrix \(A\) of Exercise~\ref{r7-dr-eigenspaces} has the eigenvalues \(-1\), \(1\), \(2\), \(3\), with the eigenvectors
\((1,0,0,1)'\), \((1,0,1,1)'\), \((1,0,1,0)'\), \((-1,-1,-1,1)'\). Without computing the matrices, find the eigenvalues, the
trace and the determinant of
\begin{enumerate}
\item \(A-2I\) and \(-3A\);
\item \(A^2\), \(A^{-1}\) and \(A^3-4A\);
\item \(A'\).
\item Give an eigenvector of \(A^3-4A\) for each of its eigenvalues, and check \(\tr(A^2)\) by computing the diagonal
entries of \(A^2\).
\end{enumerate}
\end{exercise}
\begin{solution}
(a) By Proposition~\ref{r7-prop-eig-poly}(c), \(A-2I\) has the eigenvalues \(-3\), \(-1\), \(0\), \(1\), so
\(\tr(A-2I)=-3\) and \(\det(A-2I)=0\) (Theorem~\ref{r7-thm-trace-det}); and \(-3A\) has the eigenvalues \(3\), \(-3\),
\(-6\), \(-9\), so \(\tr(-3A)=-15\) and \(\det(-3A)=-486\).

(b) \(A=P\Lambda P^{-1}\) with \(\Lambda=\diag(-1,1,2,3)\) (Exercise~\ref{r7-dr-eigenspaces}(c)), so by
Proposition~\ref{r7-prop-diag-rank}(b):
\begin{itemize}
\item \(A^2\) has the eigenvalues \(1\), \(1\), \(4\), \(9\); \(\tr(A^2)=15\), \(\det(A^2)=36\).
\item \(0\) is not an eigenvalue, so \(A\) is invertible (Theorem~\ref{r7-thm-trace-det}), and \(A^{-1}\) has the
eigenvalues \(-1\), \(1\), \(\frac12\), \(\frac13\); \(\tr(A^{-1})=-1+1+\frac12+\frac13=\frac56\), and
\(\det(A^{-1})=-\frac16\).
\item For \(f(t)=t^3-4t\), \(f(-1)=3\), \(f(1)=-3\), \(f(2)=0\), \(f(3)=15\): \(A^3-4A\) has the eigenvalues \(3\),
\(-3\), \(0\), \(15\); \(\tr(A^3-4A)=15\), \(\det(A^3-4A)=0\).
\end{itemize}
(c) By Proposition~\ref{r7-prop-eig-poly}(d), \(A'\) has the eigenvalues \(-1\), \(1\), \(2\), \(3\); \(\tr A'=5\),
\(\det A'=-6\).

(d)
\begin{itemize}
\item By Proposition~\ref{r7-prop-eig-poly}(a), \(f(A)x=f(\lambda)x\): \((1,0,0,1)'\) belongs to \(3\), \((1,0,1,1)'\) to
\(-3\), \((1,0,1,0)'\) to \(0\), and \((-1,-1,-1,1)'\) to \(15\).
\item By Theorem~\ref{r1-thm-views}(d), the \(i\)th diagonal entry of \(A^2\) is \(\row_i(A)\col_i(A)\):
\(0\cdot0+0\cdot0+2\cdot1+(-1)(-2)=4\), \(0\cdot0+3\cdot3+0\cdot0+0\cdot(-2)=9\), \(1\cdot2+0\cdot0+1\cdot1+(-1)\cdot2=1\)
and \((-2)(-1)+(-2)\cdot0+2\cdot(-1)+1\cdot1=1\); so \(\tr(A^2)=4+9+1+1=15\).
\end{itemize}
\end{solution}

\begin{exercise}\label{r7-ex-functions}
Let \(B\) be the matrix of Example~\ref{r7-exm-diag-powers}.
\begin{enumerate}
\item Without computing these matrices, find the eigenvalues and eigenvectors of \(B^3\), \(B^{-1}\), \(B+cI_4\)
(\(c\in\R\)) and \(q(B)=B^2-4B+3I_4\), and the numbers \(\tr B^3\), \(\det B^3\), \(\tr B^{-1}\), \(\det B^{-1}\).
\item Show that \(q(B)=0\). Deduce \(B^{-1}=\frac13(4I_4-B)\) and \(B^3=13B-12I_4\), and check the values of
\(\tr B^{-1}\) and \(\tr B^3\).
\end{enumerate}
\end{exercise}
\begin{solution}
(a) By Example~\ref{r7-exm-diag-powers}, \(B\) has the eigenvalue \(3\), with the eigenvectors \(p_1,p_2\), and the
eigenvalue \(1\), with the eigenvectors \(p_3,p_4\), and \(B=P\Lambda P^{-1}\) with \(\Lambda=\diag(3,3,1,1)\); so \(B\)
is invertible (Proposition~\ref{r7-prop-diag-powers}(c)). By Proposition~\ref{r7-prop-eig-poly}(a),(b), \(p_1,p_2\) and
\(p_3,p_4\) are eigenvectors of \(B^3\) belonging to \(27\) and \(1\), of \(B^{-1}\) belonging to \(\frac13\) and \(1\), of
\(B+cI_4\) belonging to \(3+c\) and \(1+c\), and of \(q(B)\) belonging to \(q(3)=9-12+3=0\) and \(q(1)=0\). By
Proposition~\ref{r7-prop-diag-rank}(b), these are the eigenvalues, listed with multiplicity (each twice), and
\(\tr B^3=56\), \(\det B^3=729\), \(\tr B^{-1}=\frac83\), \(\det B^{-1}=\frac19\). For \(B^3\), \(B^{-1}\) and \(B+cI_4\),
each of the two eigenvalues has algebraic multiplicity \(2\), so its eigenspace has dimension at most \(2\)
(Theorem~\ref{r7-thm-geo-alg}); it contains the independent \(p_1,p_2\), respectively \(p_3,p_4\), so it is spanned by them
(Theorem~\ref{thm:r4-counting}(d)). For \(q(B)\), the eigenspace for \(0\) contains \(p_1,\dots,p_4\), so it is \(\R^4\).

(b) \(q(B)\) maps the basis \(p_1,\dots,p_4\) of \(\R^4\) to \(0\), so \(q(B)x=0\) for every \(x\in\R^4\); with
\(x=e_j\), every column of \(q(B)\) is \(0\) (Theorem~\ref{r1-thm-views}(b)), and \(q(B)=0\) (multiplying out
\((B-I_4)(B-3I_4)=q(B)\) gives \(0\) as well). Hence \(B(4I_4-B)=4B-B^2=3I_4\), and \(B^{-1}=\frac13(4I_4-B)\)
(Theorem~\ref{r2-thm-invertible}):
\[
B^{-1}=\frac13\begin{pmatrix}3&-2&2&-2\\2&1&2&0\\0&2&1&2\\-2&2&-2&3\end{pmatrix}
\]
Also \(B^2=4B-3I_4\), so \(B^3=4B^2-3B=4(4B-3I_4)-3B=13B-12I_4\). With \(\tr B=8\): \(\tr B^{-1}=\frac13(16-8)=\frac83\)
and \(\tr B^3=104-48=56\).
\end{solution}

\begin{exercise}\label{r7-ex-schur-drill}
Let
\[
A=\begin{pmatrix}1&1&1&0\\1&1&0&1\\1&-1&1&2\\-1&1&2&1\end{pmatrix}
\]
and let \(Q=\frac12G\) be the orthogonal matrix of Example~\ref{r7-exm-schur}.
\begin{enumerate}
\item Show that \(\frac12\one_4\) is an eigenvector of \(A\) of length \(1\), and that \(T=Q'AQ\) is upper triangular. Read
off the eigenvalues of \(A\).
\item Use (a) to find \(\tr(A^2+I_4)\) and \(\det(A^2+I_4)\), and check the trace directly.
\item Is \(A\) diagonalizable? Is there an orthogonal \(U\) with \(U'AU\) diagonal?
\end{enumerate}
\end{exercise}
\begin{solution}
(a) Every row of \(A\) adds up to \(3\), so \(A\one_4=3\one_4\), and \(\lVert\frac12\one_4\rVert=1\). As \(G'=G\),
\[
AG=\begin{pmatrix}3&1&1&-1\\3&1&-1&1\\3&-3&1&3\\3&-3&-1&-3\end{pmatrix},\qquad
T=Q'AQ=\tfrac14G(AG)=\begin{pmatrix}3&-1&0&0\\0&2&0&0\\0&0&1&1\\0&0&0&-2\end{pmatrix}
\]
As \(Q'=Q^{-1}\) (Proposition~\ref{r6-prop-orthonormal}(b)), \(T\) is similar to \(A\), and by
Theorem~\ref{r7-thm-similar}(a) and Proposition~\ref{r7-prop-block-triangular}(b) the eigenvalues of \(A\) are \(3\),
\(2\), \(1\), \(-2\).

(b) As \(Q'Q=QQ'=I_4\), \(A=QTQ'\) and \(A^2+I_4=Q(T^2+I_4)Q'\), with
\[
T^2+I_4=\begin{pmatrix}10&-5&0&0\\0&5&0&0\\0&0&2&-1\\0&0&0&5\end{pmatrix}
\]
By Theorem~\ref{r7-thm-similar}(a), Proposition~\ref{r7-prop-block-triangular}(b) and Theorem~\ref{r7-thm-trace-det}, \(\tr(A^2+I_4)=10+5+2+5=22\) and
\(\det(A^2+I_4)=10\cdot5\cdot2\cdot5=500\), as Corollary~\ref{r7-cor-schur}(a) predicts (\(q(t)=t^2+1\) has
\(q(3)=10\), \(q(2)=5\), \(q(1)=2\), \(q(-2)=5\)). Directly,
\[
A^2+I_4=\begin{pmatrix}4&1&2&3\\1&4&3&2\\-1&1&7&3\\1&-1&3&7\end{pmatrix}
\]
with trace \(22\).

(c) \(A\) has four distinct real eigenvalues, so it is diagonalizable (Theorem~\ref{r7-thm-diag-criterion}). If
\(U'AU=\Lambda\) with \(U\) orthogonal and \(\Lambda\) diagonal, then \(A=U\Lambda U'\) would be symmetric
(Theorem~\ref{r1-thm-transpose}); but \(a_{14}=0\ne-1=a_{41}\). So no such \(U\) exists.
\end{solution}

\begin{exercise}\label{r7-dr-ab-ba}
Let
\[
A=\begin{pmatrix}0&1&0\\1&1&2\end{pmatrix}\in\R^{2\times3},\qquad B=\begin{pmatrix}-1&1\\1&1\\1&-1\end{pmatrix}\in\R^{3\times2}
\]
\begin{enumerate}
\item Compute \(AB\) and \(BA\). Find \(p_{AB}\), the eigenvalues of \(AB\), and an eigenvector for each.
\item Use Theorem~\ref{r7-thm-AB-BA} to find \(p_{BA}\) and the eigenvalues of \(BA\), and check \(p_{BA}\) by expanding
\(\det(tI-BA)\).
\item For each eigenvector \(x\) of \(AB\) from (a), check that \(Bx\) is an eigenvector of \(BA\). Find a basis of
\(\mathcal N(BA)\), and compare \(\rank(BA)\) with the number of nonzero eigenvalues of \(BA\).
\end{enumerate}
\end{exercise}
\begin{solution}
(a)
\begin{itemize}
\item Multiplying out,
\[
AB=\begin{pmatrix}1&1\\2&0\end{pmatrix},\qquad BA=\begin{pmatrix}1&0&2\\1&2&2\\-1&0&-2\end{pmatrix}
\]
\item By Definition~\ref{r7-def-charpoly}(b) (\(n=2\)), \(p_{AB}=t^2-t-2=(t-2)(t+1)\): the eigenvalues of \(AB\) are \(2\) and
\(-1\).
\item \((AB-2I)x=0\) reads \(-x_1+x_2=0\) and \(2x_1-2x_2=0\), so \(x=(1,1)'\); \((AB+I)x=0\) reads \(2x_1+x_2=0\) (twice), so
\(x=(-1,2)'\).
\end{itemize}
(b)
\begin{itemize}
\item By Theorem~\ref{r7-thm-AB-BA}(a) (\(m=2\), \(n=3\)), \(p_{BA}=t\,p_{AB}=t(t-2)(t+1)\), and by (b) of that theorem the
eigenvalues of \(BA\) are \(0\), \(2\), \(-1\).
\item Check: for \(t\in\R\), column \(2\) of \(tI-BA\) has the entry \(t-2\) in position \((2,2)\) and \(0\) elsewhere, so
\[
\begin{aligned}
\det(tI-BA)&=\det\begin{pmatrix}t-1&0&-2\\-1&t-2&-2\\1&0&t+2\end{pmatrix}=(t-2)\det\begin{pmatrix}t-1&-2\\1&t+2\end{pmatrix}\\
&=(t-2)(t^2+t)=t(t-2)(t+1)
\end{aligned}
\]
\end{itemize}
(c)
\begin{itemize}
\item \(B(1,1)'=(0,2,0)'\) and \(BA(0,2,0)'=(0,4,0)'=2(0,2,0)'\); \(B(-1,2)'=(3,1,-3)'\) and
\(BA(3,1,-3)'=(-3,-1,3)'=-(3,1,-3)'\).
\item Row-reduce \(BA\):
\[
\begin{pmatrix}1&0&2\\1&2&2\\-1&0&-2\end{pmatrix}
\xrightarrow[R_3\leftarrow R_3+R_1]{R_2\leftarrow R_2-R_1}
\begin{pmatrix}1&0&2\\0&2&0\\0&0&0\end{pmatrix}
\xrightarrow{R_2\leftarrow\frac12R_2}
\begin{pmatrix}1&0&2\\0&1&0\\0&0&0\end{pmatrix}
\]
The special solution \((-2,0,1)'\) is a basis of \(\mathcal N(BA)\), and \(\rank(BA)=2\).
\item \(BA\) has the two nonzero eigenvalues \(2\), \(-1\), as many as its rank: \(BA\) has three distinct eigenvalues, so it
is diagonalizable (Theorem~\ref{r7-thm-diag-criterion}), and Proposition~\ref{r7-prop-diag-rank}(a) applies.
\end{itemize}
\end{solution}

\begin{exercise}\label{r7-ex-ab-ba}
Let
\[
A=\begin{pmatrix}1&0\\1&1\\0&1\\1&-1\end{pmatrix},\qquad B=\begin{pmatrix}1&0&1&0\\0&1&0&1\end{pmatrix}
\]
\begin{enumerate}
\item Compute \(BA\) and \(AB\), find \(p_{BA}\), and deduce \(p_{AB}\).
\item Find eigenvectors of \(AB\) for its nonzero eigenvalues from those of \(BA\), and a basis of \(\mathcal N(AB)\).
Diagonalize \(AB\).
\item Show that \(\tr(AB)^k=2^k+(-1)^k\) for \(k\geqslant1\).
\end{enumerate}
\end{exercise}
\begin{solution}
(a)
\[
BA=\begin{pmatrix}1&1\\2&0\end{pmatrix},\qquad
AB=\begin{pmatrix}1&0&1&0\\1&1&1&1\\0&1&0&1\\1&-1&1&-1\end{pmatrix}
\]
By Definition~\ref{r7-def-charpoly}(b) (\(n=2\)), \(p_{BA}=t^2-t-2=(t-2)(t+1)\). By Theorem~\ref{r7-thm-AB-BA}(a),
applied with \(B\) and \(A\) in place of its \(A\) and \(B\) (\(m=2\), \(n=4\)), \(p_{AB}=t^2p_{BA}=t^2(t-2)(t+1)\).

(b)
\begin{itemize}
\item \(BA(1,1)'=(2,2)'\) and \(BA(1,-2)'=(-1,2)'\): \((1,1)'\) and \((1,-2)'\) are eigenvectors of \(BA\) belonging to
the simple eigenvalues \(2\) and \(-1\), so each is a basis of its eigenspace (Theorem~\ref{r7-thm-geo-alg}). By
Theorem~\ref{r7-thm-AB-BA}(c), applied as in (a),
\[
v_1=A(1,1)'=(1,2,1,0)',\qquad v_2=A(1,-2)'=(1,-1,-2,3)'
\]
are eigenvectors of \(AB\) belonging to \(2\) and \(-1\); indeed \(ABv_1=(2,4,2,0)'=2v_1\) and
\(ABv_2=(-1,1,2,-3)'=-v_2\).
\item \(B\) is in RREF, and \(Bx=0\) reads \(x_1+x_3=0\), \(x_2+x_4=0\), with the special solutions
\(n_1=(-1,0,1,0)'\) and \(n_2=(0,-1,0,1)'\) (Theorem~\ref{thm:r4-four-bases}(c)). As \(ABx=A(Bx)\),
\(\mathcal N(B)\subseteq\mathcal N(AB)\), so \(g_{AB}(0)\geqslant2\) (Theorem~\ref{thm:r4-counting}(a)); and
\(g_{AB}(0)\leqslant a_{AB}(0)=2\) (Theorem~\ref{r7-thm-geo-alg}). So \(\mathcal N(B)=\mathcal N(AB)\)
(Theorem~\ref{thm:r4-counting}(a)), with the basis \(n_1,n_2\).
\item All eigenvalues \(2\), \(-1\), \(0\), \(0\) of \(AB\) are real, and \(g_{AB}=a_{AB}\) for each of them; so \(AB\) is
diagonalizable (Theorem~\ref{r7-thm-diag-criterion}), and by the paragraph after that theorem, \(Q^{-1}(AB)Q=\Lambda\)
with \(Q=[v_1\ v_2\ n_1\ n_2]\) and \(\Lambda=\diag(2,-1,0,0)\).
\end{itemize}
(c) By (b), \(AB=Q\Lambda Q^{-1}\), so by Proposition~\ref{r7-prop-diag-rank}(b) (with \(f=t^k\)),
\(\tr(AB)^k=2^k+(-1)^k+0^k+0^k=2^k+(-1)^k\) for \(k\geqslant1\). For \(k=1\), \(\tr AB=1+1+0-1=1\).
\end{solution}

\begin{exercise}\label{r7-dr-cayley-hamilton}
Let
\[
A=\begin{pmatrix}0&-1&0&-1\\2&-1&0&-2\\-1&0&1&1\\-1&-1&0&0\end{pmatrix}
\]
\begin{enumerate}
\item Find \(p_A\).
\item Compute \(A^2\), \(A^3\) and \(A^4\), and check that \(p_A(A)=0\).
\item Use Corollary~\ref{r7-cor-cayley-hamilton}(b) to find \(A^{-1}\), and check that \(AA^{-1}=I\).
\item Use Corollary~\ref{r7-cor-cayley-hamilton}(a) to write \(A^5\) as a linear combination of \(I\), \(A\), \(A^2\),
\(A^3\), and compute \(A^5\).
\end{enumerate}
\end{exercise}
\begin{solution}
(a) For \(t\in\R\), column \(3\) of \(tI-A\) has the entry \(t-1\) in position \((3,3)\) and \(0\) elsewhere; expanding
along it, and then along row \(1\),
\[
\begin{aligned}
\det(tI-A)&=\det\begin{pmatrix}t&1&0&1\\-2&t+1&0&2\\1&0&t-1&-1\\1&1&0&t\end{pmatrix}
=(t-1)\det\begin{pmatrix}t&1&1\\-2&t+1&2\\1&1&t\end{pmatrix}\\
&=(t-1)\Bigl(t\bigl((t+1)t-2\bigr)-(-2t-2)+\bigl(-2-(t+1)\bigr)\Bigr)=(t-1)(t^3+t^2-t-1)
\end{aligned}
\]
So \(p_A=(t-1)(t^3+t^2-t-1)=t^4-2t^2+1\): \(c_0=1\), \(c_1=0\), \(c_2=-2\), \(c_3=0\), \(c_4=1\).

(b) Multiplying out,
\[
A^2=\begin{pmatrix}-1&2&0&2\\0&1&0&0\\-2&0&1&2\\-2&2&0&3\end{pmatrix},\qquad
A^3=A^2A=\begin{pmatrix}2&-3&0&-3\\2&-1&0&-2\\-3&0&1&3\\1&-3&0&-2\end{pmatrix},\qquad
A^4=A^3A=\begin{pmatrix}-3&4&0&4\\0&1&0&0\\-4&0&1&4\\-4&4&0&5\end{pmatrix}
\]
Adding entry by entry, \(p_A(A)=A^4-2A^2+I=0\), as Theorem~\ref{r7-thm-cayley-hamilton} asserts.

(c) As \(c_0=(-1)^4\det A\) (Definition~\ref{r7-def-charpoly}(b)), \(\det A=1\ne0\), and \(A\) is invertible
(Theorem~\ref{r5d-thm-product}(a)). By Corollary~\ref{r7-cor-cayley-hamilton}(b),
\(A^{-1}=-\frac1{c_0}(c_1I+c_2A+c_3A^2+c_4A^3)=2A-A^3\), that is,
\[
A^{-1}=\begin{pmatrix}-2&1&0&1\\2&-1&0&-2\\1&0&1&-1\\-3&1&0&2\end{pmatrix},\qquad
AA^{-1}=\begin{pmatrix}1&0&0&0\\0&1&0&0\\0&0&1&0\\0&0&0&1\end{pmatrix}
\]

(d)
\begin{itemize}
\item By Corollary~\ref{r7-cor-cayley-hamilton}(a), \(A^4=-(c_0I+c_1A+c_2A^2+c_3A^3)=2A^2-I\), as found in (b). Then
\(A^5=A^4A=2A^3-A\).
\item So
\[
A^5=2\begin{pmatrix}2&-3&0&-3\\2&-1&0&-2\\-3&0&1&3\\1&-3&0&-2\end{pmatrix}-\begin{pmatrix}0&-1&0&-1\\2&-1&0&-2\\-1&0&1&1\\-1&-1&0&0\end{pmatrix}
=\begin{pmatrix}4&-5&0&-5\\2&-1&0&-2\\-5&0&1&5\\3&-5&0&-4\end{pmatrix}
\]
\end{itemize}
\end{solution}

\needspace{14\baselineskip}
\subsection*{Problems}

\begin{problem}[\citeat{abadir-magnus}{Exercise~7.38}]\label{r7-ex-companion}
Let \(n\geqslant2\) and \(a_0,\dots,a_{n-1}\in\R\). Show that the companion matrix
\[
C=\begin{pmatrix}0&1&0&\cdots&0\\0&0&1&\cdots&0\\\vdots&&&\ddots&\vdots\\0&0&0&\cdots&1\\-a_0&-a_1&-a_2&\cdots&-a_{n-1}\end{pmatrix}
\in\R^{n\times n}
\]
whose entries are \(c_{i,i+1}=1\) (\(i<n\)), the last row \((-a_0,\dots,-a_{n-1})\) and \(0\) elsewhere, has the
characteristic polynomial \(p_C(t)=t^n+a_{n-1}t^{n-1}+\dots+a_1t+a_0\).
\end{problem}
\begin{solution}
Fix \(t\in\R\). For \(2\leqslant j\leqslant n\) and \(x\in\R\), let \(A_j(x)\in\R^{j\times j}\) have the entries \(t\) in
positions \((i,i)\) and \(-1\) in positions \((i,i+1)\) for \(i<j\), the last row \((a_0,a_1,\dots,a_{j-2},x)\), and \(0\)
elsewhere. Then \(tI-C=A_n(t+a_{n-1})\).
\begin{itemize}
\item Let \(j\geqslant3\). The last column of \(A_j(x)\) has the entries \(-1\) at \((j-1,j)\) and \(x\) at \((j,j)\), and \(0\)
elsewhere.
Deleting row \(j\) and column \(j\) leaves an upper triangular matrix with diagonal entries \(t\), of determinant
\(t^{j-1}\) (Theorem~\ref{r5d-thm-basic}(c)). Deleting row \(j-1\) and column \(j\) leaves \(A_{j-1}(a_{j-2})\). Expanding
along the last column (Theorem~\ref{r5d-thm-expansion}),
\[
\det A_j(x)=x\,t^{j-1}+(-1)\cdot(-1)^{(j-1)+j}\det A_{j-1}(a_{j-2})=x\,t^{j-1}+\det A_{j-1}(a_{j-2})
\]
\item For \(j=2\), \(\det A_2(x)=\det\left(\begin{smallmatrix}t&-1\\a_0&x\end{smallmatrix}\right)=tx+a_0\).
\item Hence
\[
\begin{aligned}
\det(tI-C)&=(t+a_{n-1})t^{n-1}+\det A_{n-1}(a_{n-2})=t^n+a_{n-1}t^{n-1}+a_{n-2}t^{n-2}+\det A_{n-2}(a_{n-3})\\
&=\cdots=t^n+a_{n-1}t^{n-1}+\dots+a_2t^2+\det A_2(a_1)=t^n+a_{n-1}t^{n-1}+\dots+a_1t+a_0
\end{aligned}
\]
(for \(n=2\), directly \(\det A_2(t+a_1)=t(t+a_1)+a_0\)). This holds for every real \(t\), so \(p_C\) is this polynomial
(Definition~\ref{r7-def-charpoly}(a), Proposition~\ref{r7-prop-poly}(d)).
\end{itemize}
\end{solution}

\begin{problem}[\citeat{harville-solutions}{Ch.~21, Exercise~8}]\label{r7-ex-not-similar}
Let \(n\geqslant2\). Show that for \(B\in\R^{n\times n}\) to be similar to \(A\in\R^{n\times n}\) it is not sufficient that
\(B\) have the same rank, determinant, trace and characteristic polynomial as \(A\).
\end{problem}
\begin{solution}
Take \(A=I_n\) and an upper triangular \(B\ne I_n\) with all diagonal entries \(1\), for instance \(B=I_n+E_{12}\)
(Definition~\ref{r3-def-basis}).
\begin{itemize}
\item \(\det B=1=\det I_n\) (Theorem~\ref{r5d-thm-basic}(c)), so \(B\) is invertible and \(\rank B=n=\rank I_n\)
(Theorem~\ref{r5d-thm-product}(a)); \(\tr B=n=\tr I_n\); and \(p_B=(t-1)^n=p_{I_n}\)
(Proposition~\ref{r7-prop-block-triangular}(b)).
\item If \(B\) were similar to \(I_n\), there would be an invertible \(C\) with \(B=C^{-1}I_nC=C^{-1}C=I_n\) \(\contradiction\)
\end{itemize}
\end{solution}

\begin{problem}[\citeat{harville-solutions}{Ch.~21, Exercise~13}]\label{r7-ex-sum-eigenvectors}
Let \(x_1\) and \(x_2\) be eigenvectors of \(A\in\R^{n\times n}\) belonging to \(\lambda_1\) and \(\lambda_2\), and let
\(c_1\), \(c_2\) be nonzero numbers. When is \(x=c_1x_1+c_2x_2\) an eigenvector of \(A\)?
\end{problem}
\begin{solution}
\begin{itemize}
\item \(Ax=c_1\lambda_1x_1+c_2\lambda_2x_2=\lambda_1x+(\lambda_2-\lambda_1)c_2x_2\).
\item If \(\lambda_1=\lambda_2\), then \(Ax=\lambda_1x\): \(x\) is an eigenvector unless \(x=0\), that is, unless
\(x_2=-(c_1/c_2)x_1\).
\item If \(\lambda_1\ne\lambda_2\), then \(x_1\) and \(x_2\) are independent (Theorem~\ref{r7-thm-indep}), and so are \(x\) and
\(x_2\): if \(ax+bx_2=0\), then \(ac_1x_1+(ac_2+b)x_2=0\), so \(ac_1=0\), \(a=0\), and then \(b=0\). If \(Ax=cx\) for some
\(c\), then \((\lambda_1-c)x+(\lambda_2-\lambda_1)c_2x_2=0\), so \((\lambda_2-\lambda_1)c_2=0\) \(\contradiction\) Hence \(x\)
is not an eigenvector.
\end{itemize}
So \(x\) is an eigenvector of \(A\) if and only if \(\lambda_1=\lambda_2\) and \(x\ne0\).
\end{solution}

\begin{problem}[\citeat{harville-solutions}{Ch.~21, Exercise~2}, \citeat{harville}{\S21.1}]\label{r7-ex-invariant}
Let \(A\in\R^{n\times n}\), \(B\in\R^{k\times k}\) and \(X\in\R^{n\times k}\) satisfy \(AX=XB\). A subspace \(U\) of \(\R^n\) is
\textbf{invariant} relative to \(A\) if \(Au\in U\) for every \(u\in U\).
\begin{enumerate}
\item Show that \(\mathcal C(X)\) is invariant relative to \(A\).
\item Show that if \(X\) has full column rank, then every real eigenvalue of \(B\) is an eigenvalue of \(A\).
\end{enumerate}
\end{problem}
\begin{solution}
(a) Every \(u\in\mathcal C(X)\) is \(u=Xr\) for some \(r\in\R^k\) (Definition~\ref{def:r4-spaces}), so
\(Au=AXr=XBr=X(Br)\in\mathcal C(X)\).

(b) Let \(By=\lambda y\) with \(\lambda\in\R\) and \(y\ne0\). Then \(A(Xy)=XBy=X(\lambda y)=\lambda(Xy)\). If \(X\) has full
column rank, then \(\mathcal N(X)=\{0\}\) (Theorem~\ref{thm:r4-rank-rules}(b)), so \(Xy\ne0\): \(\lambda\) is an eigenvalue of
\(A\), and \(Xy\) is an eigenvector belonging to it.
\end{solution}

\begin{problem}[\citeat{harville-solutions}{Ch.~21, Exercise~14}]\label{r7-ex-left-eigenvectors}
Let \(A\in\R^{n\times n}\), and let \(Q\in\R^{n\times n}\) be invertible with \(Q^{-1}AQ=D=\diag(d_1,\dots,d_n)\). Let \(r_i'\) be
row \(i\) of \(Q^{-1}\) (\(1\leqslant i\leqslant n\)). Show:
\begin{enumerate}
\item \((Q^{-1})'\) diagonalizes \(A'\).
\item \(d_1,\dots,d_n\) are the eigenvalues of \(A'\).
\item \(r_i\) is an eigenvector of \(A'\) belonging to \(d_i\).
\end{enumerate}
\end{problem}
\begin{solution}
(a) By Theorem~\ref{r2-thm-inverse-props}(c),(a), \((Q^{-1})'=(Q')^{-1}\), so \(\bigl((Q^{-1})'\bigr)^{-1}=Q'\). By
Theorem~\ref{r1-thm-transpose},
\(D=D'=(Q^{-1}AQ)'=Q'A'(Q^{-1})'=\bigl((Q^{-1})'\bigr)^{-1}A'(Q^{-1})'\).

(b) By (a), \(\bigl((Q^{-1})'\bigr)^{-1}A'(Q^{-1})'=D\); so by Definition~\ref{r7-def-diagonalizable}, applied to \(A'\) and
\((Q^{-1})'\), \(d_1,\dots,d_n\) are the eigenvalues of \(A'\).

(c) Column \(i\) of \((Q^{-1})'\) is the transposed row \(i\) of \(Q^{-1}\), that is, \(r_i\)
(Definition~\ref{r1-def-transpose}). By Definition~\ref{r7-def-diagonalizable}, applied as in (b), the columns of
\((Q^{-1})'\) are eigenvectors of \(A'\): \(A'r_i=d_ir_i\) with \(r_i\ne0\).
\end{solution}

\begin{problem}[\citeat{abadir-magnus}{Exercise~7.37}]\label{r7-pr-left}
Let \(A\in\R^{n\times n}\), and let \(\lambda\) be a real eigenvalue of \(A\). A vector \(w\in\R^n\), \(w\ne0\), with
\(w'A=\lambda w'\) is a \textbf{left eigenvector} of \(A\) belonging to \(\lambda\).
\begin{enumerate}
\item Show that a left eigenvector of \(A\) belonging to \(\lambda\) is an eigenvector of \(A'\) belonging to \(\lambda\).
\item Let \(\lambda\ne\mu\) be real eigenvalues of \(A\). Show that every left eigenvector of \(A\) belonging to \(\mu\) is
orthogonal to every eigenvector of \(A\) belonging to \(\lambda\).
\end{enumerate}
\end{problem}
\begin{solution}
(a) Let \(w'A=\lambda w'\) with \(w\ne0\). By Theorem~\ref{r1-thm-transpose}, \(A'w=(w'A)'=(\lambda w')'=\lambda w\); so
\(w\) is an eigenvector of \(A'\) belonging to \(\lambda\) (Definition~\ref{r7-def-eigen}).

(b) Let \(Ax=\lambda x\) with \(x\ne0\), and \(w'A=\mu w'\) with \(w\ne0\). By Theorem~\ref{r1-thm-rules},
\[
w'Ax=w'(\lambda x)=\lambda(w'x),\qquad w'Ax=(\mu w')x=\mu(w'x)
\]
So \((\lambda-\mu)\,w'x=0\), and as \(\lambda\ne\mu\), \(w'x=0\): \(w\) and \(x\) are orthogonal
(Definition~\ref{r6-def-inner}).
\end{solution}

\begin{problem}[\citeat{searle-khuri}{Exercises~7.6, 7.7}]\label{r7-ex-nilpotent}
\begin{enumerate}
\item Let \(A\in\R^{n\times n}\), \(A\ne0\), satisfy \(A^n=0\). Show that \(A\) is not diagonalizable.
\item Find a \(2\times2\) matrix that is not diagonalizable.
\end{enumerate}
\end{problem}
\begin{solution}
(a) Suppose \(P^{-1}AP=D=\diag(d_1,\dots,d_n)\) with \(P\) invertible. Then \(A=PDP^{-1}\), and by
Proposition~\ref{r7-prop-diag-powers}(a), \(0=A^n=P\diag(d_1^n,\dots,d_n^n)P^{-1}\). Multiplying by \(P^{-1}\) on the left and
by \(P\) on the right gives \(\diag(d_1^n,\dots,d_n^n)=0\); so every \(d_i=0\), \(D=0\) and \(A=PDP^{-1}=0\)
\(\contradiction\)

(b) Let \(N=\begin{pmatrix}0&0\\1&0\end{pmatrix}\). Multiplying out, \(N^2=0\), while \(N\ne0\). By (a) (with \(n=2\)),
\(N\) is not diagonalizable.
\end{solution}

\begin{problem}[\citeat{abadir-magnus}{Exercise~7.36(a), (b)}]\label{r7-pr-rank-one}
Let \(x,y\in\R^n\).
\begin{enumerate}
\item Show that \(xy'\) has \(n-1\) zero eigenvalues and one eigenvalue \(y'x\).
\item Show that \(\det(I_n+xy')=1+y'x\).
\end{enumerate}
\end{problem}
\begin{solution}
(a) For \(n=1\), \(xy'=(x_1y_1)=(y'x)\), and its only eigenvalue is \(y'x\) (Definition~\ref{r7-def-charpoly}(b)). Let
\(n\geqslant2\), and let \(\lambda_1,\dots,\lambda_n\) be the eigenvalues of \(xy'\), listed with multiplicity
(Definition~\ref{r7-def-eigenvalues}).
\begin{itemize}
\item If \(y=0\), then \(xy'=0\), \(p_{xy'}=t^n\) (Proposition~\ref{r7-prop-block-triangular}(b)), and all \(n\) eigenvalues
are \(0=y'x\).
\item Let \(y\ne0\). The only row \(y'\) of the \(1\times n\) matrix \(y'\) is nonzero, so \(\rank y'=1\)
(Proposition~\ref{r3-prop-dependence}(a), Definition~\ref{def:r4-rank}), and \(\dim\mathcal N(y')=n-1\geqslant1\)
(Theorem~\ref{thm:r4-rank-nullity}): there are \(n-1\) linearly independent vectors \(u_1,\dots,u_{n-1}\) with
\(y'u_i=0\). Each of them satisfies \(xy'u_i=x(y'u_i)=0\) (Theorem~\ref{r1-thm-rules}); so
\(\mathcal N(y')\subseteq\mathcal N(xy')\), \(g_{xy'}(0)\geqslant n-1\geqslant1\) (Theorem~\ref{thm:r4-counting}(a)),
\(0\) is an eigenvalue of \(xy'\), and \(a_{xy'}(0)\geqslant n-1\) (Theorem~\ref{r7-thm-geo-alg}).
\item So we may number the eigenvalues so that \(\lambda_1=\dots=\lambda_{n-1}=0\). By Theorem~\ref{r7-thm-trace-det}
and Theorem~\ref{r1-thm-trace}(b),
\[
\lambda_n=\lambda_1+\dots+\lambda_n=\tr(xy')=\tr(y'x)=y'x
\]
\end{itemize}
(b) By (a) and Proposition~\ref{r7-prop-eig-poly}(c) (with \(c=1\)), the eigenvalues of \(I_n+xy'\) are \(1\)
(\(n-1\) times) and \(1+y'x\). The determinant is the product of the eigenvalues (Theorem~\ref{r7-thm-trace-det}), so
\(\det(I_n+xy')=1^{n-1}(1+y'x)=1+y'x\). (This is also the case \(u=x\), \(v=y\) of Theorem~\ref{r5d-thm-det-lemma}.)
\end{solution}

\begin{problem}[\citeat{abadir-magnus}{Exercises~7.56, 7.69}]\label{r7-pr-cayley-hamilton}
\begin{enumerate}
\item Let \(A=(a_{ij})\) and \(B=(b_{ij})\) be upper triangular \(n\times n\) matrices, and let \(1\leqslant k<n\), such that
\(a_{ij}=0\) for \(1\leqslant i\leqslant k\) and \(1\leqslant j\leqslant k\), and \(b_{k+1,k+1}=0\). Show that the matrix
\(C=AB\) is upper triangular with \(c_{ij}=0\) for \(1\leqslant i\leqslant k+1\) and \(1\leqslant j\leqslant k+1\).
\item Let \(A\in\R^{n\times n}\), and let all eigenvalues \(\lambda_1,\dots,\lambda_n\) of \(A\), listed with multiplicity,
be real. Show that
\[
(\lambda_1I_n-A)(\lambda_2I_n-A)\cdots(\lambda_nI_n-A)=0
\]
and deduce that \(p_A(A)=0\).
\end{enumerate}
\end{problem}
\begin{solution}
(a) Partition \(A\), \(B\) and \(C\) with diagonal blocks of sizes \(k\) and \(n-k\) (Definition~\ref{r5b-def-partition}).
The blocks below the diagonal blocks of \(A\) and \(B\) are \(0\), as \(A\) and \(B\) are upper triangular
(Definition~\ref{r1-def-diag}), and the upper left \(k\times k\) block of \(A\) is \(0\) by assumption. By
Theorem~\ref{r5b-thm-block}(c),
\[
A=\begin{pmatrix}0&A_{12}\\0&A_{22}\end{pmatrix},\qquad B=\begin{pmatrix}B_{11}&B_{12}\\0&B_{22}\end{pmatrix},\qquad
C=AB=\begin{pmatrix}0&A_{12}B_{22}\\0&A_{22}B_{22}\end{pmatrix}
\]
\begin{itemize}
\item \(C\) is upper triangular (Proposition~\ref{r1-prop-diag}(c)), and its first \(k\) columns are \(0\). It remains to
show that column \(k+1\) of \(C\), which consists of the first columns of \(A_{12}B_{22}\) and of \(A_{22}B_{22}\), is
\(0\).
\item \(B_{22}\) is upper triangular, and its first diagonal entry is \(b_{k+1,k+1}=0\); so its first column is \(0\).
Hence for every matrix \(E\) with \(n-k\) columns, the first column of \(EB_{22}\) is \(E\cdot0=0\)
(Theorem~\ref{r1-thm-views}(a)); in particular the first columns of \(A_{12}B_{22}\) and \(A_{22}B_{22}\) are \(0\).
\end{itemize}
So the first \(k+1\) columns of \(C\) are \(0\); in particular \(c_{ij}=0\) for \(1\leqslant i\leqslant k+1\) and
\(1\leqslant j\leqslant k+1\).

(b) For \(n=1\), \(A=(a_{11})\) with \(\lambda_1=a_{11}\) (Definition~\ref{r7-def-charpoly}(b)), so
\(\lambda_1I_1-A=0\) and \(p_A(A)=A-\lambda_1I_1=0\). Let \(n\geqslant2\).
\begin{itemize}
\item By Theorem~\ref{r7-thm-schur}, \(A=QTQ'\), where \(Q\) is orthogonal and \(T\) is upper triangular with the
diagonal entries \(\lambda_1,\dots,\lambda_n\), in this order. Let \(B_i=\lambda_iI_n-T\) (\(1\leqslant i\leqslant n\)): it
is upper triangular, and its \(i\)th diagonal entry is \(0\).
\item As \(QQ'=Q'Q=I_n\) (Proposition~\ref{r6-prop-orthonormal}(b)), \(\lambda_iI_n-A=Q(\lambda_iI_n)Q'-QTQ'=QB_iQ'\)
(Theorem~\ref{r1-thm-rules}), and
\[
(\lambda_1I_n-A)\cdots(\lambda_nI_n-A)=QB_1Q'QB_2Q'\cdots QB_nQ'=Q(B_1B_2\cdots B_n)Q'
\]
\item By induction on \(j\), \(B_1\cdots B_j\) is upper triangular, and its upper left \(j\times j\) block is \(0\)
(\(1\leqslant j\leqslant n\)). For \(j=1\), the \((1,1)\) entry of \(B_1\) is \(0\). If this holds for some \(j<n\), then
(a), with \(B_1\cdots B_j\), \(B_{j+1}\) and \(j\) in place of \(A\), \(B\) and \(k\), gives it for \(j+1\), as the
\((j+1)\)th diagonal entry of \(B_{j+1}\) is \(0\). For \(j=n\), \(B_1\cdots B_n=0\); so
\((\lambda_1I_n-A)\cdots(\lambda_nI_n-A)=0\).
\item \(p_A=(t-\lambda_1)\cdots(t-\lambda_n)\) (Definition~\ref{r7-def-eigenvalues}) is a product of the real polynomials
\(t-\lambda_i\), whose values at \(A\) are \(A-\lambda_iI_n\) (Definition~\ref{r7-def-polynomial}). By
Proposition~\ref{r7-prop-poly}(b), applied \(n-1\) times, and Theorem~\ref{r1-thm-rules},
\[
p_A(A)=(A-\lambda_1I_n)\cdots(A-\lambda_nI_n)=(-1)^n(\lambda_1I_n-A)\cdots(\lambda_nI_n-A)=0
\]
\end{itemize}
\end{solution}

\begin{problem}[\citeat{prasolov}{Problem~11.3}]\label{r7-ex-pr-two-invertible}
Show that every \(A\in\R^{n\times n}\) is the sum of two invertible matrices.
\end{problem}
\begin{solution}
\begin{itemize}
\item \(p_A\) has degree \(n\) (Definition~\ref{r7-def-charpoly}(b)), so it has at most \(n\) roots
(Proposition~\ref{r7-prop-poly}(d)), and the real eigenvalues of \(A\) are among them
(Definition~\ref{r7-def-charpoly}(a)). Let \(c_1,\dots,c_k\) (\(0\leqslant k\leqslant n\)) be the distinct real
eigenvalues of \(A\), and put \(c=1+\lvert c_1\rvert+\dots+\lvert c_k\rvert\) (\(c=1\) if \(k=0\)).
\item Then \(c>\lvert c_i\rvert\) for every \(i\), so neither \(c\) nor \(-c\) is an eigenvalue of \(A\). By
Theorem~\ref{r7-thm-eig-singular}, (i)\(\Leftrightarrow\)(ii), with \(\lambda=c\) and with \(\lambda=-c\), the matrices
\(A-cI\) and \(A-(-c)I=A+cI\) are invertible.
\item Hence \(\frac12(A+cI)\) and \(\frac12(A-cI)\) are invertible, as
\(\frac12(A\pm cI)\cdot2(A\pm cI)^{-1}=I\) (Theorem~\ref{r2-thm-invertible}), and
\[
A=\tfrac12(A+cI)+\tfrac12(A-cI)
\]
\end{itemize}
\end{solution}

\begin{problem}[\citeat{prasolov}{Problem~22.1}]\label{r7-ex-pr-orthogonal-charpoly}
Let \(Q\in\R^{n\times n}\) be orthogonal, with \(p_Q(t)=c_0+c_1t+\dots+c_nt^n\), and let
\(\tilde p(t)=c_n+c_{n-1}t+\dots+c_0t^n\) be the polynomial with the coefficients of \(p_Q\) in reverse order, so that
\(\tilde p(t)=t^np_Q(1/t)\) for real \(t\ne0\). Show that \(\tilde p=p_Q\) or \(\tilde p=-p_Q\).
\end{problem}
\begin{solution}
Fix a real \(t\ne0\).
\begin{itemize}
\item \(t^np_Q(1/t)=t^n\sum_{k=0}^nc_kt^{-k}=\sum_{k=0}^nc_kt^{n-k}=\tilde p(t)\) (Definition~\ref{r7-def-polynomial}).
\item By Definition~\ref{r7-def-charpoly}(a) and Theorem~\ref{r5d-thm-product}(c),
\(t^np_Q(1/t)=t^n\det(t^{-1}I-Q)=\det\bigl(t(t^{-1}I-Q)\bigr)=\det(I-tQ)\).
\item As \(Q'Q=I\), \(I-tQ=Q'Q-tQ=(Q'-tI)Q\) (Theorem~\ref{r1-thm-rules}), so \(\det(I-tQ)=\det(Q'-tI)\cdot\det Q\)
(Theorem~\ref{r5d-thm-product}(b)).
\item By Theorem~\ref{r5d-thm-product}(c), Definition~\ref{r7-def-charpoly}(a) and
Proposition~\ref{r7-prop-eig-poly}(d),
\[
\det(Q'-tI)=\det\bigl((-1)(tI-Q')\bigr)=(-1)^n\det(tI-Q')=(-1)^np_{Q'}(t)=(-1)^np_Q(t)
\]
\end{itemize}
So \(\tilde p(t)=\varepsilon\,p_Q(t)\) for every real \(t\ne0\), where \(\varepsilon=(-1)^n\det Q\); and
\(\varepsilon=1\) or \(\varepsilon=-1\), as \(\det Q=\pm1\) (Theorem~\ref{r6-thm-orthogonal}(c)). The polynomials
\(\tilde p\) and \(\varepsilon p_Q\) have the same value at infinitely many \(t\) (Proposition~\ref{r7-prop-poly}(b)), so
\(\tilde p=\varepsilon p_Q\) (Proposition~\ref{r7-prop-poly}(d)), that is, \(c_{n-k}=\varepsilon c_k\) for
\(k=0,\dots,n\).
\end{solution}

\begin{problem}[\citeat{prasolov}{Problem~11.2}]\label{r7-ex-pr-ratio-matrix}
Let \(n\geqslant2\), let \(d_1,\dots,d_n\) be nonzero real numbers, and let \(A=(a_{ij})\in\R^{n\times n}\) with
\(a_{ij}=d_i/d_j\). Find the eigenvalues of \(A\) and the eigenvectors belonging to each of them.
\end{problem}
\begin{solution}
Let \(D=\diag(d_1,\dots,d_n)\) and \(d=(d_1,\dots,d_n)'\).
\begin{itemize}
\item \(D\) is invertible, with \(D^{-1}=\diag(d_1^{-1},\dots,d_n^{-1})\) (Definition~\ref{r2-def-inverse}). Left
multiplication by \(D\) multiplies row \(i\) by \(d_i\), and right multiplication by \(D^{-1}\) multiplies column \(j\) by
\(d_j^{-1}\) (Proposition~\ref{r1-prop-diag}(a)); as every entry of \(J_n\) is \(1\),
the \((i,j)\) entry of \(DJ_nD^{-1}\) is \(d_i\cdot1\cdot d_j^{-1}=a_{ij}\). So \(A=DJ_nD^{-1}\), that is,
\(J_n=D^{-1}AD\): \(J_n\) is similar to \(A\) (Definition~\ref{r7-def-similar}).
\item By Theorem~\ref{r7-thm-similar}(a) and Example~\ref{r7-exm-ones}, \(p_A=p_{J_n}=t^{n-1}(t-n)\): the eigenvalues of
\(A\) are \(n\), simple, and \(0\), of algebraic multiplicity \(n-1\).
\item Eigenvectors of \(J_n=\one\one'\): \(J_ny=\one(\one'y)\) (Theorem~\ref{r1-thm-rules}), a multiple of \(\one\). So
\(J_ny=0\) if and only if \(\one'y=0\), as \(\one\ne0\) (Proposition~\ref{r3-prop-dependence}(a)); and if \(J_ny=ny\),
then \(y=\frac1n(\one'y)\one\). Hence
\(\mathcal N(J_n)=\{y:\one'y=0\}\), and \(\mathcal N(J_n-nI)=\Span\{\one\}\), as \(J_n\one=n\one\).
\item By Theorem~\ref{r7-thm-similar}(b), with \(P=D\) and \(J_n=D^{-1}AD\), \(x\) is an eigenvector of \(A\) belonging
to \(c\) if and only if \(D^{-1}x\) is an eigenvector of \(J_n\) belonging to \(c\). For \(c=n\): \(D^{-1}x=s\one\) with
\(s\ne0\), that is, \(x=sD\one=sd\). For \(c=0\): \(x\ne0\) and \(\one'D^{-1}x=0\), that is,
\[
\frac{x_1}{d_1}+\frac{x_2}{d_2}+\dots+\frac{x_n}{d_n}=0
\]
\item So the eigenvectors belonging to \(n\) are the nonzero multiples of \(d\), and those belonging to \(0\) are the
nonzero vectors of \(\mathcal N(A)=\{x:x_1/d_1+\dots+x_n/d_n=0\}\). Here \(\dim\mathcal N(A)=\dim\mathcal N(J_n)=n-1\)
(Theorem~\ref{r7-thm-similar}(b); \(\mathcal N(J_n)=\mathcal N(\one')\), \(\rank\one'=1\) as its only row is nonzero, and
Theorem~\ref{thm:r4-rank-nullity}), and the
\(n-1\) vectors \(d_1e_1-d_je_j\) (\(j=2,\dots,n\)) of \(\mathcal N(A)\) are linearly independent, as only the \(j\)th of
them has a nonzero entry in position \(j\); so they form a basis of \(\mathcal N(A)\)
(Theorem~\ref{thm:r4-counting}(d)).
\end{itemize}
Check: entry \(i\) of \(Ad\) is \(\sum_j(d_i/d_j)d_j=nd_i\), and entry \(i\) of \(Ax\) is \(d_i\sum_jx_j/d_j\).
\end{solution}

\begin{problem}[\citeat{prasolov}{Problem~25.2}]\label{r7-ex-pr-idempotent-2x2}
Find all idempotent matrices \(P\in\R^{2\times2}\).
\end{problem}
\begin{solution}
Let \(P\in\R^{2\times2}\). First, \(P\) is idempotent if and only if \(P=0\), or \(P=I_2\), or \(\tr P=1\) and
\(\det P=0\).
\begin{itemize}
\item Let \(P^2=P\). Then \(\rank P\in\{0,1,2\}\) (Theorem~\ref{thm:r4-rank-rules}(a)). If \(\rank P=0\), then
\(\mathcal C(P)=\{0\}\) (Definition~\ref{def:r4-rank}, Theorem~\ref{thm:r4-dimension}), so every column of \(P\) is \(0\) (Definition~\ref{def:r4-spaces})
and \(P=0\). If \(\rank P=2\), then \(P\) is invertible (Theorem~\ref{thm:r4-rank-rules}(d)), and multiplying \(P^2=P\)
on the left by \(P^{-1}\) gives \(P=I_2\). If \(\rank P=1\), then \(\tr P=\rank P=1\) (Theorem~\ref{r6-thm-idempotent}(c)),
and \(\det P=0\), as \(P\) is not invertible (Theorem~\ref{r5d-thm-product}(a)).
\item Conversely, \(0\) and \(I_2\) are idempotent. If \(\tr P=1\) and \(\det P=0\), then
\(p_P(t)=t^2-(\tr P)t+\det P=t^2-t\) (Definition~\ref{r7-def-charpoly}(b)), and the Cayley--Hamilton theorem
(Theorem~\ref{r7-thm-cayley-hamilton}) gives \(P^2-P=0\).
\end{itemize}
Next, the matrices with \(\tr P=1\) and \(\det P=0\).
\begin{itemize}
\item Write \(P=\left(\begin{smallmatrix}\alpha&\beta\\\gamma&\delta\end{smallmatrix}\right)\). If \(\tr P=1\), then
\(\delta=1-\alpha\); with \(a=2\alpha-1\), \(b=2\beta\) and \(c=2\gamma\),
\[
P=\frac12\begin{pmatrix}1+a&b\\c&1-a\end{pmatrix},\qquad
\det P=\tfrac14(1+a)(1-a)-\tfrac14bc=\tfrac14\bigl(1-a^2-bc\bigr)
\]
\item So \(\det P=0\) if and only if \(a^2+bc=1\). Conversely, every matrix of this form with \(a,b,c\in\R\) has trace
\(\frac12(1+a)+\frac12(1-a)=1\).
\end{itemize}
Hence the idempotent \(2\times2\) matrices are \(0\), \(I_2\), and the matrices
\[
P=\frac12\begin{pmatrix}1+a&b\\c&1-a\end{pmatrix}\qquad(a,b,c\in\R,\ a^2+bc=1)
\]
\end{solution}

\begin{problem}[\citeat{prasolov}{\S26.1}]\label{r7-pr-involution}
Let \(A\in\R^{n\times n}\) satisfy \(A^2=I_n\).
\begin{enumerate}
\item Show that \(E=\frac12(I_n+A)\) is idempotent and that \(A=2E-I_n\).
\item Show that \(A\) is diagonalizable, that its eigenvalues lie in \(\{1,-1\}\), that \(\mathcal N(A-I_n)=\mathcal C(I_n+A)\)
and \(\mathcal N(A+I_n)=\mathcal C(I_n-A)\), and that \(g_A(1)+g_A(-1)=n\). Compare with
Problem~\ref{ex:r4-involutory-rank}.
\end{enumerate}
\end{problem}
\begin{solution}
(a) By Theorem~\ref{r1-thm-rules} and \(A^2=I_n\), \(E^2=\frac14(I_n+2A+A^2)=\frac14(2I_n+2A)=E\), and
\(2E-I_n=I_n+A-I_n=A\).

(b)
\begin{itemize}
\item \((A-I_n)(I_n+A)=A^2-I_n=0\) and \((A+I_n)(I_n-A)=I_n-A^2=0\) (Theorem~\ref{r1-thm-rules}). So for every
\(y\in\R^n\), \((I_n+A)y\in\mathcal N(A-I_n)\) and \((I_n-A)y\in\mathcal N(A+I_n)\): \(\mathcal C(I_n+A)\subseteq\mathcal N(A-I_n)\)
and \(\mathcal C(I_n-A)\subseteq\mathcal N(A+I_n)\) (Definition~\ref{def:r4-spaces}).
\item Conversely, if \(Ax=x\), then \(x=\frac12(x+Ax)=(I_n+A)(\frac12x)\in\mathcal C(I_n+A)\); and if \(Ax=-x\), then
\(x=\frac12(x-Ax)=(I_n-A)(\frac12x)\in\mathcal C(I_n-A)\). Hence \(\mathcal N(A-I_n)=\mathcal C(I_n+A)\) and
\(\mathcal N(A+I_n)=\mathcal C(I_n-A)\).
\item Every \(x\in\R^n\) is \(x=\frac12(x+Ax)+\frac12(x-Ax)\), where \(\frac12(x+Ax)=(I_n+A)(\frac12x)\) lies in
\(\mathcal N(A-I_n)\) and \(\frac12(x-Ax)=(I_n-A)(\frac12x)\) in \(\mathcal N(A+I_n)\), by the first item. So
\(\R^n=\mathcal N(A-I_n)+\mathcal N(A+I_n)\) (Definition~\ref{r3-def-sum}). If \(x\) lies in both null spaces, then
\(x=Ax=-x\), so \(x=0\). By Theorems~\ref{thm:r4-dim-sum} and~\ref{thm:r4-dimension} and
Definition~\ref{r7-def-multiplicity}, \(g_A(1)+g_A(-1)=\dim\mathcal N(A-I_n)+\dim\mathcal N(A+I_n)=\dim\R^n=n\).
\item If \(Ax=\lambda x\) with \(\lambda\in\R\) and \(x\ne0\), then \(x=A^2x=\lambda^2x\)
(Proposition~\ref{r7-prop-eig-poly}(a)), so \(\lambda^2=1\) (Proposition~\ref{r3-prop-dependence}(a)): the real
eigenvalues of \(A\) are among \(1\) and \(-1\). So the dimensions of the eigenspaces of the distinct real eigenvalues of
\(A\) add up to \(g_A(1)+g_A(-1)=n\) (a null space \(\mathcal N(A\mp I_n)=\{0\}\) adds \(0\)), and \(A\) is diagonalizable
(Theorem~\ref{r7-thm-diag-criterion}, (iii)\(\Rightarrow\)(i)). By Definition~\ref{r7-def-diagonalizable}, all
eigenvalues of \(A\) are then real; so they lie in \(\{1,-1\}\).
\end{itemize}
In terms of ranks, \(g_A(1)=\dim\mathcal C(I_n+A)=\rank(I_n+A)\) and \(g_A(-1)=\rank(I_n-A)\)
(Definition~\ref{def:r4-rank}), so \(\rank(I_n+A)+\rank(I_n-A)=n\), the identity of Problem~\ref{ex:r4-involutory-rank}.
\end{solution}

\clearpage
\def\unitnumber{8}
\setcounter{theorem}{0}\setcounter{problem}{0}
\section*{Chapter 8. The spectral theorem and quadratic forms}

\subsection*{Real eigenvalues and orthogonal eigenvectors}

\begin{theorem}[\citeat{bapat}{Result~3.5}]\label{r8-thm-sym-real}
Let \(A\in\R^{n\times n}\) be symmetric. Then every eigenvalue of \(A\) (every complex root of \(p_A\),
Definition~\ref{r7-def-eigenvalues}) is real, and for each eigenvalue \(\lambda\) there is a nonzero \(x\in\R^n\) with
\(Ax=\lambda x\).
\end{theorem}
\begin{proof}
Let \(\mu=\alpha+i\beta\) (\(\alpha,\beta\in\R\)) be a root of \(p_A\), and put \(Y=(A-\alpha I)^2\in\R^{n\times n}\).
\begin{itemize}
\item For every real \(s\), \(\bigl((\alpha+s)I-A\bigr)\bigl((\alpha-s)I-A\bigr)=(\alpha I-A)^2-s^2I=Y-s^2I\), as \(sI\)
commutes with every matrix. By Theorem~\ref{r5d-thm-product}(b),
\[
p_A(\alpha+s)\,p_A(\alpha-s)=\det(Y-s^2I)\qquad(s\in\R)
\]
\item Both sides are polynomials in \(s\) with real coefficients: \(p_A\) has real coefficients
(Definition~\ref{r7-def-charpoly}), and on the left we expand the powers of \(\alpha\pm s\); on the right
\(\det(Y-rI)=(-1)^np_Y(r)\) (Theorem~\ref{r5d-thm-product}(c)) with \(r=s^2\). They agree at every real \(s\), so they
are equal as polynomials (Proposition~\ref{r7-prop-poly}(d)), and the identity also holds for \(s=i\beta\)
(Proposition~\ref{r7-prop-poly}(b)): \(\det(Y+\beta^2I)=p_A(\mu)\,p_A(\alpha-i\beta)=0\).
\item \(A-\alpha I\) is symmetric, so \(Y=(A-\alpha I)'(A-\alpha I)\), and for \(x\in\R^n\)
\[
x'(Y+\beta^2I)x=\lVert(A-\alpha I)x\rVert^2+\beta^2\lVert x\rVert^2
\]
If \(\beta\ne0\), this is \(>0\) for every \(x\ne0\), so \((Y+\beta^2I)x\ne0\): \(Y+\beta^2I\) is invertible
(Theorem~\ref{r2-thm-invertible}), and \(\det(Y+\beta^2I)\ne0\) (Theorem~\ref{r5d-thm-product}(a)) \(\contradiction\)
Hence \(\beta=0\).
\item So \(\lambda=\mu=\alpha\) is real and \(\det(\lambda I-A)=0\); by Theorem~\ref{r7-thm-eig-singular},
\(Ax=\lambda x\) for some nonzero \(x\in\R^n\).\qedhere
\end{itemize}
\end{proof}

Most matrices of statistics in this book whose eigenvalues we use are symmetric, so their eigenvalues are real by
Theorem~\ref{r8-thm-sym-real}: \(A'A\) for every matrix \(A\) (Theorem~\ref{r1-thm-transpose}), the projection matrices
\(P_W\) (Theorem~\ref{r6-thm-projection-matrix}(c)), among them \(\bar J_n\) and \(C_n\)
(Proposition~\ref{r1-prop-centering}), the equicorrelation matrix of Example~\ref{r5d-exm-equicorrelation}, and the
covariance and correlation matrices of data. A matrix that is not symmetric may have eigenvalues that are not real: the
matrix of Example~\ref{r7-exm-complex} has the eigenvalues \(1\) and \(2\pm i\). The others are the matrices
\(A^{-1}B\) with \(A\) PD and \(B\) symmetric, which need not be symmetric, but whose eigenvalues are real too; this is
shown in the section on simultaneous diagonalization.

\begin{theorem}[\citeat{hoffman-kunze}{\S8.5, Theorem~15}]\label{r8-thm-sym-orth}
Let \(A\in\R^{n\times n}\) be symmetric, and let \(Ax=\lambda x\), \(Ay=\mu y\) with \(x,y\in\R^n\) and \(\lambda\ne\mu\).
Then \(x'y=0\).
\end{theorem}
\begin{proof}
\(\lambda x'y=(Ax)'y=x'A'y=x'Ay=\mu x'y\), so \((\lambda-\mu)x'y=0\), and \(\lambda-\mu\ne0\).
\end{proof}

\subsection*{The spectral theorem}

\begin{theorem}[Spectral theorem; \citeat{harville}{Theorem~21.5.7, Corollary~21.5.9}, \citeat{lang}{Ch.~VIII, Theorem~4.4}]\label{r8-thm-spectral}
Let \(A\in\R^{n\times n}\) be symmetric.
\begin{enumerate}
\item \(\R^n\) has an orthonormal basis \(q_1,\dots,q_n\) consisting of eigenvectors of \(A\).
\item Let \(Aq_i=\lambda_iq_i\) (\(i=1,\dots,n\)), \(Q=[q_1\ \cdots\ q_n]\) and \(\Lambda=\diag(\lambda_1,\dots,\lambda_n)\).
Then \(Q\) is orthogonal and
\[
Q'AQ=\Lambda,\qquad A=Q\Lambda Q'
\]
Conversely, if \(B\in\R^{n\times n}\) and \(Q'BQ\) is diagonal for an orthogonal \(Q\), then \(B\) is symmetric. The
\(q_i\) may be numbered in any order, for instance so that \(\lambda_1\geqslant\lambda_2\geqslant\dots\geqslant\lambda_n\).
\item \(p_A(t)=(t-\lambda_1)\cdots(t-\lambda_n)\): the numbers \(\lambda_1,\dots,\lambda_n\) are the eigenvalues of \(A\),
each eigenvalue \(\lambda\) listed \(a_A(\lambda)\) times.
\end{enumerate}
\end{theorem}
\begin{proof}
(a) We show by induction on \(n\) that \(Q'AQ\) is diagonal for some orthogonal \(Q\in\R^{n\times n}\).
\begin{itemize}
\item For \(n=1\), take \(Q=(1)\).
\item Let \(n\geqslant2\), and assume the claim for symmetric matrices of size \(n-1\). As \(p_A\) has \(n\geqslant1\)
roots, \(A\) has an eigenvalue \(\lambda\) (Definition~\ref{r7-def-eigenvalues}); by Theorem~\ref{r8-thm-sym-real},
\(\lambda\) is real and \(Ax=\lambda x\) for some nonzero \(x\in\R^n\). Then \(u=x/\lVert x\rVert\) has \(u'u=1\) and
\(Au=\lambda u\).
\item By Corollary~\ref{r6-cor-onb}(c), with \(p=1<n\), there is \(V\in\R^{n\times(n-1)}\) such that \([u\ V]\) is
orthogonal. Comparing blocks of \([u\ V]'[u\ V]=I_n\) (Theorem~\ref{r5b-thm-block}(b),(c)) gives \(V'u=0\) and
\(V'V=I_{n-1}\).
\item \(V'Au=\lambda V'u=0\), \(u'AV=(V'A'u)'=(V'Au)'=0\) and \(u'Au=\lambda u'u=\lambda\); so, by block multiplication
(Theorem~\ref{r5b-thm-block}(b),(c)),
\[
[u\ V]'A[u\ V]=\begin{pmatrix}u'Au&u'AV\\V'Au&V'AV\end{pmatrix}=\begin{pmatrix}\lambda&0\\0&B\end{pmatrix},\qquad
B=V'AV\in\R^{(n-1)\times(n-1)}
\]
\item \(B'=V'A'V=B\). By the induction hypothesis, \(R'BR=\Lambda_2\) is diagonal for some orthogonal
\(R\in\R^{(n-1)\times(n-1)}\). Let \(S=\diag(1,R)\) and \(Q=[u\ V]S\). By Theorems~\ref{r1-thm-transpose}
and~\ref{r5b-thm-block}(b) and Proposition~\ref{r5b-prop-block-tri}(b),
\[
\begin{aligned}
Q'Q&=S'[u\ V]'[u\ V]S=S'S=\diag(1,R'R)=I_n\\
Q'AQ&=S'\diag(\lambda,B)\,S=\diag(\lambda,R'BR)=\diag(\lambda,\Lambda_2)
\end{aligned}
\]
so \(Q\) is orthogonal and \(Q'AQ\) is diagonal.
\item Let \(\Lambda=Q'AQ\), and let \(q_1,\dots,q_n\) be the columns of \(Q\). As \(QQ'=I_n\)
(Proposition~\ref{r6-prop-orthonormal}(b)), \(AQ=QQ'AQ=Q\Lambda\), whose column \(i\) reads \(Aq_i=\lambda_iq_i\)
(Theorem~\ref{r1-thm-views}(a),(b)); and \(q_i'q_j=\delta_{ij}\), as \(Q'Q=I_n\). So \(q_1,\dots,q_n\) are orthonormal
eigenvectors of \(A\), hence independent (Proposition~\ref{r6-prop-orthogonal-list}(b)), and they form a basis of
\(\R^n\) (Theorem~\ref{thm:r4-counting}(d)).
\end{itemize}
(b)
\begin{itemize}
\item \((Q'Q)_{ij}=q_i'q_j=\delta_{ij}\), so \(Q'Q=I_n\): \(Q\) is orthogonal, and \(Q^{-1}=Q'\)
(Proposition~\ref{r6-prop-orthonormal}(b)).
\item \(AQ=[Aq_1\ \cdots\ Aq_n]=[\lambda_1q_1\ \cdots\ \lambda_nq_n]=Q\Lambda\), since
\(\col_j(Q\Lambda)=Q(\lambda_je_j)=\lambda_jq_j\) (Theorem~\ref{r1-thm-views}(a),(b)). Multiplying by \(Q'\) on the
left, or on the right, gives \(Q'AQ=\Lambda\) and \(A=Q\Lambda Q'\).
\item Conversely, if \(Q'BQ=D\) is diagonal, then \(B=QDQ'\) and \(B'=QD'Q'=QDQ'=B\)
(Theorem~\ref{r1-thm-transpose}).
\item Any renumbering of \(q_1,\dots,q_n\) is again an orthonormal basis of eigenvectors of \(A\), so the above applies to
it; numbering the \(q_i\) by decreasing \(\lambda_i\) gives \(\lambda_1\geqslant\dots\geqslant\lambda_n\).
\end{itemize}
(c) As \(Q^{-1}=Q'\), \(Q^{-1}AQ=\Lambda\); so \(p_A=p_\Lambda=(t-\lambda_1)\cdots(t-\lambda_n)\) (Theorem~\ref{r7-thm-similar}(a),
Proposition~\ref{r7-prop-block-triangular}(b)), and \(\lambda_1,\dots,\lambda_n\) are the eigenvalues of \(A\), listed with
multiplicity: this is the last sentence of Definition~\ref{r7-def-diagonalizable}.
\end{proof}

\begin{theorem}[\citeat{harville}{Lemma~14.3.1, Corollary~21.5.8}, \citeat{searle-khuri}{\S6.5.3}]\label{r8-thm-rank-eigen}
Let \(A\in\R^{n\times n}\) be symmetric. For every real \(c\), \(\dim\mathcal N(A-cI)=a_A(c)\) (which is \(0\) if \(c\) is
not an eigenvalue). Hence \(g_A(c)=a_A(c)\) for every eigenvalue \(c\); \(\rank A\) is the number of nonzero
eigenvalues of \(A\), counted with multiplicity; and \(a_A(0)=\dim\mathcal N(A)=n-\rank A\). Moreover, if
\(c_1,\dots,c_k\) are the distinct eigenvalues of \(A\), then orthonormal bases of
\(\mathcal N(A-c_1I),\dots,\mathcal N(A-c_kI)\), put side by side in any order, form an orthonormal basis of \(\R^n\)
consisting of eigenvectors of \(A\), as in Theorem~\ref{r8-thm-spectral}(a).
\end{theorem}
\begin{proof}
Let \(A=Q\Lambda Q'\) as in Theorem~\ref{r8-thm-spectral}.
\begin{itemize}
\item A diagonal matrix \(D=\diag(d_1,\dots,d_n)\) has rank equal to the number of \(i\) with \(d_i\ne0\): its nonzero
columns \(d_ie_i\) are nonzero and pairwise orthogonal, hence independent (Proposition~\ref{r6-prop-orthogonal-list}(b)),
and they span \(\mathcal C(D)\).
\item For real \(c\), \(A-cI=Q(\Lambda-cI)Q'\) with \(Q\), \(Q'\) invertible, so \(\rank(A-cI)=\rank(\Lambda-cI)\)
(Theorem~\ref{thm:r4-rank-product}(d)), the number of \(i\) with \(\lambda_i\ne c\).
\item By Theorem~\ref{thm:r4-rank-nullity}, \(\dim\mathcal N(A-cI)=n-\rank(A-cI)\) is the number of \(i\) with
\(\lambda_i=c\), which is \(a_A(c)\) by Theorem~\ref{r8-thm-spectral}(c). For an eigenvalue \(c\) this is
\(g_A(c)=a_A(c)\) (Definition~\ref{r7-def-multiplicity}), and \(c=0\) gives the statements on \(\rank A\) and
\(a_A(0)\).
\item The \(c_j\) are real (Theorem~\ref{r8-thm-sym-real}) and \(\dim\mathcal N(A-c_jI)=a_A(c_j)\geqslant1\), so each
\(\mathcal N(A-c_jI)\) has an orthonormal basis of \(a_A(c_j)\) vectors (Corollary~\ref{r6-cor-onb}(a)). Vectors from
different eigenspaces are orthogonal (Theorem~\ref{r8-thm-sym-orth}), so all these vectors together are orthonormal,
hence independent (Proposition~\ref{r6-prop-orthogonal-list}(b)); there are \(a_A(c_1)+\dots+a_A(c_k)=n\) of them
(Definition~\ref{r7-def-eigenvalues}), so they form a basis of \(\R^n\) (Theorem~\ref{thm:r4-counting}(d)).\qedhere
\end{itemize}
\end{proof}

\begin{example}[\citeat{searle-khuri}{Example~6.13}]\label{r8-exm-spectral-j3}
Find an orthogonal \(Q\) with \(Q'AQ\) diagonal for
\[
A=\begin{pmatrix}1&2&2\\2&1&2\\2&2&1\end{pmatrix}=2J_3-I_3
\]

\textit{Solution.} For real \(t\), apply \(R_1\leftarrow R_1+R_2+R_3\), take the factor \(t-5\) out of row \(1\), and
apply \(R_2\leftarrow R_2+2R_1\), \(R_3\leftarrow R_3+2R_1\) (Definition~\ref{r5d-def-det}(a),(b),
Theorem~\ref{r5d-thm-basic}(c)):
\[
\begin{aligned}
\det(tI-A)&=\det\begin{pmatrix}t-1&-2&-2\\-2&t-1&-2\\-2&-2&t-1\end{pmatrix}
=\det\begin{pmatrix}t-5&t-5&t-5\\-2&t-1&-2\\-2&-2&t-1\end{pmatrix}\\
&=(t-5)\det\begin{pmatrix}1&1&1\\-2&t-1&-2\\-2&-2&t-1\end{pmatrix}
=(t-5)\det\begin{pmatrix}1&1&1\\0&t+1&0\\0&0&t+1\end{pmatrix}=(t-5)(t+1)^2
\end{aligned}
\]
This holds for every real \(t\), hence as polynomials (Definition~\ref{r7-def-charpoly}(a), Proposition~\ref{r7-prop-poly}(d)): the eigenvalues are \(5\) with
\(a_A(5)=1\) and \(-1\) with \(a_A(-1)=2\).
\begin{itemize}
\item Every row sum of \(A\) is \(5\), so \(A\one_3=5\one_3\); as \(\dim\mathcal N(A-5I)=a_A(5)=1\)
(Theorem~\ref{r8-thm-rank-eigen}), \(\mathcal N(A-5I)=\Span\{\one_3\}\).
\item Reduce \(A+I=2J_3\):
\[
A+I=\begin{pmatrix}2&2&2\\2&2&2\\2&2&2\end{pmatrix}
\xrightarrow{R_1\leftarrow\frac12R_1}\begin{pmatrix}1&1&1\\2&2&2\\2&2&2\end{pmatrix}
\xrightarrow[R_3\leftarrow R_3-2R_1]{R_2\leftarrow R_2-2R_1}\begin{pmatrix}1&1&1\\0&0&0\\0&0&0\end{pmatrix}
\]
so \(\mathcal N(A+I)=\{x:x_1+x_2+x_3=0\}\), of dimension \(2\). It contains \(u_2=(-2,1,1)'\). A second vector of
\(\mathcal N(A+I)\) orthogonal to \(u_2\) solves \(x_1+x_2+x_3=0\) and \(u_2'x=0\):
\[
\begin{pmatrix}1&1&1\\-2&1&1\end{pmatrix}
\xrightarrow{R_2\leftarrow R_2+2R_1}
\begin{pmatrix}1&1&1\\0&3&3\end{pmatrix}
\xrightarrow{R_2\leftarrow\frac13R_2}
\begin{pmatrix}1&1&1\\0&1&1\end{pmatrix}
\xrightarrow{R_1\leftarrow R_1-R_2}
\begin{pmatrix}1&0&0\\0&1&1\end{pmatrix}
\]
so \(x_1=0\), \(x_2=-x_3\); take \(u_3=(0,-1,1)'\).
\item \(\one_3\perp u_2,u_3\) (Theorem~\ref{r8-thm-sym-orth}) and \(u_2\perp u_3\). Normalizing,
\[
Q=\begin{pmatrix}1/\sqrt3&-2/\sqrt6&0\\1/\sqrt3&1/\sqrt6&-1/\sqrt2\\1/\sqrt3&1/\sqrt6&1/\sqrt2\end{pmatrix},\qquad
Q'AQ=\diag(5,-1,-1)
\]
by Theorem~\ref{r8-thm-spectral}(b). Check: \(Au_2=(2,-1,-1)'=-u_2\) and \(Au_3=(0,1,-1)'=-u_3\).
\end{itemize}
\end{example}

\subsection*{The spectral decomposition}

\begin{theorem}[Spectral decomposition; \citeat{hoffman-kunze}{\S9.5, Theorems~9, 10}, \citeat{harville}{\S21.5f, Theorem~21.4.3}, \citeat{bapat}{\S3.2}]\label{r8-thm-spectral-decomp}
Let \(A\in\R^{n\times n}\) be symmetric with distinct eigenvalues \(c_1,\dots,c_k\), let \(W_j=\mathcal N(A-c_jI)\), and
let \(E_j=P_{W_j}\in\R^{n\times n}\) be the projection matrix for \(W_j\) (Theorem~\ref{r6-thm-projection-matrix}(a)):
\(E_jy\) is the orthogonal projection of \(y\) onto \(W_j\) (Definition~\ref{r6-def-projection}), and
\(E_j=V_j(V_j'V_j)^{-1}V_j'\) whenever the columns of \(V_j\) form a basis of \(W_j\)
(Theorem~\ref{r6-thm-projection-matrix}(b)). (Here \(E_j\) is a projection matrix, not an elementary matrix as in
Theorem~\ref{r2-thm-elementary}(c).)
\begin{enumerate}
\item \(W_i\perp W_j\) for \(i\ne j\). Every \(y\in\R^n\) can be written as \(y=y_1+\dots+y_k\) with \(y_j\in W_j\) in
exactly one way, and then \(E_jy=y_j\).
\item \(E_1+\dots+E_k=I_n\), \(E_j'=E_j=E_j^2\), \(E_iE_j=0\) for \(i\ne j\), each \(E_j\ne0\) (as \(E_jw=w\) for
\(w\in W_j\), by (a)), and
\[
A=c_1E_1+\dots+c_kE_k,\qquad A^m=c_1^mE_1+\dots+c_k^mE_k\quad(m=0,1,2,\dots)
\]
If no \(c_j\) is \(0\), then \(A\) is invertible and \(A^{-1}=c_1^{-1}E_1+\dots+c_k^{-1}E_k\).
\item If \(A=Q\Lambda Q'\) as in Theorem~\ref{r8-thm-spectral} and \(Q_j\) consists of the columns \(q_i\) of \(Q\) with
\(\lambda_i=c_j\), then \(E_j=Q_jQ_j'=\sum_{\lambda_i=c_j}q_iq_i'\), \(A=\sum_{i=1}^n\lambda_iq_iq_i'\), and
\(A^m=Q\Lambda^mQ'\) (\(m=0,1,2,\dots\)). Moreover \(\mathcal N(A)=\Span\{q_i:\lambda_i=0\}\) and
\(\mathcal C(A)=\Span\{q_i:\lambda_i\ne0\}\).
\item (Uniqueness) Let \(A=b_1F_1+\dots+b_rF_r\), where \(b_1,\dots,b_r\) are distinct numbers and
\(F_1,\dots,F_r\in\R^{n\times n}\) are nonzero, symmetric and idempotent, with \(F_1+\dots+F_r=I_n\) and \(F_lF_m=0\)
for \(l\ne m\). Then \(r=k\), and after renumbering \(b_j=c_j\) and \(F_j=E_j\) for every \(j\).
\end{enumerate}
The expression \(A=\sum_jc_jE_j\) is the \textbf{spectral decomposition} of \(A\).
\end{theorem}
\begin{proof}
(a)
\begin{itemize}
\item \(W_i\perp W_j\) by Theorem~\ref{r8-thm-sym-orth}.
\item Let \(q_1,\dots,q_n\) be an orthonormal basis of eigenvectors (Theorem~\ref{r8-thm-spectral}). Then
\(y=\sum_i(q_i'y)q_i\) (Proposition~\ref{r6-prop-orthogonal-list}(a)), and each \(q_i\) lies in the \(W_j\) with
\(c_j=\lambda_i\); grouping the terms gives \(y=\sum_jy_j\) with \(y_j\in W_j\).
\item If \(z_1+\dots+z_k=0\) with \(z_j\in W_j\), then \(0=z_i'\sum_jz_j=\sum_jz_i'z_j=\lVert z_i\rVert^2\), so every
\(z_i=0\). Applied to the differences of two such expressions of \(y\), this gives uniqueness.
\item \(y-y_j=\sum_{l\ne j}y_l\perp W_j\) and \(y_j\in W_j\), so \(y_j\) is the orthogonal projection of \(y\) onto
\(W_j\): \(E_jy=y_j\).
\end{itemize}
(b)
\begin{itemize}
\item \(\sum_jE_jy=\sum_jy_j=y\) for every \(y\), so \(\sum_jE_j=I_n\). \(E_j'=E_j=E_j^2\) by
Theorem~\ref{r6-thm-projection-matrix}(c).
\item For \(i\ne j\) and every \(y\), \(E_jy\in W_j\) is its own decomposition in (a), with zero part in \(W_i\); so
\(E_iE_jy=0\).
\item \(Ay=\sum_jAy_j=\sum_jc_jy_j=\sum_jc_jE_jy\) for every \(y\), so \(A=\sum_jc_jE_j\).
\item \(A^0=I_n=\sum_jE_j\), and if \(A^m=\sum_ic_i^mE_i\), then
\(A^{m+1}=\bigl(\sum_ic_i^mE_i\bigr)\bigl(\sum_jc_jE_j\bigr)=\sum_jc_j^{m+1}E_j\), since \(E_iE_j=0\) for \(i\ne j\) and
\(E_j^2=E_j\).
\item In the same way \(A\sum_jc_j^{-1}E_j=\sum_jE_j=I_n\), so \(A^{-1}=\sum_jc_j^{-1}E_j\)
(Theorem~\ref{r2-thm-invertible}).
\end{itemize}
(c)
\begin{itemize}
\item The columns of \(Q_j\) are orthonormal eigenvectors for \(c_j\), so they lie in \(W_j\) and are independent; their
number is \(a_A(c_j)=\dim W_j\) (Theorem~\ref{r8-thm-rank-eigen}). So they form an orthonormal basis of \(W_j\)
(Theorem~\ref{thm:r4-counting}(d)), and \(E_j=Q_jQ_j'\) (Theorem~\ref{r6-thm-projection-matrix}(b)).
\item By the outer product expansion (Theorem~\ref{r1-thm-views}(d)), \(Q_jQ_j'=\sum_{\lambda_i=c_j}q_iq_i'\), and
\(A=(Q\Lambda)Q'=\sum_i\lambda_iq_iq_i'\), as \(\col_i(Q\Lambda)=\lambda_iq_i\).
\item \(A=Q\Lambda Q^{-1}\) with \(Q^{-1}=Q'\) (Proposition~\ref{r6-prop-orthonormal}(b)), so
\(A^m=Q\Lambda^mQ^{-1}=Q\Lambda^mQ'\) (Proposition~\ref{r7-prop-diag-powers}(a)).
\item If \(\lambda_i\ne0\), then \(q_i=A(\lambda_i^{-1}q_i)\in\mathcal C(A)\); if \(\lambda_i=0\), then
\(Aq_i=\lambda_iq_i=0\). So \(\Span\{q_i:\lambda_i\ne0\}\subseteq\mathcal C(A)\) and
\(\Span\{q_i:\lambda_i=0\}\subseteq\mathcal N(A)\) (Theorem~\ref{r3-thm-span}). The \(q_i\) are independent
(Proposition~\ref{r6-prop-orthogonal-list}(b)), and there are \(\rank A=\dim\mathcal C(A)\) of them with
\(\lambda_i\ne0\) and \(a_A(0)=\dim\mathcal N(A)\) with \(\lambda_i=0\) (Theorems~\ref{r8-thm-spectral}(c)
and~\ref{r8-thm-rank-eigen}, Definition~\ref{def:r4-rank}); so both inclusions are equalities
(Theorem~\ref{thm:r4-counting}(a)).
\end{itemize}
(d)
\begin{itemize}
\item As \(F_m\ne0\), there is \(y\) with \(x=F_my\ne0\); then \(Ax=\sum_lb_lF_lF_my=b_mx\). So every \(b_m\) is an
eigenvalue of \(A\), and \(\mathcal C(F_m)\subseteq\mathcal N(A-b_mI)\).
\item Let \(Ax=cx\) with \(x\ne0\). Then \(x=\sum_lF_lx\) and \(0=Ax-cx=\sum_l(b_l-c)F_lx\). The vectors \(F_lx\) are
pairwise orthogonal (\((F_lx)'(F_mx)=x'F_lF_mx=0\)), so
\(0=\lVert\sum_l(b_l-c)F_lx\rVert^2=\sum_l(b_l-c)^2\lVert F_lx\rVert^2\). Hence \(F_lx=0\) whenever \(b_l\ne c\). As
\(x\ne0\), some \(F_lx\ne0\); so \(c\) is one of the \(b_l\), and \(x=F_lx\in\mathcal C(F_l)\) for that \(l\).
\item So \(\{b_1,\dots,b_r\}=\{c_1,\dots,c_k\}\), \(r=k\), and after renumbering \(b_j=c_j\) and
\(\mathcal C(F_j)=W_j\). \(F_j\) is symmetric and idempotent, so \(F_jy\) is the orthogonal projection of \(y\) onto
\(\mathcal C(F_j)=W_j\) (Theorem~\ref{r6-thm-idempotent}(d)): \(F_j=E_j\).\qedhere
\end{itemize}
\end{proof}

\begin{example}[cf.\ \citeat{searle-khuri}{Example~6.13}]\label{r8-exm-decomp}
For \(A=2J_3-I_3\) of Example~\ref{r8-exm-spectral-j3}, with the distinct eigenvalues \(c_1=5\) and \(c_2=-1\),
Theorem~\ref{r8-thm-spectral-decomp}(c) gives \(E_1=q_1q_1'=\frac13\one\one'=\bar J_3\), and then \(E_2=I_3-E_1=C_3\) by
(b); indeed
\[
q_2q_2'+q_3q_3'=\frac16\begin{pmatrix}4&-2&-2\\-2&1&1\\-2&1&1\end{pmatrix}+\frac12\begin{pmatrix}0&0&0\\0&1&-1\\0&-1&1\end{pmatrix}
=\frac13\begin{pmatrix}2&-1&-1\\-1&2&-1\\-1&-1&2\end{pmatrix}=C_3
\]
So \(A=5\bar J_3-C_3\), \(A^m=5^m\bar J_3+(-1)^mC_3\); for instance \(A^2=25\bar J_3+C_3=I_3+8J_3\). As no eigenvalue is
\(0\),
\[
A^{-1}=\tfrac15\bar J_3-C_3=\tfrac25J_3-I_3=\frac15\begin{pmatrix}-3&2&2\\2&-3&2\\2&2&-3\end{pmatrix}
\]
\end{example}

\begin{example}\label{r8-exm-spectral}
Find an orthogonal \(Q\) with \(Q'AQ\) diagonal, and the matrices \(E_j\) of Theorem~\ref{r8-thm-spectral-decomp}, for
\[
A=\begin{pmatrix}11&2&6&6\\2&11&6&6\\6&6&11&2\\6&6&2&11\end{pmatrix}
\]

\textit{Solution.}
\begin{itemize}
\item For real \(t\), in \(tI_4-A\) subtract row \(2\) from row \(1\) and row \(4\) from row \(3\); the new rows \(1\) and
\(3\) are \((t-9)(1,-1,0,0)\) and \((t-9)(0,0,1,-1)\). Take the two factors \(t-9\) out, add column \(1\) to column \(2\)
and column \(3\) to column \(4\) (Definition~\ref{r5d-def-det}(b),(a), Theorem~\ref{r5d-thm-transpose}), and expand along
row \(1\) and then along the row \((0,1,0)\) (Theorem~\ref{r5d-thm-expansion}):
\[
\begin{aligned}
\det(tI_4-A)&=\det\begin{pmatrix}t-11&-2&-6&-6\\-2&t-11&-6&-6\\-6&-6&t-11&-2\\-6&-6&-2&t-11\end{pmatrix}
=(t-9)^2\det\begin{pmatrix}1&-1&0&0\\-2&t-11&-6&-6\\0&0&1&-1\\-6&-6&-2&t-11\end{pmatrix}\\
&=(t-9)^2\det\begin{pmatrix}1&0&0&0\\-2&t-13&-6&-12\\0&0&1&0\\-6&-12&-2&t-13\end{pmatrix}
=(t-9)^2\det\begin{pmatrix}t-13&-6&-12\\0&1&0\\-12&-2&t-13\end{pmatrix}\\
&=(t-9)^2\bigl((t-13)^2-144\bigr)=(t-9)^2(t-25)(t-1)
\end{aligned}
\]
This holds for every real \(t\), hence as polynomials (Definition~\ref{r7-def-charpoly}(a),
Proposition~\ref{r7-prop-poly}(d)): the eigenvalues are \(c_1=25\), \(c_2=9\) and \(c_3=1\), with \(a_A(25)=a_A(1)=1\) and
\(a_A(9)=2\).
\item The row sums of \(A\) are \(25\), so \(A\one_4=25\,\one_4\); and \(v=(1,1,-1,-1)'\) has
\(Av=(11+2-12,\ 2+11-12,\ 12-11-2,\ 12-2-11)'=v\). By Theorem~\ref{r8-thm-rank-eigen},
\(\dim\mathcal N(A-25I)=\dim\mathcal N(A-I)=1\), so \(q_1=\frac12\one_4\) and \(q_4=\frac12v\) are orthonormal bases of
\(\mathcal N(A-25I)\) and \(\mathcal N(A-I)\).
\item \(\mathcal N(A-9I)\): reduce \(A-9I\),
\[
\begin{gathered}
A-9I=\begin{pmatrix}2&2&6&6\\2&2&6&6\\6&6&2&2\\6&6&2&2\end{pmatrix}
\xrightarrow[R_4\leftarrow R_4-R_3]{R_2\leftarrow R_2-R_1}
\begin{pmatrix}2&2&6&6\\0&0&0&0\\6&6&2&2\\0&0&0&0\end{pmatrix}
\xrightarrow{R_3\leftarrow R_3-3R_1}
\begin{pmatrix}2&2&6&6\\0&0&0&0\\0&0&-16&-16\\0&0&0&0\end{pmatrix}
\\
\xrightarrow{R_2\leftrightarrow R_3}
\begin{pmatrix}2&2&6&6\\0&0&-16&-16\\0&0&0&0\\0&0&0&0\end{pmatrix}
\xrightarrow[R_2\leftarrow-\frac1{16}R_2]{R_1\leftarrow\frac12R_1}
\begin{pmatrix}1&1&3&3\\0&0&1&1\\0&0&0&0\\0&0&0&0\end{pmatrix}
\xrightarrow{R_1\leftarrow R_1-3R_2}
\begin{pmatrix}1&1&0&0\\0&0&1&1\\0&0&0&0\\0&0&0&0\end{pmatrix}
\end{gathered}
\]
So \(\mathcal N(A-9I)=\{x:x_1+x_2=0,\ x_3+x_4=0\}\), with the free unknowns \(x_2\), \(x_4\) and the special
solutions \(u_2=(-1,1,0,0)'\) and \(u_3=(0,0,-1,1)'\). They are orthogonal, so \(q_2=\frac1{\sqrt2}u_2\) and
\(q_3=\frac1{\sqrt2}u_3\) form an orthonormal basis of \(\mathcal N(A-9I)\).
\item \(q_1,q_4\perp q_2,q_3\) (Theorem~\ref{r8-thm-sym-orth}), and \(q_1'q_4=0\). So \(Q=[q_1\ q_2\ q_3\ q_4]\) is
orthogonal and \(Q'AQ=\Lambda=\diag(25,9,9,1)\) (Theorem~\ref{r8-thm-spectral}(b)).
\item By Theorem~\ref{r8-thm-spectral-decomp}(c),
\[
E_1=q_1q_1'=\tfrac14J_4,\qquad
E_2=q_2q_2'+q_3q_3'=\frac12\begin{pmatrix}1&-1&0&0\\-1&1&0&0\\0&0&1&-1\\0&0&-1&1\end{pmatrix},\qquad
E_3=q_4q_4'=\tfrac14vv'
\]
\(E_1+E_2+E_3=I_4\), and \(A=25E_1+9E_2+E_3\); for instance the \((1,1)\), \((1,2)\) and \((1,3)\) entries are
\(\frac{25}4+\frac92+\frac14=11\), \(\frac{25}4-\frac92+\frac14=2\) and \(\frac{25}4-\frac14=6\).
\end{itemize}
The vectors \(\frac12(1,-1,1,-1)'\) and \(\frac12(1,-1,-1,1)'\) also form an orthonormal basis of \(\mathcal N(A-9I)\).
They give another \(Q\), with entries \(\pm\frac12\), but, by Theorem~\ref{r8-thm-spectral-decomp}(c), the same projection
matrix
\[
E_2=\tfrac14\bigl((1,-1,1,-1)'(1,-1,1,-1)+(1,-1,-1,1)'(1,-1,-1,1)\bigr)
\]
\end{example}

\subsection*{Functions of a symmetric matrix}

\begin{definition}[\citeat{hoffman-kunze}{\S9.5, (9-12)}]\label{r8-def-matrix-function}
Let \(A\in\R^{n\times n}\) be symmetric with spectral decomposition \(A=c_1E_1+\dots+c_kE_k\)
(Theorem~\ref{r8-thm-spectral-decomp}), and let \(f\) be a real-valued function defined at \(c_1,\dots,c_k\). The matrix
\[
f(A)=f(c_1)E_1+\dots+f(c_k)E_k
\]
is the \textbf{function} \(f\) of \(A\).
\end{definition}

\needspace{14\baselineskip}
\begin{proposition}[\citeat{hoffman-kunze}{\S9.5, Theorems~9, 10}]\label{r8-prop-function}
Let \(A=c_1E_1+\dots+c_kE_k\) be as in Definition~\ref{r8-def-matrix-function}, let \(W_j=\mathcal N(A-c_jI)\), and let \(f\),
\(g\) be real-valued functions defined at \(c_1,\dots,c_k\).
\begin{enumerate}
\item \(f(A)\) is symmetric, and \(f(A)x=f(c_j)x\) for every \(x\in W_j\). If \(A=Q\Lambda Q'\) as in
Theorem~\ref{r8-thm-spectral}, then \(f(A)=Q\,f(\Lambda)\,Q'\) with \(f(\Lambda)=\diag\bigl(f(\lambda_1),\dots,f(\lambda_n)\bigr)\),
and the eigenvalues of \(f(A)\), listed with multiplicity, are \(f(\lambda_1),\dots,f(\lambda_n)\).
\item \(f(A)+g(A)=(f+g)(A)\), \(af(A)=(af)(A)\) for \(a\in\R\), and \(f(A)\,g(A)=(fg)(A)\). If \(f(c_j)\ne0\) for every \(j\),
then \(f(A)\) is invertible and \(f(A)^{-1}=h(A)\) for \(h=1/f\).
\item If \(f\) is a real polynomial, then \(f(A)\) is the matrix \(f(A)\) of Definition~\ref{r7-def-polynomial}. With the
polynomials \(e_j=\prod_{i\ne j}\frac{t-c_i}{c_j-c_i}\) (\(e_1=1\) if \(k=1\)), \(E_j=e_j(A)\) for every \(j\); hence
\(f(A)=p(A)\) for the polynomial \(p=f(c_1)e_1+\dots+f(c_k)e_k\): every \(f(A)\) is a polynomial in \(A\).
\end{enumerate}
\end{proposition}
\begin{proof}
(a)
\begin{itemize}
\item \(f(A)'=\sum_jf(c_j)E_j'=f(A)\), as \(E_j'=E_j\) (Theorem~\ref{r8-thm-spectral-decomp}(b)).
\item Let \(x\in W_j\). Then \(x\) is its own decomposition in Theorem~\ref{r8-thm-spectral-decomp}(a), so \(E_jx=x\) and
\(E_ix=0\) for \(i\ne j\); hence \(f(A)x=f(c_j)x\).
\item Each \(q_i\) lies in the \(W_j\) with \(c_j=\lambda_i\), so \(f(A)q_i=f(\lambda_i)q_i\), that is, \(f(A)Q=Q\,f(\Lambda)\)
(Theorem~\ref{r1-thm-views}(a),(b)); as \(QQ'=I_n\) (Proposition~\ref{r6-prop-orthonormal}(b)), \(f(A)=Q\,f(\Lambda)\,Q'\).
\item So \(q_1,\dots,q_n\) is an orthonormal basis of eigenvectors of the symmetric matrix \(f(A)\), with
\(f(A)q_i=f(\lambda_i)q_i\), and Theorem~\ref{r8-thm-spectral}(b),(c), applied to \(f(A)\), gives its eigenvalues.
\end{itemize}
(b)
\begin{itemize}
\item The first two rules hold term by term. As \(E_iE_j=0\) for \(i\ne j\) and \(E_j^2=E_j\)
(Theorem~\ref{r8-thm-spectral-decomp}(b)),
\(f(A)\,g(A)=\sum_i\sum_jf(c_i)g(c_j)E_iE_j=\sum_jf(c_j)g(c_j)E_j=(fg)(A)\).
\item If every \(f(c_j)\ne0\), then \(h=1/f\) is defined at every \(c_j\), and \(f(A)\,h(A)=(fh)(A)=\sum_jE_j=I_n\)
(Theorem~\ref{r8-thm-spectral-decomp}(b)); so \(f(A)^{-1}=h(A)\) (Theorem~\ref{r2-thm-invertible}).
\end{itemize}
(c)
\begin{itemize}
\item Let \(f=a_0+a_1t+\dots+a_rt^r\). Since \(\sum_jc_j^mE_j=A^m\) (Theorem~\ref{r8-thm-spectral-decomp}(b)), the matrix of
Definition~\ref{r8-def-matrix-function} is \(\sum_j\sum_ma_mc_j^mE_j=\sum_ma_mA^m\), the matrix of
Definition~\ref{r7-def-polynomial}.
\item The \(c_i\) are real (Theorem~\ref{r8-thm-sym-real}), so each \(e_j\) is a real polynomial, and \(e_j(c_l)\) is \(1\) for
\(l=j\) and \(0\) for \(l\ne j\) (Proposition~\ref{r7-prop-poly}(b)). By the first item, \(e_j(A)=\sum_le_j(c_l)E_l=E_j\).
\item Hence \(f(A)=\sum_jf(c_j)E_j=\sum_jf(c_j)e_j(A)=p(A)\) (Proposition~\ref{r7-prop-poly}(b)).\qedhere
\end{itemize}
\end{proof}

\begin{example}\label{r8-exm-function-j3}
Let \(A=2J_3-I_3=5\bar J_3-C_3\), with \(c_1=5\), \(c_2=-1\), \(E_1=\bar J_3\) and \(E_2=C_3\) (Examples~\ref{r8-exm-spectral-j3}
and~\ref{r8-exm-decomp}).
\begin{itemize}
\item \(e_1=\frac{t+1}6\) and \(e_2=\frac{5-t}6\), so by Proposition~\ref{r8-prop-function}(c)
\(\bar J_3=\frac16(A+I_3)\) and \(C_3=\frac16(5I_3-A)\). Indeed \(A+I_3=2J_3\) and \(5I_3-A=6I_3-2J_3\).
\item For \(f(t)=\lvert t\rvert\),
\[
f(A)=5\bar J_3+C_3=I_3+\tfrac43J_3=\frac13\begin{pmatrix}7&4&4\\4&7&4\\4&4&7\end{pmatrix}
\]
By Proposition~\ref{r8-prop-function}, \(f(A)\) is symmetric with the eigenvalues \(5\), \(1\), \(1\); \(f(A)^2=25\bar J_3+C_3=A^2\);
and \(f(A)=p(A)\) for \(p=5e_1+e_2=\frac{2t+5}3\): indeed \(\frac13(2A+5I_3)=\frac13(4J_3+3I_3)\).
\end{itemize}
\end{example}

\begin{example}\label{r8-exm-function}
For \(A=25E_1+9E_2+E_3\) of Example~\ref{r8-exm-spectral}, Theorem~\ref{r8-thm-spectral-decomp}(b) gives
\[
A^k=25^kE_1+9^kE_2+E_3\quad(k\geqslant0),\qquad
A^{-1}=\tfrac1{25}E_1+\tfrac19E_2+E_3=\frac1{225}\begin{pmatrix}71&46&-54&-54\\46&71&-54&-54\\
-54&-54&71&46\\-54&-54&46&71\end{pmatrix}
\]
Check: row \(1\) of \(A\) times the first three columns of \(225A^{-1}\) gives \(781+92-324-324=225\),
\(506+142-324-324=0\) and \(-594-108+426+276=0\); the other entries in the same way. The polynomial
\(f(t)=(t-25)(t-9)(t-1)\) vanishes at \(25\), \(9\) and \(1\), so \(f(A)=f(25)E_1+f(9)E_2+f(1)E_3=0\)
(Definition~\ref{r8-def-matrix-function}). As \(f(t)=t^3-35t^2+259t-225\), Proposition~\ref{r8-prop-function}(c) gives
\(f(A)=A^3-35A^2+259A-225I_4\), which is \((A-25I_4)(A-9I_4)(A-I_4)\) (Proposition~\ref{r7-prop-poly}(b)); hence
\((A-25I_4)(A-9I_4)(A-I_4)=0\).
\end{example}

\subsection*{Idempotent, orthogonal, and patterned matrices}

\begin{theorem}[\citeat{harville}{Theorems~21.8.2, 21.8.3}, \citeat{hoffman-kunze}{\S6.6}, \citeat{abadir-magnus}{Exercise~8.66(b)}]\label{r8-thm-idempotent-eigen}
\begin{enumerate}
\item Let \(A\in\R^{n\times n}\) be idempotent with \(\rank A=r\). Then \(p_A(t)=t^{n-r}(t-1)^r\): the only eigenvalues of
\(A\) are \(0\) and \(1\), with \(g_A(0)=a_A(0)=n-r\) and \(g_A(1)=a_A(1)=r\) (each an eigenvalue only when this number
is positive). The eigenspaces are \(\mathcal N(A)\) and \(\mathcal N(A-I)=\mathcal C(A)\), and \(A\) is diagonalizable.
\item A symmetric matrix whose eigenvalues all lie in \(\{0,1\}\) is idempotent.
\item If \(A\in\R^{n\times n}\) is symmetric and idempotent with \(\rank A=r\), then \(Q'AQ=\diag(1,\dots,1,0,\dots,0)\)
(\(r\) entries \(1\)) for some orthogonal \(Q\); if \(r\geqslant1\), then \(A=Q_1Q_1'\), where \(Q_1\in\R^{n\times r}\)
consists of the first \(r\) columns of \(Q\), and \(Q_1'Q_1=I_r\).
\end{enumerate}
\end{theorem}
\begin{proof}
(a)
\begin{itemize}
\item If \(Ax=\lambda x\) with \(\lambda\in\R\) and \(x\ne0\), then \(\lambda x=Ax=A^2x=\lambda Ax=\lambda^2x\), so
\(\lambda^2=\lambda\): \(\lambda\) is \(0\) or \(1\).
\item \(\dim\mathcal N(A)=n-r\) (Theorem~\ref{thm:r4-rank-nullity}). \(I-A\) is idempotent, so
\(\rank(I-A)=\tr(I-A)=n-\tr A=n-r\) (Theorems~\ref{r6-thm-idempotent}(b),(c) and~\ref{r1-thm-trace}(a)), and
\(\dim\mathcal N(A-I)=n-\rank(I-A)=r\).
\item For real \(c\), \(\dim\mathcal N(A-cI)\leqslant a_A(c)\): by Theorem~\ref{r7-thm-geo-alg} if \(c\) is an
eigenvalue, and as \(\mathcal N(A-cI)=\{0\}\) otherwise (Theorem~\ref{r7-thm-eig-singular}). The multiplicities of all
eigenvalues add up to \(n\) (Definition~\ref{r7-def-eigenvalues}). Hence
\[
n=\dim\mathcal N(A)+\dim\mathcal N(A-I)\leqslant a_A(0)+a_A(1)\leqslant n
\]
so \(a_A(0)=n-r\), \(a_A(1)=r\), and \(A\) has no other eigenvalue, real or not: \(p_A(t)=t^{n-r}(t-1)^r\)
(Definition~\ref{r7-def-eigenvalues}).
\item \(Av=v\) for every \(v\in\mathcal C(A)\) (Theorem~\ref{r6-thm-idempotent}(a)), and \(v=Av\in\mathcal C(A)\) if
\(Av=v\); so \(\mathcal N(A-I)=\mathcal C(A)\). As \(\dim\mathcal N(A)+\dim\mathcal N(A-I)=n\), \(A\) is diagonalizable
(Theorem~\ref{r7-thm-diag-criterion}).
\end{itemize}
(b) By Theorem~\ref{r8-thm-spectral}, \(A=Q\Lambda Q'\) with \(Q'Q=I_n\), and the diagonal entries of \(\Lambda\) are
\(0\) or \(1\), so \(\Lambda^2=\Lambda\). Then \(A^2=Q\Lambda Q'Q\Lambda Q'=Q\Lambda^2Q'=A\).

(c) By Theorem~\ref{r8-thm-spectral}, \(Q'AQ=\Lambda\) with \(Q\) orthogonal. By (a) every \(\lambda_i\) is \(0\) or \(1\),
and exactly \(r\) of them are \(1\), as \(r=\rank A\) is the number of nonzero \(\lambda_i\)
(Theorem~\ref{r8-thm-rank-eigen}). Numbering the columns of \(Q\) so that \(\lambda_1=\dots=\lambda_r=1\)
(Theorem~\ref{r8-thm-spectral}(b)) gives \(\Lambda=\diag(1,\dots,1,0,\dots,0)\). If \(r\geqslant1\), the outer product
expansion (Theorem~\ref{r1-thm-views}(d)) gives \(A=(Q\Lambda)Q'=q_1q_1'+\dots+q_rq_r'=Q_1Q_1'\), as
\(\col_i(Q\Lambda)=\lambda_iq_i\); and \(Q_1'Q_1=I_r\), as the columns of \(Q\) are orthonormal.
\end{proof}
Without symmetry (b) fails~\citeat{abadir-magnus}{Exercise~8.56(b)}: the upper triangular matrix
\[
N=\begin{pmatrix}1&1&1\\0&1&1\\0&0&0\end{pmatrix},\qquad N^2=\begin{pmatrix}1&2&2\\0&1&1\\0&0&0\end{pmatrix}
\]
has \(p_N(t)=(t-1)^2t\) (Proposition~\ref{r7-prop-block-triangular}(b)), but \(N^2\ne N\).

\begin{example}\label{r8-exm-idempotent}
Let
\[
E=\begin{pmatrix}1&1&-1&1\\1&0&1&0\\0&-1&2&-1\\-1&-1&1&-1\end{pmatrix}
\]
Show that \(E\) is idempotent, and find an invertible \(P\) with \(P^{-1}EP\) diagonal.

\textit{Solution.}
\begin{itemize}
\item Row \(1\) of \(E\) times the columns of \(E\) gives \(1+1+0-1=1\), \(1+0+1-1=1\), \(-1+1-2+1=-1\) and
\(1+0+1-1=1\), the entries of row \(1\) of \(E\); the other rows in the same way, so \(E^2=E\). By
Theorem~\ref{r6-thm-idempotent}(c), \(\rank E=\tr E=1+0+2-1=2\).
\item The first two columns \(p_1=(1,1,0,-1)'\) and \(p_2=(1,0,-1,-1)'\) of \(E\) are not multiples of each other, so they
are linearly independent; as \(\dim\mathcal C(E)=2\), they form a basis of \(\mathcal C(E)=\mathcal N(E-I)\)
(Theorems~\ref{thm:r4-counting}(d) and~\ref{r8-thm-idempotent-eigen}(a)).
\item The vectors \(p_3=(-1,2,1,0)'\) and \(p_4=(0,-1,0,1)'\) satisfy \(Ep_3=Ep_4=0\) (multiply out), and they are not
multiples of each other, so they are linearly independent; as \(\dim\mathcal N(E)=4-2=2\)
(Theorem~\ref{thm:r4-rank-nullity}), they form a basis of \(\mathcal N(E)\) (Theorem~\ref{thm:r4-counting}(d)).
\item By Lemma~\ref{r7-lem-eigenspaces}(b), \(p_1,\dots,p_4\) are linearly independent, so
\(P=[p_1\ p_2\ p_3\ p_4]\) is invertible (Theorem~\ref{thm:r4-rank-rules}(b),(d)), and by
Definition~\ref{r7-def-diagonalizable}
\[
P^{-1}EP=\diag(1,1,0,0),\qquad P=\begin{pmatrix}1&1&-1&0\\1&0&2&-1\\0&-1&1&0\\-1&-1&0&1\end{pmatrix}
\]
Also \(p_E(t)=t^2(t-1)^2\) (Theorem~\ref{r8-thm-idempotent-eigen}(a)). Here \(\det P=1\), so \(P^{-1}\) has integer entries
(Theorem~\ref{r5d-thm-adjugate}).
\item \(E\) is not symmetric (its \((1,3)\) entry is \(-1\), its \((3,1)\) entry \(0\)), so no orthogonal \(Q\) makes
\(Q'EQ\) diagonal (Theorem~\ref{r8-thm-spectral}(b)): \(E\) is diagonalizable, but not by an orthogonal matrix.
\end{itemize}
\end{example}

\begin{proposition}[\citeat{harville}{Theorem~21.8.1}]\label{r8-prop-orthogonal-eigen}
If \(P\in\R^{n\times n}\) is orthogonal and \(\lambda\) is a real eigenvalue of \(P\), then \(\lambda=1\) or
\(\lambda=-1\).
\end{proposition}
\begin{proof}
Let \(Px=\lambda x\) with \(x\ne0\). Then \(x'x=x'P'Px=(Px)'(Px)=\lambda^2x'x\), and \(x'x>0\); so \(\lambda^2=1\).
\end{proof}
An orthogonal matrix need not have a real eigenvalue:
\(\left(\begin{smallmatrix}0&-1\\1&0\end{smallmatrix}\right)\) is orthogonal, with characteristic polynomial \(t^2+1\).

\begin{example}[\citeat{searle-khuri}{Exercise~6.15}, \citeat{abadir-magnus}{Exercise~8.74}]\label{r8-exm-patterned}
Let \(n\geqslant2\).
\begin{enumerate}
\item \(\bar J_n\) and \(C_n\) are symmetric and idempotent (Proposition~\ref{r1-prop-centering}), of ranks
\(\tr\bar J_n=1\) and \(n-1\) (Theorem~\ref{r6-thm-idempotent}(c), Corollary~\ref{r6-cor-trace-hat}). By
Theorem~\ref{r8-thm-idempotent-eigen}(a), \(\bar J_n\) has the eigenvalues \(1\) (once) and \(0\) (\(n-1\) times), and
\(C_n\) the eigenvalues \(1\) (\(n-1\) times) and \(0\) (once). Since \(\bar J_nx=\bar x\one\) and
\(C_nx=x-\bar x\one\), the eigenspaces are \(\mathcal N(\bar J_n-I_n)=\mathcal N(C_n)=\Span\{\one\}\) and
\(\mathcal N(\bar J_n)=\mathcal N(C_n-I_n)=\{x:\one'x=0\}=\Span\{\one\}^\perp\) (Theorem~\ref{r6-thm-complement}(a)).
\item \(\bar J_n+C_n=I_n\) and \(\bar J_nC_n=C_n\bar J_n=\bar J_n-\bar J_n^2=0\). So, for numbers \(a\ne b\),
\(a\bar J_n+bC_n\) satisfies the hypotheses of Theorem~\ref{r8-thm-spectral-decomp}(d): its eigenvalues are \(a\) and
\(b\). As \(a\bar J_n+bC_n-aI_n=(b-a)C_n\) and \(a\bar J_n+bC_n-bI_n=(a-b)\bar J_n\), with \(a-b\ne0\), by (a) the
eigenspace of \(a\) is \(\mathcal N(C_n)=\Span\{\one\}\) and that of \(b\) is \(\mathcal N(\bar J_n)=\Span\{\one\}^\perp\); so
\(a\) has multiplicity \(1\) and \(b\) multiplicity \(n-1\) (Theorems~\ref{r8-thm-rank-eigen}
and~\ref{r6-thm-complement}(c)). For instance
\(J_n=n\bar J_n+0\,C_n\) has the eigenvalues \(n\) (once) and \(0\) (\(n-1\) times), as found in
Example~\ref{r7-exm-ones}; and for numbers \(\alpha\), \(\beta\) with \(\beta\ne0\),
\(\alpha I_n+\beta J_n=(\alpha+n\beta)\bar J_n+\alpha C_n\) (as \(I_n=\bar J_n+C_n\) and \(J_n=n\bar J_n\)) has the
eigenvalues \(\alpha+n\beta\), with eigenspace \(\Span\{\one\}\), and \(\alpha\), of multiplicity \(n-1\).
\item The equicorrelation matrix \(R=(1-\rho)I_n+\rho J_n\) of Example~\ref{r5d-exm-equicorrelation} is, writing
\(I_n=\bar J_n+C_n\) and \(J_n=n\bar J_n\),
\[
R=\bigl(1+(n-1)\rho\bigr)\bar J_n+(1-\rho)C_n
\]
For \(\rho\ne0\), by (b), its eigenvalues are \(1+(n-1)\rho\), with eigenvector \(\one\), and \(1-\rho\), of multiplicity
\(n-1\). So \(\det R=\bigl(1+(n-1)\rho\bigr)(1-\rho)^{n-1}\) (Theorem~\ref{r7-thm-trace-det}), as in
Example~\ref{r5d-exm-equicorrelation}, and for \(\rho\ne1\), \(\rho\ne-1/(n-1)\),
\(R^{-1}=\bigl(1+(n-1)\rho\bigr)^{-1}\bar J_n+(1-\rho)^{-1}C_n\) (Theorem~\ref{r8-thm-spectral-decomp}(b)). For \(n=4\)
and \(\rho=\frac13\), \(R=2\bar J_4+\frac23C_4\) has the eigenvalues \(2\) and \(\frac23\) (three times),
\(\det R=2\cdot\bigl(\frac23\bigr)^3=\frac{16}{27}\), and \(R^{-1}=\frac12\bar J_4+\frac32C_4=\frac32I_4-\frac14J_4\), as
in Example~\ref{r5d-exm-equicorrelation}.
\item Let \(W\) be a subspace of \(\R^n\) with \(\dim W=r\). The projection matrix \(P_W\) is symmetric and idempotent,
of rank \(r\) (Theorem~\ref{r6-thm-projection-matrix}(c)). So \(P_W\) has the eigenvalue \(1\) with multiplicity \(r\)
and \(0\) with multiplicity \(n-r\) (Theorem~\ref{r8-thm-idempotent-eigen}(a)), and \(I_n-P_W=P_{W^\perp}\) the reverse;
\(P_Wy=y\) exactly for \(y\in W\), and \(P_Wy=0\) exactly for \(y\in W^\perp\) (Theorem~\ref{r6-thm-projection-matrix}(c),
Definition~\ref{r6-def-projection}). For instance, the matrix \(P\) of Example~\ref{r6-exm-projection} has the
eigenvalues \(1\), \(1\), \(1\), \(0\), as \(\rank P=\tr P=3\) (Corollary~\ref{r6-cor-trace-hat}); and for
\(W=\Span\{\one_n\}\), \(P_W=\bar J_n\) and \(I_n-P_W=C_n\) (Corollary~\ref{r6-cor-centering}(a)), which is (a).
\end{enumerate}
\end{example}

\subsection*{Quadratic forms and definiteness}

\begin{definition}[\citeat{searle-khuri}{\S5.6}, \citeat{harville}{\S14.1}, \citeat{abadir-magnus}{Exercise~8.1(a)}]\label{r8-def-quadratic-form}
The \textbf{quadratic form} of \(A\in\R^{n\times n}\) is the function \(q_A(x)=x'Ax=\sum_{i=1}^n\sum_{j=1}^na_{ij}x_ix_j\)
(\(x\in\R^n\)). Collecting the terms in \(x_i^2\) and in \(x_ix_j\) (\(i<j\)),
\[
x'Ax=\sum_{i=1}^na_{ii}x_i^2+\sum_{i<j}(a_{ij}+a_{ji})x_ix_j
\]
\end{definition}
As \(x'Ax\) is a \(1\times1\) matrix, \(x'Ax=(x'Ax)'=x'A'x\); so \(x'Ax=x'Bx\) for the symmetric matrix
\(B=\frac12(A+A')\). Different matrices can give the same form, but only one symmetric matrix does.

\begin{proposition}[\citeat{lang}{Ch.~V, \S7}]\label{r8-prop-qf-symmetric}
Let \(A,B\in\R^{n\times n}\) be symmetric.
\begin{enumerate}
\item \(x'Ay=\frac12\bigl(q_A(x+y)-q_A(x)-q_A(y)\bigr)\) for all \(x,y\in\R^n\).
\item If \(x'Ax=x'Bx\) for every \(x\in\R^n\), then \(A=B\).
\end{enumerate}
\end{proposition}
\begin{proof}
(a) \((x+y)'A(x+y)=x'Ax+x'Ay+y'Ax+y'Ay\), and \(y'Ax=(y'Ax)'=x'A'y=x'Ay\).

(b) By (a), \(x'Ay=x'By\) for all \(x,y\). For \(x=e_i\) and \(y=e_j\) this reads \(a_{ij}=b_{ij}\), since
\(e_i'Ae_j\) is entry \(i\) of \(Ae_j=\col_j(A)\) (Theorem~\ref{r1-thm-views}(b)).
\end{proof}
So every quadratic form is \(q_A\) for exactly one symmetric \(A\), the \textbf{matrix of the form}: \(a_{ii}\) is the
coefficient of \(x_i^2\), and \(a_{ij}=a_{ji}\) is half the coefficient of \(x_ix_j\) (\(i\ne j\)).

\begin{example}[cf.\ \citeat{searle-khuri}{Example~5.6, \S5.6}]\label{r8-exm-qf}
\begin{enumerate}
\item For data \(y\in\R^n\), the sum of squares about the mean \(\sum_i(y_i-\bar y)^2\) equals \(y'C_ny\)
(Example~\ref{r1-exm-centering}): it is the quadratic form of the symmetric matrix \(C_n\).
\item The matrices
\[
T=\begin{pmatrix}2&0&2&-2\\0&2&2&2\\0&0&3&0\\0&0&0&3\end{pmatrix},\qquad
T'=\begin{pmatrix}2&0&0&0\\0&2&0&0\\2&2&3&0\\-2&2&0&3\end{pmatrix},\qquad
A=\tfrac12(T+T')=\begin{pmatrix}2&0&1&-1\\0&2&1&1\\1&1&3&0\\-1&1&0&3\end{pmatrix}
\]
give the same form \(2x_1^2+2x_2^2+3x_3^2+3x_4^2+2x_1x_3-2x_1x_4+2x_2x_3+2x_2x_4\)
(Definition~\ref{r8-def-quadratic-form}: \(t_{12}+t_{21}=0\), \(t_{13}+t_{31}=2\), \(t_{14}+t_{41}=-2\),
\(t_{23}+t_{32}=2\), \(t_{24}+t_{42}=2\), \(t_{34}+t_{43}=0\)); its matrix is \(A\).
\end{enumerate}
\end{example}

\needspace{7\baselineskip}
\begin{definition}[\citeat{abadir-magnus}{Ch.~8}, \citeat{bapat}{\S3.1}]\label{r8-def-definite}
A symmetric \(A\in\R^{n\times n}\), and its quadratic form, are
\begin{itemize}
\item \textbf{positive definite} (PD) if \(x'Ax>0\) for every nonzero \(x\in\R^n\);
\item \textbf{positive semidefinite} (PSD) if \(x'Ax\geqslant0\) for every \(x\in\R^n\);
\item \textbf{negative definite} (ND) or \textbf{negative semidefinite} (NSD) if \(-A\) is PD or PSD;
\item \textbf{indefinite} if \(x'Ax\) takes both positive and negative values.
\end{itemize}
A PD matrix is PSD.
\end{definition}

\begin{theorem}[\citeat{abadir-magnus}{Exercises~7.53(a), 8.11, 8.13(a)}, \citeat{hoffman-kunze}{\S9.5, Theorem~12}, \citeat{harville}{Corollary~14.3.12}]\label{r8-thm-eigen-criterion}
Let \(A\in\R^{n\times n}\) be symmetric, with spectral decomposition \(A=c_1E_1+\dots+c_kE_k\), and let \(A=Q\Lambda Q'\) as
in Theorem~\ref{r8-thm-spectral}(b), with \(\lambda_1\geqslant\dots\geqslant\lambda_n\). Then, for \(x\in\R^n\) and
\(y=Q'x\),
\[
x'Ax=\lambda_1y_1^2+\dots+\lambda_ny_n^2=c_1\lVert E_1x\rVert^2+\dots+c_k\lVert E_kx\rVert^2
\]
\begin{enumerate}
\item \(A\) is PSD if and only if every eigenvalue of \(A\) is \(\geqslant0\), and PD if and only if every eigenvalue is
\(>0\).
\item \(A\) is NSD (ND) if and only if every eigenvalue is \(\leqslant0\) (\(<0\)), and indefinite if and only if \(A\) has a
positive and a negative eigenvalue.
\item If \(A\) is PD, then \(\det A>0\), and \(A\) is invertible; if \(A\) is PSD, then \(\det A\geqslant0\). A PSD matrix
is PD if and only if it is invertible.
\item \(\lambda_n\,x'x\leqslant x'Ax\leqslant\lambda_1\,x'x\) for every \(x\in\R^n\), with equality on the right for
\(x=q_1\) and on the left for \(x=q_n\).
\end{enumerate}
\end{theorem}
\begin{proof}
\begin{itemize}
\item Let \(x\in\R^n\) and \(y=Q'x\). Then \(x=QQ'x=Qy\) and \(y'y=x'QQ'x=x'x\) (Proposition~\ref{r6-prop-orthonormal}(b)),
and, as \(Q'AQ=\Lambda\),
\[
x'Ax=x'QQ'AQQ'x=y'\Lambda y=\lambda_1y_1^2+\dots+\lambda_ny_n^2
\]
\item Entry \(i\) of \(y=Q'x\) is \(y_i=q_i'x\) (Theorem~\ref{r1-thm-views}(d)). As \(E_j'=E_j=E_j^2\) and
\(E_j=\sum_{\lambda_i=c_j}q_iq_i'\) (Theorem~\ref{r8-thm-spectral-decomp}(b),(c)),
\(\lVert E_jx\rVert^2=x'E_j'E_jx=x'E_jx=\sum_{\lambda_i=c_j}(q_i'x)^2=\sum_{\lambda_i=c_j}y_i^2\). Grouping the terms
\(\lambda_iy_i^2\) by the value \(c_j\) of \(\lambda_i\) gives the second expression.
\end{itemize}
(a)
\begin{itemize}
\item If every \(\lambda_i\geqslant0\), then \(x'Ax\geqslant0\). If every \(\lambda_i>0\) and \(x\ne0\), then \(y\ne0\), as
\(x=Qy\); so \(x'Ax>0\).
\item Conversely, for \(x=q_i\), \(y=Q'Qe_i=e_i\) and \(x'Ax=\lambda_i\); so \(\lambda_i\geqslant0\) if \(A\) is PSD, and
\(\lambda_i>0\) if \(A\) is PD. The \(\lambda_i\) are the eigenvalues of \(A\) (Theorem~\ref{r8-thm-spectral}(c)).
\end{itemize}
(b)
\begin{itemize}
\item \(-A=\sum_j(-c_j)E_j\), so the eigenvalues of \(-A\) are \(-c_1,\dots,-c_k\)
(Theorem~\ref{r8-thm-spectral-decomp}(d)); apply (a) to \(-A\).
\item If \(c_i>0>c_l\), unit eigenvectors \(x\), \(y\) for \(c_i\), \(c_l\) give \(x'Ax=c_i>0\) and \(y'Ay=c_l<0\). If
no eigenvalue is negative (positive), then \(A\) is PSD (NSD), hence not indefinite.
\end{itemize}
(c) If \(A\) is PD, its eigenvalues \(\lambda_1,\dots,\lambda_n\) are \(>0\) by (a), so
\(\det A=\lambda_1\cdots\lambda_n>0\) (Theorem~\ref{r7-thm-trace-det}), and \(A\) is invertible
(Theorem~\ref{r5d-thm-product}(a)). If \(A\) is PSD, its eigenvalues are \(\geqslant0\) by (a), and
\(\det A=\lambda_1\cdots\lambda_n\geqslant0\) in the same way. If \(A\) is PSD and invertible, then \(0\) is not an
eigenvalue (Theorem~\ref{r7-thm-eig-singular}), so every eigenvalue is \(>0\), and \(A\) is PD by (a).

(d) As \(\lambda_n\leqslant\lambda_i\leqslant\lambda_1\) and \(y_i^2\geqslant0\), \(x'Ax=\sum_i\lambda_iy_i^2\) lies between
\(\lambda_n\sum_iy_i^2=\lambda_n\,x'x\) and \(\lambda_1\sum_iy_i^2=\lambda_1\,x'x\). Moreover \(q_1'Aq_1=\lambda_1q_1'q_1=\lambda_1\)
and \(q_n'Aq_n=\lambda_n\).
\end{proof}

\begin{example}[\citeat{abadir-magnus}{Exercise~8.74(c)}]\label{r8-exm-definite}
\begin{enumerate}
\item \(C_n\) has the eigenvalues \(1\) and \(0\) (Example~\ref{r8-exm-patterned}(a)), so it is PSD but not PD. Indeed
\(C_n=1\cdot C_n+0\cdot\bar J_n\), and the identity of Theorem~\ref{r8-thm-eigen-criterion} gives
\(\sum_i(y_i-\bar y)^2=y'C_ny=\lVert C_ny\rVert^2\geqslant0\) (Example~\ref{r1-exm-centering}), with equality exactly when
\(C_ny=0\), that is, when \(y=\bar y\one\) is constant.
\item \(A=2J_3-I_3\) (Example~\ref{r8-exm-spectral-j3}) has the eigenvalues \(5\) and \(-1\), so it is indefinite: for
\(x=\one_3\) and \(y=(1,-1,0)'\), \(x'Ax=15\) and \(y'Ay=-2\).
\item By Example~\ref{r8-exm-patterned}(c), the equicorrelation matrix \(R\) is PD if and only if \(1+(n-1)\rho>0\) and
\(1-\rho>0\), that is, \(-1/(n-1)<\rho<1\), and PSD if and only if \(-1/(n-1)\leqslant\rho\leqslant1\) (for \(\rho=0\),
\(R=I_n\)). For \(n=4\) and \(\rho=-\frac13\), \(R\) is PSD and singular (\(R\one_4=0\),
Example~\ref{r5d-exm-equicorrelation}).
\item The matrix \(A\) of Example~\ref{r8-exm-spectral} has the eigenvalues \(25\), \(9\), \(9\), \(1\), all \(>0\); so it
is PD (Theorem~\ref{r8-thm-eigen-criterion}(a)).
\end{enumerate}
\end{example}

\needspace{10\baselineskip}
\begin{proposition}[\citeat{abadir-magnus}{Exercise~8.6}, \citeat{harville}{Lemmas~14.2.2, 14.2.4, 14.2.17}, \citeat{lang}{Ch.~VIII, \S4, Exercise~5}, \citeat{shakarchi}{}]\label{r8-prop-pd-basic}
Let \(A,B\in\R^{n\times n}\) be symmetric.
\begin{enumerate}
\item A diagonal matrix \(D=\diag(d_1,\dots,d_n)\) is PD (PSD) if and only if every \(d_i>0\) (\(d_i\geqslant0\)).
\item If \(A\) and \(B\) are PSD and \(\alpha,\beta\geqslant0\), then \(\alpha A+\beta B\) is PSD; if moreover \(A\) is PD and
\(\alpha>0\), then \(\alpha A+\beta B\) is PD.
\item If \(A\) is PD, then \(A\) is invertible and \(A^{-1}\) is PD.
\item If \(A\) is idempotent, then \(A\) is PSD.
\item If \(A\) and \(-A\) are PSD, then \(A=0\).
\end{enumerate}
\end{proposition}
\begin{proof}
(a) \(x'Dx=d_1x_1^2+\dots+d_nx_n^2\). This is \(>0\) for \(x\ne0\) if every \(d_i>0\), and \(\geqslant0\) for every \(x\) if
every \(d_i\geqslant0\). Conversely \(d_i=e_i'De_i\), which is \(>0\) if \(D\) is PD and \(\geqslant0\) if \(D\) is PSD.

(b) \(\alpha A+\beta B\) is symmetric, and \(x'(\alpha A+\beta B)x=\alpha\,x'Ax+\beta\,x'Bx\). Both terms are \(\geqslant0\); if
\(A\) is PD, \(\alpha>0\) and \(x\ne0\), the first term is \(>0\).

(c)
\begin{itemize}
\item \(A\) is invertible (Theorem~\ref{r8-thm-eigen-criterion}(c)), and \((A^{-1})'=(A')^{-1}=A^{-1}\)
(Theorem~\ref{r2-thm-inverse-props}(c)).
\item Let \(x\ne0\) and \(y=A^{-1}x\). Then \(y\ne0\), as \(x=Ay\), and \(x'A^{-1}x=(Ay)'y=y'A'y=y'Ay>0\).
\end{itemize}
(d) As \(A=A'\) and \(A=A^2\), \(x'Ax=x'A'Ax=(Ax)'(Ax)=\lVert Ax\rVert^2\geqslant0\).

(e) \(x'Ax\geqslant0\) and \(-x'Ax\geqslant0\), so \(x'Ax=0=x'0x\) for every \(x\). As \(0\) is symmetric, \(A=0\)
(Proposition~\ref{r8-prop-qf-symmetric}(b)).
\end{proof}

\subsection*{Congruence, \(A'A\), and the factorization \(A=P'P\)}

\begin{theorem}[\citeat{harville}{Theorem~14.2.9, Corollary~14.2.10}]\label{r8-thm-congruence}
Let \(A\in\R^{n\times n}\) be symmetric and \(P\in\R^{n\times m}\).
\begin{enumerate}
\item If \(A\) is PSD, then \(P'AP\) is PSD.
\item If \(A\) is PD and \(\rank P=m\), then \(P'AP\) is PD.
\item If \(m=n\) and \(P\) is invertible, then \(A\) is PD (PSD) if and only if \(P'AP\) is PD (PSD).
\end{enumerate}
\end{theorem}
\begin{proof}
\(P'AP\) is symmetric, and \(x'(P'AP)x=(Px)'A(Px)\) for \(x\in\R^m\).

(a) This is \(\geqslant0\).

(b) If \(x\ne0\), then \(Px\ne0\) (Theorem~\ref{thm:r4-rank-rules}(b)), so it is \(>0\).

(c) \(\rank P=n\) (Theorem~\ref{thm:r4-rank-rules}(d)), so (a) and (b) give one direction; for the other, apply them to
\(P'AP\) and \(P^{-1}\): \(A=(P^{-1})'(P'AP)P^{-1}\) (Theorem~\ref{r2-thm-inverse-props}(c)).
\end{proof}

\begin{theorem}[\citeat{harville}{Corollary~14.2.14}, \citeat{searle-khuri}{Theorem~5.3}]\label{r8-thm-XtX}
Let \(A\in\R^{n\times p}\) have the columns \(a_1,\dots,a_p\). Then \(A'A\), the \textbf{Gram matrix} of \(a_1,\dots,a_p\),
with the entries \(a_i'a_j\) (Theorem~\ref{r1-thm-views}(d)), is PSD, and it is PD if and only if \(\rank A=p\). Likewise
\(AA'\) is PSD, and PD if and only if \(\rank A=n\).
\end{theorem}
\begin{proof}
\(A'A=A'I_nA\), and \(I_n\) is PD; so \(A'A\) is PSD, and PD if \(\rank A=p\) (Theorem~\ref{r8-thm-congruence}(a),(b)).
If \(A'A\) is PD, it is invertible (Theorem~\ref{r8-thm-eigen-criterion}(c)), so
\(\rank A=\rank(A'A)=p\) (Theorems~\ref{thm:r4-rank-ata}(b) and~\ref{thm:r4-rank-rules}(d)). For \(AA'\) apply this to
\(A'\), with \(\rank A'=\rank A\) (Theorem~\ref{thm:r4-rank-rules}(a)).
\end{proof}
For a data matrix \(X_0\in\R^{n\times k}\), the matrix of corrected sums of squares and products is
\(X_0'C_nX_0=(C_nX_0)'(C_nX_0)\) (Example~\ref{r1-exm-data}). So it is PSD, and PD exactly when \(\rank(C_nX_0)=k\)
(Theorem~\ref{r8-thm-XtX}), that is, when the centered columns are linearly independent; in Example~\ref{r1-exm-data},
\(X_0'C_4X_0=4I_2\) is PD.

\needspace{14\baselineskip}
\begin{theorem}[\citeat{harville}{Theorem~14.3.7, Corollaries~14.3.8, 14.3.13}, \citeat{hoffman-kunze}{\S9.3, Theorem~5; \S8.4, Example~24}]\label{r8-thm-psd-factor}
Let \(A\in\R^{n\times n}\).
\begin{enumerate}
\item Let \(1\leqslant r\leqslant n\). Then \(A\) is symmetric, PSD and of rank \(r\) if and only if \(A=P'P\) for some
\(P\in\R^{r\times n}\) of rank \(r\).
\item \(A\) is symmetric and PSD if and only if \(A=P'P\) for some matrix \(P\) with \(n\) columns.
\item \(A\) is symmetric and PD if and only if \(A=P'P\) for some invertible \(P\in\R^{n\times n}\). In this case
\(\det A=(\det P)^2>0\), and \(y'Ax=(Py)'(Px)\) for all \(x,y\in\R^n\).
\end{enumerate}
\end{theorem}
\begin{proof}
(a)
\begin{itemize}
\item Let \(A\) be symmetric, PSD and of rank \(r\), and let \(A=Q\Lambda Q'=\sum_{i=1}^n\lambda_iq_iq_i'\)
(Theorems~\ref{r8-thm-spectral} and~\ref{r8-thm-spectral-decomp}(c)). Exactly \(r\) of the \(\lambda_i\), say
\(\lambda_{k_1},\dots,\lambda_{k_r}\), are nonzero (Theorem~\ref{r8-thm-rank-eigen}), and they are \(>0\)
(Theorem~\ref{r8-thm-eigen-criterion}(a)).
\item Let \(P\in\R^{r\times n}\) have the rows \(\sqrt{\lambda_{k_1}}\,q_{k_1}',\dots,\sqrt{\lambda_{k_r}}\,q_{k_r}'\). By the outer
product expansion (Theorem~\ref{r1-thm-views}(d)), \(P'P=\sum_{l=1}^r\lambda_{k_l}q_{k_l}q_{k_l}'=A\), since the other terms of
\(\sum_i\lambda_iq_iq_i'\) are \(0\). Moreover \(\rank P=\rank(P'P)=\rank A=r\) (Theorem~\ref{thm:r4-rank-ata}(b)).
\item Conversely, if \(A=P'P\) with \(P\in\R^{r\times n}\) of rank \(r\), then \(A\) is symmetric
(Theorem~\ref{r1-thm-transpose}), PSD (Theorem~\ref{r8-thm-XtX}), and \(\rank A=\rank(P'P)=\rank P=r\)
(Theorem~\ref{thm:r4-rank-ata}(b)).
\end{itemize}
(b) Let \(A\) be symmetric and PSD. If \(A=0\), then \(A=P'P\) for the zero matrix \(P\in\R^{1\times n}\); otherwise (a) applies
with \(r=\rank A\geqslant1\). Conversely, \(P'P\) is symmetric and PSD (Theorems~\ref{r1-thm-transpose}
and~\ref{r8-thm-XtX}).

(c)
\begin{itemize}
\item If \(A\) is symmetric and PD, then \(A\) is invertible (Theorem~\ref{r8-thm-eigen-criterion}(c)), so \(\rank A=n\)
(Theorem~\ref{thm:r4-rank-rules}(d)). By (a) with \(r=n\), \(A=P'P\) with \(P\in\R^{n\times n}\) of rank \(n\), and \(P\) is
invertible (Theorem~\ref{thm:r4-rank-rules}(d)).
\item Conversely, if \(A=P'P\) with \(P\) invertible, then \(\rank P=n\), and \(A\) is symmetric and PD
(Theorems~\ref{r1-thm-transpose} and~\ref{r8-thm-XtX}).
\item Then \(\det A=\det P'\cdot\det P=(\det P)^2\) (Theorems~\ref{r5d-thm-product}(b) and~\ref{r5d-thm-transpose}), and
\(\det P\ne0\) (Theorem~\ref{r5d-thm-product}(a)). Finally \(y'Ax=y'P'Px=(Py)'(Px)\).\qedhere
\end{itemize}
\end{proof}

\begin{corollary}[\citeat{harville}{Corollary~14.3.11}]\label{r8-cor-psd-null}
Let \(A\in\R^{n\times n}\) be PSD and \(X\in\R^{n\times m}\). Then \(X'AX=0\) if and only if \(AX=0\). In particular
\(x'Ax=0\) if and only if \(Ax=0\): the vectors \(x\) with \(x'Ax=0\) form the subspace \(\mathcal N(A)\).
\end{corollary}
\begin{proof}
If \(AX=0\), then \(X'AX=0\). Conversely, let \(X'AX=0\), and write \(A=P'P\) (Theorem~\ref{r8-thm-psd-factor}(b)).
\begin{itemize}
\item \(X'AX=(PX)'(PX)\), so \(\tr\bigl((PX)'(PX)\bigr)=0\), and \(PX=0\) (Theorem~\ref{r1-thm-trace}(d)).
\item Hence \(AX=P'(PX)=0\). The case \(m=1\), \(X=x\), gives the last sentence.\qedhere
\end{itemize}
\end{proof}
For instance, \(\sum_i(y_i-\bar y)^2=y'C_ny=0\) only if \(C_ny=0\) (Example~\ref{r1-exm-centering}), and, for \(A\) and
\(P\) as in Theorem~\ref{r6-thm-hat}, \(y'(I_n-P)y=0\) only if \((I_n-P)y=0\): both \(C_n\) and \(I_n-P\) are symmetric
and idempotent (Proposition~\ref{r1-prop-centering}, Theorem~\ref{r6-thm-hat}(b)), hence PSD
(Proposition~\ref{r8-prop-pd-basic}(d)).

\begin{example}\label{r8-exm-psd-factor}
For \(x\in\R^4\), the sum of the squares of the two sums of three successive entries is
\[
(x_1+x_2+x_3)^2+(x_2+x_3+x_4)^2=\lVert Dx\rVert^2=x'Ax
\]
where
\[
D=\begin{pmatrix}1&1&1&0\\0&1&1&1\end{pmatrix},\qquad A=D'D=\begin{pmatrix}1&1&1&0\\1&2&2&1\\1&2&2&1\\0&1&1&1\end{pmatrix}
\]
Find the factor \(P\) of Theorem~\ref{r8-thm-psd-factor}(a) and all \(x\) with \(x'Ax=0\).

\textit{Solution.}
\begin{itemize}
\item \(A\) is symmetric and PSD, and \(\rank A=\rank D=2\) (Theorems~\ref{r8-thm-XtX} and~\ref{thm:r4-rank-ata}(b)), as the
two rows of \(D\) are independent.
\item \(A(1,2,2,1)'=(5,10,10,5)'\), \(A(1,0,0,-1)'=(1,0,0,-1)'\), and \(A(0,1,-1,0)'=A(2,-1,-1,2)'=0\). So
\(a_A(c)\geqslant1\) for \(c=5,1\), and \(a_A(0)=\dim\mathcal N(A)=4-\rank A=2\) (Theorem~\ref{r8-thm-rank-eigen}); as the
multiplicities add up to \(4\) (Definition~\ref{r7-def-eigenvalues}), the eigenvalues are \(5\), \(1\), \(0\), \(0\). The
unit eigenvectors for \(5\) and \(1\) are \(q_1=\frac1{\sqrt{10}}(1,2,2,1)'\) and \(q_2=\frac1{\sqrt2}(1,0,0,-1)'\), and the
orthogonal unit vectors \(q_3=\frac1{\sqrt2}(0,1,-1,0)'\), \(q_4=\frac1{\sqrt{10}}(2,-1,-1,2)'\) of \(\mathcal N(A)\) complete
them to an orthonormal basis of eigenvectors (Theorem~\ref{r8-thm-sym-orth}).
\item The proof of Theorem~\ref{r8-thm-psd-factor}(a) takes the rows \(\sqrt5\,q_1'=\frac1{\sqrt2}(1,2,2,1)\) and
\(\sqrt1\,q_2'=\frac1{\sqrt2}(1,0,0,-1)\):
\[
P=\frac1{\sqrt2}\begin{pmatrix}1&2&2&1\\1&0&0&-1\end{pmatrix},\qquad
P'P=\frac12\begin{pmatrix}2&2&2&0\\2&4&4&2\\2&4&4&2\\0&2&2&2\end{pmatrix}=A
\]
The matrix \(D\) is another \(2\times4\) factor of rank \(2\): the factor in Theorem~\ref{r8-thm-psd-factor}(a) is not
unique.
\item By Corollary~\ref{r8-cor-psd-null}, \(x'Ax=0\) if and only if \(x\in\mathcal N(A)=\Span\{(0,1,-1,0)',(2,-1,-1,2)'\}\)
(\(\dim\mathcal N(A)=4-2=2\), Theorem~\ref{thm:r4-rank-nullity}, and these two vectors are independent,
Theorem~\ref{thm:r4-counting}(d)), that is, \(x_1=x_4=-(x_2+x_3)\).
\end{itemize}
\end{example}

\subsection*{Principal submatrices and Schur complements}

\begin{definition}[\citeat{bapat}{\S\S3.1, 3.3}, \citeat{hoffman-kunze}{\S9.3}, \citeat{harville}{Corollary~14.2.12}]\label{r8-def-principal}
Let \(A\in\R^{n\times n}\) and \(1\leqslant i_1<\dots<i_k\leqslant n\). The \(k\times k\) matrix whose \((s,t)\) entry is
\(a_{i_si_t}\) (rows and columns \(i_1,\dots,i_k\) of \(A\)) is a \textbf{principal submatrix} of \(A\); with
\(S=[e_{i_1}\ \cdots\ e_{i_k}]\in\R^{n\times k}\) it equals \(S'AS\), as the \((s,t)\) entry of \(S'AS\) is
\(e_{i_s}'Ae_{i_t}=a_{i_si_t}\) (Theorem~\ref{r1-thm-views}(b),(d)). The \textbf{leading principal submatrix} \(A_k\)
uses rows and columns \(1,\dots,k\), and the numbers \(\Delta_k(A)=\det A_k\) (\(k=1,\dots,n\)) are the
\textbf{leading principal minors} of \(A\); we put \(\Delta_0(A)=1\).
\end{definition}

\begin{theorem}[\citeat{bapat}{Result~3.4}, \citeat{harville}{Corollaries~14.2.12, 14.2.13, 14.5.14, Theorem~14.9.11}, \citeat{abadir-magnus}{Exercises~8.7(b),(d), 8.27(c)}]\label{r8-thm-principal-sub}
If \(A\) is PD (PSD), then every principal submatrix of \(A\) is PD (PSD). In particular every diagonal entry
\(a_{ii}\) is \(>0\) (\(\geqslant0\)), and if \(A\) is PD, then \(\Delta_k(A)>0\) for \(k=1,\dots,n\). If \(A\) is PSD, then
moreover \(a_{ij}^2\leqslant a_{ii}a_{jj}\) for all \(i,j\), and if \(a_{ii}=0\), then row \(i\) and column \(i\) of \(A\)
are zero.
\end{theorem}
\begin{proof}
\begin{itemize}
\item With \(S\) as in Definition~\ref{r8-def-principal} and \(y\in\R^k\), the vector \(x=Sy\) has the entries
\(y_1,\dots,y_k\) in the positions \(i_1,\dots,i_k\) and zeros elsewhere (Theorem~\ref{r1-thm-views}(b)); so \(x\ne0\) if
\(y\ne0\), and \(y'(S'AS)y=x'Ax\). Hence \(S'AS\), which is symmetric, is PD (PSD) when \(A\) is.
\item \((a_{ii})\) is the principal submatrix with \(k=1\) and \(i_1=i\). If \(A\) is PD, then \(A_k\) is PD and
\(\Delta_k(A)=\det A_k>0\) (Theorem~\ref{r8-thm-eigen-criterion}(c)).
\item Let \(A\) be PSD and \(i\ne j\). The principal submatrix with the rows and columns \(i\) and \(j\) is PSD (first
item), and its determinant is \(a_{ii}a_{jj}-a_{ij}a_{ji}=a_{ii}a_{jj}-a_{ij}^2\) (Theorem~\ref{r5d-thm-expansion});
so \(a_{ii}a_{jj}-a_{ij}^2\geqslant0\) (Theorem~\ref{r8-thm-eigen-criterion}(c)). For \(i=j\) there is nothing to prove.
\item If \(A\) is PSD and \(a_{ii}=0\), then \(e_i'Ae_i=a_{ii}=0\) (Theorem~\ref{r1-thm-views}(b)), so \(Ae_i=0\)
(Corollary~\ref{r8-cor-psd-null}): column \(i\) of \(A\) is zero, and so is row \(i\), its transpose, as \(A'=A\).\qedhere
\end{itemize}
\end{proof}
So \(\left(\begin{smallmatrix}1&2\\2&3\end{smallmatrix}\right)\) is not PSD (\(2^2>1\cdot3\)), and neither is
\(\left(\begin{smallmatrix}0&1\\1&5\end{smallmatrix}\right)\) (\(a_{11}=0\), but row \(1\) is not zero).

\needspace{9\baselineskip}
\begin{theorem}[\citeat{bapat}{Result~3.14}, \citeat{harville}{Lemma~14.8.3, Theorems~14.8.4, 14.8.5}, \citeat{abadir-magnus}{Exercise~12.33}]\label{r8-thm-schur-pd}
Let \(M=\begin{pmatrix}A&B\\B'&D\end{pmatrix}\) be symmetric, with \(A\in\R^{p\times p}\) invertible,
\(B\in\R^{p\times q}\) and \(D\in\R^{q\times q}\), and let \(M/A=D-B'A^{-1}B\) be the Schur complement of \(A\). Then
\(M\) is PD if and only if \(A\) and \(M/A\) are PD. If \(A\) is PD, then \(M\) is PSD if and only if \(M/A\) is PSD. In
every case \(\det M=\det A\cdot\det(M/A)\).
\end{theorem}
\begin{proof}
\begin{itemize}
\item \(A\) and \(D\) are symmetric, as \(M'=M\) (Theorem~\ref{r5b-thm-block}(b)), hence so are \(A^{-1}\)
(Theorem~\ref{r2-thm-inverse-props}(c)) and \(M/A\). Let \(L=\left(\begin{smallmatrix}I_p&0\\-B'A^{-1}&I_q\end{smallmatrix}\right)\);
then \(L'=\left(\begin{smallmatrix}I_p&-A^{-1}B\\0&I_q\end{smallmatrix}\right)\), \(L\) is invertible
(Proposition~\ref{r5b-prop-block-elementary}), and by Theorem~\ref{r5b-thm-schur} (with \(C=B'\))
\[
LML'=\begin{pmatrix}A&0\\0&M/A\end{pmatrix}
\]
\item If \(M\) is PD (PSD), so is \(LML'\) (Theorem~\ref{r8-thm-congruence}(c)), and \(A\) and \(M/A\) are principal
submatrices of it, hence PD (PSD) (Theorem~\ref{r8-thm-principal-sub}).
\item Let \(A\) be PD and \(M/A\) PSD, and let \(x=\left(\begin{smallmatrix}y\\z\end{smallmatrix}\right)\) (\(y\in\R^p\),
\(z\in\R^q\)). Then \(x'(LML')x=y'Ay+z'(M/A)z\geqslant0\), and if \(M/A\) is PD and \(x\ne0\), this is \(>0\), as \(y\ne0\)
or \(z\ne0\). So \(LML'\) is PSD, and PD if \(M/A\) is PD; and so is \(M\) (Theorem~\ref{r8-thm-congruence}(c)).
\item The determinant is Theorem~\ref{r5d-thm-schur-det}(a).\qedhere
\end{itemize}
\end{proof}

\begin{theorem}[Hadamard's inequality; \citeat{abadir-magnus}{Exercise~12.26(a)}, \citeat{bapat}{Result~7.2}]\label{r8-thm-hadamard}
If \(A=(a_{ij})\in\R^{n\times n}\) is PSD, then \(\det A\leqslant a_{11}a_{22}\cdots a_{nn}\). If \(A\) is PD, equality holds
if and only if \(A\) is diagonal.
\end{theorem}
\begin{proof}
First let \(A\) be PD, with the leading principal submatrices \(A_1,\dots,A_n=A\) (Definition~\ref{r8-def-principal}).
\begin{itemize}
\item For \(2\leqslant k\leqslant n\) write
\[
A_k=\begin{pmatrix}A_{k-1}&b\\b'&a_{kk}\end{pmatrix}
\]
with \(b\in\R^{k-1}\).
\(A_{k-1}\) is PD (Theorem~\ref{r8-thm-principal-sub}), so \(\det A_{k-1}>0\) and \(A_{k-1}^{-1}\) is PD
(Theorem~\ref{r8-thm-eigen-criterion}(c), Proposition~\ref{r8-prop-pd-basic}(c)). By Theorem~\ref{r5d-thm-schur-det}(a),
\[
\det A_k=\det A_{k-1}\cdot\bigl(a_{kk}-b'A_{k-1}^{-1}b\bigr)\leqslant a_{kk}\det A_{k-1}
\]
with equality if and only if \(b'A_{k-1}^{-1}b=0\), that is, \(b=0\).
\item Every \(a_{kk}>0\) (Theorem~\ref{r8-thm-principal-sub}). Applying the first item for \(k=n,n-1,\dots,2\),
\(\det A\leqslant a_{nn}\det A_{n-1}\leqslant a_{nn}a_{n-1,n-1}\det A_{n-2}\leqslant\dots\leqslant a_{nn}\cdots a_{22}\det A_1
=a_{11}a_{22}\cdots a_{nn}\).
\item Each step multiplies one of these inequalities by positive numbers, so equality holds if and only if it holds at
every step, that is, if and only if the off-diagonal entries of the last row and column of \(A_k\) vanish for
\(k=2,\dots,n\): \(A\) is diagonal.
\item If \(A\) is PSD but not PD, then \(A\) is not invertible (Theorem~\ref{r8-thm-eigen-criterion}(c)), so \(\det A=0\)
(Theorem~\ref{r5d-thm-product}(a)), while every \(a_{ii}\geqslant0\) (Theorem~\ref{r8-thm-principal-sub}).\qedhere
\end{itemize}
\end{proof}
For the equicorrelation matrix \(R\), whose diagonal entries are \(1\) and which is PSD for \(-1/(n-1)\leqslant\rho\leqslant1\)
(Example~\ref{r8-exm-definite}(c)), Hadamard's inequality says that
\(\det R=\bigl(1+(n-1)\rho\bigr)(1-\rho)^{n-1}\leqslant1\) (Example~\ref{r5d-exm-equicorrelation}), with equality for PD \(R\) only
at \(\rho=0\).

\subsection*{Leading principal minors, LDU and Cholesky decompositions}

\begin{lemma}[\citeat{hoffman-kunze}{\S9.3, Lemma}]\label{r8-lem-triangular}
Let \(A\in\R^{n\times n}\) with \(\Delta_k(A)\ne0\) for \(k=1,\dots,n-1\). There is exactly one upper triangular
\(P\in\R^{n\times n}\) with diagonal entries \(1\) such that \(AP\) is lower triangular, and \(\Delta_k(AP)=\Delta_k(A)\)
for \(k=1,\dots,n\).
\end{lemma}
\begin{proof}
Let \(P\) be upper triangular with \(p_{kk}=1\) (\(k=1,\dots,n\)), and \(B=AP\).
\begin{itemize}
\item \(b_{jk}=\sum_ra_{jr}p_{rk}=a_{jk}+\sum_{r=1}^{k-1}a_{jr}p_{rk}\), as \(p_{rk}=0\) for \(r>k\). So \(B\) is lower
triangular (\(b_{jk}=0\) for \(j<k\)) if and only if, for \(k=2,\dots,n\),
\[
\sum_{r=1}^{k-1}a_{jr}p_{rk}=-a_{jk}\qquad(j=1,\dots,k-1)
\]
\item For fixed \(k\), this is a system of \(k-1\) equations in the \(k-1\) unknowns \(p_{1k},\dots,p_{k-1,k}\) with
coefficient matrix \(A_{k-1}\), and \(\det A_{k-1}=\Delta_{k-1}(A)\ne0\). So it has exactly one solution
(Theorems~\ref{r5d-thm-product}(a) and~\ref{r2-thm-invertible}): \(P\) exists and is unique.
\item By the first item, column \(r\) of \(B\) is column \(r\) of \(A\) plus \(\sum_{j<r}p_{jr}\) times column \(j\) of
\(A\). So \(B_k\) arises from \(A_k\) by adding to column \(r\) a combination of columns \(j<r\), for \(r=k,k-1,\dots,2\)
in turn (each step uses columns not yet changed). Such column operations do not change the determinant
(Definition~\ref{r5d-def-det}(b) and Theorem~\ref{r5d-thm-transpose}); so \(\Delta_k(B)=\Delta_k(A)\).\qedhere
\end{itemize}
\end{proof}

\needspace{12\baselineskip}
\begin{theorem}[Sylvester's criterion; \citeat{hoffman-kunze}{\S9.3, Theorem~6}]\label{r8-thm-leading-minors}
A symmetric \(A\in\R^{n\times n}\) is PD if and only if \(\Delta_k(A)>0\) for \(k=1,\dots,n\). In this case, with \(P\) as
in Lemma~\ref{r8-lem-triangular} and \(U=P^{-1}\),
\[
P'AP=D=\diag(d_1,\dots,d_n),\qquad d_k=\frac{\Delta_k(A)}{\Delta_{k-1}(A)}>0,\qquad A=U'DU
\]
and \(U\) is upper triangular with diagonal entries \(1\).
\end{theorem}
\begin{proof}
If \(A\) is PD, every \(\Delta_k(A)>0\) by Theorem~\ref{r8-thm-principal-sub}. Conversely, let every \(\Delta_k(A)>0\),
and let \(B=AP\) be lower triangular (Lemma~\ref{r8-lem-triangular}).
\begin{itemize}
\item \(D=P'AP=P'B\) is lower triangular: \((P'B)_{jk}=\sum_rp_{rj}b_{rk}\), with \(p_{rj}=0\) for \(r>j\) and
\(b_{rk}=0\) for \(r<k\), so it is \(0\) for \(j<k\). Also \(D'=P'A'P=D\). A symmetric lower triangular matrix is
diagonal.
\item As \(p_{jj}=1\), row \(j\) of \(P'B\) is row \(j\) of \(B\) plus \(\sum_{r<j}p_{rj}\) times row \(r\) of \(B\). So
\(D_k\) arises from \(B_k\) by adding to row \(j\) a combination of the rows above it, for \(j=k,\dots,2\) in turn; this
does not change the determinant (Definition~\ref{r5d-def-det}(b)), and \(\Delta_k(D)=\Delta_k(B)=\Delta_k(A)\).
\item \(\Delta_k(D)=d_1\cdots d_k\) (Theorem~\ref{r5d-thm-basic}(c)). So \(d_1\cdots d_k=\Delta_k(A)>0\) for every
\(k\), and \(d_k=\Delta_k(A)/\Delta_{k-1}(A)>0\).
\item \(\det P=1\) (Theorem~\ref{r5d-thm-basic}(c)), so \(P\) is invertible (Theorem~\ref{r5d-thm-product}(a)), and
\(A=(P')^{-1}DP^{-1}=U'DU\) (Theorem~\ref{r2-thm-inverse-props}(c)).
For \(x\ne0\), \(Ux\ne0\) and \(x'Ax=(Ux)'D(Ux)=\sum_kd_k(Ux)_k^2>0\): \(A\) is PD.
\item \(P\) is unit upper triangular, so \(U=P^{-1}\) is unit upper triangular
(Theorem~\ref{r2-thm-triangular}).\qedhere
\end{itemize}
\end{proof}

\begin{example}\label{r8-exm-leading}
Show that the matrix
\[
A=\begin{pmatrix}2&0&1&-1\\0&2&1&1\\1&1&3&0\\-1&1&0&3\end{pmatrix}
\]
of Example~\ref{r8-exm-qf} is PD, and write \(A=U'DU\).

\textit{Solution.}
\begin{itemize}
\item \(\Delta_1=2\), \(\Delta_2=\det(2I_2)=4\), and along row \(1\) (Theorem~\ref{r5d-thm-expansion}),
\(\Delta_3=2(2\cdot3-1\cdot1)-0+1(0\cdot1-2\cdot1)=10-2=8\). For \(\Delta_4=\det A\), write
\(A=\left(\begin{smallmatrix}A_2&B\\B'&3I_2\end{smallmatrix}\right)\) with \(A_2=2I_2\) and
\(B=\left(\begin{smallmatrix}1&-1\\1&1\end{smallmatrix}\right)\). As \(B'B=2I_2\), the Schur complement of \(A_2\) is
\(A/A_2=3I_2-B'(\frac12I_2)B=3I_2-I_2=2I_2\), and \(\Delta_4=\det A_2\cdot\det(2I_2)=4\cdot4=16\)
(Theorems~\ref{r5d-thm-schur-det}(a) and~\ref{r5d-thm-basic}(c)). All are \(>0\), so \(A\) is PD
(Theorem~\ref{r8-thm-leading-minors}).
\item \(P\): for \(k=2\), \(a_{11}p_{12}=-a_{12}\) gives \(p_{12}=0\); for \(k=3\),
\(A_2\left(\begin{smallmatrix}p_{13}\\p_{23}\end{smallmatrix}\right)=-\left(\begin{smallmatrix}1\\1\end{smallmatrix}\right)\)
with \(A_2=2I_2\) gives \(p_{13}=p_{23}=-\frac12\); for \(k=4\), \(A_3(p_{14},p_{24},p_{34})'=-(a_{14},a_{24},a_{34})'=(1,-1,0)'\)
reads \(2p_{14}+p_{34}=1\), \(2p_{24}+p_{34}=-1\) and \(p_{14}+p_{24}+3p_{34}=0\). The sum of the first two equations gives
\(p_{14}+p_{24}=-p_{34}\), so the third gives \(2p_{34}=0\); hence \(p_{34}=0\), \(p_{14}=\frac12\) and \(p_{24}=-\frac12\).
Then
\[
P=\begin{pmatrix}1&0&-\frac12&\frac12\\0&1&-\frac12&-\frac12\\0&0&1&0\\0&0&0&1\end{pmatrix},\qquad
AP=\begin{pmatrix}2&0&0&0\\0&2&0&0\\1&1&2&0\\-1&1&0&2\end{pmatrix},\qquad
U=P^{-1}=\begin{pmatrix}1&0&\frac12&-\frac12\\0&1&\frac12&\frac12\\0&0&1&0\\0&0&0&1\end{pmatrix}
\]
and \(D=P'(AP)=\diag(2,2,2,2)\), with \(d_k=\Delta_k/\Delta_{k-1}=2/1,\ 4/2,\ 8/4,\ 16/8\). So \(A=U'DU\), that is,
\(q_A(x)=2(x_1+\frac12x_3-\frac12x_4)^2+2(x_2+\frac12x_3+\frac12x_4)^2+2x_3^2+2x_4^2\).
\item Schur complement (Theorem~\ref{r8-thm-schur-pd}, \(p=2\)): \(A_2=2I_2\) and \(A/A_2=2I_2\) (first item) are PD
(Proposition~\ref{r8-prop-pd-basic}(a)); so again \(A\) is PD, and \(\det A=\det A_2\cdot\det(A/A_2)=16\).
\item Eigenvalues: \(A(1,1,2,0)'=(4,4,8,0)'\), \(A(-1,1,0,2)'=(-4,4,0,8)'\), \(A(1,1,-1,0)'=(1,1,-1,0)'\) and
\(A(1,-1,0,1)'=(1,-1,0,1)'\). The two vectors for \(4\) are nonzero and orthogonal, hence independent
(Proposition~\ref{r6-prop-orthogonal-list}(b)), and so are the two for \(1\). So \(a_A(c)=\dim\mathcal N(A-cI)\geqslant2\)
for \(c=4,1\) (Theorems~\ref{r8-thm-rank-eigen} and~\ref{thm:r4-counting}(c)); as the multiplicities add up to \(4\)
(Definition~\ref{r7-def-eigenvalues}), the eigenvalues are \(4\), \(4\), \(1\), \(1\), all \(>0\), as
Theorem~\ref{r8-thm-eigen-criterion}(a) requires.
\end{itemize}
\end{example}

\begin{example}\label{r8-exm-sylvester}
For which \(a\in\R\) is
\[
A=\begin{pmatrix}1&a&a&0\\a&1&0&a\\a&0&1&a\\0&a&a&1\end{pmatrix}
\]
PD?

\textit{Solution.}
\begin{itemize}
\item The first two leading principal minors are \(1\) and \(1-a^2\). Expanding along the first row
(Theorem~\ref{r5d-thm-expansion}), \(\Delta_3=(1-0)-a(a-0)+a(0-a)=1-2a^2\).
\item For \(\det A\), subtract row \(3\) from row \(2\), which becomes \((0,1,-1,0)\), and then add column \(2\) to column
\(3\) (Definition~\ref{r5d-def-det}(b), Theorem~\ref{r5d-thm-transpose}); then expand along row \(2\) and along the first
row (Theorem~\ref{r5d-thm-expansion}):
\[
\det A=\det\begin{pmatrix}1&a&2a&0\\0&1&0&0\\a&0&1&a\\0&a&2a&1\end{pmatrix}
=\det\begin{pmatrix}1&2a&0\\a&1&a\\0&2a&1\end{pmatrix}=(1-2a^2)-2a(a-0)=1-4a^2
\]
\item By Theorem~\ref{r8-thm-leading-minors}, \(A\) is PD if and only if \(1-a^2\), \(1-2a^2\) and \(1-4a^2\) are
positive, that is, \(a^2<\frac14\), or \(-\frac12<a<\frac12\).
\end{itemize}
Check with eigenvalues: the orthogonal vectors \(h_1=(1,1,1,1)'\), \(h_2=(1,1,-1,-1)'\), \(h_3=(1,-1,1,-1)'\),
\(h_4=(1,-1,-1,1)'\) satisfy
\[
Ah_1=(1+2a)h_1,\qquad Ah_2=h_2,\qquad Ah_3=h_3,\qquad Ah_4=(1-2a)h_4
\]
(for instance, row \(1\) of \(A\) times \(h_4\) is \(1-a-a+0\)). So \(A=Q\Lambda Q'\) with \(Q=\frac12[h_1\ h_2\ h_3\ h_4]\)
orthogonal and \(\Lambda=\diag(1+2a,1,1,1-2a)\) (Theorem~\ref{r8-thm-spectral}(b)), and by
Theorem~\ref{r8-thm-eigen-criterion}(a) \(A\) is PD if and only if \(1-2\lvert a\rvert>0\), the same condition. For
\(a=\pm\frac12\) the eigenvalues are \(2\), \(1\), \(1\), \(0\): \(A\) is PSD but not PD; for \(\lvert a\rvert>\frac12\),
\(A\) is indefinite (Theorem~\ref{r8-thm-eigen-criterion}(b)).
\end{example}

\needspace{12\baselineskip}
\begin{corollary}[\citeat{harville}{Corollary~14.5.7, Lemma~14.5.8}, \citeat{hoffman-kunze}{\S9.3, (9-6)}]\label{r8-cor-ldu}
Let \(A\in\R^{n\times n}\) with \(\Delta_k(A)\ne0\) for \(k=1,\dots,n-1\), and let \(P\) be as in Lemma~\ref{r8-lem-triangular}.
\begin{enumerate}
\item \(A=LDU\), where \(L\) is unit lower triangular, \(D=\diag(d_1,\dots,d_n)\) with \(d_k=\Delta_k(A)/\Delta_{k-1}(A)\), and
\(U=P^{-1}\) is unit upper triangular: an \textbf{LDU decomposition} of \(A\).
\item It is the only one: if \(A=L_1D_1U_1\) with \(L_1\) unit lower triangular, \(D_1\) diagonal and \(U_1\) unit upper
triangular, then \(L_1=L\), \(D_1=D\) and \(U_1=U\).
\item If \(A\) is symmetric, then \(L=U'\); so \(A=U'DU\) and \(P'AP=D\).
\end{enumerate}
\end{corollary}
\begin{proof}
Let \(B=AP\); it is lower triangular, and \(\Delta_k(B)=\Delta_k(A)\) for every \(k\) (Lemma~\ref{r8-lem-triangular}).

(a)
\begin{itemize}
\item As \(B\) is lower triangular, \(\Delta_k(B)=b_{11}\cdots b_{kk}\) (Theorems~\ref{r5d-thm-basic}(c)
and~\ref{r5d-thm-transpose}). So \(b_{11}\cdots b_{kk}=\Delta_k(A)\) for every \(k\), and \(b_{kk}=\Delta_k(A)/\Delta_{k-1}(A)=d_k\),
which is \(\ne0\) for \(k\leqslant n-1\).
\item Let \(L\) have the columns \(\frac1{d_k}\col_k(B)\) (\(k\leqslant n-1\)) and the last column \(e_n\). Column \(k\) of \(B\) has
zeros above the diagonal and \(d_k\) on it, so \(L\) is unit lower triangular. For every \(k\),
\(\col_k(LD)=L(d_ke_k)=d_k\col_k(L)=\col_k(B)\) (Theorem~\ref{r1-thm-views}(a),(b)); for \(k=n\), both sides are \(d_ne_n\), the
last column of \(B\). So \(B=LD\).
\item \(P\) is unit upper triangular, so \(U=P^{-1}\) is unit upper triangular (Theorem~\ref{r2-thm-triangular}), and \(A=BP^{-1}=LDU\).
\end{itemize}
(b)
\begin{itemize}
\item \(U_1\) is invertible and \(U_1^{-1}\) is unit upper triangular (Theorem~\ref{r2-thm-triangular}), and \(AU_1^{-1}=L_1D_1\)
is lower triangular, as \(\col_k(L_1D_1)=(D_1)_{kk}\col_k(L_1)\) (Theorem~\ref{r1-thm-views}(a),(b)). By the uniqueness in
Lemma~\ref{r8-lem-triangular}, \(U_1^{-1}=P\), that is, \(U_1=U\).
\item Then \(L_1D_1=AP=B=LD\). The diagonal entries of \(L_1D_1\) and of \(LD\) are those of \(D_1\) and of \(D\), so \(D_1=D\).
\item For \(k\leqslant n-1\), \(\col_k(L_1)=\frac1{d_k}\col_k(L_1D)=\frac1{d_k}\col_k(B)=\col_k(L)\), and the last columns of \(L_1\)
and \(L\) are both \(e_n\). So \(L_1=L\).
\end{itemize}
(c) If \(A'=A\), then \(A=A'=U'D'L'=U'DL'\), where \(U'\) is unit lower and \(L'\) unit upper triangular. By (b), \(L=U'\). So
\(A=U'DU\), and \(P'AP=(UP)'D(UP)=D\), as \(UP=I_n\).
\end{proof}
For a symmetric \(A\) with \(\Delta_k(A)\ne0\) (\(k\leqslant n-1\)), (c) writes the quadratic form as a combination of squares,
\(q_A(x)=(Ux)'D(Ux)=d_1(Ux)_1^2+\dots+d_n(Ux)_n^2\); this is completing the squares, as in Example~\ref{r8-exm-leading}.

\begin{proposition}[\citeat{harville}{\S14.5b}]\label{r8-prop-lu}
Let \(A\in\R^{n\times n}\). A factorization \(A=L_0U_0\) with \(L_0\) unit lower triangular and \(U_0\) upper triangular is an
\textbf{LU decomposition} of \(A\).
\begin{enumerate}
\item If \(\Delta_k(A)\ne0\) for \(k=1,\dots,n-1\), and \(A=LDU\) is the LDU decomposition of Corollary~\ref{r8-cor-ldu}, then
\(A=L(DU)\) is an LU decomposition, and the diagonal entries of \(DU\) are \(d_1,\dots,d_n\).
\item If row operations \(R_i\leftarrow R_i+cR_j\) with \(i>j\), and no others, take \(A\) to an upper triangular matrix
\(U_0\), then \(A=L_0U_0\) for a unit lower triangular \(L_0\).
\end{enumerate}
\end{proposition}
\begin{proof}
(a) \(DU\) is a product of two upper triangular matrices, so it is upper triangular, with diagonal entries \(d_k\cdot1=d_k\)
(Proposition~\ref{r1-prop-diag}(c)).

(b)
\begin{itemize}
\item By Theorem~\ref{r2-thm-elementary}(c), \(U_0=E_N\cdots E_1A\), where each \(E_k=T_{ij}(c)\) with \(i>j\) is unit lower
triangular (Definition~\ref{r2-def-elementary}).
\item Each \(E_k\) is invertible, and \(E_k^{-1}=T_{ij}(-c)\) is unit lower triangular (Theorem~\ref{r2-thm-elementary}(b)); so
\(A=(E_N\cdots E_1)^{-1}U_0=E_1^{-1}\cdots E_N^{-1}U_0\) (Theorem~\ref{r2-thm-inverse-props}(b)).
\item \(L_0=E_1^{-1}\cdots E_N^{-1}\) is a product of unit lower triangular matrices, hence unit lower triangular
(Proposition~\ref{r1-prop-diag}(c)).\qedhere
\end{itemize}
\end{proof}

\begin{example}\label{r8-exm-lu}
Find an LU decomposition of
\[
A=\begin{pmatrix}2&1&1&1\\4&3&3&3\\8&7&9&8\\4&3&5&6\end{pmatrix}
\]
by Gauss elimination.

\textit{Solution.}
\begin{itemize}
\item Elimination without interchanges:
\[
\begin{gathered}
\begin{pmatrix}2&1&1&1\\4&3&3&3\\8&7&9&8\\4&3&5&6\end{pmatrix}
\xrightarrow[R_4\leftarrow R_4-2R_1]{R_2\leftarrow R_2-2R_1,\ R_3\leftarrow R_3-4R_1}
\begin{pmatrix}2&1&1&1\\0&1&1&1\\0&3&5&4\\0&1&3&4\end{pmatrix}
\\
\xrightarrow[R_4\leftarrow R_4-R_2]{R_3\leftarrow R_3-3R_2}
\begin{pmatrix}2&1&1&1\\0&1&1&1\\0&0&2&1\\0&0&2&3\end{pmatrix}
\xrightarrow{R_4\leftarrow R_4-R_3}
\begin{pmatrix}2&1&1&1\\0&1&1&1\\0&0&2&1\\0&0&0&2\end{pmatrix}=U_0
\end{gathered}
\]
\item The elimination matrices are \(T_{21}(-2)\), \(T_{31}(-4)\), \(T_{41}(-2)\), \(T_{32}(-3)\), \(T_{42}(-1)\),
\(T_{43}(-1)\), in this order. By the proof of Proposition~\ref{r8-prop-lu}(b), \(A=L_0U_0\) with
\(L_0=T_{21}(2)\,T_{31}(4)\,T_{41}(2)\,T_{32}(3)\,T_{42}(1)\,T_{43}(1)\) (Theorem~\ref{r2-thm-elementary}(b)). Multiplying
from the left, \(T_{21}(2)\,T_{31}(4)\,T_{41}(2)\) is \(T_{41}(2)\) with \(4\) times row \(1\) added to row \(3\) and then \(2\)
times row \(1\) added to row \(2\), and \(T_{32}(3)\,T_{42}(1)\,T_{43}(1)\) is \(T_{43}(1)\) with row \(2\) added to row \(4\)
and then \(3\) times row \(2\) added to row \(3\) (Theorem~\ref{r2-thm-elementary}(a)); so
\[
L_0=\begin{pmatrix}1&0&0&0\\2&1&0&0\\4&0&1&0\\2&0&0&1\end{pmatrix}\begin{pmatrix}1&0&0&0\\0&1&0&0\\0&3&1&0\\0&1&1&1\end{pmatrix}
=\begin{pmatrix}1&0&0&0\\2&1&0&0\\4&3&1&0\\2&1&1&1\end{pmatrix}
\]
Below the diagonal of \(L_0\) stand the multipliers \(2\), \(4\), \(2\), \(3\), \(1\), \(1\) of the elimination.
\item Check by rows (Theorem~\ref{r1-thm-views}(c)): \(\row_1(L_0U_0)=(2,1,1,1)\),
\(\row_2(L_0U_0)=2\,(2,1,1,1)+(0,1,1,1)=(4,3,3,3)\), \(\row_3(L_0U_0)=4\,(2,1,1,1)+3\,(0,1,1,1)+(0,0,2,1)=(8,7,9,8)\) and
\(\row_4(L_0U_0)=2\,(2,1,1,1)+(0,1,1,1)+(0,0,2,1)+(0,0,0,2)=(4,3,5,6)\).
\item Let \(D_0=\diag(2,1,2,2)\), the diagonal of \(U_0\). Dividing the rows of \(U_0\) by \(2\), \(1\), \(2\), \(2\)
(Proposition~\ref{r1-prop-diag}(a)) gives the unit upper triangular matrix
\[
V_0=D_0^{-1}U_0=\begin{pmatrix}1&\frac12&\frac12&\frac12\\0&1&1&1\\0&0&1&\frac12\\0&0&0&1\end{pmatrix}
\]
and \(A=L_0D_0V_0\). As \(\Delta_1=2\), \(\Delta_2=2\cdot3-1\cdot4=2\) and, along row \(1\)
(Theorem~\ref{r5d-thm-expansion}), \(\Delta_3=2(3\cdot9-3\cdot7)-1(4\cdot9-3\cdot8)+1(4\cdot7-3\cdot8)=12-12+4=4\) are
nonzero, this is the LDU decomposition of Corollary~\ref{r8-cor-ldu}, by its part (b): \(d_1=2\), \(d_2=1\), \(d_3=2\),
\(d_4=2\). Indeed \(d_2=\Delta_2/\Delta_1=1\), \(d_3=\Delta_3/\Delta_2=\frac42=2\) and \(d_4=\Delta_4/\Delta_3=\frac84=2\),
with \(\Delta_4=\det A=\det L_0\det U_0=1\cdot(2\cdot1\cdot2\cdot2)=8\) (Theorem~\ref{r5d-thm-product}(b); \(\det U_0\) by
Theorem~\ref{r5d-thm-basic}(c), and \(\det L_0=\det L_0'=1\) by Theorems~\ref{r5d-thm-transpose}
and~\ref{r5d-thm-basic}(c)).
\end{itemize}
\end{example}

\begin{theorem}[Cholesky decomposition; \citeat{harville}{Theorem~14.5.11}]\label{r8-thm-cholesky}
Let \(A\in\R^{n\times n}\) be PD, and let \(A=U'DU\) as in Theorem~\ref{r8-thm-leading-minors}. There is exactly one upper
triangular \(T\in\R^{n\times n}\) with positive diagonal entries such that \(A=T'T\), namely \(T=D^{1/2}U\), where
\(D^{1/2}=\diag(\sqrt{d_1},\dots,\sqrt{d_n})\). This \(T\) is the \textbf{Cholesky factor} of \(A\), and \(A=T'T\) is the
\textbf{Cholesky decomposition} of \(A\).
\end{theorem}
\begin{proof}
\begin{itemize}
\item Every \(d_k>0\) (Theorem~\ref{r8-thm-leading-minors}), so \(D^{1/2}\) is defined. \(D^{1/2}U\) is \(U\) with row \(k\)
multiplied by \(\sqrt{d_k}\): it is upper triangular with the diagonal entries \(\sqrt{d_k}>0\), and
\((D^{1/2}U)'(D^{1/2}U)=U'D^{1/2}D^{1/2}U=U'DU=A\).
\item Let \(T\) be upper triangular with \(t_{kk}>0\) and \(A=T'T\). Put \(D_*=\diag(t_{11}^2,\dots,t_{nn}^2)\) and
\(U_*=\diag(t_{11}^{-1},\dots,t_{nn}^{-1})\,T\), which is \(T\) with row \(k\) divided by \(t_{kk}\), a unit upper triangular matrix.
Then \(T=\diag(t_{11},\dots,t_{nn})\,U_*\) and \(A=U_*'D_*U_*\).
\item This is an LDU decomposition of \(A\), and every \(\Delta_k(A)>0\) (Theorem~\ref{r8-thm-leading-minors}). By
Corollary~\ref{r8-cor-ldu}(b),(c), \(U_*=U\) and \(D_*=D\). So \(t_{kk}=\sqrt{d_k}\), and \(T=D^{1/2}U\).\qedhere
\end{itemize}
\end{proof}

\begin{remark}[\citeat{harville}{(14.5.16)--(14.5.19)}]\label{r8-rem-cholesky}
For \(i\leqslant j\), the \((i,j)\) entry of \(T'T\) is \(\sum_{k\leqslant i}t_{ki}t_{kj}\), since \(t_{ki}=0\) for
\(k>i\). Comparing with \(a_{ij}\) gives \(T\) row by row (\(i=1,\dots,n\)):
\[
t_{ii}=\Bigl(a_{ii}-\sum_{k<i}t_{ki}^2\Bigr)^{1/2},\qquad
t_{ij}=\frac1{t_{ii}}\Bigl(a_{ij}-\sum_{k<i}t_{ki}t_{kj}\Bigr)\quad(j>i)
\]
For a symmetric \(A\), this succeeds with positive numbers under all square roots exactly when \(A\) is PD: if it
succeeds, it gives an upper triangular \(T\) with positive diagonal entries and \(T'T=A\), so \(A\) is PD
(Theorems~\ref{r2-thm-triangular}(a) and~\ref{r8-thm-psd-factor}(c)); if \(A\) is PD, its Cholesky factor
(Theorem~\ref{r8-thm-cholesky}) satisfies these equations. Moreover \(A_k=T_k'T_k\) for the leading principal
submatrices, as the first \(k\) columns of \(T\) are zero below row \(k\); so
\(\Delta_k(A)=\det T_k'\det T_k=(t_{11}\cdots t_{kk})^2\) (Theorems~\ref{r5d-thm-product}(b), \ref{r5d-thm-transpose}
and~\ref{r5d-thm-basic}(c)).
\end{remark}

\begin{example}\label{r8-exm-cholesky}
Find the Cholesky factor \(T\) of
\[
A=\begin{pmatrix}1&1&1&1\\1&2&2&2\\1&2&3&3\\1&2&3&4\end{pmatrix}
\]
the matrix with entries \(a_{ij}=\min(i,j)\), and use it to compute \(A^{-1}\).

\textit{Solution.} By Remark~\ref{r8-rem-cholesky}:
\begin{itemize}
\item row \(1\): \(t_{11}=\sqrt1=1\) and \(t_{1j}=a_{1j}/1=1\) for \(j=2,3,4\);
\item row \(2\): \(t_{22}=\sqrt{2-1^2}=1\) and \(t_{2j}=(2-1\cdot1)/1=1\) for \(j=3,4\);
\item row \(3\): \(t_{33}=\sqrt{3-1-1}=1\) and \(t_{34}=(3-1-1)/1=1\);
\item row \(4\): \(t_{44}=\sqrt{4-1-1-1}=1\).
\end{itemize}
So \(T\) is the upper triangular matrix with all entries \(t_{ij}=1\) for \(i\leqslant j\); indeed the \((i,j)\) entry of
\(T'T\) is \(\col_i(T)'\col_j(T)=\min(i,j)\). Hence \(A\) is PD, and all its leading principal minors equal \(1\)
(Remark~\ref{r8-rem-cholesky}); in particular \(\det A=\Delta_4(A)=1\), and in Theorem~\ref{r8-thm-leading-minors}
\(D=I_4\) and \(U=T\) (Theorem~\ref{r8-thm-cholesky}: \(T=D^{1/2}U\)). The matrix
\[
K=\begin{pmatrix}1&0&0&0\\-1&1&0&0\\0&-1&1&0\\0&0&-1&1\end{pmatrix}
\]
has \(KT'=I_4\) (row \(i\) of \(K\) subtracts row \(i-1\) of \(T'\) from row \(i\)), so \((T')^{-1}=K\)
(Theorem~\ref{r2-thm-invertible}) and, by Theorem~\ref{r2-thm-inverse-props}(b),(c),
\[
A^{-1}=T^{-1}(T')^{-1}=K'K=\begin{pmatrix}2&-1&0&0\\-1&2&-1&0\\0&-1&2&-1\\0&0&-1&1\end{pmatrix}
\]
\end{example}

\subsection*{Sylvester's law of inertia}

\begin{theorem}[Sylvester's law of inertia; \citeat{lang}{Ch.~V, \S8; Ch.~VIII, \S4, Exercise~25}, \citeat{shakarchi}{Ch.~VIII, \S4, Exercise~25}]\label{r8-thm-inertia}
Let \(A\in\R^{n\times n}\) be symmetric, and let \(P\in\R^{n\times n}\) be invertible with \(P'AP=\diag(d_1,\dots,d_n)\).
\begin{enumerate}
\item The number \(n_0\) of indices \(i\) with \(d_i=0\) is \(\dim\mathcal N(A)=n-\rank A\).
\item The number \(n_+\) of \(i\) with \(d_i>0\) and the number \(n_-\) of \(i\) with \(d_i<0\) do not depend on the choice of
\(P\).
\item \(n_+\) and \(n_-\) are the numbers of positive and of negative eigenvalues of \(A\), counted with multiplicity.
\item There is an invertible \(R\in\R^{n\times n}\) with \(R'AR=\diag(I_{n_+},-I_{n_-},0)\), where \(0\) is the \(n_0\times n_0\)
zero matrix.
\end{enumerate}
\end{theorem}
\begin{proof}
Let \(p_1,\dots,p_n\) be the columns of \(P\). The \((i,j)\) entry of \(P'AP\) is \(p_i'Ap_j\) (Theorem~\ref{r1-thm-views}(d)), so
\(p_i'Ap_j=0\) for \(i\ne j\) and \(p_i'Ap_i=d_i\); and \(p_1,\dots,p_n\) is a basis of \(\R^n\)
(Example~\ref{r3-exm-invertible-basis}).

(a)
\begin{itemize}
\item If \(d_i=0\), then \(p_j'(Ap_i)=0\) for every \(j\), that is, \(P'(Ap_i)=0\); as \(P'\) is invertible, \(Ap_i=0\).
\item Conversely, let \(v\in\mathcal N(A)\) and \(v=\sum_lx_lp_l\). For each \(j\) with \(d_j\ne0\),
\(0=p_j'Av=\sum_lx_lp_j'Ap_l=x_jd_j\), so \(x_j=0\): \(v\) is a combination of the \(p_i\) with \(d_i=0\).
\item These \(p_i\) are independent, so they form a basis of \(\mathcal N(A)\), and \(n_0=\dim\mathcal N(A)=n-\rank A\)
(Theorem~\ref{thm:r4-rank-nullity}).
\end{itemize}
(b)
\begin{itemize}
\item Let \(R\) be another invertible matrix with \(R'AR=\diag(g_1,\dots,g_n)\), with columns \(r_1,\dots,r_n\), and let \(m_+\)
be the number of \(j\) with \(g_j>0\). Let \(\mathcal I=\{i:d_i>0\}\) (\(n_+\) indices) and \(\mathcal J=\{j:g_j\leqslant0\}\) (\(n-m_+\) indices).
\item The \(n_++n-m_+\) vectors \(p_i\) (\(i\in\mathcal I\)) and \(r_j\) (\(j\in\mathcal J\)) are linearly independent. Indeed, let
\(\sum_{i\in\mathcal I}x_ip_i+\sum_{j\in\mathcal J}y_jr_j=0\) and \(u=\sum_{i\in\mathcal I}x_ip_i=-\sum_{j\in\mathcal J}y_jr_j\). Then
\[
u'Au=\sum_{i\in\mathcal I}d_ix_i^2\geqslant0,\qquad u'Au=\sum_{j\in\mathcal J}g_jy_j^2\leqslant0
\]
so \(u'Au=0\), and every \(x_i=0\), as each \(d_i>0\) (\(i\in\mathcal I\)). Then \(\sum_{j\in\mathcal J}y_jr_j=0\), and every \(y_j=0\), as the
columns of \(R\) are independent.
\item Hence \(n_++n-m_+\leqslant n\) (Theorem~\ref{thm:r4-counting}(c)), that is, \(n_+\leqslant m_+\); exchanging the roles of \(P\)
and \(R\) gives \(m_+\leqslant n_+\).
\item By (a), \(n_-=n-n_+-n_0\) does not depend on \(P\) either.
\end{itemize}
(c) By Theorem~\ref{r8-thm-spectral}, \(Q'AQ=\Lambda=\diag(\lambda_1,\dots,\lambda_n)\) with \(Q\) orthogonal, hence invertible,
and \(\lambda_1,\dots,\lambda_n\) are the eigenvalues of \(A\), listed with multiplicity. Apply (b) with \(Q\) in place of \(P\).

(d) Number the orthonormal eigenvectors so that \(\lambda_1,\dots,\lambda_{n_+}>0\), \(\lambda_{n_++1},\dots,\lambda_{n_++n_-}<0\) and
the remaining \(\lambda_i\) are \(0\) (by (c)). Let \(R=Q\,\diag(s_1,\dots,s_n)\), with \(s_i=1/\sqrt{\lvert\lambda_i\rvert}\) if
\(\lambda_i\ne0\) and \(s_i=1\) otherwise. \(R\) is invertible (Theorem~\ref{r2-thm-inverse-props}(b)), and
\(R'AR=\diag(s_1,\dots,s_n)\,\Lambda\,\diag(s_1,\dots,s_n)=\diag(s_1^2\lambda_1,\dots,s_n^2\lambda_n)=\diag(I_{n_+},-I_{n_-},0)\).
\end{proof}

\begin{definition}[\citeat{lang}{Ch.~V, \S8}]\label{r8-def-inertia}
For a symmetric \(A\in\R^{n\times n}\), the numbers \(n_+\), \(n_-\) and \(n_0\) of Theorem~\ref{r8-thm-inertia} are the
\textbf{index of positivity}, the \textbf{index of negativity} and the \textbf{index of nullity} of \(A\) (and of its quadratic
form \(q_A\)); the triple \((n_+,n_-,n_0)\) is the \textbf{inertia} of \(A\).
\end{definition}
By Theorems~\ref{r8-thm-inertia}(c) and~\ref{r8-thm-eigen-criterion}(a),(b), \(A\) is PD if and only if \(n_+=n\), PSD if
and only if \(n_-=0\), and indefinite if and only if \(n_+\geqslant1\) and \(n_-\geqslant1\). So any diagonalization \(P'AP\)
decides the definiteness, for instance the one of Corollary~\ref{r8-cor-ldu}(c), found by completing the squares; no
eigenvalues are needed.

\begin{example}[\citeat{lang}{Ch.~V, \S8, Example}]\label{r8-exm-inertia}
\begin{enumerate}
\item Let \(C=\begin{pmatrix}-1&1\\1&-1\end{pmatrix}\), and let \(P\) have the columns \(v_1=(1,0)'\) and \(v_2=(1,1)'\). Then
\[
P=\begin{pmatrix}1&1\\0&1\end{pmatrix},\qquad
CP=\begin{pmatrix}-1&0\\1&0\end{pmatrix},\qquad
P'CP=\begin{pmatrix}-1&0\\0&0\end{pmatrix}
\]
so the inertia of \(C\) is \((0,1,1)\). Indeed \(C\) has the eigenvalues \(-2\) and \(0\), with the eigenvectors \((1,-1)'\) and \((1,1)'\). Here
\(P\) is the matrix of Lemma~\ref{r8-lem-triangular} (\(p_{12}=-c_{12}/c_{11}=1\)), and
\(P'CP=\diag(\Delta_1,\Delta_2/\Delta_1)\) is Corollary~\ref{r8-cor-ldu}(c).
\item For \(A=2J_3-I_3\) of Example~\ref{r8-exm-spectral-j3}, \(\Delta_1=1\), \(\Delta_2=1-4=-3\) and \(\Delta_3=\det A=5\cdot(-1)\cdot(-1)=5\)
(Theorem~\ref{r7-thm-trace-det}). Solving for \(P\) as in Example~\ref{r8-exm-leading} (\(p_{12}=-2\), and
\(A_2\left(\begin{smallmatrix}p_{13}\\p_{23}\end{smallmatrix}\right)=-\left(\begin{smallmatrix}2\\2\end{smallmatrix}\right)\) gives
\(p_{13}=p_{23}=-\frac23\)),
\[
P=\begin{pmatrix}1&-2&-\frac23\\0&1&-\frac23\\0&0&1\end{pmatrix},\qquad
AP=\begin{pmatrix}1&0&0\\2&-3&0\\2&-2&-\frac53\end{pmatrix},\qquad
U=P^{-1}=\begin{pmatrix}1&2&2\\0&1&\frac23\\0&0&1\end{pmatrix}
\]
and \(P'AP=\diag(1,-3,-\frac53)\) (Corollary~\ref{r8-cor-ldu}(c)). So
\(q_A(x)=(x_1+2x_2+2x_3)^2-3\bigl(x_2+\frac23x_3\bigr)^2-\frac53x_3^2\): the inertia is \((1,2,0)\), and \(A\) is indefinite, as
the eigenvalues \(5\), \(-1\), \(-1\) show (Theorem~\ref{r8-thm-inertia}(c)).
\end{enumerate}
\end{example}

\needspace{14\baselineskip}
\subsection*{The square root of a PSD matrix}

\begin{theorem}[\citeat{hoffman-kunze}{\S9.5, Theorem~13}, \citeat{harville}{Theorem~21.9.1}, \citeat{abadir-magnus}{Exercise~8.25}, \citeat{searle-khuri}{(6.41)}]\label{r8-thm-sqrt}
Let \(A\in\R^{n\times n}\) be PSD, with spectral decomposition \(A=c_1E_1+\dots+c_kE_k\), and let \(Q\), \(\Lambda\) be
as in Theorem~\ref{r8-thm-spectral}.
\begin{enumerate}
\item There is exactly one PSD matrix \(B\) with \(B^2=A\), namely
\[
A^{1/2}:=\sqrt{c_1}\,E_1+\dots+\sqrt{c_k}\,E_k=Q\Lambda^{1/2}Q',\qquad
\Lambda^{1/2}=\diag(\sqrt{\lambda_1},\dots,\sqrt{\lambda_n})
\]
the \textbf{square root} of \(A\), which is \(f(A)\) for \(f(t)=\sqrt t\) (Definition~\ref{r8-def-matrix-function}).
\item \(A^{1/2}A=AA^{1/2}\), \(\rank A^{1/2}=\rank A\), and \(x'Ax=\lVert A^{1/2}x\rVert^2\) for \(x\in\R^n\).
\item If \(A\) is PD, then \(A^{1/2}\) is PD, and \(A^{-1/2}:=(A^{1/2})^{-1}=\sum_jc_j^{-1/2}E_j=Q\Lambda^{-1/2}Q'\), with
\(\Lambda^{-1/2}=\diag(\lambda_1^{-1/2},\dots,\lambda_n^{-1/2})\), satisfies
\[
A^{-1/2}AA^{-1/2}=I_n,\qquad A^{-1/2}A^{-1/2}=A^{-1}
\]
\end{enumerate}
\end{theorem}
\begin{proof}
Every \(c_j\geqslant0\) (Theorem~\ref{r8-thm-eigen-criterion}(a)), so \(\sqrt{c_j}\) is defined. Let
\(N=\sum_j\sqrt{c_j}\,E_j\).

(a)
\begin{itemize}
\item \(N\) is symmetric, \(N^2=\sum_i\sum_j\sqrt{c_i}\sqrt{c_j}\,E_iE_j=\sum_jc_jE_j=A\), and, as in the proof of
Theorem~\ref{r8-thm-eigen-criterion}, \(x'Nx=\sum_j\sqrt{c_j}\lVert E_jx\rVert^2\geqslant0\): \(N\) is PSD. By
Theorem~\ref{r8-thm-spectral-decomp}(c), \(N=\sum_i\sqrt{\lambda_i}\,q_iq_i'=Q\Lambda^{1/2}Q'\).
\item Uniqueness: let \(B\) be PSD with \(B^2=A\), and let \(B=\sum_md_mF_m\) be its spectral decomposition. The \(d_m\)
are distinct and \(\geqslant0\) (Theorem~\ref{r8-thm-eigen-criterion}(a)), so the \(d_m^2\) are distinct, and
\(A=B^2=\sum_md_m^2F_m\) (Theorem~\ref{r8-thm-spectral-decomp}(b) for \(B\)). By
Theorem~\ref{r8-thm-spectral-decomp}(d), after renumbering \(d_m^2=c_m\) and \(F_m=E_m\); so \(d_m=\sqrt{c_m}\) and
\(B=N\).
\end{itemize}
(b)
\begin{itemize}
\item \(AN=\sum_jc_j\sqrt{c_j}\,E_j=NA\), as \(E_iE_j=0\) for \(i\ne j\).
\item \(N=f(A)\) for \(f(t)=\sqrt t\), so its eigenvalues are \(\sqrt{\lambda_1},\dots,\sqrt{\lambda_n}\)
(Proposition~\ref{r8-prop-function}(a)), and \(\sqrt{\lambda_i}\ne0\) exactly when \(\lambda_i\ne0\); hence
\(\rank N=\rank A\) (Theorem~\ref{r8-thm-rank-eigen}, for \(N\) and for \(A\)).
\item \(N\) is symmetric, so \(x'Ax=x'NNx=(Nx)'(Nx)=\lVert Nx\rVert^2\).
\end{itemize}
(c)
\begin{itemize}
\item If \(A\) is PD, every \(c_j>0\), so \(x'Nx>0\) for \(x\ne0\) (some \(E_jx\ne0\)), and
\(N\sum_jc_j^{-1/2}E_j=\sum_jE_j=I_n\) gives \(A^{-1/2}=\sum_jc_j^{-1/2}E_j\) (Theorem~\ref{r2-thm-invertible}). Then
\(A^{-1/2}AA^{-1/2}=\sum_jc_j^{-1/2}c_jc_j^{-1/2}E_j=I_n\).
\item As \(A=N^2\) with \(N\) invertible, \(A^{-1}=(N^2)^{-1}=(N^{-1})^2=A^{-1/2}A^{-1/2}\)
(Theorem~\ref{r2-thm-inverse-props}(b)); and \(A^{-1/2}=f(A)\) for \(f(t)=t^{-1/2}\), so \(A^{-1/2}=Q\Lambda^{-1/2}Q'\) (Proposition~\ref{r8-prop-function}(a)).\qedhere
\end{itemize}
\end{proof}

\begin{example}\label{r8-exm-sqrt}
For \(A=25E_1+9E_2+E_3\) of Example~\ref{r8-exm-spectral}, all eigenvalues are positive, so \(A\) is PD
(Theorem~\ref{r8-thm-eigen-criterion}(a)), and Theorem~\ref{r8-thm-sqrt}(a),(c) gives
\[
A^{1/2}=5E_1+3E_2+E_3=\begin{pmatrix}3&0&1&1\\0&3&1&1\\1&1&3&0\\1&1&0&3\end{pmatrix},\qquad
A^{-1/2}=\tfrac15E_1+\tfrac13E_2+E_3=\frac1{15}\begin{pmatrix}7&2&-3&-3\\2&7&-3&-3\\-3&-3&7&2\\-3&-3&2&7\end{pmatrix}
\]
(for instance the \((1,1)\) entry of \(A^{1/2}\) is \(\frac54+\frac32+\frac14=3\)). Check: row \(1\) of \(A^{1/2}\) times its
columns gives \(9+1+1=11\), \(0+1+1=2\), \(3+3=6\), \(3+3=6\), the first row of \(A\). The matrix
\[
G=5E_1-3E_2+E_3=\begin{pmatrix}0&3&1&1\\3&0&1&1\\1&1&0&3\\1&1&3&0\end{pmatrix}
\]
also satisfies \(G^2=25E_1+9E_2+E_3=A\), because \(E_j^2=E_j\) and \(E_jE_l=0\) for \(j\ne l\)
(Theorem~\ref{r8-thm-spectral-decomp}(b)); but \(G\) is not PSD, as its diagonal entry \(g_{11}=0\) stands in a nonzero row
(Theorem~\ref{r8-thm-principal-sub}). Among the matrices whose square is \(A\), Theorem~\ref{r8-thm-sqrt}(a) singles out
the PSD one.
\end{example}

\subsection*{The Loewner order}

\begin{definition}[\citeat{abadir-magnus}{Ch.~8}]\label{r8-def-loewner}
For symmetric \(A,B\in\R^{n\times n}\) we write \(A\succcurlyeq B\) (or \(B\preccurlyeq A\)) if \(A-B\) is PSD, and \(A\succ B\) if
\(A-B\) is PD. This is the \textbf{Loewner order}. In particular \(A\succcurlyeq0\) means that \(A\) is PSD, and \(A\succ0\) that
\(A\) is PD.
\end{definition}
\begin{example}\label{r8-exm-loewner}
\begin{enumerate}
\item Unlike numbers, two symmetric matrices need not be comparable: for \(A=\diag(1,0)\) and \(B=\diag(0,1)\), neither
\(A-B=\diag(1,-1)\) nor \(B-A\) is PSD (Proposition~\ref{r8-prop-pd-basic}(a)), so neither \(A\succcurlyeq B\) nor
\(B\succcurlyeq A\).
\item The Loewner order is not the entrywise order:
\(I_2-\frac12J_2=\frac12\left(\begin{smallmatrix}1&-1\\-1&1\end{smallmatrix}\right)\) has the eigenvalues \(0\) and \(1\),
with the eigenvectors \((1,1)'\) and \((1,-1)'\), so \(I_2\succcurlyeq\frac12J_2\) (Theorem~\ref{r8-thm-eigen-criterion}(a)),
although the \((1,2)\) entry of \(\frac12J_2\) is the larger one.
\end{enumerate}
\end{example}

\needspace{10\baselineskip}
\begin{proposition}\label{r8-prop-loewner}
Let \(A,B,D\in\R^{n\times n}\) be symmetric.
\begin{enumerate}
\item \(A\succcurlyeq A\); if \(A\succcurlyeq B\) and \(B\succcurlyeq A\), then \(A=B\); if \(A\succcurlyeq B\) and \(B\succcurlyeq D\),
then \(A\succcurlyeq D\).
\item If \(A\succcurlyeq B\), then \(A+D\succcurlyeq B+D\), and \(K'AK\succcurlyeq K'BK\) for every \(K\in\R^{n\times m}\). If
\(A\succ B\) and \(K\in\R^{n\times n}\) is invertible, then \(K'AK\succ K'BK\).
\item If \(A\succcurlyeq B\) and \(B\succ0\), then \(A\succ0\).
\item If \(A\succcurlyeq B\), then \(a_{ii}\geqslant b_{ii}\) for every \(i\), and \(\tr A\geqslant\tr B\).
\end{enumerate}
\end{proposition}
\begin{proof}
(a) \(A-A=0\) is PSD. If \(A-B\) and \(B-A=-(A-B)\) are PSD, then \(A-B=0\) (Proposition~\ref{r8-prop-pd-basic}(e)). If
\(A-B\) and \(B-D\) are PSD, so is their sum \(A-D\) (Proposition~\ref{r8-prop-pd-basic}(b)).

(b) \((A+D)-(B+D)=A-B\), and \(K'AK-K'BK=K'(A-B)K\), which is PSD if \(A-B\) is (Theorem~\ref{r8-thm-congruence}(a)), and PD if
\(A-B\) is PD and \(K\) is invertible (Theorem~\ref{r8-thm-congruence}(c)).

(c) \(A=B+(A-B)\) with \(B\) PD and \(A-B\) PSD, so \(A\) is PD (Proposition~\ref{r8-prop-pd-basic}(b)).

(d) \(A-B\) is PSD, so its diagonal entries \(a_{ii}-b_{ii}\) are \(\geqslant0\) (Theorem~\ref{r8-thm-principal-sub}); adding
over \(i\), \(\tr A-\tr B=\tr(A-B)\geqslant0\) (Theorem~\ref{r1-thm-trace}(a)).
\end{proof}

\begin{theorem}[\citeat{bapat}{Result~3.16}]\label{r8-thm-loewner-inverse}
If \(A\succcurlyeq B\succ0\), then \(B^{-1}\succcurlyeq A^{-1}\).
\end{theorem}
\begin{proof}
\(A\) is PD (Proposition~\ref{r8-prop-loewner}(c)), so \(A^{-1}\) and \(B^{-1}\) exist and are PD
(Proposition~\ref{r8-prop-pd-basic}(c)).
\begin{itemize}
\item First let \(B=I_n\), and let \(A=c_1E_1+\dots+c_kE_k\) be the spectral decomposition. By
Proposition~\ref{r8-prop-function}(c),(a) with \(f(t)=t-1\), the eigenvalues of \(A-I_n\) are \(c_j-1\), and they are \(\geqslant0\),
as \(A-I_n\) is PSD (Theorem~\ref{r8-thm-eigen-criterion}(a)). So every \(c_j\geqslant1\).
\item \(I_n-A^{-1}=g(A)\) for \(g(t)=1-1/t\) (Theorem~\ref{r8-thm-spectral-decomp}(b) and Proposition~\ref{r8-prop-function}(b)), with
the eigenvalues \(1-c_j^{-1}\geqslant0\) (Proposition~\ref{r8-prop-function}(a)). So \(I_n-A^{-1}\) is PSD
(Theorem~\ref{r8-thm-eigen-criterion}(a)), that is,
\(I_n\succcurlyeq A^{-1}\).
\item In general, let \(B^{-1/2}=(B^{1/2})^{-1}\), which is symmetric (Theorems~\ref{r8-thm-sqrt}
and~\ref{r2-thm-inverse-props}(c)), and \(M=B^{-1/2}AB^{-1/2}\). By Proposition~\ref{r8-prop-loewner}(b) with \(K=B^{-1/2}\),
\(M\succcurlyeq B^{-1/2}BB^{-1/2}=I_n\) (Theorem~\ref{r8-thm-sqrt}).
\item By the first part, \(I_n\succcurlyeq M^{-1}=B^{1/2}A^{-1}B^{1/2}\) (Theorem~\ref{r2-thm-inverse-props}(a),(b)).
\item By Proposition~\ref{r8-prop-loewner}(b) with \(K=B^{-1/2}\) again,
\(B^{-1/2}B^{-1/2}\succcurlyeq B^{-1/2}B^{1/2}A^{-1}B^{1/2}B^{-1/2}=A^{-1}\), and
\(B^{-1/2}B^{-1/2}=(B^{1/2}B^{1/2})^{-1}=B^{-1}\) (Theorem~\ref{r2-thm-inverse-props}(b)).\qedhere
\end{itemize}
\end{proof}

\needspace{9\baselineskip}
\begin{theorem}[\citeat{harville}{Theorem~18.1.6, Corollaries~18.1.7, 18.1.8}]\label{r8-thm-det-monotone}
Let \(A\in\R^{n\times n}\) be PD and \(B\in\R^{n\times n}\) PSD. Then \(\det(A+B)\geqslant\det A\), with equality if and only if
\(B=0\). Consequently, if \(A\succcurlyeq B\succ0\), then \(\det A\geqslant\det B\), with equality if and only if \(A=B\); and
\(\det A\geqslant\det B\) holds also when \(A\succcurlyeq B\succcurlyeq0\).
\end{theorem}
\begin{proof}
\begin{itemize}
\item By Theorem~\ref{r8-thm-psd-factor}(b), \(B=K'K\) for some \(K\in\R^{r\times n}\); let \(k_1',\dots,k_r'\) be the rows of
\(K\), and \(K_i\) the \(i\times n\) matrix of its first \(i\) rows. Put \(A_0=A\) and \(A_i=A+K_i'K_i\) (\(i=1,\dots,r\)), so that
\(A_r=A+B\).
\item By the outer product expansion (Theorem~\ref{r1-thm-views}(d)), \(K_i'K_i=K_{i-1}'K_{i-1}+k_ik_i'\), that is,
\(A_i=A_{i-1}+k_ik_i'\) (\(i\geqslant2\); for \(i=1\), \(A_1=A+k_1k_1'\)). Each \(K_i'K_i\) is PSD (Theorem~\ref{r8-thm-XtX}), so each
\(A_i\) is PD, and \(A_{i-1}^{-1}\) is PD (Proposition~\ref{r8-prop-pd-basic}(b),(c)).
\item By Theorem~\ref{r5d-thm-det-lemma}, \(\det A_i=\det A_{i-1}\cdot(1+k_i'A_{i-1}^{-1}k_i)\). Here \(\det A_{i-1}>0\)
(Theorem~\ref{r8-thm-eigen-criterion}(c)) and \(k_i'A_{i-1}^{-1}k_i\geqslant0\), with equality if and only if \(k_i=0\). So
\(\det A_i\geqslant\det A_{i-1}\), with equality if and only if \(k_i=0\).
\item Hence \(\det(A+B)=\det A_r\geqslant\det A_{r-1}\geqslant\dots\geqslant\det A_0=\det A\), with equality if and only if equality
holds at every step, that is, \(k_1=\dots=k_r=0\), that is, \(K=0\). And \(K=0\) if and only if \(B=K'K=0\): if \(K'K=0\), then
\(\tr(K'K)=0\) and \(K=0\) (Theorem~\ref{r1-thm-trace}(d)).
\item If \(A\succcurlyeq B\succ0\), apply this to the PD matrix \(B\) and the PSD matrix \(A-B\):
\(\det A=\det\bigl(B+(A-B)\bigr)\geqslant\det B\), with equality if and only if \(A-B=0\).
\item Let \(A\succcurlyeq B\succcurlyeq0\). Then \(A\succcurlyeq0\) (Proposition~\ref{r8-prop-loewner}(a)), so \(\det A\geqslant0\)
(Theorem~\ref{r8-thm-eigen-criterion}(c)); if \(\det B=0\), this is \(\det A\geqslant\det B\). Otherwise \(B\) is invertible
(Theorem~\ref{r5d-thm-product}(a)), hence PD (Theorem~\ref{r8-thm-eigen-criterion}(c)), and the previous item applies.\qedhere
\end{itemize}
\end{proof}

\begin{example}\label{r8-exm-loewner-inverse}
Let \(B=(\min(i,j))\in\R^{4\times4}\) be the matrix of Example~\ref{r8-exm-cholesky}, \(x=e_4\) and \(A=B+xx'\). Show that
\(A\succcurlyeq B\succ0\), and compare \(A^{-1}\) with \(B^{-1}\) and \(\det A\) with \(\det B\).

\textit{Solution.}
\begin{itemize}
\item \(xx'=e_4e_4'\) has the single nonzero entry \(1\) in position \((4,4)\), so
\[
B=\begin{pmatrix}1&1&1&1\\1&2&2&2\\1&2&3&3\\1&2&3&4\end{pmatrix},\qquad
A=\begin{pmatrix}1&1&1&1\\1&2&2&2\\1&2&3&3\\1&2&3&5\end{pmatrix}
\]
By Example~\ref{r8-exm-cholesky}, \(B\) is PD, \(\det B=1\) and \(B^{-1}=K'K\). The matrix \(A-B=xx'\) is PSD
(Theorem~\ref{r8-thm-XtX}, applied to the \(1\times4\) matrix \(x'\)); so \(A\succcurlyeq B\succ0\).
\item \(w=B^{-1}x=\col_4(B^{-1})=(0,0,-1,1)'\) (Theorem~\ref{r1-thm-views}(b)), and \(x'B^{-1}x=x'w=1\). As \(B^{-1}\) is
symmetric (Theorem~\ref{r2-thm-inverse-props}(c)), \(B^{-1}xx'B^{-1}=ww'\), and Theorem~\ref{r5d-thm-woodbury}(b) gives
\[
B^{-1}=\begin{pmatrix}2&-1&0&0\\-1&2&-1&0\\0&-1&2&-1\\0&0&-1&1\end{pmatrix},\qquad
A^{-1}=B^{-1}-\tfrac12ww'=\frac12\begin{pmatrix}4&-2&0&0\\-2&4&-2&0\\0&-2&3&-1\\0&0&-1&1\end{pmatrix}
\]
The difference \(B^{-1}-A^{-1}=\frac12ww'\), whose only nonzero entries are \(\frac12\) in positions \((3,3)\), \((4,4)\)
and \(-\frac12\) in positions \((3,4)\), \((4,3)\), is PSD (Theorem~\ref{r8-thm-XtX}), as
Theorem~\ref{r8-thm-loewner-inverse} predicts: \(B^{-1}\succcurlyeq A^{-1}\).
\item \(\det A=\det B\cdot(1+x'B^{-1}x)=1\cdot2=2\geqslant1=\det B\) (Theorems~\ref{r5d-thm-det-lemma}
and~\ref{r8-thm-det-monotone}); the inequality is strict, as \(A\ne B\).
\end{itemize}
\end{example}

\needspace{12\baselineskip}
\subsection*{Simultaneous diagonalization}

\begin{lemma}[\citeat{lang}{Ch.~VIII, Theorem~4.3}, \citeat{abadir-magnus}{Exercise~7.51}]\label{r8-lem-invariant}
Let \(A\in\R^{n\times n}\) be symmetric, and let \(W\ne\{0\}\) be a subspace of \(\R^n\) that is \textbf{invariant} under
\(A\): \(Aw\in W\) for every \(w\in W\). Then \(W\) has an orthonormal basis consisting of eigenvectors of \(A\).
\end{lemma}
\begin{proof}
\begin{itemize}
\item Let \(q_1,\dots,q_m\) be an orthonormal basis of \(W\) (Corollary~\ref{r6-cor-onb}(a)), and
\(Q_W=[q_1\ \cdots\ q_m]\in\R^{n\times m}\); so \(Q_W'Q_W=I_m\).
\item Each \(Aq_j\) lies in \(W\), so \(Aq_j=Q_Wc_j\) for some \(c_j\in\R^m\) (Theorem~\ref{r1-thm-views}(b)). With
\(M=[c_1\ \cdots\ c_m]\), \(AQ_W=Q_WM\) (Theorem~\ref{r1-thm-views}(a)); multiplying by \(Q_W'\) on the left gives
\(M=Q_W'AQ_W\), and \(M'=Q_W'A'Q_W=M\).
\item By Theorem~\ref{r8-thm-spectral}(a),(b), applied to \(M\), \(R'MR=\Lambda\) is diagonal for an orthogonal
\(R\in\R^{m\times m}\), so \(MR=RR'MR=R\Lambda\) (Proposition~\ref{r6-prop-orthonormal}(b)). Let \(U=Q_WR\). Then
\(U'U=R'Q_W'Q_WR=R'R=I_m\) and \(AU=AQ_WR=Q_WMR=Q_WR\Lambda=U\Lambda\): the \(m\) columns of \(U\) are orthonormal
eigenvectors of \(A\) (Theorem~\ref{r1-thm-views}(a),(b)).
\item They lie in \(W\), as \(\col_j(U)=Q_W\col_j(R)\), and they are independent (Proposition~\ref{r6-prop-orthogonal-list}(b));
as \(\dim W=m\), they form a basis of \(W\) (Theorem~\ref{thm:r4-counting}(d)).\qedhere
\end{itemize}
\end{proof}

\begin{theorem}[\citeat{lang}{Ch.~VIII, \S4, Exercise~20}, \citeat{shakarchi}{}, \citeat{searle-khuri}{Theorem~7.3}]\label{r8-thm-commuting}
Let \(A,B\in\R^{n\times n}\) be symmetric. There is an orthogonal \(Q\) such that \(Q'AQ\) and \(Q'BQ\) are both diagonal if and only
if \(AB=BA\).
\end{theorem}
\begin{proof}
\begin{itemize}
\item If \(Q'AQ=D_1\) and \(Q'BQ=D_2\) are diagonal, then \(D_1D_2=D_2D_1\), and \(AB=QD_1Q'QD_2Q'=QD_1D_2Q'=QD_2D_1Q'=BA\).
\item Conversely, let \(AB=BA\), let \(c_1,\dots,c_k\) be the distinct eigenvalues of \(A\), and \(W_j=\mathcal N(A-c_jI)\). For
\(x\in W_j\), \(A(Bx)=B(Ax)=c_jBx\), so \(Bx\in W_j\): \(W_j\) is invariant under \(B\).
\item By Lemma~\ref{r8-lem-invariant}, applied to \(B\) and \(W_j\), \(W_j\) has an orthonormal basis of eigenvectors of \(B\);
they are eigenvectors of \(A\) belonging to \(c_j\).
\item Together, these vectors are orthonormal, as \(W_i\perp W_j\) for \(i\ne j\) (Theorem~\ref{r8-thm-sym-orth}), and there are
\(\sum_j\dim W_j=\sum_ja_A(c_j)=n\) of them (Theorem~\ref{r8-thm-rank-eigen}, Definition~\ref{r7-def-eigenvalues}). So they are
independent (Proposition~\ref{r6-prop-orthogonal-list}(b)) and form an orthonormal basis \(q_1,\dots,q_n\) of \(\R^n\)
(Theorem~\ref{thm:r4-counting}(d)).
\item Each \(q_i\) is an eigenvector of \(A\) and of \(B\), so for \(Q=[q_1\ \cdots\ q_n]\), \(Q'AQ\) and \(Q'BQ\) are diagonal
(Theorem~\ref{r8-thm-spectral}(b), for \(A\) and for \(B\)).\qedhere
\end{itemize}
\end{proof}

\begin{example}\label{r8-exm-commuting}
Let
\[
A=\begin{pmatrix}11&2&6&6\\2&11&6&6\\6&6&11&2\\6&6&2&11\end{pmatrix},\qquad
B=\begin{pmatrix}4&2&1&0\\2&4&0&1\\1&0&4&2\\0&1&2&4\end{pmatrix}
\]
where \(A\) is the matrix of Example~\ref{r8-exm-spectral}. Find an orthogonal \(Q\) such that \(Q'AQ\) and \(Q'BQ\) are diagonal.

\textit{Solution.}
\begin{itemize}
\item By rows (Theorem~\ref{r1-thm-views}(c)),
\[
AB=\begin{pmatrix}54&36&47&38\\36&54&38&47\\47&38&54&36\\38&47&36&54\end{pmatrix}
\]
is symmetric, so \(BA=B'A'=(AB)'=AB\) (Theorem~\ref{r1-thm-transpose}). By Theorem~\ref{r8-thm-commuting}, \(Q\) exists.
\item The eigenspaces of \(A\) (Example~\ref{r8-exm-spectral}) are \(W_1=\Span\{\one_4\}\) (for \(25\)),
\(W_2=\{x:x_1+x_2=0,\ x_3+x_4=0\}=\Span\{u_2,u_3\}\) with \(u_2=(-1,1,0,0)'\) and \(u_3=(0,0,-1,1)'\) (for \(9\)), and
\(W_3=\Span\{(1,1,-1,-1)'\}\) (for \(1\)). Every row sum of \(B\) is \(7\), so \(B\one_4=7\one_4\); and
\(B(1,1,-1,-1)'=(5,5,-5,-5)'\).
\item In \(W_2\), \(x=au_2+bu_3=(-a,a,-b,b)'\) and \(Bx=(-2a-b,\ 2a+b,\ -a-2b,\ a+2b)'\). If \(Bx=\mu x\), entries \(2\) and
\(4\) give \(2a+b=\mu a\) and \(a+2b=\mu b\); their sum and their difference give \((\mu-3)(a+b)=0\) and
\((\mu-1)(a-b)=0\). So \(\mu=3\) with \(x=a(-1,1,-1,1)'\), or \(\mu=1\) with \(x=a(-1,1,1,-1)'\); entries \(1\) and \(3\)
then agree. Indeed \(B(1,-1,1,-1)'=(3,-3,3,-3)'\) and \(B(1,-1,-1,1)'=(1,-1,-1,1)'\).
\item The vectors \(\one_4\), \((1,-1,1,-1)'\), \((1,-1,-1,1)'\), \((1,1,-1,-1)'\) are eigenvectors of \(A\) and of \(B\), of
length \(2\) and pairwise orthogonal. So
\[
Q=\frac12\begin{pmatrix}1&1&1&1\\1&-1&-1&1\\1&1&-1&-1\\1&-1&1&-1\end{pmatrix},\qquad
Q'AQ=\diag(25,9,9,1),\qquad Q'BQ=\diag(7,3,1,5)
\]
by Theorem~\ref{r8-thm-spectral}(b), for \(A\) and for \(B\).
\end{itemize}
\end{example}

\needspace{12\baselineskip}
\begin{theorem}[\citeat{bapat}{Ch.~3, Exercise~17}, \citeat{searle-khuri}{Theorem~7.4}, \citeat{harville}{\S21.14}]\label{r8-thm-pd-simdiag}
Let \(A,B\in\R^{n\times n}\) be symmetric, with \(A\) PD.
\begin{enumerate}
\item There is an invertible \(E\in\R^{n\times n}\) with \(E'AE=I_n\) and \(E'BE=\Lambda=\diag(\lambda_1,\dots,\lambda_n)\), and \(E\)
can be chosen so that \(\lambda_1\geqslant\dots\geqslant\lambda_n\).
\item For every real \(\lambda\), \(\det(B-\lambda A)=\det A\cdot(\lambda_1-\lambda)\cdots(\lambda_n-\lambda)\). So
\(\lambda_1,\dots,\lambda_n\) are the roots of \(\det(B-\lambda A)=0\), and they are the eigenvalues of \(A^{-1}B\), listed with
multiplicity.
\item An \(n\times n\) matrix \(E=[f_1\ \cdots\ f_n]\) is as in (a) if and only if \(Bf_i=\lambda_iAf_i\) and \(f_i'Af_j=\delta_{ij}\)
for all \(i,j\). Moreover, if \(Bx=\lambda Ax\) and \(By=\mu Ay\) with \(\lambda\ne\mu\), then \(x'Ay=0\).
\end{enumerate}
\end{theorem}
\begin{proof}
(a)
\begin{itemize}
\item By Theorem~\ref{r8-thm-spectral}, \(P'AP=D=\diag(d_1,\dots,d_n)\) for an orthogonal \(P\), and every \(d_i>0\)
(Theorem~\ref{r8-thm-eigen-criterion}(a)). With \(D^{-1/2}=\diag(d_1^{-1/2},\dots,d_n^{-1/2})\), \(D^{-1/2}P'APD^{-1/2}=I_n\).
\item \(C=D^{-1/2}P'BPD^{-1/2}\) is symmetric (Theorem~\ref{r1-thm-transpose}), so \(Q'CQ=\Lambda\) is diagonal for an orthogonal
\(Q\), whose columns may be numbered so that \(\lambda_1\geqslant\dots\geqslant\lambda_n\) (Theorem~\ref{r8-thm-spectral}(b)).
\item \(E=PD^{-1/2}Q\) is invertible (Theorem~\ref{r2-thm-inverse-props}(b)), \(E'AE=Q'(D^{-1/2}P'APD^{-1/2})Q=Q'Q=I_n\), and
\(E'BE=Q'CQ=\Lambda\).
\end{itemize}
\needspace{4\baselineskip}
(b)
\begin{itemize}
\item For real \(\lambda\), \(E'(B-\lambda A)E=\Lambda-\lambda I_n\), so \(\det E'\cdot\det(B-\lambda A)\cdot\det E
=(\lambda_1-\lambda)\cdots(\lambda_n-\lambda)\) (Theorems~\ref{r5d-thm-product}(b) and~\ref{r5d-thm-basic}(c)).
\item From \(E'AE=I_n\), \((\det E)^2\det A=1\) (Theorems~\ref{r5d-thm-product}(b) and~\ref{r5d-thm-transpose}); multiplying by
\(\det A\) gives \(\det(B-\lambda A)=\det A\cdot(\lambda_1-\lambda)\cdots(\lambda_n-\lambda)\), and \(\det A\ne0\).
\item For real \(t\), \(\det(tI_n-A^{-1}B)=\det A^{-1}\cdot\det(tA-B)=(\det A)^{-1}(-1)^n\det(B-tA)=(t-\lambda_1)\cdots(t-\lambda_n)\)
(Theorem~\ref{r5d-thm-product}(b),(c)). By Definition~\ref{r7-def-charpoly}(a) and Proposition~\ref{r7-prop-poly}(d),
\(p_{A^{-1}B}=(t-\lambda_1)\cdots(t-\lambda_n)\), and Definition~\ref{r7-def-eigenvalues} gives the eigenvalues of \(A^{-1}B\).
\end{itemize}
(c)
\begin{itemize}
\item The \((i,j)\) entries of \(E'AE\) and \(E'BE\) are \(f_i'Af_j\) and \(f_i'Bf_j\) (Theorem~\ref{r1-thm-views}(d)). If
\(Bf_j=\lambda_jAf_j\) and \(f_i'Af_j=\delta_{ij}\) for all \(i,j\), then \(E'AE=I_n\), so \((\det E)^2\det A=1\) and \(E\) is
invertible (Theorem~\ref{r5d-thm-product}(a)), and \(f_i'Bf_j=\lambda_jf_i'Af_j=\lambda_j\delta_{ij}\): \(E'BE=\Lambda\).
\item Conversely, if \(E\) is as in (a), then \(E'AE=I_n\) gives \(AE=(E')^{-1}\), and \(BE=(E')^{-1}\Lambda=AE\Lambda\), whose
column \(i\) reads \(Bf_i=\lambda_iAf_i\) (Theorem~\ref{r1-thm-views}(a),(b)).
\item If \(Bx=\lambda Ax\) and \(By=\mu Ay\), then \(\lambda\,x'Ay=(\lambda Ax)'y=(Bx)'y=x'By=\mu\,x'Ay\), so
\((\lambda-\mu)\,x'Ay=0\).\qedhere
\end{itemize}
\end{proof}

\begin{corollary}[\citeat{abadir-magnus}{Exercises~8.36, 8.37}]\label{r8-cor-real-eigen}
Let \(A,B\in\R^{n\times n}\) be symmetric, with \(A\) PD, and let \(E\) and \(\Lambda=\diag(\lambda_1,\dots,\lambda_n)\) be as
in Theorem~\ref{r8-thm-pd-simdiag}(a).
\begin{enumerate}
\item \(A^{-1}B=E\Lambda E^{-1}\): \(A^{-1}B\) is diagonalizable, its eigenvalues \(\lambda_1,\dots,\lambda_n\) are real,
and column \(i\) of \(E\) is an eigenvector of \(A^{-1}B\) belonging to \(\lambda_i\).
\item \(A^{-1}B\) and the symmetric matrix \(A^{-1/2}BA^{-1/2}\) have the same eigenvalues, with the same multiplicities.
\end{enumerate}
\end{corollary}
\begin{proof}
(a) By Theorem~\ref{r8-thm-pd-simdiag}(c), \(Bf_i=\lambda_iAf_i\) for the columns \(f_i\) of \(E\), that is, \(BE=AE\Lambda\)
(Theorem~\ref{r1-thm-views}(a),(b)). As \(A\) is invertible (Theorem~\ref{r8-thm-eigen-criterion}(c)) and so is \(E\),
\(A^{-1}BE=E\Lambda\) and \(A^{-1}B=E\Lambda E^{-1}\); the rest is the last part of Definition~\ref{r7-def-diagonalizable}.

(b) \(A^{-1/2}\) is symmetric (Theorems~\ref{r8-thm-sqrt}(c) and~\ref{r2-thm-inverse-props}(c)), so \(A^{-1/2}BA^{-1/2}\) is
symmetric (Theorem~\ref{r1-thm-transpose}). As \(A^{-1}=A^{-1/2}A^{-1/2}\) (Theorem~\ref{r8-thm-sqrt}(c)),
\(A^{1/2}(A^{-1}B)A^{-1/2}=A^{1/2}A^{-1/2}A^{-1/2}BA^{-1/2}=A^{-1/2}BA^{-1/2}\); that is,
\(A^{-1/2}BA^{-1/2}=P^{-1}(A^{-1}B)P\) with \(P=A^{-1/2}=(A^{1/2})^{-1}\), as \(P^{-1}=A^{1/2}\)
(Theorem~\ref{r2-thm-inverse-props}(a)), and Theorem~\ref{r7-thm-similar}(a) gives the claim.
\end{proof}
This covers the matrix \(\Sigma_{11}^{-1}\Sigma_{12}\Sigma_{22}^{-1}\Sigma_{12}'\) of canonical correlation analysis,
where \(\Sigma_{11}\) and \(\Sigma_{22}\) are the covariance matrices of two random vectors and \(\Sigma_{12}\) is their
cross-covariance matrix: when \(\Sigma_{11}\) and \(\Sigma_{22}\) are PD, \(A=\Sigma_{11}\) and the symmetric matrix
\(B=\Sigma_{12}\Sigma_{22}^{-1}\Sigma_{12}'\) are as above, so its eigenvalues are real.

\needspace{10\baselineskip}
\begin{example}\label{r8-exm-simultaneous}
Let \(A=(\min(i,j))\in\R^{4\times4}\) be the matrix of Example~\ref{r8-exm-cholesky}, and
\[
B=\begin{pmatrix}3&3&3&3\\3&5&5&5\\3&5&6&6\\3&5&6&5\end{pmatrix}
\]
the matrix with entries \(b_{ij}=c_{\min(i,j)}\), where \((c_1,c_2,c_3,c_4)=(3,5,6,5)\). Find \(E\) as in
Theorem~\ref{r8-thm-pd-simdiag}(a).

\textit{Solution.}
\begin{itemize}
\item \(A\) is PD and \(\det A=1\) (Example~\ref{r8-exm-cholesky}). For real \(t\), the \((i,j)\) entry of \(B-tA\) is
\(d_{\min(i,j)}\), with \(d_m=c_m-tm\): \(d_1=3-t\), \(d_2=5-2t\), \(d_3=6-3t\), \(d_4=5-4t\).
\item Subtract row \(3\) from row \(4\), then row \(2\) from row \(3\), then row \(1\) from row \(2\)
(Definition~\ref{r5d-def-det}(b)). Row \(i\) minus row \(i-1\) has the entry \(d_j-d_j=0\) for \(j<i\) and \(d_i-d_{i-1}\)
for \(j\geqslant i\), so the result is upper triangular, and (Theorem~\ref{r5d-thm-basic}(c))
\[
\det(B-tA)=d_1(d_2-d_1)(d_3-d_2)(d_4-d_3)=(3-t)(2-t)(1-t)(-1-t)
\]
As \(\det A=1\), Theorem~\ref{r8-thm-pd-simdiag}(b) gives \(\lambda_1=3\), \(\lambda_2=2\), \(\lambda_3=1\),
\(\lambda_4=-1\).
\item Let \(f_k=e_k-e_{k-1}\) (\(k=1,\dots,4\), with \(e_0=0\) and \(d_0=0\)). Entry \(i\) of \((B-tA)f_k\) is
\(d_{\min(i,k)}-d_{\min(i,k-1)}\), which is \(0\) for \(i<k\) and \(d_k-d_{k-1}=(c_k-c_{k-1})-t\) for \(i\geqslant k\).
As \(c_k-c_{k-1}=\lambda_k\) (\(c_0=0\)), \((B-\lambda_kA)f_k=0\), that is, \(Bf_k=\lambda_kAf_k\).
\item Let \(E=[f_1\ f_2\ f_3\ f_4]\). The columns of
\[
E=\begin{pmatrix}1&-1&0&0\\0&1&-1&0\\0&0&1&-1\\0&0&0&1\end{pmatrix},\qquad
AE=\begin{pmatrix}1&0&0&0\\1&1&0&0\\1&1&1&0\\1&1&1&1\end{pmatrix}
\]
give \(E'AE=I_4\): the \((j,k)\) entry \(f_j'Af_k\) is entry \(j\) minus entry \(j-1\) of column \(k\) of \(AE\), which is
\(1\) for \(j=k\) and \(0\) otherwise. Then \(E'BE=\diag(3,2,1,-1)\), since \((E'BE)_{jk}=f_j'Bf_k=\lambda_kf_j'Af_k\).
\end{itemize}
Here \(E=K'=T^{-1}\) for the Cholesky factor \(T\) of \(A\) and \(K=(T')^{-1}\) (Example~\ref{r8-exm-cholesky},
Theorem~\ref{r2-thm-inverse-props}(c)), and \(B=T'\diag(3,2,1,-1)\,T\). In the coordinates \(y=E^{-1}x=Tx\), that is,
\(y_k=x_k+\dots+x_4\), \(x'Ax=y_1^2+y_2^2+y_3^2+y_4^2\) and \(x'Bx=3y_1^2+2y_2^2+y_3^2-y_4^2\).
\end{example}

\needspace{20\baselineskip}
\subsection*{Exercises}

\begin{exercise}\label{r8-dr-repeated}
Let
\[
A=\begin{pmatrix}2&-1&2&2\\-1&2&-2&-2\\2&-2&5&4\\2&-2&4&5\end{pmatrix}
\]
\begin{enumerate}
\item Find \(p_A\) and the eigenvalues of \(A\) with their algebraic multiplicities.
\item Find an orthogonal \(Q\) and a diagonal \(\Lambda\) with \(Q'AQ=\Lambda\).
\item Write down the spectral decomposition of \(A\), and check that the projection for the simple eigenvalue is the
polynomial in \(A\) of Proposition~\ref{r8-prop-function}(c).
\item Find \(A^{-1}\) and \(A^m\) for \(m=1,2,\dots\)
\end{enumerate}
\end{exercise}
\begin{solution}
(a) For real \(t\), apply \(R_1\leftarrow R_1+R_2\) and \(R_4\leftarrow R_4-R_3\), take the factor \(t-1\) out of rows
\(1\) and \(4\), apply \(R_2\leftarrow R_2-R_1\), \(R_3\leftarrow R_3+2R_1\), expand along column \(1\), add column \(3\) to
column \(2\), and expand along row \(3\) (Definition~\ref{r5d-def-det}(a),(b), Theorems~\ref{r5d-thm-transpose}
and~\ref{r5d-thm-expansion}):
\[
\begin{aligned}
\det(tI-A)&=\det\begin{pmatrix}t-2&1&-2&-2\\1&t-2&2&2\\-2&2&t-5&-4\\-2&2&-4&t-5\end{pmatrix}
=\det\begin{pmatrix}t-1&t-1&0&0\\1&t-2&2&2\\-2&2&t-5&-4\\0&0&1-t&t-1\end{pmatrix}\\
&=(t-1)^2\det\begin{pmatrix}1&1&0&0\\1&t-2&2&2\\-2&2&t-5&-4\\0&0&-1&1\end{pmatrix}
=(t-1)^2\det\begin{pmatrix}1&1&0&0\\0&t-3&2&2\\0&4&t-5&-4\\0&0&-1&1\end{pmatrix}\\
&=(t-1)^2\det\begin{pmatrix}t-3&2&2\\4&t-5&-4\\0&-1&1\end{pmatrix}
=(t-1)^2\det\begin{pmatrix}t-3&4&2\\4&t-9&-4\\0&0&1\end{pmatrix}\\
&=(t-1)^2\bigl((t-3)(t-9)-16\bigr)=(t-1)^2(t^2-12t+11)=(t-1)^3(t-11)
\end{aligned}
\]
So \(p_A(t)=(t-1)^3(t-11)\) (Definition~\ref{r7-def-charpoly}(a), Proposition~\ref{r7-prop-poly}(d)): the eigenvalues are
\(11\) with \(a_A(11)=1\) and \(1\) with \(a_A(1)=3\).

(b) Row-reduce \(A-11I\):
\[
\begin{gathered}
A-11I=\begin{pmatrix}-9&-1&2&2\\-1&-9&-2&-2\\2&-2&-6&4\\2&-2&4&-6\end{pmatrix}
\xrightarrow[R_3\leftarrow R_3-R_4]{R_1\leftarrow R_1+R_2}
\begin{pmatrix}-10&-10&0&0\\-1&-9&-2&-2\\0&0&-10&10\\2&-2&4&-6\end{pmatrix}
\\
\xrightarrow[R_3\leftarrow-\frac1{10}R_3]{R_1\leftarrow-\frac1{10}R_1}
\begin{pmatrix}1&1&0&0\\-1&-9&-2&-2\\0&0&1&-1\\2&-2&4&-6\end{pmatrix}
\xrightarrow[R_4\leftarrow R_4-2R_1]{R_2\leftarrow R_2+R_1}
\begin{pmatrix}1&1&0&0\\0&-8&-2&-2\\0&0&1&-1\\0&-4&4&-6\end{pmatrix}
\xrightarrow[R_4\leftarrow R_4-4R_3]{R_2\leftarrow R_2+2R_3}
\begin{pmatrix}1&1&0&0\\0&-8&0&-4\\0&0&1&-1\\0&-4&0&-2\end{pmatrix}
\\
\xrightarrow{R_4\leftarrow R_4-\frac12R_2}
\begin{pmatrix}1&1&0&0\\0&-8&0&-4\\0&0&1&-1\\0&0&0&0\end{pmatrix}
\end{gathered}
\]
So \(x_4=2\) gives \(x_3=x_4=2\), \(x_2=-\frac12x_4=-1\) and \(x_1=-x_2=1\): \(u=(1,-1,2,2)'\) spans \(\mathcal N(A-11I)\). Next,
\[
A-I=\begin{pmatrix}1&-1&2&2\\-1&1&-2&-2\\2&-2&4&4\\2&-2&4&4\end{pmatrix}
\xrightarrow[R_4\leftarrow R_4-2R_1]{R_2\leftarrow R_2+R_1,\ R_3\leftarrow R_3-2R_1}
\begin{pmatrix}1&-1&2&2\\0&0&0&0\\0&0&0&0\\0&0&0&0\end{pmatrix}
\]
so \(\mathcal N(A-I)=\{x:x_1=x_2-2x_3-2x_4\}\) has the basis \(s_1=(1,1,0,0)'\), \(s_2=(-2,0,1,0)'\), \(s_3=(-2,0,0,1)'\).
Gram--Schmidt (Theorem~\ref{r6-thm-gram-schmidt}):
\begin{itemize}
\item \(b_1=s_1\), with \(\lVert b_1\rVert^2=2\);
\item \(s_2'b_1=-2\), so \(b_2=s_2+b_1=(-1,1,1,0)'\), with \(\lVert b_2\rVert^2=3\);
\item \(s_3'b_1=-2\) and \(s_3'b_2=2\), so \(b_3=s_3+b_1-\frac23b_2=\frac13(-1,1,-2,3)'\); we use \(3b_3=(-1,1,-2,3)'\), with
\(\lVert3b_3\rVert^2=15\).
\end{itemize}
The vector \(u\) is orthogonal to \(b_1\), \(b_2\) and \(b_3\) (Theorem~\ref{r8-thm-sym-orth}). Normalizing
(Theorem~\ref{r8-thm-spectral}(b)),
\[
Q=\begin{pmatrix}1/\sqrt2&-1/\sqrt3&-1/\sqrt{15}&1/\sqrt{10}\\1/\sqrt2&1/\sqrt3&1/\sqrt{15}&-1/\sqrt{10}\\
0&1/\sqrt3&-2/\sqrt{15}&2/\sqrt{10}\\0&0&3/\sqrt{15}&2/\sqrt{10}\end{pmatrix},\qquad
\Lambda=\diag(1,1,1,11)
\]
Check: \(A(-1,1,-2,3)'=(-1,1,-2,3)'\) and \(A(1,-1,2,2)'=(11,-11,22,22)'\).

(c) With \(c_1=11\) and \(c_2=1\), Theorem~\ref{r8-thm-spectral-decomp}(b),(c) gives \(E_1=q_4q_4'=\frac1{10}uu'\) and
\(E_2=I_4-E_1\):
\[
E_1=\frac1{10}\begin{pmatrix}1&-1&2&2\\-1&1&-2&-2\\2&-2&4&4\\2&-2&4&4\end{pmatrix},\qquad
E_2=\frac1{10}\begin{pmatrix}9&1&-2&-2\\1&9&2&2\\-2&2&6&-4\\-2&2&-4&6\end{pmatrix},\qquad A=11E_1+E_2
\]
The polynomial \(e_1=\frac{t-1}{11-1}\) gives \(E_1=\frac1{10}(A-I)\), and indeed \(A-I=uu'\) is
the matrix reduced in (b).

(d) No eigenvalue is \(0\), so by Theorem~\ref{r8-thm-spectral-decomp}(b)
\[
A^{-1}=\tfrac1{11}E_1+E_2=I_4-\tfrac{10}{11}E_1=\frac1{11}\begin{pmatrix}10&1&-2&-2\\1&10&2&2\\-2&2&7&-4\\-2&2&-4&7\end{pmatrix},\qquad
A^m=11^mE_1+E_2=I_4+\tfrac{11^m-1}{10}\,uu'
\]
Check: \(\row_1(A)A^{-1}=\frac1{11}(20-1-4-4,\ 2-10+4+4,\ -4-2+14-8,\ -4-2-8+14)=(1,0,0,0)\), and likewise for rows \(2\), \(3\)
and \(4\).
\end{solution}

\begin{exercise}\label{r8-ex-projectors}
Let
\[
A=\begin{pmatrix}-1&-2&2&-2\\-2&-1&-2&2\\2&-2&-1&-2\\-2&2&-2&-1\end{pmatrix}
\]
\begin{enumerate}
\item Show that \(p_A(t)=(t-5)(t+3)^3\), find an orthonormal basis of each eigenspace, and an orthogonal \(Q\) with
\(Q'AQ\) diagonal.
\item Find the projections \(E_1\) and \(E_2\) of the spectral decomposition of \(A\) (\(c_1=5\), \(c_2=-3\)), onto
\(\mathcal N(A-5I)\) and \(\mathcal N(A+3I)\), and check that \(A=5E_1-3E_2\).
\item Use Definition~\ref{r8-def-matrix-function} to compute \(A^{-1}\) and \(A^3\).
\end{enumerate}
\end{exercise}
\begin{solution}
(a) For real \(t\), subtract rows \(2\) and \(4\) of \(tI_4-A\) from row \(1\) and add row \(3\) to it; the new row \(1\) is
\((t-5)(1,-1,1,-1)\). Take out \(t-5\), and add column \(1\) to columns \(2\) and \(4\) and subtract it from column \(3\)
(Definition~\ref{r5d-def-det}(a),(b), Theorem~\ref{r5d-thm-transpose}):
\[
\begin{aligned}
\det(tI_4-A)&=\det\begin{pmatrix}t+1&2&-2&2\\2&t+1&2&-2\\-2&2&t+1&2\\2&-2&2&t+1\end{pmatrix}
=(t-5)\det\begin{pmatrix}1&-1&1&-1\\2&t+1&2&-2\\-2&2&t+1&2\\2&-2&2&t+1\end{pmatrix}\\
&=(t-5)\det\begin{pmatrix}1&0&0&0\\2&t+3&0&0\\-2&0&t+3&0\\2&0&0&t+3\end{pmatrix}=(t-5)(t+3)^3
\end{aligned}
\]
as the last matrix is lower triangular (Theorem~\ref{r5d-thm-transpose}). So \(p_A(t)=(t-5)(t+3)^3\)
(Definition~\ref{r7-def-charpoly}(a), Proposition~\ref{r7-prop-poly}(d)).
\begin{itemize}
\item \(v=(1,-1,1,-1)'\) has \(Av=(-1+2+2+2,\ -2+1-2-2,\ 2+2-1+2,\ -2-2-2+1)'=5v\), and
\(\dim\mathcal N(A-5I)=a_A(5)=1\) (Theorem~\ref{r8-thm-rank-eigen}); so \(u_1=\frac12v\) is an orthonormal basis of
\(\mathcal N(A-5I)\).
\item \(A+3I_4=2vv'\): row \(i\) of \(A+3I_4\) is \(2v_iv'\). So \(\mathcal N(A+3I)=\{x:x_1-x_2+x_3-x_4=0\}\), of dimension
\(a_A(-3)=3\) (Theorem~\ref{r8-thm-rank-eigen}). The orthogonal vectors \((1,1,1,1)'\), \((1,1,-1,-1)'\), \((1,-1,-1,1)'\) lie
in it, so \(u_2=\frac12(1,1,1,1)'\), \(u_3=\frac12(1,1,-1,-1)'\), \(u_4=\frac12(1,-1,-1,1)'\) is an orthonormal basis of
\(\mathcal N(A+3I)\) (Proposition~\ref{r6-prop-orthogonal-list}(b), Theorem~\ref{thm:r4-counting}(d)).
\item \(Q=[u_1\ u_2\ u_3\ u_4]\) is orthogonal (\(u_1\perp u_2,u_3,u_4\) by Theorem~\ref{r8-thm-sym-orth}) and
\(Q'AQ=\diag(5,-3,-3,-3)\) (Theorem~\ref{r8-thm-spectral}(b)).
\end{itemize}
(b) By Theorem~\ref{r8-thm-spectral-decomp}(b),(c),
\[
E_1=u_1u_1'=\frac14\begin{pmatrix}1&-1&1&-1\\-1&1&-1&1\\1&-1&1&-1\\-1&1&-1&1\end{pmatrix},\qquad
E_2=I_4-E_1=\frac14\begin{pmatrix}3&1&-1&1\\1&3&1&-1\\-1&1&3&1\\1&-1&1&3\end{pmatrix}
\]
and \(5E_1-3E_2=8E_1-3I_4=2vv'-3I_4=A\): the diagonal entries are \(2-3=-1\) and the others are \(2v_iv_j=\pm2\).

(c) By Definition~\ref{r8-def-matrix-function} with \(f(t)=1/t\) and \(f(t)=t^3\), which give \(A^{-1}\) and \(A^3\)
(Proposition~\ref{r8-prop-function}(b),(c)),
\[
\begin{gathered}
A^{-1}=\tfrac15E_1-\tfrac13E_2=\tfrac8{15}E_1-\tfrac13I_4=\tfrac1{15}(2vv'-5I_4)
=\frac1{15}\begin{pmatrix}-3&-2&2&-2\\-2&-3&-2&2\\2&-2&-3&-2\\-2&2&-2&-3\end{pmatrix}\\
A^3=125E_1-27E_2=152E_1-27I_4=38vv'-27I_4
=\begin{pmatrix}11&-38&38&-38\\-38&11&-38&38\\38&-38&11&-38\\-38&38&-38&11\end{pmatrix}
\end{gathered}
\]
Check: row \(1\) of \(A\) times the first two columns of \(15A^{-1}\) gives \(3+4+4+4=15\) and \(2+6-4-4=0\). (Thus
\(A^{-1}=\frac1{15}(A-2I_4)\), that is, \(A^2-2A-15I_4=(A-5I_4)(A+3I_4)=0\).)
\end{solution}

\begin{exercise}\label{r8-dr-diagonalize}
Let
\[
B=\begin{pmatrix}3&1&0&2\\1&2&1&2\\0&1&3&2\\2&2&2&0\end{pmatrix}
\]
\begin{enumerate}
\item Find \(p_B\) and the eigenvalues of \(B\).
\item Find an orthogonal \(Q\) and a diagonal \(\Lambda\) with \(Q'BQ=\Lambda\).
\item Write down the spectral decomposition of \(B\), and use it to find \(B^{-1}\).
\end{enumerate}
\end{exercise}
\begin{solution}
(a) For real \(t\), add rows \(2\), \(3\) and \(4\) to row \(1\), take the factor \(t-6\) out of row \(1\), subtract
column \(1\) from columns \(2\), \(3\) and \(4\), and expand along row \(1\) and then along row \(3\)
(Definition~\ref{r5d-def-det}(a),(b), Theorems~\ref{r5d-thm-transpose} and~\ref{r5d-thm-expansion}):
\[
\begin{aligned}
\det(tI-B)&=\det\begin{pmatrix}t-3&-1&0&-2\\-1&t-2&-1&-2\\0&-1&t-3&-2\\-2&-2&-2&t\end{pmatrix}
=\det\begin{pmatrix}t-6&t-6&t-6&t-6\\-1&t-2&-1&-2\\0&-1&t-3&-2\\-2&-2&-2&t\end{pmatrix}\\
&=(t-6)\det\begin{pmatrix}1&1&1&1\\-1&t-2&-1&-2\\0&-1&t-3&-2\\-2&-2&-2&t\end{pmatrix}
=(t-6)\det\begin{pmatrix}1&0&0&0\\-1&t-1&0&-1\\0&-1&t-3&-2\\-2&0&0&t+2\end{pmatrix}\\
&=(t-6)\det\begin{pmatrix}t-1&0&-1\\-1&t-3&-2\\0&0&t+2\end{pmatrix}
=(t-6)(t+2)\bigl((t-1)(t-3)-0\bigr)
\end{aligned}
\]
So \(p_B(t)=(t-6)(t-3)(t-1)(t+2)\) (Definition~\ref{r7-def-charpoly}(a), Proposition~\ref{r7-prop-poly}(d)), and the
eigenvalues are \(6\), \(3\), \(1\), \(-2\).

(b) Row-reduce \(B-cI\) for \(c=6,3,1,-2\):
\[
\begin{gathered}
B-6I=\begin{pmatrix}-3&1&0&2\\1&-4&1&2\\0&1&-3&2\\2&2&2&-6\end{pmatrix}
\xrightarrow{R_4\leftarrow\frac12R_4}
\begin{pmatrix}-3&1&0&2\\1&-4&1&2\\0&1&-3&2\\1&1&1&-3\end{pmatrix}
\xrightarrow{R_1\leftrightarrow R_4}
\begin{pmatrix}1&1&1&-3\\1&-4&1&2\\0&1&-3&2\\-3&1&0&2\end{pmatrix}
\\
\xrightarrow[R_4\leftarrow R_4+3R_1]{R_2\leftarrow R_2-R_1}
\begin{pmatrix}1&1&1&-3\\0&-5&0&5\\0&1&-3&2\\0&4&3&-7\end{pmatrix}
\xrightarrow{R_2\leftarrow-\frac15R_2}
\begin{pmatrix}1&1&1&-3\\0&1&0&-1\\0&1&-3&2\\0&4&3&-7\end{pmatrix}
\xrightarrow[R_4\leftarrow R_4-4R_2]{R_3\leftarrow R_3-R_2}
\begin{pmatrix}1&1&1&-3\\0&1&0&-1\\0&0&-3&3\\0&0&3&-3\end{pmatrix}
\\
\xrightarrow{R_4\leftarrow R_4+R_3}
\begin{pmatrix}1&1&1&-3\\0&1&0&-1\\0&0&-3&3\\0&0&0&0\end{pmatrix}
\end{gathered}
\]
so \(x_4=1\) gives \(x_3=x_4=1\), \(x_2=x_4=1\) and \(x_1=-x_2-x_3+3x_4=1\);
\[
\begin{gathered}
B-3I=\begin{pmatrix}0&1&0&2\\1&-1&1&2\\0&1&0&2\\2&2&2&-3\end{pmatrix}
\xrightarrow{R_1\leftrightarrow R_2}
\begin{pmatrix}1&-1&1&2\\0&1&0&2\\0&1&0&2\\2&2&2&-3\end{pmatrix}
\xrightarrow[R_4\leftarrow R_4-2R_1]{R_3\leftarrow R_3-R_2}
\begin{pmatrix}1&-1&1&2\\0&1&0&2\\0&0&0&0\\0&4&0&-7\end{pmatrix}
\\
\xrightarrow{R_4\leftarrow R_4-4R_2}
\begin{pmatrix}1&-1&1&2\\0&1&0&2\\0&0&0&0\\0&0&0&-15\end{pmatrix}
\xrightarrow{R_3\leftrightarrow R_4}
\begin{pmatrix}1&-1&1&2\\0&1&0&2\\0&0&0&-15\\0&0&0&0\end{pmatrix}
\end{gathered}
\]
so \(x_4=0\), \(x_2=-2x_4=0\), and \(x_3=-1\) gives \(x_1=x_2-x_3-2x_4=1\);
\[
\begin{gathered}
B-I=\begin{pmatrix}2&1&0&2\\1&1&1&2\\0&1&2&2\\2&2&2&-1\end{pmatrix}
\xrightarrow{R_1\leftrightarrow R_2}
\begin{pmatrix}1&1&1&2\\2&1&0&2\\0&1&2&2\\2&2&2&-1\end{pmatrix}
\xrightarrow[R_4\leftarrow R_4-2R_1]{R_2\leftarrow R_2-2R_1}
\begin{pmatrix}1&1&1&2\\0&-1&-2&-2\\0&1&2&2\\0&0&0&-5\end{pmatrix}
\\
\xrightarrow{R_3\leftarrow R_3+R_2}
\begin{pmatrix}1&1&1&2\\0&-1&-2&-2\\0&0&0&0\\0&0&0&-5\end{pmatrix}
\xrightarrow{R_3\leftrightarrow R_4}
\begin{pmatrix}1&1&1&2\\0&-1&-2&-2\\0&0&0&-5\\0&0&0&0\end{pmatrix}
\end{gathered}
\]
so \(x_4=0\), and \(x_3=1\) gives \(x_2=-2x_3-2x_4=-2\) and \(x_1=-x_2-x_3-2x_4=1\);
\[
\begin{gathered}
B+2I=\begin{pmatrix}5&1&0&2\\1&4&1&2\\0&1&5&2\\2&2&2&2\end{pmatrix}
\xrightarrow{R_4\leftarrow\frac12R_4}
\begin{pmatrix}5&1&0&2\\1&4&1&2\\0&1&5&2\\1&1&1&1\end{pmatrix}
\xrightarrow{R_1\leftrightarrow R_4}
\begin{pmatrix}1&1&1&1\\1&4&1&2\\0&1&5&2\\5&1&0&2\end{pmatrix}
\\
\xrightarrow[R_4\leftarrow R_4-5R_1]{R_2\leftarrow R_2-R_1}
\begin{pmatrix}1&1&1&1\\0&3&0&1\\0&1&5&2\\0&-4&-5&-3\end{pmatrix}
\xrightarrow{R_2\leftrightarrow R_3}
\begin{pmatrix}1&1&1&1\\0&1&5&2\\0&3&0&1\\0&-4&-5&-3\end{pmatrix}
\xrightarrow[R_4\leftarrow R_4+4R_2]{R_3\leftarrow R_3-3R_2}
\begin{pmatrix}1&1&1&1\\0&1&5&2\\0&0&-15&-5\\0&0&15&5\end{pmatrix}
\\
\xrightarrow{R_4\leftarrow R_4+R_3}
\begin{pmatrix}1&1&1&1\\0&1&5&2\\0&0&-15&-5\\0&0&0&0\end{pmatrix}
\end{gathered}
\]
so \(x_4=-3\) gives \(x_3=-\frac13x_4=1\), \(x_2=-5x_3-2x_4=1\) and \(x_1=-x_2-x_3-x_4=1\). So the eigenspaces are spanned
by \((1,1,1,1)'\), \((1,0,-1,0)'\), \((1,-2,1,0)'\) and \((1,1,1,-3)'\). They belong to distinct eigenvalues, so they are
pairwise orthogonal (Theorem~\ref{r8-thm-sym-orth}). Normalizing them (Theorem~\ref{r8-thm-spectral}(b)),
\[
Q=\begin{pmatrix}1/2&1/\sqrt2&1/\sqrt6&1/\sqrt{12}\\1/2&0&-2/\sqrt6&1/\sqrt{12}\\1/2&-1/\sqrt2&1/\sqrt6&1/\sqrt{12}\\
1/2&0&0&-3/\sqrt{12}\end{pmatrix},\qquad\Lambda=\diag(6,3,1,-2)
\]
(c) Each eigenvalue has multiplicity \(1\), so \(E_j=q_jq_j'\) (Theorem~\ref{r8-thm-spectral-decomp}(c)), and
\(B=6E_1+3E_2+E_3-2E_4\) with
\[
\begin{gathered}
E_1=\frac14\begin{pmatrix}1&1&1&1\\1&1&1&1\\1&1&1&1\\1&1&1&1\end{pmatrix},\qquad
E_2=\frac12\begin{pmatrix}1&0&-1&0\\0&0&0&0\\-1&0&1&0\\0&0&0&0\end{pmatrix}\\
E_3=\frac16\begin{pmatrix}1&-2&1&0\\-2&4&-2&0\\1&-2&1&0\\0&0&0&0\end{pmatrix},\qquad
E_4=\frac1{12}\begin{pmatrix}1&1&1&-3\\1&1&1&-3\\1&1&1&-3\\-3&-3&-3&9\end{pmatrix}
\end{gathered}
\]
(for instance, entry \((1,1)\) of the sum is \(\frac64+\frac32+\frac16-\frac2{12}=3\)). No eigenvalue is \(0\), so by
Theorem~\ref{r8-thm-spectral-decomp}(b)
\[
B^{-1}=\tfrac16E_1+\tfrac13E_2+E_3-\tfrac12E_4=\frac16\begin{pmatrix}2&-2&0&1\\-2&4&-2&1\\0&-2&2&1\\1&1&1&-2\end{pmatrix}
\]
Check: \(\row_1(B)B^{-1}=\frac16(6-2+0+2,\ -6+4+0+2,\ 0-2+0+2,\ 3+1+0-4)=(1,0,0,0)\), and likewise for rows \(2\), \(3\)
and \(4\).
\end{solution}

\begin{exercise}\label{r8-ex-diagonalize}
Let
\[
A=\begin{pmatrix}5&3&2&0\\3&5&0&2\\2&0&5&3\\0&2&3&5\end{pmatrix}
\]
and let \(h_1=\one_4\), \(h_2=(1,1,-1,-1)'\), \(h_3=(1,-1,1,-1)'\) and \(h_4=(1,-1,-1,1)'\) (the vectors of
Example~\ref{r8-exm-sylvester}).
\begin{enumerate}
\item Show that \(h_1,\dots,h_4\) are eigenvectors of \(A\). Find an orthogonal \(Q\) and a diagonal \(\Lambda\) with
\(A=Q\Lambda Q'\) and decreasing diagonal entries, and write the spectral decomposition of \(A\).
\item Find \(\rank A\) and bases of \(\mathcal C(A)\) and \(\mathcal N(A)\). Is \(A\) PD? PSD?
\item Find the largest and the smallest value of \(x'Ax\) on the unit sphere \(x'x=1\) of \(\R^4\), and all points where
they are attained.
\end{enumerate}
\end{exercise}
\begin{solution}
(a) Multiplying the rows of \(A\) by \(h_1,\dots,h_4\):
\[
Ah_1=\begin{pmatrix}10\\10\\10\\10\end{pmatrix}=10h_1,\quad
Ah_2=\begin{pmatrix}6\\6\\-6\\-6\end{pmatrix}=6h_2,\quad
Ah_3=\begin{pmatrix}4\\-4\\4\\-4\end{pmatrix}=4h_3,\quad
Ah_4=\begin{pmatrix}5-3-2\\3-5+2\\2-5+3\\-2-3+5\end{pmatrix}=0
\]
The \(h_j\) are orthogonal with \(\lVert h_j\rVert=2\), so \(Q=\frac12[h_1\ h_2\ h_3\ h_4]\) is orthogonal, \(AQ=Q\Lambda\) with
\(\Lambda=\diag(10,6,4,0)\), and \(A=Q\Lambda Q'\); so \(10,6,4,0\) are the eigenvalues of \(A\)
(Theorem~\ref{r8-thm-spectral}(b),(c)). By
Theorem~\ref{r8-thm-spectral-decomp}(c), with \(q_j=\frac12h_j\),
\[
A=10q_1q_1'+6q_2q_2'+4q_3q_3'=\tfrac52J_4+\tfrac32h_2h_2'+h_3h_3'
\]
For instance, the entries in row \(1\) are \(\frac52+\frac32+1=5\), \(\frac52+\frac32-1=3\), \(\frac52-\frac32+1=2\) and
\(\frac52-\frac32-1=0\).

(b) \(\rank A=3\), the number of nonzero eigenvalues (Theorem~\ref{r8-thm-rank-eigen}), and
\(\mathcal C(A)=\Span\{h_1,h_2,h_3\}\), \(\mathcal N(A)=\Span\{h_4\}\) (Theorem~\ref{r8-thm-spectral-decomp}(c)); these vectors
are independent (Proposition~\ref{r6-prop-orthogonal-list}(b)), so they form bases. The eigenvalues \(10,6,4,0\) are
\(\geqslant0\) and one is \(0\): \(A\) is PSD but not PD (Theorem~\ref{r8-thm-eigen-criterion}(a)).

(c) Put \(y=Q'x\). Then \(x=QQ'x=Qy\) and \(y'y=x'x=1\) (Proposition~\ref{r6-prop-orthonormal}(b),(a)), and
\(x'Ax=y'\Lambda y=10y_1^2+6y_2^2+4y_3^2\). Hence \(x'Ax=10-4y_2^2-6y_3^2-10y_4^2\leqslant10\), with equality exactly when
\(y=\pm e_1\), that is, at \(x=\pm q_1=\pm\frac12(1,1,1,1)'\); and \(x'Ax\geqslant0\), with equality exactly when \(y=\pm e_4\),
that is, at \(x=\pm q_4=\pm\frac12(1,-1,-1,1)'\). The largest value is \(10\) and the smallest is \(0\)
(Theorem~\ref{r8-thm-eigen-criterion}(d)).
\end{solution}

\begin{exercise}\label{r8-dr-function}
Let
\[
A=\begin{pmatrix}6&5&0&5\\5&7&1&4\\0&1&2&-1\\5&4&-1&7\end{pmatrix}
\]
\begin{enumerate}
\item Check that \((1,1,0,1)'\), \((0,1,1,-1)'\), \((0,1,-2,-1)'\) and \((2,-1,0,-1)'\) are eigenvectors of \(A\), and find the
spectral decomposition of \(A\).
\item Check that the projection \(E_3\) for the eigenvalue \(c_3=1\) (with \(c_1=16\), \(c_2=4\)) equals \(\frac1{45}(A-16I)(A-4I)\)
(Proposition~\ref{r8-prop-function}(c)).
\item Find \(A^{1/2}\), \(A^{-1/2}\) and \(A^{-1}\), and check \((A^{1/2})^2=A\) and \(A^{1/2}A^{-1/2}=I\).
\end{enumerate}
\end{exercise}
\begin{solution}
(a) \(A(1,1,0,1)'=(16,16,0,16)'\), \(A(0,1,1,-1)'=(0,4,4,-4)'\), \(A(0,1,-2,-1)'=(0,1,-2,-1)'\) and
\(A(2,-1,0,-1)'=(2,-1,0,-1)'\). The last two are orthogonal, hence independent (Proposition~\ref{r6-prop-orthogonal-list}(b)),
so \(\dim\mathcal N(A-I)\geqslant2\) (Theorem~\ref{thm:r4-counting}(c)). So \(a_A(16)\geqslant1\), \(a_A(4)\geqslant1\) and
\(a_A(1)\geqslant2\) (Theorem~\ref{r8-thm-rank-eigen}); as the multiplicities add up to \(4\)
(Definition~\ref{r7-def-eigenvalues}), these are the eigenvalues, with \(a_A(16)=a_A(4)=1\) and \(a_A(1)=2\), and \(A\) is PD
(Theorem~\ref{r8-thm-eigen-criterion}(a)). The four vectors are pairwise orthogonal; with the unit vectors
\(q_1=\frac1{\sqrt3}(1,1,0,1)'\), \(q_2=\frac1{\sqrt3}(0,1,1,-1)'\), \(q_3=\frac1{\sqrt6}(0,1,-2,-1)'\) and
\(q_4=\frac1{\sqrt6}(2,-1,0,-1)'\), Theorem~\ref{r8-thm-spectral-decomp}(c), with \(c_1=16\), \(c_2=4\), \(c_3=1\), gives \(E_1=q_1q_1'\), \(E_2=q_2q_2'\),
\(E_3=q_3q_3'+q_4q_4'\):
\[
E_1=\frac13\begin{pmatrix}1&1&0&1\\1&1&0&1\\0&0&0&0\\1&1&0&1\end{pmatrix},\qquad
E_2=\frac13\begin{pmatrix}0&0&0&0\\0&1&1&-1\\0&1&1&-1\\0&-1&-1&1\end{pmatrix},\qquad
E_3=\frac16\begin{pmatrix}4&-2&0&-2\\-2&2&-2&0\\0&-2&4&2\\-2&0&2&2\end{pmatrix}
\]
and \(A=16E_1+4E_2+E_3\) (for instance, entry \((1,1)\): \(\frac{16}3+0+\frac46=6\)).

(b) \(e_3=\frac{(t-16)(t-4)}{(1-16)(1-4)}=\frac1{45}(t-16)(t-4)\), and
\[
(A-16I)(A-4I)=\begin{pmatrix}-10&5&0&5\\5&-9&1&4\\0&1&-14&-1\\5&4&-1&-9\end{pmatrix}
\begin{pmatrix}2&5&0&5\\5&3&1&4\\0&1&-2&-1\\5&4&-1&3\end{pmatrix}
=\begin{pmatrix}30&-15&0&-15\\-15&15&-15&0\\0&-15&30&15\\-15&0&15&15\end{pmatrix}=45E_3
\]
(c) By Definition~\ref{r8-def-matrix-function} with \(f(t)=\sqrt t\), \(t^{-1/2}\) and \(t^{-1}\) (Theorem~\ref{r8-thm-sqrt},
Proposition~\ref{r8-prop-function}(b)),
\[
\begin{gathered}
A^{1/2}=4E_1+2E_2+E_3=\frac13\begin{pmatrix}6&3&0&3\\3&7&1&2\\0&1&4&-1\\3&2&-1&7\end{pmatrix}\\
A^{-1/2}=\tfrac14E_1+\tfrac12E_2+E_3=\frac1{12}\begin{pmatrix}9&-3&0&-3\\-3&7&-2&-1\\0&-2&10&2\\-3&-1&2&7\end{pmatrix}\\
A^{-1}=\tfrac1{16}E_1+\tfrac14E_2+E_3=\frac1{16}\begin{pmatrix}11&-5&0&-5\\-5&7&-4&-1\\0&-4&12&4\\-5&-1&4&7\end{pmatrix}
\end{gathered}
\]
Check: \(\row_1(A^{1/2})A^{1/2}=\frac19(36+9+0+9,\ 18+21+0+6,\ 0+3+0-3,\ 18+6+0+21)=(6,5,0,5)=\row_1(A)\) and
\(\row_1(A^{1/2})A^{-1/2}=\frac1{36}(54-9+0-9,\ -18+21+0-3,\ 0-6+0+6,\ -18-3+0+21)=(1,0,0,0)\), and likewise for rows \(2\),
\(3\) and \(4\); also \(A^{-1}=(A^{-1/2})^2\), as \(\row_1(A^{-1/2})A^{-1/2}=\frac1{144}(81+9+0+9,\ -27-21+0+3,\ 0+6+0-6,\
-27+3+0-21)=\frac1{16}(11,-5,0,-5)=\row_1(A^{-1})\).
\end{solution}

\begin{exercise}\label{r8-ex-idempotent-factor}
Let \(w=(1,-1,-1,1)'\) and
\[
P=I_4-\tfrac14ww'=\frac14\begin{pmatrix}3&1&1&-1\\1&3&-1&1\\1&-1&3&1\\-1&1&1&3\end{pmatrix}
\]
\begin{enumerate}
\item Show that \(P\) is symmetric and idempotent, and find \(\rank P\).
\item Find \(Q_1\) with orthonormal columns and \(P=Q_1Q_1'\), as in Theorem~\ref{r8-thm-idempotent-eigen}(c).
\item For \(y=(2,0,0,2)'\), find the orthogonal projection \(\hat y\) of \(y\) onto \(\mathcal C(P)\), and check that
\(y-\hat y\perp\mathcal C(P)\).
\end{enumerate}
\end{exercise}
\begin{solution}
(a) \(P'=I_4-\frac14(ww')'=P\), and with \(w'w=4\), \(P^2=I_4-\frac12ww'+\frac1{16}w(w'w)w'=I_4-\frac14ww'=P\). By
Theorem~\ref{r6-thm-idempotent}(c), \(\rank P=\tr P=\frac14(3+3+3+3)=3\).

(b) With \(h_1,\dots,h_4\) of Exercise~\ref{r8-ex-diagonalize}, \(Q=\frac12[h_1\ h_2\ h_3\ h_4]\) is orthogonal and
\(\frac12h_4=\frac12w\). By Proposition~\ref{r6-prop-orthonormal}(b) and the outer product expansion
(Theorem~\ref{r1-thm-views}(d)), \(I_4=QQ'=\frac14\sum_jh_jh_j'\), so
\[
P=I_4-\tfrac14h_4h_4'=\tfrac14(h_1h_1'+h_2h_2'+h_3h_3')=Q_1Q_1',\qquad
Q_1=\frac12\begin{pmatrix}1&1&1\\1&1&-1\\1&-1&1\\1&-1&-1\end{pmatrix}
\]
and \(Q_1'Q_1=I_3\).

(c) \(P=P_{\mathcal C(P)}\) (Theorem~\ref{r6-thm-symmetric-idempotent}), so
\(\hat y=Py=y-\frac14(w'y)w=(2,0,0,2)'-(1,-1,-1,1)'=(1,1,1,1)'\). Then \(y-\hat y=w\), and \(Q_1'w=0\), as \(w=h_4\) is
orthogonal to \(h_1,h_2,h_3\). Moreover \(\mathcal C(P)=\mathcal C(Q_1Q_1')\subseteq\mathcal C(Q_1)\)
(Theorem~\ref{thm:r4-rank-product}(a)), and both have dimension \(3\) (by (a), and as the columns of \(Q_1\) are independent,
Proposition~\ref{r6-prop-orthogonal-list}(b)), so \(\mathcal C(P)=\mathcal C(Q_1)\) (Theorem~\ref{thm:r4-counting}(a)). Hence
\(y-\hat y\perp\mathcal C(P)\) (Theorem~\ref{r6-thm-complement}(a)).
\end{solution}

\begin{exercise}\label{r8-dr-classify}
For each quadratic form on \(\R^4\), write its symmetric matrix, complete the squares with Corollary~\ref{r8-cor-ldu}, and
find the inertia (Theorem~\ref{r8-thm-inertia}). Which forms are PD, PSD, indefinite?
\begin{enumerate}
\item \(q_1(x)=x_1^2+3x_2^2+7x_3^2+4x_4^2-2x_1x_2+4x_1x_3+2x_1x_4-2x_2x_4+2x_3x_4\)
\item \(q_2(x)=2x_1^2+x_2^2+x_3^2+4x_1x_2+4x_1x_3-4x_1x_4-2x_2x_3-2x_2x_4+2x_3x_4\)
\item \(q_3(x)=2x_1^2+2x_2^2+2x_3^2+x_4^2+2x_1x_3-2x_1x_4+2x_2x_3-2x_2x_4-2x_3x_4\); for this form, find also all \(x\)
with \(q_3(x)=0\).
\end{enumerate}
\end{exercise}
\begin{solution}
The matrix has the coefficient of \(x_i^2\) as entry \((i,i)\) and half the coefficient of \(x_ix_j\) as entries \((i,j)\) and
\((j,i)\) (Proposition~\ref{r8-prop-qf-symmetric} and the paragraph after it). The \(2\times2\) inverses come from
Theorem~\ref{r5d-thm-adjugate}, and column \(4\) of \(P\) solves the system of Lemma~\ref{r8-lem-triangular} with the
coefficient matrix of the first three rows and columns; then \(D=P'(AP)\) and \(\Delta_4=d_1d_2d_3d_4\)
(Corollary~\ref{r8-cor-ldu}(a)).

(a) The matrix is
\[
A=\begin{pmatrix}1&-1&2&1\\-1&3&0&-1\\2&0&7&1\\1&-1&1&4\end{pmatrix}
\]
\(\Delta_1=1\), \(\Delta_2=3-1=2\), and along row \(1\), \(\Delta_3=1(21-0)+1(-7-0)+2(0-6)=2\). Next \(p_{12}=1\), and with
\(A_2^{-1}=\frac12\left(\begin{smallmatrix}3&1\\1&1\end{smallmatrix}\right)\),
\(\left(\begin{smallmatrix}p_{13}\\p_{23}\end{smallmatrix}\right)=-A_2^{-1}\left(\begin{smallmatrix}2\\0\end{smallmatrix}\right)
=\left(\begin{smallmatrix}-3\\-1\end{smallmatrix}\right)\). Column \(4\) of \(P\) solves
\(A_3(p_{14},p_{24},p_{34})'=-(1,-1,1)'\):
\[
\left(\begin{array}{ccc|c}1&-1&2&-1\\-1&3&0&1\\2&0&7&-1\end{array}\right)
\xrightarrow[R_3\leftarrow R_3-2R_1]{R_2\leftarrow R_2+R_1}
\left(\begin{array}{ccc|c}1&-1&2&-1\\0&2&2&0\\0&2&3&1\end{array}\right)
\xrightarrow{R_3\leftarrow R_3-R_2}
\left(\begin{array}{ccc|c}1&-1&2&-1\\0&2&2&0\\0&0&1&1\end{array}\right)
\]
so \(p_{34}=1\), \(p_{24}=-p_{34}=-1\) and \(p_{14}=-1+p_{24}-2p_{34}=-4\). Hence
\[
P=\begin{pmatrix}1&1&-3&-4\\0&1&-1&-1\\0&0&1&1\\0&0&0&1\end{pmatrix},\qquad
AP=\begin{pmatrix}1&0&0&0\\-1&2&0&0\\2&2&1&0\\1&0&-1&2\end{pmatrix},\qquad
U=P^{-1}=\begin{pmatrix}1&-1&2&1\\0&1&1&0\\0&0&1&-1\\0&0&0&1\end{pmatrix}
\]
and \(D=\diag(1,2,1,2)\), so \(\Delta_4=4\). So \(q_1(x)=(x_1-x_2+2x_3+x_4)^2+2(x_2+x_3)^2+(x_3-x_4)^2+2x_4^2\): the inertia is
\((4,0,0)\), and \(q_1\) is PD.

(b) The matrix is
\[
B=\begin{pmatrix}2&2&2&-2\\2&1&-1&-1\\2&-1&1&1\\-2&-1&1&0\end{pmatrix}
\]
\(\Delta_1=2\), \(\Delta_2=2-4=-2\), and along row \(1\), \(\Delta_3=2(1-1)-2(2+2)+2(-2-2)=-16\). Next \(p_{12}=-1\), and with
\(B_2^{-1}=-\frac12\left(\begin{smallmatrix}1&-2\\-2&2\end{smallmatrix}\right)\),
\(\left(\begin{smallmatrix}p_{13}\\p_{23}\end{smallmatrix}\right)=-B_2^{-1}\left(\begin{smallmatrix}2\\-1\end{smallmatrix}\right)
=\left(\begin{smallmatrix}2\\-3\end{smallmatrix}\right)\). Column \(4\) of \(P\) solves
\(B_3(p_{14},p_{24},p_{34})'=-(-2,-1,1)'\):
\[
\left(\begin{array}{ccc|c}2&2&2&2\\2&1&-1&1\\2&-1&1&-1\end{array}\right)
\xrightarrow[R_3\leftarrow R_3-R_1]{R_2\leftarrow R_2-R_1}
\left(\begin{array}{ccc|c}2&2&2&2\\0&-1&-3&-1\\0&-3&-1&-3\end{array}\right)
\xrightarrow{R_3\leftarrow R_3-3R_2}
\left(\begin{array}{ccc|c}2&2&2&2\\0&-1&-3&-1\\0&0&8&0\end{array}\right)
\]
so \(p_{34}=0\), \(p_{24}=1-3p_{34}=1\) and \(p_{14}=\frac12(2-2p_{24}-2p_{34})=0\). Hence
\[
P=\begin{pmatrix}1&-1&2&0\\0&1&-3&1\\0&0&1&0\\0&0&0&1\end{pmatrix},\qquad
BP=\begin{pmatrix}2&0&0&0\\2&-1&0&0\\2&-3&8&0\\-2&1&0&-1\end{pmatrix},\qquad
U=P^{-1}=\begin{pmatrix}1&1&1&-1\\0&1&3&-1\\0&0&1&0\\0&0&0&1\end{pmatrix}
\]
and \(D=\diag(2,-1,8,-1)\), with \(d_k=\Delta_k/\Delta_{k-1}=2/1,\ -2/2,\ -16/(-2)\) for \(k\leqslant3\), and \(\Delta_4=16\). So
\(q_2(x)=2(x_1+x_2+x_3-x_4)^2-(x_2+3x_3-x_4)^2+8x_3^2-x_4^2\): the inertia is \((2,2,0)\), and \(q_2\) is indefinite; for
instance \(q_2(1,0,0,0)=2\) and \(q_2(-1,1,0,0)=-1\).

(c) The matrix is
\[
F=\begin{pmatrix}2&0&1&-1\\0&2&1&-1\\1&1&2&-1\\-1&-1&-1&1\end{pmatrix}
\]
\(\Delta_1=2\), \(\Delta_2=4\), and along row \(1\), \(\Delta_3=2(4-1)-0+1(0-2)=4\). Next \(p_{12}=0\),
\(\left(\begin{smallmatrix}p_{13}\\p_{23}\end{smallmatrix}\right)=-\frac12I_2\left(\begin{smallmatrix}1\\1\end{smallmatrix}\right)
=\left(\begin{smallmatrix}-1/2\\-1/2\end{smallmatrix}\right)\), and column \(4\) of \(P\) solves
\(F_3(p_{14},p_{24},p_{34})'=-(-1,-1,-1)'\):
\[
\begin{gathered}
\left(\begin{array}{ccc|c}2&0&1&1\\0&2&1&1\\1&1&2&1\end{array}\right)
\xrightarrow{R_1\leftrightarrow R_3}
\left(\begin{array}{ccc|c}1&1&2&1\\0&2&1&1\\2&0&1&1\end{array}\right)
\xrightarrow{R_3\leftarrow R_3-2R_1}
\left(\begin{array}{ccc|c}1&1&2&1\\0&2&1&1\\0&-2&-3&-1\end{array}\right)
\\
\xrightarrow{R_3\leftarrow R_3+R_2}
\left(\begin{array}{ccc|c}1&1&2&1\\0&2&1&1\\0&0&-2&0\end{array}\right)
\end{gathered}
\]
so \(p_{34}=0\), \(p_{24}=\frac12(1-p_{34})=\frac12\) and \(p_{14}=1-p_{24}-2p_{34}=\frac12\). Hence
\[
P=\begin{pmatrix}1&0&-\frac12&\frac12\\0&1&-\frac12&\frac12\\0&0&1&0\\0&0&0&1\end{pmatrix},\qquad
FP=\begin{pmatrix}2&0&0&0\\0&2&0&0\\1&1&1&0\\-1&-1&0&0\end{pmatrix},\qquad
U=P^{-1}=\begin{pmatrix}1&0&\frac12&-\frac12\\0&1&\frac12&-\frac12\\0&0&1&0\\0&0&0&1\end{pmatrix}
\]
and \(D=\diag(2,2,1,0)\), so \(\Delta_4=0\). So \(q_3(x)=2\bigl(x_1+\frac12x_3-\frac12x_4\bigr)^2+2\bigl(x_2+\frac12x_3-\frac12x_4\bigr)^2
+x_3^2\): the inertia is \((3,0,1)\), and \(q_3\) is PSD but not PD. By Corollary~\ref{r8-cor-psd-null}, \(q_3(x)=0\)
exactly for \(x\in\mathcal N(F)\):
\[
\begin{gathered}
F=\begin{pmatrix}2&0&1&-1\\0&2&1&-1\\1&1&2&-1\\-1&-1&-1&1\end{pmatrix}
\xrightarrow[R_3\leftarrow R_3+R_4]{R_1\leftarrow R_1+2R_4}
\begin{pmatrix}0&-2&-1&1\\0&2&1&-1\\0&0&1&0\\-1&-1&-1&1\end{pmatrix}
\xrightarrow[R_4\leftarrow R_4+R_2]{R_1\leftarrow R_1+R_2}
\begin{pmatrix}0&0&0&0\\0&2&1&-1\\0&0&1&0\\-1&1&0&0\end{pmatrix}
\\
\xrightarrow{R_1\leftrightarrow R_4}
\begin{pmatrix}-1&1&0&0\\0&2&1&-1\\0&0&1&0\\0&0&0&0\end{pmatrix}
\end{gathered}
\]
With \(x_4=2\), \(x_3=0\), \(x_2=\frac12(x_4-x_3)=1\) and \(x_1=x_2=1\): \(q_3(x)=0\) if and only if \(x=c(1,1,0,2)'\) for some
\(c\in\R\).
\end{solution}

\begin{exercise}[\citeat{abadir-magnus}{Exercise~8.2}]\label{r8-ex-classify-am}
Write each quadratic form with a symmetric matrix, and decide whether it is positive (negative) definite,
semidefinite, or indefinite: (a) \(x_1^2+x_2^2-x_1x_2\); (b) \(4x_1^2+5x_2^2+x_3^2+2x_1x_2+2x_2x_3\);
(c) \(-x_1^2+2x_2^2+x_1x_2\).
\end{exercise}
\begin{solution}
(a) \(\begin{pmatrix}1&-\frac12\\-\frac12&1\end{pmatrix}\):
\(x_1^2+x_2^2-x_1x_2=(x_1-\frac12x_2)^2+\frac34x_2^2>0\) for \(x\ne0\); positive definite.

(b) The matrix is
\[
\begin{pmatrix}4&1&0\\1&5&1\\0&1&1\end{pmatrix}
\]
and the form is \((x_1+x_2)^2+(x_2+x_3)^2+3x_1^2+3x_2^2\), which is \(0\) only if \(x_1=x_2=0\) and \(x_2+x_3=0\),
that is, \(x=0\); positive definite.

(c) \(\begin{pmatrix}-1&\frac12\\\frac12&2\end{pmatrix}\): \(-x_1^2+2x_2^2+x_1x_2=(\frac32x_2)^2-(x_1-\frac12x_2)^2\)
takes the value \(9\) at \(x=(1,2)'\) and \(-1\) at \(x=(1,0)'\); indefinite.
\end{solution}

\begin{exercise}\label{r8-ex-classify}
Write each quadratic form as \(x'Ax\) with \(A\) symmetric, and decide whether \(A\) is PD, PSD, negative definite,
negative semidefinite or indefinite:
(a) \(x_1^2+x_2^2+x_3^2+x_4^2+6x_1x_2+6x_2x_3+6x_3x_4+6x_1x_4\);
(b) \(x_1^2+2x_2^2+2x_3^2+x_4^2-2x_1x_2-2x_2x_3-2x_3x_4\);
(c) \(-2x_1^2-2x_2^2-2x_3^2-2x_4^2+2x_1x_2+2x_2x_3+2x_3x_4\);
(d) \(2x_1^2+2x_2^2+2x_3^2+x_4^2-2x_1x_2-2x_2x_3-2x_3x_4\).
\end{exercise}
\begin{solution}
(a) The matrix is
\[
A=\begin{pmatrix}1&3&0&3\\3&1&3&0\\0&3&1&3\\3&0&3&1\end{pmatrix}
\]
It gives \(4+24=28\) at \(\one_4\) and \(1+1-6=-4\) at \((1,-1,0,0)'\): indefinite. (The eigenvalues are \(7\), \(1\), \(1\),
\(-5\): \(A\one_4=7\one_4\), \(A(1,0,-1,0)'=(1,0,-1,0)'\), \(A(0,1,0,-1)'=(0,1,0,-1)'\) and
\(A(1,-1,1,-1)'=-5(1,-1,1,-1)'\).)

(b) The matrix is
\[
A=\begin{pmatrix}1&-1&0&0\\-1&2&-1&0\\0&-1&2&-1\\0&0&-1&1\end{pmatrix}
\]
The form equals \((x_1-x_2)^2+(x_2-x_3)^2+(x_3-x_4)^2\geqslant0\), and it is \(0\) at \(x=\one_4\ne0\): PSD but not PD.

(c) The matrix is
\[
A=-\begin{pmatrix}2&-1&0&0\\-1&2&-1&0\\0&-1&2&-1\\0&0&-1&2\end{pmatrix}
\]
Expanding along the last row (Theorem~\ref{r5d-thm-expansion}), the leading principal minors of \(-A\) satisfy
\(\Delta_k=2\Delta_{k-1}-\Delta_{k-2}\) for \(k=3,4\); so they are \(2\), \(4-1=3\), \(2\cdot3-2=4\) and \(2\cdot4-3=5\). They are
positive, so \(-A\) is PD (Theorem~\ref{r8-thm-leading-minors}) and \(A\) is negative definite.

(d) The matrix is
\[
A=\begin{pmatrix}2&-1&0&0\\-1&2&-1&0\\0&-1&2&-1\\0&0&-1&1\end{pmatrix}
\]
the matrix \(K'K\) of Example~\ref{r8-exm-cholesky}. So \(x'Ax=\lVert Kx\rVert^2=x_1^2+(x_2-x_1)^2+(x_3-x_2)^2+(x_4-x_3)^2\),
and \(A\) is PD (Theorem~\ref{r8-thm-XtX}, as \(\rank K=4\)). Its leading principal minors are \(2,3,4,1\).
\end{solution}

\begin{exercise}\label{r8-dr-psd-factor}
Let
\[
A=\begin{pmatrix}2&2&-1&1\\2&2&-1&1\\-1&-1&5&1\\1&1&1&1\end{pmatrix}
\]
\begin{enumerate}
\item Find \(p_A\), the eigenvalues, and an orthonormal basis of eigenvectors.
\item Show that \(A\) is PSD, find \(\rank A\), and find the matrix \(P\) of Theorem~\ref{r8-thm-psd-factor}(a) with \(A=P'P\).
\item Find all \(x\) with \(x'Ax=0\).
\end{enumerate}
\end{exercise}
\begin{solution}
(a) For real \(t\), apply \(R_1\leftarrow R_1-R_2\), take the factor \(t\) out of row \(1\), add column \(1\) to column \(2\),
and expand along row \(1\) and then along row \(1\) of the \(3\times3\) matrix (Definition~\ref{r5d-def-det}(a),(b),
Theorems~\ref{r5d-thm-transpose} and~\ref{r5d-thm-expansion}):
\[
\begin{aligned}
\det(tI-A)&=\det\begin{pmatrix}t-2&-2&1&-1\\-2&t-2&1&-1\\1&1&t-5&-1\\-1&-1&-1&t-1\end{pmatrix}\\
&=\det\begin{pmatrix}t&-t&0&0\\-2&t-2&1&-1\\1&1&t-5&-1\\-1&-1&-1&t-1\end{pmatrix}
=t\det\begin{pmatrix}1&-1&0&0\\-2&t-2&1&-1\\1&1&t-5&-1\\-1&-1&-1&t-1\end{pmatrix}\\
&=t\det\begin{pmatrix}1&0&0&0\\-2&t-4&1&-1\\1&2&t-5&-1\\-1&-2&-1&t-1\end{pmatrix}
=t\det\begin{pmatrix}t-4&1&-1\\2&t-5&-1\\-2&-1&t-1\end{pmatrix}\\
&=t\Bigl((t-4)\bigl((t-5)(t-1)-1\bigr)-\bigl(2(t-1)-2\bigr)-\bigl(-2+2(t-5)\bigr)\Bigr)\\
&=t(t^3-10t^2+24t)
\end{aligned}
\]
So \(p_A(t)=t^2(t-4)(t-6)\) (Definition~\ref{r7-def-charpoly}(a), Proposition~\ref{r7-prop-poly}(d)), and the eigenvalues
are \(6\), \(4\) and \(0\), with \(a_A(0)=2\). Row-reduce \(A-cI\):
\[
\begin{gathered}
A-6I=\begin{pmatrix}-4&2&-1&1\\2&-4&-1&1\\-1&-1&-1&1\\1&1&1&-5\end{pmatrix}
\xrightarrow{R_1\leftrightarrow R_4}
\begin{pmatrix}1&1&1&-5\\2&-4&-1&1\\-1&-1&-1&1\\-4&2&-1&1\end{pmatrix}
\\
\xrightarrow[R_4\leftarrow R_4+4R_1]{R_2\leftarrow R_2-2R_1,\ R_3\leftarrow R_3+R_1}
\begin{pmatrix}1&1&1&-5\\0&-6&-3&11\\0&0&0&-4\\0&6&3&-19\end{pmatrix}
\xrightarrow{R_4\leftarrow R_4+R_2}
\begin{pmatrix}1&1&1&-5\\0&-6&-3&11\\0&0&0&-4\\0&0&0&-8\end{pmatrix}
\\
\xrightarrow{R_4\leftarrow R_4-2R_3}
\begin{pmatrix}1&1&1&-5\\0&-6&-3&11\\0&0&0&-4\\0&0&0&0\end{pmatrix}
\end{gathered}
\]
so \(x_4=0\), and \(x_2=1\) gives \(x_3=-2x_2=-2\) and \(x_1=-x_2-x_3=1\);
\[
\begin{gathered}
A-4I=\begin{pmatrix}-2&2&-1&1\\2&-2&-1&1\\-1&-1&1&1\\1&1&1&-3\end{pmatrix}
\xrightarrow{R_1\leftrightarrow R_4}
\begin{pmatrix}1&1&1&-3\\2&-2&-1&1\\-1&-1&1&1\\-2&2&-1&1\end{pmatrix}
\\
\xrightarrow[R_4\leftarrow R_4+2R_1]{R_2\leftarrow R_2-2R_1,\ R_3\leftarrow R_3+R_1}
\begin{pmatrix}1&1&1&-3\\0&-4&-3&7\\0&0&2&-2\\0&4&1&-5\end{pmatrix}
\xrightarrow{R_4\leftarrow R_4+R_2}
\begin{pmatrix}1&1&1&-3\\0&-4&-3&7\\0&0&2&-2\\0&0&-2&2\end{pmatrix}
\\
\xrightarrow{R_4\leftarrow R_4+R_3}
\begin{pmatrix}1&1&1&-3\\0&-4&-3&7\\0&0&2&-2\\0&0&0&0\end{pmatrix}
\end{gathered}
\]
so \(x_4=1\) gives \(x_3=x_4=1\), \(x_2=\frac14(-3x_3+7x_4)=1\) and \(x_1=-x_2-x_3+3x_4=1\);
\[
\begin{gathered}
A=\begin{pmatrix}2&2&-1&1\\2&2&-1&1\\-1&-1&5&1\\1&1&1&1\end{pmatrix}
\xrightarrow{R_1\leftrightarrow R_4}
\begin{pmatrix}1&1&1&1\\2&2&-1&1\\-1&-1&5&1\\2&2&-1&1\end{pmatrix}
\xrightarrow[R_4\leftarrow R_4-2R_1]{R_2\leftarrow R_2-2R_1,\ R_3\leftarrow R_3+R_1}
\begin{pmatrix}1&1&1&1\\0&0&-3&-1\\0&0&6&2\\0&0&-3&-1\end{pmatrix}
\\
\xrightarrow[R_4\leftarrow R_4-R_2]{R_3\leftarrow R_3+2R_2}
\begin{pmatrix}1&1&1&1\\0&0&-3&-1\\0&0&0&0\\0&0&0&0\end{pmatrix}
\end{gathered}
\]
so \(\mathcal N(A)=\{x:x_3=-\frac13x_4,\ x_1=-x_2-x_3-x_4\}\) has the basis \(s_1=(-1,1,0,0)'\) (\(x_2=1\), \(x_4=0\)) and
\(s_2=(-2,0,-1,3)'\) (\(x_2=0\), \(x_4=3\)). So \(\mathcal N(A-6I)\) and \(\mathcal N(A-4I)\) are spanned by \((1,1,-2,0)'\) and
\((1,1,1,1)'\). In \(\mathcal N(A)\), Gram--Schmidt (Theorem~\ref{r6-thm-gram-schmidt}) gives \(b_1=s_1\), with
\(\lVert b_1\rVert^2=2\), and, as \(s_2'b_1=2\), \(b_2=s_2-b_1=(-1,-1,-1,3)'\), with \(\lVert b_2\rVert^2=12\). The eigenspaces are
pairwise orthogonal (Theorem~\ref{r8-thm-sym-orth}), so an orthonormal basis of eigenvectors is
\(q_1=\frac1{\sqrt6}(1,1,-2,0)'\), \(q_2=\frac12(1,1,1,1)'\), \(q_3=\frac1{\sqrt2}(-1,1,0,0)'\), \(q_4=\frac1{\sqrt{12}}(-1,-1,-1,3)'\).

(b) The eigenvalues are \(\geqslant0\), so \(A\) is PSD (Theorem~\ref{r8-thm-eigen-criterion}(a)), and \(\rank A=2\), the number
of nonzero eigenvalues (Theorem~\ref{r8-thm-rank-eigen}). The leading minors \(2\), \(0\), \(0\), \(0\) do not decide here, as
\(\Delta_2=0\). The proof of Theorem~\ref{r8-thm-psd-factor}(a) takes the rows \(\sqrt6\,q_1'=(1,1,-2,0)\) and
\(\sqrt4\,q_2'=(1,1,1,1)\):
\[
P=\begin{pmatrix}1&1&-2&0\\1&1&1&1\end{pmatrix},\qquad
P'P=\begin{pmatrix}1&1\\1&1\\-2&1\\0&1\end{pmatrix}\begin{pmatrix}1&1&-2&0\\1&1&1&1\end{pmatrix}
=\begin{pmatrix}2&2&-1&1\\2&2&-1&1\\-1&-1&5&1\\1&1&1&1\end{pmatrix}=A
\]
(c) By Corollary~\ref{r8-cor-psd-null}, \(x'Ax=0\) if and only if \(Ax=0\), and by the last reduction in (a) this means
\(x=c_1(-1,1,0,0)'+c_2(-2,0,-1,3)'\) for some \(c_1,c_2\in\R\). Equivalently, \(x'Ax=\lVert Px\rVert^2=(x_1+x_2-2x_3)^2
+(x_1+x_2+x_3+x_4)^2\) vanishes exactly when \(x_1+x_2-2x_3=0\) and \(x_1+x_2+x_3+x_4=0\).
\end{solution}

\begin{exercise}\label{r8-ex-leading-drill}
Let
\[
A=\begin{pmatrix}4&2&1&-3\\2&4&-1&-3\\1&-1&3&0\\-3&-3&0&9\end{pmatrix}
\]
\begin{enumerate}
\item Find the leading principal minors of \(A\), and decide whether \(A\) is PD.
\item Find \(P\), \(D\) and \(U\) of Theorem~\ref{r8-thm-leading-minors}, and write \(q_A(x)\) as a sum of squares.
\item Check (a) with Theorem~\ref{r8-thm-schur-pd}, using the Schur complement of \(A_2\).
\item Show that \((1,1,0,-2)'\), \((1,-1,2,0)'\), \((1,1,0,1)'\) and \((1,-1,-1,0)'\) are eigenvectors of \(A\), and check (a) once
more.
\end{enumerate}
\end{exercise}
\begin{solution}
(a) \(\Delta_1=4\), \(\Delta_2=16-4=12\), and along row \(1\), \(\Delta_3=4(12-1)-2(6+1)+1(-2-4)=44-14-6=24\). For
\(\Delta_4=\det A\), interchange rows \(1\) and \(3\), apply \(R_2\leftarrow R_2-2R_1\), \(R_3\leftarrow R_3-4R_1\),
\(R_4\leftarrow R_4+3R_1\), then \(R_3\leftarrow R_3-R_2\), \(R_4\leftarrow R_4+R_2\), then \(R_4\leftarrow R_4+\frac12R_3\)
(Definition~\ref{r5d-def-det}(b),(c), Theorem~\ref{r5d-thm-basic}(c)):
\[
\begin{aligned}
\det A&=-\det\begin{pmatrix}1&-1&3&0\\2&4&-1&-3\\4&2&1&-3\\-3&-3&0&9\end{pmatrix}
=-\det\begin{pmatrix}1&-1&3&0\\0&6&-7&-3\\0&6&-11&-3\\0&-6&9&9\end{pmatrix}\\
&=-\det\begin{pmatrix}1&-1&3&0\\0&6&-7&-3\\0&0&-4&0\\0&0&2&6\end{pmatrix}
=-\det\begin{pmatrix}1&-1&3&0\\0&6&-7&-3\\0&0&-4&0\\0&0&0&6\end{pmatrix}=-1\cdot6\cdot(-4)\cdot6=144
\end{aligned}
\]
All four minors are \(>0\), so \(A\) is PD (Theorem~\ref{r8-thm-leading-minors}).

(b) For \(k=2\): \(4p_{12}=-2\), so \(p_{12}=-\frac12\). For \(k=3\): with
\(A_2^{-1}=\frac1{12}\left(\begin{smallmatrix}4&-2\\-2&4\end{smallmatrix}\right)\) (Theorem~\ref{r5d-thm-adjugate}),
\(\left(\begin{smallmatrix}p_{13}\\p_{23}\end{smallmatrix}\right)=-A_2^{-1}\left(\begin{smallmatrix}1\\-1\end{smallmatrix}\right)
=-\frac1{12}\left(\begin{smallmatrix}6\\-6\end{smallmatrix}\right)=\left(\begin{smallmatrix}-1/2\\1/2\end{smallmatrix}\right)\). For
\(k=4\), \(A_3(p_{14},p_{24},p_{34})'=-(-3,-3,0)'\):
\[
\begin{gathered}
\left(\begin{array}{ccc|c}4&2&1&3\\2&4&-1&3\\1&-1&3&0\end{array}\right)
\xrightarrow{R_1\leftrightarrow R_3}
\left(\begin{array}{ccc|c}1&-1&3&0\\2&4&-1&3\\4&2&1&3\end{array}\right)
\xrightarrow[R_3\leftarrow R_3-4R_1]{R_2\leftarrow R_2-2R_1}
\left(\begin{array}{ccc|c}1&-1&3&0\\0&6&-7&3\\0&6&-11&3\end{array}\right)
\\
\xrightarrow{R_3\leftarrow R_3-R_2}
\left(\begin{array}{ccc|c}1&-1&3&0\\0&6&-7&3\\0&0&-4&0\end{array}\right)
\end{gathered}
\]
so \(p_{34}=0\), \(p_{24}=\frac16(3+7p_{34})=\frac12\) and \(p_{14}=p_{24}-3p_{34}=\frac12\). So
\[
\begin{gathered}
P=\begin{pmatrix}1&-\frac12&-\frac12&\frac12\\0&1&\frac12&\frac12\\0&0&1&0\\0&0&0&1\end{pmatrix},\qquad
AP=\begin{pmatrix}4&0&0&0\\2&3&0&0\\1&-\frac32&2&0\\-3&-\frac32&0&6\end{pmatrix}\\
D=\begin{pmatrix}4&0&0&0\\0&3&0&0\\0&0&2&0\\0&0&0&6\end{pmatrix},\qquad
U=P^{-1}=\begin{pmatrix}1&\frac12&\frac14&-\frac34\\0&1&-\frac12&-\frac12\\0&0&1&0\\0&0&0&1\end{pmatrix}
\end{gathered}
\]
with \(d_k=\Delta_k/\Delta_{k-1}=4/1,\ 12/4,\ 24/12,\ 144/24\). Hence
\(q_A(x)=(Ux)'D(Ux)=4(x_1+\frac12x_2+\frac14x_3-\frac34x_4)^2+3(x_2-\frac12x_3-\frac12x_4)^2+2x_3^2+6x_4^2\).

(c) \(A_2\) is PD (its leading minors \(4\), \(12\) are positive; Theorem~\ref{r8-thm-leading-minors}). With
\(B=\left(\begin{smallmatrix}1&-3\\-1&-3\end{smallmatrix}\right)\), the upper right \(2\times2\) block of \(A\),
\(A_2^{-1}B=\frac1{12}\left(\begin{smallmatrix}6&-6\\-6&-6\end{smallmatrix}\right)\) and
\(B'A_2^{-1}B=\left(\begin{smallmatrix}1&0\\0&3\end{smallmatrix}\right)\), so
\(A/A_2=\left(\begin{smallmatrix}3&0\\0&9\end{smallmatrix}\right)-B'A_2^{-1}B=\left(\begin{smallmatrix}2&0\\0&6\end{smallmatrix}\right)\),
which is PD (Proposition~\ref{r8-prop-pd-basic}(a)); so \(A\) is PD. Also \(\det A=12\cdot\det(A/A_2)=12\cdot12=144\).

(d) \(A(1,1,0,-2)'=(12,12,0,-24)'\), \(A(1,-1,2,0)'=(4,-4,8,0)'\), \(A(1,1,0,1)'=(3,3,0,3)'\) and \(A(1,-1,-1,0)'=(1,-1,-1,0)'\).
So \(a_A(c)\geqslant1\) for \(c=12,4,3,1\) (Theorem~\ref{r8-thm-rank-eigen}), and as the multiplicities add up to \(4\), the
eigenvalues are \(12\), \(4\), \(3\), \(1\), all \(>0\) (Theorem~\ref{r8-thm-eigen-criterion}(a)); also
\(12\cdot4\cdot3\cdot1=144=\det A\).
\end{solution}

\begin{exercise}\label{r8-ex-sylvester-parameter}
For \(a\in\R\) let
\[
A=\begin{pmatrix}4&2&2&a\\2&4&a&2\\2&a&4&2\\a&2&2&4\end{pmatrix}
\]
and let \(h_1,\dots,h_4\) be the vectors of Exercise~\ref{r8-ex-diagonalize}.
\begin{enumerate}
\item Show that \(h_1,\dots,h_4\) are eigenvectors of \(A\), and find \(\det A\).
\item For which \(a\) is \(A\) PD? Answer with Theorem~\ref{r8-thm-leading-minors}, and check with (a).
\item For which \(a\) is \(A\) PSD? Find \(\mathcal N(A)\) at the endpoints.
\end{enumerate}
\end{exercise}
\begin{solution}
(a) The row sums are \(8+a\), so \(Ah_1=(8+a)h_1\). Multiplying the rows of \(A\) by \(h_2,h_3\),
\[
Ah_2=\begin{pmatrix}4+2-2-a\\2+4-a-2\\2+a-4-2\\a+2-2-4\end{pmatrix}=(4-a)h_2,\quad
Ah_3=\begin{pmatrix}4-2+2-a\\2-4+a-2\\2-a+4-2\\a-2+2-4\end{pmatrix}=(4-a)h_3
\]
In the same way \(Ah_4=(4-2-2+a,\ 2-4-a+2,\ 2-a-4+2,\ a-2-2+4)'=ah_4\). So \(A=Q\Lambda Q'\) with
\(Q=\frac12[h_1\ h_2\ h_3\ h_4]\) and \(\Lambda=\diag(8+a,4-a,4-a,a)\), these are the eigenvalues of \(A\), and
\(\det A=a(a+8)(a-4)^2\) (Theorems~\ref{r8-thm-spectral}(b),(c) and~\ref{r7-thm-trace-det}).

(b) The leading principal minors are \(4\), \(16-4=12\),
\[
\Delta_3=4(16-a^2)-2(8-2a)+2(2a-8)=32+8a-4a^2=-4(a-4)(a+2)
\]
(expanding \(\det A_3\) along the first row), and \(\det A=a(a+8)(a-4)^2\). The third is positive exactly for \(-2<a<4\), and
the fourth exactly for \(a\ne4\) and \(a(a+8)>0\), that is, \(a>0\) or \(a<-8\). By Theorem~\ref{r8-thm-leading-minors}, \(A\)
is PD if and only if \(0<a<4\). By (a) and Theorem~\ref{r8-thm-eigen-criterion}(a), \(A\) is PD if and only if \(8+a>0\),
\(4-a>0\) and \(a>0\): the same condition.

(c) By (a) and Theorem~\ref{r8-thm-eigen-criterion}(a), \(A\) is PSD if and only if \(8+a\geqslant0\), \(4-a\geqslant0\) and
\(a\geqslant0\), that is, \(0\leqslant a\leqslant4\). By Theorem~\ref{r8-thm-spectral-decomp}(c),
\(\mathcal N(A)=\Span\{h_4\}\) for \(a=0\) (eigenvalues \(8,4,4,0\)) and \(\mathcal N(A)=\Span\{h_2,h_3\}\) for \(a=4\) (eigenvalues
\(12,0,0,4\)).
\end{solution}

\begin{exercise}\label{r8-dr-cholesky}
Let
\[
A=\begin{pmatrix}1&1&1&1\\1&2&0&2\\1&0&6&4\\1&2&4&10\end{pmatrix}
\]
\begin{enumerate}
\item Show that \(A\) is PD by its leading principal minors.
\item Find \(P\), \(D\) and \(U\) of Theorem~\ref{r8-thm-leading-minors}.
\item Find the Cholesky factor \(T\) of \(A\), check \(T'T=A\), and write \(q_A(x)\) as a sum of squares.
\end{enumerate}
\end{exercise}
\begin{solution}
(a) \(\Delta_1=1\), \(\Delta_2=2-1=1\), and along row \(1\), \(\Delta_3=1(12-0)-1(6-0)+1(0-2)=4\). For \(\Delta_4=\det A\), apply
\(R_2\leftarrow R_2-R_1\), \(R_3\leftarrow R_3-R_1\), \(R_4\leftarrow R_4-R_1\), then \(R_3\leftarrow R_3+R_2\),
\(R_4\leftarrow R_4-R_2\), then \(R_4\leftarrow R_4-R_3\) (Definition~\ref{r5d-def-det}(b), Theorem~\ref{r5d-thm-basic}(c)):
\[
\det A=\det\begin{pmatrix}1&1&1&1\\0&1&-1&1\\0&-1&5&3\\0&1&3&9\end{pmatrix}
=\det\begin{pmatrix}1&1&1&1\\0&1&-1&1\\0&0&4&4\\0&0&4&8\end{pmatrix}
=\det\begin{pmatrix}1&1&1&1\\0&1&-1&1\\0&0&4&4\\0&0&0&4\end{pmatrix}=16
\]
All four minors are \(>0\), so \(A\) is PD (Theorem~\ref{r8-thm-leading-minors}).

(b) For \(k=2\), \(p_{12}=-1\). For \(k=3\), with \(A_2^{-1}=\left(\begin{smallmatrix}2&-1\\-1&1\end{smallmatrix}\right)\)
(\(\det A_2=1\), Theorem~\ref{r5d-thm-adjugate}),
\(\left(\begin{smallmatrix}p_{13}\\p_{23}\end{smallmatrix}\right)=-A_2^{-1}\left(\begin{smallmatrix}1\\0\end{smallmatrix}\right)
=\left(\begin{smallmatrix}-2\\1\end{smallmatrix}\right)\). For \(k=4\), \(A_3(p_{14},p_{24},p_{34})'=-(1,2,4)'\):
\[
\left(\begin{array}{ccc|c}1&1&1&-1\\1&2&0&-2\\1&0&6&-4\end{array}\right)
\xrightarrow[R_3\leftarrow R_3-R_1]{R_2\leftarrow R_2-R_1}
\left(\begin{array}{ccc|c}1&1&1&-1\\0&1&-1&-1\\0&-1&5&-3\end{array}\right)
\xrightarrow{R_3\leftarrow R_3+R_2}
\left(\begin{array}{ccc|c}1&1&1&-1\\0&1&-1&-1\\0&0&4&-4\end{array}\right)
\]
so \(p_{34}=-1\), \(p_{24}=-1+p_{34}=-2\) and \(p_{14}=-1-p_{24}-p_{34}=2\). So
\[
\begin{gathered}
P=\begin{pmatrix}1&-1&-2&2\\0&1&1&-2\\0&0&1&-1\\0&0&0&1\end{pmatrix},\qquad
AP=\begin{pmatrix}1&0&0&0\\1&1&0&0\\1&-1&4&0\\1&1&4&4\end{pmatrix}\\
D=\begin{pmatrix}1&0&0&0\\0&1&0&0\\0&0&4&0\\0&0&0&4\end{pmatrix},\qquad
U=P^{-1}=\begin{pmatrix}1&1&1&1\\0&1&-1&1\\0&0&1&1\\0&0&0&1\end{pmatrix}
\end{gathered}
\]
with \(d_k=\Delta_k/\Delta_{k-1}=1/1,\ 1/1,\ 4/1,\ 16/4\).

(c) By Theorem~\ref{r8-thm-cholesky}, \(T=D^{1/2}U\) is \(U\) with rows \(3\) and \(4\) doubled:
\[
T=\begin{pmatrix}1&1&1&1\\0&1&-1&1\\0&0&2&2\\0&0&0&2\end{pmatrix},\qquad
T'T=\begin{pmatrix}1&0&0&0\\1&1&0&0\\1&-1&2&0\\1&1&2&2\end{pmatrix}\begin{pmatrix}1&1&1&1\\0&1&-1&1\\0&0&2&2\\0&0&0&2\end{pmatrix}
=\begin{pmatrix}1&1&1&1\\1&2&0&2\\1&0&6&4\\1&2&4&10\end{pmatrix}=A
\]
So \(q_A(x)=\lVert Tx\rVert^2=(x_1+x_2+x_3+x_4)^2+(x_2-x_3+x_4)^2+4(x_3+x_4)^2+4x_4^2\). Also
\(\det A=(\det T)^2=(1\cdot1\cdot2\cdot2)^2=16=\Delta_4\) (Theorem~\ref{r8-thm-psd-factor}(c)).
\end{solution}

\begin{exercise}\label{r8-ex-cholesky}
Let
\[
A=\begin{pmatrix}9&3&-6&3\\3&5&0&-1\\-6&0&9&-1\\3&-1&-1&4\end{pmatrix}
\]
(a) Find the Cholesky factor \(T\) of \(A\). (b) Find \(\det A\). (c) Solve \(Ax=b\) for \(b=(9,7,2,5)'\) by solving
\(T'y=b\) and then \(Tx=y\).
\end{exercise}
\begin{solution}
(a) By Remark~\ref{r8-rem-cholesky}, row by row:
\begin{itemize}
\item row \(1\): \(t_{11}=\sqrt9=3\), \(t_{12}=\frac33=1\), \(t_{13}=-\frac63=-2\), \(t_{14}=\frac33=1\);
\item row \(2\): \(t_{22}=\sqrt{5-1}=2\), \(t_{23}=\bigl(0-1\cdot(-2)\bigr)/2=1\), \(t_{24}=(-1-1\cdot1)/2=-1\);
\item row \(3\): \(t_{33}=\sqrt{9-4-1}=2\), \(t_{34}=\bigl(-1-(-2)\cdot1-1\cdot(-1)\bigr)/2=1\);
\item row \(4\): \(t_{44}=\sqrt{4-1-1-1}=1\).
\end{itemize}
So
\[
T=\begin{pmatrix}3&1&-2&1\\0&2&1&-1\\0&0&2&1\\0&0&0&1\end{pmatrix}
\]
and \(T'T=A\); in particular \(A\) is PD (Theorem~\ref{r8-thm-psd-factor}(c)).

(b) \(\det A=\det T'\cdot\det T=(3\cdot2\cdot2\cdot1)^2=144\) (Theorems~\ref{r5d-thm-product}(b), \ref{r5d-thm-transpose}
and~\ref{r5d-thm-basic}(c)).

(c) \(Ax=T'(Tx)\), so \(x\) solves \(Ax=b\) if \(T'y=b\) and \(Tx=y\). From the first equation downwards, \(3y_1=9\),
\(y_1+2y_2=7\), \(-2y_1+y_2+2y_3=2\), \(y_1-y_2+y_3+y_4=5\) give \(y=(3,2,3,1)'\). Back substitution in \(Tx=y\): \(x_4=1\),
\(2x_3+x_4=3\), \(2x_2+x_3-x_4=2\), \(3x_1+x_2-2x_3+x_4=3\) give \(x=(1,1,1,1)'\). Check: the row sums of \(A\) are
\(9,7,2,5\). The solution is unique, as \(A\) is invertible.
\end{solution}

\begin{exercise}\label{r8-dr-inertia}
Let
\[
A=\begin{pmatrix}1&0&2&-1\\0&1&2&-1\\2&2&-1&-1\\-1&-1&-1&5\end{pmatrix}
\]
\begin{enumerate}
\item Find the leading principal minors and the matrices \(P\), \(D\), \(U\) with \(P'AP=D\) and \(A=U'DU\)
(Corollary~\ref{r8-cor-ldu}), and write \(q_A(x)\) as a combination of squares.
\item Find the inertia of \(A\). Is \(A\) PD, PSD or indefinite?
\item Check (b) with the eigenvectors \((1,1,1,-3)'\), \((1,1,1,1)'\), \((1,-1,0,0)'\) and \((1,1,-2,0)'\) of \(A\).
\item Find an invertible \(R\) with \(R'AR=\diag(1,1,1,-1)\).
\end{enumerate}
\end{exercise}
\begin{solution}
(a) \(\Delta_1=1\), \(\Delta_2=1\), and along row \(1\), \(\Delta_3=1(-1-4)-0+2(0-2)=-9\). Next \(p_{12}=0\), and as \(A_2=I_2\),
\(\left(\begin{smallmatrix}p_{13}\\p_{23}\end{smallmatrix}\right)=-\left(\begin{smallmatrix}2\\2\end{smallmatrix}\right)\). Column
\(4\) of \(P\) solves \(A_3(p_{14},p_{24},p_{34})'=-(-1,-1,-1)'\) (Lemma~\ref{r8-lem-triangular}):
\[
\left(\begin{array}{ccc|c}1&0&2&1\\0&1&2&1\\2&2&-1&1\end{array}\right)
\xrightarrow{R_3\leftarrow R_3-2R_1}
\left(\begin{array}{ccc|c}1&0&2&1\\0&1&2&1\\0&2&-5&-1\end{array}\right)
\xrightarrow{R_3\leftarrow R_3-2R_2}
\left(\begin{array}{ccc|c}1&0&2&1\\0&1&2&1\\0&0&-9&-3\end{array}\right)
\]
so \(p_{34}=\frac13\), \(p_{24}=1-2p_{34}=\frac13\) and \(p_{14}=1-2p_{34}=\frac13\). So
\[
P=\begin{pmatrix}1&0&-2&\frac13\\0&1&-2&\frac13\\0&0&1&\frac13\\0&0&0&1\end{pmatrix},\qquad
AP=\begin{pmatrix}1&0&0&0\\0&1&0&0\\2&2&-9&0\\-1&-1&3&4\end{pmatrix},\qquad
U=P^{-1}=\begin{pmatrix}1&0&2&-1\\0&1&2&-1\\0&0&1&-\frac13\\0&0&0&1\end{pmatrix}
\]
and \(D=P'AP=\diag(1,1,-9,4)\), so \(\Delta_4=d_1d_2d_3d_4=-36\) (Corollary~\ref{r8-cor-ldu}(a)). So
\(q_A(x)=(x_1+2x_3-x_4)^2+(x_2+2x_3-x_4)^2-9\bigl(x_3-\frac13x_4\bigr)^2+4x_4^2\).

(b) Three diagonal entries of \(D\) are positive and one is negative: the inertia is \((3,1,0)\)
(Theorem~\ref{r8-thm-inertia}), and \(A\) is indefinite; for instance \(q_A(e_1)=1\) and \(q_A(e_3)=-1\).

(c) \(A(1,1,1,-3)'=(6,6,6,-18)'\), \(A(1,1,1,1)'=(2,2,2,2)'\), \(A(1,-1,0,0)'=(1,-1,0,0)'\) and \(A(1,1,-2,0)'=(-3,-3,6,0)'\).
So \(a_A(c)\geqslant1\) for \(c=6,2,1,-3\) (Theorem~\ref{r8-thm-rank-eigen}); as the multiplicities add up to \(4\), the
eigenvalues are \(6\), \(2\), \(1\), \(-3\): three positive and one negative, as Theorem~\ref{r8-thm-inertia}(c) requires.

(d) Let \(p_1,\dots,p_4\) be the columns of \(P\), so that \(p_i'Ap_j=0\) for \(i\ne j\) and \(p_i'Ap_i=d_i\)
(Theorem~\ref{r1-thm-views}(d)). Scale \(p_3\) by \(\frac13\) and \(p_4\) by \(\frac12\), and interchange them:
\(R=[p_1\ p_2\ \frac12p_4\ \frac13p_3]\). The columns of \(R\) are nonzero multiples of those of the invertible \(P\), in
another order, so they are linearly independent and \(R\) is invertible (Theorem~\ref{thm:r4-rank-rules}(b),(d)). Then
\[
R=\begin{pmatrix}1&0&\frac16&-\frac23\\0&1&\frac16&-\frac23\\0&0&\frac16&\frac13\\0&0&\frac12&0\end{pmatrix},\qquad
R'AR=\diag\bigl(1,1,\tfrac14\cdot4,\tfrac19\cdot(-9)\bigr)=\diag(1,1,1,-1)
\]
\end{solution}

\begin{exercise}\label{r8-ex-sqrt-drill}
Let
\[
A=\begin{pmatrix}6&-3&-3&0\\-3&7&-4&5\\-3&-4&7&-5\\0&5&-5&6\end{pmatrix}
\]
\begin{enumerate}
\item Find \(p_A\), the eigenvalues, and an orthonormal basis of eigenvectors.
\item Is \(A\) PD? PSD? Find \(\rank A\) and the leading principal minors.
\item Find \(A^{1/2}\).
\end{enumerate}
\end{exercise}
\begin{solution}
(a) For real \(t\), add rows \(2\) and \(3\) to row \(1\), take the factor \(t\) out of row \(1\), subtract column \(1\) from
columns \(2\) and \(3\), and expand along row \(1\); in the \(3\times3\) determinant, add row \(2\) to row \(1\), take the
factor \(t-9\) out of row \(1\), subtract column \(1\) from column \(2\), and expand along row \(1\)
(Definition~\ref{r5d-def-det}(a),(b), Theorems~\ref{r5d-thm-transpose} and~\ref{r5d-thm-expansion}):
\[
\begin{aligned}
\det(tI-A)&=\det\begin{pmatrix}t-6&3&3&0\\3&t-7&4&-5\\3&4&t-7&5\\0&-5&5&t-6\end{pmatrix}
=\det\begin{pmatrix}t&t&t&0\\3&t-7&4&-5\\3&4&t-7&5\\0&-5&5&t-6\end{pmatrix}\\
&=t\det\begin{pmatrix}1&1&1&0\\3&t-7&4&-5\\3&4&t-7&5\\0&-5&5&t-6\end{pmatrix}
=t\det\begin{pmatrix}1&0&0&0\\3&t-10&1&-5\\3&1&t-10&5\\0&-5&5&t-6\end{pmatrix}\\
&=t\det\begin{pmatrix}t-10&1&-5\\1&t-10&5\\-5&5&t-6\end{pmatrix}
=t\det\begin{pmatrix}t-9&t-9&0\\1&t-10&5\\-5&5&t-6\end{pmatrix}\\
&=t(t-9)\det\begin{pmatrix}1&1&0\\1&t-10&5\\-5&5&t-6\end{pmatrix}
=t(t-9)\det\begin{pmatrix}1&0&0\\1&t-11&5\\-5&10&t-6\end{pmatrix}\\
&=t(t-9)\bigl((t-11)(t-6)-50\bigr)=t(t-9)(t^2-17t+16)=t(t-1)(t-9)(t-16)
\end{aligned}
\]
So \(p_A(t)=t(t-1)(t-9)(t-16)\) (Definition~\ref{r7-def-charpoly}(a), Proposition~\ref{r7-prop-poly}(d)), and the
eigenvalues are \(16\), \(9\), \(1\), \(0\). Row-reduce \(A-cI\):
\[
\begin{gathered}
A-16I=\begin{pmatrix}-10&-3&-3&0\\-3&-9&-4&5\\-3&-4&-9&-5\\0&5&-5&-10\end{pmatrix}
\xrightarrow[R_2\leftarrow R_2-R_3]{R_1\leftarrow R_1-R_4}
\begin{pmatrix}-10&-8&2&10\\0&-5&5&10\\-3&-4&-9&-5\\0&5&-5&-10\end{pmatrix}
\\
\xrightarrow[R_4\leftarrow R_4+R_2]{R_1\leftarrow R_1+2R_3}
\begin{pmatrix}-16&-16&-16&0\\0&-5&5&10\\-3&-4&-9&-5\\0&0&0&0\end{pmatrix}
\xrightarrow[R_2\leftarrow\frac15R_2]{R_1\leftarrow\frac1{16}R_1}
\begin{pmatrix}-1&-1&-1&0\\0&-1&1&2\\-3&-4&-9&-5\\0&0&0&0\end{pmatrix}
\\
\xrightarrow{R_3\leftarrow R_3-3R_1}
\begin{pmatrix}-1&-1&-1&0\\0&-1&1&2\\0&-1&-6&-5\\0&0&0&0\end{pmatrix}
\xrightarrow{R_3\leftarrow R_3-R_2}
\begin{pmatrix}-1&-1&-1&0\\0&-1&1&2\\0&0&-7&-7\\0&0&0&0\end{pmatrix}
\end{gathered}
\]
so \(x_4=1\) gives \(x_3=-x_4=-1\), \(x_2=x_3+2x_4=1\) and \(x_1=-x_2-x_3=0\);
\[
\begin{gathered}
A-9I=\begin{pmatrix}-3&-3&-3&0\\-3&-2&-4&5\\-3&-4&-2&-5\\0&5&-5&-3\end{pmatrix}
\xrightarrow[R_3\leftarrow R_3-R_1]{R_2\leftarrow R_2-R_4}
\begin{pmatrix}-3&-3&-3&0\\-3&-7&1&8\\0&-1&1&-5\\0&5&-5&-3\end{pmatrix}
\xrightarrow[R_4\leftarrow R_4+R_3]{R_2\leftarrow R_2-R_1}
\begin{pmatrix}-3&-3&-3&0\\0&-4&4&8\\0&-1&1&-5\\0&4&-4&-8\end{pmatrix}
\\
\xrightarrow[R_4\leftarrow\frac14R_4]{R_2\leftarrow\frac14R_2}
\begin{pmatrix}-3&-3&-3&0\\0&-1&1&2\\0&-1&1&-5\\0&1&-1&-2\end{pmatrix}
\xrightarrow[R_4\leftarrow R_4+R_2]{R_3\leftarrow R_3-R_2}
\begin{pmatrix}-3&-3&-3&0\\0&-1&1&2\\0&0&0&-7\\0&0&0&0\end{pmatrix}
\end{gathered}
\]
so \(x_4=0\), and \(x_3=-1\) gives \(x_2=x_3+2x_4=-1\) and \(x_1=-x_2-x_3=2\);
\[
\begin{gathered}
A-I=\begin{pmatrix}5&-3&-3&0\\-3&6&-4&5\\-3&-4&6&-5\\0&5&-5&5\end{pmatrix}
\xrightarrow{R_2\leftarrow R_2-R_4}
\begin{pmatrix}5&-3&-3&0\\-3&1&1&0\\-3&-4&6&-5\\0&5&-5&5\end{pmatrix}
\xrightarrow[R_3\leftarrow R_3-R_2]{R_1\leftarrow R_1+3R_2}
\begin{pmatrix}-4&0&0&0\\-3&1&1&0\\0&-5&5&-5\\0&5&-5&5\end{pmatrix}
\\
\xrightarrow[R_4\leftarrow R_4+R_3]{R_1\leftarrow\frac14R_1}
\begin{pmatrix}-1&0&0&0\\-3&1&1&0\\0&-5&5&-5\\0&0&0&0\end{pmatrix}
\xrightarrow[R_3\leftarrow\frac15R_3]{R_2\leftarrow R_2-3R_1}
\begin{pmatrix}-1&0&0&0\\0&1&1&0\\0&-1&1&-1\\0&0&0&0\end{pmatrix}
\xrightarrow{R_3\leftarrow R_3+R_2}
\begin{pmatrix}-1&0&0&0\\0&1&1&0\\0&0&2&-1\\0&0&0&0\end{pmatrix}
\end{gathered}
\]
so \(x_1=0\), and \(x_3=-1\) gives \(x_4=2x_3=-2\) and \(x_2=-x_3=1\);
\[
\begin{gathered}
A=\begin{pmatrix}6&-3&-3&0\\-3&7&-4&5\\-3&-4&7&-5\\0&5&-5&6\end{pmatrix}
\xrightarrow[R_3\leftarrow R_3+R_4]{R_2\leftarrow R_2-R_4}
\begin{pmatrix}6&-3&-3&0\\-3&2&1&-1\\-3&1&2&1\\0&5&-5&6\end{pmatrix}
\xrightarrow[R_3\leftarrow R_3-R_2]{R_1\leftarrow R_1+R_2}
\begin{pmatrix}3&-1&-2&-1\\-3&2&1&-1\\0&-1&1&2\\0&5&-5&6\end{pmatrix}
\\
\xrightarrow[R_4\leftarrow R_4+2R_3]{R_2\leftarrow R_2+R_1}
\begin{pmatrix}3&-1&-2&-1\\0&1&-1&-2\\0&-1&1&2\\0&3&-3&10\end{pmatrix}
\xrightarrow[R_4\leftarrow R_4-3R_2]{R_3\leftarrow R_3+R_2}
\begin{pmatrix}3&-1&-2&-1\\0&1&-1&-2\\0&0&0&0\\0&0&0&16\end{pmatrix}
\\
\xrightarrow{R_3\leftrightarrow R_4}
\begin{pmatrix}3&-1&-2&-1\\0&1&-1&-2\\0&0&0&16\\0&0&0&0\end{pmatrix}
\end{gathered}
\]
so \(x_4=0\), and \(x_3=1\) gives \(x_2=x_3+2x_4=1\) and \(x_1=\frac13(x_2+2x_3+x_4)=1\). So the eigenspaces are spanned by
\((0,1,-1,1)'\), \((2,-1,-1,0)'\), \((0,1,-1,-2)'\) and \((1,1,1,0)'\), pairwise orthogonal (Theorem~\ref{r8-thm-sym-orth}); an
orthonormal basis is \(q_1=\frac1{\sqrt3}(0,1,-1,1)'\), \(q_2=\frac1{\sqrt6}(2,-1,-1,0)'\), \(q_3=\frac1{\sqrt6}(0,1,-1,-2)'\),
\(q_4=\frac1{\sqrt3}(1,1,1,0)'\).

(b) The eigenvalues are \(\geqslant0\) and one is \(0\), so \(A\) is PSD but not PD
(Theorem~\ref{r8-thm-eigen-criterion}(a)), and \(\rank A=3\) (Theorem~\ref{r8-thm-rank-eigen}). The leading minors are
\(6\), \(42-9=33\), along row \(1\) \(\Delta_3=6(49-16)+3(-21-12)-3(12+21)=0\), and \(\det A=16\cdot9\cdot1\cdot0=0\)
(Theorem~\ref{r7-thm-trace-det}).

(c) By Theorem~\ref{r8-thm-sqrt},
\[
\begin{aligned}
A^{1/2}&=4q_1q_1'+3q_2q_2'+1\,q_3q_3'+0\,q_4q_4'\\
&=\frac43\begin{pmatrix}0&0&0&0\\0&1&-1&1\\0&-1&1&-1\\0&1&-1&1\end{pmatrix}
+\frac12\begin{pmatrix}4&-2&-2&0\\-2&1&1&0\\-2&1&1&0\\0&0&0&0\end{pmatrix}
+\frac16\begin{pmatrix}0&0&0&0\\0&1&-1&-2\\0&-1&1&2\\0&-2&2&4\end{pmatrix}\\
&=\begin{pmatrix}2&-1&-1&0\\-1&2&-1&1\\-1&-1&2&-1\\0&1&-1&2\end{pmatrix}
\end{aligned}
\]
Check: \(\row_1(A^{1/2})A^{1/2}=(4+1+1+0,\ -2-2+1+0,\ -2+1-2+0,\ 0-1+1+0)=(6,-3,-3,0)=\row_1(A)\), and likewise for rows \(2\),
\(3\) and \(4\).
\end{solution}

\begin{exercise}\label{r8-ex-square-root}
Let \(A=I_4+6J_4\), the \(4\times4\) matrix with diagonal entries \(7\) and all other entries \(6\).
\begin{enumerate}
\item Find the eigenvalues of \(A\), and compute \(A^{1/2}\), \(A^{-1/2}\) and \(A^{-1}\). Check that \((A^{1/2})^2=A\).
\item Show that
\[
G=\begin{pmatrix}1&2&1&1\\2&1&1&1\\1&1&1&2\\1&1&2&1\end{pmatrix}
\]
also satisfies \(G^2=A\). Why does this not contradict Theorem~\ref{r8-thm-sqrt}(a)?
\end{enumerate}
\end{exercise}
\begin{solution}
(a) By Example~\ref{r8-exm-patterned}(b), with \(\alpha=1\) and \(\beta=6\), \(A=25\bar J_4+C_4\), and the eigenvalues are
\(25\), with \(\mathcal N(A-25I)=\Span\{\one_4\}\), and \(1\) (three times), with
\(\mathcal N(A-I)=\Span\{\one_4\}^\perp=\{x:\one'x=0\}\). The projection matrices onto these eigenspaces are
\(\bar J_4=\frac14J_4\) and \(C_4=I_4-\frac14J_4\) (Corollary~\ref{r6-cor-centering}(a)), so \(A=25\bar J_4+C_4\) is the
spectral decomposition of \(A\) (Theorem~\ref{r8-thm-spectral-decomp}(b)), and by Theorem~\ref{r8-thm-sqrt}(a),(c) and
Definition~\ref{r8-def-matrix-function}
\[
A^{1/2}=5\bar J_4+C_4=I_4+J_4,\qquad A^{-1/2}=\tfrac15\bar J_4+C_4=I_4-\tfrac15J_4,\qquad
A^{-1}=\tfrac1{25}\bar J_4+C_4=I_4-\tfrac6{25}J_4
\]
Indeed \((I_4+J_4)^2=I_4+2J_4+J_4^2=I_4+6J_4\), as \(J_4^2=4J_4\) (Proposition~\ref{r1-prop-centering}).

(b) The \((1,1)\), \((1,2)\), \((1,3)\), \((1,4)\) entries of \(G^2\) are \(1+4+1+1=7\), \(2+2+1+1=6\), \(1+2+1+2=6\),
\(1+2+2+1=6\), and the other rows follow in the same way: \(G^2=A\). But \(G\) is not PSD:
\((1,-1,0,0)G(1,-1,0,0)'=1-4+1=-2\). Theorem~\ref{r8-thm-sqrt}(a) asserts uniqueness only among PSD matrices. (With
\(h_1,\dots,h_4\) of Exercise~\ref{r8-ex-diagonalize}, \(G=\frac14(5h_1h_1'+h_2h_2'-h_3h_3'-h_4h_4')\), with the eigenvalues
\(5,1,-1,-1\).)
\end{solution}

\begin{exercise}\label{r8-dr-loewner}
Let
\[
B=\begin{pmatrix}1&-1&0&1\\-1&2&-1&-1\\0&-1&2&-1\\1&-1&-1&3\end{pmatrix},\qquad
A=\begin{pmatrix}1&-1&0&1\\-1&2&-1&-1\\0&-1&2&-1\\1&-1&-1&4\end{pmatrix}
\]
\begin{enumerate}
\item Show that \(B\) is PD and that \(A\succcurlyeq B\).
\item Find \(B^{-1}\) by Gauss--Jordan elimination, and \(A^{-1}\) by Theorem~\ref{r5d-thm-woodbury}(b).
\item Check that \(B^{-1}\succcurlyeq A^{-1}\) and \(\det A\geqslant\det B\) (Theorems~\ref{r8-thm-loewner-inverse}
and~\ref{r8-thm-det-monotone}).
\end{enumerate}
\end{exercise}
\begin{solution}
(a) \(\Delta_1(B)=1\), \(\Delta_2(B)=2-1=1\), and along row \(1\), \(\Delta_3(B)=1(4-1)+1(-2-0)=1\). For \(\det B\), apply
\(R_2\leftarrow R_2+R_1\), \(R_4\leftarrow R_4-R_1\), then \(R_3\leftarrow R_3+R_2\), then \(R_4\leftarrow R_4+R_3\)
(Definition~\ref{r5d-def-det}(b), Theorem~\ref{r5d-thm-basic}(c)):
\[
\det B=\det\begin{pmatrix}1&-1&0&1\\0&1&-1&0\\0&-1&2&-1\\0&0&-1&2\end{pmatrix}
=\det\begin{pmatrix}1&-1&0&1\\0&1&-1&0\\0&0&1&-1\\0&0&-1&2\end{pmatrix}
=\det\begin{pmatrix}1&-1&0&1\\0&1&-1&0\\0&0&1&-1\\0&0&0&1\end{pmatrix}=1
\]
So \(B\) is PD (Theorem~\ref{r8-thm-leading-minors}). \(A-B=e_4e_4'\) is PSD (Theorem~\ref{r8-thm-XtX}), so \(A\succcurlyeq B\).

(b) Row-reduce \([B\mid I_4]\) (Corollary~\ref{r2-cor-computing-inverse}(a)):
\[
\begin{gathered}
\left(\begin{array}{cccc|cccc}1&-1&0&1&1&0&0&0\\-1&2&-1&-1&0&1&0&0\\0&-1&2&-1&0&0&1&0\\1&-1&-1&3&0&0&0&1\end{array}\right)
\xrightarrow[R_4\leftarrow R_4-R_1]{R_2\leftarrow R_2+R_1}
\left(\begin{array}{cccc|cccc}1&-1&0&1&1&0&0&0\\0&1&-1&0&1&1&0&0\\0&-1&2&-1&0&0&1&0\\0&0&-1&2&-1&0&0&1\end{array}\right)
\\
\xrightarrow{R_3\leftarrow R_3+R_2}
\left(\begin{array}{cccc|cccc}1&-1&0&1&1&0&0&0\\0&1&-1&0&1&1&0&0\\0&0&1&-1&1&1&1&0\\0&0&-1&2&-1&0&0&1\end{array}\right)
\xrightarrow{R_4\leftarrow R_4+R_3}
\left(\begin{array}{cccc|cccc}1&-1&0&1&1&0&0&0\\0&1&-1&0&1&1&0&0\\0&0&1&-1&1&1&1&0\\0&0&0&1&0&1&1&1\end{array}\right)
\\
\xrightarrow[R_3\leftarrow R_3+R_4]{R_1\leftarrow R_1-R_4}
\left(\begin{array}{cccc|cccc}1&-1&0&0&1&-1&-1&-1\\0&1&-1&0&1&1&0&0\\0&0&1&0&1&2&2&1\\0&0&0&1&0&1&1&1\end{array}\right)
\xrightarrow{R_2\leftarrow R_2+R_3}
\left(\begin{array}{cccc|cccc}1&-1&0&0&1&-1&-1&-1\\0&1&0&0&2&3&2&1\\0&0&1&0&1&2&2&1\\0&0&0&1&0&1&1&1\end{array}\right)
\\
\xrightarrow{R_1\leftarrow R_1+R_2}
\left(\begin{array}{cccc|cccc}1&0&0&0&3&2&1&0\\0&1&0&0&2&3&2&1\\0&0&1&0&1&2&2&1\\0&0&0&1&0&1&1&1\end{array}\right)
\end{gathered}
\]
So
\[
B^{-1}=\begin{pmatrix}3&2&1&0\\2&3&2&1\\1&2&2&1\\0&1&1&1\end{pmatrix}
\]
Now \(A=B+e_4e_4'\), \(w=B^{-1}e_4=(0,1,1,1)'\) and \(1+e_4'B^{-1}e_4=2\), so by Theorem~\ref{r5d-thm-woodbury}(b)
\[
A^{-1}=B^{-1}-\tfrac12ww'=\frac12\begin{pmatrix}6&4&2&0\\4&5&3&1\\2&3&3&1\\0&1&1&1\end{pmatrix}
\]
Check: \(\row_1(A)A^{-1}=\frac12(6-4+0+0,\ 4-5+0+1,\ 2-3+0+1,\ 0-1+0+1)=(1,0,0,0)\), and likewise for rows \(2\), \(3\) and \(4\).

(c) \(B^{-1}-A^{-1}=\frac12ww'\) is PSD (Theorem~\ref{r8-thm-XtX}), so \(B^{-1}\succcurlyeq A^{-1}\). And
\(\det A=\det B\cdot(1+e_4'B^{-1}e_4)=2\geqslant1=\det B\) (Theorem~\ref{r5d-thm-det-lemma}).
\end{solution}

\begin{exercise}\label{r8-dr-commuting}
Let
\[
A=\begin{pmatrix}2&-1&-1&1\\-1&2&1&-1\\-1&1&2&-1\\1&-1&-1&2\end{pmatrix},\qquad
B=\begin{pmatrix}1&1&1&0\\1&-1&2&1\\1&2&-1&1\\0&1&1&1\end{pmatrix}
\]
\begin{enumerate}
\item Compute \(AB\), and conclude that \(AB=BA\).
\item Find the eigenvalues and eigenspaces of \(A\).
\item Find an orthogonal \(Q\) such that \(Q'AQ\) and \(Q'BQ\) are both diagonal (Theorem~\ref{r8-thm-commuting}).
\end{enumerate}
\end{exercise}
\begin{solution}
(a) The product is
\[
AB=\begin{pmatrix}0&2&2&-1\\2&-2&1&2\\2&1&-2&2\\-1&2&2&0\end{pmatrix}
\]
It is symmetric, so \(BA=B'A'=(AB)'=AB\) (Theorem~\ref{r1-thm-transpose}).

(b) Row-reduce \(A-5I\) and \(A-I\):
\[
\begin{gathered}
A-5I=\begin{pmatrix}-3&-1&-1&1\\-1&-3&1&-1\\-1&1&-3&-1\\1&-1&-1&-3\end{pmatrix}
\xrightarrow[R_4\leftarrow R_4+R_1]{R_2\leftarrow R_2+R_3}
\begin{pmatrix}-3&-1&-1&1\\-2&-2&-2&-2\\-1&1&-3&-1\\-2&-2&-2&-2\end{pmatrix}
\xrightarrow[R_4\leftarrow R_4-R_2]{R_1\leftarrow R_1-R_2}
\begin{pmatrix}-1&1&1&3\\-2&-2&-2&-2\\-1&1&-3&-1\\0&0&0&0\end{pmatrix}
\\
\xrightarrow[R_3\leftarrow R_3-R_1]{R_2\leftarrow R_2-2R_1}
\begin{pmatrix}-1&1&1&3\\0&-4&-4&-8\\0&0&-4&-4\\0&0&0&0\end{pmatrix}
\end{gathered}
\]
so \(x_4=1\) gives \(x_3=-x_4=-1\), \(x_2=-x_3-2x_4=-1\) and \(x_1=x_2+x_3+3x_4=1\);
\[
A-I=\begin{pmatrix}1&-1&-1&1\\-1&1&1&-1\\-1&1&1&-1\\1&-1&-1&1\end{pmatrix}
\xrightarrow[R_4\leftarrow R_4-R_1]{R_2\leftarrow R_2+R_1,\ R_3\leftarrow R_3+R_1}
\begin{pmatrix}1&-1&-1&1\\0&0&0&0\\0&0&0&0\\0&0&0&0\end{pmatrix}
\]
So \(\mathcal N(A-5I)=\Span\{v\}\) with \(v=(1,-1,-1,1)'\), and \(\mathcal N(A-I)=\{x:x_1=x_2+x_3-x_4\}\), of dimension \(3\). As
the dimensions add up to \(4\), the eigenvalues are \(5\) and \(1\), with \(a_A(1)=3\) (Theorem~\ref{r8-thm-rank-eigen}).

(c) \(\mathcal N(A-I)\) is invariant under \(B\): for \(x\) in it, \(A(Bx)=B(Ax)=Bx\). For \(x=(b+c-d,b,c,d)'\) in
\(\mathcal N(A-I)\), \(Bx=(2b+2c-d,\ 3c,\ 3b,\ b+c+d)'\). If \(Bx=\mu x\), entries \(2\) and \(3\) give \(3c=\mu b\) and
\(3b=\mu c\), so \(9b=3\mu c=\mu^2b\) and likewise \(9c=\mu^2c\); entry \(4\) gives \(b+c+d=\mu d\).
\begin{itemize}
\item If \(\mu^2\ne9\), then \(b=c=0\), and \(d=\mu d\) with \(d\ne0\) (as \(x\ne0\)): \(\mu=1\) and \(x=d(-1,0,0,1)'\).
\item If \(\mu=3\), then \(c=b\) and \(2b+d=3d\), so \(d=b\): \(x=b(1,1,1,1)'\).
\item If \(\mu=-3\), then \(c=-b\) and \(d=-3d\), so \(d=0\): \(x=b(0,1,-1,0)'\).
\end{itemize}
The first entry then agrees. Also \(Bv=(-1,1,1,-1)'=-v\). Normalizing,
\[
Q=\begin{pmatrix}1/2&1/2&0&-1/\sqrt2\\-1/2&1/2&1/\sqrt2&0\\-1/2&1/2&-1/\sqrt2&0\\1/2&1/2&0&1/\sqrt2\end{pmatrix},\qquad
Q'AQ=\diag(5,1,1,1),\qquad Q'BQ=\diag(-1,3,-3,1)
\]
\end{solution}

\begin{exercise}\label{r8-dr-simdiag}
Let
\[
A=\begin{pmatrix}2&1&1&1\\1&2&1&1\\1&1&2&1\\1&1&1&1\end{pmatrix},\qquad
B=\begin{pmatrix}0&-1&-1&-1\\-1&1&-1&-1\\-1&-1&2&-1\\-1&-1&-1&-1\end{pmatrix}
\]
\begin{enumerate}
\item Show that \(A\) is PD, and find the roots of \(\det(B-tA)=0\).
\item Find an invertible \(E\) with \(E'AE=I_4\) and \(E'BE\) diagonal (Theorem~\ref{r8-thm-pd-simdiag}).
\end{enumerate}
\end{exercise}
\begin{solution}
(a) \(\Delta_1(A)=2\), \(\Delta_2(A)=3\), and along row \(1\), \(\Delta_3(A)=2(4-1)-1(2-1)+1(1-2)=4\). Applying
\(R_1\leftarrow R_1-R_4\), \(R_2\leftarrow R_2-R_4\) and \(R_3\leftarrow R_3-R_4\) to \(A\) gives the lower triangular matrix
with rows \(e_1'\), \(e_2'\), \(e_3'\), \(\one_4'\), so \(\det A=1\) (Definition~\ref{r5d-def-det}(b),
Theorem~\ref{r5d-thm-transpose}); so \(A\) is PD (Theorem~\ref{r8-thm-leading-minors}). For real \(t\), apply the same row
operations to \(B-tA\); the result is lower triangular:
\[
\begin{aligned}
\det(B-tA)&=\det\begin{pmatrix}-2t&-1-t&-1-t&-1-t\\-1-t&1-2t&-1-t&-1-t\\-1-t&-1-t&2-2t&-1-t\\-1-t&-1-t&-1-t&-1-t\end{pmatrix}\\
&=\det\begin{pmatrix}1-t&0&0&0\\0&2-t&0&0\\0&0&3-t&0\\-1-t&-1-t&-1-t&-1-t\end{pmatrix}=(1-t)(2-t)(3-t)(-1-t)
\end{aligned}
\]
The roots are \(-1\), \(1\), \(2\), \(3\).

(b) Row-reduce \(B-cA\) for \(c=-1,1,2,3\):
\[
B+A=\begin{pmatrix}2&0&0&0\\0&3&0&0\\0&0&4&0\\0&0&0&0\end{pmatrix}
\xrightarrow[R_3\leftarrow\frac14R_3]{R_1\leftarrow\frac12R_1,\ R_2\leftarrow\frac13R_2}
\begin{pmatrix}1&0&0&0\\0&1&0&0\\0&0&1&0\\0&0&0&0\end{pmatrix}
\]
\[
B-A=\begin{pmatrix}-2&-2&-2&-2\\-2&-1&-2&-2\\-2&-2&0&-2\\-2&-2&-2&-2\end{pmatrix}
\xrightarrow[R_4\leftarrow R_4-R_1]{R_2\leftarrow R_2-R_1,\ R_3\leftarrow R_3-R_1}
\begin{pmatrix}-2&-2&-2&-2\\0&1&0&0\\0&0&2&0\\0&0&0&0\end{pmatrix}
\]
\[
\begin{gathered}
B-2A=\begin{pmatrix}-4&-3&-3&-3\\-3&-3&-3&-3\\-3&-3&-2&-3\\-3&-3&-3&-3\end{pmatrix}
\xrightarrow[R_4\leftarrow R_4-R_2]{R_1\leftarrow R_1-R_2,\ R_3\leftarrow R_3-R_2}
\begin{pmatrix}-1&0&0&0\\-3&-3&-3&-3\\0&0&1&0\\0&0&0&0\end{pmatrix}
\\
\xrightarrow{R_2\leftarrow R_2-3R_1}
\begin{pmatrix}-1&0&0&0\\0&-3&-3&-3\\0&0&1&0\\0&0&0&0\end{pmatrix}
\end{gathered}
\]
\[
\begin{gathered}
B-3A=\begin{pmatrix}-6&-4&-4&-4\\-4&-5&-4&-4\\-4&-4&-4&-4\\-4&-4&-4&-4\end{pmatrix}
\xrightarrow[R_4\leftarrow R_4-R_3]{R_1\leftarrow R_1-R_3,\ R_2\leftarrow R_2-R_3}
\begin{pmatrix}-2&0&0&0\\0&-1&0&0\\-4&-4&-4&-4\\0&0&0&0\end{pmatrix}
\\
\xrightarrow{R_3\leftarrow R_3-2R_1}
\begin{pmatrix}-2&0&0&0\\0&-1&0&0\\0&-4&-4&-4\\0&0&0&0\end{pmatrix}
\xrightarrow{R_3\leftarrow R_3-4R_2}
\begin{pmatrix}-2&0&0&0\\0&-1&0&0\\0&0&-4&-4\\0&0&0&0\end{pmatrix}
\end{gathered}
\]
So \(Bf=cAf\) for \(f_1=(0,0,0,1)'\) (\(c=-1\)), \(f_2=(1,0,0,-1)'\) (\(c=1\)), \(f_3=(0,1,0,-1)'\) (\(c=2\)) and \(f_4=(0,0,1,-1)'\)
(\(c=3\)). The roots are distinct, so \(f_i'Af_j=0\) for \(i\ne j\) (Theorem~\ref{r8-thm-pd-simdiag}(c)); and \(f_1'Af_1=1\),
\(f_2'Af_2=2-2+1=1\), \(f_3'Af_3=2-2+1=1\), \(f_4'Af_4=2-2+1=1\). By Theorem~\ref{r8-thm-pd-simdiag}(c),
\[
E=\begin{pmatrix}0&1&0&0\\0&0&1&0\\0&0&0&1\\1&-1&-1&-1\end{pmatrix},\qquad
AE=\begin{pmatrix}1&1&0&0\\1&0&1&0\\1&0&0&1\\1&0&0&0\end{pmatrix},\qquad
BE=\begin{pmatrix}-1&1&0&0\\-1&0&2&0\\-1&0&0&3\\-1&0&0&0\end{pmatrix}
\]
and \(E'AE=I_4\), \(E'BE=\diag(-1,1,2,3)\).
\end{solution}

\begin{exercise}\label{r8-ex-simultaneous}
Let
\[
A=\begin{pmatrix}1&1&0&0\\1&2&1&0\\0&1&2&1\\0&0&1&2\end{pmatrix},\qquad
B=\begin{pmatrix}4&4&0&0\\4&7&3&0\\0&3&5&2\\0&0&2&3\end{pmatrix}
\]
\begin{enumerate}
\item Show that \(A\) is PD, and find an invertible \(E\) with \(E'AE=I_4\) and \(E'BE\) diagonal.
\item Show that \(x'Ax\leqslant x'Bx\leqslant4\,x'Ax\) for every \(x\in\R^4\), that is, \(A\preccurlyeq B\preccurlyeq4A\).
\end{enumerate}
\end{exercise}
\begin{solution}
(a) Expanding \(\Delta_k(A)=\det A_k\) along its last row gives \(\Delta_k(A)=2\Delta_{k-1}(A)-\Delta_{k-2}(A)\) for \(k=3,4\)
(Theorem~\ref{r5d-thm-expansion}), so the leading principal minors are \(1\), \(2-1=1\), \(2-1=1\), \(2-1=1\), and \(A\) is PD
(Theorem~\ref{r8-thm-leading-minors}). The rows of \(B-\lambda A\) are \((4-\lambda,\,4-\lambda,\,0,\,0)\),
\((4-\lambda,\,7-2\lambda,\,3-\lambda,\,0)\), \((0,\,3-\lambda,\,5-2\lambda,\,2-\lambda)\), \((0,\,0,\,2-\lambda,\,3-2\lambda)\). Subtract
row \(1\) from row \(2\), then the new row \(2\) from row \(3\), then the new row \(3\) from row \(4\)
(Definition~\ref{r5d-def-det}(b)):
\[
\det(B-\lambda A)=\det\begin{pmatrix}4-\lambda&4-\lambda&0&0\\0&3-\lambda&3-\lambda&0\\
0&0&2-\lambda&2-\lambda\\0&0&0&1-\lambda\end{pmatrix}=(4-\lambda)(3-\lambda)(2-\lambda)(1-\lambda)
\]
(Theorem~\ref{r5d-thm-basic}(c)), so \(\lambda_1=4\), \(\lambda_2=3\), \(\lambda_3=2\), \(\lambda_4=1\)
(Theorem~\ref{r8-thm-pd-simdiag}(b), as \(\det A=1\)). The null spaces of
\[
\begin{gathered}
B-4A=\begin{pmatrix}0&0&0&0\\0&-1&-1&0\\0&-1&-3&-2\\0&0&-2&-5\end{pmatrix},\qquad
B-3A=\begin{pmatrix}1&1&0&0\\1&1&0&0\\0&0&-1&-1\\0&0&-1&-3\end{pmatrix}\\
B-2A=\begin{pmatrix}2&2&0&0\\2&3&1&0\\0&1&1&0\\0&0&0&-1\end{pmatrix},\qquad
B-A=\begin{pmatrix}3&3&0&0\\3&5&2&0\\0&2&3&1\\0&0&1&1\end{pmatrix}
\end{gathered}
\]
contain \(f_1=(1,0,0,0)'\), \(f_2=(-1,1,0,0)'\), \(f_3=(1,-1,1,0)'\) and \(f_4=(-1,1,-1,1)'\). Here \(Af_1=(1,1,0,0)'\),
\(Af_2=(0,1,1,0)'\), \(Af_3=(0,0,1,1)'\) and \(Af_4=(0,0,0,1)'\), so \(f_k'Af_k=1\), and \(f_j'Af_k=0\) for \(j\ne k\)
(Theorem~\ref{r8-thm-pd-simdiag}(c)). Hence
\[
E=\begin{pmatrix}1&-1&1&-1\\0&1&-1&1\\0&0&1&-1\\0&0&0&1\end{pmatrix}
\]
has \(E'AE=I_4\) and \(E'BE=\diag(4,3,2,1)\), since \((E'BE)_{jk}=f_j'Bf_k=\lambda_kf_j'Af_k\).

(b) \(\det E=1\), so \(E\) is invertible (Theorem~\ref{r5d-thm-product}(a)) and every \(x\) is \(x=Ey\) with \(y=E^{-1}x\);
here \(E^{-1}\) is the upper triangular matrix with ones on the diagonal and on the superdiagonal (multiply out \(E\) with
it), so \(y_k=x_k+x_{k+1}\) for \(k\leqslant3\) and \(y_4=x_4\). Then \(x'Ax=y'E'AEy=y_1^2+y_2^2+y_3^2+y_4^2\) and
\(x'Bx=y'E'BEy=4y_1^2+3y_2^2+2y_3^2+y_4^2\). Hence \(x'Ax\leqslant x'Bx\leqslant4\,x'Ax\).
\end{solution}

\needspace{14\baselineskip}
\subsection*{Problems}

\begin{problem}[\citeat{searle-khuri}{Exercise~6.17}]\label{r8-ex-zero-eigenvalues}
\begin{enumerate}
\item Show that a symmetric matrix whose eigenvalues are all \(0\) is the zero matrix.
\item Show by a \(2\times2\) example that (a) fails without symmetry, and that the same example shows that
Theorem~\ref{r8-thm-rank-eigen} needs symmetry.
\end{enumerate}
\end{problem}
\begin{solution}
(a) By Theorem~\ref{r8-thm-spectral}, \(A=Q\Lambda Q'\), and the diagonal entries of \(\Lambda\) are the eigenvalues of
\(A\). If they are all \(0\), then \(\Lambda=0\) and \(A=0\).

(b) Let \(A=\left(\begin{smallmatrix}0&1\\0&0\end{smallmatrix}\right)\). Then
\(p_A(t)=\det\left(\begin{smallmatrix}t&-1\\0&t\end{smallmatrix}\right)=t^2\), so both eigenvalues are \(0\), but
\(A\ne0\). Also \(\rank A=1\) (one nonzero row, in echelon form), although \(A\) has no nonzero eigenvalue.
\end{solution}

\begin{problem}[\citeat{abadir-magnus}{Exercise~8.62}]\label{r8-ex-idempotent-rank}
Show that a symmetric \(A\in\R^{n\times n}\) is idempotent if and only if \(\rank A+\rank(I_n-A)=n\).
\end{problem}
\begin{solution}
Let \(A=Q\Lambda Q'\) (Theorem~\ref{r8-thm-spectral}), and let \(n_0\) and \(n_1\) be the numbers of eigenvalues
\(\lambda_i\) equal to \(0\) and to \(1\).
\begin{itemize}
\item \(\rank A=n-n_0\) (Theorem~\ref{r8-thm-rank-eigen}). \(I_n-A=(-1)A+I_n\) is symmetric with the eigenvalues
\(1-\lambda_i\) (Proposition~\ref{r7-prop-eig-poly}(c), with \(c=-1\) and then \(c=1\)), of which \(n_1\) are \(0\); so \(\rank(I_n-A)=n-n_1\)
(Theorem~\ref{r8-thm-rank-eigen}).
\item Hence \(\rank A+\rank(I_n-A)=n+(n-n_0-n_1)\), which is \(n\) if and only if \(n_0+n_1=n\), that is, every
eigenvalue of \(A\) is \(0\) or \(1\).
\item This holds if and only if \(A\) is idempotent: by Theorem~\ref{r8-thm-idempotent-eigen}(a) if \(A\) is idempotent,
and by Theorem~\ref{r8-thm-idempotent-eigen}(b) conversely.\qedhere
\end{itemize}
\end{solution}

\begin{problem}[\citeat{abadir-magnus}{Exercise~8.75}]\label{r8-ex-pairwise}
Find the matrix of the quadratic form \(\sum_{i=1}^n\sum_{j=1}^n(x_i-x_j)^2\) on \(\R^n\), and its eigenvalues and
eigenvectors.
\end{problem}
\begin{solution}
\[
\begin{aligned}
\sum_{i,j}(x_i-x_j)^2&=\sum_{i,j}(x_i^2+x_j^2-2x_ix_j)=2n\sum_ix_i^2-2\Bigl(\sum_ix_i\Bigr)^2\\
&=2n\,x'x-2(\one'x)^2=2n\,x'\bigl(I_n-\tfrac1n\one\one'\bigr)x=2n\,x'C_nx
\end{aligned}
\]
So the matrix of the form is the symmetric matrix \(2nC_n\) (Proposition~\ref{r8-prop-qf-symmetric}). \(C_n\) has the
eigenvalues \(1\) (\(n-1\) times) and \(0\) (once) (Example~\ref{r8-exm-patterned}(a)), so \(2nC_n\) has the eigenvalues
(Proposition~\ref{r7-prop-eig-poly}(c))
\(2n\) (\(n-1\) times), with eigenvectors the nonzero vectors orthogonal to \(\one\), and \(0\) (once), with eigenvector
\(\one\). By Example~\ref{r1-exm-centering}, the form equals \(2n\sum_i(x_i-\bar x)^2\).
\end{solution}

\begin{problem}[\citeat{harville-solutions}{Ch.~14, Exercise~40}]\label{r8-ex-minors-counter}
Show by examples that the leading principal minors of a symmetric matrix can all be \(\geqslant0\) although the matrix
is not PSD: for a \(2\times2\) matrix, and, for every \(n\geqslant3\), for a nonsingular \(n\times n\) matrix.
\end{problem}
\begin{solution}
\begin{itemize}
\item \(\left(\begin{smallmatrix}0&0\\0&-1\end{smallmatrix}\right)\) has the leading minors \(0\), \(0\), but its
diagonal entry \(-1\) is negative, so it is not PSD (Theorem~\ref{r8-thm-principal-sub}).
\item Let
\[
A_*=\begin{pmatrix}0&0&1\\0&-1&0\\1&0&0\end{pmatrix}
\]
Its leading minors are \(0\),
\(\det\left(\begin{smallmatrix}0&0\\0&-1\end{smallmatrix}\right)=0\) and, along row \(1\),
\(\det A_*=1\cdot\det\left(\begin{smallmatrix}0&-1\\1&0\end{smallmatrix}\right)=1\). So \(A_*\) is nonsingular, but not
PSD, as \(a_{22}=-1<0\).
\item For \(n\geqslant4\) let \(A=\diag(A_*,I_{n-3})\). Its first three leading minors are those of \(A_*\), and for
\(k\geqslant4\), \(A_k=\diag(A_*,I_{k-3})\) has determinant \(\det A_*=1\) (Theorem~\ref{r5d-thm-block-tri}). So
\(A\) is nonsingular, with leading minors \(\geqslant0\), and not PSD, as \(a_{22}=-1\).\qedhere
\end{itemize}
\end{solution}

\begin{problem}[\citeat{abadir-magnus}{Exercise~12.26(b)}]\label{r8-pr-hadamard}
Let \(A=(a_{ij})\in\R^{n\times n}\). Use Hadamard's inequality (Theorem~\ref{r8-thm-hadamard}) to show that
\[
(\det A)^2\leqslant\prod_{i=1}^n\Bigl(\sum_{j=1}^na_{ij}^2\Bigr)
\]
\end{problem}
\begin{solution}
\begin{itemize}
\item If \(\det A=0\), the inequality holds, as the right side is a product of sums of squares, hence \(\geqslant0\).
\item Let \(\det A\ne0\). Then \(\rank A=n\) (Theorem~\ref{r5d-thm-product}(a)), so \(AA'\) is PD
(Theorem~\ref{r8-thm-XtX}), with \(\det(AA')=\det A\cdot\det A'=(\det A)^2\) (Theorems~\ref{r5d-thm-product}(b)
and~\ref{r5d-thm-transpose}).
\item The \(i\)th diagonal entry of \(AA'\) is \(\row_i(A)\row_i(A)'=\sum_ja_{ij}^2\) (Theorem~\ref{r1-thm-views}(d)). By
Theorem~\ref{r8-thm-hadamard},
\[
(\det A)^2=\det(AA')\leqslant\prod_{i=1}^n(AA')_{ii}=\prod_{i=1}^n\Bigl(\sum_{j=1}^na_{ij}^2\Bigr)
\]
\end{itemize}
\end{solution}

\begin{problem}[\citeat{searle-khuri}{Exercise~5.19}]\label{r8-ex-idempotent-V}
Let \(V\in\R^{n\times n}\) be PD with \(V=P'P\) for some \(P\in\R^{n\times n}\), and let \(A\in\R^{n\times n}\) be
symmetric. Show that
\begin{enumerate}
\item \(P\) is invertible;
\item \(PAP'\) is idempotent if and only if \(AV\) is idempotent;
\item if \(AV\) is idempotent, then \(\rank(PAP')=\tr(AV)\).
\end{enumerate}
\end{problem}
\begin{solution}
(a) \(\det V=\det P'\det P=(\det P)^2\) (Theorems~\ref{r5d-thm-product}(b) and~\ref{r5d-thm-transpose}), and
\(\det V\ne0\) (Theorem~\ref{r8-thm-eigen-criterion}(c)); so \(\det P\ne0\), and \(P\) is invertible
(Theorem~\ref{r5d-thm-product}(a)).

(b)
\begin{itemize}
\item If \(AV\) is idempotent, then \((AP'P)(AP'P)=AP'P\); multiplying on the right by \((P'P)^{-1}=V^{-1}\) gives
\(AP'PA=A\), and then \((PAP')(PAP')=P(AP'PA)P'=PAP'\).
\item If \(PAP'\) is idempotent, then \(P(AP'PA)P'=PAP'\); multiplying on the left by \(P^{-1}\) and on the right by
\((P')^{-1}\) gives \(AP'PA=A\), that is, \(AVA=A\), and multiplying on the right by \(V\) gives \(AVAV=AV\).
\end{itemize}
(c) By (b), \(PAP'\) is idempotent, so \(\rank(PAP')=\tr(PAP')=\tr(AP'P)=\tr(AV)\)
(Theorems~\ref{r6-thm-idempotent}(c) and~\ref{r1-thm-trace}(b)).
\end{solution}

\begin{problem}[\citeat{harville}{Theorem~14.7.6, Corollary~14.7.7}]\label{r8-pr-trace-psd}
Let \(A,B\in\R^{n\times n}\) be PSD. Show that \(\tr(AB)\geqslant0\), with equality if and only if \(AB=0\).
\end{problem}
\begin{solution}
If \(B=0\), then \(\tr(AB)=0\) and \(AB=0\). Let \(B\ne0\) and \(r=\rank B\geqslant1\).
\begin{itemize}
\item By Theorem~\ref{r8-thm-psd-factor}(a), \(B=Q'Q\) for some \(Q\in\R^{r\times n}\) of rank \(r\). Then
\(\tr(AB)=\tr(AQ'Q)=\tr(QAQ')\) (Theorem~\ref{r1-thm-trace}(b)).
\item \(M=QAQ'\in\R^{r\times r}\) is PSD (Theorem~\ref{r8-thm-congruence}(a), with \(Q'\) in place of \(P\)), so its diagonal
entries \(m_{ii}\) are \(\geqslant0\) (Theorem~\ref{r8-thm-principal-sub}) and \(\tr M=m_{11}+\dots+m_{rr}\geqslant0\). If
\(\tr M=0\), then every \(m_{ii}=0\), so every row of \(M\) is zero (Theorem~\ref{r8-thm-principal-sub}): \(M=0\). Hence
\(\tr(AB)\geqslant0\), with equality if and only if \(QAQ'=0\).
\item \(QAQ'=0\) if and only if \(BAB=0\). Indeed \(BAB=Q'(QAQ')Q\). Conversely, \(Q\) has a right inverse \(K\), \(QK=I_r\)
(Theorem~\ref{thm:r4-one-sided}(a)), so \(K'Q'=I_r\) (Theorem~\ref{r1-thm-transpose}), and if \(BAB=0\), then
\(QAQ'=K'Q'(QAQ')QK=K'(BAB)K=0\).
\item \(BAB=0\) if and only if \(AB=0\). If \(AB=0\), then \(BAB=0\). Conversely, let \(BAB=0\). By
Theorem~\ref{r8-thm-psd-factor}(b), \(A=P'P\) for some matrix \(P\) with \(n\) columns, and, as \(B'=B\),
\(0=BAB=(PB)'(PB)\). So \(\tr\bigl((PB)'(PB)\bigr)=0\), and \(PB=0\) (Theorem~\ref{r1-thm-trace}(d)); hence
\(AB=P'(PB)=0\).
\end{itemize}
So \(\tr(AB)=0\) if and only if \(QAQ'=0\), if and only if \(BAB=0\), if and only if \(AB=0\).
\end{solution}

\begin{problem}[\citeat{searle-khuri}{Exercise~9.5}]\label{r8-ex-psd-cs}
Let \(A\in\R^{n\times n}\) be PSD. Show that \((x'Ay)^2\leqslant(x'Ax)(y'Ay)\) for all \(x,y\in\R^n\).
\end{problem}
\begin{solution}
Let \(z=A^{1/2}x\) and \(w=A^{1/2}y\) (Theorem~\ref{r8-thm-sqrt}). As \(A^{1/2}\) is symmetric and
\(A^{1/2}A^{1/2}=A\), \(z'w=x'Ay\), \(z'z=x'Ax\) and \(w'w=y'Ay\). By the Cauchy--Schwarz inequality
(Theorem~\ref{r6-thm-cauchy-schwarz}(a)), \((x'Ay)^2=(z'w)^2\leqslant(z'z)(w'w)=(x'Ax)(y'Ay)\).
\end{solution}

\begin{problem}[\citeat{harville-solutions}{Ch.~14, Exercise~23}]\label{r8-ex-qr-cholesky}
Let \(A\in\R^{n\times p}\) have rank \(p\), and let \(A=QR\) as in Theorem~\ref{r6-thm-qr}. Show that \(A'A=R'R\), and that \(R\)
is the Cholesky factor of \(A'A\).
\end{problem}
\begin{solution}
\(Q'Q=I_p\), so \(A'A=R'Q'QR=R'R\). The matrix \(A'A\) is PD (Theorem~\ref{r8-thm-XtX}), and \(R\) is upper triangular with
positive diagonal entries (Theorem~\ref{r6-thm-qr}); by the uniqueness in Theorem~\ref{r8-thm-cholesky}, \(R\) is the Cholesky
factor of \(A'A\).
\end{solution}

\begin{problem}[\citeat{bapat}{Ch.~11, Exercise~49}]\label{r8-ex-rank-order}
Let \(A,B\in\R^{n\times n}\) be PSD, and let \(A-B\) be PSD. Show that \(\rank A\geqslant\rank B\).
\end{problem}
\begin{solution}
It suffices to show \(\mathcal N(A)\subseteq\mathcal N(B)\): then \(\dim\mathcal N(A)\leqslant\dim\mathcal N(B)\)
(Theorem~\ref{thm:r4-counting}(a)), and \(\rank A=n-\dim\mathcal N(A)\geqslant n-\dim\mathcal N(B)=\rank B\)
(Theorem~\ref{thm:r4-rank-nullity}).
\begin{itemize}
\item Let \(x\in\mathcal N(A)\). Then \(x'Ax=0\), and \(x'(A-B)x\geqslant0\) gives \(x'Bx\leqslant0\); as \(B\) is PSD,
\(x'Bx=0\).
\item By Corollary~\ref{r8-cor-psd-null}, \(Bx=0\): \(x\in\mathcal N(B)\).
\end{itemize}
\end{solution}

\begin{problem}[\citeat{abadir-magnus}{Exercise~12.19(b)}]\label{r8-pr-loewner-square}
Show that for symmetric matrices \(A\) and \(B\), \(A\succcurlyeq B\succcurlyeq0\) does not imply \(A^2\succcurlyeq B^2\).
\end{problem}
\begin{solution}
Let
\[
A=\begin{pmatrix}2&1\\1&1\end{pmatrix},\qquad B=\begin{pmatrix}1&0\\0&0\end{pmatrix}
\]
\begin{itemize}
\item \(A-B=\left(\begin{smallmatrix}1&1\\1&1\end{smallmatrix}\right)=\one_2\one_2'\) and \(B=e_1e_1'\) are PSD
(Theorem~\ref{r8-thm-XtX}), so \(A\succcurlyeq B\succcurlyeq0\).
\item \(A^2=\left(\begin{smallmatrix}5&3\\3&2\end{smallmatrix}\right)\) and \(B^2=B\), so
\(A^2-B^2=\left(\begin{smallmatrix}4&3\\3&2\end{smallmatrix}\right)\), with \(\det(A^2-B^2)=8-9=-1<0\). A PSD matrix has a
determinant \(\geqslant0\) (Theorem~\ref{r8-thm-eigen-criterion}(c)), so \(A^2-B^2\) is not PSD: \(A^2\succcurlyeq B^2\)
fails. For instance \(x=(2,-3)'\) gives \(x'(A^2-B^2)x=16-36+18=-2\).
\end{itemize}
\end{solution}

\begin{problem}[\citeat{abadir-magnus}{Exercise~12.18}]\label{r8-ex-loewner-ab}
Let \(A,B\in\R^{n\times k}\) with \(A'B=I_k\). Show that \(A'A\succcurlyeq(B'B)^{-1}\) and \(B'B\succcurlyeq(A'A)^{-1}\).
\end{problem}
\begin{solution}
\begin{itemize}
\item \(k=\rank I_k=\rank(A'B)\leqslant\rank B\leqslant k\) (Theorems~\ref{thm:r4-rank-product}(a) and~\ref{thm:r4-rank-rules}(a)),
so \(B\) has full column rank and \(B'B\) is invertible (Theorem~\ref{r6-thm-hat}); likewise for \(A\).
\item \(M=I_n-B(B'B)^{-1}B'\) is symmetric and idempotent (Theorem~\ref{r6-thm-hat}(b)), hence PSD
(Proposition~\ref{r8-prop-pd-basic}(d)), and \(A'MA\) is PSD (Theorem~\ref{r8-thm-congruence}(a)).
\item \(B'A=(A'B)'=I_k\), so \(A'MA=A'A-A'B(B'B)^{-1}B'A=A'A-(B'B)^{-1}\): \(A'A\succcurlyeq(B'B)^{-1}\).
\item As \(B'A=I_k\), the same argument with \(A\) and \(B\) exchanged, with \(I_n-A(A'A)^{-1}A'\), gives
\(B'B\succcurlyeq(A'A)^{-1}\).
\end{itemize}
\end{solution}

\begin{problem}[\citeat{abadir-magnus}{Exercise~8.58}]\label{r8-ex-loewner-idempotent}
Let \(A,B\in\R^{n\times n}\) be symmetric and idempotent. Show that \(A\succcurlyeq B\) if and only if \(AB=B\).
\end{problem}
\begin{solution}
\begin{itemize}
\item Let \(AB=B\). Then \(BA=B'A'=(AB)'=B'=B\), and \((A-B)^2=A^2-AB-BA+B^2=A-B-B+B=A-B\). So \(A-B\) is symmetric and
idempotent; its eigenvalues are \(0\) and \(1\) (Theorem~\ref{r8-thm-idempotent-eigen}(a)), so it is PSD
(Theorem~\ref{r8-thm-eigen-criterion}(a)): \(A\succcurlyeq B\).
\item Conversely, let \(A\succcurlyeq B\), and let \(C=(I_n-A)B\). Then \(C'C\succcurlyeq0\) (Theorem~\ref{r8-thm-XtX}).
\item \(I_n-A\) is symmetric and idempotent, so \(C'C=B(I_n-A)^2B=B(I_n-A)B=B^2-BAB=-B(A-B)B\), as \(B^2=B=B^3\); and
\(B(A-B)B=B'(A-B)B\) is PSD (Theorem~\ref{r8-thm-congruence}(a)). So \(0\succcurlyeq C'C\).
\item Hence \(C'C=0\) (Proposition~\ref{r8-prop-loewner}(a)), so \(\tr(C'C)=0\) and \(C=0\) (Theorem~\ref{r1-thm-trace}(d)): \(B=AB\).
\end{itemize}
\end{solution}

\begin{problem}[\citeat{harville-solutions}{Ch.~21, Exercise~32}]\label{r8-ex-simdiag-order}
Let \(A\in\R^{n\times n}\) be symmetric and \(B\in\R^{n\times n}\) PD. Show that \(A-B\) is PSD if and only if the \(n\) roots
of \(\det(A-\lambda B)=0\), listed with multiplicity, are all \(\geqslant1\), and PD if and only if they are all \(>1\).
\end{problem}
\begin{solution}
\begin{itemize}
\item By Theorem~\ref{r8-thm-pd-simdiag}(a),(b), applied to the PD matrix \(B\) and the symmetric matrix \(A\), there is an
invertible \(E\) with \(E'BE=I_n\) and \(E'AE=D=\diag(d_1,\dots,d_n)\), and \(\det(A-\lambda B)=\det B\cdot(d_1-\lambda)\cdots(d_n-\lambda)\):
the roots are \(d_1,\dots,d_n\).
\item So \(E'(A-B)E=D-I_n=\diag(d_1-1,\dots,d_n-1)\).
\item By Theorem~\ref{r8-thm-congruence}(c), \(A-B\) is PSD (PD) if and only if \(D-I_n\) is, that is
(Proposition~\ref{r8-prop-pd-basic}(a)), if and only if every \(d_i\geqslant1\) (\(d_i>1\)).
\end{itemize}
\end{solution}

\begin{problem}[\citeat{prasolov}{Problem~19.3}]\label{r8-ex-pr-rank-trace}
Let \(A\in\R^{n\times n}\) be symmetric, \(A\ne0\). Show that
\[
\rank A\geqslant\frac{(\tr A)^2}{\tr(A^2)}
\]
\end{problem}
\begin{solution}
Let \(A=Q\Lambda Q'\), with \(Q\) orthogonal and \(\Lambda=\diag(\lambda_1,\dots,\lambda_n)\), where
\(\lambda_1,\dots,\lambda_n\) are the eigenvalues of \(A\) (Theorem~\ref{r8-thm-spectral}).
\begin{itemize}
\item Let \(S=\{i:\lambda_i\ne0\}\), and let \(r\) be the number of its elements. Then \(r=\rank A\)
(Theorem~\ref{r8-thm-rank-eigen}), and \(r\geqslant1\), as \(A\ne0\).
\item \(\tr A=\lambda_1+\dots+\lambda_n=\sum_{i\in S}\lambda_i\) (Theorem~\ref{r7-thm-trace-det}).
\item \(A^2=Q\Lambda Q'Q\Lambda Q'=Q\Lambda^2Q'\), so
\(\tr(A^2)=\tr(\Lambda^2Q'Q)=\tr(\Lambda^2)=\sum_i\lambda_i^2=\sum_{i\in S}\lambda_i^2\) (Theorem~\ref{r1-thm-trace}(b),
Proposition~\ref{r1-prop-powers}(d)). Also \(\tr(A^2)=\tr(A'A)>0\), as \(A\ne0\) (Theorem~\ref{r1-thm-trace}(d)).
\item Let \(\mu\in\R^r\) have the entries \(\lambda_i\), \(i\in S\). By the Cauchy--Schwarz inequality
(Theorem~\ref{r6-thm-cauchy-schwarz}(a)) for \(\one_r\) and \(\mu\),
\[
(\tr A)^2=(\one_r'\mu)^2\leqslant(\one_r'\one_r)(\mu'\mu)=r\sum_{i\in S}\lambda_i^2=\rank A\cdot\tr(A^2)
\]
\end{itemize}
Dividing by \(\tr(A^2)>0\) gives the inequality.
\end{solution}

\begin{problem}[\citeat{prasolov}{Problem~33.5}]\label{r8-ex-pr-bordered}
Let \(A\in\R^{n\times n}\) be PD and \(x\in\R^n\). Show that
\[
\det\begin{pmatrix}0&x'\\x&A\end{pmatrix}\leqslant0
\]
\end{problem}
\begin{solution}
Let \(M\) be the \((n+1)\times(n+1)\) matrix in question, partitioned into the blocks \(0\in\R^{1\times1}\),
\(x'\in\R^{1\times n}\), \(x\in\R^{n\times1}\) and \(A\in\R^{n\times n}\).
\begin{itemize}
\item \(A\) is invertible and \(\det A>0\) (Theorem~\ref{r8-thm-eigen-criterion}(c)). By
Theorem~\ref{r5d-thm-schur-det}(b), with \(A\) as the invertible lower right block,
\[
\det M=\det A\cdot\det\bigl(0-x'A^{-1}x\bigr)=-\det A\cdot x'A^{-1}x
\]
as the determinant of a \(1\times1\) matrix is its entry (Definition~\ref{r5d-def-det}).
\item \(A^{-1}\) is PD (Proposition~\ref{r8-prop-pd-basic}(c)), so \(x'A^{-1}x\geqslant0\), with equality only for
\(x=0\). Hence \(\det M\leqslant0\), with equality only for \(x=0\).
\end{itemize}
\end{solution}

\begin{problem}[\citeat{prasolov}{Problem~34.3}]\label{r8-ex-pr-trace-product}
Let \(A,B\in\R^{n\times n}\) be PSD. Show that \(\lvert\tr(AB)\rvert\leqslant\tr A\cdot\tr B\).
\end{problem}
\begin{solution}
Let \(N=A^{1/2}\) and \(K=B^{1/2}\) (Theorem~\ref{r8-thm-sqrt}); they are symmetric, with \(N^2=A\) and \(K^2=B\). Let
\(G=NK\).
\begin{itemize}
\item By Theorem~\ref{r1-thm-trace}(b), applied to the factors \(NNK\) and \(K\), and by Theorem~\ref{r1-thm-transpose},
\[
\tr(AB)=\tr(NNKK)=\tr(KNNK)=\tr\bigl((NK)'(NK)\bigr)=\tr(G'G)=\sum_{i=1}^n\sum_{j=1}^ng_{ij}^2
\]
(Theorem~\ref{r1-thm-trace}(d)). In particular \(\tr(AB)\geqslant0\), so \(\lvert\tr(AB)\rvert=\tr(AB)\).
\item Let \(u_i'=\row_i(N)\) and \(v_j=\col_j(K)\). Then \(g_{ij}=u_i'v_j\) (Theorem~\ref{r1-thm-views}(d)), and
\(g_{ij}^2\leqslant\lVert u_i\rVert^2\lVert v_j\rVert^2\) (Theorem~\ref{r6-thm-cauchy-schwarz}(a)).
\item \(\sum_i\lVert u_i\rVert^2\) is the sum of the squares of all entries of \(N\), which is \(\tr(N'N)=\tr(N^2)=\tr A\);
likewise \(\sum_j\lVert v_j\rVert^2=\tr(K'K)=\tr(K^2)=\tr B\) (Theorem~\ref{r1-thm-trace}(d)).
\item Hence
\[
\tr(AB)=\sum_{i,j}g_{ij}^2\leqslant\sum_{i,j}\lVert u_i\rVert^2\lVert v_j\rVert^2
=\Bigl(\sum_i\lVert u_i\rVert^2\Bigr)\Bigl(\sum_j\lVert v_j\rVert^2\Bigr)=\tr A\cdot\tr B
\]
\end{itemize}
\end{solution}

\begin{problem}[\citeat{prasolov}{Problem~33.1}]\label{r8-ex-pr-det-sum}
Let \(A,B\in\R^{n\times n}\), \(n\geqslant2\), with \(A\) PD and \(B\) PSD. Show that
\(\det(A+B)\geqslant\det A+\det B\), with equality only for \(B=0\).
\end{problem}
\begin{solution}
\begin{itemize}
\item By Theorem~\ref{r8-thm-pd-simdiag}(a), there is an invertible \(E\in\R^{n\times n}\) with \(E'AE=I_n\) and
\(E'BE=M=\diag(\mu_1,\dots,\mu_n)\). \(M\) is PSD (Theorem~\ref{r8-thm-congruence}(a)), so every \(\mu_i\geqslant0\)
(Proposition~\ref{r8-prop-pd-basic}(a)).
\item Let \(\delta=(\det E)^2\), which is \(>0\) (Theorem~\ref{r5d-thm-product}(a)). As
\(E'(A+B)E=I_n+M=\diag(1+\mu_1,\dots,1+\mu_n)\), Theorems~\ref{r5d-thm-product}(b), \ref{r5d-thm-transpose}
and~\ref{r5d-thm-basic}(c) give
\[
\delta\det A=1,\qquad\delta\det(A+B)=(1+\mu_1)\cdots(1+\mu_n),\qquad\delta\det B=\mu_1\cdots\mu_n
\]
\item For \(k\geqslant2\) and numbers \(\mu_1,\dots,\mu_k\geqslant0\),
\[
(1+\mu_1)\cdots(1+\mu_k)\geqslant1+(\mu_1+\dots+\mu_k)+\mu_1\cdots\mu_k
\]
by induction on \(k\). For \(k=2\) both sides are \(1+\mu_1+\mu_2+\mu_1\mu_2\). If it holds for \(k\), let
\(s=\mu_1+\dots+\mu_k\) and \(\pi=\mu_1\cdots\mu_k\); multiplying by \(1+\mu_{k+1}>0\),
\[
(1+\mu_1)\cdots(1+\mu_{k+1})\geqslant(1+s+\pi)(1+\mu_{k+1})=1+(s+\mu_{k+1})+\pi\mu_{k+1}+\pi+s\mu_{k+1}
\]
which is \(\geqslant1+(s+\mu_{k+1})+\pi\mu_{k+1}\), as \(\pi\geqslant0\) and \(s\mu_{k+1}\geqslant0\).
\item With \(k=n\geqslant2\) and the three equations above,
\[
\delta\det(A+B)\geqslant1+(\mu_1+\dots+\mu_n)+\delta\det B\geqslant\delta\det A+\delta\det B
\]
as \(\delta\det A=1\) and every \(\mu_i\geqslant0\); dividing by \(\delta>0\) gives \(\det(A+B)\geqslant\det A+\det B\).
\item If equality holds, then \(1+\delta\det B=\delta\det(A+B)\geqslant1+(\mu_1+\dots+\mu_n)+\delta\det B\), so
\(\mu_1+\dots+\mu_n\leqslant0\); as every \(\mu_i\geqslant0\), every \(\mu_i=0\). So \(M=0\), and
\(B=(E')^{-1}ME^{-1}=0\) (Theorem~\ref{r2-thm-inverse-props}(c)). Conversely, if \(B=0\), then
\(\det(A+B)=\det A=\det A+\det B\), as \(\det0=0\) (Theorem~\ref{r5d-thm-basic}(b)).
\end{itemize}
\end{solution}

\clearpage
\def\unitnumber{9}
\setcounter{theorem}{0}\setcounter{problem}{0}
\section*{Chapter 9. Covariance matrices and principal components}

\subsection*{Random vectors and covariance matrices}

A \textbf{random variable} is a number whose value depends on chance; every random variable in this book is assumed to
have a finite variance. We use the following facts from probability without proof. For random variables
\(u,v\) and numbers \(a,b,c\), the \textbf{expectation} satisfies
\(\expc(au+bv+c)=a\expc(u)+b\expc(v)+c\); so the expectation of a sum of finitely many
terms is the sum of their expectations. The \textbf{variance}
\(\var(u)=\expc\bigl((u-\expc(u))^2\bigr)\) is \(\geqslant0\), with
\(\var(u)=0\) if and only if \(u=\expc(u)\) with probability \(1\); and the \textbf{covariance} is
\(\cov(u,v)=\expc\bigl((u-\expc(u))(v-\expc(v))\bigr)=\cov(v,u)\);
thus \(\var(u)=\cov(u,u)\)~\citeat{searle-khuri}{\S10.1}, \citeat{tu-clm}{Lemma~3.2}.

\begin{definition}[\citeat{bapat}{\S7.1}, \citeat{searle-khuri}{Ch.~10, introduction, \S10.1}, \citeat{tu-clm}{Lemma~3.2}]\label{r11-def-random}
A \textbf{random vector} \(z=(z_1,\dots,z_n)'\) and a \textbf{random matrix} \(Z=(z_{ij})\) have random variables as
entries, and their expectations are taken entry by entry: \(\expc(Z)=(\expc(z_{ij}))\). In
particular \(\mu=\expc(z)=(\expc(z_1),\dots,\expc(z_n))'\) is the \textbf{mean vector} of
\(z\), and
\[
\Sigma=\var(z)=\expc\bigl((z-\mu)(z-\mu)'\bigr)\in\R^{n\times n}
\]
is its \textbf{covariance matrix}. The \((i,j)\) entry of \((z-\mu)(z-\mu)'\) is \((z_i-\mu_i)(z_j-\mu_j)\), so the
\((i,j)\) entry of \(\Sigma\) is \(\sigma_{ij}=\cov(z_i,z_j)\): the diagonal entries are the variances
\(\sigma_{ii}=\var(z_i)\), and \(\Sigma\) is symmetric.
\end{definition}

\begin{definition}[\citeat{tu-clm}{Lemma~3.2}, \citeat{searle-khuri}{\S10.1}, \citeat{casella-berger}{Definition~4.5.1}]\label{r14-def-cov}
Let \(u\in\R^p\) and \(v\in\R^q\) be random vectors with means \(\mu_u=\expc(u)\) and \(\mu_v=\expc(v)\)
(Definition~\ref{r11-def-random}). The \textbf{covariance matrix of \(u\) and \(v\)} is
\[
\cov(u,v)=\expc\bigl((u-\mu_u)(v-\mu_v)'\bigr)\in\R^{p\times q}
\]
Its \((i,j)\) entry is \(\expc\bigl((u_i-\expc(u_i))(v_j-\expc(v_j))\bigr)=\cov(u_i,v_j)\).
So \(\cov(u,u)=\var(u)\), and \(\cov(v,u)=\cov(u,v)'\), since
\((v-\mu_v)(u-\mu_u)'=\bigl((u-\mu_u)(v-\mu_v)'\bigr)'\) (Theorem~\ref{r1-thm-transpose}). The vectors \(u\) and \(v\) are
\textbf{uncorrelated} if \(\cov(u,v)=0\).
\end{definition}

For the stacked vector \(z=(u',v')'\in\R^{p+q}\), whose mean is \((\mu_u',\mu_v')'\), the block product
(Theorem~\ref{r5b-thm-block}(b),(c)) gives
\[
(z-\expc(z))(z-\expc(z))'=\begin{pmatrix}u-\mu_u\\v-\mu_v\end{pmatrix}\bigl((u-\mu_u)'\ \ (v-\mu_v)'\bigr)
=\begin{pmatrix}(u-\mu_u)(u-\mu_u)'&(u-\mu_u)(v-\mu_v)'\\(v-\mu_v)(u-\mu_u)'&(v-\mu_v)(v-\mu_v)'\end{pmatrix}
\]
and, as expectations are taken entry by entry,
\[
\var(z)=\begin{pmatrix}\var(u)&\cov(u,v)\\\cov(v,u)&\var(v)\end{pmatrix}
\]

\begin{theorem}[\citeat{bapat}{\S7.1}, \citeat{tu-clm}{Lemma~3.2}, \citeat{searle-khuri}{\S10.1}]\label{r11-thm-linear}
Let \(z\in\R^n\) be a random vector with \(\expc(z)=\mu\) and \(\var(z)=\Sigma\).
\begin{enumerate}
\item For a random matrix \(Z\in\R^{m\times p}\) and constant matrices \(A\in\R^{q\times m}\), \(B\in\R^{p\times l}\) and
\(C\in\R^{q\times l}\), \(\expc(AZB+C)=A\expc(Z)B+C\).
\item For a constant matrix \(A\in\R^{m\times n}\) and a constant vector \(b\in\R^m\),
\[
\expc(Az+b)=A\mu+b,\qquad\var(Az+b)=A\Sigma A'
\]
\item \(\Sigma\) is positive semidefinite: \(a'\Sigma a=\var(a'z)\geqslant0\) for every \(a\in\R^n\).
\end{enumerate}
\end{theorem}
\begin{proof}
(a) By Definition~\ref{r1-def-product}, applied twice, the \((i,j)\) entry of \(AZB+C\) is
\(\sum_{r=1}^m\sum_{s=1}^pa_{ir}z_{rs}b_{sj}+c_{ij}\), a sum of the random variables \(z_{rs}\) times numbers, plus a
number. Its expectation is \(\sum_r\sum_sa_{ir}\expc(z_{rs})b_{sj}+c_{ij}\), the \((i,j)\) entry of
\(A\expc(Z)B+C\).

\needspace{4\baselineskip}
(b)
\begin{itemize}
\item (a) with \(Z=z\), \(B=1\) and \(C=b\) gives \(\expc(Az+b)=A\mu+b\).
\item Hence \(Az+b-\expc(Az+b)=A(z-\mu)\), whose transpose is \((z-\mu)'A'\) (Theorem~\ref{r1-thm-transpose}).
By Definition~\ref{r11-def-random} and (a),
\(\var(Az+b)=\expc\bigl(A(z-\mu)(z-\mu)'A'\bigr)=A\expc\bigl((z-\mu)(z-\mu)'\bigr)A'=A\Sigma A'\).
\end{itemize}
(c) \(\Sigma\) is symmetric (Definition~\ref{r11-def-random}). For \(a\in\R^n\), (b) with the \(1\times n\) matrix \(a'\)
and \(b=0\) gives \(\var(a'z)=a'\Sigma a\), and a variance is \(\geqslant0\). So \(\Sigma\) is positive
semidefinite (Definition~\ref{r8-def-definite}).
\end{proof}

\begin{theorem}[\citeat{tu-clm}{Lemma~3.2}, \citeat{casella-berger}{Theorem~4.5.6, Lemma~5.2.5, Theorem~4.5.3}]\label{r14-thm-cov}
Let \(u\in\R^p\) and \(v\in\R^q\) be random vectors.
\begin{enumerate}
\item For constant \(A\in\R^{k\times p}\), \(b\in\R^k\), \(C\in\R^{l\times q}\) and \(d\in\R^l\),
\(\cov(Au+b,\,Cv+d)=A\cov(u,v)C'\). In particular, if \(z\in\R^n\) is a random vector with
\(\var(z)=\Sigma\), and \(A\in\R^{k\times n}\) and \(B\in\R^{l\times n}\) are constant, then
\(\cov(Az,Bz)=A\Sigma B'\).
\item If \(u_1,u_2\in\R^p\) are random vectors, then \(\cov(u_1+u_2,v)=\cov(u_1,v)+\cov(u_2,v)\),
and likewise \(\cov(v,u_1+u_2)=\cov(v,u_1)+\cov(v,u_2)\).
\item If \(q=p\), then \(\var(u+v)=\var(u)+\cov(u,v)+\cov(v,u)+\var(v)\).
\item If \(x_1,\dots,x_m\in\R^p\) are random vectors with \(\cov(x_i,x_j)=0\) for \(i\ne j\), then
\(\var(x_1+\dots+x_m)=\var(x_1)+\dots+\var(x_m)\).
\item \(\var(u)=\expc(uu')-\mu_u\mu_u'\), that is,
\(\expc(uu')=\var(u)+\mu_u\mu_u'\), where \(\mu_u=\expc(u)\).
\end{enumerate}
\end{theorem}
\begin{proof}
(a)
\begin{itemize}
\item By Theorem~\ref{r11-thm-linear}(b), \(Au+b-\expc(Au+b)=A(u-\mu_u)\) and \(Cv+d-\expc(Cv+d)=C(v-\mu_v)\), whose
transpose is \((v-\mu_v)'C'\) (Theorem~\ref{r1-thm-transpose}).
\item By Theorem~\ref{r11-thm-linear}(a),
\(\cov(Au+b,Cv+d)=\expc\bigl(A(u-\mu_u)(v-\mu_v)'C'\bigr)=A\cov(u,v)C'\). For \(u=v=z\) this reads
\(\cov(Az,Bz)=A\var(z)B'\).
\end{itemize}
(b) The mean of \(u_1+u_2\) is \(\expc(u_1)+\expc(u_2)\), so
\((u_1+u_2-\expc(u_1+u_2))(v-\mu_v)'=(u_1-\expc(u_1))(v-\mu_v)'+(u_2-\expc(u_2))(v-\mu_v)'\); take
expectations entry by entry. The second identity follows by transposing (Definition~\ref{r14-def-cov}).

(c) By (b), applied in both arguments,
\(\var(u+v)=\cov(u+v,u+v)=\cov(u,u)+\cov(u,v)+\cov(v,u)+\cov(v,v)\).

(d) In the same way, \(\var\bigl(\sum_ix_i\bigr)=\cov\bigl(\sum_ix_i,\sum_jx_j\bigr)=\sum_i\sum_j\cov(x_i,x_j)\),
and only the \(m\) terms with \(i=j\) remain.

(e) \((u-\mu_u)(u-\mu_u)'=uu'-u\mu_u'-\mu_uu'+\mu_u\mu_u'\) (Theorems~\ref{r1-thm-rules} and~\ref{r1-thm-transpose});
taking expectations entry by entry, with \(\expc(u\mu_u')=\mu_u\mu_u'=\expc(\mu_uu')\)
(Theorem~\ref{r11-thm-linear}(a)),
\(\var(u)=\expc(uu')-\mu_u\mu_u'-\mu_u\mu_u'+\mu_u\mu_u'=\expc(uu')-\mu_u\mu_u'\).
\end{proof}

\begin{theorem}[\citeat{bapat}{\S7.1}]\label{r11-thm-cov-psd}
Let \(x\in\R^n\) be a random vector with \(\expc(x)=\mu\) and \(\var(x)=\Sigma\). Then \(\Sigma\) is
singular (not invertible) if and only if \(a'x=a'\mu\) with probability \(1\) for some nonzero \(a\in\R^n\).
\end{theorem}
\begin{proof}
\begin{itemize}
\item \(\Sigma\) is singular if and only if \(\Sigma a=0\) for some \(a\ne0\) (Theorem~\ref{r2-thm-invertible}).
\item \(\Sigma\) is PSD (Theorem~\ref{r11-thm-linear}(c)), so \(\Sigma a=0\) if and only if \(a'\Sigma a=0\)
(Corollary~\ref{r8-cor-psd-null}).
\item By Theorem~\ref{r11-thm-linear}(b), \(a'\Sigma a=\var(a'x)\) and \(\expc(a'x)=a'\mu\); and a
random variable has variance \(0\) if and only if it equals its expectation with probability \(1\).\qedhere
\end{itemize}
\end{proof}

\begin{example}\label{r11-exm-singular}
Let \(u_1,u_2,u_3\) be pairwise uncorrelated random variables with variance \(1\), and
\(x=(u_1,\ u_2,\ u_3,\ u_1+u_2+u_3)'\). Then \(x=Bu\) with \(u=(u_1,u_2,u_3)'\), \(\var(u)=I_3\)
(Definition~\ref{r11-def-random}) and
\[
B=\begin{pmatrix}1&0&0\\0&1&0\\0&0&1\\1&1&1\end{pmatrix},\qquad
\var(x)=BB'=\begin{pmatrix}1&0&0&1\\0&1&0&1\\0&0&1&1\\1&1&1&3\end{pmatrix}
\]
by Theorem~\ref{r11-thm-linear}(b). This matrix is singular: \(\var(x)a=0\) for \(a=(1,1,1,-1)'\).
Correspondingly \(a'x=u_1+u_2+u_3-(u_1+u_2+u_3)=0\), as Theorem~\ref{r11-thm-cov-psd} says. Likewise, the shares of a
budget spent on \(p\) items add up to \(1\), so their covariance matrix is singular.
\end{example}

\begin{definition}[\citeat{casella-berger}{Definitions~2.3.2, 4.5.2}, \citeat{searle-khuri}{\S10.2}]\label{r14-def-corr}
The \textbf{correlation} of random variables \(u\) and \(v\) with positive variances is
\(\rho_{uv}=\cov(u,v)/\sqrt{\var(u)\var(v)}\). Let \(z\in\R^n\) be a random vector with
\(\expc(z)=\mu\), \(\var(z)=\Sigma=(\sigma_{ij})\) and \(\sigma_{ii}>0\) for every \(i\); \(\sqrt{\sigma_{ii}}\)
is the \textbf{standard deviation} of \(z_i\). The correlation of \(z_i\) and \(z_j\) is
\(\rho_{ij}=\sigma_{ij}/\sqrt{\sigma_{ii}\sigma_{jj}}\), and with
\(D^{-1/2}=\diag\bigl(\sigma_{11}^{-1/2},\dots,\sigma_{nn}^{-1/2}\bigr)\), the \textbf{correlation matrix} of \(z\) is
\[
\cor(z)=(\rho_{ij})=D^{-1/2}\Sigma D^{-1/2}\in\R^{n\times n}
\]
by Definition~\ref{r1-def-product}: the \((i,j)\) entry of \(D^{-1/2}\Sigma D^{-1/2}\) is
\(\sigma_{ii}^{-1/2}\sigma_{ij}\sigma_{jj}^{-1/2}\).
\end{definition}

\begin{proposition}[\citeat{casella-berger}{Theorem~4.5.7(a)}, \citeat{searle-khuri}{\S10.2}, \citeat{haerdle-simar}{\S12.1}]\label{r14-prop-corr}
In Definition~\ref{r14-def-corr}:
\begin{enumerate}
\item \(-1\leqslant\rho_{ij}\leqslant1\) for all \(i,j\);
\item \(\cor(z)=\var\bigl(D^{-1/2}(z-\mu)\bigr)\) is the covariance matrix of the
\textbf{standardized variables} \((z_i-\mu_i)/\sqrt{\sigma_{ii}}\) (\(i=1,\dots,n\)), which have mean \(0\) and variance
\(1\); so \(\cor(z)\) is symmetric and PSD, with diagonal entries \(1\).
\end{enumerate}
\end{proposition}
\begin{proof}
(a)
\begin{itemize}
\item For \(t\in\R\), Theorem~\ref{r11-thm-linear}(b) for \((z_i,z_j)'\), whose covariance matrix has the entries
\(\sigma_{ii},\sigma_{ij},\sigma_{ij},\sigma_{jj}\), and the \(1\times2\) matrix \((t,1)\) gives
\(h(t)=\var(tz_i+z_j)=\sigma_{ii}t^2+2\sigma_{ij}t+\sigma_{jj}\geqslant0\).
\item A quadratic polynomial with leading coefficient \(\sigma_{ii}>0\) and two different real roots \(t_1<t_2\) equals
\(\sigma_{ii}(t-t_1)(t-t_2)<0\) for \(t_1<t<t_2\). So \(h\) has no two different real roots, and its discriminant is
\((2\sigma_{ij})^2-4\sigma_{ii}\sigma_{jj}\leqslant0\).
\item Hence \(\sigma_{ij}^2\leqslant\sigma_{ii}\sigma_{jj}\), that is, \(\rho_{ij}^2\leqslant1\).
\end{itemize}
(b) \(D^{-1/2}\) is diagonal, hence symmetric. By Theorem~\ref{r11-thm-linear}(b) with \(A=D^{-1/2}\) and
\(b=-D^{-1/2}\mu\), \(D^{-1/2}(z-\mu)\) has the mean \(D^{-1/2}\mu-D^{-1/2}\mu=0\) and
\(\var\bigl(D^{-1/2}(z-\mu)\bigr)=D^{-1/2}\Sigma D^{-1/2}=\cor(z)\). This is symmetric and PSD
(Theorem~\ref{r11-thm-linear}(c)), and its \((i,i)\) entry, the variance of \((z_i-\mu_i)/\sqrt{\sigma_{ii}}\), is
\(\sigma_{ii}/\sigma_{ii}=1\).
\end{proof}

\begin{example}\label{r14-exm-cov}
Let \(z\in\R^4\) be a random vector with
\[
\var(z)=\Sigma=\begin{pmatrix}4&0&4&0\\0&1&2&2\\4&2&9&6\\0&2&6&9\end{pmatrix}
\]
\begin{itemize}
\item The standard deviations are \(\sqrt{\sigma_{11}}=2\), \(\sqrt{\sigma_{22}}=1\), \(\sqrt{\sigma_{33}}=3\) and
\(\sqrt{\sigma_{44}}=3\), so \(\rho_{12}=\rho_{14}=0\), \(\rho_{13}=\frac4{2\cdot3}=\frac23\), \(\rho_{23}=\frac2{1\cdot3}=\frac23\),
\(\rho_{24}=\frac2{1\cdot3}=\frac23\), \(\rho_{34}=\frac6{3\cdot3}=\frac23\), and
\[
\cor(z)=\begin{pmatrix}1&0&\frac23&0\\[2pt]0&1&\frac23&\frac23\\[2pt]\frac23&\frac23&1&\frac23\\[2pt]0&\frac23&\frac23&1\end{pmatrix}
\]
So \(z_1\) is uncorrelated with \(z_2\) and with \(z_4\), and the other four pairs have the correlation \(\frac23\).
\item Let \(e=z_3-z_1-2z_2=a'z\) with \(a=(-1,-2,1,0)'\). Then
\(\Sigma a=(-4+0+4+0,\ 0-2+2+0,\ -4-4+9+0,\ 0-4+6+0)'=(0,0,1,2)'\), so by Theorem~\ref{r14-thm-cov}(a),
\(\cov(z,e)=\Sigma a=(0,0,1,2)'\): \(e\) is uncorrelated with \(z_1\) and with \(z_2\), and
\(\var(e)=a'\Sigma a=0+0+1+0=1\). Likewise \(f=z_4+2z_1+2z_2-2z_3=b'z\) with \(b=(2,2,-2,1)'\) has
\(\Sigma b=(8+0-8+0,\ 0+2-4+2,\ 8+4-18+6,\ 0+4-12+9)'=(0,0,0,1)'\): \(f\) is uncorrelated with \(z_1\), \(z_2\) and
\(z_3\), and \(\var(f)=b'\Sigma b=1\).
\item By Theorem~\ref{r14-thm-cov}(a), \(z_1+z_2+z_3+z_4=\one'z\) has the variance \(\one'\Sigma\one=51\), the sum of all
entries of \(\Sigma\). Likewise, as \((1,1,0,0)\Sigma=(4,1,6,2)\) and \((1,-1,0,0)\Sigma=(4,-1,2,-2)\), the variables
\(z_1+z_2\) and \(z_1-z_2\) have the variances \(4+1=5\) and \(4+1=5\), and the covariance \((4,1,6,2)(1,-1,0,0)'=4-1=3\).
\end{itemize}
\end{example}

\begin{theorem}[\citeat{tu-clm}{Lemma~3.6}, \citeat{searle-khuri}{(10.34)}]\label{r11-thm-quadratic}
Let \(z\in\R^n\) be a random vector with \(\expc(z)=\mu\) and \(\var(z)=\Sigma\), and let
\(A\in\R^{n\times n}\) be constant. Then
\[
\expc(z'Az)=\tr(A\Sigma)+\mu'A\mu
\]
In particular \(\expc(z'Az)=\tr(A\Sigma)\) if \(\mu=0\).
\end{theorem}
\begin{proof}
\begin{itemize}
\item \(z'Az=\sum_{i=1}^n\sum_{j=1}^na_{ij}z_iz_j\) (Definition~\ref{r8-def-quadratic-form}), so
\(\expc(z'Az)=\sum_i\sum_ja_{ij}\expc(z_iz_j)\).
\item Expanding the product,
\(\sigma_{ij}=\expc\bigl((z_i-\mu_i)(z_j-\mu_j)\bigr)=\expc(z_iz_j)-\mu_i\mu_j-\mu_j\mu_i+\mu_i\mu_j\); so
\(\expc(z_iz_j)=\sigma_{ij}+\mu_i\mu_j\), and
\(\expc(z'Az)=\sum_i\sum_ja_{ij}\sigma_{ij}+\sum_i\sum_ja_{ij}\mu_i\mu_j\).
\item \(\tr(A\Sigma)=\sum_i(A\Sigma)_{ii}=\sum_i\sum_ja_{ij}\sigma_{ji}=\sum_i\sum_ja_{ij}\sigma_{ij}\), as \(\Sigma\) is
symmetric; and \(\mu'A\mu=\sum_i\sum_ja_{ij}\mu_i\mu_j\) (Definition~\ref{r8-def-quadratic-form}).\qedhere
\end{itemize}
\end{proof}

\begin{definition}[\citeat{casella-berger}{Definitions~7.1.1, 7.3.2}]\label{r11-def-estimator}
Let observations be modeled as a random vector \(y\) whose distribution depends on unknown parameters. An
\textbf{estimator} of a parameter \(\theta\) is a random variable \(T=g(y)\), where the function \(g\) does not involve
unknown parameters; its value at the observed \(y\) is an \textbf{estimate} of \(\theta\). The estimator \(T\) is
\textbf{unbiased} if \(\expc(T)=\theta\) for all values of the parameters.
\end{definition}

\begin{example}[cf.\ \citeat{casella-berger}{Theorem~5.2.6(c)}]\label{r11-exm-quadratic-projection}
Let \(y\in\R^n\) have \(\expc(y)=\mu\) and \(\var(y)=\sigma^2I_n\), and let \(P\in\R^{n\times n}\) be
symmetric and idempotent of rank \(r\). Then \(\tr P=\rank P=r\) (Theorem~\ref{r6-thm-idempotent}(c)), and
Theorem~\ref{r11-thm-quadratic}, with \(\tr(P\sigma^2I_n)=\sigma^2\tr P\) (Theorem~\ref{r1-thm-trace}(a)), gives
\[
\expc(y'Py)=\tr(P\sigma^2I_n)+\mu'P\mu=\sigma^2r+\mu'P\mu
\]
\begin{itemize}
\item Let \(n\geqslant2\) and \(P=C_n\), which is symmetric and idempotent of rank \(n-1\)
(Proposition~\ref{r1-prop-centering}, Corollary~\ref{r6-cor-centering}(b)). By Example~\ref{r1-exm-centering},
\(y'C_ny=\sum_i(y_i-\bar y)^2\) and \(\mu'C_n\mu=\sum_i(\mu_i-\bar\mu)^2\), where \(\bar\mu=\frac1n\one'\mu\). So
\[
\expc\Bigl(\sum_{i=1}^n(y_i-\bar y)^2\Bigr)=(n-1)\sigma^2+\sum_{i=1}^n(\mu_i-\bar\mu)^2
\]
\item If the observations have a common mean, \(\mu=\mu_0\one\), then \(C_n\mu=0\) and the last sum is \(0\). So
\(s^2=\frac1{n-1}\sum_i(y_i-\bar y)^2\) is an unbiased estimator of \(\sigma^2\) (Definition~\ref{r11-def-estimator}):
\(\expc(s^2)=\sigma^2\) for all values of \(\mu_0\) and \(\sigma^2\). The estimator
\(\frac1n\sum_i(y_i-\bar y)^2\) has the expectation \(\frac{n-1}n\sigma^2\), which is smaller than \(\sigma^2\) when
\(\sigma^2>0\); so it is not unbiased.
\item If the means drift, \(s^2\) overestimates \(\sigma^2\) on average. For \(n=4\) and \(\mu=(1,3,5,7)'\),
\(\bar\mu=4\) and \(\sum_i(\mu_i-\bar\mu)^2=9+1+1+9=20\), so \(\expc(s^2)=\sigma^2+\frac{20}3\).
\end{itemize}
\end{example}

\needspace{18\baselineskip}
\subsection*{Data matrices and the sample covariance matrix}

\begin{definition}[\citeat{searle-khuri}{\S10.3}]\label{r12-def-data}
Let \(X=(x_{ij})\in\R^{n\times p}\) (\(n\geqslant2\)) be a \textbf{data matrix}: \(x_{ij}\) is the value of variable \(j\) on
unit \(i\), and row \(i\) of \(X\) is \(x_i'=(x_{i1},\dots,x_{ip})\), the \(i\)th observation, a point of \(\R^p\) (as for
\(X_0\) in Example~\ref{r1-exm-data}). With \(\one=\one_n\) and \(C_n=I_n-\bar J_n\) (Definition~\ref{r1-def-ones}):
\begin{itemize}
\item the \textbf{mean vector} is \(\bar x=\frac1nX'\one=(\bar x_1,\dots,\bar x_p)'\), where \(\bar x_j=\frac1n\sum_ix_{ij}\);
\item the \textbf{centered data matrix} is \(C_nX=X-\one\bar x'\), with entries \(x_{ij}-\bar x_j\)
(Example~\ref{r1-exm-data});
\item the \textbf{sample covariance matrix} is \(S=(s_{jl})=\frac1{n-1}X'C_nX\in\R^{p\times p}\).
\end{itemize}
For data vectors \(y,w\in\R^n\) with means \(\bar y\), \(\bar w\), the number
\(\frac1{n-1}\sum_{i=1}^n(y_i-\bar y)(w_i-\bar w)\) is their \textbf{sample covariance}, and \(y\), \(w\) are
\textbf{uncorrelated} if it is \(0\); for \(w=y\) it is the \textbf{sample variance} \(\frac1{n-1}\sum_i(y_i-\bar y)^2\) of
\(y\). For \(a\in\R^p\), the vector \(Xa=(a'x_1,\dots,a'x_n)'\) holds the values on the \(n\) units of the linear combination
of the variables with coefficient vector \(a\).
\end{definition}

\begin{proposition}[\citeat{searle-khuri}{\S10.3}, \citeat{haerdle-simar}{\S3.3, (3.22)--(3.24)}, \citeat{tu-asm}{Definition~5.19}, \citeat{bapat}{\S10.1}]\label{r12-prop-data}
Let \(X\in\R^{n\times p}\) be a data matrix, with \(\bar x\), \(C_n\) and \(S\) as in Definition~\ref{r12-def-data}.
\begin{enumerate}
\item \(X'C_nX=(C_nX)'(C_nX)=X'X-n\bar x\bar x'\), and
\(S=\frac1{n-1}\sum_{i=1}^n(x_i-\bar x)(x_i-\bar x)'\). Hence \(s_{jj}\) is the sample variance of variable \(j\) and
\(s_{jl}\) the sample covariance of variables \(j\) and \(l\); \(S\) is symmetric and PSD; \(S=0\) only if all rows of
\(X\) are equal; and \(S\) is PD if and only if \(\rank(C_nX)=p\), which requires \(n\geqslant p+1\).
\item For \(a,b\in\R^p\), \(Xa\) has mean \(a'\bar x\) and centered values \(C_n(Xa)=(C_nX)a\); the sample covariance of
\(Xa\) and \(Xb\) is \(a'Sb\), and the sample variance of \(Xa\) is \(a'Sa\).
\item Let \(A\in\R^{q\times p}\), \(b\in\R^q\) and \(Y=XA'+\one b'\), the data matrix with the rows
\(y_i'=(Ax_i+b)'\). Then \(Y\) has the mean vector \(\bar y=A\bar x+b\), the centered data matrix \(C_nY=C_nXA'\) and the
sample covariance matrix \(S_Y=ASA'\).
\end{enumerate}
\end{proposition}
\begin{proof}
(a)
\begin{itemize}
\item \(C_n'=C_n=C_n^2\) (Proposition~\ref{r1-prop-centering}), so \(X'C_nX=X'C_n'C_nX=(C_nX)'(C_nX)\); and, as
\(X'\one=n\bar x\), \(X'C_nX=X'X-\frac1n(X'\one)(\one'X)=X'X-n\bar x\bar x'\).
\item The \((j,l)\) entry of \((C_nX)'(C_nX)\) is \(\sum_i(x_{ij}-\bar x_j)(x_{il}-\bar x_l)\)
(Theorem~\ref{r1-thm-views}(d)); dividing by \(n-1\) gives the sample covariance.
\item Row \(i\) of \(C_nX=X-\one\bar x'\) is \((x_i-\bar x)'\), so by the outer product expansion
(Theorem~\ref{r1-thm-views}(d)), \((C_nX)'(C_nX)=\sum_i(x_i-\bar x)(x_i-\bar x)'\).
\item \((C_nX)'(C_nX)\) is PSD (Theorem~\ref{r8-thm-XtX}), hence so is \(S\). If \(S=0\), then every
\(s_{jj}=\frac1{n-1}\sum_i(x_{ij}-\bar x_j)^2=0\), so \(x_{ij}=\bar x_j\) for all \(i\), \(j\).
\item By Theorem~\ref{r8-thm-XtX}, \((C_nX)'(C_nX)\), and hence \(S\), is PD if and only if \(\rank(C_nX)=p\). As
\(\rank(C_nX)\leqslant\rank C_n=n-1\) (Theorem~\ref{thm:r4-rank-product}(a), Corollary~\ref{r6-cor-centering}(b)), this
requires \(p\leqslant n-1\).
\end{itemize}
(b)
\begin{itemize}
\item \(\frac1n\one'(Xa)=\bar x'a=a'\bar x\), and \(C_n(Xa)=(C_nX)a\) (Theorem~\ref{r1-thm-rules}).
\item By Example~\ref{r1-exm-centering}, \(C_n(Xa)\) and \(C_n(Xb)\) hold the deviations from the means, so the sample
covariance is \(\frac1{n-1}(C_nXa)'(C_nXb)=a'\bigl[\frac1{n-1}(C_nX)'(C_nX)\bigr]b=a'Sb\), by (a). Take \(b=a\) for the
variance.
\end{itemize}
(c) Row \(i\) of \(XA'\) is \(x_i'A'=(Ax_i)'\) (Theorems~\ref{r1-thm-views}(c) and~\ref{r1-thm-transpose}), and row \(i\)
of \(\one b'\) is \(b'\).
\begin{itemize}
\item \(Y'=AX'+b\one'\), so \(\bar y=\frac1nY'\one=A\bigl(\frac1nX'\one\bigr)+b\,\frac1n\one'\one=A\bar x+b\)
(\(\one'\one=n\), Proposition~\ref{r1-prop-centering}).
\item \(C_n\one=0\) (Proposition~\ref{r1-prop-centering}), so \(C_nY=C_nXA'+(C_n\one)b'=C_nXA'\), and by (a)
\(S_Y=\frac1{n-1}(C_nY)'(C_nY)=A\bigl[\frac1{n-1}(C_nX)'(C_nX)\bigr]A'=ASA'\).\qedhere
\end{itemize}
\end{proof}

\begin{example}\label{r12-exm-data}
Five units (\(n=5\)) are measured on four variables (\(p=4\)), for instance five students on four tests:
\[
X=\begin{pmatrix}7&3&7&3\\9&9&5&7\\5&7&9&9\\3&5&1&1\\1&1&3&5\end{pmatrix},\qquad
\one'X=(25,\ 25,\ 25,\ 25),\qquad
\bar x=\tfrac15X'\one=\begin{pmatrix}5\\5\\5\\5\end{pmatrix}
\]
Subtracting \(\bar x'=(5,5,5,5)\) from every row and using Proposition~\ref{r12-prop-data}(a),
\[
C_5X=\begin{pmatrix}2&-2&2&-2\\4&4&0&2\\0&2&4&4\\-2&0&-4&-4\\-4&-4&-2&0\end{pmatrix},\qquad
(C_5X)'(C_5X)=\begin{pmatrix}40&28&20&12\\28&40&12&20\\20&12&40&28\\12&20&28&40\end{pmatrix}
\]
and
\[
S=\frac14(C_5X)'(C_5X)=\begin{pmatrix}10&7&5&3\\7&10&3&5\\5&3&10&7\\3&5&7&10\end{pmatrix}
\]
For instance \(40=4+16+0+4+16\) is the sum of squares of column \(1\) of \(C_5X\), and \(28=-4+16+0+0+16\) is the
inner product of columns \(1\) and \(2\). So every variable has the sample variance \(10\); the pairs of variables
\(\{1,2\}\) and \(\{3,4\}\) have the sample covariance \(7\), the pairs \(\{1,3\}\) and \(\{2,4\}\) have \(5\), and the pairs
\(\{1,4\}\) and \(\{2,3\}\) have \(3\).
\end{example}

\begin{theorem}[\citeat{casella-berger}{Theorem~5.2.6, Lemma~5.2.5}, \citeat{searle-khuri}{\S10.3}]\label{r14-thm-sample}
Let \(x_1,\dots,x_n\in\R^p\) (\(n\geqslant2\)) be random vectors with \(\expc(x_i)=\mu\) and \(\var(x_i)=\Sigma\)
for every \(i\), and \(\cov(x_i,x_j)=0\) for \(i\ne j\). Let \(X\in\R^{n\times p}\) be the random data matrix with
rows \(x_1',\dots,x_n'\), with mean vector \(\bar x=\frac1nX'\one=\frac1n(x_1+\dots+x_n)\) and sample covariance matrix
\(S=\frac1{n-1}X'C_nX\) (Definition~\ref{r12-def-data}). Then
\begin{enumerate}
\item \(\expc(\bar x)=\mu\) and \(\var(\bar x)=\frac1n\Sigma\);
\item \(\expc(S)=\Sigma\): each sample variance \(s_{jj}\) and each sample covariance \(s_{jl}\) has the expectation
\(\sigma_{jj}\), respectively \(\sigma_{jl}\).
\end{enumerate}
\end{theorem}
\begin{proof}
(a)
\begin{itemize}
\item The columns of \(X'\) are \(x_1,\dots,x_n\), so \(X'\one=x_1+\dots+x_n\) (Theorem~\ref{r1-thm-views}(b)), and
\(\expc(\bar x)=\frac1n(\mu+\dots+\mu)=\mu\).
\item By Theorem~\ref{r11-thm-linear}(b) with \(A=\frac1nI_p\) and Theorem~\ref{r14-thm-cov}(d),
\(\var(\bar x)=\frac1{n^2}\var(x_1+\dots+x_n)=\frac1{n^2}\,n\Sigma=\frac1n\Sigma\).
\end{itemize}
(b)
\begin{itemize}
\item By Theorem~\ref{r14-thm-cov}(e), \(\expc(x_ix_i')=\Sigma+\mu\mu'\) and, by (a),
\(\expc(\bar x\bar x')=\frac1n\Sigma+\mu\mu'\).
\item By Proposition~\ref{r12-prop-data}(a), \(X'C_nX=X'X-n\bar x\bar x'\), and \(X'X=x_1x_1'+\dots+x_nx_n'\) by the outer
product expansion (Theorem~\ref{r1-thm-views}(d)).
\item Hence \(\expc(X'C_nX)=n(\Sigma+\mu\mu')-n\bigl(\frac1n\Sigma+\mu\mu'\bigr)=(n-1)\Sigma\), and
\(\expc(S)=\Sigma\).\qedhere
\end{itemize}
\end{proof}

With the divisor \(n\) in place of \(n-1\), the expectation would be \(\frac{n-1}n\Sigma\), not \(\Sigma\) (if \(\Sigma\ne0\)); this is
why \(S\) has the divisor \(n-1\).

\needspace{24\baselineskip}

\begin{definition}[\citeat{searle-khuri}{\S10.2, Exercise~10.7}, \citeat{tu-clm}{Definition~2.15}, \citeat{tu-asm}{Proposition~5.24}, \citeat{islr}{\S12.2.4}]\label{r12-def-correlation}
For \(u,v\in\R^n\) with means \(\bar u\), \(\bar v\), put
\[
S_{uv}=\sum_{i=1}^n(u_i-\bar u)(v_i-\bar v)=(C_nu)'(C_nv),\qquad S_{uu}=\sum_{i=1}^n(u_i-\bar u)^2=\lVert C_nu\rVert^2
\]
(\(C_nu=u-\bar u\one\) and \(C_n'C_n=C_n\); Example~\ref{r1-exm-centering} and Proposition~\ref{r1-prop-centering}). The vector
\(u\) is \textbf{constant} if \(u=\bar u\one_n\), that is, if \(C_nu=0\), or \(S_{uu}=0\). If \(u\) and \(v\) are not
constant, their \textbf{sample correlation} is
\[
\cor(u,v)=\frac{S_{uv}}{\sqrt{S_{uu}S_{vv}}}
\]
(for random variables, the correlation is \(\rho_{uv}\) of Definition~\ref{r14-def-corr}). Now let \(X\in\R^{n\times p}\) be
a data matrix whose variables all have positive sample variance, \(s_{jj}>0\) (Definition~\ref{r12-def-data}), and put
\[
D_s=\diag(s_{11},\dots,s_{pp}),\qquad D_s^{1/2}=\diag\bigl(\sqrt{s_{11}},\dots,\sqrt{s_{pp}}\bigr),\qquad
D_s^{-1/2}=\diag\Bigl(\frac1{\sqrt{s_{11}}},\dots,\frac1{\sqrt{s_{pp}}}\Bigr)
\]
The number \(\sqrt{s_{jj}}\) is the \textbf{sample standard deviation} of variable \(j\).
\begin{itemize}
\item The number \(r_{jl}=s_{jl}/\sqrt{s_{jj}s_{ll}}\) is the sample correlation of variables \(j\) and \(l\), that is,
\(\cor\) of the columns \(j\) and \(l\) of \(X\) (the divisors \(n-1\) cancel), and
\(R=(r_{jl})=D_s^{-1/2}SD_s^{-1/2}\in\R^{p\times p}\) is the \textbf{sample correlation matrix} (by
Theorem~\ref{r1-thm-views}(a),(c), the diagonal factors scale the rows and the columns of \(S\)).
\item The \textbf{standardized data matrix} is \(X_s=C_nXD_s^{-1/2}\), with the entries \((x_{ij}-\bar x_j)/\sqrt{s_{jj}}\):
each variable is centered and measured in units of its sample standard deviation.
\end{itemize}
Thus \(\cor(u,v)\) is the sample covariance of \(u\) and \(v\) (Definition~\ref{r12-def-data}) divided by the
product of the square roots of their sample variances. The numbers \(\rho_{ij}\) and the matrix \(\cor(z)\) of
Definition~\ref{r14-def-corr} are the population counterparts of the \(r_{jl}\) and of \(R\).
\end{definition}

\begin{proposition}[\citeat{tu-asm}{Proposition~5.24}, \citeat{tu-clm}{Proposition~2.17}]\label{r12-prop-correlation}
Let \(X\), \(S\), \(D_s\), \(R\) and \(X_s\) be as in Definition~\ref{r12-def-correlation}.
\begin{enumerate}
\item \(C_nX_s=X_s\), so the standardized variables have mean \(0\), and the sample covariance matrix of the data matrix
\(X_s\) is \(R\). Hence \(R\) is symmetric and PSD, \(r_{jj}=1\), \(\lvert r_{jl}\rvert\leqslant1\), and \(\tr R=p\).
\item Let \(c_1,\dots,c_p>0\) and \(\Gamma=\diag(c_1,\dots,c_p)\). The data matrix \(X\Gamma\), in which variable \(j\) is
multiplied by \(c_j\) (a change of units), has the sample covariance matrix \(\Gamma S\Gamma\), with the entries
\(c_jc_ls_{jl}\), but the same standardized data matrix \(X_s\) and the same sample correlation matrix \(R\).
\end{enumerate}
\end{proposition}
\needspace{4\baselineskip}
\begin{proof}
(a)
\begin{itemize}
\item \(C_n^2=C_n=C_n'\) and \(C_n\one=0\) (Proposition~\ref{r1-prop-centering}). So \(C_nX_s=C_n^2XD_s^{-1/2}=X_s\), and the
mean vector of \(X_s\) is \(\frac1nX_s'\one=\frac1n(C_nX_s)'\one=\frac1nX_s'C_n\one=0\).
\item By Definition~\ref{r12-def-data} and Proposition~\ref{r12-prop-data}(a), and as \((D_s^{-1/2})'=D_s^{-1/2}\), the sample
covariance matrix of \(X_s\) is
\[
\frac1{n-1}X_s'C_nX_s=\frac1{n-1}X_s'X_s=D_s^{-1/2}\Bigl(\frac1{n-1}(C_nX)'(C_nX)\Bigr)D_s^{-1/2}=D_s^{-1/2}SD_s^{-1/2}=R
\]
\item So \(R\) is symmetric and PSD (Proposition~\ref{r12-prop-data}(a) for \(X_s\)), \(r_{jj}=s_{jj}/s_{jj}=1\), and
\(\tr R=p\).
\item Let \(u_j\) be column \(j\) of \(C_nX\). Then \((n-1)s_{jl}=u_j'u_l\) (Proposition~\ref{r12-prop-data}(a),
Theorem~\ref{r1-thm-views}(d)), so \(r_{jl}=u_j'u_l/(\lVert u_j\rVert\,\lVert u_l\rVert)\), and \(\lvert r_{jl}\rvert\leqslant1\)
by the Cauchy--Schwarz inequality (Theorem~\ref{r6-thm-cauchy-schwarz}(a)).
\end{itemize}
(b)
\begin{itemize}
\item \(X\Gamma\) has the mean vector \(\frac1n\Gamma X'\one=\Gamma\bar x\), the centered data matrix \(C_nX\Gamma\), and the
sample covariance matrix \(\frac1{n-1}\Gamma X'C_nX\Gamma=\Gamma S\Gamma\), whose \((j,l)\) entry is \(c_js_{jl}c_l\)
(Theorem~\ref{r1-thm-views}(a),(c)).
\item Its diagonal entries are \(c_j^2s_{jj}\), so for \(X\Gamma\) the matrices \(D_s^{1/2}\) and \(D_s^{-1/2}\) become
\(\Gamma D_s^{1/2}\) and \(\Gamma^{-1}D_s^{-1/2}\) (\(c_j>0\)). Hence the standardized data matrix of \(X\Gamma\) is
\(C_nX\Gamma\Gamma^{-1}D_s^{-1/2}=X_s\), and by (a) its sample correlation matrix, the sample covariance matrix of \(X_s\),
is \(R\).\qedhere
\end{itemize}
\end{proof}

\subsection*{The Mahalanobis distance}

\begin{definition}[\citeat{haerdle-simar}{\S2.6, (2.32), (2.33), Theorem~14.2(a)}]\label{r8-def-mahalanobis}
Let \(S\in\R^{p\times p}\) be PD. The \textbf{Mahalanobis distance} (with respect to \(S\)) between \(x,y\in\R^p\) is
\[
d_S(x,y)=\bigl((x-y)'S^{-1}(x-y)\bigr)^{1/2}
\]
\end{definition}
\(S^{-1}\) is PD (Proposition~\ref{r8-prop-pd-basic}(c)), so \(d_S(x,y)\geqslant0\), with equality only for \(x=y\). For
\(S=I_p\) this is the distance \(\lVert x-y\rVert\) of Definition~\ref{r6-def-inner}; for
\(S=\diag(s_{11},\dots,s_{pp})\) with every \(s_{jj}>0\), \(S^{-1}=\diag(1/s_{11},\dots,1/s_{pp})\)
(Definition~\ref{r2-def-inverse}) and \(d_S(x,y)^2=\sum_j(x_j-y_j)^2/s_{jj}\): each coordinate is measured in units of
\(\sqrt{s_{jj}}\), the sample standard deviation when \(S\) is a sample covariance matrix
(Definition~\ref{r12-def-correlation}).

\begin{proposition}[\citeat{haerdle-simar}{(3.23)--(3.26), Example~6.2, Theorem~2.7}]\label{r8-prop-whitening}
Let \(X\in\R^{n\times p}\) be a data matrix with rows \(x_1',\dots,x_n'\), mean vector \(\bar x\) and a PD sample
covariance matrix \(S\) (Definition~\ref{r12-def-data}).
\begin{enumerate}
\item The vectors \(w_i=S^{-1/2}(x_i-\bar x)\) (\(i=1,\dots,n\)) have the mean vector \(0\) and the sample covariance matrix
\(I_p\). (Passing from the \(x_i\) to the \(w_i\) is called \textbf{whitening} the data.)
\item \(d_S(x_i,\bar x)=\lVert w_i\rVert\) for every \(i\), and \(\sum_{i=1}^nd_S(x_i,\bar x)^2=(n-1)p\).
\item Let \(B\in\R^{p\times p}\) be invertible, \(b\in\R^p\), and \(y_i=Bx_i+b\) (\(i=1,\dots,n\)). The \(y_i\) have the
mean vector \(\bar y=B\bar x+b\) and the PD sample covariance matrix \(S_Y=BSB'\), and \(d_{S_Y}(y_i,y_l)=d_S(x_i,x_l)\)
and \(d_{S_Y}(y_i,\bar y)=d_S(x_i,\bar x)\) for all \(i,l\).
\item Let \(S=Q\Lambda Q'\) as in Theorem~\ref{r8-thm-spectral}(b), with \(\lambda_1\geqslant\dots\geqslant\lambda_p\); then
\(\lambda_p>0\) (Theorem~\ref{r8-thm-eigen-criterion}(a)). Let \(c>0\). For \(x\in\R^p\) put \(u=Q'(x-\bar x)\), so that
\(x=\bar x+u_1q_1+\dots+u_pq_p\). Then
\[
d_S(x,\bar x)^2=\frac{u_1^2}{\lambda_1}+\dots+\frac{u_p^2}{\lambda_p}
\]
So \(\{x\in\R^p:d_S(x,\bar x)\leqslant c\}\) is an ellipse for \(p=2\) (an interval for \(p=1\), an \textbf{ellipsoid} for
\(p\geqslant3\)) with center \(\bar x\) and axes along \(q_1,\dots,q_p\): the point \(\bar x+tq_j\) lies in it if and only
if \(\lvert t\rvert\leqslant c\sqrt{\lambda_j}\).
\end{enumerate}
\end{proposition}
\begin{proof}
(a) \(S^{-1/2}\) is symmetric (Theorem~\ref{r8-thm-sqrt}(a),(c) and Theorem~\ref{r2-thm-inverse-props}(c)), so the
matrix \(X_w\) with the rows \(w_i'=(x_i-\bar x)'S^{-1/2}\) is \(X_w=C_nXS^{-1/2}=XS^{-1/2}-\one(S^{-1/2}\bar x)'\). By
Proposition~\ref{r12-prop-data}(c) with \(A=S^{-1/2}\) and \(b=-S^{-1/2}\bar x\), its mean vector is
\(S^{-1/2}\bar x-S^{-1/2}\bar x=0\) and its sample covariance matrix is \(S^{-1/2}SS^{-1/2}=I_p\)
(Theorem~\ref{r8-thm-sqrt}(c)).

(b)
\begin{itemize}
\item \(\lVert w_i\rVert^2=(x_i-\bar x)'S^{-1/2}S^{-1/2}(x_i-\bar x)=(x_i-\bar x)'S^{-1}(x_i-\bar x)\)
(Theorem~\ref{r8-thm-sqrt}(c)).
\item \((x_i-\bar x)'S^{-1}(x_i-\bar x)=\tr\bigl(S^{-1}(x_i-\bar x)(x_i-\bar x)'\bigr)\) (Theorem~\ref{r1-thm-trace}(b)); summing
over \(i\) (Theorem~\ref{r1-thm-trace}(a)) and by Proposition~\ref{r12-prop-data}(a),
\(\sum_id_S(x_i,\bar x)^2=\tr\bigl(S^{-1}(n-1)S\bigr)=\tr\bigl((n-1)I_p\bigr)=(n-1)p\).
\end{itemize}
(c)
\begin{itemize}
\item The matrix with the rows \(y_i'\) is \(Y=XB'+\one b'\). By Proposition~\ref{r12-prop-data}(c),
\(\bar y=B\bar x+b\) and \(S_Y=BSB'\), which is PD by Theorem~\ref{r8-thm-congruence}(c) with \(P=B'\).
\item \(S_Y^{-1}=(B')^{-1}S^{-1}B^{-1}\) (Theorem~\ref{r2-thm-inverse-props}(b),(c)) and \(y_i-y_l=B(x_i-x_l)\), so
\[
(y_i-y_l)'S_Y^{-1}(y_i-y_l)=(x_i-x_l)'B'(B')^{-1}S^{-1}B^{-1}B(x_i-x_l)=(x_i-x_l)'S^{-1}(x_i-x_l)
\]
and in the same way with \(y_i-\bar y=B(x_i-\bar x)\).
\end{itemize}
(d)
\begin{itemize}
\item \(QQ'=I_p\) (Proposition~\ref{r6-prop-orthonormal}(b)) gives \(x-\bar x=Qu=\sum_ju_jq_j\), and
\(S^{-1}=Q\Lambda^{-1}Q'\) (Proposition~\ref{r8-prop-function}(a),(b) with \(f(t)=t\)), so
\(d_S(x,\bar x)^2=(x-\bar x)'Q\Lambda^{-1}Q'(x-\bar x)=u'\Lambda^{-1}u=\sum_ju_j^2/\lambda_j\).
\item For \(x=\bar x+tq_j\), \(u=tQ'q_j=te_j\), so \(d_S(x,\bar x)^2=t^2/\lambda_j\), which is \(\leqslant c^2\) if and only
if \(\lvert t\rvert\leqslant c\sqrt{\lambda_j}\).\qedhere
\end{itemize}
\end{proof}
In the coordinates \(u\), the set in (d) is \(\bigl\{u:\sum_ju_j^2/(c\sqrt{\lambda_j})^2\leqslant1\bigr\}\), the standard
form of an ellipse or ellipsoid with the semi-axes \(c\sqrt{\lambda_1},\dots,c\sqrt{\lambda_p}\).

\begin{example}\label{r8-exm-mahalanobis}
Nine students take four tests. Their scores \(x_i'\) are the rows of
\[
X=\begin{pmatrix}8&9&10&11\\5&4&5&6\\3&6&5&6\\3&4&7&6\\7&8&9&8\\6&5&6&9\\4&7&6&9\\4&5&8&9\\5&6&7&8\end{pmatrix}
\]
Find \(\bar x\), \(S\) and \(R\); the sample variances of the total score and of the difference of tests \(1\) and
\(2\), and their sample covariance; \(S^{1/2}\), \(S^{-1/2}\), the whitened data and the Mahalanobis distances of the
students from \(\bar x\).

\textit{Solution.}
\begin{itemize}
\item \(\one'X=(45,\ 54,\ 63,\ 72)\), so \(\bar x=(5,6,7,8)'\), and
\[
C_9X=X-\one\bar x'=\begin{pmatrix}3&3&3&3\\0&-2&-2&-2\\-2&0&-2&-2\\-2&-2&0&-2\\2&2&2&0\\1&-1&-1&1\\-1&1&-1&1\\-1&-1&1&1\\0&0&0&0\end{pmatrix},\qquad
(C_9X)'(C_9X)=\begin{pmatrix}24&16&16&16\\16&24&16&16\\16&16&24&16\\16&16&16&24\end{pmatrix}
\]
for instance \(24=9+0+4+4+4+1+1+1+0\) and \(16=9+0+0+4+4-1-1+1+0\). Hence (Proposition~\ref{r12-prop-data}(a))
\[
S=\frac18(C_9X)'(C_9X)=\begin{pmatrix}3&2&2&2\\2&3&2&2\\2&2&3&2\\2&2&2&3\end{pmatrix}=I_4+2J_4,\qquad
R=\frac13S=\begin{pmatrix}1&\frac23&\frac23&\frac23\\[2pt]\frac23&1&\frac23&\frac23\\[2pt]\frac23&\frac23&1&\frac23\\[2pt]\frac23&\frac23&\frac23&1\end{pmatrix}
\]
as \(D_s=3I_4\): the four tests are equicorrelated with \(r_{jl}=\frac23\).
\item The total score is \(Xa\) with \(a=\one_4\), so its sample variance is \(a'Sa=\one'S\one=36\), the sum of the
entries of \(S\) (Proposition~\ref{r12-prop-data}(b)); directly, the centered totals \(12,-6,-6,-6,6,0,0,0,0\) have the
squares adding up to \(288=8\cdot36\). For the difference, \(b=(1,-1,0,0)'\): \(b'Sb=3-2\cdot2+3=2\); and
\(a'Sb=(9,9,9,9)b=0\), as every column of \(S\) sums to \(9\): total and difference are uncorrelated.
\item \(S=I_4+2J_4\) has the eigenvalues \(1+4\cdot2=9\), with eigenspace \(\Span\{\one\}\), and \(1\) (three times)
(Example~\ref{r8-exm-patterned}(b)); so \(S\) is PD (Theorem~\ref{r8-thm-eigen-criterion}(a)). Let \(B=I_4+\frac12J_4\),
which has the eigenvalues \(3\) and \(1\) in the same way, so it is PD. As \(J_4^2=4J_4\)
(Proposition~\ref{r1-prop-centering}), \(B^2=I_4+J_4+\frac14J_4^2=I_4+2J_4=S\), and by the uniqueness in
Theorem~\ref{r8-thm-sqrt}(a), \(S^{1/2}=B\). Next \(\bigl(I_4+\frac12J_4\bigr)\bigl(I_4-\frac16J_4\bigr)
=I_4+\frac12J_4-\frac16J_4-\frac13J_4=I_4\), so (Theorem~\ref{r2-thm-invertible})
\[
S^{-1/2}=B^{-1}=I_4-\tfrac16J_4,\qquad S^{-1}=\bigl(S^{-1/2}\bigr)^2=I_4-\tfrac13J_4+\tfrac4{36}J_4=I_4-\tfrac29J_4
\]
(Theorem~\ref{r8-thm-sqrt}(c)).
\item The whitened data \(w_i=S^{-1/2}(x_i-\bar x)=(x_i-\bar x)-\frac16\bigl(\one'(x_i-\bar x)\bigr)\one\) are the rows of
\[
X_w=\begin{pmatrix}1&1&1&1\\1&-1&-1&-1\\-1&1&-1&-1\\-1&-1&1&-1\\1&1&1&-1\\1&-1&-1&1\\-1&1&-1&1\\-1&-1&1&1\\0&0&0&0\end{pmatrix}
\]
(for student \(1\), \((3,3,3,3)'-2\one\); for student \(2\), \((0,-2,-2,-2)'+\one\); for students \(6\), \(7\), \(8\),
\(\one'(x_i-\bar x)=0\)). They have the mean vector \(0\) and the sample covariance matrix \(I_4\)
(Proposition~\ref{r8-prop-whitening}(a)); directly, the columns of \(X_w\) add up to \(0\) and are orthogonal with squared
length \(8\), so \(\frac18X_w'X_w=I_4\).
\item By Proposition~\ref{r8-prop-whitening}(b), \(d_S(x_i,\bar x)^2=\lVert w_i\rVert^2\), which is \(4\) for students
\(1,\dots,8\) and \(0\) for student \(9\); the sum is \(32=(n-1)p\). The squared distances \(\lVert x_i-\bar x\rVert^2\) are
\(36\), \(12\), \(12\), \(12\), \(12\), \(4\), \(4\), \(4\), \(0\). Student \(1\), three points above average on every test,
follows the strong common pattern of the tests and is no more unusual than students \(6\), \(7\), \(8\), who are one
point above or below average on each test, against the pattern.
\item \(S\) has \(\lambda_1=9\) with \(q_1=\frac12\one\) and \(\lambda_2=\lambda_3=\lambda_4=1\). The ellipsoid
\(d_S(x,\bar x)\leqslant2\) has the semi-axis \(2\cdot3=6\) along \(q_1\) and \(2\cdot1=2\) in every direction orthogonal
to \(q_1\) (Proposition~\ref{r8-prop-whitening}(d)). Student \(1\) lies at the end of the long axis:
\(x_1=\bar x+3\one=\bar x+6q_1\).
\end{itemize}
\end{example}

\subsection*{The Rayleigh quotient}

\begin{lemma}[\citeat{harville}{Corollary~21.5.9}, \citeat{bapat}{\S5.1}]\label{r9-lem-eig-order}
Let \(A\in\R^{n\times n}\) be symmetric. Its eigenvalues, listed with multiplicity, are real
(Theorem~\ref{r8-thm-sym-real}); we write them in decreasing order,
\[
\lambda_1(A)\geqslant\lambda_2(A)\geqslant\dots\geqslant\lambda_n(A)
\]
or \(\lambda_1\geqslant\dots\geqslant\lambda_n\) when \(A\) is clear.
\begin{enumerate}
\item There is an orthogonal \(Q=[q_1\ \cdots\ q_n]\in\R^{n\times n}\) with \(Q'AQ=\diag(\lambda_1,\dots,\lambda_n)\);
then \(Aq_i=\lambda_iq_i\) for every \(i\), and \(A=\lambda_1q_1q_1'+\dots+\lambda_nq_nq_n'\).
\item \(\lambda_k(-A)=-\lambda_{n+1-k}(A)\) for \(k=1,\dots,n\).
\item \(\lambda_k(A+cI_n)=\lambda_k(A)+c\) for \(k=1,\dots,n\) and every \(c\in\R\).
\end{enumerate}
\end{lemma}
\begin{proof}
(a) By Theorem~\ref{r8-thm-spectral}(a),(b), with the orthonormal eigenvectors \(q_1,\dots,q_n\) numbered so that their
eigenvalues decrease, \(Q=[q_1\ \cdots\ q_n]\) is orthogonal and \(Q'AQ\) is the diagonal matrix of these eigenvalues,
which are the eigenvalues of \(A\) listed with multiplicity (Theorem~\ref{r8-thm-spectral}(c)), that is,
\(\lambda_1,\dots,\lambda_n\). The sum is Theorem~\ref{r8-thm-spectral-decomp}(c).

(b) The eigenvalues of \(-A\) are \(-\lambda_1\leqslant\dots\leqslant-\lambda_n\) (Proposition~\ref{r7-prop-eig-poly}(c)
with \(c=-1\)); in decreasing order the \(k\)th of them is \(-\lambda_{n+1-k}\).

(c) The eigenvalues of \(A+cI_n\) are \(\lambda_1+c\geqslant\dots\geqslant\lambda_n+c\)
(Proposition~\ref{r7-prop-eig-poly}(c)), already in decreasing order; the \(k\)th of them is \(\lambda_k+c\).
\end{proof}

For symmetric \(A\in\R^{n\times n}\) and \(x\ne0\), the number \(x'Ax/x'x\) is the \textbf{Rayleigh quotient} of \(x\)~\citeat{searle-khuri}{\S9.7}, \citeat{abadir-magnus}{Exercise~7.53}.
It does not change when \(x\) is replaced by \(cx\) (\(c\ne0\)); in particular it equals the quadratic form \(u'Au\)
(Definition~\ref{r8-def-quadratic-form}) at the unit vector \(u=x/\lVert x\rVert\). So the largest and the smallest
Rayleigh quotient are also the largest and the smallest value of \(u'Au\) over the unit vectors \(u\).

\begin{theorem}[\citeat{bapat}{Result~5.1; Ch.~5, Exercise~16}, \citeat{abadir-magnus}{Exercise~12.39(a)}, \citeat{tu-asm}{Theorem~5.20}, \citeat{searle-khuri}{Lemma~9.1}]\label{r9-thm-rayleigh}
Let \(A\in\R^{n\times n}\) be symmetric, with \(\lambda_1\geqslant\dots\geqslant\lambda_n\) and
\(Q=[q_1\ \cdots\ q_n]\) as in Lemma~\ref{r9-lem-eig-order}.
\begin{enumerate}
\item For every \(x\ne0\) in \(\R^n\),
\[
\lambda_n\leqslant\frac{x'Ax}{x'x}\leqslant\lambda_1
\]
and the bounds are attained at \(x=q_n\) and \(x=q_1\). Hence \(\max_{x\ne0}x'Ax/x'x=\lambda_1\) and
\(\min_{x\ne0}x'Ax/x'x=\lambda_n\).
\item For \(x\ne0\): \(x'Ax=\lambda_1x'x\) if and only if \(Ax=\lambda_1x\), and \(x'Ax=\lambda_nx'x\) if and only if
\(Ax=\lambda_nx\). So the maximum (minimum) is attained exactly at the eigenvectors belonging to \(\lambda_1\)
(\(\lambda_n\)).
\item \(\lambda_1I_n-A\) and \(A-\lambda_nI_n\) are PSD.
\end{enumerate}
\end{theorem}
\begin{proof}
Let \(x\in\R^n\) and \(y=Q'x\). Then \(x=Qy\) and \(x'x=y'y\) (Proposition~\ref{r6-prop-orthonormal}), so \(x\ne0\)
if and only if \(y\ne0\), and
\[
x'Ax=y'Q'AQy=y'\diag(\lambda_1,\dots,\lambda_n)y=\lambda_1y_1^2+\dots+\lambda_ny_n^2
\]
\needspace{3\baselineskip}
(a) By Theorem~\ref{r8-thm-eigen-criterion}(d), \(\lambda_nx'x\leqslant x'Ax\leqslant\lambda_1x'x\), with equality on the
right for \(x=q_1\) and on the left for \(x=q_n\); divide by \(x'x>0\).

\needspace{3\baselineskip}
(b)
\begin{itemize}
\item \(\lambda_1x'x-x'Ax=\sum_i(\lambda_1-\lambda_i)y_i^2\) is a sum of terms \(\geqslant0\). It is \(0\) if and only
if \(y_i=0\) for every \(i\) with \(\lambda_i<\lambda_1\), that is, \(x=Qy=\sum_{\lambda_i=\lambda_1}y_iq_i\); and then
\(Ax=\sum_{\lambda_i=\lambda_1}y_i\lambda_1q_i=\lambda_1x\).
\item Conversely, \(Ax=\lambda_1x\) gives \(x'Ax=\lambda_1x'x\).
\item For \(\lambda_n\), apply this to \(-A\), whose largest eigenvalue is \(-\lambda_n\)
(Lemma~\ref{r9-lem-eig-order}(b)).
\end{itemize}
(c) Both matrices are symmetric, and by (a), \(x'(\lambda_1I-A)x=\lambda_1x'x-x'Ax\geqslant0\) and
\(x'(A-\lambda_nI)x=x'Ax-\lambda_nx'x\geqslant0\) for every \(x\) (Definition~\ref{r8-def-definite}).
\end{proof}

\begin{lemma}[cf.\ \citeat{harville-solutions}{Ch.~21, Exercise~17}, \citeat{abadir-magnus}{Exercise~12.46}]\label{r9-lem-compression}
Let \(A\in\R^{n\times n}\) be symmetric, let \(W\ne\{0\}\) be a subspace of \(\R^n\) of dimension \(d\), and let the
columns of \(G\in\R^{n\times d}\) form an orthonormal basis of \(W\) (Theorems~\ref{thm:r4-dimension}
and~\ref{r6-thm-gram-schmidt}). Then the nonzero vectors of \(W\) are the vectors \(x=Gz\) with \(z\in\R^d\),
\(z\ne0\), and for these
\[
\frac{x'Ax}{x'x}=\frac{z'(G'AG)z}{z'z}
\]
Hence the maximum and the minimum of \(x'Ax/x'x\) over the nonzero \(x\in W\) exist: they are \(\lambda_1(G'AG)\) and
\(\lambda_d(G'AG)\), attained at \(x=Gz\) for eigenvectors \(z\) of the symmetric \(d\times d\) matrix \(G'AG\) belonging
to these eigenvalues. Both lie between \(\lambda_n(A)\) and \(\lambda_1(A)\).
\end{lemma}
\begin{proof}
\begin{itemize}
\item The columns of \(G\) span \(W\), so \(x\in W\) if and only if \(x=Gz\) for some \(z\) (Theorem~\ref{r1-thm-views}(b));
they are independent, so \(Gz=0\) only for \(z=0\).
\item \(G'G=I_d\) (Definition~\ref{r6-def-orthonormal}), so \(x'x=z'G'Gz=z'z\); also \(x'Ax=z'(G'AG)z\), and \(G'AG\) is
symmetric.
\item Apply Theorem~\ref{r9-thm-rayleigh}(a),(b) to \(G'AG\), and Theorem~\ref{r9-thm-rayleigh}(a) to \(A\) for the
last sentence.\qedhere
\end{itemize}
\end{proof}

As before, these extrema are also the largest and the smallest value of \(u'Au\) over the unit vectors \(u\in W\).

\begin{theorem}[\citeat{bapat}{Result~5.7(ii)}, \citeat{abadir-magnus}{Exercise~12.42}, \citeat{harville-solutions}{Ch.~21, Exercise~17}]\label{r9-thm-constrained}
Let \(A\in\R^{n\times n}\) be symmetric, with \(Q=[q_1\ \cdots\ q_n]\) as in Lemma~\ref{r9-lem-eig-order}, and let
\(1\leqslant k\leqslant n\). Then
\[
\lambda_k=\max\Bigl\{\frac{x'Ax}{x'x}:x\ne0,\ x\perp q_1,\dots,q_{k-1}\Bigr\}
=\min\Bigl\{\frac{x'Ax}{x'x}:x\ne0,\ x\perp q_{k+1},\dots,q_n\Bigr\}
\]
and both are attained at \(x=q_k\). (For \(k=1\) the first set has no condition, for \(k=n\) the second.)
\end{theorem}
\begin{proof}
\begin{itemize}
\item Let \(0\leqslant j\leqslant n\), and split \(Q=[R\ T]\) with \(R=[q_1\ \cdots\ q_j]\). Then \(R'T=0\), as
\(q_i'q_l=0\) for \(i\ne l\), and \(I_n=QQ'=RR'+TT'\) (Proposition~\ref{r6-prop-orthonormal}(b) and
Theorem~\ref{r5b-thm-block}). The entries of \(R'x\) are \(q_1'x,\dots,q_j'x\). If \(R'x=0\), then \(x=TT'x\in\mathcal C(T)\);
and \(R'(Tz)=0\) for every \(z\). So \(x\perp q_1,\dots,q_j\) if and only if \(x\in\mathcal C(T)\); in the same way,
\(x\perp q_{j+1},\dots,q_n\) if and only if \(x\in\mathcal C(R)\).
\item First form (\(j=k-1\)): the \(x\perp q_1,\dots,q_{k-1}\) form \(\mathcal C(T)\), where the columns of
\(T=[q_k\ \cdots\ q_n]\) are an orthonormal basis of it. \(T'AT=\diag(\lambda_k,\dots,\lambda_n)\), the lower right block of
\(Q'AQ\), has the largest eigenvalue \(\lambda_k\) with eigenvector \(e_1\). By Lemma~\ref{r9-lem-compression}, the
maximum is \(\lambda_k\), attained at \(Te_1=q_k\).
\item Second form (\(j=k\)): the \(x\perp q_{k+1},\dots,q_n\) form \(\mathcal C(R)\), \(R=[q_1\ \cdots\ q_k]\), and
\(R'AR=\diag(\lambda_1,\dots,\lambda_k)\) has the smallest eigenvalue \(\lambda_k\) with eigenvector \(e_k\). By
Lemma~\ref{r9-lem-compression}, the minimum is \(\lambda_k\), attained at \(Re_k=q_k\).\qedhere
\end{itemize}
\end{proof}

\begin{definition}[\citeat{searle-khuri}{Definition~9.1}, \citeat{bapat}{Ch.~5, Exercise~12}]\label{r9-def-frobenius}
The \textbf{Frobenius norm} of \(A=(a_{ij})\in\R^{m\times n}\) is
\[
\lVert A\rVert_F=\bigl(\tr(A'A)\bigr)^{1/2}=\Bigl(\sum_{i=1}^m\sum_{j=1}^na_{ij}^2\Bigr)^{1/2}
\]
(Theorem~\ref{r1-thm-trace}(d)). If \(v\in\R^{mn}\) lists the entries of \(A\) column after column, then
\(\lVert A\rVert_F=\lVert v\rVert\) (Definition~\ref{r6-def-inner}); and
\(\lVert A\rVert_F^2=\lVert a_1\rVert^2+\dots+\lVert a_n\rVert^2\) for the columns \(a_j\) of \(A\). (Some books write
\(\lVert A\rVert_2\) for \(\lVert A\rVert_F\).)
\end{definition}

\begin{theorem}[Ky Fan; \citeat{bapat}{Result~5.12}, \citeat{abadir-magnus}{Exercise~12.49(a)}, \citeat{harville}{Theorem~21.12.5}, \citeat{tu-asm}{Theorem~5.22}]\label{r12-thm-kyfan}
Let \(A\in\R^{n\times n}\) be symmetric, with eigenvalues \(\lambda_1\geqslant\lambda_2\geqslant\dots\geqslant\lambda_n\) (listed
with multiplicity) and orthonormal eigenvectors \(q_1,\dots,q_n\), \(Aq_i=\lambda_iq_i\) (Lemma~\ref{r9-lem-eig-order}(a)).
Let \(1\leqslant k\leqslant n\), and let \(G\in\R^{n\times k}\) have orthonormal columns, \(G'G=I_k\).
\begin{enumerate}
\item \(\tr(G'AG)\leqslant\lambda_1+\dots+\lambda_k\), with equality for \(G=[q_1\ \cdots\ q_k]\).
\item \(\tr(G'AG)\geqslant\lambda_{n-k+1}+\dots+\lambda_n\), with equality for \(G=[q_{n-k+1}\ \cdots\ q_n]\).
\item If \(k<n\) and \(\lambda_k>\lambda_{k+1}\), then equality holds in (a) only if
\(\mathcal C(G)=\Span\{q_1,\dots,q_k\}\).
\end{enumerate}
So the largest (smallest) value of \(\tr(G'AG)\) over all \(G\in\R^{n\times k}\) with \(G'G=I_k\) is the sum of the
\(k\) largest (smallest) eigenvalues of \(A\).
\end{theorem}
\begin{proof}
(a)
\begin{itemize}
\item \(Q=[q_1\ \cdots\ q_n]\) is orthogonal and \(Q'AQ=\Lambda=\diag(\lambda_1,\dots,\lambda_n)\)
(Theorem~\ref{r8-thm-spectral}(b)). Put \(R=Q'G=(r_{ij})\in\R^{n\times k}\). As \(QQ'=I_n\)
(Proposition~\ref{r6-prop-orthonormal}(b)), \(G=QR\), and \(R'R=G'QQ'G=G'G=I_k\).
\item So \(\tr(G'AG)=\tr(R'Q'AQR)=\tr(R'\Lambda R)\). The \(j\)th diagonal entry of \(R'\Lambda R\) is
\(\col_j(R)'\Lambda\col_j(R)=\sum_i\lambda_ir_{ij}^2\) (Theorem~\ref{r1-thm-views}(d)); hence
\[
\tr(G'AG)=\sum_{j=1}^k\sum_{i=1}^n\lambda_ir_{ij}^2=\sum_{i=1}^nw_i\lambda_i,\qquad w_i=\sum_{j=1}^kr_{ij}^2
\]
\item \(w_i\geqslant0\), and \(\sum_iw_i=\sum_j\sum_ir_{ij}^2=\tr(R'R)=\tr I_k=k\) (Theorem~\ref{r1-thm-trace}(d)).
\item \(w_i\leqslant1\): \(w_i\) is the \(i\)th diagonal entry of \(RR'\) (Theorem~\ref{r1-thm-views}(d)). If \(k=n\), then
\(RR'=I_n\) (Proposition~\ref{r6-prop-orthonormal}(b)) and \(w_i=1\). If \(k<n\), Corollary~\ref{r6-cor-onb}(c) gives
\(T\in\R^{n\times(n-k)}\) with \([R\ T]\) orthogonal and \(RR'+TT'=I_n\). The \(i\)th diagonal entries give
\(w_i+\sum_lt_{il}^2=1\), so \(w_i\leqslant1\).
\item The right side of
\[
\sum_{i\leqslant k}\lambda_i-\sum_{i=1}^nw_i\lambda_i=\sum_{i\leqslant k}(\lambda_i-\lambda_k)(1-w_i)+\sum_{i>k}(\lambda_k-\lambda_i)w_i
\]
expands to \(\sum_{i\leqslant k}\lambda_i-\sum_iw_i\lambda_i-\lambda_k\bigl(k-\sum_iw_i\bigr)\), which is the left side as
\(\sum_iw_i=k\). Every term on the right is \(\geqslant0\): \(\lambda_i\geqslant\lambda_k\) and \(w_i\leqslant1\) for
\(i\leqslant k\); \(\lambda_i\leqslant\lambda_k\) and \(w_i\geqslant0\) for \(i>k\). This proves the inequality.
\item For \(G=[q_1\ \cdots\ q_k]\), \(AG=[\lambda_1q_1\ \cdots\ \lambda_kq_k]=G\diag(\lambda_1,\dots,\lambda_k)\), so
\(G'AG=\diag(\lambda_1,\dots,\lambda_k)\), with trace \(\lambda_1+\dots+\lambda_k\).
\end{itemize}
(b)
\begin{itemize}
\item \(-A\) is symmetric, with the orthonormal eigenvectors \(q_n,\dots,q_1\) for the eigenvalues
\(-\lambda_n\geqslant\dots\geqslant-\lambda_1\). By (a) for \(-A\),
\(-\tr(G'AG)=\tr(G'(-A)G)\leqslant-\lambda_n-\dots-\lambda_{n-k+1}\), with equality for \(G=[q_n\ \cdots\ q_{n-k+1}]\).
\item \(\tr(G'AG)=\sum_jg_j'Ag_j\) over the columns \(g_j\) of \(G\) (Theorem~\ref{r1-thm-views}(d)), so the order of the
columns does not matter.
\end{itemize}
\needspace{4\baselineskip}
(c)
\begin{itemize}
\item If \(\tr(G'AG)=\lambda_1+\dots+\lambda_k\), every term on the right of the identity in (a) is \(0\). For \(i>k\),
\(\lambda_k-\lambda_i\geqslant\lambda_k-\lambda_{k+1}>0\), so \(w_i=0\). Then \(\sum_{i\leqslant k}w_i=k\) with every
\(w_i\leqslant1\) gives \(w_i=1\) for \(i\leqslant k\).
\item Row \(i\) of \(R=Q'G\) is \(q_i'G\), so \(w_i=\lVert G'q_i\rVert^2\). As \(G'G=I_k\),
\(\lVert GG'q_i\rVert^2=q_i'GG'GG'q_i=w_i\), and \(GG'q_i\perp q_i-GG'q_i\), since
\((GG'q_i)'(q_i-GG'q_i)=w_i-w_i=0\). By Pythagoras (Proposition~\ref{r6-prop-pythagoras}),
\(\lVert q_i-GG'q_i\rVert^2=\lVert q_i\rVert^2-w_i=0\) for \(i\leqslant k\).
\item Hence \(q_i=G(G'q_i)\in\mathcal C(G)\) for \(i\leqslant k\), and \(\Span\{q_1,\dots,q_k\}\subseteq\mathcal C(G)\). Both
spaces have dimension \(k\) (Proposition~\ref{r6-prop-orthogonal-list}(b)), so they are equal
(Theorem~\ref{thm:r4-counting}(a)).\qedhere
\end{itemize}
\end{proof}

For \(k=1\) and \(G=x/\lVert x\rVert\) (\(x\ne0\)), (a) and (b) are the bounds
\(\lambda_n\leqslant x'Ax/x'x\leqslant\lambda_1\) of Theorem~\ref{r9-thm-rayleigh}(a).

\begin{example}\label{r9-exm-ky-fan}
Let
\[
A=\begin{pmatrix}15&5&10&2\\5&15&2&10\\10&2&15&5\\2&10&5&15\end{pmatrix}
\]
Find the maximum and the minimum of \(\tr(G'AG)\) over the matrices \(G\in\R^{4\times2}\) with \(G'G=I_2\).

\textit{Solution.}
\begin{itemize}
\item Let \(h_1=(1,1,1,1)'\), \(h_2=(1,1,-1,-1)'\), \(h_3=(1,-1,1,-1)'\) and \(h_4=(1,-1,-1,1)'\). Row by row,
\(Ah_1=32h_1\), \(Ah_2=8h_2\), \(Ah_3=18h_3\) and \(Ah_4=2h_4\); for instance \(15+5+10+2=32\), \(15+5-10-2=8\),
\(15-5+10-2=18\) and \(15-5-10+2=2\) in row \(1\). The \(h_m\) are pairwise orthogonal with \(\lVert h_m\rVert=2\), so the
eigenvalues of \(A\) are \(32,18,8,2\), with the orthonormal eigenvectors \(\frac12h_1,\frac12h_3,\frac12h_2,\frac12h_4\)
(Theorem~\ref{r8-thm-spectral}(c)).
\item By Theorem~\ref{r12-thm-kyfan}, the maximum is \(32+18=50\), attained at \(G=\frac12[h_1\ h_3]\), and the minimum
is \(8+2=10\), attained at \(G=\frac12[h_2\ h_4]\).
\item Other \(G\) give values in between: \(G=[e_1\ e_2]\) gives \(a_{11}+a_{22}=30\), and
\(G=\frac1{\sqrt2}[e_1+e_2\ \ e_3+e_4]\) gives \(\frac12(15+5+5+15)+\frac12(15+5+5+15)=40\). The matrix
\(G=\frac1{\sqrt2}[e_1+e_3\ \ e_2+e_4]\) gives \(\frac12(15+10+10+15)\cdot2=50\): its columns span the same plane \(W\) as
\(h_1=(e_1+e_3)+(e_2+e_4)\) and \(h_3=(e_1+e_3)-(e_2+e_4)\), and \(\tr(G'AG)=\tr(AGG')\) (Theorem~\ref{r1-thm-trace}(b))
depends on \(G\) only through \(GG'=P_W\) (Theorem~\ref{r6-thm-projection-matrix}(b)).
\end{itemize}
\end{example}

\begin{corollary}[\citeat{abadir-magnus}{Exercise~12.49(b)}, \citeat{searle-khuri}{Exercise~6.22}]\label{r12-cor-diagonal}
Let \(A=(a_{ij})\in\R^{n\times n}\) be symmetric, with eigenvalues \(\lambda_1\geqslant\dots\geqslant\lambda_n\)
(Lemma~\ref{r9-lem-eig-order}), and let \(1\leqslant i_1<i_2<\dots<i_k\leqslant n\). Then
\[
\lambda_{n-k+1}+\dots+\lambda_n\leqslant a_{i_1i_1}+a_{i_2i_2}+\dots+a_{i_ki_k}\leqslant\lambda_1+\dots+\lambda_k
\]
In particular \(\lambda_n\leqslant a_{ii}\leqslant\lambda_1\) for every \(i\).
\end{corollary}
\begin{proof}
\begin{itemize}
\item Let \(G=[e_{i_1}\ \cdots\ e_{i_k}]\in\R^{n\times k}\). The \((s,t)\) entry of \(G'G\) is \(e_{i_s}'e_{i_t}\)
(Theorem~\ref{r1-thm-views}(d)), which is \(1\) for \(s=t\) and \(0\) for \(s\ne t\), as \(i_s\ne i_t\); so \(G'G=I_k\).
\item The \((s,t)\) entry of \(G'AG\) is \(e_{i_s}'Ae_{i_t}=a_{i_si_t}\) (Theorem~\ref{r1-thm-views}(b),(d)): \(G'AG\) is the
submatrix of \(A\) in rows and columns \(i_1,\dots,i_k\), and \(\tr(G'AG)=a_{i_1i_1}+\dots+a_{i_ki_k}\).
\item Theorem~\ref{r12-thm-kyfan}(a),(b) bound \(\tr(G'AG)\) by \(\lambda_1+\dots+\lambda_k\) from above and by
\(\lambda_{n-k+1}+\dots+\lambda_n\) from below. For \(k=1\) and \(i_1=i\) this is the last sentence.\qedhere
\end{itemize}
\end{proof}

\subsection*{Principal components}

\begin{definition}[\citeat{tu-asm}{Definition~5.19}, \citeat{bapat}{\S10.1}]\label{r12-def-pc}
Let \(X\in\R^{n\times p}\) be a data matrix with \(S\ne0\). A \textbf{first loading vector} is a unit vector
\(v_1\in\R^p\) that maximizes \(a'Sa\) over all unit vectors \(a\in\R^p\). For \(m\geqslant2\), once \(v_1,\dots,v_{m-1}\)
are chosen, an \textbf{\(m\)th loading vector} is a unit vector \(v_m\) that maximizes \(a'Sa\) over the unit vectors
\(a\) with \(a'Sv_l=0\) for \(l=1,\dots,m-1\) (the \textbf{\(m\)th problem}). The linear combination of the variables with
coefficient vector \(v_m\) is the \textbf{\(m\)th principal component}, and the vector of its centered values,
\[
z_m=C_nXv_m\in\R^n,\qquad z_{im}=v_m'(x_i-\bar x)\quad(i=1,\dots,n)
\]
is the \textbf{\(m\)th score vector}.
\end{definition}

By Proposition~\ref{r12-prop-data}(b), \(a'Sa\) is the sample variance of \(Xa\), and \(a'Sv_l=0\) says that \(Xa\) and
\(Xv_l\) are uncorrelated. So the \(m\)th principal component has the largest sample variance among the combinations
with unit coefficient vector that are uncorrelated with the first \(m-1\) principal components.

\begin{theorem}[\citeat{tu-asm}{Theorem~5.20}, \citeat{bapat}{\S10.1}]\label{r12-thm-pca}
Let \(X\in\R^{n\times p}\) be a data matrix with \(S\ne0\), and let \(S=V\Lambda V'\), where \(V=[v_1\ \cdots\ v_p]\) is
orthogonal and \(\Lambda=\diag(\lambda_1,\dots,\lambda_p)\) with \(\lambda_1\geqslant\dots\geqslant\lambda_p\)
(Lemma~\ref{r9-lem-eig-order}(a)). Let \(r=\rank(C_nX)\) and \(z_m=C_nXv_m\).
\begin{enumerate}
\item \(\lambda_p\geqslant0\), \(r\leqslant\min\{n-1,p\}\), and \(\lambda_m>0\) if and only if \(m\leqslant r\).
\item If \(\lambda_l>0\), then for every \(a\in\R^p\), \(Xa\) and \(Xv_l\) are uncorrelated if and only if \(a\perp v_l\).
\item Let \(m\leqslant r\). If \(v_1,\dots,v_{m-1}\) are taken as the first \(m-1\) loading vectors, then \(v_m\) is an
\(m\)th loading vector, and the maximum in the \(m\)th problem is \(\lambda_m\).
\item \(z_m'z_l=(n-1)\lambda_m\) if \(l=m\) and \(0\) if \(l\ne m\): the score vectors are pairwise orthogonal and
uncorrelated, and \(z_m\) has sample variance \(\lambda_m\).
\item Let \(M\leqslant r\) and \(\lambda_1>\lambda_2>\dots>\lambda_M>\lambda_{M+1}\) (the last inequality is omitted if
\(M=p\)). Then the first \(M\) loading vectors are \(\pm v_1,\dots,\pm v_M\), and the first \(M\) score vectors are
\(\pm z_1,\dots,\pm z_M\).
\end{enumerate}
\end{theorem}
\begin{proof}
(a)
\begin{itemize}
\item \(S\) is PSD (Proposition~\ref{r12-prop-data}(a)), so \(\lambda_p\geqslant0\)
(Theorem~\ref{r8-thm-eigen-criterion}(a)).
\item \(\one'C_n=(C_n\one)'=0'\) (Proposition~\ref{r1-prop-centering}), so every column of \(C_nX\) lies in
\(\Span\{\one\}^\perp\), which has dimension \(n-1\) (Theorem~\ref{r6-thm-complement}(c)). Hence \(r\leqslant n-1\)
(Theorem~\ref{thm:r4-counting}(a)), and \(r\leqslant p\) (Theorem~\ref{thm:r4-rank-rules}(a)).
\item \(\rank S=\rank\bigl((C_nX)'(C_nX)\bigr)=r\) (Theorem~\ref{thm:r4-rank-ata}(b)), and \(\rank S=\rank\Lambda\), as \(V\) is
invertible (Theorem~\ref{thm:r4-rank-product}(d)). The nonzero columns \(\lambda_me_m\) of \(\Lambda\) are pairwise
orthogonal, hence independent (Proposition~\ref{r6-prop-orthogonal-list}(b)), and span \(\mathcal C(\Lambda)\); so
\(\rank\Lambda\) is the number of nonzero \(\lambda_m\).
\item As \(\lambda_1\geqslant\dots\geqslant\lambda_p\geqslant0\), exactly \(\lambda_1,\dots,\lambda_r\) are positive.
\end{itemize}
(b) By Proposition~\ref{r12-prop-data}(b) and \(Sv_l=\lambda_lv_l\), the sample covariance of \(Xa\) and \(Xv_l\) is
\(a'Sv_l=\lambda_la'v_l\), which is \(0\) if and only if \(a'v_l=0\), as \(\lambda_l>0\).

\needspace{4\baselineskip}
(c)
\begin{itemize}
\item Let \(a\) be a unit vector and \(c=V'a\). Then \(a=Vc\), \(c_j=v_j'a\), \(\sum_jc_j^2=\lVert c\rVert^2=\lVert a\rVert^2=1\)
(Proposition~\ref{r6-prop-orthonormal}), and \(a'Sa=c'V'SVc=c'\Lambda c=\sum_j\lambda_jc_j^2\).
\item By (b) and (a) (\(\lambda_l>0\) for \(l<m\leqslant r\)), the constraint of the \(m\)th problem says
\(c_1=\dots=c_{m-1}=0\). Then
\[
\lambda_m-a'Sa=\sum_{j\geqslant m}\lambda_mc_j^2-\sum_{j\geqslant m}\lambda_jc_j^2=\sum_{j>m}(\lambda_m-\lambda_j)c_j^2\geqslant0
\]
\item \(v_m\) satisfies the constraint, as \(v_m'Sv_l=\lambda_lv_m'v_l=0\) for \(l<m\), and \(v_m'Sv_m=\lambda_m\). So \(v_m\)
solves the \(m\)th problem, with maximum \(\lambda_m\).
\end{itemize}
(d) \(z_m'z_l=v_m'(C_nX)'(C_nX)v_l=(n-1)v_m'Sv_l=(n-1)\lambda_lv_m'v_l\) (Proposition~\ref{r12-prop-data}(a)). The score
vectors have mean \(0\)
(\(\one'C_n=0'\)), so \(\frac1{n-1}z_m'z_l\) is their sample covariance.

(e) Let \(m\leqslant M\), and suppose the first \(m-1\) loading vectors are \(\pm v_1,\dots,\pm v_{m-1}\) (nothing for
\(m=1\)).
\begin{itemize}
\item The constraint is again \(c_1=\dots=c_{m-1}=0\), and a maximizer \(a\) has \(a'Sa=\lambda_m\) by (c). By the display
in (c), \(\sum_{j>m}(\lambda_m-\lambda_j)c_j^2=0\), where \(\lambda_m-\lambda_j\geqslant\lambda_m-\lambda_{m+1}>0\) for
\(j>m\).
\item So \(c_j=0\) for \(j\ne m\), and \(a=c_mv_m\) with \(c_m^2=1\): \(a=\pm v_m\), with score vector \(\pm z_m\). By
induction on \(m\) the claim follows.\qedhere
\end{itemize}
\end{proof}

Part (c) is Theorem~\ref{r9-thm-constrained} for \(S\). The loading vectors and the maxima depend on the data only
through \(S\), and the proofs of (c) and (e) use only that \(S\) is symmetric and PSD: for a PSD matrix given without the
data, the principal components are found in the same way.

\begin{proposition}[\citeat{tu-asm}{Proposition~5.21}, \citeat{bapat}{\S10.1}]\label{r12-prop-pve}
Let \(X\), \(S\), \(\lambda_m\), \(z_m\) and \(r\) be as in Theorem~\ref{r12-thm-pca}. The \textbf{total variance}
\(s_{11}+\dots+s_{pp}\) equals \(\tr S=\lambda_1+\dots+\lambda_p\), and the \textbf{proportion of variance explained} by
the \(m\)th principal component is
\[
\mathrm{PVE}_m=\frac{\lVert z_m\rVert^2}{\lVert C_nX\rVert_F^2}=\frac{\lambda_m}{\lambda_1+\dots+\lambda_p}\qquad(m=1,\dots,p)
\]
It is positive for \(m\leqslant r\) and \(0\) for \(m>r\), and \(\mathrm{PVE}_1\geqslant\mathrm{PVE}_2\geqslant\cdots\). The
cumulative proportion \(\mathrm{PVE}_1+\dots+\mathrm{PVE}_M\) increases with \(M\leqslant r\) and equals \(1\) for
\(M\geqslant r\).
\end{proposition}
\begin{proof}
\begin{itemize}
\item \(\tr S=s_{11}+\dots+s_{pp}\), and \(\tr S=\tr(V\Lambda V')=\tr(\Lambda V'V)=\tr\Lambda\)
(Theorem~\ref{r1-thm-trace}(b)).
\item \(\lVert C_nX\rVert_F^2=\tr\bigl((C_nX)'(C_nX)\bigr)=(n-1)\tr S\) (Definition~\ref{r9-def-frobenius},
Proposition~\ref{r12-prop-data}(a)) and \(\lVert z_m\rVert^2=(n-1)\lambda_m\) (Theorem~\ref{r12-thm-pca}(d)); this gives
the second expression, and \(\tr S\geqslant\lambda_1>0\) (\(S\ne0\), so \(r\geqslant1\), and
Theorem~\ref{r12-thm-pca}(a)).
\item By Theorem~\ref{r12-thm-pca}(a), \(\lambda_m>0\) exactly for \(m\leqslant r\), and \(\lambda_m\) decreases with \(m\);
so \(\mathrm{PVE}_1+\dots+\mathrm{PVE}_r=(\lambda_1+\dots+\lambda_r)/\tr S=1\).\qedhere
\end{itemize}
\end{proof}

\begin{example}[\citeat{bapat}{Ch.~10, Exercise~1}]\label{r12-exm-bapat}
The sample covariance matrix of two variables is
\[
S=\begin{pmatrix}60&20\\20&60\end{pmatrix}
\]
Find the principal components and their sample variances.

\textit{Solution.} \(S=40I_2+20J_2\), and \(J_2\one_2=\one_2(\one_2'\one_2)=2\one_2\), \(J_2(1,-1)'=\one_2(1-1)=0\)
(Proposition~\ref{r1-prop-centering}). So
\[
S\begin{pmatrix}1\\1\end{pmatrix}=\begin{pmatrix}80\\80\end{pmatrix}=80\begin{pmatrix}1\\1\end{pmatrix},\qquad
S\begin{pmatrix}1\\-1\end{pmatrix}=\begin{pmatrix}40\\-40\end{pmatrix}=40\begin{pmatrix}1\\-1\end{pmatrix}
\]
With the orthonormal eigenvectors \(v_1=\frac1{\sqrt2}(1,1)'\), \(v_2=\frac1{\sqrt2}(1,-1)'\), the eigenvalues are
\(\lambda_1=80\), \(\lambda_2=40\) (Theorem~\ref{r8-thm-spectral}(c)). By Theorem~\ref{r12-thm-pca}, the principal components
are the sum of the two variables divided by \(\sqrt2\), with sample variance \(80\), and their difference divided by
\(\sqrt2\), with sample variance \(40\); they explain \(\frac{80}{120}=\frac23\) and \(\frac{40}{120}=\frac13\) of the total
variance (Proposition~\ref{r12-prop-pve}).
\end{example}

\begin{example}\label{r12-exm-pca}
Find the principal components, the score vectors and the proportions of variance explained for the data of
Example~\ref{r12-exm-data}.

\textit{Solution.}
\begin{itemize}
\item Let \(h_1=(1,1,1,1)'\), \(h_2=(1,1,-1,-1)'\), \(h_3=(1,-1,1,-1)'\) and \(h_4=(1,-1,-1,1)'\). Row by row,
\[
Sh_1=\begin{pmatrix}25\\25\\25\\25\end{pmatrix}=25h_1,\qquad
Sh_2=\begin{pmatrix}9\\9\\-9\\-9\end{pmatrix}=9h_2,\qquad
Sh_3=\begin{pmatrix}5\\-5\\5\\-5\end{pmatrix}=5h_3,\qquad
Sh_4=\begin{pmatrix}1\\-1\\-1\\1\end{pmatrix}=h_4
\]
for instance \(10+7+5+3=25\), \(10+7-5-3=9\), \(10-7+5-3=5\) and \(10-7-5+3=1\) in row \(1\). The \(h_m\) are pairwise
orthogonal with \(\lVert h_m\rVert=2\), so \(v_m=\frac12h_m\) (\(m=1,\dots,4\)) are orthonormal eigenvectors, and the
eigenvalues of \(S\) are \(\lambda_1=25\), \(\lambda_2=9\), \(\lambda_3=5\), \(\lambda_4=1\) (Theorem~\ref{r8-thm-spectral}(c));
with \(V=[v_1\ \cdots\ v_4]\), \(S=V\diag(25,9,5,1)V'\).
\item The eigenvalues are positive and distinct, so \(r=4\) and the loading vectors \(v_1,\dots,v_4\) are unique up to sign
(Theorem~\ref{r12-thm-pca}(a),(e)). The first principal component, with the coefficient vector \(\frac12(1,1,1,1)'\) and
sample variance \(25\), is half the sum of the four variables: an overall level. The second, \(\frac12(1,1,-1,-1)'\), with
sample variance \(9\), contrasts variables \(1\) and \(2\) with variables \(3\) and \(4\); the third,
\(\frac12(1,-1,1,-1)'\), with sample variance \(5\), contrasts variables \(1\) and \(3\) with variables \(2\) and \(4\); the
fourth, \(\frac12(1,-1,-1,1)'\), has sample variance \(1\).
\item Scores: \(C_5Xh_1=(0,10,10,-10,-10)'\), \(C_5Xh_2=(0,6,-6,6,-6)'\), \(C_5Xh_3=(8,-2,-2,-2,-2)'\) and
\(C_5Xh_4=(0,2,-2,-2,2)'\), so \(z_m=C_5Xv_m=\frac12C_5Xh_m\) gives
\[
z_1=\begin{pmatrix}0\\5\\5\\-5\\-5\end{pmatrix},\qquad
z_2=\begin{pmatrix}0\\3\\-3\\3\\-3\end{pmatrix},\qquad
z_3=\begin{pmatrix}4\\-1\\-1\\-1\\-1\end{pmatrix},\qquad
z_4=\begin{pmatrix}0\\1\\-1\\-1\\1\end{pmatrix}
\]
Check (Theorem~\ref{r12-thm-pca}(d)): \(\lVert z_1\rVert^2=100=4\cdot25\), \(\lVert z_2\rVert^2=36=4\cdot9\),
\(\lVert z_3\rVert^2=20=4\cdot5\), \(\lVert z_4\rVert^2=4=4\cdot1\), and the six products \(z_m'z_l\) (\(m<l\)) are \(0\); for
instance \(z_1'z_2=0+15-15-15+15=0\) and \(z_1'z_3=0-5-5+5+5=0\).
\item \(\tr S=10+10+10+10=40=25+9+5+1\), so (Proposition~\ref{r12-prop-pve}) \(\mathrm{PVE}_1=\frac{25}{40}=\frac58\),
\(\mathrm{PVE}_2=\frac9{40}\), \(\mathrm{PVE}_3=\frac5{40}=\frac18\), \(\mathrm{PVE}_4=\frac1{40}\): the first component explains
\(\frac58\) of the total variance, the first two together \(\frac{34}{40}=\frac{17}{20}\).
\end{itemize}
\end{example}

\begin{remark}[\citeat{islr}{\S12.2.4}, \citeat{tu-asm}{Remark after Proposition~5.24}]\label{r12-rem-number}
Usually only the first \(M\) principal components are kept, with \(M\) much smaller than \(p\); their score vectors
\(z_1,\dots,z_M\) summarize the data, for instance in a plot of the units. By Proposition~\ref{r12-prop-pve}, they explain
together the proportion \(\mathrm{PVE}_1+\dots+\mathrm{PVE}_M\) of the total variance, and
\(\mathrm{PVE}_1\geqslant\mathrm{PVE}_2\geqslant\cdots\). A common aid is the \textbf{scree plot}, the points
\((m,\mathrm{PVE}_m)\), \(m=1,\dots,p\), joined by segments: one takes the smallest \(M\) that explains a sizable part of
the total variance, typically the \(M\) after which the \(\mathrm{PVE}_m\) drop off and stay small (an \textbf{elbow} of the
plot after the \(M\)th component). There
is no objective rule; the choice depends on the data and on the purpose. Corollary~\ref{r12-cor-diagonal}, applied to
\(S\), shows what is gained: the sample variances of any \(M\) of the variables add up to at most
\(\lambda_1+\dots+\lambda_M\). So no \(M\) of the original variables explain a larger part of the total variance than the
first \(M\) principal components. More generally, by Theorem~\ref{r12-thm-kyfan}(a) for \(S\) and
Proposition~\ref{r12-prop-data}(b), the sample variances of \(M\) combinations \(Xg_1,\dots,Xg_M\) with orthonormal
coefficient vectors \(g_1,\dots,g_M\) add up to \(\tr(G'SG)\leqslant\lambda_1+\dots+\lambda_M\), where
\(G=[g_1\ \cdots\ g_M]\) (\(\tr(G'SG)=\sum_jg_j'Sg_j\), Theorem~\ref{r1-thm-views}(d)): no \(M\) orthonormal coefficient
vectors retain a larger total variance than the first \(M\) loading vectors.
\end{remark}

\begin{example}[cf.\ \citeat{islr}{\S12.2.4, Figure~12.3}]\label{r12-exm-scree}
Four tests are given to the same students: two verbal tests (variables \(1\) and \(2\)) and two quantitative tests
(variables \(3\) and \(4\)). The sample covariance matrix of the four scores is
\[
S=\begin{pmatrix}5&4&2&2\\4&5&2&2\\2&2&5&4\\2&2&4&5\end{pmatrix}
\]
Find the principal components and the proportions of variance explained, and decide how many components to keep.

\textit{Solution.}
\begin{itemize}
\item Let \(h_1=(1,1,1,1)'\), \(h_2=(1,1,-1,-1)'\), \(h_3=(1,-1,1,-1)'\) and \(h_4=(1,-1,-1,1)'\). Row by row,
\[
Sh_1=\begin{pmatrix}13\\13\\13\\13\end{pmatrix}=13h_1,\qquad
Sh_2=\begin{pmatrix}5\\5\\-5\\-5\end{pmatrix}=5h_2,\qquad
Sh_3=\begin{pmatrix}1\\-1\\1\\-1\end{pmatrix}=h_3,\qquad
Sh_4=\begin{pmatrix}1\\-1\\-1\\1\end{pmatrix}=h_4
\]
for instance \(5+4+2+2=13\) and \(5+4-2-2=5\) in row \(1\). The \(h_i\) are pairwise orthogonal with
\(\lVert h_i\rVert=2\), so \(v_i=\frac12h_i\) (\(i=1,\dots,4\)) are orthonormal eigenvectors, and the eigenvalues of \(S\)
are \(13,5,1,1\) (Theorem~\ref{r8-thm-spectral}(c)).
\item By Theorem~\ref{r12-thm-pca}(c),(e) (\(13>5>1\)), the first two loading vectors are \(\pm v_1\) and \(\pm v_2\). The
first principal component, with coefficient vector \(\frac12(1,1,1,1)'\) and sample variance \(13\), is half the total
score: an overall level. The second, with coefficient vector \(\frac12(1,1,-1,-1)'\) and sample variance \(5\), contrasts
the verbal with the quantitative scores. As \(\lambda_3=\lambda_4=1\), the third loading vector is not unique: besides
\(v_3\), also \(a=\frac1{\sqrt2}(v_3+v_4)=\frac1{\sqrt2}(1,-1,0,0)'\) is a unit vector orthogonal to \(v_1\), \(v_2\), and
\(S(1,-1,0,0)'=(1,-1,0,0)'\) gives \(a'Sa=1=\lambda_3\), the maximum of the third problem
(Theorem~\ref{r12-thm-pca}(b),(c)).
\item \(\tr S=20=13+5+1+1\), so (Proposition~\ref{r12-prop-pve})
\[
\begin{array}{lcccc}
m&1&2&3&4\\
\hline
\mathrm{PVE}_m&\frac{13}{20}&\frac5{20}=\frac14&\frac1{20}&\frac1{20}\\[2pt]
\mathrm{PVE}_1+\dots+\mathrm{PVE}_m&\frac{13}{20}&\frac{18}{20}=\frac9{10}&\frac{19}{20}&1
\end{array}
\]
The scree plot (Remark~\ref{r12-rem-number})
\begin{center}
\begin{tikzpicture}[x=13mm,y=36mm]
\draw[->] (3/5,0)--(9/2,0) node[right]{\(m\)};
\draw[->] (3/5,0)--(3/5,39/50) node[above]{\(\mathrm{PVE}_m\)};
\draw (14/25,1/4)--(16/25,1/4) node[pos=0,left]{\small\(\frac14\)};
\draw (14/25,1/2)--(16/25,1/2) node[pos=0,left]{\small\(\frac12\)};
\foreach \m in {1,2,3,4} \draw (\m,1/50)--(\m,-1/50) node[below]{\small\(\m\)};
\draw (1,13/20)--(2,1/4)--(3,1/20)--(4,1/20);
\fill (1,13/20) circle (1.6pt);
\fill (2,1/4) circle (1.6pt);
\fill (3,1/20) circle (1.6pt);
\fill (4,1/20) circle (1.6pt);
\end{tikzpicture}
\end{center}
falls steeply up to \(m=3\) and is flat afterwards: there is an elbow after the second component. Two components,
which explain \(90\%\) of the total variance, are kept.
\item Any two of the four scores have sample variances adding up to \(10\), at most \(13+5=18\), as
Corollary~\ref{r12-cor-diagonal} says: the two verbal scores carry \(\frac{10}{20}=50\%\) of the total variance, the first
two principal components \(90\%\).
\end{itemize}
\end{example}

\subsection*{Standardized variables}

The principal components of the standardized data matrix \(X_s\) (Definition~\ref{r12-def-pc}) are obtained from \(R\): by
Proposition~\ref{r12-prop-correlation}(a) and Theorem~\ref{r12-thm-pca}, the loading vectors \(v_m\) are orthonormal
eigenvectors of \(R\), the score vectors are \(C_nX_sv_m=X_sv_m\), and, as \(\tr R=p\), \(\mathrm{PVE}_m=\lambda_m(R)/p\)
(Proposition~\ref{r12-prop-pve}). By Proposition~\ref{r12-prop-correlation}(b), they do not depend on the units of the
variables, whereas those of \(S\) do: by Corollary~\ref{r12-cor-diagonal}, \(\lambda_1(S)\geqslant s_{jj}\) for every
\(j\), so a variable with a much larger sample variance than the others forces a large first eigenvalue of \(S\).
Variables measured in different units are therefore usually standardized before the analysis; variables measured in
the same units may be left as they are~\citeat{islr}{\S12.2.4}, \citeat{tu-asm}{Proposition~5.24}. For two variables with sample correlation \(r=r_{12}\), for instance,
\(R=\left(\begin{smallmatrix}1&r\\r&1\end{smallmatrix}\right)\) has \(R(1,1)'=(1+r)(1,1)'\) and \(R(1,-1)'=(1-r)(1,-1)'\):
if \(r\ne0\), the principal components of the standardized variables are their sum and their difference, divided by
\(\sqrt2\), the sum first if \(r>0\), and the first one explains the proportion \((1+\lvert r\rvert)/2\).

\begin{example}[cf.\ \citeat{islr}{\S12.2.4, Figure~12.4}]\label{r12-exm-correlation}
The sample covariance matrix of four variables is
\[
S=\begin{pmatrix}4&-2&0&0\\-2&4&0&0\\0&0&16&4\\0&0&4&16\end{pmatrix}
\]
Find the sample correlation matrix, and compare the principal components of the variables with those of the
standardized variables.

\textit{Solution.}
\begin{itemize}
\item The sample standard deviations are \(2\), \(2\), \(4\), \(4\), so \(D_s^{1/2}=\diag(2,2,4,4)\),
\(r_{12}=\frac{-2}{2\cdot2}=-\frac12\), \(r_{34}=\frac4{4\cdot4}=\frac14\), the other \(r_{jl}\) (\(j\ne l\)) are \(0\), and
\[
R=D_s^{-1/2}SD_s^{-1/2}=\begin{pmatrix}1&-\frac12&0&0\\[2pt]-\frac12&1&0&0\\[2pt]0&0&1&\frac14\\[2pt]0&0&\frac14&1\end{pmatrix}
\]
\item The variables. Row by row, \(S(0,0,1,1)'=(0,0,20,20)'\), \(S(0,0,1,-1)'=(0,0,12,-12)'\), \(S(1,-1,0,0)'=(6,-6,0,0)'\)
and \(S(1,1,0,0)'=(2,2,0,0)'\). These four vectors are pairwise orthogonal, so the eigenvalues of \(S\) are \(20,12,6,2\)
(Theorem~\ref{r8-thm-spectral}(c)). They are distinct, so the loading vectors are, up to sign,
\(\frac1{\sqrt2}(0,0,1,1)'\), \(\frac1{\sqrt2}(0,0,1,-1)'\), \(\frac1{\sqrt2}(1,-1,0,0)'\) and \(\frac1{\sqrt2}(1,1,0,0)'\)
(Theorem~\ref{r12-thm-pca}(c),(e)), with sample variances \(20\), \(12\), \(6\), \(2\); as \(\tr S=40\), the proportions of
variance explained are \(\frac{20}{40}=\frac12\), \(\frac{12}{40}=\frac3{10}\), \(\frac6{40}=\frac3{20}\) and
\(\frac2{40}=\frac1{20}\).
\item The standardized variables. In the same way \(R(1,-1,0,0)'=\frac32(1,-1,0,0)'\), \(R(0,0,1,1)'=\frac54(0,0,1,1)'\),
\(R(0,0,1,-1)'=\frac34(0,0,1,-1)'\) and \(R(1,1,0,0)'=\frac12(1,1,0,0)'\). The loading vectors are, up to sign,
\(\frac1{\sqrt2}(1,-1,0,0)'\), \(\frac1{\sqrt2}(0,0,1,1)'\), \(\frac1{\sqrt2}(0,0,1,-1)'\) and \(\frac1{\sqrt2}(1,1,0,0)'\), and, as
\(\tr R=4\), the proportions are \(\frac38\), \(\frac5{16}\), \(\frac3{16}\) and \(\frac18\).
\item Comparison. Variables \(3\) and \(4\) have the largest sample variances, \(16\), and are uncorrelated with variables
\(1\) and \(2\); the eigenvalues \(20\) and \(12\) that belong to them exceed the eigenvalues \(6\) and \(2\) that belong to
variables \(1\) and \(2\). So the first two principal components of the variables are combinations of variables \(3\) and
\(4\) only, and \(\lambda_1(S)=20\geqslant s_{33}=16\), as Corollary~\ref{r12-cor-diagonal} says. Of the standardized
variables, the first principal component is the contrast of variables \(1\) and \(2\), which are more strongly correlated
(\(r_{12}=-\frac12\)) than variables \(3\) and \(4\) (\(r_{34}=\frac14\)). A change of the units of variables \(3\) and \(4\)
would change \(S\), but not \(R\) (Proposition~\ref{r12-prop-correlation}(b)).
\end{itemize}
\end{example}

\subsection*{The closest subspaces}

\begin{definition}[\citeat{tu-asm}{Theorem~5.22}]\label{r12-def-affine}
For \(a\in\R^p\) and a subspace \(W\) of \(\R^p\), the set \(a+W=\{a+w:w\in W\}\) is an \textbf{affine subspace} of
dimension \(\dim W\): a line through \(a\) if \(\dim W=1\), a plane if \(\dim W=2\). The \textbf{distance} from \(y\in\R^p\) to
\(a+W\) is \(\operatorname{dist}(y,a+W)=\min_{w\in W}\lVert y-a-w\rVert\). By Theorem~\ref{r6-thm-best-approximation}, the
minimum is attained exactly at the orthogonal projection \(\hat w\) of \(y-a\) onto \(W\); if the columns of
\(G\in\R^{p\times M}\) form an orthonormal basis of \(W\), then \(\hat w=GG'(y-a)\) (Theorem~\ref{r6-thm-projection-matrix}(a),(b):
\(P_W=GG'\)), and
\[
\operatorname{dist}(y,a+W)^2=\bigl\lVert(I_p-GG')(y-a)\bigr\rVert^2
\]
\end{definition}

\begin{theorem}[\citeat{tu-asm}{Theorem~5.22}]\label{r12-thm-closest}
Let \(X\), \(S\), \(\lambda_m\), \(v_m\), \(z_m\) and \(r\) be as in Theorem~\ref{r12-thm-pca}, let \(1\leqslant M\leqslant p\), and
put \(V_M=[v_1\ \cdots\ v_M]\in\R^{p\times M}\) and \(Z_M=C_nXV_M=[z_1\ \cdots\ z_M]\in\R^{n\times M}\).
\begin{enumerate}
\item For all \(A\in\R^{n\times M}\) and \(B\in\R^{p\times M}\),
\(\lVert C_nX-AB'\rVert_F^2\geqslant(n-1)(\lambda_{M+1}+\dots+\lambda_p)\), with equality for \(A=Z_M\), \(B=V_M\). If
\(M\geqslant r\), then \(C_nX=Z_MV_M'\).
\item \(\lVert C_nX\rVert_F^2=\lVert z_1\rVert^2+\dots+\lVert z_M\rVert^2+\lVert C_nX-Z_MV_M'\rVert_F^2\), so
\[
\mathrm{PVE}_1+\dots+\mathrm{PVE}_M=1-\frac{\lVert C_nX-Z_MV_M'\rVert_F^2}{\lVert C_nX\rVert_F^2}
\]
\item Among all affine subspaces \(a+W\) of \(\R^p\) with \(\dim W=M\), the sum of squared distances
\(\sum_{i=1}^n\operatorname{dist}(x_i,a+W)^2\) is smallest for \(\bar x+\mathcal C(V_M)\), where it equals
\((n-1)(\lambda_{M+1}+\dots+\lambda_p)\). If \(M<p\) and \(\lambda_M>\lambda_{M+1}\), no other \(M\)-dimensional affine
subspace gives this value.
\end{enumerate}
\end{theorem}
\begin{proof}
(a)
\begin{itemize}
\item Equality: \(V_M'V_M=I_M\), so \(N=I_p-V_MV_M'\) is symmetric and idempotent, and \(C_nX-Z_MV_M'=C_nXN\). By
Definition~\ref{r9-def-frobenius}, Proposition~\ref{r12-prop-data}(a), Theorem~\ref{r1-thm-trace}(b) and
\(V_M'SV_M=\diag(\lambda_1,\dots,\lambda_M)\),
\[
\begin{aligned}
\lVert C_nXN\rVert_F^2&=\tr\bigl(N(C_nX)'(C_nX)N\bigr)=(n-1)\tr(SN^2)\\
&=(n-1)\bigl(\tr S-\tr(V_M'SV_M)\bigr)=(n-1)(\lambda_{M+1}+\dots+\lambda_p)
\end{aligned}
\]
using \(\tr S=\lambda_1+\dots+\lambda_p\) (Proposition~\ref{r12-prop-pve}). For \(M\geqslant r\) this is \(0\)
(Theorem~\ref{r12-thm-pca}(a)), so \(C_nX=Z_MV_M'\) (Theorem~\ref{r1-thm-trace}(d)).
\item The bound. Let \(W=\mathcal C(B)\), \(P=P_W\) and \(N_W=I_p-P\), and let \(d=\dim W=\rank B\); then \(d\leqslant M\)
(Definition~\ref{def:r4-rank}, Theorem~\ref{thm:r4-rank-rules}(a)). Row \(i\) of \(AB'\) is \(a_i'B'=(Ba_i)'\), where \(a_i'\)
is row \(i\) of \(A\) (Theorems~\ref{r1-thm-views}(c) and~\ref{r1-thm-transpose}), and \(Ba_i\in W\)
(Theorem~\ref{r1-thm-views}(b)); row \(i\) of \(C_nX\) is \((x_i-\bar x)'\). By Theorem~\ref{r6-thm-best-approximation} and
Theorem~\ref{r6-thm-projection-matrix}(a), \(\lVert(x_i-\bar x)-Ba_i\rVert^2\geqslant\lVert N_W(x_i-\bar x)\rVert^2\). Summing
over \(i\) (Definition~\ref{r9-def-frobenius}), and as row \(i\) of \(C_nXN_W\) is \((N_W(x_i-\bar x))'\) (\(N_W\) is
symmetric, Theorem~\ref{r6-thm-projection-matrix}(c)),
\[
\lVert C_nX-AB'\rVert_F^2\geqslant\sum_{i=1}^n\lVert N_W(x_i-\bar x)\rVert^2=\lVert C_nXN_W\rVert_F^2
\]
\item \(\tr(PS)\leqslant\lambda_1+\dots+\lambda_M\) if \(d=0\) or \(d=M\). If \(d=0\), then \(P=0\), and
\(\lambda_1\geqslant\dots\geqslant\lambda_M\geqslant0\) (Theorem~\ref{r12-thm-pca}(a)). If \(d=M\), let the columns of
\(G\in\R^{p\times M}\) form an orthonormal basis of \(W\) (Corollary~\ref{r6-cor-onb}(a)); then \(P=GG'\)
(Theorem~\ref{r6-thm-projection-matrix}(b)), and \(\tr(PS)=\tr(G'SG)\leqslant\lambda_1+\dots+\lambda_M\)
(Theorem~\ref{r1-thm-trace}(b), and Theorem~\ref{r12-thm-kyfan}(a) for \(S\)).
\item \(\tr(PS)\leqslant\lambda_1+\dots+\lambda_M\) if \(1\leqslant d<M\). Let the columns of \(Q_1\in\R^{p\times d}\) form an
orthonormal basis of \(W\) (Corollary~\ref{r6-cor-onb}(a)), so that \(P=Q_1Q_1'\)
(Theorem~\ref{r6-thm-projection-matrix}(b)). As \(d<p\), Corollary~\ref{r6-cor-onb}(c) gives \(T\) with \([Q_1\ T]\)
orthogonal; let \(Q_2\) consist of the first \(M-d\) columns of \(T\), and \(G=[Q_1\ Q_2]\in\R^{p\times M}\), so that
\(G'G=I_M\) and \(GG'=P+Q_2Q_2'\) (Theorem~\ref{r5b-thm-block}(c)). The diagonal entries \(q'Sq\) of \(Q_2'SQ_2\) are
\(\geqslant0\), as \(S\) is PSD (Proposition~\ref{r12-prop-data}(a)). So by Theorem~\ref{r1-thm-trace}(a),(b) and
Theorem~\ref{r12-thm-kyfan}(a) for \(S\),
\[
\tr(PS)=\tr(GG'S)-\tr(Q_2Q_2'S)=\tr(G'SG)-\tr(Q_2'SQ_2)\leqslant\tr(G'SG)\leqslant\lambda_1+\dots+\lambda_M
\]
\item \(N_W\) is symmetric and idempotent (Theorem~\ref{r6-thm-projection-matrix}(c)); as in the first item,
\(\lVert C_nXN_W\rVert_F^2=(n-1)\tr(SN_W)=(n-1)\bigl(\tr S-\tr(PS)\bigr)\geqslant(n-1)(\lambda_{M+1}+\dots+\lambda_p)\).
\end{itemize}
(b) \(\lVert C_nX\rVert_F^2=(n-1)\tr S=(n-1)(\lambda_1+\dots+\lambda_p)\) (Proposition~\ref{r12-prop-data}(a),
Proposition~\ref{r12-prop-pve}) and \(\lVert z_m\rVert^2=(n-1)\lambda_m\) (Theorem~\ref{r12-thm-pca}(d)); with (a), both
sides of the identity equal \((n-1)\tr S\). Divide by \(\lVert C_nX\rVert_F^2\) and use the first expression of
\(\mathrm{PVE}_m\) in Proposition~\ref{r12-prop-pve}.

(c) Let \(\dim W=M\), let the columns of \(G\in\R^{p\times M}\) form an orthonormal basis of \(W\)
(Theorems~\ref{thm:r4-dimension} and~\ref{r6-thm-gram-schmidt}), and put \(N=I_p-GG'\), symmetric and idempotent.
\begin{itemize}
\item By Definition~\ref{r12-def-affine}, \(\operatorname{dist}(x_i,a+W)^2=\lVert N(x_i-a)\rVert^2\). Write
\(x_i-a=(x_i-\bar x)+(\bar x-a)\). As \(\sum_i(x_i-\bar x)=X'\one-n\bar x=0\), expanding the squares
(Proposition~\ref{r6-prop-pythagoras}) gives
\[
\sum_{i=1}^n\operatorname{dist}(x_i,a+W)^2=\sum_{i=1}^n\lVert N(x_i-\bar x)\rVert^2+n\lVert N(\bar x-a)\rVert^2
\]
\item Row \(i\) of \(C_nXN\) is \((x_i-\bar x)'N=(N(x_i-\bar x))'\), so, as in (a), and by
Theorem~\ref{r12-thm-kyfan}(a) for \(S\),
\[
\begin{aligned}
\sum_i\lVert N(x_i-\bar x)\rVert^2&=\lVert C_nXN\rVert_F^2=(n-1)\tr(SN)=(n-1)\bigl(\tr S-\tr(G'SG)\bigr)\\
&\geqslant(n-1)(\lambda_{M+1}+\dots+\lambda_p)
\end{aligned}
\]
\item Hence the sum of squared distances is at least \((n-1)(\lambda_{M+1}+\dots+\lambda_p)\). For \(a=\bar x\) and
\(G=V_M\), the second term is \(0\) and \(\tr(V_M'SV_M)=\lambda_1+\dots+\lambda_M\): the bound is attained.
\item If the bound is attained, then \(N(\bar x-a)=0\) and \(\tr(G'SG)=\lambda_1+\dots+\lambda_M\). If \(M<p\) and
\(\lambda_M>\lambda_{M+1}\), Theorem~\ref{r12-thm-kyfan}(c) gives \(W=\mathcal C(G)=\mathcal C(V_M)\). And
\(N(\bar x-a)=0\) means \(\bar x-a=GG'(\bar x-a)\in W\), so \(a+W=\bar x+W\).\qedhere
\end{itemize}
\end{proof}

Part (a) says that the matrix \(Z_MV_M'\), of rank at most \(M\), built from the first \(M\) principal components is a best
approximation of the centered data among the products \(AB'\) with \(M\) columns in \(A\) and \(B\). Every matrix of rank
\(s\leqslant M\) is such a product (\(0=0\cdot0'\), and for \(s\geqslant1\) a full-rank factorization,
Theorem~\ref{thm:r4-full-rank-factorization}(a), with zero columns added), so \(Z_MV_M'\) is also a best approximation
among the matrices of rank at most \(M\). Part (b) says that the cumulative proportion of variance explained is the part
of \(\lVert C_nX\rVert_F^2\) that it reproduces.

\begin{example}\label{r12-exm-closest}
For the data of Example~\ref{r12-exm-data}, \(\lVert C_5X\rVert_F^2=40+40+40+40=160=4\cdot40\)
(Example~\ref{r12-exm-pca}: \(\tr S=40\)).
\begin{itemize}
\item \(M=1\): \(Z_1V_1'=z_1v_1'=\frac12(0,5,5,-5,-5)'(1,1,1,1)\) has the rows \(0'\), \(\frac52\one'\) twice and
\(-\frac52\one'\) twice, and
\[
C_5X-Z_1V_1'=\begin{pmatrix}2&-2&2&-2\\\frac32&\frac32&-\frac52&-\frac12\\[2pt]-\frac52&-\frac12&\frac32&\frac32\\[2pt]\frac12&\frac52&-\frac32&-\frac32\\[2pt]-\frac32&-\frac32&\frac12&\frac52\end{pmatrix}
\]
with the squared row lengths \(16\), \(11\), \(11\), \(11\), \(11\); so
\(\lVert C_5X-Z_1V_1'\rVert_F^2=16+11+11+11+11=60=4(9+5+1)\). By Theorem~\ref{r12-thm-closest}(c), the line \(\bar x+\Span\{(1,1,1,1)'\}\) is the line closest to the five points (the
only one, as \(25>9\)), with the sum of squared distances \(60\); and \(1-\frac{60}{160}=\frac58=\mathrm{PVE}_1\)
(Theorem~\ref{r12-thm-closest}(b)).
\item \(M=2\): \(Z_2V_2'=z_1v_1'+z_2v_2'\), with \(z_2v_2'=\frac12(0,3,-3,3,-3)'(1,1,-1,-1)\). So
\[
Z_2V_2'=\begin{pmatrix}0&0&0&0\\4&4&1&1\\1&1&4&4\\-1&-1&-4&-4\\-4&-4&-1&-1\end{pmatrix},\qquad
C_5X-Z_2V_2'=\begin{pmatrix}2&-2&2&-2\\0&0&-1&1\\-1&1&0&0\\-1&1&0&0\\0&0&-1&1\end{pmatrix}
\]
\(\lVert C_5X-Z_2V_2'\rVert_F^2=16+2+2+2+2=24=4(5+1)\), and \(1-\frac{24}{160}=\frac{17}{20}=\mathrm{PVE}_1+\mathrm{PVE}_2\).
By Theorem~\ref{r12-thm-closest}(a) with \(M=4\geqslant r\), \(C_5X-Z_2V_2'=z_3v_3'+z_4v_4'\); for instance its first row is
\(4\cdot\frac12(1,-1,1,-1)+0\).
\item The closest plane (Theorem~\ref{r12-thm-closest}(c), unique as \(9>5\)) is \(\bar x+\Span\{v_1,v_2\}\). As
\(\Span\{v_1,v_2\}=\Span\{v_3,v_4\}^\perp\) (the first lies in the second, and both have dimension \(2\):
Theorems~\ref{r6-thm-complement}(c) and~\ref{thm:r4-counting}(a)), it is
\(\{y\in\R^4:h_3'(y-\bar x)=h_4'(y-\bar x)=0\}=\{y:y_1=y_2,\ y_3=y_4\}\), since \(h_3'u+h_4'u=2(u_1-u_2)\) and
\(h_3'u-h_4'u=2(u_3-u_4)\). The point of the plane nearest to unit \(i\) is \(\bar x+V_2V_2'(x_i-\bar x)\)
(Definition~\ref{r12-def-affine}), whose transpose is \(\bar x'\) plus row \(i\) of \(Z_2V_2'\): unit \(1\), \((7,3,7,3)\), has
the nearest point \((5,5,5,5)\), at squared distance \(16\); units \(2\)--\(5\) have the nearest points \((9,9,6,6)\),
\((6,6,9,9)\), \((4,4,1,1)\), \((1,1,4,4)\), at squared distance \(2\) each.
\end{itemize}
\end{example}

\subsection*{Principal components of a random vector}

\begin{definition}[\citeat{bapat}{\S10.1}]\label{r12-def-pop-pc}
Let \(x=(x_1,\dots,x_p)'\) be a random vector with covariance matrix \(\Sigma=\var(x)\ne0\)
(Definition~\ref{r11-def-random}). A \textbf{first principal component} of \(x\) is a random variable \(a_1'x\), where the
unit vector \(a_1\in\R^p\) maximizes \(\var(a'x)\) over all unit vectors \(a\in\R^p\). For \(m\geqslant2\),
once \(a_1,\dots,a_{m-1}\) are chosen, an \textbf{\(m\)th principal component} is \(a_m'x\), where the unit vector \(a_m\)
maximizes \(\var(a'x)\) over the unit vectors \(a\) for which \(a'x\) is uncorrelated with (has covariance
\(0\) with) each of \(a_1'x,\dots,a_{m-1}'x\).
\end{definition}

\begin{theorem}[\citeat{bapat}{\S10.1}]\label{r12-thm-pop-pca}
Let \(x\) and \(\Sigma\) be as in Definition~\ref{r12-def-pop-pc}, and let \(\Sigma=V\Lambda V'\), where
\(V=[v_1\ \cdots\ v_p]\) is orthogonal and \(\Lambda=\diag(\lambda_1,\dots,\lambda_p)\) with
\(\lambda_1\geqslant\dots\geqslant\lambda_p\) (Lemma~\ref{r9-lem-eig-order}(a)).
\begin{enumerate}
\item For \(a,b\in\R^p\), \(\var(a'x)=a'\Sigma a\) and \(\cov(a'x,b'x)=a'\Sigma b\); and
\(\lambda_p\geqslant0\).
\item \(v_1'x\) is a first principal component, with variance \(\lambda_1\). Let \(m\geqslant2\) and \(\lambda_{m-1}>0\). If
\(v_1'x,\dots,v_{m-1}'x\) are taken as the first \(m-1\) principal components, then \(v_m'x\) is an \(m\)th principal
component, with variance \(\lambda_m\).
\item The principal components \(v_1'x,\dots,v_p'x\) are pairwise uncorrelated, and the \textbf{total variance} of \(x\) is
\(\var(x_1)+\dots+\var(x_p)=\tr\Sigma=\lambda_1+\dots+\lambda_p\). So the first \(M\) principal
components account for the proportion \((\lambda_1+\dots+\lambda_M)/\tr\Sigma\) of the total variance.
\end{enumerate}
\end{theorem}
\needspace{4\baselineskip}
\begin{proof}
(a)
\begin{itemize}
\item Let \(T\in\R^{2\times p}\) have the rows \(a'\) and \(b'\), so that \(Tx=(a'x,b'x)'\) (Theorem~\ref{r1-thm-views}(c)).
By Theorem~\ref{r11-thm-linear}(b), \(\var(Tx)=T\Sigma T'\), whose entries are \(a'\Sigma a\), \(a'\Sigma b\) in
the first row and \(b'\Sigma a\), \(b'\Sigma b\) in the second (Theorem~\ref{r1-thm-views}(d)).
\item By Definition~\ref{r11-def-random}, the diagonal entries of \(\var(Tx)\) are \(\var(a'x)\) and
\(\var(b'x)\), and its \((1,2)\) entry is \(\cov(a'x,b'x)\). (This is also
Theorem~\ref{r14-thm-cov}(a) for \(u=v=x\) with the rows \(A=a'\) and \(C=b'\).)
\item \(\Sigma\) is PSD (Theorem~\ref{r11-thm-linear}(c)), so \(\lambda_p\geqslant0\) (Theorem~\ref{r8-thm-eigen-criterion}(a)).
\end{itemize}
(b)
\begin{itemize}
\item By (a), the first problem is to maximize \(a'\Sigma a\) over the unit vectors \(a\), and for \(a\ne0\) the Rayleigh
quotient \(a'\Sigma a/a'a\) is the value at the unit vector \(a/\lVert a\rVert\). By Theorem~\ref{r9-thm-rayleigh}(a) the
maximum is \(\lambda_1\), attained at \(a=v_1\).
\item Let \(m\geqslant2\) and \(\lambda_{m-1}>0\); then \(\lambda_l>0\) for every \(l<m\). By (a) and \(\Sigma v_l=\lambda_lv_l\),
\(\cov(a'x,v_l'x)=a'\Sigma v_l=\lambda_la'v_l\), which is \(0\) if and only if \(a\perp v_l\). So the \(m\)th
problem is to maximize \(a'\Sigma a\) over the unit vectors \(a\perp v_1,\dots,v_{m-1}\).
\item By Theorem~\ref{r9-thm-constrained}, the largest value of \(a'\Sigma a/a'a\) over the nonzero
\(a\perp v_1,\dots,v_{m-1}\) is \(\lambda_m\), attained at \(a=v_m\); so \(v_m\) solves the \(m\)th problem, and
\(\var(v_m'x)=v_m'\Sigma v_m=\lambda_m\).
\end{itemize}
\needspace{4\baselineskip}
(c)
\begin{itemize}
\item By (a), \(\cov(v_m'x,v_l'x)=v_m'\Sigma v_l=\lambda_lv_m'v_l=0\) for \(l\ne m\).
\item \(\var(x_j)=\sigma_{jj}\) (Definition~\ref{r11-def-random}), so the total variance is
\(\tr\Sigma=\tr(V\Lambda V')=\tr(\Lambda V'V)=\tr\Lambda=\lambda_1+\dots+\lambda_p\) (Theorem~\ref{r1-thm-trace}(b)). It is
positive, so the proportions are defined: every \(\lambda_m\geqslant0\) by (a), and \(\Lambda\ne0\), as
\(\Sigma=V\Lambda V'\ne0\).\qedhere
\end{itemize}
\end{proof}

Parts (b) and (c) are Theorem~\ref{r12-thm-pca}(c),(d) and Proposition~\ref{r12-prop-pve} with \(\Sigma\) in place of
\(S\). In practice \(\Sigma\) is not known: \(n\) observations of \(x\) form the rows of a data matrix, and the principal
components of the data (Definition~\ref{r12-def-pc}), computed from \(S\), take the place of those of \(x\)~cf.\ \citeat{bapat}{\S10.1}, \citeat{searle-khuri}{\S10.3}.

\begin{example}[cf.\ \citeat{searle-khuri}{Example~10.1}]\label{r12-exm-equicorrelation}
Let \(x=(x_1,\dots,x_p)'\) (\(p\geqslant2\)) have \(\var(x_i)=\sigma^2>0\) and
\(\cov(x_i,x_j)=\rho\sigma^2\) for \(i\ne j\), where \(0<\rho<1\): \(\Sigma=\sigma^2\bigl((1-\rho)I_p+\rho J_p\bigr)\),
and \(\cor(x)\) is the equicorrelation matrix of Example~\ref{r5d-exm-equicorrelation}. Find the principal
components of \(x\).

\textit{Solution.}
\begin{itemize}
\item By Example~\ref{r8-exm-patterned}(c), \(\Sigma=\sigma^2\bigl(1+(p-1)\rho\bigr)\bar J_p+\sigma^2(1-\rho)C_p\), and by
Example~\ref{r8-exm-patterned}(b) its eigenvalues are \(\lambda_1=\sigma^2\bigl(1+(p-1)\rho\bigr)\), with eigenspace
\(\Span\{\one_p\}\), and \(\lambda_2=\dots=\lambda_p=\sigma^2(1-\rho)\), with eigenspace \(\Span\{\one_p\}^\perp\). As
\(0<\rho<1\), \(\lambda_1>\lambda_2>0\).
\item First component. With \(v_1=\frac1{\sqrt p}\one_p\), \(v_1'x=\frac1{\sqrt p}(x_1+\dots+x_p)\), a multiple of the average
of the \(p\) variables, with variance \(\sigma^2\bigl(1+(p-1)\rho\bigr)\) (Theorem~\ref{r12-thm-pop-pca}(b)). It is unique up
to sign: a unit vector \(a\) with \(a'\Sigma a=\lambda_1\) is an eigenvector for \(\lambda_1\)
(Theorem~\ref{r9-thm-rayleigh}(b)), hence \(a=\pm v_1\).
\item Second component. \(\cov(a'x,v_1'x)=\lambda_1a'v_1\) (Theorem~\ref{r12-thm-pop-pca}(a)), so the
constraint of the second problem is \(\one_p'a=0\); and for every such unit vector \(a\), \(\Sigma a=\lambda_2a\) and
\(\var(a'x)=\lambda_2\). So every \(a'x\) with \(\one_p'a=0\) and \(\lVert a\rVert=1\) is a second principal
component: for \(p\geqslant3\) the second component is not unique, even up to sign (for \(p=2\) it is
\(\pm(x_1-x_2)/\sqrt2\)).
\item The first component accounts for the proportion \(\lambda_1/\tr\Sigma=\bigl(1+(p-1)\rho\bigr)/p\) of the total variance
\(\tr\Sigma=p\sigma^2\) (Theorem~\ref{r12-thm-pop-pca}(c)); for \(p=4\) and \(\rho=\frac12\) it is \(\frac58\).
\end{itemize}
\end{example}

\needspace{20\baselineskip}
\subsection*{Exercises}

\begin{exercise}\label{r11-dr-random}
Let \(z=(z_1,z_2,z_3,z_4)'\) be a random vector, and let
\[
\mu=\expc(z)=\begin{pmatrix}1\\-1\\2\\1\end{pmatrix},\qquad
\Sigma=\var(z)=\begin{pmatrix}2&1&0&0\\1&2&1&0\\0&1&2&1\\0&0&1&2\end{pmatrix},\qquad
A=\begin{pmatrix}1&1&0&0\\0&0&1&-1\end{pmatrix},\qquad b=\begin{pmatrix}-1\\3\end{pmatrix}
\]
\begin{enumerate}
\item Find \(\expc(Az+b)\) and \(\var(Az+b)\), and \(\cov(z_1+z_2,\ z_3-z_4)\).
\item Find \(\var(z_1+z_2+z_3+z_4)\) and \(\var(z_1-z_4)\).
\item Find \(\expc(z'z)\) and \(\expc(2z_1z_2+z_3^2+2z_3z_4)\).
\end{enumerate}
\end{exercise}
\begin{solution}
(a) By Theorem~\ref{r11-thm-linear}(b),
\[
\expc(Az+b)=A\mu+b=\begin{pmatrix}1-1\\2-1\end{pmatrix}+\begin{pmatrix}-1\\3\end{pmatrix}=\begin{pmatrix}-1\\4\end{pmatrix}
\]
and, with \(A\Sigma\) computed row by row (its rows are \(\row_1(\Sigma)+\row_2(\Sigma)\) and \(\row_3(\Sigma)-\row_4(\Sigma)\),
Theorem~\ref{r1-thm-views}(c)),
\[
A\Sigma=\begin{pmatrix}3&3&1&0\\0&1&1&-1\end{pmatrix},\qquad
\var(Az+b)=A\Sigma A'=\begin{pmatrix}6&1\\1&2\end{pmatrix}
\]
With \(b=0\), Theorem~\ref{r11-thm-linear}(b) gives \(\var(Az)=A\Sigma A'\) as well. The entries of \(Az\) are
\(z_1+z_2\) and \(z_3-z_4\), so their covariance is the \((1,2)\) entry of \(A\Sigma A'\) (Definition~\ref{r11-def-random}):
\(\cov(z_1+z_2,\ z_3-z_4)=1\).

(b) By Theorem~\ref{r11-thm-linear}(b) with the \(1\times4\) matrix \(a'\), \(\var(a'z)=a'\Sigma a\). For
\(a=\one_4\) this is the sum of all entries of \(\Sigma\): \(\var(z_1+z_2+z_3+z_4)=3+4+4+3=14\), the row sums
of \(\Sigma\) being \(3,4,4,3\). For \(a=(1,0,0,-1)'\),
\(\var(z_1-z_4)=\sigma_{11}-2\sigma_{14}+\sigma_{44}=2-0+2=4\).

(c) By Theorem~\ref{r11-thm-quadratic} with \(A=I_4\), \(\expc(z'z)=\tr\Sigma+\mu'\mu=(2+2+2+2)+(1+1+4+1)=15\).
Next, \(2z_1z_2+z_3^2+2z_3z_4=z'Bz\) for the symmetric \(B\) below (Definition~\ref{r8-def-quadratic-form}), and
\[
\begin{gathered}
B=\begin{pmatrix}0&1&0&0\\1&0&0&0\\0&0&1&1\\0&0&1&0\end{pmatrix}
\\
\tr(B\Sigma)=2\sigma_{12}+\sigma_{33}+2\sigma_{34}=2+2+2=6,\qquad
\mu'B\mu=2\mu_1\mu_2+\mu_3^2+2\mu_3\mu_4=-2+4+4=6
\end{gathered}
\]
so \(\expc(2z_1z_2+z_3^2+2z_3z_4)=\tr(B\Sigma)+\mu'B\mu=12\). Check: by Theorem~\ref{r14-thm-cov}(e), whose
\((i,j)\) entry reads \(\expc(z_iz_j)=\sigma_{ij}+\mu_i\mu_j\), \(\expc(z_1z_2)=1-1=0\),
\(\expc(z_3^2)=2+4=6\) and \(\expc(z_3z_4)=1+2=3\), and \(2\cdot0+6+2\cdot3=12\).
\end{solution}

\begin{exercise}\label{r11-ex-moments-drill}
Let \(x\in\R^4\) be a random vector with \(\expc(x)=\mu=(1,0,2,-1)'\) and
\[
\var(x)=\Sigma=\begin{pmatrix}9&0&3&3\\0&1&2&1\\3&2&9&5\\3&1&5&4\end{pmatrix},\qquad
L=\begin{pmatrix}3&0&0&0\\0&1&0&0\\1&2&2&0\\1&1&1&1\end{pmatrix}
\]
\begin{enumerate}
\item Show that \(\Sigma=LL'\), and deduce that there is no nonzero \(a\in\R^4\) with \(a'x=a'\mu\) with probability \(1\).
\item Find \(\expc(y)\) and \(\var(y)\) for \(y=(x_1-x_2,\ x_3-x_4-3)'\).
\item Find \(\cov(x_1+x_3,\ x_2-x_4)\) and \(\cor(x)\).
\end{enumerate}
\end{exercise}
\begin{solution}
(a) The \((i,j)\) entry of \(LL'\) is \(\row_i(L)\row_j(L)'\) (Theorem~\ref{r1-thm-views}(d)), and \(LL'\) is symmetric
(Theorem~\ref{r1-thm-transpose}). With the rows \((3,0,0,0)\), \((0,1,0,0)\), \((1,2,2,0)\), \((1,1,1,1)\) this gives
\(9,0,3,3\) in row \(1\); \(1,2,1\) at \((2,2)\), \((2,3)\), \((2,4)\); \(1+4+4=9\) and \(1+2+2=5\) at \((3,3)\), \((3,4)\);
and \(1+1+1+1=4\) at \((4,4)\). So \(LL'=\Sigma\). As \(L\) is lower triangular, \(\det L=3\cdot1\cdot2\cdot1=6\ne0\)
(Theorem~\ref{r5d-thm-transpose}), so \(P=L'\) is invertible (Theorem~\ref{r5d-thm-product}(a)), and \(\Sigma=P'P\) is PD
(Theorem~\ref{r8-thm-psd-factor}(c)), hence invertible (Theorem~\ref{r8-thm-eigen-criterion}(c)). By
Theorem~\ref{r11-thm-cov-psd}, there is no nonzero \(a\) with \(a'x=a'\mu\) with probability \(1\).

(b) \(y=Ax+b\) with \(A=\left(\begin{smallmatrix}1&-1&0&0\\0&0&1&-1\end{smallmatrix}\right)\) and \(b=(0,-3)'\). By
Theorem~\ref{r11-thm-linear}(b), \(\expc(y)=A\mu+b=(1-0,\ 2+1-3)'=(1,0)'\), and the rows of \(A\Sigma\) are
\(\row_1(\Sigma)-\row_2(\Sigma)\) and \(\row_3(\Sigma)-\row_4(\Sigma)\):
\[
A\Sigma=\begin{pmatrix}9&-1&1&2\\0&1&4&1\end{pmatrix},\qquad
\var(y)=A\Sigma A'=\begin{pmatrix}10&-1\\-1&3\end{pmatrix}
\]
(c) \(\Sigma(0,1,0,-1)'=\col_2(\Sigma)-\col_4(\Sigma)=(-3,0,-3,-3)'\), so by Theorem~\ref{r14-thm-cov}(a)
\(\cov(x_1+x_3,\ x_2-x_4)=(1,0,1,0)(-3,0,-3,-3)'=-6\). The standard deviations are \(3,1,3,2\), so
\(D^{-1/2}=\diag(\frac13,1,\frac13,\frac12)\) and (Definition~\ref{r14-def-corr})
\[
\cor(x)=D^{-1/2}\Sigma D^{-1/2}=\begin{pmatrix}1&0&\frac13&\frac12\\[2pt]0&1&\frac23&\frac12\\[2pt]\frac13&\frac23&1&\frac56\\[2pt]\frac12&\frac12&\frac56&1\end{pmatrix}
\]
\end{solution}

\begin{exercise}\label{r14-dr-cross}
Let \(z\in\R^4\) be a random vector, and let \(u=Az\) and \(v=\one'z=z_1+z_2+z_3+z_4\), where
\[
\var(z)=\Sigma=\begin{pmatrix}4&-1&-1&-1\\-1&4&-1&-1\\-1&-1&4&-1\\-1&-1&-1&4\end{pmatrix},\qquad
A=\begin{pmatrix}1&-1&0&0\\1&1&1&2\end{pmatrix}
\]
so that \(u=(z_1-z_2,\ z_1+z_2+z_3+2z_4)'\).
\begin{enumerate}
\item Find \(\cov(u,z)\) and \(\cov(u,v)\).
\item Find \(\var(u)\), \(\var(v)\) and the covariance matrix of \((u',v)'\in\R^3\).
\item Find the correlation matrix of \(z\), and the correlation of \(u_2\) and \(v\).
\end{enumerate}
\end{exercise}
\begin{solution}
(a) By Theorem~\ref{r14-thm-cov}(a) with \(B=I_4\), computing \(A\Sigma\) row by row,
\[
\begin{aligned}
\cov(u,z)=A\Sigma&=\begin{pmatrix}4+1&-1-4&-1+1&-1+1\\4-1-1-2&-1+4-1-2&-1-1+4-2&-1-1-1+8\end{pmatrix}\\
&=\begin{pmatrix}5&-5&0&0\\0&0&0&5\end{pmatrix}
\end{aligned}
\]
and, with \(B=\one'\), \(\cov(u,v)=A\Sigma\one=(5-5+0+0,\ 0+0+0+5)'=(0,5)'\). So \(z_1-z_2\) and \(v\) are
uncorrelated.

(b) By Theorem~\ref{r11-thm-linear}(b), \(\var(u)=A\Sigma A'=(A\Sigma)A'\), and
\(\var(v)=\one'\Sigma\one\), the sum of all entries of \(\Sigma\):
\[
\var(u)=\begin{pmatrix}5+5+0+0&5-5+0+0\\0+0+0+0&0+0+0+10\end{pmatrix}=\begin{pmatrix}10&0\\0&10\end{pmatrix},\qquad
\var(v)=16-12=4
\]
By the block form of the covariance matrix of a stacked vector (after Definition~\ref{r14-def-cov}),
\[
\var\begin{pmatrix}u\\v\end{pmatrix}=\begin{pmatrix}\var(u)&\cov(u,v)\\\cov(v,u)&\var(v)\end{pmatrix}
=\begin{pmatrix}10&0&0\\0&10&5\\0&5&4\end{pmatrix}
\]
(c) Every \(\sigma_{ii}=4\), so \(D^{-1/2}=\frac12I_4\) and (Definition~\ref{r14-def-corr})
\[
\cor(z)=\frac14\Sigma=\begin{pmatrix}1&-\frac14&-\frac14&-\frac14\\[2pt]-\frac14&1&-\frac14&-\frac14\\[2pt]-\frac14&-\frac14&1&-\frac14\\[2pt]-\frac14&-\frac14&-\frac14&1\end{pmatrix}
\]
The correlation of \(u_2\) and \(v\) is
\(\cov(u_2,v)/\sqrt{\var(u_2)\var(v)}=5/\sqrt{10\cdot4}=5/(2\sqrt{10})=\sqrt{10}/4\).
\end{solution}

\begin{exercise}\label{r14-dr-corr}
Find the correlation matrix of a random vector \(z\in\R^4\) with
\[
\var(z)=\begin{pmatrix}9&3&-6&9\\3&4&2&6\\-6&2&16&-12\\9&6&-12&36\end{pmatrix}
\]
\end{exercise}
\begin{solution}
The standard deviations are \(3\), \(2\), \(4\), \(6\), so \(\rho_{12}=\frac3{3\cdot2}=\frac12\),
\(\rho_{13}=\frac{-6}{3\cdot4}=-\frac12\), \(\rho_{14}=\frac9{3\cdot6}=\frac12\), \(\rho_{23}=\frac2{2\cdot4}=\frac14\),
\(\rho_{24}=\frac6{2\cdot6}=\frac12\), \(\rho_{34}=\frac{-12}{4\cdot6}=-\frac12\), and (Definition~\ref{r14-def-corr})
\[
\cor(z)=\begin{pmatrix}1&\frac12&-\frac12&\frac12\\[2pt]\frac12&1&\frac14&\frac12\\[2pt]-\frac12&\frac14&1&-\frac12\\[2pt]\frac12&\frac12&-\frac12&1\end{pmatrix}
\]
\end{solution}

\begin{exercise}[\citeat{searle-khuri}{Exercise~10.5(b)}]\label{r12-ex-variances}
Let \(x=(x_1,x_2,x_3)'\) be a random vector with the covariance matrix
\[
\var(x)=V=\begin{pmatrix}8&-7&2\\-7&10&-3\\2&-3&1\end{pmatrix}
\]
Show that \(x_1+2x_2+3x_3\) and \(x_2+4x_3\) have the variances \(5\) and \(2\). Are they correlated?
\end{exercise}
\begin{solution}
Let \(a=(1,2,3)'\) and \(b=(0,1,4)'\). By Theorem~\ref{r12-thm-pop-pca}(a), \(\var(a'x)=a'Va\),
\(\var(b'x)=b'Vb\) and \(\cov(a'x,b'x)=a'Vb\). Here
\(Va=(8-14+6,\ -7+20-9,\ 2-6+3)'=(0,4,-1)'\) and \(Vb=(0-7+8,\ 0+10-12,\ 0-3+4)'=(1,-2,1)'\), so
\[
\var(x_1+2x_2+3x_3)=a'Va=0+8-3=5,\qquad\var(x_2+4x_3)=b'Vb=0-2+4=2
\]
and \(\cov(x_1+2x_2+3x_3,\ x_2+4x_3)=a'Vb=1-4+3=0\). So the two combinations are uncorrelated.
\end{solution}

\begin{exercise}\label{r11-ex-quadratic-blp}
Let \(x\) be as in Exercise~\ref{r11-ex-moments-drill}. Use Theorem~\ref{r11-thm-quadratic} to find
\(\expc(x_1x_3+x_2^2)\) and \(\expc(\lVert x\rVert^2)\), and check the first value directly from the means,
variances and covariances.
\end{exercise}
\begin{solution}
\(x_1x_3+x_2^2=x'Ax\) with \(A=e_1e_3'+e_2e_2'\). The rows of \(A\Sigma\) are \(\row_3(\Sigma)\), \(\row_2(\Sigma)\), \(0\),
\(0\) (Theorem~\ref{r1-thm-views}(c)), so \(\tr(A\Sigma)=\sigma_{31}+\sigma_{22}=3+1=4\); and \(\mu'A\mu=\mu_1\mu_3+\mu_2^2=2\).
Hence \(\expc(x_1x_3+x_2^2)=4+2=6\). Directly, by Theorem~\ref{r14-thm-cov}(e), whose \((i,j)\) entry reads
\(\expc(x_ix_j)=\sigma_{ij}+\mu_i\mu_j\), \(\expc(x_1x_3)=\cov(x_1,x_3)+\mu_1\mu_3=3+2=5\) and
\(\expc(x_2^2)=\var(x_2)+\mu_2^2=1\), with sum \(6\). With \(A=I_4\),
\(\expc(\lVert x\rVert^2)=\tr\Sigma+\mu'\mu=23+6=29\).
\end{solution}

\begin{exercise}\label{r8-ex-correlation-valid}
Three stocks each have correlation \(\frac23\) with a market index, and any two of the stocks have correlation \(r\). The
correlation matrix of (index, stock \(1\), stock \(2\), stock \(3\)) is
\[
R_r=\begin{pmatrix}1&\frac23&\frac23&\frac23\\[2pt]\frac23&1&r&r\\[2pt]\frac23&r&1&r\\[2pt]\frac23&r&r&1\end{pmatrix}
\]
\begin{enumerate}
\item Show that \(R_r\) is PSD if and only if \(\frac16\leqslant r\leqslant1\), and PD if and only if \(\frac16<r<1\).
\item An analyst reports the sample correlations \(\frac23,\frac23,\frac23\) of three series with a fourth, and reports the
three series as pairwise uncorrelated, all observed over the same periods. Explain why the report must contain an error.
\end{enumerate}
\end{exercise}
\begin{solution}
(a) Write \(R_r=\left(\begin{smallmatrix}1&b'\\b&D\end{smallmatrix}\right)\) with \(b=\frac23\one_3\) and
\(D=(1-r)I_3+rJ_3\). The block \((1)\) is PD, and its Schur complement is
\[
R_r/(1)=D-bb'=\begin{pmatrix}1&r&r\\r&1&r\\r&r&1\end{pmatrix}-\tfrac49J_3=(1-r)I_3+\bigl(r-\tfrac49\bigr)J_3
\]
For \(r\ne\frac49\), by Example~\ref{r8-exm-patterned}(b) its eigenvalues are \((1-r)+3\bigl(r-\frac49\bigr)=2r-\frac13\)
(eigenvector \(\one_3\)) and \(1-r\); for \(r=\frac49\) it is \(\frac59I_3\), and \(2r-\frac13=1-r=\frac59\). By
Theorem~\ref{r8-thm-schur-pd} and Theorem~\ref{r8-thm-eigen-criterion}(a), \(R_r\) is PSD if and only if
\(2r-\frac13\geqslant0\) and \(1-r\geqslant0\), that is, \(\frac16\leqslant r\leqslant1\); and PD if and only if both are
positive, that is, \(\frac16<r<1\).

(b) A sample correlation matrix is PSD (Proposition~\ref{r12-prop-correlation}(a)), and \(0<\frac16\), so by (a) the
reported matrix \(R_0\) is not; for instance \(x=(-2,1,1,1)'\) gives \(x'R_0x=4+3-2\cdot\frac23\cdot2\cdot3=-1<0\). More
generally, for any sample correlations \(r_{12},r_{13},r_{23}\) among the three series, the same \(x\) gives
\(0\leqslant7-8+2(r_{12}+r_{13}+r_{23})\): three series that each correlate \(\frac23\) with a fourth have pairwise
correlations adding up to at least \(\frac12\).
\end{solution}

\begin{exercise}\label{r8-ex-portfolio}
The monthly returns (in \%) of four funds over five months are:
\[
\begin{array}{c|ccccc}\text{month}&1&2&3&4&5\\\hline
\text{fund 1}&2&2&6&10&10\\\text{fund 2}&8&12&12&8&10\\\text{fund 3}&9&7&5&9&10\\\text{fund 4}&5&7&9&9&10\end{array}
\]
Let \(X\in\R^{5\times4}\) have the months as rows.
\begin{enumerate}
\item Compute the mean vector \(\bar x\), the sample covariance matrix \(S\) and the sample correlation matrix \(R\).
\item The equally weighted portfolio has the returns \(a'x_i\) with \(a=\frac14(1,1,1,1)'\). Compute their sample
variance as \(a'Sa\), and directly from the five portfolio returns.
\item Show that \(S\) is singular, find \(a\ne0\) with \(a'Sa=0\), and explain this from the data. What is \(\rank(C_5X)\)?
\end{enumerate}
\end{exercise}
\begin{solution}
(a) \(\bar x=\frac15(30,50,40,40)'=(6,10,8,8)'\), and (Proposition~\ref{r12-prop-data}(a))
\[
\begin{gathered}
C_5X=\begin{pmatrix}-4&-2&1&-3\\-4&2&-1&-1\\0&2&-3&1\\4&-2&1&1\\4&0&2&2\end{pmatrix}
\\
X'C_5X=\begin{pmatrix}64&-8&12&28\\-8&16&-12&4\\12&-12&16&0\\28&4&0&16\end{pmatrix},\qquad
S=\frac14X'C_5X=\begin{pmatrix}16&-2&3&7\\-2&4&-3&1\\3&-3&4&0\\7&1&0&4\end{pmatrix}
\end{gathered}
\]
(for instance \(64=16+16+0+16+16\) and \(-8=8-8+0-8+0\)). With \(D_s=\diag(16,4,4,4)\), \(r_{12}=\frac{-2}{4\cdot2}=-\frac14\),
\(r_{13}=\frac3{4\cdot2}=\frac38\), \(r_{14}=\frac7{4\cdot2}=\frac78\), \(r_{23}=\frac{-3}{2\cdot2}=-\frac34\),
\(r_{24}=\frac1{2\cdot2}=\frac14\) and \(r_{34}=0\) (Definition~\ref{r12-def-correlation}), so
\[
R=\begin{pmatrix}1&-\frac14&\frac38&\frac78\\[2pt]-\frac14&1&-\frac34&\frac14\\[2pt]\frac38&-\frac34&1&0\\[2pt]\frac78&\frac14&0&1\end{pmatrix}
\]
(b) \(a'Sa=\frac1{16}\one'S\one\), and the entries of \(S\) add up to \(40\); so \(a'Sa=\frac52\)
(Proposition~\ref{r12-prop-data}(b)). Directly, the portfolio returns are \(6,7,8,9,10\), with mean \(8\) and squared
deviations \(4,1,0,1,4\); their sample variance is \(\frac{10}4=\frac52\). Each fund alone has a sample variance of at
least \(4\); the negative correlations of fund \(2\) with funds \(1\) and \(3\) make the portfolio less variable.

(c) In every month the return of fund \(4\) is the average of those of funds \(1\) and \(2\):
\(x_{i4}=\frac12(x_{i1}+x_{i2})\). So for \(a=(1,1,0,-2)'\), every \(a'x_i\) is \(0\); their sample variance \(a'Sa\) is
\(0\) (Proposition~\ref{r12-prop-data}(b)), and \(Sa=0\) (Corollary~\ref{r8-cor-psd-null}); directly,
\(S(1,1,0,-2)'=(16-2+0-14,\ -2+4+0-2,\ 3-3+0+0,\ 7+1+0-8)'=0\). Thus \(S\) is singular and not PD, so
\(\rank(C_5X)<4\) (Proposition~\ref{r12-prop-data}(a)). The first three columns \(c_1,c_2,c_3\) of \(C_5X\) are linearly
independent: if \(b_1c_1+b_2c_2+b_3c_3=0\), rows \(3\), \(5\) and \(2\) give \(2b_2-3b_3=0\), \(4b_1+2b_3=0\) and
\(-4b_1+2b_2-b_3=0\), so \(b_2=\frac32b_3\), \(b_1=-\frac12b_3\) and \(2b_3+3b_3-b_3=4b_3=0\). So \(\rank(C_5X)=3\).
(Fund \(4\) is a \(50{:}50\) mix of funds \(1\) and \(2\), and the combination ``one unit each of funds \(1\) and \(2\),
minus two units of fund \(4\)'' has the constant return \(0\).)
\end{solution}

\begin{exercise}\label{r8-ex-outliers}
Seven students take four quizzes (scores out of \(10\)):
\[
\begin{array}{c|ccccccc}\text{student}&1&2&3&4&5&6&7\\\hline
\text{quiz 1}&9&7&3&4&9&7&3\\\text{quiz 2}&8&5&2&5&8&5&2\\\text{quiz 3}&3&8&3&8&6&8&6\\\text{quiz 4}&5&7&5&7&2&7&2\end{array}
\]
\begin{enumerate}
\item Compute \(\bar x\), \(S\), \(S^{-1}\) and the sample correlation matrix \(R\).
\item Compute the squared Mahalanobis distances \(d_i^2=d_S(x_i,\bar x)^2\) of all students, and check
\(\sum_id_i^2=(n-1)p\).
\item A screening rule flags a student with \(d_i^2>4\) as unusual. Which student is flagged? Compare with the
standardized scores \((x_{ij}-\bar x_j)/\sqrt{s_{jj}}\): is the flagged student more than one standard deviation from the
average on any quiz?
\end{enumerate}
\end{exercise}
\begin{solution}
(a) \(\bar x=\frac17(42,35,42,35)'=(6,5,6,5)'\). The rows of \(C_7X\) are \((3,3,-3,0)\), \((1,0,2,2)\), \((-3,-3,-3,0)\),
\((-2,0,2,2)\), \((3,3,0,-3)\), \((1,0,2,2)\), \((-3,-3,0,-3)\), so (Proposition~\ref{r12-prop-data}(a))
\[
X'C_7X=\begin{pmatrix}42&36&0&0\\36&36&0&0\\0&0&30&12\\0&0&12&30\end{pmatrix},\qquad
S=\frac16X'C_7X=\begin{pmatrix}7&6&0&0\\6&6&0&0\\0&0&5&2\\0&0&2&5\end{pmatrix}
\]
(for instance \(42=9+1+9+4+9+1+9\) and \(30=9+4+9+4+0+4+0\)). So \(S=\diag(S_1,S_2)\) with
\(S_1=\left(\begin{smallmatrix}7&6\\6&6\end{smallmatrix}\right)\) and \(S_2=\left(\begin{smallmatrix}5&2\\2&5\end{smallmatrix}\right)\).
As \(\det S_1=42-36=6\) and \(\det S_2=25-4=21\), the \(2\times2\) formula of Proposition~\ref{r2-prop-inverse-2x2} gives
\(S_1^{-1}=\frac16\left(\begin{smallmatrix}6&-6\\-6&7\end{smallmatrix}\right)\) and
\(S_2^{-1}=\frac1{21}\left(\begin{smallmatrix}5&-2\\-2&5\end{smallmatrix}\right)\). By Proposition~\ref{r5b-prop-block-tri}(b),
\(S^{-1}=\diag(S_1^{-1},S_2^{-1})\), that is,
\[
S^{-1}=\begin{pmatrix}1&-1&0&0\\[2pt]-1&\frac76&0&0\\[2pt]0&0&\frac5{21}&-\frac2{21}\\[2pt]0&0&-\frac2{21}&\frac5{21}\end{pmatrix}
\]
The sample standard deviations are \(\sqrt7\), \(\sqrt6\), \(\sqrt5\), \(\sqrt5\), so \(r_{12}=\frac6{\sqrt{42}}=\sqrt{\frac67}\),
\(r_{34}=\frac25\), the other \(r_{jl}\) (\(j\ne l\)) are \(0\), and (Definition~\ref{r12-def-correlation})
\[
R=\begin{pmatrix}1&\sqrt{\frac67}&0&0\\[2pt]\sqrt{\frac67}&1&0&0\\[2pt]0&0&1&\frac25\\[2pt]0&0&\frac25&1\end{pmatrix}
\]
(b) For a deviation \((u,v,w,s)=(x_i-\bar x)'\),
\(d_i^2=u^2-2uv+\frac76v^2+\frac1{21}(5w^2-4ws+5s^2)=(u-v)^2+\frac16v^2+\frac1{21}(5w^2-4ws+5s^2)\). This gives
\[
d_1^2=d_3^2=d_5^2=d_7^2=\tfrac32+\tfrac{15}7=\tfrac{51}{14},\qquad d_2^2=d_6^2=1+\tfrac87=\tfrac{15}7,\qquad
d_4^2=4+\tfrac87=\tfrac{36}7
\]
with sum \(\frac{204}{14}+\frac{30}7+\frac{36}7=\frac{168}7=24=6\cdot4\), as Proposition~\ref{r8-prop-whitening}(b) says.

(c) Only student \(4\) is flagged: \(d_4^2=\frac{36}7>4\), while all other \(d_i^2\leqslant\frac{51}{14}<4\). The
standardized scores of student \(4\) are \(-2/\sqrt7\), \(0\), \(2/\sqrt5\) and \(2/\sqrt5\), and \(2<\sqrt5<\sqrt7\): less than
one standard deviation from the average on every quiz. Students \(1,3,5,7\) have the standardized scores \(\pm3/\sqrt7\) and
\(\pm3/\sqrt6\) on quizzes \(1\) and \(2\) and \(\pm3/\sqrt5\) on one of quizzes \(3\) and \(4\), more than one standard
deviation from the average on three of the four quizzes (\(3>\sqrt7>\sqrt6>\sqrt5\)), but \(d_i^2=\frac{51}{14}\). Quizzes
\(1\) and \(2\) are strongly correlated (\(r_{12}^2=\frac67\)); students \(1,3,5,7\) follow the common pattern of these two
quizzes, while student \(4\) is below average on quiz \(1\) without being below average on quiz \(2\). The Mahalanobis
distance detects this; the coordinatewise scores do not.
\end{solution}

\begin{exercise}\label{r9-dr-shift}
Let
\[
A=\begin{pmatrix}4&1&1&1\\1&4&1&-2\\1&1&4&1\\1&-2&1&4\end{pmatrix}
\]
\begin{enumerate}
\item Check that \((1,1,1,0)'\), \((0,1,0,-1)'\), \((1,0,-1,0)'\) and \((1,-2,1,-2)'\) are eigenvectors of \(A\), and find the
eigenvalues \(\lambda_1\geqslant\lambda_2\geqslant\lambda_3\geqslant\lambda_4\) of \(A\).
\item Find the maximum and the minimum of \(x'Ax/x'x\) over \(x\ne0\), and all \(x\) at which they are attained.
\item Find the eigenvalues, in decreasing order, of \(-A\), \(A-6I_4\) and \(A+2I_4\). For which \(c\in\R\) is \(A-cI_4\)
PSD, and for which is \(cI_4-A\) PSD?
\end{enumerate}
\end{exercise}
\begin{solution}
(a) \(A(1,1,1,0)'=(6,6,6,0)'\), \(A(0,1,0,-1)'=(0,6,0,-6)'\), \(A(1,0,-1,0)'=(3,0,-3,0)'\) and
\(A(1,-2,1,-2)'=(1,-2,1,-2)'\). The first two vectors are independent eigenvectors belonging to \(6\), so
\(a_A(6)=\dim\mathcal N(A-6I)\geqslant2\), and \(a_A(3)\geqslant1\), \(a_A(1)\geqslant1\) (Theorem~\ref{r8-thm-rank-eigen});
as the multiplicities add up to \(4\) (Definition~\ref{r7-def-eigenvalues}), \(\lambda_1=\lambda_2=6\), \(\lambda_3=3\) and
\(\lambda_4=1\). Check: \(\tr A=16=6+6+3+1\).

\needspace{3\baselineskip}
(b)
\begin{itemize}
\item By Theorem~\ref{r9-thm-rayleigh}(a), the maximum is \(6\) and the minimum is \(1\).
\item By Theorem~\ref{r9-thm-rayleigh}(b), the maximum is attained exactly at the \(x\ne0\) with \(Ax=6x\). As
\(\dim\mathcal N(A-6I)=a_A(6)=2\) (Theorem~\ref{r8-thm-rank-eigen}), the two independent vectors \((1,1,1,0)'\) and
\((0,1,0,-1)'\) form a basis of \(\mathcal N(A-6I)\) (Theorem~\ref{thm:r4-counting}(d)); so these are the vectors
\(x=c_1(1,1,1,0)'+c_2(0,1,0,-1)'\ne0\): a whole plane of maximizers, as \(\lambda_1=\lambda_2\).
\item The minimum is attained exactly at the \(x\ne0\) with \(Ax=x\): \(a_A(1)=1\), so these are the vectors
\(x=c(1,-2,1,-2)'\), \(c\ne0\).
\end{itemize}
\needspace{3\baselineskip}
(c)
\begin{itemize}
\item By Lemma~\ref{r9-lem-eig-order}(b), the eigenvalues of \(-A\) are \(-1\geqslant-3\geqslant-6\geqslant-6\).
\item By Lemma~\ref{r9-lem-eig-order}(c), those of \(A-6I_4\) are \(0\geqslant0\geqslant-3\geqslant-5\), and those of
\(A+2I_4\) are \(8\geqslant8\geqslant5\geqslant3\).
\item \(A-cI_4\) is symmetric with the eigenvalues \(6-c\), \(6-c\), \(3-c\), \(1-c\), so it is PSD exactly when
\(1-c\geqslant0\), that is, \(c\leqslant1\) (Theorem~\ref{r8-thm-eigen-criterion}(a)). Likewise \(cI_4-A=-(A-cI_4)\) has the
eigenvalues \(c-1\), \(c-3\), \(c-6\), \(c-6\) and is PSD exactly when \(c\geqslant6\). For \(c=\lambda_4=1\) and
\(c=\lambda_1=6\) this is Theorem~\ref{r9-thm-rayleigh}(c).
\end{itemize}
\end{solution}

\begin{exercise}\label{r9-ex-rayleigh-drill}
Let
\[
A=\begin{pmatrix}3&1&0&-1\\1&4&2&2\\0&2&4&2\\-1&2&2&4\end{pmatrix}
\]
\begin{enumerate}
\item Check that \((0,1,1,1)'\), \((2,1,0,-1)'\), \((0,1,-2,1)'\) and \((1,-1,0,1)'\) are eigenvectors of \(A\), and find the
eigenvalues \(\lambda_1\geqslant\lambda_2\geqslant\lambda_3\geqslant\lambda_4\) of \(A\).
\item Find the maximum and the minimum of \(x'Ax\) over the unit vectors \(x\in\R^4\), and all unit vectors at which
they are attained.
\item Find the maximum of \(x'Ax/x'x\) over the \(x\ne0\) with \(x_2+x_3+x_4=0\), and the minimum over the \(x\ne0\) with
\(x_1-x_2+x_4=0\).
\end{enumerate}
\end{exercise}
\begin{solution}
(a) \(A(0,1,1,1)'=(0,8,8,8)'\), \(A(2,1,0,-1)'=(8,4,0,-4)'\), \(A(0,1,-2,1)'=(0,2,-4,2)'\) and
\(A(1,-1,0,1)'=(1,-1,0,1)'\): eigenvectors belonging to \(8\), \(4\), \(2\) and \(1\). So the algebraic multiplicities
\(\dim\mathcal N(A-cI)\) of \(c=8,4,2,1\) are at least \(1\) (Theorem~\ref{r8-thm-rank-eigen}); as all multiplicities add
up to \(4\) (Definition~\ref{r7-def-eigenvalues}), each of \(8\), \(4\), \(2\), \(1\) is simple and there are no other
eigenvalues: \(\lambda_1=8\), \(\lambda_2=4\), \(\lambda_3=2\), \(\lambda_4=1\).

(b) Each eigenspace is a line, so the unit eigenvectors are \(\pm q_1\), \(\pm q_2\), \(\pm q_3\), \(\pm q_4\), with
\(q_1=\frac1{\sqrt3}(0,1,1,1)'\), \(q_2=\frac1{\sqrt6}(2,1,0,-1)'\), \(q_3=\frac1{\sqrt6}(0,1,-2,1)'\) and
\(q_4=\frac1{\sqrt3}(1,-1,0,1)'\). By Theorem~\ref{r9-thm-rayleigh}, the maximum is \(8\), attained exactly at
\(x=\pm q_1\), and the minimum is \(1\), attained exactly at \(x=\pm q_4\).

\needspace{3\baselineskip}
(c)
\begin{itemize}
\item \(q_1,q_2,q_3,q_4\) are orthonormal (Theorem~\ref{r8-thm-sym-orth}), so \(Q=[q_1\ q_2\ q_3\ q_4]\) is as in
Lemma~\ref{r9-lem-eig-order}.
\item \(x_2+x_3+x_4=0\) means \(x\perp q_1\). By Theorem~\ref{r9-thm-constrained} with \(k=2\) (first form), the maximum
is \(\lambda_2=4\), attained at \(q_2\).
\item \(x_1-x_2+x_4=0\) means \(x\perp q_4\). By the same theorem with \(k=3\) (second form), the minimum is
\(\lambda_3=2\), attained at \(q_3\).
\end{itemize}
\end{solution}

\needspace{22\baselineskip}
\begin{exercise}\label{r9-dr-variance}
Five units are measured on four variables; row \(i\) of
\[
X=\begin{pmatrix}2&5&5&1\\6&5&1&5\\6&7&7&5\\4&1&7&7\\2&7&5&7\end{pmatrix}
\]
holds the values on unit \(i\). For \(a\in\R^4\), let \(y=Xa\).
\begin{enumerate}
\item Find the mean vector \(\bar x\), the centered data matrix \(C_5X\) and \(M=X'C_5X\).
\item Check that \((0,-1,1,1)'\), \((1,0,-1,1)'\), \((0,2,1,1)'\) and \((2,0,1,-1)'\) are eigenvectors of \(M\), and find
its eigenvalues.
\item Among the unit vectors \(a\), find those for which \(\sum_i(y_i-\bar y)^2\) is largest and smallest, these two
sums, and the corresponding sample variances.
\item Among the unit vectors \(a\) orthogonal to the maximizers of (c), find the largest \(\sum_i(y_i-\bar y)^2\) and
where it is attained.
\end{enumerate}
\end{exercise}
\begin{solution}
(a) \(\bar x=\frac15X'\one=\frac15(20,25,25,25)'=(4,5,5,5)'\), and (Proposition~\ref{r12-prop-data}(a))
\[
C_5X=X-\one\bar x'=\begin{pmatrix}-2&0&0&-4\\2&0&-4&0\\2&2&2&0\\0&-4&2&2\\-2&2&0&2\end{pmatrix},\qquad
M=(C_5X)'(C_5X)=\begin{pmatrix}16&0&-4&4\\0&24&-4&-4\\-4&-4&24&4\\4&-4&4&24\end{pmatrix}
\]
(b) \(M(0,-1,1,1)'=(0,-32,32,32)'\), \(M(1,0,-1,1)'=(24,0,-24,24)'\), \(M(0,2,1,1)'=(0,40,20,20)'\) and
\(M(2,0,1,-1)'=(24,0,12,-12)'\): eigenvectors belonging to \(32\), \(24\), \(20\) and \(12\). So \(a_M(c)\geqslant1\) for
\(c=32,24,20,12\) (Theorem~\ref{r8-thm-rank-eigen}), and as the multiplicities add up to \(4\), the eigenvalues are
\(\lambda_1=32\), \(\lambda_2=24\), \(\lambda_3=20\), \(\lambda_4=12\). Check: \(\tr M=88=32+24+20+12\).

\needspace{3\baselineskip}
(c)
\begin{itemize}
\item By Proposition~\ref{r12-prop-data}(b), \(y=Xa\) has the sample variance \(a'Sa\) with \(S=\frac14M\), so
\(\sum_i(y_i-\bar y)^2=4a'Sa=a'Ma\).
\item Each eigenspace is a line, so the unit eigenvectors are \(\pm q_1,\dots,\pm q_4\), with
\(q_1=\frac1{\sqrt3}(0,-1,1,1)'\), \(q_2=\frac1{\sqrt3}(1,0,-1,1)'\), \(q_3=\frac1{\sqrt6}(0,2,1,1)'\) and
\(q_4=\frac1{\sqrt6}(2,0,1,-1)'\). By Theorem~\ref{r9-thm-rayleigh}, the largest sum is \(32\), attained exactly at
\(a=\pm q_1\), and the smallest is \(12\), attained exactly at \(a=\pm q_4\).
\item Check: \(C_5Xq_1=\frac1{\sqrt3}(-4,-4,0,8,0)'\), whose squared length is \(\frac{96}3=32\).
\item Divided by \(n-1=4\): the sample variances \(a'Sa\) lie in \([3,8]\) for unit \(a\), with both ends attained.
\end{itemize}
(d) \(q_1,\dots,q_4\) are orthonormal (Theorem~\ref{r8-thm-sym-orth}), so \(Q=[q_1\ \cdots\ q_4]\) is as in
Lemma~\ref{r9-lem-eig-order}. By Theorem~\ref{r9-thm-constrained} (\(k=2\), first form), the largest \(a'Ma\) over the
unit vectors \(a\perp q_1\) is \(\lambda_2=24\), attained at \(a=\pm q_2\).
\end{solution}

\begin{exercise}\label{r12-ex-kyfan}
Let
\[
A=\begin{pmatrix}3&1&0&1\\1&3&0&1\\0&0&4&1\\1&1&1&2\end{pmatrix}
\]
\begin{enumerate}
\item Find the eigenvalues of \(A\) and orthonormal eigenvectors.
\item Find the largest and the smallest value of \(\tr(G'AG)\) over all \(G\in\R^{4\times2}\) with \(G'G=I_2\), and a
\(G\) attaining each.
\item For \(G=[e_1\ e_2]\), compute \(\tr(G'AG)\) and check that it lies between the two values of (b).
\end{enumerate}
\end{exercise}
\begin{solution}
(a) For real \(t\), apply \(R_1\leftarrow R_1-R_2\) to \(tI_4-A\), take the factor \(t-2\) out of row \(1\), apply
\(\col_2\leftarrow\col_2+\col_1\), and expand along row \(1\) and then along row \(1\) again
(Definition~\ref{r5d-def-det}(a),(b), Theorems~\ref{r5d-thm-transpose} and~\ref{r5d-thm-expansion}):
\[
\begin{aligned}
\det(tI_4-A)&=\det\begin{pmatrix}t-3&-1&0&-1\\-1&t-3&0&-1\\0&0&t-4&-1\\-1&-1&-1&t-2\end{pmatrix}
=(t-2)\det\begin{pmatrix}1&-1&0&0\\-1&t-3&0&-1\\0&0&t-4&-1\\-1&-1&-1&t-2\end{pmatrix}\\
&=(t-2)\det\begin{pmatrix}1&0&0&0\\-1&t-4&0&-1\\0&0&t-4&-1\\-1&-2&-1&t-2\end{pmatrix}
=(t-2)\det\begin{pmatrix}t-4&0&-1\\0&t-4&-1\\-2&-1&t-2\end{pmatrix}\\
&=(t-2)\bigl((t-4)\bigl((t-4)(t-2)-1\bigr)-2(t-4)\bigr)\\
&=(t-2)(t-4)(t^2-6t+5)=(t-5)(t-4)(t-2)(t-1)
\end{aligned}
\]
so the eigenvalues are \(5\), \(4\), \(2\), \(1\) (Definition~\ref{r7-def-charpoly}(a), Proposition~\ref{r7-prop-poly}(d)).
Eigenvectors:
\[
\begin{gathered}
A-5I=\begin{pmatrix}-2&1&0&1\\1&-2&0&1\\0&0&-1&1\\1&1&1&-3\end{pmatrix}
\xrightarrow{R_1\leftrightarrow R_4}
\begin{pmatrix}1&1&1&-3\\1&-2&0&1\\0&0&-1&1\\-2&1&0&1\end{pmatrix}
\xrightarrow[R_4\leftarrow R_4+2R_1]{R_2\leftarrow R_2-R_1}
\begin{pmatrix}1&1&1&-3\\0&-3&-1&4\\0&0&-1&1\\0&3&2&-5\end{pmatrix}
\\
\xrightarrow{R_4\leftarrow R_4+R_2}
\begin{pmatrix}1&1&1&-3\\0&-3&-1&4\\0&0&-1&1\\0&0&1&-1\end{pmatrix}
\xrightarrow{R_4\leftarrow R_4+R_3}
\begin{pmatrix}1&1&1&-3\\0&-3&-1&4\\0&0&-1&1\\0&0&0&0\end{pmatrix}
\end{gathered}
\]
so \(x_4=1\) gives \(x_3=x_4=1\), \(x_2=\frac13(4x_4-x_3)=1\) and \(x_1=-x_2-x_3+3x_4=1\); take \((1,1,1,1)'\).
\[
\begin{gathered}
A-4I=\begin{pmatrix}-1&1&0&1\\1&-1&0&1\\0&0&0&1\\1&1&1&-2\end{pmatrix}
\xrightarrow[R_4\leftarrow R_4+R_1]{R_2\leftarrow R_2+R_1}
\begin{pmatrix}-1&1&0&1\\0&0&0&2\\0&0&0&1\\0&2&1&-1\end{pmatrix}
\xrightarrow{R_2\leftrightarrow R_4}
\begin{pmatrix}-1&1&0&1\\0&2&1&-1\\0&0&0&1\\0&0&0&2\end{pmatrix}
\\
\xrightarrow{R_4\leftarrow R_4-2R_3}
\begin{pmatrix}-1&1&0&1\\0&2&1&-1\\0&0&0&1\\0&0&0&0\end{pmatrix}
\end{gathered}
\]
so \(x_4=0\), and \(x_3=-2\) gives \(x_2=\frac12(x_4-x_3)=1\) and \(x_1=x_2+x_4=1\); take \((1,1,-2,0)'\).
\[
\begin{gathered}
A-2I=\begin{pmatrix}1&1&0&1\\1&1&0&1\\0&0&2&1\\1&1&1&0\end{pmatrix}
\xrightarrow[R_4\leftarrow R_4-R_1]{R_2\leftarrow R_2-R_1}
\begin{pmatrix}1&1&0&1\\0&0&0&0\\0&0&2&1\\0&0&1&-1\end{pmatrix}
\xrightarrow{R_2\leftrightarrow R_4}
\begin{pmatrix}1&1&0&1\\0&0&1&-1\\0&0&2&1\\0&0&0&0\end{pmatrix}
\\
\xrightarrow{R_3\leftarrow R_3-2R_2}
\begin{pmatrix}1&1&0&1\\0&0&1&-1\\0&0&0&3\\0&0&0&0\end{pmatrix}
\end{gathered}
\]
so \(x_4=0\), \(x_3=x_4=0\), and \(x_2=-1\) gives \(x_1=-x_2-x_4=1\); take \((1,-1,0,0)'\).
\[
\begin{gathered}
A-I=\begin{pmatrix}2&1&0&1\\1&2&0&1\\0&0&3&1\\1&1&1&1\end{pmatrix}
\xrightarrow{R_1\leftrightarrow R_4}
\begin{pmatrix}1&1&1&1\\1&2&0&1\\0&0&3&1\\2&1&0&1\end{pmatrix}
\xrightarrow[R_4\leftarrow R_4-2R_1]{R_2\leftarrow R_2-R_1}
\begin{pmatrix}1&1&1&1\\0&1&-1&0\\0&0&3&1\\0&-1&-2&-1\end{pmatrix}
\\
\xrightarrow{R_4\leftarrow R_4+R_2}
\begin{pmatrix}1&1&1&1\\0&1&-1&0\\0&0&3&1\\0&0&-3&-1\end{pmatrix}
\xrightarrow{R_4\leftarrow R_4+R_3}
\begin{pmatrix}1&1&1&1\\0&1&-1&0\\0&0&3&1\\0&0&0&0\end{pmatrix}
\end{gathered}
\]
so \(x_4=-3\) gives \(x_3=-\frac13x_4=1\), \(x_2=x_3=1\) and \(x_1=-x_2-x_3-x_4=1\); take \((1,1,1,-3)'\). The four vectors
are pairwise orthogonal, and \(q_1=\frac12(1,1,1,1)'\), \(q_2=\frac1{\sqrt6}(1,1,-2,0)'\), \(q_3=\frac1{\sqrt2}(1,-1,0,0)'\),
\(q_4=\frac1{2\sqrt3}(1,1,1,-3)'\) are orthonormal eigenvectors for \(5\), \(4\), \(2\), \(1\).

(b) By Theorem~\ref{r12-thm-kyfan} with \(k=2\), the largest value is \(5+4=9\), attained at \(G=[q_1\ q_2]\), and the
smallest is \(2+1=3\), attained at \(G=[q_3\ q_4]\).

(c) \(G'G=I_2\), and \(G'AG\) consists of rows and columns \(1\), \(2\) of \(A\):
\(G'AG=\left(\begin{smallmatrix}3&1\\1&3\end{smallmatrix}\right)\), with trace \(6\). Indeed \(3\leqslant6\leqslant9\), as
Corollary~\ref{r12-cor-diagonal} says for \(a_{11}+a_{22}\).
\end{solution}

\begin{exercise}[cf.\ \citeat{tu-asm}{Exercise~5.9}, \citeat{islr}{\S12.6, Exercise~8}]\label{r12-ex-pca-drill}
Five units are measured on four variables:
\[
X=\begin{pmatrix}3&1&1&9\\3&4&8&4\\3&6&4&2\\5&6&6&4\\1&8&6&6\end{pmatrix}
\]
\begin{enumerate}
\item Find \(\bar x\), \(C_5X\) and \(S\).
\item Find the loading vectors, the sample variances of the principal components and the score vectors.
\item Find \(\mathrm{PVE}_1,\dots,\mathrm{PVE}_4\), once from the eigenvalues and once from the score vectors.
\end{enumerate}
\end{exercise}
\begin{solution}
(a) \(\one'X=(15,25,25,25)\), so \(\bar x=(3,5,5,5)'\), and
\[
\begin{gathered}
C_5X=\begin{pmatrix}0&-4&-4&4\\0&-1&3&-1\\0&1&-1&-3\\2&1&1&-1\\-2&3&1&1\end{pmatrix}
\\
(C_5X)'(C_5X)=\begin{pmatrix}8&-4&0&-4\\-4&28&16&-16\\0&16&28&-16\\-4&-16&-16&28\end{pmatrix},\qquad
S=\frac14(C_5X)'(C_5X)=\begin{pmatrix}2&-1&0&-1\\-1&7&4&-4\\0&4&7&-4\\-1&-4&-4&7\end{pmatrix}
\end{gathered}
\]
(b) For real \(t\), add row \(3\) of \(tI_4-S\) to row \(2\) and subtract row \(4\) from it; the new row \(2\) is
\((t-15)(0,1,1,-1)\). Take out \(t-15\), subtract column \(2\) from column \(3\) and add it to column \(4\), and expand
along row \(2\) twice (Definition~\ref{r5d-def-det}(a),(b), Theorems~\ref{r5d-thm-transpose} and~\ref{r5d-thm-expansion}):
\[
\begin{aligned}
\det(tI_4-S)&=\det\begin{pmatrix}t-2&1&0&1\\1&t-7&-4&4\\0&-4&t-7&4\\1&4&4&t-7\end{pmatrix}
=(t-15)\det\begin{pmatrix}t-2&1&0&1\\0&1&1&-1\\0&-4&t-7&4\\1&4&4&t-7\end{pmatrix}\\
&=(t-15)\det\begin{pmatrix}t-2&1&-1&2\\0&1&0&0\\0&-4&t-3&0\\1&4&0&t-3\end{pmatrix}
=(t-15)\det\begin{pmatrix}t-2&-1&2\\0&t-3&0\\1&0&t-3\end{pmatrix}\\
&=(t-15)(t-3)\bigl((t-2)(t-3)-2\bigr)\\
&=(t-15)(t-3)(t^2-5t+4)=(t-15)(t-4)(t-3)(t-1)
\end{aligned}
\]
So the eigenvalues are \(\lambda_1=15\), \(\lambda_2=4\), \(\lambda_3=3\), \(\lambda_4=1\) (Definition~\ref{r7-def-charpoly}(a),
Proposition~\ref{r7-prop-poly}(d)). Each is simple, so each eigenspace is a line (Theorem~\ref{r8-thm-rank-eigen}), spanned
by the vector found in
\[
S\begin{pmatrix}0\\1\\1\\-1\end{pmatrix}=\begin{pmatrix}0\\15\\15\\-15\end{pmatrix},\qquad
S\begin{pmatrix}1\\-1\\0\\-1\end{pmatrix}=\begin{pmatrix}4\\-4\\0\\-4\end{pmatrix},\qquad
S\begin{pmatrix}0\\1\\-2\\-1\end{pmatrix}=\begin{pmatrix}0\\3\\-6\\-3\end{pmatrix},\qquad
S\begin{pmatrix}2\\1\\0\\1\end{pmatrix}=\begin{pmatrix}2\\1\\0\\1\end{pmatrix}
\]
(the first is the row combination used for \(p_S\)). So the loading vectors are, up to sign (Theorem~\ref{r12-thm-pca}(e)),
\(v_1=\frac1{\sqrt3}(0,1,1,-1)'\), \(v_2=\frac1{\sqrt3}(1,-1,0,-1)'\), \(v_3=\frac1{\sqrt6}(0,1,-2,-1)'\) and
\(v_4=\frac1{\sqrt6}(2,1,0,1)'\) (Theorem~\ref{r8-thm-spectral}(b)), and the principal components have the sample variances
\(15\), \(4\), \(3\), \(1\). The first principal component contrasts variables \(2\) and \(3\) with variable \(4\). The score
vectors are
\[
\begin{gathered}
z_1=C_5Xv_1=\frac1{\sqrt3}\begin{pmatrix}-12\\3\\3\\3\\3\end{pmatrix},\qquad
z_2=C_5Xv_2=\frac1{\sqrt3}\begin{pmatrix}0\\2\\2\\2\\-6\end{pmatrix}
\\
z_3=C_5Xv_3=\frac1{\sqrt6}\begin{pmatrix}0\\-6\\6\\0\\0\end{pmatrix},\qquad
z_4=C_5Xv_4=\frac1{\sqrt6}\begin{pmatrix}0\\-2\\-2\\4\\0\end{pmatrix}
\end{gathered}
\]
with \(z_m'z_l=0\) for \(m\ne l\), \(\lVert z_1\rVert^2=60=4\cdot15\), \(\lVert z_2\rVert^2=16=4\cdot4\),
\(\lVert z_3\rVert^2=12=4\cdot3\) and \(\lVert z_4\rVert^2=4=4\cdot1\) (Theorem~\ref{r12-thm-pca}(d)).

(c) \(\tr S=2+7+7+7=23=15+4+3+1\), so \(\mathrm{PVE}_m=\frac{15}{23},\frac4{23},\frac3{23},\frac1{23}\). From the scores
(Proposition~\ref{r12-prop-pve}): \(\lVert C_5X\rVert_F^2=8+28+28+28=92\), and \(\frac{60}{92}=\frac{15}{23}\),
\(\frac{16}{92}=\frac4{23}\), \(\frac{12}{92}=\frac3{23}\), \(\frac4{92}=\frac1{23}\).
\end{solution}

\begin{exercise}\label{r10-ex-pca-drill}
Four variables are measured on five units: \((3,6,3,3)\), \((2,3,5,1)\), \((6,7,1,1)\), \((6,1,1,5)\), \((8,3,5,5)\).
\begin{enumerate}
\item Find \(\bar x\), the centered data matrix \(C_5X\) and \(S\).
\item Show that \((1,-1,0,1)'\), \((1,1,-1,0)'\), \((1,1,2,0)'\) and \((1,-1,0,-2)'\) are eigenvectors of \(S\), and give the
loading vectors, the variances of the components and the proportions of variance explained by the first component and
by the first two.
\item Compute the scores, and check Theorem~\ref{r12-thm-pca}(d) for them.
\item Find the sum over the units of the squared errors of the approximation of \(C_5X\) by two components.
\end{enumerate}
\end{exercise}
\begin{solution}
(a) \(\bar x=\frac15(25,20,15,15)'=(5,4,3,3)'\), and
\[
\begin{gathered}
C_5X=\begin{pmatrix}-2&2&0&0\\-3&-1&2&-2\\1&3&-2&-2\\1&-3&-2&2\\3&-1&2&2\end{pmatrix}
\\
(C_5X)'(C_5X)=\begin{pmatrix}24&-4&-4&12\\-4&24&-4&-12\\-4&-4&16&0\\12&-12&0&16\end{pmatrix},\qquad
S=\frac14(C_5X)'(C_5X)=\begin{pmatrix}6&-1&-1&3\\-1&6&-1&-3\\-1&-1&4&0\\3&-3&0&4\end{pmatrix}
\end{gathered}
\]
(b) \(S(1,-1,0,1)'=(10,-10,0,10)'\), \(S(1,1,-1,0)'=(6,6,-6,0)'\), \(S(1,1,2,0)'=(3,3,6,0)'\) and
\(S(1,-1,0,-2)'=(1,-1,0,-2)'\). The four vectors are pairwise orthogonal, so \(\lambda_1=10\), \(\lambda_2=6\),
\(\lambda_3=3\), \(\lambda_4=1\) (Theorem~\ref{r8-thm-spectral}(c)), the sample variances of the components (Theorem~\ref{r12-thm-pca}(d)), with the loading vectors \(v_1=\frac1{\sqrt3}(1,-1,0,1)'\),
\(v_2=\frac1{\sqrt3}(1,1,-1,0)'\), \(v_3=\frac1{\sqrt6}(1,1,2,0)'\), \(v_4=\frac1{\sqrt6}(1,-1,0,-2)'\), unique up to sign
(Theorem~\ref{r12-thm-pca}(e)). The total variance is \(\tr S=20\); the first component explains \(\frac{10}{20}=\frac12\),
the first two \(\frac{16}{20}=\frac45\) (Proposition~\ref{r12-prop-pve}).

(c) \(z_{im}=v_m'(x_i-\bar x)\):
\[
\begin{gathered}
(z_{i1})_i=\tfrac1{\sqrt3}(-4,-4,-4,6,6),\qquad(z_{i2})_i=\tfrac1{\sqrt3}(0,-6,6,0,0)
\\
(z_{i3})_i=\tfrac1{\sqrt6}(0,0,0,-6,6),\qquad(z_{i4})_i=\tfrac1{\sqrt6}(-4,2,2,0,0)
\end{gathered}
\]
Each has sum \(0\); the sums of squares are \(\frac{120}3=40=4\lambda_1\), \(\frac{72}3=24=4\lambda_2\),
\(\frac{72}6=12=4\lambda_3\) and \(\frac{24}6=4=4\lambda_4\); and the sums of products of two different columns are
\(\frac13(24-24)=0\), \(\frac1{\sqrt{18}}(-36+36)=0\), \(\frac1{\sqrt{18}}(16-8-8)=0\), \(0\), \(\frac1{\sqrt{18}}(-12+12)=0\) and
\(0\) for the pairs \((1,2)\), \((1,3)\), \((1,4)\), \((2,3)\), \((2,4)\), \((3,4)\).

(d) By Theorem~\ref{r12-thm-closest}(b) with \(M=2\), the sum is
\(\lVert C_5X-Z_2V_2'\rVert_F^2=\lVert C_5X\rVert_F^2-\lVert z_1\rVert^2-\lVert z_2\rVert^2=80-40-24=16\), out of
\(\lVert C_5X\rVert_F^2=(n-1)\tr S=80\); by Theorem~\ref{r12-thm-closest}(a) it is
\((n-1)(\lambda_3+\lambda_4)=4\cdot4=16=\sum_i(z_{i3}^2+z_{i4}^2)\), the smallest possible for an approximation of rank at
most \(2\) (the paragraph after Theorem~\ref{r12-thm-closest}).
\end{solution}

\begin{exercise}\label{r12-ex-ties}
\begin{enumerate}
\item The sample covariance matrix of four variables is
\[
S=\begin{pmatrix}3&1&1&1\\1&3&1&1\\1&1&3&1\\1&1&1&3\end{pmatrix}
\]
Find the eigenvalues of \(S\) and the proportions of variance explained. Show that the first loading vector is unique up
to sign, but that every unit vector orthogonal to \(\one_4\) is a second loading vector.
\item For the data matrix \(X_0\) of Example~\ref{r1-exm-data} (\(n=4\), two variables), find \(S\) and show that every unit
vector is a first loading vector.
\end{enumerate}
\end{exercise}
\begin{solution}
(a) \(S=2I_4+J_4=2(\bar J_4+C_4)+4\bar J_4=6\bar J_4+2C_4\) (Definition~\ref{r1-def-ones}). So \(S\one=6\one\), and
\(Sx=2x\) whenever \(\one'x=0\) (then \(\bar J_4x=0\), \(C_4x=x\)). With \(v_1=\frac12\one\) and the orthonormal vectors
\(\frac1{\sqrt2}(1,-1,0,0)'\), \(\frac1{\sqrt6}(1,1,-2,0)'\), \(\frac1{2\sqrt3}(1,1,1,-3)'\) orthogonal to \(\one\), the
eigenvalues are \(6,2,2,2\) (Theorem~\ref{r8-thm-spectral}(c)), and \(\mathrm{PVE}_m=\frac12,\frac16,\frac16,\frac16\)
(Proposition~\ref{r12-prop-pve}). As \(6>2\), the first loading vector is \(\pm v_1\) (Theorem~\ref{r12-thm-pca}(e),
\(M=1\)). The second problem has the constraint \(a\perp v_1\) (Theorem~\ref{r12-thm-pca}(b)), that is, \(\one'a=0\); for
every such unit vector \(a'Sa=2a'a=2\), the maximum \(\lambda_2\) (Theorem~\ref{r12-thm-pca}(c)). So every unit vector
orthogonal to \(\one_4\) is a second loading vector: with a repeated eigenvalue the loading vectors are not unique.

(b) By Example~\ref{r1-exm-data}, \(X_0'C_4X_0=4I_2\), so \(S=\frac13\cdot4I_2=\frac43I_2\). Every unit vector \(a\) gives
\(a'Sa=\frac43\), so every unit vector is a first loading vector; \(\lambda_1=\lambda_2=\frac43\), and each component
explains half of the total variance.
\end{solution}

\begin{exercise}\label{r12-ex-plane}
Five units are measured on four variables:
\[
X=\begin{pmatrix}5&5&2&4\\3&3&6&4\\3&3&3&3\\3&1&2&2\\1&3&2&2\end{pmatrix}
\]
\begin{enumerate}
\item Find \(\bar x\), \(C_5X\) and \(S\).
\item Using Exercise~\ref{r12-ex-kyfan}, find the eigenvalues of \(S\), orthonormal eigenvectors and \(r=\rank(C_5X)\).
\item Find the proportions of variance explained, and show that the five observations lie in an affine subspace of
dimension \(3\) through \(\bar x\).
\end{enumerate}
\end{exercise}
\begin{solution}
(a) \(\one'X=(15,15,15,15)\), so \(\bar x=(3,3,3,3)'\), and
\[
C_5X=\begin{pmatrix}2&2&-1&1\\0&0&3&1\\0&0&0&0\\0&-2&-1&-1\\-2&0&-1&-1\end{pmatrix},\qquad
(C_5X)'(C_5X)=\begin{pmatrix}8&4&0&4\\4&8&0&4\\0&0&12&4\\4&4&4&4\end{pmatrix},\qquad
S=\begin{pmatrix}2&1&0&1\\1&2&0&1\\0&0&3&1\\1&1&1&1\end{pmatrix}
\]
(b) \(S=A-I_4\) with \(A\) of Exercise~\ref{r12-ex-kyfan}, so \(Sq_m=Aq_m-q_m\): the vectors \(q_1,q_2,q_3,q_4\) found there
are orthonormal eigenvectors of \(S\) for \(5-1=4\), \(4-1=3\), \(2-1=1\), \(1-1=0\) (Theorem~\ref{r8-thm-spectral}(c)). By
Theorem~\ref{r12-thm-pca}(a), \(r\) is the number of positive eigenvalues: \(r=3\).

(c) \(\tr S=8\), so \(\mathrm{PVE}_m=\frac48=\frac12,\frac38,\frac18,0\) (Proposition~\ref{r12-prop-pve}): three components
explain all the variance. The fourth score vector is
\(z_4=C_5Xq_4=\frac1{2\sqrt3}C_5X(1,1,1,-3)'=\frac1{2\sqrt3}(2+2-1-3,\ 0+0+3-3,\ 0,\ 0-2-1+3,\ -2+0-1+3)'=0\); so
\((x_i-\bar x)'(1,1,1,-3)'=0\) for every \(i\), that is, all observations lie in the set \(W=\{y:y_1+y_2+y_3-3y_4=0\}\),
which contains \(\bar x=(3,3,3,3)'\): the fourth variable is the average of the other three. \(W\) is the subspace
\(\Span\{(1,1,1,-3)'\}^\perp\), of dimension \(3\) (Theorem~\ref{r6-thm-complement}(c)), and as \(\bar x\in W\), it is the
affine subspace \(\bar x+W\) (Definition~\ref{r12-def-affine}). (By Theorem~\ref{r12-thm-closest}(a) with
\(M=3\geqslant r\), \(C_5X=Z_3V_3'\) exactly.)
\end{solution}

\begin{exercise}\label{r12-dr-correlation}
Five units are measured on four variables:
\[
X=\begin{pmatrix}4&5&9&9\\7&9&1&9\\1&1&1&5\\1&9&5&1\\7&1&9&1\end{pmatrix}
\]
\begin{enumerate}
\item Find \(\bar x\), \(C_5X\) and \(S\).
\item Find \(D_s^{1/2}\), the sample correlation matrix \(R\) and the standardized data matrix \(X_s\), and check that
\(\frac14X_s'X_s=R\).
\item Find the loading vectors and the proportions of variance explained, first for the variables (from \(S\)) and then
for the standardized variables (from \(R\)).
\end{enumerate}
\end{exercise}
\begin{solution}
(a) \(\one'X=(20,25,25,25)\), so \(\bar x=(4,5,5,5)'\), and
\[
\begin{gathered}
C_5X=\begin{pmatrix}0&0&4&4\\3&4&-4&4\\-3&-4&-4&0\\-3&4&0&-4\\3&-4&4&-4\end{pmatrix}
\\
(C_5X)'(C_5X)=\begin{pmatrix}36&0&12&12\\0&64&-16&16\\12&-16&64&-16\\12&16&-16&64\end{pmatrix},\qquad
S=\begin{pmatrix}9&0&3&3\\0&16&-4&4\\3&-4&16&-4\\3&4&-4&16\end{pmatrix}
\end{gathered}
\]
(b) \(D_s^{1/2}=\diag(3,4,4,4)\), \(r_{12}=0\), \(r_{13}=r_{14}=\frac3{3\cdot4}=\frac14\), \(r_{23}=r_{34}=\frac{-4}{4\cdot4}=-\frac14\)
and \(r_{24}=\frac4{4\cdot4}=\frac14\), so
\[
\begin{gathered}
R=\begin{pmatrix}1&0&\frac14&\frac14\\[2pt]0&1&-\frac14&\frac14\\[2pt]\frac14&-\frac14&1&-\frac14\\[2pt]\frac14&\frac14&-\frac14&1\end{pmatrix}
\\
X_s=C_5XD_s^{-1/2}=\begin{pmatrix}0&0&1&1\\1&1&-1&1\\-1&-1&-1&0\\-1&1&0&-1\\1&-1&1&-1\end{pmatrix},\qquad
X_s'X_s=\begin{pmatrix}4&0&1&1\\0&4&-1&1\\1&-1&4&-1\\1&1&-1&4\end{pmatrix}
\end{gathered}
\]
For instance, column \(1\) of \(X_s\) has the sum of squares \(0+1+1+1+1=4\), and columns \(1\) and \(3\) the sum of
products \(0-1+1+0+1=1\). So \(\frac14X_s'X_s=R\) (Proposition~\ref{r12-prop-correlation}(a)).

(c) From \(S\): for real \(t\), subtract row \(3\) of \(tI_4-S\) from row \(2\) and add row \(4\) to it; the new row \(2\)
is \((t-24)(0,1,-1,1)\). Take out \(t-24\), add column \(2\) to column \(3\) and subtract it from column \(4\), and expand
along row \(2\) and then along row \(1\) (Definition~\ref{r5d-def-det}(a),(b), Theorems~\ref{r5d-thm-transpose}
and~\ref{r5d-thm-expansion}):
\[
\begin{aligned}
\det(tI_4-S)&=\det\begin{pmatrix}t-9&0&-3&-3\\0&t-16&4&-4\\-3&4&t-16&4\\-3&-4&4&t-16\end{pmatrix}
=(t-24)\det\begin{pmatrix}t-9&0&-3&-3\\0&1&-1&1\\-3&4&t-16&4\\-3&-4&4&t-16\end{pmatrix}\\
&=(t-24)\det\begin{pmatrix}t-9&0&-3&-3\\0&1&0&0\\-3&4&t-12&0\\-3&-4&0&t-12\end{pmatrix}
=(t-24)\det\begin{pmatrix}t-9&-3&-3\\-3&t-12&0\\-3&0&t-12\end{pmatrix}\\
&=(t-24)\bigl((t-9)(t-12)^2-9(t-12)-9(t-12)\bigr)\\
&=(t-24)(t-12)(t^2-21t+90)=(t-24)(t-15)(t-12)(t-6)
\end{aligned}
\]
so the eigenvalues are \(24\), \(15\), \(12\), \(6\) (Definition~\ref{r7-def-charpoly}(a), Proposition~\ref{r7-prop-poly}(d)), and \(S(0,1,-1,1)'=24(0,1,-1,1)'\), \(S(1,0,1,1)'=15(1,0,1,1)'\), \(S(0,2,1,-1)'=12(0,2,1,-1)'\),
\(S(2,0,-1,-1)'=6(2,0,-1,-1)'\). The eigenvalues are simple, so the loading vectors are, up to sign,
\(\frac1{\sqrt3}(0,1,-1,1)'\), \(\frac1{\sqrt3}(1,0,1,1)'\), \(\frac1{\sqrt6}(0,2,1,-1)'\) and \(\frac1{\sqrt6}(2,0,-1,-1)'\)
(Theorem~\ref{r12-thm-pca}(c),(e)); as \(\tr S=57\), the proportions of variance explained (Proposition~\ref{r12-prop-pve}) are
\(\frac{24}{57}=\frac8{19}\), \(\frac{15}{57}=\frac5{19}\), \(\frac{12}{57}=\frac4{19}\) and \(\frac6{57}=\frac2{19}\).

From \(R\): \(R(0,1,-1,1)'=\frac32(0,1,-1,1)'\), \(R(2,0,1,1)'=\frac54(2,0,1,1)'\), \(R(0,2,1,-1)'=\frac34(0,2,1,-1)'\) and
\(R(1,0,-1,-1)'=\frac12(1,0,-1,-1)'\). These four vectors are pairwise orthogonal, so the eigenvalues of \(R\) are
\(\frac32,\frac54,\frac34,\frac12\) (Theorem~\ref{r8-thm-spectral}(c)). The loading vectors are, up to sign,
\(\frac1{\sqrt3}(0,1,-1,1)'\), \(\frac1{\sqrt6}(2,0,1,1)'\), \(\frac1{\sqrt6}(0,2,1,-1)'\) and \(\frac1{\sqrt3}(1,0,-1,-1)'\); as
\(\tr R=4\), the proportions are \(\frac38\), \(\frac5{16}\), \(\frac3{16}\) and \(\frac18\). The first and the third loading
vectors are the same for \(S\) and for \(R\), but the second and the fourth are not: variable \(1\), whose sample standard
deviation \(3\) is the smallest, has the weight \(\frac1{\sqrt3}\) in the second loading vector of the variables and the
larger weight \(\frac2{\sqrt6}\) in that of the standardized variables.
\end{solution}

\begin{exercise}\label{r12-ex-closest-line}
Continue Exercise~\ref{r12-ex-pca-drill}.
\begin{enumerate}
\item Find a best approximation of \(C_5X\) by a matrix of rank at most \(1\), and the error.
\item Find the line in \(\R^4\) closest to the five observations and the sum of the squared distances. Which observations
lie on it?
\end{enumerate}
\end{exercise}
\begin{solution}
(a) By Theorem~\ref{r12-thm-closest}(a) with \(M=1\), \(\lVert C_5X-AB'\rVert_F^2\geqslant(5-1)(\lambda_2+\lambda_3+\lambda_4)=4\cdot8=32\)
for all \(A\in\R^{5\times1}\) and \(B\in\R^{4\times1}\), with equality for \(A=z_1\), \(B=v_1\). Every \(F\in\R^{5\times4}\) with
\(\rank F\leqslant1\) is such a product: \(F=0\) is \(0\cdot0'\), and if \(\rank F=1\), then \(F\) is a \(5\times1\) times a
\(1\times4\) matrix (Theorem~\ref{thm:r4-full-rank-factorization}(a)). So \(A_1=z_1v_1'\) is a best approximation; as
\(z_1=\sqrt3(-4,1,1,1,1)'\) and \(v_1=\frac1{\sqrt3}(0,1,1,-1)'\),
\[
\begin{gathered}
A_1=\begin{pmatrix}0&-4&-4&4\\0&1&1&-1\\0&1&1&-1\\0&1&1&-1\\0&1&1&-1\end{pmatrix},\qquad
C_5X-A_1=\begin{pmatrix}0&0&0&0\\0&-2&2&0\\0&0&-2&-2\\2&0&0&0\\-2&2&0&2\end{pmatrix}
\\
\lVert C_5X-A_1\rVert_F^2=0+8+8+4+12=32
\end{gathered}
\]
the smallest possible error, \(4(\lambda_2+\lambda_3+\lambda_4)\).

(b) By Theorem~\ref{r12-thm-closest}(c) with \(M=1\) (\(\lambda_1>\lambda_2\)), the closest line is
\(\bar x+\Span\{(0,1,1,-1)'\}=\{(3,\ 5+s,\ 5+s,\ 5-s)':s\in\R\}\), the only one, with the sum of squared distances
\((5-1)(\lambda_2+\lambda_3+\lambda_4)=32\). Row \(i\) of \(C_5X-A_1\) is \(\bigl((x_i-\bar x)-v_1v_1'(x_i-\bar x)\bigr)'\), as
\(z_{i1}=v_1'(x_i-\bar x)\); so by Definition~\ref{r12-def-affine} the squared distances of the five observations from the line
are \(0\), \(8\), \(8\), \(4\), \(12\), with sum \(32\). Observation \(1\), \((3,1,1,9)'=\bar x-4(0,1,1,-1)'\), lies on the
line; the other four do not.
\end{solution}

\begin{exercise}\label{r12-dr-pop-pca}
A random vector \(x=(x_1,x_2,x_3,x_4)'\) has the covariance matrix
\[
\Sigma=\begin{pmatrix}4&0&2&-1\\0&4&-2&-2\\2&-2&4&0\\-1&-2&0&4\end{pmatrix}
\]
\begin{enumerate}
\item Find the eigenvalues of \(\Sigma\) and orthonormal eigenvectors.
\item Find the principal components of \(x\), their variances and the proportions of the total variance. How many
components account for at least \(80\%\) of the total variance?
\item Check that the first two principal components are uncorrelated, and compare
\(\var(x_1)+\var(x_2)\) with the sum of the variances of the first two principal components.
\end{enumerate}
\end{exercise}
\begin{solution}
(a) \(\Sigma=\left(\begin{smallmatrix}4I_2&B\\B'&4I_2\end{smallmatrix}\right)\) with
\(B=\left(\begin{smallmatrix}2&-1\\-2&-2\end{smallmatrix}\right)\), and \(BB'=\left(\begin{smallmatrix}5&-2\\-2&8\end{smallmatrix}\right)\).
For real \(t\ne4\), the block \((t-4)I_2\) of \(tI_4-\Sigma=\left(\begin{smallmatrix}(t-4)I_2&-B\\-B'&(t-4)I_2\end{smallmatrix}\right)\)
is invertible, and Theorem~\ref{r5d-thm-schur-det}(b) gives, with Definition~\ref{r5d-def-det}(a) for the factor \(t-4\)
in both rows of the \(2\times2\) determinant,
\[
\begin{aligned}
\det(tI_4-\Sigma)&=(t-4)^2\det\Bigl((t-4)I_2-\frac1{t-4}BB'\Bigr)=\det\bigl((t-4)^2I_2-BB'\bigr)\\
&=\bigl((t-4)^2-5\bigr)\bigl((t-4)^2-8\bigr)-4\\
&=\bigl((t-4)^2-4\bigr)\bigl((t-4)^2-9\bigr)=(t-6)(t-2)(t-7)(t-1)
\end{aligned}
\]
This holds for every real \(t\ne4\), hence as polynomials (Definition~\ref{r7-def-charpoly}(a),
Proposition~\ref{r7-prop-poly}(d)): \(p_\Sigma(t)=(t-7)(t-6)(t-2)(t-1)\), and the eigenvalues are \(7\), \(6\), \(2\), \(1\).
Eigenvectors: row-reduce \(\Sigma-\lambda I\) and solve \((\Sigma-\lambda I)c=0\).
\[
\begin{gathered}
\Sigma-7I=\begin{pmatrix}-3&0&2&-1\\0&-3&-2&-2\\2&-2&-3&0\\-1&-2&0&-3\end{pmatrix}
\xrightarrow{R_1\leftrightarrow R_4}
\begin{pmatrix}-1&-2&0&-3\\0&-3&-2&-2\\2&-2&-3&0\\-3&0&2&-1\end{pmatrix}
\xrightarrow[R_4\leftarrow R_4-3R_1]{R_3\leftarrow R_3+2R_1}
\begin{pmatrix}-1&-2&0&-3\\0&-3&-2&-2\\0&-6&-3&-6\\0&6&2&8\end{pmatrix}
\\
\xrightarrow[R_4\leftarrow R_4+2R_2]{R_3\leftarrow R_3-2R_2}
\begin{pmatrix}-1&-2&0&-3\\0&-3&-2&-2\\0&0&1&-2\\0&0&-2&4\end{pmatrix}
\xrightarrow{R_4\leftarrow R_4+2R_3}
\begin{pmatrix}-1&-2&0&-3\\0&-3&-2&-2\\0&0&1&-2\\0&0&0&0\end{pmatrix}
\end{gathered}
\]
so \(c_4=1\) gives \(c_3=2c_4=2\), \(c_2=-\frac13(2c_3+2c_4)=-2\) and \(c_1=-2c_2-3c_4=1\); take \((1,-2,2,1)'\).
\[
\begin{gathered}
\Sigma-6I=\begin{pmatrix}-2&0&2&-1\\0&-2&-2&-2\\2&-2&-2&0\\-1&-2&0&-2\end{pmatrix}
\xrightarrow{R_1\leftrightarrow R_4}
\begin{pmatrix}-1&-2&0&-2\\0&-2&-2&-2\\2&-2&-2&0\\-2&0&2&-1\end{pmatrix}
\xrightarrow[R_4\leftarrow R_4-2R_1]{R_3\leftarrow R_3+2R_1}
\begin{pmatrix}-1&-2&0&-2\\0&-2&-2&-2\\0&-6&-2&-4\\0&4&2&3\end{pmatrix}
\\
\xrightarrow[R_4\leftarrow R_4+2R_2]{R_3\leftarrow R_3-3R_2}
\begin{pmatrix}-1&-2&0&-2\\0&-2&-2&-2\\0&0&4&2\\0&0&-2&-1\end{pmatrix}
\xrightarrow{R_4\leftarrow R_4+\frac12R_3}
\begin{pmatrix}-1&-2&0&-2\\0&-2&-2&-2\\0&0&4&2\\0&0&0&0\end{pmatrix}
\end{gathered}
\]
so \(c_4=-2\) gives \(c_3=-\frac12c_4=1\), \(c_2=-c_3-c_4=1\) and \(c_1=-2c_2-2c_4=2\); take \((2,1,1,-2)'\).
\[
\begin{gathered}
\Sigma-2I=\begin{pmatrix}2&0&2&-1\\0&2&-2&-2\\2&-2&2&0\\-1&-2&0&2\end{pmatrix}
\xrightarrow{R_1\leftrightarrow R_4}
\begin{pmatrix}-1&-2&0&2\\0&2&-2&-2\\2&-2&2&0\\2&0&2&-1\end{pmatrix}
\xrightarrow[R_4\leftarrow R_4+2R_1]{R_3\leftarrow R_3+2R_1}
\begin{pmatrix}-1&-2&0&2\\0&2&-2&-2\\0&-6&2&4\\0&-4&2&3\end{pmatrix}
\\
\xrightarrow[R_4\leftarrow R_4+2R_2]{R_3\leftarrow R_3+3R_2}
\begin{pmatrix}-1&-2&0&2\\0&2&-2&-2\\0&0&-4&-2\\0&0&-2&-1\end{pmatrix}
\xrightarrow{R_4\leftarrow R_4-\frac12R_3}
\begin{pmatrix}-1&-2&0&2\\0&2&-2&-2\\0&0&-4&-2\\0&0&0&0\end{pmatrix}
\end{gathered}
\]
so \(c_4=2\) gives \(c_3=-\frac12c_4=-1\), \(c_2=c_3+c_4=1\) and \(c_1=-2c_2+2c_4=2\); take \((2,1,-1,2)'\).
\[
\begin{gathered}
\Sigma-I=\begin{pmatrix}3&0&2&-1\\0&3&-2&-2\\2&-2&3&0\\-1&-2&0&3\end{pmatrix}
\xrightarrow{R_1\leftrightarrow R_4}
\begin{pmatrix}-1&-2&0&3\\0&3&-2&-2\\2&-2&3&0\\3&0&2&-1\end{pmatrix}
\xrightarrow[R_4\leftarrow R_4+3R_1]{R_3\leftarrow R_3+2R_1}
\begin{pmatrix}-1&-2&0&3\\0&3&-2&-2\\0&-6&3&6\\0&-6&2&8\end{pmatrix}
\\
\xrightarrow[R_4\leftarrow R_4+2R_2]{R_3\leftarrow R_3+2R_2}
\begin{pmatrix}-1&-2&0&3\\0&3&-2&-2\\0&0&-1&2\\0&0&-2&4\end{pmatrix}
\xrightarrow{R_4\leftarrow R_4-2R_3}
\begin{pmatrix}-1&-2&0&3\\0&3&-2&-2\\0&0&-1&2\\0&0&0&0\end{pmatrix}
\end{gathered}
\]
so \(c_4=1\) gives \(c_3=2c_4=2\), \(c_2=\frac13(2c_3+2c_4)=2\) and \(c_1=-2c_2+3c_4=-1\); take \((-1,2,2,1)'\). The four
vectors are pairwise orthogonal and have length \(\sqrt{10}\), so \(v_1=\frac1{\sqrt{10}}(1,-2,2,1)'\),
\(v_2=\frac1{\sqrt{10}}(2,1,1,-2)'\), \(v_3=\frac1{\sqrt{10}}(2,1,-1,2)'\) and \(v_4=\frac1{\sqrt{10}}(-1,2,2,1)'\) are orthonormal
eigenvectors for \(7\), \(6\), \(2\), \(1\). Check: \(\Sigma(1,-2,2,1)'=(7,-14,14,7)'\), \(\Sigma(2,1,1,-2)'=(12,6,6,-12)'\),
\(\Sigma(2,1,-1,2)'=(4,2,-2,4)'\), \(\Sigma(-1,2,2,1)'=(-1,2,2,1)'\).

(b) By Theorem~\ref{r12-thm-pop-pca}(b),(c), the principal components are \(v_1'x=\frac1{\sqrt{10}}(x_1-2x_2+2x_3+x_4)\),
\(v_2'x=\frac1{\sqrt{10}}(2x_1+x_2+x_3-2x_4)\), \(v_3'x=\frac1{\sqrt{10}}(2x_1+x_2-x_3+2x_4)\) and
\(v_4'x=\frac1{\sqrt{10}}(-x_1+2x_2+2x_3+x_4)\), with variances \(7\), \(6\), \(2\), \(1\). The total variance is
\(\tr\Sigma=16\), so they account for \(\frac7{16}\), \(\frac6{16}=\frac38\), \(\frac2{16}=\frac18\) and \(\frac1{16}\) of it. The
cumulative proportions are \(\frac7{16}\), \(\frac{13}{16}\), \(\frac{15}{16}\) and \(1\); as
\(\frac7{16}=\frac{35}{80}<\frac{64}{80}=\frac45\leqslant\frac{65}{80}=\frac{13}{16}\), two components account for at least
\(80\%\), one does not.

(c) \(\cov(v_1'x,v_2'x)=v_1'\Sigma v_2=\frac1{10}(1,-2,2,1)(12,6,6,-12)'=\frac1{10}(12-12+12-12)=0\)
(Theorem~\ref{r12-thm-pop-pca}(a)). And \(\var(x_1)+\var(x_2)=4+4=8\leqslant7+6=13\), as
Corollary~\ref{r12-cor-diagonal} says (every pair of the variables has the variances \(4+4=8\)).
\end{solution}

\needspace{14\baselineskip}
\subsection*{Problems}

\begin{problem}[\citeat{casella-berger}{Exercise~4.43}, \citeat{casella-berger-solutions}{Exercise~4.43}]\label{r14-ex-cov-sums}
Let \(x_1,x_2,x_3\) be uncorrelated random variables, each with mean \(\mu\) and variance \(\sigma^2\). Find
\(\cov(x_1+x_2,\ x_2+x_3)\) and \(\cov(x_1+x_2,\ x_1-x_2)\).
\end{problem}
\begin{solution}
The \((1,2)\) entry of Theorem~\ref{r14-thm-cov}(e) for the random vector \((u,v)'\) reads
\(\expc(uv)=\cov(u,v)+\expc(u)\expc(v)\), that is,
\(\cov(u,v)=\expc(uv)-\expc(u)\expc(v)\). So \(\expc(x_ix_j)=\cov(x_i,x_j)+\mu^2=\mu^2\) for
\(i\ne j\), as the \(x_i\) are uncorrelated, and \(\expc(x_i^2)=\sigma^2+\mu^2\). Hence
\[
\cov(x_1+x_2,\ x_2+x_3)=\expc(x_1x_2+x_1x_3+x_2^2+x_2x_3)-(2\mu)(2\mu)=(4\mu^2+\sigma^2)-4\mu^2=\sigma^2
\]
and \(\cov(x_1+x_2,\ x_1-x_2)=\expc(x_1^2-x_2^2)-(2\mu)\cdot0=(\sigma^2+\mu^2)-(\sigma^2+\mu^2)=0\).
\end{solution}

\begin{problem}[\citeat{searle-khuri}{Exercise~10.2}]\label{r14-ex-shift}
Let \(X\in\R^{n\times p}\) be a data matrix and \(\theta\in\R\). Show that subtracting \(\theta\) from every entry of \(X\) does not
change \(X'C_nX\), and hence does not change the sample covariance matrix \(S\).
\end{problem}
\begin{solution}
The new data matrix is \(X^*=X-\theta\one_n\one_p'\). As \(C_n\one_n=0\) and \(C_n'=C_n\) (Proposition~\ref{r1-prop-centering}),
\(C_n\one_n\one_p'=0\) and \(\one_p\one_n'C_n=(C_n\one_n\one_p')'=0\). So
\[
X^{*\prime}C_nX^*=X'C_nX-\theta X'C_n\one_n\one_p'-\theta\one_p\one_n'C_nX+\theta^2\one_p\one_n'C_n\one_n\one_p'=X'C_nX
\]
and \(S=\frac1{n-1}X'C_nX\) (Definition~\ref{r12-def-data}) is unchanged.
\end{solution}

\begin{problem}[\citeat{searle-khuri}{Example~10.1}]\label{r8-pr-equicorrelation}
Let \(p\geqslant2\) and \(\rho<-\frac1{p-1}\). Can \(p\) variables of a data matrix have all pairwise sample correlations
equal to \(\rho\)?
\end{problem}
\begin{solution}
No. Their sample correlation matrix would be the equicorrelation matrix \(R=(1-\rho)I_p+\rho J_p\) of
Example~\ref{r5d-exm-equicorrelation}. A sample correlation matrix is PSD (Proposition~\ref{r12-prop-correlation}(a)), so
its determinant is \(\geqslant0\) (Theorem~\ref{r8-thm-eigen-criterion}(c)). But for \(\rho<-\frac1{p-1}\),
\(1-\rho>0\) and \(1+(p-1)\rho<0\), so \(\det R=(1-\rho)^{p-1}\bigl(1+(p-1)\rho\bigr)<0\)
(Example~\ref{r5d-exm-equicorrelation}). For three variables, even
unequal pairwise sample correlations \(r_{12},r_{13},r_{23}\) cannot all be below \(-\frac12\): at \(x=\one_3\), the
quadratic form of their sample correlation matrix takes the value \(3+2(r_{12}+r_{13}+r_{23})\)
(Definition~\ref{r8-def-quadratic-form}), which is \(\geqslant0\); so \(r_{12}+r_{13}+r_{23}\geqslant-\frac32\).
\end{solution}

\begin{problem}[\citeat{searle-khuri}{Exercise~6.19}]\label{r9-ex-diagonal-bounds}
Let \(A=(a_{ij})\in\R^{n\times n}\) be symmetric. Show that
\(\lambda_n(A)\leqslant\frac1n\sum_{i=1}^n\sum_{j=1}^na_{ij}\leqslant\lambda_1(A)\).
\end{problem}
\begin{solution}
\(\sum_i\sum_ja_{ij}=\one_n'A\one_n\) and \(\one_n'\one_n=n\) (Proposition~\ref{r1-prop-centering}), so by
Theorem~\ref{r9-thm-rayleigh}(a), \(\lambda_n(A)\leqslant\one_n'A\one_n/\one_n'\one_n\leqslant\lambda_1(A)\).
\end{solution}

\begin{problem}[\citeat{abadir-magnus}{Exercise~12.40(a)}]\label{r9-ex-sum-bounds}
Let \(A,B\in\R^{n\times n}\) be symmetric. Show that \(\lambda_1(A+B)\leqslant\lambda_1(A)+\lambda_1(B)\) and
\(\lambda_n(A+B)\geqslant\lambda_n(A)+\lambda_n(B)\).
\end{problem}
\begin{solution}
By Theorem~\ref{r9-thm-rayleigh}(a), for every \(x\ne0\)
\[
\lambda_n(A)+\lambda_n(B)\leqslant\frac{x'(A+B)x}{x'x}=\frac{x'Ax}{x'x}+\frac{x'Bx}{x'x}\leqslant\lambda_1(A)+\lambda_1(B)
\]
For an eigenvector \(x\) of \(A+B\) belonging to \(\lambda_1(A+B)\) the middle term is \(\lambda_1(A+B)\), which gives the
first inequality; for one belonging to \(\lambda_n(A+B)\) it is \(\lambda_n(A+B)\), which gives the second.
\end{solution}

\begin{problem}[\citeat{searle-khuri}{Exercise~6.21}]\label{r9-ex-trace-bounds}
Let \(A,B\in\R^{n\times n}\) be symmetric, with \(A\) PSD. Show that
\[
\lambda_n(B)\tr A\leqslant\tr(AB)\leqslant\lambda_1(B)\tr A
\]
\end{problem}
\begin{solution}
\begin{itemize}
\item Let \(N=A^{1/2}\) (Theorem~\ref{r8-thm-sqrt}); it is symmetric and \(N^2=A\). By
Theorem~\ref{r9-thm-rayleigh}(c), \(B-\lambda_n(B)I\) and \(\lambda_1(B)I-B\) are PSD.
\item If \(M\) is PSD, so is \(NMN\): it is symmetric, and \(x'NMNx=(Nx)'M(Nx)\geqslant0\). Hence
\(NBN-\lambda_n(B)A\) and \(\lambda_1(B)A-NBN\) are PSD.
\item A PSD matrix has eigenvalues \(\geqslant0\) (Theorem~\ref{r8-thm-eigen-criterion}(a)), hence trace \(\geqslant0\)
(Theorem~\ref{r7-thm-trace-det}). By Theorem~\ref{r1-thm-trace}(a), \(\lambda_n(B)\tr A\leqslant\tr(NBN)\leqslant
\lambda_1(B)\tr A\).
\item \(\tr(NBN)=\tr(BNN)=\tr(BA)=\tr(AB)\) (Theorem~\ref{r1-thm-trace}(b)).
\end{itemize}
\end{solution}

\begin{problem}[\citeat{searle-khuri}{Exercise~6.24}]\label{r9-ex-trace-square}
Let \(A\in\R^{n\times n}\) be PSD. Show that \(\tr(A^2)\leqslant\lambda_1(A)\tr A\).
\end{problem}
\begin{solution}
With \(Q\) as in Lemma~\ref{r9-lem-eig-order}, \(Q'AQ=\diag(\lambda_1,\dots,\lambda_n)\) and \(Q^{-1}=Q'\)
(Proposition~\ref{r6-prop-orthonormal}(b)), so \(A=Q\diag(\lambda_1,\dots,\lambda_n)Q^{-1}\). By
Proposition~\ref{r7-prop-diag-rank}(b) with \(f(t)=t^2\), \(\tr(A^2)=\sum_i\lambda_i^2\). As \(A\) is PSD,
\(0\leqslant\lambda_i\leqslant\lambda_1\) (Theorem~\ref{r8-thm-eigen-criterion}(a)), so \(\lambda_i^2\leqslant\lambda_1\lambda_i\).
Hence \(\tr(A^2)\leqslant\lambda_1\sum_i\lambda_i=\lambda_1\tr A\) (Theorem~\ref{r7-thm-trace-det}).
\end{solution}

\begin{problem}[\citeat{prasolov}{Problem~38.1}]\label{r12-ex-pr-doubly-stochastic}
A matrix \(A=(a_{ij})\in\R^{n\times n}\) is \textbf{doubly stochastic} if all its entries are \(\geqslant0\) and all its row
sums and column sums are \(1\), that is, \(A\one_n=\one_n\) and \(\one_n'A=\one_n'\). Let \(A\) be doubly stochastic, and let
\(x_1\geqslant x_2\geqslant\dots\geqslant x_n\geqslant0\) and \(y_1\geqslant y_2\geqslant\dots\geqslant y_n\geqslant0\). Show
that
\[
\sum_{r=1}^n\sum_{s=1}^na_{rs}x_ry_s\leqslant\sum_{r=1}^nx_ry_r
\]
\end{problem}
\begin{solution}
All sums below run over \(1,\dots,n\) unless limits are shown, and \(\delta_{rs}\) is \(1\) for \(r=s\) and \(0\) for \(r\ne s\).
\begin{itemize}
\item Put \(x_{n+1}=y_{n+1}=0\), \(\xi_i=x_i-x_{i+1}\) and \(\eta_j=y_j-y_{j+1}\) (\(1\leqslant i,j\leqslant n\)). By the
hypotheses, \(\xi_i\geqslant0\) and \(\eta_j\geqslant0\), and the sums telescope: \(x_r=\sum_{i\geqslant r}\xi_i\) and
\(y_s=\sum_{j\geqslant s}\eta_j\).
\item As \(\sum_rx_ry_r=\sum_{r,s}\delta_{rs}x_ry_s\),
\[
\sum_rx_ry_r-\sum_{r,s}a_{rs}x_ry_s=\sum_{r,s}(\delta_{rs}-a_{rs})\sum_{i\geqslant r}\xi_i\sum_{j\geqslant s}\eta_j
=\sum_{i,j}c_{ij}\xi_i\eta_j,\qquad c_{ij}=\sum_{r\leqslant i}\sum_{s\leqslant j}(\delta_{rs}-a_{rs})
\]
since the product \((\delta_{rs}-a_{rs})\xi_i\eta_j\) occurs on both sides exactly for \(r\leqslant i\) and \(s\leqslant j\).
As every \(\xi_i\eta_j\geqslant0\), it suffices to show that every \(c_{ij}\geqslant0\).
\item Let \(i\leqslant j\). Among the pairs with \(r\leqslant i\) and \(s\leqslant j\), those with \(r=s\) are
\((1,1),\dots,(i,i)\), so \(\sum_{r\leqslant i}\sum_{s\leqslant j}\delta_{rs}=i\). As \(a_{rs}\geqslant0\) and the row sums of \(A\) are \(1\),
\(\sum_{r\leqslant i}\sum_{s\leqslant j}a_{rs}\leqslant\sum_{r\leqslant i}\sum_{s=1}^na_{rs}=i\). So \(c_{ij}\geqslant0\).
\item Let \(i\geqslant j\). Likewise the pairs with \(r=s\) are \((1,1),\dots,(j,j)\), so
\(\sum_{r\leqslant i}\sum_{s\leqslant j}\delta_{rs}=j\). As \(a_{rs}\geqslant0\) and the column sums of \(A\) are \(1\),
\(\sum_{r\leqslant i}\sum_{s\leqslant j}a_{rs}\leqslant\sum_{s\leqslant j}\sum_{r=1}^na_{rs}=j\). So \(c_{ij}\geqslant0\).
\end{itemize}
If \(P=(p_{ij})\in\R^{n\times n}\) is orthogonal, the matrix \((p_{ij}^2)\) is doubly stochastic, as the rows and the columns
of \(P\) are unit vectors (Theorem~\ref{r6-thm-orthogonal}(a)). In the proof of Theorem~\ref{r12-thm-kyfan}(a), with
\(P=[R\ T]\), the weight \(w_i\) is the sum of the first \(k\) entries of row \(i\) of \((p_{ij}^2)\).
\end{solution}

\clearpage
\def\unitnumber{10}
\setcounter{theorem}{0}\setcounter{problem}{0}
\section*{Chapter 10. The multivariate normal distribution}

\subsection*{The multivariate normal distribution}

We write \(\prb(z\in A)\) for the probability that a random vector \(z\in\R^n\) takes a value in a set
\(A\subseteq\R^n\). Two random vectors \(y,z\in\R^n\) have the \textbf{same distribution} if
\(\prb(y\in A)=\prb(z\in A)\) for every set \(A\); then also \(\expc g(y)=\expc g(z)\) for every
function \(g\) for which these expectations exist. If a random variable \(u\) takes only the values
\(a_1,\dots,a_m\), then \(\expc g(u)=g(a_1)\prb(u=a_1)+\dots+g(a_m)\prb(u=a_m)\) for every
function \(g\). Random vectors \(u_1,\dots,u_m\) are \textbf{independent} if
\(\prb(u_1\in A_1,\dots,u_m\in A_m)=\prb(u_1\in A_1)\cdots\prb(u_m\in A_m)\) for all sets
\(A_1,\dots,A_m\). The \textbf{moment generating function} (MGF) of \(z\in\R^n\) is
\[
M_z(t)=\expc\bigl(\exp(t'z)\bigr)\qquad(t\in\R^n)
\]
defined wherever this expectation is finite~\citeat{tu-clm}{\S4.1}, \citeat{casella-berger}{Definitions~1.5.8, 2.2.1, Theorem~4.2.10(a)}. (Here and below, the sets and functions are those for which the
probabilities and expectations are defined.)

\begin{definition}[\citeat{tu-clm}{Definition~4.1}]\label{r14-def-mvn}
Let \(\mu\in\R^n\), and let \(\Sigma\in\R^{n\times n}\) be symmetric and PSD. A random vector \(z\in\R^n\) has the
\textbf{multivariate normal distribution} \(N_n(\mu,\Sigma)\), written \(z\sim N_n(\mu,\Sigma)\), if
\[
M_z(t)=\exp\bigl(t'\mu+\tfrac12t'\Sigma t\bigr)\qquad\text{for all }t\in\R^n
\]
Its mean vector is \(\mu\) and its covariance matrix is \(\Sigma\): \(\expc(z)=\mu\) and \(\var(z)=\Sigma\);
and \(N_n(\mu,0)\) is the distribution of the constant vector \(\mu\) (both shown below). For \(n=1\) we write \(N(\mu,\sigma^2)\)
(\(\sigma^2\geqslant0\)), with the MGF \(\exp(t\mu+\frac12\sigma^2t^2)\); \(N(0,1)\) is the \textbf{standard normal
distribution}. Random vectors \(u\) and \(v\) are \textbf{jointly normal} if \((u',v')'\) has a multivariate normal distribution.
\end{definition}

\needspace{8\baselineskip}
We use the following facts without proof~\citeat{tu-clm}{\S4.1}.
\begin{enumerate}[label=(N\arabic*)]
\item \emph{Uniqueness.} If two random vectors have MGFs that are finite and equal on an open neighborhood of \(0\),
they have the same distribution.
\item \emph{Independence and factorization.} Random vectors \(u_1\in\R^{n_1},\dots,u_m\in\R^{n_m}\) whose MGFs are finite
near \(0\) are independent if and only if the MGF of \((u_1',\dots,u_m')'\) factors:
\(M_{(u_1',\dots,u_m')'}(t_1,\dots,t_m)=M_{u_1}(t_1)\cdots M_{u_m}(t_m)\) for all \(t_k\in\R^{n_k}\) near \(0\).
\item \emph{Functions of independent vectors.} If \(u_1,\dots,u_m\) are independent, then so are
\(g_1(u_1),\dots,g_m(u_m)\), for any functions \(g_1,\dots,g_m\).
\item \emph{Density in the nonsingular case.} This fact is not assumed here: the density of \(N_n(\mu,\Sigma)\) for PD
\(\Sigma\) is derived below from the change-of-variables formula for integrals.
\item \emph{Distributions derived from the normal.} The \textbf{chi-square distribution} \(\chi^2_r\) (\(r\geqslant1\)) is
the distribution of \(w_1^2+\dots+w_r^2\) for independent \(w_1,\dots,w_r\sim N(0,1)\).
\item \emph{Moments from the MGF.} If \(M_z\) is finite on an open neighborhood of \(0\), its partial derivatives of the first
and second order exist there, and at \(t=0\), \(\partial M_z/\partial t_i=\expc(z_i)\) and
\(\partial^2M_z/\partial t_i\partial t_j=\expc(z_iz_j)\).
\item \emph{Independence implies uncorrelatedness.} If \(u\) and \(v\) are independent random vectors, then
\(\expc(uv')=\expc(u)\expc(v)'\), that is, \(\cov(u,v)=0\) (Definition~\ref{r14-def-cov}). In
particular \(\expc(uv)=\expc(u)\expc(v)\) for independent random variables \(u\) and \(v\).
\end{enumerate}

Let \(z\sim N_n(\mu,\Sigma)\). The exponent \(q(t)=t'\mu+\frac12t'\Sigma t=\sum_kt_k\mu_k+\frac12\sum_k\sum_l\sigma_{kl}t_kt_l\) has, as
\(\sigma_{kl}=\sigma_{lk}\), the partial derivatives \(\partial q/\partial t_i=\mu_i+\sum_l\sigma_{il}t_l=\mu_i+(\Sigma t)_i\) and
\(\partial^2q/\partial t_j\partial t_i=\sigma_{ij}\). By the chain rule and the product rule,
\[
\frac{\partial M_z}{\partial t_i}=\bigl(\mu_i+(\Sigma t)_i\bigr)M_z(t),\qquad
\frac{\partial^2M_z}{\partial t_j\partial t_i}=\Bigl(\sigma_{ij}+\bigl(\mu_i+(\Sigma t)_i\bigr)\bigl(\mu_j+(\Sigma t)_j\bigr)\Bigr)M_z(t)
\]
At \(t=0\), \(M_z(0)=1\), so by (N6) \(\expc(z_i)=\mu_i\) and \(\expc(z_iz_j)=\sigma_{ij}+\mu_i\mu_j\). Hence
\(\cov(z_i,z_j)=\expc(z_iz_j)-\mu_i\mu_j=\sigma_{ij}\) (the \((i,j)\) entry of
Theorem~\ref{r14-thm-cov}(e)): \(\expc(z)=\mu\) and \(\var(z)=\Sigma\), as claimed in
Definition~\ref{r14-def-mvn}~\citeat{tu-clm}{Definition~4.1}. If \(\Sigma=0\), the constant vector \(\mu\) has the MGF \(\expc(\exp(t'\mu))=\exp(t'\mu)\); so by
(N1), \(z\sim N_n(\mu,0)\) has the distribution of the constant \(\mu\), and \(\prb(z=\mu)=1\).

\subsection*{Linear transformations and marginal distributions}

\begin{theorem}[\citeat{tu-clm}{Theorem~4.2}, \citeat{searle-khuri}{(10.28)--(10.32)}, \citeat{bapat}{\S8.1}]\label{r14-thm-mvn-linear}
Let \(z\sim N_n(\mu,\Sigma)\), and let \(A\in\R^{m\times n}\) and \(b\in\R^m\) be constant. Then
\[
Az+b\sim N_m(A\mu+b,\ A\Sigma A')
\]
\end{theorem}
\begin{proof}
\begin{itemize}
\item \(A\Sigma A'=(A')'\Sigma(A')\) is symmetric and PSD (Theorem~\ref{r8-thm-congruence}(a)), so \(N_m(A\mu+b,A\Sigma A')\) is
defined.
\item For \(t\in\R^m\), \(t'(Az+b)=t'b+(A't)'z\), so \(\exp\bigl(t'(Az+b)\bigr)=\exp(t'b)\exp\bigl((A't)'z\bigr)\), and
\(M_{Az+b}(t)=\exp(t'b)M_z(A't)\), as \(\exp(t'b)\) is a constant.
\item By Definition~\ref{r14-def-mvn},
\[
M_{Az+b}(t)=\exp\bigl(t'b+(A't)'\mu+\tfrac12(A't)'\Sigma(A't)\bigr)=\exp\bigl(t'(A\mu+b)+\tfrac12t'(A\Sigma A')t\bigr)
\]
for every \(t\in\R^m\), which is the MGF of \(N_m(A\mu+b,A\Sigma A')\).\qedhere
\end{itemize}
\end{proof}

\needspace{14\baselineskip}
\begin{corollary}[\citeat{bapat}{\S8.1}, \citeat{casella-berger}{\S4.5}]\label{r14-cor-marginal}
Let \(z\sim N_n(\mu,\Sigma)\).
\begin{enumerate}
\item Partition \(z=(z_1',z_2')'\) with \(z_1\in\R^p\), \(z_2\in\R^q\), and accordingly
\[
\mu=\begin{pmatrix}\mu_1\\\mu_2\end{pmatrix},\qquad\Sigma=\begin{pmatrix}\Sigma_{11}&\Sigma_{12}\\\Sigma_{21}&\Sigma_{22}\end{pmatrix}
\qquad(\Sigma_{11}\in\R^{p\times p})
\]
Then \(z_1\sim N_p(\mu_1,\Sigma_{11})\) and \(z_2\sim N_q(\mu_2,\Sigma_{22})\).
\item For \(1\leqslant i_1<\dots<i_k\leqslant n\), the subvector \((z_{i_1},\dots,z_{i_k})'\) is \(N_k\), with the mean
\((\mu_{i_1},\dots,\mu_{i_k})'\) and as covariance matrix the principal submatrix of \(\Sigma\) in the rows and columns
\(i_1,\dots,i_k\). In particular \(z_i\sim N(\mu_i,\sigma_{ii})\).
\item For \(a\in\R^n\), \(a'z\sim N(a'\mu,\ a'\Sigma a)\).
\end{enumerate}
\end{corollary}
\begin{proof}
(a) For \(t_1\in\R^p\), \(t_1'z_1=(t_1',0')z\), so \(M_{z_1}(t_1)=M_z\bigl((t_1',0')'\bigr)\). By the block product
(Theorem~\ref{r5b-thm-block}(c)), \((t_1',0')\mu=t_1'\mu_1\) and \((t_1',0')\Sigma(t_1',0')'=t_1'\Sigma_{11}t_1\); so
\(M_{z_1}(t_1)=\exp(t_1'\mu_1+\frac12t_1'\Sigma_{11}t_1)\) for every \(t_1\), where \(\Sigma_{11}\) is symmetric and PSD
(Theorem~\ref{r8-thm-principal-sub}). Likewise for \(z_2\), with \(t=(0',t_2')'\).

(b) With \(T=[e_{i_1}\ \cdots\ e_{i_k}]\in\R^{n\times k}\) (the matrix \(S\) of Definition~\ref{r8-def-principal}), the subvector is
\(T'z\) (Theorem~\ref{r1-thm-views}(d)), and \(T'\Sigma T\) is the principal submatrix; apply Theorem~\ref{r14-thm-mvn-linear}
with \(A=T'\) and \(b=0\).

(c) Theorem~\ref{r14-thm-mvn-linear} with \(A=a'\) and \(b=0\).
\end{proof}

\begin{example}\label{r14-exm-marginal}
Let \(z\sim N_4(\mu,\Sigma)\) with \(\mu=(1,0,2,2)'\) and \(\Sigma\) of Example~\ref{r14-exm-cov}.
\begin{itemize}
\item By Corollary~\ref{r14-cor-marginal}(b), \(z_2\sim N(0,1)\) and
\[
\begin{pmatrix}z_1\\z_3\end{pmatrix}\sim N_2\left(\begin{pmatrix}1\\2\end{pmatrix},\begin{pmatrix}4&4\\4&9\end{pmatrix}\right)
\]
\item By Corollary~\ref{r14-cor-marginal}(c) and Example~\ref{r14-exm-cov}, \(z_1+z_2+z_3+z_4\sim N(5,51)\), as
\(\one'\mu=1+0+2+2=5\); \(e=z_3-z_1-2z_2\sim N(1,1)\), as \(a'\mu=-1-0+2+0=1\) for \(a=(-1,-2,1,0)'\); and
\(f=z_4+2z_1+2z_2-2z_3\sim N(0,1)\), as \(b'\mu=2+0-4+2=0\) for \(b=(2,2,-2,1)'\).
\item With \(A=\left(\begin{smallmatrix}1&1&0&0\\1&-1&0&0\end{smallmatrix}\right)\), Theorem~\ref{r14-thm-mvn-linear} and
Example~\ref{r14-exm-cov} give
\[
\begin{pmatrix}z_1+z_2\\z_1-z_2\end{pmatrix}=Az\sim N_2\left(\begin{pmatrix}1\\1\end{pmatrix},\begin{pmatrix}5&3\\3&5\end{pmatrix}\right)
\]
\end{itemize}
\end{example}

\subsection*{Independence}

\begin{theorem}[\citeat{tu-clm}{Theorems~4.3, 4.5}, \citeat{bapat}{Results~8.1, 8.2}]\label{r14-thm-mvn-indep}
\begin{enumerate}
\item Let \(u\in\R^p\) and \(v\in\R^q\) be jointly normal, \(z=(u',v')'\sim N_{p+q}(\mu,\Sigma)\), with
\[
\Sigma=\begin{pmatrix}\Sigma_{uu}&\Sigma_{uv}\\\Sigma_{vu}&\Sigma_{vv}\end{pmatrix},\qquad\Sigma_{uu}\in\R^{p\times p}
\]
Then \(u\) and \(v\) are independent if and only if \(\Sigma_{uv}=0\), that is, if and only if \(u\) and \(v\) are
uncorrelated.
\item Let \(z=(z_1',\dots,z_m')'\sim N_n(\mu,\Sigma)\), with subvectors \(z_k\in\R^{n_k}\) and \(\Sigma=(\Sigma_{kl})\) partitioned
accordingly. Then \(z_1,\dots,z_m\) are independent if and only if \(\Sigma_{kl}=0\) for all \(k\ne l\).
\item Let \(z\sim N_n(\mu,\Sigma)\), and let \(A\in\R^{p\times n}\) and \(B\in\R^{q\times n}\) be constant. Then \(Az\) and \(Bz\) are
independent if and only if \(A\Sigma B'=0\).
\end{enumerate}
\end{theorem}
\begin{proof}
(a) Partition \(\mu=(\mu_u',\mu_v')'\). By Definition~\ref{r14-def-mvn} and the block form of \(\var(z)\) after
Definition~\ref{r14-def-cov}, \(\Sigma_{uv}=\cov(u,v)\).
\begin{itemize}
\item If \(u\) and \(v\) are independent, then \(\Sigma_{uv}=\cov(u,v)=0\) by (N7).
\item Let \(\Sigma_{uv}=0\); then \(\Sigma_{vu}=\Sigma_{uv}'=0\). For \(s\in\R^p\) and \(t\in\R^q\), the block product
(Theorem~\ref{r5b-thm-block}(c)) gives \((s',t')\mu=s'\mu_u+t'\mu_v\) and \((s',t')\Sigma(s',t')'=s'\Sigma_{uu}s+t'\Sigma_{vv}t\).
So, by Corollary~\ref{r14-cor-marginal}(a),
\[
M_z(s,t)=\exp\bigl(s'\mu_u+\tfrac12s'\Sigma_{uu}s\bigr)\exp\bigl(t'\mu_v+\tfrac12t'\Sigma_{vv}t\bigr)=M_u(s)M_v(t)
\]
and \(u\) and \(v\) are independent by (N2).
\end{itemize}
(b)
\begin{itemize}
\item If \(\Sigma_{kl}=0\) for \(k\ne l\), then for \(t=(t_1',\dots,t_m')'\) the block product gives \(t'\Sigma t=\sum_kt_k'\Sigma_{kk}t_k\),
so \(M_z(t)=\prod_k\exp(t_k'\mu_k+\frac12t_k'\Sigma_{kk}t_k)=\prod_kM_{z_k}(t_k)\), as \(z_k\sim N_{n_k}(\mu_k,\Sigma_{kk})\)
(Corollary~\ref{r14-cor-marginal}(b)). By (N2), \(z_1,\dots,z_m\) are independent.
\item Conversely, let \(z_1,\dots,z_m\) be independent, and \(k\ne l\). By (N2), \(M_z(t)=\prod_jM_{z_j}(t_j)\). Taking \(t_j=0\)
for \(j\ne k,l\), and as \(M_{z_j}(0)=1\), the MGF of \((z_k',z_l')'\) is \(M_{z_k}(t_k)M_{z_l}(t_l)\). So \(z_k\) and \(z_l\) are
independent (N2), and \(\Sigma_{kl}=\cov(z_k,z_l)=0\) by (N7).
\end{itemize}
(c)
\begin{itemize}
\item Let \(G=\left(\begin{smallmatrix}A\\B\end{smallmatrix}\right)\in\R^{(p+q)\times n}\). By Theorem~\ref{r14-thm-mvn-linear},
\(Gz=\bigl((Az)',(Bz)'\bigr)'\sim N_{p+q}(G\mu,G\Sigma G')\), and by the block product
\[
G\Sigma G'=\begin{pmatrix}A\\B\end{pmatrix}\Sigma\,(A'\ \ B')=\begin{pmatrix}A\Sigma A'&A\Sigma B'\\B\Sigma A'&B\Sigma B'\end{pmatrix}
\]
\item By (a) for \(u=Az\) and \(v=Bz\), they are independent if and only if \(A\Sigma B'=0\).\qedhere
\end{itemize}
\end{proof}

\begin{remark}[\citeat{tu-clm}{Remark~4.4}, \citeat{casella-berger}{\S4.5}]\label{r14-rem-uncorr}
Joint normality is needed in Theorem~\ref{r14-thm-mvn-indep}: uncorrelated random variables need not be independent,
even if each of them is normal.
\begin{itemize}
\item Let \(u\) take the values \(-1\), \(0\), \(1\), each with probability \(\frac13\), and let \(v=u^2\). Then
\(\expc(u)=\expc(u^3)=0\), so \(\cov(u,v)=\expc(u^3)-\expc(u)\expc(u^2)=0\);
but \(\prb(u=0,\ v=1)=0\ne\frac13\cdot\frac23=\prb(u=0)\prb(v=1)\).
\item Let \(x\sim N(0,1)\), let \(s\) be independent of \(x\) with \(\prb(s=1)=\prb(s=-1)=\frac12\), and
put \(y=sx\). Let \(\iota_+\) be \(1\) if \(s=1\) and \(0\) otherwise, and \(\iota_-=1-\iota_+\). Then
\(\expc(\iota_+)=\expc(\iota_-)=\frac12\) and \(\expc(s)=0\), and \(\iota_+\), \(\iota_-\) and
\(s\) are independent of every function of \(x\) (N3). As \(e^{ty}=\iota_+e^{tx}+\iota_-e^{-tx}\), (N7) and
Definition~\ref{r14-def-mvn} give \(M_y(t)=\frac12e^{t^2/2}+\frac12e^{t^2/2}=e^{t^2/2}\) for every \(t\): \(y\sim N(0,1)\).
Also \(\expc(x)=0\), so
\(\cov(x,y)=\expc(xy)=\expc(sx^2)=\expc(s)\expc(x^2)=0\).
\item But \(x\) and \(y\) are not independent. As \(x+y=(1+s)x\), \(e^{t(x+y)}=\iota_+e^{2tx}+\iota_-\), so
\(\expc(e^{t(x+y)})=\frac12e^{2t^2}+\frac12\). If \(x\) and \(y\) were independent, then by (N3) and (N7),
\(\expc(e^{t(x+y)})=\expc(e^{tx})\expc(e^{ty})=e^{t^2}\) for every \(t\); but at \(t=1\),
\(\frac12e^2+\frac12-e=\frac12(e-1)^2>0\). Hence, by Theorem~\ref{r14-thm-mvn-indep}(a), \(x\) and \(y\) are not jointly
normal either.
\end{itemize}
\end{remark}

\begin{example}\label{r14-exm-indep}
Let \(z\sim N_4(\mu,\Sigma)\) be as in Example~\ref{r14-exm-marginal}.
\begin{itemize}
\item By Corollary~\ref{r14-cor-marginal}(b),
\[
\begin{pmatrix}z_1\\z_2\\z_4\end{pmatrix}\sim
N_3\left(\begin{pmatrix}1\\0\\2\end{pmatrix},\begin{pmatrix}4&0&0\\0&1&2\\0&2&9\end{pmatrix}\right)
\]
and \(\sigma_{12}=\sigma_{14}=0\); so \(z_1\) and \((z_2,z_4)'\) are independent (Theorem~\ref{r14-thm-mvn-indep}(a) with
\(p=1\) and \(q=2\)).
\item Let \(A=\left(\begin{smallmatrix}1&0&0&0\\0&1&0&0\end{smallmatrix}\right)\), so that \(Az=(z_1,z_2)'\), and
\(e=a'z=z_3-z_1-2z_2\). As \(\Sigma a=(0,0,1,2)'\) (Example~\ref{r14-exm-cov}), \(A\Sigma a=(0,0)'\), and by
Theorem~\ref{r14-thm-mvn-indep}(c), \(e\) is independent of \((z_1,z_2)'\). Thus \(z_3=z_1+2z_2+e\), where \(e\sim N(1,1)\) is
independent of \(z_1\) and \(z_2\).
\item Likewise, let \(G=[I_3\ \ 0]\in\R^{3\times4}\), so that \(Gz=(z_1,z_2,z_3)'\), and \(f=b'z=z_4+2z_1+2z_2-2z_3\). As
\(\Sigma b=(0,0,0,1)'\) (Example~\ref{r14-exm-cov}), \(G\Sigma b=(0,0,0)'\), and \(f\) is independent of \((z_1,z_2,z_3)'\). Thus
\(z_4=-2z_1-2z_2+2z_3+f\), where \(f\sim N(0,1)\) (Example~\ref{r14-exm-marginal}) is independent of \(z_1\), \(z_2\) and
\(z_3\).
\item \(z_1+z_2\) and \(z_1-z_2\) are not independent: their covariance is \(3\ne0\) (Example~\ref{r14-exm-marginal}), and independent vectors are uncorrelated (N7).
\end{itemize}
\end{example}

\subsection*{The density of a normal vector}

A random vector \(z\in\R^n\) has the \textbf{density} \(f\colon\R^n\to[0,\infty)\) if
\(\prb(z\in A)=\int_Af(x)\,\mathrm{d}x\) for every set \(A\subseteq\R^n\); then \(\expc g(z)=\int_{\R^n}g(x)f(x)\,\mathrm{d}x\).
Every function \(f\geqslant0\) with \(\int_{\R^n}f(x)\,\mathrm{d}x=1\) is the density of some random vector. Random vectors
\(u_1,\dots,u_m\) with densities \(f_1,\dots,f_m\) are independent if and only if \((u_1',\dots,u_m')'\) has the density
\(f_1(u_1)\cdots f_m(u_m)\)~\citeat{casella-berger}{(4.6.2), (4.6.3), Theorem~1.6.5, Definition~4.6.5}. We also use the following facts on integrals without proof:
\begin{itemize}
\item (change of variables) if \(B\in\R^{n\times n}\) is invertible and \(c\in\R^n\), then
\(\int_{\R^n}g(x)\,\mathrm{d}x=\lvert\det B\rvert\int_{\R^n}g(c+Bu)\,\mathrm{d}u\);
\item (products) \(\int_{\R^n}g_1(u_1)\cdots g_n(u_n)\,\mathrm{d}u=\int_{\R}g_1(s)\,\mathrm{d}s\cdots\int_{\R}g_n(s)\,\mathrm{d}s\) for functions
\(g_1,\dots,g_n\geqslant0\);
\item \(\int_{\R}e^{-s^2/2}\,\mathrm{d}s=\sqrt{2\pi}\).
\end{itemize}

\needspace{14\baselineskip}
\begin{theorem}[\citeat{bapat}{\S8.1, (8.2)}, \citeat{tu-clm}{\S4.1}]\label{r14-thm-mvn-density}
Let \(\mu\in\R^n\), let \(\Sigma\in\R^{n\times n}\) be PD, and let
\[
f(x)=\frac1{(2\pi)^{n/2}(\det\Sigma)^{1/2}}\exp\Bigl(-\frac12(x-\mu)'\Sigma^{-1}(x-\mu)\Bigr)\qquad(x\in\R^n)
\]
\begin{enumerate}
\item \(f(x)>0\) for every \(x\), and \(\int_{\R^n}f(x)\,\mathrm{d}x=1\).
\item A random vector with the density \(f\) has the distribution \(N_n(\mu,\Sigma)\); in particular, \(N_n(\mu,\Sigma)\) exists.
\item Every \(z\sim N_n(\mu,\Sigma)\) has the density \(f\).
\end{enumerate}
For \(n=1\) and \(\sigma^2>0\), \(N(\mu,\sigma^2)\) has the density \((2\pi\sigma^2)^{-1/2}\exp\bigl(-(x-\mu)^2/(2\sigma^2)\bigr)\).
\end{theorem}
\begin{proof}
(a)
\begin{itemize}
\item \(\det\Sigma>0\) (Theorem~\ref{r8-thm-eigen-criterion}(c)), so \(f\) is defined, and \(f(x)>0\).
\item \(\Sigma^{1/2}\) is symmetric and PD, with \(\Sigma^{1/2}\Sigma^{1/2}=\Sigma\) (Theorem~\ref{r8-thm-sqrt}); so \(\det\Sigma^{1/2}>0\)
and \((\det\Sigma^{1/2})^2=\det\Sigma\) (Theorem~\ref{r5d-thm-product}(b)), that is, \(\det\Sigma^{1/2}=(\det\Sigma)^{1/2}\).
\item \(\Sigma^{-1}=(\Sigma^{1/2}\Sigma^{1/2})^{-1}=\Sigma^{-1/2}\Sigma^{-1/2}\) (Theorem~\ref{r2-thm-inverse-props}(b)), so
\(\Sigma^{1/2}\Sigma^{-1}\Sigma^{1/2}=I_n\), and for \(x=\mu+\Sigma^{1/2}u\),
\((x-\mu)'\Sigma^{-1}(x-\mu)=u'\Sigma^{1/2}\Sigma^{-1}\Sigma^{1/2}u=u'u=u_1^2+\dots+u_n^2\).
\item By the change of variables with \(c=\mu\) and \(B=\Sigma^{1/2}\), and then the rule for products,
\[
\int_{\R^n}f(x)\,\mathrm{d}x=\frac{(\det\Sigma)^{1/2}}{(2\pi)^{n/2}(\det\Sigma)^{1/2}}\int_{\R^n}e^{-u_1^2/2}\cdots e^{-u_n^2/2}\,\mathrm{d}u
=\frac{(\sqrt{2\pi})^n}{(2\pi)^{n/2}}=1
\]
\end{itemize}
(b)
\begin{itemize}
\item For \(t,x\in\R^n\), expanding with \(\Sigma'=\Sigma\) and \(\Sigma^{-1}\Sigma=I_n\),
\((x-\mu-\Sigma t)'\Sigma^{-1}(x-\mu-\Sigma t)=(x-\mu)'\Sigma^{-1}(x-\mu)-2t'(x-\mu)+t'\Sigma t\). Hence
\[
t'x-\tfrac12(x-\mu)'\Sigma^{-1}(x-\mu)=t'\mu+\tfrac12t'\Sigma t-\tfrac12(x-\mu-\Sigma t)'\Sigma^{-1}(x-\mu-\Sigma t)
\]
\item So \(e^{t'x}f(x)=\exp(t'\mu+\frac12t'\Sigma t)\,f_t(x)\), where \(f_t\) is the function \(f\) with \(\mu+\Sigma t\) in place of
\(\mu\). If \(y\) has the density \(f\), then by (a) for \(\mu+\Sigma t\),
\(M_y(t)=\int_{\R^n}e^{t'x}f(x)\,\mathrm{d}x=\exp(t'\mu+\frac12t'\Sigma t)\int_{\R^n}f_t(x)\,\mathrm{d}x=\exp(t'\mu+\frac12t'\Sigma t)\)
for every \(t\): \(y\sim N_n(\mu,\Sigma)\).
\item By (a), \(f\) is the density of some random vector, which has the distribution \(N_n(\mu,\Sigma)\).
\end{itemize}
(c) Let \(y\) be as in (b). Then \(M_z=M_y\), so \(z\) and \(y\) have the same distribution (N1), and
\(\prb(z\in A)=\prb(y\in A)=\int_Af(x)\,\mathrm{d}x\) for every set \(A\): \(z\) has the density \(f\).

For \(n=1\), \(\det(\sigma^2)=\sigma^2\) and the inverse of \((\sigma^2)\) is \((1/\sigma^2)\).
\end{proof}
The density \(f\) of Theorem~\ref{r14-thm-mvn-density} depends on \(x\) only through
\((x-\mu)'\Sigma^{-1}(x-\mu)=d_\Sigma(x,\mu)^2\), the squared Mahalanobis distance of Definition~\ref{r8-def-mahalanobis} with
\(\Sigma\) in place of \(S\). So \(f\) is constant on each \textbf{contour} \(\{x\in\R^n:d_\Sigma(x,\mu)=c\}\) (\(c>0\)). Let
\(\Sigma=Q\Lambda Q'\) as in Theorem~\ref{r8-thm-spectral}(b), with \(Q=[q_1\ \cdots\ q_n]\), and put \(u=Q'(x-\mu)\). Then
\(x-\mu=Qu=u_1q_1+\dots+u_nq_n\) (Proposition~\ref{r6-prop-orthonormal}(b)), every \(\lambda_j>0\)
(Theorem~\ref{r8-thm-eigen-criterion}(a)), and \(\Sigma^{-1}=Q\Lambda^{-1}Q'\) (Proposition~\ref{r8-prop-function}(a),(b) with the
function \(t\mapsto t\)), so
\[
d_\Sigma(x,\mu)^2=u'\Lambda^{-1}u=\frac{u_1^2}{\lambda_1}+\dots+\frac{u_n^2}{\lambda_n}
\]
In the coordinates \(u\), the contour is \(\bigl\{u:\sum_ju_j^2/(c\sqrt{\lambda_j})^2=1\bigr\}\). As \(x=\mu+u_1q_1+\dots+u_nq_n\),
the contour is an ellipse for \(n=2\) and an ellipsoid for \(n\geqslant3\), with center \(\mu\), axes along \(q_1,\dots,q_n\) and semi-axes
\(c\sqrt{\lambda_1},\dots,c\sqrt{\lambda_n}\)~\citeat{haerdle-simar}{\S4.4, (4.51), Example~4.14}.

\needspace{20\baselineskip}
\subsection*{Independent components and the existence of \(N_n(\mu,\Sigma)\)}

\begin{theorem}[\citeat{tu-clm}{Theorem~4.5}, \citeat{tu-asm}{Definition~P.113}, \citeat{bapat}{\S8.1}]\label{r14-thm-mvn-iid}
Let \(\mu\in\R^n\).
\begin{enumerate}
\item Let \(\sigma^2\geqslant0\). Then \(z\sim N_n(\mu,\sigma^2I_n)\) if and only if \(z_1,\dots,z_n\) are independent with
\(z_i\sim N(\mu_i,\sigma^2)\). More generally, \(z\sim N_n\bigl(\mu,\diag(\sigma_1^2,\dots,\sigma_n^2)\bigr)\) if and only if
\(z_1,\dots,z_n\) are independent with \(z_i\sim N(\mu_i,\sigma_i^2)\). Let \(n=n_1+\dots+n_m\) and
\(\mu=(\mu_1',\dots,\mu_m')'\) with \(\mu_k\in\R^{n_k}\). If random vectors \(x_1\in\R^{n_1},\dots,x_m\in\R^{n_m}\) are
independent, with \(x_k\sim N_{n_k}(\mu_k,\Sigma_k)\), then \((x_1',\dots,x_m')'\sim N_n\bigl(\mu,\diag(\Sigma_1,\dots,\Sigma_m)\bigr)\).
\item Let \(\Sigma\in\R^{n\times n}\) be PSD, and let \(w\sim N_n(0,I_n)\), which exists by
Theorem~\ref{r14-thm-mvn-density}(b). Then \(\mu+\Sigma^{1/2}w\sim N_n(\mu,\Sigma)\). So \(N_n(\mu,\Sigma)\) exists for every PSD
\(\Sigma\).
\item A random vector \(z\) has the distribution \(N_n(\mu,\Sigma)\) if and only if it has the distribution of \(\mu+Aw\) for some
\(A\in\R^{n\times k}\) with \(AA'=\Sigma\) and some \(w\sim N_k(0,I_k)\), that is (by (a)), with \(w_1,\dots,w_k\) independent
and \(N(0,1)\).
\end{enumerate}
\end{theorem}
\begin{proof}
(a)
\begin{itemize}
\item Let \(z\sim N_n(\mu,D)\) with \(D=\diag(\sigma_1^2,\dots,\sigma_n^2)\). The off-diagonal entries of \(D\) are \(0\), so
\(z_1,\dots,z_n\) are independent by Theorem~\ref{r14-thm-mvn-indep}(b) (subvectors of length \(1\)), and
\(z_i\sim N(\mu_i,\sigma_i^2)\) by Corollary~\ref{r14-cor-marginal}(b).
\item Conversely, let \(x_1,\dots,x_m\) be independent with \(x_k\sim N_{n_k}(\mu_k,\Sigma_k)\), and let
\(x=(x_1',\dots,x_m')'\) and \(D=\diag(\Sigma_1,\dots,\Sigma_m)\). For \(t=(t_1',\dots,t_m')'\) with \(t_k\in\R^{n_k}\), the
block product (Theorem~\ref{r5b-thm-block}(c)) gives \(t'\mu=\sum_kt_k'\mu_k\) and \(t'Dt=\sum_kt_k'\Sigma_kt_k\geqslant0\); so
\(D\), which is symmetric (Theorem~\ref{r5b-thm-block}(b)), is PSD. By (N2),
\(M_x(t)=\prod_kM_{x_k}(t_k)=\prod_k\exp(t_k'\mu_k+\frac12t_k'\Sigma_kt_k)=\exp(t'\mu+\frac12t'Dt)\) for \(t\) near \(0\).
\item With \(w\sim N_n(0,I_n)\) (Theorem~\ref{r14-thm-mvn-density}(b)) and the square root \(D^{1/2}\) of
Theorem~\ref{r8-thm-sqrt}, the vector \(y=\mu+D^{1/2}w\) is \(N_n(\mu,D^{1/2}I_nD^{1/2})=N_n(\mu,D)\)
(Theorem~\ref{r14-thm-mvn-linear}), so it has the same MGF as \(x\) near \(0\). Hence \(x\) and \(y\) have the same
distribution (N1) and the same MGF \(\exp(t'\mu+\frac12t'Dt)\) for every \(t\): \(x\sim N_n(\mu,D)\)
(Definition~\ref{r14-def-mvn}).
\item With \(m=n\), \(n_k=1\), \(x_k=z_k\) and \(\Sigma_k=\sigma_k^2\), this is the converse in the second statement; and
\(\sigma_1^2=\dots=\sigma_n^2=\sigma^2\) gives the first statement.
\end{itemize}
(b) \(\Sigma^{1/2}\) is symmetric and \(\Sigma^{1/2}\Sigma^{1/2}=\Sigma\) (Theorem~\ref{r8-thm-sqrt}). By
Theorem~\ref{r14-thm-mvn-linear} with \(A=\Sigma^{1/2}\) and \(b=\mu\),
\(\mu+\Sigma^{1/2}w\sim N_n(\mu,\Sigma^{1/2}I_n\Sigma^{1/2})=N_n(\mu,\Sigma)\).

(c)
\begin{itemize}
\item If \(AA'=\Sigma\) and \(w\sim N_k(0,I_k)\), then \(\mu+Aw\sim N_n(\mu,AI_kA')=N_n(\mu,\Sigma)\)
(Theorem~\ref{r14-thm-mvn-linear}). A random vector with the same distribution has the same MGF, so it is \(N_n(\mu,\Sigma)\).
\item Conversely, let \(z\sim N_n(\mu,\Sigma)\). With \(k=n\), \(A=\Sigma^{1/2}\) and \(w\) as in (b), \(\mu+Aw\sim N_n(\mu,\Sigma)\) has the
same MGF as \(z\), hence the same distribution (N1).\qedhere
\end{itemize}
\end{proof}
The multivariate normal distribution is often defined as the distribution of \(\mu+Aw\) with \(w_1,\dots,w_k\) independent
\(N(0,1)\)~\citeat{bapat}{\S8.1}, \citeat{tu-asm}{Definition~P.113}; by (c) this is the same as Definition~\ref{r14-def-mvn}. The matrix \(A\) is not unique: besides \(A=\Sigma^{1/2}\),
for PD \(\Sigma\) one may take \(A=T'\), where \(\Sigma=T'T\) is the Cholesky decomposition of
Theorem~\ref{r8-thm-cholesky}; \(T'\) is lower triangular with positive diagonal entries. For instance, for \(\Sigma\) of
Example~\ref{r14-exm-cov},
\[
\Sigma=AA',\qquad A=\begin{pmatrix}2&0&0&0\\0&1&0&0\\2&2&1&0\\0&2&2&1\end{pmatrix}
\]
where \(A'\) is upper triangular with positive diagonal entries; \(A'\) is invertible (Theorem~\ref{r2-thm-triangular}(a)),
so \(\Sigma=(A')'A'\) is PD (Theorems~\ref{r5d-thm-product}(a) and~\ref{r8-thm-XtX}), and \(A'\) is the Cholesky factor of \(\Sigma\)
(Theorem~\ref{r8-thm-cholesky}), and \(z_1=1+2w_1\), \(z_2=w_2\), \(z_3=2+2w_1+2w_2+w_3\), \(z_4=2+2w_2+2w_3+w_4\)
defines a vector \(z\sim N_4(\mu,\Sigma)\) as in Example~\ref{r14-exm-marginal}.

\begin{example}[\citeat{casella-berger}{Theorem~5.3.1, Example~5.2.8, Lemma~5.3.3}]\label{r14-exm-sample}
Let \(\mu\in\R\), \(\sigma^2>0\) and \(z\sim N_n(\mu\one_n,\sigma^2I_n)\) (\(n\geqslant2\)). By Theorem~\ref{r14-thm-mvn-iid}(a),
\(z_1,\dots,z_n\) are independent, each \(N(\mu,\sigma^2)\): a \textbf{random sample} from \(N(\mu,\sigma^2)\).
\begin{enumerate}
\item The mean \(\bar z=\frac1n\one'z\) has the distribution \(N(\mu,\sigma^2/n)\), by Corollary~\ref{r14-cor-marginal}(c) with
\(a=\frac1n\one\): \(a'\mu\one=\mu\) and \(a'(\sigma^2I_n)a=\sigma^2/n\).
\item The deviations \(C_nz=(z_1-\bar z,\dots,z_n-\bar z)'\) (Example~\ref{r1-exm-centering}) are independent of \(\bar z\): with
\(A=\frac1n\one'\) and \(B=C_n\), \(A(\sigma^2I_n)B'=\frac{\sigma^2}n(C_n\one)'=0\) (Proposition~\ref{r1-prop-centering}), and
Theorem~\ref{r14-thm-mvn-indep}(c) applies. Hence, by (N3), \(\bar z\) is independent of every function of the deviations,
such as the sample variance \(s^2=\frac1{n-1}(C_nz)'(C_nz)\).
\item By Theorem~\ref{r14-thm-mvn-linear}, \(C_nz\sim N_n(0,\sigma^2C_n)\), as \(C_n\mu\one=0\) and
\(C_n(\sigma^2I_n)C_n'=\sigma^2C_n^2=\sigma^2C_n\). This covariance matrix is singular, of rank \(n-1\)
(Corollary~\ref{r6-cor-trace-hat}): the deviations always add up to \(0\).
\end{enumerate}
\end{example}

\needspace{20\baselineskip}
\subsection*{Standardization and singular covariance matrices}

\begin{corollary}[\citeat{bapat}{\S8.1}, \citeat{tu-clm}{\S4.1}]\label{r14-cor-standard}
Let \(z\sim N_n(\mu,\Sigma)\).
\begin{enumerate}
\item If \(\Sigma\) is PD, then \(w=\Sigma^{-1/2}(z-\mu)\sim N_n(0,I_n)\): the \(w_i\) are independent and \(N(0,1)\), and
\(z=\mu+\Sigma^{1/2}w\). Moreover \((z-\mu)'\Sigma^{-1}(z-\mu)=w'w=w_1^2+\dots+w_n^2\) has the distribution \(\chi^2_n\).
\item If \(\Sigma\ne0\), let the columns of \(K\in\R^{n\times r}\) form a basis of \(\mathcal C(\Sigma)\), so that \(r=\rank\Sigma\). Then
\(\prb\bigl(z-\mu\in\mathcal C(\Sigma)\bigr)=1\): with probability \(1\), \(z\) lies in the set
\(\mu+\mathcal C(\Sigma)=\{\mu+Kc:c\in\R^r\}\). If \(\Sigma=0\), then \(\prb(z=\mu)=1\).
\end{enumerate}
\end{corollary}
\begin{proof}
(a)
\begin{itemize}
\item \(\Sigma^{-1/2}\) is symmetric and \(\Sigma^{-1/2}\Sigma\Sigma^{-1/2}=I_n\) (Theorem~\ref{r8-thm-sqrt}). By
Theorem~\ref{r14-thm-mvn-linear} with \(A=\Sigma^{-1/2}\) and \(b=-\Sigma^{-1/2}\mu\), \(w\sim N_n(0,I_n)\), and by
Theorem~\ref{r14-thm-mvn-iid}(a) the \(w_i\) are independent and \(N(0,1)\).
\item \(\Sigma^{1/2}w=\Sigma^{1/2}\Sigma^{-1/2}(z-\mu)=z-\mu\), and
\(w'w=(z-\mu)'\Sigma^{-1/2}\Sigma^{-1/2}(z-\mu)=(z-\mu)'\Sigma^{-1}(z-\mu)\), as in the proof of
Theorem~\ref{r14-thm-mvn-density}(a). By (N5), \(w_1^2+\dots+w_n^2\sim\chi^2_n\).
\end{itemize}
(b)
\begin{itemize}
\item Let \(P=K(K'K)^{-1}K'\), the matrix \(P\) of Theorem~\ref{r6-thm-hat} with \(A=K\), which has full column rank
\(r\). Every column of \(\Sigma\) lies in \(\mathcal C(K)\), so \(\Sigma=KT\) for some \(T\in\R^{r\times n}\)
(Theorem~\ref{r1-thm-views}(a),(b)), and \((I_n-P)\Sigma=(I_n-P)KT=0\) (Theorem~\ref{r6-thm-hat}(c)).
\item By Theorem~\ref{r14-thm-mvn-linear}, \((I_n-P)(z-\mu)\sim N_n\bigl(0,(I_n-P)\Sigma(I_n-P)'\bigr)=N_n(0,0)\), the
distribution of the constant \(0\) (Definition~\ref{r14-def-mvn}); so \((I_n-P)(z-\mu)=0\) with probability \(1\).
\item If \((I_n-P)(z-\mu)=0\), then \(z-\mu=P(z-\mu)=K\bigl((K'K)^{-1}K'(z-\mu)\bigr)\in\mathcal C(K)=\mathcal C(\Sigma)\).
\item The case \(\Sigma=0\) is part of Definition~\ref{r14-def-mvn}.\qedhere
\end{itemize}
\end{proof}
So a normal vector with a singular covariance matrix lies, with probability \(1\), in the affine subspace
\(\mu+\mathcal C(\Sigma)\) of dimension \(\rank\Sigma<n\). Such vectors, like the deviations \(C_nz\) of
Example~\ref{r14-exm-sample}(c), for which \(\mathcal C(\sigma^2C_n)=\Span\{\one\}^\perp\) (Example~\ref{r6-exm-centering}), are
the reason why Definition~\ref{r14-def-mvn} uses the MGF and not a density~\citeat{tu-clm}{\S4.1}.

\subsection*{Conditional distributions}

\begin{definition}[\citeat{casella-berger}{Definition~4.2.3, (4.6.4), (4.6.6)}]\label{r14-def-conditional}
Let \((u',v')'\), with \(u\in\R^p\) and \(v\in\R^q\), have the density \(f\). Then \(u\) has the density
\(f_u(a)=\int_{\R^q}f(a,b)\,\mathrm{d}b\) (\(a\in\R^p\)), the \textbf{marginal density}. For \(a\in\R^p\) with \(f_u(a)>0\), the function
\[
f(b\mid a)=\frac{f(a,b)}{f_u(a)}\qquad(b\in\R^q)
\]
is a density on \(\R^q\) (it is \(\geqslant0\), with integral \(f_u(a)/f_u(a)=1\)), the \textbf{conditional density} of \(v\) given
\(u=a\); the distribution with this density is the \textbf{conditional distribution} of \(v\) given \(u=a\).
\end{definition}

\needspace{18\baselineskip}
\begin{theorem}[\citeat{bapat}{\S8.1; Ch.~8, Exercise~13}, \citeat{searle-khuri}{\S10.4}, \citeat{casella-berger}{Example~4.5.8}]\label{r14-thm-mvn-conditional}
Let \(z\sim N_n(\mu,\Sigma)\) be partitioned as in Corollary~\ref{r14-cor-marginal}(a), with \(z_1\in\R^p\) and \(z_2\in\R^q\),
and let \(\Sigma_{11}\) be PD. Put
\[
\Sigma_{22\cdot1}=\Sigma_{22}-\Sigma_{21}\Sigma_{11}^{-1}\Sigma_{12}
\]
the Schur complement of \(\Sigma_{11}\) in \(\Sigma\) (Theorem~\ref{r5b-thm-schur}). Then the \textbf{conditional distribution of
\(z_2\) given \(z_1\)} is
\[
N_q\bigl(\mu_2+\Sigma_{21}\Sigma_{11}^{-1}(z_1-\mu_1),\ \Sigma_{22\cdot1}\bigr)
\]
in the following sense.
\begin{enumerate}
\item \(\Sigma_{22\cdot1}\) is PSD, and \(w=z_2-\mu_2-\Sigma_{21}\Sigma_{11}^{-1}(z_1-\mu_1)\) is independent of \(z_1\), with
\(w\sim N_q(0,\Sigma_{22\cdot1})\). So \(z_2\) is the sum of \(\mu_2+\Sigma_{21}\Sigma_{11}^{-1}(z_1-\mu_1)\), which is known once
\(z_1\) is known, and of \(w\), whose distribution does not depend on \(z_1\).
\item If \(\Sigma\) is PD, then \(\Sigma_{22\cdot1}\) is PD and is the inverse of the lower right \(q\times q\) block of \(\Sigma^{-1}\);
\(z_1\) has the density of \(N_p(\mu_1,\Sigma_{11})\); and for every \(a\in\R^p\), the conditional density of \(z_2\) given
\(z_1=a\) (Definition~\ref{r14-def-conditional}) is the density of
\(N_q\bigl(\mu_2+\Sigma_{21}\Sigma_{11}^{-1}(a-\mu_1),\Sigma_{22\cdot1}\bigr)\).
\end{enumerate}
\end{theorem}
\begin{proof}
Let
\[
L=\begin{pmatrix}I_p&0\\-\Sigma_{21}\Sigma_{11}^{-1}&I_q\end{pmatrix},\qquad\text{so that}\qquad
L(z-\mu)=\begin{pmatrix}z_1-\mu_1\\w\end{pmatrix}
\]
by the block product. \(\Sigma_{11}^{-1}\) is symmetric (Theorem~\ref{r2-thm-inverse-props}(c)) and \(\Sigma_{21}'=\Sigma_{12}\), so
\(L'=\left(\begin{smallmatrix}I_p&-\Sigma_{11}^{-1}\Sigma_{12}\\0&I_q\end{smallmatrix}\right)\) (Theorem~\ref{r5b-thm-block}(b)), and
Theorem~\ref{r5b-thm-schur}, with \(A=\Sigma_{11}\), \(B=\Sigma_{12}\), \(C=\Sigma_{21}\) and \(D=\Sigma_{22}\), gives
\[
L\Sigma L'=\begin{pmatrix}\Sigma_{11}&0\\0&\Sigma_{22\cdot1}\end{pmatrix}
\]
(a)
\begin{itemize}
\item By Theorem~\ref{r14-thm-mvn-linear} with \(A=L\) and \(b=-L\mu\), \(L(z-\mu)\sim N_n(0,L\Sigma L')\). By
Corollary~\ref{r14-cor-marginal}(a), \(w\sim N_q(0,\Sigma_{22\cdot1})\); so \(\Sigma_{22\cdot1}=\var(w)\) is PSD
(Theorem~\ref{r11-thm-linear}(c)).
\item The off-diagonal blocks of \(L\Sigma L'\) are \(0\), so \(z_1-\mu_1\) and \(w\) are independent
(Theorem~\ref{r14-thm-mvn-indep}(a)), and so are \(z_1=(z_1-\mu_1)+\mu_1\) and \(w\) (N3).
\end{itemize}
(b)
\begin{itemize}
\item \(\Sigma_{11}\) and \(\Sigma_{22\cdot1}\) are PD (Theorem~\ref{r8-thm-schur-pd}). \(L\) is invertible
(Proposition~\ref{r5b-prop-block-elementary}), so by Theorem~\ref{r2-thm-inverse-props}(a)--(c) and
Proposition~\ref{r5b-prop-block-tri}(b),
\[
\Sigma=L^{-1}\begin{pmatrix}\Sigma_{11}&0\\0&\Sigma_{22\cdot1}\end{pmatrix}(L')^{-1},\qquad
\Sigma^{-1}=L'\begin{pmatrix}\Sigma_{11}^{-1}&0\\0&\Sigma_{22\cdot1}^{-1}\end{pmatrix}L
\]
By Theorem~\ref{r5b-thm-partitioned-inverse}(a), the lower right block of \(\Sigma^{-1}\) is \(\Sigma_{22\cdot1}^{-1}\).
\item For \(x=(a',b')'\) with \(a\in\R^p\) and \(b\in\R^q\), put \(m(a)=\mu_2+\Sigma_{21}\Sigma_{11}^{-1}(a-\mu_1)\); then
\(L(x-\mu)=\bigl((a-\mu_1)',(b-m(a))'\bigr)'\), and
\[
\begin{aligned}
(x-\mu)'\Sigma^{-1}(x-\mu)&=\bigl(L(x-\mu)\bigr)'\begin{pmatrix}\Sigma_{11}^{-1}&0\\0&\Sigma_{22\cdot1}^{-1}\end{pmatrix}L(x-\mu)\\
&=(a-\mu_1)'\Sigma_{11}^{-1}(a-\mu_1)+(b-m(a))'\Sigma_{22\cdot1}^{-1}(b-m(a))
\end{aligned}
\]
\item \(\det\Sigma=\det\Sigma_{11}\cdot\det\Sigma_{22\cdot1}\) (Theorem~\ref{r5d-thm-schur-det}(a)) and
\((2\pi)^{n/2}=(2\pi)^{p/2}(2\pi)^{q/2}\). So the density \(f\) of \(z\) (Theorem~\ref{r14-thm-mvn-density}(c)) factors as
\(f(a,b)=f_1(a)\,g_a(b)\), where \(f_1\) is the density of \(N_p(\mu_1,\Sigma_{11})\) and \(g_a\) that of \(N_q(m(a),\Sigma_{22\cdot1})\).
\item By Theorem~\ref{r14-thm-mvn-density}(a), \(\int_{\R^q}g_a(b)\,\mathrm{d}b=1\); so the marginal density of \(z_1\) is
\(\int_{\R^q}f(a,b)\,\mathrm{d}b=f_1(a)>0\), and the conditional density of \(z_2\) given \(z_1=a\) is \(f(a,b)/f_1(a)=g_a(b)\).\qedhere
\end{itemize}
\end{proof}
The conditional mean \(\mu_2+\Sigma_{21}\Sigma_{11}^{-1}(z_1-\mu_1)\) is an affine function of \(z_1\), and the conditional covariance
matrix \(\Sigma_{22\cdot1}\) does not depend on \(z_1\). It is never larger than \(\Sigma_{22}=\var(z_2)\): the matrix
\(\Sigma_{11}^{-1}=(\Sigma_{11}^{-1/2})'\Sigma_{11}^{-1/2}\) is PSD (Theorems~\ref{r8-thm-sqrt} and~\ref{r8-thm-XtX}), so
\(\Sigma_{22}-\Sigma_{22\cdot1}=\Sigma_{12}'\Sigma_{11}^{-1}\Sigma_{12}\) is PSD (Theorem~\ref{r8-thm-congruence}(a)), and for every
\(c\in\R^q\), \(\var(c'w)=c'\Sigma_{22\cdot1}c\leqslant c'\Sigma_{22}c=\var(c'z_2)\): knowing \(z_1\) can only
reduce the uncertainty about \(z_2\). If the diagonal entries of \(\Sigma_{22\cdot1}=\var(w)\) are
positive, the correlations of the entries of \(w\) (Definition~\ref{r14-def-corr}), which are computed from the
conditional covariance matrix \(\Sigma_{22\cdot1}\), are the \textbf{partial correlations} of the entries of
\(z_2\) given \(z_1\)~\citeat{haerdle-simar}{\S5.1, Example~5.3}.

\needspace{10\baselineskip}
\begin{example}\label{r14-exm-conditional}
Let \(z\sim N_4(\mu,\Sigma)\) be as in Example~\ref{r14-exm-marginal}, with \(\mu=(1,0,2,2)'\). For the lower triangular \(A\)
with \(\Sigma=AA'\) of the paragraph after Theorem~\ref{r14-thm-mvn-iid},
\[
A^{-1}=\begin{pmatrix}\frac12&0&0&0\\0&1&0&0\\-1&-2&1&0\\2&2&-2&1\end{pmatrix},\qquad
\Sigma^{-1}=(A^{-1})'A^{-1}=\begin{pmatrix}\frac{21}4&6&-5&2\\6&9&-6&2\\-5&-6&5&-2\\2&2&-2&1\end{pmatrix}
\]
as \(AA^{-1}=I_4\) (Theorem~\ref{r2-thm-invertible}) and \(\Sigma^{-1}=(A')^{-1}A^{-1}\) (Theorem~\ref{r2-thm-inverse-props}(b),(c)); the rows \(3\) and \(4\) of
\(A^{-1}\) are \(a'\) and \(b'\) of Example~\ref{r14-exm-cov}. As \(\Sigma\) is PD (the paragraph after
Theorem~\ref{r14-thm-mvn-iid}), \(\Sigma_{11}\) is PD in (a) and (b) (Theorem~\ref{r8-thm-principal-sub}).
\begin{enumerate}
\item Given \((z_1,z_2,z_3)'\): here \(p=3\), \(q=1\), \(\Sigma_{11}\) is the leading \(3\times3\) block of \(\Sigma\),
\(\Sigma_{21}=(0,2,6)\) and \(\Sigma_{22}=9\). As
\[
(-2,-2,2)\,\Sigma_{11}=(-2,-2,2)\begin{pmatrix}4&0&4\\0&1&2\\4&2&9\end{pmatrix}=(-8+0+8,\ 0-2+4,\ -8-4+18)=(0,2,6)
\]
multiplying by \(\Sigma_{11}^{-1}\) on the right gives \(\Sigma_{21}\Sigma_{11}^{-1}=(-2,-2,2)\), and
\(\Sigma_{22\cdot1}=9-(-2,-2,2)(0,2,6)'=9-8=1\). Given \(z_1=a_1\), \(z_2=a_2\) and \(z_3=a_3\), the conditional distribution of \(z_4\) is
\[
N\bigl(2-2(a_1-1)-2a_2+2(a_3-2),\ 1\bigr)=N(2a_3-2a_1-2a_2,\ 1)
\]
for instance \(N(4,1)\) for \(a_1=3\),
\(a_2=1\), \(a_3=6\). This agrees with Example~\ref{r14-exm-indep}: \(z_4=-2z_1-2z_2+2z_3+f\), with \(f\sim N(0,1)\)
independent of \((z_1,z_2,z_3)'\). The \((4,4)\) entry of \(\Sigma^{-1}\) is \(1=\Sigma_{22\cdot1}^{-1}\), as
Theorem~\ref{r14-thm-mvn-conditional}(b) says.
\item Given \((z_1,z_2)'\): here \(p=q=2\), \(\Sigma_{11}=\diag(4,1)\), \(\Sigma_{11}^{-1}=\diag(\frac14,1)\) and
\(\Sigma_{21}=\left(\begin{smallmatrix}4&2\\0&2\end{smallmatrix}\right)\), so
\(\Sigma_{21}\Sigma_{11}^{-1}=\left(\begin{smallmatrix}1&2\\0&2\end{smallmatrix}\right)\) and
\[
\Sigma_{22\cdot1}=\begin{pmatrix}9&6\\6&9\end{pmatrix}-\begin{pmatrix}1&2\\0&2\end{pmatrix}\begin{pmatrix}4&0\\2&2\end{pmatrix}
=\begin{pmatrix}9&6\\6&9\end{pmatrix}-\begin{pmatrix}8&4\\4&4\end{pmatrix}=\begin{pmatrix}1&2\\2&5\end{pmatrix}
\]
As \((2,2)'+\left(\begin{smallmatrix}1&2\\0&2\end{smallmatrix}\right)(a_1-1,\ a_2)'=(a_1+2a_2+1,\ 2a_2+2)'\), given \(z_1=a_1\)
and \(z_2=a_2\) the vector \((z_3,z_4)'\) has the distribution \(N_2\bigl((a_1+2a_2+1,\ 2a_2+2)',\Sigma_{22\cdot1}\bigr)\). Its
first entry agrees with Example~\ref{r14-exm-indep}: \(z_3=z_1+2z_2+e\), with \(e\sim N(1,1)\) independent of
\((z_1,z_2)'\). The lower right \(2\times2\) block of \(\Sigma^{-1}\) is
\(\left(\begin{smallmatrix}5&-2\\-2&1\end{smallmatrix}\right)=\Sigma_{22\cdot1}^{-1}\), as
Theorem~\ref{r14-thm-mvn-conditional}(b) says. Knowing \(z_1\) and \(z_2\) reduces the variances of \(z_3\) and \(z_4\) from
\(9\) and \(9\) to \(1\) and \(5\).
\end{enumerate}
\end{example}

\subsection*{The bivariate normal distribution}

\begin{definition}[\citeat{casella-berger}{Definition~4.5.10}]\label{r14-def-bivariate}
Let \(\mu_u,\mu_v\in\R\), \(\sigma_u,\sigma_v>0\) and \(-1\leqslant\rho\leqslant1\). A random vector \((u,v)'\) has the
\textbf{bivariate normal distribution} with means \(\mu_u,\mu_v\), variances \(\sigma_u^2,\sigma_v^2\) and correlation \(\rho\) if
\[
\begin{pmatrix}u\\v\end{pmatrix}\sim N_2\left(\begin{pmatrix}\mu_u\\\mu_v\end{pmatrix},\
\Sigma\right),\qquad
\Sigma=\begin{pmatrix}\sigma_u^2&\rho\sigma_u\sigma_v\\\rho\sigma_u\sigma_v&\sigma_v^2\end{pmatrix}
\]
\end{definition}
This \(\Sigma\) is PSD: for \(\lvert\rho\rvert<1\) its leading principal minors \(\sigma_u^2\) and \(\sigma_u^2\sigma_v^2(1-\rho^2)\) are
positive, so it is PD (Theorem~\ref{r8-thm-leading-minors}); for \(\rho=\pm1\), \(\Sigma=cc'\) with \(c=(\sigma_u,\rho\sigma_v)'\), which is
PSD (Theorem~\ref{r8-thm-XtX}). Conversely, by Definition~\ref{r14-def-corr} and Proposition~\ref{r14-prop-corr}(a), every
\(N_2\) distribution whose covariance matrix has a positive diagonal is of this form, with \(\rho=\rho_{uv}\).

\needspace{18\baselineskip}
\begin{proposition}[\citeat{casella-berger}{\S4.5}]\label{r14-prop-bivariate}
Let \((u,v)'\) have the bivariate normal distribution of Definition~\ref{r14-def-bivariate}.
\begin{enumerate}
\item \(u\sim N(\mu_u,\sigma_u^2)\), \(v\sim N(\mu_v,\sigma_v^2)\), \(\rho_{uv}=\rho\), and for constants \(\alpha\), \(\beta\),
\(\alpha u+\beta v\sim N\bigl(\alpha\mu_u+\beta\mu_v,\ \alpha^2\sigma_u^2+2\alpha\beta\rho\sigma_u\sigma_v+\beta^2\sigma_v^2\bigr)\).
\item \(u\) and \(v\) are independent if and only if \(\rho=0\).
\item If \(\lvert\rho\rvert<1\), then \((u,v)'\) has the density
\[
f(s,t)=\frac1{2\pi\sigma_u\sigma_v\sqrt{1-\rho^2}}
\exp\biggl(-\frac1{2(1-\rho^2)}\Bigl(\Bigl(\frac{s-\mu_u}{\sigma_u}\Bigr)^2-2\rho\,\frac{s-\mu_u}{\sigma_u}\cdot\frac{t-\mu_v}{\sigma_v}
+\Bigl(\frac{t-\mu_v}{\sigma_v}\Bigr)^2\Bigr)\biggr)
\]
\item The conditional distribution of \(v\) given \(u\) is
\(N\bigl(\mu_v+\rho\frac{\sigma_v}{\sigma_u}(u-\mu_u),\ \sigma_v^2(1-\rho^2)\bigr)\). If \(\rho=\pm1\), then
\(v=\mu_v+\rho\frac{\sigma_v}{\sigma_u}(u-\mu_u)\) with probability \(1\).
\end{enumerate}
\end{proposition}
\begin{proof}
(a) The marginal distributions are given by Corollary~\ref{r14-cor-marginal}(b), and the distribution of \(\alpha u+\beta v\) by
Corollary~\ref{r14-cor-marginal}(c) with the vector \((\alpha,\beta)'\), as
\((\alpha,\beta)\Sigma(\alpha,\beta)'=\alpha^2\sigma_u^2+2\alpha\beta\rho\sigma_u\sigma_v+\beta^2\sigma_v^2\). By
Definition~\ref{r14-def-mvn}, \(\cov(u,v)=\rho\sigma_u\sigma_v\), so \(\rho_{uv}=\rho\sigma_u\sigma_v/(\sigma_u\sigma_v)=\rho\)
(Definition~\ref{r14-def-corr}).

(b) By Theorem~\ref{r14-thm-mvn-indep}(a) with \(p=q=1\), \(u\) and \(v\) are independent if and only if
\(\rho\sigma_u\sigma_v=0\), that is, \(\rho=0\).

\needspace{8\baselineskip}
(c)
\begin{itemize}
\item \(\det\Sigma=\sigma_u^2\sigma_v^2-\rho^2\sigma_u^2\sigma_v^2=\sigma_u^2\sigma_v^2(1-\rho^2)>0\), and by
Theorem~\ref{r5d-thm-adjugate},
\(\Sigma^{-1}=\frac1{\sigma_u^2\sigma_v^2(1-\rho^2)}\left(\begin{smallmatrix}\sigma_v^2&-\rho\sigma_u\sigma_v\\-\rho\sigma_u\sigma_v&\sigma_u^2\end{smallmatrix}\right)\).
\item Hence, for \(x=(s,t)'\),
\[
(x-\mu)'\Sigma^{-1}(x-\mu)=\frac{\sigma_v^2(s-\mu_u)^2-2\rho\sigma_u\sigma_v(s-\mu_u)(t-\mu_v)+\sigma_u^2(t-\mu_v)^2}{\sigma_u^2\sigma_v^2(1-\rho^2)}
\]
which is the bracket in the exponent of \(f\) divided by \(1-\rho^2\); and \((2\pi)^{2/2}(\det\Sigma)^{1/2}=2\pi\sigma_u\sigma_v\sqrt{1-\rho^2}\).
Theorem~\ref{r14-thm-mvn-density}(c) gives the claim.
\end{itemize}
(d) Theorem~\ref{r14-thm-mvn-conditional} with \(p=q=1\) and \(\Sigma_{11}=\sigma_u^2>0\): here
\(\Sigma_{21}\Sigma_{11}^{-1}=\rho\sigma_u\sigma_v/\sigma_u^2=\rho\sigma_v/\sigma_u\) and
\(\Sigma_{22\cdot1}=\sigma_v^2-\rho^2\sigma_u^2\sigma_v^2/\sigma_u^2=\sigma_v^2(1-\rho^2)\). If \(\rho=\pm1\), then
\(w=v-\mu_v-\rho\frac{\sigma_v}{\sigma_u}(u-\mu_u)\sim N(0,0)\) in (a) of that theorem, so \(\prb(w=0)=1\)
(Definition~\ref{r14-def-mvn}).
\end{proof}

\begin{example}\label{r14-exm-bivariate}
Let \((u,v)'\) be bivariate normal with \(\mu_u=5\), \(\mu_v=6\), \(\sigma_u=2\), \(\sigma_v=4\) and \(\rho=\frac12\). Then
\(\rho\sigma_u\sigma_v=4\), and
\[
\Sigma=\begin{pmatrix}4&4\\4&16\end{pmatrix},\qquad\det\Sigma=48,\qquad
\Sigma^{-1}=\frac1{48}\begin{pmatrix}16&-4\\-4&4\end{pmatrix}=\frac1{12}\begin{pmatrix}4&-1\\-1&1\end{pmatrix}
\]
\begin{itemize}
\item By Proposition~\ref{r14-prop-bivariate}(c), as \(2\pi\cdot2\cdot4\sqrt{1-\frac14}=8\sqrt3\,\pi\), the density is
\[
f(s,t)=\frac1{8\sqrt3\,\pi}\exp\Bigl(-\frac1{24}\bigl(4(s-5)^2-2(s-5)(t-6)+(t-6)^2\bigr)\Bigr)
\]
\item By Proposition~\ref{r14-prop-bivariate}(d), given \(u\) the variable \(v\) has the distribution
\(N\bigl(6+\frac12\cdot\frac42(u-5),\ 16\cdot\frac34\bigr)=N(u+1,\ 12)\); for instance \(N(10,12)\) given \(u=9\), against
\(v\sim N(6,16)\) without knowledge of \(u\). With the roles exchanged, given \(v\) the variable \(u\) has the distribution
\(N\bigl(5+\frac12\cdot\frac24(v-6),\ 4\cdot\frac34\bigr)=N\bigl(5+\frac14(v-6),\ 3\bigr)\).
\item \(v-u\) is independent of \(u\): \((1,0)\Sigma(-1,1)'=-4+4=0\) (Theorem~\ref{r14-thm-mvn-indep}(c)), and \(v-u\sim N(1,\ 4-8+16)=N(1,12)\)
by (a). So \(v=u+(v-u)\) is \(u\) plus an independent \(N(1,12)\) variable, in accordance with the conditional distribution
found above.
\end{itemize}
\end{example}

\subsection*{Maximum likelihood estimates of \(\mu\) and \(\Sigma\)}

\begin{lemma}[\citeat{abadir-magnus}{Exercise~12.20}]\label{r14-lem-logdet}
Let \(A\in\R^{p\times p}\) be PD. Then \(\log\det A\leqslant\tr A-p\), with equality only if \(A=I_p\).
\end{lemma}
\begin{proof}
\begin{itemize}
\item Let \(\varphi(\lambda)=\lambda-1-\log\lambda\) for \(\lambda>0\). Its derivative \(\varphi'(\lambda)=1-1/\lambda\) is \(<0\) for
\(\lambda<1\) and \(>0\) for \(\lambda>1\), so \(\varphi\) decreases on \((0,1]\) and increases on \([1,\infty)\), and
\(\varphi(\lambda)>\varphi(1)=0\) for \(\lambda\ne1\): \(\log\lambda\leqslant\lambda-1\), with equality only for \(\lambda=1\).
\item The eigenvalues \(\lambda_1,\dots,\lambda_p\) of \(A\) are \(>0\) (Theorem~\ref{r8-thm-eigen-criterion}(a)), and
\(\det A=\lambda_1\cdots\lambda_p\), \(\tr A=\lambda_1+\dots+\lambda_p\) (Theorem~\ref{r7-thm-trace-det}). Hence
\[
\log\det A=\sum_{i=1}^p\log\lambda_i\leqslant\sum_{i=1}^p(\lambda_i-1)=\tr A-p
\]
\item Equality forces \(\lambda_i=1\) for every \(i\); then \(A=QI_pQ'=I_p\) for an orthogonal \(Q\) (Theorem~\ref{r8-thm-spectral}(b)).\qedhere
\end{itemize}
\end{proof}

\begin{definition}[\citeat{casella-berger}{Definitions~5.1.1, 7.2.4, (7.2.3)}, \citeat{tu-mc}{\S13}]\label{r14-def-likelihood}
Let \(x_1,\dots,x_n\in\R^p\) be independent random vectors, each \(N_p(\mu,\Sigma)\) with \(\Sigma\) PD: a \textbf{random sample} from
\(N_p(\mu,\Sigma)\). Each \(x_i\) has the density \(f(\cdot\,;\mu,\Sigma)\) of Theorem~\ref{r14-thm-mvn-density}, and
\((x_1',\dots,x_n')'\) has the density \(f(x_1;\mu,\Sigma)\cdots f(x_n;\mu,\Sigma)\). For observed values \(x_1,\dots,x_n\), the function
\[
L(\mu,\Sigma)=\prod_{i=1}^nf(x_i;\mu,\Sigma)\qquad(\mu\in\R^p,\ \Sigma\in\R^{p\times p}\text{ PD})
\]
is the \textbf{likelihood}, and values \(\hat\mu\), \(\hat\Sigma\) at which \(L\) is largest are \textbf{maximum likelihood estimates}
of \(\mu\) and \(\Sigma\).
\end{definition}

\begin{theorem}[\citeat{abadir-magnus}{Exercises~13.61, 12.20}, \citeat{tu-mc}{\S13}, \citeat{casella-berger}{\S7.2.2, Example~7.2.11}]\label{r14-thm-mle}
Let \(x_1,\dots,x_n\) (\(n\geqslant2\)) be a random sample from \(N_p(\mu,\Sigma)\) (Definition~\ref{r14-def-likelihood}), let \(X\) be the
data matrix with rows \(x_1',\dots,x_n'\), and let \(\bar x\) and \(S\) be its mean vector and sample covariance matrix
(Definition~\ref{r12-def-data}), so that \(A=X'C_nX=(n-1)S\). If \(A\) is PD, which requires \(n\geqslant p+1\), then the likelihood has
exactly one maximizer,
\[
\hat\mu=\bar x,\qquad\hat\Sigma=\frac1nA=\frac1n\sum_{i=1}^n(x_i-\bar x)(x_i-\bar x)'=\frac{n-1}nS
\]
and the maximum is \(L(\bar x,\hat\Sigma)=(2\pi e)^{-np/2}(\det\hat\Sigma)^{-n/2}\). Moreover \(\expc(\bar x)=\mu\) and
\(\expc(\hat\Sigma)=\frac{n-1}n\Sigma\).
\end{theorem}
\begin{proof}
\begin{itemize}
\item The requirement: as \(A=(n-1)S\) with \(n-1>0\), \(A\) is PD if and only if \(S\) is; and \(S\) is PD if and only
if \(\rank(C_nX)=p\), which requires \(n\geqslant p+1\) (Proposition~\ref{r12-prop-data}(a)).
\item By Theorem~\ref{r14-thm-mvn-density},
\[
\log L(\mu,\Sigma)=-\frac{np}2\log(2\pi)-\frac n2\log\det\Sigma-\frac12\sum_{i=1}^n(x_i-\mu)'\Sigma^{-1}(x_i-\mu)
\]
As \(\log\) is strictly increasing, \(L\) and \(\log L\) have the same maximizers.
\item The step for \(\mu\). Write \(x_i-\mu=(x_i-\bar x)+(\bar x-\mu)\). As \(\sum_i(x_i-\bar x)=X'\one-n\bar x=0\) and \(\Sigma^{-1}\) is
symmetric, the cross terms add up to \(2(\bar x-\mu)'\Sigma^{-1}\sum_i(x_i-\bar x)=0\), and
\[
\sum_{i=1}^n(x_i-\mu)'\Sigma^{-1}(x_i-\mu)=\sum_{i=1}^n(x_i-\bar x)'\Sigma^{-1}(x_i-\bar x)+n(\bar x-\mu)'\Sigma^{-1}(\bar x-\mu)
\]
\(\Sigma^{-1}\) is PD: for \(y\ne0\), \(\Sigma^{-1}y\ne0\) and \(y'\Sigma^{-1}y=(\Sigma^{-1}y)'\Sigma(\Sigma^{-1}y)>0\). So the last term is
\(>0\) unless \(\mu=\bar x\), and \(\log L(\mu,\Sigma)\leqslant\log L(\bar x,\Sigma)\), with equality only if \(\mu=\bar x\).
\item The trace form. \((x_i-\bar x)'\Sigma^{-1}(x_i-\bar x)=\tr\bigl(\Sigma^{-1}(x_i-\bar x)(x_i-\bar x)'\bigr)\)
(Theorem~\ref{r1-thm-trace}(b)), and \(\sum_i(x_i-\bar x)(x_i-\bar x)'=(C_nX)'(C_nX)=A\), as the rows of \(C_nX\) are the
\((x_i-\bar x)'\) (Definition~\ref{r12-def-data}, Theorem~\ref{r1-thm-views}(d)). So
\(-2\log L(\bar x,\Sigma)=np\log(2\pi)+g(\Sigma)\) with \(g(\Sigma)=n\log\det\Sigma+\tr(\Sigma^{-1}A)\).
\item The step for \(\Sigma\). \(\hat\Sigma=\frac1nA\) is PD. Let \(B=\hat\Sigma^{1/2}\Sigma^{-1}\hat\Sigma^{1/2}\), which is PD
(Theorems~\ref{r8-thm-sqrt} and~\ref{r8-thm-congruence}(c)). By Theorem~\ref{r1-thm-trace}(b),
\(\tr B=\tr(\Sigma^{-1}\hat\Sigma^{1/2}\hat\Sigma^{1/2})=\tr(\Sigma^{-1}\hat\Sigma)\), and \(\det B=\det\hat\Sigma/\det\Sigma\)
(Theorem~\ref{r5d-thm-product}(b),(c)). As \(A=n\hat\Sigma\) and \(\tr(\hat\Sigma^{-1}A)=np\),
\[
g(\Sigma)-g(\hat\Sigma)=n\log\frac{\det\Sigma}{\det\hat\Sigma}+n\tr(\Sigma^{-1}\hat\Sigma)-np=n\bigl(\tr B-\log\det B-p\bigr)\geqslant0
\]
by Lemma~\ref{r14-lem-logdet}, with equality only if \(B=I_p\), that is, \(\Sigma^{-1}=\hat\Sigma^{-1/2}\hat\Sigma^{-1/2}=\hat\Sigma^{-1}\), or
\(\Sigma=\hat\Sigma\).
\item Together, \(\log L(\mu,\Sigma)\leqslant\log L(\bar x,\Sigma)\leqslant\log L(\bar x,\hat\Sigma)\), with equality throughout only for
\(\mu=\bar x\) and \(\Sigma=\hat\Sigma\). As \(g(\hat\Sigma)=n\log\det\hat\Sigma+np\),
\(\log L(\bar x,\hat\Sigma)=-\frac{np}2\log(2\pi)-\frac n2\log\det\hat\Sigma-\frac{np}2\).
\item Independent vectors are uncorrelated (N7), and \(\expc(x_i)=\mu\), \(\var(x_i)=\Sigma\)
(Definition~\ref{r14-def-mvn}); so by Theorem~\ref{r14-thm-sample}, \(\expc(\bar x)=\mu\) and
\(\expc(\hat\Sigma)=\frac{n-1}n\expc(S)=\frac{n-1}n\Sigma\).\qedhere
\end{itemize}
\end{proof}
So the maximum likelihood estimator \(\hat\Sigma\) is not unbiased (Definition~\ref{r11-def-estimator}):
\(\expc(\hat\Sigma)=\frac{n-1}n\Sigma\ne\Sigma\), while \(S\) is unbiased (Theorem~\ref{r14-thm-sample}). For \(p=1\),
Theorem~\ref{r14-thm-mle} says that for a random sample from \(N(\mu,\sigma^2)\) whose observed values \(x_1,\dots,x_n\) are
not all equal, the maximum likelihood estimates are \(\hat\mu=\bar x\) and
\(\hat\sigma^2=\frac1n\sum_{i=1}^n(x_i-\bar x)^2=\frac{n-1}ns^2\), where \(s^2\) is the sample variance.

\begin{example}\label{r14-exm-mle}
Suppose that the five observations of Example~\ref{r12-exm-data} are a random sample from \(N_4(\mu,\Sigma)\). There
\(A=(C_5X)'(C_5X)=4S\), and \(S\) has the eigenvalues \(25,9,5,1\) (Example~\ref{r12-exm-pca}), so \(A\) is PD (Theorem~\ref{r8-thm-eigen-criterion}(a)). By
Theorem~\ref{r14-thm-mle}, the maximum likelihood estimates are
\[
\hat\mu=\bar x=\begin{pmatrix}5\\5\\5\\5\end{pmatrix},\qquad
\hat\Sigma=\frac15\begin{pmatrix}40&28&20&12\\28&40&12&20\\20&12&40&28\\12&20&28&40\end{pmatrix}=\frac45S,\qquad
S=\begin{pmatrix}10&7&5&3\\7&10&3&5\\5&3&10&7\\3&5&7&10\end{pmatrix}
\]
with \(\det\hat\Sigma=\bigl(\frac45\bigr)^4\det S=\frac{256}{625}\cdot25\cdot9\cdot5\cdot1=\frac{2304}5\)
(Theorems~\ref{r5d-thm-product}(c) and~\ref{r7-thm-trace-det}). So the largest value of the likelihood is
\((2\pi e)^{-10}\bigl(\frac{2304}5\bigr)^{-5/2}\). The unbiased estimator of \(\Sigma\) is \(S\) (Theorem~\ref{r14-thm-sample}).
\end{example}

\needspace{14\baselineskip}
\subsection*{Exercises}

\needspace{22\baselineskip}
\begin{exercise}\label{r14-dr-marginal}
Let \(z\sim N_4(\mu,\Sigma)\) with
\[
\mu=\begin{pmatrix}1\\-1\\2\\0\end{pmatrix},\qquad\Sigma=\begin{pmatrix}4&2&0&0\\2&3&1&1\\0&1&2&1\\0&1&1&2\end{pmatrix}
\]
\begin{enumerate}
\item Find the distribution of \((z_1,z_3,z_4)'\). Is \(z_1\) independent of \((z_3,z_4)'\)?
\item Find the distribution of \(e=z_1-z_2+z_3\). Is \(e\) independent of \(z_2\)? Of \((z_2,z_4)'\)? Of \(z_1\)?
\item Find the distribution of \((z_1+z_2,\ z_2+z_3,\ z_3+z_4)'\).
\end{enumerate}
\end{exercise}
\begin{solution}
(a) By Corollary~\ref{r14-cor-marginal}(b),
\[
\begin{pmatrix}z_1\\z_3\\z_4\end{pmatrix}\sim N_3\left(\begin{pmatrix}1\\2\\0\end{pmatrix},\begin{pmatrix}4&0&0\\0&2&1\\0&1&2\end{pmatrix}\right)
\]
As \(\sigma_{13}=\sigma_{14}=0\), \(z_1\) and \((z_3,z_4)'\) are independent (Theorem~\ref{r14-thm-mvn-indep}(a) for this vector,
with \(p=1\) and \(q=2\)).

(b) With \(a=(1,-1,1,0)'\), \(a'\mu=1+1+2+0=4\), \(\Sigma a=(4-2+0+0,\ 2-3+1+0,\ 0-1+2+0,\ 0-1+1+0)'=(2,0,1,0)'\) and
\(a'\Sigma a=2+0+1+0=3\); so \(e\sim N(4,3)\) (Corollary~\ref{r14-cor-marginal}(c)). By Theorem~\ref{r14-thm-cov}(a),
\(\cov(z,e)=\Sigma a=(2,0,1,0)'\). So \(\cov(z_2,e)=\cov(z_4,e)=0\), and \(e\) is independent of \(z_2\) and of
\((z_2,z_4)'\) (Theorem~\ref{r14-thm-mvn-indep}(c) with \(B=a'\) and \(A=e_2'\), resp.\ the matrix \(A\) with the rows
\(e_2'\) and \(e_4'\)); but \(\cov(z_1,e)=2\ne0\), so \(e\) is not independent of \(z_1\).

(c) Let
\[
B=\begin{pmatrix}1&1&0&0\\0&1&1&0\\0&0&1&1\end{pmatrix},\qquad B\mu=\begin{pmatrix}0\\1\\2\end{pmatrix}
\]
The rows of \(B\Sigma\) are \(\row_1(\Sigma)+\row_2(\Sigma)\), \(\row_2(\Sigma)+\row_3(\Sigma)\) and \(\row_3(\Sigma)+\row_4(\Sigma)\)
(Theorem~\ref{r1-thm-views}(c)), and the columns of \(B\Sigma B'\) are the same sums of the columns of \(B\Sigma\)
(Theorem~\ref{r1-thm-views}(a),(b)):
\[
B\Sigma=\begin{pmatrix}6&5&1&1\\2&4&3&2\\0&2&3&3\end{pmatrix},\qquad
B\Sigma B'=\begin{pmatrix}6+5&5+1&1+1\\2+4&4+3&3+2\\0+2&2+3&3+3\end{pmatrix}=\begin{pmatrix}11&6&2\\6&7&5\\2&5&6\end{pmatrix}
\]
so \((z_1+z_2,\ z_2+z_3,\ z_3+z_4)'=Bz\sim N_3(B\mu,B\Sigma B')\) (Theorem~\ref{r14-thm-mvn-linear}).
\end{solution}

\begin{exercise}[\citeat{searle-khuri}{Exercise~10.10}]\label{r14-ex-mgf-probability}
The MGF of a random vector \(x=(x_1,x_2,x_3)'\) is
\[
M_x(t)=\exp\bigl(t_1-t_2+2t_3+t_1^2+\tfrac12t_2^2+2t_3^2-\tfrac12t_1t_2-t_1t_3\bigr)\qquad(t\in\R^3)
\]
Express \(\prb(x_1+x_2>x_3)\) through the distribution function \(\Phi\) of \(N(0,1)\), \(\Phi(c)=\prb(y\leqslant c)\) for
\(y\sim N(0,1)\).
\end{exercise}
\begin{solution}
\begin{itemize}
\item The exponent is \(t'\mu+\frac12t'Vt\) with \(\mu=(1,-1,2)'\) and
\[
V=\begin{pmatrix}2&-\frac12&-1\\-\frac12&1&0\\-1&0&4\end{pmatrix}
\]
since \(t'Vt=2t_1^2+t_2^2+4t_3^2-t_1t_2-2t_1t_3\) (Definition~\ref{r8-def-quadratic-form}) is twice the quadratic part. The leading
principal minors \(2\), \(2-\frac14=\frac74\) and \(\det V=6\) are positive, so \(V\) is PD (Theorem~\ref{r8-thm-leading-minors}), and
\(x\sim N_3(\mu,V)\) by Definition~\ref{r14-def-mvn}.
\item Let \(u=x_1+x_2-x_3=a'x\) with \(a=(1,1,-1)'\). By Corollary~\ref{r14-cor-marginal}(c), \(u\sim N(a'\mu,a'Va)\), with
\(a'\mu=1-1-2=-2\) and \(a'Va=2+1+4+2\bigl(-\frac12\bigr)-2(-1)-2\cdot0=8\).
\item \(z=(u+2)/\sqrt8\sim N(0,1)\) (Theorem~\ref{r14-thm-mvn-linear} with \(A=1/\sqrt8\), \(b=2/\sqrt8\)), so \(z\) has the
distribution of \(y\), and \(\prb(z>c)=1-\prb(z\leqslant c)=1-\Phi(c)\) for every \(c\). As \(2/\sqrt8=\sqrt2/2\),
\(\prb(x_1+x_2>x_3)=\prb(u>0)=\prb\bigl(z>\sqrt2/2\bigr)=1-\Phi\bigl(\sqrt2/2\bigr)\).
\end{itemize}
\end{solution}

\begin{exercise}\label{r11-ex-normal-drill}
Let \(x\sim N_4(\mu,\Sigma)\) with \(\mu=(1,2,0,-1)'\) and
\[
\Sigma=\begin{pmatrix}3&1&1&0\\1&3&1&1\\1&1&2&1\\0&1&1&1\end{pmatrix}
\]
\begin{enumerate}
\item Find the distributions of \(y=(x_1-x_2,\ x_3-x_4)'\) and of \((x_1,x_4)'\).
\item Which of the six pairs \(\{x_i,x_j\}\) are independent? Is \(x_1\) independent of \((x_3,x_4)'\)?
\item Find \(c\) such that \(x_1-cx_3\) and \(x_3\) are independent.
\item Find the conditional distribution of \((x_1,x_2)'\) given \((x_3,x_4)'=(a_3,a_4)'\), and check that \(\det\Sigma\) is the
product of the determinants of the covariance matrix of \((x_3,x_4)'\) and of this conditional covariance matrix.
\end{enumerate}
\end{exercise}
\begin{solution}
(a) \(y=Ax\) with \(A=\left(\begin{smallmatrix}1&-1&0&0\\0&0&1&-1\end{smallmatrix}\right)\). The rows of \(A\Sigma\) are
\(\row_1(\Sigma)-\row_2(\Sigma)=(2,-2,0,-1)\) and \(\row_3(\Sigma)-\row_4(\Sigma)=(1,0,1,0)\), so by
Theorem~\ref{r14-thm-mvn-linear} and Corollary~\ref{r14-cor-marginal}(b)
\[
y\sim N_2\left(\begin{pmatrix}-1\\1\end{pmatrix},\begin{pmatrix}4&1\\1&1\end{pmatrix}\right),\qquad
\begin{pmatrix}x_1\\x_4\end{pmatrix}\sim N_2\left(\begin{pmatrix}1\\-1\end{pmatrix},\begin{pmatrix}3&0\\0&1\end{pmatrix}\right)
\]
(b) Each pair forms a normal subvector of \(x\) (Corollary~\ref{r14-cor-marginal}(b)), so
Theorem~\ref{r14-thm-mvn-indep}(a) applies to it: \(x_1\) and \(x_4\) are independent, as \(\sigma_{14}=0\), and the other five
pairs are not independent, as their covariances are not \(0\). The vector \((x_1,x_3,x_4)'\) is normal, and the
cross-covariance of \(x_1\) and \((x_3,x_4)'\) is \((\sigma_{13},\sigma_{14})=(1,0)\ne0\); so \(x_1\) is not independent
of \((x_3,x_4)'\).

(c) By Theorem~\ref{r14-thm-mvn-indep}(c) with \(A=(1,0,-c,0)\) and \(B=e_3'\),
\(A\Sigma B'=\sigma_{13}-c\sigma_{33}=1-2c\), which is \(0\) exactly for \(c=\frac12\).

(d) Apply Theorem~\ref{r14-thm-mvn-conditional} to the reordered vector \(Px=(x_3,x_4,x_1,x_2)'\), where \(P\) is the
permutation matrix with the rows \(e_3',e_4',e_1',e_2'\); \(Px\sim N_4(P\mu,P\Sigma P')\) (Theorem~\ref{r14-thm-mvn-linear}),
and \(P\Sigma P'\) has the rows and columns of \(\Sigma\) in the order \(3,4,1,2\). The first block is \((x_3,x_4)'\) and the
second \((x_1,x_2)'\):
\[
\Sigma_{11}=\begin{pmatrix}2&1\\1&1\end{pmatrix},\quad\Sigma_{11}^{-1}=\begin{pmatrix}1&-1\\-1&2\end{pmatrix},\quad
\Sigma_{21}=\begin{pmatrix}1&0\\1&1\end{pmatrix},\quad\Sigma_{21}\Sigma_{11}^{-1}=\begin{pmatrix}1&-1\\0&1\end{pmatrix}
\]
with \(\Sigma_{11}\) PD (its leading principal minors are \(2\) and \(1\); Theorem~\ref{r8-thm-leading-minors}) and
\(\Sigma_{11}^{-1}\) from Proposition~\ref{r2-prop-inverse-2x2} (\(\det\Sigma_{11}=1\)), and
\[
\Sigma_{22\cdot1}=\Sigma_{22}-\Sigma_{21}\Sigma_{11}^{-1}\Sigma_{12}=\begin{pmatrix}3&1\\1&3\end{pmatrix}-\begin{pmatrix}1&-1\\0&1\end{pmatrix}
\begin{pmatrix}1&1\\0&1\end{pmatrix}=\begin{pmatrix}2&1\\1&2\end{pmatrix}
\]
The conditional mean is \((1,2)'+\Sigma_{21}\Sigma_{11}^{-1}(a_3-0,\ a_4+1)'=(a_3-a_4,\ a_4+3)'\), so given
\((x_3,x_4)'=(a_3,a_4)'\), the vector \((x_1,x_2)'\) has the distribution \(N_2\bigl((a_3-a_4,\ a_4+3)',\ \Sigma_{22\cdot1}\bigr)\).
Check: subtracting row \(4\) of \(\Sigma\) from rows \(2\) and \(3\) (Definition~\ref{r5d-def-det}(b)) gives the rows \((1,2,0,0)\)
and \((1,0,1,0)\), so that column \(4\) becomes \(e_4\); expanding along column \(4\) and then along row \(1\)
(Theorem~\ref{r5d-thm-expansion}), \(\det\Sigma=3(2\cdot1-0)-1(1\cdot1-0)+1(0-2\cdot1)=3\); and by
Theorem~\ref{r5d-thm-schur-det}(b), with \(\Sigma\) partitioned into the blocks of \((x_1,x_2)'\) and \((x_3,x_4)'\), whose
lower right block \(\Sigma_{11}\) is invertible, \(\det\Sigma=\det\Sigma_{11}\det\Sigma_{22\cdot1}=1\cdot3\).
\end{solution}

\begin{exercise}\label{r14-dr-indep}
Let \(z\sim N_4(0,\Sigma)\) with
\[
\Sigma=I_4+J_4=\begin{pmatrix}2&1&1&1\\1&2&1&1\\1&1&2&1\\1&1&1&2\end{pmatrix}
\]
\begin{enumerate}
\item Show that \(z_1+z_2+z_3+z_4\) is independent of \((z_1-z_2,\ z_2-z_3,\ z_3-z_4)'\).
\item Find all \(a\in\R^4\) for which \(a'z\) is independent of \(z_1-z_2\), of \(z_2-z_3\) and of \(z_3-z_4\).
\item Find \(b\in\R^3\) for which \(z_4-b_1z_1-b_2z_2-b_3z_3\) is independent of \((z_1,z_2,z_3)'\), and the variance of this
variable.
\end{enumerate}
\end{exercise}
\begin{solution}
(a) \(\Sigma\one=(5,5,5,5)'\), so with \(A=\one'\) and
\[
B=\begin{pmatrix}1&-1&0&0\\0&1&-1&0\\0&0&1&-1\end{pmatrix}
\]
\(A\Sigma B'=(B\Sigma\one)'=(5-5,\ 5-5,\ 5-5)=(0,0,0)\). By Theorem~\ref{r14-thm-mvn-indep}(c), \(\one'z\) and \(Bz\) are
independent.

(b) \(\Sigma(1,-1,0,0)'=(2-1,\ 1-2,\ 1-1,\ 1-1)'=(1,-1,0,0)'\), and likewise \(\Sigma(0,1,-1,0)'=(0,1,-1,0)'\) and
\(\Sigma(0,0,1,-1)'=(0,0,1,-1)'\), as \(J_4\) maps these vectors to \(0\). By Theorem~\ref{r14-thm-mvn-indep}(c), \(a'z\) is
independent of \(z_1-z_2\) if and only if \(a'\Sigma(1,-1,0,0)'=a_1-a_2=0\), of \(z_2-z_3\) if and only if \(a_2-a_3=0\), and of
\(z_3-z_4\) if and only if \(a_3-a_4=0\). So \(a=c\one\) (\(c\in\R\)): only the multiples of the total \(z_1+z_2+z_3+z_4\).

(c) Let \(a=(-b_1,-b_2,-b_3,1)'\) and \(A=[I_3\ \ 0]\in\R^{3\times4}\). Then \(A\Sigma a=(1,1,1)'-\Sigma_{11}b\), where
\(\Sigma_{11}=I_3+J_3\) is the leading \(3\times3\) block of \(\Sigma\). Its inverse is \(I_3-\frac14J_3\):
\((I_3+J_3)(I_3-\frac14J_3)=I_3+\frac34J_3-\frac14J_3^2=I_3\), as \(J_3^2=3J_3\) (Proposition~\ref{r1-prop-centering}). By
Theorem~\ref{r14-thm-mvn-indep}(c), independence holds exactly when
\(b=\Sigma_{11}^{-1}(1,1,1)'=(1,1,1)'-\frac34(1,1,1)'=\frac14(1,1,1)'\). The variance is
\(a'\Sigma a=\sigma_{44}-(1,1,1)\Sigma_{11}^{-1}(1,1,1)'=2-\frac34=\frac54\), as \(\Sigma a=(0,0,0,\ 2-\frac14-\frac14-\frac14)'\).
\end{solution}

\begin{exercise}[\citeat{searle-khuri}{Exercise~10.9}]\label{r14-ex-quadratic-form}
Let \(x=(x_1,x_2,x_3)'\sim N_3(\mu,V)\) with \(V\) PD, and suppose that
\[
Q=(x-\mu)'V^{-1}(x-\mu)=2x_1^2+3x_2^2+4x_3^2+2x_1x_2-2x_1x_3-4x_2x_3-6x_1-6x_2+10x_3+8
\]
\begin{enumerate}
\item Find \(\mu\) and \(V\).
\item Find the MGF of \(x\).
\end{enumerate}
\end{exercise}
\begin{solution}
(a)
\begin{itemize}
\item For a symmetric \(B\in\R^{3\times3}\), the partial derivatives of
\((x-\mu)'B(x-\mu)=\sum_k\sum_lb_{kl}(x_k-\mu_k)(x_l-\mu_l)\) are \(2\sum_lb_{il}(x_l-\mu_l)\), the entries of \(2B(x-\mu)\) (as in the
paragraph after the facts (N1)--(N7)). So the vector of the partial derivatives of \(Q\) is \(2V^{-1}(x-\mu)\); from the given form
of \(Q\) it is
\[
\bigl(4x_1+2x_2-2x_3-6,\ 2x_1+6x_2-4x_3-6,\ -2x_1-4x_2+8x_3+10\bigr)'
\]
\item Comparing the coefficients of \(x_1,x_2,x_3\), that is, the second partial derivatives,
\[
2V^{-1}=\begin{pmatrix}4&2&-2\\2&6&-4\\-2&-4&8\end{pmatrix},\qquad
V^{-1}=\begin{pmatrix}2&1&-1\\1&3&-2\\-1&-2&4\end{pmatrix}
\]
The leading principal minors of \(V^{-1}\) are \(2\), \(5\) and \(13\), so \(V^{-1}\) is PD (Theorem~\ref{r8-thm-leading-minors}).
\item \(\expc\bigl(V^{-1}(x-\mu)\bigr)=V^{-1}(\expc(x)-\mu)=0\) (Theorem~\ref{r11-thm-linear}(b),
Definition~\ref{r14-def-mvn}). Taking expectations of the three entries above, with \(\expc(x)=\mu\),
\[
\begin{pmatrix}4&2&-2\\2&6&-4\\-2&-4&8\end{pmatrix}\mu=\begin{pmatrix}6\\6\\-10\end{pmatrix},\qquad\text{so}\qquad\mu=\begin{pmatrix}1\\0\\-1\end{pmatrix}
\]
since \((4-0+2,\ 2-0+4,\ -2-0-8)'=(6,6,-10)'\) and the coefficient matrix is invertible (its determinant is \(8\cdot13\)). Check: the
constant term is \(\mu'V^{-1}\mu=(1,0,-1)(3,3,-5)'=8\).
\item By Theorem~\ref{r5d-thm-adjugate}, with \(\det V^{-1}=13\),
\[
V=\frac1{13}\begin{pmatrix}8&-2&1\\-2&7&3\\1&3&5\end{pmatrix}
\]
\end{itemize}
(b) By Definition~\ref{r14-def-mvn}, \(M_x(t)=\exp\bigl(t'\mu+\frac12t'Vt\bigr)\) with \(\mu\) and \(V\) of (a), that is,
\(M_x(t)=\exp\bigl(t_1-t_3+\frac1{26}(8t_1^2+7t_2^2+5t_3^2-4t_1t_2+2t_1t_3+6t_2t_3)\bigr)\).
\end{solution}

\begin{exercise}\label{r11-ex-density-drill}
A random vector \(x\in\R^4\) has the density \(f(u)=(2\pi)^{-2}\exp\bigl(-\frac12Q(u)\bigr)\), where
\[
Q(u)=3u_1^2+3u_2^2+2u_3^2+u_4^2+4u_1u_2-2u_1u_3-4u_2u_3+2u_2u_4-2u_3u_4-10u_1-8u_2+4u_3+9
\]
Show that \(x\) is normal, and find \(\mu\), \(\Sigma\), the correlation \(\rho_{12}\) of \(x_1\) and \(x_2\), and the conditional
distribution of \((x_3,x_4)'\) given \((x_1,x_2)'=(a_1,a_2)'\).
\end{exercise}
\begin{solution}
\begin{itemize}
\item The quadratic part of \(Q\) is \(u'Pu\) for the symmetric matrix \(P\) below (Definition~\ref{r8-def-quadratic-form}), and
the linear part is \(-2c'u\) with \(c=(5,4,-2,0)'\). Row reduction of \([P\mid I_4]\) (Corollary~\ref{r2-cor-computing-inverse}(a)), with the pivots taken from the last
column back:
\[
\begin{gathered}
\left(\begin{array}{cccc|cccc}3&2&-1&0&1&0&0&0\\2&3&-2&1&0&1&0&0\\-1&-2&2&-1&0&0&1&0\\0&1&-1&1&0&0&0&1\end{array}\right)
\xrightarrow[R_3\leftarrow R_3+R_4]{R_2\leftarrow R_2-R_4}
\left(\begin{array}{cccc|cccc}3&2&-1&0&1&0&0&0\\2&2&-1&0&0&1&0&-1\\-1&-1&1&0&0&0&1&1\\0&1&-1&1&0&0&0&1\end{array}\right)
\\
\xrightarrow[R_4\leftarrow R_4+R_3]{R_1\leftarrow R_1+R_3,\ R_2\leftarrow R_2+R_3}
\left(\begin{array}{cccc|cccc}2&1&0&0&1&0&1&1\\1&1&0&0&0&1&1&0\\-1&-1&1&0&0&0&1&1\\-1&0&0&1&0&0&1&2\end{array}\right)
\xrightarrow[R_3\leftarrow R_3+R_2]{R_1\leftarrow R_1-R_2}
\left(\begin{array}{cccc|cccc}1&0&0&0&1&-1&0&1\\1&1&0&0&0&1&1&0\\0&0&1&0&0&1&2&1\\-1&0&0&1&0&0&1&2\end{array}\right)
\\
\xrightarrow[R_4\leftarrow R_4+R_1]{R_2\leftarrow R_2-R_1}
\left(\begin{array}{cccc|cccc}1&0&0&0&1&-1&0&1\\0&1&0&0&-1&2&1&-1\\0&0&1&0&0&1&2&1\\0&0&0&1&1&-1&1&3\end{array}\right)
\end{gathered}
\]
So \(P\) is invertible, with
\[
\Sigma=P^{-1}=\begin{pmatrix}1&-1&0&1\\-1&2&1&-1\\0&1&2&1\\1&-1&1&3\end{pmatrix}
\]
Every arrow added multiples of rows, which does not change the determinant (Definition~\ref{r5d-def-det}(b)), and the left
half became \(I_4\); so \(\det P=1\) and \(\det\Sigma=1/\det P=1\) (Theorem~\ref{r5d-thm-product}(c)). The leading principal
minors of \(\Sigma\) are \(1\), \(2-1=1\), \(1\) and \(\det\Sigma=1\), so \(\Sigma\) is PD (Theorem~\ref{r8-thm-leading-minors}); for
the third, adding row \(1\) to row \(2\) of the leading \(3\times3\) block and then subtracting the new row \(2\) from row \(3\)
gives an upper triangular matrix with the diagonal entries \(1,1,1\) (Definition~\ref{r5d-def-det}(b),
Theorem~\ref{r5d-thm-basic}(c)).
\item Let \(\mu=\Sigma c=(5-4+0+0,\ -5+8-2+0,\ 0+4-4+0,\ 5-4-2+0)'=(1,1,0,-1)'\), so that \(P\mu=c\). Then
\((u-\mu)'P(u-\mu)=u'Pu-2c'u+\mu'c\), and \(\mu'c=5+4+0+0=9\) is the constant term of \(Q\); so
\(Q(u)=(u-\mu)'\Sigma^{-1}(u-\mu)\) and \(f(u)=(2\pi)^{-2}(\det\Sigma)^{-1/2}\exp\bigl(-\frac12(u-\mu)'\Sigma^{-1}(u-\mu)\bigr)\). By
Theorem~\ref{r14-thm-mvn-density}(b), \(x\sim N_4(\mu,\Sigma)\).
\item \(\rho_{12}=\sigma_{12}/\sqrt{\sigma_{11}\sigma_{22}}=-1/\sqrt{1\cdot2}=-\sqrt2/2\) (Definition~\ref{r14-def-corr}).
\item In Theorem~\ref{r14-thm-mvn-conditional}, the first block is \((x_1,x_2)'\):
\(\Sigma_{11}=\left(\begin{smallmatrix}1&-1\\-1&2\end{smallmatrix}\right)\) is PD as a principal submatrix of \(\Sigma\)
(Theorem~\ref{r8-thm-principal-sub}), with determinant \(1\) and the inverse
\(\left(\begin{smallmatrix}2&1\\1&1\end{smallmatrix}\right)\) (Proposition~\ref{r2-prop-inverse-2x2}), and
\(\Sigma_{21}=\left(\begin{smallmatrix}0&1\\1&-1\end{smallmatrix}\right)=\Sigma_{12}\). So
\[
\Sigma_{21}\Sigma_{11}^{-1}=\begin{pmatrix}1&1\\1&0\end{pmatrix},\qquad
\Sigma_{22\cdot1}=\begin{pmatrix}2&1\\1&3\end{pmatrix}-\begin{pmatrix}1&1\\1&0\end{pmatrix}\begin{pmatrix}0&1\\1&-1\end{pmatrix}
=\begin{pmatrix}2&1\\1&3\end{pmatrix}-\begin{pmatrix}1&0\\0&1\end{pmatrix}=\begin{pmatrix}1&1\\1&2\end{pmatrix}
\]
and the conditional mean is \((0,-1)'+\Sigma_{21}\Sigma_{11}^{-1}(a_1-1,\ a_2-1)'=(a_1+a_2-2,\ a_1-2)'\). Given
\((x_1,x_2)'=(a_1,a_2)'\), the vector \((x_3,x_4)'\) has the distribution \(N_2\bigl((a_1+a_2-2,\ a_1-2)',\ \Sigma_{22\cdot1}\bigr)\).
Check: the lower right \(2\times2\) block of \(P=\Sigma^{-1}\) is
\(\left(\begin{smallmatrix}2&-1\\-1&1\end{smallmatrix}\right)=\Sigma_{22\cdot1}^{-1}\) (Proposition~\ref{r2-prop-inverse-2x2}), as
Theorem~\ref{r14-thm-mvn-conditional}(b) says.
\end{itemize}
\end{solution}

\begin{exercise}\label{r14-dr-cov-valid}
Which of the following matrices can be the covariance matrix of a random vector?
\[
M_1=\begin{pmatrix}4&3&1&0\\3&1&0&1\\1&0&2&1\\0&1&1&3\end{pmatrix},\qquad
M_2=\begin{pmatrix}1&-\frac25&-\frac25&-\frac25\\[2pt]-\frac25&1&-\frac25&-\frac25\\[2pt]-\frac25&-\frac25&1&-\frac25\\[2pt]-\frac25&-\frac25&-\frac25&1\end{pmatrix},\qquad
M_3=\begin{pmatrix}3&-1&-1&-1\\-1&3&-1&-1\\-1&-1&3&-1\\-1&-1&-1&3\end{pmatrix}
\]
\end{exercise}
\begin{solution}
A covariance matrix is symmetric and PSD (Theorem~\ref{r11-thm-linear}(c)), and every symmetric PSD matrix is the covariance
matrix of some normal vector (Theorem~\ref{r14-thm-mvn-iid}(b)).
\begin{itemize}
\item \(M_1\): for \(x=(1,-2,0,0)'\), \(x'M_1x=4-12+4=-4<0\), so \(M_1\) is not PSD and cannot be a covariance matrix. (Its
entries would also give the correlation \(3/\sqrt{4\cdot1}=\frac32>1\), against Proposition~\ref{r14-prop-corr}(a).)
\item \(M_2\): \(\one'M_2\one\), the sum of its entries, is \(4-12\cdot\frac25=-\frac45<0\); so \(M_2\) is not PSD and cannot be a
covariance matrix.
\item \(M_3=4I_4-J_4=4C_4\) (Definition~\ref{r1-def-ones}). As \(C_4\) is symmetric and idempotent
(Proposition~\ref{r1-prop-centering}), \(x'M_3x=4x'C_4'C_4x=4\lVert C_4x\rVert^2\geqslant0\): \(M_3\) is PSD, so it is a covariance
matrix, a singular one (\(M_3\one=0\)).
\end{itemize}
\end{solution}

\begin{exercise}\label{r14-dr-random-walk}
Let \(w_1,\dots,w_4\) be independent \(N(0,1)\) variables and \(z_k=w_1+\dots+w_k\) (\(k=1,\dots,4\)), the positions of a random
walk.
\begin{enumerate}
\item Write \(z=Lw\) and find the distribution of \(z\).
\item Find \(\Sigma^{-1}\) by row reduction of \([\Sigma\mid I_4]\), and \(\det\Sigma\).
\item Write the density of \(z\), and show that \(x'\Sigma^{-1}x=x_1^2+(x_2-x_1)^2+(x_3-x_2)^2+(x_4-x_3)^2\).
\item Show that \(L'\) is the Cholesky factor of \(\Sigma\). A program produces independent \(N(0,1)\) numbers; compute
\(\mu+Lw\) for \(\mu=(1,0,-1,0)'\) and the four numbers \(w=(1,-2,2,1)'\). What is the distribution of \(\mu+Lw\) when \(w\)
consists of four such numbers?
\end{enumerate}
\end{exercise}
\begin{solution}
(a) By Theorem~\ref{r14-thm-mvn-iid}(a), \(w\sim N_4(0,I_4)\), and \(z=Lw\) with
\[
L=\begin{pmatrix}1&0&0&0\\1&1&0&0\\1&1&1&0\\1&1&1&1\end{pmatrix},\qquad
\Sigma=LL'=\begin{pmatrix}1&1&1&1\\1&2&2&2\\1&2&3&3\\1&2&3&4\end{pmatrix}
\]
So \(z\sim N_4(0,\Sigma)\) (Theorem~\ref{r14-thm-mvn-linear}). The \((i,j)\) entry of \(LL'\) is \(\row_i(L)\row_j(L)'\)
(Theorem~\ref{r1-thm-views}(d)); \(\row_i(L)\) has ones in the positions \(1,\dots,i\), so this product counts the positions
\(k\leqslant i\) with \(k\leqslant j\): \(\sigma_{ij}\) is the smaller of \(i\) and \(j\).

(b) Row-reduce \([\Sigma\mid I_4]\) (Corollary~\ref{r2-cor-computing-inverse}(a)):
\[
\begin{gathered}
\left(\begin{array}{cccc|cccc}1&1&1&1&1&0&0&0\\1&2&2&2&0&1&0&0\\1&2&3&3&0&0&1&0\\1&2&3&4&0&0&0&1\end{array}\right)
\xrightarrow[R_4\leftarrow R_4-R_1]{R_2\leftarrow R_2-R_1,\ R_3\leftarrow R_3-R_1}
\left(\begin{array}{cccc|cccc}1&1&1&1&1&0&0&0\\0&1&1&1&-1&1&0&0\\0&1&2&2&-1&0&1&0\\0&1&2&3&-1&0&0&1\end{array}\right)
\\
\xrightarrow[R_4\leftarrow R_4-R_2]{R_3\leftarrow R_3-R_2}
\left(\begin{array}{cccc|cccc}1&1&1&1&1&0&0&0\\0&1&1&1&-1&1&0&0\\0&0&1&1&0&-1&1&0\\0&0&1&2&0&-1&0&1\end{array}\right)
\xrightarrow{R_4\leftarrow R_4-R_3}
\left(\begin{array}{cccc|cccc}1&1&1&1&1&0&0&0\\0&1&1&1&-1&1&0&0\\0&0&1&1&0&-1&1&0\\0&0&0&1&0&0&-1&1\end{array}\right)
\\
\xrightarrow[R_3\leftarrow R_3-R_4]{R_1\leftarrow R_1-R_4,\ R_2\leftarrow R_2-R_4}
\left(\begin{array}{cccc|cccc}1&1&1&0&1&0&1&-1\\0&1&1&0&-1&1&1&-1\\0&0&1&0&0&-1&2&-1\\0&0&0&1&0&0&-1&1\end{array}\right)
\\
\xrightarrow[R_2\leftarrow R_2-R_3]{R_1\leftarrow R_1-R_3}
\left(\begin{array}{cccc|cccc}1&1&0&0&1&1&-1&0\\0&1&0&0&-1&2&-1&0\\0&0&1&0&0&-1&2&-1\\0&0&0&1&0&0&-1&1\end{array}\right)
\xrightarrow{R_1\leftarrow R_1-R_2}
\left(\begin{array}{cccc|cccc}1&0&0&0&2&-1&0&0\\0&1&0&0&-1&2&-1&0\\0&0&1&0&0&-1&2&-1\\0&0&0&1&0&0&-1&1\end{array}\right)
\end{gathered}
\]
So
\[
\Sigma^{-1}=\begin{pmatrix}2&-1&0&0\\-1&2&-1&0\\0&-1&2&-1\\0&0&-1&1\end{pmatrix}
\]
The first three arrows added multiples of rows, which does not change the determinant (Definition~\ref{r5d-def-det}(b)), and
gave an upper triangular left half with diagonal entries \(1,1,1,1\); so \(\det\Sigma=1\) (Theorem~\ref{r5d-thm-basic}(c)). Also
\(\det\Sigma=\det L\cdot\det L'=(\det L)^2=1\) (Theorems~\ref{r5d-thm-product}(b) and~\ref{r5d-thm-transpose}).

(c) \(L'\) is upper triangular with the diagonal entries \(1\), so \(\det L'=1\) (Theorem~\ref{r5d-thm-basic}(c)) and \(L'\) is
invertible (Theorem~\ref{r5d-thm-product}(a)); hence \(\Sigma=LL'=(L')'L'\) is PD (Theorem~\ref{r8-thm-psd-factor}(c)). By
Theorem~\ref{r14-thm-mvn-density}, \(z\) has the density \(f(x)=(2\pi)^{-2}\exp\bigl(-\frac12x'\Sigma^{-1}x\bigr)\), and
\[
x'\Sigma^{-1}x=2x_1^2+2x_2^2+2x_3^2+x_4^2-2x_1x_2-2x_2x_3-2x_3x_4=x_1^2+(x_2-x_1)^2+(x_3-x_2)^2+(x_4-x_3)^2
\]
as expanding the right side gives \(x_1^2+x_2^2-2x_1x_2+x_1^2+x_3^2-2x_2x_3+x_2^2+x_4^2-2x_3x_4+x_3^2\). The four squares are the
squared steps \(w_1=z_1\), \(w_2=z_2-z_1\), \(w_3=z_3-z_2\), \(w_4=z_4-z_3\), that is, \(L^{-1}x\).

(d) \(L'\) is upper triangular with positive diagonal entries, and \((L')'L'=LL'=\Sigma\). As \(\Sigma\) is PD by (c), the
uniqueness in Theorem~\ref{r8-thm-cholesky} shows that \(L'\) is the Cholesky factor of \(\Sigma\). If \(w\) consists of four independent \(N(0,1)\)
numbers, then \(w\sim N_4(0,I_4)\) and \(\mu+Lw\sim N_4(\mu,LL')=N_4(\mu,\Sigma)\) (Theorem~\ref{r14-thm-mvn-iid}(a),(c)): the
entries of \(Lw\) are the partial sums \(w_1+\dots+w_k\), a random walk with \(N(0,1)\) steps, shifted by \(\mu\). Here
\(Lw=(1,\ 1-2,\ 1-2+2,\ 1-2+2+1)'=(1,-1,1,2)'\), so \(\mu+Lw=(2,-1,0,2)'\).
\end{solution}

\needspace{16\baselineskip}
\begin{exercise}\label{r14-dr-triangular}
Let \(z\sim N_4(\mu,\Sigma)\) with
\[
\mu=\begin{pmatrix}1\\0\\-1\\2\end{pmatrix},\qquad\Sigma=\begin{pmatrix}4&2&-2&2\\2&2&1&0\\-2&1&6&-2\\2&0&-2&4\end{pmatrix}
\]
\begin{enumerate}
\item Find the lower triangular \(A\) with positive diagonal entries and \(AA'=\Sigma\).
\item Write \(z\) as \(\mu+Aw\) with \(w\sim N_4(0,I_4)\), and express \(w\) through \(z\).
\item Write \((z-\mu)'\Sigma^{-1}(z-\mu)\) as a sum of four squares, and give its distribution.
\end{enumerate}
\end{exercise}
\begin{solution}
(a) Comparing the entries of \(AA'\) and \(\Sigma\) column by column: \(a_{11}^2=4\), so \(a_{11}=2\); \(a_{21}a_{11}=2\),
\(a_{31}a_{11}=-2\) and \(a_{41}a_{11}=2\), so \(a_{21}=1\), \(a_{31}=-1\), \(a_{41}=1\); \(a_{21}^2+a_{22}^2=2\), so \(a_{22}=1\);
\(a_{31}a_{21}+a_{32}a_{22}=1\) and \(a_{41}a_{21}+a_{42}a_{22}=0\), so \(a_{32}=2\), \(a_{42}=-1\); \(a_{31}^2+a_{32}^2+a_{33}^2=6\), so
\(a_{33}=1\); \(a_{41}a_{31}+a_{42}a_{32}+a_{43}a_{33}=-2\), so \(a_{43}=-2+1+2=1\); \(a_{41}^2+a_{42}^2+a_{43}^2+a_{44}^2=4\), so
\(a_{44}=1\). Thus
\[
A=\begin{pmatrix}2&0&0&0\\1&1&0&0\\-1&2&1&0\\1&-1&1&1\end{pmatrix},\qquad
AA'=\begin{pmatrix}4&2&-2&2\\2&1+1&-1+2&1-1\\-2&-1+2&1+4+1&-1-2+1\\2&1-1&-1-2+1&1+1+1+1\end{pmatrix}=\Sigma
\]
\(A'\) is upper triangular with positive diagonal entries, so it is invertible (\(\det A'=2\ne0\),
Theorems~\ref{r5d-thm-basic}(c) and~\ref{r5d-thm-product}(a)),
\(\Sigma=(A')'A'\) is PD (Theorem~\ref{r8-thm-psd-factor}(c)), and \(A'\) is the Cholesky factor of \(\Sigma\)
(Theorem~\ref{r8-thm-cholesky}): \(A\) is the only such matrix.

(b) By Theorem~\ref{r14-thm-mvn-iid}(c), \(z\) has the distribution of \(\mu+Aw\) with \(w\sim N_4(0,I_4)\):
\(z_1=1+2w_1\), \(z_2=w_1+w_2\), \(z_3=-1-w_1+2w_2+w_3\), \(z_4=2+w_1-w_2+w_3+w_4\). Conversely,
\(w=A^{-1}(z-\mu)\sim N_4(0,A^{-1}\Sigma(A^{-1})')=N_4(0,I_4)\) (Theorem~\ref{r14-thm-mvn-linear}), with
\[
A^{-1}=\begin{pmatrix}\frac12&0&0&0\\-\frac12&1&0&0\\\frac32&-2&1&0\\-\frac52&3&-1&1\end{pmatrix},\qquad
AA^{-1}=\begin{pmatrix}1&0&0&0\\\frac12-\frac12&1&0&0\\-\frac12-1+\frac32&2-2&1&0\\\frac12+\frac12+\frac32-\frac52&-1-2+3&1-1&1\end{pmatrix}=I_4
\]
(Theorem~\ref{r2-thm-invertible}), that is, \(w_1=\frac12(z_1-1)\), \(w_2=z_2-\frac12(z_1-1)\), \(w_3=\frac32(z_1-1)-2z_2+(z_3+1)\) and
\(w_4=-\frac52(z_1-1)+3z_2-(z_3+1)+(z_4-2)\).

(c) \(\Sigma^{-1}=(AA')^{-1}=(A^{-1})'A^{-1}\) (Theorem~\ref{r2-thm-inverse-props}(b),(c)), so
\((z-\mu)'\Sigma^{-1}(z-\mu)=\lVert A^{-1}(z-\mu)\rVert^2=w_1^2+w_2^2+w_3^2+w_4^2\), with \(w_1,\dots,w_4\) as in (b). They are
independent \(N(0,1)\) (Theorem~\ref{r14-thm-mvn-iid}(a)), so the sum has the distribution \(\chi^2_4\) (N5).
\end{solution}

\begin{exercise}\label{r14-dr-sqrt}
Let \(z\sim N_4(\mu,\Sigma)\) with
\[
\Sigma=\begin{pmatrix}25&16&-20&-20\\16&25&-20&-20\\-20&-20&25&16\\-20&-20&16&25\end{pmatrix}
\]
\begin{enumerate}
\item Find \(\Sigma^{1/2}\), \(\Sigma^{-1/2}\) and \(\Sigma^{-1}\).
\item Find the standardized vector \(w=\Sigma^{-1/2}(z-\mu)\), and write \(z\) through \(w\).
\item Give the distribution of \((z-\mu)'\Sigma^{-1}(z-\mu)\).
\end{enumerate}
\end{exercise}
\begin{solution}
(a) Let \(h=(1,1,-1,-1)'\), \(E_1=\bar J_4=\frac14J_4\), \(E_2=\frac14hh'\) and \(E_3=I_4-E_1-E_2\). Here \(\one'h=0\) and
\(h'h=4\), so \(E_1E_2=\frac1{16}\one(\one'h)h'=0=E_2E_1\) and \(E_2^2=\frac1{16}h(h'h)h'=E_2\); also \(E_1^2=E_1\)
(Proposition~\ref{r1-prop-centering}). Hence \(E_3^2=E_3\) and \(E_1E_3=E_2E_3=0\), and \(E_1\), \(E_2\), \(E_3\) are symmetric,
idempotent and nonzero (the diagonal entries of \(E_3\) are \(1-\frac14-\frac14=\frac12\)), with \(E_1+E_2+E_3=I_4\). Comparing entries (the diagonal \(9-2+18\), the entries \((1,2)\) and
\((3,4)\) \(-2+18\), the others \(-2-18\)),
\[
\Sigma=9I_4-2J_4+18hh'=E_1+81E_2+9E_3
\]
This is the spectral decomposition of \(\Sigma\) (Theorem~\ref{r8-thm-spectral-decomp}(d)), with the eigenvalues \(1\) (on
\(\one_4\)), \(81\) (on \(h\)) and \(9\) (on the plane orthogonal to \(\one_4\) and \(h\)); they are positive, so \(\Sigma\) is PD
(Theorem~\ref{r8-thm-eigen-criterion}(a)). By Theorem~\ref{r8-thm-sqrt}(a),(c),
\[
\Sigma^{1/2}=E_1+9E_2+3E_3=3I_4-\tfrac12J_4+\tfrac32hh'=\begin{pmatrix}4&1&-2&-2\\1&4&-2&-2\\-2&-2&4&1\\-2&-2&1&4\end{pmatrix}
\]
and
\[
\Sigma^{-1/2}=E_1+\tfrac19E_2+\tfrac13E_3=\tfrac13I_4+\tfrac16J_4-\tfrac1{18}hh'
=\frac19\begin{pmatrix}4&1&2&2\\1&4&2&2\\2&2&4&1\\2&2&1&4\end{pmatrix}
\]
(the diagonal \(\frac13+\frac16-\frac1{18}=\frac49\), the entries \((1,2)\) and \((3,4)\) \(\frac16-\frac1{18}=\frac19\), the
others \(\frac16+\frac1{18}=\frac29\)), and \(\Sigma^{-1}=E_1+\frac1{81}E_2+\frac19E_3\) (Theorem~\ref{r8-thm-spectral-decomp}(b)), that is,
\[
\Sigma^{-1}=\tfrac19I_4+\tfrac29J_4-\tfrac2{81}hh'
=\frac1{81}\begin{pmatrix}25&16&20&20\\16&25&20&20\\20&20&25&16\\20&20&16&25\end{pmatrix}
\]
(the diagonal \(\frac19+\frac29-\frac2{81}=\frac{25}{81}\), the entries \((1,2)\) and \((3,4)\)
\(\frac29-\frac2{81}=\frac{16}{81}\), the others \(\frac29+\frac2{81}=\frac{20}{81}\)). Check:
\[
\begin{pmatrix}4&1&-2&-2\\1&4&-2&-2\\-2&-2&4&1\\-2&-2&1&4\end{pmatrix}^2=\begin{pmatrix}25&16&-20&-20\\16&25&-20&-20\\-20&-20&25&16\\-20&-20&16&25\end{pmatrix},\qquad
\begin{pmatrix}4&1&-2&-2\\1&4&-2&-2\\-2&-2&4&1\\-2&-2&1&4\end{pmatrix}\begin{pmatrix}4&1&2&2\\1&4&2&2\\2&2&4&1\\2&2&1&4\end{pmatrix}=9I_4
\]
and \(\Sigma\) times the integer matrix of \(\Sigma^{-1}\) is \(81I_4\); for instance \(25\cdot25+16\cdot16-20\cdot20-20\cdot20=81\) and
\(25\cdot16+16\cdot25-20\cdot20-20\cdot20=0\).

(b) By Corollary~\ref{r14-cor-standard}(a), \(w=\Sigma^{-1/2}(z-\mu)\sim N_4(0,I_4)\): with \(d=z-\mu\),
\(w_1=\frac19(4d_1+d_2+2d_3+2d_4)\), \(w_2=\frac19(d_1+4d_2+2d_3+2d_4)\), \(w_3=\frac19(2d_1+2d_2+4d_3+d_4)\) and
\(w_4=\frac19(2d_1+2d_2+d_3+4d_4)\), independent \(N(0,1)\); and \(z=\mu+\Sigma^{1/2}w\): \(z_1=\mu_1+4w_1+w_2-2w_3-2w_4\),
\(z_2=\mu_2+w_1+4w_2-2w_3-2w_4\), \(z_3=\mu_3-2w_1-2w_2+4w_3+w_4\), \(z_4=\mu_4-2w_1-2w_2+w_3+4w_4\).

(c) By Corollary~\ref{r14-cor-standard}(a), \((z-\mu)'\Sigma^{-1}(z-\mu)=w_1^2+w_2^2+w_3^2+w_4^2\sim\chi^2_4\).
\end{solution}

\needspace{16\baselineskip}
\begin{exercise}\label{r14-dr-singular}
Let \(z\sim N_4(\mu,\Sigma)\) with
\[
\mu=\begin{pmatrix}1\\2\\1\\0\end{pmatrix},\qquad\Sigma=\begin{pmatrix}1&1&0&1\\1&2&1&0\\0&1&1&-1\\1&0&-1&2\end{pmatrix}
\]
\begin{enumerate}
\item Find the RREF of \(\Sigma\), its rank, and bases of \(\mathcal N(\Sigma)\) and \(\mathcal C(\Sigma)\).
\item Show that \(z_1-z_2+z_3=0\) and \(z_2+z_4=2z_1\) with probability \(1\).
\item Find \(A\in\R^{4\times2}\) with \(AA'=\Sigma\), and check (b) on the representation \(z=\mu+Aw\), \(w\sim N_2(0,I_2)\).
\end{enumerate}
\end{exercise}
\begin{solution}
(a) Row reduction:
\[
\begin{gathered}
\Sigma=\begin{pmatrix}1&1&0&1\\1&2&1&0\\0&1&1&-1\\1&0&-1&2\end{pmatrix}
\xrightarrow[R_4\leftarrow R_4-R_1]{R_2\leftarrow R_2-R_1}
\begin{pmatrix}1&1&0&1\\0&1&1&-1\\0&1&1&-1\\0&-1&-1&1\end{pmatrix}
\\
\xrightarrow[R_4\leftarrow R_4+R_2]{R_3\leftarrow R_3-R_2}
\begin{pmatrix}1&1&0&1\\0&1&1&-1\\0&0&0&0\\0&0&0&0\end{pmatrix}
\xrightarrow{R_1\leftarrow R_1-R_2}
\begin{pmatrix}1&0&-1&2\\0&1&1&-1\\0&0&0&0\\0&0&0&0\end{pmatrix}
\end{gathered}
\]
So \(\rank\Sigma=2\), and \(\Sigma x=0\) means \(x_1=x_3-2x_4\), \(x_2=-x_3+x_4\): \(\mathcal N(\Sigma)\) has the basis
\(a=(1,-1,1,0)'\) (\(x_3=1\), \(x_4=0\)) and \(b=(-2,1,0,1)'\) (\(x_3=0\), \(x_4=1\)). The pivot columns \(1\) and \(2\) of
\(\Sigma\), \((1,1,0,1)'\) and \((1,2,1,0)'\), form a basis of \(\mathcal C(\Sigma)\) (Theorem~\ref{thm:r4-pivot-columns}(b)). As
\(\Sigma'=\Sigma\), \(\mathcal N(\Sigma)=\mathcal C(\Sigma)^\perp\) (Corollary~\ref{r6-cor-four-subspaces}): every vector \(x\) of
\(\mathcal C(\Sigma)\) is orthogonal to \(a\) and \(b\), that is, \(x_1-x_2+x_3=0\) and \(-2x_1+x_2+x_4=0\); for the basis
vectors, \(1-1+0=0\), \(1-2+1=0\), \(-2+1+1=0\) and \(-2+2+0=0\).

(b) Let \(G\in\R^{2\times4}\) have the rows \(a'\) and \(b'\). Then \(G\Sigma=0\), as \(\Sigma a=\Sigma b=0\) and
\(\Sigma'=\Sigma\), and \(G\mu=(1-2+1+0,\ -2+2+0+0)'=(0,0)'\). By Theorem~\ref{r14-thm-mvn-linear},
\(Gz=(z_1-z_2+z_3,\ -2z_1+z_2+z_4)'\sim N_2(G\mu,G\Sigma G')=N_2(0,0)\), the distribution of the constant \(0\)
(Definition~\ref{r14-def-mvn}). So both equations hold with probability \(1\), in accordance with
Corollary~\ref{r14-cor-standard}(b): with probability \(1\), \(z-\mu\in\mathcal C(\Sigma)\), so \(G(z-\mu)=0\), and \(G\mu=0\).

(c)
\[
A=\begin{pmatrix}1&0\\1&1\\0&1\\1&-1\end{pmatrix}\quad\text{gives}\quad AA'=\begin{pmatrix}1&1&0&1\\1&2&1&0\\0&1&1&-1\\1&0&-1&2\end{pmatrix}=\Sigma
\]
so \(z\) has the
distribution of \(\mu+Aw\) (Theorem~\ref{r14-thm-mvn-iid}(c)): \(z_1=1+w_1\), \(z_2=2+w_1+w_2\), \(z_3=1+w_2\), \(z_4=w_1-w_2\);
and \(z_1-z_2+z_3=(1-2+1)+(w_1-w_1-w_2+w_2)=0\) and \(-2z_1+z_2+z_4=(-2+2+0)+(-2w_1+w_1+w_2+w_1-w_2)=0\) for every \(w\).
\end{solution}

\begin{exercise}\label{r14-dr-conditional}
Let \(z\sim N_4(\mu,\Sigma)\) with
\[
\mu=\begin{pmatrix}1\\0\\-1\\2\end{pmatrix},\qquad\Sigma=\begin{pmatrix}2&1&1&3\\1&1&2&1\\1&2&7&1\\3&1&1&8\end{pmatrix}
\]
\begin{enumerate}
\item Find the conditional distribution of \((z_3,z_4)'\) given \(z_1=a_1\) and \(z_2=a_2\), and for \(a_1=2\), \(a_2=1\).
\item With
\[
L=\begin{pmatrix}1&0&0&0\\0&1&0&0\\1&-3&1&0\\-2&1&0&1\end{pmatrix}
\]
compute \(L\Sigma\) and \(L\Sigma L'\).
\item Show that \((z_3+z_1-3z_2,\ z_4-2z_1+z_2)'\) is independent of \((z_1,z_2)'\), and find its distribution.
\end{enumerate}
\end{exercise}
\begin{solution}
(a) Here \(\Sigma_{11}=\left(\begin{smallmatrix}2&1\\1&1\end{smallmatrix}\right)\), with determinant \(1\) and, by
Theorem~\ref{r5d-thm-adjugate}, \(\Sigma_{11}^{-1}=\left(\begin{smallmatrix}1&-1\\-1&2\end{smallmatrix}\right)\); it is PD (its
leading principal minors are \(2\) and \(1\); Theorem~\ref{r8-thm-leading-minors}). Also
\(\Sigma_{21}=\left(\begin{smallmatrix}1&2\\3&1\end{smallmatrix}\right)\) and
\(\Sigma_{22}=\left(\begin{smallmatrix}7&1\\1&8\end{smallmatrix}\right)\). So
\[
\begin{gathered}
\Sigma_{21}\Sigma_{11}^{-1}=\begin{pmatrix}1-2&-1+4\\3-1&-3+2\end{pmatrix}=\begin{pmatrix}-1&3\\2&-1\end{pmatrix}
\\
\Sigma_{22\cdot1}=\begin{pmatrix}7&1\\1&8\end{pmatrix}-\begin{pmatrix}-1&3\\2&-1\end{pmatrix}\begin{pmatrix}1&3\\2&1\end{pmatrix}
=\begin{pmatrix}7&1\\1&8\end{pmatrix}-\begin{pmatrix}5&0\\0&5\end{pmatrix}=\begin{pmatrix}2&1\\1&3\end{pmatrix}
\end{gathered}
\]
By Theorem~\ref{r14-thm-mvn-conditional}, given \(z_1=a_1\), \(z_2=a_2\), the vector \((z_3,z_4)'\) has the distribution
\[
N_2\left(\begin{pmatrix}-1-(a_1-1)+3(a_2-0)\\2+2(a_1-1)-(a_2-0)\end{pmatrix},\ \Sigma_{22\cdot1}\right)
=N_2\left(\begin{pmatrix}-a_1+3a_2\\2a_1-a_2\end{pmatrix},\ \begin{pmatrix}2&1\\1&3\end{pmatrix}\right)
\]
For \(a_1=2\), \(a_2=1\) the conditional distribution is \(N_2\bigl((1,3)',\Sigma_{22\cdot1}\bigr)\).

(b) The rows \(3\) and \(4\) of \(L\Sigma\) are \(\row_1(\Sigma)-3\row_2(\Sigma)+\row_3(\Sigma)=(2-3+1,\ 1-3+2,\ 1-6+7,\ 3-3+1)\) and
\(-2\row_1(\Sigma)+\row_2(\Sigma)+\row_4(\Sigma)=(-4+1+3,\ -2+1+1,\ -2+2+1,\ -6+1+8)\) (Theorem~\ref{r1-thm-views}(c)), and then
the columns \(3\) and \(4\) of \(L\Sigma L'\) are \(\col_1(L\Sigma)-3\col_2(L\Sigma)+\col_3(L\Sigma)\) and
\(-2\col_1(L\Sigma)+\col_2(L\Sigma)+\col_4(L\Sigma)\):
\[
L\Sigma=\begin{pmatrix}2&1&1&3\\1&1&2&1\\0&0&2&1\\0&0&1&3\end{pmatrix},\qquad
L\Sigma L'=\begin{pmatrix}2&1&2-3+1&-4+1+3\\1&1&1-3+2&-2+1+1\\0&0&2&1\\0&0&1&3\end{pmatrix}=\begin{pmatrix}2&1&0&0\\1&1&0&0\\0&0&2&1\\0&0&1&3\end{pmatrix}
\]
as in Theorem~\ref{r5b-thm-schur}: the lower right block is the Schur complement \(\Sigma_{22\cdot1}\).

(c) The entries \(3\) and \(4\) of \(Lz\) are \(z_1-3z_2+z_3\) and \(-2z_1+z_2+z_4\), and \(Lz\sim N_4(L\mu,L\Sigma L')\)
(Theorem~\ref{r14-thm-mvn-linear}). The off-diagonal blocks of \(L\Sigma L'\) are \(0\), so
\((z_3+z_1-3z_2,\ z_4-2z_1+z_2)'\) is independent of \((z_1,z_2)'\) (Theorem~\ref{r14-thm-mvn-indep}(a)), and it has the
distribution \(N_2\bigl((1-0-1,\ -2+0+2)',\ \Sigma_{22\cdot1}\bigr)=N_2(0,\Sigma_{22\cdot1})\) (Corollary~\ref{r14-cor-marginal}(a)).
So \((z_3,z_4)'=(-z_1+3z_2,\ 2z_1-z_2)'+(z_3+z_1-3z_2,\ z_4-2z_1+z_2)'\): the conditional mean of (a) plus an independent
\(N_2(0,\Sigma_{22\cdot1})\) vector.
\end{solution}

\begin{exercise}\label{r14-dr-markov}
Let \(z\) be the random walk of Exercise~\ref{r14-dr-random-walk}.
\begin{enumerate}
\item Find the conditional distribution of \(z_4\) given \(z_1=a_1\), \(z_2=a_2\) and \(z_3=a_3\).
\item Find the conditional distribution of \(z_2\) given \(z_1=a_1\) and \(z_3=a_3\), and given \(z_1=a_1\), \(z_3=a_3\) and
\(z_4=a_4\).
\item Compare the conditional variance in (a) with the \((4,4)\) entry of \(\Sigma^{-1}\) found in Exercise~\ref{r14-dr-random-walk}.
\end{enumerate}
\end{exercise}
\begin{solution}
(a) Here \(\Sigma_{11}\) is the leading \(3\times3\) block of \(\Sigma\), which is PD, as \(\Sigma\) is
(Exercise~\ref{r14-dr-random-walk}(c), Theorem~\ref{r8-thm-principal-sub}), and
\(\Sigma_{21}=(1,2,3)\). As
\[
(0,0,1)\,\Sigma_{11}=(0,0,1)\begin{pmatrix}1&1&1\\1&2&2\\1&2&3\end{pmatrix}=(1,2,3)
\]
multiplying by \(\Sigma_{11}^{-1}\) on the right gives \(\Sigma_{21}\Sigma_{11}^{-1}=(0,0,1)\), and
\(\Sigma_{22\cdot1}=4-(0,0,1)(1,2,3)'=4-3=1\). By Theorem~\ref{r14-thm-mvn-conditional}, given \(z_1=a_1\), \(z_2=a_2\),
\(z_3=a_3\), the variable \(z_4\) is \(N(a_3,1)\): the next position is the present one plus an independent \(N(0,1)\) step,
whatever \(a_1\) and \(a_2\) are.

(b) Reorder the coordinates: \(y=(z_1,z_3,z_2)'=Pz\) for the matrix \(P\in\R^{3\times4}\) with the rows \(e_1',e_3',e_2'\), and
\(y\sim N_3(0,P\Sigma P')\) (Theorem~\ref{r14-thm-mvn-linear}), where \(P\Sigma P'\) has the rows and columns \(1,3,2\) of
\(\Sigma\):
\[
P\Sigma P'=\begin{pmatrix}1&1&1\\1&3&2\\1&2&2\end{pmatrix},\qquad
\Sigma_{11}=\begin{pmatrix}1&1\\1&3\end{pmatrix},\qquad\Sigma_{11}^{-1}=\frac12\begin{pmatrix}3&-1\\-1&1\end{pmatrix},\qquad\Sigma_{21}=(1,2)
\]
(Theorem~\ref{r5d-thm-adjugate}); this \(\Sigma_{11}\), and the one below, are principal submatrices of \(\Sigma\), hence PD, as
in (a) (Theorem~\ref{r8-thm-principal-sub}). So \(\Sigma_{21}\Sigma_{11}^{-1}=\frac12(3-2,\ -1+2)=(\frac12,\frac12)\) and
\(\Sigma_{22\cdot1}=2-(\frac12,\frac12)(1,2)'=2-\frac32=\frac12\). Given \(z_1=a_1\), \(z_3=a_3\), the variable \(z_2\) is
\(N\bigl(\frac12(a_1+a_3),\ \frac12\bigr)\): it is centered at the midpoint of its neighbours. Given also \(z_4=a_4\), take
\(y=(z_1,z_3,z_4,z_2)'\) with the rows \(e_1',e_3',e_4',e_2'\): now \(\Sigma_{21}=(1,2,2)\) and, as
\[
\bigl(\tfrac12,\tfrac12,0\bigr)\,\Sigma_{11}=\bigl(\tfrac12,\tfrac12,0\bigr)\begin{pmatrix}1&1&1\\1&3&3\\1&3&4\end{pmatrix}=(1,2,2)
\]
\(\Sigma_{21}\Sigma_{11}^{-1}=(\frac12,\frac12,0)\) and \(\Sigma_{22\cdot1}=2-\frac12-1=\frac12\): the same distribution
\(N\bigl(\frac12(a_1+a_3),\ \frac12\bigr)\), whatever \(a_4\) is.

(c) The \((4,4)\) entry of \(\Sigma^{-1}\) is \(1\), the inverse of the conditional variance \(1\) in (a), as
Theorem~\ref{r14-thm-mvn-conditional}(b) says.
\end{solution}

\needspace{16\baselineskip}
\begin{exercise}[cf.\ \citeat{haerdle-simar}{Example~5.3}]\label{r11-ex-partial-correlation}
Let \(x\sim N_4(\mu,\Sigma)\) with
\[
\Sigma=\begin{pmatrix}4&1&2&1\\1&4&2&1\\2&2&2&1\\1&1&1&1\end{pmatrix}
\]
\begin{enumerate}
\item Show that \(\Sigma\) is PD. Find the conditional distribution of \((x_1,x_2)'\) given \((x_3,x_4)'=(a_3,a_4)'\), and
compare the partial correlation of \(x_1\) and \(x_2\) given \(x_3\) and \(x_4\) (the paragraph after
Theorem~\ref{r14-thm-mvn-conditional}) with \(\rho_{12}\).
\item Find the partial correlation of \(x_1\) and \(x_2\) given \(x_3\) alone, and given \(x_4\) alone.
\item Let \(u=(x_1,x_2)'-B(x_3,x_4)'\), with \(B\) as in (a). Compute \(\var(u)\) from Theorem~\ref{r11-thm-linear}(b), and
compare with (a).
\end{enumerate}
\end{exercise}
\begin{solution}
(a) The leading principal minors of \(\Sigma\) are \(4\), \(15\), \(6\) and \(\det\Sigma=3\), all positive
(Theorem~\ref{r8-thm-leading-minors}). Apply Theorem~\ref{r14-thm-mvn-conditional} to the reordered vector
\((x_3,x_4,x_1,x_2)'\), which is normal (Theorem~\ref{r14-thm-mvn-linear} with the permutation matrix with the rows
\(e_3',e_4',e_1',e_2'\)), with the first block \((x_3,x_4)'\): \(\Sigma_{11}=\left(\begin{smallmatrix}2&1\\1&1\end{smallmatrix}\right)\),
which is PD as a principal submatrix of \(\Sigma\) (Theorem~\ref{r8-thm-principal-sub}),
\(\Sigma_{11}^{-1}=\left(\begin{smallmatrix}1&-1\\-1&2\end{smallmatrix}\right)\) (Proposition~\ref{r2-prop-inverse-2x2}),
\(\Sigma_{21}=\left(\begin{smallmatrix}2&1\\2&1\end{smallmatrix}\right)\) and
\[
B=\Sigma_{21}\Sigma_{11}^{-1}=\begin{pmatrix}1&0\\1&0\end{pmatrix},\qquad
\Sigma_{22\cdot1}=\begin{pmatrix}4&1\\1&4\end{pmatrix}-\begin{pmatrix}1&0\\1&0\end{pmatrix}\begin{pmatrix}2&2\\1&1\end{pmatrix}
=\begin{pmatrix}2&-1\\-1&2\end{pmatrix}
\]
So given \((x_3,x_4)'=(a_3,a_4)'\), the vector \((x_1,x_2)'\) has the distribution
\(N_2\bigl((\mu_1,\mu_2)'+(a_3-\mu_3)(1,1)',\ \Sigma_{22\cdot1}\bigr)\). The partial correlation is \(-1/\sqrt{2\cdot2}=-\frac12\),
while \(\rho_{12}=1/\sqrt{4\cdot4}=\frac14\): the sign changes.

(b) Given \(x_3\) (the vector \((x_3,x_1,x_2)'\), \(\Sigma_{11}=(2)\)):
\(\left(\begin{smallmatrix}4&1\\1&4\end{smallmatrix}\right)-\frac12(2,2)'(2,2)=\left(\begin{smallmatrix}4&1\\1&4\end{smallmatrix}\right)
-\left(\begin{smallmatrix}2&2\\2&2\end{smallmatrix}\right)=\left(\begin{smallmatrix}2&-1\\-1&2\end{smallmatrix}\right)\), partial
correlation \(-\frac12\). Given \(x_4\) (\(\Sigma_{11}=(1)\)):
\(\left(\begin{smallmatrix}4&1\\1&4\end{smallmatrix}\right)-(1,1)'(1,1)=\left(\begin{smallmatrix}3&0\\0&3\end{smallmatrix}\right)\),
partial correlation \(0\).

(c) \(u=(x_1-x_3,\ x_2-x_3)'=Kx\) with \(K=\left(\begin{smallmatrix}1&0&-1&0\\0&1&-1&0\end{smallmatrix}\right)\). The rows
of \(K\Sigma\) are \((2,-1,0,0)\) and \((-1,2,0,0)\), so
\(\var(u)=K\Sigma K'=\left(\begin{smallmatrix}2&-1\\-1&2\end{smallmatrix}\right)=\Sigma_{22\cdot1}\). This agrees with
Theorem~\ref{r14-thm-mvn-conditional}(a): \(u\) differs from the vector \(w\) of that theorem (for \((x_3,x_4,x_1,x_2)'\)) by the
constant vector \((\mu_1,\mu_2)'-B(\mu_3,\mu_4)'\).
\end{solution}

\begin{exercise}\label{r14-dr-bivariate}
Let \((u,v)'\) be bivariate normal with \(\mu_u=2\), \(\mu_v=3\), \(\sigma_u=2\), \(\sigma_v=5\) and \(\rho=\frac35\). Let \(\Phi\) be the
distribution function of \(N(0,1)\): \(\Phi(c)=\prb(y\leqslant c)\) for \(y\sim N(0,1)\).
\begin{enumerate}
\item Find \(\Sigma\), \(\Sigma^{-1}\) and the density of \((u,v)'\).
\item Find the conditional distributions of \(v\) given \(u\) and of \(u\) given \(v\).
\item Express \(\prb(v>10)\) and the conditional probability \(\prb(v>10)\) given \(u=4\) through \(\Phi\). Which is larger?
\end{enumerate}
\end{exercise}
\begin{solution}
(a) \(\rho\sigma_u\sigma_v=\frac35\cdot10=6\), so
\[
\Sigma=\begin{pmatrix}4&6\\6&25\end{pmatrix},\qquad\det\Sigma=100-36=64,\qquad\Sigma^{-1}=\frac1{64}\begin{pmatrix}25&-6\\-6&4\end{pmatrix}
\]
(Theorem~\ref{r5d-thm-adjugate}). As \(2\pi\sigma_u\sigma_v\sqrt{1-\rho^2}=2\pi\cdot10\cdot\frac45=16\pi\), Proposition~\ref{r14-prop-bivariate}(c)
gives
\[
f(s,t)=\frac1{16\pi}\exp\Bigl(-\frac1{128}\bigl(25(s-2)^2-12(s-2)(t-3)+4(t-3)^2\bigr)\Bigr)
\]
(b) By Proposition~\ref{r14-prop-bivariate}(d), given \(u\), \(v\sim N\bigl(3+\frac35\cdot\frac52(u-2),\ 25\cdot\frac{16}{25}\bigr)=N\bigl(3+\frac32(u-2),\ 16\bigr)\);
with the roles exchanged, given \(v\), \(u\sim N\bigl(2+\frac35\cdot\frac25(v-3),\ 4\cdot\frac{16}{25}\bigr)=N\bigl(2+\frac6{25}(v-3),\ \frac{64}{25}\bigr)\).

(c) If \(x\sim N(m,s^2)\) with \(s>0\), then \((x-m)/s\sim N(0,1)\) (Theorem~\ref{r14-thm-mvn-linear} with \(A=1/s\), \(b=-m/s\)), and
\(\prb(x>c)=1-\Phi\bigl((c-m)/s\bigr)\). As \(v\sim N(3,25)\), \(\prb(v>10)=1-\Phi(\frac{10-3}5)=1-\Phi(\frac75)\). Given \(u=4\), \(v\) has
the distribution \(N(3+3,16)=N(6,16)\), so the conditional probability is \(1-\Phi(\frac{10-6}4)=1-\Phi(1)\). It is the larger
one: \(\Phi(\frac75)-\Phi(1)=\prb(1<y\leqslant\frac75)>0\), as \(N(0,1)\) has a positive density
(Theorem~\ref{r14-thm-mvn-density}). Knowing that \(u=4\), one standard deviation of \(u\) above its mean, makes the probability
larger.
\end{solution}

\begin{exercise}\label{r14-dr-mle}
The rows of the data matrix
\[
X=\begin{pmatrix}6&3&4&1\\2&7&2&1\\4&5&5&4\\6&7&2&3\\2&3&2&1\end{pmatrix}
\]
are a random sample of size \(n=5\) from \(N_4(\mu,\Sigma)\).
\begin{enumerate}
\item Find \(\bar x\), \(C_5X\), \(A=(C_5X)'(C_5X)\) and \(S\).
\item Find the maximum likelihood estimates \(\hat\mu\), \(\hat\Sigma\), and \(\det\hat\Sigma\).
\item Using the proof of Theorem~\ref{r14-thm-mle}, compute \(\log L(\bar x,\hat\Sigma)-\log L(\bar x,S)\).
\end{enumerate}
\end{exercise}
\begin{solution}
(a) \(\one'X=(20,25,15,10)\), so \(\bar x=(4,5,3,2)'\), and, subtracting \((4,5,3,2)\) from every row,
\[
C_5X=\begin{pmatrix}2&-2&1&-1\\-2&2&-1&-1\\0&0&2&2\\2&2&-1&1\\-2&-2&-1&-1\end{pmatrix},\qquad
A=\begin{pmatrix}16&0&4&4\\0&16&-4&4\\4&-4&8&4\\4&4&4&8\end{pmatrix},\qquad
S=\frac14A=\begin{pmatrix}4&0&1&1\\0&4&-1&1\\1&-1&2&1\\1&1&1&2\end{pmatrix}
\]
(for instance \(16=4+4+0+4+4\), \(4=2+2+0-2+2\) and \(8=1+1+4+1+1\)).

(b) By Theorem~\ref{r5d-thm-schur-det}(a) with \(S_{11}=4I_2\), \(S_{21}=\left(\begin{smallmatrix}1&-1\\1&1\end{smallmatrix}\right)=S_{12}'\)
and \(S_{22}=\left(\begin{smallmatrix}2&1\\1&2\end{smallmatrix}\right)\),
\[
S_{22}-S_{21}S_{11}^{-1}S_{12}=\begin{pmatrix}2&1\\1&2\end{pmatrix}-\frac14\begin{pmatrix}2&0\\0&2\end{pmatrix}
=\begin{pmatrix}\frac32&1\\1&\frac32\end{pmatrix},\qquad\det S=\det(4I_2)\cdot\Bigl(\frac94-1\Bigr)=16\cdot\frac54=20
\]
So \(S\), which is PSD (Proposition~\ref{r12-prop-data}(a)), is invertible (Theorem~\ref{r5d-thm-product}(a)) and hence PD
(Theorem~\ref{r8-thm-eigen-criterion}(c)), and so is \(A=4S\). By Theorem~\ref{r14-thm-mle},
\[
\hat\mu=\bar x=\begin{pmatrix}4\\5\\3\\2\end{pmatrix},\qquad\hat\Sigma=\frac15A=\frac45S,
\qquad\det\hat\Sigma=\Bigl(\frac45\Bigr)^4\det S=\frac{256}{625}\cdot20=\frac{1024}{125}
\]
(Theorem~\ref{r5d-thm-product}(c)).

(c) In the proof of Theorem~\ref{r14-thm-mle}, \(\log L(\bar x,\hat\Sigma)-\log L(\bar x,S)=\frac12\bigl(g(S)-g(\hat\Sigma)\bigr)=\frac n2(\tr B-\log\det B-p)\)
with \(B=\hat\Sigma^{1/2}S^{-1}\hat\Sigma^{1/2}\). As \(S=\frac54\hat\Sigma\), \(S^{-1}=\frac45\hat\Sigma^{-1}\), and
\(\hat\Sigma^{-1}=\hat\Sigma^{-1/2}\hat\Sigma^{-1/2}\) (Theorem~\ref{r8-thm-sqrt}), so \(B=\frac45I_4\), with \(\tr B=\frac{16}5\) and
\(\det B=\frac{256}{625}\). So the difference is
\[
\frac52\Bigl(\frac{16}5-\log\frac{256}{625}-4\Bigr)=10\log\frac54-2
\]
It is positive, as Lemma~\ref{r14-lem-logdet} requires: \(\bigl(\frac54\bigr)^5=\frac{3125}{1024}>3>e\), so \(5\log\frac54>1\).
\end{solution}

\begin{exercise}\label{r11-ex-mle-drill}
The five observations \((2,4,1,1)\), \((2,2,3,1)\), \((2,6,5,3)\), \((5,4,3,4)\), \((4,4,3,1)\) are modeled as a random sample from
\(N_4(\mu,\Sigma)\).
\begin{enumerate}
\item Find the maximum likelihood estimates \(\hat\mu\), \(\hat\Sigma\), the largest value of \(\log L\), and the correlation
matrix computed from \(\hat\Sigma\).
\item Compute \(\log L(\bar x,S)\) for the unbiased estimate \(S\), and check that it is smaller than \(\log L(\bar x,\hat\Sigma)\).
\item Find the eigenvalues \(\theta_1,\dots,\theta_4\) of \(B=\hat\Sigma^{1/2}S^{-1}\hat\Sigma^{1/2}\), and check the identity
\(g(S)-g(\hat\Sigma)=n(\tr B-\log\det B-p)\) of the proof of Theorem~\ref{r14-thm-mle}, that is,
\(n\sum_j(\theta_j-1-\log\theta_j)\) (Theorem~\ref{r7-thm-trace-det}), where
\(g(\Sigma)=n\log\det\Sigma+\tr(\Sigma^{-1}A)\).
\end{enumerate}
\end{exercise}
\begin{solution}
(a) \(\bar x=\frac15(15,20,15,10)'=(3,4,3,2)'\), and the deviations \(x_i-\bar x\) are \((-1,0,-2,-1)\), \((-1,-2,0,-1)\),
\((-1,2,2,1)\), \((2,0,0,2)\), \((1,0,0,-1)\). So
\[
A=\sum_{i=1}^5(x_i-\bar x)(x_i-\bar x)'=\begin{pmatrix}8&0&0&4\\0&8&4&4\\0&4&8&4\\4&4&4&8\end{pmatrix},\qquad
S=\frac14A=\begin{pmatrix}2&0&0&1\\0&2&1&1\\0&1&2&1\\1&1&1&2\end{pmatrix}
\]
(for instance \(8=1+1+1+4+1\) and \(4=1+1-1+4-1\)). By Theorem~\ref{r5d-thm-schur-det}(a) with \(S_{11}=2I_2\),
\(S_{21}=\left(\begin{smallmatrix}0&1\\1&1\end{smallmatrix}\right)=S_{12}'\) and \(S_{22}=\left(\begin{smallmatrix}2&1\\1&2\end{smallmatrix}\right)\),
\[
S_{22}-S_{21}S_{11}^{-1}S_{12}=\begin{pmatrix}2&1\\1&2\end{pmatrix}-\frac12\begin{pmatrix}1&1\\1&2\end{pmatrix}
=\begin{pmatrix}\frac32&\frac12\\\frac12&1\end{pmatrix},\qquad\det S=\det(2I_2)\cdot\Bigl(\frac32-\frac14\Bigr)=4\cdot\frac54=5
\]
so \(S\), which is PSD (Proposition~\ref{r12-prop-data}(a)), is invertible (Theorem~\ref{r5d-thm-product}(a)) and hence PD
(Theorem~\ref{r8-thm-eigen-criterion}(c)), and so is \(A=4S\). By Theorem~\ref{r14-thm-mle}, \(\hat\mu=\bar x=(3,4,3,2)'\),
\(\hat\Sigma=\frac15A=\frac45S\), \(\det\hat\Sigma=\bigl(\frac45\bigr)^4\cdot5=\frac{256}{125}\) (Theorem~\ref{r5d-thm-product}(c)), and the
largest value of the log-likelihood is
\[
\log L(\bar x,\hat\Sigma)=-10\log(2\pi)-\frac52\log\frac{256}{125}-10
\]
All diagonal entries of \(\hat\Sigma\) are \(\frac85\), so the correlation matrix computed from \(\hat\Sigma\)
(Definition~\ref{r14-def-corr}) is \(\frac58\hat\Sigma=\frac18A\):
\[
\frac18A=\begin{pmatrix}1&0&0&\frac12\\[2pt]0&1&\frac12&\frac12\\[2pt]0&\frac12&1&\frac12\\[2pt]\frac12&\frac12&\frac12&1\end{pmatrix}
\]
the sample correlation matrix of the data (Definition~\ref{r12-def-correlation}: the diagonal entries of \(S\) are \(2\),
and \(\frac12S=\frac18A\)).

(b) In the proof of Theorem~\ref{r14-thm-mle}, \(\log L(\bar x,\Sigma)=-10\log(2\pi)-\frac12g(\Sigma)\), with \(np=20\). Here
\(g(S)=5\log\det S+\tr(S^{-1}A)=5\log5+\tr(4I_4)=5\log5+16\), so \(\log L(\bar x,S)=-10\log(2\pi)-\frac52\log5-8\). Hence
\[
\log L(\bar x,\hat\Sigma)-\log L(\bar x,S)=\frac52\Bigl(\log5-\log\frac{256}{125}\Bigr)-2=\frac52\log\frac{625}{256}-2=10\log\frac54-2
\]
This is positive: \(\bigl(\frac54\bigr)^5=\frac{3125}{1024}>3>e\), so \(5\log\frac54>1\).

(c) \(S^{-1}=\frac45\hat\Sigma^{-1}\) and \(\hat\Sigma^{-1}=\hat\Sigma^{-1/2}\hat\Sigma^{-1/2}\) (Theorem~\ref{r8-thm-sqrt}), so
\(B=\frac45\hat\Sigma^{1/2}\hat\Sigma^{-1}\hat\Sigma^{1/2}=\frac45I_4\) and \(\theta_1=\dots=\theta_4=\frac45\). Then
\(n\sum_j(\theta_j-1-\log\theta_j)=5\cdot4\bigl(\frac45-1+\log\frac54\bigr)=20\log\frac54-4\), and
\(g(S)-g(\hat\Sigma)=5\log5+16-\bigl(5\log\frac{256}{125}+20\bigr)=5\log\frac{625}{256}-4=20\log\frac54-4\), the same number, as
\(g(\hat\Sigma)=5\log\det\hat\Sigma+\tr(5I_4)\); it is \(2\bigl(\log L(\bar x,\hat\Sigma)-\log L(\bar x,S)\bigr)\), as in the proof
of Theorem~\ref{r14-thm-mle}. The difference of (b) is the number of Exercise~\ref{r14-dr-mle}(c), for other data: for all
data, \(S=\frac n{n-1}\hat\Sigma\), so \(B=\frac{n-1}nI_p\), and the difference depends only on \(n=5\) and \(p=4\).
\end{solution}

\needspace{14\baselineskip}
\subsection*{Problems}

\begin{problem}[\citeat{casella-berger}{Exercise~4.42}, \citeat{casella-berger-solutions}{Exercise~4.42}]\label{r14-ex-corr-product}
Let \(u\) and \(v\) be independent random variables with means \(\mu_u,\mu_v\) and positive variances \(\sigma_u^2,\sigma_v^2\). Find
the correlation of \(uv\) and \(v\) in terms of these means and variances.
\end{problem}
\begin{solution}
\begin{itemize}
\item By (N3), \(u\) and \(v^2\) are independent, and so are \(u^2\) and \(v^2\); by (N7), \(\expc(uv^2)=\mu_u\expc(v^2)\),
\(\expc(uv)=\mu_u\mu_v\) and \(\expc(u^2v^2)=\expc(u^2)\expc(v^2)\), where
\(\expc(u^2)=\sigma_u^2+\mu_u^2\) and \(\expc(v^2)=\sigma_v^2+\mu_v^2\).
\item The entries \((1,2)\) and \((1,1)\) of Theorem~\ref{r14-thm-cov}(e) for the random vector \((uv,v)'\) read
\(\cov(uv,v)=\expc(uv^2)-\expc(uv)\expc(v)\) and \(\var(uv)=\expc(u^2v^2)-\bigl(\expc(uv)\bigr)^2\). So
\(\cov(uv,v)=\mu_u(\sigma_v^2+\mu_v^2)-\mu_u\mu_v^2=\mu_u\sigma_v^2\).
\item \(\var(uv)=(\sigma_u^2+\mu_u^2)(\sigma_v^2+\mu_v^2)-\mu_u^2\mu_v^2
=\sigma_u^2\sigma_v^2+\sigma_u^2\mu_v^2+\sigma_v^2\mu_u^2>0\).
\item Hence the correlation (Definition~\ref{r14-def-corr}) is
\[
\frac{\cov(uv,v)}{\sqrt{\var(uv)\var(v)}}
=\frac{\mu_u\sigma_v^2}{\sigma_v\sqrt{\sigma_u^2\sigma_v^2+\sigma_u^2\mu_v^2+\sigma_v^2\mu_u^2}}
=\frac{\mu_u\sigma_v}{\sqrt{\mu_u^2\sigma_v^2+\mu_v^2\sigma_u^2+\sigma_u^2\sigma_v^2}}
\]
\end{itemize}
\end{solution}

\begin{problem}[\citeat{casella-berger}{Exercise~4.46(b),(d)}, \citeat{casella-berger-solutions}{Exercise~4.46}]\label{r14-ex-two-factors}
Let \(w_1,w_2\) be independent \(N(0,1)\) variables, and for constants \(a_u,b_u,c_u,a_v,b_v,c_v\) let
\(u=a_uw_1+b_uw_2+c_u\) and \(v=a_vw_1+b_vw_2+c_v\). Let \(\mu_u,\mu_v\in\R\), \(\sigma_u,\sigma_v>0\) and \(-1\leqslant\rho\leqslant1\).
\begin{enumerate}
\item Let \(a_u=\sqrt{\frac{1+\rho}2}\,\sigma_u\), \(b_u=\sqrt{\frac{1-\rho}2}\,\sigma_u\), \(c_u=\mu_u\), \(a_v=\sqrt{\frac{1+\rho}2}\,\sigma_v\),
\(b_v=-\sqrt{\frac{1-\rho}2}\,\sigma_v\) and \(c_v=\mu_v\). Show that \(\expc(u)=\mu_u\), \(\var(u)=\sigma_u^2\),
\(\expc(v)=\mu_v\), \(\var(v)=\sigma_v^2\), and that the correlation of \(u\) and \(v\) is \(\rho\).
\item Show that the constants in (a) are not the only ones with these properties, by giving another choice. How many
choices are there?
\end{enumerate}
\end{problem}
\begin{solution}
Here \((u,v)'=Aw+c\) with \(A=\left(\begin{smallmatrix}a_u&b_u\\a_v&b_v\end{smallmatrix}\right)\), \(w=(w_1,w_2)'\) and \(c=(c_u,c_v)'\). As \(w\) has the mean \(0\)
and the covariance matrix \(I_2\) (Theorem~\ref{r14-thm-mvn-iid}(a), Definition~\ref{r14-def-mvn}), Theorem~\ref{r11-thm-linear}(b) gives
the mean \(c\) and the covariance matrix
\[
AA'=\begin{pmatrix}a_u^2+b_u^2&a_ua_v+b_ub_v\\a_ua_v+b_ub_v&a_v^2+b_v^2\end{pmatrix}
\]
So the requirements are \(c_u=\mu_u\), \(c_v=\mu_v\), \(a_u^2+b_u^2=\sigma_u^2\), \(a_v^2+b_v^2=\sigma_v^2\) and
\(a_ua_v+b_ub_v=\rho\sigma_u\sigma_v\); and then \((u,v)'\) is bivariate normal with these parameters
(Theorem~\ref{r14-thm-mvn-iid}(c), Definition~\ref{r14-def-bivariate}).

(a) \(a_u^2+b_u^2=\bigl(\frac{1+\rho}2+\frac{1-\rho}2\bigr)\sigma_u^2=\sigma_u^2\), likewise \(a_v^2+b_v^2=\sigma_v^2\), and
\(a_ua_v+b_ub_v=\bigl(\frac{1+\rho}2-\frac{1-\rho}2\bigr)\sigma_u\sigma_v=\rho\sigma_u\sigma_v\); so the correlation is
\(\rho\sigma_u\sigma_v/(\sigma_u\sigma_v)=\rho\).

(b)
\begin{itemize}
\item Another choice is \(a_u=\rho\sigma_u\), \(b_u=\sqrt{1-\rho^2}\,\sigma_u\), \(a_v=\sigma_v\), \(b_v=0\) (with \(c_u=\mu_u\),
\(c_v=\mu_v\)): \(a_u^2+b_u^2=\sigma_u^2\), \(a_v^2+b_v^2=\sigma_v^2\) and \(a_ua_v+b_ub_v=\rho\sigma_u\sigma_v\).
It differs from (a) if \(\rho\ne1\); for \(\rho=1\), take \(a_u=a_v=0\), \(b_u=\sigma_u\), \(b_v=\sigma_v\).
\item There are infinitely many choices. If the three equations hold, then
\((\rho\sigma_u\sigma_v-a_ua_v)^2=b_u^2b_v^2=(\sigma_u^2-a_u^2)(\sigma_v^2-a_v^2)\), that is,
\[
\sigma_u^2a_v^2-2\rho\sigma_u\sigma_va_ua_v+\sigma_v^2a_u^2=(1-\rho^2)\sigma_u^2\sigma_v^2
\]
an ellipse in the \((a_u,a_v)\)-plane if \(\lvert\rho\rvert<1\), and the line \(\sigma_ua_v=\rho\sigma_va_u\) if \(\rho=\pm1\).
\item Conversely, for every \(a_u\) with \(\lvert a_u\rvert<\sigma_u\) this quadratic equation in \(a_v\) has a real solution, as its
discriminant \(4\rho^2\sigma_u^2\sigma_v^2a_u^2-4\sigma_u^2\bigl(\sigma_v^2a_u^2-(1-\rho^2)\sigma_u^2\sigma_v^2\bigr)
=4(1-\rho^2)\sigma_u^2\sigma_v^2(\sigma_u^2-a_u^2)\) is \(\geqslant0\). Then \(b_u=\sqrt{\sigma_u^2-a_u^2}>0\) and
\(b_v=(\rho\sigma_u\sigma_v-a_ua_v)/b_u\) satisfy the three equations: \(a_ua_v+b_ub_v=\rho\sigma_u\sigma_v\), and
\(b_v^2=(\sigma_u^2-a_u^2)(\sigma_v^2-a_v^2)/b_u^2=\sigma_v^2-a_v^2\) by the displayed equation. So every \(a_u\in(-\sigma_u,\sigma_u)\)
gives a choice.
\end{itemize}
\end{solution}

\begin{problem}[\citeat{bapat}{Ch.~8, Exercise~14}]\label{r14-ex-rank-factor}
Let \(y\sim N_n(0,\Sigma)\), where \(\Sigma\) is PSD of rank \(r\geqslant1\). Show that there are \(G\in\R^{n\times r}\) with \(\Sigma=GG'\) and
\(x\sim N_r(0,I_r)\) such that \(y=Gx\) with probability \(1\).
\end{problem}
\begin{solution}
\begin{itemize}
\item By Lemma~\ref{r9-lem-eig-order}(a) and Theorem~\ref{r8-thm-rank-eigen},
\(\Sigma=P\left(\begin{smallmatrix}D&0\\0&0\end{smallmatrix}\right)P'\) with \(P\) orthogonal and
\(D=\diag(\lambda_1,\dots,\lambda_r)\), \(\lambda_1\geqslant\dots\geqslant\lambda_r>0\); the other eigenvalues are \(0\), as \(\Sigma\) is PSD
(Theorem~\ref{r8-thm-eigen-criterion}(a)). Write \(P=[P_1\ P_2]\) with \(P_1\in\R^{n\times r}\), and put
\(G=P\left(\begin{smallmatrix}D^{1/2}\\0\end{smallmatrix}\right)=P_1D^{1/2}\), where
\(D^{1/2}=\diag(\sqrt{\lambda_1},\dots,\sqrt{\lambda_r})\).
\item \(GG'=P\left(\begin{smallmatrix}D^{1/2}\\0\end{smallmatrix}\right)\bigl(D^{1/2}\ \ 0\bigr)P'
=P\left(\begin{smallmatrix}D&0\\0&0\end{smallmatrix}\right)P'=\Sigma\) (Theorem~\ref{r5b-thm-block}(c)). So, if \(x\sim N_r(0,I_r)\), then
\(Gx\sim N_n(0,GG')=N_n(0,\Sigma)\) (Theorem~\ref{r14-thm-mvn-linear}).
\item To obtain \(y=Gx\), let \(x=D^{-1/2}P_1'y\). The blocks of
\(P'\Sigma P=\left(\begin{smallmatrix}D&0\\0&0\end{smallmatrix}\right)\) are \(P_1'\Sigma P_1=D\), \(P_1'\Sigma P_2=0\) and \(P_2'\Sigma P_2=0\)
(Theorem~\ref{r5b-thm-block}(b),(c)). By Theorem~\ref{r14-thm-mvn-linear},
\(x\sim N_r(0,D^{-1/2}P_1'\Sigma P_1D^{-1/2})=N_r(0,I_r)\).
\item \(Gx=P_1P_1'y=y-P_2P_2'y\), as \(P_1P_1'+P_2P_2'=PP'=I_n\) (Theorem~\ref{r5b-thm-block}(c)). Finally
\(P_2'y\sim N_{n-r}(0,P_2'\Sigma P_2)=N_{n-r}(0,0)\) (Theorem~\ref{r14-thm-mvn-linear}), the distribution of the constant \(0\)
(Definition~\ref{r14-def-mvn}); so \(y=Gx\) with probability \(1\). (If \(r=n\), there is no \(P_2\), and \(y=Gx\).)
\end{itemize}
\end{solution}

\begin{problem}[\citeat{haerdle-simar}{Theorems~4.9, 5.7(d), \S7.1, (7.2)}]\label{r11-pr-mean-cov-independent}
Let \(x_1,\dots,x_n\) (\(n\geqslant2\)) be independent random vectors, each \(N_p(\mu,\Sigma)\), let \(X\) be the data matrix with the
rows \(x_1',\dots,x_n'\), and let \(\bar x\) and \(S\) be its mean vector and sample covariance matrix
(Definition~\ref{r12-def-data}).
\begin{enumerate}
\item Show that \(\cov(\bar x,\ x_i-\bar x)=0\) for every \(i\).
\item Show that \(\bar x\sim N_p(\mu,\frac1n\Sigma)\) and that \(\bar x\) and \(S\) are independent.
\item Let \(\Sigma\) be PD. Show that \(n(\bar x-\mu)'\Sigma^{-1}(\bar x-\mu)\sim\chi^2_p\), independently of \(S\).
\end{enumerate}
\end{problem}
\begin{solution}
(a) Let \(u=(x_1',\dots,x_n')'\in\R^{np}\). By Theorem~\ref{r14-thm-mvn-iid}(a), \(u\sim N_{np}(\mu_u,V)\) with
\(\mu_u=(\mu',\dots,\mu')'\) and \(V=\diag(\Sigma,\dots,\Sigma)\). Let \(G=[I_p\ \cdots\ I_p]\in\R^{p\times np}\), and let
\(E_i\in\R^{p\times np}\) have the block \(I_p\) in position \(i\) and zero blocks elsewhere; so \(\bar x=\frac1nGu\)
(Definition~\ref{r12-def-data}: \(\bar x=\frac1nX'\one=\frac1n(x_1+\dots+x_n)\)) and \(x_i=E_iu\). By the block product
(Theorem~\ref{r5b-thm-block}(b),(c)), \(GVE_i'=\Sigma\) and \(GVG'=n\Sigma\). By Theorem~\ref{r14-thm-cov}(a),
\[
\cov\bigl(\bar x,\ x_i-\bar x\bigr)=\frac1nGV\Bigl(E_i-\frac1nG\Bigr)'=\frac1n\Sigma-\frac1{n^2}n\Sigma=0
\]
(b) \(G\mu_u=n\mu\), so \(\bar x=\frac1nGu\sim N_p(\mu,\frac1{n^2}GVG')=N_p(\mu,\frac1n\Sigma)\) by
Theorem~\ref{r14-thm-mvn-linear}. Let \(K\in\R^{np\times np}\) have the block rows \(E_1-\frac1nG,\dots,E_n-\frac1nG\), so that
\(Ku=\bigl((x_1-\bar x)',\dots,(x_n-\bar x)'\bigr)'\). The matrix \(\frac1nGVK'\) has the blocks \(\frac1nGV(E_i-\frac1nG)'\), which
are \(0\) by (a); so \(\bar x\) and \(Ku\) are independent by Theorem~\ref{r14-thm-mvn-indep}(c), with \(A=\frac1nG\) and \(B=K\). As
\(S=\frac1{n-1}\sum_i(x_i-\bar x)(x_i-\bar x)'\) (Proposition~\ref{r12-prop-data}(a)) is a function of \(Ku\), \(\bar x\) and \(S\)
are independent (N3).

(c) \(\frac1n\Sigma\) is PD (\(y'(\frac1n\Sigma)y=\frac1ny'\Sigma y\)), with the inverse \(n\Sigma^{-1}\), so
Corollary~\ref{r14-cor-standard}(a) applied to \(\bar x\sim N_p(\mu,\frac1n\Sigma)\) gives
\((\bar x-\mu)'(n\Sigma^{-1})(\bar x-\mu)\sim\chi^2_p\). The statistic is a function of \(\bar x\), hence independent of \(S\) by
(b) and (N3).
\end{solution}

\clearpage
\def\unitnumber{11}
\setcounter{theorem}{0}\setcounter{problem}{0}
\section*{Chapter 11. Quadratic forms in normal vectors: chi-square, \(t\), and \(F\)}

\subsection*{The chi-square distribution}

This chapter uses the facts (N1)--(N7) stated after Definition~\ref{r14-def-mvn}, the facts on densities stated in the
paragraph before Theorem~\ref{r14-thm-mvn-density}, and marginal densities as in Definition~\ref{r14-def-conditional}. In
addition, it uses the following facts without proof: (D1) and (D2) from calculus, (D3) from probability~\citeat{casella-berger}{(4.3.2), (4.6.7), \S3.3, Theorem~1.5.10}. Integrals over unbounded intervals are improper
integrals, and a function \(g\) is \textbf{integrable} over a set \(\Omega\) if \(\int_\Omega\lvert g\rvert<\infty\).
\begin{enumerate}[label=(D\arabic*)]
\item (\textbf{Change of variables}) Let \(\Omega,\Omega'\subseteq\R^m\) be open, and let \(\psi\colon\Omega'\to\Omega\) be
one-to-one and onto, with continuous partial derivatives and Jacobian determinant
\(J(v)=\det\bigl(\partial\psi_i(v)/\partial v_j\bigr)\ne0\) for \(v\in\Omega'\). Then
\(\int_\Omega g(x)\,\mathrm{d}x=\int_{\Omega'}g(\psi(v))\lvert J(v)\rvert\,\mathrm{d}v\) for every \(g\) integrable over \(\Omega\). Consequently, if \(u\)
has a density \(f\) that is \(0\) outside \(\Omega\), then \(\psi^{-1}(u)\) has the density \(f(\psi(v))\lvert J(v)\rvert\) for
\(v\in\Omega'\), and \(0\) outside \(\Omega'\). For \(m=1\) this is the substitution rule of calculus; for an affine map
\(\psi(v)=c+Bv\) of \(\R^m\) onto \(\R^m\) it is the change of variables before Theorem~\ref{r14-thm-mvn-density}.
\item (\textbf{The Gamma function}) For \(c>0\), \(\Gamma(c)=\int_0^\infty x^{c-1}e^{-x}\,\mathrm{d}x\) is finite and positive,
\(\Gamma(c+1)=c\,\Gamma(c)\) and \(\Gamma(1)=1\). Substituting \(x=by\) (D1) gives
\(\int_0^\infty y^{c-1}e^{-by}\,\mathrm{d}y=\Gamma(c)\,b^{-c}\) for every \(b>0\). Substituting \(x=s^2/2\) (\(s>0\)) gives
\(\Gamma(\frac12)=\sqrt2\int_0^\infty e^{-s^2/2}\,\mathrm{d}s=\sqrt\pi\), as the integrand is even and
\(\int_\R e^{-s^2/2}\,\mathrm{d}s=\sqrt{2\pi}\) (the paragraph before Theorem~\ref{r14-thm-mvn-density}).
\item (\textbf{Distribution functions}) If random variables \(u\) and \(v\) satisfy
\(\prb(u\leqslant c)=\prb(v\leqslant c)\) for every \(c\in\R\), then \(u\) and \(v\) have the same
distribution.
\end{enumerate}
By Theorem~\ref{r14-thm-mvn-density} (the case \(n=1\)), \(N(0,1)\) has the density \(\varphi(x)=(2\pi)^{-1/2}e^{-x^2/2}\). An
identity between MGFs (defined before Definition~\ref{r14-def-mvn}) is needed only \textbf{near \(0\)}, that is, for all \(t\) in
some open interval (or ball) containing \(0\), as in (N1) and (N2).

\begin{lemma}[cf.\ \citeat{casella-berger}{Example~2.3.8}]\label{r15-lem-mgf}
\begin{enumerate}
\item Let \(w\sim N(0,1)\) and \(\delta\in\R\). For every \(t<\frac12\),
\[
\expc\bigl(e^{t(w+\delta)^2}\bigr)=(1-2t)^{-1/2}\exp\Bigl(\frac{\delta^2t}{1-2t}\Bigr)
\]
\item Let \(w\sim N_r(0,I_r)\) and \(\delta=(\delta_1,\dots,\delta_r)'\in\R^r\), and put
\(\lambda=\lVert\delta\rVert^2=\delta_1^2+\dots+\delta_r^2\). Then, near \(0\),
\[
M_{\lVert w+\delta\rVert^2}(t)=(1-2t)^{-r/2}\exp\Bigl(\frac{\lambda t}{1-2t}\Bigr)
\]
Hence the distribution of \(\lVert w+\delta\rVert^2\) depends on \(\delta\) only through \(\lambda\).
\end{enumerate}
\end{lemma}
\begin{proof}
(a)
\begin{itemize}
\item \(w\) has the density \(\varphi\), so the expectation is \(\int_\R e^{t(x+\delta)^2}\varphi(x)\,\mathrm{d}x\) (the paragraph before
Theorem~\ref{r14-thm-mvn-density}), provided this integral is finite.
\item Put \(c=1-2t>0\) and \(m=2t\delta/c\). Then
\(t(x+\delta)^2-\frac12x^2=-\frac c2x^2+2t\delta x+t\delta^2=-\frac c2(x-m)^2+\frac{t\delta^2}c\), since
\(\frac c2m^2=\frac{2t^2\delta^2}c\) and \(t\delta^2+\frac{2t^2\delta^2}c=\frac{t\delta^2(c+2t)}c=\frac{t\delta^2}c\).
\item Substituting \(x=m+y/\sqrt c\) (D1), \(\int_\R e^{-c(x-m)^2/2}\,\mathrm{d}x=c^{-1/2}\int_\R e^{-y^2/2}\,\mathrm{d}y=c^{-1/2}\sqrt{2\pi}\), as
\(\varphi\) integrates to \(1\) (Theorem~\ref{r14-thm-mvn-density}(a)).
\item Hence \(\expc\bigl(e^{t(w+\delta)^2}\bigr)=(2\pi)^{-1/2}e^{t\delta^2/c}\,c^{-1/2}\sqrt{2\pi}=c^{-1/2}e^{t\delta^2/c}\).
\end{itemize}
(b)
\begin{itemize}
\item By Theorem~\ref{r14-thm-mvn-iid}(a), \(w_1,\dots,w_r\) are independent \(N(0,1)\) variables; so the variables
\(u_i=(w_i+\delta_i)^2\) are independent (N3), and their MGFs are finite for \(t<\frac12\), by (a).
\item \(M_{u_1+\dots+u_r}(t)=\expc\bigl(e^{tu_1+\dots+tu_r}\bigr)\) is the joint MGF of \((u_1,\dots,u_r)'\) at
\((t,\dots,t)'\). By (N2) and (a), near \(0\) it equals
\(\prod_{i=1}^rM_{u_i}(t)=\prod_{i=1}^r(1-2t)^{-1/2}\exp\bigl(\delta_i^2t/(1-2t)\bigr)\), the stated expression.
\item If also \(\lVert\delta^*\rVert^2=\lambda\), then \(\lVert w+\delta^*\rVert^2\) has the same MGF near \(0\), hence the same
distribution (N1).\qedhere
\end{itemize}
\end{proof}

\begin{definition}[\citeat{tu-clm}{\S4.1}, \citeat{tu-asm}{Definition~P.115}]\label{r15-def-chisq}
Let \(r\geqslant1\), and let \(w_1,\dots,w_r\) be independent \(N(0,1)\) variables; equivalently,
\(w=(w_1,\dots,w_r)'\sim N_r(0,I_r)\) (Theorem~\ref{r14-thm-mvn-iid}(a)).
\begin{enumerate}
\item The distribution of \(w_1^2+\dots+w_r^2=\lVert w\rVert^2\) is the \textbf{chi-square distribution} \(\chi^2_r\) with \(r\)
\textbf{degrees of freedom}, as in (N5).
\item For \(\delta\in\R^r\), the distribution of \((w_1+\delta_1)^2+\dots+(w_r+\delta_r)^2=\lVert w+\delta\rVert^2\) is the
\textbf{noncentral chi-square distribution} \(\chi^2_r(\lambda)\) with \(r\) degrees of freedom and \textbf{noncentrality
parameter} \(\lambda=\delta_1^2+\dots+\delta_r^2\). By Lemma~\ref{r15-lem-mgf}(b) it depends only on \(r\) and \(\lambda\); every
\(\lambda\geqslant0\) occurs (take \(\delta_1=\sqrt\lambda\) and \(\delta_i=0\) for \(i\geqslant2\)), and \(\chi^2_r(0)=\chi^2_r\).
\end{enumerate}
We write \(u\sim\chi^2_r(\lambda)\) if \(u\) has this distribution.
\end{definition}

\begin{lemma}[\citeat{tu-asm}{Lemma~P.117}]\label{r15-lem-normal-moments}
If \(w\sim N(0,1)\), then \(\expc(w)=\expc(w^3)=0\), \(\expc(w^2)=1\) and \(\expc(w^4)=3\).
\end{lemma}
\begin{proof}
\begin{itemize}
\item \(N(0,1)\) has mean \(0\) and variance \(1\) (Definition~\ref{r14-def-mvn}), so \(\expc(w)=0\) and
\(\expc(w^2)=\var(w)=1\).
\item \(\lvert x\rvert^k\varphi(x)\) is integrable for \(k\leqslant4\), since \(x^ke^{-x^2/4}\) is bounded; so
\(\expc(w^k)=\int_\R x^k\varphi(x)\,\mathrm{d}x\) (the paragraph before Theorem~\ref{r14-thm-mvn-density}). For \(k=3\) the
integrand is odd, and \(\expc(w^3)=0\).
\item \(\varphi'(x)=-x\varphi(x)\), so, integrating by parts on \([-R,R]\),
\[
\int_{-R}^Rx^4\varphi(x)\,\mathrm{d}x=-\int_{-R}^Rx^3\varphi'(x)\,\mathrm{d}x=-2R^3\varphi(R)+3\int_{-R}^Rx^2\varphi(x)\,\mathrm{d}x
\]
\item Let \(R\to\infty\). From the exponential series, \(e^{R^2/2}\geqslant\frac12(R^2/2)^2=R^4/8\), so
\(0\leqslant R^3\varphi(R)\leqslant8(2\pi)^{-1/2}/R\to0\), while the integrals tend to \(\expc(w^4)\) and
\(\expc(w^2)=1\). Hence \(\expc(w^4)=3\).\qedhere
\end{itemize}
\end{proof}

\begin{theorem}[\citeat{casella-berger}{Lemma~5.3.2, Example~2.3.8}, \citeat{tu-asm}{Proposition~P.118}, \citeat{searle-khuri}{Corollary~10.2}]\label{r15-thm-chisq}
Let \(r,s\geqslant1\) and \(\lambda,\kappa\geqslant0\).
\begin{enumerate}
\item If \(u\sim\chi^2_r(\lambda)\), then \(M_u(t)=(1-2t)^{-r/2}\exp\bigl(\lambda t/(1-2t)\bigr)\) near \(0\); conversely, a random
variable with this MGF near \(0\) has the distribution \(\chi^2_r(\lambda)\).
\item \(\chi^2_r\) has the density
\[
f_r(x)=\frac{x^{r/2-1}e^{-x/2}}{\Gamma(r/2)\,2^{r/2}}\quad(x>0),\qquad f_r(x)=0\quad(x\leqslant0)
\]
In particular \(f_1(x)=(2\pi x)^{-1/2}e^{-x/2}\) and \(f_2(x)=\frac12e^{-x/2}\) for \(x>0\).
\item If \(u\sim\chi^2_r(\lambda)\) and \(v\sim\chi^2_s(\kappa)\) are independent, then \(u+v\sim\chi^2_{r+s}(\lambda+\kappa)\).
\item If \(u\sim\chi^2_r(\lambda)\), then \(\expc(u)=r+\lambda\) and \(\var(u)=2r+4\lambda\).
\end{enumerate}
\end{theorem}
\begin{proof}
(a) \(u\) has the distribution of \(\lVert w+\delta\rVert^2\) in Definition~\ref{r15-def-chisq}(b), whose MGF is given by
Lemma~\ref{r15-lem-mgf}(b). The converse is (N1).

\needspace{6\baselineskip}
(b)
\begin{itemize}
\item For \(t<\frac12\), by (D2) with \(c=r/2\) and \(b=(1-2t)/2\),
\[
\int_0^\infty e^{tx}f_r(x)\,\mathrm{d}x=\frac{1}{\Gamma(r/2)\,2^{r/2}}\int_0^\infty x^{r/2-1}e^{-(1-2t)x/2}\,\mathrm{d}x
=\frac{1}{2^{r/2}}\Bigl(\frac2{1-2t}\Bigr)^{r/2}=(1-2t)^{-r/2}
\]
\item For \(t=0\) this says \(\int f_r=1\). So \(f_r\geqslant0\) is the density of some random variable \(U\), and
\(M_U(t)=\int_0^\infty e^{tx}f_r(x)\,\mathrm{d}x=(1-2t)^{-r/2}\) for \(t<\frac12\) (the paragraph before
Theorem~\ref{r14-thm-mvn-density}). By (a) with \(\lambda=0\), \(U\sim\chi^2_r\); so \(\chi^2_r\) has the density \(f_r\).
\item \(\Gamma(\frac12)=\sqrt\pi\) and \(\Gamma(1)=1\) (D2) give \(f_1\) and \(f_2\).
\end{itemize}
(c) \(M_{u+v}(t)=\expc(e^{tu+tv})\) is the joint MGF of \((u,v)'\) at \((t,t)'\). By (N2) and (a), near \(0\),
\[
M_{u+v}(t)=M_u(t)M_v(t)=(1-2t)^{-(r+s)/2}\exp\Bigl(\frac{(\lambda+\kappa)t}{1-2t}\Bigr)
\]
so \(u+v\sim\chi^2_{r+s}(\lambda+\kappa)\) by (a).

(d)
\begin{itemize}
\item The mean and the variance depend only on the distribution, so we may take \(u=u_1+\dots+u_r\) with
\(u_i=(w_i+\delta_i)^2\) as in Definition~\ref{r15-def-chisq}(b).
\item By Lemma~\ref{r15-lem-normal-moments} and the linearity of the expectation,
\(\expc(u_i)=\expc(w_i^2)+2\delta_i\expc(w_i)+\delta_i^2=1+\delta_i^2\) and
\[
\expc(u_i^2)=\expc(w_i^4)+4\delta_i\expc(w_i^3)+6\delta_i^2\expc(w_i^2)+4\delta_i^3\expc(w_i)+\delta_i^4
=3+6\delta_i^2+\delta_i^4
\]
so \(\var(u_i)=\expc(u_i^2)-\bigl(\expc(u_i)\bigr)^2=2+4\delta_i^2\).
\item The \(u_i\) are independent (N3), so \(\cov(u_i,u_j)=0\) for \(i\ne j\) (N7). By
Theorem~\ref{r14-thm-cov}(d), \(\var(u)=\sum_i\var(u_i)=2r+4\lambda\); and
\(\expc(u)=\sum_i(1+\delta_i^2)=r+\lambda\).\qedhere
\end{itemize}
\end{proof}

\begin{example}\label{r15-exm-chisq}
\begin{enumerate}
\item If \(u\sim\chi^2_2\) and \(c>0\), then \(\prb(u>c)=\int_c^\infty\frac12e^{-x/2}\,\mathrm{d}x=e^{-c/2}\)
(Theorem~\ref{r15-thm-chisq}(b)). So \(\prb(u>c)=\frac1{20}\) for \(c=2\log20\).
\item Let \(z\sim N_4(\mu,I_4)\) with \(\mu=(1,2,2,4)'\). Then \(w=z-\mu\sim N_4(0,I_4)\) (Theorem~\ref{r14-thm-mvn-linear}), and
\(\lVert z\rVert^2=\lVert w+\mu\rVert^2\sim\chi^2_4(25)\), as \(\lambda=1+4+4+16=25\). By Theorem~\ref{r15-thm-chisq}(d),
\(\expc\lVert z\rVert^2=4+25=29\) and \(\var\lVert z\rVert^2=8+100=108\). Likewise \(z_1^2\sim\chi^2_1(1)\) and
\(z_2^2+z_3^2+z_4^2\sim\chi^2_3(24)\); these are independent by (N3), since \(z_1\) and \((z_2,z_3,z_4)'\) are independent
(Theorem~\ref{r14-thm-mvn-indep}(a): the off-diagonal blocks of \(I_4\) are \(0\)); and their sum is \(\chi^2_4(25)\), as
Theorem~\ref{r15-thm-chisq}(c) says.
\end{enumerate}
\end{example}

\needspace{24\baselineskip}
\subsection*{The \(t\) and \(F\) distributions}

\begin{theorem}[\citeat{casella-berger}{\S5.3.2, (5.3.6), Exercise~5.17(a)}]\label{r15-thm-t-F-density}
\begin{enumerate}
\item Let \(Z\sim N(0,1)\) and \(V\sim\chi^2_\nu\) (\(\nu\geqslant1\)) be independent. Then \(V>0\) with probability \(1\), and
\(T=Z/\sqrt{V/\nu}\) has the density
\[
f_T(x)=\frac{\Gamma\bigl(\frac{\nu+1}2\bigr)}{\Gamma\bigl(\frac\nu2\bigr)\sqrt{\nu\pi}}\Bigl(1+\frac{x^2}\nu\Bigr)^{-(\nu+1)/2}\qquad(x\in\R)
\]
\item Let \(U\sim\chi^2_a\) and \(V\sim\chi^2_b\) (\(a,b\geqslant1\)) be independent. Then \(F=(U/a)/(V/b)\) has the density
\[
f_F(x)=\frac{\Gamma\bigl(\frac{a+b}2\bigr)}{\Gamma\bigl(\frac a2\bigr)\Gamma\bigl(\frac b2\bigr)}\Bigl(\frac ab\Bigr)^{a/2}
\frac{x^{a/2-1}}{(1+ax/b)^{(a+b)/2}}\quad(x>0),\qquad f_F(x)=0\quad(x\leqslant0)
\]
and \(F\) and \(U+V\) are independent.
\end{enumerate}
\end{theorem}
\begin{proof}
(a)
\begin{itemize}
\item \(f_\nu=0\) on \((-\infty,0]\) (Theorem~\ref{r15-thm-chisq}(b)), so \(\prb(V\leqslant0)=\int_{-\infty}^0f_\nu=0\). As
\(Z\) and \(V\) are independent, \((Z,V)'\) has the density \(\varphi(z)f_\nu(v)\) (the paragraph before
Theorem~\ref{r14-thm-mvn-density}), which is \(0\) outside \(\Omega=\R\times(0,\infty)\).
\item The map \(\psi(t,w)=\bigl(t\sqrt{w/\nu},\,w\bigr)\) of \(\Omega\) onto itself is one-to-one, with inverse
\((z,v)\mapsto\bigl(z/\sqrt{v/\nu},\,v\bigr)\), and
\[
J(t,w)=\det\begin{pmatrix}\sqrt{w/\nu}&t/(2\sqrt{\nu w})\\0&1\end{pmatrix}=\sqrt{w/\nu}>0
\]
By (D1), \((T,V)=\psi^{-1}(Z,V)\) has the density \(\varphi\bigl(t\sqrt{w/\nu}\bigr)f_\nu(w)\sqrt{w/\nu}\) on \(\Omega\).
\item So \(T\) has the marginal density (Definition~\ref{r14-def-conditional})
\[
f_T(t)=\int_0^\infty\frac{e^{-t^2w/(2\nu)}}{\sqrt{2\pi}}\cdot\frac{w^{\nu/2-1}e^{-w/2}}{\Gamma(\nu/2)\,2^{\nu/2}}\cdot\frac{w^{1/2}}{\nu^{1/2}}\,\mathrm{d}w
=\frac{1}{\sqrt{2\pi}\,\Gamma(\nu/2)\,2^{\nu/2}\nu^{1/2}}\int_0^\infty w^{\frac{\nu+1}2-1}e^{-(1+t^2/\nu)w/2}\,\mathrm{d}w
\]
\item By (D2) with \(c=\frac{\nu+1}2\) and \(b=\frac12(1+t^2/\nu)\), the last integral is
\(\Gamma\bigl(\frac{\nu+1}2\bigr)2^{(\nu+1)/2}(1+t^2/\nu)^{-(\nu+1)/2}\); and \(2^{(\nu+1)/2}/(\sqrt{2\pi}\,2^{\nu/2})=1/\sqrt\pi\). This
is \(f_T(t)\).
\end{itemize}
(b) Put \(k=a/b\).
\begin{itemize}
\item As \(U\) and \(V\) are independent, \((U,V)'\) has the density \(f_a(u)f_b(v)\) (the paragraph before
Theorem~\ref{r14-thm-mvn-density}), which is \(0\) outside \(\Omega=(0,\infty)^2\).
\item The map \(\psi(x,y)=\bigl(kxy/(1+kx),\ y/(1+kx)\bigr)\) of \(\Omega\) onto itself is one-to-one, with inverse
\((u,v)\mapsto\bigl((u/a)/(v/b),\ u+v\bigr)\): for \((u,v)=\psi(x,y)\), \(u/v=kx\) and \(u+v=y\). Its Jacobian determinant is
\[
J(x,y)=\det\begin{pmatrix}ky/(1+kx)^2&kx/(1+kx)\\-ky/(1+kx)^2&1/(1+kx)\end{pmatrix}=\frac{ky(1+kx)}{(1+kx)^3}=\frac{ky}{(1+kx)^2}>0
\]
\item By (D1), \((F,U+V)=\psi^{-1}(U,V)\) has, on \(\Omega\), the density
\[
h(x,y)=f_a\Bigl(\frac{kxy}{1+kx}\Bigr)f_b\Bigl(\frac y{1+kx}\Bigr)\frac{ky}{(1+kx)^2}
=\frac{k^{a/2}x^{a/2-1}y^{(a+b)/2-1}e^{-y/2}}{\Gamma(a/2)\Gamma(b/2)\,2^{(a+b)/2}(1+kx)^{(a+b)/2}}
\]
collecting powers: \(y\) has the exponent \((\frac a2-1)+(\frac b2-1)+1\), and \(1+kx\) the exponent
\(-(\frac a2-1)-(\frac b2-1)-2\).
\item So \(h(x,y)=f_F(x)f_{a+b}(y)\) with \(f_F\geqslant0\) as stated. As \(h\) is a density and \(\int f_{a+b}=1\), the rule for
products of integrals before Theorem~\ref{r14-thm-mvn-density} gives \(1=\int\!\!\int h=\int f_F\). So \(F\) has the marginal
density \(\int_0^\infty h(x,y)\,\mathrm{d}y=f_F(x)\) (Definition~\ref{r14-def-conditional}), and \(U+V\sim\chi^2_{a+b}\) has the density
\(f_{a+b}\) (Theorem~\ref{r15-thm-chisq}(c),(b)). As \((F,U+V)'\) has the product density \(f_F(x)f_{a+b}(y)\), \(F\) and
\(U+V\) are independent (the paragraph before Theorem~\ref{r14-thm-mvn-density}).\qedhere
\end{itemize}
\end{proof}

\needspace{12\baselineskip}
\begin{definition}[\citeat{tu-clm}{\S4.1}, \citeat{tu-asm}{Definition~P.115}, \citeat{casella-berger}{Definitions~5.3.4, 5.3.6}]\label{r15-def-t-F}
Let \(\nu,a,b\geqslant1\).
\begin{enumerate}
\item If \(Z\sim N(0,1)\) and \(V\sim\chi^2_\nu\) are independent, the distribution of \(T=Z/\sqrt{V/\nu}\) is the
\textbf{\(t\) distribution} \(t_\nu\) with \(\nu\) degrees of freedom.
\item If \(U\sim\chi^2_a\) and \(V\sim\chi^2_b\) are independent, the distribution of \(F=(U/a)/(V/b)\) is the
\textbf{\(F\) distribution} \(F_{a,b}\) with \(a\) and \(b\) degrees of freedom.
\item For \(\tau\in(0,1)\), the \textbf{\(\tau\) quantiles} \(t_{\nu,\tau}\), \(F_{a,b,\tau}\) and \(\chi^2_{\nu,\tau}\) are the
numbers with \(\prb(T\leqslant t_{\nu,\tau})=\tau\) for \(T\sim t_\nu\), \(\prb(F\leqslant F_{a,b,\tau})=\tau\)
for \(F\sim F_{a,b}\), and \(\prb(u\leqslant\chi^2_{\nu,\tau})=\tau\) for \(u\sim\chi^2_\nu\).
\end{enumerate}
By Theorem~\ref{r15-thm-t-F-density}, these distributions have the densities \(f_T\) and \(f_F\), so they do not depend on
the choice of \(Z\), \(U\) and \(V\). Such variables exist: if \(g\sim N_{\nu+1}(0,I_{\nu+1})\)
(Theorem~\ref{r14-thm-mvn-density}(b)), then \(Z=g_1\sim N(0,1)\) and \(V=g_2^2+\dots+g_{\nu+1}^2\sim\chi^2_\nu\)
(Theorem~\ref{r14-thm-mvn-iid}(a), Definition~\ref{r15-def-chisq}) are independent (Theorem~\ref{r14-thm-mvn-indep}(a),
(N3)); \(U\) and \(V\) are obtained in the same way from \(g\sim N_{a+b}(0,I_{a+b})\). The quantiles exist and are unique:
\(G(x)=\prb(T\leqslant x)=\int_{-\infty}^xf_T\)
is continuous; it is strictly increasing on \(\R\), as \(f_T\) is continuous and positive there; and \(G(x)\to0\) as
\(x\to-\infty\), \(G(x)\to1\) as \(x\to\infty\), as \(\int f_T=1\). So \(G\) takes the value \(\tau\) exactly once (intermediate
value theorem). The same holds for \(F\), and for \(u\sim\chi^2_\nu\), whose density \(f_\nu\) (Theorem~\ref{r15-thm-chisq}(b))
is continuous and positive on \((0,\infty)\), with \((0,\infty)\) in place of \(\R\) and \(G(0)=0\). Hence, for
\(\alpha\in(0,1)\), \(\prb(T>t_{\nu,1-\alpha})=\alpha\), \(\prb(F>F_{a,b,1-\alpha})=\alpha\) and
\(\prb(u>\chi^2_{\nu,1-\alpha})=\alpha\).
\end{definition}

\needspace{14\baselineskip}
\begin{theorem}[\citeat{tu-clm}{\S4.1}, \citeat{tu-asm}{Corollary~P.116}, \citeat{casella-berger}{Theorem~5.3.8(b)}]\label{r15-thm-t-symmetric}
Let \(T\sim t_\nu\) and \(\alpha\in(0,1)\).
\begin{enumerate}
\item \(-T\sim t_\nu\): the \(t\) distribution is symmetric about \(0\). Hence \(\prb(T<-c)=\prb(T>c)\) for every
\(c\in\R\), and
\[
\prb(T<-t_{\nu,1-\alpha})=\alpha,\qquad
\prb\bigl(-t_{\nu,1-\alpha/2}\leqslant T\leqslant t_{\nu,1-\alpha/2}\bigr)=1-\alpha
\]
\item \(T^2\sim F_{1,\nu}\), and \(t_{\nu,1-\alpha/2}^2=F_{1,\nu,1-\alpha}\).
\end{enumerate}
\end{theorem}
\begin{proof}
Write \(T=Z/\sqrt{V/\nu}\) as in Definition~\ref{r15-def-t-F}.

(a)
\begin{itemize}
\item \(-Z\sim N(0,1)\) (Theorem~\ref{r14-thm-mvn-linear} with \(A=-1\)), and \(-Z\) is independent of \(V\) (N3). So
\(-T=(-Z)/\sqrt{V/\nu}\sim t_\nu\) (Definition~\ref{r15-def-t-F}), and
\(\prb(T<-c)=\prb(-T>c)=\prb(T>c)\).
\item For \(c=t_{\nu,1-\alpha}\), \(\prb(T>c)=\alpha\) (Definition~\ref{r15-def-t-F}); this gives the first identity.
\item With \(\alpha/2\) in place of \(\alpha\) and \(c=t_{\nu,1-\alpha/2}\),
\(\prb(-c\leqslant T\leqslant c)=1-\prb(T<-c)-\prb(T>c)=1-\frac\alpha2-\frac\alpha2\).
\end{itemize}
(b)
\begin{itemize}
\item \(T^2=(Z^2/1)/(V/\nu)\), where \(Z^2\sim\chi^2_1\) (Definition~\ref{r15-def-chisq}) is independent of \(V\) (N3); so
\(T^2\sim F_{1,\nu}\).
\item By (a), \(\prb(T<0)=\prb(T>0)\), and \(\prb(T=0)=0\), as \(T\) has a density; so
\(\prb(T\leqslant0)=\frac12<1-\frac\alpha2\), and \(c=t_{\nu,1-\alpha/2}>0\), as \(G\) in Definition~\ref{r15-def-t-F} is
increasing. Then \(\prb(T^2\leqslant c^2)=\prb(-c\leqslant T\leqslant c)=1-\alpha\) by (a), so \(c^2\) is the
\((1-\alpha)\) quantile of \(F_{1,\nu}\), which is unique (Definition~\ref{r15-def-t-F}).\qedhere
\end{itemize}
\end{proof}

\begin{proposition}[\citeat{casella-berger}{Theorem~5.3.8(a), Exercise~5.17(c)}]\label{r15-prop-F-reciprocal}
If \(F\sim F_{a,b}\), then \(1/F\sim F_{b,a}\). Hence \(F_{b,a,\alpha}=1/F_{a,b,1-\alpha}\) for \(\alpha\in(0,1)\).
\end{proposition}
\begin{proof}
\begin{itemize}
\item Write \(F=(U/a)/(V/b)\) as in Definition~\ref{r15-def-t-F}. \(U,V>0\) with probability \(1\) (as in the proof of
Theorem~\ref{r15-thm-t-F-density}(a)), and \(1/F=(V/b)/(U/a)\) with \(V\sim\chi^2_b\) and \(U\sim\chi^2_a\) independent; so
\(1/F\sim F_{b,a}\).
\item Let \(c=F_{a,b,1-\alpha}\); \(c>0\) (Definition~\ref{r15-def-t-F}). As \(F\) has a density, \(\prb(F=c)=0\), and
\(\prb(1/F\leqslant1/c)=\prb(F\geqslant c)=\prb(F>c)=\alpha\); so \(1/c\) is the \(\alpha\) quantile of
\(F_{b,a}\).\qedhere
\end{itemize}
\end{proof}

\begin{example}[\citeat{haerdle-simar}{\S4.4, (4.51), Theorem~4.7}]\label{r15-exm-quantiles}
Let \(\alpha\in(0,1)\).
\begin{itemize}
\item For \(T\sim t_4\), Theorem~\ref{r15-thm-t-symmetric}(a), with \(\alpha/2\) in place of \(\alpha\), gives
\(\prb(T<-t_{4,1-\alpha/2})=\frac\alpha2\); as \(T\) has a density, also \(\prb(T\leqslant-t_{4,1-\alpha/2})=\frac\alpha2\), so
\(t_{4,\alpha/2}=-t_{4,1-\alpha/2}\) (Definition~\ref{r15-def-t-F}(c)), and \(\prb(\lvert T\rvert\leqslant t_{4,1-\alpha/2})=1-\alpha\).
By Theorem~\ref{r15-thm-t-symmetric}(b), \(F_{1,4,1-\alpha}=t_{4,1-\alpha/2}^2\); by
Proposition~\ref{r15-prop-F-reciprocal}, \(F_{4,3,\alpha}=1/F_{3,4,1-\alpha}\).
\item Let \(u\sim\chi^2_2\) and \(\tau\in(0,1)\). For \(c>0\), \(\prb(u\leqslant c)=1-\prb(u>c)=1-e^{-c/2}\)
(Example~\ref{r15-exm-chisq}(a)), and \(1-e^{-c/2}=\tau\) exactly for \(c=-2\log(1-\tau)\), which is \(>0\). So
\(\chi^2_{2,\tau}=-2\log(1-\tau)\) (Definition~\ref{r15-def-t-F}(c)); in particular \(\chi^2_{2,1-\alpha}=2\log(1/\alpha)\), and
\(\chi^2_{2,19/20}=2\log20\).
\item Let \(x\sim N_n(\mu,\Sigma)\) with \(\Sigma\) PD. By Corollary~\ref{r14-cor-standard}(a),
\(d_\Sigma(x,\mu)^2=(x-\mu)'\Sigma^{-1}(x-\mu)\sim\chi^2_n\), with the distance \(d_\Sigma\) of the paragraph after
Theorem~\ref{r14-thm-mvn-density}. So the region \(\{u\in\R^n:d_\Sigma(u,\mu)^2\leqslant\chi^2_{n,1-\alpha}\}\), bounded by the
contour of that paragraph with \(c^2=\chi^2_{n,1-\alpha}\), contains \(x\) with probability \(1-\alpha\). For \(n=2\) it is
bounded by an ellipse, with \(c^2=2\log(1/\alpha)\).
\end{itemize}
\end{example}

\subsection*{One normal sample}

\begin{theorem}[\citeat{casella-berger}{Theorem~5.3.1}, \citeat{searle-khuri}{Examples~10.2, 10.3}]\label{r15-thm-one-sample}
Let \(n\geqslant2\), and let \(x_1,\dots,x_n\) be independent \(N(\mu,\sigma^2)\) variables (\(\sigma>0\)), with mean
\(\bar x=\frac1n\sum_ix_i\) and sample variance \(s^2=\frac1{n-1}\sum_i(x_i-\bar x)^2\) (Definition~\ref{r12-def-data}). Then:
\begin{enumerate}
\item \(\bar x\sim N(\mu,\sigma^2/n)\);
\item \((n-1)s^2/\sigma^2\sim\chi^2_{n-1}\);
\item \(\bar x\) and \(s^2\) are independent;
\item \(\sqrt n(\bar x-\mu)/s\sim t_{n-1}\).
\end{enumerate}
\end{theorem}
\begin{proof}
Let \(x=(x_1,\dots,x_n)'\). Then \(x\sim N_n(\mu\one,\sigma^2I_n)\) (Theorem~\ref{r14-thm-mvn-iid}(a)), the deviations are
\(C_nx=(x_1-\bar x,\dots,x_n-\bar x)'\) (Example~\ref{r1-exm-centering}), and \((n-1)s^2=(C_nx)'(C_nx)\).

(a) This is Example~\ref{r14-exm-sample}(a), applied to \(x\).

(b)
\begin{itemize}
\item Let \(z=(x-\mu\one)/\sigma\); then \(z\sim N_n(0,I_n)\) (Theorem~\ref{r14-thm-mvn-linear}).
\item \(C_n\) is symmetric with the eigenvalues \(1\) (\(n-1\) times) and \(0\) (once) (Example~\ref{r8-exm-patterned}(a)). By
Theorem~\ref{r8-thm-spectral}, there is an orthogonal \(P\) with \(P'C_nP=D=\diag(1,\dots,1,0)\), so \(C_n=PDP'\). (For instance
\(P=[p_2\ \cdots\ p_n\ p_1]\), with the rows \(p_k'\) of the Helmert matrix \(K_n\) of Example~\ref{r6-exm-helmert}.)
\item Let \(y=P'z\). Then \(y\sim N_n(0,P'P)=N_n(0,I_n)\) (Theorem~\ref{r14-thm-mvn-linear}), so \(y_1,\dots,y_n\) are independent
\(N(0,1)\) variables (Theorem~\ref{r14-thm-mvn-iid}(a)).
\item \(C_n\one=0\) and \(C_n'C_n=C_n\) (Proposition~\ref{r1-prop-centering}), so \(C_nx=\sigma C_nz\) and
\[
(n-1)s^2=(C_nx)'(C_nx)=\sigma^2z'C_nz=\sigma^2z'PDP'z=\sigma^2y'Dy=\sigma^2(y_1^2+\dots+y_{n-1}^2)
\]
So \((n-1)s^2/\sigma^2\sim\chi^2_{n-1}\) (Definition~\ref{r15-def-chisq}).
\end{itemize}
(c) By Example~\ref{r14-exm-sample}(b), applied to \(x\), \(\bar x\) is independent of the deviations \(C_nx\), and hence of
\(s^2=\frac1{n-1}(C_nx)'(C_nx)\).

(d)
\begin{itemize}
\item By (a) and Theorem~\ref{r14-thm-mvn-linear}, \(Z=\sqrt n(\bar x-\mu)/\sigma\sim N(0,1)\); by (b),
\(V=(n-1)s^2/\sigma^2\sim\chi^2_{n-1}\); and \(Z\) and \(V\) are independent by (c) and (N3).
\item \(V\) has a density (Theorem~\ref{r15-thm-chisq}(b)), so \(\prb(s^2=0)=\prb(V=0)=0\), and \(s>0\) with
probability \(1\). Then
\[
\frac{\sqrt n(\bar x-\mu)}s=\frac{\sqrt n(\bar x-\mu)/\sigma}{\sqrt{(n-1)s^2/\bigl((n-1)\sigma^2\bigr)}}=\frac Z{\sqrt{V/(n-1)}}
\]
which has the distribution \(t_{n-1}\) (Definition~\ref{r15-def-t-F}).\qedhere
\end{itemize}
\end{proof}

\begin{definition}[\citeat{casella-berger}{Definitions~9.1.1, 9.1.4, 9.1.5}]\label{r15-def-confidence}
Let observations be modelled as random variables whose distribution depends on parameters, and let \(\alpha\in(0,1)\).
\begin{enumerate}
\item An interval \([L,U]\) whose endpoints \(L\leqslant U\) are computed from the observations is a \textbf{confidence
interval} with level \(1-\alpha\) for a parameter \(\theta\in\R\) if \(\prb(L\leqslant\theta\leqslant U)=1-\alpha\) for all values
of the parameters.
\item A set \(C\subseteq\R^p\) computed from the observations is a \textbf{confidence region} with level \(1-\alpha\) for a
parameter \(\theta\in\R^p\) if \(\prb(\theta\in C)=1-\alpha\) for all values of the parameters.
\end{enumerate}
\end{definition}

\begin{example}[\citeat{casella-berger}{Example~9.2.10}]\label{r15-exm-one-sample}
In Theorem~\ref{r15-thm-one-sample}, let \(\alpha\in(0,1)\), \(T=\sqrt n(\bar x-\mu)/s\), \(V=(n-1)s^2/\sigma^2\) and
\(c=t_{n-1,1-\alpha/2}\); here \(s>0\) with probability \(1\) (the proof of Theorem~\ref{r15-thm-one-sample}(d)).
\begin{itemize}
\item By Theorems~\ref{r15-thm-one-sample}(d) and~\ref{r15-thm-t-symmetric}(a), \(\prb(-c\leqslant T\leqslant c)=1-\alpha\) for all
values of \(\mu\) and \(\sigma\). The inequality \(-c\leqslant T\leqslant c\) holds exactly when
\(\bar x-cs/\sqrt n\leqslant\mu\leqslant\bar x+cs/\sqrt n\). So
\(\bigl[\bar x-t_{n-1,1-\alpha/2}\,s/\sqrt n,\ \bar x+t_{n-1,1-\alpha/2}\,s/\sqrt n\bigr]\) is a confidence interval for \(\mu\) with
level \(1-\alpha\) (Definition~\ref{r15-def-confidence}(a)).
\item Let \(c_1=\chi^2_{n-1,\alpha/2}\) and \(c_2=\chi^2_{n-1,1-\alpha/2}\) (Definition~\ref{r15-def-t-F}(c)). By
Theorem~\ref{r15-thm-one-sample}(b), \(V\sim\chi^2_{n-1}\), which has a density (Theorem~\ref{r15-thm-chisq}(b)); so
\(\prb(V=c_1)=0\) and \(\prb(c_1\leqslant V\leqslant c_2)=\prb(V\leqslant c_2)-\prb(V<c_1)=(1-\frac\alpha2)-\frac\alpha2=1-\alpha\) for
all values of \(\mu\) and \(\sigma\). Also \(0<c_1<c_2\), as \(\prb(V\leqslant0)=0<\frac\alpha2<1-\frac\alpha2\). As
\(c_1\leqslant V\leqslant c_2\) holds exactly when \((n-1)s^2/c_2\leqslant\sigma^2\leqslant(n-1)s^2/c_1\),
\(\bigl[(n-1)s^2/\chi^2_{n-1,1-\alpha/2},\ (n-1)s^2/\chi^2_{n-1,\alpha/2}\bigr]\) is a confidence interval for \(\sigma^2\) with
level \(1-\alpha\).
\item Five measurements \(2,4,5,6,8\) are modelled as independent \(N(\mu,\sigma^2)\) variables. Here \(\bar x=25/5=5\), the
deviations are \(-3,-1,0,1,3\), and \(s^2=(9+1+0+1+9)/4=5\); so \(s/\sqrt5=1\) and \((n-1)s^2=20\). The two intervals are
\(\bigl[5-t_{4,1-\alpha/2},\ 5+t_{4,1-\alpha/2}\bigr]\) and \(\bigl[20/\chi^2_{4,1-\alpha/2},\ 20/\chi^2_{4,\alpha/2}\bigr]\).
\item If \(\mu=3\), then \(\sqrt5(\bar x-3)/s\sim t_4\); its value here is \(\sqrt5\cdot2/\sqrt5=2\). So \(3\) lies in the interval
for \(\mu\) exactly when \(t_{4,1-\alpha/2}\geqslant2\).
\end{itemize}
\end{example}

\begin{example}[cf.\ \citeat{haerdle-simar}{\S7.1, (7.2), Theorem~4.9}]\label{r11-exm-confidence-ellipse}
Four dimensions of each of \(n=16\) machined parts are measured; the measurement vectors \(x_1,\dots,x_{16}\in\R^4\) are
modelled as independent \(N_4(\mu,\Sigma)\), with
\[
\Sigma=2I_4+2J_4=\begin{pmatrix}4&2&2&2\\2&4&2&2\\2&2&4&2\\2&2&2&4\end{pmatrix}
\]
known from long experience and \(\mu\) unknown. So each measurement has the variance \(4\), and any two have the
correlation \(\frac12\). Let \(\alpha\in(0,1)\).
\begin{itemize}
\item By Theorem~\ref{r14-thm-mvn-iid}(a) (the block form), \(u=(x_1',\dots,x_{16}')'\sim N_{64}\bigl((\mu',\dots,\mu')',
\diag(\Sigma,\dots,\Sigma)\bigr)\), and \(\bar x=\frac1{16}(x_1+\dots+x_{16})=\frac1{16}Gu\) with \(G=[I_4\ \cdots\ I_4]\in\R^{4\times64}\).
By Theorem~\ref{r14-thm-mvn-linear} and the block product (Theorem~\ref{r5b-thm-block}(c)), as
\(G\diag(\Sigma,\dots,\Sigma)G'=16\Sigma\), \(\bar x\sim N_4\bigl(\mu,\frac1{256}\cdot16\Sigma\bigr)=N_4(\mu,\frac1{16}\Sigma)\). Each
\(\bar x_i\) has the standard deviation \(\sqrt{4/16}=\frac12\).
\item \(\Sigma\) has the eigenvalues \(2+8=10\) and \(2\) (three times) (Example~\ref{r8-exm-patterned}(b)), so it is PD
(Theorem~\ref{r8-thm-eigen-criterion}(a)), and so is \(\frac1{16}\Sigma\), with the inverse \(16\Sigma^{-1}\). Here
\(\Sigma^{-1}=\frac12(I_4-\frac15J_4)\), as \((I_4+J_4)(I_4-\frac15J_4)=I_4+J_4-\frac15J_4-\frac45J_4=I_4\) (\(J_4^2=4J_4\),
Proposition~\ref{r1-prop-centering}). By Corollary~\ref{r14-cor-standard}(a), \(q(\mu)=16(\bar x-\mu)'\Sigma^{-1}(\bar x-\mu)\sim\chi^2_4\),
so \(\prb\bigl(q(\mu)\leqslant\chi^2_{4,1-\alpha}\bigr)=1-\alpha\) for every \(\mu\) (Definition~\ref{r15-def-t-F}(c)). Thus the ellipsoid
\(\{m\in\R^4:16(\bar x-m)'\Sigma^{-1}(\bar x-m)\leqslant\chi^2_{4,1-\alpha}\}\), computed from the observations, is a confidence
region for \(\mu\) with level \(1-\alpha\) (Definition~\ref{r15-def-confidence}(b)).
\item By Theorem~\ref{r15-thm-chisq}(b), \(\chi^2_4\) has the density \(f_4(x)=\frac x4e^{-x/2}\) for \(x>0\), as
\(\Gamma(2)=\Gamma(1)=1\) (D2). For \(w\sim\chi^2_4\) and \(c>0\), integrating by parts, as \(xe^{-x/2}\to0\) for \(x\to\infty\),
\[
\prb(w>c)=\int_c^\infty\frac x4e^{-x/2}\,\mathrm{d}x=\Bigl[-\frac x2e^{-x/2}\Bigr]_c^\infty+\int_c^\infty\frac12e^{-x/2}\,\mathrm{d}x=\Bigl(1+\frac c2\Bigr)e^{-c/2}
\]
\item Let \(\alpha=\frac1{20}\), and compare two possible target values \(m_1\), \(m_2\) with \(\bar x-m_1=(1,1,1,1)'\) and
\(\bar x-m_2=(1,-1,1,-1)'\). Each coordinate of \(\bar x\) is two of its standard deviations away from the corresponding
coordinate of \(m_1\), and also of \(m_2\). With \(d=\bar x-m\), \(q(m)=16d'\Sigma^{-1}d=8\bigl(d'd-\frac15(\one'd)^2\bigr)\), so
\[
q(m_1)=8\Bigl(4-\frac{16}5\Bigr)=\frac{32}5,\qquad q(m_2)=8(4-0)=32
\]
\item From the exponential series, \(2<e<3\) (\(e=\sum_{k\geqslant0}1/k!\), and \(1/k!\leqslant1/2^{k-1}\) for \(k\geqslant1\), strictly for \(k\geqslant3\)). So
\(\prb(w>\frac{32}5)=\frac{21}5e^{-16/5}>\frac1{20}\), as \(e^{16}<3^{16}=43046721<84^5\); and \(\prb(w>32)=17e^{-16}<\frac1{20}\), as
\(e^{16}>2^{16}=65536>340\). As \(\prb(w>c)\) is strictly decreasing in \(c\) (Definition~\ref{r15-def-t-F}) and
\(\prb(w>\chi^2_{4,19/20})=\frac1{20}\), \(q(m_1)<\chi^2_{4,19/20}<q(m_2)\): \(m_1\) lies in the confidence region and \(m_2\) does
not. The measurements are positively correlated, so deviations of the same sign are common and deviations of
opposite signs are unusual; only the joint region takes this into account.
\end{itemize}
\end{example}

\needspace{14\baselineskip}
\subsection*{Quadratic forms in normal vectors}

\begin{theorem}[\citeat{tu-clm}{Theorem~4.5}, \citeat{bapat}{Result~8.3}, \citeat{searle-khuri}{Theorem~10.2}]\label{r15-thm-quadratic-chisq}
Let \(z\sim N_n(\mu,\sigma^2I_n)\) (\(\sigma>0\)), and let \(A\in\R^{n\times n}\) be symmetric and idempotent with
\(\rank A=r\geqslant1\). Then
\[
\frac{z'Az}{\sigma^2}\sim\chi^2_r\Bigl(\frac{\mu'A\mu}{\sigma^2}\Bigr)
\]
In particular \(z'Az/\sigma^2\sim\chi^2_r\) if \(A\mu=0\).
\end{theorem}
\begin{proof}
\begin{itemize}
\item By Theorem~\ref{r8-thm-idempotent-eigen}(c), \(A=P_1P_1'\) for some \(P_1\in\R^{n\times r}\) with \(P_1'P_1=I_r\).
\item Let \(u=P_1'z/\sigma\) and \(\delta=P_1'\mu/\sigma\). By Theorem~\ref{r14-thm-mvn-linear},
\(u\sim N_r\bigl(\delta,\ P_1'(\sigma^2I_n)P_1/\sigma^2\bigr)=N_r(\delta,I_r)\), and \(w=u-\delta\sim N_r(0,I_r)\).
\item \(z'Az/\sigma^2=z'P_1P_1'z/\sigma^2=\lVert u\rVert^2=\lVert w+\delta\rVert^2\). By Definition~\ref{r15-def-chisq}, this has the
distribution \(\chi^2_r(\lambda)\) with \(\lambda=\lVert\delta\rVert^2=\mu'P_1P_1'\mu/\sigma^2=\mu'A\mu/\sigma^2\).
\item If \(A\mu=0\), then \(\lambda=0\).\qedhere
\end{itemize}
\end{proof}

\needspace{18\baselineskip}
\begin{theorem}[\citeat{searle-khuri}{Theorem~10.1}]\label{r15-thm-decomposition}
Let \(y\sim N_n(\mu,\Sigma)\) with \(\Sigma\) positive definite, let \(A\in\R^{n\times n}\) be symmetric, and let
\(B=\Sigma^{1/2}A\Sigma^{1/2}\) (Theorem~\ref{r8-thm-sqrt}). Let \(\gamma_1,\dots,\gamma_k\) be the distinct nonzero eigenvalues of
\(B\), with multiplicities \(\nu_1,\dots,\nu_k\), and let the columns of \(P_i\in\R^{n\times\nu_i}\) be orthonormal eigenvectors of \(B\)
for \(\gamma_i\). Then
\[
y'Ay=\gamma_1W_1+\dots+\gamma_kW_k,\qquad W_i=\lVert P_i'\Sigma^{-1/2}y\rVert^2
\]
where \(W_1,\dots,W_k\) are independent and \(W_i\sim\chi^2_{\nu_i}(\theta_i)\) with
\(\theta_i=\mu'\Sigma^{-1/2}P_iP_i'\Sigma^{-1/2}\mu\). The \(\gamma_i\) and \(\nu_i\) are also the distinct nonzero eigenvalues of \(A\Sigma\)
and their multiplicities.
\end{theorem}
\begin{proof}
\begin{itemize}
\item \(\Sigma^{1/2}\) is symmetric and invertible, with \(\Sigma^{1/2}\Sigma^{1/2}=\Sigma\) and inverse \(\Sigma^{-1/2}\)
(Theorem~\ref{r8-thm-sqrt}). Let \(z=\Sigma^{-1/2}y\) and \(c=\Sigma^{-1/2}\mu\). By Theorem~\ref{r14-thm-mvn-linear},
\(z\sim N_n(c,\Sigma^{-1/2}\Sigma\Sigma^{-1/2})=N_n(c,I_n)\), and \(y'Ay=z'\Sigma^{1/2}A\Sigma^{1/2}z=z'Bz\).
\item \(B\) is symmetric. The \(\nu_i\) columns of \(P_i\) are orthonormal vectors of the eigenspace \(\mathcal N(B-\gamma_iI)\), whose
dimension is \(\nu_i\) (Theorem~\ref{r8-thm-rank-eigen}); so they form an orthonormal basis of it, and \(P_iP_i'\) is its
projection matrix (Theorem~\ref{r6-thm-projection-matrix}(b)). By Theorem~\ref{r8-thm-spectral-decomp}(b), with the eigenvalue
\(0\), if it occurs, contributing nothing, \(B=\sum_i\gamma_iP_iP_i'\), and \(z'Bz=\sum_i\gamma_i\lVert P_i'z\rVert^2=\sum_i\gamma_iW_i\).
\item The \(m=\nu_1+\dots+\nu_k\) columns of \(P=[P_1\ \cdots\ P_k]\) are orthonormal (Theorem~\ref{r8-thm-sym-orth} for
different eigenvalues), so \(u=P'z\sim N_m(P'c,P'P)=N_m(P'c,I_m)\) (Theorem~\ref{r14-thm-mvn-linear}). Its blocks are
\(u_i=P_i'z\sim N_{\nu_i}(P_i'c,I_{\nu_i})\), and the off-diagonal blocks of \(I_m\) are \(0\); so \(u_1,\dots,u_k\) are
independent (Theorem~\ref{r14-thm-mvn-indep}(b)).
\item By Definition~\ref{r15-def-chisq}, \(W_i=\lVert u_i\rVert^2=\lVert(u_i-P_i'c)+P_i'c\rVert^2\sim\chi^2_{\nu_i}(\theta_i)\) with
\(\theta_i=\lVert P_i'c\rVert^2=c'P_iP_i'c\); and the \(W_i\) are independent (N3).
\item \(B=\Sigma^{1/2}(A\Sigma)\Sigma^{-1/2}\) is similar to \(A\Sigma\), so the two matrices have the same characteristic
polynomial (Theorem~\ref{r7-thm-similar}), hence the same eigenvalues with the same multiplicities.\qedhere
\end{itemize}
\end{proof}

\begin{theorem}[\citeat{bapat}{Result~8.3}]\label{r15-thm-chisq-converse}
Let \(z\sim N_n(0,I_n)\), let \(A\in\R^{n\times n}\) be symmetric, and let \(r\geqslant1\). If \(z'Az\sim\chi^2_r\), then \(A\) is
idempotent and \(\rank A=r\).
\end{theorem}
\begin{proof}
\begin{itemize}
\item By Theorem~\ref{r8-thm-spectral}, \(A=Q\diag(\gamma_1,\dots,\gamma_n)Q'\) with \(Q\) orthogonal. Then \(y=Q'z\sim N_n(0,I_n)\)
(Theorem~\ref{r14-thm-mvn-linear}), its coordinates are independent \(N(0,1)\) variables (Theorem~\ref{r14-thm-mvn-iid}(a)), and
\(z'Az=y'\diag(\gamma_1,\dots,\gamma_n)y=\sum_j\gamma_jy_j^2\).
\item Let \(t_0>0\) with \(2t_0\lvert\gamma_j\rvert<1\) for every \(j\). For \(\lvert t\rvert<t_0\), Lemma~\ref{r15-lem-mgf}(a) (with
\(\delta=0\) and \(t\gamma_j\) in place of \(t\)) gives \(M_{\gamma_jy_j^2}(t)=(1-2\gamma_jt)^{-1/2}\). The \(\gamma_jy_j^2\) are
independent (N3), so, as in Lemma~\ref{r15-lem-mgf}(b), \(M_{z'Az}(t)=\prod_j(1-2\gamma_jt)^{-1/2}\) near \(0\) (N2).
\item Also \(M_{z'Az}(t)=(1-2t)^{-r/2}\) near \(0\) (Theorem~\ref{r15-thm-chisq}(a)). Squaring and inverting,
\(\prod_{j=1}^n(1-2\gamma_jt)=(1-2t)^r\) for all \(t\) in an open interval; both sides are polynomials in \(t\), so they are
equal (Proposition~\ref{r7-prop-poly}(d)).
\item The only root of the right side is \(\frac12\). If \(\gamma_j\ne0\), then \(t=1/(2\gamma_j)\) is a root of the left side
(Proposition~\ref{r7-prop-poly}(b)); so \(1/(2\gamma_j)=\frac12\), that is, \(\gamma_j=1\). Hence the left side is \((1-2t)^m\),
where \(m\) is the number of \(j\) with \(\gamma_j\ne0\), and comparing degrees (Proposition~\ref{r7-prop-poly}(a)) gives
\(m=r\).
\item So the eigenvalues of \(A\) are \(1\) (\(r\) times) and \(0\) (\(n-r\) times) (Theorem~\ref{r8-thm-spectral}(c)): \(A\) is
idempotent (Theorem~\ref{r8-thm-idempotent-eigen}(b)), and \(\rank A=r\) (Theorem~\ref{r8-thm-rank-eigen}).\qedhere
\end{itemize}
\end{proof}

\needspace{10\baselineskip}
\begin{theorem}[\citeat{searle-khuri}{Theorem~10.2, Corollary~10.1}, \citeat{bapat}{Ch.~8, Exercise~15}]\label{r15-thm-chisq-pd}
Let \(y\sim N_n(\mu,\Sigma)\) with \(\Sigma\) positive definite, let \(A\in\R^{n\times n}\) be symmetric, let \(r\geqslant1\), and let
\(B=\Sigma^{1/2}A\Sigma^{1/2}\).
\begin{enumerate}
\item \(A\Sigma\) is idempotent if and only if \(B\) is idempotent; and \(\rank(A\Sigma)=\rank B\).
\item If \(A\Sigma\) is idempotent of rank \(r\), then \(y'Ay\sim\chi^2_r(\mu'A\mu)\).
\item If \(\mu=0\) and \(y'Ay\sim\chi^2_r\), then \(A\Sigma\) is idempotent of rank \(r\).
\item \(A\Sigma\) is idempotent of rank \(r\) if and only if \(\tr(A\Sigma)=\tr\bigl((A\Sigma)^2\bigr)=\rank(A\Sigma)=r\).
\end{enumerate}
\end{theorem}
\begin{proof}
(a)
\begin{itemize}
\item \(B^2=\Sigma^{1/2}A\Sigma A\Sigma^{1/2}\) (Theorem~\ref{r8-thm-sqrt}). If \((A\Sigma)^2=A\Sigma\), multiplying on the right
by \(\Sigma^{-1}\) gives \(A\Sigma A=A\), and then \(B^2=\Sigma^{1/2}A\Sigma^{1/2}=B\).
\item If \(B^2=B\), multiplying by \(\Sigma^{-1/2}\) on both sides gives \(A\Sigma A=A\), and multiplying on the right by \(\Sigma\)
gives \((A\Sigma)^2=A\Sigma\).
\item \(B=\Sigma^{1/2}(A\Sigma)\Sigma^{-1/2}\), so \(\rank B=\rank(A\Sigma)\) (Theorem~\ref{thm:r4-rank-product}(d)).
\end{itemize}
(b) Let \(z=\Sigma^{-1/2}y\sim N_n(\Sigma^{-1/2}\mu,I_n)\) (Theorem~\ref{r14-thm-mvn-linear}); then \(y'Ay=z'Bz\). By (a), \(B\) is
symmetric and idempotent of rank \(r\), so by Theorem~\ref{r15-thm-quadratic-chisq} (with \(\sigma=1\)),
\(z'Bz\sim\chi^2_r(\mu'\Sigma^{-1/2}B\Sigma^{-1/2}\mu)=\chi^2_r(\mu'A\mu)\).

(c) Now \(z\sim N_n(0,I_n)\) and \(z'Bz=y'Ay\sim\chi^2_r\); by Theorem~\ref{r15-thm-chisq-converse}, \(B\) is idempotent of rank \(r\),
and so is \(A\Sigma\), by (a).

(d)
\begin{itemize}
\item If \(A\Sigma\) is idempotent of rank \(r\), then \(\tr(A\Sigma)=\rank(A\Sigma)=r\) (Theorem~\ref{r6-thm-idempotent}(c)), and
\((A\Sigma)^2=A\Sigma\) has the same trace.
\item Conversely, \(\tr B=\tr(A\Sigma)\) and \(\tr B^2=\tr\bigl(\Sigma^{1/2}(A\Sigma)^2\Sigma^{-1/2}\bigr)=\tr\bigl((A\Sigma)^2\bigr)\)
(Theorem~\ref{r1-thm-trace}(b)), and \(\rank B=\rank(A\Sigma)\) by (a). Write \(B=Q\diag(\beta_1,\dots,\beta_n)Q'\)
(Theorem~\ref{r8-thm-spectral}); then \(\tr B=\sum_i\beta_i\), \(\tr B^2=\sum_i\beta_i^2\) (Theorem~\ref{r1-thm-trace}(b)), and
\(\rank B\) is the number of \(i\) with \(\beta_i\ne0\) (Theorem~\ref{r8-thm-rank-eigen}).
\item If all three numbers equal \(r\), then \(\sum_{\beta_i\ne0}(\beta_i-1)^2=\sum_i\beta_i^2-2\sum_i\beta_i+r=0\), so every nonzero
\(\beta_i\) is \(1\): \(B\) is idempotent (Theorem~\ref{r8-thm-idempotent-eigen}(b)) of rank \(r\), and so is \(A\Sigma\), by
(a).\qedhere
\end{itemize}
\end{proof}

\needspace{12\baselineskip}
\subsection*{Independence of quadratic and linear forms}

\begin{lemma}\label{r15-lem-function}
Let \(A\in\R^{n\times n}\) be symmetric. There is a symmetric \(G\in\R^{n\times n}\) with \(AGA=A\); hence
\(x'Ax=(Ax)'G(Ax)\) for every \(x\in\R^n\): the number \(x'Ax\) is a function of the vector \(Ax\).
\end{lemma}
\begin{proof}
If \(A=0\), take \(G=0\). Otherwise let \(A=c_1E_1+\dots+c_kE_k\) be the spectral decomposition
(Theorem~\ref{r8-thm-spectral-decomp}), and let \(G=\sum_{c_j\ne0}c_j^{-1}E_j\), which is symmetric. As \(E_iE_j=0\) for \(i\ne j\)
and \(E_j^2=E_j\) (Theorem~\ref{r8-thm-spectral-decomp}(b)), \(AGA=\sum_{c_j\ne0}c_jc_j^{-1}c_jE_j=\sum_jc_jE_j=A\). Then
\(x'Ax=x'AGAx=(Ax)'G(Ax)\), as \(A'=A\).
\end{proof}

\needspace{14\baselineskip}
\begin{theorem}[\citeat{tu-clm}{Theorem~4.6}, \citeat{bapat}{Result~8.5}, \citeat{searle-khuri}{Theorems~10.3, 10.4}]\label{r15-thm-indep}
Let \(z\sim N_n(\mu,\sigma^2I_n)\), and let \(A,B\in\R^{n\times n}\) be symmetric.
\begin{enumerate}
\item If \(AB=0\), then \(z'Az\) and \(z'Bz\) are independent.
\item If \(C\in\R^{m\times n}\) and \(CA=0\), then \(Cz\) and \(z'Az\) are independent.
\end{enumerate}
\end{theorem}
\begin{proof}
(a)
\begin{itemize}
\item \(A(\sigma^2I_n)B'=\sigma^2AB=0\), so \(Az\) and \(Bz\) are independent (Theorem~\ref{r14-thm-mvn-indep}(c)).
\item By Lemma~\ref{r15-lem-function}, \(z'Az\) is a function of \(Az\), and \(z'Bz\) is a function of \(Bz\); so they are
independent (N3).
\end{itemize}
(b) \(C(\sigma^2I_n)A'=\sigma^2CA=0\), so \(Cz\) and \(Az\) are independent (Theorem~\ref{r14-thm-mvn-indep}(c)); and \(z'Az\) is a
function of \(Az\) (Lemma~\ref{r15-lem-function}). So \(Cz\) and \(z'Az\) are independent (N3).
\end{proof}
For instance, in Theorem~\ref{r15-thm-one-sample}, \(\bar x=\frac1n\one'x\) and \(s^2=x'C_nx/(n-1)\) are independent by (b), as
\(\frac1n\one'C_n=0\).

\needspace{14\baselineskip}
\begin{theorem}[\citeat{searle-khuri}{Theorems~10.3, 10.4}]\label{r15-thm-indep-general}
Let \(y\sim N_n(\mu,\Sigma)\) (\(\Sigma\) positive semidefinite), and let \(A,B\in\R^{n\times n}\) be symmetric and
\(C\in\R^{m\times n}\).
\begin{enumerate}
\item If \(A\Sigma B=0\), then \(y'Ay\) and \(y'By\) are independent.
\item If \(C\Sigma A=0\), then \(Cy\) and \(y'Ay\) are independent.
\end{enumerate}
\end{theorem}
\begin{proof}
(a) \(A\Sigma B'=A\Sigma B=0\), so \(Ay\) and \(By\) are independent (Theorem~\ref{r14-thm-mvn-indep}(c)); by
Lemma~\ref{r15-lem-function} and (N3), so are \(y'Ay\) and \(y'By\).

(b) \(C\Sigma A'=C\Sigma A=0\), so \(Cy\) and \(Ay\) are independent (Theorem~\ref{r14-thm-mvn-indep}(c)); by
Lemma~\ref{r15-lem-function} and (N3), so are \(Cy\) and \(y'Ay\).
\end{proof}

\begin{theorem}[\citeat{bapat}{Results~8.4, 8.5}]\label{r15-thm-indep-converse}
Let \(z\sim N_n(0,I_n)\), and let \(A_1,A_2,A\in\R^{n\times n}\) be symmetric and idempotent.
\begin{enumerate}
\item \(z'A_1z\) and \(z'A_2z\) are independent if and only if \(A_1A_2=0\).
\item For \(\ell\in\R^n\), \(z'Az\) and \(\ell'z\) are independent if and only if \(A\ell=0\).
\end{enumerate}
\end{theorem}
\begin{proof}
(a) If \(A_1A_2=0\), the two forms are independent by Theorem~\ref{r15-thm-indep}(a). Conversely, let them be independent.
If \(A_1=0\) or \(A_2=0\), then \(A_1A_2=0\); so let \(r_i=\rank A_i\geqslant1\).
\begin{itemize}
\item \(z'A_iz\sim\chi^2_{r_i}\) (Theorem~\ref{r15-thm-quadratic-chisq}), so \(z'(A_1+A_2)z\sim\chi^2_{r_1+r_2}\)
(Theorem~\ref{r15-thm-chisq}(c)). \(A_1+A_2\) is symmetric, so it is idempotent (Theorem~\ref{r15-thm-chisq-converse}).
\item Then \(A_1+A_2=(A_1+A_2)^2=A_1+A_2+A_1A_2+A_2A_1\), so \(A_1A_2+A_2A_1=0\).
\item Multiplying on the right by \(A_2\) gives \(A_1A_2+A_2A_1A_2=0\); multiplying this on the left by \(A_2\) gives
\(2A_2A_1A_2=0\). So \(A_2A_1A_2=0\), and then \(A_1A_2=0\).
\end{itemize}
(b) If \(A\ell=0\), then \(\ell'A=(A\ell)'=0\), and Theorem~\ref{r15-thm-indep}(b) (with \(C=\ell'\)) gives independence.
Conversely, let \(z'Az\) and \(\ell'z\) be independent, and let \(\ell\ne0\) (for \(\ell=0\) there is nothing to prove).
\begin{itemize}
\item \(B=\ell\ell'/\ell'\ell\) is symmetric and idempotent, and \(z'Bz=(\ell'z)^2/\ell'\ell\) is a function of \(\ell'z\); so \(z'Az\)
and \(z'Bz\) are independent (N3), and \(AB=0\) by (a).
\item Hence \(A\ell=A\ell(\ell'\ell)/\ell'\ell=AB\ell=0\).\qedhere
\end{itemize}
\end{proof}

\begin{example}[\citeat{bapat}{Ch.~8, Exercise~3}]\label{r15-exm-linear-indep}
Let \(x_1,x_2\) be independent \(N(0,1)\) variables, so that \(x=(x_1,x_2)'\sim N_2(0,I_2)\) (Theorem~\ref{r14-thm-mvn-iid}(a)). Which
linear functions \(c'x=c_1x_1+c_2x_2\) are independent of \((x_1-x_2)^2\)? With \(a=(1,-1)'\) and
\[
A=\tfrac12aa'=\frac12\begin{pmatrix}1&-1\\-1&1\end{pmatrix}
\]
which is symmetric and idempotent (\(A^2=\frac14a(a'a)a'=A\)), \((x_1-x_2)^2=(a'x)^2=2\,x'Ax\). A function of \(x'Ax\) is
independent of \(c'x\) when \(x'Ax\) is, and conversely (N3); so by Theorem~\ref{r15-thm-indep-converse}(b), \(c'x\) is
independent of \((x_1-x_2)^2\) exactly when \(Ac=\frac12a(a'c)=0\), that is, \(c_1=c_2\). These are the multiples of
\(x_1+x_2\).
\end{example}

\begin{example}[cf.\ \citeat{searle-khuri}{Example~10.1}]\label{r15-exm-equicorrelated}
Let \(n\geqslant2\), \(\sigma>0\), \(-1/(n-1)<\rho<1\), and \(x\sim N_n(\mu\one,\sigma^2R)\) with \(R=(1-\rho)I_n+\rho J_n\): equal means
and variances, and the correlation \(\rho\) between any two observations. \(R\) is positive definite
(Example~\ref{r8-exm-definite}(c)), and so is \(\Sigma=\sigma^2R\). Let \(\bar x\) and \(s^2\) be as in
Theorem~\ref{r15-thm-one-sample}.
\begin{itemize}
\item \(C_nJ_n=(C_n\one)\one'=0\), so \(C_nR=(1-\rho)C_n\). With \(A=C_n/\bigl(\sigma^2(1-\rho)\bigr)\), \(A\Sigma=C_n\), which is
idempotent of rank \(n-1\) (Proposition~\ref{r1-prop-centering}, Corollary~\ref{r6-cor-trace-hat}), and \(A(\mu\one)=0\). By
Theorem~\ref{r15-thm-chisq-pd}(b),
\((n-1)s^2/\bigl(\sigma^2(1-\rho)\bigr)=x'Ax\sim\chi^2_{n-1}\); in particular \(\expc(s^2)=\sigma^2(1-\rho)\)
(Theorem~\ref{r15-thm-chisq}(d)).
\item \(\bar x=c'x\) with \(c=\frac1n\one\). As \(\one'R=\bigl(1+(n-1)\rho\bigr)\one'\) and \(\one'C_n=0\), \(c'\Sigma A=0\); so \(\bar x\) and
\(s^2\) are independent (Theorem~\ref{r15-thm-indep-general}(b)). Also
\(\bar x\sim N\bigl(\mu,\ \sigma^2c'Rc\bigr)=N\bigl(\mu,\ \sigma^2(1+(n-1)\rho)/n\bigr)\) (Theorem~\ref{r14-thm-mvn-linear}).
\item By the two steps above and Theorem~\ref{r14-thm-mvn-linear}, \(Z=\sqrt n(\bar x-\mu)/\sqrt{\sigma^2(1+(n-1)\rho)}\sim N(0,1)\)
and \(V=(n-1)s^2/\bigl(\sigma^2(1-\rho)\bigr)\sim\chi^2_{n-1}\) are independent (N3). As in the proof of
Theorem~\ref{r15-thm-one-sample}(d), by Definition~\ref{r15-def-t-F},
\[
T=\frac{\sqrt n(\bar x-\mu)}s\sqrt{\frac{1-\rho}{1+(n-1)\rho}}=\frac Z{\sqrt{V/(n-1)}}\sim t_{n-1}
\]
For \(\rho=0\) this is Theorem~\ref{r15-thm-one-sample}(d). For \(\rho>0\) the usual statistic \(\sqrt n(\bar x-\mu)/s\) is \(T\) times
\(\sqrt{(1+(n-1)\rho)/(1-\rho)}\).
\item Let \(n=9\), \(\rho=\frac14\) and \(\alpha\in(0,1)\). The factor is \(\sqrt{(1+2)/\frac34}=2\). With \(c=t_{8,1-\alpha/2}\), which is
\(>0\) (the proof of Theorem~\ref{r15-thm-t-symmetric}(b)), the interval \(\bar x\pm cs/\sqrt9\), which has the level \(1-\alpha\) for
independent observations (Example~\ref{r15-exm-one-sample}), contains \(\mu\) exactly when \(\lvert\sqrt9(\bar x-\mu)/s\rvert\leqslant c\),
that is, when \(\lvert T\rvert\leqslant\frac c2\). By Theorem~\ref{r15-thm-t-symmetric}(a),
\(\prb(\lvert T\rvert\leqslant\frac c2)=1-2\prb(T>\frac c2)\), and \(\prb(T>\frac c2)>\prb(T>c)=\frac\alpha2\), as the distribution function of
\(t_8\) is strictly increasing (Definition~\ref{r15-def-t-F}). So the interval contains \(\mu\) with probability less than
\(1-\alpha\): with positively correlated observations it is too short.
\end{itemize}
\end{example}

\subsection*{Cochran's theorem}

\begin{theorem}[\citeat{bapat}{Result~8.6}]\label{r15-thm-cochran-algebraic}
Let \(A_1,\dots,A_p\in\R^{n\times n}\) with \(A_1+\dots+A_p=I_n\). The following are equivalent:
\begin{enumerate}
\item \(\rank A_1+\dots+\rank A_p=n\);
\item \(A_i^2=A_i\) for \(i=1,\dots,p\);
\item \(A_iA_j=0\) for \(i\ne j\).
\end{enumerate}
\end{theorem}
\begin{proof}
(a)\(\Rightarrow\)(c)
\begin{itemize}
\item Products with a zero matrix vanish, and deleting the \(A_i=0\) changes neither the sum of the matrices nor the sum of
the ranks; so let every \(r_i=\rank A_i\geqslant1\), with full-rank factorizations \(A_i=B_iC_i\), \(B_i\in\R^{n\times r_i}\),
\(C_i\in\R^{r_i\times n}\) (Theorem~\ref{thm:r4-full-rank-factorization}).
\item By the block product (Theorem~\ref{r5b-thm-block}(c)),
\[
[B_1\ \cdots\ B_p]\begin{pmatrix}C_1\\\vdots\\C_p\end{pmatrix}=B_1C_1+\dots+B_pC_p=I_n
\]
and the first factor is square, as \(r_1+\dots+r_p=n\). By Theorem~\ref{r2-thm-invertible} the two factors are inverse to each
other, so their product in the other order is \(I_n\); its \((i,j)\) block is \(C_iB_j\), so \(C_iB_j=0\) for \(i\ne j\).
\item Hence \(A_iA_j=B_i(C_iB_j)C_j=0\) for \(i\ne j\).
\end{itemize}
(c)\(\Rightarrow\)(b) \(A_j=A_jI_n=A_j(A_1+\dots+A_p)=A_j^2\), since \(A_jA_i=0\) for \(i\ne j\).

(b)\(\Rightarrow\)(a) \(\rank A_i=\tr A_i\) (Theorem~\ref{r6-thm-idempotent}(c)), so
\(\sum_i\rank A_i=\sum_i\tr A_i=\tr(A_1+\dots+A_p)=\tr I_n=n\) (Theorem~\ref{r1-thm-trace}(a)).
\end{proof}

\needspace{18\baselineskip}
\begin{theorem}[\citeat{searle-khuri}{Theorem~10.5}, \citeat{bapat}{\S8.3}]\label{r15-thm-cochran}
Let \(x\sim N_n(\mu,\sigma^2I_n)\) (\(\sigma>0\)), and let \(A_1,\dots,A_p\in\R^{n\times n}\) be symmetric with
\(A_1+\dots+A_p=I_n\) and \(r_i=\rank A_i\geqslant1\). Consider the conditions
\begin{enumerate}
\item \(r_1+\dots+r_p=n\);
\item \(x'A_ix/\sigma^2\sim\chi^2_{r_i}(\mu'A_i\mu/\sigma^2)\) for \(i=1,\dots,p\);
\item \(x'A_1x,\dots,x'A_px\) are mutually independent.
\end{enumerate}
Then (a) implies (b) and (c), and (b) implies (a).
\end{theorem}
\begin{proof}
(a)\(\Rightarrow\)(b), (c)
\begin{itemize}
\item By Theorem~\ref{r15-thm-cochran-algebraic}, every \(A_i\) is idempotent and \(A_iA_j=0\) for \(i\ne j\). So (b) holds by
Theorem~\ref{r15-thm-quadratic-chisq}.
\item Let \(u_i=A_ix\) and \(u=(u_1',\dots,u_p')'=Lx\), where \(L\in\R^{pn\times n}\) has the blocks \(A_1,\dots,A_p\) stacked. By
Theorem~\ref{r14-thm-mvn-linear}, \(u\sim N_{pn}(L\mu,\sigma^2LL')\), and the \((i,j)\) block of \(LL'\) is \(A_iA_j'=A_iA_j\), which
is \(0\) for \(i\ne j\) (Theorem~\ref{r5b-thm-block}(c)).
\item By Theorem~\ref{r14-thm-mvn-indep}(b), \(u_1,\dots,u_p\) are mutually independent, and so are the functions
\(x'A_ix=x'A_i'A_ix=\lVert u_i\rVert^2\) (N3): this is (c).
\end{itemize}
(b)\(\Rightarrow\)(a) By Theorem~\ref{r15-thm-chisq}(d), \(\expc(x'A_ix/\sigma^2)=r_i+\mu'A_i\mu/\sigma^2\); by
Theorem~\ref{r11-thm-quadratic}, \(\expc(x'A_ix)=\tr(A_i\sigma^2I_n)+\mu'A_i\mu\). So \(\tr A_i=r_i\), and
\(\sum_ir_i=\tr(A_1+\dots+A_p)=\tr I_n=n\) (Theorem~\ref{r1-thm-trace}(a)).
\end{proof}

\begin{example}[\citeat{bapat}{Ch.~8, Exercise~5}]\label{r15-exm-cochran}
Let \(x_1,\dots,x_n\) (\(n\geqslant2\)) be independent \(N(0,1)\) variables, so that \(x=(x_1,\dots,x_n)'\sim N_n(0,I_n)\)
(Theorem~\ref{r14-thm-mvn-iid}(a)), and let \(\bar x=\frac1n\one'x\). With Cochran's theorem we show that \(\bar x\) and
\(\sum_i(x_i-\bar x)^2\) are independent.
\begin{itemize}
\item Let \(A_1=\bar J_n\) and \(A_2=C_n\). They are symmetric, \(A_1+A_2=I_n\) (Definition~\ref{r1-def-ones}), and
\(\rank A_1=\tr\bar J_n=1\), \(\rank A_2=n-1\) (Proposition~\ref{r1-prop-centering}, Theorem~\ref{r6-thm-idempotent}(c),
Corollary~\ref{r6-cor-trace-hat}); so \(\rank A_1+\rank A_2=n\).
\item \(x'A_1x=\frac1n(\one'x)^2=n\bar x^2\) and \(x'A_2x=\sum_i(x_i-\bar x)^2\) (Example~\ref{r1-exm-centering}). By
Theorem~\ref{r15-thm-cochran}, (a) \(\Rightarrow\) (b), (c), with \(\mu=0\) and \(\sigma=1\): \(n\bar x^2\sim\chi^2_1\) and
\(\sum_i(x_i-\bar x)^2\sim\chi^2_{n-1}\), and they are independent.
\item The vectors \(A_1x=\bar x\one\) and \(A_2x\) are independent (Theorem~\ref{r14-thm-mvn-indep}(c)), as
\(A_1I_nA_2'=\bar J_nC_n=\frac1n\one(C_n\one)'=0\) (\(C_n\) is symmetric and \(C_n\one=0\), Proposition~\ref{r1-prop-centering}). So \(\bar x\), the first entry of \(A_1x\), and \(\sum_i(x_i-\bar x)^2=\lVert A_2x\rVert^2\)
are independent (N3).
\end{itemize}
\end{example}

\needspace{16\baselineskip}
\subsection*{Moments of quadratic forms}

\begin{theorem}[\citeat{searle-khuri}{(10.34), (10.35), Corollary~10.2}, \citeat{abadir-magnus}{Exercises~11.21, 11.22}]\label{r15-thm-moments}
Let \(y\sim N_n(\mu,\Sigma)\) (\(\Sigma\) positive semidefinite), and let \(A\in\R^{n\times n}\) be symmetric. Then
\[
\expc(y'Ay)=\tr(A\Sigma)+\mu'A\mu,\qquad\var(y'Ay)=2\tr\bigl((A\Sigma)^2\bigr)+4\mu'A\Sigma A\mu
\]
\end{theorem}
\begin{proof}
The mean is Theorem~\ref{r11-thm-quadratic}. For the variance:
\begin{itemize}
\item Let \(L=\Sigma^{1/2}\) (Theorem~\ref{r8-thm-sqrt}), so that \(L'=L\) and \(L^2=\Sigma\), and let \(w\sim N_n(0,I_n)\). Then
\(\mu+Lw\sim N_n(\mu,\Sigma)\) (Theorem~\ref{r14-thm-mvn-iid}(b)), so \(y'Ay\) has the distribution of \((\mu+Lw)'A(\mu+Lw)\), and we
may take \(y=\mu+Lw\). Then
\[
y'Ay=\mu'A\mu+c'w+w'Bw,\qquad c=2LA\mu,\quad B=LAL
\]
since the number \(w'LA\mu\) equals \(\mu'ALw\); \(B\) is symmetric.
\item Write \(B=Q\diag(\gamma_1,\dots,\gamma_n)Q'\) (Theorem~\ref{r8-thm-spectral}) and \(v=Q'w\sim N_n(0,I_n)\)
(Theorem~\ref{r14-thm-mvn-linear}); the \(v_i\) are independent \(N(0,1)\) variables (Theorem~\ref{r14-thm-mvn-iid}(a)), and
\(w'Bw=\sum_i\gamma_iv_i^2\). The \(\gamma_iv_i^2\) are independent (N3), hence uncorrelated (N7), with
\(\var(\gamma_iv_i^2)=\gamma_i^2\var(v_i^2)=2\gamma_i^2\) (Theorems~\ref{r11-thm-linear}(b) and~\ref{r15-thm-chisq}(d)). By
Theorem~\ref{r14-thm-cov}(d), \(\var(w'Bw)=2\sum_i\gamma_i^2=2\tr(B^2)\), as \(B^2=Q\diag(\gamma_1^2,\dots,\gamma_n^2)Q'\)
(Theorem~\ref{r1-thm-trace}(b)).
\item \(\var(c'w)=c'c=4\mu'ALLA\mu=4\mu'A\Sigma A\mu\) (Theorem~\ref{r11-thm-linear}(b)).
\item \(\expc(c'w)=0\), so \(\cov(w'Bw,c'w)=\expc(w'Bw\,c'w)=\sum_{i,j,l}b_{ij}c_l\expc(w_iw_jw_l)\). Each
\(\expc(w_iw_jw_l)=0\): some index \(h\) occurs once or three times among \(i,j,l\), so \(w_iw_jw_l=w_h^e\,g\) with
\(e\in\{1,3\}\) and \(g\) a function of the vector \(Ew\) of the other coordinates (the rows of \(E\) are the \(e_k'\), \(k\ne h\)).
As \(e_h'I_nE'=0\), \(w_h\) and \(Ew\) are independent (Theorem~\ref{r14-thm-mvn-indep}(c)), and so are \(w_h^e\) and \(g\) (N3); so
the expectation is \(\expc(w_h^e)\expc(g)\) (N7), and \(\expc(w_h^e)=0\) (Lemma~\ref{r15-lem-normal-moments}).
\item As \(\mu'A\mu\) is a constant, Theorem~\ref{r14-thm-cov}(a),(c) gives
\[
\var(y'Ay)=\var(w'Bw)+\var(c'w)+2\cov(w'Bw,c'w)=2\tr(B^2)+4\mu'A\Sigma A\mu
\]
and \(\tr(B^2)=\tr(LALLAL)=\tr(ALLALL)=\tr\bigl((A\Sigma)^2\bigr)\) (Theorem~\ref{r1-thm-trace}(b)).\qedhere
\end{itemize}
\end{proof}

\begin{corollary}[\citeat{searle-khuri}{Exercise~10.12}]\label{r15-cor-cov}
Let \(y\sim N_n(\mu,\Sigma)\), and let \(A,B\in\R^{n\times n}\) be symmetric. Then
\[
\cov(y'Ay,\,y'By)=2\tr(A\Sigma B\Sigma)+4\mu'A\Sigma B\mu
\]
\end{corollary}
\begin{proof}
\begin{itemize}
\item By Theorem~\ref{r14-thm-cov}(c),
\(\var\bigl(y'(A+B)y\bigr)=\var(y'Ay)+\var(y'By)+2\cov(y'Ay,y'By)\).
\item Apply Theorem~\ref{r15-thm-moments} to \(A+B\), \(A\) and \(B\). Here
\(\tr\bigl(((A+B)\Sigma)^2\bigr)=\tr\bigl((A\Sigma)^2\bigr)+\tr\bigl((B\Sigma)^2\bigr)+2\tr(A\Sigma B\Sigma)\), as
\(\tr(B\Sigma A\Sigma)=\tr(A\Sigma B\Sigma)\) (Theorem~\ref{r1-thm-trace}(b)); and
\(\mu'(A+B)\Sigma(A+B)\mu=\mu'A\Sigma A\mu+\mu'B\Sigma B\mu+2\mu'A\Sigma B\mu\), as the number \(\mu'B\Sigma A\mu\) equals its
transpose \(\mu'A\Sigma B\mu\). So
\(2\cov(y'Ay,y'By)=4\tr(A\Sigma B\Sigma)+8\mu'A\Sigma B\mu\).\qedhere
\end{itemize}
\end{proof}

\needspace{9\baselineskip}
\begin{example}[\citeat{casella-berger}{Exercise~5.10(c)}, \citeat{casella-berger-solutions}{Exercise~5.10(c)}]\label{r15-exm-var-s2}
In Theorem~\ref{r15-thm-one-sample}, \(\var(s^2)=2\sigma^4/(n-1)\). Two ways:
\begin{itemize}
\item \((n-1)s^2/\sigma^2\sim\chi^2_{n-1}\) has the variance \(2(n-1)\) (Theorem~\ref{r15-thm-chisq}(d)), so
\(\var(s^2)=\frac{\sigma^4}{(n-1)^2}\cdot2(n-1)=\frac{2\sigma^4}{n-1}\).
\item \(s^2=x'Ax\) with \(A=C_n/(n-1)\) and \(x\sim N_n(\mu\one,\sigma^2I_n)\). Here \(A\Sigma=\sigma^2C_n/(n-1)\), so
\(\tr\bigl((A\Sigma)^2\bigr)=\sigma^4\tr(C_n)/(n-1)^2=\sigma^4/(n-1)\) (Proposition~\ref{r1-prop-centering},
Corollary~\ref{r6-cor-trace-hat}), and \(A\mu\one=0\). By Theorem~\ref{r15-thm-moments},
\(\var(s^2)=2\sigma^4/(n-1)+0\).
\end{itemize}
\end{example}

\needspace{14\baselineskip}
\subsection*{Exercises}

\begin{exercise}\label{r15-dr-chisq-mgf}
Let \(u\) have the MGF \((1-2t)^{-3}\exp\bigl(4t/(1-2t)\bigr)\) near \(0\), and let \(v\sim\chi^2_2(1)\) be independent of \(u\).
\begin{enumerate}
\item Identify the distribution of \(u\), and find \(\expc(u)\) and \(\var(u)\).
\item Find the distribution of \(u+v\), its mean and its variance, and check that the means and the variances add up.
\item Let \(z\sim N_4(\mu,I_4)\) with \(\mu=(1,-1,2,1)'\). Find the distributions of \(z_1^2+z_2^2\) and of \(\lVert z\rVert^2\).
\end{enumerate}
\end{exercise}
\begin{solution}
(a) \((1-2t)^{-3}=(1-2t)^{-6/2}\), so \(u\sim\chi^2_6(4)\) (Theorem~\ref{r15-thm-chisq}(a)). By Theorem~\ref{r15-thm-chisq}(d),
\(\expc(u)=6+4=10\) and \(\var(u)=12+16=28\).

(b) By Theorem~\ref{r15-thm-chisq}(c), \(u+v\sim\chi^2_8(5)\), with mean \(8+5=13\) and variance \(16+20=36\). Also
\(\expc(v)=2+1=3\) and \(\var(v)=4+4=8\), and indeed \(10+3=13\) and \(28+8=36\).

(c) \(w=z-\mu\sim N_4(0,I_4)\) (Theorem~\ref{r14-thm-mvn-linear}), so \(w_1,\dots,w_4\) are independent \(N(0,1)\) variables
(Theorem~\ref{r14-thm-mvn-iid}(a)), and \(z_i=w_i+\mu_i\). By Definition~\ref{r15-def-chisq},
\(z_1^2+z_2^2=(w_1+1)^2+(w_2-1)^2\sim\chi^2_2(2)\) and \(\lVert z\rVert^2\sim\chi^2_4(1+1+4+1)=\chi^2_4(7)\).
\end{solution}

\begin{exercise}\label{r15-dr-chisq-tail}
\begin{enumerate}
\item Write down the densities of \(\chi^2_2\) and \(\chi^2_4\).
\item For \(u\sim\chi^2_4\) and \(c>0\), show by integration by parts that \(\prb(u>c)=(1+c/2)e^{-c/2}\). Using
\(\frac83<e<3\), show that \(8<\chi^2_{4,19/20}<10\) (Definition~\ref{r15-def-t-F}(c)).
\item For \(v\sim\chi^2_2\), find the median of \(v\), the number \(c\) with \(\prb(v\leqslant c)=\frac12\).
\end{enumerate}
\end{exercise}
\begin{solution}
(a) By Theorem~\ref{r15-thm-chisq}(b), with \(\Gamma(1)=1\) and \(\Gamma(2)=1\cdot\Gamma(1)=1\) (D2): \(f_2(x)=\frac12e^{-x/2}\) and
\(f_4(x)=\frac14xe^{-x/2}\) for \(x>0\), and both are \(0\) for \(x\leqslant0\).

(b) As \(u\) has the density \(f_4\), integrating by parts (\(xe^{-x/2}\to0\) for \(x\to\infty\)),
\[
\prb(u>c)=\int_c^\infty\frac x4e^{-x/2}\,\mathrm{d}x=\Bigl[-\frac x2e^{-x/2}\Bigr]_c^\infty+\int_c^\infty\frac12e^{-x/2}\,\mathrm{d}x=\frac c2e^{-c/2}+e^{-c/2}
\]
\begin{itemize}
\item From the series \(e=\sum_{k\geqslant0}1/k!\): \(e>1+1+\frac12+\frac16=\frac83\), and \(e<1+1+\frac12+\frac14+\dots=3\), as
\(k!\geqslant2^{k-1}\) for \(k\geqslant1\), strictly for \(k\geqslant3\).
\item \(\prb(u>8)=5e^{-4}>\frac1{20}\), as \(e^4<3^4=81<100\); and \(\prb(u>10)=6e^{-5}<\frac1{20}\), as
\(e^5>(\frac83)^5=\frac{32768}{243}>120\) (\(32768>120\cdot243=29160\)).
\item \(\prb(u>\chi^2_{4,19/20})=\frac1{20}\) (Definition~\ref{r15-def-t-F}), and \(c\mapsto\prb(u>c)\) is strictly decreasing, as the
distribution function of \(u\) is strictly increasing on \((0,\infty)\). So \(8<\chi^2_{4,19/20}<10\).
\end{itemize}
(c) \(\prb(v\leqslant c)=1-e^{-c/2}\) (Example~\ref{r15-exm-chisq}(a)), which is \(\frac12\) for \(c=2\log2\) (as in
Example~\ref{r15-exm-quantiles}, \(\chi^2_{2,1/2}=-2\log\frac12\)).
\end{solution}

\begin{exercise}\label{r15-dr-quantiles}
Let \(T\sim t_5\), \(\alpha\in(0,1)\) and \(c=t_{5,1-\alpha/2}\).
\begin{enumerate}
\item Find \(\prb(T>c)\), \(t_{5,\alpha/2}\) and \(\prb(\lvert T\rvert\leqslant c)\).
\item Express \(F_{1,5,1-\alpha}\) through \(c\).
\item Express \(F_{5,2,\alpha}\) and \(F_{2,5,\alpha}\) through quantiles \(F_{a,b,1-\alpha}\).
\end{enumerate}
\end{exercise}
\begin{solution}
(a) \(\prb(T>c)=1-(1-\frac\alpha2)=\frac\alpha2\) (Definition~\ref{r15-def-t-F}). By Theorem~\ref{r15-thm-t-symmetric}(a),
\(\prb(T<-c)=\frac\alpha2\); as \(T\) has a density, also \(\prb(T\leqslant-c)=\frac\alpha2\), so \(t_{5,\alpha/2}=-c\), and
\(\prb(\lvert T\rvert\leqslant c)=1-\frac\alpha2-\frac\alpha2=1-\alpha\).

(b) \(F_{1,5,1-\alpha}=t_{5,1-\alpha/2}^2=c^2\) (Theorem~\ref{r15-thm-t-symmetric}(b)).

(c) By Proposition~\ref{r15-prop-F-reciprocal}, \(F_{5,2,\alpha}=1/F_{2,5,1-\alpha}\) and \(F_{2,5,\alpha}=1/F_{5,2,1-\alpha}\).
\end{solution}

\begin{exercise}\label{r15-dr-one-sample}
Seven measurements \(1,3,4,5,6,7,9\) are modelled as independent \(N(\mu,\sigma^2)\) variables. Let \(\alpha\in(0,1)\) and
\(c=t_{6,1-\alpha/2}\).
\begin{enumerate}
\item Compute \(\bar x\) and \(s^2\).
\item Compute \(\sqrt7(\bar x-2)/s\), the statistic of Theorem~\ref{r15-thm-one-sample}(d) for \(\mu=2\), and find the condition on
\(c\) under which it lies in \([-c,c]\).
\item Compute the confidence interval \(\bar x\pm c\,s/\sqrt7\) for \(\mu\) with level \(1-\alpha\) (Example~\ref{r15-exm-one-sample}).
Does it contain \(2\)?
\end{enumerate}
\end{exercise}
\begin{solution}
(a) \(\bar x=35/7=5\); the deviations are \(-4,-2,-1,0,1,2,4\), so \(s^2=(16+4+1+0+1+4+16)/6=42/6=7\).

(b) \(s/\sqrt7=1\), so \(\sqrt7(\bar x-2)/s=(5-2)/1=3\), which lies in \([-c,c]\) exactly when \(c\geqslant3\). If \(\mu=2\), the
statistic is a \(t_6\) variable (Theorem~\ref{r15-thm-one-sample}(d)), and it lies outside \([-c,c]\) with probability \(\alpha\)
(Theorem~\ref{r15-thm-t-symmetric}(a)).

(c) \([5-t_{6,1-\alpha/2},\ 5+t_{6,1-\alpha/2}]\). It contains \(2\) exactly when \(5-c\leqslant2\), that is, \(t_{6,1-\alpha/2}\geqslant3\),
the condition of (b).
\end{solution}

\begin{exercise}\label{r11-ex-variance-interval}
Three measurements \(4,7,10\) are modelled as independent \(N(\mu,\sigma^2)\) variables. Let \(\alpha\in(0,1)\).
\begin{enumerate}
\item Write down \(\bar x\), \((n-1)s^2\) and the confidence interval of Example~\ref{r15-exm-one-sample} for \(\sigma^2\).
\item Write the interval of (a) in closed form with Example~\ref{r15-exm-quantiles}. Does it contain \(1\) for \(\alpha=\frac1{20}\)?
Show that its upper end then exceeds \(351\).
\item Write down the confidence interval of Example~\ref{r15-exm-one-sample} for \(\mu\).
\end{enumerate}
\end{exercise}
\begin{solution}
(a) Here \(n=3\), \(\bar x=7\) and \((n-1)s^2=9+0+9=18\). By Example~\ref{r15-exm-one-sample} the interval is
\(\bigl[18/\chi^2_{2,1-\alpha/2},\ 18/\chi^2_{2,\alpha/2}\bigr]\).

(b) By Example~\ref{r15-exm-quantiles}, \(\chi^2_{2,\tau}=-2\log(1-\tau)\); so \(\chi^2_{2,1-\alpha/2}=2\log(2/\alpha)\) and
\(\chi^2_{2,\alpha/2}=2\log\bigl(2/(2-\alpha)\bigr)\), and the interval is
\[
\Bigl[\frac9{\log(2/\alpha)},\ \frac9{\log\bigl(2/(2-\alpha)\bigr)}\Bigr],\qquad\text{for }\alpha=\tfrac1{20}:\quad
\Bigl[\frac9{\log40},\ \frac9{\log(40/39)}\Bigr]
\]
As \(e>1+1=2\) (\(e=\sum_{k\geqslant0}1/k!\)), \(e^9>2^9=512>40\), so \(\log40<9\) and the lower end exceeds \(1\): the interval does not contain \(1\). By
Lemma~\ref{r14-lem-logdet} with \(p=1\) and \(A=(\frac{40}{39})\ne I_1\), \(\log\frac{40}{39}<\frac{40}{39}-1=\frac1{39}\), so the upper
end exceeds \(9\cdot39=351\). Three observations say little about \(\sigma^2\).

(c) \(s^2=9\), \(s=3\) and \(s/\sqrt3=\sqrt3\), so the interval is
\(\bigl[7-\sqrt3\,t_{2,1-\alpha/2},\ 7+\sqrt3\,t_{2,1-\alpha/2}\bigr]\).
\end{solution}

\begin{exercise}\label{r11-ex-ellipse}
Let \(x\sim N_4(\mu,\Sigma)\) with
\[
\mu=\begin{pmatrix}1\\1\\1\\1\end{pmatrix},\qquad\Sigma=\begin{pmatrix}4&3&1&1\\3&4&1&1\\1&1&4&3\\1&1&3&4\end{pmatrix}
\]
and let \(d^2=(x-\mu)'\Sigma^{-1}(x-\mu)\).
\begin{enumerate}
\item Check that \(h_1=(1,1,1,1)'\), \(h_2=(1,1,-1,-1)'\), \(h_3=(1,-1,1,-1)'\) and \(h_4=(1,-1,-1,1)'\) are eigenvectors of \(\Sigma\),
and conclude that \(\Sigma\) is PD.
\item Show that \(\prb(d^2\leqslant c)=1-(1+c/2)e^{-c/2}\) for \(c>0\). Describe the region
\(E_\alpha=\{u\in\R^4:(u-\mu)'\Sigma^{-1}(u-\mu)\leqslant\chi^2_{4,1-\alpha}\}\), which contains \(x\) with probability \(1-\alpha\).
\item Which of the points \(a=(4,4,4,4)'\) and \(b=(3,-1,3,-1)'\) lie in \(E_{1/20}\)? For each, find the probability that
\(d^2\) exceeds its value of \((u-\mu)'\Sigma^{-1}(u-\mu)\), and compare with the Euclidean distances from \(\mu\).
\end{enumerate}
\end{exercise}
\begin{solution}
(a) Row by row,
\[
\Sigma h_1=\begin{pmatrix}9\\9\\9\\9\end{pmatrix}=9h_1,\qquad\Sigma h_2=\begin{pmatrix}5\\5\\-5\\-5\end{pmatrix}=5h_2,\qquad
\Sigma h_3=\begin{pmatrix}1\\-1\\1\\-1\end{pmatrix}=h_3,\qquad\Sigma h_4=\begin{pmatrix}1\\-1\\-1\\1\end{pmatrix}=h_4
\]
The \(h_j\) are orthogonal (\(h_i'h_j=0\) for \(i\ne j\)) with \(h_j'h_j=4\), so \(q_j=\frac12h_j\) form an orthonormal basis of
eigenvectors, and the eigenvalues of \(\Sigma\) are \(9\), \(5\), \(1\), \(1\) (Theorem~\ref{r8-thm-spectral}(b),(c)). All are \(>0\), so
\(\Sigma\) is PD (Theorem~\ref{r8-thm-eigen-criterion}(a)).

(b) By Corollary~\ref{r14-cor-standard}(a), \(d^2\sim\chi^2_4\); so, by Exercise~\ref{r15-dr-chisq-tail}(b),
\(\prb(d^2\leqslant c)=1-\prb(d^2>c)=1-(1+c/2)e^{-c/2}\) for \(c>0\), and \(\prb(x\in E_\alpha)=\prb(d^2\leqslant\chi^2_{4,1-\alpha})=1-\alpha\)
(Definition~\ref{r15-def-t-F}(c)). By the paragraph after Theorem~\ref{r14-thm-mvn-density}, with \(c=(\chi^2_{4,1-\alpha})^{1/2}\),
\(E_\alpha\) is bounded by an ellipsoid with center \(\mu\), axes along \(q_1,\dots,q_4\), and semi-axes \(3c\), \(\sqrt5\,c\), \(c\),
\(c\).

(c) In the coordinates \(v_j=q_j'(u-\mu)\) of that paragraph, \((u-\mu)'\Sigma^{-1}(u-\mu)=v_1^2/9+v_2^2/5+v_3^2+v_4^2\).
\begin{itemize}
\item For \(a\), \(a-\mu=(3,3,3,3)'=3h_1\), so \(v=(6,0,0,0)'\) and the value is \(36/9=4\). For \(b\), \(b-\mu=(2,-2,2,-2)'=2h_3\), so
\(v=(0,0,4,0)'\) and the value is \(16\).
\item By Exercise~\ref{r15-dr-chisq-tail}(b), \(8<\chi^2_{4,19/20}<10\). So \(4<\chi^2_{4,19/20}\): \(a\) lies in \(E_{1/20}\); and
\(16>\chi^2_{4,19/20}\): \(b\) does not. By (b), \(\prb(d^2>4)=3e^{-2}\) and \(\prb(d^2>16)=9e^{-8}\).
\item The Euclidean distances are \(\lVert a-\mu\rVert=6\) and \(\lVert b-\mu\rVert=4\): the point farther from \(\mu\) is the less
unusual one, because it lies in the direction \(h_1\) of the largest variance, and \(b\) in a direction of the smallest.
\end{itemize}
\end{solution}

\begin{exercise}\label{r15-dr-helmert}
Let \(C_4\) be the centering matrix (Definition~\ref{r1-def-ones}) and \(z\sim N_4(\mu,\sigma^2I_4)\) (\(\sigma>0\)), with
\(\bar z=\frac14\one'z\).
\begin{enumerate}
\item Find an orthonormal basis \(q_1,q_2,q_3\) of \(\mathcal N(C_4-I_4)\) by row reduction and Gram--Schmidt, and compare it with
the rows \(p_2'\), \(p_3'\), \(p_4'\) of the Helmert matrix \(K_4\) (Example~\ref{r6-exm-helmert}).
\item Write \(z'C_4z=u_1^2+u_2^2+u_3^2\) with \(u_i=q_i'z\), and show that \(u_1\), \(u_2\), \(u_3\) are independent with variance
\(\sigma^2\).
\item Let \(A_1=\bar J_4\) and \(A_{i+1}=q_iq_i'\) (\(i=1,2,3\)). Show that \(A_1+A_2+A_3+A_4=I_4\) and that the ranks add up to \(4\),
and conclude with Cochran's theorem that \(4\bar z^2\), \(u_1^2\), \(u_2^2\), \(u_3^2\) are independent; give their distributions
divided by \(\sigma^2\).
\item For \(\mu=(2,3,7,4)'\) and \(\sigma=1\), find the distribution of \(z'C_4z\), its mean and its variance, and show that
\(z'C_4z\) and \(\bar z\) are independent.
\end{enumerate}
\end{exercise}
\begin{solution}
(a) \(C_4-I_4=-\frac14J_4\) has the same null space as \(J_4\), which reduces as follows (the three operations use row \(1\)
and change different rows):
\[
\begin{pmatrix}1&1&1&1\\1&1&1&1\\1&1&1&1\\1&1&1&1\end{pmatrix}
\xrightarrow[R_4\leftarrow R_4-R_1]{R_2\leftarrow R_2-R_1,\ R_3\leftarrow R_3-R_1}
\begin{pmatrix}1&1&1&1\\0&0&0&0\\0&0&0&0\\0&0&0&0\end{pmatrix}
\]
So \(x_1=-x_2-x_3-x_4\), and the special solutions \(a_1=(-1,1,0,0)'\), \(a_2=(-1,0,1,0)'\), \(a_3=(-1,0,0,1)'\) form a basis
(Theorem~\ref{thm:r4-four-bases}(c)). Gram--Schmidt (Theorem~\ref{r6-thm-gram-schmidt}):
\begin{itemize}
\item \(b_1=a_1\), \(\lVert b_1\rVert^2=2\);
\item \(a_2'b_1=1\), \(b_2=a_2-\frac12b_1=(-\frac12,-\frac12,1,0)'\), \(\lVert b_2\rVert^2=\frac32\);
\item \(a_3'b_1=1\), \(a_3'b_2=\frac12\), \(b_3=a_3-\frac12b_1-\frac13b_2=(-\frac13,-\frac13,-\frac13,1)'\), \(\lVert b_3\rVert^2=\frac43\).
\end{itemize}
So \(q_1=(-1,1,0,0)'/\sqrt2\), \(q_2=(-1,-1,2,0)'/\sqrt6\) and \(q_3=(-1,-1,-1,3)'/\sqrt{12}\): these are \(-p_2\), \(-p_3\) and \(-p_4\) for
the rows \(p_2'=(1,-1,0,0)/\sqrt2\), \(p_3'=(1,1,-2,0)/\sqrt6\) and \(p_4'=(1,1,1,-3)/\sqrt{12}\) of \(K_4\).

(b) By Example~\ref{r6-exm-helmert}(b), \(C_4=p_2p_2'+p_3p_3'+p_4p_4'=q_1q_1'+q_2q_2'+q_3q_3'\), so
\(z'C_4z=(q_1'z)^2+(q_2'z)^2+(q_3'z)^2\), with \(u_1=(z_2-z_1)/\sqrt2\), \(u_2=(2z_3-z_1-z_2)/\sqrt6\) and
\(u_3=(3z_4-z_1-z_2-z_3)/\sqrt{12}\). For \(P_1=[q_1\ q_2\ q_3]\), \(P_1'P_1=I_3\), so \(u=P_1'z\sim N_3(P_1'\mu,\sigma^2I_3)\)
(Theorem~\ref{r14-thm-mvn-linear}): \(u_1\), \(u_2\), \(u_3\) are independent with variance \(\sigma^2\) (Theorem~\ref{r14-thm-mvn-iid}(a)).

(c) By (b) and Definition~\ref{r1-def-ones}, \(A_2+A_3+A_4=C_4\) and \(A_1+C_4=I_4\). The \(A_i\) are symmetric and idempotent
(Proposition~\ref{r1-prop-centering}; \((q_iq_i')^2=q_i(q_i'q_i)q_i'=q_iq_i'\)), so their ranks are their traces
(Theorem~\ref{r6-thm-idempotent}(c)): \(\tr\bar J_4=1\) and \(\tr(q_iq_i')=q_i'q_i=1\); they add up to \(4\). By
Theorem~\ref{r15-thm-cochran}, (a) \(\Rightarrow\) (b), (c), the forms \(z'A_1z=\frac14(\one'z)^2=4\bar z^2\) and
\(z'A_{i+1}z=(q_i'z)^2=u_i^2\) are independent, with \(4\bar z^2/\sigma^2\sim\chi^2_1(4\bar\mu^2/\sigma^2)\), where
\(\bar\mu=\frac14\one'\mu\), and \(u_i^2/\sigma^2\sim\chi^2_1\bigl((q_i'\mu)^2/\sigma^2\bigr)\).

(d) \(C_4\) is symmetric and idempotent of rank \(3\) (Proposition~\ref{r1-prop-centering}, Corollary~\ref{r6-cor-trace-hat}). By
Theorem~\ref{r15-thm-quadratic-chisq}, \(z'C_4z\sim\chi^2_3(\mu'C_4\mu)\), where \(\mu'C_4\mu=\sum_i(\mu_i-4)^2=4+1+9+0=14\)
(Example~\ref{r1-exm-centering}); check with (c) and Theorem~\ref{r15-thm-chisq}(c):
\((q_1'\mu)^2+(q_2'\mu)^2+(q_3'\mu)^2=\frac12+\frac{81}6+0=14\). By Theorem~\ref{r15-thm-chisq}(d), the mean is \(3+14=17\) and the
variance is \(6+56=62\). As \(\bar z=\frac14\one'z\) and \(\frac14\one'C_4=0\) (Proposition~\ref{r1-prop-centering}), \(\bar z\) and
\(z'C_4z\) are independent (Theorem~\ref{r15-thm-indep}(b)).
\end{solution}

\begin{exercise}\label{r15-dr-projections}
Let \(v=(1,1,-1,-1)'\), \(u=(1,-1,0,0)'\), \(A=\frac14vv'\) and \(B=\frac12uu'\).
\begin{enumerate}
\item Write out \(A\) and \(B\), check that \(A^2=A\), \(B^2=B\) and \(AB=0\), and find \(\rank A\) and \(\rank B\).
\item Let \(z\sim N_4(\mu,I_4)\) with \(\mu=(3,1,2,2)'\). Find the distributions of \(z'Az\) and \(z'Bz\), with their means and variances,
and explain why \(z'Az\) and \(z'Bz\) are independent.
\item Let \(\mu=0\). Find a basis of the set of \(c\in\R^4\) such that \(c'z\) is independent of \(z'Az\).
\end{enumerate}
\end{exercise}
\begin{solution}
(a)
\[
A=\frac14\begin{pmatrix}1&1&-1&-1\\1&1&-1&-1\\-1&-1&1&1\\-1&-1&1&1\end{pmatrix},\qquad
B=\frac12\begin{pmatrix}1&-1&0&0\\-1&1&0&0\\0&0&0&0\\0&0&0&0\end{pmatrix}
\]
\(A^2=\frac1{16}v(v'v)v'=\frac4{16}vv'=A\), \(B^2=\frac14u(u'u)u'=\frac24uu'=B\), and \(AB=\frac18v(v'u)u'=0\), as \(v'u=1-1+0+0=0\).
Both are symmetric and idempotent, so their ranks are their traces (Theorem~\ref{r6-thm-idempotent}(c)):
\(\rank A=\frac14v'v=1\) and \(\rank B=\frac12u'u=1\).

(b) \(\mu'A\mu=\frac14(v'\mu)^2=\frac14(3+1-2-2)^2=0\) and \(\mu'B\mu=\frac12(u'\mu)^2=\frac12(3-1)^2=2\). By
Theorem~\ref{r15-thm-quadratic-chisq}, \(z'Az\sim\chi^2_1\) and \(z'Bz\sim\chi^2_1(2)\); by Theorem~\ref{r15-thm-chisq}(d), the means are
\(1\) and \(3\), and the variances \(2\) and \(2+8=10\). They are independent by Theorem~\ref{r15-thm-indep}(a), as \(AB=0\).

(c) By Theorem~\ref{r15-thm-indep-converse}(b), \(c'z\) is independent of \(z'Az\) exactly when \(Ac=0\). Row-reduce \(4A=vv'\):
\[
\begin{pmatrix}1&1&-1&-1\\1&1&-1&-1\\-1&-1&1&1\\-1&-1&1&1\end{pmatrix}
\xrightarrow[R_4\leftarrow R_4+R_1]{R_2\leftarrow R_2-R_1,\ R_3\leftarrow R_3+R_1}
\begin{pmatrix}1&1&-1&-1\\0&0&0&0\\0&0&0&0\\0&0&0&0\end{pmatrix}
\]
So \(c_1=-c_2+c_3+c_4\), and the special solutions \((-1,1,0,0)'\), \((1,0,1,0)'\), \((1,0,0,1)'\) form a basis of \(\mathcal N(A)\)
(Theorem~\ref{thm:r4-four-bases}(c)).
\end{solution}

\begin{exercise}\label{r15-dr-pd}
Let \(y\sim N_4(\mu,\Sigma)\) with
\[
\Sigma=\begin{pmatrix}1&1&0&1\\1&2&1&2\\0&1&2&1\\1&2&1&3\end{pmatrix},\qquad\mu=\begin{pmatrix}1\\0\\-1\\1\end{pmatrix}
\]
\begin{enumerate}
\item Show that \(\Sigma\) is positive definite, and find \(\Sigma^{-1}\) by row-reducing \([\Sigma\mid I_4]\).
\item Find the distribution of \(y'\Sigma^{-1}y\), its mean and its variance.
\item Let \(A=\Sigma^{-1}-e_1e_1'\). Compute \(A\Sigma\), check the conditions of Theorem~\ref{r15-thm-chisq-pd}(d), and find the
distribution of \(y'Ay\), its mean and its variance.
\end{enumerate}
\end{exercise}
\begin{solution}
(a) Row reduction (the two operations of the first arrow use row \(1\), those of the second the new row \(2\), those of the
third row \(4\); each pair changes two different rows):
\[
\begin{gathered}
\left(\begin{array}{cccc|cccc}1&1&0&1&1&0&0&0\\1&2&1&2&0&1&0&0\\0&1&2&1&0&0&1&0\\1&2&1&3&0&0&0&1\end{array}\right)
\xrightarrow[R_4\leftarrow R_4-R_1]{R_2\leftarrow R_2-R_1}
\left(\begin{array}{cccc|cccc}1&1&0&1&1&0&0&0\\0&1&1&1&-1&1&0&0\\0&1&2&1&0&0&1&0\\0&1&1&2&-1&0&0&1\end{array}\right)
\\
\xrightarrow[R_4\leftarrow R_4-R_2]{R_3\leftarrow R_3-R_2}
\left(\begin{array}{cccc|cccc}1&1&0&1&1&0&0&0\\0&1&1&1&-1&1&0&0\\0&0&1&0&1&-1&1&0\\0&0&0&1&0&-1&0&1\end{array}\right)
\xrightarrow[R_2\leftarrow R_2-R_4]{R_1\leftarrow R_1-R_4}
\left(\begin{array}{cccc|cccc}1&1&0&0&1&1&0&-1\\0&1&1&0&-1&2&0&-1\\0&0&1&0&1&-1&1&0\\0&0&0&1&0&-1&0&1\end{array}\right)
\\
\xrightarrow{R_2\leftarrow R_2-R_3}
\left(\begin{array}{cccc|cccc}1&1&0&0&1&1&0&-1\\0&1&0&0&-2&3&-1&-1\\0&0&1&0&1&-1&1&0\\0&0&0&1&0&-1&0&1\end{array}\right)
\xrightarrow{R_1\leftarrow R_1-R_2}
\left(\begin{array}{cccc|cccc}1&0&0&0&3&-2&1&0\\0&1&0&0&-2&3&-1&-1\\0&0&1&0&1&-1&1&0\\0&0&0&1&0&-1&0&1\end{array}\right)
\end{gathered}
\]
So
\[
\Sigma^{-1}=\begin{pmatrix}3&-2&1&0\\-2&3&-1&-1\\1&-1&1&0\\0&-1&0&1\end{pmatrix}
\]
Check: \(\Sigma\Sigma^{-1}=I_4\); for instance row \(4\) of \(\Sigma\) times column \(2\) of \(\Sigma^{-1}\) is \(-2+6-1-3=0\). The leading
principal minors are \(\Delta_1(\Sigma)=1\), \(\Delta_2(\Sigma)=2-1=1\) and \(\Delta_3(\Sigma)=1\cdot(4-1)-1\cdot(2-0)=1\) (expansion along
row \(1\), Theorem~\ref{r5d-thm-expansion}); and \(\Delta_4(\Sigma)=\det\Sigma=1\), as the first two arrows add multiples of rows to
other rows (Definition~\ref{r5d-def-det}(b)) and leave the left block upper triangular with diagonal entries \(1\)
(Theorem~\ref{r5d-thm-basic}(c)). All are positive, so \(\Sigma\) is positive definite (Theorem~\ref{r8-thm-leading-minors}).

(b) For \(A=\Sigma^{-1}\), \(A\Sigma=I_4\) is idempotent of rank \(4\). By Theorem~\ref{r15-thm-chisq-pd}(b),
\(y'\Sigma^{-1}y\sim\chi^2_4(\mu'\Sigma^{-1}\mu)\); \(\Sigma^{-1}\mu=(3-1+0,\ -2+1-1,\ 1+0-1+0,\ 0+0+0+1)'=(2,-2,0,1)'\), so
\(\mu'\Sigma^{-1}\mu=2+0+0+1=3\). The mean is \(4+3=7\), and the variance is \(8+12=20\) (Theorem~\ref{r15-thm-chisq}(d)).

(c)
\[
A=\begin{pmatrix}2&-2&1&0\\-2&3&-1&-1\\1&-1&1&0\\0&-1&0&1\end{pmatrix},\qquad
A\Sigma=I_4-e_1(e_1'\Sigma)=\begin{pmatrix}0&-1&0&-1\\0&1&0&0\\0&0&1&0\\0&0&0&1\end{pmatrix}
\]
Squaring, \((A\Sigma)^2=A\Sigma\) (row \(1\) of the square is \(-1\) times row \(2\) of \(A\Sigma\) plus \(-1\) times its row \(4\)); so
\(\tr(A\Sigma)=\tr\bigl((A\Sigma)^2\bigr)=3\), and \(\rank(A\Sigma)=3\) (the first column is \(0\), the other three are independent). By
Theorem~\ref{r15-thm-chisq-pd}(d),(b), \(y'Ay\sim\chi^2_3(\mu'A\mu)\), and \(A\mu=\Sigma^{-1}\mu-\mu_1e_1=(1,-2,0,1)'\), so
\(\mu'A\mu=1+0+0+1=2\). So \(y'Ay\sim\chi^2_3(2)\), with mean \(5\) and variance \(6+8=14\) (Theorem~\ref{r15-thm-chisq}(d)).
\end{solution}

\begin{exercise}\label{r15-dr-decomposition}
Let \(z\sim N_4(\mu,I_4)\) with \(\mu=(1,1,1,1)'\), and \(A=I_4+J_4\).
\begin{enumerate}
\item Write \(A=5\bar J_4+C_4\), and give the eigenvalues of \(A\) with their multiplicities.
\item Use Theorem~\ref{r15-thm-decomposition} to write \(z'Az=5W_1+W_2\) with independent \(W_1\sim\chi^2_1(\theta_1)\) and
\(W_2\sim\chi^2_3(\theta_2)\); find \(W_1\), \(W_2\), \(\theta_1\) and \(\theta_2\).
\item Compute \(\expc(z'Az)\) and \(\var(z'Az)\) from (b), and again from Theorem~\ref{r15-thm-moments}.
\end{enumerate}
\end{exercise}
\begin{solution}
(a) \(I_4=\bar J_4+C_4\) and \(J_4=4\bar J_4\), so \(A=5\bar J_4+C_4\), a spectral decomposition (Example~\ref{r8-exm-patterned}(b)): the
eigenvalues are \(5\) (once, eigenvector \(\one\)) and \(1\) (three times, eigenspace \(\Span\{\one\}^\perp\)).

(b) Here \(\Sigma=I_4\) and \(B=A\), with \(P_1=\one/2\) for \(\gamma_1=5\) and \(P_2=[q_1\ q_2\ q_3]\) of Exercise~\ref{r15-dr-helmert} for
\(\gamma_2=1\) (\(P_1P_1'=\bar J_4\), \(P_2P_2'=C_4\)). So \(W_1=\lVert P_1'z\rVert^2=z'\bar J_4z=4\bar z^2\) and \(W_2=z'C_4z\), with
\(\theta_1=\mu'\bar J_4\mu=4\) and \(\theta_2=\mu'C_4\mu=0\).

(c) From (b) and Theorem~\ref{r15-thm-chisq}(d): \(\expc(z'Az)=5(1+4)+(3+0)=28\), and, as \(W_1\) and \(W_2\) are independent,
\(\var(z'Az)=25(2+16)+(6+0)=456\). From Theorem~\ref{r15-thm-moments}: \(\tr A=8\) and \(\mu'A\mu=\one'(I_4+J_4)\one=4+16=20\), so the
mean is \(28\); \(A^2=I_4+2J_4+J_4^2=I_4+6J_4\), so \(\tr(A^2)=28\) and \(\mu'A^2\mu=4+96=100\), and the variance is \(56+400=456\).
\end{solution}

\begin{exercise}\label{r15-dr-moments}
Let \(y\sim N_4(\mu,\Sigma)\) with \(\Sigma\) and \(\mu\) as in Exercise~\ref{r15-dr-pd}, and let
\[
A=\begin{pmatrix}1&0&0&0\\0&0&0&0\\0&0&0&0\\0&0&0&-1\end{pmatrix},\qquad B=\begin{pmatrix}0&0&0&0\\0&0&1&0\\0&1&0&0\\0&0&0&0\end{pmatrix}
\]
so that \(y'Ay=y_1^2-y_4^2\) and \(y'By=2y_2y_3\).
\begin{enumerate}
\item Compute \(\expc(y'Ay)\) and \(\var(y'Ay)\).
\item Compute \(\expc(y'By)\) and \(\var(y'By)\).
\item Compute \(\cov(y'Ay,y'By)\), and find \(\var\bigl(y'(A+B)y\bigr)\) in two ways.
\end{enumerate}
\end{exercise}
\begin{solution}
(a)
\[
A\Sigma=\begin{pmatrix}1&1&0&1\\0&0&0&0\\0&0&0&0\\-1&-2&-1&-3\end{pmatrix},\qquad
(A\Sigma)^2=\begin{pmatrix}0&-1&-1&-2\\0&0&0&0\\0&0&0&0\\2&5&3&8\end{pmatrix}
\]
so \(\tr(A\Sigma)=-2\) and \(\tr\bigl((A\Sigma)^2\bigr)=8\). Also \(\mu'A\mu=1-1=0\), \(A\mu=(1,0,0,-1)'\), \(A\Sigma A\mu=(0,0,0,2)'\) and
\(\mu'A\Sigma A\mu=2\). By Theorem~\ref{r15-thm-moments}, \(\expc(y'Ay)=-2\) and \(\var(y'Ay)=16+8=24\).

(b)
\[
B\Sigma=\begin{pmatrix}0&0&0&0\\0&1&2&1\\1&2&1&2\\0&0&0&0\end{pmatrix},\qquad
(B\Sigma)^2=\begin{pmatrix}0&0&0&0\\2&5&4&5\\1&4&5&4\\0&0&0&0\end{pmatrix}
\]
so \(\tr(B\Sigma)=2\) and \(\tr\bigl((B\Sigma)^2\bigr)=10\). Also \(\mu'B\mu=2\mu_2\mu_3=0\), \(B\mu=(0,-1,0,0)'\), \(B\Sigma B\mu=(0,-1,-2,0)'\) and
\(\mu'B\Sigma B\mu=2\). So \(\expc(y'By)=2\) and \(\var(y'By)=20+8=28\).

(c)
\[
A\Sigma B\Sigma=\begin{pmatrix}0&1&2&1\\0&0&0&0\\0&0&0&0\\-1&-4&-5&-4\end{pmatrix}
\]
has the trace \(-4\), and \(A\Sigma B\mu=A\Sigma(0,-1,0,0)'=(-1,0,0,2)'\), so \(\mu'A\Sigma B\mu=-1+2=1\). By Corollary~\ref{r15-cor-cov},
\(\cov(y'Ay,y'By)=-8+4=-4\), hence \(\var\bigl(y'(A+B)y\bigr)=24+28+2\cdot(-4)=44\). Directly,
\[
(A+B)\Sigma=\begin{pmatrix}1&1&0&1\\0&1&2&1\\1&2&1&2\\-1&-2&-1&-3\end{pmatrix},\qquad
\bigl((A+B)\Sigma\bigr)^2=\begin{pmatrix}0&0&1&-1\\1&3&3&2\\0&1&3&-1\\1&1&-2&4\end{pmatrix}
\]
with the trace \(10\), and \((A+B)\mu=(1,-1,0,-1)'\), \((A+B)\Sigma(A+B)\mu=(-1,-2,-3,4)'\), \(\mu'(A+B)\Sigma(A+B)\mu=-1+3+4=6\); so
\(\var\bigl(y'(A+B)y\bigr)=20+24=44\) (Theorem~\ref{r15-thm-moments}).
\end{solution}

\needspace{14\baselineskip}
\subsection*{Problems}

\begin{problem}[\citeat{casella-berger}{Exercise~5.13}, \citeat{casella-berger-solutions}{Exercise~5.13}]\label{r15-ex-unbiased-sigma}
Let \(x_1,\dots,x_n\) (\(n\geqslant2\)) be independent \(N(\mu,\sigma^2)\) variables (\(\sigma>0\)) with sample variance \(s^2\). Find a
constant \(c\) such that \(\expc(c\sqrt{s^2})=\sigma\).
\end{problem}
\begin{solution}
\begin{itemize}
\item \(q=(n-1)s^2/\sigma^2\sim\chi^2_{n-1}\) (Theorem~\ref{r15-thm-one-sample}(b)), and \(c\sqrt{s^2}=c\sqrt{\sigma^2/(n-1)}\,\sqrt q\).
\item \(q\) has the density \(f_{n-1}\) (Theorem~\ref{r15-thm-chisq}(b)); so, by the paragraph before
Theorem~\ref{r14-thm-mvn-density} and by (D2) with \(n/2\) and \(\frac12\) in place of \(c\) and \(b\),
\[
\expc(\sqrt q)=\int_0^\infty\sqrt x\,\frac{x^{(n-1)/2-1}e^{-x/2}}{\Gamma\bigl(\frac{n-1}2\bigr)2^{(n-1)/2}}\,\mathrm{d}x
=\frac{\Gamma\bigl(\frac n2\bigr)2^{n/2}}{\Gamma\bigl(\frac{n-1}2\bigr)2^{(n-1)/2}}=\frac{\sqrt2\,\Gamma\bigl(\frac n2\bigr)}{\Gamma\bigl(\frac{n-1}2\bigr)}
\]
\item So \(\expc(c\sqrt{s^2})=c\,\sigma\sqrt{2/(n-1)}\;\Gamma\bigl(\frac n2\bigr)/\Gamma\bigl(\frac{n-1}2\bigr)\), which is \(\sigma\) for
\[
c=\frac{\Gamma\bigl(\frac{n-1}2\bigr)\sqrt{n-1}}{\sqrt2\,\Gamma\bigl(\frac n2\bigr)}
\]
For instance, for \(n=2\), \(c=\Gamma(\frac12)/\bigl(\sqrt2\,\Gamma(1)\bigr)=\sqrt{\pi/2}\).\qedhere
\end{itemize}
\end{solution}

\begin{problem}[\citeat{casella-berger}{Exercise~5.15}, \citeat{casella-berger-solutions}{Exercise~5.15}]\label{r15-ex-recursion}
Let \(\bar x_n\) and \(s_n^2\) be the mean and the sample variance of numbers \(x_1,\dots,x_n\) (\(n\geqslant2\)), and let a further number
\(x_{n+1}\) be added. Show that
\begin{enumerate}
\item \(\bar x_{n+1}=(x_{n+1}+n\bar x_n)/(n+1)\);
\item \(ns_{n+1}^2=(n-1)s_n^2+\frac n{n+1}(x_{n+1}-\bar x_n)^2\).
\end{enumerate}
\end{problem}
\begin{solution}
(a) \(\bar x_{n+1}=\frac1{n+1}\bigl(\sum_{i\leqslant n}x_i+x_{n+1}\bigr)=\frac1{n+1}(n\bar x_n+x_{n+1})\).

(b)
\begin{itemize}
\item By (a), \(\bar x_{n+1}=\bar x_n+d\) with \(d=(x_{n+1}-\bar x_n)/(n+1)\). So
\(ns_{n+1}^2=\sum_{i=1}^{n+1}(x_i-\bar x_{n+1})^2=\sum_{i=1}^{n+1}\bigl((x_i-\bar x_n)-d\bigr)^2\).
\item Expanding, this is \(\sum_{i=1}^{n+1}(x_i-\bar x_n)^2-2d\sum_{i=1}^{n+1}(x_i-\bar x_n)+(n+1)d^2\). As
\(\sum_{i\leqslant n}(x_i-\bar x_n)=0\), the middle sum is \(x_{n+1}-\bar x_n=(n+1)d\), and the first sum is
\((n-1)s_n^2+(x_{n+1}-\bar x_n)^2\).
\item Hence \(ns_{n+1}^2=(n-1)s_n^2+(x_{n+1}-\bar x_n)^2-(n+1)d^2=(n-1)s_n^2+\bigl(1-\frac1{n+1}\bigr)(x_{n+1}-\bar x_n)^2\).\qedhere
\end{itemize}
\end{solution}

\begin{problem}[\citeat{casella-berger}{Exercise~5.16}, \citeat{casella-berger-solutions}{Exercise~5.16}]\label{r15-ex-construct}
Let \(x_1,x_2,x_3\) be independent, with \(x_i\sim N(i,i^2)\). Use them to construct a statistic with the distribution
\begin{enumerate}
\item \(\chi^2_3\);
\item \(t_2\);
\item \(F_{1,2}\).
\end{enumerate}
\end{problem}
\begin{solution}
Let \(z_i=(x_i-i)/i\). By Theorem~\ref{r14-thm-mvn-iid}(a), \(x\sim N_3(\mu,D)\) with \(\mu=(1,2,3)'\) and \(D=\diag(1,4,9)\), so
\(z=D^{-1/2}(x-\mu)\sim N_3(0,I_3)\) (Theorem~\ref{r14-thm-mvn-linear}): \(z_1,z_2,z_3\) are independent \(N(0,1)\) variables
(Theorem~\ref{r14-thm-mvn-iid}(a)), and \(z_1\) and \((z_2,z_3)'\) are independent (Theorem~\ref{r14-thm-mvn-indep}(a)).

(a) \(z_1^2+z_2^2+z_3^2=\sum_{i=1}^3\bigl((x_i-i)/i\bigr)^2\sim\chi^2_3\) (Definition~\ref{r15-def-chisq}).

(b) \(z_2^2+z_3^2\sim\chi^2_2\) is independent of \(z_1\) (N3), so
\[
\frac{z_1}{\sqrt{(z_2^2+z_3^2)/2}}=\frac{x_1-1}{\sqrt{\bigl(\frac14(x_2-2)^2+\frac19(x_3-3)^2\bigr)/2}}\sim t_2
\]
(Definition~\ref{r15-def-t-F}).

(c) The square of the statistic in (b), \(z_1^2/\bigl((z_2^2+z_3^2)/2\bigr)\), has the distribution \(F_{1,2}\)
(Theorem~\ref{r15-thm-t-symmetric}(b)).
\end{solution}

\begin{problem}[\citeat{casella-berger}{Exercise~5.17(b)}, \citeat{casella-berger-solutions}{Exercise~5.17(b)}]\label{r15-ex-F-moments}
Let \(F\sim F_{p,q}\). Show that
\[
\expc(F)=\frac q{q-2}\quad(q>2),\qquad\var(F)=2\Bigl(\frac q{q-2}\Bigr)^2\frac{p+q-2}{p(q-4)}\quad(q>4)
\]
\end{problem}
\begin{solution}
\begin{itemize}
\item Write \(F=(U/p)/(V/q)\) with \(U\sim\chi^2_p\) and \(V\sim\chi^2_q\) independent; \(V>0\) with probability \(1\), as \(f_q=0\)
on \((-\infty,0]\). For \(k,l\in\{1,2\}\), \(U^k\) and \(V^{-l}\) are independent (N3), so \(\expc(U^kV^{-l})=\expc(U^k)\expc(V^{-l})\)
(N7), where \(\expc(V^{-l})=\int_0^\infty v^{-l}f_q(v)\,\mathrm{d}v\) (the paragraph before Theorem~\ref{r14-thm-mvn-density}),
whenever these expectations are finite.
\item By (D2), for \(q>2l\), \(\int_0^\infty v^{-l}f_q(v)\,\mathrm{d}v=\frac{\Gamma(q/2-l)\,2^{q/2-l}}{\Gamma(q/2)\,2^{q/2}}\). With
\(\Gamma(q/2)=(q/2-1)\Gamma(q/2-1)\) and \(\Gamma(q/2-1)=(q/2-2)\Gamma(q/2-2)\) (D2), this gives \(\expc(V^{-1})=\frac1{q-2}\) for
\(q>2\) and \(\expc(V^{-2})=\frac1{(q-2)(q-4)}\) for \(q>4\).
\item \(\expc(U)=p\) and \(\expc(U^2)=\var(U)+p^2=2p+p^2\) (Theorem~\ref{r15-thm-chisq}(d)).
\item So \(\expc(F)=\frac qp\cdot p\cdot\frac1{q-2}=\frac q{q-2}\) and \(\expc(F^2)=\frac{q^2}{p^2}\cdot p(p+2)\cdot\frac1{(q-2)(q-4)}\),
whence
\[
\var(F)=\frac{q^2(p+2)}{p(q-2)(q-4)}-\frac{q^2}{(q-2)^2}=\frac{q^2}{(q-2)^2}\cdot\frac{(p+2)(q-2)-p(q-4)}{p(q-4)}
=2\Bigl(\frac q{q-2}\Bigr)^2\frac{p+q-2}{p(q-4)}\qedhere
\]
\end{itemize}
\end{solution}

\begin{problem}[\citeat{casella-berger}{Exercise~5.18(a)}, \citeat{casella-berger-solutions}{Exercise~5.18(a)}]\label{r15-ex-t-moments}
Let \(T\sim t_p\). Show that \(\expc(T)=0\) if \(p>1\), and \(\var(T)=p/(p-2)\) if \(p>2\).
\end{problem}
\begin{solution}
\begin{itemize}
\item Write \(T=Z/\sqrt{V/p}\) with \(Z\sim N(0,1)\) and \(V\sim\chi^2_p\) independent. As in Problem~\ref{r15-ex-F-moments},
\(Z\) and \(V^{-1/2}\) are independent (N3), so \(\expc(T)=\sqrt p\,\expc(Z)\expc(V^{-1/2})\) (N7), where
\(\expc(Z)=0\) and \(\expc(V^{-1/2})=\int_0^\infty v^{-1/2}f_p(v)\,\mathrm{d}v=\frac{\Gamma((p-1)/2)\,2^{(p-1)/2}}{\Gamma(p/2)\,2^{p/2}}\) is
finite for \(p>1\) (D2). So \(\expc(T)=0\).
\item \(T^2\sim F_{1,p}\) (Theorem~\ref{r15-thm-t-symmetric}(b)), so for \(p>2\), \(\var(T)=\expc(T^2)=p/(p-2)\), the mean of
\(F_{1,p}\) (Problem~\ref{r15-ex-F-moments}).\qedhere
\end{itemize}
\end{solution}

\begin{problem}[\citeat{casella-berger}{Exercise~5.19}, \citeat{casella-berger-solutions}{Exercise~5.19}]\label{r15-ex-stochastic-order}
A random variable \(u\) is \textbf{stochastically larger} than \(v\) if \(\prb(u>c)\geqslant\prb(v>c)\) for every \(c\),
with strict inequality for some \(c\).
\begin{enumerate}
\item Let \(p>q\geqslant1\), \(u\sim\chi^2_p\) and \(v\sim\chi^2_q\). Show that \(\prb(u>c)>\prb(v>c)\) for every \(c>0\), so
that \(\chi^2_p\) is stochastically larger than \(\chi^2_q\).
\item Let \(k_1>k_2\geqslant1\), \(F_1\sim F_{k_1,\nu}\) and \(F_2\sim F_{k_2,\nu}\). Show that \(\prb(k_1F_1>c)>\prb(k_2F_2>c)\) for every
\(c>0\), so that \(k_1F_1\) is stochastically larger than \(k_2F_2\).
\item Show that \(kF_{k,\nu,1-\alpha}>(k-1)F_{k-1,\nu,1-\alpha}\) for \(k\geqslant2\) and \(\alpha\in(0,1)\).
\end{enumerate}
\end{problem}
\begin{solution}
(a)
\begin{itemize}
\item The two probabilities depend only on the distributions, so we may take \(v=g_1^2+\dots+g_q^2\) and put
\(d=g_{q+1}^2+\dots+g_p^2\) for \(g\sim N_p(0,I_p)\) (Theorem~\ref{r14-thm-mvn-density}(b)): \(v\sim\chi^2_q\) and
\(d\sim\chi^2_{p-q}\) (Definition~\ref{r15-def-chisq}) are independent (Theorem~\ref{r14-thm-mvn-indep}(a), (N3)). Then
\(v+d\sim\chi^2_p\) (Theorem~\ref{r15-thm-chisq}(c)), so \(\prb(u>c)=\prb(v+d>c)\).
\item For \(c>0\), \(\prb(v+d>c)-\prb(v>c)=\prb(v\leqslant c<v+d)\), as \(d>0\) with probability \(1\). As \(v\)
and \(d\) are independent, this is the integral of their joint density \(f_q(x)f_{p-q}(y)\) (the paragraph before
Theorem~\ref{r14-thm-mvn-density}), which is positive on \((0,\infty)^2\), over the region \(\{(x,y):0<x\leqslant c,\ y>c-x\}\); so it
is positive. For \(c\leqslant0\) both probabilities are \(1\).
\end{itemize}
(b)
\begin{itemize}
\item Let \(g\sim N_m(0,I_m)\) with \(m=k_1+\nu\) (Theorem~\ref{r14-thm-mvn-density}(b)), and let \(v=g_1^2+\dots+g_{k_2}^2\),
\(d=g_{k_2+1}^2+\dots+g_{k_1}^2\) and \(w=g_{k_1+1}^2+\dots+g_m^2\). By Definition~\ref{r15-def-chisq}, \(v\sim\chi^2_{k_2}\),
\(d\sim\chi^2_{k_1-k_2}\), \(v+d\sim\chi^2_{k_1}\) and \(w\sim\chi^2_\nu\).
\item The three blocks of \(g\) of lengths \(k_2\), \(k_1-k_2\) and \(\nu\) are independent (Theorem~\ref{r14-thm-mvn-indep}(b)), and
so are the vectors of its first \(k_1\) and of its last \(\nu\) coordinates (Theorem~\ref{r14-thm-mvn-indep}(a)). So \(v\), \(d\),
\(w\) are independent, and so are \(v+d\) and \(w\) (N3). By
Definition~\ref{r15-def-t-F}, \(k_1F_1\) has the distribution of \((v+d)/(w/\nu)\), and \(k_2F_2\) that of \(v/(w/\nu)\).
\item For \(c>0\), \(\prb\bigl((v+d)/(w/\nu)>c\bigr)-\prb\bigl(v/(w/\nu)>c\bigr)=\prb(v\leqslant cw/\nu<v+d)>0\), as in
(a), with the density \(f_{k_2}(x)f_{k_1-k_2}(y)f_\nu(s)\) of \((v,d,w)'\) and the region \(\{(x,y,s):0<x\leqslant cs/\nu<x+y\}\); for
\(c\leqslant0\) both probabilities are \(1\).
\end{itemize}

(c) Let \(F_k\sim F_{k,\nu}\), \(F_{k-1}\sim F_{k-1,\nu}\) and \(c=(k-1)F_{k-1,\nu,1-\alpha}\); \(c>0\) (Definition~\ref{r15-def-t-F}). Then
\(\prb(kF_k>kF_{k,\nu,1-\alpha})=\alpha=\prb\bigl((k-1)F_{k-1}>c\bigr)\), and by (b) with \(k_1=k\), \(k_2=k-1\),
\(\prb(kF_k>c)>\alpha\). As
\(x\mapsto\prb(kF_k>x)\) is decreasing, \(c<kF_{k,\nu,1-\alpha}\).
\end{solution}

\begin{problem}[\citeat{casella-berger}{Exercise~5.22}, \citeat{casella-berger-solutions}{Exercise~5.22}]\label{r15-ex-min-square}
Let \(x\) and \(y\) be independent \(N(0,1)\) variables, and \(z=\min(x,y)\). Show that \(z^2\sim\chi^2_1\).
\end{problem}
\begin{solution}
Let \(\Phi(a)=\prb(x\leqslant a)=\int_{-\infty}^a\varphi\). As \(\varphi(-s)=\varphi(s)\), substituting \(s\mapsto-s\) (D1) gives
\(\prb(x\geqslant a)=\int_a^\infty\varphi=\Phi(-a)\); and \(y\) has the same distribution function. For \(c>0\), using the
independence of \(x\) and \(y\),
\[
\begin{aligned}
\prb(z^2\leqslant c)&=\prb(z\geqslant-\sqrt c)-\prb(z>\sqrt c)
=\prb(x\geqslant-\sqrt c)\prb(y\geqslant-\sqrt c)-\prb(x>\sqrt c)\prb(y>\sqrt c)\\
&=\Phi(\sqrt c)^2-\Phi(-\sqrt c)^2=\bigl(\Phi(\sqrt c)-\Phi(-\sqrt c)\bigr)\bigl(\Phi(\sqrt c)+\Phi(-\sqrt c)\bigr)=\prb(x^2\leqslant c)
\end{aligned}
\]
since \(\Phi(\sqrt c)+\Phi(-\sqrt c)=1\) and \(\Phi(\sqrt c)-\Phi(-\sqrt c)=\prb(-\sqrt c\leqslant x\leqslant\sqrt c)\). For \(c\leqslant0\) both
probabilities are \(0\). So \(z^2\) and \(x^2\) have the same distribution function, and hence the same distribution (D3). As \(x^2\sim\chi^2_1\) (Definition~\ref{r15-def-chisq}), \(z^2\sim\chi^2_1\).
\end{solution}

\begin{problem}[\citeat{searle-khuri}{Exercise~10.11}]\label{r15-ex-bilinear}
Let \(u=(x',y')'\sim N_{m+n}(\mu,V)\) with \(x\in\R^m\) and \(y\in\R^n\), partitioned accordingly as
\[
\mu=\begin{pmatrix}\mu_1\\\mu_2\end{pmatrix},\qquad V=\begin{pmatrix}V_{11}&V_{12}\\V_{12}'&V_{22}\end{pmatrix}
\]
and let \(A\in\R^{m\times n}\). Find
\begin{enumerate}
\item \(\expc(x'Ay)\);
\item \(\var(x'Ay)\).
\end{enumerate}
\end{problem}
\begin{solution}
\(x'Ay=\frac12u'Bu\) with the symmetric matrix \(B\in\R^{(m+n)\times(m+n)}\) below, since \(u'Bu=x'Ay+y'A'x=2x'Ay\)
(Theorem~\ref{r5b-thm-block}(c)); by the block product,
\[
B=\begin{pmatrix}0&A\\A'&0\end{pmatrix},\qquad BV=\begin{pmatrix}AV_{12}'&AV_{22}\\A'V_{11}&A'V_{12}\end{pmatrix},\qquad
B\mu=\begin{pmatrix}A\mu_2\\A'\mu_1\end{pmatrix}
\]
(a) By Theorem~\ref{r15-thm-moments}, \(\expc(x'Ay)=\frac12\bigl(\tr(BV)+\mu'B\mu\bigr)\). Here
\(\tr(BV)=\tr(AV_{12}')+\tr(A'V_{12})=2\tr(AV_{12}')\) (Theorem~\ref{r1-thm-trace}(a),(b)) and \(\mu'B\mu=\mu_1'A\mu_2+\mu_2'A'\mu_1=2\mu_1'A\mu_2\);
so \(\expc(x'Ay)=\mu_1'A\mu_2+\tr(AV_{12}')\).

(b) By Theorem~\ref{r15-thm-moments}, \(\var(x'Ay)=\frac14\bigl(2\tr((BV)^2)+4\mu'BVB\mu\bigr)\).
\begin{itemize}
\item The diagonal blocks of \((BV)^2\) are \(AV_{12}'AV_{12}'+AV_{22}A'V_{11}\) and \(A'V_{11}AV_{22}+A'V_{12}A'V_{12}\). As
\(\tr(A'V_{11}AV_{22})=\tr(AV_{22}A'V_{11})\) and \(\tr(A'V_{12}A'V_{12})=\tr(V_{12}'AV_{12}'A)=\tr(AV_{12}'AV_{12}')\)
(Theorem~\ref{r1-thm-trace}), \(\frac12\tr\bigl((BV)^2\bigr)=\tr(AV_{12}'AV_{12}')+\tr(AV_{22}A'V_{11})\).
\item \(\mu'BVB\mu=(B\mu)'V(B\mu)=\mu_2'A'V_{11}A\mu_2+2\mu_2'A'V_{12}A'\mu_1+\mu_1'AV_{22}A'\mu_1\), and the middle number equals its
transpose \(\mu_1'AV_{12}'A\mu_2\).
\item Hence
\(\var(x'Ay)=\tr(AV_{12}'AV_{12}')+\tr(AV_{22}A'V_{11})+\mu_1'AV_{22}A'\mu_1+\mu_2'A'V_{11}A\mu_2+2\mu_1'AV_{12}'A\mu_2\).\qedhere
\end{itemize}
\end{solution}

\begin{problem}[\citeat{bapat}{Ch.~8, Exercise~16}]\label{r15-ex-idempotent-difference}
Let \(A,B\in\R^{n\times n}\) be symmetric and idempotent, with \(A-B\) positive semidefinite.
\begin{enumerate}
\item Show that \(A-B\) is idempotent.
\item Let \(z\sim N_n(0,I_n)\) and \(A\ne B\). Show that \(z'Az-z'Bz\sim\chi^2_{\rank A-\rank B}\).
\end{enumerate}
\end{problem}
\begin{solution}
(a)
\begin{itemize}
\item Let \(r=\rank A\). If \(r=0\), then \(A=0\), and both \(-B\) and \(B=B'B\) are positive semidefinite (Theorem~\ref{r8-thm-XtX}); so
\(x'Bx=0\) for every \(x\), \(B=0\) (Proposition~\ref{r8-prop-qf-symmetric}(b)), and \(A-B=0\). If \(r=n\), then \(A=I_n\)
(Theorem~\ref{r8-thm-idempotent-eigen}(a), Theorem~\ref{r8-thm-spectral}), and \(A-B=I_n-B\) is idempotent
(Theorem~\ref{r6-thm-idempotent}(b)). Let \(1\leqslant r<n\).
\item By Theorem~\ref{r8-thm-idempotent-eigen}(c), \(P'AP=\diag(I_r,0)\) for an orthogonal \(P\), so \(A=P\diag(I_r,0)P'\). Put \(\tilde B=P'BP\): it is
symmetric, idempotent (\(\tilde B^2=P'B^2P=\tilde B\)) and positive semidefinite (\(\tilde B=\tilde B'\tilde B\), Theorem~\ref{r8-thm-XtX}), and
\(\diag(I_r,0)-\tilde B=P'(A-B)P\) is positive semidefinite (Theorem~\ref{r8-thm-congruence}(a)). Partition
\[
\tilde B=\begin{pmatrix}B_{11}&B_{12}\\B_{12}'&B_{22}\end{pmatrix},\qquad B_{11}\in\R^{r\times r}
\]
\item \(B_{22}\) and \(-B_{22}\) are principal submatrices of positive semidefinite matrices, hence positive semidefinite
(Theorem~\ref{r8-thm-principal-sub}); so \(B_{22}=0\) (Proposition~\ref{r8-prop-qf-symmetric}(b)).
\item Then \(B_{12}=0\): for \(s\in\R\), \(e\in\R^r\) and \(f\in\R^{n-r}\), the vector \(x=(se',f')'\) gives
\(0\leqslant x'\tilde Bx=s^2e'B_{11}e+2s\,e'B_{12}f\), which is negative for small \(s\) of the sign opposite to \(e'B_{12}f\) unless
\(e'B_{12}f=0\); taking for \(e\) and \(f\) the columns of identity matrices gives \(B_{12}=0\).
\item So \(\tilde B=\diag(B_{11},0)\), and \(\tilde B^2=\tilde B\) gives \(B_{11}^2=B_{11}\) (Proposition~\ref{r5b-prop-block-tri}(b)). Then
\(I_r-B_{11}\) is idempotent (Theorem~\ref{r6-thm-idempotent}(b)), hence so is \(P'(A-B)P=\diag(I_r-B_{11},0)\), and
\((A-B)^2=P\diag(I_r-B_{11},0)^2P'=A-B\).
\end{itemize}
(b) \(A-B\) is symmetric and idempotent by (a), and \(A-B\ne0\); so
\(\rank(A-B)=\tr(A-B)=\tr A-\tr B=\rank A-\rank B\geqslant1\) (Theorems~\ref{r6-thm-idempotent}(c) and~\ref{r1-thm-trace}(a)). By
Theorem~\ref{r15-thm-quadratic-chisq} (with \(\mu=0\) and \(\sigma=1\)), \(z'Az-z'Bz=z'(A-B)z\sim\chi^2_{\rank A-\rank B}\).
\end{solution}

\clearpage
\def\unitnumber{12}
\setcounter{theorem}{0}\setcounter{problem}{0}
\section*{Chapter 12. Generalized inverses}

\subsection*{Generalized inverses}

\begin{definition}[\citeat{bapat}{\S4.1}, \citeat{harville}{\S9.1, Lemma~9.1.1}]\label{r10-def-ginverse}
Let \(A\in\R^{m\times n}\). A matrix \(G\in\R^{n\times m}\) with \(AGA=A\) is a \textbf{generalized inverse} (or
\textbf{g-inverse}) of \(A\); the symbol \(A^-\) denotes any g-inverse of \(A\). If \(A\in\R^{n\times n}\) is
invertible, then \(A^{-1}\) is a g-inverse of \(A\), and it is the only one: \(AGA=A\) gives
\(G=A^{-1}(AGA)A^{-1}=A^{-1}AA^{-1}=A^{-1}\).
\end{definition}

\begin{theorem}[\citeat{bapat}{Result~4.1}, \citeat{harville}{Lemma~9.5.1}]\label{r10-thm-ginv-solves}
Let \(A\in\R^{m\times n}\) and \(G\in\R^{n\times m}\).
\begin{enumerate}
\item \(G\) is a g-inverse of \(A\) if and only if, for every \(y\in\mathcal C(A)\), \(x=Gy\) is a solution of
\(Ax=y\).
\item If \(G\) is a g-inverse of \(A\) and \(y\in\R^m\), then \(Ax=y\) is consistent if and only if \(AGy=y\).
\end{enumerate}
\end{theorem}
\begin{proof}
(a)
\begin{itemize}
\item Let \(AGA=A\) and \(y\in\mathcal C(A)\), say \(y=Az\). Then \(A(Gy)=AGAz=Az=y\).
\item Conversely, let \(AGy=y\) for every \(y\in\mathcal C(A)\). For \(z\in\R^n\), \(y=Az\) lies in \(\mathcal C(A)\), so
\(AGAz=Az\). With \(z=e_j\), column \(j\) of \(AGA\) equals column \(j\) of \(A\) (Theorem~\ref{r1-thm-views}(b)), for
\(j=1,\dots,n\); so \(AGA=A\).
\end{itemize}
(b) \(Ax=y\) is consistent if and only if \(y\in\mathcal C(A)\) (Corollary~\ref{cor:r4-ax-b}(a)). If
\(y\in\mathcal C(A)\), then \(AGy=y\) by (a). Conversely, if \(AGy=y\), then \(x=Gy\) is a solution.
\end{proof}

\begin{theorem}[\citeat{bapat}{\S4.1}, \citeat{harville}{Theorem~9.1.3, Corollary~9.1.4}]\label{r10-thm-ginv-exists}
Let \(A\in\R^{m\times n}\) have rank \(r\).
\begin{enumerate}
\item Let \(r\geqslant1\), let \(A=BC\) be a full-rank factorization (\(B\in\R^{m\times r}\), \(C\in\R^{r\times n}\)), and let
\(L\in\R^{r\times m}\) and \(K\in\R^{n\times r}\) satisfy \(LB=I_r\) and \(CK=I_r\). Then \(G=KL\) is a g-inverse of \(A\).
Such \(L\) and \(K\) exist, for instance \(L=(B'B)^{-1}B'\) and \(K=C'(CC')^{-1}\).
\item If \(r=0\), every \(G\in\R^{n\times m}\) is a g-inverse of \(A\).
\end{enumerate}
In particular, every matrix has a g-inverse.
\end{theorem}
\begin{proof}
(a) \(AGA=BC\,KL\,BC=B(CK)(LB)C=BI_rI_rC=BC=A\). The two matrices given satisfy \(LB=I_r\) and \(CK=I_r\)
(Theorem~\ref{thm:r4-full-rank-factorization}(c)).

(b) Then \(A=0\) (Theorem~\ref{thm:r4-four-bases}), and \(AGA=0=A\) for every \(G\).

Every \(A\) with \(r\geqslant1\) has a full-rank factorization (Theorem~\ref{thm:r4-full-rank-factorization}(a)), so (a)
and (b) give a g-inverse of every matrix.
\end{proof}

\begin{theorem}[\citeat{bapat}{\S4.1, Result~5.18}]\label{r10-thm-ginv-normal-form}
Let \(A\in\R^{m\times n}\) have rank \(r\geqslant1\), and let \(P\in\R^{m\times m}\) and \(Q\in\R^{n\times n}\) be invertible with
\(PAQ=N_r=\left(\begin{smallmatrix}I_r&0\\0&0\end{smallmatrix}\right)\) (Theorem~\ref{thm:r4-normal-form}). Then \(G\in\R^{n\times m}\)
is a g-inverse of \(A\) if and only if
\[
G=Q\begin{pmatrix}I_r&U\\V&W\end{pmatrix}P
\]
for some \(U\in\R^{r\times(m-r)}\), \(V\in\R^{(n-r)\times r}\) and \(W\in\R^{(n-r)\times(m-r)}\) (blocks with no rows or no
columns are absent). In particular, unless \(r=m=n\), the matrix \(A\) has infinitely many g-inverses.
\end{theorem}
\begin{proof}
Let \(G\in\R^{n\times m}\), put \(H=Q^{-1}GP^{-1}\in\R^{n\times m}\), so that \(G=QHP\), and partition
\(H=\left(\begin{smallmatrix}H_{11}&H_{12}\\H_{21}&H_{22}\end{smallmatrix}\right)\) with \(H_{11}\in\R^{r\times r}\).
\begin{itemize}
\item By the block product (Theorem~\ref{r5b-thm-block}(c)), \(N_rH=\left(\begin{smallmatrix}H_{11}&H_{12}\\0&0\end{smallmatrix}\right)\)
and
\[
N_rHN_r=\begin{pmatrix}H_{11}&0\\0&0\end{pmatrix}
\]
\item \(A=P^{-1}N_rQ^{-1}\), so \(AGA=P^{-1}N_rQ^{-1}QHPP^{-1}N_rQ^{-1}=P^{-1}N_rHN_rQ^{-1}\). Multiplying by \(P\) on the left
and by \(Q\) on the right, and back by \(P^{-1}\) and \(Q^{-1}\), shows that \(AGA=A\) if and only if \(N_rHN_r=N_r\), that is,
\(H_{11}=I_r\). Then \(U=H_{12}\), \(V=H_{21}\) and \(W=H_{22}\) are arbitrary.
\item If \(QH_1P=QH_2P\), then \(H_1=H_2\) (multiply by \(Q^{-1}\) on the left and by \(P^{-1}\) on the right); so different
blocks \(U,V,W\) give different g-inverses. If \(r<m\), the block \(U\) is present and its \((1,1)\) entry can be any number;
if \(r<n\), likewise \(V\). So \(A\) has infinitely many g-inverses unless \(r=m=n\).\qedhere
\end{itemize}
\end{proof}

\begin{corollary}[cf.\ \citeat{bapat}{\S4.1}]\label{r10-cor-ginv-rowreduction}
Let \(A\in\R^{m\times n}\) have rank \(r\geqslant1\), and let row operations take \([A\mid I_m]\) to \([R\mid P]\), where \(R\)
is the RREF of \(A\), with pivot columns \(k_1<\dots<k_r\) (Theorem~\ref{thm:r4-four-bases}). Let \(G\in\R^{n\times m}\)
have row \(k_s\) equal to row \(s\) of \(P\) (\(s=1,\dots,r\)) and all other rows zero. Then \(G\) is a g-inverse of \(A\).
\end{corollary}
\begin{proof}
\begin{itemize}
\item Let \(B=[a_{k_1}\ \cdots\ a_{k_r}]\), and let \(C\) consist of the first \(r\) rows of \(R\). Then \(A=BC\) is a
full-rank factorization (Theorem~\ref{thm:r4-full-rank-factorization}(a)).
\item \(PA=R\) (Theorem~\ref{thm:r4-four-bases}), so \(Pa_{k_s}=\col_{k_s}(R)=e_s\in\R^m\) (Definition~\ref{r2-def-rref}).
Hence \(PB=[e_1\ \cdots\ e_r]\) (Theorem~\ref{r1-thm-views}(a)), and the first \(r\) rows of \(P\) form a matrix
\(L\in\R^{r\times m}\) with \(LB=I_r\) (Theorem~\ref{r1-thm-views}(c)).
\item Let \(K=[e_{k_1}\ \cdots\ e_{k_r}]\in\R^{n\times r}\). Column \(s\) of \(CK\) is \(\col_{k_s}(C)\), the \(s\)th column of
\(I_r\) (Definition~\ref{r2-def-rref}); so \(CK=I_r\).
\item By the outer product expansion (Theorem~\ref{r1-thm-views}(d)), \(KL=e_{k_1}\row_1(L)+\dots+e_{k_r}\row_r(L)\) has
row \(k_s\) equal to \(\row_s(L)=\row_s(P)\) and all other rows zero: \(KL=G\). By
Theorem~\ref{r10-thm-ginv-exists}(a), \(G\) is a g-inverse of \(A\).\qedhere
\end{itemize}
\end{proof}

\needspace{12\baselineskip}
\begin{theorem}[\citeat{harville}{Theorem~4.4.10}, \citeat{bapat}{\S4.1}, \citeat{searle-khuri}{\S8.2.2}]\label{r10-thm-ginv-submatrix}
Let \(A\in\R^{m\times n}\) have rank \(r\geqslant1\). For rows \(i_1<\dots<i_r\) and columns \(j_1<\dots<j_r\) of \(A\), let
\(A_*\in\R^{r\times r}\) be the submatrix of \(A\) in these rows and columns: its \((s,t)\) entry is \(a_{i_sj_t}\).
\begin{enumerate}
\item \(A\) has \(r\) linearly independent rows and \(r\) linearly independent columns. If the rows \(i_1,\dots,i_r\) and
the columns \(j_1,\dots,j_r\) of \(A\) are linearly independent, then \(A_*\) is invertible.
\item Let \(A_*\) be invertible, and let \(G\in\R^{n\times m}\) have \((j_s,i_t)\) entry equal to the \((s,t)\) entry of
\(A_*^{-1}\) (\(s,t=1,\dots,r\)) and all other entries \(0\). Then \(G\) is a g-inverse of \(A\). In particular, if
\(A=\left(\begin{smallmatrix}A_{11}&A_{12}\\A_{21}&A_{22}\end{smallmatrix}\right)\) with \(A_{11}\in\R^{r\times r}\) invertible,
then \(G=\left(\begin{smallmatrix}A_{11}^{-1}&0\\0&0\end{smallmatrix}\right)\in\R^{n\times m}\) is a g-inverse of \(A\).
\end{enumerate}
\end{theorem}
\begin{proof}
(a)
\begin{itemize}
\item The pivot columns of \(A\) are \(r\) linearly independent columns (Theorem~\ref{thm:r4-four-bases}(a)). As
\(\rank A'=r\) (Theorem~\ref{thm:r4-rank-rules}(a)), the pivot columns of \(A'\) are \(r\) linearly independent columns of
\(A'\), that is, transposed rows of \(A\).
\item Let the rows \(i_1,\dots,i_r\) and the columns \(j_1,\dots,j_r\) be linearly independent. The vectors
\(\row_{i_1}(A)',\dots,\row_{i_r}(A)'\) are \(r\) linearly independent vectors of \(\mathcal C(A')\), of dimension \(r\); so
they form a basis of it (Theorem~\ref{thm:r4-counting}(d)), and
\(\row_t(A)=f_{t1}\row_{i_1}(A)+\dots+f_{tr}\row_{i_r}(A)\) (\(t=1,\dots,m\)) for some numbers \(f_{tk}\).
\item Let \(U\in\R^{m\times r}\) consist of the columns \(j_1,\dots,j_r\) of \(A\), and \(F=(f_{tk})\in\R^{m\times r}\). Entry
\(j_s\) of the last equation reads \(u_{ts}=a_{tj_s}=\sum_{k=1}^rf_{tk}a_{i_kj_s}\); that is, \(U=FA_*\).
\item The columns of \(U\) are linearly independent, so
\(r=\rank U=\rank(FA_*)\leqslant\rank A_*\leqslant r\) (Theorems~\ref{thm:r4-rank-rules}(a),(b)
and~\ref{thm:r4-rank-product}(a)), and \(A_*\) is invertible (Theorem~\ref{thm:r4-rank-rules}(d)).
\end{itemize}
(b) First the leading case \(A_*=A_{11}\) (\(i_s=j_s=s\)); blocks with no rows or no columns are absent when \(r=m\) or
\(r=n\).
\begin{itemize}
\item The columns of \(\left(\begin{smallmatrix}A_{11}\\A_{21}\end{smallmatrix}\right)\) are linearly independent:
\(\left(\begin{smallmatrix}A_{11}\\A_{21}\end{smallmatrix}\right)z=0\) gives \(A_{11}z=0\), so \(z=0\). They are
\(r=\dim\mathcal C(A)\) vectors of \(\mathcal C(A)\), hence a basis of it (Theorem~\ref{thm:r4-counting}(d)).
\item So each of the other columns is a combination of them:
\(\left(\begin{smallmatrix}A_{12}\\A_{22}\end{smallmatrix}\right)=\left(\begin{smallmatrix}A_{11}\\A_{21}\end{smallmatrix}\right)X\)
for some \(X\in\R^{r\times(n-r)}\) (Theorem~\ref{r1-thm-views}(a),(b)), that is, \(A_{12}=A_{11}X\) and \(A_{22}=A_{21}X\)
(Theorem~\ref{r5b-thm-block}(c)). Then \(X=A_{11}^{-1}A_{12}\) and \(A_{22}=A_{21}A_{11}^{-1}A_{12}\).
\item With \(F=\left(\begin{smallmatrix}A_{11}^{-1}&0\\0&0\end{smallmatrix}\right)\in\R^{n\times m}\), the block product
(Theorem~\ref{r5b-thm-block}(c)) gives
\[
AF=\begin{pmatrix}I_r&0\\A_{21}A_{11}^{-1}&0\end{pmatrix},\qquad
AFA=\begin{pmatrix}A_{11}&A_{12}\\A_{21}&A_{21}A_{11}^{-1}A_{12}\end{pmatrix}=A
\]
\end{itemize}
Now the general case.
\begin{itemize}
\item Extend \(i_1,\dots,i_r\) by the other row indices, in increasing order, to a permutation \(\varphi\) of
\(1,\dots,m\), and \(j_1,\dots,j_r\) likewise to a permutation \(\psi\) of \(1,\dots,n\); let \(P=P_\varphi\) and
\(Q=P_\psi'\) (Definition~\ref{r5d-def-permutation-matrix}). Then \(\row_s(P)=e_{i_s}'\) and \(\col_t(Q)=e_{j_t}\) for
\(s,t\leqslant r\), so the \((s,t)\) entry of \(PAQ\) is \(e_{i_s}'Ae_{j_t}=a_{i_sj_t}\) (Theorem~\ref{r1-thm-views}(c),(d)):
the leading \(r\times r\) block of \(PAQ\) is \(A_*\).
\item \(\rank(PAQ)=r\) (Theorem~\ref{thm:r4-rank-product}(d)), so by the leading case
\(F=\left(\begin{smallmatrix}A_*^{-1}&0\\0&0\end{smallmatrix}\right)\in\R^{n\times m}\) is a g-inverse of \(PAQ\).
\item \(P'P=PP'=I_m\) and \(Q'Q=QQ'=I_n\) (Proposition~\ref{r5d-prop-permutation-matrix}(c)), so \(A=P'(PAQ)Q'\) and
\[
A(QFP)A=P'(PAQ)Q'QFPP'(PAQ)Q'=P'(PAQ)F(PAQ)Q'=P'(PAQ)Q'=A
\]
\item Finally \(QFP=G\). Let \(f_{st}\) be the entries of \(F\); they are \(0\) unless \(s,t\leqslant r\). Column \(t\) of
\(QF\) is \(Q\col_t(F)=\sum_{s\leqslant r}f_{st}\col_s(Q)=\sum_{s\leqslant r}f_{st}e_{j_s}\) for \(t\leqslant r\), and \(0\) for
\(t>r\) (Theorem~\ref{r1-thm-views}(a),(b)). By the outer product expansion (Theorem~\ref{r1-thm-views}(d)),
\(QFP=\sum_{t\leqslant r}\col_t(QF)\row_t(P)=\sum_{s,t\leqslant r}f_{st}\,e_{j_s}e_{i_t}'\), and \(e_{j_s}e_{i_t}'\in\R^{n\times m}\)
has entry \(1\) in position \((j_s,i_t)\) and \(0\) elsewhere.\qedhere
\end{itemize}
\end{proof}

\begin{example}[cf.\ \citeat{harville}{\S9.5}]\label{r10-exm-ginverse}
Let \(A\) be the matrix of Example~\ref{exm:r4-four-bases}, with the full-rank factorization \(A=BC\) of
Example~\ref{exm:r4-factorization}:
\[
A=\begin{pmatrix}1&2&1&3\\1&3&0&5\\2&5&1&8\end{pmatrix},\qquad
B=\begin{pmatrix}1&2\\1&3\\2&5\end{pmatrix},\qquad
C=\begin{pmatrix}1&0&3&-1\\0&1&-1&2\end{pmatrix}
\]
Find two g-inverses of \(A\), and decide whether \(Ax=y\) is consistent for \(y=(1,2,3)'\) and for \(y=(1,1,1)'\).

\textit{Solution.}
\begin{itemize}
\item Columns \(1\), \(2\) of \(C\) are the unit vectors of \(\R^2\). Rows \(1\), \(2\) of \(B\) form
\(B_1=\left(\begin{smallmatrix}1&2\\1&3\end{smallmatrix}\right)\), and
\(B_1\left(\begin{smallmatrix}3&-2\\-1&1\end{smallmatrix}\right)=I_2\), so
\(B_1^{-1}=\left(\begin{smallmatrix}3&-2\\-1&1\end{smallmatrix}\right)\) (Theorem~\ref{r2-thm-invertible}); rows \(1\), \(3\)
form \(B_2=\left(\begin{smallmatrix}1&2\\2&5\end{smallmatrix}\right)\), and likewise
\(B_2^{-1}=\left(\begin{smallmatrix}5&-2\\-2&1\end{smallmatrix}\right)\), as
\(B_2\left(\begin{smallmatrix}5&-2\\-2&1\end{smallmatrix}\right)=I_2\). Put
\[
K=\begin{pmatrix}1&0\\0&1\\0&0\\0&0\end{pmatrix},\qquad
L_1=\begin{pmatrix}3&-2&0\\-1&1&0\end{pmatrix},\qquad
L_2=\begin{pmatrix}5&0&-2\\-2&0&1\end{pmatrix}
\]
where \(L_1\) and \(L_2\) carry the columns of \(B_1^{-1}\) and \(B_2^{-1}\) in columns \(1,2\) and \(1,3\). Then \(CK=I_2\)
(Theorem~\ref{r1-thm-views}(a),(b)), and by Theorem~\ref{r1-thm-views}(c) the rows of \(L_1B\) are
\(3(1,2)-2(1,3)=(1,0)\) and \(-(1,2)+(1,3)=(0,1)\), and those of \(L_2B\) are \(5(1,2)-2(2,5)=(1,0)\) and
\(-2(1,2)+(2,5)=(0,1)\): \(L_1B=L_2B=I_2\).
\item By Theorem~\ref{r10-thm-ginv-exists}(a), both
\[
G_1=KL_1=\begin{pmatrix}3&-2&0\\-1&1&0\\0&0&0\\0&0&0\end{pmatrix},\qquad
G_2=KL_2=\begin{pmatrix}5&0&-2\\-2&0&1\\0&0&0\\0&0&0\end{pmatrix}
\]
are g-inverses of \(A\): a g-inverse need not be unique.
\item The same two matrices come from Theorem~\ref{r10-thm-ginv-submatrix}(b): the submatrices of \(A\) in rows
\(1,2\) and in rows \(1,3\), with columns \(1,2\), are \(B_1\) and \(B_2\). Also, \(L_1\) consists of the first
two rows of the matrix \(P\) of Example~\ref{exm:r4-four-bases}, so \(G_1\) is the g-inverse of
Corollary~\ref{r10-cor-ginv-rowreduction}.
\item Since \(CK=I_2\),
\[
AG_1=BCKL_1=BL_1=\begin{pmatrix}1&2\\1&3\\2&5\end{pmatrix}\begin{pmatrix}3&-2&0\\-1&1&0\end{pmatrix}
=\begin{pmatrix}1&0&0\\0&1&0\\1&1&0\end{pmatrix}
\]
so \(AG_1y=(y_1,\ y_2,\ y_1+y_2)'\). By Theorem~\ref{r10-thm-ginv-solves}(b), \(Ax=y\) is consistent if and only if
\(y_3=y_1+y_2\); this matches \(\row_3(A)=\row_1(A)+\row_2(A)\) (Example~\ref{exm:r4-four-bases}).
\item \(y=(1,2,3)'\): consistent, and \(G_1y=(3-4,\ -1+2,\ 0,\ 0)'=(-1,1,0,0)'\) is a solution; indeed
\(A(-1,1,0,0)'=-a_1+a_2=(1,2,3)'\).
\item \(y=(1,1,1)'\): \(AG_1y=(1,1,2)'\ne y\), so \(Ax=y\) is inconsistent.
\end{itemize}
\end{example}

\needspace{14\baselineskip}
\subsection*{Solving linear systems with a generalized inverse}

\begin{theorem}[\citeat{bapat}{Result~4.4}, \citeat{harville}{Theorem~11.2.4}]\label{r10-thm-general-solution}
Let \(A\in\R^{m\times n}\), let \(G\) be a g-inverse of \(A\), and let \(y\in\mathcal C(A)\). The solutions of \(Ax=y\)
are exactly the vectors
\[
x=Gy+(I_n-GA)z\qquad(z\in\R^n)
\]
In particular (\(y=0\)), \(\mathcal N(A)=\mathcal C(I_n-GA)\).
\end{theorem}
\begin{proof}
\begin{itemize}
\item Write \(y=Av\). For every \(z\), \(A\bigl(Gy+(I-GA)z\bigr)=AGy+(A-AGA)z=AGAv+0=Av=y\).
\item Conversely, let \(Au=y\), and put \(z=u-Gy\). Since \(GAu=Gy\) and \(GAGy=G(AGy)=Gy\)
(Theorem~\ref{r10-thm-ginv-solves}(a)),
\[
Gy+(I-GA)(u-Gy)=Gy+u-Gy-GAu+GAGy=u-Gy+Gy=u
\]
\item For \(y=0\), the solutions of \(Ax=0\) are the vectors \((I-GA)z\), \(z\in\R^n\); that is,
\(\mathcal N(A)=\mathcal C(I-GA)\).\qedhere
\end{itemize}
\end{proof}

\begin{theorem}[\citeat{harville}{Theorem~11.10.1}]\label{r10-thm-invariant-combination}
Let \(A\in\R^{m\times n}\) and \(y\in\mathcal C(A)\), let \(G\) be a g-inverse of \(A\), and let \(\lambda\in\R^n\). The
following are equivalent:
\begin{enumerate}[label=(\roman*)]
\item \(\lambda'x\) has the same value for every solution \(x\) of \(Ax=y\);
\item \(\lambda\in\mathcal C(A')\), that is, \(\lambda'=a'A\) for some \(a\in\R^m\);
\item \(\lambda'GA=\lambda'\).
\end{enumerate}
If they hold, then \(\lambda'x=a'y\) for every solution \(x\).
\end{theorem}
\begin{proof}
\begin{itemize}
\item (ii)\,\(\Rightarrow\)\,(i): if \(Ax=y\), then \(\lambda'x=a'Ax=a'y\), whatever the solution \(x\).
\item (i)\,\(\Rightarrow\)\,(iii): by Theorem~\ref{r10-thm-general-solution}, \(\lambda'\bigl(Gy+(I-GA)z\bigr)=\lambda'Gy+\lambda'(I-GA)z\)
is the same for every \(z\in\R^n\). Comparing \(z=e_j\) with \(z=0\), entry \(j\) of the row vector \(\lambda'(I-GA)\) is
\(0\), for \(j=1,\dots,n\). So \(\lambda'(I-GA)=0\), that is, \(\lambda'=\lambda'GA\).
\item (iii)\,\(\Rightarrow\)\,(ii): \(\lambda'=(\lambda'G)A=a'A\) with \(a=G'\lambda\)
(Theorem~\ref{r1-thm-transpose}).\qedhere
\end{itemize}
\end{proof}

\begin{example}\label{r10-exm-general-solution}
For \(A\), \(G_1\) and \(y=(1,2,3)'\) of Example~\ref{r10-exm-ginverse}, find all solutions of \(Ax=y\), and decide
whether \(x_1\) and \(x_1+3x_3-x_4\) have the same value for all of them.

\textit{Solution.} \(G_1A=KL_1BC=KC\), as \(L_1B=I_2\); so
\[
G_1A=\begin{pmatrix}1&0&3&-1\\0&1&-1&2\\0&0&0&0\\0&0&0&0\end{pmatrix},\qquad
I_4-G_1A=\begin{pmatrix}0&0&-3&1\\0&0&1&-2\\0&0&1&0\\0&0&0&1\end{pmatrix}
\]
By Theorem~\ref{r10-thm-general-solution}, with \(G_1y=(-1,1,0,0)'\), the solutions are
\[
x=\begin{pmatrix}-1\\1\\0\\0\end{pmatrix}+z_3\begin{pmatrix}-3\\1\\1\\0\end{pmatrix}+z_4\begin{pmatrix}1\\-2\\0\\1\end{pmatrix}
=\begin{pmatrix}-1-3z_3+z_4\\1+z_3-2z_4\\z_3\\z_4\end{pmatrix}\qquad(z_3,z_4\in\R)
\]
\(z_1\) and \(z_2\) drop out, as columns \(1\), \(2\) of \(I_4-G_1A\) are \(0\). Columns \(3\), \(4\) are the special
solutions \(s_1\), \(s_2\) of Example~\ref{exm:r4-four-bases}, and \(\mathcal N(A)=\mathcal C(I_4-G_1A)=\Span\{s_1,s_2\}\).
\begin{itemize}
\item \(\lambda=(1,0,0,0)'\): \(\lambda'G_1A=\row_1(G_1A)=(1,0,3,-1)\ne\lambda'\), so \(x_1=-1-3z_3+z_4\) is not the same
for all solutions (Theorem~\ref{r10-thm-invariant-combination}).
\item \(\lambda=(1,0,3,-1)'\): \(\lambda'G_1A=\row_1(G_1A)+3\row_3(G_1A)-\row_4(G_1A)=(1,0,3,-1)=\lambda'\), so
\(x_1+3x_3-x_4\) is the same for all solutions: \((-1-3z_3+z_4)+3z_3-z_4=-1\). Here \(\lambda'\) is the first row of
the RREF of \(A\), a vector of \(\mathcal R(A)\).
\end{itemize}
\end{example}

\subsection*{Properties of generalized inverses}

\begin{proposition}[\citeat{harville}{Lemma~9.3.3, Corollary~9.3.4, Lemma~9.2.8}, \citeat{searle-khuri}{Theorem~8.2(iv)}]\label{r10-prop-ginv-basic}
Let \(A\in\R^{m\times n}\), and let \(G\) be a g-inverse of \(A\).
\begin{enumerate}
\item \(G'\) is a g-inverse of \(A'\).
\item If \(A\) is symmetric, then \(G'\) and \(\frac12(G+G')\) are g-inverses of \(A\), and \(\frac12(G+G')\) is symmetric.
So every symmetric matrix has a symmetric g-inverse.
\item If \(A\) has full column rank, its g-inverses are exactly its left inverses, the \(L\in\R^{n\times m}\) with
\(LA=I_n\). If \(A\) has full row rank, its g-inverses are exactly its right inverses, the \(K\in\R^{n\times m}\) with
\(AK=I_m\).
\end{enumerate}
\end{proposition}
\begin{proof}
(a) \(A'G'A'=(AGA)'=A'\) (Theorem~\ref{r1-thm-transpose}).

(b) If \(A'=A\), then (a) reads \(AG'A=A\). So \(A\cdot\frac12(G+G')\cdot A=\frac12(AGA+AG'A)=\frac12(A+A)=A\), and
\(\bigl(\frac12(G+G')\bigr)'=\frac12(G'+G)\).

(c)
\begin{itemize}
\item Let \(\rank A=n\). Then \(A=AI_n\) is a full-rank factorization, so \(A\) has a left inverse \(L_0\)
(Theorem~\ref{thm:r4-full-rank-factorization}(c)).
\item Every left inverse \(L\) is a g-inverse: \(ALA=AI_n=A\). Every g-inverse \(G\) is a left inverse:
\(GA=(L_0A)GA=L_0(AGA)=L_0A=I_n\).
\item Let \(\rank A=m\). Then \(A=I_mA\) is a full-rank factorization, so \(A\) has a right inverse \(K_0\). Every right
inverse \(K\) gives \(AKA=I_mA=A\), and every g-inverse \(G\) gives \(AG=AG(AK_0)=(AGA)K_0=AK_0=I_m\).\qedhere
\end{itemize}
\end{proof}

\begin{lemma}[\citeat{harville}{Lemma~9.3.5}]\label{r10-lem-range-test}
Let \(A\in\R^{m\times n}\), and let \(G\) be a g-inverse of \(A\).
\begin{enumerate}
\item For \(C\in\R^{m\times q}\): \(\mathcal C(C)\subseteq\mathcal C(A)\) if and only if \(AGC=C\).
\item For \(B\in\R^{p\times n}\): \(\mathcal R(B)\subseteq\mathcal R(A)\) if and only if \(BGA=B\).
\end{enumerate}
The inclusions do not involve \(G\); so if \(AGC=C\) (or \(BGA=B\)) holds for one g-inverse \(G\), it holds for every
g-inverse. In particular, the system \(AX=C\) (unknown \(X\in\R^{n\times q}\)) is consistent if and only if \(AGC=C\),
and then \(X=GC\) is a solution.
\end{lemma}
\begin{proof}
(a)
\begin{itemize}
\item If \(AGC=C\), then each column \(c_j=A(Gc_j)\) of \(C\) lies in \(\mathcal C(A)\) (Theorem~\ref{r1-thm-views}(a)), and
so does every combination of them: \(\mathcal C(C)\subseteq\mathcal C(A)\).
\item If \(\mathcal C(C)\subseteq\mathcal C(A)\), then \(c_j=Af_j\) for some \(f_j\in\R^n\) (\(j=1,\dots,q\)), that is, \(C=AF\)
with \(F=[f_1\ \cdots\ f_q]\) (Theorem~\ref{r1-thm-views}(a)); so \(AGC=AGAF=AF=C\).
\end{itemize}
(b)
\begin{itemize}
\item If \(BGA=B\), then each row \(\row_i(B)=(\row_i(B)G)A\) lies in \(\mathcal R(A)\) (Theorem~\ref{r1-thm-views}(c)), and
so does every combination of the rows: \(\mathcal R(B)\subseteq\mathcal R(A)\).
\item If \(\mathcal R(B)\subseteq\mathcal R(A)\), then \(\row_i(B)=l_i'A\) for some \(l_i\in\R^m\) (\(i=1,\dots,p\)), that is,
\(B=LA\), where \(L\) has the rows \(l_1',\dots,l_p'\) (Theorem~\ref{r1-thm-views}(c)); so \(BGA=LAGA=LA=B\).
\end{itemize}
Finally, \(AX=C\) is consistent if and only if every system \(Ax_j=c_j\) is (Theorem~\ref{r1-thm-views}(a)), that is,
\(c_j\in\mathcal C(A)\) for every \(j\) (Corollary~\ref{cor:r4-ax-b}(a)), that is, \(\mathcal C(C)\subseteq\mathcal C(A)\); and
then \(A(GC)=C\) by (a).
\end{proof}

\needspace{10\baselineskip}
\begin{theorem}[\citeat{harville}{Lemmas~10.2.5, 9.3.7, 10.2.6}, \citeat{bapat}{Result~4.2}]\label{r10-thm-rank-idempotent}
Let \(A\in\R^{m\times n}\) have rank \(r\), and let \(G\) be a g-inverse of \(A\).
\begin{enumerate}
\item \(AG\in\R^{m\times m}\) and \(GA\in\R^{n\times n}\) are idempotent.
\item \(\rank(AG)=\rank(GA)=r\); more precisely, \(\mathcal C(AG)=\mathcal C(A)\), \(\mathcal R(GA)=\mathcal R(A)\) and
\(\mathcal N(GA)=\mathcal N(A)\).
\item \(\tr(AG)=\tr(GA)=r\), \(\rank(I_n-GA)=n-r\) and \(\rank(I_m-AG)=m-r\).
\end{enumerate}
\end{theorem}
\begin{proof}
(a) \((AG)(AG)=(AGA)G=AG\) and \((GA)(GA)=G(AGA)=GA\).

(b)
\begin{itemize}
\item By Theorem~\ref{thm:r4-rank-product}(a), applied to the products \((AG)A\), \(AG\) and \(A(GA)\), \(GA\),
\[
r=\rank(AGA)\leqslant\rank(AG)\leqslant\rank A=r,\qquad r=\rank(AGA)\leqslant\rank(GA)\leqslant\rank A=r
\]
\item \(\mathcal C(AG)\subseteq\mathcal C(A)\) and \(\mathcal C(A)=\mathcal C((AG)A)\subseteq\mathcal C(AG)\)
(Theorem~\ref{thm:r4-rank-product}(a)); so \(\mathcal C(AG)=\mathcal C(A)\).
\item Each row \(\row_i(GA)=\row_i(G)A\) lies in \(\mathcal R(A)\), and each row \(\row_i(A)=\row_i(A)(GA)\) (as \(A=AGA\)) lies
in \(\mathcal R(GA)\) (Theorem~\ref{r1-thm-views}(c)); so \(\mathcal R(GA)=\mathcal R(A)\).
\item If \(GAx=0\), then \(Ax=AGAx=0\); and \(Ax=0\) gives \(GAx=0\). So \(\mathcal N(GA)=\mathcal N(A)\).
\end{itemize}
(c) By (a) and (b), \(\tr(AG)=\rank(AG)=r\) and \(\tr(GA)=\rank(GA)=r\) (Theorem~\ref{r6-thm-idempotent}(c)). Also
\(I_n-GA\) and \(I_m-AG\) are idempotent (Theorem~\ref{r6-thm-idempotent}(b)), so
\(\rank(I_n-GA)=\tr(I_n-GA)=n-\tr(GA)=n-r\) (Theorem~\ref{r1-thm-trace}(a)), and likewise \(\rank(I_m-AG)=m-r\).
\end{proof}

\begin{theorem}[\citeat{bapat}{\S4.1, Result~4.3}, \citeat{harville}{Lemma~20.3.2, Theorem~20.3.1}]\label{r10-thm-reflexive}
Let \(A\in\R^{m\times n}\). A g-inverse \(G\) of \(A\) with \(GAG=G\) is a \textbf{reflexive g-inverse} of \(A\); then \(A\)
is also a g-inverse of \(G\).
\begin{enumerate}
\item If \(G_1\) and \(G_2\) are g-inverses of \(A\), then \(G_1AG_2\) is a reflexive g-inverse of \(A\). In particular
\(GAG\) is one, for every g-inverse \(G\).
\item Every g-inverse \(G\) of \(A\) has \(\rank G\geqslant\rank A\), and \(\rank G=\rank A\) if and only if \(G\) is
reflexive.
\item Let \(r=\rank A\geqslant1\) and let \(A=BC\) be a full-rank factorization, with \(B\in\R^{m\times r}\) and \(C\in\R^{r\times n}\). Then \(G\) is a
reflexive g-inverse of \(A\) if and only if \(G=KL\) for some \(L\in\R^{r\times m}\) with \(LB=I_r\) and some
\(K\in\R^{n\times r}\) with \(CK=I_r\).
\end{enumerate}
\end{theorem}
\begin{proof}
(a) \(A(G_1AG_2)A=(AG_1A)G_2A=AG_2A=A\), and
\[
(G_1AG_2)A(G_1AG_2)=G_1(AG_2A)G_1AG_2=G_1(AG_1A)G_2=G_1AG_2
\]

(b)
\begin{itemize}
\item \(\rank A=\rank(AGA)\leqslant\rank G\) (Theorem~\ref{thm:r4-rank-product}(a), twice). If \(G\) is reflexive, then
also \(\rank G=\rank(GAG)\leqslant\rank A\); so \(\rank G=\rank A\).
\item Conversely, let \(\rank G=\rank A\). Then \(\mathcal C(GA)\subseteq\mathcal C(G)\) (Theorem~\ref{thm:r4-rank-product}(a)),
and \(\rank(GA)=\rank A=\rank G\) (Theorem~\ref{r10-thm-rank-idempotent}(b)); so \(\mathcal C(GA)=\mathcal C(G)\)
(Theorem~\ref{thm:r4-counting}(a)).
\item So each column of \(G\) is \(GAx_j\) for some \(x_j\in\R^n\), that is, \(G=GAX\) with \(X=[x_1\ \cdots\ x_m]\)
(Theorem~\ref{r1-thm-views}(a)), and \(GAG=GAGAX=GAX=G\).
\end{itemize}
(c)
\begin{itemize}
\item If \(G=KL\), then \(G\) is a g-inverse of \(A\) (Theorem~\ref{r10-thm-ginv-exists}(a)), and
\(GAG=KLBCKL=K(LB)(CK)L=KL=G\).
\item Conversely, let \(G\) be a reflexive g-inverse. By Theorem~\ref{thm:r4-full-rank-factorization}(c), \(B\) has a left
inverse \(L_0\) and \(C\) a right inverse \(K_0\). Multiplying \(BCGBC=AGA=A=BC\) by \(L_0\) on the left and by \(K_0\) on the
right gives \(CGB=I_r\).
\item So \(K=GB\) has \(CK=CGB=I_r\), \(L=CG\) has \(LB=CGB=I_r\), and \(G=GAG=(GB)(CG)=KL\).\qedhere
\end{itemize}
\end{proof}

\needspace{12\baselineskip}
\begin{example}\label{r10-exm-reflexive}
Let \(A\) and \(G_1\) be as in Examples~\ref{r10-exm-ginverse} and~\ref{r10-exm-general-solution}. Is \(G_1\) reflexive? For
the matrix \(U\in\R^{4\times3}\) with \((3,3)\) entry \(1\) and all other entries \(0\), show that
\(G_3=G_1+(I_4-G_1A)U\) is a g-inverse of \(A\) that is not reflexive, and find the reflexive g-inverse \(G_3AG_3\).

\textit{Solution.}
\begin{itemize}
\item The matrices \(G_1A\) and
\[
AG_1=\begin{pmatrix}1&0&0\\0&1&0\\1&1&0\end{pmatrix}
\]
are idempotent with trace \(1+1=2=\rank A\), as Theorem~\ref{r10-thm-rank-idempotent}(a),(c) says; \(AG_1\) is not symmetric.
\item \(G_1\) has the two linearly independent nonzero rows \((3,-2,0)\) and \((-1,1,0)\), so \(\rank G_1=2=\rank A\), and
\(G_1\) is reflexive (Theorem~\ref{r10-thm-reflexive}(b); also \(G_1=KL_1\), Theorem~\ref{r10-thm-reflexive}(c)).
\item \((I_4-G_1A)U\) has column \(3\) equal to column \(3\) of \(I_4-G_1A\), which is \((-3,1,1,0)'\)
(Example~\ref{r10-exm-general-solution}), and its other columns \(0\) (Theorem~\ref{r1-thm-views}(a)). So
\[
G_3=\begin{pmatrix}3&-2&-3\\-1&1&1\\0&0&1\\0&0&0\end{pmatrix}
\]
\item \(A(I_4-G_1A)=A-AG_1A=0\), so \(AG_3=AG_1+A(I_4-G_1A)U=AG_1\), and \(AG_3A=AG_1A=A\): \(G_3\) is a g-inverse of \(A\).
\item The first three rows of \(G_3\) form a matrix \(M\) with \(\det M=1\cdot(3-2)=1\) (expansion along its last row,
Theorem~\ref{r5d-thm-expansion}), so
\(M\) is invertible (Theorem~\ref{r5d-thm-product}(a)). If \(G_3x=0\), then \(Mx=0\) and \(x=0\); so the three columns of
\(G_3\) are linearly independent, and \(\rank G_3=3\ne\rank A\) (Theorem~\ref{thm:r4-rank-rules}(b)): \(G_3\) is not
reflexive (Theorem~\ref{r10-thm-reflexive}(b)).
\item \(G_3AG_3=G_3(AG_1)\). By Theorem~\ref{r1-thm-views}(a),(b), its columns are \(G_3(1,0,1)'=(3-3,\ -1+1,\ 1,\ 0)'\),
\(G_3(0,1,1)'=(-2-3,\ 1+1,\ 1,\ 0)'\) and \(0\):
\[
G_3AG_3=\begin{pmatrix}0&-5&0\\0&2&0\\1&1&0\\0&0&0\end{pmatrix}
\]
a reflexive g-inverse of \(A\) (Theorem~\ref{r10-thm-reflexive}(a)), of rank \(2\).
\end{itemize}
\end{example}

\begin{theorem}[\citeat{harville}{Theorem~9.4.1}, \citeat{bapat}{Result~6.1}]\label{r10-thm-invariance}
Let \(A\in\R^{m\times n}\), \(B\in\R^{p\times n}\) and \(C\in\R^{m\times q}\).
\begin{enumerate}
\item If \(\mathcal R(B)\subseteq\mathcal R(A)\) and \(\mathcal C(C)\subseteq\mathcal C(A)\), then \(BGC\) is the same matrix for
every g-inverse \(G\) of \(A\).
\item Conversely, let \(BGC\) be the same matrix for every g-inverse \(G\) of \(A\). If \(C\ne0\), then
\(\mathcal R(B)\subseteq\mathcal R(A)\); if \(B\ne0\), then \(\mathcal C(C)\subseteq\mathcal C(A)\).
\end{enumerate}
\end{theorem}
\begin{proof}
(a) Each row of \(B\) lies in \(\mathcal R(A)\) and each column of \(C\) in \(\mathcal C(A)\), so \(B=LA\) and \(C=AF\) for some
\(L\in\R^{p\times m}\) and \(F\in\R^{n\times q}\) (Theorem~\ref{r1-thm-views}(a),(c)). Hence \(BGC=LAGAF=LAF\) for every
g-inverse \(G\).

(b)
\begin{itemize}
\item Fix a g-inverse \(G\) of \(A\) (Theorem~\ref{r10-thm-ginv-exists}). For all \(U,V\in\R^{n\times m}\), the matrix
\(G+(I_n-GA)U+V(I_m-AG)\) is also a g-inverse of \(A\): as \(AGA=A\),
\[
A\bigl(G+(I_n-GA)U+V(I_m-AG)\bigr)A=AGA+(A-AGA)UA+AV(A-AGA)=A
\]
\item By hypothesis, \(B\bigl(G+(I_n-GA)U+V(I_m-AG)\bigr)C=BGC\). With \(V=0\) this gives \(B(I_n-GA)UC=0\) for every \(U\);
with \(U=0\), \(BV(I_m-AG)C=0\) for every \(V\).
\item Let \(C\ne0\), with a column \(c_k\ne0\), so \(c_k'c_k>0\). For \(i=1,\dots,n\), the matrix
\(U_i=\frac1{c_k'c_k}e_ic_k'\in\R^{n\times m}\) has \(U_ic_k=e_i\). Column \(k\) of \(B(I_n-GA)U_iC\) is
\(B(I_n-GA)U_ic_k=B(I_n-GA)e_i\), the \(i\)th column of \(B(I_n-GA)\) (Theorem~\ref{r1-thm-views}(a),(b)), and it is \(0\). So
\(B(I_n-GA)=0\), that is, \(BGA=B\), and \(\mathcal R(B)\subseteq\mathcal R(A)\) (Lemma~\ref{r10-lem-range-test}(b)).
\item Let \(B\ne0\), with a row \(b_k'\ne0\) (\(b_k\in\R^n\)). For \(i=1,\dots,m\), the matrix
\(V_i=\frac1{b_k'b_k}b_ke_i'\in\R^{n\times m}\) has \(b_k'V_i=e_i'\). Row \(k\) of \(BV_i(I_m-AG)C\) is
\(b_k'V_i(I_m-AG)C=e_i'(I_m-AG)C\), the \(i\)th row of \((I_m-AG)C\) (Theorem~\ref{r1-thm-views}(c)), and it is \(0\). So
\((I_m-AG)C=0\), that is, \(AGC=C\), and \(\mathcal C(C)\subseteq\mathcal C(A)\) (Lemma~\ref{r10-lem-range-test}(a)).\qedhere
\end{itemize}
\end{proof}

\needspace{14\baselineskip}
\subsection*{Generalized inverses of \(A'A\)}

\begin{theorem}[\citeat{harville}{Theorems~12.3.4, 12.3.1}, \citeat{abadir-magnus}{Exercise~10.54}]\label{r10-thm-ata}
Let \(A\in\R^{m\times n}\), let \(G\) be a g-inverse of \(A'A\), and put \(P=AGA'=A(A'A)^-A'\in\R^{m\times m}\).
\begin{enumerate}
\item \(AGA'A=A\) and \(A'AGA'=A'\). Hence \(GA'\) is a g-inverse of \(A\).
\item For every \(y\in\R^m\), \(Py\) is the orthogonal projection of \(y\) onto \(\mathcal C(A)\). Hence \(P\) does not depend on
the choice of the g-inverse \(G\), and \(P=A(A'A)^{-1}A'\), the matrix of Theorem~\ref{r6-thm-hat}, if \(\rank A=n\).
\item \(P'=P=P^2\), \(PA=A\) and \(A'P=A'\).
\item \(\rank P=\tr P=\rank A\), and \(I_m-P\) is symmetric and idempotent, of rank \(m-\rank A\).
\end{enumerate}
\end{theorem}
\begin{proof}
(a)
\begin{itemize}
\item \(A'A(GA'A-I_n)=A'AGA'A-A'A=0\), so each column of \(GA'A-I_n\) lies in \(\mathcal N(A'A)=\mathcal N(A)\)
(Theorems~\ref{r1-thm-views}(a) and~\ref{thm:r4-rank-ata}(a)). Hence \(A(GA'A-I_n)=0\), that is, \(AGA'A=A\).
\item \(A'A\) is symmetric, so \(G'\) is also a g-inverse of it (Proposition~\ref{r10-prop-ginv-basic}(b)), and the first identity
for \(G'\) reads \(AG'A'A=A\). Transposing (Theorem~\ref{r1-thm-transpose}), \(A'AGA'=A'\).
\item \(A(GA')A=AGA'A=A\).
\end{itemize}
(b)
\begin{itemize}
\item \(Py=A(GA'y)\in\mathcal C(A)\), and \(A'(y-Py)=A'y-A'AGA'y=0\) by (a). So \(y-Py\perp\mathcal C(A)\)
(Theorem~\ref{r6-thm-complement}(a)), and \(Py\) is the orthogonal projection of \(y\) onto \(\mathcal C(A)\)
(Definition~\ref{r6-def-projection}).
\item If \(G_1\), \(G_2\) are g-inverses of \(A'A\), then \(AG_1A'y=AG_2A'y\) for every \(y\in\R^m\), as the projection is unique
(Definition~\ref{r6-def-projection}). Taking for \(y\) the columns of \(I_m\), the columns of \(AG_1A'\) and \(AG_2A'\) agree
(Theorem~\ref{r1-thm-views}(b)): \(AG_1A'=AG_2A'\).
\item If \(\rank A=n\), then \(A'A\) is invertible (Theorem~\ref{r6-thm-hat}), its only g-inverse is \((A'A)^{-1}\)
(Definition~\ref{r10-def-ginverse}), and \(P=A(A'A)^{-1}A'\).
\end{itemize}
(c)
\begin{itemize}
\item \(G'\) is a g-inverse of the symmetric \(A'A\) (Proposition~\ref{r10-prop-ginv-basic}(b)), so by (b),
\(P=AG'A'=(AGA')'=P'\) (Theorem~\ref{r1-thm-transpose}).
\item \(PA=AGA'A=A\) by (a), and \(A'P=A'P'=(PA)'=A'\).
\item \(P^2=(PA)GA'=AGA'=P\).
\end{itemize}
(d)
\begin{itemize}
\item \(\mathcal C(P)=\mathcal C(A(GA'))\subseteq\mathcal C(A)\) and \(\mathcal C(A)=\mathcal C(PA)\subseteq\mathcal C(P)\)
(Theorem~\ref{thm:r4-rank-product}(a)); so \(\rank P=\rank A\), and \(\rank P=\tr P\) as \(P\) is idempotent
(Theorem~\ref{r6-thm-idempotent}(c)).
\item \(I_m-P\) is symmetric by (c) and idempotent (Theorem~\ref{r6-thm-idempotent}(b)); so
\(\rank(I_m-P)=\tr(I_m-P)=m-\tr P=m-\rank A\) (Theorems~\ref{r6-thm-idempotent}(c) and~\ref{r1-thm-trace}(a)).\qedhere
\end{itemize}
\end{proof}

\begin{corollary}[\citeat{bapat}{Result~7.1}, \citeat{harville}{Theorem~9.4.1}]\label{r10-cor-invariant-quadratic}
Let \(A\in\R^{m\times n}\), and let \(\lambda\in\R^n\), \(\lambda\ne0\). Then \(\lambda'G\lambda\) is the same number for every
g-inverse \(G\) of \(A'A\) if and only if \(\lambda\in\mathcal C(A')\). If \(\lambda'=a'A\) (\(a\in\R^m\)), then for every
g-inverse \(G\) of \(A'A\)
\[
\lambda'G\lambda=a'Pa=\lVert Pa\rVert^2\geqslant0
\]
where \(P=AGA'\) is the matrix of Theorem~\ref{r10-thm-ata}.
\end{corollary}
\begin{proof}
\begin{itemize}
\item Let \(\lambda'=a'A\). Then \(\lambda'G\lambda=a'AGA'a=a'Pa\) (Theorem~\ref{r1-thm-transpose}), and \(P\) does not depend on
\(G\) (Theorem~\ref{r10-thm-ata}(b)). As \(P\) is symmetric and idempotent (Theorem~\ref{r10-thm-ata}(c)),
\(a'Pa=a'P'Pa=\lVert Pa\rVert^2\).
\item Conversely, let \(\lambda'G\lambda\) be the same for every g-inverse \(G\) of \(A'A\). Theorem~\ref{r10-thm-invariance}(b),
for the matrix \(A'A\) with \(B=\lambda'\ne0\) and \(C=\lambda\ne0\), gives \(\mathcal R(\lambda')\subseteq\mathcal R(A'A)\):
\(\lambda'=b'A'A\) for some \(b\in\R^n\). So \(\lambda=A'(Ab)\in\mathcal C(A')\).\qedhere
\end{itemize}
\end{proof}

\needspace{12\baselineskip}
\begin{example}\label{r10-exm-ata}
Let \(A\) be the matrix of Example~\ref{r10-exm-ginverse}. Find the g-inverses \(F_1\) and \(F_2\) of \(A'A\) built from its
submatrices in rows and columns \(1,2\) and in rows and columns \(1,3\), and the matrices \(AF_1A'\) and \(AF_2A'\). Compare
\(\lambda'F_1\lambda\) and \(\lambda'F_2\lambda\) for \(\lambda=(1,0,3,-1)'\) and for \(\lambda=(1,0,0,0)'\).

\textit{Solution.}
\begin{itemize}
\item With the columns \(a_1,\dots,a_4\) of \(A\), the \((i,j)\) entry of \(A'A\) is \(a_i'a_j\) (Theorem~\ref{r1-thm-views}(d)):
\[
A=\begin{pmatrix}1&2&1&3\\1&3&0&5\\2&5&1&8\end{pmatrix},\qquad
A'A=\begin{pmatrix}6&15&3&24\\15&38&7&61\\3&7&2&11\\24&61&11&98\end{pmatrix}
\]
and \(\rank(A'A)=\rank A=2\) (Theorem~\ref{thm:r4-rank-ata}(b), Example~\ref{exm:r4-four-bases}).
\item The submatrix in rows and columns \(1,2\) is \(\left(\begin{smallmatrix}6&15\\15&38\end{smallmatrix}\right)\), with
determinant \(228-225=3\) and inverse \(\frac13\left(\begin{smallmatrix}38&-15\\-15&6\end{smallmatrix}\right)\); the one in rows and
columns \(1,3\) is \(\left(\begin{smallmatrix}6&3\\3&2\end{smallmatrix}\right)\), with determinant \(12-9=3\) and inverse
\(\frac13\left(\begin{smallmatrix}2&-3\\-3&6\end{smallmatrix}\right)\) (Theorem~\ref{r5d-thm-adjugate}). By
Theorem~\ref{r10-thm-ginv-submatrix}(b),
\[
F_1=\frac13\begin{pmatrix}38&-15&0&0\\-15&6&0&0\\0&0&0&0\\0&0&0&0\end{pmatrix},\qquad
F_2=\frac13\begin{pmatrix}2&0&-3&0\\0&0&0&0\\-3&0&6&0\\0&0&0&0\end{pmatrix}
\]
\item By Theorem~\ref{r1-thm-views}(a),(b), the columns of \(AF_1\) are \(\frac13(38a_1-15a_2)=\frac13(8,-7,1)'\),
\(\frac13(-15a_1+6a_2)=\frac13(-3,3,0)'\), \(0\) and \(0\). The rows of \(A'\) are \(a_1',\dots,a_4'\), so by the outer product
expansion (Theorem~\ref{r1-thm-views}(d))
\[
AF_1A'=\frac13\begin{pmatrix}8\\-7\\1\end{pmatrix}\begin{pmatrix}1&1&2\end{pmatrix}+\frac13\begin{pmatrix}-3\\3\\0\end{pmatrix}\begin{pmatrix}2&3&5\end{pmatrix}
=\frac13\begin{pmatrix}2&-1&1\\-1&2&1\\1&1&2\end{pmatrix}
\]
\item Likewise the columns of \(AF_2\) are \(\frac13(2a_1-3a_3)=\frac13(-1,2,1)'\), \(0\), \(\frac13(-3a_1+6a_3)=\frac13(3,-3,0)'\) and
\(0\), and
\[
AF_2A'=\frac13\begin{pmatrix}-1\\2\\1\end{pmatrix}\begin{pmatrix}1&1&2\end{pmatrix}+\frac13\begin{pmatrix}3\\-3\\0\end{pmatrix}\begin{pmatrix}1&0&1\end{pmatrix}
=\frac13\begin{pmatrix}2&-1&1\\-1&2&1\\1&1&2\end{pmatrix}
\]
\item So \(AF_1A'=AF_2A'=P\), as Theorem~\ref{r10-thm-ata}(b) says, and \(\tr P=2=\rank A\) (Theorem~\ref{r10-thm-ata}(d)). With
\(w=(1,1,-1)'\), \(P=I_3-\frac13ww'\), so \(Py=y-\frac13(w'y)w\). Here \(w'y=y_1+y_2-y_3=0\) exactly when \(y\in\mathcal C(A)\)
(Example~\ref{r10-exm-ginverse}); so \(Py=y\) for \(y\in\mathcal C(A)\), and \(Pw=0\).
\item \(\lambda=(1,0,3,-1)'\): only the entries \(\lambda_1=1\), \(\lambda_2=0\) meet the nonzero block of \(F_1\), so
\(\lambda'F_1\lambda=\frac13\cdot38=\frac{38}3\); with \(\lambda_1=1\), \(\lambda_3=3\),
\(\lambda'F_2\lambda=\frac13(2-9-9+54)=\frac{38}3\). Indeed \(\lambda'=3\row_1(A)-2\row_2(A)=a'A\) with \(a=(3,-2,0)'\), so
\(\lambda\in\mathcal C(A')\), and \(Pa=\frac13(8,-7,1)'\), \(\lVert Pa\rVert^2=\frac{64+49+1}9=\frac{38}3\), as
Corollary~\ref{r10-cor-invariant-quadratic} says.
\item \(\lambda=(1,0,0,0)'\): \(\lambda'F_1\lambda=\frac{38}3\), but \(\lambda'F_2\lambda=\frac23\). By
Corollary~\ref{r10-cor-invariant-quadratic}, \(\lambda\notin\mathcal C(A')\); correspondingly \(x_1\) is not the same for all
solutions in Example~\ref{r10-exm-general-solution}.
\end{itemize}
\end{example}

\subsection*{Minimum-norm and least-squares generalized inverses}

\begin{theorem}[\citeat{bapat}{Result~4.5}, \citeat{harville}{Theorem~20.3.6}]\label{r10-thm-min-norm}
Let \(A\in\R^{m\times n}\). A g-inverse \(G\) of \(A\) with \((GA)'=GA\) is a \textbf{minimum-norm g-inverse} of
\(A\). For \(G\in\R^{n\times m}\), the following are equivalent:
\begin{enumerate}[label=(\roman*)]
\item \(G\) is a minimum-norm g-inverse of \(A\);
\item for every \(y\in\mathcal C(A)\), \(x=Gy\) is a solution of \(Ax=y\), and \(\lVert Gy\rVert\leqslant\lVert x\rVert\) for
every solution \(x\).
\end{enumerate}
If they hold, then \(\lVert x\rVert>\lVert Gy\rVert\) for every solution \(x\ne Gy\): \(Gy\) is the only shortest solution.
\end{theorem}
\begin{proof}
For \(y=Au\in\mathcal C(A)\) and \(z\in\R^n\), Proposition~\ref{r6-prop-pythagoras} and
Theorem~\ref{r1-thm-transpose} give
\[
\lVert Gy+(I-GA)z\rVert^2=\lVert Gy\rVert^2+\lVert(I-GA)z\rVert^2+2c,\qquad c=y'G'(I-GA)z=u'(GA)'(I-GA)z
\tag{\(*\)}
\]
(i)\,\(\Rightarrow\)\,(ii)
\begin{itemize}
\item \(Gy\) is a solution, and every solution is \(x=Gy+(I-GA)z\) for some \(z\)
(Theorems~\ref{r10-thm-ginv-solves}(a) and~\ref{r10-thm-general-solution}).
\item \((GA)'=GA\), and \(GAGA=GA\) as \(AGA=A\); so \(c=u'GA(I-GA)z=u'(GA-GAGA)z=0\).
\item Hence \(\lVert x\rVert^2=\lVert Gy\rVert^2+\lVert(I-GA)z\rVert^2\geqslant\lVert Gy\rVert^2\), with equality if and
only if \((I-GA)z=0\), that is, \(x=Gy\).
\end{itemize}
(ii)\,\(\Rightarrow\)\,(i)
\begin{itemize}
\item By (ii) and Theorem~\ref{r10-thm-ginv-solves}(a), \(G\) is a g-inverse of \(A\).
\item Let \(u,z\in\R^n\), \(d=\lVert(I-GA)z\rVert^2\) and \(c=u'(GA)'(I-GA)z\). Apply (ii) to \(y=A(\alpha u)\) for
\(\alpha\in\R\) and to its solution \(Gy+(I-GA)z\) (Theorem~\ref{r10-thm-general-solution}): by \((*)\), with \(\alpha u\)
in place of \(u\), \(0\leqslant d+2\alpha c\) for every \(\alpha\in\R\).
\item If \(c\ne0\), then \(\alpha=-(d+1)/(2c)\) gives \(d+2\alpha c=-1<0\) \(\contradiction\) So \(c=0\) for all \(u,z\).
\item Taking for \(u\) and \(z\) the \(i\)th and \(j\)th columns of \(I_n\), \(c\) is the \((i,j)\) entry of \((GA)'(I-GA)\)
(Theorem~\ref{r1-thm-views}(b),(c));
so \((GA)'(I-GA)=0\), that is, \((GA)'=(GA)'GA\). This matrix is symmetric, so
\(GA=((GA)')'=(GA)'\).\qedhere
\end{itemize}
\end{proof}

\begin{theorem}[\citeat{bapat}{Result~4.6}, \citeat{harville}{Theorem~20.3.12}]\label{r10-thm-ls-ginverse}
Let \(A\in\R^{m\times n}\). A g-inverse \(G\) of \(A\) with \((AG)'=AG\) is a \textbf{least-squares g-inverse} of
\(A\). For \(G\in\R^{n\times m}\), the following are equivalent:
\begin{enumerate}[label=(\roman*)]
\item \(G\) is a least-squares g-inverse of \(A\);
\item \(\lVert AGy-y\rVert\leqslant\lVert Ax-y\rVert\) for all \(x\in\R^n\) and \(y\in\R^m\).
\end{enumerate}
So, for every \(y\), \(x=Gy\) minimizes \(\lVert Ax-y\rVert\), also when \(Ax=y\) is inconsistent.
\end{theorem}
\begin{proof}
For \(x=Gy+w\) (\(w\in\R^n\)), \(Ax-y=(AG-I)y+Aw\), and Proposition~\ref{r6-prop-pythagoras} gives
\[
\lVert Ax-y\rVert^2=\lVert(AG-I)y\rVert^2+\lVert Aw\rVert^2+2w'A'(AG-I)y
\tag{\(**\)}
\]
(i)\,\(\Rightarrow\)\,(ii) As \(AG=(AG)'\) and \(AGA=A\), \(A'(AG-I)=A'(AG)'-A'=(AGA)'-A'=0\), so \(\lVert Ax-y\rVert^2=\lVert AGy-y\rVert^2+\lVert Aw\rVert^2
\geqslant\lVert AGy-y\rVert^2\).

(ii)\,\(\Rightarrow\)\,(i)
\begin{itemize}
\item For \(x\in\R^n\) put \(y=Ax\) in (ii): \(\lVert AGAx-Ax\rVert\leqslant\lVert Ax-Ax\rVert=0\), so \(AGAx=Ax\). Taking for
\(x\) the columns of \(I_n\), \(AGA=A\) (Theorem~\ref{r1-thm-views}(b)).
\item Let \(w\in\R^n\) and \(y\in\R^m\). Apply (ii) to \(\alpha y\) (\(\alpha\in\R\)) and \(x=\alpha Gy+w\): by \((**)\),
\(\alpha^2\lVert(AG-I)y\rVert^2\leqslant\alpha^2\lVert(AG-I)y\rVert^2+\lVert Aw\rVert^2+2\alpha w'A'(AG-I)y\), that is,
\(0\leqslant d+2\alpha c\) with \(d=\lVert Aw\rVert^2\) and \(c=w'A'(AG-I)y\).
\item As in the proof of Theorem~\ref{r10-thm-min-norm}, \(c=0\) for all \(w\) and \(y\); so \(A'(AG-I)=0\), that is,
\(A'AG=A'\).
\item Transposing, \(G'A'A=A\); so \(AG=G'A'AG=(AG)'AG\), which is symmetric.\qedhere
\end{itemize}
\end{proof}
For \(A\in\R^{m\times n}\) and a g-inverse \(F\) of \(A'A\), \(G=FA'\) is a least-squares g-inverse of \(A\): with \(P=AFA'\),
\(AGA=PA=A\), and \(AG=P\) is symmetric (Theorem~\ref{r10-thm-ata}(c))~\citeat{harville}{Theorem~20.3.17}. It is even reflexive (Theorem~\ref{r10-thm-reflexive}):
\(GAG=F(A'AFA')=FA'P=FA'=G\), as \(A'P=A'\). On the other hand, for the matrix \(A\) of Example~\ref{r10-exm-ginverse}, \(G_1\) is
neither a least-squares nor a minimum-norm g-inverse: \(G_1A\) (Example~\ref{r10-exm-general-solution}) and
\[
AG_1=\begin{pmatrix}1&0&0\\0&1&0\\1&1&0\end{pmatrix}
\]
are not symmetric.

\needspace{12\baselineskip}
\subsection*{The Moore--Penrose inverse}

\begin{definition}[\citeat{bapat}{\S4.3}, \citeat{harville}{Theorem~20.1.1}]\label{r10-def-mp}
Let \(A\in\R^{m\times n}\). A matrix \(G\in\R^{n\times m}\) is a \textbf{Moore--Penrose inverse} of \(A\) if
\begin{enumerate}[label=(P\arabic*)]
\item \(AGA=A\);
\item \(GAG=G\);
\item \((AG)'=AG\);
\item \((GA)'=GA\).
\end{enumerate}
Thus a Moore--Penrose inverse of \(A\) is a reflexive g-inverse of \(A\) (Theorem~\ref{r10-thm-reflexive}) that is
both a least-squares and a minimum-norm g-inverse.
\end{definition}

\needspace{18\baselineskip}
\begin{theorem}[\citeat{bapat}{\S4.3}, \citeat{harville}{Theorem~20.1.1}]\label{r10-thm-mp-unique}
A matrix \(A\in\R^{m\times n}\) has at most one Moore--Penrose inverse. It is written \(A^+\).
\end{theorem}
\begin{proof}
Let \(G_1\) and \(G_2\) both satisfy (P1)--(P4). Each step below uses the condition shown, together with
\((BC)'=C'B'\) and \((B')'=B\) (Theorem~\ref{r1-thm-transpose}):
\[
\begin{aligned}
G_1&=G_1AG_1&&\text{(P2) for }G_1\\
&=G_1G_1'A'&&AG_1=(AG_1)'=G_1'A'\text{ by (P3)}\\
&=G_1G_1'A'G_2'A'&&A'=(AG_2A)'=A'G_2'A'\text{ by (P1)}\\
&=G_1G_1'A'AG_2&&G_2'A'=(AG_2)'=AG_2\text{ by (P3)}\\
&=G_1AG_1AG_2&&G_1'A'=(AG_1)'=AG_1\text{ by (P3)}\\
&=G_1AG_2&&AG_1A=A\text{ by (P1)}\\
&=G_1AG_2AG_2&&G_2=G_2AG_2\text{ by (P2)}\\
&=G_1AA'G_2'G_2&&G_2A=(G_2A)'=A'G_2'\text{ by (P4)}\\
&=A'G_1'A'G_2'G_2&&G_1A=(G_1A)'=A'G_1'\text{ by (P4)}\\
&=A'G_2'G_2&&A'G_1'A'=(AG_1A)'=A'\text{ by (P1)}\\
&=G_2AG_2&&A'G_2'=(G_2A)'=G_2A\text{ by (P4)}\\
&=G_2&&\text{(P2) for }G_2
\end{aligned}
\]
\end{proof}

\begin{theorem}[\citeat{bapat}{\S4.3}, \citeat{harville}{Theorem~20.1.1, \S20.2}]\label{r10-thm-mp-exists}
Every \(A\in\R^{m\times n}\) has a Moore--Penrose inverse. More precisely, let \(r=\rank A\).
\begin{enumerate}
\item If \(r=0\), then \(A^+=0\in\R^{n\times m}\).
\item If \(r\geqslant1\) and \(A=BC\) is a full-rank factorization, then \(A^+=C'(CC')^{-1}(B'B)^{-1}B'\).
\item If \(r=n\), then \(A^+=(A'A)^{-1}A'\); if \(r=m\), then \(A^+=A'(AA')^{-1}\); if \(A\) is square and invertible,
then \(A^+=A^{-1}\).
\end{enumerate}
\end{theorem}
\begin{proof}
(a) Then \(A=0\) (Theorem~\ref{thm:r4-four-bases}), and \(G=0\) satisfies (P1)--(P4), as every product in them is \(0\).

(b) Let \(G=C'(CC')^{-1}(B'B)^{-1}B'\). With \(L=(B'B)^{-1}B'\) and \(K=C'(CC')^{-1}\), \(G=KL\), \(LB=I_r\) and
\(CK=I_r\) (Theorem~\ref{thm:r4-full-rank-factorization}(c)).
\begin{itemize}
\item (P1): \(G\) is a g-inverse of \(A\) (Theorem~\ref{r10-thm-ginv-exists}(a)).
\item (P2): \(GAG=KLBCKL=K(LB)(CK)L=KL=G\).
\item (P3): \(AG=BCKL=BL=B(B'B)^{-1}B'\), the matrix \(P\) of Theorem~\ref{r6-thm-hat} for \(B\), which has full column rank;
it is symmetric
(Theorem~\ref{r6-thm-hat}(b)).
\item (P4): \(GA=KLBC=KC=C'(CC')^{-1}C\), the matrix \(P\) of Theorem~\ref{r6-thm-hat} for \(C'\in\R^{n\times r}\), which has
full column rank \(r\)
(Theorem~\ref{thm:r4-rank-rules}(a)); it is symmetric (Theorem~\ref{r6-thm-hat}(b)).
\end{itemize}
With (a), every matrix has a Moore--Penrose inverse, and it is unique (Theorem~\ref{r10-thm-mp-unique}).

(c) If \(r=n\), then \(A=AI_n\) is a full-rank factorization, and (b) gives
\(A^+=I_n(I_nI_n')^{-1}(A'A)^{-1}A'=(A'A)^{-1}A'\). If \(r=m\), then \(A=I_mA\) is one, and (b) gives
\(A^+=A'(AA')^{-1}(I_m'I_m)^{-1}I_m'=A'(AA')^{-1}\). If \(A\) is invertible, \(A^{-1}\) is its only g-inverse
(Definition~\ref{r10-def-ginverse}), so \(A^+=A^{-1}\).
\end{proof}

\needspace{9\baselineskip}
\begin{example}[\citeat{searle-khuri}{Example~8.1}]\label{r10-exm-mp}
Find \(A^+\) for
\[
A=\begin{pmatrix}1&0&-1&1\\0&2&2&2\\-1&4&5&3\end{pmatrix}
\]

\textit{Solution.} Row-reduce \(A\):
\[
A\xrightarrow{R_3\leftarrow R_3+R_1}
\begin{pmatrix}1&0&-1&1\\0&2&2&2\\0&4&4&4\end{pmatrix}
\xrightarrow{R_3\leftarrow R_3-2R_2}
\begin{pmatrix}1&0&-1&1\\0&2&2&2\\0&0&0&0\end{pmatrix}
\xrightarrow{R_2\leftarrow\frac12R_2}
\begin{pmatrix}1&0&-1&1\\0&1&1&1\\0&0&0&0\end{pmatrix}=R
\]
So \(\rank A=2\), with pivot columns \(1,2\), and Theorem~\ref{thm:r4-full-rank-factorization}(a) gives \(A=BC\) with
\[
B=\begin{pmatrix}1&0\\0&2\\-1&4\end{pmatrix},\qquad C=\begin{pmatrix}1&0&-1&1\\0&1&1&1\end{pmatrix},\qquad
B'B=\begin{pmatrix}2&-4\\-4&20\end{pmatrix},\qquad CC'=3I_2
\]
Hence \((B'B)(CC')=6\left(\begin{smallmatrix}1&-2\\-2&10\end{smallmatrix}\right)\), whose inverse is
\(\frac1{36}\left(\begin{smallmatrix}10&2\\2&1\end{smallmatrix}\right)\) (Theorem~\ref{r5d-thm-adjugate}; the determinant of
\(\left(\begin{smallmatrix}1&-2\\-2&10\end{smallmatrix}\right)\) is \(6\)). By Theorems~\ref{r10-thm-mp-exists}(b)
and~\ref{r2-thm-inverse-props}(b),
\[
\begin{aligned}
A^+&=C'(CC')^{-1}(B'B)^{-1}B'=C'\bigl((B'B)(CC')\bigr)^{-1}B'
=\frac1{36}\begin{pmatrix}1&0\\0&1\\-1&1\\1&1\end{pmatrix}\begin{pmatrix}10&2\\2&1\end{pmatrix}\begin{pmatrix}1&0&-1\\0&2&4\end{pmatrix}\\
&=\frac1{36}\begin{pmatrix}1&0\\0&1\\-1&1\\1&1\end{pmatrix}\begin{pmatrix}10&4&-2\\2&2&2\end{pmatrix}
=\frac1{36}\begin{pmatrix}10&4&-2\\2&2&2\\-8&-2&4\\12&6&0\end{pmatrix}
=\frac1{18}\begin{pmatrix}5&2&-1\\1&1&1\\-4&-1&2\\6&3&0\end{pmatrix}
\end{aligned}
\]
Check of (P3):
\[
AA^+=B(B'B)^{-1}B'=\frac16\begin{pmatrix}5&2&-1\\2&2&2\\-1&2&5\end{pmatrix}
\]
is symmetric; for instance, its first row is \(\row_1(A)A^+=\frac1{18}(5+0+4+6,\ 2+0+1+3,\ -1+0-2+0)=\frac16(5,2,-1)\).
\end{example}

\begin{proposition}[\citeat{harville}{Theorems~20.5.1, 20.5.3}, \citeat{bapat}{Result~6.3}, \citeat{abadir-magnus}{Exercises~10.25, 10.26, 10.28}]\label{r10-prop-mp-props}
Let \(A\in\R^{m\times n}\).
\begin{enumerate}
\item \(\rank A^+=\rank A\), \(\mathcal C(A^+)=\mathcal C(A')\), \(\mathcal R(A^+)=\mathcal R(A')\) and
\(\mathcal N(A^+)=\mathcal N(A')\).
\item \(AA^+\) and \(A^+A\) are symmetric and idempotent, with \(\mathcal C(AA^+)=\mathcal C(A)\) and
\(\mathcal C(A^+A)=\mathcal C(A')\). Hence \(AA^+y\) is the orthogonal projection of \(y\in\R^m\) onto \(\mathcal C(A)\), and
\(A^+Ax\) is that of \(x\in\R^n\) onto \(\mathcal C(A')\).
\item \((A')^+=(A^+)'\), \((A^+)^+=A\), and \((cA)^+=c^{-1}A^+\) for every number \(c\ne0\).
\item If \(A\) is symmetric, then \(A^+\) is symmetric and \(AA^+=A^+A\). If \(A\) is PSD, then \(A^+\) is PSD.
\end{enumerate}
\end{proposition}
\begin{proof}
(a)
\begin{itemize}
\item By (P1) and (P2), \(A^+\) is a reflexive g-inverse of \(A\); so \(\rank A^+=\rank A\)
(Theorem~\ref{r10-thm-reflexive}(b)).
\item By (P2), (P4) and Theorem~\ref{r1-thm-transpose}, \(A^+=A^+AA^+=(A^+A)'A^+=A'(A^+)'A^+\); so
\(\mathcal C(A^+)\subseteq\mathcal C(A')\) (Theorem~\ref{thm:r4-rank-product}(a)). Both spaces have dimension
\(\rank A^+=\rank A=\rank A'\) (Theorem~\ref{thm:r4-rank-rules}(a)), so they are equal (Theorem~\ref{thm:r4-counting}(a)).
\item Likewise, by (P2) and (P3), \(A^+=A^+(AA^+)'=A^+(A^+)'A'\): each row of \(A^+\) is \(\row_i(A^+(A^+)')A'\in\mathcal R(A')\)
(Theorem~\ref{r1-thm-views}(c)). So \(\mathcal R(A^+)\subseteq\mathcal R(A')\), and the dimensions agree as before; the
transposed spaces \(\mathcal C((A^+)')\subseteq\mathcal C(A)\) are therefore equal (Theorem~\ref{thm:r4-counting}(a)).
\item Hence \(\mathcal N(A^+)=\mathcal C((A^+)')^\perp=\mathcal C(A)^\perp=\mathcal N(A')\) (Corollary~\ref{r6-cor-four-subspaces}).
\end{itemize}
\needspace{5\baselineskip}
(b)
\begin{itemize}
\item \(AA^+\) and \(A^+A\) are symmetric by (P3) and (P4), and idempotent by Theorem~\ref{r10-thm-rank-idempotent}(a), as
\(A^+\) is a g-inverse of \(A\).
\item \(\mathcal C(AA^+)=\mathcal C(A)\) (Theorem~\ref{r10-thm-rank-idempotent}(b)). Also \(\mathcal C(A^+A)\subseteq\mathcal C(A^+)\)
(Theorem~\ref{thm:r4-rank-product}(a)) and \(\rank(A^+A)=\rank A=\rank A^+\) (Theorem~\ref{r10-thm-rank-idempotent}(b) and
part (a)); so \(\mathcal C(A^+A)=\mathcal C(A^+)=\mathcal C(A')\) (Theorem~\ref{thm:r4-counting}(a)).
\item The projections: Theorem~\ref{r6-thm-idempotent}(d).
\end{itemize}
(c) Each claim is checked through (P1)--(P4) and the uniqueness of Theorem~\ref{r10-thm-mp-unique}.
\begin{itemize}
\item \((A^+)'\) for \(A'\): \(A'(A^+)'A'=(AA^+A)'=A'\), \((A^+)'A'(A^+)'=(A^+AA^+)'=(A^+)'\), and
\(A'(A^+)'=(A^+A)'=A^+A\) and \((A^+)'A'=(AA^+)'=AA^+\) are symmetric.
\item \(A\) for \(A^+\): the four conditions \(A^+AA^+=A^+\), \(AA^+A=A\), \((A^+A)'=A^+A\) and \((AA^+)'=AA^+\) are (P2),
(P1), (P4) and (P3).
\item \(c^{-1}A^+\) for \(cA\): \((cA)(c^{-1}A^+)(cA)=cAA^+A=cA\), \((c^{-1}A^+)(cA)(c^{-1}A^+)=c^{-1}A^+AA^+=c^{-1}A^+\), and
\((cA)(c^{-1}A^+)=AA^+\) and \((c^{-1}A^+)(cA)=A^+A\) are symmetric.
\end{itemize}
(d)
\begin{itemize}
\item If \(A'=A\), then \((A^+)'=(A')^+=A^+\) by (c), and \(AA^+=(AA^+)'=(A^+)'A'=A^+A\) by (P3).
\item If \(A\) is PSD (so symmetric), then \(A^+\) is symmetric and \(A^+=A^+AA^+=(A^+)'AA^+\); so \(A^+\) is PSD
(Theorem~\ref{r8-thm-congruence}(a), with \(P=A^+\)).\qedhere
\end{itemize}
\end{proof}

\begin{proposition}[\citeat{abadir-magnus}{Exercise~10.35}]\label{r10-prop-mp-spectral}
Let \(A\in\R^{n\times n}\) be symmetric, with \(A=Q\Lambda Q'\) as in Theorem~\ref{r8-thm-spectral}(b): \(Q\) is orthogonal, with
columns \(q_1,\dots,q_n\), and \(\Lambda=\diag(\lambda_1,\dots,\lambda_n)\). For a number \(\lambda\) let \(\lambda^+=1/\lambda\) if
\(\lambda\ne0\) and \(0^+=0\), and let \(\Lambda^+=\diag(\lambda_1^+,\dots,\lambda_n^+)\). Then
\[
A^+=Q\Lambda^+Q'=\sum_{\lambda_i\ne0}\lambda_i^{-1}q_iq_i'
\]
\end{proposition}
\begin{proof}
\begin{itemize}
\item For every number \(\lambda\), \(\lambda\lambda^+\lambda=\lambda\) and \(\lambda^+\lambda\lambda^+=\lambda^+\) (both sides are \(0\)
if \(\lambda=0\)). A product of diagonal matrices is the diagonal matrix of the products of their diagonal entries; so
\(\Lambda\Lambda^+\Lambda=\Lambda\), \(\Lambda^+\Lambda\Lambda^+=\Lambda^+\), and \(\Lambda\Lambda^+=\Lambda^+\Lambda\) is diagonal.
\item Let \(G=Q\Lambda^+Q'\). As \(Q'Q=I_n\), \(AGA=Q\Lambda\Lambda^+\Lambda Q'=A\) and \(GAG=Q\Lambda^+\Lambda\Lambda^+Q'=G\): (P1),
(P2). Also \(AG=Q\Lambda\Lambda^+Q'=GA\), which is symmetric (Theorem~\ref{r1-thm-transpose}): (P3), (P4). So \(G=A^+\)
(Theorem~\ref{r10-thm-mp-unique}).
\item Column \(i\) of \(Q\Lambda^+\) is \(\lambda_i^+q_i\), so by the outer product expansion (Theorem~\ref{r1-thm-views}(d))
\(Q\Lambda^+Q'=\sum_i\lambda_i^+q_iq_i'\), in which the terms with \(\lambda_i=0\) are \(0\).\qedhere
\end{itemize}
\end{proof}

For a symmetric \(A\ne0\), the symmetric g-inverse \(G=\sum_{c_j\ne0}c_j^{-1}E_j\) of Lemma~\ref{r15-lem-function} is \(A^+\): by
Theorem~\ref{r8-thm-spectral-decomp}(c), \(E_j=\sum_{\lambda_i=c_j}q_iq_i'\), so \(G=\sum_{\lambda_i\ne0}\lambda_i^{-1}q_iq_i'\), which is
\(A^+\) by Proposition~\ref{r10-prop-mp-spectral}.

\needspace{14\baselineskip}
\subsection*{The minimum-norm least-squares solution}

\begin{theorem}[\citeat{harville}{Corollary~20.5.5, Theorem~20.6.1}, \citeat{abadir-magnus}{Exercise~10.29(e)}]\label{r10-thm-mp-shortest}
Let \(A\in\R^{m\times n}\) and \(y\in\R^m\).
\begin{enumerate}
\item \(A^+=(A'A)^+A'\), and \(AA^+=A(A'A)^-A'\), the matrix \(P\) of Theorem~\ref{r10-thm-ata}.
\item A vector \(x\in\R^n\) minimizes \(\lVert Ax-y\rVert\) if and only if \(Ax=AA^+y\). The vector \(A^+y\) is one of them, and
every other one has \(\lVert x\rVert>\lVert A^+y\rVert\): \(A^+y\) is the only shortest minimizer of \(\lVert Ax-y\rVert\), the
\textbf{minimum-norm least-squares solution} of \(Ax=y\). It is also the only minimizer in \(\mathcal C(A')\).
\item If \(G\in\R^{n\times m}\) and \(Gy\) is the shortest minimizer of \(\lVert Ax-y\rVert\) for every \(y\in\R^m\), then
\(G=A^+\).
\end{enumerate}
\end{theorem}
\begin{proof}
(a) Let \(F=(A'A)^+\) (Theorem~\ref{r10-thm-mp-exists}) and \(H=FA'\).
\begin{itemize}
\item \(F\) is a g-inverse of \(A'A\) by (P1), so \(AH=AFA'=P\) (Theorem~\ref{r10-thm-ata}(b)).
\item (P1): \(AHA=PA=A\); (P3): \(AH=P\) is symmetric (Theorem~\ref{r10-thm-ata}(c)).
\item (P2): \(HAH=(FA'AF)A'=FA'=H\), by (P2) for \(F\).
\item (P4): \(HA=FA'A\) is symmetric, by (P4) for \(F\), a Moore--Penrose inverse of \(A'A\).
\item So \(H=A^+\) (Theorem~\ref{r10-thm-mp-unique}), and \(AA^+=AH=P\).
\end{itemize}
(b)
\begin{itemize}
\item By Proposition~\ref{r10-prop-mp-props}(b), \(AA^+y\) is the orthogonal projection of \(y\) onto \(\mathcal C(A)\). Every
\(Ax\) lies in \(\mathcal C(A)\), so by Theorem~\ref{r6-thm-best-approximation}, \(\lVert Ax-y\rVert\geqslant\lVert AA^+y-y\rVert\),
with equality exactly when \(Ax=AA^+y\). In particular \(x=A^+y\) is a minimizer.
\item The minimizers are also the solutions of \(A'Ax=A'y\). If \(Ax=AA^+y\), then
\(A'Ax=A'(AA^+)'y=(AA^+A)'y=A'y\), by (P3), Theorem~\ref{r1-thm-transpose} and (P1). Conversely, if \(A'Ax=A'y\), then
multiplying by \((A^+)'\) on the left gives \((AA^+)'Ax=(AA^+)'y\), that is, \(AA^+Ax=AA^+y\) by (P3), and \(Ax=AA^+y\) by (P1).
\item \(A'y\in\mathcal C(A')=\mathcal C(A'A)\) (Theorem~\ref{thm:r4-rank-ata}(c)), and \(F=(A'A)^+\) satisfies (P1) and (P4) for
\(A'A\): it is a minimum-norm g-inverse of \(A'A\). By Theorem~\ref{r10-thm-min-norm}, for the matrix \(A'A\) and the right side
\(A'y\), \(FA'y\) is a solution of \(A'Ax=A'y\), and every other solution \(x\) has \(\lVert x\rVert>\lVert FA'y\rVert\). Here
\(FA'y=A^+y\) by (a).
\item If \(x\) is a minimizer, then \(A(x-A^+y)=AA^+y-AA^+y=0\), so \(x-A^+y\in\mathcal N(A)=\mathcal C(A')^\perp\)
(Corollary~\ref{r6-cor-four-subspaces}); and \(A^+y\in\mathcal C(A^+)=\mathcal C(A')\) (Proposition~\ref{r10-prop-mp-props}(a)).
If moreover \(x\in\mathcal C(A')\), then \(x-A^+y\in\mathcal C(A')\cap\mathcal C(A')^\perp=\{0\}\)
(Theorem~\ref{r6-thm-complement}(b)), that is, \(x=A^+y\).
\end{itemize}
(c) For every \(y\), \(Gy\) and \(A^+y\) are shortest minimizers of \(\lVert Ax-y\rVert\), and by (b) there is only one; so
\(Gy=A^+y\). Taking for \(y\) the columns of \(I_m\) gives \(G=A^+\) (Theorem~\ref{r1-thm-views}(b)).
\end{proof}

\begin{example}[cf.\ \citeat{searle-khuri}{Example~8.1}]\label{r10-exm-mp-shortest}
For \(A\) of Example~\ref{r10-exm-mp} and \(y=(0,0,6)'\), decide whether \(Ax=y\) is consistent, and find the minimizers of
\(\lVert Ax-y\rVert\) and the minimum-norm least-squares solution.

\textit{Solution.}
\begin{itemize}
\item By Example~\ref{r10-exm-mp},
\[
A^+=\frac1{18}\begin{pmatrix}5&2&-1\\1&1&1\\-4&-1&2\\6&3&0\end{pmatrix},\qquad
AA^+=\frac16\begin{pmatrix}5&2&-1\\2&2&2\\-1&2&5\end{pmatrix}
\]
so \(AA^+y=(-1,2,5)'\ne y\), and \(Ax=y\) is inconsistent (Theorem~\ref{r10-thm-ginv-solves}(b); \(A^+\) is a g-inverse of \(A\) by
(P1)).
\item \(y-AA^+y=(1,-2,1)'\), and \(A'(1,-2,1)'=(1-1,\ -4+4,\ -1-4+5,\ 1-4+3)'=0\): it is orthogonal to \(\mathcal C(A)\). So the
smallest value of \(\lVert Ax-y\rVert^2\) is \(\lVert(1,-2,1)'\rVert^2=6\) (Theorem~\ref{r10-thm-mp-shortest}(b)).
\item \(A^+y=\frac6{18}(-1,1,2,0)'=\frac13(-1,1,2,0)'\), with \(\lVert A^+y\rVert^2=\frac69=\frac23\); check:
\(A(-1,1,2,0)'=-a_1+a_2+2a_3=(-3,6,15)'=3(-1,2,5)'\).
\item By Theorem~\ref{r10-thm-mp-shortest}(b), the minimizers are the solutions of \(Ax=(-1,2,5)'\), that is, \(A^+y+z\) with
\(z\in\mathcal N(A)\) (Theorem~\ref{thm:r2-solution-structure}(a)). The RREF \(R\) of Example~\ref{r10-exm-mp} gives the special
solutions \(s_1=(1,-1,1,0)'\) and \(s_2=(-1,-1,0,1)'\), a basis of \(\mathcal N(A)\) (Theorem~\ref{thm:r4-four-bases}(c)); so the
minimizers are
\[
x=\frac13\begin{pmatrix}-1\\1\\2\\0\end{pmatrix}+t_1\begin{pmatrix}1\\-1\\1\\0\end{pmatrix}+t_2\begin{pmatrix}-1\\-1\\0\\1\end{pmatrix}
\qquad(t_1,t_2\in\R)
\]
\item For instance \(x_0=(-1,1,0,0)'\) (\(t_1=-\frac23\), \(t_2=0\)) has \(Ax_0=-a_1+a_2=(-1,2,5)'\), and
\(\lVert x_0\rVert^2=2>\frac23\). Here \(x_0-A^+y=-\frac23s_1\) is orthogonal to \(A^+y\), as \((-1,1,2,0)s_1=-1-1+2=0\), and
\(\lVert x_0\rVert^2=\frac23+\frac49\cdot3=2\) (Proposition~\ref{r6-prop-pythagoras}).
\item \((-1,1,2,0)=-(1,0,-1,1)+(0,1,1,1)\) is a combination of the nonzero rows of \(R\), so \(A^+y\in\mathcal C(A')\)
(Theorem~\ref{thm:r4-four-bases}(b)), as Theorem~\ref{r10-thm-mp-shortest}(b) says.
\end{itemize}
\end{example}

\needspace{14\baselineskip}
\subsection*{Pearson's statistic}

\begin{example}[\citeat{casella-berger}{Exercise~4.39}]\label{r10-exm-pearson}
In \(N\) independent trials, each trial results in exactly one of \(k\geqslant2\) categories, category \(j\) with
probability \(p_j>0\) (\(p_1+\dots+p_k=1\)). Let \(y_i\in\R^k\) record trial \(i\): \(y_i=e_j\) if trial \(i\) results in
category \(j\). Let \(c=y_1+\dots+y_N\) be the vector of counts, \(p=(p_1,\dots,p_k)'\), \(D=\diag(p_1,\dots,p_k)\),
\(V=D-pp'\) and \(u=(c-Np)/\sqrt N\). \textbf{Pearson's statistic} is
\[
X^2=\sum_{j=1}^k\frac{(c_j-Np_j)^2}{Np_j}
\]
where \(Np_j=\expc(c_j)\) is the expected count of cell \(j\). We show:
\begin{enumerate}
\item \(\expc(c)=Np\) and \(\var(c)=NV\); in particular \(\cov(c_i,c_j)=-Np_ip_j\) for \(i\ne j\);
\item \(V\one=0\), \(\mathcal N(V)=\Span\{\one\}\) and \(\rank V=k-1\), so \(\var(c)\) is singular;
\item \(D^{-1}\) is a g-inverse of \(V\);
\item \(X^2=u'D^{-1}u\), and \(u'Gu=X^2\) for every g-inverse \(G\) of \(V\);
\item let \(u_*\sim N_k(0,V)\), and take \(u_*\) as the normal approximation to \(u\) for large \(N\) (assumed, not proved
here), so that \(X^2=u'D^{-1}u\) is approximated by \(u_*'D^{-1}u_*\); then \(u_*'D^{-1}u_*\sim\chi^2_{k-1}\).
\end{enumerate}

\textit{Solution.}
(a)
\begin{itemize}
\item Entry \(j\) of \(y_i\) is \(1\) with probability \(p_j\) and \(0\) otherwise, so its expectation is \(p_j\): \(\expc(y_i)=p\).
As \(y_i\) is one of \(e_1,\dots,e_k\), \(y_iy_i'=\diag(y_i)\), so \(\expc(y_iy_i')=D\) and \(\var(y_i)=D-pp'=V\)
(Theorem~\ref{r14-thm-cov}(e)).
\item With \(Z=[y_1\ \cdots\ y_N]\in\R^{k\times N}\), \(c=Z\one_N\), and \(\expc(c)=\expc(Z)\one_N=Np\)
(Theorem~\ref{r11-thm-linear}(a)).
\item The trials are independent, so \(y_1,\dots,y_N\) are independent and \(\cov(y_i,y_j)=0\) for \(i\ne j\) (N7). By
Theorem~\ref{r14-thm-cov}(d), \(\var(c)=\var(y_1)+\dots+\var(y_N)=NV\). Its \((i,j)\) entry for \(i\ne j\) is
\(\cov(c_i,c_j)=-Np_ip_j\) (Definition~\ref{r11-def-random}).
\end{itemize}
(b) \(V\one=p-p(p'\one)=p-p=0\). If \(Vx=0\), then \(Dx=p(p'x)\), so \(p_jx_j=p_j(p'x)\) and \(x_j=p'x\) for every \(j\) (as
\(p_j>0\)): \(x=(p'x)\one\). Hence \(\mathcal N(V)=\Span\{\one\}\), and \(\rank V=k-1\) (Theorem~\ref{thm:r4-rank-nullity}). So
\(\var(c)=NV\) is singular; correspondingly \(\one'c=N\) for every outcome of the trials, a constant, as
Theorem~\ref{r11-thm-cov-psd} requires.

(c) As \(p'D^{-1}=\one'\), \(VD^{-1}=I_k-p\one'\); as \(\one'D=p'\) and \(\one'p=1\),
\[
VD^{-1}V=(I_k-p\one')(D-pp')=D-pp'-pp'+p(\one'p)p'=D-pp'=V
\]
So \(D^{-1}\) is a g-inverse of \(V\) (Definition~\ref{r10-def-ginverse}).

(d)
\begin{itemize}
\item \(u'D^{-1}u=\sum_ju_j^2/p_j=\sum_j(c_j-Np_j)^2/(Np_j)=X^2\).
\item \(\one'u=(\one'c-N\one'p)/\sqrt N=0\), so \(u\in W=\Span\{\one\}^\perp\), a subspace of dimension \(k-1\)
(Theorem~\ref{r6-thm-complement}(c)). \(W\) contains \(\mathcal C(V)\), since \(\one'Vw=(V\one)'w=0\) for every \(w\in\R^k\), and
\(\dim\mathcal C(V)=\rank V=k-1\) by (b); so \(\mathcal C(V)=W\) (Theorem~\ref{thm:r4-counting}(a)), and \(u=Vw\) for some
\(w\in\R^k\).
\item For every g-inverse \(G\) of \(V\), \(u'Gu=w'VGVw=w'Vw\), the same number for all \(G\) (this is
Theorem~\ref{r10-thm-invariance}(a) with \(B=u'\), \(C=u\)). By (c) it equals \(u'D^{-1}u=X^2\).
\end{itemize}
(e)
\begin{itemize}
\item Let \(q=(\sqrt{p_1},\dots,\sqrt{p_k})'\), so \(q'q=1\), and \(D^{-1/2}=\diag(p_1^{-1/2},\dots,p_k^{-1/2})\). Then
\(D^{-1/2}p=q\) and \(D^{-1/2}VD^{-1/2}=I_k-qq'=P\), which is symmetric and idempotent (\(P^2=I_k-2qq'+q(q'q)q'=P\)), of rank
\(\tr P=k-q'q=k-1\) (Theorems~\ref{r6-thm-idempotent}(c) and~\ref{r1-thm-trace}).
\item By Theorem~\ref{r14-thm-mvn-linear}, \(v=D^{-1/2}u_*\sim N_k(0,P)\), and for \(z\sim N_k(0,I_k)\) also \(Pz\sim N_k(0,PP')=N_k(0,P)\).
So \(v\) and \(Pz\) have the same MGF, hence the same distribution (N1), and \(v'v\) has the distribution of
\((Pz)'(Pz)=z'Pz\).
\item \(u_*'D^{-1}u_*=v'v\), and \(z'Pz\sim\chi^2_{k-1}\) (Theorem~\ref{r15-thm-quadratic-chisq} with \(\mu=0\), \(\sigma=1\)). Hence
\(u_*'D^{-1}u_*\sim\chi^2_{k-1}\).
\end{itemize}
\end{example}

\needspace{14\baselineskip}
\subsection*{Exercises}

\begin{exercise}\label{r10-ex-ginverse-drill}
Let
\[
A=\begin{pmatrix}1&2&1&0\\-1&-2&0&1\\1&2&2&1\end{pmatrix}
\]
\begin{enumerate}
\item Find a full-rank factorization \(A=BC\) from the RREF of \(A\), a left inverse \(L\) of \(B\) and a right inverse
\(K\) of \(C\) with integer entries, and the g-inverse \(G=KL\) of \(A\).
\item For which \(y\in\R^3\) is \(Ax=y\) consistent? Decide for \(y=(1,1,3)'\) and \(y=(1,0,1)'\).
\item Find all solutions of \(Ax=(1,1,3)'\).
\item Which of \(x_1+2x_2-x_4\), \(x_3+x_4\) and \(x_1+2x_2\) have the same value for all these solutions?
\end{enumerate}
\end{exercise}
\begin{solution}
(a) Row-reduce \(A\):
\[
\begin{gathered}
A\xrightarrow[R_3\leftarrow R_3-R_1]{R_2\leftarrow R_2+R_1}
\begin{pmatrix}1&2&1&0\\0&0&1&1\\0&0&1&1\end{pmatrix}
\xrightarrow{R_3\leftarrow R_3-R_2}
\begin{pmatrix}1&2&1&0\\0&0&1&1\\0&0&0&0\end{pmatrix}
\xrightarrow{R_1\leftarrow R_1-R_2}
\begin{pmatrix}1&2&0&-1\\0&0&1&1\\0&0&0&0\end{pmatrix}=R
\end{gathered}
\]
The pivot columns are \(1,3\), so (Theorem~\ref{thm:r4-full-rank-factorization}(a))
\[
B=\begin{pmatrix}1&1\\-1&0\\1&2\end{pmatrix},\qquad C=\begin{pmatrix}1&2&0&-1\\0&0&1&1\end{pmatrix}
\]
\begin{itemize}
\item Columns \(1\), \(3\) of \(C\) are the unit vectors of \(\R^2\), and rows \(1\), \(2\) of \(B\) form
\(B_1=\left(\begin{smallmatrix}1&1\\-1&0\end{smallmatrix}\right)\), with
\(B_1\left(\begin{smallmatrix}0&-1\\1&1\end{smallmatrix}\right)=I_2\), so
\(B_1^{-1}=\left(\begin{smallmatrix}0&-1\\1&1\end{smallmatrix}\right)\) (Theorem~\ref{r2-thm-invertible}). So
\[
K=\begin{pmatrix}1&0\\0&0\\0&1\\0&0\end{pmatrix},\qquad
L=[B_1^{-1}\ \ 0]=\begin{pmatrix}0&-1&0\\1&1&0\end{pmatrix}
\]
give \(CK=I_2\) and \(LB=B_1^{-1}B_1=I_2\) (Theorems~\ref{r1-thm-views}(a),(b) and~\ref{r5b-thm-block}(c)).
\item By Theorem~\ref{r10-thm-ginv-exists}(a),
\[
G=KL=\begin{pmatrix}0&-1&0\\0&0&0\\1&1&0\\0&0&0\end{pmatrix}
\]
is a g-inverse of \(A\).
\end{itemize}
(b) Since \(CK=I_2\),
\[
AG=BCKL=BL=\begin{pmatrix}1&1\\-1&0\\1&2\end{pmatrix}\begin{pmatrix}0&-1&0\\1&1&0\end{pmatrix}
=\begin{pmatrix}1&0&0\\0&1&0\\2&1&0\end{pmatrix}
\]
so \(AGy=(y_1,\ y_2,\ 2y_1+y_2)'\). By Theorem~\ref{r10-thm-ginv-solves}(b), \(Ax=y\) is consistent if and only if
\(y_3=2y_1+y_2\). For \(y=(1,1,3)'\), \(3=2+1\): consistent. For \(y=(1,0,1)'\), \(AGy=(1,0,2)'\ne y\): inconsistent.

(c) \(Gy=(-1,\ 0,\ 1+1,\ 0)'=(-1,0,2,0)'\), and \(GA=KLBC=KC\), so
\[
GA=\begin{pmatrix}1&2&0&-1\\0&0&0&0\\0&0&1&1\\0&0&0&0\end{pmatrix},\qquad
I_4-GA=\begin{pmatrix}0&-2&0&1\\0&1&0&0\\0&0&0&-1\\0&0&0&1\end{pmatrix}
\]
By Theorem~\ref{r10-thm-general-solution}, the solutions are
\[
x=\begin{pmatrix}-1\\0\\2\\0\end{pmatrix}+z_2\begin{pmatrix}-2\\1\\0\\0\end{pmatrix}+z_4\begin{pmatrix}1\\0\\-1\\1\end{pmatrix}
=\begin{pmatrix}-1-2z_2+z_4\\z_2\\2-z_4\\z_4\end{pmatrix}\qquad(z_2,z_4\in\R)
\]
Check: \(A(-1,0,2,0)'=-a_1+2a_3=(-1+2,\ 1+0,\ -1+4)'=(1,1,3)'\).

(d) Use Theorem~\ref{r10-thm-invariant-combination}(iii) with the rows of \(GA\) (Theorem~\ref{r1-thm-views}(c)).
\begin{itemize}
\item \(\lambda=(1,2,0,-1)'\): \(\lambda'GA=\row_1(GA)+2\row_2(GA)-\row_4(GA)=(1,2,0,-1)=\lambda'\); indeed
\(x_1+2x_2-x_4=(-1-2z_2+z_4)+2z_2-z_4=-1\) for all solutions.
\item \(\lambda=(0,0,1,1)'\): \(\lambda'GA=\row_3(GA)+\row_4(GA)=(0,0,1,1)=\lambda'\); indeed \(x_3+x_4=2\).
\item \(\lambda=(1,2,0,0)'\): \(\lambda'GA=\row_1(GA)+2\row_2(GA)=(1,2,0,-1)\ne\lambda'\); indeed \(x_1+2x_2=-1+z_4\) changes
with \(z_4\).
\end{itemize}
\end{solution}

\begin{exercise}\label{r10-dr-rowreduction}
Let
\[
A=\begin{pmatrix}1&-1&2&1\\2&-1&3&4\\1&0&1&3\end{pmatrix}
\]
\begin{enumerate}
\item Row-reduce \([A\mid I_3]\) to \([R\mid P]\), and write down the g-inverse \(G\) of
Corollary~\ref{r10-cor-ginv-rowreduction}.
\item Compute \(AG\), and check that \(AGA=A\).
\item Decide whether \(Ax=y\) is consistent for \(y=(1,3,2)'\) and for \(y=(2,1,0)'\), and find all solutions of the
consistent system.
\end{enumerate}
\end{exercise}
\begin{solution}
(a)
\[
\begin{gathered}
\left(\begin{array}{cccc|ccc}1&-1&2&1&1&0&0\\2&-1&3&4&0&1&0\\1&0&1&3&0&0&1\end{array}\right)
\xrightarrow[R_3\leftarrow R_3-R_1]{R_2\leftarrow R_2-2R_1}
\left(\begin{array}{cccc|ccc}1&-1&2&1&1&0&0\\0&1&-1&2&-2&1&0\\0&1&-1&2&-1&0&1\end{array}\right)
\\
\xrightarrow{R_3\leftarrow R_3-R_2}
\left(\begin{array}{cccc|ccc}1&-1&2&1&1&0&0\\0&1&-1&2&-2&1&0\\0&0&0&0&1&-1&1\end{array}\right)
\xrightarrow{R_1\leftarrow R_1+R_2}
\left(\begin{array}{cccc|ccc}1&0&1&3&-1&1&0\\0&1&-1&2&-2&1&0\\0&0&0&0&1&-1&1\end{array}\right)=[R\mid P]
\end{gathered}
\]
\(R\) is the RREF of \(A\), with pivot columns \(1,2\); so \(\rank A=2\). By Corollary~\ref{r10-cor-ginv-rowreduction},
rows \(1\), \(2\) of \(G\) are rows \(1\), \(2\) of \(P\), and rows \(3\), \(4\) are zero:
\[
G=\begin{pmatrix}-1&1&0\\-2&1&0\\0&0&0\\0&0&0\end{pmatrix}
\]
(b) By the outer product expansion (Theorem~\ref{r1-thm-views}(d)), with \(a_1=(1,2,1)'\) and \(a_2=(-1,-1,0)'\),
\[
AG=a_1\row_1(G)+a_2\row_2(G)=\begin{pmatrix}-1&1&0\\-2&2&0\\-1&1&0\end{pmatrix}+\begin{pmatrix}2&-1&0\\2&-1&0\\0&0&0\end{pmatrix}
=\begin{pmatrix}1&0&0\\0&1&0\\-1&1&0\end{pmatrix}
\]
By Theorem~\ref{r1-thm-views}(c), \(AGA\) has the rows \(\row_1(A)\), \(\row_2(A)\) and
\(-\row_1(A)+\row_2(A)=(1,0,1,3)=\row_3(A)\): \(AGA=A\).

(c) \(AGy=(y_1,\ y_2,\ -y_1+y_2)'\), so \(Ax=y\) is consistent if and only if \(y_3=y_2-y_1\)
(Theorem~\ref{r10-thm-ginv-solves}(b)).
\begin{itemize}
\item \(y=(2,1,0)'\): \(AGy=(2,1,-1)'\ne y\), so \(Ax=y\) is inconsistent.
\item \(y=(1,3,2)'\): \(AGy=(1,3,2)'=y\), so \(Ax=y\) is consistent, and \(Gy=(-1+3,\ -2+3,\ 0,\ 0)'=(2,1,0,0)'\) is a
solution.
\item Rows \(1\), \(2\) of \(GA\) are rows \(1\), \(2\) of \(PA=R\), and rows \(3\), \(4\) are zero (Theorem~\ref{r1-thm-views}(c)):
\[
GA=\begin{pmatrix}1&0&1&3\\0&1&-1&2\\0&0&0&0\\0&0&0&0\end{pmatrix},\qquad
I_4-GA=\begin{pmatrix}0&0&-1&-3\\0&0&1&-2\\0&0&1&0\\0&0&0&1\end{pmatrix}
\]
\item By Theorem~\ref{r10-thm-general-solution}, the solutions are
\[
x=\begin{pmatrix}2\\1\\0\\0\end{pmatrix}+z_3\begin{pmatrix}-1\\1\\1\\0\end{pmatrix}+z_4\begin{pmatrix}-3\\-2\\0\\1\end{pmatrix}
=\begin{pmatrix}2-z_3-3z_4\\1+z_3-2z_4\\z_3\\z_4\end{pmatrix}\qquad(z_3,z_4\in\R)
\]
(columns \(1\), \(2\) of \(I_4-GA\) are \(0\)). Check: \(A(2,1,0,0)'=2a_1+a_2=(2-1,\ 4-1,\ 2+0)'=(1,3,2)'\).
\end{itemize}
\end{solution}

\begin{exercise}[cf.\ \citeat{bapat}{Ch.~4, Exercise~2}]\label{r10-dr-submatrix}
Let
\[
A=\begin{pmatrix}1&2&1\\0&1&1\\1&3&2\end{pmatrix}
\]
\begin{enumerate}
\item Find \(\rank A\) by row reduction.
\item Find the g-inverse \(G_1\) of Theorem~\ref{r10-thm-ginv-submatrix}(b) built from the leading \(2\times2\) submatrix,
and check that \(AG_1A=A\).
\item Find the g-inverse \(G_2\) built from the submatrix in rows \(2,3\) and columns \(1,3\), and check that
\(AG_2A=A\).
\end{enumerate}
\end{exercise}
\begin{solution}
(a)
\[
A\xrightarrow{R_3\leftarrow R_3-R_1}
\begin{pmatrix}1&2&1\\0&1&1\\0&1&1\end{pmatrix}
\xrightarrow{R_3\leftarrow R_3-R_2}
\begin{pmatrix}1&2&1\\0&1&1\\0&0&0\end{pmatrix}
\xrightarrow{R_1\leftarrow R_1-2R_2}
\begin{pmatrix}1&0&-1\\0&1&1\\0&0&0\end{pmatrix}
\]
So \(\rank A=2\); the steps show \(\row_3(A)=\row_1(A)+\row_2(A)\).

(b) The leading block \(A_{11}=\left(\begin{smallmatrix}1&2\\0&1\end{smallmatrix}\right)\) is invertible, with
\(A_{11}^{-1}=\left(\begin{smallmatrix}1&-2\\0&1\end{smallmatrix}\right)\) (the product is \(I_2\)). By
Theorem~\ref{r10-thm-ginv-submatrix}(b),
\[
G_1=\begin{pmatrix}1&-2&0\\0&1&0\\0&0&0\end{pmatrix},\qquad
AG_1=\begin{pmatrix}1&0&0\\0&1&0\\1&1&0\end{pmatrix}
\]
where \(AG_1=a_1\row_1(G_1)+a_2\row_2(G_1)\) (Theorem~\ref{r1-thm-views}(d)). By Theorem~\ref{r1-thm-views}(c), \(AG_1A\) has
the rows \(\row_1(A)\), \(\row_2(A)\) and \(\row_1(A)+\row_2(A)=\row_3(A)\): \(AG_1A=A\).

(c) The submatrix in rows \(i_1=2\), \(i_2=3\) and columns \(j_1=1\), \(j_2=3\) is
\(A_*=\left(\begin{smallmatrix}0&1\\1&2\end{smallmatrix}\right)\), with \(\det A_*=-1\) and
\(A_*^{-1}=\left(\begin{smallmatrix}-2&1\\1&0\end{smallmatrix}\right)\) (Theorem~\ref{r5d-thm-adjugate}). By
Theorem~\ref{r10-thm-ginv-submatrix}(b), the \((s,t)\) entry of \(A_*^{-1}\) goes to position \((j_s,i_t)\) of \(G_2\): the
entries \(-2\), \(1\), \(1\), \(0\) go to positions \((1,2)\), \((1,3)\), \((3,2)\), \((3,3)\). So
\[
G_2=\begin{pmatrix}0&-2&1\\0&0&0\\0&1&0\end{pmatrix}
\]
and, by Theorem~\ref{r1-thm-views}(d),
\[
AG_2=a_1\row_1(G_2)+a_3\row_3(G_2)=\begin{pmatrix}0&-2&1\\0&0&0\\0&-2&1\end{pmatrix}+\begin{pmatrix}0&1&0\\0&1&0\\0&2&0\end{pmatrix}
=\begin{pmatrix}0&-1&1\\0&1&0\\0&0&1\end{pmatrix}
\]
\(AG_2A\) has the rows \(-\row_2(A)+\row_3(A)=(1,2,1)=\row_1(A)\), \(\row_2(A)\) and \(\row_3(A)\): \(AG_2A=A\). (\(G_2\ne G_1\):
a g-inverse is not unique.)
\end{solution}

\begin{exercise}\label{r10-ex-ginverse-4x4}
Let
\[
A=\begin{pmatrix}2&2&0&1\\0&-2&2&1\\1&0&1&1\\2&4&-2&0\end{pmatrix},\qquad
p_1=\begin{pmatrix}1\\-1\\0\\2\end{pmatrix},\quad p_2=\begin{pmatrix}1\\1\\1\\0\end{pmatrix},\quad
q_1=\begin{pmatrix}1\\2\\-1\\0\end{pmatrix},\quad q_2=\begin{pmatrix}1\\0\\1\\1\end{pmatrix}
\]
\begin{enumerate}
\item Show that \(A=p_1q_1'+p_2q_2'\), that this is a full-rank factorization \(A=BC\) with \(B=[p_1\ p_2]\) and the rows
\(q_1'\), \(q_2'\) of \(C\), and that \(\rank A=2\).
\item Find the g-inverses \(G_1\) and \(G_2\) of \(A\) of Theorem~\ref{r10-thm-ginv-submatrix}(b) built from the leading \(2\times2\)
submatrix and from the submatrix in rows \(1,3\) and columns \(3,4\).
\item Use \(G_2\) to find the condition on \(y\in\R^4\) for \(Ax=y\) to be consistent, and find all solutions of \(Ax=y\) for
\(y=(2,0,1,2)'\).
\end{enumerate}
\end{exercise}
\begin{solution}
(a) By the outer product expansion (Theorem~\ref{r1-thm-views}(d)), \(p_1q_1'+p_2q_2'=BC\), and
\[
p_1q_1'+p_2q_2'=\begin{pmatrix}1&2&-1&0\\-1&-2&1&0\\0&0&0&0\\2&4&-2&0\end{pmatrix}
+\begin{pmatrix}1&0&1&1\\1&0&1&1\\1&0&1&1\\0&0&0&0\end{pmatrix}=A
\]
Since \(p_1'p_2=0\) and \(q_1'q_2=0\), the nonzero vectors \(p_1,p_2\) are linearly independent, and so are \(q_1,q_2\)
(Proposition~\ref{r6-prop-orthogonal-list}(b)): \(B\) has full column rank and \(C\) full row rank
(Theorem~\ref{thm:r4-rank-rules}(b),(c)). By Theorem~\ref{thm:r4-rank-product}(b), \(\rank A=\rank C=2\), and \(A=BC\) is a
full-rank factorization.

(b)
\begin{itemize}
\item The leading submatrix \(M_1=\left(\begin{smallmatrix}2&2\\0&-2\end{smallmatrix}\right)\) has
\(M_1^{-1}=\frac14\left(\begin{smallmatrix}2&2\\0&-2\end{smallmatrix}\right)\) (the product is \(I_2\)). By
Theorem~\ref{r10-thm-ginv-submatrix}(b),
\[
G_1=\frac12\begin{pmatrix}1&1&0&0\\0&-1&0&0\\0&0&0&0\\0&0&0&0\end{pmatrix}
\]
\item The submatrix in rows \(i_1=1\), \(i_2=3\) and columns \(j_1=3\), \(j_2=4\) is \(M_2=\left(\begin{smallmatrix}0&1\\1&1\end{smallmatrix}\right)\),
with \(M_2^{-1}=\left(\begin{smallmatrix}-1&1\\1&0\end{smallmatrix}\right)\) (the product is \(I_2\)). By
Theorem~\ref{r10-thm-ginv-submatrix}(b), the \((s,t)\) entry of \(M_2^{-1}\) goes to position \((j_s,i_t)\): the entries \(-1\),
\(1\), \(1\), \(0\) go to the positions \((3,1)\), \((3,3)\), \((4,1)\), \((4,3)\):
\[
G_2=\begin{pmatrix}0&0&0&0\\0&0&0&0\\-1&0&1&0\\1&0&0&0\end{pmatrix}
\]
\end{itemize}
(c)
\begin{itemize}
\item With the columns \(a_1,\dots,a_4\) of \(A\), \(AG_2=[-a_3+a_4\ \ 0\ \ a_3\ \ 0]\) (Theorem~\ref{r1-thm-views}(a),(b)), so
\(AG_2y=y_1(1,-1,0,2)'+y_3(0,2,1,-2)'=(y_1,\ 2y_3-y_1,\ y_3,\ 2y_1-2y_3)'\). By Theorem~\ref{r10-thm-ginv-solves}(b),
\(Ax=y\) is consistent if and only if \(y_2=2y_3-y_1\) and \(y_4=2y_1-2y_3\).
\item For \(y=(2,0,1,2)'\) both conditions hold. The rows of \(G_2A\) are \(0\), \(0\), \(\row_3(A)-\row_1(A)=(-1,-2,1,0)\) and
\(\row_1(A)=(2,2,0,1)\) (Theorem~\ref{r1-thm-views}(c)), so
\[
G_2y=\begin{pmatrix}0\\0\\-1\\2\end{pmatrix},\qquad
I_4-G_2A=\begin{pmatrix}1&0&0&0\\0&1&0&0\\1&2&0&0\\-2&-2&0&0\end{pmatrix}
\]
\item By Theorem~\ref{r10-thm-general-solution}, \(\mathcal N(A)=\mathcal C(I_4-G_2A)=\Span\{(1,0,1,-2)',(0,1,2,-2)'\}\), and the
solutions are
\[
x=\begin{pmatrix}0\\0\\-1\\2\end{pmatrix}+z_1\begin{pmatrix}1\\0\\1\\-2\end{pmatrix}+z_2\begin{pmatrix}0\\1\\2\\-2\end{pmatrix}\qquad(z_1,z_2\in\R)
\]
\item The solution \(G_1y=(1,0,0,0)'=e_1\) (\(z_1=1\), \(z_2=0\)) is visible directly: \(y\) is the first column of \(A\).
\end{itemize}
\end{solution}

\needspace{16\baselineskip}
\begin{exercise}\label{r10-dr-reflexive}
Let
\[
A=\begin{pmatrix}1&2&1\\0&1&1\\1&3&2\end{pmatrix},\qquad G_1=\begin{pmatrix}1&-2&0\\0&1&0\\0&0&0\end{pmatrix}
\]
where \(\rank A=2\) and \(G_1\) is a g-inverse of \(A\) (Exercise~\ref{r10-dr-submatrix}).
\begin{enumerate}
\item Compute \(G_1A\), check that it is idempotent, and find \(\tr(G_1A)\) and \(\rank(I_3-G_1A)\).
\item Is \(G_1\) reflexive?
\item Let \(U\in\R^{3\times3}\) have \((3,3)\) entry \(1\) and all other entries \(0\). Show that \(G_3=G_1+(I_3-G_1A)U\) is an
invertible g-inverse of \(A\), and find the reflexive g-inverse \(G_3AG_3\).
\end{enumerate}
\end{exercise}
\begin{solution}
(a) By Theorem~\ref{r1-thm-views}(c), the rows of \(G_1A\) are \(\row_1(A)-2\row_2(A)=(1,0,-1)\), \(\row_2(A)=(0,1,1)\) and
\(0\):
\[
G_1A=\begin{pmatrix}1&0&-1\\0&1&1\\0&0&0\end{pmatrix},\qquad
I_3-G_1A=\begin{pmatrix}0&0&1\\0&0&-1\\0&0&1\end{pmatrix}
\]
The rows of \((G_1A)(G_1A)\) are \(\row_1(G_1A)-\row_3(G_1A)=(1,0,-1)\), \(\row_2(G_1A)+\row_3(G_1A)=(0,1,1)\) and \(0\): so
\(G_1A\) is idempotent. \(\tr(G_1A)=1+1+0=2=\rank A\), and \(I_3-G_1A\) has rank \(1=3-2\) (its columns are multiples of
\((1,-1,1)'\)), as Theorem~\ref{r10-thm-rank-idempotent}(c) says.

(b) The rows \((1,-2,0)\) and \((0,1,0)\) of \(G_1\) are linearly independent, so \(\rank G_1=2=\rank A\), and \(G_1\) is
reflexive (Theorem~\ref{r10-thm-reflexive}(b)). Check: \(G_1AG_1=(G_1A)G_1\) has the rows
\(\row_1(G_1)-\row_3(G_1)=(1,-2,0)\), \(\row_2(G_1)+\row_3(G_1)=(0,1,0)\) and \(0\), so \(G_1AG_1=G_1\).

(c)
\begin{itemize}
\item \((I_3-G_1A)U\) has column \(3\) equal to column \(3\) of \(I_3-G_1A\), that is, \((1,-1,1)'\), and its other columns
\(0\) (Theorem~\ref{r1-thm-views}(a)). So
\[
G_3=\begin{pmatrix}1&-2&1\\0&1&-1\\0&0&1\end{pmatrix}
\]
\item \(A(I_3-G_1A)=A-AG_1A=0\), so \(AG_3=AG_1\) and \(AG_3A=AG_1A=A\): \(G_3\) is a g-inverse of \(A\). It is upper
triangular with diagonal \(1,1,1\), so \(\det G_3=1\) (Theorem~\ref{r5d-thm-basic}(c)) and \(G_3\) is invertible
(Theorem~\ref{r5d-thm-product}(a)); \(\rank G_3=3>\rank A\), so \(G_3\) is not reflexive
(Theorem~\ref{r10-thm-reflexive}(b)).
\item By Theorem~\ref{r1-thm-views}(d),
\[
AG_1=a_1\row_1(G_1)+a_2\row_2(G_1)=\begin{pmatrix}1&0&0\\0&1&0\\1&1&0\end{pmatrix}
\]
and \(G_3AG_3=G_3(AG_1)\) has the columns \(G_3(1,0,1)'=(1+1,\ 0-1,\ 0+1)'\),
\(G_3(0,1,1)'=(-2+1,\ 1-1,\ 0+1)'\) and \(0\) (Theorem~\ref{r1-thm-views}(a),(b)):
\[
G_3AG_3=\begin{pmatrix}2&-1&0\\-1&0&0\\1&1&0\end{pmatrix}
\]
a reflexive g-inverse of \(A\) (Theorem~\ref{r10-thm-reflexive}(a)), of rank \(2\).
\end{itemize}
\end{solution}

\begin{exercise}\label{r10-ex-ginverse-ata}
Let
\[
A=\begin{pmatrix}1&1&0\\1&-1&2\\1&2&-1\\1&-2&3\end{pmatrix}
\]
\begin{enumerate}
\item Compute \(A'A\), show that \(\rank A=2\), and find the g-inverse \(G\) of \(A'A\) of Theorem~\ref{r10-thm-ginv-submatrix}(b)
built from its leading \(2\times2\) submatrix.
\item Check that \(AGA'A=A\), and write down the g-inverse \(GA'\) of \(A\).
\item Use \(GA'\) to decide whether \(Ax=y\) is consistent for \(y=(2,2,2,2)'\) and for \(y=e_1\). In the consistent case find
all solutions and the shortest one.
\end{enumerate}
\end{exercise}
\begin{solution}
(a) The columns are \(a_1=(1,1,1,1)'\), \(a_2=(1,-1,2,-2)'\) and \(a_3=(0,2,-1,3)'=a_1-a_2\), and, with the entries
\(a_i'a_j\) (Theorem~\ref{r1-thm-views}(d)),
\[
A'A=\begin{pmatrix}4&0&4\\0&10&-10\\4&-10&14\end{pmatrix}
\]
\(a_1\) and \(a_2\) are nonzero and orthogonal, hence linearly independent (Proposition~\ref{r6-prop-orthogonal-list}(b)), and
\(a_3\in\Span\{a_1,a_2\}\); so \(\rank A=2\), and \(\rank(A'A)=2\) (Theorem~\ref{thm:r4-rank-ata}(b)). The leading submatrix
\(\diag(4,10)\) is invertible, so \(G=\diag(\frac14,\frac1{10},0)\) is a g-inverse of \(A'A\)
(Theorem~\ref{r10-thm-ginv-submatrix}(b)).

(b) Multiplying the rows of \(A'A\) by \(\frac14,\frac1{10},0\) (Theorem~\ref{r1-thm-views}(c)),
\[
GA'A=\begin{pmatrix}1&0&1\\0&1&-1\\0&0&0\end{pmatrix}
\]
so the columns of \(A(GA'A)\) are \(a_1\), \(a_2\) and \(a_1-a_2=a_3\) (Theorem~\ref{r1-thm-views}(a),(b)): \(AGA'A=A\), as
Theorem~\ref{r10-thm-ata}(a) says. By Theorem~\ref{r10-thm-ata}(a),
\[
GA'=\begin{pmatrix}\frac14&\frac14&\frac14&\frac14\\[2pt]\frac1{10}&-\frac1{10}&\frac15&-\frac15\\[2pt]0&0&0&0\end{pmatrix}
\]
is a g-inverse of \(A\).

(c)
\begin{itemize}
\item \(A(GA')y=\frac14(a_1'y)a_1+\frac1{10}(a_2'y)a_2\) (Theorem~\ref{r1-thm-views}(d)). For \(y=(2,2,2,2)'=2a_1\), \(a_1'y=8\) and
\(a_2'y=0\), so \(AGA'y=2a_1=y\): consistent. For \(y=e_1\), \(AGA'e_1=\frac14a_1+\frac1{10}a_2=\frac1{20}(7,3,9,1)'\ne e_1\):
inconsistent (Theorem~\ref{r10-thm-ginv-solves}(b)).
\item For \(y=2a_1\), \(GA'y=(2,0,0)'\), and
\[
I_3-GA'A=\begin{pmatrix}0&0&-1\\0&0&1\\0&0&1\end{pmatrix}
\]
so the solutions are \(x=(2,0,0)'+z(-1,1,1)'\), \(z\in\R\) (Theorem~\ref{r10-thm-general-solution}).
\item As \(Ax=y\) is consistent, its solutions are the minimizers of \(\lVert Ax-y\rVert\), and the shortest one is the only one in
\(\mathcal C(A')=\mathcal N(A)^\perp\) (Theorem~\ref{r10-thm-mp-shortest}(b); Corollary~\ref{r6-cor-four-subspaces} and Theorem~\ref{r6-thm-perp-perp}(a)), where
\(\mathcal N(A)=\Span\{(-1,1,1)'\}\). The condition \((2-z)(-1)+z+z=0\) gives \(z=\frac23\) and \(x^+=\frac13(4,2,2)'\), with
\(\lVert x^+\rVert^2=\frac{24}9=\frac83\); the solution \((2,0,0)'\) has \(\lVert(2,0,0)'\rVert^2=4=\frac83+\frac43\).
\end{itemize}
\end{solution}

\begin{exercise}\label{r10-ex-mp-drill}
Let
\[
A=\begin{pmatrix}1&0&1&1\\0&1&1&2\\1&1&2&3\\2&1&3&4\end{pmatrix}
\]
\begin{enumerate}
\item Find \(A^+\) from a full-rank factorization, and check that \(AA^+\) and \(A^+A\) are symmetric.
\item Find the shortest solution of \(Ax=(0,3,3,3)'\), and check it against all solutions.
\end{enumerate}
\end{exercise}
\begin{solution}
(a) Row-reduce \(A\):
\[
\begin{pmatrix}1&0&1&1\\0&1&1&2\\1&1&2&3\\2&1&3&4\end{pmatrix}
\xrightarrow[R_4\leftarrow R_4-2R_1]{R_3\leftarrow R_3-R_1}
\begin{pmatrix}1&0&1&1\\0&1&1&2\\0&1&1&2\\0&1&1&2\end{pmatrix}
\xrightarrow[R_4\leftarrow R_4-R_2]{R_3\leftarrow R_3-R_2}
\begin{pmatrix}1&0&1&1\\0&1&1&2\\0&0&0&0\\0&0&0&0\end{pmatrix}
\]
The last matrix is the RREF; the pivot columns are \(1,2\), so \(A=BC\) with (Theorem~\ref{thm:r4-full-rank-factorization}(a))
\[
B=\begin{pmatrix}1&0\\0&1\\1&1\\2&1\end{pmatrix},\qquad C=\begin{pmatrix}1&0&1&1\\0&1&1&2\end{pmatrix}
\]
By Theorem~\ref{r5d-thm-adjugate},
\[
B'B=\begin{pmatrix}6&3\\3&3\end{pmatrix},\quad (B'B)^{-1}=\frac13\begin{pmatrix}1&-1\\-1&2\end{pmatrix},\qquad
CC'=\begin{pmatrix}3&3\\3&6\end{pmatrix},\quad (CC')^{-1}=\frac13\begin{pmatrix}2&-1\\-1&1\end{pmatrix}
\]
(both determinants are \(9\)). Then
\[
L=(B'B)^{-1}B'=\frac13\begin{pmatrix}1&-1&0&1\\-1&2&1&0\end{pmatrix},\qquad
K=C'(CC')^{-1}=\frac13\begin{pmatrix}2&-1\\-1&1\\1&0\\0&1\end{pmatrix}
\]
and by Theorem~\ref{r10-thm-mp-exists}(b)
\[
A^+=KL=\frac19\begin{pmatrix}2&-1\\-1&1\\1&0\\0&1\end{pmatrix}\begin{pmatrix}1&-1&0&1\\-1&2&1&0\end{pmatrix}
=\frac19\begin{pmatrix}3&-4&-1&2\\-2&3&1&-1\\1&-1&0&1\\-1&2&1&0\end{pmatrix}
\]
Since \(CK=I_2\) and \(LB=I_2\),
\[
AA^+=BCKL=BL=\frac13\begin{pmatrix}1&-1&0&1\\-1&2&1&0\\0&1&1&1\\1&0&1&2\end{pmatrix},\qquad
A^+A=KLBC=KC=\frac13\begin{pmatrix}2&-1&1&0\\-1&1&0&1\\1&0&1&1\\0&1&1&2\end{pmatrix}
\]
both symmetric.

(b)
\begin{itemize}
\item Let \(y=(0,3,3,3)'\). \(A^+\) is a g-inverse of \(A\), and \(AA^+y=\frac13(0-3+0+3,\ 0+6+3+0,\ 0+3+3+3,\ 0+0+3+6)'=(0,3,3,3)'=y\),
so \(Ax=y\) is consistent (Theorem~\ref{r10-thm-ginv-solves}(b)).
\item By (P1) and (P4), \(A^+\) is a minimum-norm g-inverse of \(A\), so the shortest solution is
(Theorem~\ref{r10-thm-min-norm})
\[
A^+y=\frac19\begin{pmatrix}0-12-3+6\\0+9+3-3\\0-3+0+3\\0+6+3+0\end{pmatrix}=\begin{pmatrix}-1\\1\\0\\1\end{pmatrix}
\]
check: \(A(-1,1,0,1)'=-a_1+a_2+a_4=(0,3,3,3)'\).
\item By Theorem~\ref{r10-thm-general-solution}, the solutions are \(A^+y+(I_4-A^+A)z\), where
\[
I_4-A^+A=\frac13\begin{pmatrix}1&1&-1&0\\1&2&0&-1\\-1&0&2&-1\\0&-1&-1&1\end{pmatrix}
\]
has the columns \(-\frac13s_1\), \(-\frac13s_2\), \(\frac13(2s_1-s_2)\) and \(\frac13(s_2-s_1)\), with the special solutions
\(s_1=(-1,-1,1,0)'\) and \(s_2=(-1,-2,0,1)'\) of the RREF (Theorem~\ref{thm:r4-four-bases}(c)). So the solutions are
\(x=(-1-t_1-t_2,\ 1-t_1-2t_2,\ t_1,\ 1+t_2)'\) (\(t_1,t_2\in\R\)), with
\[
\lVert x\rVert^2=(1+t_1+t_2)^2+(1-t_1-2t_2)^2+t_1^2+(1+t_2)^2=3+3(t_1+t_2)^2+3t_2^2
\]
smallest exactly for \(t_1=t_2=0\).
\end{itemize}
\end{solution}

\begin{exercise}\label{r10-ex-min-norm}
Let
\[
A=\begin{pmatrix}1&1&1&1&1\\1&-1&1&-1&1\\1&0&1&0&1\end{pmatrix},\qquad
f_1=\begin{pmatrix}1\\1\\1\end{pmatrix},\ f_2=\begin{pmatrix}1\\-1\\0\end{pmatrix},\
g_1=\begin{pmatrix}1\\0\\1\\0\\1\end{pmatrix},\ g_2=\begin{pmatrix}0\\1\\0\\1\\0\end{pmatrix}
\]
\begin{enumerate}
\item Show that \(A=f_1g_1'+f_2g_2'\), that this is a full-rank factorization \(A=BC\) with \(B=[f_1\ f_2]\) and the rows \(g_1'\),
\(g_2'\) of \(C\), and that \(B'B=CC'=\diag(3,2)\).
\item Find \(A^+\) and \(AA^+\).
\item Find the condition on \(y\in\R^3\) for \(Ax=y\) to be consistent. For \(y=(8,4,6)'\) find the shortest solution, and
compare it with the solution \((6,2,0,0,0)'\).
\end{enumerate}
\end{exercise}
\begin{solution}
(a) \(f_1g_1'\) has the three rows \((1,0,1,0,1)\), and \(f_2g_2'\) the rows \((0,1,0,1,0)\), \((0,-1,0,-1,0)\) and \(0\); their sum
is \(A\), and it is \(BC\) by the outer product expansion (Theorem~\ref{r1-thm-views}(d)). The nonzero vectors \(f_1\perp f_2\)
are linearly independent, and so are \(g_1\perp g_2\) (Proposition~\ref{r6-prop-orthogonal-list}(b)); so \(B\) has full column
rank, \(C\) has full row rank (Theorem~\ref{thm:r4-rank-rules}(b),(c)), and \(\rank A=\rank C=2\)
(Theorem~\ref{thm:r4-rank-product}(b)). With the entries \(f_i'f_j\) and \(g_i'g_j\) (Theorem~\ref{r1-thm-views}(d)),
\(B'B=\diag(3,2)\) and \(CC'=\diag(3,2)\).

(b) By Theorem~\ref{r10-thm-mp-exists}(b) and the outer product expansion,
\[
A^+=C'\diag(\tfrac13,\tfrac12)\diag(\tfrac13,\tfrac12)B'=\frac19g_1f_1'+\frac14g_2f_2'
=\frac1{36}\begin{pmatrix}4&4&4\\9&-9&0\\4&4&4\\9&-9&0\\4&4&4\end{pmatrix}
\]
and, by the proof of Theorem~\ref{r10-thm-mp-exists}(b),
\[
AA^+=B(B'B)^{-1}B'=\frac13f_1f_1'+\frac12f_2f_2'=\frac16\begin{pmatrix}5&-1&2\\-1&5&2\\2&2&2\end{pmatrix}
\]
(c)
\begin{itemize}
\item \(A^+\) is a g-inverse of \(A\), so \(Ax=y\) is consistent if and only if \(AA^+y=y\) (Theorem~\ref{r10-thm-ginv-solves}(b)),
that is, \((I_3-AA^+)y=0\). With \(w=(1,1,-2)'\),
\[
I_3-AA^+=\frac16\begin{pmatrix}1&1&-2\\1&1&-2\\-2&-2&4\end{pmatrix}=\frac16ww'
\]
so \((I_3-AA^+)y=\frac16(w'y)w\), and \(Ax=y\) is consistent if and only if \(y_1+y_2-2y_3=0\).
\item For \(y=(8,4,6)'\), \(8+4-12=0\): consistent. By (P1) and (P4), \(A^+\) is a minimum-norm g-inverse of \(A\), so the
shortest solution is (Theorem~\ref{r10-thm-min-norm}) \(x^+=A^+y=\frac1{36}(72,36,72,36,72)'=(2,1,2,1,2)'\), with
\(\lVert x^+\rVert^2=14\).
\item The solution \(x=(6,2,0,0,0)'\) (\(Ax=6a_1+2a_2=(8,4,6)'\)) has \(\lVert x\rVert^2=40=14+\lVert(4,1,-2,-1,-2)'\rVert^2=14+26\).
\end{itemize}
\end{solution}

\begin{exercise}\label{r10-ex-mp-4x4}
For the matrix \(A=p_1q_1'+p_2q_2'=BC\) of Exercise~\ref{r10-ex-ginverse-4x4}, find \(A^+\), \(AA^+\) and \(A^+A\), and the shortest
solution of \(Ax=(2,0,1,2)'\).
\end{exercise}
\begin{solution}
\begin{itemize}
\item \(p_1'p_1=q_1'q_1=6\), \(p_2'p_2=q_2'q_2=3\) and \(p_1'p_2=q_1'q_2=0\), so \(B'B=CC'=\diag(6,3)\)
(Theorem~\ref{r1-thm-views}(d)). By Theorem~\ref{r10-thm-mp-exists}(b) and the outer product expansion,
\[
A^+=C'\diag(\tfrac16,\tfrac13)\diag(\tfrac16,\tfrac13)B'=\frac1{36}q_1p_1'+\frac19q_2p_2'
=\frac1{36}\begin{pmatrix}5&3&4&2\\2&-2&0&4\\3&5&4&-2\\4&4&4&0\end{pmatrix}
\]
\item By the proof of Theorem~\ref{r10-thm-mp-exists}(b), \(AA^+=B(B'B)^{-1}B'\) and \(A^+A=C'(CC')^{-1}C\):
\[
AA^+=\frac16p_1p_1'+\frac13p_2p_2'=\frac16\begin{pmatrix}3&1&2&2\\1&3&2&-2\\2&2&2&0\\2&-2&0&4\end{pmatrix},\qquad
A^+A=\frac16q_1q_1'+\frac13q_2q_2'=\frac16\begin{pmatrix}3&2&1&2\\2&4&-2&0\\1&-2&3&2\\2&0&2&2\end{pmatrix}
\]
They are the projection matrices onto \(\mathcal C(A)\) and \(\mathcal C(A')\) (Proposition~\ref{r10-prop-mp-props}(b)); both have
trace \(2=\rank A\) (Theorem~\ref{r10-thm-rank-idempotent}(c)).
\item The system is consistent (Exercise~\ref{r10-ex-ginverse-4x4}). By (P1) and (P4), \(A^+\) is a minimum-norm g-inverse of
\(A\), so by Theorem~\ref{r10-thm-min-norm} the shortest solution is
\(x^+=A^+(2,0,1,2)'=\frac1{36}(18,12,6,12)'=\frac16(3,2,1,2)'\), with \(\lVert x^+\rVert^2=\frac{18}{36}=\frac12\). It is the first
column of \(A^+A\): \(A^+(2,0,1,2)'=A^+Ae_1\), the projection onto \(\mathcal C(A')\) of the solution \(e_1\).
\item For the solutions \(e_1\) and \((0,0,-1,2)'\) of Exercise~\ref{r10-ex-ginverse-4x4},
\(\lVert e_1-x^+\rVert^2=\frac{9+4+1+4}{36}=\frac12\) and \(\lVert(0,0,-1,2)'-x^+\rVert^2=\frac{9+4+49+100}{36}=\frac92\), and indeed
\(1=\frac12+\frac12\) and \(5=\frac12+\frac92\).
\end{itemize}
\end{solution}

\begin{exercise}\label{r10-ex-mendel}
In an experiment of Mendel's type, \(N=160\) pea seeds fall into the four categories round yellow, round green, wrinkled yellow
and wrinkled green, with the counts \(c=(93,27,31,9)'\). Mendel's theory gives the probabilities \(p=\frac1{16}(9,3,3,1)'\). Let
\(D\), \(V\), \(u\) and \(X^2\) be as in Example~\ref{r10-exm-pearson}.
\begin{enumerate}
\item Compute the expected counts \(Np_j\), the contributions \((c_j-Np_j)^2/(Np_j)\) and \(X^2\).
\item Under the approximation of Example~\ref{r10-exm-pearson}(e), the distribution of \(X^2\) is approximated by \(\chi^2_3\). Show
that \(\prb(Y>\frac8{15})>\frac12\) for \(Y\sim\chi^2_3\).
\item Compute \(V\), its leading \(3\times3\) block \(V_{11}\) and \(V_{11}^{-1}\), and check that \(u'Gu=X^2\) for the g-inverse
\(G=\left(\begin{smallmatrix}V_{11}^{-1}&0\\0&0\end{smallmatrix}\right)\) of \(V\).
\end{enumerate}
\end{exercise}
\begin{solution}
(a) The expected counts are \(160\cdot\frac9{16}=90\), \(160\cdot\frac3{16}=30\) (twice) and \(160\cdot\frac1{16}=10\), and
\(c-Np=(3,-3,1,-1)'\). The contributions are
\[
\frac9{90}=\frac1{10},\qquad\frac9{30}=\frac3{10},\qquad\frac1{30},\qquad\frac1{10}
\]
and \(X^2=\frac{3+9+1+3}{30}=\frac8{15}\).

(b) Let \(z_1,z_2,z_3\) be independent \(N(0,1)\) variables. By Definition~\ref{r15-def-chisq}(a),
\(v_3=z_1^2+z_2^2+z_3^2\sim\chi^2_3\) and \(v_2=z_1^2+z_2^2\sim\chi^2_2\), and \(v_3\geqslant v_2\). So the event \(\{v_2>c\}\) is
contained in \(\{v_3>c\}\), and \(\prb(Y>c)=\prb(v_3>c)\geqslant\prb(v_2>c)=e^{-c/2}\) for \(c>0\) (Example~\ref{r15-exm-chisq}(a)).
For \(c=\frac8{15}\), \(e^{-4/15}>\frac12\), since \(e^4<3^4=81<2^{15}\) (\(e<3\), Example~\ref{r11-exm-confidence-ellipse}). So under the
approximation a value of \(X^2\) larger than the one computed in (a) has probability more than \(\frac12\): the counts agree well with Mendel's
\(9:3:3:1\) law.

(c)
\begin{itemize}
\item \(256D=\diag(144,48,48,16)\) and \(256pp'=(9,3,3,1)'(9,3,3,1)\), so
\[
V=D-pp'=\frac1{256}\begin{pmatrix}63&-27&-27&-9\\-27&39&-9&-3\\-27&-9&39&-3\\-9&-3&-3&15\end{pmatrix},\qquad
V_{11}=\frac1{256}\begin{pmatrix}63&-27&-27\\-27&39&-9\\-27&-9&39\end{pmatrix}
\]
\item With \(D_3=\diag(p_1,p_2,p_3)\) and \(p_4=\frac1{16}\), the matrix \(D_3^{-1}+\frac1{p_4}\one_3\one_3'=\diag(\frac{16}9,\frac{16}3,\frac{16}3)+16J_3\)
is
\[
V_{11}^{-1}=\frac{16}9\begin{pmatrix}10&9&9\\9&12&9\\9&9&12\end{pmatrix}
\]
Check:
\[
\begin{pmatrix}63&-27&-27\\-27&39&-9\\-27&-9&39\end{pmatrix}\begin{pmatrix}10&9&9\\9&12&9\\9&9&12\end{pmatrix}=144I_3
\]
and \(\frac1{256}\cdot\frac{16}9\cdot144=1\); so \(V_{11}\) is invertible with this inverse (Theorem~\ref{r2-thm-invertible}).
\item \(\rank V=3\) (Example~\ref{r10-exm-pearson}(b)), so \(G\) is a g-inverse of \(V\) (Theorem~\ref{r10-thm-ginv-submatrix}(b)).
\item \(u=(c-Np)/\sqrt{160}\), so with \(u_{(3)}=(u_1,u_2,u_3)'=(3,-3,1)'/\sqrt{160}\),
\[
u'Gu=u_{(3)}'V_{11}^{-1}u_{(3)}=\frac1{160}\cdot\frac{16}9\begin{pmatrix}3&-3&1\end{pmatrix}\begin{pmatrix}12\\0\\12\end{pmatrix}
=\frac1{160}\cdot\frac{16}9\cdot48=\frac8{15}=X^2
\]
where \((12,0,12)'\) is the product of the matrix of \(V_{11}^{-1}\) (without the factor \(\frac{16}9\)) with \((3,-3,1)'\). In the
form \(D_3^{-1}+16J_3\): \(u_{(3)}'D_3^{-1}u_{(3)}+16(u_1+u_2+u_3)^2=\frac1{160}\bigl(16+48+\frac{16}3+16\bigr)=\frac8{15}\), as
\(u_1+u_2+u_3=-u_4=\frac1{\sqrt{160}}\).
\end{itemize}
\end{solution}

\needspace{14\baselineskip}
\subsection*{Problems}

\begin{problem}[\citeat{bapat}{Ch.~4, Exercise~3}, \citeat{harville}{Theorem~9.2.7}]\label{r10-ex-all-ginverses}
Let \(G\) be a g-inverse of \(A\in\R^{m\times n}\). Show that the g-inverses of \(A\) are exactly the matrices
\[
G+(I_n-GA)U+V(I_m-AG)\qquad(U,V\in\R^{n\times m})
\]
\end{problem}
\begin{solution}
\begin{itemize}
\item Each such matrix is a g-inverse: since \(AGA=A\),
\(A\bigl(G+(I-GA)U+V(I-AG)\bigr)A=AGA+(A-AGA)UA+AV(A-AGA)=A\).
\item Conversely, let \(H\) be a g-inverse of \(A\), and put \(U=HAG\), \(V=H-G\). Then
\[
\begin{aligned}
G+(I-GA)HAG+(H-G)(I-AG)&=G+HAG-GAHAG+H-HAG-G+GAG\\
&=H-G(AHA)G+GAG=H-GAG+GAG=H
\end{aligned}
\]
\end{itemize}
\end{solution}

\begin{problem}[\citeat{harville-solutions}{Ch.~14, Exercise~44}]\label{r10-ex-rank-cab}
Let \(A\in\R^{q\times p}\), \(B\in\R^{p\times n}\) and \(C\in\R^{m\times q}\), and let \(G\) be a g-inverse of \(CAB\). Show that
\begin{enumerate}
\item \(CABGC=C\) if and only if \(\rank(CAB)=\rank C\);
\item \(BGCAB=B\) if and only if \(\rank(CAB)=\rank B\).
\end{enumerate}
\end{problem}
\begin{solution}
(a)
\begin{itemize}
\item Let \(\rank(CAB)=\rank C\). As \(\mathcal C(CAB)\subseteq\mathcal C(C)\) (Theorem~\ref{thm:r4-rank-product}(a)) and the
two dimensions agree, \(\mathcal C(CAB)=\mathcal C(C)\) (Theorem~\ref{thm:r4-counting}(a)). So each column of \(C\) is
\(CABr_j\) for some \(r_j\in\R^n\), that is, \(C=CABR\) with \(R=[r_1\ \cdots\ r_q]\) (Theorem~\ref{r1-thm-views}(a)). Hence
\(CABGC=CABGCABR=CABR=C\).
\item Conversely, if \(CABGC=C\), then \(\rank C=\rank\bigl((CAB)(GC)\bigr)\leqslant\rank(CAB)\leqslant\rank C\)
(Theorem~\ref{thm:r4-rank-product}(a)), so \(\rank(CAB)=\rank C\).
\end{itemize}
(b)
\begin{itemize}
\item Let \(\rank(CAB)=\rank B\). Each row \(\row_i(CAB)=\row_i(CA)B\) lies in \(\mathcal R(B)\) (Theorem~\ref{r1-thm-views}(c)),
so \(\mathcal R(CAB)\subseteq\mathcal R(B)\); both have dimension \(\rank B\), so the spaces of the transposed vectors are
equal (Theorem~\ref{thm:r4-counting}(a)), and \(\mathcal R(CAB)=\mathcal R(B)\). So each row of \(B\) is \(l_i'CAB\) for some
\(l_i\in\R^m\), that is, \(B=LCAB\), where \(L\) has the rows \(l_1',\dots,l_p'\) (Theorem~\ref{r1-thm-views}(c)). Hence
\(BGCAB=LCABGCAB=LCAB=B\).
\item Conversely, if \(BGCAB=B\), then \(\rank B=\rank\bigl((BG)(CAB)\bigr)\leqslant\rank(CAB)\leqslant\rank B\), so
\(\rank(CAB)=\rank B\).
\end{itemize}
\end{solution}

\begin{problem}[\citeat{bapat}{Ch.~4, Exercise~12}]\label{r10-ex-min-norm-char}
Let \(A\in\R^{m\times n}\) and \(G\in\R^{n\times m}\). Show that \(G\) is a minimum-norm g-inverse of \(A\) if and only if
\(GAA'=A'\).
\end{problem}
\begin{solution}
\begin{itemize}
\item If \(AGA=A\) and \((GA)'=GA\), then \(GAA'=(GA)'A'=A'G'A'=(AGA)'=A'\).
\item Conversely, let \(GAA'=A'\). Then \(AGAA'=AA'\), that is, \((AG-I_m)AA'=0\). Transposing,
\(AA'(AG-I_m)'=0\): the columns of \((AG-I_m)'\) lie in \(\mathcal N(AA')=\mathcal N(A')\)
(Theorem~\ref{thm:r4-rank-ata}(a), applied to \(A'\)). So \(A'(AG-I_m)'=0\), and transposing again, \((AG-I_m)A=0\): \(AGA=A\).
\item Multiplying \(GAA'=A'\) by \(G'\) on the right, \(A'G'=GAA'G'=(GA)(GA)'\), which is symmetric. So \((GA)'=A'G'\) is
symmetric, and \(GA=((GA)')'=(GA)'\).
\end{itemize}
\end{solution}

\begin{problem}[\citeat{abadir-magnus}{Exercise~10.27}]\label{r10-ex-mp-idempotent}
\begin{enumerate}
\item Let \(A\in\R^{n\times n}\) be symmetric and idempotent. Show that \(A^+=A\).
\item Does the converse hold?
\item Show that \((AA^+)^+=AA^+\) and \((A^+A)^+=A^+A\) for every \(A\in\R^{m\times n}\).
\end{enumerate}
\end{problem}
\begin{solution}
(a) If \(A=A'=A^2\), then \(A^3=A\cdot A^2=A^2=A\) and \((A^2)'=A^2\). So \(G=A\) satisfies (P1)--(P4): \(AAA=A\),
\(AAA=A\), \((AA)'=AA\), \((AA)'=AA\). Hence \(A^+=A\) (Theorem~\ref{r10-thm-mp-unique}).

(b) No. For \(A=-I_n\), \(A^+=A^{-1}=-I_n=A\) (Theorem~\ref{r10-thm-mp-exists}(c)), but \(A^2=I_n\ne A\): \(A\) is not
idempotent.

(c) \(AA^+\) and \(A^+A\) are symmetric and idempotent (Proposition~\ref{r10-prop-mp-props}(b)). Apply (a).
\end{solution}

\begin{problem}[\citeat{harville-solutions}{Ch.~20, Exercise~2}]\label{r10-ex-mp-transpose}
Let \(A\in\R^{m\times n}\). Show that \(A^+=A'\) if and only if \(A'A\) is idempotent.
\end{problem}
\begin{solution}
\begin{itemize}
\item Let \(A'A\) be idempotent: \(A'A=A'AA'A\). Multiply on the left by \(A^+(A^+)'\) and on the right by \(A^+\). By
(P1)--(P4) and Theorem~\ref{r1-thm-transpose},
\[
\begin{aligned}
A^+(A^+)'A'AA^+&=A^+(AA^+)'AA^+=A^+AA^+AA^+=A^+AA^+=A^+\\
A^+(A^+)'A'AA'AA^+&=A^+(AA^+)'AA'(AA^+)'=A^+AA^+A(AA^+A)'\\
&=A^+AA'=(A^+A)'A'=(AA^+A)'=A'
\end{aligned}
\]
So \(A^+=A'\).
\item Conversely, let \(A^+=A'\). Then \(A'A=A^+A\), and \((A^+A)(A^+A)=A^+(AA^+A)=A^+A\): \(A'A\) is idempotent.
\end{itemize}
\end{solution}

\begin{problem}[\citeat{harville}{\S20.2}, \citeat{abadir-magnus}{Exercise~10.33(a)}]\label{r10-pr-mp-drill}
Verify the conditions (P1)--(P4) to show:
\begin{enumerate}
\item \(\diag(d_1,\dots,d_n)^+=\diag(d_1^+,\dots,d_n^+)\), where \(d^+=1/d\) for a number \(d\ne0\) and \(0^+=0\);
\item \((uv')^+=vu'/(\lVert u\rVert^2\lVert v\rVert^2)\) for nonzero \(u\in\R^m\) and \(v\in\R^n\); that is, \(A^+=A'/\tr(AA')\) for
\(A=uv'\).
\end{enumerate}
\end{problem}
\begin{solution}
(a) For \(D=\diag(d_1,\dots,d_n)\) and \(G=\diag(d_1^+,\dots,d_n^+)\), \(DG=GD=\diag(d_1d_1^+,\dots,d_nd_n^+)\) is symmetric, and
\(dd^+d=d\), \(d^+dd^+=d^+\) for every number \(d\) (both sides are \(0\) if \(d=0\)); so \(DGD=D\) and \(GDG=G\). Hence \(G=D^+\)
(Theorem~\ref{r10-thm-mp-unique}).

(b) Let \(A=uv'\), \(c=\lVert u\rVert^2\lVert v\rVert^2\) and \(G=vu'/c\).
\begin{itemize}
\item \(AG=uv'vu'/c=uu'/\lVert u\rVert^2\) and \(GA=vu'uv'/c=vv'/\lVert v\rVert^2\) are symmetric: (P3), (P4).
\item \(AGA=uu'uv'/\lVert u\rVert^2=uv'=A\) and \(GAG=vv'vu'/(\lVert v\rVert^2c)=vu'/c=G\): (P1), (P2). So \(G=A^+\)
(Theorem~\ref{r10-thm-mp-unique}).
\item \(A'=vu'\), and \(\tr(AA')=\tr(uv'vu')=\lVert v\rVert^2\tr(uu')=\lVert v\rVert^2u'u=c\) (Theorem~\ref{r1-thm-trace}(a),(b)); so
\(A^+=A'/\tr(AA')\).
\end{itemize}
\end{solution}

\begin{problem}[\citeat{abadir-magnus}{Exercise~10.29(c)--(e)}]\label{r10-pr-ata-ginverse}
Let \(A\in\R^{m\times n}\). Show that
\begin{enumerate}
\item \((A'A)^+=A^+(A^+)'\);
\item \((AA')^+=(A^+)'A^+\);
\item \(A^+=A'(AA')^+\).
\end{enumerate}
\end{problem}
\begin{solution}
First, by (P1)--(P4) and Theorem~\ref{r1-thm-transpose},
\[
A'AA^+=A'(AA^+)'=A'(A^+)'A'=(AA^+A)'=A',\qquad A^+(A^+)'A'=A^+(AA^+)'=A^+AA^+=A^+
\]
and likewise \(A'(A^+)'A^+=(A^+A)'A^+=A^+AA^+=A^+\).

(a) Let \(G=A^+(A^+)'\). We check (P1)--(P4) for the matrix \(A'A\), using the identities above.
\begin{itemize}
\item (P1): \(A'A\,G\,A'A=(A'AA^+)(A^+)'A'A=A'(A^+)'A'A=(AA^+A)'A=A'A\).
\item (P2): \(G\,A'A\,G=A^+(A^+)'(A'AA^+)(A^+)'=A^+(A^+)'A'(A^+)'=A^+(AA^+)'(A^+)'=A^+AA^+(A^+)'=A^+(A^+)'=G\).
\item (P3): \(A'AG=(A'AA^+)(A^+)'=A'(A^+)'=(A^+A)'=A^+A\), which is symmetric by (P4).
\item (P4): \(GA'A=(A^+(A^+)'A')A=A^+A\), which is symmetric.
\end{itemize}
So \((A'A)^+=G\) (Theorem~\ref{r10-thm-mp-unique}).

(b) By Proposition~\ref{r10-prop-mp-props}(c), \((A')^+=(A^+)'\). Part (a) for \(A'\) in place of \(A\), as \((A')'A'=AA'\), gives
\((AA')^+=(A')^+((A')^+)'=(A^+)'A^+\).

(c) By the identity \(A'(A^+)'A^+=A^+\) and (b), \(A^+=A'\bigl((A^+)'A^+\bigr)=A'(AA')^+\).
\end{solution}

\begin{problem}[\citeat{prasolov}{Problem~44.1}]\label{r10-ex-pr-two-sided}
Let \(A\in\R^{m\times n}\), \(B\in\R^{p\times q}\) and \(C\in\R^{m\times q}\), and let \(C=AX=YB\) for some
\(X\in\R^{n\times q}\) and \(Y\in\R^{m\times p}\). Show that \(C=AZB\) for some \(Z\in\R^{n\times p}\).
\end{problem}
\begin{solution}
Let \(G\) and \(H\) be g-inverses of \(A\) and \(B\) (Theorem~\ref{r10-thm-ginv-exists}).
\begin{itemize}
\item The system \(AX=C\) is consistent, so \(AGC=C\) (Lemma~\ref{r10-lem-range-test}); directly,
\(AGC=AGAX=AX=C\), as \(AGA=A\).
\item Likewise, as \(BHB=B\), \(CHB=YBHB=YB=C\).
\item Hence \(C=AGC=AG(CHB)=A(GCH)B\), that is, \(C=AZB\) with \(Z=GCH\in\R^{n\times p}\).
\end{itemize}
\end{solution}

\begin{problem}[\citeat{prasolov}{Problem~44.2}]\label{r10-ex-pr-homogeneous}
Let \(A\in\R^{m\times n}\) and \(B\in\R^{p\times q}\), and let \(G\) and \(H\) be g-inverses of \(A\) and \(B\). Show that a
matrix \(X\in\R^{n\times p}\) satisfies \(AX=0\) and \(XB=0\) if and only if
\[
X=(I_n-GA)Y(I_p-BH)\qquad\text{for some }Y\in\R^{n\times p}
\]
\end{problem}
\begin{solution}
\begin{itemize}
\item If \(X=(I_n-GA)Y(I_p-BH)\), then, as \(AGA=A\) and \(BHB=B\),
\(AX=(A-AGA)Y(I_p-BH)=0\) and \(XB=(I_n-GA)Y(B-BHB)=0\).
\item Conversely, let \(AX=0\) and \(XB=0\), and let \(r=\rank X\). If \(r=0\), then \(X=0\)
(Theorem~\ref{thm:r4-four-bases}), and \(Y=0\) will do. Let \(r\geqslant1\), and let \(X=PQ\) be a full-rank factorization,
with \(P\in\R^{n\times r}\) and \(Q\in\R^{r\times p}\) (Theorem~\ref{thm:r4-full-rank-factorization}(a)).
\item \(\mathcal C(P)=\mathcal C(X)\) and \(\mathcal N(Q)=\mathcal N(X)\) (Theorem~\ref{thm:r4-full-rank-factorization}(b)). So
each column of \(P\) is \(Xk_j\) for some \(k_j\in\R^p\), that is, \(P=XK\) with \(K=[k_1\ \cdots\ k_r]\)
(Theorem~\ref{r1-thm-views}(a)), and \(AP=AXK=0\). Also \(XB=0\) says that each column of \(B\) lies in
\(\mathcal N(X)=\mathcal N(Q)\) (Theorem~\ref{r1-thm-views}(a)), so \(QB=0\).
\item As \(AP=0\), each column of \(P\) lies in \(\mathcal N(A)=\mathcal C(I_n-GA)\) (Theorem~\ref{r10-thm-general-solution}),
so \(P=(I_n-GA)Y_1\) for some \(Y_1\in\R^{n\times r}\) (Theorem~\ref{r1-thm-views}(a)).
\item \(QB=0\) gives \(B'Q'=0\) (Theorem~\ref{r1-thm-transpose}). \(H'\) is a g-inverse of \(B'\)
(Proposition~\ref{r10-prop-ginv-basic}(a)), so each column of \(Q'\) lies in \(\mathcal N(B')=\mathcal C(I_p-H'B')\)
(Theorem~\ref{r10-thm-general-solution}), and \(Q'=(I_p-H'B')Z\) for some \(Z\in\R^{p\times r}\). Transposing,
\(Q=Z'(I_p-BH)=Y_2(I_p-BH)\) with \(Y_2=Z'\).
\item Hence \(X=PQ=(I_n-GA)Y_1Y_2(I_p-BH)\), which has the required form with \(Y=Y_1Y_2\in\R^{n\times p}\).
\end{itemize}
\end{solution}

\begin{problem}[\citeat{prasolov}{Problem~44.3}]\label{r10-ex-pr-common-solution}
Let \(A\in\R^{m\times n}\), \(B\in\R^{p\times q}\), \(C\in\R^{m\times p}\) and \(D\in\R^{n\times q}\), and let \(G\) and \(H\) be
g-inverses of \(A\) and \(B\). Show that the equations \(AX=C\) and \(XB=D\) have a common solution \(X\in\R^{n\times p}\) if
and only if each of them has a solution and \(AD=CB\); and that then \(X_0=GC+DH-GADH\) is a common solution.
\end{problem}
\begin{solution}
\begin{itemize}
\item If \(AX=C\) and \(XB=D\), then each equation has a solution, and \(AD=A(XB)=(AX)B=CB\)
(Theorem~\ref{r1-thm-rules}).
\item Conversely, let \(AX_1=C\), \(X_2B=D\) and \(AD=CB\). As \(AGA=A\) and \(BHB=B\), \(AGC=AGAX_1=AX_1=C\)
(Lemma~\ref{r10-lem-range-test}) and \(DHB=X_2BHB=X_2B=D\).
\item Using \(AGA=A\),
\[
AX_0=AGC+ADH-AGADH=C+ADH-ADH=C
\]
\item Using \(DHB=D\) and \(AD=CB\),
\[
X_0B=GCB+DHB-GADHB=GCB+D-GAD=GCB+D-GCB=D
\]
So \(X_0\) is a common solution.
\end{itemize}
\end{solution}

\clearpage
\def\unitnumber{13}
\setcounter{theorem}{0}\setcounter{problem}{0}
\section*{Chapter 13. Linear models: least squares and the Gauss--Markov theorem}

\subsection*{The linear model and least squares}

\begin{definition}[\citeat{tu-clm}{Definitions~1.1, 1.2, 2.1}, \citeat{bapat}{\S7.1}, \citeat{searle-khuri}{(12.12)}]\label{r11-def-model}
In the \textbf{linear model} \(y=X\beta+\varepsilon\), the \textbf{design matrix} \(X\in\R^{n\times(k+1)}\) is a known
constant matrix, \(\beta=(\beta_0,\dots,\beta_k)'\in\R^{k+1}\) is an unknown constant vector of \textbf{parameters}, the
\textbf{errors} \(\varepsilon=(\varepsilon_1,\dots,\varepsilon_n)'\) form a random vector, and the \textbf{response}
\(y=(y_1,\dots,y_n)'\) is observed. In \textbf{multiple regression}, \(x_{ij}\) is the value of the \(j\)th explanatory
variable on observation \(i\) (\(1\leqslant j\leqslant k\)), \(y_i=\beta_0+\beta_1x_{i1}+\dots+\beta_kx_{ik}+\varepsilon_i\), and
the \(i\)th row of \(X\) is \(x_i'=(1,x_{i1},\dots,x_{ik})\): the first column of \(X\) is \(\one_n\), the column of the
\textbf{intercept} \(\beta_0\). The \textbf{classical assumptions} are:
\begin{enumerate}[label=(A\arabic*)]
\item \(\expc(y)=X\beta\);
\item \(\expc(\varepsilon)=0\);
\item \(\var(\varepsilon)=\sigma^2I_n\) for a number \(\sigma^2>0\): the errors have the same variance
\(\sigma^2\) and are uncorrelated, \(\cov(\varepsilon_i,\varepsilon_j)=0\) for \(i\ne j\);
\item \(\rank X=k+1\): \(X\) has full column rank.
\end{enumerate}
(A1)--(A4) are the \textbf{Gauss--Markov assumptions}. Results stated for \(X\in\R^{n\times(k+1)}\) hold for every such
\(X\), with or without the column \(\one_n\) and of any rank, unless (A4) or an intercept is assumed.

A \textbf{least-squares estimate} of \(\beta\) is a vector \(\hat\beta\in\R^{k+1}\) that minimizes
\[
S(\beta)=\lVert y-X\beta\rVert^2=\sum_{i=1}^n(y_i-x_i'\beta)^2\qquad(\beta\in\R^{k+1})
\]
and the \(k+1\) equations \(X'X\beta=X'y\) are the \textbf{normal equations}.
\end{definition}
By Theorem~\ref{r11-thm-linear}(b), with \(A=I_n\) and the constant vector \(b=X\beta\),
\(\expc(y)=X\beta+\expc(\varepsilon)\) and \(\var(y)=\var(\varepsilon)\). So
(A1) and (A2) are equivalent, and under (A3), \(\var(y)=\sigma^2I_n\). Since \(\rank X\leqslant n\)
(Theorem~\ref{thm:r4-rank-rules}(a)), (A4) needs at least \(k+1\) observations. Estimators of the parameters
are functions of \(y\) that involve no unknown parameter (Definition~\ref{r11-def-estimator}).

\begin{theorem}[\citeat{harville}{Theorem~7.4.1}, \citeat{tu-clm}{Theorem~2.4}, \citeat{searle-khuri}{\S12.3}]\label{r10-thm-normal-equations}
Let \(X\in\R^{n\times(k+1)}\) and \(y\in\R^n\).
\begin{enumerate}
\item The normal equations \(X'X\beta=X'y\) (unknown \(\beta\in\R^{k+1}\)) are consistent.
\item \(X'X\) is invertible if and only if \(X\) has full column rank, \(\rank X=k+1\). Then the normal
equations have exactly one solution, \(\hat\beta=(X'X)^{-1}X'y\).
\item Let \(G\) be a g-inverse of \(X'X\). Then \(GX'y\) is a solution, and the solutions are the vectors
\(GX'y+(I_{k+1}-GX'X)z\) with \(z\in\R^{k+1}\); equivalently, \(\hat\beta+z\) with \(z\in\mathcal N(X)\), for any one
solution \(\hat\beta\). All of them give the same vector \(X\hat\beta\). If \(\rank X<k+1\), there are infinitely many
solutions.
\end{enumerate}
\end{theorem}
\begin{proof}
(a) \(X'y\in\mathcal C(X')=\mathcal C(X'X)\) (Theorem~\ref{thm:r4-rank-ata}(c)), so the system is consistent
(Corollary~\ref{cor:r4-ax-b}(a)).

(b) \(X'X\in\R^{(k+1)\times(k+1)}\) is invertible if and only if \(\rank(X'X)=k+1\)
(Theorem~\ref{thm:r4-rank-rules}(d)), that is, \(\rank X=k+1\) (Theorem~\ref{thm:r4-rank-ata}(b)). Then
multiplying \(X'X\beta=X'y\) by \((X'X)^{-1}\) gives \(\beta=(X'X)^{-1}X'y\), and this vector is a solution.

(c)
\begin{itemize}
\item By (a), \(X'y\in\mathcal C(X'X)\). Theorems~\ref{r10-thm-ginv-solves}(a) and~\ref{r10-thm-general-solution}, applied
to the matrix \(X'X\) and the right side \(X'y\), show that \(GX'y\) is a solution and that the solutions are the
vectors \(GX'y+(I-GX'X)z\).
\item By Theorem~\ref{thm:r2-solution-structure}(a), the solutions are \(\hat\beta+z\) with
\(z\in\mathcal N(X'X)=\mathcal N(X)\) (Theorem~\ref{thm:r4-rank-ata}(a)), and
\(X(\hat\beta+z)=X\hat\beta+Xz=X\hat\beta\).
\item If \(\rank X<k+1\), some \(z\ne0\) lies in \(\mathcal N(X)\) (Theorem~\ref{thm:r4-rank-rules}(b)), and the solutions
\(\hat\beta+tz\) (\(t\in\R\)) are pairwise distinct, since \(t_1z=t_2z\) forces \(t_1=t_2\).\qedhere
\end{itemize}
\end{proof}

Let \(X\in\R^{n\times p}\) and \(c>0\). The matrix \(X'X\) is PSD (Theorem~\ref{r8-thm-XtX}), so its
eigenvalues are \(\geqslant0\) (Theorem~\ref{r8-thm-eigen-criterion}(a)). By Lemma~\ref{r9-lem-eig-order}(c), every
eigenvalue of the symmetric matrix \(X'X+cI_p\) is \(\geqslant c>0\). So \(X'X+cI_p\), the matrix used in ridge
regression, is PD and invertible (Theorem~\ref{r8-thm-eigen-criterion}(a),(c)), even when \(\rank X<p\) and \(X'X\) is
singular.

\begin{theorem}[\citeat{tu-clm}{Theorems~2.4, 2.5}, \citeat{harville}{Theorem~12.4.3}]\label{r10-thm-least-squares}
Let \(X\in\R^{n\times(k+1)}\) and \(y\in\R^n\).
\begin{enumerate}
\item \(\hat\beta\) is a least-squares estimate if and only if \(X'X\hat\beta=X'y\). For every solution
\(\hat\beta\) of the normal equations,
\[
S(\beta)=S(\hat\beta)+\lVert X(\hat\beta-\beta)\rVert^2\qquad(\beta\in\R^{k+1})
\]
\item All least-squares estimates give the same vector of \textbf{fitted values} \(\hat y=X\hat\beta\), the
orthogonal projection of \(y\) onto \(\mathcal C(X)\); the vector of \textbf{residuals} \(e=y-\hat y\) lies in
\(\mathcal C(X)^\perp\).
\item If \(\rank X=k+1\), the least-squares estimate is unique, \(\hat\beta=(X'X)^{-1}X'y\), and
\(\hat y=Hy\), \(e=My\), with the \textbf{hat matrix} \(H=X(X'X)^{-1}X'\) (the matrix \(P\) of
Theorem~\ref{r6-thm-hat} for \(A=X\)) and the \textbf{residual maker} \(M=I_n-H\).
\end{enumerate}
\end{theorem}
\begin{proof}
(a)
\begin{itemize}
\item Let \(X'X\hat\beta=X'y\). Then \(X'(y-X\hat\beta)=0\), so \(y-X\hat\beta\perp\mathcal C(X)\)
(Theorem~\ref{r6-thm-complement}(a)); as \(X\hat\beta\in\mathcal C(X)\), it is the orthogonal projection of \(y\)
onto \(\mathcal C(X)\) (Definition~\ref{r6-def-projection}).
\item For \(\beta\in\R^{k+1}\), \(X\beta\in\mathcal C(X)\), and Theorem~\ref{r6-thm-best-approximation} gives
\(\lVert y-X\beta\rVert^2=\lVert y-X\hat\beta\rVert^2+\lVert X\hat\beta-X\beta\rVert^2\), the identity. Hence
\(S(\beta)\geqslant S(\hat\beta)\): \(\hat\beta\) is a least-squares estimate.
\item Solutions exist (Theorem~\ref{r10-thm-normal-equations}(a)). If \(\beta^*\) is a least-squares
estimate and \(\hat\beta\) a solution, then \(S(\beta^*)\leqslant S(\hat\beta)\), so the identity gives
\(\lVert X(\hat\beta-\beta^*)\rVert^2=0\), that is, \(X\beta^*=X\hat\beta\); then \(X'X\beta^*=X'X\hat\beta=X'y\).
\end{itemize}
(b) By the last step, every least-squares estimate \(\beta^*\) gives \(X\beta^*=X\hat\beta\), the projection
found in the first step; and \(e=y-X\hat\beta\perp\mathcal C(X)\).

(c) By Theorem~\ref{r10-thm-normal-equations}(b) and (a), \((X'X)^{-1}X'y\) is the only least-squares
estimate, and then \(\hat y=X(X'X)^{-1}X'y=Hy\) and \(e=y-Hy=My\).
\end{proof}

So least squares needs no calculus: \(\hat y\) is the projection of \(y\) onto \(\mathcal C(X)\). The formula
\(\hat\beta=(X'X)^{-1}X'y\) needs a design matrix of full column rank, \(\rank X=k+1\).
Then Corollary~\ref{r6-cor-trace-hat} with \(p=k+1\) gives \(\tr H=k+1=\rank X\) and \(\tr M=n-k-1\)~\citeat{tu-clm}{Theorem~2.9}.

\begin{theorem}[\citeat{harville}{Corollary~12.2.2, Theorems~12.3.1, 12.3.4}, \citeat{searle-khuri}{(12.13), (12.25)--(12.26)}, \citeat{tu-clm}{Definition~2.6}]\label{r10-thm-PX}\label{r11-prop-hat}
Let \(X\in\R^{n\times(k+1)}\) have rank \(r\), let \(G\) be a g-inverse of \(X'X\) (Definition~\ref{r10-def-ginverse}), and put
\[
H=XGX'=X(X'X)^-X',\qquad M=I_n-H
\]
the \textbf{hat matrix} and the \textbf{residual maker} of \(X\). For every \(y\in\R^n\), \(Hy\) is the orthogonal projection of
\(y\) onto \(\mathcal C(X)\), whatever g-inverse is used; so \(H\) does not depend on the choice of \(G\).
\begin{enumerate}
\item \(H\) and \(M\) are symmetric and idempotent, \(HX=X\), \(X'H=X'\), \(MX=0\) and \(HM=MH=0\).
\item \(\rank H=\tr H=r\) and \(\rank M=\tr M=n-r\).
\item If \(r=k+1\), then \(X'X\) is invertible, its only g-inverse is \((X'X)^{-1}\), and \(H=X(X'X)^{-1}X'\) and \(M\) are the
matrices of Theorem~\ref{r10-thm-least-squares}(c).
\item For every \(y\in\R^n\) and every least-squares estimate \(\hat\beta\), the fitted values and the residuals are
\(\hat y=X\hat\beta=Hy\) and \(e=y-\hat y=My\), and
\[
S(\hat\beta)=e'e=y'My=y'y-\hat\beta'X'y
\]
\end{enumerate}
\end{theorem}
\begin{proof}
The first statement and the independence of \(G\) are Theorem~\ref{r10-thm-ata}(b) with \(A=X\) (\(m=n\)).

(a) By Theorem~\ref{r10-thm-ata}(c),(d) with \(A=X\), \(H'=H=H^2\), \(HX=X\), \(X'H=X'\), and \(M=I_n-H\) is symmetric and
idempotent. Then \(MX=X-HX=0\) and \(HM=H-H^2=0=MH\).

(b) This is Theorem~\ref{r10-thm-ata}(d) with \(A=X\), and \(\tr M=\tr I_n-\tr H=n-r\) (Theorem~\ref{r1-thm-trace}(a)).

(c) \(X'X\) is invertible by Theorem~\ref{r10-thm-normal-equations}(b), and its only g-inverse is \((X'X)^{-1}\)
(Definition~\ref{r10-def-ginverse}); so \(H=X(X'X)^{-1}X'\) (Theorem~\ref{r10-thm-ata}(b) with \(A=X\)).

(d) Every least-squares estimate gives the fitted values \(X\hat\beta\), the orthogonal projection of \(y\) onto
\(\mathcal C(X)\) (Theorem~\ref{r10-thm-least-squares}(b)), which is \(Hy\); so \(e=y-Hy=My\). As \(M\) is symmetric and idempotent,
\(S(\hat\beta)=e'e=y'M'My=y'My\), and \(y'Hy=y'X\hat\beta=\hat\beta'X'y\) (Theorem~\ref{r1-thm-transpose}); hence
\(y'My=y'y-y'Hy=y'y-\hat\beta'X'y\).
\end{proof}

By Theorems~\ref{r10-thm-least-squares}(a),(b) and~\ref{r10-thm-normal-equations}(c), the least-squares estimates \(\hat\beta\) of
\(\beta\) are the solutions of the normal equations \(X'X\hat\beta=X'y\), for instance \(\hat\beta=(X'X)^-X'y\) for any
g-inverse. All of them give the same fitted values \(\hat y=X\hat\beta=Hy\) and residuals \(e=y-\hat y=My\), with \(H\)
and \(M\) as in Theorem~\ref{r10-thm-PX}(d). The minimum of \(\lVert y-X\beta\rVert^2\) is the \textbf{error sum of
squares} \(\mathrm{SSE}=e'e=\sum_{i=1}^n(y_i-\hat y_i)^2\)~\citeat{tu-clm}{Definition~3.1}, \citeat{searle-khuri}{(12.25)}. Under (A4), \(\hat\beta=(X'X)^{-1}X'y\) is the only
least-squares estimate.

By Theorem~\ref{r10-thm-PX}(a),(b) and Theorem~\ref{r8-thm-idempotent-eigen}(a), \(H\) has the eigenvalue \(1\) with
multiplicity \(r\) and \(0\) with multiplicity \(n-r\), and \(M\) the reverse (each an eigenvalue only when this number is
positive; both are eigenvalues when \(1\leqslant r<n\)); as \(Hy\) is the projection of \(y\) onto
\(\mathcal C(X)\), \(Hy=y\) exactly for \(y\in\mathcal C(X)\) and \(Hy=0\) exactly for \(y\in\mathcal C(X)^\perp\). For \(X=\one_n\),
\(H=\bar J_n\) and \(M=C_n\) (Example~\ref{r6-exm-centering}). Under (A4), \(X'X\) is PD (Theorem~\ref{r8-thm-XtX}); and if
\(X=QR\) as in Theorem~\ref{r6-thm-qr}, then \(X'X=R'Q'QR=R'R\) with \(R\) invertible (it is triangular with positive
diagonal entries, Theorem~\ref{r2-thm-triangular}(a)), so
\[
H=QR(R'R)^{-1}R'Q'=QRR^{-1}(R')^{-1}R'Q'=QQ'
\]

\begin{proposition}[\citeat{tu-clm}{Theorem~2.12}, \citeat{searle-khuri}{\S11.1.2}]\label{r10-prop-residuals}
Let \(\hat\beta\) be a least-squares estimate, \(\hat y=X\hat\beta\) and \(e=y-\hat y\).
\begin{enumerate}
\item \(X'e=0\): the residuals are orthogonal to every column of \(X\).
\item If \(\one_n\in\mathcal C(X)\) (for instance, if the model has an intercept), then
\(\one'e=e_1+\dots+e_n=0\); hence \(\hat y\) and \(y\) have the same mean \(\bar y\).
\item \(\hat y'e=0\).
\end{enumerate}
\end{proposition}
\begin{proof}
(a) This is \(X'(y-X\hat\beta)=0\), the normal equations (Theorem~\ref{r10-thm-least-squares}(a)).

(b) If \(\one=Xc\), then \(\one'e=c'X'e=0\) by (a), and \(\one'\hat y=\one'y-\one'e=\one'y\).

(c) \(\hat y'e=\hat\beta'X'e=0\).
\end{proof}

\begin{proposition}[\citeat{bapat}{\S8.5}]\label{r12-prop-xtx-inverse}
Let \(X=[\one_n\ X_0]\in\R^{n\times p}\) with \(X_0\in\R^{n\times k}\) and \(p=k+1\), let \(\bar x=\frac1nX_0'\one_n\) be the mean vector of
\(X_0\), and let \(S_c=X_0'C_nX_0\).
\begin{enumerate}
\item \(X'X=\begin{pmatrix}n&n\bar x'\\n\bar x&X_0'X_0\end{pmatrix}\), the Schur complement of the entry \(n\) in
\(X'X\) is \(S_c\), and \(\det(X'X)=n\det S_c\). Hence \(X\) has full column rank \(p=k+1\) if and only if \(S_c\) is
invertible. For \(n=4\) and the data matrix \(X_0\) of Example~\ref{r1-exm-data}, \(\bar x=(1,2)'\) and
\(S_c=4I_2\).
\item If \(S_c\) is invertible, then
\[
(X'X)^{-1}=\begin{pmatrix}\frac1n+\bar x'S_c^{-1}\bar x&-\bar x'S_c^{-1}\\-S_c^{-1}\bar x&S_c^{-1}\end{pmatrix}
\]
\end{enumerate}
\end{proposition}
\begin{proof}
(a) By the block product formula (Theorem~\ref{r5b-thm-block}(c)) and Theorems~\ref{r1-thm-views}(d) and~\ref{r1-thm-transpose},
\(X'X=\left(\begin{smallmatrix}\one'\one&\one'X_0\\X_0'\one&X_0'X_0\end{smallmatrix}\right)\), with \(\one'\one=n\)
and \(X_0'\one=n\bar x\). The Schur complement of \(A=(n)\) is
\(X_0'X_0-(n\bar x)\frac1n(n\bar x')=X_0'X_0-n\bar x\bar x'=X_0'C_nX_0\) (Example~\ref{r1-exm-data}).
Theorem~\ref{r5d-thm-schur-det}(a) gives \(\det(X'X)=n\det S_c\), so \(X'X\) is invertible if and only if \(S_c\) is
(Theorem~\ref{r5d-thm-product}(a)); and \(X'X\) is invertible if and only if \(\rank X=p\)
(Theorem~\ref{r10-thm-normal-equations}(b)). The numbers for Example~\ref{r1-exm-data} are computed there.

(b) Apply Theorem~\ref{r5b-thm-partitioned-inverse}(a) with \(A=(n)\), \(B=n\bar x'\), \(C=n\bar x\), \(D=X_0'X_0\)
and \(E=S_c\): here \(A^{-1}B=\bar x'\) and \(CA^{-1}=\bar x\), so \(A^{-1}+A^{-1}BE^{-1}CA^{-1}=\frac1n+\bar x'S_c^{-1}\bar x\),
\(-A^{-1}BE^{-1}=-\bar x'S_c^{-1}\) and \(-E^{-1}CA^{-1}=-S_c^{-1}\bar x\).
\end{proof}

\begin{example}[cf.\ \citeat{searle-khuri}{Example~11.1}]\label{r10-exm-least-squares}\label{r6-exm-hat}\label{r11-exm-running}
Let \(X_0\in\R^{4\times2}\) be the data matrix of Example~\ref{r1-exm-data} (\(n=4\), \(k=2\)), and let \(y=(9,3,5,3)'\) be a
response measured on the four units. The design matrix \(X=[\one_4\ X_0]\in\R^{4\times3}\) has \(i\)th row
\(x_i'=(1,x_{i1},x_{i2})\); write \(X_{\cdot1}\) and \(X_{\cdot2}\) for its other two columns, the columns of \(X_0\). By
Theorems~\ref{r1-thm-views}(d) and~\ref{r1-thm-transpose}, the \((j,l)\) entry of \(X'X\) is \(\col_j(X)'\col_l(X)\), so
\[
X=\begin{pmatrix}1&2&3\\1&0&3\\1&2&1\\1&0&1\end{pmatrix},\qquad
X'X=\begin{pmatrix}4&4&8\\4&8&8\\8&8&20\end{pmatrix},\qquad
X'y=\begin{pmatrix}9+3+5+3\\18+10\\27+9+5+3\end{pmatrix}=\begin{pmatrix}20\\28\\44\end{pmatrix}
\]
\begin{itemize}
\item By Example~\ref{r1-exm-data}, \(\bar x=(1,2)'\) and \(S_c=X_0'C_4X_0=4I_2\). So \(\det(X'X)=4\det S_c=64\ne0\), and \(X\) has
full column rank \(3\) (Proposition~\ref{r12-prop-xtx-inverse}(a)). With \(S_c^{-1}=\frac14I_2\),
\(\frac14+\bar x'S_c^{-1}\bar x=\frac14+\frac{1+4}4=\frac32\) and \(-\bar x'S_c^{-1}=(-\frac14,\ -\frac12)\),
Proposition~\ref{r12-prop-xtx-inverse}(b) gives
\[
(X'X)^{-1}=\frac14\begin{pmatrix}6&-1&-2\\-1&1&0\\-2&0&1\end{pmatrix}
\]
\item Then
\[
(X'X)^{-1}X'=\frac14\begin{pmatrix}6&-1&-2\\-1&1&0\\-2&0&1\end{pmatrix}\begin{pmatrix}1&1&1&1\\2&0&2&0\\3&3&1&1\end{pmatrix}
=\frac14\begin{pmatrix}-2&0&2&4\\1&-1&1&-1\\1&1&-1&-1\end{pmatrix}
\]
\[
H=X(X'X)^{-1}X'=\begin{pmatrix}1&2&3\\1&0&3\\1&2&1\\1&0&1\end{pmatrix}\frac14\begin{pmatrix}-2&0&2&4\\1&-1&1&-1\\1&1&-1&-1\end{pmatrix}
=\frac14\begin{pmatrix}3&1&1&-1\\1&3&-1&1\\1&-1&3&1\\-1&1&1&3\end{pmatrix}
\]
and
\[
M=I_4-H=\frac14\begin{pmatrix}1&-1&-1&1\\-1&1&1&-1\\-1&1&1&-1\\1&-1&-1&1\end{pmatrix}
\]
Hence
\[
Hy=\frac14\begin{pmatrix}27+3+5-3\\9+9-5+3\\9-3+15+3\\-9+3+5+9\end{pmatrix}=\begin{pmatrix}8\\4\\6\\2\end{pmatrix},\qquad
My=y-Hy=\begin{pmatrix}1\\-1\\-1\\1\end{pmatrix}
\]
\item The least-squares estimate is unique (Theorem~\ref{r10-thm-least-squares}(c)):
\[
\hat\beta=(X'X)^{-1}X'y=\frac14\begin{pmatrix}-2&0&2&4\\1&-1&1&-1\\1&1&-1&-1\end{pmatrix}\begin{pmatrix}9\\3\\5\\3\end{pmatrix}
=\frac14\begin{pmatrix}-18+0+10+12\\9-3+5-3\\9+3-5-3\end{pmatrix}=\begin{pmatrix}1\\2\\1\end{pmatrix}
\]
and indeed \(X'X\hat\beta=(4+8+8,\ 4+16+8,\ 8+16+20)'=(20,28,44)'=X'y\). The fitted values and residuals are
\(\hat y=X\hat\beta=Hy=(8,4,6,2)'\) and \(e=My=(1,-1,-1,1)'\), and the smallest value of \(S\) is
\(\mathrm{SSE}=S(\hat\beta)=e'e=1+1+1+1=4\). As it is positive, no \(\beta\) gives \(X\beta=y\): the system \(X\beta=y\) has no
solution.
\item As Proposition~\ref{r10-prop-residuals} says, \(X'e=(1-1-1+1,\ 2-2,\ 3-3-1+1)'=0\), \(\one'e=1-1-1+1=0\) and
\(\hat y'e=8-4-6+2=0\). By the identity of Theorem~\ref{r10-thm-least-squares}(a), for instance \(\beta=(1,2,2)'\) gives
\(S(\beta)=4+\lVert X(0,0,-1)'\rVert^2=4+(9+9+1+1)=24\).
\end{itemize}
\end{example}

\subsection*{Rank in regression}

\begin{proposition}[\citeat{tu-clm}{Remark~10.2}]\label{r11-prop-collinearity}
Let \(X\in\R^{n\times p}\) with \(p\geqslant2\). Then \(X\) has full column rank if and only if no column of
\(X\) is a linear combination of the other columns. If some column is such a combination
(\textbf{exact collinearity}), then \(Xz=0\) for some \(z\ne0\), the matrix \(X'X\) is singular, and \(\beta\) and
\(\beta+tz\) give the same \(X\beta\) for every \(t\in\R\).
\end{proposition}
\begin{proof}
By Theorem~\ref{thm:r4-rank-rules}(b), \(X\) has full column rank if and only if its columns are
independent, and \(p\geqslant2\) vectors are dependent if and only if one of them is a linear combination
of the others (Proposition~\ref{r3-prop-dependence}(c)). If they are dependent, \(Xz=0\) for some \(z\ne0\)
(Theorem~\ref{thm:r4-rank-rules}(b)), \(X'X\) is singular, as \(\rank(X'X)=\rank X<p\)
(Theorems~\ref{thm:r4-rank-ata}(b) and~\ref{thm:r4-rank-rules}(d)), and \(X(\beta+tz)=X\beta+tXz=X\beta\).
\end{proof}
For example, if a regression uses the scores \(x_{i1}\), \(x_{i2}\) of two tests and also their total
\(x_{i3}=x_{i1}+x_{i2}\), then column \(4\) of \(X\) is the sum of columns \(2\) and \(3\), and
\(X(0,1,1,-1)'=0\).

\begin{proposition}[\citeat{tu-clm}{Proposition~11.4}]\label{r11-prop-reparametrization}
Let \(X\in\R^{n\times p}\), \(T\in\R^{p\times q}\) and \(X_1=XT\in\R^{n\times q}\), with \(\rank X_1=\rank X\).
Then \(\mathcal C(X_1)=\mathcal C(X)\), the least-squares fitted values (and residuals) of the designs \(X_1\)
and \(X\) are the same, and if \(\hat\beta_1\) is a least-squares estimate for \(X_1\), then \(T\hat\beta_1\) is one
for \(X\).
\end{proposition}
\begin{proof}
\(\mathcal C(X_1)\subseteq\mathcal C(X)\) (Theorem~\ref{thm:r4-rank-product}(a)), and both have dimension
\(\rank X\), so they are equal (Theorem~\ref{thm:r4-counting}(a)). For either design, the fitted values
are the projection \(\hat y\) of \(y\) onto this subspace (Theorem~\ref{r10-thm-least-squares}(b)). Finally
\(X(T\hat\beta_1)=X_1\hat\beta_1=\hat y\), so \(\lVert y-X(T\hat\beta_1)\rVert^2=\lVert y-\hat y\rVert^2\leqslant\lVert y-X\beta\rVert^2\) for
every \(\beta\) (Theorem~\ref{r6-thm-best-approximation}).
\end{proof}
For instance, adding to the design \(X\) of Example~\ref{r11-exm-running} the total
\(X_{\cdot1}+X_{\cdot2}\) of its two variables gives \(\tilde X=[X\ \ X_{\cdot1}+X_{\cdot2}]\in\R^{4\times4}\), with
\(X=\tilde XT\) for \(T\) the first three columns of \(I_4\), and \(\rank\tilde X=\rank X=3\). So \(\tilde X\) has the
same fitted values \((8,4,6,2)'\), and \(T(1,2,1)'=(1,2,1,0)'\) is a least-squares estimate for \(\tilde X\).

\needspace{20\baselineskip}
\begin{theorem}[\citeat{searle-khuri}{\S12.1, \S12.2.1}, \citeat{tu-clm}{Remark~10.2, Proposition~11.2}]\label{r11-thm-oneway}
Let \(n_1,\dots,n_g\geqslant1\) and \(n=n_1+\dots+n_g\). Arrange \(n\) observations in \(g\) groups (first the
\(n_1\) observations of group \(1\), then the \(n_2\) of group \(2\), and so on), and let the
\textbf{indicator} \(d_l\in\R^n\) of group \(l\) have entry \(1\) in the rows of group \(l\) and \(0\) elsewhere.
The model \(y_i=\mu+\alpha_l+\varepsilon_i\) for an observation \(i\) of group \(l\) reads \(y=X\beta+\varepsilon\) with
\(\beta=(\mu,\alpha_1,\dots,\alpha_g)'\) and
\[
X=[\one_n\ d_1\ \cdots\ d_g]\in\R^{n\times(g+1)}
\]
Let \(X_1=[d_1\ \cdots\ d_g]\) (intercept dropped) and \(X_2=[\one_n\ d_2\ \cdots\ d_g]\) (indicator \(d_1\)
dropped), both in \(\R^{n\times g}\).
\begin{enumerate}
\item \(X_1\) has full column rank, and \(X=X_1S\) with \(S=[\one_g\ I_g]\in\R^{g\times(g+1)}\) of full row rank.
\item \(\rank X=g\) and \(\mathcal N(X)=\Span\{(-1,1,\dots,1)'\}\): \(X\) does not have full column rank.
\item \(X_1=XT_1\) and \(X_2=XT_2\), where \(T_1,T_2\in\R^{(g+1)\times g}\) are \(I_{g+1}\) with its first,
respectively second, column deleted. Both \(X_1\) and \(X_2\) have full column rank \(g\), and
\(\mathcal C(X_1)=\mathcal C(X_2)=\mathcal C(X)\).
\end{enumerate}
\end{theorem}
\begin{proof}
(a)
\begin{itemize}
\item No observation lies in two groups, so \(d_l'd_m=0\) for \(l\ne m\), and \(d_l'd_l=n_l\). If
\(c_1d_1+\dots+c_gd_g=0\), multiplying by \(d_m'\) gives \(c_mn_m=0\), so \(c_m=0\): \(X_1\) has full column
rank (Theorem~\ref{thm:r4-rank-rules}(b)).
\item Every observation lies in exactly one group, so \(d_1+\dots+d_g=\one_n\), that is,
\(X_1\one_g=\one_n\). By the block product (Theorem~\ref{r5b-thm-block}),
\(X_1S=[X_1\one_g\ \ X_1I_g]=[\one_n\ d_1\ \cdots\ d_g]=X\).
\item The columns of \(S\) include those of \(I_g\), so \(\mathcal C(S)=\R^g\): \(S\) has full row rank
(Theorem~\ref{thm:r4-rank-rules}(c)).
\end{itemize}
(b) By (a) and Theorem~\ref{thm:r4-rank-product}(b), \(\rank X=\rank S=g\) and \(\mathcal N(X)=\mathcal N(S)\). Now
\(S(c_0,c_1,\dots,c_g)'=c_0\one_g+(c_1,\dots,c_g)'\) is \(0\) if and only if \(c_l=-c_0\) for every \(l\), that is,
\((c_0,c_1,\dots,c_g)'=-c_0(-1,1,\dots,1)'\).

(c)
\begin{itemize}
\item The columns of \(XT_1\) are the columns of \(X\) other than the first (Theorem~\ref{r1-thm-views}(a),(b)), so
\(XT_1=X_1\);
likewise \(XT_2=X_2\).
\item \(X_1\) has full column rank by (a). For \(X_2\): as \(\one_n=d_1+\dots+d_g\) and
\(d_1=\one_n-d_2-\dots-d_g\), every column of \(X_2\) lies in \(\mathcal C(X_1)\) and every column of \(X_1\) lies in
\(\mathcal C(X_2)\). So \(\mathcal C(X_2)=\mathcal C(X_1)\), which has dimension \(g\) by (a); as \(X_2\) has \(g\) columns,
it has full column rank.
\item So \(\rank X_1=\rank X_2=g=\rank X\), and Proposition~\ref{r11-prop-reparametrization} gives
\(\mathcal C(X_1)=\mathcal C(X_2)=\mathcal C(X)\).\qedhere
\end{itemize}
\end{proof}
Using the intercept together with all \(g\) indicators (the \textbf{dummy variable trap})~\citeat{tu-clm}{Remark~10.2} thus makes \(X'X\)
singular (Theorems~\ref{thm:r4-rank-ata}(b) and~\ref{thm:r4-rank-rules}(d)); dropping the intercept or one indicator gives a
design of full column rank with the same column space, hence the same fitted values
(Proposition~\ref{r11-prop-reparametrization}).

\begin{example}[cf.\ \citeat{searle-khuri}{Examples~12.1--12.3}]\label{r11-exm-oneway}\label{r10-exm-oneway}\label{r5d-exm-oneway}\label{r3-exm-indicators}\label{r3-exm-coding}
Six plants of three types have weights \(5,7,6\) (type \(1\)), \(3,5\) (type \(2\)) and \(2\) (type \(3\)). With
\(y=(5,7,6,3,5,2)'\), the model of Theorem~\ref{r11-thm-oneway} (\(g=3\), \(n_1=3\), \(n_2=2\), \(n_3=1\)) has
\[
X=\begin{pmatrix}1&1&0&0\\1&1&0&0\\1&1&0&0\\1&0&1&0\\1&0&1&0\\1&0&0&1\end{pmatrix},\qquad
X'X=\begin{pmatrix}6&3&2&1\\3&3&0&0\\2&0&2&0\\1&0&0&1\end{pmatrix},\qquad
X'y=\begin{pmatrix}28\\18\\8\\2\end{pmatrix}
\]
The \((j,l)\) entry of \(X'X\) is \(\col_j(X)'\col_l(X)\) (Theorems~\ref{r1-thm-views}(d)
and~\ref{r1-thm-transpose}), and \(X'y\) holds the total and the
three group totals.
\begin{itemize}
\item \(\rank X=3\) and \(\mathcal N(X)=\Span\{(-1,1,1,1)'\}\) (Theorem~\ref{r11-thm-oneway}), so \(X'X\) is
singular (Theorems~\ref{thm:r4-rank-ata}(b) and~\ref{thm:r4-rank-rules}(d)). The same follows from a Schur complement:
partition \(X'X\) after its first row and column, with \(A=(6)\), \(B=(3,2,1)\), \(C=B'\) and \(D=\diag(3,2,1)\). Then
\(D^{-1}=\diag(\frac13,\frac12,1)\), \(BD^{-1}=(1,1,1)\), and the Schur complement of \(D\) is \(A-BD^{-1}C=6-(3+2+1)=0\); so
\(\rank(X'X)=3+\rank(0)=3\) (Theorem~\ref{r5d-thm-schur-rank}(b)) and \(\det(X'X)=\det D\cdot0=0\)
(Theorem~\ref{r5d-thm-schur-det}(b)).
\item \(\hat\beta=(0,6,4,2)'\) solves the normal equations: \(X'X\hat\beta=(18+8+2,\ 18,\ 8,\ 2)'=X'y\). By
Theorem~\ref{thm:r2-solution-structure}(a) and \(\mathcal N(X'X)=\mathcal N(X)\) (Theorem~\ref{thm:r4-rank-ata}(a)), all
solutions, that is, all least-squares estimates (Theorem~\ref{r10-thm-least-squares}(a)), are
\[
\begin{pmatrix}0\\6\\4\\2\end{pmatrix}+t\begin{pmatrix}-1\\1\\1\\1\end{pmatrix}
=\begin{pmatrix}-t\\6+t\\4+t\\2+t\end{pmatrix}\qquad(t\in\R)
\]
and all give \(X\hat\beta=(6,6,6,4,4,2)'\): the fitted value \(\hat\mu+\hat\alpha_l\) is the mean of group \(l\)
(\(18/3=6\), \(8/2=4\), \(2/1=2\)).
\item \(G=\diag(0,\frac13,\frac12,1)\) is a g-inverse of \(X'X\). Multiplying the columns of \(X'X\) by
\(0,\frac13,\frac12,1\) (Theorem~\ref{r1-thm-views}(a),(b)) gives
\[
X'XG=\begin{pmatrix}0&1&1&1\\0&1&0&0\\0&0&1&0\\0&0&0&1\end{pmatrix}
\]
By Theorem~\ref{r1-thm-views}(c), row \(1\) of \((X'XG)X'X\) is the sum of rows \(2,3,4\) of \(X'X\), which is \((6,3,2,1)\), and
rows \(2,3,4\) are those of \(X'X\); so \(X'XGX'X=X'X\). The solution \(GX'y=(0,\ 18/3,\ 8/2,\ 2)'=(0,6,4,2)'\) of
Theorem~\ref{r10-thm-normal-equations}(c) is the one above with \(t=0\). Multiplying the rows of \(X'X\) by \(0,\frac13,\frac12,1\),
\[
GX'X=\begin{pmatrix}0&0&0&0\\1&1&0&0\\1&0&1&0\\1&0&0&1\end{pmatrix}
\]
\item By Theorem~\ref{r10-thm-invariant-combination}, for the consistent system \(X'X\beta=X'y\), the value \(\lambda'\hat\beta\) is
the same for all least-squares estimates exactly when \(\lambda'GX'X=\lambda'\). For \(\lambda=(0,1,-1,0)'\),
\(\row_2(GX'X)-\row_3(GX'X)=(0,1,-1,0)=\lambda'\), so \(\hat\alpha_1-\hat\alpha_2=(6+t)-(4+t)=2\) for all of them. For
\(\lambda=(1,1,0,0)'\), \(\row_1(GX'X)+\row_2(GX'X)=(1,1,0,0)=\lambda'\), so \(\hat\mu+\hat\alpha_1=6\). For \(\lambda=(0,1,0,0)'\),
\(\row_2(GX'X)=(1,1,0,0)\ne\lambda'\), and indeed \(\hat\alpha_1=6+t\) changes with \(t\).
\item By Theorem~\ref{r10-thm-PX}(d), \(S(\hat\beta)=y'y-\hat\beta'X'y=148-(0+108+32+4)=4\), with
\(y'y=25+49+36+9+25+4=148\) and \(\hat\beta=(0,6,4,2)'\); indeed \(e=y-\hat y=(-1,1,0,-1,1,0)'\) and \(e'e=4\).
\item Intercept dropped: \(X_1'X_1=\diag(3,2,1)\) and \(X_1'y=(18,8,2)'\), so
\(\hat\beta=(X_1'X_1)^{-1}X_1'y=(6,4,2)'\), the three group means.
\item Indicator \(d_1\) dropped:
\[
X_2'X_2=\begin{pmatrix}6&2&1\\2&2&0\\1&0&1\end{pmatrix},\qquad
X_2'y=\begin{pmatrix}28\\8\\2\end{pmatrix},\qquad
X_2'X_2\begin{pmatrix}6\\-2\\-4\end{pmatrix}=\begin{pmatrix}36-4-4\\12-4\\6-4\end{pmatrix}=X_2'y
\]
so \(\hat\beta=(6,-2,-4)'\), the unique solution (\(X_2'X_2\) is invertible, Theorem~\ref{r6-thm-hat}): the mean \(6\) of
group \(1\) and the differences \(4-6\), \(2-6\) of the other group means from it.
\item In all three designs \(\hat y=(6,6,6,4,4,2)'\), as Proposition~\ref{r11-prop-reparametrization} says;
moreover \(T_1(6,4,2)'=(0,6,4,2)'\) and \(T_2(6,-2,-4)'=(6,0,-2,-4)'\) are the solutions above with \(t=0\)
and \(t=-6\).
\item The columns \(d_1,d_2,d_3\) of \(X_1\) form a basis of \(\mathcal C(X)\), the vectors that are constant within each group,
and so do the columns \(\one_6,d_2,d_3\) of \(X_2\) (Theorem~\ref{r11-thm-oneway}(a),(c)). The estimates \((6,4,2)'\) and
\((6,-2,-4)'\) are the coordinates of \(\hat y\) in these two bases: \(\hat y=6d_1+4d_2+2d_3=6\,\one_6-2d_2-4d_3\). In this
\textbf{baseline coding}, the coefficient of \(\one_6\) is the value of group \(1\), and the coefficients of \(d_2\) and \(d_3\)
are the differences between the other groups and group \(1\).
\item A covariate \(x=(1,2,3,1,2,1)'\) measured on the six plants is not constant within group \(1\), so \(x\notin\mathcal C(X)\),
and the four columns of \([\one_6\ d_2\ d_3\ x]\) are linearly independent (Proposition~\ref{r3-prop-adjoin}(b)): adding
the covariate gives a design of full column rank \(4\).
\end{itemize}
\end{example}

\begin{theorem}[\citeat{harville}{Corollary~20.5.5, Theorem~20.6.1}]\label{r10-thm-mnls}
Let \(X\in\R^{n\times(k+1)}\) and \(y\in\R^n\).
\begin{enumerate}
\item \(X^+=(X'X)^+X'\), and \(XX^+=H\), the hat matrix of Theorem~\ref{r10-thm-PX}.
\item \(X^+y\) is a least-squares estimate, and every other least-squares estimate \(\hat\beta\) has
\(\lVert\hat\beta\rVert>\lVert X^+y\rVert\): \(X^+y\) is the shortest least-squares estimate. It is also the only least-squares
estimate in \(\mathcal C(X')\).
\item If \(G\in\R^{(k+1)\times n}\) and \(Gy\) is the shortest least-squares estimate for every \(y\in\R^n\), then
\(G=X^+\).
\end{enumerate}
\end{theorem}
\begin{proof}
The least-squares estimates are the minimizers of \(S(\beta)=\lVert X\beta-y\rVert^2\), that is, of \(\lVert X\beta-y\rVert\)
(Definition~\ref{r11-def-model}). So (a)--(c) are Theorem~\ref{r10-thm-mp-shortest}(a)--(c) with \(A=X\), \(m=n\) and
\(x=\beta\); in (a), \(X(X'X)^-X'=H\) (Theorem~\ref{r10-thm-PX}).
\end{proof}

\needspace{14\baselineskip}
\begin{example}[cf.\ \citeat{searle-khuri}{Example~12.3}]\label{r10-exm-oneway-ginverses}\label{r10-exm-mnls}
In Example~\ref{r11-exm-oneway}, find the g-inverse \(\dot G\) of \(X'X\) built from its submatrix in rows and columns
\(1,2,4\), and the least-squares estimate \(\dot GX'y\). Compare \(\lambda'G\lambda\) and \(\lambda'\dot G\lambda\) for
\(\alpha_1-\alpha_2\), \(\mu+\alpha_1\) and \(\alpha_1\), and find the shortest least-squares estimate \(X^+y\).

\textit{Solution.}
\begin{itemize}
\item \(\rank(X'X)=\rank X=3\) (Theorem~\ref{thm:r4-rank-ata}(b)). The submatrix in rows and columns \(1,2,4\) and its
adjugate are
\[
W=\begin{pmatrix}6&3&1\\3&3&0\\1&0&1\end{pmatrix},\qquad
\operatorname{adj}W=\begin{pmatrix}3&-3&-3\\-3&5&3\\-3&3&9\end{pmatrix}
\]
and \(\det W=1\cdot(0-3)+1\cdot(18-9)=6\) (expansion along row \(3\), Theorem~\ref{r5d-thm-expansion}); so
\(W^{-1}=\frac16\operatorname{adj}W\)
(Theorem~\ref{r5d-thm-adjugate}). By Theorem~\ref{r10-thm-ginv-submatrix}(b), with \(i_s=j_s\) running through \(1,2,4\),
\[
\dot G=\frac16\begin{pmatrix}3&-3&0&-3\\-3&5&0&3\\0&0&0&0\\-3&3&0&9\end{pmatrix},\qquad
\dot GX'y=\frac16\begin{pmatrix}84-54-6\\-84+90+6\\0\\-84+54+18\end{pmatrix}=\begin{pmatrix}4\\2\\0\\-2\end{pmatrix}
\]
the estimate with \(t=-4\) in Example~\ref{r11-exm-oneway}.
\item \(\alpha_1-\alpha_2\), \(\lambda=(0,1,-1,0)'\): with \(G=\diag(0,\frac13,\frac12,1)\),
\(\lambda'G\lambda=\frac13+\frac12=\frac56\), and \(\lambda'\dot G\lambda=\frac16(5-0-0+0)=\frac56\).
\item \(\mu+\alpha_1\), \(\lambda=(1,1,0,0)'\): \(\lambda'G\lambda=0+\frac13=\frac13\), and
\(\lambda'\dot G\lambda=\frac16(3-3-3+5)=\frac13\).
\item \(\alpha_1\), \(\lambda=(0,1,0,0)'\): \(\lambda'G\lambda=\frac13\), but \(\lambda'\dot G\lambda=\frac56\).
\item This agrees with Corollary~\ref{r10-cor-invariant-quadratic} for \(A=X\): the first two \(\lambda\) lie in
\(\mathcal C(X')\), the third does not (Example~\ref{r11-exm-oneway} and Theorem~\ref{r10-thm-invariant-combination}, with
\(\mathcal C(X'X)=\mathcal C(X')\) by Theorem~\ref{thm:r4-rank-ata}(c)). For instance \((0,1,-1,0)=a'X\) with
\(a=(\frac13,\frac13,\frac13,-\frac12,-\frac12,0)'=\frac13d_1-\frac12d_2\); as \(a\in\mathcal C(X)\), its projection onto
\(\mathcal C(X)\) is \(Ha=a\) (Theorem~\ref{r10-thm-PX}, Definition~\ref{r6-def-projection}), and
\(\lambda'G\lambda=\lVert a\rVert^2=3\cdot\frac19+2\cdot\frac14=\frac56\).
\item The least-squares estimates are \(\hat\beta=(-t,\ 6+t,\ 4+t,\ 2+t)'\) (\(t\in\R\)), and
\[
\lVert\hat\beta\rVert^2=t^2+(6+t)^2+(4+t)^2+(2+t)^2=4t^2+24t+56=4(t+3)^2+20
\]
This is smallest exactly for \(t=-3\). So, by Theorem~\ref{r10-thm-mnls}(b), \(X^+y=(3,3,1,-1)'\), of squared length \(20\).
\end{itemize}
\end{example}

\begin{example}\label{exm:r4-two-factors}
Four observations are classified by two factors with two levels each; observations \(1,2,3,4\) have the levels
\((1,1)\), \((2,1)\), \((1,2)\), \((2,2)\) of (first factor, second factor). The indicators of the two levels of the
first factor, and of the second factor, are
\[
d_1=\begin{pmatrix}1\\0\\1\\0\end{pmatrix},\quad d_2=\begin{pmatrix}0\\1\\0\\1\end{pmatrix},\qquad
g_1=\begin{pmatrix}1\\1\\0\\0\end{pmatrix},\quad g_2=\begin{pmatrix}0\\0\\1\\1\end{pmatrix}
\]
Let \(U=\Span\{d_1,d_2\}\) and \(W=\Span\{g_1,g_2\}\). Find the dimensions of \(U\), \(W\), \(U+W\) and \(U\cap W\),
and a basis of \(U\cap W\).

\textit{Solution.} \(U+W=\Span\{d_1,d_2,g_1,g_2\}\) (Theorem~\ref{r3-thm-sum}(c)). Row-reduce
\([d_1\ d_2\ g_1\ g_2]\):
\[
\begin{gathered}
\begin{pmatrix}1&0&1&0\\0&1&1&0\\1&0&0&1\\0&1&0&1\end{pmatrix}
\xrightarrow[R_4\leftarrow R_4-R_2]{R_3\leftarrow R_3-R_1}
\begin{pmatrix}1&0&1&0\\0&1&1&0\\0&0&-1&1\\0&0&-1&1\end{pmatrix}
\xrightarrow{R_3\leftarrow-R_3}
\begin{pmatrix}1&0&1&0\\0&1&1&0\\0&0&1&-1\\0&0&-1&1\end{pmatrix}
\\
\xrightarrow{R_4\leftarrow R_4+R_3}
\begin{pmatrix}1&0&1&0\\0&1&1&0\\0&0&1&-1\\0&0&0&0\end{pmatrix}
\xrightarrow[R_2\leftarrow R_2-R_3]{R_1\leftarrow R_1-R_3}
\begin{pmatrix}1&0&0&1\\0&1&0&1\\0&0&1&-1\\0&0&0&0\end{pmatrix}
\end{gathered}
\]
\begin{itemize}
\item The pivot columns are \(1,2,3\), so \(d_1,d_2,g_1\) form a basis of \(U+W\) and \(\dim(U+W)=3\)
(Theorem~\ref{thm:r4-pivot-columns}(b)). Column \(4\) gives \(g_2=d_1+d_2-g_1\); check:
\((1,0,1,0)'+(0,1,0,1)'-(1,1,0,0)'=(0,0,1,1)'\).
\item Neither of \(d_1,d_2\) is a multiple of the other (each has a \(1\) where the other has a \(0\)), and the same
holds for \(g_1,g_2\); so \(\dim U=\dim W=2\) (Proposition~\ref{r3-prop-dependence}(c)).
\item By Theorem~\ref{thm:r4-dim-sum}, \(\dim(U\cap W)=2+2-3=1\). The relation of column \(4\) reads
\(d_1+d_2=g_1+g_2\), and both sides equal \(\one_4\). So \(\one_4\) lies in \(U\) and in \(W\), and this nonzero vector
is a basis of the one-dimensional space \(U\cap W\) (Theorem~\ref{thm:r4-counting}(d)): \(U\cap W=\Span\{\one_4\}\).
A combination of the indicators of the first factor that is also a combination of the indicators of the second
factor is a constant vector.
\end{itemize}
In particular, in the additive model \(y_i=\alpha_l+\gamma_m+\varepsilon_i\) for an observation \(i\) at level \(l\) of the
first factor and level \(m\) of the second, the design matrix \(X=[d_1\ d_2\ g_1\ g_2]\in\R^{4\times4}\) has rank \(3\), because
the indicators of each factor add up to \(\one_4\): \(X(1,1,-1,-1)'=0\), and \(X'X\) is singular
(Theorems~\ref{thm:r4-rank-ata}(b) and~\ref{thm:r4-rank-rules}(d)), as for the one-way layout of Theorem~\ref{r11-thm-oneway}.
\end{example}

\begin{proposition}[\citeat{tu-clm}{Definition~11.6}]\label{r11-prop-quadratic}
Let \(x=(x_1,\dots,x_n)'\in\R^n\), and let \(X\in\R^{n\times3}\) have \(i\)th row \((1,x_i,x_i^2)\): the design
matrix of the quadratic regression \(y_i=\beta_0+\beta_1x_i+\beta_2x_i^2+\varepsilon_i\) (\(k=2\), with
\(x_{i1}=x_i\) and \(x_{i2}=x_i^2\)). Then \(X\) has full column rank if and only if \(x_1,\dots,x_n\) take at
least three different values.
\end{proposition}
\begin{proof}
\begin{itemize}
\item If every \(x_i\) is one of two numbers \(a,b\) (possibly \(a=b\)), then
\((x_i-a)(x_i-b)=ab-(a+b)x_i+x_i^2=0\) for every \(i\); so \(X(ab,\,-(a+b),\,1)'=0\), and \(X\) does not
have full column rank (Theorem~\ref{thm:r4-rank-rules}(b)).
\item Let \(x_i=a\), \(x_j=b\), \(x_l=c\) be three different values. If \(Xv=0\), rows \(i,j,l\) give \(Av=0\) for
\[
A=\begin{pmatrix}1&a&a^2\\1&b&b^2\\1&c&c^2\end{pmatrix}
\]
The operations \(R_2\leftarrow R_2-R_1\), \(R_3\leftarrow R_3-R_1\) give the rows \((0,\,b-a,\,b^2-a^2)\) and
\((0,\,c-a,\,c^2-a^2)\); dividing them by \(b-a\ne0\) and \(c-a\ne0\) gives \((0,1,a+b)\) and \((0,1,a+c)\);
then \(R_3\leftarrow R_3-R_2\) gives \((0,0,c-b)\), and dividing by \(c-b\ne0\) gives an echelon form with a
leading \(1\) in every column. So \(Av=0\) only for \(v=0\) (Theorem~\ref{thm:r4-pivot-test} and
Remark~\ref{rem:r4-echelon}), and hence
\(Xv=0\) only for \(v=0\).\qedhere
\end{itemize}
\end{proof}

\begin{example}\label{r11-exm-quadratic}
(a) For \(x=(0,1,2,3)'\), four different values, \(X\) has full column rank by
Proposition~\ref{r11-prop-quadratic}. Directly:
\[
\begin{gathered}
\begin{pmatrix}1&0&0\\1&1&1\\1&2&4\\1&3&9\end{pmatrix}
\xrightarrow[R_4\leftarrow R_4-R_1]{R_2\leftarrow R_2-R_1,\ R_3\leftarrow R_3-R_1}
\begin{pmatrix}1&0&0\\0&1&1\\0&2&4\\0&3&9\end{pmatrix}
\xrightarrow[R_4\leftarrow R_4-3R_2]{R_3\leftarrow R_3-2R_2}
\begin{pmatrix}1&0&0\\0&1&1\\0&0&2\\0&0&6\end{pmatrix}
\\
\xrightarrow{R_4\leftarrow R_4-3R_3}
\begin{pmatrix}1&0&0\\0&1&1\\0&0&2\\0&0&0\end{pmatrix}
\xrightarrow{R_3\leftarrow\frac12R_3}
\begin{pmatrix}1&0&0\\0&1&1\\0&0&1\\0&0&0\end{pmatrix}
\end{gathered}
\]
an echelon form with three nonzero rows, so \(\rank X=3\).

(b) If a dose is given at only two levels, \(1\) and \(3\), say \(x=(1,3,3,1)'\), then \(x_i^2=4x_i-3\) for
every \(i\) (\(1=4-3\), \(9=12-3\)). So column \(3\) of \(X\) is \(4\) times column \(2\) minus \(3\) times
column \(1\), and \(X(3,-4,1)'=0\). Columns \(1\) and \(2\), \((1,1,1,1)'\) and \((1,3,3,1)'\), are not multiples
of each other, so \(\rank X=2\). Thus \(\beta\) and \(\beta+t(3,-4,1)'\) give the same \(X\beta\): two dose levels
cannot determine the curvature \(\beta_2\).
\end{example}

\subsection*{The Gauss--Markov theorem}

\begin{theorem}[\citeat{tu-clm}{Theorems~3.3, 3.4, 3.8, Corollary~3.10}, \citeat{searle-khuri}{\S11.2.1}]\label{r11-thm-unbiased}
In (a) and (b), assume (A4), and let \(\hat\beta=(X'X)^{-1}X'y\).
\begin{enumerate}
\item \(\hat\beta=\beta+(X'X)^{-1}X'\varepsilon\); under (A2), \(\expc(\hat\beta)=\beta\).
\item Under (A3), \(\var(\hat\beta)=\sigma^2(X'X)^{-1}\).
\item For \(X\) of any rank, with \(H\) and \(M\) of Theorem~\ref{r10-thm-PX}, under (A2) and (A3): \(e=M\varepsilon\),
\(\expc(e)=0\) and \(\var(e)=\sigma^2M\); \(\expc(\hat y)=X\beta\) and \(\var(\hat y)=\sigma^2H\); and \(\cov(\hat y,e)=0\).
Under (A4) also \(\cov(\hat\beta,e)=0\), with \(\hat\beta\) of (a).
\end{enumerate}
Only (A2)--(A4) are used; no distribution of \(\varepsilon\) is assumed.
\end{theorem}
\begin{proof}
(a) \(\hat\beta=(X'X)^{-1}X'(X\beta+\varepsilon)=(X'X)^{-1}X'X\beta+(X'X)^{-1}X'\varepsilon=\beta+(X'X)^{-1}X'\varepsilon\).
By Theorem~\ref{r11-thm-linear}(b) and (A2), \(\expc(\hat\beta)=\beta+(X'X)^{-1}X'\expc(\varepsilon)=\beta\).

(b) \((X'X)^{-1}\) is symmetric (Theorem~\ref{r2-thm-inverse-props}(c)). By (a), Theorem~\ref{r11-thm-linear}(b) and (A3),
\[
\var(\hat\beta)=(X'X)^{-1}X'(\sigma^2I_n)X(X'X)^{-1}=\sigma^2(X'X)^{-1}(X'X)(X'X)^{-1}=\sigma^2(X'X)^{-1}
\]
(c)
\begin{itemize}
\item \(e=My\) and \(\hat y=Hy\) (Theorem~\ref{r10-thm-PX}(d)), and \(MX=0\), \(HX=X\) (Theorem~\ref{r10-thm-PX}(a)), so
\(e=MX\beta+M\varepsilon=M\varepsilon\) and \(\hat y=X\beta+H\varepsilon\). As \(H\) and \(M\) are symmetric and idempotent
(Theorem~\ref{r10-thm-PX}(a)), Theorem~\ref{r11-thm-linear}(b) with (A2) and (A3) gives \(\expc(e)=0\),
\(\var(e)=M(\sigma^2I_n)M'=\sigma^2M\), \(\expc(\hat y)=X\beta\) and \(\var(\hat y)=\sigma^2H\).
\item By Theorem~\ref{r14-thm-cov}(a) and (A3),
\(\cov(\hat y,e)=\cov(H\varepsilon,M\varepsilon)=\sigma^2HM'=\sigma^2HM=0\) (Theorem~\ref{r10-thm-PX}(a)).
\item Under (A4), with \(A=(X'X)^{-1}X'\), (a) gives \(\hat\beta=\beta+A\varepsilon\), and by Theorem~\ref{r14-thm-cov}(a)
and (A3), \(\cov(\hat\beta,e)=\cov(A\varepsilon,M\varepsilon)=A(\sigma^2I_n)M'=\sigma^2(X'X)^{-1}(MX)'=0\).\qedhere
\end{itemize}
\end{proof}

\needspace{14\baselineskip}
\begin{theorem}[Gauss--Markov; \citeat{tu-clm}{Theorem~3.5}, \citeat{searle-khuri}{Theorem~11.1}]\label{r11-thm-gauss-markov}
Assume (A1)--(A4), and let \(\hat\beta=(X'X)^{-1}X'y\). Let \(C\in\R^{(k+1)\times n}\) be constant, and let
\(\tilde\beta=Cy\) be a \textbf{linear unbiased estimator} of \(\beta\): \(\expc(Cy)=\beta\) for every
\(\beta\in\R^{k+1}\). Then \(\var(\tilde\beta)-\var(\hat\beta)\) is positive semidefinite, and it is \(0\)
only if \(\tilde\beta=\hat\beta\). Consequently \(\var(a'\tilde\beta)\geqslant\var(a'\hat\beta)\) for every
\(a\in\R^{k+1}\); in particular \(\var(\tilde\beta_j)\geqslant\var(\hat\beta_j)\) for \(j=0,\dots,k\). We
say that \(\hat\beta\) is the BLUE of \(\beta\).
\end{theorem}
\begin{proof}
\begin{itemize}
\item By Theorem~\ref{r11-thm-linear}(b) and (A1), \(\expc(Cy)=CX\beta\). As \(CX\beta=\beta\) for every \(\beta\),
taking for \(\beta\) the columns of \(I_{k+1}\) gives \(CX=I_{k+1}\).
\item Write \(C=(X'X)^{-1}X'+D\). Then \(DX=CX-(X'X)^{-1}X'X=I_{k+1}-I_{k+1}=0\), and \(X'D'=(DX)'=0\).
\item By Theorem~\ref{r11-thm-linear}(b) and (A3), \(\var(\tilde\beta)=C(\sigma^2I_n)C'=\sigma^2CC'\), and
\[
CC'=(X'X)^{-1}X'X(X'X)^{-1}+(X'X)^{-1}X'D'+DX(X'X)^{-1}+DD'=(X'X)^{-1}+DD'
\]
\item By Theorem~\ref{r11-thm-unbiased}(b), \(\var(\tilde\beta)-\var(\hat\beta)=\sigma^2DD'\), which is
positive semidefinite (Theorem~\ref{r8-thm-XtX}, as \(DD'=(D')'D'\)). If it is \(0\), then \(\tr\bigl((D')'D'\bigr)=0\), so
\(D'=0\) (Theorem~\ref{r1-thm-trace}(d)), \(C=(X'X)^{-1}X'\) and \(\tilde\beta=\hat\beta\).
\item For \(a\in\R^{k+1}\), \(\var(a'\tilde\beta)=a'\var(\tilde\beta)a\) and likewise for \(\hat\beta\)
(Theorem~\ref{r11-thm-linear}(b)); so
\(\var(a'\tilde\beta)-\var(a'\hat\beta)=\sigma^2a'DD'a\geqslant0\). For \(a\) the column of \(I_{k+1}\) that
picks out \(\beta_j\), this is \(\var(\tilde\beta_j)\geqslant\var(\hat\beta_j)\).\qedhere
\end{itemize}
\end{proof}

\begin{example}\label{r11-exm-blue}
In Example~\ref{r11-exm-running}, \(\expc(y_1-y_2)=x_1'\beta-x_2'\beta=(0,2,0)\beta=2\beta_1\) by (A1). So
\(c'y=\frac12(y_1-y_2)\), with \(c=\frac12(1,-1,0,0)'\), is a linear unbiased estimator of \(\beta_1\) (\(a=(0,1,0)'\)); its
variance is \(\sigma^2c'c=\sigma^2/2\). The BLUE is \(\hat\beta_1=\frac14(y_1-y_2+y_3-y_4)\), row \(2\) of \((X'X)^{-1}X'y\)
(Example~\ref{r6-exm-hat}), with variance \(\sigma^2\cdot\frac14\), the \((2,2)\) entry of
\(\sigma^2(X'X)^{-1}\) (Theorem~\ref{r11-thm-unbiased}(b)). As in the proof of Theorem~\ref{r11-thm-gauss-markov}, write
\(c'=a'(X'X)^{-1}X'+d'\) with \(d'X=0\); then \(\var(c'y)-\var(\hat\beta_1)=\sigma^2d'd\). Here \(X'c=a\)
(unbiasedness, as in the first step of that proof), so \(X(X'X)^{-1}a=X(X'X)^{-1}X'c=Hc\) and \(d=c-Hc=Mc\):
\[
Hc=\tfrac14(1,-1,1,-1)',\qquad d=Mc=c-Hc=\tfrac14(1,-1,-1,1)',\qquad
\frac{\sigma^2}2-\frac{\sigma^2}4=\sigma^2\lVert d\rVert^2=\sigma^2\cdot\tfrac4{16}
\]
For \(y=(9,3,5,3)'\) the two estimates of \(\beta_1\) are \(\frac12(9-3)=3\) and \(\hat\beta_1=2\).
\end{example}

\begin{example}\label{r11-exm-added-observation}
For an invertible \(X'X\) and \(u=v=x\), Theorem~\ref{r5d-thm-woodbury}(b) updates \((X'X)^{-1}\) when an observation with row \(x'\) is added to the
design matrix \(X\), since then \(X'X\) becomes \(X'X+xx'\) (Theorem~\ref{r5b-thm-block}(c)): no new inversion is needed.
The condition of (b) holds: with \(w=(X'X)^{-1}x\), \(x'(X'X)^{-1}x=(X'Xw)'w=(Xw)'(Xw)\geqslant0\)
(Theorem~\ref{r1-thm-transpose}), so \(1+x'(X'X)^{-1}x\geqslant1\).

For instance, let four observations fall into a baseline group (observations \(1,2\)) and groups \(A\) (observation \(3\))
and \(B\) (observation \(4\)), with the baseline coding \(y_i=\beta_0+\beta_1a_i+\beta_2b_i+\varepsilon_i\), where \(a_i\), \(b_i\) are
the indicators of \(A\) and \(B\); add a fifth observation, in group \(A\), with row \(x'=(1,1,0)\):
\[
X=\begin{pmatrix}1&0&0\\1&0&0\\1&1&0\\1&0&1\end{pmatrix},\qquad
X'X=\begin{pmatrix}4&1&1\\1&1&0\\1&0&1\end{pmatrix},\qquad
(X'X)^{-1}=\frac12\begin{pmatrix}1&-1&-1\\-1&3&1\\-1&1&3\end{pmatrix}
\]
(\(\det(X'X)=4-1-1=2\), and \(X'X\cdot(X'X)^{-1}=I_3\), Theorem~\ref{r2-thm-invertible}).
\begin{itemize}
\item \(w=(X'X)^{-1}x=\frac12(1-1,\ -1+3,\ -1+1)'=(0,1,0)'\) and \(x'(X'X)^{-1}x=x'w=1\). The new design \(\tilde X\) has
\(\tilde X'\tilde X=X'X+xx'\), and Theorem~\ref{r5d-thm-woodbury}(b) gives
\[
\tilde X'\tilde X=\begin{pmatrix}5&2&1\\2&2&0\\1&0&1\end{pmatrix},\qquad
(\tilde X'\tilde X)^{-1}=(X'X)^{-1}-\frac1{1+1}ww'=\frac12\begin{pmatrix}1&-1&-1\\-1&2&1\\-1&1&3\end{pmatrix}
\]
\item The difference \((X'X)^{-1}-(\tilde X'\tilde X)^{-1}=\frac12ww'\) is PSD (Theorem~\ref{r8-thm-XtX}), as
Theorem~\ref{r8-thm-loewner-inverse} predicts from \(\tilde X'\tilde X\succcurlyeq X'X\succ0\); and
\(\det(\tilde X'\tilde X)=\det(X'X)\,\bigl(1+x'(X'X)^{-1}x\bigr)=2\cdot2=4\) (Theorem~\ref{r5d-thm-det-lemma}).
\item Under (A3), \(\var(\hat\beta)=\sigma^2(X'X)^{-1}\) (Theorem~\ref{r11-thm-unbiased}(b)). So the added observation
lowers \(\var(\hat\beta_1)\), the variance of the estimated difference between group \(A\) and the baseline, from
\(\frac32\sigma^2\) to \(\sigma^2\), and leaves the other variances and covariances unchanged, as \(ww'\) has only one nonzero
entry.
\end{itemize}
\end{example}

\begin{proposition}[\citeat{harville}{Theorem~14.12.20}, \citeat{tu-clm}{Exercise~B5}]\label{r12-prop-wls}
Let \(X\in\R^{n\times p}\), let \(W\in\R^{n\times n}\) be PD, and let \(y\in\R^n\). The minimizers of
\(Q(\beta)=(y-X\beta)'W(y-X\beta)\) (\(\beta\in\R^p\)) are the solutions of \(X'WX\beta=X'Wy\). If \(\rank X=p\),
then \(X'WX\) is PD, and the unique minimizer of \(Q\) is the \textbf{weighted least-squares estimate}
\[
\hat\beta_W=(X'WX)^{-1}X'Wy
\]
For \(W=\diag(w_1,\dots,w_n)\) with all \(w_i>0\), \(Q(\beta)=\sum_iw_i(y_i-x_i'\beta)^2\).
\end{proposition}
\begin{proof}
\begin{itemize}
\item Let \(T=W^{1/2}\) (Theorem~\ref{r8-thm-sqrt}): \(T\) is PD, hence symmetric and invertible, and \(T^2=W\). So
\(Q(\beta)=(y-X\beta)'T'T(y-X\beta)=\lVert Ty-TX\beta\rVert^2\): the weighted sum of squares of the response \(y\) and the
design \(X\) is the sum of squares of the response \(Ty\) and the design \(TX\).
\item By Theorem~\ref{r10-thm-least-squares}(a) for the design \(TX\) and the response \(Ty\), the minimizers of \(Q\) are the
solutions of \((TX)'TX\beta=(TX)'Ty\), that is, of \(X'WX\beta=X'Wy\), as \(T'T=T^2=W\).
\item If \(\rank X=p\), then \(\rank(TX)=\rank X=p\) (Theorem~\ref{thm:r4-rank-product}(d)), so \(X'WX=(TX)'(TX)\) is PD
(Theorem~\ref{r8-thm-XtX}), and by Theorem~\ref{r10-thm-least-squares}(c) the unique minimizer of \(Q\) is
\(\bigl((TX)'TX\bigr)^{-1}(TX)'Ty=(X'WX)^{-1}X'Wy=\hat\beta_W\).
\item A diagonal \(W\) with positive diagonal entries is PD (Proposition~\ref{r8-prop-pd-basic}(a)), and
\(Q(\beta)=\sum_iw_iu_i^2\) for \(u=y-X\beta\).\qedhere
\end{itemize}
\end{proof}

\begin{theorem}[Aitken; \citeat{tu-clm}{Exercise~B5}, \citeat{searle-khuri}{\S11.6.1}]\label{r12-thm-gls}
Let \(y=X\beta+\varepsilon\), where \(X\in\R^{n\times p}\) is known with \(\rank X=p\), \(\expc(\varepsilon)=0\) and
\(\var(\varepsilon)=\sigma^2V\) with \(V\in\R^{n\times n}\) known and PD, and \(\beta\) and \(\sigma^2>0\) are
unknown. Then \(X'V^{-1}X\) is invertible. Put
\[
\tilde\beta=(X'V^{-1}X)^{-1}X'V^{-1}y\qquad\text{(the \textbf{generalized least-squares estimator})}
\]
\begin{enumerate}
\item \(\tilde\beta\) is the unique minimizer of \((y-X\beta)'V^{-1}(y-X\beta)\), \(\expc(\tilde\beta)=\beta\) and
\(\var(\tilde\beta)=\sigma^2(X'V^{-1}X)^{-1}\).
\item \(\tilde\beta\) is the BLUE of \(\beta\): for every linear unbiased estimator \(\check\beta=Cy\) of \(\beta\), with
\(C\in\R^{p\times n}\) constant, \(\var(\check\beta)-\var(\tilde\beta)\) is PSD, and it is \(0\) only if \(\check\beta=\tilde\beta\);
so \(\var(c'\check\beta)\geqslant\var(c'\tilde\beta)\) for every \(c\in\R^p\).
\item \(\hat\beta=(X'X)^{-1}X'y\) is unbiased, and \(\var(\hat\beta)=\sigma^2(X'X)^{-1}X'VX(X'X)^{-1}\); the matrix
\(\var(\hat\beta)-\var(\tilde\beta)\) is PSD.
\end{enumerate}
\end{theorem}
\begin{proof}
\begin{itemize}
\item Let \(T=V^{-1/2}\), so that \(TVT=I_n\) and \(T^2=V^{-1}\) (Theorem~\ref{r8-thm-sqrt}(c)). Also
\(T=(V^{1/2})^{-1}\) is symmetric, as \(V^{1/2}\) is PD, hence symmetric
(Theorem~\ref{r2-thm-inverse-props}(c)).
\item Put \(y^*=Ty\), \(X^*=TX\), \(\varepsilon^*=T\varepsilon\). Then \(y^*=X^*\beta+\varepsilon^*\), \(\expc(\varepsilon^*)=0\) and
\(\var(\varepsilon^*)=T(\sigma^2V)T'=\sigma^2I_n\) (Theorem~\ref{r11-thm-linear}(b)), and \(\rank X^*=\rank X=p\)
(Theorem~\ref{thm:r4-rank-product}(d)): the transformed model is a linear model with full column rank.
\item \(X^{*\prime}X^*=X'T^2X=X'V^{-1}X\), which is therefore invertible (Theorem~\ref{r10-thm-normal-equations}(b)),
and \(X^{*\prime}y^*=X'V^{-1}y\). So \(\tilde\beta\) is the least-squares estimator of the transformed model.
\end{itemize}

(a)
\begin{itemize}
\item \(\lVert y^*-X^*\beta\rVert^2=(y-X\beta)'T^2(y-X\beta)=(y-X\beta)'V^{-1}(y-X\beta)\), whose unique minimizer is
\(\tilde\beta\) (Theorem~\ref{r10-thm-least-squares}(c)).
\item Theorem~\ref{r11-thm-unbiased}(a),(b) for the transformed model gives \(\expc(\tilde\beta)=\beta\) and
\(\var(\tilde\beta)=\sigma^2(X^{*\prime}X^*)^{-1}\).
\end{itemize}

(b)
\begin{itemize}
\item For constant \(C\in\R^{p\times n}\), \(Cy=(CT^{-1})y^*\), and \(CT^{-1}\) runs through all constant \(p\times n\) matrices
as \(C\) does. So the linear estimators of \(\beta\) of the two descriptions are the same random vectors, with the same
means and variances, and the linear unbiased estimators of \(\beta\) are the same.
\item Theorem~\ref{r11-thm-gauss-markov} for the transformed model, whose least-squares estimator is \(\tilde\beta\), gives the
claim.
\end{itemize}

(c)
\begin{itemize}
\item \(\hat\beta=Ay\) with \(A=(X'X)^{-1}X'\) and \(AX=I_p\), so \(\expc(\hat\beta)=AX\beta=\beta\) and
\(\var(\hat\beta)=A(\sigma^2V)A'\) (Theorem~\ref{r11-thm-linear}(b)).
\item So \(\hat\beta\) is a linear unbiased estimator of \(\beta\), and (b) gives that \(\var(\hat\beta)-\var(\tilde\beta)\) is
PSD.\qedhere
\end{itemize}
\end{proof}
If the errors are uncorrelated with \(\var(\varepsilon_i)=\sigma^2v_i\) and known \(v_i>0\) (\textbf{heteroscedastic}
errors), then \(V=\diag(v_1,\dots,v_n)\) and \(V^{-1}=\diag(1/v_1,\dots,1/v_n)\), and \(\tilde\beta\) is the weighted
least-squares estimate with weights \(w_i=1/v_i\) (Proposition~\ref{r12-prop-wls}).

\subsection*{Estimating \(\sigma^2\)}

\begin{theorem}[\citeat{searle-khuri}{\S12.6}, \citeat{bapat}{Result~7.5}, \citeat{tu-clm}{Theorem~3.7}]\label{r11-thm-sigma}
Assume (A2) and (A3), and let \(\rank X=r<n\). Then \(\expc(\mathrm{SSE})=(n-r)\sigma^2\). Hence
\[
\hat\sigma^2=\frac{\mathrm{SSE}}{n-r}
\]
is an unbiased estimator of \(\sigma^2\). Under (A4), \(r=k+1\), and \(\hat\sigma^2\) is the \textbf{mean squared error}
\(s^2=\mathrm{MSE}=\mathrm{SSE}/(n-k-1)\).
\end{theorem}
\begin{proof}
\begin{itemize}
\item \(e=M\varepsilon\) (Theorem~\ref{r11-thm-unbiased}(c)).
\item \(M\) is symmetric and idempotent (Theorem~\ref{r10-thm-PX}(a)), so
\(\mathrm{SSE}=e'e=\varepsilon'M'M\varepsilon=\varepsilon'M\varepsilon\).
\item By Theorem~\ref{r11-thm-quadratic} with \(z=\varepsilon\), \(\mu=0\) (A2) and \(\Sigma=\sigma^2I_n\) (A3),
\(\expc(\mathrm{SSE})=\tr(M\sigma^2I_n)=\sigma^2\tr M=\sigma^2(n-r)\) (Theorem~\ref{r10-thm-PX}(b)).
\item Dividing by \(n-r>0\) gives \(\expc(\hat\sigma^2)=\sigma^2\).\qedhere
\end{itemize}
\end{proof}

The diagonal entries \(h_{11},\dots,h_{nn}\) of the hat matrix \(H\) are the \textbf{leverages} of the \(n\) observations.
Under (A4), with \(s^2\) of Theorem~\ref{r11-thm-sigma} (so \(n>k+1\)), the numbers \(e_i/(s\sqrt{1-h_{ii}})\), defined
when \(s>0\) and \(h_{ii}<1\), are the \textbf{standardized residuals}~\citeat{tu-clm}{Definitions~2.13, 3.11}.

\begin{example}[cf.\ \citeat{searle-khuri}{Example~12.6}]\label{r11-exm-sigma}
\begin{enumerate}
\item In Example~\ref{r11-exm-running}, \(\mathrm{SSE}=4\), \(n=4\) and \(r=k+1=3\), so \(s^2=4/(4-3)=4\). Replacing
\(\sigma^2\) by \(s^2\) in Theorem~\ref{r11-thm-unbiased}(b) estimates \(\var(\hat\beta)\) by
\[
s^2(X'X)^{-1}=4\cdot\frac14\begin{pmatrix}6&-1&-2\\-1&1&0\\-2&0&1\end{pmatrix}=\begin{pmatrix}6&-1&-2\\-1&1&0\\-2&0&1\end{pmatrix}
\]
with \((X'X)^{-1}\) of Example~\ref{r11-exm-running}: for instance \(\var(\hat\beta_1)\) is estimated by \(1\).
\item In Example~\ref{r11-exm-oneway}, \(e=y-\hat y=(5-6,\ 7-6,\ 6-6,\ 3-4,\ 5-4,\ 2-2)'=(-1,1,0,-1,1,0)'\), so
\(\mathrm{SSE}=4\), and \(r=3\): \(\hat\sigma^2=4/(6-3)=\frac43\).
\end{enumerate}
\end{example}

\subsection*{Sums of squares and \(R^2\)}

\begin{definition}[\citeat{tu-clm}{Definitions~3.1, 6.6}, \citeat{searle-khuri}{\S11.1.2}]\label{r11-def-sums-of-squares}
For a least-squares fit, with fitted values \(\hat y\) and residuals \(e=y-\hat y\), the \textbf{total},
\textbf{regression} and \textbf{error sums of squares} are
\[
\mathrm{SST}=\sum_{i=1}^n(y_i-\bar y)^2,\qquad\mathrm{SSR}=\sum_{i=1}^n(\hat y_i-\bar y)^2,\qquad
\mathrm{SSE}=\sum_{i=1}^n(y_i-\hat y_i)^2=e'e
\]
If \(\one_n\in\mathcal C(X)\) and \(\mathrm{SST}>0\), the \textbf{coefficient of determination} is
\(R^2=\mathrm{SSR}/\mathrm{SST}\).
\end{definition}

\needspace{10\baselineskip}
\begin{theorem}[\citeat{tu-clm}{Theorem~6.4, Proposition~6.7}]\label{r11-thm-anova}
If \(\one_n\in\mathcal C(X)\) (for instance, if the model has an intercept), then
\(\mathrm{SST}=\mathrm{SSR}+\mathrm{SSE}\), where \(\mathrm{SST}=\lVert C_ny\rVert^2=y'C_ny\) and
\(\mathrm{SSR}=\lVert C_n\hat y\rVert^2\). Hence, if \(\mathrm{SST}>0\), then \(R^2=1-\mathrm{SSE}/\mathrm{SST}\) and
\(0\leqslant R^2\leqslant1\).
\end{theorem}
\begin{proof}
\begin{itemize}
\item For \(v\in\R^n\), \(C_nv=v-\bar v\one\), so \(\lVert C_nv\rVert^2=\sum_i(v_i-\bar v)^2\), and
\(\lVert C_nv\rVert^2=v'C_n'C_nv=v'C_nv\) (Proposition~\ref{r1-prop-centering}). With \(v=y\) this is SST; with
\(v=\hat y\), whose mean is \(\bar y\) because \(\one'\hat y=\one'y-\one'e=\one'y\) (Proposition~\ref{r10-prop-residuals}(b)),
it is SSR.
\item \(C_ne=e-\bar e\one=e\), as \(\bar e=0\) (Proposition~\ref{r10-prop-residuals}(b)). So
\(C_ny=C_n\hat y+C_ne=C_n\hat y+e\).
\item \((C_n\hat y)'e=\hat y'C_ne=\hat y'e=0\) (Proposition~\ref{r10-prop-residuals}(c)). By Pythagoras
(Proposition~\ref{r6-prop-pythagoras}),
\(\lVert C_ny\rVert^2=\lVert C_n\hat y\rVert^2+\lVert e\rVert^2\), that is, \(\mathrm{SST}=\mathrm{SSR}+\mathrm{SSE}\).
\item Dividing by \(\mathrm{SST}>0\) gives \(R^2=1-\mathrm{SSE}/\mathrm{SST}\); and
\(0\leqslant\mathrm{SSR}\leqslant\mathrm{SST}\), as \(\mathrm{SSE}\geqslant0\).\qedhere
\end{itemize}
\end{proof}
In terms of projections: \(y=\bar y\one+C_ny\) along \(\Span\{\one\}\oplus\Span\{\one\}^\perp\)
(Example~\ref{r6-exm-centering}), and \(C_ny=C_n\hat y+e\) splits the centered response into
\(C_n\hat y=\hat y-\bar y\one\in\mathcal C(X)\) and \(e\in\mathcal C(X)^\perp\); \(\mathrm{SST}=\mathrm{SSR}+\mathrm{SSE}\) is
Pythagoras for this split~\citeat{tu-clm}{Remark~6.5}. For the design \(X=\one_n\) alone, \(\hat y=\bar y\one\), so \(\mathrm{SSR}=0\) and
\(\mathrm{SSE}=\mathrm{SST}\); \(R^2\) is the fraction of SST that the other columns of \(X\) explain.

For instance, \(\mathrm{SST}=y'C_ny=0\) only if \(C_ny=0\), and \(\mathrm{SSE}=y'My=0\) for the residual maker \(M\) of
Theorem~\ref{r10-thm-PX} only if \(My=e=0\) (Corollary~\ref{r8-cor-psd-null}): both matrices are symmetric and idempotent
(Proposition~\ref{r1-prop-centering}, Theorem~\ref{r10-thm-PX}(a)), hence PSD (Proposition~\ref{r8-prop-pd-basic}(d)). Without
centering, \(y=\bar y\one_n+C_ny\) with \(\one_n'C_ny=0\) (Proposition~\ref{r1-prop-centering}), so Pythagoras
(Proposition~\ref{r6-prop-pythagoras}) and
Theorem~\ref{r11-thm-anova} give the partition of the uncorrected sum of squares~\citeat{searle-khuri}{\S12.7}
\[
y'y=\lVert\bar y\one_n\rVert^2+\lVert C_ny\rVert^2=n\bar y^2+\mathrm{SSR}+\mathrm{SSE}
\]

\begin{example}\label{r11-exm-anova}
In Example~\ref{r11-exm-running}, \(\bar y=\frac14(9+3+5+3)=5\), so
\[
C_4y=\begin{pmatrix}4\\-2\\0\\-2\end{pmatrix},\qquad
C_4\hat y=\begin{pmatrix}3\\-1\\1\\-3\end{pmatrix},\qquad
e=\begin{pmatrix}1\\-1\\-1\\1\end{pmatrix},\qquad C_4y=C_4\hat y+e
\]
Hence \(\mathrm{SST}=16+4+0+4=24\), \(\mathrm{SSR}=9+1+1+9=20\) and \(\mathrm{SSE}=1+1+1+1=4\); indeed
\(24=20+4\), and \(R^2=\frac{20}{24}=\frac56\). Without centering, \(y'y=81+9+25+9=124\) and
\(n\bar y^2=4\cdot25=100\), so \(124=100+20+4\).
\end{example}

\begin{theorem}[\citeat{tu-clm}{Theorem~6.8}]\label{r11-thm-r2-monotone}
Let \(X\in\R^{n\times p}\) with \(\one_n\in\mathcal C(X)\), let \(x\in\R^n\), let \(X_+=[X\ x]\) be \(X\) with the column
\(x\) added, and let \(y\in\R^n\) with \(\mathrm{SST}>0\). Let SSE and \(R^2\) belong to the least-squares fit with design
\(X\), and \(\mathrm{SSE}_+\) and \(R^2_+\) to the fit with design \(X_+\). Then \(\mathrm{SSE}_+\leqslant\mathrm{SSE}\) and
\(R^2_+\geqslant R^2\): adding a column cannot decrease \(R^2\).
\end{theorem}
\begin{proof}
\begin{itemize}
\item Every column of \(X\) is a column of \(X_+\), so \(\mathcal C(X)\subseteq\mathcal C(X_+)\), and both contain \(\one_n\).
\item The fitted values \(\hat y_+\) of \(X_+\) form the vector of \(\mathcal C(X_+)\) closest to \(y\)
(Theorems~\ref{r10-thm-least-squares}(b) and~\ref{r6-thm-best-approximation}), and the fitted values
\(\hat y\) of \(X\) lie in \(\mathcal C(X)\subseteq\mathcal C(X_+)\). Hence
\(\mathrm{SSE}_+=\lVert y-\hat y_+\rVert^2\leqslant\lVert y-\hat y\rVert^2=\mathrm{SSE}\).
\item SST does not depend on the design, so by Theorem~\ref{r11-thm-anova},
\(R^2_+=1-\mathrm{SSE}_+/\mathrm{SST}\geqslant1-\mathrm{SSE}/\mathrm{SST}=R^2\).\qedhere
\end{itemize}
\end{proof}
So \(R^2\) always favours the design with more columns; the adjusted \(R^2\) below charges for the number of columns.

\begin{definition}[\citeat{tu-clm}{Definition~6.6}]\label{r11-def-adjusted}
Let \(\one_n\in\mathcal C(X)\), \(\rank X=k+1<n\) and \(\mathrm{SST}>0\). The \textbf{adjusted \(R^2\)} is
\[
\bar R^2=1-\frac{\mathrm{SSE}/(n-k-1)}{\mathrm{SST}/(n-1)}=1-\frac{n-1}{n-k-1}\,(1-R^2)
\]
\end{definition}
The second form follows from \(\mathrm{SSE}/\mathrm{SST}=1-R^2\) (Theorem~\ref{r11-thm-anova}). The numerator
\(\mathrm{SSE}/(n-k-1)\) is the unbiased estimate \(s^2\) of \(\sigma^2\) (Theorem~\ref{r11-thm-sigma}), and \(\mathrm{SST}/(n-1)\)
is the sample variance of \(y\). For \(k\geqslant1\), \(\frac{n-1}{n-k-1}>1\), so \(\bar R^2<R^2\) unless \(R^2=1\), and \(\bar R^2\)
can be negative: \(R^2=0\) gives \(\bar R^2=1-\frac{n-1}{n-k-1}<0\). In Example~\ref{r11-exm-anova}, \(n=4\) and \(k=2\), so
\(\bar R^2=1-\frac{4/1}{24/3}=\frac12\), while \(R^2=\frac56\).

\begin{proposition}[\citeat{tu-clm}{Proposition~6.12}]\label{r11-prop-r2-corr}
Let \(\one_n\in\mathcal C(X)\) and \(\mathrm{SST}>0\), with \(\hat y\), SSR and \(R^2\) as in
Definition~\ref{r11-def-sums-of-squares}, and \(S_{uv}\), constant vectors and \(\cor(u,v)\) as in
Definition~\ref{r12-def-correlation}.
\begin{enumerate}
\item If \(\mathrm{SSR}>0\), then \(\hat y\) is not constant and \(R^2=\cor(y,\hat y)^2\).
\item \(\cor(y,v)^2\leqslant R^2\) for every \(v\in\mathcal C(X)\) that is not constant.
\end{enumerate}
\end{proposition}
\begin{proof}
First let \(v\in\mathcal C(X)\). Then \(C_nv=v-\bar v\one\in\mathcal C(X)\), as \(\one\in\mathcal C(X)\), so \(e'C_nv=0\)
(Theorem~\ref{r10-thm-least-squares}(b)). Moreover \(C_ne=e\), as \(\bar e=0\) (Proposition~\ref{r10-prop-residuals}(b)), so
\(C_ny=C_n\hat y+C_ne=C_n\hat y+e\). Hence
\[
S_{yv}=(C_ny)'(C_nv)=(C_n\hat y)'(C_nv)+e'C_nv=(C_n\hat y)'(C_nv)=S_{\hat yv}
\]
(a)
\begin{itemize}
\item By Theorem~\ref{r11-thm-anova}, \(S_{\hat y\hat y}=\lVert C_n\hat y\rVert^2=\mathrm{SSR}>0\), so \(\hat y\) is not constant;
and \(S_{yy}=\lVert C_ny\rVert^2=\mathrm{SST}\).
\item The first step with \(v=\hat y\in\mathcal C(X)\) gives \(S_{y\hat y}=S_{\hat y\hat y}=\mathrm{SSR}\). So
\(\cor(y,\hat y)^2=\mathrm{SSR}^2/(\mathrm{SST}\cdot\mathrm{SSR})=\mathrm{SSR}/\mathrm{SST}=R^2\).
\end{itemize}
(b) By the first step and the Cauchy--Schwarz inequality (Theorem~\ref{r6-thm-cauchy-schwarz}(a)) for \(C_n\hat y\) and
\(C_nv\), \(S_{yv}^2=\bigl((C_n\hat y)'(C_nv)\bigr)^2\leqslant\lVert C_n\hat y\rVert^2\lVert C_nv\rVert^2=\mathrm{SSR}\cdot S_{vv}\);
so \(\cor(y,v)^2=S_{yv}^2/(\mathrm{SST}\cdot S_{vv})\leqslant\mathrm{SSR}/\mathrm{SST}=R^2\).
\end{proof}
Thus, when \(\mathrm{SSR}>0\), \(\hat y\) is the vector of \(\mathcal C(X)\), that is, the linear combination of the columns
of \(X\), most highly correlated with \(y\), and \(R^2\) is the square of that correlation.

\begin{example}\label{r11-exm-r2-corr}
In Example~\ref{r11-exm-anova}, \(C_4y=(4,-2,0,-2)'\), \(C_4\hat y=(3,-1,1,-3)'\), \(\mathrm{SST}=24\) and \(\mathrm{SSR}=20\).
\begin{itemize}
\item \(S_{y\hat y}=12+2+0+6=20=\mathrm{SSR}\), so \(\cor(y,\hat y)^2=\frac{20^2}{24\cdot20}=\frac56=R^2\).
\item For the column \(v=X_{\cdot1}\) of \(X\), \(C_4v=(1,-1,1,-1)'\) (Example~\ref{r1-exm-data}), \(S_{yv}=4+2+0+2=8\) and
\(S_{vv}=4\), so \(\cor(y,v)^2=\frac{8^2}{24\cdot4}=\frac23\). For \(v=X_{\cdot2}\), \(C_4v=(1,1,-1,-1)'\),
\(S_{yv}=4-2+0+2=4\) and \(S_{vv}=4\), so \(\cor(y,v)^2=\frac{4^2}{24\cdot4}=\frac16\). Both are at most
\(R^2=\frac56\), as Proposition~\ref{r11-prop-r2-corr}(b) says.
\end{itemize}
\end{example}

\needspace{18\baselineskip}
\subsection*{Partitioned regression}

\begin{theorem}[Frisch--Waugh--Lovell; \citeat{tu-clm}{Theorem~8.1}]\label{r11-thm-fwl}
Assume (A4). Partition \(X=[X_1\ X_2]\) with \(X_1\in\R^{n\times p}\) and \(X_2\in\R^{n\times q}\) (\(p+q=k+1\)), and
accordingly \(\beta=(\gamma',\delta')'\) and \(\hat\beta=(X'X)^{-1}X'y=(\hat\gamma',\hat\delta')'\) with \(\gamma,\hat\gamma\in\R^p\),
\(\delta,\hat\delta\in\R^q\); so \(X\beta=X_1\gamma+X_2\delta\). Let \(H_2=X_2(X_2'X_2)^{-1}X_2'\) and \(M_2=I_n-H_2\) be the hat
matrix and the residual maker of \(X_2\).
\begin{enumerate}
\item \(X_1'M_2X_1\) is invertible and \(\hat\gamma=(X_1'M_2X_1)^{-1}X_1'M_2y\).
\item The least-squares fit of the response \(M_2y\) on the design \(M_2X_1\) has the estimate \(\hat\gamma\), and its residual
vector is the residual vector \(e\) of the full fit.
\end{enumerate}
With the roles of \(X_1\) and \(X_2\) exchanged, (a) and (b) hold for \(\hat\delta\) and the residual maker \(M_1\) of \(X_1\); in
particular \(\hat\delta=(X_2'M_1X_2)^{-1}X_2'M_1y\).
\end{theorem}
\begin{proof}
The columns of \(X_1\) and of \(X_2\) are columns of \(X\), hence independent (Proposition~\ref{r3-prop-dependence}(b)); so
\(H_2\), \(M_2\) and the hat matrix and the residual maker of \(X_1\) are as in Theorem~\ref{r6-thm-hat}.

(a)
\begin{itemize}
\item By the block product (Theorem~\ref{r5b-thm-block}), the normal equations \(X'X\hat\beta=X'y\) read
\[
X_1'X_1\hat\gamma+X_1'X_2\hat\delta=X_1'y,\qquad X_2'X_1\hat\gamma+X_2'X_2\hat\delta=X_2'y
\]
\item The second gives \(\hat\delta=(X_2'X_2)^{-1}X_2'(y-X_1\hat\gamma)\). Substituting into the first, with
\(X_1'X_2(X_2'X_2)^{-1}X_2'=X_1'H_2=X_1'-X_1'M_2\),
\(X_1'X_1\hat\gamma+X_1'(y-X_1\hat\gamma)-X_1'M_2(y-X_1\hat\gamma)=X_1'y\). The first two terms add up to \(X_1'y\), so
\(X_1'M_2(y-X_1\hat\gamma)=0\), that is, \(X_1'M_2X_1\hat\gamma=X_1'M_2y\).
\item \(M_2X_1\) has full column rank: if \(M_2X_1v=0\), then \(X_1v=H_2X_1v=X_2w\) with \(w=(X_2'X_2)^{-1}X_2'X_1v\), so
\(X(v',-w')'=X_1v-X_2w=0\), and (A4) gives \(v=0\) (Theorem~\ref{thm:r4-rank-rules}(b)). As \(M_2\) is symmetric and
idempotent (Theorem~\ref{r6-thm-hat}(b)), \(X_1'M_2X_1=(M_2X_1)'(M_2X_1)\), which is therefore invertible
(Theorem~\ref{r6-thm-hat}); multiplying by its inverse gives the formula for \(\hat\gamma\).
\end{itemize}
(b)
\begin{itemize}
\item By (a), \((M_2X_1)'M_2X_1=X_1'M_2X_1\) is invertible, so the fit of \(M_2y\) on \(M_2X_1\) has exactly one least-squares
estimate (Theorem~\ref{r10-thm-least-squares}(c)); as \(M_2'M_2=M_2\), it is
\[
\bigl((M_2X_1)'M_2X_1\bigr)^{-1}(M_2X_1)'M_2y=(X_1'M_2X_1)^{-1}X_1'M_2y=\hat\gamma
\]
\item Its residual vector is \(M_2y-M_2X_1\hat\gamma=M_2(y-X_1\hat\gamma)=M_2(X_2\hat\delta+e)=M_2e\), since
\(y=X_1\hat\gamma+X_2\hat\delta+e\) and \(M_2X_2=0\).
\item \(X_2'e=0\) (Proposition~\ref{r10-prop-residuals}(a)), so \(H_2e=0\) and \(M_2e=e\).
\end{itemize}
The normal equations are the same two block equations with \((X_1,\hat\gamma)\) and \((X_2,\hat\delta)\) exchanged, so the same
steps give the last statement.
\end{proof}
Here \(X_1'M_2X_1=X_1'X_1-X_1'X_2(X_2'X_2)^{-1}X_2'X_1\) is the Schur complement of \(X_2'X_2\) in
\(X'X=\left(\begin{smallmatrix}X_1'X_1&X_1'X_2\\X_2'X_1&X_2'X_2\end{smallmatrix}\right)\) (the paragraph after
Theorem~\ref{r5b-thm-schur}): solving the second block equation for \(\hat\delta\) and substituting it into the first is
block elimination. The vectors \(M_2y\) and \(M_2X_1\) are the residuals of \(y\) and of the columns of \(X_1\) after the
least-squares fit on \(X_2\) (Theorem~\ref{r10-thm-least-squares}(c)); so, by (b), \(\hat\gamma\) measures the effect of
\(X_1\) on \(y\) after both are adjusted for \(X_2\)~\citeat{tu-clm}{Remark~8.5}.

\begin{example}\label{r11-exm-fwl}
Let \(X=[\one_n\ X_0]\in\R^{n\times(k+1)}\) satisfy (A4), and let \(\bar x=\frac1nX_0'\one\) be the mean vector of the columns
of \(X_0\) (Example~\ref{r1-exm-data}). In Theorem~\ref{r11-thm-fwl} take \(X_1=\one_n\) and \(X_2=X_0\), so that
\(\gamma=\beta_0\) and \(\delta=(\beta_1,\dots,\beta_k)'\) holds the slopes. The residual maker of \(\one_n\) is \(M_1=C_n\)
(Example~\ref{r6-exm-centering}), so
\[
\hat\delta=(X_0'C_nX_0)^{-1}X_0'C_ny
\]
the least-squares estimate of the fit of the centered response \(C_ny\) on the centered data \(C_nX_0\); and \(X_0'C_nX_0\) is
the Schur complement of \(n\) in \(X'X\) (Proposition~\ref{r12-prop-xtx-inverse}(a)). The first normal equation,
\(n\hat\beta_0+\one'X_0\hat\delta=\one'y\), gives \(\hat\beta_0=\bar y-\bar x'\hat\delta\).

In Example~\ref{r11-exm-running}, \(X_0'C_4X_0=4I_2\) and \(C_4X_0\) has the columns \((1,-1,1,-1)'\) and \((1,1,-1,-1)'\)
(Example~\ref{r1-exm-data}), and \(C_4y=(4,-2,0,-2)'\) (Example~\ref{r11-exm-anova}). So
\[
X_0'C_4y=(C_4X_0)'(C_4y)=\begin{pmatrix}4+2+0+2\\4-2+0+2\end{pmatrix}=\begin{pmatrix}8\\4\end{pmatrix},\qquad
\hat\delta=\frac14\begin{pmatrix}8\\4\end{pmatrix}=\begin{pmatrix}2\\1\end{pmatrix}
\]
and \(\hat\beta_0=5-(1\cdot2+2\cdot1)=1\), with \(\bar y=5\) and \(\bar x=(1,2)'\): this is \(\hat\beta=(1,2,1)'\) of
Example~\ref{r11-exm-running}. The residuals of the fit of \(C_4y\) on \(C_4X_0\) are
\(C_4y-C_4X_0\hat\delta=(4,-2,0,-2)'-(3,-1,1,-3)'=(1,-1,-1,1)'=e\), as Theorem~\ref{r11-thm-fwl}(b) says.

The centered design \(\tilde X=[\one_n\ C_nX_0]\) gives the same fit: \(C_nX_0=X_0-\one\bar x'\) (Example~\ref{r1-exm-data}), so
\(\tilde X=XA\) with the invertible \(A=\left(\begin{smallmatrix}1&-\bar x'\\0&I_k\end{smallmatrix}\right)\) (Theorem~\ref{r5b-thm-block}),
and \(\tilde X\) has the fitted values of \(X\) (Proposition~\ref{r11-prop-reparametrization}). As \(\one'C_n=0\) and
\(C_n'C_n=C_n\) (Proposition~\ref{r1-prop-centering}), \(\tilde X'\tilde X=\left(\begin{smallmatrix}n&0\\0&X_0'C_nX_0\end{smallmatrix}\right)\).
In Example~\ref{r11-exm-running}, \(\tilde X'\tilde X=4I_3\) and \(\tilde X'y=(20,8,4)'\), so the estimate for \(\tilde X\) is
\(\frac14(20,8,4)'=(5,2,1)'=(\bar y,\hat\delta')'\), and \(A(5,2,1)'=(5-(1\cdot2+2\cdot1),\ 2,\ 1)'=(1,2,1)'=\hat\beta\).
\end{example}

\needspace{16\baselineskip}
\begin{corollary}[\citeat{tu-clm}{Corollaries~8.10, 8.12}]\label{r11-cor-vif}
Assume (A3) and (A4).
\begin{enumerate}
\item In Theorem~\ref{r11-thm-fwl}, \(\var(\hat\gamma)=\sigma^2(X_1'M_2X_1)^{-1}\) and
\(\var(\hat\delta)=\sigma^2(X_2'M_1X_2)^{-1}\).
\item Let \(X=[\one_n\ X_0]\) with \(k\geqslant2\), and let \(1\leqslant j\leqslant k\). Let \(X_{(j)}\) be \(X\) with the column
\(X_{\cdot j}\) deleted, \(M_{(j)}\) its residual maker, \(\tilde x_j=M_{(j)}X_{\cdot j}\), \(S_{jj}=\sum_{i=1}^n(x_{ij}-\bar x_j)^2\),
and \(R_j^2\) the \(R^2\) of the least-squares fit of the response \(X_{\cdot j}\) on the design \(X_{(j)}\). Then
\(\tilde x_j\ne0\), \(S_{jj}>0\), \(0\leqslant R_j^2<1\), and
\[
\var(\hat\beta_j)=\frac{\sigma^2}{\lVert\tilde x_j\rVert^2}=\frac{\sigma^2}{S_{jj}(1-R_j^2)}=\frac{\sigma^2}{S_{jj}}\,\mathrm{VIF}_j,
\qquad\text{where}\qquad\mathrm{VIF}_j=\frac1{1-R_j^2}\geqslant1
\]
\end{enumerate}
The number \(\mathrm{VIF}_j\) is the \textbf{variance inflation factor} of the \(j\)th variable.
\end{corollary}
\begin{proof}
(a) By Theorem~\ref{r11-thm-fwl}(a), \(\hat\gamma=Ly\) with \(L=(X_1'M_2X_1)^{-1}X_1'M_2\). The matrix \((X_1'M_2X_1)^{-1}\) is
symmetric (Theorem~\ref{r2-thm-inverse-props}(c)), and \(M_2M_2'=M_2\) (Theorem~\ref{r6-thm-hat}(b)); so by
Theorem~\ref{r11-thm-linear}(b) and (A3),
\[
\var(\hat\gamma)=\sigma^2LL'=\sigma^2(X_1'M_2X_1)^{-1}X_1'M_2X_1(X_1'M_2X_1)^{-1}=\sigma^2(X_1'M_2X_1)^{-1}
\]
With the roles of the blocks exchanged, the same steps give \(\var(\hat\delta)\).

(b)
\begin{itemize}
\item Let \(P\in\R^{(k+1)\times(k+1)}\) have as its columns the columns of \(I_{k+1}\), column \(j+1\) first and the others
in their order, so that \(XP=[X_{\cdot j}\ X_{(j)}]\) (Theorem~\ref{r1-thm-views}(a),(b)). The columns of \(P\) are distinct
unit vectors, so \(P'P=I_{k+1}\) (Theorem~\ref{r1-thm-views}(d)) and \(PP'=I_{k+1}\) (Theorem~\ref{r2-thm-invertible}).
\item \(XP\) has full column rank (\(\rank(XP)=\rank X\), Theorem~\ref{thm:r4-rank-product}(d)), and
\((XP)'(XP)P'\hat\beta=P'X'X\hat\beta=P'X'y=(XP)'y\). So \(P'\hat\beta\) is the least-squares estimate for \(XP\)
(Theorem~\ref{r10-thm-normal-equations}(b)); its first entry is \(\hat\beta_j\).
\item Theorem~\ref{r11-thm-fwl}(a) for \(XP\), with \(X_1=X_{\cdot j}\) and \(X_2=X_{(j)}\), gives
\(X_{\cdot j}'M_{(j)}X_{\cdot j}\ne0\) and \(\hat\beta_j=(X_{\cdot j}'M_{(j)}X_{\cdot j})^{-1}X_{\cdot j}'M_{(j)}y\). As \(M_{(j)}\) is
symmetric and idempotent, \(X_{\cdot j}'M_{(j)}X_{\cdot j}=\lVert\tilde x_j\rVert^2\). So \(\tilde x_j\ne0\), and (a) gives
\(\var(\hat\beta_j)=\sigma^2/\lVert\tilde x_j\rVert^2\).
\item If \(X_{\cdot j}\) were constant, \(X_{\cdot j}-\bar x_j\one_n=0\) would be a linear relation among the columns of \(X\),
which (A4) excludes (Theorem~\ref{thm:r4-rank-rules}(b)). So \(S_{jj}=\lVert C_nX_{\cdot j}\rVert^2>0\)
(Definition~\ref{r12-def-correlation}).
\item \(X_{(j)}\) has full column rank (Proposition~\ref{r3-prop-dependence}(b)) and contains the column \(\one_n\). In the fit
of \(X_{\cdot j}\) on \(X_{(j)}\), the
total sum of squares is \(S_{jj}\) and the residual vector is \(M_{(j)}X_{\cdot j}=\tilde x_j\)
(Theorem~\ref{r10-thm-least-squares}(c)). By Theorem~\ref{r11-thm-anova}, \(0\leqslant R_j^2=1-\lVert\tilde x_j\rVert^2/S_{jj}\), which
is \(<1\) as \(\tilde x_j\ne0\). Hence \(\lVert\tilde x_j\rVert^2=S_{jj}(1-R_j^2)\), and \(\mathrm{VIF}_j\geqslant1\).\qedhere
\end{itemize}
\end{proof}
So \(\var(\hat\beta_j)\) is large when \(\sigma^2\) is large, when the \(j\)th variable varies little (small
\(S_{jj}\)), or when it is close to a linear combination of the other columns (\(R_j^2\) close to \(1\)); \(\mathrm{VIF}_j=1\)
exactly when \(R_j^2=0\).

\begin{example}\label{r11-exm-vif}
Let \(X=[\one_4\ X_0]\), where the columns of \(X_0\) are \(X_{\cdot1}=(4,4,2,2)'\) and \(X_{\cdot2}=(3,2,2,1)'\). Then
\(\bar x=(3,2)'\) and, by Example~\ref{r1-exm-centering} and Theorem~\ref{r5d-thm-adjugate},
\[
C_4X_0=\begin{pmatrix}1&1\\1&0\\-1&0\\-1&-1\end{pmatrix},\qquad
X_0'C_4X_0=(C_4X_0)'(C_4X_0)=\begin{pmatrix}4&2\\2&2\end{pmatrix},\qquad
(X_0'C_4X_0)^{-1}=\frac14\begin{pmatrix}2&-2\\-2&4\end{pmatrix}
\]
As \(X_0'C_4X_0\) is invertible, so is \(X'X\) (Proposition~\ref{r12-prop-xtx-inverse}(a)), and \(X\) satisfies (A4)
(Theorems~\ref{thm:r4-rank-ata}(b) and~\ref{thm:r4-rank-rules}(d)).
\begin{itemize}
\item With \(X_1=\one_4\), \(X_2=X_0\) as in Example~\ref{r11-exm-fwl}, Corollary~\ref{r11-cor-vif}(a) gives
\(\var\bigl((\hat\beta_1,\hat\beta_2)'\bigr)=\sigma^2(X_0'C_4X_0)^{-1}\): \(\var(\hat\beta_1)=\frac12\sigma^2\) and
\(\var(\hat\beta_2)=\sigma^2\).
\item The diagonal of \(X_0'C_4X_0\) holds \(S_{11}=4\) and \(S_{22}=2\). By Corollary~\ref{r11-cor-vif}(b),
\(\mathrm{VIF}_j=S_{jj}\var(\hat\beta_j)/\sigma^2\): \(\mathrm{VIF}_1=4\cdot\frac12=2\) and \(\mathrm{VIF}_2=2\cdot1=2\), so
\(R_1^2=R_2^2=1-\frac12=\frac12\). The correlation between the two variables doubles the variances \(\sigma^2/S_{jj}\).
\item In Example~\ref{r11-exm-running}, by contrast, \(X_0'C_4X_0=4I_2\) (Example~\ref{r1-exm-data}), so
\(\var(\hat\beta_1)=\var(\hat\beta_2)=\frac14\sigma^2=\sigma^2/S_{jj}\): there
\(\mathrm{VIF}_1=\mathrm{VIF}_2=1\) and \(R_1^2=R_2^2=0\).
\end{itemize}
\end{example}

\needspace{12\baselineskip}
\begin{theorem}[\citeat{tu-clm}{Theorem~8.3}]\label{r11-thm-ovb}
Let \(X_1\in\R^{n\times p}\) have full column rank, let \(X_2\in\R^{n\times q}\), and suppose that
\(y=X_1\gamma+X_2\delta+\varepsilon\) with \(\expc(\varepsilon)=0\). Let \(\tilde\gamma=(X_1'X_1)^{-1}X_1'y\) be the
least-squares estimate for the design \(X_1\) alone, which omits \(X_2\). Then
\[
\expc(\tilde\gamma)=\gamma+(X_1'X_1)^{-1}X_1'X_2\delta
\]
The \textbf{omitted variable bias} \((X_1'X_1)^{-1}X_1'X_2\delta\) is \(0\) if and only if \(X_1'X_2\delta=0\), for instance when
\(X_1'X_2=0\) or \(\delta=0\); and it is \(0\) for every \(\delta\in\R^q\) if and only if \(X_1'X_2=0\).
\end{theorem}
\begin{proof}
\begin{itemize}
\item \(X_1'X_1\) is invertible (Theorem~\ref{r6-thm-hat}). Substituting \(y\),
\(\tilde\gamma=(X_1'X_1)^{-1}X_1'(X_1\gamma+X_2\delta+\varepsilon)=\gamma+(X_1'X_1)^{-1}X_1'X_2\delta+(X_1'X_1)^{-1}X_1'\varepsilon\).
\item By Theorem~\ref{r11-thm-linear}(b) and \(\expc(\varepsilon)=0\),
\(\expc(\tilde\gamma)=\gamma+(X_1'X_1)^{-1}X_1'X_2\delta\).
\item As \((X_1'X_1)^{-1}\) is invertible, the bias is \(0\) if and only if \(X_1'X_2\delta=0\).
\item If \(X_1'X_2\delta=0\) for every \(\delta\), then taking for \(\delta\) the columns of \(I_q\) gives \(X_1'X_2=0\); conversely,
\(X_1'X_2=0\) gives \(X_1'X_2\delta=0\) for every \(\delta\).\qedhere
\end{itemize}
\end{proof}

\begin{example}\label{r11-exm-ovb}
In Example~\ref{r11-exm-running}, fit the two variables without the intercept: \(X_1=X_0\), \(X_2=\one_4\),
\(\gamma=(\beta_1,\beta_2)'\) and \(\delta=\beta_0\). By Example~\ref{r1-exm-data}, \(X_0'\one=(4,8)'\) and
\[
X_0'X_0=\begin{pmatrix}8&8\\8&20\end{pmatrix},\qquad
(X_0'X_0)^{-1}=\frac1{96}\begin{pmatrix}20&-8\\-8&8\end{pmatrix},\qquad
(X_0'X_0)^{-1}X_0'\one=\frac1{96}\begin{pmatrix}80-64\\-32+64\end{pmatrix}=\begin{pmatrix}\frac16\\\frac13\end{pmatrix}
\]
(\(\det X_0'X_0=160-64=96\); Theorem~\ref{r5d-thm-adjugate}). By Theorem~\ref{r11-thm-ovb}, the slopes of the fit without the
intercept have the expectations \(\beta_1+\frac16\beta_0\) and \(\beta_2+\frac13\beta_0\): they are biased unless \(\beta_0=0\).
With centered variables, \(X_0'\one=0\), there would be no bias for any \(\beta_0\).
\end{example}

\needspace{14\baselineskip}
\subsection*{Best linear prediction}

\begin{theorem}[Best linear predictor]\label{r11-thm-blp}
Let \(y\) be a random variable and \(x\in\R^k\) a random vector, with \(\expc y=\mu_y\), \(\expc x=\mu_x\),
\(\var(y)=\sigma_{yy}\), \(\cov(x,y)=\sigma_{xy}\in\R^k\), \(\sigma_{yx}=\sigma_{xy}'\), and
\(\var(x)=\Sigma_{xx}\) PD. Put
\[
b^*=\Sigma_{xx}^{-1}\sigma_{xy},\qquad a^*=\mu_y-b^{*\prime}\mu_x
\]
Then for all \(a\in\R\) and \(b\in\R^k\),
\[
\expc(y-a-b'x)^2=\bigl(\sigma_{yy}-\sigma_{yx}\Sigma_{xx}^{-1}\sigma_{xy}\bigr)+(b-b^*)'\Sigma_{xx}(b-b^*)
+(\mu_y-a-b'\mu_x)^2
\]
So the \textbf{mean squared prediction error} \(\expc(y-a-b'x)^2\) is smallest exactly for \(a=a^*\), \(b=b^*\), with
minimum \(\sigma_{yy}-\sigma_{yx}\Sigma_{xx}^{-1}\sigma_{xy}\): the Schur complement of \(\Sigma_{xx}\) in
\(\var\bigl((x',y)'\bigr)=\left(\begin{smallmatrix}\Sigma_{xx}&\sigma_{xy}\\\sigma_{yx}&\sigma_{yy}\end{smallmatrix}\right)\)
(Theorem~\ref{r5b-thm-schur}). The random variable \(a^*+b^{*\prime}x\) is the \textbf{best linear
predictor} of \(y\) from \(x\), and the prediction error \(u=y-a^*-b^{*\prime}x\) has mean \(0\) and is
uncorrelated with \(x\).
\end{theorem}
\begin{proof}
\begin{itemize}
\item Let \(v=y-a-b'x\). By Theorem~\ref{r14-thm-cov}(e) for \(p=1\), \(\expc v^2=\var(v)+(\expc v)^2\), and
\(\expc v=\mu_y-a-b'\mu_x\) (Theorem~\ref{r11-thm-linear}(b)).
\item \(v=(1,-b')(y,x')'-a\), and \(\var\bigl((y,x')'\bigr)\) has the blocks \(\sigma_{yy}\), \(\sigma_{yx}\),
\(\sigma_{xy}\), \(\Sigma_{xx}\); by Theorem~\ref{r11-thm-linear}(b),
\(\var(v)=\sigma_{yy}-2b'\sigma_{xy}+b'\Sigma_{xx}b\).
\item As \(\Sigma_{xx}b^*=\sigma_{xy}\),
\((b-b^*)'\Sigma_{xx}(b-b^*)=b'\Sigma_{xx}b-2b'\sigma_{xy}+\sigma_{yx}\Sigma_{xx}^{-1}\sigma_{xy}\); substituting
gives the identity.
\item \(\Sigma_{xx}\) is PD, so the middle term is \(\geqslant0\), with equality only for
\(b=b^*\); the last term is \(\geqslant0\), with equality only for \(a=\mu_y-b'\mu_x\). Hence the minimum is
attained exactly at \(b=b^*\), \(a=a^*\).
\item As \(a^*=\mu_y-b^{*\prime}\mu_x\), \(u=y-a^*-b^{*\prime}x=(y-\mu_y)-b^{*\prime}(x-\mu_x)\), so \(\expc u=0\)
(Theorem~\ref{r11-thm-linear}(b)). By Theorem~\ref{r14-thm-cov}(a),(b),
\(\cov(x,u)=\cov(x,y)-\cov(x,x)b^*=\sigma_{xy}-\Sigma_{xx}b^*=0\).\qedhere
\end{itemize}
\end{proof}

Under joint normality the best linear predictor is the conditional mean. If \((x',y)'\sim N_{k+1}(\mu,\Sigma)\) with
\(\Sigma_{xx}\) PD, Theorem~\ref{r14-thm-mvn-conditional} with \(z_1=x\) and \(z_2=y\) gives the conditional distribution of \(y\)
given \(x\) as
\[
N\bigl(\mu_y+\sigma_{yx}\Sigma_{xx}^{-1}(x-\mu_x),\ \sigma_{yy}-\sigma_{yx}\Sigma_{xx}^{-1}\sigma_{xy}\bigr)
=N\bigl(a^*+b^{*\prime}x,\ \sigma_{yy}-\sigma_{yx}\Sigma_{xx}^{-1}\sigma_{xy}\bigr)
\]
in the sense of that theorem: the conditional mean is the best linear predictor of Theorem~\ref{r11-thm-blp}, and the
conditional variance is its mean squared prediction error. Likewise, for \(z\sim N_n(\mu,\Sigma)\) partitioned as in
Theorem~\ref{r14-thm-mvn-conditional}, entry \(j\) of \(w=z_2-\mu_2-\Sigma_{21}\Sigma_{11}^{-1}(z_1-\mu_1)\) is the prediction error of
the best linear predictor of entry \(j\) of \(z_2\) from \(z_1\) (row \(j\) of \(\Sigma_{21}\Sigma_{11}^{-1}\) is \(b^{*\prime}\)); so the
partial correlations of the entries of \(z_2\) given \(z_1\) (the paragraph after Theorem~\ref{r14-thm-mvn-conditional}) are the
correlations of these prediction errors. This uses only the covariance matrix.

\begin{example}\label{r11-exm-blp}
Let \((y,x_1,x_2,x_3)'\) have mean \((8,1,2,3)'\) and covariance matrix
\[
\begin{pmatrix}\sigma_{yy}&\sigma_{yx}\\\sigma_{xy}&\Sigma_{xx}\end{pmatrix}
=\left(\begin{array}{c|ccc}12&3&4&3\\\hline3&2&1&0\\4&1&2&1\\3&0&1&2\end{array}\right)
\]
which is PD (leading principal minors \(12\), \(15\), \(10\), \(8\); Theorem~\ref{r8-thm-leading-minors}); so is
\(\Sigma_{xx}\) (leading principal minors \(2\), \(3\), \(4\)), which is therefore invertible (Proposition~\ref{r8-prop-pd-basic}(c)). So \(b^*=\Sigma_{xx}^{-1}\sigma_{xy}\) is the
unique solution of \(\Sigma_{xx}b=\sigma_{xy}\). As
\[
\begin{pmatrix}2&1&0\\1&2&1\\0&1&2\end{pmatrix}\begin{pmatrix}1\\1\\1\end{pmatrix}=\begin{pmatrix}3\\4\\3\end{pmatrix}
\]
\(b^*=(1,1,1)'\) and \(a^*=8-(1+2+3)=2\). By Theorem~\ref{r11-thm-blp} the best linear predictor of \(y\) is
\(2+x_1+x_2+x_3\), with mean squared prediction error \(12-(3,4,3)(1,1,1)'=2\). With fewer predictors the minimum is larger:
\begin{itemize}
\item with \((x_1,x_3)'\), whose covariance matrix is \(2I_2\), it is \(12-(3,3)\frac12I_2(3,3)'=3\);
\item with no predictor, \(\expc(y-a)^2=\var(y)+(\mu_y-a)^2\) is smallest for \(a=\mu_y=8\), with value \(12\).
\end{itemize}
\end{example}

\subsection*{Exercises}

\begin{exercise}\label{r11-dr-hat}
Four observations of two variables and a response are
\[
\begin{array}{c|cccc}x_{i1}&2&0&1&1\\x_{i2}&1&1&0&2\\\hline y_i&7&5&3&5\end{array}
\]
Fit \(y_i=\beta_0+\beta_1x_{i1}+\beta_2x_{i2}+\varepsilon_i\).
\begin{enumerate}
\item Compute \(X'X\), \((X'X)^{-1}\) by row reduction of \([X'X\mid I_3]\), \((X'X)^{-1}X'\) and \(\hat\beta\).
\item Compute the hat matrix \(H\) and its trace, \(\hat y=Hy\) and \(e=y-\hat y\), and check that \(X'e=0\).
\item Compute SST, SSR, SSE and \(R^2\). Check that \(R^2=\cor(y,\hat y)^2\), and compare \(R^2\) with
\(\cor(y,X_{\cdot1})^2\) and \(\cor(y,X_{\cdot2})^2\).
\end{enumerate}
\end{exercise}
\begin{solution}
(a) The \(i\)th row of \(X\) is \((1,x_{i1},x_{i2})\), and the \((j,l)\) entry of \(X'X\) is \(\col_j(X)'\col_l(X)\)
(Theorems~\ref{r1-thm-views}(d) and~\ref{r1-thm-transpose}):
\[
X=\begin{pmatrix}1&2&1\\1&0&1\\1&1&0\\1&1&2\end{pmatrix},\qquad
X'X=\begin{pmatrix}4&4&4\\4&6&4\\4&4&6\end{pmatrix},\qquad
X'y=\begin{pmatrix}7+5+3+5\\14+0+3+5\\7+5+0+10\end{pmatrix}=\begin{pmatrix}20\\22\\22\end{pmatrix}
\]
Row-reduce \([X'X\mid I_3]\), postponing the scaling to the end:
\[
\begin{gathered}
\left(\begin{array}{ccc|ccc}4&4&4&1&0&0\\4&6&4&0&1&0\\4&4&6&0&0&1\end{array}\right)
\xrightarrow{R_2\leftarrow R_2-R_1,\ R_3\leftarrow R_3-R_1}
\left(\begin{array}{ccc|ccc}4&4&4&1&0&0\\0&2&0&-1&1&0\\0&0&2&-1&0&1\end{array}\right)
\\
\xrightarrow{R_1\leftarrow R_1-2R_2,\ R_1\leftarrow R_1-2R_3}
\left(\begin{array}{ccc|ccc}4&0&0&5&-2&-2\\0&2&0&-1&1&0\\0&0&2&-1&0&1\end{array}\right)
\xrightarrow[R_3\leftarrow\frac12R_3]{R_1\leftarrow\frac14R_1,\ R_2\leftarrow\frac12R_2}
\left(\begin{array}{ccc|ccc}1&0&0&\frac54&-\frac12&-\frac12\\0&1&0&-\frac12&\frac12&0\\0&0&1&-\frac12&0&\frac12\end{array}\right)
\end{gathered}
\]
So
\[
(X'X)^{-1}=\frac14\begin{pmatrix}5&-2&-2\\-2&2&0\\-2&0&2\end{pmatrix},\qquad
(X'X)^{-1}X'=\frac14\begin{pmatrix}-1&3&3&-1\\2&-2&0&0\\0&0&-2&2\end{pmatrix}
\]
and
\[
\hat\beta=(X'X)^{-1}X'y=\frac14\begin{pmatrix}-7+15+9-5\\14-10\\-6+10\end{pmatrix}=\begin{pmatrix}3\\1\\1\end{pmatrix}
\]

(b) \(H=X(X'X)^{-1}X'\) (Theorem~\ref{r6-thm-hat}); its trace is \(\frac14(3+3+3+3)=3=\rank X\)
(Corollary~\ref{r6-cor-trace-hat}):
\[
H=\frac14\begin{pmatrix}3&-1&1&1\\-1&3&1&1\\1&1&3&-1\\1&1&-1&3\end{pmatrix},\qquad
\hat y=Hy=\frac14\begin{pmatrix}21-5+3+5\\-7+15+3+5\\7+5+9-5\\7+5-3+15\end{pmatrix}=\begin{pmatrix}6\\4\\4\\6\end{pmatrix},\qquad
e=y-\hat y=\begin{pmatrix}1\\1\\-1\\-1\end{pmatrix}
\]
Also \(\hat y=X\hat\beta=(3+2+1,\ 3+0+1,\ 3+1+0,\ 3+1+2)'\). Check: \(X'e=(1+1-1-1,\ 2+0-1-1,\ 1+1+0-2)'=0\).

(c) \(\bar y=5\), and \(\hat y\) also has mean \(5\), so \(C_4y=(2,0,-2,0)'\) and \(C_4\hat y=(1,-1,-1,1)'\). Hence
\(\mathrm{SST}=4+0+4+0=8\), \(\mathrm{SSR}=1+1+1+1=4\), \(\mathrm{SSE}=e'e=1+1+1+1=4\); indeed \(8=4+4\)
(Theorem~\ref{r11-thm-anova}), and \(R^2=\frac48=\frac12\).
\begin{itemize}
\item \(S_{y\hat y}=(C_4y)'(C_4\hat y)=2+0+2+0=4\), so \(\cor(y,\hat y)^2=\frac{4^2}{8\cdot4}=\frac12=R^2\)
(Proposition~\ref{r11-prop-r2-corr}(a)).
\item \(v=X_{\cdot1}=(2,0,1,1)'\) has mean \(1\) and \(C_4v=(1,-1,0,0)'\), so \(S_{yv}=2\), \(S_{vv}=2\) and
\(\cor(y,v)^2=\frac{2^2}{8\cdot2}=\frac14\). Likewise \(v=X_{\cdot2}=(1,1,0,2)'\) has \(C_4v=(0,0,-1,1)'\),
\(S_{yv}=2\), \(S_{vv}=2\) and \(\cor(y,v)^2=\frac14\). Both are below \(R^2=\frac12\)
(Proposition~\ref{r11-prop-r2-corr}(b)).
\end{itemize}
\end{solution}

\begin{exercise}\label{r12-ex-leverages}
Let \(X=[\one_5\ X_0]\), where the columns of \(X_0\) are \(x_1=(1,2,3,4,5)'\) and \(x_2=(3,0,2,4,1)'\), and let
\(y=(9,4,4,14,9)'\).
\begin{enumerate}
\item Compute \(\bar x\), \(S_c=X_0'C_5X_0\) and \((X'X)^{-1}\) (Proposition~\ref{r12-prop-xtx-inverse}).
\item Compute the diagonal entries \(h_{11},\dots,h_{55}\) of the hat matrix \(H\), the leverages, and check that
their sum is \(\tr H=3\).
\item Compute \(\hat\beta\), \(e\), \(s^2\), and the standardized residuals \(e_i/(s\sqrt{1-h_{ii}})\).
\end{enumerate}
\end{exercise}
\begin{solution}
(a) \(\bar x=\frac15X_0'\one_5=\frac15(15,10)'=(3,2)'\), and the columns of \(C_5X_0=X_0-\one\bar x'\) (Example~\ref{r1-exm-data})
are \(a=(-2,-1,0,1,2)'\) and \(b=(1,-2,0,2,-1)'\), with \(a'a=b'b=10\) and \(a'b=-2+2+0+2-2=0\). So
\(S_c=(C_5X_0)'(C_5X_0)=10I_2\) (Proposition~\ref{r1-prop-centering}). Proposition~\ref{r12-prop-xtx-inverse}(b), with
\(\frac15+\bar x'S_c^{-1}\bar x=\frac15+\frac{9+4}{10}=\frac32\) and \(S_c^{-1}\bar x=\frac1{10}(3,2)'\), gives
\[
(X'X)^{-1}=\frac1{10}\begin{pmatrix}15&-3&-2\\-3&1&0\\-2&0&1\end{pmatrix}
\]

(b) \(h_{ii}=\row_i(X)(X'X)^{-1}\row_i(X)'\) (Theorem~\ref{r1-thm-views}(d)). For a row \((1,u,v)\),
\[
(1,u,v)(X'X)^{-1}(1,u,v)'=\tfrac1{10}\bigl(15-6u-4v+u^2+v^2\bigr)=\tfrac1{10}\bigl((u-3)^2+(v-2)^2+2\bigr)
\]
and \((u-3,v-2)\) is row \(i\) of \(C_5X_0\), that is, \((a_i,b_i)\). As \(a_i^2+b_i^2=5\) for \(i\ne3\) and \(0\) for
\(i=3\), \(h_{ii}=\frac7{10}\) for \(i\ne3\) and \(h_{33}=\frac15\). Their sum is \(4\cdot\frac7{10}+\frac15=3\), which is
\(\tr H=\rank X\) (Theorem~\ref{r10-thm-PX}(b)), as \(\rank X=3\) (Proposition~\ref{r12-prop-xtx-inverse}(a), \(\det S_c=100\)).

(c) \(X'y=(9+4+4+14+9,\ 9+8+12+56+45,\ 27+0+8+56+9)'=(40,130,100)'\), so
\[
\hat\beta=(X'X)^{-1}X'y=\frac1{10}\begin{pmatrix}600-390-200\\-120+130\\-80+100\end{pmatrix}=\begin{pmatrix}1\\1\\2\end{pmatrix}
\]
Then \(\hat y=X\hat\beta=(1+1+6,\ 1+2+0,\ 1+3+4,\ 1+4+8,\ 1+5+2)'=(8,3,8,13,8)'\), \(e=y-\hat y=(1,1,-4,1,1)'\),
\(\mathrm{SSE}=1+1+16+1+1=20\) and \(s^2=20/(5-3)=10\) (Theorem~\ref{r11-thm-sigma}). The standardized residuals are
\(1/\sqrt{10\cdot\frac3{10}}=\frac1{\sqrt3}\) for \(i\ne3\) and \(-4/\sqrt{10\cdot\frac45}=-\sqrt2\) for \(i=3\): the observation
of smallest leverage has the largest residual.
\end{solution}

\begin{exercise}\label{r10-ex-rank-deficient}
Add to the design matrix \(X\) of Example~\ref{r6-exm-hat} the sum of its two variables as a fourth column:
\(\tilde X=[\one_4\ X_{\cdot1}\ X_{\cdot2}\ X_{\cdot1}+X_{\cdot2}]\in\R^{4\times4}\), and let \(y=(9,3,5,3)'\).
\begin{enumerate}
\item Find \(\rank\tilde X\) and a basis of \(\mathcal N(\tilde X)\).
\item Solve the normal equations \(\tilde X'\tilde X\beta=\tilde X'y\) by row reduction.
\item Show that all least-squares estimates give the fitted values \(\hat y=(8,4,6,2)'\) of
Example~\ref{r10-exm-least-squares}, and find \(S(\hat\beta)\).
\end{enumerate}
\end{exercise}
\begin{solution}
(a) The columns are \(\one_4\), \((2,0,2,0)'\), \((3,3,1,1)'\) and \((5,3,3,1)'\). The fourth is the sum of the
second and the third, so \(\tilde X(0,1,1,-1)'=0\); the first three are the columns of \(X\), which are
independent (Example~\ref{r10-exm-least-squares}). So \(\rank\tilde X=3\), \(\dim\mathcal N(\tilde X)=4-3=1\)
(Theorem~\ref{thm:r4-rank-nullity}), and \((0,1,1,-1)'\) is a basis of \(\mathcal N(\tilde X)\).

(b) The entries of \(\tilde X'\tilde X\) and \(\tilde X'y\) are inner products of columns (Theorems~\ref{r1-thm-views}(d)
and~\ref{r1-thm-transpose}):
\[
\left(\begin{array}{cccc|c}4&4&8&12&20\\4&8&8&16&28\\8&8&20&28&44\\12&16&28&44&72\end{array}\right)
\xrightarrow[R_4\leftarrow R_4-3R_1]{R_2\leftarrow R_2-R_1,\ R_3\leftarrow R_3-2R_1}
\left(\begin{array}{cccc|c}4&4&8&12&20\\0&4&0&4&8\\0&0&4&4&4\\0&4&4&8&12\end{array}\right)
\]
\[
\xrightarrow{R_4\leftarrow R_4-R_2}
\left(\begin{array}{cccc|c}4&4&8&12&20\\0&4&0&4&8\\0&0&4&4&4\\0&0&4&4&4\end{array}\right)
\xrightarrow{R_4\leftarrow R_4-R_3}
\left(\begin{array}{cccc|c}4&4&8&12&20\\0&4&0&4&8\\0&0&4&4&4\\0&0&0&0&0\end{array}\right)
\]
\[
\xrightarrow[R_3\leftarrow\frac14R_3]{R_1\leftarrow\frac14R_1,\ R_2\leftarrow\frac14R_2}
\left(\begin{array}{cccc|c}1&1&2&3&5\\0&1&0&1&2\\0&0&1&1&1\\0&0&0&0&0\end{array}\right)
\]
The last matrix is an echelon form with pivots in columns \(1,2,3\) and none in the last column, so the
system is consistent and \(\beta_3=t\) is free (Theorem~\ref{thm:r2-consistency}). Back substitution:
\(\beta_2=1-t\), \(\beta_1=2-t\) and \(\beta_0=5-\beta_1-2\beta_2-3\beta_3=5-(2-t)-2(1-t)-3t=1\). So the solutions are
\[
\beta=\begin{pmatrix}1\\2-t\\1-t\\t\end{pmatrix}=\begin{pmatrix}1\\2\\1\\0\end{pmatrix}+t\begin{pmatrix}0\\-1\\-1\\1\end{pmatrix}\qquad(t\in\R)
\]
a particular solution plus \(\mathcal N(\tilde X'\tilde X)=\mathcal N(\tilde X)\), as
Theorem~\ref{r10-thm-normal-equations}(c) says.

(c) These solutions are the least-squares estimates (Theorem~\ref{r10-thm-least-squares}(a)), and each
gives \(\tilde X\beta=X(1,2,1)'+t\tilde X(0,-1,-1,1)'=(8,4,6,2)'\). So \(\hat y\) and \(e=(1,-1,-1,1)'\) are those
of Example~\ref{r10-exm-least-squares}, and \(S(\hat\beta)=e'e=4\).
\end{solution}

\needspace{20\baselineskip}
\begin{exercise}\label{r10-dr-mnls}
Let
\[
\tilde X=\begin{pmatrix}1&2&3&5\\1&0&3&3\\1&2&1&3\\1&0&1&1\end{pmatrix},\qquad y=\begin{pmatrix}9\\3\\5\\3\end{pmatrix}
\]
as in Exercise~\ref{r10-ex-rank-deficient}, where the least-squares estimates are found to be
\(\beta=(1,\ 2-t,\ 1-t,\ t)'\) (\(t\in\R\)).
\begin{enumerate}
\item Row-reduce \(\tilde X\) to its RREF, and write down the full-rank factorization \(\tilde X=BC\) of
Theorem~\ref{thm:r4-full-rank-factorization}(a).
\item Here \(B\) is the design matrix \(X\) of Example~\ref{r10-exm-least-squares}, so \((B'B)^{-1}B'y=(1,2,1)'\). Compute
\(\tilde X^+y\) by Theorem~\ref{r10-thm-mp-exists}(b).
\item Check that \(\tilde X^+y\) is the shortest least-squares estimate.
\end{enumerate}
\end{exercise}
\begin{solution}
(a)
\[
\begin{gathered}
\tilde X\xrightarrow[R_3\leftarrow R_3-R_1,\ R_4\leftarrow R_4-R_1]{R_2\leftarrow R_2-R_1}
\begin{pmatrix}1&2&3&5\\0&-2&0&-2\\0&0&-2&-2\\0&-2&-2&-4\end{pmatrix}
\xrightarrow{R_4\leftarrow R_4-R_2}
\begin{pmatrix}1&2&3&5\\0&-2&0&-2\\0&0&-2&-2\\0&0&-2&-2\end{pmatrix}
\\
\xrightarrow{R_4\leftarrow R_4-R_3}
\begin{pmatrix}1&2&3&5\\0&-2&0&-2\\0&0&-2&-2\\0&0&0&0\end{pmatrix}
\xrightarrow[R_3\leftarrow-\frac12R_3]{R_2\leftarrow-\frac12R_2}
\begin{pmatrix}1&2&3&5\\0&1&0&1\\0&0&1&1\\0&0&0&0\end{pmatrix}
\xrightarrow[R_1\leftarrow R_1-3R_3]{R_1\leftarrow R_1-2R_2}
\begin{pmatrix}1&0&0&0\\0&1&0&1\\0&0&1&1\\0&0&0&0\end{pmatrix}
\end{gathered}
\]
The pivot columns are \(1,2,3\), so \(\tilde X=BC\) with \(B\) the first three columns of \(\tilde X\), that is, the design
matrix \(X\) of Example~\ref{r10-exm-least-squares}, and
\[
C=\begin{pmatrix}1&0&0&0\\0&1&0&1\\0&0&1&1\end{pmatrix},\qquad
CC'=\begin{pmatrix}1&0&0\\0&2&1\\0&1&2\end{pmatrix}
\]
(b) \(\det(CC')=1\cdot(4-1)=3\) and, by Theorem~\ref{r5d-thm-adjugate},
\[
(CC')^{-1}=\frac13\begin{pmatrix}3&0&0\\0&2&-1\\0&-1&2\end{pmatrix},\qquad
(CC')^{-1}\begin{pmatrix}1\\2\\1\end{pmatrix}=\frac13\begin{pmatrix}3\\4-1\\-2+2\end{pmatrix}=\begin{pmatrix}1\\1\\0\end{pmatrix}
\]
By Theorem~\ref{r10-thm-mp-exists}(b), \(\tilde X^+y=C'(CC')^{-1}(B'B)^{-1}B'y=C'(1,1,0)'=\col_1(C')+\col_2(C')=(1,1,0,1)'\)
(Theorem~\ref{r1-thm-views}(b)).

(c) \(\lVert\beta\rVert^2=1+(2-t)^2+(1-t)^2+t^2=3t^2-6t+6=3(t-1)^2+3\) is smallest exactly for \(t=1\), where
\(\beta=(1,1,0,1)'=\tilde X^+y\), of squared length \(3\), as Theorem~\ref{r10-thm-mnls} says.
\end{solution}

\begin{exercise}\label{r11-dr-centered-design}
Let \(X_0\in\R^{6\times2}\) be the data of Exercise~\ref{r2-dr-centering}, with the columns \((2,2,0,0,2,0)'\) and
\((2,0,2,0,1,1)'\); let \(X=[\one_6\ X_0]\), and let \(\tilde X=[\one_6\ C_6X_0]\) be the design with the centered columns.
By Exercise~\ref{r2-dr-centering}, \(\tilde X=XQ\) for the product \(Q\) of elementary matrices found there.
\begin{enumerate}
\item Compute \(\tilde X'\tilde X\) and its inverse.
\item Find \((X'X)^{-1}\) from \(Q\) and (a), and check the first row of \((X'X)(X'X)^{-1}=I_3\).
\end{enumerate}
\end{exercise}
\begin{solution}
(a) By Exercise~\ref{r2-dr-centering}, the columns of \(\tilde X\) are \(\one_6\), \(u=(1,1,-1,-1,1,-1)'\) and
\(v=(1,-1,1,-1,0,0)'\), with \(\one'\one=6\), \(\one'u=0\), \(\one'v=0\), \(u'u=6\), \(u'v=1-1-1+1=0\) and \(v'v=4\). So
(Theorems~\ref{r1-thm-views}(d) and~\ref{r1-thm-transpose}, Definition~\ref{r2-def-inverse})
\[
\tilde X'\tilde X=\begin{pmatrix}6&0&0\\0&6&0\\0&0&4\end{pmatrix},\qquad
(\tilde X'\tilde X)^{-1}=\begin{pmatrix}\frac16&0&0\\0&\frac16&0\\0&0&\frac14\end{pmatrix}
\]

(b) \(\tilde X'\tilde X=(XQ)'XQ=Q'(X'X)Q\) (Theorem~\ref{r1-thm-transpose}), and \(Q\) is invertible
(Exercise~\ref{r2-dr-centering}). So \(X'X=(Q')^{-1}(\tilde X'\tilde X)Q^{-1}\), and by
Theorem~\ref{r2-thm-inverse-props}(a)--(c), \((X'X)^{-1}=Q(\tilde X'\tilde X)^{-1}Q'\):
\[
(X'X)^{-1}=\begin{pmatrix}1&-1&-1\\0&1&0\\0&0&1\end{pmatrix}
\begin{pmatrix}\frac16&0&0\\0&\frac16&0\\0&0&\frac14\end{pmatrix}
\begin{pmatrix}1&0&0\\-1&1&0\\-1&0&1\end{pmatrix}
=\frac1{12}\begin{pmatrix}7&-2&-3\\-2&2&0\\-3&0&3\end{pmatrix}
\]
Here the entries of \(X'X\) are \(n\), the column sums and the sums of squares and products of the columns
(Example~\ref{r1-exm-data}):
\[
X'X=\begin{pmatrix}6&6&6\\6&12&6\\6&6&10\end{pmatrix}
\]
and the first row of \((X'X)(X'X)^{-1}\) is \(\frac1{12}(42-12-18,\ -12+12,\ -18+18)=(1,0,0)\).
\end{solution}

\begin{exercise}\label{r11-ex-regression-drill}
A dose is given at the three coded levels \(-1,0,1\), two observations at each level:
\[
\begin{array}{c|cccccc}x_i&-1&-1&0&0&1&1\\\hline y_i&5&3&3&1&7&5\end{array}
\]
Fit the quadratic regression \(y_i=\beta_0+\beta_1x_i+\beta_2x_i^2+\varepsilon_i\) under (A1)--(A3).
\begin{enumerate}
\item Explain why \(X\) has full column rank, and compute \(X'X\), \(X'y\), \((X'X)^{-1}\) and \(\hat\beta\).
\item Compute \(\hat y\), \(e\), SSE and \(s^2\), and the estimate \(s^2(X'X)^{-1}\) of \(\var(\hat\beta)\).
\item Compute SST, SSR, \(R^2\) and \(\bar R^2\).
\end{enumerate}
\end{exercise}
\begin{solution}
(a) The \(x_i\) take three different values, so \(X\), with rows \((1,x_i,x_i^2)\), has full column rank
(Proposition~\ref{r11-prop-quadratic}). The columns of \(X\) are \(\one_6\), \(x=(-1,-1,0,0,1,1)'\) and \((1,1,0,0,1,1)'\), and
\[
X'X=\begin{pmatrix}6&0&4\\0&4&0\\4&0&4\end{pmatrix},\qquad
X'y=\begin{pmatrix}5+3+3+1+7+5\\-5-3+7+5\\5+3+7+5\end{pmatrix}=\begin{pmatrix}24\\4\\20\end{pmatrix}
\]
Row-reduce \([X'X\mid I_3]\), postponing the scaling to the end:
\[
\begin{gathered}
\left(\begin{array}{ccc|ccc}6&0&4&1&0&0\\0&4&0&0&1&0\\4&0&4&0&0&1\end{array}\right)
\xrightarrow{R_1\leftarrow R_1-R_3}
\left(\begin{array}{ccc|ccc}2&0&0&1&0&-1\\0&4&0&0&1&0\\4&0&4&0&0&1\end{array}\right)
\xrightarrow{R_3\leftarrow R_3-2R_1}
\left(\begin{array}{ccc|ccc}2&0&0&1&0&-1\\0&4&0&0&1&0\\0&0&4&-2&0&3\end{array}\right)
\\
\xrightarrow[R_3\leftarrow\frac14R_3]{R_1\leftarrow\frac12R_1,\ R_2\leftarrow\frac14R_2}
\left(\begin{array}{ccc|ccc}1&0&0&\frac12&0&-\frac12\\0&1&0&0&\frac14&0\\0&0&1&-\frac12&0&\frac34\end{array}\right)
\end{gathered}
\]
So
\[
(X'X)^{-1}=\frac14\begin{pmatrix}2&0&-2\\0&1&0\\-2&0&3\end{pmatrix},\qquad
\hat\beta=(X'X)^{-1}X'y=\frac14\begin{pmatrix}48-40\\4\\-48+60\end{pmatrix}=\begin{pmatrix}2\\1\\3\end{pmatrix}
\]

(b) \(\hat y_i=2+x_i+3x_i^2\) is \(4\), \(2\), \(6\) at \(x_i=-1,0,1\): \(\hat y=(4,4,2,2,6,6)'\) and
\(e=y-\hat y=(1,-1,1,-1,1,-1)'\). Check: \(X'e=(0,\ -1+1+1-1,\ 1-1+1-1)'=0\). Then \(\mathrm{SSE}=6\), and with \(n=6\),
\(k=2\), \(s^2=6/(6-3)=2\) (Theorem~\ref{r11-thm-sigma}), so
\[
s^2(X'X)^{-1}=\frac12\begin{pmatrix}2&0&-2\\0&1&0\\-2&0&3\end{pmatrix}
\]
for instance \(\var(\hat\beta_1)\) is estimated by \(\frac12\).

(c) \(\bar y=24/6=4\), so \(\mathrm{SST}=1+1+1+9+9+1=22\) and \(\mathrm{SSR}=0+0+4+4+4+4=16\); indeed \(22=16+6\)
(Theorem~\ref{r11-thm-anova}). So \(R^2=\frac{16}{22}=\frac8{11}\) and
\(\bar R^2=1-\frac{6/3}{22/5}=1-\frac{10}{22}=\frac6{11}\) (Definition~\ref{r11-def-adjusted}).
\end{solution}

\begin{exercise}\label{r6-ex-orthogonal-polynomials}
Let \(a_1=\one_5\), \(a_2=(1,2,3,4,5)'\), \(a_3=(1,4,9,16,25)'\) and \(a_4=(1,8,27,64,125)'\), the values of
\(1,t,t^2,t^3\) at \(t=1,\dots,5\), and let \(U_3=\Span\{a_1,a_2,a_3\}\) and \(U_4=\Span\{a_1,\dots,a_4\}\).
\begin{enumerate}
\item Apply the Gram--Schmidt process to \(a_1,\dots,a_4\), and show that \(b_2\), \(b_3\) and \(\frac56b_4\) are
\[
c_1=(-2,-1,0,1,2)',\qquad c_2=(2,-1,-2,-1,2)',\qquad c_3=(-1,2,0,-2,1)'
\]
\item Show that \(b_2,b_3,b_4\) are the values at \(t=1,\dots,5\) of the polynomials \(t-3\), \((t-3)^2-2\) and
\((t-3)^3-\frac{17}5(t-3)\), and that \(c_4=(1,-4,6,-4,1)'\) spans \(U_4^\perp\). (The vectors \(c_1,\dots,c_4\) are the
tabulated coefficients of the orthogonal polynomials for five equally spaced levels.)
\item For \(y=(0,7,7,5,6)'\), find the vector of \(U_3\) closest to \(y\), the polynomial of degree at most \(2\)
whose values at \(t=1,\dots,5\) it lists, and the distance from \(y\) to \(U_3\). Read the answer as the least-squares
fit of the quadratic regression \(y_i=\beta_0+\beta_1t_i+\beta_2t_i^2+\varepsilon_i\) with \(t_i=i\): find \(\hat\beta\) and SSE.
\end{enumerate}
\end{exercise}
\begin{solution}
(a) By Theorem~\ref{r6-thm-gram-schmidt}:
\begin{itemize}
\item \(b_1=\one_5\), \(\lVert b_1\rVert^2=5\); \(a_2'b_1=15\), so \(t_{12}=3\) and \(b_2=a_2-3b_1=(-2,-1,0,1,2)'=c_1\), with
\(\lVert b_2\rVert^2=10\).
\item \(a_3'b_1=55\) and \(a_3'b_2=-2-4+0+16+50=60\), so \(t_{13}=11\), \(t_{23}=6\) and
\(b_3=a_3-11b_1-6b_2=(1-11+12,\ 4-11+6,\ 9-11+0,\ 16-11-6,\ 25-11-12)'=c_2\), with \(\lVert b_3\rVert^2=14\).
\item \(a_4'b_1=225\), \(a_4'b_2=-2-8+0+64+250=304\) and \(a_4'b_3=2-8-54-64+250=126\), so \(t_{14}=45\),
\(t_{24}=\frac{152}5\), \(t_{34}=9\), and \(b_4=a_4-45b_1-\frac{152}5b_2-9b_3=\frac15(-6,12,0,-12,6)'=\frac65c_3\).
\end{itemize}
(b) The values of \(s=t-3\) at \(t=1,\dots,5\) are \(b_2=(-2,-1,0,1,2)'\); those of \(s^2-2\) are
\((4-2,\ 1-2,\ 0-2,\ 1-2,\ 4-2)'=b_3\); and those of \(s^3-\frac{17}5s\) are
\((-8+\frac{34}5,\ -1+\frac{17}5,\ 0,\ 1-\frac{17}5,\ 8-\frac{34}5)'=\frac15(-6,12,0,-12,6)'=b_4\). Next
\(c_4'a_1=1-4+6-4+1=0\), \(c_4'a_2=1-8+18-16+5=0\), \(c_4'a_3=1-16+54-64+25=0\) and
\(c_4'a_4=1-32+162-256+125=0\), so \(c_4\in U_4^\perp\) (Theorem~\ref{r6-thm-complement}(a)). The vectors
\(b_1,\dots,b_4\) are nonzero, pairwise orthogonal and span \(U_4\) (Theorem~\ref{r6-thm-gram-schmidt}), so they form a
basis of \(U_4\) (Proposition~\ref{r6-prop-orthogonal-list}(b)); hence \(\dim U_4^\perp=5-4=1\)
(Theorem~\ref{r6-thm-complement}(c)), and \(c_4\ne0\) spans \(U_4^\perp\) (Theorem~\ref{thm:r4-counting}(d)).

(c) \(b_1'y=0+7+7+5+6=25\), \(b_2'y=c_1'y=0-7+0+5+12=10\) and \(b_3'y=c_2'y=0-7-14-5+12=-14\), with
\(b_1'b_1=5\), \(b_2'b_2=10\) and \(b_3'b_3=14\) by (a). By (a), \(b_1,b_2,b_3\) are nonzero, pairwise orthogonal and span \(U_3\)
(Theorem~\ref{r6-thm-gram-schmidt}), so \(b_j/\lVert b_j\rVert\) (\(j=1,2,3\)) form an orthonormal basis of \(U_3\). By
Theorems~\ref{r6-thm-best-approximation} and~\ref{r6-thm-projection-matrix}(b), the vector of \(U_3\) closest to \(y\) is
\[
P_{U_3}y=\sum_{j=1}^3\frac{b_j'y}{b_j'b_j}\,b_j=\frac{25}5\,\one_5+\frac{10}{10}\,c_1+\frac{-14}{14}\,c_2
=5\,\one_5+c_1-c_2=(1,5,7,7,5)'
\]
and \(y-P_{U_3}y=(-1,2,0,-2,1)'=c_3\), so the distance from \(y\) to \(U_3\) is \(\sqrt{10}\). By (a) and (b),
\(P_{U_3}y=5b_1+b_2-b_3\) lists the values at \(t=1,\dots,5\) of \(5+(t-3)-\bigl((t-3)^2-2\bigr)=7+(t-3)-(t-3)^2\);
for instance \(7-2-4=1\) at \(t=1\). The columns of \(X=[a_1\ a_2\ a_3]\), with rows \((1,t_i,t_i^2)\), span \(U_3\), so
\(P_{U_3}y\) is the vector \(\hat y\) of fitted values (Theorem~\ref{r10-thm-least-squares}(b)). As
\(7+(t-3)-(t-3)^2=-5+7t-t^2\), \(\hat y=X(-5,7,-1)'\); the \(t_i\) take five different values, so \(X\) has full
column rank (Proposition~\ref{r11-prop-quadratic}), and \(\hat\beta=(-5,7,-1)'\) is the only least-squares estimate
(Theorems~\ref{r10-thm-least-squares}(a) and~\ref{r10-thm-normal-equations}(b)). Finally \(\mathrm{SSE}=\lVert y-\hat y\rVert^2=\lVert c_3\rVert^2=10\).
\end{solution}

\begin{exercise}\label{r11-dr-gauss-markov}
Let \(X\) and \(y=(7,5,3,5)'\) be as in Exercise~\ref{r11-dr-hat}, under (A1)--(A4). Let \(X_s\) consist of the first three
rows of \(X\), and \(C=[X_s^{-1}\ \ 0]\in\R^{3\times4}\): \(\tilde\beta=Cy\) solves \(X_s\beta=(y_1,y_2,y_3)'\), so it uses the
first three observations only.
\begin{enumerate}
\item Find \(X_s^{-1}\) by Theorem~\ref{r5d-thm-adjugate}, check that \(CX=I_3\), and compute \(\tilde\beta\).
\item Find \(\var(\tilde\beta)\), \(\var(\hat\beta)\) and \(D=C-(X'X)^{-1}X'\), and check that
\(\var(\tilde\beta)-\var(\hat\beta)=\sigma^2DD'\).
\item For which \(j\) is \(\var(\tilde\beta_j)=\var(\hat\beta_j)\)? Compare \(\tilde\beta_j\) with \(\hat\beta_j\) for it.
\end{enumerate}
\end{exercise}
\begin{solution}
(a) The first three rows of \(X\) are
\[
X_s=\begin{pmatrix}1&2&1\\1&0&1\\1&1&0\end{pmatrix}
\]
Expanding along row \(2\),
\(\det X_s=-(2\cdot0-1\cdot1)+0-(1\cdot1-2\cdot1)=1+1=2\), and the nine cofactors give
\[
\operatorname{adj}X_s=\begin{pmatrix}-1&1&2\\1&-1&0\\1&1&-2\end{pmatrix},\qquad
X_s^{-1}=\frac12\begin{pmatrix}-1&1&2\\1&-1&0\\1&1&-2\end{pmatrix},\qquad
C=\frac12\begin{pmatrix}-1&1&2&0\\1&-1&0&0\\1&1&-2&0\end{pmatrix}
\]
(Theorem~\ref{r5d-thm-adjugate}). By the block product, \(CX=X_s^{-1}X_s+0=I_3\), so \(\expc(Cy)=CX\beta=\beta\) by
Theorem~\ref{r11-thm-linear}(b) and (A1): \(\tilde\beta\) is a linear unbiased estimator of \(\beta\). For the data,
\(\tilde\beta=X_s^{-1}(7,5,3)'=\frac12(-7+5+6,\ 7-5+0,\ 7+5-6)'=(2,1,3)'\), whereas \(\hat\beta=(3,1,1)'\).

(b) By Theorem~\ref{r11-thm-linear}(b) and (A3), \(\var(\tilde\beta)=\sigma^2CC'=\sigma^2X_s^{-1}(X_s^{-1})'\), and
\(\var(\hat\beta)=\sigma^2(X'X)^{-1}\) (Theorem~\ref{r11-thm-unbiased}(b)), with \((X'X)^{-1}\) and \((X'X)^{-1}X'\) of
Exercise~\ref{r11-dr-hat}(a):
\[
\var(\tilde\beta)=\frac{\sigma^2}4\begin{pmatrix}6&-2&-4\\-2&2&0\\-4&0&6\end{pmatrix},\qquad
\var(\hat\beta)=\frac{\sigma^2}4\begin{pmatrix}5&-2&-2\\-2&2&0\\-2&0&2\end{pmatrix},\qquad
D=\frac14\begin{pmatrix}-1&-1&1&1\\0&0&0&0\\2&2&-2&-2\end{pmatrix}
\]
Then
\[
DD'=\frac1{16}\begin{pmatrix}4&0&-8\\0&0&0\\-8&0&16\end{pmatrix}=\frac14\begin{pmatrix}1&0&-2\\0&0&0\\-2&0&4\end{pmatrix},\qquad
\var(\tilde\beta)-\var(\hat\beta)=\frac{\sigma^2}4\begin{pmatrix}1&0&-2\\0&0&0\\-2&0&4\end{pmatrix}=\sigma^2DD'
\]
as in the proof of Theorem~\ref{r11-thm-gauss-markov}. The rows of \(D\) are multiples of \(e=(1,1,-1,-1)'\) of
Exercise~\ref{r11-dr-hat}(b), and \(X'e=0\); so \(DX=0\).

(c) Only for \(j=1\): the diagonal of \(DD'\) is \(\frac14(1,0,4)\), so \(\var(\tilde\beta_1)=\var(\hat\beta_1)=\frac12\sigma^2\),
while \(\var(\tilde\beta_0)=\frac64\sigma^2>\frac54\sigma^2\) and \(\var(\tilde\beta_2)=\frac64\sigma^2>\frac24\sigma^2\). Row \(2\)
of \(D\) is \(0\), so \(\tilde\beta_1=\hat\beta_1=\frac12(y_1-y_2)\) for every \(y\) (here \(1\)): the same estimator.
\end{solution}

\begin{exercise}\label{r11-ex-gls}
Each of six responses \(y_i\) is the mean of \(m_i\) uncorrelated determinations with variance \(\sigma^2\), where
\(m=(1,1,2,2,1,1)'\). So in the model \(y=X\beta+\varepsilon\), with \(\expc(\varepsilon)=0\), the errors are uncorrelated with
\(\var(\varepsilon_i)=\sigma^2/m_i\): \(\var(\varepsilon)=\sigma^2V\) with \(V=\diag\bigl(1,1,\frac12,\frac12,1,1\bigr)\). Let
\(X=[\one_6\ x_1\ x_2]\) and
\[
x_1=\begin{pmatrix}-1\\-1\\1\\1\\-1\\-1\end{pmatrix},\qquad x_2=\begin{pmatrix}-1\\0\\-1\\1\\0\\1\end{pmatrix},\qquad
y=\begin{pmatrix}0\\3\\7\\7\\3\\6\end{pmatrix}
\]
\begin{enumerate}
\item Compute \(X'V^{-1}X\), \(X'V^{-1}y\) and the generalized least-squares estimate \(\tilde\beta\) of Theorem~\ref{r12-thm-gls}.
\item Compute \(X'X\), \((X'X)^{-1}\), \(X'y\) and the least-squares estimate \(\hat\beta\).
\item Find \(\var(\tilde\beta)\) and \(\var(\hat\beta)\), and check that \(\var(\hat\beta)-\var(\tilde\beta)\) is PSD.
\item Compute \(Q(\tilde\beta)=(y-X\tilde\beta)'V^{-1}(y-X\tilde\beta)\) and the unbiased estimate \(Q(\tilde\beta)/(6-3)\) of
\(\sigma^2\).
\end{enumerate}
\end{exercise}
\begin{solution}
(a) \(V^{-1}=\diag(1,1,2,2,1,1)\), so \(V^{-1}X\) is \(X\) with rows \(3\) and \(4\) doubled
(Theorem~\ref{r1-thm-views}(c)), \(X'V^{-1}=(V^{-1}X)'\), and
\[
X'V^{-1}=\begin{pmatrix}1&1&2&2&1&1\\-1&-1&2&2&-1&-1\\-1&0&-2&2&0&1\end{pmatrix},\qquad
X'V^{-1}X=\begin{pmatrix}8&0&0\\0&8&0\\0&0&6\end{pmatrix},\qquad
X'V^{-1}y=\begin{pmatrix}40\\16\\6\end{pmatrix}
\]
for instance \(\one'V^{-1}y=0+3+14+14+3+6=40\) and \(x_1'V^{-1}y=0-3+14+14-3-6=16\). By Theorem~\ref{r12-thm-gls},
\(\tilde\beta=(X'V^{-1}X)^{-1}X'V^{-1}y=(\frac{40}8,\frac{16}8,\frac66)'=(5,2,1)'\).

(b) The \((j,l)\) entry of \(X'X\) is the product of columns \(j\) and \(l\) of \(X\) (Theorems~\ref{r1-thm-views}(d)
and~\ref{r1-thm-transpose}):
\[
X'X=\begin{pmatrix}6&-2&0\\-2&6&0\\0&0&4\end{pmatrix},\qquad
(X'X)^{-1}=\frac1{16}\begin{pmatrix}3&1&0\\1&3&0\\0&0&4\end{pmatrix},\qquad
X'y=\begin{pmatrix}26\\2\\6\end{pmatrix}
\]
where the product of \(X'X\) with the second matrix is \(I_3\): its first row is \(\frac1{16}(18-2,\ 6-6,\ 0)=(1,0,0)\), its
second row \(\frac1{16}(-6+6,\ -2+18,\ 0)=(0,1,0)\) and its third row \(\frac1{16}(0,0,16)\) (Definition~\ref{r2-def-inverse}).
So
\[
\hat\beta=(X'X)^{-1}X'y=\frac1{16}\begin{pmatrix}78+2\\26+6\\24\end{pmatrix}=\begin{pmatrix}5\\2\\\frac32\end{pmatrix}
\]

(c) By Theorem~\ref{r12-thm-gls}(a), \(\var(\tilde\beta)=\sigma^2(X'V^{-1}X)^{-1}=\sigma^2\diag\bigl(\frac18,\frac18,\frac16\bigr)\). By
Theorem~\ref{r12-thm-gls}(c), \(\var(\hat\beta)=\sigma^2(X'X)^{-1}X'VX(X'X)^{-1}\), with
\[
X'VX=\begin{pmatrix}5&-3&0\\-3&5&0\\0&0&3\end{pmatrix},\qquad
(X'X)^{-1}X'VX=\frac14\begin{pmatrix}3&-1&0\\-1&3&0\\0&0&3\end{pmatrix},\qquad
\var(\hat\beta)=\sigma^2\begin{pmatrix}\frac18&0&0\\0&\frac18&0\\0&0&\frac3{16}\end{pmatrix}
\]
So \(\var(\hat\beta)-\var(\tilde\beta)=\sigma^2\diag\bigl(0,0,\frac3{16}-\frac16\bigr)=\frac{\sigma^2}{48}\diag(0,0,1)\), which is
PSD (Proposition~\ref{r8-prop-pd-basic}(a)): the two estimators have the same variances for \(\beta_0\) and \(\beta_1\), and
\(\var(\hat\beta_2)=\frac3{16}\sigma^2>\frac16\sigma^2=\var(\tilde\beta_2)\).

(d) \(X\tilde\beta=(5-2-1,\ 5-2,\ 5+2-1,\ 5+2+1,\ 5-2,\ 5-2+1)'=(2,3,6,8,3,4)'\), so
\(y-X\tilde\beta=(-2,0,1,-1,0,2)'\) and, with the weights \(1,1,2,2,1,1\) of \(V^{-1}\),
\(Q(\tilde\beta)=4+0+2+2+0+4=12\). Let \(T=\diag(1,1,\sqrt2,\sqrt2,1,1)\), so that \(T^2=V^{-1}\) and \(TVT=I_6\). The model
\(Ty=(TX)\beta+T\varepsilon\) has \(\expc(T\varepsilon)=0\) and \(\var(T\varepsilon)=\sigma^2TVT=\sigma^2I_6\)
(Theorem~\ref{r11-thm-linear}(b)), and \(\rank(TX)=\rank X=3\) (\(T\) is invertible; Theorem~\ref{thm:r4-rank-product}(d)).
As \(\lVert Ty-TX\beta\rVert^2=(y-X\beta)'V^{-1}(y-X\beta)\), its least-squares estimate is \(\tilde\beta\)
(Theorem~\ref{r12-thm-gls}(a)) and its error sum of squares is \(Q(\tilde\beta)\). By Theorem~\ref{r11-thm-sigma},
\(Q(\tilde\beta)/(6-3)=12/3=4\) is an unbiased estimate of \(\sigma^2\).
\end{solution}

\begin{exercise}\label{r11-ex-fwl-drill}
Let \(X=[\one_6\ x_1\ x_2\ x_3]\in\R^{6\times4}\) and \(y\in\R^6\), with
\[
x_1=\begin{pmatrix}1\\1\\1\\-1\\-1\\-1\end{pmatrix},\quad x_2=\begin{pmatrix}1\\-1\\0\\1\\-1\\0\end{pmatrix},\quad
x_3=\begin{pmatrix}3\\3\\0\\-1\\-1\\2\end{pmatrix},\quad y=\begin{pmatrix}9\\5\\4\\-1\\-1\\2\end{pmatrix}
\]
In Theorem~\ref{r11-thm-fwl} take \(X_1=x_3\) and \(X_2=[\one_6\ x_1\ x_2]\); the design \([X_1\ X_2]\) is \(X\) with its columns
reordered, and \(\hat\gamma\) is the coefficient of \(x_3\).
\begin{enumerate}
\item Compute \(X_2'X_2\), \((X_2'X_2)^{-1}\), \(X_2'x_3\) and \(X_2'y\), and then \(H_2x_3\), \(H_2y\), \(M_2X_1\) and \(M_2y\).
\item Use Theorem~\ref{r11-thm-fwl} to find \(\hat\gamma\) and the residual vector \(e\) of the least-squares fit of \(y\) on
\(X\), and check that \(X'e=0\).
\item Check that \(X_1'M_2X_1=X_1'X_1-X_1'X_2(X_2'X_2)^{-1}X_2'X_1\), the Schur complement of \(X_2'X_2\) in \(X'X\)
(Theorem~\ref{r5b-thm-schur}, with the blocks of \(X'X\) in the order \(X_2\), \(X_1\)).
\end{enumerate}
\end{exercise}
\begin{solution}
(a) The \((j,l)\) entry of \(X_2'X_2\) is the product of columns \(j\) and \(l\) of \(X_2\) (Theorems~\ref{r1-thm-views}(d)
and~\ref{r1-thm-transpose}); these columns are orthogonal, so \(X_2'X_2\) is diagonal, and its inverse is diagonal
(Theorem~\ref{r2-thm-triangular}):
\[
X_2'X_2=\begin{pmatrix}6&0&0\\0&6&0\\0&0&4\end{pmatrix},\qquad
(X_2'X_2)^{-1}=\begin{pmatrix}\frac16&0&0\\0&\frac16&0\\0&0&\frac14\end{pmatrix},\qquad
X_2'x_3=\begin{pmatrix}6\\6\\0\end{pmatrix},\qquad X_2'y=\begin{pmatrix}18\\18\\4\end{pmatrix}
\]
for instance \(\one'x_3=3+3+0-1-1+2=6\), \(x_1'x_3=(3+3+0)-(-1-1+2)=6\) and \(x_2'y=9-5+0-1+1+0=4\). The vector
\(c=(X_2'X_2)^{-1}X_2'v\) holds the coefficients of \(H_2v=X_2c\) on the columns of \(X_2\) (Theorem~\ref{r1-thm-views}(b)):
\(c=(1,1,0)'\) for \(v=x_3\) and \(c=(3,3,1)'\) for \(v=y\). So
\[
H_2x_3=\one+x_1=\begin{pmatrix}2\\2\\2\\0\\0\\0\end{pmatrix},\qquad
H_2y=3\cdot\one+3x_1+x_2=\begin{pmatrix}7\\5\\6\\1\\-1\\0\end{pmatrix}
\]
\[
M_2X_1=x_3-H_2x_3=\begin{pmatrix}1\\1\\-2\\-1\\-1\\2\end{pmatrix},\qquad
M_2y=y-H_2y=\begin{pmatrix}2\\0\\-2\\-2\\0\\2\end{pmatrix}
\]
(b) First, \(X\) satisfies (A4): \(\one_6\), \(x_1\), \(x_2\) are nonzero and pairwise orthogonal, hence independent
(Proposition~\ref{r6-prop-orthogonal-list}(b)); \(x_3\notin\mathcal C(X_2)\), since otherwise \(H_2x_3=x_3\)
(Theorem~\ref{r6-thm-hat}(a)), while \(M_2x_3\ne0\); so the four columns of \(X\) are independent
(Proposition~\ref{r3-prop-adjoin}(b)). The design \([X_1\ X_2]=[x_3\ \one_6\ x_1\ x_2]\) has the same columns as \(X\), so its
least-squares fit has the same fitted values \(\hat y\) and residual vector \(e\) (Theorem~\ref{r10-thm-least-squares}(b)); as
these four columns are independent, they form a basis of \(\mathcal C(X)\), and the coordinates of \(\hat y\) in it are
unique (Theorem~\ref{r3-thm-unique-coef}), so the coefficient of \(x_3\) is the same in both fits.
\(M_2\) is symmetric and idempotent (Theorem~\ref{r6-thm-hat}(b)), so \(X_1'M_2X_1=(M_2X_1)'(M_2X_1)=1+1+4+1+1+4=12\) and
\(X_1'M_2y=(M_2X_1)'(M_2y)=2+0+4+2+0+4=12\). By Theorem~\ref{r11-thm-fwl}(a), \(\hat\gamma=12/12=1\). By
Theorem~\ref{r11-thm-fwl}(b), \(e\) is the residual vector of the fit of \(M_2y\) on \(M_2X_1\):
\[
e=M_2y-\hat\gamma M_2X_1=\begin{pmatrix}2\\0\\-2\\-2\\0\\2\end{pmatrix}-\begin{pmatrix}1\\1\\-2\\-1\\-1\\2\end{pmatrix}
=\begin{pmatrix}1\\-1\\0\\-1\\1\\0\end{pmatrix}
\]
Check: \(\one'e=0\), \(x_1'e=(1-1+0)-(-1+1+0)=0\), \(x_2'e=1+1+0-1-1+0=0\) and \(x_3'e=3-3+0+1-1+0=0\); so \(X'e=0\), as
Proposition~\ref{r10-prop-residuals}(a) says.

(c) \(X_1'X_1=9+9+0+1+1+4=24\) and \(X_1'X_2=(X_2'x_3)'=(6,6,0)\), so
\(X_1'X_2(X_2'X_2)^{-1}X_2'X_1=6\cdot\frac16\cdot6+6\cdot\frac16\cdot6=12\), as the third entry of \(X_2'X_1\) is \(0\)
(Theorem~\ref{r1-thm-views}(d)), and
\(X_1'X_1-X_1'X_2(X_2'X_2)^{-1}X_2'X_1=24-12=12=X_1'M_2X_1\).
\end{solution}

\begin{exercise}\label{r11-dr-vif}
Five observations of two variables and a response are
\[
\begin{array}{c|ccccc}x_{i1}&1&2&3&4&5\\x_{i2}&2&3&3&3&4\\\hline y_i&1&3&2&7&7\end{array}
\]
Fit \(y_i=\beta_0+\beta_1x_{i1}+\beta_2x_{i2}+\varepsilon_i\) under (A1)--(A4), with \(X=[\one_5\ X_0]\).
\begin{enumerate}
\item Compute \(C_5X_0\), \(X_0'C_5X_0\) and its inverse, \(C_5y\) and \(X_0'C_5y\), and find \(\hat\beta_1\), \(\hat\beta_2\) and
\(\hat\beta_0\) as in Example~\ref{r11-exm-fwl}.
\item Find \(\var(\hat\beta_1)\) and \(\var(\hat\beta_2)\), the variance inflation factors \(\mathrm{VIF}_1\),
\(\mathrm{VIF}_2\), and \(R_1^2\), \(R_2^2\) (Corollary~\ref{r11-cor-vif}).
\item Compute \(e\), SSE and \(s^2\), and the estimates of \(\var(\hat\beta_1)\) and \(\var(\hat\beta_2)\) with
\(s^2\) in place of \(\sigma^2\).
\end{enumerate}
\end{exercise}
\begin{solution}
(a) \(\bar x=(3,3)'\) and \(\bar y=4\), so (Example~\ref{r1-exm-centering})
\[
C_5X_0=\begin{pmatrix}-2&-1\\-1&0\\0&0\\1&0\\2&1\end{pmatrix},\qquad
X_0'C_5X_0=(C_5X_0)'(C_5X_0)=\begin{pmatrix}10&4\\4&2\end{pmatrix},\qquad
C_5y=\begin{pmatrix}-3\\-1\\-2\\3\\3\end{pmatrix}
\]
\(\det X_0'C_5X_0=20-16=4\), so \((X_0'C_5X_0)^{-1}=\frac14\left(\begin{smallmatrix}2&-4\\-4&10\end{smallmatrix}\right)\)
(Theorem~\ref{r5d-thm-adjugate}). Then
\[
X_0'C_5y=(C_5X_0)'(C_5y)=\begin{pmatrix}6+1+0+3+6\\3+0+0+0+3\end{pmatrix}=\begin{pmatrix}16\\6\end{pmatrix},\qquad
\begin{pmatrix}\hat\beta_1\\\hat\beta_2\end{pmatrix}=\frac14\begin{pmatrix}32-24\\-64+60\end{pmatrix}=\begin{pmatrix}2\\-1\end{pmatrix}
\]
and \(\hat\beta_0=\bar y-\bar x'(\hat\beta_1,\hat\beta_2)'=4-(6-3)=1\).

(b) By Corollary~\ref{r11-cor-vif}(a), with \(X_1=\one_5\) and \(X_2=X_0\),
\(\var\bigl((\hat\beta_1,\hat\beta_2)'\bigr)=\sigma^2(X_0'C_5X_0)^{-1}\): \(\var(\hat\beta_1)=\frac12\sigma^2\) and
\(\var(\hat\beta_2)=\frac{10}4\sigma^2=\frac52\sigma^2\). The diagonal of \(X_0'C_5X_0\) holds \(S_{11}=10\) and \(S_{22}=2\);
by Corollary~\ref{r11-cor-vif}(b), \(\mathrm{VIF}_j=S_{jj}\var(\hat\beta_j)/\sigma^2\): \(\mathrm{VIF}_1=10\cdot\frac12=5\),
\(\mathrm{VIF}_2=2\cdot\frac52=5\), and \(R_1^2=R_2^2=1-\frac15=\frac45\). Without the correlation of the two variables
(\(S_{12}=0\)), the variances would be \(\sigma^2/S_{jj}\): \(\frac1{10}\sigma^2\) and \(\frac12\sigma^2\).

(c) \(\hat y_i=1+2x_{i1}-x_{i2}\) gives \(\hat y=(1,2,4,6,7)'\) and \(e=y-\hat y=(0,1,-2,1,0)'\). Check:
\(X'e=(0+1-2+1+0,\ 0+2-6+4+0,\ 0+3-6+3+0)'=0\). So \(\mathrm{SSE}=0+1+4+1+0=6\) and \(s^2=6/(5-3)=3\)
(Theorem~\ref{r11-thm-sigma}); the estimates of the two variances are \(\frac12s^2=\frac32\) and \(\frac52s^2=\frac{15}2\).
\end{solution}

\begin{exercise}\label{r11-ex-blp-drill}
Let \(x=(x_1,x_2,x_3,x_4)'\) be the random vector of Exercise~\ref{r11-ex-moments-drill}, with \(\expc(x)=\mu=(1,0,2,-1)'\) and
\[
\var(x)=\Sigma=\begin{pmatrix}9&0&3&3\\0&1&2&1\\3&2&9&5\\3&1&5&4\end{pmatrix}
\]
\begin{enumerate}
\item Find the best linear predictor of \(x_1\) from \((x_2,x_3,x_4)'\) and its mean squared prediction error.
\item Check directly that the prediction error \(u\) of (a) has mean \(0\), that \(\cov(u,x_2)=\cov(u,x_3)=\cov(u,x_4)=0\), and
that \(\var(u)\) is the mean squared prediction error.
\end{enumerate}
\end{exercise}
\begin{solution}
(a) In Theorem~\ref{r11-thm-blp} take \(y=x_1\), with \(\mu_y=1\) and \(\sigma_{yy}=9\), and \((x_2,x_3,x_4)'\) as the predictors:
\[
\Sigma_{xx}=\begin{pmatrix}1&2&1\\2&9&5\\1&5&4\end{pmatrix},\qquad\sigma_{xy}=\begin{pmatrix}0\\3\\3\end{pmatrix},\qquad
\Sigma_{xx}\begin{pmatrix}-1\\0\\1\end{pmatrix}=\begin{pmatrix}0\\3\\3\end{pmatrix}
\]
By Exercise~\ref{r11-ex-moments-drill}(a), \(\Sigma=LL'\) with \(\det L=6\ne0\), so \(\Sigma\) is PD
(Theorem~\ref{r8-thm-psd-factor}(c) with \(P=L'\)). Its principal submatrix \(\Sigma_{xx}\) is PD
(Theorem~\ref{r8-thm-principal-sub}), hence invertible (Proposition~\ref{r8-prop-pd-basic}(c)), and
\(b^*=\Sigma_{xx}^{-1}\sigma_{xy}=(-1,0,1)'\). Then \(a^*=\mu_1-b^{*\prime}(\mu_2,\mu_3,\mu_4)'=1-(0+0-1)=2\), the best linear
predictor is \(2-x_2+x_4\), and its mean squared prediction error is \(\sigma_{yy}-\sigma_{yx}\Sigma_{xx}^{-1}\sigma_{xy}=9-(0,3,3)(-1,0,1)'=6\).

(b) \(u=x_1-2+x_2-x_4=c'x-2\) with \(c=(1,1,0,-1)'\). By Theorem~\ref{r11-thm-linear}(b),
\(\expc(u)=c'\mu-2=1+0+0+1-2=0\). By Theorem~\ref{r1-thm-views}(b),
\(\Sigma c=\col_1(\Sigma)+\col_2(\Sigma)-\col_4(\Sigma)=(9+0-3,\ 0+1-1,\ 3+2-5,\ 3+1-4)'=(6,0,0,0)'\). By
Theorem~\ref{r14-thm-cov}(a), with \(A=c'\) and \(B\) the row \(j\) of \(I_4\), \(\cov(u,x_j)=\cov(c'x,x_j)\) is entry \(j\) of
\(\Sigma c\), which is \(0\) for \(j=2,3,4\); and \(\var(u)=c'\Sigma c=6\) (Theorem~\ref{r11-thm-linear}(b)), the mean squared
prediction error of (a).
\end{solution}

\needspace{14\baselineskip}
\subsection*{Problems}

\begin{problem}[\citeat{tu-clm}{Exercise~D2}]\label{r11-ex-four-regions}
Wages \(y_i\) of \(n\) workers are regressed on years of schooling \(x_i\), an intercept, and the indicators
\(d_1,d_2,d_3,d_4\in\R^n\) of all four regions of a country (each worker lives in exactly one region):
\(X=[\one_n\ x\ d_1\ d_2\ d_3\ d_4]\in\R^{n\times6}\). (a) Show that \(X\) does not have full column rank, so
that \(\hat\beta=(X'X)^{-1}X'y\) is not defined. (b) Let \(X_2=[\one_n\ x\ d_2\ d_3\ d_4]\) (indicator \(d_1\)
dropped). Show that \(\mathcal C(X_2)=\mathcal C(X)\), so that \(\rank X_2=\rank X\).
\end{problem}
\begin{solution}
(a) Each worker lives in exactly one region, so \(d_1+d_2+d_3+d_4=\one_n\), that is,
\(X(-1,0,1,1,1,1)'=0\). So \(X\) does not have full column rank (Theorem~\ref{thm:r4-rank-rules}(b)), and
\(X'X\) is singular (Theorems~\ref{thm:r4-rank-ata}(b) and~\ref{thm:r4-rank-rules}(d)): the dummy variable trap.

(b) Every column of \(X_2\) is a column of \(X\), so \(\mathcal C(X_2)\subseteq\mathcal C(X)\). Conversely
\(d_1=\one_n-d_2-d_3-d_4\in\mathcal C(X_2)\), so every column of \(X\) lies in the subspace \(\mathcal C(X_2)\), and
therefore \(\mathcal C(X)\subseteq\mathcal C(X_2)\) (Theorem~\ref{r3-thm-span}). Hence \(\mathcal C(X_2)=\mathcal C(X)\), and
\(\rank X_2=\rank X\); by Proposition~\ref{r11-prop-reparametrization} (with \(X_2=XT\), \(T\) the columns
\(1,2,4,5,6\) of \(I_6\)), the two designs give the same fitted values.
\end{solution}

\begin{problem}[\citeat{tu-clm}{Remark~2.11, Theorem~3.7, Lemma~7.2, Remark~6.2}]\label{r11-ex-hat-identities}
Let \(X\in\R^{n\times(k+1)}\) have full column rank, with hat matrix \(H\) and residual maker \(M\).
\begin{enumerate}
\item In the model \(y=X\beta+\varepsilon\), show that \(\mathrm{SSE}=\varepsilon'M\varepsilon\).
\item Let \(\one_n\in\mathcal C(X)\). Show that \(H\one_n=\one_n\), \(H\bar J_n=\bar J_nH=\bar J_n\) and
\(MC_n=C_nM=M\).
\item Under (b), show that \(H-\bar J_n\) is symmetric and idempotent of rank \(k\), that
\((H-\bar J_n)M=0\), and that \(\mathrm{SSR}=y'(H-\bar J_n)y\).
\end{enumerate}
\end{problem}
\begin{solution}
(a) \(MX=0\) (Theorem~\ref{r10-thm-PX}(a)), so \(e=My=M(X\beta+\varepsilon)=MX\beta+M\varepsilon=M\varepsilon\)
(Theorem~\ref{r10-thm-PX}(d)). As \(M\) is symmetric and idempotent,
\(\mathrm{SSE}=e'e=\varepsilon'M'M\varepsilon=\varepsilon'M\varepsilon\).

(b)
\begin{itemize}
\item \(\one_n=Xc\) for some \(c\), so \(H\one_n=HXc=Xc=\one_n\) (Theorem~\ref{r6-thm-hat}(c)).
\item \(H\bar J_n=\frac1nH\one\one'=\frac1n\one\one'=\bar J_n\), and, as \(H\) and \(\bar J_n\) are symmetric,
\(\bar J_nH=(H\bar J_n)'=\bar J_n\).
\item \(M\bar J_n=\bar J_n-H\bar J_n=0\), so \(MC_n=M(I_n-\bar J_n)=M\), and \(C_nM=(MC_n)'=M\).
\end{itemize}
(c)
\begin{itemize}
\item \(H-\bar J_n\) is symmetric, and by (b) and \(\bar J_n^2=\bar J_n\) (Proposition~\ref{r1-prop-centering}),
\((H-\bar J_n)^2=H^2-H\bar J_n-\bar J_nH+\bar J_n^2=H-\bar J_n-\bar J_n+\bar J_n=H-\bar J_n\).
\item Its rank is its trace (Theorem~\ref{r6-thm-idempotent}(c)):
\(\tr H-\tr\bar J_n=(k+1)-1=k\) (Corollary~\ref{r6-cor-trace-hat}; the diagonal entries of \(\bar J_n\) are \(\frac1n\)).
\item \((H-\bar J_n)M=HM-\bar J_nM=0-(M\bar J_n)'=0\), by Theorem~\ref{r6-thm-hat}(c) and (b).
\item \(\hat y-\bar y\one=Hy-\bar J_ny=(H-\bar J_n)y\) (Example~\ref{r6-exm-centering}), so
\(\mathrm{SSR}=\lVert(H-\bar J_n)y\rVert^2=y'(H-\bar J_n)'(H-\bar J_n)y=y'(H-\bar J_n)y\).
\end{itemize}
\end{solution}

\begin{problem}[\citeat{tu-clm}{Theorem~2.14, Exercises~B1, D3}]\label{r11-ex-leverage}
Let \(X\in\R^{n\times(k+1)}\) have full column rank, and let \(h_{11},\dots,h_{nn}\) be the diagonal entries of its hat matrix
\(H\), the leverages of the \(n\) observations.
\begin{enumerate}
\item Show that \(0\leqslant h_{ii}\leqslant1\).
\item If \(\one_n\in\mathcal C(X)\), show that \(h_{ii}\geqslant\frac1n\).
\item In the model \(y=X\beta+\varepsilon\) under (A2) and (A3), show that \(\var(e_i)=\sigma^2(1-h_{ii})\).
\end{enumerate}
\end{problem}
\begin{solution}
(a) \(H\) is symmetric and idempotent (Theorem~\ref{r6-thm-hat}(b)), so \(H=H^2=H'H\). Comparing the \((i,i)\) entries, with
\(h_{ji}=h_{ij}\),
\(h_{ii}=\sum_{j=1}^nh_{ij}^2=h_{ii}^2+\sum_{j\ne i}h_{ij}^2\). So \(h_{ii}\geqslant0\), as a sum of squares, and
\(h_{ii}\geqslant h_{ii}^2\), that is, \(h_{ii}(1-h_{ii})\geqslant0\); together, \(h_{ii}\leqslant1\).

(b) By Problem~\ref{r11-ex-hat-identities}(c), \(H-\bar J_n\) is symmetric and idempotent, and its \((i,j)\) entry is
\(h_{ij}-\frac1n\). As in (a), \(h_{ii}-\frac1n=\sum_{j=1}^n\bigl(h_{ij}-\frac1n\bigr)^2\geqslant0\).

(c) By Theorem~\ref{r11-thm-unbiased}(c), \(\var(e)=\sigma^2M\). Its \((i,i)\) entry is \(\var(e_i)=\sigma^2(1-h_{ii})\), as
\(M=I_n-H\).
\end{solution}

\begin{problem}[\citeat{tu-clm}{Exercise~D4}]\label{r11-ex-orthogonal-regressor}
Let \(X_1\in\R^{n\times p}\) have full column rank, let \(x\in\R^n\) be nonzero with \(X_1'x=0\), and let \(X_+=[X_1\ x]\). For
\(y\in\R^n\), let \(\tilde\gamma=(X_1'X_1)^{-1}X_1'y\) and \(\mathrm{SSE}_1\) belong to the least-squares fit with design \(X_1\),
and \((\hat\gamma',\hat\theta)'\) (\(\hat\gamma\in\R^p\)) and \(\mathrm{SSE}_+\) to the fit with design \(X_+\).
\begin{enumerate}
\item Show that \(\hat\theta=x'y/x'x\).
\item Show that \(\hat\gamma=\tilde\gamma\): the coefficients of \(X_1\) do not change when \(x\) is added.
\item Show that \(\mathrm{SSE}_1-\mathrm{SSE}_+=(x'y)^2/x'x\).
\item Relate (b) to Theorem~\ref{r11-thm-ovb}.
\end{enumerate}
\end{problem}
\begin{solution}
\(X_+\) has full column rank: if \(X_1a+cx=0\), then multiplying by \(x'\) gives \(c\,x'x=0\), as \(x'X_1=0\), so \(c=0\); then
\(X_1a=0\) gives \(a=0\) (Theorem~\ref{thm:r4-rank-rules}(b)). Let \(H_1\), \(M_1\) be the hat matrix and the residual maker of
\(X_1\) (Theorem~\ref{r6-thm-hat}). As \(X_1'x=0\), \(H_1x=X_1(X_1'X_1)^{-1}X_1'x=0\), so \(M_1x=x\), \(x'M_1x=x'x>0\) and
\(x'M_1y=(M_1x)'y=x'y\).

(a) By Theorem~\ref{r11-thm-fwl} for the design \(X_+\), with the roles of the blocks exchanged (\(X_2=x\)),
\(\hat\theta=(x'M_1x)^{-1}x'M_1y=x'y/x'x\).

(b) By the block product (Theorem~\ref{r5b-thm-block}), the first block of the normal equations of \(X_+\) reads
\(X_1'X_1\hat\gamma+X_1'x\hat\theta=X_1'y\). As \(X_1'x=0\), \(\hat\gamma=(X_1'X_1)^{-1}X_1'y=\tilde\gamma\).

(c)
\begin{itemize}
\item By (b), the fitted values of \(X_+\) are \(\hat y_+=X_1\tilde\gamma+x\hat\theta=\hat y_1+x\hat\theta\), with
\(\hat y_1=X_1\tilde\gamma\); so the residual vectors satisfy \(e_+=e_1-x\hat\theta\).
\item \(x'e_1=x'(y-X_1\tilde\gamma)=x'y\), as \(x'X_1=0\). Hence
\[
\mathrm{SSE}_+=e_1'e_1-2\hat\theta\,x'e_1+\hat\theta^2x'x=\mathrm{SSE}_1-2\frac{(x'y)^2}{x'x}+\frac{(x'y)^2}{x'x}
=\mathrm{SSE}_1-\frac{(x'y)^2}{x'x}
\]
\end{itemize}
(d) Read the fit with \(X_+\) as the true model \(y=X_1\gamma+x\theta+\varepsilon\), \(\expc(\varepsilon)=0\), and the fit
with \(X_1\) as the one that omits \(x\). By Theorem~\ref{r11-thm-ovb} with \(X_2=x\), the omitted variable bias
\((X_1'X_1)^{-1}X_1'x\,\theta\) is \(0\) for every \(\theta\), as \(X_1'x=0\): omitting \(x\) does not bias \(\tilde\gamma\). Part (b)
is the exact statement about the estimates behind this statement about expectations.
\end{solution}

\begin{problem}[\citeat{tu-asm}{Exercise~5.8}, \citeat{islr}{\S12.6, Exercise~6}]\label{r12-ex-regression}
Let \(X\), \(S\), \(\lambda_m\), \(v_m=(v_{1m},\dots,v_{pm})'\), \(z_m\) and \(r\) be as in Theorem~\ref{r12-thm-pca}, let
\(M\leqslant r\), let \(V_M\) and \(Z_M\) be as in Theorem~\ref{r12-thm-closest}, and let \(c_j\) be column \(j\) of
\(C_nX\). Suppose that the first \(M\) score vectors are known, and take \(A=Z_M\) in the problem of
Theorem~\ref{r12-thm-closest}(a): minimize \(\lVert C_nX-Z_MB'\rVert_F^2\) over \(B\in\R^{p\times M}\). Show that this
problem splits into \(p\) separate least-squares problems, one for each column \(c_j\) as the response, with the
design \(Z_M\); and that their solutions are the rows of \(V_M\), so that \(B=V_M\). So the first \(M\) loading vectors
can be recovered from the first \(M\) score vectors.
\end{problem}
\begin{solution}
\begin{itemize}
\item Let \(b_j'\) be row \(j\) of \(B\). Column \(j\) of \(Z_MB'\) is \(Z_Mb_j\) (Theorem~\ref{r1-thm-views}(a)), so column
\(j\) of \(C_nX-Z_MB'\) is \(c_j-Z_Mb_j\), and
\(\lVert C_nX-Z_MB'\rVert_F^2=\sum_{j=1}^p\lVert c_j-Z_Mb_j\rVert^2\) (Definition~\ref{r9-def-frobenius}). The \(j\)th term
involves only row \(j\) of \(B\); so \(B\) minimizes the sum if and only if each \(b_j\) minimizes
\(\lVert c_j-Z_Mb_j\rVert^2\), the least-squares criterion for the design \(Z_M\) and the response \(c_j\).
\item \(Z_M'Z_M=(n-1)\diag(\lambda_1,\dots,\lambda_M)\) with \(\lambda_1,\dots,\lambda_M>0\) (Theorem~\ref{r12-thm-pca}(a),(d)) is
invertible, so \(\rank Z_M=M\) (Theorems~\ref{thm:r4-rank-ata}(b) and~\ref{thm:r4-rank-rules}(d)), and the \(j\)th problem
has exactly one solution, \(\hat b_j=(Z_M'Z_M)^{-1}Z_M'c_j\) (Theorem~\ref{r10-thm-least-squares}(c)).
\item \(SV_M=V_M\diag(\lambda_1,\dots,\lambda_M)\) and \(S'=S\), so
\(Z_M'C_nX=V_M'(C_nX)'(C_nX)=(n-1)V_M'S=(n-1)\diag(\lambda_1,\dots,\lambda_M)V_M'\). Column \(j\) of this reads
\(Z_M'c_j=(n-1)\diag(\lambda_1,\dots,\lambda_M)(v_{j1},\dots,v_{jM})'\).
\item Hence \(\hat b_j=\bigl((n-1)\diag(\lambda_1,\dots,\lambda_M)\bigr)^{-1}(n-1)\diag(\lambda_1,\dots,\lambda_M)(v_{j1},\dots,v_{jM})'
=(v_{j1},\dots,v_{jM})'\), row \(j\) of \(V_M\) transposed: the minimizing \(B\) is \(V_M\).
\end{itemize}
\end{solution}

\begin{problem}[\citeat{tu-clm}{Proposition~7.14}, \citeat{searle-khuri}{(11.40)--(11.41)}]\label{r12-pr-expected-ssr}
Assume (A1)--(A4) with \(X=[\one_n\ X_0]\in\R^{n\times(k+1)}\), \(k\geqslant1\) and \(n>k+1\), and write
\(\beta=(\beta_0,\beta_*')'\) with \(\beta_*=(\beta_1,\dots,\beta_k)'\). Show that
\[
\expc(\mathrm{SSR})=k\sigma^2+\lVert C_nX_0\beta_*\rVert^2
\]
and that \(\expc(\mathrm{SSR}/k)\geqslant\sigma^2=\expc(s^2)\), with equality if and only if \(\beta_*=0\).
\end{problem}
\begin{solution}
\begin{itemize}
\item By Problem~\ref{r11-ex-hat-identities}(c), \(\mathrm{SSR}=y'Py\) with \(P=H-\bar J_n\) symmetric and idempotent of rank
\(k\).
\item By (A1) and (A3), \(\expc(y)=X\beta\) and \(\var(y)=\sigma^2I_n\) (the paragraph after Definition~\ref{r11-def-model}). So
Example~\ref{r11-exm-quadratic-projection} with \(\mu=X\beta\) gives
\(\expc(\mathrm{SSR})=k\sigma^2+\beta'X'PX\beta=k\sigma^2+\lVert PX\beta\rVert^2\), as \(P=P'P\).
\item \(HX=X\) (Theorem~\ref{r10-thm-PX}(a)), so \((H-\bar J_n)X\beta=C_nX\beta=C_n(\beta_0\one_n+X_0\beta_*)=C_nX_0\beta_*\), as
\(C_n\one_n=0\) (Proposition~\ref{r1-prop-centering}).
\item Hence \(\expc(\mathrm{SSR}/k)=\sigma^2+\lVert C_nX_0\beta_*\rVert^2/k\geqslant\sigma^2=\expc(s^2)\)
(Theorem~\ref{r11-thm-sigma}). Equality holds if and only if
\(C_nX_0\beta_*=X_0\beta_*-(\bar x'\beta_*)\one_n=X(-\bar x'\beta_*,\ \beta_*')'\) vanishes, where
\(\bar x=\frac1nX_0'\one_n\); as \(\rank X=k+1\), this means \(\beta_*=0\) (Theorem~\ref{thm:r4-rank-rules}(b)).
\end{itemize}
So \(\mathrm{SSR}/k\) estimates \(\sigma^2\) plus a nonnegative term that vanishes exactly when \(\beta_*=0\).
\end{solution}

\clearpage
\def\unitnumber{14}
\setcounter{theorem}{0}\setcounter{problem}{0}
\section*{Chapter 14. Estimability, inference, and the analysis of variance}

\subsection*{Estimable functions}

\begin{definition}[\citeat{searle-khuri}{Definition~12.1}, \citeat{bapat}{\S7.2}]\label{r11-def-estimable}
Consider the linear model \(y=X\beta+\varepsilon\) under (A1), with \(X\in\R^{n\times(k+1)}\) of any rank, and let
\(a\in\R^{k+1}\). A \textbf{linear estimator} is a random variable \(c'y\) with a constant \(c\in\R^n\); it is a
\textbf{linear unbiased estimator} of \(a'\beta\) if \(\expc(c'y)=a'\beta\) for every \(\beta\in\R^{k+1}\). The
function \(a'\beta\) is \textbf{estimable} if it has a linear unbiased estimator. A linear unbiased estimator of
\(a'\beta\) whose variance is the smallest among all linear unbiased estimators of \(a'\beta\) is a \textbf{best linear
unbiased estimator} (\textbf{BLUE}) of \(a'\beta\).
\end{definition}

\needspace{12\baselineskip}
\begin{theorem}[\citeat{searle-khuri}{Theorems~12.1, 12.2, Corollary~12.1}, \citeat{bapat}{\S7.2}]\label{r11-thm-estimable}
Assume (A1), and let \(a\in\R^{k+1}\).
\begin{enumerate}
\item A linear estimator \(c'y\) is unbiased for \(a'\beta\) if and only if \(X'c=a\). More generally, for constant
\(A\in\R^{m\times n}\) and \(C\in\R^{m\times(k+1)}\), \(\expc(Ay)=C\beta\) for every \(\beta\in\R^{k+1}\) if and only if
\(AX=C\). Hence
\[
a'\beta\text{ is estimable}\iff a\in\mathcal C(X')\iff a\perp\mathcal N(X)\iff\rank\begin{pmatrix}X\\a'\end{pmatrix}=\rank X
\]
In particular, if \(\rank X=k+1\), every \(a'\beta\) is estimable.
\item If \(a=X'c\), then \(a'\hat\beta=c'\hat y=c'Hy\) for every least-squares estimate \(\hat\beta\): it is the same
linear estimator for all of them (for instance for \(\hat\beta=(X'X)^-X'y\) with any g-inverse), and
\(\expc(a'\hat\beta)=a'\beta\).
\item If \(a'\beta\) is not estimable, then \(a'\hat\beta\) takes every real value as \(\hat\beta\) runs through the
least-squares estimates.
\end{enumerate}
\end{theorem}
\begin{proof}
(a)
\begin{itemize}
\item By Theorem~\ref{r11-thm-linear}(b) and (A1), \(\expc(c'y)=c'X\beta\). If \(c'X\beta=a'\beta\) for every
\(\beta\), then taking for \(\beta\) the columns of \(I_{k+1}\) gives \(c'X=a'\), that is, \(X'c=a\). Conversely, if
\(X'c=a\), then \(\expc(c'y)=c'X\beta=a'\beta\) for every \(\beta\). In the same way,
\(\expc(Ay)=AX\beta\), and taking for \(\beta\) the columns of \(I_{k+1}\) shows that \(AX\beta=C\beta\) for every \(\beta\) if
and only if \(AX=C\) (Theorem~\ref{r1-thm-views}(a),(b)). So \(a'\beta\) is estimable if and only if \(a=X'c\) for some
\(c\), that is, \(a\in\mathcal C(X')\).
\item \(\mathcal N(X)=\mathcal C(X')^\perp\) (Corollary~\ref{r6-cor-four-subspaces}), so every vector of
\(\mathcal C(X')\) is orthogonal to \(\mathcal N(X)\): \(\mathcal C(X')\subseteq\mathcal N(X)^\perp\). Both have dimension
\(\rank X\): \(\dim\mathcal C(X')=\rank X'=\rank X\) (Theorem~\ref{thm:r4-rank-rules}(a)), and
\(\dim\mathcal N(X)^\perp=(k+1)-\dim\mathcal N(X)=\rank X\) (Theorems~\ref{r6-thm-complement}(c)
and~\ref{thm:r4-rank-nullity}). So \(\mathcal C(X')=\mathcal N(X)^\perp\) (Theorem~\ref{thm:r4-counting}(a)).
\item \(a\in\mathcal C(X')\) if and only if \(\rank[X'\ a]=\rank X'\) (Corollary~\ref{cor:r4-ax-b}(a)). Transposing
(Theorems~\ref{r5b-thm-block}(b) and~\ref{thm:r4-rank-rules}(a)) gives the rank condition.
\item If \(\rank X=k+1\), then \(\rank X'=k+1\), the number of rows of \(X'\), so \(\mathcal C(X')=\R^{k+1}\)
(Theorem~\ref{thm:r4-rank-rules}(a),(c)).
\end{itemize}
(b) Every least-squares estimate gives \(X\hat\beta=\hat y=Hy\) (Theorem~\ref{r10-thm-PX}(d)), so \(a'\hat\beta=c'X\hat\beta=c'Hy\). By Theorem~\ref{r11-thm-linear}(b), (A1) and
\(HX=X\) (Theorem~\ref{r10-thm-PX}(a)), \(\expc(c'Hy)=c'HX\beta=c'X\beta=a'\beta\).

(c) By (a), \(a\) is not orthogonal to \(\mathcal N(X)\): \(a'z\ne0\) for some \(z\in\mathcal N(X)\). If \(\hat\beta\) is a
least-squares estimate, so is \(\hat\beta+tz\) for every \(t\in\R\), since \(Xz=0\) gives \(X'X(\hat\beta+tz)=X'X\hat\beta=X'y\)
(Theorem~\ref{r10-thm-least-squares}(a)); and \(a'(\hat\beta+tz)=a'\hat\beta+t\,a'z\) takes every real value.
\end{proof}

\begin{corollary}[\citeat{harville}{Corollary~11.10.2}, \citeat{searle-khuri}{Theorems~12.2, 12.4}]\label{r10-cor-estimable-algebra}
Let \(X\in\R^{n\times(k+1)}\) and \(y\in\R^n\), let \(G\) be a g-inverse of \(X'X\), and let \(\lambda\in\R^{k+1}\). The
following are equivalent:
\begin{enumerate}[label=(\roman*)]
\item \(\lambda'\hat\beta\) has the same value for every least-squares estimate \(\hat\beta\);
\item \(\lambda\in\mathcal C(X')\), that is, \(\lambda'=c'X\) for some \(c\in\R^n\);
\item \(\lambda'GX'X=\lambda'\).
\end{enumerate}
If they hold, then \(\lambda'\hat\beta=c'\hat y=c'Hy\) for every least-squares estimate \(\hat\beta\). By
Theorem~\ref{r11-thm-estimable}(a), (ii) says that \(\lambda'\beta\) is estimable.
\end{corollary}
\begin{proof}
\begin{itemize}
\item The least-squares estimates are the solutions of the consistent system \(X'X\beta=X'y\)
(Theorems~\ref{r10-thm-least-squares}(a) and~\ref{r10-thm-normal-equations}(a)); so \(X'y\in\mathcal C(X'X)\)
(Corollary~\ref{cor:r4-ax-b}(a)).
\item Theorem~\ref{r10-thm-invariant-combination}, for the matrix \(X'X\) and the right side \(X'y\), shows that (i) holds
if and only if \(\lambda\in\mathcal C((X'X)')=\mathcal C(X'X)\), and if and only if \(\lambda'GX'X=\lambda'\). Here
\(\mathcal C(X'X)=\mathcal C(X')\) (Theorem~\ref{thm:r4-rank-ata}(c)).
\item If \(\lambda'=c'X\), then \(\lambda'\hat\beta=c'X\hat\beta=c'\hat y=c'Hy\)
(Theorem~\ref{r10-thm-PX}(d)).\qedhere
\end{itemize}
\end{proof}

\begin{example}[cf.\ \citeat{searle-khuri}{Example~12.1}, \citeat{bapat}{Ch.~7, Exercise~11}]\label{r11-exm-estimable}
In the one-way model of Example~\ref{r11-exm-oneway} (\(g=3\); \(n_1,n_2,n_3=3,2,1\)),
\(\mathcal N(X)=\Span\{(-1,1,1,1)'\}\). By Theorem~\ref{r11-thm-estimable}(a),
\(a'\beta=a_0\mu+a_1\alpha_1+a_2\alpha_2+a_3\alpha_3\) is estimable if and only if \(-a_0+a_1+a_2+a_3=0\).
\begin{itemize}
\item \(\mu+\alpha_1\) (\(a=(1,1,0,0)'\)) and \(\alpha_1-\alpha_2\) (\(a=(0,1,-1,0)'\)) are estimable. On the least-squares
estimates \((-t,6+t,4+t,2+t)'\) they take the values \(6\) and \(2\) for every \(t\).
\item \(\mu\) (\(a=(1,0,0,0)'\)) and \(\alpha_1\) (\(a=(0,1,0,0)'\)) are not estimable: \(-1\ne0\) and \(1\ne0\). Their
values \(-t\) and \(6+t\) change with \(t\), as Theorem~\ref{r11-thm-estimable}(c) says.
\item For every \(y\), with the group means \(\bar y_1,\bar y_2,\bar y_3\), the vector
\(\hat\beta=(0,\bar y_1,\bar y_2,\bar y_3)'\) is a least-squares estimate (Theorem~\ref{r10-thm-least-squares}(a)):
\(X'X\hat\beta=(3\bar y_1+2\bar y_2+\bar y_3,\ 3\bar y_1,\ 2\bar y_2,\ \bar y_3)'\) is the vector of the total and the group
totals, which is \(X'y\). So the estimator of \(\alpha_1-\alpha_2\) is
\[
\bar y_1-\bar y_2=\tfrac13(y_1+y_2+y_3)-\tfrac12(y_4+y_5)=c'y,\qquad c=\bigl(\tfrac13,\tfrac13,\tfrac13,-\tfrac12,-\tfrac12,0\bigr)'
\]
and indeed \(X'c=(1-1,\ 1,\ -1,\ 0)'=(0,1,-1,0)'\).
\end{itemize}
\end{example}

\needspace{12\baselineskip}
\begin{proposition}[\citeat{searle-khuri}{\S12.4}]\label{r11-prop-solution}
Let \(G\) be a g-inverse of \(X'X\), and let \(\beta^0=GX'y\), a least-squares estimate
(Theorems~\ref{r10-thm-normal-equations}(c) and~\ref{r10-thm-least-squares}(a)).
\begin{enumerate}
\item Under (A1), \(\expc(\beta^0)=GX'X\beta\). If \(\rank X<k+1\), then \(GX'X\beta\ne\beta\) for some
\(\beta\in\R^{k+1}\): \(\beta^0\) is not an unbiased estimator of \(\beta\).
\item Under (A3), \(\var(\beta^0)=\sigma^2GX'XG'\).
\end{enumerate}
\end{proposition}
\begin{proof}
(a)
\begin{itemize}
\item By Theorem~\ref{r11-thm-linear}(b) with \(A=GX'\) and (A1),
\(\expc(\beta^0)=GX'\expc(y)=GX'X\beta\).
\item If \(GX'X\beta=\beta\) for every \(\beta\), then taking for \(\beta\) the columns of \(I_{k+1}\) gives \(GX'X=I_{k+1}\), so
\(k+1=\rank(GX'X)\leqslant\rank(X'X)=\rank X\) (Theorems~\ref{thm:r4-rank-product}(a) and~\ref{thm:r4-rank-ata}(b)).
Hence \(\rank X<k+1\) gives \(GX'X\beta\ne\beta\) for some \(\beta\).
\end{itemize}
(b) By Theorem~\ref{r11-thm-linear}(b) and (A3),
\(\var(\beta^0)=GX'(\sigma^2I_n)(GX')'=\sigma^2GX'XG'\) (Theorem~\ref{r1-thm-transpose}).
\end{proof}

For \(X\) of full column rank, \(G=(X'X)^{-1}\) is the only g-inverse of \(X'X\) (Theorem~\ref{r10-thm-PX}(c)), \(GX'X=I_{k+1}\)
and \(\beta^0=\hat\beta\), and Proposition~\ref{r11-prop-solution} is Theorem~\ref{r11-thm-unbiased}(a),(b).

\begin{example}[cf.\ \citeat{searle-khuri}{Example~12.1}]\label{r11-exm-solution}
In the one-way model of Example~\ref{r11-exm-oneway} (\(g=3\); \(n_1,n_2,n_3=3,2,1\)), \(G=\diag(0,\frac13,\frac12,1)\) is a
g-inverse of \(X'X\), with
\[
GX'X=\begin{pmatrix}0&0&0&0\\1&1&0&0\\1&0&1&0\\1&0&0&1\end{pmatrix}
\]
\begin{itemize}
\item \(X'y\) holds the total and the group totals \(3\bar y_1\), \(2\bar y_2\), \(\bar y_3\), so
\(\beta^0=GX'y=(0,\bar y_1,\bar y_2,\bar y_3)'\), the least-squares estimate of Example~\ref{r11-exm-estimable}. By
Proposition~\ref{r11-prop-solution}(a),
\[
\expc(\beta^0)=GX'X\beta=(0,\ \mu+\alpha_1,\ \mu+\alpha_2,\ \mu+\alpha_3)'
\]
So \(\bar y_l\) estimates the mean \(\mu+\alpha_l\) of group \(l\), not \(\alpha_l\); for the estimable \(\alpha_1-\alpha_2\),
\(\expc(\bar y_1-\bar y_2)=\alpha_1-\alpha_2\), as Theorem~\ref{r11-thm-estimable}(b) says.
\item Multiplying the columns of \(GX'X\) by \(0,\frac13,\frac12,1\) (Theorem~\ref{r1-thm-views}(a),(b)) gives \(GX'XG'=GX'XG=\diag(0,\frac13,\frac12,1)\). So
\(\var(\beta^0)=\sigma^2\diag(0,\frac13,\frac12,1)\) by Proposition~\ref{r11-prop-solution}(b): the group means
are uncorrelated, with \(\var(\bar y_l)=\sigma^2/n_l\). For the estimable \(\alpha_1-\alpha_2\), with
\(a=(0,1,-1,0)'\), Theorem~\ref{r11-thm-linear}(b) gives
\(\var(\bar y_1-\bar y_2)=\var(a'\beta^0)=a'\var(\beta^0)a=\sigma^2(\frac13+\frac12)
=\frac56\sigma^2\), which is \(\sigma^2a'Ga\).
\end{itemize}
\end{example}

\begin{theorem}[\citeat{searle-khuri}{Theorem~12.3, Exercise~12.4}, \citeat{bapat}{Result~7.1}]\label{r11-thm-blue}
Assume (A1)--(A3), with \(X\) of any rank, and let \(a'\beta\) be estimable. Then \(a'\hat\beta\), the same for every
least-squares estimate \(\hat\beta\), is the BLUE of \(a'\beta\): it is a linear unbiased estimator of \(a'\beta\), and every
linear unbiased estimator \(c'y\) of \(a'\beta\) has \(\var(c'y)\geqslant\var(a'\hat\beta)\), with
equality only if \(c'y=a'\hat\beta\) for every \(y\). Moreover \(\var(a'\hat\beta)=\sigma^2a'(X'X)^-a\) for every
g-inverse \((X'X)^-\) of \(X'X\).
\end{theorem}
\begin{proof}
\begin{itemize}
\item Let \(c'y\) be a linear unbiased estimator of \(a'\beta\). Then \(X'c=a\), and \(a'\hat\beta=c'Hy\) is a linear
unbiased estimator of \(a'\beta\) (Theorem~\ref{r11-thm-estimable}(a),(b)).
\item By Theorem~\ref{r11-thm-linear}(b) and (A3), \(\var(c'y)=\sigma^2c'c\) and
\(\var(c'Hy)=\sigma^2c'HH'c=\sigma^2c'Hc\), as \(H\) is symmetric and idempotent
(Theorem~\ref{r10-thm-PX}(a)). So
\[
\var(c'y)-\var(a'\hat\beta)=\sigma^2c'(I_n-H)c=\sigma^2c'M'Mc=\sigma^2\lVert Mc\rVert^2\geqslant0
\]
\item Equality holds only if \(Mc=0\), that is, \(c=Hc\); then \(c'y=c'Hy=a'\hat\beta\) for every \(y\).
\item For every g-inverse, \(H=X(X'X)^-X'\) (Theorem~\ref{r10-thm-PX}), so
\(c'Hc=c'X(X'X)^-X'c=a'(X'X)^-a\), as in Corollary~\ref{r10-cor-invariant-quadratic} with \(A=X\).\qedhere
\end{itemize}
\end{proof}

The proof shows that \(\var(c'y)=\var(a'\hat\beta)+\sigma^2\lVert Mc\rVert^2\) for every linear unbiased estimator
\(c'y\) of an estimable \(a'\beta\). Under (A4), every \(a'\beta\) is estimable (Theorem~\ref{r11-thm-estimable}(a)),
\(a'\hat\beta=a'(X'X)^{-1}X'y\), and this is the scalar form of the Gauss--Markov theorem~\citeat{searle-khuri}{Theorem~11.1}: for \(X\) of full column rank,
Theorem~\ref{r11-thm-blue} is Theorem~\ref{r11-thm-gauss-markov} for one function \(a'\beta\) at a time. For
\(a'\beta=\beta_j\), \(\var(c'y)\geqslant\var(\hat\beta_j)\), with equality only if \(c'y=\hat\beta_j\) for every \(y\).

\subsection*{The normal linear model}

\begin{definition}[\citeat{tu-clm}{Definition~1.2}, \citeat{casella-berger}{Definitions~8.1.1--8.1.3, 8.3.6, 11.3.5}]\label{r12-def-inference}
The \textbf{normal linear model} is the linear model \(y=X\beta+\varepsilon\) of Definition~\ref{r11-def-model}, with
\(X\in\R^{n\times(k+1)}\) of any rank, under the assumption
\begin{enumerate}[label=(A\arabic*)]\setcounter{enumi}{4}
\item \(\varepsilon\sim N_n(0,\sigma^2I_n)\) for a number \(\sigma^2>0\).
\end{enumerate}
Then (A2) and (A3) hold (Definition~\ref{r14-def-mvn}), \(y\sim N_n(X\beta,\sigma^2I_n)\) (Theorem~\ref{r14-thm-mvn-linear}),
so (A1) holds, and \(y_1,\dots,y_n\) are independent with \(y_i\sim N(x_i'\beta,\sigma^2)\)
(Theorem~\ref{r14-thm-mvn-iid}(a)). In general, let \(\alpha\in(0,1)\), and let the distribution of a random vector \(y\)
depend on unknown parameters \(\vartheta\).
\begin{itemize}
\item The \textbf{likelihood} is the density \(f(y;\vartheta)\) of \(y\) at the observed \(y\), as a function of
\(\vartheta\), and the \textbf{log-likelihood} is its logarithm; values of \(\vartheta\) at which the likelihood is largest
are \textbf{maximum likelihood estimates} (MLE), as in Definition~\ref{r14-def-likelihood}.
\item Confidence intervals of level \(1-\alpha\) for real functions of the parameters, such as \(a'\beta\), are those of
Definition~\ref{r15-def-confidence}(a). For a further random variable \(y_0\) (a future observation), an interval
\([L,U]\) whose endpoints are computed from \(y\) is a \textbf{prediction interval} of level \(1-\alpha\) for \(y_0\) if
\(\prb(L\leqslant y_0\leqslant U)=1-\alpha\) for all values of the parameters.
\item A \textbf{hypothesis} \(H_0\) is a condition on the parameters, such as \(K\beta=m\). A \textbf{test} of \(H_0\) is
given by a \textbf{rejection region} \(\mathcal R\), a set of possible values of \(y\): \(H_0\) is rejected if the
observed \(y\) lies in \(\mathcal R\). The test has \textbf{level} \(\alpha\) if \(\prb(y\in\mathcal R)=\alpha\) for all
parameter values that satisfy \(H_0\). Usually \(\mathcal R\) is described through a \textbf{test statistic} \(T\)
computed from \(y\), as \(\{T>c\}\) or \(\{\lvert T\rvert>c\}\). When \(H_0\) is rejected at level \(\alpha\), the
departure from \(H_0\) is called \textbf{significant} at level \(\alpha\).
\item The quantiles \(t_{\nu,\tau}\), \(F_{a,b,\tau}\) and \(\chi^2_{\nu,\tau}\) (\(0<\tau<1\)) are those of
Definition~\ref{r15-def-t-F}(c).
\end{itemize}
\end{definition}

\begin{theorem}[\citeat{searle-khuri}{(12.46)--(12.47)}, \citeat{bapat}{Result~8.7}, \citeat{tu-clm}{Theorems~5.1--5.3}]\label{r12-thm-sampling}
In the normal linear model let \(\rank X=r<n\), let \(G\) be a g-inverse of \(X'X\), and let \(\beta^0=GX'y\), a
least-squares estimate (Theorems~\ref{r10-thm-normal-equations}(c) and~\ref{r10-thm-least-squares}(a)). Let
\(K\in\R^{q\times(k+1)}\) be such that every entry of \(K\beta\) is estimable.
\begin{enumerate}
\item \(K\beta^0\sim N_q(K\beta,\sigma^2KGK')\), and \(K\beta^0\) and \(KGK'\) are the same for every g-inverse \(G\) of
\(X'X\). Under (A4), \(\hat\beta=(X'X)^{-1}X'y\sim N_{k+1}(\beta,\sigma^2(X'X)^{-1})\).
\item \(\mathrm{SSE}/\sigma^2\sim\chi^2_{n-r}\); in particular \(\prb(\mathrm{SSE}=0)=0\).
\item \(K\beta^0\) and \(\mathrm{SSE}\) are independent.
\item The log-likelihood is
\[
\ell(\beta,\sigma^2)=-\frac n2\log(2\pi\sigma^2)-\frac1{2\sigma^2}\lVert y-X\beta\rVert^2\qquad(\beta\in\R^{k+1},\ \sigma^2>0)
\]
If \(\mathrm{SSE}>0\), its maximizers are the pairs \((\hat\beta,\mathrm{SSE}/n)\) with \(\hat\beta\) a least-squares
estimate.
\end{enumerate}
\end{theorem}
\begin{proof}
(a)
\begin{itemize}
\item Each row \(k_i'\) of \(K\) lies in \(\mathcal C(X')\) (Theorem~\ref{r11-thm-estimable}(a)): \(k_i'=t_i'X\) with
\(t_i\in\R^n\). So \(K=TX\), where \(T\in\R^{q\times n}\) has the rows \(t_i'\) (Theorem~\ref{r1-thm-views}(c)).
\item With \(H=XGX'\), \(K\beta^0=TXGX'y=THy\). By (A5), \(y\sim N_n(X\beta,\sigma^2I_n)\) (Definition~\ref{r12-def-inference}),
so Theorem~\ref{r14-thm-mvn-linear} gives \(K\beta^0\sim N_q(THX\beta,\ \sigma^2THH'T')\). Here \(HX=X\) and \(HH'=H\)
(Theorem~\ref{r10-thm-PX}(a)), so \(THX\beta=TX\beta=K\beta\) and \(THH'T'=THT'=TXGX'T'=KGK'\).
\item \(H\) does not depend on \(G\) (Theorem~\ref{r10-thm-PX}), so neither do \(K\beta^0=THy\) and \(KGK'=THT'\).
\item Under (A4), \(G=(X'X)^{-1}\) (Theorem~\ref{r10-thm-PX}(c)) and every \(a'\beta\) is estimable
(Theorem~\ref{r11-thm-estimable}(a)); take \(K=I_{k+1}\).
\end{itemize}
(b) By Theorem~\ref{r10-thm-PX}(a),(b),(d), \(\mathrm{SSE}=y'My\) with \(M\) symmetric and idempotent of rank \(n-r\geqslant1\),
and \(MX=0\). Theorem~\ref{r15-thm-quadratic-chisq}, for \(y\sim N_n(X\beta,\sigma^2I_n)\) and \(A=M\), gives
\(\mathrm{SSE}/\sigma^2\sim\chi^2_{n-r}\), as \(MX\beta=0\). This distribution has a density
(Theorem~\ref{r15-thm-chisq}(b)), so \(\prb(\mathrm{SSE}=0)=0\).

(c) By (a), \(K\beta^0=Cy\) with \(C=TH\), and \(CM=T(HM)=0\) (Theorem~\ref{r10-thm-PX}(a)). By
Theorem~\ref{r15-thm-indep}(b), \(Cy\) and \(y'My=\mathrm{SSE}\) are independent.

(d)
\begin{itemize}
\item \(\sigma^2I_n\) is PD with determinant \(\sigma^{2n}\) and inverse \(\sigma^{-2}I_n\), so by
Theorem~\ref{r14-thm-mvn-density}(c) the density of \(y\) is
\((2\pi\sigma^2)^{-n/2}\exp\bigl(-\lVert y-X\beta\rVert^2/(2\sigma^2)\bigr)\), whose logarithm is \(\ell\).
\item For fixed \(\sigma^2\), \(\ell\) is largest exactly when \(\lVert y-X\beta\rVert^2\) is smallest, that is, at the
least-squares estimates \(\hat\beta\); and \(\lVert y-X\hat\beta\rVert^2=\mathrm{SSE}\) for all of them
(Theorem~\ref{r10-thm-least-squares}(a), Definition~\ref{r11-def-sums-of-squares}).
\item So it remains to maximize \(h(u)=\ell(\hat\beta,u)=-\frac n2\log(2\pi)-\frac n2\log u-\mathrm{SSE}/(2u)\) over \(u>0\). Its
derivative \(h'(u)=-\frac n{2u}+\frac{\mathrm{SSE}}{2u^2}=\frac{\mathrm{SSE}-nu}{2u^2}\) vanishes only at
\(\hat u=\mathrm{SSE}/n\). As \(\mathrm{SSE}>0\), \(h(u)\to-\infty\) both as \(u\to0^+\) (the term \(-\mathrm{SSE}/(2u)\)
dominates) and as \(u\to\infty\) (the term \(-\frac n2\log u\) dominates); a differentiable function on \((0,\infty)\) with
these limits attains its largest value at a zero of its derivative, which is \(\hat u\), and nowhere else.
\item Hence \(\ell(\beta,\sigma^2)\leqslant\ell(\hat\beta,\sigma^2)\leqslant\ell(\hat\beta,\hat u)\), with equality throughout
only if \(\beta\) is a least-squares estimate and \(\sigma^2=\hat u\).\qedhere
\end{itemize}
\end{proof}

By Theorem~\ref{r11-thm-sigma}, \(\expc(\mathrm{SSE}/n)=\sigma^2(n-r)/n<\sigma^2\): the maximum likelihood estimator of
\(\sigma^2\) is biased~\citeat{tu-clm}{Theorem~5.3}, and the unbiased \(\hat\sigma^2=\mathrm{SSE}/(n-r)\) is used instead. We write \(s^2=\hat\sigma^2\)
and \(s=\sqrt{s^2}\); by Theorem~\ref{r12-thm-sampling}(b), \(s>0\) with probability \(1\).

\begin{theorem}[\citeat{searle-khuri}{\S12.11, (12.51)--(12.52)}, \citeat{tu-clm}{Theorems~7.1, 7.9, 9.1, 9.2, Proposition~7.13}]\label{r12-thm-t}
In the normal linear model let \(\rank X=r<n\), let \(G\) be a g-inverse of \(X'X\) and \(\beta^0=GX'y\), and let
\(a\in\R^{k+1}\) be nonzero with \(a'\beta\) estimable. Then \(a'Ga>0\), and we put
\(\operatorname{SE}(a'\beta^0)=s\sqrt{a'Ga}\), the \textbf{standard error} of \(a'\beta^0\), and \(t^*=t_{n-r,1-\alpha/2}\).
\begin{enumerate}
\item \(T=(a'\beta^0-a'\beta)/\operatorname{SE}(a'\beta^0)\sim t_{n-r}\).
\item \(\bigl[a'\beta^0-t^*\operatorname{SE}(a'\beta^0),\ a'\beta^0+t^*\operatorname{SE}(a'\beta^0)\bigr]\) is a confidence
interval of level \(1-\alpha\) for \(a'\beta\). For \(m_0\in\R\), the test of \(H_0\colon a'\beta=m_0\) that rejects when
\(\lvert T_0\rvert>t^*\), where \(T_0=(a'\beta^0-m_0)/\operatorname{SE}(a'\beta^0)\), has level \(\alpha\); it rejects
exactly when \(m_0\) lies outside this interval.
\item Let \(x_0\in\R^{k+1}\) with \(x_0'\beta\) estimable, and let \(y_0=x_0'\beta+\varepsilon_0\) be a future observation with
\(\varepsilon_0\sim N(0,\sigma^2)\) independent of \(\varepsilon\). Put \(w_0=s\sqrt{1+x_0'Gx_0}\). Then
\((y_0-x_0'\beta^0)/w_0\sim t_{n-r}\), and \([x_0'\beta^0-t^*w_0,\ x_0'\beta^0+t^*w_0]\) is a prediction interval of
level \(1-\alpha\) for \(y_0\).
\end{enumerate}
\(\operatorname{SE}(a'\beta^0)\) and the intervals are the same for every g-inverse \(G\) of \(X'X\)
(Theorem~\ref{r12-thm-sampling}(a) with \(K=a'\)). Under (A4), \(G=(X'X)^{-1}\), \(\beta^0=\hat\beta\) and \(r=k+1\).
\end{theorem}
\begin{proof}
(a)
\begin{itemize}
\item \(a=X'c\) for some \(c\in\R^n\) (Theorem~\ref{r11-thm-estimable}(a)), so, with \(H=XGX'\) symmetric and idempotent
(Theorem~\ref{r10-thm-PX}(a)), \(a'Ga=c'Hc=c'H'Hc=\lVert Hc\rVert^2\). As \(X'H=X'\), \(X'Hc=a\ne0\), so \(Hc\ne0\) and
\(a'Ga>0\).
\item By Theorem~\ref{r12-thm-sampling}(a) with \(K=a'\) and Theorem~\ref{r14-thm-mvn-linear},
\(z=(a'\beta^0-a'\beta)/(\sigma\sqrt{a'Ga})\sim N(0,1)\). Next \(v=\mathrm{SSE}/\sigma^2\sim\chi^2_{n-r}\)
(Theorem~\ref{r12-thm-sampling}(b)), and \(z\) and \(v\) are independent, as functions of \(a'\beta^0\) and of
\(\mathrm{SSE}\) (Theorem~\ref{r12-thm-sampling}(c), (N3)). Since \(s/\sigma=\sqrt{v/(n-r)}\),
\(T=z/\sqrt{v/(n-r)}\sim t_{n-r}\) (Definition~\ref{r15-def-t-F}(a)).
\end{itemize}
(b) By (a) and Theorem~\ref{r15-thm-t-symmetric}(a), \(\prb(\lvert T\rvert\leqslant t^*)=1-\alpha\), and
\(\lvert T\rvert\leqslant t^*\) holds exactly when \(a'\beta\) lies in the interval. If \(H_0\) holds, then \(T_0=T\),
so \(\prb(\lvert T_0\rvert>t^*)=1-(1-\alpha)=\alpha\). Finally \(\lvert T_0\rvert>t^*\) means
\(\lvert a'\beta^0-m_0\rvert>t^*\operatorname{SE}(a'\beta^0)\), that is, \(m_0\) lies outside the interval.

(c)
\begin{itemize}
\item As \(\varepsilon\sim N_n(0,\sigma^2I_n)\) and \(\varepsilon_0\sim N(0,\sigma^2)\) are independent,
\(\tilde\varepsilon=(\varepsilon',\varepsilon_0)'\sim N_{n+1}(0,\sigma^2I_{n+1})\) (Theorem~\ref{r14-thm-mvn-iid}(a)).
\item \(x_0=X'c_0\) for some \(c_0\in\R^n\) (Theorem~\ref{r11-thm-estimable}(a)), so
\(x_0'\beta^0=c_0'XGX'y=c_0'H(X\beta+\varepsilon)=x_0'\beta+c_0'H\varepsilon\), as \(HX=X\). Hence
\(y_0-x_0'\beta^0=\varepsilon_0-c_0'H\varepsilon=b'\tilde\varepsilon\), where \(b=(-c_0'H,\ 1)'\) has
\(b'b=c_0'HH'c_0+1=1+x_0'Gx_0\) (\(c_0'Hc_0=c_0'XGX'c_0=x_0'Gx_0\)). So
\(y_0-x_0'\beta^0\sim N\bigl(0,\sigma^2(1+x_0'Gx_0)\bigr)\) (Theorem~\ref{r14-thm-mvn-linear}).
\item As \(MX=0\), \(\mathrm{SSE}=y'My=\varepsilon'M\varepsilon=\tilde\varepsilon'P\tilde\varepsilon\) with \(P=\diag(M,0)\),
which is symmetric and idempotent, and \(b'P=(-c_0'HM,\ 0)=0\). By Theorem~\ref{r15-thm-indep}(b),
\(y_0-x_0'\beta^0\) and \(\mathrm{SSE}\) are independent. As in (a), \((y_0-x_0'\beta^0)/w_0\sim t_{n-r}\), and as
in (b), \(\prb(\lvert y_0-x_0'\beta^0\rvert\leqslant t^*w_0)=1-\alpha\).\qedhere
\end{itemize}
\end{proof}

For \(x_0\ne0\), the confidence interval of (b) for \(a=x_0\) has half-width \(t^*s\sqrt{x_0'Gx_0}\); the prediction
interval of (c) is wider, because the new error \(\varepsilon_0\) adds \(\sigma^2\) to the variance.

\begin{definition}[\citeat{searle-khuri}{Definition~12.2, Lemma~12.1, (12.46)--(12.47)}]\label{r11-def-testable}
Let \(\rank X=r\), let \(K\in\R^{q\times(k+1)}\) and \(m\in\R^q\). The hypothesis \(H_0\colon K\beta=m\) is \textbf{testable}
if every entry of \(K\beta\) is estimable and \(\rank K=q\) (the rows of \(K\) are linearly independent).
\end{definition}
Let \(G\) be a g-inverse of \(X'X\).
\begin{itemize}
\item \(H_0\colon K\beta=m\) is testable if and only if \(\rank K=q\) and \(K=WX'X\) for some \(W\in\R^{q\times(k+1)}\). Indeed,
if the entries of \(K\beta\) are estimable, then each row \(k_i'\) of \(K\) lies in \(\mathcal C(X')\)
(Theorem~\ref{r11-thm-estimable}(a)): \(k_i'=t_i'X\) with \(t_i\in\R^n\). So \(K=TX\), where \(T\in\R^{q\times n}\) has the
rows \(t_i'\) (Theorem~\ref{r1-thm-views}(c)), and \(X=XGX'X\) (Theorem~\ref{r10-thm-PX}(a)) gives \(K=TXGX'X=WX'X\) with
\(W=TXG\).
Conversely, if \(K=WX'X\), then \(K=TX\) with \(T=WX'\), so every row of \(K\) lies in \(\mathcal C(X')\), and the entries
of \(K\beta\) are estimable (Theorem~\ref{r11-thm-estimable}(a)).
\item Let \(H_0\) be testable, \(K=WX'X\) and \(T=WX'\), so that \(K=TX\). Then \(q=\rank K\leqslant\rank T\leqslant q\)
(Theorems~\ref{thm:r4-rank-product}(a) and~\ref{thm:r4-rank-rules}(a); \(T\) has \(q\) rows), so \(\rank T=q\) and \(TT'\) is PD (Theorem~\ref{r8-thm-XtX}).
Also \(q=\rank K\leqslant\rank X=r\) (Theorem~\ref{thm:r4-rank-product}(a)). By \(X'XGX'X=X'X\),
\[
V=KGK'=WX'XGX'XW'=WX'XW'=TT'
\]
so \(V\) is PD, hence invertible (Proposition~\ref{r8-prop-pd-basic}(c)), and it does not depend on \(G\).
\end{itemize}
Under (A4), \(H_0\colon K\beta=m\) is testable for every \(K\) of rank \(q\) (Theorem~\ref{r11-thm-estimable}(a)).

\begin{theorem}[\citeat{bapat}{Result~7.6}, \citeat{tu-clm}{Theorem~10.3, Corollary~7.6}]\label{r12-thm-cls}
Let \(X\in\R^{n\times(k+1)}\), \(y\in\R^n\), let \(G\) be a g-inverse of \(X'X\) and \(\beta^0=GX'y\), and let
\(H_0\colon K\beta=m\) be testable, with \(K\in\R^{q\times(k+1)}\), so that \(V=KGK'\) is invertible
(the paragraph after Definition~\ref{r11-def-testable}). Put
\[
\tilde\beta=\beta^0-GK'V^{-1}(K\beta^0-m)
\]
\begin{enumerate}
\item \(K\tilde\beta=m\), and for every \(\beta\) with \(K\beta=m\),
\(\lVert y-X\beta\rVert^2=\lVert y-X\tilde\beta\rVert^2+\lVert X(\tilde\beta-\beta)\rVert^2\).
\item The minimum of \(S(\beta)=\lVert y-X\beta\rVert^2\) subject to \(K\beta=m\) is attained at \(\tilde\beta\), the
\textbf{constrained least-squares estimate}; the minimizers are the \(\beta\) with \(K\beta=m\) and \(X\beta=X\tilde\beta\),
and under (A4) \(\tilde\beta\) is the only one.
\item \(\mathrm{SSE}_K-\mathrm{SSE}=(K\beta^0-m)'V^{-1}(K\beta^0-m)\), where \(\mathrm{SSE}_K=\lVert y-X\tilde\beta\rVert^2\) and
\(\mathrm{SSE}=\lVert y-X\beta^0\rVert^2\).
\end{enumerate}
\end{theorem}
\begin{proof}
Let \(K=WX'X\) and \(T=WX'\), so that \(K=TX\), \(K'=X'XW'\) and \(V=TT'\) (the paragraph after Definition~\ref{r11-def-testable}).

(a)
\begin{itemize}
\item \(K\tilde\beta=K\beta^0-KGK'V^{-1}(K\beta^0-m)=K\beta^0-(K\beta^0-m)=m\).
\item \(X'XGK'=X'XGX'XW'=X'XW'=K'\) and \(X'X\beta^0=X'y\) (Theorem~\ref{r10-thm-normal-equations}(c)), so
\(X'X\tilde\beta=X'y-K'V^{-1}(K\beta^0-m)\), that is, \(X'(y-X\tilde\beta)=K'V^{-1}(K\beta^0-m)\).
\item Let \(K\beta=m\). Then \((\tilde\beta-\beta)'X'(y-X\tilde\beta)=(K\tilde\beta-K\beta)'V^{-1}(K\beta^0-m)=0\), as
\(K\tilde\beta=m=K\beta\). So, with \(y-X\beta=(y-X\tilde\beta)+X(\tilde\beta-\beta)\),
\[
\lVert y-X\beta\rVert^2=\lVert y-X\tilde\beta\rVert^2+2(\tilde\beta-\beta)'X'(y-X\tilde\beta)+\lVert X(\tilde\beta-\beta)\rVert^2
=\lVert y-X\tilde\beta\rVert^2+\lVert X(\tilde\beta-\beta)\rVert^2
\]
\end{itemize}
(b) By (a), \(S(\beta)\geqslant S(\tilde\beta)\) for every \(\beta\) with \(K\beta=m\), with equality exactly when
\(X(\tilde\beta-\beta)=0\). Under (A4), \(X(\tilde\beta-\beta)=0\) only for \(\beta=\tilde\beta\)
(Theorem~\ref{thm:r4-rank-rules}(b)).

(c) By Theorem~\ref{r10-thm-least-squares}(a), \(S(\tilde\beta)=S(\beta^0)+\lVert X(\beta^0-\tilde\beta)\rVert^2\). Here
\(X(\beta^0-\tilde\beta)=XGK'V^{-1}d\) with \(d=K\beta^0-m\), and \(XGK'=XGX'XW'=XW'=T'\) (Theorem~\ref{r10-thm-PX}(a)); so, as
\(V=TT'\) and \(V^{-1}\) are symmetric, \(\lVert X(\beta^0-\tilde\beta)\rVert^2=d'V^{-1}TT'V^{-1}d=d'V^{-1}d\).
\end{proof}

\begin{theorem}[\citeat{searle-khuri}{(12.48)}, \citeat{bapat}{\S8.4}, \citeat{tu-clm}{Theorems~7.3, 7.4, Corollaries~7.6, 7.8, 7.15, 7.16}]\label{r12-thm-f}
In the normal linear model let \(\rank X=r<n\), let \(G\) be a g-inverse of \(X'X\) and \(\beta^0=GX'y\), let
\(H_0\colon K\beta=m\) be testable with \(K\in\R^{q\times(k+1)}\), \(V=KGK'\), and let \(\tilde\beta\) and
\(\mathrm{SSE}_K\) be as in Theorem~\ref{r12-thm-cls}.
\begin{enumerate}
\item For every \(y\) with \(\mathrm{SSE}>0\),
\[
F=\frac{(K\beta^0-m)'V^{-1}(K\beta^0-m)/q}{s^2}=\frac{(\mathrm{SSE}_K-\mathrm{SSE})/q}{\mathrm{SSE}/(n-r)}
\]
\item If \(H_0\) holds, then \(F\), which is defined with probability \(1\) (Theorem~\ref{r12-thm-sampling}(b)), has the
distribution \(F_{q,n-r}\); so the test that rejects \(H_0\) when \(F>F_{q,n-r,1-\alpha}\) has level \(\alpha\).
\item Let \(X=[X_1\ X_2]\) with \(X_1\in\R^{n\times q}\), \(\beta=(\beta_1',\beta_2')'\) accordingly, and let
\(H_0\colon\beta_1=0\) (\(K=[I_q\ 0]\), \(m=0\)) be testable. Then \(\mathrm{SSE}_K\) is the error sum of squares of the
least-squares fit with design \(X_2\). The order of the columns does not matter: reordering the columns of \(X\) and the
entries of \(\beta\) alike changes neither \(\mathcal C(X)\), \(\mathcal C(X_2)\) nor the sums of squares, so \(X_1\) may be
any \(q\) of the columns of \(X\), and \(X_2\) the others. If \(\one_n\in\mathcal C(X_2)\) and \(\mathrm{SSE}>0\), then, with \(R^2\) and
\(R^2_2\) the coefficients of determination of the fits with designs \(X\) and \(X_2\),
\[
F=\frac{(R^2-R^2_2)/q}{(1-R^2)/(n-r)}
\]
In particular, under (A4), for the \textbf{overall \(F\) test} of \(H_0\colon\beta_1=\dots=\beta_k=0\), where
\(X=[\one_n\ X_0]\) and \(k\geqslant1\), \(\mathrm{SSE}_K=\mathrm{SST}\), \(\mathrm{SSE}_K-\mathrm{SSE}=\mathrm{SSR}\) and
\[
F=\frac{\mathrm{SSR}/k}{\mathrm{SSE}/(n-k-1)}=\frac{R^2/k}{(1-R^2)/(n-k-1)}
\]
\item If \(q=1\) and \(K=a'\), then \(F=T_0^2\), where \(T_0\) is the statistic of Theorem~\ref{r12-thm-t}(b) for
\(m_0=m\); the test of (b) rejects exactly when that of Theorem~\ref{r12-thm-t}(b) does.
\end{enumerate}
Part (a) and the identities in (c) and (d) are algebraic: they hold for every \(y\) with \(\mathrm{SSE}>0\). Under (A4),
\(G=(X'X)^{-1}\), \(\beta^0=\hat\beta\) and \(r=k+1\).
\end{theorem}
\begin{proof}
(a) By Theorem~\ref{r12-thm-cls}(c), the numerator quadratic form is \(\mathrm{SSE}_K-\mathrm{SSE}\), and
\(s^2=\mathrm{SSE}/(n-r)\).

(b)
\begin{itemize}
\item Let \(K=WX'X\), \(T=WX'\), so that \(K=TX\) and \(V=TT'\) (the paragraph after Definition~\ref{r11-def-testable}); put \(H=XGX'\),
\(M=I_n-H\) and \(\varepsilon=y-X\beta\). Then \(K\beta^0=TXGX'y=THy=Ty\), as \(TH=WX'XGX'=WX'=T\) (Theorem~\ref{r10-thm-PX}(a)),
and \(XGK'=XGX'XW'=XW'=T'\).
\item So \(X\tilde\beta=Hy-T'V^{-1}(Ty-m)\). Under \(H_0\), \(m=K\beta=TX\beta\), so \(Ty-m=T\varepsilon\), and
\(y-X\tilde\beta=My+U\varepsilon\) with \(U=T'V^{-1}T\). Also \(My=M\varepsilon\), as \(MX=0\).
\item \(U\) is symmetric, \(U^2=T'V^{-1}(TT')V^{-1}T=U\), and \(\rank U=\tr U=\tr(V^{-1}TT')=\tr I_q=q\)
(Theorems~\ref{r6-thm-idempotent}(c) and~\ref{r1-thm-trace}(b)). \(MU=MXW'V^{-1}T=0\) and \(UM=(MU)'=0\).
\item Hence \(\mathrm{SSE}_K=\lVert M\varepsilon+U\varepsilon\rVert^2=\varepsilon'M\varepsilon+\varepsilon'U\varepsilon\), so
\(\mathrm{SSE}_K-\mathrm{SSE}=\varepsilon'U\varepsilon\), with \(\mathrm{SSE}=\varepsilon'M\varepsilon\).
\item By (A5) and Theorem~\ref{r15-thm-quadratic-chisq}, \(u=\varepsilon'U\varepsilon/\sigma^2\sim\chi^2_q\) and
\(v=\varepsilon'M\varepsilon/\sigma^2\sim\chi^2_{n-r}\); by Theorem~\ref{r15-thm-indep}(a), \(u\) and \(v\) are independent.
So \(F=(u/q)/(v/(n-r))\sim F_{q,n-r}\) (Definition~\ref{r15-def-t-F}(b)), and
\(\prb(F>F_{q,n-r,1-\alpha})=1-(1-\alpha)=\alpha\) (Definition~\ref{r15-def-t-F}(c)).
\end{itemize}
(c)
\begin{itemize}
\item The \(\beta\) with \(\beta_1=0\) give exactly the vectors \(X\beta=X_2\beta_2\), \(\beta_2\) arbitrary. By
Theorem~\ref{r12-thm-cls}(b), \(\mathrm{SSE}_K\) is the minimum of \(\lVert y-X\beta\rVert^2\) over these \(\beta\), that is,
the minimum of \(\lVert y-X_2\beta_2\rVert^2\): the error sum of squares of the design \(X_2\).
\item Reordering the columns of \(X\) and the entries of \(\beta\) alike leaves \(X\beta=\sum_j\beta_j\col_j(X)\)
(Theorem~\ref{r1-thm-views}(b)) and the hypothesis \(\beta_1=0\) as they are, and also \(\mathcal C(X)\) and
\(\mathcal C(X_2)\), which are spanned by the same columns as before (Definition~\ref{def:r4-spaces}). So
\(r=\dim\mathcal C(X)\) (Definition~\ref{def:r4-rank}) does not change, and neither do the fitted values \(Hy\) and
\(H_2y\), where \(H\) and \(H_2\) are the hat matrices of \(X\) and \(X_2\): they are the orthogonal projections of \(y\)
onto \(\mathcal C(X)\) and \(\mathcal C(X_2)\) (Theorem~\ref{r10-thm-PX}). Hence \(\mathrm{SSE}\), \(\mathrm{SSE}_K\) (by the first point), \(R^2\) and
\(R^2_2\) (Definition~\ref{r11-def-sums-of-squares}) do not change.
\item If \(\one_n\in\mathcal C(X_2)\), then also \(\one_n\in\mathcal C(X)\), \(\mathrm{SST}\geqslant\mathrm{SSE}>0\), and
Theorem~\ref{r11-thm-anova} gives \(\mathrm{SSE}=\mathrm{SST}(1-R^2)\) and \(\mathrm{SSE}_K=\mathrm{SST}(1-R^2_2)\);
substitute into (a) and cancel \(\mathrm{SST}\).
\item For the overall test, \(X_2=\one_n\), whose fitted values are \(\bar y\one_n\) (Example~\ref{r6-exm-centering}).
So \(\mathrm{SSE}_K=\mathrm{SST}\), \(R^2_2=0\) and \(\mathrm{SSE}_K-\mathrm{SSE}=\mathrm{SSR}\) (Theorem~\ref{r11-thm-anova});
under (A4), \(r=k+1\) and the hypothesis is testable (Definition~\ref{r11-def-testable}).
\end{itemize}
(d) For \(K=a'\), \(V=a'Ga\), and the numerator of \(F\) is \((a'\beta^0-m)^2/(a'Ga)\), so
\(F=(a'\beta^0-m)^2/(s^2a'Ga)=T_0^2\). Hence \(F>F_{1,n-r,1-\alpha}\iff T_0^2>(t^*)^2\iff\lvert T_0\rvert>t^*\), as
\((t^*)^2=F_{1,n-r,1-\alpha}\) and \(t^*>0\) (Theorem~\ref{r15-thm-t-symmetric}(b) and its proof).
\end{proof}

\begin{definition}[\citeat{searle-khuri}{\S11.9, \S12.9, \S12.13}]\label{r11-def-reduction}
Let \(X=[X_1\ X_2]\), with the columns in any order as in Theorem~\ref{r12-thm-f}(c), \(\beta=(\beta_1',\beta_2')'\)
accordingly, and let \(H\) and \(H_2\) be the hat matrices of \(X\) and
\(X_2\) (Theorem~\ref{r10-thm-PX}). The \textbf{reduction in the sum of squares} due to fitting \(X\beta\) is
\(R(\beta_1,\beta_2)=y'Hy=\lVert Hy\rVert^2\), and \(R(\beta_2)=y'H_2y\) is that due to fitting \(X_2\beta_2\) alone; the
reduction due to \(\beta_1\) after \(\beta_2\) is
\[
R(\beta_1\mid\beta_2)=R(\beta_1,\beta_2)-R(\beta_2)
\]
\end{definition}
By Theorem~\ref{r10-thm-PX}(d), \(R(\beta_1,\beta_2)=\hat\beta'X'y=y'y-\mathrm{SSE}\) for every least-squares estimate
\(\hat\beta\), and \(R(\beta_2)=y'y-\mathrm{SSE}_2\), where \(\mathrm{SSE}_2\) is the error sum of squares of the fit with design
\(X_2\). So \(R(\beta_1\mid\beta_2)=\mathrm{SSE}_2-\mathrm{SSE}\): by Theorem~\ref{r12-thm-f}(c), if \(H_0\colon\beta_1=0\) is
testable, then \(F=\bigl(R(\beta_1\mid\beta_2)/q\bigr)/s^2\). For a model with an intercept \(\mu\) (\(X_2=\one_n\)),
\(R(\mu)=n\bar y^2\) (Example~\ref{r6-exm-centering}), and \(R(\beta_1\mid\mu)=\mathrm{SST}-\mathrm{SSE}=\mathrm{SSR}\)
(Theorem~\ref{r11-thm-anova}).

\begin{example}\label{r12-exm-two-factors}
Eight plots receive a fertilizer at two levels, coded \(x_{i1}=-1,1\), and irrigation at two levels,
\(x_{i2}=-1,1\), two plots for each combination; the yields are
\[
\begin{array}{c|rrrrrrrr}
x_{i1}&-1&-1&-1&-1&1&1&1&1\\
x_{i2}&-1&-1&1&1&-1&-1&1&1\\\hline
y_i&9&3&9&7&11&13&11&17
\end{array}
\]
In the normal linear model \(y_i=\beta_0+\beta_1x_{i1}+\beta_2x_{i2}+\varepsilon_i\) (\(n=8\), \(k=2\)), find the confidence
interval of level \(1-\alpha\) for \(\beta_1\), the tests of \(\beta_2=0\) and of \(\beta_1=\beta_2=0\) at level \(\alpha\),
and the prediction interval of level \(1-\alpha\) for a plot with \(x_1=x_2=1\). Put \(t^*=t_{5,1-\alpha/2}\).

\textit{Solution.}
\begin{itemize}
\item The three columns of \(X\) are orthogonal with squared length \(8\), so \(X'X=8I_3\), \(\rank X=3\) (A4),
\(G=(X'X)^{-1}=\frac18I_3\), \(X'y=(80,24,8)'\) and \(\hat\beta=(10,3,1)'\). The fitted values are \(6,8,12,14\) for the four combinations, the residuals
\(e=(3,-3,1,-1,-1,1,-3,3)'\), \(\mathrm{SSE}=40\) and \(s^2=40/5=8\). With \(\bar y=10\),
\(\mathrm{SSR}=2(16+4+4+16)=80\), \(\mathrm{SST}=120\) and \(R^2=\frac23\).
\item \(\operatorname{SE}(\hat\beta_j)=s\sqrt{1/8}=1\) for every \(j\) (Theorem~\ref{r12-thm-t}). By Theorem~\ref{r12-thm-t}(b) the interval for
\(\beta_1\) is \([3-t^*,\ 3+t^*]\), and the test of \(\beta_2=0\) has \(T_0=1\): it rejects if and only if
\(t^*<1\); equivalently \(F=T_0^2=1\) (Theorem~\ref{r12-thm-f}(d)).
\item The overall \(F\) statistic is \(F=(80/2)/(40/5)=5=\frac{(2/3)/2}{(1/3)/5}\) (Theorem~\ref{r12-thm-f}(c)); the test
rejects if and only if \(F_{2,5,1-\alpha}<5\).
\item For \(x_0=(1,1,1)'\), \(x_0'\hat\beta=14\) and \(x_0'(X'X)^{-1}x_0=\frac38\), so \(w_0=s\sqrt{1+\frac38}=\sqrt{11}\):
the prediction interval is \([14-t^*\sqrt{11},\ 14+t^*\sqrt{11}]\) (Theorem~\ref{r12-thm-t}(c)), while the confidence
interval for \(x_0'\beta\) is \([14-t^*\sqrt3,\ 14+t^*\sqrt3]\).
\end{itemize}
\end{example}

\needspace{14\baselineskip}
\subsection*{The one-way layout and the analysis of variance}

\begin{theorem}[\citeat{tu-clm}{Proposition~11.2}, \citeat{searle-khuri}{\S12.12.1}, \citeat{bapat}{\S8.3}]\label{r11-thm-oneway-anova}
In the one-way layout of Theorem~\ref{r11-thm-oneway}, let \(\mathcal I_l\) be the set of the observations (row indices) of
group \(l\), let \(\bar y_l=n_l^{-1}\sum_{i\in\mathcal I_l}y_i\) be the \textbf{group mean} and \(\bar y\) the
\textbf{grand mean} of all \(n\) observations, and let \(H\) and \(M\) be the hat matrix and the residual maker of \(X\)
(Theorem~\ref{r10-thm-PX}).
\begin{enumerate}
\item \(H=\sum_ln_l^{-1}d_ld_l'=\diag(\bar J_{n_1},\dots,\bar J_{n_g})\), and \(\hat y=Hy=\bar y_1d_1+\dots+\bar y_gd_g\): the
fitted value of each observation is the mean of its group.
\item With \(\bar y\), \(\bar y_l-\bar y\) and \(y_i-\bar y_l\) in a row \(i\in\mathcal I_l\), the three vectors
\(\bar y\one_n\), \(Hy-\bar y\one_n\) and \(My\) are orthogonal, lie in subspaces of dimensions \(1\), \(g-1\) and \(n-g\), and add
up to \(y\); and
\[
\sum_{i=1}^n(y_i-\bar y)^2=\sum_{l=1}^gn_l(\bar y_l-\bar y)^2+\sum_{l=1}^g\sum_{i\in\mathcal I_l}(y_i-\bar y_l)^2
\]
that is, \(\mathrm{SST}=\mathrm{SSR}+\mathrm{SSE}\), with the \textbf{between-group sum of squares}
\(\mathrm{SSR}=\sum_ln_l(\bar y_l-\bar y)^2\) and the \textbf{within-group sum of squares}
\(\mathrm{SSE}=\sum_l\sum_{i\in\mathcal I_l}(y_i-\bar y_l)^2\).
\end{enumerate}
\end{theorem}
\begin{proof}
(a)
\begin{itemize}
\item No observation lies in two groups, so \(d_l'd_m=0\) for \(l\ne m\), and \(d_l'd_l=n_l\); these are the entries of
\(X_1'X_1\) (Theorems~\ref{r1-thm-views}(d) and~\ref{r1-thm-transpose}), so \(X_1'X_1=\diag(n_1,\dots,n_g)\).
\item \(\mathcal C(X_1)=\mathcal C(X)\) and \(X_1\) has full column rank (Theorem~\ref{r11-thm-oneway}(c)). For every \(y\),
\(Hy\) and the hat matrix \(H_1=X_1(X_1'X_1)^{-1}X_1'\) of \(X_1\) give the orthogonal projection of \(y\) onto the same
subspace (Theorem~\ref{r10-thm-PX}, and~(c) there for \(X_1\)), so \(H=H_1\) (Theorem~\ref{r1-thm-views}(b), with the
columns of \(I_n\) for \(y\)).
\item By the outer product expansion (Theorem~\ref{r1-thm-views}(d)),
\(H_1=X_1\diag(n_1^{-1},\dots,n_g^{-1})X_1'=\sum_ln_l^{-1}d_ld_l'\). The \((i,j)\) entry of \(d_ld_l'\) is \(1\) if
\(i,j\in\mathcal I_l\) and \(0\) otherwise, which gives the block diagonal form; and
\(n_l^{-1}d_ld_l'y=n_l^{-1}(d_l'y)d_l=\bar y_ld_l\).
\end{itemize}
(b) Let \(U=\Span\{\one_n\}\subseteq W=\mathcal C(X)\), so that \(P_W=H\). Then \(P_Uy=\bar y\one_n\)
(Corollary~\ref{r6-cor-centering}(a)), and \(P_Wy=Hy\) is given by (a). By Theorem~\ref{r6-thm-nested}(b),(c),
\(y=P_Uy+(P_W-P_U)y+(I_n-P_W)y\), with entries \(\bar y\), \(\bar y_l-\bar y\) and \(y_i-\bar y_l\) in row \(i\in\mathcal I_l\), is
a sum of three orthogonal vectors in \(U\), \(W\cap U^\perp\) and \(W^\perp\), of dimensions \(1\), \(g-1\) and \(n-g\)
(Theorems~\ref{r11-thm-oneway}(b) and~\ref{r6-thm-complement}(c)), and the squared lengths add up. Subtract
\(\lVert P_Uy\rVert^2=n\bar y^2\) and use \(\lVert y\rVert^2-n\bar y^2=\sum_i(y_i-\bar y)^2\) (Example~\ref{r1-exm-centering}).
The first sum on the right is \(\lVert Hy-\bar y\one_n\rVert^2=\mathrm{SSR}\) and the second is \(\lVert My\rVert^2=\mathrm{SSE}\)
(Definition~\ref{r11-def-sums-of-squares}).
\end{proof}

\begin{example}[\citeat{tu-clm}{Proposition~11.2}]\label{r3-exm-between-within}
In Example~\ref{r11-exm-oneway} (\(g=3\); \(n_1,n_2,n_3=3,2,1\)), let \(X_1=[d_1\ d_2\ d_3]\), so that
\(\mathcal C(X_1)=\mathcal C(X)\) is the set of vectors that are constant within each group, and let
\(W=\{x\in\R^6:d_1'x=d_2'x=d_3'x=0\}\), the vectors whose entries add up to \(0\) within each group. Then
\[
\R^6=\mathcal C(X_1)\oplus W
\]
\begin{itemize}
\item \(W\) is the set of solutions of \(X_1'x=0\), a subspace (Example~\ref{r3-exm-subspaces}(e)). The
coefficient matrix
\[
X_1'=\begin{pmatrix}1&1&1&0&0&0\\0&0&0&1&1&0\\0&0&0&0&0&1\end{pmatrix}
\]
is in RREF, with free unknowns \(x_2,x_3,x_5\). By Proposition~\ref{r2-prop-solutions}(b), \(x\in W\) if and
only if \(x_1=-x_2-x_3\), \(x_4=-x_5\) and \(x_6=0\), that is,
\[
x=x_2f_1+x_3f_2+x_5f_3,\qquad
f_1=\begin{pmatrix}-1\\1\\0\\0\\0\\0\end{pmatrix},\quad
f_2=\begin{pmatrix}-1\\0\\1\\0\\0\\0\end{pmatrix},\quad
f_3=\begin{pmatrix}0\\0\\0\\-1\\1\\0\end{pmatrix}
\]
Entries \(2,3,5\) of \(x_2f_1+x_3f_2+x_5f_3\) are \(x_2,x_3,x_5\), so \(f_1,f_2,f_3\) are
independent: they form a basis of \(W\).
\item \(\mathcal C(X_1)\cap W=\{0\}\): if \(x=c_1d_1+c_2d_2+c_3d_3\) lies in \(W\), then
\(0=d_m'x=c_md_m'd_m=c_mn_m\) for each \(m\), as \(d_m'd_l=0\) for \(l\ne m\) and \(d_m'd_m=n_m\); so every
\(c_m=0\) and \(x=0\).
\item \(\mathcal C(X_1)+W=\R^6\): for \(x\in\R^6\) let \(\bar x_l=d_l'x/n_l\), the mean of the entries of \(x\) in
group \(l\), and \(u=\bar x_1d_1+\bar x_2d_2+\bar x_3d_3\in\mathcal C(X_1)\). Then
\(d_m'(x-u)=d_m'x-\bar x_mn_m=0\) for each \(m\), so \(x-u\in W\), and \(x=u+(x-u)\).
\item Hence \(\R^6=\mathcal C(X_1)\oplus W\) (Definition~\ref{r3-def-sum}), and \(d_1,d_2,d_3,f_1,f_2,f_3\) form a
basis of \(\R^6\) (Theorem~\ref{r3-thm-direct-sum-basis}). By Theorem~\ref{r3-thm-direct-sum}, each
\(x\) is in exactly one way the sum of a vector of \(\mathcal C(X_1)\) and a vector of \(W\): its
\textbf{between-group part} \(u\), which replaces each entry by the mean of its group, and its
\textbf{within-group part} \(x-u\), the deviations from the group means. By
Theorem~\ref{r11-thm-oneway-anova}(a), \(u=Hx\) and \(x-u=Mx\). Also \(W=\mathcal N(X_1')=\mathcal C(X_1)^\perp\) by
Theorem~\ref{r6-thm-complement}(a), as the columns \(d_1,d_2,d_3\) of \(X_1\) span \(\mathcal C(X_1)\); and
\(\mathcal C(X_1)=\mathcal C(X)\) (Theorem~\ref{r11-thm-oneway}(c)), so \(W=\mathcal C(X)^\perp\).
\end{itemize}
For the responses \(y=(5,7,6,3,5,2)'\), the group means are \(6,4,2\), and
\(y=(6,6,6,4,4,2)'+(-1,1,0,-1,1,0)'\), the fitted values and the residuals of Example~\ref{r11-exm-oneway}.
With a single group, \(\mathcal C(X_1)=\Span\{\one\}\) and \(W=W_0\): this is Example~\ref{r3-exm-constant-centered}.
\end{example}

\begin{corollary}[\citeat{haerdle-simar}{\S3.6, (3.55)--(3.56), Exercise~3.22}, \citeat{bapat}{\S8.3}]\label{r11-cor-oneway-f}
In the one-way layout of Theorem~\ref{r11-thm-oneway}, use the design \(X_1=[d_1\ \cdots\ d_g]\), with
\(\beta=(\mu_1,\dots,\mu_g)'\), the group means, and assume the normal linear model with \(n>g\geqslant2\). Then
\(H_0\colon\mu_1=\dots=\mu_g\) is \(K\beta=0\) for the \((g-1)\times g\) matrix \(K\) whose row \(l\) is
\(\row_l(I_g)-\row_{l+1}(I_g)\), the hypothesis is testable, and
\[
F=\frac{\sum_ln_l(\bar y_l-\bar y)^2/(g-1)}{\sum_l\sum_{i\in\mathcal I_l}(y_i-\bar y_l)^2/(n-g)}\sim F_{g-1,\,n-g}\quad\text{under }H_0
\]
with \(\bar y_l\), \(\bar y\) and \(\mathcal I_l\) as in Theorem~\ref{r11-thm-oneway-anova}.
\end{corollary}
\begin{proof}
\begin{itemize}
\item \(K\beta=0\) says \(\mu_l=\mu_{l+1}\) for \(l=1,\dots,g-1\), that is, \(H_0\). The first \(g-1\) columns of \(K\) form
an upper triangular matrix with diagonal entries \(1\), which is invertible (Theorems~\ref{r5d-thm-basic}(c)
and~\ref{r5d-thm-product}(a)); so these \(g-1\) columns are linearly independent, and \(\rank K=g-1\), as \(K\)
has \(g-1\) rows (Theorem~\ref{thm:r4-rank-rules}(a)). \(X_1\) has full column rank (Theorem~\ref{r11-thm-oneway}(a)), so
\(H_0\) is testable (Definition~\ref{r11-def-testable}).
\item No observation lies in two groups, so \(d_l'd_m=0\) for \(l\ne m\), and \(d_l'd_l=n_l\); these are the entries of
\(X_1'X_1\) (Theorems~\ref{r1-thm-views}(d) and~\ref{r1-thm-transpose}), so \(X_1'X_1=\diag(n_1,\dots,n_g)\). And
\(X_1'y\) holds the group totals \(d_l'y\) (Theorem~\ref{r1-thm-views}(d)); so \(\hat\beta=(\bar y_1,\dots,\bar y_g)'\), each
fitted value is the mean of its group, and \(\mathrm{SSE}=\sum_l\sum_{i\in\mathcal I_l}(y_i-\bar y_l)^2\).
\item The \(\beta\) with \(K\beta=0\) are \(\beta=\mu\one_g\), with \(X_1\beta=\mu\one_n\). By Theorem~\ref{r12-thm-cls}(b),
\(\mathrm{SSE}_K=\min_\mu\lVert y-\mu\one_n\rVert^2=\mathrm{SST}\) (Example~\ref{r6-exm-centering}). As
\(\one_n\in\mathcal C(X_1)\), Theorem~\ref{r11-thm-oneway-anova}(b) gives
\(\mathrm{SSE}_K-\mathrm{SSE}=\mathrm{SSR}=\sum_ln_l(\bar y_l-\bar y)^2\). Theorem~\ref{r12-thm-f}(a),(b) with \(q=g-1\) and
\(r=g\) gives the formula and the distribution.\qedhere
\end{itemize}
\end{proof}

The computations of Corollary~\ref{r11-cor-oneway-f} are arranged in the \textbf{analysis-of-variance table}~\citeat{searle-khuri}{Tables~12.5, 12.6}, \citeat{bapat}{\S8.3}; the mean
squares are the sums of squares divided by their degrees of freedom, and \(F\) is the ratio of the two mean squares:
\[
\begin{array}{lccc}
\text{source}&\text{degrees of freedom}&\text{sum of squares}&\text{mean square}\\\hline
\text{between groups}&g-1&\mathrm{SSR}=\sum_ln_l(\bar y_l-\bar y)^2&\mathrm{SSR}/(g-1)\\
\text{within groups}&n-g&\mathrm{SSE}=\sum_l\sum_{i\in\mathcal I_l}(y_i-\bar y_l)^2&s^2=\mathrm{SSE}/(n-g)\\\hline
\text{total}&n-1&\mathrm{SST}=\sum_i(y_i-\bar y)^2&
\end{array}
\]
In the notation of Definition~\ref{r11-def-reduction}, with the intercept \(\mu\) and the effects
\(\alpha=(\alpha_1,\dots,\alpha_g)'\) of Theorem~\ref{r11-thm-oneway}, \(R(\mu)=n\bar y^2\),
\(R(\mu,\alpha)=\lVert Hy\rVert^2=\sum_ln_l\bar y_l^2\) and \(R(\alpha\mid\mu)=\mathrm{SSR}\); so
\(y'y=R(\mu)+R(\alpha\mid\mu)+\mathrm{SSE}\), the uncentered split of the paragraph after Theorem~\ref{r11-thm-anova}.

\begin{example}\label{r11-exm-oneway-f}
In Example~\ref{r11-exm-oneway}, the six weights \(y=(5,7,6,3,5,2)'\) of three types of plants, assume the normal linear
model. Test \(H_0\colon\alpha_1=\alpha_2=\alpha_3\) at level \(\alpha\), and find the confidence interval of level \(1-\alpha\)
for \(\alpha_1-\alpha_2\).

\textit{Solution.}
\begin{itemize}
\item The group means are \(6,4,2\) and \(\bar y=28/6=14/3\); the residuals are \(e=(-1,1,0,-1,1,0)'\), so \(\mathrm{SSE}=4\),
\(r=\rank X=3\) and \(s^2=4/3\) (Example~\ref{r11-exm-oneway}). By Theorem~\ref{r11-thm-oneway-anova}(b),
\[
\mathrm{SSR}=3\Bigl(6-\frac{14}3\Bigr)^2+2\Bigl(4-\frac{14}3\Bigr)^2+\Bigl(2-\frac{14}3\Bigr)^2=\frac{48+8+64}9=\frac{40}3,\qquad
\mathrm{SST}=\frac{40}3+4=\frac{52}3
\]
\item In the design \(X_1\), \(H_0\) is \(\mu_1=\mu_2=\mu_3\), and Corollary~\ref{r11-cor-oneway-f} gives
\(F=\frac{40/3}{2}\big/\frac43=5\), with \(2\) and \(3\) degrees of freedom: the test of level \(\alpha\) rejects \(H_0\) if
and only if \(F_{2,3,1-\alpha}<5\). The analysis-of-variance table:
\[
\begin{array}{lccc}
\text{source}&\text{degrees of freedom}&\text{sum of squares}&\text{mean square}\\\hline
\text{between groups}&2&40/3&20/3\\
\text{within groups}&3&4&4/3\\\hline
\text{total}&5&52/3&
\end{array}
\]
\item The same \(F\) from Theorem~\ref{r12-thm-f}(a) in the design \(X\) of rank \(3\): \(H_0\) is \(K\beta=0\) with the rows
\((0,1,-1,0)\) and \((0,0,1,-1)\) of \(K\), and with \(G=\diag(0,\frac13,\frac12,1)\) and \(\beta^0=(0,6,4,2)'\)
(Example~\ref{r11-exm-solution}), \(K\beta^0=(2,2)'\) and
\[
V=KGK'=\begin{pmatrix}\frac56&-\frac12\\-\frac12&\frac32\end{pmatrix},\qquad
V^{-1}=\begin{pmatrix}\frac32&\frac12\\\frac12&\frac56\end{pmatrix}
\]
(\(\det V=\frac54-\frac14=1\)). So \((K\beta^0)'V^{-1}K\beta^0=4\bigl(\frac32+\frac12+\frac12+\frac56\bigr)=\frac{40}3\), the same
\(\mathrm{SSE}_K-\mathrm{SSE}\), and \(F=\frac{40}3\big/(2\cdot\frac43)=5\).
\item \(\alpha_1-\alpha_2\) is estimable with BLUE \(\bar y_1-\bar y_2=2\) and \(a'Ga=\frac56\) for \(a=(0,1,-1,0)'\)
(Examples~\ref{r11-exm-estimable} and~\ref{r11-exm-solution}); so \(\operatorname{SE}=s\sqrt{5/6}=\sqrt{10/9}=\sqrt{10}/3\), and by
Theorem~\ref{r12-thm-t}(b) the interval is \(\bigl[2-t_{3,1-\alpha/2}\sqrt{10}/3,\ 2+t_{3,1-\alpha/2}\sqrt{10}/3\bigr]\). It
contains \(0\) exactly when \(t_{3,1-\alpha/2}\geqslant3\cdot2/\sqrt{10}=6/\sqrt{10}\).
\item Uncentered: \(n\bar y^2=6\cdot\frac{196}9=\frac{392}3\), and \(y'y=148=\frac{392}3+\frac{40}3+4=R(\mu)+R(\alpha\mid\mu)+\mathrm{SSE}\).
\end{itemize}
\end{example}

\begin{example}[\citeat{searle-khuri}{Exercise~10.3}, \citeat{bapat}{\S8.3}]\label{r15-exm-oneway}
In a balanced one-way layout, \(g\geqslant2\) groups have \(m\geqslant2\) observations each, \(n=gm\); \(x_{ij}\) is observation \(j\)
of group \(i\), and \(x\in\R^n\) lists group \(1\) first, then group \(2\), and so on. With the indicators \(d_1,\dots,d_g\) of the
groups (Theorem~\ref{r11-thm-oneway}), the group means are \(\bar x_i=d_i'x/m\) and the overall mean is \(\bar x=\one'x/n\). Each
observation lies in exactly one group, so \(d_i'd_l=0\) for \(i\ne l\), \(d_i'd_i=m\) and \(d_1+\dots+d_g=\one_n\). Let
\(D=d_1d_1'+\dots+d_gd_g'=\diag(J_m,\dots,J_m)\) and
\[
A_1=\bar J_n,\qquad A_2=\tfrac1mD-\bar J_n,\qquad A_3=I_n-\tfrac1mD
\]
\begin{itemize}
\item \(\frac1mDx=\sum_i\bar x_id_i\) has \(\bar x_i\) in the rows of group \(i\): it is the between-group part of \(x\), and
\(A_3x=x-\frac1mDx\) is the within-group part (Example~\ref{r3-exm-between-within}); \(\frac1mD\) is the hat matrix \(H\) of
Theorem~\ref{r11-thm-oneway-anova}(a) with \(n_1=\dots=n_g=m\). As \(\sum_i\bar x_i=\one'x/m=g\bar x\),
\[
x'A_1x=n\bar x^2,\qquad x'A_2x=m\sum_{i=1}^g\bar x_i^2-n\bar x^2=m\sum_{i=1}^g(\bar x_i-\bar x)^2=\mathrm{SS}_A
\]
and
\[
x'A_3x=x'x-m\sum_{i=1}^g\bar x_i^2=\sum_{i=1}^g\sum_{j=1}^m(x_{ij}-\bar x_i)^2=\mathrm{SS}_E
\]
using Example~\ref{r1-exm-centering} within each group. \(\mathrm{SS}_A\) and \(\mathrm{SS}_E\) are the sums of squares between and
within the groups, and \(A_1+A_2+A_3=I_n\).
\item \(D^2=\sum_{i,l}d_i(d_i'd_l)d_l'=mD\) and \(D\one=m\one\), so \((\frac1mD)^2=\frac1mD\) and
\(\frac1mD\bar J_n=\bar J_n=\bar J_n\frac1mD\). With \(\bar J_n^2=\bar J_n\) (Proposition~\ref{r1-prop-centering}),
\(A_2^2=\frac1mD-2\bar J_n+\bar J_n=A_2\) and \(A_3^2=I_n-\frac2mD+\frac1mD=A_3\): all three matrices are symmetric and idempotent.
By Theorem~\ref{r15-thm-cochran-algebraic}, their ranks add up to \(n\) and \(A_iA_j=0\) for \(i\ne j\); the ranks are the traces
\(1\), \(\frac nm-1=g-1\) and \(n-g\) (Theorem~\ref{r6-thm-idempotent}(c)).
\item Let \(x_{ij}=\mu_i+\varepsilon_{ij}\) with independent \(\varepsilon_{ij}\sim N(0,\sigma^2)\) (\(\sigma>0\)):
\(x\sim N_n(\mu,\sigma^2I_n)\) with
\(\mu=\sum_i\mu_id_i\) (Theorem~\ref{r14-thm-mvn-iid}(a)). By Theorem~\ref{r15-thm-cochran}(a), \(n\bar x^2\), \(\mathrm{SS}_A\) and
\(\mathrm{SS}_E\) are independent; \(\frac1mD\mu=\mu\) gives \(A_3\mu=0\), so \(\mathrm{SS}_E/\sigma^2\sim\chi^2_{n-g}\); and, computing
\(\mu'A_2\mu\) as \(x'A_2x\) above, \(\mathrm{SS}_A/\sigma^2\sim\chi^2_{g-1}\bigl(m\sum_i(\mu_i-\bar\mu)^2/\sigma^2\bigr)\) with
\(\bar\mu=(\mu_1+\dots+\mu_g)/g\).
\item If \(\mu_1=\dots=\mu_g\), then \(\mathrm{SS}_A/\sigma^2\sim\chi^2_{g-1}\), and by Definition~\ref{r15-def-t-F}
\[
F=\frac{\mathrm{SS}_A/(g-1)}{\mathrm{SS}_E/(n-g)}\sim F_{g-1,\,n-g}
\]
\end{itemize}
\end{example}

\subsection*{Exercises}

\begin{exercise}[cf.\ \citeat{searle-khuri}{\S12.12.2}]\label{r11-ex-two-way-estimable}
Six plants of three varieties are grown under two fertilizer treatments, one plant for each combination, and follow the
model \(y_{ij}=\mu+\alpha_i+\gamma_j+\varepsilon_{ij}\) (variety \(i\), treatment \(j\)) under (A1)--(A3), with
\(\beta=(\mu,\alpha_1,\alpha_2,\alpha_3,\gamma_1,\gamma_2)'\). With the observations ordered \(y_{11},y_{12},y_{21},y_{22},y_{31},y_{32}\),
\[
X=\begin{pmatrix}1&1&0&0&1&0\\1&1&0&0&0&1\\1&0&1&0&1&0\\1&0&1&0&0&1\\1&0&0&1&1&0\\1&0&0&1&0&1\end{pmatrix},\qquad
y=\begin{pmatrix}7\\3\\3\\3\\2\\0\end{pmatrix}
\]
\begin{enumerate}
\item Row-reduce \(X\), and find \(\rank X\) and a basis of \(\mathcal N(X)\). Check the basis against the linear relations among
the columns of \(X\).
\item Which of \(\alpha_1-\alpha_2\), \(\gamma_1-\gamma_2\), \(\mu+\alpha_1+\gamma_1\), \(\alpha_1\) and \(\mu\) are estimable?
\item Check that \(\hat\beta=(0,5,3,1,1,-1)'\) is a least-squares estimate, and find the estimates of the estimable
functions in (b).
\item Find \(\hat y\), \(e\), SSE and the unbiased estimate of \(\sigma^2\) of Theorem~\ref{r11-thm-sigma}.
\end{enumerate}
\end{exercise}
\begin{solution}
(a) Row operations do not change the solutions of \(Xw=0\):
\[
\begin{gathered}
\begin{pmatrix}1&1&0&0&1&0\\1&1&0&0&0&1\\1&0&1&0&1&0\\1&0&1&0&0&1\\1&0&0&1&1&0\\1&0&0&1&0&1\end{pmatrix}
\xrightarrow[R_5\leftarrow R_5-R_1]{R_2\leftarrow R_2-R_1,\ R_3\leftarrow R_3-R_1}
\begin{pmatrix}1&1&0&0&1&0\\0&0&0&0&-1&1\\0&-1&1&0&0&0\\1&0&1&0&0&1\\0&-1&0&1&0&0\\1&0&0&1&0&1\end{pmatrix}
\\
\xrightarrow[R_6\leftarrow R_6-R_1,\ R_6\leftarrow R_6-R_5]{R_4\leftarrow R_4-R_1,\ R_4\leftarrow R_4-R_3}
\begin{pmatrix}1&1&0&0&1&0\\0&0&0&0&-1&1\\0&-1&1&0&0&0\\0&0&0&0&-1&1\\0&-1&0&1&0&0\\0&0&0&0&-1&1\end{pmatrix}
\xrightarrow[R_6\leftarrow R_6-R_4]{R_2\leftarrow R_2+R_3}
\begin{pmatrix}1&1&0&0&1&0\\0&-1&1&0&-1&1\\0&-1&1&0&0&0\\0&0&0&0&-1&1\\0&-1&0&1&0&0\\0&0&0&0&0&0\end{pmatrix}
\\
\xrightarrow{R_3\leftarrow R_3-R_5}
\begin{pmatrix}1&1&0&0&1&0\\0&-1&1&0&-1&1\\0&0&1&-1&0&0\\0&0&0&0&-1&1\\0&-1&0&1&0&0\\0&0&0&0&0&0\end{pmatrix}
\xrightarrow[R_5\leftarrow R_5+R_4]{R_5\leftarrow R_5-R_2,\ R_5\leftarrow R_5+R_3}
\begin{pmatrix}1&1&0&0&1&0\\0&-1&1&0&-1&1\\0&0&1&-1&0&0\\0&0&0&0&-1&1\\0&0&0&0&0&0\\0&0&0&0&0&0\end{pmatrix}
\end{gathered}
\]
Multiplying rows \(2\) and \(4\) of the last matrix by \(-1\) gives an echelon form with leading entries in columns
\(1,2,3,5\); so \(\rank X=4\) (Remark~\ref{rem:r4-echelon}), and in \(Xw=0\), \(w=(w_1,\dots,w_6)'\), the unknowns \(w_4\) and
\(w_6\) are free. Back substitution: row \(4\) gives \(w_5=w_6\); row \(3\) gives \(w_3=w_4\); row \(2\) gives
\(w_2=w_3-w_5+w_6=w_4\); row \(1\) gives \(w_1=-w_2-w_5=-w_4-w_6\). So the solutions are
\[
w=w_4z_1+w_6z_2,\qquad z_1=\begin{pmatrix}-1\\1\\1\\1\\0\\0\end{pmatrix},\qquad z_2=\begin{pmatrix}-1\\0\\0\\0\\1\\1\end{pmatrix}
\]
and \(z_1,z_2\), independent (entries \(4\) and \(6\) of \(t_1z_1+t_2z_2\) are \(t_1\) and \(t_2\)), form a basis of \(\mathcal N(X)\).
Check: with the columns \(c_1,\dots,c_6\) of \(X\), each plant has one variety and one treatment, so \(c_2+c_3+c_4=c_1\) and
\(c_5+c_6=c_1\), that is, \(Xz_1=0\) and \(Xz_2=0\) (Theorem~\ref{r1-thm-views}(b)).

(b) By (a), \(a'\beta\) is estimable if and only if \(a'z_1=-a_0+a_1+a_2+a_3=0\) and \(a'z_2=-a_0+a_4+a_5=0\)
(Theorem~\ref{r11-thm-estimable}(a)), with \(a=(a_0,a_1,\dots,a_5)'\).
\begin{itemize}
\item \(\alpha_1-\alpha_2\), \(a=(0,1,-1,0,0,0)'\): \(1-1=0\) and \(0=0\); estimable.
\item \(\gamma_1-\gamma_2\), \(a=(0,0,0,0,1,-1)'\): \(0=0\) and \(1-1=0\); estimable.
\item \(\mu+\alpha_1+\gamma_1\), \(a=(1,1,0,0,1,0)'\): \(-1+1=0\) and \(-1+1=0\); estimable.
\item \(\alpha_1\), \(a=(0,1,0,0,0,0)'\): \(a'z_1=1\ne0\); not estimable. \(\mu\), \(a=(1,0,0,0,0,0)'\): \(a'z_1=-1\ne0\); not
estimable.
\end{itemize}
(c) The \((j,l)\) entry of \(X'X\) is \(c_j'c_l\) for the columns \(c_1,\dots,c_6\) of \(X\) (Theorems~\ref{r1-thm-views}(d)
and~\ref{r1-thm-transpose}), and \(X'y\) holds the total, the three variety totals and the two treatment totals:
\[
X'X=\begin{pmatrix}6&2&2&2&3&3\\2&2&0&0&1&1\\2&0&2&0&1&1\\2&0&0&2&1&1\\3&1&1&1&3&0\\3&1&1&1&0&3\end{pmatrix},\qquad
X'y=\begin{pmatrix}18\\10\\6\\2\\12\\6\end{pmatrix},\qquad
X'X\hat\beta=\begin{pmatrix}10+6+2+3-3\\10+1-1\\6+1-1\\2+1-1\\5+3+1+3\\5+3+1-3\end{pmatrix}=X'y
\]
So \(\hat\beta\) solves the normal equations and is a least-squares estimate (Theorem~\ref{r10-thm-least-squares}(a)). By
Theorem~\ref{r11-thm-estimable}(b), the estimates are \(5-3=2\) for \(\alpha_1-\alpha_2\), \(1-(-1)=2\) for \(\gamma_1-\gamma_2\), and
\(0+5+1=6\) for \(\mu+\alpha_1+\gamma_1\).

(d) Row by row, \(\hat y=X\hat\beta=(5+1,\ 5-1,\ 3+1,\ 3-1,\ 1+1,\ 1-1)'=(6,4,4,2,2,0)'\), so
\(e=y-\hat y=(1,-1,-1,1,0,0)'\) and \(\mathrm{SSE}=1+1+1+1=4\). As \(\rank X=4\) by (a),
\(\hat\sigma^2=4/(6-4)=2\). Check: \(X'e=(1-1-1+1+0+0,\ 1-1,\ -1+1,\ 0+0,\ 1-1+0,\ -1+1+0)'=0\).
\end{solution}

\begin{exercise}\label{r11-dr-solution}
Four observations in \(g=3\) groups of sizes \(2,1,1\) follow the one-way model of Theorem~\ref{r11-thm-oneway},
\(y_i=\mu+\alpha_l+\varepsilon_i\) for an observation \(i\) of group \(l\), with \(\beta=(\mu,\alpha_1,\alpha_2,\alpha_3)'\), under
(A1)--(A3). Let \(y=(3,5,7,2)'\) and
\[
G_1=\diag\bigl(0,\tfrac12,1,1\bigr),\qquad G_2=\frac12\begin{pmatrix}2&-2&-2&0\\-2&3&2&0\\-2&2&4&0\\0&0&0&0\end{pmatrix}
\]
\begin{enumerate}
\item Write down \(X\), \(X'X\) and \(X'y\), and check that \(G_1\) and \(G_2\) are g-inverses of \(X'X\).
\item Compute the least-squares estimates \(b_1=G_1X'y\) and \(b_2=G_2X'y\) and their expectations. Which functions of
\(\beta\) do their entries estimate without bias?
\item Find the estimate of \(\alpha_1-\alpha_2\) from \(b_1\) and from \(b_2\), and its variance \(\sigma^2a'Ga\) with \(G=G_1\) and
with \(G=G_2\).
\end{enumerate}
\end{exercise}
\begin{solution}
(a) The rows of \(X\) are \((1,1,0,0)\) twice, \((1,0,1,0)\) and \((1,0,0,1)\); as in Example~\ref{r11-exm-oneway},
\[
X=\begin{pmatrix}1&1&0&0\\1&1&0&0\\1&0&1&0\\1&0&0&1\end{pmatrix},\qquad
X'X=\begin{pmatrix}4&2&1&1\\2&2&0&0\\1&0&1&0\\1&0&0&1\end{pmatrix},\qquad
X'y=\begin{pmatrix}17\\8\\7\\2\end{pmatrix}
\]
(the total and the group totals). Multiplying the columns of \(X'X\) by \(0,\frac12,1,1\) (Theorem~\ref{r1-thm-views}(a),(b)),
and computing \(X'XG_2\) column by column,
\[
X'XG_1=\begin{pmatrix}0&1&1&1\\0&1&0&0\\0&0&1&0\\0&0&0&1\end{pmatrix},\qquad
X'XG_2=\begin{pmatrix}1&0&0&0\\0&1&0&0\\0&0&1&0\\1&-1&-1&0\end{pmatrix}
\]
By Theorem~\ref{r1-thm-views}(c), row \(1\) of \((X'XG_1)X'X\) is the sum of rows \(2,3,4\) of \(X'X\), which is \((4,2,1,1)\),
and rows \(2,3,4\) are those of \(X'X\); rows \(1,2,3\) of \((X'XG_2)X'X\) are those of \(X'X\), and row \(4\) is
\((4,2,1,1)-(2,2,0,0)-(1,0,1,0)=(1,0,0,1)\). So \(X'XG_1X'X=X'XG_2X'X=X'X\).

(b) By Theorems~\ref{r10-thm-normal-equations}(c) and~\ref{r10-thm-least-squares}(a), both are least-squares estimates:
\[
b_1=G_1X'y=\begin{pmatrix}0\\4\\7\\2\end{pmatrix},\qquad
b_2=G_2X'y=\frac12\begin{pmatrix}34-16-14\\-34+24+14\\-34+16+28\\0\end{pmatrix}=\begin{pmatrix}2\\2\\5\\0\end{pmatrix}
\]
and \(b_1-b_2=(-2,2,2,2)'=2(-1,1,1,1)'\in\mathcal N(X)\) (Theorem~\ref{r11-thm-oneway}(b)). Multiplying the rows of \(X'X\) by
\(0,\frac12,1,1\) (Theorem~\ref{r1-thm-views}(c)), and computing \(G_2X'X\) row by row,
\[
G_1X'X=\begin{pmatrix}0&0&0&0\\1&1&0&0\\1&0&1&0\\1&0&0&1\end{pmatrix},\qquad
G_2X'X=\begin{pmatrix}1&0&0&1\\0&1&0&-1\\0&0&1&-1\\0&0&0&0\end{pmatrix}
\]
So, by Proposition~\ref{r11-prop-solution}(a), \(\expc(b_1)=(0,\ \mu+\alpha_1,\ \mu+\alpha_2,\ \mu+\alpha_3)'\): entries
\(2,3,4\) of \(b_1\), the group means \(4,7,2\), estimate the group means \(\mu+\alpha_l\). And
\(\expc(b_2)=(\mu+\alpha_3,\ \alpha_1-\alpha_3,\ \alpha_2-\alpha_3,\ 0)'\): \(b_2\) estimates the mean of group \(3\) and the
differences of groups \(1\), \(2\) from group \(3\).

(c) For \(a=(0,1,-1,0)'\), \(a'b_1=4-7=-3\) and \(a'b_2=2-5=-3\), as Theorem~\ref{r11-thm-estimable}(b) says. Further
\(a'G_1a=\frac12+1=\frac32\) and \(a'G_2a=\frac32-2\cdot1+2=\frac32\): in both cases the estimate of \(\alpha_1-\alpha_2\)
has the variance \(\frac32\sigma^2\) (Theorem~\ref{r11-thm-blue}). The estimate is \(\bar y_1-\bar y_2=c'y\) with
\(c=(\frac12,\frac12,-1,0)'\), and indeed \(\sigma^2c'c=\sigma^2\bigl(\frac14+\frac14+1\bigr)=\frac32\sigma^2\).
\end{solution}

\begin{exercise}\label{r11-dr-blue}
Add to the design \(X\) of Example~\ref{r11-exm-running} the total of its two variables:
\[
\tilde X=[X\ \ X_{\cdot1}+X_{\cdot2}]=\begin{pmatrix}1&2&3&5\\1&0&3&3\\1&2&1&3\\1&0&1&1\end{pmatrix},\qquad
y=\begin{pmatrix}9\\3\\5\\3\end{pmatrix}
\]
and consider the model \(y=\tilde X\beta+\varepsilon\), \(\beta=(\beta_0,\beta_1,\beta_2,\beta_3)'\), under (A1)--(A3).
\begin{enumerate}
\item Row-reduce \(\tilde X\) to its RREF, and find a basis of \(\mathcal N(\tilde X)\).
\item Which of \(\beta_1+\beta_3\), \(\beta_1-\beta_2\), \(\beta_3\) and \(\beta_0\) are estimable?
\item Check that \(\hat\beta=(1,2,1,0)'\) is a least-squares estimate, and find the BLUEs of the estimable functions in (b).
\item Check that \(G=\left(\begin{smallmatrix}(X'X)^{-1}&0\\0&0\end{smallmatrix}\right)\in\R^{4\times4}\) is a g-inverse of
\(\tilde X'\tilde X\), and find the variances of these BLUEs. Compare with the linear unbiased estimator \(\frac12(y_1-y_2)\) of
\(\beta_1+\beta_3\).
\end{enumerate}
\end{exercise}
\begin{solution}
(a)
\[
\begin{gathered}
\begin{pmatrix}1&2&3&5\\1&0&3&3\\1&2&1&3\\1&0&1&1\end{pmatrix}
\xrightarrow[R_4\leftarrow R_4-R_1]{R_2\leftarrow R_2-R_1,\ R_3\leftarrow R_3-R_1}
\begin{pmatrix}1&2&3&5\\0&-2&0&-2\\0&0&-2&-2\\0&-2&-2&-4\end{pmatrix}
\xrightarrow{R_4\leftarrow R_4-R_2,\ R_4\leftarrow R_4-R_3}
\begin{pmatrix}1&2&3&5\\0&-2&0&-2\\0&0&-2&-2\\0&0&0&0\end{pmatrix}
\\
\xrightarrow{R_2\leftarrow-\frac12R_2,\ R_3\leftarrow-\frac12R_3}
\begin{pmatrix}1&2&3&5\\0&1&0&1\\0&0&1&1\\0&0&0&0\end{pmatrix}
\xrightarrow{R_1\leftarrow R_1-2R_2,\ R_1\leftarrow R_1-3R_3}
\begin{pmatrix}1&0&0&0\\0&1&0&1\\0&0&1&1\\0&0&0&0\end{pmatrix}
\end{gathered}
\]
In \(\tilde Xw=0\) the unknown \(w_4\) is free, and \(w_1=0\), \(w_2=-w_4\), \(w_3=-w_4\). So \(z=(0,-1,-1,1)'\) is a basis of
\(\mathcal N(\tilde X)\), and \(\rank\tilde X=3\).

(b) By Theorem~\ref{r11-thm-estimable}(a), \(a'\beta\) is estimable if and only if \(a'z=-a_1-a_2+a_3=0\), with
\(a=(a_0,a_1,a_2,a_3)'\). So \(\beta_1+\beta_3\) (\(-1+1=0\)), \(\beta_1-\beta_2\) (\(-1+1=0\)) and \(\beta_0\) (\(0=0\)) are
estimable, and \(\beta_3\) (\(a'z=1\)) is not.

(c) By Example~\ref{r10-exm-least-squares} (the matrix \(X'X\)) and \(\tilde X_{\cdot3}=X(0,1,1)'\),
\[
\tilde X'\tilde X=\begin{pmatrix}4&4&8&12\\4&8&8&16\\8&8&20&28\\12&16&28&44\end{pmatrix},\qquad
\tilde X'y=\begin{pmatrix}20\\28\\44\\72\end{pmatrix},\qquad
\tilde X'\tilde X\hat\beta=\begin{pmatrix}4+8+8\\4+16+8\\8+16+20\\12+32+28\end{pmatrix}=\tilde X'y
\]
so \(\hat\beta\) is a least-squares estimate (Theorem~\ref{r10-thm-least-squares}(a)). By Theorem~\ref{r11-thm-blue}, the BLUEs
have the values \(\hat\beta_1+\hat\beta_3=2\), \(\hat\beta_1-\hat\beta_2=1\) and \(\hat\beta_0=1\).

(d) With
\[
(X'X)^{-1}=\frac14\begin{pmatrix}6&-1&-2\\-1&1&0\\-2&0&1\end{pmatrix}
\]
(Example~\ref{r10-exm-least-squares}), rows \(1,2,3\) of \(\tilde X'\tilde XG\) are \((I_3\ 0)\), and row \(4\) is
\(\bigl((12,16,28)(X'X)^{-1},\,0\bigr)=(0,1,1,0)\):
\[
G=\frac14\begin{pmatrix}6&-1&-2&0\\-1&1&0&0\\-2&0&1&0\\0&0&0&0\end{pmatrix},\qquad
\tilde X'\tilde XG=\begin{pmatrix}1&0&0&0\\0&1&0&0\\0&0&1&0\\0&1&1&0\end{pmatrix}
\]
So rows \(1,2,3\) of \((\tilde X'\tilde XG)\tilde X'\tilde X\) are those of \(\tilde X'\tilde X\), and row \(4\) is the sum of rows
\(2\) and \(3\), \((12,16,28,44)\) (Theorem~\ref{r1-thm-views}(c)): \(\tilde X'\tilde XG\tilde X'\tilde X=\tilde X'\tilde X\). By
Theorem~\ref{r11-thm-blue}, the variances are \(\sigma^2a'Ga\):
\begin{itemize}
\item \(\beta_1+\beta_3\), \(a=(0,1,0,1)'\): \(a'Ga=\frac14\), variance \(\frac14\sigma^2\);
\item \(\beta_1-\beta_2\), \(a=(0,1,-1,0)'\): \(a'Ga=\frac14(1-0+1)=\frac12\), variance \(\frac12\sigma^2\);
\item \(\beta_0\), \(a=(1,0,0,0)'\): \(a'Ga=\frac64\), variance \(\frac32\sigma^2\).
\end{itemize}
By (A1), \(\expc(y_1-y_2)=(\tilde x_1-\tilde x_2)'\beta=(0,2,0,2)\beta=2(\beta_1+\beta_3)\), where \(\tilde x_i'\) is row
\(i\) of \(\tilde X\); so \(\frac12(y_1-y_2)\) is a linear unbiased estimator of \(\beta_1+\beta_3\), with variance
\(\frac14(\sigma^2+\sigma^2)=\frac12\sigma^2>\frac14\sigma^2\) (Theorem~\ref{r11-thm-linear}(b), (A3)). Its value is
\(\frac12(9-3)=3\), against the BLUE \(2\).
\end{solution}

\begin{exercise}[cf.\ \citeat{bapat}{\S7.2}]\label{r10-dr-px}
In a two-way layout with one observation in each of four cells, the additive model
\(y_{ij}=\alpha_i+\gamma_j+\varepsilon_{ij}\) (\(i,j=1,2\)) reads \(y=X\beta+\varepsilon\) with
\(y=(y_{11},y_{12},y_{21},y_{22})'\), \(\beta=(\alpha_1,\alpha_2,\gamma_1,\gamma_2)'\) and
\[
X=\begin{pmatrix}1&0&1&0\\1&0&0&1\\0&1&1&0\\0&1&0&1\end{pmatrix},\qquad y=\begin{pmatrix}5\\3\\7\\1\end{pmatrix}
\]
\begin{enumerate}
\item Find \(\rank X\) by row reduction, \(X'X\), and the g-inverse \(G\) of \(X'X\) of
Theorem~\ref{r10-thm-ginv-submatrix}(b) built from its submatrix in rows and columns \(2,3,4\).
\item Find \(H=XGX'\) and check that \(\tr H=\rank X\). Find the least-squares estimate \(GX'y\), the fitted values,
the residuals and \(S(\hat\beta)\).
\item Which of \(\alpha_1-\alpha_2\), \(\gamma_1-\gamma_2\), \(\alpha_1+\gamma_1\) and \(\alpha_1\) have the same value
\(\lambda'\hat\beta\) for all least-squares estimates \(\hat\beta\)? Give these values, and \(\lambda'G\lambda\) for
\(\alpha_1-\alpha_2\) and \(\gamma_1-\gamma_2\).
\end{enumerate}
\end{exercise}
\begin{solution}
(a)
\[
\begin{gathered}
X\xrightarrow{R_2\leftrightarrow R_3}
\begin{pmatrix}1&0&1&0\\0&1&1&0\\1&0&0&1\\0&1&0&1\end{pmatrix}
\xrightarrow[R_4\leftarrow R_4-R_2]{R_3\leftarrow R_3-R_1}
\begin{pmatrix}1&0&1&0\\0&1&1&0\\0&0&-1&1\\0&0&-1&1\end{pmatrix}
\\
\xrightarrow{R_4\leftarrow R_4-R_3}
\begin{pmatrix}1&0&1&0\\0&1&1&0\\0&0&-1&1\\0&0&0&0\end{pmatrix}
\xrightarrow{R_3\leftarrow-R_3}
\begin{pmatrix}1&0&1&0\\0&1&1&0\\0&0&1&-1\\0&0&0&0\end{pmatrix}
\end{gathered}
\]
The last matrix is an echelon form with \(3\) pivots, so \(\rank X=3\) (Definition~\ref{def:r4-rank}). The entries of
\(X'X\) are the inner products of the columns of \(X\), and \(\rank(X'X)=3\) (Theorem~\ref{thm:r4-rank-ata}(b)):
\[
X'X=\begin{pmatrix}2&0&1&1\\0&2&1&1\\1&1&2&0\\1&1&0&2\end{pmatrix},\qquad
W=\begin{pmatrix}2&1&1\\1&2&0\\1&0&2\end{pmatrix},\qquad
\operatorname{adj}W=\begin{pmatrix}4&-2&-2\\-2&3&1\\-2&1&3\end{pmatrix}
\]
where \(W\) is the submatrix in rows and columns \(2,3,4\), with \(\det W=1\cdot(0-2)+2\cdot(4-1)=4\) (expansion along
row \(3\), Theorem~\ref{r5d-thm-expansion}). So \(W^{-1}=\frac14\operatorname{adj}W\) (Theorem~\ref{r5d-thm-adjugate}), and by
Theorem~\ref{r10-thm-ginv-submatrix}(b)
\[
G=\frac14\begin{pmatrix}0&0&0&0\\0&4&-2&-2\\0&-2&3&1\\0&-2&1&3\end{pmatrix}
\]
\needspace{5\baselineskip}
(b)
\begin{itemize}
\item The rows of \(4XG\) are sums of two rows of \(4G\) (Theorem~\ref{r1-thm-views}(c)), and then
\[
4XG=\begin{pmatrix}0&-2&3&1\\0&-2&1&3\\0&2&1&-1\\0&2&-1&1\end{pmatrix},\qquad
H=XGX'=\frac14\begin{pmatrix}3&1&1&-1\\1&3&-1&1\\1&-1&3&1\\-1&1&1&3\end{pmatrix}
\]
and \(\tr H=\frac14(3+3+3+3)=3=\rank X\) (Theorem~\ref{r10-thm-PX}(b)).
\item \(X'y=(5+3,\ 7+1,\ 5+7,\ 3+1)'=(8,8,12,4)'\), and
\[
\hat\beta=GX'y=\frac14\begin{pmatrix}0\\32-24-8\\-16+36+4\\-16+12+12\end{pmatrix}=\begin{pmatrix}0\\0\\6\\2\end{pmatrix}
\]
is a least-squares estimate (Theorems~\ref{r10-thm-normal-equations}(c) and~\ref{r10-thm-least-squares}(a)).
\item \(\hat y=X\hat\beta=(\hat\alpha_1+\hat\gamma_1,\ \hat\alpha_1+\hat\gamma_2,\ \hat\alpha_2+\hat\gamma_1,\ \hat\alpha_2+\hat\gamma_2)'
=(6,2,6,2)'\); indeed \(Hy=\frac14(15+3+7-1,\ 5+9-7+1,\ 5-3+21+1,\ -5+3+7+3)'=(6,2,6,2)'\).
\item \(e=y-\hat y=(-1,1,1,-1)'\) and \(S(\hat\beta)=e'e=4\); check (Theorem~\ref{r10-thm-PX}(d)):
\(y'y-\hat\beta'X'y=(25+9+49+1)-(0+0+72+8)=84-80=4\).
\end{itemize}
(c)
\begin{itemize}
\item By Theorem~\ref{r1-thm-views}(c), \(4\row_2(GX'X)=4(0,2,1,1)-2(1,1,2,0)-2(1,1,0,2)=(-4,4,0,0)\), and likewise for
rows \(3\), \(4\):
\[
GX'X=\begin{pmatrix}0&0&0&0\\-1&1&0&0\\1&0&1&0\\1&0&0&1\end{pmatrix}
\]
\item By Corollary~\ref{r10-cor-estimable-algebra}(iii): \(\alpha_1-\alpha_2\) passes, as
\(\row_1(GX'X)-\row_2(GX'X)=(1,-1,0,0)\), with value \(\hat\alpha_1-\hat\alpha_2=0\); \(\gamma_1-\gamma_2\) passes
(\(\row_3-\row_4=(0,0,1,-1)\)), with value \(6-2=4\); \(\alpha_1+\gamma_1\) passes (\(\row_1+\row_3=(1,0,1,0)\)), with value
\(0+6=6\); \(\alpha_1\) fails, as \(\row_1(GX'X)=0\ne(1,0,0,0)\).
\item Indeed \(I_4-GX'X\) has first column \((1,1,-1,-1)'\) and its other columns \(0\), so the least-squares estimates are
\((t,\ t,\ 6-t,\ 2-t)'\) (\(t\in\R\)) (Theorem~\ref{r10-thm-normal-equations}(c)): \(\hat\alpha_1=t\) changes, the three
other values do not.
\item \(\alpha_1-\alpha_2\), \(\lambda=(1,-1,0,0)'\): \(\lambda'G\lambda=0-0-0+1=1\); \(\gamma_1-\gamma_2\),
\(\lambda=(0,0,1,-1)'\): \(\lambda'G\lambda=\frac14(3-1-1+3)=1\). Both values are the same for every g-inverse of \(X'X\)
(Corollary~\ref{r10-cor-invariant-quadratic} with \(A=X\)).
\end{itemize}
\end{solution}

\begin{exercise}[\citeat{bapat}{Ch.~7, Exercise~2}]\label{r11-ex-bapat-estimable}
Consider the model \(\expc(y_1)=2\beta_1-\beta_2-\beta_3\), \(\expc(y_2)=\beta_2-\beta_4\), \(\expc(y_3)=\beta_2+\beta_3-2\beta_4\) under
(A1)--(A3), with \(\beta=(\beta_1,\beta_2,\beta_3,\beta_4)'\). Determine the estimable functions.
\end{exercise}
\begin{solution}
The model is \(\expc(y)=X\beta\) with
\[
\begin{gathered}
X=\begin{pmatrix}2&-1&-1&0\\0&1&0&-1\\0&1&1&-2\end{pmatrix}
\xrightarrow{R_3\leftarrow R_3-R_2}
\begin{pmatrix}2&-1&-1&0\\0&1&0&-1\\0&0&1&-1\end{pmatrix}
\\
\xrightarrow{R_1\leftarrow R_1+R_2,\ R_1\leftarrow R_1+R_3}
\begin{pmatrix}2&0&0&-2\\0&1&0&-1\\0&0&1&-1\end{pmatrix}
\xrightarrow{R_1\leftarrow\frac12R_1}
\begin{pmatrix}1&0&0&-1\\0&1&0&-1\\0&0&1&-1\end{pmatrix}
\end{gathered}
\]
The last matrix is the RREF of \(X\): \(\rank X=3\), and in \(Xw=0\) the unknown \(w_4\) is free, with \(w_1=w_2=w_3=w_4\). So
\(\mathcal N(X)=\Span\{(1,1,1,1)'\}\); check: \(X(1,1,1,1)'=(2-1-1,\ 1-1,\ 1+1-2)'=0\). By
Theorem~\ref{r11-thm-estimable}(a), \(a'\beta=a_1\beta_1+a_2\beta_2+a_3\beta_3+a_4\beta_4\) is estimable if and only if
\(a\perp(1,1,1,1)'\), that is, \(a_1+a_2+a_3+a_4=0\). For instance \(\beta_1-\beta_2\) is estimable, and \(\beta_1\) is not.
\end{solution}

\begin{exercise}\label{r12-ex-cls-drill}
For the four units of Example~\ref{r10-exm-least-squares}, consider \(H_0\colon\beta_1=\beta_2\), that is, \(K\beta=0\)
with \(K=(0,1,-1)\).
\begin{enumerate}
\item Use Theorem~\ref{r12-thm-cls} to compute \(V\), \(\tilde\beta\) and \(\mathrm{SSE}_K-\mathrm{SSE}\), and check \(\mathrm{SSE}_K\)
from the residuals \(y-X\tilde\beta\).
\item Assuming the normal linear model, compute the statistic \(F\) of Theorem~\ref{r12-thm-f} and check \(F=T_0^2\).
\end{enumerate}
\end{exercise}
\begin{solution}
(a) \(X\) has full column rank, so \(H_0\) is testable (Definition~\ref{r11-def-testable}) and \(G=(X'X)^{-1}\). With
\((X'X)^{-1}\) and \(\hat\beta=(1,2,1)'\) of Example~\ref{r10-exm-least-squares}, \((X'X)^{-1}K'=\frac14(1,1,-1)'\),
\(V=K(X'X)^{-1}K'=\frac14(1+1)=\frac12\) and \(K\hat\beta=2-1=1\). So
\[
\tilde\beta=\begin{pmatrix}1\\2\\1\end{pmatrix}-\frac14\begin{pmatrix}1\\1\\-1\end{pmatrix}\cdot2\cdot1
=\begin{pmatrix}1/2\\3/2\\3/2\end{pmatrix},\qquad \mathrm{SSE}_K-\mathrm{SSE}=1\cdot2\cdot1=2
\]
and \(\mathrm{SSE}_K=4+2=6\). Check: \(K\tilde\beta=\frac32-\frac32=0\), \(X\tilde\beta=(8,5,5,2)'\), \(y-X\tilde\beta=(1,-2,0,1)'\)
and \(1+4+0+1=6\).

(b) \(n-r=1\), so \(F=(2/1)/(4/1)=\frac12\). With \(a=(0,1,-1)'\), \(\operatorname{SE}(a'\hat\beta)=s\sqrt{a'(X'X)^{-1}a}=2\sqrt{1/2}=\sqrt2\)
and \(T_0=1/\sqrt2\), so \(T_0^2=\frac12=F\) (Theorem~\ref{r12-thm-f}(d)).
\end{solution}

\begin{exercise}[\citeat{bapat}{Ch.~7, Exercise~12}]\label{r11-ex-bapat-restricted}
Consider the model \(\expc(y_1)=\beta_1+2\beta_2\), \(\expc(y_2)=2\beta_1\), \(\expc(y_3)=\beta_1+\beta_2\) under (A1)--(A3), with
\(\beta=(\beta_1,\beta_2)'\). Find the residual sum of squares subject to the restriction \(\beta_1=\beta_2\), that is,
\(\mathrm{SSE}_K\) of Theorem~\ref{r12-thm-cls} for \(K=(1,-1)\) and \(m=0\), as a function of \(y\).
\end{exercise}
\begin{solution}
\begin{itemize}
\item The design and the matrices of Theorem~\ref{r12-thm-cls} are
\[
X=\begin{pmatrix}1&2\\2&0\\1&1\end{pmatrix},\qquad X'X=\begin{pmatrix}6&3\\3&5\end{pmatrix},\qquad
(X'X)^{-1}=\frac1{21}\begin{pmatrix}5&-3\\-3&6\end{pmatrix},\qquad X'y=\begin{pmatrix}y_1+2y_2+y_3\\2y_1+y_3\end{pmatrix}
\]
with \(\det(X'X)=30-9=21\); indeed \(X'X\cdot\frac1{21}\begin{pmatrix}5&-3\\-3&6\end{pmatrix}=\frac1{21}\begin{pmatrix}30-9&-18+18\\15-15&-9+30\end{pmatrix}=I_2\). So
\(\rank X=2\), \(H_0\colon\beta_1=\beta_2\) is testable (Definition~\ref{r11-def-testable}) and \(G=(X'X)^{-1}\).
\item \(\hat\beta=(X'X)^{-1}X'y=\frac1{21}(-y_1+10y_2+2y_3,\ 9y_1-6y_2+3y_3)'\), so \(K\hat\beta=\frac1{21}(-10y_1+16y_2-y_3)\), and
\(V=K(X'X)^{-1}K'=\frac1{21}(5+3+3+6)=\frac{17}{21}\).
\item \(w=(2,1,-4)'\) has \(X'w=(2+2-4,\ 4-4)'=0\), and \(\dim\mathcal C(X)^\perp=3-2=1\) (Theorem~\ref{r6-thm-complement}(c)); so
\(\mathcal C(X)^\perp=\Span\{w\}\). The residual vector \(e\) lies in it (Theorem~\ref{r10-thm-least-squares}(b)), \(e=tw\), and
\(w'y=w'(\hat y+e)=w'e=21t\). Hence \(\mathrm{SSE}=e'e=21t^2=(2y_1+y_2-4y_3)^2/21\).
\item By Theorem~\ref{r12-thm-cls}(c),
\[
\begin{aligned}
\mathrm{SSE}_K&=\mathrm{SSE}+\frac{(K\hat\beta)^2}V=\frac{17(2y_1+y_2-4y_3)^2+(-10y_1+16y_2-y_3)^2}{357}\\
&=\frac{8y_1^2+13y_2^2+13y_3^2-12y_1y_2-12y_1y_3-8y_2y_3}{17}
\end{aligned}
\]
since the numerator of the middle fraction is
\(168y_1^2+273y_2^2+273y_3^2-252y_1y_2-252y_1y_3-168y_2y_3\), which is \(21\) times that of the last.
\end{itemize}
\end{solution}

\begin{exercise}\label{r12-ex-f-r2}
\begin{enumerate}
\item A regression with intercept and \(k=4\) further columns, of full column rank, on \(n=105\) observations has overall
\(F=20\). Find \(R^2\).
\item A regression with intercept and \(6\) further columns, of full column rank, on \(n=40\) observations has \(R^2=\frac12\);
without three of these columns, \(R^2=\frac25\). Compute the statistic \(F\) for the hypothesis that the three coefficients are \(0\),
and state when the test of level \(\alpha\) rejects it.
\end{enumerate}
\end{exercise}
\begin{solution}
(a) By Theorem~\ref{r12-thm-f}(c) (overall \(F\) test), \(F(1-R^2)/(n-k-1)=R^2/k\), so \(kF(1-R^2)=(n-k-1)R^2\) and
\(R^2=kF/(kF+n-k-1)=80/180=\frac49\).

(b) \(r=7\), \(n-r=33\), \(q=3\); under (A4) the hypothesis is testable, and by Theorem~\ref{r12-thm-f}(c),
\(F=\frac{(1/2-2/5)/3}{(1-1/2)/33}=\frac{1/30}{1/66}=\frac{11}5\). The test rejects if and only if
\(\frac{11}5>F_{3,33,1-\alpha}\) (Theorem~\ref{r12-thm-f}(b)).
\end{solution}

\begin{exercise}\label{r11-ex-two-sample}
Four plants given a new fertilizer have heights \(12,15,13,16\), and five control plants \(10,9,12,11,8\). Model them as
the one-way layout of Theorem~\ref{r11-thm-oneway} with \(g=2\) groups (\(n_1=4\), \(n_2=5\)), the design \(X_1=[d_1\ d_2]\)
and \(\beta=(\mu_1,\mu_2)'\), in the normal linear model. Compute the statistic \(T_0\) of Theorem~\ref{r12-thm-t}(b) for
\(H_0\colon\mu_1=\mu_2\) and the confidence interval of level \(1-\alpha\) for \(\mu_1-\mu_2\), and state when the test of
level \(\alpha\) rejects \(H_0\).
\end{exercise}
\begin{solution}
\begin{itemize}
\item \(X_1'X_1=\diag(4,5)\) and \(X_1'y=(56,50)'\) (the group totals), so \(\hat\beta=(14,10)'\), the group means. The
residuals are the deviations from the group means, \(-2,1,-1,2\) and \(0,-1,2,1,-2\), so
\(\mathrm{SSE}=(4+1+1+4)+(0+1+4+1+4)=20\). Here \(r=2\), \(n-r=7\) and \(s^2=\frac{20}7\).
\item For \(a=(1,-1)'\), \(a'\hat\beta=4\) and \(a'(X_1'X_1)^{-1}a=\frac14+\frac15=\frac9{20}\), so
\(\operatorname{SE}(a'\hat\beta)=\sqrt{\frac{20}7\cdot\frac9{20}}=\frac3{\sqrt7}\), and \(T_0=4\big/\frac3{\sqrt7}=\frac{4\sqrt7}3\).
\item By Theorem~\ref{r12-thm-t}(b), the interval is
\(\bigl[4-\frac3{\sqrt7}t_{7,1-\alpha/2},\ 4+\frac3{\sqrt7}t_{7,1-\alpha/2}\bigr]\), and the test rejects \(H_0\) if and only if
\(t_{7,1-\alpha/2}<\frac{4\sqrt7}3\), that is, exactly when \(0\) lies outside the interval.
\end{itemize}
\end{solution}

\begin{exercise}\label{r1-dr-group-means}
Six units form three groups, \(\{1,2\}\), \(\{3,4\}\) and \(\{5,6\}\), and \(y=(1,3,6,8,4,2)'\) is measured on them.
Let \(H\in\R^{6\times6}\) be the matrix that replaces each \(y_i\) by the mean of its group (the hat matrix of the
one-way design, Theorem~\ref{r11-thm-oneway-anova}(a)).
\begin{enumerate}
\item Write down \(H\) and compute \(Hy\) and \(\bar J_6y\).
\item Compute \(C_6y\), \((H-\bar J_6)y\) and \((I_6-H)y\), and check that \(C_6y=(H-\bar J_6)y+(I_6-H)y\).
\item Compute \(y'C_6y\), \(y'(H-\bar J_6)y\) and \(y'(I_6-H)y\), and check that the first is the sum of the other two.
\end{enumerate}
\end{exercise}
\begin{solution}
(a) Row \(i\) of \(H\) averages the two entries of the group of unit \(i\):
\[
H=\tfrac12\begin{pmatrix}1&1&0&0&0&0\\1&1&0&0&0&0\\0&0&1&1&0&0\\0&0&1&1&0&0\\0&0&0&0&1&1\\0&0&0&0&1&1\end{pmatrix},\qquad
Hy=\begin{pmatrix}2\\2\\7\\7\\3\\3\end{pmatrix},\qquad \bar J_6y=\bar y\one_6=4\,\one_6
\]
with the group means \(2\), \(7\), \(3\) and the mean \(\bar y=\frac{24}6=4\).

(b) By Example~\ref{r1-exm-centering}, \(C_6y=y-4\,\one_6\); further \((H-\bar J_6)y=Hy-4\,\one_6\) and
\((I_6-H)y=y-Hy\):
\[
C_6y=\begin{pmatrix}-3\\-1\\2\\4\\0\\-2\end{pmatrix},\qquad
(H-\bar J_6)y=\begin{pmatrix}-2\\-2\\3\\3\\-1\\-1\end{pmatrix},\qquad
(I_6-H)y=\begin{pmatrix}-1\\1\\-1\\1\\1\\-1\end{pmatrix}
\]
and the last two add up to \(C_6y\).

(c) \(H\) is symmetric, every row and every column of \(H\) adds up to \(1\), and \(H^2=H\), since each row of \(H^2\)
is the average of two equal rows of \(H\) (Theorem~\ref{r1-thm-views}(c)). So
\(H\bar J_6=\frac16(H\one)\one'=\bar J_6=\bar J_6H\), and with \(\bar J_6^2=\bar J_6\) (Proposition~\ref{r1-prop-centering}),
\((H-\bar J_6)^2=H-\bar J_6\); \(C_6\) and \(I_6-H\) are symmetric and idempotent (Propositions~\ref{r1-prop-centering}
and~\ref{r1-prop-idempotent}(b)). For a symmetric idempotent \(P\), \(y'Py=(Py)'(Py)\) is the sum of the squares of the
entries of \(Py\):
\[
y'C_6y=9+1+4+16+0+4=34,\qquad y'(H-\bar J_6)y=4+4+9+9+1+1=28,\qquad y'(I_6-H)y=6
\]
and \(34=28+6\), as Theorem~\ref{r11-thm-oneway-anova}(b) says.
\end{solution}

\begin{exercise}\label{r3-dr-grouped}
Five observations form three groups: observations \(1,2\) (group \(1\)), observations \(3,4\)
(group \(2\)) and observation \(5\) (group \(3\)), with the indicators \(d_1=(1,1,0,0,0)'\),
\(d_2=(0,0,1,1,0)'\) and \(d_3=(0,0,0,0,1)'\); let \(U=\Span\{d_1,d_2,d_3\}\) and \(y=(3,5,4,8,2)'\).
\begin{enumerate}
\item Find the between-group part \(u\) and the within-group part \(y-u\) of \(y\), as in
Example~\ref{r3-exm-between-within}, and check that the entries of \(y-u\) add up to \(0\) within each
group.
\item Find the coordinate vectors of \(u\) relative to \(\mathcal B=(d_1,d_2,d_3)\) and to
\(\tilde{\mathcal B}=(\one,d_2,d_3)\).
\item Find the change-of-basis matrix \(\hat P\) from \(\hat{\mathcal B}=(\one,d_1,d_2)\), in which group
\(3\) is the baseline, to \(\mathcal B\); check that
\[
\hat Q=\begin{pmatrix}0&0&1\\1&0&-1\\0&1&-1\end{pmatrix}
\]
is its inverse, and find \([u]_{\hat{\mathcal B}}\).
\end{enumerate}
\end{exercise}
\begin{solution}
(a) The group means of \(y\) are \((3+5)/2=4\), \((4+8)/2=6\) and \(2/1=2\). So
\(u=4d_1+6d_2+2d_3=(4,4,6,6,2)'\) and \(y-u=(-1,1,-2,2,0)'\); within the groups,
\(-1+1=0\), \(-2+2=0\) and \(0=0\).

(b) \(d_1,d_2,d_3\) are the columns of the design \(X_1\) of Theorem~\ref{r11-thm-oneway} for this layout, which
are linearly independent (Theorem~\ref{r11-thm-oneway}(a)), so \(\mathcal B\) is an
ordered basis of \(U\), and \([u]_{\mathcal B}=(4,6,2)'\). Since \(\one=d_1+d_2+d_3\), the matrix with the columns
\([\one]_{\mathcal B}=(1,1,1)'\), \([d_2]_{\mathcal B}=(0,1,0)'\), \([d_3]_{\mathcal B}=(0,0,1)'\) and its inverse are
\[
T=\begin{pmatrix}1&0&0\\1&1&0\\1&0&1\end{pmatrix},\qquad S=\begin{pmatrix}1&0&0\\-1&1&0\\-1&0&1\end{pmatrix}
\]
as \(\row_1(TS)=\row_1(S)\), \(\row_2(TS)=\row_1(S)+\row_2(S)\) and \(\row_3(TS)=\row_1(S)+\row_3(S)\)
(Theorem~\ref{r1-thm-views}(c)) give \(TS=I_3\) (Theorem~\ref{r2-thm-invertible}). So \(\tilde{\mathcal B}\) is an
ordered basis of \(U\), and \(T\) is the change-of-basis matrix from \(\tilde{\mathcal B}\) to \(\mathcal B\)
(Theorem~\ref{r3-thm-basis-from-P}); and
\[
[u]_{\tilde{\mathcal B}}=S[u]_{\mathcal B}
=\begin{pmatrix}1&0&0\\-1&1&0\\-1&0&1\end{pmatrix}\begin{pmatrix}4\\6\\2\end{pmatrix}
=\begin{pmatrix}4\\2\\-2\end{pmatrix}
\]
(Theorem~\ref{r3-thm-change-basis}(b)). Check: \(4\one+2d_2-2d_3=(4,4,6,6,2)'\).

(c) \([\one]_{\mathcal B}=(1,1,1)'\), \([d_1]_{\mathcal B}=(1,0,0)'\) and \([d_2]_{\mathcal B}=(0,1,0)'\), so
\[
\hat P=\begin{pmatrix}1&1&0\\1&0&1\\1&0&0\end{pmatrix},\qquad
\hat P\hat Q=\begin{pmatrix}1&1&0\\1&0&1\\1&0&0\end{pmatrix}\begin{pmatrix}0&0&1\\1&0&-1\\0&1&-1\end{pmatrix}
=\begin{pmatrix}1&0&0\\0&1&0\\0&0&1\end{pmatrix}
\]
So \(\hat P\) is invertible with \(\hat P^{-1}=\hat Q\) (Theorem~\ref{r2-thm-invertible}),
\(\hat{\mathcal B}\) is an ordered basis of \(U\) (Theorem~\ref{r3-thm-basis-from-P}), and
\[
[u]_{\hat{\mathcal B}}=\hat Q[u]_{\mathcal B}
=\begin{pmatrix}0&0&1\\1&0&-1\\0&1&-1\end{pmatrix}\begin{pmatrix}4\\6\\2\end{pmatrix}
=\begin{pmatrix}2\\2\\4\end{pmatrix}
\]
the value \(2\) of group \(3\) and the differences \(4-2\), \(6-2\) of groups \(1\) and \(2\) from it.
Check: \(2\one+2d_1+4d_2=(4,4,6,6,2)'\).
\end{solution}

\begin{exercise}\label{r6-ex-nested-drill}
In \(\R^4\) let \(U=\Span\{\one_4\}\), \(W=\{x:x_1+x_2=x_3+x_4\}\) and \(y=(5,1,2,0)'\).
\begin{enumerate}
\item Show that \(U\subseteq W\), and find \(W^\perp\), \(P_W\) and \(P_U\).
\item Check that \(P_WP_U=P_U\).
\item Find an orthonormal basis of \(W\cap U^\perp\), and check that the projection matrix onto \(W\cap U^\perp\) is
\(P_W-P_U\).
\item Write \(y=P_Uy+(P_W-P_U)y+(I_4-P_W)y\), and check that the squared lengths of the three terms add up to
\(\lVert y\rVert^2\).
\end{enumerate}
\end{exercise}
\begin{solution}
Let \(h=(1,1,-1,-1)'\), so that \(W=\{x:h'x=0\}\).

(a) \(h'\one_4=0\), so \(\one_4\in W\) and \(U\subseteq W\). By Theorem~\ref{r6-thm-complement}(a),
\(W=\mathcal N(h')=\Span\{h\}^\perp\), so \(W^\perp=\Span\{h\}\) (Theorem~\ref{r6-thm-perp-perp}(a)). By
Theorem~\ref{r6-thm-projection-matrix}(c) for \(\Span\{h\}\), and by Example~\ref{r6-exm-centering} and
Theorem~\ref{r6-thm-projection-matrix}(b) for \(P_{\Span\{h\}}\), \(P_W=I_4-P_{\Span\{h\}}=I_4-\frac14hh'\);
and \(P_U=\bar J_4\) (Corollary~\ref{r6-cor-centering}(a)):
\[
P_W=\frac14\begin{pmatrix}3&-1&1&1\\-1&3&1&1\\1&1&3&-1\\1&1&-1&3\end{pmatrix},\qquad
P_U=\frac14\begin{pmatrix}1&1&1&1\\1&1&1&1\\1&1&1&1\\1&1&1&1\end{pmatrix}
\]

(b) Every row of \(P_W\) has sum \(1\), that is, \(P_W\one=\one\) (as \(\one\in W\)). So
\(P_WP_U=\frac14(P_W\one)\one'=\frac14\one\one'=P_U\), as Theorem~\ref{r6-thm-nested}(a) says.

(c) \(x\in W\cap U^\perp\) means \(x_1+x_2=x_3+x_4\) and \(x_1+x_2+x_3+x_4=0\), that is, \(x_1+x_2=0=x_3+x_4\). So
\(W\cap U^\perp=\{(s,-s,t,-t)':s,t\in\R\}\), with the orthonormal basis \(u_1=\frac1{\sqrt2}(1,-1,0,0)'\),
\(u_2=\frac1{\sqrt2}(0,0,1,-1)'\). By Theorem~\ref{r6-thm-projection-matrix}(a),
\[
P_{W\cap U^\perp}=u_1u_1'+u_2u_2'=\frac12\begin{pmatrix}1&-1&0&0\\-1&1&0&0\\0&0&1&-1\\0&0&-1&1\end{pmatrix}
\]
which is \(P_W-P_U\), as Theorem~\ref{r6-thm-nested}(b) says; for instance, row \(1\) of \(P_W-P_U\) is
\(\frac14(3-1,\ -1-1,\ 1-1,\ 1-1)=\frac12(1,-1,0,0)\). The dimension of \(W\cap U^\perp\) is \(2=3-1\).

(d) \(P_Uy=\bar y\one_4=2\,\one_4\); \(h'y=5+1-2-0=4\), so \((I_4-P_W)y=P_{W^\perp}y=\frac{h'y}{h'h}h=h\) and
\(P_Wy=y-h=(4,0,3,1)'\); and \((P_W-P_U)y=(2,-2,1,-1)'\). So
\[
\begin{pmatrix}5\\1\\2\\0\end{pmatrix}=\begin{pmatrix}2\\2\\2\\2\end{pmatrix}+\begin{pmatrix}2\\-2\\1\\-1\end{pmatrix}
+\begin{pmatrix}1\\1\\-1\\-1\end{pmatrix},\qquad 30=16+10+4
\]
as Theorem~\ref{r6-thm-nested}(c) says; also \(\lVert P_Wy\rVert^2=16+0+9+1=26=16+10\).
\end{solution}

\begin{exercise}\label{r11-ex-projection-drill}
Let \(y\sim N_4(\mu,\sigma^2I_4)\) with \(\sigma>0\), and
\[
P=\frac12\begin{pmatrix}1&1&0&0\\1&1&0&0\\0&0&1&1\\0&0&1&1\end{pmatrix}
\]
\begin{enumerate}
\item Show that \(P\) and \(I_4-P\) are symmetric and idempotent of rank \(2\), that
\(y'Py=\frac12(y_1+y_2)^2+\frac12(y_3+y_4)^2\), and that \(y'(I_4-P)y=\frac12(y_1-y_2)^2+\frac12(y_3-y_4)^2\).
\item Find \(\expc(y'Py)\) and \(\expc\bigl(y'(I_4-P)y\bigr)\) for \(\mu=(1,3,5,7)'\).
\item Let \(\mu=(m_1,m_1,m_2,m_2)'\). Show that \(y'(I_4-P)y/\sigma^2\sim\chi^2_2\), that \(y'(I_4-P)y\) is
independent of \(Py\), and that \(\frac12y'(I_4-P)y\) is an unbiased estimator of \(\sigma^2\).
\end{enumerate}
\end{exercise}
\begin{solution}
(a) \(P=\frac12\diag(J_2,J_2)\) is symmetric, and \(J_2^2=2J_2\) (Proposition~\ref{r1-prop-centering}), so
\(P^2=\frac14\diag(J_2^2,J_2^2)=\frac12\diag(J_2,J_2)=P\) (Theorem~\ref{r5b-thm-block}(c)). By Theorem~\ref{r6-thm-idempotent}(b),(c), \(I_4-P\) is idempotent, \(\rank P=\tr P=2\)
and \(\rank(I_4-P)=\tr(I_4-P)=2\). As \(Py=\frac12(y_1+y_2,\ y_1+y_2,\ y_3+y_4,\ y_3+y_4)'\),
\(y'Py=\frac12(y_1+y_2)^2+\frac12(y_3+y_4)^2\); and \(y'(I_4-P)y=y'y-y'Py\), where
\(y_1^2+y_2^2-\frac12(y_1+y_2)^2=\frac12(y_1-y_2)^2\), and likewise for \(y_3,y_4\).

(b) By Theorem~\ref{r11-thm-quadratic}, with \(\tr(P\sigma^2I_4)=\sigma^2\tr P=2\sigma^2\),
\(\expc(y'Py)=2\sigma^2+\mu'P\mu=2\sigma^2+\frac12\cdot4^2+\frac12\cdot12^2=2\sigma^2+80\) and
\(\expc\bigl(y'(I_4-P)y\bigr)=2\sigma^2+\frac12\cdot2^2+\frac12\cdot2^2=2\sigma^2+4\).

(c) Let \(\varepsilon=y-\mu\sim N_4(0,\sigma^2I_4)\) (Theorem~\ref{r14-thm-mvn-linear}). Here \((I_4-P)\mu=0\), so
\((I_4-P)y=(I_4-P)\varepsilon\) and, as \(I_4-P\) is symmetric and idempotent,
\(y'(I_4-P)y=\lVert(I_4-P)y\rVert^2=\varepsilon'(I_4-P)\varepsilon\). Theorem~\ref{r15-thm-quadratic-chisq} gives
\(y'(I_4-P)y/\sigma^2\sim\chi^2_2\). As \(P(I_4-P)=P-P^2=0\), Theorem~\ref{r15-thm-indep}(b) gives that
\(P\varepsilon\) and \(\varepsilon'(I_4-P)\varepsilon\) are independent, and \(Py=P\mu+P\varepsilon\) is a function of
\(P\varepsilon\) (N3). \(\expc\bigl(y'(I_4-P)y\bigr)=2\sigma^2+0\) (Theorem~\ref{r11-thm-quadratic}), so
\(\frac12y'(I_4-P)y\) is unbiased for \(\sigma^2\) for all values of \(m_1,m_2,\sigma^2\)
(Definition~\ref{r11-def-estimator}).
\end{solution}

\begin{exercise}\label{r15-dr-oneway}
In the one-way layout of Example~\ref{r15-exm-oneway} with \(g=3\) groups of \(m=2\) observations, let \(x=(5,7,2,4,8,10)'\).
\begin{enumerate}
\item Compute the group means, \(\bar x\), \(\mathrm{SS}_A\) and \(\mathrm{SS}_E\).
\item Write out \(A_2\) and \(A_3\), compute \(A_2x\) and \(A_3x\), and check that \(x'A_2x=\mathrm{SS}_A\) and \(x'A_3x=\mathrm{SS}_E\).
\item Compute \(F=\bigl(\mathrm{SS}_A/(g-1)\bigr)/\bigl(\mathrm{SS}_E/(n-g)\bigr)\), and state when the test of level \(\alpha\) that
rejects \(H_0\colon\mu_1=\mu_2=\mu_3\) for large \(F\) rejects it.
\end{enumerate}
\end{exercise}
\begin{solution}
(a) The group means are \(6\), \(3\) and \(9\), and \(\bar x=36/6=6\). So \(\mathrm{SS}_A=2\bigl(0^2+(-3)^2+3^2\bigr)=36\) and
\(\mathrm{SS}_E=1+1+1+1+1+1=6\).

(b) Here \(\frac12D=\diag(\frac12J_2,\frac12J_2,\frac12J_2)\), and
\[
A_2=\frac16\begin{pmatrix}2&2&-1&-1&-1&-1\\2&2&-1&-1&-1&-1\\-1&-1&2&2&-1&-1\\-1&-1&2&2&-1&-1\\-1&-1&-1&-1&2&2\\-1&-1&-1&-1&2&2\end{pmatrix},\qquad
A_3=\frac12\begin{pmatrix}1&-1&0&0&0&0\\-1&1&0&0&0&0\\0&0&1&-1&0&0\\0&0&-1&1&0&0\\0&0&0&0&1&-1\\0&0&0&0&-1&1\end{pmatrix}
\]
\(\frac12Dx=(6,6,3,3,9,9)'\) and \(\bar J_6x=6\one\), so \(A_2x=(0,0,-3,-3,3,3)'\) and \(A_3x=x-\frac12Dx=(-1,1,-1,1,-1,1)'\). Then
\(x'A_2x=0+0-6-12+24+30=36\) and \(x'A_3x=-5+7-2+4-8+10=6\).

(c) \(F=(36/2)/(6/3)=9\). If the three group means \(\mu_i\) are equal, then \(F\sim F_{2,3}\) (Example~\ref{r15-exm-oneway}), and
the test that rejects when \(F>F_{2,3,1-\alpha}\) has level \(\alpha\) (Definition~\ref{r15-def-t-F}(c)): it rejects \(H_0\) if
and only if \(F_{2,3,1-\alpha}<9\).
\end{solution}

\begin{exercise}[cf.\ \citeat{bapat}{\S8.3}]\label{r15-dr-factorial}
Four observations \(y_{ij}\) (\(i,j=1,2\)) are ordered \(y=(y_{11},y_{12},y_{21},y_{22})'\). Let \(a=(1,1,-1,-1)'\), \(b=(1,-1,1,-1)'\),
\(c=(1,-1,-1,1)'\), and \(A_1=\bar J_4\), \(A_2=\frac14aa'\), \(A_3=\frac14bb'\), \(A_4=\frac14cc'\).
\begin{enumerate}
\item Check that \(A_1+A_2+A_3+A_4=I_4\) and that the ranks add up to \(4\). What does Theorem~\ref{r15-thm-cochran-algebraic} give?
\item For \(y=(7,3,4,2)'\), compute \(y'A_iy\) (\(i=1,\dots,4\)), and check that they add up to \(y'y\).
\item In the model \(y_{ij}=\mu+\alpha_i+\gamma_j+\varepsilon_{ij}\) with independent \(\varepsilon_{ij}\sim N(0,\sigma^2)\)
(\(\sigma>0\)), find the distributions of the \(y'A_iy/\sigma^2\), and show that \((y'A_2y)/(y'A_4y)\sim F_{1,1}\) when
\(\alpha_1=\alpha_2\). Compute its value for
the data of (b).
\end{enumerate}
\end{exercise}
\begin{solution}
(a) The vectors \(\one,a,b,c\) have squared length \(4\), so each \(A_i\) has the form \(\frac14qq'\) with \(q'q=4\): it is symmetric,
idempotent (\(\frac1{16}q(q'q)q'=\frac14qq'\)) and of rank \(\tr A_i=1\) (Theorem~\ref{r6-thm-idempotent}(c)). The \((k,l)\) entry of
\(4(A_1+A_2+A_3+A_4)\) is \(1+a_ka_l+b_kb_l+c_kc_l\), which is \(4\) for \(k=l\); for \(k\ne l\) it is \(0\): for \((1,2)\), \(1+1-1-1\); for
\((1,3)\), \(1-1+1-1\); for \((1,4)\), \(1-1-1+1\); for \((2,3)\), \(1-1-1+1\); for \((2,4)\), \(1-1+1-1\); for \((3,4)\), \(1+1-1-1\). So the sum is
\(I_4\), and the ranks add up to \(4\). By Theorem~\ref{r15-thm-cochran-algebraic}, \(A_iA_j=0\) for \(i\ne j\) (directly:
\(a'b=a'c=b'c=0\), and \(\one\) is orthogonal to \(a\), \(b\), \(c\)).

(b) \(\one'y=16\), \(a'y=7+3-4-2=4\), \(b'y=7-3+4-2=6\) and \(c'y=7-3-4+2=2\), so \(y'A_1y=\frac{256}4=64\),
\(y'A_2y=\frac{16}4=4\), \(y'A_3y=\frac{36}4=9\) and \(y'A_4y=\frac44=1\). Their sum is \(78=49+9+16+4=y'y\).

(c) \(y\sim N_4(m,\sigma^2I_4)\) with \(m=(\mu+\alpha_1+\gamma_1,\ \mu+\alpha_1+\gamma_2,\ \mu+\alpha_2+\gamma_1,\ \mu+\alpha_2+\gamma_2)'\)
(Theorem~\ref{r14-thm-mvn-iid}(a)). Here \(a'm=2(\alpha_1-\alpha_2)\), \(b'm=2(\gamma_1-\gamma_2)\) and \(c'm=0\). By
Theorem~\ref{r15-thm-cochran}(a), the four forms are independent, and \(y'A_1y/\sigma^2\sim\chi^2_1\bigl((\one'm)^2/(4\sigma^2)\bigr)\),
\(y'A_2y/\sigma^2\sim\chi^2_1\bigl((\alpha_1-\alpha_2)^2/\sigma^2\bigr)\), \(y'A_3y/\sigma^2\sim\chi^2_1\bigl((\gamma_1-\gamma_2)^2/\sigma^2\bigr)\) and
\(y'A_4y/\sigma^2\sim\chi^2_1\). If \(\alpha_1=\alpha_2\), then \(y'A_2y/\sigma^2\sim\chi^2_1\), and
\((y'A_2y)/(y'A_4y)=(y'A_2y/\sigma^2)/(y'A_4y/\sigma^2)\sim F_{1,1}\) (Definition~\ref{r15-def-t-F}). For the data of (b) it is \(4/1=4\).
\end{solution}

\begin{exercise}\label{r3-dr-direct-sum}
Let \(\one=(1,1,1,1)'\), \(x_A=(0,0,1,1)'\), \(x_B=(0,1,0,1)'\), \(x_{AB}=(0,0,0,1)'\) and
\(X=[\one\ x_A\ x_B\ x_{AB}]\).
\begin{enumerate}
\item Row-reduce \([X\mid y]\) for \(y=(y_1,y_2,y_3,y_4)'\), keeping the \(y_i\) as letters. Conclude that
\(\mathcal B=(\one,x_A,x_B,x_{AB})\) is an ordered basis of \(\R^4\), and find \([y]_{\mathcal B}\).
\item Find \([\mu]_{\mathcal B}\) for \(\mu=(2,5,4,9)'\), and check the answer.
\item Show that \(\R^4=\Span\{\one,x_A\}\oplus\Span\{x_B,x_{AB}\}\), and find the parts of \(\mu\) in
the two subspaces.
\end{enumerate}
\end{exercise}
\begin{solution}
(a)
\[
\begin{gathered}
\left(\begin{array}{cccc|c}1&0&0&0&y_1\\1&0&1&0&y_2\\1&1&0&0&y_3\\1&1&1&1&y_4\end{array}\right)
\xrightarrow[R_4\leftarrow R_4-R_1]{R_2\leftarrow R_2-R_1,\ R_3\leftarrow R_3-R_1}
\left(\begin{array}{cccc|c}1&0&0&0&y_1\\0&0&1&0&y_2-y_1\\0&1&0&0&y_3-y_1\\0&1&1&1&y_4-y_1\end{array}\right)
\xrightarrow{R_2\leftrightarrow R_3}
\left(\begin{array}{cccc|c}1&0&0&0&y_1\\0&1&0&0&y_3-y_1\\0&0&1&0&y_2-y_1\\0&1&1&1&y_4-y_1\end{array}\right)
\\
\xrightarrow{R_4\leftarrow R_4-R_2}
\left(\begin{array}{cccc|c}1&0&0&0&y_1\\0&1&0&0&y_3-y_1\\0&0&1&0&y_2-y_1\\0&0&1&1&y_4-y_3\end{array}\right)
\xrightarrow{R_4\leftarrow R_4-R_3}
\left(\begin{array}{cccc|c}1&0&0&0&y_1\\0&1&0&0&y_3-y_1\\0&0&1&0&y_2-y_1\\0&0&0&1&y_1-y_2-y_3+y_4\end{array}\right)
\end{gathered}
\]
After the first step the entry of row \(2\) in column \(2\) is \(0\), so Step 2 interchanges rows \(2\)
and \(3\). The left part has become \(I_4\): \(X\) is invertible (Theorem~\ref{r2-thm-invertible}), its
columns form an ordered basis \(\mathcal B\) of \(\R^4\) (Example~\ref{r3-exm-invertible-basis}), and
\(Xc=y\) has the unique solution \(c=[y]_{\mathcal B}\), read off the last column:
\[
[y]_{\mathcal B}=\begin{pmatrix}y_1\\y_3-y_1\\y_2-y_1\\y_1-y_2-y_3+y_4\end{pmatrix}
\]
(b) For \(\mu=(2,5,4,9)'\): \([\mu]_{\mathcal B}=(2,\ 4-2,\ 5-2,\ 2-5-4+9)'=(2,2,3,2)'\). Check:
\(2\one+2x_A+3x_B+2x_{AB}=(2,2,2,2)'+(0,0,2,2)'+(0,3,0,3)'+(0,0,0,2)'=(2,5,4,9)'\).

(c) The four vectors are independent, so \(\one,x_A\) and \(x_B,x_{AB}\) form bases of
\(U=\Span\{\one,x_A\}\) and \(W=\Span\{x_B,x_{AB}\}\) (Proposition~\ref{r3-prop-dependence}(b)), and the
sum \(U+W\) is direct (Theorem~\ref{r3-thm-direct-sum-basis}). Also
\(U+W=\Span\{\one,x_A,x_B,x_{AB}\}=\R^4\) (Theorem~\ref{r3-thm-sum}(c)). So \(\R^4=U\oplus W\), and
\(\mu=u+w\) with
\[
u=2\one+2x_A=(2,2,4,4)'\in U,\qquad w=3x_B+2x_{AB}=(0,3,0,5)'\in W
\]
which is the only such decomposition (Theorem~\ref{r3-thm-direct-sum}).
\end{solution}

\begin{exercise}[\citeat{searle-khuri}{\S12.12.3, Example~12.10}]\label{r11-ex-unbalanced}
Six plants of three varieties are grown under two treatments: variety \(1\) has two plants under treatment \(1\), with yields
\(3,7\), and one under treatment \(2\), with yield \(8\); variety \(2\) has one plant under each treatment, with yields \(34\)
and \(2\); variety \(3\) has one plant, under treatment \(1\), with yield \(12\). The yield \(y_{ijk}\) of plant \(k\) of
variety \(i\) under treatment \(j\) follows \(y_{ijk}=\mu+\alpha_i+\gamma_j+\varepsilon_{ijk}\). In the order given,
\(y=(y_{111},y_{112},y_{121},y_{211},y_{221},y_{311})'=(3,7,8,34,2,12)'\), and \(y=X\beta+\varepsilon\) with
\(\beta=(\mu,\alpha_1,\alpha_2,\alpha_3,\gamma_1,\gamma_2)'\) under (A1)--(A3).
\begin{enumerate}
\item Write down \(X\), \(X'X\) and \(X'y\), check that
\[
G=\frac17\begin{pmatrix}0&0&0&0&0&0\\0&5&2&4&-4&0\\0&2&5&3&-3&0\\0&4&3&13&-6&0\\0&-4&-3&-6&6&0\\0&0&0&0&0&0\end{pmatrix}
\]
is a g-inverse of \(X'X\), and compute \(\beta^0=GX'y\).
\item Compute \(GX'X\), and decide with Corollary~\ref{r10-cor-estimable-algebra}(iii) which of \(\alpha_1-\alpha_2\),
\(\gamma_1-\gamma_2\) and \(\alpha_1\) are estimable; give the BLUEs of the estimable ones.
\item Compute \(R(\mu)\), \(R(\mu,\alpha)\), \(R(\mu,\gamma)\), \(R(\mu,\alpha,\gamma)\) and \(\mathrm{SSE}\) (Definition~\ref{r11-def-reduction}),
and the two splits of \(y'y\): fitting the varieties before the treatments, and the treatments before the varieties.
\item Assuming the normal linear model, compute the statistic \(F\) of Theorem~\ref{r12-thm-f} for \(H_0\colon\gamma_1=\gamma_2\),
and state when the test of level \(\alpha\) rejects it.
\end{enumerate}
\end{exercise}
\begin{solution}
(a) The rows of \(X\) have a \(1\) in column \(1\), in the column of the variety and in the column of the treatment; the
\((j,l)\) entry of \(X'X\) is \(\col_j(X)'\col_l(X)\) (Theorems~\ref{r1-thm-views}(d) and~\ref{r1-thm-transpose}), and \(X'y\) holds
the total, the variety totals and the treatment totals:
\[
X=\begin{pmatrix}1&1&0&0&1&0\\1&1&0&0&1&0\\1&1&0&0&0&1\\1&0&1&0&1&0\\1&0&1&0&0&1\\1&0&0&1&1&0\end{pmatrix},\qquad
X'X=\begin{pmatrix}6&3&2&1&4&2\\3&3&0&0&2&1\\2&0&2&0&1&1\\1&0&0&1&1&0\\4&2&1&1&4&0\\2&1&1&0&0&2\end{pmatrix},\qquad
X'y=\begin{pmatrix}66\\18\\36\\12\\56\\10\end{pmatrix}
\]
Multiplying,
\[
X'XG=\begin{pmatrix}0&1&1&1&0&0\\0&1&0&0&0&0\\0&0&1&0&0&0\\0&0&0&1&0&0\\0&0&0&0&1&0\\0&1&1&1&-1&0\end{pmatrix}
\]
so, by Theorem~\ref{r1-thm-views}(c), rows \(2\)--\(5\) of \(X'XGX'X\) are rows \(2\)--\(5\) of \(X'X\), row \(1\) is the sum of rows
\(2,3,4\), \((6,3,2,1,4,2)\), and row \(6\) is the sum of rows \(2,3,4\) minus row \(5\), \((6,3,2,1,4,2)-(4,2,1,1,4,0)=(2,1,1,0,0,2)\):
\(X'XGX'X=X'X\). Then
\[
\beta^0=GX'y=\frac17\begin{pmatrix}0\\90+72+48-224\\36+180+36-168\\72+108+156-336\\-72-108-72+336\\0\end{pmatrix}
=\begin{pmatrix}0\\-2\\12\\0\\12\\0\end{pmatrix}
\]
a least-squares estimate (Theorems~\ref{r10-thm-normal-equations}(c) and~\ref{r10-thm-least-squares}(a)).

(b) Multiplying the rows of \(G\) by \(X'X\),
\[
GX'X=\begin{pmatrix}0&0&0&0&0&0\\1&1&0&0&0&1\\1&0&1&0&0&1\\1&0&0&1&0&1\\0&0&0&0&1&-1\\0&0&0&0&0&0\end{pmatrix}
\]
By Corollary~\ref{r10-cor-estimable-algebra}(iii), \(\lambda'\beta\) is estimable exactly when \(\lambda'GX'X=\lambda'\), and
\(\lambda'GX'X\) is the combination of the rows of \(GX'X\) with the entries of \(\lambda\) (Theorem~\ref{r1-thm-views}(c)).
\begin{itemize}
\item \(\alpha_1-\alpha_2\), \(\lambda=(0,1,-1,0,0,0)'\): \(\row_2(GX'X)-\row_3(GX'X)=(0,1,-1,0,0,0)=\lambda'\); estimable, with
BLUE \(\lambda'\beta^0=-2-12=-14\) (Theorem~\ref{r11-thm-blue}).
\item \(\gamma_1-\gamma_2\), \(\lambda=(0,0,0,0,1,-1)'\): \(\row_5(GX'X)-\row_6(GX'X)=(0,0,0,0,1,-1)=\lambda'\); estimable, with
BLUE \(\lambda'\beta^0=12-0=12\).
\item \(\alpha_1\), \(\lambda=(0,1,0,0,0,0)'\): \(\row_2(GX'X)=(1,1,0,0,0,1)\ne\lambda'\); not estimable.
\end{itemize}
(c)
\begin{itemize}
\item \(\bar y=66/6=11\), so \(R(\mu)=n\bar y^2=6\cdot121=726\) (the paragraph after Definition~\ref{r11-def-reduction}).
\item Fitting \(\mu+\alpha_i\) is the one-way layout of the varieties, with group sizes \(3,2,1\) and totals \(18,36,12\); by
Theorem~\ref{r11-thm-oneway-anova}(a), \(R(\mu,\alpha)=\lVert Hy\rVert^2=\sum_ln_l\bar y_l^2=\frac{18^2}3+\frac{36^2}2+\frac{12^2}1=108+648+144=900\).
\item Fitting \(\mu+\gamma_j\) is the one-way layout of the treatments, with design \(Z\) (columns \(1,5,6\) of \(X\)), but
Theorem~\ref{r11-thm-oneway-anova}(a) lists the observations group by group. Relabel the plants so that the four under
treatment \(1\) come first: this replaces \(Z\) and \(y\) by \(PZ\) and \(Py\) for a permutation matrix \(P\)
(Definition~\ref{r5d-def-permutation-matrix}), and \(P'P=I_6\) (Theorem~\ref{r6-thm-orthogonal}(d)), so
\((PZ)'PZ=Z'Z\) and \((PZ)'Py=Z'y\). So the normal equations \(Z'Zb=Z'y\), their solutions \(\hat b\) and the reduction
\(R(\mu,\gamma)=\hat b'Z'y\) (the paragraph after Definition~\ref{r11-def-reduction}) do not change. By
Theorem~\ref{r11-thm-oneway-anova}(a) for the relabelled plants, with group sizes \(4,2\) and totals \(56,10\),
\(R(\mu,\gamma)=\frac{56^2}4+\frac{10^2}2=784+50=834\).
\item \(R(\mu,\alpha,\gamma)=\beta^{0\prime}X'y=-2\cdot18+12\cdot36+12\cdot56=-36+432+672=1068\) (Theorem~\ref{r10-thm-PX}(d)), and
\(y'y=9+49+64+1156+4+144=1426\), so \(\mathrm{SSE}=1426-1068=358\).
\item Varieties first: \(R(\alpha\mid\mu)=900-726=174\), \(R(\gamma\mid\mu,\alpha)=1068-900=168\), and
\(1426=726+174+168+358\). Treatments first: \(R(\gamma\mid\mu)=834-726=108\), \(R(\alpha\mid\mu,\gamma)=1068-834=234\), and
\(1426=726+108+234+358\). The two splits differ because the numbers of plants in the cells are unequal.
\end{itemize}
(d)
\begin{itemize}
\item \(H_0\) is \(K\beta=0\) with \(K=(0,0,0,0,1,-1)\), estimable by (b) and of rank \(1\): testable
(Definition~\ref{r11-def-testable}).
\item The \(\beta\) with \(\gamma_1=\gamma_2=\gamma\) give
\(X\beta=(\mu+\gamma)\one_6+\alpha_1d_1+\alpha_2d_2+\alpha_3d_3\), where \(d_i\) is the indicator of variety \(i\) (the two
treatment indicators add up to \(\one_6\)); these are all vectors of \(\mathcal C([\one_6\ d_1\ d_2\ d_3])\). So, by
Theorem~\ref{r12-thm-cls}(b), \(\mathrm{SSE}_K\) is the error sum of squares of the varieties alone, \(1426-900=526\), and
\(\mathrm{SSE}_K-\mathrm{SSE}=R(\gamma\mid\mu,\alpha)=168\).
\item \(r=\rank X=4\). Column \(1\) of \(X\) is the sum of columns \(2,3,4\) and also the sum of columns \(5,6\), so
\((-1,1,1,1,0,0)'\) and \((-1,0,0,0,1,1)'\) lie in \(\mathcal N(X)\) (Theorem~\ref{r1-thm-views}(b)); they are independent
(look at entries \(2\) and \(5\)), so \(\dim\mathcal N(X)\geqslant2\) (Theorem~\ref{thm:r4-counting}(c)) and \(r\leqslant4\)
(Theorem~\ref{thm:r4-rank-nullity}). Rows \(2\)--\(5\) of \(GX'X\) in (b) lie in \(\mathcal R(GX'X)\) and are independent
(their entries \(2\)--\(5\) form \(I_4\)), so \(\rank(GX'X)\geqslant4\) (Theorem~\ref{thm:r4-counting}(c),
Definition~\ref{def:r4-rank}); and \(\rank(GX'X)\leqslant\rank X\) (Theorem~\ref{thm:r4-rank-product}(a), as
\(GX'X=(GX')X\)). So \(r=4\).
\item Theorem~\ref{r12-thm-f}(a) gives
\[
F=\frac{168/1}{358/2}=\frac{168}{179}
\]
Check: \(K\beta^0=12\) and \(V=KGK'=\frac67\), so \((K\beta^0)^2/V=144\cdot\frac76=168\). The test of level \(\alpha\) rejects
\(H_0\) if and only if \(F_{1,2,1-\alpha}<\frac{168}{179}\) (Theorem~\ref{r12-thm-f}(b)).
\end{itemize}
\end{solution}

\subsection*{Problems}

\begin{problem}[\citeat{searle-khuri}{Theorems~12.4, 12.5}]\label{r11-ex-estimable-ginverse}
Assume (A1), and let \(G\) be a g-inverse of \(X'X\). Show that for \(a\in\R^{k+1}\) the following are equivalent:
(a) \(a'\beta\) is estimable; (b) \(a'=a'GX'X\); (c) \(a'=w'GX'X\) for some \(w\in\R^{k+1}\).
\end{problem}
\begin{solution}
\begin{itemize}
\item (a)\(\Rightarrow\)(b): By Theorem~\ref{r11-thm-estimable}(a), \(a'=c'X\) for some \(c\in\R^n\). By
Theorem~\ref{r10-thm-PX}(a), \(XGX'X=HX=X\), so \(a'=c'XGX'X=a'GX'X\).
\item (b)\(\Rightarrow\)(c): take \(w=a\).
\item (c)\(\Rightarrow\)(a): \(a'=w'GX'X=(w'GX')X\), so \(a=X'(XG'w)\in\mathcal C(X')\), and \(a'\beta\) is estimable
(Theorem~\ref{r11-thm-estimable}(a)).
\end{itemize}
\end{solution}

\begin{problem}[\citeat{tu-clm}{Corollary~11.3}]\label{r11-pr-two-sample}
Let \(x_1,\dots,x_m\) be \(N(\mu_1,\sigma^2)\) and \(y_1,\dots,y_n\) be \(N(\mu_2,\sigma^2)\), all \(m+n\) independent,
with \(m,n\geqslant2\) and \(\sigma>0\), and let
\(s_p^2=\bigl(\sum_i(x_i-\bar x)^2+\sum_j(y_j-\bar y)^2\bigr)/(m+n-2)\).
\begin{enumerate}
\item Show that \(T=\dfrac{\bar x-\bar y-(\mu_1-\mu_2)}{s_p\sqrt{1/m+1/n}}\sim t_{m+n-2}\).
\item Deduce a confidence interval with level \(1-\alpha\) for \(\mu_1-\mu_2\).
\end{enumerate}
\end{problem}
\begin{solution}
(a)
\begin{itemize}
\item By Theorem~\ref{r14-thm-mvn-iid}(a), \(v=(x_1,\dots,x_m,y_1,\dots,y_n)'\) is normal with mean
\((\mu_1\one_m',\mu_2\one_n')'\) and covariance matrix \(\sigma^2I_{m+n}\); so \(\varepsilon=v-\expc v\sim N_{m+n}(0,\sigma^2I_{m+n})\)
(Theorem~\ref{r14-thm-mvn-linear}).
\item \(\bar x-\bar y-(\mu_1-\mu_2)=a'\varepsilon\) with \(a=(\frac1m\one_m',-\frac1n\one_n')'\), \(a'a=\frac1m+\frac1n\).
\item Let \(x=(x_1,\dots,x_m)'\) and \(y=(y_1,\dots,y_n)'\), so that \(\varepsilon=(\varepsilon_1',\varepsilon_2')'\) with
\(\varepsilon_1=x-\mu_1\one_m\) and \(\varepsilon_2=y-\mu_2\one_n\). By Example~\ref{r1-exm-centering},
\(\sum_{i=1}^m(x_i-\bar x)^2=(C_mx)'(C_mx)\); as \(C_m\one_m=0\) and \(C_m'C_m=C_m\) (Proposition~\ref{r1-prop-centering}),
\(C_mx=C_m\varepsilon_1\) and \(\sum_{i=1}^m(x_i-\bar x)^2=\varepsilon_1'C_m\varepsilon_1\). Likewise
\(\sum_{j=1}^n(y_j-\bar y)^2=\varepsilon_2'C_n\varepsilon_2\). So \((m+n-2)s_p^2=\varepsilon'P\varepsilon\) with
\(P=\diag(C_m,C_n)\) (Theorem~\ref{r5b-thm-block}(c)).
\item \(P\) is symmetric and idempotent (Proposition~\ref{r5b-prop-block-tri}(b)), of rank
\((m-1)+(n-1)\) (Theorem~\ref{thm:r4-partitioned}(c), Corollary~\ref{r6-cor-centering}(b)).
\item \(a'P=(\frac1m\one_m'C_m,\ -\frac1n\one_n'C_n)=0\), so \(a'\varepsilon\) and \(\varepsilon'P\varepsilon\) are
independent (Theorem~\ref{r15-thm-indep}(b)). Moreover \(u=a'\varepsilon/(\sigma\sqrt{a'a})\sim N(0,1)\) and
\(w=\varepsilon'P\varepsilon/\sigma^2\sim\chi^2_{m+n-2}\) (Theorems~\ref{r14-thm-mvn-linear} and~\ref{r15-thm-quadratic-chisq}), and \(T=u/\sqrt{w/(m+n-2)}\sim t_{m+n-2}\) by (N3) and
Definition~\ref{r15-def-t-F}(a).
\end{itemize}

(b) With \(c=t_{m+n-2,1-\alpha/2}\), Theorem~\ref{r15-thm-t-symmetric}(a) gives
\(\prb(-c\leqslant T\leqslant c)=1-\alpha\) for all values of the parameters; as in Example~\ref{r15-exm-one-sample},
\(-c\leqslant T\leqslant c\) holds exactly when \(\mu_1-\mu_2\) lies in the interval below, so \(\bar x-\bar y\pm c\,s_p\sqrt{1/m+1/n}\) is a confidence interval
with level \(1-\alpha\) for \(\mu_1-\mu_2\) (Definition~\ref{r15-def-confidence}(a)).
\end{solution}

\clearpage
\def\unitnumber{15}
\setcounter{theorem}{0}\setcounter{problem}{0}
\section*{Chapter 15. Extremal principles and the singular value decomposition}

\subsection*{The generalized Rayleigh quotient}

\needspace{12\baselineskip}
\begin{theorem}[\citeat{bapat}{Result~5.2}, \citeat{searle-khuri}{Corollary~9.3}]\label{r9-thm-gen-rayleigh}
Let \(A,B\in\R^{n\times n}\) be symmetric, with \(B\) PD, and let \(\mu_1\geqslant\dots\geqslant\mu_n\) be the
eigenvalues of the symmetric matrix \(C=B^{-1/2}AB^{-1/2}\) (Theorem~\ref{r8-thm-sqrt}). Then for every \(x\ne0\)
\[
\mu_n\leqslant\frac{x'Ax}{x'Bx}\leqslant\mu_1
\]
The maximum \(\mu_1\) is attained at \(x=B^{-1/2}y\) for every eigenvector \(y\) of \(C\) belonging to \(\mu_1\), and the
minimum \(\mu_n\) likewise. The numbers \(\mu_1,\dots,\mu_n\) are also the eigenvalues of \(B^{-1}A\) and of \(AB^{-1}\),
listed with multiplicity.
\end{theorem}
\begin{proof}
\begin{itemize}
\item \(B^{1/2}\) is symmetric and invertible, and so is \(B^{-1/2}=(B^{1/2})^{-1}\) (Theorems~\ref{r8-thm-sqrt}
and~\ref{r2-thm-inverse-props}(c)); hence \(C'=C\).
\item Put \(y=B^{1/2}x\). Then \(x=B^{-1/2}y\), \(x\ne0\) if and only if \(y\ne0\), \(x'Bx=x'B^{1/2}B^{1/2}x=y'y\) and
\(x'Ax=y'B^{-1/2}AB^{-1/2}y=y'Cy\); so
\[
\frac{x'Ax}{x'Bx}=\frac{y'Cy}{y'y}
\]
\item Theorem~\ref{r9-thm-rayleigh}(a),(b) for \(C\) gives the bounds, attained at the eigenvectors \(y\) of \(C\) belonging
to \(\mu_1\) and to \(\mu_n\), that is, at \(x=B^{-1/2}y\).
\item \(B^{-1/2}B^{-1/2}=(B^{1/2}B^{1/2})^{-1}=B^{-1}\) (Theorem~\ref{r2-thm-inverse-props}(b)). So
\(C=B^{-1/2}(AB^{-1/2})\) with \((AB^{-1/2})B^{-1/2}=AB^{-1}\), and \(B^{-1}A=B^{-1/2}(B^{-1/2}A)\) with
\((B^{-1/2}A)B^{-1/2}=C\). By Theorem~\ref{r7-thm-AB-BA}(a) (\(m=n\)), \(p_{AB^{-1}}=p_C=p_{B^{-1}A}\): both matrices
have the eigenvalues of \(C\) with the same multiplicities.\qedhere
\end{itemize}
\end{proof}

\begin{theorem}[\citeat{bapat}{Result~5.3}]\label{r9-thm-cs-max}
Let \(B\in\R^{n\times n}\) be PD and \(y\in\R^n\). Then
\[
(x'y)^2\leqslant(x'Bx)(y'B^{-1}y)\qquad(x\in\R^n)
\]
If \(y\ne0\), then \(\max_{x\ne0}(x'y)^2/x'Bx=y'B^{-1}y\), attained at \(x=B^{-1}y\).
\end{theorem}
\begin{proof}
If \(y=0\) or \(x=0\), both sides of the inequality are \(0\). Let \(y\ne0\).
\begin{itemize}
\item \((x'y)^2=x'yy'x\), so \((x'y)^2/x'Bx\) is \(x'Ax/x'Bx\) with the symmetric matrix \(A=yy'\). By
Theorem~\ref{r9-thm-gen-rayleigh}, its maximum over \(x\ne0\) is the largest eigenvalue of
\(C=B^{-1/2}yy'B^{-1/2}=ww'\), where \(w=B^{-1/2}y\ne0\). Here \(B^{-1/2}=(B^{1/2})^{-1}\) is symmetric and
\(B^{-1/2}B^{-1/2}=(B^{1/2}B^{1/2})^{-1}=B^{-1}\) (Theorems~\ref{r8-thm-sqrt} and~\ref{r2-thm-inverse-props}(b),(c)), so
\(w'w=y'B^{-1}y\) and \(B^{-1/2}w=B^{-1}y\).
\item \(\rank(ww')\leqslant\rank w=1\) (Theorem~\ref{thm:r4-rank-product}(a)) and \(Cw=w(w'w)\ne0\), so \(C\) has rank
\(1\). By Theorem~\ref{r8-thm-rank-eigen}, exactly one eigenvalue of \(C\), counted with multiplicity, is nonzero, and
it equals \(\tr C=\tr(w'w)=w'w=y'B^{-1}y>0\) (Theorems~\ref{r7-thm-trace-det} and~\ref{r1-thm-trace}(b)). So the
largest eigenvalue of \(C\) is \(y'B^{-1}y\); this gives the inequality for \(x\ne0\).
\item \(Cw=(w'w)w\): \(w\) is an eigenvector of \(C\) belonging to \(y'B^{-1}y\), and the maximum is attained at
\(x=B^{-1/2}w=B^{-1}y\) (Theorem~\ref{r9-thm-gen-rayleigh}).\qedhere
\end{itemize}
\end{proof}

In statistics, let the rows \(x_1',\dots,x_n'\) of a data matrix \(X\in\R^{n\times p}\) fall into \(g\) groups:
group \(l\) has the set \(\mathcal I_l\) of row indices, \(n_l\) rows and the mean vector \(\bar x_l\), and \(\bar x\) is the
mean vector of all \(n\) rows. For \(a\in\R^p\), the values \(Xa=(a'x_1,\dots,a'x_n)'\) of the combination \(a'x\) have the
group means \(a'\bar x_l\) and the grand mean \(a'\bar x\). By Theorem~\ref{r11-thm-oneway-anova}(b) for \(y=Xa\) (after
the rows are arranged by group; the sums of squares do not depend on the order of the rows), their total sum of squares
is the sum of the between-group and the within-group sums of squares
\[
\sum_{l=1}^gn_l(a'\bar x_l-a'\bar x)^2=a'Ba,\qquad\sum_{l=1}^g\sum_{i\in\mathcal I_l}(a'x_i-a'\bar x_l)^2=a'Wa
\]
where, as \(\bigl(a'(x_i-\bar x_l)\bigr)^2=a'(x_i-\bar x_l)(x_i-\bar x_l)'a\),
\[
B=\sum_{l=1}^gn_l(\bar x_l-\bar x)(\bar x_l-\bar x)',\qquad
W=\sum_{l=1}^g\sum_{i\in\mathcal I_l}(x_i-\bar x_l)(x_i-\bar x_l)'
\]
\textbf{Fisher's linear discriminant function} is a combination \(a'x\) that maximizes \(a'Ba/a'Wa\)~\citeat{haerdle-simar}{\S14.2, Theorem~14.4}. Let \(W\) be PD.
By Theorem~\ref{r9-thm-gen-rayleigh} for the symmetric \(B\) and the PD \(W\) (in the roles of \(A\) and \(B\) there), the
maximum is the largest eigenvalue \(\mu_1\) of \(W^{-1}B\), attained at \(a=W^{-1/2}y\) for an eigenvector \(y\) of
\(C=W^{-1/2}BW^{-1/2}\) belonging to \(\mu_1\); this \(a\) is an eigenvector of \(W^{-1}B\), as
\(W^{-1}Ba=W^{-1/2}(W^{-1/2}BW^{-1/2})y=\mu_1a\).

For two groups, let \(d=\bar x_1-\bar x_2\). As \(\bar x=(n_1\bar x_1+n_2\bar x_2)/n\) and \(n=n_1+n_2\),
\[
\bar x_1-\bar x=\frac{n_2}n\,d,\qquad\bar x_2-\bar x=-\frac{n_1}n\,d,\qquad
B=\frac{n_1n_2^2+n_2n_1^2}{n^2}\,dd'=\frac{n_1n_2}n\,dd'
\]
so that
\[
\frac{a'Ba}{a'Wa}=\frac{n_1n_2}n\cdot\frac{(a'd)^2}{a'Wa}
\]
If \(d\ne0\), Theorem~\ref{r9-thm-cs-max} gives the maximum \((n_1n_2/n)\,d'W^{-1}d\), attained at \(a=W^{-1}d\). With the
\textbf{pooled sample covariance matrix} \(S_p=W/(n-2)\), the ratio \((a'd)^2/a'S_pa\) is the squared difference of the two
group means of \(a'x\) divided by its pooled sample variance, and it is largest for \(a=S_p^{-1}d=(n-2)W^{-1}d\).

\begin{proposition}[\citeat{bapat}{\S8.5}]\label{r9-prop-multiple-correlation}
Let \(y\) be a random variable and \(x\in\R^p\) a random vector with
\[
\var\begin{pmatrix}y\\x\end{pmatrix}=V=\begin{pmatrix}\sigma^2&u'\\u&\Sigma\end{pmatrix}
\]
where \(\sigma^2>0\), \(u\in\R^p\) and \(\Sigma\in\R^{p\times p}\) is PD. For \(\alpha\ne0\) in \(\R^p\), let \(\rho(\alpha)\) be
the correlation of \(y\) and \(\alpha'x\) (Definition~\ref{r14-def-corr}).
\begin{enumerate}
\item \(\rho(\alpha)^2\leqslant u'\Sigma^{-1}u/\sigma^2\) for every \(\alpha\ne0\), with equality at \(\alpha=\Sigma^{-1}u\) if
\(u\ne0\) (and at every \(\alpha\) if \(u=0\)). The number \(r_{y(x)}\geqslant0\) with
\[
r_{y(x)}^2=\max_{\alpha\ne0}\rho(\alpha)^2=\frac{u'\Sigma^{-1}u}{\sigma^2}
\]
is the \textbf{multiple correlation coefficient} of \(y\) and \(x_1,\dots,x_p\).
\item If \(V\) is PD and \(Z=V^{-1}=(z_{ij})\), then \(r_{y(x)}^2=1-1/(\sigma^2z_{11})\).
\end{enumerate}
\end{proposition}
\begin{proof}
(a)
\begin{itemize}
\item The matrix \(T=\left(\begin{smallmatrix}1&0'\\0&\alpha'\end{smallmatrix}\right)\in\R^{2\times(p+1)}\) maps \((y,x')'\) to
\((y,\alpha'x)'\), so by Theorems~\ref{r11-thm-linear}(b) and~\ref{r5b-thm-block}
\[
\var\begin{pmatrix}y\\\alpha'x\end{pmatrix}=TVT'=\begin{pmatrix}\sigma^2&u'\alpha\\\alpha'u&\alpha'\Sigma\alpha\end{pmatrix}
\]
Here \(\var(\alpha'x)=\alpha'\Sigma\alpha>0\), as \(\Sigma\) is PD and \(\alpha\ne0\); so
\(\rho(\alpha)^2=(\alpha'u)^2/(\sigma^2\,\alpha'\Sigma\alpha)\).
\item By Theorem~\ref{r9-thm-cs-max} for the PD \(\Sigma\) and the vector \(u\),
\((\alpha'u)^2\leqslant(\alpha'\Sigma\alpha)(u'\Sigma^{-1}u)\); divide by \(\sigma^2\,\alpha'\Sigma\alpha>0\). If \(u\ne0\), the
same theorem gives equality at \(\alpha=\Sigma^{-1}u\ne0\); if \(u=0\), both sides are \(0\).
\end{itemize}
(b)
\begin{itemize}
\item \(\Sigma\) is invertible, so by Theorem~\ref{r5d-thm-schur-det}(b) (the Schur complement of the lower right block),
\(\det V=\det\Sigma\cdot(\sigma^2-u'\Sigma^{-1}u)\).
\item \(V(1|1)=\Sigma\) (Definition~\ref{r5d-def-cofactor}), so the \((1,1)\) entry of \(\operatorname{adj}V\) is
\(\det\Sigma\), and \(z_{11}=\det\Sigma/\det V\) (Theorem~\ref{r5d-thm-adjugate}). Both determinants are positive, as
\(V\) and \(\Sigma\) are PD (Theorem~\ref{r8-thm-eigen-criterion}(c)).
\item Hence \(1/z_{11}=\det V/\det\Sigma=\sigma^2-u'\Sigma^{-1}u\), and by (a)
\(r_{y(x)}^2=(\sigma^2-1/z_{11})/\sigma^2=1-1/(\sigma^2z_{11})\).\qedhere
\end{itemize}
\end{proof}
If \((x',y)'\) is normal with mean \((\mu',\tau)'\) and the PD covariance matrix
\(\left(\begin{smallmatrix}\Sigma&u\\u'&\sigma^2\end{smallmatrix}\right)\) (the entries of \(V\) rearranged), then by
Theorem~\ref{r14-thm-mvn-conditional} with \(z_1=x\) and \(z_2=y\), the conditional distribution of \(y\) given \(x\) is
\(N\bigl(\tau+u'\Sigma^{-1}(x-\mu),\ \sigma^2-u'\Sigma^{-1}u\bigr)\). Its variance is \(\sigma^2(1-r_{y(x)}^2)\), and its mean
contains \(u'\Sigma^{-1}x\), the combination \(\alpha'x\) of maximum correlation in (a); this is also the best linear
predictor of Theorem~\ref{r11-thm-blp} up to the constant term (\(b^*=\Sigma^{-1}u\)). So \(r_{y(x)}^2\) is the population
counterpart of \(R^2=\cor(y,\hat y)^2\) (Proposition~\ref{r11-prop-r2-corr}).

\subsection*{The minimax principle}

\begin{lemma}[\citeat{bapat}{Ch.~1, Exercise~19}]\label{r9-lem-intersect}
If \(U\) and \(W\) are subspaces of \(\R^n\) with \(\dim U+\dim W>n\), then \(U\cap W\) contains a unit vector.
\end{lemma}
\begin{proof}
If \(U\cap W=\{0\}\), then a basis of \(U\) together with a basis of \(W\) is linearly independent
(Theorem~\ref{r3-thm-direct-sum-basis}); these are \(\dim U+\dim W>n\) vectors of \(\R^n\) \(\contradiction\)
(Theorem~\ref{thm:r4-pivot-test}). So there is \(x\ne0\) in the subspace \(U\cap W\) (Theorem~\ref{r3-thm-sum}(a)), and
\(x/\lVert x\rVert\) is a unit vector in it.
\end{proof}

\needspace{20\baselineskip}
\begin{theorem}[Courant--Fischer; \citeat{bapat}{Result~5.8}, \citeat{abadir-magnus}{Exercises~12.43, 12.44}]\label{r9-thm-courant-fischer}
Let \(A\in\R^{n\times n}\) be symmetric, with \(Q=[q_1\ \cdots\ q_n]\) as in Lemma~\ref{r9-lem-eig-order}.
\begin{enumerate}
\item (Min--max form.) Let \(2\leqslant k\leqslant n\). For \(w_1,\dots,w_{k-1}\in\R^n\), the vectors \(u\) with
\(u\perp w_1,\dots,w_{k-1}\) form a subspace \(W\) of dimension at least \(n-k+1\), and
\[
m(w_1,\dots,w_{k-1})=\max\{u'Au:u\in W,\ \lVert u\rVert=1\}\geqslant\lambda_k
\]
with equality for \(w_i=q_i\) (\(i=1,\dots,k-1\)). That is,
\[
\lambda_k=\min_{w_1,\dots,w_{k-1}\in\R^n}\max\{u'Au:u\perp w_1,\dots,w_{k-1},\ \lVert u\rVert=1\}
\]
\item (Max--min form.) Let \(1\leqslant k\leqslant n-1\). For \(w_1,\dots,w_{n-k}\in\R^n\), the minimum of \(u'Au\) over
the unit vectors \(u\perp w_1,\dots,w_{n-k}\) exists and is \(\leqslant\lambda_k\), with equality for \(w_i=q_{k+i}\)
(\(i=1,\dots,n-k\)). That is,
\[
\lambda_k=\max_{w_1,\dots,w_{n-k}\in\R^n}\min\{u'Au:u\perp w_1,\dots,w_{n-k},\ \lVert u\rVert=1\}
\]
\end{enumerate}
\end{theorem}
\begin{proof}
(a)
\begin{itemize}
\item \(W=\mathcal N(M')\) with \(M=[w_1\ \cdots\ w_{k-1}]\in\R^{n\times(k-1)}\), since the entries of \(M'u\) are
\(w_i'u\). As \(\rank M'\leqslant k-1\) (Theorem~\ref{thm:r4-rank-rules}(a)),
\(\dim W=n-\rank M'\geqslant n-k+1\geqslant1\) (Theorem~\ref{thm:r4-rank-nullity}); so the maximum exists
(Lemma~\ref{r9-lem-compression}).
\item For \(w_i=q_i\) the maximum is \(\lambda_k\) (Theorem~\ref{r9-thm-constrained}).
\item \(T=\Span\{q_1,\dots,q_k\}\) has dimension \(k\) (orthonormal vectors are independent,
Proposition~\ref{r6-prop-orthogonal-list}(b)), and \(\dim W+\dim T\geqslant n+1\); so there is a unit vector
\(u\in W\cap T\) (Lemma~\ref{r9-lem-intersect}). Write \(u=c_1q_1+\dots+c_kq_k\). As \(q_i'q_j\) is \(1\) for \(i=j\) and
\(0\) otherwise, \(c_1^2+\dots+c_k^2=u'u=1\) and
\[
u'Au=\Bigl(\sum_{i\leqslant k}c_iq_i\Bigr)'\sum_{j\leqslant k}c_j\lambda_jq_j=\sum_{i\leqslant k}\lambda_ic_i^2
\geqslant\lambda_k\sum_{i\leqslant k}c_i^2=\lambda_k
\]
So \(m(w_1,\dots,w_{k-1})\geqslant u'Au\geqslant\lambda_k\).
\end{itemize}
\needspace{3\baselineskip}
(b)
\begin{itemize}
\item \(-A\) is symmetric, and \(\lambda_j(-A)=-\lambda_{n+1-j}(A)\) (Lemma~\ref{r9-lem-eig-order}(b)). Put
\(j=n+1-k\); then \(2\leqslant j\leqslant n\), \(j-1=n-k\) and \(\lambda_j(-A)=-\lambda_k\).
\item Let \(w_1,\dots,w_{n-k}\in\R^n\). By (a) for \(-A\) and \(j\), the maximum of \(u'(-A)u=-u'Au\) over the unit
vectors \(u\perp w_1,\dots,w_{n-k}\) exists and is \(\geqslant\lambda_j(-A)=-\lambda_k\). So the minimum of \(u'Au\) over
these \(u\) exists, it is minus that maximum, and it is \(\leqslant\lambda_k\).
\item For \(w_i=q_{k+i}\) (\(i=1,\dots,n-k\)) the minimum is \(\lambda_k\) (Theorem~\ref{r9-thm-constrained}, second
form).\qedhere
\end{itemize}
\end{proof}

\begin{theorem}[Interlacing; \citeat{bapat}{Result~5.9}]\label{r9-thm-interlacing}
Let \(A\in\R^{n\times n}\) be symmetric (\(n\geqslant2\)), let \(1\leqslant i\leqslant n\), and let
\(B=A(i|i)\in\R^{(n-1)\times(n-1)}\) be \(A\) with row \(i\) and column \(i\) deleted
(Definition~\ref{r5d-def-cofactor}). Then \(B\) is symmetric and
\[
\lambda_1(A)\geqslant\lambda_1(B)\geqslant\lambda_2(A)\geqslant\lambda_2(B)\geqslant\dots\geqslant\lambda_{n-1}(A)
\geqslant\lambda_{n-1}(B)\geqslant\lambda_n(A)
\]
that is, \(\lambda_k(A)\geqslant\lambda_k(B)\geqslant\lambda_{k+1}(A)\) for \(k=1,\dots,n-1\).
\end{theorem}
\begin{proof}
Let \(S\in\R^{n\times(n-1)}\) be \(I_n\) with column \(i\) deleted. Then \(S'S=I_{n-1}\) and \(S'AS=B\): entry \((s,t)\)
of \(S'AS\) is \(a_{jl}\), where \(j\) and \(l\) are the \(s\)th and the \(t\)th of the indices \(\ne i\). So \(B\) is
symmetric. For \(\tilde u\in\R^{n-1}\), \(u=S\tilde u\) has \(u_i=0\), \(\lVert u\rVert=\lVert\tilde u\rVert\) and
\(u'Au=\tilde u'B\tilde u\).
\begin{itemize}
\item Left inequality. Let \(q_1,\dots,q_n\) be as in Lemma~\ref{r9-lem-eig-order}. By
Theorem~\ref{r9-thm-constrained}, \(u'Au\leqslant\lambda_k(A)\) for every unit vector \(u\perp q_1,\dots,q_{k-1}\). For
\(u=S\tilde u\), \(u'q_j=\tilde u'(S'q_j)\); so \(\tilde u'B\tilde u\leqslant\lambda_k(A)\) for every unit
\(\tilde u\in\R^{n-1}\) with \(\tilde u\perp S'q_1,\dots,S'q_{k-1}\). The largest such \(\tilde u'B\tilde u\) is
\(\geqslant\lambda_k(B)\): for \(k=1\) it is \(\lambda_1(B)\) (Theorem~\ref{r9-thm-rayleigh}), and for \(k\geqslant2\)
apply Theorem~\ref{r9-thm-courant-fischer}(a) to \(B\) with \(w_j=S'q_j\). Hence \(\lambda_k(A)\geqslant\lambda_k(B)\).
\item Right inequality. \(-B\) is \(-A\) with row and column \(i\) deleted. The left inequality for \(-A\) at the index
\(n-k\in\{1,\dots,n-1\}\) gives \(\lambda_{n-k}(-A)\geqslant\lambda_{n-k}(-B)\), that is,
\(-\lambda_{k+1}(A)\geqslant-\lambda_k(B)\) (Lemma~\ref{r9-lem-eig-order}(b), for \(A\) of order \(n\) and \(B\) of order
\(n-1\)).\qedhere
\end{itemize}
\end{proof}
For a data matrix \(X_0\) with \(n\) rows, the entries of \(X_0'C_nX_0\) are \(\col_j(X_0)'C_n\col_l(X_0)\)
(Example~\ref{r1-exm-data});
deleting variable \(i\) deletes row \(i\) and column \(i\) of this matrix, so its eigenvalues interlace those of the
full matrix.

\begin{theorem}[\citeat{bapat}{Result~5.10}]\label{r9-thm-monotone}
Let \(A,B\in\R^{n\times n}\) be symmetric with \(A-B\) PSD. Then \(\lambda_k(A)\geqslant\lambda_k(B)\) for
\(k=1,\dots,n\).
\end{theorem}
\begin{proof}
Let \(u_1,\dots,u_n\) and \(v_1,\dots,v_n\) be the columns of orthogonal matrices as in Lemma~\ref{r9-lem-eig-order} for
\(A\) and for \(B\), and fix \(k\).
\begin{itemize}
\item \(T_1=\Span\{u_k,\dots,u_n\}\) and \(T_2=\Span\{v_1,\dots,v_k\}\) have dimensions \(n-k+1\) and \(k\) (orthonormal
vectors are independent, Proposition~\ref{r6-prop-orthogonal-list}(b)). By Lemma~\ref{r9-lem-intersect} there is a unit
vector \(z\in T_1\cap T_2\).
\item \(z\in T_1\), so \(z=\sum_{i\geqslant k}c_iu_i\), with \(c_i=u_i'z\), \(u_i'z=0\) for \(i<k\), and
\(\sum_{i\geqslant k}c_i^2=z'z=1\). By Lemma~\ref{r9-lem-eig-order}(a),
\[
z'Az=\sum_{i=1}^n\lambda_i(A)(u_i'z)^2=\sum_{i\geqslant k}\lambda_i(A)c_i^2\leqslant\lambda_k(A)\sum_{i\geqslant k}c_i^2
=\lambda_k(A)
\]
\item Likewise \(z=\sum_{j\leqslant k}d_jv_j\) with \(\sum_{j\leqslant k}d_j^2=1\), and
\(z'Bz=\sum_{j\leqslant k}\lambda_j(B)d_j^2\geqslant\lambda_k(B)\).
\item \(z'(A-B)z\geqslant0\), so \(\lambda_k(A)\geqslant z'Az\geqslant z'Bz\geqslant\lambda_k(B)\).\qedhere
\end{itemize}
\end{proof}
If \(X\in\R^{n\times p}\) has the rows \(x_1',\dots,x_n'\), then \(X'X=x_1x_1'+\dots+x_nx_n'\) (outer product expansion,
Theorem~\ref{r1-thm-views}(d)). So adding an observation \(x'\) as a new row turns \(X'X\) into \(X'X+xx'\) (as in
Example~\ref{r11-exm-added-observation}), and \(xx'\) is
PSD (\(v'xx'v=(x'v)^2\geqslant0\)). By Theorem~\ref{r9-thm-monotone}, no eigenvalue of \(X'X\) decreases when an
observation is added.

\subsection*{The singular value decomposition}

\begin{proposition}[\citeat{searle-khuri}{Definition~9.1}, \citeat{harville}{Lemmas~18.2.21, 8.4.2}]\label{r9-prop-frobenius}
Let \(A\in\R^{m\times n}\), with the Frobenius norm \(\lVert A\rVert_F\) of Definition~\ref{r9-def-frobenius}.
\begin{enumerate}
\item \(\lVert A\rVert_F\geqslant0\), with equality only for \(A=0\); \(\lVert cA\rVert_F=\lvert c\rvert\,\lVert A\rVert_F\)
for \(c\in\R\); and \(\lVert A+B\rVert_F\leqslant\lVert A\rVert_F+\lVert B\rVert_F\) for \(B\in\R^{m\times n}\).
\item \(\lVert AB\rVert_F\leqslant\lVert A\rVert_F\lVert B\rVert_F\) for \(B\in\R^{n\times p}\).
\item \(\lVert PAQ\rVert_F=\lVert A\rVert_F\) for orthogonal \(P\in\R^{m\times m}\) and \(Q\in\R^{n\times n}\).
\end{enumerate}
\end{proposition}
\begin{proof}
(a) The entries of \(A+B\) and \(cA\) are \(a_{ij}+b_{ij}\) and \(ca_{ij}\) (Definition~\ref{r1-def-matrix}), so if \(v\),
\(w\) list the entries of \(A\), \(B\), then \(v+w\) and \(cv\) list those of \(A+B\) and \(cA\). Apply the remarks after
Definition~\ref{r6-def-inner} and Theorem~\ref{r6-thm-cauchy-schwarz}(b) in \(\R^{mn}\).

\needspace{4\baselineskip}
(b)
\begin{itemize}
\item Column \(j\) of \(AB\) is \(Ab_j\), where \(b_j=\col_j(B)\) (Theorem~\ref{r1-thm-views}(a)), and entry \(i\) of
\(Ab_j\) is \(\row_i(A)b_j=\sum_la_{il}b_{lj}\).
\item By Cauchy--Schwarz (Theorem~\ref{r6-thm-cauchy-schwarz}(a)),
\(\lVert Ab_j\rVert^2=\sum_i(\row_i(A)b_j)^2\leqslant\sum_i\bigl(\sum_la_{il}^2\bigr)\lVert b_j\rVert^2
=\lVert A\rVert_F^2\lVert b_j\rVert^2\).
\item Summing over \(j\):
\(\lVert AB\rVert_F^2=\sum_j\lVert Ab_j\rVert^2\leqslant\lVert A\rVert_F^2\sum_j\lVert b_j\rVert^2
=\lVert A\rVert_F^2\lVert B\rVert_F^2\).
\end{itemize}
(c) \(P'P=I_m\) and \(QQ'=I_n\) (Proposition~\ref{r6-prop-orthonormal}(b)), so by Theorem~\ref{r1-thm-trace}(b)
\[
\lVert PAQ\rVert_F^2=\tr(Q'A'P'PAQ)=\tr(Q'A'AQ)=\tr(A'AQQ')=\tr(A'A)=\lVert A\rVert_F^2
\]
\end{proof}

\begin{theorem}[Singular value decomposition; \citeat{harville}{Theorem~21.12.1, Corollary~21.12.2}]\label{r9-thm-svd}
Let \(A\in\R^{m\times n}\) have rank \(r\geqslant1\). Below, \(\left(\begin{smallmatrix}D&0\\0&0\end{smallmatrix}\right)\)
denotes the matrix of the size required by the product, with \(D\) in the top left \(r\times r\) block and zeros
elsewhere; blocks with no rows or no columns are absent (when \(r=n\) there is no \(Q_2\), and when \(r=m\) no \(P_2\)).
\begin{enumerate}
\item Let \(Q=[Q_1\ Q_2]\in\R^{n\times n}\) be orthogonal, with \(Q_1\in\R^{n\times r}\), and let \(D\in\R^{r\times r}\) be
an invertible diagonal matrix with
\[
Q'A'AQ=\begin{pmatrix}D^2&0\\0&0\end{pmatrix}
\]
Put \(P_1=AQ_1D^{-1}\in\R^{m\times r}\), and let \(P=[P_1\ P_2]\), where \(P_2\in\R^{m\times(m-r)}\) is any matrix with
\(P_1'P_2=0\). Then
\[
P'AQ=\begin{pmatrix}D&0\\0&0\end{pmatrix}
\]
\item Such \(Q\) and \(D\) exist with \(D=\diag(\sigma_1,\dots,\sigma_r)\) and \(\sigma_1\geqslant\dots\geqslant\sigma_r>0\):
take \(\sigma_i^2=\lambda_i(A'A)\), the nonzero eigenvalues of \(A'A\), and \(Q\) as in Lemma~\ref{r9-lem-eig-order} for
\(A'A\). Then \(P_1'P_1=I_r\), \(P_2\) can be chosen so that \(P\) is orthogonal, and, with the columns \(p_i\) of \(P\) and
\(q_i\) of \(Q\),
\[
A=P\begin{pmatrix}D&0\\0&0\end{pmatrix}Q'=P_1DQ_1'=\sigma_1p_1q_1'+\dots+\sigma_rp_rq_r'
\]
\end{enumerate}
\end{theorem}
\begin{proof}
(a)
\begin{itemize}
\item By the block product (Theorem~\ref{r5b-thm-block}), \(Q'A'AQ\) has the blocks \(Q_1'A'AQ_1\), \(Q_1'A'AQ_2\),
\(Q_2'A'AQ_1\), \(Q_2'A'AQ_2\). So \(Q_1'A'AQ_1=D^2\), and \((AQ_2)'(AQ_2)=Q_2'A'AQ_2=0\); its trace is \(0\), so
\(AQ_2=0\) (Theorem~\ref{r1-thm-trace}(d)).
\item \(AQ_1=P_1D\), and \(P_1'=D^{-1}Q_1'A'\), as \(D\) is diagonal. Hence
\[
P'AQ=\begin{pmatrix}P_1'AQ_1&P_1'AQ_2\\P_2'AQ_1&P_2'AQ_2\end{pmatrix}
=\begin{pmatrix}D^{-1}Q_1'A'AQ_1&0\\P_2'P_1D&0\end{pmatrix}
=\begin{pmatrix}D^{-1}D^2&0\\(P_1'P_2)'D&0\end{pmatrix}
=\begin{pmatrix}D&0\\0&0\end{pmatrix}
\]
\end{itemize}
\needspace{4\baselineskip}
(b)
\begin{itemize}
\item \(A'A\) is symmetric and PSD (Theorem~\ref{r8-thm-XtX}), so its eigenvalues are \(\geqslant0\)
(Theorem~\ref{r8-thm-eigen-criterion}(a)). As \(\rank(A'A)=\rank A=r\) (Theorem~\ref{thm:r4-rank-ata}(b)), exactly \(r\)
of them are nonzero (Theorem~\ref{r8-thm-rank-eigen}):
\(\lambda_1(A'A)\geqslant\dots\geqslant\lambda_r(A'A)>0=\lambda_{r+1}(A'A)=\dots=\lambda_n(A'A)\). So \(Q\) from
Lemma~\ref{r9-lem-eig-order} and \(D\) satisfy the condition of (a).
\item \(P_1'P_1=D^{-1}Q_1'A'AQ_1D^{-1}=D^{-1}D^2D^{-1}=I_r\): the columns of \(P_1\) are orthonormal, hence independent
(Proposition~\ref{r6-prop-orthogonal-list}(b)), and \(\rank P_1'=\rank P_1=r\) (Theorem~\ref{thm:r4-rank-rules}(a),(b)).
So \(\dim\mathcal N(P_1')=m-r\) (Theorem~\ref{thm:r4-rank-nullity}).
\item If \(r<m\), let the columns of \(P_2\) form an orthonormal basis of \(\mathcal N(P_1')\)
(Theorems~\ref{thm:r4-dimension} and~\ref{r6-thm-gram-schmidt}). Then \(P_1'P_2=0\) and \(P_2'P_2=I_{m-r}\), so
\(P'P=I_m\) by blocks: \(P\) is orthogonal.
\item As \(PP'=I_m\) and \(QQ'=I_n\) (Proposition~\ref{r6-prop-orthonormal}(b)), \(A=P(P'AQ)Q'\), the first form by (a).
By blocks it is \([P_1D\ \ 0]\left(\begin{smallmatrix}Q_1'\\Q_2'\end{smallmatrix}\right)=P_1DQ_1'\), and the outer
product expansion (Theorem~\ref{r1-thm-views}(d)), with \(\col_i(P_1D)=\sigma_ip_i\) and \(\row_i(Q_1')=q_i'\), gives the
sum.\qedhere
\end{itemize}
\end{proof}
A factorization \(A=P\left(\begin{smallmatrix}D&0\\0&0\end{smallmatrix}\right)Q'\) with orthogonal
\(P\in\R^{m\times m}\) and \(Q\in\R^{n\times n}\) and \(D=\diag(\sigma_1,\dots,\sigma_r)\),
\(\sigma_1\geqslant\dots\geqslant\sigma_r>0\), is a \textbf{singular value decomposition} (SVD) of \(A\), and
\(A=P_1DQ_1'\) is the corresponding \textbf{compact SVD}. The numbers \(\sigma_1,\dots,\sigma_r\) are the
\textbf{singular values} of \(A\), and the columns of \(P\) and \(Q\) are \textbf{singular vectors} of \(A\). We put
\(\sigma_k=0\) for \(r<k\leqslant\min\{m,n\}\), and \(\sigma_1=0\) for \(A=0\). (Theorem~\ref{r9-thm-svd-props}(a)
shows that \(r=\rank A\) and that \(\sigma_1,\dots,\sigma_r\) do not depend on the choice of the SVD.)

\needspace{16\baselineskip}
\begin{theorem}[\citeat{harville}{Theorem~21.12.3, \S21.12a}]\label{r9-thm-svd-props}
Let \(A\in\R^{m\times n}\), let \(P\in\R^{m\times m}\) and \(Q\in\R^{n\times n}\) be orthogonal, and let
\(D=\diag(\sigma_1,\dots,\sigma_r)\) with \(\sigma_1\geqslant\dots\geqslant\sigma_r>0\) be such that
\(P'AQ=\left(\begin{smallmatrix}D&0\\0&0\end{smallmatrix}\right)\). Split \(P=[P_1\ P_2]\) and \(Q=[Q_1\ Q_2]\), where
\(P_1\) and \(Q_1\) have the \(r\) columns \(p_1,\dots,p_r\) and \(q_1,\dots,q_r\).
\begin{enumerate}
\item \(r=\rank A\), and \(\sigma_1^2,\dots,\sigma_r^2\) are the nonzero eigenvalues of \(A'A\) and of \(AA'\), listed with
multiplicity. So the singular values of \(A\) depend only on \(A\), not on the choice of \(P\) and \(Q\).
\item \(Q'A'AQ=\left(\begin{smallmatrix}D^2&0\\0&0\end{smallmatrix}\right)\) and
\(P'AA'P=\left(\begin{smallmatrix}D^2&0\\0&0\end{smallmatrix}\right)\): the columns of \(Q\) are orthonormal eigenvectors
of \(A'A\), and those of \(P\) of \(AA'\); \(q_i\) and \(p_i\) belong to \(\sigma_i^2\) (\(i\leqslant r\)), the other columns to
\(0\).
\item \(P_1=AQ_1D^{-1}\) and \(Q_1=A'P_1D^{-1}\); that is, \(Aq_i=\sigma_ip_i\) and \(A'p_i=\sigma_iq_i\)
(\(i=1,\dots,r\)).
\item \(\lVert A\rVert_F^2=\sigma_1^2+\dots+\sigma_r^2\).
\item If \(A\) is symmetric, its singular values are the absolute values of its nonzero eigenvalues; if \(A\) is PSD,
they are its nonzero eigenvalues. For orthogonal \(U\in\R^{m\times m}\) and \(V\in\R^{n\times n}\), the matrices \(UAV\)
and \(A'\) have the same singular values as \(A\).
\end{enumerate}
\end{theorem}
\begin{proof}
(a)
\begin{itemize}
\item \(\rank A=\rank(P'AQ)\) (Theorem~\ref{thm:r4-rank-product}(d)). The nonzero columns of \(P'AQ\) are
\(\sigma_1e_1,\dots,\sigma_re_r\), which are independent (Proposition~\ref{r6-prop-orthogonal-list}(b)) and span its
column space; so the rank is \(r\).
\item As \(PP'=I_m\), \(Q'A'AQ=Q'A'PP'AQ=(P'AQ)'(P'AQ)\), the product of an \(n\times m\) and an \(m\times n\) block matrix,
which is \(\left(\begin{smallmatrix}D^2&0\\0&0\end{smallmatrix}\right)\) (Theorem~\ref{r5b-thm-block}); likewise
\(P'AA'P=(P'AQ)(P'AQ)'=\left(\begin{smallmatrix}D^2&0\\0&0\end{smallmatrix}\right)\).
\item As \(Q^{-1}=Q'\) (Proposition~\ref{r6-prop-orthonormal}(b)), \(Q^{-1}(A'A)Q\) is diagonal with the diagonal
entries \(\sigma_1^2,\dots,\sigma_r^2,0,\dots,0\); these are the eigenvalues of \(A'A\) listed with multiplicity (Definition~\ref{r7-def-diagonalizable}), and likewise for \(AA'\) with
\(P\). Hence \(\sigma_1\geqslant\dots\geqslant\sigma_r\) are the positive square roots of the nonzero eigenvalues of
\(A'A\), in decreasing order.
\end{itemize}
(b) The two identities were shown in (a). Multiplying the first by \(Q\) on the left,
\(A'AQ=Q\left(\begin{smallmatrix}D^2&0\\0&0\end{smallmatrix}\right)\), whose columns say \(A'Aq_i=\sigma_i^2q_i\)
(\(i\leqslant r\)) and \(A'Aq_i=0\) (\(i>r\)); likewise for \(P\).

\needspace{3\baselineskip}
(c)
\begin{itemize}
\item \(PP'=P_1P_1'+P_2P_2'=I_m\), and \(P_1'AQ_1=D\) and \(P_2'AQ_1=0\) are blocks of \(P'AQ\). So
\[
P_1=P_1(P_1'AQ_1)D^{-1}=(I_m-P_2P_2')AQ_1D^{-1}=AQ_1D^{-1}-P_2(P_2'AQ_1)D^{-1}=AQ_1D^{-1}
\]
\item In the same way, \(QQ'=Q_1Q_1'+Q_2Q_2'=I_n\), \(Q_1'A'P_1=(P_1'AQ_1)'=D\) and \(Q_2'A'P_1=(P_1'AQ_2)'=0\), so
\[
Q_1=Q_1(Q_1'A'P_1)D^{-1}=(I_n-Q_2Q_2')A'P_1D^{-1}=A'P_1D^{-1}-Q_2(Q_2'A'P_1)D^{-1}=A'P_1D^{-1}
\]
\end{itemize}
(d) By Proposition~\ref{r9-prop-frobenius}(c), \(\lVert A\rVert_F=\lVert P'AQ\rVert_F\), and the squares of the entries of
\(P'AQ\) add up to \(\sigma_1^2+\dots+\sigma_r^2\).

\needspace{3\baselineskip}
(e)
\begin{itemize}
\item Let \(A\) be symmetric; we choose a new pair \(P\), \(Q\). Number an orthonormal basis of eigenvectors \(q_1,\dots,q_n\)
(Theorem~\ref{r8-thm-spectral}(a)), with eigenvalues \(d_1,\dots,d_n\), so that
\(\lvert d_1\rvert\geqslant\dots\geqslant\lvert d_n\rvert\). Then \(Q'AQ=\diag(d_1,\dots,d_n)\) for
\(Q=[q_1\ \cdots\ q_n]\), and \(d_1,\dots,d_n\) are the eigenvalues of \(A\), listed with multiplicity
(Theorem~\ref{r8-thm-spectral}(b),(c)). By Theorem~\ref{r8-thm-rank-eigen}, exactly \(\rank A\) of the \(d_i\) are
nonzero, namely \(d_1,\dots,d_{\rank A}\).
\item Let \(\Delta=\diag(\delta_1,\dots,\delta_n)\), with \(\delta_i=1\) if \(d_i\geqslant0\) and \(\delta_i=-1\) if
\(d_i<0\), and \(P=Q\Delta\). Then \(P'P=\Delta Q'Q\Delta=\Delta^2=I_n\), and
\[
P'AQ=\Delta Q'AQ=\diag(\delta_1d_1,\dots,\delta_nd_n)=\diag(\lvert d_1\rvert,\dots,\lvert d_n\rvert)
\]
is an SVD of \(A\). By (a), the singular values are the nonzero \(\lvert d_i\rvert\). If \(A\) is PSD, every
\(d_i\geqslant0\) (Theorem~\ref{r8-thm-eigen-criterion}(a)), so \(\lvert d_i\rvert=d_i\) (and \(P=Q\)).
\item In the setting of the theorem, \(UP\) and \(V'Q\) are orthogonal (\((UP)'(UP)=P'U'UP=I_m\), and likewise
\((V'Q)'(V'Q)=I_n\)), and \((UP)'(UAV)(V'Q)=P'AQ\): an SVD of \(UAV\) with the same \(D\). Also \(Q'A'P=(P'AQ)'\) is the
\(n\times m\) matrix with \(D'=D\) in the top left corner: an SVD of \(A'\) with \(Q\) and \(P\) exchanged.\qedhere
\end{itemize}
\end{proof}

\begin{corollary}[\citeat{harville-solutions}{Ch.~21, Exercise~28}]\label{r9-cor-svd-subspaces}
In Theorem~\ref{r9-thm-svd-props} let \(r\geqslant1\). Then
\[
\mathcal C(A)=\mathcal C(P_1),\qquad\mathcal N(A')=\mathcal C(P_2),\qquad\mathcal C(A')=\mathcal C(Q_1),\qquad
\mathcal N(A)=\mathcal C(Q_2)
\]
so the columns of \(P_1\), \(P_2\), \(Q_1\), \(Q_2\) are orthonormal bases of these four subspaces (an absent \(P_2\) or
\(Q_2\) means that the subspace is \(\{0\}\)). Moreover, \(P_1P_1'y\) is the orthogonal projection of \(y\in\R^m\) onto
\(\mathcal C(A)\), and if \(A\) has full column rank, then \(A(A'A)^{-1}A'=P_1P_1'\).
\end{corollary}
\begin{proof}
\begin{itemize}
\item As \(PP'=I_m\) and \(QQ'=I_n\), \(A=P(P'AQ)Q'=P\left(\begin{smallmatrix}D&0\\0&0\end{smallmatrix}\right)Q'=P_1DQ_1'\)
by blocks, so
\(Ax=P_1(DQ_1'x)\in\mathcal C(P_1)\); and \(P_1=AQ_1D^{-1}\) (Theorem~\ref{r9-thm-svd-props}(c)) gives
\(\mathcal C(P_1)\subseteq\mathcal C(A)\).
\item \(Q'Q=I_n\) gives \(Q_1'Q_2=0\), so \(AQ_2=P_1DQ_1'Q_2=0\) and \(\mathcal C(Q_2)\subseteq\mathcal N(A)\). The \(n-r\)
columns of \(Q_2\) are orthonormal, hence independent (Proposition~\ref{r6-prop-orthogonal-list}(b)), so
\(\dim\mathcal C(Q_2)=n-r=\dim\mathcal N(A)\) (Theorems~\ref{thm:r4-rank-nullity} and~\ref{r9-thm-svd-props}(a)). Hence
\(\mathcal C(Q_2)=\mathcal N(A)\) (Theorem~\ref{thm:r4-counting}(a)).
\item \(Q'A'P=(P'AQ)'\) is an SVD of \(A'\) with \(Q\) and \(P\) exchanged (Theorem~\ref{r9-thm-svd-props}(e)); the two
items above, for \(A'\), give \(\mathcal C(A')=\mathcal C(Q_1)\) and \(\mathcal N(A')=\mathcal C(P_2)\).
\item \(P_1\) has full column rank and \(P_1'P_1=I_r\), so by Theorem~\ref{r6-thm-hat}(a),
\(P_1(P_1'P_1)^{-1}P_1'y=P_1P_1'y\) is the orthogonal projection of \(y\) onto \(\mathcal C(P_1)=\mathcal C(A)\). If \(A\)
has full column rank, \(A(A'A)^{-1}A'y\) is also this projection (Theorem~\ref{r6-thm-hat}(a)), which is unique
(Definition~\ref{r6-def-projection}); so \(A(A'A)^{-1}A'y=P_1P_1'y\) for every \(y\), and with \(y=e_1,\dots,e_m\) the two
matrices have the same columns.\qedhere
\end{itemize}
\end{proof}

\begin{example}[\citeat{searle-khuri}{Example~7.3}, \citeat{abadir-magnus}{Exercise~8.39}]\label{r9-exm-svd}
Find an SVD and the compact SVD of
\[
A=\begin{pmatrix}2&0&1&1\\0&2&1&1\\1&1&1&1\end{pmatrix}
\]
and orthonormal bases of its four subspaces.

\textit{Solution.}
\begin{itemize}
\item The columns satisfy \(a_3=a_4=\frac12(a_1+a_2)\), and \(a_1\), \(a_2\) are independent; so \(\rank A=2=r\).
Hence \(A(0,0,1,-1)'=a_3-a_4=0\) and \(A(1,1,-1,-1)'=a_1+a_2-a_3-a_4=0\).
\item The matrix
\[
A'A=\begin{pmatrix}5&1&3&3\\1&5&3&3\\3&3&3&3\\3&3&3&3\end{pmatrix}
\]
has every row sum \(12\), so \(A'A\one_4=12\one_4\); and \(A'A(1,-1,0,0)'=(4,-4,0,0)'\). The two vectors of the first
item lie in \(\mathcal N(A'A)=\mathcal N(A)\) (Theorem~\ref{thm:r4-rank-ata}(a)), of dimension \(4-2=2\)
(Theorem~\ref{thm:r4-rank-nullity}). So the algebraic multiplicities \(\dim\mathcal N(A'A-cI)\)
(Theorem~\ref{r8-thm-rank-eigen}) are at least \(1\), \(1\), \(2\) for \(c=12,4,0\). All multiplicities add up to \(4\)
(Definition~\ref{r7-def-eigenvalues}); so these are exactly \(1\), \(1\), \(2\), and the eigenvalues of \(A'A\) are \(12\),
\(4\), \(0\), \(0\). Check: \(\tr(A'A)=16=12+4\).
\item The four vectors are pairwise orthogonal (Theorem~\ref{r8-thm-sym-orth}, and
\((1,1,-1,-1)(0,0,1,-1)'=0\)). Normalized: \(q_1=\frac12(1,1,1,1)'\), \(q_2=\frac1{\sqrt2}(1,-1,0,0)'\),
\(q_3=\frac12(1,1,-1,-1)'\), \(q_4=\frac1{\sqrt2}(0,0,1,-1)'\); and \(D=\diag(\sqrt{12},2)=\diag(2\sqrt3,2)\).
\item \(P_1=AQ_1D^{-1}\): \(Aq_1=(2,2,2)'\) and \(Aq_2=\frac1{\sqrt2}(2,-2,0)'\), so \(p_1=\frac1{2\sqrt3}(2,2,2)'
=\frac1{\sqrt3}(1,1,1)'\) and \(p_2=\frac12Aq_2=\frac1{\sqrt2}(1,-1,0)'\).
\item \(\mathcal N(P_1')\): \(y_1+y_2+y_3=0\) and \(y_1-y_2=0\) give \(y=t(1,1,-2)'\); so \(p_3=\frac1{\sqrt6}(1,1,-2)'\).
\item Hence \(A=P\left(\begin{smallmatrix}D&0\\0&0\end{smallmatrix}\right)Q'\) with
\[
P=\begin{pmatrix}\frac1{\sqrt3}&\frac1{\sqrt2}&\frac1{\sqrt6}\\[2pt]\frac1{\sqrt3}&-\frac1{\sqrt2}&\frac1{\sqrt6}\\[2pt]
\frac1{\sqrt3}&0&-\frac2{\sqrt6}\end{pmatrix},\qquad
\begin{pmatrix}D&0\\0&0\end{pmatrix}=\begin{pmatrix}2\sqrt3&0&0&0\\0&2&0&0\\0&0&0&0\end{pmatrix},\qquad
Q=\begin{pmatrix}\frac12&\frac1{\sqrt2}&\frac12&0\\[2pt]\frac12&-\frac1{\sqrt2}&\frac12&0\\[2pt]
\frac12&0&-\frac12&\frac1{\sqrt2}\\[2pt]\frac12&0&-\frac12&-\frac1{\sqrt2}\end{pmatrix}
\]
and the compact SVD is \(A=2\sqrt3\,p_1q_1'+2\,p_2q_2'\). Check:
\[
2\sqrt3\,p_1q_1'+2\,p_2q_2'=\begin{pmatrix}1\\1\\1\end{pmatrix}(1,1,1,1)+\begin{pmatrix}1\\-1\\0\end{pmatrix}(1,-1,0,0)
=\begin{pmatrix}1&1&1&1\\1&1&1&1\\1&1&1&1\end{pmatrix}+\begin{pmatrix}1&-1&0&0\\-1&1&0&0\\0&0&0&0\end{pmatrix}=A
\]
\item By Corollary~\ref{r9-cor-svd-subspaces}: \(p_1,p_2\) is an orthonormal basis of \(\mathcal C(A)\), \(p_3\) of
\(\mathcal N(A')\), \(q_1,q_2\) of \(\mathcal C(A')\), and \(q_3,q_4\) of \(\mathcal N(A)\).
\end{itemize}
The vector \(-p_2\) is also a unit eigenvector of \(AA'\) belonging to \(4\), but \(2\sqrt3\,p_1q_1'-2\,p_2q_2'\ne A\). So
\(P_1\) must be computed from \(Q_1\) as \(AQ_1D^{-1}\); it cannot be chosen separately among the eigenvectors of \(AA'\).
\end{example}

\subsection*{Singular values as maxima; the spectral norm}

\begin{theorem}[\citeat{bapat}{Results~5.5, 5.6}]\label{r9-thm-sigma-max}
Let \(A\in\R^{m\times n}\), \(A\ne0\), have an SVD as in Theorem~\ref{r9-thm-svd-props}, with the columns
\(p_1,\dots,p_m\) of \(P\) and \(q_1,\dots,q_n\) of \(Q\).
\begin{enumerate}
\item \(\lvert u'Av\rvert\leqslant\sigma_1\) for all unit vectors \(u\in\R^m\) and \(v\in\R^n\), and \(p_1'Aq_1=\sigma_1\).
Hence
\[
\sigma_1=\max\{u'Av:\lVert u\rVert=\lVert v\rVert=1\}
\]
\item Let \(2\leqslant k\leqslant\min\{m,n\}\). If moreover \(u\perp p_1,\dots,p_{k-1}\) and \(v\perp q_1,\dots,q_{k-1}\),
then \(\lvert u'Av\rvert\leqslant\sigma_k\); and \(p_k'Aq_k=\sigma_k\). So \(\sigma_k\) is the maximum of \(u'Av\) under
these conditions.
\end{enumerate}
\end{theorem}
\begin{proof}
Let \(u\in\R^m\) and \(v\in\R^n\) be unit vectors, \(w=P'u\) and \(z=Q'v\). Then \(\lVert w\rVert=\lVert z\rVert=1\)
(Proposition~\ref{r6-prop-orthonormal}(a); \(P'\) and \(Q'\) are orthogonal), and
\(A=P\left(\begin{smallmatrix}D&0\\0&0\end{smallmatrix}\right)Q'\) gives
\[
u'Av=w'\begin{pmatrix}D&0\\0&0\end{pmatrix}z=\sigma_1w_1z_1+\dots+\sigma_rw_rz_r
\]
\needspace{3\baselineskip}
(a)
\begin{itemize}
\item As \(0<\sigma_i\leqslant\sigma_1\),
\(\lvert u'Av\rvert\leqslant\sigma_1\sum_{i\leqslant r}\lvert w_i\rvert\lvert z_i\rvert\leqslant\sigma_1\lVert w\rVert
\lVert z\rVert=\sigma_1\), by Cauchy--Schwarz (Theorem~\ref{r6-thm-cauchy-schwarz}(a)) for the vectors
\((\lvert w_1\rvert,\dots,\lvert w_r\rvert)'\) and \((\lvert z_1\rvert,\dots,\lvert z_r\rvert)'\), whose lengths are at most
\(\lVert w\rVert\) and \(\lVert z\rVert\).
\item \(p_1'Aq_1=p_1'(\sigma_1p_1)=\sigma_1\) (Theorem~\ref{r9-thm-svd-props}(c)). As \(\sigma_1>0\), the maximum of
\(u'Av\) is \(\sigma_1\).
\end{itemize}
\needspace{3\baselineskip}
(b)
\begin{itemize}
\item The entries of \(w=P'u\) are \(w_i=p_i'u\), so \(u\perp p_1,\dots,p_{k-1}\) means \(w_1=\dots=w_{k-1}=0\); likewise
\(z_1=\dots=z_{k-1}=0\).
\item So \(u'Av=\sum_{i=k}^r\sigma_iw_iz_i\) (\(0\) if \(k>r\)), and as in (a) \(\lvert u'Av\rvert\leqslant\sigma_k\).
\item \(p_k\) and \(q_k\) satisfy the conditions, and \(p_k'Aq_k=\sigma_k\): for \(k\leqslant r\), \(Aq_k=\sigma_kp_k\); for
\(k>r\), \(Aq_k=0\) (Corollary~\ref{r9-cor-svd-subspaces}).\qedhere
\end{itemize}
\end{proof}

\begin{definition}[\citeat{searle-khuri}{(9.3)}, \citeat{bapat}{Result~5.5}]\label{r9-def-spectral-norm}
The \textbf{spectral norm} of \(A\in\R^{m\times n}\) is \(\lVert A\rVert_2=\sigma_1\), its largest singular value; and
\(\lVert0\rVert_2=0\). It is not the Frobenius norm of Definition~\ref{r9-def-frobenius}, which some books also write
\(\lVert A\rVert_2\).
\end{definition}

\begin{proposition}[\citeat{abadir-magnus}{Exercise~7.53(b)}, \citeat{bapat}{Result~5.5}]\label{r9-prop-spectral-norm}
Let \(A\in\R^{m\times n}\).
\begin{enumerate}
\item \(\lVert A\rVert_2=\bigl(\lambda_1(A'A)\bigr)^{1/2}=\max_{x\ne0}\lVert Ax\rVert/\lVert x\rVert\); so
\(\lVert Ax\rVert\leqslant\lVert A\rVert_2\lVert x\rVert\) for every \(x\in\R^n\), with equality at \(x=q_1\) if
\(A\ne0\), where \(q_1\) is the first column of \(Q\) in an SVD of \(A\) as in Theorem~\ref{r9-thm-svd-props}. If
\(\rank A=n\), then also \(\lVert Ax\rVert\geqslant\sigma_n\lVert x\rVert\) for every \(x\in\R^n\).
\item \(\lVert A\rVert_2=\max\{u'Av:\lVert u\rVert=\lVert v\rVert=1\}\).
\item \(\lVert UAV\rVert_2=\lVert A\rVert_2=\lVert A'\rVert_2\) for orthogonal \(U\in\R^{m\times m}\) and
\(V\in\R^{n\times n}\).
\item \(\lVert A\rVert_2\leqslant\lVert A\rVert_F\).
\item If \(A\) is symmetric, then \(\lVert A\rVert_2=\max_i\lvert\lambda_i(A)\rvert=\max\{\lambda_1(A),-\lambda_n(A)\}\); if
\(A\) is PSD, then \(\lVert A\rVert_2=\lambda_1(A)\).
\end{enumerate}
\end{proposition}
\begin{proof}
For \(A=0\) every statement reads \(0=0\) or \(0\leqslant0\). Let \(A\ne0\), with an SVD as in
Theorem~\ref{r9-thm-svd-props}.

\needspace{3\baselineskip}
(a)
\begin{itemize}
\item \(\lambda_1(A'A)=\sigma_1^2\) (Theorem~\ref{r9-thm-svd-props}(a)).
\item For \(x\ne0\), \(\lVert Ax\rVert^2/\lVert x\rVert^2=x'A'Ax/x'x\leqslant\lambda_1(A'A)=\sigma_1^2\)
(Theorem~\ref{r9-thm-rayleigh}(a) for \(A'A\)); take square roots.
\item \(\lVert Aq_1\rVert=\lVert\sigma_1p_1\rVert=\sigma_1=\sigma_1\lVert q_1\rVert\) (Theorem~\ref{r9-thm-svd-props}(c)).
\item Let \(\rank A=n\). Then \(A'A\in\R^{n\times n}\) has the \(n\) nonzero eigenvalues
\(\sigma_1^2\geqslant\dots\geqslant\sigma_n^2\) (Theorem~\ref{r9-thm-svd-props}(a), with \(r=n\)), so
\(\lambda_n(A'A)=\sigma_n^2\). For \(x\ne0\), \(\lVert Ax\rVert^2=x'A'Ax\geqslant\lambda_n(A'A)\,x'x=\sigma_n^2\lVert x\rVert^2\)
(Theorem~\ref{r9-thm-rayleigh}(a) for \(A'A\)); take square roots. For \(x=0\) both sides are \(0\).
\end{itemize}

(b) This is Theorem~\ref{r9-thm-sigma-max}(a). (c) By Theorem~\ref{r9-thm-svd-props}(e), the three matrices have the same
largest singular value. (d) \(\sigma_1^2\leqslant\sigma_1^2+\dots+\sigma_r^2=\lVert A\rVert_F^2\)
(Theorem~\ref{r9-thm-svd-props}(d)).

(e) By Theorem~\ref{r9-thm-svd-props}(e), \(\sigma_1\) is the largest \(\lvert\lambda_i\rvert\) with \(\lambda_i\ne0\),
which is \(\max_i\lvert\lambda_i\rvert\), as \(A\ne0\) has a nonzero eigenvalue (Theorem~\ref{r8-thm-rank-eigen}). As
\(\lambda_n\leqslant\lambda_i\leqslant\lambda_1\), \(\lvert\lambda_i\rvert\leqslant\max\{\lambda_1,-\lambda_n\}\), with equality
for \(i=1\) or \(i=n\). If \(A\) is PSD, then \(\lambda_n\geqslant0\) (Theorem~\ref{r8-thm-eigen-criterion}(a)), and the
maximum is \(\lambda_1\).
\end{proof}

\needspace{10\baselineskip}
\begin{proposition}[\citeat{knapp-basic-real-analysis}{Ch.~III, Proposition~3.2}]\label{r9-prop-spectral-norm-rules}
Let \(A,B\in\R^{m\times n}\).
\begin{enumerate}
\item \(\lVert A\rVert_2\geqslant0\), with equality only for \(A=0\); \(\lVert cA\rVert_2=\lvert c\rvert\,\lVert A\rVert_2\) for
\(c\in\R\); and \(\lVert A+B\rVert_2\leqslant\lVert A\rVert_2+\lVert B\rVert_2\).
\item \(\lVert AC\rVert_2\leqslant\lVert A\rVert_2\lVert C\rVert_2\) for \(C\in\R^{n\times p}\).
\end{enumerate}
\end{proposition}
\begin{proof}
By Proposition~\ref{r9-prop-spectral-norm}(a), every matrix \(M\in\R^{s\times t}\) has
\(\lVert M\rVert_2=\max_{x\ne0}\lVert Mx\rVert/\lVert x\rVert\) and \(\lVert Mx\rVert\leqslant\lVert M\rVert_2\lVert x\rVert\)
for every \(x\in\R^t\).

\needspace{3\baselineskip}
(a)
\begin{itemize}
\item \(\lVert A\rVert_2\) is a maximum of numbers \(\geqslant0\). If \(\lVert A\rVert_2=0\), then \(Ax=0\) for every \(x\),
so \(A=[Ae_1\ \cdots\ Ae_n]=0\) (Theorem~\ref{r1-thm-views}(b)).
\item \(\lVert cAx\rVert=\lvert c\rvert\,\lVert Ax\rVert\) for every \(x\) (the remarks after
Definition~\ref{r6-def-inner}), so the maxima for \(cA\) and for \(A\) differ by the factor \(\lvert c\rvert\).
\item For every \(x\ne0\), by Theorem~\ref{r6-thm-cauchy-schwarz}(b),
\[
\lVert(A+B)x\rVert=\lVert Ax+Bx\rVert\leqslant\lVert Ax\rVert+\lVert Bx\rVert
\leqslant\bigl(\lVert A\rVert_2+\lVert B\rVert_2\bigr)\lVert x\rVert
\]
Divide by \(\lVert x\rVert\) and take the maximum over \(x\ne0\).
\end{itemize}
(b) For every \(x\ne0\) in \(\R^p\),
\(\lVert ACx\rVert=\lVert A(Cx)\rVert\leqslant\lVert A\rVert_2\lVert Cx\rVert\leqslant\lVert A\rVert_2\lVert C\rVert_2\lVert x\rVert\).
Divide by \(\lVert x\rVert\) and take the maximum over \(x\ne0\).
\end{proof}

\subsection*{Approximation by matrices of lower rank}

\begin{theorem}[Eckart--Young; \citeat{harville}{Theorem~21.12.4}]\label{r12-thm-eckart-young}
Let \(A\in\R^{m\times n}\) have rank \(r\) and singular value decomposition
\[
A=P\begin{pmatrix}D&0\\0&0\end{pmatrix}Q'=\sigma_1p_1q_1'+\dots+\sigma_rp_rq_r',\qquad
D=\diag(\sigma_1,\dots,\sigma_r),\quad\sigma_1\geqslant\dots\geqslant\sigma_r>0
\]
(Theorem~\ref{r9-thm-svd}), let \(1\leqslant k<r\), and let \(A_k=\sigma_1p_1q_1'+\dots+\sigma_kp_kq_k'\), which has rank \(k\). For
every \(B\in\R^{m\times n}\) with \(\rank B\leqslant k\):
\begin{enumerate}
\item \(\lVert A-B\rVert_F^2\geqslant\sigma_{k+1}^2+\dots+\sigma_r^2\) (Definition~\ref{r9-def-frobenius}), with equality for \(B=A_k\);
\item \(\lVert A-B\rVert_2\geqslant\sigma_{k+1}\) (Definition~\ref{r9-def-spectral-norm}), with equality for \(B=A_k\).
\end{enumerate}
\end{theorem}
\begin{proof}
(a)
\begin{itemize}
\item Let \(s=\rank B\leqslant k\). Take an orthonormal basis \(u_1,\dots,u_s\) of \(\mathcal C(B)\) and, as
\(\dim\mathcal C(B)^\perp=m-s\geqslant k-s\) (Theorem~\ref{r6-thm-complement}(c); \(k<r\leqslant m\)), \(k-s\) vectors
\(u_{s+1},\dots,u_k\) of an orthonormal basis of \(\mathcal C(B)^\perp\) (Theorems~\ref{thm:r4-dimension}
and~\ref{r6-thm-gram-schmidt}). Then \(U=[u_1\ \cdots\ u_k]\) has \(U'U=I_k\), as \(\mathcal C(B)\perp\mathcal C(B)^\perp\), and
\(\mathcal C(B)\subseteq\mathcal C(U)\):
column \(j\) of \(B\) is \(U\ell_j\) for some \(\ell_j\in\R^k\), so \(B=UL\) with \(L=[\ell_1\ \cdots\ \ell_n]\)
(Theorem~\ref{r1-thm-views}(a)). It therefore suffices to show
\(\lVert UL-A\rVert_F^2\geqslant\sigma_{k+1}^2+\dots+\sigma_r^2\) for all \(U\in\R^{m\times k}\) with \(U'U=I_k\) and all
\(L\in\R^{k\times n}\).
\item As \(U'U=I_k\), \((UL-A)'(UL-A)=L'L-A'UL-L'U'A+A'A\), and
\(L'L-A'UL-L'U'A+A'UU'A=(L-U'A)'(L-U'A)\). By Definition~\ref{r9-def-frobenius} and Theorem~\ref{r1-thm-trace}(d),~(b),
\[
\lVert UL-A\rVert_F^2=\tr\bigl[(L-U'A)'(L-U'A)\bigr]-\tr(A'UU'A)+\tr(A'A)\geqslant\tr(A'A)-\tr(U'AA'U)
\]
\item \(Q'Q=I_n\), so \(AA'=P\diag(\sigma_1^2,\dots,\sigma_r^2,0,\dots,0)P'\) (Theorem~\ref{r5b-thm-block}(c)). Thus
\(\tr(A'A)=\tr(AA')=\sigma_1^2+\dots+\sigma_r^2\) (Theorem~\ref{r1-thm-trace}(b), \(P'P=I_m\)), and the columns of \(P\)
are orthonormal eigenvectors of the symmetric matrix \(AA'\) for the eigenvalues
\(\sigma_1^2\geqslant\dots\geqslant\sigma_r^2\geqslant0=\dots=0\) (Theorem~\ref{r8-thm-spectral}(c)).
\item By Theorem~\ref{r12-thm-kyfan}(a) for \(AA'\in\R^{m\times m}\), \(\tr(U'AA'U)\leqslant\sigma_1^2+\dots+\sigma_k^2\). Hence
\(\lVert UL-A\rVert_F^2\geqslant\sum_{i\leqslant r}\sigma_i^2-\sum_{i\leqslant k}\sigma_i^2=\sigma_{k+1}^2+\dots+\sigma_r^2\).
\item \(A-A_k=\sum_{i=k+1}^r\sigma_ip_iq_i'=PEQ'\) (Theorem~\ref{r1-thm-views}(d)), where \(E\in\R^{m\times n}\) has the
entries \(\sigma_{k+1},\dots,\sigma_r\) in positions \((k+1,k+1),\dots,(r,r)\) and zeros elsewhere. For orthogonal \(P\),
\(Q\), \(\lVert PEQ'\rVert_F^2=\tr(QE'P'PEQ')=\tr(E'EQ'Q)=\tr(E'E)=\sigma_{k+1}^2+\dots+\sigma_r^2\)
(Theorem~\ref{r1-thm-trace}(b)): equality holds.
\item Likewise \(A_k=PE_kQ'\), where \(E_k\) has \(\sigma_1,\dots,\sigma_k\) in positions \((1,1),\dots,(k,k)\); its
nonzero columns \(\sigma_ie_i\) are independent, so \(\rank A_k=\rank E_k=k\)
(Theorem~\ref{thm:r4-rank-product}(d)).
\end{itemize}
(b)
\begin{itemize}
\item Let \(\rank B\leqslant k\). Then \(\dim\mathcal N(B)=n-\rank B\geqslant n-k\) (Theorem~\ref{thm:r4-rank-nullity}), and
\(T=\Span\{q_1,\dots,q_{k+1}\}\) has dimension \(k+1\) (orthonormal vectors are independent,
Proposition~\ref{r6-prop-orthogonal-list}(b)). As \((n-k)+(k+1)>n\), there is a unit vector \(x\in\mathcal N(B)\cap T\)
(Lemma~\ref{r9-lem-intersect}).
\item Write \(x=c_1q_1+\dots+c_{k+1}q_{k+1}\); then \(c_1^2+\dots+c_{k+1}^2=x'x=1\), as the \(q_i\) are orthonormal. Since
\(k+1\leqslant r\), \(Aq_i=\sigma_ip_i\) (Theorem~\ref{r9-thm-svd-props}(c)), so \(Ax=\sum_{i\leqslant k+1}c_i\sigma_ip_i\), and, the
\(p_i\) being orthonormal,
\[
\lVert Ax\rVert^2=\sum_{i\leqslant k+1}c_i^2\sigma_i^2\geqslant\sigma_{k+1}^2\sum_{i\leqslant k+1}c_i^2=\sigma_{k+1}^2
\]
\item \(Bx=0\), so \((A-B)x=Ax\), and \(\lVert A-B\rVert_2\geqslant\lVert(A-B)x\rVert=\lVert Ax\rVert\geqslant\sigma_{k+1}\)
(Proposition~\ref{r9-prop-spectral-norm}(a)).
\item Equality: \(A-A_k=PEQ'\) with \(E\) as in (a), so \(\lVert A-A_k\rVert_2=\lVert E\rVert_2\)
(Proposition~\ref{r9-prop-spectral-norm}(c)). For \(y\in\R^n\), \(\lVert Ey\rVert^2=\sum_{i=k+1}^r\sigma_i^2y_i^2\leqslant
\sigma_{k+1}^2\lVert y\rVert^2\), with equality for \(y=e_{k+1}\); so \(\lVert E\rVert_2=\sigma_{k+1}\)
(Proposition~\ref{r9-prop-spectral-norm}(a)).\qedhere
\end{itemize}
\end{proof}
So \(A_k\) keeps the part \(\sigma_1^2+\dots+\sigma_k^2\) of \(\lVert A\rVert_F^2=\sigma_1^2+\dots+\sigma_r^2\) and loses
\(\sigma_{k+1}^2+\dots+\sigma_r^2\); no matrix \(B\) of rank at most \(k\) does better:
\(\lVert A\rVert_F^2-\lVert A-B\rVert_F^2\leqslant\sigma_1^2+\dots+\sigma_k^2\). As \(A_k=PE_kQ'\) in the proof of (a),
\(\lVert A_k\rVert_F^2=\lVert E_k\rVert_F^2=\sigma_1^2+\dots+\sigma_k^2\) (Proposition~\ref{r9-prop-frobenius}(c)): \(A_k\)
reproduces the proportion
\[
\frac{\lVert A_k\rVert_F^2}{\lVert A\rVert_F^2}=\frac{\sigma_1^2+\dots+\sigma_k^2}{\sigma_1^2+\dots+\sigma_r^2}
\]
of the sum of squares \(\lVert A\rVert_F^2=\sum_{i,j}a_{ij}^2\) of the entries of \(A\).

\begin{corollary}[\citeat{harville}{Corollary~21.12.6}]\label{r12-cor-symmetric}
Let \(A\in\R^{n\times n}\) be symmetric of rank \(r\geqslant2\). Number an orthonormal basis \(q_1,\dots,q_n\) of
eigenvectors of \(A\), \(Aq_i=d_iq_i\), so that \(\lvert d_1\rvert\geqslant\dots\geqslant\lvert d_n\rvert\)
(Theorem~\ref{r8-thm-spectral}(a)). Then \(d_1,\dots,d_r\) are the nonzero eigenvalues of \(A\), listed with
multiplicity, and \(d_{r+1}=\dots=d_n=0\) (Theorems~\ref{r8-thm-spectral}(c) and~\ref{r8-thm-rank-eigen}). Let
\(1\leqslant k<r\). For every \(B\in\R^{n\times n}\) with \(\rank B\leqslant k\),
\[
\lVert A-B\rVert_F^2\geqslant d_{k+1}^2+\dots+d_r^2
\]
with equality for \(B=d_1q_1q_1'+\dots+d_kq_kq_k'\). If \(A\) is PSD, then \(d_1\geqslant\dots\geqslant d_r>0\) are its
positive eigenvalues.
\end{corollary}
\begin{proof}
\begin{itemize}
\item \(Q=[q_1\ \cdots\ q_n]\) is orthogonal and \(Q'AQ=\diag(d_1,\dots,d_n)\) (Theorem~\ref{r8-thm-spectral}(b)). Let
\(\delta_i=-1\) if \(d_i<0\) and \(\delta_i=1\) otherwise, \(\Delta=\diag(\delta_1,\dots,\delta_n)\) and \(P=Q\Delta\), whose
columns are \(p_i=\delta_iq_i\) (Theorem~\ref{r1-thm-views}(a)). Then \(P'P=\Delta Q'Q\Delta=\Delta^2=I_n\), and, as
\(\Delta\) scales the rows (Theorem~\ref{r1-thm-views}(c)),
\[
P'AQ=\Delta Q'AQ=\diag(\delta_1d_1,\dots,\delta_nd_n)=\begin{pmatrix}D&0\\0&0\end{pmatrix},\qquad
D=\diag(\lvert d_1\rvert,\dots,\lvert d_r\rvert)
\]
\item As \(PP'=I_n\) and \(QQ'=I_n\) (Proposition~\ref{r6-prop-orthonormal}(b)), \(A=P(P'AQ)Q'\) is a singular value
decomposition of \(A\), with the singular values \(\sigma_i=\lvert d_i\rvert\) (\(i\leqslant r\)).
\item By Theorem~\ref{r12-thm-eckart-young}(a),
\(\lVert A-B\rVert_F^2\geqslant\sigma_{k+1}^2+\dots+\sigma_r^2=d_{k+1}^2+\dots+d_r^2\), with equality for
\(B=\sum_{i\leqslant k}\sigma_ip_iq_i'=\sum_{i\leqslant k}\lvert d_i\rvert\delta_iq_iq_i'=\sum_{i\leqslant k}d_iq_iq_i'\), as
\(\lvert d_i\rvert\delta_i=d_i\).
\item If \(A\) is PSD, then every \(d_i\geqslant0\) (Theorem~\ref{r8-thm-eigen-criterion}(a)), so the nonzero ones are
positive and \(\lvert d_i\rvert=d_i\).\qedhere
\end{itemize}
\end{proof}

\subsection*{The SVD in statistics}

\begin{proposition}[\citeat{tu-asm}{Proposition~5.23}]\label{r12-prop-svd}
Let \(X\in\R^{n\times p}\) be a data matrix with \(r=\rank(C_nX)\geqslant1\).
\begin{enumerate}
\item Let \(C_nX=P\begin{pmatrix}D&0\\0&0\end{pmatrix}Q'=\sigma_1p_1q_1'+\dots+\sigma_rp_rq_r'\) be a singular value
decomposition (Theorem~\ref{r9-thm-svd}). Then \(S=\frac1{n-1}Q\diag\bigl(\sigma_1^2,\dots,\sigma_r^2,0,\dots,0\bigr)Q'\).
Hence \(\lambda_m=\sigma_m^2/(n-1)\) for \(m\leqslant r\) and \(\lambda_m=0\) for \(m>r\); the columns \(q_1,\dots,q_p\) of
\(Q\) can serve as \(v_1,\dots,v_p\) in Theorem~\ref{r12-thm-pca}; with this choice \(z_m=\sigma_mp_m\) for \(m\leqslant r\);
and \(\mathrm{PVE}_m=\sigma_m^2/(\sigma_1^2+\dots+\sigma_r^2)\).
\item Conversely, let \(v_m\), \(\lambda_m\), \(z_m\) be as in Theorem~\ref{r12-thm-pca}. Then
\(C_nX=\sigma_1p_1v_1'+\dots+\sigma_rp_rv_r'\) with \(\sigma_m=\lVert z_m\rVert=\sqrt{(n-1)\lambda_m}\) and
\(p_m=z_m/\sigma_m\). The vectors \(p_1,\dots,p_r\) extend to an orthogonal \(P\in\R^{n\times n}\), and
\(C_nX=P\left(\begin{smallmatrix}D&0\\0&0\end{smallmatrix}\right)V'\), \(D=\diag(\sigma_1,\dots,\sigma_r)\), is a singular
value decomposition, with the sum above as its compact SVD.
\end{enumerate}
\end{proposition}
\needspace{4\baselineskip}
\begin{proof}
(a)
\begin{itemize}
\item \(P'P=I_n\), so \((C_nX)'(C_nX)=Q\diag(\sigma_1^2,\dots,\sigma_r^2,0,\dots,0)Q'\) (Theorem~\ref{r5b-thm-block}(c)); divide
by \(n-1\) (Proposition~\ref{r12-prop-data}(a)).
\item So \(Q\) is orthogonal and its columns are orthonormal eigenvectors of \(S\) for
\(\sigma_1^2/(n-1)\geqslant\dots\geqslant\sigma_r^2/(n-1)>0=\dots=0\). By Theorem~\ref{r8-thm-spectral}(c) these are the
eigenvalues \(\lambda_1\geqslant\dots\geqslant\lambda_p\) of \(S\), and \(V=Q\) is allowed in Theorem~\ref{r12-thm-pca}.
\item \(z_m=C_nXq_m=\sum_i\sigma_ip_i(q_i'q_m)=\sigma_mp_m\), as \(q_i'q_m=0\) for \(i\ne m\) and \(q_m'q_m=1\).
\item By Proposition~\ref{r12-prop-pve}, \(\mathrm{PVE}_m=\lambda_m/\sum_l\lambda_l=\sigma_m^2/\sum_l\sigma_l^2\).
\end{itemize}
(b)
\begin{itemize}
\item \(VV'=I_p\) (Proposition~\ref{r6-prop-orthonormal}(b)), so \(C_nX=(C_nXV)V'=\sum_{m=1}^pz_mv_m'\)
(Theorem~\ref{r1-thm-views}(d)), and \(z_m=0\) for \(m>r\), as \(\lVert z_m\rVert^2=(n-1)\lambda_m=0\)
(Theorem~\ref{r12-thm-pca}(a),(d)).
\item \((C_nX)'(C_nX)=(n-1)S=V\diag\bigl((n-1)\lambda_1,\dots,(n-1)\lambda_p\bigr)V'\) (Proposition~\ref{r12-prop-data}(a)),
with \((n-1)\lambda_1\geqslant\dots\geqslant(n-1)\lambda_p\), and its nonzero eigenvalues are
\(\sigma_1^2\geqslant\dots\geqslant\sigma_r^2\) (Theorem~\ref{r12-thm-pca}(a)). So \(Q=V\) and \(D\) are as in
Theorem~\ref{r9-thm-svd}(b) for \(A=C_nX\), and there \(P_1=C_nX[v_1\ \cdots\ v_r]D^{-1}=[p_1\ \cdots\ p_r]\), as
\(C_nXv_m/\sigma_m=z_m/\sigma_m=p_m\). Theorem~\ref{r9-thm-svd}(b) extends \(P_1\) to an orthogonal \(P\) and gives
\(C_nX=P\left(\begin{smallmatrix}D&0\\0&0\end{smallmatrix}\right)V'=\sum_{m\leqslant r}\sigma_mp_mv_m'\).\qedhere
\end{itemize}
\end{proof}
For \(A=C_nX\) and \(1\leqslant k<r\), the choice \(v_m=q_m\) in Proposition~\ref{r12-prop-svd}(a) gives, in
Theorem~\ref{r12-thm-eckart-young}, \(A_k=\sigma_1p_1q_1'+\dots+\sigma_kp_kq_k'=z_1v_1'+\dots+z_kv_k'=Z_kV_k'\)
(Theorem~\ref{r1-thm-views}(d); \(Z_k\) and \(V_k\) as in Theorem~\ref{r12-thm-closest} with \(M=k\)) and
\(\sigma_{k+1}^2+\dots+\sigma_r^2=(n-1)(\lambda_{k+1}+\dots+\lambda_p)\). So Theorem~\ref{r12-thm-eckart-young}(a) for
\(C_nX\) is Theorem~\ref{r12-thm-closest}(a) again, in the form of the paragraph after it (a best approximation among the
matrices of rank at most \(M=k\)): the same bound, attained at the same matrix \(Z_kV_k'\).

\begin{example}\label{r12-exm-svd}
For the data of Example~\ref{r12-exm-data}, Example~\ref{r12-exm-pca} and Proposition~\ref{r12-prop-svd}(b) give
\(\sigma_1=\sqrt{4\cdot25}=10\), \(\sigma_2=\sqrt{4\cdot9}=6\), \(\sigma_3=\sqrt{4\cdot5}=2\sqrt5\), \(\sigma_4=\sqrt{4\cdot1}=2\) and
\[
p_1=\frac{z_1}{\sigma_1}=\frac12\begin{pmatrix}0\\1\\1\\-1\\-1\end{pmatrix},\qquad
p_2=\frac{z_2}{\sigma_2}=\frac12\begin{pmatrix}0\\1\\-1\\1\\-1\end{pmatrix},\qquad
p_3=\frac{z_3}{\sigma_3}=\frac1{2\sqrt5}\begin{pmatrix}4\\-1\\-1\\-1\\-1\end{pmatrix},\qquad
p_4=\frac{z_4}{\sigma_4}=\frac12\begin{pmatrix}0\\1\\-1\\-1\\1\end{pmatrix}
\]
so \(C_5X=\sigma_1p_1v_1'+\sigma_2p_2v_2'+\sigma_3p_3v_3'+\sigma_4p_4v_4'\), with the loading vectors \(v_m=\frac12h_m\) of
Example~\ref{r12-exm-pca}. For instance \(\sigma_1p_1v_1'=z_1v_1'=\frac12(0,5,5,-5,-5)'(1,1,1,1)\). Check:
\(\sigma_1^2+\sigma_2^2+\sigma_3^2+\sigma_4^2=100+36+20+4=160\), which is \(\lVert C_5X\rVert_F^2=\tr\bigl((C_5X)'(C_5X)\bigr)=4\cdot40\)
(Theorem~\ref{r9-thm-svd-props}(d)).
\end{example}

\begin{theorem}[\citeat{harville}{\S21.12b}, \citeat{abadir-magnus}{Exercise~10.23}]\label{r10-thm-mp-svd}
Let \(A\in\R^{m\times n}\) have rank \(r\geqslant1\), and let \(A=P_1DQ_1'\), where \(P_1\in\R^{m\times r}\) and
\(Q_1\in\R^{n\times r}\) have orthonormal columns (\(P_1'P_1=Q_1'Q_1=I_r\)) and \(D=\diag(\sigma_1,\dots,\sigma_r)\) with
\(\sigma_1,\dots,\sigma_r>0\) (a compact SVD, Theorem~\ref{r9-thm-svd}). Then
\[
A^+=Q_1D^{-1}P_1'=\sigma_1^{-1}q_1p_1'+\dots+\sigma_r^{-1}q_rp_r'
\]
where \(p_i\) and \(q_i\) are the columns of \(P_1\) and \(Q_1\); and \(AA^+=P_1P_1'\), \(A^+A=Q_1Q_1'\). If moreover \(P=[P_1\ P_2]\in\R^{m\times m}\) and
\(Q=[Q_1\ Q_2]\in\R^{n\times n}\) are orthogonal, so that
\(A=P\left(\begin{smallmatrix}D&0\\0&0\end{smallmatrix}\right)Q'\), then
\(A^+=Q\left(\begin{smallmatrix}D^{-1}&0\\0&0\end{smallmatrix}\right)P'\).
\end{theorem}
\begin{proof}
Put \(G=Q_1D^{-1}P_1'\).
\begin{itemize}
\item \(AG=P_1DQ_1'Q_1D^{-1}P_1'=P_1P_1'\) and \(GA=Q_1D^{-1}P_1'P_1DQ_1'=Q_1Q_1'\) are symmetric: (P3), (P4) of Definition~\ref{r10-def-mp}.
\item \(AGA=P_1P_1'P_1DQ_1'=P_1DQ_1'=A\) and \(GAG=Q_1Q_1'Q_1D^{-1}P_1'=G\): (P1), (P2).
\item So \(G=A^+\) (Theorem~\ref{r10-thm-mp-unique}), and \(AA^+=AG=P_1P_1'\), \(A^+A=GA=Q_1Q_1'\). Row \(i\) of \(D^{-1}P_1'\) is \(\sigma_i^{-1}p_i'\), so the outer product
expansion (Theorem~\ref{r1-thm-views}(d)) gives the sum.
\item By the block product (Theorem~\ref{r5b-thm-block}),
\(Q\left(\begin{smallmatrix}D^{-1}&0\\0&0\end{smallmatrix}\right)P'
=[Q_1\ Q_2]\left(\begin{smallmatrix}D^{-1}P_1'\\0\end{smallmatrix}\right)=Q_1D^{-1}P_1'\).\qedhere
\end{itemize}
\end{proof}

\begin{definition}[\citeat{bapat}{\S10.2}]\label{r12-def-cca}
Let the \(p=p_1+p_2\) variables of a data matrix \(X\) be split into two groups, \(X=[X_1\ X_2]\) with
\(X_1\in\R^{n\times p_1}\), and partition \(S\) accordingly:
\[
S=\begin{pmatrix}S_{11}&S_{12}\\S_{21}&S_{22}\end{pmatrix},\qquad S_{11}\in\R^{p_1\times p_1},\quad S_{21}=S_{12}'
\]
Assume that \(S_{11}\) and \(S_{22}\) are PD. For \(a\in\R^{p_1}\) and
\(b\in\R^{p_2}\), Proposition~\ref{r12-prop-data}(b), applied to \((a',0')'\) and \((0',b')'\), shows that \(X_1a\) and
\(X_2b\) have sample variances \(a'S_{11}a\), \(b'S_{22}b\) and sample covariance \(a'S_{12}b\)
(Theorem~\ref{r5b-thm-block}(c)).
\begin{itemize}
\item A \textbf{first canonical pair} is a pair \((a_1,b_1)\) that maximizes \(a'S_{12}b\) subject to
\(a'S_{11}a=b'S_{22}b=1\); the maximum \(\rho_1\) is the \textbf{first canonical correlation}.
\item For \(k\geqslant2\), once \((a_1,b_1),\dots,(a_{k-1},b_{k-1})\) are chosen, a \textbf{\(k\)th canonical pair}
maximizes \(a'S_{12}b\) subject to \(a'S_{11}a=b'S_{22}b=1\) and \(a'S_{11}a_i=b'S_{22}b_i=0\) for \(i<k\); the maximum
\(\rho_k\) is the \textbf{\(k\)th canonical correlation}.
\end{itemize}
\end{definition}
When both variances are \(1\), \(a'S_{12}b\) is the sample correlation (Definition~\ref{r12-def-correlation}) \(a'S_{12}b/(a'S_{11}a\cdot b'S_{22}b)^{1/2}\) of
\(X_1a\) and \(X_2b\); the side conditions make \(X_1a\) uncorrelated with \(X_1a_i\), and \(X_2b\) with \(X_2b_i\) (\(i<k\)).

\begin{theorem}[\citeat{bapat}{\S10.2}]\label{r12-thm-cca}
In Definition~\ref{r12-def-cca}, let \(K=S_{11}^{-1/2}S_{12}S_{22}^{-1/2}\in\R^{p_1\times p_2}\) (Theorem~\ref{r8-thm-sqrt})
have rank \(s\geqslant1\) and compact SVD \(K=\rho_1u_1w_1'+\dots+\rho_su_sw_s'\),
\(\rho_1\geqslant\dots\geqslant\rho_s>0\) (Theorem~\ref{r9-thm-svd}), and put \(a_k=S_{11}^{-1/2}u_k\),
\(b_k=S_{22}^{-1/2}w_k\). For \(k=1,\dots,s\): if \((a_i,b_i)\) (\(i<k\)) are taken as the first \(k-1\) canonical pairs,
then \((a_k,b_k)\) is a \(k\)th canonical pair, and the \(k\)th canonical correlation is \(\rho_k\).
\end{theorem}
\begin{proof}
\begin{itemize}
\item \(S_{11}^{1/2}\) is symmetric and invertible, with inverse \(S_{11}^{-1/2}\) (Theorem~\ref{r8-thm-sqrt}); so
\(a\mapsto u=S_{11}^{1/2}a\) maps \(\R^{p_1}\) onto \(\R^{p_1}\), with inverse \(u\mapsto S_{11}^{-1/2}u\). Likewise
\(b\mapsto w=S_{22}^{1/2}b\), and \(S_{11}^{1/2}a_i=u_i\), \(S_{22}^{1/2}b_i=w_i\).
\item With \(u=S_{11}^{1/2}a\) and \(w=S_{22}^{1/2}b\): \(a'S_{11}a=a'S_{11}^{1/2}S_{11}^{1/2}a=u'u\),
\(b'S_{22}b=w'w\), \(a'S_{12}b=u'S_{11}^{-1/2}S_{12}S_{22}^{-1/2}w=u'Kw\), \(a'S_{11}a_i=u'u_i\) and \(b'S_{22}b_i=w'w_i\).
\item So the \(k\)th problem is to maximize \(u'Kw\) over unit vectors \(u\in\R^{p_1}\), \(w\in\R^{p_2}\) with
\(u\perp u_1,\dots,u_{k-1}\) and \(w\perp w_1,\dots,w_{k-1}\). By Theorem~\ref{r9-thm-sigma-max}(a),(b), the maximum is \(\rho_k\),
attained at \(u=u_k\), \(w=w_k\), that is, at \(a=a_k\), \(b=b_k\).\qedhere
\end{itemize}
\end{proof}

\begin{example}\label{r12-exm-cca}
For the data of Example~\ref{r12-exm-data}, take variables \(1\), \(2\) as the first group and variables \(3\), \(4\) as
the second:
\[
S_{11}=S_{22}=\begin{pmatrix}10&7\\7&10\end{pmatrix},\qquad S_{12}=\begin{pmatrix}5&3\\3&5\end{pmatrix}
\]
\begin{itemize}
\item \(I_2=\bar J_2+C_2\), where \(\bar J_2\) and \(C_2\) are symmetric and idempotent and
\(\bar J_2C_2=\bar J_2-\bar J_2^2=0\) (Definition~\ref{r1-def-ones}, Proposition~\ref{r1-prop-centering}); also
\(C_2\bar J_2=(\bar J_2C_2)'=0\). As \(\bar J_2=\frac12\left(\begin{smallmatrix}1&1\\1&1\end{smallmatrix}\right)\) and
\(C_2=\frac12\left(\begin{smallmatrix}1&-1\\-1&1\end{smallmatrix}\right)\),
\[
S_{11}=S_{22}=17\bar J_2+3C_2,\qquad S_{12}=8\bar J_2+2C_2
\]
\item So \((\sqrt{17}\bar J_2+\sqrt3C_2)^2=17\bar J_2+3C_2=S_{11}\), and \(\sqrt{17}\bar J_2+\sqrt3C_2\) is PSD, as
\(x'(\sqrt{17}\bar J_2+\sqrt3C_2)x=\sqrt{17}\lVert\bar J_2x\rVert^2+\sqrt3\lVert C_2x\rVert^2\). By
Theorem~\ref{r8-thm-sqrt}, \(S_{11}^{1/2}=S_{22}^{1/2}=\sqrt{17}\bar J_2+\sqrt3C_2\), and
\(S_{11}^{-1/2}=S_{22}^{-1/2}=\frac1{\sqrt{17}}\bar J_2+\frac1{\sqrt3}C_2\), since the product of the two is
\(\bar J_2+C_2=I_2\).
\item Multiplying out, with \(\bar J_2^2=\bar J_2\), \(C_2^2=C_2\) and \(\bar J_2C_2=C_2\bar J_2=0\),
\[
K=S_{11}^{-1/2}S_{12}S_{22}^{-1/2}=\frac8{17}\bar J_2+\frac23C_2=\frac1{51}\begin{pmatrix}29&-5\\-5&29\end{pmatrix}
\]
With \(u_1=w_1=\frac1{\sqrt2}(1,-1)'\) and \(u_2=w_2=\frac1{\sqrt2}(1,1)'\), \(C_2=u_1u_1'\) and \(\bar J_2=u_2u_2'\), so
\(K=\frac23u_1w_1'+\frac8{17}u_2w_2'\). As \(\frac23>\frac8{17}>0\) and \(u_1,u_2\) are orthonormal, this is a compact
SVD with \(s=2\), \(\rho_1=\frac23\) and \(\rho_2=\frac8{17}\).
\item By Theorem~\ref{r12-thm-cca}, the canonical correlations are \(\frac23\) and \(\frac8{17}\), with
\[
a_1=b_1=S_{11}^{-1/2}u_1=\frac1{\sqrt3}u_1=\frac1{\sqrt6}\begin{pmatrix}1\\-1\end{pmatrix},\qquad
a_2=b_2=S_{11}^{-1/2}u_2=\frac1{\sqrt{17}}u_2=\frac1{\sqrt{34}}\begin{pmatrix}1\\1\end{pmatrix}
\]
as \(C_2u_1=u_1\), \(\bar J_2u_1=0\), \(\bar J_2u_2=u_2\) and \(C_2u_2=0\). The first pair contrasts variable \(1\) with
variable \(2\), and variable \(3\) with variable \(4\); the second pair adds them. Check:
\(a_1'S_{11}a_1=\frac16\cdot6=1\), \(a_1'S_{12}b_1=\frac16\cdot4=\frac23\), \(a_2'S_{11}a_2=\frac1{34}\cdot34=1\),
\(a_2'S_{12}b_2=\frac1{34}\cdot16=\frac8{17}\) and \(a_2'S_{11}a_1=0\).
\end{itemize}
\end{example}

\needspace{10\baselineskip}
\begin{remark}[\citeat{bapat}{\S10.2}]\label{r12-rem-pop-cca}
Canonical correlations are also defined for two random vectors \(x\in\R^{p_1}\) and \(y\in\R^{p_2}\). Let
\[
\var\begin{pmatrix}x\\y\end{pmatrix}=\Sigma=\begin{pmatrix}\Sigma_{11}&\Sigma_{12}\\\Sigma_{21}&\Sigma_{22}\end{pmatrix},\qquad
\Sigma_{11}\in\R^{p_1\times p_1}
\]
with \(\Sigma_{11}=\var(x)\) and \(\Sigma_{22}=\var(y)\) PD (Definition~\ref{r11-def-random}). For
\(a\in\R^{p_1}\) and \(b\in\R^{p_2}\), the matrix
\(T=\left(\begin{smallmatrix}a'&0'\\0'&b'\end{smallmatrix}\right)\in\R^{2\times(p_1+p_2)}\) maps \((x',y')'\) to
\((a'x,b'y)'\), and by Theorems~\ref{r11-thm-linear}(b) and~\ref{r5b-thm-block}(b),(c)
\[
\var\begin{pmatrix}a'x\\b'y\end{pmatrix}=T\Sigma T'
=\begin{pmatrix}a'\Sigma_{11}a&a'\Sigma_{12}b\\b'\Sigma_{21}a&b'\Sigma_{22}b\end{pmatrix}
\]
So \(a'x\) and \(b'y\) have the variances \(a'\Sigma_{11}a\), \(b'\Sigma_{22}b\) and the covariance \(a'\Sigma_{12}b\); and, by
Theorem~\ref{r12-thm-pop-pca}(a) for \(x\) and for \(y\), \(\cov(a'x,a_i'x)=a'\Sigma_{11}a_i\) and
\(\cov(b'y,b_i'y)=b'\Sigma_{22}b_i\). These are the expressions of Definition~\ref{r12-def-cca} with \(\Sigma_{jl}\) in
place of \(S_{jl}\). Canonical pairs \((a_k,b_k)\) and canonical
correlations of \(x\) and \(y\) are defined as there, with the variances and covariances of \(a'x\) and \(b'y\). The proof of
Theorem~\ref{r12-thm-cca} uses only that \(S_{11}\) and \(S_{22}\) are PD. Hence the canonical correlations of \(x\) and \(y\)
are the singular values \(\rho_1\geqslant\dots\geqslant\rho_s>0\) of \(\Sigma_{11}^{-1/2}\Sigma_{12}\Sigma_{22}^{-1/2}\), attained at
\(a_k=\Sigma_{11}^{-1/2}u_k\), \(b_k=\Sigma_{22}^{-1/2}w_k\) for its singular vectors \(u_k\), \(w_k\).
With \(K=\Sigma_{11}^{-1/2}\Sigma_{12}\Sigma_{22}^{-1/2}\), \(KK'=\Sigma_{11}^{-1/2}B\Sigma_{11}^{-1/2}\) for the symmetric
\(B=\Sigma_{12}\Sigma_{22}^{-1}\Sigma_{21}\), so \(\rho_1^2,\dots,\rho_s^2\) are the nonzero eigenvalues of \(KK'\)
(Theorem~\ref{r9-thm-svd-props}(a)) and hence of \(\Sigma_{11}^{-1}\Sigma_{12}\Sigma_{22}^{-1}\Sigma_{21}\)
(Corollary~\ref{r8-cor-real-eigen}(b) with \(A=\Sigma_{11}\)), the matrix named after Corollary~\ref{r8-cor-real-eigen}.
If \(p_1=1\) and \(\Sigma_{12}\ne0\), then \(s=1\) and \(\rho_1^2=\Sigma_{12}\Sigma_{22}^{-1}\Sigma_{21}/\Sigma_{11}\): the canonical
correlation of the random variable \(x\) and \(y\) is the multiple correlation coefficient \(r_{x(y)}\) of
Proposition~\ref{r9-prop-multiple-correlation}, with \(\Sigma_{11}\), \(\Sigma_{21}\), \(\Sigma_{22}\) in the roles of
\(\sigma^2\), \(u\), \(\Sigma\).
\end{remark}

\subsection*{The condition number}

\begin{definition}[\citeat{tu-asm}{Definition~P.52}]\label{r9-def-condition}
The \textbf{condition number} of an invertible \(A\in\R^{n\times n}\) is \(\kappa(A)=\lVert A\rVert_2\lVert A^{-1}\rVert_2\).
For \(X\in\R^{n\times p}\) of full column rank with \(p<n\), with singular values \(\sigma_1\geqslant\dots\geqslant\sigma_p>0\),
put \(\kappa(X)=\sigma_1/\sigma_p\) (for \(p=n\) the first definition gives the same number, by
Proposition~\ref{r9-prop-condition}(a)).
\end{definition}

\begin{proposition}[\citeat{tu-asm}{Proposition~P.53, Corollary~P.55}]\label{r9-prop-condition}
\begin{enumerate}
\item Let \(A\in\R^{n\times n}\) be invertible. Then \(A\) has \(n\) singular values \(\sigma_1\geqslant\dots\geqslant\sigma_n>0\),
\(\lVert A^{-1}\rVert_2=1/\sigma_n\), and \(\kappa(A)=\sigma_1/\sigma_n\geqslant1\). If \(A\) is PD, with eigenvalues
\(\lambda_1\geqslant\dots\geqslant\lambda_n>0\), then \(\kappa(A)=\lambda_1/\lambda_n\).
\item If \(X\in\R^{n\times p}\) has full column rank, then \(X'X\) is PD and \(\kappa(X'X)=\kappa(X)^2\).
\end{enumerate}
\end{proposition}
\begin{proof}
(a)
\begin{itemize}
\item \(\rank A=n\) (Theorem~\ref{thm:r4-rank-rules}(d)), so \(A=PDQ'\) with \(P\), \(Q\) orthogonal and
\(D=\diag(\sigma_1,\dots,\sigma_n)\), \(\sigma_1\geqslant\dots\geqslant\sigma_n>0\) (Theorem~\ref{r9-thm-svd}(b), with \(r=n\)).
\item \(P^{-1}=P'\) and \(Q^{-1}=Q'\) (Proposition~\ref{r6-prop-orthonormal}(b)), so \(A^{-1}=QD^{-1}P'\)
(Theorem~\ref{r2-thm-inverse-props}(b)), and \(\lVert A^{-1}\rVert_2=\lVert D^{-1}\rVert_2\)
(Proposition~\ref{r9-prop-spectral-norm}(c)).
\item For \(x\in\R^n\), \(\lVert D^{-1}x\rVert^2=\sum_ix_i^2/\sigma_i^2\leqslant\lVert x\rVert^2/\sigma_n^2\), with equality for
\(x=e_n\); so \(\lVert D^{-1}\rVert_2=1/\sigma_n\) (Proposition~\ref{r9-prop-spectral-norm}(a)).
\item Hence \(\kappa(A)=\sigma_1\cdot\frac1{\sigma_n}\geqslant1\) (Definition~\ref{r9-def-spectral-norm}). If \(A\) is PD, its
singular values are its eigenvalues (Theorem~\ref{r9-thm-svd-props}(e)), so \(\kappa(A)=\lambda_1/\lambda_n\).
\end{itemize}
(b) \(X'X\) is PD (Theorem~\ref{r8-thm-XtX}). Its eigenvalues are \(\sigma_1^2\geqslant\dots\geqslant\sigma_p^2>0\)
(Theorem~\ref{r9-thm-svd-props}(a), with \(r=p\)), so by (a), \(\kappa(X'X)=\sigma_1^2/\sigma_p^2=\kappa(X)^2\).
\end{proof}

\begin{theorem}[\citeat{tu-asm}{Theorem~P.54}]\label{r9-thm-conditioning}
Let \(A\in\R^{n\times n}\) be invertible, \(Ax=b\ne0\), and \(A(x+\delta x)=b+\delta b\). Then
\[
\frac{\lVert\delta x\rVert}{\lVert x\rVert}\leqslant\kappa(A)\,\frac{\lVert\delta b\rVert}{\lVert b\rVert}
\]
\end{theorem}
\begin{proof}
\begin{itemize}
\item \(\lVert My\rVert\leqslant\lVert M\rVert_2\lVert y\rVert\) for every matrix \(M\) and vector \(y\)
(Proposition~\ref{r9-prop-spectral-norm}(a)).
\item Subtracting \(Ax=b\) from \(A(x+\delta x)=b+\delta b\) gives \(A\,\delta x=\delta b\), so \(\delta x=A^{-1}\delta b\) and
\(\lVert\delta x\rVert\leqslant\lVert A^{-1}\rVert_2\lVert\delta b\rVert\).
\item \(x\ne0\), as \(b\ne0\); and \(\lVert b\rVert=\lVert Ax\rVert\leqslant\lVert A\rVert_2\lVert x\rVert\), so
\(1/\lVert x\rVert\leqslant\lVert A\rVert_2/\lVert b\rVert\).
\item Multiplying the two bounds, both sides of each being \(\geqslant0\),
\(\lVert\delta x\rVert/\lVert x\rVert\leqslant\lVert A^{-1}\rVert_2\lVert A\rVert_2\,\lVert\delta b\rVert/\lVert b\rVert\).\qedhere
\end{itemize}
\end{proof}
A large \(\kappa(A)\) means that a small relative change of \(b\) may change the solution of \(Ax=b\) by a large relative
amount. By Proposition~\ref{r9-prop-condition}(b), forming \(X'X\) squares the condition number of \(X\). A QR decomposition
avoids this: if \(X\) has full column rank and \(X=QR\) (Theorem~\ref{r6-thm-qr}), then \(X'X=R'Q'QR=R'R\), so the normal
equations \(X'Xb=X'y\) (Theorem~\ref{r10-thm-normal-equations}) read
\(R'Rb=R'Q'y\), that is, \(Rb=Q'y\), as \(R'\) is invertible (Theorem~\ref{r2-thm-triangular}(a)); and
\(R'R=X'X\) shows that \(R\) has the same singular values as \(X\) (Theorem~\ref{r9-thm-svd-props}(a)), so
\(\kappa(R)=\kappa(X)\) (Proposition~\ref{r9-prop-condition}(a)). This is why least-squares estimates
(Theorem~\ref{r10-thm-least-squares}) are computed in practice from \(Rb=Q'y\) rather than from \(X'X\).

A large \(\kappa(X)\) also means that the columns of \(X\) are nearly linearly dependent, a near version of the exact
collinearity of Proposition~\ref{r11-prop-collinearity}. Indeed, let \(X\in\R^{n\times p}\) have full column rank
\(p\geqslant2\). A matrix \(B\in\R^{n\times p}\) has linearly dependent columns if and only if \(\rank B<p\)
(Theorem~\ref{thm:r4-rank-rules}(b)), and by Theorem~\ref{r12-thm-eckart-young}(b) for \(A=X\) (so \(r=p\)) and
\(k=p-1\), every such \(B\) has \(\lVert X-B\rVert_2\geqslant\sigma_p\), with equality for the matrix \(B=A_{p-1}\) of that
theorem. So \(\sigma_p\) is the smallest spectral norm \(\lVert E\rVert_2\) of a change \(E\) for which \(X+E\) has
linearly dependent columns, and relative to \(\lVert X\rVert_2=\sigma_1\) this distance is \(\sigma_p/\sigma_1=1/\kappa(X)\).

\begin{example}\label{r9-exm-condition}
The matrix
\[
A=\begin{pmatrix}2&0&1&-1\\0&2&1&1\\1&1&3&0\\-1&1&0&3\end{pmatrix}
\]
of Example~\ref{r8-exm-leading} is PD, with the eigenvalues \(4\), \(4\), \(1\), \(1\) and the eigenvectors \((1,1,2,0)'\),
\((-1,1,0,2)'\) for \(4\) and \((1,1,-1,0)'\), \((1,-1,0,1)'\) for \(1\) (Example~\ref{r8-exm-leading}). So
\(\kappa(A)=4/1=4\) (Proposition~\ref{r9-prop-condition}(a)). Take
\(x=(1,1,2,0)'\), so that \(b=Ax=4x=(4,4,8,0)'\), and \(\delta b=(1,1,-1,0)'\). As \(A\,\delta b=\delta b\), the change of
the solution is \(\delta x=A^{-1}\delta b=\delta b=(1,1,-1,0)'\). Then
\[
\frac{\lVert\delta x\rVert}{\lVert x\rVert}=\frac{\sqrt3}{\sqrt6},\qquad
\frac{\lVert\delta b\rVert}{\lVert b\rVert}=\frac{\sqrt3}{\sqrt{96}}=\frac14\cdot\frac{\sqrt3}{\sqrt6}
\]
so the bound of Theorem~\ref{r9-thm-conditioning} holds with equality: the relative change of \(b\) is magnified four
times in the solution.
\end{example}

\subsection*{The orthogonal Procrustes problem}

\begin{theorem}[Orthogonal Procrustes; \citeat{horn-johnson}{\S7.4.5}, \citeat{prasolov}{Problem~34.2}]\label{r10-thm-procrustes}
Let \(M\in\R^{n\times n}\), with an SVD \(M=P\left(\begin{smallmatrix}D&0\\0&0\end{smallmatrix}\right)Q'\)
(Theorem~\ref{r9-thm-svd}; for \(M=0\) take \(P=Q=I_n\)) and the singular values \(\sigma_1,\dots,\sigma_n\), where
\(\sigma_j=0\) for \(j>\rank M\).
\begin{enumerate}
\item For every orthogonal \(U\in\R^{n\times n}\), \(\tr(U'M)\leqslant\sigma_1+\dots+\sigma_n\), with equality for \(U=PQ'\).
\item Let \(B,C\in\R^{n\times k}\) and \(M=CB'\). For every orthogonal \(U\in\R^{n\times n}\),
\[
\lVert UB-C\rVert_F^2\geqslant\lVert B\rVert_F^2+\lVert C\rVert_F^2-2(\sigma_1+\dots+\sigma_n)
\]
with equality for \(U=PQ'\).
\end{enumerate}
\end{theorem}
\begin{proof}
(a)
\begin{itemize}
\item \(\Delta=\left(\begin{smallmatrix}D&0\\0&0\end{smallmatrix}\right)=\diag(\sigma_1,\dots,\sigma_n)\), and \(P\), \(Q\), \(U\)
are orthogonal: \(P'P=PP'=I_n\), and likewise for \(Q\), \(U\) (Proposition~\ref{r6-prop-orthonormal}(b)). So
\(Z=Q'U'P\) is orthogonal, as \(Z'Z=P'UQQ'U'P=I_n\).
\item By Theorem~\ref{r1-thm-trace}(b),
\[
\tr(U'M)=\tr(U'P\Delta Q')=\tr(Q'U'P\Delta)=\tr(Z\Delta)=z_{11}\sigma_1+\dots+z_{nn}\sigma_n
\]
since \(\Delta\) is diagonal.
\item Column \(j\) of \(Z\) is a unit vector (\(Z'Z=I_n\)), so \(z_{jj}\leqslant1\); and \(\sigma_j\geqslant0\). Hence
\(\tr(U'M)\leqslant\sigma_1+\dots+\sigma_n\).
\item For \(U=PQ'\), \(Z=Q'QP'P=I_n\), and equality holds.
\end{itemize}
(b)
\begin{itemize}
\item \((UB)'(UB)=B'U'UB=B'B\), and \(\tr(C'UB)=\tr\bigl((C'UB)'\bigr)=\tr(B'U'C)=\tr(U'CB')\)
(Theorem~\ref{r1-thm-trace}(a),(b)). So
\[
\lVert UB-C\rVert_F^2=\tr\bigl((UB-C)'(UB-C)\bigr)=\tr(B'B)-2\tr(U'CB')+\tr(C'C)
\]
\item By Definition~\ref{r9-def-frobenius}, \(\tr(B'B)=\lVert B\rVert_F^2\) and \(\tr(C'C)=\lVert C\rVert_F^2\); and (a) for
\(M=CB'\) gives \(\tr(U'CB')\leqslant\sigma_1+\dots+\sigma_n\), with equality for \(U=PQ'\).\qedhere
\end{itemize}
\end{proof}
In (b) the columns of \(B\) and \(C\) are two configurations of \(k\) points in \(\R^n\), and \(PQ'\) is an orthogonal
transformation (a rotation, possibly composed with a reflection) that moves the first configuration as close as possible
to the second; this is the \textbf{orthogonal Procrustes problem} of statistics (Procrustes analysis). For \(B=I_n\), (b)
says that \(PQ'\) is an orthogonal matrix nearest to \(C=M\) in the Frobenius norm.

\subsection*{Exercises}

\begin{exercise}\label{r9-ex-gen-rayleigh-drill}
Let
\[
A=\begin{pmatrix}2&4&0&2\\4&8&2&0\\0&2&2&4\\2&0&4&8\end{pmatrix},\qquad
B=\begin{pmatrix}1&0&0&0\\0&4&0&0\\0&0&1&0\\0&0&0&4\end{pmatrix},\qquad
y=\begin{pmatrix}2\\4\\1\\-4\end{pmatrix}
\]
and \(h_1=(1,1,1,1)'\), \(h_2=(1,-1,1,-1)'\), \(h_3=(1,1,-1,-1)'\), \(h_4=(1,-1,-1,1)'\).
\begin{enumerate}
\item Find \(C=B^{-1/2}AB^{-1/2}\), and check that \(h_1,\dots,h_4\) are eigenvectors of \(C\). Find the maximum and the
minimum of \(x'Ax/x'Bx\) over \(x\ne0\), and vectors at which they are attained.
\item Find the maximum of \((x'y)^2/x'Bx\) over \(x\ne0\), and a vector at which it is attained.
\end{enumerate}
\end{exercise}
\begin{solution}
(a)
\begin{itemize}
\item \(B\) is PD (its eigenvalues \(1\), \(4\), \(1\), \(4\) are positive, Theorem~\ref{r8-thm-eigen-criterion}(a)).
\(N=\diag(1,2,1,2)\) is PD with \(N^2=B\), so \(B^{1/2}=N\) (Theorem~\ref{r8-thm-sqrt}) and
\(B^{-1/2}=\diag(1,\frac12,1,\frac12)\).
\item Row \(i\) and column \(j\) of \(A\) are multiplied by the \(i\)th and the \(j\)th diagonal entry of \(B^{-1/2}\):
\[
C=B^{-1/2}AB^{-1/2}=\begin{pmatrix}2&2&0&1\\2&2&1&0\\0&1&2&2\\1&0&2&2\end{pmatrix}
\]
\item \(Ch_1=(5,5,5,5)'=5h_1\), \(Ch_2=(-1,1,-1,1)'=-h_2\), \(Ch_3=(3,3,-3,-3)'=3h_3\) and \(Ch_4=(1,-1,-1,1)'=h_4\). So the
algebraic multiplicities \(\dim\mathcal N(C-cI)\) of \(c=5,3,1,-1\) are at least \(1\) (Theorem~\ref{r8-thm-rank-eigen});
as all multiplicities add up to \(4\) (Definition~\ref{r7-def-eigenvalues}), the eigenvalues of \(C\) are \(\mu_1=5\),
\(\mu_2=3\), \(\mu_3=1\) and \(\mu_4=-1\).
\item By Theorem~\ref{r9-thm-gen-rayleigh}, the maximum is \(5\), attained at \(x=B^{-1/2}h_1=(1,\frac12,1,\frac12)'\), and so
also at \((2,1,2,1)'\); the minimum is \(-1\), attained at \(B^{-1/2}h_2=(1,-\frac12,1,-\frac12)'\) and at \((2,-1,2,-1)'\).
Check: at \(x=(2,1,2,1)'\), \(Ax=(10,20,10,20)'=5Bx\), so \(x'Ax=5x'Bx=5\cdot16=80\); at \(x=(2,-1,2,-1)'\),
\(Ax=(-2,4,-2,4)'=-Bx\), so \(x'Ax=-x'Bx=-16\).
\end{itemize}
(b) By Theorem~\ref{r9-thm-cs-max}, the maximum is \(y'B^{-1}y=4+\frac{16}4+1+\frac{16}4=13\), attained at
\(x=B^{-1}y=(2,1,1,-1)'\). Check: \((x'y)^2=(4+4+1+4)^2=13^2\) and \(x'Bx=4+4+1+4=13\).
\end{solution}

\begin{exercise}\label{r9-dr-fisher}
Let
\[
S_p=\begin{pmatrix}1&1&0&0\\1&2&1&0\\0&1&2&1\\0&0&1&2\end{pmatrix}
\]
be the pooled sample covariance matrix of four variables in two groups of units (the paragraph after
Theorem~\ref{r9-thm-cs-max}), and \(d=(1,2,1,1)'\) the difference of their mean vectors.
\begin{enumerate}
\item Show that \(S_p\) is PD.
\item Solve \(S_pa=d\) by Gauss--Jordan elimination. Find the maximum of \((a'd)^2/a'S_pa\) over \(a\ne0\) and a vector at
which it is attained.
\item Compute \((a'd)^2/a'S_pa\) for \(a=d\) and for each variable alone (\(a=e_1,e_2,e_3,e_4\)).
\end{enumerate}
\end{exercise}
\begin{solution}
(a) The leading principal minors are \(\Delta_1=1\), \(\Delta_2=2-1=1\), \(\Delta_3=1(4-1)-1(2-0)+0=1\) along row \(1\)
(Theorem~\ref{r5d-thm-expansion}), and, along row \(1\) again,
\[
\Delta_4=1\cdot\det\begin{pmatrix}2&1&0\\1&2&1\\0&1&2\end{pmatrix}-1\cdot\det\begin{pmatrix}1&1&0\\0&2&1\\0&1&2\end{pmatrix}
=\bigl(2(4-1)-1(2-0)\bigr)-1(4-1)=4-3=1
\]
All are \(>0\), so \(S_p\) is PD (Theorem~\ref{r8-thm-leading-minors}).

(b) Row-reduce the augmented matrix:
\[
\begin{gathered}
\left(\begin{array}{cccc|c}1&1&0&0&1\\1&2&1&0&2\\0&1&2&1&1\\0&0&1&2&1\end{array}\right)
\xrightarrow{R_2\leftarrow R_2-R_1}
\left(\begin{array}{cccc|c}1&1&0&0&1\\0&1&1&0&1\\0&1&2&1&1\\0&0&1&2&1\end{array}\right)
\xrightarrow{R_3\leftarrow R_3-R_2}
\left(\begin{array}{cccc|c}1&1&0&0&1\\0&1&1&0&1\\0&0&1&1&0\\0&0&1&2&1\end{array}\right)
\\
\xrightarrow{R_4\leftarrow R_4-R_3}
\left(\begin{array}{cccc|c}1&1&0&0&1\\0&1&1&0&1\\0&0&1&1&0\\0&0&0&1&1\end{array}\right)
\xrightarrow{R_3\leftarrow R_3-R_4}
\left(\begin{array}{cccc|c}1&1&0&0&1\\0&1&1&0&1\\0&0&1&0&-1\\0&0&0&1&1\end{array}\right)
\\
\xrightarrow{R_2\leftarrow R_2-R_3}
\left(\begin{array}{cccc|c}1&1&0&0&1\\0&1&0&0&2\\0&0&1&0&-1\\0&0&0&1&1\end{array}\right)
\xrightarrow{R_1\leftarrow R_1-R_2}
\left(\begin{array}{cccc|c}1&0&0&0&-1\\0&1&0&0&2\\0&0&1&0&-1\\0&0&0&1&1\end{array}\right)
\end{gathered}
\]
So \(a=S_p^{-1}d=(-1,2,-1,1)'\); check: \(S_pa=(-1+2,\ -1+4-1,\ 2-2+1,\ -1+2)'=(1,2,1,1)'\). By
Theorem~\ref{r9-thm-cs-max}, the maximum of \((a'd)^2/a'S_pa\) is \(d'S_p^{-1}d=d'a=-1+4-1+1=3\), attained at
\(a=(-1,2,-1,1)'\): there \(a'd=3\) and \(a'S_pa=a'd=3\), and \(3^2/3=3\). The best combination gives the first and the third
variable negative weights, although their mean differences are positive.

(c) For \(a=d\): \(d'd=1+4+1+1=7\), \(S_pd=(3,6,5,3)'\) and \(d'S_pd=3+12+5+3=23\); so the ratio is \(7^2/23=\frac{49}{23}\).
For \(a=e_j\) the ratio is \(d_j^2\) divided by the \(j\)th diagonal entry of \(S_p\): \(1/1=1\), \(4/2=2\), \(1/2\) and \(1/2\). All are below the maximum \(3\).
\end{solution}

\begin{exercise}[\citeat{bapat}{Ch.~8, Exercise~10}]\label{r9-ex-multiple-correlation}
The random vector \((x_1,x_2,x_3)'\) has the covariance matrix
\[
V=\begin{pmatrix}3&-1&-1\\-1&3&-1\\-1&-1&3\end{pmatrix}
\]
Find the multiple correlation coefficient between \(x_1\) and \(x_2,x_3\), and a combination \(\alpha_1x_2+\alpha_2x_3\) with
the largest correlation with \(x_1\). Check the result with the second expression of
Proposition~\ref{r9-prop-multiple-correlation}.
\end{exercise}
\begin{solution}
\begin{itemize}
\item Partition \(V\) with \(y=x_1\) and \(x=(x_2,x_3)'\): \(\sigma^2=3\), \(u=(-1,-1)'\) and
\(\Sigma=\left(\begin{smallmatrix}3&-1\\-1&3\end{smallmatrix}\right)\). \(\Sigma\) is PD: its leading principal minors are
\(3\) and \(9-1=8\) (Theorem~\ref{r8-thm-leading-minors}).
\item \(\det\Sigma=8\), so \(\Sigma^{-1}=\frac18\left(\begin{smallmatrix}3&1\\1&3\end{smallmatrix}\right)\)
(Theorem~\ref{r5d-thm-adjugate}), and \(\alpha=\Sigma^{-1}u=\frac18(-4,-4)'=-\frac12(1,1)'\).
\item \(u'\Sigma^{-1}u=u'\alpha=\frac12+\frac12=1\). By Proposition~\ref{r9-prop-multiple-correlation}, the largest squared
correlation of \(x_1\) with a combination of \(x_2\), \(x_3\) is \(u'\Sigma^{-1}u/\sigma^2=\frac13\), attained at
\(\alpha=-\frac12(1,1)'\), that is, for \(-\frac12(x_2+x_3)\) (or any nonzero multiple of it). The multiple correlation
coefficient is \(\frac1{\sqrt3}\); the covariance \(\alpha'u=1\) is positive, so \(\cor\bigl(x_1,-\frac12(x_2+x_3)\bigr)=\frac1{\sqrt3}\).
\item Second expression: along row \(1\) (Theorem~\ref{r5d-thm-expansion}),
\(\det V=3(9-1)+1(-3-1)-1(1+3)=24-4-4=16\). So the leading principal minors of \(V\) are \(3\), \(9-1=8\) and
\(\det V=16\), all positive, and \(V\) is PD (Theorem~\ref{r8-thm-leading-minors}), as
Proposition~\ref{r9-prop-multiple-correlation}(b) requires. The \((1,1)\) entry of \(Z=V^{-1}\) is
\(z_{11}=\det\Sigma/\det V=\frac8{16}=\frac12\) (Theorem~\ref{r5d-thm-adjugate}; the \((1,1)\) cofactor of \(V\) is
\(\det\Sigma\)), and \(1-1/(\sigma^2z_{11})=1-\frac1{3\cdot\frac12}=1-\frac23=\frac13\).
\end{itemize}
\end{solution}

\begin{exercise}\label{r9-dr-minimax}
Let
\[
A=\begin{pmatrix}3&3&1&-1\\3&3&-1&1\\1&-1&-2&-2\\-1&1&-2&-2\end{pmatrix}
\]
and \(w=(1,-1,0,0)'\).
\begin{enumerate}
\item Check that \((1,1,0,0)'\), \((1,-1,1,-1)'\), \((1,-1,-1,1)'\) and \((0,0,1,1)'\) are eigenvectors of \(A\), and find the
eigenvalues \(\lambda_1\geqslant\lambda_2\geqslant\lambda_3\geqslant\lambda_4\) of \(A\).
\item Find an orthonormal basis of the subspace of the vectors \(u\perp w\). Use it to find the maximum and the minimum
of \(u'Au\) over the unit vectors \(u\perp w\), and unit vectors at which they are attained.
\item Check the inequality of Theorem~\ref{r9-thm-courant-fischer}(a) for \(k=2\) and that of
Theorem~\ref{r9-thm-courant-fischer}(b) for \(k=3\), both with \(w_1=w\), and give a vector \(w_1\) for which each of them is
an equality.
\end{enumerate}
\end{exercise}
\begin{solution}
(a) \(A(1,1,0,0)'=(6,6,0,0)'\), \(A(1,-1,1,-1)'=(2,-2,2,-2)'\), \(A(1,-1,-1,1)'=(-2,2,2,-2)'\) and
\(A(0,0,1,1)'=(0,0,-4,-4)'\): eigenvectors belonging to \(6\), \(2\), \(-2\) and \(-4\). So \(\dim\mathcal N(A-cI)\geqslant1\) for
\(c=6,2,-2,-4\) (Theorem~\ref{r8-thm-rank-eigen}), and as the multiplicities add up to \(4\)
(Definition~\ref{r7-def-eigenvalues}), these are all the eigenvalues, each simple: \(\lambda_1=6\), \(\lambda_2=2\),
\(\lambda_3=-2\), \(\lambda_4=-4\). The unit eigenvectors are \(\pm q_i\), with \(q_1=\frac1{\sqrt2}(1,1,0,0)'\),
\(q_2=\frac12(1,-1,1,-1)'\), \(q_3=\frac12(1,-1,-1,1)'\) and \(q_4=\frac1{\sqrt2}(0,0,1,1)'\). Check: \(\tr A=2=6+2-2-4\).

\needspace{3\baselineskip}
(b)
\begin{itemize}
\item \(u\perp w\) means \(u_1=u_2\). The vectors \(g_1=\frac1{\sqrt2}(1,1,0,0)'\), \(g_2=\frac1{\sqrt2}(0,0,1,1)'\) and
\(g_3=\frac1{\sqrt2}(0,0,1,-1)'\) satisfy this and are orthonormal, hence linearly independent
(Proposition~\ref{r6-prop-orthogonal-list}(b)); the subspace has dimension \(3\) (Theorem~\ref{r6-thm-complement}(c)), so
they form an orthonormal basis of it (Theorem~\ref{thm:r4-counting}(d)).
\item \(g_1=q_1\) and \(g_2=q_4\) are eigenvectors, so \(Ag_1=6g_1\) and \(Ag_2=-4g_2\); and
\(Ag_3=\frac1{\sqrt2}(2,-2,0,0)'\), so \(g_3'Ag_3=\frac12(0,0,1,-1)(2,-2,0,0)'=0\). As \(A\) is symmetric,
\(g_i'Ag_j=(Ag_i)'g_j\): so \(g_1'Ag_1=6\), \(g_2'Ag_2=-4\), \(g_1'Ag_j=6g_1'g_j=0\) for \(j=2,3\), and
\(g_2'Ag_3=-4g_2'g_3=0\). With \(G=[g_1\ g_2\ g_3]\), therefore, \(G'AG=\diag(6,-4,0)\).
\item By Lemma~\ref{r9-lem-compression}, the maximum is \(\lambda_1(G'AG)=6\), attained at \(u=Ge_1=g_1\), and the
minimum is \(\lambda_3(G'AG)=-4\), attained at \(u=Ge_2=g_2\).
\end{itemize}
\needspace{3\baselineskip}
(c)
\begin{itemize}
\item Theorem~\ref{r9-thm-courant-fischer}(a) with \(k=2\), \(w_1=w\): \(m(w)=6\geqslant\lambda_2=2\).
\item Theorem~\ref{r9-thm-courant-fischer}(b) with \(k=3\), \(n-k=1\), \(w_1=w\): the minimum \(-4\leqslant\lambda_3=-2\).
\item Equality holds in (a) for \(w_1=q_1\) and in (b) for \(w_1=q_4\): the maximum of \(u'Au\) over the unit vectors
\(u\perp q_1\) is \(\lambda_2=2\), attained at \(q_2\), and the minimum over the unit vectors \(u\perp q_4\) is
\(\lambda_3=-2\), attained at \(q_3\) (Theorem~\ref{r9-thm-constrained}).
\end{itemize}
\end{solution}

\begin{exercise}\label{r9-ex-interlacing-drill}
The symmetric matrix
\[
A=\begin{pmatrix}3&1&0&-1\\1&4&2&2\\0&2&4&2\\-1&2&2&4\end{pmatrix}
\]
has the eigenvalues \(8\), \(4\), \(2\), \(1\) (Exercise~\ref{r9-ex-rayleigh-drill}). Find the eigenvalues of \(B=A(1|1)\), and
check the inequalities of Theorem~\ref{r9-thm-interlacing} for \(A\) and \(B\).
\end{exercise}
\begin{solution}
\begin{itemize}
\item Deleting row \(1\) and column \(1\) of \(A\) leaves
\[
B=\begin{pmatrix}4&2&2\\2&4&2\\2&2&4\end{pmatrix}=2I_3+2J_3
\]
\item \(2J_3=6\bar J_3\) and \(I_3=\bar J_3+C_3\) (Definition~\ref{r1-def-ones}), so \(B=8\bar J_3+2C_3\). As \(\bar J_3\one=\one\)
and \(C_3\one=0\), \(B\one=8\one\); if \(\one'x=0\), then \(\bar J_3x=0\) and \(C_3x=x\), so \(Bx=2x\). Hence
\(\frac1{\sqrt3}(1,1,1)'\), \(\frac1{\sqrt2}(1,-1,0)'\), \(\frac1{\sqrt6}(1,1,-2)'\) are orthonormal eigenvectors for \(8\), \(2\),
\(2\), and these are the eigenvalues of \(B\) (Theorem~\ref{r8-thm-spectral}(c)): \(\lambda_1(B)=8\),
\(\lambda_2(B)=\lambda_3(B)=2\).
\item With \(\lambda_1(A)=8\), \(\lambda_2(A)=4\), \(\lambda_3(A)=2\), \(\lambda_4(A)=1\), the inequalities
\(\lambda_k(A)\geqslant\lambda_k(B)\geqslant\lambda_{k+1}(A)\) (\(k=1,2,3\)) read
\[
8\geqslant8\geqslant4\geqslant2\geqslant2\geqslant2\geqslant1
\]
and they hold.
\end{itemize}
\end{solution}

\begin{exercise}\label{r9-dr-monotone}
Five units are measured on four variables; the rows of
\[
X=\begin{pmatrix}0&2&1&2\\2&0&0&2\\1&0&0&2\\0&1&2&2\\2&1&1&0\end{pmatrix}
\]
hold the values.
\begin{enumerate}
\item Find \(X'X\), check that \((1,1,1,2)'\), \((-2,1,1,0)'\), \((-2,-2,-2,3)'\) and \((0,-1,1,0)'\) are eigenvectors of it, and
find its eigenvalues.
\item A sixth unit with the values \(x'=(2,2,2,2)\) is added as a new row of \(X\). Find the new matrix \(X'X+xx'\), check
that \((2,2,2,3)'\), \((-2,1,1,0)'\), \((-1,-1,-1,2)'\) and \((0,-1,1,0)'\) are eigenvectors of it, find its eigenvalues, and
check Theorem~\ref{r9-thm-monotone}.
\item Check Theorem~\ref{r9-thm-interlacing} for \(X'X\) and its \(3\times3\) submatrix \((X'X)(4|4)\), which belongs to the
first three variables.
\end{enumerate}
\end{exercise}
\begin{solution}
(a) The entries of \(X'X\) are the inner products of the columns of \(X\) (Theorem~\ref{r1-thm-views}(d)):
\[
X'X=\begin{pmatrix}9&2&2&6\\2&6&5&6\\2&5&6&6\\6&6&6&16\end{pmatrix}
\]
Then \(X'X(1,1,1,2)'=(25,25,25,50)'\), \(X'X(-2,1,1,0)'=(-14,7,7,0)'\), \(X'X(-2,-2,-2,3)'=(-8,-8,-8,12)'\) and
\(X'X(0,-1,1,0)'=(0,-1,1,0)'\): eigenvectors belonging to \(25\), \(7\), \(4\) and \(1\). So \(\dim\mathcal N(X'X-cI)\geqslant1\) for
these four \(c\) (Theorem~\ref{r8-thm-rank-eigen}); as the multiplicities add up to \(4\) (Definition~\ref{r7-def-eigenvalues}),
the eigenvalues are \(25\), \(7\), \(4\) and \(1\).

(b) \(xx'=4J_4\), the \(4\times4\) matrix with every entry \(4\), and
\[
X'X+xx'=\begin{pmatrix}13&6&6&10\\6&10&9&10\\6&9&10&10\\10&10&10&20\end{pmatrix}
\]
which is \(X'X\) for the \(6\times4\) data matrix (Theorem~\ref{r1-thm-views}(d)). Then
\[
\begin{gathered}
(X'X+xx')(2,2,2,3)'=(80,80,80,120)',\qquad(X'X+xx')(-2,1,1,0)'=(-14,7,7,0)'
\\
(X'X+xx')(-1,-1,-1,2)'=(-5,-5,-5,10)',\qquad(X'X+xx')(0,-1,1,0)'=(0,-1,1,0)'
\end{gathered}
\]
so, as in (a), the eigenvalues are \(40\), \(7\), \(5\) and \(1\). As \(xx'\) is PSD
(\(v'xx'v=(x'v)^2\geqslant0\)), Theorem~\ref{r9-thm-monotone} applies, and indeed \(40\geqslant25\), \(7\geqslant7\),
\(5\geqslant4\) and \(1\geqslant1\).

(c) Deleting row \(4\) and column \(4\) of \(X'X\) leaves
\[
B=(X'X)(4|4)=\begin{pmatrix}9&2&2\\2&6&5\\2&5&6\end{pmatrix}
\]
the matrix of the inner products of the first three columns of \(X\). \(B(1,1,1)'=(13,13,13)'\), \(B(-2,1,1)'=(-14,7,7)'\) and
\(B(0,-1,1)'=(0,-1,1)'\), so, as in (a), the eigenvalues of \(B\) are \(13\), \(7\) and \(1\). By Theorem~\ref{r9-thm-interlacing}
(\(n=4\)), \(\lambda_k(X'X)\geqslant\lambda_k(B)\geqslant\lambda_{k+1}(X'X)\) for \(k=1,2,3\); indeed
\[
25\geqslant13\geqslant7\geqslant7\geqslant4\geqslant1\geqslant1
\]
\end{solution}

\begin{exercise}\label{r9-ex-svd-drill}
Let
\[
A=\begin{pmatrix}2&1\\1&2\\1&-1\end{pmatrix}
\]
\begin{enumerate}
\item Find an SVD and the compact SVD of \(A\).
\item Find \(\lVert A\rVert_F\) and \(\lVert A\rVert_2\).
\item Give orthonormal bases of the four subspaces of \(A\), and the matrix \(P_1P_1'\) of the projection onto
\(\mathcal C(A)\); check that it equals \(A(A'A)^{-1}A'\).
\end{enumerate}
\end{exercise}
\begin{solution}
(a)
\begin{itemize}
\item \(A'A=\left(\begin{smallmatrix}6&3\\3&6\end{smallmatrix}\right)\), with characteristic polynomial
\(t^2-12t+27=(t-9)(t-3)\) (Definition~\ref{r7-def-charpoly}(b)) and the unit eigenvectors \(q_1=\frac1{\sqrt2}(1,1)'\) for
\(9\) and \(q_2=\frac1{\sqrt2}(1,-1)'\) for \(3\). So \(r=2\), \(\sigma_1=3\), \(\sigma_2=\sqrt3\), and \(Q=[q_1\ q_2]\).
\item \(Aq_1=\frac1{\sqrt2}(3,3,0)'\) and \(Aq_2=\frac1{\sqrt2}(1,-1,2)'\), so \(p_1=\frac13Aq_1=\frac1{\sqrt2}(1,1,0)'\)
and \(p_2=\frac1{\sqrt3}Aq_2=\frac1{\sqrt6}(1,-1,2)'\) (Theorem~\ref{r9-thm-svd}).
\item \(\mathcal N(P_1')\): \(y_1+y_2=0\) and \(y_1-y_2+2y_3=0\) give \(y=t(1,-1,-1)'\); so \(p_3=\frac1{\sqrt3}(1,-1,-1)'\).
\item Hence, with \(Q'=Q=\frac1{\sqrt2}\left(\begin{smallmatrix}1&1\\1&-1\end{smallmatrix}\right)\),
\[
A=\begin{pmatrix}\frac1{\sqrt2}&\frac1{\sqrt6}&\frac1{\sqrt3}\\[2pt]\frac1{\sqrt2}&-\frac1{\sqrt6}&-\frac1{\sqrt3}\\[2pt]
0&\frac2{\sqrt6}&-\frac1{\sqrt3}\end{pmatrix}\begin{pmatrix}3&0\\0&\sqrt3\\0&0\end{pmatrix}Q',\qquad
A=3p_1q_1'+\sqrt3\,p_2q_2'=\frac32\begin{pmatrix}1&1\\1&1\\0&0\end{pmatrix}+\frac12\begin{pmatrix}1&-1\\-1&1\\2&-2\end{pmatrix}
\]
\end{itemize}
(b) \(\lVert A\rVert_F^2=4+1+1+4+1+1=12=9+3\) (Theorem~\ref{r9-thm-svd-props}(d)), so \(\lVert A\rVert_F=2\sqrt3\); and
\(\lVert A\rVert_2=\sigma_1=3\leqslant2\sqrt3\) (Proposition~\ref{r9-prop-spectral-norm}(d)).

\needspace{3\baselineskip}
(c)
\begin{itemize}
\item By Corollary~\ref{r9-cor-svd-subspaces}: \(p_1,p_2\) is an orthonormal basis of \(\mathcal C(A)\), \(p_3\) of
\(\mathcal N(A')\), \(q_1,q_2\) of \(\mathcal C(A')=\R^2\), and \(\mathcal N(A)=\{0\}\) (\(r=n=2\), no \(Q_2\)).
\item \(P_1P_1'=p_1p_1'+p_2p_2'\) (Theorem~\ref{r1-thm-views}(d)):
\[
P_1P_1'=\frac12\begin{pmatrix}1&1&0\\1&1&0\\0&0&0\end{pmatrix}+\frac16\begin{pmatrix}1&-1&2\\-1&1&-2\\2&-2&4\end{pmatrix}
=\frac13\begin{pmatrix}2&1&1\\1&2&-1\\1&-1&2\end{pmatrix}
\]
\item \(\det(A'A)=27\), so by Theorem~\ref{r5d-thm-adjugate}
\[
(A'A)^{-1}=\frac1{27}\begin{pmatrix}6&-3\\-3&6\end{pmatrix}=\frac19\begin{pmatrix}2&-1\\-1&2\end{pmatrix},\qquad
A(A'A)^{-1}=\frac13\begin{pmatrix}1&0\\0&1\\1&-1\end{pmatrix}
\]
and
\[
A(A'A)^{-1}A'=\frac13\begin{pmatrix}2&1&1\\1&2&-1\\1&-1&2\end{pmatrix}=P_1P_1'
\]
\end{itemize}
\end{solution}

\begin{exercise}[cf.\ \citeat{harville}{Ch.~21, Exercise~26}, \citeat{harville-solutions}{Ch.~21, Exercise~26}]\label{r9-ex-svd-wide}
Find an SVD and the compact SVD of \(A=\begin{pmatrix}1&0&1\\0&1&1\end{pmatrix}\) by applying Theorem~\ref{r9-thm-svd} to
\(A'\), starting from the \(2\times2\) matrix \(AA'\).
\end{exercise}
\begin{solution}
\begin{itemize}
\item \(A'\in\R^{3\times2}\) has rank \(2\), and \((A')'A'=AA'=\left(\begin{smallmatrix}2&1\\1&2\end{smallmatrix}\right)\),
with characteristic polynomial \(t^2-4t+3=(t-3)(t-1)\) and the unit eigenvectors \(p_1=\frac1{\sqrt2}(1,1)'\) for \(3\) and
\(p_2=\frac1{\sqrt2}(1,-1)'\) for \(1\). So \(D=\diag(\sqrt3,1)\).
\item Theorem~\ref{r9-thm-svd} for \(A'\), with \(P=[p_1\ p_2]\) in the role of \(Q\): the columns of \(A'PD^{-1}\) are
\(q_1=\frac1{\sqrt3}A'p_1=\frac1{\sqrt6}(1,1,2)'\) and \(q_2=A'p_2=\frac1{\sqrt2}(1,-1,0)'\). A unit vector orthogonal to
both: \(y_1+y_2+2y_3=0\) and \(y_1-y_2=0\) give \(q_3=\frac1{\sqrt3}(1,1,-1)'\).
\item So \(Q'A'P=\left(\begin{smallmatrix}D\\0\end{smallmatrix}\right)\) with \(Q=[q_1\ q_2\ q_3]\), and transposing,
\(P'AQ=[D\ \ 0]\):
\[
A=\frac1{\sqrt2}\begin{pmatrix}1&1\\1&-1\end{pmatrix}\begin{pmatrix}\sqrt3&0&0\\0&1&0\end{pmatrix}
\begin{pmatrix}\frac1{\sqrt6}&\frac1{\sqrt6}&\frac2{\sqrt6}\\[2pt]\frac1{\sqrt2}&-\frac1{\sqrt2}&0\\[2pt]
\frac1{\sqrt3}&\frac1{\sqrt3}&-\frac1{\sqrt3}\end{pmatrix}
\]
\item The compact SVD is
\(A=\sqrt3\,p_1q_1'+p_2q_2'=\frac12\left(\begin{smallmatrix}1&1&2\\1&1&2\end{smallmatrix}\right)
+\frac12\left(\begin{smallmatrix}1&-1&0\\-1&1&0\end{smallmatrix}\right)\). Also \(Aq_3=0\), so \(q_3\) spans
\(\mathcal N(A)\).
\end{itemize}
\end{solution}

\begin{exercise}\label{r9-dr-svd-rank1}
Let
\[
A=\begin{pmatrix}1&1\\2&2\\2&2\end{pmatrix}
\]
\begin{enumerate}
\item Find \(\rank A\), the eigenvalues of \(A'A\), and an SVD of \(A\).
\item Write down the compact SVD of \(A\) and orthonormal bases of the four subspaces of \(A\).
\item Find \(\lVert A\rVert_F\), \(\lVert A\rVert_2\) and the matrix \(P_1P_1'\) of the projection onto \(\mathcal C(A)\).
\end{enumerate}
\end{exercise}
\begin{solution}
(a)
\begin{itemize}
\item The two columns are equal and nonzero, so \(\rank A=1=r\).
\item \(A'A=\left(\begin{smallmatrix}9&9\\9&9\end{smallmatrix}\right)\), with characteristic polynomial
\(t^2-18t=t(t-18)\) (Definition~\ref{r7-def-charpoly}(b)), and \(A'A(1,1)'=(18,18)'\), \(A'A(1,-1)'=0\). So the
eigenvalues are \(18\) and \(0\), with the unit eigenvectors \(q_1=\frac1{\sqrt2}(1,1)'\) and \(q_2=\frac1{\sqrt2}(1,-1)'\);
\(\sigma_1=\sqrt{18}=3\sqrt2\) and \(D=(3\sqrt2)\).
\item \(P_1=AQ_1D^{-1}\) (Theorem~\ref{r9-thm-svd}): \(Aq_1=\frac1{\sqrt2}(2,4,4)'\), so
\(p_1=\frac1{3\sqrt2}\cdot\frac1{\sqrt2}(2,4,4)'=\frac13(1,2,2)'\).
\item \(\mathcal N(P_1')\) consists of the \(y\) with \(y_1+2y_2+2y_3=0\). The vectors \((2,1,-2)'\) and \((2,-2,1)'\) lie in it
and are orthogonal (\(4-2-2=0\)); normalized, \(p_2=\frac13(2,1,-2)'\) and \(p_3=\frac13(2,-2,1)'\).
\item Hence \(A=P\left(\begin{smallmatrix}D&0\\0&0\end{smallmatrix}\right)Q'\) with
\[
P=\frac13\begin{pmatrix}1&2&2\\2&1&-2\\2&-2&1\end{pmatrix},\qquad
\begin{pmatrix}D&0\\0&0\end{pmatrix}=\begin{pmatrix}3\sqrt2&0\\0&0\\0&0\end{pmatrix},\qquad
Q=\frac1{\sqrt2}\begin{pmatrix}1&1\\1&-1\end{pmatrix}
\]
\end{itemize}
(b) The compact SVD is \(A=3\sqrt2\,p_1q_1'=3\sqrt2\cdot\frac13(1,2,2)'\cdot\frac1{\sqrt2}(1,1)=(1,2,2)'(1,1)\), which is
\(A\). By Corollary~\ref{r9-cor-svd-subspaces}: \(p_1\) is an orthonormal basis of \(\mathcal C(A)\), \(p_2,p_3\) of
\(\mathcal N(A')\), \(q_1\) of \(\mathcal C(A')\), and \(q_2\) of \(\mathcal N(A)\).

(c) \(\lVert A\rVert_F^2=1+1+4+4+4+4=18=\sigma_1^2\) (Theorem~\ref{r9-thm-svd-props}(d)), so
\(\lVert A\rVert_F=3\sqrt2=\sigma_1=\lVert A\rVert_2\): the two norms agree for a matrix of rank \(1\). By
Corollary~\ref{r9-cor-svd-subspaces},
\[
P_1P_1'=p_1p_1'=\frac19\begin{pmatrix}1&2&2\\2&4&4\\2&4&4\end{pmatrix}
\]
\end{solution}

\begin{exercise}\label{r10-ex-svd-drill}
Let
\[
A=\begin{pmatrix}1&3\\3&1\\3&3\\1&1\end{pmatrix}
\]
\begin{enumerate}
\item Find a singular value decomposition \(A=P\left(\begin{smallmatrix}D\\0\end{smallmatrix}\right)Q'\) with
\(P\in\R^{4\times4}\), and the compact form.
\item Give orthonormal bases of \(\mathcal C(A)\), \(\mathcal N(A')\), \(\mathcal C(A')\) and \(\mathcal N(A)\).
\item Find \(\lVert A\rVert_2\), \(\lVert A\rVert_F\), \(\kappa(A)\) and \(\kappa(A'A)\).
\item Find \(A^+\) from the singular value decomposition, and check it with \((A'A)^{-1}A'\).
\end{enumerate}
\end{exercise}
\begin{solution}
(a)
\begin{itemize}
\item \(A'A=\left(\begin{smallmatrix}20&16\\16&20\end{smallmatrix}\right)\), with \(A'A(1,1)'=36(1,1)'\) and
\(A'A(1,-1)'=4(1,-1)'\). By Theorem~\ref{r9-thm-svd}(b): \(\sigma_1=6\), \(\sigma_2=2\), \(q_1=\frac1{\sqrt2}(1,1)'\),
\(q_2=\frac1{\sqrt2}(1,-1)'\), \(p_1=\frac16Aq_1=\frac1{6\sqrt2}(4,4,6,2)'=\frac1{3\sqrt2}(2,2,3,1)'\) and
\(p_2=\frac12Aq_2=\frac1{\sqrt2}(-1,1,0,0)'\).
\item A vector \(y\) orthogonal to \(p_1\) and \(p_2\) satisfies \(y_1=y_2\) and \(4y_1+3y_3+y_4=0\); the vectors
\((0,0,1,-3)'\) and \((5,5,-6,-2)'\) satisfy both and are orthogonal, with lengths \(\sqrt{10}\) and \(3\sqrt{10}\). Take
\(p_3=\frac1{\sqrt{10}}(0,0,1,-3)'\) and \(p_4=\frac1{3\sqrt{10}}(5,5,-6,-2)'\).
\item Then \(A=P\left(\begin{smallmatrix}D\\0\end{smallmatrix}\right)Q'\) with \(P=[p_1\ p_2\ p_3\ p_4]\),
\(D=\diag(6,2)\), \(Q=\frac1{\sqrt2}\left(\begin{smallmatrix}1&1\\1&-1\end{smallmatrix}\right)\), and
\[
A=6p_1q_1'+2p_2q_2'=\begin{pmatrix}2&2\\2&2\\3&3\\1&1\end{pmatrix}+\begin{pmatrix}-1&1\\1&-1\\0&0\\0&0\end{pmatrix}
\]
\end{itemize}
(b) By Corollary~\ref{r9-cor-svd-subspaces}: \(p_1,p_2\) is an orthonormal basis of \(\mathcal C(A)\), \(p_3,p_4\) of
\(\mathcal N(A')\), and \(q_1,q_2\) of \(\mathcal C(A')=\R^2\); \(\mathcal N(A)=\{0\}\) (\(r=n=2\), no \(Q_2\)).

(c) \(\lVert A\rVert_2=\sigma_1=6\) (Definition~\ref{r9-def-spectral-norm}); \(\lVert A\rVert_F^2=10+10+18+2=40=6^2+2^2\)
(Theorem~\ref{r9-thm-svd-props}(d)), so \(\lVert A\rVert_F=2\sqrt{10}\); \(\kappa(A)=6/2=3\)
(Definition~\ref{r9-def-condition}) and \(\kappa(A'A)=\kappa(A)^2=9\) (Proposition~\ref{r9-prop-condition}(b)).

(d) By Theorem~\ref{r10-thm-mp-svd},
\[
A^+=\frac16q_1p_1'+\frac12q_2p_2'=\frac1{36}\begin{pmatrix}2&2&3&1\\2&2&3&1\end{pmatrix}
+\frac14\begin{pmatrix}-1&1&0&0\\1&-1&0&0\end{pmatrix}
=\frac1{36}\begin{pmatrix}-7&11&3&1\\11&-7&3&1\end{pmatrix}
\]
As \(\rank A=2\), the number of columns, \(A^+=(A'A)^{-1}A'\) (Theorem~\ref{r10-thm-mp-exists}(c)). Here
\(\det(A'A)=400-256=144\), so
\((A'A)^{-1}=\frac1{144}\left(\begin{smallmatrix}20&-16\\-16&20\end{smallmatrix}\right)
=\frac1{36}\left(\begin{smallmatrix}5&-4\\-4&5\end{smallmatrix}\right)\) (Theorem~\ref{r5d-thm-adjugate}), and
\[
(A'A)^{-1}A'=\frac1{36}\begin{pmatrix}5&-4\\-4&5\end{pmatrix}\begin{pmatrix}1&3&3&1\\3&1&3&1\end{pmatrix}
=\frac1{36}\begin{pmatrix}-7&11&3&1\\11&-7&3&1\end{pmatrix}
\]
\end{solution}

\begin{exercise}\label{r10-ex-rank-one-drill}
Let
\[
A=\begin{pmatrix}2&1&2\\4&2&4\\4&2&4\\8&4&8\end{pmatrix}
\]
\begin{enumerate}
\item Write \(A=uv'\), and find a compact singular value decomposition, \(\lVert A\rVert_2\) and \(\lVert A\rVert_F\).
\item Find \(A^+\), \(AA^+\) and \(A^+A\).
\item Decide whether \(Ax=b\) is consistent for \(b=(3,6,6,12)'\) and for \(b=(1,0,0,0)'\). Find the solution of minimum
length in the consistent case, and compare it with the solution \((0,3,0)'\).
\end{enumerate}
\end{exercise}
\begin{solution}
(a)
\begin{itemize}
\item The rows of \(A\) are \(1\), \(2\), \(2\), \(4\) times \((2,1,2)\), so \(A=uv'\) with \(u=(1,2,2,4)'\), \(\lVert u\rVert=5\),
and \(v=(2,1,2)'\), \(\lVert v\rVert=3\); \(\rank A=1\).
\item \(A'A=v(u'u)v'=25vv'\) has rank \(\rank A=1\) (Theorem~\ref{thm:r4-rank-ata}(b)) and \(A'Av=25v(v'v)=225v\). By
Theorem~\ref{r8-thm-rank-eigen}, \(225\) is its only nonzero eigenvalue, and the others are \(0\).
\item By Theorem~\ref{r9-thm-svd}(b): \(\sigma_1=\sqrt{225}=15\), \(q_1=\frac13v\), and
\(p_1=\frac1{15}Aq_1=\frac1{45}u(v'v)=\frac15u\). So \(A=15p_1q_1'\) is a compact SVD.
\item \(\lVert A\rVert_2=\sigma_1=15\) (Definition~\ref{r9-def-spectral-norm}) and \(\lVert A\rVert_F^2=\sigma_1^2=225\)
(Theorem~\ref{r9-thm-svd-props}(d)), so \(\lVert A\rVert_F=15\); indeed \(\lVert A\rVert_F^2=9+36+36+144=225\).
\end{itemize}
(b) By Theorem~\ref{r10-thm-mp-svd},
\[
A^+=\frac1{15}q_1p_1'=\frac1{225}vu'=\frac1{225}\begin{pmatrix}2&4&4&8\\1&2&2&4\\2&4&4&8\end{pmatrix},\qquad
AA^+=p_1p_1'=\frac1{25}uu',\qquad A^+A=q_1q_1'=\frac19vv'
\]
(c)
\begin{itemize}
\item \(A^+\) is a g-inverse of \(A\) ((P1) of Definition~\ref{r10-def-mp}), so by Theorem~\ref{r10-thm-ginv-solves}(b),
\(Ax=b\) is consistent if and only if \(AA^+b=\frac{u'b}{25}u\) equals \(b\). For \(b=(3,6,6,12)'=3u\), \(u'b=75\) and
\(AA^+b=3u=b\): consistent. For \(b=(1,0,0,0)'\), \(AA^+b=\frac1{25}u\ne b\): there is no solution.
\item \(A^+\) is a minimum-norm g-inverse ((P1) and (P4)), so by Theorem~\ref{r10-thm-min-norm} the shortest solution for
\(b=3u\) is \(x^+=A^+b=\frac{u'b}{225}v=\frac{75}{225}v=\frac13(2,1,2)'\), with \(\lVert x^+\rVert^2=1\).
\item The solution \(x=(0,3,0)'\) (\(Ax=3\col_2(A)=3u\)) has \(\lVert x\rVert^2=9\); and \(x-x^+=\frac13(-2,8,-2)'\) is orthogonal
to \(x^+\), with \(\lVert x-x^+\rVert^2=\frac{72}9=8\), so \(9=1+8\).
\end{itemize}
\end{solution}

\begin{exercise}\label{r10-ex-svd-4x3}
Let
\[
A=\begin{pmatrix}6&0&0\\4&4&6\\5&2&6\\2&-4&0\end{pmatrix}
\]
\begin{enumerate}
\item Find a singular value decomposition of \(A\) and its compact form. (Write \(A'A=9M\) and use the vectors
\(w_1=(2,1,2)'\), \(w_2=(-2,2,1)'\), \(w_3=(1,2,-2)'\).)
\item Give orthonormal bases of the four subspaces of \(A\), and the eigenvalues of \(AA'\).
\item Find \(\lVert A\rVert_2\), \(\lVert A\rVert_F\), \(\kappa(A)\) and \(\kappa(A'A)\).
\item Find \(A^+\), once from \((A'A)^{-1}A'\) and once from the singular value decomposition, and \(AA^+\).
\item Find a matrix \(A_2\) of rank at most \(2\) nearest to \(A\) in the spectral norm, and \(\lVert A-A_2\rVert_2\).
\end{enumerate}
\end{exercise}
\begin{solution}
(a)
\begin{itemize}
\item The entries of \(A'A\) are the inner products of the columns of \(A\) (Theorem~\ref{r1-thm-views}(d)):
\[
A'A=\begin{pmatrix}81&18&54\\18&36&36\\54&36&72\end{pmatrix}=9M,\qquad
M=\begin{pmatrix}9&2&6\\2&4&4\\6&4&8\end{pmatrix}
\]
\(Mw_1=(32,16,32)'=16w_1\), \(Mw_2=(-8,8,4)'=4w_2\) and \(Mw_3=(1,2,-2)'=w_3\). So \(A'Aw_k=\lambda_kw_k\) with
\(\lambda_1=144\), \(\lambda_2=36\), \(\lambda_3=9\); as \(\dim\mathcal N(A'A-cI)\geqslant1\) for these three values and the
multiplicities add up to \(3\) (Theorem~\ref{r8-thm-rank-eigen}, Definition~\ref{r7-def-eigenvalues}), they are the
eigenvalues of \(A'A\). The \(w_k\) are orthogonal with length \(3\), so \(q_k=\frac13w_k\) gives an orthogonal
\(Q=[q_1\ q_2\ q_3]\). \(A'A\) has three nonzero eigenvalues, so \(\rank A=\rank(A'A)=3\)
(Theorems~\ref{r8-thm-rank-eigen} and~\ref{thm:r4-rank-ata}(b)), and by Theorem~\ref{r9-thm-svd}(b), \(\sigma_1=12\),
\(\sigma_2=6\), \(\sigma_3=3\).
\item \(p_k=\sigma_k^{-1}Aq_k=(3\sigma_k)^{-1}Aw_k\): \(Aw_1=(12,24,24,0)'\), \(Aw_2=(-12,6,0,-12)'\) and
\(Aw_3=(6,0,-3,-6)'\) give \(p_1=\frac13(1,2,2,0)'\), \(p_2=\frac13(-2,1,0,-2)'\) and \(p_3=\frac13(2,0,-1,-2)'\).
\item A vector \(y\) orthogonal to \(p_1,p_2,p_3\) satisfies \(y_1+2y_2+2y_3=0\), \(-2y_1+y_2-2y_4=0\) and
\(2y_1-y_3-2y_4=0\), so \(y=c\,(0,2,-2,1)'\); take \(p_4=\frac13(0,2,-2,1)'\).
\end{itemize}
Hence \(A=P\left(\begin{smallmatrix}D\\0\end{smallmatrix}\right)Q'\) with \(D=\diag(12,6,3)\),
\[
P=\frac13\begin{pmatrix}1&-2&2&0\\2&1&0&2\\2&0&-1&-2\\0&-2&-2&1\end{pmatrix},\qquad
Q=\frac13\begin{pmatrix}2&-2&1\\1&2&2\\2&1&-2\end{pmatrix}
\]
and the compact form is
\[
A=12p_1q_1'+6p_2q_2'+3p_3q_3'
=\frac43\begin{pmatrix}2&1&2\\4&2&4\\4&2&4\\0&0&0\end{pmatrix}
+\frac23\begin{pmatrix}4&-4&-2\\-2&2&1\\0&0&0\\4&-4&-2\end{pmatrix}
+\frac13\begin{pmatrix}2&4&-4\\0&0&0\\-1&-2&2\\-2&-4&4\end{pmatrix}
\]
(b) By Corollary~\ref{r9-cor-svd-subspaces}: \(p_1,p_2,p_3\) is an orthonormal basis of \(\mathcal C(A)\), \(p_4\) of
\(\mathcal N(A')\) (indeed \(A'(0,2,-2,1)'=0\)), \(q_1,q_2,q_3\) of \(\mathcal C(A')=\R^3\), and \(\mathcal N(A)=\{0\}\). By
Theorem~\ref{r9-thm-svd-props}(a),(b), the eigenvalues of the \(4\times4\) matrix \(AA'\) are \(144\), \(36\), \(9\), \(0\).

(c) \(\lVert A\rVert_2=12\); \(\lVert A\rVert_F^2=36+68+65+20=189=12^2+6^2+3^2\) (the four terms are the squared lengths of the
rows; Theorem~\ref{r9-thm-svd-props}(d)), so \(\lVert A\rVert_F=3\sqrt{21}\); \(\kappa(A)=12/3=4\)
(Definition~\ref{r9-def-condition}) and \(\kappa(A'A)=16\) (Proposition~\ref{r9-prop-condition}(b)).

(d)
\begin{itemize}
\item \(\rank A=3\), the number of columns, so \(A^+=(A'A)^{-1}A'=\frac19M^{-1}A'\) (Theorem~\ref{r10-thm-mp-exists}(c)).
Along row \(1\), \(\det M=9\cdot16-2\cdot(-8)+6\cdot(-16)=64\) (Theorem~\ref{r5d-thm-expansion}), and by
Theorem~\ref{r5d-thm-adjugate}
\[
M^{-1}=\frac1{64}\begin{pmatrix}16&8&-16\\8&36&-24\\-16&-24&32\end{pmatrix}
=\frac1{16}\begin{pmatrix}4&2&-4\\2&9&-6\\-4&-6&8\end{pmatrix}
\]
so
\[
A^+=\frac1{144}\begin{pmatrix}4&2&-4\\2&9&-6\\-4&-6&8\end{pmatrix}
\begin{pmatrix}6&4&5&2\\0&4&2&-4\\0&6&6&0\end{pmatrix}
=\frac1{36}\begin{pmatrix}6&0&0&0\\3&2&-2&-8\\-6&2&4&4\end{pmatrix}
\]
\item By Theorem~\ref{r10-thm-mp-svd}, \(A^+=\frac1{12}q_1p_1'+\frac16q_2p_2'+\frac13q_3p_3'\), the same matrix:
\[
\frac1{108}\begin{pmatrix}2&4&4&0\\1&2&2&0\\2&4&4&0\end{pmatrix}
+\frac1{54}\begin{pmatrix}4&-2&0&4\\-4&2&0&-4\\-2&1&0&-2\end{pmatrix}
+\frac1{27}\begin{pmatrix}2&0&-1&-2\\4&0&-2&-4\\-4&0&2&4\end{pmatrix}
=\frac1{36}\begin{pmatrix}6&0&0&0\\3&2&-2&-8\\-6&2&4&4\end{pmatrix}
\]
\item By Theorem~\ref{r10-thm-mp-svd}, \(AA^+=P_1P_1'\), where \(P_1=[p_1\ p_2\ p_3]\); as \(PP'=I_4\)
(Proposition~\ref{r6-prop-orthonormal}(b)), \(P_1P_1'=I_4-p_4p_4'\):
\[
AA^+=I_4-\frac19\begin{pmatrix}0&0&0&0\\0&4&-4&2\\0&-4&4&-2\\0&2&-2&1\end{pmatrix}
=\frac19\begin{pmatrix}9&0&0&0\\0&5&4&-2\\0&4&5&2\\0&-2&2&8\end{pmatrix}
\]
the matrix of the orthogonal projection onto \(\mathcal C(A)\) (Corollary~\ref{r9-cor-svd-subspaces}).
\end{itemize}
(e) By Theorem~\ref{r12-thm-eckart-young}(b) with \(k=2<r=3\), every \(B\) of rank at most \(2\) has
\(\lVert A-B\rVert_2\geqslant\sigma_3=3\), with equality for
\[
A_2=12p_1q_1'+6p_2q_2'=A-3p_3q_3'=\frac13\begin{pmatrix}16&-4&4\\12&12&18\\16&8&16\\8&-8&-4\end{pmatrix}
\]
So \(\lVert A-A_2\rVert_2=3=\lVert A\rVert_2/\kappa(A)\): the matrix \(E=A_2-A=-3p_3q_3'\), of spectral norm \(3\), makes the
columns of \(A+E=A_2\) linearly dependent (\(A_2w_3=0\)), and no \(E\) of smaller spectral norm does.
\end{solution}

\begin{exercise}\label{r10-ex-symmetric-svd}
Let
\[
A=\begin{pmatrix}1&4&0&2\\4&1&2&0\\0&2&1&4\\2&0&4&1\end{pmatrix},\qquad
h_1=\begin{pmatrix}1\\1\\1\\1\end{pmatrix},\ h_2=\begin{pmatrix}1\\-1\\1\\-1\end{pmatrix},\
h_3=\begin{pmatrix}1\\1\\-1\\-1\end{pmatrix},\ h_4=\begin{pmatrix}1\\-1\\-1\\1\end{pmatrix}
\]
\begin{enumerate}
\item Show that \(Ah_1=7h_1\), \(Ah_2=-5h_2\), \(Ah_3=3h_3\) and \(Ah_4=-h_4\).
\item Find a singular value decomposition \(A=PDQ'\) of \(A\). Why is \(P\ne Q\)?
\item Find \(\lVert A\rVert_2\), \(\lVert A\rVert_F\) and \(\kappa(A)\).
\end{enumerate}
\end{exercise}
\begin{solution}
(a) \(Ah_1=(7,7,7,7)'\) (every row sum is \(7\)), \(Ah_2=(-5,5,-5,5)'\), \(Ah_3=(3,3,-3,-3)'\) and \(Ah_4=(-1,1,1,-1)'\).

(b)
\begin{itemize}
\item The \(h_k\) are orthogonal with length \(2\), so \(Q=\frac12[h_1\ h_2\ h_3\ h_4]\) is orthogonal, and
\(Q'AQ=\Lambda=\diag(7,-5,3,-1)\) (from \(AQ=Q\Lambda\)).
\item \(A'A=A^2=Q\Lambda^2Q'\), so \(A'Aq_k=\lambda_k^2q_k\) for \(q_k=\frac12h_k\), with \(\lambda_k^2=49,25,9,1\) in decreasing
order. By Theorem~\ref{r9-thm-svd}(b), \(\sigma_1=7\), \(\sigma_2=5\), \(\sigma_3=3\), \(\sigma_4=1\), \(D=\diag(7,5,3,1)\), and
\(p_k=\sigma_k^{-1}Aq_k=(\lambda_k/\lvert\lambda_k\rvert)q_k\):
\[
P=\frac12[h_1\ \ {-h_2}\ \ h_3\ \ {-h_4}],\qquad D=\diag(7,5,3,1),\qquad Q=\frac12[h_1\ h_2\ h_3\ h_4]
\]
\item The columns of \(P\) and \(Q\) differ in sign exactly where the eigenvalue is negative: \(P=Q\Delta\) with
\(\Delta=\diag(1,-1,1,-1)\). As \(A\) has negative eigenvalues, \(P\ne Q\); the singular values are the absolute values of the
eigenvalues (Theorem~\ref{r9-thm-svd-props}(e)).
\end{itemize}
(c) \(\lVert A\rVert_2=\sigma_1=7\) (Definition~\ref{r9-def-spectral-norm}); \(\lVert A\rVert_F^2=4\cdot21=84=49+25+9+1\)
(Theorem~\ref{r9-thm-svd-props}(d)), so \(\lVert A\rVert_F=2\sqrt{21}\); \(\rank A=4\) and \(\kappa(A)=7/1=7\)
(Proposition~\ref{r9-prop-condition}(a)).
\end{solution}

\begin{exercise}\label{r9-dr-norms}
Let
\[
A=\begin{pmatrix}4&2&1&0\\1&0&4&2\\2&4&0&1\\0&1&2&4\end{pmatrix},\qquad
B=\begin{pmatrix}-4&-2&-1&0\\2&4&0&1\\1&0&4&2\\0&-1&-2&-4\end{pmatrix}
\]
and \(h_1=(1,1,1,1)'\), \(h_2=(1,-1,1,-1)'\), \(h_3=(1,1,-1,-1)'\), \(h_4=(1,-1,-1,1)'\).
\begin{enumerate}
\item Find \(A'A\) and \(A'Ah_1,\dots,A'Ah_4\). Find the singular values, \(\lVert\cdot\rVert_2\) and \(\lVert\cdot\rVert_F\) of
\(A\) and of \(B\).
\item Find \(Ah_1,\dots,Ah_4\) and the eigenvalues of \(A\), and compare their absolute values with the singular values of
\(A\).
\item Find \(\lVert A+B\rVert_2\) and \(\lVert AB\rVert_2\), and check the inequalities of
Proposition~\ref{r9-prop-spectral-norm-rules}.
\item Find unit vectors \(u\), \(v\) with \(u'Av=\lVert A\rVert_2\).
\end{enumerate}
\end{exercise}
\begin{solution}
(a)
\begin{itemize}
\item The entries of \(A'A\) are the inner products of the columns of \(A\) (Theorem~\ref{r1-thm-views}(d)):
\[
A'A=\begin{pmatrix}21&16&8&4\\16&21&4&8\\8&4&21&16\\4&8&16&21\end{pmatrix}
\]
and \(A'Ah_1=(49,49,49,49)'=49h_1\), \(A'Ah_2=(9,-9,9,-9)'=9h_2\), \(A'Ah_3=(25,25,-25,-25)'=25h_3\),
\(A'Ah_4=(1,-1,-1,1)'=h_4\). So \(\dim\mathcal N(A'A-cI)\geqslant1\) for \(c=49,25,9,1\) (Theorem~\ref{r8-thm-rank-eigen}), and
as the multiplicities add up to \(4\) (Definition~\ref{r7-def-eigenvalues}), these are the eigenvalues of \(A'A\). By
Theorem~\ref{r9-thm-svd-props}(a), \(A\) has the singular values \(7\), \(5\), \(3\), \(1\).
\item \(B=UA\) with
\[
U=\begin{pmatrix}-1&0&0&0\\0&0&1&0\\0&1&0&0\\0&0&0&-1\end{pmatrix}
\]
(row \(1\) of \(B\) is minus row \(1\) of \(A\), rows \(2\) and \(3\) are rows \(3\) and \(2\) of \(A\), row \(4\) is minus row \(4\) of
\(A\)), and \(U'U=I_4\); so \(B'B=A'U'UA=A'A\), and \(B\) has the same singular values.
\item \(\lVert A\rVert_2=\lVert B\rVert_2=7\), and \(\lVert A\rVert_F^2=49+25+9+1=84=\lVert B\rVert_F^2\)
(Theorem~\ref{r9-thm-svd-props}(d); also the sum of the squared entries, \(4\cdot21\)), so
\(\lVert A\rVert_F=\lVert B\rVert_F=2\sqrt{21}\).
\end{itemize}
(b)
\begin{itemize}
\item \(Ah_1=(7,7,7,7)'=7h_1\), \(Ah_2=(3,3,-3,-3)'=3h_3\), \(Ah_3=(5,-5,5,-5)'=5h_2\) and \(Ah_4=(1,-1,-1,1)'=h_4\). With
\(H=[h_1\ h_2\ h_3\ h_4]\), invertible as \(H'H=4I_4\), this says \(AH=HM\), that is, \(H^{-1}AH=M\), where
\[
M=\begin{pmatrix}7&0&0&0\\0&0&5&0\\0&3&0&0\\0&0&0&1\end{pmatrix}
\]
\item By Theorem~\ref{r7-thm-similar}(a), \(p_A=p_M\); and \(M\) is block diagonal, so
\(p_M(t)=(t-7)(t^2-15)(t-1)\) (Proposition~\ref{r7-prop-block-triangular}(a), applied twice; the middle block has the
characteristic polynomial \(t^2-15\)). So the eigenvalues of \(A\) are \(7\), \(\sqrt{15}\), \(-\sqrt{15}\), \(1\)
(Definition~\ref{r7-def-eigenvalues}).
\item Their absolute values \(7\), \(\sqrt{15}\), \(\sqrt{15}\), \(1\) are not the singular values \(7\), \(5\), \(3\), \(1\):
Theorem~\ref{r9-thm-svd-props}(e) needs a symmetric matrix, and \(A\) is not symmetric. (The products agree:
\(7\cdot15\cdot1=105=7\cdot5\cdot3\cdot1\), the absolute value of \(\det A\).)
\end{itemize}
\needspace{3\baselineskip}
(c)
\begin{itemize}
\item \(A+B=(0,1,1,0)'(3,4,4,3)\), the matrix with rows \(0\), \((3,4,4,3)\), \((3,4,4,3)\), \(0\). With \(w=(3,4,4,3)'\),
\((A+B)'(A+B)=w(0,1,1,0)(0,1,1,0)'w'=2ww'\), of rank \(1\), and \(2ww'w=2(w'w)w=100w\). By
Theorem~\ref{r8-thm-rank-eigen}, its eigenvalues are \(100\), \(0\), \(0\), \(0\); so
\(\lVert A+B\rVert_2=10\leqslant\lVert A\rVert_2+\lVert B\rVert_2=14\).
\item
\[
AB=\begin{pmatrix}-11&0&0&4\\0&-4&11&0\\0&11&-4&0\\4&0&0&-11\end{pmatrix}
\]
is symmetric. \(AB(1,0,0,1)'=-7(1,0,0,1)'\), \(AB(1,0,0,-1)'=-15(1,0,0,-1)'\), \(AB(0,1,1,0)'=7(0,1,1,0)'\) and
\(AB(0,1,-1,0)'=-15(0,1,-1,0)'\); as in (a), its eigenvalues are \(7\), \(-7\), \(-15\), \(-15\), and
\(\lVert AB\rVert_2=15\) (Proposition~\ref{r9-prop-spectral-norm}(e)); \(15\leqslant\lVert A\rVert_2\lVert B\rVert_2=49\).
\end{itemize}
(d) \(A'Ah_1=49h_1\), so \(q_1=\frac12h_1\), and \(p_1=\frac17Aq_1=\frac1{14}(7,7,7,7)'=\frac12(1,1,1,1)'\)
(Theorem~\ref{r9-thm-svd}). By Theorem~\ref{r9-thm-sigma-max}(a), \(u=p_1\) and \(v=q_1\) give \(u'Av=\sigma_1=7\). Check:
every row sum of \(A\) is \(7\), so \(Av=\frac12(7,7,7,7)'\) and \(u'Av=\frac14\cdot28=7\).
\end{solution}

\begin{exercise}\label{r12-ex-rank-one}
Let
\[
A=\begin{pmatrix}1&1\\1&1\\2&-2\end{pmatrix}
\]
\begin{enumerate}
\item Find a compact SVD \(A=\sigma_1p_1q_1'+\sigma_2p_2q_2'\).
\item Find a best approximation \(A_1\) of \(A\) by a matrix of rank at most \(1\), and \(\lVert A-A_1\rVert_F^2\).
\item The matrix
\[
B=\begin{pmatrix}1&1\\1&1\\0&0\end{pmatrix}
\]
also has rank \(1\). Compute \(\lVert A-B\rVert_F^2\).
\end{enumerate}
\end{exercise}
\begin{solution}
(a) \(A'A=\begin{pmatrix}6&-2\\-2&6\end{pmatrix}\), and \(\det(tI_2-A'A)=(t-6)^2-4=(t-8)(t-4)\). Indeed
\(A'A(1,-1)'=(8,-8)'\) and \(A'A(1,1)'=(4,4)'\), so \(q_1=\frac1{\sqrt2}(1,-1)'\), \(q_2=\frac1{\sqrt2}(1,1)'\),
\(\sigma_1=\sqrt8=2\sqrt2\) and \(\sigma_2=2\). Then \(p_m=Aq_m/\sigma_m\), the columns of \(P_1=AQ_1D^{-1}\)
(Theorem~\ref{r9-thm-svd}(a),(b)):
\[
p_1=\frac1{2\sqrt2}\cdot\frac1{\sqrt2}\begin{pmatrix}0\\0\\4\end{pmatrix}=\begin{pmatrix}0\\0\\1\end{pmatrix},\qquad
p_2=\frac12\cdot\frac1{\sqrt2}\begin{pmatrix}2\\2\\0\end{pmatrix}=\frac1{\sqrt2}\begin{pmatrix}1\\1\\0\end{pmatrix}
\]
\(p_1\), \(p_2\) are orthonormal, and
\[
\sigma_1p_1q_1'+\sigma_2p_2q_2'=\begin{pmatrix}0&0\\0&0\\2&-2\end{pmatrix}+\begin{pmatrix}1&1\\1&1\\0&0\end{pmatrix}=A
\]
With \(p_3=\frac1{\sqrt2}(1,-1,0)'\), \(P=[p_1\ p_2\ p_3]\) is orthogonal, and this is the compact form of the SVD
\(A=P\left(\begin{smallmatrix}D&0\\0&0\end{smallmatrix}\right)Q'\) of Theorem~\ref{r9-thm-svd}(b).

(b) By Theorem~\ref{r12-thm-eckart-young}(a) (\(k=1\), \(r=2\)),
\[
A_1=\sigma_1p_1q_1'=\begin{pmatrix}0&0\\0&0\\2&-2\end{pmatrix}
\]
and \(\lVert A-A_1\rVert_F^2=\lVert\sigma_2p_2q_2'\rVert_F^2=1+1+1+1=4=\sigma_2^2\).

(c)
\[
A-B=\begin{pmatrix}0&0\\0&0\\2&-2\end{pmatrix}
\]
so \(\lVert A-B\rVert_F^2=8=\sigma_1^2>4\): \(B=\sigma_2p_2q_2'\) keeps the
smaller singular value and loses the larger one.
\end{solution}

\begin{exercise}\label{r12-ex-symmetric}
Let \(A=2J_4-I_4\).
\begin{enumerate}
\item Show that \(A=7\bar J_4-C_4\), and find the eigenvalues of \(A\) and orthonormal eigenvectors.
\item Write \(A=PDQ'\) with \(P\), \(Q\) orthogonal and \(D=\diag(7,1,1,1)\), a singular value decomposition.
\item Find a best approximation of \(A\) by a matrix of rank at most \(1\), and the error \(\lVert A-A_1\rVert_F^2\).
\end{enumerate}
\end{exercise}
\begin{solution}
(a) \(2J_4=8\bar J_4\) and \(I_4=\bar J_4+C_4\) (Definition~\ref{r1-def-ones}), so \(A=8\bar J_4-\bar J_4-C_4=7\bar J_4-C_4\).
As \(\bar J_4\one=\one\) and \(C_4\one=0\), \(A\one=7\one\); if \(\one'x=0\), then \(\bar J_4x=0\) and \(C_4x=x\), so
\(Ax=-x\). Hence \(q_1=\frac12(1,1,1,1)'\), \(q_2=\frac12(1,-1,1,-1)'\), \(q_3=\frac12(1,1,-1,-1)'\), \(q_4=\frac12(1,-1,-1,1)'\)
are orthonormal eigenvectors for \(7\), \(-1\), \(-1\), \(-1\), and these are the eigenvalues
(Theorem~\ref{r8-thm-spectral}(c)).

(b) \(A=7q_1q_1'-q_2q_2'-q_3q_3'-q_4q_4'=7q_1q_1'+1\cdot(-q_2)q_2'+1\cdot(-q_3)q_3'+1\cdot(-q_4)q_4'\)
(Theorem~\ref{r8-thm-spectral}(b), Theorem~\ref{r1-thm-views}(d)). So \(A=PDQ'\) with \(Q=[q_1\ q_2\ q_3\ q_4]\),
\(P=[q_1\ {-q_2}\ {-q_3}\ {-q_4}]\) (both orthogonal) and \(D=\diag(7,1,1,1)\): the singular values are \(7,1,1,1\), the
absolute values of the eigenvalues.

(c) By Theorem~\ref{r12-thm-eckart-young}(a) (\(k=1\)), or by Corollary~\ref{r12-cor-symmetric} with \(d_1=7\),
\(d_2=d_3=d_4=-1\), \(A_1=7q_1q_1'=7\bar J_4=\frac74J_4\), and
\(\lVert A-A_1\rVert_F^2=\lVert C_4\rVert_F^2=\tr(C_4'C_4)=\tr C_4=4-1=3=1^2+1^2+1^2\) (\(C_4\) is symmetric and idempotent,
\(\tr\bar J_4=1\)).
\end{solution}

\begin{exercise}\label{r10-ex-low-rank}
Four stores sell three products; the sales (in hundreds) are
\[
B=\begin{pmatrix}5&2&0\\2&6&2\\0&2&7\\4&8&8\end{pmatrix}
\]
(row \(i\): store \(i\); column \(j\): product \(j\)). Let \(w_1=(1,2,2)'\), \(w_2=(2,1,-2)'\) and \(w_3=(2,-2,1)'\).
\begin{enumerate}
\item Compute \(Bw_1\), \(Bw_2\), \(Bw_3\), and deduce the singular values and a compact singular value decomposition of
\(B\).
\item Find the best approximation \(B_1\) of \(B\) by a matrix of rank at most \(1\), the errors \(\lVert B-B_1\rVert_2\) and
\(\lVert B-B_1\rVert_F\), and the fraction of \(\lVert B\rVert_F^2\) that \(B_1\) reproduces. Interpret \(B_1\).
\item Find the best approximation \(B_2\) of rank at most \(2\) and its errors.
\end{enumerate}
\end{exercise}
\begin{solution}
(a)
\begin{itemize}
\item \(Bw_1=(9,18,18,36)'=9(1,2,2,4)'\), \(Bw_2=(12,6,-12,0)'=6(2,1,-2,0)'\) and \(Bw_3=(6,-6,3,0)'=3(2,-2,1,0)'\).
\item The \(w_j\) are orthogonal with length \(3\), so the \(q_j=\frac13w_j\) form an orthogonal matrix \(Q\). The vectors
\(p_1=\frac15(1,2,2,4)'\), \(p_2=\frac13(2,1,-2,0)'\), \(p_3=\frac13(2,-2,1,0)'\) are orthonormal (lengths \(1\); inner
products \(\frac1{15}(2+2-4)\), \(\frac1{15}(2-4+2)\), \(\frac19(4-2-2)\), all \(0\)), and \(Bq_1=15p_1\), \(Bq_2=6p_2\),
\(Bq_3=3p_3\).
\item So \(BQ=P_1D\) with \(P_1=[p_1\ p_2\ p_3]\), \(D=\diag(15,6,3)\), and
\(Q'B'BQ=(P_1D)'(P_1D)=DP_1'P_1D=D^2\). By Theorem~\ref{r9-thm-svd}(a) (with \(Q_1=Q\), no \(Q_2\), and
\(P_1=BQD^{-1}\)), \(P'BQ=\left(\begin{smallmatrix}D\\0\end{smallmatrix}\right)\) for \(P=[p_1\ p_2\ p_3\ p_4]\), with any unit
vector \(p_4\) orthogonal to \(p_1,p_2,p_3\). By Theorem~\ref{r9-thm-svd-props}(a) the singular values are \(15\), \(6\), \(3\),
and \(B=BQQ'=15p_1q_1'+6p_2q_2'+3p_3q_3'\) (Theorem~\ref{r1-thm-views}(d)) is a compact SVD.
\end{itemize}
(b) By Theorem~\ref{r12-thm-eckart-young} (\(k=1\)),
\[
B_1=15p_1q_1'=(1,2,2,4)'(1,2,2)=\begin{pmatrix}1&2&2\\2&4&4\\2&4&4\\4&8&8\end{pmatrix},\qquad
B-B_1=\begin{pmatrix}4&0&-2\\0&2&-2\\-2&-2&3\\0&0&0\end{pmatrix}
\]
\(\lVert B-B_1\rVert_2=\sigma_2=6\) and \(\lVert B-B_1\rVert_F^2=36+9=45\), so \(\lVert B-B_1\rVert_F=3\sqrt5\); no matrix of rank
at most \(1\) comes closer in either norm. As \(\lVert B\rVert_F^2=225+36+9=270\) (the sum of the squared entries),
\(B_1\) reproduces \(\frac{225}{270}=\frac56\) of it. The entries of \(B_1\) are products \(a_ib_j\) with \(a=(1,2,2,4)'\) and
\(b=(1,2,2)'\): among all tables of the form ``store effect times product effect'' \(ab'\), \(B_1\) has the smallest sum of
squared differences from \(B\); store \(4\) follows it exactly.

(c) By Theorem~\ref{r12-thm-eckart-young} (\(k=2\)),
\[
B_2=B_1+6p_2q_2'=B_1+\frac23\begin{pmatrix}2\\1\\-2\\0\end{pmatrix}(2,1,-2)
=\frac13\begin{pmatrix}11&10&-2\\10&14&8\\-2&8&20\\12&24&24\end{pmatrix}
\]
with \(B-B_2=3p_3q_3'\), \(\lVert B-B_2\rVert_2=3\) and \(\lVert B-B_2\rVert_F^2=9\); \(B_2\) reproduces
\(\frac{261}{270}=\frac{29}{30}\) of \(\lVert B\rVert_F^2\).
\end{solution}

\begin{exercise}\label{r9-ex-svd-centered}
Write a compact SVD of the centered data matrix \(C_5X\), and give its singular values:
\begin{enumerate}
\item for the data of Exercise~\ref{r12-ex-pca-drill};
\item for the data of Exercise~\ref{r10-ex-pca-drill}.
\end{enumerate}
\end{exercise}
\begin{solution}
(a) In Exercise~\ref{r12-ex-pca-drill}, \(S\) has the eigenvalues \(15\), \(4\), \(3\), \(1\), with the loading vectors
\(v_1,\dots,v_4\) and the score vectors \(z_1,\dots,z_4\) found there. By Proposition~\ref{r12-prop-svd}(b), \(r=4\),
\(\sigma_m=\lVert z_m\rVert=\sqrt{4\lambda_m}\): \(\sigma_1=\sqrt{60}=2\sqrt{15}\), \(\sigma_2=\sqrt{16}=4\),
\(\sigma_3=\sqrt{12}=2\sqrt3\), \(\sigma_4=\sqrt4=2\), and \(p_m=z_m/\sigma_m\):
\[
p_1=\frac1{2\sqrt5}\begin{pmatrix}-4\\1\\1\\1\\1\end{pmatrix},\qquad
p_2=\frac1{2\sqrt3}\begin{pmatrix}0\\1\\1\\1\\-3\end{pmatrix},\qquad
p_3=\frac1{\sqrt2}\begin{pmatrix}0\\-1\\1\\0\\0\end{pmatrix},\qquad
p_4=\frac1{\sqrt6}\begin{pmatrix}0\\-1\\-1\\2\\0\end{pmatrix}
\]
So \(C_5X=2\sqrt{15}\,p_1v_1'+4p_2v_2'+2\sqrt3\,p_3v_3'+2p_4v_4'\), a compact SVD; for instance
\(2\sqrt{15}\,p_1v_1'=z_1v_1'=(-4,1,1,1,1)'(0,1,1,-1)\). Check: \(\sigma_1^2+\dots+\sigma_4^2=60+16+12+4=92=\lVert C_5X\rVert_F^2\)
(Theorem~\ref{r9-thm-svd-props}(d)).

(b) In Exercise~\ref{r10-ex-pca-drill}, \(S\) has the eigenvalues \(10\), \(6\), \(3\), \(1\). In the same way, \(r=4\),
\(\sigma_1=\sqrt{40}=2\sqrt{10}\), \(\sigma_2=\sqrt{24}=2\sqrt6\), \(\sigma_3=\sqrt{12}=2\sqrt3\), \(\sigma_4=\sqrt4=2\), and
\[
p_1=\frac1{\sqrt{30}}\begin{pmatrix}-2\\-2\\-2\\3\\3\end{pmatrix},\qquad
p_2=\frac1{\sqrt2}\begin{pmatrix}0\\-1\\1\\0\\0\end{pmatrix},\qquad
p_3=\frac1{\sqrt2}\begin{pmatrix}0\\0\\0\\-1\\1\end{pmatrix},\qquad
p_4=\frac1{\sqrt6}\begin{pmatrix}-2\\1\\1\\0\\0\end{pmatrix}
\]
from the scores of that exercise (\(z_{im}=v_m'(x_i-\bar x)\)); so \(C_5X=2\sqrt{10}\,p_1v_1'+2\sqrt6\,p_2v_2'+2\sqrt3\,p_3v_3'+2p_4v_4'\),
for instance \(2\sqrt{10}\,p_1v_1'=z_1v_1'=\frac13(-4,-4,-4,6,6)'(1,-1,0,1)\). Check:
\(40+24+12+4=80=\lVert C_5X\rVert_F^2\).
\end{solution}

\begin{exercise}\label{r10-dr-mp-svd}
The matrix
\[
A=\begin{pmatrix}2&1\\1&2\\1&-1\end{pmatrix}
\]
has the compact SVD \(A=3p_1q_1'+\sqrt3\,p_2q_2'\), with
\(p_1=\frac1{\sqrt2}(1,1,0)'\), \(p_2=\frac1{\sqrt6}(1,-1,2)'\), \(q_1=\frac1{\sqrt2}(1,1)'\) and \(q_2=\frac1{\sqrt2}(1,-1)'\)
(Exercise~\ref{r9-ex-svd-drill}).
\begin{enumerate}
\item Find \(A^+\) from the compact SVD.
\item Check the result with the formula for a matrix of full column rank.
\item For \(y=(3,0,0)'\), find \(A^+y\), the orthogonal projection \(AA^+y\) of \(y\) onto \(\mathcal C(A)\) and the residual
\(y-AA^+y\). Find also \((A')^+\) and \((2A)^+\).
\end{enumerate}
\end{exercise}
\begin{solution}
(a) By Theorem~\ref{r10-thm-mp-svd}, \(A^+=\frac13q_1p_1'+\frac1{\sqrt3}q_2p_2'\), where
\[
q_1p_1'=\frac12\begin{pmatrix}1&1&0\\1&1&0\end{pmatrix},\qquad
q_2p_2'=\frac1{\sqrt{12}}\begin{pmatrix}1&-1&2\\-1&1&-2\end{pmatrix}
\]
As \(\sqrt3\sqrt{12}=6\),
\[
A^+=\frac16\begin{pmatrix}1&1&0\\1&1&0\end{pmatrix}+\frac16\begin{pmatrix}1&-1&2\\-1&1&-2\end{pmatrix}
=\frac16\begin{pmatrix}2&0&2\\0&2&-2\end{pmatrix}=\frac13\begin{pmatrix}1&0&1\\0&1&-1\end{pmatrix}
\]
(b) \(\rank A=2\), the number of columns, so \(A^+=(A'A)^{-1}A'\) (Theorem~\ref{r10-thm-mp-exists}(c)). Here
\(A'A=\left(\begin{smallmatrix}6&3\\3&6\end{smallmatrix}\right)\), with determinant \(27\), so
\((A'A)^{-1}=\frac1{27}\left(\begin{smallmatrix}6&-3\\-3&6\end{smallmatrix}\right)=\frac19\left(\begin{smallmatrix}2&-1\\-1&2\end{smallmatrix}\right)\)
(Theorem~\ref{r5d-thm-adjugate}), and
\[
(A'A)^{-1}A'=\frac19\begin{pmatrix}2&-1\\-1&2\end{pmatrix}\begin{pmatrix}2&1&1\\1&2&-1\end{pmatrix}
=\frac19\begin{pmatrix}3&0&3\\0&3&-3\end{pmatrix}=\frac13\begin{pmatrix}1&0&1\\0&1&-1\end{pmatrix}
\]
\needspace{5\baselineskip}
(c)
\begin{itemize}
\item \(A^+y=\frac13(3,0)'=(1,0)'\), and \(AA^+y=A(1,0)'=(2,1,1)'\), the orthogonal projection of \(y\) onto \(\mathcal C(A)\)
(Proposition~\ref{r10-prop-mp-props}(b)).
\item The residual is \(y-AA^+y=(1,-1,-1)'\); check: \(A'(1,-1,-1)'=(2-1-1,\ 1-2+1)'=0\), so it is orthogonal to
\(\mathcal C(A)\).
\item By Proposition~\ref{r10-prop-mp-props}(c),
\[
(A')^+=(A^+)'=\frac13\begin{pmatrix}1&0\\0&1\\1&-1\end{pmatrix},\qquad
(2A)^+=\frac12A^+=\frac16\begin{pmatrix}1&0&1\\0&1&-1\end{pmatrix}
\]
\end{itemize}
\end{solution}

\begin{exercise}\label{r12-ex-cca}
The sample covariance matrix of four variables is
\[
S=\begin{pmatrix}4&0&2&1\\0&4&1&2\\2&1&4&0\\1&2&0&4\end{pmatrix}
\]
Find the canonical correlations and canonical pairs between the first two and the last two variables.
\end{exercise}
\begin{solution}
\(S_{11}=S_{22}=4I_2\), so \(S_{11}^{-1/2}=S_{22}^{-1/2}=\frac12I_2\) (Theorem~\ref{r8-thm-sqrt}: \(2I_2\) is PSD and
\((2I_2)^2=4I_2\)). Then
\[
K=\tfrac12S_{12}\cdot\tfrac12=\frac14\begin{pmatrix}2&1\\1&2\end{pmatrix}=\frac14(3\bar J_2+C_2)=\frac34u_1u_1'+\frac14u_2u_2',
\qquad u_1=\frac1{\sqrt2}\begin{pmatrix}1\\1\end{pmatrix},\ u_2=\frac1{\sqrt2}\begin{pmatrix}1\\-1\end{pmatrix}
\]
since \(\bar J_2=u_1u_1'\) and \(C_2=I_2-\bar J_2=u_2u_2'\). This is a singular value decomposition with
\(\rho_1=\frac34\), \(\rho_2=\frac14\) and \(w_k=u_k\). By Theorem~\ref{r12-thm-cca}, the canonical correlations are
\(\frac34\) and \(\frac14\), with \(a_1=b_1=\frac12u_1=\frac1{2\sqrt2}(1,1)'\) and \(a_2=b_2=\frac1{2\sqrt2}(1,-1)'\). Check:
\(a_1'S_{11}a_1=4\cdot\frac28=1\) and \(a_1'S_{12}b_1=\frac18(1,1)\begin{pmatrix}2&1\\1&2\end{pmatrix}\begin{pmatrix}1\\1\end{pmatrix}=\frac68=\frac34\).
\end{solution}

\begin{exercise}\label{r10-ex-procrustes}
Let \(A\) be the matrix of Exercise~\ref{r10-ex-symmetric-svd}.
\begin{enumerate}
\item Find the largest value of \(\tr(U'A)\) over the orthogonal \(U\in\R^{4\times4}\), and an orthogonal matrix \(U_0\) at
which it is attained. Compare with \(U=I_4\).
\item Find an orthogonal matrix nearest to \(A\) in the Frobenius norm, and the distance. Compare with
\(\lVert I_4-A\rVert_F\).
\end{enumerate}
\end{exercise}
\begin{solution}
(a) By Theorem~\ref{r10-thm-procrustes}(a) and the singular value decomposition \(A=PDQ'\) of
Exercise~\ref{r10-ex-symmetric-svd}, the largest value is \(7+5+3+1=16\), attained at
\[
U_0=PQ'=\sum_{k=1}^4p_kq_k'=\frac14\bigl(h_1h_1'-h_2h_2'+h_3h_3'-h_4h_4'\bigr)
=\begin{pmatrix}0&1&0&0\\1&0&0&0\\0&0&0&1\\0&0&1&0\end{pmatrix}
\]
(the \((i,j)\) entry is \(\frac14\) times the sum of the products of the entries \(i\) and \(j\) of the \(h_k\), with signs
\(+,-,+,-\)). Indeed \(\tr(U_0'A)=a_{21}+a_{12}+a_{43}+a_{34}=16\), whereas \(\tr(I_4'A)=\tr A=4\).

(b) Theorem~\ref{r10-thm-procrustes}(b) with \(B=I_4\) and \(C=A\) (so \(CB'=A\)) gives
\(\lVert U-A\rVert_F^2\geqslant\lVert I_4\rVert_F^2+\lVert A\rVert_F^2-2\cdot16=4+84-32=56\) for every orthogonal \(U\), with
equality at \(U=U_0\). So \(U_0\) is nearest, at distance \(\sqrt{56}=2\sqrt{14}\); directly, \(U_0-A\) has four entries
\(-1\), four entries \(-3\) and four entries \(-2\), and \(4+36+16=56\). By comparison,
\(\lVert I_4-A\rVert_F^2=4\cdot16+4\cdot4=80\).
\end{solution}

\subsection*{Problems}

\begin{problem}[\citeat{abadir-magnus}{Exercise~12.45}]\label{r9-ex-monotone-strict}
Let \(A,B\in\R^{n\times n}\) be symmetric, with \(B\) PD. Show that \(\lambda_k(A+B)>\lambda_k(A)\) for
\(k=1,\dots,n\).
\end{problem}
\begin{solution}
\begin{itemize}
\item \(B\) is PD, so \(\lambda_n(B)>0\) (Theorem~\ref{r8-thm-eigen-criterion}(a)), and \(u'Bu\geqslant\lambda_n(B)\) for
every unit vector \(u\) (Theorem~\ref{r9-thm-rayleigh}(a)).
\item \(k=n\): for every unit vector \(u\), \(u'(A+B)u=u'Au+u'Bu\geqslant\lambda_n(A)+\lambda_n(B)\)
(Theorem~\ref{r9-thm-rayleigh}(a)); at a unit eigenvector of \(A+B\) belonging to \(\lambda_n(A+B)\) this gives
\(\lambda_n(A+B)\geqslant\lambda_n(A)+\lambda_n(B)\).
\item \(k\leqslant n-1\): let \(w_1,\dots,w_{n-k}\in\R^n\). For every unit vector \(u\perp w_1,\dots,w_{n-k}\),
\(u'(A+B)u\geqslant u'Au+\lambda_n(B)\), so the minimum of \(u'(A+B)u\) over these \(u\) is at least the minimum of
\(u'Au\) plus \(\lambda_n(B)\). Taking the maximum over \(w_1,\dots,w_{n-k}\) on both sides,
Theorem~\ref{r9-thm-courant-fischer}(b) for \(A+B\) and for \(A\) gives \(\lambda_k(A+B)\geqslant\lambda_k(A)+\lambda_n(B)\).
\item In both cases \(\lambda_k(A+B)\geqslant\lambda_k(A)+\lambda_n(B)>\lambda_k(A)\).
\end{itemize}
\end{solution}

\begin{problem}[\citeat{bapat}{Ch.~5, Exercise~17}]\label{r9-ex-minors-sign}
Let \(A\in\R^{n\times n}\) (\(n\geqslant2\)) be symmetric with \(\det A=0\). Show that the \(n\) minors
\(\det A(1|1),\dots,\det A(n|n)\) are either all \(\geqslant0\) or all \(\leqslant0\).
\end{problem}
\begin{solution}
\begin{itemize}
\item \(\det A=\lambda_1\cdots\lambda_n=0\) (Theorem~\ref{r7-thm-trace-det}), so \(0\) is an eigenvalue of \(A\). Let \(p\)
be the number of positive eigenvalues of \(A\); then \(p\leqslant n-1\) and \(\lambda_{p+1}(A)=0\).
\item Let \(B=A(i|i)\). By Theorem~\ref{r9-thm-interlacing},
\(\lambda_k(B)\geqslant\lambda_{k+1}(A)\geqslant\lambda_{p+1}(A)=0\) for \(k\leqslant p\), and
\(\lambda_k(B)\leqslant\lambda_k(A)\leqslant\lambda_{p+1}(A)=0\) for \(p+1\leqslant k\leqslant n-1\).
\item If \(0\) has multiplicity at least \(2\), then also \(\lambda_{p+2}(A)=0\), and
\(0=\lambda_{p+1}(A)\geqslant\lambda_{p+1}(B)\geqslant\lambda_{p+2}(A)=0\). So every \(A(i|i)\) has the eigenvalue \(0\), and
\(\det A(i|i)=0\) for every \(i\).
\item Let \(0\) be a simple eigenvalue, and let \(\det B\ne0\). Then no eigenvalue of \(B\) is \(0\), so by the second item
\(B\) has exactly \(p\) positive and \(n-1-p\) negative eigenvalues, and \(\det B\) has the sign \((-1)^{n-1-p}\)
(Theorem~\ref{r7-thm-trace-det}). So all nonzero minors \(\det A(i|i)\) have this same sign.
\end{itemize}
\end{solution}

\begin{problem}[\citeat{bapat}{Ch.~5, Exercise~4}]\label{r9-ex-sum-squares-sv}
Let \(n\) be even, and let \(A\in\R^{n\times n}\) have \(n^2/2\) entries equal to \(2\) and the other entries equal to
\(4\). Show that the sum of the squares of the singular values of \(A\) is \(10n^2\).
\end{problem}
\begin{solution}
By Theorem~\ref{r9-thm-svd-props}(d) and Definition~\ref{r9-def-frobenius},
\(\sigma_1^2+\dots+\sigma_r^2=\lVert A\rVert_F^2=\sum_i\sum_ja_{ij}^2=\frac{n^2}2\cdot2^2+\frac{n^2}2\cdot4^2
=2n^2+8n^2=10n^2\).
\end{solution}

\begin{problem}[\citeat{prasolov}{Problem~34.6}]\label{r9-ex-pr-equal-sv}
Let \(A\in\R^{n\times n}\), and let all singular values \(\sigma_1,\dots,\sigma_n\) of \(A\) be equal to the same number \(c\).
Show that \(A=cU\) for an orthogonal matrix \(U\in\R^{n\times n}\).
\end{problem}
\begin{solution}
\begin{itemize}
\item The singular values are \(\geqslant0\) (the definitions after Theorem~\ref{r9-thm-svd}), so \(c\geqslant0\).
\item Let \(c=0\). Then \(\sigma_1=0\), so \(A=0\): if \(A\ne0\), then \(r=\rank A\geqslant1\) and \(\sigma_1\geqslant\sigma_r>0\)
(Theorem~\ref{r9-thm-svd}(b) and the definitions after it). So \(A=0\cdot I_n\).
\item Let \(c>0\). As \(\sigma_k=0\) for \(\rank A<k\leqslant n\) and no \(\sigma_k\) is \(0\), \(\rank A=n\). By
Theorem~\ref{r9-thm-svd-props}(a), applied to an SVD of \(A\) (Theorem~\ref{r9-thm-svd}(b)), the \(n\) eigenvalues of \(A'A\),
listed with multiplicity, are \(\sigma_1^2,\dots,\sigma_n^2\), that is, \(c^2\) (\(n\) times).
\item \(A'A\) is symmetric (Theorem~\ref{r1-thm-transpose}). By Theorem~\ref{r8-thm-spectral}, \(A'A=V\Lambda V'\) with \(V\)
orthogonal and \(\Lambda\) the diagonal matrix of the eigenvalues of \(A'A\), here \(\Lambda=c^2I_n\); so
\(A'A=c^2VV'=c^2I_n\) (Proposition~\ref{r6-prop-orthonormal}(b)).
\item Put \(U=c^{-1}A\). Then \(U'U=c^{-2}A'A=I_n\), so \(U\) is orthogonal (Definition~\ref{r6-def-orthonormal}), and
\(A=cU\).
\end{itemize}
\end{solution}

\begin{problem}[\citeat{prasolov}{Problem~34.8}]\label{r9-ex-pr-block-sv}
Let \(A\in\R^{n\times n}\) have the singular values \(\sigma_1\geqslant\dots\geqslant\sigma_n\) (with \(\sigma_k=0\) for
\(k>\rank A\), as in the definitions after Theorem~\ref{r9-thm-svd}). Show that the eigenvalues of
\[
M=\begin{pmatrix}0&A\\A'&0\end{pmatrix}\in\R^{2n\times2n}
\]
listed with multiplicity, are \(\sigma_1,\dots,\sigma_n,-\sigma_1,\dots,-\sigma_n\).
\end{problem}
\begin{solution}
\begin{itemize}
\item For \(t\in\R\), \(p_M(t)=\det(tI_{2n}-M)\) (Definition~\ref{r7-def-charpoly}(a)), and
\(tI_{2n}-M=\left(\begin{smallmatrix}tI_n&-A\\-A'&tI_n\end{smallmatrix}\right)\).
\item Let \(t\ne0\). Then \(tI_n\) is invertible with inverse \(t^{-1}I_n\), and \(\det(tI_n)=t^n\)
(Theorem~\ref{r5d-thm-basic}(c)). By Theorem~\ref{r5d-thm-schur-det}(a), and \(\det(cB)=c^n\det B\)
(Theorem~\ref{r5d-thm-product}(c)) with \(c=t^{-1}\),
\[
p_M(t)=t^n\det\bigl(tI_n-(-A')(t^{-1}I_n)(-A)\bigr)=t^n\det\bigl(t^{-1}(t^2I_n-A'A)\bigr)=t^nt^{-n}\det(t^2I_n-A'A)
\]
that is, \(p_M(t)=\det(t^2I_n-A'A)\) for every \(t\ne0\).
\item The eigenvalues of \(A'A\), listed with multiplicity, are \(\sigma_1^2,\dots,\sigma_n^2\). If \(A=0\), then \(A'A=0\),
\(p_{A'A}(t)=\det(tI_n)=t^n\) (Theorem~\ref{r5d-thm-basic}(c)), and all of them, like all \(\sigma_i\), are \(0\). If \(A\ne0\), then by Theorem~\ref{r9-thm-svd-props}(a), applied to an SVD of
\(A\) (Theorem~\ref{r9-thm-svd}(b)), \(\sigma_1^2,\dots,\sigma_r^2\) (\(r=\rank A\)) are the nonzero ones, so the other
\(n-r\) are \(0=\sigma_{r+1}^2=\dots=\sigma_n^2\). Hence, for every \(s\in\R\) (Definitions~\ref{r7-def-charpoly}(a)
and~\ref{r7-def-eigenvalues}, Proposition~\ref{r7-prop-poly}(b)),
\[
\det(sI_n-A'A)=p_{A'A}(s)=(s-\sigma_1^2)(s-\sigma_2^2)\cdots(s-\sigma_n^2)
\]
\item Take \(s=t^2\), and write \(t^2-\sigma_i^2=(t-\sigma_i)(t+\sigma_i)\). Then, for every \(t\ne0\), \(p_M(t)=g(t)\), where
\[
g=(t-\sigma_1)\cdots(t-\sigma_n)\bigl(t-(-\sigma_1)\bigr)\cdots\bigl(t-(-\sigma_n)\bigr)
\]
(Proposition~\ref{r7-prop-poly}(b)). The polynomials \(p_M\) and \(g\) have the same value at every real \(t\ne0\), so \(p_M=g\)
(Proposition~\ref{r7-prop-poly}(d)). By Definition~\ref{r7-def-eigenvalues}, the eigenvalues of \(M\), listed with
multiplicity, are \(\sigma_1,\dots,\sigma_n,-\sigma_1,\dots,-\sigma_n\).
\end{itemize}
\end{solution}

\begin{problem}[\citeat{prasolov}{Problem~35.1}]\label{r9-ex-pr-eig-norm}
Let \(A\in\R^{n\times n}\) be invertible, and let \(c\) be a real eigenvalue of \(A\). Show that
\[
\frac1{\lVert A^{-1}\rVert_2}\leqslant\lvert c\rvert\leqslant\lVert A\rVert_2
\]
\end{problem}
\begin{solution}
\begin{itemize}
\item Let \(Ax=cx\) with \(x\ne0\) (Definition~\ref{r7-def-eigen}). As \(A\) is invertible, \(c\ne0\) and \(A^{-1}x=c^{-1}x\)
(Proposition~\ref{r7-prop-eig-poly}(b)).
\item By Proposition~\ref{r9-prop-spectral-norm}(a) and \(\lVert cx\rVert=\lvert c\rvert\,\lVert x\rVert\) (the remarks after
Definition~\ref{r6-def-inner}),
\(\lVert A\rVert_2=\max_{y\ne0}\lVert Ay\rVert/\lVert y\rVert\geqslant\lVert Ax\rVert/\lVert x\rVert
=\lVert cx\rVert/\lVert x\rVert=\lvert c\rvert\).
\item Likewise
\(\lVert A^{-1}\rVert_2\geqslant\lVert A^{-1}x\rVert/\lVert x\rVert=\lVert c^{-1}x\rVert/\lVert x\rVert=1/\lvert c\rvert>0\).
Taking reciprocals, \(1/\lVert A^{-1}\rVert_2\leqslant\lvert c\rvert\).
\end{itemize}
As \(\lVert A\rVert_2=\sigma_1\) and \(\lVert A^{-1}\rVert_2=1/\sigma_n\) (Definition~\ref{r9-def-spectral-norm},
Proposition~\ref{r9-prop-condition}(a)), this says \(\sigma_n\leqslant\lvert c\rvert\leqslant\sigma_1\); and for two real
eigenvalues \(c\), \(d\) of \(A\), \(\lvert c\rvert/\lvert d\rvert\leqslant\lVert A\rVert_2\lVert A^{-1}\rVert_2=\kappa(A)\)
(Definition~\ref{r9-def-condition}).
\end{solution}

\begin{problem}[\citeat{bapat}{Ch.~5, Exercise~15}]\label{r12-ex-bapat-pd}
Let \(A\in\R^{m\times m}\) and \(B\in\R^{q\times q}\) be PD and \(C\in\R^{m\times q}\). Show that
\[
F=\begin{pmatrix}A&C\\C'&B\end{pmatrix}
\]
is PD if and only if the largest singular value \(\sigma_1\) of \(K=A^{-1/2}CB^{-1/2}\) is less than \(1\) (with
\(\sigma_1=0\) if \(C=0\)).
\end{problem}
\begin{solution}
\begin{itemize}
\item \(A^{-1/2}\), \(B^{-1/2}\) are symmetric and \(B^{-1/2}B^{-1/2}=(B^{1/2}B^{1/2})^{-1}=B^{-1}\) (Theorem~\ref{r8-thm-sqrt}), so
\(KK'=A^{-1/2}CB^{-1}C'A^{-1/2}\), and \(\sigma_1^2\) is the largest eigenvalue of \(KK'\) (Theorem~\ref{r9-thm-svd-props}(a); \(0\)
if \(C=0\)).
\item With \(KK'=Q\Lambda Q'\) (Theorem~\ref{r8-thm-spectral}(b)), \(I_m-KK'=Q(I_m-\Lambda)Q'\) has the eigenvalues \(1-\lambda_i\);
so \(\sigma_1<1\) if and only if \(I_m-KK'\) is PD (Theorem~\ref{r8-thm-eigen-criterion}(a)).
\item \(A^{1/2}(I_m-KK')A^{1/2}=A-CB^{-1}C'\), and \(x\mapsto y=A^{1/2}x\) maps the nonzero vectors onto the nonzero vectors
(\(A^{1/2}\) is invertible), with \(x'(A-CB^{-1}C')x=y'(I_m-KK')y\). So \(I_m-KK'\) is PD if and only if
\(A-CB^{-1}C'\) is PD.
\item Let \(L=\begin{pmatrix}I_m&-CB^{-1}\\0&I_q\end{pmatrix}\), invertible (Proposition~\ref{r5b-prop-block-elementary}). By
Theorem~\ref{r5b-thm-block}(c), with \((B^{-1})'=B^{-1}\),
\[
LF=\begin{pmatrix}A-CB^{-1}C'&0\\C'&B\end{pmatrix},\qquad
LFL'=\begin{pmatrix}A-CB^{-1}C'&0\\0&B\end{pmatrix}
\]
As \(x'(LFL')x=(L'x)'F(L'x)\) and \(L'\) is invertible, \(F\) is PD if and only if \(LFL'\) is PD. For
\(x=(y',z')'\), \(x'(LFL')x=y'(A-CB^{-1}C')y+z'Bz\); as \(B\) is PD, this is positive for every \(x\ne0\) if and only if
\(A-CB^{-1}C'\) is PD (take \(z=0\) for the converse).
\item Combining the three equivalences: \(F\) is PD if and only if \(\sigma_1<1\).
\end{itemize}
\end{solution}

\begin{problem}[\citeat{prasolov}{Problem~34.2}]\label{r9-ex-pr-trace-orthogonal}
Let \(U\in\R^{n\times n}\) be orthogonal and \(S\in\R^{n\times n}\) PSD.
\begin{enumerate}
\item Show that \(\lvert\tr(US)\rvert\leqslant\tr S\).
\item Let \(S\) be PD. Show that \(\tr(US)=\tr S\) only for \(U=I_n\), and \(\lvert\tr(US)\rvert=\tr S\) only for \(U=I_n\) and
\(U=-I_n\).
\end{enumerate}
\end{problem}
\begin{solution}
(a)
\begin{itemize}
\item By Theorem~\ref{r8-thm-spectral}, \(S=V\Lambda V'\) with \(V\in\R^{n\times n}\) orthogonal and
\(\Lambda=\diag(\lambda_1,\dots,\lambda_n)\), where \(\lambda_1,\dots,\lambda_n\) are the eigenvalues of \(S\); they are
\(\geqslant0\), as \(S\) is PSD (Theorem~\ref{r8-thm-eigen-criterion}(a)).
\item Let \(W=V'UV=(w_{ij})\). By Theorem~\ref{r1-thm-trace}(b),
\(\tr(US)=\tr\bigl((UV\Lambda)V'\bigr)=\tr(V'UV\Lambda)=\tr(W\Lambda)\). Column \(i\) of \(W\Lambda\) is \(\lambda_i\col_i(W)\)
(Theorem~\ref{r1-thm-views}(a),(b)), whose \(i\)th entry is \(\lambda_iw_{ii}\); so
\[
\tr(US)=\lambda_1w_{11}+\lambda_2w_{22}+\dots+\lambda_nw_{nn}
\]
\item \(W\) is orthogonal, as \(V'\), \(U\) and \(V\) are (Theorem~\ref{r6-thm-orthogonal}(b)), so its columns are unit vectors
(Theorem~\ref{r6-thm-orthogonal}(a)): \(w_{1i}^2+\dots+w_{ni}^2=1\). Hence \(w_{ii}^2\leqslant1\), that is,
\(\lvert w_{ii}\rvert\leqslant1\), for \(i=1,\dots,n\).
\item \(\tr S=\tr\bigl((V\Lambda)V'\bigr)=\tr(V'V\Lambda)=\tr\Lambda=\lambda_1+\dots+\lambda_n\)
(Theorem~\ref{r1-thm-trace}(b), \(V'V=I_n\)). As every \(\lambda_i\geqslant0\),
\[
\lvert\tr(US)\rvert\leqslant\lambda_1\lvert w_{11}\rvert+\dots+\lambda_n\lvert w_{nn}\rvert\leqslant\lambda_1+\dots+\lambda_n=\tr S
\]
\end{itemize}
(b) Now every \(\lambda_i>0\) (Theorem~\ref{r8-thm-eigen-criterion}(a)).
\begin{itemize}
\item Let \(\tr(US)=\tr S\). By (a), \(0=\tr S-\tr(US)=\lambda_1(1-w_{11})+\dots+\lambda_n(1-w_{nn})\), a sum of terms
\(\geqslant0\), as \(w_{ii}\leqslant1\). So every term is \(0\), and as \(\lambda_i>0\), \(w_{ii}=1\) for \(i=1,\dots,n\).
\item Column \(i\) of \(W\) is a unit vector with \(i\)th entry \(1\), so its other entries are \(0\): \(\col_i(W)=e_i\), and
\(W=I_n\). As \(VV'=I_n\) (Proposition~\ref{r6-prop-orthonormal}(b)), \(U=VV'UVV'=VWV'=VV'=I_n\).
\item Let \(\lvert\tr(US)\rvert=\tr S\). If \(\tr(US)=\tr S\), then \(U=I_n\). If \(\tr(US)=-\tr S\), then \(-U\) is orthogonal,
as \((-U)'(-U)=U'U=I_n\), and \(\tr\bigl((-U)S\bigr)=-\tr(US)=\tr S\) (Theorem~\ref{r1-thm-trace}(a)); so \(-U=I_n\), that is,
\(U=-I_n\).
\end{itemize}
\end{solution}

\begin{problem}[\citeat{abadir-magnus}{Exercise~8.40}]\label{r9-ex-polar}
Let \(A\in\R^{m\times n}\).
\begin{enumerate}
\item If \(A\) has full column rank \(n\), show that \(A=UV\) for some \(U\in\R^{m\times n}\) with \(U'U=I_n\) and some PD
\(V\in\R^{n\times n}\).
\item If \(A\) has full row rank \(m\), show that \(A=VU'\) for some \(U\in\R^{n\times m}\) with \(U'U=I_m\) and some PD
\(V\in\R^{m\times m}\).
\end{enumerate}
\end{problem}
\begin{solution}
(a)
\begin{itemize}
\item Here \(r=n\), so in the compact SVD \(A=P_1DQ_1'\) of Theorem~\ref{r9-thm-svd}, \(Q_1=Q\) is an orthogonal
\(n\times n\) matrix, \(P_1'P_1=I_n\), and \(D=\diag(\sigma_1,\dots,\sigma_n)\) with every \(\sigma_i>0\).
\item As \(Q'Q=I_n\), \(A=P_1(Q'Q)DQ'=(P_1Q')(QDQ')=UV\) with \(U=P_1Q'\) and \(V=QDQ'\).
\item \(U'U=QP_1'P_1Q'=QQ'=I_n\) (Proposition~\ref{r6-prop-orthonormal}(b)). \(V\) is symmetric and \(Q'VQ=D\), so its
eigenvalues are \(\sigma_1,\dots,\sigma_n>0\) (Definition~\ref{r7-def-diagonalizable}), and \(V\) is PD
(Theorem~\ref{r8-thm-eigen-criterion}(a)).
\end{itemize}
(b) \(A'\in\R^{n\times m}\) has full column rank \(m\) (Theorem~\ref{thm:r4-rank-rules}(a)). By (a), \(A'=UV\) with
\(U'U=I_m\) and \(V\) PD, so \(A=V'U'=VU'\).
\end{solution}

\clearpage
\def\unitnumber{16}
\setcounter{theorem}{0}\setcounter{problem}{0}
\section*{Chapter 16. Kronecker products, vec, and matrix derivatives}

\subsection*{The Kronecker product}

\begin{definition}[\citeat{tu-mc}{\S2.3}, \citeat{harville}{\S16.1}, \citeat{searle-khuri}{\S4.11}]\label{r13-def-kronecker}
Let \(A=(a_{ij})\in\R^{m\times n}\) and \(B\in\R^{p\times q}\). The \textbf{Kronecker product} of \(A\) and \(B\) is the
\(mp\times nq\) matrix
\[
A\otimes B=\begin{pmatrix}a_{11}B&\cdots&a_{1n}B\\\vdots&&\vdots\\a_{m1}B&\cdots&a_{mn}B\end{pmatrix}
\]
partitioned into \(m\) rows and \(n\) columns of \(p\times q\) blocks, the \((i,j)\) block being \(a_{ij}B\). Thus the
entry of \(A\otimes B\) in row \((i-1)p+r\) and column \((j-1)q+s\) is \(a_{ij}b_{rs}\) (\(1\leqslant r\leqslant p\),
\(1\leqslant s\leqslant q\)). For column vectors \(a\in\R^m\) and \(b\in\R^p\), \(a\otimes b=(a_1b',\ \dots,\ a_mb')'\in\R^{mp}\).
For a number \(c\) (a \(1\times1\) matrix), \(c\otimes B=B\otimes c=cB\).
\end{definition}
Unlike \(AB\), the product \(A\otimes B\) is defined for matrices of any sizes.

\begin{example}[\citeat{searle-khuri}{\S4.11}]\label{r13-exm-kronecker}
With \(A=(1,\ 2,\ 3)\) and \(B=\begin{pmatrix}6&7\\8&9\end{pmatrix}\),
\[
A\otimes B=\begin{pmatrix}6&7&12&14&18&21\\8&9&16&18&24&27\end{pmatrix},\qquad
B\otimes A=\begin{pmatrix}6&12&18&7&14&21\\8&16&24&9&18&27\end{pmatrix}
\]
Both contain the twelve products \(a_jb_{rs}\), in different orders, so \(A\otimes B\ne B\otimes A\). Two Kronecker products
with the summing vector (Definition~\ref{r1-def-ones}) occur in designs:
\[
I_2\otimes\one_3=\begin{pmatrix}1&0\\1&0\\1&0\\0&1\\0&1\\0&1\end{pmatrix},\qquad
\one_3\otimes I_2=\begin{pmatrix}1&0\\0&1\\1&0\\0&1\\1&0\\0&1\end{pmatrix}
\]
\end{example}

\begin{theorem}[\citeat{tu-mc}{Propositions~3, 4, Exercise~1}, \citeat{abadir-magnus}{Exercises~10.2--10.9}, \citeat{harville-solutions}{Ch.~16, Exercise~10}, \citeat{harville}{Lemma~16.1.2}]\label{r13-thm-kron-rules}
Let the sizes of the matrices be such that the operations below are defined, and let \(c\in\R\).
\begin{enumerate}
\item \((cA)\otimes B=A\otimes(cB)=c(A\otimes B)\), and \(A\otimes B=0\) if \(A=0\) or \(B=0\). For column vectors \(a\in\R^m\)
and \(b\in\R^n\), \(a\otimes b'=ab'=b'\otimes a\).
\item \((A_1+A_2)\otimes B=A_1\otimes B+A_2\otimes B\) and \(A\otimes(B_1+B_2)=A\otimes B_1+A\otimes B_2\).
\item If \(A\) is partitioned into blocks \(A_{kl}\) (Definition~\ref{r5b-def-partition}), then \(A\otimes B\) is partitioned
into the blocks \(A_{kl}\otimes B\), in the same arrangement; for instance \([A_1\ A_2]\otimes B=[A_1\otimes B\ \ A_2\otimes B]\).
\item \((A\otimes B)\otimes C=A\otimes(B\otimes C)\); this matrix is written \(A\otimes B\otimes C\).
\item \((A\otimes B)'=A'\otimes B'\).
\item (\textbf{Mixed product}) For \(A\in\R^{m\times n}\), \(B\in\R^{p\times q}\), \(C\in\R^{n\times s}\) and \(D\in\R^{q\times t}\),
\[
(A\otimes B)(C\otimes D)=(AC)\otimes(BD)
\]
\end{enumerate}
\end{theorem}
\begin{proof}
Throughout, \(A=(a_{ij})\in\R^{m\times n}\) and \(B\in\R^{p\times q}\), and we compare matrices block by block
(Definition~\ref{r5b-def-partition}).

(a)
\begin{itemize}
\item The \((i,j)\) blocks of \((cA)\otimes B\), \(A\otimes(cB)\) and \(c(A\otimes B)\) are \((ca_{ij})B\), \(a_{ij}(cB)\) and
\(c(a_{ij}B)\), which are equal. If \(A=0\) or \(B=0\), every block \(a_{ij}B\) is \(0\).
\item \(a\otimes b'\) has the rows \(a_1b',\dots,a_mb'\), which are the rows of \(ab'\) (Theorem~\ref{r1-thm-views}(c)).
\(b'\otimes a=(b_1a,\ \dots,\ b_na)\) has the columns \(b_ja\), which are the columns of \(ab'\)
(Theorem~\ref{r1-thm-views}(a)).
\end{itemize}

(b) The \((i,j)\) block of \((A_1+A_2)\otimes B\) is \((a^{(1)}_{ij}+a^{(2)}_{ij})B=a^{(1)}_{ij}B+a^{(2)}_{ij}B\), the sum of the
\((i,j)\) blocks of \(A_1\otimes B\) and \(A_2\otimes B\); and \(a_{ij}(B_1+B_2)=a_{ij}B_1+a_{ij}B_2\).

(c)
\begin{itemize}
\item Let the row groups of \(A\) have \(m_1,\dots,m_r\) rows and the column groups \(n_1,\dots,n_c\) columns. Row \(i\) of
\(A\) gives the block row of \(A\otimes B\) formed by rows \((i-1)p+1,\dots,ip\), and column \(j\) the block column formed
by columns \((j-1)q+1,\dots,jq\). So the rows of \(A\) in group \(k\) give \(m_kp\) consecutive rows of \(A\otimes B\), and the
columns in group \(l\) give \(n_lq\) consecutive columns.
\item The submatrix of \(A\otimes B\) in these rows and columns consists of the blocks \(a_{ij}B\) with row \(i\) in group
\(k\) and column \(j\) in group \(l\), arranged as the entries \(a_{ij}\) are arranged in \(A_{kl}\). By
Definition~\ref{r13-def-kronecker}, it is \(A_{kl}\otimes B\).
\end{itemize}

(d) Partition \(A\otimes B\) into its blocks \(a_{ij}B\). By (c), \((A\otimes B)\otimes C\) is partitioned into the blocks
\((a_{ij}B)\otimes C\), which equal \(a_{ij}(B\otimes C)\) by (a). These are the blocks of \(A\otimes(B\otimes C)\)
(Definition~\ref{r13-def-kronecker}).

(e) By Theorem~\ref{r5b-thm-block}(b), \((A\otimes B)'\) has \(n\) block rows and \(m\) block columns, and its \((j,i)\) block
is \((a_{ij}B)'=a_{ij}B'\). This is the \((j,i)\) block of \(A'\otimes B'\), since the \((j,i)\) entry of \(A'\) is \(a_{ij}\).

(f)
\begin{itemize}
\item \(A\otimes B\) has \(n\) block columns of width \(q\), and \(C\otimes D\) has \(n\) block rows of height \(q\): the partitions
are conformal for multiplication. By Theorem~\ref{r5b-thm-block}(c), the \((i,k)\) block of the product is
\[
\sum_{j=1}^n(a_{ij}B)(c_{jk}D)=\Bigl(\sum_{j=1}^na_{ij}c_{jk}\Bigr)BD=(AC)_{ik}\,BD
\]
\item This is the \((i,k)\) block of \((AC)\otimes(BD)\), for \(i=1,\dots,m\) and \(k=1,\dots,s\).\qedhere
\end{itemize}
\end{proof}

\begin{corollary}[\citeat{harville}{\S16.1}, \citeat{tu-mc}{Proposition~4}, \citeat{abadir-magnus}{Exercise~10.8}, \citeat{harville-solutions}{Ch.~16, Exercise~4}]\label{r13-cor-kron}
Let \(A\in\R^{m\times n}\) and \(B\in\R^{p\times q}\).
\begin{enumerate}
\item \(A\otimes B=(A\otimes I_p)(I_n\otimes B)=(I_m\otimes B)(A\otimes I_q)\); \(I_m\otimes I_p=I_{mp}\); \(I_m\otimes B\) is block
diagonal with \(m\) diagonal blocks \(B\); and for a diagonal \(D=\diag(d_1,\dots,d_m)\), \(D\otimes B=\diag(d_1B,\dots,d_mB)\).
\item If \(A\) and \(B\) are square and invertible, then \(A\otimes B\) is invertible and \((A\otimes B)^{-1}=A^{-1}\otimes B^{-1}\).
\item If \(A\) and \(B\) are square, then \(\tr(A\otimes B)=\tr A\cdot\tr B\).
\item If \(G\) is a g-inverse of \(A\) and \(F\) a g-inverse of \(B\) (Definition~\ref{r10-def-ginverse}), then \(G\otimes F\) is a
g-inverse of \(A\otimes B\).
\item If \(A\) and \(B\) are square and both symmetric, both orthogonal, or both idempotent, then \(A\otimes B\) is symmetric,
orthogonal, or idempotent, respectively.
\end{enumerate}
\end{corollary}
\begin{proof}
(a) By Theorem~\ref{r13-thm-kron-rules}(f), \((A\otimes I_p)(I_n\otimes B)=(AI_n)\otimes(I_pB)=A\otimes B\) and
\((I_m\otimes B)(A\otimes I_q)=(I_mA)\otimes(BI_q)=A\otimes B\). The \((i,j)\) block of \(D\otimes B\) is \(d_iB\) for \(i=j\) and \(0\)
for \(i\ne j\); for \(D=I_m\) this gives \(I_m\otimes B\), and for \(B=I_p\) it gives \(I_m\otimes I_p=I_{mp}\).

(b) Let \(A\in\R^{m\times m}\), \(B\in\R^{p\times p}\). By Theorem~\ref{r13-thm-kron-rules}(f) and (a),
\((A\otimes B)(A^{-1}\otimes B^{-1})=(AA^{-1})\otimes(BB^{-1})=I_m\otimes I_p=I_{mp}\); so \(A\otimes B\) is invertible with this
inverse (Theorem~\ref{r2-thm-invertible}).

(c) The diagonal blocks of \(A\otimes B\) are \(a_{11}B,\dots,a_{mm}B\), so
\(\tr(A\otimes B)=\sum_i\tr(a_{ii}B)=\bigl(\sum_ia_{ii}\bigr)\tr B=\tr A\cdot\tr B\) (Theorem~\ref{r1-thm-trace}(a)).

(d) By Theorem~\ref{r13-thm-kron-rules}(f), applied twice,
\((A\otimes B)(G\otimes F)(A\otimes B)=(AGA)\otimes(BFB)=A\otimes B\).

(e) By Theorem~\ref{r13-thm-kron-rules}(e),(f): if \(A'=A\) and \(B'=B\), then \((A\otimes B)'=A'\otimes B'=A\otimes B\); if
\(A'A=I\) and \(B'B=I\), then \((A\otimes B)'(A\otimes B)=(A'A)\otimes(B'B)=I\); if \(A^2=A\) and \(B^2=B\), then
\((A\otimes B)^2=A^2\otimes B^2=A\otimes B\).
\end{proof}

\begin{example}[cf.\ \citeat{harville}{\S16.1}, \citeat{bapat}{\S8.3}]\label{r13-exm-designs}
(a) In the one-way model of Theorem~\ref{r11-thm-oneway} with \(g\) groups of \(m\) observations each (\(n=gm\)), the
indicator of group \(l\) is \(e_l\otimes\one_m\), so \(X_1=[d_1\ \cdots\ d_g]=I_g\otimes\one_m\) and
\(X=[\one_n\ \ X_1]=[\one_g\otimes\one_m\ \ I_g\otimes\one_m]\). By Theorem~\ref{r13-thm-kron-rules}(e),(f) and
Corollary~\ref{r13-cor-kron}(a),(b),
\[
X_1'X_1=(I_g\otimes\one_m')(I_g\otimes\one_m)=I_g\otimes(\one_m'\one_m)=mI_g
\]
and the hat matrix of \(X\) (Theorem~\ref{r10-thm-PX}), the orthogonal projection matrix onto
\(\mathcal C(X)=\mathcal C(X_1)\) (Theorem~\ref{r11-thm-oneway}(c)), is (Theorem~\ref{r6-thm-hat})
\[
H=X_1(X_1'X_1)^{-1}X_1'=(I_g\otimes\one_m)\bigl(\tfrac1mI_g\bigr)(I_g\otimes\one_m')=I_g\otimes\bar J_m
\]
So the hat matrix of the balanced one-way design is block diagonal with \(g\) blocks \(\bar J_m\), and, as
\(I_n=I_g\otimes I_m=I_g\otimes(\bar J_m+C_m)\), the residual maker is \(M=I_n-H=I_g\otimes C_m\)
(Definition~\ref{r1-def-ones}, Theorem~\ref{r13-thm-kron-rules}(b)).

(b) In a \textbf{two-way layout} a response \(y_{ij}\) is observed once for each level \(i\) of a factor with \(a\) levels and
each level \(j\) of a second factor with \(b\) levels; the model without interaction is
\(y_{ij}=\mu+\alpha_i+\beta_j+\varepsilon_{ij}\). List the observations as \(y=(y_{11},\dots,y_{1b},y_{21},\dots,y_{ab})'\). The
observation \((i,j)\) is entry \((i-1)b+j\), so the column of \(\alpha_i\) is \(e_i\otimes\one_b\) and that of \(\beta_j\) is
\(\one_a\otimes e_j\) (Definition~\ref{r13-def-kronecker}), and
\[
X=[\one_a\otimes\one_b\ \ I_a\otimes\one_b\ \ \one_a\otimes I_b]\in\R^{ab\times(1+a+b)},\qquad\beta=(\mu,\alpha_1,\dots,\alpha_a,\beta_1,\dots,\beta_b)'
\]
By Theorem~\ref{r13-thm-kron-rules}(e),(f), \((\one_a\otimes\one_b)'(I_a\otimes\one_b)=\one_a'\otimes(\one_b'\one_b)=b\one_a'\),
\((I_a\otimes\one_b)'(\one_a\otimes I_b)=\one_a\otimes\one_b'=\one_a\one_b'\), \((\one_a\otimes I_b)'(\one_a\otimes I_b)=aI_b\), and
so on; with the block product (Theorem~\ref{r5b-thm-block}(c)),
\[
X'X=\begin{pmatrix}ab&b\one_a'&a\one_b'\\b\one_a&bI_a&\one_a\one_b'\\a\one_b&\one_b\one_a'&aI_b\end{pmatrix}
\]
For \(a=2\) and \(b=3\),
\[
X=\begin{pmatrix}1&1&0&1&0&0\\1&1&0&0&1&0\\1&1&0&0&0&1\\1&0&1&1&0&0\\1&0&1&0&1&0\\1&0&1&0&0&1\end{pmatrix},\qquad
X'X=\begin{pmatrix}6&3&3&2&2&2\\3&3&0&1&1&1\\3&0&3&1&1&1\\2&1&1&2&0&0\\2&1&1&0&2&0\\2&1&1&0&0&2\end{pmatrix}
\]
\end{example}

\subsection*{The vec operator}

\begin{definition}[\citeat{tu-mc}{\S2.4, Exercise~2}, \citeat{harville}{\S16.2}, \citeat{searle-khuri}{\S9.4.1}, \citeat{abadir-magnus}{Exercise~10.17}]\label{r13-def-vec}
For \(A\in\R^{m\times n}\) with columns \(a_1,\dots,a_n\), \(\operatorname{vec}A\in\R^{mn}\) is the vector obtained by stacking the
columns:
\[
\operatorname{vec}A=\begin{pmatrix}a_1\\\vdots\\a_n\end{pmatrix}
\]
Thus \(a_{ij}\) is entry \((j-1)m+i\) of \(\operatorname{vec}A\). For instance, for \(A=\begin{pmatrix}1&2&3\\4&5&6\end{pmatrix}\),
\(\operatorname{vec}A=(1,4,2,5,3,6)'\) and \(\operatorname{vec}A'=(1,2,3,4,5,6)'\).
\end{definition}
Since sums and multiples are taken entrywise, \(\operatorname{vec}(cA+dB)=c\operatorname{vec}A+d\operatorname{vec}B\) for
\(A,B\in\R^{m\times n}\) and \(c,d\in\R\); and for \(A,B\in\R^{m\times n}\), \(\operatorname{vec}A=\operatorname{vec}B\) if and only if
\(A=B\). For a column vector \(a\), \(\operatorname{vec}a=\operatorname{vec}a'=a\). The squared length of
\(\operatorname{vec}A\) is \(\sum_{i,j}a_{ij}^2=\lVert A\rVert_F^2\) (Definition~\ref{r9-def-frobenius}).

\needspace{14\baselineskip}
\begin{theorem}[\citeat{tu-mc}{Propositions~5, 6}, \citeat{harville}{\S16.2}, \citeat{abadir-magnus}{Exercises~10.18--10.20}]\label{r13-thm-vec}
\begin{enumerate}
\item For column vectors \(a\in\R^m\) and \(b\in\R^n\), \(\operatorname{vec}(ab')=b\otimes a\).
\item For \(A\in\R^{m\times n}\), \(B\in\R^{n\times p}\) and \(C\in\R^{p\times q}\),
\[
\operatorname{vec}(ABC)=(C'\otimes A)\operatorname{vec}B
\]
\item In particular, \(\operatorname{vec}(AB)=(I_p\otimes A)\operatorname{vec}B=(B'\otimes I_m)\operatorname{vec}A\) for
\(A\in\R^{m\times n}\), \(B\in\R^{n\times p}\), and \(ABx=(x'\otimes A)\operatorname{vec}B\) for \(x\in\R^p\).
\item For \(A,B\in\R^{m\times n}\), \((\operatorname{vec}A)'\operatorname{vec}B=\tr(A'B)\). More generally, for
\(A\in\R^{m\times n}\), \(B\in\R^{m\times p}\), \(C\in\R^{p\times q}\) and \(D\in\R^{n\times q}\),
\[
\tr(A'BCD')=(\operatorname{vec}A)'(D\otimes B)\operatorname{vec}C
\]
\end{enumerate}
\end{theorem}
\begin{proof}
(a) Column \(j\) of \(ab'\) is \(b_ja\) (Theorem~\ref{r1-thm-views}(a)), so
\(\operatorname{vec}(ab')=(b_1a',\ \dots,\ b_na')'=b\otimes a\) (Definition~\ref{r13-def-kronecker}).

(b)
\begin{itemize}
\item Let \(b_j\) be column \(j\) of \(B\) and \(u_j\) column \(j\) of \(I_p\). By the outer product expansion of \(B=BI_p\)
(Theorem~\ref{r1-thm-views}(d)), \(B=\sum_{j=1}^pb_ju_j'\), so \(ABC=\sum_j(Ab_j)(C'u_j)'\)
(Theorems~\ref{r1-thm-rules} and~\ref{r1-thm-transpose}).
\item By (a) and the mixed product (Theorem~\ref{r13-thm-kron-rules}(f)),
\(\operatorname{vec}\bigl((Ab_j)(C'u_j)'\bigr)=(C'u_j)\otimes(Ab_j)=(C'\otimes A)(u_j\otimes b_j)=(C'\otimes A)\operatorname{vec}(b_ju_j')\).
\item Summing over \(j\) (vec is linear):
\(\operatorname{vec}(ABC)=(C'\otimes A)\operatorname{vec}\bigl(\sum_jb_ju_j'\bigr)=(C'\otimes A)\operatorname{vec}B\).
\end{itemize}

(c) Apply (b) to \(AB\,I_p\), to \(I_m\,AB\), and to \(ABx\) (with \(C=x\in\R^{p\times1}\)), using \(\operatorname{vec}(ABx)=ABx\).

(d)
\begin{itemize}
\item The \(j\)th diagonal entry of \(A'B\) is \(\col_j(A)'\col_j(B)=\sum_ia_{ij}b_{ij}\) (Theorem~\ref{r1-thm-views}(d)), so
\(\tr(A'B)=\sum_{i,j}a_{ij}b_{ij}\); and \((\operatorname{vec}A)'\operatorname{vec}B\) is the same sum, as both vectors list the
entries in the same order.
\item Hence, by (b), \(\tr(A'BCD')=(\operatorname{vec}A)'\operatorname{vec}(BCD')=(\operatorname{vec}A)'(D\otimes B)\operatorname{vec}C\).\qedhere
\end{itemize}
\end{proof}
Applied to \(A'\) and \(B\), (d) also gives \(\tr(AB)=(\operatorname{vec}A')'\operatorname{vec}B\)
for \(A\in\R^{m\times n}\), \(B\in\R^{n\times m}\).

\begin{example}[cf.\ \citeat{bapat}{\S8.3}, \citeat{searle-khuri}{\S12.12.2}]\label{r13-exm-two-way}
In the two-way layout of Example~\ref{r13-exm-designs}(b), arrange the data in \(Y=(y_{ij})\in\R^{a\times b}\); then
\(y=\operatorname{vec}Y'\), since row \(i\) of \(Y\) is column \(i\) of \(Y'\). Let
\[
P=I_a\otimes\bar J_b+\bar J_a\otimes C_b\in\R^{ab\times ab}
\]
\begin{itemize}
\item \(P\) is symmetric (Corollary~\ref{r13-cor-kron}(e)), and \(P^2=P\): by Theorem~\ref{r13-thm-kron-rules}(b),(f) and
\(\bar J^2=\bar J\), \(C^2=C\), \(C\bar J=\bar JC=0\) (Proposition~\ref{r1-prop-centering}; \(\bar J_b\one_b=\one_b\)),
\(P^2=I_a\otimes\bar J_b+\bar J_a\otimes(\bar J_bC_b)+\bar J_a\otimes(C_b\bar J_b)+\bar J_a\otimes C_b=P\).
\item \(PX=X\): \(P(\one_a\otimes\one_b)=\one_a\otimes\one_b+\one_a\otimes C_b\one_b=\one_a\otimes\one_b\),
\(P(I_a\otimes\one_b)=I_a\otimes\one_b+\bar J_a\otimes C_b\one_b=I_a\otimes\one_b\), and
\(P(\one_a\otimes I_b)=\one_a\otimes\bar J_b+\one_a\otimes C_b=\one_a\otimes I_b\). So \(\mathcal C(X)\subseteq\mathcal C(P)\).
\item \(P=XT\) for a matrix \(T\): \(I_a\otimes\bar J_b=(I_a\otimes\one_b)(I_a\otimes\tfrac1b\one_b')\) and
\(\bar J_a\otimes C_b=(\one_a\otimes I_b)(\tfrac1a\one_a'\otimes C_b)\). So \(\mathcal C(P)\subseteq\mathcal C(X)\), and
\(\mathcal C(P)=\mathcal C(X)\).
\item By Theorem~\ref{r6-thm-idempotent}(d), \(Py\) is the orthogonal projection of \(y\) onto \(\mathcal C(X)\), that is, the
vector \(\hat y\) of fitted values (Theorem~\ref{r10-thm-least-squares}(b)).
\item By Theorem~\ref{r13-thm-vec}(b), \((I_a\otimes\bar J_b)\operatorname{vec}Y'=\operatorname{vec}(\bar J_bY')\), whose entry
\((i,j)\) is the row mean \(\bar y_{i\cdot}=\frac1b\sum_ly_{il}\); and
\((\bar J_a\otimes C_b)\operatorname{vec}Y'=\operatorname{vec}(C_bY'\bar J_a)\), whose entry \((i,j)\) is \(\bar y_{\cdot j}-\bar y_{\cdot\cdot}\),
the column mean minus the overall mean (Example~\ref{r1-exm-centering}). Hence
\[
\hat y_{ij}=\bar y_{i\cdot}+\bar y_{\cdot j}-\bar y_{\cdot\cdot}
\]
\end{itemize}
For \(a=2\), \(b=3\) and \(Y=\begin{pmatrix}6&8&4\\2&4&6\end{pmatrix}\): the row means are \(6\), \(4\), the column means
\(4\), \(6\), \(5\), and \(\bar y_{\cdot\cdot}=5\). So \(\hat y=(5,7,6,3,5,4)'\) and \(e=y-\hat y=(1,1,-2,-1,-1,2)'\); indeed
\(X'e=(0,\ 0,\ 0,\ 0,\ 0,\ 0)'\): the entries of \(e\) add up to \(0\) along each row and each column of the layout.
\end{example}

\begin{example}\label{r13-exm-multivariate}
In \textbf{multivariate regression}, \(p\) responses are measured on each of \(n\) units, \(Y\in\R^{n\times p}\) holds them
(column \(j\) is response \(j\)), and \(Y=XB+E\) with a design \(X\in\R^{n\times(k+1)}\) of full column rank, an unknown
\(B\in\R^{(k+1)\times p}\) and errors \(E\). The least-squares estimate minimizes
\(\lVert Y-XB\rVert_F^2=\lVert\operatorname{vec}Y-\operatorname{vec}(XB)\rVert^2\). By Theorem~\ref{r13-thm-vec}(c),
\(\operatorname{vec}(XB)=(I_p\otimes X)\operatorname{vec}B\): this is the least-squares problem with response \(\operatorname{vec}Y\)
and design \(I_p\otimes X\). As \((I_p\otimes X)'(I_p\otimes X)=I_p\otimes X'X\) is invertible with inverse
\(I_p\otimes(X'X)^{-1}\) (Corollary~\ref{r13-cor-kron}(b)), \(I_p\otimes X\) has full column rank
(Theorems~\ref{thm:r4-rank-ata}(b) and~\ref{thm:r4-rank-rules}(d)), and by Theorem~\ref{r10-thm-least-squares}(c) the unique
estimate is
\[
\operatorname{vec}\hat B=\bigl(I_p\otimes(X'X)^{-1}\bigr)(I_p\otimes X')\operatorname{vec}Y=\bigl(I_p\otimes(X'X)^{-1}X'\bigr)\operatorname{vec}Y
=\operatorname{vec}\bigl((X'X)^{-1}X'Y\bigr)
\]
that is, \(\hat B=(X'X)^{-1}X'Y\): column \(j\) of \(\hat B\) is the least-squares estimate for response \(j\) alone.
If \(X\) has any rank, let \(G\) be a g-inverse of \(X'X\); then \(I_p\otimes G\) is a g-inverse of
\(I_p\otimes X'X\) (Corollary~\ref{r13-cor-kron}(d)), and by Theorems~\ref{r10-thm-normal-equations}(c)
and~\ref{r10-thm-least-squares}(a) the same computation gives the least-squares estimate
\(\operatorname{vec}B^0=\bigl(I_p\otimes GX'\bigr)\operatorname{vec}Y=\operatorname{vec}(GX'Y)\), that is, \(B^0=GX'Y\): column \(j\) of
\(B^0\) is a least-squares estimate for response \(j\) alone.
\end{example}

\begin{theorem}[\citeat{abadir-magnus}{Exercise~13.65}, \citeat{tu-mc}{\S13}]\label{r13-thm-data-matrix}
Let \(Y\in\R^{n\times p}\) be a random matrix with rows \(y_1',\dots,y_n'\) such that \(\expc(y_i)=\mu\) and
\(\var(y_i)=\Sigma=(\sigma_{jl})\) for every \(i\), and \(\cov(y_{ij},y_{kl})=0\) for \(i\ne k\): the
observations on different units are uncorrelated. Then \(\expc(\operatorname{vec}Y)=\mu\otimes\one_n\) and
\(\var(\operatorname{vec}Y)=\Sigma\otimes I_n\); and \(\operatorname{vec}Y'=(y_1',\dots,y_n')'\) has
\(\expc(\operatorname{vec}Y')=\one_n\otimes\mu\) and \(\var(\operatorname{vec}Y')=I_n\otimes\Sigma\).
\end{theorem}
\begin{proof}
\begin{itemize}
\item Block \(j\) of \(\operatorname{vec}Y\) is column \(j\) of \(Y\), \((y_{1j},\dots,y_{nj})'\), with mean \(\mu_j\one_n\). So
\(\expc(\operatorname{vec}Y)\) has the blocks \(\mu_1\one_n,\dots,\mu_p\one_n\): it is \(\mu\otimes\one_n\)
(Definitions~\ref{r11-def-random} and~\ref{r13-def-kronecker}).
\item The entry of \(\var(\operatorname{vec}Y)\) in block \((j,l)\), position \((i,k)\), is
\(\cov(y_{ij},y_{kl})\) (Definition~\ref{r11-def-random}). It is \(\sigma_{jl}\) for \(i=k\) (entry \((j,l)\) of
\(\var(y_i)\)) and \(0\) for \(i\ne k\). So block \((j,l)\) is \(\sigma_{jl}I_n\), the \((j,l)\) block of
\(\Sigma\otimes I_n\).
\item Likewise block \(i\) of \(\operatorname{vec}Y'\) is \(y_i\), with mean \(\mu\), and block \((i,k)\) of its covariance
matrix, with entries \(\cov(y_{ij},y_{kl})\), is \(\Sigma\) for \(i=k\) and \(0\) for \(i\ne k\): the \((i,k)\) block
of \(I_n\otimes\Sigma\).\qedhere
\end{itemize}
\end{proof}
For \(n\geqslant2\) the rows of \(Y\) satisfy the hypotheses of Theorem~\ref{r14-thm-sample}, which gives \(\expc(\bar y)=\mu\) and
\(\var(\bar y)=\frac1n\Sigma\) for the mean vector \(\bar y=\frac1nY'\one_n\), and \(\expc(S)=\Sigma\) for the
sample covariance matrix \(S=\frac1{n-1}Y'C_nY\) (Definition~\ref{r12-def-data}). The covariance matrix of \(\bar y\) also follows
from the theorem with the mixed product: by Theorem~\ref{r13-thm-vec}(b),
\(\bar y=I_pY'(\frac1n\one_n)=(\frac1n\one_n'\otimes I_p)\operatorname{vec}Y'\), and by Theorem~\ref{r11-thm-linear}(b) and
Theorem~\ref{r13-thm-kron-rules}(e),(f),
\[
\var(\bar y)=\bigl(\tfrac1n\one_n'\otimes I_p\bigr)(I_n\otimes\Sigma)\bigl(\tfrac1n\one_n\otimes I_p\bigr)
=\tfrac1{n^2}(\one_n'\one_n)\otimes\Sigma=\tfrac1n\Sigma
\]

\subsection*{The commutation matrix}

\begin{definition}[\citeat{tu-mc}{\S2.5}, \citeat{harville}{\S16.3a}, \citeat{abadir-magnus}{Exercise~11.9(a)}]\label{r13-def-commutation}
Let \(m,n\geqslant1\). The \textbf{commutation matrix} \(K_{mn}\in\R^{mn\times mn}\) is the matrix with
\[
K_{mn}\operatorname{vec}A=\operatorname{vec}A'\qquad\text{for every }A\in\R^{m\times n}
\]
(the next theorem shows that there is exactly one such matrix). We write \(K_n=K_{nn}\) (not the Helmert matrix \(K_n\) of Example~\ref{r6-exm-helmert}). For example,
for \(A=(a_{ij})\in\R^{2\times3}\), \(\operatorname{vec}A=(a_{11},a_{21},a_{12},a_{22},a_{13},a_{23})'\) and
\(\operatorname{vec}A'=(a_{11},a_{12},a_{13},a_{21},a_{22},a_{23})'\), so
\[
K_{23}=\begin{pmatrix}1&0&0&0&0&0\\0&0&1&0&0&0\\0&0&0&0&1&0\\0&1&0&0&0&0\\0&0&0&1&0&0\\0&0&0&0&0&1\end{pmatrix}
\]
\end{definition}

\begin{theorem}[\citeat{harville}{\S16.3, Theorem~16.3.2}, \citeat{tu-mc}{Propositions~7, 8, Exercise~4}, \citeat{abadir-magnus}{Exercises~11.1--11.5}]\label{r13-thm-commutation}
Let \(m,n,p,q\geqslant1\).
\begin{enumerate}
\item There is exactly one matrix \(K_{mn}\) as in Definition~\ref{r13-def-commutation}: row \((i-1)n+j\) of \(K_{mn}\) is row
\((j-1)m+i\) of \(I_{mn}\) (\(1\leqslant i\leqslant m\), \(1\leqslant j\leqslant n\)). It is a permutation matrix.
\item \(K_{mn}'=K_{mn}^{-1}=K_{nm}\); \(K_{m1}=K_{1m}=I_m\); and \(K_n\) is symmetric.
\item \(\tr K_n=n\).
\item For \(A\in\R^{m\times n}\) and \(B\in\R^{p\times q}\),
\[
K_{pm}(A\otimes B)=(B\otimes A)K_{qn},\qquad B\otimes A=K_{pm}(A\otimes B)K_{nq}
\]
and for column vectors \(a\in\R^m\), \(b\in\R^p\), \(K_{pm}(a\otimes b)=b\otimes a\).
\end{enumerate}
\end{theorem}
\begin{proof}
(a)
\begin{itemize}
\item For \(A\in\R^{m\times n}\), \(a_{ij}\) is entry \((j-1)m+i\) of \(\operatorname{vec}A\); in \(A'\in\R^{n\times m}\) it is the \((j,i)\)
entry, so it is entry \((i-1)n+j\) of \(\operatorname{vec}A'\) (Definition~\ref{r13-def-vec}).
\item As \((i,j)\) runs over the \(mn\) pairs, both \((i-1)n+j\) and \((j-1)m+i\) run over \(1,\dots,mn\), each value once. So
\(\varphi\bigl((i-1)n+j\bigr)=(j-1)m+i\) defines a permutation \(\varphi\) of degree \(mn\) (Definition~\ref{r5d-def-permutation}),
and its permutation matrix \(P_\varphi\) has row \((i-1)n+j\) equal to row \((j-1)m+i\) of \(I_{mn}\)
(Definition~\ref{r5d-def-permutation-matrix}).
\item By Proposition~\ref{r5d-prop-permutation-matrix}(a), entry \((i-1)n+j\) of \(P_\varphi\operatorname{vec}A\) is entry
\((j-1)m+i\) of \(\operatorname{vec}A\), that is, \(a_{ij}\), which is entry \((i-1)n+j\) of \(\operatorname{vec}A'\). So
\(P_\varphi\operatorname{vec}A=\operatorname{vec}A'\) for every \(A\).
\item If also \(K\operatorname{vec}A=\operatorname{vec}A'\) for every \(A\), then \(Kv=P_\varphi v\) for every \(v\in\R^{mn}\), since
every \(v\) is \(\operatorname{vec}A\) for some \(A\); taking \(v=e_r\) gives equal columns (Theorem~\ref{r1-thm-views}(b)), so
\(K=P_\varphi\).
\end{itemize}
(b)
\begin{itemize}
\item By (a) and Proposition~\ref{r5d-prop-permutation-matrix}(c), \(K_{mn}\) is invertible with \(K_{mn}^{-1}=K_{mn}'\).
\item For every \(A\in\R^{m\times n}\), \(K_{nm}K_{mn}\operatorname{vec}A=K_{nm}\operatorname{vec}A'=\operatorname{vec}(A')'=\operatorname{vec}A\);
as in (a), \(K_{nm}K_{mn}=I_{mn}\), so \(K_{nm}=K_{mn}^{-1}\) (Theorem~\ref{r2-thm-invertible}).
\item For \(a\in\R^m\), \(\operatorname{vec}a=\operatorname{vec}a'=a\), so \(K_{m1}a=a\) and \(K_{1m}a=a\) for every \(a\): \(K_{m1}=K_{1m}=I_m\).
Finally \(K_n'=K_n^{-1}=K_n\).
\end{itemize}
(c) For \(n=1\), \(K_1=I_1\). Let \(n\geqslant2\). By (a), the diagonal entry of \(K_n\) in position \((i-1)n+j\) is \(1\) if
\((j-1)n+i=(i-1)n+j\), that is, \((j-i)(n-1)=0\), that is, \(i=j\), and \(0\) otherwise. So exactly \(n\) diagonal entries are
\(1\), and \(\tr K_n=n\).

(d)
\begin{itemize}
\item Let \(X\in\R^{q\times n}\). By Definition~\ref{r13-def-commutation} and Theorem~\ref{r13-thm-vec}(b), applied three times,
\[
K_{pm}(A\otimes B)\operatorname{vec}X=K_{pm}\operatorname{vec}(BXA')=\operatorname{vec}(AX'B')=(B\otimes A)\operatorname{vec}X'
=(B\otimes A)K_{qn}\operatorname{vec}X
\]
As \(\operatorname{vec}X\) runs over \(\R^{nq}\), this gives \(K_{pm}(A\otimes B)=(B\otimes A)K_{qn}\), as in (a).
\item Multiplying on the right by \(K_{nq}=K_{qn}^{-1}\) (by (b)) gives \(K_{pm}(A\otimes B)K_{nq}=B\otimes A\).
\item By Theorem~\ref{r13-thm-vec}(a), \(K_{pm}(a\otimes b)=K_{pm}\operatorname{vec}(ba')=\operatorname{vec}(ab')=b\otimes a\).\qedhere
\end{itemize}
\end{proof}
For the data matrix of Theorem~\ref{r13-thm-data-matrix}, \(\operatorname{vec}Y'=K_{np}\operatorname{vec}Y\), and (d) with (b) gives
again \(\var(\operatorname{vec}Y')=K_{np}(\Sigma\otimes I_n)K_{np}'=(I_n\otimes\Sigma)K_{np}K_{pn}=I_n\otimes\Sigma\)
(Theorem~\ref{r11-thm-linear}(b)).

\begin{theorem}[\citeat{harville}{(16.3.18), (16.1.26)}, \citeat{abadir-magnus}{Exercise~10.12(b)}]\label{r13-thm-kron-det}
\begin{enumerate}
\item For \(A\in\R^{m\times m}\) and \(B\in\R^{n\times n}\), \(\det(A\otimes B)=(\det A)^n(\det B)^m\).
\item For \(A\in\R^{m\times n}\) and \(B\in\R^{p\times q}\), \(\rank(A\otimes B)=\rank A\cdot\rank B\).
\item For square \(A\) and \(B\), \(A\otimes B\) is invertible if and only if \(A\) and \(B\) are invertible.
\end{enumerate}
\end{theorem}
\begin{proof}
(a)
\begin{itemize}
\item By Corollary~\ref{r13-cor-kron}(a), \(A\otimes B=(A\otimes I_n)(I_m\otimes B)\), and \(I_m\otimes B\) is block diagonal with \(m\)
blocks \(B\). Applying Theorem~\ref{r5d-thm-block-tri} (\(\det\diag(A_1,D)=\det A_1\cdot\det D\)) \(m-1\) times,
\(\det(I_m\otimes B)=(\det B)^m\); likewise \(\det(I_n\otimes A)=(\det A)^n\).
\item By Theorem~\ref{r13-thm-commutation}(d) with \(B=I_n\) (\(p=q=n\)), \(K_{nm}(A\otimes I_n)=(I_n\otimes A)K_{nm}\), so
\(A\otimes I_n=K_{mn}(I_n\otimes A)K_{nm}\) (Theorem~\ref{r13-thm-commutation}(b)).
\item By Theorem~\ref{r5d-thm-product}(b), \(\det(A\otimes I_n)=\det K_{mn}\cdot(\det A)^n\cdot\det K_{nm}\), and
\(\det K_{mn}\cdot\det K_{nm}=\det(K_{mn}K_{nm})=\det I_{mn}=1\). Hence \(\det(A\otimes B)=(\det A)^n(\det B)^m\).
\end{itemize}
(b)
\begin{itemize}
\item For any matrix \(M\) with a g-inverse \(N\), \(\rank M=\tr(MN)\) (Theorem~\ref{r10-thm-rank-idempotent}(c)).
\item Let \(G\) and \(F\) be g-inverses of \(A\) and \(B\) (Theorem~\ref{r10-thm-ginv-exists}). By
Corollary~\ref{r13-cor-kron}(d), \(G\otimes F\) is a g-inverse of \(A\otimes B\). By the first step (for \(A\otimes B\), \(A\) and
\(B\)), the mixed product and Corollary~\ref{r13-cor-kron}(c),
\[
\rank(A\otimes B)=\tr\bigl((A\otimes B)(G\otimes F)\bigr)=\tr(AG\otimes BF)=\tr(AG)\tr(BF)=\rank A\cdot\rank B
\]
\end{itemize}
(c) By (a), \(\det(A\otimes B)\ne0\) if and only if \(\det A\ne0\) and \(\det B\ne0\); apply Theorem~\ref{r5d-thm-product}(a).
\end{proof}

\needspace{12\baselineskip}
\begin{theorem}[\citeat{searle-khuri}{Theorem~6.10}, \citeat{abadir-magnus}{Exercises~10.11(a), 10.12(a)}]\label{r13-thm-kron-eigen}
Let \(A\in\R^{m\times m}\) and \(B\in\R^{n\times n}\).
\begin{enumerate}
\item If \(Ax=\lambda x\) and \(By=\mu y\) with \(x\ne0\), \(y\ne0\), then \((A\otimes B)(x\otimes y)=\lambda\mu(x\otimes y)\) and
\(x\otimes y\ne0\): \(\lambda\mu\) is an eigenvalue of \(A\otimes B\) with eigenvector \(x\otimes y\).
\item If \(P^{-1}AP=\Lambda=\diag(\lambda_1,\dots,\lambda_m)\) and \(Q^{-1}BQ=M=\diag(\mu_1,\dots,\mu_n)\) (Definition~\ref{r7-def-diagonalizable}),
then \((P\otimes Q)^{-1}(A\otimes B)(P\otimes Q)=\Lambda\otimes M\). So \(A\otimes B\) is diagonalizable, its eigenvalues, listed with
multiplicity, are the \(mn\) products \(\lambda_i\mu_j\), and \(p_i\otimes q_j\) is an eigenvector for \(\lambda_i\mu_j\), where
\(p_i=\col_i(P)\), \(q_j=\col_j(Q)\).
\item If \(A\) and \(B\) are symmetric, with \(A=P\Lambda P'\), \(B=QMQ'\) as in Theorem~\ref{r8-thm-spectral}(b), then \(P\otimes Q\) is
orthogonal and \(A\otimes B=(P\otimes Q)(\Lambda\otimes M)(P\otimes Q)'\). If \(A\) and \(B\) are PSD (PD), then \(A\otimes B\) is PSD (PD).
\end{enumerate}
\end{theorem}
\begin{proof}
(a) By the mixed product, \((A\otimes B)(x\otimes y)=(Ax)\otimes(By)=(\lambda x)\otimes(\mu y)=\lambda\mu(x\otimes y)\)
(Theorem~\ref{r13-thm-kron-rules}(a),(f)). If \(x_i\ne0\) and \(y_j\ne0\), the entry \(x_iy_j\) of \(x\otimes y\) is not \(0\).

(b)
\begin{itemize}
\item \(P\otimes Q\) is invertible with inverse \(P^{-1}\otimes Q^{-1}\) (Corollary~\ref{r13-cor-kron}(b)), and by the mixed product
\((P^{-1}\otimes Q^{-1})(A\otimes B)(P\otimes Q)=(P^{-1}AP)\otimes(Q^{-1}BQ)=\Lambda\otimes M\).
\item \(\Lambda\otimes M=\diag(\lambda_1M,\dots,\lambda_mM)\) (Corollary~\ref{r13-cor-kron}(a)) is diagonal, with \(\lambda_i\mu_j\) in
position \((i-1)n+j\); and column \((i-1)n+j\) of \(P\otimes Q\) is \(p_i\otimes q_j\) (Definition~\ref{r13-def-kronecker}). By
Definition~\ref{r7-def-diagonalizable}, these \(mn\) numbers are the eigenvalues of \(A\otimes B\), listed with multiplicity, and
\(p_i\otimes q_j\) belongs to \(\lambda_i\mu_j\).
\end{itemize}
(c)
\begin{itemize}
\item \(P\otimes Q\) is orthogonal and \(A\otimes B\) symmetric (Corollary~\ref{r13-cor-kron}(e)); as \((P\otimes Q)^{-1}=(P\otimes Q)'\),
(b) gives \(A\otimes B=(P\otimes Q)(\Lambda\otimes M)(P\otimes Q)'\).
\item If \(A\), \(B\) are PSD (PD), all \(\lambda_i,\mu_j\geqslant0\) (\(>0\)) (Theorem~\ref{r8-thm-eigen-criterion}(a)), so all eigenvalues
\(\lambda_i\mu_j\) of \(A\otimes B\) are \(\geqslant0\) (\(>0\)), and \(A\otimes B\) is PSD (PD) by the same theorem.\qedhere
\end{itemize}
\end{proof}
In (b), \(\det(A\otimes B)=\prod_{i,j}\lambda_i\mu_j=\prod_i\bigl(\lambda_i^n\prod_j\mu_j\bigr)=(\det A)^n(\det B)^m\)
(Theorem~\ref{r7-thm-trace-det}), in agreement with Theorem~\ref{r13-thm-kron-det}(a); and \(\tr(A\otimes B)=\sum_{i,j}\lambda_i\mu_j=\tr A\cdot\tr B\).

\subsection*{The duplication matrix}

\begin{definition}[\citeat{tu-mc}{\S2.6}, \citeat{harville}{\S16.4}, \citeat{searle-khuri}{\S9.4.1}]\label{r13-def-duplication}
For \(A\in\R^{n\times n}\), \(\operatorname{vech}A\in\R^{n(n+1)/2}\) is the vector obtained from \(\operatorname{vec}A\) by deleting the entries
above the diagonal:
\[
\operatorname{vech}A=(a_{11},\dots,a_{n1},a_{22},\dots,a_{n2},\dots,a_{nn})'
\]
For symmetric \(A\) it lists each distinct entry once. The \textbf{duplication matrix} \(D_n\in\R^{n^2\times n(n+1)/2}\) is the matrix with
\(D_n\operatorname{vech}A=\operatorname{vec}A\) for every symmetric \(A\in\R^{n\times n}\). For example,
\[
\operatorname{vech}\begin{pmatrix}1&7&6\\7&3&8\\6&8&2\end{pmatrix}=(1,7,6,3,8,2)',\qquad
D_2=\begin{pmatrix}1&0&0\\0&1&0\\0&1&0\\0&0&1\end{pmatrix}
\]
\end{definition}

\needspace{14\baselineskip}
\begin{proposition}[\citeat{harville}{\S16.4b, (16.4.6), (16.4.14)}, \citeat{tu-mc}{Proposition~9, Exercise~5}]\label{r13-prop-duplication}
Let \(n\geqslant1\).
\begin{enumerate}
\item There is exactly one duplication matrix \(D_n\). Its column for the position of \(a_{ij}\) (\(i\geqslant j\)) in
\(\operatorname{vech}A\) has the entry \(1\) in rows \((j-1)n+i\) and \((i-1)n+j\) (one row if \(i=j\)) and \(0\) elsewhere. The
columns of \(D_n\) are nonzero and pairwise orthogonal, \(D_n'D_n\) is diagonal with the entries \(1\) (\(i=j\)) and \(2\)
(\(i>j\)), and \(D_n\) has full column rank \(n(n+1)/2\).
\item For symmetric \(A\), \(\operatorname{vech}A=(D_n'D_n)^{-1}D_n'\operatorname{vec}A\).
\item \(K_nD_n=D_n\), and \(D_n(D_n'D_n)^{-1}D_n'=\frac12(I_{n^2}+K_n)\).
\end{enumerate}
\end{proposition}
\begin{proof}
(a)
\begin{itemize}
\item For symmetric \(A\) and \(i\geqslant j\), the entries \((j-1)n+i\) and \((i-1)n+j\) of \(\operatorname{vec}A\) are \(a_{ij}\) and
\(a_{ji}=a_{ij}\), the entry of \(\operatorname{vech}A\) for \((i,j)\). So the matrix described has \(D_n\operatorname{vech}A=\operatorname{vec}A\).
As \(A\) runs over the symmetric matrices, \(\operatorname{vech}A\) runs over all of \(\R^{n(n+1)/2}\); as in
Theorem~\ref{r13-thm-commutation}(a), \(D_n\) is unique.
\item Every row of \(D_n\) has exactly one entry \(1\), so two different columns have no nonzero entry in a common row: they are
orthogonal, and \(D_n'D_n\) is diagonal with the squared lengths \(1\) or \(2\) of the columns (Theorem~\ref{r1-thm-views}(d)).
\item \(D_n'D_n\) is invertible, so \(\rank D_n=\rank(D_n'D_n)=n(n+1)/2\) (Theorems~\ref{thm:r4-rank-ata}(b)
and~\ref{thm:r4-rank-rules}(d)).
\end{itemize}
(b) Multiply \(D_n\operatorname{vech}A=\operatorname{vec}A\) on the left by \((D_n'D_n)^{-1}D_n'\).

(c)
\begin{itemize}
\item For symmetric \(A\), \(K_nD_n\operatorname{vech}A=K_n\operatorname{vec}A=\operatorname{vec}A'=\operatorname{vec}A=D_n\operatorname{vech}A\);
as \(\operatorname{vech}A\) runs over \(\R^{n(n+1)/2}\), \(K_nD_n=D_n\).
\item Let \(M=D_n(D_n'D_n)^{-1}D_n'\) and \(N=\frac12(I_{n^2}+K_n)\). \(M\) is symmetric and idempotent. \(K_n\) is symmetric with
\(K_n^2=K_n'K_n=I\) (Theorem~\ref{r13-thm-commutation}(b)), so \(N'=N\) and \(N^2=\frac14(I+2K_n+K_n^2)=N\).
\item \(ND_n=\frac12(D_n+K_nD_n)=D_n\), so \(NM=M\) and \(MN=(NM)'=M\). Hence \(\Delta=N-M\) is symmetric with
\(\Delta^2=N-NM-MN+M=N-M=\Delta\).
\item \(\tr M=\tr\bigl((D_n'D_n)^{-1}D_n'D_n\bigr)=\frac{n(n+1)}2\) and \(\tr N=\frac12(n^2+n)\) by Theorem~\ref{r13-thm-commutation}(c)
(Theorem~\ref{r1-thm-trace}). So \(\rank\Delta=\tr\Delta=0\) (Theorem~\ref{r6-thm-idempotent}(c)), and \(\Delta=0\).\qedhere
\end{itemize}
\end{proof}

\subsection*{Derivatives and differentials}

Let \(S\subseteq\R^n\) and \(c\in S\). The \textbf{ball} \(B(c,\delta)\) of radius \(\delta>0\) is the set of \(x\in\R^n\) with
\(\lVert x-c\rVert<\delta\), and \(S\) is \textbf{open} if every \(c\in S\) has a ball \(B(c,\delta)\subseteq S\). For matrix variables
\(X\in\R^{n\times q}\) the distance is \(\lVert X-C\rVert_F=\lVert\operatorname{vec}X-\operatorname{vec}C\rVert\) (Definition~\ref{r9-def-frobenius}),
so a set of matrices is open if the corresponding set of vectors \(\operatorname{vec}X\) is. For a function \(f\) on an open set
\(S\) and \(x\in S\), the \textbf{partial derivative} \(\partial f(x)/\partial x_j\) is the derivative at \(t=0\) of the function
\(t\mapsto f(x+te_j)\) of one variable, when it exists~\citeat{harville}{\S15.1}.

\begin{definition}[\citeat{tu-mc}{\S3}, \citeat{harville}{\S15.1}, \citeat{searle-khuri}{(9.11)}]\label{r13-def-derivative}
Let \(S\subseteq\R^n\) be open and \(f=(f_1,\dots,f_m)':S\to\R^m\). If all partial derivatives of \(f_1,\dots,f_m\) exist at \(x\in S\),
the \textbf{derivative} (or \textbf{Jacobian matrix}) of \(f\) at \(x\) is the \(m\times n\) matrix
\[
\mathrm Df(x)=\frac{\partial f(x)}{\partial x'}\quad\text{with }(i,j)\text{ entry }\frac{\partial f_i(x)}{\partial x_j}
\]
Row \(i\) holds the partial derivatives of \(f_i\), column \(j\) the partial derivatives with respect to \(x_j\). For a scalar
function (\(m=1\)), \(\mathrm Df(x)\) is a row vector, and the column \(\partial f(x)/\partial x=(\mathrm Df(x))'\) is the
\textbf{gradient}. A matrix function \(F\), with values \(F(X)\in\R^{m\times p}\), of a matrix \(X\) in an open set of
\(\R^{n\times q}\) is regarded as the vector function \(\operatorname{vec}X\mapsto\operatorname{vec}F(X)\); its derivative is the
\(mp\times nq\) matrix
\[
\mathrm DF(X)=\frac{\partial\operatorname{vec}F(X)}{\partial(\operatorname{vec}X)'}
\]
For a scalar function \(\varphi(X)\) we also write \(\partial\varphi/\partial X\) for the \(n\times q\) matrix with \((i,j)\) entry
\(\partial\varphi/\partial x_{ij}\); thus \(\mathrm D\varphi(X)=(\operatorname{vec}(\partial\varphi/\partial X))'\).
\end{definition}
Some books write \(\partial y/\partial x\) for the \emph{transposed} derivative of a vector function: there
\(\partial(Ax)/\partial x=A'\). In these notes \(\mathrm D\) always has one row for each entry of the function and one column for
each variable.

We use the following facts from calculus without proof~\citeat{harville}{Lemmas~15.1.1, 15.2.2, Theorem~15.7.1}. A function on an open set is \textbf{continuously differentiable}
(\(C^1\)) if its partial derivatives exist and are continuous there, and \textbf{twice continuously differentiable} (\(C^2\)) if
moreover the partial derivatives of the partial derivatives exist and are continuous; vector and matrix functions are \(C^1\)
(\(C^2\)) if all their entries are.
\begin{itemize}
\item (\textbf{Rules}) Sums, products and quotients (with nonzero denominator) of \(C^1\) (\(C^2\)) functions are \(C^1\) (\(C^2\)), and
\(\partial(f+g)/\partial x_j=\partial f/\partial x_j+\partial g/\partial x_j\),
\(\partial(fg)/\partial x_j=(\partial f/\partial x_j)g+f\,\partial g/\partial x_j\), \(\partial(cf)/\partial x_j=c\,\partial f/\partial x_j\)
(\(c\) constant); constants and the variables are \(C^2\), so polynomials in the variables are \(C^2\). A \(C^1\) function is
continuous, and a continuous function that is not \(0\) (or is positive) at \(c\) is not \(0\) (positive) on some ball around \(c\).
\item (\textbf{Chain rule}) If \(g\) is \(C^1\) on an open \(S\subseteq\R^n\) with values in an open \(T\subseteq\R^m\), and \(f\) is \(C^1\)
(\(C^2\)) on \(T\) with values in \(\R^p\), then \(f\circ g\) is \(C^1\) on \(S\) (\(C^2\) if \(g\) is too), and
\(\mathrm D(f\circ g)(x)=\mathrm Df(g(x))\,\mathrm Dg(x)\). The functions \(\log t\) and \(t^{-1/2}\) of \(t>0\) are \(C^2\), with
derivatives \(1/t\) and \(-\frac12t^{-3/2}\).
\end{itemize}

\begin{example}[\citeat{searle-khuri}{Example~9.1, \S9.3.1, \S9.5.2}]\label{r13-exm-derivative}
(a) For \(f(x)=Ax\) with \(A\in\R^{m\times n}\) constant, \(f_i(x)=\sum_ja_{ij}x_j\), so \(\partial f_i/\partial x_j=a_{ij}\) and
\(\mathrm Df(x)=A\). For
\[
A=\begin{pmatrix}2&6&-1\\3&-2&4\\3&4&7\end{pmatrix},\qquad Ax=\begin{pmatrix}2x_1+6x_2-x_3\\3x_1-2x_2+4x_3\\3x_1+4x_2+7x_3\end{pmatrix}
\]
the partial derivatives of the first entry are \(2,6,-1\), the first row of \(\mathrm D(Ax)=A\); in the transposed notation,
\(\partial(Ax)/\partial x=A'\) has these numbers in its first column.

(b) For \(\lambda=3x_1+4x_2+9x_3=a'x\) with \(a=(3,4,9)'\), \(\mathrm D\lambda=(3,\ 4,\ 9)=a'\), and the gradient is \(a\).

(c) For \(X\in\R^{n\times q}\) and constant \(A\in\R^{q\times n}\), \(\tr(XA)=\sum_{r,s}x_{rs}a_{sr}\), so
\(\partial\tr(XA)/\partial x_{ij}=a_{ji}\): \(\partial\tr(XA)/\partial X=A'\), and \(\mathrm D\tr(XA)=(\operatorname{vec}A')'\).
\end{example}

\begin{lemma}[\citeat{tu-mc}{Proposition~2}]\label{r13-lem-forms}
\begin{enumerate}
\item If \(A,B\in\R^{m\times n}\) and \(Ax=Bx\) for every \(x\in\R^n\), then \(A=B\).
\item For \(A\in\R^{n\times n}\), \(x'Ax=0\) for every \(x\in\R^n\) if and only if \(A'=-A\) (\(A\) is skew-symmetric).
\item For \(A,B\in\R^{n\times n}\), \(x'Ax=x'Bx\) for every \(x\) if and only if \(A+A'=B+B'\). In particular, for symmetric \(A\)
and \(B\), \(x'Ax=x'Bx\) for every \(x\) only if \(A=B\).
\end{enumerate}
\end{lemma}
\begin{proof}
(a) With \(x=e_j\), \(\col_j(A)=Ae_j=Be_j=\col_j(B)\) for every \(j\) (Theorem~\ref{r1-thm-views}(b)).

(b)
\begin{itemize}
\item If \(A'=-A\), then \(x'Ax=(x'Ax)'=x'A'x=-x'Ax\), as a number is its own transpose; so \(x'Ax=0\).
\item Conversely, let \(x'Ax=0\) for every \(x\). With \(x=e_i\), \(a_{ii}=0\). With \(x=e_i+e_j\) (\(i\ne j\)),
\(0=a_{ii}+a_{ij}+a_{ji}+a_{jj}=a_{ij}+a_{ji}\). So \(A'=-A\).
\end{itemize}
(c) \(x'Ax=x'Bx\) for every \(x\) means \(x'(A-B)x=0\) for every \(x\), that is, by (b), \((A-B)'=-(A-B)\), or
\(A+A'=B+B'\). For symmetric \(A\), \(B\) this reads \(2A=2B\).
\end{proof}

\begin{definition}[\citeat{tu-mc}{\S4}, \citeat{abadir-magnus}{Ch.~13}]\label{r13-def-differential}
Let \(f\) be \(C^1\) on an open \(S\subseteq\R^n\) with values in \(\R^m\), and \(x\in S\). The \textbf{differential} of \(f\) at \(x\) is the
linear function of \(\mathrm{d}x\in\R^n\)
\[
\mathrm{d}f(x)=\mathrm Df(x)\,\mathrm{d}x
\]
where \(\mathrm{d}x\) is an arbitrary vector (an increment of \(x\)). For a \(C^1\) matrix function \(F\) of \(X\in\R^{n\times q}\), the
differential \(\mathrm{d}F(X)\) is the matrix of the size of \(F(X)\) with \(\operatorname{vec}\mathrm{d}F(X)=\mathrm DF(X)\operatorname{vec}\mathrm{d}X\),
\(\mathrm{d}X\in\R^{n\times q}\) arbitrary; its \((k,l)\) entry is \(\sum_{i,j}(\partial f_{kl}/\partial x_{ij})\,\mathrm{d}x_{ij}\).
\end{definition}
So \(\mathrm{d}f(x)\) has the size of \(f\), whatever the number of variables, while \(\mathrm Df(x)\) has one column for each variable;
this makes differentials easier to compute with. For the identity function \(f(x)=x\), \(\mathrm Df(x)=I_n\) and the differential is
\(\mathrm{d}x\) itself.

\needspace{12\baselineskip}
\begin{theorem}[First identification theorem; \citeat{tu-mc}{Proposition~14}]\label{r13-thm-identification}
\begin{enumerate}
\item Let \(f\) be \(C^1\) on an open \(S\subseteq\R^n\) with values in \(\R^m\), \(x\in S\) and \(A\in\R^{m\times n}\). Then
\(\mathrm{d}f(x)=A\,\mathrm{d}x\) for every \(\mathrm{d}x\) if and only if \(\mathrm Df(x)=A\).
\item For a \(C^1\) matrix function \(F\) of \(X\in\R^{n\times q}\) with values in \(\R^{m\times p}\) and \(A\in\R^{mp\times nq}\):
\(\operatorname{vec}\mathrm{d}F(X)=A\operatorname{vec}\mathrm{d}X\) for every \(\mathrm{d}X\) if and only if \(\mathrm DF(X)=A\).
\item For a \(C^1\) scalar function \(\varphi\) of \(X\in\R^{n\times q}\) and \(C\in\R^{n\times q}\): \(\mathrm{d}\varphi(X)=\tr(C'\mathrm{d}X)\)
for every \(\mathrm{d}X\) if and only if \(\mathrm D\varphi(X)=(\operatorname{vec}C)'\), that is, \(\partial\varphi/\partial X=C\).
\end{enumerate}
\end{theorem}
\begin{proof}
(a) If \(\mathrm Df(x)=A\), then \(\mathrm{d}f(x)=A\,\mathrm{d}x\) by Definition~\ref{r13-def-differential}. Conversely,
\(A\,\mathrm{d}x=\mathrm Df(x)\,\mathrm{d}x\) for every \(\mathrm{d}x\) gives \(\mathrm Df(x)=A\) by Lemma~\ref{r13-lem-forms}(a).

(b) This is (a) for the vector function \(\operatorname{vec}X\mapsto\operatorname{vec}F(X)\) (Definitions~\ref{r13-def-derivative}
and~\ref{r13-def-differential}).

(c) By Theorem~\ref{r13-thm-vec}(d), \(\tr(C'\mathrm{d}X)=(\operatorname{vec}C)'\operatorname{vec}\mathrm{d}X\); so this is (b) with \(m=p=1\) and
\(A=(\operatorname{vec}C)'\), and \(\mathrm D\varphi(X)=(\operatorname{vec}C)'\) means \(\partial\varphi/\partial x_{ij}=c_{ij}\) for all \(i,j\).
\end{proof}
So, to find a derivative, we compute the differential and bring it to one of these forms.

\begin{theorem}[\citeat{tu-mc}{\S4.1, \S7, Proposition~15}, \citeat{abadir-magnus}{Exercises~13.1--13.3}]\label{r13-thm-diff-rules}
Let \(F\) and \(G\) be \(C^1\) matrix functions (vector and scalar functions included) of the same variable on an open set, and
let \(A\) be a constant matrix and \(\alpha\) a number.
\begin{enumerate}
\item \(\mathrm{d}A=0\), \(\mathrm{d}(\alpha F)=\alpha\,\mathrm{d}F\), \(\mathrm{d}(F+G)=\mathrm{d}F+\mathrm{d}G\), \(\mathrm{d}(F')=(\mathrm{d}F)'\),
\(\mathrm{d}\operatorname{vec}F=\operatorname{vec}\mathrm{d}F\) and, for square \(F\), \(\mathrm{d}\tr F=\tr\mathrm{d}F\).
\item \(\mathrm{d}(FG)=(\mathrm{d}F)G+F\,\mathrm{d}G\) and \(\mathrm{d}(F\otimes G)=(\mathrm{d}F)\otimes G+F\otimes\mathrm{d}G\), whenever the products are
defined.
\item (\textbf{Chain rule}) Let \(y=g(x)\) be \(C^1\) on an open \(S\subseteq\R^n\) with values in an open \(T\subseteq\R^m\), and let
\(z=f(y)\) be \(C^1\) on \(T\), with differentials \(\mathrm{d}z=A(y)\,\mathrm{d}y\) and \(\mathrm{d}y=B(x)\,\mathrm{d}x\). Then the differential of
\(x\mapsto f(g(x))\) is \(\mathrm{d}z=A(g(x))B(x)\,\mathrm{d}x\): we may substitute \(\mathrm{d}y=B(x)\,\mathrm{d}x\) into \(\mathrm{d}z=A(y)\,\mathrm{d}y\).
The same holds for matrix functions, through \(\operatorname{vec}\).
\end{enumerate}
\end{theorem}
\begin{proof}
(a) By Definition~\ref{r13-def-differential}, the \((k,l)\) entry of \(\mathrm{d}F\) is \(\sum_{i,j}(\partial f_{kl}/\partial x_{ij})\,\mathrm{d}x_{ij}\). The
partial derivatives of a constant are \(0\), and those of \(\alpha f_{kl}\) and of \(f_{kl}+g_{kl}\) are
\(\alpha\,\partial f_{kl}/\partial x_{ij}\) and \(\partial f_{kl}/\partial x_{ij}+\partial g_{kl}/\partial x_{ij}\) (rules of calculus); this gives
the first three rules. \(F'\), \(\operatorname{vec}F\) and \(\tr F\) rearrange entries of \(F\) or add them up, so their differentials
rearrange or add up the differentials of these entries.

(b)
\begin{itemize}
\item For \(C^1\) scalar functions \(u\), \(v\), the product rule gives
\(\mathrm{d}(uv)=\sum_{i,j}\bigl((\partial u/\partial x_{ij})v+u\,\partial v/\partial x_{ij}\bigr)\mathrm{d}x_{ij}=(\mathrm{d}u)v+u\,\mathrm{d}v\).
\item The \((k,l)\) entry of \(FG\) is \(\sum_rf_{kr}g_{rl}\), so by (a) its differential is
\(\sum_r\bigl((\mathrm{d}f_{kr})g_{rl}+f_{kr}\,\mathrm{d}g_{rl}\bigr)\), the \((k,l)\) entry of \((\mathrm{d}F)G+F\,\mathrm{d}G\).
\item The entries of \(F\otimes G\) are the products \(f_{ij}g_{st}\) (Definition~\ref{r13-def-kronecker}), with differentials
\((\mathrm{d}f_{ij})g_{st}+f_{ij}\,\mathrm{d}g_{st}\), the corresponding entries of \((\mathrm{d}F)\otimes G+F\otimes\mathrm{d}G\).
\end{itemize}
(c) By the chain rule of calculus, \(\mathrm D(f\circ g)(x)=\mathrm Df(g(x))\,\mathrm Dg(x)\), and by
Theorem~\ref{r13-thm-identification}(a), \(\mathrm Df(y)=A(y)\) and \(\mathrm Dg(x)=B(x)\). So the differential of \(f\circ g\) at \(x\) is
\(A(g(x))B(x)\,\mathrm{d}x\) (Definition~\ref{r13-def-differential}).
\end{proof}
In practice we write \(f(y)\) and \(f(x)\) for the same quantity seen as a function of \(y\) or of \(x\). The rules let us compute
differentials without writing down partial derivatives.

\begin{example}[\citeat{tu-mc}{Exercises~10, 12}]\label{r13-exm-differentials}
(a) For \(f(x)=x_1^2x_2\), Theorem~\ref{r13-thm-diff-rules}(b) gives \(\mathrm{d}(x_1^2)=2x_1\,\mathrm{d}x_1\) and
\[
\mathrm{d}f=(\mathrm{d}x_1^2)x_2+x_1^2\,\mathrm{d}x_2=2x_1x_2\,\mathrm{d}x_1+x_1^2\,\mathrm{d}x_2=(2x_1x_2,\ x_1^2)\,\mathrm{d}x
\]
so \(\mathrm Df(x)=(2x_1x_2,\ x_1^2)\) (Theorem~\ref{r13-thm-identification}(a)).

(b) For \(f(x)=(x_1^2-x_2^2,\ x_1x_2x_3)'\) with \(x=(x_1,x_2,x_3)'\),
\[
\mathrm{d}f=\begin{pmatrix}2x_1\,\mathrm{d}x_1-2x_2\,\mathrm{d}x_2\\x_2x_3\,\mathrm{d}x_1+x_1x_3\,\mathrm{d}x_2+x_1x_2\,\mathrm{d}x_3\end{pmatrix}
=\begin{pmatrix}2x_1&-2x_2&0\\x_2x_3&x_1x_3&x_1x_2\end{pmatrix}\mathrm{d}x
\]
and \(\mathrm Df(x)\) is the \(2\times3\) matrix on the right.
\end{example}

\begin{theorem}[\citeat{tu-mc}{Exercise~11}, \citeat{abadir-magnus}{Exercises~13.4, 13.6(a)}]\label{r13-thm-forms}
Let \(a\in\R^n\) and \(A\in\R^{m\times n}\) be constant.
\begin{enumerate}
\item \(\mathrm{d}(a'x)=a'\,\mathrm{d}x\): \(\mathrm D(a'x)=a'\), and the gradient of \(a'x\) is \(a\).
\item \(\mathrm{d}(Ax)=A\,\mathrm{d}x\): \(\mathrm D(Ax)=A\).
\item For \(A\in\R^{n\times n}\), \(\mathrm{d}(x'Ax)=x'(A+A')\,\mathrm{d}x\): \(\mathrm D(x'Ax)=x'(A+A')\), and the gradient is \((A+A')x\), which
is \(2Ax\) if \(A\) is symmetric.
\item For \(X\in\R^{n\times q}\), \(a\in\R^n\) and \(b\in\R^q\), \(\mathrm{d}(a'Xb)=a'(\mathrm{d}X)b\):
\(\partial(a'Xb)/\partial X=ab'\), and \(\mathrm D(a'Xb)=(b\otimes a)'\).
\end{enumerate}
\end{theorem}
\begin{proof}
(a), (b) By Theorem~\ref{r13-thm-diff-rules}(a),(b), \(\mathrm{d}(Ax)=(\mathrm{d}A)x+A\,\mathrm{d}x=A\,\mathrm{d}x\); take \(A=a'\) for (a), and apply
Theorem~\ref{r13-thm-identification}(a).

(c) By the product rule, \(\mathrm{d}(x'Ax)=(\mathrm{d}x)'Ax+x'A\,\mathrm{d}x\). The number \((\mathrm{d}x)'Ax\) equals its transpose
\(x'A'\,\mathrm{d}x\); so \(\mathrm{d}(x'Ax)=x'(A+A')\,\mathrm{d}x\), and Theorem~\ref{r13-thm-identification}(a) gives the derivative.

(d) \(\mathrm{d}(a'Xb)=a'(\mathrm{d}X)b=\tr\bigl(b\,a'\,\mathrm{d}X\bigr)=\tr\bigl((ab')'\mathrm{d}X\bigr)\) (Theorem~\ref{r1-thm-trace}(b)), so
\(\partial(a'Xb)/\partial X=ab'\) by Theorem~\ref{r13-thm-identification}(c), and \(\operatorname{vec}(ab')=b\otimes a\)
(Theorem~\ref{r13-thm-vec}(a)).
\end{proof}

\begin{example}[\citeat{searle-khuri}{Example~9.2}]\label{r13-exm-quadratic}
For
\[
A=\begin{pmatrix}1&3&5\\3&4&7\\5&7&9\end{pmatrix},\qquad q=x'Ax=x_1^2+6x_1x_2+10x_1x_3+4x_2^2+14x_2x_3+9x_3^2
\]
(Definition~\ref{r8-def-quadratic-form}), and by Theorem~\ref{r13-thm-forms}(c) the gradient is
\[
2Ax=\begin{pmatrix}2x_1+6x_2+10x_3\\6x_1+8x_2+14x_3\\10x_1+14x_2+18x_3\end{pmatrix}
\]
the partial derivatives of \(q\) read off its terms. The upper triangular
\[
B=\begin{pmatrix}1&6&10\\0&4&14\\0&0&9\end{pmatrix}
\]
gives the same \(q=x'Bx\), and indeed \(B+B'=2A\) (Lemma~\ref{r13-lem-forms}(c)): the derivative \(x'(B+B')\) is the same.
\end{example}

\needspace{24\baselineskip}
\subsection*{Traces, determinants and inverses}

\begin{theorem}[\citeat{tu-mc}{\S7, (30)--(32)}, \citeat{abadir-magnus}{Exercises~13.10(a), 13.21(a), 13.32, 13.36(a)}, \citeat{searle-khuri}{(9.26), (9.33)--(9.40)}, \citeat{bapat}{Result~9.5}]\label{r13-thm-trace-det}
Let \(X\) be a variable in \(\R^{n\times n}\) (\(n\geqslant2\)) and \(A\) a constant matrix; \(c_{ij}\) are the cofactors of \(X\)
(Definition~\ref{r5d-def-cofactor}).
\begin{enumerate}
\item \(\mathrm{d}\tr(AX)=\tr(A\,\mathrm{d}X)\), so \(\partial\tr(AX)/\partial X=A'\). For \(X\in\R^{n\times q}\) and \(A\in\R^{n\times n}\),
\(\mathrm{d}\tr(X'AX)=\tr\bigl(X'(A+A')\,\mathrm{d}X\bigr)\), so \(\partial\tr(X'AX)/\partial X=(A+A')X\).
\item \(\det X\) is a polynomial in the entries of \(X\), \(\partial\det X/\partial x_{ij}=c_{ij}\), and
\(\mathrm{d}\det X=\tr(\operatorname{adj}X\,\mathrm{d}X)\). If \(X\) is invertible, then
\[
\mathrm{d}\det X=\det X\,\tr(X^{-1}\mathrm{d}X)\qquad(\textbf{Jacobi's formula}),\qquad\frac{\partial\det X}{\partial X}=\det X\,(X^{-1})'
\]
\item The invertible matrices form an open set, on which \(X\mapsto X^{-1}\) is \(C^2\), and
\[
\mathrm{d}X^{-1}=-X^{-1}(\mathrm{d}X)X^{-1},\qquad\mathrm D(X^{-1})=-(X^{-1})'\otimes X^{-1}
\]
\item On the open set \(\{X:\det X\ne0\}\), \(\mathrm{d}\log\lvert\det X\rvert=\tr(X^{-1}\mathrm{d}X)\); so
\(\partial\log\det X/\partial X=(X^{-1})'\) where \(\det X>0\).
\item (\textbf{Matrices depending on one parameter}) Let the entries of \(A(t)\in\R^{n\times n}\) be \(C^1\) functions of \(t\) on an open
interval, and let \(\frac{\mathrm{d} A}{\mathrm{d}t}\) be the matrix of their derivatives. Then
\(\frac{\mathrm{d}}{\mathrm{d}t}\det A(t)=\tr\bigl(\operatorname{adj}A(t)\,\frac{\mathrm{d}A}{\mathrm{d}t}\bigr)\), and where \(\det A(t)>0\),
\(\frac{\mathrm{d}}{\mathrm{d}t}\log\det A(t)=\tr\bigl(A(t)^{-1}\frac{\mathrm{d}A}{\mathrm{d}t}\bigr)\).
\end{enumerate}
\end{theorem}
\begin{proof}
(a)
\begin{itemize}
\item By Theorem~\ref{r13-thm-diff-rules}(a),(b), \(\mathrm{d}\tr(AX)=\tr(A\,\mathrm{d}X)=\tr((A')'\mathrm{d}X)\); apply
Theorem~\ref{r13-thm-identification}(c).
\item By the product rule, \(\mathrm{d}\tr(X'AX)=\tr((\mathrm{d}X)'AX)+\tr(X'A\,\mathrm{d}X)\). As \(\tr M=\tr M'\) (Theorem~\ref{r1-thm-trace}(a)),
\(\tr((\mathrm{d}X)'AX)=\tr(X'A'\,\mathrm{d}X)\). So \(\mathrm{d}\tr(X'AX)=\tr\bigl(((A+A')X)'\mathrm{d}X\bigr)\); apply
Theorem~\ref{r13-thm-identification}(c).
\end{itemize}
(b)
\begin{itemize}
\item \(\det X\) is a polynomial in the entries, by induction on the size: \(\det(x_{11})=x_{11}\) for a
\(1\times1\) matrix, and the expansion along row \(1\) (Theorem~\ref{r5d-thm-expansion}),
\(\det X=\sum_kx_{1k}(-1)^{1+k}\det X(1|k)\), is a sum of products of entries of \(X\) with determinants of size \(n-1\).
\item By the expansion along row \(i\) (Theorem~\ref{r5d-thm-expansion}), \(\det X=\sum_kx_{ik}c_{ik}\), where
\(c_{ik}=(-1)^{i+k}\det X(i|k)\) does not involve row \(i\) of \(X\). With all entries other than \(x_{ij}\) fixed, \(\det X\) is
\(c_{ij}x_{ij}\) plus a constant, so \(\partial\det X/\partial x_{ij}=c_{ij}\).
\item Hence \(\mathrm{d}\det X=\sum_{i,j}c_{ij}\,\mathrm{d}x_{ij}=\sum_{j}\sum_i(\operatorname{adj}X)_{ji}\,\mathrm{d}x_{ij}=\sum_j(\operatorname{adj}X\,\mathrm{d}X)_{jj}
=\tr(\operatorname{adj}X\,\mathrm{d}X)\), as \((\operatorname{adj}X)_{ji}=c_{ij}\) (Theorem~\ref{r5d-thm-adjugate}).
\item If \(X\) is invertible, \(\operatorname{adj}X=\det X\cdot X^{-1}\) (Theorem~\ref{r5d-thm-adjugate}), which gives Jacobi's formula; and
\(\partial\det X/\partial X=(c_{ij})=(\operatorname{adj}X)'=\det X\,(X^{-1})'\).
\end{itemize}
(c)
\begin{itemize}
\item \(\det\) is a polynomial, hence continuous; if \(\det C\ne0\), then \(\det X\ne0\) on a ball around \(C\) (rules of calculus): the set of
invertible matrices is open.
\item On it, \(X^{-1}=\frac1{\det X}\operatorname{adj}X\) (Theorem~\ref{r5d-thm-adjugate}), whose entries are polynomials divided by the nonzero
polynomial \(\det X\); so \(X\mapsto X^{-1}\) is \(C^2\).
\item Differentiating \(X^{-1}X=I_n\) by the product rule, \((\mathrm{d}X^{-1})X+X^{-1}\mathrm{d}X=\mathrm{d}I_n=0\); multiplying on the right by
\(X^{-1}\) gives \(\mathrm{d}X^{-1}=-X^{-1}(\mathrm{d}X)X^{-1}\).
\item By Theorem~\ref{r13-thm-vec}(b), \(\operatorname{vec}\mathrm{d}X^{-1}=-\bigl((X^{-1})'\otimes X^{-1}\bigr)\operatorname{vec}\mathrm{d}X\); apply
Theorem~\ref{r13-thm-identification}(b).
\end{itemize}
\needspace{3\baselineskip}
(d)
\begin{itemize}
\item As \(\det\) is continuous, \(\det C>0\) (\(\det C<0\)) gives \(\det X>0\) (\(\det X<0\)) on a ball around \(C\), as in (c):
the sets \(\{\det X>0\}\) and \(\{\det X<0\}\) are open, so the chain rule applies on them.
\item On \(\{\det X>0\}\), \(\log\det X=h(\det X)\) with \(h(t)=\log t\); by the chain rule (Theorem~\ref{r13-thm-diff-rules}(c)) and (b),
\(\mathrm{d}\log\det X=\frac1{\det X}\,\mathrm{d}\det X=\tr(X^{-1}\mathrm{d}X)\).
\item On \(\{\det X<0\}\), \(\log\lvert\det X\rvert=h(-\det X)\), and \(\mathrm{d}\log\lvert\det X\rvert=\frac1{-\det X}(-\mathrm{d}\det X)=\tr(X^{-1}\mathrm{d}X)\).
\item Where \(\det X>0\), Theorem~\ref{r13-thm-identification}(c) with \(C=(X^{-1})'\) gives \(\partial\log\det X/\partial X=(X^{-1})'\).
\end{itemize}
(e)
\begin{itemize}
\item By the chain rule of calculus applied to \(t\mapsto A(t)\mapsto\det A(t)\), and (b),
\(\frac{\mathrm{d}}{\mathrm{d}t}\det A=\sum_{i,j}\frac{\partial\det A}{\partial a_{ij}}\,\frac{\mathrm{d}a_{ij}}{\mathrm{d}t}=\sum_{i,j}c_{ij}\frac{\mathrm{d}a_{ij}}{\mathrm{d}t}
=\tr\bigl(\operatorname{adj}A\,\frac{\mathrm{d}A}{\mathrm{d}t}\bigr)\), as in (b).
\item Where \(\det A>0\), \(\frac{\mathrm{d}}{\mathrm{d}t}\log\det A=\frac1{\det A}\frac{\mathrm{d}}{\mathrm{d}t}\det A=\tr\bigl(\frac{\operatorname{adj}A}{\det A}\frac{\mathrm{d}A}{\mathrm{d}t}\bigr)
=\tr\bigl(A^{-1}\frac{\mathrm{d}A}{\mathrm{d}t}\bigr)\).\qedhere
\end{itemize}
\end{proof}

\begin{example}\label{r9-exm-logdet}
Let \(f(X)=\log\det X\) on the set of \(X\in\R^{4\times4}\) with \(\det X>0\), which is open, as \(\det\) is a polynomial in the
entries (Theorem~\ref{r13-thm-trace-det}(b)) and hence continuous. Evaluate \(\partial f/\partial X\) at
\[
X_0=\bigl(\min(i,j)\bigr)=\begin{pmatrix}1&1&1&1\\1&2&2&2\\1&2&3&3\\1&2&3&4\end{pmatrix}
\]
and check it along a direction in which \(f(X_0+t\,\mathrm{d}X)\) is computed exactly.

\textit{Solution.} By Theorem~\ref{r13-thm-trace-det}(d), \(\partial f/\partial X=(X^{-1})'\) where \(\det X>0\).
\begin{itemize}
\item Let \(L\in\R^{4\times4}\) have \(l_{ik}=1\) for \(k\leqslant i\) and \(l_{ik}=0\) for \(k>i\). Then
\((LL')_{ij}=\sum_kl_{ik}l_{jk}\) counts the \(k\) with \(k\leqslant i\) and \(k\leqslant j\), so \(LL'=X_0\). As \(L'\) is upper
triangular with unit diagonal, \(\det L=\det L'=1\) (Theorems~\ref{r5d-thm-transpose} and~\ref{r5d-thm-basic}(c)), and
\(\det X_0=1\) (Theorem~\ref{r5d-thm-product}(b)). So \(\det X_0>0\) and \(f(X_0)=0\).
\item Let
\[
T=\begin{pmatrix}2&-1&0&0\\-1&2&-1&0\\0&-1&2&-1\\0&0&-1&1\end{pmatrix}
\]
Column \(j\) of \(X_0T\) is \(-c_{j-1}+2c_j-c_{j+1}\) for \(j\leqslant3\) (with \(c_0=0\)) and \(c_4-c_3\) for \(j=4\), where \(c_j\) is
column \(j\) of \(X_0\); these are \(e_1,\dots,e_4\). So \(X_0T=I_4\) and \(X_0^{-1}=T\) (Theorem~\ref{r2-thm-invertible}), and, as
\(T'=T\),
\[
\frac{\partial f}{\partial X}(X_0)=(X_0^{-1})'=T,\qquad\mathrm{d}f=\tr(T\,\mathrm{d}X)\ \text{at }X_0
\]
\item Check along the rank-one direction \(\mathrm{d}X=uv'\) with \(u=e_1-e_2\) and \(v=e_1\) (adding \(t\) to \(x_{11}\) and \(-t\) to
\(x_{21}\)): by the matrix determinant lemma (Theorem~\ref{r5d-thm-det-lemma}),
\(\det(X_0+tuv')=\det X_0\,(1+t\,v'X_0^{-1}u)=1+t\bigl(2-(-1)\bigr)=1+3t\). So \(f(X_0+tuv')=\log(1+3t)\) for \(t>-\frac13\), with
derivative \(3\) at \(t=0\); and \(\tr(Tuv')=v'Tu=t_{11}-t_{12}=3\) (Theorem~\ref{r1-thm-trace}(b)).
\end{itemize}
\end{example}

\begin{example}[\citeat{tu-mc}{Exercise~22}, \citeat{abadir-magnus}{Exercises~13.22(a), 13.25(a)}]\label{r13-exm-matrix-functions}
Let \(X\) range over the invertible \(n\times n\) matrices, and let \(A\), \(B\) be constant.

(a) For \(F(X)=AX^{-1}B\), \(\mathrm{d}F=A(\mathrm{d}X^{-1})B=-AX^{-1}(\mathrm{d}X)X^{-1}B\) (Theorems~\ref{r13-thm-diff-rules}
and~\ref{r13-thm-trace-det}(c)). By Theorem~\ref{r13-thm-vec}(b),
\(\operatorname{vec}\mathrm{d}F=-\bigl((X^{-1}B)'\otimes AX^{-1}\bigr)\operatorname{vec}\mathrm{d}X\), so \(\mathrm DF(X)=-(X^{-1}B)'\otimes AX^{-1}\).

(b) For \(\varphi(X)=\tr(AX^{-1})\), \(\mathrm{d}\varphi=-\tr(AX^{-1}(\mathrm{d}X)X^{-1})=-\tr(X^{-1}AX^{-1}\mathrm{d}X)\) (Theorem~\ref{r1-thm-trace}(c)),
so \(\partial\varphi/\partial X=-(X^{-1}AX^{-1})'\) (Theorem~\ref{r13-thm-identification}(c)).
\end{example}

\begin{example}[\citeat{bapat}{Results~9.1, 9.2, 9.5}]\label{r13-exm-variance-components}
In the \textbf{one-way random effects model} the \(m\) observations of a group are \(y_j=\mu+a+\varepsilon_j\) (\(j=1,\dots,m\)), where
\(a\), \(\varepsilon_1,\dots,\varepsilon_m\) are uncorrelated random variables with means \(0\), \(\var(a)=\sigma_1^2\) and
\(\var(\varepsilon_j)=\sigma^2>0\) (the variance components). With \(y=\mu\one_m+[\one_m\ I_m](a,\varepsilon')'\),
Theorem~\ref{r11-thm-linear}(b) gives
\[
V=\var(y)=[\one_m\ I_m]\begin{pmatrix}\sigma_1^2&0\\0&\sigma^2I_m\end{pmatrix}\begin{pmatrix}\one_m'\\I_m\end{pmatrix}=\sigma_1^2J_m+\sigma^2I_m
\]
Maximum-likelihood methods for \(\sigma_1^2\), \(\sigma^2\) need the derivatives of \(\log\det V\). As \(\mathrm{d}V/\mathrm{d}\sigma_1^2=J_m\) and
\(\mathrm{d}V/\mathrm{d}\sigma^2=I_m\), Theorem~\ref{r13-thm-trace-det}(e) gives
\[
\frac{\partial\log\det V}{\partial\sigma_1^2}=\tr(V^{-1}J_m),\qquad\frac{\partial\log\det V}{\partial\sigma^2}=\tr(V^{-1})
\]
By Theorem~\ref{r5d-thm-woodbury}(b) (\(A=\sigma^2I_m\), \(u=\sigma_1^2\one_m\), \(v=\one_m\)),
\(V^{-1}=\frac1{\sigma^2}\bigl(I_m-\frac{\sigma_1^2}{\sigma^2+m\sigma_1^2}J_m\bigr)\), so \(\tr(V^{-1}J_m)=\frac m{\sigma^2+m\sigma_1^2}\). This
agrees with \(\det V=\sigma^{2(m-1)}(\sigma^2+m\sigma_1^2)\) (Theorem~\ref{r5d-thm-det-lemma}), whose logarithm has the derivative
\(\frac m{\sigma^2+m\sigma_1^2}\) in \(\sigma_1^2\). For \(m=4\), \(\sigma_1^2=1\), \(\sigma^2=2\):
\[
V=J_4+2I_4=\begin{pmatrix}3&1&1&1\\1&3&1&1\\1&1&3&1\\1&1&1&3\end{pmatrix},\qquad
V^{-1}=\frac12\Bigl(I_4-\frac16J_4\Bigr)=\frac1{12}\begin{pmatrix}5&-1&-1&-1\\-1&5&-1&-1\\-1&-1&5&-1\\-1&-1&-1&5\end{pmatrix}
\]
so \(\tr(V^{-1}J_4)=\one_4'V^{-1}\one_4=\frac{4\cdot2}{12}=\frac23=\frac m{\sigma^2+m\sigma_1^2}\) and
\(\tr(V^{-1})=\frac{4\cdot5}{12}=\frac53\); and \(\det V=2^3\cdot6=48\). Indeed
\(\frac{m-1}{\sigma^2}+\frac1{\sigma^2+m\sigma_1^2}=\frac32+\frac16=\frac53\).
\end{example}

\needspace{16\baselineskip}
\begin{theorem}[\citeat{searle-khuri}{\S\S9.5.2--9.5.4}, \citeat{abadir-magnus}{Exercises~13.25(b), 13.26(b)}]\label{r13-thm-symmetric}
Let \(X\) range over an open set of symmetric \(n\times n\) matrices, the variables being the entries \(x_{ij}\) with \(i\geqslant j\)
(that is, \(\operatorname{vech}X\)). For a function \(\varphi\) of these variables, write \(\partial\varphi/\partial X\) for the symmetric
matrix with \((i,j)\) and \((j,i)\) entries \(\partial\varphi/\partial x_{ij}\); and write \(\diag(M)\) for the diagonal part of a square
matrix \(M\).
\begin{enumerate}
\item Let \(\varphi\) be \(C^1\), with differential \(\mathrm{d}\varphi=\tr(C\,\mathrm{d}X)\) for symmetric \(\mathrm{d}X\), where \(C=(c_{ij})\) is
symmetric (and may depend on \(X\)). Then
\[
\frac{\partial\varphi}{\partial x_{ij}}=2c_{ij}\ (i>j),\qquad\frac{\partial\varphi}{\partial x_{ii}}=c_{ii},\qquad
\frac{\partial\varphi}{\partial(\operatorname{vech}X)'}=(\operatorname{vec}C)'D_n
\]
that is, \(\partial\varphi/\partial X=2C-\diag(C)\).
\item If \(X\) is PD, \(\partial\log\det X/\partial X=2X^{-1}-\diag(X^{-1})\).
\item \(\partial\tr(AX)/\partial X=A+A'-\diag(A)\) for constant \(A\in\R^{n\times n}\).
\item \(\partial(a'Xa)/\partial X=2aa'-\diag(aa')\) for constant \(a\in\R^n\).
\end{enumerate}
\end{theorem}
\begin{proof}
(a)
\begin{itemize}
\item Let \(E_{ij}=e_ie_j'\). The differential \(\mathrm{d}\varphi=\mathrm D\varphi\,\mathrm{d}\operatorname{vech}X\) at the unit vector of the
variable \(x_{ij}\) is \(\partial\varphi/\partial x_{ij}\) (Definitions~\ref{r13-def-derivative} and~\ref{r13-def-differential}). This
increment changes only \(x_{ij}\), by \(1\), that is, the \((i,j)\) and \((j,i)\) entries of \(X\): \(\mathrm{d}X=E_{ij}+E_{ji}\) for \(i>j\) and
\(\mathrm{d}X=E_{ii}\) for \(i=j\).
\item Hence \(\partial\varphi/\partial x_{ij}=\tr(CE_{ij})+\tr(CE_{ji})=c_{ji}+c_{ij}=2c_{ij}\) for \(i>j\)
(\(\tr(Ce_ie_j')=e_j'Ce_i\), Theorem~\ref{r1-thm-trace}(b)), and \(\partial\varphi/\partial x_{ii}=\tr(CE_{ii})=c_{ii}\).
\item Vector form: \(\operatorname{vec}\mathrm{d}X=D_n\,\mathrm{d}\operatorname{vech}X\) (Definition~\ref{r13-def-duplication}), so by
Theorem~\ref{r13-thm-vec}(d) (\(C'=C\))
\[
\mathrm{d}\varphi=(\operatorname{vec}C)'\operatorname{vec}\mathrm{d}X=(\operatorname{vec}C)'D_n\,\mathrm{d}\operatorname{vech}X
\]
and Theorem~\ref{r13-thm-identification}(a) applies. The entries of \((\operatorname{vec}C)'D_n\) are \(c_{ii}\) and \(c_{ij}+c_{ji}=2c_{ij}\)
(Proposition~\ref{r13-prop-duplication}(a)), as above.
\end{itemize}
(b) \(\log\det X\) is the composite of the linear map \(\operatorname{vech}X\mapsto X\) and \(\log\det\); by the chain rule
(Theorem~\ref{r13-thm-diff-rules}(c)) and Theorem~\ref{r13-thm-trace-det}(d), \(\mathrm{d}\log\det X=\tr(X^{-1}\mathrm{d}X)\) for symmetric
\(\mathrm{d}X\), and \(C=X^{-1}\) is symmetric. Apply (a).

(c) For symmetric \(\mathrm{d}X\), \(\tr(A\,\mathrm{d}X)=\tr((A\,\mathrm{d}X)')=\tr(\mathrm{d}X\,A')=\tr(A'\mathrm{d}X)\) (Theorem~\ref{r1-thm-trace}(a),(b)),
so \(\mathrm{d}\tr(AX)=\tr\bigl(\frac12(A+A')\mathrm{d}X\bigr)\) with the symmetric \(C=\frac12(A+A')\). Apply (a).

(d) \(a'Xa=\tr(Xaa')=\tr(aa'X)\) (Theorem~\ref{r1-thm-trace}(b)), so this is (c) with \(A=aa'\).
\end{proof}
So the derivative with respect to a symmetric matrix is not simply the unrestricted one: for \(\log\det X\) the unrestricted
derivative is \((X^{-1})'=X^{-1}\), while each off-diagonal variable \(x_{ij}\) (\(i>j\)) enters \(X\) twice and gets the derivative
\(2(X^{-1})_{ij}\).

\subsection*{Second differentials and the Hessian matrix}

We also use without proof~\citeat{tu-mc}{\S9}, \citeat{harville}{\S15.1}: if \(f\) is \(C^2\) on an open set, then \(\partial^2f/\partial x_i\partial x_j=\partial^2f/\partial x_j\partial x_i\), where
\(\partial^2f/\partial x_i\partial x_j\) is the partial derivative with respect to \(x_i\) of \(\partial f/\partial x_j\) (\textbf{symmetry of second
derivatives}).

\begin{definition}[\citeat{tu-mc}{\S9}, \citeat{searle-khuri}{\S9.5.7}, \citeat{harville}{\S15.1}]\label{r13-def-hessian}
Let \(f\) be a \(C^2\) scalar function on an open \(S\subseteq\R^n\). The \textbf{Hessian matrix} of \(f\) at \(x\in S\) is the \(n\times n\)
matrix
\[
\mathrm Hf(x)=\frac{\partial^2f(x)}{\partial x\,\partial x'}\quad\text{with }(i,j)\text{ entry }\frac{\partial^2f(x)}{\partial x_i\,\partial x_j}
\]
It is the derivative of the gradient, \(\mathrm Hf(x)=\mathrm D\bigl((\mathrm Df)'\bigr)(x)\), and it is symmetric. The \textbf{second
differential} \(\mathrm{d}^2f(x)\) is the differential of \(x\mapsto\mathrm{d}f(x)=\mathrm Df(x)\,\mathrm{d}x\) with \(\mathrm{d}x\) held fixed. For a
scalar function of a matrix \(X\), \(\mathrm Hf(X)=\partial^2f/\partial\operatorname{vec}X\,\partial(\operatorname{vec}X)'\).
\end{definition}
With \(a_j=\partial f/\partial x_j\), \(\mathrm{d}f(x)=\sum_ja_j(x)\,\mathrm{d}x_j\), and by Theorem~\ref{r13-thm-diff-rules},
\(\mathrm{d}^2f(x)=\sum_j(\mathrm{d}a_j)\,\mathrm{d}x_j=\sum_j\sum_i(\partial a_j/\partial x_i)\,\mathrm{d}x_i\,\mathrm{d}x_j\), that is,
\[
\mathrm{d}^2f(x)=(\mathrm{d}x)'\,\mathrm Hf(x)\,\mathrm{d}x
\]
To compute \(\mathrm{d}^2f\) we apply the rules of Theorem~\ref{r13-thm-diff-rules} to \(\mathrm{d}f\), treating \(\mathrm{d}x\) (or \(\mathrm{d}X\))
as a constant: \(\mathrm{d}(\mathrm{d}x)=0\).

\needspace{14\baselineskip}
\begin{theorem}[Second identification theorem; \citeat{tu-mc}{Propositions~16, 18}, \citeat{abadir-magnus}{Exercises~13.47(c), 13.48(c)}]\label{r13-thm-identification2}
Let \(f\) be a \(C^2\) scalar function on an open set.
\begin{enumerate}
\item For a vector variable \(x\in\R^n\) and \(B\in\R^{n\times n}\): \(\mathrm{d}^2f(x)=(\mathrm{d}x)'B\,\mathrm{d}x\) for every \(\mathrm{d}x\) if and only if
\(\mathrm Hf(x)=\frac12(B+B')\).
\item For a matrix variable \(X\in\R^{n\times q}\), \(A\in\R^{q\times q}\) and \(B\in\R^{n\times n}\):
\(\mathrm{d}^2f(X)=\tr\bigl(A(\mathrm{d}X)'B\,\mathrm{d}X\bigr)\) for every \(\mathrm{d}X\) if and only if \(\mathrm Hf(X)=\frac12(A'\otimes B+A\otimes B')\).
\item For \(X\in\R^{n\times q}\) and \(A,B\in\R^{q\times n}\): \(\mathrm{d}^2f(X)=\tr\bigl(A(\mathrm{d}X)B\,\mathrm{d}X\bigr)\) for every \(\mathrm{d}X\) if and only
if \(\mathrm Hf(X)=\frac12K_{qn}(A'\otimes B+B'\otimes A)\).
\end{enumerate}
\end{theorem}
\begin{proof}
(a) Both \(\mathrm Hf(x)\) and \(\frac12(B+B')\) are symmetric. If \(\mathrm{d}^2f(x)=(\mathrm{d}x)'B\,\mathrm{d}x\) for every \(\mathrm{d}x\), then
\((\mathrm{d}x)'\mathrm Hf(x)\,\mathrm{d}x=(\mathrm{d}x)'B\,\mathrm{d}x\), and Lemma~\ref{r13-lem-forms}(c) gives \(2\mathrm Hf(x)=B+B'\). Conversely,
\((\mathrm{d}x)'\frac12(B+B')\,\mathrm{d}x=(\mathrm{d}x)'B\,\mathrm{d}x\), as \((\mathrm{d}x)'B'\,\mathrm{d}x\) is a number equal to its transpose.

(b) By Theorems~\ref{r1-thm-trace}(b) and~\ref{r13-thm-vec}(d),(b),
\[
\tr(A(\mathrm{d}X)'B\,\mathrm{d}X)=\tr((\mathrm{d}X)'B(\mathrm{d}X)A)=(\operatorname{vec}\mathrm{d}X)'\operatorname{vec}(B(\mathrm{d}X)A)
=(\operatorname{vec}\mathrm{d}X)'(A'\otimes B)\operatorname{vec}\mathrm{d}X
\]
By (a) for the variable \(\operatorname{vec}X\),
\(\mathrm Hf(X)=\frac12\bigl(A'\otimes B+(A'\otimes B)'\bigr)=\frac12(A'\otimes B+A\otimes B')\) (Theorem~\ref{r13-thm-kron-rules}(e)).

(c)
\begin{itemize}
\item As \(\tr M=\tr M'\), \(\tr(A(\mathrm{d}X)B\,\mathrm{d}X)=\tr((\mathrm{d}X)'B'(\mathrm{d}X)'A')=(\operatorname{vec}\mathrm{d}X)'\operatorname{vec}(B'(\mathrm{d}X)'A')
=(\operatorname{vec}\mathrm{d}X)'(A\otimes B')\operatorname{vec}(\mathrm{d}X)'\), and \(\operatorname{vec}(\mathrm{d}X)'=K_{nq}\operatorname{vec}\mathrm{d}X\)
(Definition~\ref{r13-def-commutation}). So \(\mathrm{d}^2f=(\operatorname{vec}\mathrm{d}X)'(A\otimes B')K_{nq}\operatorname{vec}\mathrm{d}X\).
\item By (a), \(\mathrm Hf(X)=\frac12\bigl((A\otimes B')K_{nq}+K_{nq}'(A'\otimes B)\bigr)\), with \(K_{nq}'=K_{qn}\) (Theorem~\ref{r13-thm-commutation}(b)).
By Theorem~\ref{r13-thm-commutation}(d) for \(B'\in\R^{n\times q}\) and \(A\in\R^{q\times n}\), \(K_{qn}(B'\otimes A)=(A\otimes B')K_{nq}\). Hence
\(\mathrm Hf(X)=\frac12K_{qn}(A'\otimes B+B'\otimes A)\).\qedhere
\end{itemize}
\end{proof}

\begin{example}[\citeat{tu-mc}{Exercises~23, 25, 29}, \citeat{abadir-magnus}{Exercises~13.46, 13.47(b), 13.52}]\label{r13-exm-hessians}
(a) For \(f(x)=x'Ax\) (\(A\in\R^{n\times n}\)), \(\mathrm{d}f=(\mathrm{d}x)'Ax+x'A\,\mathrm{d}x\) and \(\mathrm{d}^2f=2(\mathrm{d}x)'A\,\mathrm{d}x\), so
\(\mathrm Hf(x)=A+A'\) (Theorem~\ref{r13-thm-identification2}(a)); for symmetric \(A\), \(\mathrm Hf=2A\).

(b) For \(f(X)=\tr(X'AX)\) with \(X\in\R^{n\times q}\), \(\mathrm{d}f=\tr(X'(A+A')\,\mathrm{d}X)\) (Theorem~\ref{r13-thm-trace-det}(a)) and
\(\mathrm{d}^2f=\tr((\mathrm{d}X)'(A+A')\,\mathrm{d}X)=\tr\bigl(I_q(\mathrm{d}X)'(A+A')\,\mathrm{d}X\bigr)\); by Theorem~\ref{r13-thm-identification2}(b),
\(\mathrm Hf(X)=I_q\otimes(A+A')\).

(c) For \(f(X)=\log\det X\) (\(\det X>0\)), \(\mathrm{d}f=\tr(X^{-1}\mathrm{d}X)\) and, by Theorem~\ref{r13-thm-trace-det}(c),
\(\mathrm{d}^2f=\tr((\mathrm{d}X^{-1})\,\mathrm{d}X)=-\tr\bigl(X^{-1}(\mathrm{d}X)X^{-1}\mathrm{d}X\bigr)\). Theorem~\ref{r13-thm-identification2}(c) with
\(q=n\), \(A=-X^{-1}\), \(B=X^{-1}\) gives \(\mathrm Hf(X)=-K_n\bigl((X^{-1})'\otimes X^{-1}\bigr)\).
\end{example}

\subsection*{Conditions for extrema}

For functions of one variable we use without proof~\citeat{abadir-magnus}{\S A.3.4}: (\textbf{Fermat's rule}) if \(h\) is differentiable on an open interval and has a
local minimum or maximum at a point \(t_0\) of it, then \(h'(t_0)=0\); and (\textbf{Taylor's formula}) if \(h\) is \(C^2\) on an open
interval that contains \([0,1]\), then \(h(1)=h(0)+h'(0)+\frac12h''(\tau)\) for some \(\tau\) with \(0<\tau<1\).

\begin{definition}[\citeat{magnus-neudecker}{\S7.2, \S7.7}, \citeat{abadir-magnus}{\S A.4.3}]\label{r13-def-extremum}
Let \(f\) be a scalar function on \(S\subseteq\R^n\) and \(c\in S\). \(f\) has a \textbf{local minimum} at \(c\) if \(f(x)\geqslant f(c)\) for all
\(x\in S\) in some ball \(B(c,\delta)\), a \textbf{strict local minimum} if \(f(x)>f(c)\) there for \(x\ne c\), and a \textbf{global minimum}
(\textbf{strict global minimum}) if \(f(x)\geqslant f(c)\) for all \(x\in S\) (\(f(x)>f(c)\) for \(x\ne c\)); maxima likewise with the
inequalities reversed. If \(S\) is open and \(f\) is \(C^1\), \(c\) is a \textbf{stationary point} of \(f\) if \(\mathrm Df(c)=0\). A set \(S\) is
\textbf{convex} if \(tx+(1-t)y\in S\) for all \(x,y\in S\) and \(0\leqslant t\leqslant1\).
\end{definition}

\needspace{16\baselineskip}
\begin{lemma}[\citeat{magnus-neudecker}{Theorems~7.2, 7.4, 7.7}]\label{r13-lem-line}
Let \(S\subseteq\R^n\) be open, \(f\) a \(C^2\) scalar function on \(S\), \(c\in S\) and \(u\in\R^n\). Let
\(I=\{t\in\R:c+tu\in S\}\) and \(h(t)=f(c+tu)\) for \(t\in I\).
\begin{enumerate}
\item \(I\) is open and contains \(0\), and \(h\) is \(C^2\) on \(I\), with \(h'(t)=\mathrm Df(c+tu)u\) and
\(h''(t)=u'\mathrm Hf(c+tu)u\).
\item Suppose that \(\lVert u\rVert<\delta\) for a ball \(B(c,\delta)\subseteq S\), or that \(S\) is convex and \(c+u\in S\). Then, for
some \(0<\tau<1\),
\[
f(c+u)=f(c)+\mathrm Df(c)u+\tfrac12u'\mathrm Hf(c+\tau u)u
\]
\end{enumerate}
\end{lemma}
\begin{proof}
(a)
\begin{itemize}
\item \(0\in I\) as \(c\in S\). \(I\) is open: if \(c+t_0u\in S\), some ball \(B(c+t_0u,\varepsilon)\) lies in \(S\), and it contains
\(c+tu\) for \(\lvert t-t_0\rvert\,\lVert u\rVert<\varepsilon\).
\item By the chain rule of calculus (\(t\mapsto c+tu\) has the derivative \(u\)),
\(h'(t)=\mathrm Df(c+tu)u=\sum_ju_j\,\partial f/\partial x_j(c+tu)\).
\item The same rule for each \(t\mapsto\partial f/\partial x_j(c+tu)\) gives
\(h''(t)=\sum_j\sum_iu_ju_i\,\partial^2f/\partial x_i\partial x_j(c+tu)=u'\mathrm Hf(c+tu)u\) (Definition~\ref{r13-def-hessian}), which is
continuous in \(t\): \(h\) is \(C^2\) on \(I\).
\end{itemize}
(b)
\begin{itemize}
\item In the first case, \(c+tu\in B(c,\delta)\subseteq S\) for \(\lvert t\rvert<\delta/\lVert u\rVert\) (for every \(t\) if \(u=0\)), and
\(\delta/\lVert u\rVert>1\): \(I\) contains an open interval that contains \([0,1]\).
\item In the second case, \(I\) is an interval: for \(t\) between two points \(t_1\), \(t_2\) of \(I\), \(c+tu\) is a convex combination of
\(c+t_1u\) and \(c+t_2u\), so it lies in \(S\). By (a) \(I\) is open, and it contains \(0\) and \(1\).
\item By Taylor's formula and (a), for some \(0<\tau<1\),
\(f(c+u)=h(1)=h(0)+h'(0)+\frac12h''(\tau)=f(c)+\mathrm Df(c)u+\frac12u'\mathrm Hf(c+\tau u)u\).\qedhere
\end{itemize}
\end{proof}

\needspace{20\baselineskip}
\begin{theorem}[\citeat{magnus-neudecker}{Theorems~7.2, 7.4, 7.7, 7.8, 7.13}]\label{r13-thm-extrema}
Let \(S\subseteq\R^n\) be open, \(f\) a \(C^2\) scalar function on \(S\), and \(c\in S\).
\begin{enumerate}
\item (\textbf{First-order condition}) If \(f\) has a local minimum or maximum at \(c\), then \(\mathrm Df(c)=0\).
\item If \(f\) has a local minimum (maximum) at \(c\), then \(\mathrm Hf(c)\) is PSD (NSD).
\item (\textbf{Second-order condition}) If \(\mathrm Df(c)=0\) and \(\mathrm Hf(c)\) is PD (ND), then \(f\) has a strict local minimum (maximum) at
\(c\). If \(\mathrm Df(c)=0\) and \(\mathrm Hf(c)\) is indefinite, then \(f\) has neither a local minimum nor a local maximum at \(c\).
\item (\textbf{Global minimum}) Let \(S\) also be convex. If \(\mathrm Df(c)=0\) and \(\mathrm Hf(x)\) is PSD for every \(x\in S\), then \(f\) has a global
minimum at \(c\); if \(\mathrm Hf(x)\) is PD for every \(x\in S\), the global minimum at \(c\) is strict.
\item (\textbf{Constraints}) Let \(S\) be convex, let \(g\) be \(C^2\) on \(S\) with values in \(\R^m\), let \(l\in\R^m\), and let
\(\psi(x)=f(x)-l'g(x)\). If \(\mathrm D\psi(c)=0\), \(g(c)=0\), and \(\mathrm H\psi(x)\) is PSD (PD) for every \(x\in S\), then \(f(x)\geqslant f(c)\)
(\(f(x)>f(c)\) for \(x\ne c\)) for every \(x\in S\) with \(g(x)=0\).
\end{enumerate}
\end{theorem}
\begin{proof}
Throughout, for a given \(u\in\R^n\), \(I\) and \(h(t)=f(c+tu)\) are as in Lemma~\ref{r13-lem-line}.

(a) Let \(f(x)\geqslant f(c)\) for \(x\in B(c,\delta)\subseteq S\) (a local minimum; \(\delta\) small, as \(S\) is open). For \(u=e_j\), \(h\) has a local minimum at \(t=0\), as
\(c+te_j\in B(c,\delta)\) for \(\lvert t\rvert<\delta\). By Fermat's rule and Lemma~\ref{r13-lem-line}(a),
\(0=h'(0)=\mathrm Df(c)e_j=\partial f/\partial x_j(c)\), for every \(j\): \(\mathrm Df(c)=0\). For a maximum apply this to \(-f\).

(b)
\begin{itemize}
\item Let \(f\) have a local minimum at \(c\), say \(f(x)\geqslant f(c)\) for \(x\in B(c,\delta)\subseteq S\) (\(\delta\) small, as in (a)), and suppose
\(u'\mathrm Hf(c)u<0\) for some \(u\). By Lemma~\ref{r13-lem-line}(a), \(h''(0)=u'\mathrm Hf(c)u<0\). As \(h''\) is continuous and \(I\) is
open, \(h''(t)<0\) for \(\lvert t\rvert\leqslant\varepsilon\), for some \(\varepsilon>0\) with \([-\varepsilon,\varepsilon]\subseteq I\).
\item Choose \(0<s_0\leqslant\varepsilon\) with \(s_0\lVert u\rVert<\delta\). By Lemma~\ref{r13-lem-line}(b) for the vector \(s_0u\) (first
case) and \(\mathrm Df(c)=0\) (by (a)), for some \(0<\tau<1\),
\[
f(c+s_0u)=f(c)+\tfrac12s_0^2u'\mathrm Hf(c+\tau s_0u)u=f(c)+\tfrac12s_0^2h''(\tau s_0)<f(c)
\]
as \(0<\tau s_0<\varepsilon\). But \(c+s_0u\in B(c,\delta)\), so \(f(c+s_0u)\geqslant f(c)\) \(\contradiction\)
\item Hence \(\mathrm Hf(c)\) is PSD. For a maximum apply this to \(-f\), whose Hessian is \(-\mathrm Hf\).
\end{itemize}
(c)
\begin{itemize}
\item Let \(\mathrm Hf(c)\) be PD with smallest eigenvalue \(\lambda>0\) (Theorem~\ref{r8-thm-eigen-criterion}(a)); then
\(u'\mathrm Hf(c)u\geqslant\lambda\lVert u\rVert^2\) for every \(u\) (Theorem~\ref{r9-thm-rayleigh}(a)).
\item The second partial derivatives are continuous, so the \(2n^2\) functions
\(\frac\lambda{2n}\pm\bigl(\partial^2f/\partial x_i\partial x_j(z)-\partial^2f/\partial x_i\partial x_j(c)\bigr)\) of \(z\), positive at \(z=c\), are positive on
a ball \(B(c,\delta)\subseteq S\) (rules of calculus). For \(z\in B(c,\delta)\), every entry of \(E=\mathrm Hf(z)-\mathrm Hf(c)\) is less than
\(\frac\lambda{2n}\) in absolute value, so \(\lVert E\rVert_F^2<n^2\frac{\lambda^2}{4n^2}\), and by Theorem~\ref{r6-thm-cauchy-schwarz}(a) and
Proposition~\ref{r9-prop-frobenius}(b), \(\lvert u'Eu\rvert\leqslant\lVert u\rVert\,\lVert Eu\rVert\leqslant\lVert E\rVert_F\lVert u\rVert^2<\frac\lambda2\lVert u\rVert^2\) for \(u\ne0\).
\item For \(0<\lVert u\rVert<\delta\), Lemma~\ref{r13-lem-line}(b) (first case) with \(\mathrm Df(c)=0\) gives
\(f(c+u)-f(c)=\frac12u'\mathrm Hf(z)u\) with \(z=c+\tau u\in B(c,\delta)\); so
\(f(c+u)-f(c)=\frac12u'\mathrm Hf(c)u+\frac12u'Eu>\frac12\lambda\lVert u\rVert^2-\frac14\lambda\lVert u\rVert^2>0\): a strict local minimum.
For ND apply this to \(-f\).
\item If \(\mathrm Hf(c)\) is indefinite, it is neither PSD nor NSD (Definition~\ref{r8-def-definite}), so by (b) \(c\) is neither a local
minimum nor a local maximum.
\end{itemize}
(d) Let \(x\in S\) and \(u=x-c\). By Lemma~\ref{r13-lem-line}(b) (\(S\) convex, \(c+u=x\in S\)) and \(\mathrm Df(c)=0\),
\(f(x)=f(c)+\frac12u'\mathrm Hf(c+\tau u)u\), where \(c+\tau u=(1-\tau)c+\tau x\in S\). This is \(\geqslant f(c)\) if \(\mathrm Hf(c+\tau u)\) is
PSD, and \(>f(c)\) for \(x\ne c\) if it is PD.

(e) \(\psi\) is \(C^2\) on \(S\) (rules of calculus), so by (d) applied to \(\psi\), \(\psi(x)\geqslant\psi(c)\) for every \(x\in S\) (\(>\) for
\(x\ne c\) in the PD case). If \(g(x)=0\), then \(f(x)=\psi(x)\), and \(\psi(c)=f(c)\) as \(g(c)=0\).
\end{proof}
To use (e) one solves \(\mathrm D\psi(x)=0\), \(g(x)=0\) for \(x\) and the \textbf{Lagrange multipliers} \(l\) (\textbf{Lagrange's method}), and
then checks the Hessian of the \textbf{Lagrangian} \(\psi\). For linear constraints \(g(x)=Rx-r\), \(\mathrm H\psi=\mathrm Hf\).

We also use without proof \textbf{Lagrange's theorem}~\citeat{magnus-neudecker}{Theorem~7.10}: let \(S\subseteq\R^n\) be open, let \(f\) be a \(C^1\) scalar function and \(g\)
a \(C^1\) function with values in \(\R^m\) (\(m<n\)) on \(S\), and let \(c\in S\) with \(g(c)=0\) and \(\rank\mathrm Dg(c)=m\). If
\(f(x)\geqslant f(c)\) for every \(x\) with \(g(x)=0\) in some ball \(B(c,\delta)\subseteq S\) (a local minimum under the constraint
\(g(x)=0\)), or \(f(x)\leqslant f(c)\) for all these \(x\) (a local maximum), then there is exactly one \(l\in\R^m\) with
\[
\mathrm Df(c)=l'\,\mathrm Dg(c)
\]
that is, \(\mathrm D\psi(c)=0\) for the Lagrangian \(\psi=f-l'g\). For \(S\), \(f\) and \(g\) as here but with \(m\leqslant n\), a point
\(c\in S\) with \(g(c)=0\), \(\rank\mathrm Dg(c)=m\) and \(\mathrm Df(c)=l'\,\mathrm Dg(c)\) for some \(l\in\R^m\) is a
\textbf{stationary point of \(f\) under the constraint \(g(x)=0\)}. So, for \(m<n\), Lagrange's method finds all the points where
\(f\) has a local minimum or maximum under the constraint and \(\rank\mathrm Dg=m\); Theorem~\ref{r13-thm-extrema}(e) shows that a
stationary point is a minimum when the Hessian of \(\psi\) is PSD on a convex \(S\).

\begin{example}[\citeat{tu-mc}{Exercises~14, 15}, \citeat{harville}{Theorem~19.1.1}]\label{r13-exm-quadratic-min}
Let \(A\in\R^{n\times n}\) be PD, \(b\in\R^n\), and \(f(x)=\frac12x'Ax-b'x\).

(a) By Theorem~\ref{r13-thm-forms}, \(\mathrm{d}f=x'A\,\mathrm{d}x-b'\,\mathrm{d}x=(Ax-b)'\mathrm{d}x\), and \(\mathrm{d}^2f=(\mathrm{d}x)'A\,\mathrm{d}x\), so
\(\mathrm Hf=A\) is PD at every \(x\). The only stationary point is \(\hat x=A^{-1}b\), and by Theorem~\ref{r13-thm-extrema}(d) it is the strict
global minimizer. Directly: \((x-\hat x)'A(x-\hat x)=x'Ax-2b'x+b'A^{-1}b=2f(x)+b'A^{-1}b\), which is smallest, \(0\), at \(x=\hat x\);
the minimum is \(f(\hat x)=-\frac12b'A^{-1}b\).

(b) Minimize \(f\) subject to \(c'x=0\) (\(c\ne0\)). The Lagrangian \(\psi(x)=\frac12x'Ax-b'x-\lambda c'x\) has
\(\mathrm{d}\psi=(Ax-b-\lambda c)'\mathrm{d}x\), so the conditions are \(Ax=b+\lambda c\) and \(c'x=0\), or, in block form,
\[
\begin{pmatrix}A&c\\c'&0\end{pmatrix}\begin{pmatrix}x\\-\lambda\end{pmatrix}=\begin{pmatrix}b\\0\end{pmatrix}
\]
From \(x=A^{-1}(b+\lambda c)\) and \(0=c'x=c'A^{-1}b+\lambda c'A^{-1}c\), with \(c'A^{-1}c>0\) (\(A^{-1}\) is PD, Proposition~\ref{r8-prop-pd-basic}(c)), we get
\(\tilde\lambda=-c'A^{-1}b/c'A^{-1}c\) and \(\tilde x=A^{-1}b-\frac{c'A^{-1}b}{c'A^{-1}c}A^{-1}c\). As \(\mathrm H\psi=A\) is PD,
\(\tilde x\) is the strict minimizer under the constraint (Theorem~\ref{r13-thm-extrema}(e)).
\end{example}

\begin{proposition}[\citeat{magnus-neudecker}{Example~7.3}]\label{r13-prop-sphere}
Let \(A\in\R^{n\times n}\) (\(n\geqslant2\)) be symmetric, \(f(x)=x'Ax\) and \(g(x)=x'x-1\) on \(\R^n\). A vector \(x\) is a stationary
point of \(f\) under the constraint \(g(x)=0\) (on the unit sphere) if and only if \(x\) is a unit eigenvector of \(A\); then
\(f(x)\) is its eigenvalue. The maximum \(\lambda_1\) and the minimum \(\lambda_n\) of \(f\) on the unit sphere (the largest and the
smallest eigenvalue of \(A\)) are attained exactly at the unit eigenvectors belonging to \(\lambda_1\) and to \(\lambda_n\).
\end{proposition}
\begin{proof}
\begin{itemize}
\item By Theorem~\ref{r13-thm-forms}(c), \(\mathrm Dg(x)=2x'\) and \(\mathrm Df(x)=2x'A\), as \(I_n\) and \(A\) are symmetric. For
\(x'x=1\), \(x\ne0\) and \(\rank\mathrm Dg(x)=1\).
\item So a unit vector \(x\) is a stationary point of \(f\) under the constraint if and only if \(2x'A=l\cdot2x'\) for some
\(l\in\R\), that is (transposing, \(A'=A\)), \(Ax=lx\): \(x\) is an eigenvector belonging to \(l\). Then
\(f(x)=x'(lx)=l\,x'x=l\).
\item On the unit sphere \(f(x)=x'Ax/x'x\), so the last statement is Theorem~\ref{r9-thm-rayleigh}(a),(b).\qedhere
\end{itemize}
\end{proof}
By Lagrange's theorem, the points where \(f\) is largest or smallest on the sphere must be stationary points under the
constraint, and the proposition finds them all; Theorem~\ref{r9-thm-rayleigh} gives the same answer without derivatives.

\begin{theorem}[Quadratic minimization under linear constraints; \citeat{harville}{Theorem~19.2.1, \S19.3a}]\label{r9-thm-equality-qp}
Let \(A\in\R^{p\times p}\) be PD, \(b\in\R^p\), \(C\in\R^{m\times p}\) of rank \(m\) and \(d\in\R^m\). Let \(q(x)=\frac12x'Ax-b'x\) and
\(\hat x=A^{-1}b\), its unconstrained minimizer (Example~\ref{r13-exm-quadratic-min}(a)). Then \(V=CA^{-1}C'\) is PD, hence
invertible, and with \(\lambda=V^{-1}(C\hat x-d)\):
\begin{enumerate}
\item the unique minimizer of \(q\) on \(\{x\in\R^p:Cx=d\}\) is
\[
x^*=\hat x-A^{-1}C'V^{-1}(C\hat x-d)=\hat x-A^{-1}C'\lambda
\]
\item \(q(x)=q(x^*)+\frac12(x-x^*)'A(x-x^*)\) for every \(x\) with \(Cx=d\);
\item \(q(x^*)-q(\hat x)=\frac12(C\hat x-d)'V^{-1}(C\hat x-d)\);
\item \((x^*,\lambda)\) is the unique solution of the \((p+m)\times(p+m)\) linear system
\[
\begin{pmatrix}A&C'\\C&0\end{pmatrix}\begin{pmatrix}x\\\lambda\end{pmatrix}=\begin{pmatrix}b\\d\end{pmatrix}
\]
Its first block row says \(\mathrm Dq(x^*)=(Ax^*-b)'=(-\lambda)'C\): \(x^*\) is a stationary point of \(q\) under the constraint
\(Cx-d=0\) (the paragraph after Theorem~\ref{r13-thm-extrema}; \(\mathrm D(Cx-d)=C\)), with multiplier vector \(-\lambda\).
\end{enumerate}
\end{theorem}
\begin{proof}
\(A^{-1}\) is PD (Proposition~\ref{r8-prop-pd-basic}(c)) and \(C'\in\R^{p\times m}\) has rank \(m\)
(Theorem~\ref{thm:r4-rank-rules}(a)), so \(V=(C')'A^{-1}C'\) is PD (Theorem~\ref{r8-thm-congruence}(b)), hence invertible
(Proposition~\ref{r8-prop-pd-basic}(c)).

(d) Let \(K\) be the block matrix. The Schur complement of \(A\) in \(K\) is \(0-CA^{-1}C'=-V\), so
\(\det K=\det A\cdot\det(-V)=(-1)^m\det A\det V\ne0\) (Theorems~\ref{r5d-thm-schur-det}(a) and~\ref{r5d-thm-product}(c)), and
\(K\) is invertible (Theorem~\ref{r5d-thm-product}(a)): the system has exactly one solution. The pair \((x^*,\lambda)\) is a
solution: \(Ax^*+C'\lambda=A\hat x-C'\lambda+C'\lambda=b\) and \(Cx^*=C\hat x-V\lambda=C\hat x-(C\hat x-d)=d\). By
Theorem~\ref{r13-thm-forms}(a),(c), \(\mathrm Dq(x)=x'A-b'=(Ax-b)'\).

(b) For \(h\in\R^p\), as \(A\) is symmetric,
\[
q(x^*+h)=\tfrac12(x^*)'Ax^*+h'Ax^*+\tfrac12h'Ah-b'x^*-b'h=q(x^*)+h'(Ax^*-b)+\tfrac12h'Ah
\]
By (d), \(Ax^*-b=-C'\lambda\). If \(Cx=d\), then \(h=x-x^*\) satisfies \(Ch=d-d=0\), so \(h'(Ax^*-b)=-(Ch)'\lambda=0\), and
\(q(x)=q(x^*)+\frac12h'Ah\).

(a) \(Cx^*=d\) by (d). By (b), \(q(x)\geqslant q(x^*)\) whenever \(Cx=d\), with equality only if \((x-x^*)'A(x-x^*)=0\), that is,
\(x=x^*\), as \(A\) is PD.

(c) By Example~\ref{r13-exm-quadratic-min}(a), \(2q(x)+b'A^{-1}b=(x-\hat x)'A(x-\hat x)\) for every \(x\), and
\(2q(\hat x)+b'A^{-1}b=0\); so \(q(x^*)=q(\hat x)+\frac12(x^*-\hat x)'A(x^*-\hat x)\), and \(x^*-\hat x=-A^{-1}C'\lambda\). So
\(q(x^*)-q(\hat x)=\frac12\lambda'CA^{-1}AA^{-1}C'\lambda=\frac12\lambda'V\lambda=\frac12(C\hat x-d)'V^{-1}(C\hat x-d)\), as \(V^{-1}\)
is symmetric.
\end{proof}
For \(A=I_p\) and \(b=a\), \(q(x)=\frac12\lVert x-a\rVert^2-\frac12\lVert a\rVert^2\), and (a) gives the point of the set
\(\{x:Cx=d\}\) closest to \(a\), namely \(x^*=a-C'(CC')^{-1}(Ca-d)\). For \(a=0\) this is \(C'(CC')^{-1}d\), the shortest solution of
\(Cx=d\), as in Theorem~\ref{r10-thm-min-norm} with the minimum-norm g-inverse \(G=C'(CC')^{-1}\) of \(C\) (\(CGC=C\), and
\(GC\) is symmetric).

\begin{example}\label{r9-exm-equality-qp}
Minimize
\[
q(x)=x_1^2+x_2^2+x_3^2+x_4^2-x_1x_2-x_2x_3-x_3x_4-2(x_1+x_2+x_3+x_4)
\]
subject to \(x_1+x_2+x_3+x_4=10\) and \(x_1-x_2+x_3-x_4=6\).

\textit{Solution.} Here \(q(x)=\frac12x'Ax-b'x\) with \(b=2\one_4\) and
\[
A=\begin{pmatrix}2&-1&0&0\\-1&2&-1&0\\0&-1&2&-1\\0&0&-1&2\end{pmatrix},\qquad
A^{-1}=\frac15\begin{pmatrix}4&3&2&1\\3&6&4&2\\2&4&6&3\\1&2&3&4\end{pmatrix}
\]
(multiply to check \(AA^{-1}=I_4\), Theorem~\ref{r2-thm-invertible}). The leading principal minors of \(A\) are \(2,3,4,5\), so
\(A\) is PD (Theorem~\ref{r8-thm-leading-minors}). The constraints are \(Cx=d\) with
\[
C=\begin{pmatrix}1&1&1&1\\1&-1&1&-1\end{pmatrix},\qquad d=\begin{pmatrix}10\\6\end{pmatrix}
\]
and \(\rank C=2\) (the rows are not multiples of each other). The unconstrained minimizer is
\(\hat x=A^{-1}b=\frac25(10,15,15,10)'=(4,6,6,4)'\), with \(q(\hat x)=-\frac12b'\hat x=-20\)
(Example~\ref{r13-exm-quadratic-min}(a)). Next
\[
A^{-1}C'=\frac15\begin{pmatrix}10&2\\15&-1\\15&1\\10&-2\end{pmatrix},\qquad
V=CA^{-1}C'=\begin{pmatrix}10&0\\0&\frac65\end{pmatrix},\qquad
C\hat x-d=\begin{pmatrix}20\\0\end{pmatrix}-\begin{pmatrix}10\\6\end{pmatrix}=\begin{pmatrix}10\\-6\end{pmatrix}
\]
so \(\lambda=V^{-1}(C\hat x-d)=(1,-5)'\). By Theorem~\ref{r9-thm-equality-qp}(a),(c),
\[
x^*=\begin{pmatrix}4\\6\\6\\4\end{pmatrix}-\frac15\begin{pmatrix}10&2\\15&-1\\15&1\\10&-2\end{pmatrix}\begin{pmatrix}1\\-5\end{pmatrix}
=\begin{pmatrix}4\\6\\6\\4\end{pmatrix}-\begin{pmatrix}0\\4\\2\\4\end{pmatrix}=\begin{pmatrix}4\\2\\4\\0\end{pmatrix},\qquad
q(x^*)-q(\hat x)=\tfrac12\bigl(10\cdot1+(-6)(-5)\bigr)=20
\]
so the constrained minimum is \(q(x^*)=0\). Checks: \(x^*\) satisfies both constraints; \(q(x^*)=16+4+16-8-8-20=0\); and
\(Ax^*-b=(6,-4,6,-4)'-(2,2,2,2)'=(4,-6,4,-6)'=-C'\lambda\), as in Theorem~\ref{r9-thm-equality-qp}(d).
\end{example}

\needspace{16\baselineskip}
\subsection*{Least squares by differentiation}

\begin{theorem}[\citeat{tu-clm}{Theorem~2.4}, \citeat{searle-khuri}{Example~9.3}, \citeat{abadir-magnus}{Exercise~13.54}]\label{r13-thm-ols}
Let \(X\in\R^{n\times(k+1)}\), \(y\in\R^n\), and \(S(\beta)=(y-X\beta)'(y-X\beta)\) for \(\beta\in\R^{k+1}\).
\begin{enumerate}
\item \(\mathrm{d}S=-2(y-X\beta)'X\,\mathrm{d}\beta\), so \(\mathrm DS(\beta)=-2(y-X\beta)'X\) and the gradient is \(2X'X\beta-2X'y\); and
\(\mathrm HS(\beta)=2X'X\), which is PSD.
\item \(\hat\beta\) minimizes \(S\) if and only if \(X'X\hat\beta=X'y\) (the normal equations).
\item If \(\rank X=k+1\), then \(\mathrm HS\) is PD and \(\hat\beta=(X'X)^{-1}X'y\) is the strict global minimizer.
\end{enumerate}
\end{theorem}
\begin{proof}
(a)
\begin{itemize}
\item With \(e=y-X\beta\), \(\mathrm{d}e=-X\,\mathrm{d}\beta\), and by the product rule \(\mathrm{d}S=(\mathrm{d}e)'e+e'\,\mathrm{d}e=2e'\,\mathrm{d}e=-2e'X\,\mathrm{d}\beta\)
(Theorems~\ref{r13-thm-diff-rules} and~\ref{r13-thm-forms}); so \(\mathrm DS(\beta)=-2(y-X\beta)'X\) (Theorem~\ref{r13-thm-identification}(a)).
Equivalently, \(S=y'y-2\beta'X'y+\beta'X'X\beta\) has the derivative \(-2y'X+2\beta'X'X\) (Theorem~\ref{r13-thm-forms}).
\item With \(\mathrm{d}\beta\) fixed, \(\mathrm{d}^2S=-2(\mathrm{d}e)'X\,\mathrm{d}\beta=2(\mathrm{d}\beta)'X'X\,\mathrm{d}\beta\), so \(\mathrm HS(\beta)=2X'X\)
(Theorem~\ref{r13-thm-identification2}(a)), which is PSD (Theorem~\ref{r8-thm-XtX}).
\end{itemize}
(b) If \(\hat\beta\) minimizes \(S\) on the open set \(\R^{k+1}\), then \(\mathrm DS(\hat\beta)=0\) (Theorem~\ref{r13-thm-extrema}(a)), that is,
\(X'(y-X\hat\beta)=0\). Conversely, if \(X'X\hat\beta=X'y\), then \(\hat\beta\) is a stationary point, and as \(\R^{k+1}\) is convex and
\(\mathrm HS\) is PSD everywhere, \(\hat\beta\) is a global minimizer (Theorem~\ref{r13-thm-extrema}(d)).

(c) If \(\rank X=k+1\), \(X'X\) is PD (Theorem~\ref{r8-thm-XtX}) and invertible, the normal equations have the single solution
\((X'X)^{-1}X'y\) (Theorem~\ref{r10-thm-normal-equations}(b)), and the minimum is strict (Theorem~\ref{r13-thm-extrema}(d)).
\end{proof}
The projection route of Chapter~13 (least squares needs no calculus: the paragraph after
Theorem~\ref{r10-thm-least-squares}) gives the same normal equations: by Theorem~\ref{r10-thm-least-squares}(a),
\(S(\beta)=S(\hat\beta)+\lVert X(\hat\beta-\beta)\rVert^2\). This identity is Lemma~\ref{r13-lem-line}(b) written out for the
quadratic function \(S\): with \(u=\beta-\hat\beta\), \(\mathrm DS(\hat\beta)u=0\) and \(\frac12u'(2X'X)u=\lVert Xu\rVert^2\), with no remainder because
the Hessian \(2X'X\) does not depend on \(\beta\). Differentiation finds the normal equations as the condition for a stationary point;
the projection explains them as \(y-X\hat\beta\perp\mathcal C(X)\) (Proposition~\ref{r10-prop-residuals}(a)).

\begin{example}\label{r13-exm-ols}
For the design \(X=[\one_4\ X_0]\in\R^{4\times3}\) (\(k=2\)) and \(y=(9,3,5,3)'\) of Example~\ref{r10-exm-least-squares},
\[
X'X=\begin{pmatrix}4&4&8\\4&8&8\\8&8&20\end{pmatrix},\qquad X'y=\begin{pmatrix}20\\28\\44\end{pmatrix}
\]
(computed there). At \(\beta=0\) the gradient \(2X'X\beta-2X'y\) is \(-2(20,28,44)'\ne0\). At
\(\hat\beta=(1,2,1)'\), \(X'X\hat\beta=(4+8+8,\ 4+16+8,\ 8+16+20)'=X'y\), so the gradient is \(0\). The leading principal minors of
\(X'X\) are \(4\), \(4\cdot8-4\cdot4=16\) and \(\det(X'X)=64\), all positive, so \(\mathrm HS=2X'X\) is PD
(Theorem~\ref{r8-thm-leading-minors}) and \(\hat\beta=(1,2,1)'\) is the strict global minimizer, with
\(S(\hat\beta)=4\) (Theorem~\ref{r13-thm-ols}(c)).
\end{example}

\subsection*{Further statistical applications}

\begin{example}[Hazel's selection index; \citeat{searle-khuri}{Example~9.4}, \citeat{bapat}{Result~5.3}]\label{r13-exm-hazel}
An animal is measured on \(p\) traits, \(x\in\R^p\) (the phenotypes), with \(x=g+e\), where \(g\) holds the unobserved genetic values
and \(e\) environmental effects uncorrelated with \(g\). Let \(G=\var(g)\), \(P=\var(x)\) (PD), and let \(a\in\R^p\)
weigh the traits by their economic value, with \(a'Ga>0\). The breeder ranks animals by an \textbf{index} \(I=b'x\) and wants it to be as
highly correlated as possible with the genetic worth \(H=a'g\).
\begin{itemize}
\item By Theorem~\ref{r11-thm-linear}(b) for \(z=(g',e')'\), with \(\var(z)=\diag(G,E)\), \(E=\var(e)\):
\(P=G+E\), and the \(2\times2\) covariance matrix of \(I=(b',b')z\) and \(H=(a',0)z\) has the off-diagonal entry
\(\cov(I,H)=b'Ga\). So the correlation is \(r(b)=\dfrac{b'Ga}{\sqrt{b'Pb}\,\sqrt{a'Ga}}\) for \(b\ne0\).
\item The derivative of \(\varphi(b)=b'Ga/\sqrt{b'Pb}\) is found from
\(\mathrm{d}\varphi=\frac{a'G\,\mathrm{d}b}{(b'Pb)^{1/2}}-\frac{(b'Ga)\,b'P\,\mathrm{d}b}{(b'Pb)^{3/2}}\) (Theorem~\ref{r13-thm-forms}, chain rule);
it is \(0\) if and only if \((b'Pb)\,Ga=(b'Ga)\,Pb\). So at a stationary point with \(b'Ga\ne0\), \(b\) is a multiple of \(P^{-1}Ga\), and
\(r\) does not change when \(b\) is multiplied by a positive number.
\item The stationary condition alone does not show a maximum. By Theorem~\ref{r9-thm-cs-max} with \(B=P\) and the vector \(Ga\),
\((b'Ga)^2\leqslant(b'Pb)(a'GP^{-1}Ga)\), with equality at \(b=P^{-1}Ga\). Hence \(r(b)\leqslant\sqrt{a'GP^{-1}Ga/a'Ga}\) for every
\(b\ne0\), and the bound is attained at \(b=P^{-1}Ga\), where \(b'Ga=a'GP^{-1}Ga\) and \(r=\sqrt{b'Ga/a'Ga}\).
\end{itemize}
For four traits with \(E=I_4\),
\[
G=\begin{pmatrix}2&1&1&0\\1&2&0&1\\1&0&2&1\\0&1&1&2\end{pmatrix},\qquad
P=G+I_4=\begin{pmatrix}3&1&1&0\\1&3&0&1\\1&0&3&1\\0&1&1&3\end{pmatrix},\qquad
a=\begin{pmatrix}3\\1\\1\\0\end{pmatrix}
\]
\begin{itemize}
\item \(G\) is PSD, as \(x'Gx=(x_1+x_2)^2+(x_1+x_3)^2+(x_2+x_4)^2+(x_3+x_4)^2\geqslant0\), so \(P=G+I_4\) is PD
(Proposition~\ref{r8-prop-pd-basic}(b)); and \(a'Ga=4^2+4^2+1^2+1^2=34>0\).
\item \(a=Pe_1\) is the first column of \(P\), and \(GP=G^2+G=PG\). So \(Ga=GPe_1=PGe_1\), and
\[
b=P^{-1}Ga=Ge_1=\begin{pmatrix}2\\1\\1\\0\end{pmatrix},\qquad
Ga=\begin{pmatrix}8\\5\\5\\2\end{pmatrix}=Pb
\]
\item Then \(b'Ga=16+5+5+0=26\), and the best index \(I=2x_1+x_2+x_3\) has \(r=\sqrt{26/34}=\sqrt{13/17}\).
\item The index \(I=a'x\) has \(a'Pa=a'Ga+a'a=34+11=45\) and only
\(r=34/\sqrt{45\cdot34}=\sqrt{34/45}\), less than \(\sqrt{13/17}\) as \(34\cdot17=578<585=13\cdot45\).
\end{itemize}
\end{example}

\begin{example}[Newton--Raphson; \citeat{searle-khuri}{\S9.2}]\label{r13-exm-newton}
To solve \(n\) nonlinear equations \(f(x)=0\) in \(n\) unknowns, with \(f\) \(C^1\) and derivative \(G(x)=\mathrm Df(x)\), one uses the
first-order approximation \(f(x_r+\Delta)\approx f(x_r)+G(x_r)\Delta\), the differential at the current guess \(x_r\). Setting the right side
equal to \(0\) gives the \textbf{Newton--Raphson} step
\[
x_{r+1}=x_r-G(x_r)^{-1}f(x_r)\qquad(r=0,1,2,\dots)
\]
provided \(G(x_r)\) is invertible; one stops when the step is negligible. To find a stationary point of a \(C^2\) scalar function
\(\varphi\), apply it to the gradient: \(f=(\mathrm D\varphi)'\), \(G=\mathrm H\varphi\), so \(x_{r+1}=x_r-\mathrm H\varphi(x_r)^{-1}(\mathrm D\varphi(x_r))'\).
If \(\varphi(x)=\frac12x'Ax-b'x\) with \(A\) symmetric and invertible, the gradient \(Ax-b\) is linear, the approximation is exact, and one
step from any \(x_0\) gives \(x_1=x_0-A^{-1}(Ax_0-b)=A^{-1}b\), the stationary point. For least squares with \(\rank X=k+1\)
(Theorem~\ref{r13-thm-ols}(c)), \(x_1=\beta_0-(2X'X)^{-1}(2X'X\beta_0-2X'y)=(X'X)^{-1}X'y\): in Example~\ref{r13-exm-ols}, one
step from \(\beta_0=0\) gives \(\hat\beta=(1,2,1)'\). For other functions several steps are needed.
\end{example}

For the stationary points of a \(C^2\) scalar function \(\varphi\), the step has a second reading. Let
\[
q_r(y)=\varphi(x_r)+\mathrm D\varphi(x_r)(y-x_r)+\tfrac12(y-x_r)'\mathrm H\varphi(x_r)(y-x_r)\qquad(y\in\R^n)
\]
the second-order Taylor polynomial of \(\varphi\) at \(x_r\). By Theorem~\ref{r13-thm-forms} and the symmetry of
\(\mathrm H\varphi(x_r)\), the gradient of \(q_r\) is \((\mathrm D\varphi(x_r))'+\mathrm H\varphi(x_r)(y-x_r)\) and its Hessian is
\(\mathrm H\varphi(x_r)\): \(q_r\) has the same value, gradient and Hessian as \(\varphi\) at \(x_r\), and its gradient is \(0\) exactly
at \(y=x_{r+1}\). So the Newton--Raphson step jumps to the stationary point of this quadratic approximation of \(\varphi\). If
\(\mathrm H\varphi(x_r)\) is PD, \(x_{r+1}\) is the unique minimizer of \(q_r\): with \(z=y-x_r\),
\(q_r(y)=\varphi(x_r)+\frac12z'\mathrm H\varphi(x_r)z-\bigl(-(\mathrm D\varphi(x_r))'\bigr)'z\), and
Example~\ref{r13-exm-quadratic-min}(a) gives the unique minimizer \(z=-\mathrm H\varphi(x_r)^{-1}(\mathrm D\varphi(x_r))'=x_{r+1}-x_r\);
the constant \(\varphi(x_r)\) does not change it.

\needspace{16\baselineskip}
\begin{theorem}[\citeat{searle-khuri}{\S9.5.5}, \citeat{abadir-magnus}{Exercises~13.37(a), 13.38(a)}]\label{r13-thm-jacobian}
\begin{enumerate}
\item If \(h\) is a one-to-one \(C^1\) map of an open \(U\subseteq\R^n\) onto an open \(W\) with a \(C^1\) inverse, then
\(\mathrm D(h^{-1})(h(x))\,\mathrm Dh(x)=I_n\); so \(\lvert\det\mathrm D(h^{-1})(h(x))\rvert=1/\lvert\det\mathrm Dh(x)\rvert\). (The
determinant \(\det\mathrm Dh(x)\) is the \textbf{Jacobian determinant} of \(h\) at \(x\).)
\item Let \(A\in\R^{n\times n}\) be invertible and \(b\in\R^n\). If \(x\) has the density \(f_x\), then \(y=Ax+b\) has the density
\[
f_y(y)=\frac{f_x\bigl(A^{-1}(y-b)\bigr)}{\lvert\det A\rvert}
\]
\item Let \(A\in\R^{m\times m}\) and \(B\in\R^{n\times n}\) be invertible. The map \(X\mapsto Y=AXB\) of \(\R^{m\times n}\) onto itself reads
\(\operatorname{vec}Y=(B'\otimes A)\operatorname{vec}X\), and \(\det(B'\otimes A)=(\det A)^n(\det B)^m\). So if the random matrix \(X\) has the
density \(f_X\) (of \(\operatorname{vec}X\)), then \(Y=AXB\) has the density
\(f_Y(Y)=f_X(A^{-1}YB^{-1})\,\lvert\det A\rvert^{-n}\lvert\det B\rvert^{-m}\).
\item On the open set of invertible \(X\in\R^{n\times n}\), the map \(X\mapsto X^{-1}\) has the Jacobian determinant
\[
\det\mathrm D(X^{-1})=(-1)^n(\det X)^{-2n}
\]
\end{enumerate}
\end{theorem}
\begin{proof}
(a) \(h^{-1}(h(x))=x\) on \(U\); by the chain rule of calculus, \(\mathrm D(h^{-1})(h(x))\,\mathrm Dh(x)=I_n\), and taking determinants
(Theorem~\ref{r5d-thm-product}(b)) gives the second statement.

(b) \(h(x)=Ax+b\) is one-to-one from \(\R^n\) onto \(\R^n\), with \(\mathrm Dh(x)=A\) (Theorem~\ref{r13-thm-forms}(b)) and the \(C^1\) inverse
\(h^{-1}(y)=A^{-1}(y-b)\). By the change of variables (D1), stated before Lemma~\ref{r15-lem-mgf}, with \(\psi=h^{-1}\),
whose Jacobian determinant is \(\det\mathrm D(h^{-1})(y)\), and (a), \(f_y(y)=f_x(A^{-1}(y-b))/\lvert\det A\rvert\).

(c) By Theorem~\ref{r13-thm-vec}(b), \(\operatorname{vec}(AXB)=(B'\otimes A)\operatorname{vec}X\). \(B'\otimes A\) is invertible
(Theorem~\ref{r13-thm-kron-det}(c)), and \(\det(B'\otimes A)=(\det B')^m(\det A)^n=(\det A)^n(\det B)^m\)
(Theorems~\ref{r13-thm-kron-det}(a) and~\ref{r5d-thm-transpose}). Apply (b) to the vector \(\operatorname{vec}X\), with
\((B'\otimes A)^{-1}\operatorname{vec}Y=\operatorname{vec}(A^{-1}YB^{-1})\) (Corollary~\ref{r13-cor-kron}(b), Theorem~\ref{r13-thm-vec}(b)).

(d) By Theorem~\ref{r13-thm-trace-det}(c), \(\mathrm D(X^{-1})=-(X^{-1})'\otimes X^{-1}\), an \(n^2\times n^2\) matrix. By
Theorem~\ref{r5d-thm-product}(c) (\(\det(cM)=c^{n^2}\det M\) for \(M\in\R^{n^2\times n^2}\)), Theorem~\ref{r13-thm-kron-det}(a)
and Theorems~\ref{r5d-thm-transpose} and~\ref{r5d-thm-product}(c),
\[
\det\mathrm D(X^{-1})=(-1)^{n^2}\bigl(\det(X^{-1})'\bigr)^n(\det X^{-1})^n=(-1)^{n^2}(\det X)^{-2n}=(-1)^n(\det X)^{-2n}
\]
as \(n^2-n=n(n-1)\) is even.
\end{proof}

\begin{example}[Aitken's integral; \citeat{searle-khuri}{\S9.5.6}]\label{r13-exm-aitken}
Let \(A\in\R^{n\times n}\) be PD. Then
\[
\int_{\R^n}e^{-\frac12x'Ax}\,\mathrm{d}x=(2\pi)^{n/2}(\det A)^{-1/2}
\]
\begin{itemize}
\item Let \(P=A^{-1/2}\) (Theorem~\ref{r8-thm-sqrt}): \(P\) is symmetric, PD as the inverse of the PD matrix \(A^{1/2}\)
(Proposition~\ref{r8-prop-pd-basic}(c)), \(P'AP=A^{-1/2}AA^{-1/2}=I_n\), and
\((\det P)^2=\det(P^2)=\det A^{-1}=1/\det A\) (Theorem~\ref{r5d-thm-product}(b),(c)); as \(\det P>0\)
(Theorem~\ref{r8-thm-eigen-criterion}(c)), \(\det P=(\det A)^{-1/2}\).
\item Change variables by \(x=Py\) (the change of variables stated before Theorem~\ref{r14-thm-mvn-density}, with \(c=0\) and
\(B=P\)): the integral is
\[
\begin{aligned}
\int e^{-\frac12y'P'APy}\det P\,\mathrm{d}y&=(\det A)^{-1/2}\int e^{-\frac12(y_1^2+\dots+y_n^2)}\,\mathrm{d}y\\
&=(\det A)^{-1/2}\prod_{i=1}^n\int_\R e^{-y_i^2/2}\,\mathrm{d}y_i=(\det A)^{-1/2}(2\pi)^{n/2}
\end{aligned}
\]
\end{itemize}
(The second line uses the rule for products and \(\int_\R e^{-s^2/2}\,\mathrm{d}s=\sqrt{2\pi}\), stated at the same place.) With
\(A=\Sigma^{-1}\) and the further change \(x\mapsto x-\mu\), this is the computation in the proof of
Theorem~\ref{r14-thm-mvn-density} that the normal density integrates to \(1\).
\end{example}

\begin{example}[\citeat{tu-mc}{\S13}, \citeat{abadir-magnus}{Exercise~13.61}]\label{r13-exm-covariance}
Let \(S\in\R^{p\times p}\) be PD (for instance a sample covariance matrix, Definition~\ref{r12-def-data}), and minimize
\[
f(\Sigma)=\log\det\Sigma+\tr(\Sigma^{-1}S)
\]
over the PD matrices \(\Sigma\in\R^{p\times p}\), with the variables \(\operatorname{vech}\Sigma\). The PD matrices form an open set (they are the
symmetric matrices whose leading principal minors, polynomials in the entries, are positive, Theorem~\ref{r8-thm-leading-minors}) and a
convex set (Proposition~\ref{r8-prop-pd-basic}(b)); \(\det\Sigma>0\) on it, and \(f\) is \(C^2\) (Theorem~\ref{r13-thm-trace-det}).

(a) By Theorems~\ref{r13-thm-diff-rules} and~\ref{r13-thm-trace-det}(c),(d) and Theorem~\ref{r1-thm-trace}(c),
\[
\mathrm{d}f=\tr(\Sigma^{-1}\mathrm{d}\Sigma)-\tr\bigl(\Sigma^{-1}(\mathrm{d}\Sigma)\Sigma^{-1}S\bigr)=\tr(M\,\mathrm{d}\Sigma),\qquad M=\Sigma^{-1}-\Sigma^{-1}S\Sigma^{-1}
\]
At a stationary point \(\tr(M\,\mathrm{d}\Sigma)=0\) for every symmetric \(\mathrm{d}\Sigma\); as \(M\) is symmetric, \(\mathrm{d}\Sigma=M\) gives
\(\tr(M'M)=0\), so \(M=0\) (Theorem~\ref{r1-thm-trace}(d)), that is, \(\Sigma^{-1}=\Sigma^{-1}S\Sigma^{-1}\), or \(\Sigma=S\). So \(S\) is the only
stationary point.

(b) With \(\mathrm{d}\Sigma\) fixed, \(\mathrm{d}^2f=\tr\bigl((\mathrm{d}M)\,\mathrm{d}\Sigma\bigr)\), where
\(\mathrm{d}M=-\Sigma^{-1}(\mathrm{d}\Sigma)\Sigma^{-1}+\Sigma^{-1}(\mathrm{d}\Sigma)\Sigma^{-1}S\Sigma^{-1}+\Sigma^{-1}S\Sigma^{-1}(\mathrm{d}\Sigma)\Sigma^{-1}\); the
first term of \(\mathrm{d}M\) contributes \(-\tr\bigl(\Sigma^{-1}(\mathrm{d}\Sigma)\Sigma^{-1}\mathrm{d}\Sigma\bigr)\), the second differential of
\(\log\det\Sigma\) (Example~\ref{r13-exm-hessians}(c)). At \(\Sigma=S\), \(\mathrm{d}M=S^{-1}(\mathrm{d}\Sigma)S^{-1}\), and with \(W=S^{-1/2}(\mathrm{d}\Sigma)S^{-1/2}\) (Theorem~\ref{r8-thm-sqrt}),
\(\mathrm{d}^2f=\tr(S^{-1}\mathrm{d}\Sigma\,S^{-1}\mathrm{d}\Sigma)=\tr(W^2)=\lVert W\rVert_F^2>0\) for \(\mathrm{d}\Sigma\ne0\). So the Hessian at \(S\) (with
respect to \(\operatorname{vech}\Sigma\)) is PD (Theorem~\ref{r13-thm-identification2}(a)), and \(S\) is a strict local minimizer
(Theorem~\ref{r13-thm-extrema}(c)).

The minimum at \(S\) is global, \(f(\Sigma)>f(S)=\log\det S+p\) for every PD \(\Sigma\ne S\): this is the step for \(\Sigma\) in
the proof of Theorem~\ref{r14-thm-mle}, with \(S\) in place of \(\hat\Sigma\) (that step uses only that \(\hat\Sigma\) is PD), by
Lemma~\ref{r14-lem-logdet}, without derivatives.
\end{example}
In multivariate statistics this criterion is used to fit a covariance matrix: among all PD matrices \(\Sigma\), the one that
makes \(\log\det\Sigma+\tr(\Sigma^{-1}S)\) smallest is the matrix \(S\) that is put in, whether it is the sample covariance matrix
\(\frac1{n-1}Y'C_nY\) of Definition~\ref{r12-def-data} or the maximum-likelihood estimate \(\frac1nY'C_nY\) of
Theorem~\ref{r14-thm-mle}.

\needspace{14\baselineskip}

\subsection*{Exercises}

\begin{exercise}\label{r13-dr-kronecker}
Let \(A=\begin{pmatrix}1&2\\0&-1\end{pmatrix}\), \(B=\begin{pmatrix}2&1&0\\1&0&-1\end{pmatrix}\) and \(C=\begin{pmatrix}2&1\\1&1\end{pmatrix}\).
\begin{enumerate}
\item Compute \(A\otimes B\) and \(B\otimes A\), and find \(\rank(A\otimes B)\).
\item Find \(A^{-1}\), \(C^{-1}\) and \((A\otimes C)^{-1}\).
\item Find \(\tr(A\otimes C)\) and \(\det(A\otimes C)\).
\end{enumerate}
\end{exercise}
\begin{solution}
(a) By Definition~\ref{r13-def-kronecker},
\[
A\otimes B=\begin{pmatrix}B&2B\\0&-B\end{pmatrix}=\begin{pmatrix}2&1&0&4&2&0\\1&0&-1&2&0&-2\\0&0&0&-2&-1&0\\0&0&0&-1&0&1\end{pmatrix},\qquad
B\otimes A=\begin{pmatrix}2&4&1&2&0&0\\0&-2&0&-1&0&0\\1&2&0&0&-1&-2\\0&-1&0&0&0&1\end{pmatrix}
\]
\(A\) is invertible (\(\det A=-1\)) and the rows of \(B\) are not proportional, so \(\rank A=\rank B=2\) and
\(\rank(A\otimes B)=2\cdot2=4\) (Theorem~\ref{r13-thm-kron-det}(b)).

(b) \(A\begin{pmatrix}1&2\\0&-1\end{pmatrix}=\begin{pmatrix}1&0\\0&1\end{pmatrix}\), so \(A^{-1}=A\); and
\(C\begin{pmatrix}1&-1\\-1&2\end{pmatrix}=I_2\), so \(C^{-1}=\begin{pmatrix}1&-1\\-1&2\end{pmatrix}\) (Theorem~\ref{r2-thm-invertible}). By
Corollary~\ref{r13-cor-kron}(b),
\[
A\otimes C=\begin{pmatrix}2&1&4&2\\1&1&2&2\\0&0&-2&-1\\0&0&-1&-1\end{pmatrix},\qquad
(A\otimes C)^{-1}=A^{-1}\otimes C^{-1}=\begin{pmatrix}1&-1&2&-2\\-1&2&-2&4\\0&0&-1&1\\0&0&1&-2\end{pmatrix}
\]
(c) \(\tr(A\otimes C)=\tr A\cdot\tr C=0\cdot3=0\) (Corollary~\ref{r13-cor-kron}(c)); directly, \(2+1-2-1=0\). By
Theorem~\ref{r13-thm-kron-det}(a) with \(m=n=2\), \(\det(A\otimes C)=(\det A)^2(\det C)^2=(-1)^2\cdot1^2=1\).
\end{solution}

\begin{exercise}\label{r13-dr-kron-eigen}
Let \(A=\begin{pmatrix}2&1\\1&2\end{pmatrix}\) and
\[
B=\begin{pmatrix}1&1&0\\1&1&0\\0&0&2\end{pmatrix}
\]
\begin{enumerate}
\item Find the eigenvalues of \(A\) and of \(B\), with eigenvectors.
\item Find the eigenvalues of \(A\otimes B\) with their multiplicities, and \(\tr(A\otimes B)\), \(\det(A\otimes B)\) and \(\rank(A\otimes B)\). Is
\(A\otimes B\) PSD?
\item Check directly that \((1,1)'\otimes(0,0,1)'\) and \((1,-1)'\otimes(1,1,0)'\) are eigenvectors of \(A\otimes B\).
\end{enumerate}
\end{exercise}
\begin{solution}
(a) \(A(1,1)'=(3,3)'\) and \(A(1,-1)'=(1,-1)'\): the eigenvalues of \(A\) are \(3\) and \(1\). \(B(1,1,0)'=(2,2,0)'\), \(B(1,-1,0)'=0\) and
\(B(0,0,1)'=(0,0,2)'\): the eigenvalues of \(B\) are \(2\), \(0\), \(2\). The eigenvectors found are orthogonal, so (after normalizing)
they give \(A=P\Lambda P'\), \(B=QMQ'\) as in Theorem~\ref{r8-thm-spectral}.

(b) By Theorem~\ref{r13-thm-kron-eigen}(b),(c), the eigenvalues of \(A\otimes B\) are the products \(3\cdot2,3\cdot0,3\cdot2,1\cdot2,1\cdot0,1\cdot2\):
\(6\) (twice), \(2\) (twice) and \(0\) (twice). Hence \(\tr(A\otimes B)=6+6+2+2=16=\tr A\cdot\tr B=4\cdot4\),
\(\det(A\otimes B)=0=(\det A)^3(\det B)^2\), and, \(A\otimes B\) being symmetric, \(\rank(A\otimes B)=4\), the number of nonzero eigenvalues
(Theorem~\ref{r8-thm-rank-eigen}), which is \(\rank A\cdot\rank B=2\cdot2\). All eigenvalues are \(\geqslant0\), so \(A\otimes B\) is PSD
(Theorem~\ref{r8-thm-eigen-criterion}(a)), but not PD.

(c)
\[
A\otimes B=\begin{pmatrix}2&2&0&1&1&0\\2&2&0&1&1&0\\0&0&4&0&0&2\\1&1&0&2&2&0\\1&1&0&2&2&0\\0&0&2&0&0&4\end{pmatrix}
\]
With \(u=(1,1)'\otimes(0,0,1)'=(0,0,1,0,0,1)'\), \((A\otimes B)u=(0,0,6,0,0,6)'=6u\); with \(w=(1,-1)'\otimes(1,1,0)'=(1,1,0,-1,-1,0)'\),
\((A\otimes B)w=(2,2,0,-2,-2,0)'=2w\), as Theorem~\ref{r13-thm-kron-eigen}(a) predicts (\(3\cdot2\) and \(1\cdot2\)).
\end{solution}

\begin{exercise}\label{r13-dr-vec}
Let \(A=\begin{pmatrix}1&2\\0&1\end{pmatrix}\), \(B=\begin{pmatrix}1&0\\2&1\end{pmatrix}\) and \(C=\begin{pmatrix}1&1\\0&2\end{pmatrix}\).
\begin{enumerate}
\item Compute \(ABC\) and \(\operatorname{vec}(ABC)\).
\item Compute \(C'\otimes A\) and check that \((C'\otimes A)\operatorname{vec}B=\operatorname{vec}(ABC)\).
\item Compute \(\tr(A'B)\) as \((\operatorname{vec}A)'\operatorname{vec}B\), and check it directly.
\end{enumerate}
\end{exercise}
\begin{solution}
(a) \(AB=\begin{pmatrix}5&2\\2&1\end{pmatrix}\) and \(ABC=\begin{pmatrix}5&9\\2&4\end{pmatrix}\), so \(\operatorname{vec}(ABC)=(5,2,9,4)'\)
(Definition~\ref{r13-def-vec}).

(b) \(C'=\begin{pmatrix}1&0\\1&2\end{pmatrix}\), so
\[
C'\otimes A=\begin{pmatrix}A&0\\A&2A\end{pmatrix}=\begin{pmatrix}1&2&0&0\\0&1&0&0\\1&2&2&4\\0&1&0&2\end{pmatrix},\qquad
(C'\otimes A)\begin{pmatrix}1\\2\\0\\1\end{pmatrix}=\begin{pmatrix}1+4\\2\\1+4+0+4\\2+2\end{pmatrix}=\begin{pmatrix}5\\2\\9\\4\end{pmatrix}
\]
with \(\operatorname{vec}B=(1,2,0,1)'\), as Theorem~\ref{r13-thm-vec}(b) says.

(c) \(\operatorname{vec}A=(1,0,2,1)'\), so \((\operatorname{vec}A)'\operatorname{vec}B=1+0+0+1=2\). Directly,
\(A'B=\begin{pmatrix}1&0\\2&1\end{pmatrix}\begin{pmatrix}1&0\\2&1\end{pmatrix}=\begin{pmatrix}1&0\\4&1\end{pmatrix}\), with trace \(2\)
(Theorem~\ref{r13-thm-vec}(d)).
\end{solution}

\begin{exercise}\label{r13-dr-commutation}
\begin{enumerate}
\item Write down \(K_{32}\) and check \(K_{32}\operatorname{vec}A=\operatorname{vec}A'\) for
\[
A=\begin{pmatrix}1&4\\2&5\\3&6\end{pmatrix}
\]
\item Check that \(K_{32}=K_{23}'\), with \(K_{23}\) of Definition~\ref{r13-def-commutation}.
\item For \(a=(1,2)'\) and \(b=(3,0,1)'\), compute \(a\otimes b\), \(b\otimes a\) and \(K_{32}(a\otimes b)\).
\end{enumerate}
\end{exercise}
\begin{solution}
(a) By Theorem~\ref{r13-thm-commutation}(a) (\(m=3\), \(n=2\)), row \((i-1)2+j\) of \(K_{32}\) is row \((j-1)3+i\) of \(I_6\): rows \(1,2,\dots,6\)
of \(K_{32}\) are rows \(1,4,2,5,3,6\) of \(I_6\),
\[
K_{32}=\begin{pmatrix}1&0&0&0&0&0\\0&0&0&1&0&0\\0&1&0&0&0&0\\0&0&0&0&1&0\\0&0&1&0&0&0\\0&0&0&0&0&1\end{pmatrix}
\]
For the given \(A\), \(\operatorname{vec}A=(1,2,3,4,5,6)'\) and \(K_{32}\operatorname{vec}A=(1,4,2,5,3,6)'=\operatorname{vec}A'\), as
\(A'=\begin{pmatrix}1&2&3\\4&5&6\end{pmatrix}\).

(b) The rows of \(K_{23}\) are rows \(1,3,5,2,4,6\) of \(I_6\); transposing, the columns of \(K_{23}'\) are \(e_1,e_3,e_5,e_2,e_4,e_6\), so the
\(1\) of column \(j\) of \(K_{23}'\) stands in the row listed \(j\)th. Hence row \(1\) of \(K_{23}'\) has its \(1\) in column \(1\), row \(2\) in
column \(4\), row \(3\) in column \(2\), row \(4\) in column \(5\), row \(5\) in column \(3\), row \(6\) in column \(6\): this is \(K_{32}\)
(Theorem~\ref{r13-thm-commutation}(b)).

(c) \(a\otimes b=(3,0,1,6,0,2)'\) and \(b\otimes a=(3,6,0,0,1,2)'\) (Definition~\ref{r13-def-kronecker}). Taking entries \(1,4,2,5,3,6\) of
\(a\otimes b\), \(K_{32}(a\otimes b)=(3,6,0,0,1,2)'=b\otimes a\) (Theorem~\ref{r13-thm-commutation}(d)).
\end{solution}

\begin{exercise}\label{r13-dr-oneway}
Three groups of two observations each give \(y=(5,7,\ 2,4,\ 6,6)'\). In the one-way model without intercept,
\(y=X_1\beta+\varepsilon\) with \(X_1=I_3\otimes\one_2\) (Example~\ref{r13-exm-designs}(a)).
\begin{enumerate}
\item Compute \(X_1'X_1\), \(X_1'y\) and the least-squares estimate \(\hat\beta\).
\item Compute the hat matrix \(H=I_3\otimes\bar J_2\), the fitted values \(Hy\), the residuals \(e=(I_3\otimes C_2)y\) and
\(\mathrm{SSE}=e'e\).
\end{enumerate}
\end{exercise}
\begin{solution}
(a) By Example~\ref{r13-exm-designs}(a), \(X_1'X_1=2I_3=\diag(2,2,2)\), and \(X_1'y=(I_3\otimes\one_2')y\) holds the three group totals,
\((12,6,12)'\). So \(\hat\beta=(X_1'X_1)^{-1}X_1'y=\frac12(12,6,12)'=(6,3,6)'\) (Theorem~\ref{r10-thm-least-squares}(c)): the group means.

(b)
\[
H=I_3\otimes\bar J_2=\frac12\begin{pmatrix}1&1&0&0&0&0\\1&1&0&0&0&0\\0&0&1&1&0&0\\0&0&1&1&0&0\\0&0&0&0&1&1\\0&0&0&0&1&1\end{pmatrix},\qquad
Hy=\begin{pmatrix}6\\6\\3\\3\\6\\6\end{pmatrix},\qquad e=y-Hy=\begin{pmatrix}-1\\1\\-1\\1\\0\\0\end{pmatrix}
\]
and \((I_3\otimes C_2)y\) gives the same \(e\) (deviations from the group means). Hence \(\mathrm{SSE}=1+1+1+1=4\).
\end{solution}

\begin{exercise}\label{r13-dr-data-cov}
Three units are measured on two variables; the rows of the random data matrix \(Y\in\R^{3\times2}\) are uncorrelated and have the same mean vector, each with
covariance matrix \(\Sigma=\begin{pmatrix}2&1\\1&3\end{pmatrix}\).
\begin{enumerate}
\item Write down \(\var(\operatorname{vec}Y)\) and \(\var(\operatorname{vec}Y')\).
\item Find \(\var(\bar y)\) and the variance of the difference \(\bar y_1-\bar y_2\) of the two variable means.
\item Find \(\var(y_{11}+y_{21})\) and \(\var(y_{11}+y_{12})\).
\end{enumerate}
\end{exercise}
\begin{solution}
(a) By Theorem~\ref{r13-thm-data-matrix},
\[
\var(\operatorname{vec}Y)=\Sigma\otimes I_3=\begin{pmatrix}2&0&0&1&0&0\\0&2&0&0&1&0\\0&0&2&0&0&1\\1&0&0&3&0&0\\0&1&0&0&3&0\\0&0&1&0&0&3\end{pmatrix},\qquad
\var(\operatorname{vec}Y')=I_3\otimes\Sigma=\begin{pmatrix}2&1&0&0&0&0\\1&3&0&0&0&0\\0&0&2&1&0&0\\0&0&1&3&0&0\\0&0&0&0&2&1\\0&0&0&0&1&3\end{pmatrix}
\]
(b) \(\var(\bar y)=\frac13\Sigma=\begin{pmatrix}2/3&1/3\\1/3&1\end{pmatrix}\) (Theorem~\ref{r14-thm-sample}(a), with the rows of \(Y\) as \(x_1,x_2,x_3\)). With \(c=(1,-1)'\),
\(\var(c'\bar y)=c'\var(\bar y)c=\frac13(2-1-1+3)=1\) (Theorem~\ref{r11-thm-linear}(b)).

(c) \(y_{11}\) and \(y_{21}\) are entries \(1\) and \(2\) of \(\operatorname{vec}Y\); with \(d=(1,1,0,0,0,0)'\),
\(\var(d'\operatorname{vec}Y)=d'(\Sigma\otimes I_3)d=2+2=4\): different units are uncorrelated. \(y_{11}\) and \(y_{12}\) are entries
\(1\) and \(2\) of \(\operatorname{vec}Y'\), so \(\var(y_{11}+y_{12})=(1,1)\Sigma(1,1)'=2+1+1+3=7\).
\end{solution}

\begin{exercise}\label{r9-ex-kronecker}
Let
\[
A=\begin{pmatrix}1&1&1\\1&0&-1\end{pmatrix}\in\R^{2\times3},\qquad
B=\begin{pmatrix}2&1\\-1&1\\1&-1\end{pmatrix}\in\R^{3\times2}
\]
\begin{enumerate}
\item Write out \(A\otimes B\in\R^{6\times6}\) and find its rank.
\item Compute \(AB\) and \(BA\), show that \((A\otimes B)(B\otimes A)=AB\otimes BA\), and find the eigenvalues and the
trace of this \(6\times6\) matrix.
\item Write the equation \(AXB=C\) (\(X\in\R^{3\times3}\), \(C\in\R^{2\times2}\)) as a linear system
\(M\operatorname{vec}X=\operatorname{vec}C\) with \(M\in\R^{4\times9}\). Show that \(\rank M=4\) and
\(MM'=B'B\otimes AA'\).
\item Show that the solution of \(AXB=C\) of least Frobenius norm is \(X^*=A'(AA')^{-1}C(B'B)^{-1}B'\). Find
\(X^*\) for \(C=\diag(18,6)\).
\end{enumerate}
\end{exercise}
\begin{solution}
(a) By Definition~\ref{r13-def-kronecker}, the blocks of \(A\otimes B\) are \(B,B,B\) in the first block row and
\(B,0,-B\) in the second:
\[
A\otimes B=\begin{pmatrix}2&1&2&1&2&1\\-1&1&-1&1&-1&1\\1&-1&1&-1&1&-1\\2&1&0&0&-2&-1\\-1&1&0&0&1&-1\\1&-1&0&0&-1&1\end{pmatrix}
\]
The two rows of \(A\) are not multiples of each other, nor are the two columns of \(B\), so \(\rank A=\rank B=2\), and
\(\rank(A\otimes B)=2\cdot2=4\) by Theorem~\ref{r13-thm-kron-det}(b). Indeed rows \(3\) and \(6\) are the negatives of
rows \(2\) and \(5\).

(b) Multiplying,
\[
AB=\begin{pmatrix}2&1\\1&2\end{pmatrix},\qquad BA=\begin{pmatrix}3&2&1\\0&-1&-2\\0&1&2\end{pmatrix}
\]
By the mixed product (Theorem~\ref{r13-thm-kron-rules}(f)), with the \(2\times3\) and \(3\times2\) factors \(A\), \(B\)
and the \(3\times2\) and \(2\times3\) factors \(B\), \(A\), \((A\otimes B)(B\otimes A)=AB\otimes BA\).
\begin{itemize}
\item \(AB(1,1)'=3(1,1)'\) and \(AB(1,-1)'=(1,-1)'\): the eigenvalues of \(AB\) are \(3\) and \(1\), and
\(P=\left(\begin{smallmatrix}1&1\\1&-1\end{smallmatrix}\right)\) gives \(P^{-1}(AB)P=\diag(3,1)\)
(Definition~\ref{r7-def-diagonalizable}).
\item Expanding along the first column (Theorem~\ref{r5d-thm-expansion}),
\(p_{BA}(t)=\det(tI_3-BA)=(t-3)\bigl((t+1)(t-2)+2\bigr)=(t-3)(t^2-t)=t(t-1)(t-3)\). The three eigenvalues \(3\), \(1\),
\(0\) are distinct, so eigenvectors belonging to them form a basis of \(\R^3\) (Theorem~\ref{r7-thm-indep}); the matrix
\(Q\) with these columns is invertible and \(Q^{-1}(BA)Q=\diag(3,1,0)\) (Definition~\ref{r7-def-diagonalizable}).
\item By Theorem~\ref{r13-thm-kron-eigen}(b), the eigenvalues of \(AB\otimes BA\), with multiplicity, are the products
\(3\cdot3,3\cdot1,3\cdot0,1\cdot3,1\cdot1,1\cdot0\), that is, \(9,3,3,1,0,0\); and the trace is
\(16=\tr(AB)\tr(BA)=4\cdot4\) (Corollary~\ref{r13-cor-kron}(c)).
\end{itemize}
(c) By Theorem~\ref{r13-thm-vec}(b), \(\operatorname{vec}(AXB)=(B'\otimes A)\operatorname{vec}X\), so
\(M=B'\otimes A\in\R^{4\times9}\), and \(\rank M=\rank B'\cdot\rank A=2\cdot2=4\) (Theorems~\ref{r13-thm-kron-det}(b)
and~\ref{thm:r4-rank-rules}(a)). By Theorem~\ref{r13-thm-kron-rules}(e),(f),
\(MM'=(B'\otimes A)(B\otimes A')=B'B\otimes AA'\). The rows of \(A\) and the columns of \(B\) are orthogonal, so
\(AA'=\diag(3,2)\), \(B'B=\diag(6,3)\) and \(MM'=\diag(18,12,9,6)\) (Corollary~\ref{r13-cor-kron}(a)).

(d)
\begin{itemize}
\item \(\lVert X\rVert_F=\lVert\operatorname{vec}X\rVert\) (Definition~\ref{r13-def-vec}). \(M\) has rank \(4\), its number of
rows, so \(\mathcal C(M)=\R^4\) (Theorem~\ref{thm:r4-rank-rules}(c)): the system \(M\operatorname{vec}X=\operatorname{vec}C\)
has solutions, and they are the minimizers of \(\lVert M\operatorname{vec}X-\operatorname{vec}C\rVert\) (minimum \(0\)). By
Theorem~\ref{r10-thm-mp-shortest}(b), the shortest of them is \(M^+\operatorname{vec}C\), and
\(M^+=M'(MM')^{-1}\) (Theorem~\ref{r10-thm-mp-exists}(c), \(\rank M=4\)).
\item By Corollary~\ref{r13-cor-kron}(b) and Theorem~\ref{r13-thm-vec}(b), as \((B'B)^{-1}\) is symmetric,
\[
(MM')^{-1}\operatorname{vec}C=\bigl((B'B)^{-1}\otimes(AA')^{-1}\bigr)\operatorname{vec}C
=\operatorname{vec}\bigl((AA')^{-1}C(B'B)^{-1}\bigr)
\]
and \(M'=B\otimes A'\) (Theorem~\ref{r13-thm-kron-rules}(e)) maps \(\operatorname{vec}Y\) to \(\operatorname{vec}(A'YB')\)
(Theorem~\ref{r13-thm-vec}(b)). Hence \(X^*=A'(AA')^{-1}C(B'B)^{-1}B'\).
\item For \(C=\diag(18,6)\), \((AA')^{-1}C(B'B)^{-1}=\diag(\frac13,\frac12)\diag(18,6)\diag(\frac16,\frac13)=I_2\), and
\[
X^*=A'B'=(BA)'=\begin{pmatrix}3&0&0\\2&-1&1\\1&-2&2\end{pmatrix},\qquad\lVert X^*\rVert_F^2=24
\]
\item Check: \(AX^*B=(AA')(B'B)=\diag(18,6)\). Other solutions are longer: \(A(1,-2,1)'=0\), so
\(X=X^*+(1,-2,1)'e_1'\) also solves \(AXB=C\), and
\[
X=\begin{pmatrix}4&0&0\\0&-1&1\\2&-2&2\end{pmatrix},\qquad\lVert X\rVert_F^2=16+1+1+4+4+4=30>24
\]
\end{itemize}
\end{solution}

\needspace{16\baselineskip}
\begin{exercise}\label{r13-dr-quadratic}
Let \(f(x)=x'Ax+b'x\) for \(x\in\R^4\), with
\[
A=\begin{pmatrix}1&2&0&0\\0&2&1&0\\0&1&3&1\\0&0&1&2\end{pmatrix},\qquad b=\begin{pmatrix}-2\\-4\\-4\\2\end{pmatrix}
\]
\begin{enumerate}
\item Find \(\mathrm{d}f\), the gradient and the Hessian of \(f\).
\item Find the stationary point by Gauss elimination.
\item Show that it is the strict global minimizer, and find the minimum value.
\end{enumerate}
\end{exercise}
\begin{solution}
(a) By Theorem~\ref{r13-thm-forms}, \(\mathrm{d}f=\bigl(x'(A+A')+b'\bigr)\mathrm{d}x\); the gradient is \((A+A')x+b\), and by
Example~\ref{r13-exm-hessians}(a) the Hessian is
\[
A+A'=\begin{pmatrix}2&2&0&0\\2&4&2&0\\0&2&6&2\\0&0&2&4\end{pmatrix}
\]
(b) Solve \((A+A')x=-b=(2,4,4,-2)'\):
\[
\left(\begin{array}{cccc|c}2&2&0&0&2\\2&4&2&0&4\\0&2&6&2&4\\0&0&2&4&-2\end{array}\right)
\xrightarrow{R_2\leftarrow R_2-R_1}
\left(\begin{array}{cccc|c}2&2&0&0&2\\0&2&2&0&2\\0&2&6&2&4\\0&0&2&4&-2\end{array}\right)
\xrightarrow{R_3\leftarrow R_3-R_2}
\left(\begin{array}{cccc|c}2&2&0&0&2\\0&2&2&0&2\\0&0&4&2&2\\0&0&2&4&-2\end{array}\right)
\]
\[
\left(\begin{array}{cccc|c}2&2&0&0&2\\0&2&2&0&2\\0&0&4&2&2\\0&0&2&4&-2\end{array}\right)
\xrightarrow{R_4\leftarrow R_4-\frac12R_3}
\left(\begin{array}{cccc|c}2&2&0&0&2\\0&2&2&0&2\\0&0&4&2&2\\0&0&0&3&-3\end{array}\right)
\]
so \(x_4=-1\); \(4x_3-2=2\) gives \(x_3=1\); \(2x_2+2=2\) gives \(x_2=0\); and \(2x_1=2\) gives \(x_1=1\): the stationary
point is \(\hat x=(1,0,1,-1)'\).

(c) The operations add multiples of a row to a later row, so the leading principal minors of \(A+A'\) are the
products of the pivots: \(2\), \(2\cdot2=4\), \(2\cdot2\cdot4=16\) and \(\det(A+A')=2\cdot2\cdot4\cdot3=48\)
(Theorem~\ref{r5d-thm-gauss}), all positive. So the Hessian is PD at every \(x\) (Theorem~\ref{r8-thm-leading-minors}), and
\(\hat x\) is the strict global minimizer (Theorem~\ref{r13-thm-extrema}(d)). As \(\hat x'(A+A')\hat x=-b'\hat x\),
\(f(\hat x)=\frac12\hat x'(A+A')\hat x+b'\hat x=\frac12b'\hat x=\frac12(-2+0-4-2)=-4\).
\end{solution}

\begin{exercise}\label{r13-dr-matrix-derivatives}
Let \(A=J_4\) and
\[
X_0=\begin{pmatrix}2&-1&0&0\\-1&2&-1&0\\0&-1&2&-1\\0&0&-1&1\end{pmatrix}
\]
\begin{enumerate}
\item Find \(\partial\det X/\partial X\) and \(\partial\log\det X/\partial X\) at \(X=X_0\) (\(X\) an unrestricted \(4\times4\) variable).
\item Write \(\mathrm D(X^{-1})\) at \(X_0\) as a Kronecker product, and find the differential of \(X^{-1}\) at \(X_0\) for the
increment \(\mathrm{d}X=E_{41}\).
\item Find \(\partial\tr(AX^{-1})/\partial X\) at \(X_0\).
\item Now let \(X\) be a symmetric variable. Find \(\partial\log\det X/\partial X\) at \(X_0\) in the sense of Theorem~\ref{r13-thm-symmetric}.
\end{enumerate}
\end{exercise}
\begin{solution}
\(X_0\) is the matrix \(T\) of Example~\ref{r9-exm-logdet}, the inverse of the matrix \(M=\bigl(\min(i,j)\bigr)\) there, and
\(\det M=1\). So \(X_0\) is invertible, \(\det X_0=1/\det M=1\) (Theorem~\ref{r5d-thm-product}(c)), and
\[
X_0^{-1}=M=\begin{pmatrix}1&1&1&1\\1&2&2&2\\1&2&3&3\\1&2&3&4\end{pmatrix}
\]
which is symmetric (check: \(X_0M=I_4\)).

(a) By Theorem~\ref{r13-thm-trace-det}(b),(d), \(\partial\det X/\partial X=\det X_0\,(X_0^{-1})'=M\), and
\(\partial\log\det X/\partial X=(X_0^{-1})'=M\) is the same matrix, as \(\det X_0=1>0\).

(b) By Theorem~\ref{r13-thm-trace-det}(c), \(\mathrm D(X^{-1})=-(X_0^{-1})'\otimes X_0^{-1}=-M\otimes M\), a
\(16\times16\) matrix, and \(\mathrm{d}X^{-1}=-X_0^{-1}(\mathrm{d}X)X_0^{-1}\). With \(\mathrm{d}X=E_{41}=e_4e_1'\)
(Proposition~\ref{r1-prop-units}), \(Me_4=(1,2,3,4)'\) is column \(4\) of \(M\) and \(e_1'M=(1,1,1,1)\) is row \(1\), so
\[
\mathrm{d}X^{-1}=-(Me_4)(e_1'M)=-\begin{pmatrix}1\\2\\3\\4\end{pmatrix}\begin{pmatrix}1&1&1&1\end{pmatrix}
=-\begin{pmatrix}1&1&1&1\\2&2&2&2\\3&3&3&3\\4&4&4&4\end{pmatrix}
\]
(c) By Example~\ref{r13-exm-matrix-functions}(b), \(\partial\tr(AX^{-1})/\partial X=-(X_0^{-1}AX_0^{-1})'\). Here
\(A=J_4=\one_4\one_4'\), so \(X_0^{-1}AX_0^{-1}=(M\one_4)(\one_4'M)=vv'\) with \(v=M\one_4=(4,7,9,10)'\), the row sums of \(M\)
(\(\one_4'M=v'\), as \(M\) is symmetric). The derivative is
\[
-vv'=-\begin{pmatrix}16&28&36&40\\28&49&63&70\\36&63&81&90\\40&70&90&100\end{pmatrix}
\]
(d) The leading principal minors of \(X_0\) are \(2\), \(4-1=3\), \(2(4-1)+(-1)\cdot2=4\) (along row \(1\),
Theorem~\ref{r5d-thm-expansion}) and \(\det X_0=1\), all positive, so \(X_0\) is PD (Theorem~\ref{r8-thm-leading-minors}). By
Theorem~\ref{r13-thm-symmetric}(b),
\[
\frac{\partial\log\det X}{\partial X}=2X_0^{-1}-\diag(X_0^{-1})=2M-\diag(1,2,3,4)
=\begin{pmatrix}1&2&2&2\\2&2&4&4\\2&4&3&6\\2&4&6&4\end{pmatrix}
\]
the off-diagonal entries double, since \(x_{ij}=x_{ji}\) is one variable.
\end{solution}

\begin{exercise}\label{r9-ex-matrix-drill}
Let
\[
X_0=\begin{pmatrix}1&1&0&0\\0&1&1&0\\0&0&1&1\\0&0&0&1\end{pmatrix},\qquad
v=\begin{pmatrix}1\\-1\\1\\-1\end{pmatrix}
\]
\begin{enumerate}
\item Show that
\[
X_0^{-1}=\begin{pmatrix}1&-1&1&-1\\0&1&-1&1\\0&0&1&-1\\0&0&0&1\end{pmatrix}
\]
and find \(\det X_0\), \(\partial\det X/\partial X\) and \(\partial\log\det X/\partial X\) at \(X_0\).
\item Find the differential of \(X^{-1}\) at \(X_0\) for the increment \(\mathrm{d}X=E_{41}\), and show that it equals
\(vv'\).
\end{enumerate}
\end{exercise}
\begin{solution}
(a) Multiplying, \(X_0\) times the given matrix is \(I_4\), so it is \(X_0^{-1}\) (Theorem~\ref{r2-thm-invertible}).
\(X_0\) is upper triangular with diagonal entries \(1\), so \(\det X_0=1\) (Theorem~\ref{r5d-thm-basic}(c)). By
Theorem~\ref{r13-thm-trace-det}(b),(d),
\[
\frac{\partial\det X}{\partial X}(X_0)=\det X_0\,(X_0^{-1})'=\frac{\partial\log\det X}{\partial X}(X_0)=(X_0^{-1})'
=\begin{pmatrix}1&0&0&0\\-1&1&0&0\\1&-1&1&0\\-1&1&-1&1\end{pmatrix}
\]
(b) By Theorem~\ref{r13-thm-trace-det}(c), \(\mathrm{d}(X^{-1})=-X_0^{-1}e_4e_1'X_0^{-1}\) (\(E_{41}=e_4e_1'\),
Proposition~\ref{r1-prop-units}). Here \(X_0^{-1}e_4=(-1,1,-1,1)'=-v\) (the fourth column) and \(e_1'X_0^{-1}=v'\) (the first
row), so
\[
\mathrm{d}(X^{-1})=vv'=\begin{pmatrix}1&-1&1&-1\\-1&1&-1&1\\1&-1&1&-1\\-1&1&-1&1\end{pmatrix}
\]
\end{solution}

\begin{exercise}\label{r13-dr-extrema}
Let \(f(x)=x_1^3+x_2^3-3x_1x_2+x_3^2+x_3x_4+x_4^2-3x_3\) for \(x\in\R^4\).
\begin{enumerate}
\item Find the gradient, the Hessian and the stationary points of \(f\).
\item Classify each stationary point by Theorem~\ref{r13-thm-extrema}(c).
\item Is the local minimum a global minimum?
\end{enumerate}
\end{exercise}
\begin{solution}
(a) The gradient is \((3x_1^2-3x_2,\ 3x_2^2-3x_1,\ 2x_3+x_4-3,\ x_3+2x_4)'\) and
\[
\mathrm Hf(x)=\begin{pmatrix}6x_1&-3&0&0\\-3&6x_2&0&0\\0&0&2&1\\0&0&1&2\end{pmatrix}
\]
The gradient is \(0\) when \(x_2=x_1^2\) and \(x_1=x_2^2=x_1^4\), that is, \(x_1(x_1^3-1)=0\): \(x_1=0\) or \(x_1=1\); and
when \(x_3=-2x_4\) and \(-4x_4+x_4=3\), that is, \(x_4=-1\), \(x_3=2\). The stationary points are \((0,0,2,-1)'\) and
\((1,1,2,-1)'\).

(b) The lower right block \(\left(\begin{smallmatrix}2&1\\1&2\end{smallmatrix}\right)\) of \(\mathrm Hf\) has the
eigenvalues \(3\) and \(1\), with the eigenvectors \((0,0,1,1)'\) and \((0,0,1,-1)'\) of \(\mathrm Hf\).
\begin{itemize}
\item At \((0,0,2,-1)'\), the upper left block is \(\left(\begin{smallmatrix}0&-3\\-3&0\end{smallmatrix}\right)\), with the
eigenvalues \(3\) and \(-3\) (eigenvectors \((1,-1,0,0)'\), \((1,1,0,0)'\)). So \(\mathrm Hf(0,0,2,-1)\) has the eigenvalues
\(3\), \(-3\), \(3\), \(1\): it is indefinite, and \((0,0,2,-1)'\) is neither a local minimum nor a local maximum.
\item At \((1,1,2,-1)'\), the upper left block is \(\left(\begin{smallmatrix}6&-3\\-3&6\end{smallmatrix}\right)\), with the
eigenvalues \(9\) and \(3\) (eigenvectors \((1,-1,0,0)'\), \((1,1,0,0)'\)). So \(\mathrm Hf(1,1,2,-1)\) has the eigenvalues
\(9\), \(3\), \(3\), \(1\): it is PD (Theorem~\ref{r8-thm-eigen-criterion}(a)), and \(f\) has a strict local minimum at
\((1,1,2,-1)'\), with \(f(1,1,2,-1)=1+1-3+4-2+1-6=-4\).
\end{itemize}
(c) No: \(f(x_1,0,2,-1)=x_1^3+4-2+1-6=x_1^3-3\) is less than \(-4\) for \(x_1<-1\). The Hessian is not PSD on all of \(\R^4\)
(for instance at \((0,0,2,-1)'\)), so Theorem~\ref{r13-thm-extrema}(d) does not apply.
\end{solution}

\begin{exercise}\label{r9-ex-classify}
Let \(M=4I_4+3J_4\) and \(f(x)=\frac14(x'x)^2-\frac12x'Mx\) on \(\R^4\).
\begin{enumerate}
\item Find the gradient and the Hessian of \(f\), and the eigenvalues and eigenvectors of \(M\).
\item Find all stationary points of \(f\), and classify them by the eigenvalues of the Hessian. Treat
\(x=(1,1,-1,-1)'\) explicitly.
\item Find the global minimum of \(f\).
\end{enumerate}
\end{exercise}
\begin{solution}
(a)
\begin{itemize}
\item By Theorems~\ref{r13-thm-forms}(c) and~\ref{r13-thm-diff-rules}, \(\mathrm{d}(x'x)=2x'\mathrm{d}x\) and
\(\mathrm{d}f=\frac12(x'x)\,2x'\mathrm{d}x-x'M\,\mathrm{d}x=\bigl((x'x)x-Mx\bigr)'\mathrm{d}x\): the gradient is
\(g(x)=(x'x)x-Mx\).
\item \(\mathrm{d}g=\bigl(\mathrm{d}(x'x)\bigr)x+(x'x)\,\mathrm{d}x-M\,\mathrm{d}x=\bigl(2xx'+(x'x)I_4-M\bigr)\mathrm{d}x\), so by
Definition~\ref{r13-def-hessian}, \(\mathrm Hf(x)=\mathrm Dg(x)=(x'x)I_4+2xx'-M\).
\item \(J_4\one_4=4\one_4\) and \(J_4v=\one_4(\one_4'v)=0\) for \(v\perp\one_4\), so \(M\one_4=16\one_4\) and \(Mv=4v\) for every
\(v\ne0\) in the \(3\)-dimensional subspace \(\one_4^\perp\): the eigenvalues of \(M\) are \(16\) and \(4\) (three times).
\end{itemize}
(b) \(g(x)=0\) means \(Mx=(x'x)x\): \(x=0\), or \(x\) is an eigenvector of \(M\) with \(x'x\) equal to its eigenvalue. The
stationary points are \(0\), \(\pm2\one_4\) (\(x'x=16\)), and every \(x\perp\one_4\) with \(x'x=4\), infinitely many.
\begin{itemize}
\item At \(\pm2\one_4\), \(\mathrm Hf=16I_4+8J_4-4I_4-3J_4=12I_4+5J_4\), with the eigenvalues \(32\) (on \(\one_4\)) and \(12\)
(three times): a strict local minimum (Theorem~\ref{r13-thm-extrema}(c)), with \(f=64-128=-64\).
\item At \(x\perp\one_4\) with \(x'x=4\), \(\mathrm Hf=2xx'-3J_4\). Then \(\mathrm Hf\,x=8x\), \(\mathrm Hf\,\one_4=-12\one_4\), and
\(\mathrm Hf\,v=0\) for \(v\perp x,\one_4\): the eigenvalues are \(8,-12,0,0\). As \(8>0>-12\), \(\mathrm Hf\) is indefinite, and
\(x\) is neither a local minimum nor a local maximum (Theorem~\ref{r13-thm-extrema}(c)), with \(f=4-8=-4\). For
\(x=(1,1,-1,-1)'\),
\[
\mathrm Hf(x)=\begin{pmatrix}-1&-1&-5&-5\\-1&-1&-5&-5\\-5&-5&-1&-1\\-5&-5&-1&-1\end{pmatrix}
\]
and \(\mathrm Hf(x)(1,-1,0,0)'=\mathrm Hf(x)(0,0,1,-1)'=0\), \(\mathrm Hf(x)\,x=8x\), \(\mathrm Hf(x)\one_4=-12\one_4\).
\item At \(0\), \(\mathrm Hf=-M\) has the eigenvalues \(-16,-4,-4,-4\): it is ND, a strict local maximum, \(f(0)=0\).
\end{itemize}
(c) By Theorem~\ref{r9-thm-rayleigh}(a) for \(M\), whose largest eigenvalue is \(16\), \(x'Mx\leqslant16s\) with \(s=x'x\),
so \(f(x)\geqslant\frac14s^2-8s=\frac14(s-16)^2-64\geqslant-64\). The global minimum is \(-64\), attained at \(\pm2\one_4\).
\end{solution}

\begin{exercise}\label{r13-dr-lagrange}
Four unbiased estimators \(t=(t_1,t_2,t_3,t_4)'\) of a number \(\theta\) have the covariance matrix
\[
\Sigma=\begin{pmatrix}1&-1&0&0\\-1&4&-1&0\\0&-1&2&-1\\0&0&-1&1\end{pmatrix}
\]
A combination \(a't\) is unbiased when \(\one'a=1\); we seek the one of least variance (as in a minimum-variance
portfolio).
\begin{enumerate}
\item Write down the Lagrangian for minimizing \(\var(a't)=a'\Sigma a\) subject to \(\one'a=1\), and the conditions
\(\mathrm D\psi=0\).
\item Solve them: find \(w=\Sigma^{-1}\one\) by Gauss elimination, then \(a\) and the multiplier.
\item Show that this \(a\) gives the smallest variance, and compare with the equal weights \(\frac14\one\).
\end{enumerate}
\end{exercise}
\begin{solution}
(a) \(\psi(a)=a'\Sigma a-2\lambda(\one'a-1)\) (multiplier \(2\lambda\)); by Theorem~\ref{r13-thm-forms},
\(\mathrm{d}\psi=2(\Sigma a-\lambda\one)'\mathrm{d}a\), so the conditions are \(\Sigma a=\lambda\one\) and \(\one'a=1\).

(b) Solve \(\Sigma w=\one\):
\[
\left(\begin{array}{cccc|c}1&-1&0&0&1\\-1&4&-1&0&1\\0&-1&2&-1&1\\0&0&-1&1&1\end{array}\right)
\xrightarrow{R_2\leftarrow R_2+R_1}
\left(\begin{array}{cccc|c}1&-1&0&0&1\\0&3&-1&0&2\\0&-1&2&-1&1\\0&0&-1&1&1\end{array}\right)
\xrightarrow{R_2\leftrightarrow R_3}
\left(\begin{array}{cccc|c}1&-1&0&0&1\\0&-1&2&-1&1\\0&3&-1&0&2\\0&0&-1&1&1\end{array}\right)
\]
\[
\left(\begin{array}{cccc|c}1&-1&0&0&1\\0&-1&2&-1&1\\0&3&-1&0&2\\0&0&-1&1&1\end{array}\right)
\xrightarrow{R_3\leftarrow R_3+3R_2}
\left(\begin{array}{cccc|c}1&-1&0&0&1\\0&-1&2&-1&1\\0&0&5&-3&5\\0&0&-1&1&1\end{array}\right)
\xrightarrow{R_3\leftrightarrow R_4}
\left(\begin{array}{cccc|c}1&-1&0&0&1\\0&-1&2&-1&1\\0&0&-1&1&1\\0&0&5&-3&5\end{array}\right)
\]
\[
\left(\begin{array}{cccc|c}1&-1&0&0&1\\0&-1&2&-1&1\\0&0&-1&1&1\\0&0&5&-3&5\end{array}\right)
\xrightarrow{R_4\leftarrow R_4+5R_3}
\left(\begin{array}{cccc|c}1&-1&0&0&1\\0&-1&2&-1&1\\0&0&-1&1&1\\0&0&0&2&10\end{array}\right)
\]
so \(w_4=5\); \(-w_3+5=1\) gives \(w_3=4\); \(-w_2+8-5=1\) gives \(w_2=2\); and \(w_1-2=1\) gives \(w_1=3\):
\(w=(3,2,4,5)'\). Then \(a=\lambda w\) with \(1=\one'a=14\lambda\), so \(\lambda=\frac1{14}\) and \(a=\frac1{14}(3,2,4,5)'\).
Check: \(\Sigma w=(3-2,\ -3+8-4,\ -2+8-5,\ -4+5)'=\one\), so \(\Sigma a=\frac1{14}\one\).

(c) The leading principal minors of \(\Sigma\) are \(1\), \(4-1=3\), \(1\cdot(8-1)+(-1)\cdot2=5\) (along row \(1\),
Theorem~\ref{r5d-thm-expansion}) and \(\det\Sigma=(-1)^2\cdot1\cdot(-1)\cdot(-1)\cdot2=2\) (Theorem~\ref{r5d-thm-gauss}, two
interchanges), all positive, so \(\Sigma\) is PD (Theorem~\ref{r8-thm-leading-minors}), and \(\mathrm H\psi=2\Sigma\) is PD
everywhere: by Theorem~\ref{r13-thm-extrema}(e), \(a\) is the strict minimizer under the constraint. Its variance is
\(a'\Sigma a=\lambda\,a'\one=\frac1{14}\). The equal weights give \(\frac1{16}\one'\Sigma\one=\frac1{16}(0+2+0+0)=\frac18>\frac1{14}\),
where \(0,2,0,0\) are the row sums of \(\Sigma\).
\end{solution}

\begin{exercise}\label{r9-ex-equality-qp}
Let \(a=(3,1,4,1)'\) and \(L=\{x\in\R^4:x_1+x_2+x_3+x_4=1,\ x_1+x_2-x_3-x_4=3\}\).
\begin{enumerate}
\item Find the point \(x^*\) of \(L\) closest to \(a\), and the distance from \(a\) to \(L\), by
Theorem~\ref{r9-thm-equality-qp} with \(A=I_4\) and \(b=a\).
\item Check Theorem~\ref{r9-thm-equality-qp}(b) for the point \(x=(2,0,0,-1)'\) of \(L\).
\end{enumerate}
\end{exercise}
\begin{solution}
(a) With \(A=I_4\) and \(b=a\), \(q(x)=\frac12x'x-a'x=\frac12\lVert x-a\rVert^2-\frac12\lVert a\rVert^2\), so
minimizing \(q\) on \(L\) is minimizing the distance to \(a\), and \(\hat x=a\). Here
\[
C=\begin{pmatrix}1&1&1&1\\1&1&-1&-1\end{pmatrix},\qquad d=\begin{pmatrix}1\\3\end{pmatrix},\qquad
V=CC'=4I_2,\qquad Ca-d=\begin{pmatrix}9\\-1\end{pmatrix}-\begin{pmatrix}1\\3\end{pmatrix}=\begin{pmatrix}8\\-4\end{pmatrix}
\]
\(\rank C=2\) (the rows are not multiples of each other), and \(\lambda=V^{-1}(Ca-d)=(2,-1)'\). By
Theorem~\ref{r9-thm-equality-qp}(a),
\[
x^*=a-C'\lambda=\begin{pmatrix}3\\1\\4\\1\end{pmatrix}-\begin{pmatrix}1\\1\\3\\3\end{pmatrix}=\begin{pmatrix}2\\0\\1\\-2\end{pmatrix}
\]
which satisfies \(2+0+1-2=1\) and \(2+0-1+2=3\). By Theorem~\ref{r9-thm-equality-qp}(c),
\(\frac12\lVert x^*-a\rVert^2=q(x^*)-q(a)=\frac12(Ca-d)'V^{-1}(Ca-d)=\frac12\cdot\frac{64+16}4=10\); so the distance is
\(\sqrt{20}=2\sqrt5\), and directly \(\lVert a-x^*\rVert^2=1+1+9+9=20\).

(b) \(x=(2,0,0,-1)'\) lies in \(L\) (\(2+0+0-1=1\), \(2+0-0+1=3\)). Here \(q(x)=\frac12\cdot5-(6-1)=-\frac52\),
\(q(x^*)=\frac12\cdot9-(6+4-2)=-\frac72\), and \(\frac12\lVert x-x^*\rVert^2=\frac12\lVert(0,0,-1,1)'\rVert^2=1\);
indeed \(-\frac52=-\frac72+1\).
\end{solution}

\begin{exercise}\label{r13-dr-least-squares}
Five units have the values \(x_{i1}=2,2,3,3,5\) and \(x_{i2}=1,5,2,2,0\) of two explanatory variables and the responses
\(y=(4,8,5,7,6)'\). In the model \(y_i=\beta_0+\beta_1x_{i1}+\beta_2x_{i2}+\varepsilon_i\) (\(k=2\)):
\begin{enumerate}
\item Write down \(X\), \(X'X\), \(X'y\) and the gradient of \(S(\beta)\) at \(\beta=0\).
\item Solve the normal equations by Gauss elimination.
\item Show that the solution minimizes \(S\), and find the residuals and \(S(\hat\beta)\).
\end{enumerate}
\end{exercise}
\begin{solution}
(a)
\[
X=\begin{pmatrix}1&2&1\\1&2&5\\1&3&2\\1&3&2\\1&5&0\end{pmatrix},\qquad
X'X=\begin{pmatrix}5&15&10\\15&51&24\\10&24&34\end{pmatrix},\qquad X'y=\begin{pmatrix}30\\90\\68\end{pmatrix}
\]
By Theorem~\ref{r13-thm-ols}(a), the gradient at \(\beta=0\) is \(-2X'y=(-60,-180,-136)'\).

(b)
\[
\left(\begin{array}{ccc|c}5&15&10&30\\15&51&24&90\\10&24&34&68\end{array}\right)
\xrightarrow[R_3\leftarrow R_3-2R_1]{R_2\leftarrow R_2-3R_1}
\left(\begin{array}{ccc|c}5&15&10&30\\0&6&-6&0\\0&-6&14&8\end{array}\right)
\xrightarrow{R_3\leftarrow R_3+R_2}
\left(\begin{array}{ccc|c}5&15&10&30\\0&6&-6&0\\0&0&8&8\end{array}\right)
\]
so \(\beta_2=1\), \(6\beta_1-6=0\) gives \(\beta_1=1\), and \(5\beta_0+15+10=30\) gives \(\beta_0=1\): \(\hat\beta=(1,1,1)'\).

(c) The pivots \(5\), \(6\), \(8\) are positive, so the leading principal minors \(5\), \(30\), \(240\) of \(X'X\) are positive and
\(\mathrm HS=2X'X\) is PD (Theorem~\ref{r8-thm-leading-minors}): \(\hat\beta\) is the strict global minimizer (Theorem~\ref{r13-thm-ols}(c)).
\(X\hat\beta=(4,8,6,6,6)'\), so \(e=y-X\hat\beta=(0,0,-1,1,0)'\) and \(S(\hat\beta)=e'e=2\); indeed \(X'e=(0,\ 0,\ 0)'\).
\end{solution}

\begin{exercise}\label{r13-dr-newton}
Let \(f(x)=\begin{pmatrix}x_1^2-x_1x_2-2x_2+2\\2x_1+x_2^2-3\end{pmatrix}\), which is \(0\) at \(x=(1,1)'\).
\begin{enumerate}
\item Find the derivative \(G(x)=\mathrm Df(x)\), and check that \(G(1,1)\) is invertible.
\item Starting from \(x_0=(-1,-1)'\), compute two Newton--Raphson steps \(x_1\) and \(x_2\) (Example~\ref{r13-exm-newton}).
\item Compare the squared distances of \(x_0\), \(x_1\), \(x_2\) from the solution \((1,1)'\).
\end{enumerate}
\end{exercise}
\begin{solution}
(a) \(G(x)=\begin{pmatrix}2x_1-x_2&-x_1-2\\2&2x_2\end{pmatrix}\), and \(G(1,1)=\begin{pmatrix}1&-3\\2&2\end{pmatrix}\) has determinant \(8\ne0\).

(b) At \(x_0=(-1,-1)'\): \(f(x_0)=(4,-4)'\) and \(G(x_0)=\begin{pmatrix}-1&-1\\2&-2\end{pmatrix}\), with determinant \(4\) and
\(G(x_0)^{-1}=\frac14\begin{pmatrix}-2&1\\-2&-1\end{pmatrix}\) (Theorem~\ref{r5d-thm-adjugate}). So
\(x_1=x_0-G(x_0)^{-1}f(x_0)=(-1,-1)'-\frac14(-12,-4)'=(2,0)'\). At \(x_1\): \(f(x_1)=(6,1)'\) and \(G(x_1)=\begin{pmatrix}4&-4\\2&0\end{pmatrix}\),
with determinant \(8\) and \(G(x_1)^{-1}=\frac18\begin{pmatrix}0&4\\-2&4\end{pmatrix}\). So \(x_2=(2,0)'-\frac18(4,-8)'=(\frac32,1)'\).

(c) The squared distances from \((1,1)'\) are \(\lVert(-2,-2)'\rVert^2=8\), \(\lVert(1,-1)'\rVert^2=2\) and
\(\lVert(\frac12,0)'\rVert^2=\frac14\). A third step gives \(x_3=(\frac{23}{22},\frac{21}{22})'\), at squared distance
\(\lVert(\frac1{22},-\frac1{22})'\rVert^2=\frac2{484}=\frac1{242}\): near the solution each step roughly squares the error.
\end{solution}

\begin{exercise}\label{r13-dr-jacobian}
Let \(A=\begin{pmatrix}2&0\\1&1\end{pmatrix}\) and \(B=\begin{pmatrix}1&1\\0&3\end{pmatrix}\), and consider \(Y=AXB\) for \(X\in\R^{2\times2}\).
\begin{enumerate}
\item Write down the matrix \(J\) with \(\operatorname{vec}Y=J\operatorname{vec}X\), and find \(\det J\) in two ways. Check \(J\operatorname{vec}X\) for
\(X=\begin{pmatrix}1&0\\2&1\end{pmatrix}\).
\item If the four entries of \(X\) have the joint density \(f_X\), what is the density of \(Y\)?
\item Compute \(\int_{\R^4}e^{-\frac12x'Px}\,\mathrm{d}x\) for the matrix
\[
P=\begin{pmatrix}2&0&1&-1\\0&2&1&1\\1&1&3&0\\-1&1&0&3\end{pmatrix}
\]
of Example~\ref{r8-exm-leading}.
\end{enumerate}
\end{exercise}
\begin{solution}
(a) By Theorem~\ref{r13-thm-jacobian}(c), \(J=B'\otimes A\) with \(B'=\begin{pmatrix}1&0\\1&3\end{pmatrix}\):
\[
J=\begin{pmatrix}A&0\\A&3A\end{pmatrix}=\begin{pmatrix}2&0&0&0\\1&1&0&0\\2&0&6&0\\1&1&3&3\end{pmatrix}
\]
\(J\) is lower triangular, so \(\det J=2\cdot1\cdot6\cdot3=36\) (Theorems~\ref{r5d-thm-basic}(c) and~\ref{r5d-thm-transpose}); and
\(\det J=(\det A)^2(\det B)^2=2^2\cdot3^2=36\) (Theorem~\ref{r13-thm-kron-det}(a)). For the given \(X\),
\(AXB=\begin{pmatrix}2&0\\3&1\end{pmatrix}B=\begin{pmatrix}2&2\\3&6\end{pmatrix}\), and \(J(1,2,0,1)'=(2,3,2,6)'=\operatorname{vec}Y\).

(b) \(f_Y(Y)=f_X(A^{-1}YB^{-1})/36\) (Theorem~\ref{r13-thm-jacobian}(c)), with \(A^{-1}=\frac12\begin{pmatrix}1&0\\-1&2\end{pmatrix}\) and
\(B^{-1}=\frac13\begin{pmatrix}3&-1\\0&1\end{pmatrix}\).

(c) By Example~\ref{r8-exm-leading}, \(P\) is PD with \(\det P=\Delta_4=16\), so by Example~\ref{r13-exm-aitken} with
\(n=4\) the integral is \((2\pi)^{4/2}\cdot16^{-1/2}=\frac{4\pi^2}4=\pi^2\).
\end{solution}

\begin{exercise}\label{r9-ex-jacobian-det}
\begin{enumerate}
\item Find the Jacobian determinant of \(X\mapsto AXB\) on \(\R^{2\times3}\) for
\[
A=\begin{pmatrix}3&1\\1&1\end{pmatrix},\qquad B=\begin{pmatrix}2&1&1\\1&2&1\\1&1&2\end{pmatrix}
\]
\item Find \(X_0^{-1}\) and the Jacobian determinant of \(X\mapsto X^{-1}\) at
\[
X_0=\begin{pmatrix}1&1&0&0\\0&1&1&0\\0&0&1&1\\-1&0&0&1\end{pmatrix}
\]
\item Find the Jacobian matrix of \(X\mapsto X'\) on \(\R^{2\times2}\) and its determinant.
\end{enumerate}
\end{exercise}
\begin{solution}
(a) \(\det A=2\) and, along row \(1\) (Theorem~\ref{r5d-thm-expansion}), \(\det B=2(4-1)-(2-1)+(1-2)=4\). By
Theorem~\ref{r13-thm-jacobian}(c) with \(m=2\) and \(n=3\), the Jacobian matrix is \(B'\otimes A\in\R^{6\times6}\), with
determinant \((\det A)^3(\det B)^2=8\cdot16=128\).

(b) Along column \(1\) (Theorem~\ref{r5d-thm-expansion}),
\[
\det X_0=1\cdot\det\begin{pmatrix}1&1&0\\0&1&1\\0&0&1\end{pmatrix}+(-1)\cdot(-1)^{4+1}\det\begin{pmatrix}1&0&0\\1&1&0\\0&1&1\end{pmatrix}
=1+1=2
\]
the two minors being triangular with diagonal entries \(1\) (Theorems~\ref{r5d-thm-basic}(c) and~\ref{r5d-thm-transpose}).
Multiplying checks \(X_0X_0^{-1}=I_4\) for
\[
X_0^{-1}=\frac12\begin{pmatrix}1&-1&1&-1\\1&1&-1&1\\-1&1&1&-1\\1&-1&1&1\end{pmatrix}
\]
(Theorem~\ref{r2-thm-invertible}). By Theorem~\ref{r13-thm-jacobian}(d) with \(n=4\), the Jacobian determinant is
\((-1)^4(\det X_0)^{-8}=2^{-8}=\frac1{256}\).

(c) The map is linear, with \(\operatorname{vec}(X')=(x_{11},x_{12},x_{21},x_{22})'\) for
\(\operatorname{vec}X=(x_{11},x_{21},x_{12},x_{22})'\): \(\operatorname{vec}(X')=K_2\operatorname{vec}X\)
(Definition~\ref{r13-def-commutation}), and \(\mathrm{d}(X')=(\mathrm{d}X)'\) gives
\(\operatorname{vec}\mathrm{d}(X')=K_2\operatorname{vec}\mathrm{d}X\). So the Jacobian matrix is
(Theorem~\ref{r13-thm-identification}(b))
\[
K_2=\begin{pmatrix}1&0&0&0\\0&0&1&0\\0&1&0&0\\0&0&0&1\end{pmatrix}
\]
obtained from \(I_4\) by interchanging rows \(2\) and \(3\), and \(\det K_2=-1\) (Theorem~\ref{r5d-thm-basic}(d)).
\end{solution}

\begin{exercise}\label{r9-ex-sphere}
Let
\[
A=\begin{pmatrix}5&1&2&0\\1&5&0&2\\2&0&5&1\\0&2&1&5\end{pmatrix}
\]
and \(h_1=(1,1,1,1)'\), \(h_2=(1,-1,1,-1)'\), \(h_3=(1,1,-1,-1)'\), \(h_4=(1,-1,-1,1)'\).
\begin{enumerate}
\item Compute \(Ah_1,\dots,Ah_4\), and show that every eigenvector of \(A\) is a multiple of one of \(h_1,\dots,h_4\).
\item Find all stationary points of \(f(x)=x'Ax\) on the unit sphere of \(\R^4\), and the maximum and the minimum of
\(f\) on the sphere.
\end{enumerate}
\end{exercise}
\begin{solution}
(a)
\begin{itemize}
\item Row by row, \(Ah_1=(8,8,8,8)'=8h_1\), \(Ah_2=(6,-6,6,-6)'=6h_2\), \(Ah_3=(4,4,-4,-4)'=4h_3\) and
\(Ah_4=(2,-2,-2,2)'=2h_4\).
\item The \(h_j\) are nonzero and pairwise orthogonal, hence linearly independent
(Proposition~\ref{r6-prop-orthogonal-list}(b)), and a basis of \(\R^4\) (Theorem~\ref{thm:r4-counting}(d)). For
\(x=\sum_jc_jh_j\), \(Ax=8c_1h_1+6c_2h_2+4c_3h_3+2c_4h_4\), so \(Ax=\lambda x\) means
\((8-\lambda)c_1=(6-\lambda)c_2=(4-\lambda)c_3=(2-\lambda)c_4=0\).
\item If \(x\ne0\), some \(c_j\ne0\), so \(\lambda\in\{8,6,4,2\}\), and then the other three coefficients vanish: \(x\) is
a multiple of one \(h_j\).
\end{itemize}
(b) Each \(\lVert h_j\rVert=2\). By Proposition~\ref{r13-prop-sphere} and (a), the stationary points are the unit
eigenvectors \(\pm\frac12h_1\) (\(f=8\)), \(\pm\frac12h_2\) (\(f=6\)), \(\pm\frac12h_3\) (\(f=4\)) and \(\pm\frac12h_4\)
(\(f=2\)). On the sphere \(f(x)=x'Ax/x'x\), so the maximum of \(f\) on the sphere is \(8\) and the minimum is \(2\)
(Theorem~\ref{r9-thm-rayleigh}(a)).
\end{solution}

\begin{exercise}\label{r9-ex-logdet-max}
Let
\[
S=\begin{pmatrix}2&1\\1&2\end{pmatrix}\otimes\begin{pmatrix}3&1\\1&3\end{pmatrix}
=\begin{pmatrix}6&2&3&1\\2&6&1&3\\3&1&6&2\\1&3&2&6\end{pmatrix}
\]
and \(g(X)=\log\det X-\tr(SX)\) for PD \(X\in\R^{4\times4}\).
\begin{enumerate}
\item Maximize \(g(cI_4)\) over \(c>0\).
\item Find the maximizer and the maximum of \(g\) over the PD \(X\in\R^{4\times4}\), and compare with (a).
\end{enumerate}
\end{exercise}
\begin{solution}
Let \(A=\left(\begin{smallmatrix}2&1\\1&2\end{smallmatrix}\right)\) and \(B=\left(\begin{smallmatrix}3&1\\1&3\end{smallmatrix}\right)\).
Their leading principal minors are \(2,3\) and \(3,8\), so \(A\) and \(B\) are PD (Theorem~\ref{r8-thm-leading-minors}), and so
is \(S=A\otimes B\) (Theorem~\ref{r13-thm-kron-eigen}(c)).

(a) \(\tr S=24\), so \(h(c)=g(cI_4)=4\log c-24c\) on the open convex set \((0,\infty)\), with derivative \(4/c-24\) and second
derivative \(-4/c^2<0\). By Theorem~\ref{r13-thm-extrema}(d) applied to \(-h\) (\(n=1\)), the stationary point \(c=\frac16\) is
the strict global maximizer, and \(g(\frac16I_4)=-4\log6-4=-2\log36-4\).

(b)
\begin{itemize}
\item The PD matrices form an open set, and we take the variables \(\operatorname{vech}X\), as in
Example~\ref{r13-exm-covariance}. By Theorems~\ref{r13-thm-diff-rules} and~\ref{r13-thm-trace-det}(a),(d),
\(\mathrm{d}g=\tr(X^{-1}\mathrm{d}X)-\tr(S\,\mathrm{d}X)=\tr\bigl((X^{-1}-S)\,\mathrm{d}X\bigr)\). At a stationary point this is \(0\)
for every symmetric \(\mathrm{d}X\); as \(X^{-1}-S\) is symmetric, \(\mathrm{d}X=X^{-1}-S\) gives
\(\tr\bigl((X^{-1}-S)'(X^{-1}-S)\bigr)=0\), so \(X^{-1}=S\) (Theorem~\ref{r1-thm-trace}(d)), as in
Example~\ref{r13-exm-covariance}(a). So the only stationary point is \(X=S^{-1}\), which is PD
(Proposition~\ref{r8-prop-pd-basic}(c)).
\item It is the global maximizer. Let \(C=S^{1/2}XS^{1/2}\) (Theorem~\ref{r8-thm-sqrt}), PD (Theorem~\ref{r8-thm-congruence}(c)).
Then \(\tr C=\tr(SX)\) (Theorem~\ref{r1-thm-trace}(b)) and \(\log\det C=\log\det S+\log\det X\)
(Theorem~\ref{r5d-thm-product}(b)), so by Lemma~\ref{r14-lem-logdet}
\[
g(X)=\log\det C-\tr C-\log\det S\leqslant-4-\log\det S
\]
with equality only if \(C=I_4\), that is, \(X=S^{-1}\).
\item By Corollary~\ref{r13-cor-kron}(b), with \(A^{-1}=\frac13\left(\begin{smallmatrix}2&-1\\-1&2\end{smallmatrix}\right)\) and
\(B^{-1}=\frac18\left(\begin{smallmatrix}3&-1\\-1&3\end{smallmatrix}\right)\),
\[
S^{-1}=A^{-1}\otimes B^{-1}=\frac1{24}\begin{pmatrix}6&-2&-3&1\\-2&6&1&-3\\-3&1&6&-2\\1&-3&-2&6\end{pmatrix}
\]
and \(\det S=(\det A)^2(\det B)^2=9\cdot64=576\) (Theorem~\ref{r13-thm-kron-det}(a)). The maximum is
\(-\log\det S-4=-\log576-4=-2\log24-4\), larger than the value \(-2\log36-4\) in (a).
\end{itemize}
\end{solution}

\needspace{14\baselineskip}
\subsection*{Problems}

\begin{problem}[\citeat{abadir-magnus}{Exercise~10.13}]\label{r13-ex-kron-nonsingular}
Let \(A\in\R^{m\times p}\) and \(B\in\R^{n\times q}\). If \(A\) and \(B\) are not square, \(A\otimes B\) may still be square. Show that \(A\otimes B\) is
singular unless both \(A\) and \(B\) are square and invertible.
\end{problem}
\begin{solution}
\begin{itemize}
\item \(A\otimes B\in\R^{mn\times pq}\) is square when \(mn=pq\). By Theorem~\ref{r13-thm-kron-det}(b) and Theorem~\ref{thm:r4-rank-rules}(a),
\(\rank(A\otimes B)=\rank A\cdot\rank B\leqslant\min\{m,p\}\min\{n,q\}\).
\item If \(A\otimes B\) is invertible, its rank is \(mn=pq\) (Theorem~\ref{thm:r4-rank-rules}(d)). From \(mn\leqslant\min\{m,p\}\min\{n,q\}\) we get
\(m\leqslant p\) and \(n\leqslant q\), and from \(pq\leqslant\min\{m,p\}\min\{n,q\}\) we get \(p\leqslant m\) and \(q\leqslant n\). So \(p=m\), \(q=n\): \(A\)
and \(B\) are square.
\item Then \(\rank A\cdot\rank B=mn\) with \(\rank A\leqslant m\) and \(\rank B\leqslant n\) forces \(\rank A=m\) and \(\rank B=n\): \(A\) and \(B\) are
invertible.
\end{itemize}
\end{solution}

\begin{problem}[\citeat{abadir-magnus}{Exercise~10.14}]\label{r13-ex-kron-loewner}
Let \(A,B\in\R^{n\times n}\) be PSD. Show that \(A\otimes A-B\otimes B\) is PSD if and only if \(A-B\) is PSD.
\end{problem}
\begin{solution}
\begin{itemize}
\item Let \(A-B\) be PSD. By Theorem~\ref{r13-thm-kron-rules}(b),
\(A\otimes A-B\otimes B=A\otimes A-A\otimes B+A\otimes B-B\otimes B=A\otimes(A-B)+(A-B)\otimes B\). Both terms are PSD by
Theorem~\ref{r13-thm-kron-eigen}(c), and a sum of PSD matrices is PSD (Proposition~\ref{r8-prop-pd-basic}(b)).
\item Conversely, let \(A\otimes A-B\otimes B\) be PSD and \(x\in\R^n\). By the mixed product, \((x\otimes x)'(A\otimes A)(x\otimes x)=(x'Ax)(x'Ax)\), so
\[
0\leqslant(x\otimes x)'(A\otimes A-B\otimes B)(x\otimes x)=(x'Ax)^2-(x'Bx)^2=x'(A+B)x\cdot x'(A-B)x
\]
\item \(x'(A+B)x\geqslant0\). If it is \(>0\), the display gives \(x'(A-B)x\geqslant0\). If it is \(0\), then \(x'Ax=x'Bx=0\) (both are \(\geqslant0\)),
and \(x'(A-B)x=0\). So \(x'(A-B)x\geqslant0\) for every \(x\): \(A-B\) is PSD.
\end{itemize}
\end{solution}

\begin{problem}[\citeat{prasolov}{Problem~27.1}]\label{r13-ex-pr-kron-vectors}
Let \(a,c\in\R^m\) and \(b,d\in\R^n\) be nonzero vectors. Show that \(a\otimes b=c\otimes d\) if and only if
\(a=\alpha c\) and \(d=\alpha b\) for some number \(\alpha\).
\end{problem}
\begin{solution}
\begin{itemize}
\item If \(a=\alpha c\) and \(d=\alpha b\), then
\(a\otimes b=(\alpha c)\otimes b=\alpha(c\otimes b)=c\otimes(\alpha b)=c\otimes d\) (Theorem~\ref{r13-thm-kron-rules}(a)).
\item Conversely, let \(a\otimes b=c\otimes d\). By Definition~\ref{r13-def-kronecker}, \(a\otimes b=(a_1b',\dots,a_mb')'\) and
\(c\otimes d=(c_1d',\dots,c_md')'\); comparing the \(i\)th blocks of \(n\) entries, \(a_ib=c_id\) for \(i=1,\dots,m\).
\item As \(a\ne0\), some \(a_k\ne0\). Then \(c_kd=a_kb\ne0\), as \(b\ne0\); so \(c_k\ne0\), and \(d=\alpha b\) with
\(\alpha=a_k/c_k\).
\item For every \(i\), \(a_ib=c_id=\alpha c_ib\), that is, \((a_i-\alpha c_i)b=0\). As \(b\) has a nonzero entry,
\(a_i=\alpha c_i\). So \(a=\alpha c\).
\end{itemize}
As \(a\otimes b=\operatorname{vec}(ba')\) (Theorem~\ref{r13-thm-vec}(a)), this also says: for nonzero vectors, \(ba'=dc'\) if
and only if \(a=\alpha c\) and \(d=\alpha b\) for some \(\alpha\).
\end{solution}

\begin{problem}[\citeat{abadir-magnus}{Exercise~10.21}]\label{r13-ex-vec-identity}
Let \(n\geqslant1\), \(\alpha\in\R\), \(v=\operatorname{vec}I_n\), and let \(A\in\R^{n\times n}\) be PD. Consider the \(n^2\times n^2\) matrices
\[
B=I_n\otimes I_n-\alpha vv',\qquad C=A\otimes A-\alpha(\operatorname{vec}A)(\operatorname{vec}A)'
\]
Show that:
\begin{enumerate}
\item the eigenvalues of \(B\) are \(1-n\alpha\) (once) and \(1\) (\(n^2-1\) times);
\item \(\det B=1-n\alpha\);
\item \(B\) is PD if and only if \(\alpha<1/n\), and PSD if and only if \(\alpha\leqslant1/n\);
\item the same holds for \(C\).
\end{enumerate}
\end{problem}
\begin{solution}
(a)
\begin{itemize}
\item \(v'v=n\), as \(v\) has \(n\) entries \(1\) and the others \(0\). \(vv'v=nv\), and \(vv'x=0\) for \(x\perp v\); an orthonormal basis of \(\R^{n^2}\)
made of \(v/\sqrt n\) and \(n^2-1\) vectors orthogonal to \(v\) (Theorems~\ref{r6-thm-complement}(c) and~\ref{r6-thm-gram-schmidt}) consists of
eigenvectors of the symmetric \(vv'\), so its eigenvalues are \(n\) (once) and \(0\) (\(n^2-1\) times) (Theorem~\ref{r8-thm-spectral}(c)).
\item \(B=I_{n^2}-\alpha vv'\) (Corollary~\ref{r13-cor-kron}(a)), so its eigenvalues are \(1-\alpha n\) and \(1-\alpha\cdot0=1\)
(Proposition~\ref{r7-prop-eig-poly}(c)).
\end{itemize}
(b) \(\det B\) is the product of the eigenvalues (Theorem~\ref{r7-thm-trace-det}): \(1-n\alpha\).

(c) \(B\) is symmetric; it is PD (PSD) if and only if all eigenvalues are \(>0\) (\(\geqslant0\)) (Theorem~\ref{r8-thm-eigen-criterion}(a)), that is,
\(1-n\alpha>0\) (\(\geqslant0\)).

(d)
\begin{itemize}
\item Let \(R=A^{1/2}\otimes A^{1/2}\) (Theorem~\ref{r8-thm-sqrt}), symmetric and invertible (Corollary~\ref{r13-cor-kron}(b),(e)). Then
\(R(I_n\otimes I_n)R=A^{1/2}A^{1/2}\otimes A^{1/2}A^{1/2}=A\otimes A\) (mixed product), and \(Rv=\operatorname{vec}(A^{1/2}I_nA^{1/2})=\operatorname{vec}A\)
(Theorem~\ref{r13-thm-vec}(b)), so \(R(vv')R=(\operatorname{vec}A)(\operatorname{vec}A)'\). Hence \(C=RBR=R'BR\).
\item By Theorem~\ref{r8-thm-congruence}(c), \(C\) is PD (PSD) if and only if \(B\) is; apply (c).
\end{itemize}
\end{solution}

\begin{problem}[\citeat{harville}{Ch.~16, Exercise~6}, \citeat{harville-solutions}{Ch.~16, Exercise~6}]\label{r13-ex-kron-projection}
For a matrix \(A\) let \(P_A=A(A'A)^-A'\), the matrix of the orthogonal projection onto \(\mathcal C(A)\) (Theorem~\ref{r10-thm-ata}(b)). Let
\(X=A\otimes B\), where \(A\in\R^{m\times n}\) and \(B\in\R^{p\times q}\). Show that \(P_X=P_A\otimes P_B\).
\end{problem}
\begin{solution}
\begin{itemize}
\item By Theorem~\ref{r13-thm-kron-rules}(e),(f), \(X'X=(A'\otimes B')(A\otimes B)=(A'A)\otimes(B'B)\), so \((A'A)^-\otimes(B'B)^-\) is a g-inverse of
\(X'X\) (Corollary~\ref{r13-cor-kron}(d)).
\item As \(P_X\) does not depend on the g-inverse (Theorem~\ref{r10-thm-ata}(b)), the mixed product gives
\[
P_X=(A\otimes B)\bigl((A'A)^-\otimes(B'B)^-\bigr)(A'\otimes B')=\bigl(A(A'A)^-A'\bigr)\otimes\bigl(B(B'B)^-B'\bigr)=P_A\otimes P_B
\]
\end{itemize}
For instance, for the balanced one-way design \(X_1=I_g\otimes\one_m\), \(P_{X_1}=I_g\otimes\bar J_m\), as in Example~\ref{r13-exm-designs}(a).
\end{solution}

\begin{problem}[\citeat{abadir-magnus}{Exercise~11.14}]\label{r13-ex-schur-inequality}
\begin{enumerate}
\item Show that \(\frac12(I_{n^2}-K_n)\) is symmetric and idempotent.
\item Show that \(\tr(A^2)\leqslant\tr(A'A)\) for every \(A\in\R^{n\times n}\), with equality if and only if \(A\) is symmetric (\textbf{Schur's inequality}).
\end{enumerate}
\end{problem}
\begin{solution}
(a) \(K_n\) is symmetric and \(K_n^2=K_n'K_n=I_{n^2}\) (Theorem~\ref{r13-thm-commutation}(b)), so \(N=\frac12(I-K_n)\) is symmetric and
\(N^2=\frac14(I-2K_n+K_n^2)=\frac12(I-K_n)=N\).

(b)
\begin{itemize}
\item By Theorem~\ref{r13-thm-vec}(d), \(\tr(A'A)=(\operatorname{vec}A)'\operatorname{vec}A\) and
\(\tr(A^2)=\tr((A')'A)=(\operatorname{vec}A')'\operatorname{vec}A=(\operatorname{vec}A)'K_n'\operatorname{vec}A=(\operatorname{vec}A)'K_n\operatorname{vec}A\).
\item So \(\tr(A'A)-\tr(A^2)=(\operatorname{vec}A)'(I-K_n)\operatorname{vec}A=2x'Nx\) with \(x=\operatorname{vec}A\), and \(x'Nx=x'N'Nx=\lVert Nx\rVert^2\geqslant0\).
\item Equality holds if and only if \(Nx=0\), that is, \(K_n\operatorname{vec}A=\operatorname{vec}A\), that is, \(\operatorname{vec}A'=\operatorname{vec}A\), or \(A'=A\).
\end{itemize}
\end{solution}

\begin{problem}[\citeat{prasolov}{Problem~32.1}]\label{r13-ex-pr-transpose-map}
Let \(m,n\geqslant1\) with \(mn>1\). Show that
\begin{enumerate}
\item there are no matrices \(A\), \(B\) with \(AXB=X'\) for every \(X\in\R^{m\times n}\);
\item there are no matrices \(A\), \(B\) with \(AX'B=X\) for every \(X\in\R^{m\times n}\).
\end{enumerate}
\end{problem}
\begin{solution}
(a)
\begin{itemize}
\item Suppose \(AXB=X'\) for every \(X\in\R^{m\times n}\). As \(AXB\) is defined and has the size \(n\times m\) of \(X'\),
\(A=(a_{ij})\in\R^{n\times m}\) and \(B=(b_{ij})\in\R^{n\times m}\).
\item Let \(1\leqslant k\leqslant m\) and \(1\leqslant s\leqslant n\), and take \(X=E_{ks}=e_ke_s'\in\R^{m\times n}\)
(Proposition~\ref{r1-prop-units}). Then \(AXB=(Ae_k)(e_s'B)=\col_k(A)\row_s(B)\) (Theorem~\ref{r1-thm-rules},
Proposition~\ref{r1-prop-units}(a)), whose \((i,j)\) entry is \(a_{ik}b_{sj}\) (Theorem~\ref{r1-thm-views}(d)); and
\(X'=e_se_k'=E_{sk}\in\R^{n\times m}\) (Theorem~\ref{r1-thm-transpose}), whose \((s,k)\) entry is \(1\) and whose other
entries are \(0\) (Proposition~\ref{r1-prop-units}(b)). So, for \(1\leqslant i,s\leqslant n\) and \(1\leqslant j,k\leqslant m\),
\[
a_{ik}b_{sj}=1\quad\text{if }(i,j)=(s,k),\qquad a_{ik}b_{sj}=0\quad\text{if }(i,j)\ne(s,k)
\]
\item With \(i=s\) and \(j=k\): \(a_{sk}b_{sk}=1\) for all \(s\) and \(k\). So every entry of \(A\) and every entry of \(B\) is
nonzero.
\item As \(mn>1\), \(n\geqslant2\) or \(m\geqslant2\). If \(n\geqslant2\), then \(i=j=k=1\) and \(s=2\) give
\(a_{11}b_{21}=0\) \(\contradiction\)
If \(m\geqslant2\), then \(i=j=s=1\) and \(k=2\) give \(a_{12}b_{11}=0\) \(\contradiction\)
\end{itemize}
(b) If \(AX'B=X\) for every \(X\in\R^{m\times n}\), then transposing gives \(B'XA'=X'\) for every \(X\in\R^{m\times n}\)
(Theorem~\ref{r1-thm-transpose}), which (a) excludes.

By Theorem~\ref{r13-thm-vec}(b) and Definition~\ref{r13-def-commutation}, \(AXB=X'\) for every \(X\) means
\((B'\otimes A)\operatorname{vec}X=K_{mn}\operatorname{vec}X\) for every \(X\), that is, \(B'\otimes A=K_{mn}\), as
\(\operatorname{vec}X\) runs over \(\R^{mn}\) (Proposition~\ref{r1-prop-units}(d)). So (a) says that for \(mn>1\) the
commutation matrix \(K_{mn}\) is not of the form \(C\otimes A\) with \(C\in\R^{m\times n}\) and \(A\in\R^{n\times m}\).
\end{solution}

\begin{problem}[\citeat{searle-khuri}{Exercise~9.1}]\label{r13-ex-jacobian-hessian}
Show that the derivative of \(f(x)=(y_1,y_2)'\), where
\(y_1=6x_1^2x_2+2x_1x_2+x_2^2\) and \(y_2=2x_1^3+x_1^2+2x_1x_2\), with respect to \(x=(x_1,x_2)'\) is the same as the Hessian matrix of
\(y=2x_1^3x_2+x_1^2x_2+x_1x_2^2\).
\end{problem}
\begin{solution}
\begin{itemize}
\item Differentiating the entries,
\[
\mathrm Df(x)=\begin{pmatrix}12x_1x_2+2x_2&6x_1^2+2x_1+2x_2\\6x_1^2+2x_1+2x_2&2x_1\end{pmatrix}
\]
\item For \(y\), \(\partial y/\partial x_1=6x_1^2x_2+2x_1x_2+x_2^2\) and \(\partial y/\partial x_2=2x_1^3+x_1^2+2x_1x_2\), and the second partial derivatives
are \(\partial^2y/\partial x_1^2=12x_1x_2+2x_2\), \(\partial^2y/\partial x_1\partial x_2=\partial^2y/\partial x_2\partial x_1=6x_1^2+2x_1+2x_2\),
\(\partial^2y/\partial x_2^2=2x_1\). So \(\mathrm Hy(x)\) is the matrix above.
\item The reason: \(y_1\) and \(y_2\) are the partial derivatives of \(y\), so \(f=(\mathrm Dy)'\) is the gradient of \(y\), and
\(\mathrm Hy=\mathrm D((\mathrm Dy)')=\mathrm Df\) (Definition~\ref{r13-def-hessian}).
\end{itemize}
\end{solution}

\begin{problem}[\citeat{abadir-magnus}{Exercise~13.24(a)}]\label{r13-ex-residual-maker}
Let \(X\in\R^{n\times k}\) have rank \(k\), and let \(M=I_n-X(X'X)^{-1}X'\) (the residual maker). Find the differential and the derivative of \(M\)
as a function of \(X\).
\end{problem}
\begin{solution}
\begin{itemize}
\item Let \(F=X(X'X)^{-1}\in\R^{n\times k}\) and \(H=I_n-M=X(X'X)^{-1}X'\). By the product rule and
\(\mathrm{d}(X'X)^{-1}=-(X'X)^{-1}\mathrm{d}(X'X)(X'X)^{-1}\) with \(\mathrm{d}(X'X)=(\mathrm{d}X)'X+X'\mathrm{d}X\)
(Theorems~\ref{r13-thm-diff-rules} and~\ref{r13-thm-trace-det}(c)),
\[
\mathrm{d}M=-(\mathrm{d}X)F'+F(\mathrm{d}X)'H+H(\mathrm{d}X)F'-F(\mathrm{d}X)'
\]
\item The first and third terms give \(-(I_n-H)(\mathrm{d}X)F'=-M(\mathrm{d}X)F'\), the second and fourth \(-F(\mathrm{d}X)'(I_n-H)=-F(\mathrm{d}X)'M\):
\(\mathrm{d}M=-M(\mathrm{d}X)(X'X)^{-1}X'-X(X'X)^{-1}(\mathrm{d}X)'M\).
\item By Theorem~\ref{r13-thm-vec}(b), \(\operatorname{vec}\bigl(M(\mathrm{d}X)F'\bigr)=(F\otimes M)\operatorname{vec}\mathrm{d}X\)
and \(\operatorname{vec}\bigl(F(\mathrm{d}X)'M\bigr)=(M\otimes F)\operatorname{vec}(\mathrm{d}X)'=(M\otimes F)K_{nk}\operatorname{vec}\mathrm{d}X\).
\item By Theorem~\ref{r13-thm-commutation}(d), \((M\otimes F)K_{nk}=K_n(F\otimes M)\). Hence
\(\operatorname{vec}\mathrm{d}M=-(I_{n^2}+K_n)(F\otimes M)\operatorname{vec}\mathrm{d}X\), and
\[
\mathrm DM(X)=-(I_{n^2}+K_n)\bigl(X(X'X)^{-1}\otimes M\bigr)
\]
\end{itemize}
\end{solution}

\begin{problem}[\citeat{tu-mc}{Exercise~16}]\label{r13-ex-lagrangian-saddle}
Show that optimizing a function subject to a constraint is not the same as optimizing its Lagrangian: maximize \(xy\) subject to
\(x+y=2\), and examine the Lagrangian \(\psi(x,y)=xy-\lambda(x+y-2)\) at the solution.
\end{problem}
\begin{solution}
\begin{itemize}
\item On the line \(x+y=2\), \(xy=x(2-x)=1-(x-1)^2\leqslant1\), with equality only at \(x=1\): the constrained maximum is attained at
\(x=y=1\).
\item \(\mathrm{d}\psi=(y-\lambda)\,\mathrm{d}x+(x-\lambda)\,\mathrm{d}y\), so \(\mathrm D\psi=0\) and \(x+y=2\) give \(x=y=\lambda=1\): Lagrange's method finds
the point \((1,1)\) with \(\lambda=1\).
\item But \(\mathrm H\psi=\begin{pmatrix}0&1\\1&0\end{pmatrix}\), with the eigenvalues \(1\) and \(-1\) (eigenvectors \((1,1)'\), \((1,-1)'\)), is
indefinite; by Theorem~\ref{r13-thm-extrema}(c), \(\psi\) has neither a maximum nor a minimum at \((1,1)\): it is a saddle point
(\(\psi=1+st\) at \((1+s,1+t)\) for \(\lambda=1\), which is \(>1\) for \(s=t\ne0\) and \(<1\) for \(s=-t\ne0\)). So the solution of the
constrained problem is not a maximum of the Lagrangian, and Theorem~\ref{r13-thm-extrema}(e) does not apply here, since
\(\mathrm H\psi\) is not PSD.
\end{itemize}
\end{solution}

\begin{problem}[\citeat{abadir-magnus}{Exercise~13.56}]\label{r13-ex-constrained-gls}
Let \(X\in\R^{n\times k}\) have rank \(k\), \(y\in\R^n\), \(\Omega\in\R^{n\times n}\) PD, \(R\in\R^{k\times r}\) of rank \(r\), and \(c\in\R^r\). Minimize
\((y-X\beta)'\Omega^{-1}(y-X\beta)\) subject to \(R'\beta=c\).
\end{problem}
\begin{solution}
\begin{itemize}
\item Take \(\psi(\beta)=\frac12(y-X\beta)'\Omega^{-1}(y-X\beta)-l'(R'\beta-c)\) with \(l\in\R^r\). Then
\(\mathrm{d}\psi=-\bigl((y-X\beta)'\Omega^{-1}X+l'R'\bigr)\mathrm{d}\beta\), so the conditions are
\(X'\Omega^{-1}X\beta-Rl=X'\Omega^{-1}y\) and \(R'\beta=c\).
\item \(\Omega^{-1}\) is PD (Proposition~\ref{r8-prop-pd-basic}(c)), so \(W=X'\Omega^{-1}X\) is PD (Theorem~\ref{r8-thm-congruence}(b)), so is
\(W^{-1}\) (the same proposition), and \(R'W^{-1}R\) is PD
(Theorem~\ref{r8-thm-congruence}(b), \(\rank R=r\)), hence invertible. Let \(\hat\beta=W^{-1}X'\Omega^{-1}y\) (the unconstrained solution, the generalized least-squares estimator of Theorem~\ref{r12-thm-gls}).
\item Multiplying the first condition by \(R'W^{-1}\) and using \(R'\beta=c\): \(c-R'W^{-1}Rl=R'\hat\beta\), so
\(\tilde l=(R'W^{-1}R)^{-1}(c-R'\hat\beta)\), and then
\[
\tilde\beta=\hat\beta+W^{-1}R\tilde l=\hat\beta+W^{-1}R(R'W^{-1}R)^{-1}(c-R'\hat\beta)
\]
\item \(\mathrm H\psi=X'\Omega^{-1}X=W\) is PD at every \(\beta\), so \(\tilde\beta\) is the strict minimizer under the constraint
(Theorem~\ref{r13-thm-extrema}(e)).
\end{itemize}
For \(\Omega=I_n\), \(W=X'X\) and \(\tilde\beta\) is the constrained least-squares estimate of Theorem~\ref{r12-thm-cls}, with
\(K=R'\), \(m=c\) and \(G=(X'X)^{-1}\), found there without derivatives.
\end{solution}

\begin{problem}[\citeat{abadir-magnus}{Exercises~8.46, 12.34(a)}]\label{r9-pr-fischer}
\begin{enumerate}
\item Let
\[
M=\begin{pmatrix}A&B\\B'&D\end{pmatrix}
\]
be PD, with square blocks \(A\) and \(D\). Show that \(\det M\leqslant\det A\cdot\det D\) (\textbf{Fischer's inequality}),
with equality if and only if \(B=0\).
\item Let \(M\) be PD and partitioned into \(k\) row and column groups, with square diagonal blocks
\(M_{11},\dots,M_{kk}\). Show that \(\det M\leqslant\det M_{11}\cdots\det M_{kk}\). For blocks of size \(1\) this is
Hadamard's inequality (Theorem~\ref{r8-thm-hadamard}) for PD \(M\).
\end{enumerate}
\end{problem}
\begin{solution}
(a)
\begin{itemize}
\item \(A\) and \(D\) are principal submatrices of \(M\), hence PD (Theorem~\ref{r8-thm-principal-sub}); in particular
\(A\) is invertible and \(\det A>0\) (Theorem~\ref{r8-thm-eigen-criterion}(c)).
\item By Theorem~\ref{r8-thm-schur-pd}, the Schur complement \(S=D-B'A^{-1}B\) is PD and \(\det M=\det A\cdot\det S\).
\item \(A^{-1}\) is PD (Proposition~\ref{r8-prop-pd-basic}(c)), so \(B'A^{-1}B\) is PSD (Theorem~\ref{r8-thm-congruence}(a)).
As \(D=S+B'A^{-1}B\) with \(S\) PD, Theorem~\ref{r8-thm-det-monotone} gives \(\det D\geqslant\det S\), with equality if and
only if \(B'A^{-1}B=0\). Multiplying by \(\det A>0\): \(\det M=\det A\cdot\det S\leqslant\det A\cdot\det D\), with equality if
and only if \(B'A^{-1}B=0\).
\item If \(B=0\), then \(B'A^{-1}B=0\). Conversely, if \(B'A^{-1}B=0\), then \((Bx)'A^{-1}(Bx)=0\) for every \(x\), so \(Bx=0\)
for every \(x\), as \(A^{-1}\) is PD; hence \(B=0\) (Proposition~\ref{r1-prop-units}(a)).
\end{itemize}
(b) Induction on \(k\); for \(k=1\) there is nothing to prove. For \(k\geqslant2\), let \(A\) be the leading block
formed by the first \(k-1\) groups and \(D=M_{kk}\). \(A\) is PD (Theorem~\ref{r8-thm-principal-sub}), with the diagonal
blocks \(M_{11},\dots,M_{k-1,k-1}\). By (a) and the induction hypothesis,
\[
\det M\leqslant\det A\cdot\det M_{kk}\leqslant\det M_{11}\cdots\det M_{k-1,k-1}\cdot\det M_{kk}
\]
\end{solution}

\clearpage
\def\unitnumber{17}
\setcounter{theorem}{0}\setcounter{problem}{0}
\section*{Chapter 17. Generalized linear models}

\subsection*{Exponential families and generalized linear models}

We use the following distributions of a real random variable \(y\). For an integer \(m\geqslant1\) and \(0<\pi<1\), \(y\)
has the \textbf{binomial distribution} \(\mathrm{Bin}(m,\pi)\) if \(\prb(y=k)=\binom mk\pi^k(1-\pi)^{m-k}\) for
\(k=0,1,\dots,m\), where \(\binom mk=\frac{m!}{k!\,(m-k)!}\); for \(m=1\), \(y\) is \textbf{Bernoulli}, with
\(\prb(y=1)=\pi=1-\prb(y=0)\). For \(\mu>0\), \(y\) has the \textbf{Poisson distribution} \(\mathrm{Poisson}(\mu)\) if
\(\prb(y=k)=e^{-\mu}\mu^k/k!\) for every integer \(k\geqslant0\). For \(\sigma^2>0\), \(N(\mu,\sigma^2)\) has the density
\((2\pi\sigma^2)^{-1/2}\exp\bigl(-(y-\mu)^2/(2\sigma^2)\bigr)\) (Theorem~\ref{r14-thm-mvn-density}), mean \(\mu\) and
variance \(\sigma^2\) (Definition~\ref{r14-def-mvn}). We use without proof that \(\expc y=m\pi\) and
\(\var y=m\pi(1-\pi)\) if \(y\sim\mathrm{Bin}(m,\pi)\), and \(\expc y=\var y=\mu\) if \(y\sim\mathrm{Poisson}(\mu)\)~\citeat{casella-berger}{\S3.2}, \citeat{tu-asm}{Propositions~7.2, 8.2}. For
\(0<\pi<1\) the \textbf{odds} are \(\pi/(1-\pi)\) and \(\operatorname{logit}\pi=\log\bigl(\pi/(1-\pi)\bigr)\); the
\textbf{logistic function} \(\operatorname{expit}t=e^t/(1+e^t)\) (\(t\in\R\)) takes values in \((0,1)\), and direct computation
gives \(1-\operatorname{expit}t=1/(1+e^t)\), \(\operatorname{logit}(\operatorname{expit}t)=t\) and
\(\operatorname{expit}(\operatorname{logit}\pi)=\pi\). To the facts of calculus stated after
Definition~\ref{r13-def-derivative} we add: \(e^t\) is \(C^2\) on \(\R\), with derivative \(e^t\).

\begin{definition}[\citeat{tu-asm}{Definitions~8.32, 8.36, 8.9, Proposition~8.35}, \citeat{shao}{\S4.4.2, (4.55)--(4.58)}, \citeat{islr}{(4.42)}]\label{r12-def-glm}
A \textbf{generalized linear model} (GLM) with canonical link consists of a known \(X\in\R^{n\times p}\) of any rank,
with rows \(x_1',\dots,x_n'\), known numbers \(o_1,\dots,o_n\), the \textbf{offsets} (often \(0\)), an unknown
\(\beta\in\R^p\), and independent responses \(y_1,\dots,y_n\). Each \(y_i\) has a density, or probabilities, of the
\textbf{one-parameter exponential family} form
\[
f_i(y_i;\theta_i,\phi)=\exp\Bigl\{\frac{y_i\theta_i-b_i(\theta_i)}\phi+c_i(y_i,\phi)\Bigr\},\qquad\theta_i=\eta_i=o_i+x_i'\beta
\]
with the \textbf{natural parameter} \(\theta_i\) equal to the \textbf{linear predictor} \(\eta_i\), the
\textbf{dispersion} \(\phi>0\), and the functions \(b_i\), \(c_i\) of one of the three families below, the same family for
all \(i\). The number \(\mu_i=b_i'(\theta_i)\) is the mean of \(y_i\) (shown after the definition), and the inverse of
\(b_i'\), which maps \(\mu_i\) to \(\eta_i\), is the \textbf{canonical link}: \(\mu\mapsto\mu\),
\(\mu\mapsto\operatorname{logit}(\mu/m_i)\) or \(\mu\mapsto\log\mu\) in the three families below.
\begin{itemize}
\item \textbf{Normal linear model}: \(y_i\sim N(\mu_i,\sigma^2)\), \(o_i=0\), \(\theta_i=\mu_i=x_i'\beta\),
\(\phi=\sigma^2\), \(b_i(\theta)=\theta^2/2\), \(c_i(y,\phi)=-y^2/(2\phi)-\frac12\log(2\pi\phi)\); this is the normal linear
model of Definition~\ref{r12-def-inference}, with \(p\) columns.
\item \textbf{Logistic regression}: \(y_i\sim\mathrm{Bin}(m_i,\pi_i)\) with known \(m_i\), \(o_i=0\),
\(\theta_i=\operatorname{logit}\pi_i=x_i'\beta\), \(\phi=1\), \(b_i(\theta)=m_i\log(1+e^\theta)\),
\(c_i(k,\phi)=\log\binom{m_i}k\). So \(\pi_i=\operatorname{expit}(x_i'\beta)\) and \(\mu_i=m_i\pi_i\); for \(m_i=1\) the
responses are binary.
\item \textbf{Poisson regression with exposure}: \(y_i\sim\mathrm{Poisson}(\mu_i)\) with \(\mu_i=t_ie^{x_i'\beta}\), that is,
\(\log\mu_i=\log t_i+x_i'\beta\), for known \textbf{exposures} \(t_i>0\) (offsets \(o_i=\log t_i\)), such as policy-years
in insurance; \(\theta_i=\log\mu_i\), \(\phi=1\), \(b_i(\theta)=e^\theta\), \(c_i(k,\phi)=-\log k!\).
\end{itemize}
The likelihood and the maximum likelihood estimates (MLE) are those of Definition~\ref{r12-def-inference}: the likelihood
is the density \(f(y;\vartheta)\) of \(y=(y_1,\dots,y_n)'\) at the observed \(y\), as a function of the parameters
\(\vartheta\) (here \(\beta\) and \(\phi\)). For a discrete \(y\),
\(f(y;\vartheta)\) is instead the probability that \(y\) takes the observed value. By independence (the paragraphs before
Definition~\ref{r14-def-mvn} and before Theorem~\ref{r14-thm-mvn-density}), it is the product
\(f_1(y_1;\theta_1,\phi)\cdots f_n(y_n;\theta_n,\phi)\).
\end{definition}

The three families have this form~\citeat{tu-asm}{Proposition~8.35}, \citeat{casella-berger}{Examples~3.4.1, 3.4.4}. For \(\mathrm{Bin}(m,\pi)\) with \(\theta=\operatorname{logit}\pi\),
\(\pi=\operatorname{expit}\theta\) and \(1-\pi=1/(1+e^\theta)\), so
\(\pi^k(1-\pi)^{m-k}=\bigl(\pi/(1-\pi)\bigr)^k(1-\pi)^m=e^{k\theta}(1+e^\theta)^{-m}\) and
\(\log\prb(y=k)=k\theta-m\log(1+e^\theta)+\log\binom mk\). For \(\mathrm{Poisson}(\mu)\) with \(\theta=\log\mu\),
\(\log(e^{-\mu}\mu^k/k!)=k\theta-e^\theta-\log k!\). For \(N(\mu,\sigma^2)\),
\(-(y-\mu)^2/(2\sigma^2)=(y\mu-\mu^2/2)/\sigma^2-y^2/(2\sigma^2)\). The functions \(b\) are \(C^2\) on \(\R\), with
\[
\begin{aligned}
&\text{binomial:}&&b'(\theta)=\frac{me^\theta}{1+e^\theta}=m\pi,&&b''(\theta)=\frac{me^\theta}{(1+e^\theta)^2}=m\pi(1-\pi)\\
&\text{Poisson:}&&b'(\theta)=e^\theta=\mu,&&b''(\theta)=e^\theta=\mu\\
&\text{normal:}&&b'(\theta)=\theta=\mu,&&b''(\theta)=1
\end{aligned}
\]
So in each family \(\expc y=b'(\theta)\) and \(\var y=\phi\,b''(\theta)>0\), by the facts in the first paragraph.

\needspace{14\baselineskip}
\begin{theorem}[\citeat{tu-asm}{Proposition~8.37, Theorem~8.40}, \citeat{shao}{\S4.4.2}]\label{r12-thm-glm}
In a GLM (Definition~\ref{r12-def-glm}), let \(\theta=o+X\beta\) with \(o=(o_1,\dots,o_n)'\), and let
\[
\mu=\begin{pmatrix}b_1'(\theta_1)\\\vdots\\b_n'(\theta_n)\end{pmatrix},\qquad
W=\diag\bigl(b_1''(\theta_1),\dots,b_n''(\theta_n)\bigr)
\]
both evaluated at \(\beta\). The log-likelihood, as a function of \(\beta\in\R^p\) with \(\phi\) fixed, is
\[
\ell(\beta)=\frac1\phi\sum_{i=1}^n\bigl(y_i\theta_i-b_i(\theta_i)\bigr)+\sum_{i=1}^nc_i(y_i,\phi)
\]
Here an \textbf{MLE} of \(\beta\) (for the given \(\phi\)) is a maximizer of \(\ell\) on \(\R^p\); for the binomial and
Poisson families \(\phi=1\), and for the normal family the maximizers do not depend on \(\phi\) and, if
\(\mathrm{SSE}>0\), they are the first components \(\hat\beta\) of the MLEs \((\hat\beta,\mathrm{SSE}/n)\) of
Theorem~\ref{r12-thm-sampling}(d).
\begin{enumerate}
\item \(\ell\) is \(C^2\) on \(\R^p\), \(\mathrm D\ell(\beta)=(y-\mu)'X/\phi\), so that the gradient, the \textbf{score}, is
\(X'(y-\mu)/\phi\), and \(\mathrm H\ell(\beta)=-X'WX/\phi\). The matrix \(X'WX\) is PSD for every \(\beta\), and PD if
\(\rank X=p\).
\item \(\hat\beta\) is an MLE if and only if it solves the \textbf{likelihood equations}
\[
X'(y-\hat\mu)=0,\qquad\hat\mu_i=b_i'(o_i+x_i'\hat\beta)
\]
which do not involve \(\phi\). If \(\rank X=p\), there is at most one MLE.
\item Let \(\rank X=p\). The Newton--Raphson step for the stationary points of \(\ell\) from \(\beta\)
(Example~\ref{r13-exm-newton}) is
\[
\beta^{\mathrm{new}}=\beta-\mathrm H\ell(\beta)^{-1}\bigl(\mathrm D\ell(\beta)\bigr)'=(X'WX)^{-1}X'Wz,\qquad
z=X\beta+W^{-1}(y-\mu)
\]
the weighted least-squares estimate (Proposition~\ref{r12-prop-wls}) for the \textbf{working response} \(z\) and the weight
matrix \(W\); it is the unique maximizer of the second-order Taylor polynomial of \(\ell\) at \(\beta\). Repeating the step
is \textbf{iteratively reweighted least squares} (IRLS).
\end{enumerate}
\end{theorem}
\begin{proof}
By Definition~\ref{r12-def-glm} the likelihood is \(\prod_if_i(y_i;\theta_i,\phi)\), and its logarithm is \(\ell\).

(a)
\begin{itemize}
\item Each \(b_i\) is \(C^2\) on \(\R\) (the paragraph after Definition~\ref{r12-def-glm}), and \(\theta_i=o_i+x_i'\beta\) is a
polynomial in \(\beta\); so \(\ell\) is \(C^2\) on \(\R^p\) (rules and chain rule of calculus, after
Definition~\ref{r13-def-derivative}).
\item By Theorems~\ref{r13-thm-forms}(a) and~\ref{r13-thm-diff-rules}(a),(c), \(\mathrm{d}\theta_i=x_i'\,\mathrm{d}\beta\) and
\(\mathrm{d}\,b_i(\theta_i)=b_i'(\theta_i)\,x_i'\,\mathrm{d}\beta=\mu_ix_i'\,\mathrm{d}\beta\). By Theorem~\ref{r1-thm-views}(c), the row
\((y-\mu)'X\) is \(\sum_i(y_i-\mu_i)x_i'\), so
\[
\mathrm{d}\ell=\frac1\phi\sum_{i=1}^n(y_i-\mu_i)\,x_i'\,\mathrm{d}\beta=\frac1\phi(y-\mu)'X\,\mathrm{d}\beta
\]
and \(\mathrm D\ell(\beta)=(y-\mu)'X/\phi\) (Theorem~\ref{r13-thm-identification}(a)); the gradient is its transpose
\(X'(y-\mu)/\phi\).
\item With \(\mathrm{d}\beta\) fixed, \(\mathrm{d}\mu_i=b_i''(\theta_i)\,x_i'\,\mathrm{d}\beta\), a number equal to its transpose
\(b_i''(\theta_i)(\mathrm{d}\beta)'x_i\). By Theorem~\ref{r1-thm-views}(d), \(X'WX=\sum_ix_i\,b_i''(\theta_i)x_i'\) (the
columns of \(X'\) are the \(x_i\), the rows of \(WX\) are the \(b_i''(\theta_i)x_i'\)). So
\[
\mathrm{d}^2\ell=-\frac1\phi\sum_{i=1}^n(\mathrm{d}\mu_i)\,x_i'\,\mathrm{d}\beta
=-\frac1\phi(\mathrm{d}\beta)'\Bigl(\sum_{i=1}^nb_i''(\theta_i)x_ix_i'\Bigr)\mathrm{d}\beta=-\frac1\phi(\mathrm{d}\beta)'X'WX\,\mathrm{d}\beta
\]
\(X'WX\) is symmetric, so \(\mathrm H\ell(\beta)=-X'WX/\phi\)
(Theorem~\ref{r13-thm-identification2}(a)).
\item \(W\) is diagonal with positive diagonal entries (the paragraph after Definition~\ref{r12-def-glm}), hence PD
(Proposition~\ref{r8-prop-pd-basic}(a)). By Theorem~\ref{r8-thm-congruence}(a),(b) with \(A=W\) and \(P=X\), \(X'WX\) is
PSD, and PD if \(\rank X=p\).
\end{itemize}
(b)
\begin{itemize}
\item \(\R^p\) is open and convex. If \(\hat\beta\) is an MLE, \(\ell\) has a global, hence a local, maximum at
\(\hat\beta\), so \(\mathrm D\ell(\hat\beta)=0\) (Theorem~\ref{r13-thm-extrema}(a)); by (a) and \(\phi>0\), this is
\(X'(y-\hat\mu)=0\).
\item Conversely, let \(X'(y-\hat\mu)=0\). The function \(-\ell\) is \(C^2\), \(\mathrm D(-\ell)(\hat\beta)=0\), and its
Hessian \(X'WX/\phi\) is PSD at every \(\beta\) by (a). By Theorem~\ref{r13-thm-extrema}(d), \(-\ell\) has a global
minimum at \(\hat\beta\): \(\ell(\beta)\leqslant\ell(\hat\beta)\) for every \(\beta\), and \(\hat\beta\) is an MLE.
\item If \(\rank X=p\), the Hessian of \(-\ell\) is PD at every \(\beta\), and the global minimum is strict
(Theorem~\ref{r13-thm-extrema}(d)): \(\ell(\beta)<\ell(\hat\beta)\) for \(\beta\ne\hat\beta\), so no other \(\beta\) is an
MLE.
\end{itemize}
(c)
\begin{itemize}
\item By (a), \(X'WX\) is PD, hence invertible (Theorem~\ref{r8-thm-eigen-criterion}(c)), and so is
\(\mathrm H\ell(\beta)=-X'WX/\phi\). The step of Example~\ref{r13-exm-newton} for the function \(\ell\) is
\(\beta-(-X'WX/\phi)^{-1}X'(y-\mu)/\phi=\beta+(X'WX)^{-1}X'(y-\mu)\), in which \(\phi\) cancels.
\item As \(\beta=(X'WX)^{-1}X'WX\beta\), this is \((X'WX)^{-1}X'W\bigl(X\beta+W^{-1}(y-\mu)\bigr)=(X'WX)^{-1}X'Wz\), and by
Proposition~\ref{r12-prop-wls} (\(W\) PD, \(\rank X=p\)) it is the unique minimizer of \((z-Xb)'W(z-Xb)\).
\item The Hessian \(X'WX/\phi\) of \(-\ell\) at \(\beta\) is PD, so by the paragraph after Example~\ref{r13-exm-newton},
applied to \(-\ell\), \(\beta^{\mathrm{new}}\) is the unique minimizer of the second-order Taylor polynomial \(q\) of \(-\ell\)
at \(\beta\); and \(-q\) is the second-order Taylor polynomial of \(\ell\).\qedhere
\end{itemize}
\end{proof}

The theorem does not assert that an MLE exists. For instance, in a logistic regression whose design has the column
\(\one_n\) and whose responses are all \(0\), the likelihood equation \(\one_n'(y-\hat\mu)=0\) would require
\(\sum_im_i\operatorname{expit}(x_i'\hat\beta)=0\), which is impossible, as \(\operatorname{expit}>0\). Nor does the theorem
assert that the IRLS iterates converge. For the normal family, \(b_i''=1\), so \(W=I_n\) and \(\mu=X\beta\): the likelihood
equations \(X'(y-X\hat\beta)=0\) are the normal equations (Definition~\ref{r11-def-model}), as in
Theorem~\ref{r12-thm-sampling}(d); and \(z=X\beta+(y-X\beta)=y\), so that for \(\rank X=p\) one IRLS step from any \(\beta\)
gives the least-squares estimate \((X'X)^{-1}X'y\), as in Example~\ref{r13-exm-newton}~\citeat{shao}{Example~4.36}.

\needspace{18\baselineskip}
\subsection*{The rank-general case and weighted projections}

\begin{proposition}\label{r12-prop-glm-fitted}
In a GLM (Definition~\ref{r12-def-glm}), let \(\rank X=r\geqslant1\), and let \(X=X_1C\) be a full-rank factorization
(Theorem~\ref{thm:r4-full-rank-factorization}): \(X_1\in\R^{n\times r}\) has full column rank and \(C\in\R^{r\times p}\) full
row rank. Let \(\ell_1(\gamma)\) (\(\gamma\in\R^r\)) be the log-likelihood of the GLM with the same responses, family,
offsets and \(\phi\) and the design \(X_1\).
\begin{enumerate}
\item \(\ell(\beta)=\ell_1(C\beta)\), and \(\beta\) is an MLE if and only if \(C\beta\) is an MLE of the design \(X_1\), of
which there is at most one. If \(\beta^0\) is an MLE, the MLEs are the vectors \(\beta^0+u\) with \(u\in\mathcal N(X)\), and
all of them give the same \(X\hat\beta\), the same linear predictors \(\hat\eta=o+X\hat\beta\) and the same fitted means
\(\hat\mu\).
\item If \(a\in\mathcal C(X')\), then \(a'\hat\beta\) is the same number for all MLEs \(\hat\beta\).
\item At any \(\beta\), with \(W\) and \(z\) as in Theorem~\ref{r12-thm-glm}, the system \(X'WXb=X'Wz\) (unknown
\(b\in\R^p\)) is consistent, and its solutions are the minimizers of \((z-Xb)'W(z-Xb)\); one of them is \(GX'Wz\) for any
g-inverse \(G\) of \(X'WX\). All of them give the same \(Xb=X_1\gamma^{\mathrm{new}}\), where \(\gamma^{\mathrm{new}}\) is the
IRLS step of Theorem~\ref{r12-thm-glm}(c) for the design \(X_1\) from \(\gamma=C\beta\).
\end{enumerate}
\end{proposition}
\begin{proof}
(a)
\begin{itemize}
\item \(\theta=o+X\beta=o+X_1(C\beta)\), and \(\ell\) depends on \(\beta\) only through \(\theta\)
(Theorem~\ref{r12-thm-glm}); so \(\ell(\beta)=\ell_1(C\beta)\). As \(C\) has full row rank, \(\mathcal C(C)=\R^r\)
(Theorem~\ref{thm:r4-rank-rules}(c)): every \(\gamma\in\R^r\) is \(C\beta'\) for some \(\beta'\in\R^p\).
\item If \(C\beta\) maximizes \(\ell_1\), then \(\ell(\beta')=\ell_1(C\beta')\leqslant\ell_1(C\beta)=\ell(\beta)\) for every
\(\beta'\): \(\beta\) is an MLE. Conversely, if \(\beta\) is an MLE and \(\gamma=C\beta'\in\R^r\), then
\(\ell_1(\gamma)=\ell(\beta')\leqslant\ell(\beta)=\ell_1(C\beta)\): \(C\beta\) maximizes \(\ell_1\).
\item \(X_1\) has full column rank \(r\), so \(\ell_1\) has at most one maximizer \(\hat\gamma\)
(Theorem~\ref{r12-thm-glm}(b)). If \(\beta^0\) is an MLE, \(\hat\gamma=C\beta^0\), and the MLEs are the solutions of
\(C\beta=C\beta^0\), that is, the vectors \(\beta^0+u\) with \(u\in\mathcal N(C)=\mathcal N(X)\)
(Theorem~\ref{thm:r4-full-rank-factorization}(b)). Each of them gives \(X\beta=X_1C\beta=X_1\hat\gamma\), hence the same
\(\hat\eta=o+X\hat\beta\) and \(\hat\mu_i=b_i'(\hat\eta_i)\).
\end{itemize}
(b) \(a=X'c\) for some \(c\in\R^n\), so \(a'\hat\beta=c'X\hat\beta=c'X_1\hat\gamma\) for every MLE \(\hat\beta\), by (a).

(c)
\begin{itemize}
\item By Proposition~\ref{r12-prop-wls} (\(W\) PD, \(X\) of any rank), the minimizers of \(Q(b)=(z-Xb)'W(z-Xb)\) are the
solutions of \(X'WXb=X'Wz\). \(W\), \(\mu\) and \(z\) depend on \(\beta\) only through \(X\beta=X_1\gamma\), so they are also
those of the design \(X_1\) at \(\gamma=C\beta\); and \(Q(b)=Q_1(Cb)\) with \(Q_1(h)=(z-X_1h)'W(z-X_1h)\) (\(h\in\R^r\)).
\item By Proposition~\ref{r12-prop-wls} for \(X_1\) (rank \(r\)), \(Q_1\) has exactly one minimizer, which is
\(\gamma^{\mathrm{new}}=(X_1'WX_1)^{-1}X_1'Wz\), the step of Theorem~\ref{r12-thm-glm}(c) for \(X_1\). As in (a), \(b\) minimizes
\(Q\) if and only if \(Cb\) minimizes \(Q_1\), that is, \(Cb=\gamma^{\mathrm{new}}\); such \(b\) exist (\(\mathcal C(C)=\R^r\)), so
the system is consistent, and \(Xb=X_1Cb=X_1\gamma^{\mathrm{new}}\) for all of them.
\item As the system is consistent, \(X'Wz\in\mathcal C(X'WX)\), and \(GX'Wz\) is a solution
(Theorem~\ref{r10-thm-ginv-solves}(a)).\qedhere
\end{itemize}
\end{proof}

\needspace{14\baselineskip}
\begin{proposition}[\citeat{harville}{Theorems~14.12.11, 14.12.3, Corollary~14.12.4}]\label{r12-prop-weighted-projection}
Let \(X\in\R^{n\times p}\), let \(W\in\R^{n\times n}\) be PD, let \(G\) be a g-inverse of \(X'WX\), and put
\[
P_{X,W}=XGX'W=X(X'WX)^-X'W
\]
\begin{enumerate}
\item Let \(T=W^{1/2}\) (Theorem~\ref{r8-thm-sqrt}), and let \(P=TXG(TX)'\), the orthogonal projection matrix onto
\(\mathcal C(TX)\) of Theorem~\ref{r10-thm-ata} for \(A=TX\). Then \(P_{X,W}=T^{-1}PT\), and \(P_{X,W}\) does not depend on the
choice of \(G\).
\item \(P_{X,W}X=X\), \(P_{X,W}^2=P_{X,W}\), and \(WP_{X,W}\) is symmetric.
\item \(\mathcal C(P_{X,W})=\mathcal C(X)\), \(\rank P_{X,W}=\rank X\), and \(\rank(I_n-P_{X,W})=n-\rank X\).
\item For every \(z\in\R^n\), \(P_{X,W}z=Xb\) for every minimizer \(b\) of \((z-Xb)'W(z-Xb)\). In particular, in a GLM the
IRLS step of Proposition~\ref{r12-prop-glm-fitted}(c) gives \(Xb=P_{X,W}z\).
\end{enumerate}
\end{proposition}
\begin{proof}
\begin{itemize}
\item \(T\) is PD, hence symmetric and invertible, and \(T^2=W\) (Theorem~\ref{r8-thm-sqrt}(a),(c)); so
\((TX)'(TX)=X'T^2X=X'WX\), and \(G\) is a g-inverse of \((TX)'(TX)\).
\item By Theorem~\ref{r10-thm-ata} for \(A=TX\): \(P=TXGX'T\) does not depend on \(G\), \(P'=P=P^2\), \(P(TX)=TX\), and
\(I_n-P\) has rank \(n-\rank(TX)\), where \(\rank(TX)=\rank X\) (Theorem~\ref{thm:r4-rank-product}(d)).
\end{itemize}
(a) \(T^{-1}PT=T^{-1}TXGX'TT=XGX'W=P_{X,W}\). As \(P\) does not depend on \(G\), neither does \(P_{X,W}\).

(b) \(P_{X,W}X=T^{-1}P(TX)=T^{-1}TX=X\); \(P_{X,W}^2=T^{-1}P^2T=T^{-1}PT=P_{X,W}\); and \(WP_{X,W}=T^2T^{-1}PT=TPT\), with
\((TPT)'=T'P'T'=TPT\).

(c) \(\mathcal C(P_{X,W})=\mathcal C\bigl(X(GX'W)\bigr)\subseteq\mathcal C(X)\) and
\(\mathcal C(X)=\mathcal C(P_{X,W}X)\subseteq\mathcal C(P_{X,W})\) (Theorem~\ref{thm:r4-rank-product}(a)); so the two are equal
and have the same dimension. \(I_n-P_{X,W}=T^{-1}(I_n-P)T\) has the rank \(n-\rank X\) of \(I_n-P\)
(Theorem~\ref{thm:r4-rank-product}(d)).

(d) By Proposition~\ref{r12-prop-wls}, the minimizers \(b\) are the solutions of \(X'WXb=X'Wz\), that is, of
\((TX)'(TX)b=(TX)'(Tz)\), the normal equations of the design \(TX\) and the response \(Tz\). By
Theorem~\ref{r10-thm-normal-equations}(c), \(G(TX)'Tz=GX'Wz\) is one of them, and all of them give the same \(TXb\), hence
the same \(Xb\) (\(T\) is invertible), which is \(XGX'Wz=P_{X,W}z\). In a GLM, the solutions of
Proposition~\ref{r12-prop-glm-fitted}(c) are these minimizers.
\end{proof}
For \(W=I_n\), \(P_{X,W}=X(X'X)^-X'\) is the hat matrix \(H\) of Theorem~\ref{r10-thm-PX}. For other \(W\) it is idempotent
but in general not symmetric. By (d), \(P_{X,W}z\) is the point \(Xb\) of \(\mathcal C(X)\) nearest to \(z\) when the length of
a vector \(v\) is measured by \(\lVert Tv\rVert=(v'Wv)^{1/2}\); so each IRLS step replaces the working response \(z\) by this
projection \(P_{X,W}z\) of it.

\subsection*{Grouped data and examples}

In logistic regression the odds are \(\pi_i/(1-\pi_i)=e^{x_i'\beta}\): increasing \(x_{ij}\) by \(1\), the other entries of
\(x_i\) fixed, multiplies the odds by \(e^{\beta_j}\), the \textbf{odds ratio}. In Poisson regression with exposure, the
expected count per unit of exposure is \(\mu_i/t_i=e^{x_i'\beta}\); increasing \(x_{ij}\) by \(1\) multiplies it by
\(e^{\beta_j}\), the \textbf{rate ratio}~\citeat{tu-asm}{Theorem~7.11, Proposition~8.10}, \citeat{islr}{\S4.3, \S4.6.2}. When column \(j\) of \(X\) is the indicator of a group, \(e^{\beta_j}\) compares that
group with the observations that have \(x_{ij}=0\) and the same other entries: in the baseline coding
\([\one_n\ d_2\ \cdots\ d_g]\) of Theorem~\ref{r11-thm-oneway}, the observations of group \(1\); in
\([\one_n\ d_1\ \cdots\ d_g]\) there are no such observations, as every row has exactly one indicator entry \(1\).

\needspace{16\baselineskip}
\begin{proposition}[\citeat{tu-asm}{Corollaries~8.12, 7.16, 8.14}]\label{r12-prop-glm-groups}
In a GLM (Definition~\ref{r12-def-glm}), let the \(n\) observations fall into \(g\) nonempty groups, and let
\(d_l\in\R^n\) be the \textbf{indicator} of group \(l\), with entry \(1\) in the rows of group \(l\) and \(0\) elsewhere (as in
Theorem~\ref{r11-thm-oneway}); \(\sum_{i\in l}\) is the sum over the observations \(i\) of group \(l\).
\begin{enumerate}
\item If \(\one_n\in\mathcal C(X)\), every MLE satisfies \(\sum_{i=1}^n(y_i-\hat\mu_i)=0\).
\item If \(d_l\in\mathcal C(X)\), every MLE satisfies \(\sum_{i\in l}\hat\mu_i=\sum_{i\in l}y_i\).
\item Let \(\mathcal C(X)=\mathcal C(D)\) with \(D=[d_1\ \cdots\ d_g]\), for instance \(X=[\one_n\ d_1\ \cdots\ d_g]\),
\([\one_n\ d_2\ \cdots\ d_g]\) or \(D\) (Theorem~\ref{r11-thm-oneway}).
\begin{itemize}
\item In logistic regression, let \(s_l=\sum_{i\in l}y_i\) and \(n_l=\sum_{i\in l}m_i\). An MLE exists if and only if
\(0<s_l<n_l\) for every \(l\); the MLEs are then the \(\beta\) with \(x_i'\beta=\operatorname{logit}(s_l/n_l)\) for every \(i\)
of every group \(l\), and \(\hat\pi_i=s_l/n_l\).
\item In Poisson regression with exposures, let \(Y_l=\sum_{i\in l}y_i\) and \(\tau_l=\sum_{i\in l}t_i\). An MLE exists if and
only if \(Y_l>0\) for every \(l\); the MLEs are then the \(\beta\) with \(x_i'\beta=\log(Y_l/\tau_l)\) for every \(i\) of every
group \(l\), and \(\hat\mu_i=t_iY_l/\tau_l\).
\end{itemize}
\end{enumerate}
\end{proposition}
\begin{proof}
(a) If \(\one_n=Xa\), then by Theorem~\ref{r12-thm-glm}(b),
\(\sum_i(y_i-\hat\mu_i)=\one_n'(y-\hat\mu)=a'X'(y-\hat\mu)=0\).

(b) In the same way, \(d_l=Xa_l\) gives \(\sum_{i\in l}(y_i-\hat\mu_i)=d_l'(y-\hat\mu)=a_l'X'(y-\hat\mu)=0\).

(c)
\begin{itemize}
\item Every column of \(X\) lies in \(\mathcal C(D)\), so \(X=DF\) for some \(F\in\R^{g\times p}\)
(Theorem~\ref{r1-thm-views}(a),(b)); row \(i\) of \(D\) is \(e_l'\) for \(i\) in group \(l\), so \(x_i'\beta=(F\beta)_l\) is the
same number \(c_l\) for all \(i\) of group \(l\), and \(X\beta=Dc\) with \(c=F\beta\).
\item Logistic. \(\mu_i=m_i\operatorname{expit}(c_l)\) for \(i\) in group \(l\). If \(\beta\) is an MLE, (b) gives
\(n_l\operatorname{expit}(c_l)=s_l\); as \(0<\operatorname{expit}<1\), this has a solution only if \(0<s_l<n_l\), and then
\(c_l=\operatorname{logit}(s_l/n_l)\). Conversely, let \(0<s_l<n_l\) for every \(l\), let
\(c_l=\operatorname{logit}(s_l/n_l)\), and let \(\beta\) satisfy \(X\beta=Dc\) (there is such a \(\beta\), as
\(Dc\in\mathcal C(D)=\mathcal C(X)\)). Then \(\hat\pi_i=s_l/n_l\), the entries of \(D'(y-\hat\mu)\) are
\(s_l-n_l\,s_l/n_l=0\), and \(X'(y-\hat\mu)=F'D'(y-\hat\mu)=0\): \(\beta\) is an MLE (Theorem~\ref{r12-thm-glm}(b)).
\item Poisson. \(\mu_i=t_ie^{c_l}\) for \(i\) in group \(l\), and (b) reads \(\tau_le^{c_l}=Y_l\), which has a solution if and
only if \(Y_l>0\), namely \(c_l=\log(Y_l/\tau_l)\); the converse is as in the logistic case.\qedhere
\end{itemize}
\end{proof}

\needspace{20\baselineskip}
\begin{example}\label{r12-exm-glm-factorial}
Four batches of \(8\) insects are exposed to an insecticide. Batch \(i\) has the dose \(a_i\in\{0,1\}\) (low, high) and the
strain \(b_i\in\{0,1\}\), with \((a_i,b_i)=(0,0),(1,0),(0,1),(1,1)\), and \(y=(2,4,4,6)'\) of the insects die. Fit the
logistic regression \(y_i\sim\mathrm{Bin}(8,\pi_i)\), \(\operatorname{logit}\pi_i=\beta_0+\beta_1a_i+\beta_2b_i\): carry
out one IRLS step from \(\beta=0\), and find the MLE.

\textit{Solution.} With \(X\) and \(X'X\) below, \(\det(X'X)=4\), and
\((X'X)^{-1}=\frac14\operatorname{adj}(X'X)\) (Theorem~\ref{r5d-thm-adjugate}):
\[
X=\begin{pmatrix}1&0&0\\1&1&0\\1&0&1\\1&1&1\end{pmatrix},\qquad
X'X=\begin{pmatrix}4&2&2\\2&2&1\\2&1&2\end{pmatrix},\qquad
(X'X)^{-1}=\frac14\begin{pmatrix}3&-2&-2\\-2&4&0\\-2&0&4\end{pmatrix}
\]
so \(\rank X=3\) (Theorem~\ref{r10-thm-normal-equations}(b)). By the paragraph after Definition~\ref{r12-def-glm},
\(\mu_i=8\pi_i\) and
\(W=\diag\bigl(8\pi_i(1-\pi_i)\bigr)\), and \(\phi=1\).
\begin{itemize}
\item At \(\beta=0\): \(\pi_i=\frac12\), \(\mu=4\one_4\), \(W=2I_4\), \(y-\mu=(-2,0,0,2)'\) and
\(z=X\beta+W^{-1}(y-\mu)=(-1,0,0,1)'\). The score \(X'(y-\mu)=(0,2,2)'\) is not \(0\), \(X'Wz=X'(y-\mu)\), and
\[
X'WX=2X'X=\begin{pmatrix}8&4&4\\4&4&2\\4&2&4\end{pmatrix},\qquad
\beta^{(1)}=(X'WX)^{-1}X'Wz=\frac18\begin{pmatrix}3&-2&-2\\-2&4&0\\-2&0&4\end{pmatrix}\begin{pmatrix}0\\2\\2\end{pmatrix}
=\begin{pmatrix}-1\\1\\1\end{pmatrix}
\]
by Theorem~\ref{r12-thm-glm}(c).
\item The likelihood equations \(X'(y-\hat\mu)=0\) (Theorem~\ref{r12-thm-glm}(b)) hold if \(\hat\mu=y\), that is,
\(\hat\pi=\frac18y=(\frac14,\frac12,\frac12,\frac34)'\), with \(\operatorname{logit}\hat\pi=(-\log3,0,0,\log3)'\). The
equations \(\beta_0=-\log3\), \(\beta_0+\beta_1=0\), \(\beta_0+\beta_2=0\), \(\beta_0+\beta_1+\beta_2=\log3\) are solved by
\(\hat\beta=(-\log3,\log3,\log3)'\). As \(\rank X=3\), it is the only MLE. For each strain, the high dose multiplies the
odds of death by \(e^{\hat\beta_1}=3\), and for each dose, the strain \(b=1\) multiplies them by \(e^{\hat\beta_2}=3\).
\item At \(\hat\beta\), \(\hat W=\diag\bigl(8\hat\pi_i(1-\hat\pi_i)\bigr)=\diag(\frac32,2,2,\frac32)\), and
\(\mathrm H\ell(\hat\beta)=-X'\hat WX\) with
\[
X'\hat WX=\begin{pmatrix}7&\frac72&\frac72\\\frac72&\frac72&\frac32\\\frac72&\frac32&\frac72\end{pmatrix}
\]
which is PD (Theorem~\ref{r12-thm-glm}(a)).
\end{itemize}
\end{example}

\needspace{18\baselineskip}
\begin{example}\label{r12-exm-glm-oneway}
Three groups of \(36\) patients receive the treatments \(1\), \(2\), \(3\), and \(12\), \(24\), \(18\) of them recover. With
one row for each group, \(y_l\sim\mathrm{Bin}(36,\pi_l)\) (\(l=1,2,3\)), so \(n=3\) and \(d_l=e_l\). Fit the one-way
logistic model \(\operatorname{logit}\pi_l=\beta_0+\alpha_l\), that is, \(\beta=(\beta_0,\alpha_1,\alpha_2,\alpha_3)'\) and
\[
X=[\one_3\ d_1\ d_2\ d_3]=\begin{pmatrix}1&1&0&0\\1&0&1&0\\1&0&0&1\end{pmatrix}
\]
(Theorem~\ref{r11-thm-oneway}): find the MLEs, the fitted probabilities and the odds ratios of the treatments.

\textit{Solution.}
\begin{itemize}
\item \(\mathcal C(X)=\mathcal C(D)\) with \(D=[d_1\ d_2\ d_3]=I_3\) (Theorem~\ref{r11-thm-oneway}(c)), and
\(0<s_l<36\) for each group. By Proposition~\ref{r12-prop-glm-groups}(c), the MLEs exist, \(\hat\pi=(\frac13,\frac23,\frac12)'\),
and they are the \(\beta\) with
\[
X\beta=\begin{pmatrix}\operatorname{logit}\frac13\\\operatorname{logit}\frac23\\\operatorname{logit}\frac12\end{pmatrix}
=\begin{pmatrix}-\log2\\\log2\\0\end{pmatrix}
\]
\item \(\rank X=3<4\) and \(\mathcal N(X)=\Span\{(-1,1,1,1)'\}\) (Theorem~\ref{r11-thm-oneway}(b)). One MLE is
\(\beta^0=(0,-\log2,\log2,0)'\), so by Proposition~\ref{r12-prop-glm-fitted}(a) the MLEs are
\(\beta^0+t(-1,1,1,1)'=(-t,\ t-\log2,\ t+\log2,\ t)'\) (\(t\in\R\)), all with the same \(\hat\pi\). Their first entry
\(-t\) takes every real value: the data do not determine \(\beta_0\).
\item \(\alpha_2-\alpha_1=a'\beta\) with \(a=(0,-1,1,0)'=X'(-1,1,0)'\in\mathcal C(X')\), and likewise
\(\alpha_3-\alpha_1\), with \(a=X'(-1,0,1)'\). By Proposition~\ref{r12-prop-glm-fitted}(b) they have the same value for all
MLEs: \(\hat\alpha_2-\hat\alpha_1=2\log2\) and \(\hat\alpha_3-\hat\alpha_1=\log2\). So the estimated odds ratios of
treatments \(2\) and \(3\) to treatment \(1\) are \(4\) and \(2\): the estimated odds of recovery are \(\frac{24}{12}=2\),
\(\frac{12}{24}=\frac12\) and \(1\) for treatments \(2\), \(1\), \(3\).
\end{itemize}
\end{example}

We state without proof, as it needs limit theorems that are not in this book: in a GLM with \(\rank X=p\) and \(\phi\) known
(\(\phi=1\) for the binomial and Poisson families), under conditions on \(X\) (for instance, bounded rows \(x_i\) and
\(X'WX/n\) converging to a PD matrix), for large \(n\) an MLE exists with probability close to \(1\), and \(\hat\beta\) is
approximately \(N_p\bigl(\beta,\phi(X'\hat WX)^{-1}\bigr)\)~\citeat{tu-asm}{Theorem~8.41, Corollary~7.21}, \citeat{shao}{Theorem~4.18}, where \(\hat W\) is \(W\) at \(\hat\beta\). Binomial rows are
covered, with \(\sum_im_i\) in place of \(n\): a row with \(m_i\) trials and response \(y_i\) contributes
\((y_i-m_i\pi_i)x_i\) to the score and \(m_i\pi_i(1-\pi_i)x_ix_i'\) to \(X'WX\) (the sums in the proof of
Theorem~\ref{r12-thm-glm}(a)), as do \(m_i\) Bernoulli rows with the same \(x_i\) and responses of sum \(y_i\). Accordingly,
the \textbf{standard error} of \(c'\hat\beta\) is taken to be \(\operatorname{SE}(c'\hat\beta)=\sqrt{\phi\,c'(X'\hat WX)^{-1}c}\),
with \(\phi(X'\hat WX)^{-1}\) in place of the \(\sigma^2(X'X)^{-1}\) of Theorem~\ref{r12-thm-sampling}(a); for the normal
family, where \(W=I_n\) and \(\phi=\sigma^2\) is known, the normal distribution of \(\hat\beta\) is exact. For
\(H_0\colon K\beta=k_0\) with \(K\in\R^{q\times p}\) of rank \(q\) and \(k_0\in\R^q\), the \textbf{Wald statistic}
\((K\hat\beta-k_0)'\bigl(\phi K(X'\hat WX)^{-1}K'\bigr)^{-1}(K\hat\beta-k_0)\) and the \textbf{likelihood-ratio statistic}
\(2\bigl(\ell(\hat\beta)-\ell(\tilde\beta)\bigr)\), where \(\tilde\beta\) maximizes \(\ell\) subject to \(K\beta=k_0\), are both
approximately \(\chi^2_q\) under \(H_0\), and \(H_0\) is rejected at approximate level \(\alpha\) when the statistic exceeds
\(\chi^2_{q,1-\alpha}\) (Definition~\ref{r15-def-t-F}(c)).

\needspace{14\baselineskip}

\subsection*{Exercises}

\begin{exercise}\label{r12-ex-logistic-groups}
Use the data of Example~\ref{r12-exm-glm-oneway} (three groups of \(36\) patients, \(12\), \(24\), \(18\) recoveries, one row
for each group).
\begin{enumerate}
\item In the baseline coding \(X_2=[\one_3\ d_2\ d_3]\) (Example~\ref{r11-exm-oneway}),
\(\operatorname{logit}\pi_l=\beta_0+\beta_2d_{l2}+\beta_3d_{l3}\)
(\(d_{l2}\), \(d_{l3}\) the entries of \(d_2\), \(d_3\)), find the MLE \(\hat\beta=(\hat\beta_0,\hat\beta_2,\hat\beta_3)'\) and
the odds ratios \(e^{\hat\beta_2}\), \(e^{\hat\beta_3}\), and compare them with Example~\ref{r12-exm-glm-oneway}.
\item At the MLE, find \(\hat W\), \(X_2'\hat WX_2\) and its inverse, and \(\operatorname{SE}(\hat\beta_2)\),
\(\operatorname{SE}(\hat\beta_3)\).
\item In the coding \(X_1=[d_1\ d_2\ d_3]=I_3\) without intercept (Theorem~\ref{r11-thm-oneway}), carry out one IRLS step
from \(\gamma=0\).
\item In the coding \(X=[\one_3\ d_1\ d_2\ d_3]\) of Example~\ref{r12-exm-glm-oneway}, at \(\beta=0\), find \(X'WX\), check
that \(G_0=\frac19\diag(0,1,1,1)\) is a g-inverse of it, compute \(b=G_0X'Wz\), and compare \(Xb\) with the step of (c).
\end{enumerate}
\end{exercise}
\begin{solution}
(a) \(\mathcal C(X_2)=\mathcal C(I_3)\) (Theorem~\ref{r11-thm-oneway}(c)), so by
Proposition~\ref{r12-prop-glm-groups}(c) the MLE satisfies \(X_2\hat\beta=(-\log2,\log2,0)'\), as in
Example~\ref{r12-exm-glm-oneway}: \(\hat\beta_0=-\log2\), \(\hat\beta_0+\hat\beta_2=\log2\), \(\hat\beta_0+\hat\beta_3=0\). So
\(\hat\beta=(-\log2,\ 2\log2,\ \log2)'\), the only MLE (\(\rank X_2=3\), Theorem~\ref{r12-thm-glm}(b)), and
\(e^{\hat\beta_2}=4\), \(e^{\hat\beta_3}=2\). These are the odds ratios of Example~\ref{r12-exm-glm-oneway}:
\(\hat\beta_2\), \(\hat\beta_3\) equal \(\hat\alpha_2-\hat\alpha_1\), \(\hat\alpha_3-\hat\alpha_1\), as both codings give the same
fitted probabilities \(\hat\pi=(\frac13,\frac23,\frac12)'\) (Proposition~\ref{r12-prop-glm-groups}(c)).

(b) \(\hat W=\diag\bigl(36\hat\pi_l(1-\hat\pi_l)\bigr)=\diag(36\cdot\frac13\cdot\frac23,\ 36\cdot\frac23\cdot\frac13,\
36\cdot\frac12\cdot\frac12)=\diag(8,8,9)\), and \(X_2\) is invertible, with the inverse below (\(X_2X_2^{-1}=I_3\) by
multiplication):
\[
X_2=\begin{pmatrix}1&0&0\\1&1&0\\1&0&1\end{pmatrix},\qquad
X_2'\hat WX_2=\begin{pmatrix}25&8&9\\8&8&0\\9&0&9\end{pmatrix},\qquad
X_2^{-1}=\begin{pmatrix}1&0&0\\-1&1&0\\-1&0&1\end{pmatrix}
\]
By Theorem~\ref{r2-thm-inverse-props}(b),(c), \((X_2'\hat WX_2)^{-1}=X_2^{-1}\hat W^{-1}(X_2^{-1})'\), with
\[
X_2^{-1}\hat W^{-1}=\begin{pmatrix}\frac18&0&0\\-\frac18&\frac18&0\\-\frac18&0&\frac19\end{pmatrix},\qquad
(X_2'\hat WX_2)^{-1}=\begin{pmatrix}\frac18&-\frac18&-\frac18\\-\frac18&\frac14&\frac18\\-\frac18&\frac18&\frac{17}{72}\end{pmatrix}
\]
With \(\phi=1\) and the paragraph after Example~\ref{r12-exm-glm-oneway},
\(\operatorname{SE}(\hat\beta_2)^2=\frac14=\frac18+\frac18\) and \(\operatorname{SE}(\hat\beta_3)^2=\frac{17}{72}=\frac18+\frac19\): so
\(\operatorname{SE}(\hat\beta_2)=\frac12\) and \(\operatorname{SE}(\hat\beta_3)=\frac{\sqrt{34}}{12}\).

(c) At \(\gamma=0\): \(\pi_l=\frac12\), \(\mu=18\one_3\), \(W=36\cdot\frac12\cdot\frac12I_3=9I_3\), \(y-\mu=(-6,6,0)'\) and
\(z=\gamma+W^{-1}(y-\mu)=(-\frac23,\frac23,0)'\). By Theorem~\ref{r12-thm-glm}(c),
\(\gamma^{\mathrm{new}}=(X_1'WX_1)^{-1}X_1'Wz=(9I_3)^{-1}\,9z=z=(-\frac23,\frac23,0)'\) (the MLE is
\(\hat\gamma=(-\log2,\log2,0)'\)).

(d) At \(\beta=0\), \(X\beta=0\), so \(\pi\), \(\mu\), \(W=9I_3\) and \(z\) are those of (c), and \(X'Wz=X'(y-\mu)=(0,-6,6,0)'\).
\[
X'WX=9X'X=\begin{pmatrix}27&9&9&9\\9&9&0&0\\9&0&9&0\\9&0&0&9\end{pmatrix},\qquad
X'WX\,G_0=\begin{pmatrix}0&1&1&1\\0&1&0&0\\0&0&1&0\\0&0&0&1\end{pmatrix}
\]
and \(X'WX\,G_0\,X'WX=X'WX\): the first row of the product is the sum of rows \(2\), \(3\), \(4\) of \(X'WX\), which is
\((27,9,9,9)\), and the other rows are rows \(2\), \(3\), \(4\) themselves. So \(G_0\) is a g-inverse of \(X'WX\)
(Definition~\ref{r10-def-ginverse}). Then
\[
b=G_0X'Wz=\frac19\begin{pmatrix}0\\-6\\6\\0\end{pmatrix}=\begin{pmatrix}0\\-\frac23\\\frac23\\0\end{pmatrix},\qquad
Xb=\begin{pmatrix}-\frac23\\\frac23\\0\end{pmatrix}=\gamma^{\mathrm{new}}
\]
as Proposition~\ref{r12-prop-glm-fitted}(c) asserts for the full-rank factorization \(X=X_1S\), \(S=[\one_3\ I_3]\)
(Theorem~\ref{r11-thm-oneway}(a)), although \(b\) itself depends on the g-inverse.
\end{solution}

\begin{exercise}\label{r12-ex-poisson-classes}
A motor portfolio has three rating classes A, B, C. Policy \(i\) is observed for \(t_i\) years (its exposure) and has \(y_i\)
claims. Class A has total exposure \(1000\) policy-years and \(24\) claims, class B \(500\) policy-years and \(48\) claims,
class C \(750\) policy-years and \(36\) claims. Assume the Poisson regression
\(\log\mu_i=\log t_i+\beta_0+\beta_1d_{Bi}+\beta_2d_{Ci}\), where \(d_B\), \(d_C\) are the indicators of classes B and C.
\begin{enumerate}
\item Show that, apart from a term free of \(\beta\), the log-likelihood depends on the data only through the class
totals; find the MLE and the rate ratios of classes B and C to class A.
\item Find \(X'\hat WX\) at the MLE, its inverse, and \(\operatorname{SE}(\hat\beta_1)\), \(\operatorname{SE}(\hat\beta_2)\).
\end{enumerate}
\end{exercise}
\begin{solution}
(a) Let \(Y_l\), \(\tau_l\) be the total count and exposure of class \(l\), and \(\eta_A=\beta_0\), \(\eta_B=\beta_0+\beta_1\),
\(\eta_C=\beta_0+\beta_2\), so that \(x_i'\beta=\eta_l\) for \(i\) in class \(l\). By Theorem~\ref{r12-thm-glm} with
\(b_i(\theta)=e^\theta\), \(c_i(k,\phi)=-\log k!\) and \(\phi=1\),
\[
\ell(\beta)=\sum_i\bigl(y_i(\log t_i+x_i'\beta)-t_ie^{x_i'\beta}-\log y_i!\bigr)
=\sum_{l=A,B,C}\bigl(Y_l\eta_l-\tau_le^{\eta_l}\bigr)+\sum_i(y_i\log t_i-\log y_i!)
\]
\(\mathcal C(X)=\mathcal C([d_A\ d_B\ d_C])\) (Theorem~\ref{r11-thm-oneway}(c)) and every \(Y_l>0\), so by
Proposition~\ref{r12-prop-glm-groups}(c), \(e^{\hat\eta_l}=Y_l/\tau_l\): \(\frac{24}{1000}=\frac3{125}\),
\(\frac{48}{500}=\frac{12}{125}\), \(\frac{36}{750}=\frac6{125}\). Hence \(\hat\beta_0=\log\frac3{125}\),
\(\hat\beta_1=\log\bigl(\frac{12}{125}\big/\frac3{125}\bigr)=\log4\) and \(\hat\beta_2=\log\bigl(\frac6{125}\big/\frac3{125}\bigr)=\log2\),
the only MLE (\(X=[\one_n\ d_B\ d_C]\) has rank \(3\), Theorem~\ref{r11-thm-oneway}(c)). So the rate ratios of classes B and C
to class A are \(e^{\hat\beta_1}=4\) and \(e^{\hat\beta_2}=2\).

(b) Here \(\hat W=\diag(\hat\mu_1,\dots,\hat\mu_n)\) (\(b_i''(\theta)=e^\theta\)). By Theorem~\ref{r1-thm-views}(d),
\(X'\hat WX=\sum_i\hat\mu_ix_ix_i'\); \(x_i\) is \(x_A=(1,0,0)'\), \(x_B=(1,1,0)'\) or \(x_C=(1,0,1)'\) according to the class,
and \(\sum_{i\in l}\hat\mu_i=Y_l\) (Proposition~\ref{r12-prop-glm-groups}(b)). The matrix with rows \(x_A'\), \(x_B'\), \(x_C'\)
is \(X_2\) of Exercise~\ref{r12-ex-logistic-groups}(b), with the inverse computed there; so, by
Theorem~\ref{r1-thm-views}(d) again, \(X'\hat WX=\sum_lY_lx_lx_l'=X_2'\diag(24,48,36)X_2\):
\[
X_2=\begin{pmatrix}1&0&0\\1&1&0\\1&0&1\end{pmatrix},\qquad
X'\hat WX=\begin{pmatrix}108&48&36\\48&48&0\\36&0&36\end{pmatrix},\qquad
X_2^{-1}=\begin{pmatrix}1&0&0\\-1&1&0\\-1&0&1\end{pmatrix}
\]
By Theorem~\ref{r2-thm-inverse-props}(b),(c),
\((X'\hat WX)^{-1}=X_2^{-1}\diag(\frac1{24},\frac1{48},\frac1{36})(X_2^{-1})'\), with
\[
X_2^{-1}\diag\bigl(\tfrac1{24},\tfrac1{48},\tfrac1{36}\bigr)=\begin{pmatrix}\frac1{24}&0&0\\-\frac1{24}&\frac1{48}&0\\-\frac1{24}&0&\frac1{36}\end{pmatrix},\qquad
(X'\hat WX)^{-1}=\begin{pmatrix}\frac1{24}&-\frac1{24}&-\frac1{24}\\-\frac1{24}&\frac1{16}&\frac1{24}\\-\frac1{24}&\frac1{24}&\frac5{72}\end{pmatrix}
\]
So \(\operatorname{SE}(\hat\beta_1)^2=\frac1{16}=\frac1{24}+\frac1{48}\) and \(\operatorname{SE}(\hat\beta_2)^2=\frac5{72}=\frac1{24}+\frac1{36}\)
(\(\phi=1\); the paragraph after Example~\ref{r12-exm-glm-oneway}): \(\operatorname{SE}(\hat\beta_1)=\frac14\) and
\(\operatorname{SE}(\hat\beta_2)=\frac{\sqrt{10}}{12}\).
\end{solution}

\begin{exercise}\label{r12-ex-weighted-projection}
For \(X\) and \(\hat W\) of Example~\ref{r12-exm-glm-factorial}, compute \(P_{X,\hat W}=X(X'\hat WX)^{-1}X'\hat W\).
Check that \(P_{X,\hat W}^2=P_{X,\hat W}\), that \(\hat WP_{X,\hat W}\) is symmetric while \(P_{X,\hat W}\) is not, that
\(\tr P_{X,\hat W}=\rank X\), and that \(I_4-P_{X,\hat W}\) has rank \(1\) (Proposition~\ref{r12-prop-weighted-projection}).
\end{exercise}
\begin{solution}
By Example~\ref{r12-exm-glm-factorial}, \(X'\hat WX=\frac12B\) with \(B\) below. Expanding along the first row,
\(\det B=14(49-9)-7(49-21)+7(21-49)=560-196-196=168\), and by Theorem~\ref{r5d-thm-adjugate}
\[
B=\begin{pmatrix}14&7&7\\7&7&3\\7&3&7\end{pmatrix},\qquad
\operatorname{adj}B=\begin{pmatrix}40&-28&-28\\-28&49&7\\-28&7&49\end{pmatrix}
\]
so \((X'\hat WX)^{-1}=2B^{-1}=\frac2{168}\operatorname{adj}B=\frac1{84}\operatorname{adj}B\). Then
\[
X(X'\hat WX)^{-1}=\frac1{84}\begin{pmatrix}40&-28&-28\\12&21&-21\\12&-21&21\\-16&28&28\end{pmatrix},\qquad
X'\hat W=\begin{pmatrix}\frac32&2&2&\frac32\\0&2&0&\frac32\\0&0&2&\frac32\end{pmatrix}
\]
and \(P_{X,\hat W}=\frac1{14}M\) with
\[
M=\begin{pmatrix}10&4&4&-4\\3&11&-3&3\\3&-3&11&3\\-4&4&4&10\end{pmatrix},\qquad
M^2=\begin{pmatrix}140&56&56&-56\\42&154&-42&42\\42&-42&154&42\\-56&56&56&140\end{pmatrix}=14M
\]
so \(P_{X,\hat W}^2=\frac1{196}M^2=\frac1{14}M=P_{X,\hat W}\). Further,
\[
\hat WP_{X,\hat W}=\frac1{28}\begin{pmatrix}30&12&12&-12\\12&44&-12&12\\12&-12&44&12\\-12&12&12&30\end{pmatrix}
\]
is symmetric, while \(P_{X,\hat W}\) is not: its \((1,2)\) entry is \(\frac27\) and its \((2,1)\) entry \(\frac3{14}\).
\(\tr P_{X,\hat W}=\frac{10+11+11+10}{14}=3=\rank X\) (Example~\ref{r12-exm-glm-factorial}), as it must be:
\(P_{X,\hat W}\) is idempotent of rank \(\rank X\) (Proposition~\ref{r12-prop-weighted-projection}(b),(c)), and an
idempotent matrix has trace equal to its rank (Theorem~\ref{r6-thm-idempotent}(c)). Finally,
\[
I_4-P_{X,\hat W}=\frac1{14}\begin{pmatrix}4&-4&-4&4\\-3&3&3&-3\\-3&3&3&-3\\4&-4&-4&4\end{pmatrix}
=\frac1{14}\begin{pmatrix}4\\-3\\-3\\4\end{pmatrix}(1,\ -1,\ -1,\ 1)
\]
is a nonzero matrix whose columns are multiples of one vector, so its rank is \(1=4-\rank X\).
\end{solution}

\subsection*{Problems}

\begin{problem}[\citeat{tu-asm}{Exercise~8.3}]\label{r12-pr-poisson-residuals}
In a Poisson regression (Definition~\ref{r12-def-glm}) whose design \(X\) has the column \(\one_n\) (an intercept), let
\(\hat\beta\) be an MLE and \(e=y-\hat\mu\) the vector of residuals. Show that
\begin{enumerate}
\item \(e_1+\dots+e_n=0\);
\item the residuals and each column \(u\) of \(X\) have sample covariance \(0\) (Definition~\ref{r12-def-data}), so that their
sample correlation is \(0\) whenever neither \(e\) nor \(u\) is constant (Definition~\ref{r12-def-correlation}).
\end{enumerate}
\end{problem}
\begin{solution}
(a) This is Proposition~\ref{r12-prop-glm-groups}(a), as \(\one_n\in\mathcal C(X)\).

(b) By Theorem~\ref{r12-thm-glm}(b), \(X'e=0\), so \(u'e=\sum_iu_ie_i=0\) for each column \(u\) of \(X\). By (a),
\(\bar e=0\), so
\[
\sum_{i=1}^n(u_i-\bar u)(e_i-\bar e)=\sum_{i=1}^nu_ie_i-\bar u\sum_{i=1}^ne_i=0
\]
and the sample covariance \(\frac1{n-1}\sum_i(u_i-\bar u)(e_i-\bar e)\) is \(0\). The sample correlation
\(\cor(u,e)=S_{ue}/\sqrt{S_{uu}S_{ee}}\) has the numerator \(S_{ue}=0\), so it is \(0\) when \(u\) and \(e\) are not constant
(for the column \(\one_n\) itself it is not defined).
\end{solution}

\clearpage
\def\runningtitle{Bibliography}
\renewcommand{\refname}{Bibliography}
\begin{thebibliography}{10}

\bibitem{hoffman-kunze}
Kenneth Hoffman and Ray Kunze.
\newblock {\em {Linear Algebra}}.
\newblock Prentice-Hall, Englewood Cliffs, NJ, second edition, 1971.

\bibitem{searle-khuri}
Shayle~R. Searle and Andr{\'e}~I. Khuri.
\newblock {\em {Matrix Algebra Useful for Statistics}}.
\newblock John Wiley \& Sons, Hoboken, NJ, second edition, 2017.

\bibitem{bapat}
R.~B. Bapat.
\newblock {\em {Linear Algebra and Linear Models}}.
\newblock Universitext. Springer, London, third edition, 2012.
\newblock \url{https://doi.org/10.1007/978-1-4471-2739-0}.

\bibitem{lang}
Serge Lang.
\newblock {\em {Linear Algebra}}.
\newblock Undergraduate Texts in Mathematics. Springer, New York, third
  edition, 1987.
\newblock \url{https://doi.org/10.1007/978-1-4757-1949-9}.

\bibitem{harville}
David~A. Harville.
\newblock {\em {Matrix Algebra From a Statistician's Perspective}}.
\newblock Springer, New York, 1997.
\newblock \url{https://doi.org/10.1007/b98818}.

\bibitem{abadir-magnus}
Karim~M. Abadir and Jan~R. Magnus.
\newblock {\em {Matrix Algebra}}, volume~1 of {\em Econometric Exercises}.
\newblock Cambridge University Press, Cambridge, 2005.

\bibitem{shakarchi}
Rami Shakarchi.
\newblock {\em {Solutions Manual for Lang's Linear Algebra}}.
\newblock Springer, New York, 1996.
\newblock \url{https://doi.org/10.1007/978-1-4612-0755-9}.

\bibitem{harville-solutions}
David~A. Harville.
\newblock {\em {Matrix Algebra: Exercises and Solutions}}.
\newblock Springer, New York, 2001.
\newblock \url{https://doi.org/10.1007/978-1-4613-0181-3}.

\bibitem{tu-asm}
Chang-Ye Tu.
\newblock {Actuarial Statistical Models}.
\newblock Lecture notes (Mathematical Prerequisites and Units 1--14),
  \url{https://github.com/chang-ye-tu/asm}.

\bibitem{prasolov}
V.~V. Prasolov.
\newblock {\em {Problems and Theorems in Linear Algebra}}, volume 134 of {\em
  Translations of Mathematical Monographs}.
\newblock American Mathematical Society, Providence, RI, 1994.

\bibitem{hefferon}
Jim Hefferon.
\newblock {Linear Algebra}.
\newblock Fourth edition, \url{https://hefferon.net/linearalgebra}, 2021.

\bibitem{tu-clm}
Chang-Ye Tu.
\newblock {Notes on the Classical Theory of Linear Models}.
\newblock Lecture notes,
  \url{https://github.com/chang-ye-tu/mva/blob/master/note/clm.pdf}.

\bibitem{lang-ila}
Serge Lang.
\newblock {\em {Introduction to Linear Algebra}}.
\newblock Undergraduate Texts in Mathematics. Springer, New York, second
  edition, 1986.
\newblock \url{https://doi.org/10.1007/978-1-4612-1070-2}.

\bibitem{knapp-basic-algebra}
Anthony~W. Knapp.
\newblock {\em {Basic Algebra}}.
\newblock Anthony W. Knapp, East Setauket, NY, digital second edition, 2016.
\newblock Freely downloadable from the author,
  \url{https://www.math.stonybrook.edu/~aknapp/download.html}.

\bibitem{zhang}
Fuzhen Zhang.
\newblock {\em {Matrix Theory: Basic Results and Techniques}}.
\newblock Universitext. Springer, New York, third edition, 2026.

\bibitem{horn-johnson}
Roger~A. Horn and Charles~R. Johnson.
\newblock {\em {Matrix Analysis}}.
\newblock Cambridge University Press, Cambridge, second edition, 2013.

\bibitem{casella-berger}
George Casella and Roger~L. Berger.
\newblock {\em {Statistical Inference}}.
\newblock Chapman and Hall/CRC, Boca Raton, FL, second edition, 2024.
\newblock Reissue of the 2002 Duxbury edition,
  \url{https://doi.org/10.1201/9781003456285}.

\bibitem{haerdle-simar}
Wolfgang~Karl H{\"a}rdle and L{\'e}opold Simar.
\newblock {\em {Applied Multivariate Statistical Analysis}}.
\newblock Springer, Berlin, fourth edition, 2015.
\newblock \url{https://doi.org/10.1007/978-3-662-45171-7}.

\bibitem{islr}
Gareth James, Daniela Witten, Trevor Hastie, and Robert Tibshirani.
\newblock {\em {An Introduction to Statistical Learning with Applications in
  R}}.
\newblock Springer, New York, second edition, 2021.
\newblock Freely downloadable from the authors,
  \url{https://www.statlearning.com}.

\bibitem{casella-berger-solutions}
George Casella, Roger~L. Berger, and Damaris Santana.
\newblock {Solutions Manual for Statistical Inference, Second Edition}.
\newblock University of Florida and North Carolina State University.

\bibitem{tu-mc}
Chang-Ye Tu.
\newblock {Introduction to Matrix Calculus}.
\newblock Lecture notes, expanded from Jan R. Magnus, ``A Gentle Introduction
  to Matrix Calculus'',
  \url{https://github.com/chang-ye-tu/mva/blob/master/note/mc.pdf}.

\bibitem{knapp-basic-real-analysis}
Anthony~W. Knapp.
\newblock {\em {Basic Real Analysis}}.
\newblock Anthony W. Knapp, East Setauket, NY, digital second edition, 2016.
\newblock Freely downloadable from the author,
  \url{https://www.math.stonybrook.edu/~aknapp/download.html}.

\bibitem{magnus-neudecker}
Jan~R. Magnus and Heinz Neudecker.
\newblock {\em {Matrix Differential Calculus with Applications in Statistics
  and Econometrics}}.
\newblock Wiley Series in Probability and Statistics. John Wiley \& Sons,
  Hoboken, NJ, third edition, 2019.

\bibitem{shao}
Jun Shao.
\newblock {\em {Mathematical Statistics}}.
\newblock Springer Texts in Statistics. Springer, New York, second edition,
  2003.
\newblock \url{https://doi.org/10.1007/b97553}.

\end{thebibliography}
\end{document}
